%% Versión en español del manuscrito (REVTeX 4.2, doble columna). Traducción de manuscrito_aps.tex.
%% La Fig. 1 se dibuja directamente con TikZ; las demás figuras PGF se comparten con la versión en inglés.
\documentclass[aps,pra,twocolumn,groupedaddress,amsmath,amssymb,floatfix]{revtex4-2}

\usepackage[spanish,es-nodecimaldot]{babel}
\usepackage{graphicx}
\usepackage{pgf}
\usepackage{tikz}
\usetikzlibrary{arrows.meta,babel}
\usepackage{xcolor}
\usepackage{bm}
\usepackage{lmodern}
\usepackage[T1]{fontenc}

\newcommand{\ket}[1]{\lvert #1\rangle}
\newcommand{\bra}[1]{\langle #1\rvert}
\newcommand{\wb}{\omega_b}
\newcommand{\Dxx}{\Delta_{xx}}

\begin{document}

\title{Memoria para un qubit fon\'onico basada en estados oscuros de tres puntos cu\'anticos acoplados}

\author{J. S. Garc\'ia Barrera}
\email{jhgarciaba@unal.edu.co}
\affiliation{Departamento de F\'isica, Facultad de Ciencias Exactas y Naturales,
\\Universidad Nacional de Colombia, Sede Manizales, Manizales, Colombia}

\date{\today}

\begin{abstract}
Una memoria para un qubit fon\'onico debe conservar no solo la excitaci\'on sino tambi\'en la fase entre sus dos intervalos temporales. Proponemos almacenar las dos componentes en los dos estados colectivos \'opticamente oscuros de tres puntos cu\'anticos acoplados por transferencia de F\"orster, que una deformaci\'on no uniforme conecta con un estado brillante. Pulsos Raman selectivos en frecuencia escriben y leen cada componente, y el tercer punto coloca ambas en la misma variedad de un ex\'citon, de modo que el ruido de energ\'ia de modo com\'un deja invariante su fase relativa. Con carga y lectura a trav\'es de gu\'ias de onda en cascada, comparamos la memoria con dos modos ac\'usticos pasivos con los mismos puertos. En el l\'imite de ruido perfectamente com\'un, la probabilidad de recuperaci\'on de la memoria, cercana a $0.7$, y su fidelidad incondicional son casi independientes de la espera hasta ocho tiempos de vida del modo, mientras que los modos pasivos pierden casi toda la excitaci\'on. La fidelidad condicional es menor que la pasiva ($0.83$--$0.88$ frente a $0.95$), de modo que la ganancia est\'a en la probabilidad de recuperar el qubit. El resultado exige una correlaci\'on del ruido cercana a la perfecta, que se endurece con la espera, puertos mucho m\'as r\'apidos que la p\'erdida intr\'inseca, un acoplo por deformaci\'on por encima de un umbral y un desplazamiento de dos excitones entre puntos grande.
\end{abstract}

\maketitle

\section{Introducci\'on}
\label{sec:intro}

Una memoria cu\'antica sincroniza operaciones probabil\'isticas y desacopla la generaci\'on, la transmisi\'on y el procesamiento de un estado cu\'antico, lo que la hace central para la comunicaci\'on a larga distancia~\cite{duan2001}. Se han construido memorias \'opticas con polaritones de estado oscuro, transferencia Raman y peines de frecuencia at\'omicos~\cite{fleischhauer2002,gorshkov2007,afzelius2009}. Su desempe\~no se juzga por el tiempo de almacenamiento, la probabilidad de recuperar la excitaci\'on y la fidelidad del estado recuperado, porque ni una vida media larga ni una salida intensa demuestran que se haya preservado un qubit~\cite{hedges2010,ortu2022,wang2019}. Para un qubit codificado en dos intervalos temporales, la fase relativa entre las componentes temprana y tard\'ia tambi\'en debe sobrevivir a la espera~\cite{ortu2022}. Los mismos criterios se aplican a un qubit fon\'onico.

Una memoria de fonones extender\'ia esa funci\'on a los circuitos integrados, donde los fonones pueden transportar informaci\'on entre componentes~\cite{habraken2012,schuetz2015}. Se han controlado fonones individuales, transferido entre nodos y enviado en dos intervalos temporales por gu\'ias de onda fon\'onicas~\cite{satzinger2018,bienfait2019,dumur2021,zivariWaveguide2022,zivariTimebin2022}. Osciladores mec\'anicos acoplados a circuitos superconductores han almacenado campos de microondas y fonones individuales~\cite{oconnell2010,chu2017,palomaki2013,liu2023,reed2017}. Dispositivos optomec\'anicos han almacenado fonones y luz de banda ancha, y la ingenier\'ia de bandas fon\'onicas ha extendido el tiempo de vida de modos ac\'usticos~\cite{galland2014,wallucks2020,england2013,macCabe2020,taylor2022,bozkurt2025}. En general, sin embargo, la excitaci\'on almacenada sigue siendo un fon\'on durante la espera, de modo que la vida media del modo limita el almacenamiento, y por eso la ingenier\'ia de bandas busca alargarla. Estas arquitecturas deben juzgarse frente al almacenamiento pasivo, que incluye la captura y la lectura y no solo la espera; ese es el criterio que usamos aqu\'i.

Los puntos cu\'anticos (QD) semiconductores acoplan luz y movimiento: el ex\'citon responde a la deformaci\'on y puede controlarse \'opticamente, y su acoplo al movimiento mec\'anico y a resonadores se ha observado hasta el r\'egimen de pocos fonones~\cite{kettler2021,spinnler2024}. La teor\'ia ha usado este acoplo para \emph{emitir} fonones: Bin \emph{et al.}\ propusieron resonancias Stokes de un solo punto para producir paquetes de fonones~\cite{bin2020}, se mostr\'o que el ex\'citon oscuro de esp\'in de un solo punto admite oscilaciones de Rabi gigantes~\cite{vargas2022}, y tratamientos microsc\'opicos del ba\~no muestran que el continuo de fonones modifica la estad\'istica de emisi\'on~\cite{cosacchi2022}. Sin embargo, emitir no es almacenar: estas propuestas no absorben un qubit fon\'onico entrante, no preservan su fase ni lo devuelven. Almacenar requiere estados internos de vida larga a los que puedan acceder tanto el fon\'on como un l\'aser, y preguntamos si un emisor de estado s\'olido controlado \'opticamente puede ofrecerlos.

Una mol\'ecula de QD es un candidato. El acoplo entre emisores crea estados excit\'onicos colectivos, y se han observado la mezcla de excitones brillantes y oscuros, el acoplo controlado, los efectos optofon\'onicos y el almacenamiento de informaci\'on excit\'onica~\cite{ortner2003,krenner2005,gerardot2005,skold2013,kerfoot2014,boyerdelagiroday2011,carlson2019}. Aqu\'i ``oscuro'' significa oscuro por interferencia de los dipolos \'opticos, lo que difiere de los excitones oscuros de esp\'in de un punto aislado~\cite{lueker2017}. Se ha mostrado te\'oricamente que el acoplo vibr\'onico transfiere poblaci\'on entre estados colectivos brillantes y oscuros de emisores moleculares~\cite{holzinger2022}; es un an\'alogo f\'isico cercano, pero ese trabajo no implement\'o un canal de memoria fon\'onica de entrada y salida.

La oscuridad \'optica, sin embargo, no basta. El continuo de fonones a\~nade transferencia incoherente y desfase~\cite{rozbicki2008,muljarov2005,sirkina2025,daniels2013}, y la fase almacenada depende de las correlaciones del ruido: los subespacios libres de decoherencia protegen contra el ruido colectivo~\cite{lidar1998}, se ha propuesto una codificaci\'on distinta de cuadrupolo de carga en tres puntos para suprimir las fluctuaciones uniformes del campo el\'ectrico~\cite{friesen2017}, y la coherencia de un ba\~no de fonones compartido puede reducir el desfase de puntos acoplados~\cite{hall2025}. Falta, por tanto, un protocolo que absorba un qubit fon\'onico en estados oscuros, lo conserve durante una espera bajo estos mecanismos y lo devuelva, y que se compare con el almacenamiento pasivo.

Proponemos una memoria que almacena las dos componentes temporales de un qubit fon\'onico en los dos estados oscuros colectivos de tres puntos acoplados, de modo que la excitaci\'on no se mantenga en un estado \'opticamente brillante durante la espera. La transferencia es posible porque, con la transferencia de F\"orster y una deformaci\'on no uniforme, aparecen elementos de acoplo entre autoestados colectivos, mientras que el acoplo longitudinal de un solo punto no intercambia una excitaci\'on entre dos niveles. Cada componente se escribe mediante una resonancia Raman de dos fotones seleccionada por un pulso l\'aser y se lee invirtiendo el proceso; el l\'aser est\'a apagado durante la espera, pero el acoplo de la cavidad a la gu\'ia de onda debe variarse durante la carga y la lectura. Comparamos la memoria con dos modos ac\'usticos pasivos que tienen las mismas eficiencias de entrada y salida, y analizamos el efecto de las correlaciones del ruido y de los puertos. En el r\'egimen que estudiamos, la memoria mantiene una probabilidad sustancial de recuperar el qubit durante una espera en la que la referencia pasiva pierde casi toda la excitaci\'on, aunque la fidelidad condicionada a la recuperaci\'on es menor.

El art\'iculo se organiza as\'i. La Sec.~\ref{sec:model} introduce el modelo y el subespacio de almacenamiento oscuro, la Sec.~\ref{sec:raman} deriva el proceso de escritura y lectura Raman, la Sec.~\ref{sec:noise} trata las p\'erdidas, el ruido y los puertos, y la Sec.~\ref{sec:results} presenta la comparaci\'on con el almacenamiento pasivo. La Sec.~\ref{sec:platform} discute los requisitos de plataforma, que son condiciones del resultado y no una implementaci\'on, y la Sec.~\ref{sec:conclusions} concluye.

\section{Modelo y subespacio de almacenamiento oscuro}
\label{sec:model}

Consideramos tres puntos cu\'anticos semiconductores con operadores de descenso $\sigma_i$ y ocupaciones $n_i=\sigma_i^\dagger\sigma_i$ ($i=1,2,3$), acoplados por transferencia de F\"orster y situados dentro de una cavidad ac\'ustica con un solo modo de frecuencia $\wb$ y operador de descenso $b$. Un campo l\'aser com\'un excita los excitones, y la cavidad se acopla a gu\'ias de onda de entrada y salida durante la carga y la lectura. La Figura~\ref{fig:scheme} muestra las transiciones Raman internas. Tomamos $\hbar=1$ y expresamos energ\'ias y tasas en unidades de $\wb$.

\begin{figure}[t]
	\centering
	\begin{tikzpicture}[x=1cm,y=1cm,>=Latex,
		level/.style={line width=1.1pt,line cap=round},
		coupling/.style={{Latex[length=2mm]}-{Latex[length=2mm]},line width=0.8pt}]
		% Niveles esquemáticos: cada ranura se pone en resonancia con un láser distinto.
		\definecolor{schemegold}{RGB}{169,111,28}
		\definecolor{schemered}{RGB}{186,42,38}
		\definecolor{schemeblue}{RGB}{35,96,170}
		\definecolor{schemegreen}{RGB}{38,126,78}
		\draw[level] (0.15,0.55)--(1.15,0.55);
		\node[above,font=\small] at (0.65,0.55) {$\ket{G,1}$};
		\draw[level,schemered] (2.35,2.55)--(3.75,2.55);
		\node[above,schemered,font=\small] at (3.05,2.55) {$\ket{B,1}$};
		\draw[level,schemeblue] (2.35,1.48)--(3.75,1.48);
		\node[below,schemeblue,font=\small] at (3.05,1.48) {$\ket{D_1,0}$};
		\draw[level,schemegreen] (2.35,0.42)--(3.75,0.42);
		\node[below,schemegreen,font=\small] at (3.05,0.42) {$\ket{D_2,0}$};
		\draw[coupling,schemegold] (1.15,0.66)--(2.35,2.44);
		\node[schemegold,font=\small,anchor=east,align=right] at (1.75,1.90)
		{$\Omega_s(t)$\\$\omega_{L,s}$};
		\draw[coupling,schemeblue] (2.57,2.43)--(2.57,1.60);
		\node[schemeblue,font=\small,anchor=west] at (2.70,2.02) {$g_1$};
		\draw[coupling,schemegreen] (3.60,2.43)--(3.60,0.54);
		\node[schemegreen,font=\small,anchor=west] at (3.70,2.02) {$g_2$};
	\end{tikzpicture}
	\caption{(En color en l\'inea) Transiciones Raman usadas para escribir y leer las dos ranuras oscuras. Un l\'aser de frecuencia $\omega_{L,s}$ acopla $\ket{G,1}$ fuera de resonancia con el intermedio brillante virtual $\ket{B,1}$, y la deformaci\'on no uniforme acopla este \'ultimo con $\ket{D_s,0}$ mediante el elemento de matriz $g_s$. Las posiciones de los niveles son esquem\'aticas.}
	\label{fig:scheme}
\end{figure}

El hamiltoniano es
\begin{align}
H(t)={}&\wb\, b^\dagger b+\sum_{ij}h_{ij}\,\sigma_i^\dagger\sigma_j
+\Dxx\sum_{i<j}n_in_j\nonumber\\
&+\sum_i\lambda_i\,n_i\,(b+b^\dagger)+H_L(t),\label{eq:H}\\
H_L(t)={}&\frac{\Omega(t)}{\sqrt3}\sum_i\bigl(\sigma_i\,e^{i\omega_Lt}+\sigma_i^\dagger\,e^{-i\omega_Lt}\bigr),\nonumber\\
&h_{ii}=\varepsilon_i,\quad h_{i\neq j}=J_{ij}.\nonumber
\end{align}
Aqu\'i $\varepsilon_i$ son las energ\'ias de excit\'on de los puntos, $J_{ij}$ los acoplos de F\"orster, $\Dxx$ la energ\'ia de interacci\'on de cada par de excitones, $\lambda_i$ el acoplo longitudinal por potencial de deformaci\'on del punto $i$ con el modo ac\'ustico y $\Omega(t)$ la amplitud de un l\'aser de frecuencia $\omega_L$ con la misma fase en los tres puntos. Las mol\'eculas de puntos reales tambi\'en pueden tener acoplo t\'unel y transiciones asistidas por fonones~\cite{daniels2013}; estos procesos no se representan con t\'erminos coherentes de salto adicionales en la Ec.~\eqref{eq:H}. El modelo \emph{completo} conserva los ocho estados electr\'onicos de los tres puntos y permite comprobar si los pulsos pueblan biexcitones. El modelo \emph{reducido} restringe el espacio electr\'onico al vac\'io y a la variedad de un ex\'citon, lo cual es exacto para $|\Dxx|\to\infty$.

En la base de sitio $\ket{i}=\sigma_i^\dagger\ket{G}$, donde $\ket G$ es el estado fundamental electr\'onico, el l\'aser com\'un solo excita
\begin{equation}
\ket{B}=\frac{\ket1+\ket2+\ket3}{\sqrt3},\qquad
\bra{B}\sum_i\sigma_i^\dagger\ket G=\sqrt3 .\label{eq:bright}
\end{equation}
Dos vectores ortogonales son
\begin{equation}
\ket{D_1}=\frac{\ket2-\ket3}{\sqrt2},\qquad
\ket{D_2}=\frac{2\ket1-\ket2-\ket3}{\sqrt6}.\label{eq:dark}
\end{equation}
La ortogonalidad con $\ket B$ los hace oscuros para un campo colectivo y para un canal radiativo colectivo en el l\'imite de Dicke de fases \'opticas indistinguibles~\cite{dicke1954,mokhlespour2012}. Para que adem\'as sean autoestados estacionarios de $h$, $\ket B$ debe ser un autovector, de modo que las tres sumas de fila de $h$ deben ser iguales. Esta es una condici\'on sobre las energ\'ias de sitio y los acoplos, no una propiedad autom\'atica de tres emisores. Fijando $E_B=0$, $E_{D_1}=-G_1$ y $E_{D_2}=-G_2$ con $0<G_1<G_2$, la descomposici\'on espectral $h=\sum_aE_a\ket a\bra a$ da
\begin{align}
\varepsilon_1&=-\frac{2G_2}{3},&
\varepsilon_2=\varepsilon_3&=-\frac{G_1}{2}-\frac{G_2}{6},\nonumber\\
J_{12}=J_{13}&=\frac{G_2}{3},&
J_{23}&=\frac{G_1}{2}-\frac{G_2}{6},\label{eq:inverse}
\end{align}
de modo que $\varepsilon_1+2J_{12}=\varepsilon_2+J_{12}+J_{23}=0$.

Con $\Lambda=\mathrm{diag}(\lambda_1,\lambda_2,\lambda_3)$ en la variedad de un ex\'citon, los elementos de matriz que conectan el estado brillante con los dos oscuros son
\begin{align}
g_1&=\bra B\Lambda\ket{D_1}=\frac{\lambda_2-\lambda_3}{\sqrt6},\nonumber\\
g_2&=\bra B\Lambda\ket{D_2}=\frac{2\lambda_1-\lambda_2-\lambda_3}{\sqrt{18}}.\label{eq:g}
\end{align}
Una componente de $\lambda_i$ igual en los tres puntos desplaza las energ\'ias pero no contribuye a $g_s$.

La entrada que se almacena es un qubit de un fon\'on en dos intervalos temporales, entrelazado con una referencia $R$ que no se somete al protocolo:
\begin{equation}
\ket{\Psi_{\rm in}}=\frac{\ket0_R\ket{1_{t_1},0_{t_2}}+\ket1_R\ket{0_{t_1},1_{t_2}}}{\sqrt2}.
\label{eq:input}
\end{equation}
La primera componente se escribe en $D_1$ y la segunda en $D_2$. Durante la espera, ambas ramas ocupan la misma variedad de un ex\'citon, $\mathcal H_L=\mathrm{span}\{\ket{D_1,0},\ket{D_2,0}\}$. Hacen falta tres puntos porque con dos emisores solo hay un estado oscuro colectivo, de modo que no caben dos componentes en la misma variedad; una codificaci\'on $\ket G/\ket D$ comparar\'ia en cambio los sectores de cero y un ex\'citon, cuya fase relativa no est\'a protegida contra el ruido com\'un.

\section{Escritura y lectura Raman}
\label{sec:raman}

Cada componente temporal se escribe en una ranura oscura con un pulso l\'aser de frecuencia seleccionada. Las dos ranuras tienen energ\'ias distintas, de modo que cada una se direcciona solo con la frecuencia del l\'aser y no hace falta conmutar ninguna interacci\'on, y la lectura invierte el proceso empezando por la ranura escrita en \'ultimo lugar. Para la ranura $s$, la transici\'on Raman de dos fotones va de $\ket{G,1}$ (puntos vac\'ios, un fon\'on) a $\ket{D_s,0}$ (ex\'citon en la ranura $s$, sin fon\'on), con el estado brillante $\ket B$ como nivel intermedio muy desintonizado. En el sistema rotante a la frecuencia del l\'aser, $\delta_a=E_a-\omega_L$, los estados $\ket{G,1}$, $\ket{B,1}$ y $\ket{D_s,0}$ tienen energ\'ias $\wb$, $\delta_B+\wb$ y $\delta_s$. El l\'aser acopla los dos primeros con amplitud $\Omega$, y la deformaci\'on acopla los dos \'ultimos con amplitud $g_s$ (Ec.~\eqref{eq:g}), conservando solo la parte que conserva la excitaci\'on, $\ket B\bra{D_s}\,b^\dagger$. La resonancia de dos fotones es $\delta_s=\wb+\delta_2$, con $\delta_2$ peque\~no, y la desintonizaci\'on del nivel intermedio es $\Delta_{I,s}\equiv\delta_B\simeq\wb+G_s$, que siempre es grande porque los estados oscuros est\'an por debajo de $\ket B$. Eliminando $\ket{B,1}$ para $|\Omega|,|g_s|\ll\Delta_{I,s}$ se obtiene, en la base $\{\ket{G,1},\ket{D_s,0}\}$,
\begin{align}
H_{\rm eff}^{(s)}&=
\begin{pmatrix}
\wb-\Omega^2/\Delta_{I,s}&-\Omega_{{\rm eff},s}\\
-\Omega_{{\rm eff},s}&\wb+\delta_2-g_s^2/\Delta_{I,s}
\end{pmatrix},\nonumber\\
\Omega_{{\rm eff},s}&=\frac{\Omega g_s}{\Delta_{I,s}}.\label{eq:heff}
\end{align}
La eliminaci\'on es adiab\'atica cuando la envolvente del pulso es suave en la escala de tiempo $1/\Delta_I$~\cite{brion2007,reiter2012}; un apagado abrupto deja una mezcla residual de $\ket B$. La transferencia completa ocurre cuando las energ\'ias diagonales son iguales, $\delta_{2,s}^{\star}=(g_s^2-\Omega^2)/\Delta_{I,s}$, tras el tiempo $\pi$
\begin{equation}
t_{\pi,s}=\frac{\pi}{2\Omega_{{\rm eff},s}}=\frac{\pi\Delta_{I,s}}{2\Omega g_s}.\label{eq:tpi}
\end{equation}

La Ecuaci\'on~\eqref{eq:heff} omite que $\ket{D_s,0}$ tambi\'en est\'a acoplado por $\Lambda$ a todos los estados de un fon\'on de la variedad de un ex\'citon, $\ket{c,1}$ con $c\in\{B,D_1,D_2\}$, incluyendo $c=s$ a trav\'es del elemento diagonal $\bra{D_s}\Lambda\ket{D_s}$. A segundo orden en $\Lambda$ la energ\'ia de $\ket{D_s,0}$ se desplaza en $-\sum_c|\Lambda_{cs}|^2/(\wb+E_c-E_s)$, y la condici\'on de resonancia pasa a ser
\begin{align}
\delta_{2,s}^{\star}&\simeq-\frac{\Omega^2}{\Delta_{I,s}}+\sum_{c\in\{B,D_1,D_2\}}\frac{|\Lambda_{cs}|^2}{\wb+E_c-E_s},\nonumber\\
\Lambda_{cs}&=\bra c\Lambda\ket{D_s}.\label{eq:stark}
\end{align}
El t\'ermino $c=B$ reproduce el desplazamiento de la Ec.~\eqref{eq:heff}. El t\'ermino diagonal $c=s$ es el desplazamiento polar\'onico, y domina para $s=2$ porque $\bra{D_2}\Lambda\ket{D_2}=0.157\,\wb$ para los par\'ametros de la Fig.~\ref{fig:raman_map}. Los denominadores deben mantenerse lejos de resonancias adicionales.

Mientras act\'ua el pulso, el fon\'on permanece en el modo con probabilidad $\cos^2(\Omega_{\rm eff}t)$ y decae con la tasa de la cavidad $\kappa$, de modo que la eficiencia de escritura es, a primer orden en $\kappa$,
\begin{equation}
\eta_w\simeq1-\kappa\int_0^{t_\pi}\!\cos^2(\Omega_{\rm eff}t)\,\mathrm dt=1-\frac{\kappa t_\pi}{2}.\label{eq:etaw}
\end{equation}
Reducir $\Delta_I$ acorta $t_\pi$ pero aumenta la poblaci\'on virtual del estado brillante, $p_B\sim(\Omega/\Delta_I)^2$, que desfasa y se relaja. A $t_\pi$ fijo, $\Omega_{\rm eff}=\pi/2t_\pi$ y el acoplo requerido para un $p_{B,\max}$ dado es
\begin{align}
g_{\min}&=\frac{\Omega_{\rm eff}}{\sqrt{p_{B,\max}}}=\frac{\pi}{2\,t_\pi\sqrt{p_{B,\max}}},\nonumber\\
\Omega&=\frac{\Omega_{\rm eff}\Delta_I}{g}\propto\frac1g.\label{eq:gmin}
\end{align}
Para $t_\pi\simeq190/\wb$ esto da $g_{\min}=0.026\,\wb$ si $p_{B,\max}=0.1$ y $0.019\,\wb$ si $p_{B,\max}=0.2$. Un acoplo por deformaci\'on menor requiere por tanto un pulso proporcionalmente m\'as intenso de duraci\'on casi igual.

La ranura no direccionada, $s'$, est\'a desintonizada en frecuencia de dos fotones por $\Delta_{ss'}=|E_{D_s}-E_{D_{s'}}|$ y se excita con amplitud efectiva $\Omega_{\rm eff}'=\Omega g_{s'}/\Delta_I$. La poblaci\'on m\'axima que alcanza es la de una oscilaci\'on de Rabi desintonizada,
\begin{equation}
P_{s'}^{\max}=\frac{(2\Omega_{\rm eff}')^2}{(2\Omega_{\rm eff}')^2+\Delta_{ss'}^2}\simeq\frac{4\Omega_{\rm eff}'^2}{\Delta_{ss'}^2},\label{eq:leak}
\end{equation}
que da $5\times10^{-3}$ y $7\times10^{-3}$ para los pulsos sobre $D_1$ y $D_2$ ($\Omega_{\rm eff}'\simeq0.009$ y $0.010\,\wb$, $\Delta_{ss'}=0.25\,\wb$).

\begin{figure}[t]
	\centering
	\resizebox{\columnwidth}{!}{%% Creator: Matplotlib, PGF backend
%%
%% To include the figure in your LaTeX document, write
%%   \input{<filename>.pgf}
%%
%% Make sure the required packages are loaded in your preamble
%%   \usepackage{pgf}
%%
%% Also ensure that all the required font packages are loaded; for instance,
%% the lmodern package is sometimes necessary when using math font.
%%   \usepackage{lmodern}
%%
%% Figures using additional raster images can only be included by \input if
%% they are in the same directory as the main LaTeX file. For loading figures
%% from other directories you can use the `import` package
%%   \usepackage{import}
%%
%% and then include the figures with
%%   \import{<path to file>}{<filename>.pgf}
%%
%% Matplotlib used the following preamble
%%   \def\mathdefault#1{#1}
%%   \everymath=\expandafter{\the\everymath\displaystyle}
%%   \IfFileExists{scrextend.sty}{
%%     \usepackage[fontsize=8.000000pt]{scrextend}
%%   }{
%%     \renewcommand{\normalsize}{\fontsize{8.000000}{9.600000}\selectfont}
%%     \normalsize
%%   }
%%   \usepackage{amsmath,amssymb}
%%   \makeatletter\@ifpackageloaded{underscore}{}{\usepackage[strings]{underscore}}\makeatother
%%
\begingroup%
\makeatletter%
\begin{pgfpicture}%
\pgfpathrectangle{\pgfpointorigin}{\pgfqpoint{5.007798in}{2.380989in}}%
\pgfusepath{use as bounding box, clip}%
\begin{pgfscope}%
\pgfsetbuttcap%
\pgfsetmiterjoin%
\definecolor{currentfill}{rgb}{1.000000,1.000000,1.000000}%
\pgfsetfillcolor{currentfill}%
\pgfsetlinewidth{0.000000pt}%
\definecolor{currentstroke}{rgb}{1.000000,1.000000,1.000000}%
\pgfsetstrokecolor{currentstroke}%
\pgfsetdash{}{0pt}%
\pgfpathmoveto{\pgfqpoint{0.000000in}{0.000000in}}%
\pgfpathlineto{\pgfqpoint{5.007798in}{0.000000in}}%
\pgfpathlineto{\pgfqpoint{5.007798in}{2.380989in}}%
\pgfpathlineto{\pgfqpoint{0.000000in}{2.380989in}}%
\pgfpathlineto{\pgfqpoint{0.000000in}{0.000000in}}%
\pgfpathclose%
\pgfusepath{fill}%
\end{pgfscope}%
\begin{pgfscope}%
\pgfsetbuttcap%
\pgfsetmiterjoin%
\definecolor{currentfill}{rgb}{1.000000,1.000000,1.000000}%
\pgfsetfillcolor{currentfill}%
\pgfsetlinewidth{0.000000pt}%
\definecolor{currentstroke}{rgb}{0.000000,0.000000,0.000000}%
\pgfsetstrokecolor{currentstroke}%
\pgfsetstrokeopacity{0.000000}%
\pgfsetdash{}{0pt}%
\pgfpathmoveto{\pgfqpoint{0.540975in}{0.476543in}}%
\pgfpathlineto{\pgfqpoint{2.361363in}{0.476543in}}%
\pgfpathlineto{\pgfqpoint{2.361363in}{2.132543in}}%
\pgfpathlineto{\pgfqpoint{0.540975in}{2.132543in}}%
\pgfpathlineto{\pgfqpoint{0.540975in}{0.476543in}}%
\pgfpathclose%
\pgfusepath{fill}%
\end{pgfscope}%
\begin{pgfscope}%
\pgfpathrectangle{\pgfqpoint{0.540975in}{0.476543in}}{\pgfqpoint{1.820388in}{1.656000in}}%
\pgfusepath{clip}%
\pgfsetbuttcap%
\pgfsetroundjoin%
\definecolor{currentfill}{rgb}{0.207536,0.260208,0.608151}%
\pgfsetfillcolor{currentfill}%
\pgfsetlinewidth{0.000000pt}%
\definecolor{currentstroke}{rgb}{0.000000,0.000000,0.000000}%
\pgfsetstrokecolor{currentstroke}%
\pgfsetdash{}{0pt}%
\pgfpathmoveto{\pgfqpoint{0.566258in}{0.476543in}}%
\pgfpathlineto{\pgfqpoint{0.591541in}{0.476543in}}%
\pgfpathlineto{\pgfqpoint{0.616824in}{0.476543in}}%
\pgfpathlineto{\pgfqpoint{0.642107in}{0.476543in}}%
\pgfpathlineto{\pgfqpoint{0.667391in}{0.476543in}}%
\pgfpathlineto{\pgfqpoint{0.692674in}{0.476543in}}%
\pgfpathlineto{\pgfqpoint{0.717957in}{0.476543in}}%
\pgfpathlineto{\pgfqpoint{0.743240in}{0.476543in}}%
\pgfpathlineto{\pgfqpoint{0.768523in}{0.476543in}}%
\pgfpathlineto{\pgfqpoint{0.793806in}{0.476543in}}%
\pgfpathlineto{\pgfqpoint{0.819090in}{0.476543in}}%
\pgfpathlineto{\pgfqpoint{0.844373in}{0.476543in}}%
\pgfpathlineto{\pgfqpoint{0.869656in}{0.476543in}}%
\pgfpathlineto{\pgfqpoint{0.894939in}{0.476543in}}%
\pgfpathlineto{\pgfqpoint{0.920222in}{0.476543in}}%
\pgfpathlineto{\pgfqpoint{0.945505in}{0.476543in}}%
\pgfpathlineto{\pgfqpoint{0.970789in}{0.476543in}}%
\pgfpathlineto{\pgfqpoint{0.996072in}{0.476543in}}%
\pgfpathlineto{\pgfqpoint{1.021355in}{0.476543in}}%
\pgfpathlineto{\pgfqpoint{1.046638in}{0.476543in}}%
\pgfpathlineto{\pgfqpoint{1.071921in}{0.476543in}}%
\pgfpathlineto{\pgfqpoint{1.097204in}{0.476543in}}%
\pgfpathlineto{\pgfqpoint{1.122488in}{0.476543in}}%
\pgfpathlineto{\pgfqpoint{1.147771in}{0.476543in}}%
\pgfpathlineto{\pgfqpoint{1.173054in}{0.476543in}}%
\pgfpathlineto{\pgfqpoint{1.198337in}{0.476543in}}%
\pgfpathlineto{\pgfqpoint{1.223620in}{0.476543in}}%
\pgfpathlineto{\pgfqpoint{1.248904in}{0.476543in}}%
\pgfpathlineto{\pgfqpoint{1.274187in}{0.476543in}}%
\pgfpathlineto{\pgfqpoint{1.299470in}{0.476543in}}%
\pgfpathlineto{\pgfqpoint{1.324753in}{0.476543in}}%
\pgfpathlineto{\pgfqpoint{1.350036in}{0.476543in}}%
\pgfpathlineto{\pgfqpoint{1.375319in}{0.476543in}}%
\pgfpathlineto{\pgfqpoint{1.400603in}{0.476543in}}%
\pgfpathlineto{\pgfqpoint{1.425886in}{0.476543in}}%
\pgfpathlineto{\pgfqpoint{1.451169in}{0.476543in}}%
\pgfpathlineto{\pgfqpoint{1.476452in}{0.476543in}}%
\pgfpathlineto{\pgfqpoint{1.501735in}{0.476543in}}%
\pgfpathlineto{\pgfqpoint{1.527018in}{0.476543in}}%
\pgfpathlineto{\pgfqpoint{1.552302in}{0.476543in}}%
\pgfpathlineto{\pgfqpoint{1.577585in}{0.476543in}}%
\pgfpathlineto{\pgfqpoint{1.602868in}{0.476543in}}%
\pgfpathlineto{\pgfqpoint{1.628151in}{0.476543in}}%
\pgfpathlineto{\pgfqpoint{1.653434in}{0.476543in}}%
\pgfpathlineto{\pgfqpoint{1.678717in}{0.476543in}}%
\pgfpathlineto{\pgfqpoint{1.704001in}{0.476543in}}%
\pgfpathlineto{\pgfqpoint{1.729284in}{0.476543in}}%
\pgfpathlineto{\pgfqpoint{1.754567in}{0.476543in}}%
\pgfpathlineto{\pgfqpoint{1.779850in}{0.476543in}}%
\pgfpathlineto{\pgfqpoint{1.805133in}{0.476543in}}%
\pgfpathlineto{\pgfqpoint{1.830416in}{0.476543in}}%
\pgfpathlineto{\pgfqpoint{1.855700in}{0.476543in}}%
\pgfpathlineto{\pgfqpoint{1.880983in}{0.476543in}}%
\pgfpathlineto{\pgfqpoint{1.906266in}{0.476543in}}%
\pgfpathlineto{\pgfqpoint{1.931549in}{0.476543in}}%
\pgfpathlineto{\pgfqpoint{1.956832in}{0.476543in}}%
\pgfpathlineto{\pgfqpoint{1.982115in}{0.476543in}}%
\pgfpathlineto{\pgfqpoint{2.007399in}{0.476543in}}%
\pgfpathlineto{\pgfqpoint{2.032682in}{0.476543in}}%
\pgfpathlineto{\pgfqpoint{2.057965in}{0.476543in}}%
\pgfpathlineto{\pgfqpoint{2.083248in}{0.476543in}}%
\pgfpathlineto{\pgfqpoint{2.108531in}{0.476543in}}%
\pgfpathlineto{\pgfqpoint{2.133815in}{0.476543in}}%
\pgfpathlineto{\pgfqpoint{2.159098in}{0.476543in}}%
\pgfpathlineto{\pgfqpoint{2.184381in}{0.476543in}}%
\pgfpathlineto{\pgfqpoint{2.209664in}{0.476543in}}%
\pgfpathlineto{\pgfqpoint{2.234947in}{0.476543in}}%
\pgfpathlineto{\pgfqpoint{2.260230in}{0.476543in}}%
\pgfpathlineto{\pgfqpoint{2.285514in}{0.476543in}}%
\pgfpathlineto{\pgfqpoint{2.310797in}{0.476543in}}%
\pgfpathlineto{\pgfqpoint{2.336080in}{0.476543in}}%
\pgfpathlineto{\pgfqpoint{2.361363in}{0.476543in}}%
\pgfpathlineto{\pgfqpoint{2.361363in}{0.489791in}}%
\pgfpathlineto{\pgfqpoint{2.361363in}{0.503039in}}%
\pgfpathlineto{\pgfqpoint{2.361363in}{0.516287in}}%
\pgfpathlineto{\pgfqpoint{2.361363in}{0.529535in}}%
\pgfpathlineto{\pgfqpoint{2.361363in}{0.542783in}}%
\pgfpathlineto{\pgfqpoint{2.361363in}{0.556031in}}%
\pgfpathlineto{\pgfqpoint{2.361363in}{0.569279in}}%
\pgfpathlineto{\pgfqpoint{2.361363in}{0.582527in}}%
\pgfpathlineto{\pgfqpoint{2.361363in}{0.595775in}}%
\pgfpathlineto{\pgfqpoint{2.361363in}{0.609023in}}%
\pgfpathlineto{\pgfqpoint{2.361363in}{0.622271in}}%
\pgfpathlineto{\pgfqpoint{2.361363in}{0.635519in}}%
\pgfpathlineto{\pgfqpoint{2.361363in}{0.648767in}}%
\pgfpathlineto{\pgfqpoint{2.361363in}{0.662015in}}%
\pgfpathlineto{\pgfqpoint{2.361363in}{0.675263in}}%
\pgfpathlineto{\pgfqpoint{2.361363in}{0.688511in}}%
\pgfpathlineto{\pgfqpoint{2.361363in}{0.701759in}}%
\pgfpathlineto{\pgfqpoint{2.361363in}{0.715007in}}%
\pgfpathlineto{\pgfqpoint{2.361363in}{0.728255in}}%
\pgfpathlineto{\pgfqpoint{2.361363in}{0.741503in}}%
\pgfpathlineto{\pgfqpoint{2.361363in}{0.754751in}}%
\pgfpathlineto{\pgfqpoint{2.361363in}{0.767999in}}%
\pgfpathlineto{\pgfqpoint{2.361363in}{0.781247in}}%
\pgfpathlineto{\pgfqpoint{2.361363in}{0.794495in}}%
\pgfpathlineto{\pgfqpoint{2.361363in}{0.807743in}}%
\pgfpathlineto{\pgfqpoint{2.361363in}{0.820991in}}%
\pgfpathlineto{\pgfqpoint{2.361363in}{0.834239in}}%
\pgfpathlineto{\pgfqpoint{2.361363in}{0.847487in}}%
\pgfpathlineto{\pgfqpoint{2.361363in}{0.860735in}}%
\pgfpathlineto{\pgfqpoint{2.361363in}{0.873983in}}%
\pgfpathlineto{\pgfqpoint{2.361363in}{0.887231in}}%
\pgfpathlineto{\pgfqpoint{2.361363in}{0.900479in}}%
\pgfpathlineto{\pgfqpoint{2.361363in}{0.913727in}}%
\pgfpathlineto{\pgfqpoint{2.361363in}{0.926975in}}%
\pgfpathlineto{\pgfqpoint{2.361363in}{0.940223in}}%
\pgfpathlineto{\pgfqpoint{2.361363in}{0.953471in}}%
\pgfpathlineto{\pgfqpoint{2.361363in}{0.966719in}}%
\pgfpathlineto{\pgfqpoint{2.361363in}{0.979967in}}%
\pgfpathlineto{\pgfqpoint{2.361363in}{0.993215in}}%
\pgfpathlineto{\pgfqpoint{2.361363in}{1.006463in}}%
\pgfpathlineto{\pgfqpoint{2.361363in}{1.019711in}}%
\pgfpathlineto{\pgfqpoint{2.361363in}{1.032959in}}%
\pgfpathlineto{\pgfqpoint{2.361363in}{1.046207in}}%
\pgfpathlineto{\pgfqpoint{2.361363in}{1.059455in}}%
\pgfpathlineto{\pgfqpoint{2.361363in}{1.072703in}}%
\pgfpathlineto{\pgfqpoint{2.361363in}{1.085951in}}%
\pgfpathlineto{\pgfqpoint{2.361363in}{1.099199in}}%
\pgfpathlineto{\pgfqpoint{2.361363in}{1.112447in}}%
\pgfpathlineto{\pgfqpoint{2.361363in}{1.125695in}}%
\pgfpathlineto{\pgfqpoint{2.361363in}{1.138943in}}%
\pgfpathlineto{\pgfqpoint{2.361363in}{1.152191in}}%
\pgfpathlineto{\pgfqpoint{2.361363in}{1.165439in}}%
\pgfpathlineto{\pgfqpoint{2.361363in}{1.178687in}}%
\pgfpathlineto{\pgfqpoint{2.361363in}{1.191935in}}%
\pgfpathlineto{\pgfqpoint{2.361363in}{1.205183in}}%
\pgfpathlineto{\pgfqpoint{2.361363in}{1.218431in}}%
\pgfpathlineto{\pgfqpoint{2.361363in}{1.231679in}}%
\pgfpathlineto{\pgfqpoint{2.361363in}{1.244927in}}%
\pgfpathlineto{\pgfqpoint{2.361363in}{1.258175in}}%
\pgfpathlineto{\pgfqpoint{2.361363in}{1.271423in}}%
\pgfpathlineto{\pgfqpoint{2.361363in}{1.284671in}}%
\pgfpathlineto{\pgfqpoint{2.361363in}{1.297919in}}%
\pgfpathlineto{\pgfqpoint{2.361363in}{1.311167in}}%
\pgfpathlineto{\pgfqpoint{2.361363in}{1.324415in}}%
\pgfpathlineto{\pgfqpoint{2.361363in}{1.337663in}}%
\pgfpathlineto{\pgfqpoint{2.361363in}{1.350911in}}%
\pgfpathlineto{\pgfqpoint{2.361363in}{1.364159in}}%
\pgfpathlineto{\pgfqpoint{2.361363in}{1.377407in}}%
\pgfpathlineto{\pgfqpoint{2.361363in}{1.390655in}}%
\pgfpathlineto{\pgfqpoint{2.361363in}{1.403903in}}%
\pgfpathlineto{\pgfqpoint{2.361363in}{1.417151in}}%
\pgfpathlineto{\pgfqpoint{2.361363in}{1.430399in}}%
\pgfpathlineto{\pgfqpoint{2.361363in}{1.443647in}}%
\pgfpathlineto{\pgfqpoint{2.361363in}{1.456895in}}%
\pgfpathlineto{\pgfqpoint{2.361363in}{1.470143in}}%
\pgfpathlineto{\pgfqpoint{2.361363in}{1.483391in}}%
\pgfpathlineto{\pgfqpoint{2.361363in}{1.496639in}}%
\pgfpathlineto{\pgfqpoint{2.361363in}{1.509887in}}%
\pgfpathlineto{\pgfqpoint{2.361363in}{1.523135in}}%
\pgfpathlineto{\pgfqpoint{2.361363in}{1.536383in}}%
\pgfpathlineto{\pgfqpoint{2.361363in}{1.549631in}}%
\pgfpathlineto{\pgfqpoint{2.361363in}{1.562879in}}%
\pgfpathlineto{\pgfqpoint{2.361363in}{1.576127in}}%
\pgfpathlineto{\pgfqpoint{2.361363in}{1.589375in}}%
\pgfpathlineto{\pgfqpoint{2.361363in}{1.602623in}}%
\pgfpathlineto{\pgfqpoint{2.361363in}{1.615871in}}%
\pgfpathlineto{\pgfqpoint{2.361363in}{1.629119in}}%
\pgfpathlineto{\pgfqpoint{2.361363in}{1.642367in}}%
\pgfpathlineto{\pgfqpoint{2.361363in}{1.655615in}}%
\pgfpathlineto{\pgfqpoint{2.361363in}{1.668863in}}%
\pgfpathlineto{\pgfqpoint{2.361363in}{1.682111in}}%
\pgfpathlineto{\pgfqpoint{2.361363in}{1.695359in}}%
\pgfpathlineto{\pgfqpoint{2.361363in}{1.708607in}}%
\pgfpathlineto{\pgfqpoint{2.361363in}{1.721855in}}%
\pgfpathlineto{\pgfqpoint{2.361363in}{1.735103in}}%
\pgfpathlineto{\pgfqpoint{2.361363in}{1.748351in}}%
\pgfpathlineto{\pgfqpoint{2.361363in}{1.761599in}}%
\pgfpathlineto{\pgfqpoint{2.361363in}{1.774847in}}%
\pgfpathlineto{\pgfqpoint{2.361363in}{1.788095in}}%
\pgfpathlineto{\pgfqpoint{2.361363in}{1.801343in}}%
\pgfpathlineto{\pgfqpoint{2.361363in}{1.814591in}}%
\pgfpathlineto{\pgfqpoint{2.361363in}{1.827839in}}%
\pgfpathlineto{\pgfqpoint{2.361363in}{1.841087in}}%
\pgfpathlineto{\pgfqpoint{2.361363in}{1.854335in}}%
\pgfpathlineto{\pgfqpoint{2.361363in}{1.867583in}}%
\pgfpathlineto{\pgfqpoint{2.361363in}{1.880831in}}%
\pgfpathlineto{\pgfqpoint{2.361363in}{1.894079in}}%
\pgfpathlineto{\pgfqpoint{2.361363in}{1.907327in}}%
\pgfpathlineto{\pgfqpoint{2.361363in}{1.920575in}}%
\pgfpathlineto{\pgfqpoint{2.361363in}{1.933823in}}%
\pgfpathlineto{\pgfqpoint{2.361363in}{1.947071in}}%
\pgfpathlineto{\pgfqpoint{2.361363in}{1.960319in}}%
\pgfpathlineto{\pgfqpoint{2.361363in}{1.973567in}}%
\pgfpathlineto{\pgfqpoint{2.361363in}{1.986815in}}%
\pgfpathlineto{\pgfqpoint{2.361363in}{2.000063in}}%
\pgfpathlineto{\pgfqpoint{2.361363in}{2.013311in}}%
\pgfpathlineto{\pgfqpoint{2.361363in}{2.026559in}}%
\pgfpathlineto{\pgfqpoint{2.361363in}{2.039807in}}%
\pgfpathlineto{\pgfqpoint{2.361363in}{2.053055in}}%
\pgfpathlineto{\pgfqpoint{2.361363in}{2.066303in}}%
\pgfpathlineto{\pgfqpoint{2.361363in}{2.079551in}}%
\pgfpathlineto{\pgfqpoint{2.361363in}{2.092799in}}%
\pgfpathlineto{\pgfqpoint{2.361363in}{2.106047in}}%
\pgfpathlineto{\pgfqpoint{2.361363in}{2.119295in}}%
\pgfpathlineto{\pgfqpoint{2.361363in}{2.132543in}}%
\pgfpathlineto{\pgfqpoint{2.336080in}{2.132543in}}%
\pgfpathlineto{\pgfqpoint{2.310797in}{2.132543in}}%
\pgfpathlineto{\pgfqpoint{2.285514in}{2.132543in}}%
\pgfpathlineto{\pgfqpoint{2.260230in}{2.132543in}}%
\pgfpathlineto{\pgfqpoint{2.234947in}{2.132543in}}%
\pgfpathlineto{\pgfqpoint{2.209664in}{2.132543in}}%
\pgfpathlineto{\pgfqpoint{2.184381in}{2.132543in}}%
\pgfpathlineto{\pgfqpoint{2.159098in}{2.132543in}}%
\pgfpathlineto{\pgfqpoint{2.133815in}{2.132543in}}%
\pgfpathlineto{\pgfqpoint{2.130056in}{2.132543in}}%
\pgfpathlineto{\pgfqpoint{2.133815in}{2.127619in}}%
\pgfpathlineto{\pgfqpoint{2.139886in}{2.119295in}}%
\pgfpathlineto{\pgfqpoint{2.140685in}{2.106047in}}%
\pgfpathlineto{\pgfqpoint{2.147849in}{2.092799in}}%
\pgfpathlineto{\pgfqpoint{2.151981in}{2.079551in}}%
\pgfpathlineto{\pgfqpoint{2.157215in}{2.066303in}}%
\pgfpathlineto{\pgfqpoint{2.159098in}{2.063619in}}%
\pgfpathlineto{\pgfqpoint{2.166170in}{2.053055in}}%
\pgfpathlineto{\pgfqpoint{2.167258in}{2.039807in}}%
\pgfpathlineto{\pgfqpoint{2.177449in}{2.026559in}}%
\pgfpathlineto{\pgfqpoint{2.177048in}{2.013311in}}%
\pgfpathlineto{\pgfqpoint{2.184381in}{2.004413in}}%
\pgfpathlineto{\pgfqpoint{2.187554in}{2.000063in}}%
\pgfpathlineto{\pgfqpoint{2.190274in}{1.986815in}}%
\pgfpathlineto{\pgfqpoint{2.194875in}{1.973567in}}%
\pgfpathlineto{\pgfqpoint{2.202517in}{1.960319in}}%
\pgfpathlineto{\pgfqpoint{2.203370in}{1.947071in}}%
\pgfpathlineto{\pgfqpoint{2.209664in}{1.941337in}}%
\pgfpathlineto{\pgfqpoint{2.216829in}{1.933823in}}%
\pgfpathlineto{\pgfqpoint{2.216067in}{1.920575in}}%
\pgfpathlineto{\pgfqpoint{2.227710in}{1.907327in}}%
\pgfpathlineto{\pgfqpoint{2.228967in}{1.894079in}}%
\pgfpathlineto{\pgfqpoint{2.234947in}{1.885945in}}%
\pgfpathlineto{\pgfqpoint{2.237871in}{1.880831in}}%
\pgfpathlineto{\pgfqpoint{2.245388in}{1.867583in}}%
\pgfpathlineto{\pgfqpoint{2.245948in}{1.854335in}}%
\pgfpathlineto{\pgfqpoint{2.259683in}{1.841087in}}%
\pgfpathlineto{\pgfqpoint{2.253459in}{1.827839in}}%
\pgfpathlineto{\pgfqpoint{2.260230in}{1.823734in}}%
\pgfpathlineto{\pgfqpoint{2.270944in}{1.814591in}}%
\pgfpathlineto{\pgfqpoint{2.267862in}{1.801343in}}%
\pgfpathlineto{\pgfqpoint{2.277762in}{1.788095in}}%
\pgfpathlineto{\pgfqpoint{2.282230in}{1.774847in}}%
\pgfpathlineto{\pgfqpoint{2.280714in}{1.761599in}}%
\pgfpathlineto{\pgfqpoint{2.285514in}{1.759372in}}%
\pgfpathlineto{\pgfqpoint{2.299775in}{1.748351in}}%
\pgfpathlineto{\pgfqpoint{2.286110in}{1.735103in}}%
\pgfpathlineto{\pgfqpoint{2.310797in}{1.724238in}}%
\pgfpathlineto{\pgfqpoint{2.314082in}{1.721855in}}%
\pgfpathlineto{\pgfqpoint{2.310797in}{1.718389in}}%
\pgfpathlineto{\pgfqpoint{2.285514in}{1.713124in}}%
\pgfpathlineto{\pgfqpoint{2.260230in}{1.719119in}}%
\pgfpathlineto{\pgfqpoint{2.254824in}{1.721855in}}%
\pgfpathlineto{\pgfqpoint{2.260230in}{1.732965in}}%
\pgfpathlineto{\pgfqpoint{2.278974in}{1.735103in}}%
\pgfpathlineto{\pgfqpoint{2.260230in}{1.735913in}}%
\pgfpathlineto{\pgfqpoint{2.234947in}{1.747798in}}%
\pgfpathlineto{\pgfqpoint{2.234263in}{1.748351in}}%
\pgfpathlineto{\pgfqpoint{2.230605in}{1.761599in}}%
\pgfpathlineto{\pgfqpoint{2.221595in}{1.774847in}}%
\pgfpathlineto{\pgfqpoint{2.209664in}{1.780903in}}%
\pgfpathlineto{\pgfqpoint{2.200085in}{1.788095in}}%
\pgfpathlineto{\pgfqpoint{2.198333in}{1.801343in}}%
\pgfpathlineto{\pgfqpoint{2.184381in}{1.808585in}}%
\pgfpathlineto{\pgfqpoint{2.176612in}{1.814591in}}%
\pgfpathlineto{\pgfqpoint{2.176837in}{1.827839in}}%
\pgfpathlineto{\pgfqpoint{2.159098in}{1.840784in}}%
\pgfpathlineto{\pgfqpoint{2.158803in}{1.841087in}}%
\pgfpathlineto{\pgfqpoint{2.152932in}{1.854335in}}%
\pgfpathlineto{\pgfqpoint{2.145257in}{1.867583in}}%
\pgfpathlineto{\pgfqpoint{2.133815in}{1.879145in}}%
\pgfpathlineto{\pgfqpoint{2.132469in}{1.880831in}}%
\pgfpathlineto{\pgfqpoint{2.131909in}{1.894079in}}%
\pgfpathlineto{\pgfqpoint{2.116061in}{1.907327in}}%
\pgfpathlineto{\pgfqpoint{2.114191in}{1.920575in}}%
\pgfpathlineto{\pgfqpoint{2.108531in}{1.925347in}}%
\pgfpathlineto{\pgfqpoint{2.101135in}{1.933823in}}%
\pgfpathlineto{\pgfqpoint{2.096165in}{1.947071in}}%
\pgfpathlineto{\pgfqpoint{2.089367in}{1.960319in}}%
\pgfpathlineto{\pgfqpoint{2.083248in}{1.966865in}}%
\pgfpathlineto{\pgfqpoint{2.078006in}{1.973567in}}%
\pgfpathlineto{\pgfqpoint{2.075306in}{1.986815in}}%
\pgfpathlineto{\pgfqpoint{2.063287in}{2.000063in}}%
\pgfpathlineto{\pgfqpoint{2.062029in}{2.013311in}}%
\pgfpathlineto{\pgfqpoint{2.057965in}{2.017542in}}%
\pgfpathlineto{\pgfqpoint{2.051135in}{2.026559in}}%
\pgfpathlineto{\pgfqpoint{2.046689in}{2.039807in}}%
\pgfpathlineto{\pgfqpoint{2.040166in}{2.053055in}}%
\pgfpathlineto{\pgfqpoint{2.032682in}{2.065114in}}%
\pgfpathlineto{\pgfqpoint{2.031998in}{2.066303in}}%
\pgfpathlineto{\pgfqpoint{2.030174in}{2.079551in}}%
\pgfpathlineto{\pgfqpoint{2.019515in}{2.092799in}}%
\pgfpathlineto{\pgfqpoint{2.017666in}{2.106047in}}%
\pgfpathlineto{\pgfqpoint{2.007863in}{2.119295in}}%
\pgfpathlineto{\pgfqpoint{2.007399in}{2.120728in}}%
\pgfpathlineto{\pgfqpoint{2.003774in}{2.132543in}}%
\pgfpathlineto{\pgfqpoint{1.982115in}{2.132543in}}%
\pgfpathlineto{\pgfqpoint{1.956832in}{2.132543in}}%
\pgfpathlineto{\pgfqpoint{1.931549in}{2.132543in}}%
\pgfpathlineto{\pgfqpoint{1.906266in}{2.132543in}}%
\pgfpathlineto{\pgfqpoint{1.893542in}{2.132543in}}%
\pgfpathlineto{\pgfqpoint{1.899271in}{2.119295in}}%
\pgfpathlineto{\pgfqpoint{1.900889in}{2.106047in}}%
\pgfpathlineto{\pgfqpoint{1.904646in}{2.092799in}}%
\pgfpathlineto{\pgfqpoint{1.906266in}{2.086741in}}%
\pgfpathlineto{\pgfqpoint{1.909056in}{2.079551in}}%
\pgfpathlineto{\pgfqpoint{1.912460in}{2.066303in}}%
\pgfpathlineto{\pgfqpoint{1.919137in}{2.053055in}}%
\pgfpathlineto{\pgfqpoint{1.920563in}{2.039807in}}%
\pgfpathlineto{\pgfqpoint{1.926466in}{2.026559in}}%
\pgfpathlineto{\pgfqpoint{1.928072in}{2.013311in}}%
\pgfpathlineto{\pgfqpoint{1.931549in}{2.003957in}}%
\pgfpathlineto{\pgfqpoint{1.933497in}{2.000063in}}%
\pgfpathlineto{\pgfqpoint{1.937856in}{1.986815in}}%
\pgfpathlineto{\pgfqpoint{1.940959in}{1.973567in}}%
\pgfpathlineto{\pgfqpoint{1.947010in}{1.960319in}}%
\pgfpathlineto{\pgfqpoint{1.948586in}{1.947071in}}%
\pgfpathlineto{\pgfqpoint{1.955901in}{1.933823in}}%
\pgfpathlineto{\pgfqpoint{1.956832in}{1.927040in}}%
\pgfpathlineto{\pgfqpoint{1.958047in}{1.920575in}}%
\pgfpathlineto{\pgfqpoint{1.964978in}{1.907327in}}%
\pgfpathlineto{\pgfqpoint{1.969120in}{1.894079in}}%
\pgfpathlineto{\pgfqpoint{1.973047in}{1.880831in}}%
\pgfpathlineto{\pgfqpoint{1.979455in}{1.867583in}}%
\pgfpathlineto{\pgfqpoint{1.981075in}{1.854335in}}%
\pgfpathlineto{\pgfqpoint{1.982115in}{1.852436in}}%
\pgfpathlineto{\pgfqpoint{1.989456in}{1.841087in}}%
\pgfpathlineto{\pgfqpoint{1.991083in}{1.827839in}}%
\pgfpathlineto{\pgfqpoint{1.999093in}{1.814591in}}%
\pgfpathlineto{\pgfqpoint{2.002766in}{1.801343in}}%
\pgfpathlineto{\pgfqpoint{2.007200in}{1.788095in}}%
\pgfpathlineto{\pgfqpoint{2.007399in}{1.787693in}}%
\pgfpathlineto{\pgfqpoint{2.014926in}{1.774847in}}%
\pgfpathlineto{\pgfqpoint{2.017104in}{1.761599in}}%
\pgfpathlineto{\pgfqpoint{2.026943in}{1.748351in}}%
\pgfpathlineto{\pgfqpoint{2.028638in}{1.735103in}}%
\pgfpathlineto{\pgfqpoint{2.032682in}{1.728630in}}%
\pgfpathlineto{\pgfqpoint{2.037398in}{1.721855in}}%
\pgfpathlineto{\pgfqpoint{2.041199in}{1.708607in}}%
\pgfpathlineto{\pgfqpoint{2.047821in}{1.695359in}}%
\pgfpathlineto{\pgfqpoint{2.054974in}{1.682111in}}%
\pgfpathlineto{\pgfqpoint{2.057401in}{1.668863in}}%
\pgfpathlineto{\pgfqpoint{2.057965in}{1.668078in}}%
\pgfpathlineto{\pgfqpoint{2.067702in}{1.655615in}}%
\pgfpathlineto{\pgfqpoint{2.068409in}{1.642367in}}%
\pgfpathlineto{\pgfqpoint{2.079866in}{1.629119in}}%
\pgfpathlineto{\pgfqpoint{2.083248in}{1.615911in}}%
\pgfpathlineto{\pgfqpoint{2.083260in}{1.615871in}}%
\pgfpathlineto{\pgfqpoint{2.091608in}{1.602623in}}%
\pgfpathlineto{\pgfqpoint{2.099428in}{1.589375in}}%
\pgfpathlineto{\pgfqpoint{2.103258in}{1.576127in}}%
\pgfpathlineto{\pgfqpoint{2.108531in}{1.570408in}}%
\pgfpathlineto{\pgfqpoint{2.115953in}{1.562879in}}%
\pgfpathlineto{\pgfqpoint{2.116858in}{1.549631in}}%
\pgfpathlineto{\pgfqpoint{2.128697in}{1.536383in}}%
\pgfpathlineto{\pgfqpoint{2.130134in}{1.523135in}}%
\pgfpathlineto{\pgfqpoint{2.133815in}{1.518589in}}%
\pgfpathlineto{\pgfqpoint{2.140883in}{1.509887in}}%
\pgfpathlineto{\pgfqpoint{2.148185in}{1.496639in}}%
\pgfpathlineto{\pgfqpoint{2.153161in}{1.483391in}}%
\pgfpathlineto{\pgfqpoint{2.159098in}{1.477334in}}%
\pgfpathlineto{\pgfqpoint{2.166404in}{1.470143in}}%
\pgfpathlineto{\pgfqpoint{2.166592in}{1.456895in}}%
\pgfpathlineto{\pgfqpoint{2.184364in}{1.443647in}}%
\pgfpathlineto{\pgfqpoint{2.184381in}{1.443445in}}%
\pgfpathlineto{\pgfqpoint{2.185608in}{1.430399in}}%
\pgfpathlineto{\pgfqpoint{2.198000in}{1.417151in}}%
\pgfpathlineto{\pgfqpoint{2.204091in}{1.403903in}}%
\pgfpathlineto{\pgfqpoint{2.209664in}{1.392395in}}%
\pgfpathlineto{\pgfqpoint{2.210488in}{1.390655in}}%
\pgfpathlineto{\pgfqpoint{2.224906in}{1.377407in}}%
\pgfpathlineto{\pgfqpoint{2.224555in}{1.364159in}}%
\pgfpathlineto{\pgfqpoint{2.234947in}{1.356847in}}%
\pgfpathlineto{\pgfqpoint{2.243114in}{1.350911in}}%
\pgfpathlineto{\pgfqpoint{2.238815in}{1.337663in}}%
\pgfpathlineto{\pgfqpoint{2.259068in}{1.324415in}}%
\pgfpathlineto{\pgfqpoint{2.260230in}{1.319119in}}%
\pgfpathlineto{\pgfqpoint{2.262097in}{1.311167in}}%
\pgfpathlineto{\pgfqpoint{2.270262in}{1.297919in}}%
\pgfpathlineto{\pgfqpoint{2.285372in}{1.284671in}}%
\pgfpathlineto{\pgfqpoint{2.280643in}{1.271423in}}%
\pgfpathlineto{\pgfqpoint{2.285514in}{1.269218in}}%
\pgfpathlineto{\pgfqpoint{2.308649in}{1.258175in}}%
\pgfpathlineto{\pgfqpoint{2.297316in}{1.244927in}}%
\pgfpathlineto{\pgfqpoint{2.310797in}{1.238369in}}%
\pgfpathlineto{\pgfqpoint{2.322778in}{1.231679in}}%
\pgfpathlineto{\pgfqpoint{2.310797in}{1.221588in}}%
\pgfpathlineto{\pgfqpoint{2.301973in}{1.218431in}}%
\pgfpathlineto{\pgfqpoint{2.310797in}{1.213344in}}%
\pgfpathlineto{\pgfqpoint{2.320757in}{1.205183in}}%
\pgfpathlineto{\pgfqpoint{2.310797in}{1.194379in}}%
\pgfpathlineto{\pgfqpoint{2.285514in}{1.192834in}}%
\pgfpathlineto{\pgfqpoint{2.260230in}{1.197008in}}%
\pgfpathlineto{\pgfqpoint{2.234947in}{1.203057in}}%
\pgfpathlineto{\pgfqpoint{2.228501in}{1.205183in}}%
\pgfpathlineto{\pgfqpoint{2.228703in}{1.218431in}}%
\pgfpathlineto{\pgfqpoint{2.209664in}{1.221881in}}%
\pgfpathlineto{\pgfqpoint{2.184381in}{1.228224in}}%
\pgfpathlineto{\pgfqpoint{2.173760in}{1.231679in}}%
\pgfpathlineto{\pgfqpoint{2.172106in}{1.244927in}}%
\pgfpathlineto{\pgfqpoint{2.159098in}{1.249387in}}%
\pgfpathlineto{\pgfqpoint{2.139835in}{1.258175in}}%
\pgfpathlineto{\pgfqpoint{2.133815in}{1.266122in}}%
\pgfpathlineto{\pgfqpoint{2.130159in}{1.271423in}}%
\pgfpathlineto{\pgfqpoint{2.114786in}{1.284671in}}%
\pgfpathlineto{\pgfqpoint{2.108531in}{1.287913in}}%
\pgfpathlineto{\pgfqpoint{2.091488in}{1.297919in}}%
\pgfpathlineto{\pgfqpoint{2.084942in}{1.311167in}}%
\pgfpathlineto{\pgfqpoint{2.083248in}{1.311882in}}%
\pgfpathlineto{\pgfqpoint{2.059218in}{1.324415in}}%
\pgfpathlineto{\pgfqpoint{2.057965in}{1.328587in}}%
\pgfpathlineto{\pgfqpoint{2.055460in}{1.337663in}}%
\pgfpathlineto{\pgfqpoint{2.033058in}{1.350911in}}%
\pgfpathlineto{\pgfqpoint{2.032682in}{1.351411in}}%
\pgfpathlineto{\pgfqpoint{2.023763in}{1.364159in}}%
\pgfpathlineto{\pgfqpoint{2.011594in}{1.377407in}}%
\pgfpathlineto{\pgfqpoint{2.007399in}{1.380967in}}%
\pgfpathlineto{\pgfqpoint{1.997122in}{1.390655in}}%
\pgfpathlineto{\pgfqpoint{1.993705in}{1.403903in}}%
\pgfpathlineto{\pgfqpoint{1.982115in}{1.411560in}}%
\pgfpathlineto{\pgfqpoint{1.974645in}{1.417151in}}%
\pgfpathlineto{\pgfqpoint{1.971431in}{1.430399in}}%
\pgfpathlineto{\pgfqpoint{1.956832in}{1.443006in}}%
\pgfpathlineto{\pgfqpoint{1.956162in}{1.443647in}}%
\pgfpathlineto{\pgfqpoint{1.949225in}{1.456895in}}%
\pgfpathlineto{\pgfqpoint{1.939873in}{1.470143in}}%
\pgfpathlineto{\pgfqpoint{1.931549in}{1.478726in}}%
\pgfpathlineto{\pgfqpoint{1.927183in}{1.483391in}}%
\pgfpathlineto{\pgfqpoint{1.922571in}{1.496639in}}%
\pgfpathlineto{\pgfqpoint{1.909082in}{1.509887in}}%
\pgfpathlineto{\pgfqpoint{1.906589in}{1.523135in}}%
\pgfpathlineto{\pgfqpoint{1.906266in}{1.523435in}}%
\pgfpathlineto{\pgfqpoint{1.893798in}{1.536383in}}%
\pgfpathlineto{\pgfqpoint{1.888290in}{1.549631in}}%
\pgfpathlineto{\pgfqpoint{1.880983in}{1.562077in}}%
\pgfpathlineto{\pgfqpoint{1.880511in}{1.562879in}}%
\pgfpathlineto{\pgfqpoint{1.871515in}{1.576127in}}%
\pgfpathlineto{\pgfqpoint{1.868157in}{1.589375in}}%
\pgfpathlineto{\pgfqpoint{1.856616in}{1.602623in}}%
\pgfpathlineto{\pgfqpoint{1.855700in}{1.607667in}}%
\pgfpathlineto{\pgfqpoint{1.854071in}{1.615871in}}%
\pgfpathlineto{\pgfqpoint{1.843641in}{1.629119in}}%
\pgfpathlineto{\pgfqpoint{1.839612in}{1.642367in}}%
\pgfpathlineto{\pgfqpoint{1.832964in}{1.655615in}}%
\pgfpathlineto{\pgfqpoint{1.830416in}{1.659774in}}%
\pgfpathlineto{\pgfqpoint{1.824435in}{1.668863in}}%
\pgfpathlineto{\pgfqpoint{1.820734in}{1.682111in}}%
\pgfpathlineto{\pgfqpoint{1.811959in}{1.695359in}}%
\pgfpathlineto{\pgfqpoint{1.810123in}{1.708607in}}%
\pgfpathlineto{\pgfqpoint{1.805133in}{1.716306in}}%
\pgfpathlineto{\pgfqpoint{1.801216in}{1.721855in}}%
\pgfpathlineto{\pgfqpoint{1.797644in}{1.735103in}}%
\pgfpathlineto{\pgfqpoint{1.791534in}{1.748351in}}%
\pgfpathlineto{\pgfqpoint{1.786084in}{1.761599in}}%
\pgfpathlineto{\pgfqpoint{1.783405in}{1.774847in}}%
\pgfpathlineto{\pgfqpoint{1.779850in}{1.780391in}}%
\pgfpathlineto{\pgfqpoint{1.774209in}{1.788095in}}%
\pgfpathlineto{\pgfqpoint{1.771847in}{1.801343in}}%
\pgfpathlineto{\pgfqpoint{1.764279in}{1.814591in}}%
\pgfpathlineto{\pgfqpoint{1.761603in}{1.827839in}}%
\pgfpathlineto{\pgfqpoint{1.756778in}{1.841087in}}%
\pgfpathlineto{\pgfqpoint{1.754567in}{1.847063in}}%
\pgfpathlineto{\pgfqpoint{1.751091in}{1.854335in}}%
\pgfpathlineto{\pgfqpoint{1.747900in}{1.867583in}}%
\pgfpathlineto{\pgfqpoint{1.741363in}{1.880831in}}%
\pgfpathlineto{\pgfqpoint{1.740079in}{1.894079in}}%
\pgfpathlineto{\pgfqpoint{1.733505in}{1.907327in}}%
\pgfpathlineto{\pgfqpoint{1.731021in}{1.920575in}}%
\pgfpathlineto{\pgfqpoint{1.729284in}{1.925412in}}%
\pgfpathlineto{\pgfqpoint{1.725192in}{1.933823in}}%
\pgfpathlineto{\pgfqpoint{1.720440in}{1.947071in}}%
\pgfpathlineto{\pgfqpoint{1.717781in}{1.960319in}}%
\pgfpathlineto{\pgfqpoint{1.711780in}{1.973567in}}%
\pgfpathlineto{\pgfqpoint{1.710450in}{1.986815in}}%
\pgfpathlineto{\pgfqpoint{1.705251in}{2.000063in}}%
\pgfpathlineto{\pgfqpoint{1.704001in}{2.010557in}}%
\pgfpathlineto{\pgfqpoint{1.703466in}{2.013311in}}%
\pgfpathlineto{\pgfqpoint{1.697734in}{2.026559in}}%
\pgfpathlineto{\pgfqpoint{1.693241in}{2.039807in}}%
\pgfpathlineto{\pgfqpoint{1.690427in}{2.053055in}}%
\pgfpathlineto{\pgfqpoint{1.685859in}{2.066303in}}%
\pgfpathlineto{\pgfqpoint{1.684755in}{2.079551in}}%
\pgfpathlineto{\pgfqpoint{1.679901in}{2.092799in}}%
\pgfpathlineto{\pgfqpoint{1.678717in}{2.102507in}}%
\pgfpathlineto{\pgfqpoint{1.677932in}{2.106047in}}%
\pgfpathlineto{\pgfqpoint{1.671863in}{2.119295in}}%
\pgfpathlineto{\pgfqpoint{1.668380in}{2.132543in}}%
\pgfpathlineto{\pgfqpoint{1.653434in}{2.132543in}}%
\pgfpathlineto{\pgfqpoint{1.628151in}{2.132543in}}%
\pgfpathlineto{\pgfqpoint{1.616784in}{2.132543in}}%
\pgfpathlineto{\pgfqpoint{1.620673in}{2.119295in}}%
\pgfpathlineto{\pgfqpoint{1.621477in}{2.106047in}}%
\pgfpathlineto{\pgfqpoint{1.624048in}{2.092799in}}%
\pgfpathlineto{\pgfqpoint{1.625886in}{2.079551in}}%
\pgfpathlineto{\pgfqpoint{1.627524in}{2.066303in}}%
\pgfpathlineto{\pgfqpoint{1.628151in}{2.063571in}}%
\pgfpathlineto{\pgfqpoint{1.633427in}{2.053055in}}%
\pgfpathlineto{\pgfqpoint{1.634973in}{2.039807in}}%
\pgfpathlineto{\pgfqpoint{1.639835in}{2.026559in}}%
\pgfpathlineto{\pgfqpoint{1.641054in}{2.013311in}}%
\pgfpathlineto{\pgfqpoint{1.644676in}{2.000063in}}%
\pgfpathlineto{\pgfqpoint{1.646616in}{1.986815in}}%
\pgfpathlineto{\pgfqpoint{1.648377in}{1.973567in}}%
\pgfpathlineto{\pgfqpoint{1.651260in}{1.960319in}}%
\pgfpathlineto{\pgfqpoint{1.652205in}{1.947071in}}%
\pgfpathlineto{\pgfqpoint{1.653434in}{1.942717in}}%
\pgfpathlineto{\pgfqpoint{1.658476in}{1.933823in}}%
\pgfpathlineto{\pgfqpoint{1.660276in}{1.920575in}}%
\pgfpathlineto{\pgfqpoint{1.664765in}{1.907327in}}%
\pgfpathlineto{\pgfqpoint{1.667334in}{1.894079in}}%
\pgfpathlineto{\pgfqpoint{1.669788in}{1.880831in}}%
\pgfpathlineto{\pgfqpoint{1.673404in}{1.867583in}}%
\pgfpathlineto{\pgfqpoint{1.674454in}{1.854335in}}%
\pgfpathlineto{\pgfqpoint{1.677996in}{1.841087in}}%
\pgfpathlineto{\pgfqpoint{1.678717in}{1.831910in}}%
\pgfpathlineto{\pgfqpoint{1.679385in}{1.827839in}}%
\pgfpathlineto{\pgfqpoint{1.685781in}{1.814591in}}%
\pgfpathlineto{\pgfqpoint{1.688840in}{1.801343in}}%
\pgfpathlineto{\pgfqpoint{1.691808in}{1.788095in}}%
\pgfpathlineto{\pgfqpoint{1.695825in}{1.774847in}}%
\pgfpathlineto{\pgfqpoint{1.697166in}{1.761599in}}%
\pgfpathlineto{\pgfqpoint{1.701882in}{1.748351in}}%
\pgfpathlineto{\pgfqpoint{1.703010in}{1.735103in}}%
\pgfpathlineto{\pgfqpoint{1.704001in}{1.731601in}}%
\pgfpathlineto{\pgfqpoint{1.708755in}{1.721855in}}%
\pgfpathlineto{\pgfqpoint{1.712156in}{1.708607in}}%
\pgfpathlineto{\pgfqpoint{1.715951in}{1.695359in}}%
\pgfpathlineto{\pgfqpoint{1.720578in}{1.682111in}}%
\pgfpathlineto{\pgfqpoint{1.722100in}{1.668863in}}%
\pgfpathlineto{\pgfqpoint{1.726793in}{1.655615in}}%
\pgfpathlineto{\pgfqpoint{1.727910in}{1.642367in}}%
\pgfpathlineto{\pgfqpoint{1.729284in}{1.638610in}}%
\pgfpathlineto{\pgfqpoint{1.734811in}{1.629119in}}%
\pgfpathlineto{\pgfqpoint{1.738165in}{1.615871in}}%
\pgfpathlineto{\pgfqpoint{1.742374in}{1.602623in}}%
\pgfpathlineto{\pgfqpoint{1.747243in}{1.589375in}}%
\pgfpathlineto{\pgfqpoint{1.749210in}{1.576127in}}%
\pgfpathlineto{\pgfqpoint{1.754567in}{1.563999in}}%
\pgfpathlineto{\pgfqpoint{1.755338in}{1.562879in}}%
\pgfpathlineto{\pgfqpoint{1.757301in}{1.549631in}}%
\pgfpathlineto{\pgfqpoint{1.764269in}{1.536383in}}%
\pgfpathlineto{\pgfqpoint{1.767158in}{1.523135in}}%
\pgfpathlineto{\pgfqpoint{1.771771in}{1.509887in}}%
\pgfpathlineto{\pgfqpoint{1.776673in}{1.496639in}}%
\pgfpathlineto{\pgfqpoint{1.778929in}{1.483391in}}%
\pgfpathlineto{\pgfqpoint{1.779850in}{1.481462in}}%
\pgfpathlineto{\pgfqpoint{1.787435in}{1.470143in}}%
\pgfpathlineto{\pgfqpoint{1.789472in}{1.456895in}}%
\pgfpathlineto{\pgfqpoint{1.797718in}{1.443647in}}%
\pgfpathlineto{\pgfqpoint{1.800513in}{1.430399in}}%
\pgfpathlineto{\pgfqpoint{1.805104in}{1.417151in}}%
\pgfpathlineto{\pgfqpoint{1.805133in}{1.417073in}}%
\pgfpathlineto{\pgfqpoint{1.812031in}{1.403903in}}%
\pgfpathlineto{\pgfqpoint{1.815203in}{1.390655in}}%
\pgfpathlineto{\pgfqpoint{1.823838in}{1.377407in}}%
\pgfpathlineto{\pgfqpoint{1.825594in}{1.364159in}}%
\pgfpathlineto{\pgfqpoint{1.830416in}{1.355912in}}%
\pgfpathlineto{\pgfqpoint{1.834246in}{1.350911in}}%
\pgfpathlineto{\pgfqpoint{1.837853in}{1.337663in}}%
\pgfpathlineto{\pgfqpoint{1.845759in}{1.324415in}}%
\pgfpathlineto{\pgfqpoint{1.851975in}{1.311167in}}%
\pgfpathlineto{\pgfqpoint{1.855420in}{1.297919in}}%
\pgfpathlineto{\pgfqpoint{1.855700in}{1.297490in}}%
\pgfpathlineto{\pgfqpoint{1.866158in}{1.284671in}}%
\pgfpathlineto{\pgfqpoint{1.868199in}{1.271423in}}%
\pgfpathlineto{\pgfqpoint{1.879288in}{1.258175in}}%
\pgfpathlineto{\pgfqpoint{1.880983in}{1.250541in}}%
\pgfpathlineto{\pgfqpoint{1.882653in}{1.244927in}}%
\pgfpathlineto{\pgfqpoint{1.892105in}{1.231679in}}%
\pgfpathlineto{\pgfqpoint{1.898780in}{1.218431in}}%
\pgfpathlineto{\pgfqpoint{1.903644in}{1.205183in}}%
\pgfpathlineto{\pgfqpoint{1.906266in}{1.201892in}}%
\pgfpathlineto{\pgfqpoint{1.915896in}{1.191935in}}%
\pgfpathlineto{\pgfqpoint{1.918358in}{1.178687in}}%
\pgfpathlineto{\pgfqpoint{1.930656in}{1.165439in}}%
\pgfpathlineto{\pgfqpoint{1.931549in}{1.161486in}}%
\pgfpathlineto{\pgfqpoint{1.934236in}{1.152191in}}%
\pgfpathlineto{\pgfqpoint{1.947622in}{1.138943in}}%
\pgfpathlineto{\pgfqpoint{1.955333in}{1.125695in}}%
\pgfpathlineto{\pgfqpoint{1.956832in}{1.122806in}}%
\pgfpathlineto{\pgfqpoint{1.963177in}{1.112447in}}%
\pgfpathlineto{\pgfqpoint{1.976091in}{1.099199in}}%
\pgfpathlineto{\pgfqpoint{1.978555in}{1.085951in}}%
\pgfpathlineto{\pgfqpoint{1.982115in}{1.083062in}}%
\pgfpathlineto{\pgfqpoint{1.996392in}{1.072703in}}%
\pgfpathlineto{\pgfqpoint{1.998878in}{1.059455in}}%
\pgfpathlineto{\pgfqpoint{2.007399in}{1.051248in}}%
\pgfpathlineto{\pgfqpoint{2.013185in}{1.046207in}}%
\pgfpathlineto{\pgfqpoint{2.021038in}{1.032959in}}%
\pgfpathlineto{\pgfqpoint{2.031632in}{1.019711in}}%
\pgfpathlineto{\pgfqpoint{2.032682in}{1.018783in}}%
\pgfpathlineto{\pgfqpoint{2.048556in}{1.006463in}}%
\pgfpathlineto{\pgfqpoint{2.053620in}{0.993215in}}%
\pgfpathlineto{\pgfqpoint{2.057965in}{0.990297in}}%
\pgfpathlineto{\pgfqpoint{2.074534in}{0.979967in}}%
\pgfpathlineto{\pgfqpoint{2.076170in}{0.966719in}}%
\pgfpathlineto{\pgfqpoint{2.083248in}{0.962503in}}%
\pgfpathlineto{\pgfqpoint{2.099237in}{0.953471in}}%
\pgfpathlineto{\pgfqpoint{2.105938in}{0.940223in}}%
\pgfpathlineto{\pgfqpoint{2.108531in}{0.937669in}}%
\pgfpathlineto{\pgfqpoint{2.120196in}{0.926975in}}%
\pgfpathlineto{\pgfqpoint{2.133815in}{0.915371in}}%
\pgfpathlineto{\pgfqpoint{2.135980in}{0.913727in}}%
\pgfpathlineto{\pgfqpoint{2.141700in}{0.900479in}}%
\pgfpathlineto{\pgfqpoint{2.159098in}{0.892042in}}%
\pgfpathlineto{\pgfqpoint{2.169803in}{0.887231in}}%
\pgfpathlineto{\pgfqpoint{2.170717in}{0.873983in}}%
\pgfpathlineto{\pgfqpoint{2.184381in}{0.867391in}}%
\pgfpathlineto{\pgfqpoint{2.198753in}{0.860735in}}%
\pgfpathlineto{\pgfqpoint{2.200322in}{0.847487in}}%
\pgfpathlineto{\pgfqpoint{2.209664in}{0.842533in}}%
\pgfpathlineto{\pgfqpoint{2.225822in}{0.834239in}}%
\pgfpathlineto{\pgfqpoint{2.234947in}{0.826348in}}%
\pgfpathlineto{\pgfqpoint{2.242072in}{0.820991in}}%
\pgfpathlineto{\pgfqpoint{2.251861in}{0.807743in}}%
\pgfpathlineto{\pgfqpoint{2.260230in}{0.803765in}}%
\pgfpathlineto{\pgfqpoint{2.281567in}{0.794495in}}%
\pgfpathlineto{\pgfqpoint{2.273094in}{0.781247in}}%
\pgfpathlineto{\pgfqpoint{2.285514in}{0.777306in}}%
\pgfpathlineto{\pgfqpoint{2.310797in}{0.769761in}}%
\pgfpathlineto{\pgfqpoint{2.317005in}{0.767999in}}%
\pgfpathlineto{\pgfqpoint{2.310797in}{0.762061in}}%
\pgfpathlineto{\pgfqpoint{2.302260in}{0.754751in}}%
\pgfpathlineto{\pgfqpoint{2.310797in}{0.750833in}}%
\pgfpathlineto{\pgfqpoint{2.332261in}{0.741503in}}%
\pgfpathlineto{\pgfqpoint{2.315785in}{0.728255in}}%
\pgfpathlineto{\pgfqpoint{2.323477in}{0.715007in}}%
\pgfpathlineto{\pgfqpoint{2.330586in}{0.701759in}}%
\pgfpathlineto{\pgfqpoint{2.310797in}{0.694865in}}%
\pgfpathlineto{\pgfqpoint{2.289972in}{0.688511in}}%
\pgfpathlineto{\pgfqpoint{2.310797in}{0.675492in}}%
\pgfpathlineto{\pgfqpoint{2.311236in}{0.675263in}}%
\pgfpathlineto{\pgfqpoint{2.310797in}{0.675213in}}%
\pgfpathlineto{\pgfqpoint{2.285514in}{0.672452in}}%
\pgfpathlineto{\pgfqpoint{2.260230in}{0.669945in}}%
\pgfpathlineto{\pgfqpoint{2.234947in}{0.667659in}}%
\pgfpathlineto{\pgfqpoint{2.209664in}{0.665569in}}%
\pgfpathlineto{\pgfqpoint{2.184381in}{0.663650in}}%
\pgfpathlineto{\pgfqpoint{2.161055in}{0.662015in}}%
\pgfpathlineto{\pgfqpoint{2.179282in}{0.648767in}}%
\pgfpathlineto{\pgfqpoint{2.159098in}{0.647573in}}%
\pgfpathlineto{\pgfqpoint{2.133815in}{0.646193in}}%
\pgfpathlineto{\pgfqpoint{2.108531in}{0.644928in}}%
\pgfpathlineto{\pgfqpoint{2.083248in}{0.643766in}}%
\pgfpathlineto{\pgfqpoint{2.057965in}{0.642697in}}%
\pgfpathlineto{\pgfqpoint{2.032682in}{0.641714in}}%
\pgfpathlineto{\pgfqpoint{2.007399in}{0.640807in}}%
\pgfpathlineto{\pgfqpoint{1.982115in}{0.639972in}}%
\pgfpathlineto{\pgfqpoint{1.956832in}{0.639201in}}%
\pgfpathlineto{\pgfqpoint{1.931549in}{0.638490in}}%
\pgfpathlineto{\pgfqpoint{1.906266in}{0.637833in}}%
\pgfpathlineto{\pgfqpoint{1.880983in}{0.637227in}}%
\pgfpathlineto{\pgfqpoint{1.855700in}{0.636667in}}%
\pgfpathlineto{\pgfqpoint{1.830416in}{0.636151in}}%
\pgfpathlineto{\pgfqpoint{1.805133in}{0.635675in}}%
\pgfpathlineto{\pgfqpoint{1.796204in}{0.635519in}}%
\pgfpathlineto{\pgfqpoint{1.779850in}{0.634255in}}%
\pgfpathlineto{\pgfqpoint{1.754567in}{0.632445in}}%
\pgfpathlineto{\pgfqpoint{1.729284in}{0.630780in}}%
\pgfpathlineto{\pgfqpoint{1.704001in}{0.629249in}}%
\pgfpathlineto{\pgfqpoint{1.678717in}{0.627847in}}%
\pgfpathlineto{\pgfqpoint{1.653434in}{0.626565in}}%
\pgfpathlineto{\pgfqpoint{1.628151in}{0.625398in}}%
\pgfpathlineto{\pgfqpoint{1.602868in}{0.624342in}}%
\pgfpathlineto{\pgfqpoint{1.577585in}{0.623393in}}%
\pgfpathlineto{\pgfqpoint{1.552302in}{0.622548in}}%
\pgfpathlineto{\pgfqpoint{1.542886in}{0.622271in}}%
\pgfpathlineto{\pgfqpoint{1.527018in}{0.622165in}}%
\pgfpathlineto{\pgfqpoint{1.501735in}{0.622020in}}%
\pgfpathlineto{\pgfqpoint{1.476452in}{0.621900in}}%
\pgfpathlineto{\pgfqpoint{1.451169in}{0.621806in}}%
\pgfpathlineto{\pgfqpoint{1.425886in}{0.621735in}}%
\pgfpathlineto{\pgfqpoint{1.400603in}{0.621688in}}%
\pgfpathlineto{\pgfqpoint{1.375319in}{0.621665in}}%
\pgfpathlineto{\pgfqpoint{1.350036in}{0.621664in}}%
\pgfpathlineto{\pgfqpoint{1.324753in}{0.621687in}}%
\pgfpathlineto{\pgfqpoint{1.299470in}{0.621733in}}%
\pgfpathlineto{\pgfqpoint{1.274187in}{0.621801in}}%
\pgfpathlineto{\pgfqpoint{1.248904in}{0.621893in}}%
\pgfpathlineto{\pgfqpoint{1.223620in}{0.622008in}}%
\pgfpathlineto{\pgfqpoint{1.198337in}{0.622147in}}%
\pgfpathlineto{\pgfqpoint{1.178998in}{0.622271in}}%
\pgfpathlineto{\pgfqpoint{1.173054in}{0.622470in}}%
\pgfpathlineto{\pgfqpoint{1.147771in}{0.623458in}}%
\pgfpathlineto{\pgfqpoint{1.122488in}{0.624610in}}%
\pgfpathlineto{\pgfqpoint{1.097204in}{0.625938in}}%
\pgfpathlineto{\pgfqpoint{1.071921in}{0.627455in}}%
\pgfpathlineto{\pgfqpoint{1.046638in}{0.629175in}}%
\pgfpathlineto{\pgfqpoint{1.021355in}{0.631116in}}%
\pgfpathlineto{\pgfqpoint{0.996072in}{0.633295in}}%
\pgfpathlineto{\pgfqpoint{0.972946in}{0.635519in}}%
\pgfpathlineto{\pgfqpoint{0.970789in}{0.635566in}}%
\pgfpathlineto{\pgfqpoint{0.945505in}{0.636171in}}%
\pgfpathlineto{\pgfqpoint{0.920222in}{0.636855in}}%
\pgfpathlineto{\pgfqpoint{0.894939in}{0.637628in}}%
\pgfpathlineto{\pgfqpoint{0.869656in}{0.638505in}}%
\pgfpathlineto{\pgfqpoint{0.844373in}{0.639500in}}%
\pgfpathlineto{\pgfqpoint{0.819090in}{0.640635in}}%
\pgfpathlineto{\pgfqpoint{0.793806in}{0.641932in}}%
\pgfpathlineto{\pgfqpoint{0.768523in}{0.643424in}}%
\pgfpathlineto{\pgfqpoint{0.743240in}{0.645148in}}%
\pgfpathlineto{\pgfqpoint{0.717957in}{0.647154in}}%
\pgfpathlineto{\pgfqpoint{0.700238in}{0.648767in}}%
\pgfpathlineto{\pgfqpoint{0.692674in}{0.649374in}}%
\pgfpathlineto{\pgfqpoint{0.667391in}{0.651493in}}%
\pgfpathlineto{\pgfqpoint{0.642107in}{0.653737in}}%
\pgfpathlineto{\pgfqpoint{0.616824in}{0.656121in}}%
\pgfpathlineto{\pgfqpoint{0.591541in}{0.658656in}}%
\pgfpathlineto{\pgfqpoint{0.566258in}{0.661360in}}%
\pgfpathlineto{\pgfqpoint{0.560402in}{0.662015in}}%
\pgfpathlineto{\pgfqpoint{0.566258in}{0.664037in}}%
\pgfpathlineto{\pgfqpoint{0.591541in}{0.675172in}}%
\pgfpathlineto{\pgfqpoint{0.591709in}{0.675263in}}%
\pgfpathlineto{\pgfqpoint{0.591541in}{0.675278in}}%
\pgfpathlineto{\pgfqpoint{0.566258in}{0.677589in}}%
\pgfpathlineto{\pgfqpoint{0.540975in}{0.680137in}}%
\pgfpathlineto{\pgfqpoint{0.540975in}{0.675263in}}%
\pgfpathlineto{\pgfqpoint{0.540975in}{0.662015in}}%
\pgfpathlineto{\pgfqpoint{0.540975in}{0.648767in}}%
\pgfpathlineto{\pgfqpoint{0.540975in}{0.635519in}}%
\pgfpathlineto{\pgfqpoint{0.540975in}{0.622271in}}%
\pgfpathlineto{\pgfqpoint{0.540975in}{0.609023in}}%
\pgfpathlineto{\pgfqpoint{0.540975in}{0.595775in}}%
\pgfpathlineto{\pgfqpoint{0.540975in}{0.582527in}}%
\pgfpathlineto{\pgfqpoint{0.540975in}{0.569279in}}%
\pgfpathlineto{\pgfqpoint{0.540975in}{0.556031in}}%
\pgfpathlineto{\pgfqpoint{0.540975in}{0.542783in}}%
\pgfpathlineto{\pgfqpoint{0.540975in}{0.529535in}}%
\pgfpathlineto{\pgfqpoint{0.540975in}{0.516287in}}%
\pgfpathlineto{\pgfqpoint{0.540975in}{0.503039in}}%
\pgfpathlineto{\pgfqpoint{0.540975in}{0.489791in}}%
\pgfpathlineto{\pgfqpoint{0.540975in}{0.476543in}}%
\pgfpathlineto{\pgfqpoint{0.566258in}{0.476543in}}%
\pgfpathclose%
\pgfpathmoveto{\pgfqpoint{2.284420in}{1.695359in}}%
\pgfpathlineto{\pgfqpoint{2.285514in}{1.696742in}}%
\pgfpathlineto{\pgfqpoint{2.310797in}{1.698535in}}%
\pgfpathlineto{\pgfqpoint{2.316086in}{1.695359in}}%
\pgfpathlineto{\pgfqpoint{2.310797in}{1.691110in}}%
\pgfpathlineto{\pgfqpoint{2.285514in}{1.694723in}}%
\pgfpathlineto{\pgfqpoint{2.284420in}{1.695359in}}%
\pgfpathclose%
\pgfpathmoveto{\pgfqpoint{0.558099in}{0.741503in}}%
\pgfpathlineto{\pgfqpoint{0.540975in}{0.751581in}}%
\pgfpathlineto{\pgfqpoint{0.540975in}{0.741503in}}%
\pgfpathlineto{\pgfqpoint{0.540975in}{0.734579in}}%
\pgfpathlineto{\pgfqpoint{0.558099in}{0.741503in}}%
\pgfpathclose%
\pgfpathmoveto{\pgfqpoint{0.566258in}{0.761233in}}%
\pgfpathlineto{\pgfqpoint{0.591541in}{0.766981in}}%
\pgfpathlineto{\pgfqpoint{0.595777in}{0.767999in}}%
\pgfpathlineto{\pgfqpoint{0.591541in}{0.771843in}}%
\pgfpathlineto{\pgfqpoint{0.582229in}{0.781247in}}%
\pgfpathlineto{\pgfqpoint{0.591541in}{0.784655in}}%
\pgfpathlineto{\pgfqpoint{0.616824in}{0.794040in}}%
\pgfpathlineto{\pgfqpoint{0.618020in}{0.794495in}}%
\pgfpathlineto{\pgfqpoint{0.624906in}{0.807743in}}%
\pgfpathlineto{\pgfqpoint{0.633268in}{0.820991in}}%
\pgfpathlineto{\pgfqpoint{0.642107in}{0.823955in}}%
\pgfpathlineto{\pgfqpoint{0.667391in}{0.833330in}}%
\pgfpathlineto{\pgfqpoint{0.669599in}{0.834239in}}%
\pgfpathlineto{\pgfqpoint{0.668390in}{0.847487in}}%
\pgfpathlineto{\pgfqpoint{0.692674in}{0.856197in}}%
\pgfpathlineto{\pgfqpoint{0.704955in}{0.860735in}}%
\pgfpathlineto{\pgfqpoint{0.704643in}{0.873983in}}%
\pgfpathlineto{\pgfqpoint{0.717957in}{0.880341in}}%
\pgfpathlineto{\pgfqpoint{0.732316in}{0.887231in}}%
\pgfpathlineto{\pgfqpoint{0.743240in}{0.897206in}}%
\pgfpathlineto{\pgfqpoint{0.746278in}{0.900479in}}%
\pgfpathlineto{\pgfqpoint{0.757720in}{0.913727in}}%
\pgfpathlineto{\pgfqpoint{0.768523in}{0.919801in}}%
\pgfpathlineto{\pgfqpoint{0.780081in}{0.926975in}}%
\pgfpathlineto{\pgfqpoint{0.779536in}{0.940223in}}%
\pgfpathlineto{\pgfqpoint{0.793806in}{0.946720in}}%
\pgfpathlineto{\pgfqpoint{0.808057in}{0.953471in}}%
\pgfpathlineto{\pgfqpoint{0.810833in}{0.966719in}}%
\pgfpathlineto{\pgfqpoint{0.819090in}{0.972006in}}%
\pgfpathlineto{\pgfqpoint{0.831025in}{0.979967in}}%
\pgfpathlineto{\pgfqpoint{0.840901in}{0.993215in}}%
\pgfpathlineto{\pgfqpoint{0.844373in}{0.998067in}}%
\pgfpathlineto{\pgfqpoint{0.849847in}{1.006463in}}%
\pgfpathlineto{\pgfqpoint{0.869656in}{1.019705in}}%
\pgfpathlineto{\pgfqpoint{0.869664in}{1.019711in}}%
\pgfpathlineto{\pgfqpoint{0.872527in}{1.032959in}}%
\pgfpathlineto{\pgfqpoint{0.892175in}{1.046207in}}%
\pgfpathlineto{\pgfqpoint{0.894046in}{1.059455in}}%
\pgfpathlineto{\pgfqpoint{0.894939in}{1.060153in}}%
\pgfpathlineto{\pgfqpoint{0.909863in}{1.072703in}}%
\pgfpathlineto{\pgfqpoint{0.918866in}{1.085951in}}%
\pgfpathlineto{\pgfqpoint{0.920222in}{1.087727in}}%
\pgfpathlineto{\pgfqpoint{0.927911in}{1.099199in}}%
\pgfpathlineto{\pgfqpoint{0.940757in}{1.112447in}}%
\pgfpathlineto{\pgfqpoint{0.943422in}{1.125695in}}%
\pgfpathlineto{\pgfqpoint{0.945505in}{1.127238in}}%
\pgfpathlineto{\pgfqpoint{0.959355in}{1.138943in}}%
\pgfpathlineto{\pgfqpoint{0.961464in}{1.152191in}}%
\pgfpathlineto{\pgfqpoint{0.970789in}{1.160907in}}%
\pgfpathlineto{\pgfqpoint{0.975048in}{1.165439in}}%
\pgfpathlineto{\pgfqpoint{0.979896in}{1.178687in}}%
\pgfpathlineto{\pgfqpoint{0.987576in}{1.191935in}}%
\pgfpathlineto{\pgfqpoint{0.996072in}{1.200981in}}%
\pgfpathlineto{\pgfqpoint{0.999247in}{1.205183in}}%
\pgfpathlineto{\pgfqpoint{1.003103in}{1.218431in}}%
\pgfpathlineto{\pgfqpoint{1.015483in}{1.231679in}}%
\pgfpathlineto{\pgfqpoint{1.017294in}{1.244927in}}%
\pgfpathlineto{\pgfqpoint{1.021355in}{1.249299in}}%
\pgfpathlineto{\pgfqpoint{1.028191in}{1.258175in}}%
\pgfpathlineto{\pgfqpoint{1.032011in}{1.271423in}}%
\pgfpathlineto{\pgfqpoint{1.039691in}{1.284671in}}%
\pgfpathlineto{\pgfqpoint{1.046638in}{1.295818in}}%
\pgfpathlineto{\pgfqpoint{1.047629in}{1.297919in}}%
\pgfpathlineto{\pgfqpoint{1.050250in}{1.311167in}}%
\pgfpathlineto{\pgfqpoint{1.060560in}{1.324415in}}%
\pgfpathlineto{\pgfqpoint{1.063074in}{1.337663in}}%
\pgfpathlineto{\pgfqpoint{1.071921in}{1.348089in}}%
\pgfpathlineto{\pgfqpoint{1.073812in}{1.350911in}}%
\pgfpathlineto{\pgfqpoint{1.076237in}{1.364159in}}%
\pgfpathlineto{\pgfqpoint{1.082348in}{1.377407in}}%
\pgfpathlineto{\pgfqpoint{1.088552in}{1.390655in}}%
\pgfpathlineto{\pgfqpoint{1.092617in}{1.403903in}}%
\pgfpathlineto{\pgfqpoint{1.097204in}{1.410314in}}%
\pgfpathlineto{\pgfqpoint{1.100764in}{1.417151in}}%
\pgfpathlineto{\pgfqpoint{1.101980in}{1.430399in}}%
\pgfpathlineto{\pgfqpoint{1.110074in}{1.443647in}}%
\pgfpathlineto{\pgfqpoint{1.112734in}{1.456895in}}%
\pgfpathlineto{\pgfqpoint{1.120559in}{1.470143in}}%
\pgfpathlineto{\pgfqpoint{1.122488in}{1.475022in}}%
\pgfpathlineto{\pgfqpoint{1.124636in}{1.483391in}}%
\pgfpathlineto{\pgfqpoint{1.127520in}{1.496639in}}%
\pgfpathlineto{\pgfqpoint{1.133993in}{1.509887in}}%
\pgfpathlineto{\pgfqpoint{1.135878in}{1.523135in}}%
\pgfpathlineto{\pgfqpoint{1.144326in}{1.536383in}}%
\pgfpathlineto{\pgfqpoint{1.146413in}{1.549631in}}%
\pgfpathlineto{\pgfqpoint{1.147771in}{1.552028in}}%
\pgfpathlineto{\pgfqpoint{1.151814in}{1.562879in}}%
\pgfpathlineto{\pgfqpoint{1.155047in}{1.576127in}}%
\pgfpathlineto{\pgfqpoint{1.158735in}{1.589375in}}%
\pgfpathlineto{\pgfqpoint{1.163934in}{1.602623in}}%
\pgfpathlineto{\pgfqpoint{1.165633in}{1.615871in}}%
\pgfpathlineto{\pgfqpoint{1.173054in}{1.627885in}}%
\pgfpathlineto{\pgfqpoint{1.173548in}{1.629119in}}%
\pgfpathlineto{\pgfqpoint{1.174799in}{1.642367in}}%
\pgfpathlineto{\pgfqpoint{1.179858in}{1.655615in}}%
\pgfpathlineto{\pgfqpoint{1.182178in}{1.668863in}}%
\pgfpathlineto{\pgfqpoint{1.185784in}{1.682111in}}%
\pgfpathlineto{\pgfqpoint{1.190956in}{1.695359in}}%
\pgfpathlineto{\pgfqpoint{1.193518in}{1.708607in}}%
\pgfpathlineto{\pgfqpoint{1.198337in}{1.717308in}}%
\pgfpathlineto{\pgfqpoint{1.199824in}{1.721855in}}%
\pgfpathlineto{\pgfqpoint{1.200577in}{1.735103in}}%
\pgfpathlineto{\pgfqpoint{1.204768in}{1.748351in}}%
\pgfpathlineto{\pgfqpoint{1.206652in}{1.761599in}}%
\pgfpathlineto{\pgfqpoint{1.210621in}{1.774847in}}%
\pgfpathlineto{\pgfqpoint{1.214294in}{1.788095in}}%
\pgfpathlineto{\pgfqpoint{1.216795in}{1.801343in}}%
\pgfpathlineto{\pgfqpoint{1.223620in}{1.814331in}}%
\pgfpathlineto{\pgfqpoint{1.223697in}{1.814591in}}%
\pgfpathlineto{\pgfqpoint{1.224683in}{1.827839in}}%
\pgfpathlineto{\pgfqpoint{1.228731in}{1.841087in}}%
\pgfpathlineto{\pgfqpoint{1.229805in}{1.854335in}}%
\pgfpathlineto{\pgfqpoint{1.233062in}{1.867583in}}%
\pgfpathlineto{\pgfqpoint{1.235949in}{1.880831in}}%
\pgfpathlineto{\pgfqpoint{1.238727in}{1.894079in}}%
\pgfpathlineto{\pgfqpoint{1.243693in}{1.907327in}}%
\pgfpathlineto{\pgfqpoint{1.245393in}{1.920575in}}%
\pgfpathlineto{\pgfqpoint{1.248904in}{1.927414in}}%
\pgfpathlineto{\pgfqpoint{1.250549in}{1.933823in}}%
\pgfpathlineto{\pgfqpoint{1.251528in}{1.947071in}}%
\pgfpathlineto{\pgfqpoint{1.254906in}{1.960319in}}%
\pgfpathlineto{\pgfqpoint{1.256574in}{1.973567in}}%
\pgfpathlineto{\pgfqpoint{1.258901in}{1.986815in}}%
\pgfpathlineto{\pgfqpoint{1.262510in}{2.000063in}}%
\pgfpathlineto{\pgfqpoint{1.264369in}{2.013311in}}%
\pgfpathlineto{\pgfqpoint{1.270076in}{2.026559in}}%
\pgfpathlineto{\pgfqpoint{1.271623in}{2.039807in}}%
\pgfpathlineto{\pgfqpoint{1.274187in}{2.045413in}}%
\pgfpathlineto{\pgfqpoint{1.275763in}{2.053055in}}%
\pgfpathlineto{\pgfqpoint{1.277052in}{2.066303in}}%
\pgfpathlineto{\pgfqpoint{1.279433in}{2.079551in}}%
\pgfpathlineto{\pgfqpoint{1.281649in}{2.092799in}}%
\pgfpathlineto{\pgfqpoint{1.283140in}{2.106047in}}%
\pgfpathlineto{\pgfqpoint{1.286817in}{2.119295in}}%
\pgfpathlineto{\pgfqpoint{1.288075in}{2.132543in}}%
\pgfpathlineto{\pgfqpoint{1.274187in}{2.132543in}}%
\pgfpathlineto{\pgfqpoint{1.248904in}{2.132543in}}%
\pgfpathlineto{\pgfqpoint{1.237050in}{2.132543in}}%
\pgfpathlineto{\pgfqpoint{1.235144in}{2.119295in}}%
\pgfpathlineto{\pgfqpoint{1.228725in}{2.106047in}}%
\pgfpathlineto{\pgfqpoint{1.226888in}{2.092799in}}%
\pgfpathlineto{\pgfqpoint{1.223620in}{2.087068in}}%
\pgfpathlineto{\pgfqpoint{1.220954in}{2.079551in}}%
\pgfpathlineto{\pgfqpoint{1.219357in}{2.066303in}}%
\pgfpathlineto{\pgfqpoint{1.216086in}{2.053055in}}%
\pgfpathlineto{\pgfqpoint{1.212190in}{2.039807in}}%
\pgfpathlineto{\pgfqpoint{1.209551in}{2.026559in}}%
\pgfpathlineto{\pgfqpoint{1.203538in}{2.013311in}}%
\pgfpathlineto{\pgfqpoint{1.201996in}{2.000063in}}%
\pgfpathlineto{\pgfqpoint{1.198337in}{1.993974in}}%
\pgfpathlineto{\pgfqpoint{1.195338in}{1.986815in}}%
\pgfpathlineto{\pgfqpoint{1.193570in}{1.973567in}}%
\pgfpathlineto{\pgfqpoint{1.188729in}{1.960319in}}%
\pgfpathlineto{\pgfqpoint{1.185436in}{1.947071in}}%
\pgfpathlineto{\pgfqpoint{1.182357in}{1.933823in}}%
\pgfpathlineto{\pgfqpoint{1.175208in}{1.920575in}}%
\pgfpathlineto{\pgfqpoint{1.173054in}{1.910052in}}%
\pgfpathlineto{\pgfqpoint{1.172667in}{1.907327in}}%
\pgfpathlineto{\pgfqpoint{1.166181in}{1.894079in}}%
\pgfpathlineto{\pgfqpoint{1.164815in}{1.880831in}}%
\pgfpathlineto{\pgfqpoint{1.159069in}{1.867583in}}%
\pgfpathlineto{\pgfqpoint{1.155310in}{1.854335in}}%
\pgfpathlineto{\pgfqpoint{1.150293in}{1.841087in}}%
\pgfpathlineto{\pgfqpoint{1.147771in}{1.836250in}}%
\pgfpathlineto{\pgfqpoint{1.144362in}{1.827839in}}%
\pgfpathlineto{\pgfqpoint{1.142386in}{1.814591in}}%
\pgfpathlineto{\pgfqpoint{1.134793in}{1.801343in}}%
\pgfpathlineto{\pgfqpoint{1.132626in}{1.788095in}}%
\pgfpathlineto{\pgfqpoint{1.123406in}{1.774847in}}%
\pgfpathlineto{\pgfqpoint{1.122488in}{1.771539in}}%
\pgfpathlineto{\pgfqpoint{1.120336in}{1.761599in}}%
\pgfpathlineto{\pgfqpoint{1.115403in}{1.748351in}}%
\pgfpathlineto{\pgfqpoint{1.109173in}{1.735103in}}%
\pgfpathlineto{\pgfqpoint{1.105930in}{1.721855in}}%
\pgfpathlineto{\pgfqpoint{1.097204in}{1.709558in}}%
\pgfpathlineto{\pgfqpoint{1.096576in}{1.708607in}}%
\pgfpathlineto{\pgfqpoint{1.094930in}{1.695359in}}%
\pgfpathlineto{\pgfqpoint{1.086029in}{1.682111in}}%
\pgfpathlineto{\pgfqpoint{1.082846in}{1.668863in}}%
\pgfpathlineto{\pgfqpoint{1.074003in}{1.655615in}}%
\pgfpathlineto{\pgfqpoint{1.071921in}{1.651497in}}%
\pgfpathlineto{\pgfqpoint{1.067856in}{1.642367in}}%
\pgfpathlineto{\pgfqpoint{1.063886in}{1.629119in}}%
\pgfpathlineto{\pgfqpoint{1.053851in}{1.615871in}}%
\pgfpathlineto{\pgfqpoint{1.051440in}{1.602623in}}%
\pgfpathlineto{\pgfqpoint{1.046638in}{1.597392in}}%
\pgfpathlineto{\pgfqpoint{1.039213in}{1.589375in}}%
\pgfpathlineto{\pgfqpoint{1.036611in}{1.576127in}}%
\pgfpathlineto{\pgfqpoint{1.026546in}{1.562879in}}%
\pgfpathlineto{\pgfqpoint{1.021355in}{1.552416in}}%
\pgfpathlineto{\pgfqpoint{1.020061in}{1.549631in}}%
\pgfpathlineto{\pgfqpoint{1.012875in}{1.536383in}}%
\pgfpathlineto{\pgfqpoint{1.002220in}{1.523135in}}%
\pgfpathlineto{\pgfqpoint{0.999726in}{1.509887in}}%
\pgfpathlineto{\pgfqpoint{0.996072in}{1.506352in}}%
\pgfpathlineto{\pgfqpoint{0.985476in}{1.496639in}}%
\pgfpathlineto{\pgfqpoint{0.982975in}{1.483391in}}%
\pgfpathlineto{\pgfqpoint{0.970789in}{1.472531in}}%
\pgfpathlineto{\pgfqpoint{0.967736in}{1.470143in}}%
\pgfpathlineto{\pgfqpoint{0.961835in}{1.456895in}}%
\pgfpathlineto{\pgfqpoint{0.951828in}{1.443647in}}%
\pgfpathlineto{\pgfqpoint{0.945505in}{1.436947in}}%
\pgfpathlineto{\pgfqpoint{0.939163in}{1.430399in}}%
\pgfpathlineto{\pgfqpoint{0.933954in}{1.417151in}}%
\pgfpathlineto{\pgfqpoint{0.920222in}{1.406867in}}%
\pgfpathlineto{\pgfqpoint{0.915806in}{1.403903in}}%
\pgfpathlineto{\pgfqpoint{0.914681in}{1.390655in}}%
\pgfpathlineto{\pgfqpoint{0.896250in}{1.377407in}}%
\pgfpathlineto{\pgfqpoint{0.894939in}{1.374002in}}%
\pgfpathlineto{\pgfqpoint{0.891160in}{1.364159in}}%
\pgfpathlineto{\pgfqpoint{0.874112in}{1.350911in}}%
\pgfpathlineto{\pgfqpoint{0.869656in}{1.346375in}}%
\pgfpathlineto{\pgfqpoint{0.860736in}{1.337663in}}%
\pgfpathlineto{\pgfqpoint{0.852798in}{1.324415in}}%
\pgfpathlineto{\pgfqpoint{0.844373in}{1.319418in}}%
\pgfpathlineto{\pgfqpoint{0.828797in}{1.311167in}}%
\pgfpathlineto{\pgfqpoint{0.826761in}{1.297919in}}%
\pgfpathlineto{\pgfqpoint{0.819090in}{1.294139in}}%
\pgfpathlineto{\pgfqpoint{0.795991in}{1.284671in}}%
\pgfpathlineto{\pgfqpoint{0.795026in}{1.271423in}}%
\pgfpathlineto{\pgfqpoint{0.793806in}{1.270612in}}%
\pgfpathlineto{\pgfqpoint{0.770160in}{1.258175in}}%
\pgfpathlineto{\pgfqpoint{0.768523in}{1.256544in}}%
\pgfpathlineto{\pgfqpoint{0.756503in}{1.244927in}}%
\pgfpathlineto{\pgfqpoint{0.743240in}{1.234973in}}%
\pgfpathlineto{\pgfqpoint{0.737332in}{1.231679in}}%
\pgfpathlineto{\pgfqpoint{0.717957in}{1.223684in}}%
\pgfpathlineto{\pgfqpoint{0.704069in}{1.218431in}}%
\pgfpathlineto{\pgfqpoint{0.704701in}{1.205183in}}%
\pgfpathlineto{\pgfqpoint{0.692674in}{1.201496in}}%
\pgfpathlineto{\pgfqpoint{0.667391in}{1.195348in}}%
\pgfpathlineto{\pgfqpoint{0.649277in}{1.191935in}}%
\pgfpathlineto{\pgfqpoint{0.661769in}{1.178687in}}%
\pgfpathlineto{\pgfqpoint{0.642107in}{1.173984in}}%
\pgfpathlineto{\pgfqpoint{0.616824in}{1.169829in}}%
\pgfpathlineto{\pgfqpoint{0.591541in}{1.167506in}}%
\pgfpathlineto{\pgfqpoint{0.566258in}{1.167298in}}%
\pgfpathlineto{\pgfqpoint{0.540975in}{1.170021in}}%
\pgfpathlineto{\pgfqpoint{0.540975in}{1.165439in}}%
\pgfpathlineto{\pgfqpoint{0.540975in}{1.158551in}}%
\pgfpathlineto{\pgfqpoint{0.566258in}{1.160485in}}%
\pgfpathlineto{\pgfqpoint{0.591541in}{1.152563in}}%
\pgfpathlineto{\pgfqpoint{0.591847in}{1.152191in}}%
\pgfpathlineto{\pgfqpoint{0.591541in}{1.152082in}}%
\pgfpathlineto{\pgfqpoint{0.566258in}{1.147018in}}%
\pgfpathlineto{\pgfqpoint{0.540975in}{1.145891in}}%
\pgfpathlineto{\pgfqpoint{0.540975in}{1.138943in}}%
\pgfpathlineto{\pgfqpoint{0.540975in}{1.125695in}}%
\pgfpathlineto{\pgfqpoint{0.540975in}{1.112447in}}%
\pgfpathlineto{\pgfqpoint{0.540975in}{1.099199in}}%
\pgfpathlineto{\pgfqpoint{0.540975in}{1.085951in}}%
\pgfpathlineto{\pgfqpoint{0.540975in}{1.072703in}}%
\pgfpathlineto{\pgfqpoint{0.540975in}{1.059455in}}%
\pgfpathlineto{\pgfqpoint{0.540975in}{1.046207in}}%
\pgfpathlineto{\pgfqpoint{0.540975in}{1.032959in}}%
\pgfpathlineto{\pgfqpoint{0.540975in}{1.019711in}}%
\pgfpathlineto{\pgfqpoint{0.540975in}{1.006463in}}%
\pgfpathlineto{\pgfqpoint{0.540975in}{0.993215in}}%
\pgfpathlineto{\pgfqpoint{0.540975in}{0.979967in}}%
\pgfpathlineto{\pgfqpoint{0.540975in}{0.966719in}}%
\pgfpathlineto{\pgfqpoint{0.540975in}{0.953471in}}%
\pgfpathlineto{\pgfqpoint{0.540975in}{0.940223in}}%
\pgfpathlineto{\pgfqpoint{0.540975in}{0.926975in}}%
\pgfpathlineto{\pgfqpoint{0.540975in}{0.913727in}}%
\pgfpathlineto{\pgfqpoint{0.540975in}{0.900479in}}%
\pgfpathlineto{\pgfqpoint{0.540975in}{0.887231in}}%
\pgfpathlineto{\pgfqpoint{0.540975in}{0.873983in}}%
\pgfpathlineto{\pgfqpoint{0.540975in}{0.860735in}}%
\pgfpathlineto{\pgfqpoint{0.540975in}{0.847487in}}%
\pgfpathlineto{\pgfqpoint{0.540975in}{0.834239in}}%
\pgfpathlineto{\pgfqpoint{0.540975in}{0.820991in}}%
\pgfpathlineto{\pgfqpoint{0.540975in}{0.807743in}}%
\pgfpathlineto{\pgfqpoint{0.540975in}{0.794495in}}%
\pgfpathlineto{\pgfqpoint{0.540975in}{0.781247in}}%
\pgfpathlineto{\pgfqpoint{0.540975in}{0.767999in}}%
\pgfpathlineto{\pgfqpoint{0.540975in}{0.755797in}}%
\pgfpathlineto{\pgfqpoint{0.566258in}{0.761233in}}%
\pgfpathclose%
\pgfpathmoveto{\pgfqpoint{0.555755in}{1.191935in}}%
\pgfpathlineto{\pgfqpoint{0.544722in}{1.205183in}}%
\pgfpathlineto{\pgfqpoint{0.560732in}{1.218431in}}%
\pgfpathlineto{\pgfqpoint{0.566258in}{1.222329in}}%
\pgfpathlineto{\pgfqpoint{0.578353in}{1.231679in}}%
\pgfpathlineto{\pgfqpoint{0.566258in}{1.244475in}}%
\pgfpathlineto{\pgfqpoint{0.565932in}{1.244927in}}%
\pgfpathlineto{\pgfqpoint{0.566258in}{1.245076in}}%
\pgfpathlineto{\pgfqpoint{0.591541in}{1.256050in}}%
\pgfpathlineto{\pgfqpoint{0.596630in}{1.258175in}}%
\pgfpathlineto{\pgfqpoint{0.591541in}{1.263823in}}%
\pgfpathlineto{\pgfqpoint{0.586611in}{1.271423in}}%
\pgfpathlineto{\pgfqpoint{0.591541in}{1.274236in}}%
\pgfpathlineto{\pgfqpoint{0.611795in}{1.284671in}}%
\pgfpathlineto{\pgfqpoint{0.612928in}{1.297919in}}%
\pgfpathlineto{\pgfqpoint{0.616824in}{1.304541in}}%
\pgfpathlineto{\pgfqpoint{0.621106in}{1.311167in}}%
\pgfpathlineto{\pgfqpoint{0.638711in}{1.324415in}}%
\pgfpathlineto{\pgfqpoint{0.639563in}{1.337663in}}%
\pgfpathlineto{\pgfqpoint{0.642107in}{1.339137in}}%
\pgfpathlineto{\pgfqpoint{0.662558in}{1.350911in}}%
\pgfpathlineto{\pgfqpoint{0.658255in}{1.364159in}}%
\pgfpathlineto{\pgfqpoint{0.667391in}{1.371033in}}%
\pgfpathlineto{\pgfqpoint{0.676506in}{1.377407in}}%
\pgfpathlineto{\pgfqpoint{0.679995in}{1.390655in}}%
\pgfpathlineto{\pgfqpoint{0.690385in}{1.403903in}}%
\pgfpathlineto{\pgfqpoint{0.692674in}{1.406112in}}%
\pgfpathlineto{\pgfqpoint{0.703247in}{1.417151in}}%
\pgfpathlineto{\pgfqpoint{0.703379in}{1.430399in}}%
\pgfpathlineto{\pgfqpoint{0.717957in}{1.440746in}}%
\pgfpathlineto{\pgfqpoint{0.722024in}{1.443647in}}%
\pgfpathlineto{\pgfqpoint{0.721476in}{1.456895in}}%
\pgfpathlineto{\pgfqpoint{0.737862in}{1.470143in}}%
\pgfpathlineto{\pgfqpoint{0.740410in}{1.483391in}}%
\pgfpathlineto{\pgfqpoint{0.743240in}{1.488072in}}%
\pgfpathlineto{\pgfqpoint{0.748387in}{1.496639in}}%
\pgfpathlineto{\pgfqpoint{0.758919in}{1.509887in}}%
\pgfpathlineto{\pgfqpoint{0.762067in}{1.523135in}}%
\pgfpathlineto{\pgfqpoint{0.768523in}{1.528607in}}%
\pgfpathlineto{\pgfqpoint{0.777275in}{1.536383in}}%
\pgfpathlineto{\pgfqpoint{0.777252in}{1.549631in}}%
\pgfpathlineto{\pgfqpoint{0.790805in}{1.562879in}}%
\pgfpathlineto{\pgfqpoint{0.793806in}{1.575096in}}%
\pgfpathlineto{\pgfqpoint{0.794018in}{1.576127in}}%
\pgfpathlineto{\pgfqpoint{0.802751in}{1.589375in}}%
\pgfpathlineto{\pgfqpoint{0.810007in}{1.602623in}}%
\pgfpathlineto{\pgfqpoint{0.812239in}{1.615871in}}%
\pgfpathlineto{\pgfqpoint{0.819090in}{1.622986in}}%
\pgfpathlineto{\pgfqpoint{0.824470in}{1.629119in}}%
\pgfpathlineto{\pgfqpoint{0.825461in}{1.642367in}}%
\pgfpathlineto{\pgfqpoint{0.837858in}{1.655615in}}%
\pgfpathlineto{\pgfqpoint{0.840593in}{1.668863in}}%
\pgfpathlineto{\pgfqpoint{0.844373in}{1.674942in}}%
\pgfpathlineto{\pgfqpoint{0.848486in}{1.682111in}}%
\pgfpathlineto{\pgfqpoint{0.855088in}{1.695359in}}%
\pgfpathlineto{\pgfqpoint{0.858833in}{1.708607in}}%
\pgfpathlineto{\pgfqpoint{0.868937in}{1.721855in}}%
\pgfpathlineto{\pgfqpoint{0.869476in}{1.735103in}}%
\pgfpathlineto{\pgfqpoint{0.869656in}{1.735326in}}%
\pgfpathlineto{\pgfqpoint{0.878906in}{1.748351in}}%
\pgfpathlineto{\pgfqpoint{0.880939in}{1.761599in}}%
\pgfpathlineto{\pgfqpoint{0.889213in}{1.774847in}}%
\pgfpathlineto{\pgfqpoint{0.894939in}{1.787786in}}%
\pgfpathlineto{\pgfqpoint{0.895048in}{1.788095in}}%
\pgfpathlineto{\pgfqpoint{0.898416in}{1.801343in}}%
\pgfpathlineto{\pgfqpoint{0.906429in}{1.814591in}}%
\pgfpathlineto{\pgfqpoint{0.907847in}{1.827839in}}%
\pgfpathlineto{\pgfqpoint{0.917659in}{1.841087in}}%
\pgfpathlineto{\pgfqpoint{0.919050in}{1.854335in}}%
\pgfpathlineto{\pgfqpoint{0.920222in}{1.856386in}}%
\pgfpathlineto{\pgfqpoint{0.925509in}{1.867583in}}%
\pgfpathlineto{\pgfqpoint{0.929580in}{1.880831in}}%
\pgfpathlineto{\pgfqpoint{0.933822in}{1.894079in}}%
\pgfpathlineto{\pgfqpoint{0.940998in}{1.907327in}}%
\pgfpathlineto{\pgfqpoint{0.942660in}{1.920575in}}%
\pgfpathlineto{\pgfqpoint{0.945505in}{1.924784in}}%
\pgfpathlineto{\pgfqpoint{0.950322in}{1.933823in}}%
\pgfpathlineto{\pgfqpoint{0.951523in}{1.947071in}}%
\pgfpathlineto{\pgfqpoint{0.958374in}{1.960319in}}%
\pgfpathlineto{\pgfqpoint{0.961266in}{1.973567in}}%
\pgfpathlineto{\pgfqpoint{0.965576in}{1.986815in}}%
\pgfpathlineto{\pgfqpoint{0.970789in}{1.996469in}}%
\pgfpathlineto{\pgfqpoint{0.972219in}{2.000063in}}%
\pgfpathlineto{\pgfqpoint{0.974068in}{2.013311in}}%
\pgfpathlineto{\pgfqpoint{0.980612in}{2.026559in}}%
\pgfpathlineto{\pgfqpoint{0.981623in}{2.039807in}}%
\pgfpathlineto{\pgfqpoint{0.987550in}{2.053055in}}%
\pgfpathlineto{\pgfqpoint{0.990163in}{2.066303in}}%
\pgfpathlineto{\pgfqpoint{0.995514in}{2.079551in}}%
\pgfpathlineto{\pgfqpoint{0.996072in}{2.080826in}}%
\pgfpathlineto{\pgfqpoint{0.999556in}{2.092799in}}%
\pgfpathlineto{\pgfqpoint{1.001277in}{2.106047in}}%
\pgfpathlineto{\pgfqpoint{1.007010in}{2.119295in}}%
\pgfpathlineto{\pgfqpoint{1.008376in}{2.132543in}}%
\pgfpathlineto{\pgfqpoint{0.996072in}{2.132543in}}%
\pgfpathlineto{\pgfqpoint{0.970789in}{2.132543in}}%
\pgfpathlineto{\pgfqpoint{0.945505in}{2.132543in}}%
\pgfpathlineto{\pgfqpoint{0.920222in}{2.132543in}}%
\pgfpathlineto{\pgfqpoint{0.904060in}{2.132543in}}%
\pgfpathlineto{\pgfqpoint{0.900114in}{2.119295in}}%
\pgfpathlineto{\pgfqpoint{0.894939in}{2.112069in}}%
\pgfpathlineto{\pgfqpoint{0.890420in}{2.106047in}}%
\pgfpathlineto{\pgfqpoint{0.888639in}{2.092799in}}%
\pgfpathlineto{\pgfqpoint{0.878146in}{2.079551in}}%
\pgfpathlineto{\pgfqpoint{0.877094in}{2.066303in}}%
\pgfpathlineto{\pgfqpoint{0.869656in}{2.055365in}}%
\pgfpathlineto{\pgfqpoint{0.867746in}{2.053055in}}%
\pgfpathlineto{\pgfqpoint{0.862411in}{2.039807in}}%
\pgfpathlineto{\pgfqpoint{0.855714in}{2.026559in}}%
\pgfpathlineto{\pgfqpoint{0.847093in}{2.013311in}}%
\pgfpathlineto{\pgfqpoint{0.845673in}{2.000063in}}%
\pgfpathlineto{\pgfqpoint{0.844373in}{1.998618in}}%
\pgfpathlineto{\pgfqpoint{0.832402in}{1.986815in}}%
\pgfpathlineto{\pgfqpoint{0.831719in}{1.973567in}}%
\pgfpathlineto{\pgfqpoint{0.819745in}{1.960319in}}%
\pgfpathlineto{\pgfqpoint{0.819090in}{1.957709in}}%
\pgfpathlineto{\pgfqpoint{0.816257in}{1.947071in}}%
\pgfpathlineto{\pgfqpoint{0.807991in}{1.933823in}}%
\pgfpathlineto{\pgfqpoint{0.798185in}{1.920575in}}%
\pgfpathlineto{\pgfqpoint{0.793806in}{1.907363in}}%
\pgfpathlineto{\pgfqpoint{0.793791in}{1.907327in}}%
\pgfpathlineto{\pgfqpoint{0.777990in}{1.894079in}}%
\pgfpathlineto{\pgfqpoint{0.778633in}{1.880831in}}%
\pgfpathlineto{\pgfqpoint{0.768523in}{1.871454in}}%
\pgfpathlineto{\pgfqpoint{0.762769in}{1.867583in}}%
\pgfpathlineto{\pgfqpoint{0.761066in}{1.854335in}}%
\pgfpathlineto{\pgfqpoint{0.748453in}{1.841087in}}%
\pgfpathlineto{\pgfqpoint{0.743240in}{1.831675in}}%
\pgfpathlineto{\pgfqpoint{0.740714in}{1.827839in}}%
\pgfpathlineto{\pgfqpoint{0.735500in}{1.814591in}}%
\pgfpathlineto{\pgfqpoint{0.718416in}{1.801343in}}%
\pgfpathlineto{\pgfqpoint{0.717957in}{1.793979in}}%
\pgfpathlineto{\pgfqpoint{0.717481in}{1.788095in}}%
\pgfpathlineto{\pgfqpoint{0.693756in}{1.774847in}}%
\pgfpathlineto{\pgfqpoint{0.695652in}{1.761599in}}%
\pgfpathlineto{\pgfqpoint{0.692674in}{1.758799in}}%
\pgfpathlineto{\pgfqpoint{0.676045in}{1.748351in}}%
\pgfpathlineto{\pgfqpoint{0.670611in}{1.735103in}}%
\pgfpathlineto{\pgfqpoint{0.667391in}{1.730426in}}%
\pgfpathlineto{\pgfqpoint{0.657843in}{1.721855in}}%
\pgfpathlineto{\pgfqpoint{0.642107in}{1.710509in}}%
\pgfpathlineto{\pgfqpoint{0.638070in}{1.708607in}}%
\pgfpathlineto{\pgfqpoint{0.642107in}{1.696069in}}%
\pgfpathlineto{\pgfqpoint{0.642213in}{1.695359in}}%
\pgfpathlineto{\pgfqpoint{0.642107in}{1.695301in}}%
\pgfpathlineto{\pgfqpoint{0.616824in}{1.686111in}}%
\pgfpathlineto{\pgfqpoint{0.594304in}{1.682111in}}%
\pgfpathlineto{\pgfqpoint{0.615193in}{1.668863in}}%
\pgfpathlineto{\pgfqpoint{0.591541in}{1.661379in}}%
\pgfpathlineto{\pgfqpoint{0.566258in}{1.659851in}}%
\pgfpathlineto{\pgfqpoint{0.544615in}{1.668863in}}%
\pgfpathlineto{\pgfqpoint{0.566258in}{1.678671in}}%
\pgfpathlineto{\pgfqpoint{0.587213in}{1.682111in}}%
\pgfpathlineto{\pgfqpoint{0.566258in}{1.688865in}}%
\pgfpathlineto{\pgfqpoint{0.562179in}{1.695359in}}%
\pgfpathlineto{\pgfqpoint{0.566258in}{1.700893in}}%
\pgfpathlineto{\pgfqpoint{0.575568in}{1.708607in}}%
\pgfpathlineto{\pgfqpoint{0.583969in}{1.721855in}}%
\pgfpathlineto{\pgfqpoint{0.578273in}{1.735103in}}%
\pgfpathlineto{\pgfqpoint{0.591541in}{1.744231in}}%
\pgfpathlineto{\pgfqpoint{0.599413in}{1.748351in}}%
\pgfpathlineto{\pgfqpoint{0.591541in}{1.759383in}}%
\pgfpathlineto{\pgfqpoint{0.590654in}{1.761599in}}%
\pgfpathlineto{\pgfqpoint{0.591541in}{1.762282in}}%
\pgfpathlineto{\pgfqpoint{0.612282in}{1.774847in}}%
\pgfpathlineto{\pgfqpoint{0.610327in}{1.788095in}}%
\pgfpathlineto{\pgfqpoint{0.616824in}{1.799270in}}%
\pgfpathlineto{\pgfqpoint{0.618478in}{1.801343in}}%
\pgfpathlineto{\pgfqpoint{0.627902in}{1.814591in}}%
\pgfpathlineto{\pgfqpoint{0.628401in}{1.827839in}}%
\pgfpathlineto{\pgfqpoint{0.642107in}{1.839573in}}%
\pgfpathlineto{\pgfqpoint{0.644230in}{1.841087in}}%
\pgfpathlineto{\pgfqpoint{0.642107in}{1.846798in}}%
\pgfpathlineto{\pgfqpoint{0.640197in}{1.854335in}}%
\pgfpathlineto{\pgfqpoint{0.642107in}{1.856245in}}%
\pgfpathlineto{\pgfqpoint{0.655762in}{1.867583in}}%
\pgfpathlineto{\pgfqpoint{0.656074in}{1.880831in}}%
\pgfpathlineto{\pgfqpoint{0.665842in}{1.894079in}}%
\pgfpathlineto{\pgfqpoint{0.667391in}{1.896812in}}%
\pgfpathlineto{\pgfqpoint{0.673401in}{1.907327in}}%
\pgfpathlineto{\pgfqpoint{0.674523in}{1.920575in}}%
\pgfpathlineto{\pgfqpoint{0.686600in}{1.933823in}}%
\pgfpathlineto{\pgfqpoint{0.684971in}{1.947071in}}%
\pgfpathlineto{\pgfqpoint{0.692674in}{1.954793in}}%
\pgfpathlineto{\pgfqpoint{0.698998in}{1.960319in}}%
\pgfpathlineto{\pgfqpoint{0.699333in}{1.973567in}}%
\pgfpathlineto{\pgfqpoint{0.707530in}{1.986815in}}%
\pgfpathlineto{\pgfqpoint{0.713956in}{2.000063in}}%
\pgfpathlineto{\pgfqpoint{0.717583in}{2.013311in}}%
\pgfpathlineto{\pgfqpoint{0.717957in}{2.013755in}}%
\pgfpathlineto{\pgfqpoint{0.728729in}{2.026559in}}%
\pgfpathlineto{\pgfqpoint{0.728018in}{2.039807in}}%
\pgfpathlineto{\pgfqpoint{0.738250in}{2.053055in}}%
\pgfpathlineto{\pgfqpoint{0.738586in}{2.066303in}}%
\pgfpathlineto{\pgfqpoint{0.743240in}{2.073501in}}%
\pgfpathlineto{\pgfqpoint{0.747472in}{2.079551in}}%
\pgfpathlineto{\pgfqpoint{0.752087in}{2.092799in}}%
\pgfpathlineto{\pgfqpoint{0.755277in}{2.106047in}}%
\pgfpathlineto{\pgfqpoint{0.764351in}{2.119295in}}%
\pgfpathlineto{\pgfqpoint{0.765275in}{2.132543in}}%
\pgfpathlineto{\pgfqpoint{0.743240in}{2.132543in}}%
\pgfpathlineto{\pgfqpoint{0.717957in}{2.132543in}}%
\pgfpathlineto{\pgfqpoint{0.692674in}{2.132543in}}%
\pgfpathlineto{\pgfqpoint{0.667391in}{2.132543in}}%
\pgfpathlineto{\pgfqpoint{0.642107in}{2.132543in}}%
\pgfpathlineto{\pgfqpoint{0.616824in}{2.132543in}}%
\pgfpathlineto{\pgfqpoint{0.591541in}{2.132543in}}%
\pgfpathlineto{\pgfqpoint{0.568649in}{2.132543in}}%
\pgfpathlineto{\pgfqpoint{0.566258in}{2.129593in}}%
\pgfpathlineto{\pgfqpoint{0.562233in}{2.132543in}}%
\pgfpathlineto{\pgfqpoint{0.540975in}{2.132543in}}%
\pgfpathlineto{\pgfqpoint{0.540975in}{2.119295in}}%
\pgfpathlineto{\pgfqpoint{0.540975in}{2.106047in}}%
\pgfpathlineto{\pgfqpoint{0.540975in}{2.092799in}}%
\pgfpathlineto{\pgfqpoint{0.540975in}{2.079551in}}%
\pgfpathlineto{\pgfqpoint{0.540975in}{2.066303in}}%
\pgfpathlineto{\pgfqpoint{0.540975in}{2.053055in}}%
\pgfpathlineto{\pgfqpoint{0.540975in}{2.039807in}}%
\pgfpathlineto{\pgfqpoint{0.540975in}{2.026559in}}%
\pgfpathlineto{\pgfqpoint{0.540975in}{2.013311in}}%
\pgfpathlineto{\pgfqpoint{0.540975in}{2.000063in}}%
\pgfpathlineto{\pgfqpoint{0.540975in}{1.986815in}}%
\pgfpathlineto{\pgfqpoint{0.540975in}{1.973567in}}%
\pgfpathlineto{\pgfqpoint{0.540975in}{1.960319in}}%
\pgfpathlineto{\pgfqpoint{0.540975in}{1.947071in}}%
\pgfpathlineto{\pgfqpoint{0.540975in}{1.933823in}}%
\pgfpathlineto{\pgfqpoint{0.540975in}{1.920575in}}%
\pgfpathlineto{\pgfqpoint{0.540975in}{1.907327in}}%
\pgfpathlineto{\pgfqpoint{0.540975in}{1.894079in}}%
\pgfpathlineto{\pgfqpoint{0.540975in}{1.880831in}}%
\pgfpathlineto{\pgfqpoint{0.540975in}{1.867583in}}%
\pgfpathlineto{\pgfqpoint{0.540975in}{1.854335in}}%
\pgfpathlineto{\pgfqpoint{0.540975in}{1.841087in}}%
\pgfpathlineto{\pgfqpoint{0.540975in}{1.827839in}}%
\pgfpathlineto{\pgfqpoint{0.540975in}{1.814591in}}%
\pgfpathlineto{\pgfqpoint{0.540975in}{1.801343in}}%
\pgfpathlineto{\pgfqpoint{0.540975in}{1.788095in}}%
\pgfpathlineto{\pgfqpoint{0.540975in}{1.774847in}}%
\pgfpathlineto{\pgfqpoint{0.540975in}{1.761599in}}%
\pgfpathlineto{\pgfqpoint{0.540975in}{1.748351in}}%
\pgfpathlineto{\pgfqpoint{0.540975in}{1.735103in}}%
\pgfpathlineto{\pgfqpoint{0.540975in}{1.721855in}}%
\pgfpathlineto{\pgfqpoint{0.540975in}{1.708607in}}%
\pgfpathlineto{\pgfqpoint{0.540975in}{1.695359in}}%
\pgfpathlineto{\pgfqpoint{0.540975in}{1.682111in}}%
\pgfpathlineto{\pgfqpoint{0.540975in}{1.668863in}}%
\pgfpathlineto{\pgfqpoint{0.540975in}{1.655615in}}%
\pgfpathlineto{\pgfqpoint{0.540975in}{1.645215in}}%
\pgfpathlineto{\pgfqpoint{0.566258in}{1.648286in}}%
\pgfpathlineto{\pgfqpoint{0.575922in}{1.642367in}}%
\pgfpathlineto{\pgfqpoint{0.566258in}{1.635875in}}%
\pgfpathlineto{\pgfqpoint{0.540975in}{1.635700in}}%
\pgfpathlineto{\pgfqpoint{0.540975in}{1.629119in}}%
\pgfpathlineto{\pgfqpoint{0.540975in}{1.615871in}}%
\pgfpathlineto{\pgfqpoint{0.540975in}{1.602623in}}%
\pgfpathlineto{\pgfqpoint{0.540975in}{1.589375in}}%
\pgfpathlineto{\pgfqpoint{0.540975in}{1.576127in}}%
\pgfpathlineto{\pgfqpoint{0.540975in}{1.562879in}}%
\pgfpathlineto{\pgfqpoint{0.540975in}{1.549631in}}%
\pgfpathlineto{\pgfqpoint{0.540975in}{1.536383in}}%
\pgfpathlineto{\pgfqpoint{0.540975in}{1.523135in}}%
\pgfpathlineto{\pgfqpoint{0.540975in}{1.509887in}}%
\pgfpathlineto{\pgfqpoint{0.540975in}{1.496639in}}%
\pgfpathlineto{\pgfqpoint{0.540975in}{1.483391in}}%
\pgfpathlineto{\pgfqpoint{0.540975in}{1.470143in}}%
\pgfpathlineto{\pgfqpoint{0.540975in}{1.456895in}}%
\pgfpathlineto{\pgfqpoint{0.540975in}{1.443647in}}%
\pgfpathlineto{\pgfqpoint{0.540975in}{1.430399in}}%
\pgfpathlineto{\pgfqpoint{0.540975in}{1.417151in}}%
\pgfpathlineto{\pgfqpoint{0.540975in}{1.403903in}}%
\pgfpathlineto{\pgfqpoint{0.540975in}{1.390655in}}%
\pgfpathlineto{\pgfqpoint{0.540975in}{1.377407in}}%
\pgfpathlineto{\pgfqpoint{0.540975in}{1.364159in}}%
\pgfpathlineto{\pgfqpoint{0.540975in}{1.350911in}}%
\pgfpathlineto{\pgfqpoint{0.540975in}{1.337663in}}%
\pgfpathlineto{\pgfqpoint{0.540975in}{1.324415in}}%
\pgfpathlineto{\pgfqpoint{0.540975in}{1.311167in}}%
\pgfpathlineto{\pgfqpoint{0.540975in}{1.297919in}}%
\pgfpathlineto{\pgfqpoint{0.540975in}{1.284671in}}%
\pgfpathlineto{\pgfqpoint{0.540975in}{1.271423in}}%
\pgfpathlineto{\pgfqpoint{0.540975in}{1.258175in}}%
\pgfpathlineto{\pgfqpoint{0.540975in}{1.244927in}}%
\pgfpathlineto{\pgfqpoint{0.540975in}{1.231679in}}%
\pgfpathlineto{\pgfqpoint{0.540975in}{1.218431in}}%
\pgfpathlineto{\pgfqpoint{0.540975in}{1.205183in}}%
\pgfpathlineto{\pgfqpoint{0.540975in}{1.191935in}}%
\pgfpathlineto{\pgfqpoint{0.540975in}{1.188094in}}%
\pgfpathlineto{\pgfqpoint{0.555755in}{1.191935in}}%
\pgfpathclose%
\pgfusepath{fill}%
\end{pgfscope}%
\begin{pgfscope}%
\pgfpathrectangle{\pgfqpoint{0.540975in}{0.476543in}}{\pgfqpoint{1.820388in}{1.656000in}}%
\pgfusepath{clip}%
\pgfsetbuttcap%
\pgfsetroundjoin%
\definecolor{currentfill}{rgb}{0.241369,0.366782,0.660592}%
\pgfsetfillcolor{currentfill}%
\pgfsetlinewidth{0.000000pt}%
\definecolor{currentstroke}{rgb}{0.000000,0.000000,0.000000}%
\pgfsetstrokecolor{currentstroke}%
\pgfsetdash{}{0pt}%
\pgfpathmoveto{\pgfqpoint{1.198337in}{0.622147in}}%
\pgfpathlineto{\pgfqpoint{1.223620in}{0.622008in}}%
\pgfpathlineto{\pgfqpoint{1.248904in}{0.621893in}}%
\pgfpathlineto{\pgfqpoint{1.274187in}{0.621801in}}%
\pgfpathlineto{\pgfqpoint{1.299470in}{0.621733in}}%
\pgfpathlineto{\pgfqpoint{1.324753in}{0.621687in}}%
\pgfpathlineto{\pgfqpoint{1.350036in}{0.621664in}}%
\pgfpathlineto{\pgfqpoint{1.375319in}{0.621665in}}%
\pgfpathlineto{\pgfqpoint{1.400603in}{0.621688in}}%
\pgfpathlineto{\pgfqpoint{1.425886in}{0.621735in}}%
\pgfpathlineto{\pgfqpoint{1.451169in}{0.621806in}}%
\pgfpathlineto{\pgfqpoint{1.476452in}{0.621900in}}%
\pgfpathlineto{\pgfqpoint{1.501735in}{0.622020in}}%
\pgfpathlineto{\pgfqpoint{1.527018in}{0.622165in}}%
\pgfpathlineto{\pgfqpoint{1.542886in}{0.622271in}}%
\pgfpathlineto{\pgfqpoint{1.552302in}{0.622548in}}%
\pgfpathlineto{\pgfqpoint{1.577585in}{0.623393in}}%
\pgfpathlineto{\pgfqpoint{1.602868in}{0.624342in}}%
\pgfpathlineto{\pgfqpoint{1.628151in}{0.625398in}}%
\pgfpathlineto{\pgfqpoint{1.653434in}{0.626565in}}%
\pgfpathlineto{\pgfqpoint{1.678717in}{0.627847in}}%
\pgfpathlineto{\pgfqpoint{1.704001in}{0.629249in}}%
\pgfpathlineto{\pgfqpoint{1.729284in}{0.630780in}}%
\pgfpathlineto{\pgfqpoint{1.754567in}{0.632445in}}%
\pgfpathlineto{\pgfqpoint{1.779850in}{0.634255in}}%
\pgfpathlineto{\pgfqpoint{1.796204in}{0.635519in}}%
\pgfpathlineto{\pgfqpoint{1.805133in}{0.635675in}}%
\pgfpathlineto{\pgfqpoint{1.830416in}{0.636151in}}%
\pgfpathlineto{\pgfqpoint{1.855700in}{0.636667in}}%
\pgfpathlineto{\pgfqpoint{1.880983in}{0.637227in}}%
\pgfpathlineto{\pgfqpoint{1.906266in}{0.637833in}}%
\pgfpathlineto{\pgfqpoint{1.931549in}{0.638490in}}%
\pgfpathlineto{\pgfqpoint{1.956832in}{0.639201in}}%
\pgfpathlineto{\pgfqpoint{1.982115in}{0.639972in}}%
\pgfpathlineto{\pgfqpoint{2.007399in}{0.640807in}}%
\pgfpathlineto{\pgfqpoint{2.032682in}{0.641714in}}%
\pgfpathlineto{\pgfqpoint{2.057965in}{0.642697in}}%
\pgfpathlineto{\pgfqpoint{2.083248in}{0.643766in}}%
\pgfpathlineto{\pgfqpoint{2.108531in}{0.644928in}}%
\pgfpathlineto{\pgfqpoint{2.133815in}{0.646193in}}%
\pgfpathlineto{\pgfqpoint{2.159098in}{0.647573in}}%
\pgfpathlineto{\pgfqpoint{2.179282in}{0.648767in}}%
\pgfpathlineto{\pgfqpoint{2.161055in}{0.662015in}}%
\pgfpathlineto{\pgfqpoint{2.184381in}{0.663650in}}%
\pgfpathlineto{\pgfqpoint{2.209664in}{0.665569in}}%
\pgfpathlineto{\pgfqpoint{2.234947in}{0.667659in}}%
\pgfpathlineto{\pgfqpoint{2.260230in}{0.669945in}}%
\pgfpathlineto{\pgfqpoint{2.285514in}{0.672452in}}%
\pgfpathlineto{\pgfqpoint{2.310797in}{0.675213in}}%
\pgfpathlineto{\pgfqpoint{2.311236in}{0.675263in}}%
\pgfpathlineto{\pgfqpoint{2.310797in}{0.675492in}}%
\pgfpathlineto{\pgfqpoint{2.289972in}{0.688511in}}%
\pgfpathlineto{\pgfqpoint{2.310797in}{0.694865in}}%
\pgfpathlineto{\pgfqpoint{2.330586in}{0.701759in}}%
\pgfpathlineto{\pgfqpoint{2.323477in}{0.715007in}}%
\pgfpathlineto{\pgfqpoint{2.315785in}{0.728255in}}%
\pgfpathlineto{\pgfqpoint{2.332261in}{0.741503in}}%
\pgfpathlineto{\pgfqpoint{2.310797in}{0.750833in}}%
\pgfpathlineto{\pgfqpoint{2.302260in}{0.754751in}}%
\pgfpathlineto{\pgfqpoint{2.310797in}{0.762061in}}%
\pgfpathlineto{\pgfqpoint{2.317005in}{0.767999in}}%
\pgfpathlineto{\pgfqpoint{2.310797in}{0.769761in}}%
\pgfpathlineto{\pgfqpoint{2.285514in}{0.777306in}}%
\pgfpathlineto{\pgfqpoint{2.273094in}{0.781247in}}%
\pgfpathlineto{\pgfqpoint{2.281567in}{0.794495in}}%
\pgfpathlineto{\pgfqpoint{2.260230in}{0.803765in}}%
\pgfpathlineto{\pgfqpoint{2.251861in}{0.807743in}}%
\pgfpathlineto{\pgfqpoint{2.242072in}{0.820991in}}%
\pgfpathlineto{\pgfqpoint{2.234947in}{0.826348in}}%
\pgfpathlineto{\pgfqpoint{2.225822in}{0.834239in}}%
\pgfpathlineto{\pgfqpoint{2.209664in}{0.842533in}}%
\pgfpathlineto{\pgfqpoint{2.200322in}{0.847487in}}%
\pgfpathlineto{\pgfqpoint{2.198753in}{0.860735in}}%
\pgfpathlineto{\pgfqpoint{2.184381in}{0.867391in}}%
\pgfpathlineto{\pgfqpoint{2.170717in}{0.873983in}}%
\pgfpathlineto{\pgfqpoint{2.169803in}{0.887231in}}%
\pgfpathlineto{\pgfqpoint{2.159098in}{0.892042in}}%
\pgfpathlineto{\pgfqpoint{2.141700in}{0.900479in}}%
\pgfpathlineto{\pgfqpoint{2.135980in}{0.913727in}}%
\pgfpathlineto{\pgfqpoint{2.133815in}{0.915371in}}%
\pgfpathlineto{\pgfqpoint{2.120196in}{0.926975in}}%
\pgfpathlineto{\pgfqpoint{2.108531in}{0.937669in}}%
\pgfpathlineto{\pgfqpoint{2.105938in}{0.940223in}}%
\pgfpathlineto{\pgfqpoint{2.099237in}{0.953471in}}%
\pgfpathlineto{\pgfqpoint{2.083248in}{0.962503in}}%
\pgfpathlineto{\pgfqpoint{2.076170in}{0.966719in}}%
\pgfpathlineto{\pgfqpoint{2.074534in}{0.979967in}}%
\pgfpathlineto{\pgfqpoint{2.057965in}{0.990297in}}%
\pgfpathlineto{\pgfqpoint{2.053620in}{0.993215in}}%
\pgfpathlineto{\pgfqpoint{2.048556in}{1.006463in}}%
\pgfpathlineto{\pgfqpoint{2.032682in}{1.018783in}}%
\pgfpathlineto{\pgfqpoint{2.031632in}{1.019711in}}%
\pgfpathlineto{\pgfqpoint{2.021038in}{1.032959in}}%
\pgfpathlineto{\pgfqpoint{2.013185in}{1.046207in}}%
\pgfpathlineto{\pgfqpoint{2.007399in}{1.051248in}}%
\pgfpathlineto{\pgfqpoint{1.998878in}{1.059455in}}%
\pgfpathlineto{\pgfqpoint{1.996392in}{1.072703in}}%
\pgfpathlineto{\pgfqpoint{1.982115in}{1.083062in}}%
\pgfpathlineto{\pgfqpoint{1.978555in}{1.085951in}}%
\pgfpathlineto{\pgfqpoint{1.976091in}{1.099199in}}%
\pgfpathlineto{\pgfqpoint{1.963177in}{1.112447in}}%
\pgfpathlineto{\pgfqpoint{1.956832in}{1.122806in}}%
\pgfpathlineto{\pgfqpoint{1.955333in}{1.125695in}}%
\pgfpathlineto{\pgfqpoint{1.947622in}{1.138943in}}%
\pgfpathlineto{\pgfqpoint{1.934236in}{1.152191in}}%
\pgfpathlineto{\pgfqpoint{1.931549in}{1.161486in}}%
\pgfpathlineto{\pgfqpoint{1.930656in}{1.165439in}}%
\pgfpathlineto{\pgfqpoint{1.918358in}{1.178687in}}%
\pgfpathlineto{\pgfqpoint{1.915896in}{1.191935in}}%
\pgfpathlineto{\pgfqpoint{1.906266in}{1.201892in}}%
\pgfpathlineto{\pgfqpoint{1.903644in}{1.205183in}}%
\pgfpathlineto{\pgfqpoint{1.898780in}{1.218431in}}%
\pgfpathlineto{\pgfqpoint{1.892105in}{1.231679in}}%
\pgfpathlineto{\pgfqpoint{1.882653in}{1.244927in}}%
\pgfpathlineto{\pgfqpoint{1.880983in}{1.250541in}}%
\pgfpathlineto{\pgfqpoint{1.879288in}{1.258175in}}%
\pgfpathlineto{\pgfqpoint{1.868199in}{1.271423in}}%
\pgfpathlineto{\pgfqpoint{1.866158in}{1.284671in}}%
\pgfpathlineto{\pgfqpoint{1.855700in}{1.297490in}}%
\pgfpathlineto{\pgfqpoint{1.855420in}{1.297919in}}%
\pgfpathlineto{\pgfqpoint{1.851975in}{1.311167in}}%
\pgfpathlineto{\pgfqpoint{1.845759in}{1.324415in}}%
\pgfpathlineto{\pgfqpoint{1.837853in}{1.337663in}}%
\pgfpathlineto{\pgfqpoint{1.834246in}{1.350911in}}%
\pgfpathlineto{\pgfqpoint{1.830416in}{1.355912in}}%
\pgfpathlineto{\pgfqpoint{1.825594in}{1.364159in}}%
\pgfpathlineto{\pgfqpoint{1.823838in}{1.377407in}}%
\pgfpathlineto{\pgfqpoint{1.815203in}{1.390655in}}%
\pgfpathlineto{\pgfqpoint{1.812031in}{1.403903in}}%
\pgfpathlineto{\pgfqpoint{1.805133in}{1.417073in}}%
\pgfpathlineto{\pgfqpoint{1.805104in}{1.417151in}}%
\pgfpathlineto{\pgfqpoint{1.800513in}{1.430399in}}%
\pgfpathlineto{\pgfqpoint{1.797718in}{1.443647in}}%
\pgfpathlineto{\pgfqpoint{1.789472in}{1.456895in}}%
\pgfpathlineto{\pgfqpoint{1.787435in}{1.470143in}}%
\pgfpathlineto{\pgfqpoint{1.779850in}{1.481462in}}%
\pgfpathlineto{\pgfqpoint{1.778929in}{1.483391in}}%
\pgfpathlineto{\pgfqpoint{1.776673in}{1.496639in}}%
\pgfpathlineto{\pgfqpoint{1.771771in}{1.509887in}}%
\pgfpathlineto{\pgfqpoint{1.767158in}{1.523135in}}%
\pgfpathlineto{\pgfqpoint{1.764269in}{1.536383in}}%
\pgfpathlineto{\pgfqpoint{1.757301in}{1.549631in}}%
\pgfpathlineto{\pgfqpoint{1.755338in}{1.562879in}}%
\pgfpathlineto{\pgfqpoint{1.754567in}{1.563999in}}%
\pgfpathlineto{\pgfqpoint{1.749210in}{1.576127in}}%
\pgfpathlineto{\pgfqpoint{1.747243in}{1.589375in}}%
\pgfpathlineto{\pgfqpoint{1.742374in}{1.602623in}}%
\pgfpathlineto{\pgfqpoint{1.738165in}{1.615871in}}%
\pgfpathlineto{\pgfqpoint{1.734811in}{1.629119in}}%
\pgfpathlineto{\pgfqpoint{1.729284in}{1.638610in}}%
\pgfpathlineto{\pgfqpoint{1.727910in}{1.642367in}}%
\pgfpathlineto{\pgfqpoint{1.726793in}{1.655615in}}%
\pgfpathlineto{\pgfqpoint{1.722100in}{1.668863in}}%
\pgfpathlineto{\pgfqpoint{1.720578in}{1.682111in}}%
\pgfpathlineto{\pgfqpoint{1.715951in}{1.695359in}}%
\pgfpathlineto{\pgfqpoint{1.712156in}{1.708607in}}%
\pgfpathlineto{\pgfqpoint{1.708755in}{1.721855in}}%
\pgfpathlineto{\pgfqpoint{1.704001in}{1.731601in}}%
\pgfpathlineto{\pgfqpoint{1.703010in}{1.735103in}}%
\pgfpathlineto{\pgfqpoint{1.701882in}{1.748351in}}%
\pgfpathlineto{\pgfqpoint{1.697166in}{1.761599in}}%
\pgfpathlineto{\pgfqpoint{1.695825in}{1.774847in}}%
\pgfpathlineto{\pgfqpoint{1.691808in}{1.788095in}}%
\pgfpathlineto{\pgfqpoint{1.688840in}{1.801343in}}%
\pgfpathlineto{\pgfqpoint{1.685781in}{1.814591in}}%
\pgfpathlineto{\pgfqpoint{1.679385in}{1.827839in}}%
\pgfpathlineto{\pgfqpoint{1.678717in}{1.831910in}}%
\pgfpathlineto{\pgfqpoint{1.677996in}{1.841087in}}%
\pgfpathlineto{\pgfqpoint{1.674454in}{1.854335in}}%
\pgfpathlineto{\pgfqpoint{1.673404in}{1.867583in}}%
\pgfpathlineto{\pgfqpoint{1.669788in}{1.880831in}}%
\pgfpathlineto{\pgfqpoint{1.667334in}{1.894079in}}%
\pgfpathlineto{\pgfqpoint{1.664765in}{1.907327in}}%
\pgfpathlineto{\pgfqpoint{1.660276in}{1.920575in}}%
\pgfpathlineto{\pgfqpoint{1.658476in}{1.933823in}}%
\pgfpathlineto{\pgfqpoint{1.653434in}{1.942717in}}%
\pgfpathlineto{\pgfqpoint{1.652205in}{1.947071in}}%
\pgfpathlineto{\pgfqpoint{1.651260in}{1.960319in}}%
\pgfpathlineto{\pgfqpoint{1.648377in}{1.973567in}}%
\pgfpathlineto{\pgfqpoint{1.646616in}{1.986815in}}%
\pgfpathlineto{\pgfqpoint{1.644676in}{2.000063in}}%
\pgfpathlineto{\pgfqpoint{1.641054in}{2.013311in}}%
\pgfpathlineto{\pgfqpoint{1.639835in}{2.026559in}}%
\pgfpathlineto{\pgfqpoint{1.634973in}{2.039807in}}%
\pgfpathlineto{\pgfqpoint{1.633427in}{2.053055in}}%
\pgfpathlineto{\pgfqpoint{1.628151in}{2.063571in}}%
\pgfpathlineto{\pgfqpoint{1.627524in}{2.066303in}}%
\pgfpathlineto{\pgfqpoint{1.625886in}{2.079551in}}%
\pgfpathlineto{\pgfqpoint{1.624048in}{2.092799in}}%
\pgfpathlineto{\pgfqpoint{1.621477in}{2.106047in}}%
\pgfpathlineto{\pgfqpoint{1.620673in}{2.119295in}}%
\pgfpathlineto{\pgfqpoint{1.616784in}{2.132543in}}%
\pgfpathlineto{\pgfqpoint{1.602868in}{2.132543in}}%
\pgfpathlineto{\pgfqpoint{1.597931in}{2.132543in}}%
\pgfpathlineto{\pgfqpoint{1.601411in}{2.119295in}}%
\pgfpathlineto{\pgfqpoint{1.602347in}{2.106047in}}%
\pgfpathlineto{\pgfqpoint{1.602868in}{2.103526in}}%
\pgfpathlineto{\pgfqpoint{1.606396in}{2.092799in}}%
\pgfpathlineto{\pgfqpoint{1.609350in}{2.079551in}}%
\pgfpathlineto{\pgfqpoint{1.611982in}{2.066303in}}%
\pgfpathlineto{\pgfqpoint{1.616037in}{2.053055in}}%
\pgfpathlineto{\pgfqpoint{1.617236in}{2.039807in}}%
\pgfpathlineto{\pgfqpoint{1.621089in}{2.026559in}}%
\pgfpathlineto{\pgfqpoint{1.622144in}{2.013311in}}%
\pgfpathlineto{\pgfqpoint{1.625296in}{2.000063in}}%
\pgfpathlineto{\pgfqpoint{1.627152in}{1.986815in}}%
\pgfpathlineto{\pgfqpoint{1.628151in}{1.979310in}}%
\pgfpathlineto{\pgfqpoint{1.629402in}{1.973567in}}%
\pgfpathlineto{\pgfqpoint{1.634008in}{1.960319in}}%
\pgfpathlineto{\pgfqpoint{1.635471in}{1.947071in}}%
\pgfpathlineto{\pgfqpoint{1.640515in}{1.933823in}}%
\pgfpathlineto{\pgfqpoint{1.641822in}{1.920575in}}%
\pgfpathlineto{\pgfqpoint{1.645477in}{1.907327in}}%
\pgfpathlineto{\pgfqpoint{1.647650in}{1.894079in}}%
\pgfpathlineto{\pgfqpoint{1.649876in}{1.880831in}}%
\pgfpathlineto{\pgfqpoint{1.653162in}{1.867583in}}%
\pgfpathlineto{\pgfqpoint{1.653434in}{1.864192in}}%
\pgfpathlineto{\pgfqpoint{1.654703in}{1.854335in}}%
\pgfpathlineto{\pgfqpoint{1.660204in}{1.841087in}}%
\pgfpathlineto{\pgfqpoint{1.661559in}{1.827839in}}%
\pgfpathlineto{\pgfqpoint{1.666270in}{1.814591in}}%
\pgfpathlineto{\pgfqpoint{1.668604in}{1.801343in}}%
\pgfpathlineto{\pgfqpoint{1.671214in}{1.788095in}}%
\pgfpathlineto{\pgfqpoint{1.674815in}{1.774847in}}%
\pgfpathlineto{\pgfqpoint{1.676061in}{1.761599in}}%
\pgfpathlineto{\pgfqpoint{1.678717in}{1.753844in}}%
\pgfpathlineto{\pgfqpoint{1.681469in}{1.748351in}}%
\pgfpathlineto{\pgfqpoint{1.683042in}{1.735103in}}%
\pgfpathlineto{\pgfqpoint{1.688259in}{1.721855in}}%
\pgfpathlineto{\pgfqpoint{1.690816in}{1.708607in}}%
\pgfpathlineto{\pgfqpoint{1.694182in}{1.695359in}}%
\pgfpathlineto{\pgfqpoint{1.698159in}{1.682111in}}%
\pgfpathlineto{\pgfqpoint{1.699569in}{1.668863in}}%
\pgfpathlineto{\pgfqpoint{1.704001in}{1.656044in}}%
\pgfpathlineto{\pgfqpoint{1.704212in}{1.655615in}}%
\pgfpathlineto{\pgfqpoint{1.705579in}{1.642367in}}%
\pgfpathlineto{\pgfqpoint{1.712055in}{1.629119in}}%
\pgfpathlineto{\pgfqpoint{1.714677in}{1.615871in}}%
\pgfpathlineto{\pgfqpoint{1.718587in}{1.602623in}}%
\pgfpathlineto{\pgfqpoint{1.722977in}{1.589375in}}%
\pgfpathlineto{\pgfqpoint{1.724894in}{1.576127in}}%
\pgfpathlineto{\pgfqpoint{1.729284in}{1.565786in}}%
\pgfpathlineto{\pgfqpoint{1.730961in}{1.562879in}}%
\pgfpathlineto{\pgfqpoint{1.732366in}{1.549631in}}%
\pgfpathlineto{\pgfqpoint{1.738754in}{1.536383in}}%
\pgfpathlineto{\pgfqpoint{1.741005in}{1.523135in}}%
\pgfpathlineto{\pgfqpoint{1.745515in}{1.509887in}}%
\pgfpathlineto{\pgfqpoint{1.749934in}{1.496639in}}%
\pgfpathlineto{\pgfqpoint{1.752169in}{1.483391in}}%
\pgfpathlineto{\pgfqpoint{1.754567in}{1.478161in}}%
\pgfpathlineto{\pgfqpoint{1.759317in}{1.470143in}}%
\pgfpathlineto{\pgfqpoint{1.760867in}{1.456895in}}%
\pgfpathlineto{\pgfqpoint{1.768857in}{1.443647in}}%
\pgfpathlineto{\pgfqpoint{1.771139in}{1.430399in}}%
\pgfpathlineto{\pgfqpoint{1.775998in}{1.417151in}}%
\pgfpathlineto{\pgfqpoint{1.779850in}{1.405585in}}%
\pgfpathlineto{\pgfqpoint{1.780583in}{1.403903in}}%
\pgfpathlineto{\pgfqpoint{1.783691in}{1.390655in}}%
\pgfpathlineto{\pgfqpoint{1.791853in}{1.377407in}}%
\pgfpathlineto{\pgfqpoint{1.793270in}{1.364159in}}%
\pgfpathlineto{\pgfqpoint{1.801090in}{1.350911in}}%
\pgfpathlineto{\pgfqpoint{1.803032in}{1.337663in}}%
\pgfpathlineto{\pgfqpoint{1.805133in}{1.333447in}}%
\pgfpathlineto{\pgfqpoint{1.810564in}{1.324415in}}%
\pgfpathlineto{\pgfqpoint{1.816044in}{1.311167in}}%
\pgfpathlineto{\pgfqpoint{1.819628in}{1.297919in}}%
\pgfpathlineto{\pgfqpoint{1.827753in}{1.284671in}}%
\pgfpathlineto{\pgfqpoint{1.828978in}{1.271423in}}%
\pgfpathlineto{\pgfqpoint{1.830416in}{1.269462in}}%
\pgfpathlineto{\pgfqpoint{1.840055in}{1.258175in}}%
\pgfpathlineto{\pgfqpoint{1.842181in}{1.244927in}}%
\pgfpathlineto{\pgfqpoint{1.850186in}{1.231679in}}%
\pgfpathlineto{\pgfqpoint{1.854790in}{1.218431in}}%
\pgfpathlineto{\pgfqpoint{1.855700in}{1.215689in}}%
\pgfpathlineto{\pgfqpoint{1.859800in}{1.205183in}}%
\pgfpathlineto{\pgfqpoint{1.869486in}{1.191935in}}%
\pgfpathlineto{\pgfqpoint{1.870896in}{1.178687in}}%
\pgfpathlineto{\pgfqpoint{1.880983in}{1.166560in}}%
\pgfpathlineto{\pgfqpoint{1.882040in}{1.165439in}}%
\pgfpathlineto{\pgfqpoint{1.883539in}{1.152191in}}%
\pgfpathlineto{\pgfqpoint{1.895652in}{1.138943in}}%
\pgfpathlineto{\pgfqpoint{1.900869in}{1.125695in}}%
\pgfpathlineto{\pgfqpoint{1.906266in}{1.114459in}}%
\pgfpathlineto{\pgfqpoint{1.907358in}{1.112447in}}%
\pgfpathlineto{\pgfqpoint{1.917790in}{1.099199in}}%
\pgfpathlineto{\pgfqpoint{1.919184in}{1.085951in}}%
\pgfpathlineto{\pgfqpoint{1.931549in}{1.074191in}}%
\pgfpathlineto{\pgfqpoint{1.933289in}{1.072703in}}%
\pgfpathlineto{\pgfqpoint{1.933163in}{1.059455in}}%
\pgfpathlineto{\pgfqpoint{1.946307in}{1.046207in}}%
\pgfpathlineto{\pgfqpoint{1.949486in}{1.032959in}}%
\pgfpathlineto{\pgfqpoint{1.956832in}{1.022681in}}%
\pgfpathlineto{\pgfqpoint{1.959135in}{1.019711in}}%
\pgfpathlineto{\pgfqpoint{1.970154in}{1.006463in}}%
\pgfpathlineto{\pgfqpoint{1.973111in}{0.993215in}}%
\pgfpathlineto{\pgfqpoint{1.982115in}{0.985642in}}%
\pgfpathlineto{\pgfqpoint{1.989484in}{0.979967in}}%
\pgfpathlineto{\pgfqpoint{1.986747in}{0.966719in}}%
\pgfpathlineto{\pgfqpoint{2.005941in}{0.953471in}}%
\pgfpathlineto{\pgfqpoint{2.005689in}{0.940223in}}%
\pgfpathlineto{\pgfqpoint{2.007399in}{0.938244in}}%
\pgfpathlineto{\pgfqpoint{2.017691in}{0.926975in}}%
\pgfpathlineto{\pgfqpoint{2.024544in}{0.913727in}}%
\pgfpathlineto{\pgfqpoint{2.027024in}{0.900479in}}%
\pgfpathlineto{\pgfqpoint{2.032682in}{0.895867in}}%
\pgfpathlineto{\pgfqpoint{2.044489in}{0.887231in}}%
\pgfpathlineto{\pgfqpoint{2.037170in}{0.873983in}}%
\pgfpathlineto{\pgfqpoint{2.057631in}{0.860735in}}%
\pgfpathlineto{\pgfqpoint{2.043893in}{0.847487in}}%
\pgfpathlineto{\pgfqpoint{2.057965in}{0.837946in}}%
\pgfpathlineto{\pgfqpoint{2.063686in}{0.834239in}}%
\pgfpathlineto{\pgfqpoint{2.058037in}{0.820991in}}%
\pgfpathlineto{\pgfqpoint{2.059983in}{0.807743in}}%
\pgfpathlineto{\pgfqpoint{2.066394in}{0.794495in}}%
\pgfpathlineto{\pgfqpoint{2.057965in}{0.789896in}}%
\pgfpathlineto{\pgfqpoint{2.040728in}{0.781247in}}%
\pgfpathlineto{\pgfqpoint{2.057965in}{0.768921in}}%
\pgfpathlineto{\pgfqpoint{2.059421in}{0.767999in}}%
\pgfpathlineto{\pgfqpoint{2.057965in}{0.767588in}}%
\pgfpathlineto{\pgfqpoint{2.032682in}{0.761037in}}%
\pgfpathlineto{\pgfqpoint{2.007399in}{0.755447in}}%
\pgfpathlineto{\pgfqpoint{2.003954in}{0.754751in}}%
\pgfpathlineto{\pgfqpoint{2.007399in}{0.748853in}}%
\pgfpathlineto{\pgfqpoint{2.012066in}{0.741503in}}%
\pgfpathlineto{\pgfqpoint{2.007399in}{0.740723in}}%
\pgfpathlineto{\pgfqpoint{1.982115in}{0.736883in}}%
\pgfpathlineto{\pgfqpoint{1.956832in}{0.733580in}}%
\pgfpathlineto{\pgfqpoint{1.931549in}{0.730722in}}%
\pgfpathlineto{\pgfqpoint{1.906484in}{0.728255in}}%
\pgfpathlineto{\pgfqpoint{1.906266in}{0.728107in}}%
\pgfpathlineto{\pgfqpoint{1.884945in}{0.715007in}}%
\pgfpathlineto{\pgfqpoint{1.880983in}{0.714527in}}%
\pgfpathlineto{\pgfqpoint{1.855700in}{0.711767in}}%
\pgfpathlineto{\pgfqpoint{1.830416in}{0.709405in}}%
\pgfpathlineto{\pgfqpoint{1.805133in}{0.707378in}}%
\pgfpathlineto{\pgfqpoint{1.779850in}{0.705634in}}%
\pgfpathlineto{\pgfqpoint{1.754567in}{0.704133in}}%
\pgfpathlineto{\pgfqpoint{1.729284in}{0.702841in}}%
\pgfpathlineto{\pgfqpoint{1.704664in}{0.701759in}}%
\pgfpathlineto{\pgfqpoint{1.704001in}{0.701706in}}%
\pgfpathlineto{\pgfqpoint{1.678717in}{0.699826in}}%
\pgfpathlineto{\pgfqpoint{1.653434in}{0.698094in}}%
\pgfpathlineto{\pgfqpoint{1.628151in}{0.696501in}}%
\pgfpathlineto{\pgfqpoint{1.602868in}{0.695043in}}%
\pgfpathlineto{\pgfqpoint{1.577585in}{0.693713in}}%
\pgfpathlineto{\pgfqpoint{1.552302in}{0.692510in}}%
\pgfpathlineto{\pgfqpoint{1.527018in}{0.691434in}}%
\pgfpathlineto{\pgfqpoint{1.501735in}{0.690484in}}%
\pgfpathlineto{\pgfqpoint{1.476452in}{0.689667in}}%
\pgfpathlineto{\pgfqpoint{1.451169in}{0.688987in}}%
\pgfpathlineto{\pgfqpoint{1.428463in}{0.688511in}}%
\pgfpathlineto{\pgfqpoint{1.425886in}{0.688485in}}%
\pgfpathlineto{\pgfqpoint{1.400603in}{0.688320in}}%
\pgfpathlineto{\pgfqpoint{1.375319in}{0.688254in}}%
\pgfpathlineto{\pgfqpoint{1.350036in}{0.688279in}}%
\pgfpathlineto{\pgfqpoint{1.324753in}{0.688392in}}%
\pgfpathlineto{\pgfqpoint{1.309631in}{0.688511in}}%
\pgfpathlineto{\pgfqpoint{1.299470in}{0.688744in}}%
\pgfpathlineto{\pgfqpoint{1.274187in}{0.689662in}}%
\pgfpathlineto{\pgfqpoint{1.248904in}{0.691032in}}%
\pgfpathlineto{\pgfqpoint{1.223620in}{0.692986in}}%
\pgfpathlineto{\pgfqpoint{1.198337in}{0.695717in}}%
\pgfpathlineto{\pgfqpoint{1.173054in}{0.699518in}}%
\pgfpathlineto{\pgfqpoint{1.161611in}{0.701759in}}%
\pgfpathlineto{\pgfqpoint{1.147771in}{0.702242in}}%
\pgfpathlineto{\pgfqpoint{1.122488in}{0.703225in}}%
\pgfpathlineto{\pgfqpoint{1.097204in}{0.704330in}}%
\pgfpathlineto{\pgfqpoint{1.071921in}{0.705569in}}%
\pgfpathlineto{\pgfqpoint{1.046638in}{0.706959in}}%
\pgfpathlineto{\pgfqpoint{1.021355in}{0.708518in}}%
\pgfpathlineto{\pgfqpoint{0.996072in}{0.710270in}}%
\pgfpathlineto{\pgfqpoint{0.970789in}{0.712242in}}%
\pgfpathlineto{\pgfqpoint{0.945505in}{0.714467in}}%
\pgfpathlineto{\pgfqpoint{0.939877in}{0.715007in}}%
\pgfpathlineto{\pgfqpoint{0.929364in}{0.728255in}}%
\pgfpathlineto{\pgfqpoint{0.920222in}{0.729951in}}%
\pgfpathlineto{\pgfqpoint{0.894939in}{0.735756in}}%
\pgfpathlineto{\pgfqpoint{0.875974in}{0.741503in}}%
\pgfpathlineto{\pgfqpoint{0.869656in}{0.743187in}}%
\pgfpathlineto{\pgfqpoint{0.844373in}{0.750359in}}%
\pgfpathlineto{\pgfqpoint{0.830023in}{0.754751in}}%
\pgfpathlineto{\pgfqpoint{0.844373in}{0.762473in}}%
\pgfpathlineto{\pgfqpoint{0.853032in}{0.767999in}}%
\pgfpathlineto{\pgfqpoint{0.844373in}{0.770277in}}%
\pgfpathlineto{\pgfqpoint{0.819090in}{0.777685in}}%
\pgfpathlineto{\pgfqpoint{0.808266in}{0.781247in}}%
\pgfpathlineto{\pgfqpoint{0.819090in}{0.786655in}}%
\pgfpathlineto{\pgfqpoint{0.833880in}{0.794495in}}%
\pgfpathlineto{\pgfqpoint{0.819090in}{0.801995in}}%
\pgfpathlineto{\pgfqpoint{0.809535in}{0.807743in}}%
\pgfpathlineto{\pgfqpoint{0.816637in}{0.820991in}}%
\pgfpathlineto{\pgfqpoint{0.819090in}{0.823364in}}%
\pgfpathlineto{\pgfqpoint{0.828637in}{0.834239in}}%
\pgfpathlineto{\pgfqpoint{0.820522in}{0.847487in}}%
\pgfpathlineto{\pgfqpoint{0.844373in}{0.859538in}}%
\pgfpathlineto{\pgfqpoint{0.846541in}{0.860735in}}%
\pgfpathlineto{\pgfqpoint{0.844373in}{0.862836in}}%
\pgfpathlineto{\pgfqpoint{0.833691in}{0.873983in}}%
\pgfpathlineto{\pgfqpoint{0.844373in}{0.879638in}}%
\pgfpathlineto{\pgfqpoint{0.858212in}{0.887231in}}%
\pgfpathlineto{\pgfqpoint{0.858081in}{0.900479in}}%
\pgfpathlineto{\pgfqpoint{0.868417in}{0.913727in}}%
\pgfpathlineto{\pgfqpoint{0.869656in}{0.914903in}}%
\pgfpathlineto{\pgfqpoint{0.880852in}{0.926975in}}%
\pgfpathlineto{\pgfqpoint{0.878121in}{0.940223in}}%
\pgfpathlineto{\pgfqpoint{0.894939in}{0.949831in}}%
\pgfpathlineto{\pgfqpoint{0.900847in}{0.953471in}}%
\pgfpathlineto{\pgfqpoint{0.898058in}{0.966719in}}%
\pgfpathlineto{\pgfqpoint{0.916527in}{0.979967in}}%
\pgfpathlineto{\pgfqpoint{0.919088in}{0.993215in}}%
\pgfpathlineto{\pgfqpoint{0.920222in}{0.994839in}}%
\pgfpathlineto{\pgfqpoint{0.927815in}{1.006463in}}%
\pgfpathlineto{\pgfqpoint{0.940494in}{1.019711in}}%
\pgfpathlineto{\pgfqpoint{0.942258in}{1.032959in}}%
\pgfpathlineto{\pgfqpoint{0.945505in}{1.035434in}}%
\pgfpathlineto{\pgfqpoint{0.958412in}{1.046207in}}%
\pgfpathlineto{\pgfqpoint{0.957307in}{1.059455in}}%
\pgfpathlineto{\pgfqpoint{0.970789in}{1.070667in}}%
\pgfpathlineto{\pgfqpoint{0.973084in}{1.072703in}}%
\pgfpathlineto{\pgfqpoint{0.977021in}{1.085951in}}%
\pgfpathlineto{\pgfqpoint{0.986173in}{1.099199in}}%
\pgfpathlineto{\pgfqpoint{0.995533in}{1.112447in}}%
\pgfpathlineto{\pgfqpoint{0.996072in}{1.115847in}}%
\pgfpathlineto{\pgfqpoint{0.997420in}{1.125695in}}%
\pgfpathlineto{\pgfqpoint{1.010679in}{1.138943in}}%
\pgfpathlineto{\pgfqpoint{1.010937in}{1.152191in}}%
\pgfpathlineto{\pgfqpoint{1.021355in}{1.162831in}}%
\pgfpathlineto{\pgfqpoint{1.023658in}{1.165439in}}%
\pgfpathlineto{\pgfqpoint{1.026168in}{1.178687in}}%
\pgfpathlineto{\pgfqpoint{1.033853in}{1.191935in}}%
\pgfpathlineto{\pgfqpoint{1.042409in}{1.205183in}}%
\pgfpathlineto{\pgfqpoint{1.046638in}{1.218404in}}%
\pgfpathlineto{\pgfqpoint{1.046646in}{1.218431in}}%
\pgfpathlineto{\pgfqpoint{1.056630in}{1.231679in}}%
\pgfpathlineto{\pgfqpoint{1.057255in}{1.244927in}}%
\pgfpathlineto{\pgfqpoint{1.068020in}{1.258175in}}%
\pgfpathlineto{\pgfqpoint{1.070350in}{1.271423in}}%
\pgfpathlineto{\pgfqpoint{1.071921in}{1.274001in}}%
\pgfpathlineto{\pgfqpoint{1.077451in}{1.284671in}}%
\pgfpathlineto{\pgfqpoint{1.082725in}{1.297919in}}%
\pgfpathlineto{\pgfqpoint{1.085827in}{1.311167in}}%
\pgfpathlineto{\pgfqpoint{1.095671in}{1.324415in}}%
\pgfpathlineto{\pgfqpoint{1.097204in}{1.334612in}}%
\pgfpathlineto{\pgfqpoint{1.097562in}{1.337663in}}%
\pgfpathlineto{\pgfqpoint{1.106174in}{1.350911in}}%
\pgfpathlineto{\pgfqpoint{1.107762in}{1.364159in}}%
\pgfpathlineto{\pgfqpoint{1.114568in}{1.377407in}}%
\pgfpathlineto{\pgfqpoint{1.119507in}{1.390655in}}%
\pgfpathlineto{\pgfqpoint{1.122488in}{1.400028in}}%
\pgfpathlineto{\pgfqpoint{1.123460in}{1.403903in}}%
\pgfpathlineto{\pgfqpoint{1.129838in}{1.417151in}}%
\pgfpathlineto{\pgfqpoint{1.130966in}{1.430399in}}%
\pgfpathlineto{\pgfqpoint{1.138999in}{1.443647in}}%
\pgfpathlineto{\pgfqpoint{1.140764in}{1.456895in}}%
\pgfpathlineto{\pgfqpoint{1.147771in}{1.469268in}}%
\pgfpathlineto{\pgfqpoint{1.148164in}{1.470143in}}%
\pgfpathlineto{\pgfqpoint{1.150945in}{1.483391in}}%
\pgfpathlineto{\pgfqpoint{1.154173in}{1.496639in}}%
\pgfpathlineto{\pgfqpoint{1.159906in}{1.509887in}}%
\pgfpathlineto{\pgfqpoint{1.161669in}{1.523135in}}%
\pgfpathlineto{\pgfqpoint{1.169288in}{1.536383in}}%
\pgfpathlineto{\pgfqpoint{1.170707in}{1.549631in}}%
\pgfpathlineto{\pgfqpoint{1.173054in}{1.554087in}}%
\pgfpathlineto{\pgfqpoint{1.176547in}{1.562879in}}%
\pgfpathlineto{\pgfqpoint{1.179048in}{1.576127in}}%
\pgfpathlineto{\pgfqpoint{1.182763in}{1.589375in}}%
\pgfpathlineto{\pgfqpoint{1.187094in}{1.602623in}}%
\pgfpathlineto{\pgfqpoint{1.188697in}{1.615871in}}%
\pgfpathlineto{\pgfqpoint{1.195425in}{1.629119in}}%
\pgfpathlineto{\pgfqpoint{1.196896in}{1.642367in}}%
\pgfpathlineto{\pgfqpoint{1.198337in}{1.645058in}}%
\pgfpathlineto{\pgfqpoint{1.202413in}{1.655615in}}%
\pgfpathlineto{\pgfqpoint{1.204300in}{1.668863in}}%
\pgfpathlineto{\pgfqpoint{1.207949in}{1.682111in}}%
\pgfpathlineto{\pgfqpoint{1.211905in}{1.695359in}}%
\pgfpathlineto{\pgfqpoint{1.214118in}{1.708607in}}%
\pgfpathlineto{\pgfqpoint{1.219783in}{1.721855in}}%
\pgfpathlineto{\pgfqpoint{1.220801in}{1.735103in}}%
\pgfpathlineto{\pgfqpoint{1.223620in}{1.741047in}}%
\pgfpathlineto{\pgfqpoint{1.226028in}{1.748351in}}%
\pgfpathlineto{\pgfqpoint{1.227458in}{1.761599in}}%
\pgfpathlineto{\pgfqpoint{1.231209in}{1.774847in}}%
\pgfpathlineto{\pgfqpoint{1.234088in}{1.788095in}}%
\pgfpathlineto{\pgfqpoint{1.236342in}{1.801343in}}%
\pgfpathlineto{\pgfqpoint{1.241184in}{1.814591in}}%
\pgfpathlineto{\pgfqpoint{1.242575in}{1.827839in}}%
\pgfpathlineto{\pgfqpoint{1.248747in}{1.841087in}}%
\pgfpathlineto{\pgfqpoint{1.248904in}{1.842462in}}%
\pgfpathlineto{\pgfqpoint{1.249730in}{1.854335in}}%
\pgfpathlineto{\pgfqpoint{1.253042in}{1.867583in}}%
\pgfpathlineto{\pgfqpoint{1.255303in}{1.880831in}}%
\pgfpathlineto{\pgfqpoint{1.257773in}{1.894079in}}%
\pgfpathlineto{\pgfqpoint{1.261426in}{1.907327in}}%
\pgfpathlineto{\pgfqpoint{1.262708in}{1.920575in}}%
\pgfpathlineto{\pgfqpoint{1.267500in}{1.933823in}}%
\pgfpathlineto{\pgfqpoint{1.268758in}{1.947071in}}%
\pgfpathlineto{\pgfqpoint{1.274187in}{1.960220in}}%
\pgfpathlineto{\pgfqpoint{1.274214in}{1.960319in}}%
\pgfpathlineto{\pgfqpoint{1.275705in}{1.973567in}}%
\pgfpathlineto{\pgfqpoint{1.278071in}{1.986815in}}%
\pgfpathlineto{\pgfqpoint{1.281055in}{2.000063in}}%
\pgfpathlineto{\pgfqpoint{1.282632in}{2.013311in}}%
\pgfpathlineto{\pgfqpoint{1.286706in}{2.026559in}}%
\pgfpathlineto{\pgfqpoint{1.287671in}{2.039807in}}%
\pgfpathlineto{\pgfqpoint{1.291850in}{2.053055in}}%
\pgfpathlineto{\pgfqpoint{1.293542in}{2.066303in}}%
\pgfpathlineto{\pgfqpoint{1.297588in}{2.079551in}}%
\pgfpathlineto{\pgfqpoint{1.299470in}{2.086572in}}%
\pgfpathlineto{\pgfqpoint{1.300473in}{2.092799in}}%
\pgfpathlineto{\pgfqpoint{1.302017in}{2.106047in}}%
\pgfpathlineto{\pgfqpoint{1.305341in}{2.119295in}}%
\pgfpathlineto{\pgfqpoint{1.306419in}{2.132543in}}%
\pgfpathlineto{\pgfqpoint{1.299470in}{2.132543in}}%
\pgfpathlineto{\pgfqpoint{1.288075in}{2.132543in}}%
\pgfpathlineto{\pgfqpoint{1.286817in}{2.119295in}}%
\pgfpathlineto{\pgfqpoint{1.283140in}{2.106047in}}%
\pgfpathlineto{\pgfqpoint{1.281649in}{2.092799in}}%
\pgfpathlineto{\pgfqpoint{1.279433in}{2.079551in}}%
\pgfpathlineto{\pgfqpoint{1.277052in}{2.066303in}}%
\pgfpathlineto{\pgfqpoint{1.275763in}{2.053055in}}%
\pgfpathlineto{\pgfqpoint{1.274187in}{2.045413in}}%
\pgfpathlineto{\pgfqpoint{1.271623in}{2.039807in}}%
\pgfpathlineto{\pgfqpoint{1.270076in}{2.026559in}}%
\pgfpathlineto{\pgfqpoint{1.264369in}{2.013311in}}%
\pgfpathlineto{\pgfqpoint{1.262510in}{2.000063in}}%
\pgfpathlineto{\pgfqpoint{1.258901in}{1.986815in}}%
\pgfpathlineto{\pgfqpoint{1.256574in}{1.973567in}}%
\pgfpathlineto{\pgfqpoint{1.254906in}{1.960319in}}%
\pgfpathlineto{\pgfqpoint{1.251528in}{1.947071in}}%
\pgfpathlineto{\pgfqpoint{1.250549in}{1.933823in}}%
\pgfpathlineto{\pgfqpoint{1.248904in}{1.927414in}}%
\pgfpathlineto{\pgfqpoint{1.245393in}{1.920575in}}%
\pgfpathlineto{\pgfqpoint{1.243693in}{1.907327in}}%
\pgfpathlineto{\pgfqpoint{1.238727in}{1.894079in}}%
\pgfpathlineto{\pgfqpoint{1.235949in}{1.880831in}}%
\pgfpathlineto{\pgfqpoint{1.233062in}{1.867583in}}%
\pgfpathlineto{\pgfqpoint{1.229805in}{1.854335in}}%
\pgfpathlineto{\pgfqpoint{1.228731in}{1.841087in}}%
\pgfpathlineto{\pgfqpoint{1.224683in}{1.827839in}}%
\pgfpathlineto{\pgfqpoint{1.223697in}{1.814591in}}%
\pgfpathlineto{\pgfqpoint{1.223620in}{1.814331in}}%
\pgfpathlineto{\pgfqpoint{1.216795in}{1.801343in}}%
\pgfpathlineto{\pgfqpoint{1.214294in}{1.788095in}}%
\pgfpathlineto{\pgfqpoint{1.210621in}{1.774847in}}%
\pgfpathlineto{\pgfqpoint{1.206652in}{1.761599in}}%
\pgfpathlineto{\pgfqpoint{1.204768in}{1.748351in}}%
\pgfpathlineto{\pgfqpoint{1.200577in}{1.735103in}}%
\pgfpathlineto{\pgfqpoint{1.199824in}{1.721855in}}%
\pgfpathlineto{\pgfqpoint{1.198337in}{1.717308in}}%
\pgfpathlineto{\pgfqpoint{1.193518in}{1.708607in}}%
\pgfpathlineto{\pgfqpoint{1.190956in}{1.695359in}}%
\pgfpathlineto{\pgfqpoint{1.185784in}{1.682111in}}%
\pgfpathlineto{\pgfqpoint{1.182178in}{1.668863in}}%
\pgfpathlineto{\pgfqpoint{1.179858in}{1.655615in}}%
\pgfpathlineto{\pgfqpoint{1.174799in}{1.642367in}}%
\pgfpathlineto{\pgfqpoint{1.173548in}{1.629119in}}%
\pgfpathlineto{\pgfqpoint{1.173054in}{1.627885in}}%
\pgfpathlineto{\pgfqpoint{1.165633in}{1.615871in}}%
\pgfpathlineto{\pgfqpoint{1.163934in}{1.602623in}}%
\pgfpathlineto{\pgfqpoint{1.158735in}{1.589375in}}%
\pgfpathlineto{\pgfqpoint{1.155047in}{1.576127in}}%
\pgfpathlineto{\pgfqpoint{1.151814in}{1.562879in}}%
\pgfpathlineto{\pgfqpoint{1.147771in}{1.552028in}}%
\pgfpathlineto{\pgfqpoint{1.146413in}{1.549631in}}%
\pgfpathlineto{\pgfqpoint{1.144326in}{1.536383in}}%
\pgfpathlineto{\pgfqpoint{1.135878in}{1.523135in}}%
\pgfpathlineto{\pgfqpoint{1.133993in}{1.509887in}}%
\pgfpathlineto{\pgfqpoint{1.127520in}{1.496639in}}%
\pgfpathlineto{\pgfqpoint{1.124636in}{1.483391in}}%
\pgfpathlineto{\pgfqpoint{1.122488in}{1.475022in}}%
\pgfpathlineto{\pgfqpoint{1.120559in}{1.470143in}}%
\pgfpathlineto{\pgfqpoint{1.112734in}{1.456895in}}%
\pgfpathlineto{\pgfqpoint{1.110074in}{1.443647in}}%
\pgfpathlineto{\pgfqpoint{1.101980in}{1.430399in}}%
\pgfpathlineto{\pgfqpoint{1.100764in}{1.417151in}}%
\pgfpathlineto{\pgfqpoint{1.097204in}{1.410314in}}%
\pgfpathlineto{\pgfqpoint{1.092617in}{1.403903in}}%
\pgfpathlineto{\pgfqpoint{1.088552in}{1.390655in}}%
\pgfpathlineto{\pgfqpoint{1.082348in}{1.377407in}}%
\pgfpathlineto{\pgfqpoint{1.076237in}{1.364159in}}%
\pgfpathlineto{\pgfqpoint{1.073812in}{1.350911in}}%
\pgfpathlineto{\pgfqpoint{1.071921in}{1.348089in}}%
\pgfpathlineto{\pgfqpoint{1.063074in}{1.337663in}}%
\pgfpathlineto{\pgfqpoint{1.060560in}{1.324415in}}%
\pgfpathlineto{\pgfqpoint{1.050250in}{1.311167in}}%
\pgfpathlineto{\pgfqpoint{1.047629in}{1.297919in}}%
\pgfpathlineto{\pgfqpoint{1.046638in}{1.295818in}}%
\pgfpathlineto{\pgfqpoint{1.039691in}{1.284671in}}%
\pgfpathlineto{\pgfqpoint{1.032011in}{1.271423in}}%
\pgfpathlineto{\pgfqpoint{1.028191in}{1.258175in}}%
\pgfpathlineto{\pgfqpoint{1.021355in}{1.249299in}}%
\pgfpathlineto{\pgfqpoint{1.017294in}{1.244927in}}%
\pgfpathlineto{\pgfqpoint{1.015483in}{1.231679in}}%
\pgfpathlineto{\pgfqpoint{1.003103in}{1.218431in}}%
\pgfpathlineto{\pgfqpoint{0.999247in}{1.205183in}}%
\pgfpathlineto{\pgfqpoint{0.996072in}{1.200981in}}%
\pgfpathlineto{\pgfqpoint{0.987576in}{1.191935in}}%
\pgfpathlineto{\pgfqpoint{0.979896in}{1.178687in}}%
\pgfpathlineto{\pgfqpoint{0.975048in}{1.165439in}}%
\pgfpathlineto{\pgfqpoint{0.970789in}{1.160907in}}%
\pgfpathlineto{\pgfqpoint{0.961464in}{1.152191in}}%
\pgfpathlineto{\pgfqpoint{0.959355in}{1.138943in}}%
\pgfpathlineto{\pgfqpoint{0.945505in}{1.127238in}}%
\pgfpathlineto{\pgfqpoint{0.943422in}{1.125695in}}%
\pgfpathlineto{\pgfqpoint{0.940757in}{1.112447in}}%
\pgfpathlineto{\pgfqpoint{0.927911in}{1.099199in}}%
\pgfpathlineto{\pgfqpoint{0.920222in}{1.087727in}}%
\pgfpathlineto{\pgfqpoint{0.918866in}{1.085951in}}%
\pgfpathlineto{\pgfqpoint{0.909863in}{1.072703in}}%
\pgfpathlineto{\pgfqpoint{0.894939in}{1.060153in}}%
\pgfpathlineto{\pgfqpoint{0.894046in}{1.059455in}}%
\pgfpathlineto{\pgfqpoint{0.892175in}{1.046207in}}%
\pgfpathlineto{\pgfqpoint{0.872527in}{1.032959in}}%
\pgfpathlineto{\pgfqpoint{0.869664in}{1.019711in}}%
\pgfpathlineto{\pgfqpoint{0.869656in}{1.019705in}}%
\pgfpathlineto{\pgfqpoint{0.849847in}{1.006463in}}%
\pgfpathlineto{\pgfqpoint{0.844373in}{0.998067in}}%
\pgfpathlineto{\pgfqpoint{0.840901in}{0.993215in}}%
\pgfpathlineto{\pgfqpoint{0.831025in}{0.979967in}}%
\pgfpathlineto{\pgfqpoint{0.819090in}{0.972006in}}%
\pgfpathlineto{\pgfqpoint{0.810833in}{0.966719in}}%
\pgfpathlineto{\pgfqpoint{0.808057in}{0.953471in}}%
\pgfpathlineto{\pgfqpoint{0.793806in}{0.946720in}}%
\pgfpathlineto{\pgfqpoint{0.779536in}{0.940223in}}%
\pgfpathlineto{\pgfqpoint{0.780081in}{0.926975in}}%
\pgfpathlineto{\pgfqpoint{0.768523in}{0.919801in}}%
\pgfpathlineto{\pgfqpoint{0.757720in}{0.913727in}}%
\pgfpathlineto{\pgfqpoint{0.746278in}{0.900479in}}%
\pgfpathlineto{\pgfqpoint{0.743240in}{0.897206in}}%
\pgfpathlineto{\pgfqpoint{0.732316in}{0.887231in}}%
\pgfpathlineto{\pgfqpoint{0.717957in}{0.880341in}}%
\pgfpathlineto{\pgfqpoint{0.704643in}{0.873983in}}%
\pgfpathlineto{\pgfqpoint{0.704955in}{0.860735in}}%
\pgfpathlineto{\pgfqpoint{0.692674in}{0.856197in}}%
\pgfpathlineto{\pgfqpoint{0.668390in}{0.847487in}}%
\pgfpathlineto{\pgfqpoint{0.669599in}{0.834239in}}%
\pgfpathlineto{\pgfqpoint{0.667391in}{0.833330in}}%
\pgfpathlineto{\pgfqpoint{0.642107in}{0.823955in}}%
\pgfpathlineto{\pgfqpoint{0.633268in}{0.820991in}}%
\pgfpathlineto{\pgfqpoint{0.624906in}{0.807743in}}%
\pgfpathlineto{\pgfqpoint{0.618020in}{0.794495in}}%
\pgfpathlineto{\pgfqpoint{0.616824in}{0.794040in}}%
\pgfpathlineto{\pgfqpoint{0.591541in}{0.784655in}}%
\pgfpathlineto{\pgfqpoint{0.582229in}{0.781247in}}%
\pgfpathlineto{\pgfqpoint{0.591541in}{0.771843in}}%
\pgfpathlineto{\pgfqpoint{0.595777in}{0.767999in}}%
\pgfpathlineto{\pgfqpoint{0.591541in}{0.766981in}}%
\pgfpathlineto{\pgfqpoint{0.566258in}{0.761233in}}%
\pgfpathlineto{\pgfqpoint{0.540975in}{0.755797in}}%
\pgfpathlineto{\pgfqpoint{0.540975in}{0.754751in}}%
\pgfpathlineto{\pgfqpoint{0.540975in}{0.751581in}}%
\pgfpathlineto{\pgfqpoint{0.558099in}{0.741503in}}%
\pgfpathlineto{\pgfqpoint{0.540975in}{0.734579in}}%
\pgfpathlineto{\pgfqpoint{0.540975in}{0.728255in}}%
\pgfpathlineto{\pgfqpoint{0.540975in}{0.715007in}}%
\pgfpathlineto{\pgfqpoint{0.540975in}{0.701759in}}%
\pgfpathlineto{\pgfqpoint{0.540975in}{0.688511in}}%
\pgfpathlineto{\pgfqpoint{0.540975in}{0.680137in}}%
\pgfpathlineto{\pgfqpoint{0.566258in}{0.677589in}}%
\pgfpathlineto{\pgfqpoint{0.591541in}{0.675278in}}%
\pgfpathlineto{\pgfqpoint{0.591709in}{0.675263in}}%
\pgfpathlineto{\pgfqpoint{0.591541in}{0.675172in}}%
\pgfpathlineto{\pgfqpoint{0.566258in}{0.664037in}}%
\pgfpathlineto{\pgfqpoint{0.560402in}{0.662015in}}%
\pgfpathlineto{\pgfqpoint{0.566258in}{0.661360in}}%
\pgfpathlineto{\pgfqpoint{0.591541in}{0.658656in}}%
\pgfpathlineto{\pgfqpoint{0.616824in}{0.656121in}}%
\pgfpathlineto{\pgfqpoint{0.642107in}{0.653737in}}%
\pgfpathlineto{\pgfqpoint{0.667391in}{0.651493in}}%
\pgfpathlineto{\pgfqpoint{0.692674in}{0.649374in}}%
\pgfpathlineto{\pgfqpoint{0.700238in}{0.648767in}}%
\pgfpathlineto{\pgfqpoint{0.717957in}{0.647154in}}%
\pgfpathlineto{\pgfqpoint{0.743240in}{0.645148in}}%
\pgfpathlineto{\pgfqpoint{0.768523in}{0.643424in}}%
\pgfpathlineto{\pgfqpoint{0.793806in}{0.641932in}}%
\pgfpathlineto{\pgfqpoint{0.819090in}{0.640635in}}%
\pgfpathlineto{\pgfqpoint{0.844373in}{0.639500in}}%
\pgfpathlineto{\pgfqpoint{0.869656in}{0.638505in}}%
\pgfpathlineto{\pgfqpoint{0.894939in}{0.637628in}}%
\pgfpathlineto{\pgfqpoint{0.920222in}{0.636855in}}%
\pgfpathlineto{\pgfqpoint{0.945505in}{0.636171in}}%
\pgfpathlineto{\pgfqpoint{0.970789in}{0.635566in}}%
\pgfpathlineto{\pgfqpoint{0.972946in}{0.635519in}}%
\pgfpathlineto{\pgfqpoint{0.996072in}{0.633295in}}%
\pgfpathlineto{\pgfqpoint{1.021355in}{0.631116in}}%
\pgfpathlineto{\pgfqpoint{1.046638in}{0.629175in}}%
\pgfpathlineto{\pgfqpoint{1.071921in}{0.627455in}}%
\pgfpathlineto{\pgfqpoint{1.097204in}{0.625938in}}%
\pgfpathlineto{\pgfqpoint{1.122488in}{0.624610in}}%
\pgfpathlineto{\pgfqpoint{1.147771in}{0.623458in}}%
\pgfpathlineto{\pgfqpoint{1.173054in}{0.622470in}}%
\pgfpathlineto{\pgfqpoint{1.178998in}{0.622271in}}%
\pgfpathlineto{\pgfqpoint{1.198337in}{0.622147in}}%
\pgfpathclose%
\pgfpathmoveto{\pgfqpoint{0.566258in}{1.147018in}}%
\pgfpathlineto{\pgfqpoint{0.591541in}{1.152082in}}%
\pgfpathlineto{\pgfqpoint{0.591847in}{1.152191in}}%
\pgfpathlineto{\pgfqpoint{0.591541in}{1.152563in}}%
\pgfpathlineto{\pgfqpoint{0.566258in}{1.160485in}}%
\pgfpathlineto{\pgfqpoint{0.540975in}{1.158551in}}%
\pgfpathlineto{\pgfqpoint{0.540975in}{1.152191in}}%
\pgfpathlineto{\pgfqpoint{0.540975in}{1.145891in}}%
\pgfpathlineto{\pgfqpoint{0.566258in}{1.147018in}}%
\pgfpathclose%
\pgfpathmoveto{\pgfqpoint{0.566258in}{1.167298in}}%
\pgfpathlineto{\pgfqpoint{0.591541in}{1.167506in}}%
\pgfpathlineto{\pgfqpoint{0.616824in}{1.169829in}}%
\pgfpathlineto{\pgfqpoint{0.642107in}{1.173984in}}%
\pgfpathlineto{\pgfqpoint{0.661769in}{1.178687in}}%
\pgfpathlineto{\pgfqpoint{0.649277in}{1.191935in}}%
\pgfpathlineto{\pgfqpoint{0.667391in}{1.195348in}}%
\pgfpathlineto{\pgfqpoint{0.692674in}{1.201496in}}%
\pgfpathlineto{\pgfqpoint{0.704701in}{1.205183in}}%
\pgfpathlineto{\pgfqpoint{0.704069in}{1.218431in}}%
\pgfpathlineto{\pgfqpoint{0.717957in}{1.223684in}}%
\pgfpathlineto{\pgfqpoint{0.737332in}{1.231679in}}%
\pgfpathlineto{\pgfqpoint{0.743240in}{1.234973in}}%
\pgfpathlineto{\pgfqpoint{0.756503in}{1.244927in}}%
\pgfpathlineto{\pgfqpoint{0.768523in}{1.256544in}}%
\pgfpathlineto{\pgfqpoint{0.770160in}{1.258175in}}%
\pgfpathlineto{\pgfqpoint{0.793806in}{1.270612in}}%
\pgfpathlineto{\pgfqpoint{0.795026in}{1.271423in}}%
\pgfpathlineto{\pgfqpoint{0.795991in}{1.284671in}}%
\pgfpathlineto{\pgfqpoint{0.819090in}{1.294139in}}%
\pgfpathlineto{\pgfqpoint{0.826761in}{1.297919in}}%
\pgfpathlineto{\pgfqpoint{0.828797in}{1.311167in}}%
\pgfpathlineto{\pgfqpoint{0.844373in}{1.319418in}}%
\pgfpathlineto{\pgfqpoint{0.852798in}{1.324415in}}%
\pgfpathlineto{\pgfqpoint{0.860736in}{1.337663in}}%
\pgfpathlineto{\pgfqpoint{0.869656in}{1.346375in}}%
\pgfpathlineto{\pgfqpoint{0.874112in}{1.350911in}}%
\pgfpathlineto{\pgfqpoint{0.891160in}{1.364159in}}%
\pgfpathlineto{\pgfqpoint{0.894939in}{1.374002in}}%
\pgfpathlineto{\pgfqpoint{0.896250in}{1.377407in}}%
\pgfpathlineto{\pgfqpoint{0.914681in}{1.390655in}}%
\pgfpathlineto{\pgfqpoint{0.915806in}{1.403903in}}%
\pgfpathlineto{\pgfqpoint{0.920222in}{1.406867in}}%
\pgfpathlineto{\pgfqpoint{0.933954in}{1.417151in}}%
\pgfpathlineto{\pgfqpoint{0.939163in}{1.430399in}}%
\pgfpathlineto{\pgfqpoint{0.945505in}{1.436947in}}%
\pgfpathlineto{\pgfqpoint{0.951828in}{1.443647in}}%
\pgfpathlineto{\pgfqpoint{0.961835in}{1.456895in}}%
\pgfpathlineto{\pgfqpoint{0.967736in}{1.470143in}}%
\pgfpathlineto{\pgfqpoint{0.970789in}{1.472531in}}%
\pgfpathlineto{\pgfqpoint{0.982975in}{1.483391in}}%
\pgfpathlineto{\pgfqpoint{0.985476in}{1.496639in}}%
\pgfpathlineto{\pgfqpoint{0.996072in}{1.506352in}}%
\pgfpathlineto{\pgfqpoint{0.999726in}{1.509887in}}%
\pgfpathlineto{\pgfqpoint{1.002220in}{1.523135in}}%
\pgfpathlineto{\pgfqpoint{1.012875in}{1.536383in}}%
\pgfpathlineto{\pgfqpoint{1.020061in}{1.549631in}}%
\pgfpathlineto{\pgfqpoint{1.021355in}{1.552416in}}%
\pgfpathlineto{\pgfqpoint{1.026546in}{1.562879in}}%
\pgfpathlineto{\pgfqpoint{1.036611in}{1.576127in}}%
\pgfpathlineto{\pgfqpoint{1.039213in}{1.589375in}}%
\pgfpathlineto{\pgfqpoint{1.046638in}{1.597392in}}%
\pgfpathlineto{\pgfqpoint{1.051440in}{1.602623in}}%
\pgfpathlineto{\pgfqpoint{1.053851in}{1.615871in}}%
\pgfpathlineto{\pgfqpoint{1.063886in}{1.629119in}}%
\pgfpathlineto{\pgfqpoint{1.067856in}{1.642367in}}%
\pgfpathlineto{\pgfqpoint{1.071921in}{1.651497in}}%
\pgfpathlineto{\pgfqpoint{1.074003in}{1.655615in}}%
\pgfpathlineto{\pgfqpoint{1.082846in}{1.668863in}}%
\pgfpathlineto{\pgfqpoint{1.086029in}{1.682111in}}%
\pgfpathlineto{\pgfqpoint{1.094930in}{1.695359in}}%
\pgfpathlineto{\pgfqpoint{1.096576in}{1.708607in}}%
\pgfpathlineto{\pgfqpoint{1.097204in}{1.709558in}}%
\pgfpathlineto{\pgfqpoint{1.105930in}{1.721855in}}%
\pgfpathlineto{\pgfqpoint{1.109173in}{1.735103in}}%
\pgfpathlineto{\pgfqpoint{1.115403in}{1.748351in}}%
\pgfpathlineto{\pgfqpoint{1.120336in}{1.761599in}}%
\pgfpathlineto{\pgfqpoint{1.122488in}{1.771539in}}%
\pgfpathlineto{\pgfqpoint{1.123406in}{1.774847in}}%
\pgfpathlineto{\pgfqpoint{1.132626in}{1.788095in}}%
\pgfpathlineto{\pgfqpoint{1.134793in}{1.801343in}}%
\pgfpathlineto{\pgfqpoint{1.142386in}{1.814591in}}%
\pgfpathlineto{\pgfqpoint{1.144362in}{1.827839in}}%
\pgfpathlineto{\pgfqpoint{1.147771in}{1.836250in}}%
\pgfpathlineto{\pgfqpoint{1.150293in}{1.841087in}}%
\pgfpathlineto{\pgfqpoint{1.155310in}{1.854335in}}%
\pgfpathlineto{\pgfqpoint{1.159069in}{1.867583in}}%
\pgfpathlineto{\pgfqpoint{1.164815in}{1.880831in}}%
\pgfpathlineto{\pgfqpoint{1.166181in}{1.894079in}}%
\pgfpathlineto{\pgfqpoint{1.172667in}{1.907327in}}%
\pgfpathlineto{\pgfqpoint{1.173054in}{1.910052in}}%
\pgfpathlineto{\pgfqpoint{1.175208in}{1.920575in}}%
\pgfpathlineto{\pgfqpoint{1.182357in}{1.933823in}}%
\pgfpathlineto{\pgfqpoint{1.185436in}{1.947071in}}%
\pgfpathlineto{\pgfqpoint{1.188729in}{1.960319in}}%
\pgfpathlineto{\pgfqpoint{1.193570in}{1.973567in}}%
\pgfpathlineto{\pgfqpoint{1.195338in}{1.986815in}}%
\pgfpathlineto{\pgfqpoint{1.198337in}{1.993974in}}%
\pgfpathlineto{\pgfqpoint{1.201996in}{2.000063in}}%
\pgfpathlineto{\pgfqpoint{1.203538in}{2.013311in}}%
\pgfpathlineto{\pgfqpoint{1.209551in}{2.026559in}}%
\pgfpathlineto{\pgfqpoint{1.212190in}{2.039807in}}%
\pgfpathlineto{\pgfqpoint{1.216086in}{2.053055in}}%
\pgfpathlineto{\pgfqpoint{1.219357in}{2.066303in}}%
\pgfpathlineto{\pgfqpoint{1.220954in}{2.079551in}}%
\pgfpathlineto{\pgfqpoint{1.223620in}{2.087068in}}%
\pgfpathlineto{\pgfqpoint{1.226888in}{2.092799in}}%
\pgfpathlineto{\pgfqpoint{1.228725in}{2.106047in}}%
\pgfpathlineto{\pgfqpoint{1.235144in}{2.119295in}}%
\pgfpathlineto{\pgfqpoint{1.237050in}{2.132543in}}%
\pgfpathlineto{\pgfqpoint{1.223620in}{2.132543in}}%
\pgfpathlineto{\pgfqpoint{1.214142in}{2.132543in}}%
\pgfpathlineto{\pgfqpoint{1.211964in}{2.119295in}}%
\pgfpathlineto{\pgfqpoint{1.205926in}{2.106047in}}%
\pgfpathlineto{\pgfqpoint{1.204584in}{2.092799in}}%
\pgfpathlineto{\pgfqpoint{1.198337in}{2.080520in}}%
\pgfpathlineto{\pgfqpoint{1.197861in}{2.079551in}}%
\pgfpathlineto{\pgfqpoint{1.196241in}{2.066303in}}%
\pgfpathlineto{\pgfqpoint{1.191611in}{2.053055in}}%
\pgfpathlineto{\pgfqpoint{1.187718in}{2.039807in}}%
\pgfpathlineto{\pgfqpoint{1.184505in}{2.026559in}}%
\pgfpathlineto{\pgfqpoint{1.178554in}{2.013311in}}%
\pgfpathlineto{\pgfqpoint{1.177252in}{2.000063in}}%
\pgfpathlineto{\pgfqpoint{1.173054in}{1.992618in}}%
\pgfpathlineto{\pgfqpoint{1.169789in}{1.986815in}}%
\pgfpathlineto{\pgfqpoint{1.168266in}{1.973567in}}%
\pgfpathlineto{\pgfqpoint{1.161864in}{1.960319in}}%
\pgfpathlineto{\pgfqpoint{1.158689in}{1.947071in}}%
\pgfpathlineto{\pgfqpoint{1.154670in}{1.933823in}}%
\pgfpathlineto{\pgfqpoint{1.147926in}{1.920575in}}%
\pgfpathlineto{\pgfqpoint{1.147771in}{1.919756in}}%
\pgfpathlineto{\pgfqpoint{1.145434in}{1.907327in}}%
\pgfpathlineto{\pgfqpoint{1.136974in}{1.894079in}}%
\pgfpathlineto{\pgfqpoint{1.135953in}{1.880831in}}%
\pgfpathlineto{\pgfqpoint{1.128320in}{1.867583in}}%
\pgfpathlineto{\pgfqpoint{1.125421in}{1.854335in}}%
\pgfpathlineto{\pgfqpoint{1.122488in}{1.847614in}}%
\pgfpathlineto{\pgfqpoint{1.119360in}{1.841087in}}%
\pgfpathlineto{\pgfqpoint{1.113310in}{1.827839in}}%
\pgfpathlineto{\pgfqpoint{1.110540in}{1.814591in}}%
\pgfpathlineto{\pgfqpoint{1.101211in}{1.801343in}}%
\pgfpathlineto{\pgfqpoint{1.099459in}{1.788095in}}%
\pgfpathlineto{\pgfqpoint{1.097204in}{1.785150in}}%
\pgfpathlineto{\pgfqpoint{1.088207in}{1.774847in}}%
\pgfpathlineto{\pgfqpoint{1.086072in}{1.761599in}}%
\pgfpathlineto{\pgfqpoint{1.078430in}{1.748351in}}%
\pgfpathlineto{\pgfqpoint{1.072491in}{1.735103in}}%
\pgfpathlineto{\pgfqpoint{1.071921in}{1.733338in}}%
\pgfpathlineto{\pgfqpoint{1.067598in}{1.721855in}}%
\pgfpathlineto{\pgfqpoint{1.057128in}{1.708607in}}%
\pgfpathlineto{\pgfqpoint{1.055766in}{1.695359in}}%
\pgfpathlineto{\pgfqpoint{1.046638in}{1.685422in}}%
\pgfpathlineto{\pgfqpoint{1.042777in}{1.682111in}}%
\pgfpathlineto{\pgfqpoint{1.041093in}{1.668863in}}%
\pgfpathlineto{\pgfqpoint{1.028383in}{1.655615in}}%
\pgfpathlineto{\pgfqpoint{1.023652in}{1.642367in}}%
\pgfpathlineto{\pgfqpoint{1.021355in}{1.637338in}}%
\pgfpathlineto{\pgfqpoint{1.016512in}{1.629119in}}%
\pgfpathlineto{\pgfqpoint{1.004393in}{1.615871in}}%
\pgfpathlineto{\pgfqpoint{1.001905in}{1.602623in}}%
\pgfpathlineto{\pgfqpoint{0.996072in}{1.597670in}}%
\pgfpathlineto{\pgfqpoint{0.983282in}{1.589375in}}%
\pgfpathlineto{\pgfqpoint{0.983176in}{1.576127in}}%
\pgfpathlineto{\pgfqpoint{0.970789in}{1.566390in}}%
\pgfpathlineto{\pgfqpoint{0.964020in}{1.562879in}}%
\pgfpathlineto{\pgfqpoint{0.961172in}{1.549631in}}%
\pgfpathlineto{\pgfqpoint{0.947123in}{1.536383in}}%
\pgfpathlineto{\pgfqpoint{0.945505in}{1.534361in}}%
\pgfpathlineto{\pgfqpoint{0.934006in}{1.523135in}}%
\pgfpathlineto{\pgfqpoint{0.930932in}{1.509887in}}%
\pgfpathlineto{\pgfqpoint{0.920222in}{1.503483in}}%
\pgfpathlineto{\pgfqpoint{0.903235in}{1.496639in}}%
\pgfpathlineto{\pgfqpoint{0.907701in}{1.483391in}}%
\pgfpathlineto{\pgfqpoint{0.894939in}{1.478203in}}%
\pgfpathlineto{\pgfqpoint{0.869656in}{1.473090in}}%
\pgfpathlineto{\pgfqpoint{0.844373in}{1.474302in}}%
\pgfpathlineto{\pgfqpoint{0.825644in}{1.483391in}}%
\pgfpathlineto{\pgfqpoint{0.840588in}{1.496639in}}%
\pgfpathlineto{\pgfqpoint{0.837494in}{1.509887in}}%
\pgfpathlineto{\pgfqpoint{0.839741in}{1.523135in}}%
\pgfpathlineto{\pgfqpoint{0.844373in}{1.527851in}}%
\pgfpathlineto{\pgfqpoint{0.854201in}{1.536383in}}%
\pgfpathlineto{\pgfqpoint{0.844872in}{1.549631in}}%
\pgfpathlineto{\pgfqpoint{0.863001in}{1.562879in}}%
\pgfpathlineto{\pgfqpoint{0.856014in}{1.576127in}}%
\pgfpathlineto{\pgfqpoint{0.867224in}{1.589375in}}%
\pgfpathlineto{\pgfqpoint{0.867822in}{1.602623in}}%
\pgfpathlineto{\pgfqpoint{0.869609in}{1.615871in}}%
\pgfpathlineto{\pgfqpoint{0.869656in}{1.615933in}}%
\pgfpathlineto{\pgfqpoint{0.880304in}{1.629119in}}%
\pgfpathlineto{\pgfqpoint{0.878186in}{1.642367in}}%
\pgfpathlineto{\pgfqpoint{0.891561in}{1.655615in}}%
\pgfpathlineto{\pgfqpoint{0.889395in}{1.668863in}}%
\pgfpathlineto{\pgfqpoint{0.894939in}{1.676657in}}%
\pgfpathlineto{\pgfqpoint{0.899416in}{1.682111in}}%
\pgfpathlineto{\pgfqpoint{0.901864in}{1.695359in}}%
\pgfpathlineto{\pgfqpoint{0.906069in}{1.708607in}}%
\pgfpathlineto{\pgfqpoint{0.913550in}{1.721855in}}%
\pgfpathlineto{\pgfqpoint{0.912336in}{1.735103in}}%
\pgfpathlineto{\pgfqpoint{0.920222in}{1.745207in}}%
\pgfpathlineto{\pgfqpoint{0.922857in}{1.748351in}}%
\pgfpathlineto{\pgfqpoint{0.921869in}{1.761599in}}%
\pgfpathlineto{\pgfqpoint{0.931261in}{1.774847in}}%
\pgfpathlineto{\pgfqpoint{0.933457in}{1.788095in}}%
\pgfpathlineto{\pgfqpoint{0.937582in}{1.801343in}}%
\pgfpathlineto{\pgfqpoint{0.944230in}{1.814591in}}%
\pgfpathlineto{\pgfqpoint{0.944908in}{1.827839in}}%
\pgfpathlineto{\pgfqpoint{0.945505in}{1.828685in}}%
\pgfpathlineto{\pgfqpoint{0.954118in}{1.841087in}}%
\pgfpathlineto{\pgfqpoint{0.953440in}{1.854335in}}%
\pgfpathlineto{\pgfqpoint{0.961187in}{1.867583in}}%
\pgfpathlineto{\pgfqpoint{0.963237in}{1.880831in}}%
\pgfpathlineto{\pgfqpoint{0.968243in}{1.894079in}}%
\pgfpathlineto{\pgfqpoint{0.970789in}{1.900041in}}%
\pgfpathlineto{\pgfqpoint{0.973656in}{1.907327in}}%
\pgfpathlineto{\pgfqpoint{0.974640in}{1.920575in}}%
\pgfpathlineto{\pgfqpoint{0.981958in}{1.933823in}}%
\pgfpathlineto{\pgfqpoint{0.981883in}{1.947071in}}%
\pgfpathlineto{\pgfqpoint{0.989396in}{1.960319in}}%
\pgfpathlineto{\pgfqpoint{0.990763in}{1.973567in}}%
\pgfpathlineto{\pgfqpoint{0.995626in}{1.986815in}}%
\pgfpathlineto{\pgfqpoint{0.996072in}{1.987920in}}%
\pgfpathlineto{\pgfqpoint{1.000351in}{2.000063in}}%
\pgfpathlineto{\pgfqpoint{1.002389in}{2.013311in}}%
\pgfpathlineto{\pgfqpoint{1.008861in}{2.026559in}}%
\pgfpathlineto{\pgfqpoint{1.009027in}{2.039807in}}%
\pgfpathlineto{\pgfqpoint{1.015294in}{2.053055in}}%
\pgfpathlineto{\pgfqpoint{1.016345in}{2.066303in}}%
\pgfpathlineto{\pgfqpoint{1.021355in}{2.079044in}}%
\pgfpathlineto{\pgfqpoint{1.021540in}{2.079551in}}%
\pgfpathlineto{\pgfqpoint{1.024830in}{2.092799in}}%
\pgfpathlineto{\pgfqpoint{1.026801in}{2.106047in}}%
\pgfpathlineto{\pgfqpoint{1.032406in}{2.119295in}}%
\pgfpathlineto{\pgfqpoint{1.033272in}{2.132543in}}%
\pgfpathlineto{\pgfqpoint{1.021355in}{2.132543in}}%
\pgfpathlineto{\pgfqpoint{1.008376in}{2.132543in}}%
\pgfpathlineto{\pgfqpoint{1.007010in}{2.119295in}}%
\pgfpathlineto{\pgfqpoint{1.001277in}{2.106047in}}%
\pgfpathlineto{\pgfqpoint{0.999556in}{2.092799in}}%
\pgfpathlineto{\pgfqpoint{0.996072in}{2.080826in}}%
\pgfpathlineto{\pgfqpoint{0.995514in}{2.079551in}}%
\pgfpathlineto{\pgfqpoint{0.990163in}{2.066303in}}%
\pgfpathlineto{\pgfqpoint{0.987550in}{2.053055in}}%
\pgfpathlineto{\pgfqpoint{0.981623in}{2.039807in}}%
\pgfpathlineto{\pgfqpoint{0.980612in}{2.026559in}}%
\pgfpathlineto{\pgfqpoint{0.974068in}{2.013311in}}%
\pgfpathlineto{\pgfqpoint{0.972219in}{2.000063in}}%
\pgfpathlineto{\pgfqpoint{0.970789in}{1.996469in}}%
\pgfpathlineto{\pgfqpoint{0.965576in}{1.986815in}}%
\pgfpathlineto{\pgfqpoint{0.961266in}{1.973567in}}%
\pgfpathlineto{\pgfqpoint{0.958374in}{1.960319in}}%
\pgfpathlineto{\pgfqpoint{0.951523in}{1.947071in}}%
\pgfpathlineto{\pgfqpoint{0.950322in}{1.933823in}}%
\pgfpathlineto{\pgfqpoint{0.945505in}{1.924784in}}%
\pgfpathlineto{\pgfqpoint{0.942660in}{1.920575in}}%
\pgfpathlineto{\pgfqpoint{0.940998in}{1.907327in}}%
\pgfpathlineto{\pgfqpoint{0.933822in}{1.894079in}}%
\pgfpathlineto{\pgfqpoint{0.929580in}{1.880831in}}%
\pgfpathlineto{\pgfqpoint{0.925509in}{1.867583in}}%
\pgfpathlineto{\pgfqpoint{0.920222in}{1.856386in}}%
\pgfpathlineto{\pgfqpoint{0.919050in}{1.854335in}}%
\pgfpathlineto{\pgfqpoint{0.917659in}{1.841087in}}%
\pgfpathlineto{\pgfqpoint{0.907847in}{1.827839in}}%
\pgfpathlineto{\pgfqpoint{0.906429in}{1.814591in}}%
\pgfpathlineto{\pgfqpoint{0.898416in}{1.801343in}}%
\pgfpathlineto{\pgfqpoint{0.895048in}{1.788095in}}%
\pgfpathlineto{\pgfqpoint{0.894939in}{1.787786in}}%
\pgfpathlineto{\pgfqpoint{0.889213in}{1.774847in}}%
\pgfpathlineto{\pgfqpoint{0.880939in}{1.761599in}}%
\pgfpathlineto{\pgfqpoint{0.878906in}{1.748351in}}%
\pgfpathlineto{\pgfqpoint{0.869656in}{1.735326in}}%
\pgfpathlineto{\pgfqpoint{0.869476in}{1.735103in}}%
\pgfpathlineto{\pgfqpoint{0.868937in}{1.721855in}}%
\pgfpathlineto{\pgfqpoint{0.858833in}{1.708607in}}%
\pgfpathlineto{\pgfqpoint{0.855088in}{1.695359in}}%
\pgfpathlineto{\pgfqpoint{0.848486in}{1.682111in}}%
\pgfpathlineto{\pgfqpoint{0.844373in}{1.674942in}}%
\pgfpathlineto{\pgfqpoint{0.840593in}{1.668863in}}%
\pgfpathlineto{\pgfqpoint{0.837858in}{1.655615in}}%
\pgfpathlineto{\pgfqpoint{0.825461in}{1.642367in}}%
\pgfpathlineto{\pgfqpoint{0.824470in}{1.629119in}}%
\pgfpathlineto{\pgfqpoint{0.819090in}{1.622986in}}%
\pgfpathlineto{\pgfqpoint{0.812239in}{1.615871in}}%
\pgfpathlineto{\pgfqpoint{0.810007in}{1.602623in}}%
\pgfpathlineto{\pgfqpoint{0.802751in}{1.589375in}}%
\pgfpathlineto{\pgfqpoint{0.794018in}{1.576127in}}%
\pgfpathlineto{\pgfqpoint{0.793806in}{1.575096in}}%
\pgfpathlineto{\pgfqpoint{0.790805in}{1.562879in}}%
\pgfpathlineto{\pgfqpoint{0.777252in}{1.549631in}}%
\pgfpathlineto{\pgfqpoint{0.777275in}{1.536383in}}%
\pgfpathlineto{\pgfqpoint{0.768523in}{1.528607in}}%
\pgfpathlineto{\pgfqpoint{0.762067in}{1.523135in}}%
\pgfpathlineto{\pgfqpoint{0.758919in}{1.509887in}}%
\pgfpathlineto{\pgfqpoint{0.748387in}{1.496639in}}%
\pgfpathlineto{\pgfqpoint{0.743240in}{1.488072in}}%
\pgfpathlineto{\pgfqpoint{0.740410in}{1.483391in}}%
\pgfpathlineto{\pgfqpoint{0.737862in}{1.470143in}}%
\pgfpathlineto{\pgfqpoint{0.721476in}{1.456895in}}%
\pgfpathlineto{\pgfqpoint{0.722024in}{1.443647in}}%
\pgfpathlineto{\pgfqpoint{0.717957in}{1.440746in}}%
\pgfpathlineto{\pgfqpoint{0.703379in}{1.430399in}}%
\pgfpathlineto{\pgfqpoint{0.703247in}{1.417151in}}%
\pgfpathlineto{\pgfqpoint{0.692674in}{1.406112in}}%
\pgfpathlineto{\pgfqpoint{0.690385in}{1.403903in}}%
\pgfpathlineto{\pgfqpoint{0.679995in}{1.390655in}}%
\pgfpathlineto{\pgfqpoint{0.676506in}{1.377407in}}%
\pgfpathlineto{\pgfqpoint{0.667391in}{1.371033in}}%
\pgfpathlineto{\pgfqpoint{0.658255in}{1.364159in}}%
\pgfpathlineto{\pgfqpoint{0.662558in}{1.350911in}}%
\pgfpathlineto{\pgfqpoint{0.642107in}{1.339137in}}%
\pgfpathlineto{\pgfqpoint{0.639563in}{1.337663in}}%
\pgfpathlineto{\pgfqpoint{0.638711in}{1.324415in}}%
\pgfpathlineto{\pgfqpoint{0.621106in}{1.311167in}}%
\pgfpathlineto{\pgfqpoint{0.616824in}{1.304541in}}%
\pgfpathlineto{\pgfqpoint{0.612928in}{1.297919in}}%
\pgfpathlineto{\pgfqpoint{0.611795in}{1.284671in}}%
\pgfpathlineto{\pgfqpoint{0.591541in}{1.274236in}}%
\pgfpathlineto{\pgfqpoint{0.586611in}{1.271423in}}%
\pgfpathlineto{\pgfqpoint{0.591541in}{1.263823in}}%
\pgfpathlineto{\pgfqpoint{0.596630in}{1.258175in}}%
\pgfpathlineto{\pgfqpoint{0.591541in}{1.256050in}}%
\pgfpathlineto{\pgfqpoint{0.566258in}{1.245076in}}%
\pgfpathlineto{\pgfqpoint{0.565932in}{1.244927in}}%
\pgfpathlineto{\pgfqpoint{0.566258in}{1.244475in}}%
\pgfpathlineto{\pgfqpoint{0.578353in}{1.231679in}}%
\pgfpathlineto{\pgfqpoint{0.566258in}{1.222329in}}%
\pgfpathlineto{\pgfqpoint{0.560732in}{1.218431in}}%
\pgfpathlineto{\pgfqpoint{0.544722in}{1.205183in}}%
\pgfpathlineto{\pgfqpoint{0.555755in}{1.191935in}}%
\pgfpathlineto{\pgfqpoint{0.540975in}{1.188094in}}%
\pgfpathlineto{\pgfqpoint{0.540975in}{1.178687in}}%
\pgfpathlineto{\pgfqpoint{0.540975in}{1.170021in}}%
\pgfpathlineto{\pgfqpoint{0.566258in}{1.167298in}}%
\pgfpathclose%
\pgfpathmoveto{\pgfqpoint{0.826355in}{1.456895in}}%
\pgfpathlineto{\pgfqpoint{0.844373in}{1.463051in}}%
\pgfpathlineto{\pgfqpoint{0.869656in}{1.458273in}}%
\pgfpathlineto{\pgfqpoint{0.870598in}{1.456895in}}%
\pgfpathlineto{\pgfqpoint{0.869656in}{1.456461in}}%
\pgfpathlineto{\pgfqpoint{0.844373in}{1.452879in}}%
\pgfpathlineto{\pgfqpoint{0.826355in}{1.456895in}}%
\pgfpathclose%
\pgfpathmoveto{\pgfqpoint{2.234947in}{1.203057in}}%
\pgfpathlineto{\pgfqpoint{2.260230in}{1.197008in}}%
\pgfpathlineto{\pgfqpoint{2.285514in}{1.192834in}}%
\pgfpathlineto{\pgfqpoint{2.310797in}{1.194379in}}%
\pgfpathlineto{\pgfqpoint{2.320757in}{1.205183in}}%
\pgfpathlineto{\pgfqpoint{2.310797in}{1.213344in}}%
\pgfpathlineto{\pgfqpoint{2.301973in}{1.218431in}}%
\pgfpathlineto{\pgfqpoint{2.310797in}{1.221588in}}%
\pgfpathlineto{\pgfqpoint{2.322778in}{1.231679in}}%
\pgfpathlineto{\pgfqpoint{2.310797in}{1.238369in}}%
\pgfpathlineto{\pgfqpoint{2.297316in}{1.244927in}}%
\pgfpathlineto{\pgfqpoint{2.308649in}{1.258175in}}%
\pgfpathlineto{\pgfqpoint{2.285514in}{1.269218in}}%
\pgfpathlineto{\pgfqpoint{2.280643in}{1.271423in}}%
\pgfpathlineto{\pgfqpoint{2.285372in}{1.284671in}}%
\pgfpathlineto{\pgfqpoint{2.270262in}{1.297919in}}%
\pgfpathlineto{\pgfqpoint{2.262097in}{1.311167in}}%
\pgfpathlineto{\pgfqpoint{2.260230in}{1.319119in}}%
\pgfpathlineto{\pgfqpoint{2.259068in}{1.324415in}}%
\pgfpathlineto{\pgfqpoint{2.238815in}{1.337663in}}%
\pgfpathlineto{\pgfqpoint{2.243114in}{1.350911in}}%
\pgfpathlineto{\pgfqpoint{2.234947in}{1.356847in}}%
\pgfpathlineto{\pgfqpoint{2.224555in}{1.364159in}}%
\pgfpathlineto{\pgfqpoint{2.224906in}{1.377407in}}%
\pgfpathlineto{\pgfqpoint{2.210488in}{1.390655in}}%
\pgfpathlineto{\pgfqpoint{2.209664in}{1.392395in}}%
\pgfpathlineto{\pgfqpoint{2.204091in}{1.403903in}}%
\pgfpathlineto{\pgfqpoint{2.198000in}{1.417151in}}%
\pgfpathlineto{\pgfqpoint{2.185608in}{1.430399in}}%
\pgfpathlineto{\pgfqpoint{2.184381in}{1.443445in}}%
\pgfpathlineto{\pgfqpoint{2.184364in}{1.443647in}}%
\pgfpathlineto{\pgfqpoint{2.166592in}{1.456895in}}%
\pgfpathlineto{\pgfqpoint{2.166404in}{1.470143in}}%
\pgfpathlineto{\pgfqpoint{2.159098in}{1.477334in}}%
\pgfpathlineto{\pgfqpoint{2.153161in}{1.483391in}}%
\pgfpathlineto{\pgfqpoint{2.148185in}{1.496639in}}%
\pgfpathlineto{\pgfqpoint{2.140883in}{1.509887in}}%
\pgfpathlineto{\pgfqpoint{2.133815in}{1.518589in}}%
\pgfpathlineto{\pgfqpoint{2.130134in}{1.523135in}}%
\pgfpathlineto{\pgfqpoint{2.128697in}{1.536383in}}%
\pgfpathlineto{\pgfqpoint{2.116858in}{1.549631in}}%
\pgfpathlineto{\pgfqpoint{2.115953in}{1.562879in}}%
\pgfpathlineto{\pgfqpoint{2.108531in}{1.570408in}}%
\pgfpathlineto{\pgfqpoint{2.103258in}{1.576127in}}%
\pgfpathlineto{\pgfqpoint{2.099428in}{1.589375in}}%
\pgfpathlineto{\pgfqpoint{2.091608in}{1.602623in}}%
\pgfpathlineto{\pgfqpoint{2.083260in}{1.615871in}}%
\pgfpathlineto{\pgfqpoint{2.083248in}{1.615911in}}%
\pgfpathlineto{\pgfqpoint{2.079866in}{1.629119in}}%
\pgfpathlineto{\pgfqpoint{2.068409in}{1.642367in}}%
\pgfpathlineto{\pgfqpoint{2.067702in}{1.655615in}}%
\pgfpathlineto{\pgfqpoint{2.057965in}{1.668078in}}%
\pgfpathlineto{\pgfqpoint{2.057401in}{1.668863in}}%
\pgfpathlineto{\pgfqpoint{2.054974in}{1.682111in}}%
\pgfpathlineto{\pgfqpoint{2.047821in}{1.695359in}}%
\pgfpathlineto{\pgfqpoint{2.041199in}{1.708607in}}%
\pgfpathlineto{\pgfqpoint{2.037398in}{1.721855in}}%
\pgfpathlineto{\pgfqpoint{2.032682in}{1.728630in}}%
\pgfpathlineto{\pgfqpoint{2.028638in}{1.735103in}}%
\pgfpathlineto{\pgfqpoint{2.026943in}{1.748351in}}%
\pgfpathlineto{\pgfqpoint{2.017104in}{1.761599in}}%
\pgfpathlineto{\pgfqpoint{2.014926in}{1.774847in}}%
\pgfpathlineto{\pgfqpoint{2.007399in}{1.787693in}}%
\pgfpathlineto{\pgfqpoint{2.007200in}{1.788095in}}%
\pgfpathlineto{\pgfqpoint{2.002766in}{1.801343in}}%
\pgfpathlineto{\pgfqpoint{1.999093in}{1.814591in}}%
\pgfpathlineto{\pgfqpoint{1.991083in}{1.827839in}}%
\pgfpathlineto{\pgfqpoint{1.989456in}{1.841087in}}%
\pgfpathlineto{\pgfqpoint{1.982115in}{1.852436in}}%
\pgfpathlineto{\pgfqpoint{1.981075in}{1.854335in}}%
\pgfpathlineto{\pgfqpoint{1.979455in}{1.867583in}}%
\pgfpathlineto{\pgfqpoint{1.973047in}{1.880831in}}%
\pgfpathlineto{\pgfqpoint{1.969120in}{1.894079in}}%
\pgfpathlineto{\pgfqpoint{1.964978in}{1.907327in}}%
\pgfpathlineto{\pgfqpoint{1.958047in}{1.920575in}}%
\pgfpathlineto{\pgfqpoint{1.956832in}{1.927040in}}%
\pgfpathlineto{\pgfqpoint{1.955901in}{1.933823in}}%
\pgfpathlineto{\pgfqpoint{1.948586in}{1.947071in}}%
\pgfpathlineto{\pgfqpoint{1.947010in}{1.960319in}}%
\pgfpathlineto{\pgfqpoint{1.940959in}{1.973567in}}%
\pgfpathlineto{\pgfqpoint{1.937856in}{1.986815in}}%
\pgfpathlineto{\pgfqpoint{1.933497in}{2.000063in}}%
\pgfpathlineto{\pgfqpoint{1.931549in}{2.003957in}}%
\pgfpathlineto{\pgfqpoint{1.928072in}{2.013311in}}%
\pgfpathlineto{\pgfqpoint{1.926466in}{2.026559in}}%
\pgfpathlineto{\pgfqpoint{1.920563in}{2.039807in}}%
\pgfpathlineto{\pgfqpoint{1.919137in}{2.053055in}}%
\pgfpathlineto{\pgfqpoint{1.912460in}{2.066303in}}%
\pgfpathlineto{\pgfqpoint{1.909056in}{2.079551in}}%
\pgfpathlineto{\pgfqpoint{1.906266in}{2.086741in}}%
\pgfpathlineto{\pgfqpoint{1.904646in}{2.092799in}}%
\pgfpathlineto{\pgfqpoint{1.900889in}{2.106047in}}%
\pgfpathlineto{\pgfqpoint{1.899271in}{2.119295in}}%
\pgfpathlineto{\pgfqpoint{1.893542in}{2.132543in}}%
\pgfpathlineto{\pgfqpoint{1.880983in}{2.132543in}}%
\pgfpathlineto{\pgfqpoint{1.867829in}{2.132543in}}%
\pgfpathlineto{\pgfqpoint{1.873531in}{2.119295in}}%
\pgfpathlineto{\pgfqpoint{1.874464in}{2.106047in}}%
\pgfpathlineto{\pgfqpoint{1.878747in}{2.092799in}}%
\pgfpathlineto{\pgfqpoint{1.880983in}{2.083179in}}%
\pgfpathlineto{\pgfqpoint{1.881973in}{2.079551in}}%
\pgfpathlineto{\pgfqpoint{1.885118in}{2.066303in}}%
\pgfpathlineto{\pgfqpoint{1.890949in}{2.053055in}}%
\pgfpathlineto{\pgfqpoint{1.891906in}{2.039807in}}%
\pgfpathlineto{\pgfqpoint{1.898197in}{2.026559in}}%
\pgfpathlineto{\pgfqpoint{1.898798in}{2.013311in}}%
\pgfpathlineto{\pgfqpoint{1.904392in}{2.000063in}}%
\pgfpathlineto{\pgfqpoint{1.906266in}{1.990326in}}%
\pgfpathlineto{\pgfqpoint{1.907032in}{1.986815in}}%
\pgfpathlineto{\pgfqpoint{1.910006in}{1.973567in}}%
\pgfpathlineto{\pgfqpoint{1.915431in}{1.960319in}}%
\pgfpathlineto{\pgfqpoint{1.916391in}{1.947071in}}%
\pgfpathlineto{\pgfqpoint{1.923991in}{1.933823in}}%
\pgfpathlineto{\pgfqpoint{1.924582in}{1.920575in}}%
\pgfpathlineto{\pgfqpoint{1.931186in}{1.907327in}}%
\pgfpathlineto{\pgfqpoint{1.931549in}{1.905267in}}%
\pgfpathlineto{\pgfqpoint{1.933677in}{1.894079in}}%
\pgfpathlineto{\pgfqpoint{1.937958in}{1.880831in}}%
\pgfpathlineto{\pgfqpoint{1.943543in}{1.867583in}}%
\pgfpathlineto{\pgfqpoint{1.944469in}{1.854335in}}%
\pgfpathlineto{\pgfqpoint{1.952099in}{1.841087in}}%
\pgfpathlineto{\pgfqpoint{1.951658in}{1.827839in}}%
\pgfpathlineto{\pgfqpoint{1.956832in}{1.820058in}}%
\pgfpathlineto{\pgfqpoint{1.960324in}{1.814591in}}%
\pgfpathlineto{\pgfqpoint{1.961758in}{1.801343in}}%
\pgfpathlineto{\pgfqpoint{1.966535in}{1.788095in}}%
\pgfpathlineto{\pgfqpoint{1.971667in}{1.774847in}}%
\pgfpathlineto{\pgfqpoint{1.972771in}{1.761599in}}%
\pgfpathlineto{\pgfqpoint{1.982115in}{1.749050in}}%
\pgfpathlineto{\pgfqpoint{1.982621in}{1.748351in}}%
\pgfpathlineto{\pgfqpoint{1.982115in}{1.738312in}}%
\pgfpathlineto{\pgfqpoint{1.981921in}{1.735103in}}%
\pgfpathlineto{\pgfqpoint{1.982115in}{1.734860in}}%
\pgfpathlineto{\pgfqpoint{1.991251in}{1.721855in}}%
\pgfpathlineto{\pgfqpoint{1.991159in}{1.708607in}}%
\pgfpathlineto{\pgfqpoint{1.998656in}{1.695359in}}%
\pgfpathlineto{\pgfqpoint{2.003116in}{1.682111in}}%
\pgfpathlineto{\pgfqpoint{2.004073in}{1.668863in}}%
\pgfpathlineto{\pgfqpoint{2.007399in}{1.664228in}}%
\pgfpathlineto{\pgfqpoint{2.012985in}{1.655615in}}%
\pgfpathlineto{\pgfqpoint{2.009039in}{1.642367in}}%
\pgfpathlineto{\pgfqpoint{2.021562in}{1.629119in}}%
\pgfpathlineto{\pgfqpoint{2.018482in}{1.615871in}}%
\pgfpathlineto{\pgfqpoint{2.028621in}{1.602623in}}%
\pgfpathlineto{\pgfqpoint{2.031088in}{1.589375in}}%
\pgfpathlineto{\pgfqpoint{2.032682in}{1.582861in}}%
\pgfpathlineto{\pgfqpoint{2.033923in}{1.576127in}}%
\pgfpathlineto{\pgfqpoint{2.043585in}{1.562879in}}%
\pgfpathlineto{\pgfqpoint{2.035956in}{1.549631in}}%
\pgfpathlineto{\pgfqpoint{2.050745in}{1.536383in}}%
\pgfpathlineto{\pgfqpoint{2.032682in}{1.523908in}}%
\pgfpathlineto{\pgfqpoint{2.026336in}{1.523135in}}%
\pgfpathlineto{\pgfqpoint{2.032682in}{1.521851in}}%
\pgfpathlineto{\pgfqpoint{2.049406in}{1.509887in}}%
\pgfpathlineto{\pgfqpoint{2.032682in}{1.499924in}}%
\pgfpathlineto{\pgfqpoint{2.007399in}{1.502682in}}%
\pgfpathlineto{\pgfqpoint{1.989723in}{1.509887in}}%
\pgfpathlineto{\pgfqpoint{2.000096in}{1.523135in}}%
\pgfpathlineto{\pgfqpoint{1.982115in}{1.526919in}}%
\pgfpathlineto{\pgfqpoint{1.961523in}{1.536383in}}%
\pgfpathlineto{\pgfqpoint{1.958570in}{1.549631in}}%
\pgfpathlineto{\pgfqpoint{1.956832in}{1.550494in}}%
\pgfpathlineto{\pgfqpoint{1.943161in}{1.562879in}}%
\pgfpathlineto{\pgfqpoint{1.931549in}{1.575713in}}%
\pgfpathlineto{\pgfqpoint{1.931259in}{1.576127in}}%
\pgfpathlineto{\pgfqpoint{1.927412in}{1.589375in}}%
\pgfpathlineto{\pgfqpoint{1.910250in}{1.602623in}}%
\pgfpathlineto{\pgfqpoint{1.909418in}{1.615871in}}%
\pgfpathlineto{\pgfqpoint{1.906266in}{1.617834in}}%
\pgfpathlineto{\pgfqpoint{1.893764in}{1.629119in}}%
\pgfpathlineto{\pgfqpoint{1.891019in}{1.642367in}}%
\pgfpathlineto{\pgfqpoint{1.880983in}{1.652789in}}%
\pgfpathlineto{\pgfqpoint{1.878954in}{1.655615in}}%
\pgfpathlineto{\pgfqpoint{1.870429in}{1.668863in}}%
\pgfpathlineto{\pgfqpoint{1.865410in}{1.682111in}}%
\pgfpathlineto{\pgfqpoint{1.855700in}{1.692192in}}%
\pgfpathlineto{\pgfqpoint{1.853171in}{1.695359in}}%
\pgfpathlineto{\pgfqpoint{1.852133in}{1.708607in}}%
\pgfpathlineto{\pgfqpoint{1.840506in}{1.721855in}}%
\pgfpathlineto{\pgfqpoint{1.837967in}{1.735103in}}%
\pgfpathlineto{\pgfqpoint{1.830416in}{1.746267in}}%
\pgfpathlineto{\pgfqpoint{1.829237in}{1.748351in}}%
\pgfpathlineto{\pgfqpoint{1.823327in}{1.761599in}}%
\pgfpathlineto{\pgfqpoint{1.819741in}{1.774847in}}%
\pgfpathlineto{\pgfqpoint{1.809266in}{1.788095in}}%
\pgfpathlineto{\pgfqpoint{1.806994in}{1.801343in}}%
\pgfpathlineto{\pgfqpoint{1.805133in}{1.803571in}}%
\pgfpathlineto{\pgfqpoint{1.797373in}{1.814591in}}%
\pgfpathlineto{\pgfqpoint{1.794910in}{1.827839in}}%
\pgfpathlineto{\pgfqpoint{1.788078in}{1.841087in}}%
\pgfpathlineto{\pgfqpoint{1.782593in}{1.854335in}}%
\pgfpathlineto{\pgfqpoint{1.779850in}{1.864221in}}%
\pgfpathlineto{\pgfqpoint{1.778977in}{1.867583in}}%
\pgfpathlineto{\pgfqpoint{1.771632in}{1.880831in}}%
\pgfpathlineto{\pgfqpoint{1.770464in}{1.894079in}}%
\pgfpathlineto{\pgfqpoint{1.762063in}{1.907327in}}%
\pgfpathlineto{\pgfqpoint{1.759586in}{1.920575in}}%
\pgfpathlineto{\pgfqpoint{1.754567in}{1.930756in}}%
\pgfpathlineto{\pgfqpoint{1.753110in}{1.933823in}}%
\pgfpathlineto{\pgfqpoint{1.748482in}{1.947071in}}%
\pgfpathlineto{\pgfqpoint{1.745407in}{1.960319in}}%
\pgfpathlineto{\pgfqpoint{1.738600in}{1.973567in}}%
\pgfpathlineto{\pgfqpoint{1.737095in}{1.986815in}}%
\pgfpathlineto{\pgfqpoint{1.730271in}{2.000063in}}%
\pgfpathlineto{\pgfqpoint{1.729284in}{2.007482in}}%
\pgfpathlineto{\pgfqpoint{1.728444in}{2.013311in}}%
\pgfpathlineto{\pgfqpoint{1.722869in}{2.026559in}}%
\pgfpathlineto{\pgfqpoint{1.718775in}{2.039807in}}%
\pgfpathlineto{\pgfqpoint{1.715645in}{2.053055in}}%
\pgfpathlineto{\pgfqpoint{1.710265in}{2.066303in}}%
\pgfpathlineto{\pgfqpoint{1.709015in}{2.079551in}}%
\pgfpathlineto{\pgfqpoint{1.704001in}{2.089999in}}%
\pgfpathlineto{\pgfqpoint{1.702545in}{2.092799in}}%
\pgfpathlineto{\pgfqpoint{1.700544in}{2.106047in}}%
\pgfpathlineto{\pgfqpoint{1.695260in}{2.119295in}}%
\pgfpathlineto{\pgfqpoint{1.691907in}{2.132543in}}%
\pgfpathlineto{\pgfqpoint{1.678717in}{2.132543in}}%
\pgfpathlineto{\pgfqpoint{1.668380in}{2.132543in}}%
\pgfpathlineto{\pgfqpoint{1.671863in}{2.119295in}}%
\pgfpathlineto{\pgfqpoint{1.677932in}{2.106047in}}%
\pgfpathlineto{\pgfqpoint{1.678717in}{2.102507in}}%
\pgfpathlineto{\pgfqpoint{1.679901in}{2.092799in}}%
\pgfpathlineto{\pgfqpoint{1.684755in}{2.079551in}}%
\pgfpathlineto{\pgfqpoint{1.685859in}{2.066303in}}%
\pgfpathlineto{\pgfqpoint{1.690427in}{2.053055in}}%
\pgfpathlineto{\pgfqpoint{1.693241in}{2.039807in}}%
\pgfpathlineto{\pgfqpoint{1.697734in}{2.026559in}}%
\pgfpathlineto{\pgfqpoint{1.703466in}{2.013311in}}%
\pgfpathlineto{\pgfqpoint{1.704001in}{2.010557in}}%
\pgfpathlineto{\pgfqpoint{1.705251in}{2.000063in}}%
\pgfpathlineto{\pgfqpoint{1.710450in}{1.986815in}}%
\pgfpathlineto{\pgfqpoint{1.711780in}{1.973567in}}%
\pgfpathlineto{\pgfqpoint{1.717781in}{1.960319in}}%
\pgfpathlineto{\pgfqpoint{1.720440in}{1.947071in}}%
\pgfpathlineto{\pgfqpoint{1.725192in}{1.933823in}}%
\pgfpathlineto{\pgfqpoint{1.729284in}{1.925412in}}%
\pgfpathlineto{\pgfqpoint{1.731021in}{1.920575in}}%
\pgfpathlineto{\pgfqpoint{1.733505in}{1.907327in}}%
\pgfpathlineto{\pgfqpoint{1.740079in}{1.894079in}}%
\pgfpathlineto{\pgfqpoint{1.741363in}{1.880831in}}%
\pgfpathlineto{\pgfqpoint{1.747900in}{1.867583in}}%
\pgfpathlineto{\pgfqpoint{1.751091in}{1.854335in}}%
\pgfpathlineto{\pgfqpoint{1.754567in}{1.847063in}}%
\pgfpathlineto{\pgfqpoint{1.756778in}{1.841087in}}%
\pgfpathlineto{\pgfqpoint{1.761603in}{1.827839in}}%
\pgfpathlineto{\pgfqpoint{1.764279in}{1.814591in}}%
\pgfpathlineto{\pgfqpoint{1.771847in}{1.801343in}}%
\pgfpathlineto{\pgfqpoint{1.774209in}{1.788095in}}%
\pgfpathlineto{\pgfqpoint{1.779850in}{1.780391in}}%
\pgfpathlineto{\pgfqpoint{1.783405in}{1.774847in}}%
\pgfpathlineto{\pgfqpoint{1.786084in}{1.761599in}}%
\pgfpathlineto{\pgfqpoint{1.791534in}{1.748351in}}%
\pgfpathlineto{\pgfqpoint{1.797644in}{1.735103in}}%
\pgfpathlineto{\pgfqpoint{1.801216in}{1.721855in}}%
\pgfpathlineto{\pgfqpoint{1.805133in}{1.716306in}}%
\pgfpathlineto{\pgfqpoint{1.810123in}{1.708607in}}%
\pgfpathlineto{\pgfqpoint{1.811959in}{1.695359in}}%
\pgfpathlineto{\pgfqpoint{1.820734in}{1.682111in}}%
\pgfpathlineto{\pgfqpoint{1.824435in}{1.668863in}}%
\pgfpathlineto{\pgfqpoint{1.830416in}{1.659774in}}%
\pgfpathlineto{\pgfqpoint{1.832964in}{1.655615in}}%
\pgfpathlineto{\pgfqpoint{1.839612in}{1.642367in}}%
\pgfpathlineto{\pgfqpoint{1.843641in}{1.629119in}}%
\pgfpathlineto{\pgfqpoint{1.854071in}{1.615871in}}%
\pgfpathlineto{\pgfqpoint{1.855700in}{1.607667in}}%
\pgfpathlineto{\pgfqpoint{1.856616in}{1.602623in}}%
\pgfpathlineto{\pgfqpoint{1.868157in}{1.589375in}}%
\pgfpathlineto{\pgfqpoint{1.871515in}{1.576127in}}%
\pgfpathlineto{\pgfqpoint{1.880511in}{1.562879in}}%
\pgfpathlineto{\pgfqpoint{1.880983in}{1.562077in}}%
\pgfpathlineto{\pgfqpoint{1.888290in}{1.549631in}}%
\pgfpathlineto{\pgfqpoint{1.893798in}{1.536383in}}%
\pgfpathlineto{\pgfqpoint{1.906266in}{1.523435in}}%
\pgfpathlineto{\pgfqpoint{1.906589in}{1.523135in}}%
\pgfpathlineto{\pgfqpoint{1.909082in}{1.509887in}}%
\pgfpathlineto{\pgfqpoint{1.922571in}{1.496639in}}%
\pgfpathlineto{\pgfqpoint{1.927183in}{1.483391in}}%
\pgfpathlineto{\pgfqpoint{1.931549in}{1.478726in}}%
\pgfpathlineto{\pgfqpoint{1.939873in}{1.470143in}}%
\pgfpathlineto{\pgfqpoint{1.949225in}{1.456895in}}%
\pgfpathlineto{\pgfqpoint{1.956162in}{1.443647in}}%
\pgfpathlineto{\pgfqpoint{1.956832in}{1.443006in}}%
\pgfpathlineto{\pgfqpoint{1.971431in}{1.430399in}}%
\pgfpathlineto{\pgfqpoint{1.974645in}{1.417151in}}%
\pgfpathlineto{\pgfqpoint{1.982115in}{1.411560in}}%
\pgfpathlineto{\pgfqpoint{1.993705in}{1.403903in}}%
\pgfpathlineto{\pgfqpoint{1.997122in}{1.390655in}}%
\pgfpathlineto{\pgfqpoint{2.007399in}{1.380967in}}%
\pgfpathlineto{\pgfqpoint{2.011594in}{1.377407in}}%
\pgfpathlineto{\pgfqpoint{2.023763in}{1.364159in}}%
\pgfpathlineto{\pgfqpoint{2.032682in}{1.351411in}}%
\pgfpathlineto{\pgfqpoint{2.033058in}{1.350911in}}%
\pgfpathlineto{\pgfqpoint{2.055460in}{1.337663in}}%
\pgfpathlineto{\pgfqpoint{2.057965in}{1.328587in}}%
\pgfpathlineto{\pgfqpoint{2.059218in}{1.324415in}}%
\pgfpathlineto{\pgfqpoint{2.083248in}{1.311882in}}%
\pgfpathlineto{\pgfqpoint{2.084942in}{1.311167in}}%
\pgfpathlineto{\pgfqpoint{2.091488in}{1.297919in}}%
\pgfpathlineto{\pgfqpoint{2.108531in}{1.287913in}}%
\pgfpathlineto{\pgfqpoint{2.114786in}{1.284671in}}%
\pgfpathlineto{\pgfqpoint{2.130159in}{1.271423in}}%
\pgfpathlineto{\pgfqpoint{2.133815in}{1.266122in}}%
\pgfpathlineto{\pgfqpoint{2.139835in}{1.258175in}}%
\pgfpathlineto{\pgfqpoint{2.159098in}{1.249387in}}%
\pgfpathlineto{\pgfqpoint{2.172106in}{1.244927in}}%
\pgfpathlineto{\pgfqpoint{2.173760in}{1.231679in}}%
\pgfpathlineto{\pgfqpoint{2.184381in}{1.228224in}}%
\pgfpathlineto{\pgfqpoint{2.209664in}{1.221881in}}%
\pgfpathlineto{\pgfqpoint{2.228703in}{1.218431in}}%
\pgfpathlineto{\pgfqpoint{2.228501in}{1.205183in}}%
\pgfpathlineto{\pgfqpoint{2.234947in}{1.203057in}}%
\pgfpathclose%
\pgfpathmoveto{\pgfqpoint{0.566258in}{1.635875in}}%
\pgfpathlineto{\pgfqpoint{0.575922in}{1.642367in}}%
\pgfpathlineto{\pgfqpoint{0.566258in}{1.648286in}}%
\pgfpathlineto{\pgfqpoint{0.540975in}{1.645215in}}%
\pgfpathlineto{\pgfqpoint{0.540975in}{1.642367in}}%
\pgfpathlineto{\pgfqpoint{0.540975in}{1.635700in}}%
\pgfpathlineto{\pgfqpoint{0.566258in}{1.635875in}}%
\pgfpathclose%
\pgfpathmoveto{\pgfqpoint{0.566258in}{1.659851in}}%
\pgfpathlineto{\pgfqpoint{0.591541in}{1.661379in}}%
\pgfpathlineto{\pgfqpoint{0.615193in}{1.668863in}}%
\pgfpathlineto{\pgfqpoint{0.594304in}{1.682111in}}%
\pgfpathlineto{\pgfqpoint{0.616824in}{1.686111in}}%
\pgfpathlineto{\pgfqpoint{0.642107in}{1.695301in}}%
\pgfpathlineto{\pgfqpoint{0.642213in}{1.695359in}}%
\pgfpathlineto{\pgfqpoint{0.642107in}{1.696069in}}%
\pgfpathlineto{\pgfqpoint{0.638070in}{1.708607in}}%
\pgfpathlineto{\pgfqpoint{0.642107in}{1.710509in}}%
\pgfpathlineto{\pgfqpoint{0.657843in}{1.721855in}}%
\pgfpathlineto{\pgfqpoint{0.667391in}{1.730426in}}%
\pgfpathlineto{\pgfqpoint{0.670611in}{1.735103in}}%
\pgfpathlineto{\pgfqpoint{0.676045in}{1.748351in}}%
\pgfpathlineto{\pgfqpoint{0.692674in}{1.758799in}}%
\pgfpathlineto{\pgfqpoint{0.695652in}{1.761599in}}%
\pgfpathlineto{\pgfqpoint{0.693756in}{1.774847in}}%
\pgfpathlineto{\pgfqpoint{0.717481in}{1.788095in}}%
\pgfpathlineto{\pgfqpoint{0.717957in}{1.793979in}}%
\pgfpathlineto{\pgfqpoint{0.718416in}{1.801343in}}%
\pgfpathlineto{\pgfqpoint{0.735500in}{1.814591in}}%
\pgfpathlineto{\pgfqpoint{0.740714in}{1.827839in}}%
\pgfpathlineto{\pgfqpoint{0.743240in}{1.831675in}}%
\pgfpathlineto{\pgfqpoint{0.748453in}{1.841087in}}%
\pgfpathlineto{\pgfqpoint{0.761066in}{1.854335in}}%
\pgfpathlineto{\pgfqpoint{0.762769in}{1.867583in}}%
\pgfpathlineto{\pgfqpoint{0.768523in}{1.871454in}}%
\pgfpathlineto{\pgfqpoint{0.778633in}{1.880831in}}%
\pgfpathlineto{\pgfqpoint{0.777990in}{1.894079in}}%
\pgfpathlineto{\pgfqpoint{0.793791in}{1.907327in}}%
\pgfpathlineto{\pgfqpoint{0.793806in}{1.907363in}}%
\pgfpathlineto{\pgfqpoint{0.798185in}{1.920575in}}%
\pgfpathlineto{\pgfqpoint{0.807991in}{1.933823in}}%
\pgfpathlineto{\pgfqpoint{0.816257in}{1.947071in}}%
\pgfpathlineto{\pgfqpoint{0.819090in}{1.957709in}}%
\pgfpathlineto{\pgfqpoint{0.819745in}{1.960319in}}%
\pgfpathlineto{\pgfqpoint{0.831719in}{1.973567in}}%
\pgfpathlineto{\pgfqpoint{0.832402in}{1.986815in}}%
\pgfpathlineto{\pgfqpoint{0.844373in}{1.998618in}}%
\pgfpathlineto{\pgfqpoint{0.845673in}{2.000063in}}%
\pgfpathlineto{\pgfqpoint{0.847093in}{2.013311in}}%
\pgfpathlineto{\pgfqpoint{0.855714in}{2.026559in}}%
\pgfpathlineto{\pgfqpoint{0.862411in}{2.039807in}}%
\pgfpathlineto{\pgfqpoint{0.867746in}{2.053055in}}%
\pgfpathlineto{\pgfqpoint{0.869656in}{2.055365in}}%
\pgfpathlineto{\pgfqpoint{0.877094in}{2.066303in}}%
\pgfpathlineto{\pgfqpoint{0.878146in}{2.079551in}}%
\pgfpathlineto{\pgfqpoint{0.888639in}{2.092799in}}%
\pgfpathlineto{\pgfqpoint{0.890420in}{2.106047in}}%
\pgfpathlineto{\pgfqpoint{0.894939in}{2.112069in}}%
\pgfpathlineto{\pgfqpoint{0.900114in}{2.119295in}}%
\pgfpathlineto{\pgfqpoint{0.904060in}{2.132543in}}%
\pgfpathlineto{\pgfqpoint{0.894939in}{2.132543in}}%
\pgfpathlineto{\pgfqpoint{0.869656in}{2.132543in}}%
\pgfpathlineto{\pgfqpoint{0.844373in}{2.132543in}}%
\pgfpathlineto{\pgfqpoint{0.819090in}{2.132543in}}%
\pgfpathlineto{\pgfqpoint{0.793806in}{2.132543in}}%
\pgfpathlineto{\pgfqpoint{0.768523in}{2.132543in}}%
\pgfpathlineto{\pgfqpoint{0.765275in}{2.132543in}}%
\pgfpathlineto{\pgfqpoint{0.764351in}{2.119295in}}%
\pgfpathlineto{\pgfqpoint{0.755277in}{2.106047in}}%
\pgfpathlineto{\pgfqpoint{0.752087in}{2.092799in}}%
\pgfpathlineto{\pgfqpoint{0.747472in}{2.079551in}}%
\pgfpathlineto{\pgfqpoint{0.743240in}{2.073501in}}%
\pgfpathlineto{\pgfqpoint{0.738586in}{2.066303in}}%
\pgfpathlineto{\pgfqpoint{0.738250in}{2.053055in}}%
\pgfpathlineto{\pgfqpoint{0.728018in}{2.039807in}}%
\pgfpathlineto{\pgfqpoint{0.728729in}{2.026559in}}%
\pgfpathlineto{\pgfqpoint{0.717957in}{2.013755in}}%
\pgfpathlineto{\pgfqpoint{0.717583in}{2.013311in}}%
\pgfpathlineto{\pgfqpoint{0.713956in}{2.000063in}}%
\pgfpathlineto{\pgfqpoint{0.707530in}{1.986815in}}%
\pgfpathlineto{\pgfqpoint{0.699333in}{1.973567in}}%
\pgfpathlineto{\pgfqpoint{0.698998in}{1.960319in}}%
\pgfpathlineto{\pgfqpoint{0.692674in}{1.954793in}}%
\pgfpathlineto{\pgfqpoint{0.684971in}{1.947071in}}%
\pgfpathlineto{\pgfqpoint{0.686600in}{1.933823in}}%
\pgfpathlineto{\pgfqpoint{0.674523in}{1.920575in}}%
\pgfpathlineto{\pgfqpoint{0.673401in}{1.907327in}}%
\pgfpathlineto{\pgfqpoint{0.667391in}{1.896812in}}%
\pgfpathlineto{\pgfqpoint{0.665842in}{1.894079in}}%
\pgfpathlineto{\pgfqpoint{0.656074in}{1.880831in}}%
\pgfpathlineto{\pgfqpoint{0.655762in}{1.867583in}}%
\pgfpathlineto{\pgfqpoint{0.642107in}{1.856245in}}%
\pgfpathlineto{\pgfqpoint{0.640197in}{1.854335in}}%
\pgfpathlineto{\pgfqpoint{0.642107in}{1.846798in}}%
\pgfpathlineto{\pgfqpoint{0.644230in}{1.841087in}}%
\pgfpathlineto{\pgfqpoint{0.642107in}{1.839573in}}%
\pgfpathlineto{\pgfqpoint{0.628401in}{1.827839in}}%
\pgfpathlineto{\pgfqpoint{0.627902in}{1.814591in}}%
\pgfpathlineto{\pgfqpoint{0.618478in}{1.801343in}}%
\pgfpathlineto{\pgfqpoint{0.616824in}{1.799270in}}%
\pgfpathlineto{\pgfqpoint{0.610327in}{1.788095in}}%
\pgfpathlineto{\pgfqpoint{0.612282in}{1.774847in}}%
\pgfpathlineto{\pgfqpoint{0.591541in}{1.762282in}}%
\pgfpathlineto{\pgfqpoint{0.590654in}{1.761599in}}%
\pgfpathlineto{\pgfqpoint{0.591541in}{1.759383in}}%
\pgfpathlineto{\pgfqpoint{0.599413in}{1.748351in}}%
\pgfpathlineto{\pgfqpoint{0.591541in}{1.744231in}}%
\pgfpathlineto{\pgfqpoint{0.578273in}{1.735103in}}%
\pgfpathlineto{\pgfqpoint{0.583969in}{1.721855in}}%
\pgfpathlineto{\pgfqpoint{0.575568in}{1.708607in}}%
\pgfpathlineto{\pgfqpoint{0.566258in}{1.700893in}}%
\pgfpathlineto{\pgfqpoint{0.562179in}{1.695359in}}%
\pgfpathlineto{\pgfqpoint{0.566258in}{1.688865in}}%
\pgfpathlineto{\pgfqpoint{0.587213in}{1.682111in}}%
\pgfpathlineto{\pgfqpoint{0.566258in}{1.678671in}}%
\pgfpathlineto{\pgfqpoint{0.544615in}{1.668863in}}%
\pgfpathlineto{\pgfqpoint{0.566258in}{1.659851in}}%
\pgfpathclose%
\pgfpathmoveto{\pgfqpoint{2.285514in}{1.694723in}}%
\pgfpathlineto{\pgfqpoint{2.310797in}{1.691110in}}%
\pgfpathlineto{\pgfqpoint{2.316086in}{1.695359in}}%
\pgfpathlineto{\pgfqpoint{2.310797in}{1.698535in}}%
\pgfpathlineto{\pgfqpoint{2.285514in}{1.696742in}}%
\pgfpathlineto{\pgfqpoint{2.284420in}{1.695359in}}%
\pgfpathlineto{\pgfqpoint{2.285514in}{1.694723in}}%
\pgfpathclose%
\pgfpathmoveto{\pgfqpoint{2.260230in}{1.719119in}}%
\pgfpathlineto{\pgfqpoint{2.285514in}{1.713124in}}%
\pgfpathlineto{\pgfqpoint{2.310797in}{1.718389in}}%
\pgfpathlineto{\pgfqpoint{2.314082in}{1.721855in}}%
\pgfpathlineto{\pgfqpoint{2.310797in}{1.724238in}}%
\pgfpathlineto{\pgfqpoint{2.286110in}{1.735103in}}%
\pgfpathlineto{\pgfqpoint{2.299775in}{1.748351in}}%
\pgfpathlineto{\pgfqpoint{2.285514in}{1.759372in}}%
\pgfpathlineto{\pgfqpoint{2.280714in}{1.761599in}}%
\pgfpathlineto{\pgfqpoint{2.282230in}{1.774847in}}%
\pgfpathlineto{\pgfqpoint{2.277762in}{1.788095in}}%
\pgfpathlineto{\pgfqpoint{2.267862in}{1.801343in}}%
\pgfpathlineto{\pgfqpoint{2.270944in}{1.814591in}}%
\pgfpathlineto{\pgfqpoint{2.260230in}{1.823734in}}%
\pgfpathlineto{\pgfqpoint{2.253459in}{1.827839in}}%
\pgfpathlineto{\pgfqpoint{2.259683in}{1.841087in}}%
\pgfpathlineto{\pgfqpoint{2.245948in}{1.854335in}}%
\pgfpathlineto{\pgfqpoint{2.245388in}{1.867583in}}%
\pgfpathlineto{\pgfqpoint{2.237871in}{1.880831in}}%
\pgfpathlineto{\pgfqpoint{2.234947in}{1.885945in}}%
\pgfpathlineto{\pgfqpoint{2.228967in}{1.894079in}}%
\pgfpathlineto{\pgfqpoint{2.227710in}{1.907327in}}%
\pgfpathlineto{\pgfqpoint{2.216067in}{1.920575in}}%
\pgfpathlineto{\pgfqpoint{2.216829in}{1.933823in}}%
\pgfpathlineto{\pgfqpoint{2.209664in}{1.941337in}}%
\pgfpathlineto{\pgfqpoint{2.203370in}{1.947071in}}%
\pgfpathlineto{\pgfqpoint{2.202517in}{1.960319in}}%
\pgfpathlineto{\pgfqpoint{2.194875in}{1.973567in}}%
\pgfpathlineto{\pgfqpoint{2.190274in}{1.986815in}}%
\pgfpathlineto{\pgfqpoint{2.187554in}{2.000063in}}%
\pgfpathlineto{\pgfqpoint{2.184381in}{2.004413in}}%
\pgfpathlineto{\pgfqpoint{2.177048in}{2.013311in}}%
\pgfpathlineto{\pgfqpoint{2.177449in}{2.026559in}}%
\pgfpathlineto{\pgfqpoint{2.167258in}{2.039807in}}%
\pgfpathlineto{\pgfqpoint{2.166170in}{2.053055in}}%
\pgfpathlineto{\pgfqpoint{2.159098in}{2.063619in}}%
\pgfpathlineto{\pgfqpoint{2.157215in}{2.066303in}}%
\pgfpathlineto{\pgfqpoint{2.151981in}{2.079551in}}%
\pgfpathlineto{\pgfqpoint{2.147849in}{2.092799in}}%
\pgfpathlineto{\pgfqpoint{2.140685in}{2.106047in}}%
\pgfpathlineto{\pgfqpoint{2.139886in}{2.119295in}}%
\pgfpathlineto{\pgfqpoint{2.133815in}{2.127619in}}%
\pgfpathlineto{\pgfqpoint{2.130056in}{2.132543in}}%
\pgfpathlineto{\pgfqpoint{2.108531in}{2.132543in}}%
\pgfpathlineto{\pgfqpoint{2.083248in}{2.132543in}}%
\pgfpathlineto{\pgfqpoint{2.057965in}{2.132543in}}%
\pgfpathlineto{\pgfqpoint{2.032682in}{2.132543in}}%
\pgfpathlineto{\pgfqpoint{2.007399in}{2.132543in}}%
\pgfpathlineto{\pgfqpoint{2.003774in}{2.132543in}}%
\pgfpathlineto{\pgfqpoint{2.007399in}{2.120728in}}%
\pgfpathlineto{\pgfqpoint{2.007863in}{2.119295in}}%
\pgfpathlineto{\pgfqpoint{2.017666in}{2.106047in}}%
\pgfpathlineto{\pgfqpoint{2.019515in}{2.092799in}}%
\pgfpathlineto{\pgfqpoint{2.030174in}{2.079551in}}%
\pgfpathlineto{\pgfqpoint{2.031998in}{2.066303in}}%
\pgfpathlineto{\pgfqpoint{2.032682in}{2.065114in}}%
\pgfpathlineto{\pgfqpoint{2.040166in}{2.053055in}}%
\pgfpathlineto{\pgfqpoint{2.046689in}{2.039807in}}%
\pgfpathlineto{\pgfqpoint{2.051135in}{2.026559in}}%
\pgfpathlineto{\pgfqpoint{2.057965in}{2.017542in}}%
\pgfpathlineto{\pgfqpoint{2.062029in}{2.013311in}}%
\pgfpathlineto{\pgfqpoint{2.063287in}{2.000063in}}%
\pgfpathlineto{\pgfqpoint{2.075306in}{1.986815in}}%
\pgfpathlineto{\pgfqpoint{2.078006in}{1.973567in}}%
\pgfpathlineto{\pgfqpoint{2.083248in}{1.966865in}}%
\pgfpathlineto{\pgfqpoint{2.089367in}{1.960319in}}%
\pgfpathlineto{\pgfqpoint{2.096165in}{1.947071in}}%
\pgfpathlineto{\pgfqpoint{2.101135in}{1.933823in}}%
\pgfpathlineto{\pgfqpoint{2.108531in}{1.925347in}}%
\pgfpathlineto{\pgfqpoint{2.114191in}{1.920575in}}%
\pgfpathlineto{\pgfqpoint{2.116061in}{1.907327in}}%
\pgfpathlineto{\pgfqpoint{2.131909in}{1.894079in}}%
\pgfpathlineto{\pgfqpoint{2.132469in}{1.880831in}}%
\pgfpathlineto{\pgfqpoint{2.133815in}{1.879145in}}%
\pgfpathlineto{\pgfqpoint{2.145257in}{1.867583in}}%
\pgfpathlineto{\pgfqpoint{2.152932in}{1.854335in}}%
\pgfpathlineto{\pgfqpoint{2.158803in}{1.841087in}}%
\pgfpathlineto{\pgfqpoint{2.159098in}{1.840784in}}%
\pgfpathlineto{\pgfqpoint{2.176837in}{1.827839in}}%
\pgfpathlineto{\pgfqpoint{2.176612in}{1.814591in}}%
\pgfpathlineto{\pgfqpoint{2.184381in}{1.808585in}}%
\pgfpathlineto{\pgfqpoint{2.198333in}{1.801343in}}%
\pgfpathlineto{\pgfqpoint{2.200085in}{1.788095in}}%
\pgfpathlineto{\pgfqpoint{2.209664in}{1.780903in}}%
\pgfpathlineto{\pgfqpoint{2.221595in}{1.774847in}}%
\pgfpathlineto{\pgfqpoint{2.230605in}{1.761599in}}%
\pgfpathlineto{\pgfqpoint{2.234263in}{1.748351in}}%
\pgfpathlineto{\pgfqpoint{2.234947in}{1.747798in}}%
\pgfpathlineto{\pgfqpoint{2.260230in}{1.735913in}}%
\pgfpathlineto{\pgfqpoint{2.278974in}{1.735103in}}%
\pgfpathlineto{\pgfqpoint{2.260230in}{1.732965in}}%
\pgfpathlineto{\pgfqpoint{2.254824in}{1.721855in}}%
\pgfpathlineto{\pgfqpoint{2.260230in}{1.719119in}}%
\pgfpathclose%
\pgfpathmoveto{\pgfqpoint{0.566258in}{2.129593in}}%
\pgfpathlineto{\pgfqpoint{0.568649in}{2.132543in}}%
\pgfpathlineto{\pgfqpoint{0.566258in}{2.132543in}}%
\pgfpathlineto{\pgfqpoint{0.562233in}{2.132543in}}%
\pgfpathlineto{\pgfqpoint{0.566258in}{2.129593in}}%
\pgfpathclose%
\pgfusepath{fill}%
\end{pgfscope}%
\begin{pgfscope}%
\pgfpathrectangle{\pgfqpoint{0.540975in}{0.476543in}}{\pgfqpoint{1.820388in}{1.656000in}}%
\pgfusepath{clip}%
\pgfsetbuttcap%
\pgfsetroundjoin%
\definecolor{currentfill}{rgb}{0.274202,0.463130,0.708112}%
\pgfsetfillcolor{currentfill}%
\pgfsetlinewidth{0.000000pt}%
\definecolor{currentstroke}{rgb}{0.000000,0.000000,0.000000}%
\pgfsetstrokecolor{currentstroke}%
\pgfsetdash{}{0pt}%
\pgfpathmoveto{\pgfqpoint{1.324753in}{0.688392in}}%
\pgfpathlineto{\pgfqpoint{1.350036in}{0.688279in}}%
\pgfpathlineto{\pgfqpoint{1.375319in}{0.688254in}}%
\pgfpathlineto{\pgfqpoint{1.400603in}{0.688320in}}%
\pgfpathlineto{\pgfqpoint{1.425886in}{0.688485in}}%
\pgfpathlineto{\pgfqpoint{1.428463in}{0.688511in}}%
\pgfpathlineto{\pgfqpoint{1.451169in}{0.688987in}}%
\pgfpathlineto{\pgfqpoint{1.476452in}{0.689667in}}%
\pgfpathlineto{\pgfqpoint{1.501735in}{0.690484in}}%
\pgfpathlineto{\pgfqpoint{1.527018in}{0.691434in}}%
\pgfpathlineto{\pgfqpoint{1.552302in}{0.692510in}}%
\pgfpathlineto{\pgfqpoint{1.577585in}{0.693713in}}%
\pgfpathlineto{\pgfqpoint{1.602868in}{0.695043in}}%
\pgfpathlineto{\pgfqpoint{1.628151in}{0.696501in}}%
\pgfpathlineto{\pgfqpoint{1.653434in}{0.698094in}}%
\pgfpathlineto{\pgfqpoint{1.678717in}{0.699826in}}%
\pgfpathlineto{\pgfqpoint{1.704001in}{0.701706in}}%
\pgfpathlineto{\pgfqpoint{1.704664in}{0.701759in}}%
\pgfpathlineto{\pgfqpoint{1.729284in}{0.702841in}}%
\pgfpathlineto{\pgfqpoint{1.754567in}{0.704133in}}%
\pgfpathlineto{\pgfqpoint{1.779850in}{0.705634in}}%
\pgfpathlineto{\pgfqpoint{1.805133in}{0.707378in}}%
\pgfpathlineto{\pgfqpoint{1.830416in}{0.709405in}}%
\pgfpathlineto{\pgfqpoint{1.855700in}{0.711767in}}%
\pgfpathlineto{\pgfqpoint{1.880983in}{0.714527in}}%
\pgfpathlineto{\pgfqpoint{1.884945in}{0.715007in}}%
\pgfpathlineto{\pgfqpoint{1.906266in}{0.728107in}}%
\pgfpathlineto{\pgfqpoint{1.906484in}{0.728255in}}%
\pgfpathlineto{\pgfqpoint{1.931549in}{0.730722in}}%
\pgfpathlineto{\pgfqpoint{1.956832in}{0.733580in}}%
\pgfpathlineto{\pgfqpoint{1.982115in}{0.736883in}}%
\pgfpathlineto{\pgfqpoint{2.007399in}{0.740723in}}%
\pgfpathlineto{\pgfqpoint{2.012066in}{0.741503in}}%
\pgfpathlineto{\pgfqpoint{2.007399in}{0.748853in}}%
\pgfpathlineto{\pgfqpoint{2.003954in}{0.754751in}}%
\pgfpathlineto{\pgfqpoint{2.007399in}{0.755447in}}%
\pgfpathlineto{\pgfqpoint{2.032682in}{0.761037in}}%
\pgfpathlineto{\pgfqpoint{2.057965in}{0.767588in}}%
\pgfpathlineto{\pgfqpoint{2.059421in}{0.767999in}}%
\pgfpathlineto{\pgfqpoint{2.057965in}{0.768921in}}%
\pgfpathlineto{\pgfqpoint{2.040728in}{0.781247in}}%
\pgfpathlineto{\pgfqpoint{2.057965in}{0.789896in}}%
\pgfpathlineto{\pgfqpoint{2.066394in}{0.794495in}}%
\pgfpathlineto{\pgfqpoint{2.059983in}{0.807743in}}%
\pgfpathlineto{\pgfqpoint{2.058037in}{0.820991in}}%
\pgfpathlineto{\pgfqpoint{2.063686in}{0.834239in}}%
\pgfpathlineto{\pgfqpoint{2.057965in}{0.837946in}}%
\pgfpathlineto{\pgfqpoint{2.043893in}{0.847487in}}%
\pgfpathlineto{\pgfqpoint{2.057631in}{0.860735in}}%
\pgfpathlineto{\pgfqpoint{2.037170in}{0.873983in}}%
\pgfpathlineto{\pgfqpoint{2.044489in}{0.887231in}}%
\pgfpathlineto{\pgfqpoint{2.032682in}{0.895867in}}%
\pgfpathlineto{\pgfqpoint{2.027024in}{0.900479in}}%
\pgfpathlineto{\pgfqpoint{2.024544in}{0.913727in}}%
\pgfpathlineto{\pgfqpoint{2.017691in}{0.926975in}}%
\pgfpathlineto{\pgfqpoint{2.007399in}{0.938244in}}%
\pgfpathlineto{\pgfqpoint{2.005689in}{0.940223in}}%
\pgfpathlineto{\pgfqpoint{2.005941in}{0.953471in}}%
\pgfpathlineto{\pgfqpoint{1.986747in}{0.966719in}}%
\pgfpathlineto{\pgfqpoint{1.989484in}{0.979967in}}%
\pgfpathlineto{\pgfqpoint{1.982115in}{0.985642in}}%
\pgfpathlineto{\pgfqpoint{1.973111in}{0.993215in}}%
\pgfpathlineto{\pgfqpoint{1.970154in}{1.006463in}}%
\pgfpathlineto{\pgfqpoint{1.959135in}{1.019711in}}%
\pgfpathlineto{\pgfqpoint{1.956832in}{1.022681in}}%
\pgfpathlineto{\pgfqpoint{1.949486in}{1.032959in}}%
\pgfpathlineto{\pgfqpoint{1.946307in}{1.046207in}}%
\pgfpathlineto{\pgfqpoint{1.933163in}{1.059455in}}%
\pgfpathlineto{\pgfqpoint{1.933289in}{1.072703in}}%
\pgfpathlineto{\pgfqpoint{1.931549in}{1.074191in}}%
\pgfpathlineto{\pgfqpoint{1.919184in}{1.085951in}}%
\pgfpathlineto{\pgfqpoint{1.917790in}{1.099199in}}%
\pgfpathlineto{\pgfqpoint{1.907358in}{1.112447in}}%
\pgfpathlineto{\pgfqpoint{1.906266in}{1.114459in}}%
\pgfpathlineto{\pgfqpoint{1.900869in}{1.125695in}}%
\pgfpathlineto{\pgfqpoint{1.895652in}{1.138943in}}%
\pgfpathlineto{\pgfqpoint{1.883539in}{1.152191in}}%
\pgfpathlineto{\pgfqpoint{1.882040in}{1.165439in}}%
\pgfpathlineto{\pgfqpoint{1.880983in}{1.166560in}}%
\pgfpathlineto{\pgfqpoint{1.870896in}{1.178687in}}%
\pgfpathlineto{\pgfqpoint{1.869486in}{1.191935in}}%
\pgfpathlineto{\pgfqpoint{1.859800in}{1.205183in}}%
\pgfpathlineto{\pgfqpoint{1.855700in}{1.215689in}}%
\pgfpathlineto{\pgfqpoint{1.854790in}{1.218431in}}%
\pgfpathlineto{\pgfqpoint{1.850186in}{1.231679in}}%
\pgfpathlineto{\pgfqpoint{1.842181in}{1.244927in}}%
\pgfpathlineto{\pgfqpoint{1.840055in}{1.258175in}}%
\pgfpathlineto{\pgfqpoint{1.830416in}{1.269462in}}%
\pgfpathlineto{\pgfqpoint{1.828978in}{1.271423in}}%
\pgfpathlineto{\pgfqpoint{1.827753in}{1.284671in}}%
\pgfpathlineto{\pgfqpoint{1.819628in}{1.297919in}}%
\pgfpathlineto{\pgfqpoint{1.816044in}{1.311167in}}%
\pgfpathlineto{\pgfqpoint{1.810564in}{1.324415in}}%
\pgfpathlineto{\pgfqpoint{1.805133in}{1.333447in}}%
\pgfpathlineto{\pgfqpoint{1.803032in}{1.337663in}}%
\pgfpathlineto{\pgfqpoint{1.801090in}{1.350911in}}%
\pgfpathlineto{\pgfqpoint{1.793270in}{1.364159in}}%
\pgfpathlineto{\pgfqpoint{1.791853in}{1.377407in}}%
\pgfpathlineto{\pgfqpoint{1.783691in}{1.390655in}}%
\pgfpathlineto{\pgfqpoint{1.780583in}{1.403903in}}%
\pgfpathlineto{\pgfqpoint{1.779850in}{1.405585in}}%
\pgfpathlineto{\pgfqpoint{1.775998in}{1.417151in}}%
\pgfpathlineto{\pgfqpoint{1.771139in}{1.430399in}}%
\pgfpathlineto{\pgfqpoint{1.768857in}{1.443647in}}%
\pgfpathlineto{\pgfqpoint{1.760867in}{1.456895in}}%
\pgfpathlineto{\pgfqpoint{1.759317in}{1.470143in}}%
\pgfpathlineto{\pgfqpoint{1.754567in}{1.478161in}}%
\pgfpathlineto{\pgfqpoint{1.752169in}{1.483391in}}%
\pgfpathlineto{\pgfqpoint{1.749934in}{1.496639in}}%
\pgfpathlineto{\pgfqpoint{1.745515in}{1.509887in}}%
\pgfpathlineto{\pgfqpoint{1.741005in}{1.523135in}}%
\pgfpathlineto{\pgfqpoint{1.738754in}{1.536383in}}%
\pgfpathlineto{\pgfqpoint{1.732366in}{1.549631in}}%
\pgfpathlineto{\pgfqpoint{1.730961in}{1.562879in}}%
\pgfpathlineto{\pgfqpoint{1.729284in}{1.565786in}}%
\pgfpathlineto{\pgfqpoint{1.724894in}{1.576127in}}%
\pgfpathlineto{\pgfqpoint{1.722977in}{1.589375in}}%
\pgfpathlineto{\pgfqpoint{1.718587in}{1.602623in}}%
\pgfpathlineto{\pgfqpoint{1.714677in}{1.615871in}}%
\pgfpathlineto{\pgfqpoint{1.712055in}{1.629119in}}%
\pgfpathlineto{\pgfqpoint{1.705579in}{1.642367in}}%
\pgfpathlineto{\pgfqpoint{1.704212in}{1.655615in}}%
\pgfpathlineto{\pgfqpoint{1.704001in}{1.656044in}}%
\pgfpathlineto{\pgfqpoint{1.699569in}{1.668863in}}%
\pgfpathlineto{\pgfqpoint{1.698159in}{1.682111in}}%
\pgfpathlineto{\pgfqpoint{1.694182in}{1.695359in}}%
\pgfpathlineto{\pgfqpoint{1.690816in}{1.708607in}}%
\pgfpathlineto{\pgfqpoint{1.688259in}{1.721855in}}%
\pgfpathlineto{\pgfqpoint{1.683042in}{1.735103in}}%
\pgfpathlineto{\pgfqpoint{1.681469in}{1.748351in}}%
\pgfpathlineto{\pgfqpoint{1.678717in}{1.753844in}}%
\pgfpathlineto{\pgfqpoint{1.676061in}{1.761599in}}%
\pgfpathlineto{\pgfqpoint{1.674815in}{1.774847in}}%
\pgfpathlineto{\pgfqpoint{1.671214in}{1.788095in}}%
\pgfpathlineto{\pgfqpoint{1.668604in}{1.801343in}}%
\pgfpathlineto{\pgfqpoint{1.666270in}{1.814591in}}%
\pgfpathlineto{\pgfqpoint{1.661559in}{1.827839in}}%
\pgfpathlineto{\pgfqpoint{1.660204in}{1.841087in}}%
\pgfpathlineto{\pgfqpoint{1.654703in}{1.854335in}}%
\pgfpathlineto{\pgfqpoint{1.653434in}{1.864192in}}%
\pgfpathlineto{\pgfqpoint{1.653162in}{1.867583in}}%
\pgfpathlineto{\pgfqpoint{1.649876in}{1.880831in}}%
\pgfpathlineto{\pgfqpoint{1.647650in}{1.894079in}}%
\pgfpathlineto{\pgfqpoint{1.645477in}{1.907327in}}%
\pgfpathlineto{\pgfqpoint{1.641822in}{1.920575in}}%
\pgfpathlineto{\pgfqpoint{1.640515in}{1.933823in}}%
\pgfpathlineto{\pgfqpoint{1.635471in}{1.947071in}}%
\pgfpathlineto{\pgfqpoint{1.634008in}{1.960319in}}%
\pgfpathlineto{\pgfqpoint{1.629402in}{1.973567in}}%
\pgfpathlineto{\pgfqpoint{1.628151in}{1.979310in}}%
\pgfpathlineto{\pgfqpoint{1.627152in}{1.986815in}}%
\pgfpathlineto{\pgfqpoint{1.625296in}{2.000063in}}%
\pgfpathlineto{\pgfqpoint{1.622144in}{2.013311in}}%
\pgfpathlineto{\pgfqpoint{1.621089in}{2.026559in}}%
\pgfpathlineto{\pgfqpoint{1.617236in}{2.039807in}}%
\pgfpathlineto{\pgfqpoint{1.616037in}{2.053055in}}%
\pgfpathlineto{\pgfqpoint{1.611982in}{2.066303in}}%
\pgfpathlineto{\pgfqpoint{1.609350in}{2.079551in}}%
\pgfpathlineto{\pgfqpoint{1.606396in}{2.092799in}}%
\pgfpathlineto{\pgfqpoint{1.602868in}{2.103526in}}%
\pgfpathlineto{\pgfqpoint{1.602347in}{2.106047in}}%
\pgfpathlineto{\pgfqpoint{1.601411in}{2.119295in}}%
\pgfpathlineto{\pgfqpoint{1.597931in}{2.132543in}}%
\pgfpathlineto{\pgfqpoint{1.585955in}{2.132543in}}%
\pgfpathlineto{\pgfqpoint{1.590229in}{2.119295in}}%
\pgfpathlineto{\pgfqpoint{1.591264in}{2.106047in}}%
\pgfpathlineto{\pgfqpoint{1.594525in}{2.092799in}}%
\pgfpathlineto{\pgfqpoint{1.596627in}{2.079551in}}%
\pgfpathlineto{\pgfqpoint{1.598802in}{2.066303in}}%
\pgfpathlineto{\pgfqpoint{1.602089in}{2.053055in}}%
\pgfpathlineto{\pgfqpoint{1.602868in}{2.043433in}}%
\pgfpathlineto{\pgfqpoint{1.603297in}{2.039807in}}%
\pgfpathlineto{\pgfqpoint{1.608287in}{2.026559in}}%
\pgfpathlineto{\pgfqpoint{1.609490in}{2.013311in}}%
\pgfpathlineto{\pgfqpoint{1.613474in}{2.000063in}}%
\pgfpathlineto{\pgfqpoint{1.615560in}{1.986815in}}%
\pgfpathlineto{\pgfqpoint{1.617756in}{1.973567in}}%
\pgfpathlineto{\pgfqpoint{1.621072in}{1.960319in}}%
\pgfpathlineto{\pgfqpoint{1.622218in}{1.947071in}}%
\pgfpathlineto{\pgfqpoint{1.626279in}{1.933823in}}%
\pgfpathlineto{\pgfqpoint{1.627318in}{1.920575in}}%
\pgfpathlineto{\pgfqpoint{1.628151in}{1.917335in}}%
\pgfpathlineto{\pgfqpoint{1.631686in}{1.907327in}}%
\pgfpathlineto{\pgfqpoint{1.634208in}{1.894079in}}%
\pgfpathlineto{\pgfqpoint{1.637054in}{1.880831in}}%
\pgfpathlineto{\pgfqpoint{1.640865in}{1.867583in}}%
\pgfpathlineto{\pgfqpoint{1.642125in}{1.854335in}}%
\pgfpathlineto{\pgfqpoint{1.646354in}{1.841087in}}%
\pgfpathlineto{\pgfqpoint{1.647313in}{1.827839in}}%
\pgfpathlineto{\pgfqpoint{1.651270in}{1.814591in}}%
\pgfpathlineto{\pgfqpoint{1.653119in}{1.801343in}}%
\pgfpathlineto{\pgfqpoint{1.653434in}{1.799631in}}%
\pgfpathlineto{\pgfqpoint{1.656342in}{1.788095in}}%
\pgfpathlineto{\pgfqpoint{1.660606in}{1.774847in}}%
\pgfpathlineto{\pgfqpoint{1.662102in}{1.761599in}}%
\pgfpathlineto{\pgfqpoint{1.667323in}{1.748351in}}%
\pgfpathlineto{\pgfqpoint{1.668419in}{1.735103in}}%
\pgfpathlineto{\pgfqpoint{1.672791in}{1.721855in}}%
\pgfpathlineto{\pgfqpoint{1.674742in}{1.708607in}}%
\pgfpathlineto{\pgfqpoint{1.677770in}{1.695359in}}%
\pgfpathlineto{\pgfqpoint{1.678717in}{1.691666in}}%
\pgfpathlineto{\pgfqpoint{1.682032in}{1.682111in}}%
\pgfpathlineto{\pgfqpoint{1.683716in}{1.668863in}}%
\pgfpathlineto{\pgfqpoint{1.689159in}{1.655615in}}%
\pgfpathlineto{\pgfqpoint{1.690003in}{1.642367in}}%
\pgfpathlineto{\pgfqpoint{1.695390in}{1.629119in}}%
\pgfpathlineto{\pgfqpoint{1.697355in}{1.615871in}}%
\pgfpathlineto{\pgfqpoint{1.700899in}{1.602623in}}%
\pgfpathlineto{\pgfqpoint{1.704001in}{1.591674in}}%
\pgfpathlineto{\pgfqpoint{1.704861in}{1.589375in}}%
\pgfpathlineto{\pgfqpoint{1.707198in}{1.576127in}}%
\pgfpathlineto{\pgfqpoint{1.713657in}{1.562879in}}%
\pgfpathlineto{\pgfqpoint{1.714581in}{1.549631in}}%
\pgfpathlineto{\pgfqpoint{1.720185in}{1.536383in}}%
\pgfpathlineto{\pgfqpoint{1.721769in}{1.523135in}}%
\pgfpathlineto{\pgfqpoint{1.725949in}{1.509887in}}%
\pgfpathlineto{\pgfqpoint{1.729284in}{1.498015in}}%
\pgfpathlineto{\pgfqpoint{1.729783in}{1.496639in}}%
\pgfpathlineto{\pgfqpoint{1.732434in}{1.483391in}}%
\pgfpathlineto{\pgfqpoint{1.739212in}{1.470143in}}%
\pgfpathlineto{\pgfqpoint{1.740261in}{1.456895in}}%
\pgfpathlineto{\pgfqpoint{1.747478in}{1.443647in}}%
\pgfpathlineto{\pgfqpoint{1.749110in}{1.430399in}}%
\pgfpathlineto{\pgfqpoint{1.753824in}{1.417151in}}%
\pgfpathlineto{\pgfqpoint{1.754567in}{1.414491in}}%
\pgfpathlineto{\pgfqpoint{1.758257in}{1.403903in}}%
\pgfpathlineto{\pgfqpoint{1.761069in}{1.390655in}}%
\pgfpathlineto{\pgfqpoint{1.768152in}{1.377407in}}%
\pgfpathlineto{\pgfqpoint{1.769130in}{1.364159in}}%
\pgfpathlineto{\pgfqpoint{1.776542in}{1.350911in}}%
\pgfpathlineto{\pgfqpoint{1.777692in}{1.337663in}}%
\pgfpathlineto{\pgfqpoint{1.779850in}{1.333216in}}%
\pgfpathlineto{\pgfqpoint{1.784811in}{1.324415in}}%
\pgfpathlineto{\pgfqpoint{1.789076in}{1.311167in}}%
\pgfpathlineto{\pgfqpoint{1.792461in}{1.297919in}}%
\pgfpathlineto{\pgfqpoint{1.799813in}{1.284671in}}%
\pgfpathlineto{\pgfqpoint{1.800592in}{1.271423in}}%
\pgfpathlineto{\pgfqpoint{1.805133in}{1.264829in}}%
\pgfpathlineto{\pgfqpoint{1.810355in}{1.258175in}}%
\pgfpathlineto{\pgfqpoint{1.811471in}{1.244927in}}%
\pgfpathlineto{\pgfqpoint{1.819407in}{1.231679in}}%
\pgfpathlineto{\pgfqpoint{1.822706in}{1.218431in}}%
\pgfpathlineto{\pgfqpoint{1.826962in}{1.205183in}}%
\pgfpathlineto{\pgfqpoint{1.830416in}{1.198761in}}%
\pgfpathlineto{\pgfqpoint{1.834679in}{1.191935in}}%
\pgfpathlineto{\pgfqpoint{1.835373in}{1.178687in}}%
\pgfpathlineto{\pgfqpoint{1.845890in}{1.165439in}}%
\pgfpathlineto{\pgfqpoint{1.845714in}{1.152191in}}%
\pgfpathlineto{\pgfqpoint{1.855700in}{1.139738in}}%
\pgfpathlineto{\pgfqpoint{1.856397in}{1.138943in}}%
\pgfpathlineto{\pgfqpoint{1.859670in}{1.125695in}}%
\pgfpathlineto{\pgfqpoint{1.865982in}{1.112447in}}%
\pgfpathlineto{\pgfqpoint{1.873542in}{1.099199in}}%
\pgfpathlineto{\pgfqpoint{1.874010in}{1.085951in}}%
\pgfpathlineto{\pgfqpoint{1.880983in}{1.077932in}}%
\pgfpathlineto{\pgfqpoint{1.886008in}{1.072703in}}%
\pgfpathlineto{\pgfqpoint{1.883643in}{1.059455in}}%
\pgfpathlineto{\pgfqpoint{1.895418in}{1.046207in}}%
\pgfpathlineto{\pgfqpoint{1.894979in}{1.032959in}}%
\pgfpathlineto{\pgfqpoint{1.903716in}{1.019711in}}%
\pgfpathlineto{\pgfqpoint{1.906266in}{1.013924in}}%
\pgfpathlineto{\pgfqpoint{1.910014in}{1.006463in}}%
\pgfpathlineto{\pgfqpoint{1.911230in}{0.993215in}}%
\pgfpathlineto{\pgfqpoint{1.923718in}{0.979967in}}%
\pgfpathlineto{\pgfqpoint{1.917517in}{0.966719in}}%
\pgfpathlineto{\pgfqpoint{1.931549in}{0.955082in}}%
\pgfpathlineto{\pgfqpoint{1.933616in}{0.953471in}}%
\pgfpathlineto{\pgfqpoint{1.931549in}{0.949018in}}%
\pgfpathlineto{\pgfqpoint{1.927236in}{0.940223in}}%
\pgfpathlineto{\pgfqpoint{1.931549in}{0.934432in}}%
\pgfpathlineto{\pgfqpoint{1.937355in}{0.926975in}}%
\pgfpathlineto{\pgfqpoint{1.935791in}{0.913727in}}%
\pgfpathlineto{\pgfqpoint{1.935790in}{0.900479in}}%
\pgfpathlineto{\pgfqpoint{1.943812in}{0.887231in}}%
\pgfpathlineto{\pgfqpoint{1.931549in}{0.876016in}}%
\pgfpathlineto{\pgfqpoint{1.929165in}{0.873983in}}%
\pgfpathlineto{\pgfqpoint{1.931549in}{0.871349in}}%
\pgfpathlineto{\pgfqpoint{1.941997in}{0.860735in}}%
\pgfpathlineto{\pgfqpoint{1.931549in}{0.855518in}}%
\pgfpathlineto{\pgfqpoint{1.913157in}{0.847487in}}%
\pgfpathlineto{\pgfqpoint{1.927048in}{0.834239in}}%
\pgfpathlineto{\pgfqpoint{1.906266in}{0.824309in}}%
\pgfpathlineto{\pgfqpoint{1.897894in}{0.820991in}}%
\pgfpathlineto{\pgfqpoint{1.892317in}{0.807743in}}%
\pgfpathlineto{\pgfqpoint{1.880983in}{0.801152in}}%
\pgfpathlineto{\pgfqpoint{1.866649in}{0.794495in}}%
\pgfpathlineto{\pgfqpoint{1.855700in}{0.791138in}}%
\pgfpathlineto{\pgfqpoint{1.830416in}{0.784019in}}%
\pgfpathlineto{\pgfqpoint{1.819715in}{0.781247in}}%
\pgfpathlineto{\pgfqpoint{1.805133in}{0.771933in}}%
\pgfpathlineto{\pgfqpoint{1.797157in}{0.767999in}}%
\pgfpathlineto{\pgfqpoint{1.779850in}{0.765891in}}%
\pgfpathlineto{\pgfqpoint{1.754567in}{0.763111in}}%
\pgfpathlineto{\pgfqpoint{1.729284in}{0.760620in}}%
\pgfpathlineto{\pgfqpoint{1.704001in}{0.758385in}}%
\pgfpathlineto{\pgfqpoint{1.678717in}{0.756379in}}%
\pgfpathlineto{\pgfqpoint{1.655896in}{0.754751in}}%
\pgfpathlineto{\pgfqpoint{1.653434in}{0.754103in}}%
\pgfpathlineto{\pgfqpoint{1.628151in}{0.748579in}}%
\pgfpathlineto{\pgfqpoint{1.602868in}{0.744388in}}%
\pgfpathlineto{\pgfqpoint{1.579795in}{0.741503in}}%
\pgfpathlineto{\pgfqpoint{1.577585in}{0.741433in}}%
\pgfpathlineto{\pgfqpoint{1.552302in}{0.740764in}}%
\pgfpathlineto{\pgfqpoint{1.527018in}{0.740236in}}%
\pgfpathlineto{\pgfqpoint{1.501735in}{0.739843in}}%
\pgfpathlineto{\pgfqpoint{1.476452in}{0.739584in}}%
\pgfpathlineto{\pgfqpoint{1.451169in}{0.739458in}}%
\pgfpathlineto{\pgfqpoint{1.425886in}{0.739467in}}%
\pgfpathlineto{\pgfqpoint{1.400603in}{0.739614in}}%
\pgfpathlineto{\pgfqpoint{1.375319in}{0.739905in}}%
\pgfpathlineto{\pgfqpoint{1.350036in}{0.740348in}}%
\pgfpathlineto{\pgfqpoint{1.324753in}{0.740954in}}%
\pgfpathlineto{\pgfqpoint{1.306859in}{0.741503in}}%
\pgfpathlineto{\pgfqpoint{1.299470in}{0.742019in}}%
\pgfpathlineto{\pgfqpoint{1.274187in}{0.744073in}}%
\pgfpathlineto{\pgfqpoint{1.248904in}{0.746399in}}%
\pgfpathlineto{\pgfqpoint{1.223620in}{0.749000in}}%
\pgfpathlineto{\pgfqpoint{1.198337in}{0.751885in}}%
\pgfpathlineto{\pgfqpoint{1.175560in}{0.754751in}}%
\pgfpathlineto{\pgfqpoint{1.173054in}{0.755096in}}%
\pgfpathlineto{\pgfqpoint{1.147771in}{0.759452in}}%
\pgfpathlineto{\pgfqpoint{1.122488in}{0.765568in}}%
\pgfpathlineto{\pgfqpoint{1.114660in}{0.767999in}}%
\pgfpathlineto{\pgfqpoint{1.097204in}{0.770210in}}%
\pgfpathlineto{\pgfqpoint{1.071921in}{0.773723in}}%
\pgfpathlineto{\pgfqpoint{1.046638in}{0.777622in}}%
\pgfpathlineto{\pgfqpoint{1.025419in}{0.781247in}}%
\pgfpathlineto{\pgfqpoint{1.034081in}{0.794495in}}%
\pgfpathlineto{\pgfqpoint{1.021355in}{0.796779in}}%
\pgfpathlineto{\pgfqpoint{0.996072in}{0.802002in}}%
\pgfpathlineto{\pgfqpoint{0.972846in}{0.807743in}}%
\pgfpathlineto{\pgfqpoint{0.976444in}{0.820991in}}%
\pgfpathlineto{\pgfqpoint{0.970789in}{0.826122in}}%
\pgfpathlineto{\pgfqpoint{0.963464in}{0.834239in}}%
\pgfpathlineto{\pgfqpoint{0.948664in}{0.847487in}}%
\pgfpathlineto{\pgfqpoint{0.962470in}{0.860735in}}%
\pgfpathlineto{\pgfqpoint{0.945505in}{0.870049in}}%
\pgfpathlineto{\pgfqpoint{0.939026in}{0.873983in}}%
\pgfpathlineto{\pgfqpoint{0.945505in}{0.878101in}}%
\pgfpathlineto{\pgfqpoint{0.959116in}{0.887231in}}%
\pgfpathlineto{\pgfqpoint{0.947766in}{0.900479in}}%
\pgfpathlineto{\pgfqpoint{0.956379in}{0.913727in}}%
\pgfpathlineto{\pgfqpoint{0.960136in}{0.926975in}}%
\pgfpathlineto{\pgfqpoint{0.955619in}{0.940223in}}%
\pgfpathlineto{\pgfqpoint{0.970789in}{0.951650in}}%
\pgfpathlineto{\pgfqpoint{0.972995in}{0.953471in}}%
\pgfpathlineto{\pgfqpoint{0.970789in}{0.957511in}}%
\pgfpathlineto{\pgfqpoint{0.965899in}{0.966719in}}%
\pgfpathlineto{\pgfqpoint{0.970789in}{0.970514in}}%
\pgfpathlineto{\pgfqpoint{0.982377in}{0.979967in}}%
\pgfpathlineto{\pgfqpoint{0.979190in}{0.993215in}}%
\pgfpathlineto{\pgfqpoint{0.987589in}{1.006463in}}%
\pgfpathlineto{\pgfqpoint{0.995153in}{1.019711in}}%
\pgfpathlineto{\pgfqpoint{0.995853in}{1.032959in}}%
\pgfpathlineto{\pgfqpoint{0.996072in}{1.033160in}}%
\pgfpathlineto{\pgfqpoint{1.009062in}{1.046207in}}%
\pgfpathlineto{\pgfqpoint{1.005900in}{1.059455in}}%
\pgfpathlineto{\pgfqpoint{1.020620in}{1.072703in}}%
\pgfpathlineto{\pgfqpoint{1.021355in}{1.082197in}}%
\pgfpathlineto{\pgfqpoint{1.021604in}{1.085951in}}%
\pgfpathlineto{\pgfqpoint{1.030302in}{1.099199in}}%
\pgfpathlineto{\pgfqpoint{1.036285in}{1.112447in}}%
\pgfpathlineto{\pgfqpoint{1.037816in}{1.125695in}}%
\pgfpathlineto{\pgfqpoint{1.046638in}{1.135469in}}%
\pgfpathlineto{\pgfqpoint{1.049425in}{1.138943in}}%
\pgfpathlineto{\pgfqpoint{1.048490in}{1.152191in}}%
\pgfpathlineto{\pgfqpoint{1.059745in}{1.165439in}}%
\pgfpathlineto{\pgfqpoint{1.060535in}{1.178687in}}%
\pgfpathlineto{\pgfqpoint{1.068656in}{1.191935in}}%
\pgfpathlineto{\pgfqpoint{1.071921in}{1.198369in}}%
\pgfpathlineto{\pgfqpoint{1.074844in}{1.205183in}}%
\pgfpathlineto{\pgfqpoint{1.078375in}{1.218431in}}%
\pgfpathlineto{\pgfqpoint{1.087427in}{1.231679in}}%
\pgfpathlineto{\pgfqpoint{1.087358in}{1.244927in}}%
\pgfpathlineto{\pgfqpoint{1.097204in}{1.257564in}}%
\pgfpathlineto{\pgfqpoint{1.097628in}{1.258175in}}%
\pgfpathlineto{\pgfqpoint{1.098426in}{1.271423in}}%
\pgfpathlineto{\pgfqpoint{1.105358in}{1.284671in}}%
\pgfpathlineto{\pgfqpoint{1.109411in}{1.297919in}}%
\pgfpathlineto{\pgfqpoint{1.112573in}{1.311167in}}%
\pgfpathlineto{\pgfqpoint{1.121060in}{1.324415in}}%
\pgfpathlineto{\pgfqpoint{1.122488in}{1.337630in}}%
\pgfpathlineto{\pgfqpoint{1.122491in}{1.337663in}}%
\pgfpathlineto{\pgfqpoint{1.130672in}{1.350911in}}%
\pgfpathlineto{\pgfqpoint{1.131463in}{1.364159in}}%
\pgfpathlineto{\pgfqpoint{1.138117in}{1.377407in}}%
\pgfpathlineto{\pgfqpoint{1.141725in}{1.390655in}}%
\pgfpathlineto{\pgfqpoint{1.145774in}{1.403903in}}%
\pgfpathlineto{\pgfqpoint{1.147771in}{1.407711in}}%
\pgfpathlineto{\pgfqpoint{1.151790in}{1.417151in}}%
\pgfpathlineto{\pgfqpoint{1.152753in}{1.430399in}}%
\pgfpathlineto{\pgfqpoint{1.160035in}{1.443647in}}%
\pgfpathlineto{\pgfqpoint{1.161073in}{1.456895in}}%
\pgfpathlineto{\pgfqpoint{1.167833in}{1.470143in}}%
\pgfpathlineto{\pgfqpoint{1.170472in}{1.483391in}}%
\pgfpathlineto{\pgfqpoint{1.173054in}{1.491754in}}%
\pgfpathlineto{\pgfqpoint{1.174271in}{1.496639in}}%
\pgfpathlineto{\pgfqpoint{1.179088in}{1.509887in}}%
\pgfpathlineto{\pgfqpoint{1.180666in}{1.523135in}}%
\pgfpathlineto{\pgfqpoint{1.187210in}{1.536383in}}%
\pgfpathlineto{\pgfqpoint{1.188126in}{1.549631in}}%
\pgfpathlineto{\pgfqpoint{1.194455in}{1.562879in}}%
\pgfpathlineto{\pgfqpoint{1.196894in}{1.576127in}}%
\pgfpathlineto{\pgfqpoint{1.198337in}{1.580311in}}%
\pgfpathlineto{\pgfqpoint{1.200781in}{1.589375in}}%
\pgfpathlineto{\pgfqpoint{1.204281in}{1.602623in}}%
\pgfpathlineto{\pgfqpoint{1.205733in}{1.615871in}}%
\pgfpathlineto{\pgfqpoint{1.211224in}{1.629119in}}%
\pgfpathlineto{\pgfqpoint{1.212211in}{1.642367in}}%
\pgfpathlineto{\pgfqpoint{1.218337in}{1.655615in}}%
\pgfpathlineto{\pgfqpoint{1.220278in}{1.668863in}}%
\pgfpathlineto{\pgfqpoint{1.223620in}{1.678311in}}%
\pgfpathlineto{\pgfqpoint{1.224654in}{1.682111in}}%
\pgfpathlineto{\pgfqpoint{1.227773in}{1.695359in}}%
\pgfpathlineto{\pgfqpoint{1.229739in}{1.708607in}}%
\pgfpathlineto{\pgfqpoint{1.234207in}{1.721855in}}%
\pgfpathlineto{\pgfqpoint{1.234943in}{1.735103in}}%
\pgfpathlineto{\pgfqpoint{1.240087in}{1.748351in}}%
\pgfpathlineto{\pgfqpoint{1.241503in}{1.761599in}}%
\pgfpathlineto{\pgfqpoint{1.246278in}{1.774847in}}%
\pgfpathlineto{\pgfqpoint{1.248904in}{1.785466in}}%
\pgfpathlineto{\pgfqpoint{1.249366in}{1.788095in}}%
\pgfpathlineto{\pgfqpoint{1.251435in}{1.801343in}}%
\pgfpathlineto{\pgfqpoint{1.255267in}{1.814591in}}%
\pgfpathlineto{\pgfqpoint{1.256360in}{1.827839in}}%
\pgfpathlineto{\pgfqpoint{1.261041in}{1.841087in}}%
\pgfpathlineto{\pgfqpoint{1.261975in}{1.854335in}}%
\pgfpathlineto{\pgfqpoint{1.266208in}{1.867583in}}%
\pgfpathlineto{\pgfqpoint{1.268691in}{1.880831in}}%
\pgfpathlineto{\pgfqpoint{1.271943in}{1.894079in}}%
\pgfpathlineto{\pgfqpoint{1.274187in}{1.900937in}}%
\pgfpathlineto{\pgfqpoint{1.275660in}{1.907327in}}%
\pgfpathlineto{\pgfqpoint{1.276777in}{1.920575in}}%
\pgfpathlineto{\pgfqpoint{1.280628in}{1.933823in}}%
\pgfpathlineto{\pgfqpoint{1.281478in}{1.947071in}}%
\pgfpathlineto{\pgfqpoint{1.285579in}{1.960319in}}%
\pgfpathlineto{\pgfqpoint{1.287155in}{1.973567in}}%
\pgfpathlineto{\pgfqpoint{1.290261in}{1.986815in}}%
\pgfpathlineto{\pgfqpoint{1.293750in}{2.000063in}}%
\pgfpathlineto{\pgfqpoint{1.295837in}{2.013311in}}%
\pgfpathlineto{\pgfqpoint{1.299470in}{2.022879in}}%
\pgfpathlineto{\pgfqpoint{1.300472in}{2.026559in}}%
\pgfpathlineto{\pgfqpoint{1.301261in}{2.039807in}}%
\pgfpathlineto{\pgfqpoint{1.304776in}{2.053055in}}%
\pgfpathlineto{\pgfqpoint{1.305920in}{2.066303in}}%
\pgfpathlineto{\pgfqpoint{1.309063in}{2.079551in}}%
\pgfpathlineto{\pgfqpoint{1.311382in}{2.092799in}}%
\pgfpathlineto{\pgfqpoint{1.313420in}{2.106047in}}%
\pgfpathlineto{\pgfqpoint{1.317429in}{2.119295in}}%
\pgfpathlineto{\pgfqpoint{1.318794in}{2.132543in}}%
\pgfpathlineto{\pgfqpoint{1.306419in}{2.132543in}}%
\pgfpathlineto{\pgfqpoint{1.305341in}{2.119295in}}%
\pgfpathlineto{\pgfqpoint{1.302017in}{2.106047in}}%
\pgfpathlineto{\pgfqpoint{1.300473in}{2.092799in}}%
\pgfpathlineto{\pgfqpoint{1.299470in}{2.086572in}}%
\pgfpathlineto{\pgfqpoint{1.297588in}{2.079551in}}%
\pgfpathlineto{\pgfqpoint{1.293542in}{2.066303in}}%
\pgfpathlineto{\pgfqpoint{1.291850in}{2.053055in}}%
\pgfpathlineto{\pgfqpoint{1.287671in}{2.039807in}}%
\pgfpathlineto{\pgfqpoint{1.286706in}{2.026559in}}%
\pgfpathlineto{\pgfqpoint{1.282632in}{2.013311in}}%
\pgfpathlineto{\pgfqpoint{1.281055in}{2.000063in}}%
\pgfpathlineto{\pgfqpoint{1.278071in}{1.986815in}}%
\pgfpathlineto{\pgfqpoint{1.275705in}{1.973567in}}%
\pgfpathlineto{\pgfqpoint{1.274214in}{1.960319in}}%
\pgfpathlineto{\pgfqpoint{1.274187in}{1.960220in}}%
\pgfpathlineto{\pgfqpoint{1.268758in}{1.947071in}}%
\pgfpathlineto{\pgfqpoint{1.267500in}{1.933823in}}%
\pgfpathlineto{\pgfqpoint{1.262708in}{1.920575in}}%
\pgfpathlineto{\pgfqpoint{1.261426in}{1.907327in}}%
\pgfpathlineto{\pgfqpoint{1.257773in}{1.894079in}}%
\pgfpathlineto{\pgfqpoint{1.255303in}{1.880831in}}%
\pgfpathlineto{\pgfqpoint{1.253042in}{1.867583in}}%
\pgfpathlineto{\pgfqpoint{1.249730in}{1.854335in}}%
\pgfpathlineto{\pgfqpoint{1.248904in}{1.842462in}}%
\pgfpathlineto{\pgfqpoint{1.248747in}{1.841087in}}%
\pgfpathlineto{\pgfqpoint{1.242575in}{1.827839in}}%
\pgfpathlineto{\pgfqpoint{1.241184in}{1.814591in}}%
\pgfpathlineto{\pgfqpoint{1.236342in}{1.801343in}}%
\pgfpathlineto{\pgfqpoint{1.234088in}{1.788095in}}%
\pgfpathlineto{\pgfqpoint{1.231209in}{1.774847in}}%
\pgfpathlineto{\pgfqpoint{1.227458in}{1.761599in}}%
\pgfpathlineto{\pgfqpoint{1.226028in}{1.748351in}}%
\pgfpathlineto{\pgfqpoint{1.223620in}{1.741047in}}%
\pgfpathlineto{\pgfqpoint{1.220801in}{1.735103in}}%
\pgfpathlineto{\pgfqpoint{1.219783in}{1.721855in}}%
\pgfpathlineto{\pgfqpoint{1.214118in}{1.708607in}}%
\pgfpathlineto{\pgfqpoint{1.211905in}{1.695359in}}%
\pgfpathlineto{\pgfqpoint{1.207949in}{1.682111in}}%
\pgfpathlineto{\pgfqpoint{1.204300in}{1.668863in}}%
\pgfpathlineto{\pgfqpoint{1.202413in}{1.655615in}}%
\pgfpathlineto{\pgfqpoint{1.198337in}{1.645058in}}%
\pgfpathlineto{\pgfqpoint{1.196896in}{1.642367in}}%
\pgfpathlineto{\pgfqpoint{1.195425in}{1.629119in}}%
\pgfpathlineto{\pgfqpoint{1.188697in}{1.615871in}}%
\pgfpathlineto{\pgfqpoint{1.187094in}{1.602623in}}%
\pgfpathlineto{\pgfqpoint{1.182763in}{1.589375in}}%
\pgfpathlineto{\pgfqpoint{1.179048in}{1.576127in}}%
\pgfpathlineto{\pgfqpoint{1.176547in}{1.562879in}}%
\pgfpathlineto{\pgfqpoint{1.173054in}{1.554087in}}%
\pgfpathlineto{\pgfqpoint{1.170707in}{1.549631in}}%
\pgfpathlineto{\pgfqpoint{1.169288in}{1.536383in}}%
\pgfpathlineto{\pgfqpoint{1.161669in}{1.523135in}}%
\pgfpathlineto{\pgfqpoint{1.159906in}{1.509887in}}%
\pgfpathlineto{\pgfqpoint{1.154173in}{1.496639in}}%
\pgfpathlineto{\pgfqpoint{1.150945in}{1.483391in}}%
\pgfpathlineto{\pgfqpoint{1.148164in}{1.470143in}}%
\pgfpathlineto{\pgfqpoint{1.147771in}{1.469268in}}%
\pgfpathlineto{\pgfqpoint{1.140764in}{1.456895in}}%
\pgfpathlineto{\pgfqpoint{1.138999in}{1.443647in}}%
\pgfpathlineto{\pgfqpoint{1.130966in}{1.430399in}}%
\pgfpathlineto{\pgfqpoint{1.129838in}{1.417151in}}%
\pgfpathlineto{\pgfqpoint{1.123460in}{1.403903in}}%
\pgfpathlineto{\pgfqpoint{1.122488in}{1.400028in}}%
\pgfpathlineto{\pgfqpoint{1.119507in}{1.390655in}}%
\pgfpathlineto{\pgfqpoint{1.114568in}{1.377407in}}%
\pgfpathlineto{\pgfqpoint{1.107762in}{1.364159in}}%
\pgfpathlineto{\pgfqpoint{1.106174in}{1.350911in}}%
\pgfpathlineto{\pgfqpoint{1.097562in}{1.337663in}}%
\pgfpathlineto{\pgfqpoint{1.097204in}{1.334612in}}%
\pgfpathlineto{\pgfqpoint{1.095671in}{1.324415in}}%
\pgfpathlineto{\pgfqpoint{1.085827in}{1.311167in}}%
\pgfpathlineto{\pgfqpoint{1.082725in}{1.297919in}}%
\pgfpathlineto{\pgfqpoint{1.077451in}{1.284671in}}%
\pgfpathlineto{\pgfqpoint{1.071921in}{1.274001in}}%
\pgfpathlineto{\pgfqpoint{1.070350in}{1.271423in}}%
\pgfpathlineto{\pgfqpoint{1.068020in}{1.258175in}}%
\pgfpathlineto{\pgfqpoint{1.057255in}{1.244927in}}%
\pgfpathlineto{\pgfqpoint{1.056630in}{1.231679in}}%
\pgfpathlineto{\pgfqpoint{1.046646in}{1.218431in}}%
\pgfpathlineto{\pgfqpoint{1.046638in}{1.218404in}}%
\pgfpathlineto{\pgfqpoint{1.042409in}{1.205183in}}%
\pgfpathlineto{\pgfqpoint{1.033853in}{1.191935in}}%
\pgfpathlineto{\pgfqpoint{1.026168in}{1.178687in}}%
\pgfpathlineto{\pgfqpoint{1.023658in}{1.165439in}}%
\pgfpathlineto{\pgfqpoint{1.021355in}{1.162831in}}%
\pgfpathlineto{\pgfqpoint{1.010937in}{1.152191in}}%
\pgfpathlineto{\pgfqpoint{1.010679in}{1.138943in}}%
\pgfpathlineto{\pgfqpoint{0.997420in}{1.125695in}}%
\pgfpathlineto{\pgfqpoint{0.996072in}{1.115847in}}%
\pgfpathlineto{\pgfqpoint{0.995533in}{1.112447in}}%
\pgfpathlineto{\pgfqpoint{0.986173in}{1.099199in}}%
\pgfpathlineto{\pgfqpoint{0.977021in}{1.085951in}}%
\pgfpathlineto{\pgfqpoint{0.973084in}{1.072703in}}%
\pgfpathlineto{\pgfqpoint{0.970789in}{1.070667in}}%
\pgfpathlineto{\pgfqpoint{0.957307in}{1.059455in}}%
\pgfpathlineto{\pgfqpoint{0.958412in}{1.046207in}}%
\pgfpathlineto{\pgfqpoint{0.945505in}{1.035434in}}%
\pgfpathlineto{\pgfqpoint{0.942258in}{1.032959in}}%
\pgfpathlineto{\pgfqpoint{0.940494in}{1.019711in}}%
\pgfpathlineto{\pgfqpoint{0.927815in}{1.006463in}}%
\pgfpathlineto{\pgfqpoint{0.920222in}{0.994839in}}%
\pgfpathlineto{\pgfqpoint{0.919088in}{0.993215in}}%
\pgfpathlineto{\pgfqpoint{0.916527in}{0.979967in}}%
\pgfpathlineto{\pgfqpoint{0.898058in}{0.966719in}}%
\pgfpathlineto{\pgfqpoint{0.900847in}{0.953471in}}%
\pgfpathlineto{\pgfqpoint{0.894939in}{0.949831in}}%
\pgfpathlineto{\pgfqpoint{0.878121in}{0.940223in}}%
\pgfpathlineto{\pgfqpoint{0.880852in}{0.926975in}}%
\pgfpathlineto{\pgfqpoint{0.869656in}{0.914903in}}%
\pgfpathlineto{\pgfqpoint{0.868417in}{0.913727in}}%
\pgfpathlineto{\pgfqpoint{0.858081in}{0.900479in}}%
\pgfpathlineto{\pgfqpoint{0.858212in}{0.887231in}}%
\pgfpathlineto{\pgfqpoint{0.844373in}{0.879638in}}%
\pgfpathlineto{\pgfqpoint{0.833691in}{0.873983in}}%
\pgfpathlineto{\pgfqpoint{0.844373in}{0.862836in}}%
\pgfpathlineto{\pgfqpoint{0.846541in}{0.860735in}}%
\pgfpathlineto{\pgfqpoint{0.844373in}{0.859538in}}%
\pgfpathlineto{\pgfqpoint{0.820522in}{0.847487in}}%
\pgfpathlineto{\pgfqpoint{0.828637in}{0.834239in}}%
\pgfpathlineto{\pgfqpoint{0.819090in}{0.823364in}}%
\pgfpathlineto{\pgfqpoint{0.816637in}{0.820991in}}%
\pgfpathlineto{\pgfqpoint{0.809535in}{0.807743in}}%
\pgfpathlineto{\pgfqpoint{0.819090in}{0.801995in}}%
\pgfpathlineto{\pgfqpoint{0.833880in}{0.794495in}}%
\pgfpathlineto{\pgfqpoint{0.819090in}{0.786655in}}%
\pgfpathlineto{\pgfqpoint{0.808266in}{0.781247in}}%
\pgfpathlineto{\pgfqpoint{0.819090in}{0.777685in}}%
\pgfpathlineto{\pgfqpoint{0.844373in}{0.770277in}}%
\pgfpathlineto{\pgfqpoint{0.853032in}{0.767999in}}%
\pgfpathlineto{\pgfqpoint{0.844373in}{0.762473in}}%
\pgfpathlineto{\pgfqpoint{0.830023in}{0.754751in}}%
\pgfpathlineto{\pgfqpoint{0.844373in}{0.750359in}}%
\pgfpathlineto{\pgfqpoint{0.869656in}{0.743187in}}%
\pgfpathlineto{\pgfqpoint{0.875974in}{0.741503in}}%
\pgfpathlineto{\pgfqpoint{0.894939in}{0.735756in}}%
\pgfpathlineto{\pgfqpoint{0.920222in}{0.729951in}}%
\pgfpathlineto{\pgfqpoint{0.929364in}{0.728255in}}%
\pgfpathlineto{\pgfqpoint{0.939877in}{0.715007in}}%
\pgfpathlineto{\pgfqpoint{0.945505in}{0.714467in}}%
\pgfpathlineto{\pgfqpoint{0.970789in}{0.712242in}}%
\pgfpathlineto{\pgfqpoint{0.996072in}{0.710270in}}%
\pgfpathlineto{\pgfqpoint{1.021355in}{0.708518in}}%
\pgfpathlineto{\pgfqpoint{1.046638in}{0.706959in}}%
\pgfpathlineto{\pgfqpoint{1.071921in}{0.705569in}}%
\pgfpathlineto{\pgfqpoint{1.097204in}{0.704330in}}%
\pgfpathlineto{\pgfqpoint{1.122488in}{0.703225in}}%
\pgfpathlineto{\pgfqpoint{1.147771in}{0.702242in}}%
\pgfpathlineto{\pgfqpoint{1.161611in}{0.701759in}}%
\pgfpathlineto{\pgfqpoint{1.173054in}{0.699518in}}%
\pgfpathlineto{\pgfqpoint{1.198337in}{0.695717in}}%
\pgfpathlineto{\pgfqpoint{1.223620in}{0.692986in}}%
\pgfpathlineto{\pgfqpoint{1.248904in}{0.691032in}}%
\pgfpathlineto{\pgfqpoint{1.274187in}{0.689662in}}%
\pgfpathlineto{\pgfqpoint{1.299470in}{0.688744in}}%
\pgfpathlineto{\pgfqpoint{1.309631in}{0.688511in}}%
\pgfpathlineto{\pgfqpoint{1.324753in}{0.688392in}}%
\pgfpathclose%
\pgfpathmoveto{\pgfqpoint{0.844373in}{1.452879in}}%
\pgfpathlineto{\pgfqpoint{0.869656in}{1.456461in}}%
\pgfpathlineto{\pgfqpoint{0.870598in}{1.456895in}}%
\pgfpathlineto{\pgfqpoint{0.869656in}{1.458273in}}%
\pgfpathlineto{\pgfqpoint{0.844373in}{1.463051in}}%
\pgfpathlineto{\pgfqpoint{0.826355in}{1.456895in}}%
\pgfpathlineto{\pgfqpoint{0.844373in}{1.452879in}}%
\pgfpathclose%
\pgfpathmoveto{\pgfqpoint{0.844373in}{1.474302in}}%
\pgfpathlineto{\pgfqpoint{0.869656in}{1.473090in}}%
\pgfpathlineto{\pgfqpoint{0.894939in}{1.478203in}}%
\pgfpathlineto{\pgfqpoint{0.907701in}{1.483391in}}%
\pgfpathlineto{\pgfqpoint{0.903235in}{1.496639in}}%
\pgfpathlineto{\pgfqpoint{0.920222in}{1.503483in}}%
\pgfpathlineto{\pgfqpoint{0.930932in}{1.509887in}}%
\pgfpathlineto{\pgfqpoint{0.934006in}{1.523135in}}%
\pgfpathlineto{\pgfqpoint{0.945505in}{1.534361in}}%
\pgfpathlineto{\pgfqpoint{0.947123in}{1.536383in}}%
\pgfpathlineto{\pgfqpoint{0.961172in}{1.549631in}}%
\pgfpathlineto{\pgfqpoint{0.964020in}{1.562879in}}%
\pgfpathlineto{\pgfqpoint{0.970789in}{1.566390in}}%
\pgfpathlineto{\pgfqpoint{0.983176in}{1.576127in}}%
\pgfpathlineto{\pgfqpoint{0.983282in}{1.589375in}}%
\pgfpathlineto{\pgfqpoint{0.996072in}{1.597670in}}%
\pgfpathlineto{\pgfqpoint{1.001905in}{1.602623in}}%
\pgfpathlineto{\pgfqpoint{1.004393in}{1.615871in}}%
\pgfpathlineto{\pgfqpoint{1.016512in}{1.629119in}}%
\pgfpathlineto{\pgfqpoint{1.021355in}{1.637338in}}%
\pgfpathlineto{\pgfqpoint{1.023652in}{1.642367in}}%
\pgfpathlineto{\pgfqpoint{1.028383in}{1.655615in}}%
\pgfpathlineto{\pgfqpoint{1.041093in}{1.668863in}}%
\pgfpathlineto{\pgfqpoint{1.042777in}{1.682111in}}%
\pgfpathlineto{\pgfqpoint{1.046638in}{1.685422in}}%
\pgfpathlineto{\pgfqpoint{1.055766in}{1.695359in}}%
\pgfpathlineto{\pgfqpoint{1.057128in}{1.708607in}}%
\pgfpathlineto{\pgfqpoint{1.067598in}{1.721855in}}%
\pgfpathlineto{\pgfqpoint{1.071921in}{1.733338in}}%
\pgfpathlineto{\pgfqpoint{1.072491in}{1.735103in}}%
\pgfpathlineto{\pgfqpoint{1.078430in}{1.748351in}}%
\pgfpathlineto{\pgfqpoint{1.086072in}{1.761599in}}%
\pgfpathlineto{\pgfqpoint{1.088207in}{1.774847in}}%
\pgfpathlineto{\pgfqpoint{1.097204in}{1.785150in}}%
\pgfpathlineto{\pgfqpoint{1.099459in}{1.788095in}}%
\pgfpathlineto{\pgfqpoint{1.101211in}{1.801343in}}%
\pgfpathlineto{\pgfqpoint{1.110540in}{1.814591in}}%
\pgfpathlineto{\pgfqpoint{1.113310in}{1.827839in}}%
\pgfpathlineto{\pgfqpoint{1.119360in}{1.841087in}}%
\pgfpathlineto{\pgfqpoint{1.122488in}{1.847614in}}%
\pgfpathlineto{\pgfqpoint{1.125421in}{1.854335in}}%
\pgfpathlineto{\pgfqpoint{1.128320in}{1.867583in}}%
\pgfpathlineto{\pgfqpoint{1.135953in}{1.880831in}}%
\pgfpathlineto{\pgfqpoint{1.136974in}{1.894079in}}%
\pgfpathlineto{\pgfqpoint{1.145434in}{1.907327in}}%
\pgfpathlineto{\pgfqpoint{1.147771in}{1.919756in}}%
\pgfpathlineto{\pgfqpoint{1.147926in}{1.920575in}}%
\pgfpathlineto{\pgfqpoint{1.154670in}{1.933823in}}%
\pgfpathlineto{\pgfqpoint{1.158689in}{1.947071in}}%
\pgfpathlineto{\pgfqpoint{1.161864in}{1.960319in}}%
\pgfpathlineto{\pgfqpoint{1.168266in}{1.973567in}}%
\pgfpathlineto{\pgfqpoint{1.169789in}{1.986815in}}%
\pgfpathlineto{\pgfqpoint{1.173054in}{1.992618in}}%
\pgfpathlineto{\pgfqpoint{1.177252in}{2.000063in}}%
\pgfpathlineto{\pgfqpoint{1.178554in}{2.013311in}}%
\pgfpathlineto{\pgfqpoint{1.184505in}{2.026559in}}%
\pgfpathlineto{\pgfqpoint{1.187718in}{2.039807in}}%
\pgfpathlineto{\pgfqpoint{1.191611in}{2.053055in}}%
\pgfpathlineto{\pgfqpoint{1.196241in}{2.066303in}}%
\pgfpathlineto{\pgfqpoint{1.197861in}{2.079551in}}%
\pgfpathlineto{\pgfqpoint{1.198337in}{2.080520in}}%
\pgfpathlineto{\pgfqpoint{1.204584in}{2.092799in}}%
\pgfpathlineto{\pgfqpoint{1.205926in}{2.106047in}}%
\pgfpathlineto{\pgfqpoint{1.211964in}{2.119295in}}%
\pgfpathlineto{\pgfqpoint{1.214142in}{2.132543in}}%
\pgfpathlineto{\pgfqpoint{1.198337in}{2.132543in}}%
\pgfpathlineto{\pgfqpoint{1.195362in}{2.132543in}}%
\pgfpathlineto{\pgfqpoint{1.192580in}{2.119295in}}%
\pgfpathlineto{\pgfqpoint{1.186008in}{2.106047in}}%
\pgfpathlineto{\pgfqpoint{1.184854in}{2.092799in}}%
\pgfpathlineto{\pgfqpoint{1.177693in}{2.079551in}}%
\pgfpathlineto{\pgfqpoint{1.176578in}{2.066303in}}%
\pgfpathlineto{\pgfqpoint{1.173054in}{2.058273in}}%
\pgfpathlineto{\pgfqpoint{1.170294in}{2.053055in}}%
\pgfpathlineto{\pgfqpoint{1.166329in}{2.039807in}}%
\pgfpathlineto{\pgfqpoint{1.162025in}{2.026559in}}%
\pgfpathlineto{\pgfqpoint{1.155753in}{2.013311in}}%
\pgfpathlineto{\pgfqpoint{1.154606in}{2.000063in}}%
\pgfpathlineto{\pgfqpoint{1.147771in}{1.988899in}}%
\pgfpathlineto{\pgfqpoint{1.146154in}{1.986815in}}%
\pgfpathlineto{\pgfqpoint{1.145249in}{1.973567in}}%
\pgfpathlineto{\pgfqpoint{1.136344in}{1.960319in}}%
\pgfpathlineto{\pgfqpoint{1.133795in}{1.947071in}}%
\pgfpathlineto{\pgfqpoint{1.128354in}{1.933823in}}%
\pgfpathlineto{\pgfqpoint{1.122488in}{1.921891in}}%
\pgfpathlineto{\pgfqpoint{1.121662in}{1.920575in}}%
\pgfpathlineto{\pgfqpoint{1.118135in}{1.907327in}}%
\pgfpathlineto{\pgfqpoint{1.107170in}{1.894079in}}%
\pgfpathlineto{\pgfqpoint{1.107209in}{1.880831in}}%
\pgfpathlineto{\pgfqpoint{1.097204in}{1.867595in}}%
\pgfpathlineto{\pgfqpoint{1.097190in}{1.867583in}}%
\pgfpathlineto{\pgfqpoint{1.095494in}{1.854335in}}%
\pgfpathlineto{\pgfqpoint{1.085395in}{1.841087in}}%
\pgfpathlineto{\pgfqpoint{1.079832in}{1.827839in}}%
\pgfpathlineto{\pgfqpoint{1.076064in}{1.814591in}}%
\pgfpathlineto{\pgfqpoint{1.071921in}{1.809475in}}%
\pgfpathlineto{\pgfqpoint{1.062332in}{1.801343in}}%
\pgfpathlineto{\pgfqpoint{1.061689in}{1.788095in}}%
\pgfpathlineto{\pgfqpoint{1.046638in}{1.775859in}}%
\pgfpathlineto{\pgfqpoint{1.044189in}{1.774847in}}%
\pgfpathlineto{\pgfqpoint{1.046638in}{1.761747in}}%
\pgfpathlineto{\pgfqpoint{1.046650in}{1.761599in}}%
\pgfpathlineto{\pgfqpoint{1.046638in}{1.761585in}}%
\pgfpathlineto{\pgfqpoint{1.027811in}{1.748351in}}%
\pgfpathlineto{\pgfqpoint{1.025088in}{1.735103in}}%
\pgfpathlineto{\pgfqpoint{1.021355in}{1.729018in}}%
\pgfpathlineto{\pgfqpoint{1.010185in}{1.721855in}}%
\pgfpathlineto{\pgfqpoint{0.996072in}{1.711113in}}%
\pgfpathlineto{\pgfqpoint{0.973208in}{1.721855in}}%
\pgfpathlineto{\pgfqpoint{0.970789in}{1.723174in}}%
\pgfpathlineto{\pgfqpoint{0.964688in}{1.735103in}}%
\pgfpathlineto{\pgfqpoint{0.970789in}{1.742724in}}%
\pgfpathlineto{\pgfqpoint{0.978595in}{1.748351in}}%
\pgfpathlineto{\pgfqpoint{0.970789in}{1.754360in}}%
\pgfpathlineto{\pgfqpoint{0.966765in}{1.761599in}}%
\pgfpathlineto{\pgfqpoint{0.970789in}{1.766530in}}%
\pgfpathlineto{\pgfqpoint{0.981349in}{1.774847in}}%
\pgfpathlineto{\pgfqpoint{0.974633in}{1.788095in}}%
\pgfpathlineto{\pgfqpoint{0.980598in}{1.801343in}}%
\pgfpathlineto{\pgfqpoint{0.983584in}{1.814591in}}%
\pgfpathlineto{\pgfqpoint{0.982688in}{1.827839in}}%
\pgfpathlineto{\pgfqpoint{0.991384in}{1.841087in}}%
\pgfpathlineto{\pgfqpoint{0.987476in}{1.854335in}}%
\pgfpathlineto{\pgfqpoint{0.996072in}{1.867267in}}%
\pgfpathlineto{\pgfqpoint{0.996358in}{1.867583in}}%
\pgfpathlineto{\pgfqpoint{0.996072in}{1.869532in}}%
\pgfpathlineto{\pgfqpoint{0.994996in}{1.880831in}}%
\pgfpathlineto{\pgfqpoint{0.996072in}{1.883431in}}%
\pgfpathlineto{\pgfqpoint{1.001627in}{1.894079in}}%
\pgfpathlineto{\pgfqpoint{1.004634in}{1.907327in}}%
\pgfpathlineto{\pgfqpoint{1.004860in}{1.920575in}}%
\pgfpathlineto{\pgfqpoint{1.011716in}{1.933823in}}%
\pgfpathlineto{\pgfqpoint{1.009804in}{1.947071in}}%
\pgfpathlineto{\pgfqpoint{1.017716in}{1.960319in}}%
\pgfpathlineto{\pgfqpoint{1.016930in}{1.973567in}}%
\pgfpathlineto{\pgfqpoint{1.021355in}{1.984769in}}%
\pgfpathlineto{\pgfqpoint{1.022316in}{1.986815in}}%
\pgfpathlineto{\pgfqpoint{1.025727in}{2.000063in}}%
\pgfpathlineto{\pgfqpoint{1.027943in}{2.013311in}}%
\pgfpathlineto{\pgfqpoint{1.034005in}{2.026559in}}%
\pgfpathlineto{\pgfqpoint{1.033038in}{2.039807in}}%
\pgfpathlineto{\pgfqpoint{1.039456in}{2.053055in}}%
\pgfpathlineto{\pgfqpoint{1.038869in}{2.066303in}}%
\pgfpathlineto{\pgfqpoint{1.044151in}{2.079551in}}%
\pgfpathlineto{\pgfqpoint{1.046297in}{2.092799in}}%
\pgfpathlineto{\pgfqpoint{1.046638in}{2.094869in}}%
\pgfpathlineto{\pgfqpoint{1.048567in}{2.106047in}}%
\pgfpathlineto{\pgfqpoint{1.053963in}{2.119295in}}%
\pgfpathlineto{\pgfqpoint{1.054292in}{2.132543in}}%
\pgfpathlineto{\pgfqpoint{1.046638in}{2.132543in}}%
\pgfpathlineto{\pgfqpoint{1.033272in}{2.132543in}}%
\pgfpathlineto{\pgfqpoint{1.032406in}{2.119295in}}%
\pgfpathlineto{\pgfqpoint{1.026801in}{2.106047in}}%
\pgfpathlineto{\pgfqpoint{1.024830in}{2.092799in}}%
\pgfpathlineto{\pgfqpoint{1.021540in}{2.079551in}}%
\pgfpathlineto{\pgfqpoint{1.021355in}{2.079044in}}%
\pgfpathlineto{\pgfqpoint{1.016345in}{2.066303in}}%
\pgfpathlineto{\pgfqpoint{1.015294in}{2.053055in}}%
\pgfpathlineto{\pgfqpoint{1.009027in}{2.039807in}}%
\pgfpathlineto{\pgfqpoint{1.008861in}{2.026559in}}%
\pgfpathlineto{\pgfqpoint{1.002389in}{2.013311in}}%
\pgfpathlineto{\pgfqpoint{1.000351in}{2.000063in}}%
\pgfpathlineto{\pgfqpoint{0.996072in}{1.987920in}}%
\pgfpathlineto{\pgfqpoint{0.995626in}{1.986815in}}%
\pgfpathlineto{\pgfqpoint{0.990763in}{1.973567in}}%
\pgfpathlineto{\pgfqpoint{0.989396in}{1.960319in}}%
\pgfpathlineto{\pgfqpoint{0.981883in}{1.947071in}}%
\pgfpathlineto{\pgfqpoint{0.981958in}{1.933823in}}%
\pgfpathlineto{\pgfqpoint{0.974640in}{1.920575in}}%
\pgfpathlineto{\pgfqpoint{0.973656in}{1.907327in}}%
\pgfpathlineto{\pgfqpoint{0.970789in}{1.900041in}}%
\pgfpathlineto{\pgfqpoint{0.968243in}{1.894079in}}%
\pgfpathlineto{\pgfqpoint{0.963237in}{1.880831in}}%
\pgfpathlineto{\pgfqpoint{0.961187in}{1.867583in}}%
\pgfpathlineto{\pgfqpoint{0.953440in}{1.854335in}}%
\pgfpathlineto{\pgfqpoint{0.954118in}{1.841087in}}%
\pgfpathlineto{\pgfqpoint{0.945505in}{1.828685in}}%
\pgfpathlineto{\pgfqpoint{0.944908in}{1.827839in}}%
\pgfpathlineto{\pgfqpoint{0.944230in}{1.814591in}}%
\pgfpathlineto{\pgfqpoint{0.937582in}{1.801343in}}%
\pgfpathlineto{\pgfqpoint{0.933457in}{1.788095in}}%
\pgfpathlineto{\pgfqpoint{0.931261in}{1.774847in}}%
\pgfpathlineto{\pgfqpoint{0.921869in}{1.761599in}}%
\pgfpathlineto{\pgfqpoint{0.922857in}{1.748351in}}%
\pgfpathlineto{\pgfqpoint{0.920222in}{1.745207in}}%
\pgfpathlineto{\pgfqpoint{0.912336in}{1.735103in}}%
\pgfpathlineto{\pgfqpoint{0.913550in}{1.721855in}}%
\pgfpathlineto{\pgfqpoint{0.906069in}{1.708607in}}%
\pgfpathlineto{\pgfqpoint{0.901864in}{1.695359in}}%
\pgfpathlineto{\pgfqpoint{0.899416in}{1.682111in}}%
\pgfpathlineto{\pgfqpoint{0.894939in}{1.676657in}}%
\pgfpathlineto{\pgfqpoint{0.889395in}{1.668863in}}%
\pgfpathlineto{\pgfqpoint{0.891561in}{1.655615in}}%
\pgfpathlineto{\pgfqpoint{0.878186in}{1.642367in}}%
\pgfpathlineto{\pgfqpoint{0.880304in}{1.629119in}}%
\pgfpathlineto{\pgfqpoint{0.869656in}{1.615933in}}%
\pgfpathlineto{\pgfqpoint{0.869609in}{1.615871in}}%
\pgfpathlineto{\pgfqpoint{0.867822in}{1.602623in}}%
\pgfpathlineto{\pgfqpoint{0.867224in}{1.589375in}}%
\pgfpathlineto{\pgfqpoint{0.856014in}{1.576127in}}%
\pgfpathlineto{\pgfqpoint{0.863001in}{1.562879in}}%
\pgfpathlineto{\pgfqpoint{0.844872in}{1.549631in}}%
\pgfpathlineto{\pgfqpoint{0.854201in}{1.536383in}}%
\pgfpathlineto{\pgfqpoint{0.844373in}{1.527851in}}%
\pgfpathlineto{\pgfqpoint{0.839741in}{1.523135in}}%
\pgfpathlineto{\pgfqpoint{0.837494in}{1.509887in}}%
\pgfpathlineto{\pgfqpoint{0.840588in}{1.496639in}}%
\pgfpathlineto{\pgfqpoint{0.825644in}{1.483391in}}%
\pgfpathlineto{\pgfqpoint{0.844373in}{1.474302in}}%
\pgfpathclose%
\pgfpathmoveto{\pgfqpoint{0.968439in}{1.695359in}}%
\pgfpathlineto{\pgfqpoint{0.970789in}{1.700892in}}%
\pgfpathlineto{\pgfqpoint{0.996072in}{1.695457in}}%
\pgfpathlineto{\pgfqpoint{0.996087in}{1.695359in}}%
\pgfpathlineto{\pgfqpoint{0.996072in}{1.695351in}}%
\pgfpathlineto{\pgfqpoint{0.970789in}{1.693901in}}%
\pgfpathlineto{\pgfqpoint{0.968439in}{1.695359in}}%
\pgfpathclose%
\pgfpathmoveto{\pgfqpoint{2.007399in}{1.502682in}}%
\pgfpathlineto{\pgfqpoint{2.032682in}{1.499924in}}%
\pgfpathlineto{\pgfqpoint{2.049406in}{1.509887in}}%
\pgfpathlineto{\pgfqpoint{2.032682in}{1.521851in}}%
\pgfpathlineto{\pgfqpoint{2.026336in}{1.523135in}}%
\pgfpathlineto{\pgfqpoint{2.032682in}{1.523908in}}%
\pgfpathlineto{\pgfqpoint{2.050745in}{1.536383in}}%
\pgfpathlineto{\pgfqpoint{2.035956in}{1.549631in}}%
\pgfpathlineto{\pgfqpoint{2.043585in}{1.562879in}}%
\pgfpathlineto{\pgfqpoint{2.033923in}{1.576127in}}%
\pgfpathlineto{\pgfqpoint{2.032682in}{1.582861in}}%
\pgfpathlineto{\pgfqpoint{2.031088in}{1.589375in}}%
\pgfpathlineto{\pgfqpoint{2.028621in}{1.602623in}}%
\pgfpathlineto{\pgfqpoint{2.018482in}{1.615871in}}%
\pgfpathlineto{\pgfqpoint{2.021562in}{1.629119in}}%
\pgfpathlineto{\pgfqpoint{2.009039in}{1.642367in}}%
\pgfpathlineto{\pgfqpoint{2.012985in}{1.655615in}}%
\pgfpathlineto{\pgfqpoint{2.007399in}{1.664228in}}%
\pgfpathlineto{\pgfqpoint{2.004073in}{1.668863in}}%
\pgfpathlineto{\pgfqpoint{2.003116in}{1.682111in}}%
\pgfpathlineto{\pgfqpoint{1.998656in}{1.695359in}}%
\pgfpathlineto{\pgfqpoint{1.991159in}{1.708607in}}%
\pgfpathlineto{\pgfqpoint{1.991251in}{1.721855in}}%
\pgfpathlineto{\pgfqpoint{1.982115in}{1.734860in}}%
\pgfpathlineto{\pgfqpoint{1.981921in}{1.735103in}}%
\pgfpathlineto{\pgfqpoint{1.982115in}{1.738312in}}%
\pgfpathlineto{\pgfqpoint{1.982621in}{1.748351in}}%
\pgfpathlineto{\pgfqpoint{1.982115in}{1.749050in}}%
\pgfpathlineto{\pgfqpoint{1.972771in}{1.761599in}}%
\pgfpathlineto{\pgfqpoint{1.971667in}{1.774847in}}%
\pgfpathlineto{\pgfqpoint{1.966535in}{1.788095in}}%
\pgfpathlineto{\pgfqpoint{1.961758in}{1.801343in}}%
\pgfpathlineto{\pgfqpoint{1.960324in}{1.814591in}}%
\pgfpathlineto{\pgfqpoint{1.956832in}{1.820058in}}%
\pgfpathlineto{\pgfqpoint{1.951658in}{1.827839in}}%
\pgfpathlineto{\pgfqpoint{1.952099in}{1.841087in}}%
\pgfpathlineto{\pgfqpoint{1.944469in}{1.854335in}}%
\pgfpathlineto{\pgfqpoint{1.943543in}{1.867583in}}%
\pgfpathlineto{\pgfqpoint{1.937958in}{1.880831in}}%
\pgfpathlineto{\pgfqpoint{1.933677in}{1.894079in}}%
\pgfpathlineto{\pgfqpoint{1.931549in}{1.905267in}}%
\pgfpathlineto{\pgfqpoint{1.931186in}{1.907327in}}%
\pgfpathlineto{\pgfqpoint{1.924582in}{1.920575in}}%
\pgfpathlineto{\pgfqpoint{1.923991in}{1.933823in}}%
\pgfpathlineto{\pgfqpoint{1.916391in}{1.947071in}}%
\pgfpathlineto{\pgfqpoint{1.915431in}{1.960319in}}%
\pgfpathlineto{\pgfqpoint{1.910006in}{1.973567in}}%
\pgfpathlineto{\pgfqpoint{1.907032in}{1.986815in}}%
\pgfpathlineto{\pgfqpoint{1.906266in}{1.990326in}}%
\pgfpathlineto{\pgfqpoint{1.904392in}{2.000063in}}%
\pgfpathlineto{\pgfqpoint{1.898798in}{2.013311in}}%
\pgfpathlineto{\pgfqpoint{1.898197in}{2.026559in}}%
\pgfpathlineto{\pgfqpoint{1.891906in}{2.039807in}}%
\pgfpathlineto{\pgfqpoint{1.890949in}{2.053055in}}%
\pgfpathlineto{\pgfqpoint{1.885118in}{2.066303in}}%
\pgfpathlineto{\pgfqpoint{1.881973in}{2.079551in}}%
\pgfpathlineto{\pgfqpoint{1.880983in}{2.083179in}}%
\pgfpathlineto{\pgfqpoint{1.878747in}{2.092799in}}%
\pgfpathlineto{\pgfqpoint{1.874464in}{2.106047in}}%
\pgfpathlineto{\pgfqpoint{1.873531in}{2.119295in}}%
\pgfpathlineto{\pgfqpoint{1.867829in}{2.132543in}}%
\pgfpathlineto{\pgfqpoint{1.855700in}{2.132543in}}%
\pgfpathlineto{\pgfqpoint{1.845773in}{2.132543in}}%
\pgfpathlineto{\pgfqpoint{1.851665in}{2.119295in}}%
\pgfpathlineto{\pgfqpoint{1.851444in}{2.106047in}}%
\pgfpathlineto{\pgfqpoint{1.855700in}{2.094889in}}%
\pgfpathlineto{\pgfqpoint{1.856442in}{2.092799in}}%
\pgfpathlineto{\pgfqpoint{1.858364in}{2.079551in}}%
\pgfpathlineto{\pgfqpoint{1.861347in}{2.066303in}}%
\pgfpathlineto{\pgfqpoint{1.866204in}{2.053055in}}%
\pgfpathlineto{\pgfqpoint{1.866352in}{2.039807in}}%
\pgfpathlineto{\pgfqpoint{1.872891in}{2.026559in}}%
\pgfpathlineto{\pgfqpoint{1.871748in}{2.013311in}}%
\pgfpathlineto{\pgfqpoint{1.878296in}{2.000063in}}%
\pgfpathlineto{\pgfqpoint{1.879141in}{1.986815in}}%
\pgfpathlineto{\pgfqpoint{1.880983in}{1.979236in}}%
\pgfpathlineto{\pgfqpoint{1.882148in}{1.973567in}}%
\pgfpathlineto{\pgfqpoint{1.886122in}{1.960319in}}%
\pgfpathlineto{\pgfqpoint{1.886003in}{1.947071in}}%
\pgfpathlineto{\pgfqpoint{1.893379in}{1.933823in}}%
\pgfpathlineto{\pgfqpoint{1.891395in}{1.920575in}}%
\pgfpathlineto{\pgfqpoint{1.898993in}{1.907327in}}%
\pgfpathlineto{\pgfqpoint{1.897936in}{1.894079in}}%
\pgfpathlineto{\pgfqpoint{1.903428in}{1.880831in}}%
\pgfpathlineto{\pgfqpoint{1.906266in}{1.871979in}}%
\pgfpathlineto{\pgfqpoint{1.907432in}{1.867583in}}%
\pgfpathlineto{\pgfqpoint{1.906801in}{1.854335in}}%
\pgfpathlineto{\pgfqpoint{1.913784in}{1.841087in}}%
\pgfpathlineto{\pgfqpoint{1.908940in}{1.827839in}}%
\pgfpathlineto{\pgfqpoint{1.918604in}{1.814591in}}%
\pgfpathlineto{\pgfqpoint{1.913117in}{1.801343in}}%
\pgfpathlineto{\pgfqpoint{1.918627in}{1.788095in}}%
\pgfpathlineto{\pgfqpoint{1.916377in}{1.774847in}}%
\pgfpathlineto{\pgfqpoint{1.912625in}{1.761599in}}%
\pgfpathlineto{\pgfqpoint{1.922570in}{1.748351in}}%
\pgfpathlineto{\pgfqpoint{1.906266in}{1.740005in}}%
\pgfpathlineto{\pgfqpoint{1.881189in}{1.748351in}}%
\pgfpathlineto{\pgfqpoint{1.880983in}{1.748633in}}%
\pgfpathlineto{\pgfqpoint{1.877628in}{1.761599in}}%
\pgfpathlineto{\pgfqpoint{1.870577in}{1.774847in}}%
\pgfpathlineto{\pgfqpoint{1.855700in}{1.783845in}}%
\pgfpathlineto{\pgfqpoint{1.851654in}{1.788095in}}%
\pgfpathlineto{\pgfqpoint{1.850855in}{1.801343in}}%
\pgfpathlineto{\pgfqpoint{1.834197in}{1.814591in}}%
\pgfpathlineto{\pgfqpoint{1.833610in}{1.827839in}}%
\pgfpathlineto{\pgfqpoint{1.830416in}{1.830263in}}%
\pgfpathlineto{\pgfqpoint{1.822459in}{1.841087in}}%
\pgfpathlineto{\pgfqpoint{1.817537in}{1.854335in}}%
\pgfpathlineto{\pgfqpoint{1.812171in}{1.867583in}}%
\pgfpathlineto{\pgfqpoint{1.805133in}{1.875994in}}%
\pgfpathlineto{\pgfqpoint{1.802220in}{1.880831in}}%
\pgfpathlineto{\pgfqpoint{1.801784in}{1.894079in}}%
\pgfpathlineto{\pgfqpoint{1.790852in}{1.907327in}}%
\pgfpathlineto{\pgfqpoint{1.789300in}{1.920575in}}%
\pgfpathlineto{\pgfqpoint{1.779850in}{1.933366in}}%
\pgfpathlineto{\pgfqpoint{1.779598in}{1.933823in}}%
\pgfpathlineto{\pgfqpoint{1.775482in}{1.947071in}}%
\pgfpathlineto{\pgfqpoint{1.771702in}{1.960319in}}%
\pgfpathlineto{\pgfqpoint{1.763845in}{1.973567in}}%
\pgfpathlineto{\pgfqpoint{1.762327in}{1.986815in}}%
\pgfpathlineto{\pgfqpoint{1.754567in}{1.997633in}}%
\pgfpathlineto{\pgfqpoint{1.753161in}{2.000063in}}%
\pgfpathlineto{\pgfqpoint{1.752019in}{2.013311in}}%
\pgfpathlineto{\pgfqpoint{1.745491in}{2.026559in}}%
\pgfpathlineto{\pgfqpoint{1.741706in}{2.039807in}}%
\pgfpathlineto{\pgfqpoint{1.737815in}{2.053055in}}%
\pgfpathlineto{\pgfqpoint{1.731328in}{2.066303in}}%
\pgfpathlineto{\pgfqpoint{1.730030in}{2.079551in}}%
\pgfpathlineto{\pgfqpoint{1.729284in}{2.080665in}}%
\pgfpathlineto{\pgfqpoint{1.722622in}{2.092799in}}%
\pgfpathlineto{\pgfqpoint{1.721184in}{2.106047in}}%
\pgfpathlineto{\pgfqpoint{1.715192in}{2.119295in}}%
\pgfpathlineto{\pgfqpoint{1.711835in}{2.132543in}}%
\pgfpathlineto{\pgfqpoint{1.704001in}{2.132543in}}%
\pgfpathlineto{\pgfqpoint{1.691907in}{2.132543in}}%
\pgfpathlineto{\pgfqpoint{1.695260in}{2.119295in}}%
\pgfpathlineto{\pgfqpoint{1.700544in}{2.106047in}}%
\pgfpathlineto{\pgfqpoint{1.702545in}{2.092799in}}%
\pgfpathlineto{\pgfqpoint{1.704001in}{2.089999in}}%
\pgfpathlineto{\pgfqpoint{1.709015in}{2.079551in}}%
\pgfpathlineto{\pgfqpoint{1.710265in}{2.066303in}}%
\pgfpathlineto{\pgfqpoint{1.715645in}{2.053055in}}%
\pgfpathlineto{\pgfqpoint{1.718775in}{2.039807in}}%
\pgfpathlineto{\pgfqpoint{1.722869in}{2.026559in}}%
\pgfpathlineto{\pgfqpoint{1.728444in}{2.013311in}}%
\pgfpathlineto{\pgfqpoint{1.729284in}{2.007482in}}%
\pgfpathlineto{\pgfqpoint{1.730271in}{2.000063in}}%
\pgfpathlineto{\pgfqpoint{1.737095in}{1.986815in}}%
\pgfpathlineto{\pgfqpoint{1.738600in}{1.973567in}}%
\pgfpathlineto{\pgfqpoint{1.745407in}{1.960319in}}%
\pgfpathlineto{\pgfqpoint{1.748482in}{1.947071in}}%
\pgfpathlineto{\pgfqpoint{1.753110in}{1.933823in}}%
\pgfpathlineto{\pgfqpoint{1.754567in}{1.930756in}}%
\pgfpathlineto{\pgfqpoint{1.759586in}{1.920575in}}%
\pgfpathlineto{\pgfqpoint{1.762063in}{1.907327in}}%
\pgfpathlineto{\pgfqpoint{1.770464in}{1.894079in}}%
\pgfpathlineto{\pgfqpoint{1.771632in}{1.880831in}}%
\pgfpathlineto{\pgfqpoint{1.778977in}{1.867583in}}%
\pgfpathlineto{\pgfqpoint{1.779850in}{1.864221in}}%
\pgfpathlineto{\pgfqpoint{1.782593in}{1.854335in}}%
\pgfpathlineto{\pgfqpoint{1.788078in}{1.841087in}}%
\pgfpathlineto{\pgfqpoint{1.794910in}{1.827839in}}%
\pgfpathlineto{\pgfqpoint{1.797373in}{1.814591in}}%
\pgfpathlineto{\pgfqpoint{1.805133in}{1.803571in}}%
\pgfpathlineto{\pgfqpoint{1.806994in}{1.801343in}}%
\pgfpathlineto{\pgfqpoint{1.809266in}{1.788095in}}%
\pgfpathlineto{\pgfqpoint{1.819741in}{1.774847in}}%
\pgfpathlineto{\pgfqpoint{1.823327in}{1.761599in}}%
\pgfpathlineto{\pgfqpoint{1.829237in}{1.748351in}}%
\pgfpathlineto{\pgfqpoint{1.830416in}{1.746267in}}%
\pgfpathlineto{\pgfqpoint{1.837967in}{1.735103in}}%
\pgfpathlineto{\pgfqpoint{1.840506in}{1.721855in}}%
\pgfpathlineto{\pgfqpoint{1.852133in}{1.708607in}}%
\pgfpathlineto{\pgfqpoint{1.853171in}{1.695359in}}%
\pgfpathlineto{\pgfqpoint{1.855700in}{1.692192in}}%
\pgfpathlineto{\pgfqpoint{1.865410in}{1.682111in}}%
\pgfpathlineto{\pgfqpoint{1.870429in}{1.668863in}}%
\pgfpathlineto{\pgfqpoint{1.878954in}{1.655615in}}%
\pgfpathlineto{\pgfqpoint{1.880983in}{1.652789in}}%
\pgfpathlineto{\pgfqpoint{1.891019in}{1.642367in}}%
\pgfpathlineto{\pgfqpoint{1.893764in}{1.629119in}}%
\pgfpathlineto{\pgfqpoint{1.906266in}{1.617834in}}%
\pgfpathlineto{\pgfqpoint{1.909418in}{1.615871in}}%
\pgfpathlineto{\pgfqpoint{1.910250in}{1.602623in}}%
\pgfpathlineto{\pgfqpoint{1.927412in}{1.589375in}}%
\pgfpathlineto{\pgfqpoint{1.931259in}{1.576127in}}%
\pgfpathlineto{\pgfqpoint{1.931549in}{1.575713in}}%
\pgfpathlineto{\pgfqpoint{1.943161in}{1.562879in}}%
\pgfpathlineto{\pgfqpoint{1.956832in}{1.550494in}}%
\pgfpathlineto{\pgfqpoint{1.958570in}{1.549631in}}%
\pgfpathlineto{\pgfqpoint{1.961523in}{1.536383in}}%
\pgfpathlineto{\pgfqpoint{1.982115in}{1.526919in}}%
\pgfpathlineto{\pgfqpoint{2.000096in}{1.523135in}}%
\pgfpathlineto{\pgfqpoint{1.989723in}{1.509887in}}%
\pgfpathlineto{\pgfqpoint{2.007399in}{1.502682in}}%
\pgfpathclose%
\pgfpathmoveto{\pgfqpoint{1.903353in}{1.721855in}}%
\pgfpathlineto{\pgfqpoint{1.906266in}{1.725492in}}%
\pgfpathlineto{\pgfqpoint{1.916527in}{1.721855in}}%
\pgfpathlineto{\pgfqpoint{1.906266in}{1.720517in}}%
\pgfpathlineto{\pgfqpoint{1.903353in}{1.721855in}}%
\pgfpathclose%
\pgfusepath{fill}%
\end{pgfscope}%
\begin{pgfscope}%
\pgfpathrectangle{\pgfqpoint{0.540975in}{0.476543in}}{\pgfqpoint{1.820388in}{1.656000in}}%
\pgfusepath{clip}%
\pgfsetbuttcap%
\pgfsetroundjoin%
\definecolor{currentfill}{rgb}{0.353710,0.557862,0.757170}%
\pgfsetfillcolor{currentfill}%
\pgfsetlinewidth{0.000000pt}%
\definecolor{currentstroke}{rgb}{0.000000,0.000000,0.000000}%
\pgfsetstrokecolor{currentstroke}%
\pgfsetdash{}{0pt}%
\pgfpathmoveto{\pgfqpoint{1.324753in}{0.740954in}}%
\pgfpathlineto{\pgfqpoint{1.350036in}{0.740348in}}%
\pgfpathlineto{\pgfqpoint{1.375319in}{0.739905in}}%
\pgfpathlineto{\pgfqpoint{1.400603in}{0.739614in}}%
\pgfpathlineto{\pgfqpoint{1.425886in}{0.739467in}}%
\pgfpathlineto{\pgfqpoint{1.451169in}{0.739458in}}%
\pgfpathlineto{\pgfqpoint{1.476452in}{0.739584in}}%
\pgfpathlineto{\pgfqpoint{1.501735in}{0.739843in}}%
\pgfpathlineto{\pgfqpoint{1.527018in}{0.740236in}}%
\pgfpathlineto{\pgfqpoint{1.552302in}{0.740764in}}%
\pgfpathlineto{\pgfqpoint{1.577585in}{0.741433in}}%
\pgfpathlineto{\pgfqpoint{1.579795in}{0.741503in}}%
\pgfpathlineto{\pgfqpoint{1.602868in}{0.744388in}}%
\pgfpathlineto{\pgfqpoint{1.628151in}{0.748579in}}%
\pgfpathlineto{\pgfqpoint{1.653434in}{0.754103in}}%
\pgfpathlineto{\pgfqpoint{1.655896in}{0.754751in}}%
\pgfpathlineto{\pgfqpoint{1.678717in}{0.756379in}}%
\pgfpathlineto{\pgfqpoint{1.704001in}{0.758385in}}%
\pgfpathlineto{\pgfqpoint{1.729284in}{0.760620in}}%
\pgfpathlineto{\pgfqpoint{1.754567in}{0.763111in}}%
\pgfpathlineto{\pgfqpoint{1.779850in}{0.765891in}}%
\pgfpathlineto{\pgfqpoint{1.797157in}{0.767999in}}%
\pgfpathlineto{\pgfqpoint{1.805133in}{0.771933in}}%
\pgfpathlineto{\pgfqpoint{1.819715in}{0.781247in}}%
\pgfpathlineto{\pgfqpoint{1.830416in}{0.784019in}}%
\pgfpathlineto{\pgfqpoint{1.855700in}{0.791138in}}%
\pgfpathlineto{\pgfqpoint{1.866649in}{0.794495in}}%
\pgfpathlineto{\pgfqpoint{1.880983in}{0.801152in}}%
\pgfpathlineto{\pgfqpoint{1.892317in}{0.807743in}}%
\pgfpathlineto{\pgfqpoint{1.897894in}{0.820991in}}%
\pgfpathlineto{\pgfqpoint{1.906266in}{0.824309in}}%
\pgfpathlineto{\pgfqpoint{1.927048in}{0.834239in}}%
\pgfpathlineto{\pgfqpoint{1.913157in}{0.847487in}}%
\pgfpathlineto{\pgfqpoint{1.931549in}{0.855518in}}%
\pgfpathlineto{\pgfqpoint{1.941997in}{0.860735in}}%
\pgfpathlineto{\pgfqpoint{1.931549in}{0.871349in}}%
\pgfpathlineto{\pgfqpoint{1.929165in}{0.873983in}}%
\pgfpathlineto{\pgfqpoint{1.931549in}{0.876016in}}%
\pgfpathlineto{\pgfqpoint{1.943812in}{0.887231in}}%
\pgfpathlineto{\pgfqpoint{1.935790in}{0.900479in}}%
\pgfpathlineto{\pgfqpoint{1.935791in}{0.913727in}}%
\pgfpathlineto{\pgfqpoint{1.937355in}{0.926975in}}%
\pgfpathlineto{\pgfqpoint{1.931549in}{0.934432in}}%
\pgfpathlineto{\pgfqpoint{1.927236in}{0.940223in}}%
\pgfpathlineto{\pgfqpoint{1.931549in}{0.949018in}}%
\pgfpathlineto{\pgfqpoint{1.933616in}{0.953471in}}%
\pgfpathlineto{\pgfqpoint{1.931549in}{0.955082in}}%
\pgfpathlineto{\pgfqpoint{1.917517in}{0.966719in}}%
\pgfpathlineto{\pgfqpoint{1.923718in}{0.979967in}}%
\pgfpathlineto{\pgfqpoint{1.911230in}{0.993215in}}%
\pgfpathlineto{\pgfqpoint{1.910014in}{1.006463in}}%
\pgfpathlineto{\pgfqpoint{1.906266in}{1.013924in}}%
\pgfpathlineto{\pgfqpoint{1.903716in}{1.019711in}}%
\pgfpathlineto{\pgfqpoint{1.894979in}{1.032959in}}%
\pgfpathlineto{\pgfqpoint{1.895418in}{1.046207in}}%
\pgfpathlineto{\pgfqpoint{1.883643in}{1.059455in}}%
\pgfpathlineto{\pgfqpoint{1.886008in}{1.072703in}}%
\pgfpathlineto{\pgfqpoint{1.880983in}{1.077932in}}%
\pgfpathlineto{\pgfqpoint{1.874010in}{1.085951in}}%
\pgfpathlineto{\pgfqpoint{1.873542in}{1.099199in}}%
\pgfpathlineto{\pgfqpoint{1.865982in}{1.112447in}}%
\pgfpathlineto{\pgfqpoint{1.859670in}{1.125695in}}%
\pgfpathlineto{\pgfqpoint{1.856397in}{1.138943in}}%
\pgfpathlineto{\pgfqpoint{1.855700in}{1.139738in}}%
\pgfpathlineto{\pgfqpoint{1.845714in}{1.152191in}}%
\pgfpathlineto{\pgfqpoint{1.845890in}{1.165439in}}%
\pgfpathlineto{\pgfqpoint{1.835373in}{1.178687in}}%
\pgfpathlineto{\pgfqpoint{1.834679in}{1.191935in}}%
\pgfpathlineto{\pgfqpoint{1.830416in}{1.198761in}}%
\pgfpathlineto{\pgfqpoint{1.826962in}{1.205183in}}%
\pgfpathlineto{\pgfqpoint{1.822706in}{1.218431in}}%
\pgfpathlineto{\pgfqpoint{1.819407in}{1.231679in}}%
\pgfpathlineto{\pgfqpoint{1.811471in}{1.244927in}}%
\pgfpathlineto{\pgfqpoint{1.810355in}{1.258175in}}%
\pgfpathlineto{\pgfqpoint{1.805133in}{1.264829in}}%
\pgfpathlineto{\pgfqpoint{1.800592in}{1.271423in}}%
\pgfpathlineto{\pgfqpoint{1.799813in}{1.284671in}}%
\pgfpathlineto{\pgfqpoint{1.792461in}{1.297919in}}%
\pgfpathlineto{\pgfqpoint{1.789076in}{1.311167in}}%
\pgfpathlineto{\pgfqpoint{1.784811in}{1.324415in}}%
\pgfpathlineto{\pgfqpoint{1.779850in}{1.333216in}}%
\pgfpathlineto{\pgfqpoint{1.777692in}{1.337663in}}%
\pgfpathlineto{\pgfqpoint{1.776542in}{1.350911in}}%
\pgfpathlineto{\pgfqpoint{1.769130in}{1.364159in}}%
\pgfpathlineto{\pgfqpoint{1.768152in}{1.377407in}}%
\pgfpathlineto{\pgfqpoint{1.761069in}{1.390655in}}%
\pgfpathlineto{\pgfqpoint{1.758257in}{1.403903in}}%
\pgfpathlineto{\pgfqpoint{1.754567in}{1.414491in}}%
\pgfpathlineto{\pgfqpoint{1.753824in}{1.417151in}}%
\pgfpathlineto{\pgfqpoint{1.749110in}{1.430399in}}%
\pgfpathlineto{\pgfqpoint{1.747478in}{1.443647in}}%
\pgfpathlineto{\pgfqpoint{1.740261in}{1.456895in}}%
\pgfpathlineto{\pgfqpoint{1.739212in}{1.470143in}}%
\pgfpathlineto{\pgfqpoint{1.732434in}{1.483391in}}%
\pgfpathlineto{\pgfqpoint{1.729783in}{1.496639in}}%
\pgfpathlineto{\pgfqpoint{1.729284in}{1.498015in}}%
\pgfpathlineto{\pgfqpoint{1.725949in}{1.509887in}}%
\pgfpathlineto{\pgfqpoint{1.721769in}{1.523135in}}%
\pgfpathlineto{\pgfqpoint{1.720185in}{1.536383in}}%
\pgfpathlineto{\pgfqpoint{1.714581in}{1.549631in}}%
\pgfpathlineto{\pgfqpoint{1.713657in}{1.562879in}}%
\pgfpathlineto{\pgfqpoint{1.707198in}{1.576127in}}%
\pgfpathlineto{\pgfqpoint{1.704861in}{1.589375in}}%
\pgfpathlineto{\pgfqpoint{1.704001in}{1.591674in}}%
\pgfpathlineto{\pgfqpoint{1.700899in}{1.602623in}}%
\pgfpathlineto{\pgfqpoint{1.697355in}{1.615871in}}%
\pgfpathlineto{\pgfqpoint{1.695390in}{1.629119in}}%
\pgfpathlineto{\pgfqpoint{1.690003in}{1.642367in}}%
\pgfpathlineto{\pgfqpoint{1.689159in}{1.655615in}}%
\pgfpathlineto{\pgfqpoint{1.683716in}{1.668863in}}%
\pgfpathlineto{\pgfqpoint{1.682032in}{1.682111in}}%
\pgfpathlineto{\pgfqpoint{1.678717in}{1.691666in}}%
\pgfpathlineto{\pgfqpoint{1.677770in}{1.695359in}}%
\pgfpathlineto{\pgfqpoint{1.674742in}{1.708607in}}%
\pgfpathlineto{\pgfqpoint{1.672791in}{1.721855in}}%
\pgfpathlineto{\pgfqpoint{1.668419in}{1.735103in}}%
\pgfpathlineto{\pgfqpoint{1.667323in}{1.748351in}}%
\pgfpathlineto{\pgfqpoint{1.662102in}{1.761599in}}%
\pgfpathlineto{\pgfqpoint{1.660606in}{1.774847in}}%
\pgfpathlineto{\pgfqpoint{1.656342in}{1.788095in}}%
\pgfpathlineto{\pgfqpoint{1.653434in}{1.799631in}}%
\pgfpathlineto{\pgfqpoint{1.653119in}{1.801343in}}%
\pgfpathlineto{\pgfqpoint{1.651270in}{1.814591in}}%
\pgfpathlineto{\pgfqpoint{1.647313in}{1.827839in}}%
\pgfpathlineto{\pgfqpoint{1.646354in}{1.841087in}}%
\pgfpathlineto{\pgfqpoint{1.642125in}{1.854335in}}%
\pgfpathlineto{\pgfqpoint{1.640865in}{1.867583in}}%
\pgfpathlineto{\pgfqpoint{1.637054in}{1.880831in}}%
\pgfpathlineto{\pgfqpoint{1.634208in}{1.894079in}}%
\pgfpathlineto{\pgfqpoint{1.631686in}{1.907327in}}%
\pgfpathlineto{\pgfqpoint{1.628151in}{1.917335in}}%
\pgfpathlineto{\pgfqpoint{1.627318in}{1.920575in}}%
\pgfpathlineto{\pgfqpoint{1.626279in}{1.933823in}}%
\pgfpathlineto{\pgfqpoint{1.622218in}{1.947071in}}%
\pgfpathlineto{\pgfqpoint{1.621072in}{1.960319in}}%
\pgfpathlineto{\pgfqpoint{1.617756in}{1.973567in}}%
\pgfpathlineto{\pgfqpoint{1.615560in}{1.986815in}}%
\pgfpathlineto{\pgfqpoint{1.613474in}{2.000063in}}%
\pgfpathlineto{\pgfqpoint{1.609490in}{2.013311in}}%
\pgfpathlineto{\pgfqpoint{1.608287in}{2.026559in}}%
\pgfpathlineto{\pgfqpoint{1.603297in}{2.039807in}}%
\pgfpathlineto{\pgfqpoint{1.602868in}{2.043433in}}%
\pgfpathlineto{\pgfqpoint{1.602089in}{2.053055in}}%
\pgfpathlineto{\pgfqpoint{1.598802in}{2.066303in}}%
\pgfpathlineto{\pgfqpoint{1.596627in}{2.079551in}}%
\pgfpathlineto{\pgfqpoint{1.594525in}{2.092799in}}%
\pgfpathlineto{\pgfqpoint{1.591264in}{2.106047in}}%
\pgfpathlineto{\pgfqpoint{1.590229in}{2.119295in}}%
\pgfpathlineto{\pgfqpoint{1.585955in}{2.132543in}}%
\pgfpathlineto{\pgfqpoint{1.577585in}{2.132543in}}%
\pgfpathlineto{\pgfqpoint{1.574938in}{2.132543in}}%
\pgfpathlineto{\pgfqpoint{1.577585in}{2.123307in}}%
\pgfpathlineto{\pgfqpoint{1.579048in}{2.119295in}}%
\pgfpathlineto{\pgfqpoint{1.580180in}{2.106047in}}%
\pgfpathlineto{\pgfqpoint{1.584099in}{2.092799in}}%
\pgfpathlineto{\pgfqpoint{1.586363in}{2.079551in}}%
\pgfpathlineto{\pgfqpoint{1.588969in}{2.066303in}}%
\pgfpathlineto{\pgfqpoint{1.592558in}{2.053055in}}%
\pgfpathlineto{\pgfqpoint{1.593776in}{2.039807in}}%
\pgfpathlineto{\pgfqpoint{1.597702in}{2.026559in}}%
\pgfpathlineto{\pgfqpoint{1.598582in}{2.013311in}}%
\pgfpathlineto{\pgfqpoint{1.601985in}{2.000063in}}%
\pgfpathlineto{\pgfqpoint{1.602868in}{1.993174in}}%
\pgfpathlineto{\pgfqpoint{1.603968in}{1.986815in}}%
\pgfpathlineto{\pgfqpoint{1.606627in}{1.973567in}}%
\pgfpathlineto{\pgfqpoint{1.610353in}{1.960319in}}%
\pgfpathlineto{\pgfqpoint{1.611672in}{1.947071in}}%
\pgfpathlineto{\pgfqpoint{1.616246in}{1.933823in}}%
\pgfpathlineto{\pgfqpoint{1.617277in}{1.920575in}}%
\pgfpathlineto{\pgfqpoint{1.621048in}{1.907327in}}%
\pgfpathlineto{\pgfqpoint{1.622892in}{1.894079in}}%
\pgfpathlineto{\pgfqpoint{1.625327in}{1.880831in}}%
\pgfpathlineto{\pgfqpoint{1.628151in}{1.868907in}}%
\pgfpathlineto{\pgfqpoint{1.628568in}{1.867583in}}%
\pgfpathlineto{\pgfqpoint{1.630040in}{1.854335in}}%
\pgfpathlineto{\pgfqpoint{1.634925in}{1.841087in}}%
\pgfpathlineto{\pgfqpoint{1.635871in}{1.827839in}}%
\pgfpathlineto{\pgfqpoint{1.640428in}{1.814591in}}%
\pgfpathlineto{\pgfqpoint{1.642254in}{1.801343in}}%
\pgfpathlineto{\pgfqpoint{1.644987in}{1.788095in}}%
\pgfpathlineto{\pgfqpoint{1.648295in}{1.774847in}}%
\pgfpathlineto{\pgfqpoint{1.649543in}{1.761599in}}%
\pgfpathlineto{\pgfqpoint{1.653434in}{1.750018in}}%
\pgfpathlineto{\pgfqpoint{1.654153in}{1.748351in}}%
\pgfpathlineto{\pgfqpoint{1.655273in}{1.735103in}}%
\pgfpathlineto{\pgfqpoint{1.660381in}{1.721855in}}%
\pgfpathlineto{\pgfqpoint{1.662317in}{1.708607in}}%
\pgfpathlineto{\pgfqpoint{1.665839in}{1.695359in}}%
\pgfpathlineto{\pgfqpoint{1.669330in}{1.682111in}}%
\pgfpathlineto{\pgfqpoint{1.670703in}{1.668863in}}%
\pgfpathlineto{\pgfqpoint{1.675265in}{1.655615in}}%
\pgfpathlineto{\pgfqpoint{1.675822in}{1.642367in}}%
\pgfpathlineto{\pgfqpoint{1.678717in}{1.634506in}}%
\pgfpathlineto{\pgfqpoint{1.681162in}{1.629119in}}%
\pgfpathlineto{\pgfqpoint{1.683075in}{1.615871in}}%
\pgfpathlineto{\pgfqpoint{1.687190in}{1.602623in}}%
\pgfpathlineto{\pgfqpoint{1.691061in}{1.589375in}}%
\pgfpathlineto{\pgfqpoint{1.693037in}{1.576127in}}%
\pgfpathlineto{\pgfqpoint{1.698445in}{1.562879in}}%
\pgfpathlineto{\pgfqpoint{1.699061in}{1.549631in}}%
\pgfpathlineto{\pgfqpoint{1.704001in}{1.536771in}}%
\pgfpathlineto{\pgfqpoint{1.704181in}{1.536383in}}%
\pgfpathlineto{\pgfqpoint{1.705539in}{1.523135in}}%
\pgfpathlineto{\pgfqpoint{1.710348in}{1.509887in}}%
\pgfpathlineto{\pgfqpoint{1.714055in}{1.496639in}}%
\pgfpathlineto{\pgfqpoint{1.716326in}{1.483391in}}%
\pgfpathlineto{\pgfqpoint{1.722109in}{1.470143in}}%
\pgfpathlineto{\pgfqpoint{1.722790in}{1.456895in}}%
\pgfpathlineto{\pgfqpoint{1.729284in}{1.443820in}}%
\pgfpathlineto{\pgfqpoint{1.729385in}{1.443647in}}%
\pgfpathlineto{\pgfqpoint{1.730688in}{1.430399in}}%
\pgfpathlineto{\pgfqpoint{1.736070in}{1.417151in}}%
\pgfpathlineto{\pgfqpoint{1.739606in}{1.403903in}}%
\pgfpathlineto{\pgfqpoint{1.742160in}{1.390655in}}%
\pgfpathlineto{\pgfqpoint{1.748272in}{1.377407in}}%
\pgfpathlineto{\pgfqpoint{1.748874in}{1.364159in}}%
\pgfpathlineto{\pgfqpoint{1.754567in}{1.353285in}}%
\pgfpathlineto{\pgfqpoint{1.755994in}{1.350911in}}%
\pgfpathlineto{\pgfqpoint{1.756503in}{1.337663in}}%
\pgfpathlineto{\pgfqpoint{1.763468in}{1.324415in}}%
\pgfpathlineto{\pgfqpoint{1.766587in}{1.311167in}}%
\pgfpathlineto{\pgfqpoint{1.769724in}{1.297919in}}%
\pgfpathlineto{\pgfqpoint{1.776151in}{1.284671in}}%
\pgfpathlineto{\pgfqpoint{1.776528in}{1.271423in}}%
\pgfpathlineto{\pgfqpoint{1.779850in}{1.266245in}}%
\pgfpathlineto{\pgfqpoint{1.785645in}{1.258175in}}%
\pgfpathlineto{\pgfqpoint{1.785814in}{1.244927in}}%
\pgfpathlineto{\pgfqpoint{1.793413in}{1.231679in}}%
\pgfpathlineto{\pgfqpoint{1.795377in}{1.218431in}}%
\pgfpathlineto{\pgfqpoint{1.799442in}{1.205183in}}%
\pgfpathlineto{\pgfqpoint{1.805133in}{1.192391in}}%
\pgfpathlineto{\pgfqpoint{1.805366in}{1.191935in}}%
\pgfpathlineto{\pgfqpoint{1.805381in}{1.178687in}}%
\pgfpathlineto{\pgfqpoint{1.814985in}{1.165439in}}%
\pgfpathlineto{\pgfqpoint{1.813440in}{1.152191in}}%
\pgfpathlineto{\pgfqpoint{1.823689in}{1.138943in}}%
\pgfpathlineto{\pgfqpoint{1.824764in}{1.125695in}}%
\pgfpathlineto{\pgfqpoint{1.830306in}{1.112447in}}%
\pgfpathlineto{\pgfqpoint{1.830416in}{1.112147in}}%
\pgfpathlineto{\pgfqpoint{1.835807in}{1.099199in}}%
\pgfpathlineto{\pgfqpoint{1.835501in}{1.085951in}}%
\pgfpathlineto{\pgfqpoint{1.845471in}{1.072703in}}%
\pgfpathlineto{\pgfqpoint{1.841233in}{1.059455in}}%
\pgfpathlineto{\pgfqpoint{1.851596in}{1.046207in}}%
\pgfpathlineto{\pgfqpoint{1.847825in}{1.032959in}}%
\pgfpathlineto{\pgfqpoint{1.855700in}{1.019793in}}%
\pgfpathlineto{\pgfqpoint{1.855751in}{1.019711in}}%
\pgfpathlineto{\pgfqpoint{1.857806in}{1.006463in}}%
\pgfpathlineto{\pgfqpoint{1.857379in}{0.993215in}}%
\pgfpathlineto{\pgfqpoint{1.865917in}{0.979967in}}%
\pgfpathlineto{\pgfqpoint{1.856316in}{0.966719in}}%
\pgfpathlineto{\pgfqpoint{1.869505in}{0.953471in}}%
\pgfpathlineto{\pgfqpoint{1.856940in}{0.940223in}}%
\pgfpathlineto{\pgfqpoint{1.865350in}{0.926975in}}%
\pgfpathlineto{\pgfqpoint{1.855700in}{0.914761in}}%
\pgfpathlineto{\pgfqpoint{1.854770in}{0.913727in}}%
\pgfpathlineto{\pgfqpoint{1.852291in}{0.900479in}}%
\pgfpathlineto{\pgfqpoint{1.849859in}{0.887231in}}%
\pgfpathlineto{\pgfqpoint{1.830416in}{0.875697in}}%
\pgfpathlineto{\pgfqpoint{1.827245in}{0.873983in}}%
\pgfpathlineto{\pgfqpoint{1.830194in}{0.860735in}}%
\pgfpathlineto{\pgfqpoint{1.805133in}{0.853269in}}%
\pgfpathlineto{\pgfqpoint{1.782333in}{0.847487in}}%
\pgfpathlineto{\pgfqpoint{1.787895in}{0.834239in}}%
\pgfpathlineto{\pgfqpoint{1.779850in}{0.832241in}}%
\pgfpathlineto{\pgfqpoint{1.754567in}{0.826813in}}%
\pgfpathlineto{\pgfqpoint{1.729284in}{0.822353in}}%
\pgfpathlineto{\pgfqpoint{1.720309in}{0.820991in}}%
\pgfpathlineto{\pgfqpoint{1.704001in}{0.808143in}}%
\pgfpathlineto{\pgfqpoint{1.703435in}{0.807743in}}%
\pgfpathlineto{\pgfqpoint{1.678717in}{0.803819in}}%
\pgfpathlineto{\pgfqpoint{1.653434in}{0.800621in}}%
\pgfpathlineto{\pgfqpoint{1.628151in}{0.798084in}}%
\pgfpathlineto{\pgfqpoint{1.602868in}{0.796086in}}%
\pgfpathlineto{\pgfqpoint{1.577585in}{0.794536in}}%
\pgfpathlineto{\pgfqpoint{1.576748in}{0.794495in}}%
\pgfpathlineto{\pgfqpoint{1.552302in}{0.792159in}}%
\pgfpathlineto{\pgfqpoint{1.527018in}{0.790154in}}%
\pgfpathlineto{\pgfqpoint{1.501735in}{0.788570in}}%
\pgfpathlineto{\pgfqpoint{1.476452in}{0.787431in}}%
\pgfpathlineto{\pgfqpoint{1.451169in}{0.786778in}}%
\pgfpathlineto{\pgfqpoint{1.425886in}{0.786675in}}%
\pgfpathlineto{\pgfqpoint{1.400603in}{0.787216in}}%
\pgfpathlineto{\pgfqpoint{1.375319in}{0.788536in}}%
\pgfpathlineto{\pgfqpoint{1.350036in}{0.790835in}}%
\pgfpathlineto{\pgfqpoint{1.324753in}{0.794407in}}%
\pgfpathlineto{\pgfqpoint{1.324284in}{0.794495in}}%
\pgfpathlineto{\pgfqpoint{1.299470in}{0.795696in}}%
\pgfpathlineto{\pgfqpoint{1.274187in}{0.797162in}}%
\pgfpathlineto{\pgfqpoint{1.248904in}{0.798891in}}%
\pgfpathlineto{\pgfqpoint{1.223620in}{0.800913in}}%
\pgfpathlineto{\pgfqpoint{1.198337in}{0.803264in}}%
\pgfpathlineto{\pgfqpoint{1.173054in}{0.805992in}}%
\pgfpathlineto{\pgfqpoint{1.158822in}{0.807743in}}%
\pgfpathlineto{\pgfqpoint{1.153790in}{0.820991in}}%
\pgfpathlineto{\pgfqpoint{1.147771in}{0.822133in}}%
\pgfpathlineto{\pgfqpoint{1.122488in}{0.827937in}}%
\pgfpathlineto{\pgfqpoint{1.101257in}{0.834239in}}%
\pgfpathlineto{\pgfqpoint{1.097204in}{0.836367in}}%
\pgfpathlineto{\pgfqpoint{1.077725in}{0.847487in}}%
\pgfpathlineto{\pgfqpoint{1.074481in}{0.860735in}}%
\pgfpathlineto{\pgfqpoint{1.071921in}{0.861585in}}%
\pgfpathlineto{\pgfqpoint{1.046638in}{0.870865in}}%
\pgfpathlineto{\pgfqpoint{1.039044in}{0.873983in}}%
\pgfpathlineto{\pgfqpoint{1.046638in}{0.880871in}}%
\pgfpathlineto{\pgfqpoint{1.053196in}{0.887231in}}%
\pgfpathlineto{\pgfqpoint{1.046638in}{0.890605in}}%
\pgfpathlineto{\pgfqpoint{1.030082in}{0.900479in}}%
\pgfpathlineto{\pgfqpoint{1.036972in}{0.913727in}}%
\pgfpathlineto{\pgfqpoint{1.031438in}{0.926975in}}%
\pgfpathlineto{\pgfqpoint{1.025219in}{0.940223in}}%
\pgfpathlineto{\pgfqpoint{1.036875in}{0.953471in}}%
\pgfpathlineto{\pgfqpoint{1.026010in}{0.966719in}}%
\pgfpathlineto{\pgfqpoint{1.040362in}{0.979967in}}%
\pgfpathlineto{\pgfqpoint{1.031906in}{0.993215in}}%
\pgfpathlineto{\pgfqpoint{1.039636in}{1.006463in}}%
\pgfpathlineto{\pgfqpoint{1.042338in}{1.019711in}}%
\pgfpathlineto{\pgfqpoint{1.042095in}{1.032959in}}%
\pgfpathlineto{\pgfqpoint{1.046638in}{1.038332in}}%
\pgfpathlineto{\pgfqpoint{1.052695in}{1.046207in}}%
\pgfpathlineto{\pgfqpoint{1.047834in}{1.059455in}}%
\pgfpathlineto{\pgfqpoint{1.061055in}{1.072703in}}%
\pgfpathlineto{\pgfqpoint{1.059309in}{1.085951in}}%
\pgfpathlineto{\pgfqpoint{1.068045in}{1.099199in}}%
\pgfpathlineto{\pgfqpoint{1.071529in}{1.112447in}}%
\pgfpathlineto{\pgfqpoint{1.071921in}{1.116797in}}%
\pgfpathlineto{\pgfqpoint{1.072658in}{1.125695in}}%
\pgfpathlineto{\pgfqpoint{1.081885in}{1.138943in}}%
\pgfpathlineto{\pgfqpoint{1.079910in}{1.152191in}}%
\pgfpathlineto{\pgfqpoint{1.090641in}{1.165439in}}%
\pgfpathlineto{\pgfqpoint{1.089823in}{1.178687in}}%
\pgfpathlineto{\pgfqpoint{1.097204in}{1.190775in}}%
\pgfpathlineto{\pgfqpoint{1.097864in}{1.191935in}}%
\pgfpathlineto{\pgfqpoint{1.102046in}{1.205183in}}%
\pgfpathlineto{\pgfqpoint{1.105464in}{1.218431in}}%
\pgfpathlineto{\pgfqpoint{1.113316in}{1.231679in}}%
\pgfpathlineto{\pgfqpoint{1.112666in}{1.244927in}}%
\pgfpathlineto{\pgfqpoint{1.122227in}{1.258175in}}%
\pgfpathlineto{\pgfqpoint{1.122040in}{1.271423in}}%
\pgfpathlineto{\pgfqpoint{1.122488in}{1.272204in}}%
\pgfpathlineto{\pgfqpoint{1.128933in}{1.284671in}}%
\pgfpathlineto{\pgfqpoint{1.131816in}{1.297919in}}%
\pgfpathlineto{\pgfqpoint{1.134919in}{1.311167in}}%
\pgfpathlineto{\pgfqpoint{1.142060in}{1.324415in}}%
\pgfpathlineto{\pgfqpoint{1.143036in}{1.337663in}}%
\pgfpathlineto{\pgfqpoint{1.147771in}{1.344815in}}%
\pgfpathlineto{\pgfqpoint{1.151307in}{1.350911in}}%
\pgfpathlineto{\pgfqpoint{1.151446in}{1.364159in}}%
\pgfpathlineto{\pgfqpoint{1.157811in}{1.377407in}}%
\pgfpathlineto{\pgfqpoint{1.160365in}{1.390655in}}%
\pgfpathlineto{\pgfqpoint{1.164177in}{1.403903in}}%
\pgfpathlineto{\pgfqpoint{1.169906in}{1.417151in}}%
\pgfpathlineto{\pgfqpoint{1.170881in}{1.430399in}}%
\pgfpathlineto{\pgfqpoint{1.173054in}{1.434081in}}%
\pgfpathlineto{\pgfqpoint{1.177849in}{1.443647in}}%
\pgfpathlineto{\pgfqpoint{1.178368in}{1.456895in}}%
\pgfpathlineto{\pgfqpoint{1.184507in}{1.470143in}}%
\pgfpathlineto{\pgfqpoint{1.186263in}{1.483391in}}%
\pgfpathlineto{\pgfqpoint{1.190027in}{1.496639in}}%
\pgfpathlineto{\pgfqpoint{1.194946in}{1.509887in}}%
\pgfpathlineto{\pgfqpoint{1.196725in}{1.523135in}}%
\pgfpathlineto{\pgfqpoint{1.198337in}{1.526160in}}%
\pgfpathlineto{\pgfqpoint{1.202838in}{1.536383in}}%
\pgfpathlineto{\pgfqpoint{1.203435in}{1.549631in}}%
\pgfpathlineto{\pgfqpoint{1.209043in}{1.562879in}}%
\pgfpathlineto{\pgfqpoint{1.210658in}{1.576127in}}%
\pgfpathlineto{\pgfqpoint{1.214731in}{1.589375in}}%
\pgfpathlineto{\pgfqpoint{1.218275in}{1.602623in}}%
\pgfpathlineto{\pgfqpoint{1.219979in}{1.615871in}}%
\pgfpathlineto{\pgfqpoint{1.223620in}{1.623901in}}%
\pgfpathlineto{\pgfqpoint{1.225524in}{1.629119in}}%
\pgfpathlineto{\pgfqpoint{1.226252in}{1.642367in}}%
\pgfpathlineto{\pgfqpoint{1.231547in}{1.655615in}}%
\pgfpathlineto{\pgfqpoint{1.232880in}{1.668863in}}%
\pgfpathlineto{\pgfqpoint{1.236932in}{1.682111in}}%
\pgfpathlineto{\pgfqpoint{1.240070in}{1.695359in}}%
\pgfpathlineto{\pgfqpoint{1.242388in}{1.708607in}}%
\pgfpathlineto{\pgfqpoint{1.247313in}{1.721855in}}%
\pgfpathlineto{\pgfqpoint{1.248099in}{1.735103in}}%
\pgfpathlineto{\pgfqpoint{1.248904in}{1.736970in}}%
\pgfpathlineto{\pgfqpoint{1.252777in}{1.748351in}}%
\pgfpathlineto{\pgfqpoint{1.253756in}{1.761599in}}%
\pgfpathlineto{\pgfqpoint{1.257779in}{1.774847in}}%
\pgfpathlineto{\pgfqpoint{1.260085in}{1.788095in}}%
\pgfpathlineto{\pgfqpoint{1.262538in}{1.801343in}}%
\pgfpathlineto{\pgfqpoint{1.266660in}{1.814591in}}%
\pgfpathlineto{\pgfqpoint{1.267897in}{1.827839in}}%
\pgfpathlineto{\pgfqpoint{1.273274in}{1.841087in}}%
\pgfpathlineto{\pgfqpoint{1.274187in}{1.853866in}}%
\pgfpathlineto{\pgfqpoint{1.274212in}{1.854335in}}%
\pgfpathlineto{\pgfqpoint{1.277929in}{1.867583in}}%
\pgfpathlineto{\pgfqpoint{1.279755in}{1.880831in}}%
\pgfpathlineto{\pgfqpoint{1.282471in}{1.894079in}}%
\pgfpathlineto{\pgfqpoint{1.285645in}{1.907327in}}%
\pgfpathlineto{\pgfqpoint{1.286950in}{1.920575in}}%
\pgfpathlineto{\pgfqpoint{1.291266in}{1.933823in}}%
\pgfpathlineto{\pgfqpoint{1.292123in}{1.947071in}}%
\pgfpathlineto{\pgfqpoint{1.296944in}{1.960319in}}%
\pgfpathlineto{\pgfqpoint{1.298605in}{1.973567in}}%
\pgfpathlineto{\pgfqpoint{1.299470in}{1.976694in}}%
\pgfpathlineto{\pgfqpoint{1.301629in}{1.986815in}}%
\pgfpathlineto{\pgfqpoint{1.304354in}{2.000063in}}%
\pgfpathlineto{\pgfqpoint{1.306112in}{2.013311in}}%
\pgfpathlineto{\pgfqpoint{1.309936in}{2.026559in}}%
\pgfpathlineto{\pgfqpoint{1.310793in}{2.039807in}}%
\pgfpathlineto{\pgfqpoint{1.314859in}{2.053055in}}%
\pgfpathlineto{\pgfqpoint{1.315990in}{2.066303in}}%
\pgfpathlineto{\pgfqpoint{1.319766in}{2.079551in}}%
\pgfpathlineto{\pgfqpoint{1.322291in}{2.092799in}}%
\pgfpathlineto{\pgfqpoint{1.324753in}{2.105695in}}%
\pgfpathlineto{\pgfqpoint{1.324806in}{2.106047in}}%
\pgfpathlineto{\pgfqpoint{1.328198in}{2.119295in}}%
\pgfpathlineto{\pgfqpoint{1.329357in}{2.132543in}}%
\pgfpathlineto{\pgfqpoint{1.324753in}{2.132543in}}%
\pgfpathlineto{\pgfqpoint{1.318794in}{2.132543in}}%
\pgfpathlineto{\pgfqpoint{1.317429in}{2.119295in}}%
\pgfpathlineto{\pgfqpoint{1.313420in}{2.106047in}}%
\pgfpathlineto{\pgfqpoint{1.311382in}{2.092799in}}%
\pgfpathlineto{\pgfqpoint{1.309063in}{2.079551in}}%
\pgfpathlineto{\pgfqpoint{1.305920in}{2.066303in}}%
\pgfpathlineto{\pgfqpoint{1.304776in}{2.053055in}}%
\pgfpathlineto{\pgfqpoint{1.301261in}{2.039807in}}%
\pgfpathlineto{\pgfqpoint{1.300472in}{2.026559in}}%
\pgfpathlineto{\pgfqpoint{1.299470in}{2.022879in}}%
\pgfpathlineto{\pgfqpoint{1.295837in}{2.013311in}}%
\pgfpathlineto{\pgfqpoint{1.293750in}{2.000063in}}%
\pgfpathlineto{\pgfqpoint{1.290261in}{1.986815in}}%
\pgfpathlineto{\pgfqpoint{1.287155in}{1.973567in}}%
\pgfpathlineto{\pgfqpoint{1.285579in}{1.960319in}}%
\pgfpathlineto{\pgfqpoint{1.281478in}{1.947071in}}%
\pgfpathlineto{\pgfqpoint{1.280628in}{1.933823in}}%
\pgfpathlineto{\pgfqpoint{1.276777in}{1.920575in}}%
\pgfpathlineto{\pgfqpoint{1.275660in}{1.907327in}}%
\pgfpathlineto{\pgfqpoint{1.274187in}{1.900937in}}%
\pgfpathlineto{\pgfqpoint{1.271943in}{1.894079in}}%
\pgfpathlineto{\pgfqpoint{1.268691in}{1.880831in}}%
\pgfpathlineto{\pgfqpoint{1.266208in}{1.867583in}}%
\pgfpathlineto{\pgfqpoint{1.261975in}{1.854335in}}%
\pgfpathlineto{\pgfqpoint{1.261041in}{1.841087in}}%
\pgfpathlineto{\pgfqpoint{1.256360in}{1.827839in}}%
\pgfpathlineto{\pgfqpoint{1.255267in}{1.814591in}}%
\pgfpathlineto{\pgfqpoint{1.251435in}{1.801343in}}%
\pgfpathlineto{\pgfqpoint{1.249366in}{1.788095in}}%
\pgfpathlineto{\pgfqpoint{1.248904in}{1.785466in}}%
\pgfpathlineto{\pgfqpoint{1.246278in}{1.774847in}}%
\pgfpathlineto{\pgfqpoint{1.241503in}{1.761599in}}%
\pgfpathlineto{\pgfqpoint{1.240087in}{1.748351in}}%
\pgfpathlineto{\pgfqpoint{1.234943in}{1.735103in}}%
\pgfpathlineto{\pgfqpoint{1.234207in}{1.721855in}}%
\pgfpathlineto{\pgfqpoint{1.229739in}{1.708607in}}%
\pgfpathlineto{\pgfqpoint{1.227773in}{1.695359in}}%
\pgfpathlineto{\pgfqpoint{1.224654in}{1.682111in}}%
\pgfpathlineto{\pgfqpoint{1.223620in}{1.678311in}}%
\pgfpathlineto{\pgfqpoint{1.220278in}{1.668863in}}%
\pgfpathlineto{\pgfqpoint{1.218337in}{1.655615in}}%
\pgfpathlineto{\pgfqpoint{1.212211in}{1.642367in}}%
\pgfpathlineto{\pgfqpoint{1.211224in}{1.629119in}}%
\pgfpathlineto{\pgfqpoint{1.205733in}{1.615871in}}%
\pgfpathlineto{\pgfqpoint{1.204281in}{1.602623in}}%
\pgfpathlineto{\pgfqpoint{1.200781in}{1.589375in}}%
\pgfpathlineto{\pgfqpoint{1.198337in}{1.580311in}}%
\pgfpathlineto{\pgfqpoint{1.196894in}{1.576127in}}%
\pgfpathlineto{\pgfqpoint{1.194455in}{1.562879in}}%
\pgfpathlineto{\pgfqpoint{1.188126in}{1.549631in}}%
\pgfpathlineto{\pgfqpoint{1.187210in}{1.536383in}}%
\pgfpathlineto{\pgfqpoint{1.180666in}{1.523135in}}%
\pgfpathlineto{\pgfqpoint{1.179088in}{1.509887in}}%
\pgfpathlineto{\pgfqpoint{1.174271in}{1.496639in}}%
\pgfpathlineto{\pgfqpoint{1.173054in}{1.491754in}}%
\pgfpathlineto{\pgfqpoint{1.170472in}{1.483391in}}%
\pgfpathlineto{\pgfqpoint{1.167833in}{1.470143in}}%
\pgfpathlineto{\pgfqpoint{1.161073in}{1.456895in}}%
\pgfpathlineto{\pgfqpoint{1.160035in}{1.443647in}}%
\pgfpathlineto{\pgfqpoint{1.152753in}{1.430399in}}%
\pgfpathlineto{\pgfqpoint{1.151790in}{1.417151in}}%
\pgfpathlineto{\pgfqpoint{1.147771in}{1.407711in}}%
\pgfpathlineto{\pgfqpoint{1.145774in}{1.403903in}}%
\pgfpathlineto{\pgfqpoint{1.141725in}{1.390655in}}%
\pgfpathlineto{\pgfqpoint{1.138117in}{1.377407in}}%
\pgfpathlineto{\pgfqpoint{1.131463in}{1.364159in}}%
\pgfpathlineto{\pgfqpoint{1.130672in}{1.350911in}}%
\pgfpathlineto{\pgfqpoint{1.122491in}{1.337663in}}%
\pgfpathlineto{\pgfqpoint{1.122488in}{1.337630in}}%
\pgfpathlineto{\pgfqpoint{1.121060in}{1.324415in}}%
\pgfpathlineto{\pgfqpoint{1.112573in}{1.311167in}}%
\pgfpathlineto{\pgfqpoint{1.109411in}{1.297919in}}%
\pgfpathlineto{\pgfqpoint{1.105358in}{1.284671in}}%
\pgfpathlineto{\pgfqpoint{1.098426in}{1.271423in}}%
\pgfpathlineto{\pgfqpoint{1.097628in}{1.258175in}}%
\pgfpathlineto{\pgfqpoint{1.097204in}{1.257564in}}%
\pgfpathlineto{\pgfqpoint{1.087358in}{1.244927in}}%
\pgfpathlineto{\pgfqpoint{1.087427in}{1.231679in}}%
\pgfpathlineto{\pgfqpoint{1.078375in}{1.218431in}}%
\pgfpathlineto{\pgfqpoint{1.074844in}{1.205183in}}%
\pgfpathlineto{\pgfqpoint{1.071921in}{1.198369in}}%
\pgfpathlineto{\pgfqpoint{1.068656in}{1.191935in}}%
\pgfpathlineto{\pgfqpoint{1.060535in}{1.178687in}}%
\pgfpathlineto{\pgfqpoint{1.059745in}{1.165439in}}%
\pgfpathlineto{\pgfqpoint{1.048490in}{1.152191in}}%
\pgfpathlineto{\pgfqpoint{1.049425in}{1.138943in}}%
\pgfpathlineto{\pgfqpoint{1.046638in}{1.135469in}}%
\pgfpathlineto{\pgfqpoint{1.037816in}{1.125695in}}%
\pgfpathlineto{\pgfqpoint{1.036285in}{1.112447in}}%
\pgfpathlineto{\pgfqpoint{1.030302in}{1.099199in}}%
\pgfpathlineto{\pgfqpoint{1.021604in}{1.085951in}}%
\pgfpathlineto{\pgfqpoint{1.021355in}{1.082197in}}%
\pgfpathlineto{\pgfqpoint{1.020620in}{1.072703in}}%
\pgfpathlineto{\pgfqpoint{1.005900in}{1.059455in}}%
\pgfpathlineto{\pgfqpoint{1.009062in}{1.046207in}}%
\pgfpathlineto{\pgfqpoint{0.996072in}{1.033160in}}%
\pgfpathlineto{\pgfqpoint{0.995853in}{1.032959in}}%
\pgfpathlineto{\pgfqpoint{0.995153in}{1.019711in}}%
\pgfpathlineto{\pgfqpoint{0.987589in}{1.006463in}}%
\pgfpathlineto{\pgfqpoint{0.979190in}{0.993215in}}%
\pgfpathlineto{\pgfqpoint{0.982377in}{0.979967in}}%
\pgfpathlineto{\pgfqpoint{0.970789in}{0.970514in}}%
\pgfpathlineto{\pgfqpoint{0.965899in}{0.966719in}}%
\pgfpathlineto{\pgfqpoint{0.970789in}{0.957511in}}%
\pgfpathlineto{\pgfqpoint{0.972995in}{0.953471in}}%
\pgfpathlineto{\pgfqpoint{0.970789in}{0.951650in}}%
\pgfpathlineto{\pgfqpoint{0.955619in}{0.940223in}}%
\pgfpathlineto{\pgfqpoint{0.960136in}{0.926975in}}%
\pgfpathlineto{\pgfqpoint{0.956379in}{0.913727in}}%
\pgfpathlineto{\pgfqpoint{0.947766in}{0.900479in}}%
\pgfpathlineto{\pgfqpoint{0.959116in}{0.887231in}}%
\pgfpathlineto{\pgfqpoint{0.945505in}{0.878101in}}%
\pgfpathlineto{\pgfqpoint{0.939026in}{0.873983in}}%
\pgfpathlineto{\pgfqpoint{0.945505in}{0.870049in}}%
\pgfpathlineto{\pgfqpoint{0.962470in}{0.860735in}}%
\pgfpathlineto{\pgfqpoint{0.948664in}{0.847487in}}%
\pgfpathlineto{\pgfqpoint{0.963464in}{0.834239in}}%
\pgfpathlineto{\pgfqpoint{0.970789in}{0.826122in}}%
\pgfpathlineto{\pgfqpoint{0.976444in}{0.820991in}}%
\pgfpathlineto{\pgfqpoint{0.972846in}{0.807743in}}%
\pgfpathlineto{\pgfqpoint{0.996072in}{0.802002in}}%
\pgfpathlineto{\pgfqpoint{1.021355in}{0.796779in}}%
\pgfpathlineto{\pgfqpoint{1.034081in}{0.794495in}}%
\pgfpathlineto{\pgfqpoint{1.025419in}{0.781247in}}%
\pgfpathlineto{\pgfqpoint{1.046638in}{0.777622in}}%
\pgfpathlineto{\pgfqpoint{1.071921in}{0.773723in}}%
\pgfpathlineto{\pgfqpoint{1.097204in}{0.770210in}}%
\pgfpathlineto{\pgfqpoint{1.114660in}{0.767999in}}%
\pgfpathlineto{\pgfqpoint{1.122488in}{0.765568in}}%
\pgfpathlineto{\pgfqpoint{1.147771in}{0.759452in}}%
\pgfpathlineto{\pgfqpoint{1.173054in}{0.755096in}}%
\pgfpathlineto{\pgfqpoint{1.175560in}{0.754751in}}%
\pgfpathlineto{\pgfqpoint{1.198337in}{0.751885in}}%
\pgfpathlineto{\pgfqpoint{1.223620in}{0.749000in}}%
\pgfpathlineto{\pgfqpoint{1.248904in}{0.746399in}}%
\pgfpathlineto{\pgfqpoint{1.274187in}{0.744073in}}%
\pgfpathlineto{\pgfqpoint{1.299470in}{0.742019in}}%
\pgfpathlineto{\pgfqpoint{1.306859in}{0.741503in}}%
\pgfpathlineto{\pgfqpoint{1.324753in}{0.740954in}}%
\pgfpathclose%
\pgfpathmoveto{\pgfqpoint{0.970789in}{1.693901in}}%
\pgfpathlineto{\pgfqpoint{0.996072in}{1.695351in}}%
\pgfpathlineto{\pgfqpoint{0.996087in}{1.695359in}}%
\pgfpathlineto{\pgfqpoint{0.996072in}{1.695457in}}%
\pgfpathlineto{\pgfqpoint{0.970789in}{1.700892in}}%
\pgfpathlineto{\pgfqpoint{0.968439in}{1.695359in}}%
\pgfpathlineto{\pgfqpoint{0.970789in}{1.693901in}}%
\pgfpathclose%
\pgfpathmoveto{\pgfqpoint{0.996072in}{1.711113in}}%
\pgfpathlineto{\pgfqpoint{1.010185in}{1.721855in}}%
\pgfpathlineto{\pgfqpoint{1.021355in}{1.729018in}}%
\pgfpathlineto{\pgfqpoint{1.025088in}{1.735103in}}%
\pgfpathlineto{\pgfqpoint{1.027811in}{1.748351in}}%
\pgfpathlineto{\pgfqpoint{1.046638in}{1.761585in}}%
\pgfpathlineto{\pgfqpoint{1.046650in}{1.761599in}}%
\pgfpathlineto{\pgfqpoint{1.046638in}{1.761747in}}%
\pgfpathlineto{\pgfqpoint{1.044189in}{1.774847in}}%
\pgfpathlineto{\pgfqpoint{1.046638in}{1.775859in}}%
\pgfpathlineto{\pgfqpoint{1.061689in}{1.788095in}}%
\pgfpathlineto{\pgfqpoint{1.062332in}{1.801343in}}%
\pgfpathlineto{\pgfqpoint{1.071921in}{1.809475in}}%
\pgfpathlineto{\pgfqpoint{1.076064in}{1.814591in}}%
\pgfpathlineto{\pgfqpoint{1.079832in}{1.827839in}}%
\pgfpathlineto{\pgfqpoint{1.085395in}{1.841087in}}%
\pgfpathlineto{\pgfqpoint{1.095494in}{1.854335in}}%
\pgfpathlineto{\pgfqpoint{1.097190in}{1.867583in}}%
\pgfpathlineto{\pgfqpoint{1.097204in}{1.867595in}}%
\pgfpathlineto{\pgfqpoint{1.107209in}{1.880831in}}%
\pgfpathlineto{\pgfqpoint{1.107170in}{1.894079in}}%
\pgfpathlineto{\pgfqpoint{1.118135in}{1.907327in}}%
\pgfpathlineto{\pgfqpoint{1.121662in}{1.920575in}}%
\pgfpathlineto{\pgfqpoint{1.122488in}{1.921891in}}%
\pgfpathlineto{\pgfqpoint{1.128354in}{1.933823in}}%
\pgfpathlineto{\pgfqpoint{1.133795in}{1.947071in}}%
\pgfpathlineto{\pgfqpoint{1.136344in}{1.960319in}}%
\pgfpathlineto{\pgfqpoint{1.145249in}{1.973567in}}%
\pgfpathlineto{\pgfqpoint{1.146154in}{1.986815in}}%
\pgfpathlineto{\pgfqpoint{1.147771in}{1.988899in}}%
\pgfpathlineto{\pgfqpoint{1.154606in}{2.000063in}}%
\pgfpathlineto{\pgfqpoint{1.155753in}{2.013311in}}%
\pgfpathlineto{\pgfqpoint{1.162025in}{2.026559in}}%
\pgfpathlineto{\pgfqpoint{1.166329in}{2.039807in}}%
\pgfpathlineto{\pgfqpoint{1.170294in}{2.053055in}}%
\pgfpathlineto{\pgfqpoint{1.173054in}{2.058273in}}%
\pgfpathlineto{\pgfqpoint{1.176578in}{2.066303in}}%
\pgfpathlineto{\pgfqpoint{1.177693in}{2.079551in}}%
\pgfpathlineto{\pgfqpoint{1.184854in}{2.092799in}}%
\pgfpathlineto{\pgfqpoint{1.186008in}{2.106047in}}%
\pgfpathlineto{\pgfqpoint{1.192580in}{2.119295in}}%
\pgfpathlineto{\pgfqpoint{1.195362in}{2.132543in}}%
\pgfpathlineto{\pgfqpoint{1.176175in}{2.132543in}}%
\pgfpathlineto{\pgfqpoint{1.173083in}{2.119295in}}%
\pgfpathlineto{\pgfqpoint{1.173054in}{2.119237in}}%
\pgfpathlineto{\pgfqpoint{1.164483in}{2.106047in}}%
\pgfpathlineto{\pgfqpoint{1.163478in}{2.092799in}}%
\pgfpathlineto{\pgfqpoint{1.154251in}{2.079551in}}%
\pgfpathlineto{\pgfqpoint{1.154255in}{2.066303in}}%
\pgfpathlineto{\pgfqpoint{1.147771in}{2.055038in}}%
\pgfpathlineto{\pgfqpoint{1.145813in}{2.053055in}}%
\pgfpathlineto{\pgfqpoint{1.142056in}{2.039807in}}%
\pgfpathlineto{\pgfqpoint{1.134906in}{2.026559in}}%
\pgfpathlineto{\pgfqpoint{1.128071in}{2.013311in}}%
\pgfpathlineto{\pgfqpoint{1.127245in}{2.000063in}}%
\pgfpathlineto{\pgfqpoint{1.122488in}{1.993608in}}%
\pgfpathlineto{\pgfqpoint{1.113480in}{1.986815in}}%
\pgfpathlineto{\pgfqpoint{1.115659in}{1.973567in}}%
\pgfpathlineto{\pgfqpoint{1.100789in}{1.960319in}}%
\pgfpathlineto{\pgfqpoint{1.101376in}{1.947071in}}%
\pgfpathlineto{\pgfqpoint{1.097204in}{1.940615in}}%
\pgfpathlineto{\pgfqpoint{1.084902in}{1.933823in}}%
\pgfpathlineto{\pgfqpoint{1.077206in}{1.920575in}}%
\pgfpathlineto{\pgfqpoint{1.071921in}{1.907871in}}%
\pgfpathlineto{\pgfqpoint{1.050834in}{1.920575in}}%
\pgfpathlineto{\pgfqpoint{1.057443in}{1.933823in}}%
\pgfpathlineto{\pgfqpoint{1.046638in}{1.941753in}}%
\pgfpathlineto{\pgfqpoint{1.044271in}{1.947071in}}%
\pgfpathlineto{\pgfqpoint{1.046638in}{1.950730in}}%
\pgfpathlineto{\pgfqpoint{1.058177in}{1.960319in}}%
\pgfpathlineto{\pgfqpoint{1.047599in}{1.973567in}}%
\pgfpathlineto{\pgfqpoint{1.057449in}{1.986815in}}%
\pgfpathlineto{\pgfqpoint{1.056020in}{2.000063in}}%
\pgfpathlineto{\pgfqpoint{1.058448in}{2.013311in}}%
\pgfpathlineto{\pgfqpoint{1.063098in}{2.026559in}}%
\pgfpathlineto{\pgfqpoint{1.059735in}{2.039807in}}%
\pgfpathlineto{\pgfqpoint{1.066524in}{2.053055in}}%
\pgfpathlineto{\pgfqpoint{1.063173in}{2.066303in}}%
\pgfpathlineto{\pgfqpoint{1.069098in}{2.079551in}}%
\pgfpathlineto{\pgfqpoint{1.068954in}{2.092799in}}%
\pgfpathlineto{\pgfqpoint{1.071076in}{2.106047in}}%
\pgfpathlineto{\pgfqpoint{1.071921in}{2.108695in}}%
\pgfpathlineto{\pgfqpoint{1.076017in}{2.119295in}}%
\pgfpathlineto{\pgfqpoint{1.075415in}{2.132543in}}%
\pgfpathlineto{\pgfqpoint{1.071921in}{2.132543in}}%
\pgfpathlineto{\pgfqpoint{1.054292in}{2.132543in}}%
\pgfpathlineto{\pgfqpoint{1.053963in}{2.119295in}}%
\pgfpathlineto{\pgfqpoint{1.048567in}{2.106047in}}%
\pgfpathlineto{\pgfqpoint{1.046638in}{2.094869in}}%
\pgfpathlineto{\pgfqpoint{1.046297in}{2.092799in}}%
\pgfpathlineto{\pgfqpoint{1.044151in}{2.079551in}}%
\pgfpathlineto{\pgfqpoint{1.038869in}{2.066303in}}%
\pgfpathlineto{\pgfqpoint{1.039456in}{2.053055in}}%
\pgfpathlineto{\pgfqpoint{1.033038in}{2.039807in}}%
\pgfpathlineto{\pgfqpoint{1.034005in}{2.026559in}}%
\pgfpathlineto{\pgfqpoint{1.027943in}{2.013311in}}%
\pgfpathlineto{\pgfqpoint{1.025727in}{2.000063in}}%
\pgfpathlineto{\pgfqpoint{1.022316in}{1.986815in}}%
\pgfpathlineto{\pgfqpoint{1.021355in}{1.984769in}}%
\pgfpathlineto{\pgfqpoint{1.016930in}{1.973567in}}%
\pgfpathlineto{\pgfqpoint{1.017716in}{1.960319in}}%
\pgfpathlineto{\pgfqpoint{1.009804in}{1.947071in}}%
\pgfpathlineto{\pgfqpoint{1.011716in}{1.933823in}}%
\pgfpathlineto{\pgfqpoint{1.004860in}{1.920575in}}%
\pgfpathlineto{\pgfqpoint{1.004634in}{1.907327in}}%
\pgfpathlineto{\pgfqpoint{1.001627in}{1.894079in}}%
\pgfpathlineto{\pgfqpoint{0.996072in}{1.883431in}}%
\pgfpathlineto{\pgfqpoint{0.994996in}{1.880831in}}%
\pgfpathlineto{\pgfqpoint{0.996072in}{1.869532in}}%
\pgfpathlineto{\pgfqpoint{0.996358in}{1.867583in}}%
\pgfpathlineto{\pgfqpoint{0.996072in}{1.867267in}}%
\pgfpathlineto{\pgfqpoint{0.987476in}{1.854335in}}%
\pgfpathlineto{\pgfqpoint{0.991384in}{1.841087in}}%
\pgfpathlineto{\pgfqpoint{0.982688in}{1.827839in}}%
\pgfpathlineto{\pgfqpoint{0.983584in}{1.814591in}}%
\pgfpathlineto{\pgfqpoint{0.980598in}{1.801343in}}%
\pgfpathlineto{\pgfqpoint{0.974633in}{1.788095in}}%
\pgfpathlineto{\pgfqpoint{0.981349in}{1.774847in}}%
\pgfpathlineto{\pgfqpoint{0.970789in}{1.766530in}}%
\pgfpathlineto{\pgfqpoint{0.966765in}{1.761599in}}%
\pgfpathlineto{\pgfqpoint{0.970789in}{1.754360in}}%
\pgfpathlineto{\pgfqpoint{0.978595in}{1.748351in}}%
\pgfpathlineto{\pgfqpoint{0.970789in}{1.742724in}}%
\pgfpathlineto{\pgfqpoint{0.964688in}{1.735103in}}%
\pgfpathlineto{\pgfqpoint{0.970789in}{1.723174in}}%
\pgfpathlineto{\pgfqpoint{0.973208in}{1.721855in}}%
\pgfpathlineto{\pgfqpoint{0.996072in}{1.711113in}}%
\pgfpathclose%
\pgfpathmoveto{\pgfqpoint{1.906266in}{1.720517in}}%
\pgfpathlineto{\pgfqpoint{1.916527in}{1.721855in}}%
\pgfpathlineto{\pgfqpoint{1.906266in}{1.725492in}}%
\pgfpathlineto{\pgfqpoint{1.903353in}{1.721855in}}%
\pgfpathlineto{\pgfqpoint{1.906266in}{1.720517in}}%
\pgfpathclose%
\pgfpathmoveto{\pgfqpoint{1.906266in}{1.740005in}}%
\pgfpathlineto{\pgfqpoint{1.922570in}{1.748351in}}%
\pgfpathlineto{\pgfqpoint{1.912625in}{1.761599in}}%
\pgfpathlineto{\pgfqpoint{1.916377in}{1.774847in}}%
\pgfpathlineto{\pgfqpoint{1.918627in}{1.788095in}}%
\pgfpathlineto{\pgfqpoint{1.913117in}{1.801343in}}%
\pgfpathlineto{\pgfqpoint{1.918604in}{1.814591in}}%
\pgfpathlineto{\pgfqpoint{1.908940in}{1.827839in}}%
\pgfpathlineto{\pgfqpoint{1.913784in}{1.841087in}}%
\pgfpathlineto{\pgfqpoint{1.906801in}{1.854335in}}%
\pgfpathlineto{\pgfqpoint{1.907432in}{1.867583in}}%
\pgfpathlineto{\pgfqpoint{1.906266in}{1.871979in}}%
\pgfpathlineto{\pgfqpoint{1.903428in}{1.880831in}}%
\pgfpathlineto{\pgfqpoint{1.897936in}{1.894079in}}%
\pgfpathlineto{\pgfqpoint{1.898993in}{1.907327in}}%
\pgfpathlineto{\pgfqpoint{1.891395in}{1.920575in}}%
\pgfpathlineto{\pgfqpoint{1.893379in}{1.933823in}}%
\pgfpathlineto{\pgfqpoint{1.886003in}{1.947071in}}%
\pgfpathlineto{\pgfqpoint{1.886122in}{1.960319in}}%
\pgfpathlineto{\pgfqpoint{1.882148in}{1.973567in}}%
\pgfpathlineto{\pgfqpoint{1.880983in}{1.979236in}}%
\pgfpathlineto{\pgfqpoint{1.879141in}{1.986815in}}%
\pgfpathlineto{\pgfqpoint{1.878296in}{2.000063in}}%
\pgfpathlineto{\pgfqpoint{1.871748in}{2.013311in}}%
\pgfpathlineto{\pgfqpoint{1.872891in}{2.026559in}}%
\pgfpathlineto{\pgfqpoint{1.866352in}{2.039807in}}%
\pgfpathlineto{\pgfqpoint{1.866204in}{2.053055in}}%
\pgfpathlineto{\pgfqpoint{1.861347in}{2.066303in}}%
\pgfpathlineto{\pgfqpoint{1.858364in}{2.079551in}}%
\pgfpathlineto{\pgfqpoint{1.856442in}{2.092799in}}%
\pgfpathlineto{\pgfqpoint{1.855700in}{2.094889in}}%
\pgfpathlineto{\pgfqpoint{1.851444in}{2.106047in}}%
\pgfpathlineto{\pgfqpoint{1.851665in}{2.119295in}}%
\pgfpathlineto{\pgfqpoint{1.845773in}{2.132543in}}%
\pgfpathlineto{\pgfqpoint{1.830416in}{2.132543in}}%
\pgfpathlineto{\pgfqpoint{1.822412in}{2.132543in}}%
\pgfpathlineto{\pgfqpoint{1.829251in}{2.119295in}}%
\pgfpathlineto{\pgfqpoint{1.826187in}{2.106047in}}%
\pgfpathlineto{\pgfqpoint{1.830416in}{2.097918in}}%
\pgfpathlineto{\pgfqpoint{1.832279in}{2.092799in}}%
\pgfpathlineto{\pgfqpoint{1.831739in}{2.079551in}}%
\pgfpathlineto{\pgfqpoint{1.834739in}{2.066303in}}%
\pgfpathlineto{\pgfqpoint{1.838021in}{2.053055in}}%
\pgfpathlineto{\pgfqpoint{1.836351in}{2.039807in}}%
\pgfpathlineto{\pgfqpoint{1.843074in}{2.026559in}}%
\pgfpathlineto{\pgfqpoint{1.837160in}{2.013311in}}%
\pgfpathlineto{\pgfqpoint{1.845411in}{2.000063in}}%
\pgfpathlineto{\pgfqpoint{1.839201in}{1.986815in}}%
\pgfpathlineto{\pgfqpoint{1.842410in}{1.973567in}}%
\pgfpathlineto{\pgfqpoint{1.839829in}{1.960319in}}%
\pgfpathlineto{\pgfqpoint{1.830944in}{1.947071in}}%
\pgfpathlineto{\pgfqpoint{1.837429in}{1.933823in}}%
\pgfpathlineto{\pgfqpoint{1.830416in}{1.931778in}}%
\pgfpathlineto{\pgfqpoint{1.827464in}{1.933823in}}%
\pgfpathlineto{\pgfqpoint{1.829821in}{1.947071in}}%
\pgfpathlineto{\pgfqpoint{1.815682in}{1.960319in}}%
\pgfpathlineto{\pgfqpoint{1.805133in}{1.965314in}}%
\pgfpathlineto{\pgfqpoint{1.799407in}{1.973567in}}%
\pgfpathlineto{\pgfqpoint{1.798839in}{1.986815in}}%
\pgfpathlineto{\pgfqpoint{1.782779in}{2.000063in}}%
\pgfpathlineto{\pgfqpoint{1.784471in}{2.013311in}}%
\pgfpathlineto{\pgfqpoint{1.779850in}{2.016322in}}%
\pgfpathlineto{\pgfqpoint{1.772784in}{2.026559in}}%
\pgfpathlineto{\pgfqpoint{1.769977in}{2.039807in}}%
\pgfpathlineto{\pgfqpoint{1.764120in}{2.053055in}}%
\pgfpathlineto{\pgfqpoint{1.755220in}{2.066303in}}%
\pgfpathlineto{\pgfqpoint{1.754567in}{2.074740in}}%
\pgfpathlineto{\pgfqpoint{1.754340in}{2.079551in}}%
\pgfpathlineto{\pgfqpoint{1.745606in}{2.092799in}}%
\pgfpathlineto{\pgfqpoint{1.745008in}{2.106047in}}%
\pgfpathlineto{\pgfqpoint{1.736484in}{2.119295in}}%
\pgfpathlineto{\pgfqpoint{1.732874in}{2.132543in}}%
\pgfpathlineto{\pgfqpoint{1.729284in}{2.132543in}}%
\pgfpathlineto{\pgfqpoint{1.711835in}{2.132543in}}%
\pgfpathlineto{\pgfqpoint{1.715192in}{2.119295in}}%
\pgfpathlineto{\pgfqpoint{1.721184in}{2.106047in}}%
\pgfpathlineto{\pgfqpoint{1.722622in}{2.092799in}}%
\pgfpathlineto{\pgfqpoint{1.729284in}{2.080665in}}%
\pgfpathlineto{\pgfqpoint{1.730030in}{2.079551in}}%
\pgfpathlineto{\pgfqpoint{1.731328in}{2.066303in}}%
\pgfpathlineto{\pgfqpoint{1.737815in}{2.053055in}}%
\pgfpathlineto{\pgfqpoint{1.741706in}{2.039807in}}%
\pgfpathlineto{\pgfqpoint{1.745491in}{2.026559in}}%
\pgfpathlineto{\pgfqpoint{1.752019in}{2.013311in}}%
\pgfpathlineto{\pgfqpoint{1.753161in}{2.000063in}}%
\pgfpathlineto{\pgfqpoint{1.754567in}{1.997633in}}%
\pgfpathlineto{\pgfqpoint{1.762327in}{1.986815in}}%
\pgfpathlineto{\pgfqpoint{1.763845in}{1.973567in}}%
\pgfpathlineto{\pgfqpoint{1.771702in}{1.960319in}}%
\pgfpathlineto{\pgfqpoint{1.775482in}{1.947071in}}%
\pgfpathlineto{\pgfqpoint{1.779598in}{1.933823in}}%
\pgfpathlineto{\pgfqpoint{1.779850in}{1.933366in}}%
\pgfpathlineto{\pgfqpoint{1.789300in}{1.920575in}}%
\pgfpathlineto{\pgfqpoint{1.790852in}{1.907327in}}%
\pgfpathlineto{\pgfqpoint{1.801784in}{1.894079in}}%
\pgfpathlineto{\pgfqpoint{1.802220in}{1.880831in}}%
\pgfpathlineto{\pgfqpoint{1.805133in}{1.875994in}}%
\pgfpathlineto{\pgfqpoint{1.812171in}{1.867583in}}%
\pgfpathlineto{\pgfqpoint{1.817537in}{1.854335in}}%
\pgfpathlineto{\pgfqpoint{1.822459in}{1.841087in}}%
\pgfpathlineto{\pgfqpoint{1.830416in}{1.830263in}}%
\pgfpathlineto{\pgfqpoint{1.833610in}{1.827839in}}%
\pgfpathlineto{\pgfqpoint{1.834197in}{1.814591in}}%
\pgfpathlineto{\pgfqpoint{1.850855in}{1.801343in}}%
\pgfpathlineto{\pgfqpoint{1.851654in}{1.788095in}}%
\pgfpathlineto{\pgfqpoint{1.855700in}{1.783845in}}%
\pgfpathlineto{\pgfqpoint{1.870577in}{1.774847in}}%
\pgfpathlineto{\pgfqpoint{1.877628in}{1.761599in}}%
\pgfpathlineto{\pgfqpoint{1.880983in}{1.748633in}}%
\pgfpathlineto{\pgfqpoint{1.881189in}{1.748351in}}%
\pgfpathlineto{\pgfqpoint{1.906266in}{1.740005in}}%
\pgfpathclose%
\pgfusepath{fill}%
\end{pgfscope}%
\begin{pgfscope}%
\pgfpathrectangle{\pgfqpoint{0.540975in}{0.476543in}}{\pgfqpoint{1.820388in}{1.656000in}}%
\pgfusepath{clip}%
\pgfsetbuttcap%
\pgfsetroundjoin%
\definecolor{currentfill}{rgb}{0.433218,0.652595,0.806228}%
\pgfsetfillcolor{currentfill}%
\pgfsetlinewidth{0.000000pt}%
\definecolor{currentstroke}{rgb}{0.000000,0.000000,0.000000}%
\pgfsetstrokecolor{currentstroke}%
\pgfsetdash{}{0pt}%
\pgfpathmoveto{\pgfqpoint{1.324753in}{0.794407in}}%
\pgfpathlineto{\pgfqpoint{1.350036in}{0.790835in}}%
\pgfpathlineto{\pgfqpoint{1.375319in}{0.788536in}}%
\pgfpathlineto{\pgfqpoint{1.400603in}{0.787216in}}%
\pgfpathlineto{\pgfqpoint{1.425886in}{0.786675in}}%
\pgfpathlineto{\pgfqpoint{1.451169in}{0.786778in}}%
\pgfpathlineto{\pgfqpoint{1.476452in}{0.787431in}}%
\pgfpathlineto{\pgfqpoint{1.501735in}{0.788570in}}%
\pgfpathlineto{\pgfqpoint{1.527018in}{0.790154in}}%
\pgfpathlineto{\pgfqpoint{1.552302in}{0.792159in}}%
\pgfpathlineto{\pgfqpoint{1.576748in}{0.794495in}}%
\pgfpathlineto{\pgfqpoint{1.577585in}{0.794536in}}%
\pgfpathlineto{\pgfqpoint{1.602868in}{0.796086in}}%
\pgfpathlineto{\pgfqpoint{1.628151in}{0.798084in}}%
\pgfpathlineto{\pgfqpoint{1.653434in}{0.800621in}}%
\pgfpathlineto{\pgfqpoint{1.678717in}{0.803819in}}%
\pgfpathlineto{\pgfqpoint{1.703435in}{0.807743in}}%
\pgfpathlineto{\pgfqpoint{1.704001in}{0.808143in}}%
\pgfpathlineto{\pgfqpoint{1.720309in}{0.820991in}}%
\pgfpathlineto{\pgfqpoint{1.729284in}{0.822353in}}%
\pgfpathlineto{\pgfqpoint{1.754567in}{0.826813in}}%
\pgfpathlineto{\pgfqpoint{1.779850in}{0.832241in}}%
\pgfpathlineto{\pgfqpoint{1.787895in}{0.834239in}}%
\pgfpathlineto{\pgfqpoint{1.782333in}{0.847487in}}%
\pgfpathlineto{\pgfqpoint{1.805133in}{0.853269in}}%
\pgfpathlineto{\pgfqpoint{1.830194in}{0.860735in}}%
\pgfpathlineto{\pgfqpoint{1.827245in}{0.873983in}}%
\pgfpathlineto{\pgfqpoint{1.830416in}{0.875697in}}%
\pgfpathlineto{\pgfqpoint{1.849859in}{0.887231in}}%
\pgfpathlineto{\pgfqpoint{1.852291in}{0.900479in}}%
\pgfpathlineto{\pgfqpoint{1.854770in}{0.913727in}}%
\pgfpathlineto{\pgfqpoint{1.855700in}{0.914761in}}%
\pgfpathlineto{\pgfqpoint{1.865350in}{0.926975in}}%
\pgfpathlineto{\pgfqpoint{1.856940in}{0.940223in}}%
\pgfpathlineto{\pgfqpoint{1.869505in}{0.953471in}}%
\pgfpathlineto{\pgfqpoint{1.856316in}{0.966719in}}%
\pgfpathlineto{\pgfqpoint{1.865917in}{0.979967in}}%
\pgfpathlineto{\pgfqpoint{1.857379in}{0.993215in}}%
\pgfpathlineto{\pgfqpoint{1.857806in}{1.006463in}}%
\pgfpathlineto{\pgfqpoint{1.855751in}{1.019711in}}%
\pgfpathlineto{\pgfqpoint{1.855700in}{1.019793in}}%
\pgfpathlineto{\pgfqpoint{1.847825in}{1.032959in}}%
\pgfpathlineto{\pgfqpoint{1.851596in}{1.046207in}}%
\pgfpathlineto{\pgfqpoint{1.841233in}{1.059455in}}%
\pgfpathlineto{\pgfqpoint{1.845471in}{1.072703in}}%
\pgfpathlineto{\pgfqpoint{1.835501in}{1.085951in}}%
\pgfpathlineto{\pgfqpoint{1.835807in}{1.099199in}}%
\pgfpathlineto{\pgfqpoint{1.830416in}{1.112147in}}%
\pgfpathlineto{\pgfqpoint{1.830306in}{1.112447in}}%
\pgfpathlineto{\pgfqpoint{1.824764in}{1.125695in}}%
\pgfpathlineto{\pgfqpoint{1.823689in}{1.138943in}}%
\pgfpathlineto{\pgfqpoint{1.813440in}{1.152191in}}%
\pgfpathlineto{\pgfqpoint{1.814985in}{1.165439in}}%
\pgfpathlineto{\pgfqpoint{1.805381in}{1.178687in}}%
\pgfpathlineto{\pgfqpoint{1.805366in}{1.191935in}}%
\pgfpathlineto{\pgfqpoint{1.805133in}{1.192391in}}%
\pgfpathlineto{\pgfqpoint{1.799442in}{1.205183in}}%
\pgfpathlineto{\pgfqpoint{1.795377in}{1.218431in}}%
\pgfpathlineto{\pgfqpoint{1.793413in}{1.231679in}}%
\pgfpathlineto{\pgfqpoint{1.785814in}{1.244927in}}%
\pgfpathlineto{\pgfqpoint{1.785645in}{1.258175in}}%
\pgfpathlineto{\pgfqpoint{1.779850in}{1.266245in}}%
\pgfpathlineto{\pgfqpoint{1.776528in}{1.271423in}}%
\pgfpathlineto{\pgfqpoint{1.776151in}{1.284671in}}%
\pgfpathlineto{\pgfqpoint{1.769724in}{1.297919in}}%
\pgfpathlineto{\pgfqpoint{1.766587in}{1.311167in}}%
\pgfpathlineto{\pgfqpoint{1.763468in}{1.324415in}}%
\pgfpathlineto{\pgfqpoint{1.756503in}{1.337663in}}%
\pgfpathlineto{\pgfqpoint{1.755994in}{1.350911in}}%
\pgfpathlineto{\pgfqpoint{1.754567in}{1.353285in}}%
\pgfpathlineto{\pgfqpoint{1.748874in}{1.364159in}}%
\pgfpathlineto{\pgfqpoint{1.748272in}{1.377407in}}%
\pgfpathlineto{\pgfqpoint{1.742160in}{1.390655in}}%
\pgfpathlineto{\pgfqpoint{1.739606in}{1.403903in}}%
\pgfpathlineto{\pgfqpoint{1.736070in}{1.417151in}}%
\pgfpathlineto{\pgfqpoint{1.730688in}{1.430399in}}%
\pgfpathlineto{\pgfqpoint{1.729385in}{1.443647in}}%
\pgfpathlineto{\pgfqpoint{1.729284in}{1.443820in}}%
\pgfpathlineto{\pgfqpoint{1.722790in}{1.456895in}}%
\pgfpathlineto{\pgfqpoint{1.722109in}{1.470143in}}%
\pgfpathlineto{\pgfqpoint{1.716326in}{1.483391in}}%
\pgfpathlineto{\pgfqpoint{1.714055in}{1.496639in}}%
\pgfpathlineto{\pgfqpoint{1.710348in}{1.509887in}}%
\pgfpathlineto{\pgfqpoint{1.705539in}{1.523135in}}%
\pgfpathlineto{\pgfqpoint{1.704181in}{1.536383in}}%
\pgfpathlineto{\pgfqpoint{1.704001in}{1.536771in}}%
\pgfpathlineto{\pgfqpoint{1.699061in}{1.549631in}}%
\pgfpathlineto{\pgfqpoint{1.698445in}{1.562879in}}%
\pgfpathlineto{\pgfqpoint{1.693037in}{1.576127in}}%
\pgfpathlineto{\pgfqpoint{1.691061in}{1.589375in}}%
\pgfpathlineto{\pgfqpoint{1.687190in}{1.602623in}}%
\pgfpathlineto{\pgfqpoint{1.683075in}{1.615871in}}%
\pgfpathlineto{\pgfqpoint{1.681162in}{1.629119in}}%
\pgfpathlineto{\pgfqpoint{1.678717in}{1.634506in}}%
\pgfpathlineto{\pgfqpoint{1.675822in}{1.642367in}}%
\pgfpathlineto{\pgfqpoint{1.675265in}{1.655615in}}%
\pgfpathlineto{\pgfqpoint{1.670703in}{1.668863in}}%
\pgfpathlineto{\pgfqpoint{1.669330in}{1.682111in}}%
\pgfpathlineto{\pgfqpoint{1.665839in}{1.695359in}}%
\pgfpathlineto{\pgfqpoint{1.662317in}{1.708607in}}%
\pgfpathlineto{\pgfqpoint{1.660381in}{1.721855in}}%
\pgfpathlineto{\pgfqpoint{1.655273in}{1.735103in}}%
\pgfpathlineto{\pgfqpoint{1.654153in}{1.748351in}}%
\pgfpathlineto{\pgfqpoint{1.653434in}{1.750018in}}%
\pgfpathlineto{\pgfqpoint{1.649543in}{1.761599in}}%
\pgfpathlineto{\pgfqpoint{1.648295in}{1.774847in}}%
\pgfpathlineto{\pgfqpoint{1.644987in}{1.788095in}}%
\pgfpathlineto{\pgfqpoint{1.642254in}{1.801343in}}%
\pgfpathlineto{\pgfqpoint{1.640428in}{1.814591in}}%
\pgfpathlineto{\pgfqpoint{1.635871in}{1.827839in}}%
\pgfpathlineto{\pgfqpoint{1.634925in}{1.841087in}}%
\pgfpathlineto{\pgfqpoint{1.630040in}{1.854335in}}%
\pgfpathlineto{\pgfqpoint{1.628568in}{1.867583in}}%
\pgfpathlineto{\pgfqpoint{1.628151in}{1.868907in}}%
\pgfpathlineto{\pgfqpoint{1.625327in}{1.880831in}}%
\pgfpathlineto{\pgfqpoint{1.622892in}{1.894079in}}%
\pgfpathlineto{\pgfqpoint{1.621048in}{1.907327in}}%
\pgfpathlineto{\pgfqpoint{1.617277in}{1.920575in}}%
\pgfpathlineto{\pgfqpoint{1.616246in}{1.933823in}}%
\pgfpathlineto{\pgfqpoint{1.611672in}{1.947071in}}%
\pgfpathlineto{\pgfqpoint{1.610353in}{1.960319in}}%
\pgfpathlineto{\pgfqpoint{1.606627in}{1.973567in}}%
\pgfpathlineto{\pgfqpoint{1.603968in}{1.986815in}}%
\pgfpathlineto{\pgfqpoint{1.602868in}{1.993174in}}%
\pgfpathlineto{\pgfqpoint{1.601985in}{2.000063in}}%
\pgfpathlineto{\pgfqpoint{1.598582in}{2.013311in}}%
\pgfpathlineto{\pgfqpoint{1.597702in}{2.026559in}}%
\pgfpathlineto{\pgfqpoint{1.593776in}{2.039807in}}%
\pgfpathlineto{\pgfqpoint{1.592558in}{2.053055in}}%
\pgfpathlineto{\pgfqpoint{1.588969in}{2.066303in}}%
\pgfpathlineto{\pgfqpoint{1.586363in}{2.079551in}}%
\pgfpathlineto{\pgfqpoint{1.584099in}{2.092799in}}%
\pgfpathlineto{\pgfqpoint{1.580180in}{2.106047in}}%
\pgfpathlineto{\pgfqpoint{1.579048in}{2.119295in}}%
\pgfpathlineto{\pgfqpoint{1.577585in}{2.123307in}}%
\pgfpathlineto{\pgfqpoint{1.574938in}{2.132543in}}%
\pgfpathlineto{\pgfqpoint{1.566151in}{2.132543in}}%
\pgfpathlineto{\pgfqpoint{1.570220in}{2.119295in}}%
\pgfpathlineto{\pgfqpoint{1.571072in}{2.106047in}}%
\pgfpathlineto{\pgfqpoint{1.574528in}{2.092799in}}%
\pgfpathlineto{\pgfqpoint{1.576397in}{2.079551in}}%
\pgfpathlineto{\pgfqpoint{1.577585in}{2.073211in}}%
\pgfpathlineto{\pgfqpoint{1.579136in}{2.066303in}}%
\pgfpathlineto{\pgfqpoint{1.583027in}{2.053055in}}%
\pgfpathlineto{\pgfqpoint{1.584396in}{2.039807in}}%
\pgfpathlineto{\pgfqpoint{1.588743in}{2.026559in}}%
\pgfpathlineto{\pgfqpoint{1.589591in}{2.013311in}}%
\pgfpathlineto{\pgfqpoint{1.593395in}{2.000063in}}%
\pgfpathlineto{\pgfqpoint{1.595065in}{1.986815in}}%
\pgfpathlineto{\pgfqpoint{1.597338in}{1.973567in}}%
\pgfpathlineto{\pgfqpoint{1.600370in}{1.960319in}}%
\pgfpathlineto{\pgfqpoint{1.601516in}{1.947071in}}%
\pgfpathlineto{\pgfqpoint{1.602868in}{1.942682in}}%
\pgfpathlineto{\pgfqpoint{1.606213in}{1.933823in}}%
\pgfpathlineto{\pgfqpoint{1.607235in}{1.920575in}}%
\pgfpathlineto{\pgfqpoint{1.611496in}{1.907327in}}%
\pgfpathlineto{\pgfqpoint{1.613319in}{1.894079in}}%
\pgfpathlineto{\pgfqpoint{1.616084in}{1.880831in}}%
\pgfpathlineto{\pgfqpoint{1.619325in}{1.867583in}}%
\pgfpathlineto{\pgfqpoint{1.620531in}{1.854335in}}%
\pgfpathlineto{\pgfqpoint{1.624575in}{1.841087in}}%
\pgfpathlineto{\pgfqpoint{1.625269in}{1.827839in}}%
\pgfpathlineto{\pgfqpoint{1.628151in}{1.818422in}}%
\pgfpathlineto{\pgfqpoint{1.629586in}{1.814591in}}%
\pgfpathlineto{\pgfqpoint{1.631390in}{1.801343in}}%
\pgfpathlineto{\pgfqpoint{1.634486in}{1.788095in}}%
\pgfpathlineto{\pgfqpoint{1.637917in}{1.774847in}}%
\pgfpathlineto{\pgfqpoint{1.639276in}{1.761599in}}%
\pgfpathlineto{\pgfqpoint{1.644002in}{1.748351in}}%
\pgfpathlineto{\pgfqpoint{1.644790in}{1.735103in}}%
\pgfpathlineto{\pgfqpoint{1.649173in}{1.721855in}}%
\pgfpathlineto{\pgfqpoint{1.650630in}{1.708607in}}%
\pgfpathlineto{\pgfqpoint{1.653434in}{1.696982in}}%
\pgfpathlineto{\pgfqpoint{1.653909in}{1.695359in}}%
\pgfpathlineto{\pgfqpoint{1.657515in}{1.682111in}}%
\pgfpathlineto{\pgfqpoint{1.658997in}{1.668863in}}%
\pgfpathlineto{\pgfqpoint{1.663869in}{1.655615in}}%
\pgfpathlineto{\pgfqpoint{1.664275in}{1.642367in}}%
\pgfpathlineto{\pgfqpoint{1.669487in}{1.629119in}}%
\pgfpathlineto{\pgfqpoint{1.670832in}{1.615871in}}%
\pgfpathlineto{\pgfqpoint{1.674508in}{1.602623in}}%
\pgfpathlineto{\pgfqpoint{1.677699in}{1.589375in}}%
\pgfpathlineto{\pgfqpoint{1.678717in}{1.581874in}}%
\pgfpathlineto{\pgfqpoint{1.679654in}{1.576127in}}%
\pgfpathlineto{\pgfqpoint{1.685377in}{1.562879in}}%
\pgfpathlineto{\pgfqpoint{1.685829in}{1.549631in}}%
\pgfpathlineto{\pgfqpoint{1.691346in}{1.536383in}}%
\pgfpathlineto{\pgfqpoint{1.692115in}{1.523135in}}%
\pgfpathlineto{\pgfqpoint{1.696427in}{1.509887in}}%
\pgfpathlineto{\pgfqpoint{1.699369in}{1.496639in}}%
\pgfpathlineto{\pgfqpoint{1.701411in}{1.483391in}}%
\pgfpathlineto{\pgfqpoint{1.704001in}{1.476689in}}%
\pgfpathlineto{\pgfqpoint{1.706968in}{1.470143in}}%
\pgfpathlineto{\pgfqpoint{1.707431in}{1.456895in}}%
\pgfpathlineto{\pgfqpoint{1.714442in}{1.443647in}}%
\pgfpathlineto{\pgfqpoint{1.715103in}{1.430399in}}%
\pgfpathlineto{\pgfqpoint{1.720041in}{1.417151in}}%
\pgfpathlineto{\pgfqpoint{1.722754in}{1.403903in}}%
\pgfpathlineto{\pgfqpoint{1.725130in}{1.390655in}}%
\pgfpathlineto{\pgfqpoint{1.729284in}{1.380344in}}%
\pgfpathlineto{\pgfqpoint{1.730654in}{1.377407in}}%
\pgfpathlineto{\pgfqpoint{1.730982in}{1.364159in}}%
\pgfpathlineto{\pgfqpoint{1.738249in}{1.350911in}}%
\pgfpathlineto{\pgfqpoint{1.738012in}{1.337663in}}%
\pgfpathlineto{\pgfqpoint{1.744454in}{1.324415in}}%
\pgfpathlineto{\pgfqpoint{1.746581in}{1.311167in}}%
\pgfpathlineto{\pgfqpoint{1.749499in}{1.297919in}}%
\pgfpathlineto{\pgfqpoint{1.754567in}{1.285781in}}%
\pgfpathlineto{\pgfqpoint{1.755094in}{1.284671in}}%
\pgfpathlineto{\pgfqpoint{1.755121in}{1.271423in}}%
\pgfpathlineto{\pgfqpoint{1.763745in}{1.258175in}}%
\pgfpathlineto{\pgfqpoint{1.763065in}{1.244927in}}%
\pgfpathlineto{\pgfqpoint{1.770270in}{1.231679in}}%
\pgfpathlineto{\pgfqpoint{1.770955in}{1.218431in}}%
\pgfpathlineto{\pgfqpoint{1.774802in}{1.205183in}}%
\pgfpathlineto{\pgfqpoint{1.779374in}{1.191935in}}%
\pgfpathlineto{\pgfqpoint{1.778821in}{1.178687in}}%
\pgfpathlineto{\pgfqpoint{1.779850in}{1.176943in}}%
\pgfpathlineto{\pgfqpoint{1.787256in}{1.165439in}}%
\pgfpathlineto{\pgfqpoint{1.784443in}{1.152191in}}%
\pgfpathlineto{\pgfqpoint{1.794138in}{1.138943in}}%
\pgfpathlineto{\pgfqpoint{1.793318in}{1.125695in}}%
\pgfpathlineto{\pgfqpoint{1.798553in}{1.112447in}}%
\pgfpathlineto{\pgfqpoint{1.801598in}{1.099199in}}%
\pgfpathlineto{\pgfqpoint{1.800607in}{1.085951in}}%
\pgfpathlineto{\pgfqpoint{1.805133in}{1.078025in}}%
\pgfpathlineto{\pgfqpoint{1.808452in}{1.072703in}}%
\pgfpathlineto{\pgfqpoint{1.805133in}{1.065440in}}%
\pgfpathlineto{\pgfqpoint{1.802333in}{1.059455in}}%
\pgfpathlineto{\pgfqpoint{1.805133in}{1.055320in}}%
\pgfpathlineto{\pgfqpoint{1.811633in}{1.046207in}}%
\pgfpathlineto{\pgfqpoint{1.805133in}{1.034227in}}%
\pgfpathlineto{\pgfqpoint{1.804398in}{1.032959in}}%
\pgfpathlineto{\pgfqpoint{1.805133in}{1.031577in}}%
\pgfpathlineto{\pgfqpoint{1.811671in}{1.019711in}}%
\pgfpathlineto{\pgfqpoint{1.809312in}{1.006463in}}%
\pgfpathlineto{\pgfqpoint{1.807310in}{0.993215in}}%
\pgfpathlineto{\pgfqpoint{1.811417in}{0.979967in}}%
\pgfpathlineto{\pgfqpoint{1.805133in}{0.973329in}}%
\pgfpathlineto{\pgfqpoint{1.798464in}{0.966719in}}%
\pgfpathlineto{\pgfqpoint{1.805133in}{0.957486in}}%
\pgfpathlineto{\pgfqpoint{1.808213in}{0.953471in}}%
\pgfpathlineto{\pgfqpoint{1.805133in}{0.951153in}}%
\pgfpathlineto{\pgfqpoint{1.788838in}{0.940223in}}%
\pgfpathlineto{\pgfqpoint{1.795308in}{0.926975in}}%
\pgfpathlineto{\pgfqpoint{1.779850in}{0.917052in}}%
\pgfpathlineto{\pgfqpoint{1.773693in}{0.913727in}}%
\pgfpathlineto{\pgfqpoint{1.768359in}{0.900479in}}%
\pgfpathlineto{\pgfqpoint{1.754567in}{0.889056in}}%
\pgfpathlineto{\pgfqpoint{1.751865in}{0.887231in}}%
\pgfpathlineto{\pgfqpoint{1.729284in}{0.878032in}}%
\pgfpathlineto{\pgfqpoint{1.718166in}{0.873983in}}%
\pgfpathlineto{\pgfqpoint{1.704133in}{0.860735in}}%
\pgfpathlineto{\pgfqpoint{1.704001in}{0.860709in}}%
\pgfpathlineto{\pgfqpoint{1.678717in}{0.856146in}}%
\pgfpathlineto{\pgfqpoint{1.653434in}{0.852200in}}%
\pgfpathlineto{\pgfqpoint{1.628151in}{0.848780in}}%
\pgfpathlineto{\pgfqpoint{1.617206in}{0.847487in}}%
\pgfpathlineto{\pgfqpoint{1.602868in}{0.835810in}}%
\pgfpathlineto{\pgfqpoint{1.600293in}{0.834239in}}%
\pgfpathlineto{\pgfqpoint{1.577585in}{0.832323in}}%
\pgfpathlineto{\pgfqpoint{1.552302in}{0.830614in}}%
\pgfpathlineto{\pgfqpoint{1.527018in}{0.829292in}}%
\pgfpathlineto{\pgfqpoint{1.501735in}{0.828329in}}%
\pgfpathlineto{\pgfqpoint{1.476452in}{0.827704in}}%
\pgfpathlineto{\pgfqpoint{1.451169in}{0.827407in}}%
\pgfpathlineto{\pgfqpoint{1.425886in}{0.827436in}}%
\pgfpathlineto{\pgfqpoint{1.400603in}{0.827799in}}%
\pgfpathlineto{\pgfqpoint{1.375319in}{0.828514in}}%
\pgfpathlineto{\pgfqpoint{1.350036in}{0.829609in}}%
\pgfpathlineto{\pgfqpoint{1.324753in}{0.831126in}}%
\pgfpathlineto{\pgfqpoint{1.299470in}{0.833123in}}%
\pgfpathlineto{\pgfqpoint{1.288142in}{0.834239in}}%
\pgfpathlineto{\pgfqpoint{1.274187in}{0.837757in}}%
\pgfpathlineto{\pgfqpoint{1.248904in}{0.845033in}}%
\pgfpathlineto{\pgfqpoint{1.241356in}{0.847487in}}%
\pgfpathlineto{\pgfqpoint{1.223620in}{0.852179in}}%
\pgfpathlineto{\pgfqpoint{1.200733in}{0.860735in}}%
\pgfpathlineto{\pgfqpoint{1.198337in}{0.861229in}}%
\pgfpathlineto{\pgfqpoint{1.173054in}{0.867012in}}%
\pgfpathlineto{\pgfqpoint{1.147771in}{0.873606in}}%
\pgfpathlineto{\pgfqpoint{1.146457in}{0.873983in}}%
\pgfpathlineto{\pgfqpoint{1.147771in}{0.877399in}}%
\pgfpathlineto{\pgfqpoint{1.151281in}{0.887231in}}%
\pgfpathlineto{\pgfqpoint{1.147771in}{0.888253in}}%
\pgfpathlineto{\pgfqpoint{1.122488in}{0.896764in}}%
\pgfpathlineto{\pgfqpoint{1.113102in}{0.900479in}}%
\pgfpathlineto{\pgfqpoint{1.117682in}{0.913727in}}%
\pgfpathlineto{\pgfqpoint{1.100883in}{0.926975in}}%
\pgfpathlineto{\pgfqpoint{1.097204in}{0.932832in}}%
\pgfpathlineto{\pgfqpoint{1.092750in}{0.940223in}}%
\pgfpathlineto{\pgfqpoint{1.097204in}{0.951337in}}%
\pgfpathlineto{\pgfqpoint{1.097977in}{0.953471in}}%
\pgfpathlineto{\pgfqpoint{1.097204in}{0.954111in}}%
\pgfpathlineto{\pgfqpoint{1.083108in}{0.966719in}}%
\pgfpathlineto{\pgfqpoint{1.094963in}{0.979967in}}%
\pgfpathlineto{\pgfqpoint{1.081168in}{0.993215in}}%
\pgfpathlineto{\pgfqpoint{1.088036in}{1.006463in}}%
\pgfpathlineto{\pgfqpoint{1.085978in}{1.019711in}}%
\pgfpathlineto{\pgfqpoint{1.084858in}{1.032959in}}%
\pgfpathlineto{\pgfqpoint{1.092584in}{1.046207in}}%
\pgfpathlineto{\pgfqpoint{1.085921in}{1.059455in}}%
\pgfpathlineto{\pgfqpoint{1.097204in}{1.071609in}}%
\pgfpathlineto{\pgfqpoint{1.098168in}{1.072703in}}%
\pgfpathlineto{\pgfqpoint{1.097204in}{1.075592in}}%
\pgfpathlineto{\pgfqpoint{1.093881in}{1.085951in}}%
\pgfpathlineto{\pgfqpoint{1.097204in}{1.091140in}}%
\pgfpathlineto{\pgfqpoint{1.102161in}{1.099199in}}%
\pgfpathlineto{\pgfqpoint{1.102978in}{1.112447in}}%
\pgfpathlineto{\pgfqpoint{1.103750in}{1.125695in}}%
\pgfpathlineto{\pgfqpoint{1.111248in}{1.138943in}}%
\pgfpathlineto{\pgfqpoint{1.108341in}{1.152191in}}%
\pgfpathlineto{\pgfqpoint{1.118317in}{1.165439in}}%
\pgfpathlineto{\pgfqpoint{1.116016in}{1.178687in}}%
\pgfpathlineto{\pgfqpoint{1.122488in}{1.189568in}}%
\pgfpathlineto{\pgfqpoint{1.123817in}{1.191935in}}%
\pgfpathlineto{\pgfqpoint{1.126440in}{1.205183in}}%
\pgfpathlineto{\pgfqpoint{1.129710in}{1.218431in}}%
\pgfpathlineto{\pgfqpoint{1.136334in}{1.231679in}}%
\pgfpathlineto{\pgfqpoint{1.135191in}{1.244927in}}%
\pgfpathlineto{\pgfqpoint{1.143921in}{1.258175in}}%
\pgfpathlineto{\pgfqpoint{1.142797in}{1.271423in}}%
\pgfpathlineto{\pgfqpoint{1.147771in}{1.280391in}}%
\pgfpathlineto{\pgfqpoint{1.149959in}{1.284671in}}%
\pgfpathlineto{\pgfqpoint{1.151793in}{1.297919in}}%
\pgfpathlineto{\pgfqpoint{1.154825in}{1.311167in}}%
\pgfpathlineto{\pgfqpoint{1.160806in}{1.324415in}}%
\pgfpathlineto{\pgfqpoint{1.161425in}{1.337663in}}%
\pgfpathlineto{\pgfqpoint{1.169307in}{1.350911in}}%
\pgfpathlineto{\pgfqpoint{1.168833in}{1.364159in}}%
\pgfpathlineto{\pgfqpoint{1.173054in}{1.372259in}}%
\pgfpathlineto{\pgfqpoint{1.175469in}{1.377407in}}%
\pgfpathlineto{\pgfqpoint{1.177205in}{1.390655in}}%
\pgfpathlineto{\pgfqpoint{1.180856in}{1.403903in}}%
\pgfpathlineto{\pgfqpoint{1.185581in}{1.417151in}}%
\pgfpathlineto{\pgfqpoint{1.186369in}{1.430399in}}%
\pgfpathlineto{\pgfqpoint{1.193249in}{1.443647in}}%
\pgfpathlineto{\pgfqpoint{1.193379in}{1.456895in}}%
\pgfpathlineto{\pgfqpoint{1.198337in}{1.466767in}}%
\pgfpathlineto{\pgfqpoint{1.199827in}{1.470143in}}%
\pgfpathlineto{\pgfqpoint{1.200971in}{1.483391in}}%
\pgfpathlineto{\pgfqpoint{1.204511in}{1.496639in}}%
\pgfpathlineto{\pgfqpoint{1.208465in}{1.509887in}}%
\pgfpathlineto{\pgfqpoint{1.210023in}{1.523135in}}%
\pgfpathlineto{\pgfqpoint{1.216100in}{1.536383in}}%
\pgfpathlineto{\pgfqpoint{1.216523in}{1.549631in}}%
\pgfpathlineto{\pgfqpoint{1.222713in}{1.562879in}}%
\pgfpathlineto{\pgfqpoint{1.223620in}{1.571737in}}%
\pgfpathlineto{\pgfqpoint{1.223983in}{1.576127in}}%
\pgfpathlineto{\pgfqpoint{1.227730in}{1.589375in}}%
\pgfpathlineto{\pgfqpoint{1.230507in}{1.602623in}}%
\pgfpathlineto{\pgfqpoint{1.232025in}{1.615871in}}%
\pgfpathlineto{\pgfqpoint{1.236980in}{1.629119in}}%
\pgfpathlineto{\pgfqpoint{1.237630in}{1.642367in}}%
\pgfpathlineto{\pgfqpoint{1.243410in}{1.655615in}}%
\pgfpathlineto{\pgfqpoint{1.244589in}{1.668863in}}%
\pgfpathlineto{\pgfqpoint{1.248904in}{1.681268in}}%
\pgfpathlineto{\pgfqpoint{1.249151in}{1.682111in}}%
\pgfpathlineto{\pgfqpoint{1.251639in}{1.695359in}}%
\pgfpathlineto{\pgfqpoint{1.253717in}{1.708607in}}%
\pgfpathlineto{\pgfqpoint{1.257697in}{1.721855in}}%
\pgfpathlineto{\pgfqpoint{1.258303in}{1.735103in}}%
\pgfpathlineto{\pgfqpoint{1.263163in}{1.748351in}}%
\pgfpathlineto{\pgfqpoint{1.264013in}{1.761599in}}%
\pgfpathlineto{\pgfqpoint{1.268528in}{1.774847in}}%
\pgfpathlineto{\pgfqpoint{1.270805in}{1.788095in}}%
\pgfpathlineto{\pgfqpoint{1.273641in}{1.801343in}}%
\pgfpathlineto{\pgfqpoint{1.274187in}{1.803018in}}%
\pgfpathlineto{\pgfqpoint{1.277171in}{1.814591in}}%
\pgfpathlineto{\pgfqpoint{1.278221in}{1.827839in}}%
\pgfpathlineto{\pgfqpoint{1.282633in}{1.841087in}}%
\pgfpathlineto{\pgfqpoint{1.283273in}{1.854335in}}%
\pgfpathlineto{\pgfqpoint{1.287430in}{1.867583in}}%
\pgfpathlineto{\pgfqpoint{1.289200in}{1.880831in}}%
\pgfpathlineto{\pgfqpoint{1.292314in}{1.894079in}}%
\pgfpathlineto{\pgfqpoint{1.295631in}{1.907327in}}%
\pgfpathlineto{\pgfqpoint{1.297124in}{1.920575in}}%
\pgfpathlineto{\pgfqpoint{1.299470in}{1.927225in}}%
\pgfpathlineto{\pgfqpoint{1.301351in}{1.933823in}}%
\pgfpathlineto{\pgfqpoint{1.302008in}{1.947071in}}%
\pgfpathlineto{\pgfqpoint{1.306096in}{1.960319in}}%
\pgfpathlineto{\pgfqpoint{1.307274in}{1.973567in}}%
\pgfpathlineto{\pgfqpoint{1.310462in}{1.986815in}}%
\pgfpathlineto{\pgfqpoint{1.313245in}{2.000063in}}%
\pgfpathlineto{\pgfqpoint{1.315274in}{2.013311in}}%
\pgfpathlineto{\pgfqpoint{1.319401in}{2.026559in}}%
\pgfpathlineto{\pgfqpoint{1.320325in}{2.039807in}}%
\pgfpathlineto{\pgfqpoint{1.324753in}{2.052544in}}%
\pgfpathlineto{\pgfqpoint{1.324903in}{2.053055in}}%
\pgfpathlineto{\pgfqpoint{1.325786in}{2.066303in}}%
\pgfpathlineto{\pgfqpoint{1.329197in}{2.079551in}}%
\pgfpathlineto{\pgfqpoint{1.331161in}{2.092799in}}%
\pgfpathlineto{\pgfqpoint{1.333384in}{2.106047in}}%
\pgfpathlineto{\pgfqpoint{1.336937in}{2.119295in}}%
\pgfpathlineto{\pgfqpoint{1.338237in}{2.132543in}}%
\pgfpathlineto{\pgfqpoint{1.329357in}{2.132543in}}%
\pgfpathlineto{\pgfqpoint{1.328198in}{2.119295in}}%
\pgfpathlineto{\pgfqpoint{1.324806in}{2.106047in}}%
\pgfpathlineto{\pgfqpoint{1.324753in}{2.105695in}}%
\pgfpathlineto{\pgfqpoint{1.322291in}{2.092799in}}%
\pgfpathlineto{\pgfqpoint{1.319766in}{2.079551in}}%
\pgfpathlineto{\pgfqpoint{1.315990in}{2.066303in}}%
\pgfpathlineto{\pgfqpoint{1.314859in}{2.053055in}}%
\pgfpathlineto{\pgfqpoint{1.310793in}{2.039807in}}%
\pgfpathlineto{\pgfqpoint{1.309936in}{2.026559in}}%
\pgfpathlineto{\pgfqpoint{1.306112in}{2.013311in}}%
\pgfpathlineto{\pgfqpoint{1.304354in}{2.000063in}}%
\pgfpathlineto{\pgfqpoint{1.301629in}{1.986815in}}%
\pgfpathlineto{\pgfqpoint{1.299470in}{1.976694in}}%
\pgfpathlineto{\pgfqpoint{1.298605in}{1.973567in}}%
\pgfpathlineto{\pgfqpoint{1.296944in}{1.960319in}}%
\pgfpathlineto{\pgfqpoint{1.292123in}{1.947071in}}%
\pgfpathlineto{\pgfqpoint{1.291266in}{1.933823in}}%
\pgfpathlineto{\pgfqpoint{1.286950in}{1.920575in}}%
\pgfpathlineto{\pgfqpoint{1.285645in}{1.907327in}}%
\pgfpathlineto{\pgfqpoint{1.282471in}{1.894079in}}%
\pgfpathlineto{\pgfqpoint{1.279755in}{1.880831in}}%
\pgfpathlineto{\pgfqpoint{1.277929in}{1.867583in}}%
\pgfpathlineto{\pgfqpoint{1.274212in}{1.854335in}}%
\pgfpathlineto{\pgfqpoint{1.274187in}{1.853866in}}%
\pgfpathlineto{\pgfqpoint{1.273274in}{1.841087in}}%
\pgfpathlineto{\pgfqpoint{1.267897in}{1.827839in}}%
\pgfpathlineto{\pgfqpoint{1.266660in}{1.814591in}}%
\pgfpathlineto{\pgfqpoint{1.262538in}{1.801343in}}%
\pgfpathlineto{\pgfqpoint{1.260085in}{1.788095in}}%
\pgfpathlineto{\pgfqpoint{1.257779in}{1.774847in}}%
\pgfpathlineto{\pgfqpoint{1.253756in}{1.761599in}}%
\pgfpathlineto{\pgfqpoint{1.252777in}{1.748351in}}%
\pgfpathlineto{\pgfqpoint{1.248904in}{1.736970in}}%
\pgfpathlineto{\pgfqpoint{1.248099in}{1.735103in}}%
\pgfpathlineto{\pgfqpoint{1.247313in}{1.721855in}}%
\pgfpathlineto{\pgfqpoint{1.242388in}{1.708607in}}%
\pgfpathlineto{\pgfqpoint{1.240070in}{1.695359in}}%
\pgfpathlineto{\pgfqpoint{1.236932in}{1.682111in}}%
\pgfpathlineto{\pgfqpoint{1.232880in}{1.668863in}}%
\pgfpathlineto{\pgfqpoint{1.231547in}{1.655615in}}%
\pgfpathlineto{\pgfqpoint{1.226252in}{1.642367in}}%
\pgfpathlineto{\pgfqpoint{1.225524in}{1.629119in}}%
\pgfpathlineto{\pgfqpoint{1.223620in}{1.623901in}}%
\pgfpathlineto{\pgfqpoint{1.219979in}{1.615871in}}%
\pgfpathlineto{\pgfqpoint{1.218275in}{1.602623in}}%
\pgfpathlineto{\pgfqpoint{1.214731in}{1.589375in}}%
\pgfpathlineto{\pgfqpoint{1.210658in}{1.576127in}}%
\pgfpathlineto{\pgfqpoint{1.209043in}{1.562879in}}%
\pgfpathlineto{\pgfqpoint{1.203435in}{1.549631in}}%
\pgfpathlineto{\pgfqpoint{1.202838in}{1.536383in}}%
\pgfpathlineto{\pgfqpoint{1.198337in}{1.526160in}}%
\pgfpathlineto{\pgfqpoint{1.196725in}{1.523135in}}%
\pgfpathlineto{\pgfqpoint{1.194946in}{1.509887in}}%
\pgfpathlineto{\pgfqpoint{1.190027in}{1.496639in}}%
\pgfpathlineto{\pgfqpoint{1.186263in}{1.483391in}}%
\pgfpathlineto{\pgfqpoint{1.184507in}{1.470143in}}%
\pgfpathlineto{\pgfqpoint{1.178368in}{1.456895in}}%
\pgfpathlineto{\pgfqpoint{1.177849in}{1.443647in}}%
\pgfpathlineto{\pgfqpoint{1.173054in}{1.434081in}}%
\pgfpathlineto{\pgfqpoint{1.170881in}{1.430399in}}%
\pgfpathlineto{\pgfqpoint{1.169906in}{1.417151in}}%
\pgfpathlineto{\pgfqpoint{1.164177in}{1.403903in}}%
\pgfpathlineto{\pgfqpoint{1.160365in}{1.390655in}}%
\pgfpathlineto{\pgfqpoint{1.157811in}{1.377407in}}%
\pgfpathlineto{\pgfqpoint{1.151446in}{1.364159in}}%
\pgfpathlineto{\pgfqpoint{1.151307in}{1.350911in}}%
\pgfpathlineto{\pgfqpoint{1.147771in}{1.344815in}}%
\pgfpathlineto{\pgfqpoint{1.143036in}{1.337663in}}%
\pgfpathlineto{\pgfqpoint{1.142060in}{1.324415in}}%
\pgfpathlineto{\pgfqpoint{1.134919in}{1.311167in}}%
\pgfpathlineto{\pgfqpoint{1.131816in}{1.297919in}}%
\pgfpathlineto{\pgfqpoint{1.128933in}{1.284671in}}%
\pgfpathlineto{\pgfqpoint{1.122488in}{1.272204in}}%
\pgfpathlineto{\pgfqpoint{1.122040in}{1.271423in}}%
\pgfpathlineto{\pgfqpoint{1.122227in}{1.258175in}}%
\pgfpathlineto{\pgfqpoint{1.112666in}{1.244927in}}%
\pgfpathlineto{\pgfqpoint{1.113316in}{1.231679in}}%
\pgfpathlineto{\pgfqpoint{1.105464in}{1.218431in}}%
\pgfpathlineto{\pgfqpoint{1.102046in}{1.205183in}}%
\pgfpathlineto{\pgfqpoint{1.097864in}{1.191935in}}%
\pgfpathlineto{\pgfqpoint{1.097204in}{1.190775in}}%
\pgfpathlineto{\pgfqpoint{1.089823in}{1.178687in}}%
\pgfpathlineto{\pgfqpoint{1.090641in}{1.165439in}}%
\pgfpathlineto{\pgfqpoint{1.079910in}{1.152191in}}%
\pgfpathlineto{\pgfqpoint{1.081885in}{1.138943in}}%
\pgfpathlineto{\pgfqpoint{1.072658in}{1.125695in}}%
\pgfpathlineto{\pgfqpoint{1.071921in}{1.116797in}}%
\pgfpathlineto{\pgfqpoint{1.071529in}{1.112447in}}%
\pgfpathlineto{\pgfqpoint{1.068045in}{1.099199in}}%
\pgfpathlineto{\pgfqpoint{1.059309in}{1.085951in}}%
\pgfpathlineto{\pgfqpoint{1.061055in}{1.072703in}}%
\pgfpathlineto{\pgfqpoint{1.047834in}{1.059455in}}%
\pgfpathlineto{\pgfqpoint{1.052695in}{1.046207in}}%
\pgfpathlineto{\pgfqpoint{1.046638in}{1.038332in}}%
\pgfpathlineto{\pgfqpoint{1.042095in}{1.032959in}}%
\pgfpathlineto{\pgfqpoint{1.042338in}{1.019711in}}%
\pgfpathlineto{\pgfqpoint{1.039636in}{1.006463in}}%
\pgfpathlineto{\pgfqpoint{1.031906in}{0.993215in}}%
\pgfpathlineto{\pgfqpoint{1.040362in}{0.979967in}}%
\pgfpathlineto{\pgfqpoint{1.026010in}{0.966719in}}%
\pgfpathlineto{\pgfqpoint{1.036875in}{0.953471in}}%
\pgfpathlineto{\pgfqpoint{1.025219in}{0.940223in}}%
\pgfpathlineto{\pgfqpoint{1.031438in}{0.926975in}}%
\pgfpathlineto{\pgfqpoint{1.036972in}{0.913727in}}%
\pgfpathlineto{\pgfqpoint{1.030082in}{0.900479in}}%
\pgfpathlineto{\pgfqpoint{1.046638in}{0.890605in}}%
\pgfpathlineto{\pgfqpoint{1.053196in}{0.887231in}}%
\pgfpathlineto{\pgfqpoint{1.046638in}{0.880871in}}%
\pgfpathlineto{\pgfqpoint{1.039044in}{0.873983in}}%
\pgfpathlineto{\pgfqpoint{1.046638in}{0.870865in}}%
\pgfpathlineto{\pgfqpoint{1.071921in}{0.861585in}}%
\pgfpathlineto{\pgfqpoint{1.074481in}{0.860735in}}%
\pgfpathlineto{\pgfqpoint{1.077725in}{0.847487in}}%
\pgfpathlineto{\pgfqpoint{1.097204in}{0.836367in}}%
\pgfpathlineto{\pgfqpoint{1.101257in}{0.834239in}}%
\pgfpathlineto{\pgfqpoint{1.122488in}{0.827937in}}%
\pgfpathlineto{\pgfqpoint{1.147771in}{0.822133in}}%
\pgfpathlineto{\pgfqpoint{1.153790in}{0.820991in}}%
\pgfpathlineto{\pgfqpoint{1.158822in}{0.807743in}}%
\pgfpathlineto{\pgfqpoint{1.173054in}{0.805992in}}%
\pgfpathlineto{\pgfqpoint{1.198337in}{0.803264in}}%
\pgfpathlineto{\pgfqpoint{1.223620in}{0.800913in}}%
\pgfpathlineto{\pgfqpoint{1.248904in}{0.798891in}}%
\pgfpathlineto{\pgfqpoint{1.274187in}{0.797162in}}%
\pgfpathlineto{\pgfqpoint{1.299470in}{0.795696in}}%
\pgfpathlineto{\pgfqpoint{1.324284in}{0.794495in}}%
\pgfpathlineto{\pgfqpoint{1.324753in}{0.794407in}}%
\pgfpathclose%
\pgfpathmoveto{\pgfqpoint{1.071921in}{1.907871in}}%
\pgfpathlineto{\pgfqpoint{1.077206in}{1.920575in}}%
\pgfpathlineto{\pgfqpoint{1.084902in}{1.933823in}}%
\pgfpathlineto{\pgfqpoint{1.097204in}{1.940615in}}%
\pgfpathlineto{\pgfqpoint{1.101376in}{1.947071in}}%
\pgfpathlineto{\pgfqpoint{1.100789in}{1.960319in}}%
\pgfpathlineto{\pgfqpoint{1.115659in}{1.973567in}}%
\pgfpathlineto{\pgfqpoint{1.113480in}{1.986815in}}%
\pgfpathlineto{\pgfqpoint{1.122488in}{1.993608in}}%
\pgfpathlineto{\pgfqpoint{1.127245in}{2.000063in}}%
\pgfpathlineto{\pgfqpoint{1.128071in}{2.013311in}}%
\pgfpathlineto{\pgfqpoint{1.134906in}{2.026559in}}%
\pgfpathlineto{\pgfqpoint{1.142056in}{2.039807in}}%
\pgfpathlineto{\pgfqpoint{1.145813in}{2.053055in}}%
\pgfpathlineto{\pgfqpoint{1.147771in}{2.055038in}}%
\pgfpathlineto{\pgfqpoint{1.154255in}{2.066303in}}%
\pgfpathlineto{\pgfqpoint{1.154251in}{2.079551in}}%
\pgfpathlineto{\pgfqpoint{1.163478in}{2.092799in}}%
\pgfpathlineto{\pgfqpoint{1.164483in}{2.106047in}}%
\pgfpathlineto{\pgfqpoint{1.173054in}{2.119237in}}%
\pgfpathlineto{\pgfqpoint{1.173083in}{2.119295in}}%
\pgfpathlineto{\pgfqpoint{1.176175in}{2.132543in}}%
\pgfpathlineto{\pgfqpoint{1.173054in}{2.132543in}}%
\pgfpathlineto{\pgfqpoint{1.149378in}{2.132543in}}%
\pgfpathlineto{\pgfqpoint{1.147771in}{2.127648in}}%
\pgfpathlineto{\pgfqpoint{1.138776in}{2.119295in}}%
\pgfpathlineto{\pgfqpoint{1.126656in}{2.106047in}}%
\pgfpathlineto{\pgfqpoint{1.127644in}{2.092799in}}%
\pgfpathlineto{\pgfqpoint{1.122488in}{2.088463in}}%
\pgfpathlineto{\pgfqpoint{1.110426in}{2.092799in}}%
\pgfpathlineto{\pgfqpoint{1.115163in}{2.106047in}}%
\pgfpathlineto{\pgfqpoint{1.111320in}{2.119295in}}%
\pgfpathlineto{\pgfqpoint{1.105823in}{2.132543in}}%
\pgfpathlineto{\pgfqpoint{1.097204in}{2.132543in}}%
\pgfpathlineto{\pgfqpoint{1.075415in}{2.132543in}}%
\pgfpathlineto{\pgfqpoint{1.076017in}{2.119295in}}%
\pgfpathlineto{\pgfqpoint{1.071921in}{2.108695in}}%
\pgfpathlineto{\pgfqpoint{1.071076in}{2.106047in}}%
\pgfpathlineto{\pgfqpoint{1.068954in}{2.092799in}}%
\pgfpathlineto{\pgfqpoint{1.069098in}{2.079551in}}%
\pgfpathlineto{\pgfqpoint{1.063173in}{2.066303in}}%
\pgfpathlineto{\pgfqpoint{1.066524in}{2.053055in}}%
\pgfpathlineto{\pgfqpoint{1.059735in}{2.039807in}}%
\pgfpathlineto{\pgfqpoint{1.063098in}{2.026559in}}%
\pgfpathlineto{\pgfqpoint{1.058448in}{2.013311in}}%
\pgfpathlineto{\pgfqpoint{1.056020in}{2.000063in}}%
\pgfpathlineto{\pgfqpoint{1.057449in}{1.986815in}}%
\pgfpathlineto{\pgfqpoint{1.047599in}{1.973567in}}%
\pgfpathlineto{\pgfqpoint{1.058177in}{1.960319in}}%
\pgfpathlineto{\pgfqpoint{1.046638in}{1.950730in}}%
\pgfpathlineto{\pgfqpoint{1.044271in}{1.947071in}}%
\pgfpathlineto{\pgfqpoint{1.046638in}{1.941753in}}%
\pgfpathlineto{\pgfqpoint{1.057443in}{1.933823in}}%
\pgfpathlineto{\pgfqpoint{1.050834in}{1.920575in}}%
\pgfpathlineto{\pgfqpoint{1.071921in}{1.907871in}}%
\pgfpathclose%
\pgfpathmoveto{\pgfqpoint{1.830416in}{1.931778in}}%
\pgfpathlineto{\pgfqpoint{1.837429in}{1.933823in}}%
\pgfpathlineto{\pgfqpoint{1.830944in}{1.947071in}}%
\pgfpathlineto{\pgfqpoint{1.839829in}{1.960319in}}%
\pgfpathlineto{\pgfqpoint{1.842410in}{1.973567in}}%
\pgfpathlineto{\pgfqpoint{1.839201in}{1.986815in}}%
\pgfpathlineto{\pgfqpoint{1.845411in}{2.000063in}}%
\pgfpathlineto{\pgfqpoint{1.837160in}{2.013311in}}%
\pgfpathlineto{\pgfqpoint{1.843074in}{2.026559in}}%
\pgfpathlineto{\pgfqpoint{1.836351in}{2.039807in}}%
\pgfpathlineto{\pgfqpoint{1.838021in}{2.053055in}}%
\pgfpathlineto{\pgfqpoint{1.834739in}{2.066303in}}%
\pgfpathlineto{\pgfqpoint{1.831739in}{2.079551in}}%
\pgfpathlineto{\pgfqpoint{1.832279in}{2.092799in}}%
\pgfpathlineto{\pgfqpoint{1.830416in}{2.097918in}}%
\pgfpathlineto{\pgfqpoint{1.826187in}{2.106047in}}%
\pgfpathlineto{\pgfqpoint{1.829251in}{2.119295in}}%
\pgfpathlineto{\pgfqpoint{1.822412in}{2.132543in}}%
\pgfpathlineto{\pgfqpoint{1.805133in}{2.132543in}}%
\pgfpathlineto{\pgfqpoint{1.781130in}{2.132543in}}%
\pgfpathlineto{\pgfqpoint{1.786670in}{2.119295in}}%
\pgfpathlineto{\pgfqpoint{1.779850in}{2.115676in}}%
\pgfpathlineto{\pgfqpoint{1.774838in}{2.119295in}}%
\pgfpathlineto{\pgfqpoint{1.777927in}{2.132543in}}%
\pgfpathlineto{\pgfqpoint{1.754567in}{2.132543in}}%
\pgfpathlineto{\pgfqpoint{1.732874in}{2.132543in}}%
\pgfpathlineto{\pgfqpoint{1.736484in}{2.119295in}}%
\pgfpathlineto{\pgfqpoint{1.745008in}{2.106047in}}%
\pgfpathlineto{\pgfqpoint{1.745606in}{2.092799in}}%
\pgfpathlineto{\pgfqpoint{1.754340in}{2.079551in}}%
\pgfpathlineto{\pgfqpoint{1.754567in}{2.074740in}}%
\pgfpathlineto{\pgfqpoint{1.755220in}{2.066303in}}%
\pgfpathlineto{\pgfqpoint{1.764120in}{2.053055in}}%
\pgfpathlineto{\pgfqpoint{1.769977in}{2.039807in}}%
\pgfpathlineto{\pgfqpoint{1.772784in}{2.026559in}}%
\pgfpathlineto{\pgfqpoint{1.779850in}{2.016322in}}%
\pgfpathlineto{\pgfqpoint{1.784471in}{2.013311in}}%
\pgfpathlineto{\pgfqpoint{1.782779in}{2.000063in}}%
\pgfpathlineto{\pgfqpoint{1.798839in}{1.986815in}}%
\pgfpathlineto{\pgfqpoint{1.799407in}{1.973567in}}%
\pgfpathlineto{\pgfqpoint{1.805133in}{1.965314in}}%
\pgfpathlineto{\pgfqpoint{1.815682in}{1.960319in}}%
\pgfpathlineto{\pgfqpoint{1.829821in}{1.947071in}}%
\pgfpathlineto{\pgfqpoint{1.827464in}{1.933823in}}%
\pgfpathlineto{\pgfqpoint{1.830416in}{1.931778in}}%
\pgfpathclose%
\pgfusepath{fill}%
\end{pgfscope}%
\begin{pgfscope}%
\pgfpathrectangle{\pgfqpoint{0.540975in}{0.476543in}}{\pgfqpoint{1.820388in}{1.656000in}}%
\pgfusepath{clip}%
\pgfsetbuttcap%
\pgfsetroundjoin%
\definecolor{currentfill}{rgb}{0.514110,0.725798,0.845444}%
\pgfsetfillcolor{currentfill}%
\pgfsetlinewidth{0.000000pt}%
\definecolor{currentstroke}{rgb}{0.000000,0.000000,0.000000}%
\pgfsetstrokecolor{currentstroke}%
\pgfsetdash{}{0pt}%
\pgfpathmoveto{\pgfqpoint{1.299470in}{0.833123in}}%
\pgfpathlineto{\pgfqpoint{1.324753in}{0.831126in}}%
\pgfpathlineto{\pgfqpoint{1.350036in}{0.829609in}}%
\pgfpathlineto{\pgfqpoint{1.375319in}{0.828514in}}%
\pgfpathlineto{\pgfqpoint{1.400603in}{0.827799in}}%
\pgfpathlineto{\pgfqpoint{1.425886in}{0.827436in}}%
\pgfpathlineto{\pgfqpoint{1.451169in}{0.827407in}}%
\pgfpathlineto{\pgfqpoint{1.476452in}{0.827704in}}%
\pgfpathlineto{\pgfqpoint{1.501735in}{0.828329in}}%
\pgfpathlineto{\pgfqpoint{1.527018in}{0.829292in}}%
\pgfpathlineto{\pgfqpoint{1.552302in}{0.830614in}}%
\pgfpathlineto{\pgfqpoint{1.577585in}{0.832323in}}%
\pgfpathlineto{\pgfqpoint{1.600293in}{0.834239in}}%
\pgfpathlineto{\pgfqpoint{1.602868in}{0.835810in}}%
\pgfpathlineto{\pgfqpoint{1.617206in}{0.847487in}}%
\pgfpathlineto{\pgfqpoint{1.628151in}{0.848780in}}%
\pgfpathlineto{\pgfqpoint{1.653434in}{0.852200in}}%
\pgfpathlineto{\pgfqpoint{1.678717in}{0.856146in}}%
\pgfpathlineto{\pgfqpoint{1.704001in}{0.860709in}}%
\pgfpathlineto{\pgfqpoint{1.704133in}{0.860735in}}%
\pgfpathlineto{\pgfqpoint{1.718166in}{0.873983in}}%
\pgfpathlineto{\pgfqpoint{1.729284in}{0.878032in}}%
\pgfpathlineto{\pgfqpoint{1.751865in}{0.887231in}}%
\pgfpathlineto{\pgfqpoint{1.754567in}{0.889056in}}%
\pgfpathlineto{\pgfqpoint{1.768359in}{0.900479in}}%
\pgfpathlineto{\pgfqpoint{1.773693in}{0.913727in}}%
\pgfpathlineto{\pgfqpoint{1.779850in}{0.917052in}}%
\pgfpathlineto{\pgfqpoint{1.795308in}{0.926975in}}%
\pgfpathlineto{\pgfqpoint{1.788838in}{0.940223in}}%
\pgfpathlineto{\pgfqpoint{1.805133in}{0.951153in}}%
\pgfpathlineto{\pgfqpoint{1.808213in}{0.953471in}}%
\pgfpathlineto{\pgfqpoint{1.805133in}{0.957486in}}%
\pgfpathlineto{\pgfqpoint{1.798464in}{0.966719in}}%
\pgfpathlineto{\pgfqpoint{1.805133in}{0.973329in}}%
\pgfpathlineto{\pgfqpoint{1.811417in}{0.979967in}}%
\pgfpathlineto{\pgfqpoint{1.807310in}{0.993215in}}%
\pgfpathlineto{\pgfqpoint{1.809312in}{1.006463in}}%
\pgfpathlineto{\pgfqpoint{1.811671in}{1.019711in}}%
\pgfpathlineto{\pgfqpoint{1.805133in}{1.031577in}}%
\pgfpathlineto{\pgfqpoint{1.804398in}{1.032959in}}%
\pgfpathlineto{\pgfqpoint{1.805133in}{1.034227in}}%
\pgfpathlineto{\pgfqpoint{1.811633in}{1.046207in}}%
\pgfpathlineto{\pgfqpoint{1.805133in}{1.055320in}}%
\pgfpathlineto{\pgfqpoint{1.802333in}{1.059455in}}%
\pgfpathlineto{\pgfqpoint{1.805133in}{1.065440in}}%
\pgfpathlineto{\pgfqpoint{1.808452in}{1.072703in}}%
\pgfpathlineto{\pgfqpoint{1.805133in}{1.078025in}}%
\pgfpathlineto{\pgfqpoint{1.800607in}{1.085951in}}%
\pgfpathlineto{\pgfqpoint{1.801598in}{1.099199in}}%
\pgfpathlineto{\pgfqpoint{1.798553in}{1.112447in}}%
\pgfpathlineto{\pgfqpoint{1.793318in}{1.125695in}}%
\pgfpathlineto{\pgfqpoint{1.794138in}{1.138943in}}%
\pgfpathlineto{\pgfqpoint{1.784443in}{1.152191in}}%
\pgfpathlineto{\pgfqpoint{1.787256in}{1.165439in}}%
\pgfpathlineto{\pgfqpoint{1.779850in}{1.176943in}}%
\pgfpathlineto{\pgfqpoint{1.778821in}{1.178687in}}%
\pgfpathlineto{\pgfqpoint{1.779374in}{1.191935in}}%
\pgfpathlineto{\pgfqpoint{1.774802in}{1.205183in}}%
\pgfpathlineto{\pgfqpoint{1.770955in}{1.218431in}}%
\pgfpathlineto{\pgfqpoint{1.770270in}{1.231679in}}%
\pgfpathlineto{\pgfqpoint{1.763065in}{1.244927in}}%
\pgfpathlineto{\pgfqpoint{1.763745in}{1.258175in}}%
\pgfpathlineto{\pgfqpoint{1.755121in}{1.271423in}}%
\pgfpathlineto{\pgfqpoint{1.755094in}{1.284671in}}%
\pgfpathlineto{\pgfqpoint{1.754567in}{1.285781in}}%
\pgfpathlineto{\pgfqpoint{1.749499in}{1.297919in}}%
\pgfpathlineto{\pgfqpoint{1.746581in}{1.311167in}}%
\pgfpathlineto{\pgfqpoint{1.744454in}{1.324415in}}%
\pgfpathlineto{\pgfqpoint{1.738012in}{1.337663in}}%
\pgfpathlineto{\pgfqpoint{1.738249in}{1.350911in}}%
\pgfpathlineto{\pgfqpoint{1.730982in}{1.364159in}}%
\pgfpathlineto{\pgfqpoint{1.730654in}{1.377407in}}%
\pgfpathlineto{\pgfqpoint{1.729284in}{1.380344in}}%
\pgfpathlineto{\pgfqpoint{1.725130in}{1.390655in}}%
\pgfpathlineto{\pgfqpoint{1.722754in}{1.403903in}}%
\pgfpathlineto{\pgfqpoint{1.720041in}{1.417151in}}%
\pgfpathlineto{\pgfqpoint{1.715103in}{1.430399in}}%
\pgfpathlineto{\pgfqpoint{1.714442in}{1.443647in}}%
\pgfpathlineto{\pgfqpoint{1.707431in}{1.456895in}}%
\pgfpathlineto{\pgfqpoint{1.706968in}{1.470143in}}%
\pgfpathlineto{\pgfqpoint{1.704001in}{1.476689in}}%
\pgfpathlineto{\pgfqpoint{1.701411in}{1.483391in}}%
\pgfpathlineto{\pgfqpoint{1.699369in}{1.496639in}}%
\pgfpathlineto{\pgfqpoint{1.696427in}{1.509887in}}%
\pgfpathlineto{\pgfqpoint{1.692115in}{1.523135in}}%
\pgfpathlineto{\pgfqpoint{1.691346in}{1.536383in}}%
\pgfpathlineto{\pgfqpoint{1.685829in}{1.549631in}}%
\pgfpathlineto{\pgfqpoint{1.685377in}{1.562879in}}%
\pgfpathlineto{\pgfqpoint{1.679654in}{1.576127in}}%
\pgfpathlineto{\pgfqpoint{1.678717in}{1.581874in}}%
\pgfpathlineto{\pgfqpoint{1.677699in}{1.589375in}}%
\pgfpathlineto{\pgfqpoint{1.674508in}{1.602623in}}%
\pgfpathlineto{\pgfqpoint{1.670832in}{1.615871in}}%
\pgfpathlineto{\pgfqpoint{1.669487in}{1.629119in}}%
\pgfpathlineto{\pgfqpoint{1.664275in}{1.642367in}}%
\pgfpathlineto{\pgfqpoint{1.663869in}{1.655615in}}%
\pgfpathlineto{\pgfqpoint{1.658997in}{1.668863in}}%
\pgfpathlineto{\pgfqpoint{1.657515in}{1.682111in}}%
\pgfpathlineto{\pgfqpoint{1.653909in}{1.695359in}}%
\pgfpathlineto{\pgfqpoint{1.653434in}{1.696982in}}%
\pgfpathlineto{\pgfqpoint{1.650630in}{1.708607in}}%
\pgfpathlineto{\pgfqpoint{1.649173in}{1.721855in}}%
\pgfpathlineto{\pgfqpoint{1.644790in}{1.735103in}}%
\pgfpathlineto{\pgfqpoint{1.644002in}{1.748351in}}%
\pgfpathlineto{\pgfqpoint{1.639276in}{1.761599in}}%
\pgfpathlineto{\pgfqpoint{1.637917in}{1.774847in}}%
\pgfpathlineto{\pgfqpoint{1.634486in}{1.788095in}}%
\pgfpathlineto{\pgfqpoint{1.631390in}{1.801343in}}%
\pgfpathlineto{\pgfqpoint{1.629586in}{1.814591in}}%
\pgfpathlineto{\pgfqpoint{1.628151in}{1.818422in}}%
\pgfpathlineto{\pgfqpoint{1.625269in}{1.827839in}}%
\pgfpathlineto{\pgfqpoint{1.624575in}{1.841087in}}%
\pgfpathlineto{\pgfqpoint{1.620531in}{1.854335in}}%
\pgfpathlineto{\pgfqpoint{1.619325in}{1.867583in}}%
\pgfpathlineto{\pgfqpoint{1.616084in}{1.880831in}}%
\pgfpathlineto{\pgfqpoint{1.613319in}{1.894079in}}%
\pgfpathlineto{\pgfqpoint{1.611496in}{1.907327in}}%
\pgfpathlineto{\pgfqpoint{1.607235in}{1.920575in}}%
\pgfpathlineto{\pgfqpoint{1.606213in}{1.933823in}}%
\pgfpathlineto{\pgfqpoint{1.602868in}{1.942682in}}%
\pgfpathlineto{\pgfqpoint{1.601516in}{1.947071in}}%
\pgfpathlineto{\pgfqpoint{1.600370in}{1.960319in}}%
\pgfpathlineto{\pgfqpoint{1.597338in}{1.973567in}}%
\pgfpathlineto{\pgfqpoint{1.595065in}{1.986815in}}%
\pgfpathlineto{\pgfqpoint{1.593395in}{2.000063in}}%
\pgfpathlineto{\pgfqpoint{1.589591in}{2.013311in}}%
\pgfpathlineto{\pgfqpoint{1.588743in}{2.026559in}}%
\pgfpathlineto{\pgfqpoint{1.584396in}{2.039807in}}%
\pgfpathlineto{\pgfqpoint{1.583027in}{2.053055in}}%
\pgfpathlineto{\pgfqpoint{1.579136in}{2.066303in}}%
\pgfpathlineto{\pgfqpoint{1.577585in}{2.073211in}}%
\pgfpathlineto{\pgfqpoint{1.576397in}{2.079551in}}%
\pgfpathlineto{\pgfqpoint{1.574528in}{2.092799in}}%
\pgfpathlineto{\pgfqpoint{1.571072in}{2.106047in}}%
\pgfpathlineto{\pgfqpoint{1.570220in}{2.119295in}}%
\pgfpathlineto{\pgfqpoint{1.566151in}{2.132543in}}%
\pgfpathlineto{\pgfqpoint{1.557363in}{2.132543in}}%
\pgfpathlineto{\pgfqpoint{1.561747in}{2.119295in}}%
\pgfpathlineto{\pgfqpoint{1.562567in}{2.106047in}}%
\pgfpathlineto{\pgfqpoint{1.566381in}{2.092799in}}%
\pgfpathlineto{\pgfqpoint{1.568194in}{2.079551in}}%
\pgfpathlineto{\pgfqpoint{1.570915in}{2.066303in}}%
\pgfpathlineto{\pgfqpoint{1.574203in}{2.053055in}}%
\pgfpathlineto{\pgfqpoint{1.575452in}{2.039807in}}%
\pgfpathlineto{\pgfqpoint{1.577585in}{2.032822in}}%
\pgfpathlineto{\pgfqpoint{1.579784in}{2.026559in}}%
\pgfpathlineto{\pgfqpoint{1.580600in}{2.013311in}}%
\pgfpathlineto{\pgfqpoint{1.584806in}{2.000063in}}%
\pgfpathlineto{\pgfqpoint{1.586444in}{1.986815in}}%
\pgfpathlineto{\pgfqpoint{1.588988in}{1.973567in}}%
\pgfpathlineto{\pgfqpoint{1.592090in}{1.960319in}}%
\pgfpathlineto{\pgfqpoint{1.593331in}{1.947071in}}%
\pgfpathlineto{\pgfqpoint{1.597538in}{1.933823in}}%
\pgfpathlineto{\pgfqpoint{1.598307in}{1.920575in}}%
\pgfpathlineto{\pgfqpoint{1.602113in}{1.907327in}}%
\pgfpathlineto{\pgfqpoint{1.602868in}{1.900529in}}%
\pgfpathlineto{\pgfqpoint{1.603746in}{1.894079in}}%
\pgfpathlineto{\pgfqpoint{1.606841in}{1.880831in}}%
\pgfpathlineto{\pgfqpoint{1.610188in}{1.867583in}}%
\pgfpathlineto{\pgfqpoint{1.611498in}{1.854335in}}%
\pgfpathlineto{\pgfqpoint{1.615796in}{1.841087in}}%
\pgfpathlineto{\pgfqpoint{1.616409in}{1.827839in}}%
\pgfpathlineto{\pgfqpoint{1.620723in}{1.814591in}}%
\pgfpathlineto{\pgfqpoint{1.622038in}{1.801343in}}%
\pgfpathlineto{\pgfqpoint{1.624782in}{1.788095in}}%
\pgfpathlineto{\pgfqpoint{1.627646in}{1.774847in}}%
\pgfpathlineto{\pgfqpoint{1.628151in}{1.769369in}}%
\pgfpathlineto{\pgfqpoint{1.629010in}{1.761599in}}%
\pgfpathlineto{\pgfqpoint{1.634026in}{1.748351in}}%
\pgfpathlineto{\pgfqpoint{1.634739in}{1.735103in}}%
\pgfpathlineto{\pgfqpoint{1.639493in}{1.721855in}}%
\pgfpathlineto{\pgfqpoint{1.640792in}{1.708607in}}%
\pgfpathlineto{\pgfqpoint{1.644264in}{1.695359in}}%
\pgfpathlineto{\pgfqpoint{1.647118in}{1.682111in}}%
\pgfpathlineto{\pgfqpoint{1.648400in}{1.668863in}}%
\pgfpathlineto{\pgfqpoint{1.652631in}{1.655615in}}%
\pgfpathlineto{\pgfqpoint{1.652842in}{1.642367in}}%
\pgfpathlineto{\pgfqpoint{1.653434in}{1.640753in}}%
\pgfpathlineto{\pgfqpoint{1.658342in}{1.629119in}}%
\pgfpathlineto{\pgfqpoint{1.659484in}{1.615871in}}%
\pgfpathlineto{\pgfqpoint{1.663488in}{1.602623in}}%
\pgfpathlineto{\pgfqpoint{1.666554in}{1.589375in}}%
\pgfpathlineto{\pgfqpoint{1.668456in}{1.576127in}}%
\pgfpathlineto{\pgfqpoint{1.673324in}{1.562879in}}%
\pgfpathlineto{\pgfqpoint{1.673548in}{1.549631in}}%
\pgfpathlineto{\pgfqpoint{1.678572in}{1.536383in}}%
\pgfpathlineto{\pgfqpoint{1.678717in}{1.531163in}}%
\pgfpathlineto{\pgfqpoint{1.678984in}{1.523135in}}%
\pgfpathlineto{\pgfqpoint{1.683659in}{1.509887in}}%
\pgfpathlineto{\pgfqpoint{1.686366in}{1.496639in}}%
\pgfpathlineto{\pgfqpoint{1.688503in}{1.483391in}}%
\pgfpathlineto{\pgfqpoint{1.693602in}{1.470143in}}%
\pgfpathlineto{\pgfqpoint{1.693778in}{1.456895in}}%
\pgfpathlineto{\pgfqpoint{1.700092in}{1.443647in}}%
\pgfpathlineto{\pgfqpoint{1.700271in}{1.430399in}}%
\pgfpathlineto{\pgfqpoint{1.704001in}{1.419891in}}%
\pgfpathlineto{\pgfqpoint{1.705080in}{1.417151in}}%
\pgfpathlineto{\pgfqpoint{1.707393in}{1.403903in}}%
\pgfpathlineto{\pgfqpoint{1.709910in}{1.390655in}}%
\pgfpathlineto{\pgfqpoint{1.715079in}{1.377407in}}%
\pgfpathlineto{\pgfqpoint{1.715087in}{1.364159in}}%
\pgfpathlineto{\pgfqpoint{1.721652in}{1.350911in}}%
\pgfpathlineto{\pgfqpoint{1.720792in}{1.337663in}}%
\pgfpathlineto{\pgfqpoint{1.726850in}{1.324415in}}%
\pgfpathlineto{\pgfqpoint{1.728110in}{1.311167in}}%
\pgfpathlineto{\pgfqpoint{1.729284in}{1.305521in}}%
\pgfpathlineto{\pgfqpoint{1.731006in}{1.297919in}}%
\pgfpathlineto{\pgfqpoint{1.736090in}{1.284671in}}%
\pgfpathlineto{\pgfqpoint{1.735769in}{1.271423in}}%
\pgfpathlineto{\pgfqpoint{1.743533in}{1.258175in}}%
\pgfpathlineto{\pgfqpoint{1.742076in}{1.244927in}}%
\pgfpathlineto{\pgfqpoint{1.748913in}{1.231679in}}%
\pgfpathlineto{\pgfqpoint{1.748371in}{1.218431in}}%
\pgfpathlineto{\pgfqpoint{1.752013in}{1.205183in}}%
\pgfpathlineto{\pgfqpoint{1.754567in}{1.194620in}}%
\pgfpathlineto{\pgfqpoint{1.755293in}{1.191935in}}%
\pgfpathlineto{\pgfqpoint{1.754567in}{1.183470in}}%
\pgfpathlineto{\pgfqpoint{1.754168in}{1.178687in}}%
\pgfpathlineto{\pgfqpoint{1.754567in}{1.177913in}}%
\pgfpathlineto{\pgfqpoint{1.761509in}{1.165439in}}%
\pgfpathlineto{\pgfqpoint{1.757497in}{1.152191in}}%
\pgfpathlineto{\pgfqpoint{1.766545in}{1.138943in}}%
\pgfpathlineto{\pgfqpoint{1.763809in}{1.125695in}}%
\pgfpathlineto{\pgfqpoint{1.768684in}{1.112447in}}%
\pgfpathlineto{\pgfqpoint{1.769445in}{1.099199in}}%
\pgfpathlineto{\pgfqpoint{1.767772in}{1.085951in}}%
\pgfpathlineto{\pgfqpoint{1.773461in}{1.072703in}}%
\pgfpathlineto{\pgfqpoint{1.765489in}{1.059455in}}%
\pgfpathlineto{\pgfqpoint{1.773442in}{1.046207in}}%
\pgfpathlineto{\pgfqpoint{1.762728in}{1.032959in}}%
\pgfpathlineto{\pgfqpoint{1.769085in}{1.019711in}}%
\pgfpathlineto{\pgfqpoint{1.761945in}{1.006463in}}%
\pgfpathlineto{\pgfqpoint{1.758325in}{0.993215in}}%
\pgfpathlineto{\pgfqpoint{1.757225in}{0.979967in}}%
\pgfpathlineto{\pgfqpoint{1.754567in}{0.977733in}}%
\pgfpathlineto{\pgfqpoint{1.740355in}{0.966719in}}%
\pgfpathlineto{\pgfqpoint{1.745889in}{0.953471in}}%
\pgfpathlineto{\pgfqpoint{1.729284in}{0.945177in}}%
\pgfpathlineto{\pgfqpoint{1.717804in}{0.940223in}}%
\pgfpathlineto{\pgfqpoint{1.721833in}{0.926975in}}%
\pgfpathlineto{\pgfqpoint{1.704001in}{0.920053in}}%
\pgfpathlineto{\pgfqpoint{1.683723in}{0.913727in}}%
\pgfpathlineto{\pgfqpoint{1.678717in}{0.906338in}}%
\pgfpathlineto{\pgfqpoint{1.674328in}{0.900479in}}%
\pgfpathlineto{\pgfqpoint{1.653434in}{0.893243in}}%
\pgfpathlineto{\pgfqpoint{1.630116in}{0.887231in}}%
\pgfpathlineto{\pgfqpoint{1.628151in}{0.886701in}}%
\pgfpathlineto{\pgfqpoint{1.602868in}{0.880692in}}%
\pgfpathlineto{\pgfqpoint{1.577585in}{0.875457in}}%
\pgfpathlineto{\pgfqpoint{1.569315in}{0.873983in}}%
\pgfpathlineto{\pgfqpoint{1.552302in}{0.870555in}}%
\pgfpathlineto{\pgfqpoint{1.527018in}{0.866983in}}%
\pgfpathlineto{\pgfqpoint{1.501735in}{0.864699in}}%
\pgfpathlineto{\pgfqpoint{1.476452in}{0.863428in}}%
\pgfpathlineto{\pgfqpoint{1.451169in}{0.862987in}}%
\pgfpathlineto{\pgfqpoint{1.425886in}{0.863249in}}%
\pgfpathlineto{\pgfqpoint{1.400603in}{0.864132in}}%
\pgfpathlineto{\pgfqpoint{1.375319in}{0.865586in}}%
\pgfpathlineto{\pgfqpoint{1.350036in}{0.867583in}}%
\pgfpathlineto{\pgfqpoint{1.324753in}{0.870118in}}%
\pgfpathlineto{\pgfqpoint{1.299470in}{0.873204in}}%
\pgfpathlineto{\pgfqpoint{1.294090in}{0.873983in}}%
\pgfpathlineto{\pgfqpoint{1.274187in}{0.886235in}}%
\pgfpathlineto{\pgfqpoint{1.273030in}{0.887231in}}%
\pgfpathlineto{\pgfqpoint{1.248904in}{0.891275in}}%
\pgfpathlineto{\pgfqpoint{1.223620in}{0.896375in}}%
\pgfpathlineto{\pgfqpoint{1.206476in}{0.900479in}}%
\pgfpathlineto{\pgfqpoint{1.207459in}{0.913727in}}%
\pgfpathlineto{\pgfqpoint{1.198337in}{0.916727in}}%
\pgfpathlineto{\pgfqpoint{1.173799in}{0.926975in}}%
\pgfpathlineto{\pgfqpoint{1.173054in}{0.927860in}}%
\pgfpathlineto{\pgfqpoint{1.163338in}{0.940223in}}%
\pgfpathlineto{\pgfqpoint{1.160149in}{0.953471in}}%
\pgfpathlineto{\pgfqpoint{1.147771in}{0.961552in}}%
\pgfpathlineto{\pgfqpoint{1.140558in}{0.966719in}}%
\pgfpathlineto{\pgfqpoint{1.147771in}{0.977639in}}%
\pgfpathlineto{\pgfqpoint{1.149239in}{0.979967in}}%
\pgfpathlineto{\pgfqpoint{1.147771in}{0.980835in}}%
\pgfpathlineto{\pgfqpoint{1.129476in}{0.993215in}}%
\pgfpathlineto{\pgfqpoint{1.135315in}{1.006463in}}%
\pgfpathlineto{\pgfqpoint{1.128157in}{1.019711in}}%
\pgfpathlineto{\pgfqpoint{1.126136in}{1.032959in}}%
\pgfpathlineto{\pgfqpoint{1.130659in}{1.046207in}}%
\pgfpathlineto{\pgfqpoint{1.122488in}{1.059273in}}%
\pgfpathlineto{\pgfqpoint{1.122378in}{1.059455in}}%
\pgfpathlineto{\pgfqpoint{1.122488in}{1.059584in}}%
\pgfpathlineto{\pgfqpoint{1.133134in}{1.072703in}}%
\pgfpathlineto{\pgfqpoint{1.126422in}{1.085951in}}%
\pgfpathlineto{\pgfqpoint{1.134298in}{1.099199in}}%
\pgfpathlineto{\pgfqpoint{1.132571in}{1.112447in}}%
\pgfpathlineto{\pgfqpoint{1.133007in}{1.125695in}}%
\pgfpathlineto{\pgfqpoint{1.138659in}{1.138943in}}%
\pgfpathlineto{\pgfqpoint{1.134888in}{1.152191in}}%
\pgfpathlineto{\pgfqpoint{1.143999in}{1.165439in}}%
\pgfpathlineto{\pgfqpoint{1.140299in}{1.178687in}}%
\pgfpathlineto{\pgfqpoint{1.147771in}{1.191665in}}%
\pgfpathlineto{\pgfqpoint{1.147920in}{1.191935in}}%
\pgfpathlineto{\pgfqpoint{1.149028in}{1.205183in}}%
\pgfpathlineto{\pgfqpoint{1.152155in}{1.218431in}}%
\pgfpathlineto{\pgfqpoint{1.157586in}{1.231679in}}%
\pgfpathlineto{\pgfqpoint{1.156006in}{1.244927in}}%
\pgfpathlineto{\pgfqpoint{1.163966in}{1.258175in}}%
\pgfpathlineto{\pgfqpoint{1.162031in}{1.271423in}}%
\pgfpathlineto{\pgfqpoint{1.169045in}{1.284671in}}%
\pgfpathlineto{\pgfqpoint{1.169972in}{1.297919in}}%
\pgfpathlineto{\pgfqpoint{1.173054in}{1.310485in}}%
\pgfpathlineto{\pgfqpoint{1.173209in}{1.311167in}}%
\pgfpathlineto{\pgfqpoint{1.178221in}{1.324415in}}%
\pgfpathlineto{\pgfqpoint{1.178552in}{1.337663in}}%
\pgfpathlineto{\pgfqpoint{1.185647in}{1.350911in}}%
\pgfpathlineto{\pgfqpoint{1.184647in}{1.364159in}}%
\pgfpathlineto{\pgfqpoint{1.191123in}{1.377407in}}%
\pgfpathlineto{\pgfqpoint{1.192263in}{1.390655in}}%
\pgfpathlineto{\pgfqpoint{1.196234in}{1.403903in}}%
\pgfpathlineto{\pgfqpoint{1.198337in}{1.410038in}}%
\pgfpathlineto{\pgfqpoint{1.200454in}{1.417151in}}%
\pgfpathlineto{\pgfqpoint{1.201117in}{1.430399in}}%
\pgfpathlineto{\pgfqpoint{1.207235in}{1.443647in}}%
\pgfpathlineto{\pgfqpoint{1.206982in}{1.456895in}}%
\pgfpathlineto{\pgfqpoint{1.213019in}{1.470143in}}%
\pgfpathlineto{\pgfqpoint{1.213712in}{1.483391in}}%
\pgfpathlineto{\pgfqpoint{1.217575in}{1.496639in}}%
\pgfpathlineto{\pgfqpoint{1.221348in}{1.509887in}}%
\pgfpathlineto{\pgfqpoint{1.223012in}{1.523135in}}%
\pgfpathlineto{\pgfqpoint{1.223620in}{1.524428in}}%
\pgfpathlineto{\pgfqpoint{1.228487in}{1.536383in}}%
\pgfpathlineto{\pgfqpoint{1.228683in}{1.549631in}}%
\pgfpathlineto{\pgfqpoint{1.234204in}{1.562879in}}%
\pgfpathlineto{\pgfqpoint{1.234995in}{1.576127in}}%
\pgfpathlineto{\pgfqpoint{1.239058in}{1.589375in}}%
\pgfpathlineto{\pgfqpoint{1.241649in}{1.602623in}}%
\pgfpathlineto{\pgfqpoint{1.243317in}{1.615871in}}%
\pgfpathlineto{\pgfqpoint{1.248436in}{1.629119in}}%
\pgfpathlineto{\pgfqpoint{1.248904in}{1.639936in}}%
\pgfpathlineto{\pgfqpoint{1.248991in}{1.642367in}}%
\pgfpathlineto{\pgfqpoint{1.254152in}{1.655615in}}%
\pgfpathlineto{\pgfqpoint{1.254942in}{1.668863in}}%
\pgfpathlineto{\pgfqpoint{1.259031in}{1.682111in}}%
\pgfpathlineto{\pgfqpoint{1.261351in}{1.695359in}}%
\pgfpathlineto{\pgfqpoint{1.263642in}{1.708607in}}%
\pgfpathlineto{\pgfqpoint{1.267704in}{1.721855in}}%
\pgfpathlineto{\pgfqpoint{1.268315in}{1.735103in}}%
\pgfpathlineto{\pgfqpoint{1.273549in}{1.748351in}}%
\pgfpathlineto{\pgfqpoint{1.274187in}{1.760059in}}%
\pgfpathlineto{\pgfqpoint{1.274255in}{1.761599in}}%
\pgfpathlineto{\pgfqpoint{1.278324in}{1.774847in}}%
\pgfpathlineto{\pgfqpoint{1.280053in}{1.788095in}}%
\pgfpathlineto{\pgfqpoint{1.282556in}{1.801343in}}%
\pgfpathlineto{\pgfqpoint{1.285965in}{1.814591in}}%
\pgfpathlineto{\pgfqpoint{1.287092in}{1.827839in}}%
\pgfpathlineto{\pgfqpoint{1.291760in}{1.841087in}}%
\pgfpathlineto{\pgfqpoint{1.292334in}{1.854335in}}%
\pgfpathlineto{\pgfqpoint{1.296930in}{1.867583in}}%
\pgfpathlineto{\pgfqpoint{1.298644in}{1.880831in}}%
\pgfpathlineto{\pgfqpoint{1.299470in}{1.884044in}}%
\pgfpathlineto{\pgfqpoint{1.301665in}{1.894079in}}%
\pgfpathlineto{\pgfqpoint{1.304377in}{1.907327in}}%
\pgfpathlineto{\pgfqpoint{1.305700in}{1.920575in}}%
\pgfpathlineto{\pgfqpoint{1.309577in}{1.933823in}}%
\pgfpathlineto{\pgfqpoint{1.310200in}{1.947071in}}%
\pgfpathlineto{\pgfqpoint{1.314617in}{1.960319in}}%
\pgfpathlineto{\pgfqpoint{1.315715in}{1.973567in}}%
\pgfpathlineto{\pgfqpoint{1.319295in}{1.986815in}}%
\pgfpathlineto{\pgfqpoint{1.322135in}{2.000063in}}%
\pgfpathlineto{\pgfqpoint{1.324437in}{2.013311in}}%
\pgfpathlineto{\pgfqpoint{1.324753in}{2.014286in}}%
\pgfpathlineto{\pgfqpoint{1.328125in}{2.026559in}}%
\pgfpathlineto{\pgfqpoint{1.328931in}{2.039807in}}%
\pgfpathlineto{\pgfqpoint{1.332950in}{2.053055in}}%
\pgfpathlineto{\pgfqpoint{1.333751in}{2.066303in}}%
\pgfpathlineto{\pgfqpoint{1.337520in}{2.079551in}}%
\pgfpathlineto{\pgfqpoint{1.339437in}{2.092799in}}%
\pgfpathlineto{\pgfqpoint{1.341962in}{2.106047in}}%
\pgfpathlineto{\pgfqpoint{1.345676in}{2.119295in}}%
\pgfpathlineto{\pgfqpoint{1.347117in}{2.132543in}}%
\pgfpathlineto{\pgfqpoint{1.338237in}{2.132543in}}%
\pgfpathlineto{\pgfqpoint{1.336937in}{2.119295in}}%
\pgfpathlineto{\pgfqpoint{1.333384in}{2.106047in}}%
\pgfpathlineto{\pgfqpoint{1.331161in}{2.092799in}}%
\pgfpathlineto{\pgfqpoint{1.329197in}{2.079551in}}%
\pgfpathlineto{\pgfqpoint{1.325786in}{2.066303in}}%
\pgfpathlineto{\pgfqpoint{1.324903in}{2.053055in}}%
\pgfpathlineto{\pgfqpoint{1.324753in}{2.052544in}}%
\pgfpathlineto{\pgfqpoint{1.320325in}{2.039807in}}%
\pgfpathlineto{\pgfqpoint{1.319401in}{2.026559in}}%
\pgfpathlineto{\pgfqpoint{1.315274in}{2.013311in}}%
\pgfpathlineto{\pgfqpoint{1.313245in}{2.000063in}}%
\pgfpathlineto{\pgfqpoint{1.310462in}{1.986815in}}%
\pgfpathlineto{\pgfqpoint{1.307274in}{1.973567in}}%
\pgfpathlineto{\pgfqpoint{1.306096in}{1.960319in}}%
\pgfpathlineto{\pgfqpoint{1.302008in}{1.947071in}}%
\pgfpathlineto{\pgfqpoint{1.301351in}{1.933823in}}%
\pgfpathlineto{\pgfqpoint{1.299470in}{1.927225in}}%
\pgfpathlineto{\pgfqpoint{1.297124in}{1.920575in}}%
\pgfpathlineto{\pgfqpoint{1.295631in}{1.907327in}}%
\pgfpathlineto{\pgfqpoint{1.292314in}{1.894079in}}%
\pgfpathlineto{\pgfqpoint{1.289200in}{1.880831in}}%
\pgfpathlineto{\pgfqpoint{1.287430in}{1.867583in}}%
\pgfpathlineto{\pgfqpoint{1.283273in}{1.854335in}}%
\pgfpathlineto{\pgfqpoint{1.282633in}{1.841087in}}%
\pgfpathlineto{\pgfqpoint{1.278221in}{1.827839in}}%
\pgfpathlineto{\pgfqpoint{1.277171in}{1.814591in}}%
\pgfpathlineto{\pgfqpoint{1.274187in}{1.803018in}}%
\pgfpathlineto{\pgfqpoint{1.273641in}{1.801343in}}%
\pgfpathlineto{\pgfqpoint{1.270805in}{1.788095in}}%
\pgfpathlineto{\pgfqpoint{1.268528in}{1.774847in}}%
\pgfpathlineto{\pgfqpoint{1.264013in}{1.761599in}}%
\pgfpathlineto{\pgfqpoint{1.263163in}{1.748351in}}%
\pgfpathlineto{\pgfqpoint{1.258303in}{1.735103in}}%
\pgfpathlineto{\pgfqpoint{1.257697in}{1.721855in}}%
\pgfpathlineto{\pgfqpoint{1.253717in}{1.708607in}}%
\pgfpathlineto{\pgfqpoint{1.251639in}{1.695359in}}%
\pgfpathlineto{\pgfqpoint{1.249151in}{1.682111in}}%
\pgfpathlineto{\pgfqpoint{1.248904in}{1.681268in}}%
\pgfpathlineto{\pgfqpoint{1.244589in}{1.668863in}}%
\pgfpathlineto{\pgfqpoint{1.243410in}{1.655615in}}%
\pgfpathlineto{\pgfqpoint{1.237630in}{1.642367in}}%
\pgfpathlineto{\pgfqpoint{1.236980in}{1.629119in}}%
\pgfpathlineto{\pgfqpoint{1.232025in}{1.615871in}}%
\pgfpathlineto{\pgfqpoint{1.230507in}{1.602623in}}%
\pgfpathlineto{\pgfqpoint{1.227730in}{1.589375in}}%
\pgfpathlineto{\pgfqpoint{1.223983in}{1.576127in}}%
\pgfpathlineto{\pgfqpoint{1.223620in}{1.571737in}}%
\pgfpathlineto{\pgfqpoint{1.222713in}{1.562879in}}%
\pgfpathlineto{\pgfqpoint{1.216523in}{1.549631in}}%
\pgfpathlineto{\pgfqpoint{1.216100in}{1.536383in}}%
\pgfpathlineto{\pgfqpoint{1.210023in}{1.523135in}}%
\pgfpathlineto{\pgfqpoint{1.208465in}{1.509887in}}%
\pgfpathlineto{\pgfqpoint{1.204511in}{1.496639in}}%
\pgfpathlineto{\pgfqpoint{1.200971in}{1.483391in}}%
\pgfpathlineto{\pgfqpoint{1.199827in}{1.470143in}}%
\pgfpathlineto{\pgfqpoint{1.198337in}{1.466767in}}%
\pgfpathlineto{\pgfqpoint{1.193379in}{1.456895in}}%
\pgfpathlineto{\pgfqpoint{1.193249in}{1.443647in}}%
\pgfpathlineto{\pgfqpoint{1.186369in}{1.430399in}}%
\pgfpathlineto{\pgfqpoint{1.185581in}{1.417151in}}%
\pgfpathlineto{\pgfqpoint{1.180856in}{1.403903in}}%
\pgfpathlineto{\pgfqpoint{1.177205in}{1.390655in}}%
\pgfpathlineto{\pgfqpoint{1.175469in}{1.377407in}}%
\pgfpathlineto{\pgfqpoint{1.173054in}{1.372259in}}%
\pgfpathlineto{\pgfqpoint{1.168833in}{1.364159in}}%
\pgfpathlineto{\pgfqpoint{1.169307in}{1.350911in}}%
\pgfpathlineto{\pgfqpoint{1.161425in}{1.337663in}}%
\pgfpathlineto{\pgfqpoint{1.160806in}{1.324415in}}%
\pgfpathlineto{\pgfqpoint{1.154825in}{1.311167in}}%
\pgfpathlineto{\pgfqpoint{1.151793in}{1.297919in}}%
\pgfpathlineto{\pgfqpoint{1.149959in}{1.284671in}}%
\pgfpathlineto{\pgfqpoint{1.147771in}{1.280391in}}%
\pgfpathlineto{\pgfqpoint{1.142797in}{1.271423in}}%
\pgfpathlineto{\pgfqpoint{1.143921in}{1.258175in}}%
\pgfpathlineto{\pgfqpoint{1.135191in}{1.244927in}}%
\pgfpathlineto{\pgfqpoint{1.136334in}{1.231679in}}%
\pgfpathlineto{\pgfqpoint{1.129710in}{1.218431in}}%
\pgfpathlineto{\pgfqpoint{1.126440in}{1.205183in}}%
\pgfpathlineto{\pgfqpoint{1.123817in}{1.191935in}}%
\pgfpathlineto{\pgfqpoint{1.122488in}{1.189568in}}%
\pgfpathlineto{\pgfqpoint{1.116016in}{1.178687in}}%
\pgfpathlineto{\pgfqpoint{1.118317in}{1.165439in}}%
\pgfpathlineto{\pgfqpoint{1.108341in}{1.152191in}}%
\pgfpathlineto{\pgfqpoint{1.111248in}{1.138943in}}%
\pgfpathlineto{\pgfqpoint{1.103750in}{1.125695in}}%
\pgfpathlineto{\pgfqpoint{1.102978in}{1.112447in}}%
\pgfpathlineto{\pgfqpoint{1.102161in}{1.099199in}}%
\pgfpathlineto{\pgfqpoint{1.097204in}{1.091140in}}%
\pgfpathlineto{\pgfqpoint{1.093881in}{1.085951in}}%
\pgfpathlineto{\pgfqpoint{1.097204in}{1.075592in}}%
\pgfpathlineto{\pgfqpoint{1.098168in}{1.072703in}}%
\pgfpathlineto{\pgfqpoint{1.097204in}{1.071609in}}%
\pgfpathlineto{\pgfqpoint{1.085921in}{1.059455in}}%
\pgfpathlineto{\pgfqpoint{1.092584in}{1.046207in}}%
\pgfpathlineto{\pgfqpoint{1.084858in}{1.032959in}}%
\pgfpathlineto{\pgfqpoint{1.085978in}{1.019711in}}%
\pgfpathlineto{\pgfqpoint{1.088036in}{1.006463in}}%
\pgfpathlineto{\pgfqpoint{1.081168in}{0.993215in}}%
\pgfpathlineto{\pgfqpoint{1.094963in}{0.979967in}}%
\pgfpathlineto{\pgfqpoint{1.083108in}{0.966719in}}%
\pgfpathlineto{\pgfqpoint{1.097204in}{0.954111in}}%
\pgfpathlineto{\pgfqpoint{1.097977in}{0.953471in}}%
\pgfpathlineto{\pgfqpoint{1.097204in}{0.951337in}}%
\pgfpathlineto{\pgfqpoint{1.092750in}{0.940223in}}%
\pgfpathlineto{\pgfqpoint{1.097204in}{0.932832in}}%
\pgfpathlineto{\pgfqpoint{1.100883in}{0.926975in}}%
\pgfpathlineto{\pgfqpoint{1.117682in}{0.913727in}}%
\pgfpathlineto{\pgfqpoint{1.113102in}{0.900479in}}%
\pgfpathlineto{\pgfqpoint{1.122488in}{0.896764in}}%
\pgfpathlineto{\pgfqpoint{1.147771in}{0.888253in}}%
\pgfpathlineto{\pgfqpoint{1.151281in}{0.887231in}}%
\pgfpathlineto{\pgfqpoint{1.147771in}{0.877399in}}%
\pgfpathlineto{\pgfqpoint{1.146457in}{0.873983in}}%
\pgfpathlineto{\pgfqpoint{1.147771in}{0.873606in}}%
\pgfpathlineto{\pgfqpoint{1.173054in}{0.867012in}}%
\pgfpathlineto{\pgfqpoint{1.198337in}{0.861229in}}%
\pgfpathlineto{\pgfqpoint{1.200733in}{0.860735in}}%
\pgfpathlineto{\pgfqpoint{1.223620in}{0.852179in}}%
\pgfpathlineto{\pgfqpoint{1.241356in}{0.847487in}}%
\pgfpathlineto{\pgfqpoint{1.248904in}{0.845033in}}%
\pgfpathlineto{\pgfqpoint{1.274187in}{0.837757in}}%
\pgfpathlineto{\pgfqpoint{1.288142in}{0.834239in}}%
\pgfpathlineto{\pgfqpoint{1.299470in}{0.833123in}}%
\pgfpathclose%
\pgfpathmoveto{\pgfqpoint{1.122488in}{2.088463in}}%
\pgfpathlineto{\pgfqpoint{1.127644in}{2.092799in}}%
\pgfpathlineto{\pgfqpoint{1.126656in}{2.106047in}}%
\pgfpathlineto{\pgfqpoint{1.138776in}{2.119295in}}%
\pgfpathlineto{\pgfqpoint{1.147771in}{2.127648in}}%
\pgfpathlineto{\pgfqpoint{1.149378in}{2.132543in}}%
\pgfpathlineto{\pgfqpoint{1.147771in}{2.132543in}}%
\pgfpathlineto{\pgfqpoint{1.122488in}{2.132543in}}%
\pgfpathlineto{\pgfqpoint{1.105823in}{2.132543in}}%
\pgfpathlineto{\pgfqpoint{1.111320in}{2.119295in}}%
\pgfpathlineto{\pgfqpoint{1.115163in}{2.106047in}}%
\pgfpathlineto{\pgfqpoint{1.110426in}{2.092799in}}%
\pgfpathlineto{\pgfqpoint{1.122488in}{2.088463in}}%
\pgfpathclose%
\pgfpathmoveto{\pgfqpoint{1.779850in}{2.115676in}}%
\pgfpathlineto{\pgfqpoint{1.786670in}{2.119295in}}%
\pgfpathlineto{\pgfqpoint{1.781130in}{2.132543in}}%
\pgfpathlineto{\pgfqpoint{1.779850in}{2.132543in}}%
\pgfpathlineto{\pgfqpoint{1.777927in}{2.132543in}}%
\pgfpathlineto{\pgfqpoint{1.774838in}{2.119295in}}%
\pgfpathlineto{\pgfqpoint{1.779850in}{2.115676in}}%
\pgfpathclose%
\pgfusepath{fill}%
\end{pgfscope}%
\begin{pgfscope}%
\pgfpathrectangle{\pgfqpoint{0.540975in}{0.476543in}}{\pgfqpoint{1.820388in}{1.656000in}}%
\pgfusepath{clip}%
\pgfsetbuttcap%
\pgfsetroundjoin%
\definecolor{currentfill}{rgb}{0.607151,0.800231,0.886044}%
\pgfsetfillcolor{currentfill}%
\pgfsetlinewidth{0.000000pt}%
\definecolor{currentstroke}{rgb}{0.000000,0.000000,0.000000}%
\pgfsetstrokecolor{currentstroke}%
\pgfsetdash{}{0pt}%
\pgfpathmoveto{\pgfqpoint{1.299470in}{0.873204in}}%
\pgfpathlineto{\pgfqpoint{1.324753in}{0.870118in}}%
\pgfpathlineto{\pgfqpoint{1.350036in}{0.867583in}}%
\pgfpathlineto{\pgfqpoint{1.375319in}{0.865586in}}%
\pgfpathlineto{\pgfqpoint{1.400603in}{0.864132in}}%
\pgfpathlineto{\pgfqpoint{1.425886in}{0.863249in}}%
\pgfpathlineto{\pgfqpoint{1.451169in}{0.862987in}}%
\pgfpathlineto{\pgfqpoint{1.476452in}{0.863428in}}%
\pgfpathlineto{\pgfqpoint{1.501735in}{0.864699in}}%
\pgfpathlineto{\pgfqpoint{1.527018in}{0.866983in}}%
\pgfpathlineto{\pgfqpoint{1.552302in}{0.870555in}}%
\pgfpathlineto{\pgfqpoint{1.569315in}{0.873983in}}%
\pgfpathlineto{\pgfqpoint{1.577585in}{0.875457in}}%
\pgfpathlineto{\pgfqpoint{1.602868in}{0.880692in}}%
\pgfpathlineto{\pgfqpoint{1.628151in}{0.886701in}}%
\pgfpathlineto{\pgfqpoint{1.630116in}{0.887231in}}%
\pgfpathlineto{\pgfqpoint{1.653434in}{0.893243in}}%
\pgfpathlineto{\pgfqpoint{1.674328in}{0.900479in}}%
\pgfpathlineto{\pgfqpoint{1.678717in}{0.906338in}}%
\pgfpathlineto{\pgfqpoint{1.683723in}{0.913727in}}%
\pgfpathlineto{\pgfqpoint{1.704001in}{0.920053in}}%
\pgfpathlineto{\pgfqpoint{1.721833in}{0.926975in}}%
\pgfpathlineto{\pgfqpoint{1.717804in}{0.940223in}}%
\pgfpathlineto{\pgfqpoint{1.729284in}{0.945177in}}%
\pgfpathlineto{\pgfqpoint{1.745889in}{0.953471in}}%
\pgfpathlineto{\pgfqpoint{1.740355in}{0.966719in}}%
\pgfpathlineto{\pgfqpoint{1.754567in}{0.977733in}}%
\pgfpathlineto{\pgfqpoint{1.757225in}{0.979967in}}%
\pgfpathlineto{\pgfqpoint{1.758325in}{0.993215in}}%
\pgfpathlineto{\pgfqpoint{1.761945in}{1.006463in}}%
\pgfpathlineto{\pgfqpoint{1.769085in}{1.019711in}}%
\pgfpathlineto{\pgfqpoint{1.762728in}{1.032959in}}%
\pgfpathlineto{\pgfqpoint{1.773442in}{1.046207in}}%
\pgfpathlineto{\pgfqpoint{1.765489in}{1.059455in}}%
\pgfpathlineto{\pgfqpoint{1.773461in}{1.072703in}}%
\pgfpathlineto{\pgfqpoint{1.767772in}{1.085951in}}%
\pgfpathlineto{\pgfqpoint{1.769445in}{1.099199in}}%
\pgfpathlineto{\pgfqpoint{1.768684in}{1.112447in}}%
\pgfpathlineto{\pgfqpoint{1.763809in}{1.125695in}}%
\pgfpathlineto{\pgfqpoint{1.766545in}{1.138943in}}%
\pgfpathlineto{\pgfqpoint{1.757497in}{1.152191in}}%
\pgfpathlineto{\pgfqpoint{1.761509in}{1.165439in}}%
\pgfpathlineto{\pgfqpoint{1.754567in}{1.177913in}}%
\pgfpathlineto{\pgfqpoint{1.754168in}{1.178687in}}%
\pgfpathlineto{\pgfqpoint{1.754567in}{1.183470in}}%
\pgfpathlineto{\pgfqpoint{1.755293in}{1.191935in}}%
\pgfpathlineto{\pgfqpoint{1.754567in}{1.194620in}}%
\pgfpathlineto{\pgfqpoint{1.752013in}{1.205183in}}%
\pgfpathlineto{\pgfqpoint{1.748371in}{1.218431in}}%
\pgfpathlineto{\pgfqpoint{1.748913in}{1.231679in}}%
\pgfpathlineto{\pgfqpoint{1.742076in}{1.244927in}}%
\pgfpathlineto{\pgfqpoint{1.743533in}{1.258175in}}%
\pgfpathlineto{\pgfqpoint{1.735769in}{1.271423in}}%
\pgfpathlineto{\pgfqpoint{1.736090in}{1.284671in}}%
\pgfpathlineto{\pgfqpoint{1.731006in}{1.297919in}}%
\pgfpathlineto{\pgfqpoint{1.729284in}{1.305521in}}%
\pgfpathlineto{\pgfqpoint{1.728110in}{1.311167in}}%
\pgfpathlineto{\pgfqpoint{1.726850in}{1.324415in}}%
\pgfpathlineto{\pgfqpoint{1.720792in}{1.337663in}}%
\pgfpathlineto{\pgfqpoint{1.721652in}{1.350911in}}%
\pgfpathlineto{\pgfqpoint{1.715087in}{1.364159in}}%
\pgfpathlineto{\pgfqpoint{1.715079in}{1.377407in}}%
\pgfpathlineto{\pgfqpoint{1.709910in}{1.390655in}}%
\pgfpathlineto{\pgfqpoint{1.707393in}{1.403903in}}%
\pgfpathlineto{\pgfqpoint{1.705080in}{1.417151in}}%
\pgfpathlineto{\pgfqpoint{1.704001in}{1.419891in}}%
\pgfpathlineto{\pgfqpoint{1.700271in}{1.430399in}}%
\pgfpathlineto{\pgfqpoint{1.700092in}{1.443647in}}%
\pgfpathlineto{\pgfqpoint{1.693778in}{1.456895in}}%
\pgfpathlineto{\pgfqpoint{1.693602in}{1.470143in}}%
\pgfpathlineto{\pgfqpoint{1.688503in}{1.483391in}}%
\pgfpathlineto{\pgfqpoint{1.686366in}{1.496639in}}%
\pgfpathlineto{\pgfqpoint{1.683659in}{1.509887in}}%
\pgfpathlineto{\pgfqpoint{1.678984in}{1.523135in}}%
\pgfpathlineto{\pgfqpoint{1.678717in}{1.531163in}}%
\pgfpathlineto{\pgfqpoint{1.678572in}{1.536383in}}%
\pgfpathlineto{\pgfqpoint{1.673548in}{1.549631in}}%
\pgfpathlineto{\pgfqpoint{1.673324in}{1.562879in}}%
\pgfpathlineto{\pgfqpoint{1.668456in}{1.576127in}}%
\pgfpathlineto{\pgfqpoint{1.666554in}{1.589375in}}%
\pgfpathlineto{\pgfqpoint{1.663488in}{1.602623in}}%
\pgfpathlineto{\pgfqpoint{1.659484in}{1.615871in}}%
\pgfpathlineto{\pgfqpoint{1.658342in}{1.629119in}}%
\pgfpathlineto{\pgfqpoint{1.653434in}{1.640753in}}%
\pgfpathlineto{\pgfqpoint{1.652842in}{1.642367in}}%
\pgfpathlineto{\pgfqpoint{1.652631in}{1.655615in}}%
\pgfpathlineto{\pgfqpoint{1.648400in}{1.668863in}}%
\pgfpathlineto{\pgfqpoint{1.647118in}{1.682111in}}%
\pgfpathlineto{\pgfqpoint{1.644264in}{1.695359in}}%
\pgfpathlineto{\pgfqpoint{1.640792in}{1.708607in}}%
\pgfpathlineto{\pgfqpoint{1.639493in}{1.721855in}}%
\pgfpathlineto{\pgfqpoint{1.634739in}{1.735103in}}%
\pgfpathlineto{\pgfqpoint{1.634026in}{1.748351in}}%
\pgfpathlineto{\pgfqpoint{1.629010in}{1.761599in}}%
\pgfpathlineto{\pgfqpoint{1.628151in}{1.769369in}}%
\pgfpathlineto{\pgfqpoint{1.627646in}{1.774847in}}%
\pgfpathlineto{\pgfqpoint{1.624782in}{1.788095in}}%
\pgfpathlineto{\pgfqpoint{1.622038in}{1.801343in}}%
\pgfpathlineto{\pgfqpoint{1.620723in}{1.814591in}}%
\pgfpathlineto{\pgfqpoint{1.616409in}{1.827839in}}%
\pgfpathlineto{\pgfqpoint{1.615796in}{1.841087in}}%
\pgfpathlineto{\pgfqpoint{1.611498in}{1.854335in}}%
\pgfpathlineto{\pgfqpoint{1.610188in}{1.867583in}}%
\pgfpathlineto{\pgfqpoint{1.606841in}{1.880831in}}%
\pgfpathlineto{\pgfqpoint{1.603746in}{1.894079in}}%
\pgfpathlineto{\pgfqpoint{1.602868in}{1.900529in}}%
\pgfpathlineto{\pgfqpoint{1.602113in}{1.907327in}}%
\pgfpathlineto{\pgfqpoint{1.598307in}{1.920575in}}%
\pgfpathlineto{\pgfqpoint{1.597538in}{1.933823in}}%
\pgfpathlineto{\pgfqpoint{1.593331in}{1.947071in}}%
\pgfpathlineto{\pgfqpoint{1.592090in}{1.960319in}}%
\pgfpathlineto{\pgfqpoint{1.588988in}{1.973567in}}%
\pgfpathlineto{\pgfqpoint{1.586444in}{1.986815in}}%
\pgfpathlineto{\pgfqpoint{1.584806in}{2.000063in}}%
\pgfpathlineto{\pgfqpoint{1.580600in}{2.013311in}}%
\pgfpathlineto{\pgfqpoint{1.579784in}{2.026559in}}%
\pgfpathlineto{\pgfqpoint{1.577585in}{2.032822in}}%
\pgfpathlineto{\pgfqpoint{1.575452in}{2.039807in}}%
\pgfpathlineto{\pgfqpoint{1.574203in}{2.053055in}}%
\pgfpathlineto{\pgfqpoint{1.570915in}{2.066303in}}%
\pgfpathlineto{\pgfqpoint{1.568194in}{2.079551in}}%
\pgfpathlineto{\pgfqpoint{1.566381in}{2.092799in}}%
\pgfpathlineto{\pgfqpoint{1.562567in}{2.106047in}}%
\pgfpathlineto{\pgfqpoint{1.561747in}{2.119295in}}%
\pgfpathlineto{\pgfqpoint{1.557363in}{2.132543in}}%
\pgfpathlineto{\pgfqpoint{1.552302in}{2.132543in}}%
\pgfpathlineto{\pgfqpoint{1.548903in}{2.132543in}}%
\pgfpathlineto{\pgfqpoint{1.552302in}{2.122117in}}%
\pgfpathlineto{\pgfqpoint{1.553274in}{2.119295in}}%
\pgfpathlineto{\pgfqpoint{1.554063in}{2.106047in}}%
\pgfpathlineto{\pgfqpoint{1.558233in}{2.092799in}}%
\pgfpathlineto{\pgfqpoint{1.559992in}{2.079551in}}%
\pgfpathlineto{\pgfqpoint{1.562997in}{2.066303in}}%
\pgfpathlineto{\pgfqpoint{1.566318in}{2.053055in}}%
\pgfpathlineto{\pgfqpoint{1.567667in}{2.039807in}}%
\pgfpathlineto{\pgfqpoint{1.571839in}{2.026559in}}%
\pgfpathlineto{\pgfqpoint{1.572478in}{2.013311in}}%
\pgfpathlineto{\pgfqpoint{1.576399in}{2.000063in}}%
\pgfpathlineto{\pgfqpoint{1.577585in}{1.988768in}}%
\pgfpathlineto{\pgfqpoint{1.577822in}{1.986815in}}%
\pgfpathlineto{\pgfqpoint{1.580637in}{1.973567in}}%
\pgfpathlineto{\pgfqpoint{1.583810in}{1.960319in}}%
\pgfpathlineto{\pgfqpoint{1.585145in}{1.947071in}}%
\pgfpathlineto{\pgfqpoint{1.589541in}{1.933823in}}%
\pgfpathlineto{\pgfqpoint{1.590236in}{1.920575in}}%
\pgfpathlineto{\pgfqpoint{1.594301in}{1.907327in}}%
\pgfpathlineto{\pgfqpoint{1.595649in}{1.894079in}}%
\pgfpathlineto{\pgfqpoint{1.598462in}{1.880831in}}%
\pgfpathlineto{\pgfqpoint{1.601321in}{1.867583in}}%
\pgfpathlineto{\pgfqpoint{1.602525in}{1.854335in}}%
\pgfpathlineto{\pgfqpoint{1.602868in}{1.853196in}}%
\pgfpathlineto{\pgfqpoint{1.607017in}{1.841087in}}%
\pgfpathlineto{\pgfqpoint{1.607549in}{1.827839in}}%
\pgfpathlineto{\pgfqpoint{1.612162in}{1.814591in}}%
\pgfpathlineto{\pgfqpoint{1.613328in}{1.801343in}}%
\pgfpathlineto{\pgfqpoint{1.616287in}{1.788095in}}%
\pgfpathlineto{\pgfqpoint{1.619071in}{1.774847in}}%
\pgfpathlineto{\pgfqpoint{1.620355in}{1.761599in}}%
\pgfpathlineto{\pgfqpoint{1.624682in}{1.748351in}}%
\pgfpathlineto{\pgfqpoint{1.625209in}{1.735103in}}%
\pgfpathlineto{\pgfqpoint{1.628151in}{1.726262in}}%
\pgfpathlineto{\pgfqpoint{1.629813in}{1.721855in}}%
\pgfpathlineto{\pgfqpoint{1.630953in}{1.708607in}}%
\pgfpathlineto{\pgfqpoint{1.634715in}{1.695359in}}%
\pgfpathlineto{\pgfqpoint{1.637471in}{1.682111in}}%
\pgfpathlineto{\pgfqpoint{1.638806in}{1.668863in}}%
\pgfpathlineto{\pgfqpoint{1.643094in}{1.655615in}}%
\pgfpathlineto{\pgfqpoint{1.643147in}{1.642367in}}%
\pgfpathlineto{\pgfqpoint{1.648111in}{1.629119in}}%
\pgfpathlineto{\pgfqpoint{1.648872in}{1.615871in}}%
\pgfpathlineto{\pgfqpoint{1.652594in}{1.602623in}}%
\pgfpathlineto{\pgfqpoint{1.653434in}{1.598235in}}%
\pgfpathlineto{\pgfqpoint{1.655409in}{1.589375in}}%
\pgfpathlineto{\pgfqpoint{1.657421in}{1.576127in}}%
\pgfpathlineto{\pgfqpoint{1.662325in}{1.562879in}}%
\pgfpathlineto{\pgfqpoint{1.662370in}{1.549631in}}%
\pgfpathlineto{\pgfqpoint{1.667590in}{1.536383in}}%
\pgfpathlineto{\pgfqpoint{1.667602in}{1.523135in}}%
\pgfpathlineto{\pgfqpoint{1.671892in}{1.509887in}}%
\pgfpathlineto{\pgfqpoint{1.673989in}{1.496639in}}%
\pgfpathlineto{\pgfqpoint{1.675952in}{1.483391in}}%
\pgfpathlineto{\pgfqpoint{1.678717in}{1.475235in}}%
\pgfpathlineto{\pgfqpoint{1.680670in}{1.470143in}}%
\pgfpathlineto{\pgfqpoint{1.680614in}{1.456895in}}%
\pgfpathlineto{\pgfqpoint{1.687080in}{1.443647in}}%
\pgfpathlineto{\pgfqpoint{1.686799in}{1.430399in}}%
\pgfpathlineto{\pgfqpoint{1.691705in}{1.417151in}}%
\pgfpathlineto{\pgfqpoint{1.693277in}{1.403903in}}%
\pgfpathlineto{\pgfqpoint{1.695632in}{1.390655in}}%
\pgfpathlineto{\pgfqpoint{1.700059in}{1.377407in}}%
\pgfpathlineto{\pgfqpoint{1.699803in}{1.364159in}}%
\pgfpathlineto{\pgfqpoint{1.704001in}{1.354990in}}%
\pgfpathlineto{\pgfqpoint{1.706037in}{1.350911in}}%
\pgfpathlineto{\pgfqpoint{1.704491in}{1.337663in}}%
\pgfpathlineto{\pgfqpoint{1.710700in}{1.324415in}}%
\pgfpathlineto{\pgfqpoint{1.711145in}{1.311167in}}%
\pgfpathlineto{\pgfqpoint{1.713917in}{1.297919in}}%
\pgfpathlineto{\pgfqpoint{1.718002in}{1.284671in}}%
\pgfpathlineto{\pgfqpoint{1.717377in}{1.271423in}}%
\pgfpathlineto{\pgfqpoint{1.724408in}{1.258175in}}%
\pgfpathlineto{\pgfqpoint{1.722207in}{1.244927in}}%
\pgfpathlineto{\pgfqpoint{1.728732in}{1.231679in}}%
\pgfpathlineto{\pgfqpoint{1.726973in}{1.218431in}}%
\pgfpathlineto{\pgfqpoint{1.729284in}{1.209582in}}%
\pgfpathlineto{\pgfqpoint{1.730485in}{1.205183in}}%
\pgfpathlineto{\pgfqpoint{1.732429in}{1.191935in}}%
\pgfpathlineto{\pgfqpoint{1.730771in}{1.178687in}}%
\pgfpathlineto{\pgfqpoint{1.736948in}{1.165439in}}%
\pgfpathlineto{\pgfqpoint{1.731762in}{1.152191in}}%
\pgfpathlineto{\pgfqpoint{1.740109in}{1.138943in}}%
\pgfpathlineto{\pgfqpoint{1.735394in}{1.125695in}}%
\pgfpathlineto{\pgfqpoint{1.739877in}{1.112447in}}%
\pgfpathlineto{\pgfqpoint{1.738157in}{1.099199in}}%
\pgfpathlineto{\pgfqpoint{1.735815in}{1.085951in}}%
\pgfpathlineto{\pgfqpoint{1.739202in}{1.072703in}}%
\pgfpathlineto{\pgfqpoint{1.729284in}{1.059515in}}%
\pgfpathlineto{\pgfqpoint{1.729236in}{1.059455in}}%
\pgfpathlineto{\pgfqpoint{1.729284in}{1.059354in}}%
\pgfpathlineto{\pgfqpoint{1.735737in}{1.046207in}}%
\pgfpathlineto{\pgfqpoint{1.729284in}{1.040025in}}%
\pgfpathlineto{\pgfqpoint{1.721005in}{1.032959in}}%
\pgfpathlineto{\pgfqpoint{1.726461in}{1.019711in}}%
\pgfpathlineto{\pgfqpoint{1.713525in}{1.006463in}}%
\pgfpathlineto{\pgfqpoint{1.708083in}{0.993215in}}%
\pgfpathlineto{\pgfqpoint{1.704001in}{0.985872in}}%
\pgfpathlineto{\pgfqpoint{1.700068in}{0.979967in}}%
\pgfpathlineto{\pgfqpoint{1.678717in}{0.966944in}}%
\pgfpathlineto{\pgfqpoint{1.678306in}{0.966719in}}%
\pgfpathlineto{\pgfqpoint{1.677723in}{0.953471in}}%
\pgfpathlineto{\pgfqpoint{1.653434in}{0.945269in}}%
\pgfpathlineto{\pgfqpoint{1.635555in}{0.940223in}}%
\pgfpathlineto{\pgfqpoint{1.635308in}{0.926975in}}%
\pgfpathlineto{\pgfqpoint{1.628151in}{0.925252in}}%
\pgfpathlineto{\pgfqpoint{1.602868in}{0.920232in}}%
\pgfpathlineto{\pgfqpoint{1.577585in}{0.916320in}}%
\pgfpathlineto{\pgfqpoint{1.555911in}{0.913727in}}%
\pgfpathlineto{\pgfqpoint{1.552302in}{0.911241in}}%
\pgfpathlineto{\pgfqpoint{1.532931in}{0.900479in}}%
\pgfpathlineto{\pgfqpoint{1.527018in}{0.899801in}}%
\pgfpathlineto{\pgfqpoint{1.501735in}{0.897740in}}%
\pgfpathlineto{\pgfqpoint{1.476452in}{0.896495in}}%
\pgfpathlineto{\pgfqpoint{1.451169in}{0.895966in}}%
\pgfpathlineto{\pgfqpoint{1.425886in}{0.896091in}}%
\pgfpathlineto{\pgfqpoint{1.400603in}{0.896834in}}%
\pgfpathlineto{\pgfqpoint{1.375319in}{0.898185in}}%
\pgfpathlineto{\pgfqpoint{1.350036in}{0.900161in}}%
\pgfpathlineto{\pgfqpoint{1.346946in}{0.900479in}}%
\pgfpathlineto{\pgfqpoint{1.335293in}{0.913727in}}%
\pgfpathlineto{\pgfqpoint{1.324753in}{0.914909in}}%
\pgfpathlineto{\pgfqpoint{1.299470in}{0.918621in}}%
\pgfpathlineto{\pgfqpoint{1.274187in}{0.923523in}}%
\pgfpathlineto{\pgfqpoint{1.260109in}{0.926975in}}%
\pgfpathlineto{\pgfqpoint{1.248904in}{0.936605in}}%
\pgfpathlineto{\pgfqpoint{1.245119in}{0.940223in}}%
\pgfpathlineto{\pgfqpoint{1.228275in}{0.953471in}}%
\pgfpathlineto{\pgfqpoint{1.223620in}{0.955610in}}%
\pgfpathlineto{\pgfqpoint{1.202173in}{0.966719in}}%
\pgfpathlineto{\pgfqpoint{1.206542in}{0.979967in}}%
\pgfpathlineto{\pgfqpoint{1.198337in}{0.983504in}}%
\pgfpathlineto{\pgfqpoint{1.179125in}{0.993215in}}%
\pgfpathlineto{\pgfqpoint{1.183641in}{1.006463in}}%
\pgfpathlineto{\pgfqpoint{1.173054in}{1.016713in}}%
\pgfpathlineto{\pgfqpoint{1.170410in}{1.019711in}}%
\pgfpathlineto{\pgfqpoint{1.167430in}{1.032959in}}%
\pgfpathlineto{\pgfqpoint{1.168283in}{1.046207in}}%
\pgfpathlineto{\pgfqpoint{1.158215in}{1.059455in}}%
\pgfpathlineto{\pgfqpoint{1.167430in}{1.072703in}}%
\pgfpathlineto{\pgfqpoint{1.158074in}{1.085951in}}%
\pgfpathlineto{\pgfqpoint{1.165565in}{1.099199in}}%
\pgfpathlineto{\pgfqpoint{1.161147in}{1.112447in}}%
\pgfpathlineto{\pgfqpoint{1.161237in}{1.125695in}}%
\pgfpathlineto{\pgfqpoint{1.164929in}{1.138943in}}%
\pgfpathlineto{\pgfqpoint{1.160320in}{1.152191in}}%
\pgfpathlineto{\pgfqpoint{1.168500in}{1.165439in}}%
\pgfpathlineto{\pgfqpoint{1.163431in}{1.178687in}}%
\pgfpathlineto{\pgfqpoint{1.170761in}{1.191935in}}%
\pgfpathlineto{\pgfqpoint{1.170323in}{1.205183in}}%
\pgfpathlineto{\pgfqpoint{1.173054in}{1.216782in}}%
\pgfpathlineto{\pgfqpoint{1.173427in}{1.218431in}}%
\pgfpathlineto{\pgfqpoint{1.177702in}{1.231679in}}%
\pgfpathlineto{\pgfqpoint{1.175709in}{1.244927in}}%
\pgfpathlineto{\pgfqpoint{1.182984in}{1.258175in}}%
\pgfpathlineto{\pgfqpoint{1.180298in}{1.271423in}}%
\pgfpathlineto{\pgfqpoint{1.187017in}{1.284671in}}%
\pgfpathlineto{\pgfqpoint{1.186920in}{1.297919in}}%
\pgfpathlineto{\pgfqpoint{1.190035in}{1.311167in}}%
\pgfpathlineto{\pgfqpoint{1.194420in}{1.324415in}}%
\pgfpathlineto{\pgfqpoint{1.194504in}{1.337663in}}%
\pgfpathlineto{\pgfqpoint{1.198337in}{1.344859in}}%
\pgfpathlineto{\pgfqpoint{1.201289in}{1.350911in}}%
\pgfpathlineto{\pgfqpoint{1.199823in}{1.364159in}}%
\pgfpathlineto{\pgfqpoint{1.205986in}{1.377407in}}%
\pgfpathlineto{\pgfqpoint{1.206419in}{1.390655in}}%
\pgfpathlineto{\pgfqpoint{1.210210in}{1.403903in}}%
\pgfpathlineto{\pgfqpoint{1.213808in}{1.417151in}}%
\pgfpathlineto{\pgfqpoint{1.214422in}{1.430399in}}%
\pgfpathlineto{\pgfqpoint{1.220524in}{1.443647in}}%
\pgfpathlineto{\pgfqpoint{1.219891in}{1.456895in}}%
\pgfpathlineto{\pgfqpoint{1.223620in}{1.464783in}}%
\pgfpathlineto{\pgfqpoint{1.225934in}{1.470143in}}%
\pgfpathlineto{\pgfqpoint{1.226133in}{1.483391in}}%
\pgfpathlineto{\pgfqpoint{1.229801in}{1.496639in}}%
\pgfpathlineto{\pgfqpoint{1.232824in}{1.509887in}}%
\pgfpathlineto{\pgfqpoint{1.234331in}{1.523135in}}%
\pgfpathlineto{\pgfqpoint{1.239729in}{1.536383in}}%
\pgfpathlineto{\pgfqpoint{1.239743in}{1.549631in}}%
\pgfpathlineto{\pgfqpoint{1.245541in}{1.562879in}}%
\pgfpathlineto{\pgfqpoint{1.246008in}{1.576127in}}%
\pgfpathlineto{\pgfqpoint{1.248904in}{1.584973in}}%
\pgfpathlineto{\pgfqpoint{1.250199in}{1.589375in}}%
\pgfpathlineto{\pgfqpoint{1.252257in}{1.602623in}}%
\pgfpathlineto{\pgfqpoint{1.253811in}{1.615871in}}%
\pgfpathlineto{\pgfqpoint{1.258168in}{1.629119in}}%
\pgfpathlineto{\pgfqpoint{1.258563in}{1.642367in}}%
\pgfpathlineto{\pgfqpoint{1.263926in}{1.655615in}}%
\pgfpathlineto{\pgfqpoint{1.264503in}{1.668863in}}%
\pgfpathlineto{\pgfqpoint{1.268911in}{1.682111in}}%
\pgfpathlineto{\pgfqpoint{1.271063in}{1.695359in}}%
\pgfpathlineto{\pgfqpoint{1.273568in}{1.708607in}}%
\pgfpathlineto{\pgfqpoint{1.274187in}{1.710599in}}%
\pgfpathlineto{\pgfqpoint{1.277177in}{1.721855in}}%
\pgfpathlineto{\pgfqpoint{1.277698in}{1.735103in}}%
\pgfpathlineto{\pgfqpoint{1.282281in}{1.748351in}}%
\pgfpathlineto{\pgfqpoint{1.282729in}{1.761599in}}%
\pgfpathlineto{\pgfqpoint{1.287062in}{1.774847in}}%
\pgfpathlineto{\pgfqpoint{1.288624in}{1.788095in}}%
\pgfpathlineto{\pgfqpoint{1.291357in}{1.801343in}}%
\pgfpathlineto{\pgfqpoint{1.294760in}{1.814591in}}%
\pgfpathlineto{\pgfqpoint{1.295963in}{1.827839in}}%
\pgfpathlineto{\pgfqpoint{1.299470in}{1.837350in}}%
\pgfpathlineto{\pgfqpoint{1.300677in}{1.841087in}}%
\pgfpathlineto{\pgfqpoint{1.301104in}{1.854335in}}%
\pgfpathlineto{\pgfqpoint{1.305277in}{1.867583in}}%
\pgfpathlineto{\pgfqpoint{1.306577in}{1.880831in}}%
\pgfpathlineto{\pgfqpoint{1.309707in}{1.894079in}}%
\pgfpathlineto{\pgfqpoint{1.312350in}{1.907327in}}%
\pgfpathlineto{\pgfqpoint{1.313797in}{1.920575in}}%
\pgfpathlineto{\pgfqpoint{1.317802in}{1.933823in}}%
\pgfpathlineto{\pgfqpoint{1.318391in}{1.947071in}}%
\pgfpathlineto{\pgfqpoint{1.323137in}{1.960319in}}%
\pgfpathlineto{\pgfqpoint{1.324157in}{1.973567in}}%
\pgfpathlineto{\pgfqpoint{1.324753in}{1.975635in}}%
\pgfpathlineto{\pgfqpoint{1.327660in}{1.986815in}}%
\pgfpathlineto{\pgfqpoint{1.330053in}{2.000063in}}%
\pgfpathlineto{\pgfqpoint{1.332201in}{2.013311in}}%
\pgfpathlineto{\pgfqpoint{1.335886in}{2.026559in}}%
\pgfpathlineto{\pgfqpoint{1.336733in}{2.039807in}}%
\pgfpathlineto{\pgfqpoint{1.340997in}{2.053055in}}%
\pgfpathlineto{\pgfqpoint{1.341716in}{2.066303in}}%
\pgfpathlineto{\pgfqpoint{1.345842in}{2.079551in}}%
\pgfpathlineto{\pgfqpoint{1.347713in}{2.092799in}}%
\pgfpathlineto{\pgfqpoint{1.350036in}{2.103752in}}%
\pgfpathlineto{\pgfqpoint{1.350501in}{2.106047in}}%
\pgfpathlineto{\pgfqpoint{1.353973in}{2.119295in}}%
\pgfpathlineto{\pgfqpoint{1.355397in}{2.132543in}}%
\pgfpathlineto{\pgfqpoint{1.350036in}{2.132543in}}%
\pgfpathlineto{\pgfqpoint{1.347117in}{2.132543in}}%
\pgfpathlineto{\pgfqpoint{1.345676in}{2.119295in}}%
\pgfpathlineto{\pgfqpoint{1.341962in}{2.106047in}}%
\pgfpathlineto{\pgfqpoint{1.339437in}{2.092799in}}%
\pgfpathlineto{\pgfqpoint{1.337520in}{2.079551in}}%
\pgfpathlineto{\pgfqpoint{1.333751in}{2.066303in}}%
\pgfpathlineto{\pgfqpoint{1.332950in}{2.053055in}}%
\pgfpathlineto{\pgfqpoint{1.328931in}{2.039807in}}%
\pgfpathlineto{\pgfqpoint{1.328125in}{2.026559in}}%
\pgfpathlineto{\pgfqpoint{1.324753in}{2.014286in}}%
\pgfpathlineto{\pgfqpoint{1.324437in}{2.013311in}}%
\pgfpathlineto{\pgfqpoint{1.322135in}{2.000063in}}%
\pgfpathlineto{\pgfqpoint{1.319295in}{1.986815in}}%
\pgfpathlineto{\pgfqpoint{1.315715in}{1.973567in}}%
\pgfpathlineto{\pgfqpoint{1.314617in}{1.960319in}}%
\pgfpathlineto{\pgfqpoint{1.310200in}{1.947071in}}%
\pgfpathlineto{\pgfqpoint{1.309577in}{1.933823in}}%
\pgfpathlineto{\pgfqpoint{1.305700in}{1.920575in}}%
\pgfpathlineto{\pgfqpoint{1.304377in}{1.907327in}}%
\pgfpathlineto{\pgfqpoint{1.301665in}{1.894079in}}%
\pgfpathlineto{\pgfqpoint{1.299470in}{1.884044in}}%
\pgfpathlineto{\pgfqpoint{1.298644in}{1.880831in}}%
\pgfpathlineto{\pgfqpoint{1.296930in}{1.867583in}}%
\pgfpathlineto{\pgfqpoint{1.292334in}{1.854335in}}%
\pgfpathlineto{\pgfqpoint{1.291760in}{1.841087in}}%
\pgfpathlineto{\pgfqpoint{1.287092in}{1.827839in}}%
\pgfpathlineto{\pgfqpoint{1.285965in}{1.814591in}}%
\pgfpathlineto{\pgfqpoint{1.282556in}{1.801343in}}%
\pgfpathlineto{\pgfqpoint{1.280053in}{1.788095in}}%
\pgfpathlineto{\pgfqpoint{1.278324in}{1.774847in}}%
\pgfpathlineto{\pgfqpoint{1.274255in}{1.761599in}}%
\pgfpathlineto{\pgfqpoint{1.274187in}{1.760059in}}%
\pgfpathlineto{\pgfqpoint{1.273549in}{1.748351in}}%
\pgfpathlineto{\pgfqpoint{1.268315in}{1.735103in}}%
\pgfpathlineto{\pgfqpoint{1.267704in}{1.721855in}}%
\pgfpathlineto{\pgfqpoint{1.263642in}{1.708607in}}%
\pgfpathlineto{\pgfqpoint{1.261351in}{1.695359in}}%
\pgfpathlineto{\pgfqpoint{1.259031in}{1.682111in}}%
\pgfpathlineto{\pgfqpoint{1.254942in}{1.668863in}}%
\pgfpathlineto{\pgfqpoint{1.254152in}{1.655615in}}%
\pgfpathlineto{\pgfqpoint{1.248991in}{1.642367in}}%
\pgfpathlineto{\pgfqpoint{1.248904in}{1.639936in}}%
\pgfpathlineto{\pgfqpoint{1.248436in}{1.629119in}}%
\pgfpathlineto{\pgfqpoint{1.243317in}{1.615871in}}%
\pgfpathlineto{\pgfqpoint{1.241649in}{1.602623in}}%
\pgfpathlineto{\pgfqpoint{1.239058in}{1.589375in}}%
\pgfpathlineto{\pgfqpoint{1.234995in}{1.576127in}}%
\pgfpathlineto{\pgfqpoint{1.234204in}{1.562879in}}%
\pgfpathlineto{\pgfqpoint{1.228683in}{1.549631in}}%
\pgfpathlineto{\pgfqpoint{1.228487in}{1.536383in}}%
\pgfpathlineto{\pgfqpoint{1.223620in}{1.524428in}}%
\pgfpathlineto{\pgfqpoint{1.223012in}{1.523135in}}%
\pgfpathlineto{\pgfqpoint{1.221348in}{1.509887in}}%
\pgfpathlineto{\pgfqpoint{1.217575in}{1.496639in}}%
\pgfpathlineto{\pgfqpoint{1.213712in}{1.483391in}}%
\pgfpathlineto{\pgfqpoint{1.213019in}{1.470143in}}%
\pgfpathlineto{\pgfqpoint{1.206982in}{1.456895in}}%
\pgfpathlineto{\pgfqpoint{1.207235in}{1.443647in}}%
\pgfpathlineto{\pgfqpoint{1.201117in}{1.430399in}}%
\pgfpathlineto{\pgfqpoint{1.200454in}{1.417151in}}%
\pgfpathlineto{\pgfqpoint{1.198337in}{1.410038in}}%
\pgfpathlineto{\pgfqpoint{1.196234in}{1.403903in}}%
\pgfpathlineto{\pgfqpoint{1.192263in}{1.390655in}}%
\pgfpathlineto{\pgfqpoint{1.191123in}{1.377407in}}%
\pgfpathlineto{\pgfqpoint{1.184647in}{1.364159in}}%
\pgfpathlineto{\pgfqpoint{1.185647in}{1.350911in}}%
\pgfpathlineto{\pgfqpoint{1.178552in}{1.337663in}}%
\pgfpathlineto{\pgfqpoint{1.178221in}{1.324415in}}%
\pgfpathlineto{\pgfqpoint{1.173209in}{1.311167in}}%
\pgfpathlineto{\pgfqpoint{1.173054in}{1.310485in}}%
\pgfpathlineto{\pgfqpoint{1.169972in}{1.297919in}}%
\pgfpathlineto{\pgfqpoint{1.169045in}{1.284671in}}%
\pgfpathlineto{\pgfqpoint{1.162031in}{1.271423in}}%
\pgfpathlineto{\pgfqpoint{1.163966in}{1.258175in}}%
\pgfpathlineto{\pgfqpoint{1.156006in}{1.244927in}}%
\pgfpathlineto{\pgfqpoint{1.157586in}{1.231679in}}%
\pgfpathlineto{\pgfqpoint{1.152155in}{1.218431in}}%
\pgfpathlineto{\pgfqpoint{1.149028in}{1.205183in}}%
\pgfpathlineto{\pgfqpoint{1.147920in}{1.191935in}}%
\pgfpathlineto{\pgfqpoint{1.147771in}{1.191665in}}%
\pgfpathlineto{\pgfqpoint{1.140299in}{1.178687in}}%
\pgfpathlineto{\pgfqpoint{1.143999in}{1.165439in}}%
\pgfpathlineto{\pgfqpoint{1.134888in}{1.152191in}}%
\pgfpathlineto{\pgfqpoint{1.138659in}{1.138943in}}%
\pgfpathlineto{\pgfqpoint{1.133007in}{1.125695in}}%
\pgfpathlineto{\pgfqpoint{1.132571in}{1.112447in}}%
\pgfpathlineto{\pgfqpoint{1.134298in}{1.099199in}}%
\pgfpathlineto{\pgfqpoint{1.126422in}{1.085951in}}%
\pgfpathlineto{\pgfqpoint{1.133134in}{1.072703in}}%
\pgfpathlineto{\pgfqpoint{1.122488in}{1.059584in}}%
\pgfpathlineto{\pgfqpoint{1.122378in}{1.059455in}}%
\pgfpathlineto{\pgfqpoint{1.122488in}{1.059273in}}%
\pgfpathlineto{\pgfqpoint{1.130659in}{1.046207in}}%
\pgfpathlineto{\pgfqpoint{1.126136in}{1.032959in}}%
\pgfpathlineto{\pgfqpoint{1.128157in}{1.019711in}}%
\pgfpathlineto{\pgfqpoint{1.135315in}{1.006463in}}%
\pgfpathlineto{\pgfqpoint{1.129476in}{0.993215in}}%
\pgfpathlineto{\pgfqpoint{1.147771in}{0.980835in}}%
\pgfpathlineto{\pgfqpoint{1.149239in}{0.979967in}}%
\pgfpathlineto{\pgfqpoint{1.147771in}{0.977639in}}%
\pgfpathlineto{\pgfqpoint{1.140558in}{0.966719in}}%
\pgfpathlineto{\pgfqpoint{1.147771in}{0.961552in}}%
\pgfpathlineto{\pgfqpoint{1.160149in}{0.953471in}}%
\pgfpathlineto{\pgfqpoint{1.163338in}{0.940223in}}%
\pgfpathlineto{\pgfqpoint{1.173054in}{0.927860in}}%
\pgfpathlineto{\pgfqpoint{1.173799in}{0.926975in}}%
\pgfpathlineto{\pgfqpoint{1.198337in}{0.916727in}}%
\pgfpathlineto{\pgfqpoint{1.207459in}{0.913727in}}%
\pgfpathlineto{\pgfqpoint{1.206476in}{0.900479in}}%
\pgfpathlineto{\pgfqpoint{1.223620in}{0.896375in}}%
\pgfpathlineto{\pgfqpoint{1.248904in}{0.891275in}}%
\pgfpathlineto{\pgfqpoint{1.273030in}{0.887231in}}%
\pgfpathlineto{\pgfqpoint{1.274187in}{0.886235in}}%
\pgfpathlineto{\pgfqpoint{1.294090in}{0.873983in}}%
\pgfpathlineto{\pgfqpoint{1.299470in}{0.873204in}}%
\pgfpathclose%
\pgfusepath{fill}%
\end{pgfscope}%
\begin{pgfscope}%
\pgfpathrectangle{\pgfqpoint{0.540975in}{0.476543in}}{\pgfqpoint{1.820388in}{1.656000in}}%
\pgfusepath{clip}%
\pgfsetbuttcap%
\pgfsetroundjoin%
\definecolor{currentfill}{rgb}{0.699116,0.864975,0.921799}%
\pgfsetfillcolor{currentfill}%
\pgfsetlinewidth{0.000000pt}%
\definecolor{currentstroke}{rgb}{0.000000,0.000000,0.000000}%
\pgfsetstrokecolor{currentstroke}%
\pgfsetdash{}{0pt}%
\pgfpathmoveto{\pgfqpoint{1.350036in}{0.900161in}}%
\pgfpathlineto{\pgfqpoint{1.375319in}{0.898185in}}%
\pgfpathlineto{\pgfqpoint{1.400603in}{0.896834in}}%
\pgfpathlineto{\pgfqpoint{1.425886in}{0.896091in}}%
\pgfpathlineto{\pgfqpoint{1.451169in}{0.895966in}}%
\pgfpathlineto{\pgfqpoint{1.476452in}{0.896495in}}%
\pgfpathlineto{\pgfqpoint{1.501735in}{0.897740in}}%
\pgfpathlineto{\pgfqpoint{1.527018in}{0.899801in}}%
\pgfpathlineto{\pgfqpoint{1.532931in}{0.900479in}}%
\pgfpathlineto{\pgfqpoint{1.552302in}{0.911241in}}%
\pgfpathlineto{\pgfqpoint{1.555911in}{0.913727in}}%
\pgfpathlineto{\pgfqpoint{1.577585in}{0.916320in}}%
\pgfpathlineto{\pgfqpoint{1.602868in}{0.920232in}}%
\pgfpathlineto{\pgfqpoint{1.628151in}{0.925252in}}%
\pgfpathlineto{\pgfqpoint{1.635308in}{0.926975in}}%
\pgfpathlineto{\pgfqpoint{1.635555in}{0.940223in}}%
\pgfpathlineto{\pgfqpoint{1.653434in}{0.945269in}}%
\pgfpathlineto{\pgfqpoint{1.677723in}{0.953471in}}%
\pgfpathlineto{\pgfqpoint{1.678306in}{0.966719in}}%
\pgfpathlineto{\pgfqpoint{1.678717in}{0.966944in}}%
\pgfpathlineto{\pgfqpoint{1.700068in}{0.979967in}}%
\pgfpathlineto{\pgfqpoint{1.704001in}{0.985872in}}%
\pgfpathlineto{\pgfqpoint{1.708083in}{0.993215in}}%
\pgfpathlineto{\pgfqpoint{1.713525in}{1.006463in}}%
\pgfpathlineto{\pgfqpoint{1.726461in}{1.019711in}}%
\pgfpathlineto{\pgfqpoint{1.721005in}{1.032959in}}%
\pgfpathlineto{\pgfqpoint{1.729284in}{1.040025in}}%
\pgfpathlineto{\pgfqpoint{1.735737in}{1.046207in}}%
\pgfpathlineto{\pgfqpoint{1.729284in}{1.059354in}}%
\pgfpathlineto{\pgfqpoint{1.729236in}{1.059455in}}%
\pgfpathlineto{\pgfqpoint{1.729284in}{1.059515in}}%
\pgfpathlineto{\pgfqpoint{1.739202in}{1.072703in}}%
\pgfpathlineto{\pgfqpoint{1.735815in}{1.085951in}}%
\pgfpathlineto{\pgfqpoint{1.738157in}{1.099199in}}%
\pgfpathlineto{\pgfqpoint{1.739877in}{1.112447in}}%
\pgfpathlineto{\pgfqpoint{1.735394in}{1.125695in}}%
\pgfpathlineto{\pgfqpoint{1.740109in}{1.138943in}}%
\pgfpathlineto{\pgfqpoint{1.731762in}{1.152191in}}%
\pgfpathlineto{\pgfqpoint{1.736948in}{1.165439in}}%
\pgfpathlineto{\pgfqpoint{1.730771in}{1.178687in}}%
\pgfpathlineto{\pgfqpoint{1.732429in}{1.191935in}}%
\pgfpathlineto{\pgfqpoint{1.730485in}{1.205183in}}%
\pgfpathlineto{\pgfqpoint{1.729284in}{1.209582in}}%
\pgfpathlineto{\pgfqpoint{1.726973in}{1.218431in}}%
\pgfpathlineto{\pgfqpoint{1.728732in}{1.231679in}}%
\pgfpathlineto{\pgfqpoint{1.722207in}{1.244927in}}%
\pgfpathlineto{\pgfqpoint{1.724408in}{1.258175in}}%
\pgfpathlineto{\pgfqpoint{1.717377in}{1.271423in}}%
\pgfpathlineto{\pgfqpoint{1.718002in}{1.284671in}}%
\pgfpathlineto{\pgfqpoint{1.713917in}{1.297919in}}%
\pgfpathlineto{\pgfqpoint{1.711145in}{1.311167in}}%
\pgfpathlineto{\pgfqpoint{1.710700in}{1.324415in}}%
\pgfpathlineto{\pgfqpoint{1.704491in}{1.337663in}}%
\pgfpathlineto{\pgfqpoint{1.706037in}{1.350911in}}%
\pgfpathlineto{\pgfqpoint{1.704001in}{1.354990in}}%
\pgfpathlineto{\pgfqpoint{1.699803in}{1.364159in}}%
\pgfpathlineto{\pgfqpoint{1.700059in}{1.377407in}}%
\pgfpathlineto{\pgfqpoint{1.695632in}{1.390655in}}%
\pgfpathlineto{\pgfqpoint{1.693277in}{1.403903in}}%
\pgfpathlineto{\pgfqpoint{1.691705in}{1.417151in}}%
\pgfpathlineto{\pgfqpoint{1.686799in}{1.430399in}}%
\pgfpathlineto{\pgfqpoint{1.687080in}{1.443647in}}%
\pgfpathlineto{\pgfqpoint{1.680614in}{1.456895in}}%
\pgfpathlineto{\pgfqpoint{1.680670in}{1.470143in}}%
\pgfpathlineto{\pgfqpoint{1.678717in}{1.475235in}}%
\pgfpathlineto{\pgfqpoint{1.675952in}{1.483391in}}%
\pgfpathlineto{\pgfqpoint{1.673989in}{1.496639in}}%
\pgfpathlineto{\pgfqpoint{1.671892in}{1.509887in}}%
\pgfpathlineto{\pgfqpoint{1.667602in}{1.523135in}}%
\pgfpathlineto{\pgfqpoint{1.667590in}{1.536383in}}%
\pgfpathlineto{\pgfqpoint{1.662370in}{1.549631in}}%
\pgfpathlineto{\pgfqpoint{1.662325in}{1.562879in}}%
\pgfpathlineto{\pgfqpoint{1.657421in}{1.576127in}}%
\pgfpathlineto{\pgfqpoint{1.655409in}{1.589375in}}%
\pgfpathlineto{\pgfqpoint{1.653434in}{1.598235in}}%
\pgfpathlineto{\pgfqpoint{1.652594in}{1.602623in}}%
\pgfpathlineto{\pgfqpoint{1.648872in}{1.615871in}}%
\pgfpathlineto{\pgfqpoint{1.648111in}{1.629119in}}%
\pgfpathlineto{\pgfqpoint{1.643147in}{1.642367in}}%
\pgfpathlineto{\pgfqpoint{1.643094in}{1.655615in}}%
\pgfpathlineto{\pgfqpoint{1.638806in}{1.668863in}}%
\pgfpathlineto{\pgfqpoint{1.637471in}{1.682111in}}%
\pgfpathlineto{\pgfqpoint{1.634715in}{1.695359in}}%
\pgfpathlineto{\pgfqpoint{1.630953in}{1.708607in}}%
\pgfpathlineto{\pgfqpoint{1.629813in}{1.721855in}}%
\pgfpathlineto{\pgfqpoint{1.628151in}{1.726262in}}%
\pgfpathlineto{\pgfqpoint{1.625209in}{1.735103in}}%
\pgfpathlineto{\pgfqpoint{1.624682in}{1.748351in}}%
\pgfpathlineto{\pgfqpoint{1.620355in}{1.761599in}}%
\pgfpathlineto{\pgfqpoint{1.619071in}{1.774847in}}%
\pgfpathlineto{\pgfqpoint{1.616287in}{1.788095in}}%
\pgfpathlineto{\pgfqpoint{1.613328in}{1.801343in}}%
\pgfpathlineto{\pgfqpoint{1.612162in}{1.814591in}}%
\pgfpathlineto{\pgfqpoint{1.607549in}{1.827839in}}%
\pgfpathlineto{\pgfqpoint{1.607017in}{1.841087in}}%
\pgfpathlineto{\pgfqpoint{1.602868in}{1.853196in}}%
\pgfpathlineto{\pgfqpoint{1.602525in}{1.854335in}}%
\pgfpathlineto{\pgfqpoint{1.601321in}{1.867583in}}%
\pgfpathlineto{\pgfqpoint{1.598462in}{1.880831in}}%
\pgfpathlineto{\pgfqpoint{1.595649in}{1.894079in}}%
\pgfpathlineto{\pgfqpoint{1.594301in}{1.907327in}}%
\pgfpathlineto{\pgfqpoint{1.590236in}{1.920575in}}%
\pgfpathlineto{\pgfqpoint{1.589541in}{1.933823in}}%
\pgfpathlineto{\pgfqpoint{1.585145in}{1.947071in}}%
\pgfpathlineto{\pgfqpoint{1.583810in}{1.960319in}}%
\pgfpathlineto{\pgfqpoint{1.580637in}{1.973567in}}%
\pgfpathlineto{\pgfqpoint{1.577822in}{1.986815in}}%
\pgfpathlineto{\pgfqpoint{1.577585in}{1.988768in}}%
\pgfpathlineto{\pgfqpoint{1.576399in}{2.000063in}}%
\pgfpathlineto{\pgfqpoint{1.572478in}{2.013311in}}%
\pgfpathlineto{\pgfqpoint{1.571839in}{2.026559in}}%
\pgfpathlineto{\pgfqpoint{1.567667in}{2.039807in}}%
\pgfpathlineto{\pgfqpoint{1.566318in}{2.053055in}}%
\pgfpathlineto{\pgfqpoint{1.562997in}{2.066303in}}%
\pgfpathlineto{\pgfqpoint{1.559992in}{2.079551in}}%
\pgfpathlineto{\pgfqpoint{1.558233in}{2.092799in}}%
\pgfpathlineto{\pgfqpoint{1.554063in}{2.106047in}}%
\pgfpathlineto{\pgfqpoint{1.553274in}{2.119295in}}%
\pgfpathlineto{\pgfqpoint{1.552302in}{2.122117in}}%
\pgfpathlineto{\pgfqpoint{1.548903in}{2.132543in}}%
\pgfpathlineto{\pgfqpoint{1.540887in}{2.132543in}}%
\pgfpathlineto{\pgfqpoint{1.545315in}{2.119295in}}%
\pgfpathlineto{\pgfqpoint{1.546000in}{2.106047in}}%
\pgfpathlineto{\pgfqpoint{1.550207in}{2.092799in}}%
\pgfpathlineto{\pgfqpoint{1.551811in}{2.079551in}}%
\pgfpathlineto{\pgfqpoint{1.552302in}{2.077547in}}%
\pgfpathlineto{\pgfqpoint{1.555078in}{2.066303in}}%
\pgfpathlineto{\pgfqpoint{1.558433in}{2.053055in}}%
\pgfpathlineto{\pgfqpoint{1.559881in}{2.039807in}}%
\pgfpathlineto{\pgfqpoint{1.564223in}{2.026559in}}%
\pgfpathlineto{\pgfqpoint{1.564794in}{2.013311in}}%
\pgfpathlineto{\pgfqpoint{1.568956in}{2.000063in}}%
\pgfpathlineto{\pgfqpoint{1.570228in}{1.986815in}}%
\pgfpathlineto{\pgfqpoint{1.572917in}{1.973567in}}%
\pgfpathlineto{\pgfqpoint{1.575739in}{1.960319in}}%
\pgfpathlineto{\pgfqpoint{1.577024in}{1.947071in}}%
\pgfpathlineto{\pgfqpoint{1.577585in}{1.945303in}}%
\pgfpathlineto{\pgfqpoint{1.581545in}{1.933823in}}%
\pgfpathlineto{\pgfqpoint{1.582164in}{1.920575in}}%
\pgfpathlineto{\pgfqpoint{1.586489in}{1.907327in}}%
\pgfpathlineto{\pgfqpoint{1.587701in}{1.894079in}}%
\pgfpathlineto{\pgfqpoint{1.590735in}{1.880831in}}%
\pgfpathlineto{\pgfqpoint{1.593535in}{1.867583in}}%
\pgfpathlineto{\pgfqpoint{1.594817in}{1.854335in}}%
\pgfpathlineto{\pgfqpoint{1.598839in}{1.841087in}}%
\pgfpathlineto{\pgfqpoint{1.599221in}{1.827839in}}%
\pgfpathlineto{\pgfqpoint{1.602868in}{1.816628in}}%
\pgfpathlineto{\pgfqpoint{1.603602in}{1.814591in}}%
\pgfpathlineto{\pgfqpoint{1.604618in}{1.801343in}}%
\pgfpathlineto{\pgfqpoint{1.607793in}{1.788095in}}%
\pgfpathlineto{\pgfqpoint{1.610496in}{1.774847in}}%
\pgfpathlineto{\pgfqpoint{1.611847in}{1.761599in}}%
\pgfpathlineto{\pgfqpoint{1.616244in}{1.748351in}}%
\pgfpathlineto{\pgfqpoint{1.616672in}{1.735103in}}%
\pgfpathlineto{\pgfqpoint{1.621238in}{1.721855in}}%
\pgfpathlineto{\pgfqpoint{1.622029in}{1.708607in}}%
\pgfpathlineto{\pgfqpoint{1.625533in}{1.695359in}}%
\pgfpathlineto{\pgfqpoint{1.627859in}{1.682111in}}%
\pgfpathlineto{\pgfqpoint{1.628151in}{1.678996in}}%
\pgfpathlineto{\pgfqpoint{1.629212in}{1.668863in}}%
\pgfpathlineto{\pgfqpoint{1.633558in}{1.655615in}}%
\pgfpathlineto{\pgfqpoint{1.633453in}{1.642367in}}%
\pgfpathlineto{\pgfqpoint{1.638600in}{1.629119in}}%
\pgfpathlineto{\pgfqpoint{1.639101in}{1.615871in}}%
\pgfpathlineto{\pgfqpoint{1.643023in}{1.602623in}}%
\pgfpathlineto{\pgfqpoint{1.645353in}{1.589375in}}%
\pgfpathlineto{\pgfqpoint{1.647204in}{1.576127in}}%
\pgfpathlineto{\pgfqpoint{1.651540in}{1.562879in}}%
\pgfpathlineto{\pgfqpoint{1.651418in}{1.549631in}}%
\pgfpathlineto{\pgfqpoint{1.653434in}{1.544205in}}%
\pgfpathlineto{\pgfqpoint{1.656608in}{1.536383in}}%
\pgfpathlineto{\pgfqpoint{1.656255in}{1.523135in}}%
\pgfpathlineto{\pgfqpoint{1.660757in}{1.509887in}}%
\pgfpathlineto{\pgfqpoint{1.662507in}{1.496639in}}%
\pgfpathlineto{\pgfqpoint{1.664514in}{1.483391in}}%
\pgfpathlineto{\pgfqpoint{1.668828in}{1.470143in}}%
\pgfpathlineto{\pgfqpoint{1.668556in}{1.456895in}}%
\pgfpathlineto{\pgfqpoint{1.674463in}{1.443647in}}%
\pgfpathlineto{\pgfqpoint{1.673763in}{1.430399in}}%
\pgfpathlineto{\pgfqpoint{1.678473in}{1.417151in}}%
\pgfpathlineto{\pgfqpoint{1.678717in}{1.413762in}}%
\pgfpathlineto{\pgfqpoint{1.679515in}{1.403903in}}%
\pgfpathlineto{\pgfqpoint{1.681953in}{1.390655in}}%
\pgfpathlineto{\pgfqpoint{1.686020in}{1.377407in}}%
\pgfpathlineto{\pgfqpoint{1.685488in}{1.364159in}}%
\pgfpathlineto{\pgfqpoint{1.691440in}{1.350911in}}%
\pgfpathlineto{\pgfqpoint{1.689335in}{1.337663in}}%
\pgfpathlineto{\pgfqpoint{1.695152in}{1.324415in}}%
\pgfpathlineto{\pgfqpoint{1.694737in}{1.311167in}}%
\pgfpathlineto{\pgfqpoint{1.697351in}{1.297919in}}%
\pgfpathlineto{\pgfqpoint{1.700551in}{1.284671in}}%
\pgfpathlineto{\pgfqpoint{1.699639in}{1.271423in}}%
\pgfpathlineto{\pgfqpoint{1.704001in}{1.262320in}}%
\pgfpathlineto{\pgfqpoint{1.706117in}{1.258175in}}%
\pgfpathlineto{\pgfqpoint{1.704001in}{1.248880in}}%
\pgfpathlineto{\pgfqpoint{1.703087in}{1.244927in}}%
\pgfpathlineto{\pgfqpoint{1.704001in}{1.242979in}}%
\pgfpathlineto{\pgfqpoint{1.709493in}{1.231679in}}%
\pgfpathlineto{\pgfqpoint{1.706386in}{1.218431in}}%
\pgfpathlineto{\pgfqpoint{1.709754in}{1.205183in}}%
\pgfpathlineto{\pgfqpoint{1.710192in}{1.191935in}}%
\pgfpathlineto{\pgfqpoint{1.708043in}{1.178687in}}%
\pgfpathlineto{\pgfqpoint{1.712988in}{1.165439in}}%
\pgfpathlineto{\pgfqpoint{1.706605in}{1.152191in}}%
\pgfpathlineto{\pgfqpoint{1.714203in}{1.138943in}}%
\pgfpathlineto{\pgfqpoint{1.707378in}{1.125695in}}%
\pgfpathlineto{\pgfqpoint{1.711441in}{1.112447in}}%
\pgfpathlineto{\pgfqpoint{1.706963in}{1.099199in}}%
\pgfpathlineto{\pgfqpoint{1.704001in}{1.086135in}}%
\pgfpathlineto{\pgfqpoint{1.703957in}{1.085951in}}%
\pgfpathlineto{\pgfqpoint{1.704001in}{1.085131in}}%
\pgfpathlineto{\pgfqpoint{1.704700in}{1.072703in}}%
\pgfpathlineto{\pgfqpoint{1.704001in}{1.071892in}}%
\pgfpathlineto{\pgfqpoint{1.692363in}{1.059455in}}%
\pgfpathlineto{\pgfqpoint{1.697147in}{1.046207in}}%
\pgfpathlineto{\pgfqpoint{1.678717in}{1.033574in}}%
\pgfpathlineto{\pgfqpoint{1.677679in}{1.032959in}}%
\pgfpathlineto{\pgfqpoint{1.678717in}{1.029759in}}%
\pgfpathlineto{\pgfqpoint{1.681925in}{1.019711in}}%
\pgfpathlineto{\pgfqpoint{1.678717in}{1.017308in}}%
\pgfpathlineto{\pgfqpoint{1.661214in}{1.006463in}}%
\pgfpathlineto{\pgfqpoint{1.653519in}{0.993215in}}%
\pgfpathlineto{\pgfqpoint{1.653434in}{0.993141in}}%
\pgfpathlineto{\pgfqpoint{1.634335in}{0.979967in}}%
\pgfpathlineto{\pgfqpoint{1.628151in}{0.976961in}}%
\pgfpathlineto{\pgfqpoint{1.604157in}{0.966719in}}%
\pgfpathlineto{\pgfqpoint{1.602868in}{0.965118in}}%
\pgfpathlineto{\pgfqpoint{1.590378in}{0.953471in}}%
\pgfpathlineto{\pgfqpoint{1.577585in}{0.950746in}}%
\pgfpathlineto{\pgfqpoint{1.552302in}{0.946365in}}%
\pgfpathlineto{\pgfqpoint{1.527018in}{0.942925in}}%
\pgfpathlineto{\pgfqpoint{1.501735in}{0.940344in}}%
\pgfpathlineto{\pgfqpoint{1.499974in}{0.940223in}}%
\pgfpathlineto{\pgfqpoint{1.476452in}{0.931068in}}%
\pgfpathlineto{\pgfqpoint{1.451169in}{0.927790in}}%
\pgfpathlineto{\pgfqpoint{1.425886in}{0.929147in}}%
\pgfpathlineto{\pgfqpoint{1.400603in}{0.934056in}}%
\pgfpathlineto{\pgfqpoint{1.381063in}{0.940223in}}%
\pgfpathlineto{\pgfqpoint{1.375319in}{0.940747in}}%
\pgfpathlineto{\pgfqpoint{1.350036in}{0.944500in}}%
\pgfpathlineto{\pgfqpoint{1.324753in}{0.950397in}}%
\pgfpathlineto{\pgfqpoint{1.315353in}{0.953471in}}%
\pgfpathlineto{\pgfqpoint{1.299470in}{0.958292in}}%
\pgfpathlineto{\pgfqpoint{1.275781in}{0.966719in}}%
\pgfpathlineto{\pgfqpoint{1.274187in}{0.971718in}}%
\pgfpathlineto{\pgfqpoint{1.272141in}{0.979967in}}%
\pgfpathlineto{\pgfqpoint{1.248904in}{0.987304in}}%
\pgfpathlineto{\pgfqpoint{1.233360in}{0.993215in}}%
\pgfpathlineto{\pgfqpoint{1.235904in}{1.006463in}}%
\pgfpathlineto{\pgfqpoint{1.223620in}{1.013498in}}%
\pgfpathlineto{\pgfqpoint{1.214731in}{1.019711in}}%
\pgfpathlineto{\pgfqpoint{1.210579in}{1.032959in}}%
\pgfpathlineto{\pgfqpoint{1.206826in}{1.046207in}}%
\pgfpathlineto{\pgfqpoint{1.198337in}{1.055209in}}%
\pgfpathlineto{\pgfqpoint{1.194631in}{1.059455in}}%
\pgfpathlineto{\pgfqpoint{1.198337in}{1.065996in}}%
\pgfpathlineto{\pgfqpoint{1.202009in}{1.072703in}}%
\pgfpathlineto{\pgfqpoint{1.198337in}{1.076371in}}%
\pgfpathlineto{\pgfqpoint{1.189737in}{1.085951in}}%
\pgfpathlineto{\pgfqpoint{1.196728in}{1.099199in}}%
\pgfpathlineto{\pgfqpoint{1.189389in}{1.112447in}}%
\pgfpathlineto{\pgfqpoint{1.189118in}{1.125695in}}%
\pgfpathlineto{\pgfqpoint{1.190686in}{1.138943in}}%
\pgfpathlineto{\pgfqpoint{1.185228in}{1.152191in}}%
\pgfpathlineto{\pgfqpoint{1.192410in}{1.165439in}}%
\pgfpathlineto{\pgfqpoint{1.185939in}{1.178687in}}%
\pgfpathlineto{\pgfqpoint{1.192973in}{1.191935in}}%
\pgfpathlineto{\pgfqpoint{1.190929in}{1.205183in}}%
\pgfpathlineto{\pgfqpoint{1.193859in}{1.218431in}}%
\pgfpathlineto{\pgfqpoint{1.197016in}{1.231679in}}%
\pgfpathlineto{\pgfqpoint{1.194539in}{1.244927in}}%
\pgfpathlineto{\pgfqpoint{1.198337in}{1.252169in}}%
\pgfpathlineto{\pgfqpoint{1.201308in}{1.258175in}}%
\pgfpathlineto{\pgfqpoint{1.198337in}{1.269542in}}%
\pgfpathlineto{\pgfqpoint{1.197851in}{1.271423in}}%
\pgfpathlineto{\pgfqpoint{1.198337in}{1.272371in}}%
\pgfpathlineto{\pgfqpoint{1.204403in}{1.284671in}}%
\pgfpathlineto{\pgfqpoint{1.203346in}{1.297919in}}%
\pgfpathlineto{\pgfqpoint{1.206394in}{1.311167in}}%
\pgfpathlineto{\pgfqpoint{1.209801in}{1.324415in}}%
\pgfpathlineto{\pgfqpoint{1.209638in}{1.337663in}}%
\pgfpathlineto{\pgfqpoint{1.215897in}{1.350911in}}%
\pgfpathlineto{\pgfqpoint{1.213865in}{1.364159in}}%
\pgfpathlineto{\pgfqpoint{1.220171in}{1.377407in}}%
\pgfpathlineto{\pgfqpoint{1.219965in}{1.390655in}}%
\pgfpathlineto{\pgfqpoint{1.223620in}{1.402780in}}%
\pgfpathlineto{\pgfqpoint{1.223941in}{1.403903in}}%
\pgfpathlineto{\pgfqpoint{1.226884in}{1.417151in}}%
\pgfpathlineto{\pgfqpoint{1.227404in}{1.430399in}}%
\pgfpathlineto{\pgfqpoint{1.232862in}{1.443647in}}%
\pgfpathlineto{\pgfqpoint{1.231934in}{1.456895in}}%
\pgfpathlineto{\pgfqpoint{1.237718in}{1.470143in}}%
\pgfpathlineto{\pgfqpoint{1.237435in}{1.483391in}}%
\pgfpathlineto{\pgfqpoint{1.241303in}{1.496639in}}%
\pgfpathlineto{\pgfqpoint{1.243999in}{1.509887in}}%
\pgfpathlineto{\pgfqpoint{1.245568in}{1.523135in}}%
\pgfpathlineto{\pgfqpoint{1.248904in}{1.531317in}}%
\pgfpathlineto{\pgfqpoint{1.250766in}{1.536383in}}%
\pgfpathlineto{\pgfqpoint{1.250615in}{1.549631in}}%
\pgfpathlineto{\pgfqpoint{1.255977in}{1.562879in}}%
\pgfpathlineto{\pgfqpoint{1.256065in}{1.576127in}}%
\pgfpathlineto{\pgfqpoint{1.260102in}{1.589375in}}%
\pgfpathlineto{\pgfqpoint{1.261867in}{1.602623in}}%
\pgfpathlineto{\pgfqpoint{1.263526in}{1.615871in}}%
\pgfpathlineto{\pgfqpoint{1.267826in}{1.629119in}}%
\pgfpathlineto{\pgfqpoint{1.268135in}{1.642367in}}%
\pgfpathlineto{\pgfqpoint{1.273701in}{1.655615in}}%
\pgfpathlineto{\pgfqpoint{1.274065in}{1.668863in}}%
\pgfpathlineto{\pgfqpoint{1.274187in}{1.669217in}}%
\pgfpathlineto{\pgfqpoint{1.278233in}{1.682111in}}%
\pgfpathlineto{\pgfqpoint{1.279905in}{1.695359in}}%
\pgfpathlineto{\pgfqpoint{1.282233in}{1.708607in}}%
\pgfpathlineto{\pgfqpoint{1.285667in}{1.721855in}}%
\pgfpathlineto{\pgfqpoint{1.286187in}{1.735103in}}%
\pgfpathlineto{\pgfqpoint{1.290905in}{1.748351in}}%
\pgfpathlineto{\pgfqpoint{1.291202in}{1.761599in}}%
\pgfpathlineto{\pgfqpoint{1.295799in}{1.774847in}}%
\pgfpathlineto{\pgfqpoint{1.297194in}{1.788095in}}%
\pgfpathlineto{\pgfqpoint{1.299470in}{1.798331in}}%
\pgfpathlineto{\pgfqpoint{1.300079in}{1.801343in}}%
\pgfpathlineto{\pgfqpoint{1.303022in}{1.814591in}}%
\pgfpathlineto{\pgfqpoint{1.304132in}{1.827839in}}%
\pgfpathlineto{\pgfqpoint{1.308444in}{1.841087in}}%
\pgfpathlineto{\pgfqpoint{1.308796in}{1.854335in}}%
\pgfpathlineto{\pgfqpoint{1.313203in}{1.867583in}}%
\pgfpathlineto{\pgfqpoint{1.314364in}{1.880831in}}%
\pgfpathlineto{\pgfqpoint{1.317748in}{1.894079in}}%
\pgfpathlineto{\pgfqpoint{1.320323in}{1.907327in}}%
\pgfpathlineto{\pgfqpoint{1.321894in}{1.920575in}}%
\pgfpathlineto{\pgfqpoint{1.324753in}{1.929783in}}%
\pgfpathlineto{\pgfqpoint{1.325892in}{1.933823in}}%
\pgfpathlineto{\pgfqpoint{1.326389in}{1.947071in}}%
\pgfpathlineto{\pgfqpoint{1.330813in}{1.960319in}}%
\pgfpathlineto{\pgfqpoint{1.331585in}{1.973567in}}%
\pgfpathlineto{\pgfqpoint{1.335269in}{1.986815in}}%
\pgfpathlineto{\pgfqpoint{1.337564in}{2.000063in}}%
\pgfpathlineto{\pgfqpoint{1.339916in}{2.013311in}}%
\pgfpathlineto{\pgfqpoint{1.343646in}{2.026559in}}%
\pgfpathlineto{\pgfqpoint{1.344536in}{2.039807in}}%
\pgfpathlineto{\pgfqpoint{1.349045in}{2.053055in}}%
\pgfpathlineto{\pgfqpoint{1.349681in}{2.066303in}}%
\pgfpathlineto{\pgfqpoint{1.350036in}{2.067396in}}%
\pgfpathlineto{\pgfqpoint{1.353903in}{2.079551in}}%
\pgfpathlineto{\pgfqpoint{1.355538in}{2.092799in}}%
\pgfpathlineto{\pgfqpoint{1.358405in}{2.106047in}}%
\pgfpathlineto{\pgfqpoint{1.361831in}{2.119295in}}%
\pgfpathlineto{\pgfqpoint{1.363382in}{2.132543in}}%
\pgfpathlineto{\pgfqpoint{1.355397in}{2.132543in}}%
\pgfpathlineto{\pgfqpoint{1.353973in}{2.119295in}}%
\pgfpathlineto{\pgfqpoint{1.350501in}{2.106047in}}%
\pgfpathlineto{\pgfqpoint{1.350036in}{2.103752in}}%
\pgfpathlineto{\pgfqpoint{1.347713in}{2.092799in}}%
\pgfpathlineto{\pgfqpoint{1.345842in}{2.079551in}}%
\pgfpathlineto{\pgfqpoint{1.341716in}{2.066303in}}%
\pgfpathlineto{\pgfqpoint{1.340997in}{2.053055in}}%
\pgfpathlineto{\pgfqpoint{1.336733in}{2.039807in}}%
\pgfpathlineto{\pgfqpoint{1.335886in}{2.026559in}}%
\pgfpathlineto{\pgfqpoint{1.332201in}{2.013311in}}%
\pgfpathlineto{\pgfqpoint{1.330053in}{2.000063in}}%
\pgfpathlineto{\pgfqpoint{1.327660in}{1.986815in}}%
\pgfpathlineto{\pgfqpoint{1.324753in}{1.975635in}}%
\pgfpathlineto{\pgfqpoint{1.324157in}{1.973567in}}%
\pgfpathlineto{\pgfqpoint{1.323137in}{1.960319in}}%
\pgfpathlineto{\pgfqpoint{1.318391in}{1.947071in}}%
\pgfpathlineto{\pgfqpoint{1.317802in}{1.933823in}}%
\pgfpathlineto{\pgfqpoint{1.313797in}{1.920575in}}%
\pgfpathlineto{\pgfqpoint{1.312350in}{1.907327in}}%
\pgfpathlineto{\pgfqpoint{1.309707in}{1.894079in}}%
\pgfpathlineto{\pgfqpoint{1.306577in}{1.880831in}}%
\pgfpathlineto{\pgfqpoint{1.305277in}{1.867583in}}%
\pgfpathlineto{\pgfqpoint{1.301104in}{1.854335in}}%
\pgfpathlineto{\pgfqpoint{1.300677in}{1.841087in}}%
\pgfpathlineto{\pgfqpoint{1.299470in}{1.837350in}}%
\pgfpathlineto{\pgfqpoint{1.295963in}{1.827839in}}%
\pgfpathlineto{\pgfqpoint{1.294760in}{1.814591in}}%
\pgfpathlineto{\pgfqpoint{1.291357in}{1.801343in}}%
\pgfpathlineto{\pgfqpoint{1.288624in}{1.788095in}}%
\pgfpathlineto{\pgfqpoint{1.287062in}{1.774847in}}%
\pgfpathlineto{\pgfqpoint{1.282729in}{1.761599in}}%
\pgfpathlineto{\pgfqpoint{1.282281in}{1.748351in}}%
\pgfpathlineto{\pgfqpoint{1.277698in}{1.735103in}}%
\pgfpathlineto{\pgfqpoint{1.277177in}{1.721855in}}%
\pgfpathlineto{\pgfqpoint{1.274187in}{1.710599in}}%
\pgfpathlineto{\pgfqpoint{1.273568in}{1.708607in}}%
\pgfpathlineto{\pgfqpoint{1.271063in}{1.695359in}}%
\pgfpathlineto{\pgfqpoint{1.268911in}{1.682111in}}%
\pgfpathlineto{\pgfqpoint{1.264503in}{1.668863in}}%
\pgfpathlineto{\pgfqpoint{1.263926in}{1.655615in}}%
\pgfpathlineto{\pgfqpoint{1.258563in}{1.642367in}}%
\pgfpathlineto{\pgfqpoint{1.258168in}{1.629119in}}%
\pgfpathlineto{\pgfqpoint{1.253811in}{1.615871in}}%
\pgfpathlineto{\pgfqpoint{1.252257in}{1.602623in}}%
\pgfpathlineto{\pgfqpoint{1.250199in}{1.589375in}}%
\pgfpathlineto{\pgfqpoint{1.248904in}{1.584973in}}%
\pgfpathlineto{\pgfqpoint{1.246008in}{1.576127in}}%
\pgfpathlineto{\pgfqpoint{1.245541in}{1.562879in}}%
\pgfpathlineto{\pgfqpoint{1.239743in}{1.549631in}}%
\pgfpathlineto{\pgfqpoint{1.239729in}{1.536383in}}%
\pgfpathlineto{\pgfqpoint{1.234331in}{1.523135in}}%
\pgfpathlineto{\pgfqpoint{1.232824in}{1.509887in}}%
\pgfpathlineto{\pgfqpoint{1.229801in}{1.496639in}}%
\pgfpathlineto{\pgfqpoint{1.226133in}{1.483391in}}%
\pgfpathlineto{\pgfqpoint{1.225934in}{1.470143in}}%
\pgfpathlineto{\pgfqpoint{1.223620in}{1.464783in}}%
\pgfpathlineto{\pgfqpoint{1.219891in}{1.456895in}}%
\pgfpathlineto{\pgfqpoint{1.220524in}{1.443647in}}%
\pgfpathlineto{\pgfqpoint{1.214422in}{1.430399in}}%
\pgfpathlineto{\pgfqpoint{1.213808in}{1.417151in}}%
\pgfpathlineto{\pgfqpoint{1.210210in}{1.403903in}}%
\pgfpathlineto{\pgfqpoint{1.206419in}{1.390655in}}%
\pgfpathlineto{\pgfqpoint{1.205986in}{1.377407in}}%
\pgfpathlineto{\pgfqpoint{1.199823in}{1.364159in}}%
\pgfpathlineto{\pgfqpoint{1.201289in}{1.350911in}}%
\pgfpathlineto{\pgfqpoint{1.198337in}{1.344859in}}%
\pgfpathlineto{\pgfqpoint{1.194504in}{1.337663in}}%
\pgfpathlineto{\pgfqpoint{1.194420in}{1.324415in}}%
\pgfpathlineto{\pgfqpoint{1.190035in}{1.311167in}}%
\pgfpathlineto{\pgfqpoint{1.186920in}{1.297919in}}%
\pgfpathlineto{\pgfqpoint{1.187017in}{1.284671in}}%
\pgfpathlineto{\pgfqpoint{1.180298in}{1.271423in}}%
\pgfpathlineto{\pgfqpoint{1.182984in}{1.258175in}}%
\pgfpathlineto{\pgfqpoint{1.175709in}{1.244927in}}%
\pgfpathlineto{\pgfqpoint{1.177702in}{1.231679in}}%
\pgfpathlineto{\pgfqpoint{1.173427in}{1.218431in}}%
\pgfpathlineto{\pgfqpoint{1.173054in}{1.216782in}}%
\pgfpathlineto{\pgfqpoint{1.170323in}{1.205183in}}%
\pgfpathlineto{\pgfqpoint{1.170761in}{1.191935in}}%
\pgfpathlineto{\pgfqpoint{1.163431in}{1.178687in}}%
\pgfpathlineto{\pgfqpoint{1.168500in}{1.165439in}}%
\pgfpathlineto{\pgfqpoint{1.160320in}{1.152191in}}%
\pgfpathlineto{\pgfqpoint{1.164929in}{1.138943in}}%
\pgfpathlineto{\pgfqpoint{1.161237in}{1.125695in}}%
\pgfpathlineto{\pgfqpoint{1.161147in}{1.112447in}}%
\pgfpathlineto{\pgfqpoint{1.165565in}{1.099199in}}%
\pgfpathlineto{\pgfqpoint{1.158074in}{1.085951in}}%
\pgfpathlineto{\pgfqpoint{1.167430in}{1.072703in}}%
\pgfpathlineto{\pgfqpoint{1.158215in}{1.059455in}}%
\pgfpathlineto{\pgfqpoint{1.168283in}{1.046207in}}%
\pgfpathlineto{\pgfqpoint{1.167430in}{1.032959in}}%
\pgfpathlineto{\pgfqpoint{1.170410in}{1.019711in}}%
\pgfpathlineto{\pgfqpoint{1.173054in}{1.016713in}}%
\pgfpathlineto{\pgfqpoint{1.183641in}{1.006463in}}%
\pgfpathlineto{\pgfqpoint{1.179125in}{0.993215in}}%
\pgfpathlineto{\pgfqpoint{1.198337in}{0.983504in}}%
\pgfpathlineto{\pgfqpoint{1.206542in}{0.979967in}}%
\pgfpathlineto{\pgfqpoint{1.202173in}{0.966719in}}%
\pgfpathlineto{\pgfqpoint{1.223620in}{0.955610in}}%
\pgfpathlineto{\pgfqpoint{1.228275in}{0.953471in}}%
\pgfpathlineto{\pgfqpoint{1.245119in}{0.940223in}}%
\pgfpathlineto{\pgfqpoint{1.248904in}{0.936605in}}%
\pgfpathlineto{\pgfqpoint{1.260109in}{0.926975in}}%
\pgfpathlineto{\pgfqpoint{1.274187in}{0.923523in}}%
\pgfpathlineto{\pgfqpoint{1.299470in}{0.918621in}}%
\pgfpathlineto{\pgfqpoint{1.324753in}{0.914909in}}%
\pgfpathlineto{\pgfqpoint{1.335293in}{0.913727in}}%
\pgfpathlineto{\pgfqpoint{1.346946in}{0.900479in}}%
\pgfpathlineto{\pgfqpoint{1.350036in}{0.900161in}}%
\pgfpathclose%
\pgfusepath{fill}%
\end{pgfscope}%
\begin{pgfscope}%
\pgfpathrectangle{\pgfqpoint{0.540975in}{0.476543in}}{\pgfqpoint{1.820388in}{1.656000in}}%
\pgfusepath{clip}%
\pgfsetbuttcap%
\pgfsetroundjoin%
\definecolor{currentfill}{rgb}{0.780623,0.904960,0.944867}%
\pgfsetfillcolor{currentfill}%
\pgfsetlinewidth{0.000000pt}%
\definecolor{currentstroke}{rgb}{0.000000,0.000000,0.000000}%
\pgfsetstrokecolor{currentstroke}%
\pgfsetdash{}{0pt}%
\pgfpathmoveto{\pgfqpoint{1.400603in}{0.934056in}}%
\pgfpathlineto{\pgfqpoint{1.425886in}{0.929147in}}%
\pgfpathlineto{\pgfqpoint{1.451169in}{0.927790in}}%
\pgfpathlineto{\pgfqpoint{1.476452in}{0.931068in}}%
\pgfpathlineto{\pgfqpoint{1.499974in}{0.940223in}}%
\pgfpathlineto{\pgfqpoint{1.501735in}{0.940344in}}%
\pgfpathlineto{\pgfqpoint{1.527018in}{0.942925in}}%
\pgfpathlineto{\pgfqpoint{1.552302in}{0.946365in}}%
\pgfpathlineto{\pgfqpoint{1.577585in}{0.950746in}}%
\pgfpathlineto{\pgfqpoint{1.590378in}{0.953471in}}%
\pgfpathlineto{\pgfqpoint{1.602868in}{0.965118in}}%
\pgfpathlineto{\pgfqpoint{1.604157in}{0.966719in}}%
\pgfpathlineto{\pgfqpoint{1.628151in}{0.976961in}}%
\pgfpathlineto{\pgfqpoint{1.634335in}{0.979967in}}%
\pgfpathlineto{\pgfqpoint{1.653434in}{0.993141in}}%
\pgfpathlineto{\pgfqpoint{1.653519in}{0.993215in}}%
\pgfpathlineto{\pgfqpoint{1.661214in}{1.006463in}}%
\pgfpathlineto{\pgfqpoint{1.678717in}{1.017308in}}%
\pgfpathlineto{\pgfqpoint{1.681925in}{1.019711in}}%
\pgfpathlineto{\pgfqpoint{1.678717in}{1.029759in}}%
\pgfpathlineto{\pgfqpoint{1.677679in}{1.032959in}}%
\pgfpathlineto{\pgfqpoint{1.678717in}{1.033574in}}%
\pgfpathlineto{\pgfqpoint{1.697147in}{1.046207in}}%
\pgfpathlineto{\pgfqpoint{1.692363in}{1.059455in}}%
\pgfpathlineto{\pgfqpoint{1.704001in}{1.071892in}}%
\pgfpathlineto{\pgfqpoint{1.704700in}{1.072703in}}%
\pgfpathlineto{\pgfqpoint{1.704001in}{1.085131in}}%
\pgfpathlineto{\pgfqpoint{1.703957in}{1.085951in}}%
\pgfpathlineto{\pgfqpoint{1.704001in}{1.086135in}}%
\pgfpathlineto{\pgfqpoint{1.706963in}{1.099199in}}%
\pgfpathlineto{\pgfqpoint{1.711441in}{1.112447in}}%
\pgfpathlineto{\pgfqpoint{1.707378in}{1.125695in}}%
\pgfpathlineto{\pgfqpoint{1.714203in}{1.138943in}}%
\pgfpathlineto{\pgfqpoint{1.706605in}{1.152191in}}%
\pgfpathlineto{\pgfqpoint{1.712988in}{1.165439in}}%
\pgfpathlineto{\pgfqpoint{1.708043in}{1.178687in}}%
\pgfpathlineto{\pgfqpoint{1.710192in}{1.191935in}}%
\pgfpathlineto{\pgfqpoint{1.709754in}{1.205183in}}%
\pgfpathlineto{\pgfqpoint{1.706386in}{1.218431in}}%
\pgfpathlineto{\pgfqpoint{1.709493in}{1.231679in}}%
\pgfpathlineto{\pgfqpoint{1.704001in}{1.242979in}}%
\pgfpathlineto{\pgfqpoint{1.703087in}{1.244927in}}%
\pgfpathlineto{\pgfqpoint{1.704001in}{1.248880in}}%
\pgfpathlineto{\pgfqpoint{1.706117in}{1.258175in}}%
\pgfpathlineto{\pgfqpoint{1.704001in}{1.262320in}}%
\pgfpathlineto{\pgfqpoint{1.699639in}{1.271423in}}%
\pgfpathlineto{\pgfqpoint{1.700551in}{1.284671in}}%
\pgfpathlineto{\pgfqpoint{1.697351in}{1.297919in}}%
\pgfpathlineto{\pgfqpoint{1.694737in}{1.311167in}}%
\pgfpathlineto{\pgfqpoint{1.695152in}{1.324415in}}%
\pgfpathlineto{\pgfqpoint{1.689335in}{1.337663in}}%
\pgfpathlineto{\pgfqpoint{1.691440in}{1.350911in}}%
\pgfpathlineto{\pgfqpoint{1.685488in}{1.364159in}}%
\pgfpathlineto{\pgfqpoint{1.686020in}{1.377407in}}%
\pgfpathlineto{\pgfqpoint{1.681953in}{1.390655in}}%
\pgfpathlineto{\pgfqpoint{1.679515in}{1.403903in}}%
\pgfpathlineto{\pgfqpoint{1.678717in}{1.413762in}}%
\pgfpathlineto{\pgfqpoint{1.678473in}{1.417151in}}%
\pgfpathlineto{\pgfqpoint{1.673763in}{1.430399in}}%
\pgfpathlineto{\pgfqpoint{1.674463in}{1.443647in}}%
\pgfpathlineto{\pgfqpoint{1.668556in}{1.456895in}}%
\pgfpathlineto{\pgfqpoint{1.668828in}{1.470143in}}%
\pgfpathlineto{\pgfqpoint{1.664514in}{1.483391in}}%
\pgfpathlineto{\pgfqpoint{1.662507in}{1.496639in}}%
\pgfpathlineto{\pgfqpoint{1.660757in}{1.509887in}}%
\pgfpathlineto{\pgfqpoint{1.656255in}{1.523135in}}%
\pgfpathlineto{\pgfqpoint{1.656608in}{1.536383in}}%
\pgfpathlineto{\pgfqpoint{1.653434in}{1.544205in}}%
\pgfpathlineto{\pgfqpoint{1.651418in}{1.549631in}}%
\pgfpathlineto{\pgfqpoint{1.651540in}{1.562879in}}%
\pgfpathlineto{\pgfqpoint{1.647204in}{1.576127in}}%
\pgfpathlineto{\pgfqpoint{1.645353in}{1.589375in}}%
\pgfpathlineto{\pgfqpoint{1.643023in}{1.602623in}}%
\pgfpathlineto{\pgfqpoint{1.639101in}{1.615871in}}%
\pgfpathlineto{\pgfqpoint{1.638600in}{1.629119in}}%
\pgfpathlineto{\pgfqpoint{1.633453in}{1.642367in}}%
\pgfpathlineto{\pgfqpoint{1.633558in}{1.655615in}}%
\pgfpathlineto{\pgfqpoint{1.629212in}{1.668863in}}%
\pgfpathlineto{\pgfqpoint{1.628151in}{1.678996in}}%
\pgfpathlineto{\pgfqpoint{1.627859in}{1.682111in}}%
\pgfpathlineto{\pgfqpoint{1.625533in}{1.695359in}}%
\pgfpathlineto{\pgfqpoint{1.622029in}{1.708607in}}%
\pgfpathlineto{\pgfqpoint{1.621238in}{1.721855in}}%
\pgfpathlineto{\pgfqpoint{1.616672in}{1.735103in}}%
\pgfpathlineto{\pgfqpoint{1.616244in}{1.748351in}}%
\pgfpathlineto{\pgfqpoint{1.611847in}{1.761599in}}%
\pgfpathlineto{\pgfqpoint{1.610496in}{1.774847in}}%
\pgfpathlineto{\pgfqpoint{1.607793in}{1.788095in}}%
\pgfpathlineto{\pgfqpoint{1.604618in}{1.801343in}}%
\pgfpathlineto{\pgfqpoint{1.603602in}{1.814591in}}%
\pgfpathlineto{\pgfqpoint{1.602868in}{1.816628in}}%
\pgfpathlineto{\pgfqpoint{1.599221in}{1.827839in}}%
\pgfpathlineto{\pgfqpoint{1.598839in}{1.841087in}}%
\pgfpathlineto{\pgfqpoint{1.594817in}{1.854335in}}%
\pgfpathlineto{\pgfqpoint{1.593535in}{1.867583in}}%
\pgfpathlineto{\pgfqpoint{1.590735in}{1.880831in}}%
\pgfpathlineto{\pgfqpoint{1.587701in}{1.894079in}}%
\pgfpathlineto{\pgfqpoint{1.586489in}{1.907327in}}%
\pgfpathlineto{\pgfqpoint{1.582164in}{1.920575in}}%
\pgfpathlineto{\pgfqpoint{1.581545in}{1.933823in}}%
\pgfpathlineto{\pgfqpoint{1.577585in}{1.945303in}}%
\pgfpathlineto{\pgfqpoint{1.577024in}{1.947071in}}%
\pgfpathlineto{\pgfqpoint{1.575739in}{1.960319in}}%
\pgfpathlineto{\pgfqpoint{1.572917in}{1.973567in}}%
\pgfpathlineto{\pgfqpoint{1.570228in}{1.986815in}}%
\pgfpathlineto{\pgfqpoint{1.568956in}{2.000063in}}%
\pgfpathlineto{\pgfqpoint{1.564794in}{2.013311in}}%
\pgfpathlineto{\pgfqpoint{1.564223in}{2.026559in}}%
\pgfpathlineto{\pgfqpoint{1.559881in}{2.039807in}}%
\pgfpathlineto{\pgfqpoint{1.558433in}{2.053055in}}%
\pgfpathlineto{\pgfqpoint{1.555078in}{2.066303in}}%
\pgfpathlineto{\pgfqpoint{1.552302in}{2.077547in}}%
\pgfpathlineto{\pgfqpoint{1.551811in}{2.079551in}}%
\pgfpathlineto{\pgfqpoint{1.550207in}{2.092799in}}%
\pgfpathlineto{\pgfqpoint{1.546000in}{2.106047in}}%
\pgfpathlineto{\pgfqpoint{1.545315in}{2.119295in}}%
\pgfpathlineto{\pgfqpoint{1.540887in}{2.132543in}}%
\pgfpathlineto{\pgfqpoint{1.532871in}{2.132543in}}%
\pgfpathlineto{\pgfqpoint{1.537423in}{2.119295in}}%
\pgfpathlineto{\pgfqpoint{1.538054in}{2.106047in}}%
\pgfpathlineto{\pgfqpoint{1.542505in}{2.092799in}}%
\pgfpathlineto{\pgfqpoint{1.543957in}{2.079551in}}%
\pgfpathlineto{\pgfqpoint{1.547364in}{2.066303in}}%
\pgfpathlineto{\pgfqpoint{1.550586in}{2.053055in}}%
\pgfpathlineto{\pgfqpoint{1.552100in}{2.039807in}}%
\pgfpathlineto{\pgfqpoint{1.552302in}{2.039215in}}%
\pgfpathlineto{\pgfqpoint{1.556608in}{2.026559in}}%
\pgfpathlineto{\pgfqpoint{1.557111in}{2.013311in}}%
\pgfpathlineto{\pgfqpoint{1.561513in}{2.000063in}}%
\pgfpathlineto{\pgfqpoint{1.562663in}{1.986815in}}%
\pgfpathlineto{\pgfqpoint{1.565561in}{1.973567in}}%
\pgfpathlineto{\pgfqpoint{1.568304in}{1.960319in}}%
\pgfpathlineto{\pgfqpoint{1.569674in}{1.947071in}}%
\pgfpathlineto{\pgfqpoint{1.573890in}{1.933823in}}%
\pgfpathlineto{\pgfqpoint{1.574379in}{1.920575in}}%
\pgfpathlineto{\pgfqpoint{1.577585in}{1.910565in}}%
\pgfpathlineto{\pgfqpoint{1.578678in}{1.907327in}}%
\pgfpathlineto{\pgfqpoint{1.579754in}{1.894079in}}%
\pgfpathlineto{\pgfqpoint{1.583008in}{1.880831in}}%
\pgfpathlineto{\pgfqpoint{1.585749in}{1.867583in}}%
\pgfpathlineto{\pgfqpoint{1.587108in}{1.854335in}}%
\pgfpathlineto{\pgfqpoint{1.591200in}{1.841087in}}%
\pgfpathlineto{\pgfqpoint{1.591491in}{1.827839in}}%
\pgfpathlineto{\pgfqpoint{1.595940in}{1.814591in}}%
\pgfpathlineto{\pgfqpoint{1.596653in}{1.801343in}}%
\pgfpathlineto{\pgfqpoint{1.599662in}{1.788095in}}%
\pgfpathlineto{\pgfqpoint{1.602004in}{1.774847in}}%
\pgfpathlineto{\pgfqpoint{1.602868in}{1.766023in}}%
\pgfpathlineto{\pgfqpoint{1.603339in}{1.761599in}}%
\pgfpathlineto{\pgfqpoint{1.607806in}{1.748351in}}%
\pgfpathlineto{\pgfqpoint{1.608134in}{1.735103in}}%
\pgfpathlineto{\pgfqpoint{1.612892in}{1.721855in}}%
\pgfpathlineto{\pgfqpoint{1.613470in}{1.708607in}}%
\pgfpathlineto{\pgfqpoint{1.617159in}{1.695359in}}%
\pgfpathlineto{\pgfqpoint{1.619280in}{1.682111in}}%
\pgfpathlineto{\pgfqpoint{1.620552in}{1.668863in}}%
\pgfpathlineto{\pgfqpoint{1.624414in}{1.655615in}}%
\pgfpathlineto{\pgfqpoint{1.624174in}{1.642367in}}%
\pgfpathlineto{\pgfqpoint{1.628151in}{1.631487in}}%
\pgfpathlineto{\pgfqpoint{1.629089in}{1.629119in}}%
\pgfpathlineto{\pgfqpoint{1.629331in}{1.615871in}}%
\pgfpathlineto{\pgfqpoint{1.633452in}{1.602623in}}%
\pgfpathlineto{\pgfqpoint{1.635531in}{1.589375in}}%
\pgfpathlineto{\pgfqpoint{1.637449in}{1.576127in}}%
\pgfpathlineto{\pgfqpoint{1.641657in}{1.562879in}}%
\pgfpathlineto{\pgfqpoint{1.641364in}{1.549631in}}%
\pgfpathlineto{\pgfqpoint{1.646318in}{1.536383in}}%
\pgfpathlineto{\pgfqpoint{1.645634in}{1.523135in}}%
\pgfpathlineto{\pgfqpoint{1.649920in}{1.509887in}}%
\pgfpathlineto{\pgfqpoint{1.651191in}{1.496639in}}%
\pgfpathlineto{\pgfqpoint{1.653101in}{1.483391in}}%
\pgfpathlineto{\pgfqpoint{1.653434in}{1.482218in}}%
\pgfpathlineto{\pgfqpoint{1.657181in}{1.470143in}}%
\pgfpathlineto{\pgfqpoint{1.656684in}{1.456895in}}%
\pgfpathlineto{\pgfqpoint{1.662556in}{1.443647in}}%
\pgfpathlineto{\pgfqpoint{1.661384in}{1.430399in}}%
\pgfpathlineto{\pgfqpoint{1.666200in}{1.417151in}}%
\pgfpathlineto{\pgfqpoint{1.666604in}{1.403903in}}%
\pgfpathlineto{\pgfqpoint{1.668938in}{1.390655in}}%
\pgfpathlineto{\pgfqpoint{1.672332in}{1.377407in}}%
\pgfpathlineto{\pgfqpoint{1.671569in}{1.364159in}}%
\pgfpathlineto{\pgfqpoint{1.677066in}{1.350911in}}%
\pgfpathlineto{\pgfqpoint{1.674394in}{1.337663in}}%
\pgfpathlineto{\pgfqpoint{1.678717in}{1.327409in}}%
\pgfpathlineto{\pgfqpoint{1.680032in}{1.324415in}}%
\pgfpathlineto{\pgfqpoint{1.678734in}{1.311167in}}%
\pgfpathlineto{\pgfqpoint{1.681321in}{1.297919in}}%
\pgfpathlineto{\pgfqpoint{1.683698in}{1.284671in}}%
\pgfpathlineto{\pgfqpoint{1.682479in}{1.271423in}}%
\pgfpathlineto{\pgfqpoint{1.688339in}{1.258175in}}%
\pgfpathlineto{\pgfqpoint{1.684529in}{1.244927in}}%
\pgfpathlineto{\pgfqpoint{1.690566in}{1.231679in}}%
\pgfpathlineto{\pgfqpoint{1.686135in}{1.218431in}}%
\pgfpathlineto{\pgfqpoint{1.689304in}{1.205183in}}%
\pgfpathlineto{\pgfqpoint{1.688166in}{1.191935in}}%
\pgfpathlineto{\pgfqpoint{1.685522in}{1.178687in}}%
\pgfpathlineto{\pgfqpoint{1.689143in}{1.165439in}}%
\pgfpathlineto{\pgfqpoint{1.681487in}{1.152191in}}%
\pgfpathlineto{\pgfqpoint{1.688266in}{1.138943in}}%
\pgfpathlineto{\pgfqpoint{1.679107in}{1.125695in}}%
\pgfpathlineto{\pgfqpoint{1.682713in}{1.112447in}}%
\pgfpathlineto{\pgfqpoint{1.678717in}{1.105181in}}%
\pgfpathlineto{\pgfqpoint{1.674913in}{1.099199in}}%
\pgfpathlineto{\pgfqpoint{1.671098in}{1.085951in}}%
\pgfpathlineto{\pgfqpoint{1.668525in}{1.072703in}}%
\pgfpathlineto{\pgfqpoint{1.653712in}{1.059455in}}%
\pgfpathlineto{\pgfqpoint{1.656101in}{1.046207in}}%
\pgfpathlineto{\pgfqpoint{1.653434in}{1.044648in}}%
\pgfpathlineto{\pgfqpoint{1.629874in}{1.032959in}}%
\pgfpathlineto{\pgfqpoint{1.632521in}{1.019711in}}%
\pgfpathlineto{\pgfqpoint{1.628151in}{1.017556in}}%
\pgfpathlineto{\pgfqpoint{1.602868in}{1.007629in}}%
\pgfpathlineto{\pgfqpoint{1.599050in}{1.006463in}}%
\pgfpathlineto{\pgfqpoint{1.586924in}{0.993215in}}%
\pgfpathlineto{\pgfqpoint{1.577585in}{0.989595in}}%
\pgfpathlineto{\pgfqpoint{1.552302in}{0.982534in}}%
\pgfpathlineto{\pgfqpoint{1.538637in}{0.979967in}}%
\pgfpathlineto{\pgfqpoint{1.527018in}{0.977024in}}%
\pgfpathlineto{\pgfqpoint{1.501735in}{0.972117in}}%
\pgfpathlineto{\pgfqpoint{1.476452in}{0.968653in}}%
\pgfpathlineto{\pgfqpoint{1.451169in}{0.966800in}}%
\pgfpathlineto{\pgfqpoint{1.425886in}{0.966922in}}%
\pgfpathlineto{\pgfqpoint{1.400603in}{0.969716in}}%
\pgfpathlineto{\pgfqpoint{1.375319in}{0.976528in}}%
\pgfpathlineto{\pgfqpoint{1.367898in}{0.979967in}}%
\pgfpathlineto{\pgfqpoint{1.350036in}{0.982513in}}%
\pgfpathlineto{\pgfqpoint{1.324753in}{0.987167in}}%
\pgfpathlineto{\pgfqpoint{1.299470in}{0.993048in}}%
\pgfpathlineto{\pgfqpoint{1.298870in}{0.993215in}}%
\pgfpathlineto{\pgfqpoint{1.297833in}{1.006463in}}%
\pgfpathlineto{\pgfqpoint{1.274187in}{1.014968in}}%
\pgfpathlineto{\pgfqpoint{1.263903in}{1.019711in}}%
\pgfpathlineto{\pgfqpoint{1.258150in}{1.032959in}}%
\pgfpathlineto{\pgfqpoint{1.248904in}{1.044444in}}%
\pgfpathlineto{\pgfqpoint{1.247759in}{1.046207in}}%
\pgfpathlineto{\pgfqpoint{1.233218in}{1.059455in}}%
\pgfpathlineto{\pgfqpoint{1.238159in}{1.072703in}}%
\pgfpathlineto{\pgfqpoint{1.223620in}{1.084651in}}%
\pgfpathlineto{\pgfqpoint{1.222227in}{1.085951in}}%
\pgfpathlineto{\pgfqpoint{1.223620in}{1.088799in}}%
\pgfpathlineto{\pgfqpoint{1.228790in}{1.099199in}}%
\pgfpathlineto{\pgfqpoint{1.223620in}{1.105039in}}%
\pgfpathlineto{\pgfqpoint{1.217954in}{1.112447in}}%
\pgfpathlineto{\pgfqpoint{1.217301in}{1.125695in}}%
\pgfpathlineto{\pgfqpoint{1.216488in}{1.138943in}}%
\pgfpathlineto{\pgfqpoint{1.210130in}{1.152191in}}%
\pgfpathlineto{\pgfqpoint{1.216217in}{1.165439in}}%
\pgfpathlineto{\pgfqpoint{1.208236in}{1.178687in}}%
\pgfpathlineto{\pgfqpoint{1.214983in}{1.191935in}}%
\pgfpathlineto{\pgfqpoint{1.211260in}{1.205183in}}%
\pgfpathlineto{\pgfqpoint{1.214039in}{1.218431in}}%
\pgfpathlineto{\pgfqpoint{1.215851in}{1.231679in}}%
\pgfpathlineto{\pgfqpoint{1.212993in}{1.244927in}}%
\pgfpathlineto{\pgfqpoint{1.219065in}{1.258175in}}%
\pgfpathlineto{\pgfqpoint{1.214776in}{1.271423in}}%
\pgfpathlineto{\pgfqpoint{1.221276in}{1.284671in}}%
\pgfpathlineto{\pgfqpoint{1.219188in}{1.297919in}}%
\pgfpathlineto{\pgfqpoint{1.222297in}{1.311167in}}%
\pgfpathlineto{\pgfqpoint{1.223620in}{1.317684in}}%
\pgfpathlineto{\pgfqpoint{1.224875in}{1.324415in}}%
\pgfpathlineto{\pgfqpoint{1.224483in}{1.337663in}}%
\pgfpathlineto{\pgfqpoint{1.230180in}{1.350911in}}%
\pgfpathlineto{\pgfqpoint{1.227705in}{1.364159in}}%
\pgfpathlineto{\pgfqpoint{1.233739in}{1.377407in}}%
\pgfpathlineto{\pgfqpoint{1.232896in}{1.390655in}}%
\pgfpathlineto{\pgfqpoint{1.236775in}{1.403903in}}%
\pgfpathlineto{\pgfqpoint{1.239192in}{1.417151in}}%
\pgfpathlineto{\pgfqpoint{1.239662in}{1.430399in}}%
\pgfpathlineto{\pgfqpoint{1.244911in}{1.443647in}}%
\pgfpathlineto{\pgfqpoint{1.243625in}{1.456895in}}%
\pgfpathlineto{\pgfqpoint{1.248904in}{1.468805in}}%
\pgfpathlineto{\pgfqpoint{1.249465in}{1.470143in}}%
\pgfpathlineto{\pgfqpoint{1.248904in}{1.480430in}}%
\pgfpathlineto{\pgfqpoint{1.248738in}{1.483391in}}%
\pgfpathlineto{\pgfqpoint{1.248904in}{1.483938in}}%
\pgfpathlineto{\pgfqpoint{1.252527in}{1.496639in}}%
\pgfpathlineto{\pgfqpoint{1.254659in}{1.509887in}}%
\pgfpathlineto{\pgfqpoint{1.256147in}{1.523135in}}%
\pgfpathlineto{\pgfqpoint{1.260900in}{1.536383in}}%
\pgfpathlineto{\pgfqpoint{1.260579in}{1.549631in}}%
\pgfpathlineto{\pgfqpoint{1.266032in}{1.562879in}}%
\pgfpathlineto{\pgfqpoint{1.265782in}{1.576127in}}%
\pgfpathlineto{\pgfqpoint{1.270006in}{1.589375in}}%
\pgfpathlineto{\pgfqpoint{1.271477in}{1.602623in}}%
\pgfpathlineto{\pgfqpoint{1.273240in}{1.615871in}}%
\pgfpathlineto{\pgfqpoint{1.274187in}{1.618814in}}%
\pgfpathlineto{\pgfqpoint{1.277181in}{1.629119in}}%
\pgfpathlineto{\pgfqpoint{1.277383in}{1.642367in}}%
\pgfpathlineto{\pgfqpoint{1.282482in}{1.655615in}}%
\pgfpathlineto{\pgfqpoint{1.282578in}{1.668863in}}%
\pgfpathlineto{\pgfqpoint{1.286913in}{1.682111in}}%
\pgfpathlineto{\pgfqpoint{1.288335in}{1.695359in}}%
\pgfpathlineto{\pgfqpoint{1.290813in}{1.708607in}}%
\pgfpathlineto{\pgfqpoint{1.294156in}{1.721855in}}%
\pgfpathlineto{\pgfqpoint{1.294677in}{1.735103in}}%
\pgfpathlineto{\pgfqpoint{1.299470in}{1.748191in}}%
\pgfpathlineto{\pgfqpoint{1.299524in}{1.748351in}}%
\pgfpathlineto{\pgfqpoint{1.299657in}{1.761599in}}%
\pgfpathlineto{\pgfqpoint{1.304029in}{1.774847in}}%
\pgfpathlineto{\pgfqpoint{1.305076in}{1.788095in}}%
\pgfpathlineto{\pgfqpoint{1.307873in}{1.801343in}}%
\pgfpathlineto{\pgfqpoint{1.310671in}{1.814591in}}%
\pgfpathlineto{\pgfqpoint{1.311841in}{1.827839in}}%
\pgfpathlineto{\pgfqpoint{1.316212in}{1.841087in}}%
\pgfpathlineto{\pgfqpoint{1.316487in}{1.854335in}}%
\pgfpathlineto{\pgfqpoint{1.321128in}{1.867583in}}%
\pgfpathlineto{\pgfqpoint{1.322152in}{1.880831in}}%
\pgfpathlineto{\pgfqpoint{1.324753in}{1.890390in}}%
\pgfpathlineto{\pgfqpoint{1.325713in}{1.894079in}}%
\pgfpathlineto{\pgfqpoint{1.327983in}{1.907327in}}%
\pgfpathlineto{\pgfqpoint{1.329529in}{1.920575in}}%
\pgfpathlineto{\pgfqpoint{1.333242in}{1.933823in}}%
\pgfpathlineto{\pgfqpoint{1.333710in}{1.947071in}}%
\pgfpathlineto{\pgfqpoint{1.338291in}{1.960319in}}%
\pgfpathlineto{\pgfqpoint{1.338936in}{1.973567in}}%
\pgfpathlineto{\pgfqpoint{1.342878in}{1.986815in}}%
\pgfpathlineto{\pgfqpoint{1.345076in}{2.000063in}}%
\pgfpathlineto{\pgfqpoint{1.347630in}{2.013311in}}%
\pgfpathlineto{\pgfqpoint{1.350036in}{2.021769in}}%
\pgfpathlineto{\pgfqpoint{1.351357in}{2.026559in}}%
\pgfpathlineto{\pgfqpoint{1.352261in}{2.039807in}}%
\pgfpathlineto{\pgfqpoint{1.356731in}{2.053055in}}%
\pgfpathlineto{\pgfqpoint{1.357245in}{2.066303in}}%
\pgfpathlineto{\pgfqpoint{1.361699in}{2.079551in}}%
\pgfpathlineto{\pgfqpoint{1.363187in}{2.092799in}}%
\pgfpathlineto{\pgfqpoint{1.366310in}{2.106047in}}%
\pgfpathlineto{\pgfqpoint{1.369689in}{2.119295in}}%
\pgfpathlineto{\pgfqpoint{1.371367in}{2.132543in}}%
\pgfpathlineto{\pgfqpoint{1.363382in}{2.132543in}}%
\pgfpathlineto{\pgfqpoint{1.361831in}{2.119295in}}%
\pgfpathlineto{\pgfqpoint{1.358405in}{2.106047in}}%
\pgfpathlineto{\pgfqpoint{1.355538in}{2.092799in}}%
\pgfpathlineto{\pgfqpoint{1.353903in}{2.079551in}}%
\pgfpathlineto{\pgfqpoint{1.350036in}{2.067396in}}%
\pgfpathlineto{\pgfqpoint{1.349681in}{2.066303in}}%
\pgfpathlineto{\pgfqpoint{1.349045in}{2.053055in}}%
\pgfpathlineto{\pgfqpoint{1.344536in}{2.039807in}}%
\pgfpathlineto{\pgfqpoint{1.343646in}{2.026559in}}%
\pgfpathlineto{\pgfqpoint{1.339916in}{2.013311in}}%
\pgfpathlineto{\pgfqpoint{1.337564in}{2.000063in}}%
\pgfpathlineto{\pgfqpoint{1.335269in}{1.986815in}}%
\pgfpathlineto{\pgfqpoint{1.331585in}{1.973567in}}%
\pgfpathlineto{\pgfqpoint{1.330813in}{1.960319in}}%
\pgfpathlineto{\pgfqpoint{1.326389in}{1.947071in}}%
\pgfpathlineto{\pgfqpoint{1.325892in}{1.933823in}}%
\pgfpathlineto{\pgfqpoint{1.324753in}{1.929783in}}%
\pgfpathlineto{\pgfqpoint{1.321894in}{1.920575in}}%
\pgfpathlineto{\pgfqpoint{1.320323in}{1.907327in}}%
\pgfpathlineto{\pgfqpoint{1.317748in}{1.894079in}}%
\pgfpathlineto{\pgfqpoint{1.314364in}{1.880831in}}%
\pgfpathlineto{\pgfqpoint{1.313203in}{1.867583in}}%
\pgfpathlineto{\pgfqpoint{1.308796in}{1.854335in}}%
\pgfpathlineto{\pgfqpoint{1.308444in}{1.841087in}}%
\pgfpathlineto{\pgfqpoint{1.304132in}{1.827839in}}%
\pgfpathlineto{\pgfqpoint{1.303022in}{1.814591in}}%
\pgfpathlineto{\pgfqpoint{1.300079in}{1.801343in}}%
\pgfpathlineto{\pgfqpoint{1.299470in}{1.798331in}}%
\pgfpathlineto{\pgfqpoint{1.297194in}{1.788095in}}%
\pgfpathlineto{\pgfqpoint{1.295799in}{1.774847in}}%
\pgfpathlineto{\pgfqpoint{1.291202in}{1.761599in}}%
\pgfpathlineto{\pgfqpoint{1.290905in}{1.748351in}}%
\pgfpathlineto{\pgfqpoint{1.286187in}{1.735103in}}%
\pgfpathlineto{\pgfqpoint{1.285667in}{1.721855in}}%
\pgfpathlineto{\pgfqpoint{1.282233in}{1.708607in}}%
\pgfpathlineto{\pgfqpoint{1.279905in}{1.695359in}}%
\pgfpathlineto{\pgfqpoint{1.278233in}{1.682111in}}%
\pgfpathlineto{\pgfqpoint{1.274187in}{1.669217in}}%
\pgfpathlineto{\pgfqpoint{1.274065in}{1.668863in}}%
\pgfpathlineto{\pgfqpoint{1.273701in}{1.655615in}}%
\pgfpathlineto{\pgfqpoint{1.268135in}{1.642367in}}%
\pgfpathlineto{\pgfqpoint{1.267826in}{1.629119in}}%
\pgfpathlineto{\pgfqpoint{1.263526in}{1.615871in}}%
\pgfpathlineto{\pgfqpoint{1.261867in}{1.602623in}}%
\pgfpathlineto{\pgfqpoint{1.260102in}{1.589375in}}%
\pgfpathlineto{\pgfqpoint{1.256065in}{1.576127in}}%
\pgfpathlineto{\pgfqpoint{1.255977in}{1.562879in}}%
\pgfpathlineto{\pgfqpoint{1.250615in}{1.549631in}}%
\pgfpathlineto{\pgfqpoint{1.250766in}{1.536383in}}%
\pgfpathlineto{\pgfqpoint{1.248904in}{1.531317in}}%
\pgfpathlineto{\pgfqpoint{1.245568in}{1.523135in}}%
\pgfpathlineto{\pgfqpoint{1.243999in}{1.509887in}}%
\pgfpathlineto{\pgfqpoint{1.241303in}{1.496639in}}%
\pgfpathlineto{\pgfqpoint{1.237435in}{1.483391in}}%
\pgfpathlineto{\pgfqpoint{1.237718in}{1.470143in}}%
\pgfpathlineto{\pgfqpoint{1.231934in}{1.456895in}}%
\pgfpathlineto{\pgfqpoint{1.232862in}{1.443647in}}%
\pgfpathlineto{\pgfqpoint{1.227404in}{1.430399in}}%
\pgfpathlineto{\pgfqpoint{1.226884in}{1.417151in}}%
\pgfpathlineto{\pgfqpoint{1.223941in}{1.403903in}}%
\pgfpathlineto{\pgfqpoint{1.223620in}{1.402780in}}%
\pgfpathlineto{\pgfqpoint{1.219965in}{1.390655in}}%
\pgfpathlineto{\pgfqpoint{1.220171in}{1.377407in}}%
\pgfpathlineto{\pgfqpoint{1.213865in}{1.364159in}}%
\pgfpathlineto{\pgfqpoint{1.215897in}{1.350911in}}%
\pgfpathlineto{\pgfqpoint{1.209638in}{1.337663in}}%
\pgfpathlineto{\pgfqpoint{1.209801in}{1.324415in}}%
\pgfpathlineto{\pgfqpoint{1.206394in}{1.311167in}}%
\pgfpathlineto{\pgfqpoint{1.203346in}{1.297919in}}%
\pgfpathlineto{\pgfqpoint{1.204403in}{1.284671in}}%
\pgfpathlineto{\pgfqpoint{1.198337in}{1.272371in}}%
\pgfpathlineto{\pgfqpoint{1.197851in}{1.271423in}}%
\pgfpathlineto{\pgfqpoint{1.198337in}{1.269542in}}%
\pgfpathlineto{\pgfqpoint{1.201308in}{1.258175in}}%
\pgfpathlineto{\pgfqpoint{1.198337in}{1.252169in}}%
\pgfpathlineto{\pgfqpoint{1.194539in}{1.244927in}}%
\pgfpathlineto{\pgfqpoint{1.197016in}{1.231679in}}%
\pgfpathlineto{\pgfqpoint{1.193859in}{1.218431in}}%
\pgfpathlineto{\pgfqpoint{1.190929in}{1.205183in}}%
\pgfpathlineto{\pgfqpoint{1.192973in}{1.191935in}}%
\pgfpathlineto{\pgfqpoint{1.185939in}{1.178687in}}%
\pgfpathlineto{\pgfqpoint{1.192410in}{1.165439in}}%
\pgfpathlineto{\pgfqpoint{1.185228in}{1.152191in}}%
\pgfpathlineto{\pgfqpoint{1.190686in}{1.138943in}}%
\pgfpathlineto{\pgfqpoint{1.189118in}{1.125695in}}%
\pgfpathlineto{\pgfqpoint{1.189389in}{1.112447in}}%
\pgfpathlineto{\pgfqpoint{1.196728in}{1.099199in}}%
\pgfpathlineto{\pgfqpoint{1.189737in}{1.085951in}}%
\pgfpathlineto{\pgfqpoint{1.198337in}{1.076371in}}%
\pgfpathlineto{\pgfqpoint{1.202009in}{1.072703in}}%
\pgfpathlineto{\pgfqpoint{1.198337in}{1.065996in}}%
\pgfpathlineto{\pgfqpoint{1.194631in}{1.059455in}}%
\pgfpathlineto{\pgfqpoint{1.198337in}{1.055209in}}%
\pgfpathlineto{\pgfqpoint{1.206826in}{1.046207in}}%
\pgfpathlineto{\pgfqpoint{1.210579in}{1.032959in}}%
\pgfpathlineto{\pgfqpoint{1.214731in}{1.019711in}}%
\pgfpathlineto{\pgfqpoint{1.223620in}{1.013498in}}%
\pgfpathlineto{\pgfqpoint{1.235904in}{1.006463in}}%
\pgfpathlineto{\pgfqpoint{1.233360in}{0.993215in}}%
\pgfpathlineto{\pgfqpoint{1.248904in}{0.987304in}}%
\pgfpathlineto{\pgfqpoint{1.272141in}{0.979967in}}%
\pgfpathlineto{\pgfqpoint{1.274187in}{0.971718in}}%
\pgfpathlineto{\pgfqpoint{1.275781in}{0.966719in}}%
\pgfpathlineto{\pgfqpoint{1.299470in}{0.958292in}}%
\pgfpathlineto{\pgfqpoint{1.315353in}{0.953471in}}%
\pgfpathlineto{\pgfqpoint{1.324753in}{0.950397in}}%
\pgfpathlineto{\pgfqpoint{1.350036in}{0.944500in}}%
\pgfpathlineto{\pgfqpoint{1.375319in}{0.940747in}}%
\pgfpathlineto{\pgfqpoint{1.381063in}{0.940223in}}%
\pgfpathlineto{\pgfqpoint{1.400603in}{0.934056in}}%
\pgfpathclose%
\pgfusepath{fill}%
\end{pgfscope}%
\begin{pgfscope}%
\pgfpathrectangle{\pgfqpoint{0.540975in}{0.476543in}}{\pgfqpoint{1.820388in}{1.656000in}}%
\pgfusepath{clip}%
\pgfsetbuttcap%
\pgfsetroundjoin%
\definecolor{currentfill}{rgb}{0.870281,0.948943,0.970242}%
\pgfsetfillcolor{currentfill}%
\pgfsetlinewidth{0.000000pt}%
\definecolor{currentstroke}{rgb}{0.000000,0.000000,0.000000}%
\pgfsetstrokecolor{currentstroke}%
\pgfsetdash{}{0pt}%
\pgfpathmoveto{\pgfqpoint{1.375319in}{0.976528in}}%
\pgfpathlineto{\pgfqpoint{1.400603in}{0.969716in}}%
\pgfpathlineto{\pgfqpoint{1.425886in}{0.966922in}}%
\pgfpathlineto{\pgfqpoint{1.451169in}{0.966800in}}%
\pgfpathlineto{\pgfqpoint{1.476452in}{0.968653in}}%
\pgfpathlineto{\pgfqpoint{1.501735in}{0.972117in}}%
\pgfpathlineto{\pgfqpoint{1.527018in}{0.977024in}}%
\pgfpathlineto{\pgfqpoint{1.538637in}{0.979967in}}%
\pgfpathlineto{\pgfqpoint{1.552302in}{0.982534in}}%
\pgfpathlineto{\pgfqpoint{1.577585in}{0.989595in}}%
\pgfpathlineto{\pgfqpoint{1.586924in}{0.993215in}}%
\pgfpathlineto{\pgfqpoint{1.599050in}{1.006463in}}%
\pgfpathlineto{\pgfqpoint{1.602868in}{1.007629in}}%
\pgfpathlineto{\pgfqpoint{1.628151in}{1.017556in}}%
\pgfpathlineto{\pgfqpoint{1.632521in}{1.019711in}}%
\pgfpathlineto{\pgfqpoint{1.629874in}{1.032959in}}%
\pgfpathlineto{\pgfqpoint{1.653434in}{1.044648in}}%
\pgfpathlineto{\pgfqpoint{1.656101in}{1.046207in}}%
\pgfpathlineto{\pgfqpoint{1.653712in}{1.059455in}}%
\pgfpathlineto{\pgfqpoint{1.668525in}{1.072703in}}%
\pgfpathlineto{\pgfqpoint{1.671098in}{1.085951in}}%
\pgfpathlineto{\pgfqpoint{1.674913in}{1.099199in}}%
\pgfpathlineto{\pgfqpoint{1.678717in}{1.105181in}}%
\pgfpathlineto{\pgfqpoint{1.682713in}{1.112447in}}%
\pgfpathlineto{\pgfqpoint{1.679107in}{1.125695in}}%
\pgfpathlineto{\pgfqpoint{1.688266in}{1.138943in}}%
\pgfpathlineto{\pgfqpoint{1.681487in}{1.152191in}}%
\pgfpathlineto{\pgfqpoint{1.689143in}{1.165439in}}%
\pgfpathlineto{\pgfqpoint{1.685522in}{1.178687in}}%
\pgfpathlineto{\pgfqpoint{1.688166in}{1.191935in}}%
\pgfpathlineto{\pgfqpoint{1.689304in}{1.205183in}}%
\pgfpathlineto{\pgfqpoint{1.686135in}{1.218431in}}%
\pgfpathlineto{\pgfqpoint{1.690566in}{1.231679in}}%
\pgfpathlineto{\pgfqpoint{1.684529in}{1.244927in}}%
\pgfpathlineto{\pgfqpoint{1.688339in}{1.258175in}}%
\pgfpathlineto{\pgfqpoint{1.682479in}{1.271423in}}%
\pgfpathlineto{\pgfqpoint{1.683698in}{1.284671in}}%
\pgfpathlineto{\pgfqpoint{1.681321in}{1.297919in}}%
\pgfpathlineto{\pgfqpoint{1.678734in}{1.311167in}}%
\pgfpathlineto{\pgfqpoint{1.680032in}{1.324415in}}%
\pgfpathlineto{\pgfqpoint{1.678717in}{1.327409in}}%
\pgfpathlineto{\pgfqpoint{1.674394in}{1.337663in}}%
\pgfpathlineto{\pgfqpoint{1.677066in}{1.350911in}}%
\pgfpathlineto{\pgfqpoint{1.671569in}{1.364159in}}%
\pgfpathlineto{\pgfqpoint{1.672332in}{1.377407in}}%
\pgfpathlineto{\pgfqpoint{1.668938in}{1.390655in}}%
\pgfpathlineto{\pgfqpoint{1.666604in}{1.403903in}}%
\pgfpathlineto{\pgfqpoint{1.666200in}{1.417151in}}%
\pgfpathlineto{\pgfqpoint{1.661384in}{1.430399in}}%
\pgfpathlineto{\pgfqpoint{1.662556in}{1.443647in}}%
\pgfpathlineto{\pgfqpoint{1.656684in}{1.456895in}}%
\pgfpathlineto{\pgfqpoint{1.657181in}{1.470143in}}%
\pgfpathlineto{\pgfqpoint{1.653434in}{1.482218in}}%
\pgfpathlineto{\pgfqpoint{1.653101in}{1.483391in}}%
\pgfpathlineto{\pgfqpoint{1.651191in}{1.496639in}}%
\pgfpathlineto{\pgfqpoint{1.649920in}{1.509887in}}%
\pgfpathlineto{\pgfqpoint{1.645634in}{1.523135in}}%
\pgfpathlineto{\pgfqpoint{1.646318in}{1.536383in}}%
\pgfpathlineto{\pgfqpoint{1.641364in}{1.549631in}}%
\pgfpathlineto{\pgfqpoint{1.641657in}{1.562879in}}%
\pgfpathlineto{\pgfqpoint{1.637449in}{1.576127in}}%
\pgfpathlineto{\pgfqpoint{1.635531in}{1.589375in}}%
\pgfpathlineto{\pgfqpoint{1.633452in}{1.602623in}}%
\pgfpathlineto{\pgfqpoint{1.629331in}{1.615871in}}%
\pgfpathlineto{\pgfqpoint{1.629089in}{1.629119in}}%
\pgfpathlineto{\pgfqpoint{1.628151in}{1.631487in}}%
\pgfpathlineto{\pgfqpoint{1.624174in}{1.642367in}}%
\pgfpathlineto{\pgfqpoint{1.624414in}{1.655615in}}%
\pgfpathlineto{\pgfqpoint{1.620552in}{1.668863in}}%
\pgfpathlineto{\pgfqpoint{1.619280in}{1.682111in}}%
\pgfpathlineto{\pgfqpoint{1.617159in}{1.695359in}}%
\pgfpathlineto{\pgfqpoint{1.613470in}{1.708607in}}%
\pgfpathlineto{\pgfqpoint{1.612892in}{1.721855in}}%
\pgfpathlineto{\pgfqpoint{1.608134in}{1.735103in}}%
\pgfpathlineto{\pgfqpoint{1.607806in}{1.748351in}}%
\pgfpathlineto{\pgfqpoint{1.603339in}{1.761599in}}%
\pgfpathlineto{\pgfqpoint{1.602868in}{1.766023in}}%
\pgfpathlineto{\pgfqpoint{1.602004in}{1.774847in}}%
\pgfpathlineto{\pgfqpoint{1.599662in}{1.788095in}}%
\pgfpathlineto{\pgfqpoint{1.596653in}{1.801343in}}%
\pgfpathlineto{\pgfqpoint{1.595940in}{1.814591in}}%
\pgfpathlineto{\pgfqpoint{1.591491in}{1.827839in}}%
\pgfpathlineto{\pgfqpoint{1.591200in}{1.841087in}}%
\pgfpathlineto{\pgfqpoint{1.587108in}{1.854335in}}%
\pgfpathlineto{\pgfqpoint{1.585749in}{1.867583in}}%
\pgfpathlineto{\pgfqpoint{1.583008in}{1.880831in}}%
\pgfpathlineto{\pgfqpoint{1.579754in}{1.894079in}}%
\pgfpathlineto{\pgfqpoint{1.578678in}{1.907327in}}%
\pgfpathlineto{\pgfqpoint{1.577585in}{1.910565in}}%
\pgfpathlineto{\pgfqpoint{1.574379in}{1.920575in}}%
\pgfpathlineto{\pgfqpoint{1.573890in}{1.933823in}}%
\pgfpathlineto{\pgfqpoint{1.569674in}{1.947071in}}%
\pgfpathlineto{\pgfqpoint{1.568304in}{1.960319in}}%
\pgfpathlineto{\pgfqpoint{1.565561in}{1.973567in}}%
\pgfpathlineto{\pgfqpoint{1.562663in}{1.986815in}}%
\pgfpathlineto{\pgfqpoint{1.561513in}{2.000063in}}%
\pgfpathlineto{\pgfqpoint{1.557111in}{2.013311in}}%
\pgfpathlineto{\pgfqpoint{1.556608in}{2.026559in}}%
\pgfpathlineto{\pgfqpoint{1.552302in}{2.039215in}}%
\pgfpathlineto{\pgfqpoint{1.552100in}{2.039807in}}%
\pgfpathlineto{\pgfqpoint{1.550586in}{2.053055in}}%
\pgfpathlineto{\pgfqpoint{1.547364in}{2.066303in}}%
\pgfpathlineto{\pgfqpoint{1.543957in}{2.079551in}}%
\pgfpathlineto{\pgfqpoint{1.542505in}{2.092799in}}%
\pgfpathlineto{\pgfqpoint{1.538054in}{2.106047in}}%
\pgfpathlineto{\pgfqpoint{1.537423in}{2.119295in}}%
\pgfpathlineto{\pgfqpoint{1.532871in}{2.132543in}}%
\pgfpathlineto{\pgfqpoint{1.527018in}{2.132543in}}%
\pgfpathlineto{\pgfqpoint{1.524611in}{2.132543in}}%
\pgfpathlineto{\pgfqpoint{1.527018in}{2.126464in}}%
\pgfpathlineto{\pgfqpoint{1.529530in}{2.119295in}}%
\pgfpathlineto{\pgfqpoint{1.530107in}{2.106047in}}%
\pgfpathlineto{\pgfqpoint{1.534803in}{2.092799in}}%
\pgfpathlineto{\pgfqpoint{1.536103in}{2.079551in}}%
\pgfpathlineto{\pgfqpoint{1.539759in}{2.066303in}}%
\pgfpathlineto{\pgfqpoint{1.542867in}{2.053055in}}%
\pgfpathlineto{\pgfqpoint{1.544492in}{2.039807in}}%
\pgfpathlineto{\pgfqpoint{1.549006in}{2.026559in}}%
\pgfpathlineto{\pgfqpoint{1.549441in}{2.013311in}}%
\pgfpathlineto{\pgfqpoint{1.552302in}{2.005209in}}%
\pgfpathlineto{\pgfqpoint{1.554070in}{2.000063in}}%
\pgfpathlineto{\pgfqpoint{1.555098in}{1.986815in}}%
\pgfpathlineto{\pgfqpoint{1.558204in}{1.973567in}}%
\pgfpathlineto{\pgfqpoint{1.560868in}{1.960319in}}%
\pgfpathlineto{\pgfqpoint{1.562323in}{1.947071in}}%
\pgfpathlineto{\pgfqpoint{1.566569in}{1.933823in}}%
\pgfpathlineto{\pgfqpoint{1.566969in}{1.920575in}}%
\pgfpathlineto{\pgfqpoint{1.571336in}{1.907327in}}%
\pgfpathlineto{\pgfqpoint{1.572163in}{1.894079in}}%
\pgfpathlineto{\pgfqpoint{1.575414in}{1.880831in}}%
\pgfpathlineto{\pgfqpoint{1.577585in}{1.869444in}}%
\pgfpathlineto{\pgfqpoint{1.577964in}{1.867583in}}%
\pgfpathlineto{\pgfqpoint{1.579400in}{1.854335in}}%
\pgfpathlineto{\pgfqpoint{1.583561in}{1.841087in}}%
\pgfpathlineto{\pgfqpoint{1.583761in}{1.827839in}}%
\pgfpathlineto{\pgfqpoint{1.588363in}{1.814591in}}%
\pgfpathlineto{\pgfqpoint{1.588875in}{1.801343in}}%
\pgfpathlineto{\pgfqpoint{1.592031in}{1.788095in}}%
\pgfpathlineto{\pgfqpoint{1.594184in}{1.774847in}}%
\pgfpathlineto{\pgfqpoint{1.595536in}{1.761599in}}%
\pgfpathlineto{\pgfqpoint{1.599622in}{1.748351in}}%
\pgfpathlineto{\pgfqpoint{1.599832in}{1.735103in}}%
\pgfpathlineto{\pgfqpoint{1.602868in}{1.726415in}}%
\pgfpathlineto{\pgfqpoint{1.604546in}{1.721855in}}%
\pgfpathlineto{\pgfqpoint{1.604910in}{1.708607in}}%
\pgfpathlineto{\pgfqpoint{1.608785in}{1.695359in}}%
\pgfpathlineto{\pgfqpoint{1.610700in}{1.682111in}}%
\pgfpathlineto{\pgfqpoint{1.612009in}{1.668863in}}%
\pgfpathlineto{\pgfqpoint{1.615783in}{1.655615in}}%
\pgfpathlineto{\pgfqpoint{1.615396in}{1.642367in}}%
\pgfpathlineto{\pgfqpoint{1.620285in}{1.629119in}}%
\pgfpathlineto{\pgfqpoint{1.620226in}{1.615871in}}%
\pgfpathlineto{\pgfqpoint{1.624185in}{1.602623in}}%
\pgfpathlineto{\pgfqpoint{1.625857in}{1.589375in}}%
\pgfpathlineto{\pgfqpoint{1.627721in}{1.576127in}}%
\pgfpathlineto{\pgfqpoint{1.628151in}{1.574626in}}%
\pgfpathlineto{\pgfqpoint{1.631774in}{1.562879in}}%
\pgfpathlineto{\pgfqpoint{1.631309in}{1.549631in}}%
\pgfpathlineto{\pgfqpoint{1.636309in}{1.536383in}}%
\pgfpathlineto{\pgfqpoint{1.635252in}{1.523135in}}%
\pgfpathlineto{\pgfqpoint{1.639656in}{1.509887in}}%
\pgfpathlineto{\pgfqpoint{1.640503in}{1.496639in}}%
\pgfpathlineto{\pgfqpoint{1.642434in}{1.483391in}}%
\pgfpathlineto{\pgfqpoint{1.645963in}{1.470143in}}%
\pgfpathlineto{\pgfqpoint{1.645282in}{1.456895in}}%
\pgfpathlineto{\pgfqpoint{1.650767in}{1.443647in}}%
\pgfpathlineto{\pgfqpoint{1.649184in}{1.430399in}}%
\pgfpathlineto{\pgfqpoint{1.653434in}{1.418487in}}%
\pgfpathlineto{\pgfqpoint{1.653927in}{1.417151in}}%
\pgfpathlineto{\pgfqpoint{1.653745in}{1.403903in}}%
\pgfpathlineto{\pgfqpoint{1.656128in}{1.390655in}}%
\pgfpathlineto{\pgfqpoint{1.659027in}{1.377407in}}%
\pgfpathlineto{\pgfqpoint{1.658005in}{1.364159in}}%
\pgfpathlineto{\pgfqpoint{1.663211in}{1.350911in}}%
\pgfpathlineto{\pgfqpoint{1.659876in}{1.337663in}}%
\pgfpathlineto{\pgfqpoint{1.665349in}{1.324415in}}%
\pgfpathlineto{\pgfqpoint{1.663147in}{1.311167in}}%
\pgfpathlineto{\pgfqpoint{1.665612in}{1.297919in}}%
\pgfpathlineto{\pgfqpoint{1.667006in}{1.284671in}}%
\pgfpathlineto{\pgfqpoint{1.665519in}{1.271423in}}%
\pgfpathlineto{\pgfqpoint{1.670658in}{1.258175in}}%
\pgfpathlineto{\pgfqpoint{1.666023in}{1.244927in}}%
\pgfpathlineto{\pgfqpoint{1.671736in}{1.231679in}}%
\pgfpathlineto{\pgfqpoint{1.665840in}{1.218431in}}%
\pgfpathlineto{\pgfqpoint{1.668851in}{1.205183in}}%
\pgfpathlineto{\pgfqpoint{1.665970in}{1.191935in}}%
\pgfpathlineto{\pgfqpoint{1.662807in}{1.178687in}}%
\pgfpathlineto{\pgfqpoint{1.664934in}{1.165439in}}%
\pgfpathlineto{\pgfqpoint{1.655863in}{1.152191in}}%
\pgfpathlineto{\pgfqpoint{1.661706in}{1.138943in}}%
\pgfpathlineto{\pgfqpoint{1.653434in}{1.129456in}}%
\pgfpathlineto{\pgfqpoint{1.649659in}{1.125695in}}%
\pgfpathlineto{\pgfqpoint{1.652919in}{1.112447in}}%
\pgfpathlineto{\pgfqpoint{1.640744in}{1.099199in}}%
\pgfpathlineto{\pgfqpoint{1.636160in}{1.085951in}}%
\pgfpathlineto{\pgfqpoint{1.629326in}{1.072703in}}%
\pgfpathlineto{\pgfqpoint{1.628151in}{1.071782in}}%
\pgfpathlineto{\pgfqpoint{1.610303in}{1.059455in}}%
\pgfpathlineto{\pgfqpoint{1.609166in}{1.046207in}}%
\pgfpathlineto{\pgfqpoint{1.602868in}{1.043551in}}%
\pgfpathlineto{\pgfqpoint{1.577585in}{1.034776in}}%
\pgfpathlineto{\pgfqpoint{1.571072in}{1.032959in}}%
\pgfpathlineto{\pgfqpoint{1.570952in}{1.019711in}}%
\pgfpathlineto{\pgfqpoint{1.552302in}{1.014947in}}%
\pgfpathlineto{\pgfqpoint{1.527018in}{1.010283in}}%
\pgfpathlineto{\pgfqpoint{1.501735in}{1.007156in}}%
\pgfpathlineto{\pgfqpoint{1.492315in}{1.006463in}}%
\pgfpathlineto{\pgfqpoint{1.476452in}{1.000641in}}%
\pgfpathlineto{\pgfqpoint{1.451169in}{0.996300in}}%
\pgfpathlineto{\pgfqpoint{1.425886in}{0.997847in}}%
\pgfpathlineto{\pgfqpoint{1.403746in}{1.006463in}}%
\pgfpathlineto{\pgfqpoint{1.400603in}{1.006661in}}%
\pgfpathlineto{\pgfqpoint{1.375319in}{1.009463in}}%
\pgfpathlineto{\pgfqpoint{1.350036in}{1.013695in}}%
\pgfpathlineto{\pgfqpoint{1.324753in}{1.019686in}}%
\pgfpathlineto{\pgfqpoint{1.324669in}{1.019711in}}%
\pgfpathlineto{\pgfqpoint{1.316528in}{1.032959in}}%
\pgfpathlineto{\pgfqpoint{1.299470in}{1.042551in}}%
\pgfpathlineto{\pgfqpoint{1.294639in}{1.046207in}}%
\pgfpathlineto{\pgfqpoint{1.275970in}{1.059455in}}%
\pgfpathlineto{\pgfqpoint{1.277474in}{1.072703in}}%
\pgfpathlineto{\pgfqpoint{1.274187in}{1.074558in}}%
\pgfpathlineto{\pgfqpoint{1.257098in}{1.085951in}}%
\pgfpathlineto{\pgfqpoint{1.262917in}{1.099199in}}%
\pgfpathlineto{\pgfqpoint{1.248904in}{1.111120in}}%
\pgfpathlineto{\pgfqpoint{1.247590in}{1.112447in}}%
\pgfpathlineto{\pgfqpoint{1.246519in}{1.125695in}}%
\pgfpathlineto{\pgfqpoint{1.242922in}{1.138943in}}%
\pgfpathlineto{\pgfqpoint{1.235574in}{1.152191in}}%
\pgfpathlineto{\pgfqpoint{1.240396in}{1.165439in}}%
\pgfpathlineto{\pgfqpoint{1.230710in}{1.178687in}}%
\pgfpathlineto{\pgfqpoint{1.237168in}{1.191935in}}%
\pgfpathlineto{\pgfqpoint{1.231589in}{1.205183in}}%
\pgfpathlineto{\pgfqpoint{1.234251in}{1.218431in}}%
\pgfpathlineto{\pgfqpoint{1.234649in}{1.231679in}}%
\pgfpathlineto{\pgfqpoint{1.231370in}{1.244927in}}%
\pgfpathlineto{\pgfqpoint{1.236700in}{1.258175in}}%
\pgfpathlineto{\pgfqpoint{1.231621in}{1.271423in}}%
\pgfpathlineto{\pgfqpoint{1.237911in}{1.284671in}}%
\pgfpathlineto{\pgfqpoint{1.234794in}{1.297919in}}%
\pgfpathlineto{\pgfqpoint{1.237851in}{1.311167in}}%
\pgfpathlineto{\pgfqpoint{1.239474in}{1.324415in}}%
\pgfpathlineto{\pgfqpoint{1.238847in}{1.337663in}}%
\pgfpathlineto{\pgfqpoint{1.244097in}{1.350911in}}%
\pgfpathlineto{\pgfqpoint{1.241085in}{1.364159in}}%
\pgfpathlineto{\pgfqpoint{1.247109in}{1.377407in}}%
\pgfpathlineto{\pgfqpoint{1.245599in}{1.390655in}}%
\pgfpathlineto{\pgfqpoint{1.248904in}{1.401595in}}%
\pgfpathlineto{\pgfqpoint{1.249588in}{1.403903in}}%
\pgfpathlineto{\pgfqpoint{1.251402in}{1.417151in}}%
\pgfpathlineto{\pgfqpoint{1.251807in}{1.430399in}}%
\pgfpathlineto{\pgfqpoint{1.256553in}{1.443647in}}%
\pgfpathlineto{\pgfqpoint{1.254994in}{1.456895in}}%
\pgfpathlineto{\pgfqpoint{1.260523in}{1.470143in}}%
\pgfpathlineto{\pgfqpoint{1.259313in}{1.483391in}}%
\pgfpathlineto{\pgfqpoint{1.263207in}{1.496639in}}%
\pgfpathlineto{\pgfqpoint{1.264918in}{1.509887in}}%
\pgfpathlineto{\pgfqpoint{1.266447in}{1.523135in}}%
\pgfpathlineto{\pgfqpoint{1.271034in}{1.536383in}}%
\pgfpathlineto{\pgfqpoint{1.270543in}{1.549631in}}%
\pgfpathlineto{\pgfqpoint{1.274187in}{1.558364in}}%
\pgfpathlineto{\pgfqpoint{1.275978in}{1.562879in}}%
\pgfpathlineto{\pgfqpoint{1.275419in}{1.576127in}}%
\pgfpathlineto{\pgfqpoint{1.279524in}{1.589375in}}%
\pgfpathlineto{\pgfqpoint{1.280562in}{1.602623in}}%
\pgfpathlineto{\pgfqpoint{1.282276in}{1.615871in}}%
\pgfpathlineto{\pgfqpoint{1.285951in}{1.629119in}}%
\pgfpathlineto{\pgfqpoint{1.286077in}{1.642367in}}%
\pgfpathlineto{\pgfqpoint{1.291212in}{1.655615in}}%
\pgfpathlineto{\pgfqpoint{1.291077in}{1.668863in}}%
\pgfpathlineto{\pgfqpoint{1.295594in}{1.682111in}}%
\pgfpathlineto{\pgfqpoint{1.296766in}{1.695359in}}%
\pgfpathlineto{\pgfqpoint{1.299394in}{1.708607in}}%
\pgfpathlineto{\pgfqpoint{1.299470in}{1.708913in}}%
\pgfpathlineto{\pgfqpoint{1.302411in}{1.721855in}}%
\pgfpathlineto{\pgfqpoint{1.302900in}{1.735103in}}%
\pgfpathlineto{\pgfqpoint{1.307393in}{1.748351in}}%
\pgfpathlineto{\pgfqpoint{1.307378in}{1.761599in}}%
\pgfpathlineto{\pgfqpoint{1.311891in}{1.774847in}}%
\pgfpathlineto{\pgfqpoint{1.312709in}{1.788095in}}%
\pgfpathlineto{\pgfqpoint{1.315667in}{1.801343in}}%
\pgfpathlineto{\pgfqpoint{1.318320in}{1.814591in}}%
\pgfpathlineto{\pgfqpoint{1.319551in}{1.827839in}}%
\pgfpathlineto{\pgfqpoint{1.323980in}{1.841087in}}%
\pgfpathlineto{\pgfqpoint{1.324179in}{1.854335in}}%
\pgfpathlineto{\pgfqpoint{1.324753in}{1.855937in}}%
\pgfpathlineto{\pgfqpoint{1.328785in}{1.867583in}}%
\pgfpathlineto{\pgfqpoint{1.329585in}{1.880831in}}%
\pgfpathlineto{\pgfqpoint{1.333158in}{1.894079in}}%
\pgfpathlineto{\pgfqpoint{1.335251in}{1.907327in}}%
\pgfpathlineto{\pgfqpoint{1.336912in}{1.920575in}}%
\pgfpathlineto{\pgfqpoint{1.340591in}{1.933823in}}%
\pgfpathlineto{\pgfqpoint{1.341032in}{1.947071in}}%
\pgfpathlineto{\pgfqpoint{1.345768in}{1.960319in}}%
\pgfpathlineto{\pgfqpoint{1.346288in}{1.973567in}}%
\pgfpathlineto{\pgfqpoint{1.350036in}{1.985437in}}%
\pgfpathlineto{\pgfqpoint{1.350484in}{1.986815in}}%
\pgfpathlineto{\pgfqpoint{1.352543in}{2.000063in}}%
\pgfpathlineto{\pgfqpoint{1.355251in}{2.013311in}}%
\pgfpathlineto{\pgfqpoint{1.358835in}{2.026559in}}%
\pgfpathlineto{\pgfqpoint{1.359804in}{2.039807in}}%
\pgfpathlineto{\pgfqpoint{1.364367in}{2.053055in}}%
\pgfpathlineto{\pgfqpoint{1.364790in}{2.066303in}}%
\pgfpathlineto{\pgfqpoint{1.369495in}{2.079551in}}%
\pgfpathlineto{\pgfqpoint{1.370835in}{2.092799in}}%
\pgfpathlineto{\pgfqpoint{1.374214in}{2.106047in}}%
\pgfpathlineto{\pgfqpoint{1.375319in}{2.110424in}}%
\pgfpathlineto{\pgfqpoint{1.377748in}{2.119295in}}%
\pgfpathlineto{\pgfqpoint{1.379741in}{2.132543in}}%
\pgfpathlineto{\pgfqpoint{1.375319in}{2.132543in}}%
\pgfpathlineto{\pgfqpoint{1.371367in}{2.132543in}}%
\pgfpathlineto{\pgfqpoint{1.369689in}{2.119295in}}%
\pgfpathlineto{\pgfqpoint{1.366310in}{2.106047in}}%
\pgfpathlineto{\pgfqpoint{1.363187in}{2.092799in}}%
\pgfpathlineto{\pgfqpoint{1.361699in}{2.079551in}}%
\pgfpathlineto{\pgfqpoint{1.357245in}{2.066303in}}%
\pgfpathlineto{\pgfqpoint{1.356731in}{2.053055in}}%
\pgfpathlineto{\pgfqpoint{1.352261in}{2.039807in}}%
\pgfpathlineto{\pgfqpoint{1.351357in}{2.026559in}}%
\pgfpathlineto{\pgfqpoint{1.350036in}{2.021769in}}%
\pgfpathlineto{\pgfqpoint{1.347630in}{2.013311in}}%
\pgfpathlineto{\pgfqpoint{1.345076in}{2.000063in}}%
\pgfpathlineto{\pgfqpoint{1.342878in}{1.986815in}}%
\pgfpathlineto{\pgfqpoint{1.338936in}{1.973567in}}%
\pgfpathlineto{\pgfqpoint{1.338291in}{1.960319in}}%
\pgfpathlineto{\pgfqpoint{1.333710in}{1.947071in}}%
\pgfpathlineto{\pgfqpoint{1.333242in}{1.933823in}}%
\pgfpathlineto{\pgfqpoint{1.329529in}{1.920575in}}%
\pgfpathlineto{\pgfqpoint{1.327983in}{1.907327in}}%
\pgfpathlineto{\pgfqpoint{1.325713in}{1.894079in}}%
\pgfpathlineto{\pgfqpoint{1.324753in}{1.890390in}}%
\pgfpathlineto{\pgfqpoint{1.322152in}{1.880831in}}%
\pgfpathlineto{\pgfqpoint{1.321128in}{1.867583in}}%
\pgfpathlineto{\pgfqpoint{1.316487in}{1.854335in}}%
\pgfpathlineto{\pgfqpoint{1.316212in}{1.841087in}}%
\pgfpathlineto{\pgfqpoint{1.311841in}{1.827839in}}%
\pgfpathlineto{\pgfqpoint{1.310671in}{1.814591in}}%
\pgfpathlineto{\pgfqpoint{1.307873in}{1.801343in}}%
\pgfpathlineto{\pgfqpoint{1.305076in}{1.788095in}}%
\pgfpathlineto{\pgfqpoint{1.304029in}{1.774847in}}%
\pgfpathlineto{\pgfqpoint{1.299657in}{1.761599in}}%
\pgfpathlineto{\pgfqpoint{1.299524in}{1.748351in}}%
\pgfpathlineto{\pgfqpoint{1.299470in}{1.748191in}}%
\pgfpathlineto{\pgfqpoint{1.294677in}{1.735103in}}%
\pgfpathlineto{\pgfqpoint{1.294156in}{1.721855in}}%
\pgfpathlineto{\pgfqpoint{1.290813in}{1.708607in}}%
\pgfpathlineto{\pgfqpoint{1.288335in}{1.695359in}}%
\pgfpathlineto{\pgfqpoint{1.286913in}{1.682111in}}%
\pgfpathlineto{\pgfqpoint{1.282578in}{1.668863in}}%
\pgfpathlineto{\pgfqpoint{1.282482in}{1.655615in}}%
\pgfpathlineto{\pgfqpoint{1.277383in}{1.642367in}}%
\pgfpathlineto{\pgfqpoint{1.277181in}{1.629119in}}%
\pgfpathlineto{\pgfqpoint{1.274187in}{1.618814in}}%
\pgfpathlineto{\pgfqpoint{1.273240in}{1.615871in}}%
\pgfpathlineto{\pgfqpoint{1.271477in}{1.602623in}}%
\pgfpathlineto{\pgfqpoint{1.270006in}{1.589375in}}%
\pgfpathlineto{\pgfqpoint{1.265782in}{1.576127in}}%
\pgfpathlineto{\pgfqpoint{1.266032in}{1.562879in}}%
\pgfpathlineto{\pgfqpoint{1.260579in}{1.549631in}}%
\pgfpathlineto{\pgfqpoint{1.260900in}{1.536383in}}%
\pgfpathlineto{\pgfqpoint{1.256147in}{1.523135in}}%
\pgfpathlineto{\pgfqpoint{1.254659in}{1.509887in}}%
\pgfpathlineto{\pgfqpoint{1.252527in}{1.496639in}}%
\pgfpathlineto{\pgfqpoint{1.248904in}{1.483938in}}%
\pgfpathlineto{\pgfqpoint{1.248738in}{1.483391in}}%
\pgfpathlineto{\pgfqpoint{1.248904in}{1.480430in}}%
\pgfpathlineto{\pgfqpoint{1.249465in}{1.470143in}}%
\pgfpathlineto{\pgfqpoint{1.248904in}{1.468805in}}%
\pgfpathlineto{\pgfqpoint{1.243625in}{1.456895in}}%
\pgfpathlineto{\pgfqpoint{1.244911in}{1.443647in}}%
\pgfpathlineto{\pgfqpoint{1.239662in}{1.430399in}}%
\pgfpathlineto{\pgfqpoint{1.239192in}{1.417151in}}%
\pgfpathlineto{\pgfqpoint{1.236775in}{1.403903in}}%
\pgfpathlineto{\pgfqpoint{1.232896in}{1.390655in}}%
\pgfpathlineto{\pgfqpoint{1.233739in}{1.377407in}}%
\pgfpathlineto{\pgfqpoint{1.227705in}{1.364159in}}%
\pgfpathlineto{\pgfqpoint{1.230180in}{1.350911in}}%
\pgfpathlineto{\pgfqpoint{1.224483in}{1.337663in}}%
\pgfpathlineto{\pgfqpoint{1.224875in}{1.324415in}}%
\pgfpathlineto{\pgfqpoint{1.223620in}{1.317684in}}%
\pgfpathlineto{\pgfqpoint{1.222297in}{1.311167in}}%
\pgfpathlineto{\pgfqpoint{1.219188in}{1.297919in}}%
\pgfpathlineto{\pgfqpoint{1.221276in}{1.284671in}}%
\pgfpathlineto{\pgfqpoint{1.214776in}{1.271423in}}%
\pgfpathlineto{\pgfqpoint{1.219065in}{1.258175in}}%
\pgfpathlineto{\pgfqpoint{1.212993in}{1.244927in}}%
\pgfpathlineto{\pgfqpoint{1.215851in}{1.231679in}}%
\pgfpathlineto{\pgfqpoint{1.214039in}{1.218431in}}%
\pgfpathlineto{\pgfqpoint{1.211260in}{1.205183in}}%
\pgfpathlineto{\pgfqpoint{1.214983in}{1.191935in}}%
\pgfpathlineto{\pgfqpoint{1.208236in}{1.178687in}}%
\pgfpathlineto{\pgfqpoint{1.216217in}{1.165439in}}%
\pgfpathlineto{\pgfqpoint{1.210130in}{1.152191in}}%
\pgfpathlineto{\pgfqpoint{1.216488in}{1.138943in}}%
\pgfpathlineto{\pgfqpoint{1.217301in}{1.125695in}}%
\pgfpathlineto{\pgfqpoint{1.217954in}{1.112447in}}%
\pgfpathlineto{\pgfqpoint{1.223620in}{1.105039in}}%
\pgfpathlineto{\pgfqpoint{1.228790in}{1.099199in}}%
\pgfpathlineto{\pgfqpoint{1.223620in}{1.088799in}}%
\pgfpathlineto{\pgfqpoint{1.222227in}{1.085951in}}%
\pgfpathlineto{\pgfqpoint{1.223620in}{1.084651in}}%
\pgfpathlineto{\pgfqpoint{1.238159in}{1.072703in}}%
\pgfpathlineto{\pgfqpoint{1.233218in}{1.059455in}}%
\pgfpathlineto{\pgfqpoint{1.247759in}{1.046207in}}%
\pgfpathlineto{\pgfqpoint{1.248904in}{1.044444in}}%
\pgfpathlineto{\pgfqpoint{1.258150in}{1.032959in}}%
\pgfpathlineto{\pgfqpoint{1.263903in}{1.019711in}}%
\pgfpathlineto{\pgfqpoint{1.274187in}{1.014968in}}%
\pgfpathlineto{\pgfqpoint{1.297833in}{1.006463in}}%
\pgfpathlineto{\pgfqpoint{1.298870in}{0.993215in}}%
\pgfpathlineto{\pgfqpoint{1.299470in}{0.993048in}}%
\pgfpathlineto{\pgfqpoint{1.324753in}{0.987167in}}%
\pgfpathlineto{\pgfqpoint{1.350036in}{0.982513in}}%
\pgfpathlineto{\pgfqpoint{1.367898in}{0.979967in}}%
\pgfpathlineto{\pgfqpoint{1.375319in}{0.976528in}}%
\pgfpathclose%
\pgfusepath{fill}%
\end{pgfscope}%
\begin{pgfscope}%
\pgfpathrectangle{\pgfqpoint{0.540975in}{0.476543in}}{\pgfqpoint{1.820388in}{1.656000in}}%
\pgfusepath{clip}%
\pgfsetbuttcap%
\pgfsetroundjoin%
\definecolor{currentfill}{rgb}{0.926105,0.971396,0.884890}%
\pgfsetfillcolor{currentfill}%
\pgfsetlinewidth{0.000000pt}%
\definecolor{currentstroke}{rgb}{0.000000,0.000000,0.000000}%
\pgfsetstrokecolor{currentstroke}%
\pgfsetdash{}{0pt}%
\pgfpathmoveto{\pgfqpoint{1.425886in}{0.997847in}}%
\pgfpathlineto{\pgfqpoint{1.451169in}{0.996300in}}%
\pgfpathlineto{\pgfqpoint{1.476452in}{1.000641in}}%
\pgfpathlineto{\pgfqpoint{1.492315in}{1.006463in}}%
\pgfpathlineto{\pgfqpoint{1.501735in}{1.007156in}}%
\pgfpathlineto{\pgfqpoint{1.527018in}{1.010283in}}%
\pgfpathlineto{\pgfqpoint{1.552302in}{1.014947in}}%
\pgfpathlineto{\pgfqpoint{1.570952in}{1.019711in}}%
\pgfpathlineto{\pgfqpoint{1.571072in}{1.032959in}}%
\pgfpathlineto{\pgfqpoint{1.577585in}{1.034776in}}%
\pgfpathlineto{\pgfqpoint{1.602868in}{1.043551in}}%
\pgfpathlineto{\pgfqpoint{1.609166in}{1.046207in}}%
\pgfpathlineto{\pgfqpoint{1.610303in}{1.059455in}}%
\pgfpathlineto{\pgfqpoint{1.628151in}{1.071782in}}%
\pgfpathlineto{\pgfqpoint{1.629326in}{1.072703in}}%
\pgfpathlineto{\pgfqpoint{1.636160in}{1.085951in}}%
\pgfpathlineto{\pgfqpoint{1.640744in}{1.099199in}}%
\pgfpathlineto{\pgfqpoint{1.652919in}{1.112447in}}%
\pgfpathlineto{\pgfqpoint{1.649659in}{1.125695in}}%
\pgfpathlineto{\pgfqpoint{1.653434in}{1.129456in}}%
\pgfpathlineto{\pgfqpoint{1.661706in}{1.138943in}}%
\pgfpathlineto{\pgfqpoint{1.655863in}{1.152191in}}%
\pgfpathlineto{\pgfqpoint{1.664934in}{1.165439in}}%
\pgfpathlineto{\pgfqpoint{1.662807in}{1.178687in}}%
\pgfpathlineto{\pgfqpoint{1.665970in}{1.191935in}}%
\pgfpathlineto{\pgfqpoint{1.668851in}{1.205183in}}%
\pgfpathlineto{\pgfqpoint{1.665840in}{1.218431in}}%
\pgfpathlineto{\pgfqpoint{1.671736in}{1.231679in}}%
\pgfpathlineto{\pgfqpoint{1.666023in}{1.244927in}}%
\pgfpathlineto{\pgfqpoint{1.670658in}{1.258175in}}%
\pgfpathlineto{\pgfqpoint{1.665519in}{1.271423in}}%
\pgfpathlineto{\pgfqpoint{1.667006in}{1.284671in}}%
\pgfpathlineto{\pgfqpoint{1.665612in}{1.297919in}}%
\pgfpathlineto{\pgfqpoint{1.663147in}{1.311167in}}%
\pgfpathlineto{\pgfqpoint{1.665349in}{1.324415in}}%
\pgfpathlineto{\pgfqpoint{1.659876in}{1.337663in}}%
\pgfpathlineto{\pgfqpoint{1.663211in}{1.350911in}}%
\pgfpathlineto{\pgfqpoint{1.658005in}{1.364159in}}%
\pgfpathlineto{\pgfqpoint{1.659027in}{1.377407in}}%
\pgfpathlineto{\pgfqpoint{1.656128in}{1.390655in}}%
\pgfpathlineto{\pgfqpoint{1.653745in}{1.403903in}}%
\pgfpathlineto{\pgfqpoint{1.653927in}{1.417151in}}%
\pgfpathlineto{\pgfqpoint{1.653434in}{1.418487in}}%
\pgfpathlineto{\pgfqpoint{1.649184in}{1.430399in}}%
\pgfpathlineto{\pgfqpoint{1.650767in}{1.443647in}}%
\pgfpathlineto{\pgfqpoint{1.645282in}{1.456895in}}%
\pgfpathlineto{\pgfqpoint{1.645963in}{1.470143in}}%
\pgfpathlineto{\pgfqpoint{1.642434in}{1.483391in}}%
\pgfpathlineto{\pgfqpoint{1.640503in}{1.496639in}}%
\pgfpathlineto{\pgfqpoint{1.639656in}{1.509887in}}%
\pgfpathlineto{\pgfqpoint{1.635252in}{1.523135in}}%
\pgfpathlineto{\pgfqpoint{1.636309in}{1.536383in}}%
\pgfpathlineto{\pgfqpoint{1.631309in}{1.549631in}}%
\pgfpathlineto{\pgfqpoint{1.631774in}{1.562879in}}%
\pgfpathlineto{\pgfqpoint{1.628151in}{1.574626in}}%
\pgfpathlineto{\pgfqpoint{1.627721in}{1.576127in}}%
\pgfpathlineto{\pgfqpoint{1.625857in}{1.589375in}}%
\pgfpathlineto{\pgfqpoint{1.624185in}{1.602623in}}%
\pgfpathlineto{\pgfqpoint{1.620226in}{1.615871in}}%
\pgfpathlineto{\pgfqpoint{1.620285in}{1.629119in}}%
\pgfpathlineto{\pgfqpoint{1.615396in}{1.642367in}}%
\pgfpathlineto{\pgfqpoint{1.615783in}{1.655615in}}%
\pgfpathlineto{\pgfqpoint{1.612009in}{1.668863in}}%
\pgfpathlineto{\pgfqpoint{1.610700in}{1.682111in}}%
\pgfpathlineto{\pgfqpoint{1.608785in}{1.695359in}}%
\pgfpathlineto{\pgfqpoint{1.604910in}{1.708607in}}%
\pgfpathlineto{\pgfqpoint{1.604546in}{1.721855in}}%
\pgfpathlineto{\pgfqpoint{1.602868in}{1.726415in}}%
\pgfpathlineto{\pgfqpoint{1.599832in}{1.735103in}}%
\pgfpathlineto{\pgfqpoint{1.599622in}{1.748351in}}%
\pgfpathlineto{\pgfqpoint{1.595536in}{1.761599in}}%
\pgfpathlineto{\pgfqpoint{1.594184in}{1.774847in}}%
\pgfpathlineto{\pgfqpoint{1.592031in}{1.788095in}}%
\pgfpathlineto{\pgfqpoint{1.588875in}{1.801343in}}%
\pgfpathlineto{\pgfqpoint{1.588363in}{1.814591in}}%
\pgfpathlineto{\pgfqpoint{1.583761in}{1.827839in}}%
\pgfpathlineto{\pgfqpoint{1.583561in}{1.841087in}}%
\pgfpathlineto{\pgfqpoint{1.579400in}{1.854335in}}%
\pgfpathlineto{\pgfqpoint{1.577964in}{1.867583in}}%
\pgfpathlineto{\pgfqpoint{1.577585in}{1.869444in}}%
\pgfpathlineto{\pgfqpoint{1.575414in}{1.880831in}}%
\pgfpathlineto{\pgfqpoint{1.572163in}{1.894079in}}%
\pgfpathlineto{\pgfqpoint{1.571336in}{1.907327in}}%
\pgfpathlineto{\pgfqpoint{1.566969in}{1.920575in}}%
\pgfpathlineto{\pgfqpoint{1.566569in}{1.933823in}}%
\pgfpathlineto{\pgfqpoint{1.562323in}{1.947071in}}%
\pgfpathlineto{\pgfqpoint{1.560868in}{1.960319in}}%
\pgfpathlineto{\pgfqpoint{1.558204in}{1.973567in}}%
\pgfpathlineto{\pgfqpoint{1.555098in}{1.986815in}}%
\pgfpathlineto{\pgfqpoint{1.554070in}{2.000063in}}%
\pgfpathlineto{\pgfqpoint{1.552302in}{2.005209in}}%
\pgfpathlineto{\pgfqpoint{1.549441in}{2.013311in}}%
\pgfpathlineto{\pgfqpoint{1.549006in}{2.026559in}}%
\pgfpathlineto{\pgfqpoint{1.544492in}{2.039807in}}%
\pgfpathlineto{\pgfqpoint{1.542867in}{2.053055in}}%
\pgfpathlineto{\pgfqpoint{1.539759in}{2.066303in}}%
\pgfpathlineto{\pgfqpoint{1.536103in}{2.079551in}}%
\pgfpathlineto{\pgfqpoint{1.534803in}{2.092799in}}%
\pgfpathlineto{\pgfqpoint{1.530107in}{2.106047in}}%
\pgfpathlineto{\pgfqpoint{1.529530in}{2.119295in}}%
\pgfpathlineto{\pgfqpoint{1.527018in}{2.126464in}}%
\pgfpathlineto{\pgfqpoint{1.524611in}{2.132543in}}%
\pgfpathlineto{\pgfqpoint{1.515693in}{2.132543in}}%
\pgfpathlineto{\pgfqpoint{1.520922in}{2.119295in}}%
\pgfpathlineto{\pgfqpoint{1.521535in}{2.106047in}}%
\pgfpathlineto{\pgfqpoint{1.527018in}{2.093026in}}%
\pgfpathlineto{\pgfqpoint{1.527100in}{2.092799in}}%
\pgfpathlineto{\pgfqpoint{1.528250in}{2.079551in}}%
\pgfpathlineto{\pgfqpoint{1.532155in}{2.066303in}}%
\pgfpathlineto{\pgfqpoint{1.535148in}{2.053055in}}%
\pgfpathlineto{\pgfqpoint{1.536885in}{2.039807in}}%
\pgfpathlineto{\pgfqpoint{1.541422in}{2.026559in}}%
\pgfpathlineto{\pgfqpoint{1.541794in}{2.013311in}}%
\pgfpathlineto{\pgfqpoint{1.546593in}{2.000063in}}%
\pgfpathlineto{\pgfqpoint{1.547472in}{1.986815in}}%
\pgfpathlineto{\pgfqpoint{1.550827in}{1.973567in}}%
\pgfpathlineto{\pgfqpoint{1.552302in}{1.966082in}}%
\pgfpathlineto{\pgfqpoint{1.553433in}{1.960319in}}%
\pgfpathlineto{\pgfqpoint{1.554973in}{1.947071in}}%
\pgfpathlineto{\pgfqpoint{1.559249in}{1.933823in}}%
\pgfpathlineto{\pgfqpoint{1.559559in}{1.920575in}}%
\pgfpathlineto{\pgfqpoint{1.564070in}{1.907327in}}%
\pgfpathlineto{\pgfqpoint{1.564706in}{1.894079in}}%
\pgfpathlineto{\pgfqpoint{1.568134in}{1.880831in}}%
\pgfpathlineto{\pgfqpoint{1.570508in}{1.867583in}}%
\pgfpathlineto{\pgfqpoint{1.571962in}{1.854335in}}%
\pgfpathlineto{\pgfqpoint{1.575973in}{1.841087in}}%
\pgfpathlineto{\pgfqpoint{1.576079in}{1.827839in}}%
\pgfpathlineto{\pgfqpoint{1.577585in}{1.823568in}}%
\pgfpathlineto{\pgfqpoint{1.580786in}{1.814591in}}%
\pgfpathlineto{\pgfqpoint{1.581097in}{1.801343in}}%
\pgfpathlineto{\pgfqpoint{1.584400in}{1.788095in}}%
\pgfpathlineto{\pgfqpoint{1.586363in}{1.774847in}}%
\pgfpathlineto{\pgfqpoint{1.587774in}{1.761599in}}%
\pgfpathlineto{\pgfqpoint{1.591797in}{1.748351in}}%
\pgfpathlineto{\pgfqpoint{1.591910in}{1.735103in}}%
\pgfpathlineto{\pgfqpoint{1.596602in}{1.721855in}}%
\pgfpathlineto{\pgfqpoint{1.596712in}{1.708607in}}%
\pgfpathlineto{\pgfqpoint{1.600533in}{1.695359in}}%
\pgfpathlineto{\pgfqpoint{1.602150in}{1.682111in}}%
\pgfpathlineto{\pgfqpoint{1.602868in}{1.674765in}}%
\pgfpathlineto{\pgfqpoint{1.603466in}{1.668863in}}%
\pgfpathlineto{\pgfqpoint{1.607152in}{1.655615in}}%
\pgfpathlineto{\pgfqpoint{1.606618in}{1.642367in}}%
\pgfpathlineto{\pgfqpoint{1.611558in}{1.629119in}}%
\pgfpathlineto{\pgfqpoint{1.611213in}{1.615871in}}%
\pgfpathlineto{\pgfqpoint{1.615294in}{1.602623in}}%
\pgfpathlineto{\pgfqpoint{1.616629in}{1.589375in}}%
\pgfpathlineto{\pgfqpoint{1.618542in}{1.576127in}}%
\pgfpathlineto{\pgfqpoint{1.622179in}{1.562879in}}%
\pgfpathlineto{\pgfqpoint{1.621579in}{1.549631in}}%
\pgfpathlineto{\pgfqpoint{1.626366in}{1.536383in}}%
\pgfpathlineto{\pgfqpoint{1.624983in}{1.523135in}}%
\pgfpathlineto{\pgfqpoint{1.628151in}{1.513550in}}%
\pgfpathlineto{\pgfqpoint{1.629391in}{1.509887in}}%
\pgfpathlineto{\pgfqpoint{1.629816in}{1.496639in}}%
\pgfpathlineto{\pgfqpoint{1.631766in}{1.483391in}}%
\pgfpathlineto{\pgfqpoint{1.634949in}{1.470143in}}%
\pgfpathlineto{\pgfqpoint{1.634059in}{1.456895in}}%
\pgfpathlineto{\pgfqpoint{1.639364in}{1.443647in}}%
\pgfpathlineto{\pgfqpoint{1.637305in}{1.430399in}}%
\pgfpathlineto{\pgfqpoint{1.642048in}{1.417151in}}%
\pgfpathlineto{\pgfqpoint{1.641225in}{1.403903in}}%
\pgfpathlineto{\pgfqpoint{1.643577in}{1.390655in}}%
\pgfpathlineto{\pgfqpoint{1.645839in}{1.377407in}}%
\pgfpathlineto{\pgfqpoint{1.644590in}{1.364159in}}%
\pgfpathlineto{\pgfqpoint{1.649376in}{1.350911in}}%
\pgfpathlineto{\pgfqpoint{1.645405in}{1.337663in}}%
\pgfpathlineto{\pgfqpoint{1.650702in}{1.324415in}}%
\pgfpathlineto{\pgfqpoint{1.647524in}{1.311167in}}%
\pgfpathlineto{\pgfqpoint{1.649939in}{1.297919in}}%
\pgfpathlineto{\pgfqpoint{1.650327in}{1.284671in}}%
\pgfpathlineto{\pgfqpoint{1.648542in}{1.271423in}}%
\pgfpathlineto{\pgfqpoint{1.653003in}{1.258175in}}%
\pgfpathlineto{\pgfqpoint{1.647397in}{1.244927in}}%
\pgfpathlineto{\pgfqpoint{1.652842in}{1.231679in}}%
\pgfpathlineto{\pgfqpoint{1.645220in}{1.218431in}}%
\pgfpathlineto{\pgfqpoint{1.648111in}{1.205183in}}%
\pgfpathlineto{\pgfqpoint{1.643182in}{1.191935in}}%
\pgfpathlineto{\pgfqpoint{1.639447in}{1.178687in}}%
\pgfpathlineto{\pgfqpoint{1.639787in}{1.165439in}}%
\pgfpathlineto{\pgfqpoint{1.629080in}{1.152191in}}%
\pgfpathlineto{\pgfqpoint{1.633779in}{1.138943in}}%
\pgfpathlineto{\pgfqpoint{1.628151in}{1.133723in}}%
\pgfpathlineto{\pgfqpoint{1.617915in}{1.125695in}}%
\pgfpathlineto{\pgfqpoint{1.620558in}{1.112447in}}%
\pgfpathlineto{\pgfqpoint{1.602908in}{1.099199in}}%
\pgfpathlineto{\pgfqpoint{1.602868in}{1.099100in}}%
\pgfpathlineto{\pgfqpoint{1.596882in}{1.085951in}}%
\pgfpathlineto{\pgfqpoint{1.582518in}{1.072703in}}%
\pgfpathlineto{\pgfqpoint{1.577585in}{1.069795in}}%
\pgfpathlineto{\pgfqpoint{1.556773in}{1.059455in}}%
\pgfpathlineto{\pgfqpoint{1.552302in}{1.052012in}}%
\pgfpathlineto{\pgfqpoint{1.547358in}{1.046207in}}%
\pgfpathlineto{\pgfqpoint{1.527018in}{1.041646in}}%
\pgfpathlineto{\pgfqpoint{1.501735in}{1.037553in}}%
\pgfpathlineto{\pgfqpoint{1.476452in}{1.034927in}}%
\pgfpathlineto{\pgfqpoint{1.451169in}{1.033748in}}%
\pgfpathlineto{\pgfqpoint{1.425886in}{1.034119in}}%
\pgfpathlineto{\pgfqpoint{1.400603in}{1.036293in}}%
\pgfpathlineto{\pgfqpoint{1.375319in}{1.040750in}}%
\pgfpathlineto{\pgfqpoint{1.356557in}{1.046207in}}%
\pgfpathlineto{\pgfqpoint{1.350036in}{1.048908in}}%
\pgfpathlineto{\pgfqpoint{1.329459in}{1.059455in}}%
\pgfpathlineto{\pgfqpoint{1.324753in}{1.070203in}}%
\pgfpathlineto{\pgfqpoint{1.323955in}{1.072703in}}%
\pgfpathlineto{\pgfqpoint{1.299470in}{1.084045in}}%
\pgfpathlineto{\pgfqpoint{1.296073in}{1.085951in}}%
\pgfpathlineto{\pgfqpoint{1.299470in}{1.095617in}}%
\pgfpathlineto{\pgfqpoint{1.300825in}{1.099199in}}%
\pgfpathlineto{\pgfqpoint{1.299470in}{1.099883in}}%
\pgfpathlineto{\pgfqpoint{1.279749in}{1.112447in}}%
\pgfpathlineto{\pgfqpoint{1.278073in}{1.125695in}}%
\pgfpathlineto{\pgfqpoint{1.274187in}{1.132037in}}%
\pgfpathlineto{\pgfqpoint{1.270719in}{1.138943in}}%
\pgfpathlineto{\pgfqpoint{1.262246in}{1.152191in}}%
\pgfpathlineto{\pgfqpoint{1.265516in}{1.165439in}}%
\pgfpathlineto{\pgfqpoint{1.253811in}{1.178687in}}%
\pgfpathlineto{\pgfqpoint{1.259952in}{1.191935in}}%
\pgfpathlineto{\pgfqpoint{1.252195in}{1.205183in}}%
\pgfpathlineto{\pgfqpoint{1.254772in}{1.218431in}}%
\pgfpathlineto{\pgfqpoint{1.253575in}{1.231679in}}%
\pgfpathlineto{\pgfqpoint{1.249798in}{1.244927in}}%
\pgfpathlineto{\pgfqpoint{1.254402in}{1.258175in}}%
\pgfpathlineto{\pgfqpoint{1.248904in}{1.270233in}}%
\pgfpathlineto{\pgfqpoint{1.248396in}{1.271423in}}%
\pgfpathlineto{\pgfqpoint{1.248904in}{1.272512in}}%
\pgfpathlineto{\pgfqpoint{1.254590in}{1.284671in}}%
\pgfpathlineto{\pgfqpoint{1.250326in}{1.297919in}}%
\pgfpathlineto{\pgfqpoint{1.253413in}{1.311167in}}%
\pgfpathlineto{\pgfqpoint{1.254070in}{1.324415in}}%
\pgfpathlineto{\pgfqpoint{1.253214in}{1.337663in}}%
\pgfpathlineto{\pgfqpoint{1.257904in}{1.350911in}}%
\pgfpathlineto{\pgfqpoint{1.254406in}{1.364159in}}%
\pgfpathlineto{\pgfqpoint{1.260239in}{1.377407in}}%
\pgfpathlineto{\pgfqpoint{1.258081in}{1.390655in}}%
\pgfpathlineto{\pgfqpoint{1.262062in}{1.403903in}}%
\pgfpathlineto{\pgfqpoint{1.263245in}{1.417151in}}%
\pgfpathlineto{\pgfqpoint{1.263607in}{1.430399in}}%
\pgfpathlineto{\pgfqpoint{1.267992in}{1.443647in}}%
\pgfpathlineto{\pgfqpoint{1.266100in}{1.456895in}}%
\pgfpathlineto{\pgfqpoint{1.271582in}{1.470143in}}%
\pgfpathlineto{\pgfqpoint{1.269876in}{1.483391in}}%
\pgfpathlineto{\pgfqpoint{1.273887in}{1.496639in}}%
\pgfpathlineto{\pgfqpoint{1.274187in}{1.499625in}}%
\pgfpathlineto{\pgfqpoint{1.275145in}{1.509887in}}%
\pgfpathlineto{\pgfqpoint{1.276665in}{1.523135in}}%
\pgfpathlineto{\pgfqpoint{1.280848in}{1.536383in}}%
\pgfpathlineto{\pgfqpoint{1.280224in}{1.549631in}}%
\pgfpathlineto{\pgfqpoint{1.285454in}{1.562879in}}%
\pgfpathlineto{\pgfqpoint{1.284553in}{1.576127in}}%
\pgfpathlineto{\pgfqpoint{1.288762in}{1.589375in}}%
\pgfpathlineto{\pgfqpoint{1.289441in}{1.602623in}}%
\pgfpathlineto{\pgfqpoint{1.291239in}{1.615871in}}%
\pgfpathlineto{\pgfqpoint{1.294722in}{1.629119in}}%
\pgfpathlineto{\pgfqpoint{1.294770in}{1.642367in}}%
\pgfpathlineto{\pgfqpoint{1.299470in}{1.654412in}}%
\pgfpathlineto{\pgfqpoint{1.299923in}{1.655615in}}%
\pgfpathlineto{\pgfqpoint{1.299572in}{1.668863in}}%
\pgfpathlineto{\pgfqpoint{1.304033in}{1.682111in}}%
\pgfpathlineto{\pgfqpoint{1.304864in}{1.695359in}}%
\pgfpathlineto{\pgfqpoint{1.307462in}{1.708607in}}%
\pgfpathlineto{\pgfqpoint{1.310274in}{1.721855in}}%
\pgfpathlineto{\pgfqpoint{1.310776in}{1.735103in}}%
\pgfpathlineto{\pgfqpoint{1.315261in}{1.748351in}}%
\pgfpathlineto{\pgfqpoint{1.315099in}{1.761599in}}%
\pgfpathlineto{\pgfqpoint{1.319752in}{1.774847in}}%
\pgfpathlineto{\pgfqpoint{1.320342in}{1.788095in}}%
\pgfpathlineto{\pgfqpoint{1.323460in}{1.801343in}}%
\pgfpathlineto{\pgfqpoint{1.324753in}{1.808109in}}%
\pgfpathlineto{\pgfqpoint{1.325925in}{1.814591in}}%
\pgfpathlineto{\pgfqpoint{1.327175in}{1.827839in}}%
\pgfpathlineto{\pgfqpoint{1.331395in}{1.841087in}}%
\pgfpathlineto{\pgfqpoint{1.331519in}{1.854335in}}%
\pgfpathlineto{\pgfqpoint{1.336215in}{1.867583in}}%
\pgfpathlineto{\pgfqpoint{1.336841in}{1.880831in}}%
\pgfpathlineto{\pgfqpoint{1.340603in}{1.894079in}}%
\pgfpathlineto{\pgfqpoint{1.342520in}{1.907327in}}%
\pgfpathlineto{\pgfqpoint{1.344295in}{1.920575in}}%
\pgfpathlineto{\pgfqpoint{1.347941in}{1.933823in}}%
\pgfpathlineto{\pgfqpoint{1.348353in}{1.947071in}}%
\pgfpathlineto{\pgfqpoint{1.350036in}{1.951692in}}%
\pgfpathlineto{\pgfqpoint{1.353266in}{1.960319in}}%
\pgfpathlineto{\pgfqpoint{1.353653in}{1.973567in}}%
\pgfpathlineto{\pgfqpoint{1.358059in}{1.986815in}}%
\pgfpathlineto{\pgfqpoint{1.359925in}{2.000063in}}%
\pgfpathlineto{\pgfqpoint{1.362830in}{2.013311in}}%
\pgfpathlineto{\pgfqpoint{1.366312in}{2.026559in}}%
\pgfpathlineto{\pgfqpoint{1.367346in}{2.039807in}}%
\pgfpathlineto{\pgfqpoint{1.372003in}{2.053055in}}%
\pgfpathlineto{\pgfqpoint{1.372335in}{2.066303in}}%
\pgfpathlineto{\pgfqpoint{1.375319in}{2.074385in}}%
\pgfpathlineto{\pgfqpoint{1.377528in}{2.079551in}}%
\pgfpathlineto{\pgfqpoint{1.378837in}{2.092799in}}%
\pgfpathlineto{\pgfqpoint{1.382882in}{2.106047in}}%
\pgfpathlineto{\pgfqpoint{1.386316in}{2.119295in}}%
\pgfpathlineto{\pgfqpoint{1.388499in}{2.132543in}}%
\pgfpathlineto{\pgfqpoint{1.379741in}{2.132543in}}%
\pgfpathlineto{\pgfqpoint{1.377748in}{2.119295in}}%
\pgfpathlineto{\pgfqpoint{1.375319in}{2.110424in}}%
\pgfpathlineto{\pgfqpoint{1.374214in}{2.106047in}}%
\pgfpathlineto{\pgfqpoint{1.370835in}{2.092799in}}%
\pgfpathlineto{\pgfqpoint{1.369495in}{2.079551in}}%
\pgfpathlineto{\pgfqpoint{1.364790in}{2.066303in}}%
\pgfpathlineto{\pgfqpoint{1.364367in}{2.053055in}}%
\pgfpathlineto{\pgfqpoint{1.359804in}{2.039807in}}%
\pgfpathlineto{\pgfqpoint{1.358835in}{2.026559in}}%
\pgfpathlineto{\pgfqpoint{1.355251in}{2.013311in}}%
\pgfpathlineto{\pgfqpoint{1.352543in}{2.000063in}}%
\pgfpathlineto{\pgfqpoint{1.350484in}{1.986815in}}%
\pgfpathlineto{\pgfqpoint{1.350036in}{1.985437in}}%
\pgfpathlineto{\pgfqpoint{1.346288in}{1.973567in}}%
\pgfpathlineto{\pgfqpoint{1.345768in}{1.960319in}}%
\pgfpathlineto{\pgfqpoint{1.341032in}{1.947071in}}%
\pgfpathlineto{\pgfqpoint{1.340591in}{1.933823in}}%
\pgfpathlineto{\pgfqpoint{1.336912in}{1.920575in}}%
\pgfpathlineto{\pgfqpoint{1.335251in}{1.907327in}}%
\pgfpathlineto{\pgfqpoint{1.333158in}{1.894079in}}%
\pgfpathlineto{\pgfqpoint{1.329585in}{1.880831in}}%
\pgfpathlineto{\pgfqpoint{1.328785in}{1.867583in}}%
\pgfpathlineto{\pgfqpoint{1.324753in}{1.855937in}}%
\pgfpathlineto{\pgfqpoint{1.324179in}{1.854335in}}%
\pgfpathlineto{\pgfqpoint{1.323980in}{1.841087in}}%
\pgfpathlineto{\pgfqpoint{1.319551in}{1.827839in}}%
\pgfpathlineto{\pgfqpoint{1.318320in}{1.814591in}}%
\pgfpathlineto{\pgfqpoint{1.315667in}{1.801343in}}%
\pgfpathlineto{\pgfqpoint{1.312709in}{1.788095in}}%
\pgfpathlineto{\pgfqpoint{1.311891in}{1.774847in}}%
\pgfpathlineto{\pgfqpoint{1.307378in}{1.761599in}}%
\pgfpathlineto{\pgfqpoint{1.307393in}{1.748351in}}%
\pgfpathlineto{\pgfqpoint{1.302900in}{1.735103in}}%
\pgfpathlineto{\pgfqpoint{1.302411in}{1.721855in}}%
\pgfpathlineto{\pgfqpoint{1.299470in}{1.708913in}}%
\pgfpathlineto{\pgfqpoint{1.299394in}{1.708607in}}%
\pgfpathlineto{\pgfqpoint{1.296766in}{1.695359in}}%
\pgfpathlineto{\pgfqpoint{1.295594in}{1.682111in}}%
\pgfpathlineto{\pgfqpoint{1.291077in}{1.668863in}}%
\pgfpathlineto{\pgfqpoint{1.291212in}{1.655615in}}%
\pgfpathlineto{\pgfqpoint{1.286077in}{1.642367in}}%
\pgfpathlineto{\pgfqpoint{1.285951in}{1.629119in}}%
\pgfpathlineto{\pgfqpoint{1.282276in}{1.615871in}}%
\pgfpathlineto{\pgfqpoint{1.280562in}{1.602623in}}%
\pgfpathlineto{\pgfqpoint{1.279524in}{1.589375in}}%
\pgfpathlineto{\pgfqpoint{1.275419in}{1.576127in}}%
\pgfpathlineto{\pgfqpoint{1.275978in}{1.562879in}}%
\pgfpathlineto{\pgfqpoint{1.274187in}{1.558364in}}%
\pgfpathlineto{\pgfqpoint{1.270543in}{1.549631in}}%
\pgfpathlineto{\pgfqpoint{1.271034in}{1.536383in}}%
\pgfpathlineto{\pgfqpoint{1.266447in}{1.523135in}}%
\pgfpathlineto{\pgfqpoint{1.264918in}{1.509887in}}%
\pgfpathlineto{\pgfqpoint{1.263207in}{1.496639in}}%
\pgfpathlineto{\pgfqpoint{1.259313in}{1.483391in}}%
\pgfpathlineto{\pgfqpoint{1.260523in}{1.470143in}}%
\pgfpathlineto{\pgfqpoint{1.254994in}{1.456895in}}%
\pgfpathlineto{\pgfqpoint{1.256553in}{1.443647in}}%
\pgfpathlineto{\pgfqpoint{1.251807in}{1.430399in}}%
\pgfpathlineto{\pgfqpoint{1.251402in}{1.417151in}}%
\pgfpathlineto{\pgfqpoint{1.249588in}{1.403903in}}%
\pgfpathlineto{\pgfqpoint{1.248904in}{1.401595in}}%
\pgfpathlineto{\pgfqpoint{1.245599in}{1.390655in}}%
\pgfpathlineto{\pgfqpoint{1.247109in}{1.377407in}}%
\pgfpathlineto{\pgfqpoint{1.241085in}{1.364159in}}%
\pgfpathlineto{\pgfqpoint{1.244097in}{1.350911in}}%
\pgfpathlineto{\pgfqpoint{1.238847in}{1.337663in}}%
\pgfpathlineto{\pgfqpoint{1.239474in}{1.324415in}}%
\pgfpathlineto{\pgfqpoint{1.237851in}{1.311167in}}%
\pgfpathlineto{\pgfqpoint{1.234794in}{1.297919in}}%
\pgfpathlineto{\pgfqpoint{1.237911in}{1.284671in}}%
\pgfpathlineto{\pgfqpoint{1.231621in}{1.271423in}}%
\pgfpathlineto{\pgfqpoint{1.236700in}{1.258175in}}%
\pgfpathlineto{\pgfqpoint{1.231370in}{1.244927in}}%
\pgfpathlineto{\pgfqpoint{1.234649in}{1.231679in}}%
\pgfpathlineto{\pgfqpoint{1.234251in}{1.218431in}}%
\pgfpathlineto{\pgfqpoint{1.231589in}{1.205183in}}%
\pgfpathlineto{\pgfqpoint{1.237168in}{1.191935in}}%
\pgfpathlineto{\pgfqpoint{1.230710in}{1.178687in}}%
\pgfpathlineto{\pgfqpoint{1.240396in}{1.165439in}}%
\pgfpathlineto{\pgfqpoint{1.235574in}{1.152191in}}%
\pgfpathlineto{\pgfqpoint{1.242922in}{1.138943in}}%
\pgfpathlineto{\pgfqpoint{1.246519in}{1.125695in}}%
\pgfpathlineto{\pgfqpoint{1.247590in}{1.112447in}}%
\pgfpathlineto{\pgfqpoint{1.248904in}{1.111120in}}%
\pgfpathlineto{\pgfqpoint{1.262917in}{1.099199in}}%
\pgfpathlineto{\pgfqpoint{1.257098in}{1.085951in}}%
\pgfpathlineto{\pgfqpoint{1.274187in}{1.074558in}}%
\pgfpathlineto{\pgfqpoint{1.277474in}{1.072703in}}%
\pgfpathlineto{\pgfqpoint{1.275970in}{1.059455in}}%
\pgfpathlineto{\pgfqpoint{1.294639in}{1.046207in}}%
\pgfpathlineto{\pgfqpoint{1.299470in}{1.042551in}}%
\pgfpathlineto{\pgfqpoint{1.316528in}{1.032959in}}%
\pgfpathlineto{\pgfqpoint{1.324669in}{1.019711in}}%
\pgfpathlineto{\pgfqpoint{1.324753in}{1.019686in}}%
\pgfpathlineto{\pgfqpoint{1.350036in}{1.013695in}}%
\pgfpathlineto{\pgfqpoint{1.375319in}{1.009463in}}%
\pgfpathlineto{\pgfqpoint{1.400603in}{1.006661in}}%
\pgfpathlineto{\pgfqpoint{1.403746in}{1.006463in}}%
\pgfpathlineto{\pgfqpoint{1.425886in}{0.997847in}}%
\pgfpathclose%
\pgfusepath{fill}%
\end{pgfscope}%
\begin{pgfscope}%
\pgfpathrectangle{\pgfqpoint{0.540975in}{0.476543in}}{\pgfqpoint{1.820388in}{1.656000in}}%
\pgfusepath{clip}%
\pgfsetbuttcap%
\pgfsetroundjoin%
\definecolor{currentfill}{rgb}{0.973779,0.989850,0.797232}%
\pgfsetfillcolor{currentfill}%
\pgfsetlinewidth{0.000000pt}%
\definecolor{currentstroke}{rgb}{0.000000,0.000000,0.000000}%
\pgfsetstrokecolor{currentstroke}%
\pgfsetdash{}{0pt}%
\pgfpathmoveto{\pgfqpoint{1.375319in}{1.040750in}}%
\pgfpathlineto{\pgfqpoint{1.400603in}{1.036293in}}%
\pgfpathlineto{\pgfqpoint{1.425886in}{1.034119in}}%
\pgfpathlineto{\pgfqpoint{1.451169in}{1.033748in}}%
\pgfpathlineto{\pgfqpoint{1.476452in}{1.034927in}}%
\pgfpathlineto{\pgfqpoint{1.501735in}{1.037553in}}%
\pgfpathlineto{\pgfqpoint{1.527018in}{1.041646in}}%
\pgfpathlineto{\pgfqpoint{1.547358in}{1.046207in}}%
\pgfpathlineto{\pgfqpoint{1.552302in}{1.052012in}}%
\pgfpathlineto{\pgfqpoint{1.556773in}{1.059455in}}%
\pgfpathlineto{\pgfqpoint{1.577585in}{1.069795in}}%
\pgfpathlineto{\pgfqpoint{1.582518in}{1.072703in}}%
\pgfpathlineto{\pgfqpoint{1.596882in}{1.085951in}}%
\pgfpathlineto{\pgfqpoint{1.602868in}{1.099100in}}%
\pgfpathlineto{\pgfqpoint{1.602908in}{1.099199in}}%
\pgfpathlineto{\pgfqpoint{1.620558in}{1.112447in}}%
\pgfpathlineto{\pgfqpoint{1.617915in}{1.125695in}}%
\pgfpathlineto{\pgfqpoint{1.628151in}{1.133723in}}%
\pgfpathlineto{\pgfqpoint{1.633779in}{1.138943in}}%
\pgfpathlineto{\pgfqpoint{1.629080in}{1.152191in}}%
\pgfpathlineto{\pgfqpoint{1.639787in}{1.165439in}}%
\pgfpathlineto{\pgfqpoint{1.639447in}{1.178687in}}%
\pgfpathlineto{\pgfqpoint{1.643182in}{1.191935in}}%
\pgfpathlineto{\pgfqpoint{1.648111in}{1.205183in}}%
\pgfpathlineto{\pgfqpoint{1.645220in}{1.218431in}}%
\pgfpathlineto{\pgfqpoint{1.652842in}{1.231679in}}%
\pgfpathlineto{\pgfqpoint{1.647397in}{1.244927in}}%
\pgfpathlineto{\pgfqpoint{1.653003in}{1.258175in}}%
\pgfpathlineto{\pgfqpoint{1.648542in}{1.271423in}}%
\pgfpathlineto{\pgfqpoint{1.650327in}{1.284671in}}%
\pgfpathlineto{\pgfqpoint{1.649939in}{1.297919in}}%
\pgfpathlineto{\pgfqpoint{1.647524in}{1.311167in}}%
\pgfpathlineto{\pgfqpoint{1.650702in}{1.324415in}}%
\pgfpathlineto{\pgfqpoint{1.645405in}{1.337663in}}%
\pgfpathlineto{\pgfqpoint{1.649376in}{1.350911in}}%
\pgfpathlineto{\pgfqpoint{1.644590in}{1.364159in}}%
\pgfpathlineto{\pgfqpoint{1.645839in}{1.377407in}}%
\pgfpathlineto{\pgfqpoint{1.643577in}{1.390655in}}%
\pgfpathlineto{\pgfqpoint{1.641225in}{1.403903in}}%
\pgfpathlineto{\pgfqpoint{1.642048in}{1.417151in}}%
\pgfpathlineto{\pgfqpoint{1.637305in}{1.430399in}}%
\pgfpathlineto{\pgfqpoint{1.639364in}{1.443647in}}%
\pgfpathlineto{\pgfqpoint{1.634059in}{1.456895in}}%
\pgfpathlineto{\pgfqpoint{1.634949in}{1.470143in}}%
\pgfpathlineto{\pgfqpoint{1.631766in}{1.483391in}}%
\pgfpathlineto{\pgfqpoint{1.629816in}{1.496639in}}%
\pgfpathlineto{\pgfqpoint{1.629391in}{1.509887in}}%
\pgfpathlineto{\pgfqpoint{1.628151in}{1.513550in}}%
\pgfpathlineto{\pgfqpoint{1.624983in}{1.523135in}}%
\pgfpathlineto{\pgfqpoint{1.626366in}{1.536383in}}%
\pgfpathlineto{\pgfqpoint{1.621579in}{1.549631in}}%
\pgfpathlineto{\pgfqpoint{1.622179in}{1.562879in}}%
\pgfpathlineto{\pgfqpoint{1.618542in}{1.576127in}}%
\pgfpathlineto{\pgfqpoint{1.616629in}{1.589375in}}%
\pgfpathlineto{\pgfqpoint{1.615294in}{1.602623in}}%
\pgfpathlineto{\pgfqpoint{1.611213in}{1.615871in}}%
\pgfpathlineto{\pgfqpoint{1.611558in}{1.629119in}}%
\pgfpathlineto{\pgfqpoint{1.606618in}{1.642367in}}%
\pgfpathlineto{\pgfqpoint{1.607152in}{1.655615in}}%
\pgfpathlineto{\pgfqpoint{1.603466in}{1.668863in}}%
\pgfpathlineto{\pgfqpoint{1.602868in}{1.674765in}}%
\pgfpathlineto{\pgfqpoint{1.602150in}{1.682111in}}%
\pgfpathlineto{\pgfqpoint{1.600533in}{1.695359in}}%
\pgfpathlineto{\pgfqpoint{1.596712in}{1.708607in}}%
\pgfpathlineto{\pgfqpoint{1.596602in}{1.721855in}}%
\pgfpathlineto{\pgfqpoint{1.591910in}{1.735103in}}%
\pgfpathlineto{\pgfqpoint{1.591797in}{1.748351in}}%
\pgfpathlineto{\pgfqpoint{1.587774in}{1.761599in}}%
\pgfpathlineto{\pgfqpoint{1.586363in}{1.774847in}}%
\pgfpathlineto{\pgfqpoint{1.584400in}{1.788095in}}%
\pgfpathlineto{\pgfqpoint{1.581097in}{1.801343in}}%
\pgfpathlineto{\pgfqpoint{1.580786in}{1.814591in}}%
\pgfpathlineto{\pgfqpoint{1.577585in}{1.823568in}}%
\pgfpathlineto{\pgfqpoint{1.576079in}{1.827839in}}%
\pgfpathlineto{\pgfqpoint{1.575973in}{1.841087in}}%
\pgfpathlineto{\pgfqpoint{1.571962in}{1.854335in}}%
\pgfpathlineto{\pgfqpoint{1.570508in}{1.867583in}}%
\pgfpathlineto{\pgfqpoint{1.568134in}{1.880831in}}%
\pgfpathlineto{\pgfqpoint{1.564706in}{1.894079in}}%
\pgfpathlineto{\pgfqpoint{1.564070in}{1.907327in}}%
\pgfpathlineto{\pgfqpoint{1.559559in}{1.920575in}}%
\pgfpathlineto{\pgfqpoint{1.559249in}{1.933823in}}%
\pgfpathlineto{\pgfqpoint{1.554973in}{1.947071in}}%
\pgfpathlineto{\pgfqpoint{1.553433in}{1.960319in}}%
\pgfpathlineto{\pgfqpoint{1.552302in}{1.966082in}}%
\pgfpathlineto{\pgfqpoint{1.550827in}{1.973567in}}%
\pgfpathlineto{\pgfqpoint{1.547472in}{1.986815in}}%
\pgfpathlineto{\pgfqpoint{1.546593in}{2.000063in}}%
\pgfpathlineto{\pgfqpoint{1.541794in}{2.013311in}}%
\pgfpathlineto{\pgfqpoint{1.541422in}{2.026559in}}%
\pgfpathlineto{\pgfqpoint{1.536885in}{2.039807in}}%
\pgfpathlineto{\pgfqpoint{1.535148in}{2.053055in}}%
\pgfpathlineto{\pgfqpoint{1.532155in}{2.066303in}}%
\pgfpathlineto{\pgfqpoint{1.528250in}{2.079551in}}%
\pgfpathlineto{\pgfqpoint{1.527100in}{2.092799in}}%
\pgfpathlineto{\pgfqpoint{1.527018in}{2.093026in}}%
\pgfpathlineto{\pgfqpoint{1.521535in}{2.106047in}}%
\pgfpathlineto{\pgfqpoint{1.520922in}{2.119295in}}%
\pgfpathlineto{\pgfqpoint{1.515693in}{2.132543in}}%
\pgfpathlineto{\pgfqpoint{1.506774in}{2.132543in}}%
\pgfpathlineto{\pgfqpoint{1.511979in}{2.119295in}}%
\pgfpathlineto{\pgfqpoint{1.512565in}{2.106047in}}%
\pgfpathlineto{\pgfqpoint{1.518334in}{2.092799in}}%
\pgfpathlineto{\pgfqpoint{1.519413in}{2.079551in}}%
\pgfpathlineto{\pgfqpoint{1.524186in}{2.066303in}}%
\pgfpathlineto{\pgfqpoint{1.527018in}{2.054922in}}%
\pgfpathlineto{\pgfqpoint{1.527429in}{2.053055in}}%
\pgfpathlineto{\pgfqpoint{1.529277in}{2.039807in}}%
\pgfpathlineto{\pgfqpoint{1.533839in}{2.026559in}}%
\pgfpathlineto{\pgfqpoint{1.534147in}{2.013311in}}%
\pgfpathlineto{\pgfqpoint{1.539105in}{2.000063in}}%
\pgfpathlineto{\pgfqpoint{1.539809in}{1.986815in}}%
\pgfpathlineto{\pgfqpoint{1.543368in}{1.973567in}}%
\pgfpathlineto{\pgfqpoint{1.545814in}{1.960319in}}%
\pgfpathlineto{\pgfqpoint{1.547504in}{1.947071in}}%
\pgfpathlineto{\pgfqpoint{1.551912in}{1.933823in}}%
\pgfpathlineto{\pgfqpoint{1.552143in}{1.920575in}}%
\pgfpathlineto{\pgfqpoint{1.552302in}{1.920150in}}%
\pgfpathlineto{\pgfqpoint{1.556804in}{1.907327in}}%
\pgfpathlineto{\pgfqpoint{1.557248in}{1.894079in}}%
\pgfpathlineto{\pgfqpoint{1.560853in}{1.880831in}}%
\pgfpathlineto{\pgfqpoint{1.563068in}{1.867583in}}%
\pgfpathlineto{\pgfqpoint{1.564607in}{1.854335in}}%
\pgfpathlineto{\pgfqpoint{1.568568in}{1.841087in}}%
\pgfpathlineto{\pgfqpoint{1.568594in}{1.827839in}}%
\pgfpathlineto{\pgfqpoint{1.573297in}{1.814591in}}%
\pgfpathlineto{\pgfqpoint{1.573383in}{1.801343in}}%
\pgfpathlineto{\pgfqpoint{1.576778in}{1.788095in}}%
\pgfpathlineto{\pgfqpoint{1.577585in}{1.781922in}}%
\pgfpathlineto{\pgfqpoint{1.578543in}{1.774847in}}%
\pgfpathlineto{\pgfqpoint{1.580012in}{1.761599in}}%
\pgfpathlineto{\pgfqpoint{1.583972in}{1.748351in}}%
\pgfpathlineto{\pgfqpoint{1.583987in}{1.735103in}}%
\pgfpathlineto{\pgfqpoint{1.588759in}{1.721855in}}%
\pgfpathlineto{\pgfqpoint{1.588627in}{1.708607in}}%
\pgfpathlineto{\pgfqpoint{1.592575in}{1.695359in}}%
\pgfpathlineto{\pgfqpoint{1.593908in}{1.682111in}}%
\pgfpathlineto{\pgfqpoint{1.595240in}{1.668863in}}%
\pgfpathlineto{\pgfqpoint{1.598636in}{1.655615in}}%
\pgfpathlineto{\pgfqpoint{1.597985in}{1.642367in}}%
\pgfpathlineto{\pgfqpoint{1.602831in}{1.629119in}}%
\pgfpathlineto{\pgfqpoint{1.602210in}{1.615871in}}%
\pgfpathlineto{\pgfqpoint{1.602868in}{1.613790in}}%
\pgfpathlineto{\pgfqpoint{1.606403in}{1.602623in}}%
\pgfpathlineto{\pgfqpoint{1.607402in}{1.589375in}}%
\pgfpathlineto{\pgfqpoint{1.609363in}{1.576127in}}%
\pgfpathlineto{\pgfqpoint{1.612751in}{1.562879in}}%
\pgfpathlineto{\pgfqpoint{1.611996in}{1.549631in}}%
\pgfpathlineto{\pgfqpoint{1.616713in}{1.536383in}}%
\pgfpathlineto{\pgfqpoint{1.614958in}{1.523135in}}%
\pgfpathlineto{\pgfqpoint{1.619378in}{1.509887in}}%
\pgfpathlineto{\pgfqpoint{1.619316in}{1.496639in}}%
\pgfpathlineto{\pgfqpoint{1.621239in}{1.483391in}}%
\pgfpathlineto{\pgfqpoint{1.623968in}{1.470143in}}%
\pgfpathlineto{\pgfqpoint{1.622886in}{1.456895in}}%
\pgfpathlineto{\pgfqpoint{1.627961in}{1.443647in}}%
\pgfpathlineto{\pgfqpoint{1.625418in}{1.430399in}}%
\pgfpathlineto{\pgfqpoint{1.628151in}{1.422819in}}%
\pgfpathlineto{\pgfqpoint{1.630186in}{1.417151in}}%
\pgfpathlineto{\pgfqpoint{1.628714in}{1.403903in}}%
\pgfpathlineto{\pgfqpoint{1.631096in}{1.390655in}}%
\pgfpathlineto{\pgfqpoint{1.632735in}{1.377407in}}%
\pgfpathlineto{\pgfqpoint{1.631252in}{1.364159in}}%
\pgfpathlineto{\pgfqpoint{1.635593in}{1.350911in}}%
\pgfpathlineto{\pgfqpoint{1.630974in}{1.337663in}}%
\pgfpathlineto{\pgfqpoint{1.636031in}{1.324415in}}%
\pgfpathlineto{\pgfqpoint{1.631839in}{1.311167in}}%
\pgfpathlineto{\pgfqpoint{1.634169in}{1.297919in}}%
\pgfpathlineto{\pgfqpoint{1.633408in}{1.284671in}}%
\pgfpathlineto{\pgfqpoint{1.631381in}{1.271423in}}%
\pgfpathlineto{\pgfqpoint{1.634893in}{1.258175in}}%
\pgfpathlineto{\pgfqpoint{1.628462in}{1.244927in}}%
\pgfpathlineto{\pgfqpoint{1.633344in}{1.231679in}}%
\pgfpathlineto{\pgfqpoint{1.628151in}{1.223933in}}%
\pgfpathlineto{\pgfqpoint{1.623905in}{1.218431in}}%
\pgfpathlineto{\pgfqpoint{1.626710in}{1.205183in}}%
\pgfpathlineto{\pgfqpoint{1.619215in}{1.191935in}}%
\pgfpathlineto{\pgfqpoint{1.614835in}{1.178687in}}%
\pgfpathlineto{\pgfqpoint{1.612882in}{1.165439in}}%
\pgfpathlineto{\pgfqpoint{1.602868in}{1.154883in}}%
\pgfpathlineto{\pgfqpoint{1.599924in}{1.152191in}}%
\pgfpathlineto{\pgfqpoint{1.602868in}{1.141047in}}%
\pgfpathlineto{\pgfqpoint{1.603402in}{1.138943in}}%
\pgfpathlineto{\pgfqpoint{1.602868in}{1.138541in}}%
\pgfpathlineto{\pgfqpoint{1.582092in}{1.125695in}}%
\pgfpathlineto{\pgfqpoint{1.583869in}{1.112447in}}%
\pgfpathlineto{\pgfqpoint{1.577585in}{1.108810in}}%
\pgfpathlineto{\pgfqpoint{1.555331in}{1.099199in}}%
\pgfpathlineto{\pgfqpoint{1.552302in}{1.093409in}}%
\pgfpathlineto{\pgfqpoint{1.547528in}{1.085951in}}%
\pgfpathlineto{\pgfqpoint{1.527018in}{1.076865in}}%
\pgfpathlineto{\pgfqpoint{1.511564in}{1.072703in}}%
\pgfpathlineto{\pgfqpoint{1.501735in}{1.069727in}}%
\pgfpathlineto{\pgfqpoint{1.476452in}{1.064667in}}%
\pgfpathlineto{\pgfqpoint{1.451169in}{1.062295in}}%
\pgfpathlineto{\pgfqpoint{1.425886in}{1.063396in}}%
\pgfpathlineto{\pgfqpoint{1.400603in}{1.069736in}}%
\pgfpathlineto{\pgfqpoint{1.394847in}{1.072703in}}%
\pgfpathlineto{\pgfqpoint{1.375319in}{1.076798in}}%
\pgfpathlineto{\pgfqpoint{1.350036in}{1.084186in}}%
\pgfpathlineto{\pgfqpoint{1.345367in}{1.085951in}}%
\pgfpathlineto{\pgfqpoint{1.348207in}{1.099199in}}%
\pgfpathlineto{\pgfqpoint{1.324753in}{1.108227in}}%
\pgfpathlineto{\pgfqpoint{1.316421in}{1.112447in}}%
\pgfpathlineto{\pgfqpoint{1.314112in}{1.125695in}}%
\pgfpathlineto{\pgfqpoint{1.301120in}{1.138943in}}%
\pgfpathlineto{\pgfqpoint{1.299470in}{1.140909in}}%
\pgfpathlineto{\pgfqpoint{1.291161in}{1.152191in}}%
\pgfpathlineto{\pgfqpoint{1.292377in}{1.165439in}}%
\pgfpathlineto{\pgfqpoint{1.278173in}{1.178687in}}%
\pgfpathlineto{\pgfqpoint{1.283917in}{1.191935in}}%
\pgfpathlineto{\pgfqpoint{1.274187in}{1.204164in}}%
\pgfpathlineto{\pgfqpoint{1.273496in}{1.205183in}}%
\pgfpathlineto{\pgfqpoint{1.274187in}{1.209040in}}%
\pgfpathlineto{\pgfqpoint{1.275968in}{1.218431in}}%
\pgfpathlineto{\pgfqpoint{1.274187in}{1.225635in}}%
\pgfpathlineto{\pgfqpoint{1.272906in}{1.231679in}}%
\pgfpathlineto{\pgfqpoint{1.268786in}{1.244927in}}%
\pgfpathlineto{\pgfqpoint{1.272354in}{1.258175in}}%
\pgfpathlineto{\pgfqpoint{1.265552in}{1.271423in}}%
\pgfpathlineto{\pgfqpoint{1.271430in}{1.284671in}}%
\pgfpathlineto{\pgfqpoint{1.266025in}{1.297919in}}%
\pgfpathlineto{\pgfqpoint{1.269071in}{1.311167in}}%
\pgfpathlineto{\pgfqpoint{1.268659in}{1.324415in}}%
\pgfpathlineto{\pgfqpoint{1.267584in}{1.337663in}}%
\pgfpathlineto{\pgfqpoint{1.271651in}{1.350911in}}%
\pgfpathlineto{\pgfqpoint{1.267646in}{1.364159in}}%
\pgfpathlineto{\pgfqpoint{1.273332in}{1.377407in}}%
\pgfpathlineto{\pgfqpoint{1.270485in}{1.390655in}}%
\pgfpathlineto{\pgfqpoint{1.274187in}{1.402768in}}%
\pgfpathlineto{\pgfqpoint{1.274540in}{1.403903in}}%
\pgfpathlineto{\pgfqpoint{1.275092in}{1.417151in}}%
\pgfpathlineto{\pgfqpoint{1.275414in}{1.430399in}}%
\pgfpathlineto{\pgfqpoint{1.279398in}{1.443647in}}%
\pgfpathlineto{\pgfqpoint{1.277193in}{1.456895in}}%
\pgfpathlineto{\pgfqpoint{1.282511in}{1.470143in}}%
\pgfpathlineto{\pgfqpoint{1.280333in}{1.483391in}}%
\pgfpathlineto{\pgfqpoint{1.284329in}{1.496639in}}%
\pgfpathlineto{\pgfqpoint{1.285080in}{1.509887in}}%
\pgfpathlineto{\pgfqpoint{1.286636in}{1.523135in}}%
\pgfpathlineto{\pgfqpoint{1.290517in}{1.536383in}}%
\pgfpathlineto{\pgfqpoint{1.289741in}{1.549631in}}%
\pgfpathlineto{\pgfqpoint{1.294929in}{1.562879in}}%
\pgfpathlineto{\pgfqpoint{1.293686in}{1.576127in}}%
\pgfpathlineto{\pgfqpoint{1.297999in}{1.589375in}}%
\pgfpathlineto{\pgfqpoint{1.298320in}{1.602623in}}%
\pgfpathlineto{\pgfqpoint{1.299470in}{1.610746in}}%
\pgfpathlineto{\pgfqpoint{1.300192in}{1.615871in}}%
\pgfpathlineto{\pgfqpoint{1.303385in}{1.629119in}}%
\pgfpathlineto{\pgfqpoint{1.303364in}{1.642367in}}%
\pgfpathlineto{\pgfqpoint{1.308311in}{1.655615in}}%
\pgfpathlineto{\pgfqpoint{1.307726in}{1.668863in}}%
\pgfpathlineto{\pgfqpoint{1.312277in}{1.682111in}}%
\pgfpathlineto{\pgfqpoint{1.312806in}{1.695359in}}%
\pgfpathlineto{\pgfqpoint{1.315525in}{1.708607in}}%
\pgfpathlineto{\pgfqpoint{1.318136in}{1.721855in}}%
\pgfpathlineto{\pgfqpoint{1.318652in}{1.735103in}}%
\pgfpathlineto{\pgfqpoint{1.323130in}{1.748351in}}%
\pgfpathlineto{\pgfqpoint{1.322820in}{1.761599in}}%
\pgfpathlineto{\pgfqpoint{1.324753in}{1.767000in}}%
\pgfpathlineto{\pgfqpoint{1.327578in}{1.774847in}}%
\pgfpathlineto{\pgfqpoint{1.327915in}{1.788095in}}%
\pgfpathlineto{\pgfqpoint{1.331108in}{1.801343in}}%
\pgfpathlineto{\pgfqpoint{1.333301in}{1.814591in}}%
\pgfpathlineto{\pgfqpoint{1.334622in}{1.827839in}}%
\pgfpathlineto{\pgfqpoint{1.338771in}{1.841087in}}%
\pgfpathlineto{\pgfqpoint{1.338832in}{1.854335in}}%
\pgfpathlineto{\pgfqpoint{1.343645in}{1.867583in}}%
\pgfpathlineto{\pgfqpoint{1.344097in}{1.880831in}}%
\pgfpathlineto{\pgfqpoint{1.348048in}{1.894079in}}%
\pgfpathlineto{\pgfqpoint{1.349789in}{1.907327in}}%
\pgfpathlineto{\pgfqpoint{1.350036in}{1.909082in}}%
\pgfpathlineto{\pgfqpoint{1.351735in}{1.920575in}}%
\pgfpathlineto{\pgfqpoint{1.355382in}{1.933823in}}%
\pgfpathlineto{\pgfqpoint{1.355797in}{1.947071in}}%
\pgfpathlineto{\pgfqpoint{1.360790in}{1.960319in}}%
\pgfpathlineto{\pgfqpoint{1.361032in}{1.973567in}}%
\pgfpathlineto{\pgfqpoint{1.365633in}{1.986815in}}%
\pgfpathlineto{\pgfqpoint{1.367306in}{2.000063in}}%
\pgfpathlineto{\pgfqpoint{1.370410in}{2.013311in}}%
\pgfpathlineto{\pgfqpoint{1.373789in}{2.026559in}}%
\pgfpathlineto{\pgfqpoint{1.374888in}{2.039807in}}%
\pgfpathlineto{\pgfqpoint{1.375319in}{2.041023in}}%
\pgfpathlineto{\pgfqpoint{1.380182in}{2.053055in}}%
\pgfpathlineto{\pgfqpoint{1.380476in}{2.066303in}}%
\pgfpathlineto{\pgfqpoint{1.386264in}{2.079551in}}%
\pgfpathlineto{\pgfqpoint{1.387338in}{2.092799in}}%
\pgfpathlineto{\pgfqpoint{1.391674in}{2.106047in}}%
\pgfpathlineto{\pgfqpoint{1.394885in}{2.119295in}}%
\pgfpathlineto{\pgfqpoint{1.397256in}{2.132543in}}%
\pgfpathlineto{\pgfqpoint{1.388499in}{2.132543in}}%
\pgfpathlineto{\pgfqpoint{1.386316in}{2.119295in}}%
\pgfpathlineto{\pgfqpoint{1.382882in}{2.106047in}}%
\pgfpathlineto{\pgfqpoint{1.378837in}{2.092799in}}%
\pgfpathlineto{\pgfqpoint{1.377528in}{2.079551in}}%
\pgfpathlineto{\pgfqpoint{1.375319in}{2.074385in}}%
\pgfpathlineto{\pgfqpoint{1.372335in}{2.066303in}}%
\pgfpathlineto{\pgfqpoint{1.372003in}{2.053055in}}%
\pgfpathlineto{\pgfqpoint{1.367346in}{2.039807in}}%
\pgfpathlineto{\pgfqpoint{1.366312in}{2.026559in}}%
\pgfpathlineto{\pgfqpoint{1.362830in}{2.013311in}}%
\pgfpathlineto{\pgfqpoint{1.359925in}{2.000063in}}%
\pgfpathlineto{\pgfqpoint{1.358059in}{1.986815in}}%
\pgfpathlineto{\pgfqpoint{1.353653in}{1.973567in}}%
\pgfpathlineto{\pgfqpoint{1.353266in}{1.960319in}}%
\pgfpathlineto{\pgfqpoint{1.350036in}{1.951692in}}%
\pgfpathlineto{\pgfqpoint{1.348353in}{1.947071in}}%
\pgfpathlineto{\pgfqpoint{1.347941in}{1.933823in}}%
\pgfpathlineto{\pgfqpoint{1.344295in}{1.920575in}}%
\pgfpathlineto{\pgfqpoint{1.342520in}{1.907327in}}%
\pgfpathlineto{\pgfqpoint{1.340603in}{1.894079in}}%
\pgfpathlineto{\pgfqpoint{1.336841in}{1.880831in}}%
\pgfpathlineto{\pgfqpoint{1.336215in}{1.867583in}}%
\pgfpathlineto{\pgfqpoint{1.331519in}{1.854335in}}%
\pgfpathlineto{\pgfqpoint{1.331395in}{1.841087in}}%
\pgfpathlineto{\pgfqpoint{1.327175in}{1.827839in}}%
\pgfpathlineto{\pgfqpoint{1.325925in}{1.814591in}}%
\pgfpathlineto{\pgfqpoint{1.324753in}{1.808109in}}%
\pgfpathlineto{\pgfqpoint{1.323460in}{1.801343in}}%
\pgfpathlineto{\pgfqpoint{1.320342in}{1.788095in}}%
\pgfpathlineto{\pgfqpoint{1.319752in}{1.774847in}}%
\pgfpathlineto{\pgfqpoint{1.315099in}{1.761599in}}%
\pgfpathlineto{\pgfqpoint{1.315261in}{1.748351in}}%
\pgfpathlineto{\pgfqpoint{1.310776in}{1.735103in}}%
\pgfpathlineto{\pgfqpoint{1.310274in}{1.721855in}}%
\pgfpathlineto{\pgfqpoint{1.307462in}{1.708607in}}%
\pgfpathlineto{\pgfqpoint{1.304864in}{1.695359in}}%
\pgfpathlineto{\pgfqpoint{1.304033in}{1.682111in}}%
\pgfpathlineto{\pgfqpoint{1.299572in}{1.668863in}}%
\pgfpathlineto{\pgfqpoint{1.299923in}{1.655615in}}%
\pgfpathlineto{\pgfqpoint{1.299470in}{1.654412in}}%
\pgfpathlineto{\pgfqpoint{1.294770in}{1.642367in}}%
\pgfpathlineto{\pgfqpoint{1.294722in}{1.629119in}}%
\pgfpathlineto{\pgfqpoint{1.291239in}{1.615871in}}%
\pgfpathlineto{\pgfqpoint{1.289441in}{1.602623in}}%
\pgfpathlineto{\pgfqpoint{1.288762in}{1.589375in}}%
\pgfpathlineto{\pgfqpoint{1.284553in}{1.576127in}}%
\pgfpathlineto{\pgfqpoint{1.285454in}{1.562879in}}%
\pgfpathlineto{\pgfqpoint{1.280224in}{1.549631in}}%
\pgfpathlineto{\pgfqpoint{1.280848in}{1.536383in}}%
\pgfpathlineto{\pgfqpoint{1.276665in}{1.523135in}}%
\pgfpathlineto{\pgfqpoint{1.275145in}{1.509887in}}%
\pgfpathlineto{\pgfqpoint{1.274187in}{1.499625in}}%
\pgfpathlineto{\pgfqpoint{1.273887in}{1.496639in}}%
\pgfpathlineto{\pgfqpoint{1.269876in}{1.483391in}}%
\pgfpathlineto{\pgfqpoint{1.271582in}{1.470143in}}%
\pgfpathlineto{\pgfqpoint{1.266100in}{1.456895in}}%
\pgfpathlineto{\pgfqpoint{1.267992in}{1.443647in}}%
\pgfpathlineto{\pgfqpoint{1.263607in}{1.430399in}}%
\pgfpathlineto{\pgfqpoint{1.263245in}{1.417151in}}%
\pgfpathlineto{\pgfqpoint{1.262062in}{1.403903in}}%
\pgfpathlineto{\pgfqpoint{1.258081in}{1.390655in}}%
\pgfpathlineto{\pgfqpoint{1.260239in}{1.377407in}}%
\pgfpathlineto{\pgfqpoint{1.254406in}{1.364159in}}%
\pgfpathlineto{\pgfqpoint{1.257904in}{1.350911in}}%
\pgfpathlineto{\pgfqpoint{1.253214in}{1.337663in}}%
\pgfpathlineto{\pgfqpoint{1.254070in}{1.324415in}}%
\pgfpathlineto{\pgfqpoint{1.253413in}{1.311167in}}%
\pgfpathlineto{\pgfqpoint{1.250326in}{1.297919in}}%
\pgfpathlineto{\pgfqpoint{1.254590in}{1.284671in}}%
\pgfpathlineto{\pgfqpoint{1.248904in}{1.272512in}}%
\pgfpathlineto{\pgfqpoint{1.248396in}{1.271423in}}%
\pgfpathlineto{\pgfqpoint{1.248904in}{1.270233in}}%
\pgfpathlineto{\pgfqpoint{1.254402in}{1.258175in}}%
\pgfpathlineto{\pgfqpoint{1.249798in}{1.244927in}}%
\pgfpathlineto{\pgfqpoint{1.253575in}{1.231679in}}%
\pgfpathlineto{\pgfqpoint{1.254772in}{1.218431in}}%
\pgfpathlineto{\pgfqpoint{1.252195in}{1.205183in}}%
\pgfpathlineto{\pgfqpoint{1.259952in}{1.191935in}}%
\pgfpathlineto{\pgfqpoint{1.253811in}{1.178687in}}%
\pgfpathlineto{\pgfqpoint{1.265516in}{1.165439in}}%
\pgfpathlineto{\pgfqpoint{1.262246in}{1.152191in}}%
\pgfpathlineto{\pgfqpoint{1.270719in}{1.138943in}}%
\pgfpathlineto{\pgfqpoint{1.274187in}{1.132037in}}%
\pgfpathlineto{\pgfqpoint{1.278073in}{1.125695in}}%
\pgfpathlineto{\pgfqpoint{1.279749in}{1.112447in}}%
\pgfpathlineto{\pgfqpoint{1.299470in}{1.099883in}}%
\pgfpathlineto{\pgfqpoint{1.300825in}{1.099199in}}%
\pgfpathlineto{\pgfqpoint{1.299470in}{1.095617in}}%
\pgfpathlineto{\pgfqpoint{1.296073in}{1.085951in}}%
\pgfpathlineto{\pgfqpoint{1.299470in}{1.084045in}}%
\pgfpathlineto{\pgfqpoint{1.323955in}{1.072703in}}%
\pgfpathlineto{\pgfqpoint{1.324753in}{1.070203in}}%
\pgfpathlineto{\pgfqpoint{1.329459in}{1.059455in}}%
\pgfpathlineto{\pgfqpoint{1.350036in}{1.048908in}}%
\pgfpathlineto{\pgfqpoint{1.356557in}{1.046207in}}%
\pgfpathlineto{\pgfqpoint{1.375319in}{1.040750in}}%
\pgfpathclose%
\pgfusepath{fill}%
\end{pgfscope}%
\begin{pgfscope}%
\pgfpathrectangle{\pgfqpoint{0.540975in}{0.476543in}}{\pgfqpoint{1.820388in}{1.656000in}}%
\pgfusepath{clip}%
\pgfsetbuttcap%
\pgfsetroundjoin%
\definecolor{currentfill}{rgb}{0.999154,0.973779,0.709266}%
\pgfsetfillcolor{currentfill}%
\pgfsetlinewidth{0.000000pt}%
\definecolor{currentstroke}{rgb}{0.000000,0.000000,0.000000}%
\pgfsetstrokecolor{currentstroke}%
\pgfsetdash{}{0pt}%
\pgfpathmoveto{\pgfqpoint{1.400603in}{1.069736in}}%
\pgfpathlineto{\pgfqpoint{1.425886in}{1.063396in}}%
\pgfpathlineto{\pgfqpoint{1.451169in}{1.062295in}}%
\pgfpathlineto{\pgfqpoint{1.476452in}{1.064667in}}%
\pgfpathlineto{\pgfqpoint{1.501735in}{1.069727in}}%
\pgfpathlineto{\pgfqpoint{1.511564in}{1.072703in}}%
\pgfpathlineto{\pgfqpoint{1.527018in}{1.076865in}}%
\pgfpathlineto{\pgfqpoint{1.547528in}{1.085951in}}%
\pgfpathlineto{\pgfqpoint{1.552302in}{1.093409in}}%
\pgfpathlineto{\pgfqpoint{1.555331in}{1.099199in}}%
\pgfpathlineto{\pgfqpoint{1.577585in}{1.108810in}}%
\pgfpathlineto{\pgfqpoint{1.583869in}{1.112447in}}%
\pgfpathlineto{\pgfqpoint{1.582092in}{1.125695in}}%
\pgfpathlineto{\pgfqpoint{1.602868in}{1.138541in}}%
\pgfpathlineto{\pgfqpoint{1.603402in}{1.138943in}}%
\pgfpathlineto{\pgfqpoint{1.602868in}{1.141047in}}%
\pgfpathlineto{\pgfqpoint{1.599924in}{1.152191in}}%
\pgfpathlineto{\pgfqpoint{1.602868in}{1.154883in}}%
\pgfpathlineto{\pgfqpoint{1.612882in}{1.165439in}}%
\pgfpathlineto{\pgfqpoint{1.614835in}{1.178687in}}%
\pgfpathlineto{\pgfqpoint{1.619215in}{1.191935in}}%
\pgfpathlineto{\pgfqpoint{1.626710in}{1.205183in}}%
\pgfpathlineto{\pgfqpoint{1.623905in}{1.218431in}}%
\pgfpathlineto{\pgfqpoint{1.628151in}{1.223933in}}%
\pgfpathlineto{\pgfqpoint{1.633344in}{1.231679in}}%
\pgfpathlineto{\pgfqpoint{1.628462in}{1.244927in}}%
\pgfpathlineto{\pgfqpoint{1.634893in}{1.258175in}}%
\pgfpathlineto{\pgfqpoint{1.631381in}{1.271423in}}%
\pgfpathlineto{\pgfqpoint{1.633408in}{1.284671in}}%
\pgfpathlineto{\pgfqpoint{1.634169in}{1.297919in}}%
\pgfpathlineto{\pgfqpoint{1.631839in}{1.311167in}}%
\pgfpathlineto{\pgfqpoint{1.636031in}{1.324415in}}%
\pgfpathlineto{\pgfqpoint{1.630974in}{1.337663in}}%
\pgfpathlineto{\pgfqpoint{1.635593in}{1.350911in}}%
\pgfpathlineto{\pgfqpoint{1.631252in}{1.364159in}}%
\pgfpathlineto{\pgfqpoint{1.632735in}{1.377407in}}%
\pgfpathlineto{\pgfqpoint{1.631096in}{1.390655in}}%
\pgfpathlineto{\pgfqpoint{1.628714in}{1.403903in}}%
\pgfpathlineto{\pgfqpoint{1.630186in}{1.417151in}}%
\pgfpathlineto{\pgfqpoint{1.628151in}{1.422819in}}%
\pgfpathlineto{\pgfqpoint{1.625418in}{1.430399in}}%
\pgfpathlineto{\pgfqpoint{1.627961in}{1.443647in}}%
\pgfpathlineto{\pgfqpoint{1.622886in}{1.456895in}}%
\pgfpathlineto{\pgfqpoint{1.623968in}{1.470143in}}%
\pgfpathlineto{\pgfqpoint{1.621239in}{1.483391in}}%
\pgfpathlineto{\pgfqpoint{1.619316in}{1.496639in}}%
\pgfpathlineto{\pgfqpoint{1.619378in}{1.509887in}}%
\pgfpathlineto{\pgfqpoint{1.614958in}{1.523135in}}%
\pgfpathlineto{\pgfqpoint{1.616713in}{1.536383in}}%
\pgfpathlineto{\pgfqpoint{1.611996in}{1.549631in}}%
\pgfpathlineto{\pgfqpoint{1.612751in}{1.562879in}}%
\pgfpathlineto{\pgfqpoint{1.609363in}{1.576127in}}%
\pgfpathlineto{\pgfqpoint{1.607402in}{1.589375in}}%
\pgfpathlineto{\pgfqpoint{1.606403in}{1.602623in}}%
\pgfpathlineto{\pgfqpoint{1.602868in}{1.613790in}}%
\pgfpathlineto{\pgfqpoint{1.602210in}{1.615871in}}%
\pgfpathlineto{\pgfqpoint{1.602831in}{1.629119in}}%
\pgfpathlineto{\pgfqpoint{1.597985in}{1.642367in}}%
\pgfpathlineto{\pgfqpoint{1.598636in}{1.655615in}}%
\pgfpathlineto{\pgfqpoint{1.595240in}{1.668863in}}%
\pgfpathlineto{\pgfqpoint{1.593908in}{1.682111in}}%
\pgfpathlineto{\pgfqpoint{1.592575in}{1.695359in}}%
\pgfpathlineto{\pgfqpoint{1.588627in}{1.708607in}}%
\pgfpathlineto{\pgfqpoint{1.588759in}{1.721855in}}%
\pgfpathlineto{\pgfqpoint{1.583987in}{1.735103in}}%
\pgfpathlineto{\pgfqpoint{1.583972in}{1.748351in}}%
\pgfpathlineto{\pgfqpoint{1.580012in}{1.761599in}}%
\pgfpathlineto{\pgfqpoint{1.578543in}{1.774847in}}%
\pgfpathlineto{\pgfqpoint{1.577585in}{1.781922in}}%
\pgfpathlineto{\pgfqpoint{1.576778in}{1.788095in}}%
\pgfpathlineto{\pgfqpoint{1.573383in}{1.801343in}}%
\pgfpathlineto{\pgfqpoint{1.573297in}{1.814591in}}%
\pgfpathlineto{\pgfqpoint{1.568594in}{1.827839in}}%
\pgfpathlineto{\pgfqpoint{1.568568in}{1.841087in}}%
\pgfpathlineto{\pgfqpoint{1.564607in}{1.854335in}}%
\pgfpathlineto{\pgfqpoint{1.563068in}{1.867583in}}%
\pgfpathlineto{\pgfqpoint{1.560853in}{1.880831in}}%
\pgfpathlineto{\pgfqpoint{1.557248in}{1.894079in}}%
\pgfpathlineto{\pgfqpoint{1.556804in}{1.907327in}}%
\pgfpathlineto{\pgfqpoint{1.552302in}{1.920150in}}%
\pgfpathlineto{\pgfqpoint{1.552143in}{1.920575in}}%
\pgfpathlineto{\pgfqpoint{1.551912in}{1.933823in}}%
\pgfpathlineto{\pgfqpoint{1.547504in}{1.947071in}}%
\pgfpathlineto{\pgfqpoint{1.545814in}{1.960319in}}%
\pgfpathlineto{\pgfqpoint{1.543368in}{1.973567in}}%
\pgfpathlineto{\pgfqpoint{1.539809in}{1.986815in}}%
\pgfpathlineto{\pgfqpoint{1.539105in}{2.000063in}}%
\pgfpathlineto{\pgfqpoint{1.534147in}{2.013311in}}%
\pgfpathlineto{\pgfqpoint{1.533839in}{2.026559in}}%
\pgfpathlineto{\pgfqpoint{1.529277in}{2.039807in}}%
\pgfpathlineto{\pgfqpoint{1.527429in}{2.053055in}}%
\pgfpathlineto{\pgfqpoint{1.527018in}{2.054922in}}%
\pgfpathlineto{\pgfqpoint{1.524186in}{2.066303in}}%
\pgfpathlineto{\pgfqpoint{1.519413in}{2.079551in}}%
\pgfpathlineto{\pgfqpoint{1.518334in}{2.092799in}}%
\pgfpathlineto{\pgfqpoint{1.512565in}{2.106047in}}%
\pgfpathlineto{\pgfqpoint{1.511979in}{2.119295in}}%
\pgfpathlineto{\pgfqpoint{1.506774in}{2.132543in}}%
\pgfpathlineto{\pgfqpoint{1.501735in}{2.132543in}}%
\pgfpathlineto{\pgfqpoint{1.496095in}{2.132543in}}%
\pgfpathlineto{\pgfqpoint{1.501735in}{2.122614in}}%
\pgfpathlineto{\pgfqpoint{1.503035in}{2.119295in}}%
\pgfpathlineto{\pgfqpoint{1.503596in}{2.106047in}}%
\pgfpathlineto{\pgfqpoint{1.509557in}{2.092799in}}%
\pgfpathlineto{\pgfqpoint{1.510394in}{2.079551in}}%
\pgfpathlineto{\pgfqpoint{1.515457in}{2.066303in}}%
\pgfpathlineto{\pgfqpoint{1.518482in}{2.053055in}}%
\pgfpathlineto{\pgfqpoint{1.520815in}{2.039807in}}%
\pgfpathlineto{\pgfqpoint{1.526117in}{2.026559in}}%
\pgfpathlineto{\pgfqpoint{1.526410in}{2.013311in}}%
\pgfpathlineto{\pgfqpoint{1.527018in}{2.011995in}}%
\pgfpathlineto{\pgfqpoint{1.531618in}{2.000063in}}%
\pgfpathlineto{\pgfqpoint{1.532147in}{1.986815in}}%
\pgfpathlineto{\pgfqpoint{1.535908in}{1.973567in}}%
\pgfpathlineto{\pgfqpoint{1.538162in}{1.960319in}}%
\pgfpathlineto{\pgfqpoint{1.539967in}{1.947071in}}%
\pgfpathlineto{\pgfqpoint{1.544279in}{1.933823in}}%
\pgfpathlineto{\pgfqpoint{1.544430in}{1.920575in}}%
\pgfpathlineto{\pgfqpoint{1.549392in}{1.907327in}}%
\pgfpathlineto{\pgfqpoint{1.549646in}{1.894079in}}%
\pgfpathlineto{\pgfqpoint{1.552302in}{1.885356in}}%
\pgfpathlineto{\pgfqpoint{1.553573in}{1.880831in}}%
\pgfpathlineto{\pgfqpoint{1.555629in}{1.867583in}}%
\pgfpathlineto{\pgfqpoint{1.557252in}{1.854335in}}%
\pgfpathlineto{\pgfqpoint{1.561162in}{1.841087in}}%
\pgfpathlineto{\pgfqpoint{1.561108in}{1.827839in}}%
\pgfpathlineto{\pgfqpoint{1.565871in}{1.814591in}}%
\pgfpathlineto{\pgfqpoint{1.565723in}{1.801343in}}%
\pgfpathlineto{\pgfqpoint{1.569232in}{1.788095in}}%
\pgfpathlineto{\pgfqpoint{1.570718in}{1.774847in}}%
\pgfpathlineto{\pgfqpoint{1.572257in}{1.761599in}}%
\pgfpathlineto{\pgfqpoint{1.576129in}{1.748351in}}%
\pgfpathlineto{\pgfqpoint{1.576047in}{1.735103in}}%
\pgfpathlineto{\pgfqpoint{1.577585in}{1.730980in}}%
\pgfpathlineto{\pgfqpoint{1.580917in}{1.721855in}}%
\pgfpathlineto{\pgfqpoint{1.580542in}{1.708607in}}%
\pgfpathlineto{\pgfqpoint{1.584617in}{1.695359in}}%
\pgfpathlineto{\pgfqpoint{1.585667in}{1.682111in}}%
\pgfpathlineto{\pgfqpoint{1.587038in}{1.668863in}}%
\pgfpathlineto{\pgfqpoint{1.590235in}{1.655615in}}%
\pgfpathlineto{\pgfqpoint{1.589459in}{1.642367in}}%
\pgfpathlineto{\pgfqpoint{1.594254in}{1.629119in}}%
\pgfpathlineto{\pgfqpoint{1.593329in}{1.615871in}}%
\pgfpathlineto{\pgfqpoint{1.597557in}{1.602623in}}%
\pgfpathlineto{\pgfqpoint{1.598173in}{1.589375in}}%
\pgfpathlineto{\pgfqpoint{1.600182in}{1.576127in}}%
\pgfpathlineto{\pgfqpoint{1.602868in}{1.564755in}}%
\pgfpathlineto{\pgfqpoint{1.603322in}{1.562879in}}%
\pgfpathlineto{\pgfqpoint{1.602868in}{1.556204in}}%
\pgfpathlineto{\pgfqpoint{1.602408in}{1.549631in}}%
\pgfpathlineto{\pgfqpoint{1.602868in}{1.548326in}}%
\pgfpathlineto{\pgfqpoint{1.607061in}{1.536383in}}%
\pgfpathlineto{\pgfqpoint{1.604933in}{1.523135in}}%
\pgfpathlineto{\pgfqpoint{1.609400in}{1.509887in}}%
\pgfpathlineto{\pgfqpoint{1.608851in}{1.496639in}}%
\pgfpathlineto{\pgfqpoint{1.610784in}{1.483391in}}%
\pgfpathlineto{\pgfqpoint{1.613039in}{1.470143in}}%
\pgfpathlineto{\pgfqpoint{1.611770in}{1.456895in}}%
\pgfpathlineto{\pgfqpoint{1.616531in}{1.443647in}}%
\pgfpathlineto{\pgfqpoint{1.613509in}{1.430399in}}%
\pgfpathlineto{\pgfqpoint{1.618232in}{1.417151in}}%
\pgfpathlineto{\pgfqpoint{1.616031in}{1.403903in}}%
\pgfpathlineto{\pgfqpoint{1.618469in}{1.390655in}}%
\pgfpathlineto{\pgfqpoint{1.619418in}{1.377407in}}%
\pgfpathlineto{\pgfqpoint{1.617689in}{1.364159in}}%
\pgfpathlineto{\pgfqpoint{1.621595in}{1.350911in}}%
\pgfpathlineto{\pgfqpoint{1.616180in}{1.337663in}}%
\pgfpathlineto{\pgfqpoint{1.621093in}{1.324415in}}%
\pgfpathlineto{\pgfqpoint{1.615648in}{1.311167in}}%
\pgfpathlineto{\pgfqpoint{1.617976in}{1.297919in}}%
\pgfpathlineto{\pgfqpoint{1.615873in}{1.284671in}}%
\pgfpathlineto{\pgfqpoint{1.613539in}{1.271423in}}%
\pgfpathlineto{\pgfqpoint{1.616091in}{1.258175in}}%
\pgfpathlineto{\pgfqpoint{1.608462in}{1.244927in}}%
\pgfpathlineto{\pgfqpoint{1.612897in}{1.231679in}}%
\pgfpathlineto{\pgfqpoint{1.602868in}{1.219959in}}%
\pgfpathlineto{\pgfqpoint{1.601308in}{1.218431in}}%
\pgfpathlineto{\pgfqpoint{1.602868in}{1.210868in}}%
\pgfpathlineto{\pgfqpoint{1.603944in}{1.205183in}}%
\pgfpathlineto{\pgfqpoint{1.602868in}{1.203602in}}%
\pgfpathlineto{\pgfqpoint{1.593121in}{1.191935in}}%
\pgfpathlineto{\pgfqpoint{1.588000in}{1.178687in}}%
\pgfpathlineto{\pgfqpoint{1.582868in}{1.165439in}}%
\pgfpathlineto{\pgfqpoint{1.577585in}{1.160710in}}%
\pgfpathlineto{\pgfqpoint{1.566260in}{1.152191in}}%
\pgfpathlineto{\pgfqpoint{1.567251in}{1.138943in}}%
\pgfpathlineto{\pgfqpoint{1.552302in}{1.131701in}}%
\pgfpathlineto{\pgfqpoint{1.535788in}{1.125695in}}%
\pgfpathlineto{\pgfqpoint{1.536208in}{1.112447in}}%
\pgfpathlineto{\pgfqpoint{1.527018in}{1.109448in}}%
\pgfpathlineto{\pgfqpoint{1.501735in}{1.103929in}}%
\pgfpathlineto{\pgfqpoint{1.476452in}{1.100897in}}%
\pgfpathlineto{\pgfqpoint{1.451169in}{1.099960in}}%
\pgfpathlineto{\pgfqpoint{1.425886in}{1.100956in}}%
\pgfpathlineto{\pgfqpoint{1.400603in}{1.103914in}}%
\pgfpathlineto{\pgfqpoint{1.375319in}{1.109059in}}%
\pgfpathlineto{\pgfqpoint{1.363956in}{1.112447in}}%
\pgfpathlineto{\pgfqpoint{1.360198in}{1.125695in}}%
\pgfpathlineto{\pgfqpoint{1.350036in}{1.130371in}}%
\pgfpathlineto{\pgfqpoint{1.337134in}{1.138943in}}%
\pgfpathlineto{\pgfqpoint{1.324753in}{1.151340in}}%
\pgfpathlineto{\pgfqpoint{1.324031in}{1.152191in}}%
\pgfpathlineto{\pgfqpoint{1.322285in}{1.165439in}}%
\pgfpathlineto{\pgfqpoint{1.304847in}{1.178687in}}%
\pgfpathlineto{\pgfqpoint{1.310005in}{1.191935in}}%
\pgfpathlineto{\pgfqpoint{1.299470in}{1.201620in}}%
\pgfpathlineto{\pgfqpoint{1.296289in}{1.205183in}}%
\pgfpathlineto{\pgfqpoint{1.298533in}{1.218431in}}%
\pgfpathlineto{\pgfqpoint{1.293339in}{1.231679in}}%
\pgfpathlineto{\pgfqpoint{1.288671in}{1.244927in}}%
\pgfpathlineto{\pgfqpoint{1.291174in}{1.258175in}}%
\pgfpathlineto{\pgfqpoint{1.283219in}{1.271423in}}%
\pgfpathlineto{\pgfqpoint{1.288965in}{1.284671in}}%
\pgfpathlineto{\pgfqpoint{1.282087in}{1.297919in}}%
\pgfpathlineto{\pgfqpoint{1.285205in}{1.311167in}}%
\pgfpathlineto{\pgfqpoint{1.283578in}{1.324415in}}%
\pgfpathlineto{\pgfqpoint{1.282253in}{1.337663in}}%
\pgfpathlineto{\pgfqpoint{1.285690in}{1.350911in}}%
\pgfpathlineto{\pgfqpoint{1.281082in}{1.364159in}}%
\pgfpathlineto{\pgfqpoint{1.286657in}{1.377407in}}%
\pgfpathlineto{\pgfqpoint{1.283041in}{1.390655in}}%
\pgfpathlineto{\pgfqpoint{1.287181in}{1.403903in}}%
\pgfpathlineto{\pgfqpoint{1.286991in}{1.417151in}}%
\pgfpathlineto{\pgfqpoint{1.287287in}{1.430399in}}%
\pgfpathlineto{\pgfqpoint{1.290763in}{1.443647in}}%
\pgfpathlineto{\pgfqpoint{1.288252in}{1.456895in}}%
\pgfpathlineto{\pgfqpoint{1.293400in}{1.470143in}}%
\pgfpathlineto{\pgfqpoint{1.290715in}{1.483391in}}%
\pgfpathlineto{\pgfqpoint{1.294764in}{1.496639in}}%
\pgfpathlineto{\pgfqpoint{1.295015in}{1.509887in}}%
\pgfpathlineto{\pgfqpoint{1.296606in}{1.523135in}}%
\pgfpathlineto{\pgfqpoint{1.299470in}{1.533667in}}%
\pgfpathlineto{\pgfqpoint{1.300193in}{1.536383in}}%
\pgfpathlineto{\pgfqpoint{1.299470in}{1.546580in}}%
\pgfpathlineto{\pgfqpoint{1.299259in}{1.549631in}}%
\pgfpathlineto{\pgfqpoint{1.299470in}{1.550172in}}%
\pgfpathlineto{\pgfqpoint{1.304405in}{1.562879in}}%
\pgfpathlineto{\pgfqpoint{1.302819in}{1.576127in}}%
\pgfpathlineto{\pgfqpoint{1.307185in}{1.589375in}}%
\pgfpathlineto{\pgfqpoint{1.307095in}{1.602623in}}%
\pgfpathlineto{\pgfqpoint{1.309032in}{1.615871in}}%
\pgfpathlineto{\pgfqpoint{1.311920in}{1.629119in}}%
\pgfpathlineto{\pgfqpoint{1.311841in}{1.642367in}}%
\pgfpathlineto{\pgfqpoint{1.316698in}{1.655615in}}%
\pgfpathlineto{\pgfqpoint{1.315880in}{1.668863in}}%
\pgfpathlineto{\pgfqpoint{1.320521in}{1.682111in}}%
\pgfpathlineto{\pgfqpoint{1.320748in}{1.695359in}}%
\pgfpathlineto{\pgfqpoint{1.323589in}{1.708607in}}%
\pgfpathlineto{\pgfqpoint{1.324753in}{1.714924in}}%
\pgfpathlineto{\pgfqpoint{1.326008in}{1.721855in}}%
\pgfpathlineto{\pgfqpoint{1.326549in}{1.735103in}}%
\pgfpathlineto{\pgfqpoint{1.330978in}{1.748351in}}%
\pgfpathlineto{\pgfqpoint{1.330531in}{1.761599in}}%
\pgfpathlineto{\pgfqpoint{1.335340in}{1.774847in}}%
\pgfpathlineto{\pgfqpoint{1.335406in}{1.788095in}}%
\pgfpathlineto{\pgfqpoint{1.338726in}{1.801343in}}%
\pgfpathlineto{\pgfqpoint{1.340677in}{1.814591in}}%
\pgfpathlineto{\pgfqpoint{1.342069in}{1.827839in}}%
\pgfpathlineto{\pgfqpoint{1.346147in}{1.841087in}}%
\pgfpathlineto{\pgfqpoint{1.346144in}{1.854335in}}%
\pgfpathlineto{\pgfqpoint{1.350036in}{1.864828in}}%
\pgfpathlineto{\pgfqpoint{1.351127in}{1.867583in}}%
\pgfpathlineto{\pgfqpoint{1.351416in}{1.880831in}}%
\pgfpathlineto{\pgfqpoint{1.355731in}{1.894079in}}%
\pgfpathlineto{\pgfqpoint{1.357280in}{1.907327in}}%
\pgfpathlineto{\pgfqpoint{1.359373in}{1.920575in}}%
\pgfpathlineto{\pgfqpoint{1.362860in}{1.933823in}}%
\pgfpathlineto{\pgfqpoint{1.363277in}{1.947071in}}%
\pgfpathlineto{\pgfqpoint{1.368314in}{1.960319in}}%
\pgfpathlineto{\pgfqpoint{1.368411in}{1.973567in}}%
\pgfpathlineto{\pgfqpoint{1.373207in}{1.986815in}}%
\pgfpathlineto{\pgfqpoint{1.374687in}{2.000063in}}%
\pgfpathlineto{\pgfqpoint{1.375319in}{2.002654in}}%
\pgfpathlineto{\pgfqpoint{1.378406in}{2.013311in}}%
\pgfpathlineto{\pgfqpoint{1.382073in}{2.026559in}}%
\pgfpathlineto{\pgfqpoint{1.383466in}{2.039807in}}%
\pgfpathlineto{\pgfqpoint{1.388780in}{2.053055in}}%
\pgfpathlineto{\pgfqpoint{1.389007in}{2.066303in}}%
\pgfpathlineto{\pgfqpoint{1.395000in}{2.079551in}}%
\pgfpathlineto{\pgfqpoint{1.395840in}{2.092799in}}%
\pgfpathlineto{\pgfqpoint{1.400467in}{2.106047in}}%
\pgfpathlineto{\pgfqpoint{1.400603in}{2.106636in}}%
\pgfpathlineto{\pgfqpoint{1.404588in}{2.119295in}}%
\pgfpathlineto{\pgfqpoint{1.408287in}{2.132543in}}%
\pgfpathlineto{\pgfqpoint{1.400603in}{2.132543in}}%
\pgfpathlineto{\pgfqpoint{1.397256in}{2.132543in}}%
\pgfpathlineto{\pgfqpoint{1.394885in}{2.119295in}}%
\pgfpathlineto{\pgfqpoint{1.391674in}{2.106047in}}%
\pgfpathlineto{\pgfqpoint{1.387338in}{2.092799in}}%
\pgfpathlineto{\pgfqpoint{1.386264in}{2.079551in}}%
\pgfpathlineto{\pgfqpoint{1.380476in}{2.066303in}}%
\pgfpathlineto{\pgfqpoint{1.380182in}{2.053055in}}%
\pgfpathlineto{\pgfqpoint{1.375319in}{2.041023in}}%
\pgfpathlineto{\pgfqpoint{1.374888in}{2.039807in}}%
\pgfpathlineto{\pgfqpoint{1.373789in}{2.026559in}}%
\pgfpathlineto{\pgfqpoint{1.370410in}{2.013311in}}%
\pgfpathlineto{\pgfqpoint{1.367306in}{2.000063in}}%
\pgfpathlineto{\pgfqpoint{1.365633in}{1.986815in}}%
\pgfpathlineto{\pgfqpoint{1.361032in}{1.973567in}}%
\pgfpathlineto{\pgfqpoint{1.360790in}{1.960319in}}%
\pgfpathlineto{\pgfqpoint{1.355797in}{1.947071in}}%
\pgfpathlineto{\pgfqpoint{1.355382in}{1.933823in}}%
\pgfpathlineto{\pgfqpoint{1.351735in}{1.920575in}}%
\pgfpathlineto{\pgfqpoint{1.350036in}{1.909082in}}%
\pgfpathlineto{\pgfqpoint{1.349789in}{1.907327in}}%
\pgfpathlineto{\pgfqpoint{1.348048in}{1.894079in}}%
\pgfpathlineto{\pgfqpoint{1.344097in}{1.880831in}}%
\pgfpathlineto{\pgfqpoint{1.343645in}{1.867583in}}%
\pgfpathlineto{\pgfqpoint{1.338832in}{1.854335in}}%
\pgfpathlineto{\pgfqpoint{1.338771in}{1.841087in}}%
\pgfpathlineto{\pgfqpoint{1.334622in}{1.827839in}}%
\pgfpathlineto{\pgfqpoint{1.333301in}{1.814591in}}%
\pgfpathlineto{\pgfqpoint{1.331108in}{1.801343in}}%
\pgfpathlineto{\pgfqpoint{1.327915in}{1.788095in}}%
\pgfpathlineto{\pgfqpoint{1.327578in}{1.774847in}}%
\pgfpathlineto{\pgfqpoint{1.324753in}{1.767000in}}%
\pgfpathlineto{\pgfqpoint{1.322820in}{1.761599in}}%
\pgfpathlineto{\pgfqpoint{1.323130in}{1.748351in}}%
\pgfpathlineto{\pgfqpoint{1.318652in}{1.735103in}}%
\pgfpathlineto{\pgfqpoint{1.318136in}{1.721855in}}%
\pgfpathlineto{\pgfqpoint{1.315525in}{1.708607in}}%
\pgfpathlineto{\pgfqpoint{1.312806in}{1.695359in}}%
\pgfpathlineto{\pgfqpoint{1.312277in}{1.682111in}}%
\pgfpathlineto{\pgfqpoint{1.307726in}{1.668863in}}%
\pgfpathlineto{\pgfqpoint{1.308311in}{1.655615in}}%
\pgfpathlineto{\pgfqpoint{1.303364in}{1.642367in}}%
\pgfpathlineto{\pgfqpoint{1.303385in}{1.629119in}}%
\pgfpathlineto{\pgfqpoint{1.300192in}{1.615871in}}%
\pgfpathlineto{\pgfqpoint{1.299470in}{1.610746in}}%
\pgfpathlineto{\pgfqpoint{1.298320in}{1.602623in}}%
\pgfpathlineto{\pgfqpoint{1.297999in}{1.589375in}}%
\pgfpathlineto{\pgfqpoint{1.293686in}{1.576127in}}%
\pgfpathlineto{\pgfqpoint{1.294929in}{1.562879in}}%
\pgfpathlineto{\pgfqpoint{1.289741in}{1.549631in}}%
\pgfpathlineto{\pgfqpoint{1.290517in}{1.536383in}}%
\pgfpathlineto{\pgfqpoint{1.286636in}{1.523135in}}%
\pgfpathlineto{\pgfqpoint{1.285080in}{1.509887in}}%
\pgfpathlineto{\pgfqpoint{1.284329in}{1.496639in}}%
\pgfpathlineto{\pgfqpoint{1.280333in}{1.483391in}}%
\pgfpathlineto{\pgfqpoint{1.282511in}{1.470143in}}%
\pgfpathlineto{\pgfqpoint{1.277193in}{1.456895in}}%
\pgfpathlineto{\pgfqpoint{1.279398in}{1.443647in}}%
\pgfpathlineto{\pgfqpoint{1.275414in}{1.430399in}}%
\pgfpathlineto{\pgfqpoint{1.275092in}{1.417151in}}%
\pgfpathlineto{\pgfqpoint{1.274540in}{1.403903in}}%
\pgfpathlineto{\pgfqpoint{1.274187in}{1.402768in}}%
\pgfpathlineto{\pgfqpoint{1.270485in}{1.390655in}}%
\pgfpathlineto{\pgfqpoint{1.273332in}{1.377407in}}%
\pgfpathlineto{\pgfqpoint{1.267646in}{1.364159in}}%
\pgfpathlineto{\pgfqpoint{1.271651in}{1.350911in}}%
\pgfpathlineto{\pgfqpoint{1.267584in}{1.337663in}}%
\pgfpathlineto{\pgfqpoint{1.268659in}{1.324415in}}%
\pgfpathlineto{\pgfqpoint{1.269071in}{1.311167in}}%
\pgfpathlineto{\pgfqpoint{1.266025in}{1.297919in}}%
\pgfpathlineto{\pgfqpoint{1.271430in}{1.284671in}}%
\pgfpathlineto{\pgfqpoint{1.265552in}{1.271423in}}%
\pgfpathlineto{\pgfqpoint{1.272354in}{1.258175in}}%
\pgfpathlineto{\pgfqpoint{1.268786in}{1.244927in}}%
\pgfpathlineto{\pgfqpoint{1.272906in}{1.231679in}}%
\pgfpathlineto{\pgfqpoint{1.274187in}{1.225635in}}%
\pgfpathlineto{\pgfqpoint{1.275968in}{1.218431in}}%
\pgfpathlineto{\pgfqpoint{1.274187in}{1.209040in}}%
\pgfpathlineto{\pgfqpoint{1.273496in}{1.205183in}}%
\pgfpathlineto{\pgfqpoint{1.274187in}{1.204164in}}%
\pgfpathlineto{\pgfqpoint{1.283917in}{1.191935in}}%
\pgfpathlineto{\pgfqpoint{1.278173in}{1.178687in}}%
\pgfpathlineto{\pgfqpoint{1.292377in}{1.165439in}}%
\pgfpathlineto{\pgfqpoint{1.291161in}{1.152191in}}%
\pgfpathlineto{\pgfqpoint{1.299470in}{1.140909in}}%
\pgfpathlineto{\pgfqpoint{1.301120in}{1.138943in}}%
\pgfpathlineto{\pgfqpoint{1.314112in}{1.125695in}}%
\pgfpathlineto{\pgfqpoint{1.316421in}{1.112447in}}%
\pgfpathlineto{\pgfqpoint{1.324753in}{1.108227in}}%
\pgfpathlineto{\pgfqpoint{1.348207in}{1.099199in}}%
\pgfpathlineto{\pgfqpoint{1.345367in}{1.085951in}}%
\pgfpathlineto{\pgfqpoint{1.350036in}{1.084186in}}%
\pgfpathlineto{\pgfqpoint{1.375319in}{1.076798in}}%
\pgfpathlineto{\pgfqpoint{1.394847in}{1.072703in}}%
\pgfpathlineto{\pgfqpoint{1.400603in}{1.069736in}}%
\pgfpathclose%
\pgfusepath{fill}%
\end{pgfscope}%
\begin{pgfscope}%
\pgfpathrectangle{\pgfqpoint{0.540975in}{0.476543in}}{\pgfqpoint{1.820388in}{1.656000in}}%
\pgfusepath{clip}%
\pgfsetbuttcap%
\pgfsetroundjoin%
\definecolor{currentfill}{rgb}{0.997463,0.921338,0.629758}%
\pgfsetfillcolor{currentfill}%
\pgfsetlinewidth{0.000000pt}%
\definecolor{currentstroke}{rgb}{0.000000,0.000000,0.000000}%
\pgfsetstrokecolor{currentstroke}%
\pgfsetdash{}{0pt}%
\pgfpathmoveto{\pgfqpoint{1.375319in}{1.109059in}}%
\pgfpathlineto{\pgfqpoint{1.400603in}{1.103914in}}%
\pgfpathlineto{\pgfqpoint{1.425886in}{1.100956in}}%
\pgfpathlineto{\pgfqpoint{1.451169in}{1.099960in}}%
\pgfpathlineto{\pgfqpoint{1.476452in}{1.100897in}}%
\pgfpathlineto{\pgfqpoint{1.501735in}{1.103929in}}%
\pgfpathlineto{\pgfqpoint{1.527018in}{1.109448in}}%
\pgfpathlineto{\pgfqpoint{1.536208in}{1.112447in}}%
\pgfpathlineto{\pgfqpoint{1.535788in}{1.125695in}}%
\pgfpathlineto{\pgfqpoint{1.552302in}{1.131701in}}%
\pgfpathlineto{\pgfqpoint{1.567251in}{1.138943in}}%
\pgfpathlineto{\pgfqpoint{1.566260in}{1.152191in}}%
\pgfpathlineto{\pgfqpoint{1.577585in}{1.160710in}}%
\pgfpathlineto{\pgfqpoint{1.582868in}{1.165439in}}%
\pgfpathlineto{\pgfqpoint{1.588000in}{1.178687in}}%
\pgfpathlineto{\pgfqpoint{1.593121in}{1.191935in}}%
\pgfpathlineto{\pgfqpoint{1.602868in}{1.203602in}}%
\pgfpathlineto{\pgfqpoint{1.603944in}{1.205183in}}%
\pgfpathlineto{\pgfqpoint{1.602868in}{1.210868in}}%
\pgfpathlineto{\pgfqpoint{1.601308in}{1.218431in}}%
\pgfpathlineto{\pgfqpoint{1.602868in}{1.219959in}}%
\pgfpathlineto{\pgfqpoint{1.612897in}{1.231679in}}%
\pgfpathlineto{\pgfqpoint{1.608462in}{1.244927in}}%
\pgfpathlineto{\pgfqpoint{1.616091in}{1.258175in}}%
\pgfpathlineto{\pgfqpoint{1.613539in}{1.271423in}}%
\pgfpathlineto{\pgfqpoint{1.615873in}{1.284671in}}%
\pgfpathlineto{\pgfqpoint{1.617976in}{1.297919in}}%
\pgfpathlineto{\pgfqpoint{1.615648in}{1.311167in}}%
\pgfpathlineto{\pgfqpoint{1.621093in}{1.324415in}}%
\pgfpathlineto{\pgfqpoint{1.616180in}{1.337663in}}%
\pgfpathlineto{\pgfqpoint{1.621595in}{1.350911in}}%
\pgfpathlineto{\pgfqpoint{1.617689in}{1.364159in}}%
\pgfpathlineto{\pgfqpoint{1.619418in}{1.377407in}}%
\pgfpathlineto{\pgfqpoint{1.618469in}{1.390655in}}%
\pgfpathlineto{\pgfqpoint{1.616031in}{1.403903in}}%
\pgfpathlineto{\pgfqpoint{1.618232in}{1.417151in}}%
\pgfpathlineto{\pgfqpoint{1.613509in}{1.430399in}}%
\pgfpathlineto{\pgfqpoint{1.616531in}{1.443647in}}%
\pgfpathlineto{\pgfqpoint{1.611770in}{1.456895in}}%
\pgfpathlineto{\pgfqpoint{1.613039in}{1.470143in}}%
\pgfpathlineto{\pgfqpoint{1.610784in}{1.483391in}}%
\pgfpathlineto{\pgfqpoint{1.608851in}{1.496639in}}%
\pgfpathlineto{\pgfqpoint{1.609400in}{1.509887in}}%
\pgfpathlineto{\pgfqpoint{1.604933in}{1.523135in}}%
\pgfpathlineto{\pgfqpoint{1.607061in}{1.536383in}}%
\pgfpathlineto{\pgfqpoint{1.602868in}{1.548326in}}%
\pgfpathlineto{\pgfqpoint{1.602408in}{1.549631in}}%
\pgfpathlineto{\pgfqpoint{1.602868in}{1.556204in}}%
\pgfpathlineto{\pgfqpoint{1.603322in}{1.562879in}}%
\pgfpathlineto{\pgfqpoint{1.602868in}{1.564755in}}%
\pgfpathlineto{\pgfqpoint{1.600182in}{1.576127in}}%
\pgfpathlineto{\pgfqpoint{1.598173in}{1.589375in}}%
\pgfpathlineto{\pgfqpoint{1.597557in}{1.602623in}}%
\pgfpathlineto{\pgfqpoint{1.593329in}{1.615871in}}%
\pgfpathlineto{\pgfqpoint{1.594254in}{1.629119in}}%
\pgfpathlineto{\pgfqpoint{1.589459in}{1.642367in}}%
\pgfpathlineto{\pgfqpoint{1.590235in}{1.655615in}}%
\pgfpathlineto{\pgfqpoint{1.587038in}{1.668863in}}%
\pgfpathlineto{\pgfqpoint{1.585667in}{1.682111in}}%
\pgfpathlineto{\pgfqpoint{1.584617in}{1.695359in}}%
\pgfpathlineto{\pgfqpoint{1.580542in}{1.708607in}}%
\pgfpathlineto{\pgfqpoint{1.580917in}{1.721855in}}%
\pgfpathlineto{\pgfqpoint{1.577585in}{1.730980in}}%
\pgfpathlineto{\pgfqpoint{1.576047in}{1.735103in}}%
\pgfpathlineto{\pgfqpoint{1.576129in}{1.748351in}}%
\pgfpathlineto{\pgfqpoint{1.572257in}{1.761599in}}%
\pgfpathlineto{\pgfqpoint{1.570718in}{1.774847in}}%
\pgfpathlineto{\pgfqpoint{1.569232in}{1.788095in}}%
\pgfpathlineto{\pgfqpoint{1.565723in}{1.801343in}}%
\pgfpathlineto{\pgfqpoint{1.565871in}{1.814591in}}%
\pgfpathlineto{\pgfqpoint{1.561108in}{1.827839in}}%
\pgfpathlineto{\pgfqpoint{1.561162in}{1.841087in}}%
\pgfpathlineto{\pgfqpoint{1.557252in}{1.854335in}}%
\pgfpathlineto{\pgfqpoint{1.555629in}{1.867583in}}%
\pgfpathlineto{\pgfqpoint{1.553573in}{1.880831in}}%
\pgfpathlineto{\pgfqpoint{1.552302in}{1.885356in}}%
\pgfpathlineto{\pgfqpoint{1.549646in}{1.894079in}}%
\pgfpathlineto{\pgfqpoint{1.549392in}{1.907327in}}%
\pgfpathlineto{\pgfqpoint{1.544430in}{1.920575in}}%
\pgfpathlineto{\pgfqpoint{1.544279in}{1.933823in}}%
\pgfpathlineto{\pgfqpoint{1.539967in}{1.947071in}}%
\pgfpathlineto{\pgfqpoint{1.538162in}{1.960319in}}%
\pgfpathlineto{\pgfqpoint{1.535908in}{1.973567in}}%
\pgfpathlineto{\pgfqpoint{1.532147in}{1.986815in}}%
\pgfpathlineto{\pgfqpoint{1.531618in}{2.000063in}}%
\pgfpathlineto{\pgfqpoint{1.527018in}{2.011995in}}%
\pgfpathlineto{\pgfqpoint{1.526410in}{2.013311in}}%
\pgfpathlineto{\pgfqpoint{1.526117in}{2.026559in}}%
\pgfpathlineto{\pgfqpoint{1.520815in}{2.039807in}}%
\pgfpathlineto{\pgfqpoint{1.518482in}{2.053055in}}%
\pgfpathlineto{\pgfqpoint{1.515457in}{2.066303in}}%
\pgfpathlineto{\pgfqpoint{1.510394in}{2.079551in}}%
\pgfpathlineto{\pgfqpoint{1.509557in}{2.092799in}}%
\pgfpathlineto{\pgfqpoint{1.503596in}{2.106047in}}%
\pgfpathlineto{\pgfqpoint{1.503035in}{2.119295in}}%
\pgfpathlineto{\pgfqpoint{1.501735in}{2.122614in}}%
\pgfpathlineto{\pgfqpoint{1.496095in}{2.132543in}}%
\pgfpathlineto{\pgfqpoint{1.483127in}{2.132543in}}%
\pgfpathlineto{\pgfqpoint{1.490344in}{2.119295in}}%
\pgfpathlineto{\pgfqpoint{1.491308in}{2.106047in}}%
\pgfpathlineto{\pgfqpoint{1.500316in}{2.092799in}}%
\pgfpathlineto{\pgfqpoint{1.501198in}{2.079551in}}%
\pgfpathlineto{\pgfqpoint{1.501735in}{2.078686in}}%
\pgfpathlineto{\pgfqpoint{1.506728in}{2.066303in}}%
\pgfpathlineto{\pgfqpoint{1.509465in}{2.053055in}}%
\pgfpathlineto{\pgfqpoint{1.511993in}{2.039807in}}%
\pgfpathlineto{\pgfqpoint{1.517154in}{2.026559in}}%
\pgfpathlineto{\pgfqpoint{1.517432in}{2.013311in}}%
\pgfpathlineto{\pgfqpoint{1.523593in}{2.000063in}}%
\pgfpathlineto{\pgfqpoint{1.524004in}{1.986815in}}%
\pgfpathlineto{\pgfqpoint{1.527018in}{1.978430in}}%
\pgfpathlineto{\pgfqpoint{1.528449in}{1.973567in}}%
\pgfpathlineto{\pgfqpoint{1.530511in}{1.960319in}}%
\pgfpathlineto{\pgfqpoint{1.532430in}{1.947071in}}%
\pgfpathlineto{\pgfqpoint{1.536645in}{1.933823in}}%
\pgfpathlineto{\pgfqpoint{1.536717in}{1.920575in}}%
\pgfpathlineto{\pgfqpoint{1.541740in}{1.907327in}}%
\pgfpathlineto{\pgfqpoint{1.541756in}{1.894079in}}%
\pgfpathlineto{\pgfqpoint{1.545928in}{1.880831in}}%
\pgfpathlineto{\pgfqpoint{1.547890in}{1.867583in}}%
\pgfpathlineto{\pgfqpoint{1.549734in}{1.854335in}}%
\pgfpathlineto{\pgfqpoint{1.552302in}{1.846062in}}%
\pgfpathlineto{\pgfqpoint{1.553757in}{1.841087in}}%
\pgfpathlineto{\pgfqpoint{1.553622in}{1.827839in}}%
\pgfpathlineto{\pgfqpoint{1.558444in}{1.814591in}}%
\pgfpathlineto{\pgfqpoint{1.558062in}{1.801343in}}%
\pgfpathlineto{\pgfqpoint{1.561685in}{1.788095in}}%
\pgfpathlineto{\pgfqpoint{1.562893in}{1.774847in}}%
\pgfpathlineto{\pgfqpoint{1.564505in}{1.761599in}}%
\pgfpathlineto{\pgfqpoint{1.568201in}{1.748351in}}%
\pgfpathlineto{\pgfqpoint{1.568038in}{1.735103in}}%
\pgfpathlineto{\pgfqpoint{1.572971in}{1.721855in}}%
\pgfpathlineto{\pgfqpoint{1.572329in}{1.708607in}}%
\pgfpathlineto{\pgfqpoint{1.576630in}{1.695359in}}%
\pgfpathlineto{\pgfqpoint{1.577419in}{1.682111in}}%
\pgfpathlineto{\pgfqpoint{1.577585in}{1.680619in}}%
\pgfpathlineto{\pgfqpoint{1.578836in}{1.668863in}}%
\pgfpathlineto{\pgfqpoint{1.581834in}{1.655615in}}%
\pgfpathlineto{\pgfqpoint{1.580934in}{1.642367in}}%
\pgfpathlineto{\pgfqpoint{1.585676in}{1.629119in}}%
\pgfpathlineto{\pgfqpoint{1.584447in}{1.615871in}}%
\pgfpathlineto{\pgfqpoint{1.588742in}{1.602623in}}%
\pgfpathlineto{\pgfqpoint{1.588943in}{1.589375in}}%
\pgfpathlineto{\pgfqpoint{1.590998in}{1.576127in}}%
\pgfpathlineto{\pgfqpoint{1.593777in}{1.562879in}}%
\pgfpathlineto{\pgfqpoint{1.592730in}{1.549631in}}%
\pgfpathlineto{\pgfqpoint{1.597293in}{1.536383in}}%
\pgfpathlineto{\pgfqpoint{1.594746in}{1.523135in}}%
\pgfpathlineto{\pgfqpoint{1.599331in}{1.509887in}}%
\pgfpathlineto{\pgfqpoint{1.598241in}{1.496639in}}%
\pgfpathlineto{\pgfqpoint{1.600248in}{1.483391in}}%
\pgfpathlineto{\pgfqpoint{1.602078in}{1.470143in}}%
\pgfpathlineto{\pgfqpoint{1.600566in}{1.456895in}}%
\pgfpathlineto{\pgfqpoint{1.602868in}{1.450208in}}%
\pgfpathlineto{\pgfqpoint{1.605101in}{1.443647in}}%
\pgfpathlineto{\pgfqpoint{1.602868in}{1.435072in}}%
\pgfpathlineto{\pgfqpoint{1.601536in}{1.430399in}}%
\pgfpathlineto{\pgfqpoint{1.602868in}{1.426781in}}%
\pgfpathlineto{\pgfqpoint{1.606259in}{1.417151in}}%
\pgfpathlineto{\pgfqpoint{1.603340in}{1.403903in}}%
\pgfpathlineto{\pgfqpoint{1.605798in}{1.390655in}}%
\pgfpathlineto{\pgfqpoint{1.605986in}{1.377407in}}%
\pgfpathlineto{\pgfqpoint{1.604056in}{1.364159in}}%
\pgfpathlineto{\pgfqpoint{1.607346in}{1.350911in}}%
\pgfpathlineto{\pgfqpoint{1.602868in}{1.340992in}}%
\pgfpathlineto{\pgfqpoint{1.601182in}{1.337663in}}%
\pgfpathlineto{\pgfqpoint{1.602868in}{1.333018in}}%
\pgfpathlineto{\pgfqpoint{1.605846in}{1.324415in}}%
\pgfpathlineto{\pgfqpoint{1.602868in}{1.318157in}}%
\pgfpathlineto{\pgfqpoint{1.599000in}{1.311167in}}%
\pgfpathlineto{\pgfqpoint{1.601408in}{1.297919in}}%
\pgfpathlineto{\pgfqpoint{1.597605in}{1.284671in}}%
\pgfpathlineto{\pgfqpoint{1.594885in}{1.271423in}}%
\pgfpathlineto{\pgfqpoint{1.596265in}{1.258175in}}%
\pgfpathlineto{\pgfqpoint{1.587056in}{1.244927in}}%
\pgfpathlineto{\pgfqpoint{1.590957in}{1.231679in}}%
\pgfpathlineto{\pgfqpoint{1.577585in}{1.219359in}}%
\pgfpathlineto{\pgfqpoint{1.576331in}{1.218431in}}%
\pgfpathlineto{\pgfqpoint{1.577585in}{1.212294in}}%
\pgfpathlineto{\pgfqpoint{1.578862in}{1.205183in}}%
\pgfpathlineto{\pgfqpoint{1.577585in}{1.203826in}}%
\pgfpathlineto{\pgfqpoint{1.563130in}{1.191935in}}%
\pgfpathlineto{\pgfqpoint{1.557180in}{1.178687in}}%
\pgfpathlineto{\pgfqpoint{1.552302in}{1.171804in}}%
\pgfpathlineto{\pgfqpoint{1.545969in}{1.165439in}}%
\pgfpathlineto{\pgfqpoint{1.527018in}{1.154028in}}%
\pgfpathlineto{\pgfqpoint{1.522983in}{1.152191in}}%
\pgfpathlineto{\pgfqpoint{1.517680in}{1.138943in}}%
\pgfpathlineto{\pgfqpoint{1.501735in}{1.134732in}}%
\pgfpathlineto{\pgfqpoint{1.476452in}{1.130713in}}%
\pgfpathlineto{\pgfqpoint{1.451169in}{1.129228in}}%
\pgfpathlineto{\pgfqpoint{1.425886in}{1.130377in}}%
\pgfpathlineto{\pgfqpoint{1.400603in}{1.134603in}}%
\pgfpathlineto{\pgfqpoint{1.386674in}{1.138943in}}%
\pgfpathlineto{\pgfqpoint{1.375319in}{1.146032in}}%
\pgfpathlineto{\pgfqpoint{1.367876in}{1.152191in}}%
\pgfpathlineto{\pgfqpoint{1.359445in}{1.165439in}}%
\pgfpathlineto{\pgfqpoint{1.350036in}{1.169960in}}%
\pgfpathlineto{\pgfqpoint{1.335840in}{1.178687in}}%
\pgfpathlineto{\pgfqpoint{1.339969in}{1.191935in}}%
\pgfpathlineto{\pgfqpoint{1.324753in}{1.202230in}}%
\pgfpathlineto{\pgfqpoint{1.321327in}{1.205183in}}%
\pgfpathlineto{\pgfqpoint{1.323481in}{1.218431in}}%
\pgfpathlineto{\pgfqpoint{1.315289in}{1.231679in}}%
\pgfpathlineto{\pgfqpoint{1.309923in}{1.244927in}}%
\pgfpathlineto{\pgfqpoint{1.311088in}{1.258175in}}%
\pgfpathlineto{\pgfqpoint{1.301588in}{1.271423in}}%
\pgfpathlineto{\pgfqpoint{1.307287in}{1.284671in}}%
\pgfpathlineto{\pgfqpoint{1.299470in}{1.296300in}}%
\pgfpathlineto{\pgfqpoint{1.298541in}{1.297919in}}%
\pgfpathlineto{\pgfqpoint{1.299470in}{1.301968in}}%
\pgfpathlineto{\pgfqpoint{1.301754in}{1.311167in}}%
\pgfpathlineto{\pgfqpoint{1.299470in}{1.320644in}}%
\pgfpathlineto{\pgfqpoint{1.298698in}{1.324415in}}%
\pgfpathlineto{\pgfqpoint{1.297177in}{1.337663in}}%
\pgfpathlineto{\pgfqpoint{1.299470in}{1.349185in}}%
\pgfpathlineto{\pgfqpoint{1.299817in}{1.350911in}}%
\pgfpathlineto{\pgfqpoint{1.299470in}{1.351731in}}%
\pgfpathlineto{\pgfqpoint{1.294710in}{1.364159in}}%
\pgfpathlineto{\pgfqpoint{1.299470in}{1.376057in}}%
\pgfpathlineto{\pgfqpoint{1.300033in}{1.377407in}}%
\pgfpathlineto{\pgfqpoint{1.299470in}{1.378946in}}%
\pgfpathlineto{\pgfqpoint{1.295663in}{1.390655in}}%
\pgfpathlineto{\pgfqpoint{1.299470in}{1.402782in}}%
\pgfpathlineto{\pgfqpoint{1.299844in}{1.403903in}}%
\pgfpathlineto{\pgfqpoint{1.299470in}{1.408728in}}%
\pgfpathlineto{\pgfqpoint{1.298891in}{1.417151in}}%
\pgfpathlineto{\pgfqpoint{1.299159in}{1.430399in}}%
\pgfpathlineto{\pgfqpoint{1.299470in}{1.431733in}}%
\pgfpathlineto{\pgfqpoint{1.302243in}{1.443647in}}%
\pgfpathlineto{\pgfqpoint{1.299470in}{1.456128in}}%
\pgfpathlineto{\pgfqpoint{1.299311in}{1.456895in}}%
\pgfpathlineto{\pgfqpoint{1.299470in}{1.457312in}}%
\pgfpathlineto{\pgfqpoint{1.304460in}{1.470143in}}%
\pgfpathlineto{\pgfqpoint{1.301156in}{1.483391in}}%
\pgfpathlineto{\pgfqpoint{1.305372in}{1.496639in}}%
\pgfpathlineto{\pgfqpoint{1.305076in}{1.509887in}}%
\pgfpathlineto{\pgfqpoint{1.306742in}{1.523135in}}%
\pgfpathlineto{\pgfqpoint{1.309959in}{1.536383in}}%
\pgfpathlineto{\pgfqpoint{1.308897in}{1.549631in}}%
\pgfpathlineto{\pgfqpoint{1.313881in}{1.562879in}}%
\pgfpathlineto{\pgfqpoint{1.311953in}{1.576127in}}%
\pgfpathlineto{\pgfqpoint{1.316360in}{1.589375in}}%
\pgfpathlineto{\pgfqpoint{1.315853in}{1.602623in}}%
\pgfpathlineto{\pgfqpoint{1.317872in}{1.615871in}}%
\pgfpathlineto{\pgfqpoint{1.320456in}{1.629119in}}%
\pgfpathlineto{\pgfqpoint{1.320319in}{1.642367in}}%
\pgfpathlineto{\pgfqpoint{1.324753in}{1.654681in}}%
\pgfpathlineto{\pgfqpoint{1.325098in}{1.655615in}}%
\pgfpathlineto{\pgfqpoint{1.324753in}{1.659726in}}%
\pgfpathlineto{\pgfqpoint{1.324034in}{1.668863in}}%
\pgfpathlineto{\pgfqpoint{1.324753in}{1.670896in}}%
\pgfpathlineto{\pgfqpoint{1.328871in}{1.682111in}}%
\pgfpathlineto{\pgfqpoint{1.328773in}{1.695359in}}%
\pgfpathlineto{\pgfqpoint{1.331789in}{1.708607in}}%
\pgfpathlineto{\pgfqpoint{1.333931in}{1.721855in}}%
\pgfpathlineto{\pgfqpoint{1.334516in}{1.735103in}}%
\pgfpathlineto{\pgfqpoint{1.338820in}{1.748351in}}%
\pgfpathlineto{\pgfqpoint{1.338240in}{1.761599in}}%
\pgfpathlineto{\pgfqpoint{1.343103in}{1.774847in}}%
\pgfpathlineto{\pgfqpoint{1.342897in}{1.788095in}}%
\pgfpathlineto{\pgfqpoint{1.346344in}{1.801343in}}%
\pgfpathlineto{\pgfqpoint{1.348053in}{1.814591in}}%
\pgfpathlineto{\pgfqpoint{1.349516in}{1.827839in}}%
\pgfpathlineto{\pgfqpoint{1.350036in}{1.829546in}}%
\pgfpathlineto{\pgfqpoint{1.353731in}{1.841087in}}%
\pgfpathlineto{\pgfqpoint{1.353677in}{1.854335in}}%
\pgfpathlineto{\pgfqpoint{1.358934in}{1.867583in}}%
\pgfpathlineto{\pgfqpoint{1.359024in}{1.880831in}}%
\pgfpathlineto{\pgfqpoint{1.363500in}{1.894079in}}%
\pgfpathlineto{\pgfqpoint{1.364779in}{1.907327in}}%
\pgfpathlineto{\pgfqpoint{1.367011in}{1.920575in}}%
\pgfpathlineto{\pgfqpoint{1.370337in}{1.933823in}}%
\pgfpathlineto{\pgfqpoint{1.370757in}{1.947071in}}%
\pgfpathlineto{\pgfqpoint{1.375319in}{1.958973in}}%
\pgfpathlineto{\pgfqpoint{1.375927in}{1.960319in}}%
\pgfpathlineto{\pgfqpoint{1.375872in}{1.973567in}}%
\pgfpathlineto{\pgfqpoint{1.381677in}{1.986815in}}%
\pgfpathlineto{\pgfqpoint{1.383099in}{2.000063in}}%
\pgfpathlineto{\pgfqpoint{1.387169in}{2.013311in}}%
\pgfpathlineto{\pgfqpoint{1.390565in}{2.026559in}}%
\pgfpathlineto{\pgfqpoint{1.392106in}{2.039807in}}%
\pgfpathlineto{\pgfqpoint{1.397377in}{2.053055in}}%
\pgfpathlineto{\pgfqpoint{1.397538in}{2.066303in}}%
\pgfpathlineto{\pgfqpoint{1.400603in}{2.072932in}}%
\pgfpathlineto{\pgfqpoint{1.405081in}{2.079551in}}%
\pgfpathlineto{\pgfqpoint{1.405921in}{2.092799in}}%
\pgfpathlineto{\pgfqpoint{1.412991in}{2.106047in}}%
\pgfpathlineto{\pgfqpoint{1.416568in}{2.119295in}}%
\pgfpathlineto{\pgfqpoint{1.420725in}{2.132543in}}%
\pgfpathlineto{\pgfqpoint{1.408287in}{2.132543in}}%
\pgfpathlineto{\pgfqpoint{1.404588in}{2.119295in}}%
\pgfpathlineto{\pgfqpoint{1.400603in}{2.106636in}}%
\pgfpathlineto{\pgfqpoint{1.400467in}{2.106047in}}%
\pgfpathlineto{\pgfqpoint{1.395840in}{2.092799in}}%
\pgfpathlineto{\pgfqpoint{1.395000in}{2.079551in}}%
\pgfpathlineto{\pgfqpoint{1.389007in}{2.066303in}}%
\pgfpathlineto{\pgfqpoint{1.388780in}{2.053055in}}%
\pgfpathlineto{\pgfqpoint{1.383466in}{2.039807in}}%
\pgfpathlineto{\pgfqpoint{1.382073in}{2.026559in}}%
\pgfpathlineto{\pgfqpoint{1.378406in}{2.013311in}}%
\pgfpathlineto{\pgfqpoint{1.375319in}{2.002654in}}%
\pgfpathlineto{\pgfqpoint{1.374687in}{2.000063in}}%
\pgfpathlineto{\pgfqpoint{1.373207in}{1.986815in}}%
\pgfpathlineto{\pgfqpoint{1.368411in}{1.973567in}}%
\pgfpathlineto{\pgfqpoint{1.368314in}{1.960319in}}%
\pgfpathlineto{\pgfqpoint{1.363277in}{1.947071in}}%
\pgfpathlineto{\pgfqpoint{1.362860in}{1.933823in}}%
\pgfpathlineto{\pgfqpoint{1.359373in}{1.920575in}}%
\pgfpathlineto{\pgfqpoint{1.357280in}{1.907327in}}%
\pgfpathlineto{\pgfqpoint{1.355731in}{1.894079in}}%
\pgfpathlineto{\pgfqpoint{1.351416in}{1.880831in}}%
\pgfpathlineto{\pgfqpoint{1.351127in}{1.867583in}}%
\pgfpathlineto{\pgfqpoint{1.350036in}{1.864828in}}%
\pgfpathlineto{\pgfqpoint{1.346144in}{1.854335in}}%
\pgfpathlineto{\pgfqpoint{1.346147in}{1.841087in}}%
\pgfpathlineto{\pgfqpoint{1.342069in}{1.827839in}}%
\pgfpathlineto{\pgfqpoint{1.340677in}{1.814591in}}%
\pgfpathlineto{\pgfqpoint{1.338726in}{1.801343in}}%
\pgfpathlineto{\pgfqpoint{1.335406in}{1.788095in}}%
\pgfpathlineto{\pgfqpoint{1.335340in}{1.774847in}}%
\pgfpathlineto{\pgfqpoint{1.330531in}{1.761599in}}%
\pgfpathlineto{\pgfqpoint{1.330978in}{1.748351in}}%
\pgfpathlineto{\pgfqpoint{1.326549in}{1.735103in}}%
\pgfpathlineto{\pgfqpoint{1.326008in}{1.721855in}}%
\pgfpathlineto{\pgfqpoint{1.324753in}{1.714924in}}%
\pgfpathlineto{\pgfqpoint{1.323589in}{1.708607in}}%
\pgfpathlineto{\pgfqpoint{1.320748in}{1.695359in}}%
\pgfpathlineto{\pgfqpoint{1.320521in}{1.682111in}}%
\pgfpathlineto{\pgfqpoint{1.315880in}{1.668863in}}%
\pgfpathlineto{\pgfqpoint{1.316698in}{1.655615in}}%
\pgfpathlineto{\pgfqpoint{1.311841in}{1.642367in}}%
\pgfpathlineto{\pgfqpoint{1.311920in}{1.629119in}}%
\pgfpathlineto{\pgfqpoint{1.309032in}{1.615871in}}%
\pgfpathlineto{\pgfqpoint{1.307095in}{1.602623in}}%
\pgfpathlineto{\pgfqpoint{1.307185in}{1.589375in}}%
\pgfpathlineto{\pgfqpoint{1.302819in}{1.576127in}}%
\pgfpathlineto{\pgfqpoint{1.304405in}{1.562879in}}%
\pgfpathlineto{\pgfqpoint{1.299470in}{1.550172in}}%
\pgfpathlineto{\pgfqpoint{1.299259in}{1.549631in}}%
\pgfpathlineto{\pgfqpoint{1.299470in}{1.546580in}}%
\pgfpathlineto{\pgfqpoint{1.300193in}{1.536383in}}%
\pgfpathlineto{\pgfqpoint{1.299470in}{1.533667in}}%
\pgfpathlineto{\pgfqpoint{1.296606in}{1.523135in}}%
\pgfpathlineto{\pgfqpoint{1.295015in}{1.509887in}}%
\pgfpathlineto{\pgfqpoint{1.294764in}{1.496639in}}%
\pgfpathlineto{\pgfqpoint{1.290715in}{1.483391in}}%
\pgfpathlineto{\pgfqpoint{1.293400in}{1.470143in}}%
\pgfpathlineto{\pgfqpoint{1.288252in}{1.456895in}}%
\pgfpathlineto{\pgfqpoint{1.290763in}{1.443647in}}%
\pgfpathlineto{\pgfqpoint{1.287287in}{1.430399in}}%
\pgfpathlineto{\pgfqpoint{1.286991in}{1.417151in}}%
\pgfpathlineto{\pgfqpoint{1.287181in}{1.403903in}}%
\pgfpathlineto{\pgfqpoint{1.283041in}{1.390655in}}%
\pgfpathlineto{\pgfqpoint{1.286657in}{1.377407in}}%
\pgfpathlineto{\pgfqpoint{1.281082in}{1.364159in}}%
\pgfpathlineto{\pgfqpoint{1.285690in}{1.350911in}}%
\pgfpathlineto{\pgfqpoint{1.282253in}{1.337663in}}%
\pgfpathlineto{\pgfqpoint{1.283578in}{1.324415in}}%
\pgfpathlineto{\pgfqpoint{1.285205in}{1.311167in}}%
\pgfpathlineto{\pgfqpoint{1.282087in}{1.297919in}}%
\pgfpathlineto{\pgfqpoint{1.288965in}{1.284671in}}%
\pgfpathlineto{\pgfqpoint{1.283219in}{1.271423in}}%
\pgfpathlineto{\pgfqpoint{1.291174in}{1.258175in}}%
\pgfpathlineto{\pgfqpoint{1.288671in}{1.244927in}}%
\pgfpathlineto{\pgfqpoint{1.293339in}{1.231679in}}%
\pgfpathlineto{\pgfqpoint{1.298533in}{1.218431in}}%
\pgfpathlineto{\pgfqpoint{1.296289in}{1.205183in}}%
\pgfpathlineto{\pgfqpoint{1.299470in}{1.201620in}}%
\pgfpathlineto{\pgfqpoint{1.310005in}{1.191935in}}%
\pgfpathlineto{\pgfqpoint{1.304847in}{1.178687in}}%
\pgfpathlineto{\pgfqpoint{1.322285in}{1.165439in}}%
\pgfpathlineto{\pgfqpoint{1.324031in}{1.152191in}}%
\pgfpathlineto{\pgfqpoint{1.324753in}{1.151340in}}%
\pgfpathlineto{\pgfqpoint{1.337134in}{1.138943in}}%
\pgfpathlineto{\pgfqpoint{1.350036in}{1.130371in}}%
\pgfpathlineto{\pgfqpoint{1.360198in}{1.125695in}}%
\pgfpathlineto{\pgfqpoint{1.363956in}{1.112447in}}%
\pgfpathlineto{\pgfqpoint{1.375319in}{1.109059in}}%
\pgfpathclose%
\pgfusepath{fill}%
\end{pgfscope}%
\begin{pgfscope}%
\pgfpathrectangle{\pgfqpoint{0.540975in}{0.476543in}}{\pgfqpoint{1.820388in}{1.656000in}}%
\pgfusepath{clip}%
\pgfsetbuttcap%
\pgfsetroundjoin%
\definecolor{currentfill}{rgb}{0.995925,0.870742,0.557478}%
\pgfsetfillcolor{currentfill}%
\pgfsetlinewidth{0.000000pt}%
\definecolor{currentstroke}{rgb}{0.000000,0.000000,0.000000}%
\pgfsetstrokecolor{currentstroke}%
\pgfsetdash{}{0pt}%
\pgfpathmoveto{\pgfqpoint{1.400603in}{1.134603in}}%
\pgfpathlineto{\pgfqpoint{1.425886in}{1.130377in}}%
\pgfpathlineto{\pgfqpoint{1.451169in}{1.129228in}}%
\pgfpathlineto{\pgfqpoint{1.476452in}{1.130713in}}%
\pgfpathlineto{\pgfqpoint{1.501735in}{1.134732in}}%
\pgfpathlineto{\pgfqpoint{1.517680in}{1.138943in}}%
\pgfpathlineto{\pgfqpoint{1.522983in}{1.152191in}}%
\pgfpathlineto{\pgfqpoint{1.527018in}{1.154028in}}%
\pgfpathlineto{\pgfqpoint{1.545969in}{1.165439in}}%
\pgfpathlineto{\pgfqpoint{1.552302in}{1.171804in}}%
\pgfpathlineto{\pgfqpoint{1.557180in}{1.178687in}}%
\pgfpathlineto{\pgfqpoint{1.563130in}{1.191935in}}%
\pgfpathlineto{\pgfqpoint{1.577585in}{1.203826in}}%
\pgfpathlineto{\pgfqpoint{1.578862in}{1.205183in}}%
\pgfpathlineto{\pgfqpoint{1.577585in}{1.212294in}}%
\pgfpathlineto{\pgfqpoint{1.576331in}{1.218431in}}%
\pgfpathlineto{\pgfqpoint{1.577585in}{1.219359in}}%
\pgfpathlineto{\pgfqpoint{1.590957in}{1.231679in}}%
\pgfpathlineto{\pgfqpoint{1.587056in}{1.244927in}}%
\pgfpathlineto{\pgfqpoint{1.596265in}{1.258175in}}%
\pgfpathlineto{\pgfqpoint{1.594885in}{1.271423in}}%
\pgfpathlineto{\pgfqpoint{1.597605in}{1.284671in}}%
\pgfpathlineto{\pgfqpoint{1.601408in}{1.297919in}}%
\pgfpathlineto{\pgfqpoint{1.599000in}{1.311167in}}%
\pgfpathlineto{\pgfqpoint{1.602868in}{1.318157in}}%
\pgfpathlineto{\pgfqpoint{1.605846in}{1.324415in}}%
\pgfpathlineto{\pgfqpoint{1.602868in}{1.333018in}}%
\pgfpathlineto{\pgfqpoint{1.601182in}{1.337663in}}%
\pgfpathlineto{\pgfqpoint{1.602868in}{1.340992in}}%
\pgfpathlineto{\pgfqpoint{1.607346in}{1.350911in}}%
\pgfpathlineto{\pgfqpoint{1.604056in}{1.364159in}}%
\pgfpathlineto{\pgfqpoint{1.605986in}{1.377407in}}%
\pgfpathlineto{\pgfqpoint{1.605798in}{1.390655in}}%
\pgfpathlineto{\pgfqpoint{1.603340in}{1.403903in}}%
\pgfpathlineto{\pgfqpoint{1.606259in}{1.417151in}}%
\pgfpathlineto{\pgfqpoint{1.602868in}{1.426781in}}%
\pgfpathlineto{\pgfqpoint{1.601536in}{1.430399in}}%
\pgfpathlineto{\pgfqpoint{1.602868in}{1.435072in}}%
\pgfpathlineto{\pgfqpoint{1.605101in}{1.443647in}}%
\pgfpathlineto{\pgfqpoint{1.602868in}{1.450208in}}%
\pgfpathlineto{\pgfqpoint{1.600566in}{1.456895in}}%
\pgfpathlineto{\pgfqpoint{1.602078in}{1.470143in}}%
\pgfpathlineto{\pgfqpoint{1.600248in}{1.483391in}}%
\pgfpathlineto{\pgfqpoint{1.598241in}{1.496639in}}%
\pgfpathlineto{\pgfqpoint{1.599331in}{1.509887in}}%
\pgfpathlineto{\pgfqpoint{1.594746in}{1.523135in}}%
\pgfpathlineto{\pgfqpoint{1.597293in}{1.536383in}}%
\pgfpathlineto{\pgfqpoint{1.592730in}{1.549631in}}%
\pgfpathlineto{\pgfqpoint{1.593777in}{1.562879in}}%
\pgfpathlineto{\pgfqpoint{1.590998in}{1.576127in}}%
\pgfpathlineto{\pgfqpoint{1.588943in}{1.589375in}}%
\pgfpathlineto{\pgfqpoint{1.588742in}{1.602623in}}%
\pgfpathlineto{\pgfqpoint{1.584447in}{1.615871in}}%
\pgfpathlineto{\pgfqpoint{1.585676in}{1.629119in}}%
\pgfpathlineto{\pgfqpoint{1.580934in}{1.642367in}}%
\pgfpathlineto{\pgfqpoint{1.581834in}{1.655615in}}%
\pgfpathlineto{\pgfqpoint{1.578836in}{1.668863in}}%
\pgfpathlineto{\pgfqpoint{1.577585in}{1.680619in}}%
\pgfpathlineto{\pgfqpoint{1.577419in}{1.682111in}}%
\pgfpathlineto{\pgfqpoint{1.576630in}{1.695359in}}%
\pgfpathlineto{\pgfqpoint{1.572329in}{1.708607in}}%
\pgfpathlineto{\pgfqpoint{1.572971in}{1.721855in}}%
\pgfpathlineto{\pgfqpoint{1.568038in}{1.735103in}}%
\pgfpathlineto{\pgfqpoint{1.568201in}{1.748351in}}%
\pgfpathlineto{\pgfqpoint{1.564505in}{1.761599in}}%
\pgfpathlineto{\pgfqpoint{1.562893in}{1.774847in}}%
\pgfpathlineto{\pgfqpoint{1.561685in}{1.788095in}}%
\pgfpathlineto{\pgfqpoint{1.558062in}{1.801343in}}%
\pgfpathlineto{\pgfqpoint{1.558444in}{1.814591in}}%
\pgfpathlineto{\pgfqpoint{1.553622in}{1.827839in}}%
\pgfpathlineto{\pgfqpoint{1.553757in}{1.841087in}}%
\pgfpathlineto{\pgfqpoint{1.552302in}{1.846062in}}%
\pgfpathlineto{\pgfqpoint{1.549734in}{1.854335in}}%
\pgfpathlineto{\pgfqpoint{1.547890in}{1.867583in}}%
\pgfpathlineto{\pgfqpoint{1.545928in}{1.880831in}}%
\pgfpathlineto{\pgfqpoint{1.541756in}{1.894079in}}%
\pgfpathlineto{\pgfqpoint{1.541740in}{1.907327in}}%
\pgfpathlineto{\pgfqpoint{1.536717in}{1.920575in}}%
\pgfpathlineto{\pgfqpoint{1.536645in}{1.933823in}}%
\pgfpathlineto{\pgfqpoint{1.532430in}{1.947071in}}%
\pgfpathlineto{\pgfqpoint{1.530511in}{1.960319in}}%
\pgfpathlineto{\pgfqpoint{1.528449in}{1.973567in}}%
\pgfpathlineto{\pgfqpoint{1.527018in}{1.978430in}}%
\pgfpathlineto{\pgfqpoint{1.524004in}{1.986815in}}%
\pgfpathlineto{\pgfqpoint{1.523593in}{2.000063in}}%
\pgfpathlineto{\pgfqpoint{1.517432in}{2.013311in}}%
\pgfpathlineto{\pgfqpoint{1.517154in}{2.026559in}}%
\pgfpathlineto{\pgfqpoint{1.511993in}{2.039807in}}%
\pgfpathlineto{\pgfqpoint{1.509465in}{2.053055in}}%
\pgfpathlineto{\pgfqpoint{1.506728in}{2.066303in}}%
\pgfpathlineto{\pgfqpoint{1.501735in}{2.078686in}}%
\pgfpathlineto{\pgfqpoint{1.501198in}{2.079551in}}%
\pgfpathlineto{\pgfqpoint{1.500316in}{2.092799in}}%
\pgfpathlineto{\pgfqpoint{1.491308in}{2.106047in}}%
\pgfpathlineto{\pgfqpoint{1.490344in}{2.119295in}}%
\pgfpathlineto{\pgfqpoint{1.483127in}{2.132543in}}%
\pgfpathlineto{\pgfqpoint{1.476452in}{2.132543in}}%
\pgfpathlineto{\pgfqpoint{1.457955in}{2.132543in}}%
\pgfpathlineto{\pgfqpoint{1.476452in}{2.120356in}}%
\pgfpathlineto{\pgfqpoint{1.477015in}{2.119295in}}%
\pgfpathlineto{\pgfqpoint{1.478151in}{2.106047in}}%
\pgfpathlineto{\pgfqpoint{1.487288in}{2.092799in}}%
\pgfpathlineto{\pgfqpoint{1.487751in}{2.079551in}}%
\pgfpathlineto{\pgfqpoint{1.496198in}{2.066303in}}%
\pgfpathlineto{\pgfqpoint{1.499783in}{2.053055in}}%
\pgfpathlineto{\pgfqpoint{1.501735in}{2.046868in}}%
\pgfpathlineto{\pgfqpoint{1.503172in}{2.039807in}}%
\pgfpathlineto{\pgfqpoint{1.508191in}{2.026559in}}%
\pgfpathlineto{\pgfqpoint{1.508453in}{2.013311in}}%
\pgfpathlineto{\pgfqpoint{1.514710in}{2.000063in}}%
\pgfpathlineto{\pgfqpoint{1.514891in}{1.986815in}}%
\pgfpathlineto{\pgfqpoint{1.519860in}{1.973567in}}%
\pgfpathlineto{\pgfqpoint{1.522010in}{1.960319in}}%
\pgfpathlineto{\pgfqpoint{1.524481in}{1.947071in}}%
\pgfpathlineto{\pgfqpoint{1.527018in}{1.940193in}}%
\pgfpathlineto{\pgfqpoint{1.529012in}{1.933823in}}%
\pgfpathlineto{\pgfqpoint{1.529004in}{1.920575in}}%
\pgfpathlineto{\pgfqpoint{1.534089in}{1.907327in}}%
\pgfpathlineto{\pgfqpoint{1.533867in}{1.894079in}}%
\pgfpathlineto{\pgfqpoint{1.538207in}{1.880831in}}%
\pgfpathlineto{\pgfqpoint{1.539908in}{1.867583in}}%
\pgfpathlineto{\pgfqpoint{1.541880in}{1.854335in}}%
\pgfpathlineto{\pgfqpoint{1.545851in}{1.841087in}}%
\pgfpathlineto{\pgfqpoint{1.545650in}{1.827839in}}%
\pgfpathlineto{\pgfqpoint{1.550901in}{1.814591in}}%
\pgfpathlineto{\pgfqpoint{1.550225in}{1.801343in}}%
\pgfpathlineto{\pgfqpoint{1.552302in}{1.794656in}}%
\pgfpathlineto{\pgfqpoint{1.554138in}{1.788095in}}%
\pgfpathlineto{\pgfqpoint{1.555068in}{1.774847in}}%
\pgfpathlineto{\pgfqpoint{1.556753in}{1.761599in}}%
\pgfpathlineto{\pgfqpoint{1.560274in}{1.748351in}}%
\pgfpathlineto{\pgfqpoint{1.560029in}{1.735103in}}%
\pgfpathlineto{\pgfqpoint{1.564950in}{1.721855in}}%
\pgfpathlineto{\pgfqpoint{1.564042in}{1.708607in}}%
\pgfpathlineto{\pgfqpoint{1.568431in}{1.695359in}}%
\pgfpathlineto{\pgfqpoint{1.568855in}{1.682111in}}%
\pgfpathlineto{\pgfqpoint{1.570381in}{1.668863in}}%
\pgfpathlineto{\pgfqpoint{1.573227in}{1.655615in}}%
\pgfpathlineto{\pgfqpoint{1.572177in}{1.642367in}}%
\pgfpathlineto{\pgfqpoint{1.577070in}{1.629119in}}%
\pgfpathlineto{\pgfqpoint{1.575451in}{1.615871in}}%
\pgfpathlineto{\pgfqpoint{1.577585in}{1.609763in}}%
\pgfpathlineto{\pgfqpoint{1.579927in}{1.602623in}}%
\pgfpathlineto{\pgfqpoint{1.579714in}{1.589375in}}%
\pgfpathlineto{\pgfqpoint{1.581814in}{1.576127in}}%
\pgfpathlineto{\pgfqpoint{1.584226in}{1.562879in}}%
\pgfpathlineto{\pgfqpoint{1.583051in}{1.549631in}}%
\pgfpathlineto{\pgfqpoint{1.587435in}{1.536383in}}%
\pgfpathlineto{\pgfqpoint{1.584517in}{1.523135in}}%
\pgfpathlineto{\pgfqpoint{1.589088in}{1.509887in}}%
\pgfpathlineto{\pgfqpoint{1.587437in}{1.496639in}}%
\pgfpathlineto{\pgfqpoint{1.589456in}{1.483391in}}%
\pgfpathlineto{\pgfqpoint{1.590668in}{1.470143in}}%
\pgfpathlineto{\pgfqpoint{1.589002in}{1.456895in}}%
\pgfpathlineto{\pgfqpoint{1.593184in}{1.443647in}}%
\pgfpathlineto{\pgfqpoint{1.589022in}{1.430399in}}%
\pgfpathlineto{\pgfqpoint{1.593789in}{1.417151in}}%
\pgfpathlineto{\pgfqpoint{1.589896in}{1.403903in}}%
\pgfpathlineto{\pgfqpoint{1.592527in}{1.390655in}}%
\pgfpathlineto{\pgfqpoint{1.591816in}{1.377407in}}%
\pgfpathlineto{\pgfqpoint{1.589595in}{1.364159in}}%
\pgfpathlineto{\pgfqpoint{1.592322in}{1.350911in}}%
\pgfpathlineto{\pgfqpoint{1.585194in}{1.337663in}}%
\pgfpathlineto{\pgfqpoint{1.589579in}{1.324415in}}%
\pgfpathlineto{\pgfqpoint{1.581270in}{1.311167in}}%
\pgfpathlineto{\pgfqpoint{1.583554in}{1.297919in}}%
\pgfpathlineto{\pgfqpoint{1.578103in}{1.284671in}}%
\pgfpathlineto{\pgfqpoint{1.577585in}{1.282231in}}%
\pgfpathlineto{\pgfqpoint{1.574957in}{1.271423in}}%
\pgfpathlineto{\pgfqpoint{1.574693in}{1.258175in}}%
\pgfpathlineto{\pgfqpoint{1.563285in}{1.244927in}}%
\pgfpathlineto{\pgfqpoint{1.566406in}{1.231679in}}%
\pgfpathlineto{\pgfqpoint{1.552302in}{1.221534in}}%
\pgfpathlineto{\pgfqpoint{1.546615in}{1.218431in}}%
\pgfpathlineto{\pgfqpoint{1.549367in}{1.205183in}}%
\pgfpathlineto{\pgfqpoint{1.527018in}{1.192725in}}%
\pgfpathlineto{\pgfqpoint{1.524849in}{1.191935in}}%
\pgfpathlineto{\pgfqpoint{1.515955in}{1.178687in}}%
\pgfpathlineto{\pgfqpoint{1.501735in}{1.170715in}}%
\pgfpathlineto{\pgfqpoint{1.483135in}{1.165439in}}%
\pgfpathlineto{\pgfqpoint{1.476452in}{1.163503in}}%
\pgfpathlineto{\pgfqpoint{1.451169in}{1.160803in}}%
\pgfpathlineto{\pgfqpoint{1.425886in}{1.163730in}}%
\pgfpathlineto{\pgfqpoint{1.421584in}{1.165439in}}%
\pgfpathlineto{\pgfqpoint{1.400603in}{1.170031in}}%
\pgfpathlineto{\pgfqpoint{1.375707in}{1.178687in}}%
\pgfpathlineto{\pgfqpoint{1.378372in}{1.191935in}}%
\pgfpathlineto{\pgfqpoint{1.375319in}{1.192946in}}%
\pgfpathlineto{\pgfqpoint{1.350376in}{1.205183in}}%
\pgfpathlineto{\pgfqpoint{1.352949in}{1.218431in}}%
\pgfpathlineto{\pgfqpoint{1.350036in}{1.220691in}}%
\pgfpathlineto{\pgfqpoint{1.339829in}{1.231679in}}%
\pgfpathlineto{\pgfqpoint{1.333554in}{1.244927in}}%
\pgfpathlineto{\pgfqpoint{1.332799in}{1.258175in}}%
\pgfpathlineto{\pgfqpoint{1.324753in}{1.267223in}}%
\pgfpathlineto{\pgfqpoint{1.321598in}{1.271423in}}%
\pgfpathlineto{\pgfqpoint{1.324753in}{1.279732in}}%
\pgfpathlineto{\pgfqpoint{1.326818in}{1.284671in}}%
\pgfpathlineto{\pgfqpoint{1.324753in}{1.286895in}}%
\pgfpathlineto{\pgfqpoint{1.316409in}{1.297919in}}%
\pgfpathlineto{\pgfqpoint{1.319561in}{1.311167in}}%
\pgfpathlineto{\pgfqpoint{1.314961in}{1.324415in}}%
\pgfpathlineto{\pgfqpoint{1.313156in}{1.337663in}}%
\pgfpathlineto{\pgfqpoint{1.314908in}{1.350911in}}%
\pgfpathlineto{\pgfqpoint{1.309009in}{1.364159in}}%
\pgfpathlineto{\pgfqpoint{1.314232in}{1.377407in}}%
\pgfpathlineto{\pgfqpoint{1.308857in}{1.390655in}}%
\pgfpathlineto{\pgfqpoint{1.313251in}{1.403903in}}%
\pgfpathlineto{\pgfqpoint{1.311388in}{1.417151in}}%
\pgfpathlineto{\pgfqpoint{1.311671in}{1.430399in}}%
\pgfpathlineto{\pgfqpoint{1.314093in}{1.443647in}}%
\pgfpathlineto{\pgfqpoint{1.310883in}{1.456895in}}%
\pgfpathlineto{\pgfqpoint{1.315735in}{1.470143in}}%
\pgfpathlineto{\pgfqpoint{1.311911in}{1.483391in}}%
\pgfpathlineto{\pgfqpoint{1.316121in}{1.496639in}}%
\pgfpathlineto{\pgfqpoint{1.315238in}{1.509887in}}%
\pgfpathlineto{\pgfqpoint{1.316946in}{1.523135in}}%
\pgfpathlineto{\pgfqpoint{1.319725in}{1.536383in}}%
\pgfpathlineto{\pgfqpoint{1.318537in}{1.549631in}}%
\pgfpathlineto{\pgfqpoint{1.323357in}{1.562879in}}%
\pgfpathlineto{\pgfqpoint{1.321087in}{1.576127in}}%
\pgfpathlineto{\pgfqpoint{1.324753in}{1.587053in}}%
\pgfpathlineto{\pgfqpoint{1.325585in}{1.589375in}}%
\pgfpathlineto{\pgfqpoint{1.324753in}{1.600522in}}%
\pgfpathlineto{\pgfqpoint{1.324612in}{1.602623in}}%
\pgfpathlineto{\pgfqpoint{1.324753in}{1.603520in}}%
\pgfpathlineto{\pgfqpoint{1.326826in}{1.615871in}}%
\pgfpathlineto{\pgfqpoint{1.329185in}{1.629119in}}%
\pgfpathlineto{\pgfqpoint{1.328998in}{1.642367in}}%
\pgfpathlineto{\pgfqpoint{1.333779in}{1.655615in}}%
\pgfpathlineto{\pgfqpoint{1.332458in}{1.668863in}}%
\pgfpathlineto{\pgfqpoint{1.337334in}{1.682111in}}%
\pgfpathlineto{\pgfqpoint{1.336882in}{1.695359in}}%
\pgfpathlineto{\pgfqpoint{1.340013in}{1.708607in}}%
\pgfpathlineto{\pgfqpoint{1.341853in}{1.721855in}}%
\pgfpathlineto{\pgfqpoint{1.342482in}{1.735103in}}%
\pgfpathlineto{\pgfqpoint{1.346662in}{1.748351in}}%
\pgfpathlineto{\pgfqpoint{1.345948in}{1.761599in}}%
\pgfpathlineto{\pgfqpoint{1.350036in}{1.772626in}}%
\pgfpathlineto{\pgfqpoint{1.350940in}{1.774847in}}%
\pgfpathlineto{\pgfqpoint{1.350419in}{1.788095in}}%
\pgfpathlineto{\pgfqpoint{1.354292in}{1.801343in}}%
\pgfpathlineto{\pgfqpoint{1.355816in}{1.814591in}}%
\pgfpathlineto{\pgfqpoint{1.357492in}{1.827839in}}%
\pgfpathlineto{\pgfqpoint{1.361547in}{1.841087in}}%
\pgfpathlineto{\pgfqpoint{1.361461in}{1.854335in}}%
\pgfpathlineto{\pgfqpoint{1.366741in}{1.867583in}}%
\pgfpathlineto{\pgfqpoint{1.366631in}{1.880831in}}%
\pgfpathlineto{\pgfqpoint{1.371268in}{1.894079in}}%
\pgfpathlineto{\pgfqpoint{1.372278in}{1.907327in}}%
\pgfpathlineto{\pgfqpoint{1.374650in}{1.920575in}}%
\pgfpathlineto{\pgfqpoint{1.375319in}{1.923332in}}%
\pgfpathlineto{\pgfqpoint{1.378258in}{1.933823in}}%
\pgfpathlineto{\pgfqpoint{1.378789in}{1.947071in}}%
\pgfpathlineto{\pgfqpoint{1.384731in}{1.960319in}}%
\pgfpathlineto{\pgfqpoint{1.384527in}{1.973567in}}%
\pgfpathlineto{\pgfqpoint{1.390493in}{1.986815in}}%
\pgfpathlineto{\pgfqpoint{1.391608in}{2.000063in}}%
\pgfpathlineto{\pgfqpoint{1.395933in}{2.013311in}}%
\pgfpathlineto{\pgfqpoint{1.399056in}{2.026559in}}%
\pgfpathlineto{\pgfqpoint{1.400603in}{2.038702in}}%
\pgfpathlineto{\pgfqpoint{1.400811in}{2.039807in}}%
\pgfpathlineto{\pgfqpoint{1.408251in}{2.053055in}}%
\pgfpathlineto{\pgfqpoint{1.408497in}{2.066303in}}%
\pgfpathlineto{\pgfqpoint{1.417568in}{2.079551in}}%
\pgfpathlineto{\pgfqpoint{1.418016in}{2.092799in}}%
\pgfpathlineto{\pgfqpoint{1.425574in}{2.106047in}}%
\pgfpathlineto{\pgfqpoint{1.425886in}{2.107377in}}%
\pgfpathlineto{\pgfqpoint{1.432506in}{2.119295in}}%
\pgfpathlineto{\pgfqpoint{1.445271in}{2.132543in}}%
\pgfpathlineto{\pgfqpoint{1.425886in}{2.132543in}}%
\pgfpathlineto{\pgfqpoint{1.420725in}{2.132543in}}%
\pgfpathlineto{\pgfqpoint{1.416568in}{2.119295in}}%
\pgfpathlineto{\pgfqpoint{1.412991in}{2.106047in}}%
\pgfpathlineto{\pgfqpoint{1.405921in}{2.092799in}}%
\pgfpathlineto{\pgfqpoint{1.405081in}{2.079551in}}%
\pgfpathlineto{\pgfqpoint{1.400603in}{2.072932in}}%
\pgfpathlineto{\pgfqpoint{1.397538in}{2.066303in}}%
\pgfpathlineto{\pgfqpoint{1.397377in}{2.053055in}}%
\pgfpathlineto{\pgfqpoint{1.392106in}{2.039807in}}%
\pgfpathlineto{\pgfqpoint{1.390565in}{2.026559in}}%
\pgfpathlineto{\pgfqpoint{1.387169in}{2.013311in}}%
\pgfpathlineto{\pgfqpoint{1.383099in}{2.000063in}}%
\pgfpathlineto{\pgfqpoint{1.381677in}{1.986815in}}%
\pgfpathlineto{\pgfqpoint{1.375872in}{1.973567in}}%
\pgfpathlineto{\pgfqpoint{1.375927in}{1.960319in}}%
\pgfpathlineto{\pgfqpoint{1.375319in}{1.958973in}}%
\pgfpathlineto{\pgfqpoint{1.370757in}{1.947071in}}%
\pgfpathlineto{\pgfqpoint{1.370337in}{1.933823in}}%
\pgfpathlineto{\pgfqpoint{1.367011in}{1.920575in}}%
\pgfpathlineto{\pgfqpoint{1.364779in}{1.907327in}}%
\pgfpathlineto{\pgfqpoint{1.363500in}{1.894079in}}%
\pgfpathlineto{\pgfqpoint{1.359024in}{1.880831in}}%
\pgfpathlineto{\pgfqpoint{1.358934in}{1.867583in}}%
\pgfpathlineto{\pgfqpoint{1.353677in}{1.854335in}}%
\pgfpathlineto{\pgfqpoint{1.353731in}{1.841087in}}%
\pgfpathlineto{\pgfqpoint{1.350036in}{1.829546in}}%
\pgfpathlineto{\pgfqpoint{1.349516in}{1.827839in}}%
\pgfpathlineto{\pgfqpoint{1.348053in}{1.814591in}}%
\pgfpathlineto{\pgfqpoint{1.346344in}{1.801343in}}%
\pgfpathlineto{\pgfqpoint{1.342897in}{1.788095in}}%
\pgfpathlineto{\pgfqpoint{1.343103in}{1.774847in}}%
\pgfpathlineto{\pgfqpoint{1.338240in}{1.761599in}}%
\pgfpathlineto{\pgfqpoint{1.338820in}{1.748351in}}%
\pgfpathlineto{\pgfqpoint{1.334516in}{1.735103in}}%
\pgfpathlineto{\pgfqpoint{1.333931in}{1.721855in}}%
\pgfpathlineto{\pgfqpoint{1.331789in}{1.708607in}}%
\pgfpathlineto{\pgfqpoint{1.328773in}{1.695359in}}%
\pgfpathlineto{\pgfqpoint{1.328871in}{1.682111in}}%
\pgfpathlineto{\pgfqpoint{1.324753in}{1.670896in}}%
\pgfpathlineto{\pgfqpoint{1.324034in}{1.668863in}}%
\pgfpathlineto{\pgfqpoint{1.324753in}{1.659726in}}%
\pgfpathlineto{\pgfqpoint{1.325098in}{1.655615in}}%
\pgfpathlineto{\pgfqpoint{1.324753in}{1.654681in}}%
\pgfpathlineto{\pgfqpoint{1.320319in}{1.642367in}}%
\pgfpathlineto{\pgfqpoint{1.320456in}{1.629119in}}%
\pgfpathlineto{\pgfqpoint{1.317872in}{1.615871in}}%
\pgfpathlineto{\pgfqpoint{1.315853in}{1.602623in}}%
\pgfpathlineto{\pgfqpoint{1.316360in}{1.589375in}}%
\pgfpathlineto{\pgfqpoint{1.311953in}{1.576127in}}%
\pgfpathlineto{\pgfqpoint{1.313881in}{1.562879in}}%
\pgfpathlineto{\pgfqpoint{1.308897in}{1.549631in}}%
\pgfpathlineto{\pgfqpoint{1.309959in}{1.536383in}}%
\pgfpathlineto{\pgfqpoint{1.306742in}{1.523135in}}%
\pgfpathlineto{\pgfqpoint{1.305076in}{1.509887in}}%
\pgfpathlineto{\pgfqpoint{1.305372in}{1.496639in}}%
\pgfpathlineto{\pgfqpoint{1.301156in}{1.483391in}}%
\pgfpathlineto{\pgfqpoint{1.304460in}{1.470143in}}%
\pgfpathlineto{\pgfqpoint{1.299470in}{1.457312in}}%
\pgfpathlineto{\pgfqpoint{1.299311in}{1.456895in}}%
\pgfpathlineto{\pgfqpoint{1.299470in}{1.456128in}}%
\pgfpathlineto{\pgfqpoint{1.302243in}{1.443647in}}%
\pgfpathlineto{\pgfqpoint{1.299470in}{1.431733in}}%
\pgfpathlineto{\pgfqpoint{1.299159in}{1.430399in}}%
\pgfpathlineto{\pgfqpoint{1.298891in}{1.417151in}}%
\pgfpathlineto{\pgfqpoint{1.299470in}{1.408728in}}%
\pgfpathlineto{\pgfqpoint{1.299844in}{1.403903in}}%
\pgfpathlineto{\pgfqpoint{1.299470in}{1.402782in}}%
\pgfpathlineto{\pgfqpoint{1.295663in}{1.390655in}}%
\pgfpathlineto{\pgfqpoint{1.299470in}{1.378946in}}%
\pgfpathlineto{\pgfqpoint{1.300033in}{1.377407in}}%
\pgfpathlineto{\pgfqpoint{1.299470in}{1.376057in}}%
\pgfpathlineto{\pgfqpoint{1.294710in}{1.364159in}}%
\pgfpathlineto{\pgfqpoint{1.299470in}{1.351731in}}%
\pgfpathlineto{\pgfqpoint{1.299817in}{1.350911in}}%
\pgfpathlineto{\pgfqpoint{1.299470in}{1.349185in}}%
\pgfpathlineto{\pgfqpoint{1.297177in}{1.337663in}}%
\pgfpathlineto{\pgfqpoint{1.298698in}{1.324415in}}%
\pgfpathlineto{\pgfqpoint{1.299470in}{1.320644in}}%
\pgfpathlineto{\pgfqpoint{1.301754in}{1.311167in}}%
\pgfpathlineto{\pgfqpoint{1.299470in}{1.301968in}}%
\pgfpathlineto{\pgfqpoint{1.298541in}{1.297919in}}%
\pgfpathlineto{\pgfqpoint{1.299470in}{1.296300in}}%
\pgfpathlineto{\pgfqpoint{1.307287in}{1.284671in}}%
\pgfpathlineto{\pgfqpoint{1.301588in}{1.271423in}}%
\pgfpathlineto{\pgfqpoint{1.311088in}{1.258175in}}%
\pgfpathlineto{\pgfqpoint{1.309923in}{1.244927in}}%
\pgfpathlineto{\pgfqpoint{1.315289in}{1.231679in}}%
\pgfpathlineto{\pgfqpoint{1.323481in}{1.218431in}}%
\pgfpathlineto{\pgfqpoint{1.321327in}{1.205183in}}%
\pgfpathlineto{\pgfqpoint{1.324753in}{1.202230in}}%
\pgfpathlineto{\pgfqpoint{1.339969in}{1.191935in}}%
\pgfpathlineto{\pgfqpoint{1.335840in}{1.178687in}}%
\pgfpathlineto{\pgfqpoint{1.350036in}{1.169960in}}%
\pgfpathlineto{\pgfqpoint{1.359445in}{1.165439in}}%
\pgfpathlineto{\pgfqpoint{1.367876in}{1.152191in}}%
\pgfpathlineto{\pgfqpoint{1.375319in}{1.146032in}}%
\pgfpathlineto{\pgfqpoint{1.386674in}{1.138943in}}%
\pgfpathlineto{\pgfqpoint{1.400603in}{1.134603in}}%
\pgfpathclose%
\pgfusepath{fill}%
\end{pgfscope}%
\begin{pgfscope}%
\pgfpathrectangle{\pgfqpoint{0.540975in}{0.476543in}}{\pgfqpoint{1.820388in}{1.656000in}}%
\pgfusepath{clip}%
\pgfsetbuttcap%
\pgfsetroundjoin%
\definecolor{currentfill}{rgb}{0.994233,0.786159,0.477970}%
\pgfsetfillcolor{currentfill}%
\pgfsetlinewidth{0.000000pt}%
\definecolor{currentstroke}{rgb}{0.000000,0.000000,0.000000}%
\pgfsetstrokecolor{currentstroke}%
\pgfsetdash{}{0pt}%
\pgfpathmoveto{\pgfqpoint{1.425886in}{1.163730in}}%
\pgfpathlineto{\pgfqpoint{1.451169in}{1.160803in}}%
\pgfpathlineto{\pgfqpoint{1.476452in}{1.163503in}}%
\pgfpathlineto{\pgfqpoint{1.483135in}{1.165439in}}%
\pgfpathlineto{\pgfqpoint{1.501735in}{1.170715in}}%
\pgfpathlineto{\pgfqpoint{1.515955in}{1.178687in}}%
\pgfpathlineto{\pgfqpoint{1.524849in}{1.191935in}}%
\pgfpathlineto{\pgfqpoint{1.527018in}{1.192725in}}%
\pgfpathlineto{\pgfqpoint{1.549367in}{1.205183in}}%
\pgfpathlineto{\pgfqpoint{1.546615in}{1.218431in}}%
\pgfpathlineto{\pgfqpoint{1.552302in}{1.221534in}}%
\pgfpathlineto{\pgfqpoint{1.566406in}{1.231679in}}%
\pgfpathlineto{\pgfqpoint{1.563285in}{1.244927in}}%
\pgfpathlineto{\pgfqpoint{1.574693in}{1.258175in}}%
\pgfpathlineto{\pgfqpoint{1.574957in}{1.271423in}}%
\pgfpathlineto{\pgfqpoint{1.577585in}{1.282231in}}%
\pgfpathlineto{\pgfqpoint{1.578103in}{1.284671in}}%
\pgfpathlineto{\pgfqpoint{1.583554in}{1.297919in}}%
\pgfpathlineto{\pgfqpoint{1.581270in}{1.311167in}}%
\pgfpathlineto{\pgfqpoint{1.589579in}{1.324415in}}%
\pgfpathlineto{\pgfqpoint{1.585194in}{1.337663in}}%
\pgfpathlineto{\pgfqpoint{1.592322in}{1.350911in}}%
\pgfpathlineto{\pgfqpoint{1.589595in}{1.364159in}}%
\pgfpathlineto{\pgfqpoint{1.591816in}{1.377407in}}%
\pgfpathlineto{\pgfqpoint{1.592527in}{1.390655in}}%
\pgfpathlineto{\pgfqpoint{1.589896in}{1.403903in}}%
\pgfpathlineto{\pgfqpoint{1.593789in}{1.417151in}}%
\pgfpathlineto{\pgfqpoint{1.589022in}{1.430399in}}%
\pgfpathlineto{\pgfqpoint{1.593184in}{1.443647in}}%
\pgfpathlineto{\pgfqpoint{1.589002in}{1.456895in}}%
\pgfpathlineto{\pgfqpoint{1.590668in}{1.470143in}}%
\pgfpathlineto{\pgfqpoint{1.589456in}{1.483391in}}%
\pgfpathlineto{\pgfqpoint{1.587437in}{1.496639in}}%
\pgfpathlineto{\pgfqpoint{1.589088in}{1.509887in}}%
\pgfpathlineto{\pgfqpoint{1.584517in}{1.523135in}}%
\pgfpathlineto{\pgfqpoint{1.587435in}{1.536383in}}%
\pgfpathlineto{\pgfqpoint{1.583051in}{1.549631in}}%
\pgfpathlineto{\pgfqpoint{1.584226in}{1.562879in}}%
\pgfpathlineto{\pgfqpoint{1.581814in}{1.576127in}}%
\pgfpathlineto{\pgfqpoint{1.579714in}{1.589375in}}%
\pgfpathlineto{\pgfqpoint{1.579927in}{1.602623in}}%
\pgfpathlineto{\pgfqpoint{1.577585in}{1.609763in}}%
\pgfpathlineto{\pgfqpoint{1.575451in}{1.615871in}}%
\pgfpathlineto{\pgfqpoint{1.577070in}{1.629119in}}%
\pgfpathlineto{\pgfqpoint{1.572177in}{1.642367in}}%
\pgfpathlineto{\pgfqpoint{1.573227in}{1.655615in}}%
\pgfpathlineto{\pgfqpoint{1.570381in}{1.668863in}}%
\pgfpathlineto{\pgfqpoint{1.568855in}{1.682111in}}%
\pgfpathlineto{\pgfqpoint{1.568431in}{1.695359in}}%
\pgfpathlineto{\pgfqpoint{1.564042in}{1.708607in}}%
\pgfpathlineto{\pgfqpoint{1.564950in}{1.721855in}}%
\pgfpathlineto{\pgfqpoint{1.560029in}{1.735103in}}%
\pgfpathlineto{\pgfqpoint{1.560274in}{1.748351in}}%
\pgfpathlineto{\pgfqpoint{1.556753in}{1.761599in}}%
\pgfpathlineto{\pgfqpoint{1.555068in}{1.774847in}}%
\pgfpathlineto{\pgfqpoint{1.554138in}{1.788095in}}%
\pgfpathlineto{\pgfqpoint{1.552302in}{1.794656in}}%
\pgfpathlineto{\pgfqpoint{1.550225in}{1.801343in}}%
\pgfpathlineto{\pgfqpoint{1.550901in}{1.814591in}}%
\pgfpathlineto{\pgfqpoint{1.545650in}{1.827839in}}%
\pgfpathlineto{\pgfqpoint{1.545851in}{1.841087in}}%
\pgfpathlineto{\pgfqpoint{1.541880in}{1.854335in}}%
\pgfpathlineto{\pgfqpoint{1.539908in}{1.867583in}}%
\pgfpathlineto{\pgfqpoint{1.538207in}{1.880831in}}%
\pgfpathlineto{\pgfqpoint{1.533867in}{1.894079in}}%
\pgfpathlineto{\pgfqpoint{1.534089in}{1.907327in}}%
\pgfpathlineto{\pgfqpoint{1.529004in}{1.920575in}}%
\pgfpathlineto{\pgfqpoint{1.529012in}{1.933823in}}%
\pgfpathlineto{\pgfqpoint{1.527018in}{1.940193in}}%
\pgfpathlineto{\pgfqpoint{1.524481in}{1.947071in}}%
\pgfpathlineto{\pgfqpoint{1.522010in}{1.960319in}}%
\pgfpathlineto{\pgfqpoint{1.519860in}{1.973567in}}%
\pgfpathlineto{\pgfqpoint{1.514891in}{1.986815in}}%
\pgfpathlineto{\pgfqpoint{1.514710in}{2.000063in}}%
\pgfpathlineto{\pgfqpoint{1.508453in}{2.013311in}}%
\pgfpathlineto{\pgfqpoint{1.508191in}{2.026559in}}%
\pgfpathlineto{\pgfqpoint{1.503172in}{2.039807in}}%
\pgfpathlineto{\pgfqpoint{1.501735in}{2.046868in}}%
\pgfpathlineto{\pgfqpoint{1.499783in}{2.053055in}}%
\pgfpathlineto{\pgfqpoint{1.496198in}{2.066303in}}%
\pgfpathlineto{\pgfqpoint{1.487751in}{2.079551in}}%
\pgfpathlineto{\pgfqpoint{1.487288in}{2.092799in}}%
\pgfpathlineto{\pgfqpoint{1.478151in}{2.106047in}}%
\pgfpathlineto{\pgfqpoint{1.477015in}{2.119295in}}%
\pgfpathlineto{\pgfqpoint{1.476452in}{2.120356in}}%
\pgfpathlineto{\pgfqpoint{1.457955in}{2.132543in}}%
\pgfpathlineto{\pgfqpoint{1.451169in}{2.132543in}}%
\pgfpathlineto{\pgfqpoint{1.445271in}{2.132543in}}%
\pgfpathlineto{\pgfqpoint{1.432506in}{2.119295in}}%
\pgfpathlineto{\pgfqpoint{1.425886in}{2.107377in}}%
\pgfpathlineto{\pgfqpoint{1.425574in}{2.106047in}}%
\pgfpathlineto{\pgfqpoint{1.418016in}{2.092799in}}%
\pgfpathlineto{\pgfqpoint{1.417568in}{2.079551in}}%
\pgfpathlineto{\pgfqpoint{1.408497in}{2.066303in}}%
\pgfpathlineto{\pgfqpoint{1.408251in}{2.053055in}}%
\pgfpathlineto{\pgfqpoint{1.400811in}{2.039807in}}%
\pgfpathlineto{\pgfqpoint{1.400603in}{2.038702in}}%
\pgfpathlineto{\pgfqpoint{1.399056in}{2.026559in}}%
\pgfpathlineto{\pgfqpoint{1.395933in}{2.013311in}}%
\pgfpathlineto{\pgfqpoint{1.391608in}{2.000063in}}%
\pgfpathlineto{\pgfqpoint{1.390493in}{1.986815in}}%
\pgfpathlineto{\pgfqpoint{1.384527in}{1.973567in}}%
\pgfpathlineto{\pgfqpoint{1.384731in}{1.960319in}}%
\pgfpathlineto{\pgfqpoint{1.378789in}{1.947071in}}%
\pgfpathlineto{\pgfqpoint{1.378258in}{1.933823in}}%
\pgfpathlineto{\pgfqpoint{1.375319in}{1.923332in}}%
\pgfpathlineto{\pgfqpoint{1.374650in}{1.920575in}}%
\pgfpathlineto{\pgfqpoint{1.372278in}{1.907327in}}%
\pgfpathlineto{\pgfqpoint{1.371268in}{1.894079in}}%
\pgfpathlineto{\pgfqpoint{1.366631in}{1.880831in}}%
\pgfpathlineto{\pgfqpoint{1.366741in}{1.867583in}}%
\pgfpathlineto{\pgfqpoint{1.361461in}{1.854335in}}%
\pgfpathlineto{\pgfqpoint{1.361547in}{1.841087in}}%
\pgfpathlineto{\pgfqpoint{1.357492in}{1.827839in}}%
\pgfpathlineto{\pgfqpoint{1.355816in}{1.814591in}}%
\pgfpathlineto{\pgfqpoint{1.354292in}{1.801343in}}%
\pgfpathlineto{\pgfqpoint{1.350419in}{1.788095in}}%
\pgfpathlineto{\pgfqpoint{1.350940in}{1.774847in}}%
\pgfpathlineto{\pgfqpoint{1.350036in}{1.772626in}}%
\pgfpathlineto{\pgfqpoint{1.345948in}{1.761599in}}%
\pgfpathlineto{\pgfqpoint{1.346662in}{1.748351in}}%
\pgfpathlineto{\pgfqpoint{1.342482in}{1.735103in}}%
\pgfpathlineto{\pgfqpoint{1.341853in}{1.721855in}}%
\pgfpathlineto{\pgfqpoint{1.340013in}{1.708607in}}%
\pgfpathlineto{\pgfqpoint{1.336882in}{1.695359in}}%
\pgfpathlineto{\pgfqpoint{1.337334in}{1.682111in}}%
\pgfpathlineto{\pgfqpoint{1.332458in}{1.668863in}}%
\pgfpathlineto{\pgfqpoint{1.333779in}{1.655615in}}%
\pgfpathlineto{\pgfqpoint{1.328998in}{1.642367in}}%
\pgfpathlineto{\pgfqpoint{1.329185in}{1.629119in}}%
\pgfpathlineto{\pgfqpoint{1.326826in}{1.615871in}}%
\pgfpathlineto{\pgfqpoint{1.324753in}{1.603520in}}%
\pgfpathlineto{\pgfqpoint{1.324612in}{1.602623in}}%
\pgfpathlineto{\pgfqpoint{1.324753in}{1.600522in}}%
\pgfpathlineto{\pgfqpoint{1.325585in}{1.589375in}}%
\pgfpathlineto{\pgfqpoint{1.324753in}{1.587053in}}%
\pgfpathlineto{\pgfqpoint{1.321087in}{1.576127in}}%
\pgfpathlineto{\pgfqpoint{1.323357in}{1.562879in}}%
\pgfpathlineto{\pgfqpoint{1.318537in}{1.549631in}}%
\pgfpathlineto{\pgfqpoint{1.319725in}{1.536383in}}%
\pgfpathlineto{\pgfqpoint{1.316946in}{1.523135in}}%
\pgfpathlineto{\pgfqpoint{1.315238in}{1.509887in}}%
\pgfpathlineto{\pgfqpoint{1.316121in}{1.496639in}}%
\pgfpathlineto{\pgfqpoint{1.311911in}{1.483391in}}%
\pgfpathlineto{\pgfqpoint{1.315735in}{1.470143in}}%
\pgfpathlineto{\pgfqpoint{1.310883in}{1.456895in}}%
\pgfpathlineto{\pgfqpoint{1.314093in}{1.443647in}}%
\pgfpathlineto{\pgfqpoint{1.311671in}{1.430399in}}%
\pgfpathlineto{\pgfqpoint{1.311388in}{1.417151in}}%
\pgfpathlineto{\pgfqpoint{1.313251in}{1.403903in}}%
\pgfpathlineto{\pgfqpoint{1.308857in}{1.390655in}}%
\pgfpathlineto{\pgfqpoint{1.314232in}{1.377407in}}%
\pgfpathlineto{\pgfqpoint{1.309009in}{1.364159in}}%
\pgfpathlineto{\pgfqpoint{1.314908in}{1.350911in}}%
\pgfpathlineto{\pgfqpoint{1.313156in}{1.337663in}}%
\pgfpathlineto{\pgfqpoint{1.314961in}{1.324415in}}%
\pgfpathlineto{\pgfqpoint{1.319561in}{1.311167in}}%
\pgfpathlineto{\pgfqpoint{1.316409in}{1.297919in}}%
\pgfpathlineto{\pgfqpoint{1.324753in}{1.286895in}}%
\pgfpathlineto{\pgfqpoint{1.326818in}{1.284671in}}%
\pgfpathlineto{\pgfqpoint{1.324753in}{1.279732in}}%
\pgfpathlineto{\pgfqpoint{1.321598in}{1.271423in}}%
\pgfpathlineto{\pgfqpoint{1.324753in}{1.267223in}}%
\pgfpathlineto{\pgfqpoint{1.332799in}{1.258175in}}%
\pgfpathlineto{\pgfqpoint{1.333554in}{1.244927in}}%
\pgfpathlineto{\pgfqpoint{1.339829in}{1.231679in}}%
\pgfpathlineto{\pgfqpoint{1.350036in}{1.220691in}}%
\pgfpathlineto{\pgfqpoint{1.352949in}{1.218431in}}%
\pgfpathlineto{\pgfqpoint{1.350376in}{1.205183in}}%
\pgfpathlineto{\pgfqpoint{1.375319in}{1.192946in}}%
\pgfpathlineto{\pgfqpoint{1.378372in}{1.191935in}}%
\pgfpathlineto{\pgfqpoint{1.375707in}{1.178687in}}%
\pgfpathlineto{\pgfqpoint{1.400603in}{1.170031in}}%
\pgfpathlineto{\pgfqpoint{1.421584in}{1.165439in}}%
\pgfpathlineto{\pgfqpoint{1.425886in}{1.163730in}}%
\pgfpathclose%
\pgfpathmoveto{\pgfqpoint{1.391241in}{1.205183in}}%
\pgfpathlineto{\pgfqpoint{1.393987in}{1.218431in}}%
\pgfpathlineto{\pgfqpoint{1.375319in}{1.227266in}}%
\pgfpathlineto{\pgfqpoint{1.369305in}{1.231679in}}%
\pgfpathlineto{\pgfqpoint{1.361858in}{1.244927in}}%
\pgfpathlineto{\pgfqpoint{1.357891in}{1.258175in}}%
\pgfpathlineto{\pgfqpoint{1.350036in}{1.265077in}}%
\pgfpathlineto{\pgfqpoint{1.344242in}{1.271423in}}%
\pgfpathlineto{\pgfqpoint{1.348852in}{1.284671in}}%
\pgfpathlineto{\pgfqpoint{1.335766in}{1.297919in}}%
\pgfpathlineto{\pgfqpoint{1.339173in}{1.311167in}}%
\pgfpathlineto{\pgfqpoint{1.332163in}{1.324415in}}%
\pgfpathlineto{\pgfqpoint{1.329968in}{1.337663in}}%
\pgfpathlineto{\pgfqpoint{1.330653in}{1.350911in}}%
\pgfpathlineto{\pgfqpoint{1.324753in}{1.361834in}}%
\pgfpathlineto{\pgfqpoint{1.323669in}{1.364159in}}%
\pgfpathlineto{\pgfqpoint{1.324753in}{1.367103in}}%
\pgfpathlineto{\pgfqpoint{1.328874in}{1.377407in}}%
\pgfpathlineto{\pgfqpoint{1.324753in}{1.385175in}}%
\pgfpathlineto{\pgfqpoint{1.322297in}{1.390655in}}%
\pgfpathlineto{\pgfqpoint{1.324753in}{1.398109in}}%
\pgfpathlineto{\pgfqpoint{1.326884in}{1.403903in}}%
\pgfpathlineto{\pgfqpoint{1.324753in}{1.412910in}}%
\pgfpathlineto{\pgfqpoint{1.323915in}{1.417151in}}%
\pgfpathlineto{\pgfqpoint{1.324201in}{1.430399in}}%
\pgfpathlineto{\pgfqpoint{1.324753in}{1.434439in}}%
\pgfpathlineto{\pgfqpoint{1.326065in}{1.443647in}}%
\pgfpathlineto{\pgfqpoint{1.324753in}{1.448110in}}%
\pgfpathlineto{\pgfqpoint{1.322462in}{1.456895in}}%
\pgfpathlineto{\pgfqpoint{1.324753in}{1.463481in}}%
\pgfpathlineto{\pgfqpoint{1.327227in}{1.470143in}}%
\pgfpathlineto{\pgfqpoint{1.324753in}{1.476868in}}%
\pgfpathlineto{\pgfqpoint{1.322665in}{1.483391in}}%
\pgfpathlineto{\pgfqpoint{1.324753in}{1.489967in}}%
\pgfpathlineto{\pgfqpoint{1.327068in}{1.496639in}}%
\pgfpathlineto{\pgfqpoint{1.325457in}{1.509887in}}%
\pgfpathlineto{\pgfqpoint{1.327362in}{1.523135in}}%
\pgfpathlineto{\pgfqpoint{1.329845in}{1.536383in}}%
\pgfpathlineto{\pgfqpoint{1.328453in}{1.549631in}}%
\pgfpathlineto{\pgfqpoint{1.333370in}{1.562879in}}%
\pgfpathlineto{\pgfqpoint{1.330599in}{1.576127in}}%
\pgfpathlineto{\pgfqpoint{1.335331in}{1.589375in}}%
\pgfpathlineto{\pgfqpoint{1.333864in}{1.602623in}}%
\pgfpathlineto{\pgfqpoint{1.336181in}{1.615871in}}%
\pgfpathlineto{\pgfqpoint{1.338110in}{1.629119in}}%
\pgfpathlineto{\pgfqpoint{1.337898in}{1.642367in}}%
\pgfpathlineto{\pgfqpoint{1.342461in}{1.655615in}}%
\pgfpathlineto{\pgfqpoint{1.340907in}{1.668863in}}%
\pgfpathlineto{\pgfqpoint{1.345796in}{1.682111in}}%
\pgfpathlineto{\pgfqpoint{1.344992in}{1.695359in}}%
\pgfpathlineto{\pgfqpoint{1.348236in}{1.708607in}}%
\pgfpathlineto{\pgfqpoint{1.349775in}{1.721855in}}%
\pgfpathlineto{\pgfqpoint{1.350036in}{1.727007in}}%
\pgfpathlineto{\pgfqpoint{1.350495in}{1.735103in}}%
\pgfpathlineto{\pgfqpoint{1.354936in}{1.748351in}}%
\pgfpathlineto{\pgfqpoint{1.354027in}{1.761599in}}%
\pgfpathlineto{\pgfqpoint{1.359404in}{1.774847in}}%
\pgfpathlineto{\pgfqpoint{1.358563in}{1.788095in}}%
\pgfpathlineto{\pgfqpoint{1.362550in}{1.801343in}}%
\pgfpathlineto{\pgfqpoint{1.363721in}{1.814591in}}%
\pgfpathlineto{\pgfqpoint{1.365507in}{1.827839in}}%
\pgfpathlineto{\pgfqpoint{1.369362in}{1.841087in}}%
\pgfpathlineto{\pgfqpoint{1.369245in}{1.854335in}}%
\pgfpathlineto{\pgfqpoint{1.374548in}{1.867583in}}%
\pgfpathlineto{\pgfqpoint{1.374238in}{1.880831in}}%
\pgfpathlineto{\pgfqpoint{1.375319in}{1.883865in}}%
\pgfpathlineto{\pgfqpoint{1.379784in}{1.894079in}}%
\pgfpathlineto{\pgfqpoint{1.380622in}{1.907327in}}%
\pgfpathlineto{\pgfqpoint{1.383664in}{1.920575in}}%
\pgfpathlineto{\pgfqpoint{1.387063in}{1.933823in}}%
\pgfpathlineto{\pgfqpoint{1.387681in}{1.947071in}}%
\pgfpathlineto{\pgfqpoint{1.393534in}{1.960319in}}%
\pgfpathlineto{\pgfqpoint{1.393183in}{1.973567in}}%
\pgfpathlineto{\pgfqpoint{1.399309in}{1.986815in}}%
\pgfpathlineto{\pgfqpoint{1.400116in}{2.000063in}}%
\pgfpathlineto{\pgfqpoint{1.400603in}{2.001509in}}%
\pgfpathlineto{\pgfqpoint{1.406589in}{2.013311in}}%
\pgfpathlineto{\pgfqpoint{1.410520in}{2.026559in}}%
\pgfpathlineto{\pgfqpoint{1.413401in}{2.039807in}}%
\pgfpathlineto{\pgfqpoint{1.420491in}{2.053055in}}%
\pgfpathlineto{\pgfqpoint{1.420817in}{2.066303in}}%
\pgfpathlineto{\pgfqpoint{1.425886in}{2.073616in}}%
\pgfpathlineto{\pgfqpoint{1.436690in}{2.079551in}}%
\pgfpathlineto{\pgfqpoint{1.436805in}{2.092799in}}%
\pgfpathlineto{\pgfqpoint{1.451169in}{2.101734in}}%
\pgfpathlineto{\pgfqpoint{1.469675in}{2.092799in}}%
\pgfpathlineto{\pgfqpoint{1.469801in}{2.079551in}}%
\pgfpathlineto{\pgfqpoint{1.476452in}{2.076466in}}%
\pgfpathlineto{\pgfqpoint{1.483260in}{2.066303in}}%
\pgfpathlineto{\pgfqpoint{1.486095in}{2.053055in}}%
\pgfpathlineto{\pgfqpoint{1.490737in}{2.039807in}}%
\pgfpathlineto{\pgfqpoint{1.497892in}{2.026559in}}%
\pgfpathlineto{\pgfqpoint{1.498338in}{2.013311in}}%
\pgfpathlineto{\pgfqpoint{1.501735in}{2.008631in}}%
\pgfpathlineto{\pgfqpoint{1.505828in}{2.000063in}}%
\pgfpathlineto{\pgfqpoint{1.505777in}{1.986815in}}%
\pgfpathlineto{\pgfqpoint{1.511002in}{1.973567in}}%
\pgfpathlineto{\pgfqpoint{1.512797in}{1.960319in}}%
\pgfpathlineto{\pgfqpoint{1.515485in}{1.947071in}}%
\pgfpathlineto{\pgfqpoint{1.520162in}{1.933823in}}%
\pgfpathlineto{\pgfqpoint{1.520106in}{1.920575in}}%
\pgfpathlineto{\pgfqpoint{1.526308in}{1.907327in}}%
\pgfpathlineto{\pgfqpoint{1.525745in}{1.894079in}}%
\pgfpathlineto{\pgfqpoint{1.527018in}{1.891071in}}%
\pgfpathlineto{\pgfqpoint{1.530486in}{1.880831in}}%
\pgfpathlineto{\pgfqpoint{1.531927in}{1.867583in}}%
\pgfpathlineto{\pgfqpoint{1.534026in}{1.854335in}}%
\pgfpathlineto{\pgfqpoint{1.537823in}{1.841087in}}%
\pgfpathlineto{\pgfqpoint{1.537575in}{1.827839in}}%
\pgfpathlineto{\pgfqpoint{1.542796in}{1.814591in}}%
\pgfpathlineto{\pgfqpoint{1.541851in}{1.801343in}}%
\pgfpathlineto{\pgfqpoint{1.546041in}{1.788095in}}%
\pgfpathlineto{\pgfqpoint{1.546702in}{1.774847in}}%
\pgfpathlineto{\pgfqpoint{1.548664in}{1.761599in}}%
\pgfpathlineto{\pgfqpoint{1.552302in}{1.748527in}}%
\pgfpathlineto{\pgfqpoint{1.552347in}{1.748351in}}%
\pgfpathlineto{\pgfqpoint{1.552302in}{1.746493in}}%
\pgfpathlineto{\pgfqpoint{1.551989in}{1.735103in}}%
\pgfpathlineto{\pgfqpoint{1.552302in}{1.734343in}}%
\pgfpathlineto{\pgfqpoint{1.556928in}{1.721855in}}%
\pgfpathlineto{\pgfqpoint{1.555755in}{1.708607in}}%
\pgfpathlineto{\pgfqpoint{1.560232in}{1.695359in}}%
\pgfpathlineto{\pgfqpoint{1.560292in}{1.682111in}}%
\pgfpathlineto{\pgfqpoint{1.561880in}{1.668863in}}%
\pgfpathlineto{\pgfqpoint{1.564409in}{1.655615in}}%
\pgfpathlineto{\pgfqpoint{1.563270in}{1.642367in}}%
\pgfpathlineto{\pgfqpoint{1.568005in}{1.629119in}}%
\pgfpathlineto{\pgfqpoint{1.566066in}{1.615871in}}%
\pgfpathlineto{\pgfqpoint{1.570705in}{1.602623in}}%
\pgfpathlineto{\pgfqpoint{1.569983in}{1.589375in}}%
\pgfpathlineto{\pgfqpoint{1.572285in}{1.576127in}}%
\pgfpathlineto{\pgfqpoint{1.574435in}{1.562879in}}%
\pgfpathlineto{\pgfqpoint{1.573050in}{1.549631in}}%
\pgfpathlineto{\pgfqpoint{1.577576in}{1.536383in}}%
\pgfpathlineto{\pgfqpoint{1.574007in}{1.523135in}}%
\pgfpathlineto{\pgfqpoint{1.577585in}{1.513546in}}%
\pgfpathlineto{\pgfqpoint{1.578845in}{1.509887in}}%
\pgfpathlineto{\pgfqpoint{1.577585in}{1.502167in}}%
\pgfpathlineto{\pgfqpoint{1.576542in}{1.496639in}}%
\pgfpathlineto{\pgfqpoint{1.577585in}{1.490436in}}%
\pgfpathlineto{\pgfqpoint{1.578664in}{1.483391in}}%
\pgfpathlineto{\pgfqpoint{1.579257in}{1.470143in}}%
\pgfpathlineto{\pgfqpoint{1.577585in}{1.457953in}}%
\pgfpathlineto{\pgfqpoint{1.577423in}{1.456895in}}%
\pgfpathlineto{\pgfqpoint{1.577585in}{1.456350in}}%
\pgfpathlineto{\pgfqpoint{1.581149in}{1.443647in}}%
\pgfpathlineto{\pgfqpoint{1.577585in}{1.433384in}}%
\pgfpathlineto{\pgfqpoint{1.576389in}{1.430399in}}%
\pgfpathlineto{\pgfqpoint{1.577585in}{1.427276in}}%
\pgfpathlineto{\pgfqpoint{1.581121in}{1.417151in}}%
\pgfpathlineto{\pgfqpoint{1.577585in}{1.407030in}}%
\pgfpathlineto{\pgfqpoint{1.576283in}{1.403903in}}%
\pgfpathlineto{\pgfqpoint{1.577585in}{1.398104in}}%
\pgfpathlineto{\pgfqpoint{1.579075in}{1.390655in}}%
\pgfpathlineto{\pgfqpoint{1.577585in}{1.378632in}}%
\pgfpathlineto{\pgfqpoint{1.577401in}{1.377407in}}%
\pgfpathlineto{\pgfqpoint{1.574745in}{1.364159in}}%
\pgfpathlineto{\pgfqpoint{1.576854in}{1.350911in}}%
\pgfpathlineto{\pgfqpoint{1.568129in}{1.337663in}}%
\pgfpathlineto{\pgfqpoint{1.572437in}{1.324415in}}%
\pgfpathlineto{\pgfqpoint{1.561597in}{1.311167in}}%
\pgfpathlineto{\pgfqpoint{1.564055in}{1.297919in}}%
\pgfpathlineto{\pgfqpoint{1.555772in}{1.284671in}}%
\pgfpathlineto{\pgfqpoint{1.552604in}{1.271423in}}%
\pgfpathlineto{\pgfqpoint{1.552302in}{1.269620in}}%
\pgfpathlineto{\pgfqpoint{1.549703in}{1.258175in}}%
\pgfpathlineto{\pgfqpoint{1.535040in}{1.244927in}}%
\pgfpathlineto{\pgfqpoint{1.536741in}{1.231679in}}%
\pgfpathlineto{\pgfqpoint{1.527018in}{1.226388in}}%
\pgfpathlineto{\pgfqpoint{1.506017in}{1.218431in}}%
\pgfpathlineto{\pgfqpoint{1.508847in}{1.205183in}}%
\pgfpathlineto{\pgfqpoint{1.501735in}{1.202497in}}%
\pgfpathlineto{\pgfqpoint{1.476452in}{1.197064in}}%
\pgfpathlineto{\pgfqpoint{1.451169in}{1.195366in}}%
\pgfpathlineto{\pgfqpoint{1.425886in}{1.196991in}}%
\pgfpathlineto{\pgfqpoint{1.400603in}{1.201993in}}%
\pgfpathlineto{\pgfqpoint{1.391241in}{1.205183in}}%
\pgfpathclose%
\pgfusepath{fill}%
\end{pgfscope}%
\begin{pgfscope}%
\pgfpathrectangle{\pgfqpoint{0.540975in}{0.476543in}}{\pgfqpoint{1.820388in}{1.656000in}}%
\pgfusepath{clip}%
\pgfsetbuttcap%
\pgfsetroundjoin%
\definecolor{currentfill}{rgb}{0.992541,0.701576,0.398462}%
\pgfsetfillcolor{currentfill}%
\pgfsetlinewidth{0.000000pt}%
\definecolor{currentstroke}{rgb}{0.000000,0.000000,0.000000}%
\pgfsetstrokecolor{currentstroke}%
\pgfsetdash{}{0pt}%
\pgfpathmoveto{\pgfqpoint{1.400603in}{1.201993in}}%
\pgfpathlineto{\pgfqpoint{1.425886in}{1.196991in}}%
\pgfpathlineto{\pgfqpoint{1.451169in}{1.195366in}}%
\pgfpathlineto{\pgfqpoint{1.476452in}{1.197064in}}%
\pgfpathlineto{\pgfqpoint{1.501735in}{1.202497in}}%
\pgfpathlineto{\pgfqpoint{1.508847in}{1.205183in}}%
\pgfpathlineto{\pgfqpoint{1.506017in}{1.218431in}}%
\pgfpathlineto{\pgfqpoint{1.527018in}{1.226388in}}%
\pgfpathlineto{\pgfqpoint{1.536741in}{1.231679in}}%
\pgfpathlineto{\pgfqpoint{1.535040in}{1.244927in}}%
\pgfpathlineto{\pgfqpoint{1.549703in}{1.258175in}}%
\pgfpathlineto{\pgfqpoint{1.552302in}{1.269620in}}%
\pgfpathlineto{\pgfqpoint{1.552604in}{1.271423in}}%
\pgfpathlineto{\pgfqpoint{1.555772in}{1.284671in}}%
\pgfpathlineto{\pgfqpoint{1.564055in}{1.297919in}}%
\pgfpathlineto{\pgfqpoint{1.561597in}{1.311167in}}%
\pgfpathlineto{\pgfqpoint{1.572437in}{1.324415in}}%
\pgfpathlineto{\pgfqpoint{1.568129in}{1.337663in}}%
\pgfpathlineto{\pgfqpoint{1.576854in}{1.350911in}}%
\pgfpathlineto{\pgfqpoint{1.574745in}{1.364159in}}%
\pgfpathlineto{\pgfqpoint{1.577401in}{1.377407in}}%
\pgfpathlineto{\pgfqpoint{1.577585in}{1.378632in}}%
\pgfpathlineto{\pgfqpoint{1.579075in}{1.390655in}}%
\pgfpathlineto{\pgfqpoint{1.577585in}{1.398104in}}%
\pgfpathlineto{\pgfqpoint{1.576283in}{1.403903in}}%
\pgfpathlineto{\pgfqpoint{1.577585in}{1.407030in}}%
\pgfpathlineto{\pgfqpoint{1.581121in}{1.417151in}}%
\pgfpathlineto{\pgfqpoint{1.577585in}{1.427276in}}%
\pgfpathlineto{\pgfqpoint{1.576389in}{1.430399in}}%
\pgfpathlineto{\pgfqpoint{1.577585in}{1.433384in}}%
\pgfpathlineto{\pgfqpoint{1.581149in}{1.443647in}}%
\pgfpathlineto{\pgfqpoint{1.577585in}{1.456350in}}%
\pgfpathlineto{\pgfqpoint{1.577423in}{1.456895in}}%
\pgfpathlineto{\pgfqpoint{1.577585in}{1.457953in}}%
\pgfpathlineto{\pgfqpoint{1.579257in}{1.470143in}}%
\pgfpathlineto{\pgfqpoint{1.578664in}{1.483391in}}%
\pgfpathlineto{\pgfqpoint{1.577585in}{1.490436in}}%
\pgfpathlineto{\pgfqpoint{1.576542in}{1.496639in}}%
\pgfpathlineto{\pgfqpoint{1.577585in}{1.502167in}}%
\pgfpathlineto{\pgfqpoint{1.578845in}{1.509887in}}%
\pgfpathlineto{\pgfqpoint{1.577585in}{1.513546in}}%
\pgfpathlineto{\pgfqpoint{1.574007in}{1.523135in}}%
\pgfpathlineto{\pgfqpoint{1.577576in}{1.536383in}}%
\pgfpathlineto{\pgfqpoint{1.573050in}{1.549631in}}%
\pgfpathlineto{\pgfqpoint{1.574435in}{1.562879in}}%
\pgfpathlineto{\pgfqpoint{1.572285in}{1.576127in}}%
\pgfpathlineto{\pgfqpoint{1.569983in}{1.589375in}}%
\pgfpathlineto{\pgfqpoint{1.570705in}{1.602623in}}%
\pgfpathlineto{\pgfqpoint{1.566066in}{1.615871in}}%
\pgfpathlineto{\pgfqpoint{1.568005in}{1.629119in}}%
\pgfpathlineto{\pgfqpoint{1.563270in}{1.642367in}}%
\pgfpathlineto{\pgfqpoint{1.564409in}{1.655615in}}%
\pgfpathlineto{\pgfqpoint{1.561880in}{1.668863in}}%
\pgfpathlineto{\pgfqpoint{1.560292in}{1.682111in}}%
\pgfpathlineto{\pgfqpoint{1.560232in}{1.695359in}}%
\pgfpathlineto{\pgfqpoint{1.555755in}{1.708607in}}%
\pgfpathlineto{\pgfqpoint{1.556928in}{1.721855in}}%
\pgfpathlineto{\pgfqpoint{1.552302in}{1.734343in}}%
\pgfpathlineto{\pgfqpoint{1.551989in}{1.735103in}}%
\pgfpathlineto{\pgfqpoint{1.552302in}{1.746493in}}%
\pgfpathlineto{\pgfqpoint{1.552347in}{1.748351in}}%
\pgfpathlineto{\pgfqpoint{1.552302in}{1.748527in}}%
\pgfpathlineto{\pgfqpoint{1.548664in}{1.761599in}}%
\pgfpathlineto{\pgfqpoint{1.546702in}{1.774847in}}%
\pgfpathlineto{\pgfqpoint{1.546041in}{1.788095in}}%
\pgfpathlineto{\pgfqpoint{1.541851in}{1.801343in}}%
\pgfpathlineto{\pgfqpoint{1.542796in}{1.814591in}}%
\pgfpathlineto{\pgfqpoint{1.537575in}{1.827839in}}%
\pgfpathlineto{\pgfqpoint{1.537823in}{1.841087in}}%
\pgfpathlineto{\pgfqpoint{1.534026in}{1.854335in}}%
\pgfpathlineto{\pgfqpoint{1.531927in}{1.867583in}}%
\pgfpathlineto{\pgfqpoint{1.530486in}{1.880831in}}%
\pgfpathlineto{\pgfqpoint{1.527018in}{1.891071in}}%
\pgfpathlineto{\pgfqpoint{1.525745in}{1.894079in}}%
\pgfpathlineto{\pgfqpoint{1.526308in}{1.907327in}}%
\pgfpathlineto{\pgfqpoint{1.520106in}{1.920575in}}%
\pgfpathlineto{\pgfqpoint{1.520162in}{1.933823in}}%
\pgfpathlineto{\pgfqpoint{1.515485in}{1.947071in}}%
\pgfpathlineto{\pgfqpoint{1.512797in}{1.960319in}}%
\pgfpathlineto{\pgfqpoint{1.511002in}{1.973567in}}%
\pgfpathlineto{\pgfqpoint{1.505777in}{1.986815in}}%
\pgfpathlineto{\pgfqpoint{1.505828in}{2.000063in}}%
\pgfpathlineto{\pgfqpoint{1.501735in}{2.008631in}}%
\pgfpathlineto{\pgfqpoint{1.498338in}{2.013311in}}%
\pgfpathlineto{\pgfqpoint{1.497892in}{2.026559in}}%
\pgfpathlineto{\pgfqpoint{1.490737in}{2.039807in}}%
\pgfpathlineto{\pgfqpoint{1.486095in}{2.053055in}}%
\pgfpathlineto{\pgfqpoint{1.483260in}{2.066303in}}%
\pgfpathlineto{\pgfqpoint{1.476452in}{2.076466in}}%
\pgfpathlineto{\pgfqpoint{1.469801in}{2.079551in}}%
\pgfpathlineto{\pgfqpoint{1.469675in}{2.092799in}}%
\pgfpathlineto{\pgfqpoint{1.451169in}{2.101734in}}%
\pgfpathlineto{\pgfqpoint{1.436805in}{2.092799in}}%
\pgfpathlineto{\pgfqpoint{1.436690in}{2.079551in}}%
\pgfpathlineto{\pgfqpoint{1.425886in}{2.073616in}}%
\pgfpathlineto{\pgfqpoint{1.420817in}{2.066303in}}%
\pgfpathlineto{\pgfqpoint{1.420491in}{2.053055in}}%
\pgfpathlineto{\pgfqpoint{1.413401in}{2.039807in}}%
\pgfpathlineto{\pgfqpoint{1.410520in}{2.026559in}}%
\pgfpathlineto{\pgfqpoint{1.406589in}{2.013311in}}%
\pgfpathlineto{\pgfqpoint{1.400603in}{2.001509in}}%
\pgfpathlineto{\pgfqpoint{1.400116in}{2.000063in}}%
\pgfpathlineto{\pgfqpoint{1.399309in}{1.986815in}}%
\pgfpathlineto{\pgfqpoint{1.393183in}{1.973567in}}%
\pgfpathlineto{\pgfqpoint{1.393534in}{1.960319in}}%
\pgfpathlineto{\pgfqpoint{1.387681in}{1.947071in}}%
\pgfpathlineto{\pgfqpoint{1.387063in}{1.933823in}}%
\pgfpathlineto{\pgfqpoint{1.383664in}{1.920575in}}%
\pgfpathlineto{\pgfqpoint{1.380622in}{1.907327in}}%
\pgfpathlineto{\pgfqpoint{1.379784in}{1.894079in}}%
\pgfpathlineto{\pgfqpoint{1.375319in}{1.883865in}}%
\pgfpathlineto{\pgfqpoint{1.374238in}{1.880831in}}%
\pgfpathlineto{\pgfqpoint{1.374548in}{1.867583in}}%
\pgfpathlineto{\pgfqpoint{1.369245in}{1.854335in}}%
\pgfpathlineto{\pgfqpoint{1.369362in}{1.841087in}}%
\pgfpathlineto{\pgfqpoint{1.365507in}{1.827839in}}%
\pgfpathlineto{\pgfqpoint{1.363721in}{1.814591in}}%
\pgfpathlineto{\pgfqpoint{1.362550in}{1.801343in}}%
\pgfpathlineto{\pgfqpoint{1.358563in}{1.788095in}}%
\pgfpathlineto{\pgfqpoint{1.359404in}{1.774847in}}%
\pgfpathlineto{\pgfqpoint{1.354027in}{1.761599in}}%
\pgfpathlineto{\pgfqpoint{1.354936in}{1.748351in}}%
\pgfpathlineto{\pgfqpoint{1.350495in}{1.735103in}}%
\pgfpathlineto{\pgfqpoint{1.350036in}{1.727007in}}%
\pgfpathlineto{\pgfqpoint{1.349775in}{1.721855in}}%
\pgfpathlineto{\pgfqpoint{1.348236in}{1.708607in}}%
\pgfpathlineto{\pgfqpoint{1.344992in}{1.695359in}}%
\pgfpathlineto{\pgfqpoint{1.345796in}{1.682111in}}%
\pgfpathlineto{\pgfqpoint{1.340907in}{1.668863in}}%
\pgfpathlineto{\pgfqpoint{1.342461in}{1.655615in}}%
\pgfpathlineto{\pgfqpoint{1.337898in}{1.642367in}}%
\pgfpathlineto{\pgfqpoint{1.338110in}{1.629119in}}%
\pgfpathlineto{\pgfqpoint{1.336181in}{1.615871in}}%
\pgfpathlineto{\pgfqpoint{1.333864in}{1.602623in}}%
\pgfpathlineto{\pgfqpoint{1.335331in}{1.589375in}}%
\pgfpathlineto{\pgfqpoint{1.330599in}{1.576127in}}%
\pgfpathlineto{\pgfqpoint{1.333370in}{1.562879in}}%
\pgfpathlineto{\pgfqpoint{1.328453in}{1.549631in}}%
\pgfpathlineto{\pgfqpoint{1.329845in}{1.536383in}}%
\pgfpathlineto{\pgfqpoint{1.327362in}{1.523135in}}%
\pgfpathlineto{\pgfqpoint{1.325457in}{1.509887in}}%
\pgfpathlineto{\pgfqpoint{1.327068in}{1.496639in}}%
\pgfpathlineto{\pgfqpoint{1.324753in}{1.489967in}}%
\pgfpathlineto{\pgfqpoint{1.322665in}{1.483391in}}%
\pgfpathlineto{\pgfqpoint{1.324753in}{1.476868in}}%
\pgfpathlineto{\pgfqpoint{1.327227in}{1.470143in}}%
\pgfpathlineto{\pgfqpoint{1.324753in}{1.463481in}}%
\pgfpathlineto{\pgfqpoint{1.322462in}{1.456895in}}%
\pgfpathlineto{\pgfqpoint{1.324753in}{1.448110in}}%
\pgfpathlineto{\pgfqpoint{1.326065in}{1.443647in}}%
\pgfpathlineto{\pgfqpoint{1.324753in}{1.434439in}}%
\pgfpathlineto{\pgfqpoint{1.324201in}{1.430399in}}%
\pgfpathlineto{\pgfqpoint{1.323915in}{1.417151in}}%
\pgfpathlineto{\pgfqpoint{1.324753in}{1.412910in}}%
\pgfpathlineto{\pgfqpoint{1.326884in}{1.403903in}}%
\pgfpathlineto{\pgfqpoint{1.324753in}{1.398109in}}%
\pgfpathlineto{\pgfqpoint{1.322297in}{1.390655in}}%
\pgfpathlineto{\pgfqpoint{1.324753in}{1.385175in}}%
\pgfpathlineto{\pgfqpoint{1.328874in}{1.377407in}}%
\pgfpathlineto{\pgfqpoint{1.324753in}{1.367103in}}%
\pgfpathlineto{\pgfqpoint{1.323669in}{1.364159in}}%
\pgfpathlineto{\pgfqpoint{1.324753in}{1.361834in}}%
\pgfpathlineto{\pgfqpoint{1.330653in}{1.350911in}}%
\pgfpathlineto{\pgfqpoint{1.329968in}{1.337663in}}%
\pgfpathlineto{\pgfqpoint{1.332163in}{1.324415in}}%
\pgfpathlineto{\pgfqpoint{1.339173in}{1.311167in}}%
\pgfpathlineto{\pgfqpoint{1.335766in}{1.297919in}}%
\pgfpathlineto{\pgfqpoint{1.348852in}{1.284671in}}%
\pgfpathlineto{\pgfqpoint{1.344242in}{1.271423in}}%
\pgfpathlineto{\pgfqpoint{1.350036in}{1.265077in}}%
\pgfpathlineto{\pgfqpoint{1.357891in}{1.258175in}}%
\pgfpathlineto{\pgfqpoint{1.361858in}{1.244927in}}%
\pgfpathlineto{\pgfqpoint{1.369305in}{1.231679in}}%
\pgfpathlineto{\pgfqpoint{1.375319in}{1.227266in}}%
\pgfpathlineto{\pgfqpoint{1.393987in}{1.218431in}}%
\pgfpathlineto{\pgfqpoint{1.391241in}{1.205183in}}%
\pgfpathlineto{\pgfqpoint{1.400603in}{1.201993in}}%
\pgfpathclose%
\pgfpathmoveto{\pgfqpoint{1.415378in}{1.231679in}}%
\pgfpathlineto{\pgfqpoint{1.401504in}{1.244927in}}%
\pgfpathlineto{\pgfqpoint{1.400603in}{1.245513in}}%
\pgfpathlineto{\pgfqpoint{1.390549in}{1.258175in}}%
\pgfpathlineto{\pgfqpoint{1.375319in}{1.268058in}}%
\pgfpathlineto{\pgfqpoint{1.371468in}{1.271423in}}%
\pgfpathlineto{\pgfqpoint{1.375319in}{1.284295in}}%
\pgfpathlineto{\pgfqpoint{1.375461in}{1.284671in}}%
\pgfpathlineto{\pgfqpoint{1.375319in}{1.284746in}}%
\pgfpathlineto{\pgfqpoint{1.357760in}{1.297919in}}%
\pgfpathlineto{\pgfqpoint{1.361627in}{1.311167in}}%
\pgfpathlineto{\pgfqpoint{1.350815in}{1.324415in}}%
\pgfpathlineto{\pgfqpoint{1.350036in}{1.328157in}}%
\pgfpathlineto{\pgfqpoint{1.348409in}{1.337663in}}%
\pgfpathlineto{\pgfqpoint{1.347626in}{1.350911in}}%
\pgfpathlineto{\pgfqpoint{1.340145in}{1.364159in}}%
\pgfpathlineto{\pgfqpoint{1.344785in}{1.377407in}}%
\pgfpathlineto{\pgfqpoint{1.337092in}{1.390655in}}%
\pgfpathlineto{\pgfqpoint{1.341885in}{1.403903in}}%
\pgfpathlineto{\pgfqpoint{1.337745in}{1.417151in}}%
\pgfpathlineto{\pgfqpoint{1.338117in}{1.430399in}}%
\pgfpathlineto{\pgfqpoint{1.339124in}{1.443647in}}%
\pgfpathlineto{\pgfqpoint{1.335058in}{1.456895in}}%
\pgfpathlineto{\pgfqpoint{1.339588in}{1.470143in}}%
\pgfpathlineto{\pgfqpoint{1.334281in}{1.483391in}}%
\pgfpathlineto{\pgfqpoint{1.338820in}{1.496639in}}%
\pgfpathlineto{\pgfqpoint{1.336504in}{1.509887in}}%
\pgfpathlineto{\pgfqpoint{1.338473in}{1.523135in}}%
\pgfpathlineto{\pgfqpoint{1.340342in}{1.536383in}}%
\pgfpathlineto{\pgfqpoint{1.338868in}{1.549631in}}%
\pgfpathlineto{\pgfqpoint{1.343475in}{1.562879in}}%
\pgfpathlineto{\pgfqpoint{1.340365in}{1.576127in}}%
\pgfpathlineto{\pgfqpoint{1.345077in}{1.589375in}}%
\pgfpathlineto{\pgfqpoint{1.343123in}{1.602623in}}%
\pgfpathlineto{\pgfqpoint{1.345535in}{1.615871in}}%
\pgfpathlineto{\pgfqpoint{1.347035in}{1.629119in}}%
\pgfpathlineto{\pgfqpoint{1.346798in}{1.642367in}}%
\pgfpathlineto{\pgfqpoint{1.350036in}{1.652180in}}%
\pgfpathlineto{\pgfqpoint{1.351283in}{1.655615in}}%
\pgfpathlineto{\pgfqpoint{1.350036in}{1.663739in}}%
\pgfpathlineto{\pgfqpoint{1.349357in}{1.668863in}}%
\pgfpathlineto{\pgfqpoint{1.350036in}{1.670701in}}%
\pgfpathlineto{\pgfqpoint{1.354769in}{1.682111in}}%
\pgfpathlineto{\pgfqpoint{1.353462in}{1.695359in}}%
\pgfpathlineto{\pgfqpoint{1.357217in}{1.708607in}}%
\pgfpathlineto{\pgfqpoint{1.358506in}{1.721855in}}%
\pgfpathlineto{\pgfqpoint{1.359360in}{1.735103in}}%
\pgfpathlineto{\pgfqpoint{1.363536in}{1.748351in}}%
\pgfpathlineto{\pgfqpoint{1.362524in}{1.761599in}}%
\pgfpathlineto{\pgfqpoint{1.367867in}{1.774847in}}%
\pgfpathlineto{\pgfqpoint{1.366707in}{1.788095in}}%
\pgfpathlineto{\pgfqpoint{1.370808in}{1.801343in}}%
\pgfpathlineto{\pgfqpoint{1.371627in}{1.814591in}}%
\pgfpathlineto{\pgfqpoint{1.373523in}{1.827839in}}%
\pgfpathlineto{\pgfqpoint{1.375319in}{1.834267in}}%
\pgfpathlineto{\pgfqpoint{1.377567in}{1.841087in}}%
\pgfpathlineto{\pgfqpoint{1.377407in}{1.854335in}}%
\pgfpathlineto{\pgfqpoint{1.383788in}{1.867583in}}%
\pgfpathlineto{\pgfqpoint{1.383189in}{1.880831in}}%
\pgfpathlineto{\pgfqpoint{1.389115in}{1.894079in}}%
\pgfpathlineto{\pgfqpoint{1.389542in}{1.907327in}}%
\pgfpathlineto{\pgfqpoint{1.392812in}{1.920575in}}%
\pgfpathlineto{\pgfqpoint{1.395868in}{1.933823in}}%
\pgfpathlineto{\pgfqpoint{1.396574in}{1.947071in}}%
\pgfpathlineto{\pgfqpoint{1.400603in}{1.956302in}}%
\pgfpathlineto{\pgfqpoint{1.403139in}{1.960319in}}%
\pgfpathlineto{\pgfqpoint{1.402428in}{1.973567in}}%
\pgfpathlineto{\pgfqpoint{1.411610in}{1.986815in}}%
\pgfpathlineto{\pgfqpoint{1.412225in}{2.000063in}}%
\pgfpathlineto{\pgfqpoint{1.419405in}{2.013311in}}%
\pgfpathlineto{\pgfqpoint{1.422646in}{2.026559in}}%
\pgfpathlineto{\pgfqpoint{1.425886in}{2.039404in}}%
\pgfpathlineto{\pgfqpoint{1.426171in}{2.039807in}}%
\pgfpathlineto{\pgfqpoint{1.443016in}{2.053055in}}%
\pgfpathlineto{\pgfqpoint{1.445241in}{2.066303in}}%
\pgfpathlineto{\pgfqpoint{1.451169in}{2.069562in}}%
\pgfpathlineto{\pgfqpoint{1.458137in}{2.066303in}}%
\pgfpathlineto{\pgfqpoint{1.463152in}{2.053055in}}%
\pgfpathlineto{\pgfqpoint{1.476452in}{2.042833in}}%
\pgfpathlineto{\pgfqpoint{1.477601in}{2.039807in}}%
\pgfpathlineto{\pgfqpoint{1.484155in}{2.026559in}}%
\pgfpathlineto{\pgfqpoint{1.484841in}{2.013311in}}%
\pgfpathlineto{\pgfqpoint{1.494428in}{2.000063in}}%
\pgfpathlineto{\pgfqpoint{1.494019in}{1.986815in}}%
\pgfpathlineto{\pgfqpoint{1.501735in}{1.974583in}}%
\pgfpathlineto{\pgfqpoint{1.502145in}{1.973567in}}%
\pgfpathlineto{\pgfqpoint{1.503584in}{1.960319in}}%
\pgfpathlineto{\pgfqpoint{1.506488in}{1.947071in}}%
\pgfpathlineto{\pgfqpoint{1.510882in}{1.933823in}}%
\pgfpathlineto{\pgfqpoint{1.510799in}{1.920575in}}%
\pgfpathlineto{\pgfqpoint{1.516955in}{1.907327in}}%
\pgfpathlineto{\pgfqpoint{1.516094in}{1.894079in}}%
\pgfpathlineto{\pgfqpoint{1.521810in}{1.880831in}}%
\pgfpathlineto{\pgfqpoint{1.523214in}{1.867583in}}%
\pgfpathlineto{\pgfqpoint{1.525979in}{1.854335in}}%
\pgfpathlineto{\pgfqpoint{1.527018in}{1.851186in}}%
\pgfpathlineto{\pgfqpoint{1.529795in}{1.841087in}}%
\pgfpathlineto{\pgfqpoint{1.529499in}{1.827839in}}%
\pgfpathlineto{\pgfqpoint{1.534690in}{1.814591in}}%
\pgfpathlineto{\pgfqpoint{1.533478in}{1.801343in}}%
\pgfpathlineto{\pgfqpoint{1.537767in}{1.788095in}}%
\pgfpathlineto{\pgfqpoint{1.538039in}{1.774847in}}%
\pgfpathlineto{\pgfqpoint{1.540121in}{1.761599in}}%
\pgfpathlineto{\pgfqpoint{1.543490in}{1.748351in}}%
\pgfpathlineto{\pgfqpoint{1.543086in}{1.735103in}}%
\pgfpathlineto{\pgfqpoint{1.548481in}{1.721855in}}%
\pgfpathlineto{\pgfqpoint{1.546868in}{1.708607in}}%
\pgfpathlineto{\pgfqpoint{1.551997in}{1.695359in}}%
\pgfpathlineto{\pgfqpoint{1.551649in}{1.682111in}}%
\pgfpathlineto{\pgfqpoint{1.552302in}{1.677533in}}%
\pgfpathlineto{\pgfqpoint{1.553379in}{1.668863in}}%
\pgfpathlineto{\pgfqpoint{1.555590in}{1.655615in}}%
\pgfpathlineto{\pgfqpoint{1.554364in}{1.642367in}}%
\pgfpathlineto{\pgfqpoint{1.558940in}{1.629119in}}%
\pgfpathlineto{\pgfqpoint{1.556681in}{1.615871in}}%
\pgfpathlineto{\pgfqpoint{1.561336in}{1.602623in}}%
\pgfpathlineto{\pgfqpoint{1.560102in}{1.589375in}}%
\pgfpathlineto{\pgfqpoint{1.562463in}{1.576127in}}%
\pgfpathlineto{\pgfqpoint{1.564097in}{1.562879in}}%
\pgfpathlineto{\pgfqpoint{1.562631in}{1.549631in}}%
\pgfpathlineto{\pgfqpoint{1.566841in}{1.536383in}}%
\pgfpathlineto{\pgfqpoint{1.562910in}{1.523135in}}%
\pgfpathlineto{\pgfqpoint{1.567779in}{1.509887in}}%
\pgfpathlineto{\pgfqpoint{1.564701in}{1.496639in}}%
\pgfpathlineto{\pgfqpoint{1.566953in}{1.483391in}}%
\pgfpathlineto{\pgfqpoint{1.566802in}{1.470143in}}%
\pgfpathlineto{\pgfqpoint{1.564695in}{1.456895in}}%
\pgfpathlineto{\pgfqpoint{1.568139in}{1.443647in}}%
\pgfpathlineto{\pgfqpoint{1.562499in}{1.430399in}}%
\pgfpathlineto{\pgfqpoint{1.567375in}{1.417151in}}%
\pgfpathlineto{\pgfqpoint{1.561191in}{1.403903in}}%
\pgfpathlineto{\pgfqpoint{1.564196in}{1.390655in}}%
\pgfpathlineto{\pgfqpoint{1.561135in}{1.377407in}}%
\pgfpathlineto{\pgfqpoint{1.558429in}{1.364159in}}%
\pgfpathlineto{\pgfqpoint{1.559364in}{1.350911in}}%
\pgfpathlineto{\pgfqpoint{1.552302in}{1.340751in}}%
\pgfpathlineto{\pgfqpoint{1.549636in}{1.337663in}}%
\pgfpathlineto{\pgfqpoint{1.552302in}{1.329252in}}%
\pgfpathlineto{\pgfqpoint{1.553629in}{1.324415in}}%
\pgfpathlineto{\pgfqpoint{1.552302in}{1.322883in}}%
\pgfpathlineto{\pgfqpoint{1.539082in}{1.311167in}}%
\pgfpathlineto{\pgfqpoint{1.541904in}{1.297919in}}%
\pgfpathlineto{\pgfqpoint{1.529076in}{1.284671in}}%
\pgfpathlineto{\pgfqpoint{1.527018in}{1.276931in}}%
\pgfpathlineto{\pgfqpoint{1.525100in}{1.271423in}}%
\pgfpathlineto{\pgfqpoint{1.516668in}{1.258175in}}%
\pgfpathlineto{\pgfqpoint{1.501735in}{1.248225in}}%
\pgfpathlineto{\pgfqpoint{1.494006in}{1.244927in}}%
\pgfpathlineto{\pgfqpoint{1.490811in}{1.231679in}}%
\pgfpathlineto{\pgfqpoint{1.476452in}{1.228391in}}%
\pgfpathlineto{\pgfqpoint{1.451169in}{1.226529in}}%
\pgfpathlineto{\pgfqpoint{1.425886in}{1.228799in}}%
\pgfpathlineto{\pgfqpoint{1.415378in}{1.231679in}}%
\pgfpathclose%
\pgfusepath{fill}%
\end{pgfscope}%
\begin{pgfscope}%
\pgfpathrectangle{\pgfqpoint{0.540975in}{0.476543in}}{\pgfqpoint{1.820388in}{1.656000in}}%
\pgfusepath{clip}%
\pgfsetbuttcap%
\pgfsetroundjoin%
\definecolor{currentfill}{rgb}{0.981776,0.607382,0.345790}%
\pgfsetfillcolor{currentfill}%
\pgfsetlinewidth{0.000000pt}%
\definecolor{currentstroke}{rgb}{0.000000,0.000000,0.000000}%
\pgfsetstrokecolor{currentstroke}%
\pgfsetdash{}{0pt}%
\pgfpathmoveto{\pgfqpoint{1.425886in}{1.228799in}}%
\pgfpathlineto{\pgfqpoint{1.451169in}{1.226529in}}%
\pgfpathlineto{\pgfqpoint{1.476452in}{1.228391in}}%
\pgfpathlineto{\pgfqpoint{1.490811in}{1.231679in}}%
\pgfpathlineto{\pgfqpoint{1.494006in}{1.244927in}}%
\pgfpathlineto{\pgfqpoint{1.501735in}{1.248225in}}%
\pgfpathlineto{\pgfqpoint{1.516668in}{1.258175in}}%
\pgfpathlineto{\pgfqpoint{1.525100in}{1.271423in}}%
\pgfpathlineto{\pgfqpoint{1.527018in}{1.276931in}}%
\pgfpathlineto{\pgfqpoint{1.529076in}{1.284671in}}%
\pgfpathlineto{\pgfqpoint{1.541904in}{1.297919in}}%
\pgfpathlineto{\pgfqpoint{1.539082in}{1.311167in}}%
\pgfpathlineto{\pgfqpoint{1.552302in}{1.322883in}}%
\pgfpathlineto{\pgfqpoint{1.553629in}{1.324415in}}%
\pgfpathlineto{\pgfqpoint{1.552302in}{1.329252in}}%
\pgfpathlineto{\pgfqpoint{1.549636in}{1.337663in}}%
\pgfpathlineto{\pgfqpoint{1.552302in}{1.340751in}}%
\pgfpathlineto{\pgfqpoint{1.559364in}{1.350911in}}%
\pgfpathlineto{\pgfqpoint{1.558429in}{1.364159in}}%
\pgfpathlineto{\pgfqpoint{1.561135in}{1.377407in}}%
\pgfpathlineto{\pgfqpoint{1.564196in}{1.390655in}}%
\pgfpathlineto{\pgfqpoint{1.561191in}{1.403903in}}%
\pgfpathlineto{\pgfqpoint{1.567375in}{1.417151in}}%
\pgfpathlineto{\pgfqpoint{1.562499in}{1.430399in}}%
\pgfpathlineto{\pgfqpoint{1.568139in}{1.443647in}}%
\pgfpathlineto{\pgfqpoint{1.564695in}{1.456895in}}%
\pgfpathlineto{\pgfqpoint{1.566802in}{1.470143in}}%
\pgfpathlineto{\pgfqpoint{1.566953in}{1.483391in}}%
\pgfpathlineto{\pgfqpoint{1.564701in}{1.496639in}}%
\pgfpathlineto{\pgfqpoint{1.567779in}{1.509887in}}%
\pgfpathlineto{\pgfqpoint{1.562910in}{1.523135in}}%
\pgfpathlineto{\pgfqpoint{1.566841in}{1.536383in}}%
\pgfpathlineto{\pgfqpoint{1.562631in}{1.549631in}}%
\pgfpathlineto{\pgfqpoint{1.564097in}{1.562879in}}%
\pgfpathlineto{\pgfqpoint{1.562463in}{1.576127in}}%
\pgfpathlineto{\pgfqpoint{1.560102in}{1.589375in}}%
\pgfpathlineto{\pgfqpoint{1.561336in}{1.602623in}}%
\pgfpathlineto{\pgfqpoint{1.556681in}{1.615871in}}%
\pgfpathlineto{\pgfqpoint{1.558940in}{1.629119in}}%
\pgfpathlineto{\pgfqpoint{1.554364in}{1.642367in}}%
\pgfpathlineto{\pgfqpoint{1.555590in}{1.655615in}}%
\pgfpathlineto{\pgfqpoint{1.553379in}{1.668863in}}%
\pgfpathlineto{\pgfqpoint{1.552302in}{1.677533in}}%
\pgfpathlineto{\pgfqpoint{1.551649in}{1.682111in}}%
\pgfpathlineto{\pgfqpoint{1.551997in}{1.695359in}}%
\pgfpathlineto{\pgfqpoint{1.546868in}{1.708607in}}%
\pgfpathlineto{\pgfqpoint{1.548481in}{1.721855in}}%
\pgfpathlineto{\pgfqpoint{1.543086in}{1.735103in}}%
\pgfpathlineto{\pgfqpoint{1.543490in}{1.748351in}}%
\pgfpathlineto{\pgfqpoint{1.540121in}{1.761599in}}%
\pgfpathlineto{\pgfqpoint{1.538039in}{1.774847in}}%
\pgfpathlineto{\pgfqpoint{1.537767in}{1.788095in}}%
\pgfpathlineto{\pgfqpoint{1.533478in}{1.801343in}}%
\pgfpathlineto{\pgfqpoint{1.534690in}{1.814591in}}%
\pgfpathlineto{\pgfqpoint{1.529499in}{1.827839in}}%
\pgfpathlineto{\pgfqpoint{1.529795in}{1.841087in}}%
\pgfpathlineto{\pgfqpoint{1.527018in}{1.851186in}}%
\pgfpathlineto{\pgfqpoint{1.525979in}{1.854335in}}%
\pgfpathlineto{\pgfqpoint{1.523214in}{1.867583in}}%
\pgfpathlineto{\pgfqpoint{1.521810in}{1.880831in}}%
\pgfpathlineto{\pgfqpoint{1.516094in}{1.894079in}}%
\pgfpathlineto{\pgfqpoint{1.516955in}{1.907327in}}%
\pgfpathlineto{\pgfqpoint{1.510799in}{1.920575in}}%
\pgfpathlineto{\pgfqpoint{1.510882in}{1.933823in}}%
\pgfpathlineto{\pgfqpoint{1.506488in}{1.947071in}}%
\pgfpathlineto{\pgfqpoint{1.503584in}{1.960319in}}%
\pgfpathlineto{\pgfqpoint{1.502145in}{1.973567in}}%
\pgfpathlineto{\pgfqpoint{1.501735in}{1.974583in}}%
\pgfpathlineto{\pgfqpoint{1.494019in}{1.986815in}}%
\pgfpathlineto{\pgfqpoint{1.494428in}{2.000063in}}%
\pgfpathlineto{\pgfqpoint{1.484841in}{2.013311in}}%
\pgfpathlineto{\pgfqpoint{1.484155in}{2.026559in}}%
\pgfpathlineto{\pgfqpoint{1.477601in}{2.039807in}}%
\pgfpathlineto{\pgfqpoint{1.476452in}{2.042833in}}%
\pgfpathlineto{\pgfqpoint{1.463152in}{2.053055in}}%
\pgfpathlineto{\pgfqpoint{1.458137in}{2.066303in}}%
\pgfpathlineto{\pgfqpoint{1.451169in}{2.069562in}}%
\pgfpathlineto{\pgfqpoint{1.445241in}{2.066303in}}%
\pgfpathlineto{\pgfqpoint{1.443016in}{2.053055in}}%
\pgfpathlineto{\pgfqpoint{1.426171in}{2.039807in}}%
\pgfpathlineto{\pgfqpoint{1.425886in}{2.039404in}}%
\pgfpathlineto{\pgfqpoint{1.422646in}{2.026559in}}%
\pgfpathlineto{\pgfqpoint{1.419405in}{2.013311in}}%
\pgfpathlineto{\pgfqpoint{1.412225in}{2.000063in}}%
\pgfpathlineto{\pgfqpoint{1.411610in}{1.986815in}}%
\pgfpathlineto{\pgfqpoint{1.402428in}{1.973567in}}%
\pgfpathlineto{\pgfqpoint{1.403139in}{1.960319in}}%
\pgfpathlineto{\pgfqpoint{1.400603in}{1.956302in}}%
\pgfpathlineto{\pgfqpoint{1.396574in}{1.947071in}}%
\pgfpathlineto{\pgfqpoint{1.395868in}{1.933823in}}%
\pgfpathlineto{\pgfqpoint{1.392812in}{1.920575in}}%
\pgfpathlineto{\pgfqpoint{1.389542in}{1.907327in}}%
\pgfpathlineto{\pgfqpoint{1.389115in}{1.894079in}}%
\pgfpathlineto{\pgfqpoint{1.383189in}{1.880831in}}%
\pgfpathlineto{\pgfqpoint{1.383788in}{1.867583in}}%
\pgfpathlineto{\pgfqpoint{1.377407in}{1.854335in}}%
\pgfpathlineto{\pgfqpoint{1.377567in}{1.841087in}}%
\pgfpathlineto{\pgfqpoint{1.375319in}{1.834267in}}%
\pgfpathlineto{\pgfqpoint{1.373523in}{1.827839in}}%
\pgfpathlineto{\pgfqpoint{1.371627in}{1.814591in}}%
\pgfpathlineto{\pgfqpoint{1.370808in}{1.801343in}}%
\pgfpathlineto{\pgfqpoint{1.366707in}{1.788095in}}%
\pgfpathlineto{\pgfqpoint{1.367867in}{1.774847in}}%
\pgfpathlineto{\pgfqpoint{1.362524in}{1.761599in}}%
\pgfpathlineto{\pgfqpoint{1.363536in}{1.748351in}}%
\pgfpathlineto{\pgfqpoint{1.359360in}{1.735103in}}%
\pgfpathlineto{\pgfqpoint{1.358506in}{1.721855in}}%
\pgfpathlineto{\pgfqpoint{1.357217in}{1.708607in}}%
\pgfpathlineto{\pgfqpoint{1.353462in}{1.695359in}}%
\pgfpathlineto{\pgfqpoint{1.354769in}{1.682111in}}%
\pgfpathlineto{\pgfqpoint{1.350036in}{1.670701in}}%
\pgfpathlineto{\pgfqpoint{1.349357in}{1.668863in}}%
\pgfpathlineto{\pgfqpoint{1.350036in}{1.663739in}}%
\pgfpathlineto{\pgfqpoint{1.351283in}{1.655615in}}%
\pgfpathlineto{\pgfqpoint{1.350036in}{1.652180in}}%
\pgfpathlineto{\pgfqpoint{1.346798in}{1.642367in}}%
\pgfpathlineto{\pgfqpoint{1.347035in}{1.629119in}}%
\pgfpathlineto{\pgfqpoint{1.345535in}{1.615871in}}%
\pgfpathlineto{\pgfqpoint{1.343123in}{1.602623in}}%
\pgfpathlineto{\pgfqpoint{1.345077in}{1.589375in}}%
\pgfpathlineto{\pgfqpoint{1.340365in}{1.576127in}}%
\pgfpathlineto{\pgfqpoint{1.343475in}{1.562879in}}%
\pgfpathlineto{\pgfqpoint{1.338868in}{1.549631in}}%
\pgfpathlineto{\pgfqpoint{1.340342in}{1.536383in}}%
\pgfpathlineto{\pgfqpoint{1.338473in}{1.523135in}}%
\pgfpathlineto{\pgfqpoint{1.336504in}{1.509887in}}%
\pgfpathlineto{\pgfqpoint{1.338820in}{1.496639in}}%
\pgfpathlineto{\pgfqpoint{1.334281in}{1.483391in}}%
\pgfpathlineto{\pgfqpoint{1.339588in}{1.470143in}}%
\pgfpathlineto{\pgfqpoint{1.335058in}{1.456895in}}%
\pgfpathlineto{\pgfqpoint{1.339124in}{1.443647in}}%
\pgfpathlineto{\pgfqpoint{1.338117in}{1.430399in}}%
\pgfpathlineto{\pgfqpoint{1.337745in}{1.417151in}}%
\pgfpathlineto{\pgfqpoint{1.341885in}{1.403903in}}%
\pgfpathlineto{\pgfqpoint{1.337092in}{1.390655in}}%
\pgfpathlineto{\pgfqpoint{1.344785in}{1.377407in}}%
\pgfpathlineto{\pgfqpoint{1.340145in}{1.364159in}}%
\pgfpathlineto{\pgfqpoint{1.347626in}{1.350911in}}%
\pgfpathlineto{\pgfqpoint{1.348409in}{1.337663in}}%
\pgfpathlineto{\pgfqpoint{1.350036in}{1.328157in}}%
\pgfpathlineto{\pgfqpoint{1.350815in}{1.324415in}}%
\pgfpathlineto{\pgfqpoint{1.361627in}{1.311167in}}%
\pgfpathlineto{\pgfqpoint{1.357760in}{1.297919in}}%
\pgfpathlineto{\pgfqpoint{1.375319in}{1.284746in}}%
\pgfpathlineto{\pgfqpoint{1.375461in}{1.284671in}}%
\pgfpathlineto{\pgfqpoint{1.375319in}{1.284295in}}%
\pgfpathlineto{\pgfqpoint{1.371468in}{1.271423in}}%
\pgfpathlineto{\pgfqpoint{1.375319in}{1.268058in}}%
\pgfpathlineto{\pgfqpoint{1.390549in}{1.258175in}}%
\pgfpathlineto{\pgfqpoint{1.400603in}{1.245513in}}%
\pgfpathlineto{\pgfqpoint{1.401504in}{1.244927in}}%
\pgfpathlineto{\pgfqpoint{1.415378in}{1.231679in}}%
\pgfpathlineto{\pgfqpoint{1.425886in}{1.228799in}}%
\pgfpathclose%
\pgfpathmoveto{\pgfqpoint{1.415206in}{1.271423in}}%
\pgfpathlineto{\pgfqpoint{1.417960in}{1.284671in}}%
\pgfpathlineto{\pgfqpoint{1.400603in}{1.290074in}}%
\pgfpathlineto{\pgfqpoint{1.385677in}{1.297919in}}%
\pgfpathlineto{\pgfqpoint{1.390395in}{1.311167in}}%
\pgfpathlineto{\pgfqpoint{1.375319in}{1.322140in}}%
\pgfpathlineto{\pgfqpoint{1.373232in}{1.324415in}}%
\pgfpathlineto{\pgfqpoint{1.370572in}{1.337663in}}%
\pgfpathlineto{\pgfqpoint{1.367516in}{1.350911in}}%
\pgfpathlineto{\pgfqpoint{1.358219in}{1.364159in}}%
\pgfpathlineto{\pgfqpoint{1.362804in}{1.377407in}}%
\pgfpathlineto{\pgfqpoint{1.352630in}{1.390655in}}%
\pgfpathlineto{\pgfqpoint{1.358249in}{1.403903in}}%
\pgfpathlineto{\pgfqpoint{1.351979in}{1.417151in}}%
\pgfpathlineto{\pgfqpoint{1.352506in}{1.430399in}}%
\pgfpathlineto{\pgfqpoint{1.352572in}{1.443647in}}%
\pgfpathlineto{\pgfqpoint{1.350036in}{1.450239in}}%
\pgfpathlineto{\pgfqpoint{1.347904in}{1.456895in}}%
\pgfpathlineto{\pgfqpoint{1.350036in}{1.463750in}}%
\pgfpathlineto{\pgfqpoint{1.352289in}{1.470143in}}%
\pgfpathlineto{\pgfqpoint{1.350036in}{1.474351in}}%
\pgfpathlineto{\pgfqpoint{1.346104in}{1.483391in}}%
\pgfpathlineto{\pgfqpoint{1.350036in}{1.495042in}}%
\pgfpathlineto{\pgfqpoint{1.350667in}{1.496639in}}%
\pgfpathlineto{\pgfqpoint{1.350036in}{1.498871in}}%
\pgfpathlineto{\pgfqpoint{1.347551in}{1.509887in}}%
\pgfpathlineto{\pgfqpoint{1.349583in}{1.523135in}}%
\pgfpathlineto{\pgfqpoint{1.350036in}{1.527743in}}%
\pgfpathlineto{\pgfqpoint{1.350966in}{1.536383in}}%
\pgfpathlineto{\pgfqpoint{1.350036in}{1.543192in}}%
\pgfpathlineto{\pgfqpoint{1.349283in}{1.549631in}}%
\pgfpathlineto{\pgfqpoint{1.350036in}{1.551895in}}%
\pgfpathlineto{\pgfqpoint{1.354120in}{1.562879in}}%
\pgfpathlineto{\pgfqpoint{1.350145in}{1.576127in}}%
\pgfpathlineto{\pgfqpoint{1.355545in}{1.589375in}}%
\pgfpathlineto{\pgfqpoint{1.352729in}{1.602623in}}%
\pgfpathlineto{\pgfqpoint{1.355619in}{1.615871in}}%
\pgfpathlineto{\pgfqpoint{1.356768in}{1.629119in}}%
\pgfpathlineto{\pgfqpoint{1.356518in}{1.642367in}}%
\pgfpathlineto{\pgfqpoint{1.361069in}{1.655615in}}%
\pgfpathlineto{\pgfqpoint{1.358835in}{1.668863in}}%
\pgfpathlineto{\pgfqpoint{1.364255in}{1.682111in}}%
\pgfpathlineto{\pgfqpoint{1.362526in}{1.695359in}}%
\pgfpathlineto{\pgfqpoint{1.366409in}{1.708607in}}%
\pgfpathlineto{\pgfqpoint{1.367264in}{1.721855in}}%
\pgfpathlineto{\pgfqpoint{1.368225in}{1.735103in}}%
\pgfpathlineto{\pgfqpoint{1.372135in}{1.748351in}}%
\pgfpathlineto{\pgfqpoint{1.371021in}{1.761599in}}%
\pgfpathlineto{\pgfqpoint{1.375319in}{1.772318in}}%
\pgfpathlineto{\pgfqpoint{1.376568in}{1.774847in}}%
\pgfpathlineto{\pgfqpoint{1.375319in}{1.783786in}}%
\pgfpathlineto{\pgfqpoint{1.374851in}{1.788095in}}%
\pgfpathlineto{\pgfqpoint{1.375319in}{1.789587in}}%
\pgfpathlineto{\pgfqpoint{1.379944in}{1.801343in}}%
\pgfpathlineto{\pgfqpoint{1.380462in}{1.814591in}}%
\pgfpathlineto{\pgfqpoint{1.382975in}{1.827839in}}%
\pgfpathlineto{\pgfqpoint{1.387020in}{1.841087in}}%
\pgfpathlineto{\pgfqpoint{1.386916in}{1.854335in}}%
\pgfpathlineto{\pgfqpoint{1.393185in}{1.867583in}}%
\pgfpathlineto{\pgfqpoint{1.392362in}{1.880831in}}%
\pgfpathlineto{\pgfqpoint{1.398446in}{1.894079in}}%
\pgfpathlineto{\pgfqpoint{1.398462in}{1.907327in}}%
\pgfpathlineto{\pgfqpoint{1.400603in}{1.915517in}}%
\pgfpathlineto{\pgfqpoint{1.402637in}{1.920575in}}%
\pgfpathlineto{\pgfqpoint{1.406563in}{1.933823in}}%
\pgfpathlineto{\pgfqpoint{1.407882in}{1.947071in}}%
\pgfpathlineto{\pgfqpoint{1.416004in}{1.960319in}}%
\pgfpathlineto{\pgfqpoint{1.415214in}{1.973567in}}%
\pgfpathlineto{\pgfqpoint{1.424508in}{1.986815in}}%
\pgfpathlineto{\pgfqpoint{1.424552in}{2.000063in}}%
\pgfpathlineto{\pgfqpoint{1.425886in}{2.002441in}}%
\pgfpathlineto{\pgfqpoint{1.442853in}{2.013311in}}%
\pgfpathlineto{\pgfqpoint{1.447926in}{2.026559in}}%
\pgfpathlineto{\pgfqpoint{1.451169in}{2.030398in}}%
\pgfpathlineto{\pgfqpoint{1.456152in}{2.026559in}}%
\pgfpathlineto{\pgfqpoint{1.461037in}{2.013311in}}%
\pgfpathlineto{\pgfqpoint{1.476452in}{2.006199in}}%
\pgfpathlineto{\pgfqpoint{1.480877in}{2.000063in}}%
\pgfpathlineto{\pgfqpoint{1.480154in}{1.986815in}}%
\pgfpathlineto{\pgfqpoint{1.488967in}{1.973567in}}%
\pgfpathlineto{\pgfqpoint{1.490391in}{1.960319in}}%
\pgfpathlineto{\pgfqpoint{1.495334in}{1.947071in}}%
\pgfpathlineto{\pgfqpoint{1.501527in}{1.933823in}}%
\pgfpathlineto{\pgfqpoint{1.501364in}{1.920575in}}%
\pgfpathlineto{\pgfqpoint{1.501735in}{1.920044in}}%
\pgfpathlineto{\pgfqpoint{1.507602in}{1.907327in}}%
\pgfpathlineto{\pgfqpoint{1.506443in}{1.894079in}}%
\pgfpathlineto{\pgfqpoint{1.512356in}{1.880831in}}%
\pgfpathlineto{\pgfqpoint{1.513331in}{1.867583in}}%
\pgfpathlineto{\pgfqpoint{1.516345in}{1.854335in}}%
\pgfpathlineto{\pgfqpoint{1.520463in}{1.841087in}}%
\pgfpathlineto{\pgfqpoint{1.520110in}{1.827839in}}%
\pgfpathlineto{\pgfqpoint{1.526476in}{1.814591in}}%
\pgfpathlineto{\pgfqpoint{1.524628in}{1.801343in}}%
\pgfpathlineto{\pgfqpoint{1.527018in}{1.795605in}}%
\pgfpathlineto{\pgfqpoint{1.529494in}{1.788095in}}%
\pgfpathlineto{\pgfqpoint{1.529376in}{1.774847in}}%
\pgfpathlineto{\pgfqpoint{1.531577in}{1.761599in}}%
\pgfpathlineto{\pgfqpoint{1.534627in}{1.748351in}}%
\pgfpathlineto{\pgfqpoint{1.534184in}{1.735103in}}%
\pgfpathlineto{\pgfqpoint{1.539453in}{1.721855in}}%
\pgfpathlineto{\pgfqpoint{1.537551in}{1.708607in}}%
\pgfpathlineto{\pgfqpoint{1.542733in}{1.695359in}}%
\pgfpathlineto{\pgfqpoint{1.541905in}{1.682111in}}%
\pgfpathlineto{\pgfqpoint{1.543896in}{1.668863in}}%
\pgfpathlineto{\pgfqpoint{1.545958in}{1.655615in}}%
\pgfpathlineto{\pgfqpoint{1.544517in}{1.642367in}}%
\pgfpathlineto{\pgfqpoint{1.549507in}{1.629119in}}%
\pgfpathlineto{\pgfqpoint{1.546554in}{1.615871in}}%
\pgfpathlineto{\pgfqpoint{1.551915in}{1.602623in}}%
\pgfpathlineto{\pgfqpoint{1.549884in}{1.589375in}}%
\pgfpathlineto{\pgfqpoint{1.552302in}{1.577996in}}%
\pgfpathlineto{\pgfqpoint{1.552641in}{1.576127in}}%
\pgfpathlineto{\pgfqpoint{1.553759in}{1.562879in}}%
\pgfpathlineto{\pgfqpoint{1.552302in}{1.550390in}}%
\pgfpathlineto{\pgfqpoint{1.552198in}{1.549631in}}%
\pgfpathlineto{\pgfqpoint{1.552302in}{1.549319in}}%
\pgfpathlineto{\pgfqpoint{1.556106in}{1.536383in}}%
\pgfpathlineto{\pgfqpoint{1.552302in}{1.524603in}}%
\pgfpathlineto{\pgfqpoint{1.551728in}{1.523135in}}%
\pgfpathlineto{\pgfqpoint{1.552302in}{1.521770in}}%
\pgfpathlineto{\pgfqpoint{1.556598in}{1.509887in}}%
\pgfpathlineto{\pgfqpoint{1.552860in}{1.496639in}}%
\pgfpathlineto{\pgfqpoint{1.555141in}{1.483391in}}%
\pgfpathlineto{\pgfqpoint{1.554168in}{1.470143in}}%
\pgfpathlineto{\pgfqpoint{1.552302in}{1.458895in}}%
\pgfpathlineto{\pgfqpoint{1.551906in}{1.456895in}}%
\pgfpathlineto{\pgfqpoint{1.552302in}{1.455210in}}%
\pgfpathlineto{\pgfqpoint{1.554720in}{1.443647in}}%
\pgfpathlineto{\pgfqpoint{1.552302in}{1.438294in}}%
\pgfpathlineto{\pgfqpoint{1.547907in}{1.430399in}}%
\pgfpathlineto{\pgfqpoint{1.552302in}{1.419729in}}%
\pgfpathlineto{\pgfqpoint{1.553211in}{1.417151in}}%
\pgfpathlineto{\pgfqpoint{1.552302in}{1.415361in}}%
\pgfpathlineto{\pgfqpoint{1.544851in}{1.403903in}}%
\pgfpathlineto{\pgfqpoint{1.548511in}{1.390655in}}%
\pgfpathlineto{\pgfqpoint{1.543293in}{1.377407in}}%
\pgfpathlineto{\pgfqpoint{1.540091in}{1.364159in}}%
\pgfpathlineto{\pgfqpoint{1.539594in}{1.350911in}}%
\pgfpathlineto{\pgfqpoint{1.527919in}{1.337663in}}%
\pgfpathlineto{\pgfqpoint{1.531016in}{1.324415in}}%
\pgfpathlineto{\pgfqpoint{1.527018in}{1.321026in}}%
\pgfpathlineto{\pgfqpoint{1.510364in}{1.311167in}}%
\pgfpathlineto{\pgfqpoint{1.513969in}{1.297919in}}%
\pgfpathlineto{\pgfqpoint{1.501735in}{1.289960in}}%
\pgfpathlineto{\pgfqpoint{1.486574in}{1.284671in}}%
\pgfpathlineto{\pgfqpoint{1.481822in}{1.271423in}}%
\pgfpathlineto{\pgfqpoint{1.476452in}{1.268025in}}%
\pgfpathlineto{\pgfqpoint{1.451169in}{1.263291in}}%
\pgfpathlineto{\pgfqpoint{1.425886in}{1.267027in}}%
\pgfpathlineto{\pgfqpoint{1.415206in}{1.271423in}}%
\pgfpathclose%
\pgfusepath{fill}%
\end{pgfscope}%
\begin{pgfscope}%
\pgfpathrectangle{\pgfqpoint{0.540975in}{0.476543in}}{\pgfqpoint{1.820388in}{1.656000in}}%
\pgfusepath{clip}%
\pgfsetbuttcap%
\pgfsetroundjoin%
\definecolor{currentfill}{rgb}{0.966551,0.497424,0.295040}%
\pgfsetfillcolor{currentfill}%
\pgfsetlinewidth{0.000000pt}%
\definecolor{currentstroke}{rgb}{0.000000,0.000000,0.000000}%
\pgfsetstrokecolor{currentstroke}%
\pgfsetdash{}{0pt}%
\pgfpathmoveto{\pgfqpoint{1.425886in}{1.267027in}}%
\pgfpathlineto{\pgfqpoint{1.451169in}{1.263291in}}%
\pgfpathlineto{\pgfqpoint{1.476452in}{1.268025in}}%
\pgfpathlineto{\pgfqpoint{1.481822in}{1.271423in}}%
\pgfpathlineto{\pgfqpoint{1.486574in}{1.284671in}}%
\pgfpathlineto{\pgfqpoint{1.501735in}{1.289960in}}%
\pgfpathlineto{\pgfqpoint{1.513969in}{1.297919in}}%
\pgfpathlineto{\pgfqpoint{1.510364in}{1.311167in}}%
\pgfpathlineto{\pgfqpoint{1.527018in}{1.321026in}}%
\pgfpathlineto{\pgfqpoint{1.531016in}{1.324415in}}%
\pgfpathlineto{\pgfqpoint{1.527919in}{1.337663in}}%
\pgfpathlineto{\pgfqpoint{1.539594in}{1.350911in}}%
\pgfpathlineto{\pgfqpoint{1.540091in}{1.364159in}}%
\pgfpathlineto{\pgfqpoint{1.543293in}{1.377407in}}%
\pgfpathlineto{\pgfqpoint{1.548511in}{1.390655in}}%
\pgfpathlineto{\pgfqpoint{1.544851in}{1.403903in}}%
\pgfpathlineto{\pgfqpoint{1.552302in}{1.415361in}}%
\pgfpathlineto{\pgfqpoint{1.553211in}{1.417151in}}%
\pgfpathlineto{\pgfqpoint{1.552302in}{1.419729in}}%
\pgfpathlineto{\pgfqpoint{1.547907in}{1.430399in}}%
\pgfpathlineto{\pgfqpoint{1.552302in}{1.438294in}}%
\pgfpathlineto{\pgfqpoint{1.554720in}{1.443647in}}%
\pgfpathlineto{\pgfqpoint{1.552302in}{1.455210in}}%
\pgfpathlineto{\pgfqpoint{1.551906in}{1.456895in}}%
\pgfpathlineto{\pgfqpoint{1.552302in}{1.458895in}}%
\pgfpathlineto{\pgfqpoint{1.554168in}{1.470143in}}%
\pgfpathlineto{\pgfqpoint{1.555141in}{1.483391in}}%
\pgfpathlineto{\pgfqpoint{1.552860in}{1.496639in}}%
\pgfpathlineto{\pgfqpoint{1.556598in}{1.509887in}}%
\pgfpathlineto{\pgfqpoint{1.552302in}{1.521770in}}%
\pgfpathlineto{\pgfqpoint{1.551728in}{1.523135in}}%
\pgfpathlineto{\pgfqpoint{1.552302in}{1.524603in}}%
\pgfpathlineto{\pgfqpoint{1.556106in}{1.536383in}}%
\pgfpathlineto{\pgfqpoint{1.552302in}{1.549319in}}%
\pgfpathlineto{\pgfqpoint{1.552198in}{1.549631in}}%
\pgfpathlineto{\pgfqpoint{1.552302in}{1.550390in}}%
\pgfpathlineto{\pgfqpoint{1.553759in}{1.562879in}}%
\pgfpathlineto{\pgfqpoint{1.552641in}{1.576127in}}%
\pgfpathlineto{\pgfqpoint{1.552302in}{1.577996in}}%
\pgfpathlineto{\pgfqpoint{1.549884in}{1.589375in}}%
\pgfpathlineto{\pgfqpoint{1.551915in}{1.602623in}}%
\pgfpathlineto{\pgfqpoint{1.546554in}{1.615871in}}%
\pgfpathlineto{\pgfqpoint{1.549507in}{1.629119in}}%
\pgfpathlineto{\pgfqpoint{1.544517in}{1.642367in}}%
\pgfpathlineto{\pgfqpoint{1.545958in}{1.655615in}}%
\pgfpathlineto{\pgfqpoint{1.543896in}{1.668863in}}%
\pgfpathlineto{\pgfqpoint{1.541905in}{1.682111in}}%
\pgfpathlineto{\pgfqpoint{1.542733in}{1.695359in}}%
\pgfpathlineto{\pgfqpoint{1.537551in}{1.708607in}}%
\pgfpathlineto{\pgfqpoint{1.539453in}{1.721855in}}%
\pgfpathlineto{\pgfqpoint{1.534184in}{1.735103in}}%
\pgfpathlineto{\pgfqpoint{1.534627in}{1.748351in}}%
\pgfpathlineto{\pgfqpoint{1.531577in}{1.761599in}}%
\pgfpathlineto{\pgfqpoint{1.529376in}{1.774847in}}%
\pgfpathlineto{\pgfqpoint{1.529494in}{1.788095in}}%
\pgfpathlineto{\pgfqpoint{1.527018in}{1.795605in}}%
\pgfpathlineto{\pgfqpoint{1.524628in}{1.801343in}}%
\pgfpathlineto{\pgfqpoint{1.526476in}{1.814591in}}%
\pgfpathlineto{\pgfqpoint{1.520110in}{1.827839in}}%
\pgfpathlineto{\pgfqpoint{1.520463in}{1.841087in}}%
\pgfpathlineto{\pgfqpoint{1.516345in}{1.854335in}}%
\pgfpathlineto{\pgfqpoint{1.513331in}{1.867583in}}%
\pgfpathlineto{\pgfqpoint{1.512356in}{1.880831in}}%
\pgfpathlineto{\pgfqpoint{1.506443in}{1.894079in}}%
\pgfpathlineto{\pgfqpoint{1.507602in}{1.907327in}}%
\pgfpathlineto{\pgfqpoint{1.501735in}{1.920044in}}%
\pgfpathlineto{\pgfqpoint{1.501364in}{1.920575in}}%
\pgfpathlineto{\pgfqpoint{1.501527in}{1.933823in}}%
\pgfpathlineto{\pgfqpoint{1.495334in}{1.947071in}}%
\pgfpathlineto{\pgfqpoint{1.490391in}{1.960319in}}%
\pgfpathlineto{\pgfqpoint{1.488967in}{1.973567in}}%
\pgfpathlineto{\pgfqpoint{1.480154in}{1.986815in}}%
\pgfpathlineto{\pgfqpoint{1.480877in}{2.000063in}}%
\pgfpathlineto{\pgfqpoint{1.476452in}{2.006199in}}%
\pgfpathlineto{\pgfqpoint{1.461037in}{2.013311in}}%
\pgfpathlineto{\pgfqpoint{1.456152in}{2.026559in}}%
\pgfpathlineto{\pgfqpoint{1.451169in}{2.030398in}}%
\pgfpathlineto{\pgfqpoint{1.447926in}{2.026559in}}%
\pgfpathlineto{\pgfqpoint{1.442853in}{2.013311in}}%
\pgfpathlineto{\pgfqpoint{1.425886in}{2.002441in}}%
\pgfpathlineto{\pgfqpoint{1.424552in}{2.000063in}}%
\pgfpathlineto{\pgfqpoint{1.424508in}{1.986815in}}%
\pgfpathlineto{\pgfqpoint{1.415214in}{1.973567in}}%
\pgfpathlineto{\pgfqpoint{1.416004in}{1.960319in}}%
\pgfpathlineto{\pgfqpoint{1.407882in}{1.947071in}}%
\pgfpathlineto{\pgfqpoint{1.406563in}{1.933823in}}%
\pgfpathlineto{\pgfqpoint{1.402637in}{1.920575in}}%
\pgfpathlineto{\pgfqpoint{1.400603in}{1.915517in}}%
\pgfpathlineto{\pgfqpoint{1.398462in}{1.907327in}}%
\pgfpathlineto{\pgfqpoint{1.398446in}{1.894079in}}%
\pgfpathlineto{\pgfqpoint{1.392362in}{1.880831in}}%
\pgfpathlineto{\pgfqpoint{1.393185in}{1.867583in}}%
\pgfpathlineto{\pgfqpoint{1.386916in}{1.854335in}}%
\pgfpathlineto{\pgfqpoint{1.387020in}{1.841087in}}%
\pgfpathlineto{\pgfqpoint{1.382975in}{1.827839in}}%
\pgfpathlineto{\pgfqpoint{1.380462in}{1.814591in}}%
\pgfpathlineto{\pgfqpoint{1.379944in}{1.801343in}}%
\pgfpathlineto{\pgfqpoint{1.375319in}{1.789587in}}%
\pgfpathlineto{\pgfqpoint{1.374851in}{1.788095in}}%
\pgfpathlineto{\pgfqpoint{1.375319in}{1.783786in}}%
\pgfpathlineto{\pgfqpoint{1.376568in}{1.774847in}}%
\pgfpathlineto{\pgfqpoint{1.375319in}{1.772318in}}%
\pgfpathlineto{\pgfqpoint{1.371021in}{1.761599in}}%
\pgfpathlineto{\pgfqpoint{1.372135in}{1.748351in}}%
\pgfpathlineto{\pgfqpoint{1.368225in}{1.735103in}}%
\pgfpathlineto{\pgfqpoint{1.367264in}{1.721855in}}%
\pgfpathlineto{\pgfqpoint{1.366409in}{1.708607in}}%
\pgfpathlineto{\pgfqpoint{1.362526in}{1.695359in}}%
\pgfpathlineto{\pgfqpoint{1.364255in}{1.682111in}}%
\pgfpathlineto{\pgfqpoint{1.358835in}{1.668863in}}%
\pgfpathlineto{\pgfqpoint{1.361069in}{1.655615in}}%
\pgfpathlineto{\pgfqpoint{1.356518in}{1.642367in}}%
\pgfpathlineto{\pgfqpoint{1.356768in}{1.629119in}}%
\pgfpathlineto{\pgfqpoint{1.355619in}{1.615871in}}%
\pgfpathlineto{\pgfqpoint{1.352729in}{1.602623in}}%
\pgfpathlineto{\pgfqpoint{1.355545in}{1.589375in}}%
\pgfpathlineto{\pgfqpoint{1.350145in}{1.576127in}}%
\pgfpathlineto{\pgfqpoint{1.354120in}{1.562879in}}%
\pgfpathlineto{\pgfqpoint{1.350036in}{1.551895in}}%
\pgfpathlineto{\pgfqpoint{1.349283in}{1.549631in}}%
\pgfpathlineto{\pgfqpoint{1.350036in}{1.543192in}}%
\pgfpathlineto{\pgfqpoint{1.350966in}{1.536383in}}%
\pgfpathlineto{\pgfqpoint{1.350036in}{1.527743in}}%
\pgfpathlineto{\pgfqpoint{1.349583in}{1.523135in}}%
\pgfpathlineto{\pgfqpoint{1.347551in}{1.509887in}}%
\pgfpathlineto{\pgfqpoint{1.350036in}{1.498871in}}%
\pgfpathlineto{\pgfqpoint{1.350667in}{1.496639in}}%
\pgfpathlineto{\pgfqpoint{1.350036in}{1.495042in}}%
\pgfpathlineto{\pgfqpoint{1.346104in}{1.483391in}}%
\pgfpathlineto{\pgfqpoint{1.350036in}{1.474351in}}%
\pgfpathlineto{\pgfqpoint{1.352289in}{1.470143in}}%
\pgfpathlineto{\pgfqpoint{1.350036in}{1.463750in}}%
\pgfpathlineto{\pgfqpoint{1.347904in}{1.456895in}}%
\pgfpathlineto{\pgfqpoint{1.350036in}{1.450239in}}%
\pgfpathlineto{\pgfqpoint{1.352572in}{1.443647in}}%
\pgfpathlineto{\pgfqpoint{1.352506in}{1.430399in}}%
\pgfpathlineto{\pgfqpoint{1.351979in}{1.417151in}}%
\pgfpathlineto{\pgfqpoint{1.358249in}{1.403903in}}%
\pgfpathlineto{\pgfqpoint{1.352630in}{1.390655in}}%
\pgfpathlineto{\pgfqpoint{1.362804in}{1.377407in}}%
\pgfpathlineto{\pgfqpoint{1.358219in}{1.364159in}}%
\pgfpathlineto{\pgfqpoint{1.367516in}{1.350911in}}%
\pgfpathlineto{\pgfqpoint{1.370572in}{1.337663in}}%
\pgfpathlineto{\pgfqpoint{1.373232in}{1.324415in}}%
\pgfpathlineto{\pgfqpoint{1.375319in}{1.322140in}}%
\pgfpathlineto{\pgfqpoint{1.390395in}{1.311167in}}%
\pgfpathlineto{\pgfqpoint{1.385677in}{1.297919in}}%
\pgfpathlineto{\pgfqpoint{1.400603in}{1.290074in}}%
\pgfpathlineto{\pgfqpoint{1.417960in}{1.284671in}}%
\pgfpathlineto{\pgfqpoint{1.415206in}{1.271423in}}%
\pgfpathlineto{\pgfqpoint{1.425886in}{1.267027in}}%
\pgfpathclose%
\pgfpathmoveto{\pgfqpoint{1.446587in}{1.297919in}}%
\pgfpathlineto{\pgfqpoint{1.451169in}{1.300877in}}%
\pgfpathlineto{\pgfqpoint{1.454861in}{1.297919in}}%
\pgfpathlineto{\pgfqpoint{1.451169in}{1.297460in}}%
\pgfpathlineto{\pgfqpoint{1.446587in}{1.297919in}}%
\pgfpathclose%
\pgfpathmoveto{\pgfqpoint{1.403586in}{1.324415in}}%
\pgfpathlineto{\pgfqpoint{1.400603in}{1.332605in}}%
\pgfpathlineto{\pgfqpoint{1.399429in}{1.337663in}}%
\pgfpathlineto{\pgfqpoint{1.391932in}{1.350911in}}%
\pgfpathlineto{\pgfqpoint{1.379517in}{1.364159in}}%
\pgfpathlineto{\pgfqpoint{1.383966in}{1.377407in}}%
\pgfpathlineto{\pgfqpoint{1.375319in}{1.385103in}}%
\pgfpathlineto{\pgfqpoint{1.370816in}{1.390655in}}%
\pgfpathlineto{\pgfqpoint{1.375319in}{1.401645in}}%
\pgfpathlineto{\pgfqpoint{1.376535in}{1.403903in}}%
\pgfpathlineto{\pgfqpoint{1.375319in}{1.405375in}}%
\pgfpathlineto{\pgfqpoint{1.368564in}{1.417151in}}%
\pgfpathlineto{\pgfqpoint{1.369265in}{1.430399in}}%
\pgfpathlineto{\pgfqpoint{1.368001in}{1.443647in}}%
\pgfpathlineto{\pgfqpoint{1.362833in}{1.456895in}}%
\pgfpathlineto{\pgfqpoint{1.366846in}{1.470143in}}%
\pgfpathlineto{\pgfqpoint{1.359387in}{1.483391in}}%
\pgfpathlineto{\pgfqpoint{1.364506in}{1.496639in}}%
\pgfpathlineto{\pgfqpoint{1.360070in}{1.509887in}}%
\pgfpathlineto{\pgfqpoint{1.362570in}{1.523135in}}%
\pgfpathlineto{\pgfqpoint{1.363135in}{1.536383in}}%
\pgfpathlineto{\pgfqpoint{1.361332in}{1.549631in}}%
\pgfpathlineto{\pgfqpoint{1.365761in}{1.562879in}}%
\pgfpathlineto{\pgfqpoint{1.361460in}{1.576127in}}%
\pgfpathlineto{\pgfqpoint{1.366760in}{1.589375in}}%
\pgfpathlineto{\pgfqpoint{1.363355in}{1.602623in}}%
\pgfpathlineto{\pgfqpoint{1.366379in}{1.615871in}}%
\pgfpathlineto{\pgfqpoint{1.366912in}{1.629119in}}%
\pgfpathlineto{\pgfqpoint{1.366706in}{1.642367in}}%
\pgfpathlineto{\pgfqpoint{1.370855in}{1.655615in}}%
\pgfpathlineto{\pgfqpoint{1.368403in}{1.668863in}}%
\pgfpathlineto{\pgfqpoint{1.373741in}{1.682111in}}%
\pgfpathlineto{\pgfqpoint{1.371591in}{1.695359in}}%
\pgfpathlineto{\pgfqpoint{1.375319in}{1.707685in}}%
\pgfpathlineto{\pgfqpoint{1.375676in}{1.708607in}}%
\pgfpathlineto{\pgfqpoint{1.376195in}{1.721855in}}%
\pgfpathlineto{\pgfqpoint{1.377553in}{1.735103in}}%
\pgfpathlineto{\pgfqpoint{1.382030in}{1.748351in}}%
\pgfpathlineto{\pgfqpoint{1.380577in}{1.761599in}}%
\pgfpathlineto{\pgfqpoint{1.387028in}{1.774847in}}%
\pgfpathlineto{\pgfqpoint{1.384811in}{1.788095in}}%
\pgfpathlineto{\pgfqpoint{1.390137in}{1.801343in}}%
\pgfpathlineto{\pgfqpoint{1.390113in}{1.814591in}}%
\pgfpathlineto{\pgfqpoint{1.392843in}{1.827839in}}%
\pgfpathlineto{\pgfqpoint{1.396472in}{1.841087in}}%
\pgfpathlineto{\pgfqpoint{1.396425in}{1.854335in}}%
\pgfpathlineto{\pgfqpoint{1.400603in}{1.863290in}}%
\pgfpathlineto{\pgfqpoint{1.403545in}{1.867583in}}%
\pgfpathlineto{\pgfqpoint{1.402000in}{1.880831in}}%
\pgfpathlineto{\pgfqpoint{1.411314in}{1.894079in}}%
\pgfpathlineto{\pgfqpoint{1.410613in}{1.907327in}}%
\pgfpathlineto{\pgfqpoint{1.416356in}{1.920575in}}%
\pgfpathlineto{\pgfqpoint{1.419453in}{1.933823in}}%
\pgfpathlineto{\pgfqpoint{1.421191in}{1.947071in}}%
\pgfpathlineto{\pgfqpoint{1.425886in}{1.955064in}}%
\pgfpathlineto{\pgfqpoint{1.433554in}{1.960319in}}%
\pgfpathlineto{\pgfqpoint{1.431614in}{1.973567in}}%
\pgfpathlineto{\pgfqpoint{1.451169in}{1.984074in}}%
\pgfpathlineto{\pgfqpoint{1.473838in}{1.973567in}}%
\pgfpathlineto{\pgfqpoint{1.475636in}{1.960319in}}%
\pgfpathlineto{\pgfqpoint{1.476452in}{1.959744in}}%
\pgfpathlineto{\pgfqpoint{1.481763in}{1.947071in}}%
\pgfpathlineto{\pgfqpoint{1.487112in}{1.933823in}}%
\pgfpathlineto{\pgfqpoint{1.487196in}{1.920575in}}%
\pgfpathlineto{\pgfqpoint{1.496316in}{1.907327in}}%
\pgfpathlineto{\pgfqpoint{1.494094in}{1.894079in}}%
\pgfpathlineto{\pgfqpoint{1.501735in}{1.883405in}}%
\pgfpathlineto{\pgfqpoint{1.502903in}{1.880831in}}%
\pgfpathlineto{\pgfqpoint{1.503449in}{1.867583in}}%
\pgfpathlineto{\pgfqpoint{1.506710in}{1.854335in}}%
\pgfpathlineto{\pgfqpoint{1.510441in}{1.841087in}}%
\pgfpathlineto{\pgfqpoint{1.510139in}{1.827839in}}%
\pgfpathlineto{\pgfqpoint{1.516328in}{1.814591in}}%
\pgfpathlineto{\pgfqpoint{1.514171in}{1.801343in}}%
\pgfpathlineto{\pgfqpoint{1.519762in}{1.788095in}}%
\pgfpathlineto{\pgfqpoint{1.519053in}{1.774847in}}%
\pgfpathlineto{\pgfqpoint{1.522029in}{1.761599in}}%
\pgfpathlineto{\pgfqpoint{1.525421in}{1.748351in}}%
\pgfpathlineto{\pgfqpoint{1.524831in}{1.735103in}}%
\pgfpathlineto{\pgfqpoint{1.527018in}{1.730589in}}%
\pgfpathlineto{\pgfqpoint{1.530426in}{1.721855in}}%
\pgfpathlineto{\pgfqpoint{1.528234in}{1.708607in}}%
\pgfpathlineto{\pgfqpoint{1.533470in}{1.695359in}}%
\pgfpathlineto{\pgfqpoint{1.532160in}{1.682111in}}%
\pgfpathlineto{\pgfqpoint{1.534270in}{1.668863in}}%
\pgfpathlineto{\pgfqpoint{1.535842in}{1.655615in}}%
\pgfpathlineto{\pgfqpoint{1.534388in}{1.642367in}}%
\pgfpathlineto{\pgfqpoint{1.539064in}{1.629119in}}%
\pgfpathlineto{\pgfqpoint{1.535779in}{1.615871in}}%
\pgfpathlineto{\pgfqpoint{1.541091in}{1.602623in}}%
\pgfpathlineto{\pgfqpoint{1.538399in}{1.589375in}}%
\pgfpathlineto{\pgfqpoint{1.541311in}{1.576127in}}%
\pgfpathlineto{\pgfqpoint{1.541872in}{1.562879in}}%
\pgfpathlineto{\pgfqpoint{1.540075in}{1.549631in}}%
\pgfpathlineto{\pgfqpoint{1.544129in}{1.536383in}}%
\pgfpathlineto{\pgfqpoint{1.538739in}{1.523135in}}%
\pgfpathlineto{\pgfqpoint{1.544188in}{1.509887in}}%
\pgfpathlineto{\pgfqpoint{1.538971in}{1.496639in}}%
\pgfpathlineto{\pgfqpoint{1.541730in}{1.483391in}}%
\pgfpathlineto{\pgfqpoint{1.539451in}{1.470143in}}%
\pgfpathlineto{\pgfqpoint{1.536864in}{1.456895in}}%
\pgfpathlineto{\pgfqpoint{1.539095in}{1.443647in}}%
\pgfpathlineto{\pgfqpoint{1.531374in}{1.430399in}}%
\pgfpathlineto{\pgfqpoint{1.536382in}{1.417151in}}%
\pgfpathlineto{\pgfqpoint{1.527018in}{1.404283in}}%
\pgfpathlineto{\pgfqpoint{1.526626in}{1.403903in}}%
\pgfpathlineto{\pgfqpoint{1.527018in}{1.402864in}}%
\pgfpathlineto{\pgfqpoint{1.530464in}{1.390655in}}%
\pgfpathlineto{\pgfqpoint{1.527018in}{1.383733in}}%
\pgfpathlineto{\pgfqpoint{1.522385in}{1.377407in}}%
\pgfpathlineto{\pgfqpoint{1.518453in}{1.364159in}}%
\pgfpathlineto{\pgfqpoint{1.515172in}{1.350911in}}%
\pgfpathlineto{\pgfqpoint{1.501735in}{1.339608in}}%
\pgfpathlineto{\pgfqpoint{1.498094in}{1.337663in}}%
\pgfpathlineto{\pgfqpoint{1.501395in}{1.324415in}}%
\pgfpathlineto{\pgfqpoint{1.476452in}{1.315713in}}%
\pgfpathlineto{\pgfqpoint{1.451169in}{1.312746in}}%
\pgfpathlineto{\pgfqpoint{1.425886in}{1.315633in}}%
\pgfpathlineto{\pgfqpoint{1.403586in}{1.324415in}}%
\pgfpathclose%
\pgfusepath{fill}%
\end{pgfscope}%
\begin{pgfscope}%
\pgfpathrectangle{\pgfqpoint{0.540975in}{0.476543in}}{\pgfqpoint{1.820388in}{1.656000in}}%
\pgfusepath{clip}%
\pgfsetbuttcap%
\pgfsetroundjoin%
\definecolor{currentfill}{rgb}{0.939023,0.389927,0.245521}%
\pgfsetfillcolor{currentfill}%
\pgfsetlinewidth{0.000000pt}%
\definecolor{currentstroke}{rgb}{0.000000,0.000000,0.000000}%
\pgfsetstrokecolor{currentstroke}%
\pgfsetdash{}{0pt}%
\pgfpathmoveto{\pgfqpoint{1.451169in}{1.297460in}}%
\pgfpathlineto{\pgfqpoint{1.454861in}{1.297919in}}%
\pgfpathlineto{\pgfqpoint{1.451169in}{1.300877in}}%
\pgfpathlineto{\pgfqpoint{1.446587in}{1.297919in}}%
\pgfpathlineto{\pgfqpoint{1.451169in}{1.297460in}}%
\pgfpathclose%
\pgfpathmoveto{\pgfqpoint{1.425886in}{1.315633in}}%
\pgfpathlineto{\pgfqpoint{1.451169in}{1.312746in}}%
\pgfpathlineto{\pgfqpoint{1.476452in}{1.315713in}}%
\pgfpathlineto{\pgfqpoint{1.501395in}{1.324415in}}%
\pgfpathlineto{\pgfqpoint{1.498094in}{1.337663in}}%
\pgfpathlineto{\pgfqpoint{1.501735in}{1.339608in}}%
\pgfpathlineto{\pgfqpoint{1.515172in}{1.350911in}}%
\pgfpathlineto{\pgfqpoint{1.518453in}{1.364159in}}%
\pgfpathlineto{\pgfqpoint{1.522385in}{1.377407in}}%
\pgfpathlineto{\pgfqpoint{1.527018in}{1.383733in}}%
\pgfpathlineto{\pgfqpoint{1.530464in}{1.390655in}}%
\pgfpathlineto{\pgfqpoint{1.527018in}{1.402864in}}%
\pgfpathlineto{\pgfqpoint{1.526626in}{1.403903in}}%
\pgfpathlineto{\pgfqpoint{1.527018in}{1.404283in}}%
\pgfpathlineto{\pgfqpoint{1.536382in}{1.417151in}}%
\pgfpathlineto{\pgfqpoint{1.531374in}{1.430399in}}%
\pgfpathlineto{\pgfqpoint{1.539095in}{1.443647in}}%
\pgfpathlineto{\pgfqpoint{1.536864in}{1.456895in}}%
\pgfpathlineto{\pgfqpoint{1.539451in}{1.470143in}}%
\pgfpathlineto{\pgfqpoint{1.541730in}{1.483391in}}%
\pgfpathlineto{\pgfqpoint{1.538971in}{1.496639in}}%
\pgfpathlineto{\pgfqpoint{1.544188in}{1.509887in}}%
\pgfpathlineto{\pgfqpoint{1.538739in}{1.523135in}}%
\pgfpathlineto{\pgfqpoint{1.544129in}{1.536383in}}%
\pgfpathlineto{\pgfqpoint{1.540075in}{1.549631in}}%
\pgfpathlineto{\pgfqpoint{1.541872in}{1.562879in}}%
\pgfpathlineto{\pgfqpoint{1.541311in}{1.576127in}}%
\pgfpathlineto{\pgfqpoint{1.538399in}{1.589375in}}%
\pgfpathlineto{\pgfqpoint{1.541091in}{1.602623in}}%
\pgfpathlineto{\pgfqpoint{1.535779in}{1.615871in}}%
\pgfpathlineto{\pgfqpoint{1.539064in}{1.629119in}}%
\pgfpathlineto{\pgfqpoint{1.534388in}{1.642367in}}%
\pgfpathlineto{\pgfqpoint{1.535842in}{1.655615in}}%
\pgfpathlineto{\pgfqpoint{1.534270in}{1.668863in}}%
\pgfpathlineto{\pgfqpoint{1.532160in}{1.682111in}}%
\pgfpathlineto{\pgfqpoint{1.533470in}{1.695359in}}%
\pgfpathlineto{\pgfqpoint{1.528234in}{1.708607in}}%
\pgfpathlineto{\pgfqpoint{1.530426in}{1.721855in}}%
\pgfpathlineto{\pgfqpoint{1.527018in}{1.730589in}}%
\pgfpathlineto{\pgfqpoint{1.524831in}{1.735103in}}%
\pgfpathlineto{\pgfqpoint{1.525421in}{1.748351in}}%
\pgfpathlineto{\pgfqpoint{1.522029in}{1.761599in}}%
\pgfpathlineto{\pgfqpoint{1.519053in}{1.774847in}}%
\pgfpathlineto{\pgfqpoint{1.519762in}{1.788095in}}%
\pgfpathlineto{\pgfqpoint{1.514171in}{1.801343in}}%
\pgfpathlineto{\pgfqpoint{1.516328in}{1.814591in}}%
\pgfpathlineto{\pgfqpoint{1.510139in}{1.827839in}}%
\pgfpathlineto{\pgfqpoint{1.510441in}{1.841087in}}%
\pgfpathlineto{\pgfqpoint{1.506710in}{1.854335in}}%
\pgfpathlineto{\pgfqpoint{1.503449in}{1.867583in}}%
\pgfpathlineto{\pgfqpoint{1.502903in}{1.880831in}}%
\pgfpathlineto{\pgfqpoint{1.501735in}{1.883405in}}%
\pgfpathlineto{\pgfqpoint{1.494094in}{1.894079in}}%
\pgfpathlineto{\pgfqpoint{1.496316in}{1.907327in}}%
\pgfpathlineto{\pgfqpoint{1.487196in}{1.920575in}}%
\pgfpathlineto{\pgfqpoint{1.487112in}{1.933823in}}%
\pgfpathlineto{\pgfqpoint{1.481763in}{1.947071in}}%
\pgfpathlineto{\pgfqpoint{1.476452in}{1.959744in}}%
\pgfpathlineto{\pgfqpoint{1.475636in}{1.960319in}}%
\pgfpathlineto{\pgfqpoint{1.473838in}{1.973567in}}%
\pgfpathlineto{\pgfqpoint{1.451169in}{1.984074in}}%
\pgfpathlineto{\pgfqpoint{1.431614in}{1.973567in}}%
\pgfpathlineto{\pgfqpoint{1.433554in}{1.960319in}}%
\pgfpathlineto{\pgfqpoint{1.425886in}{1.955064in}}%
\pgfpathlineto{\pgfqpoint{1.421191in}{1.947071in}}%
\pgfpathlineto{\pgfqpoint{1.419453in}{1.933823in}}%
\pgfpathlineto{\pgfqpoint{1.416356in}{1.920575in}}%
\pgfpathlineto{\pgfqpoint{1.410613in}{1.907327in}}%
\pgfpathlineto{\pgfqpoint{1.411314in}{1.894079in}}%
\pgfpathlineto{\pgfqpoint{1.402000in}{1.880831in}}%
\pgfpathlineto{\pgfqpoint{1.403545in}{1.867583in}}%
\pgfpathlineto{\pgfqpoint{1.400603in}{1.863290in}}%
\pgfpathlineto{\pgfqpoint{1.396425in}{1.854335in}}%
\pgfpathlineto{\pgfqpoint{1.396472in}{1.841087in}}%
\pgfpathlineto{\pgfqpoint{1.392843in}{1.827839in}}%
\pgfpathlineto{\pgfqpoint{1.390113in}{1.814591in}}%
\pgfpathlineto{\pgfqpoint{1.390137in}{1.801343in}}%
\pgfpathlineto{\pgfqpoint{1.384811in}{1.788095in}}%
\pgfpathlineto{\pgfqpoint{1.387028in}{1.774847in}}%
\pgfpathlineto{\pgfqpoint{1.380577in}{1.761599in}}%
\pgfpathlineto{\pgfqpoint{1.382030in}{1.748351in}}%
\pgfpathlineto{\pgfqpoint{1.377553in}{1.735103in}}%
\pgfpathlineto{\pgfqpoint{1.376195in}{1.721855in}}%
\pgfpathlineto{\pgfqpoint{1.375676in}{1.708607in}}%
\pgfpathlineto{\pgfqpoint{1.375319in}{1.707685in}}%
\pgfpathlineto{\pgfqpoint{1.371591in}{1.695359in}}%
\pgfpathlineto{\pgfqpoint{1.373741in}{1.682111in}}%
\pgfpathlineto{\pgfqpoint{1.368403in}{1.668863in}}%
\pgfpathlineto{\pgfqpoint{1.370855in}{1.655615in}}%
\pgfpathlineto{\pgfqpoint{1.366706in}{1.642367in}}%
\pgfpathlineto{\pgfqpoint{1.366912in}{1.629119in}}%
\pgfpathlineto{\pgfqpoint{1.366379in}{1.615871in}}%
\pgfpathlineto{\pgfqpoint{1.363355in}{1.602623in}}%
\pgfpathlineto{\pgfqpoint{1.366760in}{1.589375in}}%
\pgfpathlineto{\pgfqpoint{1.361460in}{1.576127in}}%
\pgfpathlineto{\pgfqpoint{1.365761in}{1.562879in}}%
\pgfpathlineto{\pgfqpoint{1.361332in}{1.549631in}}%
\pgfpathlineto{\pgfqpoint{1.363135in}{1.536383in}}%
\pgfpathlineto{\pgfqpoint{1.362570in}{1.523135in}}%
\pgfpathlineto{\pgfqpoint{1.360070in}{1.509887in}}%
\pgfpathlineto{\pgfqpoint{1.364506in}{1.496639in}}%
\pgfpathlineto{\pgfqpoint{1.359387in}{1.483391in}}%
\pgfpathlineto{\pgfqpoint{1.366846in}{1.470143in}}%
\pgfpathlineto{\pgfqpoint{1.362833in}{1.456895in}}%
\pgfpathlineto{\pgfqpoint{1.368001in}{1.443647in}}%
\pgfpathlineto{\pgfqpoint{1.369265in}{1.430399in}}%
\pgfpathlineto{\pgfqpoint{1.368564in}{1.417151in}}%
\pgfpathlineto{\pgfqpoint{1.375319in}{1.405375in}}%
\pgfpathlineto{\pgfqpoint{1.376535in}{1.403903in}}%
\pgfpathlineto{\pgfqpoint{1.375319in}{1.401645in}}%
\pgfpathlineto{\pgfqpoint{1.370816in}{1.390655in}}%
\pgfpathlineto{\pgfqpoint{1.375319in}{1.385103in}}%
\pgfpathlineto{\pgfqpoint{1.383966in}{1.377407in}}%
\pgfpathlineto{\pgfqpoint{1.379517in}{1.364159in}}%
\pgfpathlineto{\pgfqpoint{1.391932in}{1.350911in}}%
\pgfpathlineto{\pgfqpoint{1.399429in}{1.337663in}}%
\pgfpathlineto{\pgfqpoint{1.400603in}{1.332605in}}%
\pgfpathlineto{\pgfqpoint{1.403586in}{1.324415in}}%
\pgfpathlineto{\pgfqpoint{1.425886in}{1.315633in}}%
\pgfpathclose%
\pgfpathmoveto{\pgfqpoint{1.435071in}{1.350911in}}%
\pgfpathlineto{\pgfqpoint{1.425886in}{1.353546in}}%
\pgfpathlineto{\pgfqpoint{1.410910in}{1.364159in}}%
\pgfpathlineto{\pgfqpoint{1.414138in}{1.377407in}}%
\pgfpathlineto{\pgfqpoint{1.400603in}{1.384576in}}%
\pgfpathlineto{\pgfqpoint{1.393605in}{1.390655in}}%
\pgfpathlineto{\pgfqpoint{1.400428in}{1.403903in}}%
\pgfpathlineto{\pgfqpoint{1.388278in}{1.417151in}}%
\pgfpathlineto{\pgfqpoint{1.389591in}{1.430399in}}%
\pgfpathlineto{\pgfqpoint{1.385916in}{1.443647in}}%
\pgfpathlineto{\pgfqpoint{1.379124in}{1.456895in}}%
\pgfpathlineto{\pgfqpoint{1.383257in}{1.470143in}}%
\pgfpathlineto{\pgfqpoint{1.375319in}{1.480116in}}%
\pgfpathlineto{\pgfqpoint{1.373397in}{1.483391in}}%
\pgfpathlineto{\pgfqpoint{1.375319in}{1.488500in}}%
\pgfpathlineto{\pgfqpoint{1.379284in}{1.496639in}}%
\pgfpathlineto{\pgfqpoint{1.375319in}{1.503946in}}%
\pgfpathlineto{\pgfqpoint{1.373017in}{1.509887in}}%
\pgfpathlineto{\pgfqpoint{1.375319in}{1.521542in}}%
\pgfpathlineto{\pgfqpoint{1.375737in}{1.523135in}}%
\pgfpathlineto{\pgfqpoint{1.375319in}{1.535709in}}%
\pgfpathlineto{\pgfqpoint{1.375304in}{1.536383in}}%
\pgfpathlineto{\pgfqpoint{1.373507in}{1.549631in}}%
\pgfpathlineto{\pgfqpoint{1.375319in}{1.555646in}}%
\pgfpathlineto{\pgfqpoint{1.377990in}{1.562879in}}%
\pgfpathlineto{\pgfqpoint{1.375319in}{1.568750in}}%
\pgfpathlineto{\pgfqpoint{1.372775in}{1.576127in}}%
\pgfpathlineto{\pgfqpoint{1.375319in}{1.582580in}}%
\pgfpathlineto{\pgfqpoint{1.378734in}{1.589375in}}%
\pgfpathlineto{\pgfqpoint{1.375319in}{1.598022in}}%
\pgfpathlineto{\pgfqpoint{1.373980in}{1.602623in}}%
\pgfpathlineto{\pgfqpoint{1.375319in}{1.608279in}}%
\pgfpathlineto{\pgfqpoint{1.377670in}{1.615871in}}%
\pgfpathlineto{\pgfqpoint{1.377530in}{1.629119in}}%
\pgfpathlineto{\pgfqpoint{1.377349in}{1.642367in}}%
\pgfpathlineto{\pgfqpoint{1.382044in}{1.655615in}}%
\pgfpathlineto{\pgfqpoint{1.378703in}{1.668863in}}%
\pgfpathlineto{\pgfqpoint{1.385281in}{1.682111in}}%
\pgfpathlineto{\pgfqpoint{1.382039in}{1.695359in}}%
\pgfpathlineto{\pgfqpoint{1.387275in}{1.708607in}}%
\pgfpathlineto{\pgfqpoint{1.387113in}{1.721855in}}%
\pgfpathlineto{\pgfqpoint{1.388732in}{1.735103in}}%
\pgfpathlineto{\pgfqpoint{1.392686in}{1.748351in}}%
\pgfpathlineto{\pgfqpoint{1.391218in}{1.761599in}}%
\pgfpathlineto{\pgfqpoint{1.397489in}{1.774847in}}%
\pgfpathlineto{\pgfqpoint{1.394882in}{1.788095in}}%
\pgfpathlineto{\pgfqpoint{1.400331in}{1.801343in}}%
\pgfpathlineto{\pgfqpoint{1.399764in}{1.814591in}}%
\pgfpathlineto{\pgfqpoint{1.400603in}{1.818422in}}%
\pgfpathlineto{\pgfqpoint{1.403829in}{1.827839in}}%
\pgfpathlineto{\pgfqpoint{1.408518in}{1.841087in}}%
\pgfpathlineto{\pgfqpoint{1.408712in}{1.854335in}}%
\pgfpathlineto{\pgfqpoint{1.417510in}{1.867583in}}%
\pgfpathlineto{\pgfqpoint{1.415752in}{1.880831in}}%
\pgfpathlineto{\pgfqpoint{1.425245in}{1.894079in}}%
\pgfpathlineto{\pgfqpoint{1.423784in}{1.907327in}}%
\pgfpathlineto{\pgfqpoint{1.425886in}{1.911874in}}%
\pgfpathlineto{\pgfqpoint{1.437474in}{1.920575in}}%
\pgfpathlineto{\pgfqpoint{1.442317in}{1.933823in}}%
\pgfpathlineto{\pgfqpoint{1.449983in}{1.947071in}}%
\pgfpathlineto{\pgfqpoint{1.451169in}{1.947917in}}%
\pgfpathlineto{\pgfqpoint{1.452426in}{1.947071in}}%
\pgfpathlineto{\pgfqpoint{1.464034in}{1.933823in}}%
\pgfpathlineto{\pgfqpoint{1.466312in}{1.920575in}}%
\pgfpathlineto{\pgfqpoint{1.476452in}{1.915311in}}%
\pgfpathlineto{\pgfqpoint{1.481780in}{1.907327in}}%
\pgfpathlineto{\pgfqpoint{1.479176in}{1.894079in}}%
\pgfpathlineto{\pgfqpoint{1.488943in}{1.880831in}}%
\pgfpathlineto{\pgfqpoint{1.488923in}{1.867583in}}%
\pgfpathlineto{\pgfqpoint{1.494585in}{1.854335in}}%
\pgfpathlineto{\pgfqpoint{1.499656in}{1.841087in}}%
\pgfpathlineto{\pgfqpoint{1.499321in}{1.827839in}}%
\pgfpathlineto{\pgfqpoint{1.501735in}{1.824339in}}%
\pgfpathlineto{\pgfqpoint{1.506180in}{1.814591in}}%
\pgfpathlineto{\pgfqpoint{1.503715in}{1.801343in}}%
\pgfpathlineto{\pgfqpoint{1.509408in}{1.788095in}}%
\pgfpathlineto{\pgfqpoint{1.508110in}{1.774847in}}%
\pgfpathlineto{\pgfqpoint{1.511332in}{1.761599in}}%
\pgfpathlineto{\pgfqpoint{1.514129in}{1.748351in}}%
\pgfpathlineto{\pgfqpoint{1.513622in}{1.735103in}}%
\pgfpathlineto{\pgfqpoint{1.519827in}{1.721855in}}%
\pgfpathlineto{\pgfqpoint{1.516708in}{1.708607in}}%
\pgfpathlineto{\pgfqpoint{1.523420in}{1.695359in}}%
\pgfpathlineto{\pgfqpoint{1.521090in}{1.682111in}}%
\pgfpathlineto{\pgfqpoint{1.523989in}{1.668863in}}%
\pgfpathlineto{\pgfqpoint{1.525341in}{1.655615in}}%
\pgfpathlineto{\pgfqpoint{1.523493in}{1.642367in}}%
\pgfpathlineto{\pgfqpoint{1.527018in}{1.633893in}}%
\pgfpathlineto{\pgfqpoint{1.528621in}{1.629119in}}%
\pgfpathlineto{\pgfqpoint{1.527018in}{1.623148in}}%
\pgfpathlineto{\pgfqpoint{1.524423in}{1.615871in}}%
\pgfpathlineto{\pgfqpoint{1.527018in}{1.610786in}}%
\pgfpathlineto{\pgfqpoint{1.530268in}{1.602623in}}%
\pgfpathlineto{\pgfqpoint{1.527018in}{1.589761in}}%
\pgfpathlineto{\pgfqpoint{1.526883in}{1.589375in}}%
\pgfpathlineto{\pgfqpoint{1.527018in}{1.588924in}}%
\pgfpathlineto{\pgfqpoint{1.529927in}{1.576127in}}%
\pgfpathlineto{\pgfqpoint{1.529732in}{1.562879in}}%
\pgfpathlineto{\pgfqpoint{1.527953in}{1.549631in}}%
\pgfpathlineto{\pgfqpoint{1.531470in}{1.536383in}}%
\pgfpathlineto{\pgfqpoint{1.527018in}{1.526015in}}%
\pgfpathlineto{\pgfqpoint{1.525364in}{1.523135in}}%
\pgfpathlineto{\pgfqpoint{1.527018in}{1.519907in}}%
\pgfpathlineto{\pgfqpoint{1.531012in}{1.509887in}}%
\pgfpathlineto{\pgfqpoint{1.527018in}{1.500940in}}%
\pgfpathlineto{\pgfqpoint{1.524333in}{1.496639in}}%
\pgfpathlineto{\pgfqpoint{1.527018in}{1.487123in}}%
\pgfpathlineto{\pgfqpoint{1.527814in}{1.483391in}}%
\pgfpathlineto{\pgfqpoint{1.527018in}{1.480137in}}%
\pgfpathlineto{\pgfqpoint{1.523492in}{1.470143in}}%
\pgfpathlineto{\pgfqpoint{1.520210in}{1.456895in}}%
\pgfpathlineto{\pgfqpoint{1.521611in}{1.443647in}}%
\pgfpathlineto{\pgfqpoint{1.510959in}{1.430399in}}%
\pgfpathlineto{\pgfqpoint{1.516797in}{1.417151in}}%
\pgfpathlineto{\pgfqpoint{1.502475in}{1.403903in}}%
\pgfpathlineto{\pgfqpoint{1.507647in}{1.390655in}}%
\pgfpathlineto{\pgfqpoint{1.501735in}{1.383686in}}%
\pgfpathlineto{\pgfqpoint{1.491988in}{1.377407in}}%
\pgfpathlineto{\pgfqpoint{1.487330in}{1.364159in}}%
\pgfpathlineto{\pgfqpoint{1.476452in}{1.351348in}}%
\pgfpathlineto{\pgfqpoint{1.474954in}{1.350911in}}%
\pgfpathlineto{\pgfqpoint{1.451169in}{1.346761in}}%
\pgfpathlineto{\pgfqpoint{1.435071in}{1.350911in}}%
\pgfpathclose%
\pgfusepath{fill}%
\end{pgfscope}%
\begin{pgfscope}%
\pgfpathrectangle{\pgfqpoint{0.540975in}{0.476543in}}{\pgfqpoint{1.820388in}{1.656000in}}%
\pgfusepath{clip}%
\pgfsetbuttcap%
\pgfsetroundjoin%
\definecolor{currentfill}{rgb}{0.894425,0.296117,0.202461}%
\pgfsetfillcolor{currentfill}%
\pgfsetlinewidth{0.000000pt}%
\definecolor{currentstroke}{rgb}{0.000000,0.000000,0.000000}%
\pgfsetstrokecolor{currentstroke}%
\pgfsetdash{}{0pt}%
\pgfpathmoveto{\pgfqpoint{1.451169in}{1.346761in}}%
\pgfpathlineto{\pgfqpoint{1.474954in}{1.350911in}}%
\pgfpathlineto{\pgfqpoint{1.476452in}{1.351348in}}%
\pgfpathlineto{\pgfqpoint{1.487330in}{1.364159in}}%
\pgfpathlineto{\pgfqpoint{1.491988in}{1.377407in}}%
\pgfpathlineto{\pgfqpoint{1.501735in}{1.383686in}}%
\pgfpathlineto{\pgfqpoint{1.507647in}{1.390655in}}%
\pgfpathlineto{\pgfqpoint{1.502475in}{1.403903in}}%
\pgfpathlineto{\pgfqpoint{1.516797in}{1.417151in}}%
\pgfpathlineto{\pgfqpoint{1.510959in}{1.430399in}}%
\pgfpathlineto{\pgfqpoint{1.521611in}{1.443647in}}%
\pgfpathlineto{\pgfqpoint{1.520210in}{1.456895in}}%
\pgfpathlineto{\pgfqpoint{1.523492in}{1.470143in}}%
\pgfpathlineto{\pgfqpoint{1.527018in}{1.480137in}}%
\pgfpathlineto{\pgfqpoint{1.527814in}{1.483391in}}%
\pgfpathlineto{\pgfqpoint{1.527018in}{1.487123in}}%
\pgfpathlineto{\pgfqpoint{1.524333in}{1.496639in}}%
\pgfpathlineto{\pgfqpoint{1.527018in}{1.500940in}}%
\pgfpathlineto{\pgfqpoint{1.531012in}{1.509887in}}%
\pgfpathlineto{\pgfqpoint{1.527018in}{1.519907in}}%
\pgfpathlineto{\pgfqpoint{1.525364in}{1.523135in}}%
\pgfpathlineto{\pgfqpoint{1.527018in}{1.526015in}}%
\pgfpathlineto{\pgfqpoint{1.531470in}{1.536383in}}%
\pgfpathlineto{\pgfqpoint{1.527953in}{1.549631in}}%
\pgfpathlineto{\pgfqpoint{1.529732in}{1.562879in}}%
\pgfpathlineto{\pgfqpoint{1.529927in}{1.576127in}}%
\pgfpathlineto{\pgfqpoint{1.527018in}{1.588924in}}%
\pgfpathlineto{\pgfqpoint{1.526883in}{1.589375in}}%
\pgfpathlineto{\pgfqpoint{1.527018in}{1.589761in}}%
\pgfpathlineto{\pgfqpoint{1.530268in}{1.602623in}}%
\pgfpathlineto{\pgfqpoint{1.527018in}{1.610786in}}%
\pgfpathlineto{\pgfqpoint{1.524423in}{1.615871in}}%
\pgfpathlineto{\pgfqpoint{1.527018in}{1.623148in}}%
\pgfpathlineto{\pgfqpoint{1.528621in}{1.629119in}}%
\pgfpathlineto{\pgfqpoint{1.527018in}{1.633893in}}%
\pgfpathlineto{\pgfqpoint{1.523493in}{1.642367in}}%
\pgfpathlineto{\pgfqpoint{1.525341in}{1.655615in}}%
\pgfpathlineto{\pgfqpoint{1.523989in}{1.668863in}}%
\pgfpathlineto{\pgfqpoint{1.521090in}{1.682111in}}%
\pgfpathlineto{\pgfqpoint{1.523420in}{1.695359in}}%
\pgfpathlineto{\pgfqpoint{1.516708in}{1.708607in}}%
\pgfpathlineto{\pgfqpoint{1.519827in}{1.721855in}}%
\pgfpathlineto{\pgfqpoint{1.513622in}{1.735103in}}%
\pgfpathlineto{\pgfqpoint{1.514129in}{1.748351in}}%
\pgfpathlineto{\pgfqpoint{1.511332in}{1.761599in}}%
\pgfpathlineto{\pgfqpoint{1.508110in}{1.774847in}}%
\pgfpathlineto{\pgfqpoint{1.509408in}{1.788095in}}%
\pgfpathlineto{\pgfqpoint{1.503715in}{1.801343in}}%
\pgfpathlineto{\pgfqpoint{1.506180in}{1.814591in}}%
\pgfpathlineto{\pgfqpoint{1.501735in}{1.824339in}}%
\pgfpathlineto{\pgfqpoint{1.499321in}{1.827839in}}%
\pgfpathlineto{\pgfqpoint{1.499656in}{1.841087in}}%
\pgfpathlineto{\pgfqpoint{1.494585in}{1.854335in}}%
\pgfpathlineto{\pgfqpoint{1.488923in}{1.867583in}}%
\pgfpathlineto{\pgfqpoint{1.488943in}{1.880831in}}%
\pgfpathlineto{\pgfqpoint{1.479176in}{1.894079in}}%
\pgfpathlineto{\pgfqpoint{1.481780in}{1.907327in}}%
\pgfpathlineto{\pgfqpoint{1.476452in}{1.915311in}}%
\pgfpathlineto{\pgfqpoint{1.466312in}{1.920575in}}%
\pgfpathlineto{\pgfqpoint{1.464034in}{1.933823in}}%
\pgfpathlineto{\pgfqpoint{1.452426in}{1.947071in}}%
\pgfpathlineto{\pgfqpoint{1.451169in}{1.947917in}}%
\pgfpathlineto{\pgfqpoint{1.449983in}{1.947071in}}%
\pgfpathlineto{\pgfqpoint{1.442317in}{1.933823in}}%
\pgfpathlineto{\pgfqpoint{1.437474in}{1.920575in}}%
\pgfpathlineto{\pgfqpoint{1.425886in}{1.911874in}}%
\pgfpathlineto{\pgfqpoint{1.423784in}{1.907327in}}%
\pgfpathlineto{\pgfqpoint{1.425245in}{1.894079in}}%
\pgfpathlineto{\pgfqpoint{1.415752in}{1.880831in}}%
\pgfpathlineto{\pgfqpoint{1.417510in}{1.867583in}}%
\pgfpathlineto{\pgfqpoint{1.408712in}{1.854335in}}%
\pgfpathlineto{\pgfqpoint{1.408518in}{1.841087in}}%
\pgfpathlineto{\pgfqpoint{1.403829in}{1.827839in}}%
\pgfpathlineto{\pgfqpoint{1.400603in}{1.818422in}}%
\pgfpathlineto{\pgfqpoint{1.399764in}{1.814591in}}%
\pgfpathlineto{\pgfqpoint{1.400331in}{1.801343in}}%
\pgfpathlineto{\pgfqpoint{1.394882in}{1.788095in}}%
\pgfpathlineto{\pgfqpoint{1.397489in}{1.774847in}}%
\pgfpathlineto{\pgfqpoint{1.391218in}{1.761599in}}%
\pgfpathlineto{\pgfqpoint{1.392686in}{1.748351in}}%
\pgfpathlineto{\pgfqpoint{1.388732in}{1.735103in}}%
\pgfpathlineto{\pgfqpoint{1.387113in}{1.721855in}}%
\pgfpathlineto{\pgfqpoint{1.387275in}{1.708607in}}%
\pgfpathlineto{\pgfqpoint{1.382039in}{1.695359in}}%
\pgfpathlineto{\pgfqpoint{1.385281in}{1.682111in}}%
\pgfpathlineto{\pgfqpoint{1.378703in}{1.668863in}}%
\pgfpathlineto{\pgfqpoint{1.382044in}{1.655615in}}%
\pgfpathlineto{\pgfqpoint{1.377349in}{1.642367in}}%
\pgfpathlineto{\pgfqpoint{1.377530in}{1.629119in}}%
\pgfpathlineto{\pgfqpoint{1.377670in}{1.615871in}}%
\pgfpathlineto{\pgfqpoint{1.375319in}{1.608279in}}%
\pgfpathlineto{\pgfqpoint{1.373980in}{1.602623in}}%
\pgfpathlineto{\pgfqpoint{1.375319in}{1.598022in}}%
\pgfpathlineto{\pgfqpoint{1.378734in}{1.589375in}}%
\pgfpathlineto{\pgfqpoint{1.375319in}{1.582580in}}%
\pgfpathlineto{\pgfqpoint{1.372775in}{1.576127in}}%
\pgfpathlineto{\pgfqpoint{1.375319in}{1.568750in}}%
\pgfpathlineto{\pgfqpoint{1.377990in}{1.562879in}}%
\pgfpathlineto{\pgfqpoint{1.375319in}{1.555646in}}%
\pgfpathlineto{\pgfqpoint{1.373507in}{1.549631in}}%
\pgfpathlineto{\pgfqpoint{1.375304in}{1.536383in}}%
\pgfpathlineto{\pgfqpoint{1.375319in}{1.535709in}}%
\pgfpathlineto{\pgfqpoint{1.375737in}{1.523135in}}%
\pgfpathlineto{\pgfqpoint{1.375319in}{1.521542in}}%
\pgfpathlineto{\pgfqpoint{1.373017in}{1.509887in}}%
\pgfpathlineto{\pgfqpoint{1.375319in}{1.503946in}}%
\pgfpathlineto{\pgfqpoint{1.379284in}{1.496639in}}%
\pgfpathlineto{\pgfqpoint{1.375319in}{1.488500in}}%
\pgfpathlineto{\pgfqpoint{1.373397in}{1.483391in}}%
\pgfpathlineto{\pgfqpoint{1.375319in}{1.480116in}}%
\pgfpathlineto{\pgfqpoint{1.383257in}{1.470143in}}%
\pgfpathlineto{\pgfqpoint{1.379124in}{1.456895in}}%
\pgfpathlineto{\pgfqpoint{1.385916in}{1.443647in}}%
\pgfpathlineto{\pgfqpoint{1.389591in}{1.430399in}}%
\pgfpathlineto{\pgfqpoint{1.388278in}{1.417151in}}%
\pgfpathlineto{\pgfqpoint{1.400428in}{1.403903in}}%
\pgfpathlineto{\pgfqpoint{1.393605in}{1.390655in}}%
\pgfpathlineto{\pgfqpoint{1.400603in}{1.384576in}}%
\pgfpathlineto{\pgfqpoint{1.414138in}{1.377407in}}%
\pgfpathlineto{\pgfqpoint{1.410910in}{1.364159in}}%
\pgfpathlineto{\pgfqpoint{1.425886in}{1.353546in}}%
\pgfpathlineto{\pgfqpoint{1.435071in}{1.350911in}}%
\pgfpathlineto{\pgfqpoint{1.451169in}{1.346761in}}%
\pgfpathclose%
\pgfpathmoveto{\pgfqpoint{1.434548in}{1.390655in}}%
\pgfpathlineto{\pgfqpoint{1.451169in}{1.398129in}}%
\pgfpathlineto{\pgfqpoint{1.466246in}{1.390655in}}%
\pgfpathlineto{\pgfqpoint{1.451169in}{1.387654in}}%
\pgfpathlineto{\pgfqpoint{1.434548in}{1.390655in}}%
\pgfpathclose%
\pgfpathmoveto{\pgfqpoint{1.415731in}{1.417151in}}%
\pgfpathlineto{\pgfqpoint{1.419031in}{1.430399in}}%
\pgfpathlineto{\pgfqpoint{1.409162in}{1.443647in}}%
\pgfpathlineto{\pgfqpoint{1.400603in}{1.454791in}}%
\pgfpathlineto{\pgfqpoint{1.399555in}{1.456895in}}%
\pgfpathlineto{\pgfqpoint{1.400603in}{1.461814in}}%
\pgfpathlineto{\pgfqpoint{1.403187in}{1.470143in}}%
\pgfpathlineto{\pgfqpoint{1.400603in}{1.472092in}}%
\pgfpathlineto{\pgfqpoint{1.391288in}{1.483391in}}%
\pgfpathlineto{\pgfqpoint{1.397416in}{1.496639in}}%
\pgfpathlineto{\pgfqpoint{1.389207in}{1.509887in}}%
\pgfpathlineto{\pgfqpoint{1.392899in}{1.523135in}}%
\pgfpathlineto{\pgfqpoint{1.390979in}{1.536383in}}%
\pgfpathlineto{\pgfqpoint{1.388872in}{1.549631in}}%
\pgfpathlineto{\pgfqpoint{1.392913in}{1.562879in}}%
\pgfpathlineto{\pgfqpoint{1.386693in}{1.576127in}}%
\pgfpathlineto{\pgfqpoint{1.393151in}{1.589375in}}%
\pgfpathlineto{\pgfqpoint{1.387249in}{1.602623in}}%
\pgfpathlineto{\pgfqpoint{1.391565in}{1.615871in}}%
\pgfpathlineto{\pgfqpoint{1.390440in}{1.629119in}}%
\pgfpathlineto{\pgfqpoint{1.390482in}{1.642367in}}%
\pgfpathlineto{\pgfqpoint{1.394409in}{1.655615in}}%
\pgfpathlineto{\pgfqpoint{1.390911in}{1.668863in}}%
\pgfpathlineto{\pgfqpoint{1.397232in}{1.682111in}}%
\pgfpathlineto{\pgfqpoint{1.393454in}{1.695359in}}%
\pgfpathlineto{\pgfqpoint{1.398874in}{1.708607in}}%
\pgfpathlineto{\pgfqpoint{1.398030in}{1.721855in}}%
\pgfpathlineto{\pgfqpoint{1.399911in}{1.735103in}}%
\pgfpathlineto{\pgfqpoint{1.400603in}{1.737674in}}%
\pgfpathlineto{\pgfqpoint{1.404745in}{1.748351in}}%
\pgfpathlineto{\pgfqpoint{1.402550in}{1.761599in}}%
\pgfpathlineto{\pgfqpoint{1.411736in}{1.774847in}}%
\pgfpathlineto{\pgfqpoint{1.407239in}{1.788095in}}%
\pgfpathlineto{\pgfqpoint{1.415722in}{1.801343in}}%
\pgfpathlineto{\pgfqpoint{1.413829in}{1.814591in}}%
\pgfpathlineto{\pgfqpoint{1.418925in}{1.827839in}}%
\pgfpathlineto{\pgfqpoint{1.422577in}{1.841087in}}%
\pgfpathlineto{\pgfqpoint{1.423177in}{1.854335in}}%
\pgfpathlineto{\pgfqpoint{1.425886in}{1.858558in}}%
\pgfpathlineto{\pgfqpoint{1.440309in}{1.867583in}}%
\pgfpathlineto{\pgfqpoint{1.435640in}{1.880831in}}%
\pgfpathlineto{\pgfqpoint{1.451169in}{1.888847in}}%
\pgfpathlineto{\pgfqpoint{1.469981in}{1.880831in}}%
\pgfpathlineto{\pgfqpoint{1.466506in}{1.867583in}}%
\pgfpathlineto{\pgfqpoint{1.476452in}{1.861445in}}%
\pgfpathlineto{\pgfqpoint{1.479799in}{1.854335in}}%
\pgfpathlineto{\pgfqpoint{1.483816in}{1.841087in}}%
\pgfpathlineto{\pgfqpoint{1.483967in}{1.827839in}}%
\pgfpathlineto{\pgfqpoint{1.492741in}{1.814591in}}%
\pgfpathlineto{\pgfqpoint{1.488513in}{1.801343in}}%
\pgfpathlineto{\pgfqpoint{1.497546in}{1.788095in}}%
\pgfpathlineto{\pgfqpoint{1.494505in}{1.774847in}}%
\pgfpathlineto{\pgfqpoint{1.500030in}{1.761599in}}%
\pgfpathlineto{\pgfqpoint{1.501735in}{1.754804in}}%
\pgfpathlineto{\pgfqpoint{1.502838in}{1.748351in}}%
\pgfpathlineto{\pgfqpoint{1.502413in}{1.735103in}}%
\pgfpathlineto{\pgfqpoint{1.508276in}{1.721855in}}%
\pgfpathlineto{\pgfqpoint{1.504851in}{1.708607in}}%
\pgfpathlineto{\pgfqpoint{1.511567in}{1.695359in}}%
\pgfpathlineto{\pgfqpoint{1.508537in}{1.682111in}}%
\pgfpathlineto{\pgfqpoint{1.511706in}{1.668863in}}%
\pgfpathlineto{\pgfqpoint{1.512220in}{1.655615in}}%
\pgfpathlineto{\pgfqpoint{1.510558in}{1.642367in}}%
\pgfpathlineto{\pgfqpoint{1.515537in}{1.629119in}}%
\pgfpathlineto{\pgfqpoint{1.510545in}{1.615871in}}%
\pgfpathlineto{\pgfqpoint{1.517187in}{1.602623in}}%
\pgfpathlineto{\pgfqpoint{1.511891in}{1.589375in}}%
\pgfpathlineto{\pgfqpoint{1.516020in}{1.576127in}}%
\pgfpathlineto{\pgfqpoint{1.514570in}{1.562879in}}%
\pgfpathlineto{\pgfqpoint{1.512489in}{1.549631in}}%
\pgfpathlineto{\pgfqpoint{1.516158in}{1.536383in}}%
\pgfpathlineto{\pgfqpoint{1.508425in}{1.523135in}}%
\pgfpathlineto{\pgfqpoint{1.514933in}{1.509887in}}%
\pgfpathlineto{\pgfqpoint{1.505901in}{1.496639in}}%
\pgfpathlineto{\pgfqpoint{1.509833in}{1.483391in}}%
\pgfpathlineto{\pgfqpoint{1.503394in}{1.470143in}}%
\pgfpathlineto{\pgfqpoint{1.501735in}{1.462608in}}%
\pgfpathlineto{\pgfqpoint{1.499789in}{1.456895in}}%
\pgfpathlineto{\pgfqpoint{1.498895in}{1.443647in}}%
\pgfpathlineto{\pgfqpoint{1.481733in}{1.430399in}}%
\pgfpathlineto{\pgfqpoint{1.489202in}{1.417151in}}%
\pgfpathlineto{\pgfqpoint{1.476452in}{1.410315in}}%
\pgfpathlineto{\pgfqpoint{1.451169in}{1.406028in}}%
\pgfpathlineto{\pgfqpoint{1.425886in}{1.410959in}}%
\pgfpathlineto{\pgfqpoint{1.415731in}{1.417151in}}%
\pgfpathclose%
\pgfusepath{fill}%
\end{pgfscope}%
\begin{pgfscope}%
\pgfpathrectangle{\pgfqpoint{0.540975in}{0.476543in}}{\pgfqpoint{1.820388in}{1.656000in}}%
\pgfusepath{clip}%
\pgfsetbuttcap%
\pgfsetroundjoin%
\definecolor{currentfill}{rgb}{0.845367,0.192926,0.155094}%
\pgfsetfillcolor{currentfill}%
\pgfsetlinewidth{0.000000pt}%
\definecolor{currentstroke}{rgb}{0.000000,0.000000,0.000000}%
\pgfsetstrokecolor{currentstroke}%
\pgfsetdash{}{0pt}%
\pgfpathmoveto{\pgfqpoint{1.451169in}{1.387654in}}%
\pgfpathlineto{\pgfqpoint{1.466246in}{1.390655in}}%
\pgfpathlineto{\pgfqpoint{1.451169in}{1.398129in}}%
\pgfpathlineto{\pgfqpoint{1.434548in}{1.390655in}}%
\pgfpathlineto{\pgfqpoint{1.451169in}{1.387654in}}%
\pgfpathclose%
\pgfpathmoveto{\pgfqpoint{1.425886in}{1.410959in}}%
\pgfpathlineto{\pgfqpoint{1.451169in}{1.406028in}}%
\pgfpathlineto{\pgfqpoint{1.476452in}{1.410315in}}%
\pgfpathlineto{\pgfqpoint{1.489202in}{1.417151in}}%
\pgfpathlineto{\pgfqpoint{1.481733in}{1.430399in}}%
\pgfpathlineto{\pgfqpoint{1.498895in}{1.443647in}}%
\pgfpathlineto{\pgfqpoint{1.499789in}{1.456895in}}%
\pgfpathlineto{\pgfqpoint{1.501735in}{1.462608in}}%
\pgfpathlineto{\pgfqpoint{1.503394in}{1.470143in}}%
\pgfpathlineto{\pgfqpoint{1.509833in}{1.483391in}}%
\pgfpathlineto{\pgfqpoint{1.505901in}{1.496639in}}%
\pgfpathlineto{\pgfqpoint{1.514933in}{1.509887in}}%
\pgfpathlineto{\pgfqpoint{1.508425in}{1.523135in}}%
\pgfpathlineto{\pgfqpoint{1.516158in}{1.536383in}}%
\pgfpathlineto{\pgfqpoint{1.512489in}{1.549631in}}%
\pgfpathlineto{\pgfqpoint{1.514570in}{1.562879in}}%
\pgfpathlineto{\pgfqpoint{1.516020in}{1.576127in}}%
\pgfpathlineto{\pgfqpoint{1.511891in}{1.589375in}}%
\pgfpathlineto{\pgfqpoint{1.517187in}{1.602623in}}%
\pgfpathlineto{\pgfqpoint{1.510545in}{1.615871in}}%
\pgfpathlineto{\pgfqpoint{1.515537in}{1.629119in}}%
\pgfpathlineto{\pgfqpoint{1.510558in}{1.642367in}}%
\pgfpathlineto{\pgfqpoint{1.512220in}{1.655615in}}%
\pgfpathlineto{\pgfqpoint{1.511706in}{1.668863in}}%
\pgfpathlineto{\pgfqpoint{1.508537in}{1.682111in}}%
\pgfpathlineto{\pgfqpoint{1.511567in}{1.695359in}}%
\pgfpathlineto{\pgfqpoint{1.504851in}{1.708607in}}%
\pgfpathlineto{\pgfqpoint{1.508276in}{1.721855in}}%
\pgfpathlineto{\pgfqpoint{1.502413in}{1.735103in}}%
\pgfpathlineto{\pgfqpoint{1.502838in}{1.748351in}}%
\pgfpathlineto{\pgfqpoint{1.501735in}{1.754804in}}%
\pgfpathlineto{\pgfqpoint{1.500030in}{1.761599in}}%
\pgfpathlineto{\pgfqpoint{1.494505in}{1.774847in}}%
\pgfpathlineto{\pgfqpoint{1.497546in}{1.788095in}}%
\pgfpathlineto{\pgfqpoint{1.488513in}{1.801343in}}%
\pgfpathlineto{\pgfqpoint{1.492741in}{1.814591in}}%
\pgfpathlineto{\pgfqpoint{1.483967in}{1.827839in}}%
\pgfpathlineto{\pgfqpoint{1.483816in}{1.841087in}}%
\pgfpathlineto{\pgfqpoint{1.479799in}{1.854335in}}%
\pgfpathlineto{\pgfqpoint{1.476452in}{1.861445in}}%
\pgfpathlineto{\pgfqpoint{1.466506in}{1.867583in}}%
\pgfpathlineto{\pgfqpoint{1.469981in}{1.880831in}}%
\pgfpathlineto{\pgfqpoint{1.451169in}{1.888847in}}%
\pgfpathlineto{\pgfqpoint{1.435640in}{1.880831in}}%
\pgfpathlineto{\pgfqpoint{1.440309in}{1.867583in}}%
\pgfpathlineto{\pgfqpoint{1.425886in}{1.858558in}}%
\pgfpathlineto{\pgfqpoint{1.423177in}{1.854335in}}%
\pgfpathlineto{\pgfqpoint{1.422577in}{1.841087in}}%
\pgfpathlineto{\pgfqpoint{1.418925in}{1.827839in}}%
\pgfpathlineto{\pgfqpoint{1.413829in}{1.814591in}}%
\pgfpathlineto{\pgfqpoint{1.415722in}{1.801343in}}%
\pgfpathlineto{\pgfqpoint{1.407239in}{1.788095in}}%
\pgfpathlineto{\pgfqpoint{1.411736in}{1.774847in}}%
\pgfpathlineto{\pgfqpoint{1.402550in}{1.761599in}}%
\pgfpathlineto{\pgfqpoint{1.404745in}{1.748351in}}%
\pgfpathlineto{\pgfqpoint{1.400603in}{1.737674in}}%
\pgfpathlineto{\pgfqpoint{1.399911in}{1.735103in}}%
\pgfpathlineto{\pgfqpoint{1.398030in}{1.721855in}}%
\pgfpathlineto{\pgfqpoint{1.398874in}{1.708607in}}%
\pgfpathlineto{\pgfqpoint{1.393454in}{1.695359in}}%
\pgfpathlineto{\pgfqpoint{1.397232in}{1.682111in}}%
\pgfpathlineto{\pgfqpoint{1.390911in}{1.668863in}}%
\pgfpathlineto{\pgfqpoint{1.394409in}{1.655615in}}%
\pgfpathlineto{\pgfqpoint{1.390482in}{1.642367in}}%
\pgfpathlineto{\pgfqpoint{1.390440in}{1.629119in}}%
\pgfpathlineto{\pgfqpoint{1.391565in}{1.615871in}}%
\pgfpathlineto{\pgfqpoint{1.387249in}{1.602623in}}%
\pgfpathlineto{\pgfqpoint{1.393151in}{1.589375in}}%
\pgfpathlineto{\pgfqpoint{1.386693in}{1.576127in}}%
\pgfpathlineto{\pgfqpoint{1.392913in}{1.562879in}}%
\pgfpathlineto{\pgfqpoint{1.388872in}{1.549631in}}%
\pgfpathlineto{\pgfqpoint{1.390979in}{1.536383in}}%
\pgfpathlineto{\pgfqpoint{1.392899in}{1.523135in}}%
\pgfpathlineto{\pgfqpoint{1.389207in}{1.509887in}}%
\pgfpathlineto{\pgfqpoint{1.397416in}{1.496639in}}%
\pgfpathlineto{\pgfqpoint{1.391288in}{1.483391in}}%
\pgfpathlineto{\pgfqpoint{1.400603in}{1.472092in}}%
\pgfpathlineto{\pgfqpoint{1.403187in}{1.470143in}}%
\pgfpathlineto{\pgfqpoint{1.400603in}{1.461814in}}%
\pgfpathlineto{\pgfqpoint{1.399555in}{1.456895in}}%
\pgfpathlineto{\pgfqpoint{1.400603in}{1.454791in}}%
\pgfpathlineto{\pgfqpoint{1.409162in}{1.443647in}}%
\pgfpathlineto{\pgfqpoint{1.419031in}{1.430399in}}%
\pgfpathlineto{\pgfqpoint{1.415731in}{1.417151in}}%
\pgfpathlineto{\pgfqpoint{1.425886in}{1.410959in}}%
\pgfpathclose%
\pgfpathmoveto{\pgfqpoint{1.446154in}{1.456895in}}%
\pgfpathlineto{\pgfqpoint{1.444765in}{1.470143in}}%
\pgfpathlineto{\pgfqpoint{1.425886in}{1.475501in}}%
\pgfpathlineto{\pgfqpoint{1.415435in}{1.483391in}}%
\pgfpathlineto{\pgfqpoint{1.424349in}{1.496639in}}%
\pgfpathlineto{\pgfqpoint{1.409288in}{1.509887in}}%
\pgfpathlineto{\pgfqpoint{1.415793in}{1.523135in}}%
\pgfpathlineto{\pgfqpoint{1.409991in}{1.536383in}}%
\pgfpathlineto{\pgfqpoint{1.407314in}{1.549631in}}%
\pgfpathlineto{\pgfqpoint{1.411755in}{1.562879in}}%
\pgfpathlineto{\pgfqpoint{1.401811in}{1.576127in}}%
\pgfpathlineto{\pgfqpoint{1.411463in}{1.589375in}}%
\pgfpathlineto{\pgfqpoint{1.401064in}{1.602623in}}%
\pgfpathlineto{\pgfqpoint{1.408289in}{1.615871in}}%
\pgfpathlineto{\pgfqpoint{1.404847in}{1.629119in}}%
\pgfpathlineto{\pgfqpoint{1.405389in}{1.642367in}}%
\pgfpathlineto{\pgfqpoint{1.410055in}{1.655615in}}%
\pgfpathlineto{\pgfqpoint{1.404548in}{1.668863in}}%
\pgfpathlineto{\pgfqpoint{1.413757in}{1.682111in}}%
\pgfpathlineto{\pgfqpoint{1.407168in}{1.695359in}}%
\pgfpathlineto{\pgfqpoint{1.415881in}{1.708607in}}%
\pgfpathlineto{\pgfqpoint{1.413280in}{1.721855in}}%
\pgfpathlineto{\pgfqpoint{1.416938in}{1.735103in}}%
\pgfpathlineto{\pgfqpoint{1.420859in}{1.748351in}}%
\pgfpathlineto{\pgfqpoint{1.419038in}{1.761599in}}%
\pgfpathlineto{\pgfqpoint{1.425886in}{1.772123in}}%
\pgfpathlineto{\pgfqpoint{1.430393in}{1.774847in}}%
\pgfpathlineto{\pgfqpoint{1.425886in}{1.779280in}}%
\pgfpathlineto{\pgfqpoint{1.422602in}{1.788095in}}%
\pgfpathlineto{\pgfqpoint{1.425886in}{1.793157in}}%
\pgfpathlineto{\pgfqpoint{1.440625in}{1.801343in}}%
\pgfpathlineto{\pgfqpoint{1.432219in}{1.814591in}}%
\pgfpathlineto{\pgfqpoint{1.448953in}{1.827839in}}%
\pgfpathlineto{\pgfqpoint{1.451169in}{1.833975in}}%
\pgfpathlineto{\pgfqpoint{1.453497in}{1.827839in}}%
\pgfpathlineto{\pgfqpoint{1.476452in}{1.815037in}}%
\pgfpathlineto{\pgfqpoint{1.476737in}{1.814591in}}%
\pgfpathlineto{\pgfqpoint{1.476452in}{1.813745in}}%
\pgfpathlineto{\pgfqpoint{1.463482in}{1.801343in}}%
\pgfpathlineto{\pgfqpoint{1.476452in}{1.795230in}}%
\pgfpathlineto{\pgfqpoint{1.481370in}{1.788095in}}%
\pgfpathlineto{\pgfqpoint{1.477187in}{1.774847in}}%
\pgfpathlineto{\pgfqpoint{1.483447in}{1.761599in}}%
\pgfpathlineto{\pgfqpoint{1.485449in}{1.748351in}}%
\pgfpathlineto{\pgfqpoint{1.485366in}{1.735103in}}%
\pgfpathlineto{\pgfqpoint{1.493713in}{1.721855in}}%
\pgfpathlineto{\pgfqpoint{1.487963in}{1.708607in}}%
\pgfpathlineto{\pgfqpoint{1.498523in}{1.695359in}}%
\pgfpathlineto{\pgfqpoint{1.492514in}{1.682111in}}%
\pgfpathlineto{\pgfqpoint{1.498105in}{1.668863in}}%
\pgfpathlineto{\pgfqpoint{1.497474in}{1.655615in}}%
\pgfpathlineto{\pgfqpoint{1.495320in}{1.642367in}}%
\pgfpathlineto{\pgfqpoint{1.501735in}{1.629813in}}%
\pgfpathlineto{\pgfqpoint{1.501974in}{1.629119in}}%
\pgfpathlineto{\pgfqpoint{1.501735in}{1.628511in}}%
\pgfpathlineto{\pgfqpoint{1.493751in}{1.615871in}}%
\pgfpathlineto{\pgfqpoint{1.501735in}{1.605468in}}%
\pgfpathlineto{\pgfqpoint{1.503138in}{1.602623in}}%
\pgfpathlineto{\pgfqpoint{1.501735in}{1.599492in}}%
\pgfpathlineto{\pgfqpoint{1.493947in}{1.589375in}}%
\pgfpathlineto{\pgfqpoint{1.500961in}{1.576127in}}%
\pgfpathlineto{\pgfqpoint{1.496491in}{1.562879in}}%
\pgfpathlineto{\pgfqpoint{1.493856in}{1.549631in}}%
\pgfpathlineto{\pgfqpoint{1.497918in}{1.536383in}}%
\pgfpathlineto{\pgfqpoint{1.485509in}{1.523135in}}%
\pgfpathlineto{\pgfqpoint{1.495058in}{1.509887in}}%
\pgfpathlineto{\pgfqpoint{1.478792in}{1.496639in}}%
\pgfpathlineto{\pgfqpoint{1.485667in}{1.483391in}}%
\pgfpathlineto{\pgfqpoint{1.476452in}{1.474394in}}%
\pgfpathlineto{\pgfqpoint{1.459986in}{1.470143in}}%
\pgfpathlineto{\pgfqpoint{1.455181in}{1.456895in}}%
\pgfpathlineto{\pgfqpoint{1.451169in}{1.453939in}}%
\pgfpathlineto{\pgfqpoint{1.446154in}{1.456895in}}%
\pgfpathclose%
\pgfusepath{fill}%
\end{pgfscope}%
\begin{pgfscope}%
\pgfpathrectangle{\pgfqpoint{0.540975in}{0.476543in}}{\pgfqpoint{1.820388in}{1.656000in}}%
\pgfusepath{clip}%
\pgfsetbuttcap%
\pgfsetroundjoin%
\definecolor{currentfill}{rgb}{0.762399,0.110727,0.151326}%
\pgfsetfillcolor{currentfill}%
\pgfsetlinewidth{0.000000pt}%
\definecolor{currentstroke}{rgb}{0.000000,0.000000,0.000000}%
\pgfsetstrokecolor{currentstroke}%
\pgfsetdash{}{0pt}%
\pgfpathmoveto{\pgfqpoint{1.451169in}{1.453939in}}%
\pgfpathlineto{\pgfqpoint{1.455181in}{1.456895in}}%
\pgfpathlineto{\pgfqpoint{1.459986in}{1.470143in}}%
\pgfpathlineto{\pgfqpoint{1.476452in}{1.474394in}}%
\pgfpathlineto{\pgfqpoint{1.485667in}{1.483391in}}%
\pgfpathlineto{\pgfqpoint{1.478792in}{1.496639in}}%
\pgfpathlineto{\pgfqpoint{1.495058in}{1.509887in}}%
\pgfpathlineto{\pgfqpoint{1.485509in}{1.523135in}}%
\pgfpathlineto{\pgfqpoint{1.497918in}{1.536383in}}%
\pgfpathlineto{\pgfqpoint{1.493856in}{1.549631in}}%
\pgfpathlineto{\pgfqpoint{1.496491in}{1.562879in}}%
\pgfpathlineto{\pgfqpoint{1.500961in}{1.576127in}}%
\pgfpathlineto{\pgfqpoint{1.493947in}{1.589375in}}%
\pgfpathlineto{\pgfqpoint{1.501735in}{1.599492in}}%
\pgfpathlineto{\pgfqpoint{1.503138in}{1.602623in}}%
\pgfpathlineto{\pgfqpoint{1.501735in}{1.605468in}}%
\pgfpathlineto{\pgfqpoint{1.493751in}{1.615871in}}%
\pgfpathlineto{\pgfqpoint{1.501735in}{1.628511in}}%
\pgfpathlineto{\pgfqpoint{1.501974in}{1.629119in}}%
\pgfpathlineto{\pgfqpoint{1.501735in}{1.629813in}}%
\pgfpathlineto{\pgfqpoint{1.495320in}{1.642367in}}%
\pgfpathlineto{\pgfqpoint{1.497474in}{1.655615in}}%
\pgfpathlineto{\pgfqpoint{1.498105in}{1.668863in}}%
\pgfpathlineto{\pgfqpoint{1.492514in}{1.682111in}}%
\pgfpathlineto{\pgfqpoint{1.498523in}{1.695359in}}%
\pgfpathlineto{\pgfqpoint{1.487963in}{1.708607in}}%
\pgfpathlineto{\pgfqpoint{1.493713in}{1.721855in}}%
\pgfpathlineto{\pgfqpoint{1.485366in}{1.735103in}}%
\pgfpathlineto{\pgfqpoint{1.485449in}{1.748351in}}%
\pgfpathlineto{\pgfqpoint{1.483447in}{1.761599in}}%
\pgfpathlineto{\pgfqpoint{1.477187in}{1.774847in}}%
\pgfpathlineto{\pgfqpoint{1.481370in}{1.788095in}}%
\pgfpathlineto{\pgfqpoint{1.476452in}{1.795230in}}%
\pgfpathlineto{\pgfqpoint{1.463482in}{1.801343in}}%
\pgfpathlineto{\pgfqpoint{1.476452in}{1.813745in}}%
\pgfpathlineto{\pgfqpoint{1.476737in}{1.814591in}}%
\pgfpathlineto{\pgfqpoint{1.476452in}{1.815037in}}%
\pgfpathlineto{\pgfqpoint{1.453497in}{1.827839in}}%
\pgfpathlineto{\pgfqpoint{1.451169in}{1.833975in}}%
\pgfpathlineto{\pgfqpoint{1.448953in}{1.827839in}}%
\pgfpathlineto{\pgfqpoint{1.432219in}{1.814591in}}%
\pgfpathlineto{\pgfqpoint{1.440625in}{1.801343in}}%
\pgfpathlineto{\pgfqpoint{1.425886in}{1.793157in}}%
\pgfpathlineto{\pgfqpoint{1.422602in}{1.788095in}}%
\pgfpathlineto{\pgfqpoint{1.425886in}{1.779280in}}%
\pgfpathlineto{\pgfqpoint{1.430393in}{1.774847in}}%
\pgfpathlineto{\pgfqpoint{1.425886in}{1.772123in}}%
\pgfpathlineto{\pgfqpoint{1.419038in}{1.761599in}}%
\pgfpathlineto{\pgfqpoint{1.420859in}{1.748351in}}%
\pgfpathlineto{\pgfqpoint{1.416938in}{1.735103in}}%
\pgfpathlineto{\pgfqpoint{1.413280in}{1.721855in}}%
\pgfpathlineto{\pgfqpoint{1.415881in}{1.708607in}}%
\pgfpathlineto{\pgfqpoint{1.407168in}{1.695359in}}%
\pgfpathlineto{\pgfqpoint{1.413757in}{1.682111in}}%
\pgfpathlineto{\pgfqpoint{1.404548in}{1.668863in}}%
\pgfpathlineto{\pgfqpoint{1.410055in}{1.655615in}}%
\pgfpathlineto{\pgfqpoint{1.405389in}{1.642367in}}%
\pgfpathlineto{\pgfqpoint{1.404847in}{1.629119in}}%
\pgfpathlineto{\pgfqpoint{1.408289in}{1.615871in}}%
\pgfpathlineto{\pgfqpoint{1.401064in}{1.602623in}}%
\pgfpathlineto{\pgfqpoint{1.411463in}{1.589375in}}%
\pgfpathlineto{\pgfqpoint{1.401811in}{1.576127in}}%
\pgfpathlineto{\pgfqpoint{1.411755in}{1.562879in}}%
\pgfpathlineto{\pgfqpoint{1.407314in}{1.549631in}}%
\pgfpathlineto{\pgfqpoint{1.409991in}{1.536383in}}%
\pgfpathlineto{\pgfqpoint{1.415793in}{1.523135in}}%
\pgfpathlineto{\pgfqpoint{1.409288in}{1.509887in}}%
\pgfpathlineto{\pgfqpoint{1.424349in}{1.496639in}}%
\pgfpathlineto{\pgfqpoint{1.415435in}{1.483391in}}%
\pgfpathlineto{\pgfqpoint{1.425886in}{1.475501in}}%
\pgfpathlineto{\pgfqpoint{1.444765in}{1.470143in}}%
\pgfpathlineto{\pgfqpoint{1.446154in}{1.456895in}}%
\pgfpathlineto{\pgfqpoint{1.451169in}{1.453939in}}%
\pgfpathclose%
\pgfpathmoveto{\pgfqpoint{1.447845in}{1.536383in}}%
\pgfpathlineto{\pgfqpoint{1.446777in}{1.549631in}}%
\pgfpathlineto{\pgfqpoint{1.448805in}{1.562879in}}%
\pgfpathlineto{\pgfqpoint{1.425886in}{1.574989in}}%
\pgfpathlineto{\pgfqpoint{1.425043in}{1.576127in}}%
\pgfpathlineto{\pgfqpoint{1.425886in}{1.577346in}}%
\pgfpathlineto{\pgfqpoint{1.447832in}{1.589375in}}%
\pgfpathlineto{\pgfqpoint{1.425886in}{1.598476in}}%
\pgfpathlineto{\pgfqpoint{1.422402in}{1.602623in}}%
\pgfpathlineto{\pgfqpoint{1.425886in}{1.608580in}}%
\pgfpathlineto{\pgfqpoint{1.438834in}{1.615871in}}%
\pgfpathlineto{\pgfqpoint{1.425886in}{1.626249in}}%
\pgfpathlineto{\pgfqpoint{1.424784in}{1.629119in}}%
\pgfpathlineto{\pgfqpoint{1.425886in}{1.639177in}}%
\pgfpathlineto{\pgfqpoint{1.426997in}{1.642367in}}%
\pgfpathlineto{\pgfqpoint{1.434002in}{1.655615in}}%
\pgfpathlineto{\pgfqpoint{1.425886in}{1.663406in}}%
\pgfpathlineto{\pgfqpoint{1.423685in}{1.668863in}}%
\pgfpathlineto{\pgfqpoint{1.425886in}{1.672227in}}%
\pgfpathlineto{\pgfqpoint{1.442289in}{1.682111in}}%
\pgfpathlineto{\pgfqpoint{1.425886in}{1.693203in}}%
\pgfpathlineto{\pgfqpoint{1.424731in}{1.695359in}}%
\pgfpathlineto{\pgfqpoint{1.425886in}{1.697072in}}%
\pgfpathlineto{\pgfqpoint{1.448063in}{1.708607in}}%
\pgfpathlineto{\pgfqpoint{1.436261in}{1.721855in}}%
\pgfpathlineto{\pgfqpoint{1.450497in}{1.735103in}}%
\pgfpathlineto{\pgfqpoint{1.451169in}{1.736945in}}%
\pgfpathlineto{\pgfqpoint{1.451838in}{1.735103in}}%
\pgfpathlineto{\pgfqpoint{1.472354in}{1.721855in}}%
\pgfpathlineto{\pgfqpoint{1.454684in}{1.708607in}}%
\pgfpathlineto{\pgfqpoint{1.476452in}{1.699450in}}%
\pgfpathlineto{\pgfqpoint{1.479683in}{1.695359in}}%
\pgfpathlineto{\pgfqpoint{1.476452in}{1.689278in}}%
\pgfpathlineto{\pgfqpoint{1.463184in}{1.682111in}}%
\pgfpathlineto{\pgfqpoint{1.476452in}{1.673879in}}%
\pgfpathlineto{\pgfqpoint{1.478825in}{1.668863in}}%
\pgfpathlineto{\pgfqpoint{1.476452in}{1.656481in}}%
\pgfpathlineto{\pgfqpoint{1.475841in}{1.655615in}}%
\pgfpathlineto{\pgfqpoint{1.472772in}{1.642367in}}%
\pgfpathlineto{\pgfqpoint{1.476452in}{1.638779in}}%
\pgfpathlineto{\pgfqpoint{1.480272in}{1.629119in}}%
\pgfpathlineto{\pgfqpoint{1.476452in}{1.623083in}}%
\pgfpathlineto{\pgfqpoint{1.463225in}{1.615871in}}%
\pgfpathlineto{\pgfqpoint{1.476452in}{1.609503in}}%
\pgfpathlineto{\pgfqpoint{1.481519in}{1.602623in}}%
\pgfpathlineto{\pgfqpoint{1.476452in}{1.596657in}}%
\pgfpathlineto{\pgfqpoint{1.455346in}{1.589375in}}%
\pgfpathlineto{\pgfqpoint{1.476452in}{1.577999in}}%
\pgfpathlineto{\pgfqpoint{1.477514in}{1.576127in}}%
\pgfpathlineto{\pgfqpoint{1.476452in}{1.574009in}}%
\pgfpathlineto{\pgfqpoint{1.454792in}{1.562879in}}%
\pgfpathlineto{\pgfqpoint{1.455208in}{1.549631in}}%
\pgfpathlineto{\pgfqpoint{1.456064in}{1.536383in}}%
\pgfpathlineto{\pgfqpoint{1.451169in}{1.534574in}}%
\pgfpathlineto{\pgfqpoint{1.447845in}{1.536383in}}%
\pgfpathclose%
\pgfusepath{fill}%
\end{pgfscope}%
\begin{pgfscope}%
\pgfpathrectangle{\pgfqpoint{0.540975in}{0.476543in}}{\pgfqpoint{1.820388in}{1.656000in}}%
\pgfusepath{clip}%
\pgfsetbuttcap%
\pgfsetroundjoin%
\definecolor{currentfill}{rgb}{0.685506,0.036909,0.149789}%
\pgfsetfillcolor{currentfill}%
\pgfsetlinewidth{0.000000pt}%
\definecolor{currentstroke}{rgb}{0.000000,0.000000,0.000000}%
\pgfsetstrokecolor{currentstroke}%
\pgfsetdash{}{0pt}%
\pgfpathmoveto{\pgfqpoint{1.451169in}{1.534574in}}%
\pgfpathlineto{\pgfqpoint{1.456064in}{1.536383in}}%
\pgfpathlineto{\pgfqpoint{1.455208in}{1.549631in}}%
\pgfpathlineto{\pgfqpoint{1.454792in}{1.562879in}}%
\pgfpathlineto{\pgfqpoint{1.476452in}{1.574009in}}%
\pgfpathlineto{\pgfqpoint{1.477514in}{1.576127in}}%
\pgfpathlineto{\pgfqpoint{1.476452in}{1.577999in}}%
\pgfpathlineto{\pgfqpoint{1.455346in}{1.589375in}}%
\pgfpathlineto{\pgfqpoint{1.476452in}{1.596657in}}%
\pgfpathlineto{\pgfqpoint{1.481519in}{1.602623in}}%
\pgfpathlineto{\pgfqpoint{1.476452in}{1.609503in}}%
\pgfpathlineto{\pgfqpoint{1.463225in}{1.615871in}}%
\pgfpathlineto{\pgfqpoint{1.476452in}{1.623083in}}%
\pgfpathlineto{\pgfqpoint{1.480272in}{1.629119in}}%
\pgfpathlineto{\pgfqpoint{1.476452in}{1.638779in}}%
\pgfpathlineto{\pgfqpoint{1.472772in}{1.642367in}}%
\pgfpathlineto{\pgfqpoint{1.475841in}{1.655615in}}%
\pgfpathlineto{\pgfqpoint{1.476452in}{1.656481in}}%
\pgfpathlineto{\pgfqpoint{1.478825in}{1.668863in}}%
\pgfpathlineto{\pgfqpoint{1.476452in}{1.673879in}}%
\pgfpathlineto{\pgfqpoint{1.463184in}{1.682111in}}%
\pgfpathlineto{\pgfqpoint{1.476452in}{1.689278in}}%
\pgfpathlineto{\pgfqpoint{1.479683in}{1.695359in}}%
\pgfpathlineto{\pgfqpoint{1.476452in}{1.699450in}}%
\pgfpathlineto{\pgfqpoint{1.454684in}{1.708607in}}%
\pgfpathlineto{\pgfqpoint{1.472354in}{1.721855in}}%
\pgfpathlineto{\pgfqpoint{1.451838in}{1.735103in}}%
\pgfpathlineto{\pgfqpoint{1.451169in}{1.736945in}}%
\pgfpathlineto{\pgfqpoint{1.450497in}{1.735103in}}%
\pgfpathlineto{\pgfqpoint{1.436261in}{1.721855in}}%
\pgfpathlineto{\pgfqpoint{1.448063in}{1.708607in}}%
\pgfpathlineto{\pgfqpoint{1.425886in}{1.697072in}}%
\pgfpathlineto{\pgfqpoint{1.424731in}{1.695359in}}%
\pgfpathlineto{\pgfqpoint{1.425886in}{1.693203in}}%
\pgfpathlineto{\pgfqpoint{1.442289in}{1.682111in}}%
\pgfpathlineto{\pgfqpoint{1.425886in}{1.672227in}}%
\pgfpathlineto{\pgfqpoint{1.423685in}{1.668863in}}%
\pgfpathlineto{\pgfqpoint{1.425886in}{1.663406in}}%
\pgfpathlineto{\pgfqpoint{1.434002in}{1.655615in}}%
\pgfpathlineto{\pgfqpoint{1.426997in}{1.642367in}}%
\pgfpathlineto{\pgfqpoint{1.425886in}{1.639177in}}%
\pgfpathlineto{\pgfqpoint{1.424784in}{1.629119in}}%
\pgfpathlineto{\pgfqpoint{1.425886in}{1.626249in}}%
\pgfpathlineto{\pgfqpoint{1.438834in}{1.615871in}}%
\pgfpathlineto{\pgfqpoint{1.425886in}{1.608580in}}%
\pgfpathlineto{\pgfqpoint{1.422402in}{1.602623in}}%
\pgfpathlineto{\pgfqpoint{1.425886in}{1.598476in}}%
\pgfpathlineto{\pgfqpoint{1.447832in}{1.589375in}}%
\pgfpathlineto{\pgfqpoint{1.425886in}{1.577346in}}%
\pgfpathlineto{\pgfqpoint{1.425043in}{1.576127in}}%
\pgfpathlineto{\pgfqpoint{1.425886in}{1.574989in}}%
\pgfpathlineto{\pgfqpoint{1.448805in}{1.562879in}}%
\pgfpathlineto{\pgfqpoint{1.446777in}{1.549631in}}%
\pgfpathlineto{\pgfqpoint{1.447845in}{1.536383in}}%
\pgfpathlineto{\pgfqpoint{1.451169in}{1.534574in}}%
\pgfpathclose%
\pgfusepath{fill}%
\end{pgfscope}%
\begin{pgfscope}%
\pgfsetbuttcap%
\pgfsetroundjoin%
\definecolor{currentfill}{rgb}{0.000000,0.000000,0.000000}%
\pgfsetfillcolor{currentfill}%
\pgfsetlinewidth{0.803000pt}%
\definecolor{currentstroke}{rgb}{0.000000,0.000000,0.000000}%
\pgfsetstrokecolor{currentstroke}%
\pgfsetdash{}{0pt}%
\pgfsys@defobject{currentmarker}{\pgfqpoint{0.000000in}{-0.048611in}}{\pgfqpoint{0.000000in}{0.000000in}}{%
\pgfpathmoveto{\pgfqpoint{0.000000in}{0.000000in}}%
\pgfpathlineto{\pgfqpoint{0.000000in}{-0.048611in}}%
\pgfusepath{stroke,fill}%
}%
\begin{pgfscope}%
\pgfsys@transformshift{0.642107in}{0.476543in}%
\pgfsys@useobject{currentmarker}{}%
\end{pgfscope}%
\end{pgfscope}%
\begin{pgfscope}%
\definecolor{textcolor}{rgb}{0.000000,0.000000,0.000000}%
\pgfsetstrokecolor{textcolor}%
\pgfsetfillcolor{textcolor}%
\pgftext[x=0.642107in,y=0.379321in,,top]{\color{textcolor}{\rmfamily\fontsize{8.000000}{9.600000}\selectfont\catcode`\^=\active\def^{\ifmmode\sp\else\^{}\fi}\catcode`\%=\active\def%{\%}\ensuremath{-}0.10}}%
\end{pgfscope}%
\begin{pgfscope}%
\pgfsetbuttcap%
\pgfsetroundjoin%
\definecolor{currentfill}{rgb}{0.000000,0.000000,0.000000}%
\pgfsetfillcolor{currentfill}%
\pgfsetlinewidth{0.803000pt}%
\definecolor{currentstroke}{rgb}{0.000000,0.000000,0.000000}%
\pgfsetstrokecolor{currentstroke}%
\pgfsetdash{}{0pt}%
\pgfsys@defobject{currentmarker}{\pgfqpoint{0.000000in}{-0.048611in}}{\pgfqpoint{0.000000in}{0.000000in}}{%
\pgfpathmoveto{\pgfqpoint{0.000000in}{0.000000in}}%
\pgfpathlineto{\pgfqpoint{0.000000in}{-0.048611in}}%
\pgfusepath{stroke,fill}%
}%
\begin{pgfscope}%
\pgfsys@transformshift{1.147771in}{0.476543in}%
\pgfsys@useobject{currentmarker}{}%
\end{pgfscope}%
\end{pgfscope}%
\begin{pgfscope}%
\definecolor{textcolor}{rgb}{0.000000,0.000000,0.000000}%
\pgfsetstrokecolor{textcolor}%
\pgfsetfillcolor{textcolor}%
\pgftext[x=1.147771in,y=0.379321in,,top]{\color{textcolor}{\rmfamily\fontsize{8.000000}{9.600000}\selectfont\catcode`\^=\active\def^{\ifmmode\sp\else\^{}\fi}\catcode`\%=\active\def%{\%}\ensuremath{-}0.05}}%
\end{pgfscope}%
\begin{pgfscope}%
\pgfsetbuttcap%
\pgfsetroundjoin%
\definecolor{currentfill}{rgb}{0.000000,0.000000,0.000000}%
\pgfsetfillcolor{currentfill}%
\pgfsetlinewidth{0.803000pt}%
\definecolor{currentstroke}{rgb}{0.000000,0.000000,0.000000}%
\pgfsetstrokecolor{currentstroke}%
\pgfsetdash{}{0pt}%
\pgfsys@defobject{currentmarker}{\pgfqpoint{0.000000in}{-0.048611in}}{\pgfqpoint{0.000000in}{0.000000in}}{%
\pgfpathmoveto{\pgfqpoint{0.000000in}{0.000000in}}%
\pgfpathlineto{\pgfqpoint{0.000000in}{-0.048611in}}%
\pgfusepath{stroke,fill}%
}%
\begin{pgfscope}%
\pgfsys@transformshift{1.653434in}{0.476543in}%
\pgfsys@useobject{currentmarker}{}%
\end{pgfscope}%
\end{pgfscope}%
\begin{pgfscope}%
\definecolor{textcolor}{rgb}{0.000000,0.000000,0.000000}%
\pgfsetstrokecolor{textcolor}%
\pgfsetfillcolor{textcolor}%
\pgftext[x=1.653434in,y=0.379321in,,top]{\color{textcolor}{\rmfamily\fontsize{8.000000}{9.600000}\selectfont\catcode`\^=\active\def^{\ifmmode\sp\else\^{}\fi}\catcode`\%=\active\def%{\%}0.00}}%
\end{pgfscope}%
\begin{pgfscope}%
\pgfsetbuttcap%
\pgfsetroundjoin%
\definecolor{currentfill}{rgb}{0.000000,0.000000,0.000000}%
\pgfsetfillcolor{currentfill}%
\pgfsetlinewidth{0.803000pt}%
\definecolor{currentstroke}{rgb}{0.000000,0.000000,0.000000}%
\pgfsetstrokecolor{currentstroke}%
\pgfsetdash{}{0pt}%
\pgfsys@defobject{currentmarker}{\pgfqpoint{0.000000in}{-0.048611in}}{\pgfqpoint{0.000000in}{0.000000in}}{%
\pgfpathmoveto{\pgfqpoint{0.000000in}{0.000000in}}%
\pgfpathlineto{\pgfqpoint{0.000000in}{-0.048611in}}%
\pgfusepath{stroke,fill}%
}%
\begin{pgfscope}%
\pgfsys@transformshift{2.159098in}{0.476543in}%
\pgfsys@useobject{currentmarker}{}%
\end{pgfscope}%
\end{pgfscope}%
\begin{pgfscope}%
\definecolor{textcolor}{rgb}{0.000000,0.000000,0.000000}%
\pgfsetstrokecolor{textcolor}%
\pgfsetfillcolor{textcolor}%
\pgftext[x=2.159098in,y=0.379321in,,top]{\color{textcolor}{\rmfamily\fontsize{8.000000}{9.600000}\selectfont\catcode`\^=\active\def^{\ifmmode\sp\else\^{}\fi}\catcode`\%=\active\def%{\%}0.05}}%
\end{pgfscope}%
\begin{pgfscope}%
\definecolor{textcolor}{rgb}{0.000000,0.000000,0.000000}%
\pgfsetstrokecolor{textcolor}%
\pgfsetfillcolor{textcolor}%
\pgftext[x=1.451169in,y=0.225000in,,top]{\color{textcolor}{\rmfamily\fontsize{9.000000}{10.800000}\selectfont\catcode`\^=\active\def^{\ifmmode\sp\else\^{}\fi}\catcode`\%=\active\def%{\%}$\delta_2/\omega_b$}}%
\end{pgfscope}%
\begin{pgfscope}%
\pgfsetbuttcap%
\pgfsetroundjoin%
\definecolor{currentfill}{rgb}{0.000000,0.000000,0.000000}%
\pgfsetfillcolor{currentfill}%
\pgfsetlinewidth{0.803000pt}%
\definecolor{currentstroke}{rgb}{0.000000,0.000000,0.000000}%
\pgfsetstrokecolor{currentstroke}%
\pgfsetdash{}{0pt}%
\pgfsys@defobject{currentmarker}{\pgfqpoint{-0.048611in}{0.000000in}}{\pgfqpoint{-0.000000in}{0.000000in}}{%
\pgfpathmoveto{\pgfqpoint{-0.000000in}{0.000000in}}%
\pgfpathlineto{\pgfqpoint{-0.048611in}{0.000000in}}%
\pgfusepath{stroke,fill}%
}%
\begin{pgfscope}%
\pgfsys@transformshift{0.540975in}{0.476543in}%
\pgfsys@useobject{currentmarker}{}%
\end{pgfscope}%
\end{pgfscope}%
\begin{pgfscope}%
\definecolor{textcolor}{rgb}{0.000000,0.000000,0.000000}%
\pgfsetstrokecolor{textcolor}%
\pgfsetfillcolor{textcolor}%
\pgftext[x=0.384724in, y=0.437963in, left, base]{\color{textcolor}{\rmfamily\fontsize{8.000000}{9.600000}\selectfont\catcode`\^=\active\def^{\ifmmode\sp\else\^{}\fi}\catcode`\%=\active\def%{\%}0}}%
\end{pgfscope}%
\begin{pgfscope}%
\pgfsetbuttcap%
\pgfsetroundjoin%
\definecolor{currentfill}{rgb}{0.000000,0.000000,0.000000}%
\pgfsetfillcolor{currentfill}%
\pgfsetlinewidth{0.803000pt}%
\definecolor{currentstroke}{rgb}{0.000000,0.000000,0.000000}%
\pgfsetstrokecolor{currentstroke}%
\pgfsetdash{}{0pt}%
\pgfsys@defobject{currentmarker}{\pgfqpoint{-0.048611in}{0.000000in}}{\pgfqpoint{-0.000000in}{0.000000in}}{%
\pgfpathmoveto{\pgfqpoint{-0.000000in}{0.000000in}}%
\pgfpathlineto{\pgfqpoint{-0.048611in}{0.000000in}}%
\pgfusepath{stroke,fill}%
}%
\begin{pgfscope}%
\pgfsys@transformshift{0.540975in}{0.807743in}%
\pgfsys@useobject{currentmarker}{}%
\end{pgfscope}%
\end{pgfscope}%
\begin{pgfscope}%
\definecolor{textcolor}{rgb}{0.000000,0.000000,0.000000}%
\pgfsetstrokecolor{textcolor}%
\pgfsetfillcolor{textcolor}%
\pgftext[x=0.325695in, y=0.769163in, left, base]{\color{textcolor}{\rmfamily\fontsize{8.000000}{9.600000}\selectfont\catcode`\^=\active\def^{\ifmmode\sp\else\^{}\fi}\catcode`\%=\active\def%{\%}50}}%
\end{pgfscope}%
\begin{pgfscope}%
\pgfsetbuttcap%
\pgfsetroundjoin%
\definecolor{currentfill}{rgb}{0.000000,0.000000,0.000000}%
\pgfsetfillcolor{currentfill}%
\pgfsetlinewidth{0.803000pt}%
\definecolor{currentstroke}{rgb}{0.000000,0.000000,0.000000}%
\pgfsetstrokecolor{currentstroke}%
\pgfsetdash{}{0pt}%
\pgfsys@defobject{currentmarker}{\pgfqpoint{-0.048611in}{0.000000in}}{\pgfqpoint{-0.000000in}{0.000000in}}{%
\pgfpathmoveto{\pgfqpoint{-0.000000in}{0.000000in}}%
\pgfpathlineto{\pgfqpoint{-0.048611in}{0.000000in}}%
\pgfusepath{stroke,fill}%
}%
\begin{pgfscope}%
\pgfsys@transformshift{0.540975in}{1.138943in}%
\pgfsys@useobject{currentmarker}{}%
\end{pgfscope}%
\end{pgfscope}%
\begin{pgfscope}%
\definecolor{textcolor}{rgb}{0.000000,0.000000,0.000000}%
\pgfsetstrokecolor{textcolor}%
\pgfsetfillcolor{textcolor}%
\pgftext[x=0.266667in, y=1.100363in, left, base]{\color{textcolor}{\rmfamily\fontsize{8.000000}{9.600000}\selectfont\catcode`\^=\active\def^{\ifmmode\sp\else\^{}\fi}\catcode`\%=\active\def%{\%}100}}%
\end{pgfscope}%
\begin{pgfscope}%
\pgfsetbuttcap%
\pgfsetroundjoin%
\definecolor{currentfill}{rgb}{0.000000,0.000000,0.000000}%
\pgfsetfillcolor{currentfill}%
\pgfsetlinewidth{0.803000pt}%
\definecolor{currentstroke}{rgb}{0.000000,0.000000,0.000000}%
\pgfsetstrokecolor{currentstroke}%
\pgfsetdash{}{0pt}%
\pgfsys@defobject{currentmarker}{\pgfqpoint{-0.048611in}{0.000000in}}{\pgfqpoint{-0.000000in}{0.000000in}}{%
\pgfpathmoveto{\pgfqpoint{-0.000000in}{0.000000in}}%
\pgfpathlineto{\pgfqpoint{-0.048611in}{0.000000in}}%
\pgfusepath{stroke,fill}%
}%
\begin{pgfscope}%
\pgfsys@transformshift{0.540975in}{1.470143in}%
\pgfsys@useobject{currentmarker}{}%
\end{pgfscope}%
\end{pgfscope}%
\begin{pgfscope}%
\definecolor{textcolor}{rgb}{0.000000,0.000000,0.000000}%
\pgfsetstrokecolor{textcolor}%
\pgfsetfillcolor{textcolor}%
\pgftext[x=0.266667in, y=1.431563in, left, base]{\color{textcolor}{\rmfamily\fontsize{8.000000}{9.600000}\selectfont\catcode`\^=\active\def^{\ifmmode\sp\else\^{}\fi}\catcode`\%=\active\def%{\%}150}}%
\end{pgfscope}%
\begin{pgfscope}%
\pgfsetbuttcap%
\pgfsetroundjoin%
\definecolor{currentfill}{rgb}{0.000000,0.000000,0.000000}%
\pgfsetfillcolor{currentfill}%
\pgfsetlinewidth{0.803000pt}%
\definecolor{currentstroke}{rgb}{0.000000,0.000000,0.000000}%
\pgfsetstrokecolor{currentstroke}%
\pgfsetdash{}{0pt}%
\pgfsys@defobject{currentmarker}{\pgfqpoint{-0.048611in}{0.000000in}}{\pgfqpoint{-0.000000in}{0.000000in}}{%
\pgfpathmoveto{\pgfqpoint{-0.000000in}{0.000000in}}%
\pgfpathlineto{\pgfqpoint{-0.048611in}{0.000000in}}%
\pgfusepath{stroke,fill}%
}%
\begin{pgfscope}%
\pgfsys@transformshift{0.540975in}{1.801343in}%
\pgfsys@useobject{currentmarker}{}%
\end{pgfscope}%
\end{pgfscope}%
\begin{pgfscope}%
\definecolor{textcolor}{rgb}{0.000000,0.000000,0.000000}%
\pgfsetstrokecolor{textcolor}%
\pgfsetfillcolor{textcolor}%
\pgftext[x=0.266667in, y=1.762763in, left, base]{\color{textcolor}{\rmfamily\fontsize{8.000000}{9.600000}\selectfont\catcode`\^=\active\def^{\ifmmode\sp\else\^{}\fi}\catcode`\%=\active\def%{\%}200}}%
\end{pgfscope}%
\begin{pgfscope}%
\pgfsetbuttcap%
\pgfsetroundjoin%
\definecolor{currentfill}{rgb}{0.000000,0.000000,0.000000}%
\pgfsetfillcolor{currentfill}%
\pgfsetlinewidth{0.803000pt}%
\definecolor{currentstroke}{rgb}{0.000000,0.000000,0.000000}%
\pgfsetstrokecolor{currentstroke}%
\pgfsetdash{}{0pt}%
\pgfsys@defobject{currentmarker}{\pgfqpoint{-0.048611in}{0.000000in}}{\pgfqpoint{-0.000000in}{0.000000in}}{%
\pgfpathmoveto{\pgfqpoint{-0.000000in}{0.000000in}}%
\pgfpathlineto{\pgfqpoint{-0.048611in}{0.000000in}}%
\pgfusepath{stroke,fill}%
}%
\begin{pgfscope}%
\pgfsys@transformshift{0.540975in}{2.132543in}%
\pgfsys@useobject{currentmarker}{}%
\end{pgfscope}%
\end{pgfscope}%
\begin{pgfscope}%
\definecolor{textcolor}{rgb}{0.000000,0.000000,0.000000}%
\pgfsetstrokecolor{textcolor}%
\pgfsetfillcolor{textcolor}%
\pgftext[x=0.266667in, y=2.093963in, left, base]{\color{textcolor}{\rmfamily\fontsize{8.000000}{9.600000}\selectfont\catcode`\^=\active\def^{\ifmmode\sp\else\^{}\fi}\catcode`\%=\active\def%{\%}250}}%
\end{pgfscope}%
\begin{pgfscope}%
\definecolor{textcolor}{rgb}{0.000000,0.000000,0.000000}%
\pgfsetstrokecolor{textcolor}%
\pgfsetfillcolor{textcolor}%
\pgftext[x=0.211111in,y=1.304543in,,bottom,rotate=90.000000]{\color{textcolor}{\rmfamily\fontsize{9.000000}{10.800000}\selectfont\catcode`\^=\active\def^{\ifmmode\sp\else\^{}\fi}\catcode`\%=\active\def%{\%}$\omega_b\,t$}}%
\end{pgfscope}%
\begin{pgfscope}%
\pgfpathrectangle{\pgfqpoint{0.540975in}{0.476543in}}{\pgfqpoint{1.820388in}{1.656000in}}%
\pgfusepath{clip}%
\pgfsetbuttcap%
\pgfsetroundjoin%
\pgfsetlinewidth{0.501875pt}%
\definecolor{currentstroke}{rgb}{1.000000,1.000000,1.000000}%
\pgfsetstrokecolor{currentstroke}%
\pgfsetstrokeopacity{0.800000}%
\pgfsetdash{}{0pt}%
\pgfpathmoveto{\pgfqpoint{1.341138in}{2.132543in}}%
\pgfpathlineto{\pgfqpoint{1.339792in}{2.119295in}}%
\pgfpathlineto{\pgfqpoint{1.336187in}{2.106047in}}%
\pgfpathlineto{\pgfqpoint{1.324753in}{2.043587in}}%
\pgfpathlineto{\pgfqpoint{1.323439in}{2.039807in}}%
\pgfpathlineto{\pgfqpoint{1.322493in}{2.026559in}}%
\pgfpathlineto{\pgfqpoint{1.318268in}{2.013311in}}%
\pgfpathlineto{\pgfqpoint{1.313348in}{1.986815in}}%
\pgfpathlineto{\pgfqpoint{1.310032in}{1.973567in}}%
\pgfpathlineto{\pgfqpoint{1.308880in}{1.960319in}}%
\pgfpathlineto{\pgfqpoint{1.304684in}{1.947071in}}%
\pgfpathlineto{\pgfqpoint{1.304039in}{1.933823in}}%
\pgfpathlineto{\pgfqpoint{1.300249in}{1.920575in}}%
\pgfpathlineto{\pgfqpoint{1.298893in}{1.907327in}}%
\pgfpathlineto{\pgfqpoint{1.292285in}{1.880831in}}%
\pgfpathlineto{\pgfqpoint{1.290534in}{1.867583in}}%
\pgfpathlineto{\pgfqpoint{1.286233in}{1.854335in}}%
\pgfpathlineto{\pgfqpoint{1.285615in}{1.841087in}}%
\pgfpathlineto{\pgfqpoint{1.281119in}{1.827839in}}%
\pgfpathlineto{\pgfqpoint{1.280044in}{1.814591in}}%
\pgfpathlineto{\pgfqpoint{1.274283in}{1.788095in}}%
\pgfpathlineto{\pgfqpoint{1.272039in}{1.774847in}}%
\pgfpathlineto{\pgfqpoint{1.267364in}{1.761599in}}%
\pgfpathlineto{\pgfqpoint{1.266556in}{1.748351in}}%
\pgfpathlineto{\pgfqpoint{1.261574in}{1.735103in}}%
\pgfpathlineto{\pgfqpoint{1.260966in}{1.721855in}}%
\pgfpathlineto{\pgfqpoint{1.256960in}{1.708607in}}%
\pgfpathlineto{\pgfqpoint{1.252379in}{1.682111in}}%
\pgfpathlineto{\pgfqpoint{1.248414in}{1.668863in}}%
\pgfpathlineto{\pgfqpoint{1.247286in}{1.655615in}}%
\pgfpathlineto{\pgfqpoint{1.241347in}{1.642367in}}%
\pgfpathlineto{\pgfqpoint{1.240723in}{1.629119in}}%
\pgfpathlineto{\pgfqpoint{1.235714in}{1.615871in}}%
\pgfpathlineto{\pgfqpoint{1.234147in}{1.602623in}}%
\pgfpathlineto{\pgfqpoint{1.223620in}{1.555633in}}%
\pgfpathlineto{\pgfqpoint{1.220799in}{1.549631in}}%
\pgfpathlineto{\pgfqpoint{1.220433in}{1.536383in}}%
\pgfpathlineto{\pgfqpoint{1.214267in}{1.523135in}}%
\pgfpathlineto{\pgfqpoint{1.212674in}{1.509887in}}%
\pgfpathlineto{\pgfqpoint{1.205134in}{1.483391in}}%
\pgfpathlineto{\pgfqpoint{1.204137in}{1.470143in}}%
\pgfpathlineto{\pgfqpoint{1.198283in}{1.456895in}}%
\pgfpathlineto{\pgfqpoint{1.198280in}{1.443647in}}%
\pgfpathlineto{\pgfqpoint{1.191312in}{1.430399in}}%
\pgfpathlineto{\pgfqpoint{1.190534in}{1.417151in}}%
\pgfpathlineto{\pgfqpoint{1.182125in}{1.390655in}}%
\pgfpathlineto{\pgfqpoint{1.180584in}{1.377407in}}%
\pgfpathlineto{\pgfqpoint{1.174339in}{1.364159in}}%
\pgfpathlineto{\pgfqpoint{1.174939in}{1.350911in}}%
\pgfpathlineto{\pgfqpoint{1.167222in}{1.337663in}}%
\pgfpathlineto{\pgfqpoint{1.166682in}{1.324415in}}%
\pgfpathlineto{\pgfqpoint{1.160836in}{1.311167in}}%
\pgfpathlineto{\pgfqpoint{1.157732in}{1.297919in}}%
\pgfpathlineto{\pgfqpoint{1.156195in}{1.284671in}}%
\pgfpathlineto{\pgfqpoint{1.149393in}{1.271423in}}%
\pgfpathlineto{\pgfqpoint{1.150706in}{1.258175in}}%
\pgfpathlineto{\pgfqpoint{1.142217in}{1.244927in}}%
\pgfpathlineto{\pgfqpoint{1.143545in}{1.231679in}}%
\pgfpathlineto{\pgfqpoint{1.137141in}{1.218431in}}%
\pgfpathlineto{\pgfqpoint{1.131694in}{1.191935in}}%
\pgfpathlineto{\pgfqpoint{1.124237in}{1.178687in}}%
\pgfpathlineto{\pgfqpoint{1.126885in}{1.165439in}}%
\pgfpathlineto{\pgfqpoint{1.117231in}{1.152191in}}%
\pgfpathlineto{\pgfqpoint{1.120488in}{1.138943in}}%
\pgfpathlineto{\pgfqpoint{1.113449in}{1.125695in}}%
\pgfpathlineto{\pgfqpoint{1.112782in}{1.099199in}}%
\pgfpathlineto{\pgfqpoint{1.104634in}{1.085951in}}%
\pgfpathlineto{\pgfqpoint{1.109671in}{1.072703in}}%
\pgfpathlineto{\pgfqpoint{1.098026in}{1.059455in}}%
\pgfpathlineto{\pgfqpoint{1.105141in}{1.046207in}}%
\pgfpathlineto{\pgfqpoint{1.098472in}{1.032959in}}%
\pgfpathlineto{\pgfqpoint{1.099872in}{1.019711in}}%
\pgfpathlineto{\pgfqpoint{1.103513in}{1.006463in}}%
\pgfpathlineto{\pgfqpoint{1.097204in}{0.992976in}}%
\pgfpathlineto{\pgfqpoint{1.112600in}{0.979967in}}%
\pgfpathlineto{\pgfqpoint{1.101677in}{0.966719in}}%
\pgfpathlineto{\pgfqpoint{1.117968in}{0.953471in}}%
\pgfpathlineto{\pgfqpoint{1.115204in}{0.940223in}}%
\pgfpathlineto{\pgfqpoint{1.124004in}{0.926975in}}%
\pgfpathlineto{\pgfqpoint{1.145354in}{0.913727in}}%
\pgfpathlineto{\pgfqpoint{1.141748in}{0.900479in}}%
\pgfpathlineto{\pgfqpoint{1.147771in}{0.898376in}}%
\pgfpathlineto{\pgfqpoint{1.186791in}{0.887231in}}%
\pgfpathlineto{\pgfqpoint{1.186913in}{0.873983in}}%
\pgfpathlineto{\pgfqpoint{1.223620in}{0.865945in}}%
\pgfpathlineto{\pgfqpoint{1.324753in}{0.848114in}}%
\pgfpathlineto{\pgfqpoint{1.332143in}{0.847487in}}%
\pgfpathlineto{\pgfqpoint{1.375319in}{0.841071in}}%
\pgfpathlineto{\pgfqpoint{1.425886in}{0.837388in}}%
\pgfpathlineto{\pgfqpoint{1.451169in}{0.837298in}}%
\pgfpathlineto{\pgfqpoint{1.476452in}{0.838669in}}%
\pgfpathlineto{\pgfqpoint{1.501735in}{0.841888in}}%
\pgfpathlineto{\pgfqpoint{1.527018in}{0.847506in}}%
\pgfpathlineto{\pgfqpoint{1.577585in}{0.851525in}}%
\pgfpathlineto{\pgfqpoint{1.628151in}{0.857225in}}%
\pgfpathlineto{\pgfqpoint{1.653434in}{0.860973in}}%
\pgfpathlineto{\pgfqpoint{1.678717in}{0.874313in}}%
\pgfpathlineto{\pgfqpoint{1.716359in}{0.887231in}}%
\pgfpathlineto{\pgfqpoint{1.729284in}{0.893917in}}%
\pgfpathlineto{\pgfqpoint{1.739432in}{0.900479in}}%
\pgfpathlineto{\pgfqpoint{1.745863in}{0.913727in}}%
\pgfpathlineto{\pgfqpoint{1.754567in}{0.917606in}}%
\pgfpathlineto{\pgfqpoint{1.771994in}{0.926975in}}%
\pgfpathlineto{\pgfqpoint{1.766227in}{0.940223in}}%
\pgfpathlineto{\pgfqpoint{1.788153in}{0.953471in}}%
\pgfpathlineto{\pgfqpoint{1.779850in}{0.966856in}}%
\pgfpathlineto{\pgfqpoint{1.793803in}{0.979967in}}%
\pgfpathlineto{\pgfqpoint{1.791303in}{0.993215in}}%
\pgfpathlineto{\pgfqpoint{1.793817in}{1.006463in}}%
\pgfpathlineto{\pgfqpoint{1.797661in}{1.019711in}}%
\pgfpathlineto{\pgfqpoint{1.790729in}{1.032959in}}%
\pgfpathlineto{\pgfqpoint{1.799017in}{1.046207in}}%
\pgfpathlineto{\pgfqpoint{1.790199in}{1.059455in}}%
\pgfpathlineto{\pgfqpoint{1.796900in}{1.072703in}}%
\pgfpathlineto{\pgfqpoint{1.789758in}{1.085951in}}%
\pgfpathlineto{\pgfqpoint{1.790989in}{1.099199in}}%
\pgfpathlineto{\pgfqpoint{1.788645in}{1.112447in}}%
\pgfpathlineto{\pgfqpoint{1.783438in}{1.125695in}}%
\pgfpathlineto{\pgfqpoint{1.784898in}{1.138943in}}%
\pgfpathlineto{\pgfqpoint{1.775457in}{1.152191in}}%
\pgfpathlineto{\pgfqpoint{1.778542in}{1.165439in}}%
\pgfpathlineto{\pgfqpoint{1.770760in}{1.178687in}}%
\pgfpathlineto{\pgfqpoint{1.771506in}{1.191935in}}%
\pgfpathlineto{\pgfqpoint{1.763433in}{1.218431in}}%
\pgfpathlineto{\pgfqpoint{1.763147in}{1.231679in}}%
\pgfpathlineto{\pgfqpoint{1.755843in}{1.244927in}}%
\pgfpathlineto{\pgfqpoint{1.756805in}{1.258175in}}%
\pgfpathlineto{\pgfqpoint{1.748761in}{1.271423in}}%
\pgfpathlineto{\pgfqpoint{1.748850in}{1.284671in}}%
\pgfpathlineto{\pgfqpoint{1.743457in}{1.297919in}}%
\pgfpathlineto{\pgfqpoint{1.738618in}{1.324415in}}%
\pgfpathlineto{\pgfqpoint{1.732059in}{1.337663in}}%
\pgfpathlineto{\pgfqpoint{1.732519in}{1.350911in}}%
\pgfpathlineto{\pgfqpoint{1.725648in}{1.364159in}}%
\pgfpathlineto{\pgfqpoint{1.725449in}{1.377407in}}%
\pgfpathlineto{\pgfqpoint{1.720157in}{1.390655in}}%
\pgfpathlineto{\pgfqpoint{1.715153in}{1.417151in}}%
\pgfpathlineto{\pgfqpoint{1.710087in}{1.430399in}}%
\pgfpathlineto{\pgfqpoint{1.709566in}{1.443647in}}%
\pgfpathlineto{\pgfqpoint{1.702640in}{1.456895in}}%
\pgfpathlineto{\pgfqpoint{1.702310in}{1.470143in}}%
\pgfpathlineto{\pgfqpoint{1.697194in}{1.483391in}}%
\pgfpathlineto{\pgfqpoint{1.692256in}{1.509887in}}%
\pgfpathlineto{\pgfqpoint{1.687825in}{1.523135in}}%
\pgfpathlineto{\pgfqpoint{1.687165in}{1.536383in}}%
\pgfpathlineto{\pgfqpoint{1.681506in}{1.549631in}}%
\pgfpathlineto{\pgfqpoint{1.681108in}{1.562879in}}%
\pgfpathlineto{\pgfqpoint{1.675885in}{1.576127in}}%
\pgfpathlineto{\pgfqpoint{1.674058in}{1.589375in}}%
\pgfpathlineto{\pgfqpoint{1.667124in}{1.615871in}}%
\pgfpathlineto{\pgfqpoint{1.665846in}{1.629119in}}%
\pgfpathlineto{\pgfqpoint{1.660503in}{1.642367in}}%
\pgfpathlineto{\pgfqpoint{1.660146in}{1.655615in}}%
\pgfpathlineto{\pgfqpoint{1.655173in}{1.668863in}}%
\pgfpathlineto{\pgfqpoint{1.653434in}{1.682919in}}%
\pgfpathlineto{\pgfqpoint{1.647416in}{1.708607in}}%
\pgfpathlineto{\pgfqpoint{1.646010in}{1.721855in}}%
\pgfpathlineto{\pgfqpoint{1.641506in}{1.735103in}}%
\pgfpathlineto{\pgfqpoint{1.640743in}{1.748351in}}%
\pgfpathlineto{\pgfqpoint{1.635922in}{1.761599in}}%
\pgfpathlineto{\pgfqpoint{1.634526in}{1.774847in}}%
\pgfpathlineto{\pgfqpoint{1.627902in}{1.801343in}}%
\pgfpathlineto{\pgfqpoint{1.626487in}{1.814591in}}%
\pgfpathlineto{\pgfqpoint{1.622374in}{1.827839in}}%
\pgfpathlineto{\pgfqpoint{1.621707in}{1.841087in}}%
\pgfpathlineto{\pgfqpoint{1.617580in}{1.854335in}}%
\pgfpathlineto{\pgfqpoint{1.616340in}{1.867583in}}%
\pgfpathlineto{\pgfqpoint{1.610191in}{1.894079in}}%
\pgfpathlineto{\pgfqpoint{1.608376in}{1.907327in}}%
\pgfpathlineto{\pgfqpoint{1.603955in}{1.920575in}}%
\pgfpathlineto{\pgfqpoint{1.602868in}{1.934000in}}%
\pgfpathlineto{\pgfqpoint{1.598841in}{1.947071in}}%
\pgfpathlineto{\pgfqpoint{1.597665in}{1.960319in}}%
\pgfpathlineto{\pgfqpoint{1.592248in}{1.986815in}}%
\pgfpathlineto{\pgfqpoint{1.590589in}{2.000063in}}%
\pgfpathlineto{\pgfqpoint{1.586653in}{2.013311in}}%
\pgfpathlineto{\pgfqpoint{1.585816in}{2.026559in}}%
\pgfpathlineto{\pgfqpoint{1.581331in}{2.039807in}}%
\pgfpathlineto{\pgfqpoint{1.579913in}{2.053055in}}%
\pgfpathlineto{\pgfqpoint{1.576247in}{2.066303in}}%
\pgfpathlineto{\pgfqpoint{1.568293in}{2.106047in}}%
\pgfpathlineto{\pgfqpoint{1.567452in}{2.119295in}}%
\pgfpathlineto{\pgfqpoint{1.563280in}{2.132543in}}%
\pgfpathlineto{\pgfqpoint{1.563280in}{2.132543in}}%
\pgfusepath{stroke}%
\end{pgfscope}%
\begin{pgfscope}%
\pgfpathrectangle{\pgfqpoint{0.540975in}{0.476543in}}{\pgfqpoint{1.820388in}{1.656000in}}%
\pgfusepath{clip}%
\pgfsetbuttcap%
\pgfsetroundjoin%
\pgfsetlinewidth{0.501875pt}%
\definecolor{currentstroke}{rgb}{1.000000,1.000000,1.000000}%
\pgfsetstrokecolor{currentstroke}%
\pgfsetstrokeopacity{0.800000}%
\pgfsetdash{}{0pt}%
\pgfpathmoveto{\pgfqpoint{1.385464in}{2.132543in}}%
\pgfpathlineto{\pgfqpoint{1.383347in}{2.119295in}}%
\pgfpathlineto{\pgfqpoint{1.375319in}{2.087270in}}%
\pgfpathlineto{\pgfqpoint{1.374589in}{2.079551in}}%
\pgfpathlineto{\pgfqpoint{1.369720in}{2.066303in}}%
\pgfpathlineto{\pgfqpoint{1.369356in}{2.053055in}}%
\pgfpathlineto{\pgfqpoint{1.364732in}{2.039807in}}%
\pgfpathlineto{\pgfqpoint{1.363721in}{2.026559in}}%
\pgfpathlineto{\pgfqpoint{1.357366in}{2.000063in}}%
\pgfpathlineto{\pgfqpoint{1.355434in}{1.986815in}}%
\pgfpathlineto{\pgfqpoint{1.351096in}{1.973567in}}%
\pgfpathlineto{\pgfqpoint{1.350036in}{1.958657in}}%
\pgfpathlineto{\pgfqpoint{1.345815in}{1.947071in}}%
\pgfpathlineto{\pgfqpoint{1.345394in}{1.933823in}}%
\pgfpathlineto{\pgfqpoint{1.341736in}{1.920575in}}%
\pgfpathlineto{\pgfqpoint{1.338023in}{1.894079in}}%
\pgfpathlineto{\pgfqpoint{1.334326in}{1.880831in}}%
\pgfpathlineto{\pgfqpoint{1.333640in}{1.867583in}}%
\pgfpathlineto{\pgfqpoint{1.328985in}{1.854335in}}%
\pgfpathlineto{\pgfqpoint{1.328839in}{1.841087in}}%
\pgfpathlineto{\pgfqpoint{1.324589in}{1.827839in}}%
\pgfpathlineto{\pgfqpoint{1.323318in}{1.814591in}}%
\pgfpathlineto{\pgfqpoint{1.317696in}{1.788095in}}%
\pgfpathlineto{\pgfqpoint{1.317028in}{1.774847in}}%
\pgfpathlineto{\pgfqpoint{1.312423in}{1.761599in}}%
\pgfpathlineto{\pgfqpoint{1.312534in}{1.748351in}}%
\pgfpathlineto{\pgfqpoint{1.308046in}{1.735103in}}%
\pgfpathlineto{\pgfqpoint{1.307549in}{1.721855in}}%
\pgfpathlineto{\pgfqpoint{1.302112in}{1.695359in}}%
\pgfpathlineto{\pgfqpoint{1.301176in}{1.682111in}}%
\pgfpathlineto{\pgfqpoint{1.296631in}{1.668863in}}%
\pgfpathlineto{\pgfqpoint{1.296916in}{1.655615in}}%
\pgfpathlineto{\pgfqpoint{1.291757in}{1.642367in}}%
\pgfpathlineto{\pgfqpoint{1.291682in}{1.629119in}}%
\pgfpathlineto{\pgfqpoint{1.288133in}{1.615871in}}%
\pgfpathlineto{\pgfqpoint{1.285560in}{1.589375in}}%
\pgfpathlineto{\pgfqpoint{1.281387in}{1.576127in}}%
\pgfpathlineto{\pgfqpoint{1.282170in}{1.562879in}}%
\pgfpathlineto{\pgfqpoint{1.276925in}{1.549631in}}%
\pgfpathlineto{\pgfqpoint{1.277497in}{1.536383in}}%
\pgfpathlineto{\pgfqpoint{1.273177in}{1.523135in}}%
\pgfpathlineto{\pgfqpoint{1.270185in}{1.496639in}}%
\pgfpathlineto{\pgfqpoint{1.266215in}{1.483391in}}%
\pgfpathlineto{\pgfqpoint{1.267750in}{1.470143in}}%
\pgfpathlineto{\pgfqpoint{1.262251in}{1.456895in}}%
\pgfpathlineto{\pgfqpoint{1.264028in}{1.443647in}}%
\pgfpathlineto{\pgfqpoint{1.259518in}{1.430399in}}%
\pgfpathlineto{\pgfqpoint{1.257739in}{1.403903in}}%
\pgfpathlineto{\pgfqpoint{1.253782in}{1.390655in}}%
\pgfpathlineto{\pgfqpoint{1.255702in}{1.377407in}}%
\pgfpathlineto{\pgfqpoint{1.249818in}{1.364159in}}%
\pgfpathlineto{\pgfqpoint{1.253139in}{1.350911in}}%
\pgfpathlineto{\pgfqpoint{1.248234in}{1.337663in}}%
\pgfpathlineto{\pgfqpoint{1.248904in}{1.322901in}}%
\pgfpathlineto{\pgfqpoint{1.247994in}{1.311167in}}%
\pgfpathlineto{\pgfqpoint{1.244932in}{1.297919in}}%
\pgfpathlineto{\pgfqpoint{1.248755in}{1.284671in}}%
\pgfpathlineto{\pgfqpoint{1.242582in}{1.271423in}}%
\pgfpathlineto{\pgfqpoint{1.248195in}{1.258175in}}%
\pgfpathlineto{\pgfqpoint{1.243393in}{1.244927in}}%
\pgfpathlineto{\pgfqpoint{1.246932in}{1.231679in}}%
\pgfpathlineto{\pgfqpoint{1.247520in}{1.218431in}}%
\pgfpathlineto{\pgfqpoint{1.244971in}{1.205183in}}%
\pgfpathlineto{\pgfqpoint{1.251880in}{1.191935in}}%
\pgfpathlineto{\pgfqpoint{1.245645in}{1.178687in}}%
\pgfpathlineto{\pgfqpoint{1.256667in}{1.165439in}}%
\pgfpathlineto{\pgfqpoint{1.252721in}{1.152191in}}%
\pgfpathlineto{\pgfqpoint{1.260970in}{1.138943in}}%
\pgfpathlineto{\pgfqpoint{1.266865in}{1.125695in}}%
\pgfpathlineto{\pgfqpoint{1.268265in}{1.112447in}}%
\pgfpathlineto{\pgfqpoint{1.274187in}{1.107717in}}%
\pgfpathlineto{\pgfqpoint{1.287216in}{1.099199in}}%
\pgfpathlineto{\pgfqpoint{1.281954in}{1.085951in}}%
\pgfpathlineto{\pgfqpoint{1.306846in}{1.072703in}}%
\pgfpathlineto{\pgfqpoint{1.309171in}{1.059455in}}%
\pgfpathlineto{\pgfqpoint{1.332052in}{1.046207in}}%
\pgfpathlineto{\pgfqpoint{1.350036in}{1.038878in}}%
\pgfpathlineto{\pgfqpoint{1.372528in}{1.032959in}}%
\pgfpathlineto{\pgfqpoint{1.375319in}{1.029940in}}%
\pgfpathlineto{\pgfqpoint{1.388377in}{1.019711in}}%
\pgfpathlineto{\pgfqpoint{1.400603in}{1.018203in}}%
\pgfpathlineto{\pgfqpoint{1.425886in}{1.016438in}}%
\pgfpathlineto{\pgfqpoint{1.451169in}{1.015962in}}%
\pgfpathlineto{\pgfqpoint{1.476452in}{1.016751in}}%
\pgfpathlineto{\pgfqpoint{1.507882in}{1.019711in}}%
\pgfpathlineto{\pgfqpoint{1.514201in}{1.032959in}}%
\pgfpathlineto{\pgfqpoint{1.527018in}{1.035029in}}%
\pgfpathlineto{\pgfqpoint{1.552302in}{1.040600in}}%
\pgfpathlineto{\pgfqpoint{1.571841in}{1.046207in}}%
\pgfpathlineto{\pgfqpoint{1.577585in}{1.059456in}}%
\pgfpathlineto{\pgfqpoint{1.602868in}{1.075463in}}%
\pgfpathlineto{\pgfqpoint{1.611180in}{1.085951in}}%
\pgfpathlineto{\pgfqpoint{1.616440in}{1.099199in}}%
\pgfpathlineto{\pgfqpoint{1.632181in}{1.112447in}}%
\pgfpathlineto{\pgfqpoint{1.629444in}{1.125695in}}%
\pgfpathlineto{\pgfqpoint{1.643622in}{1.138943in}}%
\pgfpathlineto{\pgfqpoint{1.638402in}{1.152191in}}%
\pgfpathlineto{\pgfqpoint{1.648709in}{1.165439in}}%
\pgfpathlineto{\pgfqpoint{1.647674in}{1.178687in}}%
\pgfpathlineto{\pgfqpoint{1.655422in}{1.205183in}}%
\pgfpathlineto{\pgfqpoint{1.652515in}{1.218431in}}%
\pgfpathlineto{\pgfqpoint{1.659402in}{1.231679in}}%
\pgfpathlineto{\pgfqpoint{1.653947in}{1.244927in}}%
\pgfpathlineto{\pgfqpoint{1.659129in}{1.258175in}}%
\pgfpathlineto{\pgfqpoint{1.654473in}{1.271423in}}%
\pgfpathlineto{\pgfqpoint{1.656142in}{1.284671in}}%
\pgfpathlineto{\pgfqpoint{1.655389in}{1.297919in}}%
\pgfpathlineto{\pgfqpoint{1.653434in}{1.313271in}}%
\pgfpathlineto{\pgfqpoint{1.655782in}{1.324415in}}%
\pgfpathlineto{\pgfqpoint{1.650407in}{1.337663in}}%
\pgfpathlineto{\pgfqpoint{1.654157in}{1.350911in}}%
\pgfpathlineto{\pgfqpoint{1.649213in}{1.364159in}}%
\pgfpathlineto{\pgfqpoint{1.650380in}{1.377407in}}%
\pgfpathlineto{\pgfqpoint{1.645561in}{1.403903in}}%
\pgfpathlineto{\pgfqpoint{1.646159in}{1.417151in}}%
\pgfpathlineto{\pgfqpoint{1.641422in}{1.430399in}}%
\pgfpathlineto{\pgfqpoint{1.643316in}{1.443647in}}%
\pgfpathlineto{\pgfqpoint{1.637948in}{1.456895in}}%
\pgfpathlineto{\pgfqpoint{1.638766in}{1.470143in}}%
\pgfpathlineto{\pgfqpoint{1.635463in}{1.483391in}}%
\pgfpathlineto{\pgfqpoint{1.633520in}{1.496639in}}%
\pgfpathlineto{\pgfqpoint{1.632948in}{1.509887in}}%
\pgfpathlineto{\pgfqpoint{1.628468in}{1.523135in}}%
\pgfpathlineto{\pgfqpoint{1.629768in}{1.536383in}}%
\pgfpathlineto{\pgfqpoint{1.624900in}{1.549631in}}%
\pgfpathlineto{\pgfqpoint{1.625446in}{1.562879in}}%
\pgfpathlineto{\pgfqpoint{1.621723in}{1.576127in}}%
\pgfpathlineto{\pgfqpoint{1.618375in}{1.602623in}}%
\pgfpathlineto{\pgfqpoint{1.614337in}{1.615871in}}%
\pgfpathlineto{\pgfqpoint{1.614582in}{1.629119in}}%
\pgfpathlineto{\pgfqpoint{1.609660in}{1.642367in}}%
\pgfpathlineto{\pgfqpoint{1.610143in}{1.655615in}}%
\pgfpathlineto{\pgfqpoint{1.606427in}{1.668863in}}%
\pgfpathlineto{\pgfqpoint{1.602868in}{1.696866in}}%
\pgfpathlineto{\pgfqpoint{1.599514in}{1.708607in}}%
\pgfpathlineto{\pgfqpoint{1.599320in}{1.721855in}}%
\pgfpathlineto{\pgfqpoint{1.594655in}{1.735103in}}%
\pgfpathlineto{\pgfqpoint{1.594509in}{1.748351in}}%
\pgfpathlineto{\pgfqpoint{1.590464in}{1.761599in}}%
\pgfpathlineto{\pgfqpoint{1.587045in}{1.788095in}}%
\pgfpathlineto{\pgfqpoint{1.583792in}{1.801343in}}%
\pgfpathlineto{\pgfqpoint{1.583412in}{1.814591in}}%
\pgfpathlineto{\pgfqpoint{1.578709in}{1.827839in}}%
\pgfpathlineto{\pgfqpoint{1.577585in}{1.844210in}}%
\pgfpathlineto{\pgfqpoint{1.574511in}{1.854335in}}%
\pgfpathlineto{\pgfqpoint{1.570657in}{1.880831in}}%
\pgfpathlineto{\pgfqpoint{1.567290in}{1.894079in}}%
\pgfpathlineto{\pgfqpoint{1.566588in}{1.907327in}}%
\pgfpathlineto{\pgfqpoint{1.562127in}{1.920575in}}%
\pgfpathlineto{\pgfqpoint{1.561786in}{1.933823in}}%
\pgfpathlineto{\pgfqpoint{1.557520in}{1.947071in}}%
\pgfpathlineto{\pgfqpoint{1.553397in}{1.973567in}}%
\pgfpathlineto{\pgfqpoint{1.550127in}{1.986815in}}%
\pgfpathlineto{\pgfqpoint{1.549188in}{2.000063in}}%
\pgfpathlineto{\pgfqpoint{1.544444in}{2.013311in}}%
\pgfpathlineto{\pgfqpoint{1.544051in}{2.026559in}}%
\pgfpathlineto{\pgfqpoint{1.539521in}{2.039807in}}%
\pgfpathlineto{\pgfqpoint{1.537823in}{2.053055in}}%
\pgfpathlineto{\pgfqpoint{1.524644in}{2.106047in}}%
\pgfpathlineto{\pgfqpoint{1.524021in}{2.119295in}}%
\pgfpathlineto{\pgfqpoint{1.518783in}{2.132543in}}%
\pgfpathlineto{\pgfqpoint{1.518783in}{2.132543in}}%
\pgfusepath{stroke}%
\end{pgfscope}%
\begin{pgfscope}%
\pgfpathrectangle{\pgfqpoint{0.540975in}{0.476543in}}{\pgfqpoint{1.820388in}{1.656000in}}%
\pgfusepath{clip}%
\pgfsetbuttcap%
\pgfsetroundjoin%
\pgfsetlinewidth{0.501875pt}%
\definecolor{currentstroke}{rgb}{1.000000,1.000000,1.000000}%
\pgfsetstrokecolor{currentstroke}%
\pgfsetstrokeopacity{0.800000}%
\pgfsetdash{}{0pt}%
\pgfpathmoveto{\pgfqpoint{1.400603in}{1.201655in}}%
\pgfpathlineto{\pgfqpoint{1.425886in}{1.196660in}}%
\pgfpathlineto{\pgfqpoint{1.451169in}{1.195035in}}%
\pgfpathlineto{\pgfqpoint{1.476452in}{1.196722in}}%
\pgfpathlineto{\pgfqpoint{1.501735in}{1.202135in}}%
\pgfpathlineto{\pgfqpoint{1.509806in}{1.205183in}}%
\pgfpathlineto{\pgfqpoint{1.506959in}{1.218431in}}%
\pgfpathlineto{\pgfqpoint{1.527018in}{1.226031in}}%
\pgfpathlineto{\pgfqpoint{1.537396in}{1.231679in}}%
\pgfpathlineto{\pgfqpoint{1.535649in}{1.244927in}}%
\pgfpathlineto{\pgfqpoint{1.550304in}{1.258175in}}%
\pgfpathlineto{\pgfqpoint{1.552302in}{1.266975in}}%
\pgfpathlineto{\pgfqpoint{1.553048in}{1.271423in}}%
\pgfpathlineto{\pgfqpoint{1.556217in}{1.284671in}}%
\pgfpathlineto{\pgfqpoint{1.564459in}{1.297919in}}%
\pgfpathlineto{\pgfqpoint{1.561998in}{1.311167in}}%
\pgfpathlineto{\pgfqpoint{1.572810in}{1.324415in}}%
\pgfpathlineto{\pgfqpoint{1.568487in}{1.337663in}}%
\pgfpathlineto{\pgfqpoint{1.577201in}{1.350911in}}%
\pgfpathlineto{\pgfqpoint{1.575069in}{1.364159in}}%
\pgfpathlineto{\pgfqpoint{1.577585in}{1.376710in}}%
\pgfpathlineto{\pgfqpoint{1.577708in}{1.377407in}}%
\pgfpathlineto{\pgfqpoint{1.579342in}{1.390655in}}%
\pgfpathlineto{\pgfqpoint{1.577585in}{1.399439in}}%
\pgfpathlineto{\pgfqpoint{1.576583in}{1.403903in}}%
\pgfpathlineto{\pgfqpoint{1.577585in}{1.406310in}}%
\pgfpathlineto{\pgfqpoint{1.581373in}{1.417151in}}%
\pgfpathlineto{\pgfqpoint{1.577585in}{1.427997in}}%
\pgfpathlineto{\pgfqpoint{1.576665in}{1.430399in}}%
\pgfpathlineto{\pgfqpoint{1.577585in}{1.432695in}}%
\pgfpathlineto{\pgfqpoint{1.581388in}{1.443647in}}%
\pgfpathlineto{\pgfqpoint{1.577667in}{1.456895in}}%
\pgfpathlineto{\pgfqpoint{1.579484in}{1.470143in}}%
\pgfpathlineto{\pgfqpoint{1.578878in}{1.483391in}}%
\pgfpathlineto{\pgfqpoint{1.577585in}{1.491835in}}%
\pgfpathlineto{\pgfqpoint{1.576777in}{1.496639in}}%
\pgfpathlineto{\pgfqpoint{1.577585in}{1.500920in}}%
\pgfpathlineto{\pgfqpoint{1.579048in}{1.509887in}}%
\pgfpathlineto{\pgfqpoint{1.577585in}{1.514137in}}%
\pgfpathlineto{\pgfqpoint{1.574228in}{1.523135in}}%
\pgfpathlineto{\pgfqpoint{1.577585in}{1.535600in}}%
\pgfpathlineto{\pgfqpoint{1.577772in}{1.536383in}}%
\pgfpathlineto{\pgfqpoint{1.577585in}{1.536963in}}%
\pgfpathlineto{\pgfqpoint{1.573257in}{1.549631in}}%
\pgfpathlineto{\pgfqpoint{1.574640in}{1.562879in}}%
\pgfpathlineto{\pgfqpoint{1.572480in}{1.576127in}}%
\pgfpathlineto{\pgfqpoint{1.570179in}{1.589375in}}%
\pgfpathlineto{\pgfqpoint{1.570891in}{1.602623in}}%
\pgfpathlineto{\pgfqpoint{1.566253in}{1.615871in}}%
\pgfpathlineto{\pgfqpoint{1.568185in}{1.629119in}}%
\pgfpathlineto{\pgfqpoint{1.563447in}{1.642367in}}%
\pgfpathlineto{\pgfqpoint{1.564584in}{1.655615in}}%
\pgfpathlineto{\pgfqpoint{1.562049in}{1.668863in}}%
\pgfpathlineto{\pgfqpoint{1.560462in}{1.682111in}}%
\pgfpathlineto{\pgfqpoint{1.560394in}{1.695359in}}%
\pgfpathlineto{\pgfqpoint{1.555920in}{1.708607in}}%
\pgfpathlineto{\pgfqpoint{1.557088in}{1.721855in}}%
\pgfpathlineto{\pgfqpoint{1.552302in}{1.734773in}}%
\pgfpathlineto{\pgfqpoint{1.552166in}{1.735103in}}%
\pgfpathlineto{\pgfqpoint{1.552302in}{1.740055in}}%
\pgfpathlineto{\pgfqpoint{1.552504in}{1.748351in}}%
\pgfpathlineto{\pgfqpoint{1.552302in}{1.749137in}}%
\pgfpathlineto{\pgfqpoint{1.548834in}{1.761599in}}%
\pgfpathlineto{\pgfqpoint{1.546874in}{1.774847in}}%
\pgfpathlineto{\pgfqpoint{1.546205in}{1.788095in}}%
\pgfpathlineto{\pgfqpoint{1.542018in}{1.801343in}}%
\pgfpathlineto{\pgfqpoint{1.542957in}{1.814591in}}%
\pgfpathlineto{\pgfqpoint{1.537735in}{1.827839in}}%
\pgfpathlineto{\pgfqpoint{1.537983in}{1.841087in}}%
\pgfpathlineto{\pgfqpoint{1.534181in}{1.854335in}}%
\pgfpathlineto{\pgfqpoint{1.532086in}{1.867583in}}%
\pgfpathlineto{\pgfqpoint{1.530639in}{1.880831in}}%
\pgfpathlineto{\pgfqpoint{1.527018in}{1.891524in}}%
\pgfpathlineto{\pgfqpoint{1.525937in}{1.894079in}}%
\pgfpathlineto{\pgfqpoint{1.526493in}{1.907327in}}%
\pgfpathlineto{\pgfqpoint{1.520291in}{1.920575in}}%
\pgfpathlineto{\pgfqpoint{1.520346in}{1.933823in}}%
\pgfpathlineto{\pgfqpoint{1.515664in}{1.947071in}}%
\pgfpathlineto{\pgfqpoint{1.512980in}{1.960319in}}%
\pgfpathlineto{\pgfqpoint{1.511178in}{1.973567in}}%
\pgfpathlineto{\pgfqpoint{1.505958in}{1.986815in}}%
\pgfpathlineto{\pgfqpoint{1.506005in}{2.000063in}}%
\pgfpathlineto{\pgfqpoint{1.501735in}{2.009000in}}%
\pgfpathlineto{\pgfqpoint{1.498606in}{2.013311in}}%
\pgfpathlineto{\pgfqpoint{1.498165in}{2.026559in}}%
\pgfpathlineto{\pgfqpoint{1.490998in}{2.039807in}}%
\pgfpathlineto{\pgfqpoint{1.486367in}{2.053055in}}%
\pgfpathlineto{\pgfqpoint{1.483517in}{2.066303in}}%
\pgfpathlineto{\pgfqpoint{1.476452in}{2.076850in}}%
\pgfpathlineto{\pgfqpoint{1.470628in}{2.079551in}}%
\pgfpathlineto{\pgfqpoint{1.470474in}{2.092799in}}%
\pgfpathlineto{\pgfqpoint{1.451169in}{2.102120in}}%
\pgfpathlineto{\pgfqpoint{1.436184in}{2.092799in}}%
\pgfpathlineto{\pgfqpoint{1.436048in}{2.079551in}}%
\pgfpathlineto{\pgfqpoint{1.425886in}{2.073969in}}%
\pgfpathlineto{\pgfqpoint{1.420573in}{2.066303in}}%
\pgfpathlineto{\pgfqpoint{1.420248in}{2.053055in}}%
\pgfpathlineto{\pgfqpoint{1.413151in}{2.039807in}}%
\pgfpathlineto{\pgfqpoint{1.410280in}{2.026559in}}%
\pgfpathlineto{\pgfqpoint{1.406335in}{2.013311in}}%
\pgfpathlineto{\pgfqpoint{1.400603in}{2.002010in}}%
\pgfpathlineto{\pgfqpoint{1.399947in}{2.000063in}}%
\pgfpathlineto{\pgfqpoint{1.399134in}{1.986815in}}%
\pgfpathlineto{\pgfqpoint{1.393011in}{1.973567in}}%
\pgfpathlineto{\pgfqpoint{1.393359in}{1.960319in}}%
\pgfpathlineto{\pgfqpoint{1.387505in}{1.947071in}}%
\pgfpathlineto{\pgfqpoint{1.386888in}{1.933823in}}%
\pgfpathlineto{\pgfqpoint{1.383483in}{1.920575in}}%
\pgfpathlineto{\pgfqpoint{1.380444in}{1.907327in}}%
\pgfpathlineto{\pgfqpoint{1.379599in}{1.894079in}}%
\pgfpathlineto{\pgfqpoint{1.375319in}{1.884289in}}%
\pgfpathlineto{\pgfqpoint{1.374087in}{1.880831in}}%
\pgfpathlineto{\pgfqpoint{1.374393in}{1.867583in}}%
\pgfpathlineto{\pgfqpoint{1.369090in}{1.854335in}}%
\pgfpathlineto{\pgfqpoint{1.369207in}{1.841087in}}%
\pgfpathlineto{\pgfqpoint{1.365348in}{1.827839in}}%
\pgfpathlineto{\pgfqpoint{1.363564in}{1.814591in}}%
\pgfpathlineto{\pgfqpoint{1.362386in}{1.801343in}}%
\pgfpathlineto{\pgfqpoint{1.358401in}{1.788095in}}%
\pgfpathlineto{\pgfqpoint{1.359235in}{1.774847in}}%
\pgfpathlineto{\pgfqpoint{1.353858in}{1.761599in}}%
\pgfpathlineto{\pgfqpoint{1.354765in}{1.748351in}}%
\pgfpathlineto{\pgfqpoint{1.350319in}{1.735103in}}%
\pgfpathlineto{\pgfqpoint{1.350036in}{1.730113in}}%
\pgfpathlineto{\pgfqpoint{1.349618in}{1.721855in}}%
\pgfpathlineto{\pgfqpoint{1.348073in}{1.708607in}}%
\pgfpathlineto{\pgfqpoint{1.344831in}{1.695359in}}%
\pgfpathlineto{\pgfqpoint{1.345628in}{1.682111in}}%
\pgfpathlineto{\pgfqpoint{1.340740in}{1.668863in}}%
\pgfpathlineto{\pgfqpoint{1.342288in}{1.655615in}}%
\pgfpathlineto{\pgfqpoint{1.337721in}{1.642367in}}%
\pgfpathlineto{\pgfqpoint{1.337933in}{1.629119in}}%
\pgfpathlineto{\pgfqpoint{1.335995in}{1.615871in}}%
\pgfpathlineto{\pgfqpoint{1.333680in}{1.602623in}}%
\pgfpathlineto{\pgfqpoint{1.335137in}{1.589375in}}%
\pgfpathlineto{\pgfqpoint{1.330405in}{1.576127in}}%
\pgfpathlineto{\pgfqpoint{1.333169in}{1.562879in}}%
\pgfpathlineto{\pgfqpoint{1.328246in}{1.549631in}}%
\pgfpathlineto{\pgfqpoint{1.329637in}{1.536383in}}%
\pgfpathlineto{\pgfqpoint{1.327141in}{1.523135in}}%
\pgfpathlineto{\pgfqpoint{1.325237in}{1.509887in}}%
\pgfpathlineto{\pgfqpoint{1.326835in}{1.496639in}}%
\pgfpathlineto{\pgfqpoint{1.324753in}{1.490640in}}%
\pgfpathlineto{\pgfqpoint{1.322452in}{1.483391in}}%
\pgfpathlineto{\pgfqpoint{1.324753in}{1.476201in}}%
\pgfpathlineto{\pgfqpoint{1.326982in}{1.470143in}}%
\pgfpathlineto{\pgfqpoint{1.324753in}{1.464142in}}%
\pgfpathlineto{\pgfqpoint{1.322232in}{1.456895in}}%
\pgfpathlineto{\pgfqpoint{1.324753in}{1.447228in}}%
\pgfpathlineto{\pgfqpoint{1.325806in}{1.443647in}}%
\pgfpathlineto{\pgfqpoint{1.324753in}{1.436259in}}%
\pgfpathlineto{\pgfqpoint{1.323952in}{1.430399in}}%
\pgfpathlineto{\pgfqpoint{1.323666in}{1.417151in}}%
\pgfpathlineto{\pgfqpoint{1.324753in}{1.411651in}}%
\pgfpathlineto{\pgfqpoint{1.326586in}{1.403903in}}%
\pgfpathlineto{\pgfqpoint{1.324753in}{1.398919in}}%
\pgfpathlineto{\pgfqpoint{1.322030in}{1.390655in}}%
\pgfpathlineto{\pgfqpoint{1.324753in}{1.384579in}}%
\pgfpathlineto{\pgfqpoint{1.328558in}{1.377407in}}%
\pgfpathlineto{\pgfqpoint{1.324753in}{1.367893in}}%
\pgfpathlineto{\pgfqpoint{1.323377in}{1.364159in}}%
\pgfpathlineto{\pgfqpoint{1.324753in}{1.361211in}}%
\pgfpathlineto{\pgfqpoint{1.330316in}{1.350911in}}%
\pgfpathlineto{\pgfqpoint{1.329601in}{1.337663in}}%
\pgfpathlineto{\pgfqpoint{1.331795in}{1.324415in}}%
\pgfpathlineto{\pgfqpoint{1.338769in}{1.311167in}}%
\pgfpathlineto{\pgfqpoint{1.335358in}{1.297919in}}%
\pgfpathlineto{\pgfqpoint{1.348414in}{1.284671in}}%
\pgfpathlineto{\pgfqpoint{1.343782in}{1.271423in}}%
\pgfpathlineto{\pgfqpoint{1.350036in}{1.264574in}}%
\pgfpathlineto{\pgfqpoint{1.357318in}{1.258175in}}%
\pgfpathlineto{\pgfqpoint{1.361218in}{1.244927in}}%
\pgfpathlineto{\pgfqpoint{1.368675in}{1.231679in}}%
\pgfpathlineto{\pgfqpoint{1.375319in}{1.226804in}}%
\pgfpathlineto{\pgfqpoint{1.393009in}{1.218431in}}%
\pgfpathlineto{\pgfqpoint{1.390248in}{1.205183in}}%
\pgfpathlineto{\pgfqpoint{1.400603in}{1.201655in}}%
\pgfpathclose%
\pgfusepath{stroke}%
\end{pgfscope}%
\begin{pgfscope}%
\pgfpathrectangle{\pgfqpoint{0.540975in}{0.476543in}}{\pgfqpoint{1.820388in}{1.656000in}}%
\pgfusepath{clip}%
\pgfsetbuttcap%
\pgfsetroundjoin%
\pgfsetlinewidth{0.501875pt}%
\definecolor{currentstroke}{rgb}{1.000000,1.000000,1.000000}%
\pgfsetstrokecolor{currentstroke}%
\pgfsetstrokeopacity{0.800000}%
\pgfsetdash{}{0pt}%
\pgfpathmoveto{\pgfqpoint{1.451169in}{1.412153in}}%
\pgfpathlineto{\pgfqpoint{1.476452in}{1.416425in}}%
\pgfpathlineto{\pgfqpoint{1.477806in}{1.417151in}}%
\pgfpathlineto{\pgfqpoint{1.476452in}{1.419681in}}%
\pgfpathlineto{\pgfqpoint{1.461421in}{1.430399in}}%
\pgfpathlineto{\pgfqpoint{1.476452in}{1.434695in}}%
\pgfpathlineto{\pgfqpoint{1.487944in}{1.443647in}}%
\pgfpathlineto{\pgfqpoint{1.490246in}{1.456895in}}%
\pgfpathlineto{\pgfqpoint{1.494340in}{1.470143in}}%
\pgfpathlineto{\pgfqpoint{1.501735in}{1.478359in}}%
\pgfpathlineto{\pgfqpoint{1.504240in}{1.483391in}}%
\pgfpathlineto{\pgfqpoint{1.501735in}{1.491730in}}%
\pgfpathlineto{\pgfqpoint{1.499338in}{1.496639in}}%
\pgfpathlineto{\pgfqpoint{1.501735in}{1.498642in}}%
\pgfpathlineto{\pgfqpoint{1.509612in}{1.509887in}}%
\pgfpathlineto{\pgfqpoint{1.503227in}{1.523135in}}%
\pgfpathlineto{\pgfqpoint{1.511017in}{1.536383in}}%
\pgfpathlineto{\pgfqpoint{1.507658in}{1.549631in}}%
\pgfpathlineto{\pgfqpoint{1.509650in}{1.562879in}}%
\pgfpathlineto{\pgfqpoint{1.511487in}{1.576127in}}%
\pgfpathlineto{\pgfqpoint{1.507290in}{1.589375in}}%
\pgfpathlineto{\pgfqpoint{1.512876in}{1.602623in}}%
\pgfpathlineto{\pgfqpoint{1.506287in}{1.615871in}}%
\pgfpathlineto{\pgfqpoint{1.511375in}{1.629119in}}%
\pgfpathlineto{\pgfqpoint{1.506589in}{1.642367in}}%
\pgfpathlineto{\pgfqpoint{1.508193in}{1.655615in}}%
\pgfpathlineto{\pgfqpoint{1.507936in}{1.668863in}}%
\pgfpathlineto{\pgfqpoint{1.504686in}{1.682111in}}%
\pgfpathlineto{\pgfqpoint{1.507930in}{1.695359in}}%
\pgfpathlineto{\pgfqpoint{1.501735in}{1.707578in}}%
\pgfpathlineto{\pgfqpoint{1.500912in}{1.708607in}}%
\pgfpathlineto{\pgfqpoint{1.501735in}{1.710530in}}%
\pgfpathlineto{\pgfqpoint{1.504731in}{1.721855in}}%
\pgfpathlineto{\pgfqpoint{1.501735in}{1.728650in}}%
\pgfpathlineto{\pgfqpoint{1.497443in}{1.735103in}}%
\pgfpathlineto{\pgfqpoint{1.497959in}{1.748351in}}%
\pgfpathlineto{\pgfqpoint{1.494942in}{1.761599in}}%
\pgfpathlineto{\pgfqpoint{1.489191in}{1.774847in}}%
\pgfpathlineto{\pgfqpoint{1.492582in}{1.788095in}}%
\pgfpathlineto{\pgfqpoint{1.483508in}{1.801343in}}%
\pgfpathlineto{\pgfqpoint{1.487830in}{1.814591in}}%
\pgfpathlineto{\pgfqpoint{1.479256in}{1.827839in}}%
\pgfpathlineto{\pgfqpoint{1.478955in}{1.841087in}}%
\pgfpathlineto{\pgfqpoint{1.476452in}{1.849865in}}%
\pgfpathlineto{\pgfqpoint{1.472955in}{1.854335in}}%
\pgfpathlineto{\pgfqpoint{1.451169in}{1.867439in}}%
\pgfpathlineto{\pgfqpoint{1.430756in}{1.854335in}}%
\pgfpathlineto{\pgfqpoint{1.428451in}{1.841087in}}%
\pgfpathlineto{\pgfqpoint{1.425886in}{1.836890in}}%
\pgfpathlineto{\pgfqpoint{1.423558in}{1.827839in}}%
\pgfpathlineto{\pgfqpoint{1.418274in}{1.814591in}}%
\pgfpathlineto{\pgfqpoint{1.420489in}{1.801343in}}%
\pgfpathlineto{\pgfqpoint{1.411953in}{1.788095in}}%
\pgfpathlineto{\pgfqpoint{1.416601in}{1.774847in}}%
\pgfpathlineto{\pgfqpoint{1.407609in}{1.761599in}}%
\pgfpathlineto{\pgfqpoint{1.409690in}{1.748351in}}%
\pgfpathlineto{\pgfqpoint{1.404868in}{1.735103in}}%
\pgfpathlineto{\pgfqpoint{1.401784in}{1.721855in}}%
\pgfpathlineto{\pgfqpoint{1.403436in}{1.708607in}}%
\pgfpathlineto{\pgfqpoint{1.400603in}{1.704226in}}%
\pgfpathlineto{\pgfqpoint{1.396957in}{1.695359in}}%
\pgfpathlineto{\pgfqpoint{1.400603in}{1.683067in}}%
\pgfpathlineto{\pgfqpoint{1.401058in}{1.682111in}}%
\pgfpathlineto{\pgfqpoint{1.400603in}{1.681468in}}%
\pgfpathlineto{\pgfqpoint{1.394658in}{1.668863in}}%
\pgfpathlineto{\pgfqpoint{1.398203in}{1.655615in}}%
\pgfpathlineto{\pgfqpoint{1.394512in}{1.642367in}}%
\pgfpathlineto{\pgfqpoint{1.394402in}{1.629119in}}%
\pgfpathlineto{\pgfqpoint{1.395828in}{1.615871in}}%
\pgfpathlineto{\pgfqpoint{1.391437in}{1.602623in}}%
\pgfpathlineto{\pgfqpoint{1.397575in}{1.589375in}}%
\pgfpathlineto{\pgfqpoint{1.391195in}{1.576127in}}%
\pgfpathlineto{\pgfqpoint{1.397492in}{1.562879in}}%
\pgfpathlineto{\pgfqpoint{1.393758in}{1.549631in}}%
\pgfpathlineto{\pgfqpoint{1.395791in}{1.536383in}}%
\pgfpathlineto{\pgfqpoint{1.398166in}{1.523135in}}%
\pgfpathlineto{\pgfqpoint{1.394390in}{1.509887in}}%
\pgfpathlineto{\pgfqpoint{1.400603in}{1.500121in}}%
\pgfpathlineto{\pgfqpoint{1.404380in}{1.496639in}}%
\pgfpathlineto{\pgfqpoint{1.400603in}{1.491336in}}%
\pgfpathlineto{\pgfqpoint{1.396967in}{1.483391in}}%
\pgfpathlineto{\pgfqpoint{1.400603in}{1.478982in}}%
\pgfpathlineto{\pgfqpoint{1.412323in}{1.470143in}}%
\pgfpathlineto{\pgfqpoint{1.409123in}{1.456895in}}%
\pgfpathlineto{\pgfqpoint{1.418838in}{1.443647in}}%
\pgfpathlineto{\pgfqpoint{1.425886in}{1.434982in}}%
\pgfpathlineto{\pgfqpoint{1.439253in}{1.430399in}}%
\pgfpathlineto{\pgfqpoint{1.427225in}{1.417151in}}%
\pgfpathlineto{\pgfqpoint{1.451169in}{1.412153in}}%
\pgfpathclose%
\pgfpathmoveto{\pgfqpoint{1.451169in}{1.868204in}}%
\pgfpathlineto{\pgfqpoint{1.456201in}{1.880831in}}%
\pgfpathlineto{\pgfqpoint{1.451169in}{1.882976in}}%
\pgfpathlineto{\pgfqpoint{1.447015in}{1.880831in}}%
\pgfpathlineto{\pgfqpoint{1.451169in}{1.868204in}}%
\pgfpathclose%
\pgfusepath{stroke}%
\end{pgfscope}%
\begin{pgfscope}%
\pgfpathrectangle{\pgfqpoint{0.540975in}{0.476543in}}{\pgfqpoint{1.820388in}{1.656000in}}%
\pgfusepath{clip}%
\pgfsetbuttcap%
\pgfsetbeveljoin%
\definecolor{currentfill}{rgb}{1.000000,1.000000,1.000000}%
\pgfsetfillcolor{currentfill}%
\pgfsetlinewidth{0.501875pt}%
\definecolor{currentstroke}{rgb}{0.000000,0.000000,0.000000}%
\pgfsetstrokecolor{currentstroke}%
\pgfsetdash{}{0pt}%
\pgfsys@defobject{currentmarker}{\pgfqpoint{-0.052836in}{-0.044945in}}{\pgfqpoint{0.052836in}{0.055556in}}{%
\pgfpathmoveto{\pgfqpoint{0.000000in}{0.055556in}}%
\pgfpathlineto{\pgfqpoint{-0.012473in}{0.017168in}}%
\pgfpathlineto{\pgfqpoint{-0.052836in}{0.017168in}}%
\pgfpathlineto{\pgfqpoint{-0.020182in}{-0.006557in}}%
\pgfpathlineto{\pgfqpoint{-0.032655in}{-0.044945in}}%
\pgfpathlineto{\pgfqpoint{-0.000000in}{-0.021220in}}%
\pgfpathlineto{\pgfqpoint{0.032655in}{-0.044945in}}%
\pgfpathlineto{\pgfqpoint{0.020182in}{-0.006557in}}%
\pgfpathlineto{\pgfqpoint{0.052836in}{0.017168in}}%
\pgfpathlineto{\pgfqpoint{0.012473in}{0.017168in}}%
\pgfpathlineto{\pgfqpoint{0.000000in}{0.055556in}}%
\pgfpathclose%
\pgfusepath{stroke,fill}%
}%
\begin{pgfscope}%
\pgfsys@transformshift{1.451169in}{1.602623in}%
\pgfsys@useobject{currentmarker}{}%
\end{pgfscope}%
\end{pgfscope}%
\begin{pgfscope}%
\pgfsetrectcap%
\pgfsetmiterjoin%
\pgfsetlinewidth{0.803000pt}%
\definecolor{currentstroke}{rgb}{0.000000,0.000000,0.000000}%
\pgfsetstrokecolor{currentstroke}%
\pgfsetdash{}{0pt}%
\pgfpathmoveto{\pgfqpoint{0.540975in}{0.476543in}}%
\pgfpathlineto{\pgfqpoint{0.540975in}{2.132543in}}%
\pgfusepath{stroke}%
\end{pgfscope}%
\begin{pgfscope}%
\pgfsetrectcap%
\pgfsetmiterjoin%
\pgfsetlinewidth{0.803000pt}%
\definecolor{currentstroke}{rgb}{0.000000,0.000000,0.000000}%
\pgfsetstrokecolor{currentstroke}%
\pgfsetdash{}{0pt}%
\pgfpathmoveto{\pgfqpoint{2.361363in}{0.476543in}}%
\pgfpathlineto{\pgfqpoint{2.361363in}{2.132543in}}%
\pgfusepath{stroke}%
\end{pgfscope}%
\begin{pgfscope}%
\pgfsetrectcap%
\pgfsetmiterjoin%
\pgfsetlinewidth{0.803000pt}%
\definecolor{currentstroke}{rgb}{0.000000,0.000000,0.000000}%
\pgfsetstrokecolor{currentstroke}%
\pgfsetdash{}{0pt}%
\pgfpathmoveto{\pgfqpoint{0.540975in}{0.476543in}}%
\pgfpathlineto{\pgfqpoint{2.361363in}{0.476543in}}%
\pgfusepath{stroke}%
\end{pgfscope}%
\begin{pgfscope}%
\pgfsetrectcap%
\pgfsetmiterjoin%
\pgfsetlinewidth{0.803000pt}%
\definecolor{currentstroke}{rgb}{0.000000,0.000000,0.000000}%
\pgfsetstrokecolor{currentstroke}%
\pgfsetdash{}{0pt}%
\pgfpathmoveto{\pgfqpoint{0.540975in}{2.132543in}}%
\pgfpathlineto{\pgfqpoint{2.361363in}{2.132543in}}%
\pgfusepath{stroke}%
\end{pgfscope}%
\begin{pgfscope}%
\definecolor{textcolor}{rgb}{1.000000,1.000000,1.000000}%
\pgfsetstrokecolor{textcolor}%
\pgfsetfillcolor{textcolor}%
\pgftext[x=0.595586in,y=2.049743in,left,top]{\color{textcolor}{\rmfamily\fontsize{9.000000}{10.800000}\selectfont\catcode`\^=\active\def^{\ifmmode\sp\else\^{}\fi}\catcode`\%=\active\def%{\%}$(a)$}}%
\end{pgfscope}%
\begin{pgfscope}%
\definecolor{textcolor}{rgb}{0.200000,0.200000,0.200000}%
\pgfsetstrokecolor{textcolor}%
\pgfsetfillcolor{textcolor}%
\pgftext[x=1.451169in,y=2.182223in,,bottom]{\color{textcolor}{\rmfamily\fontsize{8.000000}{9.600000}\selectfont\catcode`\^=\active\def^{\ifmmode\sp\else\^{}\fi}\catcode`\%=\active\def%{\%}$D_1$}}%
\end{pgfscope}%
\begin{pgfscope}%
\pgfsetbuttcap%
\pgfsetmiterjoin%
\definecolor{currentfill}{rgb}{1.000000,1.000000,1.000000}%
\pgfsetfillcolor{currentfill}%
\pgfsetlinewidth{0.000000pt}%
\definecolor{currentstroke}{rgb}{0.000000,0.000000,0.000000}%
\pgfsetstrokecolor{currentstroke}%
\pgfsetstrokeopacity{0.000000}%
\pgfsetdash{}{0pt}%
\pgfpathmoveto{\pgfqpoint{2.470586in}{0.476543in}}%
\pgfpathlineto{\pgfqpoint{4.290975in}{0.476543in}}%
\pgfpathlineto{\pgfqpoint{4.290975in}{2.132543in}}%
\pgfpathlineto{\pgfqpoint{2.470586in}{2.132543in}}%
\pgfpathlineto{\pgfqpoint{2.470586in}{0.476543in}}%
\pgfpathclose%
\pgfusepath{fill}%
\end{pgfscope}%
\begin{pgfscope}%
\pgfpathrectangle{\pgfqpoint{2.470586in}{0.476543in}}{\pgfqpoint{1.820388in}{1.656000in}}%
\pgfusepath{clip}%
\pgfsetbuttcap%
\pgfsetroundjoin%
\definecolor{currentfill}{rgb}{0.207536,0.260208,0.608151}%
\pgfsetfillcolor{currentfill}%
\pgfsetlinewidth{0.000000pt}%
\definecolor{currentstroke}{rgb}{0.000000,0.000000,0.000000}%
\pgfsetstrokecolor{currentstroke}%
\pgfsetdash{}{0pt}%
\pgfpathmoveto{\pgfqpoint{2.495870in}{0.476543in}}%
\pgfpathlineto{\pgfqpoint{2.521153in}{0.476543in}}%
\pgfpathlineto{\pgfqpoint{2.546436in}{0.476543in}}%
\pgfpathlineto{\pgfqpoint{2.571719in}{0.476543in}}%
\pgfpathlineto{\pgfqpoint{2.597002in}{0.476543in}}%
\pgfpathlineto{\pgfqpoint{2.622285in}{0.476543in}}%
\pgfpathlineto{\pgfqpoint{2.647569in}{0.476543in}}%
\pgfpathlineto{\pgfqpoint{2.672852in}{0.476543in}}%
\pgfpathlineto{\pgfqpoint{2.698135in}{0.476543in}}%
\pgfpathlineto{\pgfqpoint{2.723418in}{0.476543in}}%
\pgfpathlineto{\pgfqpoint{2.748701in}{0.476543in}}%
\pgfpathlineto{\pgfqpoint{2.773984in}{0.476543in}}%
\pgfpathlineto{\pgfqpoint{2.799268in}{0.476543in}}%
\pgfpathlineto{\pgfqpoint{2.824551in}{0.476543in}}%
\pgfpathlineto{\pgfqpoint{2.849834in}{0.476543in}}%
\pgfpathlineto{\pgfqpoint{2.875117in}{0.476543in}}%
\pgfpathlineto{\pgfqpoint{2.900400in}{0.476543in}}%
\pgfpathlineto{\pgfqpoint{2.925683in}{0.476543in}}%
\pgfpathlineto{\pgfqpoint{2.950967in}{0.476543in}}%
\pgfpathlineto{\pgfqpoint{2.976250in}{0.476543in}}%
\pgfpathlineto{\pgfqpoint{3.001533in}{0.476543in}}%
\pgfpathlineto{\pgfqpoint{3.026816in}{0.476543in}}%
\pgfpathlineto{\pgfqpoint{3.052099in}{0.476543in}}%
\pgfpathlineto{\pgfqpoint{3.077382in}{0.476543in}}%
\pgfpathlineto{\pgfqpoint{3.102666in}{0.476543in}}%
\pgfpathlineto{\pgfqpoint{3.127949in}{0.476543in}}%
\pgfpathlineto{\pgfqpoint{3.153232in}{0.476543in}}%
\pgfpathlineto{\pgfqpoint{3.178515in}{0.476543in}}%
\pgfpathlineto{\pgfqpoint{3.203798in}{0.476543in}}%
\pgfpathlineto{\pgfqpoint{3.229082in}{0.476543in}}%
\pgfpathlineto{\pgfqpoint{3.254365in}{0.476543in}}%
\pgfpathlineto{\pgfqpoint{3.279648in}{0.476543in}}%
\pgfpathlineto{\pgfqpoint{3.304931in}{0.476543in}}%
\pgfpathlineto{\pgfqpoint{3.330214in}{0.476543in}}%
\pgfpathlineto{\pgfqpoint{3.355497in}{0.476543in}}%
\pgfpathlineto{\pgfqpoint{3.380781in}{0.476543in}}%
\pgfpathlineto{\pgfqpoint{3.406064in}{0.476543in}}%
\pgfpathlineto{\pgfqpoint{3.431347in}{0.476543in}}%
\pgfpathlineto{\pgfqpoint{3.456630in}{0.476543in}}%
\pgfpathlineto{\pgfqpoint{3.481913in}{0.476543in}}%
\pgfpathlineto{\pgfqpoint{3.507196in}{0.476543in}}%
\pgfpathlineto{\pgfqpoint{3.532480in}{0.476543in}}%
\pgfpathlineto{\pgfqpoint{3.557763in}{0.476543in}}%
\pgfpathlineto{\pgfqpoint{3.583046in}{0.476543in}}%
\pgfpathlineto{\pgfqpoint{3.608329in}{0.476543in}}%
\pgfpathlineto{\pgfqpoint{3.633612in}{0.476543in}}%
\pgfpathlineto{\pgfqpoint{3.658895in}{0.476543in}}%
\pgfpathlineto{\pgfqpoint{3.684179in}{0.476543in}}%
\pgfpathlineto{\pgfqpoint{3.709462in}{0.476543in}}%
\pgfpathlineto{\pgfqpoint{3.734745in}{0.476543in}}%
\pgfpathlineto{\pgfqpoint{3.760028in}{0.476543in}}%
\pgfpathlineto{\pgfqpoint{3.785311in}{0.476543in}}%
\pgfpathlineto{\pgfqpoint{3.810594in}{0.476543in}}%
\pgfpathlineto{\pgfqpoint{3.835878in}{0.476543in}}%
\pgfpathlineto{\pgfqpoint{3.861161in}{0.476543in}}%
\pgfpathlineto{\pgfqpoint{3.886444in}{0.476543in}}%
\pgfpathlineto{\pgfqpoint{3.911727in}{0.476543in}}%
\pgfpathlineto{\pgfqpoint{3.937010in}{0.476543in}}%
\pgfpathlineto{\pgfqpoint{3.962293in}{0.476543in}}%
\pgfpathlineto{\pgfqpoint{3.987577in}{0.476543in}}%
\pgfpathlineto{\pgfqpoint{4.012860in}{0.476543in}}%
\pgfpathlineto{\pgfqpoint{4.038143in}{0.476543in}}%
\pgfpathlineto{\pgfqpoint{4.063426in}{0.476543in}}%
\pgfpathlineto{\pgfqpoint{4.088709in}{0.476543in}}%
\pgfpathlineto{\pgfqpoint{4.113993in}{0.476543in}}%
\pgfpathlineto{\pgfqpoint{4.139276in}{0.476543in}}%
\pgfpathlineto{\pgfqpoint{4.164559in}{0.476543in}}%
\pgfpathlineto{\pgfqpoint{4.189842in}{0.476543in}}%
\pgfpathlineto{\pgfqpoint{4.215125in}{0.476543in}}%
\pgfpathlineto{\pgfqpoint{4.240408in}{0.476543in}}%
\pgfpathlineto{\pgfqpoint{4.265692in}{0.476543in}}%
\pgfpathlineto{\pgfqpoint{4.290975in}{0.476543in}}%
\pgfpathlineto{\pgfqpoint{4.290975in}{0.489791in}}%
\pgfpathlineto{\pgfqpoint{4.290975in}{0.503039in}}%
\pgfpathlineto{\pgfqpoint{4.290975in}{0.516287in}}%
\pgfpathlineto{\pgfqpoint{4.290975in}{0.529535in}}%
\pgfpathlineto{\pgfqpoint{4.290975in}{0.542783in}}%
\pgfpathlineto{\pgfqpoint{4.290975in}{0.556031in}}%
\pgfpathlineto{\pgfqpoint{4.290975in}{0.569279in}}%
\pgfpathlineto{\pgfqpoint{4.290975in}{0.582527in}}%
\pgfpathlineto{\pgfqpoint{4.290975in}{0.595775in}}%
\pgfpathlineto{\pgfqpoint{4.290975in}{0.609023in}}%
\pgfpathlineto{\pgfqpoint{4.290975in}{0.613212in}}%
\pgfpathlineto{\pgfqpoint{4.265692in}{0.611933in}}%
\pgfpathlineto{\pgfqpoint{4.240408in}{0.610784in}}%
\pgfpathlineto{\pgfqpoint{4.215125in}{0.609749in}}%
\pgfpathlineto{\pgfqpoint{4.195664in}{0.609023in}}%
\pgfpathlineto{\pgfqpoint{4.189842in}{0.608667in}}%
\pgfpathlineto{\pgfqpoint{4.164559in}{0.607176in}}%
\pgfpathlineto{\pgfqpoint{4.139276in}{0.605749in}}%
\pgfpathlineto{\pgfqpoint{4.113993in}{0.604384in}}%
\pgfpathlineto{\pgfqpoint{4.088709in}{0.603075in}}%
\pgfpathlineto{\pgfqpoint{4.063426in}{0.601819in}}%
\pgfpathlineto{\pgfqpoint{4.038143in}{0.600612in}}%
\pgfpathlineto{\pgfqpoint{4.012860in}{0.599452in}}%
\pgfpathlineto{\pgfqpoint{3.987577in}{0.598336in}}%
\pgfpathlineto{\pgfqpoint{3.962293in}{0.597260in}}%
\pgfpathlineto{\pgfqpoint{3.937010in}{0.596222in}}%
\pgfpathlineto{\pgfqpoint{3.925760in}{0.595775in}}%
\pgfpathlineto{\pgfqpoint{3.911727in}{0.595558in}}%
\pgfpathlineto{\pgfqpoint{3.886444in}{0.595186in}}%
\pgfpathlineto{\pgfqpoint{3.861161in}{0.594834in}}%
\pgfpathlineto{\pgfqpoint{3.835878in}{0.594503in}}%
\pgfpathlineto{\pgfqpoint{3.810594in}{0.594191in}}%
\pgfpathlineto{\pgfqpoint{3.785311in}{0.593897in}}%
\pgfpathlineto{\pgfqpoint{3.760028in}{0.593621in}}%
\pgfpathlineto{\pgfqpoint{3.734745in}{0.593361in}}%
\pgfpathlineto{\pgfqpoint{3.709462in}{0.593118in}}%
\pgfpathlineto{\pgfqpoint{3.684179in}{0.592890in}}%
\pgfpathlineto{\pgfqpoint{3.658895in}{0.592677in}}%
\pgfpathlineto{\pgfqpoint{3.633612in}{0.592478in}}%
\pgfpathlineto{\pgfqpoint{3.608329in}{0.592293in}}%
\pgfpathlineto{\pgfqpoint{3.583046in}{0.592121in}}%
\pgfpathlineto{\pgfqpoint{3.557763in}{0.591962in}}%
\pgfpathlineto{\pgfqpoint{3.532480in}{0.591816in}}%
\pgfpathlineto{\pgfqpoint{3.507196in}{0.591682in}}%
\pgfpathlineto{\pgfqpoint{3.481913in}{0.591560in}}%
\pgfpathlineto{\pgfqpoint{3.456630in}{0.591449in}}%
\pgfpathlineto{\pgfqpoint{3.431347in}{0.591350in}}%
\pgfpathlineto{\pgfqpoint{3.406064in}{0.591262in}}%
\pgfpathlineto{\pgfqpoint{3.380781in}{0.591185in}}%
\pgfpathlineto{\pgfqpoint{3.355497in}{0.591119in}}%
\pgfpathlineto{\pgfqpoint{3.330214in}{0.591063in}}%
\pgfpathlineto{\pgfqpoint{3.304931in}{0.591018in}}%
\pgfpathlineto{\pgfqpoint{3.279648in}{0.590983in}}%
\pgfpathlineto{\pgfqpoint{3.254365in}{0.590959in}}%
\pgfpathlineto{\pgfqpoint{3.229082in}{0.590945in}}%
\pgfpathlineto{\pgfqpoint{3.203798in}{0.590941in}}%
\pgfpathlineto{\pgfqpoint{3.178515in}{0.590948in}}%
\pgfpathlineto{\pgfqpoint{3.153232in}{0.590966in}}%
\pgfpathlineto{\pgfqpoint{3.127949in}{0.590993in}}%
\pgfpathlineto{\pgfqpoint{3.102666in}{0.591032in}}%
\pgfpathlineto{\pgfqpoint{3.077382in}{0.591081in}}%
\pgfpathlineto{\pgfqpoint{3.052099in}{0.591141in}}%
\pgfpathlineto{\pgfqpoint{3.026816in}{0.591212in}}%
\pgfpathlineto{\pgfqpoint{3.001533in}{0.591294in}}%
\pgfpathlineto{\pgfqpoint{2.976250in}{0.591388in}}%
\pgfpathlineto{\pgfqpoint{2.950967in}{0.591494in}}%
\pgfpathlineto{\pgfqpoint{2.925683in}{0.591612in}}%
\pgfpathlineto{\pgfqpoint{2.900400in}{0.591743in}}%
\pgfpathlineto{\pgfqpoint{2.875117in}{0.591887in}}%
\pgfpathlineto{\pgfqpoint{2.849834in}{0.592044in}}%
\pgfpathlineto{\pgfqpoint{2.824551in}{0.592215in}}%
\pgfpathlineto{\pgfqpoint{2.799268in}{0.592401in}}%
\pgfpathlineto{\pgfqpoint{2.773984in}{0.592602in}}%
\pgfpathlineto{\pgfqpoint{2.748701in}{0.592818in}}%
\pgfpathlineto{\pgfqpoint{2.723418in}{0.593051in}}%
\pgfpathlineto{\pgfqpoint{2.698135in}{0.593302in}}%
\pgfpathlineto{\pgfqpoint{2.672852in}{0.593571in}}%
\pgfpathlineto{\pgfqpoint{2.647569in}{0.593858in}}%
\pgfpathlineto{\pgfqpoint{2.622285in}{0.594167in}}%
\pgfpathlineto{\pgfqpoint{2.597002in}{0.594496in}}%
\pgfpathlineto{\pgfqpoint{2.571719in}{0.594849in}}%
\pgfpathlineto{\pgfqpoint{2.546436in}{0.595225in}}%
\pgfpathlineto{\pgfqpoint{2.521153in}{0.595627in}}%
\pgfpathlineto{\pgfqpoint{2.512340in}{0.595775in}}%
\pgfpathlineto{\pgfqpoint{2.495870in}{0.597305in}}%
\pgfpathlineto{\pgfqpoint{2.470586in}{0.599811in}}%
\pgfpathlineto{\pgfqpoint{2.470586in}{0.595775in}}%
\pgfpathlineto{\pgfqpoint{2.470586in}{0.582527in}}%
\pgfpathlineto{\pgfqpoint{2.470586in}{0.569279in}}%
\pgfpathlineto{\pgfqpoint{2.470586in}{0.556031in}}%
\pgfpathlineto{\pgfqpoint{2.470586in}{0.542783in}}%
\pgfpathlineto{\pgfqpoint{2.470586in}{0.529535in}}%
\pgfpathlineto{\pgfqpoint{2.470586in}{0.516287in}}%
\pgfpathlineto{\pgfqpoint{2.470586in}{0.503039in}}%
\pgfpathlineto{\pgfqpoint{2.470586in}{0.489791in}}%
\pgfpathlineto{\pgfqpoint{2.470586in}{0.476543in}}%
\pgfpathlineto{\pgfqpoint{2.495870in}{0.476543in}}%
\pgfpathclose%
\pgfpathmoveto{\pgfqpoint{4.290975in}{0.780126in}}%
\pgfpathlineto{\pgfqpoint{4.290975in}{0.781247in}}%
\pgfpathlineto{\pgfqpoint{4.290975in}{0.794495in}}%
\pgfpathlineto{\pgfqpoint{4.290975in}{0.807743in}}%
\pgfpathlineto{\pgfqpoint{4.290975in}{0.820991in}}%
\pgfpathlineto{\pgfqpoint{4.290975in}{0.834239in}}%
\pgfpathlineto{\pgfqpoint{4.290975in}{0.847487in}}%
\pgfpathlineto{\pgfqpoint{4.290975in}{0.860735in}}%
\pgfpathlineto{\pgfqpoint{4.290975in}{0.873983in}}%
\pgfpathlineto{\pgfqpoint{4.290975in}{0.887231in}}%
\pgfpathlineto{\pgfqpoint{4.290975in}{0.900479in}}%
\pgfpathlineto{\pgfqpoint{4.290975in}{0.913727in}}%
\pgfpathlineto{\pgfqpoint{4.290975in}{0.926975in}}%
\pgfpathlineto{\pgfqpoint{4.290975in}{0.940223in}}%
\pgfpathlineto{\pgfqpoint{4.290975in}{0.953471in}}%
\pgfpathlineto{\pgfqpoint{4.290975in}{0.966719in}}%
\pgfpathlineto{\pgfqpoint{4.290975in}{0.979967in}}%
\pgfpathlineto{\pgfqpoint{4.290975in}{0.993215in}}%
\pgfpathlineto{\pgfqpoint{4.290975in}{1.006463in}}%
\pgfpathlineto{\pgfqpoint{4.290975in}{1.019711in}}%
\pgfpathlineto{\pgfqpoint{4.290975in}{1.032959in}}%
\pgfpathlineto{\pgfqpoint{4.290975in}{1.046207in}}%
\pgfpathlineto{\pgfqpoint{4.290975in}{1.059455in}}%
\pgfpathlineto{\pgfqpoint{4.290975in}{1.061804in}}%
\pgfpathlineto{\pgfqpoint{4.265692in}{1.066254in}}%
\pgfpathlineto{\pgfqpoint{4.240408in}{1.071413in}}%
\pgfpathlineto{\pgfqpoint{4.234961in}{1.072703in}}%
\pgfpathlineto{\pgfqpoint{4.219278in}{1.085951in}}%
\pgfpathlineto{\pgfqpoint{4.215125in}{1.087175in}}%
\pgfpathlineto{\pgfqpoint{4.189842in}{1.095745in}}%
\pgfpathlineto{\pgfqpoint{4.180881in}{1.099199in}}%
\pgfpathlineto{\pgfqpoint{4.177530in}{1.112447in}}%
\pgfpathlineto{\pgfqpoint{4.164559in}{1.116231in}}%
\pgfpathlineto{\pgfqpoint{4.139276in}{1.124771in}}%
\pgfpathlineto{\pgfqpoint{4.136891in}{1.125695in}}%
\pgfpathlineto{\pgfqpoint{4.127624in}{1.138943in}}%
\pgfpathlineto{\pgfqpoint{4.113993in}{1.143133in}}%
\pgfpathlineto{\pgfqpoint{4.088709in}{1.151649in}}%
\pgfpathlineto{\pgfqpoint{4.087235in}{1.152191in}}%
\pgfpathlineto{\pgfqpoint{4.073090in}{1.165439in}}%
\pgfpathlineto{\pgfqpoint{4.063426in}{1.170471in}}%
\pgfpathlineto{\pgfqpoint{4.048827in}{1.178687in}}%
\pgfpathlineto{\pgfqpoint{4.038143in}{1.191767in}}%
\pgfpathlineto{\pgfqpoint{4.038013in}{1.191935in}}%
\pgfpathlineto{\pgfqpoint{4.017019in}{1.205183in}}%
\pgfpathlineto{\pgfqpoint{4.012860in}{1.209545in}}%
\pgfpathlineto{\pgfqpoint{4.005075in}{1.218431in}}%
\pgfpathlineto{\pgfqpoint{3.987577in}{1.229687in}}%
\pgfpathlineto{\pgfqpoint{3.984582in}{1.231679in}}%
\pgfpathlineto{\pgfqpoint{3.966885in}{1.244927in}}%
\pgfpathlineto{\pgfqpoint{3.962293in}{1.248858in}}%
\pgfpathlineto{\pgfqpoint{3.951598in}{1.258175in}}%
\pgfpathlineto{\pgfqpoint{3.937010in}{1.268786in}}%
\pgfpathlineto{\pgfqpoint{3.933556in}{1.271423in}}%
\pgfpathlineto{\pgfqpoint{3.923291in}{1.284671in}}%
\pgfpathlineto{\pgfqpoint{3.911727in}{1.296434in}}%
\pgfpathlineto{\pgfqpoint{3.910313in}{1.297919in}}%
\pgfpathlineto{\pgfqpoint{3.902774in}{1.311167in}}%
\pgfpathlineto{\pgfqpoint{3.886444in}{1.323292in}}%
\pgfpathlineto{\pgfqpoint{3.885009in}{1.324415in}}%
\pgfpathlineto{\pgfqpoint{3.876856in}{1.337663in}}%
\pgfpathlineto{\pgfqpoint{3.861161in}{1.348106in}}%
\pgfpathlineto{\pgfqpoint{3.857025in}{1.350911in}}%
\pgfpathlineto{\pgfqpoint{3.851799in}{1.364159in}}%
\pgfpathlineto{\pgfqpoint{3.837588in}{1.377407in}}%
\pgfpathlineto{\pgfqpoint{3.835878in}{1.383153in}}%
\pgfpathlineto{\pgfqpoint{3.833453in}{1.390655in}}%
\pgfpathlineto{\pgfqpoint{3.817844in}{1.403903in}}%
\pgfpathlineto{\pgfqpoint{3.812781in}{1.417151in}}%
\pgfpathlineto{\pgfqpoint{3.810594in}{1.418870in}}%
\pgfpathlineto{\pgfqpoint{3.795923in}{1.430399in}}%
\pgfpathlineto{\pgfqpoint{3.790690in}{1.443647in}}%
\pgfpathlineto{\pgfqpoint{3.785311in}{1.449290in}}%
\pgfpathlineto{\pgfqpoint{3.777600in}{1.456895in}}%
\pgfpathlineto{\pgfqpoint{3.772185in}{1.470143in}}%
\pgfpathlineto{\pgfqpoint{3.762968in}{1.483391in}}%
\pgfpathlineto{\pgfqpoint{3.760028in}{1.493004in}}%
\pgfpathlineto{\pgfqpoint{3.758723in}{1.496639in}}%
\pgfpathlineto{\pgfqpoint{3.747628in}{1.509887in}}%
\pgfpathlineto{\pgfqpoint{3.740021in}{1.523135in}}%
\pgfpathlineto{\pgfqpoint{3.734745in}{1.530743in}}%
\pgfpathlineto{\pgfqpoint{3.730122in}{1.536383in}}%
\pgfpathlineto{\pgfqpoint{3.720893in}{1.549631in}}%
\pgfpathlineto{\pgfqpoint{3.715885in}{1.562879in}}%
\pgfpathlineto{\pgfqpoint{3.709462in}{1.576030in}}%
\pgfpathlineto{\pgfqpoint{3.709402in}{1.576127in}}%
\pgfpathlineto{\pgfqpoint{3.702960in}{1.589375in}}%
\pgfpathlineto{\pgfqpoint{3.693799in}{1.602623in}}%
\pgfpathlineto{\pgfqpoint{3.689093in}{1.615871in}}%
\pgfpathlineto{\pgfqpoint{3.684179in}{1.622559in}}%
\pgfpathlineto{\pgfqpoint{3.678135in}{1.629119in}}%
\pgfpathlineto{\pgfqpoint{3.674997in}{1.642367in}}%
\pgfpathlineto{\pgfqpoint{3.666239in}{1.655615in}}%
\pgfpathlineto{\pgfqpoint{3.663758in}{1.668863in}}%
\pgfpathlineto{\pgfqpoint{3.658895in}{1.677431in}}%
\pgfpathlineto{\pgfqpoint{3.655452in}{1.682111in}}%
\pgfpathlineto{\pgfqpoint{3.652448in}{1.695359in}}%
\pgfpathlineto{\pgfqpoint{3.641734in}{1.708607in}}%
\pgfpathlineto{\pgfqpoint{3.638571in}{1.721855in}}%
\pgfpathlineto{\pgfqpoint{3.633612in}{1.730024in}}%
\pgfpathlineto{\pgfqpoint{3.629214in}{1.735103in}}%
\pgfpathlineto{\pgfqpoint{3.625616in}{1.748351in}}%
\pgfpathlineto{\pgfqpoint{3.619325in}{1.761599in}}%
\pgfpathlineto{\pgfqpoint{3.616739in}{1.774847in}}%
\pgfpathlineto{\pgfqpoint{3.610093in}{1.788095in}}%
\pgfpathlineto{\pgfqpoint{3.608329in}{1.794402in}}%
\pgfpathlineto{\pgfqpoint{3.605197in}{1.801343in}}%
\pgfpathlineto{\pgfqpoint{3.597357in}{1.814591in}}%
\pgfpathlineto{\pgfqpoint{3.592469in}{1.827839in}}%
\pgfpathlineto{\pgfqpoint{3.588818in}{1.841087in}}%
\pgfpathlineto{\pgfqpoint{3.584875in}{1.854335in}}%
\pgfpathlineto{\pgfqpoint{3.583046in}{1.860813in}}%
\pgfpathlineto{\pgfqpoint{3.579859in}{1.867583in}}%
\pgfpathlineto{\pgfqpoint{3.573168in}{1.880831in}}%
\pgfpathlineto{\pgfqpoint{3.569018in}{1.894079in}}%
\pgfpathlineto{\pgfqpoint{3.562220in}{1.907327in}}%
\pgfpathlineto{\pgfqpoint{3.559754in}{1.920575in}}%
\pgfpathlineto{\pgfqpoint{3.557763in}{1.925624in}}%
\pgfpathlineto{\pgfqpoint{3.551965in}{1.933823in}}%
\pgfpathlineto{\pgfqpoint{3.550278in}{1.947071in}}%
\pgfpathlineto{\pgfqpoint{3.544237in}{1.960319in}}%
\pgfpathlineto{\pgfqpoint{3.542198in}{1.973567in}}%
\pgfpathlineto{\pgfqpoint{3.535446in}{1.986815in}}%
\pgfpathlineto{\pgfqpoint{3.533553in}{2.000063in}}%
\pgfpathlineto{\pgfqpoint{3.532480in}{2.002470in}}%
\pgfpathlineto{\pgfqpoint{3.523795in}{2.013311in}}%
\pgfpathlineto{\pgfqpoint{3.521979in}{2.026559in}}%
\pgfpathlineto{\pgfqpoint{3.516295in}{2.039807in}}%
\pgfpathlineto{\pgfqpoint{3.514529in}{2.053055in}}%
\pgfpathlineto{\pgfqpoint{3.509423in}{2.066303in}}%
\pgfpathlineto{\pgfqpoint{3.507451in}{2.079551in}}%
\pgfpathlineto{\pgfqpoint{3.507196in}{2.080203in}}%
\pgfpathlineto{\pgfqpoint{3.498113in}{2.092799in}}%
\pgfpathlineto{\pgfqpoint{3.494257in}{2.106047in}}%
\pgfpathlineto{\pgfqpoint{3.489901in}{2.119295in}}%
\pgfpathlineto{\pgfqpoint{3.487284in}{2.132543in}}%
\pgfpathlineto{\pgfqpoint{3.481913in}{2.132543in}}%
\pgfpathlineto{\pgfqpoint{3.456630in}{2.132543in}}%
\pgfpathlineto{\pgfqpoint{3.431347in}{2.132543in}}%
\pgfpathlineto{\pgfqpoint{3.430617in}{2.132543in}}%
\pgfpathlineto{\pgfqpoint{3.431347in}{2.128980in}}%
\pgfpathlineto{\pgfqpoint{3.434645in}{2.119295in}}%
\pgfpathlineto{\pgfqpoint{3.443466in}{2.106047in}}%
\pgfpathlineto{\pgfqpoint{3.446062in}{2.092799in}}%
\pgfpathlineto{\pgfqpoint{3.450583in}{2.079551in}}%
\pgfpathlineto{\pgfqpoint{3.452924in}{2.066303in}}%
\pgfpathlineto{\pgfqpoint{3.456630in}{2.057029in}}%
\pgfpathlineto{\pgfqpoint{3.459471in}{2.053055in}}%
\pgfpathlineto{\pgfqpoint{3.464980in}{2.039807in}}%
\pgfpathlineto{\pgfqpoint{3.470927in}{2.026559in}}%
\pgfpathlineto{\pgfqpoint{3.474949in}{2.013311in}}%
\pgfpathlineto{\pgfqpoint{3.477658in}{2.000063in}}%
\pgfpathlineto{\pgfqpoint{3.480688in}{1.986815in}}%
\pgfpathlineto{\pgfqpoint{3.481913in}{1.980429in}}%
\pgfpathlineto{\pgfqpoint{3.484511in}{1.973567in}}%
\pgfpathlineto{\pgfqpoint{3.490463in}{1.960319in}}%
\pgfpathlineto{\pgfqpoint{3.494359in}{1.947071in}}%
\pgfpathlineto{\pgfqpoint{3.500577in}{1.933823in}}%
\pgfpathlineto{\pgfqpoint{3.502751in}{1.920575in}}%
\pgfpathlineto{\pgfqpoint{3.507142in}{1.907327in}}%
\pgfpathlineto{\pgfqpoint{3.507196in}{1.906737in}}%
\pgfpathlineto{\pgfqpoint{3.509664in}{1.894079in}}%
\pgfpathlineto{\pgfqpoint{3.516661in}{1.880831in}}%
\pgfpathlineto{\pgfqpoint{3.519021in}{1.867583in}}%
\pgfpathlineto{\pgfqpoint{3.525677in}{1.854335in}}%
\pgfpathlineto{\pgfqpoint{3.527445in}{1.841087in}}%
\pgfpathlineto{\pgfqpoint{3.532480in}{1.828197in}}%
\pgfpathlineto{\pgfqpoint{3.532780in}{1.827839in}}%
\pgfpathlineto{\pgfqpoint{3.535146in}{1.814591in}}%
\pgfpathlineto{\pgfqpoint{3.542159in}{1.801343in}}%
\pgfpathlineto{\pgfqpoint{3.544135in}{1.788095in}}%
\pgfpathlineto{\pgfqpoint{3.549948in}{1.774847in}}%
\pgfpathlineto{\pgfqpoint{3.552191in}{1.761599in}}%
\pgfpathlineto{\pgfqpoint{3.557524in}{1.748351in}}%
\pgfpathlineto{\pgfqpoint{3.557763in}{1.747220in}}%
\pgfpathlineto{\pgfqpoint{3.562794in}{1.735103in}}%
\pgfpathlineto{\pgfqpoint{3.568614in}{1.721855in}}%
\pgfpathlineto{\pgfqpoint{3.572028in}{1.708607in}}%
\pgfpathlineto{\pgfqpoint{3.575588in}{1.695359in}}%
\pgfpathlineto{\pgfqpoint{3.578872in}{1.682111in}}%
\pgfpathlineto{\pgfqpoint{3.582876in}{1.668863in}}%
\pgfpathlineto{\pgfqpoint{3.583046in}{1.668386in}}%
\pgfpathlineto{\pgfqpoint{3.590920in}{1.655615in}}%
\pgfpathlineto{\pgfqpoint{3.594736in}{1.642367in}}%
\pgfpathlineto{\pgfqpoint{3.600429in}{1.629119in}}%
\pgfpathlineto{\pgfqpoint{3.602272in}{1.615871in}}%
\pgfpathlineto{\pgfqpoint{3.607096in}{1.602623in}}%
\pgfpathlineto{\pgfqpoint{3.608329in}{1.596141in}}%
\pgfpathlineto{\pgfqpoint{3.610627in}{1.589375in}}%
\pgfpathlineto{\pgfqpoint{3.620034in}{1.576127in}}%
\pgfpathlineto{\pgfqpoint{3.622787in}{1.562879in}}%
\pgfpathlineto{\pgfqpoint{3.630589in}{1.549631in}}%
\pgfpathlineto{\pgfqpoint{3.632417in}{1.536383in}}%
\pgfpathlineto{\pgfqpoint{3.633612in}{1.533559in}}%
\pgfpathlineto{\pgfqpoint{3.640377in}{1.523135in}}%
\pgfpathlineto{\pgfqpoint{3.642355in}{1.509887in}}%
\pgfpathlineto{\pgfqpoint{3.650071in}{1.496639in}}%
\pgfpathlineto{\pgfqpoint{3.652707in}{1.483391in}}%
\pgfpathlineto{\pgfqpoint{3.658895in}{1.473220in}}%
\pgfpathlineto{\pgfqpoint{3.661751in}{1.470143in}}%
\pgfpathlineto{\pgfqpoint{3.666779in}{1.456895in}}%
\pgfpathlineto{\pgfqpoint{3.673992in}{1.443647in}}%
\pgfpathlineto{\pgfqpoint{3.677505in}{1.430399in}}%
\pgfpathlineto{\pgfqpoint{3.683089in}{1.417151in}}%
\pgfpathlineto{\pgfqpoint{3.684179in}{1.413782in}}%
\pgfpathlineto{\pgfqpoint{3.688920in}{1.403903in}}%
\pgfpathlineto{\pgfqpoint{3.697100in}{1.390655in}}%
\pgfpathlineto{\pgfqpoint{3.704135in}{1.377407in}}%
\pgfpathlineto{\pgfqpoint{3.708970in}{1.364159in}}%
\pgfpathlineto{\pgfqpoint{3.709462in}{1.363223in}}%
\pgfpathlineto{\pgfqpoint{3.718454in}{1.350911in}}%
\pgfpathlineto{\pgfqpoint{3.722669in}{1.337663in}}%
\pgfpathlineto{\pgfqpoint{3.729665in}{1.324415in}}%
\pgfpathlineto{\pgfqpoint{3.733174in}{1.311167in}}%
\pgfpathlineto{\pgfqpoint{3.734745in}{1.309044in}}%
\pgfpathlineto{\pgfqpoint{3.745308in}{1.297919in}}%
\pgfpathlineto{\pgfqpoint{3.750175in}{1.284671in}}%
\pgfpathlineto{\pgfqpoint{3.760028in}{1.273767in}}%
\pgfpathlineto{\pgfqpoint{3.762768in}{1.271423in}}%
\pgfpathlineto{\pgfqpoint{3.766932in}{1.258175in}}%
\pgfpathlineto{\pgfqpoint{3.777532in}{1.244927in}}%
\pgfpathlineto{\pgfqpoint{3.780116in}{1.231679in}}%
\pgfpathlineto{\pgfqpoint{3.785311in}{1.225809in}}%
\pgfpathlineto{\pgfqpoint{3.793229in}{1.218431in}}%
\pgfpathlineto{\pgfqpoint{3.798140in}{1.205183in}}%
\pgfpathlineto{\pgfqpoint{3.810594in}{1.192911in}}%
\pgfpathlineto{\pgfqpoint{3.811820in}{1.191935in}}%
\pgfpathlineto{\pgfqpoint{3.818311in}{1.178687in}}%
\pgfpathlineto{\pgfqpoint{3.829979in}{1.165439in}}%
\pgfpathlineto{\pgfqpoint{3.835528in}{1.152191in}}%
\pgfpathlineto{\pgfqpoint{3.835878in}{1.151727in}}%
\pgfpathlineto{\pgfqpoint{3.847352in}{1.138943in}}%
\pgfpathlineto{\pgfqpoint{3.854399in}{1.125695in}}%
\pgfpathlineto{\pgfqpoint{3.861161in}{1.117851in}}%
\pgfpathlineto{\pgfqpoint{3.866817in}{1.112447in}}%
\pgfpathlineto{\pgfqpoint{3.880476in}{1.099199in}}%
\pgfpathlineto{\pgfqpoint{3.886444in}{1.091657in}}%
\pgfpathlineto{\pgfqpoint{3.891904in}{1.085951in}}%
\pgfpathlineto{\pgfqpoint{3.906355in}{1.072703in}}%
\pgfpathlineto{\pgfqpoint{3.911727in}{1.060920in}}%
\pgfpathlineto{\pgfqpoint{3.912534in}{1.059455in}}%
\pgfpathlineto{\pgfqpoint{3.926041in}{1.046207in}}%
\pgfpathlineto{\pgfqpoint{3.932517in}{1.032959in}}%
\pgfpathlineto{\pgfqpoint{3.937010in}{1.030014in}}%
\pgfpathlineto{\pgfqpoint{3.953933in}{1.019711in}}%
\pgfpathlineto{\pgfqpoint{3.961322in}{1.006463in}}%
\pgfpathlineto{\pgfqpoint{3.962293in}{1.005889in}}%
\pgfpathlineto{\pgfqpoint{3.985544in}{0.993215in}}%
\pgfpathlineto{\pgfqpoint{3.987577in}{0.988429in}}%
\pgfpathlineto{\pgfqpoint{3.991806in}{0.979967in}}%
\pgfpathlineto{\pgfqpoint{4.012860in}{0.966852in}}%
\pgfpathlineto{\pgfqpoint{4.013086in}{0.966719in}}%
\pgfpathlineto{\pgfqpoint{4.018954in}{0.953471in}}%
\pgfpathlineto{\pgfqpoint{4.038143in}{0.943176in}}%
\pgfpathlineto{\pgfqpoint{4.043906in}{0.940223in}}%
\pgfpathlineto{\pgfqpoint{4.051140in}{0.926975in}}%
\pgfpathlineto{\pgfqpoint{4.063426in}{0.921237in}}%
\pgfpathlineto{\pgfqpoint{4.080712in}{0.913727in}}%
\pgfpathlineto{\pgfqpoint{4.088709in}{0.906342in}}%
\pgfpathlineto{\pgfqpoint{4.095724in}{0.900479in}}%
\pgfpathlineto{\pgfqpoint{4.113993in}{0.890243in}}%
\pgfpathlineto{\pgfqpoint{4.119760in}{0.887231in}}%
\pgfpathlineto{\pgfqpoint{4.131068in}{0.873983in}}%
\pgfpathlineto{\pgfqpoint{4.139276in}{0.867984in}}%
\pgfpathlineto{\pgfqpoint{4.150049in}{0.860735in}}%
\pgfpathlineto{\pgfqpoint{4.164559in}{0.848303in}}%
\pgfpathlineto{\pgfqpoint{4.165568in}{0.847487in}}%
\pgfpathlineto{\pgfqpoint{4.188883in}{0.834239in}}%
\pgfpathlineto{\pgfqpoint{4.189842in}{0.833868in}}%
\pgfpathlineto{\pgfqpoint{4.215125in}{0.824521in}}%
\pgfpathlineto{\pgfqpoint{4.224935in}{0.820991in}}%
\pgfpathlineto{\pgfqpoint{4.240408in}{0.809527in}}%
\pgfpathlineto{\pgfqpoint{4.243112in}{0.807743in}}%
\pgfpathlineto{\pgfqpoint{4.265692in}{0.799110in}}%
\pgfpathlineto{\pgfqpoint{4.278255in}{0.794495in}}%
\pgfpathlineto{\pgfqpoint{4.288057in}{0.781247in}}%
\pgfpathlineto{\pgfqpoint{4.290975in}{0.780126in}}%
\pgfpathclose%
\pgfpathmoveto{\pgfqpoint{2.495870in}{0.847100in}}%
\pgfpathlineto{\pgfqpoint{2.496986in}{0.847487in}}%
\pgfpathlineto{\pgfqpoint{2.502188in}{0.860735in}}%
\pgfpathlineto{\pgfqpoint{2.521153in}{0.869236in}}%
\pgfpathlineto{\pgfqpoint{2.531149in}{0.873983in}}%
\pgfpathlineto{\pgfqpoint{2.537743in}{0.887231in}}%
\pgfpathlineto{\pgfqpoint{2.546436in}{0.890784in}}%
\pgfpathlineto{\pgfqpoint{2.568098in}{0.900479in}}%
\pgfpathlineto{\pgfqpoint{2.571719in}{0.903943in}}%
\pgfpathlineto{\pgfqpoint{2.581438in}{0.913727in}}%
\pgfpathlineto{\pgfqpoint{2.597002in}{0.921264in}}%
\pgfpathlineto{\pgfqpoint{2.607300in}{0.926975in}}%
\pgfpathlineto{\pgfqpoint{2.622285in}{0.936058in}}%
\pgfpathlineto{\pgfqpoint{2.628914in}{0.940223in}}%
\pgfpathlineto{\pgfqpoint{2.645105in}{0.953471in}}%
\pgfpathlineto{\pgfqpoint{2.647569in}{0.955559in}}%
\pgfpathlineto{\pgfqpoint{2.660066in}{0.966719in}}%
\pgfpathlineto{\pgfqpoint{2.669618in}{0.979967in}}%
\pgfpathlineto{\pgfqpoint{2.672852in}{0.982549in}}%
\pgfpathlineto{\pgfqpoint{2.685737in}{0.993215in}}%
\pgfpathlineto{\pgfqpoint{2.696134in}{1.006463in}}%
\pgfpathlineto{\pgfqpoint{2.698135in}{1.007461in}}%
\pgfpathlineto{\pgfqpoint{2.721768in}{1.019711in}}%
\pgfpathlineto{\pgfqpoint{2.723418in}{1.022717in}}%
\pgfpathlineto{\pgfqpoint{2.727939in}{1.032959in}}%
\pgfpathlineto{\pgfqpoint{2.748272in}{1.046207in}}%
\pgfpathlineto{\pgfqpoint{2.748701in}{1.047655in}}%
\pgfpathlineto{\pgfqpoint{2.751503in}{1.059455in}}%
\pgfpathlineto{\pgfqpoint{2.769672in}{1.072703in}}%
\pgfpathlineto{\pgfqpoint{2.773984in}{1.084972in}}%
\pgfpathlineto{\pgfqpoint{2.774258in}{1.085951in}}%
\pgfpathlineto{\pgfqpoint{2.795226in}{1.099199in}}%
\pgfpathlineto{\pgfqpoint{2.799268in}{1.109721in}}%
\pgfpathlineto{\pgfqpoint{2.800097in}{1.112447in}}%
\pgfpathlineto{\pgfqpoint{2.818205in}{1.125695in}}%
\pgfpathlineto{\pgfqpoint{2.824551in}{1.138745in}}%
\pgfpathlineto{\pgfqpoint{2.824629in}{1.138943in}}%
\pgfpathlineto{\pgfqpoint{2.836129in}{1.152191in}}%
\pgfpathlineto{\pgfqpoint{2.839297in}{1.165439in}}%
\pgfpathlineto{\pgfqpoint{2.849834in}{1.176038in}}%
\pgfpathlineto{\pgfqpoint{2.851936in}{1.178687in}}%
\pgfpathlineto{\pgfqpoint{2.858646in}{1.191935in}}%
\pgfpathlineto{\pgfqpoint{2.871719in}{1.205183in}}%
\pgfpathlineto{\pgfqpoint{2.875117in}{1.208796in}}%
\pgfpathlineto{\pgfqpoint{2.882486in}{1.218431in}}%
\pgfpathlineto{\pgfqpoint{2.889367in}{1.231679in}}%
\pgfpathlineto{\pgfqpoint{2.899090in}{1.244927in}}%
\pgfpathlineto{\pgfqpoint{2.900400in}{1.247631in}}%
\pgfpathlineto{\pgfqpoint{2.904049in}{1.258175in}}%
\pgfpathlineto{\pgfqpoint{2.913292in}{1.271423in}}%
\pgfpathlineto{\pgfqpoint{2.918455in}{1.284671in}}%
\pgfpathlineto{\pgfqpoint{2.925683in}{1.291457in}}%
\pgfpathlineto{\pgfqpoint{2.931081in}{1.297919in}}%
\pgfpathlineto{\pgfqpoint{2.934296in}{1.311167in}}%
\pgfpathlineto{\pgfqpoint{2.946309in}{1.324415in}}%
\pgfpathlineto{\pgfqpoint{2.949543in}{1.337663in}}%
\pgfpathlineto{\pgfqpoint{2.950967in}{1.339453in}}%
\pgfpathlineto{\pgfqpoint{2.957787in}{1.350911in}}%
\pgfpathlineto{\pgfqpoint{2.959697in}{1.364159in}}%
\pgfpathlineto{\pgfqpoint{2.971487in}{1.377407in}}%
\pgfpathlineto{\pgfqpoint{2.976250in}{1.389375in}}%
\pgfpathlineto{\pgfqpoint{2.976592in}{1.390655in}}%
\pgfpathlineto{\pgfqpoint{2.985776in}{1.403903in}}%
\pgfpathlineto{\pgfqpoint{2.990026in}{1.417151in}}%
\pgfpathlineto{\pgfqpoint{2.997037in}{1.430399in}}%
\pgfpathlineto{\pgfqpoint{3.001533in}{1.443108in}}%
\pgfpathlineto{\pgfqpoint{3.001660in}{1.443647in}}%
\pgfpathlineto{\pgfqpoint{3.007074in}{1.456895in}}%
\pgfpathlineto{\pgfqpoint{3.011952in}{1.470143in}}%
\pgfpathlineto{\pgfqpoint{3.018151in}{1.483391in}}%
\pgfpathlineto{\pgfqpoint{3.026816in}{1.496238in}}%
\pgfpathlineto{\pgfqpoint{3.027001in}{1.496639in}}%
\pgfpathlineto{\pgfqpoint{3.030195in}{1.509887in}}%
\pgfpathlineto{\pgfqpoint{3.036372in}{1.523135in}}%
\pgfpathlineto{\pgfqpoint{3.039196in}{1.536383in}}%
\pgfpathlineto{\pgfqpoint{3.046300in}{1.549631in}}%
\pgfpathlineto{\pgfqpoint{3.049957in}{1.562879in}}%
\pgfpathlineto{\pgfqpoint{3.052099in}{1.565229in}}%
\pgfpathlineto{\pgfqpoint{3.058303in}{1.576127in}}%
\pgfpathlineto{\pgfqpoint{3.060680in}{1.589375in}}%
\pgfpathlineto{\pgfqpoint{3.068226in}{1.602623in}}%
\pgfpathlineto{\pgfqpoint{3.070004in}{1.615871in}}%
\pgfpathlineto{\pgfqpoint{3.077382in}{1.628880in}}%
\pgfpathlineto{\pgfqpoint{3.077467in}{1.629119in}}%
\pgfpathlineto{\pgfqpoint{3.079330in}{1.642367in}}%
\pgfpathlineto{\pgfqpoint{3.085391in}{1.655615in}}%
\pgfpathlineto{\pgfqpoint{3.088660in}{1.668863in}}%
\pgfpathlineto{\pgfqpoint{3.095435in}{1.682111in}}%
\pgfpathlineto{\pgfqpoint{3.101328in}{1.695359in}}%
\pgfpathlineto{\pgfqpoint{3.102666in}{1.698247in}}%
\pgfpathlineto{\pgfqpoint{3.105233in}{1.708607in}}%
\pgfpathlineto{\pgfqpoint{3.108362in}{1.721855in}}%
\pgfpathlineto{\pgfqpoint{3.111375in}{1.735103in}}%
\pgfpathlineto{\pgfqpoint{3.115815in}{1.748351in}}%
\pgfpathlineto{\pgfqpoint{3.120006in}{1.761599in}}%
\pgfpathlineto{\pgfqpoint{3.127949in}{1.772612in}}%
\pgfpathlineto{\pgfqpoint{3.128895in}{1.774847in}}%
\pgfpathlineto{\pgfqpoint{3.130628in}{1.788095in}}%
\pgfpathlineto{\pgfqpoint{3.135067in}{1.801343in}}%
\pgfpathlineto{\pgfqpoint{3.136463in}{1.814591in}}%
\pgfpathlineto{\pgfqpoint{3.142114in}{1.827839in}}%
\pgfpathlineto{\pgfqpoint{3.144924in}{1.841087in}}%
\pgfpathlineto{\pgfqpoint{3.153232in}{1.851835in}}%
\pgfpathlineto{\pgfqpoint{3.154320in}{1.854335in}}%
\pgfpathlineto{\pgfqpoint{3.156143in}{1.867583in}}%
\pgfpathlineto{\pgfqpoint{3.160898in}{1.880831in}}%
\pgfpathlineto{\pgfqpoint{3.163253in}{1.894079in}}%
\pgfpathlineto{\pgfqpoint{3.167472in}{1.907327in}}%
\pgfpathlineto{\pgfqpoint{3.170135in}{1.920575in}}%
\pgfpathlineto{\pgfqpoint{3.176963in}{1.933823in}}%
\pgfpathlineto{\pgfqpoint{3.178515in}{1.936771in}}%
\pgfpathlineto{\pgfqpoint{3.181161in}{1.947071in}}%
\pgfpathlineto{\pgfqpoint{3.184965in}{1.960319in}}%
\pgfpathlineto{\pgfqpoint{3.190102in}{1.973567in}}%
\pgfpathlineto{\pgfqpoint{3.192673in}{1.986815in}}%
\pgfpathlineto{\pgfqpoint{3.198031in}{2.000063in}}%
\pgfpathlineto{\pgfqpoint{3.201642in}{2.013311in}}%
\pgfpathlineto{\pgfqpoint{3.203798in}{2.016253in}}%
\pgfpathlineto{\pgfqpoint{3.207861in}{2.026559in}}%
\pgfpathlineto{\pgfqpoint{3.209991in}{2.039807in}}%
\pgfpathlineto{\pgfqpoint{3.216768in}{2.053055in}}%
\pgfpathlineto{\pgfqpoint{3.219051in}{2.066303in}}%
\pgfpathlineto{\pgfqpoint{3.226708in}{2.079551in}}%
\pgfpathlineto{\pgfqpoint{3.229082in}{2.088315in}}%
\pgfpathlineto{\pgfqpoint{3.229766in}{2.092799in}}%
\pgfpathlineto{\pgfqpoint{3.234934in}{2.106047in}}%
\pgfpathlineto{\pgfqpoint{3.237662in}{2.119295in}}%
\pgfpathlineto{\pgfqpoint{3.245277in}{2.132543in}}%
\pgfpathlineto{\pgfqpoint{3.229082in}{2.132543in}}%
\pgfpathlineto{\pgfqpoint{3.203798in}{2.132543in}}%
\pgfpathlineto{\pgfqpoint{3.187113in}{2.132543in}}%
\pgfpathlineto{\pgfqpoint{3.179741in}{2.119295in}}%
\pgfpathlineto{\pgfqpoint{3.178515in}{2.115046in}}%
\pgfpathlineto{\pgfqpoint{3.177408in}{2.106047in}}%
\pgfpathlineto{\pgfqpoint{3.171919in}{2.092799in}}%
\pgfpathlineto{\pgfqpoint{3.169635in}{2.079551in}}%
\pgfpathlineto{\pgfqpoint{3.162907in}{2.066303in}}%
\pgfpathlineto{\pgfqpoint{3.160420in}{2.053055in}}%
\pgfpathlineto{\pgfqpoint{3.153442in}{2.039807in}}%
\pgfpathlineto{\pgfqpoint{3.153232in}{2.039134in}}%
\pgfpathlineto{\pgfqpoint{3.151418in}{2.026559in}}%
\pgfpathlineto{\pgfqpoint{3.146649in}{2.013311in}}%
\pgfpathlineto{\pgfqpoint{3.143138in}{2.000063in}}%
\pgfpathlineto{\pgfqpoint{3.137405in}{1.986815in}}%
\pgfpathlineto{\pgfqpoint{3.130709in}{1.973567in}}%
\pgfpathlineto{\pgfqpoint{3.127949in}{1.967065in}}%
\pgfpathlineto{\pgfqpoint{3.126419in}{1.960319in}}%
\pgfpathlineto{\pgfqpoint{3.122574in}{1.947071in}}%
\pgfpathlineto{\pgfqpoint{3.120247in}{1.933823in}}%
\pgfpathlineto{\pgfqpoint{3.115200in}{1.920575in}}%
\pgfpathlineto{\pgfqpoint{3.111689in}{1.907327in}}%
\pgfpathlineto{\pgfqpoint{3.102666in}{1.896167in}}%
\pgfpathlineto{\pgfqpoint{3.101539in}{1.894079in}}%
\pgfpathlineto{\pgfqpoint{3.099390in}{1.880831in}}%
\pgfpathlineto{\pgfqpoint{3.093342in}{1.867583in}}%
\pgfpathlineto{\pgfqpoint{3.091453in}{1.854335in}}%
\pgfpathlineto{\pgfqpoint{3.084491in}{1.841087in}}%
\pgfpathlineto{\pgfqpoint{3.081306in}{1.827839in}}%
\pgfpathlineto{\pgfqpoint{3.077382in}{1.823305in}}%
\pgfpathlineto{\pgfqpoint{3.072603in}{1.814591in}}%
\pgfpathlineto{\pgfqpoint{3.069499in}{1.801343in}}%
\pgfpathlineto{\pgfqpoint{3.062644in}{1.788095in}}%
\pgfpathlineto{\pgfqpoint{3.057803in}{1.774847in}}%
\pgfpathlineto{\pgfqpoint{3.052099in}{1.763482in}}%
\pgfpathlineto{\pgfqpoint{3.051449in}{1.761599in}}%
\pgfpathlineto{\pgfqpoint{3.048285in}{1.748351in}}%
\pgfpathlineto{\pgfqpoint{3.043434in}{1.735103in}}%
\pgfpathlineto{\pgfqpoint{3.036710in}{1.721855in}}%
\pgfpathlineto{\pgfqpoint{3.030262in}{1.708607in}}%
\pgfpathlineto{\pgfqpoint{3.026816in}{1.704506in}}%
\pgfpathlineto{\pgfqpoint{3.020956in}{1.695359in}}%
\pgfpathlineto{\pgfqpoint{3.018054in}{1.682111in}}%
\pgfpathlineto{\pgfqpoint{3.010637in}{1.668863in}}%
\pgfpathlineto{\pgfqpoint{3.007546in}{1.655615in}}%
\pgfpathlineto{\pgfqpoint{3.001533in}{1.648241in}}%
\pgfpathlineto{\pgfqpoint{2.997596in}{1.642367in}}%
\pgfpathlineto{\pgfqpoint{2.995131in}{1.629119in}}%
\pgfpathlineto{\pgfqpoint{2.983780in}{1.615871in}}%
\pgfpathlineto{\pgfqpoint{2.980127in}{1.602623in}}%
\pgfpathlineto{\pgfqpoint{2.976250in}{1.598029in}}%
\pgfpathlineto{\pgfqpoint{2.970145in}{1.589375in}}%
\pgfpathlineto{\pgfqpoint{2.966734in}{1.576127in}}%
\pgfpathlineto{\pgfqpoint{2.957449in}{1.562879in}}%
\pgfpathlineto{\pgfqpoint{2.953235in}{1.549631in}}%
\pgfpathlineto{\pgfqpoint{2.950967in}{1.547273in}}%
\pgfpathlineto{\pgfqpoint{2.942483in}{1.536383in}}%
\pgfpathlineto{\pgfqpoint{2.934316in}{1.523135in}}%
\pgfpathlineto{\pgfqpoint{2.925683in}{1.511548in}}%
\pgfpathlineto{\pgfqpoint{2.924624in}{1.509887in}}%
\pgfpathlineto{\pgfqpoint{2.916598in}{1.496639in}}%
\pgfpathlineto{\pgfqpoint{2.911582in}{1.483391in}}%
\pgfpathlineto{\pgfqpoint{2.902987in}{1.470143in}}%
\pgfpathlineto{\pgfqpoint{2.900400in}{1.465481in}}%
\pgfpathlineto{\pgfqpoint{2.896492in}{1.456895in}}%
\pgfpathlineto{\pgfqpoint{2.883200in}{1.443647in}}%
\pgfpathlineto{\pgfqpoint{2.877527in}{1.430399in}}%
\pgfpathlineto{\pgfqpoint{2.875117in}{1.428437in}}%
\pgfpathlineto{\pgfqpoint{2.861485in}{1.417151in}}%
\pgfpathlineto{\pgfqpoint{2.856992in}{1.403903in}}%
\pgfpathlineto{\pgfqpoint{2.849834in}{1.397597in}}%
\pgfpathlineto{\pgfqpoint{2.841679in}{1.390655in}}%
\pgfpathlineto{\pgfqpoint{2.838908in}{1.377407in}}%
\pgfpathlineto{\pgfqpoint{2.824551in}{1.366117in}}%
\pgfpathlineto{\pgfqpoint{2.821981in}{1.364159in}}%
\pgfpathlineto{\pgfqpoint{2.818006in}{1.350911in}}%
\pgfpathlineto{\pgfqpoint{2.799268in}{1.338859in}}%
\pgfpathlineto{\pgfqpoint{2.797363in}{1.337663in}}%
\pgfpathlineto{\pgfqpoint{2.789546in}{1.324415in}}%
\pgfpathlineto{\pgfqpoint{2.773984in}{1.311321in}}%
\pgfpathlineto{\pgfqpoint{2.773800in}{1.311167in}}%
\pgfpathlineto{\pgfqpoint{2.767467in}{1.297919in}}%
\pgfpathlineto{\pgfqpoint{2.754210in}{1.284671in}}%
\pgfpathlineto{\pgfqpoint{2.748701in}{1.278084in}}%
\pgfpathlineto{\pgfqpoint{2.742656in}{1.271423in}}%
\pgfpathlineto{\pgfqpoint{2.725649in}{1.258175in}}%
\pgfpathlineto{\pgfqpoint{2.723418in}{1.256607in}}%
\pgfpathlineto{\pgfqpoint{2.704912in}{1.244927in}}%
\pgfpathlineto{\pgfqpoint{2.698135in}{1.237451in}}%
\pgfpathlineto{\pgfqpoint{2.693140in}{1.231679in}}%
\pgfpathlineto{\pgfqpoint{2.672852in}{1.219884in}}%
\pgfpathlineto{\pgfqpoint{2.669916in}{1.218431in}}%
\pgfpathlineto{\pgfqpoint{2.662841in}{1.205183in}}%
\pgfpathlineto{\pgfqpoint{2.647569in}{1.195997in}}%
\pgfpathlineto{\pgfqpoint{2.639598in}{1.191935in}}%
\pgfpathlineto{\pgfqpoint{2.635937in}{1.178687in}}%
\pgfpathlineto{\pgfqpoint{2.622285in}{1.172697in}}%
\pgfpathlineto{\pgfqpoint{2.602872in}{1.165439in}}%
\pgfpathlineto{\pgfqpoint{2.597002in}{1.155226in}}%
\pgfpathlineto{\pgfqpoint{2.595241in}{1.152191in}}%
\pgfpathlineto{\pgfqpoint{2.571719in}{1.144563in}}%
\pgfpathlineto{\pgfqpoint{2.551836in}{1.138943in}}%
\pgfpathlineto{\pgfqpoint{2.546436in}{1.126319in}}%
\pgfpathlineto{\pgfqpoint{2.546156in}{1.125695in}}%
\pgfpathlineto{\pgfqpoint{2.521153in}{1.116219in}}%
\pgfpathlineto{\pgfqpoint{2.510001in}{1.112447in}}%
\pgfpathlineto{\pgfqpoint{2.507477in}{1.099199in}}%
\pgfpathlineto{\pgfqpoint{2.495870in}{1.095398in}}%
\pgfpathlineto{\pgfqpoint{2.470586in}{1.087826in}}%
\pgfpathlineto{\pgfqpoint{2.470586in}{1.085951in}}%
\pgfpathlineto{\pgfqpoint{2.470586in}{1.072703in}}%
\pgfpathlineto{\pgfqpoint{2.470586in}{1.059455in}}%
\pgfpathlineto{\pgfqpoint{2.470586in}{1.046207in}}%
\pgfpathlineto{\pgfqpoint{2.470586in}{1.032959in}}%
\pgfpathlineto{\pgfqpoint{2.470586in}{1.019711in}}%
\pgfpathlineto{\pgfqpoint{2.470586in}{1.006463in}}%
\pgfpathlineto{\pgfqpoint{2.470586in}{0.993215in}}%
\pgfpathlineto{\pgfqpoint{2.470586in}{0.979967in}}%
\pgfpathlineto{\pgfqpoint{2.470586in}{0.966719in}}%
\pgfpathlineto{\pgfqpoint{2.470586in}{0.953471in}}%
\pgfpathlineto{\pgfqpoint{2.470586in}{0.940223in}}%
\pgfpathlineto{\pgfqpoint{2.470586in}{0.926975in}}%
\pgfpathlineto{\pgfqpoint{2.470586in}{0.913727in}}%
\pgfpathlineto{\pgfqpoint{2.470586in}{0.900479in}}%
\pgfpathlineto{\pgfqpoint{2.470586in}{0.887231in}}%
\pgfpathlineto{\pgfqpoint{2.470586in}{0.873983in}}%
\pgfpathlineto{\pgfqpoint{2.470586in}{0.860735in}}%
\pgfpathlineto{\pgfqpoint{2.470586in}{0.847487in}}%
\pgfpathlineto{\pgfqpoint{2.470586in}{0.839033in}}%
\pgfpathlineto{\pgfqpoint{2.495870in}{0.847100in}}%
\pgfpathclose%
\pgfpathmoveto{\pgfqpoint{4.290975in}{1.224152in}}%
\pgfpathlineto{\pgfqpoint{4.290975in}{1.231679in}}%
\pgfpathlineto{\pgfqpoint{4.290975in}{1.244927in}}%
\pgfpathlineto{\pgfqpoint{4.290975in}{1.258175in}}%
\pgfpathlineto{\pgfqpoint{4.290975in}{1.271423in}}%
\pgfpathlineto{\pgfqpoint{4.290975in}{1.284671in}}%
\pgfpathlineto{\pgfqpoint{4.290975in}{1.297919in}}%
\pgfpathlineto{\pgfqpoint{4.290975in}{1.311167in}}%
\pgfpathlineto{\pgfqpoint{4.290975in}{1.324415in}}%
\pgfpathlineto{\pgfqpoint{4.290975in}{1.337663in}}%
\pgfpathlineto{\pgfqpoint{4.290975in}{1.350911in}}%
\pgfpathlineto{\pgfqpoint{4.290975in}{1.364159in}}%
\pgfpathlineto{\pgfqpoint{4.290975in}{1.377407in}}%
\pgfpathlineto{\pgfqpoint{4.290975in}{1.390655in}}%
\pgfpathlineto{\pgfqpoint{4.290975in}{1.403903in}}%
\pgfpathlineto{\pgfqpoint{4.290975in}{1.417151in}}%
\pgfpathlineto{\pgfqpoint{4.290975in}{1.430399in}}%
\pgfpathlineto{\pgfqpoint{4.290975in}{1.443647in}}%
\pgfpathlineto{\pgfqpoint{4.290975in}{1.456895in}}%
\pgfpathlineto{\pgfqpoint{4.290975in}{1.470143in}}%
\pgfpathlineto{\pgfqpoint{4.290975in}{1.483391in}}%
\pgfpathlineto{\pgfqpoint{4.290975in}{1.496639in}}%
\pgfpathlineto{\pgfqpoint{4.290975in}{1.505144in}}%
\pgfpathlineto{\pgfqpoint{4.284667in}{1.509887in}}%
\pgfpathlineto{\pgfqpoint{4.265692in}{1.519832in}}%
\pgfpathlineto{\pgfqpoint{4.260665in}{1.523135in}}%
\pgfpathlineto{\pgfqpoint{4.245883in}{1.536383in}}%
\pgfpathlineto{\pgfqpoint{4.240408in}{1.539704in}}%
\pgfpathlineto{\pgfqpoint{4.226688in}{1.549631in}}%
\pgfpathlineto{\pgfqpoint{4.219486in}{1.562879in}}%
\pgfpathlineto{\pgfqpoint{4.215125in}{1.566090in}}%
\pgfpathlineto{\pgfqpoint{4.204073in}{1.576127in}}%
\pgfpathlineto{\pgfqpoint{4.197270in}{1.589375in}}%
\pgfpathlineto{\pgfqpoint{4.189842in}{1.593245in}}%
\pgfpathlineto{\pgfqpoint{4.175335in}{1.602623in}}%
\pgfpathlineto{\pgfqpoint{4.169419in}{1.615871in}}%
\pgfpathlineto{\pgfqpoint{4.164559in}{1.618570in}}%
\pgfpathlineto{\pgfqpoint{4.148305in}{1.629119in}}%
\pgfpathlineto{\pgfqpoint{4.143799in}{1.642367in}}%
\pgfpathlineto{\pgfqpoint{4.139276in}{1.645663in}}%
\pgfpathlineto{\pgfqpoint{4.127488in}{1.655615in}}%
\pgfpathlineto{\pgfqpoint{4.123876in}{1.668863in}}%
\pgfpathlineto{\pgfqpoint{4.113993in}{1.676230in}}%
\pgfpathlineto{\pgfqpoint{4.107056in}{1.682111in}}%
\pgfpathlineto{\pgfqpoint{4.102901in}{1.695359in}}%
\pgfpathlineto{\pgfqpoint{4.088709in}{1.705150in}}%
\pgfpathlineto{\pgfqpoint{4.083994in}{1.708607in}}%
\pgfpathlineto{\pgfqpoint{4.076425in}{1.721855in}}%
\pgfpathlineto{\pgfqpoint{4.063426in}{1.734724in}}%
\pgfpathlineto{\pgfqpoint{4.063045in}{1.735103in}}%
\pgfpathlineto{\pgfqpoint{4.057348in}{1.748351in}}%
\pgfpathlineto{\pgfqpoint{4.047828in}{1.761599in}}%
\pgfpathlineto{\pgfqpoint{4.041870in}{1.774847in}}%
\pgfpathlineto{\pgfqpoint{4.038143in}{1.779081in}}%
\pgfpathlineto{\pgfqpoint{4.030319in}{1.788095in}}%
\pgfpathlineto{\pgfqpoint{4.020531in}{1.801343in}}%
\pgfpathlineto{\pgfqpoint{4.012860in}{1.813727in}}%
\pgfpathlineto{\pgfqpoint{4.012285in}{1.814591in}}%
\pgfpathlineto{\pgfqpoint{4.001946in}{1.827839in}}%
\pgfpathlineto{\pgfqpoint{3.996436in}{1.841087in}}%
\pgfpathlineto{\pgfqpoint{3.987636in}{1.854335in}}%
\pgfpathlineto{\pgfqpoint{3.987577in}{1.854480in}}%
\pgfpathlineto{\pgfqpoint{3.981862in}{1.867583in}}%
\pgfpathlineto{\pgfqpoint{3.970916in}{1.880831in}}%
\pgfpathlineto{\pgfqpoint{3.966475in}{1.894079in}}%
\pgfpathlineto{\pgfqpoint{3.962293in}{1.898277in}}%
\pgfpathlineto{\pgfqpoint{3.953136in}{1.907327in}}%
\pgfpathlineto{\pgfqpoint{3.949328in}{1.920575in}}%
\pgfpathlineto{\pgfqpoint{3.940187in}{1.933823in}}%
\pgfpathlineto{\pgfqpoint{3.938303in}{1.947071in}}%
\pgfpathlineto{\pgfqpoint{3.937010in}{1.948772in}}%
\pgfpathlineto{\pgfqpoint{3.928152in}{1.960319in}}%
\pgfpathlineto{\pgfqpoint{3.924845in}{1.973567in}}%
\pgfpathlineto{\pgfqpoint{3.913687in}{1.986815in}}%
\pgfpathlineto{\pgfqpoint{3.911727in}{1.992957in}}%
\pgfpathlineto{\pgfqpoint{3.909007in}{2.000063in}}%
\pgfpathlineto{\pgfqpoint{3.899767in}{2.013311in}}%
\pgfpathlineto{\pgfqpoint{3.895924in}{2.026559in}}%
\pgfpathlineto{\pgfqpoint{3.889083in}{2.039807in}}%
\pgfpathlineto{\pgfqpoint{3.886444in}{2.048919in}}%
\pgfpathlineto{\pgfqpoint{3.884945in}{2.053055in}}%
\pgfpathlineto{\pgfqpoint{3.877017in}{2.066303in}}%
\pgfpathlineto{\pgfqpoint{3.871441in}{2.079551in}}%
\pgfpathlineto{\pgfqpoint{3.865721in}{2.092799in}}%
\pgfpathlineto{\pgfqpoint{3.861161in}{2.102371in}}%
\pgfpathlineto{\pgfqpoint{3.858856in}{2.106047in}}%
\pgfpathlineto{\pgfqpoint{3.854166in}{2.119295in}}%
\pgfpathlineto{\pgfqpoint{3.848450in}{2.132543in}}%
\pgfpathlineto{\pgfqpoint{3.835878in}{2.132543in}}%
\pgfpathlineto{\pgfqpoint{3.810594in}{2.132543in}}%
\pgfpathlineto{\pgfqpoint{3.785311in}{2.132543in}}%
\pgfpathlineto{\pgfqpoint{3.777992in}{2.132543in}}%
\pgfpathlineto{\pgfqpoint{3.779757in}{2.119295in}}%
\pgfpathlineto{\pgfqpoint{3.785241in}{2.106047in}}%
\pgfpathlineto{\pgfqpoint{3.785311in}{2.105411in}}%
\pgfpathlineto{\pgfqpoint{3.787916in}{2.092799in}}%
\pgfpathlineto{\pgfqpoint{3.794444in}{2.079551in}}%
\pgfpathlineto{\pgfqpoint{3.796181in}{2.066303in}}%
\pgfpathlineto{\pgfqpoint{3.802740in}{2.053055in}}%
\pgfpathlineto{\pgfqpoint{3.805272in}{2.039807in}}%
\pgfpathlineto{\pgfqpoint{3.810594in}{2.027642in}}%
\pgfpathlineto{\pgfqpoint{3.811407in}{2.026559in}}%
\pgfpathlineto{\pgfqpoint{3.815991in}{2.013311in}}%
\pgfpathlineto{\pgfqpoint{3.821760in}{2.000063in}}%
\pgfpathlineto{\pgfqpoint{3.825110in}{1.986815in}}%
\pgfpathlineto{\pgfqpoint{3.829372in}{1.973567in}}%
\pgfpathlineto{\pgfqpoint{3.832651in}{1.960319in}}%
\pgfpathlineto{\pgfqpoint{3.835878in}{1.950115in}}%
\pgfpathlineto{\pgfqpoint{3.837426in}{1.947071in}}%
\pgfpathlineto{\pgfqpoint{3.845200in}{1.933823in}}%
\pgfpathlineto{\pgfqpoint{3.849612in}{1.920575in}}%
\pgfpathlineto{\pgfqpoint{3.855781in}{1.907327in}}%
\pgfpathlineto{\pgfqpoint{3.858191in}{1.894079in}}%
\pgfpathlineto{\pgfqpoint{3.861161in}{1.886538in}}%
\pgfpathlineto{\pgfqpoint{3.864412in}{1.880831in}}%
\pgfpathlineto{\pgfqpoint{3.868088in}{1.867583in}}%
\pgfpathlineto{\pgfqpoint{3.877211in}{1.854335in}}%
\pgfpathlineto{\pgfqpoint{3.880138in}{1.841087in}}%
\pgfpathlineto{\pgfqpoint{3.886444in}{1.830672in}}%
\pgfpathlineto{\pgfqpoint{3.888884in}{1.827839in}}%
\pgfpathlineto{\pgfqpoint{3.891459in}{1.814591in}}%
\pgfpathlineto{\pgfqpoint{3.899577in}{1.801343in}}%
\pgfpathlineto{\pgfqpoint{3.901769in}{1.788095in}}%
\pgfpathlineto{\pgfqpoint{3.909735in}{1.774847in}}%
\pgfpathlineto{\pgfqpoint{3.911727in}{1.765308in}}%
\pgfpathlineto{\pgfqpoint{3.912829in}{1.761599in}}%
\pgfpathlineto{\pgfqpoint{3.923886in}{1.748351in}}%
\pgfpathlineto{\pgfqpoint{3.928952in}{1.735103in}}%
\pgfpathlineto{\pgfqpoint{3.936065in}{1.721855in}}%
\pgfpathlineto{\pgfqpoint{3.937010in}{1.718836in}}%
\pgfpathlineto{\pgfqpoint{3.941308in}{1.708607in}}%
\pgfpathlineto{\pgfqpoint{3.948118in}{1.695359in}}%
\pgfpathlineto{\pgfqpoint{3.953337in}{1.682111in}}%
\pgfpathlineto{\pgfqpoint{3.960390in}{1.668863in}}%
\pgfpathlineto{\pgfqpoint{3.962293in}{1.665873in}}%
\pgfpathlineto{\pgfqpoint{3.970263in}{1.655615in}}%
\pgfpathlineto{\pgfqpoint{3.975985in}{1.642367in}}%
\pgfpathlineto{\pgfqpoint{3.984594in}{1.629119in}}%
\pgfpathlineto{\pgfqpoint{3.987577in}{1.617445in}}%
\pgfpathlineto{\pgfqpoint{3.988094in}{1.615871in}}%
\pgfpathlineto{\pgfqpoint{3.998179in}{1.602623in}}%
\pgfpathlineto{\pgfqpoint{4.002449in}{1.589375in}}%
\pgfpathlineto{\pgfqpoint{4.012860in}{1.577934in}}%
\pgfpathlineto{\pgfqpoint{4.014767in}{1.576127in}}%
\pgfpathlineto{\pgfqpoint{4.019329in}{1.562879in}}%
\pgfpathlineto{\pgfqpoint{4.033633in}{1.549631in}}%
\pgfpathlineto{\pgfqpoint{4.037468in}{1.536383in}}%
\pgfpathlineto{\pgfqpoint{4.038143in}{1.535539in}}%
\pgfpathlineto{\pgfqpoint{4.049252in}{1.523135in}}%
\pgfpathlineto{\pgfqpoint{4.051465in}{1.509887in}}%
\pgfpathlineto{\pgfqpoint{4.063426in}{1.496705in}}%
\pgfpathlineto{\pgfqpoint{4.063493in}{1.496639in}}%
\pgfpathlineto{\pgfqpoint{4.068425in}{1.483391in}}%
\pgfpathlineto{\pgfqpoint{4.083769in}{1.470143in}}%
\pgfpathlineto{\pgfqpoint{4.088709in}{1.461374in}}%
\pgfpathlineto{\pgfqpoint{4.091626in}{1.456895in}}%
\pgfpathlineto{\pgfqpoint{4.103057in}{1.443647in}}%
\pgfpathlineto{\pgfqpoint{4.109666in}{1.430399in}}%
\pgfpathlineto{\pgfqpoint{4.113993in}{1.424526in}}%
\pgfpathlineto{\pgfqpoint{4.120046in}{1.417151in}}%
\pgfpathlineto{\pgfqpoint{4.129357in}{1.403903in}}%
\pgfpathlineto{\pgfqpoint{4.139276in}{1.390858in}}%
\pgfpathlineto{\pgfqpoint{4.139455in}{1.390655in}}%
\pgfpathlineto{\pgfqpoint{4.154600in}{1.377407in}}%
\pgfpathlineto{\pgfqpoint{4.162302in}{1.364159in}}%
\pgfpathlineto{\pgfqpoint{4.164559in}{1.362220in}}%
\pgfpathlineto{\pgfqpoint{4.178395in}{1.350911in}}%
\pgfpathlineto{\pgfqpoint{4.183438in}{1.337663in}}%
\pgfpathlineto{\pgfqpoint{4.189842in}{1.331466in}}%
\pgfpathlineto{\pgfqpoint{4.197342in}{1.324415in}}%
\pgfpathlineto{\pgfqpoint{4.200026in}{1.311167in}}%
\pgfpathlineto{\pgfqpoint{4.215125in}{1.302241in}}%
\pgfpathlineto{\pgfqpoint{4.222459in}{1.297919in}}%
\pgfpathlineto{\pgfqpoint{4.228806in}{1.284671in}}%
\pgfpathlineto{\pgfqpoint{4.240408in}{1.278323in}}%
\pgfpathlineto{\pgfqpoint{4.253450in}{1.271423in}}%
\pgfpathlineto{\pgfqpoint{4.258476in}{1.258175in}}%
\pgfpathlineto{\pgfqpoint{4.265692in}{1.252948in}}%
\pgfpathlineto{\pgfqpoint{4.277002in}{1.244927in}}%
\pgfpathlineto{\pgfqpoint{4.278162in}{1.231679in}}%
\pgfpathlineto{\pgfqpoint{4.290975in}{1.224152in}}%
\pgfpathclose%
\pgfpathmoveto{\pgfqpoint{2.472947in}{1.324415in}}%
\pgfpathlineto{\pgfqpoint{2.475033in}{1.337663in}}%
\pgfpathlineto{\pgfqpoint{2.492138in}{1.350911in}}%
\pgfpathlineto{\pgfqpoint{2.493011in}{1.364159in}}%
\pgfpathlineto{\pgfqpoint{2.495870in}{1.365887in}}%
\pgfpathlineto{\pgfqpoint{2.513335in}{1.377407in}}%
\pgfpathlineto{\pgfqpoint{2.521153in}{1.389261in}}%
\pgfpathlineto{\pgfqpoint{2.522016in}{1.390655in}}%
\pgfpathlineto{\pgfqpoint{2.539097in}{1.403903in}}%
\pgfpathlineto{\pgfqpoint{2.546436in}{1.414339in}}%
\pgfpathlineto{\pgfqpoint{2.548260in}{1.417151in}}%
\pgfpathlineto{\pgfqpoint{2.557239in}{1.430399in}}%
\pgfpathlineto{\pgfqpoint{2.564507in}{1.443647in}}%
\pgfpathlineto{\pgfqpoint{2.571719in}{1.453718in}}%
\pgfpathlineto{\pgfqpoint{2.573578in}{1.456895in}}%
\pgfpathlineto{\pgfqpoint{2.583871in}{1.470143in}}%
\pgfpathlineto{\pgfqpoint{2.591578in}{1.483391in}}%
\pgfpathlineto{\pgfqpoint{2.597002in}{1.488102in}}%
\pgfpathlineto{\pgfqpoint{2.606172in}{1.496639in}}%
\pgfpathlineto{\pgfqpoint{2.611305in}{1.509887in}}%
\pgfpathlineto{\pgfqpoint{2.622285in}{1.520002in}}%
\pgfpathlineto{\pgfqpoint{2.625483in}{1.523135in}}%
\pgfpathlineto{\pgfqpoint{2.628126in}{1.536383in}}%
\pgfpathlineto{\pgfqpoint{2.640075in}{1.549631in}}%
\pgfpathlineto{\pgfqpoint{2.642761in}{1.562879in}}%
\pgfpathlineto{\pgfqpoint{2.647569in}{1.566470in}}%
\pgfpathlineto{\pgfqpoint{2.658978in}{1.576127in}}%
\pgfpathlineto{\pgfqpoint{2.662421in}{1.589375in}}%
\pgfpathlineto{\pgfqpoint{2.672852in}{1.599108in}}%
\pgfpathlineto{\pgfqpoint{2.676059in}{1.602623in}}%
\pgfpathlineto{\pgfqpoint{2.677637in}{1.615871in}}%
\pgfpathlineto{\pgfqpoint{2.686763in}{1.629119in}}%
\pgfpathlineto{\pgfqpoint{2.689745in}{1.642367in}}%
\pgfpathlineto{\pgfqpoint{2.698135in}{1.650971in}}%
\pgfpathlineto{\pgfqpoint{2.701757in}{1.655615in}}%
\pgfpathlineto{\pgfqpoint{2.706452in}{1.668863in}}%
\pgfpathlineto{\pgfqpoint{2.716398in}{1.682111in}}%
\pgfpathlineto{\pgfqpoint{2.723418in}{1.693636in}}%
\pgfpathlineto{\pgfqpoint{2.724276in}{1.695359in}}%
\pgfpathlineto{\pgfqpoint{2.730389in}{1.708607in}}%
\pgfpathlineto{\pgfqpoint{2.736193in}{1.721855in}}%
\pgfpathlineto{\pgfqpoint{2.741162in}{1.735103in}}%
\pgfpathlineto{\pgfqpoint{2.748701in}{1.747768in}}%
\pgfpathlineto{\pgfqpoint{2.748979in}{1.748351in}}%
\pgfpathlineto{\pgfqpoint{2.753349in}{1.761599in}}%
\pgfpathlineto{\pgfqpoint{2.763317in}{1.774847in}}%
\pgfpathlineto{\pgfqpoint{2.766444in}{1.788095in}}%
\pgfpathlineto{\pgfqpoint{2.773984in}{1.798488in}}%
\pgfpathlineto{\pgfqpoint{2.775602in}{1.801343in}}%
\pgfpathlineto{\pgfqpoint{2.777096in}{1.814591in}}%
\pgfpathlineto{\pgfqpoint{2.784881in}{1.827839in}}%
\pgfpathlineto{\pgfqpoint{2.787585in}{1.841087in}}%
\pgfpathlineto{\pgfqpoint{2.799268in}{1.853997in}}%
\pgfpathlineto{\pgfqpoint{2.799500in}{1.854335in}}%
\pgfpathlineto{\pgfqpoint{2.801672in}{1.867583in}}%
\pgfpathlineto{\pgfqpoint{2.809394in}{1.880831in}}%
\pgfpathlineto{\pgfqpoint{2.812007in}{1.894079in}}%
\pgfpathlineto{\pgfqpoint{2.818917in}{1.907327in}}%
\pgfpathlineto{\pgfqpoint{2.821499in}{1.920575in}}%
\pgfpathlineto{\pgfqpoint{2.824551in}{1.925294in}}%
\pgfpathlineto{\pgfqpoint{2.828021in}{1.933823in}}%
\pgfpathlineto{\pgfqpoint{2.831598in}{1.947071in}}%
\pgfpathlineto{\pgfqpoint{2.837804in}{1.960319in}}%
\pgfpathlineto{\pgfqpoint{2.844290in}{1.973567in}}%
\pgfpathlineto{\pgfqpoint{2.849188in}{1.986815in}}%
\pgfpathlineto{\pgfqpoint{2.849834in}{1.988216in}}%
\pgfpathlineto{\pgfqpoint{2.853260in}{2.000063in}}%
\pgfpathlineto{\pgfqpoint{2.855822in}{2.013311in}}%
\pgfpathlineto{\pgfqpoint{2.861575in}{2.026559in}}%
\pgfpathlineto{\pgfqpoint{2.864810in}{2.039807in}}%
\pgfpathlineto{\pgfqpoint{2.873866in}{2.053055in}}%
\pgfpathlineto{\pgfqpoint{2.875117in}{2.059060in}}%
\pgfpathlineto{\pgfqpoint{2.875924in}{2.066303in}}%
\pgfpathlineto{\pgfqpoint{2.881105in}{2.079551in}}%
\pgfpathlineto{\pgfqpoint{2.882425in}{2.092799in}}%
\pgfpathlineto{\pgfqpoint{2.887734in}{2.106047in}}%
\pgfpathlineto{\pgfqpoint{2.889315in}{2.119295in}}%
\pgfpathlineto{\pgfqpoint{2.897696in}{2.132543in}}%
\pgfpathlineto{\pgfqpoint{2.875117in}{2.132543in}}%
\pgfpathlineto{\pgfqpoint{2.849834in}{2.132543in}}%
\pgfpathlineto{\pgfqpoint{2.825494in}{2.132543in}}%
\pgfpathlineto{\pgfqpoint{2.824551in}{2.131204in}}%
\pgfpathlineto{\pgfqpoint{2.818272in}{2.119295in}}%
\pgfpathlineto{\pgfqpoint{2.816182in}{2.106047in}}%
\pgfpathlineto{\pgfqpoint{2.806852in}{2.092799in}}%
\pgfpathlineto{\pgfqpoint{2.802080in}{2.079551in}}%
\pgfpathlineto{\pgfqpoint{2.799268in}{2.075246in}}%
\pgfpathlineto{\pgfqpoint{2.794708in}{2.066303in}}%
\pgfpathlineto{\pgfqpoint{2.790921in}{2.053055in}}%
\pgfpathlineto{\pgfqpoint{2.785251in}{2.039807in}}%
\pgfpathlineto{\pgfqpoint{2.779796in}{2.026559in}}%
\pgfpathlineto{\pgfqpoint{2.773984in}{2.016828in}}%
\pgfpathlineto{\pgfqpoint{2.772312in}{2.013311in}}%
\pgfpathlineto{\pgfqpoint{2.764285in}{2.000063in}}%
\pgfpathlineto{\pgfqpoint{2.759147in}{1.986815in}}%
\pgfpathlineto{\pgfqpoint{2.748701in}{1.973684in}}%
\pgfpathlineto{\pgfqpoint{2.748609in}{1.973567in}}%
\pgfpathlineto{\pgfqpoint{2.745406in}{1.960319in}}%
\pgfpathlineto{\pgfqpoint{2.737433in}{1.947071in}}%
\pgfpathlineto{\pgfqpoint{2.735568in}{1.933823in}}%
\pgfpathlineto{\pgfqpoint{2.723903in}{1.920575in}}%
\pgfpathlineto{\pgfqpoint{2.723418in}{1.918990in}}%
\pgfpathlineto{\pgfqpoint{2.720509in}{1.907327in}}%
\pgfpathlineto{\pgfqpoint{2.706600in}{1.894079in}}%
\pgfpathlineto{\pgfqpoint{2.703127in}{1.880831in}}%
\pgfpathlineto{\pgfqpoint{2.698135in}{1.875040in}}%
\pgfpathlineto{\pgfqpoint{2.691897in}{1.867583in}}%
\pgfpathlineto{\pgfqpoint{2.689249in}{1.854335in}}%
\pgfpathlineto{\pgfqpoint{2.677673in}{1.841087in}}%
\pgfpathlineto{\pgfqpoint{2.673039in}{1.827839in}}%
\pgfpathlineto{\pgfqpoint{2.672852in}{1.827660in}}%
\pgfpathlineto{\pgfqpoint{2.660863in}{1.814591in}}%
\pgfpathlineto{\pgfqpoint{2.653170in}{1.801343in}}%
\pgfpathlineto{\pgfqpoint{2.647569in}{1.793698in}}%
\pgfpathlineto{\pgfqpoint{2.643724in}{1.788095in}}%
\pgfpathlineto{\pgfqpoint{2.635132in}{1.774847in}}%
\pgfpathlineto{\pgfqpoint{2.628516in}{1.761599in}}%
\pgfpathlineto{\pgfqpoint{2.622285in}{1.752031in}}%
\pgfpathlineto{\pgfqpoint{2.619677in}{1.748351in}}%
\pgfpathlineto{\pgfqpoint{2.612939in}{1.735103in}}%
\pgfpathlineto{\pgfqpoint{2.597002in}{1.722314in}}%
\pgfpathlineto{\pgfqpoint{2.596322in}{1.721855in}}%
\pgfpathlineto{\pgfqpoint{2.589793in}{1.708607in}}%
\pgfpathlineto{\pgfqpoint{2.571719in}{1.695955in}}%
\pgfpathlineto{\pgfqpoint{2.570714in}{1.695359in}}%
\pgfpathlineto{\pgfqpoint{2.568807in}{1.682111in}}%
\pgfpathlineto{\pgfqpoint{2.553599in}{1.668863in}}%
\pgfpathlineto{\pgfqpoint{2.552089in}{1.655615in}}%
\pgfpathlineto{\pgfqpoint{2.546436in}{1.651721in}}%
\pgfpathlineto{\pgfqpoint{2.531407in}{1.642367in}}%
\pgfpathlineto{\pgfqpoint{2.528996in}{1.629119in}}%
\pgfpathlineto{\pgfqpoint{2.521153in}{1.624143in}}%
\pgfpathlineto{\pgfqpoint{2.507360in}{1.615871in}}%
\pgfpathlineto{\pgfqpoint{2.501320in}{1.602623in}}%
\pgfpathlineto{\pgfqpoint{2.495870in}{1.598464in}}%
\pgfpathlineto{\pgfqpoint{2.482952in}{1.589375in}}%
\pgfpathlineto{\pgfqpoint{2.477135in}{1.576127in}}%
\pgfpathlineto{\pgfqpoint{2.470586in}{1.570943in}}%
\pgfpathlineto{\pgfqpoint{2.470586in}{1.562879in}}%
\pgfpathlineto{\pgfqpoint{2.470586in}{1.549631in}}%
\pgfpathlineto{\pgfqpoint{2.470586in}{1.536383in}}%
\pgfpathlineto{\pgfqpoint{2.470586in}{1.523135in}}%
\pgfpathlineto{\pgfqpoint{2.470586in}{1.509887in}}%
\pgfpathlineto{\pgfqpoint{2.470586in}{1.496639in}}%
\pgfpathlineto{\pgfqpoint{2.470586in}{1.483391in}}%
\pgfpathlineto{\pgfqpoint{2.470586in}{1.470143in}}%
\pgfpathlineto{\pgfqpoint{2.470586in}{1.456895in}}%
\pgfpathlineto{\pgfqpoint{2.470586in}{1.443647in}}%
\pgfpathlineto{\pgfqpoint{2.470586in}{1.430399in}}%
\pgfpathlineto{\pgfqpoint{2.470586in}{1.417151in}}%
\pgfpathlineto{\pgfqpoint{2.470586in}{1.403903in}}%
\pgfpathlineto{\pgfqpoint{2.470586in}{1.390655in}}%
\pgfpathlineto{\pgfqpoint{2.470586in}{1.377407in}}%
\pgfpathlineto{\pgfqpoint{2.470586in}{1.364159in}}%
\pgfpathlineto{\pgfqpoint{2.470586in}{1.350911in}}%
\pgfpathlineto{\pgfqpoint{2.470586in}{1.337663in}}%
\pgfpathlineto{\pgfqpoint{2.470586in}{1.324415in}}%
\pgfpathlineto{\pgfqpoint{2.470586in}{1.323066in}}%
\pgfpathlineto{\pgfqpoint{2.472947in}{1.324415in}}%
\pgfpathclose%
\pgfpathmoveto{\pgfqpoint{4.290975in}{1.667261in}}%
\pgfpathlineto{\pgfqpoint{4.290975in}{1.668863in}}%
\pgfpathlineto{\pgfqpoint{4.290975in}{1.682111in}}%
\pgfpathlineto{\pgfqpoint{4.290975in}{1.695359in}}%
\pgfpathlineto{\pgfqpoint{4.290975in}{1.708607in}}%
\pgfpathlineto{\pgfqpoint{4.290975in}{1.721855in}}%
\pgfpathlineto{\pgfqpoint{4.290975in}{1.735103in}}%
\pgfpathlineto{\pgfqpoint{4.290975in}{1.748351in}}%
\pgfpathlineto{\pgfqpoint{4.290975in}{1.761599in}}%
\pgfpathlineto{\pgfqpoint{4.290975in}{1.774847in}}%
\pgfpathlineto{\pgfqpoint{4.290975in}{1.788095in}}%
\pgfpathlineto{\pgfqpoint{4.290975in}{1.801343in}}%
\pgfpathlineto{\pgfqpoint{4.290975in}{1.814591in}}%
\pgfpathlineto{\pgfqpoint{4.290975in}{1.827839in}}%
\pgfpathlineto{\pgfqpoint{4.290975in}{1.841087in}}%
\pgfpathlineto{\pgfqpoint{4.290975in}{1.854335in}}%
\pgfpathlineto{\pgfqpoint{4.290975in}{1.867583in}}%
\pgfpathlineto{\pgfqpoint{4.290975in}{1.880831in}}%
\pgfpathlineto{\pgfqpoint{4.290975in}{1.894079in}}%
\pgfpathlineto{\pgfqpoint{4.290975in}{1.907327in}}%
\pgfpathlineto{\pgfqpoint{4.290975in}{1.920575in}}%
\pgfpathlineto{\pgfqpoint{4.290975in}{1.933823in}}%
\pgfpathlineto{\pgfqpoint{4.290975in}{1.947071in}}%
\pgfpathlineto{\pgfqpoint{4.290975in}{1.948728in}}%
\pgfpathlineto{\pgfqpoint{4.276446in}{1.960319in}}%
\pgfpathlineto{\pgfqpoint{4.271710in}{1.973567in}}%
\pgfpathlineto{\pgfqpoint{4.265692in}{1.977165in}}%
\pgfpathlineto{\pgfqpoint{4.252883in}{1.986815in}}%
\pgfpathlineto{\pgfqpoint{4.245573in}{2.000063in}}%
\pgfpathlineto{\pgfqpoint{4.240408in}{2.004992in}}%
\pgfpathlineto{\pgfqpoint{4.233058in}{2.013311in}}%
\pgfpathlineto{\pgfqpoint{4.227802in}{2.026559in}}%
\pgfpathlineto{\pgfqpoint{4.216280in}{2.039807in}}%
\pgfpathlineto{\pgfqpoint{4.215125in}{2.041927in}}%
\pgfpathlineto{\pgfqpoint{4.210121in}{2.053055in}}%
\pgfpathlineto{\pgfqpoint{4.200604in}{2.066303in}}%
\pgfpathlineto{\pgfqpoint{4.190030in}{2.079551in}}%
\pgfpathlineto{\pgfqpoint{4.189842in}{2.079839in}}%
\pgfpathlineto{\pgfqpoint{4.182340in}{2.092799in}}%
\pgfpathlineto{\pgfqpoint{4.171951in}{2.106047in}}%
\pgfpathlineto{\pgfqpoint{4.166412in}{2.119295in}}%
\pgfpathlineto{\pgfqpoint{4.164559in}{2.121787in}}%
\pgfpathlineto{\pgfqpoint{4.157773in}{2.132543in}}%
\pgfpathlineto{\pgfqpoint{4.139276in}{2.132543in}}%
\pgfpathlineto{\pgfqpoint{4.113993in}{2.132543in}}%
\pgfpathlineto{\pgfqpoint{4.088709in}{2.132543in}}%
\pgfpathlineto{\pgfqpoint{4.063426in}{2.132543in}}%
\pgfpathlineto{\pgfqpoint{4.038143in}{2.132543in}}%
\pgfpathlineto{\pgfqpoint{4.037538in}{2.132543in}}%
\pgfpathlineto{\pgfqpoint{4.038143in}{2.129418in}}%
\pgfpathlineto{\pgfqpoint{4.040738in}{2.119295in}}%
\pgfpathlineto{\pgfqpoint{4.051123in}{2.106047in}}%
\pgfpathlineto{\pgfqpoint{4.053665in}{2.092799in}}%
\pgfpathlineto{\pgfqpoint{4.060532in}{2.079551in}}%
\pgfpathlineto{\pgfqpoint{4.062486in}{2.066303in}}%
\pgfpathlineto{\pgfqpoint{4.063426in}{2.064943in}}%
\pgfpathlineto{\pgfqpoint{4.072715in}{2.053055in}}%
\pgfpathlineto{\pgfqpoint{4.076666in}{2.039807in}}%
\pgfpathlineto{\pgfqpoint{4.085780in}{2.026559in}}%
\pgfpathlineto{\pgfqpoint{4.088709in}{2.018908in}}%
\pgfpathlineto{\pgfqpoint{4.091327in}{2.013311in}}%
\pgfpathlineto{\pgfqpoint{4.098769in}{2.000063in}}%
\pgfpathlineto{\pgfqpoint{4.103791in}{1.986815in}}%
\pgfpathlineto{\pgfqpoint{4.109575in}{1.973567in}}%
\pgfpathlineto{\pgfqpoint{4.113993in}{1.962490in}}%
\pgfpathlineto{\pgfqpoint{4.114992in}{1.960319in}}%
\pgfpathlineto{\pgfqpoint{4.121332in}{1.947071in}}%
\pgfpathlineto{\pgfqpoint{4.131526in}{1.933823in}}%
\pgfpathlineto{\pgfqpoint{4.136724in}{1.920575in}}%
\pgfpathlineto{\pgfqpoint{4.139276in}{1.917185in}}%
\pgfpathlineto{\pgfqpoint{4.147366in}{1.907327in}}%
\pgfpathlineto{\pgfqpoint{4.150576in}{1.894079in}}%
\pgfpathlineto{\pgfqpoint{4.159453in}{1.880831in}}%
\pgfpathlineto{\pgfqpoint{4.162389in}{1.867583in}}%
\pgfpathlineto{\pgfqpoint{4.164559in}{1.865512in}}%
\pgfpathlineto{\pgfqpoint{4.176360in}{1.854335in}}%
\pgfpathlineto{\pgfqpoint{4.179845in}{1.841087in}}%
\pgfpathlineto{\pgfqpoint{4.189842in}{1.831043in}}%
\pgfpathlineto{\pgfqpoint{4.193310in}{1.827839in}}%
\pgfpathlineto{\pgfqpoint{4.196306in}{1.814591in}}%
\pgfpathlineto{\pgfqpoint{4.207208in}{1.801343in}}%
\pgfpathlineto{\pgfqpoint{4.209377in}{1.788095in}}%
\pgfpathlineto{\pgfqpoint{4.215125in}{1.781887in}}%
\pgfpathlineto{\pgfqpoint{4.221790in}{1.774847in}}%
\pgfpathlineto{\pgfqpoint{4.226000in}{1.761599in}}%
\pgfpathlineto{\pgfqpoint{4.239446in}{1.748351in}}%
\pgfpathlineto{\pgfqpoint{4.240408in}{1.746973in}}%
\pgfpathlineto{\pgfqpoint{4.248746in}{1.735103in}}%
\pgfpathlineto{\pgfqpoint{4.258731in}{1.721855in}}%
\pgfpathlineto{\pgfqpoint{4.265692in}{1.708676in}}%
\pgfpathlineto{\pgfqpoint{4.265730in}{1.708607in}}%
\pgfpathlineto{\pgfqpoint{4.272678in}{1.695359in}}%
\pgfpathlineto{\pgfqpoint{4.281256in}{1.682111in}}%
\pgfpathlineto{\pgfqpoint{4.288884in}{1.668863in}}%
\pgfpathlineto{\pgfqpoint{4.290975in}{1.667261in}}%
\pgfpathclose%
\pgfpathmoveto{\pgfqpoint{2.477160in}{1.827839in}}%
\pgfpathlineto{\pgfqpoint{2.479175in}{1.841087in}}%
\pgfpathlineto{\pgfqpoint{2.495175in}{1.854335in}}%
\pgfpathlineto{\pgfqpoint{2.495870in}{1.856704in}}%
\pgfpathlineto{\pgfqpoint{2.498654in}{1.867583in}}%
\pgfpathlineto{\pgfqpoint{2.510833in}{1.880831in}}%
\pgfpathlineto{\pgfqpoint{2.515308in}{1.894079in}}%
\pgfpathlineto{\pgfqpoint{2.521153in}{1.902494in}}%
\pgfpathlineto{\pgfqpoint{2.524032in}{1.907327in}}%
\pgfpathlineto{\pgfqpoint{2.526776in}{1.920575in}}%
\pgfpathlineto{\pgfqpoint{2.534342in}{1.933823in}}%
\pgfpathlineto{\pgfqpoint{2.541647in}{1.947071in}}%
\pgfpathlineto{\pgfqpoint{2.546436in}{1.953818in}}%
\pgfpathlineto{\pgfqpoint{2.550105in}{1.960319in}}%
\pgfpathlineto{\pgfqpoint{2.559967in}{1.973567in}}%
\pgfpathlineto{\pgfqpoint{2.564079in}{1.986815in}}%
\pgfpathlineto{\pgfqpoint{2.571719in}{1.998570in}}%
\pgfpathlineto{\pgfqpoint{2.572620in}{2.000063in}}%
\pgfpathlineto{\pgfqpoint{2.575218in}{2.013311in}}%
\pgfpathlineto{\pgfqpoint{2.584971in}{2.026559in}}%
\pgfpathlineto{\pgfqpoint{2.587732in}{2.039807in}}%
\pgfpathlineto{\pgfqpoint{2.597002in}{2.049543in}}%
\pgfpathlineto{\pgfqpoint{2.599930in}{2.053055in}}%
\pgfpathlineto{\pgfqpoint{2.601700in}{2.066303in}}%
\pgfpathlineto{\pgfqpoint{2.611159in}{2.079551in}}%
\pgfpathlineto{\pgfqpoint{2.612638in}{2.092799in}}%
\pgfpathlineto{\pgfqpoint{2.622168in}{2.106047in}}%
\pgfpathlineto{\pgfqpoint{2.622285in}{2.106980in}}%
\pgfpathlineto{\pgfqpoint{2.623400in}{2.119295in}}%
\pgfpathlineto{\pgfqpoint{2.632421in}{2.132543in}}%
\pgfpathlineto{\pgfqpoint{2.622285in}{2.132543in}}%
\pgfpathlineto{\pgfqpoint{2.597002in}{2.132543in}}%
\pgfpathlineto{\pgfqpoint{2.571719in}{2.132543in}}%
\pgfpathlineto{\pgfqpoint{2.546436in}{2.132543in}}%
\pgfpathlineto{\pgfqpoint{2.521153in}{2.132543in}}%
\pgfpathlineto{\pgfqpoint{2.520208in}{2.132543in}}%
\pgfpathlineto{\pgfqpoint{2.507946in}{2.119295in}}%
\pgfpathlineto{\pgfqpoint{2.505137in}{2.106047in}}%
\pgfpathlineto{\pgfqpoint{2.495870in}{2.097402in}}%
\pgfpathlineto{\pgfqpoint{2.490721in}{2.092799in}}%
\pgfpathlineto{\pgfqpoint{2.483946in}{2.079551in}}%
\pgfpathlineto{\pgfqpoint{2.474607in}{2.066303in}}%
\pgfpathlineto{\pgfqpoint{2.470586in}{2.056822in}}%
\pgfpathlineto{\pgfqpoint{2.470586in}{2.053055in}}%
\pgfpathlineto{\pgfqpoint{2.470586in}{2.039807in}}%
\pgfpathlineto{\pgfqpoint{2.470586in}{2.026559in}}%
\pgfpathlineto{\pgfqpoint{2.470586in}{2.013311in}}%
\pgfpathlineto{\pgfqpoint{2.470586in}{2.000063in}}%
\pgfpathlineto{\pgfqpoint{2.470586in}{1.986815in}}%
\pgfpathlineto{\pgfqpoint{2.470586in}{1.973567in}}%
\pgfpathlineto{\pgfqpoint{2.470586in}{1.960319in}}%
\pgfpathlineto{\pgfqpoint{2.470586in}{1.947071in}}%
\pgfpathlineto{\pgfqpoint{2.470586in}{1.933823in}}%
\pgfpathlineto{\pgfqpoint{2.470586in}{1.920575in}}%
\pgfpathlineto{\pgfqpoint{2.470586in}{1.907327in}}%
\pgfpathlineto{\pgfqpoint{2.470586in}{1.894079in}}%
\pgfpathlineto{\pgfqpoint{2.470586in}{1.880831in}}%
\pgfpathlineto{\pgfqpoint{2.470586in}{1.867583in}}%
\pgfpathlineto{\pgfqpoint{2.470586in}{1.854335in}}%
\pgfpathlineto{\pgfqpoint{2.470586in}{1.841087in}}%
\pgfpathlineto{\pgfqpoint{2.470586in}{1.827839in}}%
\pgfpathlineto{\pgfqpoint{2.470586in}{1.821666in}}%
\pgfpathlineto{\pgfqpoint{2.477160in}{1.827839in}}%
\pgfpathclose%
\pgfpathmoveto{\pgfqpoint{4.290975in}{2.114802in}}%
\pgfpathlineto{\pgfqpoint{4.290975in}{2.119295in}}%
\pgfpathlineto{\pgfqpoint{4.290975in}{2.132543in}}%
\pgfpathlineto{\pgfqpoint{4.283458in}{2.132543in}}%
\pgfpathlineto{\pgfqpoint{4.286195in}{2.119295in}}%
\pgfpathlineto{\pgfqpoint{4.290975in}{2.114802in}}%
\pgfpathclose%
\pgfusepath{fill}%
\end{pgfscope}%
\begin{pgfscope}%
\pgfpathrectangle{\pgfqpoint{2.470586in}{0.476543in}}{\pgfqpoint{1.820388in}{1.656000in}}%
\pgfusepath{clip}%
\pgfsetbuttcap%
\pgfsetroundjoin%
\definecolor{currentfill}{rgb}{0.241369,0.366782,0.660592}%
\pgfsetfillcolor{currentfill}%
\pgfsetlinewidth{0.000000pt}%
\definecolor{currentstroke}{rgb}{0.000000,0.000000,0.000000}%
\pgfsetstrokecolor{currentstroke}%
\pgfsetdash{}{0pt}%
\pgfpathmoveto{\pgfqpoint{2.521153in}{0.595627in}}%
\pgfpathlineto{\pgfqpoint{2.546436in}{0.595225in}}%
\pgfpathlineto{\pgfqpoint{2.571719in}{0.594849in}}%
\pgfpathlineto{\pgfqpoint{2.597002in}{0.594496in}}%
\pgfpathlineto{\pgfqpoint{2.622285in}{0.594167in}}%
\pgfpathlineto{\pgfqpoint{2.647569in}{0.593858in}}%
\pgfpathlineto{\pgfqpoint{2.672852in}{0.593571in}}%
\pgfpathlineto{\pgfqpoint{2.698135in}{0.593302in}}%
\pgfpathlineto{\pgfqpoint{2.723418in}{0.593051in}}%
\pgfpathlineto{\pgfqpoint{2.748701in}{0.592818in}}%
\pgfpathlineto{\pgfqpoint{2.773984in}{0.592602in}}%
\pgfpathlineto{\pgfqpoint{2.799268in}{0.592401in}}%
\pgfpathlineto{\pgfqpoint{2.824551in}{0.592215in}}%
\pgfpathlineto{\pgfqpoint{2.849834in}{0.592044in}}%
\pgfpathlineto{\pgfqpoint{2.875117in}{0.591887in}}%
\pgfpathlineto{\pgfqpoint{2.900400in}{0.591743in}}%
\pgfpathlineto{\pgfqpoint{2.925683in}{0.591612in}}%
\pgfpathlineto{\pgfqpoint{2.950967in}{0.591494in}}%
\pgfpathlineto{\pgfqpoint{2.976250in}{0.591388in}}%
\pgfpathlineto{\pgfqpoint{3.001533in}{0.591294in}}%
\pgfpathlineto{\pgfqpoint{3.026816in}{0.591212in}}%
\pgfpathlineto{\pgfqpoint{3.052099in}{0.591141in}}%
\pgfpathlineto{\pgfqpoint{3.077382in}{0.591081in}}%
\pgfpathlineto{\pgfqpoint{3.102666in}{0.591032in}}%
\pgfpathlineto{\pgfqpoint{3.127949in}{0.590993in}}%
\pgfpathlineto{\pgfqpoint{3.153232in}{0.590966in}}%
\pgfpathlineto{\pgfqpoint{3.178515in}{0.590948in}}%
\pgfpathlineto{\pgfqpoint{3.203798in}{0.590941in}}%
\pgfpathlineto{\pgfqpoint{3.229082in}{0.590945in}}%
\pgfpathlineto{\pgfqpoint{3.254365in}{0.590959in}}%
\pgfpathlineto{\pgfqpoint{3.279648in}{0.590983in}}%
\pgfpathlineto{\pgfqpoint{3.304931in}{0.591018in}}%
\pgfpathlineto{\pgfqpoint{3.330214in}{0.591063in}}%
\pgfpathlineto{\pgfqpoint{3.355497in}{0.591119in}}%
\pgfpathlineto{\pgfqpoint{3.380781in}{0.591185in}}%
\pgfpathlineto{\pgfqpoint{3.406064in}{0.591262in}}%
\pgfpathlineto{\pgfqpoint{3.431347in}{0.591350in}}%
\pgfpathlineto{\pgfqpoint{3.456630in}{0.591449in}}%
\pgfpathlineto{\pgfqpoint{3.481913in}{0.591560in}}%
\pgfpathlineto{\pgfqpoint{3.507196in}{0.591682in}}%
\pgfpathlineto{\pgfqpoint{3.532480in}{0.591816in}}%
\pgfpathlineto{\pgfqpoint{3.557763in}{0.591962in}}%
\pgfpathlineto{\pgfqpoint{3.583046in}{0.592121in}}%
\pgfpathlineto{\pgfqpoint{3.608329in}{0.592293in}}%
\pgfpathlineto{\pgfqpoint{3.633612in}{0.592478in}}%
\pgfpathlineto{\pgfqpoint{3.658895in}{0.592677in}}%
\pgfpathlineto{\pgfqpoint{3.684179in}{0.592890in}}%
\pgfpathlineto{\pgfqpoint{3.709462in}{0.593118in}}%
\pgfpathlineto{\pgfqpoint{3.734745in}{0.593361in}}%
\pgfpathlineto{\pgfqpoint{3.760028in}{0.593621in}}%
\pgfpathlineto{\pgfqpoint{3.785311in}{0.593897in}}%
\pgfpathlineto{\pgfqpoint{3.810594in}{0.594191in}}%
\pgfpathlineto{\pgfqpoint{3.835878in}{0.594503in}}%
\pgfpathlineto{\pgfqpoint{3.861161in}{0.594834in}}%
\pgfpathlineto{\pgfqpoint{3.886444in}{0.595186in}}%
\pgfpathlineto{\pgfqpoint{3.911727in}{0.595558in}}%
\pgfpathlineto{\pgfqpoint{3.925760in}{0.595775in}}%
\pgfpathlineto{\pgfqpoint{3.937010in}{0.596222in}}%
\pgfpathlineto{\pgfqpoint{3.962293in}{0.597260in}}%
\pgfpathlineto{\pgfqpoint{3.987577in}{0.598336in}}%
\pgfpathlineto{\pgfqpoint{4.012860in}{0.599452in}}%
\pgfpathlineto{\pgfqpoint{4.038143in}{0.600612in}}%
\pgfpathlineto{\pgfqpoint{4.063426in}{0.601819in}}%
\pgfpathlineto{\pgfqpoint{4.088709in}{0.603075in}}%
\pgfpathlineto{\pgfqpoint{4.113993in}{0.604384in}}%
\pgfpathlineto{\pgfqpoint{4.139276in}{0.605749in}}%
\pgfpathlineto{\pgfqpoint{4.164559in}{0.607176in}}%
\pgfpathlineto{\pgfqpoint{4.189842in}{0.608667in}}%
\pgfpathlineto{\pgfqpoint{4.195664in}{0.609023in}}%
\pgfpathlineto{\pgfqpoint{4.215125in}{0.609749in}}%
\pgfpathlineto{\pgfqpoint{4.240408in}{0.610784in}}%
\pgfpathlineto{\pgfqpoint{4.265692in}{0.611933in}}%
\pgfpathlineto{\pgfqpoint{4.290975in}{0.613212in}}%
\pgfpathlineto{\pgfqpoint{4.290975in}{0.622271in}}%
\pgfpathlineto{\pgfqpoint{4.290975in}{0.635519in}}%
\pgfpathlineto{\pgfqpoint{4.290975in}{0.648767in}}%
\pgfpathlineto{\pgfqpoint{4.290975in}{0.662015in}}%
\pgfpathlineto{\pgfqpoint{4.290975in}{0.675263in}}%
\pgfpathlineto{\pgfqpoint{4.290975in}{0.688511in}}%
\pgfpathlineto{\pgfqpoint{4.290975in}{0.701759in}}%
\pgfpathlineto{\pgfqpoint{4.290975in}{0.715007in}}%
\pgfpathlineto{\pgfqpoint{4.290975in}{0.728255in}}%
\pgfpathlineto{\pgfqpoint{4.290975in}{0.741503in}}%
\pgfpathlineto{\pgfqpoint{4.290975in}{0.754751in}}%
\pgfpathlineto{\pgfqpoint{4.290975in}{0.767999in}}%
\pgfpathlineto{\pgfqpoint{4.290975in}{0.780126in}}%
\pgfpathlineto{\pgfqpoint{4.288057in}{0.781247in}}%
\pgfpathlineto{\pgfqpoint{4.278255in}{0.794495in}}%
\pgfpathlineto{\pgfqpoint{4.265692in}{0.799110in}}%
\pgfpathlineto{\pgfqpoint{4.243112in}{0.807743in}}%
\pgfpathlineto{\pgfqpoint{4.240408in}{0.809527in}}%
\pgfpathlineto{\pgfqpoint{4.224935in}{0.820991in}}%
\pgfpathlineto{\pgfqpoint{4.215125in}{0.824521in}}%
\pgfpathlineto{\pgfqpoint{4.189842in}{0.833868in}}%
\pgfpathlineto{\pgfqpoint{4.188883in}{0.834239in}}%
\pgfpathlineto{\pgfqpoint{4.165568in}{0.847487in}}%
\pgfpathlineto{\pgfqpoint{4.164559in}{0.848303in}}%
\pgfpathlineto{\pgfqpoint{4.150049in}{0.860735in}}%
\pgfpathlineto{\pgfqpoint{4.139276in}{0.867984in}}%
\pgfpathlineto{\pgfqpoint{4.131068in}{0.873983in}}%
\pgfpathlineto{\pgfqpoint{4.119760in}{0.887231in}}%
\pgfpathlineto{\pgfqpoint{4.113993in}{0.890243in}}%
\pgfpathlineto{\pgfqpoint{4.095724in}{0.900479in}}%
\pgfpathlineto{\pgfqpoint{4.088709in}{0.906342in}}%
\pgfpathlineto{\pgfqpoint{4.080712in}{0.913727in}}%
\pgfpathlineto{\pgfqpoint{4.063426in}{0.921237in}}%
\pgfpathlineto{\pgfqpoint{4.051140in}{0.926975in}}%
\pgfpathlineto{\pgfqpoint{4.043906in}{0.940223in}}%
\pgfpathlineto{\pgfqpoint{4.038143in}{0.943176in}}%
\pgfpathlineto{\pgfqpoint{4.018954in}{0.953471in}}%
\pgfpathlineto{\pgfqpoint{4.013086in}{0.966719in}}%
\pgfpathlineto{\pgfqpoint{4.012860in}{0.966852in}}%
\pgfpathlineto{\pgfqpoint{3.991806in}{0.979967in}}%
\pgfpathlineto{\pgfqpoint{3.987577in}{0.988429in}}%
\pgfpathlineto{\pgfqpoint{3.985544in}{0.993215in}}%
\pgfpathlineto{\pgfqpoint{3.962293in}{1.005889in}}%
\pgfpathlineto{\pgfqpoint{3.961322in}{1.006463in}}%
\pgfpathlineto{\pgfqpoint{3.953933in}{1.019711in}}%
\pgfpathlineto{\pgfqpoint{3.937010in}{1.030014in}}%
\pgfpathlineto{\pgfqpoint{3.932517in}{1.032959in}}%
\pgfpathlineto{\pgfqpoint{3.926041in}{1.046207in}}%
\pgfpathlineto{\pgfqpoint{3.912534in}{1.059455in}}%
\pgfpathlineto{\pgfqpoint{3.911727in}{1.060920in}}%
\pgfpathlineto{\pgfqpoint{3.906355in}{1.072703in}}%
\pgfpathlineto{\pgfqpoint{3.891904in}{1.085951in}}%
\pgfpathlineto{\pgfqpoint{3.886444in}{1.091657in}}%
\pgfpathlineto{\pgfqpoint{3.880476in}{1.099199in}}%
\pgfpathlineto{\pgfqpoint{3.866817in}{1.112447in}}%
\pgfpathlineto{\pgfqpoint{3.861161in}{1.117851in}}%
\pgfpathlineto{\pgfqpoint{3.854399in}{1.125695in}}%
\pgfpathlineto{\pgfqpoint{3.847352in}{1.138943in}}%
\pgfpathlineto{\pgfqpoint{3.835878in}{1.151727in}}%
\pgfpathlineto{\pgfqpoint{3.835528in}{1.152191in}}%
\pgfpathlineto{\pgfqpoint{3.829979in}{1.165439in}}%
\pgfpathlineto{\pgfqpoint{3.818311in}{1.178687in}}%
\pgfpathlineto{\pgfqpoint{3.811820in}{1.191935in}}%
\pgfpathlineto{\pgfqpoint{3.810594in}{1.192911in}}%
\pgfpathlineto{\pgfqpoint{3.798140in}{1.205183in}}%
\pgfpathlineto{\pgfqpoint{3.793229in}{1.218431in}}%
\pgfpathlineto{\pgfqpoint{3.785311in}{1.225809in}}%
\pgfpathlineto{\pgfqpoint{3.780116in}{1.231679in}}%
\pgfpathlineto{\pgfqpoint{3.777532in}{1.244927in}}%
\pgfpathlineto{\pgfqpoint{3.766932in}{1.258175in}}%
\pgfpathlineto{\pgfqpoint{3.762768in}{1.271423in}}%
\pgfpathlineto{\pgfqpoint{3.760028in}{1.273767in}}%
\pgfpathlineto{\pgfqpoint{3.750175in}{1.284671in}}%
\pgfpathlineto{\pgfqpoint{3.745308in}{1.297919in}}%
\pgfpathlineto{\pgfqpoint{3.734745in}{1.309044in}}%
\pgfpathlineto{\pgfqpoint{3.733174in}{1.311167in}}%
\pgfpathlineto{\pgfqpoint{3.729665in}{1.324415in}}%
\pgfpathlineto{\pgfqpoint{3.722669in}{1.337663in}}%
\pgfpathlineto{\pgfqpoint{3.718454in}{1.350911in}}%
\pgfpathlineto{\pgfqpoint{3.709462in}{1.363223in}}%
\pgfpathlineto{\pgfqpoint{3.708970in}{1.364159in}}%
\pgfpathlineto{\pgfqpoint{3.704135in}{1.377407in}}%
\pgfpathlineto{\pgfqpoint{3.697100in}{1.390655in}}%
\pgfpathlineto{\pgfqpoint{3.688920in}{1.403903in}}%
\pgfpathlineto{\pgfqpoint{3.684179in}{1.413782in}}%
\pgfpathlineto{\pgfqpoint{3.683089in}{1.417151in}}%
\pgfpathlineto{\pgfqpoint{3.677505in}{1.430399in}}%
\pgfpathlineto{\pgfqpoint{3.673992in}{1.443647in}}%
\pgfpathlineto{\pgfqpoint{3.666779in}{1.456895in}}%
\pgfpathlineto{\pgfqpoint{3.661751in}{1.470143in}}%
\pgfpathlineto{\pgfqpoint{3.658895in}{1.473220in}}%
\pgfpathlineto{\pgfqpoint{3.652707in}{1.483391in}}%
\pgfpathlineto{\pgfqpoint{3.650071in}{1.496639in}}%
\pgfpathlineto{\pgfqpoint{3.642355in}{1.509887in}}%
\pgfpathlineto{\pgfqpoint{3.640377in}{1.523135in}}%
\pgfpathlineto{\pgfqpoint{3.633612in}{1.533559in}}%
\pgfpathlineto{\pgfqpoint{3.632417in}{1.536383in}}%
\pgfpathlineto{\pgfqpoint{3.630589in}{1.549631in}}%
\pgfpathlineto{\pgfqpoint{3.622787in}{1.562879in}}%
\pgfpathlineto{\pgfqpoint{3.620034in}{1.576127in}}%
\pgfpathlineto{\pgfqpoint{3.610627in}{1.589375in}}%
\pgfpathlineto{\pgfqpoint{3.608329in}{1.596141in}}%
\pgfpathlineto{\pgfqpoint{3.607096in}{1.602623in}}%
\pgfpathlineto{\pgfqpoint{3.602272in}{1.615871in}}%
\pgfpathlineto{\pgfqpoint{3.600429in}{1.629119in}}%
\pgfpathlineto{\pgfqpoint{3.594736in}{1.642367in}}%
\pgfpathlineto{\pgfqpoint{3.590920in}{1.655615in}}%
\pgfpathlineto{\pgfqpoint{3.583046in}{1.668386in}}%
\pgfpathlineto{\pgfqpoint{3.582876in}{1.668863in}}%
\pgfpathlineto{\pgfqpoint{3.578872in}{1.682111in}}%
\pgfpathlineto{\pgfqpoint{3.575588in}{1.695359in}}%
\pgfpathlineto{\pgfqpoint{3.572028in}{1.708607in}}%
\pgfpathlineto{\pgfqpoint{3.568614in}{1.721855in}}%
\pgfpathlineto{\pgfqpoint{3.562794in}{1.735103in}}%
\pgfpathlineto{\pgfqpoint{3.557763in}{1.747220in}}%
\pgfpathlineto{\pgfqpoint{3.557524in}{1.748351in}}%
\pgfpathlineto{\pgfqpoint{3.552191in}{1.761599in}}%
\pgfpathlineto{\pgfqpoint{3.549948in}{1.774847in}}%
\pgfpathlineto{\pgfqpoint{3.544135in}{1.788095in}}%
\pgfpathlineto{\pgfqpoint{3.542159in}{1.801343in}}%
\pgfpathlineto{\pgfqpoint{3.535146in}{1.814591in}}%
\pgfpathlineto{\pgfqpoint{3.532780in}{1.827839in}}%
\pgfpathlineto{\pgfqpoint{3.532480in}{1.828197in}}%
\pgfpathlineto{\pgfqpoint{3.527445in}{1.841087in}}%
\pgfpathlineto{\pgfqpoint{3.525677in}{1.854335in}}%
\pgfpathlineto{\pgfqpoint{3.519021in}{1.867583in}}%
\pgfpathlineto{\pgfqpoint{3.516661in}{1.880831in}}%
\pgfpathlineto{\pgfqpoint{3.509664in}{1.894079in}}%
\pgfpathlineto{\pgfqpoint{3.507196in}{1.906737in}}%
\pgfpathlineto{\pgfqpoint{3.507142in}{1.907327in}}%
\pgfpathlineto{\pgfqpoint{3.502751in}{1.920575in}}%
\pgfpathlineto{\pgfqpoint{3.500577in}{1.933823in}}%
\pgfpathlineto{\pgfqpoint{3.494359in}{1.947071in}}%
\pgfpathlineto{\pgfqpoint{3.490463in}{1.960319in}}%
\pgfpathlineto{\pgfqpoint{3.484511in}{1.973567in}}%
\pgfpathlineto{\pgfqpoint{3.481913in}{1.980429in}}%
\pgfpathlineto{\pgfqpoint{3.480688in}{1.986815in}}%
\pgfpathlineto{\pgfqpoint{3.477658in}{2.000063in}}%
\pgfpathlineto{\pgfqpoint{3.474949in}{2.013311in}}%
\pgfpathlineto{\pgfqpoint{3.470927in}{2.026559in}}%
\pgfpathlineto{\pgfqpoint{3.464980in}{2.039807in}}%
\pgfpathlineto{\pgfqpoint{3.459471in}{2.053055in}}%
\pgfpathlineto{\pgfqpoint{3.456630in}{2.057029in}}%
\pgfpathlineto{\pgfqpoint{3.452924in}{2.066303in}}%
\pgfpathlineto{\pgfqpoint{3.450583in}{2.079551in}}%
\pgfpathlineto{\pgfqpoint{3.446062in}{2.092799in}}%
\pgfpathlineto{\pgfqpoint{3.443466in}{2.106047in}}%
\pgfpathlineto{\pgfqpoint{3.434645in}{2.119295in}}%
\pgfpathlineto{\pgfqpoint{3.431347in}{2.128980in}}%
\pgfpathlineto{\pgfqpoint{3.430617in}{2.132543in}}%
\pgfpathlineto{\pgfqpoint{3.406064in}{2.132543in}}%
\pgfpathlineto{\pgfqpoint{3.398955in}{2.132543in}}%
\pgfpathlineto{\pgfqpoint{3.403616in}{2.119295in}}%
\pgfpathlineto{\pgfqpoint{3.406064in}{2.115989in}}%
\pgfpathlineto{\pgfqpoint{3.412943in}{2.106047in}}%
\pgfpathlineto{\pgfqpoint{3.416723in}{2.092799in}}%
\pgfpathlineto{\pgfqpoint{3.423274in}{2.079551in}}%
\pgfpathlineto{\pgfqpoint{3.426823in}{2.066303in}}%
\pgfpathlineto{\pgfqpoint{3.431347in}{2.057525in}}%
\pgfpathlineto{\pgfqpoint{3.434004in}{2.053055in}}%
\pgfpathlineto{\pgfqpoint{3.439026in}{2.039807in}}%
\pgfpathlineto{\pgfqpoint{3.444852in}{2.026559in}}%
\pgfpathlineto{\pgfqpoint{3.449324in}{2.013311in}}%
\pgfpathlineto{\pgfqpoint{3.453169in}{2.000063in}}%
\pgfpathlineto{\pgfqpoint{3.456630in}{1.988496in}}%
\pgfpathlineto{\pgfqpoint{3.457293in}{1.986815in}}%
\pgfpathlineto{\pgfqpoint{3.461734in}{1.973567in}}%
\pgfpathlineto{\pgfqpoint{3.466929in}{1.960319in}}%
\pgfpathlineto{\pgfqpoint{3.470471in}{1.947071in}}%
\pgfpathlineto{\pgfqpoint{3.476552in}{1.933823in}}%
\pgfpathlineto{\pgfqpoint{3.479283in}{1.920575in}}%
\pgfpathlineto{\pgfqpoint{3.481913in}{1.913895in}}%
\pgfpathlineto{\pgfqpoint{3.485449in}{1.907327in}}%
\pgfpathlineto{\pgfqpoint{3.487877in}{1.894079in}}%
\pgfpathlineto{\pgfqpoint{3.493902in}{1.880831in}}%
\pgfpathlineto{\pgfqpoint{3.496182in}{1.867583in}}%
\pgfpathlineto{\pgfqpoint{3.502572in}{1.854335in}}%
\pgfpathlineto{\pgfqpoint{3.504485in}{1.841087in}}%
\pgfpathlineto{\pgfqpoint{3.507196in}{1.834595in}}%
\pgfpathlineto{\pgfqpoint{3.511170in}{1.827839in}}%
\pgfpathlineto{\pgfqpoint{3.513214in}{1.814591in}}%
\pgfpathlineto{\pgfqpoint{3.519348in}{1.801343in}}%
\pgfpathlineto{\pgfqpoint{3.521393in}{1.788095in}}%
\pgfpathlineto{\pgfqpoint{3.527087in}{1.774847in}}%
\pgfpathlineto{\pgfqpoint{3.529369in}{1.761599in}}%
\pgfpathlineto{\pgfqpoint{3.532480in}{1.753981in}}%
\pgfpathlineto{\pgfqpoint{3.535778in}{1.748351in}}%
\pgfpathlineto{\pgfqpoint{3.539993in}{1.735103in}}%
\pgfpathlineto{\pgfqpoint{3.545277in}{1.721855in}}%
\pgfpathlineto{\pgfqpoint{3.548518in}{1.708607in}}%
\pgfpathlineto{\pgfqpoint{3.552373in}{1.695359in}}%
\pgfpathlineto{\pgfqpoint{3.555795in}{1.682111in}}%
\pgfpathlineto{\pgfqpoint{3.557763in}{1.675515in}}%
\pgfpathlineto{\pgfqpoint{3.560651in}{1.668863in}}%
\pgfpathlineto{\pgfqpoint{3.567042in}{1.655615in}}%
\pgfpathlineto{\pgfqpoint{3.570536in}{1.642367in}}%
\pgfpathlineto{\pgfqpoint{3.575819in}{1.629119in}}%
\pgfpathlineto{\pgfqpoint{3.578145in}{1.615871in}}%
\pgfpathlineto{\pgfqpoint{3.583046in}{1.603469in}}%
\pgfpathlineto{\pgfqpoint{3.583509in}{1.602623in}}%
\pgfpathlineto{\pgfqpoint{3.587079in}{1.589375in}}%
\pgfpathlineto{\pgfqpoint{3.595147in}{1.576127in}}%
\pgfpathlineto{\pgfqpoint{3.597616in}{1.562879in}}%
\pgfpathlineto{\pgfqpoint{3.604968in}{1.549631in}}%
\pgfpathlineto{\pgfqpoint{3.607060in}{1.536383in}}%
\pgfpathlineto{\pgfqpoint{3.608329in}{1.533554in}}%
\pgfpathlineto{\pgfqpoint{3.614511in}{1.523135in}}%
\pgfpathlineto{\pgfqpoint{3.616379in}{1.509887in}}%
\pgfpathlineto{\pgfqpoint{3.623832in}{1.496639in}}%
\pgfpathlineto{\pgfqpoint{3.626402in}{1.483391in}}%
\pgfpathlineto{\pgfqpoint{3.633612in}{1.470726in}}%
\pgfpathlineto{\pgfqpoint{3.634056in}{1.470143in}}%
\pgfpathlineto{\pgfqpoint{3.638575in}{1.456895in}}%
\pgfpathlineto{\pgfqpoint{3.645620in}{1.443647in}}%
\pgfpathlineto{\pgfqpoint{3.649238in}{1.430399in}}%
\pgfpathlineto{\pgfqpoint{3.655179in}{1.417151in}}%
\pgfpathlineto{\pgfqpoint{3.658895in}{1.406074in}}%
\pgfpathlineto{\pgfqpoint{3.659852in}{1.403903in}}%
\pgfpathlineto{\pgfqpoint{3.667068in}{1.390655in}}%
\pgfpathlineto{\pgfqpoint{3.673940in}{1.377407in}}%
\pgfpathlineto{\pgfqpoint{3.678589in}{1.364159in}}%
\pgfpathlineto{\pgfqpoint{3.684179in}{1.353196in}}%
\pgfpathlineto{\pgfqpoint{3.685658in}{1.350911in}}%
\pgfpathlineto{\pgfqpoint{3.690118in}{1.337663in}}%
\pgfpathlineto{\pgfqpoint{3.697716in}{1.324415in}}%
\pgfpathlineto{\pgfqpoint{3.700904in}{1.311167in}}%
\pgfpathlineto{\pgfqpoint{3.709462in}{1.299243in}}%
\pgfpathlineto{\pgfqpoint{3.710611in}{1.297919in}}%
\pgfpathlineto{\pgfqpoint{3.715061in}{1.284671in}}%
\pgfpathlineto{\pgfqpoint{3.726971in}{1.271423in}}%
\pgfpathlineto{\pgfqpoint{3.729933in}{1.258175in}}%
\pgfpathlineto{\pgfqpoint{3.734745in}{1.251212in}}%
\pgfpathlineto{\pgfqpoint{3.739896in}{1.244927in}}%
\pgfpathlineto{\pgfqpoint{3.742388in}{1.231679in}}%
\pgfpathlineto{\pgfqpoint{3.754147in}{1.218431in}}%
\pgfpathlineto{\pgfqpoint{3.757733in}{1.205183in}}%
\pgfpathlineto{\pgfqpoint{3.760028in}{1.202426in}}%
\pgfpathlineto{\pgfqpoint{3.770332in}{1.191935in}}%
\pgfpathlineto{\pgfqpoint{3.775242in}{1.178687in}}%
\pgfpathlineto{\pgfqpoint{3.784978in}{1.165439in}}%
\pgfpathlineto{\pgfqpoint{3.785311in}{1.164587in}}%
\pgfpathlineto{\pgfqpoint{3.791048in}{1.152191in}}%
\pgfpathlineto{\pgfqpoint{3.800947in}{1.138943in}}%
\pgfpathlineto{\pgfqpoint{3.807353in}{1.125695in}}%
\pgfpathlineto{\pgfqpoint{3.810594in}{1.120855in}}%
\pgfpathlineto{\pgfqpoint{3.817259in}{1.112447in}}%
\pgfpathlineto{\pgfqpoint{3.829354in}{1.099199in}}%
\pgfpathlineto{\pgfqpoint{3.835878in}{1.088653in}}%
\pgfpathlineto{\pgfqpoint{3.837850in}{1.085951in}}%
\pgfpathlineto{\pgfqpoint{3.850679in}{1.072703in}}%
\pgfpathlineto{\pgfqpoint{3.855351in}{1.059455in}}%
\pgfpathlineto{\pgfqpoint{3.861161in}{1.052839in}}%
\pgfpathlineto{\pgfqpoint{3.867538in}{1.046207in}}%
\pgfpathlineto{\pgfqpoint{3.871641in}{1.032959in}}%
\pgfpathlineto{\pgfqpoint{3.886444in}{1.022295in}}%
\pgfpathlineto{\pgfqpoint{3.890294in}{1.019711in}}%
\pgfpathlineto{\pgfqpoint{3.894905in}{1.006463in}}%
\pgfpathlineto{\pgfqpoint{3.911727in}{0.995018in}}%
\pgfpathlineto{\pgfqpoint{3.914600in}{0.993215in}}%
\pgfpathlineto{\pgfqpoint{3.918449in}{0.979967in}}%
\pgfpathlineto{\pgfqpoint{3.937010in}{0.966786in}}%
\pgfpathlineto{\pgfqpoint{3.937111in}{0.966719in}}%
\pgfpathlineto{\pgfqpoint{3.939937in}{0.953471in}}%
\pgfpathlineto{\pgfqpoint{3.960550in}{0.940223in}}%
\pgfpathlineto{\pgfqpoint{3.962293in}{0.933873in}}%
\pgfpathlineto{\pgfqpoint{3.964421in}{0.926975in}}%
\pgfpathlineto{\pgfqpoint{3.986062in}{0.913727in}}%
\pgfpathlineto{\pgfqpoint{3.987577in}{0.911741in}}%
\pgfpathlineto{\pgfqpoint{3.996982in}{0.900479in}}%
\pgfpathlineto{\pgfqpoint{4.012860in}{0.887755in}}%
\pgfpathlineto{\pgfqpoint{4.013582in}{0.887231in}}%
\pgfpathlineto{\pgfqpoint{4.020484in}{0.873983in}}%
\pgfpathlineto{\pgfqpoint{4.031398in}{0.860735in}}%
\pgfpathlineto{\pgfqpoint{4.038143in}{0.853720in}}%
\pgfpathlineto{\pgfqpoint{4.044457in}{0.847487in}}%
\pgfpathlineto{\pgfqpoint{4.054443in}{0.834239in}}%
\pgfpathlineto{\pgfqpoint{4.063426in}{0.829788in}}%
\pgfpathlineto{\pgfqpoint{4.082126in}{0.820991in}}%
\pgfpathlineto{\pgfqpoint{4.085272in}{0.807743in}}%
\pgfpathlineto{\pgfqpoint{4.088709in}{0.805793in}}%
\pgfpathlineto{\pgfqpoint{4.109897in}{0.794495in}}%
\pgfpathlineto{\pgfqpoint{4.102329in}{0.781247in}}%
\pgfpathlineto{\pgfqpoint{4.113993in}{0.775166in}}%
\pgfpathlineto{\pgfqpoint{4.128532in}{0.767999in}}%
\pgfpathlineto{\pgfqpoint{4.113993in}{0.760158in}}%
\pgfpathlineto{\pgfqpoint{4.102476in}{0.754751in}}%
\pgfpathlineto{\pgfqpoint{4.113993in}{0.749754in}}%
\pgfpathlineto{\pgfqpoint{4.135102in}{0.741503in}}%
\pgfpathlineto{\pgfqpoint{4.113993in}{0.731720in}}%
\pgfpathlineto{\pgfqpoint{4.105499in}{0.728255in}}%
\pgfpathlineto{\pgfqpoint{4.113993in}{0.723486in}}%
\pgfpathlineto{\pgfqpoint{4.131292in}{0.715007in}}%
\pgfpathlineto{\pgfqpoint{4.113993in}{0.709523in}}%
\pgfpathlineto{\pgfqpoint{4.088709in}{0.702642in}}%
\pgfpathlineto{\pgfqpoint{4.085119in}{0.701759in}}%
\pgfpathlineto{\pgfqpoint{4.069709in}{0.688511in}}%
\pgfpathlineto{\pgfqpoint{4.063426in}{0.687614in}}%
\pgfpathlineto{\pgfqpoint{4.038143in}{0.684194in}}%
\pgfpathlineto{\pgfqpoint{4.012860in}{0.681032in}}%
\pgfpathlineto{\pgfqpoint{3.987577in}{0.678096in}}%
\pgfpathlineto{\pgfqpoint{3.962293in}{0.675361in}}%
\pgfpathlineto{\pgfqpoint{3.961346in}{0.675263in}}%
\pgfpathlineto{\pgfqpoint{3.937010in}{0.671850in}}%
\pgfpathlineto{\pgfqpoint{3.911727in}{0.668871in}}%
\pgfpathlineto{\pgfqpoint{3.886444in}{0.666359in}}%
\pgfpathlineto{\pgfqpoint{3.861161in}{0.664227in}}%
\pgfpathlineto{\pgfqpoint{3.835878in}{0.662410in}}%
\pgfpathlineto{\pgfqpoint{3.829741in}{0.662015in}}%
\pgfpathlineto{\pgfqpoint{3.810594in}{0.659844in}}%
\pgfpathlineto{\pgfqpoint{3.785311in}{0.657187in}}%
\pgfpathlineto{\pgfqpoint{3.760028in}{0.654733in}}%
\pgfpathlineto{\pgfqpoint{3.734745in}{0.652465in}}%
\pgfpathlineto{\pgfqpoint{3.709462in}{0.650364in}}%
\pgfpathlineto{\pgfqpoint{3.688790in}{0.648767in}}%
\pgfpathlineto{\pgfqpoint{3.684179in}{0.648602in}}%
\pgfpathlineto{\pgfqpoint{3.658895in}{0.647768in}}%
\pgfpathlineto{\pgfqpoint{3.633612in}{0.647019in}}%
\pgfpathlineto{\pgfqpoint{3.608329in}{0.646348in}}%
\pgfpathlineto{\pgfqpoint{3.583046in}{0.645750in}}%
\pgfpathlineto{\pgfqpoint{3.557763in}{0.645217in}}%
\pgfpathlineto{\pgfqpoint{3.532480in}{0.644748in}}%
\pgfpathlineto{\pgfqpoint{3.507196in}{0.644336in}}%
\pgfpathlineto{\pgfqpoint{3.481913in}{0.643980in}}%
\pgfpathlineto{\pgfqpoint{3.456630in}{0.643675in}}%
\pgfpathlineto{\pgfqpoint{3.431347in}{0.643419in}}%
\pgfpathlineto{\pgfqpoint{3.406064in}{0.643210in}}%
\pgfpathlineto{\pgfqpoint{3.380781in}{0.643045in}}%
\pgfpathlineto{\pgfqpoint{3.355497in}{0.642924in}}%
\pgfpathlineto{\pgfqpoint{3.330214in}{0.642845in}}%
\pgfpathlineto{\pgfqpoint{3.304931in}{0.642806in}}%
\pgfpathlineto{\pgfqpoint{3.279648in}{0.642807in}}%
\pgfpathlineto{\pgfqpoint{3.254365in}{0.642846in}}%
\pgfpathlineto{\pgfqpoint{3.229082in}{0.642925in}}%
\pgfpathlineto{\pgfqpoint{3.203798in}{0.643041in}}%
\pgfpathlineto{\pgfqpoint{3.178515in}{0.643196in}}%
\pgfpathlineto{\pgfqpoint{3.153232in}{0.643389in}}%
\pgfpathlineto{\pgfqpoint{3.127949in}{0.643621in}}%
\pgfpathlineto{\pgfqpoint{3.102666in}{0.643893in}}%
\pgfpathlineto{\pgfqpoint{3.077382in}{0.644205in}}%
\pgfpathlineto{\pgfqpoint{3.052099in}{0.644558in}}%
\pgfpathlineto{\pgfqpoint{3.026816in}{0.644954in}}%
\pgfpathlineto{\pgfqpoint{3.001533in}{0.645394in}}%
\pgfpathlineto{\pgfqpoint{2.976250in}{0.645880in}}%
\pgfpathlineto{\pgfqpoint{2.950967in}{0.646414in}}%
\pgfpathlineto{\pgfqpoint{2.925683in}{0.646999in}}%
\pgfpathlineto{\pgfqpoint{2.900400in}{0.647637in}}%
\pgfpathlineto{\pgfqpoint{2.875117in}{0.648332in}}%
\pgfpathlineto{\pgfqpoint{2.860434in}{0.648767in}}%
\pgfpathlineto{\pgfqpoint{2.849834in}{0.651007in}}%
\pgfpathlineto{\pgfqpoint{2.824551in}{0.657700in}}%
\pgfpathlineto{\pgfqpoint{2.811743in}{0.662015in}}%
\pgfpathlineto{\pgfqpoint{2.799268in}{0.662488in}}%
\pgfpathlineto{\pgfqpoint{2.773984in}{0.663516in}}%
\pgfpathlineto{\pgfqpoint{2.748701in}{0.664633in}}%
\pgfpathlineto{\pgfqpoint{2.723418in}{0.665849in}}%
\pgfpathlineto{\pgfqpoint{2.698135in}{0.667173in}}%
\pgfpathlineto{\pgfqpoint{2.672852in}{0.668620in}}%
\pgfpathlineto{\pgfqpoint{2.647569in}{0.670203in}}%
\pgfpathlineto{\pgfqpoint{2.622285in}{0.671940in}}%
\pgfpathlineto{\pgfqpoint{2.597002in}{0.673850in}}%
\pgfpathlineto{\pgfqpoint{2.579801in}{0.675263in}}%
\pgfpathlineto{\pgfqpoint{2.597002in}{0.685847in}}%
\pgfpathlineto{\pgfqpoint{2.600808in}{0.688511in}}%
\pgfpathlineto{\pgfqpoint{2.597002in}{0.688840in}}%
\pgfpathlineto{\pgfqpoint{2.571719in}{0.691163in}}%
\pgfpathlineto{\pgfqpoint{2.546436in}{0.693740in}}%
\pgfpathlineto{\pgfqpoint{2.521153in}{0.696615in}}%
\pgfpathlineto{\pgfqpoint{2.495870in}{0.699841in}}%
\pgfpathlineto{\pgfqpoint{2.482168in}{0.701759in}}%
\pgfpathlineto{\pgfqpoint{2.495870in}{0.705596in}}%
\pgfpathlineto{\pgfqpoint{2.521153in}{0.713368in}}%
\pgfpathlineto{\pgfqpoint{2.526063in}{0.715007in}}%
\pgfpathlineto{\pgfqpoint{2.521153in}{0.716062in}}%
\pgfpathlineto{\pgfqpoint{2.495870in}{0.721982in}}%
\pgfpathlineto{\pgfqpoint{2.472762in}{0.728255in}}%
\pgfpathlineto{\pgfqpoint{2.495870in}{0.732941in}}%
\pgfpathlineto{\pgfqpoint{2.521153in}{0.738477in}}%
\pgfpathlineto{\pgfqpoint{2.533995in}{0.741503in}}%
\pgfpathlineto{\pgfqpoint{2.521153in}{0.749143in}}%
\pgfpathlineto{\pgfqpoint{2.512766in}{0.754751in}}%
\pgfpathlineto{\pgfqpoint{2.521153in}{0.757180in}}%
\pgfpathlineto{\pgfqpoint{2.546436in}{0.765110in}}%
\pgfpathlineto{\pgfqpoint{2.554851in}{0.767999in}}%
\pgfpathlineto{\pgfqpoint{2.546436in}{0.771964in}}%
\pgfpathlineto{\pgfqpoint{2.528472in}{0.781247in}}%
\pgfpathlineto{\pgfqpoint{2.546436in}{0.788103in}}%
\pgfpathlineto{\pgfqpoint{2.561726in}{0.794495in}}%
\pgfpathlineto{\pgfqpoint{2.549179in}{0.807743in}}%
\pgfpathlineto{\pgfqpoint{2.571719in}{0.815030in}}%
\pgfpathlineto{\pgfqpoint{2.588239in}{0.820991in}}%
\pgfpathlineto{\pgfqpoint{2.589533in}{0.834239in}}%
\pgfpathlineto{\pgfqpoint{2.597002in}{0.837660in}}%
\pgfpathlineto{\pgfqpoint{2.616099in}{0.847487in}}%
\pgfpathlineto{\pgfqpoint{2.619304in}{0.860735in}}%
\pgfpathlineto{\pgfqpoint{2.622285in}{0.862576in}}%
\pgfpathlineto{\pgfqpoint{2.638907in}{0.873983in}}%
\pgfpathlineto{\pgfqpoint{2.646610in}{0.887231in}}%
\pgfpathlineto{\pgfqpoint{2.647569in}{0.887854in}}%
\pgfpathlineto{\pgfqpoint{2.664528in}{0.900479in}}%
\pgfpathlineto{\pgfqpoint{2.672852in}{0.908279in}}%
\pgfpathlineto{\pgfqpoint{2.678430in}{0.913727in}}%
\pgfpathlineto{\pgfqpoint{2.693125in}{0.926975in}}%
\pgfpathlineto{\pgfqpoint{2.698135in}{0.930199in}}%
\pgfpathlineto{\pgfqpoint{2.713136in}{0.940223in}}%
\pgfpathlineto{\pgfqpoint{2.721950in}{0.953471in}}%
\pgfpathlineto{\pgfqpoint{2.723418in}{0.954708in}}%
\pgfpathlineto{\pgfqpoint{2.736968in}{0.966719in}}%
\pgfpathlineto{\pgfqpoint{2.740364in}{0.979967in}}%
\pgfpathlineto{\pgfqpoint{2.748701in}{0.986257in}}%
\pgfpathlineto{\pgfqpoint{2.757552in}{0.993215in}}%
\pgfpathlineto{\pgfqpoint{2.762033in}{1.006463in}}%
\pgfpathlineto{\pgfqpoint{2.773984in}{1.013222in}}%
\pgfpathlineto{\pgfqpoint{2.784837in}{1.019711in}}%
\pgfpathlineto{\pgfqpoint{2.787680in}{1.032959in}}%
\pgfpathlineto{\pgfqpoint{2.799268in}{1.041102in}}%
\pgfpathlineto{\pgfqpoint{2.806022in}{1.046207in}}%
\pgfpathlineto{\pgfqpoint{2.807669in}{1.059455in}}%
\pgfpathlineto{\pgfqpoint{2.824551in}{1.072031in}}%
\pgfpathlineto{\pgfqpoint{2.825384in}{1.072703in}}%
\pgfpathlineto{\pgfqpoint{2.828021in}{1.085951in}}%
\pgfpathlineto{\pgfqpoint{2.845693in}{1.099199in}}%
\pgfpathlineto{\pgfqpoint{2.849729in}{1.112447in}}%
\pgfpathlineto{\pgfqpoint{2.849834in}{1.112530in}}%
\pgfpathlineto{\pgfqpoint{2.864268in}{1.125695in}}%
\pgfpathlineto{\pgfqpoint{2.869941in}{1.138943in}}%
\pgfpathlineto{\pgfqpoint{2.875117in}{1.145019in}}%
\pgfpathlineto{\pgfqpoint{2.880289in}{1.152191in}}%
\pgfpathlineto{\pgfqpoint{2.884051in}{1.165439in}}%
\pgfpathlineto{\pgfqpoint{2.893465in}{1.178687in}}%
\pgfpathlineto{\pgfqpoint{2.900400in}{1.189422in}}%
\pgfpathlineto{\pgfqpoint{2.901816in}{1.191935in}}%
\pgfpathlineto{\pgfqpoint{2.910605in}{1.205183in}}%
\pgfpathlineto{\pgfqpoint{2.921597in}{1.218431in}}%
\pgfpathlineto{\pgfqpoint{2.925683in}{1.227850in}}%
\pgfpathlineto{\pgfqpoint{2.926997in}{1.231679in}}%
\pgfpathlineto{\pgfqpoint{2.935688in}{1.244927in}}%
\pgfpathlineto{\pgfqpoint{2.939516in}{1.258175in}}%
\pgfpathlineto{\pgfqpoint{2.950527in}{1.271423in}}%
\pgfpathlineto{\pgfqpoint{2.950967in}{1.272997in}}%
\pgfpathlineto{\pgfqpoint{2.953429in}{1.284671in}}%
\pgfpathlineto{\pgfqpoint{2.964252in}{1.297919in}}%
\pgfpathlineto{\pgfqpoint{2.966827in}{1.311167in}}%
\pgfpathlineto{\pgfqpoint{2.976250in}{1.321709in}}%
\pgfpathlineto{\pgfqpoint{2.978270in}{1.324415in}}%
\pgfpathlineto{\pgfqpoint{2.980411in}{1.337663in}}%
\pgfpathlineto{\pgfqpoint{2.988877in}{1.350911in}}%
\pgfpathlineto{\pgfqpoint{2.990964in}{1.364159in}}%
\pgfpathlineto{\pgfqpoint{3.001533in}{1.376830in}}%
\pgfpathlineto{\pgfqpoint{3.001920in}{1.377407in}}%
\pgfpathlineto{\pgfqpoint{3.005598in}{1.390655in}}%
\pgfpathlineto{\pgfqpoint{3.013997in}{1.403903in}}%
\pgfpathlineto{\pgfqpoint{3.018387in}{1.417151in}}%
\pgfpathlineto{\pgfqpoint{3.024642in}{1.430399in}}%
\pgfpathlineto{\pgfqpoint{3.026816in}{1.436141in}}%
\pgfpathlineto{\pgfqpoint{3.028987in}{1.443647in}}%
\pgfpathlineto{\pgfqpoint{3.033901in}{1.456895in}}%
\pgfpathlineto{\pgfqpoint{3.039486in}{1.470143in}}%
\pgfpathlineto{\pgfqpoint{3.044207in}{1.483391in}}%
\pgfpathlineto{\pgfqpoint{3.052099in}{1.496097in}}%
\pgfpathlineto{\pgfqpoint{3.052362in}{1.496639in}}%
\pgfpathlineto{\pgfqpoint{3.055228in}{1.509887in}}%
\pgfpathlineto{\pgfqpoint{3.061946in}{1.523135in}}%
\pgfpathlineto{\pgfqpoint{3.064341in}{1.536383in}}%
\pgfpathlineto{\pgfqpoint{3.071592in}{1.549631in}}%
\pgfpathlineto{\pgfqpoint{3.074166in}{1.562879in}}%
\pgfpathlineto{\pgfqpoint{3.077382in}{1.567103in}}%
\pgfpathlineto{\pgfqpoint{3.082502in}{1.576127in}}%
\pgfpathlineto{\pgfqpoint{3.084651in}{1.589375in}}%
\pgfpathlineto{\pgfqpoint{3.091513in}{1.602623in}}%
\pgfpathlineto{\pgfqpoint{3.093276in}{1.615871in}}%
\pgfpathlineto{\pgfqpoint{3.099787in}{1.629119in}}%
\pgfpathlineto{\pgfqpoint{3.102666in}{1.640358in}}%
\pgfpathlineto{\pgfqpoint{3.103024in}{1.642367in}}%
\pgfpathlineto{\pgfqpoint{3.108693in}{1.655615in}}%
\pgfpathlineto{\pgfqpoint{3.112122in}{1.668863in}}%
\pgfpathlineto{\pgfqpoint{3.117479in}{1.682111in}}%
\pgfpathlineto{\pgfqpoint{3.122611in}{1.695359in}}%
\pgfpathlineto{\pgfqpoint{3.127389in}{1.708607in}}%
\pgfpathlineto{\pgfqpoint{3.127949in}{1.709947in}}%
\pgfpathlineto{\pgfqpoint{3.131445in}{1.721855in}}%
\pgfpathlineto{\pgfqpoint{3.134089in}{1.735103in}}%
\pgfpathlineto{\pgfqpoint{3.138927in}{1.748351in}}%
\pgfpathlineto{\pgfqpoint{3.141934in}{1.761599in}}%
\pgfpathlineto{\pgfqpoint{3.149575in}{1.774847in}}%
\pgfpathlineto{\pgfqpoint{3.152076in}{1.788095in}}%
\pgfpathlineto{\pgfqpoint{3.153232in}{1.790237in}}%
\pgfpathlineto{\pgfqpoint{3.157495in}{1.801343in}}%
\pgfpathlineto{\pgfqpoint{3.158999in}{1.814591in}}%
\pgfpathlineto{\pgfqpoint{3.164971in}{1.827839in}}%
\pgfpathlineto{\pgfqpoint{3.167438in}{1.841087in}}%
\pgfpathlineto{\pgfqpoint{3.175004in}{1.854335in}}%
\pgfpathlineto{\pgfqpoint{3.178094in}{1.867583in}}%
\pgfpathlineto{\pgfqpoint{3.178515in}{1.868379in}}%
\pgfpathlineto{\pgfqpoint{3.183113in}{1.880831in}}%
\pgfpathlineto{\pgfqpoint{3.185922in}{1.894079in}}%
\pgfpathlineto{\pgfqpoint{3.190145in}{1.907327in}}%
\pgfpathlineto{\pgfqpoint{3.193374in}{1.920575in}}%
\pgfpathlineto{\pgfqpoint{3.198169in}{1.933823in}}%
\pgfpathlineto{\pgfqpoint{3.203798in}{1.946074in}}%
\pgfpathlineto{\pgfqpoint{3.204154in}{1.947071in}}%
\pgfpathlineto{\pgfqpoint{3.207899in}{1.960319in}}%
\pgfpathlineto{\pgfqpoint{3.213678in}{1.973567in}}%
\pgfpathlineto{\pgfqpoint{3.216176in}{1.986815in}}%
\pgfpathlineto{\pgfqpoint{3.222184in}{2.000063in}}%
\pgfpathlineto{\pgfqpoint{3.224970in}{2.013311in}}%
\pgfpathlineto{\pgfqpoint{3.229082in}{2.019813in}}%
\pgfpathlineto{\pgfqpoint{3.232673in}{2.026559in}}%
\pgfpathlineto{\pgfqpoint{3.234963in}{2.039807in}}%
\pgfpathlineto{\pgfqpoint{3.242497in}{2.053055in}}%
\pgfpathlineto{\pgfqpoint{3.245188in}{2.066303in}}%
\pgfpathlineto{\pgfqpoint{3.253060in}{2.079551in}}%
\pgfpathlineto{\pgfqpoint{3.254365in}{2.084128in}}%
\pgfpathlineto{\pgfqpoint{3.256651in}{2.092799in}}%
\pgfpathlineto{\pgfqpoint{3.263832in}{2.106047in}}%
\pgfpathlineto{\pgfqpoint{3.268675in}{2.119295in}}%
\pgfpathlineto{\pgfqpoint{3.277289in}{2.132543in}}%
\pgfpathlineto{\pgfqpoint{3.254365in}{2.132543in}}%
\pgfpathlineto{\pgfqpoint{3.245277in}{2.132543in}}%
\pgfpathlineto{\pgfqpoint{3.237662in}{2.119295in}}%
\pgfpathlineto{\pgfqpoint{3.234934in}{2.106047in}}%
\pgfpathlineto{\pgfqpoint{3.229766in}{2.092799in}}%
\pgfpathlineto{\pgfqpoint{3.229082in}{2.088315in}}%
\pgfpathlineto{\pgfqpoint{3.226708in}{2.079551in}}%
\pgfpathlineto{\pgfqpoint{3.219051in}{2.066303in}}%
\pgfpathlineto{\pgfqpoint{3.216768in}{2.053055in}}%
\pgfpathlineto{\pgfqpoint{3.209991in}{2.039807in}}%
\pgfpathlineto{\pgfqpoint{3.207861in}{2.026559in}}%
\pgfpathlineto{\pgfqpoint{3.203798in}{2.016253in}}%
\pgfpathlineto{\pgfqpoint{3.201642in}{2.013311in}}%
\pgfpathlineto{\pgfqpoint{3.198031in}{2.000063in}}%
\pgfpathlineto{\pgfqpoint{3.192673in}{1.986815in}}%
\pgfpathlineto{\pgfqpoint{3.190102in}{1.973567in}}%
\pgfpathlineto{\pgfqpoint{3.184965in}{1.960319in}}%
\pgfpathlineto{\pgfqpoint{3.181161in}{1.947071in}}%
\pgfpathlineto{\pgfqpoint{3.178515in}{1.936771in}}%
\pgfpathlineto{\pgfqpoint{3.176963in}{1.933823in}}%
\pgfpathlineto{\pgfqpoint{3.170135in}{1.920575in}}%
\pgfpathlineto{\pgfqpoint{3.167472in}{1.907327in}}%
\pgfpathlineto{\pgfqpoint{3.163253in}{1.894079in}}%
\pgfpathlineto{\pgfqpoint{3.160898in}{1.880831in}}%
\pgfpathlineto{\pgfqpoint{3.156143in}{1.867583in}}%
\pgfpathlineto{\pgfqpoint{3.154320in}{1.854335in}}%
\pgfpathlineto{\pgfqpoint{3.153232in}{1.851835in}}%
\pgfpathlineto{\pgfqpoint{3.144924in}{1.841087in}}%
\pgfpathlineto{\pgfqpoint{3.142114in}{1.827839in}}%
\pgfpathlineto{\pgfqpoint{3.136463in}{1.814591in}}%
\pgfpathlineto{\pgfqpoint{3.135067in}{1.801343in}}%
\pgfpathlineto{\pgfqpoint{3.130628in}{1.788095in}}%
\pgfpathlineto{\pgfqpoint{3.128895in}{1.774847in}}%
\pgfpathlineto{\pgfqpoint{3.127949in}{1.772612in}}%
\pgfpathlineto{\pgfqpoint{3.120006in}{1.761599in}}%
\pgfpathlineto{\pgfqpoint{3.115815in}{1.748351in}}%
\pgfpathlineto{\pgfqpoint{3.111375in}{1.735103in}}%
\pgfpathlineto{\pgfqpoint{3.108362in}{1.721855in}}%
\pgfpathlineto{\pgfqpoint{3.105233in}{1.708607in}}%
\pgfpathlineto{\pgfqpoint{3.102666in}{1.698247in}}%
\pgfpathlineto{\pgfqpoint{3.101328in}{1.695359in}}%
\pgfpathlineto{\pgfqpoint{3.095435in}{1.682111in}}%
\pgfpathlineto{\pgfqpoint{3.088660in}{1.668863in}}%
\pgfpathlineto{\pgfqpoint{3.085391in}{1.655615in}}%
\pgfpathlineto{\pgfqpoint{3.079330in}{1.642367in}}%
\pgfpathlineto{\pgfqpoint{3.077467in}{1.629119in}}%
\pgfpathlineto{\pgfqpoint{3.077382in}{1.628880in}}%
\pgfpathlineto{\pgfqpoint{3.070004in}{1.615871in}}%
\pgfpathlineto{\pgfqpoint{3.068226in}{1.602623in}}%
\pgfpathlineto{\pgfqpoint{3.060680in}{1.589375in}}%
\pgfpathlineto{\pgfqpoint{3.058303in}{1.576127in}}%
\pgfpathlineto{\pgfqpoint{3.052099in}{1.565229in}}%
\pgfpathlineto{\pgfqpoint{3.049957in}{1.562879in}}%
\pgfpathlineto{\pgfqpoint{3.046300in}{1.549631in}}%
\pgfpathlineto{\pgfqpoint{3.039196in}{1.536383in}}%
\pgfpathlineto{\pgfqpoint{3.036372in}{1.523135in}}%
\pgfpathlineto{\pgfqpoint{3.030195in}{1.509887in}}%
\pgfpathlineto{\pgfqpoint{3.027001in}{1.496639in}}%
\pgfpathlineto{\pgfqpoint{3.026816in}{1.496238in}}%
\pgfpathlineto{\pgfqpoint{3.018151in}{1.483391in}}%
\pgfpathlineto{\pgfqpoint{3.011952in}{1.470143in}}%
\pgfpathlineto{\pgfqpoint{3.007074in}{1.456895in}}%
\pgfpathlineto{\pgfqpoint{3.001660in}{1.443647in}}%
\pgfpathlineto{\pgfqpoint{3.001533in}{1.443108in}}%
\pgfpathlineto{\pgfqpoint{2.997037in}{1.430399in}}%
\pgfpathlineto{\pgfqpoint{2.990026in}{1.417151in}}%
\pgfpathlineto{\pgfqpoint{2.985776in}{1.403903in}}%
\pgfpathlineto{\pgfqpoint{2.976592in}{1.390655in}}%
\pgfpathlineto{\pgfqpoint{2.976250in}{1.389375in}}%
\pgfpathlineto{\pgfqpoint{2.971487in}{1.377407in}}%
\pgfpathlineto{\pgfqpoint{2.959697in}{1.364159in}}%
\pgfpathlineto{\pgfqpoint{2.957787in}{1.350911in}}%
\pgfpathlineto{\pgfqpoint{2.950967in}{1.339453in}}%
\pgfpathlineto{\pgfqpoint{2.949543in}{1.337663in}}%
\pgfpathlineto{\pgfqpoint{2.946309in}{1.324415in}}%
\pgfpathlineto{\pgfqpoint{2.934296in}{1.311167in}}%
\pgfpathlineto{\pgfqpoint{2.931081in}{1.297919in}}%
\pgfpathlineto{\pgfqpoint{2.925683in}{1.291457in}}%
\pgfpathlineto{\pgfqpoint{2.918455in}{1.284671in}}%
\pgfpathlineto{\pgfqpoint{2.913292in}{1.271423in}}%
\pgfpathlineto{\pgfqpoint{2.904049in}{1.258175in}}%
\pgfpathlineto{\pgfqpoint{2.900400in}{1.247631in}}%
\pgfpathlineto{\pgfqpoint{2.899090in}{1.244927in}}%
\pgfpathlineto{\pgfqpoint{2.889367in}{1.231679in}}%
\pgfpathlineto{\pgfqpoint{2.882486in}{1.218431in}}%
\pgfpathlineto{\pgfqpoint{2.875117in}{1.208796in}}%
\pgfpathlineto{\pgfqpoint{2.871719in}{1.205183in}}%
\pgfpathlineto{\pgfqpoint{2.858646in}{1.191935in}}%
\pgfpathlineto{\pgfqpoint{2.851936in}{1.178687in}}%
\pgfpathlineto{\pgfqpoint{2.849834in}{1.176038in}}%
\pgfpathlineto{\pgfqpoint{2.839297in}{1.165439in}}%
\pgfpathlineto{\pgfqpoint{2.836129in}{1.152191in}}%
\pgfpathlineto{\pgfqpoint{2.824629in}{1.138943in}}%
\pgfpathlineto{\pgfqpoint{2.824551in}{1.138745in}}%
\pgfpathlineto{\pgfqpoint{2.818205in}{1.125695in}}%
\pgfpathlineto{\pgfqpoint{2.800097in}{1.112447in}}%
\pgfpathlineto{\pgfqpoint{2.799268in}{1.109721in}}%
\pgfpathlineto{\pgfqpoint{2.795226in}{1.099199in}}%
\pgfpathlineto{\pgfqpoint{2.774258in}{1.085951in}}%
\pgfpathlineto{\pgfqpoint{2.773984in}{1.084972in}}%
\pgfpathlineto{\pgfqpoint{2.769672in}{1.072703in}}%
\pgfpathlineto{\pgfqpoint{2.751503in}{1.059455in}}%
\pgfpathlineto{\pgfqpoint{2.748701in}{1.047655in}}%
\pgfpathlineto{\pgfqpoint{2.748272in}{1.046207in}}%
\pgfpathlineto{\pgfqpoint{2.727939in}{1.032959in}}%
\pgfpathlineto{\pgfqpoint{2.723418in}{1.022717in}}%
\pgfpathlineto{\pgfqpoint{2.721768in}{1.019711in}}%
\pgfpathlineto{\pgfqpoint{2.698135in}{1.007461in}}%
\pgfpathlineto{\pgfqpoint{2.696134in}{1.006463in}}%
\pgfpathlineto{\pgfqpoint{2.685737in}{0.993215in}}%
\pgfpathlineto{\pgfqpoint{2.672852in}{0.982549in}}%
\pgfpathlineto{\pgfqpoint{2.669618in}{0.979967in}}%
\pgfpathlineto{\pgfqpoint{2.660066in}{0.966719in}}%
\pgfpathlineto{\pgfqpoint{2.647569in}{0.955559in}}%
\pgfpathlineto{\pgfqpoint{2.645105in}{0.953471in}}%
\pgfpathlineto{\pgfqpoint{2.628914in}{0.940223in}}%
\pgfpathlineto{\pgfqpoint{2.622285in}{0.936058in}}%
\pgfpathlineto{\pgfqpoint{2.607300in}{0.926975in}}%
\pgfpathlineto{\pgfqpoint{2.597002in}{0.921264in}}%
\pgfpathlineto{\pgfqpoint{2.581438in}{0.913727in}}%
\pgfpathlineto{\pgfqpoint{2.571719in}{0.903943in}}%
\pgfpathlineto{\pgfqpoint{2.568098in}{0.900479in}}%
\pgfpathlineto{\pgfqpoint{2.546436in}{0.890784in}}%
\pgfpathlineto{\pgfqpoint{2.537743in}{0.887231in}}%
\pgfpathlineto{\pgfqpoint{2.531149in}{0.873983in}}%
\pgfpathlineto{\pgfqpoint{2.521153in}{0.869236in}}%
\pgfpathlineto{\pgfqpoint{2.502188in}{0.860735in}}%
\pgfpathlineto{\pgfqpoint{2.496986in}{0.847487in}}%
\pgfpathlineto{\pgfqpoint{2.495870in}{0.847100in}}%
\pgfpathlineto{\pgfqpoint{2.470586in}{0.839033in}}%
\pgfpathlineto{\pgfqpoint{2.470586in}{0.834239in}}%
\pgfpathlineto{\pgfqpoint{2.470586in}{0.820991in}}%
\pgfpathlineto{\pgfqpoint{2.470586in}{0.807743in}}%
\pgfpathlineto{\pgfqpoint{2.470586in}{0.794495in}}%
\pgfpathlineto{\pgfqpoint{2.470586in}{0.781247in}}%
\pgfpathlineto{\pgfqpoint{2.470586in}{0.767999in}}%
\pgfpathlineto{\pgfqpoint{2.470586in}{0.754751in}}%
\pgfpathlineto{\pgfqpoint{2.470586in}{0.741503in}}%
\pgfpathlineto{\pgfqpoint{2.470586in}{0.728255in}}%
\pgfpathlineto{\pgfqpoint{2.470586in}{0.715007in}}%
\pgfpathlineto{\pgfqpoint{2.470586in}{0.701759in}}%
\pgfpathlineto{\pgfqpoint{2.470586in}{0.688511in}}%
\pgfpathlineto{\pgfqpoint{2.470586in}{0.675263in}}%
\pgfpathlineto{\pgfqpoint{2.470586in}{0.662015in}}%
\pgfpathlineto{\pgfqpoint{2.470586in}{0.648767in}}%
\pgfpathlineto{\pgfqpoint{2.470586in}{0.635519in}}%
\pgfpathlineto{\pgfqpoint{2.470586in}{0.622271in}}%
\pgfpathlineto{\pgfqpoint{2.470586in}{0.609023in}}%
\pgfpathlineto{\pgfqpoint{2.470586in}{0.599811in}}%
\pgfpathlineto{\pgfqpoint{2.495870in}{0.597305in}}%
\pgfpathlineto{\pgfqpoint{2.512340in}{0.595775in}}%
\pgfpathlineto{\pgfqpoint{2.521153in}{0.595627in}}%
\pgfpathclose%
\pgfpathmoveto{\pgfqpoint{4.240408in}{1.071413in}}%
\pgfpathlineto{\pgfqpoint{4.265692in}{1.066254in}}%
\pgfpathlineto{\pgfqpoint{4.290975in}{1.061804in}}%
\pgfpathlineto{\pgfqpoint{4.290975in}{1.072703in}}%
\pgfpathlineto{\pgfqpoint{4.290975in}{1.085951in}}%
\pgfpathlineto{\pgfqpoint{4.290975in}{1.099199in}}%
\pgfpathlineto{\pgfqpoint{4.290975in}{1.112447in}}%
\pgfpathlineto{\pgfqpoint{4.290975in}{1.125695in}}%
\pgfpathlineto{\pgfqpoint{4.290975in}{1.138943in}}%
\pgfpathlineto{\pgfqpoint{4.290975in}{1.152191in}}%
\pgfpathlineto{\pgfqpoint{4.290975in}{1.165439in}}%
\pgfpathlineto{\pgfqpoint{4.290975in}{1.178687in}}%
\pgfpathlineto{\pgfqpoint{4.290975in}{1.191935in}}%
\pgfpathlineto{\pgfqpoint{4.290975in}{1.205183in}}%
\pgfpathlineto{\pgfqpoint{4.290975in}{1.218431in}}%
\pgfpathlineto{\pgfqpoint{4.290975in}{1.224152in}}%
\pgfpathlineto{\pgfqpoint{4.278162in}{1.231679in}}%
\pgfpathlineto{\pgfqpoint{4.277002in}{1.244927in}}%
\pgfpathlineto{\pgfqpoint{4.265692in}{1.252948in}}%
\pgfpathlineto{\pgfqpoint{4.258476in}{1.258175in}}%
\pgfpathlineto{\pgfqpoint{4.253450in}{1.271423in}}%
\pgfpathlineto{\pgfqpoint{4.240408in}{1.278323in}}%
\pgfpathlineto{\pgfqpoint{4.228806in}{1.284671in}}%
\pgfpathlineto{\pgfqpoint{4.222459in}{1.297919in}}%
\pgfpathlineto{\pgfqpoint{4.215125in}{1.302241in}}%
\pgfpathlineto{\pgfqpoint{4.200026in}{1.311167in}}%
\pgfpathlineto{\pgfqpoint{4.197342in}{1.324415in}}%
\pgfpathlineto{\pgfqpoint{4.189842in}{1.331466in}}%
\pgfpathlineto{\pgfqpoint{4.183438in}{1.337663in}}%
\pgfpathlineto{\pgfqpoint{4.178395in}{1.350911in}}%
\pgfpathlineto{\pgfqpoint{4.164559in}{1.362220in}}%
\pgfpathlineto{\pgfqpoint{4.162302in}{1.364159in}}%
\pgfpathlineto{\pgfqpoint{4.154600in}{1.377407in}}%
\pgfpathlineto{\pgfqpoint{4.139455in}{1.390655in}}%
\pgfpathlineto{\pgfqpoint{4.139276in}{1.390858in}}%
\pgfpathlineto{\pgfqpoint{4.129357in}{1.403903in}}%
\pgfpathlineto{\pgfqpoint{4.120046in}{1.417151in}}%
\pgfpathlineto{\pgfqpoint{4.113993in}{1.424526in}}%
\pgfpathlineto{\pgfqpoint{4.109666in}{1.430399in}}%
\pgfpathlineto{\pgfqpoint{4.103057in}{1.443647in}}%
\pgfpathlineto{\pgfqpoint{4.091626in}{1.456895in}}%
\pgfpathlineto{\pgfqpoint{4.088709in}{1.461374in}}%
\pgfpathlineto{\pgfqpoint{4.083769in}{1.470143in}}%
\pgfpathlineto{\pgfqpoint{4.068425in}{1.483391in}}%
\pgfpathlineto{\pgfqpoint{4.063493in}{1.496639in}}%
\pgfpathlineto{\pgfqpoint{4.063426in}{1.496705in}}%
\pgfpathlineto{\pgfqpoint{4.051465in}{1.509887in}}%
\pgfpathlineto{\pgfqpoint{4.049252in}{1.523135in}}%
\pgfpathlineto{\pgfqpoint{4.038143in}{1.535539in}}%
\pgfpathlineto{\pgfqpoint{4.037468in}{1.536383in}}%
\pgfpathlineto{\pgfqpoint{4.033633in}{1.549631in}}%
\pgfpathlineto{\pgfqpoint{4.019329in}{1.562879in}}%
\pgfpathlineto{\pgfqpoint{4.014767in}{1.576127in}}%
\pgfpathlineto{\pgfqpoint{4.012860in}{1.577934in}}%
\pgfpathlineto{\pgfqpoint{4.002449in}{1.589375in}}%
\pgfpathlineto{\pgfqpoint{3.998179in}{1.602623in}}%
\pgfpathlineto{\pgfqpoint{3.988094in}{1.615871in}}%
\pgfpathlineto{\pgfqpoint{3.987577in}{1.617445in}}%
\pgfpathlineto{\pgfqpoint{3.984594in}{1.629119in}}%
\pgfpathlineto{\pgfqpoint{3.975985in}{1.642367in}}%
\pgfpathlineto{\pgfqpoint{3.970263in}{1.655615in}}%
\pgfpathlineto{\pgfqpoint{3.962293in}{1.665873in}}%
\pgfpathlineto{\pgfqpoint{3.960390in}{1.668863in}}%
\pgfpathlineto{\pgfqpoint{3.953337in}{1.682111in}}%
\pgfpathlineto{\pgfqpoint{3.948118in}{1.695359in}}%
\pgfpathlineto{\pgfqpoint{3.941308in}{1.708607in}}%
\pgfpathlineto{\pgfqpoint{3.937010in}{1.718836in}}%
\pgfpathlineto{\pgfqpoint{3.936065in}{1.721855in}}%
\pgfpathlineto{\pgfqpoint{3.928952in}{1.735103in}}%
\pgfpathlineto{\pgfqpoint{3.923886in}{1.748351in}}%
\pgfpathlineto{\pgfqpoint{3.912829in}{1.761599in}}%
\pgfpathlineto{\pgfqpoint{3.911727in}{1.765308in}}%
\pgfpathlineto{\pgfqpoint{3.909735in}{1.774847in}}%
\pgfpathlineto{\pgfqpoint{3.901769in}{1.788095in}}%
\pgfpathlineto{\pgfqpoint{3.899577in}{1.801343in}}%
\pgfpathlineto{\pgfqpoint{3.891459in}{1.814591in}}%
\pgfpathlineto{\pgfqpoint{3.888884in}{1.827839in}}%
\pgfpathlineto{\pgfqpoint{3.886444in}{1.830672in}}%
\pgfpathlineto{\pgfqpoint{3.880138in}{1.841087in}}%
\pgfpathlineto{\pgfqpoint{3.877211in}{1.854335in}}%
\pgfpathlineto{\pgfqpoint{3.868088in}{1.867583in}}%
\pgfpathlineto{\pgfqpoint{3.864412in}{1.880831in}}%
\pgfpathlineto{\pgfqpoint{3.861161in}{1.886538in}}%
\pgfpathlineto{\pgfqpoint{3.858191in}{1.894079in}}%
\pgfpathlineto{\pgfqpoint{3.855781in}{1.907327in}}%
\pgfpathlineto{\pgfqpoint{3.849612in}{1.920575in}}%
\pgfpathlineto{\pgfqpoint{3.845200in}{1.933823in}}%
\pgfpathlineto{\pgfqpoint{3.837426in}{1.947071in}}%
\pgfpathlineto{\pgfqpoint{3.835878in}{1.950115in}}%
\pgfpathlineto{\pgfqpoint{3.832651in}{1.960319in}}%
\pgfpathlineto{\pgfqpoint{3.829372in}{1.973567in}}%
\pgfpathlineto{\pgfqpoint{3.825110in}{1.986815in}}%
\pgfpathlineto{\pgfqpoint{3.821760in}{2.000063in}}%
\pgfpathlineto{\pgfqpoint{3.815991in}{2.013311in}}%
\pgfpathlineto{\pgfqpoint{3.811407in}{2.026559in}}%
\pgfpathlineto{\pgfqpoint{3.810594in}{2.027642in}}%
\pgfpathlineto{\pgfqpoint{3.805272in}{2.039807in}}%
\pgfpathlineto{\pgfqpoint{3.802740in}{2.053055in}}%
\pgfpathlineto{\pgfqpoint{3.796181in}{2.066303in}}%
\pgfpathlineto{\pgfqpoint{3.794444in}{2.079551in}}%
\pgfpathlineto{\pgfqpoint{3.787916in}{2.092799in}}%
\pgfpathlineto{\pgfqpoint{3.785311in}{2.105411in}}%
\pgfpathlineto{\pgfqpoint{3.785241in}{2.106047in}}%
\pgfpathlineto{\pgfqpoint{3.779757in}{2.119295in}}%
\pgfpathlineto{\pgfqpoint{3.777992in}{2.132543in}}%
\pgfpathlineto{\pgfqpoint{3.760028in}{2.132543in}}%
\pgfpathlineto{\pgfqpoint{3.754889in}{2.132543in}}%
\pgfpathlineto{\pgfqpoint{3.756556in}{2.119295in}}%
\pgfpathlineto{\pgfqpoint{3.760028in}{2.111227in}}%
\pgfpathlineto{\pgfqpoint{3.762986in}{2.106047in}}%
\pgfpathlineto{\pgfqpoint{3.764975in}{2.092799in}}%
\pgfpathlineto{\pgfqpoint{3.770628in}{2.079551in}}%
\pgfpathlineto{\pgfqpoint{3.772178in}{2.066303in}}%
\pgfpathlineto{\pgfqpoint{3.778483in}{2.053055in}}%
\pgfpathlineto{\pgfqpoint{3.780962in}{2.039807in}}%
\pgfpathlineto{\pgfqpoint{3.785311in}{2.029658in}}%
\pgfpathlineto{\pgfqpoint{3.787048in}{2.026559in}}%
\pgfpathlineto{\pgfqpoint{3.790623in}{2.013311in}}%
\pgfpathlineto{\pgfqpoint{3.795762in}{2.000063in}}%
\pgfpathlineto{\pgfqpoint{3.799018in}{1.986815in}}%
\pgfpathlineto{\pgfqpoint{3.803305in}{1.973567in}}%
\pgfpathlineto{\pgfqpoint{3.806850in}{1.960319in}}%
\pgfpathlineto{\pgfqpoint{3.810594in}{1.947264in}}%
\pgfpathlineto{\pgfqpoint{3.810666in}{1.947071in}}%
\pgfpathlineto{\pgfqpoint{3.817591in}{1.933823in}}%
\pgfpathlineto{\pgfqpoint{3.821418in}{1.920575in}}%
\pgfpathlineto{\pgfqpoint{3.827666in}{1.907327in}}%
\pgfpathlineto{\pgfqpoint{3.830019in}{1.894079in}}%
\pgfpathlineto{\pgfqpoint{3.835858in}{1.880831in}}%
\pgfpathlineto{\pgfqpoint{3.835878in}{1.880708in}}%
\pgfpathlineto{\pgfqpoint{3.838509in}{1.867583in}}%
\pgfpathlineto{\pgfqpoint{3.847557in}{1.854335in}}%
\pgfpathlineto{\pgfqpoint{3.849824in}{1.841087in}}%
\pgfpathlineto{\pgfqpoint{3.857813in}{1.827839in}}%
\pgfpathlineto{\pgfqpoint{3.859392in}{1.814591in}}%
\pgfpathlineto{\pgfqpoint{3.861161in}{1.811393in}}%
\pgfpathlineto{\pgfqpoint{3.867354in}{1.801343in}}%
\pgfpathlineto{\pgfqpoint{3.869117in}{1.788095in}}%
\pgfpathlineto{\pgfqpoint{3.877468in}{1.774847in}}%
\pgfpathlineto{\pgfqpoint{3.880016in}{1.761599in}}%
\pgfpathlineto{\pgfqpoint{3.886444in}{1.751413in}}%
\pgfpathlineto{\pgfqpoint{3.888645in}{1.748351in}}%
\pgfpathlineto{\pgfqpoint{3.893801in}{1.735103in}}%
\pgfpathlineto{\pgfqpoint{3.900899in}{1.721855in}}%
\pgfpathlineto{\pgfqpoint{3.904942in}{1.708607in}}%
\pgfpathlineto{\pgfqpoint{3.910288in}{1.695359in}}%
\pgfpathlineto{\pgfqpoint{3.911727in}{1.691756in}}%
\pgfpathlineto{\pgfqpoint{3.915824in}{1.682111in}}%
\pgfpathlineto{\pgfqpoint{3.921582in}{1.668863in}}%
\pgfpathlineto{\pgfqpoint{3.930434in}{1.655615in}}%
\pgfpathlineto{\pgfqpoint{3.934197in}{1.642367in}}%
\pgfpathlineto{\pgfqpoint{3.937010in}{1.638010in}}%
\pgfpathlineto{\pgfqpoint{3.942993in}{1.629119in}}%
\pgfpathlineto{\pgfqpoint{3.945418in}{1.615871in}}%
\pgfpathlineto{\pgfqpoint{3.955222in}{1.602623in}}%
\pgfpathlineto{\pgfqpoint{3.957120in}{1.589375in}}%
\pgfpathlineto{\pgfqpoint{3.962293in}{1.583987in}}%
\pgfpathlineto{\pgfqpoint{3.969882in}{1.576127in}}%
\pgfpathlineto{\pgfqpoint{3.971785in}{1.562879in}}%
\pgfpathlineto{\pgfqpoint{3.985066in}{1.549631in}}%
\pgfpathlineto{\pgfqpoint{3.987443in}{1.536383in}}%
\pgfpathlineto{\pgfqpoint{3.987577in}{1.536218in}}%
\pgfpathlineto{\pgfqpoint{3.998116in}{1.523135in}}%
\pgfpathlineto{\pgfqpoint{3.997654in}{1.509887in}}%
\pgfpathlineto{\pgfqpoint{4.009291in}{1.496639in}}%
\pgfpathlineto{\pgfqpoint{4.012593in}{1.483391in}}%
\pgfpathlineto{\pgfqpoint{4.012860in}{1.483103in}}%
\pgfpathlineto{\pgfqpoint{4.025042in}{1.470143in}}%
\pgfpathlineto{\pgfqpoint{4.031612in}{1.456895in}}%
\pgfpathlineto{\pgfqpoint{4.038143in}{1.446787in}}%
\pgfpathlineto{\pgfqpoint{4.040221in}{1.443647in}}%
\pgfpathlineto{\pgfqpoint{4.045506in}{1.430399in}}%
\pgfpathlineto{\pgfqpoint{4.052078in}{1.417151in}}%
\pgfpathlineto{\pgfqpoint{4.061852in}{1.403903in}}%
\pgfpathlineto{\pgfqpoint{4.063426in}{1.397101in}}%
\pgfpathlineto{\pgfqpoint{4.064985in}{1.390655in}}%
\pgfpathlineto{\pgfqpoint{4.079894in}{1.377407in}}%
\pgfpathlineto{\pgfqpoint{4.080308in}{1.364159in}}%
\pgfpathlineto{\pgfqpoint{4.088709in}{1.357564in}}%
\pgfpathlineto{\pgfqpoint{4.096425in}{1.350911in}}%
\pgfpathlineto{\pgfqpoint{4.091317in}{1.337663in}}%
\pgfpathlineto{\pgfqpoint{4.106474in}{1.324415in}}%
\pgfpathlineto{\pgfqpoint{4.090272in}{1.311167in}}%
\pgfpathlineto{\pgfqpoint{4.113993in}{1.300573in}}%
\pgfpathlineto{\pgfqpoint{4.119181in}{1.297919in}}%
\pgfpathlineto{\pgfqpoint{4.113993in}{1.288508in}}%
\pgfpathlineto{\pgfqpoint{4.109743in}{1.284671in}}%
\pgfpathlineto{\pgfqpoint{4.113993in}{1.282818in}}%
\pgfpathlineto{\pgfqpoint{4.136066in}{1.271423in}}%
\pgfpathlineto{\pgfqpoint{4.113993in}{1.260887in}}%
\pgfpathlineto{\pgfqpoint{4.088709in}{1.259338in}}%
\pgfpathlineto{\pgfqpoint{4.063426in}{1.263068in}}%
\pgfpathlineto{\pgfqpoint{4.038143in}{1.270202in}}%
\pgfpathlineto{\pgfqpoint{4.035237in}{1.271423in}}%
\pgfpathlineto{\pgfqpoint{4.028257in}{1.284671in}}%
\pgfpathlineto{\pgfqpoint{4.012860in}{1.290632in}}%
\pgfpathlineto{\pgfqpoint{3.999867in}{1.297919in}}%
\pgfpathlineto{\pgfqpoint{3.999541in}{1.311167in}}%
\pgfpathlineto{\pgfqpoint{3.987577in}{1.314312in}}%
\pgfpathlineto{\pgfqpoint{3.962293in}{1.324104in}}%
\pgfpathlineto{\pgfqpoint{3.961683in}{1.324415in}}%
\pgfpathlineto{\pgfqpoint{3.955931in}{1.337663in}}%
\pgfpathlineto{\pgfqpoint{3.937010in}{1.345344in}}%
\pgfpathlineto{\pgfqpoint{3.925964in}{1.350911in}}%
\pgfpathlineto{\pgfqpoint{3.921531in}{1.364159in}}%
\pgfpathlineto{\pgfqpoint{3.911727in}{1.369555in}}%
\pgfpathlineto{\pgfqpoint{3.900187in}{1.377407in}}%
\pgfpathlineto{\pgfqpoint{3.896817in}{1.390655in}}%
\pgfpathlineto{\pgfqpoint{3.886444in}{1.396195in}}%
\pgfpathlineto{\pgfqpoint{3.874599in}{1.403903in}}%
\pgfpathlineto{\pgfqpoint{3.869560in}{1.417151in}}%
\pgfpathlineto{\pgfqpoint{3.861161in}{1.422263in}}%
\pgfpathlineto{\pgfqpoint{3.849549in}{1.430399in}}%
\pgfpathlineto{\pgfqpoint{3.843026in}{1.443647in}}%
\pgfpathlineto{\pgfqpoint{3.835878in}{1.449489in}}%
\pgfpathlineto{\pgfqpoint{3.827799in}{1.456895in}}%
\pgfpathlineto{\pgfqpoint{3.821446in}{1.470143in}}%
\pgfpathlineto{\pgfqpoint{3.810594in}{1.481783in}}%
\pgfpathlineto{\pgfqpoint{3.809221in}{1.483391in}}%
\pgfpathlineto{\pgfqpoint{3.803954in}{1.496639in}}%
\pgfpathlineto{\pgfqpoint{3.792514in}{1.509887in}}%
\pgfpathlineto{\pgfqpoint{3.785311in}{1.518851in}}%
\pgfpathlineto{\pgfqpoint{3.782214in}{1.523135in}}%
\pgfpathlineto{\pgfqpoint{3.772205in}{1.536383in}}%
\pgfpathlineto{\pgfqpoint{3.761776in}{1.549631in}}%
\pgfpathlineto{\pgfqpoint{3.760028in}{1.553287in}}%
\pgfpathlineto{\pgfqpoint{3.755563in}{1.562879in}}%
\pgfpathlineto{\pgfqpoint{3.746744in}{1.576127in}}%
\pgfpathlineto{\pgfqpoint{3.740964in}{1.589375in}}%
\pgfpathlineto{\pgfqpoint{3.734745in}{1.596524in}}%
\pgfpathlineto{\pgfqpoint{3.729796in}{1.602623in}}%
\pgfpathlineto{\pgfqpoint{3.725183in}{1.615871in}}%
\pgfpathlineto{\pgfqpoint{3.713716in}{1.629119in}}%
\pgfpathlineto{\pgfqpoint{3.710462in}{1.642367in}}%
\pgfpathlineto{\pgfqpoint{3.709462in}{1.643581in}}%
\pgfpathlineto{\pgfqpoint{3.700008in}{1.655615in}}%
\pgfpathlineto{\pgfqpoint{3.696971in}{1.668863in}}%
\pgfpathlineto{\pgfqpoint{3.687290in}{1.682111in}}%
\pgfpathlineto{\pgfqpoint{3.684731in}{1.695359in}}%
\pgfpathlineto{\pgfqpoint{3.684179in}{1.695977in}}%
\pgfpathlineto{\pgfqpoint{3.673118in}{1.708607in}}%
\pgfpathlineto{\pgfqpoint{3.669235in}{1.721855in}}%
\pgfpathlineto{\pgfqpoint{3.659935in}{1.735103in}}%
\pgfpathlineto{\pgfqpoint{3.658895in}{1.739253in}}%
\pgfpathlineto{\pgfqpoint{3.656373in}{1.748351in}}%
\pgfpathlineto{\pgfqpoint{3.648912in}{1.761599in}}%
\pgfpathlineto{\pgfqpoint{3.645714in}{1.774847in}}%
\pgfpathlineto{\pgfqpoint{3.637966in}{1.788095in}}%
\pgfpathlineto{\pgfqpoint{3.633612in}{1.800679in}}%
\pgfpathlineto{\pgfqpoint{3.633352in}{1.801343in}}%
\pgfpathlineto{\pgfqpoint{3.626234in}{1.814591in}}%
\pgfpathlineto{\pgfqpoint{3.620912in}{1.827839in}}%
\pgfpathlineto{\pgfqpoint{3.616174in}{1.841087in}}%
\pgfpathlineto{\pgfqpoint{3.611107in}{1.854335in}}%
\pgfpathlineto{\pgfqpoint{3.608329in}{1.862400in}}%
\pgfpathlineto{\pgfqpoint{3.606207in}{1.867583in}}%
\pgfpathlineto{\pgfqpoint{3.599698in}{1.880831in}}%
\pgfpathlineto{\pgfqpoint{3.595721in}{1.894079in}}%
\pgfpathlineto{\pgfqpoint{3.588357in}{1.907327in}}%
\pgfpathlineto{\pgfqpoint{3.585233in}{1.920575in}}%
\pgfpathlineto{\pgfqpoint{3.583046in}{1.925149in}}%
\pgfpathlineto{\pgfqpoint{3.578136in}{1.933823in}}%
\pgfpathlineto{\pgfqpoint{3.576016in}{1.947071in}}%
\pgfpathlineto{\pgfqpoint{3.569658in}{1.960319in}}%
\pgfpathlineto{\pgfqpoint{3.567451in}{1.973567in}}%
\pgfpathlineto{\pgfqpoint{3.560069in}{1.986815in}}%
\pgfpathlineto{\pgfqpoint{3.557777in}{2.000063in}}%
\pgfpathlineto{\pgfqpoint{3.557763in}{2.000091in}}%
\pgfpathlineto{\pgfqpoint{3.549635in}{2.013311in}}%
\pgfpathlineto{\pgfqpoint{3.547384in}{2.026559in}}%
\pgfpathlineto{\pgfqpoint{3.541195in}{2.039807in}}%
\pgfpathlineto{\pgfqpoint{3.538947in}{2.053055in}}%
\pgfpathlineto{\pgfqpoint{3.533089in}{2.066303in}}%
\pgfpathlineto{\pgfqpoint{3.532480in}{2.069610in}}%
\pgfpathlineto{\pgfqpoint{3.529994in}{2.079551in}}%
\pgfpathlineto{\pgfqpoint{3.523164in}{2.092799in}}%
\pgfpathlineto{\pgfqpoint{3.519471in}{2.106047in}}%
\pgfpathlineto{\pgfqpoint{3.514749in}{2.119295in}}%
\pgfpathlineto{\pgfqpoint{3.511745in}{2.132543in}}%
\pgfpathlineto{\pgfqpoint{3.507196in}{2.132543in}}%
\pgfpathlineto{\pgfqpoint{3.487284in}{2.132543in}}%
\pgfpathlineto{\pgfqpoint{3.489901in}{2.119295in}}%
\pgfpathlineto{\pgfqpoint{3.494257in}{2.106047in}}%
\pgfpathlineto{\pgfqpoint{3.498113in}{2.092799in}}%
\pgfpathlineto{\pgfqpoint{3.507196in}{2.080203in}}%
\pgfpathlineto{\pgfqpoint{3.507451in}{2.079551in}}%
\pgfpathlineto{\pgfqpoint{3.509423in}{2.066303in}}%
\pgfpathlineto{\pgfqpoint{3.514529in}{2.053055in}}%
\pgfpathlineto{\pgfqpoint{3.516295in}{2.039807in}}%
\pgfpathlineto{\pgfqpoint{3.521979in}{2.026559in}}%
\pgfpathlineto{\pgfqpoint{3.523795in}{2.013311in}}%
\pgfpathlineto{\pgfqpoint{3.532480in}{2.002470in}}%
\pgfpathlineto{\pgfqpoint{3.533553in}{2.000063in}}%
\pgfpathlineto{\pgfqpoint{3.535446in}{1.986815in}}%
\pgfpathlineto{\pgfqpoint{3.542198in}{1.973567in}}%
\pgfpathlineto{\pgfqpoint{3.544237in}{1.960319in}}%
\pgfpathlineto{\pgfqpoint{3.550278in}{1.947071in}}%
\pgfpathlineto{\pgfqpoint{3.551965in}{1.933823in}}%
\pgfpathlineto{\pgfqpoint{3.557763in}{1.925624in}}%
\pgfpathlineto{\pgfqpoint{3.559754in}{1.920575in}}%
\pgfpathlineto{\pgfqpoint{3.562220in}{1.907327in}}%
\pgfpathlineto{\pgfqpoint{3.569018in}{1.894079in}}%
\pgfpathlineto{\pgfqpoint{3.573168in}{1.880831in}}%
\pgfpathlineto{\pgfqpoint{3.579859in}{1.867583in}}%
\pgfpathlineto{\pgfqpoint{3.583046in}{1.860813in}}%
\pgfpathlineto{\pgfqpoint{3.584875in}{1.854335in}}%
\pgfpathlineto{\pgfqpoint{3.588818in}{1.841087in}}%
\pgfpathlineto{\pgfqpoint{3.592469in}{1.827839in}}%
\pgfpathlineto{\pgfqpoint{3.597357in}{1.814591in}}%
\pgfpathlineto{\pgfqpoint{3.605197in}{1.801343in}}%
\pgfpathlineto{\pgfqpoint{3.608329in}{1.794402in}}%
\pgfpathlineto{\pgfqpoint{3.610093in}{1.788095in}}%
\pgfpathlineto{\pgfqpoint{3.616739in}{1.774847in}}%
\pgfpathlineto{\pgfqpoint{3.619325in}{1.761599in}}%
\pgfpathlineto{\pgfqpoint{3.625616in}{1.748351in}}%
\pgfpathlineto{\pgfqpoint{3.629214in}{1.735103in}}%
\pgfpathlineto{\pgfqpoint{3.633612in}{1.730024in}}%
\pgfpathlineto{\pgfqpoint{3.638571in}{1.721855in}}%
\pgfpathlineto{\pgfqpoint{3.641734in}{1.708607in}}%
\pgfpathlineto{\pgfqpoint{3.652448in}{1.695359in}}%
\pgfpathlineto{\pgfqpoint{3.655452in}{1.682111in}}%
\pgfpathlineto{\pgfqpoint{3.658895in}{1.677431in}}%
\pgfpathlineto{\pgfqpoint{3.663758in}{1.668863in}}%
\pgfpathlineto{\pgfqpoint{3.666239in}{1.655615in}}%
\pgfpathlineto{\pgfqpoint{3.674997in}{1.642367in}}%
\pgfpathlineto{\pgfqpoint{3.678135in}{1.629119in}}%
\pgfpathlineto{\pgfqpoint{3.684179in}{1.622559in}}%
\pgfpathlineto{\pgfqpoint{3.689093in}{1.615871in}}%
\pgfpathlineto{\pgfqpoint{3.693799in}{1.602623in}}%
\pgfpathlineto{\pgfqpoint{3.702960in}{1.589375in}}%
\pgfpathlineto{\pgfqpoint{3.709402in}{1.576127in}}%
\pgfpathlineto{\pgfqpoint{3.709462in}{1.576030in}}%
\pgfpathlineto{\pgfqpoint{3.715885in}{1.562879in}}%
\pgfpathlineto{\pgfqpoint{3.720893in}{1.549631in}}%
\pgfpathlineto{\pgfqpoint{3.730122in}{1.536383in}}%
\pgfpathlineto{\pgfqpoint{3.734745in}{1.530743in}}%
\pgfpathlineto{\pgfqpoint{3.740021in}{1.523135in}}%
\pgfpathlineto{\pgfqpoint{3.747628in}{1.509887in}}%
\pgfpathlineto{\pgfqpoint{3.758723in}{1.496639in}}%
\pgfpathlineto{\pgfqpoint{3.760028in}{1.493004in}}%
\pgfpathlineto{\pgfqpoint{3.762968in}{1.483391in}}%
\pgfpathlineto{\pgfqpoint{3.772185in}{1.470143in}}%
\pgfpathlineto{\pgfqpoint{3.777600in}{1.456895in}}%
\pgfpathlineto{\pgfqpoint{3.785311in}{1.449290in}}%
\pgfpathlineto{\pgfqpoint{3.790690in}{1.443647in}}%
\pgfpathlineto{\pgfqpoint{3.795923in}{1.430399in}}%
\pgfpathlineto{\pgfqpoint{3.810594in}{1.418870in}}%
\pgfpathlineto{\pgfqpoint{3.812781in}{1.417151in}}%
\pgfpathlineto{\pgfqpoint{3.817844in}{1.403903in}}%
\pgfpathlineto{\pgfqpoint{3.833453in}{1.390655in}}%
\pgfpathlineto{\pgfqpoint{3.835878in}{1.383153in}}%
\pgfpathlineto{\pgfqpoint{3.837588in}{1.377407in}}%
\pgfpathlineto{\pgfqpoint{3.851799in}{1.364159in}}%
\pgfpathlineto{\pgfqpoint{3.857025in}{1.350911in}}%
\pgfpathlineto{\pgfqpoint{3.861161in}{1.348106in}}%
\pgfpathlineto{\pgfqpoint{3.876856in}{1.337663in}}%
\pgfpathlineto{\pgfqpoint{3.885009in}{1.324415in}}%
\pgfpathlineto{\pgfqpoint{3.886444in}{1.323292in}}%
\pgfpathlineto{\pgfqpoint{3.902774in}{1.311167in}}%
\pgfpathlineto{\pgfqpoint{3.910313in}{1.297919in}}%
\pgfpathlineto{\pgfqpoint{3.911727in}{1.296434in}}%
\pgfpathlineto{\pgfqpoint{3.923291in}{1.284671in}}%
\pgfpathlineto{\pgfqpoint{3.933556in}{1.271423in}}%
\pgfpathlineto{\pgfqpoint{3.937010in}{1.268786in}}%
\pgfpathlineto{\pgfqpoint{3.951598in}{1.258175in}}%
\pgfpathlineto{\pgfqpoint{3.962293in}{1.248858in}}%
\pgfpathlineto{\pgfqpoint{3.966885in}{1.244927in}}%
\pgfpathlineto{\pgfqpoint{3.984582in}{1.231679in}}%
\pgfpathlineto{\pgfqpoint{3.987577in}{1.229687in}}%
\pgfpathlineto{\pgfqpoint{4.005075in}{1.218431in}}%
\pgfpathlineto{\pgfqpoint{4.012860in}{1.209545in}}%
\pgfpathlineto{\pgfqpoint{4.017019in}{1.205183in}}%
\pgfpathlineto{\pgfqpoint{4.038013in}{1.191935in}}%
\pgfpathlineto{\pgfqpoint{4.038143in}{1.191767in}}%
\pgfpathlineto{\pgfqpoint{4.048827in}{1.178687in}}%
\pgfpathlineto{\pgfqpoint{4.063426in}{1.170471in}}%
\pgfpathlineto{\pgfqpoint{4.073090in}{1.165439in}}%
\pgfpathlineto{\pgfqpoint{4.087235in}{1.152191in}}%
\pgfpathlineto{\pgfqpoint{4.088709in}{1.151649in}}%
\pgfpathlineto{\pgfqpoint{4.113993in}{1.143133in}}%
\pgfpathlineto{\pgfqpoint{4.127624in}{1.138943in}}%
\pgfpathlineto{\pgfqpoint{4.136891in}{1.125695in}}%
\pgfpathlineto{\pgfqpoint{4.139276in}{1.124771in}}%
\pgfpathlineto{\pgfqpoint{4.164559in}{1.116231in}}%
\pgfpathlineto{\pgfqpoint{4.177530in}{1.112447in}}%
\pgfpathlineto{\pgfqpoint{4.180881in}{1.099199in}}%
\pgfpathlineto{\pgfqpoint{4.189842in}{1.095745in}}%
\pgfpathlineto{\pgfqpoint{4.215125in}{1.087175in}}%
\pgfpathlineto{\pgfqpoint{4.219278in}{1.085951in}}%
\pgfpathlineto{\pgfqpoint{4.234961in}{1.072703in}}%
\pgfpathlineto{\pgfqpoint{4.240408in}{1.071413in}}%
\pgfpathclose%
\pgfpathmoveto{\pgfqpoint{2.495870in}{1.095398in}}%
\pgfpathlineto{\pgfqpoint{2.507477in}{1.099199in}}%
\pgfpathlineto{\pgfqpoint{2.510001in}{1.112447in}}%
\pgfpathlineto{\pgfqpoint{2.521153in}{1.116219in}}%
\pgfpathlineto{\pgfqpoint{2.546156in}{1.125695in}}%
\pgfpathlineto{\pgfqpoint{2.546436in}{1.126319in}}%
\pgfpathlineto{\pgfqpoint{2.551836in}{1.138943in}}%
\pgfpathlineto{\pgfqpoint{2.571719in}{1.144563in}}%
\pgfpathlineto{\pgfqpoint{2.595241in}{1.152191in}}%
\pgfpathlineto{\pgfqpoint{2.597002in}{1.155226in}}%
\pgfpathlineto{\pgfqpoint{2.602872in}{1.165439in}}%
\pgfpathlineto{\pgfqpoint{2.622285in}{1.172697in}}%
\pgfpathlineto{\pgfqpoint{2.635937in}{1.178687in}}%
\pgfpathlineto{\pgfqpoint{2.639598in}{1.191935in}}%
\pgfpathlineto{\pgfqpoint{2.647569in}{1.195997in}}%
\pgfpathlineto{\pgfqpoint{2.662841in}{1.205183in}}%
\pgfpathlineto{\pgfqpoint{2.669916in}{1.218431in}}%
\pgfpathlineto{\pgfqpoint{2.672852in}{1.219884in}}%
\pgfpathlineto{\pgfqpoint{2.693140in}{1.231679in}}%
\pgfpathlineto{\pgfqpoint{2.698135in}{1.237451in}}%
\pgfpathlineto{\pgfqpoint{2.704912in}{1.244927in}}%
\pgfpathlineto{\pgfqpoint{2.723418in}{1.256607in}}%
\pgfpathlineto{\pgfqpoint{2.725649in}{1.258175in}}%
\pgfpathlineto{\pgfqpoint{2.742656in}{1.271423in}}%
\pgfpathlineto{\pgfqpoint{2.748701in}{1.278084in}}%
\pgfpathlineto{\pgfqpoint{2.754210in}{1.284671in}}%
\pgfpathlineto{\pgfqpoint{2.767467in}{1.297919in}}%
\pgfpathlineto{\pgfqpoint{2.773800in}{1.311167in}}%
\pgfpathlineto{\pgfqpoint{2.773984in}{1.311321in}}%
\pgfpathlineto{\pgfqpoint{2.789546in}{1.324415in}}%
\pgfpathlineto{\pgfqpoint{2.797363in}{1.337663in}}%
\pgfpathlineto{\pgfqpoint{2.799268in}{1.338859in}}%
\pgfpathlineto{\pgfqpoint{2.818006in}{1.350911in}}%
\pgfpathlineto{\pgfqpoint{2.821981in}{1.364159in}}%
\pgfpathlineto{\pgfqpoint{2.824551in}{1.366117in}}%
\pgfpathlineto{\pgfqpoint{2.838908in}{1.377407in}}%
\pgfpathlineto{\pgfqpoint{2.841679in}{1.390655in}}%
\pgfpathlineto{\pgfqpoint{2.849834in}{1.397597in}}%
\pgfpathlineto{\pgfqpoint{2.856992in}{1.403903in}}%
\pgfpathlineto{\pgfqpoint{2.861485in}{1.417151in}}%
\pgfpathlineto{\pgfqpoint{2.875117in}{1.428437in}}%
\pgfpathlineto{\pgfqpoint{2.877527in}{1.430399in}}%
\pgfpathlineto{\pgfqpoint{2.883200in}{1.443647in}}%
\pgfpathlineto{\pgfqpoint{2.896492in}{1.456895in}}%
\pgfpathlineto{\pgfqpoint{2.900400in}{1.465481in}}%
\pgfpathlineto{\pgfqpoint{2.902987in}{1.470143in}}%
\pgfpathlineto{\pgfqpoint{2.911582in}{1.483391in}}%
\pgfpathlineto{\pgfqpoint{2.916598in}{1.496639in}}%
\pgfpathlineto{\pgfqpoint{2.924624in}{1.509887in}}%
\pgfpathlineto{\pgfqpoint{2.925683in}{1.511548in}}%
\pgfpathlineto{\pgfqpoint{2.934316in}{1.523135in}}%
\pgfpathlineto{\pgfqpoint{2.942483in}{1.536383in}}%
\pgfpathlineto{\pgfqpoint{2.950967in}{1.547273in}}%
\pgfpathlineto{\pgfqpoint{2.953235in}{1.549631in}}%
\pgfpathlineto{\pgfqpoint{2.957449in}{1.562879in}}%
\pgfpathlineto{\pgfqpoint{2.966734in}{1.576127in}}%
\pgfpathlineto{\pgfqpoint{2.970145in}{1.589375in}}%
\pgfpathlineto{\pgfqpoint{2.976250in}{1.598029in}}%
\pgfpathlineto{\pgfqpoint{2.980127in}{1.602623in}}%
\pgfpathlineto{\pgfqpoint{2.983780in}{1.615871in}}%
\pgfpathlineto{\pgfqpoint{2.995131in}{1.629119in}}%
\pgfpathlineto{\pgfqpoint{2.997596in}{1.642367in}}%
\pgfpathlineto{\pgfqpoint{3.001533in}{1.648241in}}%
\pgfpathlineto{\pgfqpoint{3.007546in}{1.655615in}}%
\pgfpathlineto{\pgfqpoint{3.010637in}{1.668863in}}%
\pgfpathlineto{\pgfqpoint{3.018054in}{1.682111in}}%
\pgfpathlineto{\pgfqpoint{3.020956in}{1.695359in}}%
\pgfpathlineto{\pgfqpoint{3.026816in}{1.704506in}}%
\pgfpathlineto{\pgfqpoint{3.030262in}{1.708607in}}%
\pgfpathlineto{\pgfqpoint{3.036710in}{1.721855in}}%
\pgfpathlineto{\pgfqpoint{3.043434in}{1.735103in}}%
\pgfpathlineto{\pgfqpoint{3.048285in}{1.748351in}}%
\pgfpathlineto{\pgfqpoint{3.051449in}{1.761599in}}%
\pgfpathlineto{\pgfqpoint{3.052099in}{1.763482in}}%
\pgfpathlineto{\pgfqpoint{3.057803in}{1.774847in}}%
\pgfpathlineto{\pgfqpoint{3.062644in}{1.788095in}}%
\pgfpathlineto{\pgfqpoint{3.069499in}{1.801343in}}%
\pgfpathlineto{\pgfqpoint{3.072603in}{1.814591in}}%
\pgfpathlineto{\pgfqpoint{3.077382in}{1.823305in}}%
\pgfpathlineto{\pgfqpoint{3.081306in}{1.827839in}}%
\pgfpathlineto{\pgfqpoint{3.084491in}{1.841087in}}%
\pgfpathlineto{\pgfqpoint{3.091453in}{1.854335in}}%
\pgfpathlineto{\pgfqpoint{3.093342in}{1.867583in}}%
\pgfpathlineto{\pgfqpoint{3.099390in}{1.880831in}}%
\pgfpathlineto{\pgfqpoint{3.101539in}{1.894079in}}%
\pgfpathlineto{\pgfqpoint{3.102666in}{1.896167in}}%
\pgfpathlineto{\pgfqpoint{3.111689in}{1.907327in}}%
\pgfpathlineto{\pgfqpoint{3.115200in}{1.920575in}}%
\pgfpathlineto{\pgfqpoint{3.120247in}{1.933823in}}%
\pgfpathlineto{\pgfqpoint{3.122574in}{1.947071in}}%
\pgfpathlineto{\pgfqpoint{3.126419in}{1.960319in}}%
\pgfpathlineto{\pgfqpoint{3.127949in}{1.967065in}}%
\pgfpathlineto{\pgfqpoint{3.130709in}{1.973567in}}%
\pgfpathlineto{\pgfqpoint{3.137405in}{1.986815in}}%
\pgfpathlineto{\pgfqpoint{3.143138in}{2.000063in}}%
\pgfpathlineto{\pgfqpoint{3.146649in}{2.013311in}}%
\pgfpathlineto{\pgfqpoint{3.151418in}{2.026559in}}%
\pgfpathlineto{\pgfqpoint{3.153232in}{2.039134in}}%
\pgfpathlineto{\pgfqpoint{3.153442in}{2.039807in}}%
\pgfpathlineto{\pgfqpoint{3.160420in}{2.053055in}}%
\pgfpathlineto{\pgfqpoint{3.162907in}{2.066303in}}%
\pgfpathlineto{\pgfqpoint{3.169635in}{2.079551in}}%
\pgfpathlineto{\pgfqpoint{3.171919in}{2.092799in}}%
\pgfpathlineto{\pgfqpoint{3.177408in}{2.106047in}}%
\pgfpathlineto{\pgfqpoint{3.178515in}{2.115046in}}%
\pgfpathlineto{\pgfqpoint{3.179741in}{2.119295in}}%
\pgfpathlineto{\pgfqpoint{3.187113in}{2.132543in}}%
\pgfpathlineto{\pgfqpoint{3.178515in}{2.132543in}}%
\pgfpathlineto{\pgfqpoint{3.163670in}{2.132543in}}%
\pgfpathlineto{\pgfqpoint{3.157141in}{2.119295in}}%
\pgfpathlineto{\pgfqpoint{3.154792in}{2.106047in}}%
\pgfpathlineto{\pgfqpoint{3.153232in}{2.103599in}}%
\pgfpathlineto{\pgfqpoint{3.148093in}{2.092799in}}%
\pgfpathlineto{\pgfqpoint{3.145764in}{2.079551in}}%
\pgfpathlineto{\pgfqpoint{3.139520in}{2.066303in}}%
\pgfpathlineto{\pgfqpoint{3.137068in}{2.053055in}}%
\pgfpathlineto{\pgfqpoint{3.131429in}{2.039807in}}%
\pgfpathlineto{\pgfqpoint{3.127949in}{2.027961in}}%
\pgfpathlineto{\pgfqpoint{3.127628in}{2.026559in}}%
\pgfpathlineto{\pgfqpoint{3.122463in}{2.013311in}}%
\pgfpathlineto{\pgfqpoint{3.118169in}{2.000063in}}%
\pgfpathlineto{\pgfqpoint{3.113385in}{1.986815in}}%
\pgfpathlineto{\pgfqpoint{3.107127in}{1.973567in}}%
\pgfpathlineto{\pgfqpoint{3.103034in}{1.960319in}}%
\pgfpathlineto{\pgfqpoint{3.102666in}{1.959575in}}%
\pgfpathlineto{\pgfqpoint{3.097613in}{1.947071in}}%
\pgfpathlineto{\pgfqpoint{3.095430in}{1.933823in}}%
\pgfpathlineto{\pgfqpoint{3.089132in}{1.920575in}}%
\pgfpathlineto{\pgfqpoint{3.086354in}{1.907327in}}%
\pgfpathlineto{\pgfqpoint{3.077382in}{1.894901in}}%
\pgfpathlineto{\pgfqpoint{3.076844in}{1.894079in}}%
\pgfpathlineto{\pgfqpoint{3.074575in}{1.880831in}}%
\pgfpathlineto{\pgfqpoint{3.067308in}{1.867583in}}%
\pgfpathlineto{\pgfqpoint{3.065116in}{1.854335in}}%
\pgfpathlineto{\pgfqpoint{3.057336in}{1.841087in}}%
\pgfpathlineto{\pgfqpoint{3.054054in}{1.827839in}}%
\pgfpathlineto{\pgfqpoint{3.052099in}{1.825025in}}%
\pgfpathlineto{\pgfqpoint{3.045864in}{1.814591in}}%
\pgfpathlineto{\pgfqpoint{3.041966in}{1.801343in}}%
\pgfpathlineto{\pgfqpoint{3.035170in}{1.788095in}}%
\pgfpathlineto{\pgfqpoint{3.029590in}{1.774847in}}%
\pgfpathlineto{\pgfqpoint{3.026816in}{1.768469in}}%
\pgfpathlineto{\pgfqpoint{3.024176in}{1.761599in}}%
\pgfpathlineto{\pgfqpoint{3.019064in}{1.748351in}}%
\pgfpathlineto{\pgfqpoint{3.014348in}{1.735103in}}%
\pgfpathlineto{\pgfqpoint{3.005668in}{1.721855in}}%
\pgfpathlineto{\pgfqpoint{3.001533in}{1.711505in}}%
\pgfpathlineto{\pgfqpoint{3.000486in}{1.708607in}}%
\pgfpathlineto{\pgfqpoint{2.989888in}{1.695359in}}%
\pgfpathlineto{\pgfqpoint{2.987293in}{1.682111in}}%
\pgfpathlineto{\pgfqpoint{2.977611in}{1.668863in}}%
\pgfpathlineto{\pgfqpoint{2.976250in}{1.662078in}}%
\pgfpathlineto{\pgfqpoint{2.975053in}{1.655615in}}%
\pgfpathlineto{\pgfqpoint{2.963892in}{1.642367in}}%
\pgfpathlineto{\pgfqpoint{2.961516in}{1.629119in}}%
\pgfpathlineto{\pgfqpoint{2.950967in}{1.617836in}}%
\pgfpathlineto{\pgfqpoint{2.949071in}{1.615871in}}%
\pgfpathlineto{\pgfqpoint{2.945374in}{1.602623in}}%
\pgfpathlineto{\pgfqpoint{2.934255in}{1.589375in}}%
\pgfpathlineto{\pgfqpoint{2.929527in}{1.576127in}}%
\pgfpathlineto{\pgfqpoint{2.925683in}{1.571120in}}%
\pgfpathlineto{\pgfqpoint{2.919272in}{1.562879in}}%
\pgfpathlineto{\pgfqpoint{2.913675in}{1.549631in}}%
\pgfpathlineto{\pgfqpoint{2.902964in}{1.536383in}}%
\pgfpathlineto{\pgfqpoint{2.900400in}{1.533127in}}%
\pgfpathlineto{\pgfqpoint{2.891702in}{1.523135in}}%
\pgfpathlineto{\pgfqpoint{2.883674in}{1.509887in}}%
\pgfpathlineto{\pgfqpoint{2.875117in}{1.499703in}}%
\pgfpathlineto{\pgfqpoint{2.872169in}{1.496639in}}%
\pgfpathlineto{\pgfqpoint{2.868144in}{1.483391in}}%
\pgfpathlineto{\pgfqpoint{2.855554in}{1.470143in}}%
\pgfpathlineto{\pgfqpoint{2.850789in}{1.456895in}}%
\pgfpathlineto{\pgfqpoint{2.849834in}{1.456132in}}%
\pgfpathlineto{\pgfqpoint{2.831834in}{1.443647in}}%
\pgfpathlineto{\pgfqpoint{2.828451in}{1.430399in}}%
\pgfpathlineto{\pgfqpoint{2.824551in}{1.427713in}}%
\pgfpathlineto{\pgfqpoint{2.806729in}{1.417151in}}%
\pgfpathlineto{\pgfqpoint{2.802888in}{1.403903in}}%
\pgfpathlineto{\pgfqpoint{2.799268in}{1.401380in}}%
\pgfpathlineto{\pgfqpoint{2.781277in}{1.390655in}}%
\pgfpathlineto{\pgfqpoint{2.779872in}{1.377407in}}%
\pgfpathlineto{\pgfqpoint{2.773984in}{1.373442in}}%
\pgfpathlineto{\pgfqpoint{2.758405in}{1.364159in}}%
\pgfpathlineto{\pgfqpoint{2.753175in}{1.350911in}}%
\pgfpathlineto{\pgfqpoint{2.748701in}{1.348290in}}%
\pgfpathlineto{\pgfqpoint{2.728670in}{1.337663in}}%
\pgfpathlineto{\pgfqpoint{2.723418in}{1.331629in}}%
\pgfpathlineto{\pgfqpoint{2.715423in}{1.324415in}}%
\pgfpathlineto{\pgfqpoint{2.699195in}{1.311167in}}%
\pgfpathlineto{\pgfqpoint{2.698135in}{1.309906in}}%
\pgfpathlineto{\pgfqpoint{2.684962in}{1.297919in}}%
\pgfpathlineto{\pgfqpoint{2.672852in}{1.284672in}}%
\pgfpathlineto{\pgfqpoint{2.672851in}{1.284671in}}%
\pgfpathlineto{\pgfqpoint{2.647569in}{1.273343in}}%
\pgfpathlineto{\pgfqpoint{2.640732in}{1.271423in}}%
\pgfpathlineto{\pgfqpoint{2.629863in}{1.258175in}}%
\pgfpathlineto{\pgfqpoint{2.622285in}{1.255735in}}%
\pgfpathlineto{\pgfqpoint{2.597002in}{1.250326in}}%
\pgfpathlineto{\pgfqpoint{2.571719in}{1.247551in}}%
\pgfpathlineto{\pgfqpoint{2.546436in}{1.247815in}}%
\pgfpathlineto{\pgfqpoint{2.521153in}{1.252366in}}%
\pgfpathlineto{\pgfqpoint{2.507224in}{1.258175in}}%
\pgfpathlineto{\pgfqpoint{2.521153in}{1.264341in}}%
\pgfpathlineto{\pgfqpoint{2.538469in}{1.271423in}}%
\pgfpathlineto{\pgfqpoint{2.521153in}{1.280577in}}%
\pgfpathlineto{\pgfqpoint{2.516228in}{1.284671in}}%
\pgfpathlineto{\pgfqpoint{2.521153in}{1.287132in}}%
\pgfpathlineto{\pgfqpoint{2.543046in}{1.297919in}}%
\pgfpathlineto{\pgfqpoint{2.538279in}{1.311167in}}%
\pgfpathlineto{\pgfqpoint{2.546436in}{1.316096in}}%
\pgfpathlineto{\pgfqpoint{2.560334in}{1.324415in}}%
\pgfpathlineto{\pgfqpoint{2.556961in}{1.337663in}}%
\pgfpathlineto{\pgfqpoint{2.569822in}{1.350911in}}%
\pgfpathlineto{\pgfqpoint{2.569508in}{1.364159in}}%
\pgfpathlineto{\pgfqpoint{2.571719in}{1.365979in}}%
\pgfpathlineto{\pgfqpoint{2.584811in}{1.377407in}}%
\pgfpathlineto{\pgfqpoint{2.594878in}{1.390655in}}%
\pgfpathlineto{\pgfqpoint{2.597002in}{1.393367in}}%
\pgfpathlineto{\pgfqpoint{2.604573in}{1.403903in}}%
\pgfpathlineto{\pgfqpoint{2.613937in}{1.417151in}}%
\pgfpathlineto{\pgfqpoint{2.616925in}{1.430399in}}%
\pgfpathlineto{\pgfqpoint{2.622285in}{1.437831in}}%
\pgfpathlineto{\pgfqpoint{2.626843in}{1.443647in}}%
\pgfpathlineto{\pgfqpoint{2.629175in}{1.456895in}}%
\pgfpathlineto{\pgfqpoint{2.642428in}{1.470143in}}%
\pgfpathlineto{\pgfqpoint{2.643944in}{1.483391in}}%
\pgfpathlineto{\pgfqpoint{2.647569in}{1.486424in}}%
\pgfpathlineto{\pgfqpoint{2.660272in}{1.496639in}}%
\pgfpathlineto{\pgfqpoint{2.661629in}{1.509887in}}%
\pgfpathlineto{\pgfqpoint{2.672852in}{1.519855in}}%
\pgfpathlineto{\pgfqpoint{2.676684in}{1.523135in}}%
\pgfpathlineto{\pgfqpoint{2.675868in}{1.536383in}}%
\pgfpathlineto{\pgfqpoint{2.688755in}{1.549631in}}%
\pgfpathlineto{\pgfqpoint{2.688743in}{1.562879in}}%
\pgfpathlineto{\pgfqpoint{2.698135in}{1.570790in}}%
\pgfpathlineto{\pgfqpoint{2.704269in}{1.576127in}}%
\pgfpathlineto{\pgfqpoint{2.706232in}{1.589375in}}%
\pgfpathlineto{\pgfqpoint{2.717535in}{1.602623in}}%
\pgfpathlineto{\pgfqpoint{2.718641in}{1.615871in}}%
\pgfpathlineto{\pgfqpoint{2.723418in}{1.622755in}}%
\pgfpathlineto{\pgfqpoint{2.727499in}{1.629119in}}%
\pgfpathlineto{\pgfqpoint{2.730784in}{1.642367in}}%
\pgfpathlineto{\pgfqpoint{2.739850in}{1.655615in}}%
\pgfpathlineto{\pgfqpoint{2.745774in}{1.668863in}}%
\pgfpathlineto{\pgfqpoint{2.748701in}{1.673799in}}%
\pgfpathlineto{\pgfqpoint{2.752962in}{1.682111in}}%
\pgfpathlineto{\pgfqpoint{2.760630in}{1.695359in}}%
\pgfpathlineto{\pgfqpoint{2.765348in}{1.708607in}}%
\pgfpathlineto{\pgfqpoint{2.772715in}{1.721855in}}%
\pgfpathlineto{\pgfqpoint{2.773984in}{1.727692in}}%
\pgfpathlineto{\pgfqpoint{2.775312in}{1.735103in}}%
\pgfpathlineto{\pgfqpoint{2.782871in}{1.748351in}}%
\pgfpathlineto{\pgfqpoint{2.785558in}{1.761599in}}%
\pgfpathlineto{\pgfqpoint{2.796340in}{1.774847in}}%
\pgfpathlineto{\pgfqpoint{2.797989in}{1.788095in}}%
\pgfpathlineto{\pgfqpoint{2.799268in}{1.789939in}}%
\pgfpathlineto{\pgfqpoint{2.806311in}{1.801343in}}%
\pgfpathlineto{\pgfqpoint{2.807128in}{1.814591in}}%
\pgfpathlineto{\pgfqpoint{2.816145in}{1.827839in}}%
\pgfpathlineto{\pgfqpoint{2.818222in}{1.841087in}}%
\pgfpathlineto{\pgfqpoint{2.824551in}{1.849245in}}%
\pgfpathlineto{\pgfqpoint{2.827945in}{1.854335in}}%
\pgfpathlineto{\pgfqpoint{2.830010in}{1.867583in}}%
\pgfpathlineto{\pgfqpoint{2.837376in}{1.880831in}}%
\pgfpathlineto{\pgfqpoint{2.840179in}{1.894079in}}%
\pgfpathlineto{\pgfqpoint{2.846255in}{1.907327in}}%
\pgfpathlineto{\pgfqpoint{2.849178in}{1.920575in}}%
\pgfpathlineto{\pgfqpoint{2.849834in}{1.921970in}}%
\pgfpathlineto{\pgfqpoint{2.854136in}{1.933823in}}%
\pgfpathlineto{\pgfqpoint{2.858602in}{1.947071in}}%
\pgfpathlineto{\pgfqpoint{2.863559in}{1.960319in}}%
\pgfpathlineto{\pgfqpoint{2.869983in}{1.973567in}}%
\pgfpathlineto{\pgfqpoint{2.873153in}{1.986815in}}%
\pgfpathlineto{\pgfqpoint{2.875117in}{1.991074in}}%
\pgfpathlineto{\pgfqpoint{2.878504in}{2.000063in}}%
\pgfpathlineto{\pgfqpoint{2.880517in}{2.013311in}}%
\pgfpathlineto{\pgfqpoint{2.887004in}{2.026559in}}%
\pgfpathlineto{\pgfqpoint{2.888966in}{2.039807in}}%
\pgfpathlineto{\pgfqpoint{2.896500in}{2.053055in}}%
\pgfpathlineto{\pgfqpoint{2.898008in}{2.066303in}}%
\pgfpathlineto{\pgfqpoint{2.900400in}{2.070528in}}%
\pgfpathlineto{\pgfqpoint{2.904376in}{2.079551in}}%
\pgfpathlineto{\pgfqpoint{2.905603in}{2.092799in}}%
\pgfpathlineto{\pgfqpoint{2.911068in}{2.106047in}}%
\pgfpathlineto{\pgfqpoint{2.912488in}{2.119295in}}%
\pgfpathlineto{\pgfqpoint{2.918983in}{2.132543in}}%
\pgfpathlineto{\pgfqpoint{2.900400in}{2.132543in}}%
\pgfpathlineto{\pgfqpoint{2.897696in}{2.132543in}}%
\pgfpathlineto{\pgfqpoint{2.889315in}{2.119295in}}%
\pgfpathlineto{\pgfqpoint{2.887734in}{2.106047in}}%
\pgfpathlineto{\pgfqpoint{2.882425in}{2.092799in}}%
\pgfpathlineto{\pgfqpoint{2.881105in}{2.079551in}}%
\pgfpathlineto{\pgfqpoint{2.875924in}{2.066303in}}%
\pgfpathlineto{\pgfqpoint{2.875117in}{2.059060in}}%
\pgfpathlineto{\pgfqpoint{2.873866in}{2.053055in}}%
\pgfpathlineto{\pgfqpoint{2.864810in}{2.039807in}}%
\pgfpathlineto{\pgfqpoint{2.861575in}{2.026559in}}%
\pgfpathlineto{\pgfqpoint{2.855822in}{2.013311in}}%
\pgfpathlineto{\pgfqpoint{2.853260in}{2.000063in}}%
\pgfpathlineto{\pgfqpoint{2.849834in}{1.988216in}}%
\pgfpathlineto{\pgfqpoint{2.849188in}{1.986815in}}%
\pgfpathlineto{\pgfqpoint{2.844290in}{1.973567in}}%
\pgfpathlineto{\pgfqpoint{2.837804in}{1.960319in}}%
\pgfpathlineto{\pgfqpoint{2.831598in}{1.947071in}}%
\pgfpathlineto{\pgfqpoint{2.828021in}{1.933823in}}%
\pgfpathlineto{\pgfqpoint{2.824551in}{1.925294in}}%
\pgfpathlineto{\pgfqpoint{2.821499in}{1.920575in}}%
\pgfpathlineto{\pgfqpoint{2.818917in}{1.907327in}}%
\pgfpathlineto{\pgfqpoint{2.812007in}{1.894079in}}%
\pgfpathlineto{\pgfqpoint{2.809394in}{1.880831in}}%
\pgfpathlineto{\pgfqpoint{2.801672in}{1.867583in}}%
\pgfpathlineto{\pgfqpoint{2.799500in}{1.854335in}}%
\pgfpathlineto{\pgfqpoint{2.799268in}{1.853997in}}%
\pgfpathlineto{\pgfqpoint{2.787585in}{1.841087in}}%
\pgfpathlineto{\pgfqpoint{2.784881in}{1.827839in}}%
\pgfpathlineto{\pgfqpoint{2.777096in}{1.814591in}}%
\pgfpathlineto{\pgfqpoint{2.775602in}{1.801343in}}%
\pgfpathlineto{\pgfqpoint{2.773984in}{1.798488in}}%
\pgfpathlineto{\pgfqpoint{2.766444in}{1.788095in}}%
\pgfpathlineto{\pgfqpoint{2.763317in}{1.774847in}}%
\pgfpathlineto{\pgfqpoint{2.753349in}{1.761599in}}%
\pgfpathlineto{\pgfqpoint{2.748979in}{1.748351in}}%
\pgfpathlineto{\pgfqpoint{2.748701in}{1.747768in}}%
\pgfpathlineto{\pgfqpoint{2.741162in}{1.735103in}}%
\pgfpathlineto{\pgfqpoint{2.736193in}{1.721855in}}%
\pgfpathlineto{\pgfqpoint{2.730389in}{1.708607in}}%
\pgfpathlineto{\pgfqpoint{2.724276in}{1.695359in}}%
\pgfpathlineto{\pgfqpoint{2.723418in}{1.693636in}}%
\pgfpathlineto{\pgfqpoint{2.716398in}{1.682111in}}%
\pgfpathlineto{\pgfqpoint{2.706452in}{1.668863in}}%
\pgfpathlineto{\pgfqpoint{2.701757in}{1.655615in}}%
\pgfpathlineto{\pgfqpoint{2.698135in}{1.650971in}}%
\pgfpathlineto{\pgfqpoint{2.689745in}{1.642367in}}%
\pgfpathlineto{\pgfqpoint{2.686763in}{1.629119in}}%
\pgfpathlineto{\pgfqpoint{2.677637in}{1.615871in}}%
\pgfpathlineto{\pgfqpoint{2.676059in}{1.602623in}}%
\pgfpathlineto{\pgfqpoint{2.672852in}{1.599108in}}%
\pgfpathlineto{\pgfqpoint{2.662421in}{1.589375in}}%
\pgfpathlineto{\pgfqpoint{2.658978in}{1.576127in}}%
\pgfpathlineto{\pgfqpoint{2.647569in}{1.566470in}}%
\pgfpathlineto{\pgfqpoint{2.642761in}{1.562879in}}%
\pgfpathlineto{\pgfqpoint{2.640075in}{1.549631in}}%
\pgfpathlineto{\pgfqpoint{2.628126in}{1.536383in}}%
\pgfpathlineto{\pgfqpoint{2.625483in}{1.523135in}}%
\pgfpathlineto{\pgfqpoint{2.622285in}{1.520002in}}%
\pgfpathlineto{\pgfqpoint{2.611305in}{1.509887in}}%
\pgfpathlineto{\pgfqpoint{2.606172in}{1.496639in}}%
\pgfpathlineto{\pgfqpoint{2.597002in}{1.488102in}}%
\pgfpathlineto{\pgfqpoint{2.591578in}{1.483391in}}%
\pgfpathlineto{\pgfqpoint{2.583871in}{1.470143in}}%
\pgfpathlineto{\pgfqpoint{2.573578in}{1.456895in}}%
\pgfpathlineto{\pgfqpoint{2.571719in}{1.453718in}}%
\pgfpathlineto{\pgfqpoint{2.564507in}{1.443647in}}%
\pgfpathlineto{\pgfqpoint{2.557239in}{1.430399in}}%
\pgfpathlineto{\pgfqpoint{2.548260in}{1.417151in}}%
\pgfpathlineto{\pgfqpoint{2.546436in}{1.414339in}}%
\pgfpathlineto{\pgfqpoint{2.539097in}{1.403903in}}%
\pgfpathlineto{\pgfqpoint{2.522016in}{1.390655in}}%
\pgfpathlineto{\pgfqpoint{2.521153in}{1.389261in}}%
\pgfpathlineto{\pgfqpoint{2.513335in}{1.377407in}}%
\pgfpathlineto{\pgfqpoint{2.495870in}{1.365887in}}%
\pgfpathlineto{\pgfqpoint{2.493011in}{1.364159in}}%
\pgfpathlineto{\pgfqpoint{2.492138in}{1.350911in}}%
\pgfpathlineto{\pgfqpoint{2.475033in}{1.337663in}}%
\pgfpathlineto{\pgfqpoint{2.472947in}{1.324415in}}%
\pgfpathlineto{\pgfqpoint{2.470586in}{1.323066in}}%
\pgfpathlineto{\pgfqpoint{2.470586in}{1.311167in}}%
\pgfpathlineto{\pgfqpoint{2.470586in}{1.297919in}}%
\pgfpathlineto{\pgfqpoint{2.470586in}{1.284671in}}%
\pgfpathlineto{\pgfqpoint{2.470586in}{1.271423in}}%
\pgfpathlineto{\pgfqpoint{2.470586in}{1.258175in}}%
\pgfpathlineto{\pgfqpoint{2.470586in}{1.244927in}}%
\pgfpathlineto{\pgfqpoint{2.470586in}{1.231679in}}%
\pgfpathlineto{\pgfqpoint{2.470586in}{1.218431in}}%
\pgfpathlineto{\pgfqpoint{2.470586in}{1.205183in}}%
\pgfpathlineto{\pgfqpoint{2.470586in}{1.191935in}}%
\pgfpathlineto{\pgfqpoint{2.470586in}{1.178687in}}%
\pgfpathlineto{\pgfqpoint{2.470586in}{1.165439in}}%
\pgfpathlineto{\pgfqpoint{2.470586in}{1.152191in}}%
\pgfpathlineto{\pgfqpoint{2.470586in}{1.138943in}}%
\pgfpathlineto{\pgfqpoint{2.470586in}{1.125695in}}%
\pgfpathlineto{\pgfqpoint{2.470586in}{1.112447in}}%
\pgfpathlineto{\pgfqpoint{2.470586in}{1.099199in}}%
\pgfpathlineto{\pgfqpoint{2.470586in}{1.087826in}}%
\pgfpathlineto{\pgfqpoint{2.495870in}{1.095398in}}%
\pgfpathclose%
\pgfpathmoveto{\pgfqpoint{2.511867in}{1.231679in}}%
\pgfpathlineto{\pgfqpoint{2.521153in}{1.234459in}}%
\pgfpathlineto{\pgfqpoint{2.546436in}{1.238086in}}%
\pgfpathlineto{\pgfqpoint{2.571719in}{1.232135in}}%
\pgfpathlineto{\pgfqpoint{2.572148in}{1.231679in}}%
\pgfpathlineto{\pgfqpoint{2.571719in}{1.231586in}}%
\pgfpathlineto{\pgfqpoint{2.546436in}{1.229062in}}%
\pgfpathlineto{\pgfqpoint{2.521153in}{1.229832in}}%
\pgfpathlineto{\pgfqpoint{2.511867in}{1.231679in}}%
\pgfpathclose%
\pgfpathmoveto{\pgfqpoint{4.290975in}{1.505144in}}%
\pgfpathlineto{\pgfqpoint{4.290975in}{1.509887in}}%
\pgfpathlineto{\pgfqpoint{4.290975in}{1.523135in}}%
\pgfpathlineto{\pgfqpoint{4.290975in}{1.536383in}}%
\pgfpathlineto{\pgfqpoint{4.290975in}{1.549631in}}%
\pgfpathlineto{\pgfqpoint{4.290975in}{1.562879in}}%
\pgfpathlineto{\pgfqpoint{4.290975in}{1.576127in}}%
\pgfpathlineto{\pgfqpoint{4.290975in}{1.589375in}}%
\pgfpathlineto{\pgfqpoint{4.290975in}{1.602623in}}%
\pgfpathlineto{\pgfqpoint{4.290975in}{1.615871in}}%
\pgfpathlineto{\pgfqpoint{4.290975in}{1.629119in}}%
\pgfpathlineto{\pgfqpoint{4.290975in}{1.642367in}}%
\pgfpathlineto{\pgfqpoint{4.290975in}{1.655615in}}%
\pgfpathlineto{\pgfqpoint{4.290975in}{1.667261in}}%
\pgfpathlineto{\pgfqpoint{4.288884in}{1.668863in}}%
\pgfpathlineto{\pgfqpoint{4.281256in}{1.682111in}}%
\pgfpathlineto{\pgfqpoint{4.272678in}{1.695359in}}%
\pgfpathlineto{\pgfqpoint{4.265730in}{1.708607in}}%
\pgfpathlineto{\pgfqpoint{4.265692in}{1.708676in}}%
\pgfpathlineto{\pgfqpoint{4.258731in}{1.721855in}}%
\pgfpathlineto{\pgfqpoint{4.248746in}{1.735103in}}%
\pgfpathlineto{\pgfqpoint{4.240408in}{1.746973in}}%
\pgfpathlineto{\pgfqpoint{4.239446in}{1.748351in}}%
\pgfpathlineto{\pgfqpoint{4.226000in}{1.761599in}}%
\pgfpathlineto{\pgfqpoint{4.221790in}{1.774847in}}%
\pgfpathlineto{\pgfqpoint{4.215125in}{1.781887in}}%
\pgfpathlineto{\pgfqpoint{4.209377in}{1.788095in}}%
\pgfpathlineto{\pgfqpoint{4.207208in}{1.801343in}}%
\pgfpathlineto{\pgfqpoint{4.196306in}{1.814591in}}%
\pgfpathlineto{\pgfqpoint{4.193310in}{1.827839in}}%
\pgfpathlineto{\pgfqpoint{4.189842in}{1.831043in}}%
\pgfpathlineto{\pgfqpoint{4.179845in}{1.841087in}}%
\pgfpathlineto{\pgfqpoint{4.176360in}{1.854335in}}%
\pgfpathlineto{\pgfqpoint{4.164559in}{1.865512in}}%
\pgfpathlineto{\pgfqpoint{4.162389in}{1.867583in}}%
\pgfpathlineto{\pgfqpoint{4.159453in}{1.880831in}}%
\pgfpathlineto{\pgfqpoint{4.150576in}{1.894079in}}%
\pgfpathlineto{\pgfqpoint{4.147366in}{1.907327in}}%
\pgfpathlineto{\pgfqpoint{4.139276in}{1.917185in}}%
\pgfpathlineto{\pgfqpoint{4.136724in}{1.920575in}}%
\pgfpathlineto{\pgfqpoint{4.131526in}{1.933823in}}%
\pgfpathlineto{\pgfqpoint{4.121332in}{1.947071in}}%
\pgfpathlineto{\pgfqpoint{4.114992in}{1.960319in}}%
\pgfpathlineto{\pgfqpoint{4.113993in}{1.962490in}}%
\pgfpathlineto{\pgfqpoint{4.109575in}{1.973567in}}%
\pgfpathlineto{\pgfqpoint{4.103791in}{1.986815in}}%
\pgfpathlineto{\pgfqpoint{4.098769in}{2.000063in}}%
\pgfpathlineto{\pgfqpoint{4.091327in}{2.013311in}}%
\pgfpathlineto{\pgfqpoint{4.088709in}{2.018908in}}%
\pgfpathlineto{\pgfqpoint{4.085780in}{2.026559in}}%
\pgfpathlineto{\pgfqpoint{4.076666in}{2.039807in}}%
\pgfpathlineto{\pgfqpoint{4.072715in}{2.053055in}}%
\pgfpathlineto{\pgfqpoint{4.063426in}{2.064943in}}%
\pgfpathlineto{\pgfqpoint{4.062486in}{2.066303in}}%
\pgfpathlineto{\pgfqpoint{4.060532in}{2.079551in}}%
\pgfpathlineto{\pgfqpoint{4.053665in}{2.092799in}}%
\pgfpathlineto{\pgfqpoint{4.051123in}{2.106047in}}%
\pgfpathlineto{\pgfqpoint{4.040738in}{2.119295in}}%
\pgfpathlineto{\pgfqpoint{4.038143in}{2.129418in}}%
\pgfpathlineto{\pgfqpoint{4.037538in}{2.132543in}}%
\pgfpathlineto{\pgfqpoint{4.012860in}{2.132543in}}%
\pgfpathlineto{\pgfqpoint{4.000772in}{2.132543in}}%
\pgfpathlineto{\pgfqpoint{4.001850in}{2.119295in}}%
\pgfpathlineto{\pgfqpoint{4.011878in}{2.106047in}}%
\pgfpathlineto{\pgfqpoint{4.012860in}{2.095538in}}%
\pgfpathlineto{\pgfqpoint{4.013096in}{2.092799in}}%
\pgfpathlineto{\pgfqpoint{4.020295in}{2.079551in}}%
\pgfpathlineto{\pgfqpoint{4.020591in}{2.066303in}}%
\pgfpathlineto{\pgfqpoint{4.029906in}{2.053055in}}%
\pgfpathlineto{\pgfqpoint{4.033154in}{2.039807in}}%
\pgfpathlineto{\pgfqpoint{4.038143in}{2.032034in}}%
\pgfpathlineto{\pgfqpoint{4.041286in}{2.026559in}}%
\pgfpathlineto{\pgfqpoint{4.045710in}{2.013311in}}%
\pgfpathlineto{\pgfqpoint{4.050773in}{2.000063in}}%
\pgfpathlineto{\pgfqpoint{4.055224in}{1.986815in}}%
\pgfpathlineto{\pgfqpoint{4.058340in}{1.973567in}}%
\pgfpathlineto{\pgfqpoint{4.063426in}{1.964323in}}%
\pgfpathlineto{\pgfqpoint{4.065056in}{1.960319in}}%
\pgfpathlineto{\pgfqpoint{4.065544in}{1.947071in}}%
\pgfpathlineto{\pgfqpoint{4.077899in}{1.933823in}}%
\pgfpathlineto{\pgfqpoint{4.078541in}{1.920575in}}%
\pgfpathlineto{\pgfqpoint{4.088709in}{1.910659in}}%
\pgfpathlineto{\pgfqpoint{4.091259in}{1.907327in}}%
\pgfpathlineto{\pgfqpoint{4.088709in}{1.897034in}}%
\pgfpathlineto{\pgfqpoint{4.086988in}{1.894079in}}%
\pgfpathlineto{\pgfqpoint{4.088709in}{1.892885in}}%
\pgfpathlineto{\pgfqpoint{4.098437in}{1.880831in}}%
\pgfpathlineto{\pgfqpoint{4.090589in}{1.867583in}}%
\pgfpathlineto{\pgfqpoint{4.110443in}{1.854335in}}%
\pgfpathlineto{\pgfqpoint{4.098649in}{1.841087in}}%
\pgfpathlineto{\pgfqpoint{4.113993in}{1.832443in}}%
\pgfpathlineto{\pgfqpoint{4.118413in}{1.827839in}}%
\pgfpathlineto{\pgfqpoint{4.113993in}{1.821483in}}%
\pgfpathlineto{\pgfqpoint{4.088709in}{1.818969in}}%
\pgfpathlineto{\pgfqpoint{4.073157in}{1.827839in}}%
\pgfpathlineto{\pgfqpoint{4.072509in}{1.841087in}}%
\pgfpathlineto{\pgfqpoint{4.063426in}{1.843367in}}%
\pgfpathlineto{\pgfqpoint{4.049560in}{1.854335in}}%
\pgfpathlineto{\pgfqpoint{4.051943in}{1.867583in}}%
\pgfpathlineto{\pgfqpoint{4.038143in}{1.872724in}}%
\pgfpathlineto{\pgfqpoint{4.027458in}{1.880831in}}%
\pgfpathlineto{\pgfqpoint{4.026620in}{1.894079in}}%
\pgfpathlineto{\pgfqpoint{4.012860in}{1.901194in}}%
\pgfpathlineto{\pgfqpoint{4.005275in}{1.907327in}}%
\pgfpathlineto{\pgfqpoint{4.003243in}{1.920575in}}%
\pgfpathlineto{\pgfqpoint{3.987988in}{1.933823in}}%
\pgfpathlineto{\pgfqpoint{3.989200in}{1.947071in}}%
\pgfpathlineto{\pgfqpoint{3.987577in}{1.948054in}}%
\pgfpathlineto{\pgfqpoint{3.975425in}{1.960319in}}%
\pgfpathlineto{\pgfqpoint{3.972632in}{1.973567in}}%
\pgfpathlineto{\pgfqpoint{3.962293in}{1.981509in}}%
\pgfpathlineto{\pgfqpoint{3.957130in}{1.986815in}}%
\pgfpathlineto{\pgfqpoint{3.952532in}{2.000063in}}%
\pgfpathlineto{\pgfqpoint{3.941994in}{2.013311in}}%
\pgfpathlineto{\pgfqpoint{3.937010in}{2.025580in}}%
\pgfpathlineto{\pgfqpoint{3.936697in}{2.026559in}}%
\pgfpathlineto{\pgfqpoint{3.928395in}{2.039807in}}%
\pgfpathlineto{\pgfqpoint{3.923655in}{2.053055in}}%
\pgfpathlineto{\pgfqpoint{3.915976in}{2.066303in}}%
\pgfpathlineto{\pgfqpoint{3.911727in}{2.072875in}}%
\pgfpathlineto{\pgfqpoint{3.908395in}{2.079551in}}%
\pgfpathlineto{\pgfqpoint{3.902792in}{2.092799in}}%
\pgfpathlineto{\pgfqpoint{3.894520in}{2.106047in}}%
\pgfpathlineto{\pgfqpoint{3.890058in}{2.119295in}}%
\pgfpathlineto{\pgfqpoint{3.886444in}{2.125163in}}%
\pgfpathlineto{\pgfqpoint{3.882761in}{2.132543in}}%
\pgfpathlineto{\pgfqpoint{3.861161in}{2.132543in}}%
\pgfpathlineto{\pgfqpoint{3.848450in}{2.132543in}}%
\pgfpathlineto{\pgfqpoint{3.854166in}{2.119295in}}%
\pgfpathlineto{\pgfqpoint{3.858856in}{2.106047in}}%
\pgfpathlineto{\pgfqpoint{3.861161in}{2.102371in}}%
\pgfpathlineto{\pgfqpoint{3.865721in}{2.092799in}}%
\pgfpathlineto{\pgfqpoint{3.871441in}{2.079551in}}%
\pgfpathlineto{\pgfqpoint{3.877017in}{2.066303in}}%
\pgfpathlineto{\pgfqpoint{3.884945in}{2.053055in}}%
\pgfpathlineto{\pgfqpoint{3.886444in}{2.048919in}}%
\pgfpathlineto{\pgfqpoint{3.889083in}{2.039807in}}%
\pgfpathlineto{\pgfqpoint{3.895924in}{2.026559in}}%
\pgfpathlineto{\pgfqpoint{3.899767in}{2.013311in}}%
\pgfpathlineto{\pgfqpoint{3.909007in}{2.000063in}}%
\pgfpathlineto{\pgfqpoint{3.911727in}{1.992957in}}%
\pgfpathlineto{\pgfqpoint{3.913687in}{1.986815in}}%
\pgfpathlineto{\pgfqpoint{3.924845in}{1.973567in}}%
\pgfpathlineto{\pgfqpoint{3.928152in}{1.960319in}}%
\pgfpathlineto{\pgfqpoint{3.937010in}{1.948772in}}%
\pgfpathlineto{\pgfqpoint{3.938303in}{1.947071in}}%
\pgfpathlineto{\pgfqpoint{3.940187in}{1.933823in}}%
\pgfpathlineto{\pgfqpoint{3.949328in}{1.920575in}}%
\pgfpathlineto{\pgfqpoint{3.953136in}{1.907327in}}%
\pgfpathlineto{\pgfqpoint{3.962293in}{1.898277in}}%
\pgfpathlineto{\pgfqpoint{3.966475in}{1.894079in}}%
\pgfpathlineto{\pgfqpoint{3.970916in}{1.880831in}}%
\pgfpathlineto{\pgfqpoint{3.981862in}{1.867583in}}%
\pgfpathlineto{\pgfqpoint{3.987577in}{1.854480in}}%
\pgfpathlineto{\pgfqpoint{3.987636in}{1.854335in}}%
\pgfpathlineto{\pgfqpoint{3.996436in}{1.841087in}}%
\pgfpathlineto{\pgfqpoint{4.001946in}{1.827839in}}%
\pgfpathlineto{\pgfqpoint{4.012285in}{1.814591in}}%
\pgfpathlineto{\pgfqpoint{4.012860in}{1.813727in}}%
\pgfpathlineto{\pgfqpoint{4.020531in}{1.801343in}}%
\pgfpathlineto{\pgfqpoint{4.030319in}{1.788095in}}%
\pgfpathlineto{\pgfqpoint{4.038143in}{1.779081in}}%
\pgfpathlineto{\pgfqpoint{4.041870in}{1.774847in}}%
\pgfpathlineto{\pgfqpoint{4.047828in}{1.761599in}}%
\pgfpathlineto{\pgfqpoint{4.057348in}{1.748351in}}%
\pgfpathlineto{\pgfqpoint{4.063045in}{1.735103in}}%
\pgfpathlineto{\pgfqpoint{4.063426in}{1.734724in}}%
\pgfpathlineto{\pgfqpoint{4.076425in}{1.721855in}}%
\pgfpathlineto{\pgfqpoint{4.083994in}{1.708607in}}%
\pgfpathlineto{\pgfqpoint{4.088709in}{1.705150in}}%
\pgfpathlineto{\pgfqpoint{4.102901in}{1.695359in}}%
\pgfpathlineto{\pgfqpoint{4.107056in}{1.682111in}}%
\pgfpathlineto{\pgfqpoint{4.113993in}{1.676230in}}%
\pgfpathlineto{\pgfqpoint{4.123876in}{1.668863in}}%
\pgfpathlineto{\pgfqpoint{4.127488in}{1.655615in}}%
\pgfpathlineto{\pgfqpoint{4.139276in}{1.645663in}}%
\pgfpathlineto{\pgfqpoint{4.143799in}{1.642367in}}%
\pgfpathlineto{\pgfqpoint{4.148305in}{1.629119in}}%
\pgfpathlineto{\pgfqpoint{4.164559in}{1.618570in}}%
\pgfpathlineto{\pgfqpoint{4.169419in}{1.615871in}}%
\pgfpathlineto{\pgfqpoint{4.175335in}{1.602623in}}%
\pgfpathlineto{\pgfqpoint{4.189842in}{1.593245in}}%
\pgfpathlineto{\pgfqpoint{4.197270in}{1.589375in}}%
\pgfpathlineto{\pgfqpoint{4.204073in}{1.576127in}}%
\pgfpathlineto{\pgfqpoint{4.215125in}{1.566090in}}%
\pgfpathlineto{\pgfqpoint{4.219486in}{1.562879in}}%
\pgfpathlineto{\pgfqpoint{4.226688in}{1.549631in}}%
\pgfpathlineto{\pgfqpoint{4.240408in}{1.539704in}}%
\pgfpathlineto{\pgfqpoint{4.245883in}{1.536383in}}%
\pgfpathlineto{\pgfqpoint{4.260665in}{1.523135in}}%
\pgfpathlineto{\pgfqpoint{4.265692in}{1.519832in}}%
\pgfpathlineto{\pgfqpoint{4.284667in}{1.509887in}}%
\pgfpathlineto{\pgfqpoint{4.290975in}{1.505144in}}%
\pgfpathclose%
\pgfpathmoveto{\pgfqpoint{4.112518in}{1.801343in}}%
\pgfpathlineto{\pgfqpoint{4.113993in}{1.802518in}}%
\pgfpathlineto{\pgfqpoint{4.114950in}{1.801343in}}%
\pgfpathlineto{\pgfqpoint{4.113993in}{1.800873in}}%
\pgfpathlineto{\pgfqpoint{4.112518in}{1.801343in}}%
\pgfpathclose%
\pgfpathmoveto{\pgfqpoint{2.477135in}{1.576127in}}%
\pgfpathlineto{\pgfqpoint{2.482952in}{1.589375in}}%
\pgfpathlineto{\pgfqpoint{2.495870in}{1.598464in}}%
\pgfpathlineto{\pgfqpoint{2.501320in}{1.602623in}}%
\pgfpathlineto{\pgfqpoint{2.507360in}{1.615871in}}%
\pgfpathlineto{\pgfqpoint{2.521153in}{1.624143in}}%
\pgfpathlineto{\pgfqpoint{2.528996in}{1.629119in}}%
\pgfpathlineto{\pgfqpoint{2.531407in}{1.642367in}}%
\pgfpathlineto{\pgfqpoint{2.546436in}{1.651721in}}%
\pgfpathlineto{\pgfqpoint{2.552089in}{1.655615in}}%
\pgfpathlineto{\pgfqpoint{2.553599in}{1.668863in}}%
\pgfpathlineto{\pgfqpoint{2.568807in}{1.682111in}}%
\pgfpathlineto{\pgfqpoint{2.570714in}{1.695359in}}%
\pgfpathlineto{\pgfqpoint{2.571719in}{1.695955in}}%
\pgfpathlineto{\pgfqpoint{2.589793in}{1.708607in}}%
\pgfpathlineto{\pgfqpoint{2.596322in}{1.721855in}}%
\pgfpathlineto{\pgfqpoint{2.597002in}{1.722314in}}%
\pgfpathlineto{\pgfqpoint{2.612939in}{1.735103in}}%
\pgfpathlineto{\pgfqpoint{2.619677in}{1.748351in}}%
\pgfpathlineto{\pgfqpoint{2.622285in}{1.752031in}}%
\pgfpathlineto{\pgfqpoint{2.628516in}{1.761599in}}%
\pgfpathlineto{\pgfqpoint{2.635132in}{1.774847in}}%
\pgfpathlineto{\pgfqpoint{2.643724in}{1.788095in}}%
\pgfpathlineto{\pgfqpoint{2.647569in}{1.793698in}}%
\pgfpathlineto{\pgfqpoint{2.653170in}{1.801343in}}%
\pgfpathlineto{\pgfqpoint{2.660863in}{1.814591in}}%
\pgfpathlineto{\pgfqpoint{2.672852in}{1.827660in}}%
\pgfpathlineto{\pgfqpoint{2.673039in}{1.827839in}}%
\pgfpathlineto{\pgfqpoint{2.677673in}{1.841087in}}%
\pgfpathlineto{\pgfqpoint{2.689249in}{1.854335in}}%
\pgfpathlineto{\pgfqpoint{2.691897in}{1.867583in}}%
\pgfpathlineto{\pgfqpoint{2.698135in}{1.875040in}}%
\pgfpathlineto{\pgfqpoint{2.703127in}{1.880831in}}%
\pgfpathlineto{\pgfqpoint{2.706600in}{1.894079in}}%
\pgfpathlineto{\pgfqpoint{2.720509in}{1.907327in}}%
\pgfpathlineto{\pgfqpoint{2.723418in}{1.918990in}}%
\pgfpathlineto{\pgfqpoint{2.723903in}{1.920575in}}%
\pgfpathlineto{\pgfqpoint{2.735568in}{1.933823in}}%
\pgfpathlineto{\pgfqpoint{2.737433in}{1.947071in}}%
\pgfpathlineto{\pgfqpoint{2.745406in}{1.960319in}}%
\pgfpathlineto{\pgfqpoint{2.748609in}{1.973567in}}%
\pgfpathlineto{\pgfqpoint{2.748701in}{1.973684in}}%
\pgfpathlineto{\pgfqpoint{2.759147in}{1.986815in}}%
\pgfpathlineto{\pgfqpoint{2.764285in}{2.000063in}}%
\pgfpathlineto{\pgfqpoint{2.772312in}{2.013311in}}%
\pgfpathlineto{\pgfqpoint{2.773984in}{2.016828in}}%
\pgfpathlineto{\pgfqpoint{2.779796in}{2.026559in}}%
\pgfpathlineto{\pgfqpoint{2.785251in}{2.039807in}}%
\pgfpathlineto{\pgfqpoint{2.790921in}{2.053055in}}%
\pgfpathlineto{\pgfqpoint{2.794708in}{2.066303in}}%
\pgfpathlineto{\pgfqpoint{2.799268in}{2.075246in}}%
\pgfpathlineto{\pgfqpoint{2.802080in}{2.079551in}}%
\pgfpathlineto{\pgfqpoint{2.806852in}{2.092799in}}%
\pgfpathlineto{\pgfqpoint{2.816182in}{2.106047in}}%
\pgfpathlineto{\pgfqpoint{2.818272in}{2.119295in}}%
\pgfpathlineto{\pgfqpoint{2.824551in}{2.131204in}}%
\pgfpathlineto{\pgfqpoint{2.825494in}{2.132543in}}%
\pgfpathlineto{\pgfqpoint{2.824551in}{2.132543in}}%
\pgfpathlineto{\pgfqpoint{2.799268in}{2.132543in}}%
\pgfpathlineto{\pgfqpoint{2.793432in}{2.132543in}}%
\pgfpathlineto{\pgfqpoint{2.784698in}{2.119295in}}%
\pgfpathlineto{\pgfqpoint{2.782037in}{2.106047in}}%
\pgfpathlineto{\pgfqpoint{2.773984in}{2.094719in}}%
\pgfpathlineto{\pgfqpoint{2.772402in}{2.092799in}}%
\pgfpathlineto{\pgfqpoint{2.766390in}{2.079551in}}%
\pgfpathlineto{\pgfqpoint{2.759665in}{2.066303in}}%
\pgfpathlineto{\pgfqpoint{2.753874in}{2.053055in}}%
\pgfpathlineto{\pgfqpoint{2.748701in}{2.040086in}}%
\pgfpathlineto{\pgfqpoint{2.748571in}{2.039807in}}%
\pgfpathlineto{\pgfqpoint{2.738813in}{2.026559in}}%
\pgfpathlineto{\pgfqpoint{2.734215in}{2.013311in}}%
\pgfpathlineto{\pgfqpoint{2.723418in}{2.001253in}}%
\pgfpathlineto{\pgfqpoint{2.721878in}{2.000063in}}%
\pgfpathlineto{\pgfqpoint{2.718684in}{1.986815in}}%
\pgfpathlineto{\pgfqpoint{2.704017in}{1.973567in}}%
\pgfpathlineto{\pgfqpoint{2.703686in}{1.960319in}}%
\pgfpathlineto{\pgfqpoint{2.698135in}{1.953819in}}%
\pgfpathlineto{\pgfqpoint{2.690206in}{1.947071in}}%
\pgfpathlineto{\pgfqpoint{2.691991in}{1.933823in}}%
\pgfpathlineto{\pgfqpoint{2.676151in}{1.920575in}}%
\pgfpathlineto{\pgfqpoint{2.674117in}{1.907327in}}%
\pgfpathlineto{\pgfqpoint{2.672852in}{1.906299in}}%
\pgfpathlineto{\pgfqpoint{2.653460in}{1.894079in}}%
\pgfpathlineto{\pgfqpoint{2.651122in}{1.880831in}}%
\pgfpathlineto{\pgfqpoint{2.647569in}{1.877076in}}%
\pgfpathlineto{\pgfqpoint{2.635240in}{1.867583in}}%
\pgfpathlineto{\pgfqpoint{2.631865in}{1.854335in}}%
\pgfpathlineto{\pgfqpoint{2.622285in}{1.843937in}}%
\pgfpathlineto{\pgfqpoint{2.618503in}{1.841087in}}%
\pgfpathlineto{\pgfqpoint{2.608204in}{1.827839in}}%
\pgfpathlineto{\pgfqpoint{2.600244in}{1.814591in}}%
\pgfpathlineto{\pgfqpoint{2.597002in}{1.811325in}}%
\pgfpathlineto{\pgfqpoint{2.577963in}{1.801343in}}%
\pgfpathlineto{\pgfqpoint{2.576431in}{1.788095in}}%
\pgfpathlineto{\pgfqpoint{2.571719in}{1.784636in}}%
\pgfpathlineto{\pgfqpoint{2.546436in}{1.777583in}}%
\pgfpathlineto{\pgfqpoint{2.521153in}{1.787250in}}%
\pgfpathlineto{\pgfqpoint{2.520514in}{1.788095in}}%
\pgfpathlineto{\pgfqpoint{2.521153in}{1.788647in}}%
\pgfpathlineto{\pgfqpoint{2.542407in}{1.801343in}}%
\pgfpathlineto{\pgfqpoint{2.529898in}{1.814591in}}%
\pgfpathlineto{\pgfqpoint{2.545805in}{1.827839in}}%
\pgfpathlineto{\pgfqpoint{2.543554in}{1.841087in}}%
\pgfpathlineto{\pgfqpoint{2.546436in}{1.844154in}}%
\pgfpathlineto{\pgfqpoint{2.557675in}{1.854335in}}%
\pgfpathlineto{\pgfqpoint{2.560436in}{1.867583in}}%
\pgfpathlineto{\pgfqpoint{2.567809in}{1.880831in}}%
\pgfpathlineto{\pgfqpoint{2.571719in}{1.892523in}}%
\pgfpathlineto{\pgfqpoint{2.572513in}{1.894079in}}%
\pgfpathlineto{\pgfqpoint{2.575588in}{1.907327in}}%
\pgfpathlineto{\pgfqpoint{2.580370in}{1.920575in}}%
\pgfpathlineto{\pgfqpoint{2.582448in}{1.933823in}}%
\pgfpathlineto{\pgfqpoint{2.593155in}{1.947071in}}%
\pgfpathlineto{\pgfqpoint{2.595876in}{1.960319in}}%
\pgfpathlineto{\pgfqpoint{2.597002in}{1.961590in}}%
\pgfpathlineto{\pgfqpoint{2.610082in}{1.973567in}}%
\pgfpathlineto{\pgfqpoint{2.608709in}{1.986815in}}%
\pgfpathlineto{\pgfqpoint{2.618998in}{2.000063in}}%
\pgfpathlineto{\pgfqpoint{2.617234in}{2.013311in}}%
\pgfpathlineto{\pgfqpoint{2.622285in}{2.018910in}}%
\pgfpathlineto{\pgfqpoint{2.630448in}{2.026559in}}%
\pgfpathlineto{\pgfqpoint{2.628889in}{2.039807in}}%
\pgfpathlineto{\pgfqpoint{2.640677in}{2.053055in}}%
\pgfpathlineto{\pgfqpoint{2.640341in}{2.066303in}}%
\pgfpathlineto{\pgfqpoint{2.647569in}{2.075865in}}%
\pgfpathlineto{\pgfqpoint{2.650724in}{2.079551in}}%
\pgfpathlineto{\pgfqpoint{2.650769in}{2.092799in}}%
\pgfpathlineto{\pgfqpoint{2.659246in}{2.106047in}}%
\pgfpathlineto{\pgfqpoint{2.660089in}{2.119295in}}%
\pgfpathlineto{\pgfqpoint{2.668560in}{2.132543in}}%
\pgfpathlineto{\pgfqpoint{2.647569in}{2.132543in}}%
\pgfpathlineto{\pgfqpoint{2.632421in}{2.132543in}}%
\pgfpathlineto{\pgfqpoint{2.623400in}{2.119295in}}%
\pgfpathlineto{\pgfqpoint{2.622285in}{2.106980in}}%
\pgfpathlineto{\pgfqpoint{2.622168in}{2.106047in}}%
\pgfpathlineto{\pgfqpoint{2.612638in}{2.092799in}}%
\pgfpathlineto{\pgfqpoint{2.611159in}{2.079551in}}%
\pgfpathlineto{\pgfqpoint{2.601700in}{2.066303in}}%
\pgfpathlineto{\pgfqpoint{2.599930in}{2.053055in}}%
\pgfpathlineto{\pgfqpoint{2.597002in}{2.049543in}}%
\pgfpathlineto{\pgfqpoint{2.587732in}{2.039807in}}%
\pgfpathlineto{\pgfqpoint{2.584971in}{2.026559in}}%
\pgfpathlineto{\pgfqpoint{2.575218in}{2.013311in}}%
\pgfpathlineto{\pgfqpoint{2.572620in}{2.000063in}}%
\pgfpathlineto{\pgfqpoint{2.571719in}{1.998570in}}%
\pgfpathlineto{\pgfqpoint{2.564079in}{1.986815in}}%
\pgfpathlineto{\pgfqpoint{2.559967in}{1.973567in}}%
\pgfpathlineto{\pgfqpoint{2.550105in}{1.960319in}}%
\pgfpathlineto{\pgfqpoint{2.546436in}{1.953818in}}%
\pgfpathlineto{\pgfqpoint{2.541647in}{1.947071in}}%
\pgfpathlineto{\pgfqpoint{2.534342in}{1.933823in}}%
\pgfpathlineto{\pgfqpoint{2.526776in}{1.920575in}}%
\pgfpathlineto{\pgfqpoint{2.524032in}{1.907327in}}%
\pgfpathlineto{\pgfqpoint{2.521153in}{1.902494in}}%
\pgfpathlineto{\pgfqpoint{2.515308in}{1.894079in}}%
\pgfpathlineto{\pgfqpoint{2.510833in}{1.880831in}}%
\pgfpathlineto{\pgfqpoint{2.498654in}{1.867583in}}%
\pgfpathlineto{\pgfqpoint{2.495870in}{1.856704in}}%
\pgfpathlineto{\pgfqpoint{2.495175in}{1.854335in}}%
\pgfpathlineto{\pgfqpoint{2.479175in}{1.841087in}}%
\pgfpathlineto{\pgfqpoint{2.477160in}{1.827839in}}%
\pgfpathlineto{\pgfqpoint{2.470586in}{1.821666in}}%
\pgfpathlineto{\pgfqpoint{2.470586in}{1.814591in}}%
\pgfpathlineto{\pgfqpoint{2.470586in}{1.801343in}}%
\pgfpathlineto{\pgfqpoint{2.470586in}{1.788095in}}%
\pgfpathlineto{\pgfqpoint{2.470586in}{1.774847in}}%
\pgfpathlineto{\pgfqpoint{2.470586in}{1.761599in}}%
\pgfpathlineto{\pgfqpoint{2.470586in}{1.748351in}}%
\pgfpathlineto{\pgfqpoint{2.470586in}{1.735103in}}%
\pgfpathlineto{\pgfqpoint{2.470586in}{1.721855in}}%
\pgfpathlineto{\pgfqpoint{2.470586in}{1.708607in}}%
\pgfpathlineto{\pgfqpoint{2.470586in}{1.695359in}}%
\pgfpathlineto{\pgfqpoint{2.470586in}{1.682111in}}%
\pgfpathlineto{\pgfqpoint{2.470586in}{1.668863in}}%
\pgfpathlineto{\pgfqpoint{2.470586in}{1.655615in}}%
\pgfpathlineto{\pgfqpoint{2.470586in}{1.642367in}}%
\pgfpathlineto{\pgfqpoint{2.470586in}{1.629119in}}%
\pgfpathlineto{\pgfqpoint{2.470586in}{1.615871in}}%
\pgfpathlineto{\pgfqpoint{2.470586in}{1.602623in}}%
\pgfpathlineto{\pgfqpoint{2.470586in}{1.589375in}}%
\pgfpathlineto{\pgfqpoint{2.470586in}{1.576127in}}%
\pgfpathlineto{\pgfqpoint{2.470586in}{1.570943in}}%
\pgfpathlineto{\pgfqpoint{2.477135in}{1.576127in}}%
\pgfpathclose%
\pgfpathmoveto{\pgfqpoint{2.496470in}{1.735103in}}%
\pgfpathlineto{\pgfqpoint{2.521153in}{1.740715in}}%
\pgfpathlineto{\pgfqpoint{2.531072in}{1.735103in}}%
\pgfpathlineto{\pgfqpoint{2.521153in}{1.732757in}}%
\pgfpathlineto{\pgfqpoint{2.496470in}{1.735103in}}%
\pgfpathclose%
\pgfpathmoveto{\pgfqpoint{2.507482in}{1.761599in}}%
\pgfpathlineto{\pgfqpoint{2.521153in}{1.767636in}}%
\pgfpathlineto{\pgfqpoint{2.546436in}{1.770712in}}%
\pgfpathlineto{\pgfqpoint{2.555641in}{1.761599in}}%
\pgfpathlineto{\pgfqpoint{2.546436in}{1.756598in}}%
\pgfpathlineto{\pgfqpoint{2.521153in}{1.754060in}}%
\pgfpathlineto{\pgfqpoint{2.507482in}{1.761599in}}%
\pgfpathclose%
\pgfpathmoveto{\pgfqpoint{4.290975in}{1.948728in}}%
\pgfpathlineto{\pgfqpoint{4.290975in}{1.960319in}}%
\pgfpathlineto{\pgfqpoint{4.290975in}{1.973567in}}%
\pgfpathlineto{\pgfqpoint{4.290975in}{1.986815in}}%
\pgfpathlineto{\pgfqpoint{4.290975in}{2.000063in}}%
\pgfpathlineto{\pgfqpoint{4.290975in}{2.013311in}}%
\pgfpathlineto{\pgfqpoint{4.290975in}{2.026559in}}%
\pgfpathlineto{\pgfqpoint{4.290975in}{2.039807in}}%
\pgfpathlineto{\pgfqpoint{4.290975in}{2.053055in}}%
\pgfpathlineto{\pgfqpoint{4.290975in}{2.066303in}}%
\pgfpathlineto{\pgfqpoint{4.290975in}{2.079551in}}%
\pgfpathlineto{\pgfqpoint{4.290975in}{2.092799in}}%
\pgfpathlineto{\pgfqpoint{4.290975in}{2.106047in}}%
\pgfpathlineto{\pgfqpoint{4.290975in}{2.114802in}}%
\pgfpathlineto{\pgfqpoint{4.286195in}{2.119295in}}%
\pgfpathlineto{\pgfqpoint{4.283458in}{2.132543in}}%
\pgfpathlineto{\pgfqpoint{4.265692in}{2.132543in}}%
\pgfpathlineto{\pgfqpoint{4.240408in}{2.132543in}}%
\pgfpathlineto{\pgfqpoint{4.215125in}{2.132543in}}%
\pgfpathlineto{\pgfqpoint{4.189842in}{2.132543in}}%
\pgfpathlineto{\pgfqpoint{4.164559in}{2.132543in}}%
\pgfpathlineto{\pgfqpoint{4.157773in}{2.132543in}}%
\pgfpathlineto{\pgfqpoint{4.164559in}{2.121787in}}%
\pgfpathlineto{\pgfqpoint{4.166412in}{2.119295in}}%
\pgfpathlineto{\pgfqpoint{4.171951in}{2.106047in}}%
\pgfpathlineto{\pgfqpoint{4.182340in}{2.092799in}}%
\pgfpathlineto{\pgfqpoint{4.189842in}{2.079839in}}%
\pgfpathlineto{\pgfqpoint{4.190030in}{2.079551in}}%
\pgfpathlineto{\pgfqpoint{4.200604in}{2.066303in}}%
\pgfpathlineto{\pgfqpoint{4.210121in}{2.053055in}}%
\pgfpathlineto{\pgfqpoint{4.215125in}{2.041927in}}%
\pgfpathlineto{\pgfqpoint{4.216280in}{2.039807in}}%
\pgfpathlineto{\pgfqpoint{4.227802in}{2.026559in}}%
\pgfpathlineto{\pgfqpoint{4.233058in}{2.013311in}}%
\pgfpathlineto{\pgfqpoint{4.240408in}{2.004992in}}%
\pgfpathlineto{\pgfqpoint{4.245573in}{2.000063in}}%
\pgfpathlineto{\pgfqpoint{4.252883in}{1.986815in}}%
\pgfpathlineto{\pgfqpoint{4.265692in}{1.977165in}}%
\pgfpathlineto{\pgfqpoint{4.271710in}{1.973567in}}%
\pgfpathlineto{\pgfqpoint{4.276446in}{1.960319in}}%
\pgfpathlineto{\pgfqpoint{4.290975in}{1.948728in}}%
\pgfpathclose%
\pgfpathmoveto{\pgfqpoint{2.474607in}{2.066303in}}%
\pgfpathlineto{\pgfqpoint{2.483946in}{2.079551in}}%
\pgfpathlineto{\pgfqpoint{2.490721in}{2.092799in}}%
\pgfpathlineto{\pgfqpoint{2.495870in}{2.097402in}}%
\pgfpathlineto{\pgfqpoint{2.505137in}{2.106047in}}%
\pgfpathlineto{\pgfqpoint{2.507946in}{2.119295in}}%
\pgfpathlineto{\pgfqpoint{2.520208in}{2.132543in}}%
\pgfpathlineto{\pgfqpoint{2.495870in}{2.132543in}}%
\pgfpathlineto{\pgfqpoint{2.470586in}{2.132543in}}%
\pgfpathlineto{\pgfqpoint{2.470586in}{2.119295in}}%
\pgfpathlineto{\pgfqpoint{2.470586in}{2.106047in}}%
\pgfpathlineto{\pgfqpoint{2.470586in}{2.092799in}}%
\pgfpathlineto{\pgfqpoint{2.470586in}{2.079551in}}%
\pgfpathlineto{\pgfqpoint{2.470586in}{2.066303in}}%
\pgfpathlineto{\pgfqpoint{2.470586in}{2.056822in}}%
\pgfpathlineto{\pgfqpoint{2.474607in}{2.066303in}}%
\pgfpathclose%
\pgfusepath{fill}%
\end{pgfscope}%
\begin{pgfscope}%
\pgfpathrectangle{\pgfqpoint{2.470586in}{0.476543in}}{\pgfqpoint{1.820388in}{1.656000in}}%
\pgfusepath{clip}%
\pgfsetbuttcap%
\pgfsetroundjoin%
\definecolor{currentfill}{rgb}{0.274202,0.463130,0.708112}%
\pgfsetfillcolor{currentfill}%
\pgfsetlinewidth{0.000000pt}%
\definecolor{currentstroke}{rgb}{0.000000,0.000000,0.000000}%
\pgfsetstrokecolor{currentstroke}%
\pgfsetdash{}{0pt}%
\pgfpathmoveto{\pgfqpoint{2.875117in}{0.648332in}}%
\pgfpathlineto{\pgfqpoint{2.900400in}{0.647637in}}%
\pgfpathlineto{\pgfqpoint{2.925683in}{0.646999in}}%
\pgfpathlineto{\pgfqpoint{2.950967in}{0.646414in}}%
\pgfpathlineto{\pgfqpoint{2.976250in}{0.645880in}}%
\pgfpathlineto{\pgfqpoint{3.001533in}{0.645394in}}%
\pgfpathlineto{\pgfqpoint{3.026816in}{0.644954in}}%
\pgfpathlineto{\pgfqpoint{3.052099in}{0.644558in}}%
\pgfpathlineto{\pgfqpoint{3.077382in}{0.644205in}}%
\pgfpathlineto{\pgfqpoint{3.102666in}{0.643893in}}%
\pgfpathlineto{\pgfqpoint{3.127949in}{0.643621in}}%
\pgfpathlineto{\pgfqpoint{3.153232in}{0.643389in}}%
\pgfpathlineto{\pgfqpoint{3.178515in}{0.643196in}}%
\pgfpathlineto{\pgfqpoint{3.203798in}{0.643041in}}%
\pgfpathlineto{\pgfqpoint{3.229082in}{0.642925in}}%
\pgfpathlineto{\pgfqpoint{3.254365in}{0.642846in}}%
\pgfpathlineto{\pgfqpoint{3.279648in}{0.642807in}}%
\pgfpathlineto{\pgfqpoint{3.304931in}{0.642806in}}%
\pgfpathlineto{\pgfqpoint{3.330214in}{0.642845in}}%
\pgfpathlineto{\pgfqpoint{3.355497in}{0.642924in}}%
\pgfpathlineto{\pgfqpoint{3.380781in}{0.643045in}}%
\pgfpathlineto{\pgfqpoint{3.406064in}{0.643210in}}%
\pgfpathlineto{\pgfqpoint{3.431347in}{0.643419in}}%
\pgfpathlineto{\pgfqpoint{3.456630in}{0.643675in}}%
\pgfpathlineto{\pgfqpoint{3.481913in}{0.643980in}}%
\pgfpathlineto{\pgfqpoint{3.507196in}{0.644336in}}%
\pgfpathlineto{\pgfqpoint{3.532480in}{0.644748in}}%
\pgfpathlineto{\pgfqpoint{3.557763in}{0.645217in}}%
\pgfpathlineto{\pgfqpoint{3.583046in}{0.645750in}}%
\pgfpathlineto{\pgfqpoint{3.608329in}{0.646348in}}%
\pgfpathlineto{\pgfqpoint{3.633612in}{0.647019in}}%
\pgfpathlineto{\pgfqpoint{3.658895in}{0.647768in}}%
\pgfpathlineto{\pgfqpoint{3.684179in}{0.648602in}}%
\pgfpathlineto{\pgfqpoint{3.688790in}{0.648767in}}%
\pgfpathlineto{\pgfqpoint{3.709462in}{0.650364in}}%
\pgfpathlineto{\pgfqpoint{3.734745in}{0.652465in}}%
\pgfpathlineto{\pgfqpoint{3.760028in}{0.654733in}}%
\pgfpathlineto{\pgfqpoint{3.785311in}{0.657187in}}%
\pgfpathlineto{\pgfqpoint{3.810594in}{0.659844in}}%
\pgfpathlineto{\pgfqpoint{3.829741in}{0.662015in}}%
\pgfpathlineto{\pgfqpoint{3.835878in}{0.662410in}}%
\pgfpathlineto{\pgfqpoint{3.861161in}{0.664227in}}%
\pgfpathlineto{\pgfqpoint{3.886444in}{0.666359in}}%
\pgfpathlineto{\pgfqpoint{3.911727in}{0.668871in}}%
\pgfpathlineto{\pgfqpoint{3.937010in}{0.671850in}}%
\pgfpathlineto{\pgfqpoint{3.961346in}{0.675263in}}%
\pgfpathlineto{\pgfqpoint{3.962293in}{0.675361in}}%
\pgfpathlineto{\pgfqpoint{3.987577in}{0.678096in}}%
\pgfpathlineto{\pgfqpoint{4.012860in}{0.681032in}}%
\pgfpathlineto{\pgfqpoint{4.038143in}{0.684194in}}%
\pgfpathlineto{\pgfqpoint{4.063426in}{0.687614in}}%
\pgfpathlineto{\pgfqpoint{4.069709in}{0.688511in}}%
\pgfpathlineto{\pgfqpoint{4.085119in}{0.701759in}}%
\pgfpathlineto{\pgfqpoint{4.088709in}{0.702642in}}%
\pgfpathlineto{\pgfqpoint{4.113993in}{0.709523in}}%
\pgfpathlineto{\pgfqpoint{4.131292in}{0.715007in}}%
\pgfpathlineto{\pgfqpoint{4.113993in}{0.723486in}}%
\pgfpathlineto{\pgfqpoint{4.105499in}{0.728255in}}%
\pgfpathlineto{\pgfqpoint{4.113993in}{0.731720in}}%
\pgfpathlineto{\pgfqpoint{4.135102in}{0.741503in}}%
\pgfpathlineto{\pgfqpoint{4.113993in}{0.749754in}}%
\pgfpathlineto{\pgfqpoint{4.102476in}{0.754751in}}%
\pgfpathlineto{\pgfqpoint{4.113993in}{0.760158in}}%
\pgfpathlineto{\pgfqpoint{4.128532in}{0.767999in}}%
\pgfpathlineto{\pgfqpoint{4.113993in}{0.775166in}}%
\pgfpathlineto{\pgfqpoint{4.102329in}{0.781247in}}%
\pgfpathlineto{\pgfqpoint{4.109897in}{0.794495in}}%
\pgfpathlineto{\pgfqpoint{4.088709in}{0.805793in}}%
\pgfpathlineto{\pgfqpoint{4.085272in}{0.807743in}}%
\pgfpathlineto{\pgfqpoint{4.082126in}{0.820991in}}%
\pgfpathlineto{\pgfqpoint{4.063426in}{0.829788in}}%
\pgfpathlineto{\pgfqpoint{4.054443in}{0.834239in}}%
\pgfpathlineto{\pgfqpoint{4.044457in}{0.847487in}}%
\pgfpathlineto{\pgfqpoint{4.038143in}{0.853720in}}%
\pgfpathlineto{\pgfqpoint{4.031398in}{0.860735in}}%
\pgfpathlineto{\pgfqpoint{4.020484in}{0.873983in}}%
\pgfpathlineto{\pgfqpoint{4.013582in}{0.887231in}}%
\pgfpathlineto{\pgfqpoint{4.012860in}{0.887755in}}%
\pgfpathlineto{\pgfqpoint{3.996982in}{0.900479in}}%
\pgfpathlineto{\pgfqpoint{3.987577in}{0.911741in}}%
\pgfpathlineto{\pgfqpoint{3.986062in}{0.913727in}}%
\pgfpathlineto{\pgfqpoint{3.964421in}{0.926975in}}%
\pgfpathlineto{\pgfqpoint{3.962293in}{0.933873in}}%
\pgfpathlineto{\pgfqpoint{3.960550in}{0.940223in}}%
\pgfpathlineto{\pgfqpoint{3.939937in}{0.953471in}}%
\pgfpathlineto{\pgfqpoint{3.937111in}{0.966719in}}%
\pgfpathlineto{\pgfqpoint{3.937010in}{0.966786in}}%
\pgfpathlineto{\pgfqpoint{3.918449in}{0.979967in}}%
\pgfpathlineto{\pgfqpoint{3.914600in}{0.993215in}}%
\pgfpathlineto{\pgfqpoint{3.911727in}{0.995018in}}%
\pgfpathlineto{\pgfqpoint{3.894905in}{1.006463in}}%
\pgfpathlineto{\pgfqpoint{3.890294in}{1.019711in}}%
\pgfpathlineto{\pgfqpoint{3.886444in}{1.022295in}}%
\pgfpathlineto{\pgfqpoint{3.871641in}{1.032959in}}%
\pgfpathlineto{\pgfqpoint{3.867538in}{1.046207in}}%
\pgfpathlineto{\pgfqpoint{3.861161in}{1.052839in}}%
\pgfpathlineto{\pgfqpoint{3.855351in}{1.059455in}}%
\pgfpathlineto{\pgfqpoint{3.850679in}{1.072703in}}%
\pgfpathlineto{\pgfqpoint{3.837850in}{1.085951in}}%
\pgfpathlineto{\pgfqpoint{3.835878in}{1.088653in}}%
\pgfpathlineto{\pgfqpoint{3.829354in}{1.099199in}}%
\pgfpathlineto{\pgfqpoint{3.817259in}{1.112447in}}%
\pgfpathlineto{\pgfqpoint{3.810594in}{1.120855in}}%
\pgfpathlineto{\pgfqpoint{3.807353in}{1.125695in}}%
\pgfpathlineto{\pgfqpoint{3.800947in}{1.138943in}}%
\pgfpathlineto{\pgfqpoint{3.791048in}{1.152191in}}%
\pgfpathlineto{\pgfqpoint{3.785311in}{1.164587in}}%
\pgfpathlineto{\pgfqpoint{3.784978in}{1.165439in}}%
\pgfpathlineto{\pgfqpoint{3.775242in}{1.178687in}}%
\pgfpathlineto{\pgfqpoint{3.770332in}{1.191935in}}%
\pgfpathlineto{\pgfqpoint{3.760028in}{1.202426in}}%
\pgfpathlineto{\pgfqpoint{3.757733in}{1.205183in}}%
\pgfpathlineto{\pgfqpoint{3.754147in}{1.218431in}}%
\pgfpathlineto{\pgfqpoint{3.742388in}{1.231679in}}%
\pgfpathlineto{\pgfqpoint{3.739896in}{1.244927in}}%
\pgfpathlineto{\pgfqpoint{3.734745in}{1.251212in}}%
\pgfpathlineto{\pgfqpoint{3.729933in}{1.258175in}}%
\pgfpathlineto{\pgfqpoint{3.726971in}{1.271423in}}%
\pgfpathlineto{\pgfqpoint{3.715061in}{1.284671in}}%
\pgfpathlineto{\pgfqpoint{3.710611in}{1.297919in}}%
\pgfpathlineto{\pgfqpoint{3.709462in}{1.299243in}}%
\pgfpathlineto{\pgfqpoint{3.700904in}{1.311167in}}%
\pgfpathlineto{\pgfqpoint{3.697716in}{1.324415in}}%
\pgfpathlineto{\pgfqpoint{3.690118in}{1.337663in}}%
\pgfpathlineto{\pgfqpoint{3.685658in}{1.350911in}}%
\pgfpathlineto{\pgfqpoint{3.684179in}{1.353196in}}%
\pgfpathlineto{\pgfqpoint{3.678589in}{1.364159in}}%
\pgfpathlineto{\pgfqpoint{3.673940in}{1.377407in}}%
\pgfpathlineto{\pgfqpoint{3.667068in}{1.390655in}}%
\pgfpathlineto{\pgfqpoint{3.659852in}{1.403903in}}%
\pgfpathlineto{\pgfqpoint{3.658895in}{1.406074in}}%
\pgfpathlineto{\pgfqpoint{3.655179in}{1.417151in}}%
\pgfpathlineto{\pgfqpoint{3.649238in}{1.430399in}}%
\pgfpathlineto{\pgfqpoint{3.645620in}{1.443647in}}%
\pgfpathlineto{\pgfqpoint{3.638575in}{1.456895in}}%
\pgfpathlineto{\pgfqpoint{3.634056in}{1.470143in}}%
\pgfpathlineto{\pgfqpoint{3.633612in}{1.470726in}}%
\pgfpathlineto{\pgfqpoint{3.626402in}{1.483391in}}%
\pgfpathlineto{\pgfqpoint{3.623832in}{1.496639in}}%
\pgfpathlineto{\pgfqpoint{3.616379in}{1.509887in}}%
\pgfpathlineto{\pgfqpoint{3.614511in}{1.523135in}}%
\pgfpathlineto{\pgfqpoint{3.608329in}{1.533554in}}%
\pgfpathlineto{\pgfqpoint{3.607060in}{1.536383in}}%
\pgfpathlineto{\pgfqpoint{3.604968in}{1.549631in}}%
\pgfpathlineto{\pgfqpoint{3.597616in}{1.562879in}}%
\pgfpathlineto{\pgfqpoint{3.595147in}{1.576127in}}%
\pgfpathlineto{\pgfqpoint{3.587079in}{1.589375in}}%
\pgfpathlineto{\pgfqpoint{3.583509in}{1.602623in}}%
\pgfpathlineto{\pgfqpoint{3.583046in}{1.603469in}}%
\pgfpathlineto{\pgfqpoint{3.578145in}{1.615871in}}%
\pgfpathlineto{\pgfqpoint{3.575819in}{1.629119in}}%
\pgfpathlineto{\pgfqpoint{3.570536in}{1.642367in}}%
\pgfpathlineto{\pgfqpoint{3.567042in}{1.655615in}}%
\pgfpathlineto{\pgfqpoint{3.560651in}{1.668863in}}%
\pgfpathlineto{\pgfqpoint{3.557763in}{1.675515in}}%
\pgfpathlineto{\pgfqpoint{3.555795in}{1.682111in}}%
\pgfpathlineto{\pgfqpoint{3.552373in}{1.695359in}}%
\pgfpathlineto{\pgfqpoint{3.548518in}{1.708607in}}%
\pgfpathlineto{\pgfqpoint{3.545277in}{1.721855in}}%
\pgfpathlineto{\pgfqpoint{3.539993in}{1.735103in}}%
\pgfpathlineto{\pgfqpoint{3.535778in}{1.748351in}}%
\pgfpathlineto{\pgfqpoint{3.532480in}{1.753981in}}%
\pgfpathlineto{\pgfqpoint{3.529369in}{1.761599in}}%
\pgfpathlineto{\pgfqpoint{3.527087in}{1.774847in}}%
\pgfpathlineto{\pgfqpoint{3.521393in}{1.788095in}}%
\pgfpathlineto{\pgfqpoint{3.519348in}{1.801343in}}%
\pgfpathlineto{\pgfqpoint{3.513214in}{1.814591in}}%
\pgfpathlineto{\pgfqpoint{3.511170in}{1.827839in}}%
\pgfpathlineto{\pgfqpoint{3.507196in}{1.834595in}}%
\pgfpathlineto{\pgfqpoint{3.504485in}{1.841087in}}%
\pgfpathlineto{\pgfqpoint{3.502572in}{1.854335in}}%
\pgfpathlineto{\pgfqpoint{3.496182in}{1.867583in}}%
\pgfpathlineto{\pgfqpoint{3.493902in}{1.880831in}}%
\pgfpathlineto{\pgfqpoint{3.487877in}{1.894079in}}%
\pgfpathlineto{\pgfqpoint{3.485449in}{1.907327in}}%
\pgfpathlineto{\pgfqpoint{3.481913in}{1.913895in}}%
\pgfpathlineto{\pgfqpoint{3.479283in}{1.920575in}}%
\pgfpathlineto{\pgfqpoint{3.476552in}{1.933823in}}%
\pgfpathlineto{\pgfqpoint{3.470471in}{1.947071in}}%
\pgfpathlineto{\pgfqpoint{3.466929in}{1.960319in}}%
\pgfpathlineto{\pgfqpoint{3.461734in}{1.973567in}}%
\pgfpathlineto{\pgfqpoint{3.457293in}{1.986815in}}%
\pgfpathlineto{\pgfqpoint{3.456630in}{1.988496in}}%
\pgfpathlineto{\pgfqpoint{3.453169in}{2.000063in}}%
\pgfpathlineto{\pgfqpoint{3.449324in}{2.013311in}}%
\pgfpathlineto{\pgfqpoint{3.444852in}{2.026559in}}%
\pgfpathlineto{\pgfqpoint{3.439026in}{2.039807in}}%
\pgfpathlineto{\pgfqpoint{3.434004in}{2.053055in}}%
\pgfpathlineto{\pgfqpoint{3.431347in}{2.057525in}}%
\pgfpathlineto{\pgfqpoint{3.426823in}{2.066303in}}%
\pgfpathlineto{\pgfqpoint{3.423274in}{2.079551in}}%
\pgfpathlineto{\pgfqpoint{3.416723in}{2.092799in}}%
\pgfpathlineto{\pgfqpoint{3.412943in}{2.106047in}}%
\pgfpathlineto{\pgfqpoint{3.406064in}{2.115989in}}%
\pgfpathlineto{\pgfqpoint{3.403616in}{2.119295in}}%
\pgfpathlineto{\pgfqpoint{3.398955in}{2.132543in}}%
\pgfpathlineto{\pgfqpoint{3.380781in}{2.132543in}}%
\pgfpathlineto{\pgfqpoint{3.363401in}{2.132543in}}%
\pgfpathlineto{\pgfqpoint{3.373326in}{2.119295in}}%
\pgfpathlineto{\pgfqpoint{3.380781in}{2.113321in}}%
\pgfpathlineto{\pgfqpoint{3.386723in}{2.106047in}}%
\pgfpathlineto{\pgfqpoint{3.392120in}{2.092799in}}%
\pgfpathlineto{\pgfqpoint{3.400713in}{2.079551in}}%
\pgfpathlineto{\pgfqpoint{3.406002in}{2.066303in}}%
\pgfpathlineto{\pgfqpoint{3.406064in}{2.066212in}}%
\pgfpathlineto{\pgfqpoint{3.413990in}{2.053055in}}%
\pgfpathlineto{\pgfqpoint{3.419500in}{2.039807in}}%
\pgfpathlineto{\pgfqpoint{3.425388in}{2.026559in}}%
\pgfpathlineto{\pgfqpoint{3.430617in}{2.013311in}}%
\pgfpathlineto{\pgfqpoint{3.431347in}{2.011168in}}%
\pgfpathlineto{\pgfqpoint{3.435417in}{2.000063in}}%
\pgfpathlineto{\pgfqpoint{3.440689in}{1.986815in}}%
\pgfpathlineto{\pgfqpoint{3.444547in}{1.973567in}}%
\pgfpathlineto{\pgfqpoint{3.449784in}{1.960319in}}%
\pgfpathlineto{\pgfqpoint{3.453071in}{1.947071in}}%
\pgfpathlineto{\pgfqpoint{3.456630in}{1.939876in}}%
\pgfpathlineto{\pgfqpoint{3.460053in}{1.933823in}}%
\pgfpathlineto{\pgfqpoint{3.463423in}{1.920575in}}%
\pgfpathlineto{\pgfqpoint{3.469762in}{1.907327in}}%
\pgfpathlineto{\pgfqpoint{3.471845in}{1.894079in}}%
\pgfpathlineto{\pgfqpoint{3.477510in}{1.880831in}}%
\pgfpathlineto{\pgfqpoint{3.479587in}{1.867583in}}%
\pgfpathlineto{\pgfqpoint{3.481913in}{1.862851in}}%
\pgfpathlineto{\pgfqpoint{3.486967in}{1.854335in}}%
\pgfpathlineto{\pgfqpoint{3.489269in}{1.841087in}}%
\pgfpathlineto{\pgfqpoint{3.495723in}{1.827839in}}%
\pgfpathlineto{\pgfqpoint{3.497498in}{1.814591in}}%
\pgfpathlineto{\pgfqpoint{3.502998in}{1.801343in}}%
\pgfpathlineto{\pgfqpoint{3.504929in}{1.788095in}}%
\pgfpathlineto{\pgfqpoint{3.507196in}{1.782651in}}%
\pgfpathlineto{\pgfqpoint{3.511309in}{1.774847in}}%
\pgfpathlineto{\pgfqpoint{3.514216in}{1.761599in}}%
\pgfpathlineto{\pgfqpoint{3.520433in}{1.748351in}}%
\pgfpathlineto{\pgfqpoint{3.524019in}{1.735103in}}%
\pgfpathlineto{\pgfqpoint{3.528593in}{1.721855in}}%
\pgfpathlineto{\pgfqpoint{3.531593in}{1.708607in}}%
\pgfpathlineto{\pgfqpoint{3.532480in}{1.705382in}}%
\pgfpathlineto{\pgfqpoint{3.536094in}{1.695359in}}%
\pgfpathlineto{\pgfqpoint{3.540530in}{1.682111in}}%
\pgfpathlineto{\pgfqpoint{3.544921in}{1.668863in}}%
\pgfpathlineto{\pgfqpoint{3.550325in}{1.655615in}}%
\pgfpathlineto{\pgfqpoint{3.553241in}{1.642367in}}%
\pgfpathlineto{\pgfqpoint{3.557763in}{1.630036in}}%
\pgfpathlineto{\pgfqpoint{3.558199in}{1.629119in}}%
\pgfpathlineto{\pgfqpoint{3.561130in}{1.615871in}}%
\pgfpathlineto{\pgfqpoint{3.567787in}{1.602623in}}%
\pgfpathlineto{\pgfqpoint{3.570453in}{1.589375in}}%
\pgfpathlineto{\pgfqpoint{3.577413in}{1.576127in}}%
\pgfpathlineto{\pgfqpoint{3.579451in}{1.562879in}}%
\pgfpathlineto{\pgfqpoint{3.583046in}{1.555956in}}%
\pgfpathlineto{\pgfqpoint{3.587171in}{1.549631in}}%
\pgfpathlineto{\pgfqpoint{3.589668in}{1.536383in}}%
\pgfpathlineto{\pgfqpoint{3.596508in}{1.523135in}}%
\pgfpathlineto{\pgfqpoint{3.598021in}{1.509887in}}%
\pgfpathlineto{\pgfqpoint{3.604675in}{1.496639in}}%
\pgfpathlineto{\pgfqpoint{3.607022in}{1.483391in}}%
\pgfpathlineto{\pgfqpoint{3.608329in}{1.480927in}}%
\pgfpathlineto{\pgfqpoint{3.615436in}{1.470143in}}%
\pgfpathlineto{\pgfqpoint{3.619277in}{1.456895in}}%
\pgfpathlineto{\pgfqpoint{3.625401in}{1.443647in}}%
\pgfpathlineto{\pgfqpoint{3.628768in}{1.430399in}}%
\pgfpathlineto{\pgfqpoint{3.633612in}{1.418686in}}%
\pgfpathlineto{\pgfqpoint{3.634407in}{1.417151in}}%
\pgfpathlineto{\pgfqpoint{3.639906in}{1.403903in}}%
\pgfpathlineto{\pgfqpoint{3.645742in}{1.390655in}}%
\pgfpathlineto{\pgfqpoint{3.652036in}{1.377407in}}%
\pgfpathlineto{\pgfqpoint{3.656007in}{1.364159in}}%
\pgfpathlineto{\pgfqpoint{3.658895in}{1.358385in}}%
\pgfpathlineto{\pgfqpoint{3.663371in}{1.350911in}}%
\pgfpathlineto{\pgfqpoint{3.667067in}{1.337663in}}%
\pgfpathlineto{\pgfqpoint{3.674341in}{1.324415in}}%
\pgfpathlineto{\pgfqpoint{3.676830in}{1.311167in}}%
\pgfpathlineto{\pgfqpoint{3.684179in}{1.300649in}}%
\pgfpathlineto{\pgfqpoint{3.686396in}{1.297919in}}%
\pgfpathlineto{\pgfqpoint{3.689988in}{1.284671in}}%
\pgfpathlineto{\pgfqpoint{3.701042in}{1.271423in}}%
\pgfpathlineto{\pgfqpoint{3.703492in}{1.258175in}}%
\pgfpathlineto{\pgfqpoint{3.709462in}{1.249217in}}%
\pgfpathlineto{\pgfqpoint{3.712762in}{1.244927in}}%
\pgfpathlineto{\pgfqpoint{3.714703in}{1.231679in}}%
\pgfpathlineto{\pgfqpoint{3.725853in}{1.218431in}}%
\pgfpathlineto{\pgfqpoint{3.728973in}{1.205183in}}%
\pgfpathlineto{\pgfqpoint{3.734745in}{1.197481in}}%
\pgfpathlineto{\pgfqpoint{3.739571in}{1.191935in}}%
\pgfpathlineto{\pgfqpoint{3.744066in}{1.178687in}}%
\pgfpathlineto{\pgfqpoint{3.753159in}{1.165439in}}%
\pgfpathlineto{\pgfqpoint{3.758120in}{1.152191in}}%
\pgfpathlineto{\pgfqpoint{3.760028in}{1.148946in}}%
\pgfpathlineto{\pgfqpoint{3.766810in}{1.138943in}}%
\pgfpathlineto{\pgfqpoint{3.773231in}{1.125695in}}%
\pgfpathlineto{\pgfqpoint{3.780725in}{1.112447in}}%
\pgfpathlineto{\pgfqpoint{3.785311in}{1.106786in}}%
\pgfpathlineto{\pgfqpoint{3.792073in}{1.099199in}}%
\pgfpathlineto{\pgfqpoint{3.798849in}{1.085951in}}%
\pgfpathlineto{\pgfqpoint{3.809783in}{1.072703in}}%
\pgfpathlineto{\pgfqpoint{3.810594in}{1.068664in}}%
\pgfpathlineto{\pgfqpoint{3.812754in}{1.059455in}}%
\pgfpathlineto{\pgfqpoint{3.824517in}{1.046207in}}%
\pgfpathlineto{\pgfqpoint{3.826230in}{1.032959in}}%
\pgfpathlineto{\pgfqpoint{3.835878in}{1.025297in}}%
\pgfpathlineto{\pgfqpoint{3.843404in}{1.019711in}}%
\pgfpathlineto{\pgfqpoint{3.845401in}{1.006463in}}%
\pgfpathlineto{\pgfqpoint{3.861161in}{0.994154in}}%
\pgfpathlineto{\pgfqpoint{3.862462in}{0.993215in}}%
\pgfpathlineto{\pgfqpoint{3.864026in}{0.979967in}}%
\pgfpathlineto{\pgfqpoint{3.880811in}{0.966719in}}%
\pgfpathlineto{\pgfqpoint{3.880987in}{0.953471in}}%
\pgfpathlineto{\pgfqpoint{3.886444in}{0.948937in}}%
\pgfpathlineto{\pgfqpoint{3.897917in}{0.940223in}}%
\pgfpathlineto{\pgfqpoint{3.899331in}{0.926975in}}%
\pgfpathlineto{\pgfqpoint{3.911727in}{0.916367in}}%
\pgfpathlineto{\pgfqpoint{3.915159in}{0.913727in}}%
\pgfpathlineto{\pgfqpoint{3.922842in}{0.900479in}}%
\pgfpathlineto{\pgfqpoint{3.933399in}{0.887231in}}%
\pgfpathlineto{\pgfqpoint{3.936658in}{0.873983in}}%
\pgfpathlineto{\pgfqpoint{3.937010in}{0.872269in}}%
\pgfpathlineto{\pgfqpoint{3.939722in}{0.860735in}}%
\pgfpathlineto{\pgfqpoint{3.950556in}{0.847487in}}%
\pgfpathlineto{\pgfqpoint{3.949401in}{0.834239in}}%
\pgfpathlineto{\pgfqpoint{3.962293in}{0.826151in}}%
\pgfpathlineto{\pgfqpoint{3.971032in}{0.820991in}}%
\pgfpathlineto{\pgfqpoint{3.962293in}{0.809664in}}%
\pgfpathlineto{\pgfqpoint{3.960634in}{0.807743in}}%
\pgfpathlineto{\pgfqpoint{3.962293in}{0.806362in}}%
\pgfpathlineto{\pgfqpoint{3.977533in}{0.794495in}}%
\pgfpathlineto{\pgfqpoint{3.962293in}{0.786297in}}%
\pgfpathlineto{\pgfqpoint{3.951407in}{0.781247in}}%
\pgfpathlineto{\pgfqpoint{3.962293in}{0.771917in}}%
\pgfpathlineto{\pgfqpoint{3.967242in}{0.767999in}}%
\pgfpathlineto{\pgfqpoint{3.962293in}{0.766552in}}%
\pgfpathlineto{\pgfqpoint{3.937010in}{0.759861in}}%
\pgfpathlineto{\pgfqpoint{3.914468in}{0.754751in}}%
\pgfpathlineto{\pgfqpoint{3.922241in}{0.741503in}}%
\pgfpathlineto{\pgfqpoint{3.911727in}{0.738976in}}%
\pgfpathlineto{\pgfqpoint{3.886444in}{0.733444in}}%
\pgfpathlineto{\pgfqpoint{3.861161in}{0.728554in}}%
\pgfpathlineto{\pgfqpoint{3.859498in}{0.728255in}}%
\pgfpathlineto{\pgfqpoint{3.840712in}{0.715007in}}%
\pgfpathlineto{\pgfqpoint{3.835878in}{0.714314in}}%
\pgfpathlineto{\pgfqpoint{3.810594in}{0.710944in}}%
\pgfpathlineto{\pgfqpoint{3.785311in}{0.707896in}}%
\pgfpathlineto{\pgfqpoint{3.760028in}{0.705128in}}%
\pgfpathlineto{\pgfqpoint{3.734745in}{0.702605in}}%
\pgfpathlineto{\pgfqpoint{3.725623in}{0.701759in}}%
\pgfpathlineto{\pgfqpoint{3.709462in}{0.699989in}}%
\pgfpathlineto{\pgfqpoint{3.684179in}{0.697589in}}%
\pgfpathlineto{\pgfqpoint{3.658895in}{0.695555in}}%
\pgfpathlineto{\pgfqpoint{3.633612in}{0.693833in}}%
\pgfpathlineto{\pgfqpoint{3.608329in}{0.692377in}}%
\pgfpathlineto{\pgfqpoint{3.583046in}{0.691151in}}%
\pgfpathlineto{\pgfqpoint{3.557763in}{0.690123in}}%
\pgfpathlineto{\pgfqpoint{3.532480in}{0.689270in}}%
\pgfpathlineto{\pgfqpoint{3.507196in}{0.688571in}}%
\pgfpathlineto{\pgfqpoint{3.504591in}{0.688511in}}%
\pgfpathlineto{\pgfqpoint{3.481913in}{0.687767in}}%
\pgfpathlineto{\pgfqpoint{3.456630in}{0.687047in}}%
\pgfpathlineto{\pgfqpoint{3.431347in}{0.686439in}}%
\pgfpathlineto{\pgfqpoint{3.406064in}{0.685946in}}%
\pgfpathlineto{\pgfqpoint{3.380781in}{0.685574in}}%
\pgfpathlineto{\pgfqpoint{3.355497in}{0.685328in}}%
\pgfpathlineto{\pgfqpoint{3.330214in}{0.685219in}}%
\pgfpathlineto{\pgfqpoint{3.304931in}{0.685259in}}%
\pgfpathlineto{\pgfqpoint{3.279648in}{0.685461in}}%
\pgfpathlineto{\pgfqpoint{3.254365in}{0.685845in}}%
\pgfpathlineto{\pgfqpoint{3.229082in}{0.686433in}}%
\pgfpathlineto{\pgfqpoint{3.203798in}{0.687254in}}%
\pgfpathlineto{\pgfqpoint{3.178515in}{0.688344in}}%
\pgfpathlineto{\pgfqpoint{3.175358in}{0.688511in}}%
\pgfpathlineto{\pgfqpoint{3.153232in}{0.688943in}}%
\pgfpathlineto{\pgfqpoint{3.127949in}{0.689499in}}%
\pgfpathlineto{\pgfqpoint{3.102666in}{0.690121in}}%
\pgfpathlineto{\pgfqpoint{3.077382in}{0.690808in}}%
\pgfpathlineto{\pgfqpoint{3.052099in}{0.691564in}}%
\pgfpathlineto{\pgfqpoint{3.026816in}{0.692391in}}%
\pgfpathlineto{\pgfqpoint{3.001533in}{0.693293in}}%
\pgfpathlineto{\pgfqpoint{2.976250in}{0.694275in}}%
\pgfpathlineto{\pgfqpoint{2.950967in}{0.695342in}}%
\pgfpathlineto{\pgfqpoint{2.925683in}{0.696499in}}%
\pgfpathlineto{\pgfqpoint{2.900400in}{0.697754in}}%
\pgfpathlineto{\pgfqpoint{2.875117in}{0.699115in}}%
\pgfpathlineto{\pgfqpoint{2.849834in}{0.700592in}}%
\pgfpathlineto{\pgfqpoint{2.831339in}{0.701759in}}%
\pgfpathlineto{\pgfqpoint{2.824551in}{0.708331in}}%
\pgfpathlineto{\pgfqpoint{2.819374in}{0.715007in}}%
\pgfpathlineto{\pgfqpoint{2.799268in}{0.717266in}}%
\pgfpathlineto{\pgfqpoint{2.773984in}{0.720340in}}%
\pgfpathlineto{\pgfqpoint{2.748701in}{0.723700in}}%
\pgfpathlineto{\pgfqpoint{2.723418in}{0.727388in}}%
\pgfpathlineto{\pgfqpoint{2.717846in}{0.728255in}}%
\pgfpathlineto{\pgfqpoint{2.723418in}{0.730445in}}%
\pgfpathlineto{\pgfqpoint{2.747170in}{0.741503in}}%
\pgfpathlineto{\pgfqpoint{2.723418in}{0.747601in}}%
\pgfpathlineto{\pgfqpoint{2.699337in}{0.754751in}}%
\pgfpathlineto{\pgfqpoint{2.717936in}{0.767999in}}%
\pgfpathlineto{\pgfqpoint{2.698135in}{0.774005in}}%
\pgfpathlineto{\pgfqpoint{2.676661in}{0.781247in}}%
\pgfpathlineto{\pgfqpoint{2.693320in}{0.794495in}}%
\pgfpathlineto{\pgfqpoint{2.674942in}{0.807743in}}%
\pgfpathlineto{\pgfqpoint{2.697782in}{0.820991in}}%
\pgfpathlineto{\pgfqpoint{2.694452in}{0.834239in}}%
\pgfpathlineto{\pgfqpoint{2.698135in}{0.837299in}}%
\pgfpathlineto{\pgfqpoint{2.708800in}{0.847487in}}%
\pgfpathlineto{\pgfqpoint{2.709929in}{0.860735in}}%
\pgfpathlineto{\pgfqpoint{2.721027in}{0.873983in}}%
\pgfpathlineto{\pgfqpoint{2.723418in}{0.878238in}}%
\pgfpathlineto{\pgfqpoint{2.728275in}{0.887231in}}%
\pgfpathlineto{\pgfqpoint{2.737186in}{0.900479in}}%
\pgfpathlineto{\pgfqpoint{2.748701in}{0.911793in}}%
\pgfpathlineto{\pgfqpoint{2.750596in}{0.913727in}}%
\pgfpathlineto{\pgfqpoint{2.757743in}{0.926975in}}%
\pgfpathlineto{\pgfqpoint{2.773984in}{0.938296in}}%
\pgfpathlineto{\pgfqpoint{2.776644in}{0.940223in}}%
\pgfpathlineto{\pgfqpoint{2.779636in}{0.953471in}}%
\pgfpathlineto{\pgfqpoint{2.794094in}{0.966719in}}%
\pgfpathlineto{\pgfqpoint{2.793216in}{0.979967in}}%
\pgfpathlineto{\pgfqpoint{2.799268in}{0.984646in}}%
\pgfpathlineto{\pgfqpoint{2.809871in}{0.993215in}}%
\pgfpathlineto{\pgfqpoint{2.811091in}{1.006463in}}%
\pgfpathlineto{\pgfqpoint{2.824551in}{1.014859in}}%
\pgfpathlineto{\pgfqpoint{2.831860in}{1.019711in}}%
\pgfpathlineto{\pgfqpoint{2.832402in}{1.032959in}}%
\pgfpathlineto{\pgfqpoint{2.849366in}{1.046207in}}%
\pgfpathlineto{\pgfqpoint{2.849813in}{1.059455in}}%
\pgfpathlineto{\pgfqpoint{2.849834in}{1.059472in}}%
\pgfpathlineto{\pgfqpoint{2.865366in}{1.072703in}}%
\pgfpathlineto{\pgfqpoint{2.867280in}{1.085951in}}%
\pgfpathlineto{\pgfqpoint{2.875117in}{1.092400in}}%
\pgfpathlineto{\pgfqpoint{2.882457in}{1.099199in}}%
\pgfpathlineto{\pgfqpoint{2.885630in}{1.112447in}}%
\pgfpathlineto{\pgfqpoint{2.898694in}{1.125695in}}%
\pgfpathlineto{\pgfqpoint{2.900400in}{1.129748in}}%
\pgfpathlineto{\pgfqpoint{2.903824in}{1.138943in}}%
\pgfpathlineto{\pgfqpoint{2.912161in}{1.152191in}}%
\pgfpathlineto{\pgfqpoint{2.916381in}{1.165439in}}%
\pgfpathlineto{\pgfqpoint{2.924252in}{1.178687in}}%
\pgfpathlineto{\pgfqpoint{2.925683in}{1.180883in}}%
\pgfpathlineto{\pgfqpoint{2.932132in}{1.191935in}}%
\pgfpathlineto{\pgfqpoint{2.939347in}{1.205183in}}%
\pgfpathlineto{\pgfqpoint{2.950199in}{1.218431in}}%
\pgfpathlineto{\pgfqpoint{2.950967in}{1.220808in}}%
\pgfpathlineto{\pgfqpoint{2.953829in}{1.231679in}}%
\pgfpathlineto{\pgfqpoint{2.962797in}{1.244927in}}%
\pgfpathlineto{\pgfqpoint{2.965613in}{1.258175in}}%
\pgfpathlineto{\pgfqpoint{2.976250in}{1.271261in}}%
\pgfpathlineto{\pgfqpoint{2.976366in}{1.271423in}}%
\pgfpathlineto{\pgfqpoint{2.978285in}{1.284671in}}%
\pgfpathlineto{\pgfqpoint{2.988511in}{1.297919in}}%
\pgfpathlineto{\pgfqpoint{2.990461in}{1.311167in}}%
\pgfpathlineto{\pgfqpoint{3.001434in}{1.324415in}}%
\pgfpathlineto{\pgfqpoint{3.001533in}{1.324968in}}%
\pgfpathlineto{\pgfqpoint{3.003363in}{1.337663in}}%
\pgfpathlineto{\pgfqpoint{3.011322in}{1.350911in}}%
\pgfpathlineto{\pgfqpoint{3.013305in}{1.364159in}}%
\pgfpathlineto{\pgfqpoint{3.022835in}{1.377407in}}%
\pgfpathlineto{\pgfqpoint{3.026816in}{1.388890in}}%
\pgfpathlineto{\pgfqpoint{3.027316in}{1.390655in}}%
\pgfpathlineto{\pgfqpoint{3.034545in}{1.403903in}}%
\pgfpathlineto{\pgfqpoint{3.038667in}{1.417151in}}%
\pgfpathlineto{\pgfqpoint{3.043732in}{1.430399in}}%
\pgfpathlineto{\pgfqpoint{3.048723in}{1.443647in}}%
\pgfpathlineto{\pgfqpoint{3.052099in}{1.451982in}}%
\pgfpathlineto{\pgfqpoint{3.053639in}{1.456895in}}%
\pgfpathlineto{\pgfqpoint{3.059118in}{1.470143in}}%
\pgfpathlineto{\pgfqpoint{3.062690in}{1.483391in}}%
\pgfpathlineto{\pgfqpoint{3.070014in}{1.496639in}}%
\pgfpathlineto{\pgfqpoint{3.073089in}{1.509887in}}%
\pgfpathlineto{\pgfqpoint{3.077382in}{1.516788in}}%
\pgfpathlineto{\pgfqpoint{3.080615in}{1.523135in}}%
\pgfpathlineto{\pgfqpoint{3.082515in}{1.536383in}}%
\pgfpathlineto{\pgfqpoint{3.089095in}{1.549631in}}%
\pgfpathlineto{\pgfqpoint{3.091020in}{1.562879in}}%
\pgfpathlineto{\pgfqpoint{3.099793in}{1.576127in}}%
\pgfpathlineto{\pgfqpoint{3.102335in}{1.589375in}}%
\pgfpathlineto{\pgfqpoint{3.102666in}{1.589934in}}%
\pgfpathlineto{\pgfqpoint{3.108468in}{1.602623in}}%
\pgfpathlineto{\pgfqpoint{3.110099in}{1.615871in}}%
\pgfpathlineto{\pgfqpoint{3.115570in}{1.629119in}}%
\pgfpathlineto{\pgfqpoint{3.118583in}{1.642367in}}%
\pgfpathlineto{\pgfqpoint{3.125175in}{1.655615in}}%
\pgfpathlineto{\pgfqpoint{3.127949in}{1.663848in}}%
\pgfpathlineto{\pgfqpoint{3.129259in}{1.668863in}}%
\pgfpathlineto{\pgfqpoint{3.133666in}{1.682111in}}%
\pgfpathlineto{\pgfqpoint{3.138307in}{1.695359in}}%
\pgfpathlineto{\pgfqpoint{3.141891in}{1.708607in}}%
\pgfpathlineto{\pgfqpoint{3.147004in}{1.721855in}}%
\pgfpathlineto{\pgfqpoint{3.149920in}{1.735103in}}%
\pgfpathlineto{\pgfqpoint{3.153232in}{1.742117in}}%
\pgfpathlineto{\pgfqpoint{3.155567in}{1.748351in}}%
\pgfpathlineto{\pgfqpoint{3.157960in}{1.761599in}}%
\pgfpathlineto{\pgfqpoint{3.164548in}{1.774847in}}%
\pgfpathlineto{\pgfqpoint{3.166403in}{1.788095in}}%
\pgfpathlineto{\pgfqpoint{3.172744in}{1.801343in}}%
\pgfpathlineto{\pgfqpoint{3.174648in}{1.814591in}}%
\pgfpathlineto{\pgfqpoint{3.178515in}{1.821472in}}%
\pgfpathlineto{\pgfqpoint{3.181415in}{1.827839in}}%
\pgfpathlineto{\pgfqpoint{3.183700in}{1.841087in}}%
\pgfpathlineto{\pgfqpoint{3.189993in}{1.854335in}}%
\pgfpathlineto{\pgfqpoint{3.192662in}{1.867583in}}%
\pgfpathlineto{\pgfqpoint{3.198426in}{1.880831in}}%
\pgfpathlineto{\pgfqpoint{3.202308in}{1.894079in}}%
\pgfpathlineto{\pgfqpoint{3.203798in}{1.897983in}}%
\pgfpathlineto{\pgfqpoint{3.206652in}{1.907327in}}%
\pgfpathlineto{\pgfqpoint{3.210136in}{1.920575in}}%
\pgfpathlineto{\pgfqpoint{3.214167in}{1.933823in}}%
\pgfpathlineto{\pgfqpoint{3.219926in}{1.947071in}}%
\pgfpathlineto{\pgfqpoint{3.224133in}{1.960319in}}%
\pgfpathlineto{\pgfqpoint{3.229082in}{1.969215in}}%
\pgfpathlineto{\pgfqpoint{3.231323in}{1.973567in}}%
\pgfpathlineto{\pgfqpoint{3.233732in}{1.986815in}}%
\pgfpathlineto{\pgfqpoint{3.240259in}{2.000063in}}%
\pgfpathlineto{\pgfqpoint{3.242818in}{2.013311in}}%
\pgfpathlineto{\pgfqpoint{3.251313in}{2.026559in}}%
\pgfpathlineto{\pgfqpoint{3.254072in}{2.039807in}}%
\pgfpathlineto{\pgfqpoint{3.254365in}{2.040254in}}%
\pgfpathlineto{\pgfqpoint{3.262872in}{2.053055in}}%
\pgfpathlineto{\pgfqpoint{3.266329in}{2.066303in}}%
\pgfpathlineto{\pgfqpoint{3.274999in}{2.079551in}}%
\pgfpathlineto{\pgfqpoint{3.279648in}{2.092006in}}%
\pgfpathlineto{\pgfqpoint{3.280031in}{2.092799in}}%
\pgfpathlineto{\pgfqpoint{3.290043in}{2.106047in}}%
\pgfpathlineto{\pgfqpoint{3.298614in}{2.119295in}}%
\pgfpathlineto{\pgfqpoint{3.304931in}{2.127166in}}%
\pgfpathlineto{\pgfqpoint{3.313212in}{2.132543in}}%
\pgfpathlineto{\pgfqpoint{3.304931in}{2.132543in}}%
\pgfpathlineto{\pgfqpoint{3.279648in}{2.132543in}}%
\pgfpathlineto{\pgfqpoint{3.277289in}{2.132543in}}%
\pgfpathlineto{\pgfqpoint{3.268675in}{2.119295in}}%
\pgfpathlineto{\pgfqpoint{3.263832in}{2.106047in}}%
\pgfpathlineto{\pgfqpoint{3.256651in}{2.092799in}}%
\pgfpathlineto{\pgfqpoint{3.254365in}{2.084128in}}%
\pgfpathlineto{\pgfqpoint{3.253060in}{2.079551in}}%
\pgfpathlineto{\pgfqpoint{3.245188in}{2.066303in}}%
\pgfpathlineto{\pgfqpoint{3.242497in}{2.053055in}}%
\pgfpathlineto{\pgfqpoint{3.234963in}{2.039807in}}%
\pgfpathlineto{\pgfqpoint{3.232673in}{2.026559in}}%
\pgfpathlineto{\pgfqpoint{3.229082in}{2.019813in}}%
\pgfpathlineto{\pgfqpoint{3.224970in}{2.013311in}}%
\pgfpathlineto{\pgfqpoint{3.222184in}{2.000063in}}%
\pgfpathlineto{\pgfqpoint{3.216176in}{1.986815in}}%
\pgfpathlineto{\pgfqpoint{3.213678in}{1.973567in}}%
\pgfpathlineto{\pgfqpoint{3.207899in}{1.960319in}}%
\pgfpathlineto{\pgfqpoint{3.204154in}{1.947071in}}%
\pgfpathlineto{\pgfqpoint{3.203798in}{1.946074in}}%
\pgfpathlineto{\pgfqpoint{3.198169in}{1.933823in}}%
\pgfpathlineto{\pgfqpoint{3.193374in}{1.920575in}}%
\pgfpathlineto{\pgfqpoint{3.190145in}{1.907327in}}%
\pgfpathlineto{\pgfqpoint{3.185922in}{1.894079in}}%
\pgfpathlineto{\pgfqpoint{3.183113in}{1.880831in}}%
\pgfpathlineto{\pgfqpoint{3.178515in}{1.868379in}}%
\pgfpathlineto{\pgfqpoint{3.178094in}{1.867583in}}%
\pgfpathlineto{\pgfqpoint{3.175004in}{1.854335in}}%
\pgfpathlineto{\pgfqpoint{3.167438in}{1.841087in}}%
\pgfpathlineto{\pgfqpoint{3.164971in}{1.827839in}}%
\pgfpathlineto{\pgfqpoint{3.158999in}{1.814591in}}%
\pgfpathlineto{\pgfqpoint{3.157495in}{1.801343in}}%
\pgfpathlineto{\pgfqpoint{3.153232in}{1.790237in}}%
\pgfpathlineto{\pgfqpoint{3.152076in}{1.788095in}}%
\pgfpathlineto{\pgfqpoint{3.149575in}{1.774847in}}%
\pgfpathlineto{\pgfqpoint{3.141934in}{1.761599in}}%
\pgfpathlineto{\pgfqpoint{3.138927in}{1.748351in}}%
\pgfpathlineto{\pgfqpoint{3.134089in}{1.735103in}}%
\pgfpathlineto{\pgfqpoint{3.131445in}{1.721855in}}%
\pgfpathlineto{\pgfqpoint{3.127949in}{1.709947in}}%
\pgfpathlineto{\pgfqpoint{3.127389in}{1.708607in}}%
\pgfpathlineto{\pgfqpoint{3.122611in}{1.695359in}}%
\pgfpathlineto{\pgfqpoint{3.117479in}{1.682111in}}%
\pgfpathlineto{\pgfqpoint{3.112122in}{1.668863in}}%
\pgfpathlineto{\pgfqpoint{3.108693in}{1.655615in}}%
\pgfpathlineto{\pgfqpoint{3.103024in}{1.642367in}}%
\pgfpathlineto{\pgfqpoint{3.102666in}{1.640358in}}%
\pgfpathlineto{\pgfqpoint{3.099787in}{1.629119in}}%
\pgfpathlineto{\pgfqpoint{3.093276in}{1.615871in}}%
\pgfpathlineto{\pgfqpoint{3.091513in}{1.602623in}}%
\pgfpathlineto{\pgfqpoint{3.084651in}{1.589375in}}%
\pgfpathlineto{\pgfqpoint{3.082502in}{1.576127in}}%
\pgfpathlineto{\pgfqpoint{3.077382in}{1.567103in}}%
\pgfpathlineto{\pgfqpoint{3.074166in}{1.562879in}}%
\pgfpathlineto{\pgfqpoint{3.071592in}{1.549631in}}%
\pgfpathlineto{\pgfqpoint{3.064341in}{1.536383in}}%
\pgfpathlineto{\pgfqpoint{3.061946in}{1.523135in}}%
\pgfpathlineto{\pgfqpoint{3.055228in}{1.509887in}}%
\pgfpathlineto{\pgfqpoint{3.052362in}{1.496639in}}%
\pgfpathlineto{\pgfqpoint{3.052099in}{1.496097in}}%
\pgfpathlineto{\pgfqpoint{3.044207in}{1.483391in}}%
\pgfpathlineto{\pgfqpoint{3.039486in}{1.470143in}}%
\pgfpathlineto{\pgfqpoint{3.033901in}{1.456895in}}%
\pgfpathlineto{\pgfqpoint{3.028987in}{1.443647in}}%
\pgfpathlineto{\pgfqpoint{3.026816in}{1.436141in}}%
\pgfpathlineto{\pgfqpoint{3.024642in}{1.430399in}}%
\pgfpathlineto{\pgfqpoint{3.018387in}{1.417151in}}%
\pgfpathlineto{\pgfqpoint{3.013997in}{1.403903in}}%
\pgfpathlineto{\pgfqpoint{3.005598in}{1.390655in}}%
\pgfpathlineto{\pgfqpoint{3.001920in}{1.377407in}}%
\pgfpathlineto{\pgfqpoint{3.001533in}{1.376830in}}%
\pgfpathlineto{\pgfqpoint{2.990964in}{1.364159in}}%
\pgfpathlineto{\pgfqpoint{2.988877in}{1.350911in}}%
\pgfpathlineto{\pgfqpoint{2.980411in}{1.337663in}}%
\pgfpathlineto{\pgfqpoint{2.978270in}{1.324415in}}%
\pgfpathlineto{\pgfqpoint{2.976250in}{1.321709in}}%
\pgfpathlineto{\pgfqpoint{2.966827in}{1.311167in}}%
\pgfpathlineto{\pgfqpoint{2.964252in}{1.297919in}}%
\pgfpathlineto{\pgfqpoint{2.953429in}{1.284671in}}%
\pgfpathlineto{\pgfqpoint{2.950967in}{1.272997in}}%
\pgfpathlineto{\pgfqpoint{2.950527in}{1.271423in}}%
\pgfpathlineto{\pgfqpoint{2.939516in}{1.258175in}}%
\pgfpathlineto{\pgfqpoint{2.935688in}{1.244927in}}%
\pgfpathlineto{\pgfqpoint{2.926997in}{1.231679in}}%
\pgfpathlineto{\pgfqpoint{2.925683in}{1.227850in}}%
\pgfpathlineto{\pgfqpoint{2.921597in}{1.218431in}}%
\pgfpathlineto{\pgfqpoint{2.910605in}{1.205183in}}%
\pgfpathlineto{\pgfqpoint{2.901816in}{1.191935in}}%
\pgfpathlineto{\pgfqpoint{2.900400in}{1.189422in}}%
\pgfpathlineto{\pgfqpoint{2.893465in}{1.178687in}}%
\pgfpathlineto{\pgfqpoint{2.884051in}{1.165439in}}%
\pgfpathlineto{\pgfqpoint{2.880289in}{1.152191in}}%
\pgfpathlineto{\pgfqpoint{2.875117in}{1.145019in}}%
\pgfpathlineto{\pgfqpoint{2.869941in}{1.138943in}}%
\pgfpathlineto{\pgfqpoint{2.864268in}{1.125695in}}%
\pgfpathlineto{\pgfqpoint{2.849834in}{1.112530in}}%
\pgfpathlineto{\pgfqpoint{2.849729in}{1.112447in}}%
\pgfpathlineto{\pgfqpoint{2.845693in}{1.099199in}}%
\pgfpathlineto{\pgfqpoint{2.828021in}{1.085951in}}%
\pgfpathlineto{\pgfqpoint{2.825384in}{1.072703in}}%
\pgfpathlineto{\pgfqpoint{2.824551in}{1.072031in}}%
\pgfpathlineto{\pgfqpoint{2.807669in}{1.059455in}}%
\pgfpathlineto{\pgfqpoint{2.806022in}{1.046207in}}%
\pgfpathlineto{\pgfqpoint{2.799268in}{1.041102in}}%
\pgfpathlineto{\pgfqpoint{2.787680in}{1.032959in}}%
\pgfpathlineto{\pgfqpoint{2.784837in}{1.019711in}}%
\pgfpathlineto{\pgfqpoint{2.773984in}{1.013222in}}%
\pgfpathlineto{\pgfqpoint{2.762033in}{1.006463in}}%
\pgfpathlineto{\pgfqpoint{2.757552in}{0.993215in}}%
\pgfpathlineto{\pgfqpoint{2.748701in}{0.986257in}}%
\pgfpathlineto{\pgfqpoint{2.740364in}{0.979967in}}%
\pgfpathlineto{\pgfqpoint{2.736968in}{0.966719in}}%
\pgfpathlineto{\pgfqpoint{2.723418in}{0.954708in}}%
\pgfpathlineto{\pgfqpoint{2.721950in}{0.953471in}}%
\pgfpathlineto{\pgfqpoint{2.713136in}{0.940223in}}%
\pgfpathlineto{\pgfqpoint{2.698135in}{0.930199in}}%
\pgfpathlineto{\pgfqpoint{2.693125in}{0.926975in}}%
\pgfpathlineto{\pgfqpoint{2.678430in}{0.913727in}}%
\pgfpathlineto{\pgfqpoint{2.672852in}{0.908279in}}%
\pgfpathlineto{\pgfqpoint{2.664528in}{0.900479in}}%
\pgfpathlineto{\pgfqpoint{2.647569in}{0.887854in}}%
\pgfpathlineto{\pgfqpoint{2.646610in}{0.887231in}}%
\pgfpathlineto{\pgfqpoint{2.638907in}{0.873983in}}%
\pgfpathlineto{\pgfqpoint{2.622285in}{0.862576in}}%
\pgfpathlineto{\pgfqpoint{2.619304in}{0.860735in}}%
\pgfpathlineto{\pgfqpoint{2.616099in}{0.847487in}}%
\pgfpathlineto{\pgfqpoint{2.597002in}{0.837660in}}%
\pgfpathlineto{\pgfqpoint{2.589533in}{0.834239in}}%
\pgfpathlineto{\pgfqpoint{2.588239in}{0.820991in}}%
\pgfpathlineto{\pgfqpoint{2.571719in}{0.815030in}}%
\pgfpathlineto{\pgfqpoint{2.549179in}{0.807743in}}%
\pgfpathlineto{\pgfqpoint{2.561726in}{0.794495in}}%
\pgfpathlineto{\pgfqpoint{2.546436in}{0.788103in}}%
\pgfpathlineto{\pgfqpoint{2.528472in}{0.781247in}}%
\pgfpathlineto{\pgfqpoint{2.546436in}{0.771964in}}%
\pgfpathlineto{\pgfqpoint{2.554851in}{0.767999in}}%
\pgfpathlineto{\pgfqpoint{2.546436in}{0.765110in}}%
\pgfpathlineto{\pgfqpoint{2.521153in}{0.757180in}}%
\pgfpathlineto{\pgfqpoint{2.512766in}{0.754751in}}%
\pgfpathlineto{\pgfqpoint{2.521153in}{0.749143in}}%
\pgfpathlineto{\pgfqpoint{2.533995in}{0.741503in}}%
\pgfpathlineto{\pgfqpoint{2.521153in}{0.738477in}}%
\pgfpathlineto{\pgfqpoint{2.495870in}{0.732941in}}%
\pgfpathlineto{\pgfqpoint{2.472762in}{0.728255in}}%
\pgfpathlineto{\pgfqpoint{2.495870in}{0.721982in}}%
\pgfpathlineto{\pgfqpoint{2.521153in}{0.716062in}}%
\pgfpathlineto{\pgfqpoint{2.526063in}{0.715007in}}%
\pgfpathlineto{\pgfqpoint{2.521153in}{0.713368in}}%
\pgfpathlineto{\pgfqpoint{2.495870in}{0.705596in}}%
\pgfpathlineto{\pgfqpoint{2.482168in}{0.701759in}}%
\pgfpathlineto{\pgfqpoint{2.495870in}{0.699841in}}%
\pgfpathlineto{\pgfqpoint{2.521153in}{0.696615in}}%
\pgfpathlineto{\pgfqpoint{2.546436in}{0.693740in}}%
\pgfpathlineto{\pgfqpoint{2.571719in}{0.691163in}}%
\pgfpathlineto{\pgfqpoint{2.597002in}{0.688840in}}%
\pgfpathlineto{\pgfqpoint{2.600808in}{0.688511in}}%
\pgfpathlineto{\pgfqpoint{2.597002in}{0.685847in}}%
\pgfpathlineto{\pgfqpoint{2.579801in}{0.675263in}}%
\pgfpathlineto{\pgfqpoint{2.597002in}{0.673850in}}%
\pgfpathlineto{\pgfqpoint{2.622285in}{0.671940in}}%
\pgfpathlineto{\pgfqpoint{2.647569in}{0.670203in}}%
\pgfpathlineto{\pgfqpoint{2.672852in}{0.668620in}}%
\pgfpathlineto{\pgfqpoint{2.698135in}{0.667173in}}%
\pgfpathlineto{\pgfqpoint{2.723418in}{0.665849in}}%
\pgfpathlineto{\pgfqpoint{2.748701in}{0.664633in}}%
\pgfpathlineto{\pgfqpoint{2.773984in}{0.663516in}}%
\pgfpathlineto{\pgfqpoint{2.799268in}{0.662488in}}%
\pgfpathlineto{\pgfqpoint{2.811743in}{0.662015in}}%
\pgfpathlineto{\pgfqpoint{2.824551in}{0.657700in}}%
\pgfpathlineto{\pgfqpoint{2.849834in}{0.651007in}}%
\pgfpathlineto{\pgfqpoint{2.860434in}{0.648767in}}%
\pgfpathlineto{\pgfqpoint{2.875117in}{0.648332in}}%
\pgfpathclose%
\pgfpathmoveto{\pgfqpoint{2.521153in}{1.229832in}}%
\pgfpathlineto{\pgfqpoint{2.546436in}{1.229062in}}%
\pgfpathlineto{\pgfqpoint{2.571719in}{1.231586in}}%
\pgfpathlineto{\pgfqpoint{2.572148in}{1.231679in}}%
\pgfpathlineto{\pgfqpoint{2.571719in}{1.232135in}}%
\pgfpathlineto{\pgfqpoint{2.546436in}{1.238086in}}%
\pgfpathlineto{\pgfqpoint{2.521153in}{1.234459in}}%
\pgfpathlineto{\pgfqpoint{2.511867in}{1.231679in}}%
\pgfpathlineto{\pgfqpoint{2.521153in}{1.229832in}}%
\pgfpathclose%
\pgfpathmoveto{\pgfqpoint{2.521153in}{1.252366in}}%
\pgfpathlineto{\pgfqpoint{2.546436in}{1.247815in}}%
\pgfpathlineto{\pgfqpoint{2.571719in}{1.247551in}}%
\pgfpathlineto{\pgfqpoint{2.597002in}{1.250326in}}%
\pgfpathlineto{\pgfqpoint{2.622285in}{1.255735in}}%
\pgfpathlineto{\pgfqpoint{2.629863in}{1.258175in}}%
\pgfpathlineto{\pgfqpoint{2.640732in}{1.271423in}}%
\pgfpathlineto{\pgfqpoint{2.647569in}{1.273343in}}%
\pgfpathlineto{\pgfqpoint{2.672851in}{1.284671in}}%
\pgfpathlineto{\pgfqpoint{2.672852in}{1.284672in}}%
\pgfpathlineto{\pgfqpoint{2.684962in}{1.297919in}}%
\pgfpathlineto{\pgfqpoint{2.698135in}{1.309906in}}%
\pgfpathlineto{\pgfqpoint{2.699195in}{1.311167in}}%
\pgfpathlineto{\pgfqpoint{2.715423in}{1.324415in}}%
\pgfpathlineto{\pgfqpoint{2.723418in}{1.331629in}}%
\pgfpathlineto{\pgfqpoint{2.728670in}{1.337663in}}%
\pgfpathlineto{\pgfqpoint{2.748701in}{1.348290in}}%
\pgfpathlineto{\pgfqpoint{2.753175in}{1.350911in}}%
\pgfpathlineto{\pgfqpoint{2.758405in}{1.364159in}}%
\pgfpathlineto{\pgfqpoint{2.773984in}{1.373442in}}%
\pgfpathlineto{\pgfqpoint{2.779872in}{1.377407in}}%
\pgfpathlineto{\pgfqpoint{2.781277in}{1.390655in}}%
\pgfpathlineto{\pgfqpoint{2.799268in}{1.401380in}}%
\pgfpathlineto{\pgfqpoint{2.802888in}{1.403903in}}%
\pgfpathlineto{\pgfqpoint{2.806729in}{1.417151in}}%
\pgfpathlineto{\pgfqpoint{2.824551in}{1.427713in}}%
\pgfpathlineto{\pgfqpoint{2.828451in}{1.430399in}}%
\pgfpathlineto{\pgfqpoint{2.831834in}{1.443647in}}%
\pgfpathlineto{\pgfqpoint{2.849834in}{1.456132in}}%
\pgfpathlineto{\pgfqpoint{2.850789in}{1.456895in}}%
\pgfpathlineto{\pgfqpoint{2.855554in}{1.470143in}}%
\pgfpathlineto{\pgfqpoint{2.868144in}{1.483391in}}%
\pgfpathlineto{\pgfqpoint{2.872169in}{1.496639in}}%
\pgfpathlineto{\pgfqpoint{2.875117in}{1.499703in}}%
\pgfpathlineto{\pgfqpoint{2.883674in}{1.509887in}}%
\pgfpathlineto{\pgfqpoint{2.891702in}{1.523135in}}%
\pgfpathlineto{\pgfqpoint{2.900400in}{1.533127in}}%
\pgfpathlineto{\pgfqpoint{2.902964in}{1.536383in}}%
\pgfpathlineto{\pgfqpoint{2.913675in}{1.549631in}}%
\pgfpathlineto{\pgfqpoint{2.919272in}{1.562879in}}%
\pgfpathlineto{\pgfqpoint{2.925683in}{1.571120in}}%
\pgfpathlineto{\pgfqpoint{2.929527in}{1.576127in}}%
\pgfpathlineto{\pgfqpoint{2.934255in}{1.589375in}}%
\pgfpathlineto{\pgfqpoint{2.945374in}{1.602623in}}%
\pgfpathlineto{\pgfqpoint{2.949071in}{1.615871in}}%
\pgfpathlineto{\pgfqpoint{2.950967in}{1.617836in}}%
\pgfpathlineto{\pgfqpoint{2.961516in}{1.629119in}}%
\pgfpathlineto{\pgfqpoint{2.963892in}{1.642367in}}%
\pgfpathlineto{\pgfqpoint{2.975053in}{1.655615in}}%
\pgfpathlineto{\pgfqpoint{2.976250in}{1.662078in}}%
\pgfpathlineto{\pgfqpoint{2.977611in}{1.668863in}}%
\pgfpathlineto{\pgfqpoint{2.987293in}{1.682111in}}%
\pgfpathlineto{\pgfqpoint{2.989888in}{1.695359in}}%
\pgfpathlineto{\pgfqpoint{3.000486in}{1.708607in}}%
\pgfpathlineto{\pgfqpoint{3.001533in}{1.711505in}}%
\pgfpathlineto{\pgfqpoint{3.005668in}{1.721855in}}%
\pgfpathlineto{\pgfqpoint{3.014348in}{1.735103in}}%
\pgfpathlineto{\pgfqpoint{3.019064in}{1.748351in}}%
\pgfpathlineto{\pgfqpoint{3.024176in}{1.761599in}}%
\pgfpathlineto{\pgfqpoint{3.026816in}{1.768469in}}%
\pgfpathlineto{\pgfqpoint{3.029590in}{1.774847in}}%
\pgfpathlineto{\pgfqpoint{3.035170in}{1.788095in}}%
\pgfpathlineto{\pgfqpoint{3.041966in}{1.801343in}}%
\pgfpathlineto{\pgfqpoint{3.045864in}{1.814591in}}%
\pgfpathlineto{\pgfqpoint{3.052099in}{1.825025in}}%
\pgfpathlineto{\pgfqpoint{3.054054in}{1.827839in}}%
\pgfpathlineto{\pgfqpoint{3.057336in}{1.841087in}}%
\pgfpathlineto{\pgfqpoint{3.065116in}{1.854335in}}%
\pgfpathlineto{\pgfqpoint{3.067308in}{1.867583in}}%
\pgfpathlineto{\pgfqpoint{3.074575in}{1.880831in}}%
\pgfpathlineto{\pgfqpoint{3.076844in}{1.894079in}}%
\pgfpathlineto{\pgfqpoint{3.077382in}{1.894901in}}%
\pgfpathlineto{\pgfqpoint{3.086354in}{1.907327in}}%
\pgfpathlineto{\pgfqpoint{3.089132in}{1.920575in}}%
\pgfpathlineto{\pgfqpoint{3.095430in}{1.933823in}}%
\pgfpathlineto{\pgfqpoint{3.097613in}{1.947071in}}%
\pgfpathlineto{\pgfqpoint{3.102666in}{1.959575in}}%
\pgfpathlineto{\pgfqpoint{3.103034in}{1.960319in}}%
\pgfpathlineto{\pgfqpoint{3.107127in}{1.973567in}}%
\pgfpathlineto{\pgfqpoint{3.113385in}{1.986815in}}%
\pgfpathlineto{\pgfqpoint{3.118169in}{2.000063in}}%
\pgfpathlineto{\pgfqpoint{3.122463in}{2.013311in}}%
\pgfpathlineto{\pgfqpoint{3.127628in}{2.026559in}}%
\pgfpathlineto{\pgfqpoint{3.127949in}{2.027961in}}%
\pgfpathlineto{\pgfqpoint{3.131429in}{2.039807in}}%
\pgfpathlineto{\pgfqpoint{3.137068in}{2.053055in}}%
\pgfpathlineto{\pgfqpoint{3.139520in}{2.066303in}}%
\pgfpathlineto{\pgfqpoint{3.145764in}{2.079551in}}%
\pgfpathlineto{\pgfqpoint{3.148093in}{2.092799in}}%
\pgfpathlineto{\pgfqpoint{3.153232in}{2.103599in}}%
\pgfpathlineto{\pgfqpoint{3.154792in}{2.106047in}}%
\pgfpathlineto{\pgfqpoint{3.157141in}{2.119295in}}%
\pgfpathlineto{\pgfqpoint{3.163670in}{2.132543in}}%
\pgfpathlineto{\pgfqpoint{3.153232in}{2.132543in}}%
\pgfpathlineto{\pgfqpoint{3.146686in}{2.132543in}}%
\pgfpathlineto{\pgfqpoint{3.140565in}{2.119295in}}%
\pgfpathlineto{\pgfqpoint{3.138719in}{2.106047in}}%
\pgfpathlineto{\pgfqpoint{3.131201in}{2.092799in}}%
\pgfpathlineto{\pgfqpoint{3.128309in}{2.079551in}}%
\pgfpathlineto{\pgfqpoint{3.127949in}{2.078881in}}%
\pgfpathlineto{\pgfqpoint{3.122131in}{2.066303in}}%
\pgfpathlineto{\pgfqpoint{3.119551in}{2.053055in}}%
\pgfpathlineto{\pgfqpoint{3.114473in}{2.039807in}}%
\pgfpathlineto{\pgfqpoint{3.110573in}{2.026559in}}%
\pgfpathlineto{\pgfqpoint{3.104788in}{2.013311in}}%
\pgfpathlineto{\pgfqpoint{3.102666in}{2.008224in}}%
\pgfpathlineto{\pgfqpoint{3.099474in}{2.000063in}}%
\pgfpathlineto{\pgfqpoint{3.095161in}{1.986815in}}%
\pgfpathlineto{\pgfqpoint{3.088618in}{1.973567in}}%
\pgfpathlineto{\pgfqpoint{3.085410in}{1.960319in}}%
\pgfpathlineto{\pgfqpoint{3.078508in}{1.947071in}}%
\pgfpathlineto{\pgfqpoint{3.077382in}{1.939949in}}%
\pgfpathlineto{\pgfqpoint{3.076512in}{1.933823in}}%
\pgfpathlineto{\pgfqpoint{3.069222in}{1.920575in}}%
\pgfpathlineto{\pgfqpoint{3.066815in}{1.907327in}}%
\pgfpathlineto{\pgfqpoint{3.057300in}{1.894079in}}%
\pgfpathlineto{\pgfqpoint{3.054924in}{1.880831in}}%
\pgfpathlineto{\pgfqpoint{3.052099in}{1.876394in}}%
\pgfpathlineto{\pgfqpoint{3.046460in}{1.867583in}}%
\pgfpathlineto{\pgfqpoint{3.044029in}{1.854335in}}%
\pgfpathlineto{\pgfqpoint{3.035841in}{1.841087in}}%
\pgfpathlineto{\pgfqpoint{3.032401in}{1.827839in}}%
\pgfpathlineto{\pgfqpoint{3.026816in}{1.819069in}}%
\pgfpathlineto{\pgfqpoint{3.023940in}{1.814591in}}%
\pgfpathlineto{\pgfqpoint{3.019112in}{1.801343in}}%
\pgfpathlineto{\pgfqpoint{3.012471in}{1.788095in}}%
\pgfpathlineto{\pgfqpoint{3.006021in}{1.774847in}}%
\pgfpathlineto{\pgfqpoint{3.001533in}{1.762983in}}%
\pgfpathlineto{\pgfqpoint{3.000995in}{1.761599in}}%
\pgfpathlineto{\pgfqpoint{2.994000in}{1.748351in}}%
\pgfpathlineto{\pgfqpoint{2.989778in}{1.735103in}}%
\pgfpathlineto{\pgfqpoint{2.979251in}{1.721855in}}%
\pgfpathlineto{\pgfqpoint{2.976250in}{1.712585in}}%
\pgfpathlineto{\pgfqpoint{2.974895in}{1.708607in}}%
\pgfpathlineto{\pgfqpoint{2.961629in}{1.695359in}}%
\pgfpathlineto{\pgfqpoint{2.959923in}{1.682111in}}%
\pgfpathlineto{\pgfqpoint{2.950967in}{1.671493in}}%
\pgfpathlineto{\pgfqpoint{2.948315in}{1.668863in}}%
\pgfpathlineto{\pgfqpoint{2.946342in}{1.655615in}}%
\pgfpathlineto{\pgfqpoint{2.932706in}{1.642367in}}%
\pgfpathlineto{\pgfqpoint{2.930731in}{1.629119in}}%
\pgfpathlineto{\pgfqpoint{2.925683in}{1.623971in}}%
\pgfpathlineto{\pgfqpoint{2.916581in}{1.615871in}}%
\pgfpathlineto{\pgfqpoint{2.911958in}{1.602623in}}%
\pgfpathlineto{\pgfqpoint{2.900400in}{1.590145in}}%
\pgfpathlineto{\pgfqpoint{2.899540in}{1.589375in}}%
\pgfpathlineto{\pgfqpoint{2.892580in}{1.576127in}}%
\pgfpathlineto{\pgfqpoint{2.881966in}{1.562879in}}%
\pgfpathlineto{\pgfqpoint{2.875117in}{1.551479in}}%
\pgfpathlineto{\pgfqpoint{2.873616in}{1.549631in}}%
\pgfpathlineto{\pgfqpoint{2.862301in}{1.536383in}}%
\pgfpathlineto{\pgfqpoint{2.849834in}{1.526004in}}%
\pgfpathlineto{\pgfqpoint{2.844820in}{1.523135in}}%
\pgfpathlineto{\pgfqpoint{2.838597in}{1.509887in}}%
\pgfpathlineto{\pgfqpoint{2.824551in}{1.498996in}}%
\pgfpathlineto{\pgfqpoint{2.819948in}{1.496639in}}%
\pgfpathlineto{\pgfqpoint{2.820930in}{1.483391in}}%
\pgfpathlineto{\pgfqpoint{2.799268in}{1.470775in}}%
\pgfpathlineto{\pgfqpoint{2.797442in}{1.470143in}}%
\pgfpathlineto{\pgfqpoint{2.797966in}{1.456895in}}%
\pgfpathlineto{\pgfqpoint{2.773984in}{1.447722in}}%
\pgfpathlineto{\pgfqpoint{2.755930in}{1.443647in}}%
\pgfpathlineto{\pgfqpoint{2.762086in}{1.430399in}}%
\pgfpathlineto{\pgfqpoint{2.748701in}{1.426136in}}%
\pgfpathlineto{\pgfqpoint{2.723418in}{1.422488in}}%
\pgfpathlineto{\pgfqpoint{2.698135in}{1.425520in}}%
\pgfpathlineto{\pgfqpoint{2.689733in}{1.430399in}}%
\pgfpathlineto{\pgfqpoint{2.698135in}{1.437677in}}%
\pgfpathlineto{\pgfqpoint{2.710561in}{1.443647in}}%
\pgfpathlineto{\pgfqpoint{2.698135in}{1.448182in}}%
\pgfpathlineto{\pgfqpoint{2.689169in}{1.456895in}}%
\pgfpathlineto{\pgfqpoint{2.698135in}{1.463973in}}%
\pgfpathlineto{\pgfqpoint{2.708644in}{1.470143in}}%
\pgfpathlineto{\pgfqpoint{2.698135in}{1.481479in}}%
\pgfpathlineto{\pgfqpoint{2.697109in}{1.483391in}}%
\pgfpathlineto{\pgfqpoint{2.698135in}{1.484140in}}%
\pgfpathlineto{\pgfqpoint{2.718740in}{1.496639in}}%
\pgfpathlineto{\pgfqpoint{2.712493in}{1.509887in}}%
\pgfpathlineto{\pgfqpoint{2.723418in}{1.518801in}}%
\pgfpathlineto{\pgfqpoint{2.729813in}{1.523135in}}%
\pgfpathlineto{\pgfqpoint{2.723418in}{1.532197in}}%
\pgfpathlineto{\pgfqpoint{2.721429in}{1.536383in}}%
\pgfpathlineto{\pgfqpoint{2.723418in}{1.538440in}}%
\pgfpathlineto{\pgfqpoint{2.735289in}{1.549631in}}%
\pgfpathlineto{\pgfqpoint{2.731737in}{1.562879in}}%
\pgfpathlineto{\pgfqpoint{2.746043in}{1.576127in}}%
\pgfpathlineto{\pgfqpoint{2.746159in}{1.589375in}}%
\pgfpathlineto{\pgfqpoint{2.748701in}{1.592921in}}%
\pgfpathlineto{\pgfqpoint{2.755815in}{1.602623in}}%
\pgfpathlineto{\pgfqpoint{2.755856in}{1.615871in}}%
\pgfpathlineto{\pgfqpoint{2.762872in}{1.629119in}}%
\pgfpathlineto{\pgfqpoint{2.766562in}{1.642367in}}%
\pgfpathlineto{\pgfqpoint{2.772958in}{1.655615in}}%
\pgfpathlineto{\pgfqpoint{2.773984in}{1.657768in}}%
\pgfpathlineto{\pgfqpoint{2.779708in}{1.668863in}}%
\pgfpathlineto{\pgfqpoint{2.784182in}{1.682111in}}%
\pgfpathlineto{\pgfqpoint{2.792875in}{1.695359in}}%
\pgfpathlineto{\pgfqpoint{2.795034in}{1.708607in}}%
\pgfpathlineto{\pgfqpoint{2.799268in}{1.715882in}}%
\pgfpathlineto{\pgfqpoint{2.802962in}{1.721855in}}%
\pgfpathlineto{\pgfqpoint{2.802932in}{1.735103in}}%
\pgfpathlineto{\pgfqpoint{2.811587in}{1.748351in}}%
\pgfpathlineto{\pgfqpoint{2.812400in}{1.761599in}}%
\pgfpathlineto{\pgfqpoint{2.823211in}{1.774847in}}%
\pgfpathlineto{\pgfqpoint{2.823257in}{1.788095in}}%
\pgfpathlineto{\pgfqpoint{2.824551in}{1.790030in}}%
\pgfpathlineto{\pgfqpoint{2.832017in}{1.801343in}}%
\pgfpathlineto{\pgfqpoint{2.831890in}{1.814591in}}%
\pgfpathlineto{\pgfqpoint{2.841012in}{1.827839in}}%
\pgfpathlineto{\pgfqpoint{2.842388in}{1.841087in}}%
\pgfpathlineto{\pgfqpoint{2.849834in}{1.852188in}}%
\pgfpathlineto{\pgfqpoint{2.851204in}{1.854335in}}%
\pgfpathlineto{\pgfqpoint{2.853082in}{1.867583in}}%
\pgfpathlineto{\pgfqpoint{2.859763in}{1.880831in}}%
\pgfpathlineto{\pgfqpoint{2.862548in}{1.894079in}}%
\pgfpathlineto{\pgfqpoint{2.867321in}{1.907327in}}%
\pgfpathlineto{\pgfqpoint{2.870210in}{1.920575in}}%
\pgfpathlineto{\pgfqpoint{2.874557in}{1.933823in}}%
\pgfpathlineto{\pgfqpoint{2.875117in}{1.935176in}}%
\pgfpathlineto{\pgfqpoint{2.879744in}{1.947071in}}%
\pgfpathlineto{\pgfqpoint{2.883624in}{1.960319in}}%
\pgfpathlineto{\pgfqpoint{2.889985in}{1.973567in}}%
\pgfpathlineto{\pgfqpoint{2.891812in}{1.986815in}}%
\pgfpathlineto{\pgfqpoint{2.897849in}{2.000063in}}%
\pgfpathlineto{\pgfqpoint{2.899319in}{2.013311in}}%
\pgfpathlineto{\pgfqpoint{2.900400in}{2.015195in}}%
\pgfpathlineto{\pgfqpoint{2.906303in}{2.026559in}}%
\pgfpathlineto{\pgfqpoint{2.907381in}{2.039807in}}%
\pgfpathlineto{\pgfqpoint{2.914153in}{2.053055in}}%
\pgfpathlineto{\pgfqpoint{2.914968in}{2.066303in}}%
\pgfpathlineto{\pgfqpoint{2.921558in}{2.079551in}}%
\pgfpathlineto{\pgfqpoint{2.922743in}{2.092799in}}%
\pgfpathlineto{\pgfqpoint{2.925683in}{2.098967in}}%
\pgfpathlineto{\pgfqpoint{2.928580in}{2.106047in}}%
\pgfpathlineto{\pgfqpoint{2.929890in}{2.119295in}}%
\pgfpathlineto{\pgfqpoint{2.935477in}{2.132543in}}%
\pgfpathlineto{\pgfqpoint{2.925683in}{2.132543in}}%
\pgfpathlineto{\pgfqpoint{2.918983in}{2.132543in}}%
\pgfpathlineto{\pgfqpoint{2.912488in}{2.119295in}}%
\pgfpathlineto{\pgfqpoint{2.911068in}{2.106047in}}%
\pgfpathlineto{\pgfqpoint{2.905603in}{2.092799in}}%
\pgfpathlineto{\pgfqpoint{2.904376in}{2.079551in}}%
\pgfpathlineto{\pgfqpoint{2.900400in}{2.070528in}}%
\pgfpathlineto{\pgfqpoint{2.898008in}{2.066303in}}%
\pgfpathlineto{\pgfqpoint{2.896500in}{2.053055in}}%
\pgfpathlineto{\pgfqpoint{2.888966in}{2.039807in}}%
\pgfpathlineto{\pgfqpoint{2.887004in}{2.026559in}}%
\pgfpathlineto{\pgfqpoint{2.880517in}{2.013311in}}%
\pgfpathlineto{\pgfqpoint{2.878504in}{2.000063in}}%
\pgfpathlineto{\pgfqpoint{2.875117in}{1.991074in}}%
\pgfpathlineto{\pgfqpoint{2.873153in}{1.986815in}}%
\pgfpathlineto{\pgfqpoint{2.869983in}{1.973567in}}%
\pgfpathlineto{\pgfqpoint{2.863559in}{1.960319in}}%
\pgfpathlineto{\pgfqpoint{2.858602in}{1.947071in}}%
\pgfpathlineto{\pgfqpoint{2.854136in}{1.933823in}}%
\pgfpathlineto{\pgfqpoint{2.849834in}{1.921970in}}%
\pgfpathlineto{\pgfqpoint{2.849178in}{1.920575in}}%
\pgfpathlineto{\pgfqpoint{2.846255in}{1.907327in}}%
\pgfpathlineto{\pgfqpoint{2.840179in}{1.894079in}}%
\pgfpathlineto{\pgfqpoint{2.837376in}{1.880831in}}%
\pgfpathlineto{\pgfqpoint{2.830010in}{1.867583in}}%
\pgfpathlineto{\pgfqpoint{2.827945in}{1.854335in}}%
\pgfpathlineto{\pgfqpoint{2.824551in}{1.849245in}}%
\pgfpathlineto{\pgfqpoint{2.818222in}{1.841087in}}%
\pgfpathlineto{\pgfqpoint{2.816145in}{1.827839in}}%
\pgfpathlineto{\pgfqpoint{2.807128in}{1.814591in}}%
\pgfpathlineto{\pgfqpoint{2.806311in}{1.801343in}}%
\pgfpathlineto{\pgfqpoint{2.799268in}{1.789939in}}%
\pgfpathlineto{\pgfqpoint{2.797989in}{1.788095in}}%
\pgfpathlineto{\pgfqpoint{2.796340in}{1.774847in}}%
\pgfpathlineto{\pgfqpoint{2.785558in}{1.761599in}}%
\pgfpathlineto{\pgfqpoint{2.782871in}{1.748351in}}%
\pgfpathlineto{\pgfqpoint{2.775312in}{1.735103in}}%
\pgfpathlineto{\pgfqpoint{2.773984in}{1.727692in}}%
\pgfpathlineto{\pgfqpoint{2.772715in}{1.721855in}}%
\pgfpathlineto{\pgfqpoint{2.765348in}{1.708607in}}%
\pgfpathlineto{\pgfqpoint{2.760630in}{1.695359in}}%
\pgfpathlineto{\pgfqpoint{2.752962in}{1.682111in}}%
\pgfpathlineto{\pgfqpoint{2.748701in}{1.673799in}}%
\pgfpathlineto{\pgfqpoint{2.745774in}{1.668863in}}%
\pgfpathlineto{\pgfqpoint{2.739850in}{1.655615in}}%
\pgfpathlineto{\pgfqpoint{2.730784in}{1.642367in}}%
\pgfpathlineto{\pgfqpoint{2.727499in}{1.629119in}}%
\pgfpathlineto{\pgfqpoint{2.723418in}{1.622755in}}%
\pgfpathlineto{\pgfqpoint{2.718641in}{1.615871in}}%
\pgfpathlineto{\pgfqpoint{2.717535in}{1.602623in}}%
\pgfpathlineto{\pgfqpoint{2.706232in}{1.589375in}}%
\pgfpathlineto{\pgfqpoint{2.704269in}{1.576127in}}%
\pgfpathlineto{\pgfqpoint{2.698135in}{1.570790in}}%
\pgfpathlineto{\pgfqpoint{2.688743in}{1.562879in}}%
\pgfpathlineto{\pgfqpoint{2.688755in}{1.549631in}}%
\pgfpathlineto{\pgfqpoint{2.675868in}{1.536383in}}%
\pgfpathlineto{\pgfqpoint{2.676684in}{1.523135in}}%
\pgfpathlineto{\pgfqpoint{2.672852in}{1.519855in}}%
\pgfpathlineto{\pgfqpoint{2.661629in}{1.509887in}}%
\pgfpathlineto{\pgfqpoint{2.660272in}{1.496639in}}%
\pgfpathlineto{\pgfqpoint{2.647569in}{1.486424in}}%
\pgfpathlineto{\pgfqpoint{2.643944in}{1.483391in}}%
\pgfpathlineto{\pgfqpoint{2.642428in}{1.470143in}}%
\pgfpathlineto{\pgfqpoint{2.629175in}{1.456895in}}%
\pgfpathlineto{\pgfqpoint{2.626843in}{1.443647in}}%
\pgfpathlineto{\pgfqpoint{2.622285in}{1.437831in}}%
\pgfpathlineto{\pgfqpoint{2.616925in}{1.430399in}}%
\pgfpathlineto{\pgfqpoint{2.613937in}{1.417151in}}%
\pgfpathlineto{\pgfqpoint{2.604573in}{1.403903in}}%
\pgfpathlineto{\pgfqpoint{2.597002in}{1.393367in}}%
\pgfpathlineto{\pgfqpoint{2.594878in}{1.390655in}}%
\pgfpathlineto{\pgfqpoint{2.584811in}{1.377407in}}%
\pgfpathlineto{\pgfqpoint{2.571719in}{1.365979in}}%
\pgfpathlineto{\pgfqpoint{2.569508in}{1.364159in}}%
\pgfpathlineto{\pgfqpoint{2.569822in}{1.350911in}}%
\pgfpathlineto{\pgfqpoint{2.556961in}{1.337663in}}%
\pgfpathlineto{\pgfqpoint{2.560334in}{1.324415in}}%
\pgfpathlineto{\pgfqpoint{2.546436in}{1.316096in}}%
\pgfpathlineto{\pgfqpoint{2.538279in}{1.311167in}}%
\pgfpathlineto{\pgfqpoint{2.543046in}{1.297919in}}%
\pgfpathlineto{\pgfqpoint{2.521153in}{1.287132in}}%
\pgfpathlineto{\pgfqpoint{2.516228in}{1.284671in}}%
\pgfpathlineto{\pgfqpoint{2.521153in}{1.280577in}}%
\pgfpathlineto{\pgfqpoint{2.538469in}{1.271423in}}%
\pgfpathlineto{\pgfqpoint{2.521153in}{1.264341in}}%
\pgfpathlineto{\pgfqpoint{2.507224in}{1.258175in}}%
\pgfpathlineto{\pgfqpoint{2.521153in}{1.252366in}}%
\pgfpathclose%
\pgfpathmoveto{\pgfqpoint{4.038143in}{1.270202in}}%
\pgfpathlineto{\pgfqpoint{4.063426in}{1.263068in}}%
\pgfpathlineto{\pgfqpoint{4.088709in}{1.259338in}}%
\pgfpathlineto{\pgfqpoint{4.113993in}{1.260887in}}%
\pgfpathlineto{\pgfqpoint{4.136066in}{1.271423in}}%
\pgfpathlineto{\pgfqpoint{4.113993in}{1.282818in}}%
\pgfpathlineto{\pgfqpoint{4.109743in}{1.284671in}}%
\pgfpathlineto{\pgfqpoint{4.113993in}{1.288508in}}%
\pgfpathlineto{\pgfqpoint{4.119181in}{1.297919in}}%
\pgfpathlineto{\pgfqpoint{4.113993in}{1.300573in}}%
\pgfpathlineto{\pgfqpoint{4.090272in}{1.311167in}}%
\pgfpathlineto{\pgfqpoint{4.106474in}{1.324415in}}%
\pgfpathlineto{\pgfqpoint{4.091317in}{1.337663in}}%
\pgfpathlineto{\pgfqpoint{4.096425in}{1.350911in}}%
\pgfpathlineto{\pgfqpoint{4.088709in}{1.357564in}}%
\pgfpathlineto{\pgfqpoint{4.080308in}{1.364159in}}%
\pgfpathlineto{\pgfqpoint{4.079894in}{1.377407in}}%
\pgfpathlineto{\pgfqpoint{4.064985in}{1.390655in}}%
\pgfpathlineto{\pgfqpoint{4.063426in}{1.397101in}}%
\pgfpathlineto{\pgfqpoint{4.061852in}{1.403903in}}%
\pgfpathlineto{\pgfqpoint{4.052078in}{1.417151in}}%
\pgfpathlineto{\pgfqpoint{4.045506in}{1.430399in}}%
\pgfpathlineto{\pgfqpoint{4.040221in}{1.443647in}}%
\pgfpathlineto{\pgfqpoint{4.038143in}{1.446787in}}%
\pgfpathlineto{\pgfqpoint{4.031612in}{1.456895in}}%
\pgfpathlineto{\pgfqpoint{4.025042in}{1.470143in}}%
\pgfpathlineto{\pgfqpoint{4.012860in}{1.483103in}}%
\pgfpathlineto{\pgfqpoint{4.012593in}{1.483391in}}%
\pgfpathlineto{\pgfqpoint{4.009291in}{1.496639in}}%
\pgfpathlineto{\pgfqpoint{3.997654in}{1.509887in}}%
\pgfpathlineto{\pgfqpoint{3.998116in}{1.523135in}}%
\pgfpathlineto{\pgfqpoint{3.987577in}{1.536218in}}%
\pgfpathlineto{\pgfqpoint{3.987443in}{1.536383in}}%
\pgfpathlineto{\pgfqpoint{3.985066in}{1.549631in}}%
\pgfpathlineto{\pgfqpoint{3.971785in}{1.562879in}}%
\pgfpathlineto{\pgfqpoint{3.969882in}{1.576127in}}%
\pgfpathlineto{\pgfqpoint{3.962293in}{1.583987in}}%
\pgfpathlineto{\pgfqpoint{3.957120in}{1.589375in}}%
\pgfpathlineto{\pgfqpoint{3.955222in}{1.602623in}}%
\pgfpathlineto{\pgfqpoint{3.945418in}{1.615871in}}%
\pgfpathlineto{\pgfqpoint{3.942993in}{1.629119in}}%
\pgfpathlineto{\pgfqpoint{3.937010in}{1.638010in}}%
\pgfpathlineto{\pgfqpoint{3.934197in}{1.642367in}}%
\pgfpathlineto{\pgfqpoint{3.930434in}{1.655615in}}%
\pgfpathlineto{\pgfqpoint{3.921582in}{1.668863in}}%
\pgfpathlineto{\pgfqpoint{3.915824in}{1.682111in}}%
\pgfpathlineto{\pgfqpoint{3.911727in}{1.691756in}}%
\pgfpathlineto{\pgfqpoint{3.910288in}{1.695359in}}%
\pgfpathlineto{\pgfqpoint{3.904942in}{1.708607in}}%
\pgfpathlineto{\pgfqpoint{3.900899in}{1.721855in}}%
\pgfpathlineto{\pgfqpoint{3.893801in}{1.735103in}}%
\pgfpathlineto{\pgfqpoint{3.888645in}{1.748351in}}%
\pgfpathlineto{\pgfqpoint{3.886444in}{1.751413in}}%
\pgfpathlineto{\pgfqpoint{3.880016in}{1.761599in}}%
\pgfpathlineto{\pgfqpoint{3.877468in}{1.774847in}}%
\pgfpathlineto{\pgfqpoint{3.869117in}{1.788095in}}%
\pgfpathlineto{\pgfqpoint{3.867354in}{1.801343in}}%
\pgfpathlineto{\pgfqpoint{3.861161in}{1.811393in}}%
\pgfpathlineto{\pgfqpoint{3.859392in}{1.814591in}}%
\pgfpathlineto{\pgfqpoint{3.857813in}{1.827839in}}%
\pgfpathlineto{\pgfqpoint{3.849824in}{1.841087in}}%
\pgfpathlineto{\pgfqpoint{3.847557in}{1.854335in}}%
\pgfpathlineto{\pgfqpoint{3.838509in}{1.867583in}}%
\pgfpathlineto{\pgfqpoint{3.835878in}{1.880708in}}%
\pgfpathlineto{\pgfqpoint{3.835858in}{1.880831in}}%
\pgfpathlineto{\pgfqpoint{3.830019in}{1.894079in}}%
\pgfpathlineto{\pgfqpoint{3.827666in}{1.907327in}}%
\pgfpathlineto{\pgfqpoint{3.821418in}{1.920575in}}%
\pgfpathlineto{\pgfqpoint{3.817591in}{1.933823in}}%
\pgfpathlineto{\pgfqpoint{3.810666in}{1.947071in}}%
\pgfpathlineto{\pgfqpoint{3.810594in}{1.947264in}}%
\pgfpathlineto{\pgfqpoint{3.806850in}{1.960319in}}%
\pgfpathlineto{\pgfqpoint{3.803305in}{1.973567in}}%
\pgfpathlineto{\pgfqpoint{3.799018in}{1.986815in}}%
\pgfpathlineto{\pgfqpoint{3.795762in}{2.000063in}}%
\pgfpathlineto{\pgfqpoint{3.790623in}{2.013311in}}%
\pgfpathlineto{\pgfqpoint{3.787048in}{2.026559in}}%
\pgfpathlineto{\pgfqpoint{3.785311in}{2.029658in}}%
\pgfpathlineto{\pgfqpoint{3.780962in}{2.039807in}}%
\pgfpathlineto{\pgfqpoint{3.778483in}{2.053055in}}%
\pgfpathlineto{\pgfqpoint{3.772178in}{2.066303in}}%
\pgfpathlineto{\pgfqpoint{3.770628in}{2.079551in}}%
\pgfpathlineto{\pgfqpoint{3.764975in}{2.092799in}}%
\pgfpathlineto{\pgfqpoint{3.762986in}{2.106047in}}%
\pgfpathlineto{\pgfqpoint{3.760028in}{2.111227in}}%
\pgfpathlineto{\pgfqpoint{3.756556in}{2.119295in}}%
\pgfpathlineto{\pgfqpoint{3.754889in}{2.132543in}}%
\pgfpathlineto{\pgfqpoint{3.738176in}{2.132543in}}%
\pgfpathlineto{\pgfqpoint{3.739781in}{2.119295in}}%
\pgfpathlineto{\pgfqpoint{3.746095in}{2.106047in}}%
\pgfpathlineto{\pgfqpoint{3.747462in}{2.092799in}}%
\pgfpathlineto{\pgfqpoint{3.752570in}{2.079551in}}%
\pgfpathlineto{\pgfqpoint{3.753814in}{2.066303in}}%
\pgfpathlineto{\pgfqpoint{3.759919in}{2.053055in}}%
\pgfpathlineto{\pgfqpoint{3.760028in}{2.052453in}}%
\pgfpathlineto{\pgfqpoint{3.762629in}{2.039807in}}%
\pgfpathlineto{\pgfqpoint{3.768538in}{2.026559in}}%
\pgfpathlineto{\pgfqpoint{3.771491in}{2.013311in}}%
\pgfpathlineto{\pgfqpoint{3.775859in}{2.000063in}}%
\pgfpathlineto{\pgfqpoint{3.778939in}{1.986815in}}%
\pgfpathlineto{\pgfqpoint{3.782750in}{1.973567in}}%
\pgfpathlineto{\pgfqpoint{3.785311in}{1.964564in}}%
\pgfpathlineto{\pgfqpoint{3.786619in}{1.960319in}}%
\pgfpathlineto{\pgfqpoint{3.789979in}{1.947071in}}%
\pgfpathlineto{\pgfqpoint{3.796418in}{1.933823in}}%
\pgfpathlineto{\pgfqpoint{3.799321in}{1.920575in}}%
\pgfpathlineto{\pgfqpoint{3.805543in}{1.907327in}}%
\pgfpathlineto{\pgfqpoint{3.807160in}{1.894079in}}%
\pgfpathlineto{\pgfqpoint{3.810594in}{1.887073in}}%
\pgfpathlineto{\pgfqpoint{3.813738in}{1.880831in}}%
\pgfpathlineto{\pgfqpoint{3.815073in}{1.867583in}}%
\pgfpathlineto{\pgfqpoint{3.823761in}{1.854335in}}%
\pgfpathlineto{\pgfqpoint{3.825006in}{1.841087in}}%
\pgfpathlineto{\pgfqpoint{3.832861in}{1.827839in}}%
\pgfpathlineto{\pgfqpoint{3.833648in}{1.814591in}}%
\pgfpathlineto{\pgfqpoint{3.835878in}{1.810742in}}%
\pgfpathlineto{\pgfqpoint{3.841310in}{1.801343in}}%
\pgfpathlineto{\pgfqpoint{3.842121in}{1.788095in}}%
\pgfpathlineto{\pgfqpoint{3.849958in}{1.774847in}}%
\pgfpathlineto{\pgfqpoint{3.851976in}{1.761599in}}%
\pgfpathlineto{\pgfqpoint{3.859793in}{1.748351in}}%
\pgfpathlineto{\pgfqpoint{3.861161in}{1.744579in}}%
\pgfpathlineto{\pgfqpoint{3.864513in}{1.735103in}}%
\pgfpathlineto{\pgfqpoint{3.870434in}{1.721855in}}%
\pgfpathlineto{\pgfqpoint{3.873931in}{1.708607in}}%
\pgfpathlineto{\pgfqpoint{3.877594in}{1.695359in}}%
\pgfpathlineto{\pgfqpoint{3.883308in}{1.682111in}}%
\pgfpathlineto{\pgfqpoint{3.886444in}{1.669963in}}%
\pgfpathlineto{\pgfqpoint{3.886731in}{1.668863in}}%
\pgfpathlineto{\pgfqpoint{3.895743in}{1.655615in}}%
\pgfpathlineto{\pgfqpoint{3.896814in}{1.642367in}}%
\pgfpathlineto{\pgfqpoint{3.905649in}{1.629119in}}%
\pgfpathlineto{\pgfqpoint{3.904962in}{1.615871in}}%
\pgfpathlineto{\pgfqpoint{3.911727in}{1.608300in}}%
\pgfpathlineto{\pgfqpoint{3.916175in}{1.602623in}}%
\pgfpathlineto{\pgfqpoint{3.913973in}{1.589375in}}%
\pgfpathlineto{\pgfqpoint{3.926857in}{1.576127in}}%
\pgfpathlineto{\pgfqpoint{3.924718in}{1.562879in}}%
\pgfpathlineto{\pgfqpoint{3.937010in}{1.551453in}}%
\pgfpathlineto{\pgfqpoint{3.938769in}{1.549631in}}%
\pgfpathlineto{\pgfqpoint{3.937010in}{1.536954in}}%
\pgfpathlineto{\pgfqpoint{3.936892in}{1.536383in}}%
\pgfpathlineto{\pgfqpoint{3.937010in}{1.536236in}}%
\pgfpathlineto{\pgfqpoint{3.945776in}{1.523135in}}%
\pgfpathlineto{\pgfqpoint{3.938505in}{1.509887in}}%
\pgfpathlineto{\pgfqpoint{3.946615in}{1.496639in}}%
\pgfpathlineto{\pgfqpoint{3.945922in}{1.483391in}}%
\pgfpathlineto{\pgfqpoint{3.952280in}{1.470143in}}%
\pgfpathlineto{\pgfqpoint{3.955238in}{1.456895in}}%
\pgfpathlineto{\pgfqpoint{3.937010in}{1.443871in}}%
\pgfpathlineto{\pgfqpoint{3.911727in}{1.449902in}}%
\pgfpathlineto{\pgfqpoint{3.897915in}{1.456895in}}%
\pgfpathlineto{\pgfqpoint{3.886444in}{1.468634in}}%
\pgfpathlineto{\pgfqpoint{3.885585in}{1.470143in}}%
\pgfpathlineto{\pgfqpoint{3.869743in}{1.483391in}}%
\pgfpathlineto{\pgfqpoint{3.861161in}{1.493325in}}%
\pgfpathlineto{\pgfqpoint{3.859392in}{1.496639in}}%
\pgfpathlineto{\pgfqpoint{3.846609in}{1.509887in}}%
\pgfpathlineto{\pgfqpoint{3.835878in}{1.516513in}}%
\pgfpathlineto{\pgfqpoint{3.828738in}{1.523135in}}%
\pgfpathlineto{\pgfqpoint{3.817730in}{1.536383in}}%
\pgfpathlineto{\pgfqpoint{3.810594in}{1.541802in}}%
\pgfpathlineto{\pgfqpoint{3.802974in}{1.549631in}}%
\pgfpathlineto{\pgfqpoint{3.797359in}{1.562879in}}%
\pgfpathlineto{\pgfqpoint{3.785311in}{1.574495in}}%
\pgfpathlineto{\pgfqpoint{3.783995in}{1.576127in}}%
\pgfpathlineto{\pgfqpoint{3.779251in}{1.589375in}}%
\pgfpathlineto{\pgfqpoint{3.764183in}{1.602623in}}%
\pgfpathlineto{\pgfqpoint{3.760159in}{1.615871in}}%
\pgfpathlineto{\pgfqpoint{3.760028in}{1.615969in}}%
\pgfpathlineto{\pgfqpoint{3.746330in}{1.629119in}}%
\pgfpathlineto{\pgfqpoint{3.743382in}{1.642367in}}%
\pgfpathlineto{\pgfqpoint{3.734745in}{1.650622in}}%
\pgfpathlineto{\pgfqpoint{3.730383in}{1.655615in}}%
\pgfpathlineto{\pgfqpoint{3.727778in}{1.668863in}}%
\pgfpathlineto{\pgfqpoint{3.715807in}{1.682111in}}%
\pgfpathlineto{\pgfqpoint{3.713362in}{1.695359in}}%
\pgfpathlineto{\pgfqpoint{3.709462in}{1.698933in}}%
\pgfpathlineto{\pgfqpoint{3.700405in}{1.708607in}}%
\pgfpathlineto{\pgfqpoint{3.696084in}{1.721855in}}%
\pgfpathlineto{\pgfqpoint{3.685342in}{1.735103in}}%
\pgfpathlineto{\pgfqpoint{3.684179in}{1.739146in}}%
\pgfpathlineto{\pgfqpoint{3.681720in}{1.748351in}}%
\pgfpathlineto{\pgfqpoint{3.673851in}{1.761599in}}%
\pgfpathlineto{\pgfqpoint{3.670192in}{1.774847in}}%
\pgfpathlineto{\pgfqpoint{3.661678in}{1.788095in}}%
\pgfpathlineto{\pgfqpoint{3.658895in}{1.794372in}}%
\pgfpathlineto{\pgfqpoint{3.656018in}{1.801343in}}%
\pgfpathlineto{\pgfqpoint{3.649471in}{1.814591in}}%
\pgfpathlineto{\pgfqpoint{3.643512in}{1.827839in}}%
\pgfpathlineto{\pgfqpoint{3.638386in}{1.841087in}}%
\pgfpathlineto{\pgfqpoint{3.633612in}{1.851025in}}%
\pgfpathlineto{\pgfqpoint{3.632070in}{1.854335in}}%
\pgfpathlineto{\pgfqpoint{3.627235in}{1.867583in}}%
\pgfpathlineto{\pgfqpoint{3.620335in}{1.880831in}}%
\pgfpathlineto{\pgfqpoint{3.616693in}{1.894079in}}%
\pgfpathlineto{\pgfqpoint{3.608351in}{1.907327in}}%
\pgfpathlineto{\pgfqpoint{3.608329in}{1.907415in}}%
\pgfpathlineto{\pgfqpoint{3.604865in}{1.920575in}}%
\pgfpathlineto{\pgfqpoint{3.597533in}{1.933823in}}%
\pgfpathlineto{\pgfqpoint{3.595718in}{1.947071in}}%
\pgfpathlineto{\pgfqpoint{3.588884in}{1.960319in}}%
\pgfpathlineto{\pgfqpoint{3.586725in}{1.973567in}}%
\pgfpathlineto{\pgfqpoint{3.583046in}{1.979301in}}%
\pgfpathlineto{\pgfqpoint{3.578083in}{1.986815in}}%
\pgfpathlineto{\pgfqpoint{3.575474in}{2.000063in}}%
\pgfpathlineto{\pgfqpoint{3.568130in}{2.013311in}}%
\pgfpathlineto{\pgfqpoint{3.565856in}{2.026559in}}%
\pgfpathlineto{\pgfqpoint{3.559368in}{2.039807in}}%
\pgfpathlineto{\pgfqpoint{3.557763in}{2.048326in}}%
\pgfpathlineto{\pgfqpoint{3.556750in}{2.053055in}}%
\pgfpathlineto{\pgfqpoint{3.549707in}{2.066303in}}%
\pgfpathlineto{\pgfqpoint{3.546783in}{2.079551in}}%
\pgfpathlineto{\pgfqpoint{3.540938in}{2.092799in}}%
\pgfpathlineto{\pgfqpoint{3.537218in}{2.106047in}}%
\pgfpathlineto{\pgfqpoint{3.532516in}{2.119295in}}%
\pgfpathlineto{\pgfqpoint{3.532480in}{2.119436in}}%
\pgfpathlineto{\pgfqpoint{3.528569in}{2.132543in}}%
\pgfpathlineto{\pgfqpoint{3.511745in}{2.132543in}}%
\pgfpathlineto{\pgfqpoint{3.514749in}{2.119295in}}%
\pgfpathlineto{\pgfqpoint{3.519471in}{2.106047in}}%
\pgfpathlineto{\pgfqpoint{3.523164in}{2.092799in}}%
\pgfpathlineto{\pgfqpoint{3.529994in}{2.079551in}}%
\pgfpathlineto{\pgfqpoint{3.532480in}{2.069610in}}%
\pgfpathlineto{\pgfqpoint{3.533089in}{2.066303in}}%
\pgfpathlineto{\pgfqpoint{3.538947in}{2.053055in}}%
\pgfpathlineto{\pgfqpoint{3.541195in}{2.039807in}}%
\pgfpathlineto{\pgfqpoint{3.547384in}{2.026559in}}%
\pgfpathlineto{\pgfqpoint{3.549635in}{2.013311in}}%
\pgfpathlineto{\pgfqpoint{3.557763in}{2.000091in}}%
\pgfpathlineto{\pgfqpoint{3.557777in}{2.000063in}}%
\pgfpathlineto{\pgfqpoint{3.560069in}{1.986815in}}%
\pgfpathlineto{\pgfqpoint{3.567451in}{1.973567in}}%
\pgfpathlineto{\pgfqpoint{3.569658in}{1.960319in}}%
\pgfpathlineto{\pgfqpoint{3.576016in}{1.947071in}}%
\pgfpathlineto{\pgfqpoint{3.578136in}{1.933823in}}%
\pgfpathlineto{\pgfqpoint{3.583046in}{1.925149in}}%
\pgfpathlineto{\pgfqpoint{3.585233in}{1.920575in}}%
\pgfpathlineto{\pgfqpoint{3.588357in}{1.907327in}}%
\pgfpathlineto{\pgfqpoint{3.595721in}{1.894079in}}%
\pgfpathlineto{\pgfqpoint{3.599698in}{1.880831in}}%
\pgfpathlineto{\pgfqpoint{3.606207in}{1.867583in}}%
\pgfpathlineto{\pgfqpoint{3.608329in}{1.862400in}}%
\pgfpathlineto{\pgfqpoint{3.611107in}{1.854335in}}%
\pgfpathlineto{\pgfqpoint{3.616174in}{1.841087in}}%
\pgfpathlineto{\pgfqpoint{3.620912in}{1.827839in}}%
\pgfpathlineto{\pgfqpoint{3.626234in}{1.814591in}}%
\pgfpathlineto{\pgfqpoint{3.633352in}{1.801343in}}%
\pgfpathlineto{\pgfqpoint{3.633612in}{1.800679in}}%
\pgfpathlineto{\pgfqpoint{3.637966in}{1.788095in}}%
\pgfpathlineto{\pgfqpoint{3.645714in}{1.774847in}}%
\pgfpathlineto{\pgfqpoint{3.648912in}{1.761599in}}%
\pgfpathlineto{\pgfqpoint{3.656373in}{1.748351in}}%
\pgfpathlineto{\pgfqpoint{3.658895in}{1.739253in}}%
\pgfpathlineto{\pgfqpoint{3.659935in}{1.735103in}}%
\pgfpathlineto{\pgfqpoint{3.669235in}{1.721855in}}%
\pgfpathlineto{\pgfqpoint{3.673118in}{1.708607in}}%
\pgfpathlineto{\pgfqpoint{3.684179in}{1.695977in}}%
\pgfpathlineto{\pgfqpoint{3.684731in}{1.695359in}}%
\pgfpathlineto{\pgfqpoint{3.687290in}{1.682111in}}%
\pgfpathlineto{\pgfqpoint{3.696971in}{1.668863in}}%
\pgfpathlineto{\pgfqpoint{3.700008in}{1.655615in}}%
\pgfpathlineto{\pgfqpoint{3.709462in}{1.643581in}}%
\pgfpathlineto{\pgfqpoint{3.710462in}{1.642367in}}%
\pgfpathlineto{\pgfqpoint{3.713716in}{1.629119in}}%
\pgfpathlineto{\pgfqpoint{3.725183in}{1.615871in}}%
\pgfpathlineto{\pgfqpoint{3.729796in}{1.602623in}}%
\pgfpathlineto{\pgfqpoint{3.734745in}{1.596524in}}%
\pgfpathlineto{\pgfqpoint{3.740964in}{1.589375in}}%
\pgfpathlineto{\pgfqpoint{3.746744in}{1.576127in}}%
\pgfpathlineto{\pgfqpoint{3.755563in}{1.562879in}}%
\pgfpathlineto{\pgfqpoint{3.760028in}{1.553287in}}%
\pgfpathlineto{\pgfqpoint{3.761776in}{1.549631in}}%
\pgfpathlineto{\pgfqpoint{3.772205in}{1.536383in}}%
\pgfpathlineto{\pgfqpoint{3.782214in}{1.523135in}}%
\pgfpathlineto{\pgfqpoint{3.785311in}{1.518851in}}%
\pgfpathlineto{\pgfqpoint{3.792514in}{1.509887in}}%
\pgfpathlineto{\pgfqpoint{3.803954in}{1.496639in}}%
\pgfpathlineto{\pgfqpoint{3.809221in}{1.483391in}}%
\pgfpathlineto{\pgfqpoint{3.810594in}{1.481783in}}%
\pgfpathlineto{\pgfqpoint{3.821446in}{1.470143in}}%
\pgfpathlineto{\pgfqpoint{3.827799in}{1.456895in}}%
\pgfpathlineto{\pgfqpoint{3.835878in}{1.449489in}}%
\pgfpathlineto{\pgfqpoint{3.843026in}{1.443647in}}%
\pgfpathlineto{\pgfqpoint{3.849549in}{1.430399in}}%
\pgfpathlineto{\pgfqpoint{3.861161in}{1.422263in}}%
\pgfpathlineto{\pgfqpoint{3.869560in}{1.417151in}}%
\pgfpathlineto{\pgfqpoint{3.874599in}{1.403903in}}%
\pgfpathlineto{\pgfqpoint{3.886444in}{1.396195in}}%
\pgfpathlineto{\pgfqpoint{3.896817in}{1.390655in}}%
\pgfpathlineto{\pgfqpoint{3.900187in}{1.377407in}}%
\pgfpathlineto{\pgfqpoint{3.911727in}{1.369555in}}%
\pgfpathlineto{\pgfqpoint{3.921531in}{1.364159in}}%
\pgfpathlineto{\pgfqpoint{3.925964in}{1.350911in}}%
\pgfpathlineto{\pgfqpoint{3.937010in}{1.345344in}}%
\pgfpathlineto{\pgfqpoint{3.955931in}{1.337663in}}%
\pgfpathlineto{\pgfqpoint{3.961683in}{1.324415in}}%
\pgfpathlineto{\pgfqpoint{3.962293in}{1.324104in}}%
\pgfpathlineto{\pgfqpoint{3.987577in}{1.314312in}}%
\pgfpathlineto{\pgfqpoint{3.999541in}{1.311167in}}%
\pgfpathlineto{\pgfqpoint{3.999867in}{1.297919in}}%
\pgfpathlineto{\pgfqpoint{4.012860in}{1.290632in}}%
\pgfpathlineto{\pgfqpoint{4.028257in}{1.284671in}}%
\pgfpathlineto{\pgfqpoint{4.035237in}{1.271423in}}%
\pgfpathlineto{\pgfqpoint{4.038143in}{1.270202in}}%
\pgfpathclose%
\pgfpathmoveto{\pgfqpoint{2.521153in}{1.732757in}}%
\pgfpathlineto{\pgfqpoint{2.531072in}{1.735103in}}%
\pgfpathlineto{\pgfqpoint{2.521153in}{1.740715in}}%
\pgfpathlineto{\pgfqpoint{2.496470in}{1.735103in}}%
\pgfpathlineto{\pgfqpoint{2.521153in}{1.732757in}}%
\pgfpathclose%
\pgfpathmoveto{\pgfqpoint{2.521153in}{1.754060in}}%
\pgfpathlineto{\pgfqpoint{2.546436in}{1.756598in}}%
\pgfpathlineto{\pgfqpoint{2.555641in}{1.761599in}}%
\pgfpathlineto{\pgfqpoint{2.546436in}{1.770712in}}%
\pgfpathlineto{\pgfqpoint{2.521153in}{1.767636in}}%
\pgfpathlineto{\pgfqpoint{2.507482in}{1.761599in}}%
\pgfpathlineto{\pgfqpoint{2.521153in}{1.754060in}}%
\pgfpathclose%
\pgfpathmoveto{\pgfqpoint{2.521153in}{1.787250in}}%
\pgfpathlineto{\pgfqpoint{2.546436in}{1.777583in}}%
\pgfpathlineto{\pgfqpoint{2.571719in}{1.784636in}}%
\pgfpathlineto{\pgfqpoint{2.576431in}{1.788095in}}%
\pgfpathlineto{\pgfqpoint{2.577963in}{1.801343in}}%
\pgfpathlineto{\pgfqpoint{2.597002in}{1.811325in}}%
\pgfpathlineto{\pgfqpoint{2.600244in}{1.814591in}}%
\pgfpathlineto{\pgfqpoint{2.608204in}{1.827839in}}%
\pgfpathlineto{\pgfqpoint{2.618503in}{1.841087in}}%
\pgfpathlineto{\pgfqpoint{2.622285in}{1.843937in}}%
\pgfpathlineto{\pgfqpoint{2.631865in}{1.854335in}}%
\pgfpathlineto{\pgfqpoint{2.635240in}{1.867583in}}%
\pgfpathlineto{\pgfqpoint{2.647569in}{1.877076in}}%
\pgfpathlineto{\pgfqpoint{2.651122in}{1.880831in}}%
\pgfpathlineto{\pgfqpoint{2.653460in}{1.894079in}}%
\pgfpathlineto{\pgfqpoint{2.672852in}{1.906299in}}%
\pgfpathlineto{\pgfqpoint{2.674117in}{1.907327in}}%
\pgfpathlineto{\pgfqpoint{2.676151in}{1.920575in}}%
\pgfpathlineto{\pgfqpoint{2.691991in}{1.933823in}}%
\pgfpathlineto{\pgfqpoint{2.690206in}{1.947071in}}%
\pgfpathlineto{\pgfqpoint{2.698135in}{1.953819in}}%
\pgfpathlineto{\pgfqpoint{2.703686in}{1.960319in}}%
\pgfpathlineto{\pgfqpoint{2.704017in}{1.973567in}}%
\pgfpathlineto{\pgfqpoint{2.718684in}{1.986815in}}%
\pgfpathlineto{\pgfqpoint{2.721878in}{2.000063in}}%
\pgfpathlineto{\pgfqpoint{2.723418in}{2.001253in}}%
\pgfpathlineto{\pgfqpoint{2.734215in}{2.013311in}}%
\pgfpathlineto{\pgfqpoint{2.738813in}{2.026559in}}%
\pgfpathlineto{\pgfqpoint{2.748571in}{2.039807in}}%
\pgfpathlineto{\pgfqpoint{2.748701in}{2.040086in}}%
\pgfpathlineto{\pgfqpoint{2.753874in}{2.053055in}}%
\pgfpathlineto{\pgfqpoint{2.759665in}{2.066303in}}%
\pgfpathlineto{\pgfqpoint{2.766390in}{2.079551in}}%
\pgfpathlineto{\pgfqpoint{2.772402in}{2.092799in}}%
\pgfpathlineto{\pgfqpoint{2.773984in}{2.094719in}}%
\pgfpathlineto{\pgfqpoint{2.782037in}{2.106047in}}%
\pgfpathlineto{\pgfqpoint{2.784698in}{2.119295in}}%
\pgfpathlineto{\pgfqpoint{2.793432in}{2.132543in}}%
\pgfpathlineto{\pgfqpoint{2.773984in}{2.132543in}}%
\pgfpathlineto{\pgfqpoint{2.752312in}{2.132543in}}%
\pgfpathlineto{\pgfqpoint{2.748701in}{2.126917in}}%
\pgfpathlineto{\pgfqpoint{2.737215in}{2.119295in}}%
\pgfpathlineto{\pgfqpoint{2.730690in}{2.106047in}}%
\pgfpathlineto{\pgfqpoint{2.723418in}{2.094505in}}%
\pgfpathlineto{\pgfqpoint{2.713936in}{2.106047in}}%
\pgfpathlineto{\pgfqpoint{2.709939in}{2.119295in}}%
\pgfpathlineto{\pgfqpoint{2.714117in}{2.132543in}}%
\pgfpathlineto{\pgfqpoint{2.698135in}{2.132543in}}%
\pgfpathlineto{\pgfqpoint{2.672852in}{2.132543in}}%
\pgfpathlineto{\pgfqpoint{2.668560in}{2.132543in}}%
\pgfpathlineto{\pgfqpoint{2.660089in}{2.119295in}}%
\pgfpathlineto{\pgfqpoint{2.659246in}{2.106047in}}%
\pgfpathlineto{\pgfqpoint{2.650769in}{2.092799in}}%
\pgfpathlineto{\pgfqpoint{2.650724in}{2.079551in}}%
\pgfpathlineto{\pgfqpoint{2.647569in}{2.075865in}}%
\pgfpathlineto{\pgfqpoint{2.640341in}{2.066303in}}%
\pgfpathlineto{\pgfqpoint{2.640677in}{2.053055in}}%
\pgfpathlineto{\pgfqpoint{2.628889in}{2.039807in}}%
\pgfpathlineto{\pgfqpoint{2.630448in}{2.026559in}}%
\pgfpathlineto{\pgfqpoint{2.622285in}{2.018910in}}%
\pgfpathlineto{\pgfqpoint{2.617234in}{2.013311in}}%
\pgfpathlineto{\pgfqpoint{2.618998in}{2.000063in}}%
\pgfpathlineto{\pgfqpoint{2.608709in}{1.986815in}}%
\pgfpathlineto{\pgfqpoint{2.610082in}{1.973567in}}%
\pgfpathlineto{\pgfqpoint{2.597002in}{1.961590in}}%
\pgfpathlineto{\pgfqpoint{2.595876in}{1.960319in}}%
\pgfpathlineto{\pgfqpoint{2.593155in}{1.947071in}}%
\pgfpathlineto{\pgfqpoint{2.582448in}{1.933823in}}%
\pgfpathlineto{\pgfqpoint{2.580370in}{1.920575in}}%
\pgfpathlineto{\pgfqpoint{2.575588in}{1.907327in}}%
\pgfpathlineto{\pgfqpoint{2.572513in}{1.894079in}}%
\pgfpathlineto{\pgfqpoint{2.571719in}{1.892523in}}%
\pgfpathlineto{\pgfqpoint{2.567809in}{1.880831in}}%
\pgfpathlineto{\pgfqpoint{2.560436in}{1.867583in}}%
\pgfpathlineto{\pgfqpoint{2.557675in}{1.854335in}}%
\pgfpathlineto{\pgfqpoint{2.546436in}{1.844154in}}%
\pgfpathlineto{\pgfqpoint{2.543554in}{1.841087in}}%
\pgfpathlineto{\pgfqpoint{2.545805in}{1.827839in}}%
\pgfpathlineto{\pgfqpoint{2.529898in}{1.814591in}}%
\pgfpathlineto{\pgfqpoint{2.542407in}{1.801343in}}%
\pgfpathlineto{\pgfqpoint{2.521153in}{1.788647in}}%
\pgfpathlineto{\pgfqpoint{2.520514in}{1.788095in}}%
\pgfpathlineto{\pgfqpoint{2.521153in}{1.787250in}}%
\pgfpathclose%
\pgfpathmoveto{\pgfqpoint{4.113993in}{1.800873in}}%
\pgfpathlineto{\pgfqpoint{4.114950in}{1.801343in}}%
\pgfpathlineto{\pgfqpoint{4.113993in}{1.802518in}}%
\pgfpathlineto{\pgfqpoint{4.112518in}{1.801343in}}%
\pgfpathlineto{\pgfqpoint{4.113993in}{1.800873in}}%
\pgfpathclose%
\pgfpathmoveto{\pgfqpoint{4.088709in}{1.818969in}}%
\pgfpathlineto{\pgfqpoint{4.113993in}{1.821483in}}%
\pgfpathlineto{\pgfqpoint{4.118413in}{1.827839in}}%
\pgfpathlineto{\pgfqpoint{4.113993in}{1.832443in}}%
\pgfpathlineto{\pgfqpoint{4.098649in}{1.841087in}}%
\pgfpathlineto{\pgfqpoint{4.110443in}{1.854335in}}%
\pgfpathlineto{\pgfqpoint{4.090589in}{1.867583in}}%
\pgfpathlineto{\pgfqpoint{4.098437in}{1.880831in}}%
\pgfpathlineto{\pgfqpoint{4.088709in}{1.892885in}}%
\pgfpathlineto{\pgfqpoint{4.086988in}{1.894079in}}%
\pgfpathlineto{\pgfqpoint{4.088709in}{1.897034in}}%
\pgfpathlineto{\pgfqpoint{4.091259in}{1.907327in}}%
\pgfpathlineto{\pgfqpoint{4.088709in}{1.910659in}}%
\pgfpathlineto{\pgfqpoint{4.078541in}{1.920575in}}%
\pgfpathlineto{\pgfqpoint{4.077899in}{1.933823in}}%
\pgfpathlineto{\pgfqpoint{4.065544in}{1.947071in}}%
\pgfpathlineto{\pgfqpoint{4.065056in}{1.960319in}}%
\pgfpathlineto{\pgfqpoint{4.063426in}{1.964323in}}%
\pgfpathlineto{\pgfqpoint{4.058340in}{1.973567in}}%
\pgfpathlineto{\pgfqpoint{4.055224in}{1.986815in}}%
\pgfpathlineto{\pgfqpoint{4.050773in}{2.000063in}}%
\pgfpathlineto{\pgfqpoint{4.045710in}{2.013311in}}%
\pgfpathlineto{\pgfqpoint{4.041286in}{2.026559in}}%
\pgfpathlineto{\pgfqpoint{4.038143in}{2.032034in}}%
\pgfpathlineto{\pgfqpoint{4.033154in}{2.039807in}}%
\pgfpathlineto{\pgfqpoint{4.029906in}{2.053055in}}%
\pgfpathlineto{\pgfqpoint{4.020591in}{2.066303in}}%
\pgfpathlineto{\pgfqpoint{4.020295in}{2.079551in}}%
\pgfpathlineto{\pgfqpoint{4.013096in}{2.092799in}}%
\pgfpathlineto{\pgfqpoint{4.012860in}{2.095538in}}%
\pgfpathlineto{\pgfqpoint{4.011878in}{2.106047in}}%
\pgfpathlineto{\pgfqpoint{4.001850in}{2.119295in}}%
\pgfpathlineto{\pgfqpoint{4.000772in}{2.132543in}}%
\pgfpathlineto{\pgfqpoint{3.987577in}{2.132543in}}%
\pgfpathlineto{\pgfqpoint{3.962293in}{2.132543in}}%
\pgfpathlineto{\pgfqpoint{3.938211in}{2.132543in}}%
\pgfpathlineto{\pgfqpoint{3.937010in}{2.131913in}}%
\pgfpathlineto{\pgfqpoint{3.936369in}{2.132543in}}%
\pgfpathlineto{\pgfqpoint{3.911727in}{2.132543in}}%
\pgfpathlineto{\pgfqpoint{3.886444in}{2.132543in}}%
\pgfpathlineto{\pgfqpoint{3.882761in}{2.132543in}}%
\pgfpathlineto{\pgfqpoint{3.886444in}{2.125163in}}%
\pgfpathlineto{\pgfqpoint{3.890058in}{2.119295in}}%
\pgfpathlineto{\pgfqpoint{3.894520in}{2.106047in}}%
\pgfpathlineto{\pgfqpoint{3.902792in}{2.092799in}}%
\pgfpathlineto{\pgfqpoint{3.908395in}{2.079551in}}%
\pgfpathlineto{\pgfqpoint{3.911727in}{2.072875in}}%
\pgfpathlineto{\pgfqpoint{3.915976in}{2.066303in}}%
\pgfpathlineto{\pgfqpoint{3.923655in}{2.053055in}}%
\pgfpathlineto{\pgfqpoint{3.928395in}{2.039807in}}%
\pgfpathlineto{\pgfqpoint{3.936697in}{2.026559in}}%
\pgfpathlineto{\pgfqpoint{3.937010in}{2.025580in}}%
\pgfpathlineto{\pgfqpoint{3.941994in}{2.013311in}}%
\pgfpathlineto{\pgfqpoint{3.952532in}{2.000063in}}%
\pgfpathlineto{\pgfqpoint{3.957130in}{1.986815in}}%
\pgfpathlineto{\pgfqpoint{3.962293in}{1.981509in}}%
\pgfpathlineto{\pgfqpoint{3.972632in}{1.973567in}}%
\pgfpathlineto{\pgfqpoint{3.975425in}{1.960319in}}%
\pgfpathlineto{\pgfqpoint{3.987577in}{1.948054in}}%
\pgfpathlineto{\pgfqpoint{3.989200in}{1.947071in}}%
\pgfpathlineto{\pgfqpoint{3.987988in}{1.933823in}}%
\pgfpathlineto{\pgfqpoint{4.003243in}{1.920575in}}%
\pgfpathlineto{\pgfqpoint{4.005275in}{1.907327in}}%
\pgfpathlineto{\pgfqpoint{4.012860in}{1.901194in}}%
\pgfpathlineto{\pgfqpoint{4.026620in}{1.894079in}}%
\pgfpathlineto{\pgfqpoint{4.027458in}{1.880831in}}%
\pgfpathlineto{\pgfqpoint{4.038143in}{1.872724in}}%
\pgfpathlineto{\pgfqpoint{4.051943in}{1.867583in}}%
\pgfpathlineto{\pgfqpoint{4.049560in}{1.854335in}}%
\pgfpathlineto{\pgfqpoint{4.063426in}{1.843367in}}%
\pgfpathlineto{\pgfqpoint{4.072509in}{1.841087in}}%
\pgfpathlineto{\pgfqpoint{4.073157in}{1.827839in}}%
\pgfpathlineto{\pgfqpoint{4.088709in}{1.818969in}}%
\pgfpathclose%
\pgfusepath{fill}%
\end{pgfscope}%
\begin{pgfscope}%
\pgfpathrectangle{\pgfqpoint{2.470586in}{0.476543in}}{\pgfqpoint{1.820388in}{1.656000in}}%
\pgfusepath{clip}%
\pgfsetbuttcap%
\pgfsetroundjoin%
\definecolor{currentfill}{rgb}{0.353710,0.557862,0.757170}%
\pgfsetfillcolor{currentfill}%
\pgfsetlinewidth{0.000000pt}%
\definecolor{currentstroke}{rgb}{0.000000,0.000000,0.000000}%
\pgfsetstrokecolor{currentstroke}%
\pgfsetdash{}{0pt}%
\pgfpathmoveto{\pgfqpoint{3.178515in}{0.688344in}}%
\pgfpathlineto{\pgfqpoint{3.203798in}{0.687254in}}%
\pgfpathlineto{\pgfqpoint{3.229082in}{0.686433in}}%
\pgfpathlineto{\pgfqpoint{3.254365in}{0.685845in}}%
\pgfpathlineto{\pgfqpoint{3.279648in}{0.685461in}}%
\pgfpathlineto{\pgfqpoint{3.304931in}{0.685259in}}%
\pgfpathlineto{\pgfqpoint{3.330214in}{0.685219in}}%
\pgfpathlineto{\pgfqpoint{3.355497in}{0.685328in}}%
\pgfpathlineto{\pgfqpoint{3.380781in}{0.685574in}}%
\pgfpathlineto{\pgfqpoint{3.406064in}{0.685946in}}%
\pgfpathlineto{\pgfqpoint{3.431347in}{0.686439in}}%
\pgfpathlineto{\pgfqpoint{3.456630in}{0.687047in}}%
\pgfpathlineto{\pgfqpoint{3.481913in}{0.687767in}}%
\pgfpathlineto{\pgfqpoint{3.504591in}{0.688511in}}%
\pgfpathlineto{\pgfqpoint{3.507196in}{0.688571in}}%
\pgfpathlineto{\pgfqpoint{3.532480in}{0.689270in}}%
\pgfpathlineto{\pgfqpoint{3.557763in}{0.690123in}}%
\pgfpathlineto{\pgfqpoint{3.583046in}{0.691151in}}%
\pgfpathlineto{\pgfqpoint{3.608329in}{0.692377in}}%
\pgfpathlineto{\pgfqpoint{3.633612in}{0.693833in}}%
\pgfpathlineto{\pgfqpoint{3.658895in}{0.695555in}}%
\pgfpathlineto{\pgfqpoint{3.684179in}{0.697589in}}%
\pgfpathlineto{\pgfqpoint{3.709462in}{0.699989in}}%
\pgfpathlineto{\pgfqpoint{3.725623in}{0.701759in}}%
\pgfpathlineto{\pgfqpoint{3.734745in}{0.702605in}}%
\pgfpathlineto{\pgfqpoint{3.760028in}{0.705128in}}%
\pgfpathlineto{\pgfqpoint{3.785311in}{0.707896in}}%
\pgfpathlineto{\pgfqpoint{3.810594in}{0.710944in}}%
\pgfpathlineto{\pgfqpoint{3.835878in}{0.714314in}}%
\pgfpathlineto{\pgfqpoint{3.840712in}{0.715007in}}%
\pgfpathlineto{\pgfqpoint{3.859498in}{0.728255in}}%
\pgfpathlineto{\pgfqpoint{3.861161in}{0.728554in}}%
\pgfpathlineto{\pgfqpoint{3.886444in}{0.733444in}}%
\pgfpathlineto{\pgfqpoint{3.911727in}{0.738976in}}%
\pgfpathlineto{\pgfqpoint{3.922241in}{0.741503in}}%
\pgfpathlineto{\pgfqpoint{3.914468in}{0.754751in}}%
\pgfpathlineto{\pgfqpoint{3.937010in}{0.759861in}}%
\pgfpathlineto{\pgfqpoint{3.962293in}{0.766552in}}%
\pgfpathlineto{\pgfqpoint{3.967242in}{0.767999in}}%
\pgfpathlineto{\pgfqpoint{3.962293in}{0.771917in}}%
\pgfpathlineto{\pgfqpoint{3.951407in}{0.781247in}}%
\pgfpathlineto{\pgfqpoint{3.962293in}{0.786297in}}%
\pgfpathlineto{\pgfqpoint{3.977533in}{0.794495in}}%
\pgfpathlineto{\pgfqpoint{3.962293in}{0.806362in}}%
\pgfpathlineto{\pgfqpoint{3.960634in}{0.807743in}}%
\pgfpathlineto{\pgfqpoint{3.962293in}{0.809664in}}%
\pgfpathlineto{\pgfqpoint{3.971032in}{0.820991in}}%
\pgfpathlineto{\pgfqpoint{3.962293in}{0.826151in}}%
\pgfpathlineto{\pgfqpoint{3.949401in}{0.834239in}}%
\pgfpathlineto{\pgfqpoint{3.950556in}{0.847487in}}%
\pgfpathlineto{\pgfqpoint{3.939722in}{0.860735in}}%
\pgfpathlineto{\pgfqpoint{3.937010in}{0.872269in}}%
\pgfpathlineto{\pgfqpoint{3.936658in}{0.873983in}}%
\pgfpathlineto{\pgfqpoint{3.933399in}{0.887231in}}%
\pgfpathlineto{\pgfqpoint{3.922842in}{0.900479in}}%
\pgfpathlineto{\pgfqpoint{3.915159in}{0.913727in}}%
\pgfpathlineto{\pgfqpoint{3.911727in}{0.916367in}}%
\pgfpathlineto{\pgfqpoint{3.899331in}{0.926975in}}%
\pgfpathlineto{\pgfqpoint{3.897917in}{0.940223in}}%
\pgfpathlineto{\pgfqpoint{3.886444in}{0.948937in}}%
\pgfpathlineto{\pgfqpoint{3.880987in}{0.953471in}}%
\pgfpathlineto{\pgfqpoint{3.880811in}{0.966719in}}%
\pgfpathlineto{\pgfqpoint{3.864026in}{0.979967in}}%
\pgfpathlineto{\pgfqpoint{3.862462in}{0.993215in}}%
\pgfpathlineto{\pgfqpoint{3.861161in}{0.994154in}}%
\pgfpathlineto{\pgfqpoint{3.845401in}{1.006463in}}%
\pgfpathlineto{\pgfqpoint{3.843404in}{1.019711in}}%
\pgfpathlineto{\pgfqpoint{3.835878in}{1.025297in}}%
\pgfpathlineto{\pgfqpoint{3.826230in}{1.032959in}}%
\pgfpathlineto{\pgfqpoint{3.824517in}{1.046207in}}%
\pgfpathlineto{\pgfqpoint{3.812754in}{1.059455in}}%
\pgfpathlineto{\pgfqpoint{3.810594in}{1.068664in}}%
\pgfpathlineto{\pgfqpoint{3.809783in}{1.072703in}}%
\pgfpathlineto{\pgfqpoint{3.798849in}{1.085951in}}%
\pgfpathlineto{\pgfqpoint{3.792073in}{1.099199in}}%
\pgfpathlineto{\pgfqpoint{3.785311in}{1.106786in}}%
\pgfpathlineto{\pgfqpoint{3.780725in}{1.112447in}}%
\pgfpathlineto{\pgfqpoint{3.773231in}{1.125695in}}%
\pgfpathlineto{\pgfqpoint{3.766810in}{1.138943in}}%
\pgfpathlineto{\pgfqpoint{3.760028in}{1.148946in}}%
\pgfpathlineto{\pgfqpoint{3.758120in}{1.152191in}}%
\pgfpathlineto{\pgfqpoint{3.753159in}{1.165439in}}%
\pgfpathlineto{\pgfqpoint{3.744066in}{1.178687in}}%
\pgfpathlineto{\pgfqpoint{3.739571in}{1.191935in}}%
\pgfpathlineto{\pgfqpoint{3.734745in}{1.197481in}}%
\pgfpathlineto{\pgfqpoint{3.728973in}{1.205183in}}%
\pgfpathlineto{\pgfqpoint{3.725853in}{1.218431in}}%
\pgfpathlineto{\pgfqpoint{3.714703in}{1.231679in}}%
\pgfpathlineto{\pgfqpoint{3.712762in}{1.244927in}}%
\pgfpathlineto{\pgfqpoint{3.709462in}{1.249217in}}%
\pgfpathlineto{\pgfqpoint{3.703492in}{1.258175in}}%
\pgfpathlineto{\pgfqpoint{3.701042in}{1.271423in}}%
\pgfpathlineto{\pgfqpoint{3.689988in}{1.284671in}}%
\pgfpathlineto{\pgfqpoint{3.686396in}{1.297919in}}%
\pgfpathlineto{\pgfqpoint{3.684179in}{1.300649in}}%
\pgfpathlineto{\pgfqpoint{3.676830in}{1.311167in}}%
\pgfpathlineto{\pgfqpoint{3.674341in}{1.324415in}}%
\pgfpathlineto{\pgfqpoint{3.667067in}{1.337663in}}%
\pgfpathlineto{\pgfqpoint{3.663371in}{1.350911in}}%
\pgfpathlineto{\pgfqpoint{3.658895in}{1.358385in}}%
\pgfpathlineto{\pgfqpoint{3.656007in}{1.364159in}}%
\pgfpathlineto{\pgfqpoint{3.652036in}{1.377407in}}%
\pgfpathlineto{\pgfqpoint{3.645742in}{1.390655in}}%
\pgfpathlineto{\pgfqpoint{3.639906in}{1.403903in}}%
\pgfpathlineto{\pgfqpoint{3.634407in}{1.417151in}}%
\pgfpathlineto{\pgfqpoint{3.633612in}{1.418686in}}%
\pgfpathlineto{\pgfqpoint{3.628768in}{1.430399in}}%
\pgfpathlineto{\pgfqpoint{3.625401in}{1.443647in}}%
\pgfpathlineto{\pgfqpoint{3.619277in}{1.456895in}}%
\pgfpathlineto{\pgfqpoint{3.615436in}{1.470143in}}%
\pgfpathlineto{\pgfqpoint{3.608329in}{1.480927in}}%
\pgfpathlineto{\pgfqpoint{3.607022in}{1.483391in}}%
\pgfpathlineto{\pgfqpoint{3.604675in}{1.496639in}}%
\pgfpathlineto{\pgfqpoint{3.598021in}{1.509887in}}%
\pgfpathlineto{\pgfqpoint{3.596508in}{1.523135in}}%
\pgfpathlineto{\pgfqpoint{3.589668in}{1.536383in}}%
\pgfpathlineto{\pgfqpoint{3.587171in}{1.549631in}}%
\pgfpathlineto{\pgfqpoint{3.583046in}{1.555956in}}%
\pgfpathlineto{\pgfqpoint{3.579451in}{1.562879in}}%
\pgfpathlineto{\pgfqpoint{3.577413in}{1.576127in}}%
\pgfpathlineto{\pgfqpoint{3.570453in}{1.589375in}}%
\pgfpathlineto{\pgfqpoint{3.567787in}{1.602623in}}%
\pgfpathlineto{\pgfqpoint{3.561130in}{1.615871in}}%
\pgfpathlineto{\pgfqpoint{3.558199in}{1.629119in}}%
\pgfpathlineto{\pgfqpoint{3.557763in}{1.630036in}}%
\pgfpathlineto{\pgfqpoint{3.553241in}{1.642367in}}%
\pgfpathlineto{\pgfqpoint{3.550325in}{1.655615in}}%
\pgfpathlineto{\pgfqpoint{3.544921in}{1.668863in}}%
\pgfpathlineto{\pgfqpoint{3.540530in}{1.682111in}}%
\pgfpathlineto{\pgfqpoint{3.536094in}{1.695359in}}%
\pgfpathlineto{\pgfqpoint{3.532480in}{1.705382in}}%
\pgfpathlineto{\pgfqpoint{3.531593in}{1.708607in}}%
\pgfpathlineto{\pgfqpoint{3.528593in}{1.721855in}}%
\pgfpathlineto{\pgfqpoint{3.524019in}{1.735103in}}%
\pgfpathlineto{\pgfqpoint{3.520433in}{1.748351in}}%
\pgfpathlineto{\pgfqpoint{3.514216in}{1.761599in}}%
\pgfpathlineto{\pgfqpoint{3.511309in}{1.774847in}}%
\pgfpathlineto{\pgfqpoint{3.507196in}{1.782651in}}%
\pgfpathlineto{\pgfqpoint{3.504929in}{1.788095in}}%
\pgfpathlineto{\pgfqpoint{3.502998in}{1.801343in}}%
\pgfpathlineto{\pgfqpoint{3.497498in}{1.814591in}}%
\pgfpathlineto{\pgfqpoint{3.495723in}{1.827839in}}%
\pgfpathlineto{\pgfqpoint{3.489269in}{1.841087in}}%
\pgfpathlineto{\pgfqpoint{3.486967in}{1.854335in}}%
\pgfpathlineto{\pgfqpoint{3.481913in}{1.862851in}}%
\pgfpathlineto{\pgfqpoint{3.479587in}{1.867583in}}%
\pgfpathlineto{\pgfqpoint{3.477510in}{1.880831in}}%
\pgfpathlineto{\pgfqpoint{3.471845in}{1.894079in}}%
\pgfpathlineto{\pgfqpoint{3.469762in}{1.907327in}}%
\pgfpathlineto{\pgfqpoint{3.463423in}{1.920575in}}%
\pgfpathlineto{\pgfqpoint{3.460053in}{1.933823in}}%
\pgfpathlineto{\pgfqpoint{3.456630in}{1.939876in}}%
\pgfpathlineto{\pgfqpoint{3.453071in}{1.947071in}}%
\pgfpathlineto{\pgfqpoint{3.449784in}{1.960319in}}%
\pgfpathlineto{\pgfqpoint{3.444547in}{1.973567in}}%
\pgfpathlineto{\pgfqpoint{3.440689in}{1.986815in}}%
\pgfpathlineto{\pgfqpoint{3.435417in}{2.000063in}}%
\pgfpathlineto{\pgfqpoint{3.431347in}{2.011168in}}%
\pgfpathlineto{\pgfqpoint{3.430617in}{2.013311in}}%
\pgfpathlineto{\pgfqpoint{3.425388in}{2.026559in}}%
\pgfpathlineto{\pgfqpoint{3.419500in}{2.039807in}}%
\pgfpathlineto{\pgfqpoint{3.413990in}{2.053055in}}%
\pgfpathlineto{\pgfqpoint{3.406064in}{2.066212in}}%
\pgfpathlineto{\pgfqpoint{3.406002in}{2.066303in}}%
\pgfpathlineto{\pgfqpoint{3.400713in}{2.079551in}}%
\pgfpathlineto{\pgfqpoint{3.392120in}{2.092799in}}%
\pgfpathlineto{\pgfqpoint{3.386723in}{2.106047in}}%
\pgfpathlineto{\pgfqpoint{3.380781in}{2.113321in}}%
\pgfpathlineto{\pgfqpoint{3.373326in}{2.119295in}}%
\pgfpathlineto{\pgfqpoint{3.363401in}{2.132543in}}%
\pgfpathlineto{\pgfqpoint{3.355497in}{2.132543in}}%
\pgfpathlineto{\pgfqpoint{3.330214in}{2.132543in}}%
\pgfpathlineto{\pgfqpoint{3.313212in}{2.132543in}}%
\pgfpathlineto{\pgfqpoint{3.304931in}{2.127166in}}%
\pgfpathlineto{\pgfqpoint{3.298614in}{2.119295in}}%
\pgfpathlineto{\pgfqpoint{3.290043in}{2.106047in}}%
\pgfpathlineto{\pgfqpoint{3.280031in}{2.092799in}}%
\pgfpathlineto{\pgfqpoint{3.279648in}{2.092006in}}%
\pgfpathlineto{\pgfqpoint{3.274999in}{2.079551in}}%
\pgfpathlineto{\pgfqpoint{3.266329in}{2.066303in}}%
\pgfpathlineto{\pgfqpoint{3.262872in}{2.053055in}}%
\pgfpathlineto{\pgfqpoint{3.254365in}{2.040254in}}%
\pgfpathlineto{\pgfqpoint{3.254072in}{2.039807in}}%
\pgfpathlineto{\pgfqpoint{3.251313in}{2.026559in}}%
\pgfpathlineto{\pgfqpoint{3.242818in}{2.013311in}}%
\pgfpathlineto{\pgfqpoint{3.240259in}{2.000063in}}%
\pgfpathlineto{\pgfqpoint{3.233732in}{1.986815in}}%
\pgfpathlineto{\pgfqpoint{3.231323in}{1.973567in}}%
\pgfpathlineto{\pgfqpoint{3.229082in}{1.969215in}}%
\pgfpathlineto{\pgfqpoint{3.224133in}{1.960319in}}%
\pgfpathlineto{\pgfqpoint{3.219926in}{1.947071in}}%
\pgfpathlineto{\pgfqpoint{3.214167in}{1.933823in}}%
\pgfpathlineto{\pgfqpoint{3.210136in}{1.920575in}}%
\pgfpathlineto{\pgfqpoint{3.206652in}{1.907327in}}%
\pgfpathlineto{\pgfqpoint{3.203798in}{1.897983in}}%
\pgfpathlineto{\pgfqpoint{3.202308in}{1.894079in}}%
\pgfpathlineto{\pgfqpoint{3.198426in}{1.880831in}}%
\pgfpathlineto{\pgfqpoint{3.192662in}{1.867583in}}%
\pgfpathlineto{\pgfqpoint{3.189993in}{1.854335in}}%
\pgfpathlineto{\pgfqpoint{3.183700in}{1.841087in}}%
\pgfpathlineto{\pgfqpoint{3.181415in}{1.827839in}}%
\pgfpathlineto{\pgfqpoint{3.178515in}{1.821472in}}%
\pgfpathlineto{\pgfqpoint{3.174648in}{1.814591in}}%
\pgfpathlineto{\pgfqpoint{3.172744in}{1.801343in}}%
\pgfpathlineto{\pgfqpoint{3.166403in}{1.788095in}}%
\pgfpathlineto{\pgfqpoint{3.164548in}{1.774847in}}%
\pgfpathlineto{\pgfqpoint{3.157960in}{1.761599in}}%
\pgfpathlineto{\pgfqpoint{3.155567in}{1.748351in}}%
\pgfpathlineto{\pgfqpoint{3.153232in}{1.742117in}}%
\pgfpathlineto{\pgfqpoint{3.149920in}{1.735103in}}%
\pgfpathlineto{\pgfqpoint{3.147004in}{1.721855in}}%
\pgfpathlineto{\pgfqpoint{3.141891in}{1.708607in}}%
\pgfpathlineto{\pgfqpoint{3.138307in}{1.695359in}}%
\pgfpathlineto{\pgfqpoint{3.133666in}{1.682111in}}%
\pgfpathlineto{\pgfqpoint{3.129259in}{1.668863in}}%
\pgfpathlineto{\pgfqpoint{3.127949in}{1.663848in}}%
\pgfpathlineto{\pgfqpoint{3.125175in}{1.655615in}}%
\pgfpathlineto{\pgfqpoint{3.118583in}{1.642367in}}%
\pgfpathlineto{\pgfqpoint{3.115570in}{1.629119in}}%
\pgfpathlineto{\pgfqpoint{3.110099in}{1.615871in}}%
\pgfpathlineto{\pgfqpoint{3.108468in}{1.602623in}}%
\pgfpathlineto{\pgfqpoint{3.102666in}{1.589934in}}%
\pgfpathlineto{\pgfqpoint{3.102335in}{1.589375in}}%
\pgfpathlineto{\pgfqpoint{3.099793in}{1.576127in}}%
\pgfpathlineto{\pgfqpoint{3.091020in}{1.562879in}}%
\pgfpathlineto{\pgfqpoint{3.089095in}{1.549631in}}%
\pgfpathlineto{\pgfqpoint{3.082515in}{1.536383in}}%
\pgfpathlineto{\pgfqpoint{3.080615in}{1.523135in}}%
\pgfpathlineto{\pgfqpoint{3.077382in}{1.516788in}}%
\pgfpathlineto{\pgfqpoint{3.073089in}{1.509887in}}%
\pgfpathlineto{\pgfqpoint{3.070014in}{1.496639in}}%
\pgfpathlineto{\pgfqpoint{3.062690in}{1.483391in}}%
\pgfpathlineto{\pgfqpoint{3.059118in}{1.470143in}}%
\pgfpathlineto{\pgfqpoint{3.053639in}{1.456895in}}%
\pgfpathlineto{\pgfqpoint{3.052099in}{1.451982in}}%
\pgfpathlineto{\pgfqpoint{3.048723in}{1.443647in}}%
\pgfpathlineto{\pgfqpoint{3.043732in}{1.430399in}}%
\pgfpathlineto{\pgfqpoint{3.038667in}{1.417151in}}%
\pgfpathlineto{\pgfqpoint{3.034545in}{1.403903in}}%
\pgfpathlineto{\pgfqpoint{3.027316in}{1.390655in}}%
\pgfpathlineto{\pgfqpoint{3.026816in}{1.388890in}}%
\pgfpathlineto{\pgfqpoint{3.022835in}{1.377407in}}%
\pgfpathlineto{\pgfqpoint{3.013305in}{1.364159in}}%
\pgfpathlineto{\pgfqpoint{3.011322in}{1.350911in}}%
\pgfpathlineto{\pgfqpoint{3.003363in}{1.337663in}}%
\pgfpathlineto{\pgfqpoint{3.001533in}{1.324968in}}%
\pgfpathlineto{\pgfqpoint{3.001434in}{1.324415in}}%
\pgfpathlineto{\pgfqpoint{2.990461in}{1.311167in}}%
\pgfpathlineto{\pgfqpoint{2.988511in}{1.297919in}}%
\pgfpathlineto{\pgfqpoint{2.978285in}{1.284671in}}%
\pgfpathlineto{\pgfqpoint{2.976366in}{1.271423in}}%
\pgfpathlineto{\pgfqpoint{2.976250in}{1.271261in}}%
\pgfpathlineto{\pgfqpoint{2.965613in}{1.258175in}}%
\pgfpathlineto{\pgfqpoint{2.962797in}{1.244927in}}%
\pgfpathlineto{\pgfqpoint{2.953829in}{1.231679in}}%
\pgfpathlineto{\pgfqpoint{2.950967in}{1.220808in}}%
\pgfpathlineto{\pgfqpoint{2.950199in}{1.218431in}}%
\pgfpathlineto{\pgfqpoint{2.939347in}{1.205183in}}%
\pgfpathlineto{\pgfqpoint{2.932132in}{1.191935in}}%
\pgfpathlineto{\pgfqpoint{2.925683in}{1.180883in}}%
\pgfpathlineto{\pgfqpoint{2.924252in}{1.178687in}}%
\pgfpathlineto{\pgfqpoint{2.916381in}{1.165439in}}%
\pgfpathlineto{\pgfqpoint{2.912161in}{1.152191in}}%
\pgfpathlineto{\pgfqpoint{2.903824in}{1.138943in}}%
\pgfpathlineto{\pgfqpoint{2.900400in}{1.129748in}}%
\pgfpathlineto{\pgfqpoint{2.898694in}{1.125695in}}%
\pgfpathlineto{\pgfqpoint{2.885630in}{1.112447in}}%
\pgfpathlineto{\pgfqpoint{2.882457in}{1.099199in}}%
\pgfpathlineto{\pgfqpoint{2.875117in}{1.092400in}}%
\pgfpathlineto{\pgfqpoint{2.867280in}{1.085951in}}%
\pgfpathlineto{\pgfqpoint{2.865366in}{1.072703in}}%
\pgfpathlineto{\pgfqpoint{2.849834in}{1.059472in}}%
\pgfpathlineto{\pgfqpoint{2.849813in}{1.059455in}}%
\pgfpathlineto{\pgfqpoint{2.849366in}{1.046207in}}%
\pgfpathlineto{\pgfqpoint{2.832402in}{1.032959in}}%
\pgfpathlineto{\pgfqpoint{2.831860in}{1.019711in}}%
\pgfpathlineto{\pgfqpoint{2.824551in}{1.014859in}}%
\pgfpathlineto{\pgfqpoint{2.811091in}{1.006463in}}%
\pgfpathlineto{\pgfqpoint{2.809871in}{0.993215in}}%
\pgfpathlineto{\pgfqpoint{2.799268in}{0.984646in}}%
\pgfpathlineto{\pgfqpoint{2.793216in}{0.979967in}}%
\pgfpathlineto{\pgfqpoint{2.794094in}{0.966719in}}%
\pgfpathlineto{\pgfqpoint{2.779636in}{0.953471in}}%
\pgfpathlineto{\pgfqpoint{2.776644in}{0.940223in}}%
\pgfpathlineto{\pgfqpoint{2.773984in}{0.938296in}}%
\pgfpathlineto{\pgfqpoint{2.757743in}{0.926975in}}%
\pgfpathlineto{\pgfqpoint{2.750596in}{0.913727in}}%
\pgfpathlineto{\pgfqpoint{2.748701in}{0.911793in}}%
\pgfpathlineto{\pgfqpoint{2.737186in}{0.900479in}}%
\pgfpathlineto{\pgfqpoint{2.728275in}{0.887231in}}%
\pgfpathlineto{\pgfqpoint{2.723418in}{0.878238in}}%
\pgfpathlineto{\pgfqpoint{2.721027in}{0.873983in}}%
\pgfpathlineto{\pgfqpoint{2.709929in}{0.860735in}}%
\pgfpathlineto{\pgfqpoint{2.708800in}{0.847487in}}%
\pgfpathlineto{\pgfqpoint{2.698135in}{0.837299in}}%
\pgfpathlineto{\pgfqpoint{2.694452in}{0.834239in}}%
\pgfpathlineto{\pgfqpoint{2.697782in}{0.820991in}}%
\pgfpathlineto{\pgfqpoint{2.674942in}{0.807743in}}%
\pgfpathlineto{\pgfqpoint{2.693320in}{0.794495in}}%
\pgfpathlineto{\pgfqpoint{2.676661in}{0.781247in}}%
\pgfpathlineto{\pgfqpoint{2.698135in}{0.774005in}}%
\pgfpathlineto{\pgfqpoint{2.717936in}{0.767999in}}%
\pgfpathlineto{\pgfqpoint{2.699337in}{0.754751in}}%
\pgfpathlineto{\pgfqpoint{2.723418in}{0.747601in}}%
\pgfpathlineto{\pgfqpoint{2.747170in}{0.741503in}}%
\pgfpathlineto{\pgfqpoint{2.723418in}{0.730445in}}%
\pgfpathlineto{\pgfqpoint{2.717846in}{0.728255in}}%
\pgfpathlineto{\pgfqpoint{2.723418in}{0.727388in}}%
\pgfpathlineto{\pgfqpoint{2.748701in}{0.723700in}}%
\pgfpathlineto{\pgfqpoint{2.773984in}{0.720340in}}%
\pgfpathlineto{\pgfqpoint{2.799268in}{0.717266in}}%
\pgfpathlineto{\pgfqpoint{2.819374in}{0.715007in}}%
\pgfpathlineto{\pgfqpoint{2.824551in}{0.708331in}}%
\pgfpathlineto{\pgfqpoint{2.831339in}{0.701759in}}%
\pgfpathlineto{\pgfqpoint{2.849834in}{0.700592in}}%
\pgfpathlineto{\pgfqpoint{2.875117in}{0.699115in}}%
\pgfpathlineto{\pgfqpoint{2.900400in}{0.697754in}}%
\pgfpathlineto{\pgfqpoint{2.925683in}{0.696499in}}%
\pgfpathlineto{\pgfqpoint{2.950967in}{0.695342in}}%
\pgfpathlineto{\pgfqpoint{2.976250in}{0.694275in}}%
\pgfpathlineto{\pgfqpoint{3.001533in}{0.693293in}}%
\pgfpathlineto{\pgfqpoint{3.026816in}{0.692391in}}%
\pgfpathlineto{\pgfqpoint{3.052099in}{0.691564in}}%
\pgfpathlineto{\pgfqpoint{3.077382in}{0.690808in}}%
\pgfpathlineto{\pgfqpoint{3.102666in}{0.690121in}}%
\pgfpathlineto{\pgfqpoint{3.127949in}{0.689499in}}%
\pgfpathlineto{\pgfqpoint{3.153232in}{0.688943in}}%
\pgfpathlineto{\pgfqpoint{3.175358in}{0.688511in}}%
\pgfpathlineto{\pgfqpoint{3.178515in}{0.688344in}}%
\pgfpathclose%
\pgfpathmoveto{\pgfqpoint{3.029846in}{0.728255in}}%
\pgfpathlineto{\pgfqpoint{3.026816in}{0.728883in}}%
\pgfpathlineto{\pgfqpoint{3.001533in}{0.735848in}}%
\pgfpathlineto{\pgfqpoint{2.987296in}{0.741503in}}%
\pgfpathlineto{\pgfqpoint{2.976250in}{0.742814in}}%
\pgfpathlineto{\pgfqpoint{2.950967in}{0.746101in}}%
\pgfpathlineto{\pgfqpoint{2.925683in}{0.749764in}}%
\pgfpathlineto{\pgfqpoint{2.900400in}{0.753858in}}%
\pgfpathlineto{\pgfqpoint{2.895319in}{0.754751in}}%
\pgfpathlineto{\pgfqpoint{2.878600in}{0.767999in}}%
\pgfpathlineto{\pgfqpoint{2.875117in}{0.768694in}}%
\pgfpathlineto{\pgfqpoint{2.849834in}{0.774065in}}%
\pgfpathlineto{\pgfqpoint{2.824551in}{0.779983in}}%
\pgfpathlineto{\pgfqpoint{2.819497in}{0.781247in}}%
\pgfpathlineto{\pgfqpoint{2.814896in}{0.794495in}}%
\pgfpathlineto{\pgfqpoint{2.799268in}{0.802723in}}%
\pgfpathlineto{\pgfqpoint{2.790274in}{0.807743in}}%
\pgfpathlineto{\pgfqpoint{2.796345in}{0.820991in}}%
\pgfpathlineto{\pgfqpoint{2.788426in}{0.834239in}}%
\pgfpathlineto{\pgfqpoint{2.791084in}{0.847487in}}%
\pgfpathlineto{\pgfqpoint{2.790199in}{0.860735in}}%
\pgfpathlineto{\pgfqpoint{2.792459in}{0.873983in}}%
\pgfpathlineto{\pgfqpoint{2.799073in}{0.887231in}}%
\pgfpathlineto{\pgfqpoint{2.799268in}{0.890018in}}%
\pgfpathlineto{\pgfqpoint{2.799893in}{0.900479in}}%
\pgfpathlineto{\pgfqpoint{2.812159in}{0.913727in}}%
\pgfpathlineto{\pgfqpoint{2.813221in}{0.926975in}}%
\pgfpathlineto{\pgfqpoint{2.824551in}{0.935428in}}%
\pgfpathlineto{\pgfqpoint{2.830731in}{0.940223in}}%
\pgfpathlineto{\pgfqpoint{2.828993in}{0.953471in}}%
\pgfpathlineto{\pgfqpoint{2.842321in}{0.966719in}}%
\pgfpathlineto{\pgfqpoint{2.837876in}{0.979967in}}%
\pgfpathlineto{\pgfqpoint{2.849834in}{0.989758in}}%
\pgfpathlineto{\pgfqpoint{2.853865in}{0.993215in}}%
\pgfpathlineto{\pgfqpoint{2.852569in}{1.006463in}}%
\pgfpathlineto{\pgfqpoint{2.871430in}{1.019711in}}%
\pgfpathlineto{\pgfqpoint{2.870103in}{1.032959in}}%
\pgfpathlineto{\pgfqpoint{2.875117in}{1.037128in}}%
\pgfpathlineto{\pgfqpoint{2.885261in}{1.046207in}}%
\pgfpathlineto{\pgfqpoint{2.884616in}{1.059455in}}%
\pgfpathlineto{\pgfqpoint{2.899246in}{1.072703in}}%
\pgfpathlineto{\pgfqpoint{2.900400in}{1.085673in}}%
\pgfpathlineto{\pgfqpoint{2.900422in}{1.085951in}}%
\pgfpathlineto{\pgfqpoint{2.913153in}{1.099199in}}%
\pgfpathlineto{\pgfqpoint{2.916014in}{1.112447in}}%
\pgfpathlineto{\pgfqpoint{2.925683in}{1.123781in}}%
\pgfpathlineto{\pgfqpoint{2.927123in}{1.125695in}}%
\pgfpathlineto{\pgfqpoint{2.931933in}{1.138943in}}%
\pgfpathlineto{\pgfqpoint{2.938857in}{1.152191in}}%
\pgfpathlineto{\pgfqpoint{2.943217in}{1.165439in}}%
\pgfpathlineto{\pgfqpoint{2.949468in}{1.178687in}}%
\pgfpathlineto{\pgfqpoint{2.950967in}{1.180981in}}%
\pgfpathlineto{\pgfqpoint{2.957510in}{1.191935in}}%
\pgfpathlineto{\pgfqpoint{2.963241in}{1.205183in}}%
\pgfpathlineto{\pgfqpoint{2.973605in}{1.218431in}}%
\pgfpathlineto{\pgfqpoint{2.976250in}{1.230586in}}%
\pgfpathlineto{\pgfqpoint{2.976449in}{1.231679in}}%
\pgfpathlineto{\pgfqpoint{2.985298in}{1.244927in}}%
\pgfpathlineto{\pgfqpoint{2.987208in}{1.258175in}}%
\pgfpathlineto{\pgfqpoint{2.997345in}{1.271423in}}%
\pgfpathlineto{\pgfqpoint{2.998741in}{1.284671in}}%
\pgfpathlineto{\pgfqpoint{3.001533in}{1.288051in}}%
\pgfpathlineto{\pgfqpoint{3.008683in}{1.297919in}}%
\pgfpathlineto{\pgfqpoint{3.010147in}{1.311167in}}%
\pgfpathlineto{\pgfqpoint{3.019995in}{1.324415in}}%
\pgfpathlineto{\pgfqpoint{3.021966in}{1.337663in}}%
\pgfpathlineto{\pgfqpoint{3.026816in}{1.345088in}}%
\pgfpathlineto{\pgfqpoint{3.030076in}{1.350911in}}%
\pgfpathlineto{\pgfqpoint{3.031942in}{1.364159in}}%
\pgfpathlineto{\pgfqpoint{3.040194in}{1.377407in}}%
\pgfpathlineto{\pgfqpoint{3.044393in}{1.390655in}}%
\pgfpathlineto{\pgfqpoint{3.051943in}{1.403903in}}%
\pgfpathlineto{\pgfqpoint{3.052099in}{1.404350in}}%
\pgfpathlineto{\pgfqpoint{3.055868in}{1.417151in}}%
\pgfpathlineto{\pgfqpoint{3.060047in}{1.430399in}}%
\pgfpathlineto{\pgfqpoint{3.064775in}{1.443647in}}%
\pgfpathlineto{\pgfqpoint{3.068912in}{1.456895in}}%
\pgfpathlineto{\pgfqpoint{3.075417in}{1.470143in}}%
\pgfpathlineto{\pgfqpoint{3.077382in}{1.477514in}}%
\pgfpathlineto{\pgfqpoint{3.078624in}{1.483391in}}%
\pgfpathlineto{\pgfqpoint{3.085363in}{1.496639in}}%
\pgfpathlineto{\pgfqpoint{3.087746in}{1.509887in}}%
\pgfpathlineto{\pgfqpoint{3.095236in}{1.523135in}}%
\pgfpathlineto{\pgfqpoint{3.097103in}{1.536383in}}%
\pgfpathlineto{\pgfqpoint{3.102666in}{1.546136in}}%
\pgfpathlineto{\pgfqpoint{3.104344in}{1.549631in}}%
\pgfpathlineto{\pgfqpoint{3.105915in}{1.562879in}}%
\pgfpathlineto{\pgfqpoint{3.113439in}{1.576127in}}%
\pgfpathlineto{\pgfqpoint{3.115500in}{1.589375in}}%
\pgfpathlineto{\pgfqpoint{3.122161in}{1.602623in}}%
\pgfpathlineto{\pgfqpoint{3.124113in}{1.615871in}}%
\pgfpathlineto{\pgfqpoint{3.127949in}{1.624149in}}%
\pgfpathlineto{\pgfqpoint{3.129830in}{1.629119in}}%
\pgfpathlineto{\pgfqpoint{3.132688in}{1.642367in}}%
\pgfpathlineto{\pgfqpoint{3.138054in}{1.655615in}}%
\pgfpathlineto{\pgfqpoint{3.142039in}{1.668863in}}%
\pgfpathlineto{\pgfqpoint{3.146683in}{1.682111in}}%
\pgfpathlineto{\pgfqpoint{3.152284in}{1.695359in}}%
\pgfpathlineto{\pgfqpoint{3.153232in}{1.698627in}}%
\pgfpathlineto{\pgfqpoint{3.155523in}{1.708607in}}%
\pgfpathlineto{\pgfqpoint{3.160273in}{1.721855in}}%
\pgfpathlineto{\pgfqpoint{3.162477in}{1.735103in}}%
\pgfpathlineto{\pgfqpoint{3.168123in}{1.748351in}}%
\pgfpathlineto{\pgfqpoint{3.170598in}{1.761599in}}%
\pgfpathlineto{\pgfqpoint{3.178295in}{1.774847in}}%
\pgfpathlineto{\pgfqpoint{3.178515in}{1.776304in}}%
\pgfpathlineto{\pgfqpoint{3.179951in}{1.788095in}}%
\pgfpathlineto{\pgfqpoint{3.185773in}{1.801343in}}%
\pgfpathlineto{\pgfqpoint{3.187465in}{1.814591in}}%
\pgfpathlineto{\pgfqpoint{3.194092in}{1.827839in}}%
\pgfpathlineto{\pgfqpoint{3.196770in}{1.841087in}}%
\pgfpathlineto{\pgfqpoint{3.203798in}{1.854300in}}%
\pgfpathlineto{\pgfqpoint{3.203815in}{1.854335in}}%
\pgfpathlineto{\pgfqpoint{3.206504in}{1.867583in}}%
\pgfpathlineto{\pgfqpoint{3.211684in}{1.880831in}}%
\pgfpathlineto{\pgfqpoint{3.215481in}{1.894079in}}%
\pgfpathlineto{\pgfqpoint{3.219762in}{1.907327in}}%
\pgfpathlineto{\pgfqpoint{3.224161in}{1.920575in}}%
\pgfpathlineto{\pgfqpoint{3.228508in}{1.933823in}}%
\pgfpathlineto{\pgfqpoint{3.229082in}{1.934976in}}%
\pgfpathlineto{\pgfqpoint{3.234881in}{1.947071in}}%
\pgfpathlineto{\pgfqpoint{3.238737in}{1.960319in}}%
\pgfpathlineto{\pgfqpoint{3.246353in}{1.973567in}}%
\pgfpathlineto{\pgfqpoint{3.248869in}{1.986815in}}%
\pgfpathlineto{\pgfqpoint{3.254365in}{1.996236in}}%
\pgfpathlineto{\pgfqpoint{3.256804in}{2.000063in}}%
\pgfpathlineto{\pgfqpoint{3.259654in}{2.013311in}}%
\pgfpathlineto{\pgfqpoint{3.269287in}{2.026559in}}%
\pgfpathlineto{\pgfqpoint{3.272337in}{2.039807in}}%
\pgfpathlineto{\pgfqpoint{3.279648in}{2.049740in}}%
\pgfpathlineto{\pgfqpoint{3.282822in}{2.053055in}}%
\pgfpathlineto{\pgfqpoint{3.288447in}{2.066303in}}%
\pgfpathlineto{\pgfqpoint{3.300126in}{2.079551in}}%
\pgfpathlineto{\pgfqpoint{3.304931in}{2.087673in}}%
\pgfpathlineto{\pgfqpoint{3.311706in}{2.092799in}}%
\pgfpathlineto{\pgfqpoint{3.330214in}{2.105234in}}%
\pgfpathlineto{\pgfqpoint{3.355497in}{2.100088in}}%
\pgfpathlineto{\pgfqpoint{3.361287in}{2.092799in}}%
\pgfpathlineto{\pgfqpoint{3.375443in}{2.079551in}}%
\pgfpathlineto{\pgfqpoint{3.380781in}{2.073133in}}%
\pgfpathlineto{\pgfqpoint{3.384229in}{2.066303in}}%
\pgfpathlineto{\pgfqpoint{3.393391in}{2.053055in}}%
\pgfpathlineto{\pgfqpoint{3.400868in}{2.039807in}}%
\pgfpathlineto{\pgfqpoint{3.406064in}{2.030026in}}%
\pgfpathlineto{\pgfqpoint{3.407703in}{2.026559in}}%
\pgfpathlineto{\pgfqpoint{3.414121in}{2.013311in}}%
\pgfpathlineto{\pgfqpoint{3.418908in}{2.000063in}}%
\pgfpathlineto{\pgfqpoint{3.424864in}{1.986815in}}%
\pgfpathlineto{\pgfqpoint{3.428825in}{1.973567in}}%
\pgfpathlineto{\pgfqpoint{3.431347in}{1.968318in}}%
\pgfpathlineto{\pgfqpoint{3.435109in}{1.960319in}}%
\pgfpathlineto{\pgfqpoint{3.438639in}{1.947071in}}%
\pgfpathlineto{\pgfqpoint{3.446086in}{1.933823in}}%
\pgfpathlineto{\pgfqpoint{3.449130in}{1.920575in}}%
\pgfpathlineto{\pgfqpoint{3.455455in}{1.907327in}}%
\pgfpathlineto{\pgfqpoint{3.456630in}{1.899670in}}%
\pgfpathlineto{\pgfqpoint{3.457622in}{1.894079in}}%
\pgfpathlineto{\pgfqpoint{3.464403in}{1.880831in}}%
\pgfpathlineto{\pgfqpoint{3.466662in}{1.867583in}}%
\pgfpathlineto{\pgfqpoint{3.473748in}{1.854335in}}%
\pgfpathlineto{\pgfqpoint{3.475745in}{1.841087in}}%
\pgfpathlineto{\pgfqpoint{3.481691in}{1.827839in}}%
\pgfpathlineto{\pgfqpoint{3.481913in}{1.826158in}}%
\pgfpathlineto{\pgfqpoint{3.483753in}{1.814591in}}%
\pgfpathlineto{\pgfqpoint{3.490201in}{1.801343in}}%
\pgfpathlineto{\pgfqpoint{3.492458in}{1.788095in}}%
\pgfpathlineto{\pgfqpoint{3.498426in}{1.774847in}}%
\pgfpathlineto{\pgfqpoint{3.501021in}{1.761599in}}%
\pgfpathlineto{\pgfqpoint{3.506440in}{1.748351in}}%
\pgfpathlineto{\pgfqpoint{3.507196in}{1.745492in}}%
\pgfpathlineto{\pgfqpoint{3.510524in}{1.735103in}}%
\pgfpathlineto{\pgfqpoint{3.515683in}{1.721855in}}%
\pgfpathlineto{\pgfqpoint{3.519191in}{1.708607in}}%
\pgfpathlineto{\pgfqpoint{3.523034in}{1.695359in}}%
\pgfpathlineto{\pgfqpoint{3.527019in}{1.682111in}}%
\pgfpathlineto{\pgfqpoint{3.530678in}{1.668863in}}%
\pgfpathlineto{\pgfqpoint{3.532480in}{1.664241in}}%
\pgfpathlineto{\pgfqpoint{3.536523in}{1.655615in}}%
\pgfpathlineto{\pgfqpoint{3.539655in}{1.642367in}}%
\pgfpathlineto{\pgfqpoint{3.545265in}{1.629119in}}%
\pgfpathlineto{\pgfqpoint{3.547584in}{1.615871in}}%
\pgfpathlineto{\pgfqpoint{3.553523in}{1.602623in}}%
\pgfpathlineto{\pgfqpoint{3.555693in}{1.589375in}}%
\pgfpathlineto{\pgfqpoint{3.557763in}{1.585223in}}%
\pgfpathlineto{\pgfqpoint{3.563166in}{1.576127in}}%
\pgfpathlineto{\pgfqpoint{3.565282in}{1.562879in}}%
\pgfpathlineto{\pgfqpoint{3.572782in}{1.549631in}}%
\pgfpathlineto{\pgfqpoint{3.574828in}{1.536383in}}%
\pgfpathlineto{\pgfqpoint{3.580922in}{1.523135in}}%
\pgfpathlineto{\pgfqpoint{3.582188in}{1.509887in}}%
\pgfpathlineto{\pgfqpoint{3.583046in}{1.508074in}}%
\pgfpathlineto{\pgfqpoint{3.589498in}{1.496639in}}%
\pgfpathlineto{\pgfqpoint{3.592138in}{1.483391in}}%
\pgfpathlineto{\pgfqpoint{3.599492in}{1.470143in}}%
\pgfpathlineto{\pgfqpoint{3.602888in}{1.456895in}}%
\pgfpathlineto{\pgfqpoint{3.608299in}{1.443647in}}%
\pgfpathlineto{\pgfqpoint{3.608329in}{1.443523in}}%
\pgfpathlineto{\pgfqpoint{3.612180in}{1.430399in}}%
\pgfpathlineto{\pgfqpoint{3.618089in}{1.417151in}}%
\pgfpathlineto{\pgfqpoint{3.623058in}{1.403903in}}%
\pgfpathlineto{\pgfqpoint{3.627842in}{1.390655in}}%
\pgfpathlineto{\pgfqpoint{3.633612in}{1.377544in}}%
\pgfpathlineto{\pgfqpoint{3.633683in}{1.377407in}}%
\pgfpathlineto{\pgfqpoint{3.637714in}{1.364159in}}%
\pgfpathlineto{\pgfqpoint{3.645019in}{1.350911in}}%
\pgfpathlineto{\pgfqpoint{3.647965in}{1.337663in}}%
\pgfpathlineto{\pgfqpoint{3.654828in}{1.324415in}}%
\pgfpathlineto{\pgfqpoint{3.656676in}{1.311167in}}%
\pgfpathlineto{\pgfqpoint{3.658895in}{1.307907in}}%
\pgfpathlineto{\pgfqpoint{3.666564in}{1.297919in}}%
\pgfpathlineto{\pgfqpoint{3.669396in}{1.284671in}}%
\pgfpathlineto{\pgfqpoint{3.679526in}{1.271423in}}%
\pgfpathlineto{\pgfqpoint{3.681423in}{1.258175in}}%
\pgfpathlineto{\pgfqpoint{3.684179in}{1.253841in}}%
\pgfpathlineto{\pgfqpoint{3.690582in}{1.244927in}}%
\pgfpathlineto{\pgfqpoint{3.691963in}{1.231679in}}%
\pgfpathlineto{\pgfqpoint{3.702196in}{1.218431in}}%
\pgfpathlineto{\pgfqpoint{3.704835in}{1.205183in}}%
\pgfpathlineto{\pgfqpoint{3.709462in}{1.198257in}}%
\pgfpathlineto{\pgfqpoint{3.714294in}{1.191935in}}%
\pgfpathlineto{\pgfqpoint{3.718288in}{1.178687in}}%
\pgfpathlineto{\pgfqpoint{3.726420in}{1.165439in}}%
\pgfpathlineto{\pgfqpoint{3.731117in}{1.152191in}}%
\pgfpathlineto{\pgfqpoint{3.734745in}{1.145061in}}%
\pgfpathlineto{\pgfqpoint{3.738304in}{1.138943in}}%
\pgfpathlineto{\pgfqpoint{3.744491in}{1.125695in}}%
\pgfpathlineto{\pgfqpoint{3.750440in}{1.112447in}}%
\pgfpathlineto{\pgfqpoint{3.760028in}{1.100126in}}%
\pgfpathlineto{\pgfqpoint{3.760812in}{1.099199in}}%
\pgfpathlineto{\pgfqpoint{3.765946in}{1.085951in}}%
\pgfpathlineto{\pgfqpoint{3.776348in}{1.072703in}}%
\pgfpathlineto{\pgfqpoint{3.777360in}{1.059455in}}%
\pgfpathlineto{\pgfqpoint{3.785311in}{1.049546in}}%
\pgfpathlineto{\pgfqpoint{3.788175in}{1.046207in}}%
\pgfpathlineto{\pgfqpoint{3.787894in}{1.032959in}}%
\pgfpathlineto{\pgfqpoint{3.803808in}{1.019711in}}%
\pgfpathlineto{\pgfqpoint{3.803656in}{1.006463in}}%
\pgfpathlineto{\pgfqpoint{3.810594in}{1.000168in}}%
\pgfpathlineto{\pgfqpoint{3.818864in}{0.993215in}}%
\pgfpathlineto{\pgfqpoint{3.818328in}{0.979967in}}%
\pgfpathlineto{\pgfqpoint{3.832860in}{0.966719in}}%
\pgfpathlineto{\pgfqpoint{3.830796in}{0.953471in}}%
\pgfpathlineto{\pgfqpoint{3.835878in}{0.948094in}}%
\pgfpathlineto{\pgfqpoint{3.844072in}{0.940223in}}%
\pgfpathlineto{\pgfqpoint{3.843372in}{0.926975in}}%
\pgfpathlineto{\pgfqpoint{3.854223in}{0.913727in}}%
\pgfpathlineto{\pgfqpoint{3.859055in}{0.900479in}}%
\pgfpathlineto{\pgfqpoint{3.861161in}{0.894109in}}%
\pgfpathlineto{\pgfqpoint{3.863725in}{0.887231in}}%
\pgfpathlineto{\pgfqpoint{3.864046in}{0.873983in}}%
\pgfpathlineto{\pgfqpoint{3.861161in}{0.866317in}}%
\pgfpathlineto{\pgfqpoint{3.858840in}{0.860735in}}%
\pgfpathlineto{\pgfqpoint{3.861161in}{0.857047in}}%
\pgfpathlineto{\pgfqpoint{3.867413in}{0.847487in}}%
\pgfpathlineto{\pgfqpoint{3.861161in}{0.840291in}}%
\pgfpathlineto{\pgfqpoint{3.855182in}{0.834239in}}%
\pgfpathlineto{\pgfqpoint{3.861161in}{0.829060in}}%
\pgfpathlineto{\pgfqpoint{3.871092in}{0.820991in}}%
\pgfpathlineto{\pgfqpoint{3.861161in}{0.815246in}}%
\pgfpathlineto{\pgfqpoint{3.846079in}{0.807743in}}%
\pgfpathlineto{\pgfqpoint{3.855198in}{0.794495in}}%
\pgfpathlineto{\pgfqpoint{3.835878in}{0.788631in}}%
\pgfpathlineto{\pgfqpoint{3.810594in}{0.782162in}}%
\pgfpathlineto{\pgfqpoint{3.806579in}{0.781247in}}%
\pgfpathlineto{\pgfqpoint{3.807821in}{0.767999in}}%
\pgfpathlineto{\pgfqpoint{3.785311in}{0.764149in}}%
\pgfpathlineto{\pgfqpoint{3.760028in}{0.760301in}}%
\pgfpathlineto{\pgfqpoint{3.734745in}{0.756877in}}%
\pgfpathlineto{\pgfqpoint{3.717379in}{0.754751in}}%
\pgfpathlineto{\pgfqpoint{3.709462in}{0.751130in}}%
\pgfpathlineto{\pgfqpoint{3.684179in}{0.742080in}}%
\pgfpathlineto{\pgfqpoint{3.682180in}{0.741503in}}%
\pgfpathlineto{\pgfqpoint{3.658895in}{0.738710in}}%
\pgfpathlineto{\pgfqpoint{3.633612in}{0.735953in}}%
\pgfpathlineto{\pgfqpoint{3.608329in}{0.733447in}}%
\pgfpathlineto{\pgfqpoint{3.583046in}{0.731166in}}%
\pgfpathlineto{\pgfqpoint{3.557763in}{0.729089in}}%
\pgfpathlineto{\pgfqpoint{3.546652in}{0.728255in}}%
\pgfpathlineto{\pgfqpoint{3.532480in}{0.726810in}}%
\pgfpathlineto{\pgfqpoint{3.507196in}{0.724654in}}%
\pgfpathlineto{\pgfqpoint{3.481913in}{0.722911in}}%
\pgfpathlineto{\pgfqpoint{3.456630in}{0.721521in}}%
\pgfpathlineto{\pgfqpoint{3.431347in}{0.720436in}}%
\pgfpathlineto{\pgfqpoint{3.406064in}{0.719619in}}%
\pgfpathlineto{\pgfqpoint{3.380781in}{0.719038in}}%
\pgfpathlineto{\pgfqpoint{3.355497in}{0.718670in}}%
\pgfpathlineto{\pgfqpoint{3.330214in}{0.718495in}}%
\pgfpathlineto{\pgfqpoint{3.304931in}{0.718497in}}%
\pgfpathlineto{\pgfqpoint{3.279648in}{0.718665in}}%
\pgfpathlineto{\pgfqpoint{3.254365in}{0.718990in}}%
\pgfpathlineto{\pgfqpoint{3.229082in}{0.719464in}}%
\pgfpathlineto{\pgfqpoint{3.203798in}{0.720083in}}%
\pgfpathlineto{\pgfqpoint{3.178515in}{0.720844in}}%
\pgfpathlineto{\pgfqpoint{3.153232in}{0.721746in}}%
\pgfpathlineto{\pgfqpoint{3.127949in}{0.722790in}}%
\pgfpathlineto{\pgfqpoint{3.102666in}{0.723978in}}%
\pgfpathlineto{\pgfqpoint{3.077382in}{0.725314in}}%
\pgfpathlineto{\pgfqpoint{3.052099in}{0.726803in}}%
\pgfpathlineto{\pgfqpoint{3.029846in}{0.728255in}}%
\pgfpathclose%
\pgfpathmoveto{\pgfqpoint{2.698135in}{1.425520in}}%
\pgfpathlineto{\pgfqpoint{2.723418in}{1.422488in}}%
\pgfpathlineto{\pgfqpoint{2.748701in}{1.426136in}}%
\pgfpathlineto{\pgfqpoint{2.762086in}{1.430399in}}%
\pgfpathlineto{\pgfqpoint{2.755930in}{1.443647in}}%
\pgfpathlineto{\pgfqpoint{2.773984in}{1.447722in}}%
\pgfpathlineto{\pgfqpoint{2.797966in}{1.456895in}}%
\pgfpathlineto{\pgfqpoint{2.797442in}{1.470143in}}%
\pgfpathlineto{\pgfqpoint{2.799268in}{1.470775in}}%
\pgfpathlineto{\pgfqpoint{2.820930in}{1.483391in}}%
\pgfpathlineto{\pgfqpoint{2.819948in}{1.496639in}}%
\pgfpathlineto{\pgfqpoint{2.824551in}{1.498996in}}%
\pgfpathlineto{\pgfqpoint{2.838597in}{1.509887in}}%
\pgfpathlineto{\pgfqpoint{2.844820in}{1.523135in}}%
\pgfpathlineto{\pgfqpoint{2.849834in}{1.526004in}}%
\pgfpathlineto{\pgfqpoint{2.862301in}{1.536383in}}%
\pgfpathlineto{\pgfqpoint{2.873616in}{1.549631in}}%
\pgfpathlineto{\pgfqpoint{2.875117in}{1.551479in}}%
\pgfpathlineto{\pgfqpoint{2.881966in}{1.562879in}}%
\pgfpathlineto{\pgfqpoint{2.892580in}{1.576127in}}%
\pgfpathlineto{\pgfqpoint{2.899540in}{1.589375in}}%
\pgfpathlineto{\pgfqpoint{2.900400in}{1.590145in}}%
\pgfpathlineto{\pgfqpoint{2.911958in}{1.602623in}}%
\pgfpathlineto{\pgfqpoint{2.916581in}{1.615871in}}%
\pgfpathlineto{\pgfqpoint{2.925683in}{1.623971in}}%
\pgfpathlineto{\pgfqpoint{2.930731in}{1.629119in}}%
\pgfpathlineto{\pgfqpoint{2.932706in}{1.642367in}}%
\pgfpathlineto{\pgfqpoint{2.946342in}{1.655615in}}%
\pgfpathlineto{\pgfqpoint{2.948315in}{1.668863in}}%
\pgfpathlineto{\pgfqpoint{2.950967in}{1.671493in}}%
\pgfpathlineto{\pgfqpoint{2.959923in}{1.682111in}}%
\pgfpathlineto{\pgfqpoint{2.961629in}{1.695359in}}%
\pgfpathlineto{\pgfqpoint{2.974895in}{1.708607in}}%
\pgfpathlineto{\pgfqpoint{2.976250in}{1.712585in}}%
\pgfpathlineto{\pgfqpoint{2.979251in}{1.721855in}}%
\pgfpathlineto{\pgfqpoint{2.989778in}{1.735103in}}%
\pgfpathlineto{\pgfqpoint{2.994000in}{1.748351in}}%
\pgfpathlineto{\pgfqpoint{3.000995in}{1.761599in}}%
\pgfpathlineto{\pgfqpoint{3.001533in}{1.762983in}}%
\pgfpathlineto{\pgfqpoint{3.006021in}{1.774847in}}%
\pgfpathlineto{\pgfqpoint{3.012471in}{1.788095in}}%
\pgfpathlineto{\pgfqpoint{3.019112in}{1.801343in}}%
\pgfpathlineto{\pgfqpoint{3.023940in}{1.814591in}}%
\pgfpathlineto{\pgfqpoint{3.026816in}{1.819069in}}%
\pgfpathlineto{\pgfqpoint{3.032401in}{1.827839in}}%
\pgfpathlineto{\pgfqpoint{3.035841in}{1.841087in}}%
\pgfpathlineto{\pgfqpoint{3.044029in}{1.854335in}}%
\pgfpathlineto{\pgfqpoint{3.046460in}{1.867583in}}%
\pgfpathlineto{\pgfqpoint{3.052099in}{1.876394in}}%
\pgfpathlineto{\pgfqpoint{3.054924in}{1.880831in}}%
\pgfpathlineto{\pgfqpoint{3.057300in}{1.894079in}}%
\pgfpathlineto{\pgfqpoint{3.066815in}{1.907327in}}%
\pgfpathlineto{\pgfqpoint{3.069222in}{1.920575in}}%
\pgfpathlineto{\pgfqpoint{3.076512in}{1.933823in}}%
\pgfpathlineto{\pgfqpoint{3.077382in}{1.939949in}}%
\pgfpathlineto{\pgfqpoint{3.078508in}{1.947071in}}%
\pgfpathlineto{\pgfqpoint{3.085410in}{1.960319in}}%
\pgfpathlineto{\pgfqpoint{3.088618in}{1.973567in}}%
\pgfpathlineto{\pgfqpoint{3.095161in}{1.986815in}}%
\pgfpathlineto{\pgfqpoint{3.099474in}{2.000063in}}%
\pgfpathlineto{\pgfqpoint{3.102666in}{2.008224in}}%
\pgfpathlineto{\pgfqpoint{3.104788in}{2.013311in}}%
\pgfpathlineto{\pgfqpoint{3.110573in}{2.026559in}}%
\pgfpathlineto{\pgfqpoint{3.114473in}{2.039807in}}%
\pgfpathlineto{\pgfqpoint{3.119551in}{2.053055in}}%
\pgfpathlineto{\pgfqpoint{3.122131in}{2.066303in}}%
\pgfpathlineto{\pgfqpoint{3.127949in}{2.078881in}}%
\pgfpathlineto{\pgfqpoint{3.128309in}{2.079551in}}%
\pgfpathlineto{\pgfqpoint{3.131201in}{2.092799in}}%
\pgfpathlineto{\pgfqpoint{3.138719in}{2.106047in}}%
\pgfpathlineto{\pgfqpoint{3.140565in}{2.119295in}}%
\pgfpathlineto{\pgfqpoint{3.146686in}{2.132543in}}%
\pgfpathlineto{\pgfqpoint{3.132271in}{2.132543in}}%
\pgfpathlineto{\pgfqpoint{3.127949in}{2.124299in}}%
\pgfpathlineto{\pgfqpoint{3.125373in}{2.119295in}}%
\pgfpathlineto{\pgfqpoint{3.123582in}{2.106047in}}%
\pgfpathlineto{\pgfqpoint{3.116257in}{2.092799in}}%
\pgfpathlineto{\pgfqpoint{3.113511in}{2.079551in}}%
\pgfpathlineto{\pgfqpoint{3.107122in}{2.066303in}}%
\pgfpathlineto{\pgfqpoint{3.104018in}{2.053055in}}%
\pgfpathlineto{\pgfqpoint{3.102666in}{2.049761in}}%
\pgfpathlineto{\pgfqpoint{3.098654in}{2.039807in}}%
\pgfpathlineto{\pgfqpoint{3.094107in}{2.026559in}}%
\pgfpathlineto{\pgfqpoint{3.088895in}{2.013311in}}%
\pgfpathlineto{\pgfqpoint{3.082900in}{2.000063in}}%
\pgfpathlineto{\pgfqpoint{3.078849in}{1.986815in}}%
\pgfpathlineto{\pgfqpoint{3.077382in}{1.984172in}}%
\pgfpathlineto{\pgfqpoint{3.070989in}{1.973567in}}%
\pgfpathlineto{\pgfqpoint{3.068165in}{1.960319in}}%
\pgfpathlineto{\pgfqpoint{3.060629in}{1.947071in}}%
\pgfpathlineto{\pgfqpoint{3.059335in}{1.933823in}}%
\pgfpathlineto{\pgfqpoint{3.052099in}{1.922321in}}%
\pgfpathlineto{\pgfqpoint{3.050837in}{1.920575in}}%
\pgfpathlineto{\pgfqpoint{3.048505in}{1.907327in}}%
\pgfpathlineto{\pgfqpoint{3.037733in}{1.894079in}}%
\pgfpathlineto{\pgfqpoint{3.035565in}{1.880831in}}%
\pgfpathlineto{\pgfqpoint{3.026816in}{1.867754in}}%
\pgfpathlineto{\pgfqpoint{3.026682in}{1.867583in}}%
\pgfpathlineto{\pgfqpoint{3.023712in}{1.854335in}}%
\pgfpathlineto{\pgfqpoint{3.014397in}{1.841087in}}%
\pgfpathlineto{\pgfqpoint{3.010390in}{1.827839in}}%
\pgfpathlineto{\pgfqpoint{3.002427in}{1.814591in}}%
\pgfpathlineto{\pgfqpoint{3.001533in}{1.812415in}}%
\pgfpathlineto{\pgfqpoint{2.995968in}{1.801343in}}%
\pgfpathlineto{\pgfqpoint{2.989272in}{1.788095in}}%
\pgfpathlineto{\pgfqpoint{2.981093in}{1.774847in}}%
\pgfpathlineto{\pgfqpoint{2.977101in}{1.761599in}}%
\pgfpathlineto{\pgfqpoint{2.976250in}{1.760170in}}%
\pgfpathlineto{\pgfqpoint{2.967167in}{1.748351in}}%
\pgfpathlineto{\pgfqpoint{2.963945in}{1.735103in}}%
\pgfpathlineto{\pgfqpoint{2.950967in}{1.722483in}}%
\pgfpathlineto{\pgfqpoint{2.950023in}{1.721855in}}%
\pgfpathlineto{\pgfqpoint{2.946283in}{1.708607in}}%
\pgfpathlineto{\pgfqpoint{2.928736in}{1.695359in}}%
\pgfpathlineto{\pgfqpoint{2.929113in}{1.682111in}}%
\pgfpathlineto{\pgfqpoint{2.925683in}{1.678720in}}%
\pgfpathlineto{\pgfqpoint{2.911791in}{1.668863in}}%
\pgfpathlineto{\pgfqpoint{2.911757in}{1.655615in}}%
\pgfpathlineto{\pgfqpoint{2.900400in}{1.646611in}}%
\pgfpathlineto{\pgfqpoint{2.892429in}{1.642367in}}%
\pgfpathlineto{\pgfqpoint{2.891290in}{1.629119in}}%
\pgfpathlineto{\pgfqpoint{2.875117in}{1.616835in}}%
\pgfpathlineto{\pgfqpoint{2.873001in}{1.615871in}}%
\pgfpathlineto{\pgfqpoint{2.862191in}{1.602623in}}%
\pgfpathlineto{\pgfqpoint{2.849834in}{1.593128in}}%
\pgfpathlineto{\pgfqpoint{2.838811in}{1.589375in}}%
\pgfpathlineto{\pgfqpoint{2.824551in}{1.581427in}}%
\pgfpathlineto{\pgfqpoint{2.810614in}{1.589375in}}%
\pgfpathlineto{\pgfqpoint{2.809198in}{1.602623in}}%
\pgfpathlineto{\pgfqpoint{2.803788in}{1.615871in}}%
\pgfpathlineto{\pgfqpoint{2.805355in}{1.629119in}}%
\pgfpathlineto{\pgfqpoint{2.810536in}{1.642367in}}%
\pgfpathlineto{\pgfqpoint{2.810992in}{1.655615in}}%
\pgfpathlineto{\pgfqpoint{2.818962in}{1.668863in}}%
\pgfpathlineto{\pgfqpoint{2.818329in}{1.682111in}}%
\pgfpathlineto{\pgfqpoint{2.824551in}{1.690504in}}%
\pgfpathlineto{\pgfqpoint{2.829455in}{1.695359in}}%
\pgfpathlineto{\pgfqpoint{2.826214in}{1.708607in}}%
\pgfpathlineto{\pgfqpoint{2.835873in}{1.721855in}}%
\pgfpathlineto{\pgfqpoint{2.831328in}{1.735103in}}%
\pgfpathlineto{\pgfqpoint{2.841146in}{1.748351in}}%
\pgfpathlineto{\pgfqpoint{2.839329in}{1.761599in}}%
\pgfpathlineto{\pgfqpoint{2.849834in}{1.774281in}}%
\pgfpathlineto{\pgfqpoint{2.850391in}{1.774847in}}%
\pgfpathlineto{\pgfqpoint{2.849834in}{1.777731in}}%
\pgfpathlineto{\pgfqpoint{2.848319in}{1.788095in}}%
\pgfpathlineto{\pgfqpoint{2.849834in}{1.790445in}}%
\pgfpathlineto{\pgfqpoint{2.857592in}{1.801343in}}%
\pgfpathlineto{\pgfqpoint{2.856005in}{1.814591in}}%
\pgfpathlineto{\pgfqpoint{2.865222in}{1.827839in}}%
\pgfpathlineto{\pgfqpoint{2.865808in}{1.841087in}}%
\pgfpathlineto{\pgfqpoint{2.873336in}{1.854335in}}%
\pgfpathlineto{\pgfqpoint{2.874834in}{1.867583in}}%
\pgfpathlineto{\pgfqpoint{2.875117in}{1.868203in}}%
\pgfpathlineto{\pgfqpoint{2.880920in}{1.880831in}}%
\pgfpathlineto{\pgfqpoint{2.883822in}{1.894079in}}%
\pgfpathlineto{\pgfqpoint{2.887326in}{1.907327in}}%
\pgfpathlineto{\pgfqpoint{2.890338in}{1.920575in}}%
\pgfpathlineto{\pgfqpoint{2.893130in}{1.933823in}}%
\pgfpathlineto{\pgfqpoint{2.898941in}{1.947071in}}%
\pgfpathlineto{\pgfqpoint{2.900400in}{1.953540in}}%
\pgfpathlineto{\pgfqpoint{2.901877in}{1.960319in}}%
\pgfpathlineto{\pgfqpoint{2.908814in}{1.973567in}}%
\pgfpathlineto{\pgfqpoint{2.909476in}{1.986815in}}%
\pgfpathlineto{\pgfqpoint{2.915618in}{2.000063in}}%
\pgfpathlineto{\pgfqpoint{2.915955in}{2.013311in}}%
\pgfpathlineto{\pgfqpoint{2.923318in}{2.026559in}}%
\pgfpathlineto{\pgfqpoint{2.923593in}{2.039807in}}%
\pgfpathlineto{\pgfqpoint{2.925683in}{2.043700in}}%
\pgfpathlineto{\pgfqpoint{2.930567in}{2.053055in}}%
\pgfpathlineto{\pgfqpoint{2.930903in}{2.066303in}}%
\pgfpathlineto{\pgfqpoint{2.937147in}{2.079551in}}%
\pgfpathlineto{\pgfqpoint{2.937929in}{2.092799in}}%
\pgfpathlineto{\pgfqpoint{2.943288in}{2.106047in}}%
\pgfpathlineto{\pgfqpoint{2.944598in}{2.119295in}}%
\pgfpathlineto{\pgfqpoint{2.950234in}{2.132543in}}%
\pgfpathlineto{\pgfqpoint{2.935477in}{2.132543in}}%
\pgfpathlineto{\pgfqpoint{2.929890in}{2.119295in}}%
\pgfpathlineto{\pgfqpoint{2.928580in}{2.106047in}}%
\pgfpathlineto{\pgfqpoint{2.925683in}{2.098967in}}%
\pgfpathlineto{\pgfqpoint{2.922743in}{2.092799in}}%
\pgfpathlineto{\pgfqpoint{2.921558in}{2.079551in}}%
\pgfpathlineto{\pgfqpoint{2.914968in}{2.066303in}}%
\pgfpathlineto{\pgfqpoint{2.914153in}{2.053055in}}%
\pgfpathlineto{\pgfqpoint{2.907381in}{2.039807in}}%
\pgfpathlineto{\pgfqpoint{2.906303in}{2.026559in}}%
\pgfpathlineto{\pgfqpoint{2.900400in}{2.015195in}}%
\pgfpathlineto{\pgfqpoint{2.899319in}{2.013311in}}%
\pgfpathlineto{\pgfqpoint{2.897849in}{2.000063in}}%
\pgfpathlineto{\pgfqpoint{2.891812in}{1.986815in}}%
\pgfpathlineto{\pgfqpoint{2.889985in}{1.973567in}}%
\pgfpathlineto{\pgfqpoint{2.883624in}{1.960319in}}%
\pgfpathlineto{\pgfqpoint{2.879744in}{1.947071in}}%
\pgfpathlineto{\pgfqpoint{2.875117in}{1.935176in}}%
\pgfpathlineto{\pgfqpoint{2.874557in}{1.933823in}}%
\pgfpathlineto{\pgfqpoint{2.870210in}{1.920575in}}%
\pgfpathlineto{\pgfqpoint{2.867321in}{1.907327in}}%
\pgfpathlineto{\pgfqpoint{2.862548in}{1.894079in}}%
\pgfpathlineto{\pgfqpoint{2.859763in}{1.880831in}}%
\pgfpathlineto{\pgfqpoint{2.853082in}{1.867583in}}%
\pgfpathlineto{\pgfqpoint{2.851204in}{1.854335in}}%
\pgfpathlineto{\pgfqpoint{2.849834in}{1.852188in}}%
\pgfpathlineto{\pgfqpoint{2.842388in}{1.841087in}}%
\pgfpathlineto{\pgfqpoint{2.841012in}{1.827839in}}%
\pgfpathlineto{\pgfqpoint{2.831890in}{1.814591in}}%
\pgfpathlineto{\pgfqpoint{2.832017in}{1.801343in}}%
\pgfpathlineto{\pgfqpoint{2.824551in}{1.790030in}}%
\pgfpathlineto{\pgfqpoint{2.823257in}{1.788095in}}%
\pgfpathlineto{\pgfqpoint{2.823211in}{1.774847in}}%
\pgfpathlineto{\pgfqpoint{2.812400in}{1.761599in}}%
\pgfpathlineto{\pgfqpoint{2.811587in}{1.748351in}}%
\pgfpathlineto{\pgfqpoint{2.802932in}{1.735103in}}%
\pgfpathlineto{\pgfqpoint{2.802962in}{1.721855in}}%
\pgfpathlineto{\pgfqpoint{2.799268in}{1.715882in}}%
\pgfpathlineto{\pgfqpoint{2.795034in}{1.708607in}}%
\pgfpathlineto{\pgfqpoint{2.792875in}{1.695359in}}%
\pgfpathlineto{\pgfqpoint{2.784182in}{1.682111in}}%
\pgfpathlineto{\pgfqpoint{2.779708in}{1.668863in}}%
\pgfpathlineto{\pgfqpoint{2.773984in}{1.657768in}}%
\pgfpathlineto{\pgfqpoint{2.772958in}{1.655615in}}%
\pgfpathlineto{\pgfqpoint{2.766562in}{1.642367in}}%
\pgfpathlineto{\pgfqpoint{2.762872in}{1.629119in}}%
\pgfpathlineto{\pgfqpoint{2.755856in}{1.615871in}}%
\pgfpathlineto{\pgfqpoint{2.755815in}{1.602623in}}%
\pgfpathlineto{\pgfqpoint{2.748701in}{1.592921in}}%
\pgfpathlineto{\pgfqpoint{2.746159in}{1.589375in}}%
\pgfpathlineto{\pgfqpoint{2.746043in}{1.576127in}}%
\pgfpathlineto{\pgfqpoint{2.731737in}{1.562879in}}%
\pgfpathlineto{\pgfqpoint{2.735289in}{1.549631in}}%
\pgfpathlineto{\pgfqpoint{2.723418in}{1.538440in}}%
\pgfpathlineto{\pgfqpoint{2.721429in}{1.536383in}}%
\pgfpathlineto{\pgfqpoint{2.723418in}{1.532197in}}%
\pgfpathlineto{\pgfqpoint{2.729813in}{1.523135in}}%
\pgfpathlineto{\pgfqpoint{2.723418in}{1.518801in}}%
\pgfpathlineto{\pgfqpoint{2.712493in}{1.509887in}}%
\pgfpathlineto{\pgfqpoint{2.718740in}{1.496639in}}%
\pgfpathlineto{\pgfqpoint{2.698135in}{1.484140in}}%
\pgfpathlineto{\pgfqpoint{2.697109in}{1.483391in}}%
\pgfpathlineto{\pgfqpoint{2.698135in}{1.481479in}}%
\pgfpathlineto{\pgfqpoint{2.708644in}{1.470143in}}%
\pgfpathlineto{\pgfqpoint{2.698135in}{1.463973in}}%
\pgfpathlineto{\pgfqpoint{2.689169in}{1.456895in}}%
\pgfpathlineto{\pgfqpoint{2.698135in}{1.448182in}}%
\pgfpathlineto{\pgfqpoint{2.710561in}{1.443647in}}%
\pgfpathlineto{\pgfqpoint{2.698135in}{1.437677in}}%
\pgfpathlineto{\pgfqpoint{2.689733in}{1.430399in}}%
\pgfpathlineto{\pgfqpoint{2.698135in}{1.425520in}}%
\pgfpathclose%
\pgfpathmoveto{\pgfqpoint{3.911727in}{1.449902in}}%
\pgfpathlineto{\pgfqpoint{3.937010in}{1.443871in}}%
\pgfpathlineto{\pgfqpoint{3.955238in}{1.456895in}}%
\pgfpathlineto{\pgfqpoint{3.952280in}{1.470143in}}%
\pgfpathlineto{\pgfqpoint{3.945922in}{1.483391in}}%
\pgfpathlineto{\pgfqpoint{3.946615in}{1.496639in}}%
\pgfpathlineto{\pgfqpoint{3.938505in}{1.509887in}}%
\pgfpathlineto{\pgfqpoint{3.945776in}{1.523135in}}%
\pgfpathlineto{\pgfqpoint{3.937010in}{1.536236in}}%
\pgfpathlineto{\pgfqpoint{3.936892in}{1.536383in}}%
\pgfpathlineto{\pgfqpoint{3.937010in}{1.536954in}}%
\pgfpathlineto{\pgfqpoint{3.938769in}{1.549631in}}%
\pgfpathlineto{\pgfqpoint{3.937010in}{1.551453in}}%
\pgfpathlineto{\pgfqpoint{3.924718in}{1.562879in}}%
\pgfpathlineto{\pgfqpoint{3.926857in}{1.576127in}}%
\pgfpathlineto{\pgfqpoint{3.913973in}{1.589375in}}%
\pgfpathlineto{\pgfqpoint{3.916175in}{1.602623in}}%
\pgfpathlineto{\pgfqpoint{3.911727in}{1.608300in}}%
\pgfpathlineto{\pgfqpoint{3.904962in}{1.615871in}}%
\pgfpathlineto{\pgfqpoint{3.905649in}{1.629119in}}%
\pgfpathlineto{\pgfqpoint{3.896814in}{1.642367in}}%
\pgfpathlineto{\pgfqpoint{3.895743in}{1.655615in}}%
\pgfpathlineto{\pgfqpoint{3.886731in}{1.668863in}}%
\pgfpathlineto{\pgfqpoint{3.886444in}{1.669963in}}%
\pgfpathlineto{\pgfqpoint{3.883308in}{1.682111in}}%
\pgfpathlineto{\pgfqpoint{3.877594in}{1.695359in}}%
\pgfpathlineto{\pgfqpoint{3.873931in}{1.708607in}}%
\pgfpathlineto{\pgfqpoint{3.870434in}{1.721855in}}%
\pgfpathlineto{\pgfqpoint{3.864513in}{1.735103in}}%
\pgfpathlineto{\pgfqpoint{3.861161in}{1.744579in}}%
\pgfpathlineto{\pgfqpoint{3.859793in}{1.748351in}}%
\pgfpathlineto{\pgfqpoint{3.851976in}{1.761599in}}%
\pgfpathlineto{\pgfqpoint{3.849958in}{1.774847in}}%
\pgfpathlineto{\pgfqpoint{3.842121in}{1.788095in}}%
\pgfpathlineto{\pgfqpoint{3.841310in}{1.801343in}}%
\pgfpathlineto{\pgfqpoint{3.835878in}{1.810742in}}%
\pgfpathlineto{\pgfqpoint{3.833648in}{1.814591in}}%
\pgfpathlineto{\pgfqpoint{3.832861in}{1.827839in}}%
\pgfpathlineto{\pgfqpoint{3.825006in}{1.841087in}}%
\pgfpathlineto{\pgfqpoint{3.823761in}{1.854335in}}%
\pgfpathlineto{\pgfqpoint{3.815073in}{1.867583in}}%
\pgfpathlineto{\pgfqpoint{3.813738in}{1.880831in}}%
\pgfpathlineto{\pgfqpoint{3.810594in}{1.887073in}}%
\pgfpathlineto{\pgfqpoint{3.807160in}{1.894079in}}%
\pgfpathlineto{\pgfqpoint{3.805543in}{1.907327in}}%
\pgfpathlineto{\pgfqpoint{3.799321in}{1.920575in}}%
\pgfpathlineto{\pgfqpoint{3.796418in}{1.933823in}}%
\pgfpathlineto{\pgfqpoint{3.789979in}{1.947071in}}%
\pgfpathlineto{\pgfqpoint{3.786619in}{1.960319in}}%
\pgfpathlineto{\pgfqpoint{3.785311in}{1.964564in}}%
\pgfpathlineto{\pgfqpoint{3.782750in}{1.973567in}}%
\pgfpathlineto{\pgfqpoint{3.778939in}{1.986815in}}%
\pgfpathlineto{\pgfqpoint{3.775859in}{2.000063in}}%
\pgfpathlineto{\pgfqpoint{3.771491in}{2.013311in}}%
\pgfpathlineto{\pgfqpoint{3.768538in}{2.026559in}}%
\pgfpathlineto{\pgfqpoint{3.762629in}{2.039807in}}%
\pgfpathlineto{\pgfqpoint{3.760028in}{2.052453in}}%
\pgfpathlineto{\pgfqpoint{3.759919in}{2.053055in}}%
\pgfpathlineto{\pgfqpoint{3.753814in}{2.066303in}}%
\pgfpathlineto{\pgfqpoint{3.752570in}{2.079551in}}%
\pgfpathlineto{\pgfqpoint{3.747462in}{2.092799in}}%
\pgfpathlineto{\pgfqpoint{3.746095in}{2.106047in}}%
\pgfpathlineto{\pgfqpoint{3.739781in}{2.119295in}}%
\pgfpathlineto{\pgfqpoint{3.738176in}{2.132543in}}%
\pgfpathlineto{\pgfqpoint{3.734745in}{2.132543in}}%
\pgfpathlineto{\pgfqpoint{3.723441in}{2.132543in}}%
\pgfpathlineto{\pgfqpoint{3.724483in}{2.119295in}}%
\pgfpathlineto{\pgfqpoint{3.730417in}{2.106047in}}%
\pgfpathlineto{\pgfqpoint{3.731390in}{2.092799in}}%
\pgfpathlineto{\pgfqpoint{3.734745in}{2.084312in}}%
\pgfpathlineto{\pgfqpoint{3.736694in}{2.079551in}}%
\pgfpathlineto{\pgfqpoint{3.737701in}{2.066303in}}%
\pgfpathlineto{\pgfqpoint{3.743840in}{2.053055in}}%
\pgfpathlineto{\pgfqpoint{3.746078in}{2.039807in}}%
\pgfpathlineto{\pgfqpoint{3.751073in}{2.026559in}}%
\pgfpathlineto{\pgfqpoint{3.753807in}{2.013311in}}%
\pgfpathlineto{\pgfqpoint{3.757777in}{2.000063in}}%
\pgfpathlineto{\pgfqpoint{3.760028in}{1.990690in}}%
\pgfpathlineto{\pgfqpoint{3.760949in}{1.986815in}}%
\pgfpathlineto{\pgfqpoint{3.764069in}{1.973567in}}%
\pgfpathlineto{\pgfqpoint{3.767936in}{1.960319in}}%
\pgfpathlineto{\pgfqpoint{3.769976in}{1.947071in}}%
\pgfpathlineto{\pgfqpoint{3.776591in}{1.933823in}}%
\pgfpathlineto{\pgfqpoint{3.778545in}{1.920575in}}%
\pgfpathlineto{\pgfqpoint{3.785245in}{1.907327in}}%
\pgfpathlineto{\pgfqpoint{3.785311in}{1.905065in}}%
\pgfpathlineto{\pgfqpoint{3.785629in}{1.894079in}}%
\pgfpathlineto{\pgfqpoint{3.792253in}{1.880831in}}%
\pgfpathlineto{\pgfqpoint{3.792067in}{1.867583in}}%
\pgfpathlineto{\pgfqpoint{3.800825in}{1.854335in}}%
\pgfpathlineto{\pgfqpoint{3.800630in}{1.841087in}}%
\pgfpathlineto{\pgfqpoint{3.808603in}{1.827839in}}%
\pgfpathlineto{\pgfqpoint{3.807896in}{1.814591in}}%
\pgfpathlineto{\pgfqpoint{3.810594in}{1.810194in}}%
\pgfpathlineto{\pgfqpoint{3.815368in}{1.801343in}}%
\pgfpathlineto{\pgfqpoint{3.814738in}{1.788095in}}%
\pgfpathlineto{\pgfqpoint{3.821597in}{1.774847in}}%
\pgfpathlineto{\pgfqpoint{3.822678in}{1.761599in}}%
\pgfpathlineto{\pgfqpoint{3.829203in}{1.748351in}}%
\pgfpathlineto{\pgfqpoint{3.834012in}{1.735103in}}%
\pgfpathlineto{\pgfqpoint{3.835878in}{1.729422in}}%
\pgfpathlineto{\pgfqpoint{3.838041in}{1.721855in}}%
\pgfpathlineto{\pgfqpoint{3.840387in}{1.708607in}}%
\pgfpathlineto{\pgfqpoint{3.839836in}{1.695359in}}%
\pgfpathlineto{\pgfqpoint{3.845705in}{1.682111in}}%
\pgfpathlineto{\pgfqpoint{3.842713in}{1.668863in}}%
\pgfpathlineto{\pgfqpoint{3.854080in}{1.655615in}}%
\pgfpathlineto{\pgfqpoint{3.845151in}{1.642367in}}%
\pgfpathlineto{\pgfqpoint{3.857937in}{1.629119in}}%
\pgfpathlineto{\pgfqpoint{3.835878in}{1.616732in}}%
\pgfpathlineto{\pgfqpoint{3.810594in}{1.619850in}}%
\pgfpathlineto{\pgfqpoint{3.790775in}{1.629119in}}%
\pgfpathlineto{\pgfqpoint{3.790891in}{1.642367in}}%
\pgfpathlineto{\pgfqpoint{3.785311in}{1.644307in}}%
\pgfpathlineto{\pgfqpoint{3.769072in}{1.655615in}}%
\pgfpathlineto{\pgfqpoint{3.768613in}{1.668863in}}%
\pgfpathlineto{\pgfqpoint{3.760028in}{1.673572in}}%
\pgfpathlineto{\pgfqpoint{3.750254in}{1.682111in}}%
\pgfpathlineto{\pgfqpoint{3.748906in}{1.695359in}}%
\pgfpathlineto{\pgfqpoint{3.734745in}{1.705034in}}%
\pgfpathlineto{\pgfqpoint{3.730903in}{1.708607in}}%
\pgfpathlineto{\pgfqpoint{3.726097in}{1.721855in}}%
\pgfpathlineto{\pgfqpoint{3.713273in}{1.735103in}}%
\pgfpathlineto{\pgfqpoint{3.709462in}{1.746475in}}%
\pgfpathlineto{\pgfqpoint{3.708976in}{1.748351in}}%
\pgfpathlineto{\pgfqpoint{3.700555in}{1.761599in}}%
\pgfpathlineto{\pgfqpoint{3.696133in}{1.774847in}}%
\pgfpathlineto{\pgfqpoint{3.686463in}{1.788095in}}%
\pgfpathlineto{\pgfqpoint{3.684179in}{1.791794in}}%
\pgfpathlineto{\pgfqpoint{3.679644in}{1.801343in}}%
\pgfpathlineto{\pgfqpoint{3.673342in}{1.814591in}}%
\pgfpathlineto{\pgfqpoint{3.665919in}{1.827839in}}%
\pgfpathlineto{\pgfqpoint{3.660216in}{1.841087in}}%
\pgfpathlineto{\pgfqpoint{3.658895in}{1.843076in}}%
\pgfpathlineto{\pgfqpoint{3.652882in}{1.854335in}}%
\pgfpathlineto{\pgfqpoint{3.648539in}{1.867583in}}%
\pgfpathlineto{\pgfqpoint{3.640235in}{1.880831in}}%
\pgfpathlineto{\pgfqpoint{3.636731in}{1.894079in}}%
\pgfpathlineto{\pgfqpoint{3.633612in}{1.897882in}}%
\pgfpathlineto{\pgfqpoint{3.627057in}{1.907327in}}%
\pgfpathlineto{\pgfqpoint{3.624044in}{1.920575in}}%
\pgfpathlineto{\pgfqpoint{3.615984in}{1.933823in}}%
\pgfpathlineto{\pgfqpoint{3.614449in}{1.947071in}}%
\pgfpathlineto{\pgfqpoint{3.608329in}{1.956716in}}%
\pgfpathlineto{\pgfqpoint{3.606269in}{1.960319in}}%
\pgfpathlineto{\pgfqpoint{3.604227in}{1.973567in}}%
\pgfpathlineto{\pgfqpoint{3.595219in}{1.986815in}}%
\pgfpathlineto{\pgfqpoint{3.592705in}{2.000063in}}%
\pgfpathlineto{\pgfqpoint{3.584832in}{2.013311in}}%
\pgfpathlineto{\pgfqpoint{3.583046in}{2.022457in}}%
\pgfpathlineto{\pgfqpoint{3.582259in}{2.026559in}}%
\pgfpathlineto{\pgfqpoint{3.575026in}{2.039807in}}%
\pgfpathlineto{\pgfqpoint{3.572331in}{2.053055in}}%
\pgfpathlineto{\pgfqpoint{3.565629in}{2.066303in}}%
\pgfpathlineto{\pgfqpoint{3.562517in}{2.079551in}}%
\pgfpathlineto{\pgfqpoint{3.557763in}{2.090085in}}%
\pgfpathlineto{\pgfqpoint{3.556492in}{2.092799in}}%
\pgfpathlineto{\pgfqpoint{3.552121in}{2.106047in}}%
\pgfpathlineto{\pgfqpoint{3.547182in}{2.119295in}}%
\pgfpathlineto{\pgfqpoint{3.543355in}{2.132543in}}%
\pgfpathlineto{\pgfqpoint{3.532480in}{2.132543in}}%
\pgfpathlineto{\pgfqpoint{3.528569in}{2.132543in}}%
\pgfpathlineto{\pgfqpoint{3.532480in}{2.119436in}}%
\pgfpathlineto{\pgfqpoint{3.532516in}{2.119295in}}%
\pgfpathlineto{\pgfqpoint{3.537218in}{2.106047in}}%
\pgfpathlineto{\pgfqpoint{3.540938in}{2.092799in}}%
\pgfpathlineto{\pgfqpoint{3.546783in}{2.079551in}}%
\pgfpathlineto{\pgfqpoint{3.549707in}{2.066303in}}%
\pgfpathlineto{\pgfqpoint{3.556750in}{2.053055in}}%
\pgfpathlineto{\pgfqpoint{3.557763in}{2.048326in}}%
\pgfpathlineto{\pgfqpoint{3.559368in}{2.039807in}}%
\pgfpathlineto{\pgfqpoint{3.565856in}{2.026559in}}%
\pgfpathlineto{\pgfqpoint{3.568130in}{2.013311in}}%
\pgfpathlineto{\pgfqpoint{3.575474in}{2.000063in}}%
\pgfpathlineto{\pgfqpoint{3.578083in}{1.986815in}}%
\pgfpathlineto{\pgfqpoint{3.583046in}{1.979301in}}%
\pgfpathlineto{\pgfqpoint{3.586725in}{1.973567in}}%
\pgfpathlineto{\pgfqpoint{3.588884in}{1.960319in}}%
\pgfpathlineto{\pgfqpoint{3.595718in}{1.947071in}}%
\pgfpathlineto{\pgfqpoint{3.597533in}{1.933823in}}%
\pgfpathlineto{\pgfqpoint{3.604865in}{1.920575in}}%
\pgfpathlineto{\pgfqpoint{3.608329in}{1.907415in}}%
\pgfpathlineto{\pgfqpoint{3.608351in}{1.907327in}}%
\pgfpathlineto{\pgfqpoint{3.616693in}{1.894079in}}%
\pgfpathlineto{\pgfqpoint{3.620335in}{1.880831in}}%
\pgfpathlineto{\pgfqpoint{3.627235in}{1.867583in}}%
\pgfpathlineto{\pgfqpoint{3.632070in}{1.854335in}}%
\pgfpathlineto{\pgfqpoint{3.633612in}{1.851025in}}%
\pgfpathlineto{\pgfqpoint{3.638386in}{1.841087in}}%
\pgfpathlineto{\pgfqpoint{3.643512in}{1.827839in}}%
\pgfpathlineto{\pgfqpoint{3.649471in}{1.814591in}}%
\pgfpathlineto{\pgfqpoint{3.656018in}{1.801343in}}%
\pgfpathlineto{\pgfqpoint{3.658895in}{1.794372in}}%
\pgfpathlineto{\pgfqpoint{3.661678in}{1.788095in}}%
\pgfpathlineto{\pgfqpoint{3.670192in}{1.774847in}}%
\pgfpathlineto{\pgfqpoint{3.673851in}{1.761599in}}%
\pgfpathlineto{\pgfqpoint{3.681720in}{1.748351in}}%
\pgfpathlineto{\pgfqpoint{3.684179in}{1.739146in}}%
\pgfpathlineto{\pgfqpoint{3.685342in}{1.735103in}}%
\pgfpathlineto{\pgfqpoint{3.696084in}{1.721855in}}%
\pgfpathlineto{\pgfqpoint{3.700405in}{1.708607in}}%
\pgfpathlineto{\pgfqpoint{3.709462in}{1.698933in}}%
\pgfpathlineto{\pgfqpoint{3.713362in}{1.695359in}}%
\pgfpathlineto{\pgfqpoint{3.715807in}{1.682111in}}%
\pgfpathlineto{\pgfqpoint{3.727778in}{1.668863in}}%
\pgfpathlineto{\pgfqpoint{3.730383in}{1.655615in}}%
\pgfpathlineto{\pgfqpoint{3.734745in}{1.650622in}}%
\pgfpathlineto{\pgfqpoint{3.743382in}{1.642367in}}%
\pgfpathlineto{\pgfqpoint{3.746330in}{1.629119in}}%
\pgfpathlineto{\pgfqpoint{3.760028in}{1.615969in}}%
\pgfpathlineto{\pgfqpoint{3.760159in}{1.615871in}}%
\pgfpathlineto{\pgfqpoint{3.764183in}{1.602623in}}%
\pgfpathlineto{\pgfqpoint{3.779251in}{1.589375in}}%
\pgfpathlineto{\pgfqpoint{3.783995in}{1.576127in}}%
\pgfpathlineto{\pgfqpoint{3.785311in}{1.574495in}}%
\pgfpathlineto{\pgfqpoint{3.797359in}{1.562879in}}%
\pgfpathlineto{\pgfqpoint{3.802974in}{1.549631in}}%
\pgfpathlineto{\pgfqpoint{3.810594in}{1.541802in}}%
\pgfpathlineto{\pgfqpoint{3.817730in}{1.536383in}}%
\pgfpathlineto{\pgfqpoint{3.828738in}{1.523135in}}%
\pgfpathlineto{\pgfqpoint{3.835878in}{1.516513in}}%
\pgfpathlineto{\pgfqpoint{3.846609in}{1.509887in}}%
\pgfpathlineto{\pgfqpoint{3.859392in}{1.496639in}}%
\pgfpathlineto{\pgfqpoint{3.861161in}{1.493325in}}%
\pgfpathlineto{\pgfqpoint{3.869743in}{1.483391in}}%
\pgfpathlineto{\pgfqpoint{3.885585in}{1.470143in}}%
\pgfpathlineto{\pgfqpoint{3.886444in}{1.468634in}}%
\pgfpathlineto{\pgfqpoint{3.897915in}{1.456895in}}%
\pgfpathlineto{\pgfqpoint{3.911727in}{1.449902in}}%
\pgfpathclose%
\pgfpathmoveto{\pgfqpoint{3.820646in}{1.602623in}}%
\pgfpathlineto{\pgfqpoint{3.835878in}{1.613956in}}%
\pgfpathlineto{\pgfqpoint{3.858119in}{1.602623in}}%
\pgfpathlineto{\pgfqpoint{3.835878in}{1.598428in}}%
\pgfpathlineto{\pgfqpoint{3.820646in}{1.602623in}}%
\pgfpathclose%
\pgfpathmoveto{\pgfqpoint{2.723418in}{2.094505in}}%
\pgfpathlineto{\pgfqpoint{2.730690in}{2.106047in}}%
\pgfpathlineto{\pgfqpoint{2.737215in}{2.119295in}}%
\pgfpathlineto{\pgfqpoint{2.748701in}{2.126917in}}%
\pgfpathlineto{\pgfqpoint{2.752312in}{2.132543in}}%
\pgfpathlineto{\pgfqpoint{2.748701in}{2.132543in}}%
\pgfpathlineto{\pgfqpoint{2.723418in}{2.132543in}}%
\pgfpathlineto{\pgfqpoint{2.714117in}{2.132543in}}%
\pgfpathlineto{\pgfqpoint{2.709939in}{2.119295in}}%
\pgfpathlineto{\pgfqpoint{2.713936in}{2.106047in}}%
\pgfpathlineto{\pgfqpoint{2.723418in}{2.094505in}}%
\pgfpathclose%
\pgfpathmoveto{\pgfqpoint{3.937010in}{2.131913in}}%
\pgfpathlineto{\pgfqpoint{3.938211in}{2.132543in}}%
\pgfpathlineto{\pgfqpoint{3.937010in}{2.132543in}}%
\pgfpathlineto{\pgfqpoint{3.936369in}{2.132543in}}%
\pgfpathlineto{\pgfqpoint{3.937010in}{2.131913in}}%
\pgfpathclose%
\pgfusepath{fill}%
\end{pgfscope}%
\begin{pgfscope}%
\pgfpathrectangle{\pgfqpoint{2.470586in}{0.476543in}}{\pgfqpoint{1.820388in}{1.656000in}}%
\pgfusepath{clip}%
\pgfsetbuttcap%
\pgfsetroundjoin%
\definecolor{currentfill}{rgb}{0.433218,0.652595,0.806228}%
\pgfsetfillcolor{currentfill}%
\pgfsetlinewidth{0.000000pt}%
\definecolor{currentstroke}{rgb}{0.000000,0.000000,0.000000}%
\pgfsetstrokecolor{currentstroke}%
\pgfsetdash{}{0pt}%
\pgfpathmoveto{\pgfqpoint{3.052099in}{0.726803in}}%
\pgfpathlineto{\pgfqpoint{3.077382in}{0.725314in}}%
\pgfpathlineto{\pgfqpoint{3.102666in}{0.723978in}}%
\pgfpathlineto{\pgfqpoint{3.127949in}{0.722790in}}%
\pgfpathlineto{\pgfqpoint{3.153232in}{0.721746in}}%
\pgfpathlineto{\pgfqpoint{3.178515in}{0.720844in}}%
\pgfpathlineto{\pgfqpoint{3.203798in}{0.720083in}}%
\pgfpathlineto{\pgfqpoint{3.229082in}{0.719464in}}%
\pgfpathlineto{\pgfqpoint{3.254365in}{0.718990in}}%
\pgfpathlineto{\pgfqpoint{3.279648in}{0.718665in}}%
\pgfpathlineto{\pgfqpoint{3.304931in}{0.718497in}}%
\pgfpathlineto{\pgfqpoint{3.330214in}{0.718495in}}%
\pgfpathlineto{\pgfqpoint{3.355497in}{0.718670in}}%
\pgfpathlineto{\pgfqpoint{3.380781in}{0.719038in}}%
\pgfpathlineto{\pgfqpoint{3.406064in}{0.719619in}}%
\pgfpathlineto{\pgfqpoint{3.431347in}{0.720436in}}%
\pgfpathlineto{\pgfqpoint{3.456630in}{0.721521in}}%
\pgfpathlineto{\pgfqpoint{3.481913in}{0.722911in}}%
\pgfpathlineto{\pgfqpoint{3.507196in}{0.724654in}}%
\pgfpathlineto{\pgfqpoint{3.532480in}{0.726810in}}%
\pgfpathlineto{\pgfqpoint{3.546652in}{0.728255in}}%
\pgfpathlineto{\pgfqpoint{3.557763in}{0.729089in}}%
\pgfpathlineto{\pgfqpoint{3.583046in}{0.731166in}}%
\pgfpathlineto{\pgfqpoint{3.608329in}{0.733447in}}%
\pgfpathlineto{\pgfqpoint{3.633612in}{0.735953in}}%
\pgfpathlineto{\pgfqpoint{3.658895in}{0.738710in}}%
\pgfpathlineto{\pgfqpoint{3.682180in}{0.741503in}}%
\pgfpathlineto{\pgfqpoint{3.684179in}{0.742080in}}%
\pgfpathlineto{\pgfqpoint{3.709462in}{0.751130in}}%
\pgfpathlineto{\pgfqpoint{3.717379in}{0.754751in}}%
\pgfpathlineto{\pgfqpoint{3.734745in}{0.756877in}}%
\pgfpathlineto{\pgfqpoint{3.760028in}{0.760301in}}%
\pgfpathlineto{\pgfqpoint{3.785311in}{0.764149in}}%
\pgfpathlineto{\pgfqpoint{3.807821in}{0.767999in}}%
\pgfpathlineto{\pgfqpoint{3.806579in}{0.781247in}}%
\pgfpathlineto{\pgfqpoint{3.810594in}{0.782162in}}%
\pgfpathlineto{\pgfqpoint{3.835878in}{0.788631in}}%
\pgfpathlineto{\pgfqpoint{3.855198in}{0.794495in}}%
\pgfpathlineto{\pgfqpoint{3.846079in}{0.807743in}}%
\pgfpathlineto{\pgfqpoint{3.861161in}{0.815246in}}%
\pgfpathlineto{\pgfqpoint{3.871092in}{0.820991in}}%
\pgfpathlineto{\pgfqpoint{3.861161in}{0.829060in}}%
\pgfpathlineto{\pgfqpoint{3.855182in}{0.834239in}}%
\pgfpathlineto{\pgfqpoint{3.861161in}{0.840291in}}%
\pgfpathlineto{\pgfqpoint{3.867413in}{0.847487in}}%
\pgfpathlineto{\pgfqpoint{3.861161in}{0.857047in}}%
\pgfpathlineto{\pgfqpoint{3.858840in}{0.860735in}}%
\pgfpathlineto{\pgfqpoint{3.861161in}{0.866317in}}%
\pgfpathlineto{\pgfqpoint{3.864046in}{0.873983in}}%
\pgfpathlineto{\pgfqpoint{3.863725in}{0.887231in}}%
\pgfpathlineto{\pgfqpoint{3.861161in}{0.894109in}}%
\pgfpathlineto{\pgfqpoint{3.859055in}{0.900479in}}%
\pgfpathlineto{\pgfqpoint{3.854223in}{0.913727in}}%
\pgfpathlineto{\pgfqpoint{3.843372in}{0.926975in}}%
\pgfpathlineto{\pgfqpoint{3.844072in}{0.940223in}}%
\pgfpathlineto{\pgfqpoint{3.835878in}{0.948094in}}%
\pgfpathlineto{\pgfqpoint{3.830796in}{0.953471in}}%
\pgfpathlineto{\pgfqpoint{3.832860in}{0.966719in}}%
\pgfpathlineto{\pgfqpoint{3.818328in}{0.979967in}}%
\pgfpathlineto{\pgfqpoint{3.818864in}{0.993215in}}%
\pgfpathlineto{\pgfqpoint{3.810594in}{1.000168in}}%
\pgfpathlineto{\pgfqpoint{3.803656in}{1.006463in}}%
\pgfpathlineto{\pgfqpoint{3.803808in}{1.019711in}}%
\pgfpathlineto{\pgfqpoint{3.787894in}{1.032959in}}%
\pgfpathlineto{\pgfqpoint{3.788175in}{1.046207in}}%
\pgfpathlineto{\pgfqpoint{3.785311in}{1.049546in}}%
\pgfpathlineto{\pgfqpoint{3.777360in}{1.059455in}}%
\pgfpathlineto{\pgfqpoint{3.776348in}{1.072703in}}%
\pgfpathlineto{\pgfqpoint{3.765946in}{1.085951in}}%
\pgfpathlineto{\pgfqpoint{3.760812in}{1.099199in}}%
\pgfpathlineto{\pgfqpoint{3.760028in}{1.100126in}}%
\pgfpathlineto{\pgfqpoint{3.750440in}{1.112447in}}%
\pgfpathlineto{\pgfqpoint{3.744491in}{1.125695in}}%
\pgfpathlineto{\pgfqpoint{3.738304in}{1.138943in}}%
\pgfpathlineto{\pgfqpoint{3.734745in}{1.145061in}}%
\pgfpathlineto{\pgfqpoint{3.731117in}{1.152191in}}%
\pgfpathlineto{\pgfqpoint{3.726420in}{1.165439in}}%
\pgfpathlineto{\pgfqpoint{3.718288in}{1.178687in}}%
\pgfpathlineto{\pgfqpoint{3.714294in}{1.191935in}}%
\pgfpathlineto{\pgfqpoint{3.709462in}{1.198257in}}%
\pgfpathlineto{\pgfqpoint{3.704835in}{1.205183in}}%
\pgfpathlineto{\pgfqpoint{3.702196in}{1.218431in}}%
\pgfpathlineto{\pgfqpoint{3.691963in}{1.231679in}}%
\pgfpathlineto{\pgfqpoint{3.690582in}{1.244927in}}%
\pgfpathlineto{\pgfqpoint{3.684179in}{1.253841in}}%
\pgfpathlineto{\pgfqpoint{3.681423in}{1.258175in}}%
\pgfpathlineto{\pgfqpoint{3.679526in}{1.271423in}}%
\pgfpathlineto{\pgfqpoint{3.669396in}{1.284671in}}%
\pgfpathlineto{\pgfqpoint{3.666564in}{1.297919in}}%
\pgfpathlineto{\pgfqpoint{3.658895in}{1.307907in}}%
\pgfpathlineto{\pgfqpoint{3.656676in}{1.311167in}}%
\pgfpathlineto{\pgfqpoint{3.654828in}{1.324415in}}%
\pgfpathlineto{\pgfqpoint{3.647965in}{1.337663in}}%
\pgfpathlineto{\pgfqpoint{3.645019in}{1.350911in}}%
\pgfpathlineto{\pgfqpoint{3.637714in}{1.364159in}}%
\pgfpathlineto{\pgfqpoint{3.633683in}{1.377407in}}%
\pgfpathlineto{\pgfqpoint{3.633612in}{1.377544in}}%
\pgfpathlineto{\pgfqpoint{3.627842in}{1.390655in}}%
\pgfpathlineto{\pgfqpoint{3.623058in}{1.403903in}}%
\pgfpathlineto{\pgfqpoint{3.618089in}{1.417151in}}%
\pgfpathlineto{\pgfqpoint{3.612180in}{1.430399in}}%
\pgfpathlineto{\pgfqpoint{3.608329in}{1.443523in}}%
\pgfpathlineto{\pgfqpoint{3.608299in}{1.443647in}}%
\pgfpathlineto{\pgfqpoint{3.602888in}{1.456895in}}%
\pgfpathlineto{\pgfqpoint{3.599492in}{1.470143in}}%
\pgfpathlineto{\pgfqpoint{3.592138in}{1.483391in}}%
\pgfpathlineto{\pgfqpoint{3.589498in}{1.496639in}}%
\pgfpathlineto{\pgfqpoint{3.583046in}{1.508074in}}%
\pgfpathlineto{\pgfqpoint{3.582188in}{1.509887in}}%
\pgfpathlineto{\pgfqpoint{3.580922in}{1.523135in}}%
\pgfpathlineto{\pgfqpoint{3.574828in}{1.536383in}}%
\pgfpathlineto{\pgfqpoint{3.572782in}{1.549631in}}%
\pgfpathlineto{\pgfqpoint{3.565282in}{1.562879in}}%
\pgfpathlineto{\pgfqpoint{3.563166in}{1.576127in}}%
\pgfpathlineto{\pgfqpoint{3.557763in}{1.585223in}}%
\pgfpathlineto{\pgfqpoint{3.555693in}{1.589375in}}%
\pgfpathlineto{\pgfqpoint{3.553523in}{1.602623in}}%
\pgfpathlineto{\pgfqpoint{3.547584in}{1.615871in}}%
\pgfpathlineto{\pgfqpoint{3.545265in}{1.629119in}}%
\pgfpathlineto{\pgfqpoint{3.539655in}{1.642367in}}%
\pgfpathlineto{\pgfqpoint{3.536523in}{1.655615in}}%
\pgfpathlineto{\pgfqpoint{3.532480in}{1.664241in}}%
\pgfpathlineto{\pgfqpoint{3.530678in}{1.668863in}}%
\pgfpathlineto{\pgfqpoint{3.527019in}{1.682111in}}%
\pgfpathlineto{\pgfqpoint{3.523034in}{1.695359in}}%
\pgfpathlineto{\pgfqpoint{3.519191in}{1.708607in}}%
\pgfpathlineto{\pgfqpoint{3.515683in}{1.721855in}}%
\pgfpathlineto{\pgfqpoint{3.510524in}{1.735103in}}%
\pgfpathlineto{\pgfqpoint{3.507196in}{1.745492in}}%
\pgfpathlineto{\pgfqpoint{3.506440in}{1.748351in}}%
\pgfpathlineto{\pgfqpoint{3.501021in}{1.761599in}}%
\pgfpathlineto{\pgfqpoint{3.498426in}{1.774847in}}%
\pgfpathlineto{\pgfqpoint{3.492458in}{1.788095in}}%
\pgfpathlineto{\pgfqpoint{3.490201in}{1.801343in}}%
\pgfpathlineto{\pgfqpoint{3.483753in}{1.814591in}}%
\pgfpathlineto{\pgfqpoint{3.481913in}{1.826158in}}%
\pgfpathlineto{\pgfqpoint{3.481691in}{1.827839in}}%
\pgfpathlineto{\pgfqpoint{3.475745in}{1.841087in}}%
\pgfpathlineto{\pgfqpoint{3.473748in}{1.854335in}}%
\pgfpathlineto{\pgfqpoint{3.466662in}{1.867583in}}%
\pgfpathlineto{\pgfqpoint{3.464403in}{1.880831in}}%
\pgfpathlineto{\pgfqpoint{3.457622in}{1.894079in}}%
\pgfpathlineto{\pgfqpoint{3.456630in}{1.899670in}}%
\pgfpathlineto{\pgfqpoint{3.455455in}{1.907327in}}%
\pgfpathlineto{\pgfqpoint{3.449130in}{1.920575in}}%
\pgfpathlineto{\pgfqpoint{3.446086in}{1.933823in}}%
\pgfpathlineto{\pgfqpoint{3.438639in}{1.947071in}}%
\pgfpathlineto{\pgfqpoint{3.435109in}{1.960319in}}%
\pgfpathlineto{\pgfqpoint{3.431347in}{1.968318in}}%
\pgfpathlineto{\pgfqpoint{3.428825in}{1.973567in}}%
\pgfpathlineto{\pgfqpoint{3.424864in}{1.986815in}}%
\pgfpathlineto{\pgfqpoint{3.418908in}{2.000063in}}%
\pgfpathlineto{\pgfqpoint{3.414121in}{2.013311in}}%
\pgfpathlineto{\pgfqpoint{3.407703in}{2.026559in}}%
\pgfpathlineto{\pgfqpoint{3.406064in}{2.030026in}}%
\pgfpathlineto{\pgfqpoint{3.400868in}{2.039807in}}%
\pgfpathlineto{\pgfqpoint{3.393391in}{2.053055in}}%
\pgfpathlineto{\pgfqpoint{3.384229in}{2.066303in}}%
\pgfpathlineto{\pgfqpoint{3.380781in}{2.073133in}}%
\pgfpathlineto{\pgfqpoint{3.375443in}{2.079551in}}%
\pgfpathlineto{\pgfqpoint{3.361287in}{2.092799in}}%
\pgfpathlineto{\pgfqpoint{3.355497in}{2.100088in}}%
\pgfpathlineto{\pgfqpoint{3.330214in}{2.105234in}}%
\pgfpathlineto{\pgfqpoint{3.311706in}{2.092799in}}%
\pgfpathlineto{\pgfqpoint{3.304931in}{2.087673in}}%
\pgfpathlineto{\pgfqpoint{3.300126in}{2.079551in}}%
\pgfpathlineto{\pgfqpoint{3.288447in}{2.066303in}}%
\pgfpathlineto{\pgfqpoint{3.282822in}{2.053055in}}%
\pgfpathlineto{\pgfqpoint{3.279648in}{2.049740in}}%
\pgfpathlineto{\pgfqpoint{3.272337in}{2.039807in}}%
\pgfpathlineto{\pgfqpoint{3.269287in}{2.026559in}}%
\pgfpathlineto{\pgfqpoint{3.259654in}{2.013311in}}%
\pgfpathlineto{\pgfqpoint{3.256804in}{2.000063in}}%
\pgfpathlineto{\pgfqpoint{3.254365in}{1.996236in}}%
\pgfpathlineto{\pgfqpoint{3.248869in}{1.986815in}}%
\pgfpathlineto{\pgfqpoint{3.246353in}{1.973567in}}%
\pgfpathlineto{\pgfqpoint{3.238737in}{1.960319in}}%
\pgfpathlineto{\pgfqpoint{3.234881in}{1.947071in}}%
\pgfpathlineto{\pgfqpoint{3.229082in}{1.934976in}}%
\pgfpathlineto{\pgfqpoint{3.228508in}{1.933823in}}%
\pgfpathlineto{\pgfqpoint{3.224161in}{1.920575in}}%
\pgfpathlineto{\pgfqpoint{3.219762in}{1.907327in}}%
\pgfpathlineto{\pgfqpoint{3.215481in}{1.894079in}}%
\pgfpathlineto{\pgfqpoint{3.211684in}{1.880831in}}%
\pgfpathlineto{\pgfqpoint{3.206504in}{1.867583in}}%
\pgfpathlineto{\pgfqpoint{3.203815in}{1.854335in}}%
\pgfpathlineto{\pgfqpoint{3.203798in}{1.854300in}}%
\pgfpathlineto{\pgfqpoint{3.196770in}{1.841087in}}%
\pgfpathlineto{\pgfqpoint{3.194092in}{1.827839in}}%
\pgfpathlineto{\pgfqpoint{3.187465in}{1.814591in}}%
\pgfpathlineto{\pgfqpoint{3.185773in}{1.801343in}}%
\pgfpathlineto{\pgfqpoint{3.179951in}{1.788095in}}%
\pgfpathlineto{\pgfqpoint{3.178515in}{1.776304in}}%
\pgfpathlineto{\pgfqpoint{3.178295in}{1.774847in}}%
\pgfpathlineto{\pgfqpoint{3.170598in}{1.761599in}}%
\pgfpathlineto{\pgfqpoint{3.168123in}{1.748351in}}%
\pgfpathlineto{\pgfqpoint{3.162477in}{1.735103in}}%
\pgfpathlineto{\pgfqpoint{3.160273in}{1.721855in}}%
\pgfpathlineto{\pgfqpoint{3.155523in}{1.708607in}}%
\pgfpathlineto{\pgfqpoint{3.153232in}{1.698627in}}%
\pgfpathlineto{\pgfqpoint{3.152284in}{1.695359in}}%
\pgfpathlineto{\pgfqpoint{3.146683in}{1.682111in}}%
\pgfpathlineto{\pgfqpoint{3.142039in}{1.668863in}}%
\pgfpathlineto{\pgfqpoint{3.138054in}{1.655615in}}%
\pgfpathlineto{\pgfqpoint{3.132688in}{1.642367in}}%
\pgfpathlineto{\pgfqpoint{3.129830in}{1.629119in}}%
\pgfpathlineto{\pgfqpoint{3.127949in}{1.624149in}}%
\pgfpathlineto{\pgfqpoint{3.124113in}{1.615871in}}%
\pgfpathlineto{\pgfqpoint{3.122161in}{1.602623in}}%
\pgfpathlineto{\pgfqpoint{3.115500in}{1.589375in}}%
\pgfpathlineto{\pgfqpoint{3.113439in}{1.576127in}}%
\pgfpathlineto{\pgfqpoint{3.105915in}{1.562879in}}%
\pgfpathlineto{\pgfqpoint{3.104344in}{1.549631in}}%
\pgfpathlineto{\pgfqpoint{3.102666in}{1.546136in}}%
\pgfpathlineto{\pgfqpoint{3.097103in}{1.536383in}}%
\pgfpathlineto{\pgfqpoint{3.095236in}{1.523135in}}%
\pgfpathlineto{\pgfqpoint{3.087746in}{1.509887in}}%
\pgfpathlineto{\pgfqpoint{3.085363in}{1.496639in}}%
\pgfpathlineto{\pgfqpoint{3.078624in}{1.483391in}}%
\pgfpathlineto{\pgfqpoint{3.077382in}{1.477514in}}%
\pgfpathlineto{\pgfqpoint{3.075417in}{1.470143in}}%
\pgfpathlineto{\pgfqpoint{3.068912in}{1.456895in}}%
\pgfpathlineto{\pgfqpoint{3.064775in}{1.443647in}}%
\pgfpathlineto{\pgfqpoint{3.060047in}{1.430399in}}%
\pgfpathlineto{\pgfqpoint{3.055868in}{1.417151in}}%
\pgfpathlineto{\pgfqpoint{3.052099in}{1.404350in}}%
\pgfpathlineto{\pgfqpoint{3.051943in}{1.403903in}}%
\pgfpathlineto{\pgfqpoint{3.044393in}{1.390655in}}%
\pgfpathlineto{\pgfqpoint{3.040194in}{1.377407in}}%
\pgfpathlineto{\pgfqpoint{3.031942in}{1.364159in}}%
\pgfpathlineto{\pgfqpoint{3.030076in}{1.350911in}}%
\pgfpathlineto{\pgfqpoint{3.026816in}{1.345088in}}%
\pgfpathlineto{\pgfqpoint{3.021966in}{1.337663in}}%
\pgfpathlineto{\pgfqpoint{3.019995in}{1.324415in}}%
\pgfpathlineto{\pgfqpoint{3.010147in}{1.311167in}}%
\pgfpathlineto{\pgfqpoint{3.008683in}{1.297919in}}%
\pgfpathlineto{\pgfqpoint{3.001533in}{1.288051in}}%
\pgfpathlineto{\pgfqpoint{2.998741in}{1.284671in}}%
\pgfpathlineto{\pgfqpoint{2.997345in}{1.271423in}}%
\pgfpathlineto{\pgfqpoint{2.987208in}{1.258175in}}%
\pgfpathlineto{\pgfqpoint{2.985298in}{1.244927in}}%
\pgfpathlineto{\pgfqpoint{2.976449in}{1.231679in}}%
\pgfpathlineto{\pgfqpoint{2.976250in}{1.230586in}}%
\pgfpathlineto{\pgfqpoint{2.973605in}{1.218431in}}%
\pgfpathlineto{\pgfqpoint{2.963241in}{1.205183in}}%
\pgfpathlineto{\pgfqpoint{2.957510in}{1.191935in}}%
\pgfpathlineto{\pgfqpoint{2.950967in}{1.180981in}}%
\pgfpathlineto{\pgfqpoint{2.949468in}{1.178687in}}%
\pgfpathlineto{\pgfqpoint{2.943217in}{1.165439in}}%
\pgfpathlineto{\pgfqpoint{2.938857in}{1.152191in}}%
\pgfpathlineto{\pgfqpoint{2.931933in}{1.138943in}}%
\pgfpathlineto{\pgfqpoint{2.927123in}{1.125695in}}%
\pgfpathlineto{\pgfqpoint{2.925683in}{1.123781in}}%
\pgfpathlineto{\pgfqpoint{2.916014in}{1.112447in}}%
\pgfpathlineto{\pgfqpoint{2.913153in}{1.099199in}}%
\pgfpathlineto{\pgfqpoint{2.900422in}{1.085951in}}%
\pgfpathlineto{\pgfqpoint{2.900400in}{1.085673in}}%
\pgfpathlineto{\pgfqpoint{2.899246in}{1.072703in}}%
\pgfpathlineto{\pgfqpoint{2.884616in}{1.059455in}}%
\pgfpathlineto{\pgfqpoint{2.885261in}{1.046207in}}%
\pgfpathlineto{\pgfqpoint{2.875117in}{1.037128in}}%
\pgfpathlineto{\pgfqpoint{2.870103in}{1.032959in}}%
\pgfpathlineto{\pgfqpoint{2.871430in}{1.019711in}}%
\pgfpathlineto{\pgfqpoint{2.852569in}{1.006463in}}%
\pgfpathlineto{\pgfqpoint{2.853865in}{0.993215in}}%
\pgfpathlineto{\pgfqpoint{2.849834in}{0.989758in}}%
\pgfpathlineto{\pgfqpoint{2.837876in}{0.979967in}}%
\pgfpathlineto{\pgfqpoint{2.842321in}{0.966719in}}%
\pgfpathlineto{\pgfqpoint{2.828993in}{0.953471in}}%
\pgfpathlineto{\pgfqpoint{2.830731in}{0.940223in}}%
\pgfpathlineto{\pgfqpoint{2.824551in}{0.935428in}}%
\pgfpathlineto{\pgfqpoint{2.813221in}{0.926975in}}%
\pgfpathlineto{\pgfqpoint{2.812159in}{0.913727in}}%
\pgfpathlineto{\pgfqpoint{2.799893in}{0.900479in}}%
\pgfpathlineto{\pgfqpoint{2.799268in}{0.890018in}}%
\pgfpathlineto{\pgfqpoint{2.799073in}{0.887231in}}%
\pgfpathlineto{\pgfqpoint{2.792459in}{0.873983in}}%
\pgfpathlineto{\pgfqpoint{2.790199in}{0.860735in}}%
\pgfpathlineto{\pgfqpoint{2.791084in}{0.847487in}}%
\pgfpathlineto{\pgfqpoint{2.788426in}{0.834239in}}%
\pgfpathlineto{\pgfqpoint{2.796345in}{0.820991in}}%
\pgfpathlineto{\pgfqpoint{2.790274in}{0.807743in}}%
\pgfpathlineto{\pgfqpoint{2.799268in}{0.802723in}}%
\pgfpathlineto{\pgfqpoint{2.814896in}{0.794495in}}%
\pgfpathlineto{\pgfqpoint{2.819497in}{0.781247in}}%
\pgfpathlineto{\pgfqpoint{2.824551in}{0.779983in}}%
\pgfpathlineto{\pgfqpoint{2.849834in}{0.774065in}}%
\pgfpathlineto{\pgfqpoint{2.875117in}{0.768694in}}%
\pgfpathlineto{\pgfqpoint{2.878600in}{0.767999in}}%
\pgfpathlineto{\pgfqpoint{2.895319in}{0.754751in}}%
\pgfpathlineto{\pgfqpoint{2.900400in}{0.753858in}}%
\pgfpathlineto{\pgfqpoint{2.925683in}{0.749764in}}%
\pgfpathlineto{\pgfqpoint{2.950967in}{0.746101in}}%
\pgfpathlineto{\pgfqpoint{2.976250in}{0.742814in}}%
\pgfpathlineto{\pgfqpoint{2.987296in}{0.741503in}}%
\pgfpathlineto{\pgfqpoint{3.001533in}{0.735848in}}%
\pgfpathlineto{\pgfqpoint{3.026816in}{0.728883in}}%
\pgfpathlineto{\pgfqpoint{3.029846in}{0.728255in}}%
\pgfpathlineto{\pgfqpoint{3.052099in}{0.726803in}}%
\pgfpathclose%
\pgfpathmoveto{\pgfqpoint{3.230460in}{0.754751in}}%
\pgfpathlineto{\pgfqpoint{3.229082in}{0.754784in}}%
\pgfpathlineto{\pgfqpoint{3.203798in}{0.755654in}}%
\pgfpathlineto{\pgfqpoint{3.178515in}{0.756874in}}%
\pgfpathlineto{\pgfqpoint{3.153232in}{0.758515in}}%
\pgfpathlineto{\pgfqpoint{3.127949in}{0.760675in}}%
\pgfpathlineto{\pgfqpoint{3.102666in}{0.763484in}}%
\pgfpathlineto{\pgfqpoint{3.077382in}{0.767129in}}%
\pgfpathlineto{\pgfqpoint{3.072361in}{0.767999in}}%
\pgfpathlineto{\pgfqpoint{3.052099in}{0.770535in}}%
\pgfpathlineto{\pgfqpoint{3.026816in}{0.773974in}}%
\pgfpathlineto{\pgfqpoint{3.001533in}{0.777720in}}%
\pgfpathlineto{\pgfqpoint{2.979664in}{0.781247in}}%
\pgfpathlineto{\pgfqpoint{2.976250in}{0.782124in}}%
\pgfpathlineto{\pgfqpoint{2.950967in}{0.790262in}}%
\pgfpathlineto{\pgfqpoint{2.940855in}{0.794495in}}%
\pgfpathlineto{\pgfqpoint{2.925683in}{0.800305in}}%
\pgfpathlineto{\pgfqpoint{2.907675in}{0.807743in}}%
\pgfpathlineto{\pgfqpoint{2.900400in}{0.813725in}}%
\pgfpathlineto{\pgfqpoint{2.893358in}{0.820991in}}%
\pgfpathlineto{\pgfqpoint{2.880053in}{0.834239in}}%
\pgfpathlineto{\pgfqpoint{2.875117in}{0.839985in}}%
\pgfpathlineto{\pgfqpoint{2.869699in}{0.847487in}}%
\pgfpathlineto{\pgfqpoint{2.866618in}{0.860735in}}%
\pgfpathlineto{\pgfqpoint{2.859296in}{0.873983in}}%
\pgfpathlineto{\pgfqpoint{2.864724in}{0.887231in}}%
\pgfpathlineto{\pgfqpoint{2.857472in}{0.900479in}}%
\pgfpathlineto{\pgfqpoint{2.868678in}{0.913727in}}%
\pgfpathlineto{\pgfqpoint{2.863918in}{0.926975in}}%
\pgfpathlineto{\pgfqpoint{2.875117in}{0.936056in}}%
\pgfpathlineto{\pgfqpoint{2.880064in}{0.940223in}}%
\pgfpathlineto{\pgfqpoint{2.875117in}{0.950444in}}%
\pgfpathlineto{\pgfqpoint{2.873760in}{0.953471in}}%
\pgfpathlineto{\pgfqpoint{2.875117in}{0.954898in}}%
\pgfpathlineto{\pgfqpoint{2.885819in}{0.966719in}}%
\pgfpathlineto{\pgfqpoint{2.878226in}{0.979967in}}%
\pgfpathlineto{\pgfqpoint{2.893039in}{0.993215in}}%
\pgfpathlineto{\pgfqpoint{2.889458in}{1.006463in}}%
\pgfpathlineto{\pgfqpoint{2.900400in}{1.014681in}}%
\pgfpathlineto{\pgfqpoint{2.906669in}{1.019711in}}%
\pgfpathlineto{\pgfqpoint{2.903794in}{1.032959in}}%
\pgfpathlineto{\pgfqpoint{2.917522in}{1.046207in}}%
\pgfpathlineto{\pgfqpoint{2.915909in}{1.059455in}}%
\pgfpathlineto{\pgfqpoint{2.925683in}{1.069185in}}%
\pgfpathlineto{\pgfqpoint{2.928924in}{1.072703in}}%
\pgfpathlineto{\pgfqpoint{2.929358in}{1.085951in}}%
\pgfpathlineto{\pgfqpoint{2.940382in}{1.099199in}}%
\pgfpathlineto{\pgfqpoint{2.942912in}{1.112447in}}%
\pgfpathlineto{\pgfqpoint{2.950967in}{1.123642in}}%
\pgfpathlineto{\pgfqpoint{2.952277in}{1.125695in}}%
\pgfpathlineto{\pgfqpoint{2.956925in}{1.138943in}}%
\pgfpathlineto{\pgfqpoint{2.962419in}{1.152191in}}%
\pgfpathlineto{\pgfqpoint{2.966787in}{1.165439in}}%
\pgfpathlineto{\pgfqpoint{2.971550in}{1.178687in}}%
\pgfpathlineto{\pgfqpoint{2.976250in}{1.185896in}}%
\pgfpathlineto{\pgfqpoint{2.979914in}{1.191935in}}%
\pgfpathlineto{\pgfqpoint{2.984342in}{1.205183in}}%
\pgfpathlineto{\pgfqpoint{2.994185in}{1.218431in}}%
\pgfpathlineto{\pgfqpoint{2.995932in}{1.231679in}}%
\pgfpathlineto{\pgfqpoint{3.001533in}{1.239419in}}%
\pgfpathlineto{\pgfqpoint{3.005195in}{1.244927in}}%
\pgfpathlineto{\pgfqpoint{3.006342in}{1.258175in}}%
\pgfpathlineto{\pgfqpoint{3.015845in}{1.271423in}}%
\pgfpathlineto{\pgfqpoint{3.016511in}{1.284671in}}%
\pgfpathlineto{\pgfqpoint{3.026622in}{1.297919in}}%
\pgfpathlineto{\pgfqpoint{3.026816in}{1.299932in}}%
\pgfpathlineto{\pgfqpoint{3.027742in}{1.311167in}}%
\pgfpathlineto{\pgfqpoint{3.036717in}{1.324415in}}%
\pgfpathlineto{\pgfqpoint{3.038377in}{1.337663in}}%
\pgfpathlineto{\pgfqpoint{3.046019in}{1.350911in}}%
\pgfpathlineto{\pgfqpoint{3.048076in}{1.364159in}}%
\pgfpathlineto{\pgfqpoint{3.052099in}{1.370406in}}%
\pgfpathlineto{\pgfqpoint{3.055981in}{1.377407in}}%
\pgfpathlineto{\pgfqpoint{3.059914in}{1.390655in}}%
\pgfpathlineto{\pgfqpoint{3.066136in}{1.403903in}}%
\pgfpathlineto{\pgfqpoint{3.070458in}{1.417151in}}%
\pgfpathlineto{\pgfqpoint{3.074575in}{1.430399in}}%
\pgfpathlineto{\pgfqpoint{3.077382in}{1.437342in}}%
\pgfpathlineto{\pgfqpoint{3.079619in}{1.443647in}}%
\pgfpathlineto{\pgfqpoint{3.082959in}{1.456895in}}%
\pgfpathlineto{\pgfqpoint{3.088925in}{1.470143in}}%
\pgfpathlineto{\pgfqpoint{3.091531in}{1.483391in}}%
\pgfpathlineto{\pgfqpoint{3.099062in}{1.496639in}}%
\pgfpathlineto{\pgfqpoint{3.101389in}{1.509887in}}%
\pgfpathlineto{\pgfqpoint{3.102666in}{1.512005in}}%
\pgfpathlineto{\pgfqpoint{3.108508in}{1.523135in}}%
\pgfpathlineto{\pgfqpoint{3.109932in}{1.536383in}}%
\pgfpathlineto{\pgfqpoint{3.116628in}{1.549631in}}%
\pgfpathlineto{\pgfqpoint{3.118203in}{1.562879in}}%
\pgfpathlineto{\pgfqpoint{3.126357in}{1.576127in}}%
\pgfpathlineto{\pgfqpoint{3.127949in}{1.585652in}}%
\pgfpathlineto{\pgfqpoint{3.128467in}{1.589375in}}%
\pgfpathlineto{\pgfqpoint{3.134276in}{1.602623in}}%
\pgfpathlineto{\pgfqpoint{3.136062in}{1.615871in}}%
\pgfpathlineto{\pgfqpoint{3.141274in}{1.629119in}}%
\pgfpathlineto{\pgfqpoint{3.144592in}{1.642367in}}%
\pgfpathlineto{\pgfqpoint{3.150203in}{1.655615in}}%
\pgfpathlineto{\pgfqpoint{3.153232in}{1.664457in}}%
\pgfpathlineto{\pgfqpoint{3.154540in}{1.668863in}}%
\pgfpathlineto{\pgfqpoint{3.158443in}{1.682111in}}%
\pgfpathlineto{\pgfqpoint{3.163564in}{1.695359in}}%
\pgfpathlineto{\pgfqpoint{3.166585in}{1.708607in}}%
\pgfpathlineto{\pgfqpoint{3.172012in}{1.721855in}}%
\pgfpathlineto{\pgfqpoint{3.174167in}{1.735103in}}%
\pgfpathlineto{\pgfqpoint{3.178515in}{1.744155in}}%
\pgfpathlineto{\pgfqpoint{3.180317in}{1.748351in}}%
\pgfpathlineto{\pgfqpoint{3.182409in}{1.761599in}}%
\pgfpathlineto{\pgfqpoint{3.189346in}{1.774847in}}%
\pgfpathlineto{\pgfqpoint{3.191000in}{1.788095in}}%
\pgfpathlineto{\pgfqpoint{3.197449in}{1.801343in}}%
\pgfpathlineto{\pgfqpoint{3.199353in}{1.814591in}}%
\pgfpathlineto{\pgfqpoint{3.203798in}{1.822734in}}%
\pgfpathlineto{\pgfqpoint{3.206354in}{1.827839in}}%
\pgfpathlineto{\pgfqpoint{3.208970in}{1.841087in}}%
\pgfpathlineto{\pgfqpoint{3.215247in}{1.854335in}}%
\pgfpathlineto{\pgfqpoint{3.218375in}{1.867583in}}%
\pgfpathlineto{\pgfqpoint{3.223832in}{1.880831in}}%
\pgfpathlineto{\pgfqpoint{3.228332in}{1.894079in}}%
\pgfpathlineto{\pgfqpoint{3.229082in}{1.896302in}}%
\pgfpathlineto{\pgfqpoint{3.232557in}{1.907327in}}%
\pgfpathlineto{\pgfqpoint{3.237367in}{1.920575in}}%
\pgfpathlineto{\pgfqpoint{3.241345in}{1.933823in}}%
\pgfpathlineto{\pgfqpoint{3.248705in}{1.947071in}}%
\pgfpathlineto{\pgfqpoint{3.252626in}{1.960319in}}%
\pgfpathlineto{\pgfqpoint{3.254365in}{1.963119in}}%
\pgfpathlineto{\pgfqpoint{3.261610in}{1.973567in}}%
\pgfpathlineto{\pgfqpoint{3.264242in}{1.986815in}}%
\pgfpathlineto{\pgfqpoint{3.273340in}{2.000063in}}%
\pgfpathlineto{\pgfqpoint{3.276424in}{2.013311in}}%
\pgfpathlineto{\pgfqpoint{3.279648in}{2.017475in}}%
\pgfpathlineto{\pgfqpoint{3.289018in}{2.026559in}}%
\pgfpathlineto{\pgfqpoint{3.293512in}{2.039807in}}%
\pgfpathlineto{\pgfqpoint{3.304931in}{2.051431in}}%
\pgfpathlineto{\pgfqpoint{3.308156in}{2.053055in}}%
\pgfpathlineto{\pgfqpoint{3.323808in}{2.066303in}}%
\pgfpathlineto{\pgfqpoint{3.330214in}{2.069900in}}%
\pgfpathlineto{\pgfqpoint{3.346916in}{2.066303in}}%
\pgfpathlineto{\pgfqpoint{3.355497in}{2.064313in}}%
\pgfpathlineto{\pgfqpoint{3.367101in}{2.053055in}}%
\pgfpathlineto{\pgfqpoint{3.380781in}{2.040159in}}%
\pgfpathlineto{\pgfqpoint{3.381002in}{2.039807in}}%
\pgfpathlineto{\pgfqpoint{3.387962in}{2.026559in}}%
\pgfpathlineto{\pgfqpoint{3.396583in}{2.013311in}}%
\pgfpathlineto{\pgfqpoint{3.402271in}{2.000063in}}%
\pgfpathlineto{\pgfqpoint{3.406064in}{1.994093in}}%
\pgfpathlineto{\pgfqpoint{3.409809in}{1.986815in}}%
\pgfpathlineto{\pgfqpoint{3.413837in}{1.973567in}}%
\pgfpathlineto{\pgfqpoint{3.420950in}{1.960319in}}%
\pgfpathlineto{\pgfqpoint{3.424436in}{1.947071in}}%
\pgfpathlineto{\pgfqpoint{3.431347in}{1.936195in}}%
\pgfpathlineto{\pgfqpoint{3.432782in}{1.933823in}}%
\pgfpathlineto{\pgfqpoint{3.436011in}{1.920575in}}%
\pgfpathlineto{\pgfqpoint{3.443145in}{1.907327in}}%
\pgfpathlineto{\pgfqpoint{3.445215in}{1.894079in}}%
\pgfpathlineto{\pgfqpoint{3.451923in}{1.880831in}}%
\pgfpathlineto{\pgfqpoint{3.454054in}{1.867583in}}%
\pgfpathlineto{\pgfqpoint{3.456630in}{1.862959in}}%
\pgfpathlineto{\pgfqpoint{3.461671in}{1.854335in}}%
\pgfpathlineto{\pgfqpoint{3.463804in}{1.841087in}}%
\pgfpathlineto{\pgfqpoint{3.470290in}{1.827839in}}%
\pgfpathlineto{\pgfqpoint{3.472203in}{1.814591in}}%
\pgfpathlineto{\pgfqpoint{3.478125in}{1.801343in}}%
\pgfpathlineto{\pgfqpoint{3.480286in}{1.788095in}}%
\pgfpathlineto{\pgfqpoint{3.481913in}{1.784351in}}%
\pgfpathlineto{\pgfqpoint{3.486564in}{1.774847in}}%
\pgfpathlineto{\pgfqpoint{3.489516in}{1.761599in}}%
\pgfpathlineto{\pgfqpoint{3.495284in}{1.748351in}}%
\pgfpathlineto{\pgfqpoint{3.499113in}{1.735103in}}%
\pgfpathlineto{\pgfqpoint{3.503575in}{1.721855in}}%
\pgfpathlineto{\pgfqpoint{3.506859in}{1.708607in}}%
\pgfpathlineto{\pgfqpoint{3.507196in}{1.707314in}}%
\pgfpathlineto{\pgfqpoint{3.510892in}{1.695359in}}%
\pgfpathlineto{\pgfqpoint{3.515466in}{1.682111in}}%
\pgfpathlineto{\pgfqpoint{3.519202in}{1.668863in}}%
\pgfpathlineto{\pgfqpoint{3.524753in}{1.655615in}}%
\pgfpathlineto{\pgfqpoint{3.527307in}{1.642367in}}%
\pgfpathlineto{\pgfqpoint{3.532456in}{1.629119in}}%
\pgfpathlineto{\pgfqpoint{3.532480in}{1.628964in}}%
\pgfpathlineto{\pgfqpoint{3.534894in}{1.615871in}}%
\pgfpathlineto{\pgfqpoint{3.541611in}{1.602623in}}%
\pgfpathlineto{\pgfqpoint{3.543737in}{1.589375in}}%
\pgfpathlineto{\pgfqpoint{3.550693in}{1.576127in}}%
\pgfpathlineto{\pgfqpoint{3.552392in}{1.562879in}}%
\pgfpathlineto{\pgfqpoint{3.557763in}{1.552403in}}%
\pgfpathlineto{\pgfqpoint{3.559421in}{1.549631in}}%
\pgfpathlineto{\pgfqpoint{3.561559in}{1.536383in}}%
\pgfpathlineto{\pgfqpoint{3.568218in}{1.523135in}}%
\pgfpathlineto{\pgfqpoint{3.569463in}{1.509887in}}%
\pgfpathlineto{\pgfqpoint{3.575947in}{1.496639in}}%
\pgfpathlineto{\pgfqpoint{3.578306in}{1.483391in}}%
\pgfpathlineto{\pgfqpoint{3.583046in}{1.473750in}}%
\pgfpathlineto{\pgfqpoint{3.585136in}{1.470143in}}%
\pgfpathlineto{\pgfqpoint{3.588841in}{1.456895in}}%
\pgfpathlineto{\pgfqpoint{3.594498in}{1.443647in}}%
\pgfpathlineto{\pgfqpoint{3.597966in}{1.430399in}}%
\pgfpathlineto{\pgfqpoint{3.603016in}{1.417151in}}%
\pgfpathlineto{\pgfqpoint{3.607666in}{1.403903in}}%
\pgfpathlineto{\pgfqpoint{3.608329in}{1.401693in}}%
\pgfpathlineto{\pgfqpoint{3.612275in}{1.390655in}}%
\pgfpathlineto{\pgfqpoint{3.618640in}{1.377407in}}%
\pgfpathlineto{\pgfqpoint{3.621802in}{1.364159in}}%
\pgfpathlineto{\pgfqpoint{3.628523in}{1.350911in}}%
\pgfpathlineto{\pgfqpoint{3.630863in}{1.337663in}}%
\pgfpathlineto{\pgfqpoint{3.633612in}{1.332201in}}%
\pgfpathlineto{\pgfqpoint{3.638001in}{1.324415in}}%
\pgfpathlineto{\pgfqpoint{3.639465in}{1.311167in}}%
\pgfpathlineto{\pgfqpoint{3.648898in}{1.297919in}}%
\pgfpathlineto{\pgfqpoint{3.651121in}{1.284671in}}%
\pgfpathlineto{\pgfqpoint{3.658895in}{1.273713in}}%
\pgfpathlineto{\pgfqpoint{3.660715in}{1.271423in}}%
\pgfpathlineto{\pgfqpoint{3.662287in}{1.258175in}}%
\pgfpathlineto{\pgfqpoint{3.671029in}{1.244927in}}%
\pgfpathlineto{\pgfqpoint{3.671924in}{1.231679in}}%
\pgfpathlineto{\pgfqpoint{3.681262in}{1.218431in}}%
\pgfpathlineto{\pgfqpoint{3.683470in}{1.205183in}}%
\pgfpathlineto{\pgfqpoint{3.684179in}{1.203973in}}%
\pgfpathlineto{\pgfqpoint{3.692143in}{1.191935in}}%
\pgfpathlineto{\pgfqpoint{3.695671in}{1.178687in}}%
\pgfpathlineto{\pgfqpoint{3.702798in}{1.165439in}}%
\pgfpathlineto{\pgfqpoint{3.707197in}{1.152191in}}%
\pgfpathlineto{\pgfqpoint{3.709462in}{1.146819in}}%
\pgfpathlineto{\pgfqpoint{3.713225in}{1.138943in}}%
\pgfpathlineto{\pgfqpoint{3.719104in}{1.125695in}}%
\pgfpathlineto{\pgfqpoint{3.723517in}{1.112447in}}%
\pgfpathlineto{\pgfqpoint{3.733442in}{1.099199in}}%
\pgfpathlineto{\pgfqpoint{3.734745in}{1.093640in}}%
\pgfpathlineto{\pgfqpoint{3.736794in}{1.085951in}}%
\pgfpathlineto{\pgfqpoint{3.746509in}{1.072703in}}%
\pgfpathlineto{\pgfqpoint{3.745743in}{1.059455in}}%
\pgfpathlineto{\pgfqpoint{3.756008in}{1.046207in}}%
\pgfpathlineto{\pgfqpoint{3.753897in}{1.032959in}}%
\pgfpathlineto{\pgfqpoint{3.760028in}{1.027183in}}%
\pgfpathlineto{\pgfqpoint{3.768466in}{1.019711in}}%
\pgfpathlineto{\pgfqpoint{3.766337in}{1.006463in}}%
\pgfpathlineto{\pgfqpoint{3.779637in}{0.993215in}}%
\pgfpathlineto{\pgfqpoint{3.777203in}{0.979967in}}%
\pgfpathlineto{\pgfqpoint{3.785311in}{0.971146in}}%
\pgfpathlineto{\pgfqpoint{3.789690in}{0.966719in}}%
\pgfpathlineto{\pgfqpoint{3.785338in}{0.953471in}}%
\pgfpathlineto{\pgfqpoint{3.795008in}{0.940223in}}%
\pgfpathlineto{\pgfqpoint{3.792328in}{0.926975in}}%
\pgfpathlineto{\pgfqpoint{3.798243in}{0.913727in}}%
\pgfpathlineto{\pgfqpoint{3.800580in}{0.900479in}}%
\pgfpathlineto{\pgfqpoint{3.798891in}{0.887231in}}%
\pgfpathlineto{\pgfqpoint{3.796319in}{0.873983in}}%
\pgfpathlineto{\pgfqpoint{3.785311in}{0.863510in}}%
\pgfpathlineto{\pgfqpoint{3.781958in}{0.860735in}}%
\pgfpathlineto{\pgfqpoint{3.785311in}{0.853467in}}%
\pgfpathlineto{\pgfqpoint{3.788153in}{0.847487in}}%
\pgfpathlineto{\pgfqpoint{3.785311in}{0.845740in}}%
\pgfpathlineto{\pgfqpoint{3.763389in}{0.834239in}}%
\pgfpathlineto{\pgfqpoint{3.772964in}{0.820991in}}%
\pgfpathlineto{\pgfqpoint{3.760028in}{0.816528in}}%
\pgfpathlineto{\pgfqpoint{3.734745in}{0.809131in}}%
\pgfpathlineto{\pgfqpoint{3.729296in}{0.807743in}}%
\pgfpathlineto{\pgfqpoint{3.728607in}{0.794495in}}%
\pgfpathlineto{\pgfqpoint{3.709462in}{0.790857in}}%
\pgfpathlineto{\pgfqpoint{3.684179in}{0.786665in}}%
\pgfpathlineto{\pgfqpoint{3.658895in}{0.783041in}}%
\pgfpathlineto{\pgfqpoint{3.644760in}{0.781247in}}%
\pgfpathlineto{\pgfqpoint{3.633612in}{0.775497in}}%
\pgfpathlineto{\pgfqpoint{3.614958in}{0.767999in}}%
\pgfpathlineto{\pgfqpoint{3.608329in}{0.767350in}}%
\pgfpathlineto{\pgfqpoint{3.583046in}{0.765110in}}%
\pgfpathlineto{\pgfqpoint{3.557763in}{0.763117in}}%
\pgfpathlineto{\pgfqpoint{3.532480in}{0.761349in}}%
\pgfpathlineto{\pgfqpoint{3.507196in}{0.759789in}}%
\pgfpathlineto{\pgfqpoint{3.481913in}{0.758423in}}%
\pgfpathlineto{\pgfqpoint{3.456630in}{0.757239in}}%
\pgfpathlineto{\pgfqpoint{3.431347in}{0.756232in}}%
\pgfpathlineto{\pgfqpoint{3.406064in}{0.755397in}}%
\pgfpathlineto{\pgfqpoint{3.381498in}{0.754751in}}%
\pgfpathlineto{\pgfqpoint{3.380781in}{0.754722in}}%
\pgfpathlineto{\pgfqpoint{3.355497in}{0.753995in}}%
\pgfpathlineto{\pgfqpoint{3.330214in}{0.753593in}}%
\pgfpathlineto{\pgfqpoint{3.304931in}{0.753488in}}%
\pgfpathlineto{\pgfqpoint{3.279648in}{0.753660in}}%
\pgfpathlineto{\pgfqpoint{3.254365in}{0.754097in}}%
\pgfpathlineto{\pgfqpoint{3.230460in}{0.754751in}}%
\pgfpathclose%
\pgfpathmoveto{\pgfqpoint{2.824551in}{1.581427in}}%
\pgfpathlineto{\pgfqpoint{2.838811in}{1.589375in}}%
\pgfpathlineto{\pgfqpoint{2.849834in}{1.593128in}}%
\pgfpathlineto{\pgfqpoint{2.862191in}{1.602623in}}%
\pgfpathlineto{\pgfqpoint{2.873001in}{1.615871in}}%
\pgfpathlineto{\pgfqpoint{2.875117in}{1.616835in}}%
\pgfpathlineto{\pgfqpoint{2.891290in}{1.629119in}}%
\pgfpathlineto{\pgfqpoint{2.892429in}{1.642367in}}%
\pgfpathlineto{\pgfqpoint{2.900400in}{1.646611in}}%
\pgfpathlineto{\pgfqpoint{2.911757in}{1.655615in}}%
\pgfpathlineto{\pgfqpoint{2.911791in}{1.668863in}}%
\pgfpathlineto{\pgfqpoint{2.925683in}{1.678720in}}%
\pgfpathlineto{\pgfqpoint{2.929113in}{1.682111in}}%
\pgfpathlineto{\pgfqpoint{2.928736in}{1.695359in}}%
\pgfpathlineto{\pgfqpoint{2.946283in}{1.708607in}}%
\pgfpathlineto{\pgfqpoint{2.950023in}{1.721855in}}%
\pgfpathlineto{\pgfqpoint{2.950967in}{1.722483in}}%
\pgfpathlineto{\pgfqpoint{2.963945in}{1.735103in}}%
\pgfpathlineto{\pgfqpoint{2.967167in}{1.748351in}}%
\pgfpathlineto{\pgfqpoint{2.976250in}{1.760170in}}%
\pgfpathlineto{\pgfqpoint{2.977101in}{1.761599in}}%
\pgfpathlineto{\pgfqpoint{2.981093in}{1.774847in}}%
\pgfpathlineto{\pgfqpoint{2.989272in}{1.788095in}}%
\pgfpathlineto{\pgfqpoint{2.995968in}{1.801343in}}%
\pgfpathlineto{\pgfqpoint{3.001533in}{1.812415in}}%
\pgfpathlineto{\pgfqpoint{3.002427in}{1.814591in}}%
\pgfpathlineto{\pgfqpoint{3.010390in}{1.827839in}}%
\pgfpathlineto{\pgfqpoint{3.014397in}{1.841087in}}%
\pgfpathlineto{\pgfqpoint{3.023712in}{1.854335in}}%
\pgfpathlineto{\pgfqpoint{3.026682in}{1.867583in}}%
\pgfpathlineto{\pgfqpoint{3.026816in}{1.867754in}}%
\pgfpathlineto{\pgfqpoint{3.035565in}{1.880831in}}%
\pgfpathlineto{\pgfqpoint{3.037733in}{1.894079in}}%
\pgfpathlineto{\pgfqpoint{3.048505in}{1.907327in}}%
\pgfpathlineto{\pgfqpoint{3.050837in}{1.920575in}}%
\pgfpathlineto{\pgfqpoint{3.052099in}{1.922321in}}%
\pgfpathlineto{\pgfqpoint{3.059335in}{1.933823in}}%
\pgfpathlineto{\pgfqpoint{3.060629in}{1.947071in}}%
\pgfpathlineto{\pgfqpoint{3.068165in}{1.960319in}}%
\pgfpathlineto{\pgfqpoint{3.070989in}{1.973567in}}%
\pgfpathlineto{\pgfqpoint{3.077382in}{1.984172in}}%
\pgfpathlineto{\pgfqpoint{3.078849in}{1.986815in}}%
\pgfpathlineto{\pgfqpoint{3.082900in}{2.000063in}}%
\pgfpathlineto{\pgfqpoint{3.088895in}{2.013311in}}%
\pgfpathlineto{\pgfqpoint{3.094107in}{2.026559in}}%
\pgfpathlineto{\pgfqpoint{3.098654in}{2.039807in}}%
\pgfpathlineto{\pgfqpoint{3.102666in}{2.049761in}}%
\pgfpathlineto{\pgfqpoint{3.104018in}{2.053055in}}%
\pgfpathlineto{\pgfqpoint{3.107122in}{2.066303in}}%
\pgfpathlineto{\pgfqpoint{3.113511in}{2.079551in}}%
\pgfpathlineto{\pgfqpoint{3.116257in}{2.092799in}}%
\pgfpathlineto{\pgfqpoint{3.123582in}{2.106047in}}%
\pgfpathlineto{\pgfqpoint{3.125373in}{2.119295in}}%
\pgfpathlineto{\pgfqpoint{3.127949in}{2.124299in}}%
\pgfpathlineto{\pgfqpoint{3.132271in}{2.132543in}}%
\pgfpathlineto{\pgfqpoint{3.127949in}{2.132543in}}%
\pgfpathlineto{\pgfqpoint{3.118216in}{2.132543in}}%
\pgfpathlineto{\pgfqpoint{3.111219in}{2.119295in}}%
\pgfpathlineto{\pgfqpoint{3.109452in}{2.106047in}}%
\pgfpathlineto{\pgfqpoint{3.102666in}{2.094300in}}%
\pgfpathlineto{\pgfqpoint{3.101705in}{2.092799in}}%
\pgfpathlineto{\pgfqpoint{3.098507in}{2.079551in}}%
\pgfpathlineto{\pgfqpoint{3.091996in}{2.066303in}}%
\pgfpathlineto{\pgfqpoint{3.088486in}{2.053055in}}%
\pgfpathlineto{\pgfqpoint{3.083410in}{2.039807in}}%
\pgfpathlineto{\pgfqpoint{3.078150in}{2.026559in}}%
\pgfpathlineto{\pgfqpoint{3.077382in}{2.024520in}}%
\pgfpathlineto{\pgfqpoint{3.072774in}{2.013311in}}%
\pgfpathlineto{\pgfqpoint{3.065370in}{2.000063in}}%
\pgfpathlineto{\pgfqpoint{3.061935in}{1.986815in}}%
\pgfpathlineto{\pgfqpoint{3.053265in}{1.973567in}}%
\pgfpathlineto{\pgfqpoint{3.052099in}{1.966357in}}%
\pgfpathlineto{\pgfqpoint{3.050958in}{1.960319in}}%
\pgfpathlineto{\pgfqpoint{3.041298in}{1.947071in}}%
\pgfpathlineto{\pgfqpoint{3.040874in}{1.933823in}}%
\pgfpathlineto{\pgfqpoint{3.031087in}{1.920575in}}%
\pgfpathlineto{\pgfqpoint{3.029353in}{1.907327in}}%
\pgfpathlineto{\pgfqpoint{3.026816in}{1.904274in}}%
\pgfpathlineto{\pgfqpoint{3.015972in}{1.894079in}}%
\pgfpathlineto{\pgfqpoint{3.013883in}{1.880831in}}%
\pgfpathlineto{\pgfqpoint{3.003967in}{1.867583in}}%
\pgfpathlineto{\pgfqpoint{3.001533in}{1.856183in}}%
\pgfpathlineto{\pgfqpoint{3.000958in}{1.854335in}}%
\pgfpathlineto{\pgfqpoint{2.989371in}{1.841087in}}%
\pgfpathlineto{\pgfqpoint{2.984099in}{1.827839in}}%
\pgfpathlineto{\pgfqpoint{2.976739in}{1.814591in}}%
\pgfpathlineto{\pgfqpoint{2.976250in}{1.813664in}}%
\pgfpathlineto{\pgfqpoint{2.965989in}{1.801343in}}%
\pgfpathlineto{\pgfqpoint{2.960170in}{1.788095in}}%
\pgfpathlineto{\pgfqpoint{2.950967in}{1.777748in}}%
\pgfpathlineto{\pgfqpoint{2.945606in}{1.774847in}}%
\pgfpathlineto{\pgfqpoint{2.944288in}{1.761599in}}%
\pgfpathlineto{\pgfqpoint{2.927103in}{1.748351in}}%
\pgfpathlineto{\pgfqpoint{2.928486in}{1.735103in}}%
\pgfpathlineto{\pgfqpoint{2.925683in}{1.733233in}}%
\pgfpathlineto{\pgfqpoint{2.900400in}{1.725909in}}%
\pgfpathlineto{\pgfqpoint{2.875117in}{1.730512in}}%
\pgfpathlineto{\pgfqpoint{2.869798in}{1.735103in}}%
\pgfpathlineto{\pgfqpoint{2.875117in}{1.741404in}}%
\pgfpathlineto{\pgfqpoint{2.886725in}{1.748351in}}%
\pgfpathlineto{\pgfqpoint{2.875117in}{1.757004in}}%
\pgfpathlineto{\pgfqpoint{2.872708in}{1.761599in}}%
\pgfpathlineto{\pgfqpoint{2.875117in}{1.764328in}}%
\pgfpathlineto{\pgfqpoint{2.888956in}{1.774847in}}%
\pgfpathlineto{\pgfqpoint{2.879286in}{1.788095in}}%
\pgfpathlineto{\pgfqpoint{2.889338in}{1.801343in}}%
\pgfpathlineto{\pgfqpoint{2.884321in}{1.814591in}}%
\pgfpathlineto{\pgfqpoint{2.893223in}{1.827839in}}%
\pgfpathlineto{\pgfqpoint{2.892384in}{1.841087in}}%
\pgfpathlineto{\pgfqpoint{2.898128in}{1.854335in}}%
\pgfpathlineto{\pgfqpoint{2.898977in}{1.867583in}}%
\pgfpathlineto{\pgfqpoint{2.900400in}{1.872207in}}%
\pgfpathlineto{\pgfqpoint{2.903512in}{1.880831in}}%
\pgfpathlineto{\pgfqpoint{2.906760in}{1.894079in}}%
\pgfpathlineto{\pgfqpoint{2.908280in}{1.907327in}}%
\pgfpathlineto{\pgfqpoint{2.911701in}{1.920575in}}%
\pgfpathlineto{\pgfqpoint{2.912698in}{1.933823in}}%
\pgfpathlineto{\pgfqpoint{2.919238in}{1.947071in}}%
\pgfpathlineto{\pgfqpoint{2.920414in}{1.960319in}}%
\pgfpathlineto{\pgfqpoint{2.925683in}{1.970286in}}%
\pgfpathlineto{\pgfqpoint{2.927780in}{1.973567in}}%
\pgfpathlineto{\pgfqpoint{2.926808in}{1.986815in}}%
\pgfpathlineto{\pgfqpoint{2.933921in}{2.000063in}}%
\pgfpathlineto{\pgfqpoint{2.932945in}{2.013311in}}%
\pgfpathlineto{\pgfqpoint{2.940552in}{2.026559in}}%
\pgfpathlineto{\pgfqpoint{2.939715in}{2.039807in}}%
\pgfpathlineto{\pgfqpoint{2.946329in}{2.053055in}}%
\pgfpathlineto{\pgfqpoint{2.946069in}{2.066303in}}%
\pgfpathlineto{\pgfqpoint{2.950967in}{2.076819in}}%
\pgfpathlineto{\pgfqpoint{2.952340in}{2.079551in}}%
\pgfpathlineto{\pgfqpoint{2.952830in}{2.092799in}}%
\pgfpathlineto{\pgfqpoint{2.958042in}{2.106047in}}%
\pgfpathlineto{\pgfqpoint{2.959256in}{2.119295in}}%
\pgfpathlineto{\pgfqpoint{2.964146in}{2.132543in}}%
\pgfpathlineto{\pgfqpoint{2.950967in}{2.132543in}}%
\pgfpathlineto{\pgfqpoint{2.950234in}{2.132543in}}%
\pgfpathlineto{\pgfqpoint{2.944598in}{2.119295in}}%
\pgfpathlineto{\pgfqpoint{2.943288in}{2.106047in}}%
\pgfpathlineto{\pgfqpoint{2.937929in}{2.092799in}}%
\pgfpathlineto{\pgfqpoint{2.937147in}{2.079551in}}%
\pgfpathlineto{\pgfqpoint{2.930903in}{2.066303in}}%
\pgfpathlineto{\pgfqpoint{2.930567in}{2.053055in}}%
\pgfpathlineto{\pgfqpoint{2.925683in}{2.043700in}}%
\pgfpathlineto{\pgfqpoint{2.923593in}{2.039807in}}%
\pgfpathlineto{\pgfqpoint{2.923318in}{2.026559in}}%
\pgfpathlineto{\pgfqpoint{2.915955in}{2.013311in}}%
\pgfpathlineto{\pgfqpoint{2.915618in}{2.000063in}}%
\pgfpathlineto{\pgfqpoint{2.909476in}{1.986815in}}%
\pgfpathlineto{\pgfqpoint{2.908814in}{1.973567in}}%
\pgfpathlineto{\pgfqpoint{2.901877in}{1.960319in}}%
\pgfpathlineto{\pgfqpoint{2.900400in}{1.953540in}}%
\pgfpathlineto{\pgfqpoint{2.898941in}{1.947071in}}%
\pgfpathlineto{\pgfqpoint{2.893130in}{1.933823in}}%
\pgfpathlineto{\pgfqpoint{2.890338in}{1.920575in}}%
\pgfpathlineto{\pgfqpoint{2.887326in}{1.907327in}}%
\pgfpathlineto{\pgfqpoint{2.883822in}{1.894079in}}%
\pgfpathlineto{\pgfqpoint{2.880920in}{1.880831in}}%
\pgfpathlineto{\pgfqpoint{2.875117in}{1.868203in}}%
\pgfpathlineto{\pgfqpoint{2.874834in}{1.867583in}}%
\pgfpathlineto{\pgfqpoint{2.873336in}{1.854335in}}%
\pgfpathlineto{\pgfqpoint{2.865808in}{1.841087in}}%
\pgfpathlineto{\pgfqpoint{2.865222in}{1.827839in}}%
\pgfpathlineto{\pgfqpoint{2.856005in}{1.814591in}}%
\pgfpathlineto{\pgfqpoint{2.857592in}{1.801343in}}%
\pgfpathlineto{\pgfqpoint{2.849834in}{1.790445in}}%
\pgfpathlineto{\pgfqpoint{2.848319in}{1.788095in}}%
\pgfpathlineto{\pgfqpoint{2.849834in}{1.777731in}}%
\pgfpathlineto{\pgfqpoint{2.850391in}{1.774847in}}%
\pgfpathlineto{\pgfqpoint{2.849834in}{1.774281in}}%
\pgfpathlineto{\pgfqpoint{2.839329in}{1.761599in}}%
\pgfpathlineto{\pgfqpoint{2.841146in}{1.748351in}}%
\pgfpathlineto{\pgfqpoint{2.831328in}{1.735103in}}%
\pgfpathlineto{\pgfqpoint{2.835873in}{1.721855in}}%
\pgfpathlineto{\pgfqpoint{2.826214in}{1.708607in}}%
\pgfpathlineto{\pgfqpoint{2.829455in}{1.695359in}}%
\pgfpathlineto{\pgfqpoint{2.824551in}{1.690504in}}%
\pgfpathlineto{\pgfqpoint{2.818329in}{1.682111in}}%
\pgfpathlineto{\pgfqpoint{2.818962in}{1.668863in}}%
\pgfpathlineto{\pgfqpoint{2.810992in}{1.655615in}}%
\pgfpathlineto{\pgfqpoint{2.810536in}{1.642367in}}%
\pgfpathlineto{\pgfqpoint{2.805355in}{1.629119in}}%
\pgfpathlineto{\pgfqpoint{2.803788in}{1.615871in}}%
\pgfpathlineto{\pgfqpoint{2.809198in}{1.602623in}}%
\pgfpathlineto{\pgfqpoint{2.810614in}{1.589375in}}%
\pgfpathlineto{\pgfqpoint{2.824551in}{1.581427in}}%
\pgfpathclose%
\pgfpathmoveto{\pgfqpoint{3.835878in}{1.598428in}}%
\pgfpathlineto{\pgfqpoint{3.858119in}{1.602623in}}%
\pgfpathlineto{\pgfqpoint{3.835878in}{1.613956in}}%
\pgfpathlineto{\pgfqpoint{3.820646in}{1.602623in}}%
\pgfpathlineto{\pgfqpoint{3.835878in}{1.598428in}}%
\pgfpathclose%
\pgfpathmoveto{\pgfqpoint{3.810594in}{1.619850in}}%
\pgfpathlineto{\pgfqpoint{3.835878in}{1.616732in}}%
\pgfpathlineto{\pgfqpoint{3.857937in}{1.629119in}}%
\pgfpathlineto{\pgfqpoint{3.845151in}{1.642367in}}%
\pgfpathlineto{\pgfqpoint{3.854080in}{1.655615in}}%
\pgfpathlineto{\pgfqpoint{3.842713in}{1.668863in}}%
\pgfpathlineto{\pgfqpoint{3.845705in}{1.682111in}}%
\pgfpathlineto{\pgfqpoint{3.839836in}{1.695359in}}%
\pgfpathlineto{\pgfqpoint{3.840387in}{1.708607in}}%
\pgfpathlineto{\pgfqpoint{3.838041in}{1.721855in}}%
\pgfpathlineto{\pgfqpoint{3.835878in}{1.729422in}}%
\pgfpathlineto{\pgfqpoint{3.834012in}{1.735103in}}%
\pgfpathlineto{\pgfqpoint{3.829203in}{1.748351in}}%
\pgfpathlineto{\pgfqpoint{3.822678in}{1.761599in}}%
\pgfpathlineto{\pgfqpoint{3.821597in}{1.774847in}}%
\pgfpathlineto{\pgfqpoint{3.814738in}{1.788095in}}%
\pgfpathlineto{\pgfqpoint{3.815368in}{1.801343in}}%
\pgfpathlineto{\pgfqpoint{3.810594in}{1.810194in}}%
\pgfpathlineto{\pgfqpoint{3.807896in}{1.814591in}}%
\pgfpathlineto{\pgfqpoint{3.808603in}{1.827839in}}%
\pgfpathlineto{\pgfqpoint{3.800630in}{1.841087in}}%
\pgfpathlineto{\pgfqpoint{3.800825in}{1.854335in}}%
\pgfpathlineto{\pgfqpoint{3.792067in}{1.867583in}}%
\pgfpathlineto{\pgfqpoint{3.792253in}{1.880831in}}%
\pgfpathlineto{\pgfqpoint{3.785629in}{1.894079in}}%
\pgfpathlineto{\pgfqpoint{3.785311in}{1.905065in}}%
\pgfpathlineto{\pgfqpoint{3.785245in}{1.907327in}}%
\pgfpathlineto{\pgfqpoint{3.778545in}{1.920575in}}%
\pgfpathlineto{\pgfqpoint{3.776591in}{1.933823in}}%
\pgfpathlineto{\pgfqpoint{3.769976in}{1.947071in}}%
\pgfpathlineto{\pgfqpoint{3.767936in}{1.960319in}}%
\pgfpathlineto{\pgfqpoint{3.764069in}{1.973567in}}%
\pgfpathlineto{\pgfqpoint{3.760949in}{1.986815in}}%
\pgfpathlineto{\pgfqpoint{3.760028in}{1.990690in}}%
\pgfpathlineto{\pgfqpoint{3.757777in}{2.000063in}}%
\pgfpathlineto{\pgfqpoint{3.753807in}{2.013311in}}%
\pgfpathlineto{\pgfqpoint{3.751073in}{2.026559in}}%
\pgfpathlineto{\pgfqpoint{3.746078in}{2.039807in}}%
\pgfpathlineto{\pgfqpoint{3.743840in}{2.053055in}}%
\pgfpathlineto{\pgfqpoint{3.737701in}{2.066303in}}%
\pgfpathlineto{\pgfqpoint{3.736694in}{2.079551in}}%
\pgfpathlineto{\pgfqpoint{3.734745in}{2.084312in}}%
\pgfpathlineto{\pgfqpoint{3.731390in}{2.092799in}}%
\pgfpathlineto{\pgfqpoint{3.730417in}{2.106047in}}%
\pgfpathlineto{\pgfqpoint{3.724483in}{2.119295in}}%
\pgfpathlineto{\pgfqpoint{3.723441in}{2.132543in}}%
\pgfpathlineto{\pgfqpoint{3.709462in}{2.132543in}}%
\pgfpathlineto{\pgfqpoint{3.709201in}{2.132543in}}%
\pgfpathlineto{\pgfqpoint{3.709462in}{2.127053in}}%
\pgfpathlineto{\pgfqpoint{3.709819in}{2.119295in}}%
\pgfpathlineto{\pgfqpoint{3.715682in}{2.106047in}}%
\pgfpathlineto{\pgfqpoint{3.716055in}{2.092799in}}%
\pgfpathlineto{\pgfqpoint{3.721002in}{2.079551in}}%
\pgfpathlineto{\pgfqpoint{3.721505in}{2.066303in}}%
\pgfpathlineto{\pgfqpoint{3.727308in}{2.053055in}}%
\pgfpathlineto{\pgfqpoint{3.729381in}{2.039807in}}%
\pgfpathlineto{\pgfqpoint{3.734070in}{2.026559in}}%
\pgfpathlineto{\pgfqpoint{3.734745in}{2.023343in}}%
\pgfpathlineto{\pgfqpoint{3.736603in}{2.013311in}}%
\pgfpathlineto{\pgfqpoint{3.739605in}{2.000063in}}%
\pgfpathlineto{\pgfqpoint{3.742451in}{1.986815in}}%
\pgfpathlineto{\pgfqpoint{3.744129in}{1.973567in}}%
\pgfpathlineto{\pgfqpoint{3.748144in}{1.960319in}}%
\pgfpathlineto{\pgfqpoint{3.748525in}{1.947071in}}%
\pgfpathlineto{\pgfqpoint{3.756219in}{1.933823in}}%
\pgfpathlineto{\pgfqpoint{3.756600in}{1.920575in}}%
\pgfpathlineto{\pgfqpoint{3.760028in}{1.914768in}}%
\pgfpathlineto{\pgfqpoint{3.763548in}{1.907327in}}%
\pgfpathlineto{\pgfqpoint{3.761631in}{1.894079in}}%
\pgfpathlineto{\pgfqpoint{3.768502in}{1.880831in}}%
\pgfpathlineto{\pgfqpoint{3.765651in}{1.867583in}}%
\pgfpathlineto{\pgfqpoint{3.774871in}{1.854335in}}%
\pgfpathlineto{\pgfqpoint{3.771520in}{1.841087in}}%
\pgfpathlineto{\pgfqpoint{3.779443in}{1.827839in}}%
\pgfpathlineto{\pgfqpoint{3.774369in}{1.814591in}}%
\pgfpathlineto{\pgfqpoint{3.783037in}{1.801343in}}%
\pgfpathlineto{\pgfqpoint{3.775815in}{1.788095in}}%
\pgfpathlineto{\pgfqpoint{3.781417in}{1.774847in}}%
\pgfpathlineto{\pgfqpoint{3.775798in}{1.761599in}}%
\pgfpathlineto{\pgfqpoint{3.764256in}{1.748351in}}%
\pgfpathlineto{\pgfqpoint{3.760028in}{1.743514in}}%
\pgfpathlineto{\pgfqpoint{3.759149in}{1.748351in}}%
\pgfpathlineto{\pgfqpoint{3.745619in}{1.761599in}}%
\pgfpathlineto{\pgfqpoint{3.734745in}{1.771039in}}%
\pgfpathlineto{\pgfqpoint{3.733123in}{1.774847in}}%
\pgfpathlineto{\pgfqpoint{3.723589in}{1.788095in}}%
\pgfpathlineto{\pgfqpoint{3.709462in}{1.801189in}}%
\pgfpathlineto{\pgfqpoint{3.709368in}{1.801343in}}%
\pgfpathlineto{\pgfqpoint{3.704007in}{1.814591in}}%
\pgfpathlineto{\pgfqpoint{3.693343in}{1.827839in}}%
\pgfpathlineto{\pgfqpoint{3.687247in}{1.841087in}}%
\pgfpathlineto{\pgfqpoint{3.684179in}{1.843902in}}%
\pgfpathlineto{\pgfqpoint{3.677164in}{1.854335in}}%
\pgfpathlineto{\pgfqpoint{3.673702in}{1.867583in}}%
\pgfpathlineto{\pgfqpoint{3.662012in}{1.880831in}}%
\pgfpathlineto{\pgfqpoint{3.658895in}{1.893749in}}%
\pgfpathlineto{\pgfqpoint{3.658834in}{1.894079in}}%
\pgfpathlineto{\pgfqpoint{3.647822in}{1.907327in}}%
\pgfpathlineto{\pgfqpoint{3.645220in}{1.920575in}}%
\pgfpathlineto{\pgfqpoint{3.634819in}{1.933823in}}%
\pgfpathlineto{\pgfqpoint{3.633792in}{1.947071in}}%
\pgfpathlineto{\pgfqpoint{3.633612in}{1.947271in}}%
\pgfpathlineto{\pgfqpoint{3.624991in}{1.960319in}}%
\pgfpathlineto{\pgfqpoint{3.623260in}{1.973567in}}%
\pgfpathlineto{\pgfqpoint{3.612672in}{1.986815in}}%
\pgfpathlineto{\pgfqpoint{3.609848in}{2.000063in}}%
\pgfpathlineto{\pgfqpoint{3.608329in}{2.002120in}}%
\pgfpathlineto{\pgfqpoint{3.601543in}{2.013311in}}%
\pgfpathlineto{\pgfqpoint{3.598994in}{2.026559in}}%
\pgfpathlineto{\pgfqpoint{3.591252in}{2.039807in}}%
\pgfpathlineto{\pgfqpoint{3.588235in}{2.053055in}}%
\pgfpathlineto{\pgfqpoint{3.583046in}{2.062131in}}%
\pgfpathlineto{\pgfqpoint{3.580896in}{2.066303in}}%
\pgfpathlineto{\pgfqpoint{3.577518in}{2.079551in}}%
\pgfpathlineto{\pgfqpoint{3.571737in}{2.092799in}}%
\pgfpathlineto{\pgfqpoint{3.567001in}{2.106047in}}%
\pgfpathlineto{\pgfqpoint{3.562030in}{2.119295in}}%
\pgfpathlineto{\pgfqpoint{3.557763in}{2.131784in}}%
\pgfpathlineto{\pgfqpoint{3.557523in}{2.132543in}}%
\pgfpathlineto{\pgfqpoint{3.543355in}{2.132543in}}%
\pgfpathlineto{\pgfqpoint{3.547182in}{2.119295in}}%
\pgfpathlineto{\pgfqpoint{3.552121in}{2.106047in}}%
\pgfpathlineto{\pgfqpoint{3.556492in}{2.092799in}}%
\pgfpathlineto{\pgfqpoint{3.557763in}{2.090085in}}%
\pgfpathlineto{\pgfqpoint{3.562517in}{2.079551in}}%
\pgfpathlineto{\pgfqpoint{3.565629in}{2.066303in}}%
\pgfpathlineto{\pgfqpoint{3.572331in}{2.053055in}}%
\pgfpathlineto{\pgfqpoint{3.575026in}{2.039807in}}%
\pgfpathlineto{\pgfqpoint{3.582259in}{2.026559in}}%
\pgfpathlineto{\pgfqpoint{3.583046in}{2.022457in}}%
\pgfpathlineto{\pgfqpoint{3.584832in}{2.013311in}}%
\pgfpathlineto{\pgfqpoint{3.592705in}{2.000063in}}%
\pgfpathlineto{\pgfqpoint{3.595219in}{1.986815in}}%
\pgfpathlineto{\pgfqpoint{3.604227in}{1.973567in}}%
\pgfpathlineto{\pgfqpoint{3.606269in}{1.960319in}}%
\pgfpathlineto{\pgfqpoint{3.608329in}{1.956716in}}%
\pgfpathlineto{\pgfqpoint{3.614449in}{1.947071in}}%
\pgfpathlineto{\pgfqpoint{3.615984in}{1.933823in}}%
\pgfpathlineto{\pgfqpoint{3.624044in}{1.920575in}}%
\pgfpathlineto{\pgfqpoint{3.627057in}{1.907327in}}%
\pgfpathlineto{\pgfqpoint{3.633612in}{1.897882in}}%
\pgfpathlineto{\pgfqpoint{3.636731in}{1.894079in}}%
\pgfpathlineto{\pgfqpoint{3.640235in}{1.880831in}}%
\pgfpathlineto{\pgfqpoint{3.648539in}{1.867583in}}%
\pgfpathlineto{\pgfqpoint{3.652882in}{1.854335in}}%
\pgfpathlineto{\pgfqpoint{3.658895in}{1.843076in}}%
\pgfpathlineto{\pgfqpoint{3.660216in}{1.841087in}}%
\pgfpathlineto{\pgfqpoint{3.665919in}{1.827839in}}%
\pgfpathlineto{\pgfqpoint{3.673342in}{1.814591in}}%
\pgfpathlineto{\pgfqpoint{3.679644in}{1.801343in}}%
\pgfpathlineto{\pgfqpoint{3.684179in}{1.791794in}}%
\pgfpathlineto{\pgfqpoint{3.686463in}{1.788095in}}%
\pgfpathlineto{\pgfqpoint{3.696133in}{1.774847in}}%
\pgfpathlineto{\pgfqpoint{3.700555in}{1.761599in}}%
\pgfpathlineto{\pgfqpoint{3.708976in}{1.748351in}}%
\pgfpathlineto{\pgfqpoint{3.709462in}{1.746475in}}%
\pgfpathlineto{\pgfqpoint{3.713273in}{1.735103in}}%
\pgfpathlineto{\pgfqpoint{3.726097in}{1.721855in}}%
\pgfpathlineto{\pgfqpoint{3.730903in}{1.708607in}}%
\pgfpathlineto{\pgfqpoint{3.734745in}{1.705034in}}%
\pgfpathlineto{\pgfqpoint{3.748906in}{1.695359in}}%
\pgfpathlineto{\pgfqpoint{3.750254in}{1.682111in}}%
\pgfpathlineto{\pgfqpoint{3.760028in}{1.673572in}}%
\pgfpathlineto{\pgfqpoint{3.768613in}{1.668863in}}%
\pgfpathlineto{\pgfqpoint{3.769072in}{1.655615in}}%
\pgfpathlineto{\pgfqpoint{3.785311in}{1.644307in}}%
\pgfpathlineto{\pgfqpoint{3.790891in}{1.642367in}}%
\pgfpathlineto{\pgfqpoint{3.790775in}{1.629119in}}%
\pgfpathlineto{\pgfqpoint{3.810594in}{1.619850in}}%
\pgfpathclose%
\pgfusepath{fill}%
\end{pgfscope}%
\begin{pgfscope}%
\pgfpathrectangle{\pgfqpoint{2.470586in}{0.476543in}}{\pgfqpoint{1.820388in}{1.656000in}}%
\pgfusepath{clip}%
\pgfsetbuttcap%
\pgfsetroundjoin%
\definecolor{currentfill}{rgb}{0.514110,0.725798,0.845444}%
\pgfsetfillcolor{currentfill}%
\pgfsetlinewidth{0.000000pt}%
\definecolor{currentstroke}{rgb}{0.000000,0.000000,0.000000}%
\pgfsetstrokecolor{currentstroke}%
\pgfsetdash{}{0pt}%
\pgfpathmoveto{\pgfqpoint{3.254365in}{0.754097in}}%
\pgfpathlineto{\pgfqpoint{3.279648in}{0.753660in}}%
\pgfpathlineto{\pgfqpoint{3.304931in}{0.753488in}}%
\pgfpathlineto{\pgfqpoint{3.330214in}{0.753593in}}%
\pgfpathlineto{\pgfqpoint{3.355497in}{0.753995in}}%
\pgfpathlineto{\pgfqpoint{3.380781in}{0.754722in}}%
\pgfpathlineto{\pgfqpoint{3.381498in}{0.754751in}}%
\pgfpathlineto{\pgfqpoint{3.406064in}{0.755397in}}%
\pgfpathlineto{\pgfqpoint{3.431347in}{0.756232in}}%
\pgfpathlineto{\pgfqpoint{3.456630in}{0.757239in}}%
\pgfpathlineto{\pgfqpoint{3.481913in}{0.758423in}}%
\pgfpathlineto{\pgfqpoint{3.507196in}{0.759789in}}%
\pgfpathlineto{\pgfqpoint{3.532480in}{0.761349in}}%
\pgfpathlineto{\pgfqpoint{3.557763in}{0.763117in}}%
\pgfpathlineto{\pgfqpoint{3.583046in}{0.765110in}}%
\pgfpathlineto{\pgfqpoint{3.608329in}{0.767350in}}%
\pgfpathlineto{\pgfqpoint{3.614958in}{0.767999in}}%
\pgfpathlineto{\pgfqpoint{3.633612in}{0.775497in}}%
\pgfpathlineto{\pgfqpoint{3.644760in}{0.781247in}}%
\pgfpathlineto{\pgfqpoint{3.658895in}{0.783041in}}%
\pgfpathlineto{\pgfqpoint{3.684179in}{0.786665in}}%
\pgfpathlineto{\pgfqpoint{3.709462in}{0.790857in}}%
\pgfpathlineto{\pgfqpoint{3.728607in}{0.794495in}}%
\pgfpathlineto{\pgfqpoint{3.729296in}{0.807743in}}%
\pgfpathlineto{\pgfqpoint{3.734745in}{0.809131in}}%
\pgfpathlineto{\pgfqpoint{3.760028in}{0.816528in}}%
\pgfpathlineto{\pgfqpoint{3.772964in}{0.820991in}}%
\pgfpathlineto{\pgfqpoint{3.763389in}{0.834239in}}%
\pgfpathlineto{\pgfqpoint{3.785311in}{0.845740in}}%
\pgfpathlineto{\pgfqpoint{3.788153in}{0.847487in}}%
\pgfpathlineto{\pgfqpoint{3.785311in}{0.853467in}}%
\pgfpathlineto{\pgfqpoint{3.781958in}{0.860735in}}%
\pgfpathlineto{\pgfqpoint{3.785311in}{0.863510in}}%
\pgfpathlineto{\pgfqpoint{3.796319in}{0.873983in}}%
\pgfpathlineto{\pgfqpoint{3.798891in}{0.887231in}}%
\pgfpathlineto{\pgfqpoint{3.800580in}{0.900479in}}%
\pgfpathlineto{\pgfqpoint{3.798243in}{0.913727in}}%
\pgfpathlineto{\pgfqpoint{3.792328in}{0.926975in}}%
\pgfpathlineto{\pgfqpoint{3.795008in}{0.940223in}}%
\pgfpathlineto{\pgfqpoint{3.785338in}{0.953471in}}%
\pgfpathlineto{\pgfqpoint{3.789690in}{0.966719in}}%
\pgfpathlineto{\pgfqpoint{3.785311in}{0.971146in}}%
\pgfpathlineto{\pgfqpoint{3.777203in}{0.979967in}}%
\pgfpathlineto{\pgfqpoint{3.779637in}{0.993215in}}%
\pgfpathlineto{\pgfqpoint{3.766337in}{1.006463in}}%
\pgfpathlineto{\pgfqpoint{3.768466in}{1.019711in}}%
\pgfpathlineto{\pgfqpoint{3.760028in}{1.027183in}}%
\pgfpathlineto{\pgfqpoint{3.753897in}{1.032959in}}%
\pgfpathlineto{\pgfqpoint{3.756008in}{1.046207in}}%
\pgfpathlineto{\pgfqpoint{3.745743in}{1.059455in}}%
\pgfpathlineto{\pgfqpoint{3.746509in}{1.072703in}}%
\pgfpathlineto{\pgfqpoint{3.736794in}{1.085951in}}%
\pgfpathlineto{\pgfqpoint{3.734745in}{1.093640in}}%
\pgfpathlineto{\pgfqpoint{3.733442in}{1.099199in}}%
\pgfpathlineto{\pgfqpoint{3.723517in}{1.112447in}}%
\pgfpathlineto{\pgfqpoint{3.719104in}{1.125695in}}%
\pgfpathlineto{\pgfqpoint{3.713225in}{1.138943in}}%
\pgfpathlineto{\pgfqpoint{3.709462in}{1.146819in}}%
\pgfpathlineto{\pgfqpoint{3.707197in}{1.152191in}}%
\pgfpathlineto{\pgfqpoint{3.702798in}{1.165439in}}%
\pgfpathlineto{\pgfqpoint{3.695671in}{1.178687in}}%
\pgfpathlineto{\pgfqpoint{3.692143in}{1.191935in}}%
\pgfpathlineto{\pgfqpoint{3.684179in}{1.203973in}}%
\pgfpathlineto{\pgfqpoint{3.683470in}{1.205183in}}%
\pgfpathlineto{\pgfqpoint{3.681262in}{1.218431in}}%
\pgfpathlineto{\pgfqpoint{3.671924in}{1.231679in}}%
\pgfpathlineto{\pgfqpoint{3.671029in}{1.244927in}}%
\pgfpathlineto{\pgfqpoint{3.662287in}{1.258175in}}%
\pgfpathlineto{\pgfqpoint{3.660715in}{1.271423in}}%
\pgfpathlineto{\pgfqpoint{3.658895in}{1.273713in}}%
\pgfpathlineto{\pgfqpoint{3.651121in}{1.284671in}}%
\pgfpathlineto{\pgfqpoint{3.648898in}{1.297919in}}%
\pgfpathlineto{\pgfqpoint{3.639465in}{1.311167in}}%
\pgfpathlineto{\pgfqpoint{3.638001in}{1.324415in}}%
\pgfpathlineto{\pgfqpoint{3.633612in}{1.332201in}}%
\pgfpathlineto{\pgfqpoint{3.630863in}{1.337663in}}%
\pgfpathlineto{\pgfqpoint{3.628523in}{1.350911in}}%
\pgfpathlineto{\pgfqpoint{3.621802in}{1.364159in}}%
\pgfpathlineto{\pgfqpoint{3.618640in}{1.377407in}}%
\pgfpathlineto{\pgfqpoint{3.612275in}{1.390655in}}%
\pgfpathlineto{\pgfqpoint{3.608329in}{1.401693in}}%
\pgfpathlineto{\pgfqpoint{3.607666in}{1.403903in}}%
\pgfpathlineto{\pgfqpoint{3.603016in}{1.417151in}}%
\pgfpathlineto{\pgfqpoint{3.597966in}{1.430399in}}%
\pgfpathlineto{\pgfqpoint{3.594498in}{1.443647in}}%
\pgfpathlineto{\pgfqpoint{3.588841in}{1.456895in}}%
\pgfpathlineto{\pgfqpoint{3.585136in}{1.470143in}}%
\pgfpathlineto{\pgfqpoint{3.583046in}{1.473750in}}%
\pgfpathlineto{\pgfqpoint{3.578306in}{1.483391in}}%
\pgfpathlineto{\pgfqpoint{3.575947in}{1.496639in}}%
\pgfpathlineto{\pgfqpoint{3.569463in}{1.509887in}}%
\pgfpathlineto{\pgfqpoint{3.568218in}{1.523135in}}%
\pgfpathlineto{\pgfqpoint{3.561559in}{1.536383in}}%
\pgfpathlineto{\pgfqpoint{3.559421in}{1.549631in}}%
\pgfpathlineto{\pgfqpoint{3.557763in}{1.552403in}}%
\pgfpathlineto{\pgfqpoint{3.552392in}{1.562879in}}%
\pgfpathlineto{\pgfqpoint{3.550693in}{1.576127in}}%
\pgfpathlineto{\pgfqpoint{3.543737in}{1.589375in}}%
\pgfpathlineto{\pgfqpoint{3.541611in}{1.602623in}}%
\pgfpathlineto{\pgfqpoint{3.534894in}{1.615871in}}%
\pgfpathlineto{\pgfqpoint{3.532480in}{1.628964in}}%
\pgfpathlineto{\pgfqpoint{3.532456in}{1.629119in}}%
\pgfpathlineto{\pgfqpoint{3.527307in}{1.642367in}}%
\pgfpathlineto{\pgfqpoint{3.524753in}{1.655615in}}%
\pgfpathlineto{\pgfqpoint{3.519202in}{1.668863in}}%
\pgfpathlineto{\pgfqpoint{3.515466in}{1.682111in}}%
\pgfpathlineto{\pgfqpoint{3.510892in}{1.695359in}}%
\pgfpathlineto{\pgfqpoint{3.507196in}{1.707314in}}%
\pgfpathlineto{\pgfqpoint{3.506859in}{1.708607in}}%
\pgfpathlineto{\pgfqpoint{3.503575in}{1.721855in}}%
\pgfpathlineto{\pgfqpoint{3.499113in}{1.735103in}}%
\pgfpathlineto{\pgfqpoint{3.495284in}{1.748351in}}%
\pgfpathlineto{\pgfqpoint{3.489516in}{1.761599in}}%
\pgfpathlineto{\pgfqpoint{3.486564in}{1.774847in}}%
\pgfpathlineto{\pgfqpoint{3.481913in}{1.784351in}}%
\pgfpathlineto{\pgfqpoint{3.480286in}{1.788095in}}%
\pgfpathlineto{\pgfqpoint{3.478125in}{1.801343in}}%
\pgfpathlineto{\pgfqpoint{3.472203in}{1.814591in}}%
\pgfpathlineto{\pgfqpoint{3.470290in}{1.827839in}}%
\pgfpathlineto{\pgfqpoint{3.463804in}{1.841087in}}%
\pgfpathlineto{\pgfqpoint{3.461671in}{1.854335in}}%
\pgfpathlineto{\pgfqpoint{3.456630in}{1.862959in}}%
\pgfpathlineto{\pgfqpoint{3.454054in}{1.867583in}}%
\pgfpathlineto{\pgfqpoint{3.451923in}{1.880831in}}%
\pgfpathlineto{\pgfqpoint{3.445215in}{1.894079in}}%
\pgfpathlineto{\pgfqpoint{3.443145in}{1.907327in}}%
\pgfpathlineto{\pgfqpoint{3.436011in}{1.920575in}}%
\pgfpathlineto{\pgfqpoint{3.432782in}{1.933823in}}%
\pgfpathlineto{\pgfqpoint{3.431347in}{1.936195in}}%
\pgfpathlineto{\pgfqpoint{3.424436in}{1.947071in}}%
\pgfpathlineto{\pgfqpoint{3.420950in}{1.960319in}}%
\pgfpathlineto{\pgfqpoint{3.413837in}{1.973567in}}%
\pgfpathlineto{\pgfqpoint{3.409809in}{1.986815in}}%
\pgfpathlineto{\pgfqpoint{3.406064in}{1.994093in}}%
\pgfpathlineto{\pgfqpoint{3.402271in}{2.000063in}}%
\pgfpathlineto{\pgfqpoint{3.396583in}{2.013311in}}%
\pgfpathlineto{\pgfqpoint{3.387962in}{2.026559in}}%
\pgfpathlineto{\pgfqpoint{3.381002in}{2.039807in}}%
\pgfpathlineto{\pgfqpoint{3.380781in}{2.040159in}}%
\pgfpathlineto{\pgfqpoint{3.367101in}{2.053055in}}%
\pgfpathlineto{\pgfqpoint{3.355497in}{2.064313in}}%
\pgfpathlineto{\pgfqpoint{3.346916in}{2.066303in}}%
\pgfpathlineto{\pgfqpoint{3.330214in}{2.069900in}}%
\pgfpathlineto{\pgfqpoint{3.323808in}{2.066303in}}%
\pgfpathlineto{\pgfqpoint{3.308156in}{2.053055in}}%
\pgfpathlineto{\pgfqpoint{3.304931in}{2.051431in}}%
\pgfpathlineto{\pgfqpoint{3.293512in}{2.039807in}}%
\pgfpathlineto{\pgfqpoint{3.289018in}{2.026559in}}%
\pgfpathlineto{\pgfqpoint{3.279648in}{2.017475in}}%
\pgfpathlineto{\pgfqpoint{3.276424in}{2.013311in}}%
\pgfpathlineto{\pgfqpoint{3.273340in}{2.000063in}}%
\pgfpathlineto{\pgfqpoint{3.264242in}{1.986815in}}%
\pgfpathlineto{\pgfqpoint{3.261610in}{1.973567in}}%
\pgfpathlineto{\pgfqpoint{3.254365in}{1.963119in}}%
\pgfpathlineto{\pgfqpoint{3.252626in}{1.960319in}}%
\pgfpathlineto{\pgfqpoint{3.248705in}{1.947071in}}%
\pgfpathlineto{\pgfqpoint{3.241345in}{1.933823in}}%
\pgfpathlineto{\pgfqpoint{3.237367in}{1.920575in}}%
\pgfpathlineto{\pgfqpoint{3.232557in}{1.907327in}}%
\pgfpathlineto{\pgfqpoint{3.229082in}{1.896302in}}%
\pgfpathlineto{\pgfqpoint{3.228332in}{1.894079in}}%
\pgfpathlineto{\pgfqpoint{3.223832in}{1.880831in}}%
\pgfpathlineto{\pgfqpoint{3.218375in}{1.867583in}}%
\pgfpathlineto{\pgfqpoint{3.215247in}{1.854335in}}%
\pgfpathlineto{\pgfqpoint{3.208970in}{1.841087in}}%
\pgfpathlineto{\pgfqpoint{3.206354in}{1.827839in}}%
\pgfpathlineto{\pgfqpoint{3.203798in}{1.822734in}}%
\pgfpathlineto{\pgfqpoint{3.199353in}{1.814591in}}%
\pgfpathlineto{\pgfqpoint{3.197449in}{1.801343in}}%
\pgfpathlineto{\pgfqpoint{3.191000in}{1.788095in}}%
\pgfpathlineto{\pgfqpoint{3.189346in}{1.774847in}}%
\pgfpathlineto{\pgfqpoint{3.182409in}{1.761599in}}%
\pgfpathlineto{\pgfqpoint{3.180317in}{1.748351in}}%
\pgfpathlineto{\pgfqpoint{3.178515in}{1.744155in}}%
\pgfpathlineto{\pgfqpoint{3.174167in}{1.735103in}}%
\pgfpathlineto{\pgfqpoint{3.172012in}{1.721855in}}%
\pgfpathlineto{\pgfqpoint{3.166585in}{1.708607in}}%
\pgfpathlineto{\pgfqpoint{3.163564in}{1.695359in}}%
\pgfpathlineto{\pgfqpoint{3.158443in}{1.682111in}}%
\pgfpathlineto{\pgfqpoint{3.154540in}{1.668863in}}%
\pgfpathlineto{\pgfqpoint{3.153232in}{1.664457in}}%
\pgfpathlineto{\pgfqpoint{3.150203in}{1.655615in}}%
\pgfpathlineto{\pgfqpoint{3.144592in}{1.642367in}}%
\pgfpathlineto{\pgfqpoint{3.141274in}{1.629119in}}%
\pgfpathlineto{\pgfqpoint{3.136062in}{1.615871in}}%
\pgfpathlineto{\pgfqpoint{3.134276in}{1.602623in}}%
\pgfpathlineto{\pgfqpoint{3.128467in}{1.589375in}}%
\pgfpathlineto{\pgfqpoint{3.127949in}{1.585652in}}%
\pgfpathlineto{\pgfqpoint{3.126357in}{1.576127in}}%
\pgfpathlineto{\pgfqpoint{3.118203in}{1.562879in}}%
\pgfpathlineto{\pgfqpoint{3.116628in}{1.549631in}}%
\pgfpathlineto{\pgfqpoint{3.109932in}{1.536383in}}%
\pgfpathlineto{\pgfqpoint{3.108508in}{1.523135in}}%
\pgfpathlineto{\pgfqpoint{3.102666in}{1.512005in}}%
\pgfpathlineto{\pgfqpoint{3.101389in}{1.509887in}}%
\pgfpathlineto{\pgfqpoint{3.099062in}{1.496639in}}%
\pgfpathlineto{\pgfqpoint{3.091531in}{1.483391in}}%
\pgfpathlineto{\pgfqpoint{3.088925in}{1.470143in}}%
\pgfpathlineto{\pgfqpoint{3.082959in}{1.456895in}}%
\pgfpathlineto{\pgfqpoint{3.079619in}{1.443647in}}%
\pgfpathlineto{\pgfqpoint{3.077382in}{1.437342in}}%
\pgfpathlineto{\pgfqpoint{3.074575in}{1.430399in}}%
\pgfpathlineto{\pgfqpoint{3.070458in}{1.417151in}}%
\pgfpathlineto{\pgfqpoint{3.066136in}{1.403903in}}%
\pgfpathlineto{\pgfqpoint{3.059914in}{1.390655in}}%
\pgfpathlineto{\pgfqpoint{3.055981in}{1.377407in}}%
\pgfpathlineto{\pgfqpoint{3.052099in}{1.370406in}}%
\pgfpathlineto{\pgfqpoint{3.048076in}{1.364159in}}%
\pgfpathlineto{\pgfqpoint{3.046019in}{1.350911in}}%
\pgfpathlineto{\pgfqpoint{3.038377in}{1.337663in}}%
\pgfpathlineto{\pgfqpoint{3.036717in}{1.324415in}}%
\pgfpathlineto{\pgfqpoint{3.027742in}{1.311167in}}%
\pgfpathlineto{\pgfqpoint{3.026816in}{1.299932in}}%
\pgfpathlineto{\pgfqpoint{3.026622in}{1.297919in}}%
\pgfpathlineto{\pgfqpoint{3.016511in}{1.284671in}}%
\pgfpathlineto{\pgfqpoint{3.015845in}{1.271423in}}%
\pgfpathlineto{\pgfqpoint{3.006342in}{1.258175in}}%
\pgfpathlineto{\pgfqpoint{3.005195in}{1.244927in}}%
\pgfpathlineto{\pgfqpoint{3.001533in}{1.239419in}}%
\pgfpathlineto{\pgfqpoint{2.995932in}{1.231679in}}%
\pgfpathlineto{\pgfqpoint{2.994185in}{1.218431in}}%
\pgfpathlineto{\pgfqpoint{2.984342in}{1.205183in}}%
\pgfpathlineto{\pgfqpoint{2.979914in}{1.191935in}}%
\pgfpathlineto{\pgfqpoint{2.976250in}{1.185896in}}%
\pgfpathlineto{\pgfqpoint{2.971550in}{1.178687in}}%
\pgfpathlineto{\pgfqpoint{2.966787in}{1.165439in}}%
\pgfpathlineto{\pgfqpoint{2.962419in}{1.152191in}}%
\pgfpathlineto{\pgfqpoint{2.956925in}{1.138943in}}%
\pgfpathlineto{\pgfqpoint{2.952277in}{1.125695in}}%
\pgfpathlineto{\pgfqpoint{2.950967in}{1.123642in}}%
\pgfpathlineto{\pgfqpoint{2.942912in}{1.112447in}}%
\pgfpathlineto{\pgfqpoint{2.940382in}{1.099199in}}%
\pgfpathlineto{\pgfqpoint{2.929358in}{1.085951in}}%
\pgfpathlineto{\pgfqpoint{2.928924in}{1.072703in}}%
\pgfpathlineto{\pgfqpoint{2.925683in}{1.069185in}}%
\pgfpathlineto{\pgfqpoint{2.915909in}{1.059455in}}%
\pgfpathlineto{\pgfqpoint{2.917522in}{1.046207in}}%
\pgfpathlineto{\pgfqpoint{2.903794in}{1.032959in}}%
\pgfpathlineto{\pgfqpoint{2.906669in}{1.019711in}}%
\pgfpathlineto{\pgfqpoint{2.900400in}{1.014681in}}%
\pgfpathlineto{\pgfqpoint{2.889458in}{1.006463in}}%
\pgfpathlineto{\pgfqpoint{2.893039in}{0.993215in}}%
\pgfpathlineto{\pgfqpoint{2.878226in}{0.979967in}}%
\pgfpathlineto{\pgfqpoint{2.885819in}{0.966719in}}%
\pgfpathlineto{\pgfqpoint{2.875117in}{0.954898in}}%
\pgfpathlineto{\pgfqpoint{2.873760in}{0.953471in}}%
\pgfpathlineto{\pgfqpoint{2.875117in}{0.950444in}}%
\pgfpathlineto{\pgfqpoint{2.880064in}{0.940223in}}%
\pgfpathlineto{\pgfqpoint{2.875117in}{0.936056in}}%
\pgfpathlineto{\pgfqpoint{2.863918in}{0.926975in}}%
\pgfpathlineto{\pgfqpoint{2.868678in}{0.913727in}}%
\pgfpathlineto{\pgfqpoint{2.857472in}{0.900479in}}%
\pgfpathlineto{\pgfqpoint{2.864724in}{0.887231in}}%
\pgfpathlineto{\pgfqpoint{2.859296in}{0.873983in}}%
\pgfpathlineto{\pgfqpoint{2.866618in}{0.860735in}}%
\pgfpathlineto{\pgfqpoint{2.869699in}{0.847487in}}%
\pgfpathlineto{\pgfqpoint{2.875117in}{0.839985in}}%
\pgfpathlineto{\pgfqpoint{2.880053in}{0.834239in}}%
\pgfpathlineto{\pgfqpoint{2.893358in}{0.820991in}}%
\pgfpathlineto{\pgfqpoint{2.900400in}{0.813725in}}%
\pgfpathlineto{\pgfqpoint{2.907675in}{0.807743in}}%
\pgfpathlineto{\pgfqpoint{2.925683in}{0.800305in}}%
\pgfpathlineto{\pgfqpoint{2.940855in}{0.794495in}}%
\pgfpathlineto{\pgfqpoint{2.950967in}{0.790262in}}%
\pgfpathlineto{\pgfqpoint{2.976250in}{0.782124in}}%
\pgfpathlineto{\pgfqpoint{2.979664in}{0.781247in}}%
\pgfpathlineto{\pgfqpoint{3.001533in}{0.777720in}}%
\pgfpathlineto{\pgfqpoint{3.026816in}{0.773974in}}%
\pgfpathlineto{\pgfqpoint{3.052099in}{0.770535in}}%
\pgfpathlineto{\pgfqpoint{3.072361in}{0.767999in}}%
\pgfpathlineto{\pgfqpoint{3.077382in}{0.767129in}}%
\pgfpathlineto{\pgfqpoint{3.102666in}{0.763484in}}%
\pgfpathlineto{\pgfqpoint{3.127949in}{0.760675in}}%
\pgfpathlineto{\pgfqpoint{3.153232in}{0.758515in}}%
\pgfpathlineto{\pgfqpoint{3.178515in}{0.756874in}}%
\pgfpathlineto{\pgfqpoint{3.203798in}{0.755654in}}%
\pgfpathlineto{\pgfqpoint{3.229082in}{0.754784in}}%
\pgfpathlineto{\pgfqpoint{3.230460in}{0.754751in}}%
\pgfpathlineto{\pgfqpoint{3.254365in}{0.754097in}}%
\pgfpathclose%
\pgfpathmoveto{\pgfqpoint{3.095288in}{0.794495in}}%
\pgfpathlineto{\pgfqpoint{3.077382in}{0.798796in}}%
\pgfpathlineto{\pgfqpoint{3.052099in}{0.805497in}}%
\pgfpathlineto{\pgfqpoint{3.044347in}{0.807743in}}%
\pgfpathlineto{\pgfqpoint{3.026816in}{0.811803in}}%
\pgfpathlineto{\pgfqpoint{3.001533in}{0.819523in}}%
\pgfpathlineto{\pgfqpoint{2.997618in}{0.820991in}}%
\pgfpathlineto{\pgfqpoint{2.976745in}{0.834239in}}%
\pgfpathlineto{\pgfqpoint{2.976250in}{0.834433in}}%
\pgfpathlineto{\pgfqpoint{2.950967in}{0.846671in}}%
\pgfpathlineto{\pgfqpoint{2.949597in}{0.847487in}}%
\pgfpathlineto{\pgfqpoint{2.943906in}{0.860735in}}%
\pgfpathlineto{\pgfqpoint{2.925683in}{0.873480in}}%
\pgfpathlineto{\pgfqpoint{2.925082in}{0.873983in}}%
\pgfpathlineto{\pgfqpoint{2.925683in}{0.875993in}}%
\pgfpathlineto{\pgfqpoint{2.928985in}{0.887231in}}%
\pgfpathlineto{\pgfqpoint{2.925683in}{0.889651in}}%
\pgfpathlineto{\pgfqpoint{2.913077in}{0.900479in}}%
\pgfpathlineto{\pgfqpoint{2.922863in}{0.913727in}}%
\pgfpathlineto{\pgfqpoint{2.912217in}{0.926975in}}%
\pgfpathlineto{\pgfqpoint{2.925683in}{0.939091in}}%
\pgfpathlineto{\pgfqpoint{2.926898in}{0.940223in}}%
\pgfpathlineto{\pgfqpoint{2.925683in}{0.941557in}}%
\pgfpathlineto{\pgfqpoint{2.915891in}{0.953471in}}%
\pgfpathlineto{\pgfqpoint{2.925683in}{0.965428in}}%
\pgfpathlineto{\pgfqpoint{2.926692in}{0.966719in}}%
\pgfpathlineto{\pgfqpoint{2.925683in}{0.967872in}}%
\pgfpathlineto{\pgfqpoint{2.915819in}{0.979967in}}%
\pgfpathlineto{\pgfqpoint{2.925683in}{0.989340in}}%
\pgfpathlineto{\pgfqpoint{2.929574in}{0.993215in}}%
\pgfpathlineto{\pgfqpoint{2.925683in}{1.002316in}}%
\pgfpathlineto{\pgfqpoint{2.923969in}{1.006463in}}%
\pgfpathlineto{\pgfqpoint{2.925683in}{1.007855in}}%
\pgfpathlineto{\pgfqpoint{2.939361in}{1.019711in}}%
\pgfpathlineto{\pgfqpoint{2.934919in}{1.032959in}}%
\pgfpathlineto{\pgfqpoint{2.947357in}{1.046207in}}%
\pgfpathlineto{\pgfqpoint{2.944843in}{1.059455in}}%
\pgfpathlineto{\pgfqpoint{2.950967in}{1.066302in}}%
\pgfpathlineto{\pgfqpoint{2.956218in}{1.072703in}}%
\pgfpathlineto{\pgfqpoint{2.956040in}{1.085951in}}%
\pgfpathlineto{\pgfqpoint{2.965336in}{1.099199in}}%
\pgfpathlineto{\pgfqpoint{2.967545in}{1.112447in}}%
\pgfpathlineto{\pgfqpoint{2.975328in}{1.125695in}}%
\pgfpathlineto{\pgfqpoint{2.976250in}{1.128244in}}%
\pgfpathlineto{\pgfqpoint{2.979865in}{1.138943in}}%
\pgfpathlineto{\pgfqpoint{2.983983in}{1.152191in}}%
\pgfpathlineto{\pgfqpoint{2.988312in}{1.165439in}}%
\pgfpathlineto{\pgfqpoint{2.991749in}{1.178687in}}%
\pgfpathlineto{\pgfqpoint{3.000294in}{1.191935in}}%
\pgfpathlineto{\pgfqpoint{3.001533in}{1.196357in}}%
\pgfpathlineto{\pgfqpoint{3.003689in}{1.205183in}}%
\pgfpathlineto{\pgfqpoint{3.013090in}{1.218431in}}%
\pgfpathlineto{\pgfqpoint{3.013890in}{1.231679in}}%
\pgfpathlineto{\pgfqpoint{3.023089in}{1.244927in}}%
\pgfpathlineto{\pgfqpoint{3.023621in}{1.258175in}}%
\pgfpathlineto{\pgfqpoint{3.026816in}{1.262463in}}%
\pgfpathlineto{\pgfqpoint{3.032936in}{1.271423in}}%
\pgfpathlineto{\pgfqpoint{3.033032in}{1.284671in}}%
\pgfpathlineto{\pgfqpoint{3.042198in}{1.297919in}}%
\pgfpathlineto{\pgfqpoint{3.043071in}{1.311167in}}%
\pgfpathlineto{\pgfqpoint{3.052099in}{1.324022in}}%
\pgfpathlineto{\pgfqpoint{3.052349in}{1.324415in}}%
\pgfpathlineto{\pgfqpoint{3.053785in}{1.337663in}}%
\pgfpathlineto{\pgfqpoint{3.060605in}{1.350911in}}%
\pgfpathlineto{\pgfqpoint{3.062512in}{1.364159in}}%
\pgfpathlineto{\pgfqpoint{3.069871in}{1.377407in}}%
\pgfpathlineto{\pgfqpoint{3.074154in}{1.390655in}}%
\pgfpathlineto{\pgfqpoint{3.077382in}{1.397597in}}%
\pgfpathlineto{\pgfqpoint{3.079901in}{1.403903in}}%
\pgfpathlineto{\pgfqpoint{3.083950in}{1.417151in}}%
\pgfpathlineto{\pgfqpoint{3.087263in}{1.430399in}}%
\pgfpathlineto{\pgfqpoint{3.092363in}{1.443647in}}%
\pgfpathlineto{\pgfqpoint{3.095479in}{1.456895in}}%
\pgfpathlineto{\pgfqpoint{3.102051in}{1.470143in}}%
\pgfpathlineto{\pgfqpoint{3.102666in}{1.473512in}}%
\pgfpathlineto{\pgfqpoint{3.104181in}{1.483391in}}%
\pgfpathlineto{\pgfqpoint{3.111148in}{1.496639in}}%
\pgfpathlineto{\pgfqpoint{3.112959in}{1.509887in}}%
\pgfpathlineto{\pgfqpoint{3.120387in}{1.523135in}}%
\pgfpathlineto{\pgfqpoint{3.121676in}{1.536383in}}%
\pgfpathlineto{\pgfqpoint{3.127949in}{1.547935in}}%
\pgfpathlineto{\pgfqpoint{3.128769in}{1.549631in}}%
\pgfpathlineto{\pgfqpoint{3.130107in}{1.562879in}}%
\pgfpathlineto{\pgfqpoint{3.137318in}{1.576127in}}%
\pgfpathlineto{\pgfqpoint{3.139231in}{1.589375in}}%
\pgfpathlineto{\pgfqpoint{3.145236in}{1.602623in}}%
\pgfpathlineto{\pgfqpoint{3.147234in}{1.615871in}}%
\pgfpathlineto{\pgfqpoint{3.152719in}{1.629119in}}%
\pgfpathlineto{\pgfqpoint{3.153232in}{1.630981in}}%
\pgfpathlineto{\pgfqpoint{3.156011in}{1.642367in}}%
\pgfpathlineto{\pgfqpoint{3.160820in}{1.655615in}}%
\pgfpathlineto{\pgfqpoint{3.165066in}{1.668863in}}%
\pgfpathlineto{\pgfqpoint{3.168931in}{1.682111in}}%
\pgfpathlineto{\pgfqpoint{3.174646in}{1.695359in}}%
\pgfpathlineto{\pgfqpoint{3.177647in}{1.708607in}}%
\pgfpathlineto{\pgfqpoint{3.178515in}{1.710589in}}%
\pgfpathlineto{\pgfqpoint{3.183008in}{1.721855in}}%
\pgfpathlineto{\pgfqpoint{3.184750in}{1.735103in}}%
\pgfpathlineto{\pgfqpoint{3.190766in}{1.748351in}}%
\pgfpathlineto{\pgfqpoint{3.192832in}{1.761599in}}%
\pgfpathlineto{\pgfqpoint{3.200353in}{1.774847in}}%
\pgfpathlineto{\pgfqpoint{3.202048in}{1.788095in}}%
\pgfpathlineto{\pgfqpoint{3.203798in}{1.791509in}}%
\pgfpathlineto{\pgfqpoint{3.208531in}{1.801343in}}%
\pgfpathlineto{\pgfqpoint{3.210404in}{1.814591in}}%
\pgfpathlineto{\pgfqpoint{3.217256in}{1.827839in}}%
\pgfpathlineto{\pgfqpoint{3.220158in}{1.841087in}}%
\pgfpathlineto{\pgfqpoint{3.226680in}{1.854335in}}%
\pgfpathlineto{\pgfqpoint{3.229082in}{1.863367in}}%
\pgfpathlineto{\pgfqpoint{3.230206in}{1.867583in}}%
\pgfpathlineto{\pgfqpoint{3.235585in}{1.880831in}}%
\pgfpathlineto{\pgfqpoint{3.240450in}{1.894079in}}%
\pgfpathlineto{\pgfqpoint{3.244575in}{1.907327in}}%
\pgfpathlineto{\pgfqpoint{3.250129in}{1.920575in}}%
\pgfpathlineto{\pgfqpoint{3.254119in}{1.933823in}}%
\pgfpathlineto{\pgfqpoint{3.254365in}{1.934241in}}%
\pgfpathlineto{\pgfqpoint{3.263068in}{1.947071in}}%
\pgfpathlineto{\pgfqpoint{3.267145in}{1.960319in}}%
\pgfpathlineto{\pgfqpoint{3.277125in}{1.973567in}}%
\pgfpathlineto{\pgfqpoint{3.279648in}{1.986298in}}%
\pgfpathlineto{\pgfqpoint{3.279782in}{1.986815in}}%
\pgfpathlineto{\pgfqpoint{3.292790in}{2.000063in}}%
\pgfpathlineto{\pgfqpoint{3.297159in}{2.013311in}}%
\pgfpathlineto{\pgfqpoint{3.304931in}{2.020660in}}%
\pgfpathlineto{\pgfqpoint{3.317681in}{2.026559in}}%
\pgfpathlineto{\pgfqpoint{3.329349in}{2.039807in}}%
\pgfpathlineto{\pgfqpoint{3.330214in}{2.040243in}}%
\pgfpathlineto{\pgfqpoint{3.332787in}{2.039807in}}%
\pgfpathlineto{\pgfqpoint{3.355497in}{2.033110in}}%
\pgfpathlineto{\pgfqpoint{3.360423in}{2.026559in}}%
\pgfpathlineto{\pgfqpoint{3.376537in}{2.013311in}}%
\pgfpathlineto{\pgfqpoint{3.380781in}{2.006893in}}%
\pgfpathlineto{\pgfqpoint{3.383691in}{2.000063in}}%
\pgfpathlineto{\pgfqpoint{3.393053in}{1.986815in}}%
\pgfpathlineto{\pgfqpoint{3.397599in}{1.973567in}}%
\pgfpathlineto{\pgfqpoint{3.406064in}{1.961891in}}%
\pgfpathlineto{\pgfqpoint{3.406969in}{1.960319in}}%
\pgfpathlineto{\pgfqpoint{3.410466in}{1.947071in}}%
\pgfpathlineto{\pgfqpoint{3.419695in}{1.933823in}}%
\pgfpathlineto{\pgfqpoint{3.422938in}{1.920575in}}%
\pgfpathlineto{\pgfqpoint{3.430830in}{1.907327in}}%
\pgfpathlineto{\pgfqpoint{3.431347in}{1.904122in}}%
\pgfpathlineto{\pgfqpoint{3.432945in}{1.894079in}}%
\pgfpathlineto{\pgfqpoint{3.440356in}{1.880831in}}%
\pgfpathlineto{\pgfqpoint{3.442544in}{1.867583in}}%
\pgfpathlineto{\pgfqpoint{3.450201in}{1.854335in}}%
\pgfpathlineto{\pgfqpoint{3.452255in}{1.841087in}}%
\pgfpathlineto{\pgfqpoint{3.456630in}{1.832231in}}%
\pgfpathlineto{\pgfqpoint{3.458888in}{1.827839in}}%
\pgfpathlineto{\pgfqpoint{3.460993in}{1.814591in}}%
\pgfpathlineto{\pgfqpoint{3.467374in}{1.801343in}}%
\pgfpathlineto{\pgfqpoint{3.469752in}{1.788095in}}%
\pgfpathlineto{\pgfqpoint{3.475648in}{1.774847in}}%
\pgfpathlineto{\pgfqpoint{3.478513in}{1.761599in}}%
\pgfpathlineto{\pgfqpoint{3.481913in}{1.753243in}}%
\pgfpathlineto{\pgfqpoint{3.484128in}{1.748351in}}%
\pgfpathlineto{\pgfqpoint{3.488383in}{1.735103in}}%
\pgfpathlineto{\pgfqpoint{3.493006in}{1.721855in}}%
\pgfpathlineto{\pgfqpoint{3.496615in}{1.708607in}}%
\pgfpathlineto{\pgfqpoint{3.500059in}{1.695359in}}%
\pgfpathlineto{\pgfqpoint{3.504399in}{1.682111in}}%
\pgfpathlineto{\pgfqpoint{3.507196in}{1.670712in}}%
\pgfpathlineto{\pgfqpoint{3.507725in}{1.668863in}}%
\pgfpathlineto{\pgfqpoint{3.513826in}{1.655615in}}%
\pgfpathlineto{\pgfqpoint{3.516345in}{1.642367in}}%
\pgfpathlineto{\pgfqpoint{3.521963in}{1.629119in}}%
\pgfpathlineto{\pgfqpoint{3.523943in}{1.615871in}}%
\pgfpathlineto{\pgfqpoint{3.530127in}{1.602623in}}%
\pgfpathlineto{\pgfqpoint{3.531883in}{1.589375in}}%
\pgfpathlineto{\pgfqpoint{3.532480in}{1.588198in}}%
\pgfpathlineto{\pgfqpoint{3.539303in}{1.576127in}}%
\pgfpathlineto{\pgfqpoint{3.540948in}{1.562879in}}%
\pgfpathlineto{\pgfqpoint{3.548095in}{1.549631in}}%
\pgfpathlineto{\pgfqpoint{3.549863in}{1.536383in}}%
\pgfpathlineto{\pgfqpoint{3.555847in}{1.523135in}}%
\pgfpathlineto{\pgfqpoint{3.556886in}{1.509887in}}%
\pgfpathlineto{\pgfqpoint{3.557763in}{1.507957in}}%
\pgfpathlineto{\pgfqpoint{3.563598in}{1.496639in}}%
\pgfpathlineto{\pgfqpoint{3.566125in}{1.483391in}}%
\pgfpathlineto{\pgfqpoint{3.572771in}{1.470143in}}%
\pgfpathlineto{\pgfqpoint{3.576081in}{1.456895in}}%
\pgfpathlineto{\pgfqpoint{3.581026in}{1.443647in}}%
\pgfpathlineto{\pgfqpoint{3.583046in}{1.435471in}}%
\pgfpathlineto{\pgfqpoint{3.584471in}{1.430399in}}%
\pgfpathlineto{\pgfqpoint{3.589576in}{1.417151in}}%
\pgfpathlineto{\pgfqpoint{3.594612in}{1.403903in}}%
\pgfpathlineto{\pgfqpoint{3.598384in}{1.390655in}}%
\pgfpathlineto{\pgfqpoint{3.604248in}{1.377407in}}%
\pgfpathlineto{\pgfqpoint{3.606797in}{1.364159in}}%
\pgfpathlineto{\pgfqpoint{3.608329in}{1.361036in}}%
\pgfpathlineto{\pgfqpoint{3.613815in}{1.350911in}}%
\pgfpathlineto{\pgfqpoint{3.615841in}{1.337663in}}%
\pgfpathlineto{\pgfqpoint{3.622835in}{1.324415in}}%
\pgfpathlineto{\pgfqpoint{3.623693in}{1.311167in}}%
\pgfpathlineto{\pgfqpoint{3.632512in}{1.297919in}}%
\pgfpathlineto{\pgfqpoint{3.633612in}{1.289474in}}%
\pgfpathlineto{\pgfqpoint{3.634330in}{1.284671in}}%
\pgfpathlineto{\pgfqpoint{3.644040in}{1.271423in}}%
\pgfpathlineto{\pgfqpoint{3.645033in}{1.258175in}}%
\pgfpathlineto{\pgfqpoint{3.653037in}{1.244927in}}%
\pgfpathlineto{\pgfqpoint{3.653518in}{1.231679in}}%
\pgfpathlineto{\pgfqpoint{3.658895in}{1.223331in}}%
\pgfpathlineto{\pgfqpoint{3.662385in}{1.218431in}}%
\pgfpathlineto{\pgfqpoint{3.664376in}{1.205183in}}%
\pgfpathlineto{\pgfqpoint{3.671894in}{1.191935in}}%
\pgfpathlineto{\pgfqpoint{3.675022in}{1.178687in}}%
\pgfpathlineto{\pgfqpoint{3.681192in}{1.165439in}}%
\pgfpathlineto{\pgfqpoint{3.684179in}{1.155833in}}%
\pgfpathlineto{\pgfqpoint{3.685406in}{1.152191in}}%
\pgfpathlineto{\pgfqpoint{3.690309in}{1.138943in}}%
\pgfpathlineto{\pgfqpoint{3.695878in}{1.125695in}}%
\pgfpathlineto{\pgfqpoint{3.698813in}{1.112447in}}%
\pgfpathlineto{\pgfqpoint{3.708340in}{1.099199in}}%
\pgfpathlineto{\pgfqpoint{3.709462in}{1.090398in}}%
\pgfpathlineto{\pgfqpoint{3.710100in}{1.085951in}}%
\pgfpathlineto{\pgfqpoint{3.719081in}{1.072703in}}%
\pgfpathlineto{\pgfqpoint{3.716530in}{1.059455in}}%
\pgfpathlineto{\pgfqpoint{3.726329in}{1.046207in}}%
\pgfpathlineto{\pgfqpoint{3.722385in}{1.032959in}}%
\pgfpathlineto{\pgfqpoint{3.734745in}{1.020493in}}%
\pgfpathlineto{\pgfqpoint{3.735568in}{1.019711in}}%
\pgfpathlineto{\pgfqpoint{3.734745in}{1.016893in}}%
\pgfpathlineto{\pgfqpoint{3.731621in}{1.006463in}}%
\pgfpathlineto{\pgfqpoint{3.734745in}{1.002709in}}%
\pgfpathlineto{\pgfqpoint{3.743205in}{0.993215in}}%
\pgfpathlineto{\pgfqpoint{3.738808in}{0.979967in}}%
\pgfpathlineto{\pgfqpoint{3.749161in}{0.966719in}}%
\pgfpathlineto{\pgfqpoint{3.742692in}{0.953471in}}%
\pgfpathlineto{\pgfqpoint{3.748374in}{0.940223in}}%
\pgfpathlineto{\pgfqpoint{3.743749in}{0.926975in}}%
\pgfpathlineto{\pgfqpoint{3.744436in}{0.913727in}}%
\pgfpathlineto{\pgfqpoint{3.744233in}{0.900479in}}%
\pgfpathlineto{\pgfqpoint{3.735356in}{0.887231in}}%
\pgfpathlineto{\pgfqpoint{3.734745in}{0.885713in}}%
\pgfpathlineto{\pgfqpoint{3.729746in}{0.873983in}}%
\pgfpathlineto{\pgfqpoint{3.709462in}{0.863072in}}%
\pgfpathlineto{\pgfqpoint{3.704279in}{0.860735in}}%
\pgfpathlineto{\pgfqpoint{3.707714in}{0.847487in}}%
\pgfpathlineto{\pgfqpoint{3.684179in}{0.839309in}}%
\pgfpathlineto{\pgfqpoint{3.666423in}{0.834239in}}%
\pgfpathlineto{\pgfqpoint{3.667725in}{0.820991in}}%
\pgfpathlineto{\pgfqpoint{3.658895in}{0.819024in}}%
\pgfpathlineto{\pgfqpoint{3.633612in}{0.814083in}}%
\pgfpathlineto{\pgfqpoint{3.608329in}{0.809913in}}%
\pgfpathlineto{\pgfqpoint{3.593091in}{0.807743in}}%
\pgfpathlineto{\pgfqpoint{3.583046in}{0.800263in}}%
\pgfpathlineto{\pgfqpoint{3.573543in}{0.794495in}}%
\pgfpathlineto{\pgfqpoint{3.557763in}{0.792841in}}%
\pgfpathlineto{\pgfqpoint{3.532480in}{0.790526in}}%
\pgfpathlineto{\pgfqpoint{3.507196in}{0.788535in}}%
\pgfpathlineto{\pgfqpoint{3.481913in}{0.786836in}}%
\pgfpathlineto{\pgfqpoint{3.456630in}{0.785403in}}%
\pgfpathlineto{\pgfqpoint{3.431347in}{0.784218in}}%
\pgfpathlineto{\pgfqpoint{3.406064in}{0.783267in}}%
\pgfpathlineto{\pgfqpoint{3.380781in}{0.782542in}}%
\pgfpathlineto{\pgfqpoint{3.355497in}{0.782040in}}%
\pgfpathlineto{\pgfqpoint{3.330214in}{0.781762in}}%
\pgfpathlineto{\pgfqpoint{3.304931in}{0.781716in}}%
\pgfpathlineto{\pgfqpoint{3.279648in}{0.781914in}}%
\pgfpathlineto{\pgfqpoint{3.254365in}{0.782379in}}%
\pgfpathlineto{\pgfqpoint{3.229082in}{0.783139in}}%
\pgfpathlineto{\pgfqpoint{3.203798in}{0.784234in}}%
\pgfpathlineto{\pgfqpoint{3.178515in}{0.785718in}}%
\pgfpathlineto{\pgfqpoint{3.153232in}{0.787665in}}%
\pgfpathlineto{\pgfqpoint{3.127949in}{0.790171in}}%
\pgfpathlineto{\pgfqpoint{3.102666in}{0.793371in}}%
\pgfpathlineto{\pgfqpoint{3.095288in}{0.794495in}}%
\pgfpathclose%
\pgfpathmoveto{\pgfqpoint{2.875117in}{1.730512in}}%
\pgfpathlineto{\pgfqpoint{2.900400in}{1.725909in}}%
\pgfpathlineto{\pgfqpoint{2.925683in}{1.733233in}}%
\pgfpathlineto{\pgfqpoint{2.928486in}{1.735103in}}%
\pgfpathlineto{\pgfqpoint{2.927103in}{1.748351in}}%
\pgfpathlineto{\pgfqpoint{2.944288in}{1.761599in}}%
\pgfpathlineto{\pgfqpoint{2.945606in}{1.774847in}}%
\pgfpathlineto{\pgfqpoint{2.950967in}{1.777748in}}%
\pgfpathlineto{\pgfqpoint{2.960170in}{1.788095in}}%
\pgfpathlineto{\pgfqpoint{2.965989in}{1.801343in}}%
\pgfpathlineto{\pgfqpoint{2.976250in}{1.813664in}}%
\pgfpathlineto{\pgfqpoint{2.976739in}{1.814591in}}%
\pgfpathlineto{\pgfqpoint{2.984099in}{1.827839in}}%
\pgfpathlineto{\pgfqpoint{2.989371in}{1.841087in}}%
\pgfpathlineto{\pgfqpoint{3.000958in}{1.854335in}}%
\pgfpathlineto{\pgfqpoint{3.001533in}{1.856183in}}%
\pgfpathlineto{\pgfqpoint{3.003967in}{1.867583in}}%
\pgfpathlineto{\pgfqpoint{3.013883in}{1.880831in}}%
\pgfpathlineto{\pgfqpoint{3.015972in}{1.894079in}}%
\pgfpathlineto{\pgfqpoint{3.026816in}{1.904274in}}%
\pgfpathlineto{\pgfqpoint{3.029353in}{1.907327in}}%
\pgfpathlineto{\pgfqpoint{3.031087in}{1.920575in}}%
\pgfpathlineto{\pgfqpoint{3.040874in}{1.933823in}}%
\pgfpathlineto{\pgfqpoint{3.041298in}{1.947071in}}%
\pgfpathlineto{\pgfqpoint{3.050958in}{1.960319in}}%
\pgfpathlineto{\pgfqpoint{3.052099in}{1.966357in}}%
\pgfpathlineto{\pgfqpoint{3.053265in}{1.973567in}}%
\pgfpathlineto{\pgfqpoint{3.061935in}{1.986815in}}%
\pgfpathlineto{\pgfqpoint{3.065370in}{2.000063in}}%
\pgfpathlineto{\pgfqpoint{3.072774in}{2.013311in}}%
\pgfpathlineto{\pgfqpoint{3.077382in}{2.024520in}}%
\pgfpathlineto{\pgfqpoint{3.078150in}{2.026559in}}%
\pgfpathlineto{\pgfqpoint{3.083410in}{2.039807in}}%
\pgfpathlineto{\pgfqpoint{3.088486in}{2.053055in}}%
\pgfpathlineto{\pgfqpoint{3.091996in}{2.066303in}}%
\pgfpathlineto{\pgfqpoint{3.098507in}{2.079551in}}%
\pgfpathlineto{\pgfqpoint{3.101705in}{2.092799in}}%
\pgfpathlineto{\pgfqpoint{3.102666in}{2.094300in}}%
\pgfpathlineto{\pgfqpoint{3.109452in}{2.106047in}}%
\pgfpathlineto{\pgfqpoint{3.111219in}{2.119295in}}%
\pgfpathlineto{\pgfqpoint{3.118216in}{2.132543in}}%
\pgfpathlineto{\pgfqpoint{3.104314in}{2.132543in}}%
\pgfpathlineto{\pgfqpoint{3.102666in}{2.129489in}}%
\pgfpathlineto{\pgfqpoint{3.096169in}{2.119295in}}%
\pgfpathlineto{\pgfqpoint{3.094254in}{2.106047in}}%
\pgfpathlineto{\pgfqpoint{3.086053in}{2.092799in}}%
\pgfpathlineto{\pgfqpoint{3.082777in}{2.079551in}}%
\pgfpathlineto{\pgfqpoint{3.077382in}{2.067597in}}%
\pgfpathlineto{\pgfqpoint{3.076650in}{2.066303in}}%
\pgfpathlineto{\pgfqpoint{3.071714in}{2.053055in}}%
\pgfpathlineto{\pgfqpoint{3.066292in}{2.039807in}}%
\pgfpathlineto{\pgfqpoint{3.059457in}{2.026559in}}%
\pgfpathlineto{\pgfqpoint{3.055347in}{2.013311in}}%
\pgfpathlineto{\pgfqpoint{3.052099in}{2.007818in}}%
\pgfpathlineto{\pgfqpoint{3.045360in}{2.000063in}}%
\pgfpathlineto{\pgfqpoint{3.042567in}{1.986815in}}%
\pgfpathlineto{\pgfqpoint{3.030922in}{1.973567in}}%
\pgfpathlineto{\pgfqpoint{3.030228in}{1.960319in}}%
\pgfpathlineto{\pgfqpoint{3.026816in}{1.955643in}}%
\pgfpathlineto{\pgfqpoint{3.016974in}{1.947071in}}%
\pgfpathlineto{\pgfqpoint{3.018857in}{1.933823in}}%
\pgfpathlineto{\pgfqpoint{3.005344in}{1.920575in}}%
\pgfpathlineto{\pgfqpoint{3.004441in}{1.907327in}}%
\pgfpathlineto{\pgfqpoint{3.001533in}{1.904476in}}%
\pgfpathlineto{\pgfqpoint{2.983562in}{1.894079in}}%
\pgfpathlineto{\pgfqpoint{2.982346in}{1.880831in}}%
\pgfpathlineto{\pgfqpoint{2.976250in}{1.874265in}}%
\pgfpathlineto{\pgfqpoint{2.959902in}{1.867583in}}%
\pgfpathlineto{\pgfqpoint{2.953225in}{1.854335in}}%
\pgfpathlineto{\pgfqpoint{2.950967in}{1.852116in}}%
\pgfpathlineto{\pgfqpoint{2.948343in}{1.854335in}}%
\pgfpathlineto{\pgfqpoint{2.943504in}{1.867583in}}%
\pgfpathlineto{\pgfqpoint{2.937769in}{1.880831in}}%
\pgfpathlineto{\pgfqpoint{2.941042in}{1.894079in}}%
\pgfpathlineto{\pgfqpoint{2.934967in}{1.907327in}}%
\pgfpathlineto{\pgfqpoint{2.939199in}{1.920575in}}%
\pgfpathlineto{\pgfqpoint{2.935853in}{1.933823in}}%
\pgfpathlineto{\pgfqpoint{2.944364in}{1.947071in}}%
\pgfpathlineto{\pgfqpoint{2.942445in}{1.960319in}}%
\pgfpathlineto{\pgfqpoint{2.950179in}{1.973567in}}%
\pgfpathlineto{\pgfqpoint{2.946652in}{1.986815in}}%
\pgfpathlineto{\pgfqpoint{2.950967in}{1.995378in}}%
\pgfpathlineto{\pgfqpoint{2.954287in}{2.000063in}}%
\pgfpathlineto{\pgfqpoint{2.950967in}{2.011483in}}%
\pgfpathlineto{\pgfqpoint{2.950630in}{2.013311in}}%
\pgfpathlineto{\pgfqpoint{2.950967in}{2.013918in}}%
\pgfpathlineto{\pgfqpoint{2.959723in}{2.026559in}}%
\pgfpathlineto{\pgfqpoint{2.957014in}{2.039807in}}%
\pgfpathlineto{\pgfqpoint{2.963983in}{2.053055in}}%
\pgfpathlineto{\pgfqpoint{2.962788in}{2.066303in}}%
\pgfpathlineto{\pgfqpoint{2.968706in}{2.079551in}}%
\pgfpathlineto{\pgfqpoint{2.968615in}{2.092799in}}%
\pgfpathlineto{\pgfqpoint{2.972845in}{2.106047in}}%
\pgfpathlineto{\pgfqpoint{2.973876in}{2.119295in}}%
\pgfpathlineto{\pgfqpoint{2.976250in}{2.126726in}}%
\pgfpathlineto{\pgfqpoint{2.978335in}{2.132543in}}%
\pgfpathlineto{\pgfqpoint{2.976250in}{2.132543in}}%
\pgfpathlineto{\pgfqpoint{2.964146in}{2.132543in}}%
\pgfpathlineto{\pgfqpoint{2.959256in}{2.119295in}}%
\pgfpathlineto{\pgfqpoint{2.958042in}{2.106047in}}%
\pgfpathlineto{\pgfqpoint{2.952830in}{2.092799in}}%
\pgfpathlineto{\pgfqpoint{2.952340in}{2.079551in}}%
\pgfpathlineto{\pgfqpoint{2.950967in}{2.076819in}}%
\pgfpathlineto{\pgfqpoint{2.946069in}{2.066303in}}%
\pgfpathlineto{\pgfqpoint{2.946329in}{2.053055in}}%
\pgfpathlineto{\pgfqpoint{2.939715in}{2.039807in}}%
\pgfpathlineto{\pgfqpoint{2.940552in}{2.026559in}}%
\pgfpathlineto{\pgfqpoint{2.932945in}{2.013311in}}%
\pgfpathlineto{\pgfqpoint{2.933921in}{2.000063in}}%
\pgfpathlineto{\pgfqpoint{2.926808in}{1.986815in}}%
\pgfpathlineto{\pgfqpoint{2.927780in}{1.973567in}}%
\pgfpathlineto{\pgfqpoint{2.925683in}{1.970286in}}%
\pgfpathlineto{\pgfqpoint{2.920414in}{1.960319in}}%
\pgfpathlineto{\pgfqpoint{2.919238in}{1.947071in}}%
\pgfpathlineto{\pgfqpoint{2.912698in}{1.933823in}}%
\pgfpathlineto{\pgfqpoint{2.911701in}{1.920575in}}%
\pgfpathlineto{\pgfqpoint{2.908280in}{1.907327in}}%
\pgfpathlineto{\pgfqpoint{2.906760in}{1.894079in}}%
\pgfpathlineto{\pgfqpoint{2.903512in}{1.880831in}}%
\pgfpathlineto{\pgfqpoint{2.900400in}{1.872207in}}%
\pgfpathlineto{\pgfqpoint{2.898977in}{1.867583in}}%
\pgfpathlineto{\pgfqpoint{2.898128in}{1.854335in}}%
\pgfpathlineto{\pgfqpoint{2.892384in}{1.841087in}}%
\pgfpathlineto{\pgfqpoint{2.893223in}{1.827839in}}%
\pgfpathlineto{\pgfqpoint{2.884321in}{1.814591in}}%
\pgfpathlineto{\pgfqpoint{2.889338in}{1.801343in}}%
\pgfpathlineto{\pgfqpoint{2.879286in}{1.788095in}}%
\pgfpathlineto{\pgfqpoint{2.888956in}{1.774847in}}%
\pgfpathlineto{\pgfqpoint{2.875117in}{1.764328in}}%
\pgfpathlineto{\pgfqpoint{2.872708in}{1.761599in}}%
\pgfpathlineto{\pgfqpoint{2.875117in}{1.757004in}}%
\pgfpathlineto{\pgfqpoint{2.886725in}{1.748351in}}%
\pgfpathlineto{\pgfqpoint{2.875117in}{1.741404in}}%
\pgfpathlineto{\pgfqpoint{2.869798in}{1.735103in}}%
\pgfpathlineto{\pgfqpoint{2.875117in}{1.730512in}}%
\pgfpathclose%
\pgfpathmoveto{\pgfqpoint{3.760028in}{1.743514in}}%
\pgfpathlineto{\pgfqpoint{3.764256in}{1.748351in}}%
\pgfpathlineto{\pgfqpoint{3.775798in}{1.761599in}}%
\pgfpathlineto{\pgfqpoint{3.781417in}{1.774847in}}%
\pgfpathlineto{\pgfqpoint{3.775815in}{1.788095in}}%
\pgfpathlineto{\pgfqpoint{3.783037in}{1.801343in}}%
\pgfpathlineto{\pgfqpoint{3.774369in}{1.814591in}}%
\pgfpathlineto{\pgfqpoint{3.779443in}{1.827839in}}%
\pgfpathlineto{\pgfqpoint{3.771520in}{1.841087in}}%
\pgfpathlineto{\pgfqpoint{3.774871in}{1.854335in}}%
\pgfpathlineto{\pgfqpoint{3.765651in}{1.867583in}}%
\pgfpathlineto{\pgfqpoint{3.768502in}{1.880831in}}%
\pgfpathlineto{\pgfqpoint{3.761631in}{1.894079in}}%
\pgfpathlineto{\pgfqpoint{3.763548in}{1.907327in}}%
\pgfpathlineto{\pgfqpoint{3.760028in}{1.914768in}}%
\pgfpathlineto{\pgfqpoint{3.756600in}{1.920575in}}%
\pgfpathlineto{\pgfqpoint{3.756219in}{1.933823in}}%
\pgfpathlineto{\pgfqpoint{3.748525in}{1.947071in}}%
\pgfpathlineto{\pgfqpoint{3.748144in}{1.960319in}}%
\pgfpathlineto{\pgfqpoint{3.744129in}{1.973567in}}%
\pgfpathlineto{\pgfqpoint{3.742451in}{1.986815in}}%
\pgfpathlineto{\pgfqpoint{3.739605in}{2.000063in}}%
\pgfpathlineto{\pgfqpoint{3.736603in}{2.013311in}}%
\pgfpathlineto{\pgfqpoint{3.734745in}{2.023343in}}%
\pgfpathlineto{\pgfqpoint{3.734070in}{2.026559in}}%
\pgfpathlineto{\pgfqpoint{3.729381in}{2.039807in}}%
\pgfpathlineto{\pgfqpoint{3.727308in}{2.053055in}}%
\pgfpathlineto{\pgfqpoint{3.721505in}{2.066303in}}%
\pgfpathlineto{\pgfqpoint{3.721002in}{2.079551in}}%
\pgfpathlineto{\pgfqpoint{3.716055in}{2.092799in}}%
\pgfpathlineto{\pgfqpoint{3.715682in}{2.106047in}}%
\pgfpathlineto{\pgfqpoint{3.709819in}{2.119295in}}%
\pgfpathlineto{\pgfqpoint{3.709462in}{2.127053in}}%
\pgfpathlineto{\pgfqpoint{3.709201in}{2.132543in}}%
\pgfpathlineto{\pgfqpoint{3.694092in}{2.132543in}}%
\pgfpathlineto{\pgfqpoint{3.693892in}{2.119295in}}%
\pgfpathlineto{\pgfqpoint{3.699715in}{2.106047in}}%
\pgfpathlineto{\pgfqpoint{3.699177in}{2.092799in}}%
\pgfpathlineto{\pgfqpoint{3.704324in}{2.079551in}}%
\pgfpathlineto{\pgfqpoint{3.704144in}{2.066303in}}%
\pgfpathlineto{\pgfqpoint{3.709462in}{2.054934in}}%
\pgfpathlineto{\pgfqpoint{3.710189in}{2.053055in}}%
\pgfpathlineto{\pgfqpoint{3.711672in}{2.039807in}}%
\pgfpathlineto{\pgfqpoint{3.714573in}{2.026559in}}%
\pgfpathlineto{\pgfqpoint{3.716588in}{2.013311in}}%
\pgfpathlineto{\pgfqpoint{3.718020in}{2.000063in}}%
\pgfpathlineto{\pgfqpoint{3.720730in}{1.986815in}}%
\pgfpathlineto{\pgfqpoint{3.719492in}{1.973567in}}%
\pgfpathlineto{\pgfqpoint{3.724123in}{1.960319in}}%
\pgfpathlineto{\pgfqpoint{3.720023in}{1.947071in}}%
\pgfpathlineto{\pgfqpoint{3.730646in}{1.933823in}}%
\pgfpathlineto{\pgfqpoint{3.724736in}{1.920575in}}%
\pgfpathlineto{\pgfqpoint{3.734745in}{1.909024in}}%
\pgfpathlineto{\pgfqpoint{3.735571in}{1.907327in}}%
\pgfpathlineto{\pgfqpoint{3.734745in}{1.905694in}}%
\pgfpathlineto{\pgfqpoint{3.715281in}{1.894079in}}%
\pgfpathlineto{\pgfqpoint{3.734231in}{1.880831in}}%
\pgfpathlineto{\pgfqpoint{3.709462in}{1.874805in}}%
\pgfpathlineto{\pgfqpoint{3.699264in}{1.880831in}}%
\pgfpathlineto{\pgfqpoint{3.703970in}{1.894079in}}%
\pgfpathlineto{\pgfqpoint{3.684179in}{1.900381in}}%
\pgfpathlineto{\pgfqpoint{3.675932in}{1.907327in}}%
\pgfpathlineto{\pgfqpoint{3.675326in}{1.920575in}}%
\pgfpathlineto{\pgfqpoint{3.658895in}{1.933445in}}%
\pgfpathlineto{\pgfqpoint{3.658607in}{1.933823in}}%
\pgfpathlineto{\pgfqpoint{3.658895in}{1.937893in}}%
\pgfpathlineto{\pgfqpoint{3.660404in}{1.947071in}}%
\pgfpathlineto{\pgfqpoint{3.658895in}{1.947921in}}%
\pgfpathlineto{\pgfqpoint{3.648201in}{1.960319in}}%
\pgfpathlineto{\pgfqpoint{3.647114in}{1.973567in}}%
\pgfpathlineto{\pgfqpoint{3.633612in}{1.985318in}}%
\pgfpathlineto{\pgfqpoint{3.632431in}{1.986815in}}%
\pgfpathlineto{\pgfqpoint{3.630201in}{2.000063in}}%
\pgfpathlineto{\pgfqpoint{3.621018in}{2.013311in}}%
\pgfpathlineto{\pgfqpoint{3.618073in}{2.026559in}}%
\pgfpathlineto{\pgfqpoint{3.608329in}{2.039371in}}%
\pgfpathlineto{\pgfqpoint{3.608075in}{2.039807in}}%
\pgfpathlineto{\pgfqpoint{3.605175in}{2.053055in}}%
\pgfpathlineto{\pgfqpoint{3.598306in}{2.066303in}}%
\pgfpathlineto{\pgfqpoint{3.594257in}{2.079551in}}%
\pgfpathlineto{\pgfqpoint{3.587941in}{2.092799in}}%
\pgfpathlineto{\pgfqpoint{3.583046in}{2.102924in}}%
\pgfpathlineto{\pgfqpoint{3.581868in}{2.106047in}}%
\pgfpathlineto{\pgfqpoint{3.577353in}{2.119295in}}%
\pgfpathlineto{\pgfqpoint{3.572592in}{2.132543in}}%
\pgfpathlineto{\pgfqpoint{3.557763in}{2.132543in}}%
\pgfpathlineto{\pgfqpoint{3.557523in}{2.132543in}}%
\pgfpathlineto{\pgfqpoint{3.557763in}{2.131784in}}%
\pgfpathlineto{\pgfqpoint{3.562030in}{2.119295in}}%
\pgfpathlineto{\pgfqpoint{3.567001in}{2.106047in}}%
\pgfpathlineto{\pgfqpoint{3.571737in}{2.092799in}}%
\pgfpathlineto{\pgfqpoint{3.577518in}{2.079551in}}%
\pgfpathlineto{\pgfqpoint{3.580896in}{2.066303in}}%
\pgfpathlineto{\pgfqpoint{3.583046in}{2.062131in}}%
\pgfpathlineto{\pgfqpoint{3.588235in}{2.053055in}}%
\pgfpathlineto{\pgfqpoint{3.591252in}{2.039807in}}%
\pgfpathlineto{\pgfqpoint{3.598994in}{2.026559in}}%
\pgfpathlineto{\pgfqpoint{3.601543in}{2.013311in}}%
\pgfpathlineto{\pgfqpoint{3.608329in}{2.002120in}}%
\pgfpathlineto{\pgfqpoint{3.609848in}{2.000063in}}%
\pgfpathlineto{\pgfqpoint{3.612672in}{1.986815in}}%
\pgfpathlineto{\pgfqpoint{3.623260in}{1.973567in}}%
\pgfpathlineto{\pgfqpoint{3.624991in}{1.960319in}}%
\pgfpathlineto{\pgfqpoint{3.633612in}{1.947271in}}%
\pgfpathlineto{\pgfqpoint{3.633792in}{1.947071in}}%
\pgfpathlineto{\pgfqpoint{3.634819in}{1.933823in}}%
\pgfpathlineto{\pgfqpoint{3.645220in}{1.920575in}}%
\pgfpathlineto{\pgfqpoint{3.647822in}{1.907327in}}%
\pgfpathlineto{\pgfqpoint{3.658834in}{1.894079in}}%
\pgfpathlineto{\pgfqpoint{3.658895in}{1.893749in}}%
\pgfpathlineto{\pgfqpoint{3.662012in}{1.880831in}}%
\pgfpathlineto{\pgfqpoint{3.673702in}{1.867583in}}%
\pgfpathlineto{\pgfqpoint{3.677164in}{1.854335in}}%
\pgfpathlineto{\pgfqpoint{3.684179in}{1.843902in}}%
\pgfpathlineto{\pgfqpoint{3.687247in}{1.841087in}}%
\pgfpathlineto{\pgfqpoint{3.693343in}{1.827839in}}%
\pgfpathlineto{\pgfqpoint{3.704007in}{1.814591in}}%
\pgfpathlineto{\pgfqpoint{3.709368in}{1.801343in}}%
\pgfpathlineto{\pgfqpoint{3.709462in}{1.801189in}}%
\pgfpathlineto{\pgfqpoint{3.723589in}{1.788095in}}%
\pgfpathlineto{\pgfqpoint{3.733123in}{1.774847in}}%
\pgfpathlineto{\pgfqpoint{3.734745in}{1.771039in}}%
\pgfpathlineto{\pgfqpoint{3.745619in}{1.761599in}}%
\pgfpathlineto{\pgfqpoint{3.759149in}{1.748351in}}%
\pgfpathlineto{\pgfqpoint{3.760028in}{1.743514in}}%
\pgfpathclose%
\pgfusepath{fill}%
\end{pgfscope}%
\begin{pgfscope}%
\pgfpathrectangle{\pgfqpoint{2.470586in}{0.476543in}}{\pgfqpoint{1.820388in}{1.656000in}}%
\pgfusepath{clip}%
\pgfsetbuttcap%
\pgfsetroundjoin%
\definecolor{currentfill}{rgb}{0.607151,0.800231,0.886044}%
\pgfsetfillcolor{currentfill}%
\pgfsetlinewidth{0.000000pt}%
\definecolor{currentstroke}{rgb}{0.000000,0.000000,0.000000}%
\pgfsetstrokecolor{currentstroke}%
\pgfsetdash{}{0pt}%
\pgfpathmoveto{\pgfqpoint{3.102666in}{0.793371in}}%
\pgfpathlineto{\pgfqpoint{3.127949in}{0.790171in}}%
\pgfpathlineto{\pgfqpoint{3.153232in}{0.787665in}}%
\pgfpathlineto{\pgfqpoint{3.178515in}{0.785718in}}%
\pgfpathlineto{\pgfqpoint{3.203798in}{0.784234in}}%
\pgfpathlineto{\pgfqpoint{3.229082in}{0.783139in}}%
\pgfpathlineto{\pgfqpoint{3.254365in}{0.782379in}}%
\pgfpathlineto{\pgfqpoint{3.279648in}{0.781914in}}%
\pgfpathlineto{\pgfqpoint{3.304931in}{0.781716in}}%
\pgfpathlineto{\pgfqpoint{3.330214in}{0.781762in}}%
\pgfpathlineto{\pgfqpoint{3.355497in}{0.782040in}}%
\pgfpathlineto{\pgfqpoint{3.380781in}{0.782542in}}%
\pgfpathlineto{\pgfqpoint{3.406064in}{0.783267in}}%
\pgfpathlineto{\pgfqpoint{3.431347in}{0.784218in}}%
\pgfpathlineto{\pgfqpoint{3.456630in}{0.785403in}}%
\pgfpathlineto{\pgfqpoint{3.481913in}{0.786836in}}%
\pgfpathlineto{\pgfqpoint{3.507196in}{0.788535in}}%
\pgfpathlineto{\pgfqpoint{3.532480in}{0.790526in}}%
\pgfpathlineto{\pgfqpoint{3.557763in}{0.792841in}}%
\pgfpathlineto{\pgfqpoint{3.573543in}{0.794495in}}%
\pgfpathlineto{\pgfqpoint{3.583046in}{0.800263in}}%
\pgfpathlineto{\pgfqpoint{3.593091in}{0.807743in}}%
\pgfpathlineto{\pgfqpoint{3.608329in}{0.809913in}}%
\pgfpathlineto{\pgfqpoint{3.633612in}{0.814083in}}%
\pgfpathlineto{\pgfqpoint{3.658895in}{0.819024in}}%
\pgfpathlineto{\pgfqpoint{3.667725in}{0.820991in}}%
\pgfpathlineto{\pgfqpoint{3.666423in}{0.834239in}}%
\pgfpathlineto{\pgfqpoint{3.684179in}{0.839309in}}%
\pgfpathlineto{\pgfqpoint{3.707714in}{0.847487in}}%
\pgfpathlineto{\pgfqpoint{3.704279in}{0.860735in}}%
\pgfpathlineto{\pgfqpoint{3.709462in}{0.863072in}}%
\pgfpathlineto{\pgfqpoint{3.729746in}{0.873983in}}%
\pgfpathlineto{\pgfqpoint{3.734745in}{0.885713in}}%
\pgfpathlineto{\pgfqpoint{3.735356in}{0.887231in}}%
\pgfpathlineto{\pgfqpoint{3.744233in}{0.900479in}}%
\pgfpathlineto{\pgfqpoint{3.744436in}{0.913727in}}%
\pgfpathlineto{\pgfqpoint{3.743749in}{0.926975in}}%
\pgfpathlineto{\pgfqpoint{3.748374in}{0.940223in}}%
\pgfpathlineto{\pgfqpoint{3.742692in}{0.953471in}}%
\pgfpathlineto{\pgfqpoint{3.749161in}{0.966719in}}%
\pgfpathlineto{\pgfqpoint{3.738808in}{0.979967in}}%
\pgfpathlineto{\pgfqpoint{3.743205in}{0.993215in}}%
\pgfpathlineto{\pgfqpoint{3.734745in}{1.002709in}}%
\pgfpathlineto{\pgfqpoint{3.731621in}{1.006463in}}%
\pgfpathlineto{\pgfqpoint{3.734745in}{1.016893in}}%
\pgfpathlineto{\pgfqpoint{3.735568in}{1.019711in}}%
\pgfpathlineto{\pgfqpoint{3.734745in}{1.020493in}}%
\pgfpathlineto{\pgfqpoint{3.722385in}{1.032959in}}%
\pgfpathlineto{\pgfqpoint{3.726329in}{1.046207in}}%
\pgfpathlineto{\pgfqpoint{3.716530in}{1.059455in}}%
\pgfpathlineto{\pgfqpoint{3.719081in}{1.072703in}}%
\pgfpathlineto{\pgfqpoint{3.710100in}{1.085951in}}%
\pgfpathlineto{\pgfqpoint{3.709462in}{1.090398in}}%
\pgfpathlineto{\pgfqpoint{3.708340in}{1.099199in}}%
\pgfpathlineto{\pgfqpoint{3.698813in}{1.112447in}}%
\pgfpathlineto{\pgfqpoint{3.695878in}{1.125695in}}%
\pgfpathlineto{\pgfqpoint{3.690309in}{1.138943in}}%
\pgfpathlineto{\pgfqpoint{3.685406in}{1.152191in}}%
\pgfpathlineto{\pgfqpoint{3.684179in}{1.155833in}}%
\pgfpathlineto{\pgfqpoint{3.681192in}{1.165439in}}%
\pgfpathlineto{\pgfqpoint{3.675022in}{1.178687in}}%
\pgfpathlineto{\pgfqpoint{3.671894in}{1.191935in}}%
\pgfpathlineto{\pgfqpoint{3.664376in}{1.205183in}}%
\pgfpathlineto{\pgfqpoint{3.662385in}{1.218431in}}%
\pgfpathlineto{\pgfqpoint{3.658895in}{1.223331in}}%
\pgfpathlineto{\pgfqpoint{3.653518in}{1.231679in}}%
\pgfpathlineto{\pgfqpoint{3.653037in}{1.244927in}}%
\pgfpathlineto{\pgfqpoint{3.645033in}{1.258175in}}%
\pgfpathlineto{\pgfqpoint{3.644040in}{1.271423in}}%
\pgfpathlineto{\pgfqpoint{3.634330in}{1.284671in}}%
\pgfpathlineto{\pgfqpoint{3.633612in}{1.289474in}}%
\pgfpathlineto{\pgfqpoint{3.632512in}{1.297919in}}%
\pgfpathlineto{\pgfqpoint{3.623693in}{1.311167in}}%
\pgfpathlineto{\pgfqpoint{3.622835in}{1.324415in}}%
\pgfpathlineto{\pgfqpoint{3.615841in}{1.337663in}}%
\pgfpathlineto{\pgfqpoint{3.613815in}{1.350911in}}%
\pgfpathlineto{\pgfqpoint{3.608329in}{1.361036in}}%
\pgfpathlineto{\pgfqpoint{3.606797in}{1.364159in}}%
\pgfpathlineto{\pgfqpoint{3.604248in}{1.377407in}}%
\pgfpathlineto{\pgfqpoint{3.598384in}{1.390655in}}%
\pgfpathlineto{\pgfqpoint{3.594612in}{1.403903in}}%
\pgfpathlineto{\pgfqpoint{3.589576in}{1.417151in}}%
\pgfpathlineto{\pgfqpoint{3.584471in}{1.430399in}}%
\pgfpathlineto{\pgfqpoint{3.583046in}{1.435471in}}%
\pgfpathlineto{\pgfqpoint{3.581026in}{1.443647in}}%
\pgfpathlineto{\pgfqpoint{3.576081in}{1.456895in}}%
\pgfpathlineto{\pgfqpoint{3.572771in}{1.470143in}}%
\pgfpathlineto{\pgfqpoint{3.566125in}{1.483391in}}%
\pgfpathlineto{\pgfqpoint{3.563598in}{1.496639in}}%
\pgfpathlineto{\pgfqpoint{3.557763in}{1.507957in}}%
\pgfpathlineto{\pgfqpoint{3.556886in}{1.509887in}}%
\pgfpathlineto{\pgfqpoint{3.555847in}{1.523135in}}%
\pgfpathlineto{\pgfqpoint{3.549863in}{1.536383in}}%
\pgfpathlineto{\pgfqpoint{3.548095in}{1.549631in}}%
\pgfpathlineto{\pgfqpoint{3.540948in}{1.562879in}}%
\pgfpathlineto{\pgfqpoint{3.539303in}{1.576127in}}%
\pgfpathlineto{\pgfqpoint{3.532480in}{1.588198in}}%
\pgfpathlineto{\pgfqpoint{3.531883in}{1.589375in}}%
\pgfpathlineto{\pgfqpoint{3.530127in}{1.602623in}}%
\pgfpathlineto{\pgfqpoint{3.523943in}{1.615871in}}%
\pgfpathlineto{\pgfqpoint{3.521963in}{1.629119in}}%
\pgfpathlineto{\pgfqpoint{3.516345in}{1.642367in}}%
\pgfpathlineto{\pgfqpoint{3.513826in}{1.655615in}}%
\pgfpathlineto{\pgfqpoint{3.507725in}{1.668863in}}%
\pgfpathlineto{\pgfqpoint{3.507196in}{1.670712in}}%
\pgfpathlineto{\pgfqpoint{3.504399in}{1.682111in}}%
\pgfpathlineto{\pgfqpoint{3.500059in}{1.695359in}}%
\pgfpathlineto{\pgfqpoint{3.496615in}{1.708607in}}%
\pgfpathlineto{\pgfqpoint{3.493006in}{1.721855in}}%
\pgfpathlineto{\pgfqpoint{3.488383in}{1.735103in}}%
\pgfpathlineto{\pgfqpoint{3.484128in}{1.748351in}}%
\pgfpathlineto{\pgfqpoint{3.481913in}{1.753243in}}%
\pgfpathlineto{\pgfqpoint{3.478513in}{1.761599in}}%
\pgfpathlineto{\pgfqpoint{3.475648in}{1.774847in}}%
\pgfpathlineto{\pgfqpoint{3.469752in}{1.788095in}}%
\pgfpathlineto{\pgfqpoint{3.467374in}{1.801343in}}%
\pgfpathlineto{\pgfqpoint{3.460993in}{1.814591in}}%
\pgfpathlineto{\pgfqpoint{3.458888in}{1.827839in}}%
\pgfpathlineto{\pgfqpoint{3.456630in}{1.832231in}}%
\pgfpathlineto{\pgfqpoint{3.452255in}{1.841087in}}%
\pgfpathlineto{\pgfqpoint{3.450201in}{1.854335in}}%
\pgfpathlineto{\pgfqpoint{3.442544in}{1.867583in}}%
\pgfpathlineto{\pgfqpoint{3.440356in}{1.880831in}}%
\pgfpathlineto{\pgfqpoint{3.432945in}{1.894079in}}%
\pgfpathlineto{\pgfqpoint{3.431347in}{1.904122in}}%
\pgfpathlineto{\pgfqpoint{3.430830in}{1.907327in}}%
\pgfpathlineto{\pgfqpoint{3.422938in}{1.920575in}}%
\pgfpathlineto{\pgfqpoint{3.419695in}{1.933823in}}%
\pgfpathlineto{\pgfqpoint{3.410466in}{1.947071in}}%
\pgfpathlineto{\pgfqpoint{3.406969in}{1.960319in}}%
\pgfpathlineto{\pgfqpoint{3.406064in}{1.961891in}}%
\pgfpathlineto{\pgfqpoint{3.397599in}{1.973567in}}%
\pgfpathlineto{\pgfqpoint{3.393053in}{1.986815in}}%
\pgfpathlineto{\pgfqpoint{3.383691in}{2.000063in}}%
\pgfpathlineto{\pgfqpoint{3.380781in}{2.006893in}}%
\pgfpathlineto{\pgfqpoint{3.376537in}{2.013311in}}%
\pgfpathlineto{\pgfqpoint{3.360423in}{2.026559in}}%
\pgfpathlineto{\pgfqpoint{3.355497in}{2.033110in}}%
\pgfpathlineto{\pgfqpoint{3.332787in}{2.039807in}}%
\pgfpathlineto{\pgfqpoint{3.330214in}{2.040243in}}%
\pgfpathlineto{\pgfqpoint{3.329349in}{2.039807in}}%
\pgfpathlineto{\pgfqpoint{3.317681in}{2.026559in}}%
\pgfpathlineto{\pgfqpoint{3.304931in}{2.020660in}}%
\pgfpathlineto{\pgfqpoint{3.297159in}{2.013311in}}%
\pgfpathlineto{\pgfqpoint{3.292790in}{2.000063in}}%
\pgfpathlineto{\pgfqpoint{3.279782in}{1.986815in}}%
\pgfpathlineto{\pgfqpoint{3.279648in}{1.986298in}}%
\pgfpathlineto{\pgfqpoint{3.277125in}{1.973567in}}%
\pgfpathlineto{\pgfqpoint{3.267145in}{1.960319in}}%
\pgfpathlineto{\pgfqpoint{3.263068in}{1.947071in}}%
\pgfpathlineto{\pgfqpoint{3.254365in}{1.934241in}}%
\pgfpathlineto{\pgfqpoint{3.254119in}{1.933823in}}%
\pgfpathlineto{\pgfqpoint{3.250129in}{1.920575in}}%
\pgfpathlineto{\pgfqpoint{3.244575in}{1.907327in}}%
\pgfpathlineto{\pgfqpoint{3.240450in}{1.894079in}}%
\pgfpathlineto{\pgfqpoint{3.235585in}{1.880831in}}%
\pgfpathlineto{\pgfqpoint{3.230206in}{1.867583in}}%
\pgfpathlineto{\pgfqpoint{3.229082in}{1.863367in}}%
\pgfpathlineto{\pgfqpoint{3.226680in}{1.854335in}}%
\pgfpathlineto{\pgfqpoint{3.220158in}{1.841087in}}%
\pgfpathlineto{\pgfqpoint{3.217256in}{1.827839in}}%
\pgfpathlineto{\pgfqpoint{3.210404in}{1.814591in}}%
\pgfpathlineto{\pgfqpoint{3.208531in}{1.801343in}}%
\pgfpathlineto{\pgfqpoint{3.203798in}{1.791509in}}%
\pgfpathlineto{\pgfqpoint{3.202048in}{1.788095in}}%
\pgfpathlineto{\pgfqpoint{3.200353in}{1.774847in}}%
\pgfpathlineto{\pgfqpoint{3.192832in}{1.761599in}}%
\pgfpathlineto{\pgfqpoint{3.190766in}{1.748351in}}%
\pgfpathlineto{\pgfqpoint{3.184750in}{1.735103in}}%
\pgfpathlineto{\pgfqpoint{3.183008in}{1.721855in}}%
\pgfpathlineto{\pgfqpoint{3.178515in}{1.710589in}}%
\pgfpathlineto{\pgfqpoint{3.177647in}{1.708607in}}%
\pgfpathlineto{\pgfqpoint{3.174646in}{1.695359in}}%
\pgfpathlineto{\pgfqpoint{3.168931in}{1.682111in}}%
\pgfpathlineto{\pgfqpoint{3.165066in}{1.668863in}}%
\pgfpathlineto{\pgfqpoint{3.160820in}{1.655615in}}%
\pgfpathlineto{\pgfqpoint{3.156011in}{1.642367in}}%
\pgfpathlineto{\pgfqpoint{3.153232in}{1.630981in}}%
\pgfpathlineto{\pgfqpoint{3.152719in}{1.629119in}}%
\pgfpathlineto{\pgfqpoint{3.147234in}{1.615871in}}%
\pgfpathlineto{\pgfqpoint{3.145236in}{1.602623in}}%
\pgfpathlineto{\pgfqpoint{3.139231in}{1.589375in}}%
\pgfpathlineto{\pgfqpoint{3.137318in}{1.576127in}}%
\pgfpathlineto{\pgfqpoint{3.130107in}{1.562879in}}%
\pgfpathlineto{\pgfqpoint{3.128769in}{1.549631in}}%
\pgfpathlineto{\pgfqpoint{3.127949in}{1.547935in}}%
\pgfpathlineto{\pgfqpoint{3.121676in}{1.536383in}}%
\pgfpathlineto{\pgfqpoint{3.120387in}{1.523135in}}%
\pgfpathlineto{\pgfqpoint{3.112959in}{1.509887in}}%
\pgfpathlineto{\pgfqpoint{3.111148in}{1.496639in}}%
\pgfpathlineto{\pgfqpoint{3.104181in}{1.483391in}}%
\pgfpathlineto{\pgfqpoint{3.102666in}{1.473512in}}%
\pgfpathlineto{\pgfqpoint{3.102051in}{1.470143in}}%
\pgfpathlineto{\pgfqpoint{3.095479in}{1.456895in}}%
\pgfpathlineto{\pgfqpoint{3.092363in}{1.443647in}}%
\pgfpathlineto{\pgfqpoint{3.087263in}{1.430399in}}%
\pgfpathlineto{\pgfqpoint{3.083950in}{1.417151in}}%
\pgfpathlineto{\pgfqpoint{3.079901in}{1.403903in}}%
\pgfpathlineto{\pgfqpoint{3.077382in}{1.397597in}}%
\pgfpathlineto{\pgfqpoint{3.074154in}{1.390655in}}%
\pgfpathlineto{\pgfqpoint{3.069871in}{1.377407in}}%
\pgfpathlineto{\pgfqpoint{3.062512in}{1.364159in}}%
\pgfpathlineto{\pgfqpoint{3.060605in}{1.350911in}}%
\pgfpathlineto{\pgfqpoint{3.053785in}{1.337663in}}%
\pgfpathlineto{\pgfqpoint{3.052349in}{1.324415in}}%
\pgfpathlineto{\pgfqpoint{3.052099in}{1.324022in}}%
\pgfpathlineto{\pgfqpoint{3.043071in}{1.311167in}}%
\pgfpathlineto{\pgfqpoint{3.042198in}{1.297919in}}%
\pgfpathlineto{\pgfqpoint{3.033032in}{1.284671in}}%
\pgfpathlineto{\pgfqpoint{3.032936in}{1.271423in}}%
\pgfpathlineto{\pgfqpoint{3.026816in}{1.262463in}}%
\pgfpathlineto{\pgfqpoint{3.023621in}{1.258175in}}%
\pgfpathlineto{\pgfqpoint{3.023089in}{1.244927in}}%
\pgfpathlineto{\pgfqpoint{3.013890in}{1.231679in}}%
\pgfpathlineto{\pgfqpoint{3.013090in}{1.218431in}}%
\pgfpathlineto{\pgfqpoint{3.003689in}{1.205183in}}%
\pgfpathlineto{\pgfqpoint{3.001533in}{1.196357in}}%
\pgfpathlineto{\pgfqpoint{3.000294in}{1.191935in}}%
\pgfpathlineto{\pgfqpoint{2.991749in}{1.178687in}}%
\pgfpathlineto{\pgfqpoint{2.988312in}{1.165439in}}%
\pgfpathlineto{\pgfqpoint{2.983983in}{1.152191in}}%
\pgfpathlineto{\pgfqpoint{2.979865in}{1.138943in}}%
\pgfpathlineto{\pgfqpoint{2.976250in}{1.128244in}}%
\pgfpathlineto{\pgfqpoint{2.975328in}{1.125695in}}%
\pgfpathlineto{\pgfqpoint{2.967545in}{1.112447in}}%
\pgfpathlineto{\pgfqpoint{2.965336in}{1.099199in}}%
\pgfpathlineto{\pgfqpoint{2.956040in}{1.085951in}}%
\pgfpathlineto{\pgfqpoint{2.956218in}{1.072703in}}%
\pgfpathlineto{\pgfqpoint{2.950967in}{1.066302in}}%
\pgfpathlineto{\pgfqpoint{2.944843in}{1.059455in}}%
\pgfpathlineto{\pgfqpoint{2.947357in}{1.046207in}}%
\pgfpathlineto{\pgfqpoint{2.934919in}{1.032959in}}%
\pgfpathlineto{\pgfqpoint{2.939361in}{1.019711in}}%
\pgfpathlineto{\pgfqpoint{2.925683in}{1.007855in}}%
\pgfpathlineto{\pgfqpoint{2.923969in}{1.006463in}}%
\pgfpathlineto{\pgfqpoint{2.925683in}{1.002316in}}%
\pgfpathlineto{\pgfqpoint{2.929574in}{0.993215in}}%
\pgfpathlineto{\pgfqpoint{2.925683in}{0.989340in}}%
\pgfpathlineto{\pgfqpoint{2.915819in}{0.979967in}}%
\pgfpathlineto{\pgfqpoint{2.925683in}{0.967872in}}%
\pgfpathlineto{\pgfqpoint{2.926692in}{0.966719in}}%
\pgfpathlineto{\pgfqpoint{2.925683in}{0.965428in}}%
\pgfpathlineto{\pgfqpoint{2.915891in}{0.953471in}}%
\pgfpathlineto{\pgfqpoint{2.925683in}{0.941557in}}%
\pgfpathlineto{\pgfqpoint{2.926898in}{0.940223in}}%
\pgfpathlineto{\pgfqpoint{2.925683in}{0.939091in}}%
\pgfpathlineto{\pgfqpoint{2.912217in}{0.926975in}}%
\pgfpathlineto{\pgfqpoint{2.922863in}{0.913727in}}%
\pgfpathlineto{\pgfqpoint{2.913077in}{0.900479in}}%
\pgfpathlineto{\pgfqpoint{2.925683in}{0.889651in}}%
\pgfpathlineto{\pgfqpoint{2.928985in}{0.887231in}}%
\pgfpathlineto{\pgfqpoint{2.925683in}{0.875993in}}%
\pgfpathlineto{\pgfqpoint{2.925082in}{0.873983in}}%
\pgfpathlineto{\pgfqpoint{2.925683in}{0.873480in}}%
\pgfpathlineto{\pgfqpoint{2.943906in}{0.860735in}}%
\pgfpathlineto{\pgfqpoint{2.949597in}{0.847487in}}%
\pgfpathlineto{\pgfqpoint{2.950967in}{0.846671in}}%
\pgfpathlineto{\pgfqpoint{2.976250in}{0.834433in}}%
\pgfpathlineto{\pgfqpoint{2.976745in}{0.834239in}}%
\pgfpathlineto{\pgfqpoint{2.997618in}{0.820991in}}%
\pgfpathlineto{\pgfqpoint{3.001533in}{0.819523in}}%
\pgfpathlineto{\pgfqpoint{3.026816in}{0.811803in}}%
\pgfpathlineto{\pgfqpoint{3.044347in}{0.807743in}}%
\pgfpathlineto{\pgfqpoint{3.052099in}{0.805497in}}%
\pgfpathlineto{\pgfqpoint{3.077382in}{0.798796in}}%
\pgfpathlineto{\pgfqpoint{3.095288in}{0.794495in}}%
\pgfpathlineto{\pgfqpoint{3.102666in}{0.793371in}}%
\pgfpathclose%
\pgfpathmoveto{\pgfqpoint{3.128912in}{0.820991in}}%
\pgfpathlineto{\pgfqpoint{3.127949in}{0.821292in}}%
\pgfpathlineto{\pgfqpoint{3.102666in}{0.830298in}}%
\pgfpathlineto{\pgfqpoint{3.092977in}{0.834239in}}%
\pgfpathlineto{\pgfqpoint{3.077382in}{0.837183in}}%
\pgfpathlineto{\pgfqpoint{3.052099in}{0.843154in}}%
\pgfpathlineto{\pgfqpoint{3.037372in}{0.847487in}}%
\pgfpathlineto{\pgfqpoint{3.027734in}{0.860735in}}%
\pgfpathlineto{\pgfqpoint{3.026816in}{0.861014in}}%
\pgfpathlineto{\pgfqpoint{3.001533in}{0.870268in}}%
\pgfpathlineto{\pgfqpoint{2.993357in}{0.873983in}}%
\pgfpathlineto{\pgfqpoint{2.994995in}{0.887231in}}%
\pgfpathlineto{\pgfqpoint{2.976250in}{0.896245in}}%
\pgfpathlineto{\pgfqpoint{2.968897in}{0.900479in}}%
\pgfpathlineto{\pgfqpoint{2.976250in}{0.912649in}}%
\pgfpathlineto{\pgfqpoint{2.976887in}{0.913727in}}%
\pgfpathlineto{\pgfqpoint{2.976250in}{0.914158in}}%
\pgfpathlineto{\pgfqpoint{2.959831in}{0.926975in}}%
\pgfpathlineto{\pgfqpoint{2.972847in}{0.940223in}}%
\pgfpathlineto{\pgfqpoint{2.956948in}{0.953471in}}%
\pgfpathlineto{\pgfqpoint{2.966090in}{0.966719in}}%
\pgfpathlineto{\pgfqpoint{2.952103in}{0.979967in}}%
\pgfpathlineto{\pgfqpoint{2.964461in}{0.993215in}}%
\pgfpathlineto{\pgfqpoint{2.956820in}{1.006463in}}%
\pgfpathlineto{\pgfqpoint{2.970561in}{1.019711in}}%
\pgfpathlineto{\pgfqpoint{2.964601in}{1.032959in}}%
\pgfpathlineto{\pgfqpoint{2.975593in}{1.046207in}}%
\pgfpathlineto{\pgfqpoint{2.972224in}{1.059455in}}%
\pgfpathlineto{\pgfqpoint{2.976250in}{1.064655in}}%
\pgfpathlineto{\pgfqpoint{2.981973in}{1.072703in}}%
\pgfpathlineto{\pgfqpoint{2.981200in}{1.085951in}}%
\pgfpathlineto{\pgfqpoint{2.988783in}{1.099199in}}%
\pgfpathlineto{\pgfqpoint{2.990689in}{1.112447in}}%
\pgfpathlineto{\pgfqpoint{2.996859in}{1.125695in}}%
\pgfpathlineto{\pgfqpoint{3.001398in}{1.138943in}}%
\pgfpathlineto{\pgfqpoint{3.001533in}{1.139511in}}%
\pgfpathlineto{\pgfqpoint{3.004205in}{1.152191in}}%
\pgfpathlineto{\pgfqpoint{3.008493in}{1.165439in}}%
\pgfpathlineto{\pgfqpoint{3.010735in}{1.178687in}}%
\pgfpathlineto{\pgfqpoint{3.019068in}{1.191935in}}%
\pgfpathlineto{\pgfqpoint{3.021458in}{1.205183in}}%
\pgfpathlineto{\pgfqpoint{3.026816in}{1.212511in}}%
\pgfpathlineto{\pgfqpoint{3.030894in}{1.218431in}}%
\pgfpathlineto{\pgfqpoint{3.030869in}{1.231679in}}%
\pgfpathlineto{\pgfqpoint{3.039644in}{1.244927in}}%
\pgfpathlineto{\pgfqpoint{3.039510in}{1.258175in}}%
\pgfpathlineto{\pgfqpoint{3.048772in}{1.271423in}}%
\pgfpathlineto{\pgfqpoint{3.048346in}{1.284671in}}%
\pgfpathlineto{\pgfqpoint{3.052099in}{1.290009in}}%
\pgfpathlineto{\pgfqpoint{3.057188in}{1.297919in}}%
\pgfpathlineto{\pgfqpoint{3.057767in}{1.311167in}}%
\pgfpathlineto{\pgfqpoint{3.066180in}{1.324415in}}%
\pgfpathlineto{\pgfqpoint{3.067551in}{1.337663in}}%
\pgfpathlineto{\pgfqpoint{3.074353in}{1.350911in}}%
\pgfpathlineto{\pgfqpoint{3.076384in}{1.364159in}}%
\pgfpathlineto{\pgfqpoint{3.077382in}{1.365955in}}%
\pgfpathlineto{\pgfqpoint{3.083020in}{1.377407in}}%
\pgfpathlineto{\pgfqpoint{3.087057in}{1.390655in}}%
\pgfpathlineto{\pgfqpoint{3.092129in}{1.403903in}}%
\pgfpathlineto{\pgfqpoint{3.096450in}{1.417151in}}%
\pgfpathlineto{\pgfqpoint{3.099510in}{1.430399in}}%
\pgfpathlineto{\pgfqpoint{3.102666in}{1.437998in}}%
\pgfpathlineto{\pgfqpoint{3.104823in}{1.443647in}}%
\pgfpathlineto{\pgfqpoint{3.107326in}{1.456895in}}%
\pgfpathlineto{\pgfqpoint{3.113449in}{1.470143in}}%
\pgfpathlineto{\pgfqpoint{3.115212in}{1.483391in}}%
\pgfpathlineto{\pgfqpoint{3.122659in}{1.496639in}}%
\pgfpathlineto{\pgfqpoint{3.124315in}{1.509887in}}%
\pgfpathlineto{\pgfqpoint{3.127949in}{1.516091in}}%
\pgfpathlineto{\pgfqpoint{3.131722in}{1.523135in}}%
\pgfpathlineto{\pgfqpoint{3.132710in}{1.536383in}}%
\pgfpathlineto{\pgfqpoint{3.139228in}{1.549631in}}%
\pgfpathlineto{\pgfqpoint{3.140540in}{1.562879in}}%
\pgfpathlineto{\pgfqpoint{3.148004in}{1.576127in}}%
\pgfpathlineto{\pgfqpoint{3.149995in}{1.589375in}}%
\pgfpathlineto{\pgfqpoint{3.153232in}{1.596350in}}%
\pgfpathlineto{\pgfqpoint{3.155828in}{1.602623in}}%
\pgfpathlineto{\pgfqpoint{3.157759in}{1.615871in}}%
\pgfpathlineto{\pgfqpoint{3.162593in}{1.629119in}}%
\pgfpathlineto{\pgfqpoint{3.166148in}{1.642367in}}%
\pgfpathlineto{\pgfqpoint{3.170928in}{1.655615in}}%
\pgfpathlineto{\pgfqpoint{3.175592in}{1.668863in}}%
\pgfpathlineto{\pgfqpoint{3.178515in}{1.678974in}}%
\pgfpathlineto{\pgfqpoint{3.179324in}{1.682111in}}%
\pgfpathlineto{\pgfqpoint{3.184901in}{1.695359in}}%
\pgfpathlineto{\pgfqpoint{3.187406in}{1.708607in}}%
\pgfpathlineto{\pgfqpoint{3.193082in}{1.721855in}}%
\pgfpathlineto{\pgfqpoint{3.194677in}{1.735103in}}%
\pgfpathlineto{\pgfqpoint{3.201216in}{1.748351in}}%
\pgfpathlineto{\pgfqpoint{3.203256in}{1.761599in}}%
\pgfpathlineto{\pgfqpoint{3.203798in}{1.762532in}}%
\pgfpathlineto{\pgfqpoint{3.210696in}{1.774847in}}%
\pgfpathlineto{\pgfqpoint{3.212262in}{1.788095in}}%
\pgfpathlineto{\pgfqpoint{3.218906in}{1.801343in}}%
\pgfpathlineto{\pgfqpoint{3.220956in}{1.814591in}}%
\pgfpathlineto{\pgfqpoint{3.228158in}{1.827839in}}%
\pgfpathlineto{\pgfqpoint{3.229082in}{1.831748in}}%
\pgfpathlineto{\pgfqpoint{3.231317in}{1.841087in}}%
\pgfpathlineto{\pgfqpoint{3.237788in}{1.854335in}}%
\pgfpathlineto{\pgfqpoint{3.241667in}{1.867583in}}%
\pgfpathlineto{\pgfqpoint{3.247037in}{1.880831in}}%
\pgfpathlineto{\pgfqpoint{3.252523in}{1.894079in}}%
\pgfpathlineto{\pgfqpoint{3.254365in}{1.900059in}}%
\pgfpathlineto{\pgfqpoint{3.256793in}{1.907327in}}%
\pgfpathlineto{\pgfqpoint{3.263646in}{1.920575in}}%
\pgfpathlineto{\pgfqpoint{3.267829in}{1.933823in}}%
\pgfpathlineto{\pgfqpoint{3.277804in}{1.947071in}}%
\pgfpathlineto{\pgfqpoint{3.279648in}{1.953227in}}%
\pgfpathlineto{\pgfqpoint{3.282451in}{1.960319in}}%
\pgfpathlineto{\pgfqpoint{3.296663in}{1.973567in}}%
\pgfpathlineto{\pgfqpoint{3.300132in}{1.986815in}}%
\pgfpathlineto{\pgfqpoint{3.304931in}{1.991517in}}%
\pgfpathlineto{\pgfqpoint{3.323439in}{2.000063in}}%
\pgfpathlineto{\pgfqpoint{3.330214in}{2.008604in}}%
\pgfpathlineto{\pgfqpoint{3.355497in}{2.000412in}}%
\pgfpathlineto{\pgfqpoint{3.355714in}{2.000063in}}%
\pgfpathlineto{\pgfqpoint{3.372823in}{1.986815in}}%
\pgfpathlineto{\pgfqpoint{3.379542in}{1.973567in}}%
\pgfpathlineto{\pgfqpoint{3.380781in}{1.972675in}}%
\pgfpathlineto{\pgfqpoint{3.390712in}{1.960319in}}%
\pgfpathlineto{\pgfqpoint{3.394736in}{1.947071in}}%
\pgfpathlineto{\pgfqpoint{3.406064in}{1.934618in}}%
\pgfpathlineto{\pgfqpoint{3.406634in}{1.933823in}}%
\pgfpathlineto{\pgfqpoint{3.409889in}{1.920575in}}%
\pgfpathlineto{\pgfqpoint{3.418389in}{1.907327in}}%
\pgfpathlineto{\pgfqpoint{3.420480in}{1.894079in}}%
\pgfpathlineto{\pgfqpoint{3.428700in}{1.880831in}}%
\pgfpathlineto{\pgfqpoint{3.431024in}{1.867583in}}%
\pgfpathlineto{\pgfqpoint{3.431347in}{1.867095in}}%
\pgfpathlineto{\pgfqpoint{3.439165in}{1.854335in}}%
\pgfpathlineto{\pgfqpoint{3.441296in}{1.841087in}}%
\pgfpathlineto{\pgfqpoint{3.448062in}{1.827839in}}%
\pgfpathlineto{\pgfqpoint{3.450211in}{1.814591in}}%
\pgfpathlineto{\pgfqpoint{3.456624in}{1.801343in}}%
\pgfpathlineto{\pgfqpoint{3.456630in}{1.801310in}}%
\pgfpathlineto{\pgfqpoint{3.459218in}{1.788095in}}%
\pgfpathlineto{\pgfqpoint{3.465344in}{1.774847in}}%
\pgfpathlineto{\pgfqpoint{3.468488in}{1.761599in}}%
\pgfpathlineto{\pgfqpoint{3.473946in}{1.748351in}}%
\pgfpathlineto{\pgfqpoint{3.478095in}{1.735103in}}%
\pgfpathlineto{\pgfqpoint{3.481913in}{1.723327in}}%
\pgfpathlineto{\pgfqpoint{3.482438in}{1.721855in}}%
\pgfpathlineto{\pgfqpoint{3.486371in}{1.708607in}}%
\pgfpathlineto{\pgfqpoint{3.489798in}{1.695359in}}%
\pgfpathlineto{\pgfqpoint{3.494555in}{1.682111in}}%
\pgfpathlineto{\pgfqpoint{3.497712in}{1.668863in}}%
\pgfpathlineto{\pgfqpoint{3.503429in}{1.655615in}}%
\pgfpathlineto{\pgfqpoint{3.505588in}{1.642367in}}%
\pgfpathlineto{\pgfqpoint{3.507196in}{1.638542in}}%
\pgfpathlineto{\pgfqpoint{3.511471in}{1.629119in}}%
\pgfpathlineto{\pgfqpoint{3.513401in}{1.615871in}}%
\pgfpathlineto{\pgfqpoint{3.520053in}{1.602623in}}%
\pgfpathlineto{\pgfqpoint{3.521668in}{1.589375in}}%
\pgfpathlineto{\pgfqpoint{3.528505in}{1.576127in}}%
\pgfpathlineto{\pgfqpoint{3.529875in}{1.562879in}}%
\pgfpathlineto{\pgfqpoint{3.532480in}{1.557773in}}%
\pgfpathlineto{\pgfqpoint{3.537057in}{1.549631in}}%
\pgfpathlineto{\pgfqpoint{3.538798in}{1.536383in}}%
\pgfpathlineto{\pgfqpoint{3.545026in}{1.523135in}}%
\pgfpathlineto{\pgfqpoint{3.546004in}{1.509887in}}%
\pgfpathlineto{\pgfqpoint{3.552083in}{1.496639in}}%
\pgfpathlineto{\pgfqpoint{3.554422in}{1.483391in}}%
\pgfpathlineto{\pgfqpoint{3.557763in}{1.475996in}}%
\pgfpathlineto{\pgfqpoint{3.560747in}{1.470143in}}%
\pgfpathlineto{\pgfqpoint{3.564226in}{1.456895in}}%
\pgfpathlineto{\pgfqpoint{3.569156in}{1.443647in}}%
\pgfpathlineto{\pgfqpoint{3.572606in}{1.430399in}}%
\pgfpathlineto{\pgfqpoint{3.576966in}{1.417151in}}%
\pgfpathlineto{\pgfqpoint{3.581726in}{1.403903in}}%
\pgfpathlineto{\pgfqpoint{3.583046in}{1.398213in}}%
\pgfpathlineto{\pgfqpoint{3.585063in}{1.390655in}}%
\pgfpathlineto{\pgfqpoint{3.591254in}{1.377407in}}%
\pgfpathlineto{\pgfqpoint{3.593483in}{1.364159in}}%
\pgfpathlineto{\pgfqpoint{3.600180in}{1.350911in}}%
\pgfpathlineto{\pgfqpoint{3.601625in}{1.337663in}}%
\pgfpathlineto{\pgfqpoint{3.608262in}{1.324415in}}%
\pgfpathlineto{\pgfqpoint{3.608329in}{1.321941in}}%
\pgfpathlineto{\pgfqpoint{3.608663in}{1.311167in}}%
\pgfpathlineto{\pgfqpoint{3.617860in}{1.297919in}}%
\pgfpathlineto{\pgfqpoint{3.619273in}{1.284671in}}%
\pgfpathlineto{\pgfqpoint{3.628163in}{1.271423in}}%
\pgfpathlineto{\pgfqpoint{3.628672in}{1.258175in}}%
\pgfpathlineto{\pgfqpoint{3.633612in}{1.249367in}}%
\pgfpathlineto{\pgfqpoint{3.636326in}{1.244927in}}%
\pgfpathlineto{\pgfqpoint{3.636447in}{1.231679in}}%
\pgfpathlineto{\pgfqpoint{3.644857in}{1.218431in}}%
\pgfpathlineto{\pgfqpoint{3.646436in}{1.205183in}}%
\pgfpathlineto{\pgfqpoint{3.652889in}{1.191935in}}%
\pgfpathlineto{\pgfqpoint{3.655684in}{1.178687in}}%
\pgfpathlineto{\pgfqpoint{3.658895in}{1.170512in}}%
\pgfpathlineto{\pgfqpoint{3.661096in}{1.165439in}}%
\pgfpathlineto{\pgfqpoint{3.665198in}{1.152191in}}%
\pgfpathlineto{\pgfqpoint{3.668792in}{1.138943in}}%
\pgfpathlineto{\pgfqpoint{3.674081in}{1.125695in}}%
\pgfpathlineto{\pgfqpoint{3.675593in}{1.112447in}}%
\pgfpathlineto{\pgfqpoint{3.684179in}{1.100011in}}%
\pgfpathlineto{\pgfqpoint{3.684767in}{1.099199in}}%
\pgfpathlineto{\pgfqpoint{3.685011in}{1.085951in}}%
\pgfpathlineto{\pgfqpoint{3.693255in}{1.072703in}}%
\pgfpathlineto{\pgfqpoint{3.688907in}{1.059455in}}%
\pgfpathlineto{\pgfqpoint{3.698156in}{1.046207in}}%
\pgfpathlineto{\pgfqpoint{3.692388in}{1.032959in}}%
\pgfpathlineto{\pgfqpoint{3.704396in}{1.019711in}}%
\pgfpathlineto{\pgfqpoint{3.698561in}{1.006463in}}%
\pgfpathlineto{\pgfqpoint{3.708255in}{0.993215in}}%
\pgfpathlineto{\pgfqpoint{3.701973in}{0.979967in}}%
\pgfpathlineto{\pgfqpoint{3.709462in}{0.967413in}}%
\pgfpathlineto{\pgfqpoint{3.709907in}{0.966719in}}%
\pgfpathlineto{\pgfqpoint{3.709462in}{0.966007in}}%
\pgfpathlineto{\pgfqpoint{3.701189in}{0.953471in}}%
\pgfpathlineto{\pgfqpoint{3.702554in}{0.940223in}}%
\pgfpathlineto{\pgfqpoint{3.695887in}{0.926975in}}%
\pgfpathlineto{\pgfqpoint{3.690696in}{0.913727in}}%
\pgfpathlineto{\pgfqpoint{3.687797in}{0.900479in}}%
\pgfpathlineto{\pgfqpoint{3.684179in}{0.897418in}}%
\pgfpathlineto{\pgfqpoint{3.669913in}{0.887231in}}%
\pgfpathlineto{\pgfqpoint{3.660925in}{0.873983in}}%
\pgfpathlineto{\pgfqpoint{3.658895in}{0.873177in}}%
\pgfpathlineto{\pgfqpoint{3.633612in}{0.864537in}}%
\pgfpathlineto{\pgfqpoint{3.619970in}{0.860735in}}%
\pgfpathlineto{\pgfqpoint{3.619507in}{0.847487in}}%
\pgfpathlineto{\pgfqpoint{3.608329in}{0.844809in}}%
\pgfpathlineto{\pgfqpoint{3.583046in}{0.839624in}}%
\pgfpathlineto{\pgfqpoint{3.557763in}{0.835342in}}%
\pgfpathlineto{\pgfqpoint{3.550160in}{0.834239in}}%
\pgfpathlineto{\pgfqpoint{3.535575in}{0.820991in}}%
\pgfpathlineto{\pgfqpoint{3.532480in}{0.820604in}}%
\pgfpathlineto{\pgfqpoint{3.507196in}{0.817862in}}%
\pgfpathlineto{\pgfqpoint{3.481913in}{0.815571in}}%
\pgfpathlineto{\pgfqpoint{3.456630in}{0.813681in}}%
\pgfpathlineto{\pgfqpoint{3.431347in}{0.812154in}}%
\pgfpathlineto{\pgfqpoint{3.406064in}{0.810963in}}%
\pgfpathlineto{\pgfqpoint{3.380781in}{0.810087in}}%
\pgfpathlineto{\pgfqpoint{3.355497in}{0.809516in}}%
\pgfpathlineto{\pgfqpoint{3.330214in}{0.809248in}}%
\pgfpathlineto{\pgfqpoint{3.304931in}{0.809287in}}%
\pgfpathlineto{\pgfqpoint{3.279648in}{0.809649in}}%
\pgfpathlineto{\pgfqpoint{3.254365in}{0.810359in}}%
\pgfpathlineto{\pgfqpoint{3.229082in}{0.811454in}}%
\pgfpathlineto{\pgfqpoint{3.203798in}{0.812989in}}%
\pgfpathlineto{\pgfqpoint{3.178515in}{0.815038in}}%
\pgfpathlineto{\pgfqpoint{3.153232in}{0.817704in}}%
\pgfpathlineto{\pgfqpoint{3.128912in}{0.820991in}}%
\pgfpathclose%
\pgfpathmoveto{\pgfqpoint{2.950967in}{1.852116in}}%
\pgfpathlineto{\pgfqpoint{2.953225in}{1.854335in}}%
\pgfpathlineto{\pgfqpoint{2.959902in}{1.867583in}}%
\pgfpathlineto{\pgfqpoint{2.976250in}{1.874265in}}%
\pgfpathlineto{\pgfqpoint{2.982346in}{1.880831in}}%
\pgfpathlineto{\pgfqpoint{2.983562in}{1.894079in}}%
\pgfpathlineto{\pgfqpoint{3.001533in}{1.904476in}}%
\pgfpathlineto{\pgfqpoint{3.004441in}{1.907327in}}%
\pgfpathlineto{\pgfqpoint{3.005344in}{1.920575in}}%
\pgfpathlineto{\pgfqpoint{3.018857in}{1.933823in}}%
\pgfpathlineto{\pgfqpoint{3.016974in}{1.947071in}}%
\pgfpathlineto{\pgfqpoint{3.026816in}{1.955643in}}%
\pgfpathlineto{\pgfqpoint{3.030228in}{1.960319in}}%
\pgfpathlineto{\pgfqpoint{3.030922in}{1.973567in}}%
\pgfpathlineto{\pgfqpoint{3.042567in}{1.986815in}}%
\pgfpathlineto{\pgfqpoint{3.045360in}{2.000063in}}%
\pgfpathlineto{\pgfqpoint{3.052099in}{2.007818in}}%
\pgfpathlineto{\pgfqpoint{3.055347in}{2.013311in}}%
\pgfpathlineto{\pgfqpoint{3.059457in}{2.026559in}}%
\pgfpathlineto{\pgfqpoint{3.066292in}{2.039807in}}%
\pgfpathlineto{\pgfqpoint{3.071714in}{2.053055in}}%
\pgfpathlineto{\pgfqpoint{3.076650in}{2.066303in}}%
\pgfpathlineto{\pgfqpoint{3.077382in}{2.067597in}}%
\pgfpathlineto{\pgfqpoint{3.082777in}{2.079551in}}%
\pgfpathlineto{\pgfqpoint{3.086053in}{2.092799in}}%
\pgfpathlineto{\pgfqpoint{3.094254in}{2.106047in}}%
\pgfpathlineto{\pgfqpoint{3.096169in}{2.119295in}}%
\pgfpathlineto{\pgfqpoint{3.102666in}{2.129489in}}%
\pgfpathlineto{\pgfqpoint{3.104314in}{2.132543in}}%
\pgfpathlineto{\pgfqpoint{3.102666in}{2.132543in}}%
\pgfpathlineto{\pgfqpoint{3.087579in}{2.132543in}}%
\pgfpathlineto{\pgfqpoint{3.079750in}{2.119295in}}%
\pgfpathlineto{\pgfqpoint{3.078068in}{2.106047in}}%
\pgfpathlineto{\pgfqpoint{3.077382in}{2.104898in}}%
\pgfpathlineto{\pgfqpoint{3.067360in}{2.092799in}}%
\pgfpathlineto{\pgfqpoint{3.062854in}{2.079551in}}%
\pgfpathlineto{\pgfqpoint{3.056917in}{2.066303in}}%
\pgfpathlineto{\pgfqpoint{3.052099in}{2.053744in}}%
\pgfpathlineto{\pgfqpoint{3.051571in}{2.053055in}}%
\pgfpathlineto{\pgfqpoint{3.044965in}{2.039807in}}%
\pgfpathlineto{\pgfqpoint{3.033204in}{2.026559in}}%
\pgfpathlineto{\pgfqpoint{3.031066in}{2.013311in}}%
\pgfpathlineto{\pgfqpoint{3.026816in}{2.008369in}}%
\pgfpathlineto{\pgfqpoint{3.007985in}{2.000063in}}%
\pgfpathlineto{\pgfqpoint{3.011839in}{1.986815in}}%
\pgfpathlineto{\pgfqpoint{3.001533in}{1.980604in}}%
\pgfpathlineto{\pgfqpoint{2.976250in}{1.985484in}}%
\pgfpathlineto{\pgfqpoint{2.975224in}{1.986815in}}%
\pgfpathlineto{\pgfqpoint{2.976250in}{1.988598in}}%
\pgfpathlineto{\pgfqpoint{2.993754in}{2.000063in}}%
\pgfpathlineto{\pgfqpoint{2.976250in}{2.010607in}}%
\pgfpathlineto{\pgfqpoint{2.974916in}{2.013311in}}%
\pgfpathlineto{\pgfqpoint{2.976250in}{2.015680in}}%
\pgfpathlineto{\pgfqpoint{2.987184in}{2.026559in}}%
\pgfpathlineto{\pgfqpoint{2.977779in}{2.039807in}}%
\pgfpathlineto{\pgfqpoint{2.986302in}{2.053055in}}%
\pgfpathlineto{\pgfqpoint{2.982567in}{2.066303in}}%
\pgfpathlineto{\pgfqpoint{2.988854in}{2.079551in}}%
\pgfpathlineto{\pgfqpoint{2.987626in}{2.092799in}}%
\pgfpathlineto{\pgfqpoint{2.990893in}{2.106047in}}%
\pgfpathlineto{\pgfqpoint{2.991733in}{2.119295in}}%
\pgfpathlineto{\pgfqpoint{2.994736in}{2.132543in}}%
\pgfpathlineto{\pgfqpoint{2.978335in}{2.132543in}}%
\pgfpathlineto{\pgfqpoint{2.976250in}{2.126726in}}%
\pgfpathlineto{\pgfqpoint{2.973876in}{2.119295in}}%
\pgfpathlineto{\pgfqpoint{2.972845in}{2.106047in}}%
\pgfpathlineto{\pgfqpoint{2.968615in}{2.092799in}}%
\pgfpathlineto{\pgfqpoint{2.968706in}{2.079551in}}%
\pgfpathlineto{\pgfqpoint{2.962788in}{2.066303in}}%
\pgfpathlineto{\pgfqpoint{2.963983in}{2.053055in}}%
\pgfpathlineto{\pgfqpoint{2.957014in}{2.039807in}}%
\pgfpathlineto{\pgfqpoint{2.959723in}{2.026559in}}%
\pgfpathlineto{\pgfqpoint{2.950967in}{2.013918in}}%
\pgfpathlineto{\pgfqpoint{2.950630in}{2.013311in}}%
\pgfpathlineto{\pgfqpoint{2.950967in}{2.011483in}}%
\pgfpathlineto{\pgfqpoint{2.954287in}{2.000063in}}%
\pgfpathlineto{\pgfqpoint{2.950967in}{1.995378in}}%
\pgfpathlineto{\pgfqpoint{2.946652in}{1.986815in}}%
\pgfpathlineto{\pgfqpoint{2.950179in}{1.973567in}}%
\pgfpathlineto{\pgfqpoint{2.942445in}{1.960319in}}%
\pgfpathlineto{\pgfqpoint{2.944364in}{1.947071in}}%
\pgfpathlineto{\pgfqpoint{2.935853in}{1.933823in}}%
\pgfpathlineto{\pgfqpoint{2.939199in}{1.920575in}}%
\pgfpathlineto{\pgfqpoint{2.934967in}{1.907327in}}%
\pgfpathlineto{\pgfqpoint{2.941042in}{1.894079in}}%
\pgfpathlineto{\pgfqpoint{2.937769in}{1.880831in}}%
\pgfpathlineto{\pgfqpoint{2.943504in}{1.867583in}}%
\pgfpathlineto{\pgfqpoint{2.948343in}{1.854335in}}%
\pgfpathlineto{\pgfqpoint{2.950967in}{1.852116in}}%
\pgfpathclose%
\pgfpathmoveto{\pgfqpoint{3.709462in}{1.874805in}}%
\pgfpathlineto{\pgfqpoint{3.734231in}{1.880831in}}%
\pgfpathlineto{\pgfqpoint{3.715281in}{1.894079in}}%
\pgfpathlineto{\pgfqpoint{3.734745in}{1.905694in}}%
\pgfpathlineto{\pgfqpoint{3.735571in}{1.907327in}}%
\pgfpathlineto{\pgfqpoint{3.734745in}{1.909024in}}%
\pgfpathlineto{\pgfqpoint{3.724736in}{1.920575in}}%
\pgfpathlineto{\pgfqpoint{3.730646in}{1.933823in}}%
\pgfpathlineto{\pgfqpoint{3.720023in}{1.947071in}}%
\pgfpathlineto{\pgfqpoint{3.724123in}{1.960319in}}%
\pgfpathlineto{\pgfqpoint{3.719492in}{1.973567in}}%
\pgfpathlineto{\pgfqpoint{3.720730in}{1.986815in}}%
\pgfpathlineto{\pgfqpoint{3.718020in}{2.000063in}}%
\pgfpathlineto{\pgfqpoint{3.716588in}{2.013311in}}%
\pgfpathlineto{\pgfqpoint{3.714573in}{2.026559in}}%
\pgfpathlineto{\pgfqpoint{3.711672in}{2.039807in}}%
\pgfpathlineto{\pgfqpoint{3.710189in}{2.053055in}}%
\pgfpathlineto{\pgfqpoint{3.709462in}{2.054934in}}%
\pgfpathlineto{\pgfqpoint{3.704144in}{2.066303in}}%
\pgfpathlineto{\pgfqpoint{3.704324in}{2.079551in}}%
\pgfpathlineto{\pgfqpoint{3.699177in}{2.092799in}}%
\pgfpathlineto{\pgfqpoint{3.699715in}{2.106047in}}%
\pgfpathlineto{\pgfqpoint{3.693892in}{2.119295in}}%
\pgfpathlineto{\pgfqpoint{3.694092in}{2.132543in}}%
\pgfpathlineto{\pgfqpoint{3.684179in}{2.132543in}}%
\pgfpathlineto{\pgfqpoint{3.676915in}{2.132543in}}%
\pgfpathlineto{\pgfqpoint{3.675133in}{2.119295in}}%
\pgfpathlineto{\pgfqpoint{3.682078in}{2.106047in}}%
\pgfpathlineto{\pgfqpoint{3.679122in}{2.092799in}}%
\pgfpathlineto{\pgfqpoint{3.684179in}{2.081681in}}%
\pgfpathlineto{\pgfqpoint{3.684809in}{2.079551in}}%
\pgfpathlineto{\pgfqpoint{3.684179in}{2.073110in}}%
\pgfpathlineto{\pgfqpoint{3.682775in}{2.066303in}}%
\pgfpathlineto{\pgfqpoint{3.684179in}{2.063214in}}%
\pgfpathlineto{\pgfqpoint{3.686832in}{2.053055in}}%
\pgfpathlineto{\pgfqpoint{3.687321in}{2.039807in}}%
\pgfpathlineto{\pgfqpoint{3.686828in}{2.026559in}}%
\pgfpathlineto{\pgfqpoint{3.687627in}{2.013311in}}%
\pgfpathlineto{\pgfqpoint{3.684179in}{2.003584in}}%
\pgfpathlineto{\pgfqpoint{3.658895in}{2.004738in}}%
\pgfpathlineto{\pgfqpoint{3.650782in}{2.013311in}}%
\pgfpathlineto{\pgfqpoint{3.647249in}{2.026559in}}%
\pgfpathlineto{\pgfqpoint{3.633612in}{2.037063in}}%
\pgfpathlineto{\pgfqpoint{3.631783in}{2.039807in}}%
\pgfpathlineto{\pgfqpoint{3.628500in}{2.053055in}}%
\pgfpathlineto{\pgfqpoint{3.620903in}{2.066303in}}%
\pgfpathlineto{\pgfqpoint{3.614535in}{2.079551in}}%
\pgfpathlineto{\pgfqpoint{3.608329in}{2.089465in}}%
\pgfpathlineto{\pgfqpoint{3.606999in}{2.092799in}}%
\pgfpathlineto{\pgfqpoint{3.600651in}{2.106047in}}%
\pgfpathlineto{\pgfqpoint{3.596262in}{2.119295in}}%
\pgfpathlineto{\pgfqpoint{3.589598in}{2.132543in}}%
\pgfpathlineto{\pgfqpoint{3.583046in}{2.132543in}}%
\pgfpathlineto{\pgfqpoint{3.572592in}{2.132543in}}%
\pgfpathlineto{\pgfqpoint{3.577353in}{2.119295in}}%
\pgfpathlineto{\pgfqpoint{3.581868in}{2.106047in}}%
\pgfpathlineto{\pgfqpoint{3.583046in}{2.102924in}}%
\pgfpathlineto{\pgfqpoint{3.587941in}{2.092799in}}%
\pgfpathlineto{\pgfqpoint{3.594257in}{2.079551in}}%
\pgfpathlineto{\pgfqpoint{3.598306in}{2.066303in}}%
\pgfpathlineto{\pgfqpoint{3.605175in}{2.053055in}}%
\pgfpathlineto{\pgfqpoint{3.608075in}{2.039807in}}%
\pgfpathlineto{\pgfqpoint{3.608329in}{2.039371in}}%
\pgfpathlineto{\pgfqpoint{3.618073in}{2.026559in}}%
\pgfpathlineto{\pgfqpoint{3.621018in}{2.013311in}}%
\pgfpathlineto{\pgfqpoint{3.630201in}{2.000063in}}%
\pgfpathlineto{\pgfqpoint{3.632431in}{1.986815in}}%
\pgfpathlineto{\pgfqpoint{3.633612in}{1.985318in}}%
\pgfpathlineto{\pgfqpoint{3.647114in}{1.973567in}}%
\pgfpathlineto{\pgfqpoint{3.648201in}{1.960319in}}%
\pgfpathlineto{\pgfqpoint{3.658895in}{1.947921in}}%
\pgfpathlineto{\pgfqpoint{3.660404in}{1.947071in}}%
\pgfpathlineto{\pgfqpoint{3.658895in}{1.937893in}}%
\pgfpathlineto{\pgfqpoint{3.658607in}{1.933823in}}%
\pgfpathlineto{\pgfqpoint{3.658895in}{1.933445in}}%
\pgfpathlineto{\pgfqpoint{3.675326in}{1.920575in}}%
\pgfpathlineto{\pgfqpoint{3.675932in}{1.907327in}}%
\pgfpathlineto{\pgfqpoint{3.684179in}{1.900381in}}%
\pgfpathlineto{\pgfqpoint{3.703970in}{1.894079in}}%
\pgfpathlineto{\pgfqpoint{3.699264in}{1.880831in}}%
\pgfpathlineto{\pgfqpoint{3.709462in}{1.874805in}}%
\pgfpathclose%
\pgfusepath{fill}%
\end{pgfscope}%
\begin{pgfscope}%
\pgfpathrectangle{\pgfqpoint{2.470586in}{0.476543in}}{\pgfqpoint{1.820388in}{1.656000in}}%
\pgfusepath{clip}%
\pgfsetbuttcap%
\pgfsetroundjoin%
\definecolor{currentfill}{rgb}{0.699116,0.864975,0.921799}%
\pgfsetfillcolor{currentfill}%
\pgfsetlinewidth{0.000000pt}%
\definecolor{currentstroke}{rgb}{0.000000,0.000000,0.000000}%
\pgfsetstrokecolor{currentstroke}%
\pgfsetdash{}{0pt}%
\pgfpathmoveto{\pgfqpoint{3.153232in}{0.817704in}}%
\pgfpathlineto{\pgfqpoint{3.178515in}{0.815038in}}%
\pgfpathlineto{\pgfqpoint{3.203798in}{0.812989in}}%
\pgfpathlineto{\pgfqpoint{3.229082in}{0.811454in}}%
\pgfpathlineto{\pgfqpoint{3.254365in}{0.810359in}}%
\pgfpathlineto{\pgfqpoint{3.279648in}{0.809649in}}%
\pgfpathlineto{\pgfqpoint{3.304931in}{0.809287in}}%
\pgfpathlineto{\pgfqpoint{3.330214in}{0.809248in}}%
\pgfpathlineto{\pgfqpoint{3.355497in}{0.809516in}}%
\pgfpathlineto{\pgfqpoint{3.380781in}{0.810087in}}%
\pgfpathlineto{\pgfqpoint{3.406064in}{0.810963in}}%
\pgfpathlineto{\pgfqpoint{3.431347in}{0.812154in}}%
\pgfpathlineto{\pgfqpoint{3.456630in}{0.813681in}}%
\pgfpathlineto{\pgfqpoint{3.481913in}{0.815571in}}%
\pgfpathlineto{\pgfqpoint{3.507196in}{0.817862in}}%
\pgfpathlineto{\pgfqpoint{3.532480in}{0.820604in}}%
\pgfpathlineto{\pgfqpoint{3.535575in}{0.820991in}}%
\pgfpathlineto{\pgfqpoint{3.550160in}{0.834239in}}%
\pgfpathlineto{\pgfqpoint{3.557763in}{0.835342in}}%
\pgfpathlineto{\pgfqpoint{3.583046in}{0.839624in}}%
\pgfpathlineto{\pgfqpoint{3.608329in}{0.844809in}}%
\pgfpathlineto{\pgfqpoint{3.619507in}{0.847487in}}%
\pgfpathlineto{\pgfqpoint{3.619970in}{0.860735in}}%
\pgfpathlineto{\pgfqpoint{3.633612in}{0.864537in}}%
\pgfpathlineto{\pgfqpoint{3.658895in}{0.873177in}}%
\pgfpathlineto{\pgfqpoint{3.660925in}{0.873983in}}%
\pgfpathlineto{\pgfqpoint{3.669913in}{0.887231in}}%
\pgfpathlineto{\pgfqpoint{3.684179in}{0.897418in}}%
\pgfpathlineto{\pgfqpoint{3.687797in}{0.900479in}}%
\pgfpathlineto{\pgfqpoint{3.690696in}{0.913727in}}%
\pgfpathlineto{\pgfqpoint{3.695887in}{0.926975in}}%
\pgfpathlineto{\pgfqpoint{3.702554in}{0.940223in}}%
\pgfpathlineto{\pgfqpoint{3.701189in}{0.953471in}}%
\pgfpathlineto{\pgfqpoint{3.709462in}{0.966007in}}%
\pgfpathlineto{\pgfqpoint{3.709907in}{0.966719in}}%
\pgfpathlineto{\pgfqpoint{3.709462in}{0.967413in}}%
\pgfpathlineto{\pgfqpoint{3.701973in}{0.979967in}}%
\pgfpathlineto{\pgfqpoint{3.708255in}{0.993215in}}%
\pgfpathlineto{\pgfqpoint{3.698561in}{1.006463in}}%
\pgfpathlineto{\pgfqpoint{3.704396in}{1.019711in}}%
\pgfpathlineto{\pgfqpoint{3.692388in}{1.032959in}}%
\pgfpathlineto{\pgfqpoint{3.698156in}{1.046207in}}%
\pgfpathlineto{\pgfqpoint{3.688907in}{1.059455in}}%
\pgfpathlineto{\pgfqpoint{3.693255in}{1.072703in}}%
\pgfpathlineto{\pgfqpoint{3.685011in}{1.085951in}}%
\pgfpathlineto{\pgfqpoint{3.684767in}{1.099199in}}%
\pgfpathlineto{\pgfqpoint{3.684179in}{1.100011in}}%
\pgfpathlineto{\pgfqpoint{3.675593in}{1.112447in}}%
\pgfpathlineto{\pgfqpoint{3.674081in}{1.125695in}}%
\pgfpathlineto{\pgfqpoint{3.668792in}{1.138943in}}%
\pgfpathlineto{\pgfqpoint{3.665198in}{1.152191in}}%
\pgfpathlineto{\pgfqpoint{3.661096in}{1.165439in}}%
\pgfpathlineto{\pgfqpoint{3.658895in}{1.170512in}}%
\pgfpathlineto{\pgfqpoint{3.655684in}{1.178687in}}%
\pgfpathlineto{\pgfqpoint{3.652889in}{1.191935in}}%
\pgfpathlineto{\pgfqpoint{3.646436in}{1.205183in}}%
\pgfpathlineto{\pgfqpoint{3.644857in}{1.218431in}}%
\pgfpathlineto{\pgfqpoint{3.636447in}{1.231679in}}%
\pgfpathlineto{\pgfqpoint{3.636326in}{1.244927in}}%
\pgfpathlineto{\pgfqpoint{3.633612in}{1.249367in}}%
\pgfpathlineto{\pgfqpoint{3.628672in}{1.258175in}}%
\pgfpathlineto{\pgfqpoint{3.628163in}{1.271423in}}%
\pgfpathlineto{\pgfqpoint{3.619273in}{1.284671in}}%
\pgfpathlineto{\pgfqpoint{3.617860in}{1.297919in}}%
\pgfpathlineto{\pgfqpoint{3.608663in}{1.311167in}}%
\pgfpathlineto{\pgfqpoint{3.608329in}{1.321941in}}%
\pgfpathlineto{\pgfqpoint{3.608262in}{1.324415in}}%
\pgfpathlineto{\pgfqpoint{3.601625in}{1.337663in}}%
\pgfpathlineto{\pgfqpoint{3.600180in}{1.350911in}}%
\pgfpathlineto{\pgfqpoint{3.593483in}{1.364159in}}%
\pgfpathlineto{\pgfqpoint{3.591254in}{1.377407in}}%
\pgfpathlineto{\pgfqpoint{3.585063in}{1.390655in}}%
\pgfpathlineto{\pgfqpoint{3.583046in}{1.398213in}}%
\pgfpathlineto{\pgfqpoint{3.581726in}{1.403903in}}%
\pgfpathlineto{\pgfqpoint{3.576966in}{1.417151in}}%
\pgfpathlineto{\pgfqpoint{3.572606in}{1.430399in}}%
\pgfpathlineto{\pgfqpoint{3.569156in}{1.443647in}}%
\pgfpathlineto{\pgfqpoint{3.564226in}{1.456895in}}%
\pgfpathlineto{\pgfqpoint{3.560747in}{1.470143in}}%
\pgfpathlineto{\pgfqpoint{3.557763in}{1.475996in}}%
\pgfpathlineto{\pgfqpoint{3.554422in}{1.483391in}}%
\pgfpathlineto{\pgfqpoint{3.552083in}{1.496639in}}%
\pgfpathlineto{\pgfqpoint{3.546004in}{1.509887in}}%
\pgfpathlineto{\pgfqpoint{3.545026in}{1.523135in}}%
\pgfpathlineto{\pgfqpoint{3.538798in}{1.536383in}}%
\pgfpathlineto{\pgfqpoint{3.537057in}{1.549631in}}%
\pgfpathlineto{\pgfqpoint{3.532480in}{1.557773in}}%
\pgfpathlineto{\pgfqpoint{3.529875in}{1.562879in}}%
\pgfpathlineto{\pgfqpoint{3.528505in}{1.576127in}}%
\pgfpathlineto{\pgfqpoint{3.521668in}{1.589375in}}%
\pgfpathlineto{\pgfqpoint{3.520053in}{1.602623in}}%
\pgfpathlineto{\pgfqpoint{3.513401in}{1.615871in}}%
\pgfpathlineto{\pgfqpoint{3.511471in}{1.629119in}}%
\pgfpathlineto{\pgfqpoint{3.507196in}{1.638542in}}%
\pgfpathlineto{\pgfqpoint{3.505588in}{1.642367in}}%
\pgfpathlineto{\pgfqpoint{3.503429in}{1.655615in}}%
\pgfpathlineto{\pgfqpoint{3.497712in}{1.668863in}}%
\pgfpathlineto{\pgfqpoint{3.494555in}{1.682111in}}%
\pgfpathlineto{\pgfqpoint{3.489798in}{1.695359in}}%
\pgfpathlineto{\pgfqpoint{3.486371in}{1.708607in}}%
\pgfpathlineto{\pgfqpoint{3.482438in}{1.721855in}}%
\pgfpathlineto{\pgfqpoint{3.481913in}{1.723327in}}%
\pgfpathlineto{\pgfqpoint{3.478095in}{1.735103in}}%
\pgfpathlineto{\pgfqpoint{3.473946in}{1.748351in}}%
\pgfpathlineto{\pgfqpoint{3.468488in}{1.761599in}}%
\pgfpathlineto{\pgfqpoint{3.465344in}{1.774847in}}%
\pgfpathlineto{\pgfqpoint{3.459218in}{1.788095in}}%
\pgfpathlineto{\pgfqpoint{3.456630in}{1.801310in}}%
\pgfpathlineto{\pgfqpoint{3.456624in}{1.801343in}}%
\pgfpathlineto{\pgfqpoint{3.450211in}{1.814591in}}%
\pgfpathlineto{\pgfqpoint{3.448062in}{1.827839in}}%
\pgfpathlineto{\pgfqpoint{3.441296in}{1.841087in}}%
\pgfpathlineto{\pgfqpoint{3.439165in}{1.854335in}}%
\pgfpathlineto{\pgfqpoint{3.431347in}{1.867095in}}%
\pgfpathlineto{\pgfqpoint{3.431024in}{1.867583in}}%
\pgfpathlineto{\pgfqpoint{3.428700in}{1.880831in}}%
\pgfpathlineto{\pgfqpoint{3.420480in}{1.894079in}}%
\pgfpathlineto{\pgfqpoint{3.418389in}{1.907327in}}%
\pgfpathlineto{\pgfqpoint{3.409889in}{1.920575in}}%
\pgfpathlineto{\pgfqpoint{3.406634in}{1.933823in}}%
\pgfpathlineto{\pgfqpoint{3.406064in}{1.934618in}}%
\pgfpathlineto{\pgfqpoint{3.394736in}{1.947071in}}%
\pgfpathlineto{\pgfqpoint{3.390712in}{1.960319in}}%
\pgfpathlineto{\pgfqpoint{3.380781in}{1.972675in}}%
\pgfpathlineto{\pgfqpoint{3.379542in}{1.973567in}}%
\pgfpathlineto{\pgfqpoint{3.372823in}{1.986815in}}%
\pgfpathlineto{\pgfqpoint{3.355714in}{2.000063in}}%
\pgfpathlineto{\pgfqpoint{3.355497in}{2.000412in}}%
\pgfpathlineto{\pgfqpoint{3.330214in}{2.008604in}}%
\pgfpathlineto{\pgfqpoint{3.323439in}{2.000063in}}%
\pgfpathlineto{\pgfqpoint{3.304931in}{1.991517in}}%
\pgfpathlineto{\pgfqpoint{3.300132in}{1.986815in}}%
\pgfpathlineto{\pgfqpoint{3.296663in}{1.973567in}}%
\pgfpathlineto{\pgfqpoint{3.282451in}{1.960319in}}%
\pgfpathlineto{\pgfqpoint{3.279648in}{1.953227in}}%
\pgfpathlineto{\pgfqpoint{3.277804in}{1.947071in}}%
\pgfpathlineto{\pgfqpoint{3.267829in}{1.933823in}}%
\pgfpathlineto{\pgfqpoint{3.263646in}{1.920575in}}%
\pgfpathlineto{\pgfqpoint{3.256793in}{1.907327in}}%
\pgfpathlineto{\pgfqpoint{3.254365in}{1.900059in}}%
\pgfpathlineto{\pgfqpoint{3.252523in}{1.894079in}}%
\pgfpathlineto{\pgfqpoint{3.247037in}{1.880831in}}%
\pgfpathlineto{\pgfqpoint{3.241667in}{1.867583in}}%
\pgfpathlineto{\pgfqpoint{3.237788in}{1.854335in}}%
\pgfpathlineto{\pgfqpoint{3.231317in}{1.841087in}}%
\pgfpathlineto{\pgfqpoint{3.229082in}{1.831748in}}%
\pgfpathlineto{\pgfqpoint{3.228158in}{1.827839in}}%
\pgfpathlineto{\pgfqpoint{3.220956in}{1.814591in}}%
\pgfpathlineto{\pgfqpoint{3.218906in}{1.801343in}}%
\pgfpathlineto{\pgfqpoint{3.212262in}{1.788095in}}%
\pgfpathlineto{\pgfqpoint{3.210696in}{1.774847in}}%
\pgfpathlineto{\pgfqpoint{3.203798in}{1.762532in}}%
\pgfpathlineto{\pgfqpoint{3.203256in}{1.761599in}}%
\pgfpathlineto{\pgfqpoint{3.201216in}{1.748351in}}%
\pgfpathlineto{\pgfqpoint{3.194677in}{1.735103in}}%
\pgfpathlineto{\pgfqpoint{3.193082in}{1.721855in}}%
\pgfpathlineto{\pgfqpoint{3.187406in}{1.708607in}}%
\pgfpathlineto{\pgfqpoint{3.184901in}{1.695359in}}%
\pgfpathlineto{\pgfqpoint{3.179324in}{1.682111in}}%
\pgfpathlineto{\pgfqpoint{3.178515in}{1.678974in}}%
\pgfpathlineto{\pgfqpoint{3.175592in}{1.668863in}}%
\pgfpathlineto{\pgfqpoint{3.170928in}{1.655615in}}%
\pgfpathlineto{\pgfqpoint{3.166148in}{1.642367in}}%
\pgfpathlineto{\pgfqpoint{3.162593in}{1.629119in}}%
\pgfpathlineto{\pgfqpoint{3.157759in}{1.615871in}}%
\pgfpathlineto{\pgfqpoint{3.155828in}{1.602623in}}%
\pgfpathlineto{\pgfqpoint{3.153232in}{1.596350in}}%
\pgfpathlineto{\pgfqpoint{3.149995in}{1.589375in}}%
\pgfpathlineto{\pgfqpoint{3.148004in}{1.576127in}}%
\pgfpathlineto{\pgfqpoint{3.140540in}{1.562879in}}%
\pgfpathlineto{\pgfqpoint{3.139228in}{1.549631in}}%
\pgfpathlineto{\pgfqpoint{3.132710in}{1.536383in}}%
\pgfpathlineto{\pgfqpoint{3.131722in}{1.523135in}}%
\pgfpathlineto{\pgfqpoint{3.127949in}{1.516091in}}%
\pgfpathlineto{\pgfqpoint{3.124315in}{1.509887in}}%
\pgfpathlineto{\pgfqpoint{3.122659in}{1.496639in}}%
\pgfpathlineto{\pgfqpoint{3.115212in}{1.483391in}}%
\pgfpathlineto{\pgfqpoint{3.113449in}{1.470143in}}%
\pgfpathlineto{\pgfqpoint{3.107326in}{1.456895in}}%
\pgfpathlineto{\pgfqpoint{3.104823in}{1.443647in}}%
\pgfpathlineto{\pgfqpoint{3.102666in}{1.437998in}}%
\pgfpathlineto{\pgfqpoint{3.099510in}{1.430399in}}%
\pgfpathlineto{\pgfqpoint{3.096450in}{1.417151in}}%
\pgfpathlineto{\pgfqpoint{3.092129in}{1.403903in}}%
\pgfpathlineto{\pgfqpoint{3.087057in}{1.390655in}}%
\pgfpathlineto{\pgfqpoint{3.083020in}{1.377407in}}%
\pgfpathlineto{\pgfqpoint{3.077382in}{1.365955in}}%
\pgfpathlineto{\pgfqpoint{3.076384in}{1.364159in}}%
\pgfpathlineto{\pgfqpoint{3.074353in}{1.350911in}}%
\pgfpathlineto{\pgfqpoint{3.067551in}{1.337663in}}%
\pgfpathlineto{\pgfqpoint{3.066180in}{1.324415in}}%
\pgfpathlineto{\pgfqpoint{3.057767in}{1.311167in}}%
\pgfpathlineto{\pgfqpoint{3.057188in}{1.297919in}}%
\pgfpathlineto{\pgfqpoint{3.052099in}{1.290009in}}%
\pgfpathlineto{\pgfqpoint{3.048346in}{1.284671in}}%
\pgfpathlineto{\pgfqpoint{3.048772in}{1.271423in}}%
\pgfpathlineto{\pgfqpoint{3.039510in}{1.258175in}}%
\pgfpathlineto{\pgfqpoint{3.039644in}{1.244927in}}%
\pgfpathlineto{\pgfqpoint{3.030869in}{1.231679in}}%
\pgfpathlineto{\pgfqpoint{3.030894in}{1.218431in}}%
\pgfpathlineto{\pgfqpoint{3.026816in}{1.212511in}}%
\pgfpathlineto{\pgfqpoint{3.021458in}{1.205183in}}%
\pgfpathlineto{\pgfqpoint{3.019068in}{1.191935in}}%
\pgfpathlineto{\pgfqpoint{3.010735in}{1.178687in}}%
\pgfpathlineto{\pgfqpoint{3.008493in}{1.165439in}}%
\pgfpathlineto{\pgfqpoint{3.004205in}{1.152191in}}%
\pgfpathlineto{\pgfqpoint{3.001533in}{1.139511in}}%
\pgfpathlineto{\pgfqpoint{3.001398in}{1.138943in}}%
\pgfpathlineto{\pgfqpoint{2.996859in}{1.125695in}}%
\pgfpathlineto{\pgfqpoint{2.990689in}{1.112447in}}%
\pgfpathlineto{\pgfqpoint{2.988783in}{1.099199in}}%
\pgfpathlineto{\pgfqpoint{2.981200in}{1.085951in}}%
\pgfpathlineto{\pgfqpoint{2.981973in}{1.072703in}}%
\pgfpathlineto{\pgfqpoint{2.976250in}{1.064655in}}%
\pgfpathlineto{\pgfqpoint{2.972224in}{1.059455in}}%
\pgfpathlineto{\pgfqpoint{2.975593in}{1.046207in}}%
\pgfpathlineto{\pgfqpoint{2.964601in}{1.032959in}}%
\pgfpathlineto{\pgfqpoint{2.970561in}{1.019711in}}%
\pgfpathlineto{\pgfqpoint{2.956820in}{1.006463in}}%
\pgfpathlineto{\pgfqpoint{2.964461in}{0.993215in}}%
\pgfpathlineto{\pgfqpoint{2.952103in}{0.979967in}}%
\pgfpathlineto{\pgfqpoint{2.966090in}{0.966719in}}%
\pgfpathlineto{\pgfqpoint{2.956948in}{0.953471in}}%
\pgfpathlineto{\pgfqpoint{2.972847in}{0.940223in}}%
\pgfpathlineto{\pgfqpoint{2.959831in}{0.926975in}}%
\pgfpathlineto{\pgfqpoint{2.976250in}{0.914158in}}%
\pgfpathlineto{\pgfqpoint{2.976887in}{0.913727in}}%
\pgfpathlineto{\pgfqpoint{2.976250in}{0.912649in}}%
\pgfpathlineto{\pgfqpoint{2.968897in}{0.900479in}}%
\pgfpathlineto{\pgfqpoint{2.976250in}{0.896245in}}%
\pgfpathlineto{\pgfqpoint{2.994995in}{0.887231in}}%
\pgfpathlineto{\pgfqpoint{2.993357in}{0.873983in}}%
\pgfpathlineto{\pgfqpoint{3.001533in}{0.870268in}}%
\pgfpathlineto{\pgfqpoint{3.026816in}{0.861014in}}%
\pgfpathlineto{\pgfqpoint{3.027734in}{0.860735in}}%
\pgfpathlineto{\pgfqpoint{3.037372in}{0.847487in}}%
\pgfpathlineto{\pgfqpoint{3.052099in}{0.843154in}}%
\pgfpathlineto{\pgfqpoint{3.077382in}{0.837183in}}%
\pgfpathlineto{\pgfqpoint{3.092977in}{0.834239in}}%
\pgfpathlineto{\pgfqpoint{3.102666in}{0.830298in}}%
\pgfpathlineto{\pgfqpoint{3.127949in}{0.821292in}}%
\pgfpathlineto{\pgfqpoint{3.128912in}{0.820991in}}%
\pgfpathlineto{\pgfqpoint{3.153232in}{0.817704in}}%
\pgfpathclose%
\pgfpathmoveto{\pgfqpoint{3.148815in}{0.847487in}}%
\pgfpathlineto{\pgfqpoint{3.131551in}{0.860735in}}%
\pgfpathlineto{\pgfqpoint{3.127949in}{0.861274in}}%
\pgfpathlineto{\pgfqpoint{3.102666in}{0.865845in}}%
\pgfpathlineto{\pgfqpoint{3.077382in}{0.871696in}}%
\pgfpathlineto{\pgfqpoint{3.069193in}{0.873983in}}%
\pgfpathlineto{\pgfqpoint{3.067251in}{0.887231in}}%
\pgfpathlineto{\pgfqpoint{3.052099in}{0.891615in}}%
\pgfpathlineto{\pgfqpoint{3.027363in}{0.900479in}}%
\pgfpathlineto{\pgfqpoint{3.033171in}{0.913727in}}%
\pgfpathlineto{\pgfqpoint{3.026816in}{0.916720in}}%
\pgfpathlineto{\pgfqpoint{3.008357in}{0.926975in}}%
\pgfpathlineto{\pgfqpoint{3.019489in}{0.940223in}}%
\pgfpathlineto{\pgfqpoint{3.001533in}{0.951005in}}%
\pgfpathlineto{\pgfqpoint{2.997953in}{0.953471in}}%
\pgfpathlineto{\pgfqpoint{3.001533in}{0.959842in}}%
\pgfpathlineto{\pgfqpoint{3.005244in}{0.966719in}}%
\pgfpathlineto{\pgfqpoint{3.001533in}{0.969327in}}%
\pgfpathlineto{\pgfqpoint{2.987833in}{0.979967in}}%
\pgfpathlineto{\pgfqpoint{2.998699in}{0.993215in}}%
\pgfpathlineto{\pgfqpoint{2.988960in}{1.006463in}}%
\pgfpathlineto{\pgfqpoint{3.000912in}{1.019711in}}%
\pgfpathlineto{\pgfqpoint{2.993444in}{1.032959in}}%
\pgfpathlineto{\pgfqpoint{3.001533in}{1.044266in}}%
\pgfpathlineto{\pgfqpoint{3.002834in}{1.046207in}}%
\pgfpathlineto{\pgfqpoint{3.001533in}{1.050278in}}%
\pgfpathlineto{\pgfqpoint{2.998656in}{1.059455in}}%
\pgfpathlineto{\pgfqpoint{3.001533in}{1.063928in}}%
\pgfpathlineto{\pgfqpoint{3.006729in}{1.072703in}}%
\pgfpathlineto{\pgfqpoint{3.005363in}{1.085951in}}%
\pgfpathlineto{\pgfqpoint{3.011240in}{1.099199in}}%
\pgfpathlineto{\pgfqpoint{3.012863in}{1.112447in}}%
\pgfpathlineto{\pgfqpoint{3.017491in}{1.125695in}}%
\pgfpathlineto{\pgfqpoint{3.021790in}{1.138943in}}%
\pgfpathlineto{\pgfqpoint{3.023321in}{1.152191in}}%
\pgfpathlineto{\pgfqpoint{3.026816in}{1.162589in}}%
\pgfpathlineto{\pgfqpoint{3.027738in}{1.165439in}}%
\pgfpathlineto{\pgfqpoint{3.028871in}{1.178687in}}%
\pgfpathlineto{\pgfqpoint{3.037074in}{1.191935in}}%
\pgfpathlineto{\pgfqpoint{3.038393in}{1.205183in}}%
\pgfpathlineto{\pgfqpoint{3.047634in}{1.218431in}}%
\pgfpathlineto{\pgfqpoint{3.046782in}{1.231679in}}%
\pgfpathlineto{\pgfqpoint{3.052099in}{1.239508in}}%
\pgfpathlineto{\pgfqpoint{3.055595in}{1.244927in}}%
\pgfpathlineto{\pgfqpoint{3.054875in}{1.258175in}}%
\pgfpathlineto{\pgfqpoint{3.063556in}{1.271423in}}%
\pgfpathlineto{\pgfqpoint{3.062674in}{1.284671in}}%
\pgfpathlineto{\pgfqpoint{3.071196in}{1.297919in}}%
\pgfpathlineto{\pgfqpoint{3.071554in}{1.311167in}}%
\pgfpathlineto{\pgfqpoint{3.077382in}{1.320305in}}%
\pgfpathlineto{\pgfqpoint{3.079796in}{1.324415in}}%
\pgfpathlineto{\pgfqpoint{3.080988in}{1.337663in}}%
\pgfpathlineto{\pgfqpoint{3.087044in}{1.350911in}}%
\pgfpathlineto{\pgfqpoint{3.088963in}{1.364159in}}%
\pgfpathlineto{\pgfqpoint{3.095299in}{1.377407in}}%
\pgfpathlineto{\pgfqpoint{3.099568in}{1.390655in}}%
\pgfpathlineto{\pgfqpoint{3.102666in}{1.399156in}}%
\pgfpathlineto{\pgfqpoint{3.104200in}{1.403903in}}%
\pgfpathlineto{\pgfqpoint{3.108334in}{1.417151in}}%
\pgfpathlineto{\pgfqpoint{3.110770in}{1.430399in}}%
\pgfpathlineto{\pgfqpoint{3.116082in}{1.443647in}}%
\pgfpathlineto{\pgfqpoint{3.118266in}{1.456895in}}%
\pgfpathlineto{\pgfqpoint{3.124763in}{1.470143in}}%
\pgfpathlineto{\pgfqpoint{3.126242in}{1.483391in}}%
\pgfpathlineto{\pgfqpoint{3.127949in}{1.486340in}}%
\pgfpathlineto{\pgfqpoint{3.133521in}{1.496639in}}%
\pgfpathlineto{\pgfqpoint{3.134829in}{1.509887in}}%
\pgfpathlineto{\pgfqpoint{3.142105in}{1.523135in}}%
\pgfpathlineto{\pgfqpoint{3.142929in}{1.536383in}}%
\pgfpathlineto{\pgfqpoint{3.149687in}{1.549631in}}%
\pgfpathlineto{\pgfqpoint{3.150973in}{1.562879in}}%
\pgfpathlineto{\pgfqpoint{3.153232in}{1.566823in}}%
\pgfpathlineto{\pgfqpoint{3.158118in}{1.576127in}}%
\pgfpathlineto{\pgfqpoint{3.159958in}{1.589375in}}%
\pgfpathlineto{\pgfqpoint{3.165428in}{1.602623in}}%
\pgfpathlineto{\pgfqpoint{3.167533in}{1.615871in}}%
\pgfpathlineto{\pgfqpoint{3.172394in}{1.629119in}}%
\pgfpathlineto{\pgfqpoint{3.176285in}{1.642367in}}%
\pgfpathlineto{\pgfqpoint{3.178515in}{1.648578in}}%
\pgfpathlineto{\pgfqpoint{3.180818in}{1.655615in}}%
\pgfpathlineto{\pgfqpoint{3.185416in}{1.668863in}}%
\pgfpathlineto{\pgfqpoint{3.188710in}{1.682111in}}%
\pgfpathlineto{\pgfqpoint{3.194712in}{1.695359in}}%
\pgfpathlineto{\pgfqpoint{3.197055in}{1.708607in}}%
\pgfpathlineto{\pgfqpoint{3.203155in}{1.721855in}}%
\pgfpathlineto{\pgfqpoint{3.203798in}{1.727693in}}%
\pgfpathlineto{\pgfqpoint{3.204561in}{1.735103in}}%
\pgfpathlineto{\pgfqpoint{3.211149in}{1.748351in}}%
\pgfpathlineto{\pgfqpoint{3.213003in}{1.761599in}}%
\pgfpathlineto{\pgfqpoint{3.220735in}{1.774847in}}%
\pgfpathlineto{\pgfqpoint{3.222317in}{1.788095in}}%
\pgfpathlineto{\pgfqpoint{3.229082in}{1.800977in}}%
\pgfpathlineto{\pgfqpoint{3.229281in}{1.801343in}}%
\pgfpathlineto{\pgfqpoint{3.231536in}{1.814591in}}%
\pgfpathlineto{\pgfqpoint{3.238933in}{1.827839in}}%
\pgfpathlineto{\pgfqpoint{3.242369in}{1.841087in}}%
\pgfpathlineto{\pgfqpoint{3.248811in}{1.854335in}}%
\pgfpathlineto{\pgfqpoint{3.253128in}{1.867583in}}%
\pgfpathlineto{\pgfqpoint{3.254365in}{1.870637in}}%
\pgfpathlineto{\pgfqpoint{3.258930in}{1.880831in}}%
\pgfpathlineto{\pgfqpoint{3.265711in}{1.894079in}}%
\pgfpathlineto{\pgfqpoint{3.269890in}{1.907327in}}%
\pgfpathlineto{\pgfqpoint{3.277537in}{1.920575in}}%
\pgfpathlineto{\pgfqpoint{3.279648in}{1.927492in}}%
\pgfpathlineto{\pgfqpoint{3.282218in}{1.933823in}}%
\pgfpathlineto{\pgfqpoint{3.296954in}{1.947071in}}%
\pgfpathlineto{\pgfqpoint{3.301894in}{1.960319in}}%
\pgfpathlineto{\pgfqpoint{3.304931in}{1.963081in}}%
\pgfpathlineto{\pgfqpoint{3.329783in}{1.973567in}}%
\pgfpathlineto{\pgfqpoint{3.330214in}{1.974301in}}%
\pgfpathlineto{\pgfqpoint{3.331824in}{1.973567in}}%
\pgfpathlineto{\pgfqpoint{3.355497in}{1.970791in}}%
\pgfpathlineto{\pgfqpoint{3.370441in}{1.960319in}}%
\pgfpathlineto{\pgfqpoint{3.376652in}{1.947071in}}%
\pgfpathlineto{\pgfqpoint{3.380781in}{1.944549in}}%
\pgfpathlineto{\pgfqpoint{3.391143in}{1.933823in}}%
\pgfpathlineto{\pgfqpoint{3.394974in}{1.920575in}}%
\pgfpathlineto{\pgfqpoint{3.405922in}{1.907327in}}%
\pgfpathlineto{\pgfqpoint{3.406064in}{1.906568in}}%
\pgfpathlineto{\pgfqpoint{3.407986in}{1.894079in}}%
\pgfpathlineto{\pgfqpoint{3.416729in}{1.880831in}}%
\pgfpathlineto{\pgfqpoint{3.419095in}{1.867583in}}%
\pgfpathlineto{\pgfqpoint{3.427942in}{1.854335in}}%
\pgfpathlineto{\pgfqpoint{3.430280in}{1.841087in}}%
\pgfpathlineto{\pgfqpoint{3.431347in}{1.839227in}}%
\pgfpathlineto{\pgfqpoint{3.437378in}{1.827839in}}%
\pgfpathlineto{\pgfqpoint{3.439704in}{1.814591in}}%
\pgfpathlineto{\pgfqpoint{3.446327in}{1.801343in}}%
\pgfpathlineto{\pgfqpoint{3.449021in}{1.788095in}}%
\pgfpathlineto{\pgfqpoint{3.455074in}{1.774847in}}%
\pgfpathlineto{\pgfqpoint{3.456630in}{1.768783in}}%
\pgfpathlineto{\pgfqpoint{3.458463in}{1.761599in}}%
\pgfpathlineto{\pgfqpoint{3.464005in}{1.748351in}}%
\pgfpathlineto{\pgfqpoint{3.468478in}{1.735103in}}%
\pgfpathlineto{\pgfqpoint{3.472725in}{1.721855in}}%
\pgfpathlineto{\pgfqpoint{3.476598in}{1.708607in}}%
\pgfpathlineto{\pgfqpoint{3.479696in}{1.695359in}}%
\pgfpathlineto{\pgfqpoint{3.481913in}{1.689413in}}%
\pgfpathlineto{\pgfqpoint{3.484712in}{1.682111in}}%
\pgfpathlineto{\pgfqpoint{3.487769in}{1.668863in}}%
\pgfpathlineto{\pgfqpoint{3.493850in}{1.655615in}}%
\pgfpathlineto{\pgfqpoint{3.495866in}{1.642367in}}%
\pgfpathlineto{\pgfqpoint{3.501624in}{1.629119in}}%
\pgfpathlineto{\pgfqpoint{3.503275in}{1.615871in}}%
\pgfpathlineto{\pgfqpoint{3.507196in}{1.607943in}}%
\pgfpathlineto{\pgfqpoint{3.509978in}{1.602623in}}%
\pgfpathlineto{\pgfqpoint{3.511454in}{1.589375in}}%
\pgfpathlineto{\pgfqpoint{3.518591in}{1.576127in}}%
\pgfpathlineto{\pgfqpoint{3.519858in}{1.562879in}}%
\pgfpathlineto{\pgfqpoint{3.526727in}{1.549631in}}%
\pgfpathlineto{\pgfqpoint{3.528232in}{1.536383in}}%
\pgfpathlineto{\pgfqpoint{3.532480in}{1.526726in}}%
\pgfpathlineto{\pgfqpoint{3.534206in}{1.523135in}}%
\pgfpathlineto{\pgfqpoint{3.535122in}{1.509887in}}%
\pgfpathlineto{\pgfqpoint{3.541316in}{1.496639in}}%
\pgfpathlineto{\pgfqpoint{3.543769in}{1.483391in}}%
\pgfpathlineto{\pgfqpoint{3.549715in}{1.470143in}}%
\pgfpathlineto{\pgfqpoint{3.552940in}{1.456895in}}%
\pgfpathlineto{\pgfqpoint{3.557331in}{1.443647in}}%
\pgfpathlineto{\pgfqpoint{3.557763in}{1.441944in}}%
\pgfpathlineto{\pgfqpoint{3.560934in}{1.430399in}}%
\pgfpathlineto{\pgfqpoint{3.565139in}{1.417151in}}%
\pgfpathlineto{\pgfqpoint{3.570148in}{1.403903in}}%
\pgfpathlineto{\pgfqpoint{3.572894in}{1.390655in}}%
\pgfpathlineto{\pgfqpoint{3.578712in}{1.377407in}}%
\pgfpathlineto{\pgfqpoint{3.580418in}{1.364159in}}%
\pgfpathlineto{\pgfqpoint{3.583046in}{1.358802in}}%
\pgfpathlineto{\pgfqpoint{3.587184in}{1.350911in}}%
\pgfpathlineto{\pgfqpoint{3.588216in}{1.337663in}}%
\pgfpathlineto{\pgfqpoint{3.595103in}{1.324415in}}%
\pgfpathlineto{\pgfqpoint{3.594997in}{1.311167in}}%
\pgfpathlineto{\pgfqpoint{3.603624in}{1.297919in}}%
\pgfpathlineto{\pgfqpoint{3.604605in}{1.284671in}}%
\pgfpathlineto{\pgfqpoint{3.608329in}{1.278736in}}%
\pgfpathlineto{\pgfqpoint{3.613271in}{1.271423in}}%
\pgfpathlineto{\pgfqpoint{3.613341in}{1.258175in}}%
\pgfpathlineto{\pgfqpoint{3.620703in}{1.244927in}}%
\pgfpathlineto{\pgfqpoint{3.620456in}{1.231679in}}%
\pgfpathlineto{\pgfqpoint{3.628010in}{1.218431in}}%
\pgfpathlineto{\pgfqpoint{3.629243in}{1.205183in}}%
\pgfpathlineto{\pgfqpoint{3.633612in}{1.194621in}}%
\pgfpathlineto{\pgfqpoint{3.634834in}{1.191935in}}%
\pgfpathlineto{\pgfqpoint{3.637480in}{1.178687in}}%
\pgfpathlineto{\pgfqpoint{3.642069in}{1.165439in}}%
\pgfpathlineto{\pgfqpoint{3.645865in}{1.152191in}}%
\pgfpathlineto{\pgfqpoint{3.648197in}{1.138943in}}%
\pgfpathlineto{\pgfqpoint{3.653251in}{1.125695in}}%
\pgfpathlineto{\pgfqpoint{3.653359in}{1.112447in}}%
\pgfpathlineto{\pgfqpoint{3.658895in}{1.104097in}}%
\pgfpathlineto{\pgfqpoint{3.662268in}{1.099199in}}%
\pgfpathlineto{\pgfqpoint{3.660932in}{1.085951in}}%
\pgfpathlineto{\pgfqpoint{3.668454in}{1.072703in}}%
\pgfpathlineto{\pgfqpoint{3.662266in}{1.059455in}}%
\pgfpathlineto{\pgfqpoint{3.670903in}{1.046207in}}%
\pgfpathlineto{\pgfqpoint{3.663286in}{1.032959in}}%
\pgfpathlineto{\pgfqpoint{3.674016in}{1.019711in}}%
\pgfpathlineto{\pgfqpoint{3.666219in}{1.006463in}}%
\pgfpathlineto{\pgfqpoint{3.674061in}{0.993215in}}%
\pgfpathlineto{\pgfqpoint{3.665712in}{0.979967in}}%
\pgfpathlineto{\pgfqpoint{3.670941in}{0.966719in}}%
\pgfpathlineto{\pgfqpoint{3.659806in}{0.953471in}}%
\pgfpathlineto{\pgfqpoint{3.658895in}{0.949848in}}%
\pgfpathlineto{\pgfqpoint{3.656102in}{0.940223in}}%
\pgfpathlineto{\pgfqpoint{3.647128in}{0.926975in}}%
\pgfpathlineto{\pgfqpoint{3.634792in}{0.913727in}}%
\pgfpathlineto{\pgfqpoint{3.633612in}{0.910982in}}%
\pgfpathlineto{\pgfqpoint{3.628677in}{0.900479in}}%
\pgfpathlineto{\pgfqpoint{3.608329in}{0.891010in}}%
\pgfpathlineto{\pgfqpoint{3.598214in}{0.887231in}}%
\pgfpathlineto{\pgfqpoint{3.584610in}{0.873983in}}%
\pgfpathlineto{\pgfqpoint{3.583046in}{0.873574in}}%
\pgfpathlineto{\pgfqpoint{3.557763in}{0.867844in}}%
\pgfpathlineto{\pgfqpoint{3.532480in}{0.863227in}}%
\pgfpathlineto{\pgfqpoint{3.515873in}{0.860735in}}%
\pgfpathlineto{\pgfqpoint{3.507867in}{0.847487in}}%
\pgfpathlineto{\pgfqpoint{3.507196in}{0.847393in}}%
\pgfpathlineto{\pgfqpoint{3.481913in}{0.844350in}}%
\pgfpathlineto{\pgfqpoint{3.456630in}{0.841881in}}%
\pgfpathlineto{\pgfqpoint{3.431347in}{0.839917in}}%
\pgfpathlineto{\pgfqpoint{3.406064in}{0.838407in}}%
\pgfpathlineto{\pgfqpoint{3.380781in}{0.837317in}}%
\pgfpathlineto{\pgfqpoint{3.355497in}{0.836624in}}%
\pgfpathlineto{\pgfqpoint{3.330214in}{0.836316in}}%
\pgfpathlineto{\pgfqpoint{3.304931in}{0.836396in}}%
\pgfpathlineto{\pgfqpoint{3.279648in}{0.836876in}}%
\pgfpathlineto{\pgfqpoint{3.254365in}{0.837782in}}%
\pgfpathlineto{\pgfqpoint{3.229082in}{0.839155in}}%
\pgfpathlineto{\pgfqpoint{3.203798in}{0.841056in}}%
\pgfpathlineto{\pgfqpoint{3.178515in}{0.843567in}}%
\pgfpathlineto{\pgfqpoint{3.153232in}{0.846803in}}%
\pgfpathlineto{\pgfqpoint{3.148815in}{0.847487in}}%
\pgfpathclose%
\pgfpathmoveto{\pgfqpoint{2.976250in}{1.985484in}}%
\pgfpathlineto{\pgfqpoint{3.001533in}{1.980604in}}%
\pgfpathlineto{\pgfqpoint{3.011839in}{1.986815in}}%
\pgfpathlineto{\pgfqpoint{3.007985in}{2.000063in}}%
\pgfpathlineto{\pgfqpoint{3.026816in}{2.008369in}}%
\pgfpathlineto{\pgfqpoint{3.031066in}{2.013311in}}%
\pgfpathlineto{\pgfqpoint{3.033204in}{2.026559in}}%
\pgfpathlineto{\pgfqpoint{3.044965in}{2.039807in}}%
\pgfpathlineto{\pgfqpoint{3.051571in}{2.053055in}}%
\pgfpathlineto{\pgfqpoint{3.052099in}{2.053744in}}%
\pgfpathlineto{\pgfqpoint{3.056917in}{2.066303in}}%
\pgfpathlineto{\pgfqpoint{3.062854in}{2.079551in}}%
\pgfpathlineto{\pgfqpoint{3.067360in}{2.092799in}}%
\pgfpathlineto{\pgfqpoint{3.077382in}{2.104898in}}%
\pgfpathlineto{\pgfqpoint{3.078068in}{2.106047in}}%
\pgfpathlineto{\pgfqpoint{3.079750in}{2.119295in}}%
\pgfpathlineto{\pgfqpoint{3.087579in}{2.132543in}}%
\pgfpathlineto{\pgfqpoint{3.077382in}{2.132543in}}%
\pgfpathlineto{\pgfqpoint{3.064553in}{2.132543in}}%
\pgfpathlineto{\pgfqpoint{3.054502in}{2.119295in}}%
\pgfpathlineto{\pgfqpoint{3.052756in}{2.106047in}}%
\pgfpathlineto{\pgfqpoint{3.052099in}{2.105070in}}%
\pgfpathlineto{\pgfqpoint{3.031070in}{2.092799in}}%
\pgfpathlineto{\pgfqpoint{3.026816in}{2.087055in}}%
\pgfpathlineto{\pgfqpoint{3.023033in}{2.092799in}}%
\pgfpathlineto{\pgfqpoint{3.018664in}{2.106047in}}%
\pgfpathlineto{\pgfqpoint{3.018572in}{2.119295in}}%
\pgfpathlineto{\pgfqpoint{3.017784in}{2.132543in}}%
\pgfpathlineto{\pgfqpoint{3.001533in}{2.132543in}}%
\pgfpathlineto{\pgfqpoint{2.994736in}{2.132543in}}%
\pgfpathlineto{\pgfqpoint{2.991733in}{2.119295in}}%
\pgfpathlineto{\pgfqpoint{2.990893in}{2.106047in}}%
\pgfpathlineto{\pgfqpoint{2.987626in}{2.092799in}}%
\pgfpathlineto{\pgfqpoint{2.988854in}{2.079551in}}%
\pgfpathlineto{\pgfqpoint{2.982567in}{2.066303in}}%
\pgfpathlineto{\pgfqpoint{2.986302in}{2.053055in}}%
\pgfpathlineto{\pgfqpoint{2.977779in}{2.039807in}}%
\pgfpathlineto{\pgfqpoint{2.987184in}{2.026559in}}%
\pgfpathlineto{\pgfqpoint{2.976250in}{2.015680in}}%
\pgfpathlineto{\pgfqpoint{2.974916in}{2.013311in}}%
\pgfpathlineto{\pgfqpoint{2.976250in}{2.010607in}}%
\pgfpathlineto{\pgfqpoint{2.993754in}{2.000063in}}%
\pgfpathlineto{\pgfqpoint{2.976250in}{1.988598in}}%
\pgfpathlineto{\pgfqpoint{2.975224in}{1.986815in}}%
\pgfpathlineto{\pgfqpoint{2.976250in}{1.985484in}}%
\pgfpathclose%
\pgfpathmoveto{\pgfqpoint{3.658895in}{2.004738in}}%
\pgfpathlineto{\pgfqpoint{3.684179in}{2.003584in}}%
\pgfpathlineto{\pgfqpoint{3.687627in}{2.013311in}}%
\pgfpathlineto{\pgfqpoint{3.686828in}{2.026559in}}%
\pgfpathlineto{\pgfqpoint{3.687321in}{2.039807in}}%
\pgfpathlineto{\pgfqpoint{3.686832in}{2.053055in}}%
\pgfpathlineto{\pgfqpoint{3.684179in}{2.063214in}}%
\pgfpathlineto{\pgfqpoint{3.682775in}{2.066303in}}%
\pgfpathlineto{\pgfqpoint{3.684179in}{2.073110in}}%
\pgfpathlineto{\pgfqpoint{3.684809in}{2.079551in}}%
\pgfpathlineto{\pgfqpoint{3.684179in}{2.081681in}}%
\pgfpathlineto{\pgfqpoint{3.679122in}{2.092799in}}%
\pgfpathlineto{\pgfqpoint{3.682078in}{2.106047in}}%
\pgfpathlineto{\pgfqpoint{3.675133in}{2.119295in}}%
\pgfpathlineto{\pgfqpoint{3.676915in}{2.132543in}}%
\pgfpathlineto{\pgfqpoint{3.658895in}{2.132543in}}%
\pgfpathlineto{\pgfqpoint{3.650509in}{2.132543in}}%
\pgfpathlineto{\pgfqpoint{3.637436in}{2.119295in}}%
\pgfpathlineto{\pgfqpoint{3.640182in}{2.106047in}}%
\pgfpathlineto{\pgfqpoint{3.633612in}{2.103895in}}%
\pgfpathlineto{\pgfqpoint{3.631816in}{2.106047in}}%
\pgfpathlineto{\pgfqpoint{3.630815in}{2.119295in}}%
\pgfpathlineto{\pgfqpoint{3.615726in}{2.132543in}}%
\pgfpathlineto{\pgfqpoint{3.608329in}{2.132543in}}%
\pgfpathlineto{\pgfqpoint{3.589598in}{2.132543in}}%
\pgfpathlineto{\pgfqpoint{3.596262in}{2.119295in}}%
\pgfpathlineto{\pgfqpoint{3.600651in}{2.106047in}}%
\pgfpathlineto{\pgfqpoint{3.606999in}{2.092799in}}%
\pgfpathlineto{\pgfqpoint{3.608329in}{2.089465in}}%
\pgfpathlineto{\pgfqpoint{3.614535in}{2.079551in}}%
\pgfpathlineto{\pgfqpoint{3.620903in}{2.066303in}}%
\pgfpathlineto{\pgfqpoint{3.628500in}{2.053055in}}%
\pgfpathlineto{\pgfqpoint{3.631783in}{2.039807in}}%
\pgfpathlineto{\pgfqpoint{3.633612in}{2.037063in}}%
\pgfpathlineto{\pgfqpoint{3.647249in}{2.026559in}}%
\pgfpathlineto{\pgfqpoint{3.650782in}{2.013311in}}%
\pgfpathlineto{\pgfqpoint{3.658895in}{2.004738in}}%
\pgfpathclose%
\pgfusepath{fill}%
\end{pgfscope}%
\begin{pgfscope}%
\pgfpathrectangle{\pgfqpoint{2.470586in}{0.476543in}}{\pgfqpoint{1.820388in}{1.656000in}}%
\pgfusepath{clip}%
\pgfsetbuttcap%
\pgfsetroundjoin%
\definecolor{currentfill}{rgb}{0.780623,0.904960,0.944867}%
\pgfsetfillcolor{currentfill}%
\pgfsetlinewidth{0.000000pt}%
\definecolor{currentstroke}{rgb}{0.000000,0.000000,0.000000}%
\pgfsetstrokecolor{currentstroke}%
\pgfsetdash{}{0pt}%
\pgfpathmoveto{\pgfqpoint{3.153232in}{0.846803in}}%
\pgfpathlineto{\pgfqpoint{3.178515in}{0.843567in}}%
\pgfpathlineto{\pgfqpoint{3.203798in}{0.841056in}}%
\pgfpathlineto{\pgfqpoint{3.229082in}{0.839155in}}%
\pgfpathlineto{\pgfqpoint{3.254365in}{0.837782in}}%
\pgfpathlineto{\pgfqpoint{3.279648in}{0.836876in}}%
\pgfpathlineto{\pgfqpoint{3.304931in}{0.836396in}}%
\pgfpathlineto{\pgfqpoint{3.330214in}{0.836316in}}%
\pgfpathlineto{\pgfqpoint{3.355497in}{0.836624in}}%
\pgfpathlineto{\pgfqpoint{3.380781in}{0.837317in}}%
\pgfpathlineto{\pgfqpoint{3.406064in}{0.838407in}}%
\pgfpathlineto{\pgfqpoint{3.431347in}{0.839917in}}%
\pgfpathlineto{\pgfqpoint{3.456630in}{0.841881in}}%
\pgfpathlineto{\pgfqpoint{3.481913in}{0.844350in}}%
\pgfpathlineto{\pgfqpoint{3.507196in}{0.847393in}}%
\pgfpathlineto{\pgfqpoint{3.507867in}{0.847487in}}%
\pgfpathlineto{\pgfqpoint{3.515873in}{0.860735in}}%
\pgfpathlineto{\pgfqpoint{3.532480in}{0.863227in}}%
\pgfpathlineto{\pgfqpoint{3.557763in}{0.867844in}}%
\pgfpathlineto{\pgfqpoint{3.583046in}{0.873574in}}%
\pgfpathlineto{\pgfqpoint{3.584610in}{0.873983in}}%
\pgfpathlineto{\pgfqpoint{3.598214in}{0.887231in}}%
\pgfpathlineto{\pgfqpoint{3.608329in}{0.891010in}}%
\pgfpathlineto{\pgfqpoint{3.628677in}{0.900479in}}%
\pgfpathlineto{\pgfqpoint{3.633612in}{0.910982in}}%
\pgfpathlineto{\pgfqpoint{3.634792in}{0.913727in}}%
\pgfpathlineto{\pgfqpoint{3.647128in}{0.926975in}}%
\pgfpathlineto{\pgfqpoint{3.656102in}{0.940223in}}%
\pgfpathlineto{\pgfqpoint{3.658895in}{0.949848in}}%
\pgfpathlineto{\pgfqpoint{3.659806in}{0.953471in}}%
\pgfpathlineto{\pgfqpoint{3.670941in}{0.966719in}}%
\pgfpathlineto{\pgfqpoint{3.665712in}{0.979967in}}%
\pgfpathlineto{\pgfqpoint{3.674061in}{0.993215in}}%
\pgfpathlineto{\pgfqpoint{3.666219in}{1.006463in}}%
\pgfpathlineto{\pgfqpoint{3.674016in}{1.019711in}}%
\pgfpathlineto{\pgfqpoint{3.663286in}{1.032959in}}%
\pgfpathlineto{\pgfqpoint{3.670903in}{1.046207in}}%
\pgfpathlineto{\pgfqpoint{3.662266in}{1.059455in}}%
\pgfpathlineto{\pgfqpoint{3.668454in}{1.072703in}}%
\pgfpathlineto{\pgfqpoint{3.660932in}{1.085951in}}%
\pgfpathlineto{\pgfqpoint{3.662268in}{1.099199in}}%
\pgfpathlineto{\pgfqpoint{3.658895in}{1.104097in}}%
\pgfpathlineto{\pgfqpoint{3.653359in}{1.112447in}}%
\pgfpathlineto{\pgfqpoint{3.653251in}{1.125695in}}%
\pgfpathlineto{\pgfqpoint{3.648197in}{1.138943in}}%
\pgfpathlineto{\pgfqpoint{3.645865in}{1.152191in}}%
\pgfpathlineto{\pgfqpoint{3.642069in}{1.165439in}}%
\pgfpathlineto{\pgfqpoint{3.637480in}{1.178687in}}%
\pgfpathlineto{\pgfqpoint{3.634834in}{1.191935in}}%
\pgfpathlineto{\pgfqpoint{3.633612in}{1.194621in}}%
\pgfpathlineto{\pgfqpoint{3.629243in}{1.205183in}}%
\pgfpathlineto{\pgfqpoint{3.628010in}{1.218431in}}%
\pgfpathlineto{\pgfqpoint{3.620456in}{1.231679in}}%
\pgfpathlineto{\pgfqpoint{3.620703in}{1.244927in}}%
\pgfpathlineto{\pgfqpoint{3.613341in}{1.258175in}}%
\pgfpathlineto{\pgfqpoint{3.613271in}{1.271423in}}%
\pgfpathlineto{\pgfqpoint{3.608329in}{1.278736in}}%
\pgfpathlineto{\pgfqpoint{3.604605in}{1.284671in}}%
\pgfpathlineto{\pgfqpoint{3.603624in}{1.297919in}}%
\pgfpathlineto{\pgfqpoint{3.594997in}{1.311167in}}%
\pgfpathlineto{\pgfqpoint{3.595103in}{1.324415in}}%
\pgfpathlineto{\pgfqpoint{3.588216in}{1.337663in}}%
\pgfpathlineto{\pgfqpoint{3.587184in}{1.350911in}}%
\pgfpathlineto{\pgfqpoint{3.583046in}{1.358802in}}%
\pgfpathlineto{\pgfqpoint{3.580418in}{1.364159in}}%
\pgfpathlineto{\pgfqpoint{3.578712in}{1.377407in}}%
\pgfpathlineto{\pgfqpoint{3.572894in}{1.390655in}}%
\pgfpathlineto{\pgfqpoint{3.570148in}{1.403903in}}%
\pgfpathlineto{\pgfqpoint{3.565139in}{1.417151in}}%
\pgfpathlineto{\pgfqpoint{3.560934in}{1.430399in}}%
\pgfpathlineto{\pgfqpoint{3.557763in}{1.441944in}}%
\pgfpathlineto{\pgfqpoint{3.557331in}{1.443647in}}%
\pgfpathlineto{\pgfqpoint{3.552940in}{1.456895in}}%
\pgfpathlineto{\pgfqpoint{3.549715in}{1.470143in}}%
\pgfpathlineto{\pgfqpoint{3.543769in}{1.483391in}}%
\pgfpathlineto{\pgfqpoint{3.541316in}{1.496639in}}%
\pgfpathlineto{\pgfqpoint{3.535122in}{1.509887in}}%
\pgfpathlineto{\pgfqpoint{3.534206in}{1.523135in}}%
\pgfpathlineto{\pgfqpoint{3.532480in}{1.526726in}}%
\pgfpathlineto{\pgfqpoint{3.528232in}{1.536383in}}%
\pgfpathlineto{\pgfqpoint{3.526727in}{1.549631in}}%
\pgfpathlineto{\pgfqpoint{3.519858in}{1.562879in}}%
\pgfpathlineto{\pgfqpoint{3.518591in}{1.576127in}}%
\pgfpathlineto{\pgfqpoint{3.511454in}{1.589375in}}%
\pgfpathlineto{\pgfqpoint{3.509978in}{1.602623in}}%
\pgfpathlineto{\pgfqpoint{3.507196in}{1.607943in}}%
\pgfpathlineto{\pgfqpoint{3.503275in}{1.615871in}}%
\pgfpathlineto{\pgfqpoint{3.501624in}{1.629119in}}%
\pgfpathlineto{\pgfqpoint{3.495866in}{1.642367in}}%
\pgfpathlineto{\pgfqpoint{3.493850in}{1.655615in}}%
\pgfpathlineto{\pgfqpoint{3.487769in}{1.668863in}}%
\pgfpathlineto{\pgfqpoint{3.484712in}{1.682111in}}%
\pgfpathlineto{\pgfqpoint{3.481913in}{1.689413in}}%
\pgfpathlineto{\pgfqpoint{3.479696in}{1.695359in}}%
\pgfpathlineto{\pgfqpoint{3.476598in}{1.708607in}}%
\pgfpathlineto{\pgfqpoint{3.472725in}{1.721855in}}%
\pgfpathlineto{\pgfqpoint{3.468478in}{1.735103in}}%
\pgfpathlineto{\pgfqpoint{3.464005in}{1.748351in}}%
\pgfpathlineto{\pgfqpoint{3.458463in}{1.761599in}}%
\pgfpathlineto{\pgfqpoint{3.456630in}{1.768783in}}%
\pgfpathlineto{\pgfqpoint{3.455074in}{1.774847in}}%
\pgfpathlineto{\pgfqpoint{3.449021in}{1.788095in}}%
\pgfpathlineto{\pgfqpoint{3.446327in}{1.801343in}}%
\pgfpathlineto{\pgfqpoint{3.439704in}{1.814591in}}%
\pgfpathlineto{\pgfqpoint{3.437378in}{1.827839in}}%
\pgfpathlineto{\pgfqpoint{3.431347in}{1.839227in}}%
\pgfpathlineto{\pgfqpoint{3.430280in}{1.841087in}}%
\pgfpathlineto{\pgfqpoint{3.427942in}{1.854335in}}%
\pgfpathlineto{\pgfqpoint{3.419095in}{1.867583in}}%
\pgfpathlineto{\pgfqpoint{3.416729in}{1.880831in}}%
\pgfpathlineto{\pgfqpoint{3.407986in}{1.894079in}}%
\pgfpathlineto{\pgfqpoint{3.406064in}{1.906568in}}%
\pgfpathlineto{\pgfqpoint{3.405922in}{1.907327in}}%
\pgfpathlineto{\pgfqpoint{3.394974in}{1.920575in}}%
\pgfpathlineto{\pgfqpoint{3.391143in}{1.933823in}}%
\pgfpathlineto{\pgfqpoint{3.380781in}{1.944549in}}%
\pgfpathlineto{\pgfqpoint{3.376652in}{1.947071in}}%
\pgfpathlineto{\pgfqpoint{3.370441in}{1.960319in}}%
\pgfpathlineto{\pgfqpoint{3.355497in}{1.970791in}}%
\pgfpathlineto{\pgfqpoint{3.331824in}{1.973567in}}%
\pgfpathlineto{\pgfqpoint{3.330214in}{1.974301in}}%
\pgfpathlineto{\pgfqpoint{3.329783in}{1.973567in}}%
\pgfpathlineto{\pgfqpoint{3.304931in}{1.963081in}}%
\pgfpathlineto{\pgfqpoint{3.301894in}{1.960319in}}%
\pgfpathlineto{\pgfqpoint{3.296954in}{1.947071in}}%
\pgfpathlineto{\pgfqpoint{3.282218in}{1.933823in}}%
\pgfpathlineto{\pgfqpoint{3.279648in}{1.927492in}}%
\pgfpathlineto{\pgfqpoint{3.277537in}{1.920575in}}%
\pgfpathlineto{\pgfqpoint{3.269890in}{1.907327in}}%
\pgfpathlineto{\pgfqpoint{3.265711in}{1.894079in}}%
\pgfpathlineto{\pgfqpoint{3.258930in}{1.880831in}}%
\pgfpathlineto{\pgfqpoint{3.254365in}{1.870637in}}%
\pgfpathlineto{\pgfqpoint{3.253128in}{1.867583in}}%
\pgfpathlineto{\pgfqpoint{3.248811in}{1.854335in}}%
\pgfpathlineto{\pgfqpoint{3.242369in}{1.841087in}}%
\pgfpathlineto{\pgfqpoint{3.238933in}{1.827839in}}%
\pgfpathlineto{\pgfqpoint{3.231536in}{1.814591in}}%
\pgfpathlineto{\pgfqpoint{3.229281in}{1.801343in}}%
\pgfpathlineto{\pgfqpoint{3.229082in}{1.800977in}}%
\pgfpathlineto{\pgfqpoint{3.222317in}{1.788095in}}%
\pgfpathlineto{\pgfqpoint{3.220735in}{1.774847in}}%
\pgfpathlineto{\pgfqpoint{3.213003in}{1.761599in}}%
\pgfpathlineto{\pgfqpoint{3.211149in}{1.748351in}}%
\pgfpathlineto{\pgfqpoint{3.204561in}{1.735103in}}%
\pgfpathlineto{\pgfqpoint{3.203798in}{1.727693in}}%
\pgfpathlineto{\pgfqpoint{3.203155in}{1.721855in}}%
\pgfpathlineto{\pgfqpoint{3.197055in}{1.708607in}}%
\pgfpathlineto{\pgfqpoint{3.194712in}{1.695359in}}%
\pgfpathlineto{\pgfqpoint{3.188710in}{1.682111in}}%
\pgfpathlineto{\pgfqpoint{3.185416in}{1.668863in}}%
\pgfpathlineto{\pgfqpoint{3.180818in}{1.655615in}}%
\pgfpathlineto{\pgfqpoint{3.178515in}{1.648578in}}%
\pgfpathlineto{\pgfqpoint{3.176285in}{1.642367in}}%
\pgfpathlineto{\pgfqpoint{3.172394in}{1.629119in}}%
\pgfpathlineto{\pgfqpoint{3.167533in}{1.615871in}}%
\pgfpathlineto{\pgfqpoint{3.165428in}{1.602623in}}%
\pgfpathlineto{\pgfqpoint{3.159958in}{1.589375in}}%
\pgfpathlineto{\pgfqpoint{3.158118in}{1.576127in}}%
\pgfpathlineto{\pgfqpoint{3.153232in}{1.566823in}}%
\pgfpathlineto{\pgfqpoint{3.150973in}{1.562879in}}%
\pgfpathlineto{\pgfqpoint{3.149687in}{1.549631in}}%
\pgfpathlineto{\pgfqpoint{3.142929in}{1.536383in}}%
\pgfpathlineto{\pgfqpoint{3.142105in}{1.523135in}}%
\pgfpathlineto{\pgfqpoint{3.134829in}{1.509887in}}%
\pgfpathlineto{\pgfqpoint{3.133521in}{1.496639in}}%
\pgfpathlineto{\pgfqpoint{3.127949in}{1.486340in}}%
\pgfpathlineto{\pgfqpoint{3.126242in}{1.483391in}}%
\pgfpathlineto{\pgfqpoint{3.124763in}{1.470143in}}%
\pgfpathlineto{\pgfqpoint{3.118266in}{1.456895in}}%
\pgfpathlineto{\pgfqpoint{3.116082in}{1.443647in}}%
\pgfpathlineto{\pgfqpoint{3.110770in}{1.430399in}}%
\pgfpathlineto{\pgfqpoint{3.108334in}{1.417151in}}%
\pgfpathlineto{\pgfqpoint{3.104200in}{1.403903in}}%
\pgfpathlineto{\pgfqpoint{3.102666in}{1.399156in}}%
\pgfpathlineto{\pgfqpoint{3.099568in}{1.390655in}}%
\pgfpathlineto{\pgfqpoint{3.095299in}{1.377407in}}%
\pgfpathlineto{\pgfqpoint{3.088963in}{1.364159in}}%
\pgfpathlineto{\pgfqpoint{3.087044in}{1.350911in}}%
\pgfpathlineto{\pgfqpoint{3.080988in}{1.337663in}}%
\pgfpathlineto{\pgfqpoint{3.079796in}{1.324415in}}%
\pgfpathlineto{\pgfqpoint{3.077382in}{1.320305in}}%
\pgfpathlineto{\pgfqpoint{3.071554in}{1.311167in}}%
\pgfpathlineto{\pgfqpoint{3.071196in}{1.297919in}}%
\pgfpathlineto{\pgfqpoint{3.062674in}{1.284671in}}%
\pgfpathlineto{\pgfqpoint{3.063556in}{1.271423in}}%
\pgfpathlineto{\pgfqpoint{3.054875in}{1.258175in}}%
\pgfpathlineto{\pgfqpoint{3.055595in}{1.244927in}}%
\pgfpathlineto{\pgfqpoint{3.052099in}{1.239508in}}%
\pgfpathlineto{\pgfqpoint{3.046782in}{1.231679in}}%
\pgfpathlineto{\pgfqpoint{3.047634in}{1.218431in}}%
\pgfpathlineto{\pgfqpoint{3.038393in}{1.205183in}}%
\pgfpathlineto{\pgfqpoint{3.037074in}{1.191935in}}%
\pgfpathlineto{\pgfqpoint{3.028871in}{1.178687in}}%
\pgfpathlineto{\pgfqpoint{3.027738in}{1.165439in}}%
\pgfpathlineto{\pgfqpoint{3.026816in}{1.162589in}}%
\pgfpathlineto{\pgfqpoint{3.023321in}{1.152191in}}%
\pgfpathlineto{\pgfqpoint{3.021790in}{1.138943in}}%
\pgfpathlineto{\pgfqpoint{3.017491in}{1.125695in}}%
\pgfpathlineto{\pgfqpoint{3.012863in}{1.112447in}}%
\pgfpathlineto{\pgfqpoint{3.011240in}{1.099199in}}%
\pgfpathlineto{\pgfqpoint{3.005363in}{1.085951in}}%
\pgfpathlineto{\pgfqpoint{3.006729in}{1.072703in}}%
\pgfpathlineto{\pgfqpoint{3.001533in}{1.063928in}}%
\pgfpathlineto{\pgfqpoint{2.998656in}{1.059455in}}%
\pgfpathlineto{\pgfqpoint{3.001533in}{1.050278in}}%
\pgfpathlineto{\pgfqpoint{3.002834in}{1.046207in}}%
\pgfpathlineto{\pgfqpoint{3.001533in}{1.044266in}}%
\pgfpathlineto{\pgfqpoint{2.993444in}{1.032959in}}%
\pgfpathlineto{\pgfqpoint{3.000912in}{1.019711in}}%
\pgfpathlineto{\pgfqpoint{2.988960in}{1.006463in}}%
\pgfpathlineto{\pgfqpoint{2.998699in}{0.993215in}}%
\pgfpathlineto{\pgfqpoint{2.987833in}{0.979967in}}%
\pgfpathlineto{\pgfqpoint{3.001533in}{0.969327in}}%
\pgfpathlineto{\pgfqpoint{3.005244in}{0.966719in}}%
\pgfpathlineto{\pgfqpoint{3.001533in}{0.959842in}}%
\pgfpathlineto{\pgfqpoint{2.997953in}{0.953471in}}%
\pgfpathlineto{\pgfqpoint{3.001533in}{0.951005in}}%
\pgfpathlineto{\pgfqpoint{3.019489in}{0.940223in}}%
\pgfpathlineto{\pgfqpoint{3.008357in}{0.926975in}}%
\pgfpathlineto{\pgfqpoint{3.026816in}{0.916720in}}%
\pgfpathlineto{\pgfqpoint{3.033171in}{0.913727in}}%
\pgfpathlineto{\pgfqpoint{3.027363in}{0.900479in}}%
\pgfpathlineto{\pgfqpoint{3.052099in}{0.891615in}}%
\pgfpathlineto{\pgfqpoint{3.067251in}{0.887231in}}%
\pgfpathlineto{\pgfqpoint{3.069193in}{0.873983in}}%
\pgfpathlineto{\pgfqpoint{3.077382in}{0.871696in}}%
\pgfpathlineto{\pgfqpoint{3.102666in}{0.865845in}}%
\pgfpathlineto{\pgfqpoint{3.127949in}{0.861274in}}%
\pgfpathlineto{\pgfqpoint{3.131551in}{0.860735in}}%
\pgfpathlineto{\pgfqpoint{3.148815in}{0.847487in}}%
\pgfpathlineto{\pgfqpoint{3.153232in}{0.846803in}}%
\pgfpathclose%
\pgfpathmoveto{\pgfqpoint{3.166437in}{0.873983in}}%
\pgfpathlineto{\pgfqpoint{3.155985in}{0.887231in}}%
\pgfpathlineto{\pgfqpoint{3.153232in}{0.887645in}}%
\pgfpathlineto{\pgfqpoint{3.127949in}{0.892193in}}%
\pgfpathlineto{\pgfqpoint{3.102666in}{0.897860in}}%
\pgfpathlineto{\pgfqpoint{3.092882in}{0.900479in}}%
\pgfpathlineto{\pgfqpoint{3.095137in}{0.913727in}}%
\pgfpathlineto{\pgfqpoint{3.077382in}{0.919772in}}%
\pgfpathlineto{\pgfqpoint{3.059863in}{0.926975in}}%
\pgfpathlineto{\pgfqpoint{3.068604in}{0.940223in}}%
\pgfpathlineto{\pgfqpoint{3.052099in}{0.947410in}}%
\pgfpathlineto{\pgfqpoint{3.040234in}{0.953471in}}%
\pgfpathlineto{\pgfqpoint{3.045179in}{0.966719in}}%
\pgfpathlineto{\pgfqpoint{3.026816in}{0.977964in}}%
\pgfpathlineto{\pgfqpoint{3.023884in}{0.979967in}}%
\pgfpathlineto{\pgfqpoint{3.026816in}{0.984100in}}%
\pgfpathlineto{\pgfqpoint{3.033024in}{0.993215in}}%
\pgfpathlineto{\pgfqpoint{3.026816in}{0.999810in}}%
\pgfpathlineto{\pgfqpoint{3.021031in}{1.006463in}}%
\pgfpathlineto{\pgfqpoint{3.026816in}{1.013955in}}%
\pgfpathlineto{\pgfqpoint{3.031007in}{1.019711in}}%
\pgfpathlineto{\pgfqpoint{3.026816in}{1.025655in}}%
\pgfpathlineto{\pgfqpoint{3.021972in}{1.032959in}}%
\pgfpathlineto{\pgfqpoint{3.026816in}{1.041130in}}%
\pgfpathlineto{\pgfqpoint{3.029646in}{1.046207in}}%
\pgfpathlineto{\pgfqpoint{3.026816in}{1.053567in}}%
\pgfpathlineto{\pgfqpoint{3.024625in}{1.059455in}}%
\pgfpathlineto{\pgfqpoint{3.026816in}{1.063842in}}%
\pgfpathlineto{\pgfqpoint{3.030906in}{1.072703in}}%
\pgfpathlineto{\pgfqpoint{3.028932in}{1.085951in}}%
\pgfpathlineto{\pgfqpoint{3.033075in}{1.099199in}}%
\pgfpathlineto{\pgfqpoint{3.034430in}{1.112447in}}%
\pgfpathlineto{\pgfqpoint{3.037547in}{1.125695in}}%
\pgfpathlineto{\pgfqpoint{3.041648in}{1.138943in}}%
\pgfpathlineto{\pgfqpoint{3.041775in}{1.152191in}}%
\pgfpathlineto{\pgfqpoint{3.046130in}{1.165439in}}%
\pgfpathlineto{\pgfqpoint{3.046157in}{1.178687in}}%
\pgfpathlineto{\pgfqpoint{3.052099in}{1.188043in}}%
\pgfpathlineto{\pgfqpoint{3.054511in}{1.191935in}}%
\pgfpathlineto{\pgfqpoint{3.054844in}{1.205183in}}%
\pgfpathlineto{\pgfqpoint{3.063726in}{1.218431in}}%
\pgfpathlineto{\pgfqpoint{3.062102in}{1.231679in}}%
\pgfpathlineto{\pgfqpoint{3.070706in}{1.244927in}}%
\pgfpathlineto{\pgfqpoint{3.069360in}{1.258175in}}%
\pgfpathlineto{\pgfqpoint{3.077382in}{1.270392in}}%
\pgfpathlineto{\pgfqpoint{3.078023in}{1.271423in}}%
\pgfpathlineto{\pgfqpoint{3.077382in}{1.277820in}}%
\pgfpathlineto{\pgfqpoint{3.076681in}{1.284671in}}%
\pgfpathlineto{\pgfqpoint{3.077382in}{1.285762in}}%
\pgfpathlineto{\pgfqpoint{3.084678in}{1.297919in}}%
\pgfpathlineto{\pgfqpoint{3.084804in}{1.311167in}}%
\pgfpathlineto{\pgfqpoint{3.092493in}{1.324415in}}%
\pgfpathlineto{\pgfqpoint{3.093603in}{1.337663in}}%
\pgfpathlineto{\pgfqpoint{3.099435in}{1.350911in}}%
\pgfpathlineto{\pgfqpoint{3.101441in}{1.364159in}}%
\pgfpathlineto{\pgfqpoint{3.102666in}{1.366770in}}%
\pgfpathlineto{\pgfqpoint{3.107203in}{1.377407in}}%
\pgfpathlineto{\pgfqpoint{3.111321in}{1.390655in}}%
\pgfpathlineto{\pgfqpoint{3.115292in}{1.403903in}}%
\pgfpathlineto{\pgfqpoint{3.119609in}{1.417151in}}%
\pgfpathlineto{\pgfqpoint{3.121689in}{1.430399in}}%
\pgfpathlineto{\pgfqpoint{3.127342in}{1.443647in}}%
\pgfpathlineto{\pgfqpoint{3.127949in}{1.447878in}}%
\pgfpathlineto{\pgfqpoint{3.129109in}{1.456895in}}%
\pgfpathlineto{\pgfqpoint{3.135367in}{1.470143in}}%
\pgfpathlineto{\pgfqpoint{3.136417in}{1.483391in}}%
\pgfpathlineto{\pgfqpoint{3.143833in}{1.496639in}}%
\pgfpathlineto{\pgfqpoint{3.144947in}{1.509887in}}%
\pgfpathlineto{\pgfqpoint{3.152488in}{1.523135in}}%
\pgfpathlineto{\pgfqpoint{3.153148in}{1.536383in}}%
\pgfpathlineto{\pgfqpoint{3.153232in}{1.536547in}}%
\pgfpathlineto{\pgfqpoint{3.159549in}{1.549631in}}%
\pgfpathlineto{\pgfqpoint{3.160697in}{1.562879in}}%
\pgfpathlineto{\pgfqpoint{3.167684in}{1.576127in}}%
\pgfpathlineto{\pgfqpoint{3.169577in}{1.589375in}}%
\pgfpathlineto{\pgfqpoint{3.175027in}{1.602623in}}%
\pgfpathlineto{\pgfqpoint{3.177307in}{1.615871in}}%
\pgfpathlineto{\pgfqpoint{3.178515in}{1.619152in}}%
\pgfpathlineto{\pgfqpoint{3.181944in}{1.629119in}}%
\pgfpathlineto{\pgfqpoint{3.185857in}{1.642367in}}%
\pgfpathlineto{\pgfqpoint{3.190052in}{1.655615in}}%
\pgfpathlineto{\pgfqpoint{3.194971in}{1.668863in}}%
\pgfpathlineto{\pgfqpoint{3.198095in}{1.682111in}}%
\pgfpathlineto{\pgfqpoint{3.203798in}{1.693925in}}%
\pgfpathlineto{\pgfqpoint{3.204503in}{1.695359in}}%
\pgfpathlineto{\pgfqpoint{3.206600in}{1.708607in}}%
\pgfpathlineto{\pgfqpoint{3.212780in}{1.721855in}}%
\pgfpathlineto{\pgfqpoint{3.213973in}{1.735103in}}%
\pgfpathlineto{\pgfqpoint{3.220912in}{1.748351in}}%
\pgfpathlineto{\pgfqpoint{3.222712in}{1.761599in}}%
\pgfpathlineto{\pgfqpoint{3.229082in}{1.772138in}}%
\pgfpathlineto{\pgfqpoint{3.230817in}{1.774847in}}%
\pgfpathlineto{\pgfqpoint{3.232465in}{1.788095in}}%
\pgfpathlineto{\pgfqpoint{3.239738in}{1.801343in}}%
\pgfpathlineto{\pgfqpoint{3.242211in}{1.814591in}}%
\pgfpathlineto{\pgfqpoint{3.249696in}{1.827839in}}%
\pgfpathlineto{\pgfqpoint{3.253420in}{1.841087in}}%
\pgfpathlineto{\pgfqpoint{3.254365in}{1.843035in}}%
\pgfpathlineto{\pgfqpoint{3.260492in}{1.854335in}}%
\pgfpathlineto{\pgfqpoint{3.265905in}{1.867583in}}%
\pgfpathlineto{\pgfqpoint{3.271608in}{1.880831in}}%
\pgfpathlineto{\pgfqpoint{3.279099in}{1.894079in}}%
\pgfpathlineto{\pgfqpoint{3.279648in}{1.895913in}}%
\pgfpathlineto{\pgfqpoint{3.284159in}{1.907327in}}%
\pgfpathlineto{\pgfqpoint{3.295653in}{1.920575in}}%
\pgfpathlineto{\pgfqpoint{3.300701in}{1.933823in}}%
\pgfpathlineto{\pgfqpoint{3.304931in}{1.937496in}}%
\pgfpathlineto{\pgfqpoint{3.330214in}{1.947066in}}%
\pgfpathlineto{\pgfqpoint{3.355497in}{1.944026in}}%
\pgfpathlineto{\pgfqpoint{3.372336in}{1.933823in}}%
\pgfpathlineto{\pgfqpoint{3.378368in}{1.920575in}}%
\pgfpathlineto{\pgfqpoint{3.380781in}{1.918898in}}%
\pgfpathlineto{\pgfqpoint{3.390757in}{1.907327in}}%
\pgfpathlineto{\pgfqpoint{3.393176in}{1.894079in}}%
\pgfpathlineto{\pgfqpoint{3.404446in}{1.880831in}}%
\pgfpathlineto{\pgfqpoint{3.406064in}{1.873655in}}%
\pgfpathlineto{\pgfqpoint{3.407165in}{1.867583in}}%
\pgfpathlineto{\pgfqpoint{3.416264in}{1.854335in}}%
\pgfpathlineto{\pgfqpoint{3.418708in}{1.841087in}}%
\pgfpathlineto{\pgfqpoint{3.426350in}{1.827839in}}%
\pgfpathlineto{\pgfqpoint{3.429050in}{1.814591in}}%
\pgfpathlineto{\pgfqpoint{3.431347in}{1.810480in}}%
\pgfpathlineto{\pgfqpoint{3.436031in}{1.801343in}}%
\pgfpathlineto{\pgfqpoint{3.438935in}{1.788095in}}%
\pgfpathlineto{\pgfqpoint{3.444996in}{1.774847in}}%
\pgfpathlineto{\pgfqpoint{3.448628in}{1.761599in}}%
\pgfpathlineto{\pgfqpoint{3.454080in}{1.748351in}}%
\pgfpathlineto{\pgfqpoint{3.456630in}{1.741373in}}%
\pgfpathlineto{\pgfqpoint{3.458861in}{1.735103in}}%
\pgfpathlineto{\pgfqpoint{3.463056in}{1.721855in}}%
\pgfpathlineto{\pgfqpoint{3.467187in}{1.708607in}}%
\pgfpathlineto{\pgfqpoint{3.470122in}{1.695359in}}%
\pgfpathlineto{\pgfqpoint{3.475308in}{1.682111in}}%
\pgfpathlineto{\pgfqpoint{3.478038in}{1.668863in}}%
\pgfpathlineto{\pgfqpoint{3.481913in}{1.660577in}}%
\pgfpathlineto{\pgfqpoint{3.484271in}{1.655615in}}%
\pgfpathlineto{\pgfqpoint{3.486144in}{1.642367in}}%
\pgfpathlineto{\pgfqpoint{3.492222in}{1.629119in}}%
\pgfpathlineto{\pgfqpoint{3.493745in}{1.615871in}}%
\pgfpathlineto{\pgfqpoint{3.500511in}{1.602623in}}%
\pgfpathlineto{\pgfqpoint{3.501705in}{1.589375in}}%
\pgfpathlineto{\pgfqpoint{3.507196in}{1.578828in}}%
\pgfpathlineto{\pgfqpoint{3.508677in}{1.576127in}}%
\pgfpathlineto{\pgfqpoint{3.509842in}{1.562879in}}%
\pgfpathlineto{\pgfqpoint{3.516900in}{1.549631in}}%
\pgfpathlineto{\pgfqpoint{3.518330in}{1.536383in}}%
\pgfpathlineto{\pgfqpoint{3.524203in}{1.523135in}}%
\pgfpathlineto{\pgfqpoint{3.524974in}{1.509887in}}%
\pgfpathlineto{\pgfqpoint{3.530693in}{1.496639in}}%
\pgfpathlineto{\pgfqpoint{3.532480in}{1.486697in}}%
\pgfpathlineto{\pgfqpoint{3.533115in}{1.483391in}}%
\pgfpathlineto{\pgfqpoint{3.539011in}{1.470143in}}%
\pgfpathlineto{\pgfqpoint{3.542336in}{1.456895in}}%
\pgfpathlineto{\pgfqpoint{3.546556in}{1.443647in}}%
\pgfpathlineto{\pgfqpoint{3.550019in}{1.430399in}}%
\pgfpathlineto{\pgfqpoint{3.553652in}{1.417151in}}%
\pgfpathlineto{\pgfqpoint{3.557763in}{1.405975in}}%
\pgfpathlineto{\pgfqpoint{3.558570in}{1.403903in}}%
\pgfpathlineto{\pgfqpoint{3.560930in}{1.390655in}}%
\pgfpathlineto{\pgfqpoint{3.566948in}{1.377407in}}%
\pgfpathlineto{\pgfqpoint{3.568262in}{1.364159in}}%
\pgfpathlineto{\pgfqpoint{3.574888in}{1.350911in}}%
\pgfpathlineto{\pgfqpoint{3.575403in}{1.337663in}}%
\pgfpathlineto{\pgfqpoint{3.582015in}{1.324415in}}%
\pgfpathlineto{\pgfqpoint{3.581464in}{1.311167in}}%
\pgfpathlineto{\pgfqpoint{3.583046in}{1.308657in}}%
\pgfpathlineto{\pgfqpoint{3.590161in}{1.297919in}}%
\pgfpathlineto{\pgfqpoint{3.590781in}{1.284671in}}%
\pgfpathlineto{\pgfqpoint{3.599027in}{1.271423in}}%
\pgfpathlineto{\pgfqpoint{3.598651in}{1.258175in}}%
\pgfpathlineto{\pgfqpoint{3.605445in}{1.244927in}}%
\pgfpathlineto{\pgfqpoint{3.604862in}{1.231679in}}%
\pgfpathlineto{\pgfqpoint{3.608329in}{1.224876in}}%
\pgfpathlineto{\pgfqpoint{3.611867in}{1.218431in}}%
\pgfpathlineto{\pgfqpoint{3.612828in}{1.205183in}}%
\pgfpathlineto{\pgfqpoint{3.617627in}{1.191935in}}%
\pgfpathlineto{\pgfqpoint{3.619971in}{1.178687in}}%
\pgfpathlineto{\pgfqpoint{3.623581in}{1.165439in}}%
\pgfpathlineto{\pgfqpoint{3.627131in}{1.152191in}}%
\pgfpathlineto{\pgfqpoint{3.628213in}{1.138943in}}%
\pgfpathlineto{\pgfqpoint{3.633079in}{1.125695in}}%
\pgfpathlineto{\pgfqpoint{3.631761in}{1.112447in}}%
\pgfpathlineto{\pgfqpoint{3.633612in}{1.109537in}}%
\pgfpathlineto{\pgfqpoint{3.640365in}{1.099199in}}%
\pgfpathlineto{\pgfqpoint{3.637418in}{1.085951in}}%
\pgfpathlineto{\pgfqpoint{3.644230in}{1.072703in}}%
\pgfpathlineto{\pgfqpoint{3.636116in}{1.059455in}}%
\pgfpathlineto{\pgfqpoint{3.644085in}{1.046207in}}%
\pgfpathlineto{\pgfqpoint{3.634545in}{1.032959in}}%
\pgfpathlineto{\pgfqpoint{3.643850in}{1.019711in}}%
\pgfpathlineto{\pgfqpoint{3.634004in}{1.006463in}}%
\pgfpathlineto{\pgfqpoint{3.639766in}{0.993215in}}%
\pgfpathlineto{\pgfqpoint{3.633612in}{0.985271in}}%
\pgfpathlineto{\pgfqpoint{3.629166in}{0.979967in}}%
\pgfpathlineto{\pgfqpoint{3.631229in}{0.966719in}}%
\pgfpathlineto{\pgfqpoint{3.617182in}{0.953471in}}%
\pgfpathlineto{\pgfqpoint{3.608329in}{0.941497in}}%
\pgfpathlineto{\pgfqpoint{3.607195in}{0.940223in}}%
\pgfpathlineto{\pgfqpoint{3.595429in}{0.926975in}}%
\pgfpathlineto{\pgfqpoint{3.583046in}{0.918985in}}%
\pgfpathlineto{\pgfqpoint{3.572944in}{0.913727in}}%
\pgfpathlineto{\pgfqpoint{3.562820in}{0.900479in}}%
\pgfpathlineto{\pgfqpoint{3.557763in}{0.898812in}}%
\pgfpathlineto{\pgfqpoint{3.532480in}{0.891935in}}%
\pgfpathlineto{\pgfqpoint{3.509970in}{0.887231in}}%
\pgfpathlineto{\pgfqpoint{3.507196in}{0.885423in}}%
\pgfpathlineto{\pgfqpoint{3.486908in}{0.873983in}}%
\pgfpathlineto{\pgfqpoint{3.481913in}{0.873265in}}%
\pgfpathlineto{\pgfqpoint{3.456630in}{0.870258in}}%
\pgfpathlineto{\pgfqpoint{3.431347in}{0.867906in}}%
\pgfpathlineto{\pgfqpoint{3.406064in}{0.866128in}}%
\pgfpathlineto{\pgfqpoint{3.380781in}{0.864866in}}%
\pgfpathlineto{\pgfqpoint{3.355497in}{0.864081in}}%
\pgfpathlineto{\pgfqpoint{3.330214in}{0.863753in}}%
\pgfpathlineto{\pgfqpoint{3.304931in}{0.863875in}}%
\pgfpathlineto{\pgfqpoint{3.279648in}{0.864459in}}%
\pgfpathlineto{\pgfqpoint{3.254365in}{0.865530in}}%
\pgfpathlineto{\pgfqpoint{3.229082in}{0.867134in}}%
\pgfpathlineto{\pgfqpoint{3.203798in}{0.869341in}}%
\pgfpathlineto{\pgfqpoint{3.178515in}{0.872246in}}%
\pgfpathlineto{\pgfqpoint{3.166437in}{0.873983in}}%
\pgfpathclose%
\pgfpathmoveto{\pgfqpoint{3.026816in}{2.087055in}}%
\pgfpathlineto{\pgfqpoint{3.031070in}{2.092799in}}%
\pgfpathlineto{\pgfqpoint{3.052099in}{2.105070in}}%
\pgfpathlineto{\pgfqpoint{3.052756in}{2.106047in}}%
\pgfpathlineto{\pgfqpoint{3.054502in}{2.119295in}}%
\pgfpathlineto{\pgfqpoint{3.064553in}{2.132543in}}%
\pgfpathlineto{\pgfqpoint{3.052099in}{2.132543in}}%
\pgfpathlineto{\pgfqpoint{3.026816in}{2.132543in}}%
\pgfpathlineto{\pgfqpoint{3.017784in}{2.132543in}}%
\pgfpathlineto{\pgfqpoint{3.018572in}{2.119295in}}%
\pgfpathlineto{\pgfqpoint{3.018664in}{2.106047in}}%
\pgfpathlineto{\pgfqpoint{3.023033in}{2.092799in}}%
\pgfpathlineto{\pgfqpoint{3.026816in}{2.087055in}}%
\pgfpathclose%
\pgfpathmoveto{\pgfqpoint{3.633612in}{2.103895in}}%
\pgfpathlineto{\pgfqpoint{3.640182in}{2.106047in}}%
\pgfpathlineto{\pgfqpoint{3.637436in}{2.119295in}}%
\pgfpathlineto{\pgfqpoint{3.650509in}{2.132543in}}%
\pgfpathlineto{\pgfqpoint{3.633612in}{2.132543in}}%
\pgfpathlineto{\pgfqpoint{3.615726in}{2.132543in}}%
\pgfpathlineto{\pgfqpoint{3.630815in}{2.119295in}}%
\pgfpathlineto{\pgfqpoint{3.631816in}{2.106047in}}%
\pgfpathlineto{\pgfqpoint{3.633612in}{2.103895in}}%
\pgfpathclose%
\pgfusepath{fill}%
\end{pgfscope}%
\begin{pgfscope}%
\pgfpathrectangle{\pgfqpoint{2.470586in}{0.476543in}}{\pgfqpoint{1.820388in}{1.656000in}}%
\pgfusepath{clip}%
\pgfsetbuttcap%
\pgfsetroundjoin%
\definecolor{currentfill}{rgb}{0.870281,0.948943,0.970242}%
\pgfsetfillcolor{currentfill}%
\pgfsetlinewidth{0.000000pt}%
\definecolor{currentstroke}{rgb}{0.000000,0.000000,0.000000}%
\pgfsetstrokecolor{currentstroke}%
\pgfsetdash{}{0pt}%
\pgfpathmoveto{\pgfqpoint{3.178515in}{0.872246in}}%
\pgfpathlineto{\pgfqpoint{3.203798in}{0.869341in}}%
\pgfpathlineto{\pgfqpoint{3.229082in}{0.867134in}}%
\pgfpathlineto{\pgfqpoint{3.254365in}{0.865530in}}%
\pgfpathlineto{\pgfqpoint{3.279648in}{0.864459in}}%
\pgfpathlineto{\pgfqpoint{3.304931in}{0.863875in}}%
\pgfpathlineto{\pgfqpoint{3.330214in}{0.863753in}}%
\pgfpathlineto{\pgfqpoint{3.355497in}{0.864081in}}%
\pgfpathlineto{\pgfqpoint{3.380781in}{0.864866in}}%
\pgfpathlineto{\pgfqpoint{3.406064in}{0.866128in}}%
\pgfpathlineto{\pgfqpoint{3.431347in}{0.867906in}}%
\pgfpathlineto{\pgfqpoint{3.456630in}{0.870258in}}%
\pgfpathlineto{\pgfqpoint{3.481913in}{0.873265in}}%
\pgfpathlineto{\pgfqpoint{3.486908in}{0.873983in}}%
\pgfpathlineto{\pgfqpoint{3.507196in}{0.885423in}}%
\pgfpathlineto{\pgfqpoint{3.509970in}{0.887231in}}%
\pgfpathlineto{\pgfqpoint{3.532480in}{0.891935in}}%
\pgfpathlineto{\pgfqpoint{3.557763in}{0.898812in}}%
\pgfpathlineto{\pgfqpoint{3.562820in}{0.900479in}}%
\pgfpathlineto{\pgfqpoint{3.572944in}{0.913727in}}%
\pgfpathlineto{\pgfqpoint{3.583046in}{0.918985in}}%
\pgfpathlineto{\pgfqpoint{3.595429in}{0.926975in}}%
\pgfpathlineto{\pgfqpoint{3.607195in}{0.940223in}}%
\pgfpathlineto{\pgfqpoint{3.608329in}{0.941497in}}%
\pgfpathlineto{\pgfqpoint{3.617182in}{0.953471in}}%
\pgfpathlineto{\pgfqpoint{3.631229in}{0.966719in}}%
\pgfpathlineto{\pgfqpoint{3.629166in}{0.979967in}}%
\pgfpathlineto{\pgfqpoint{3.633612in}{0.985271in}}%
\pgfpathlineto{\pgfqpoint{3.639766in}{0.993215in}}%
\pgfpathlineto{\pgfqpoint{3.634004in}{1.006463in}}%
\pgfpathlineto{\pgfqpoint{3.643850in}{1.019711in}}%
\pgfpathlineto{\pgfqpoint{3.634545in}{1.032959in}}%
\pgfpathlineto{\pgfqpoint{3.644085in}{1.046207in}}%
\pgfpathlineto{\pgfqpoint{3.636116in}{1.059455in}}%
\pgfpathlineto{\pgfqpoint{3.644230in}{1.072703in}}%
\pgfpathlineto{\pgfqpoint{3.637418in}{1.085951in}}%
\pgfpathlineto{\pgfqpoint{3.640365in}{1.099199in}}%
\pgfpathlineto{\pgfqpoint{3.633612in}{1.109537in}}%
\pgfpathlineto{\pgfqpoint{3.631761in}{1.112447in}}%
\pgfpathlineto{\pgfqpoint{3.633079in}{1.125695in}}%
\pgfpathlineto{\pgfqpoint{3.628213in}{1.138943in}}%
\pgfpathlineto{\pgfqpoint{3.627131in}{1.152191in}}%
\pgfpathlineto{\pgfqpoint{3.623581in}{1.165439in}}%
\pgfpathlineto{\pgfqpoint{3.619971in}{1.178687in}}%
\pgfpathlineto{\pgfqpoint{3.617627in}{1.191935in}}%
\pgfpathlineto{\pgfqpoint{3.612828in}{1.205183in}}%
\pgfpathlineto{\pgfqpoint{3.611867in}{1.218431in}}%
\pgfpathlineto{\pgfqpoint{3.608329in}{1.224876in}}%
\pgfpathlineto{\pgfqpoint{3.604862in}{1.231679in}}%
\pgfpathlineto{\pgfqpoint{3.605445in}{1.244927in}}%
\pgfpathlineto{\pgfqpoint{3.598651in}{1.258175in}}%
\pgfpathlineto{\pgfqpoint{3.599027in}{1.271423in}}%
\pgfpathlineto{\pgfqpoint{3.590781in}{1.284671in}}%
\pgfpathlineto{\pgfqpoint{3.590161in}{1.297919in}}%
\pgfpathlineto{\pgfqpoint{3.583046in}{1.308657in}}%
\pgfpathlineto{\pgfqpoint{3.581464in}{1.311167in}}%
\pgfpathlineto{\pgfqpoint{3.582015in}{1.324415in}}%
\pgfpathlineto{\pgfqpoint{3.575403in}{1.337663in}}%
\pgfpathlineto{\pgfqpoint{3.574888in}{1.350911in}}%
\pgfpathlineto{\pgfqpoint{3.568262in}{1.364159in}}%
\pgfpathlineto{\pgfqpoint{3.566948in}{1.377407in}}%
\pgfpathlineto{\pgfqpoint{3.560930in}{1.390655in}}%
\pgfpathlineto{\pgfqpoint{3.558570in}{1.403903in}}%
\pgfpathlineto{\pgfqpoint{3.557763in}{1.405975in}}%
\pgfpathlineto{\pgfqpoint{3.553652in}{1.417151in}}%
\pgfpathlineto{\pgfqpoint{3.550019in}{1.430399in}}%
\pgfpathlineto{\pgfqpoint{3.546556in}{1.443647in}}%
\pgfpathlineto{\pgfqpoint{3.542336in}{1.456895in}}%
\pgfpathlineto{\pgfqpoint{3.539011in}{1.470143in}}%
\pgfpathlineto{\pgfqpoint{3.533115in}{1.483391in}}%
\pgfpathlineto{\pgfqpoint{3.532480in}{1.486697in}}%
\pgfpathlineto{\pgfqpoint{3.530693in}{1.496639in}}%
\pgfpathlineto{\pgfqpoint{3.524974in}{1.509887in}}%
\pgfpathlineto{\pgfqpoint{3.524203in}{1.523135in}}%
\pgfpathlineto{\pgfqpoint{3.518330in}{1.536383in}}%
\pgfpathlineto{\pgfqpoint{3.516900in}{1.549631in}}%
\pgfpathlineto{\pgfqpoint{3.509842in}{1.562879in}}%
\pgfpathlineto{\pgfqpoint{3.508677in}{1.576127in}}%
\pgfpathlineto{\pgfqpoint{3.507196in}{1.578828in}}%
\pgfpathlineto{\pgfqpoint{3.501705in}{1.589375in}}%
\pgfpathlineto{\pgfqpoint{3.500511in}{1.602623in}}%
\pgfpathlineto{\pgfqpoint{3.493745in}{1.615871in}}%
\pgfpathlineto{\pgfqpoint{3.492222in}{1.629119in}}%
\pgfpathlineto{\pgfqpoint{3.486144in}{1.642367in}}%
\pgfpathlineto{\pgfqpoint{3.484271in}{1.655615in}}%
\pgfpathlineto{\pgfqpoint{3.481913in}{1.660577in}}%
\pgfpathlineto{\pgfqpoint{3.478038in}{1.668863in}}%
\pgfpathlineto{\pgfqpoint{3.475308in}{1.682111in}}%
\pgfpathlineto{\pgfqpoint{3.470122in}{1.695359in}}%
\pgfpathlineto{\pgfqpoint{3.467187in}{1.708607in}}%
\pgfpathlineto{\pgfqpoint{3.463056in}{1.721855in}}%
\pgfpathlineto{\pgfqpoint{3.458861in}{1.735103in}}%
\pgfpathlineto{\pgfqpoint{3.456630in}{1.741373in}}%
\pgfpathlineto{\pgfqpoint{3.454080in}{1.748351in}}%
\pgfpathlineto{\pgfqpoint{3.448628in}{1.761599in}}%
\pgfpathlineto{\pgfqpoint{3.444996in}{1.774847in}}%
\pgfpathlineto{\pgfqpoint{3.438935in}{1.788095in}}%
\pgfpathlineto{\pgfqpoint{3.436031in}{1.801343in}}%
\pgfpathlineto{\pgfqpoint{3.431347in}{1.810480in}}%
\pgfpathlineto{\pgfqpoint{3.429050in}{1.814591in}}%
\pgfpathlineto{\pgfqpoint{3.426350in}{1.827839in}}%
\pgfpathlineto{\pgfqpoint{3.418708in}{1.841087in}}%
\pgfpathlineto{\pgfqpoint{3.416264in}{1.854335in}}%
\pgfpathlineto{\pgfqpoint{3.407165in}{1.867583in}}%
\pgfpathlineto{\pgfqpoint{3.406064in}{1.873655in}}%
\pgfpathlineto{\pgfqpoint{3.404446in}{1.880831in}}%
\pgfpathlineto{\pgfqpoint{3.393176in}{1.894079in}}%
\pgfpathlineto{\pgfqpoint{3.390757in}{1.907327in}}%
\pgfpathlineto{\pgfqpoint{3.380781in}{1.918898in}}%
\pgfpathlineto{\pgfqpoint{3.378368in}{1.920575in}}%
\pgfpathlineto{\pgfqpoint{3.372336in}{1.933823in}}%
\pgfpathlineto{\pgfqpoint{3.355497in}{1.944026in}}%
\pgfpathlineto{\pgfqpoint{3.330214in}{1.947066in}}%
\pgfpathlineto{\pgfqpoint{3.304931in}{1.937496in}}%
\pgfpathlineto{\pgfqpoint{3.300701in}{1.933823in}}%
\pgfpathlineto{\pgfqpoint{3.295653in}{1.920575in}}%
\pgfpathlineto{\pgfqpoint{3.284159in}{1.907327in}}%
\pgfpathlineto{\pgfqpoint{3.279648in}{1.895913in}}%
\pgfpathlineto{\pgfqpoint{3.279099in}{1.894079in}}%
\pgfpathlineto{\pgfqpoint{3.271608in}{1.880831in}}%
\pgfpathlineto{\pgfqpoint{3.265905in}{1.867583in}}%
\pgfpathlineto{\pgfqpoint{3.260492in}{1.854335in}}%
\pgfpathlineto{\pgfqpoint{3.254365in}{1.843035in}}%
\pgfpathlineto{\pgfqpoint{3.253420in}{1.841087in}}%
\pgfpathlineto{\pgfqpoint{3.249696in}{1.827839in}}%
\pgfpathlineto{\pgfqpoint{3.242211in}{1.814591in}}%
\pgfpathlineto{\pgfqpoint{3.239738in}{1.801343in}}%
\pgfpathlineto{\pgfqpoint{3.232465in}{1.788095in}}%
\pgfpathlineto{\pgfqpoint{3.230817in}{1.774847in}}%
\pgfpathlineto{\pgfqpoint{3.229082in}{1.772138in}}%
\pgfpathlineto{\pgfqpoint{3.222712in}{1.761599in}}%
\pgfpathlineto{\pgfqpoint{3.220912in}{1.748351in}}%
\pgfpathlineto{\pgfqpoint{3.213973in}{1.735103in}}%
\pgfpathlineto{\pgfqpoint{3.212780in}{1.721855in}}%
\pgfpathlineto{\pgfqpoint{3.206600in}{1.708607in}}%
\pgfpathlineto{\pgfqpoint{3.204503in}{1.695359in}}%
\pgfpathlineto{\pgfqpoint{3.203798in}{1.693925in}}%
\pgfpathlineto{\pgfqpoint{3.198095in}{1.682111in}}%
\pgfpathlineto{\pgfqpoint{3.194971in}{1.668863in}}%
\pgfpathlineto{\pgfqpoint{3.190052in}{1.655615in}}%
\pgfpathlineto{\pgfqpoint{3.185857in}{1.642367in}}%
\pgfpathlineto{\pgfqpoint{3.181944in}{1.629119in}}%
\pgfpathlineto{\pgfqpoint{3.178515in}{1.619152in}}%
\pgfpathlineto{\pgfqpoint{3.177307in}{1.615871in}}%
\pgfpathlineto{\pgfqpoint{3.175027in}{1.602623in}}%
\pgfpathlineto{\pgfqpoint{3.169577in}{1.589375in}}%
\pgfpathlineto{\pgfqpoint{3.167684in}{1.576127in}}%
\pgfpathlineto{\pgfqpoint{3.160697in}{1.562879in}}%
\pgfpathlineto{\pgfqpoint{3.159549in}{1.549631in}}%
\pgfpathlineto{\pgfqpoint{3.153232in}{1.536547in}}%
\pgfpathlineto{\pgfqpoint{3.153148in}{1.536383in}}%
\pgfpathlineto{\pgfqpoint{3.152488in}{1.523135in}}%
\pgfpathlineto{\pgfqpoint{3.144947in}{1.509887in}}%
\pgfpathlineto{\pgfqpoint{3.143833in}{1.496639in}}%
\pgfpathlineto{\pgfqpoint{3.136417in}{1.483391in}}%
\pgfpathlineto{\pgfqpoint{3.135367in}{1.470143in}}%
\pgfpathlineto{\pgfqpoint{3.129109in}{1.456895in}}%
\pgfpathlineto{\pgfqpoint{3.127949in}{1.447878in}}%
\pgfpathlineto{\pgfqpoint{3.127342in}{1.443647in}}%
\pgfpathlineto{\pgfqpoint{3.121689in}{1.430399in}}%
\pgfpathlineto{\pgfqpoint{3.119609in}{1.417151in}}%
\pgfpathlineto{\pgfqpoint{3.115292in}{1.403903in}}%
\pgfpathlineto{\pgfqpoint{3.111321in}{1.390655in}}%
\pgfpathlineto{\pgfqpoint{3.107203in}{1.377407in}}%
\pgfpathlineto{\pgfqpoint{3.102666in}{1.366770in}}%
\pgfpathlineto{\pgfqpoint{3.101441in}{1.364159in}}%
\pgfpathlineto{\pgfqpoint{3.099435in}{1.350911in}}%
\pgfpathlineto{\pgfqpoint{3.093603in}{1.337663in}}%
\pgfpathlineto{\pgfqpoint{3.092493in}{1.324415in}}%
\pgfpathlineto{\pgfqpoint{3.084804in}{1.311167in}}%
\pgfpathlineto{\pgfqpoint{3.084678in}{1.297919in}}%
\pgfpathlineto{\pgfqpoint{3.077382in}{1.285762in}}%
\pgfpathlineto{\pgfqpoint{3.076681in}{1.284671in}}%
\pgfpathlineto{\pgfqpoint{3.077382in}{1.277820in}}%
\pgfpathlineto{\pgfqpoint{3.078023in}{1.271423in}}%
\pgfpathlineto{\pgfqpoint{3.077382in}{1.270392in}}%
\pgfpathlineto{\pgfqpoint{3.069360in}{1.258175in}}%
\pgfpathlineto{\pgfqpoint{3.070706in}{1.244927in}}%
\pgfpathlineto{\pgfqpoint{3.062102in}{1.231679in}}%
\pgfpathlineto{\pgfqpoint{3.063726in}{1.218431in}}%
\pgfpathlineto{\pgfqpoint{3.054844in}{1.205183in}}%
\pgfpathlineto{\pgfqpoint{3.054511in}{1.191935in}}%
\pgfpathlineto{\pgfqpoint{3.052099in}{1.188043in}}%
\pgfpathlineto{\pgfqpoint{3.046157in}{1.178687in}}%
\pgfpathlineto{\pgfqpoint{3.046130in}{1.165439in}}%
\pgfpathlineto{\pgfqpoint{3.041775in}{1.152191in}}%
\pgfpathlineto{\pgfqpoint{3.041648in}{1.138943in}}%
\pgfpathlineto{\pgfqpoint{3.037547in}{1.125695in}}%
\pgfpathlineto{\pgfqpoint{3.034430in}{1.112447in}}%
\pgfpathlineto{\pgfqpoint{3.033075in}{1.099199in}}%
\pgfpathlineto{\pgfqpoint{3.028932in}{1.085951in}}%
\pgfpathlineto{\pgfqpoint{3.030906in}{1.072703in}}%
\pgfpathlineto{\pgfqpoint{3.026816in}{1.063842in}}%
\pgfpathlineto{\pgfqpoint{3.024625in}{1.059455in}}%
\pgfpathlineto{\pgfqpoint{3.026816in}{1.053567in}}%
\pgfpathlineto{\pgfqpoint{3.029646in}{1.046207in}}%
\pgfpathlineto{\pgfqpoint{3.026816in}{1.041130in}}%
\pgfpathlineto{\pgfqpoint{3.021972in}{1.032959in}}%
\pgfpathlineto{\pgfqpoint{3.026816in}{1.025655in}}%
\pgfpathlineto{\pgfqpoint{3.031007in}{1.019711in}}%
\pgfpathlineto{\pgfqpoint{3.026816in}{1.013955in}}%
\pgfpathlineto{\pgfqpoint{3.021031in}{1.006463in}}%
\pgfpathlineto{\pgfqpoint{3.026816in}{0.999810in}}%
\pgfpathlineto{\pgfqpoint{3.033024in}{0.993215in}}%
\pgfpathlineto{\pgfqpoint{3.026816in}{0.984100in}}%
\pgfpathlineto{\pgfqpoint{3.023884in}{0.979967in}}%
\pgfpathlineto{\pgfqpoint{3.026816in}{0.977964in}}%
\pgfpathlineto{\pgfqpoint{3.045179in}{0.966719in}}%
\pgfpathlineto{\pgfqpoint{3.040234in}{0.953471in}}%
\pgfpathlineto{\pgfqpoint{3.052099in}{0.947410in}}%
\pgfpathlineto{\pgfqpoint{3.068604in}{0.940223in}}%
\pgfpathlineto{\pgfqpoint{3.059863in}{0.926975in}}%
\pgfpathlineto{\pgfqpoint{3.077382in}{0.919772in}}%
\pgfpathlineto{\pgfqpoint{3.095137in}{0.913727in}}%
\pgfpathlineto{\pgfqpoint{3.092882in}{0.900479in}}%
\pgfpathlineto{\pgfqpoint{3.102666in}{0.897860in}}%
\pgfpathlineto{\pgfqpoint{3.127949in}{0.892193in}}%
\pgfpathlineto{\pgfqpoint{3.153232in}{0.887645in}}%
\pgfpathlineto{\pgfqpoint{3.155985in}{0.887231in}}%
\pgfpathlineto{\pgfqpoint{3.166437in}{0.873983in}}%
\pgfpathlineto{\pgfqpoint{3.178515in}{0.872246in}}%
\pgfpathclose%
\pgfpathmoveto{\pgfqpoint{3.175588in}{0.900479in}}%
\pgfpathlineto{\pgfqpoint{3.171168in}{0.913727in}}%
\pgfpathlineto{\pgfqpoint{3.153232in}{0.917448in}}%
\pgfpathlineto{\pgfqpoint{3.127949in}{0.923901in}}%
\pgfpathlineto{\pgfqpoint{3.117903in}{0.926975in}}%
\pgfpathlineto{\pgfqpoint{3.123091in}{0.940223in}}%
\pgfpathlineto{\pgfqpoint{3.102666in}{0.946782in}}%
\pgfpathlineto{\pgfqpoint{3.085322in}{0.953471in}}%
\pgfpathlineto{\pgfqpoint{3.087261in}{0.966719in}}%
\pgfpathlineto{\pgfqpoint{3.077382in}{0.971376in}}%
\pgfpathlineto{\pgfqpoint{3.061345in}{0.979967in}}%
\pgfpathlineto{\pgfqpoint{3.068272in}{0.993215in}}%
\pgfpathlineto{\pgfqpoint{3.053687in}{1.006463in}}%
\pgfpathlineto{\pgfqpoint{3.061501in}{1.019711in}}%
\pgfpathlineto{\pgfqpoint{3.052099in}{1.031116in}}%
\pgfpathlineto{\pgfqpoint{3.050688in}{1.032959in}}%
\pgfpathlineto{\pgfqpoint{3.052099in}{1.036042in}}%
\pgfpathlineto{\pgfqpoint{3.056497in}{1.046207in}}%
\pgfpathlineto{\pgfqpoint{3.052099in}{1.055904in}}%
\pgfpathlineto{\pgfqpoint{3.050562in}{1.059455in}}%
\pgfpathlineto{\pgfqpoint{3.052099in}{1.063958in}}%
\pgfpathlineto{\pgfqpoint{3.054870in}{1.072703in}}%
\pgfpathlineto{\pgfqpoint{3.052252in}{1.085951in}}%
\pgfpathlineto{\pgfqpoint{3.054575in}{1.099199in}}%
\pgfpathlineto{\pgfqpoint{3.055673in}{1.112447in}}%
\pgfpathlineto{\pgfqpoint{3.057251in}{1.125695in}}%
\pgfpathlineto{\pgfqpoint{3.061199in}{1.138943in}}%
\pgfpathlineto{\pgfqpoint{3.059908in}{1.152191in}}%
\pgfpathlineto{\pgfqpoint{3.064203in}{1.165439in}}%
\pgfpathlineto{\pgfqpoint{3.063096in}{1.178687in}}%
\pgfpathlineto{\pgfqpoint{3.071373in}{1.191935in}}%
\pgfpathlineto{\pgfqpoint{3.070703in}{1.205183in}}%
\pgfpathlineto{\pgfqpoint{3.077382in}{1.215149in}}%
\pgfpathlineto{\pgfqpoint{3.079521in}{1.218431in}}%
\pgfpathlineto{\pgfqpoint{3.077382in}{1.230209in}}%
\pgfpathlineto{\pgfqpoint{3.077122in}{1.231679in}}%
\pgfpathlineto{\pgfqpoint{3.077382in}{1.232078in}}%
\pgfpathlineto{\pgfqpoint{3.085473in}{1.244927in}}%
\pgfpathlineto{\pgfqpoint{3.083574in}{1.258175in}}%
\pgfpathlineto{\pgfqpoint{3.091731in}{1.271423in}}%
\pgfpathlineto{\pgfqpoint{3.089957in}{1.284671in}}%
\pgfpathlineto{\pgfqpoint{3.097744in}{1.297919in}}%
\pgfpathlineto{\pgfqpoint{3.097660in}{1.311167in}}%
\pgfpathlineto{\pgfqpoint{3.102666in}{1.319937in}}%
\pgfpathlineto{\pgfqpoint{3.105068in}{1.324415in}}%
\pgfpathlineto{\pgfqpoint{3.106045in}{1.337663in}}%
\pgfpathlineto{\pgfqpoint{3.111259in}{1.350911in}}%
\pgfpathlineto{\pgfqpoint{3.113210in}{1.364159in}}%
\pgfpathlineto{\pgfqpoint{3.118547in}{1.377407in}}%
\pgfpathlineto{\pgfqpoint{3.122825in}{1.390655in}}%
\pgfpathlineto{\pgfqpoint{3.126384in}{1.403903in}}%
\pgfpathlineto{\pgfqpoint{3.127949in}{1.408560in}}%
\pgfpathlineto{\pgfqpoint{3.130722in}{1.417151in}}%
\pgfpathlineto{\pgfqpoint{3.132314in}{1.430399in}}%
\pgfpathlineto{\pgfqpoint{3.137855in}{1.443647in}}%
\pgfpathlineto{\pgfqpoint{3.139207in}{1.456895in}}%
\pgfpathlineto{\pgfqpoint{3.145693in}{1.470143in}}%
\pgfpathlineto{\pgfqpoint{3.146434in}{1.483391in}}%
\pgfpathlineto{\pgfqpoint{3.153232in}{1.495112in}}%
\pgfpathlineto{\pgfqpoint{3.154097in}{1.496639in}}%
\pgfpathlineto{\pgfqpoint{3.154966in}{1.509887in}}%
\pgfpathlineto{\pgfqpoint{3.162205in}{1.523135in}}%
\pgfpathlineto{\pgfqpoint{3.162652in}{1.536383in}}%
\pgfpathlineto{\pgfqpoint{3.169104in}{1.549631in}}%
\pgfpathlineto{\pgfqpoint{3.170226in}{1.562879in}}%
\pgfpathlineto{\pgfqpoint{3.177251in}{1.576127in}}%
\pgfpathlineto{\pgfqpoint{3.178515in}{1.584759in}}%
\pgfpathlineto{\pgfqpoint{3.179169in}{1.589375in}}%
\pgfpathlineto{\pgfqpoint{3.184295in}{1.602623in}}%
\pgfpathlineto{\pgfqpoint{3.186627in}{1.615871in}}%
\pgfpathlineto{\pgfqpoint{3.191077in}{1.629119in}}%
\pgfpathlineto{\pgfqpoint{3.195270in}{1.642367in}}%
\pgfpathlineto{\pgfqpoint{3.199287in}{1.655615in}}%
\pgfpathlineto{\pgfqpoint{3.203798in}{1.667078in}}%
\pgfpathlineto{\pgfqpoint{3.204519in}{1.668863in}}%
\pgfpathlineto{\pgfqpoint{3.207412in}{1.682111in}}%
\pgfpathlineto{\pgfqpoint{3.214051in}{1.695359in}}%
\pgfpathlineto{\pgfqpoint{3.215904in}{1.708607in}}%
\pgfpathlineto{\pgfqpoint{3.222374in}{1.721855in}}%
\pgfpathlineto{\pgfqpoint{3.223385in}{1.735103in}}%
\pgfpathlineto{\pgfqpoint{3.229082in}{1.745537in}}%
\pgfpathlineto{\pgfqpoint{3.230742in}{1.748351in}}%
\pgfpathlineto{\pgfqpoint{3.232566in}{1.761599in}}%
\pgfpathlineto{\pgfqpoint{3.241106in}{1.774847in}}%
\pgfpathlineto{\pgfqpoint{3.242799in}{1.788095in}}%
\pgfpathlineto{\pgfqpoint{3.250194in}{1.801343in}}%
\pgfpathlineto{\pgfqpoint{3.252886in}{1.814591in}}%
\pgfpathlineto{\pgfqpoint{3.254365in}{1.817195in}}%
\pgfpathlineto{\pgfqpoint{3.261301in}{1.827839in}}%
\pgfpathlineto{\pgfqpoint{3.265933in}{1.841087in}}%
\pgfpathlineto{\pgfqpoint{3.272842in}{1.854335in}}%
\pgfpathlineto{\pgfqpoint{3.278841in}{1.867583in}}%
\pgfpathlineto{\pgfqpoint{3.279648in}{1.869513in}}%
\pgfpathlineto{\pgfqpoint{3.285960in}{1.880831in}}%
\pgfpathlineto{\pgfqpoint{3.297288in}{1.894079in}}%
\pgfpathlineto{\pgfqpoint{3.301853in}{1.907327in}}%
\pgfpathlineto{\pgfqpoint{3.304931in}{1.910705in}}%
\pgfpathlineto{\pgfqpoint{3.325630in}{1.920575in}}%
\pgfpathlineto{\pgfqpoint{3.330214in}{1.926750in}}%
\pgfpathlineto{\pgfqpoint{3.343876in}{1.920575in}}%
\pgfpathlineto{\pgfqpoint{3.355497in}{1.918923in}}%
\pgfpathlineto{\pgfqpoint{3.372252in}{1.907327in}}%
\pgfpathlineto{\pgfqpoint{3.376176in}{1.894079in}}%
\pgfpathlineto{\pgfqpoint{3.380781in}{1.890930in}}%
\pgfpathlineto{\pgfqpoint{3.389626in}{1.880831in}}%
\pgfpathlineto{\pgfqpoint{3.392741in}{1.867583in}}%
\pgfpathlineto{\pgfqpoint{3.404205in}{1.854335in}}%
\pgfpathlineto{\pgfqpoint{3.406064in}{1.846689in}}%
\pgfpathlineto{\pgfqpoint{3.407136in}{1.841087in}}%
\pgfpathlineto{\pgfqpoint{3.414874in}{1.827839in}}%
\pgfpathlineto{\pgfqpoint{3.417829in}{1.814591in}}%
\pgfpathlineto{\pgfqpoint{3.425238in}{1.801343in}}%
\pgfpathlineto{\pgfqpoint{3.428641in}{1.788095in}}%
\pgfpathlineto{\pgfqpoint{3.431347in}{1.782644in}}%
\pgfpathlineto{\pgfqpoint{3.434918in}{1.774847in}}%
\pgfpathlineto{\pgfqpoint{3.438836in}{1.761599in}}%
\pgfpathlineto{\pgfqpoint{3.444200in}{1.748351in}}%
\pgfpathlineto{\pgfqpoint{3.449274in}{1.735103in}}%
\pgfpathlineto{\pgfqpoint{3.453349in}{1.721855in}}%
\pgfpathlineto{\pgfqpoint{3.456630in}{1.712138in}}%
\pgfpathlineto{\pgfqpoint{3.457776in}{1.708607in}}%
\pgfpathlineto{\pgfqpoint{3.460547in}{1.695359in}}%
\pgfpathlineto{\pgfqpoint{3.466078in}{1.682111in}}%
\pgfpathlineto{\pgfqpoint{3.468610in}{1.668863in}}%
\pgfpathlineto{\pgfqpoint{3.475008in}{1.655615in}}%
\pgfpathlineto{\pgfqpoint{3.476623in}{1.642367in}}%
\pgfpathlineto{\pgfqpoint{3.481913in}{1.631050in}}%
\pgfpathlineto{\pgfqpoint{3.482819in}{1.629119in}}%
\pgfpathlineto{\pgfqpoint{3.484216in}{1.615871in}}%
\pgfpathlineto{\pgfqpoint{3.491275in}{1.602623in}}%
\pgfpathlineto{\pgfqpoint{3.492288in}{1.589375in}}%
\pgfpathlineto{\pgfqpoint{3.499304in}{1.576127in}}%
\pgfpathlineto{\pgfqpoint{3.500283in}{1.562879in}}%
\pgfpathlineto{\pgfqpoint{3.507078in}{1.549631in}}%
\pgfpathlineto{\pgfqpoint{3.507196in}{1.548411in}}%
\pgfpathlineto{\pgfqpoint{3.508428in}{1.536383in}}%
\pgfpathlineto{\pgfqpoint{3.514355in}{1.523135in}}%
\pgfpathlineto{\pgfqpoint{3.515061in}{1.509887in}}%
\pgfpathlineto{\pgfqpoint{3.520730in}{1.496639in}}%
\pgfpathlineto{\pgfqpoint{3.523200in}{1.483391in}}%
\pgfpathlineto{\pgfqpoint{3.528558in}{1.470143in}}%
\pgfpathlineto{\pgfqpoint{3.531775in}{1.456895in}}%
\pgfpathlineto{\pgfqpoint{3.532480in}{1.454418in}}%
\pgfpathlineto{\pgfqpoint{3.535781in}{1.443647in}}%
\pgfpathlineto{\pgfqpoint{3.539385in}{1.430399in}}%
\pgfpathlineto{\pgfqpoint{3.542729in}{1.417151in}}%
\pgfpathlineto{\pgfqpoint{3.547754in}{1.403903in}}%
\pgfpathlineto{\pgfqpoint{3.549509in}{1.390655in}}%
\pgfpathlineto{\pgfqpoint{3.555328in}{1.377407in}}%
\pgfpathlineto{\pgfqpoint{3.556188in}{1.364159in}}%
\pgfpathlineto{\pgfqpoint{3.557763in}{1.360976in}}%
\pgfpathlineto{\pgfqpoint{3.562921in}{1.350911in}}%
\pgfpathlineto{\pgfqpoint{3.562965in}{1.337663in}}%
\pgfpathlineto{\pgfqpoint{3.569692in}{1.324415in}}%
\pgfpathlineto{\pgfqpoint{3.568655in}{1.311167in}}%
\pgfpathlineto{\pgfqpoint{3.577014in}{1.297919in}}%
\pgfpathlineto{\pgfqpoint{3.577245in}{1.284671in}}%
\pgfpathlineto{\pgfqpoint{3.583046in}{1.274752in}}%
\pgfpathlineto{\pgfqpoint{3.585104in}{1.271423in}}%
\pgfpathlineto{\pgfqpoint{3.584273in}{1.258175in}}%
\pgfpathlineto{\pgfqpoint{3.590815in}{1.244927in}}%
\pgfpathlineto{\pgfqpoint{3.589876in}{1.231679in}}%
\pgfpathlineto{\pgfqpoint{3.596203in}{1.218431in}}%
\pgfpathlineto{\pgfqpoint{3.596853in}{1.205183in}}%
\pgfpathlineto{\pgfqpoint{3.600699in}{1.191935in}}%
\pgfpathlineto{\pgfqpoint{3.602801in}{1.178687in}}%
\pgfpathlineto{\pgfqpoint{3.605475in}{1.165439in}}%
\pgfpathlineto{\pgfqpoint{3.608329in}{1.154170in}}%
\pgfpathlineto{\pgfqpoint{3.608843in}{1.152191in}}%
\pgfpathlineto{\pgfqpoint{3.608628in}{1.138943in}}%
\pgfpathlineto{\pgfqpoint{3.613375in}{1.125695in}}%
\pgfpathlineto{\pgfqpoint{3.610522in}{1.112447in}}%
\pgfpathlineto{\pgfqpoint{3.618749in}{1.099199in}}%
\pgfpathlineto{\pgfqpoint{3.614105in}{1.085951in}}%
\pgfpathlineto{\pgfqpoint{3.620208in}{1.072703in}}%
\pgfpathlineto{\pgfqpoint{3.610018in}{1.059455in}}%
\pgfpathlineto{\pgfqpoint{3.617249in}{1.046207in}}%
\pgfpathlineto{\pgfqpoint{3.608329in}{1.035847in}}%
\pgfpathlineto{\pgfqpoint{3.605602in}{1.032959in}}%
\pgfpathlineto{\pgfqpoint{3.608329in}{1.028179in}}%
\pgfpathlineto{\pgfqpoint{3.613348in}{1.019711in}}%
\pgfpathlineto{\pgfqpoint{3.608329in}{1.013978in}}%
\pgfpathlineto{\pgfqpoint{3.601114in}{1.006463in}}%
\pgfpathlineto{\pgfqpoint{3.604555in}{0.993215in}}%
\pgfpathlineto{\pgfqpoint{3.591329in}{0.979967in}}%
\pgfpathlineto{\pgfqpoint{3.589307in}{0.966719in}}%
\pgfpathlineto{\pgfqpoint{3.583046in}{0.961669in}}%
\pgfpathlineto{\pgfqpoint{3.571753in}{0.953471in}}%
\pgfpathlineto{\pgfqpoint{3.557763in}{0.943174in}}%
\pgfpathlineto{\pgfqpoint{3.552693in}{0.940223in}}%
\pgfpathlineto{\pgfqpoint{3.537340in}{0.926975in}}%
\pgfpathlineto{\pgfqpoint{3.532480in}{0.924911in}}%
\pgfpathlineto{\pgfqpoint{3.507196in}{0.916347in}}%
\pgfpathlineto{\pgfqpoint{3.497086in}{0.913727in}}%
\pgfpathlineto{\pgfqpoint{3.481913in}{0.901971in}}%
\pgfpathlineto{\pgfqpoint{3.479622in}{0.900479in}}%
\pgfpathlineto{\pgfqpoint{3.456630in}{0.896838in}}%
\pgfpathlineto{\pgfqpoint{3.431347in}{0.893821in}}%
\pgfpathlineto{\pgfqpoint{3.406064in}{0.891643in}}%
\pgfpathlineto{\pgfqpoint{3.380781in}{0.890185in}}%
\pgfpathlineto{\pgfqpoint{3.355497in}{0.889364in}}%
\pgfpathlineto{\pgfqpoint{3.330214in}{0.889127in}}%
\pgfpathlineto{\pgfqpoint{3.304931in}{0.889446in}}%
\pgfpathlineto{\pgfqpoint{3.279648in}{0.890316in}}%
\pgfpathlineto{\pgfqpoint{3.254365in}{0.891752in}}%
\pgfpathlineto{\pgfqpoint{3.229082in}{0.893794in}}%
\pgfpathlineto{\pgfqpoint{3.203798in}{0.896508in}}%
\pgfpathlineto{\pgfqpoint{3.178515in}{0.899991in}}%
\pgfpathlineto{\pgfqpoint{3.175588in}{0.900479in}}%
\pgfpathclose%
\pgfusepath{fill}%
\end{pgfscope}%
\begin{pgfscope}%
\pgfpathrectangle{\pgfqpoint{2.470586in}{0.476543in}}{\pgfqpoint{1.820388in}{1.656000in}}%
\pgfusepath{clip}%
\pgfsetbuttcap%
\pgfsetroundjoin%
\definecolor{currentfill}{rgb}{0.926105,0.971396,0.884890}%
\pgfsetfillcolor{currentfill}%
\pgfsetlinewidth{0.000000pt}%
\definecolor{currentstroke}{rgb}{0.000000,0.000000,0.000000}%
\pgfsetstrokecolor{currentstroke}%
\pgfsetdash{}{0pt}%
\pgfpathmoveto{\pgfqpoint{3.178515in}{0.899991in}}%
\pgfpathlineto{\pgfqpoint{3.203798in}{0.896508in}}%
\pgfpathlineto{\pgfqpoint{3.229082in}{0.893794in}}%
\pgfpathlineto{\pgfqpoint{3.254365in}{0.891752in}}%
\pgfpathlineto{\pgfqpoint{3.279648in}{0.890316in}}%
\pgfpathlineto{\pgfqpoint{3.304931in}{0.889446in}}%
\pgfpathlineto{\pgfqpoint{3.330214in}{0.889127in}}%
\pgfpathlineto{\pgfqpoint{3.355497in}{0.889364in}}%
\pgfpathlineto{\pgfqpoint{3.380781in}{0.890185in}}%
\pgfpathlineto{\pgfqpoint{3.406064in}{0.891643in}}%
\pgfpathlineto{\pgfqpoint{3.431347in}{0.893821in}}%
\pgfpathlineto{\pgfqpoint{3.456630in}{0.896838in}}%
\pgfpathlineto{\pgfqpoint{3.479622in}{0.900479in}}%
\pgfpathlineto{\pgfqpoint{3.481913in}{0.901971in}}%
\pgfpathlineto{\pgfqpoint{3.497086in}{0.913727in}}%
\pgfpathlineto{\pgfqpoint{3.507196in}{0.916347in}}%
\pgfpathlineto{\pgfqpoint{3.532480in}{0.924911in}}%
\pgfpathlineto{\pgfqpoint{3.537340in}{0.926975in}}%
\pgfpathlineto{\pgfqpoint{3.552693in}{0.940223in}}%
\pgfpathlineto{\pgfqpoint{3.557763in}{0.943174in}}%
\pgfpathlineto{\pgfqpoint{3.571753in}{0.953471in}}%
\pgfpathlineto{\pgfqpoint{3.583046in}{0.961669in}}%
\pgfpathlineto{\pgfqpoint{3.589307in}{0.966719in}}%
\pgfpathlineto{\pgfqpoint{3.591329in}{0.979967in}}%
\pgfpathlineto{\pgfqpoint{3.604555in}{0.993215in}}%
\pgfpathlineto{\pgfqpoint{3.601114in}{1.006463in}}%
\pgfpathlineto{\pgfqpoint{3.608329in}{1.013978in}}%
\pgfpathlineto{\pgfqpoint{3.613348in}{1.019711in}}%
\pgfpathlineto{\pgfqpoint{3.608329in}{1.028179in}}%
\pgfpathlineto{\pgfqpoint{3.605602in}{1.032959in}}%
\pgfpathlineto{\pgfqpoint{3.608329in}{1.035847in}}%
\pgfpathlineto{\pgfqpoint{3.617249in}{1.046207in}}%
\pgfpathlineto{\pgfqpoint{3.610018in}{1.059455in}}%
\pgfpathlineto{\pgfqpoint{3.620208in}{1.072703in}}%
\pgfpathlineto{\pgfqpoint{3.614105in}{1.085951in}}%
\pgfpathlineto{\pgfqpoint{3.618749in}{1.099199in}}%
\pgfpathlineto{\pgfqpoint{3.610522in}{1.112447in}}%
\pgfpathlineto{\pgfqpoint{3.613375in}{1.125695in}}%
\pgfpathlineto{\pgfqpoint{3.608628in}{1.138943in}}%
\pgfpathlineto{\pgfqpoint{3.608843in}{1.152191in}}%
\pgfpathlineto{\pgfqpoint{3.608329in}{1.154170in}}%
\pgfpathlineto{\pgfqpoint{3.605475in}{1.165439in}}%
\pgfpathlineto{\pgfqpoint{3.602801in}{1.178687in}}%
\pgfpathlineto{\pgfqpoint{3.600699in}{1.191935in}}%
\pgfpathlineto{\pgfqpoint{3.596853in}{1.205183in}}%
\pgfpathlineto{\pgfqpoint{3.596203in}{1.218431in}}%
\pgfpathlineto{\pgfqpoint{3.589876in}{1.231679in}}%
\pgfpathlineto{\pgfqpoint{3.590815in}{1.244927in}}%
\pgfpathlineto{\pgfqpoint{3.584273in}{1.258175in}}%
\pgfpathlineto{\pgfqpoint{3.585104in}{1.271423in}}%
\pgfpathlineto{\pgfqpoint{3.583046in}{1.274752in}}%
\pgfpathlineto{\pgfqpoint{3.577245in}{1.284671in}}%
\pgfpathlineto{\pgfqpoint{3.577014in}{1.297919in}}%
\pgfpathlineto{\pgfqpoint{3.568655in}{1.311167in}}%
\pgfpathlineto{\pgfqpoint{3.569692in}{1.324415in}}%
\pgfpathlineto{\pgfqpoint{3.562965in}{1.337663in}}%
\pgfpathlineto{\pgfqpoint{3.562921in}{1.350911in}}%
\pgfpathlineto{\pgfqpoint{3.557763in}{1.360976in}}%
\pgfpathlineto{\pgfqpoint{3.556188in}{1.364159in}}%
\pgfpathlineto{\pgfqpoint{3.555328in}{1.377407in}}%
\pgfpathlineto{\pgfqpoint{3.549509in}{1.390655in}}%
\pgfpathlineto{\pgfqpoint{3.547754in}{1.403903in}}%
\pgfpathlineto{\pgfqpoint{3.542729in}{1.417151in}}%
\pgfpathlineto{\pgfqpoint{3.539385in}{1.430399in}}%
\pgfpathlineto{\pgfqpoint{3.535781in}{1.443647in}}%
\pgfpathlineto{\pgfqpoint{3.532480in}{1.454418in}}%
\pgfpathlineto{\pgfqpoint{3.531775in}{1.456895in}}%
\pgfpathlineto{\pgfqpoint{3.528558in}{1.470143in}}%
\pgfpathlineto{\pgfqpoint{3.523200in}{1.483391in}}%
\pgfpathlineto{\pgfqpoint{3.520730in}{1.496639in}}%
\pgfpathlineto{\pgfqpoint{3.515061in}{1.509887in}}%
\pgfpathlineto{\pgfqpoint{3.514355in}{1.523135in}}%
\pgfpathlineto{\pgfqpoint{3.508428in}{1.536383in}}%
\pgfpathlineto{\pgfqpoint{3.507196in}{1.548411in}}%
\pgfpathlineto{\pgfqpoint{3.507078in}{1.549631in}}%
\pgfpathlineto{\pgfqpoint{3.500283in}{1.562879in}}%
\pgfpathlineto{\pgfqpoint{3.499304in}{1.576127in}}%
\pgfpathlineto{\pgfqpoint{3.492288in}{1.589375in}}%
\pgfpathlineto{\pgfqpoint{3.491275in}{1.602623in}}%
\pgfpathlineto{\pgfqpoint{3.484216in}{1.615871in}}%
\pgfpathlineto{\pgfqpoint{3.482819in}{1.629119in}}%
\pgfpathlineto{\pgfqpoint{3.481913in}{1.631050in}}%
\pgfpathlineto{\pgfqpoint{3.476623in}{1.642367in}}%
\pgfpathlineto{\pgfqpoint{3.475008in}{1.655615in}}%
\pgfpathlineto{\pgfqpoint{3.468610in}{1.668863in}}%
\pgfpathlineto{\pgfqpoint{3.466078in}{1.682111in}}%
\pgfpathlineto{\pgfqpoint{3.460547in}{1.695359in}}%
\pgfpathlineto{\pgfqpoint{3.457776in}{1.708607in}}%
\pgfpathlineto{\pgfqpoint{3.456630in}{1.712138in}}%
\pgfpathlineto{\pgfqpoint{3.453349in}{1.721855in}}%
\pgfpathlineto{\pgfqpoint{3.449274in}{1.735103in}}%
\pgfpathlineto{\pgfqpoint{3.444200in}{1.748351in}}%
\pgfpathlineto{\pgfqpoint{3.438836in}{1.761599in}}%
\pgfpathlineto{\pgfqpoint{3.434918in}{1.774847in}}%
\pgfpathlineto{\pgfqpoint{3.431347in}{1.782644in}}%
\pgfpathlineto{\pgfqpoint{3.428641in}{1.788095in}}%
\pgfpathlineto{\pgfqpoint{3.425238in}{1.801343in}}%
\pgfpathlineto{\pgfqpoint{3.417829in}{1.814591in}}%
\pgfpathlineto{\pgfqpoint{3.414874in}{1.827839in}}%
\pgfpathlineto{\pgfqpoint{3.407136in}{1.841087in}}%
\pgfpathlineto{\pgfqpoint{3.406064in}{1.846689in}}%
\pgfpathlineto{\pgfqpoint{3.404205in}{1.854335in}}%
\pgfpathlineto{\pgfqpoint{3.392741in}{1.867583in}}%
\pgfpathlineto{\pgfqpoint{3.389626in}{1.880831in}}%
\pgfpathlineto{\pgfqpoint{3.380781in}{1.890930in}}%
\pgfpathlineto{\pgfqpoint{3.376176in}{1.894079in}}%
\pgfpathlineto{\pgfqpoint{3.372252in}{1.907327in}}%
\pgfpathlineto{\pgfqpoint{3.355497in}{1.918923in}}%
\pgfpathlineto{\pgfqpoint{3.343876in}{1.920575in}}%
\pgfpathlineto{\pgfqpoint{3.330214in}{1.926750in}}%
\pgfpathlineto{\pgfqpoint{3.325630in}{1.920575in}}%
\pgfpathlineto{\pgfqpoint{3.304931in}{1.910705in}}%
\pgfpathlineto{\pgfqpoint{3.301853in}{1.907327in}}%
\pgfpathlineto{\pgfqpoint{3.297288in}{1.894079in}}%
\pgfpathlineto{\pgfqpoint{3.285960in}{1.880831in}}%
\pgfpathlineto{\pgfqpoint{3.279648in}{1.869513in}}%
\pgfpathlineto{\pgfqpoint{3.278841in}{1.867583in}}%
\pgfpathlineto{\pgfqpoint{3.272842in}{1.854335in}}%
\pgfpathlineto{\pgfqpoint{3.265933in}{1.841087in}}%
\pgfpathlineto{\pgfqpoint{3.261301in}{1.827839in}}%
\pgfpathlineto{\pgfqpoint{3.254365in}{1.817195in}}%
\pgfpathlineto{\pgfqpoint{3.252886in}{1.814591in}}%
\pgfpathlineto{\pgfqpoint{3.250194in}{1.801343in}}%
\pgfpathlineto{\pgfqpoint{3.242799in}{1.788095in}}%
\pgfpathlineto{\pgfqpoint{3.241106in}{1.774847in}}%
\pgfpathlineto{\pgfqpoint{3.232566in}{1.761599in}}%
\pgfpathlineto{\pgfqpoint{3.230742in}{1.748351in}}%
\pgfpathlineto{\pgfqpoint{3.229082in}{1.745537in}}%
\pgfpathlineto{\pgfqpoint{3.223385in}{1.735103in}}%
\pgfpathlineto{\pgfqpoint{3.222374in}{1.721855in}}%
\pgfpathlineto{\pgfqpoint{3.215904in}{1.708607in}}%
\pgfpathlineto{\pgfqpoint{3.214051in}{1.695359in}}%
\pgfpathlineto{\pgfqpoint{3.207412in}{1.682111in}}%
\pgfpathlineto{\pgfqpoint{3.204519in}{1.668863in}}%
\pgfpathlineto{\pgfqpoint{3.203798in}{1.667078in}}%
\pgfpathlineto{\pgfqpoint{3.199287in}{1.655615in}}%
\pgfpathlineto{\pgfqpoint{3.195270in}{1.642367in}}%
\pgfpathlineto{\pgfqpoint{3.191077in}{1.629119in}}%
\pgfpathlineto{\pgfqpoint{3.186627in}{1.615871in}}%
\pgfpathlineto{\pgfqpoint{3.184295in}{1.602623in}}%
\pgfpathlineto{\pgfqpoint{3.179169in}{1.589375in}}%
\pgfpathlineto{\pgfqpoint{3.178515in}{1.584759in}}%
\pgfpathlineto{\pgfqpoint{3.177251in}{1.576127in}}%
\pgfpathlineto{\pgfqpoint{3.170226in}{1.562879in}}%
\pgfpathlineto{\pgfqpoint{3.169104in}{1.549631in}}%
\pgfpathlineto{\pgfqpoint{3.162652in}{1.536383in}}%
\pgfpathlineto{\pgfqpoint{3.162205in}{1.523135in}}%
\pgfpathlineto{\pgfqpoint{3.154966in}{1.509887in}}%
\pgfpathlineto{\pgfqpoint{3.154097in}{1.496639in}}%
\pgfpathlineto{\pgfqpoint{3.153232in}{1.495112in}}%
\pgfpathlineto{\pgfqpoint{3.146434in}{1.483391in}}%
\pgfpathlineto{\pgfqpoint{3.145693in}{1.470143in}}%
\pgfpathlineto{\pgfqpoint{3.139207in}{1.456895in}}%
\pgfpathlineto{\pgfqpoint{3.137855in}{1.443647in}}%
\pgfpathlineto{\pgfqpoint{3.132314in}{1.430399in}}%
\pgfpathlineto{\pgfqpoint{3.130722in}{1.417151in}}%
\pgfpathlineto{\pgfqpoint{3.127949in}{1.408560in}}%
\pgfpathlineto{\pgfqpoint{3.126384in}{1.403903in}}%
\pgfpathlineto{\pgfqpoint{3.122825in}{1.390655in}}%
\pgfpathlineto{\pgfqpoint{3.118547in}{1.377407in}}%
\pgfpathlineto{\pgfqpoint{3.113210in}{1.364159in}}%
\pgfpathlineto{\pgfqpoint{3.111259in}{1.350911in}}%
\pgfpathlineto{\pgfqpoint{3.106045in}{1.337663in}}%
\pgfpathlineto{\pgfqpoint{3.105068in}{1.324415in}}%
\pgfpathlineto{\pgfqpoint{3.102666in}{1.319937in}}%
\pgfpathlineto{\pgfqpoint{3.097660in}{1.311167in}}%
\pgfpathlineto{\pgfqpoint{3.097744in}{1.297919in}}%
\pgfpathlineto{\pgfqpoint{3.089957in}{1.284671in}}%
\pgfpathlineto{\pgfqpoint{3.091731in}{1.271423in}}%
\pgfpathlineto{\pgfqpoint{3.083574in}{1.258175in}}%
\pgfpathlineto{\pgfqpoint{3.085473in}{1.244927in}}%
\pgfpathlineto{\pgfqpoint{3.077382in}{1.232078in}}%
\pgfpathlineto{\pgfqpoint{3.077122in}{1.231679in}}%
\pgfpathlineto{\pgfqpoint{3.077382in}{1.230209in}}%
\pgfpathlineto{\pgfqpoint{3.079521in}{1.218431in}}%
\pgfpathlineto{\pgfqpoint{3.077382in}{1.215149in}}%
\pgfpathlineto{\pgfqpoint{3.070703in}{1.205183in}}%
\pgfpathlineto{\pgfqpoint{3.071373in}{1.191935in}}%
\pgfpathlineto{\pgfqpoint{3.063096in}{1.178687in}}%
\pgfpathlineto{\pgfqpoint{3.064203in}{1.165439in}}%
\pgfpathlineto{\pgfqpoint{3.059908in}{1.152191in}}%
\pgfpathlineto{\pgfqpoint{3.061199in}{1.138943in}}%
\pgfpathlineto{\pgfqpoint{3.057251in}{1.125695in}}%
\pgfpathlineto{\pgfqpoint{3.055673in}{1.112447in}}%
\pgfpathlineto{\pgfqpoint{3.054575in}{1.099199in}}%
\pgfpathlineto{\pgfqpoint{3.052252in}{1.085951in}}%
\pgfpathlineto{\pgfqpoint{3.054870in}{1.072703in}}%
\pgfpathlineto{\pgfqpoint{3.052099in}{1.063958in}}%
\pgfpathlineto{\pgfqpoint{3.050562in}{1.059455in}}%
\pgfpathlineto{\pgfqpoint{3.052099in}{1.055904in}}%
\pgfpathlineto{\pgfqpoint{3.056497in}{1.046207in}}%
\pgfpathlineto{\pgfqpoint{3.052099in}{1.036042in}}%
\pgfpathlineto{\pgfqpoint{3.050688in}{1.032959in}}%
\pgfpathlineto{\pgfqpoint{3.052099in}{1.031116in}}%
\pgfpathlineto{\pgfqpoint{3.061501in}{1.019711in}}%
\pgfpathlineto{\pgfqpoint{3.053687in}{1.006463in}}%
\pgfpathlineto{\pgfqpoint{3.068272in}{0.993215in}}%
\pgfpathlineto{\pgfqpoint{3.061345in}{0.979967in}}%
\pgfpathlineto{\pgfqpoint{3.077382in}{0.971376in}}%
\pgfpathlineto{\pgfqpoint{3.087261in}{0.966719in}}%
\pgfpathlineto{\pgfqpoint{3.085322in}{0.953471in}}%
\pgfpathlineto{\pgfqpoint{3.102666in}{0.946782in}}%
\pgfpathlineto{\pgfqpoint{3.123091in}{0.940223in}}%
\pgfpathlineto{\pgfqpoint{3.117903in}{0.926975in}}%
\pgfpathlineto{\pgfqpoint{3.127949in}{0.923901in}}%
\pgfpathlineto{\pgfqpoint{3.153232in}{0.917448in}}%
\pgfpathlineto{\pgfqpoint{3.171168in}{0.913727in}}%
\pgfpathlineto{\pgfqpoint{3.175588in}{0.900479in}}%
\pgfpathlineto{\pgfqpoint{3.178515in}{0.899991in}}%
\pgfpathclose%
\pgfpathmoveto{\pgfqpoint{3.192104in}{0.926975in}}%
\pgfpathlineto{\pgfqpoint{3.190514in}{0.940223in}}%
\pgfpathlineto{\pgfqpoint{3.178515in}{0.942567in}}%
\pgfpathlineto{\pgfqpoint{3.153232in}{0.948571in}}%
\pgfpathlineto{\pgfqpoint{3.136237in}{0.953471in}}%
\pgfpathlineto{\pgfqpoint{3.133540in}{0.966719in}}%
\pgfpathlineto{\pgfqpoint{3.127949in}{0.968741in}}%
\pgfpathlineto{\pgfqpoint{3.102666in}{0.979311in}}%
\pgfpathlineto{\pgfqpoint{3.101284in}{0.979967in}}%
\pgfpathlineto{\pgfqpoint{3.102666in}{0.984283in}}%
\pgfpathlineto{\pgfqpoint{3.105450in}{0.993215in}}%
\pgfpathlineto{\pgfqpoint{3.102666in}{0.995128in}}%
\pgfpathlineto{\pgfqpoint{3.087961in}{1.006463in}}%
\pgfpathlineto{\pgfqpoint{3.093039in}{1.019711in}}%
\pgfpathlineto{\pgfqpoint{3.080249in}{1.032959in}}%
\pgfpathlineto{\pgfqpoint{3.083867in}{1.046207in}}%
\pgfpathlineto{\pgfqpoint{3.077382in}{1.058496in}}%
\pgfpathlineto{\pgfqpoint{3.076906in}{1.059455in}}%
\pgfpathlineto{\pgfqpoint{3.077382in}{1.062341in}}%
\pgfpathlineto{\pgfqpoint{3.078979in}{1.072703in}}%
\pgfpathlineto{\pgfqpoint{3.077382in}{1.078982in}}%
\pgfpathlineto{\pgfqpoint{3.075695in}{1.085951in}}%
\pgfpathlineto{\pgfqpoint{3.076014in}{1.099199in}}%
\pgfpathlineto{\pgfqpoint{3.076846in}{1.112447in}}%
\pgfpathlineto{\pgfqpoint{3.076781in}{1.125695in}}%
\pgfpathlineto{\pgfqpoint{3.077382in}{1.127772in}}%
\pgfpathlineto{\pgfqpoint{3.080613in}{1.138943in}}%
\pgfpathlineto{\pgfqpoint{3.077815in}{1.152191in}}%
\pgfpathlineto{\pgfqpoint{3.082105in}{1.165439in}}%
\pgfpathlineto{\pgfqpoint{3.079834in}{1.178687in}}%
\pgfpathlineto{\pgfqpoint{3.088069in}{1.191935in}}%
\pgfpathlineto{\pgfqpoint{3.086386in}{1.205183in}}%
\pgfpathlineto{\pgfqpoint{3.094932in}{1.218431in}}%
\pgfpathlineto{\pgfqpoint{3.091699in}{1.231679in}}%
\pgfpathlineto{\pgfqpoint{3.099968in}{1.244927in}}%
\pgfpathlineto{\pgfqpoint{3.097451in}{1.258175in}}%
\pgfpathlineto{\pgfqpoint{3.102666in}{1.266786in}}%
\pgfpathlineto{\pgfqpoint{3.105368in}{1.271423in}}%
\pgfpathlineto{\pgfqpoint{3.103182in}{1.284671in}}%
\pgfpathlineto{\pgfqpoint{3.110519in}{1.297919in}}%
\pgfpathlineto{\pgfqpoint{3.110241in}{1.311167in}}%
\pgfpathlineto{\pgfqpoint{3.117152in}{1.324415in}}%
\pgfpathlineto{\pgfqpoint{3.118044in}{1.337663in}}%
\pgfpathlineto{\pgfqpoint{3.122882in}{1.350911in}}%
\pgfpathlineto{\pgfqpoint{3.124903in}{1.364159in}}%
\pgfpathlineto{\pgfqpoint{3.127949in}{1.372155in}}%
\pgfpathlineto{\pgfqpoint{3.129818in}{1.377407in}}%
\pgfpathlineto{\pgfqpoint{3.134069in}{1.390655in}}%
\pgfpathlineto{\pgfqpoint{3.136992in}{1.403903in}}%
\pgfpathlineto{\pgfqpoint{3.141377in}{1.417151in}}%
\pgfpathlineto{\pgfqpoint{3.142545in}{1.430399in}}%
\pgfpathlineto{\pgfqpoint{3.148325in}{1.443647in}}%
\pgfpathlineto{\pgfqpoint{3.149306in}{1.456895in}}%
\pgfpathlineto{\pgfqpoint{3.153232in}{1.464715in}}%
\pgfpathlineto{\pgfqpoint{3.155914in}{1.470143in}}%
\pgfpathlineto{\pgfqpoint{3.156322in}{1.483391in}}%
\pgfpathlineto{\pgfqpoint{3.163880in}{1.496639in}}%
\pgfpathlineto{\pgfqpoint{3.164538in}{1.509887in}}%
\pgfpathlineto{\pgfqpoint{3.171872in}{1.523135in}}%
\pgfpathlineto{\pgfqpoint{3.172151in}{1.536383in}}%
\pgfpathlineto{\pgfqpoint{3.178515in}{1.549339in}}%
\pgfpathlineto{\pgfqpoint{3.178656in}{1.549631in}}%
\pgfpathlineto{\pgfqpoint{3.179725in}{1.562879in}}%
\pgfpathlineto{\pgfqpoint{3.186489in}{1.576127in}}%
\pgfpathlineto{\pgfqpoint{3.188412in}{1.589375in}}%
\pgfpathlineto{\pgfqpoint{3.193375in}{1.602623in}}%
\pgfpathlineto{\pgfqpoint{3.195882in}{1.615871in}}%
\pgfpathlineto{\pgfqpoint{3.200210in}{1.629119in}}%
\pgfpathlineto{\pgfqpoint{3.203798in}{1.639809in}}%
\pgfpathlineto{\pgfqpoint{3.204690in}{1.642367in}}%
\pgfpathlineto{\pgfqpoint{3.208498in}{1.655615in}}%
\pgfpathlineto{\pgfqpoint{3.213987in}{1.668863in}}%
\pgfpathlineto{\pgfqpoint{3.216622in}{1.682111in}}%
\pgfpathlineto{\pgfqpoint{3.223599in}{1.695359in}}%
\pgfpathlineto{\pgfqpoint{3.225209in}{1.708607in}}%
\pgfpathlineto{\pgfqpoint{3.229082in}{1.716298in}}%
\pgfpathlineto{\pgfqpoint{3.232126in}{1.721855in}}%
\pgfpathlineto{\pgfqpoint{3.233000in}{1.735103in}}%
\pgfpathlineto{\pgfqpoint{3.240915in}{1.748351in}}%
\pgfpathlineto{\pgfqpoint{3.242695in}{1.761599in}}%
\pgfpathlineto{\pgfqpoint{3.251395in}{1.774847in}}%
\pgfpathlineto{\pgfqpoint{3.253134in}{1.788095in}}%
\pgfpathlineto{\pgfqpoint{3.254365in}{1.790286in}}%
\pgfpathlineto{\pgfqpoint{3.261621in}{1.801343in}}%
\pgfpathlineto{\pgfqpoint{3.265081in}{1.814591in}}%
\pgfpathlineto{\pgfqpoint{3.273552in}{1.827839in}}%
\pgfpathlineto{\pgfqpoint{3.278584in}{1.841087in}}%
\pgfpathlineto{\pgfqpoint{3.279648in}{1.843178in}}%
\pgfpathlineto{\pgfqpoint{3.287226in}{1.854335in}}%
\pgfpathlineto{\pgfqpoint{3.296526in}{1.867583in}}%
\pgfpathlineto{\pgfqpoint{3.303216in}{1.880831in}}%
\pgfpathlineto{\pgfqpoint{3.304931in}{1.882757in}}%
\pgfpathlineto{\pgfqpoint{3.328266in}{1.894079in}}%
\pgfpathlineto{\pgfqpoint{3.330214in}{1.897422in}}%
\pgfpathlineto{\pgfqpoint{3.336013in}{1.894079in}}%
\pgfpathlineto{\pgfqpoint{3.355497in}{1.891316in}}%
\pgfpathlineto{\pgfqpoint{3.370840in}{1.880831in}}%
\pgfpathlineto{\pgfqpoint{3.376362in}{1.867583in}}%
\pgfpathlineto{\pgfqpoint{3.380781in}{1.864437in}}%
\pgfpathlineto{\pgfqpoint{3.389514in}{1.854335in}}%
\pgfpathlineto{\pgfqpoint{3.392994in}{1.841087in}}%
\pgfpathlineto{\pgfqpoint{3.402685in}{1.827839in}}%
\pgfpathlineto{\pgfqpoint{3.406064in}{1.816880in}}%
\pgfpathlineto{\pgfqpoint{3.406608in}{1.814591in}}%
\pgfpathlineto{\pgfqpoint{3.414029in}{1.801343in}}%
\pgfpathlineto{\pgfqpoint{3.417717in}{1.788095in}}%
\pgfpathlineto{\pgfqpoint{3.424164in}{1.774847in}}%
\pgfpathlineto{\pgfqpoint{3.428819in}{1.761599in}}%
\pgfpathlineto{\pgfqpoint{3.431347in}{1.755789in}}%
\pgfpathlineto{\pgfqpoint{3.434321in}{1.748351in}}%
\pgfpathlineto{\pgfqpoint{3.439697in}{1.735103in}}%
\pgfpathlineto{\pgfqpoint{3.443564in}{1.721855in}}%
\pgfpathlineto{\pgfqpoint{3.448257in}{1.708607in}}%
\pgfpathlineto{\pgfqpoint{3.450831in}{1.695359in}}%
\pgfpathlineto{\pgfqpoint{3.456630in}{1.682621in}}%
\pgfpathlineto{\pgfqpoint{3.456848in}{1.682111in}}%
\pgfpathlineto{\pgfqpoint{3.459182in}{1.668863in}}%
\pgfpathlineto{\pgfqpoint{3.465847in}{1.655615in}}%
\pgfpathlineto{\pgfqpoint{3.467257in}{1.642367in}}%
\pgfpathlineto{\pgfqpoint{3.473664in}{1.629119in}}%
\pgfpathlineto{\pgfqpoint{3.474861in}{1.615871in}}%
\pgfpathlineto{\pgfqpoint{3.481913in}{1.602858in}}%
\pgfpathlineto{\pgfqpoint{3.482039in}{1.602623in}}%
\pgfpathlineto{\pgfqpoint{3.482871in}{1.589375in}}%
\pgfpathlineto{\pgfqpoint{3.490027in}{1.576127in}}%
\pgfpathlineto{\pgfqpoint{3.490889in}{1.562879in}}%
\pgfpathlineto{\pgfqpoint{3.497730in}{1.549631in}}%
\pgfpathlineto{\pgfqpoint{3.498932in}{1.536383in}}%
\pgfpathlineto{\pgfqpoint{3.504596in}{1.523135in}}%
\pgfpathlineto{\pgfqpoint{3.505219in}{1.509887in}}%
\pgfpathlineto{\pgfqpoint{3.507196in}{1.505044in}}%
\pgfpathlineto{\pgfqpoint{3.510768in}{1.496639in}}%
\pgfpathlineto{\pgfqpoint{3.513332in}{1.483391in}}%
\pgfpathlineto{\pgfqpoint{3.518500in}{1.470143in}}%
\pgfpathlineto{\pgfqpoint{3.521783in}{1.456895in}}%
\pgfpathlineto{\pgfqpoint{3.525351in}{1.443647in}}%
\pgfpathlineto{\pgfqpoint{3.528917in}{1.430399in}}%
\pgfpathlineto{\pgfqpoint{3.531829in}{1.417151in}}%
\pgfpathlineto{\pgfqpoint{3.532480in}{1.415455in}}%
\pgfpathlineto{\pgfqpoint{3.536995in}{1.403903in}}%
\pgfpathlineto{\pgfqpoint{3.538285in}{1.390655in}}%
\pgfpathlineto{\pgfqpoint{3.544220in}{1.377407in}}%
\pgfpathlineto{\pgfqpoint{3.544639in}{1.364159in}}%
\pgfpathlineto{\pgfqpoint{3.551245in}{1.350911in}}%
\pgfpathlineto{\pgfqpoint{3.550801in}{1.337663in}}%
\pgfpathlineto{\pgfqpoint{3.557380in}{1.324415in}}%
\pgfpathlineto{\pgfqpoint{3.555901in}{1.311167in}}%
\pgfpathlineto{\pgfqpoint{3.557763in}{1.308139in}}%
\pgfpathlineto{\pgfqpoint{3.564221in}{1.297919in}}%
\pgfpathlineto{\pgfqpoint{3.564074in}{1.284671in}}%
\pgfpathlineto{\pgfqpoint{3.571607in}{1.271423in}}%
\pgfpathlineto{\pgfqpoint{3.570349in}{1.258175in}}%
\pgfpathlineto{\pgfqpoint{3.576348in}{1.244927in}}%
\pgfpathlineto{\pgfqpoint{3.575089in}{1.231679in}}%
\pgfpathlineto{\pgfqpoint{3.580704in}{1.218431in}}%
\pgfpathlineto{\pgfqpoint{3.581070in}{1.205183in}}%
\pgfpathlineto{\pgfqpoint{3.583046in}{1.196141in}}%
\pgfpathlineto{\pgfqpoint{3.584027in}{1.191935in}}%
\pgfpathlineto{\pgfqpoint{3.585942in}{1.178687in}}%
\pgfpathlineto{\pgfqpoint{3.587640in}{1.165439in}}%
\pgfpathlineto{\pgfqpoint{3.590835in}{1.152191in}}%
\pgfpathlineto{\pgfqpoint{3.589190in}{1.138943in}}%
\pgfpathlineto{\pgfqpoint{3.593739in}{1.125695in}}%
\pgfpathlineto{\pgfqpoint{3.589281in}{1.112447in}}%
\pgfpathlineto{\pgfqpoint{3.597155in}{1.099199in}}%
\pgfpathlineto{\pgfqpoint{3.590654in}{1.085951in}}%
\pgfpathlineto{\pgfqpoint{3.596024in}{1.072703in}}%
\pgfpathlineto{\pgfqpoint{3.583532in}{1.059455in}}%
\pgfpathlineto{\pgfqpoint{3.589924in}{1.046207in}}%
\pgfpathlineto{\pgfqpoint{3.583046in}{1.039391in}}%
\pgfpathlineto{\pgfqpoint{3.575824in}{1.032959in}}%
\pgfpathlineto{\pgfqpoint{3.581838in}{1.019711in}}%
\pgfpathlineto{\pgfqpoint{3.566788in}{1.006463in}}%
\pgfpathlineto{\pgfqpoint{3.567240in}{0.993215in}}%
\pgfpathlineto{\pgfqpoint{3.557763in}{0.985183in}}%
\pgfpathlineto{\pgfqpoint{3.550863in}{0.979967in}}%
\pgfpathlineto{\pgfqpoint{3.543054in}{0.966719in}}%
\pgfpathlineto{\pgfqpoint{3.532480in}{0.960065in}}%
\pgfpathlineto{\pgfqpoint{3.520551in}{0.953471in}}%
\pgfpathlineto{\pgfqpoint{3.507196in}{0.947445in}}%
\pgfpathlineto{\pgfqpoint{3.485621in}{0.940223in}}%
\pgfpathlineto{\pgfqpoint{3.481913in}{0.937767in}}%
\pgfpathlineto{\pgfqpoint{3.463021in}{0.926975in}}%
\pgfpathlineto{\pgfqpoint{3.456630in}{0.925524in}}%
\pgfpathlineto{\pgfqpoint{3.431347in}{0.921106in}}%
\pgfpathlineto{\pgfqpoint{3.406064in}{0.918004in}}%
\pgfpathlineto{\pgfqpoint{3.380781in}{0.915981in}}%
\pgfpathlineto{\pgfqpoint{3.355497in}{0.914879in}}%
\pgfpathlineto{\pgfqpoint{3.330214in}{0.914597in}}%
\pgfpathlineto{\pgfqpoint{3.304931in}{0.915080in}}%
\pgfpathlineto{\pgfqpoint{3.279648in}{0.916308in}}%
\pgfpathlineto{\pgfqpoint{3.254365in}{0.918299in}}%
\pgfpathlineto{\pgfqpoint{3.229082in}{0.921105in}}%
\pgfpathlineto{\pgfqpoint{3.203798in}{0.924818in}}%
\pgfpathlineto{\pgfqpoint{3.192104in}{0.926975in}}%
\pgfpathclose%
\pgfusepath{fill}%
\end{pgfscope}%
\begin{pgfscope}%
\pgfpathrectangle{\pgfqpoint{2.470586in}{0.476543in}}{\pgfqpoint{1.820388in}{1.656000in}}%
\pgfusepath{clip}%
\pgfsetbuttcap%
\pgfsetroundjoin%
\definecolor{currentfill}{rgb}{0.973779,0.989850,0.797232}%
\pgfsetfillcolor{currentfill}%
\pgfsetlinewidth{0.000000pt}%
\definecolor{currentstroke}{rgb}{0.000000,0.000000,0.000000}%
\pgfsetstrokecolor{currentstroke}%
\pgfsetdash{}{0pt}%
\pgfpathmoveto{\pgfqpoint{3.203798in}{0.924818in}}%
\pgfpathlineto{\pgfqpoint{3.229082in}{0.921105in}}%
\pgfpathlineto{\pgfqpoint{3.254365in}{0.918299in}}%
\pgfpathlineto{\pgfqpoint{3.279648in}{0.916308in}}%
\pgfpathlineto{\pgfqpoint{3.304931in}{0.915080in}}%
\pgfpathlineto{\pgfqpoint{3.330214in}{0.914597in}}%
\pgfpathlineto{\pgfqpoint{3.355497in}{0.914879in}}%
\pgfpathlineto{\pgfqpoint{3.380781in}{0.915981in}}%
\pgfpathlineto{\pgfqpoint{3.406064in}{0.918004in}}%
\pgfpathlineto{\pgfqpoint{3.431347in}{0.921106in}}%
\pgfpathlineto{\pgfqpoint{3.456630in}{0.925524in}}%
\pgfpathlineto{\pgfqpoint{3.463021in}{0.926975in}}%
\pgfpathlineto{\pgfqpoint{3.481913in}{0.937767in}}%
\pgfpathlineto{\pgfqpoint{3.485621in}{0.940223in}}%
\pgfpathlineto{\pgfqpoint{3.507196in}{0.947445in}}%
\pgfpathlineto{\pgfqpoint{3.520551in}{0.953471in}}%
\pgfpathlineto{\pgfqpoint{3.532480in}{0.960065in}}%
\pgfpathlineto{\pgfqpoint{3.543054in}{0.966719in}}%
\pgfpathlineto{\pgfqpoint{3.550863in}{0.979967in}}%
\pgfpathlineto{\pgfqpoint{3.557763in}{0.985183in}}%
\pgfpathlineto{\pgfqpoint{3.567240in}{0.993215in}}%
\pgfpathlineto{\pgfqpoint{3.566788in}{1.006463in}}%
\pgfpathlineto{\pgfqpoint{3.581838in}{1.019711in}}%
\pgfpathlineto{\pgfqpoint{3.575824in}{1.032959in}}%
\pgfpathlineto{\pgfqpoint{3.583046in}{1.039391in}}%
\pgfpathlineto{\pgfqpoint{3.589924in}{1.046207in}}%
\pgfpathlineto{\pgfqpoint{3.583532in}{1.059455in}}%
\pgfpathlineto{\pgfqpoint{3.596024in}{1.072703in}}%
\pgfpathlineto{\pgfqpoint{3.590654in}{1.085951in}}%
\pgfpathlineto{\pgfqpoint{3.597155in}{1.099199in}}%
\pgfpathlineto{\pgfqpoint{3.589281in}{1.112447in}}%
\pgfpathlineto{\pgfqpoint{3.593739in}{1.125695in}}%
\pgfpathlineto{\pgfqpoint{3.589190in}{1.138943in}}%
\pgfpathlineto{\pgfqpoint{3.590835in}{1.152191in}}%
\pgfpathlineto{\pgfqpoint{3.587640in}{1.165439in}}%
\pgfpathlineto{\pgfqpoint{3.585942in}{1.178687in}}%
\pgfpathlineto{\pgfqpoint{3.584027in}{1.191935in}}%
\pgfpathlineto{\pgfqpoint{3.583046in}{1.196141in}}%
\pgfpathlineto{\pgfqpoint{3.581070in}{1.205183in}}%
\pgfpathlineto{\pgfqpoint{3.580704in}{1.218431in}}%
\pgfpathlineto{\pgfqpoint{3.575089in}{1.231679in}}%
\pgfpathlineto{\pgfqpoint{3.576348in}{1.244927in}}%
\pgfpathlineto{\pgfqpoint{3.570349in}{1.258175in}}%
\pgfpathlineto{\pgfqpoint{3.571607in}{1.271423in}}%
\pgfpathlineto{\pgfqpoint{3.564074in}{1.284671in}}%
\pgfpathlineto{\pgfqpoint{3.564221in}{1.297919in}}%
\pgfpathlineto{\pgfqpoint{3.557763in}{1.308139in}}%
\pgfpathlineto{\pgfqpoint{3.555901in}{1.311167in}}%
\pgfpathlineto{\pgfqpoint{3.557380in}{1.324415in}}%
\pgfpathlineto{\pgfqpoint{3.550801in}{1.337663in}}%
\pgfpathlineto{\pgfqpoint{3.551245in}{1.350911in}}%
\pgfpathlineto{\pgfqpoint{3.544639in}{1.364159in}}%
\pgfpathlineto{\pgfqpoint{3.544220in}{1.377407in}}%
\pgfpathlineto{\pgfqpoint{3.538285in}{1.390655in}}%
\pgfpathlineto{\pgfqpoint{3.536995in}{1.403903in}}%
\pgfpathlineto{\pgfqpoint{3.532480in}{1.415455in}}%
\pgfpathlineto{\pgfqpoint{3.531829in}{1.417151in}}%
\pgfpathlineto{\pgfqpoint{3.528917in}{1.430399in}}%
\pgfpathlineto{\pgfqpoint{3.525351in}{1.443647in}}%
\pgfpathlineto{\pgfqpoint{3.521783in}{1.456895in}}%
\pgfpathlineto{\pgfqpoint{3.518500in}{1.470143in}}%
\pgfpathlineto{\pgfqpoint{3.513332in}{1.483391in}}%
\pgfpathlineto{\pgfqpoint{3.510768in}{1.496639in}}%
\pgfpathlineto{\pgfqpoint{3.507196in}{1.505044in}}%
\pgfpathlineto{\pgfqpoint{3.505219in}{1.509887in}}%
\pgfpathlineto{\pgfqpoint{3.504596in}{1.523135in}}%
\pgfpathlineto{\pgfqpoint{3.498932in}{1.536383in}}%
\pgfpathlineto{\pgfqpoint{3.497730in}{1.549631in}}%
\pgfpathlineto{\pgfqpoint{3.490889in}{1.562879in}}%
\pgfpathlineto{\pgfqpoint{3.490027in}{1.576127in}}%
\pgfpathlineto{\pgfqpoint{3.482871in}{1.589375in}}%
\pgfpathlineto{\pgfqpoint{3.482039in}{1.602623in}}%
\pgfpathlineto{\pgfqpoint{3.481913in}{1.602858in}}%
\pgfpathlineto{\pgfqpoint{3.474861in}{1.615871in}}%
\pgfpathlineto{\pgfqpoint{3.473664in}{1.629119in}}%
\pgfpathlineto{\pgfqpoint{3.467257in}{1.642367in}}%
\pgfpathlineto{\pgfqpoint{3.465847in}{1.655615in}}%
\pgfpathlineto{\pgfqpoint{3.459182in}{1.668863in}}%
\pgfpathlineto{\pgfqpoint{3.456848in}{1.682111in}}%
\pgfpathlineto{\pgfqpoint{3.456630in}{1.682621in}}%
\pgfpathlineto{\pgfqpoint{3.450831in}{1.695359in}}%
\pgfpathlineto{\pgfqpoint{3.448257in}{1.708607in}}%
\pgfpathlineto{\pgfqpoint{3.443564in}{1.721855in}}%
\pgfpathlineto{\pgfqpoint{3.439697in}{1.735103in}}%
\pgfpathlineto{\pgfqpoint{3.434321in}{1.748351in}}%
\pgfpathlineto{\pgfqpoint{3.431347in}{1.755789in}}%
\pgfpathlineto{\pgfqpoint{3.428819in}{1.761599in}}%
\pgfpathlineto{\pgfqpoint{3.424164in}{1.774847in}}%
\pgfpathlineto{\pgfqpoint{3.417717in}{1.788095in}}%
\pgfpathlineto{\pgfqpoint{3.414029in}{1.801343in}}%
\pgfpathlineto{\pgfqpoint{3.406608in}{1.814591in}}%
\pgfpathlineto{\pgfqpoint{3.406064in}{1.816880in}}%
\pgfpathlineto{\pgfqpoint{3.402685in}{1.827839in}}%
\pgfpathlineto{\pgfqpoint{3.392994in}{1.841087in}}%
\pgfpathlineto{\pgfqpoint{3.389514in}{1.854335in}}%
\pgfpathlineto{\pgfqpoint{3.380781in}{1.864437in}}%
\pgfpathlineto{\pgfqpoint{3.376362in}{1.867583in}}%
\pgfpathlineto{\pgfqpoint{3.370840in}{1.880831in}}%
\pgfpathlineto{\pgfqpoint{3.355497in}{1.891316in}}%
\pgfpathlineto{\pgfqpoint{3.336013in}{1.894079in}}%
\pgfpathlineto{\pgfqpoint{3.330214in}{1.897422in}}%
\pgfpathlineto{\pgfqpoint{3.328266in}{1.894079in}}%
\pgfpathlineto{\pgfqpoint{3.304931in}{1.882757in}}%
\pgfpathlineto{\pgfqpoint{3.303216in}{1.880831in}}%
\pgfpathlineto{\pgfqpoint{3.296526in}{1.867583in}}%
\pgfpathlineto{\pgfqpoint{3.287226in}{1.854335in}}%
\pgfpathlineto{\pgfqpoint{3.279648in}{1.843178in}}%
\pgfpathlineto{\pgfqpoint{3.278584in}{1.841087in}}%
\pgfpathlineto{\pgfqpoint{3.273552in}{1.827839in}}%
\pgfpathlineto{\pgfqpoint{3.265081in}{1.814591in}}%
\pgfpathlineto{\pgfqpoint{3.261621in}{1.801343in}}%
\pgfpathlineto{\pgfqpoint{3.254365in}{1.790286in}}%
\pgfpathlineto{\pgfqpoint{3.253134in}{1.788095in}}%
\pgfpathlineto{\pgfqpoint{3.251395in}{1.774847in}}%
\pgfpathlineto{\pgfqpoint{3.242695in}{1.761599in}}%
\pgfpathlineto{\pgfqpoint{3.240915in}{1.748351in}}%
\pgfpathlineto{\pgfqpoint{3.233000in}{1.735103in}}%
\pgfpathlineto{\pgfqpoint{3.232126in}{1.721855in}}%
\pgfpathlineto{\pgfqpoint{3.229082in}{1.716298in}}%
\pgfpathlineto{\pgfqpoint{3.225209in}{1.708607in}}%
\pgfpathlineto{\pgfqpoint{3.223599in}{1.695359in}}%
\pgfpathlineto{\pgfqpoint{3.216622in}{1.682111in}}%
\pgfpathlineto{\pgfqpoint{3.213987in}{1.668863in}}%
\pgfpathlineto{\pgfqpoint{3.208498in}{1.655615in}}%
\pgfpathlineto{\pgfqpoint{3.204690in}{1.642367in}}%
\pgfpathlineto{\pgfqpoint{3.203798in}{1.639809in}}%
\pgfpathlineto{\pgfqpoint{3.200210in}{1.629119in}}%
\pgfpathlineto{\pgfqpoint{3.195882in}{1.615871in}}%
\pgfpathlineto{\pgfqpoint{3.193375in}{1.602623in}}%
\pgfpathlineto{\pgfqpoint{3.188412in}{1.589375in}}%
\pgfpathlineto{\pgfqpoint{3.186489in}{1.576127in}}%
\pgfpathlineto{\pgfqpoint{3.179725in}{1.562879in}}%
\pgfpathlineto{\pgfqpoint{3.178656in}{1.549631in}}%
\pgfpathlineto{\pgfqpoint{3.178515in}{1.549339in}}%
\pgfpathlineto{\pgfqpoint{3.172151in}{1.536383in}}%
\pgfpathlineto{\pgfqpoint{3.171872in}{1.523135in}}%
\pgfpathlineto{\pgfqpoint{3.164538in}{1.509887in}}%
\pgfpathlineto{\pgfqpoint{3.163880in}{1.496639in}}%
\pgfpathlineto{\pgfqpoint{3.156322in}{1.483391in}}%
\pgfpathlineto{\pgfqpoint{3.155914in}{1.470143in}}%
\pgfpathlineto{\pgfqpoint{3.153232in}{1.464715in}}%
\pgfpathlineto{\pgfqpoint{3.149306in}{1.456895in}}%
\pgfpathlineto{\pgfqpoint{3.148325in}{1.443647in}}%
\pgfpathlineto{\pgfqpoint{3.142545in}{1.430399in}}%
\pgfpathlineto{\pgfqpoint{3.141377in}{1.417151in}}%
\pgfpathlineto{\pgfqpoint{3.136992in}{1.403903in}}%
\pgfpathlineto{\pgfqpoint{3.134069in}{1.390655in}}%
\pgfpathlineto{\pgfqpoint{3.129818in}{1.377407in}}%
\pgfpathlineto{\pgfqpoint{3.127949in}{1.372155in}}%
\pgfpathlineto{\pgfqpoint{3.124903in}{1.364159in}}%
\pgfpathlineto{\pgfqpoint{3.122882in}{1.350911in}}%
\pgfpathlineto{\pgfqpoint{3.118044in}{1.337663in}}%
\pgfpathlineto{\pgfqpoint{3.117152in}{1.324415in}}%
\pgfpathlineto{\pgfqpoint{3.110241in}{1.311167in}}%
\pgfpathlineto{\pgfqpoint{3.110519in}{1.297919in}}%
\pgfpathlineto{\pgfqpoint{3.103182in}{1.284671in}}%
\pgfpathlineto{\pgfqpoint{3.105368in}{1.271423in}}%
\pgfpathlineto{\pgfqpoint{3.102666in}{1.266786in}}%
\pgfpathlineto{\pgfqpoint{3.097451in}{1.258175in}}%
\pgfpathlineto{\pgfqpoint{3.099968in}{1.244927in}}%
\pgfpathlineto{\pgfqpoint{3.091699in}{1.231679in}}%
\pgfpathlineto{\pgfqpoint{3.094932in}{1.218431in}}%
\pgfpathlineto{\pgfqpoint{3.086386in}{1.205183in}}%
\pgfpathlineto{\pgfqpoint{3.088069in}{1.191935in}}%
\pgfpathlineto{\pgfqpoint{3.079834in}{1.178687in}}%
\pgfpathlineto{\pgfqpoint{3.082105in}{1.165439in}}%
\pgfpathlineto{\pgfqpoint{3.077815in}{1.152191in}}%
\pgfpathlineto{\pgfqpoint{3.080613in}{1.138943in}}%
\pgfpathlineto{\pgfqpoint{3.077382in}{1.127772in}}%
\pgfpathlineto{\pgfqpoint{3.076781in}{1.125695in}}%
\pgfpathlineto{\pgfqpoint{3.076846in}{1.112447in}}%
\pgfpathlineto{\pgfqpoint{3.076014in}{1.099199in}}%
\pgfpathlineto{\pgfqpoint{3.075695in}{1.085951in}}%
\pgfpathlineto{\pgfqpoint{3.077382in}{1.078982in}}%
\pgfpathlineto{\pgfqpoint{3.078979in}{1.072703in}}%
\pgfpathlineto{\pgfqpoint{3.077382in}{1.062341in}}%
\pgfpathlineto{\pgfqpoint{3.076906in}{1.059455in}}%
\pgfpathlineto{\pgfqpoint{3.077382in}{1.058496in}}%
\pgfpathlineto{\pgfqpoint{3.083867in}{1.046207in}}%
\pgfpathlineto{\pgfqpoint{3.080249in}{1.032959in}}%
\pgfpathlineto{\pgfqpoint{3.093039in}{1.019711in}}%
\pgfpathlineto{\pgfqpoint{3.087961in}{1.006463in}}%
\pgfpathlineto{\pgfqpoint{3.102666in}{0.995128in}}%
\pgfpathlineto{\pgfqpoint{3.105450in}{0.993215in}}%
\pgfpathlineto{\pgfqpoint{3.102666in}{0.984283in}}%
\pgfpathlineto{\pgfqpoint{3.101284in}{0.979967in}}%
\pgfpathlineto{\pgfqpoint{3.102666in}{0.979311in}}%
\pgfpathlineto{\pgfqpoint{3.127949in}{0.968741in}}%
\pgfpathlineto{\pgfqpoint{3.133540in}{0.966719in}}%
\pgfpathlineto{\pgfqpoint{3.136237in}{0.953471in}}%
\pgfpathlineto{\pgfqpoint{3.153232in}{0.948571in}}%
\pgfpathlineto{\pgfqpoint{3.178515in}{0.942567in}}%
\pgfpathlineto{\pgfqpoint{3.190514in}{0.940223in}}%
\pgfpathlineto{\pgfqpoint{3.192104in}{0.926975in}}%
\pgfpathlineto{\pgfqpoint{3.203798in}{0.924818in}}%
\pgfpathclose%
\pgfpathmoveto{\pgfqpoint{3.200540in}{0.953471in}}%
\pgfpathlineto{\pgfqpoint{3.189082in}{0.966719in}}%
\pgfpathlineto{\pgfqpoint{3.178515in}{0.969549in}}%
\pgfpathlineto{\pgfqpoint{3.153232in}{0.977551in}}%
\pgfpathlineto{\pgfqpoint{3.146664in}{0.979967in}}%
\pgfpathlineto{\pgfqpoint{3.146621in}{0.993215in}}%
\pgfpathlineto{\pgfqpoint{3.127949in}{1.004325in}}%
\pgfpathlineto{\pgfqpoint{3.124789in}{1.006463in}}%
\pgfpathlineto{\pgfqpoint{3.126441in}{1.019711in}}%
\pgfpathlineto{\pgfqpoint{3.111473in}{1.032959in}}%
\pgfpathlineto{\pgfqpoint{3.112346in}{1.046207in}}%
\pgfpathlineto{\pgfqpoint{3.104240in}{1.059455in}}%
\pgfpathlineto{\pgfqpoint{3.103642in}{1.072703in}}%
\pgfpathlineto{\pgfqpoint{3.102666in}{1.075781in}}%
\pgfpathlineto{\pgfqpoint{3.099648in}{1.085951in}}%
\pgfpathlineto{\pgfqpoint{3.097767in}{1.099199in}}%
\pgfpathlineto{\pgfqpoint{3.098334in}{1.112447in}}%
\pgfpathlineto{\pgfqpoint{3.096463in}{1.125695in}}%
\pgfpathlineto{\pgfqpoint{3.100117in}{1.138943in}}%
\pgfpathlineto{\pgfqpoint{3.095716in}{1.152191in}}%
\pgfpathlineto{\pgfqpoint{3.099952in}{1.165439in}}%
\pgfpathlineto{\pgfqpoint{3.096461in}{1.178687in}}%
\pgfpathlineto{\pgfqpoint{3.102666in}{1.188694in}}%
\pgfpathlineto{\pgfqpoint{3.104691in}{1.191935in}}%
\pgfpathlineto{\pgfqpoint{3.102666in}{1.201501in}}%
\pgfpathlineto{\pgfqpoint{3.101942in}{1.205183in}}%
\pgfpathlineto{\pgfqpoint{3.102666in}{1.206315in}}%
\pgfpathlineto{\pgfqpoint{3.110324in}{1.218431in}}%
\pgfpathlineto{\pgfqpoint{3.106257in}{1.231679in}}%
\pgfpathlineto{\pgfqpoint{3.114302in}{1.244927in}}%
\pgfpathlineto{\pgfqpoint{3.111213in}{1.258175in}}%
\pgfpathlineto{\pgfqpoint{3.118723in}{1.271423in}}%
\pgfpathlineto{\pgfqpoint{3.116096in}{1.284671in}}%
\pgfpathlineto{\pgfqpoint{3.123117in}{1.297919in}}%
\pgfpathlineto{\pgfqpoint{3.122647in}{1.311167in}}%
\pgfpathlineto{\pgfqpoint{3.127949in}{1.321773in}}%
\pgfpathlineto{\pgfqpoint{3.129218in}{1.324415in}}%
\pgfpathlineto{\pgfqpoint{3.130016in}{1.337663in}}%
\pgfpathlineto{\pgfqpoint{3.134339in}{1.350911in}}%
\pgfpathlineto{\pgfqpoint{3.136374in}{1.364159in}}%
\pgfpathlineto{\pgfqpoint{3.140743in}{1.377407in}}%
\pgfpathlineto{\pgfqpoint{3.145106in}{1.390655in}}%
\pgfpathlineto{\pgfqpoint{3.147521in}{1.403903in}}%
\pgfpathlineto{\pgfqpoint{3.152032in}{1.417151in}}%
\pgfpathlineto{\pgfqpoint{3.152777in}{1.430399in}}%
\pgfpathlineto{\pgfqpoint{3.153232in}{1.431423in}}%
\pgfpathlineto{\pgfqpoint{3.158663in}{1.443647in}}%
\pgfpathlineto{\pgfqpoint{3.159229in}{1.456895in}}%
\pgfpathlineto{\pgfqpoint{3.165851in}{1.470143in}}%
\pgfpathlineto{\pgfqpoint{3.165938in}{1.483391in}}%
\pgfpathlineto{\pgfqpoint{3.173663in}{1.496639in}}%
\pgfpathlineto{\pgfqpoint{3.174111in}{1.509887in}}%
\pgfpathlineto{\pgfqpoint{3.178515in}{1.517774in}}%
\pgfpathlineto{\pgfqpoint{3.181504in}{1.523135in}}%
\pgfpathlineto{\pgfqpoint{3.181617in}{1.536383in}}%
\pgfpathlineto{\pgfqpoint{3.187971in}{1.549631in}}%
\pgfpathlineto{\pgfqpoint{3.189030in}{1.562879in}}%
\pgfpathlineto{\pgfqpoint{3.195678in}{1.576127in}}%
\pgfpathlineto{\pgfqpoint{3.197655in}{1.589375in}}%
\pgfpathlineto{\pgfqpoint{3.202455in}{1.602623in}}%
\pgfpathlineto{\pgfqpoint{3.203798in}{1.609321in}}%
\pgfpathlineto{\pgfqpoint{3.205168in}{1.615871in}}%
\pgfpathlineto{\pgfqpoint{3.209394in}{1.629119in}}%
\pgfpathlineto{\pgfqpoint{3.214175in}{1.642367in}}%
\pgfpathlineto{\pgfqpoint{3.217688in}{1.655615in}}%
\pgfpathlineto{\pgfqpoint{3.223455in}{1.668863in}}%
\pgfpathlineto{\pgfqpoint{3.225832in}{1.682111in}}%
\pgfpathlineto{\pgfqpoint{3.229082in}{1.688114in}}%
\pgfpathlineto{\pgfqpoint{3.233439in}{1.695359in}}%
\pgfpathlineto{\pgfqpoint{3.234872in}{1.708607in}}%
\pgfpathlineto{\pgfqpoint{3.242245in}{1.721855in}}%
\pgfpathlineto{\pgfqpoint{3.242926in}{1.735103in}}%
\pgfpathlineto{\pgfqpoint{3.251088in}{1.748351in}}%
\pgfpathlineto{\pgfqpoint{3.252824in}{1.761599in}}%
\pgfpathlineto{\pgfqpoint{3.254365in}{1.763932in}}%
\pgfpathlineto{\pgfqpoint{3.262904in}{1.774847in}}%
\pgfpathlineto{\pgfqpoint{3.265078in}{1.788095in}}%
\pgfpathlineto{\pgfqpoint{3.273691in}{1.801343in}}%
\pgfpathlineto{\pgfqpoint{3.277519in}{1.814591in}}%
\pgfpathlineto{\pgfqpoint{3.279648in}{1.817957in}}%
\pgfpathlineto{\pgfqpoint{3.288155in}{1.827839in}}%
\pgfpathlineto{\pgfqpoint{3.295907in}{1.841087in}}%
\pgfpathlineto{\pgfqpoint{3.304104in}{1.854335in}}%
\pgfpathlineto{\pgfqpoint{3.304931in}{1.855451in}}%
\pgfpathlineto{\pgfqpoint{3.326091in}{1.867583in}}%
\pgfpathlineto{\pgfqpoint{3.330214in}{1.872108in}}%
\pgfpathlineto{\pgfqpoint{3.341646in}{1.867583in}}%
\pgfpathlineto{\pgfqpoint{3.355497in}{1.865326in}}%
\pgfpathlineto{\pgfqpoint{3.370712in}{1.854335in}}%
\pgfpathlineto{\pgfqpoint{3.377191in}{1.841087in}}%
\pgfpathlineto{\pgfqpoint{3.380781in}{1.838021in}}%
\pgfpathlineto{\pgfqpoint{3.388133in}{1.827839in}}%
\pgfpathlineto{\pgfqpoint{3.392721in}{1.814591in}}%
\pgfpathlineto{\pgfqpoint{3.401921in}{1.801343in}}%
\pgfpathlineto{\pgfqpoint{3.406064in}{1.790580in}}%
\pgfpathlineto{\pgfqpoint{3.406793in}{1.788095in}}%
\pgfpathlineto{\pgfqpoint{3.413039in}{1.774847in}}%
\pgfpathlineto{\pgfqpoint{3.418072in}{1.761599in}}%
\pgfpathlineto{\pgfqpoint{3.423649in}{1.748351in}}%
\pgfpathlineto{\pgfqpoint{3.429980in}{1.735103in}}%
\pgfpathlineto{\pgfqpoint{3.431347in}{1.730598in}}%
\pgfpathlineto{\pgfqpoint{3.433780in}{1.721855in}}%
\pgfpathlineto{\pgfqpoint{3.438723in}{1.708607in}}%
\pgfpathlineto{\pgfqpoint{3.441017in}{1.695359in}}%
\pgfpathlineto{\pgfqpoint{3.447370in}{1.682111in}}%
\pgfpathlineto{\pgfqpoint{3.449510in}{1.668863in}}%
\pgfpathlineto{\pgfqpoint{3.456630in}{1.655727in}}%
\pgfpathlineto{\pgfqpoint{3.456687in}{1.655615in}}%
\pgfpathlineto{\pgfqpoint{3.457891in}{1.642367in}}%
\pgfpathlineto{\pgfqpoint{3.464535in}{1.629119in}}%
\pgfpathlineto{\pgfqpoint{3.465562in}{1.615871in}}%
\pgfpathlineto{\pgfqpoint{3.472916in}{1.602623in}}%
\pgfpathlineto{\pgfqpoint{3.473540in}{1.589375in}}%
\pgfpathlineto{\pgfqpoint{3.480746in}{1.576127in}}%
\pgfpathlineto{\pgfqpoint{3.481494in}{1.562879in}}%
\pgfpathlineto{\pgfqpoint{3.481913in}{1.562078in}}%
\pgfpathlineto{\pgfqpoint{3.488382in}{1.549631in}}%
\pgfpathlineto{\pgfqpoint{3.489492in}{1.536383in}}%
\pgfpathlineto{\pgfqpoint{3.495073in}{1.523135in}}%
\pgfpathlineto{\pgfqpoint{3.495647in}{1.509887in}}%
\pgfpathlineto{\pgfqpoint{3.500942in}{1.496639in}}%
\pgfpathlineto{\pgfqpoint{3.503546in}{1.483391in}}%
\pgfpathlineto{\pgfqpoint{3.507196in}{1.473410in}}%
\pgfpathlineto{\pgfqpoint{3.508442in}{1.470143in}}%
\pgfpathlineto{\pgfqpoint{3.511791in}{1.456895in}}%
\pgfpathlineto{\pgfqpoint{3.515073in}{1.443647in}}%
\pgfpathlineto{\pgfqpoint{3.518753in}{1.430399in}}%
\pgfpathlineto{\pgfqpoint{3.521271in}{1.417151in}}%
\pgfpathlineto{\pgfqpoint{3.526415in}{1.403903in}}%
\pgfpathlineto{\pgfqpoint{3.527177in}{1.390655in}}%
\pgfpathlineto{\pgfqpoint{3.532480in}{1.378804in}}%
\pgfpathlineto{\pgfqpoint{3.533112in}{1.377407in}}%
\pgfpathlineto{\pgfqpoint{3.533089in}{1.364159in}}%
\pgfpathlineto{\pgfqpoint{3.539790in}{1.350911in}}%
\pgfpathlineto{\pgfqpoint{3.538834in}{1.337663in}}%
\pgfpathlineto{\pgfqpoint{3.545422in}{1.324415in}}%
\pgfpathlineto{\pgfqpoint{3.543456in}{1.311167in}}%
\pgfpathlineto{\pgfqpoint{3.551543in}{1.297919in}}%
\pgfpathlineto{\pgfqpoint{3.551018in}{1.284671in}}%
\pgfpathlineto{\pgfqpoint{3.557763in}{1.272177in}}%
\pgfpathlineto{\pgfqpoint{3.558185in}{1.271423in}}%
\pgfpathlineto{\pgfqpoint{3.557763in}{1.268086in}}%
\pgfpathlineto{\pgfqpoint{3.556475in}{1.258175in}}%
\pgfpathlineto{\pgfqpoint{3.557763in}{1.255041in}}%
\pgfpathlineto{\pgfqpoint{3.562065in}{1.244927in}}%
\pgfpathlineto{\pgfqpoint{3.560467in}{1.231679in}}%
\pgfpathlineto{\pgfqpoint{3.565374in}{1.218431in}}%
\pgfpathlineto{\pgfqpoint{3.565464in}{1.205183in}}%
\pgfpathlineto{\pgfqpoint{3.567422in}{1.191935in}}%
\pgfpathlineto{\pgfqpoint{3.569138in}{1.178687in}}%
\pgfpathlineto{\pgfqpoint{3.569750in}{1.165439in}}%
\pgfpathlineto{\pgfqpoint{3.572775in}{1.152191in}}%
\pgfpathlineto{\pgfqpoint{3.569581in}{1.138943in}}%
\pgfpathlineto{\pgfqpoint{3.573984in}{1.125695in}}%
\pgfpathlineto{\pgfqpoint{3.567745in}{1.112447in}}%
\pgfpathlineto{\pgfqpoint{3.575314in}{1.099199in}}%
\pgfpathlineto{\pgfqpoint{3.566702in}{1.085951in}}%
\pgfpathlineto{\pgfqpoint{3.571279in}{1.072703in}}%
\pgfpathlineto{\pgfqpoint{3.557763in}{1.060763in}}%
\pgfpathlineto{\pgfqpoint{3.556068in}{1.059455in}}%
\pgfpathlineto{\pgfqpoint{3.557763in}{1.055378in}}%
\pgfpathlineto{\pgfqpoint{3.561550in}{1.046207in}}%
\pgfpathlineto{\pgfqpoint{3.557763in}{1.042980in}}%
\pgfpathlineto{\pgfqpoint{3.544479in}{1.032959in}}%
\pgfpathlineto{\pgfqpoint{3.548114in}{1.019711in}}%
\pgfpathlineto{\pgfqpoint{3.532480in}{1.008173in}}%
\pgfpathlineto{\pgfqpoint{3.529831in}{1.006463in}}%
\pgfpathlineto{\pgfqpoint{3.526167in}{0.993215in}}%
\pgfpathlineto{\pgfqpoint{3.507196in}{0.980951in}}%
\pgfpathlineto{\pgfqpoint{3.505441in}{0.979967in}}%
\pgfpathlineto{\pgfqpoint{3.487837in}{0.966719in}}%
\pgfpathlineto{\pgfqpoint{3.481913in}{0.963885in}}%
\pgfpathlineto{\pgfqpoint{3.456630in}{0.953604in}}%
\pgfpathlineto{\pgfqpoint{3.456238in}{0.953471in}}%
\pgfpathlineto{\pgfqpoint{3.431347in}{0.948103in}}%
\pgfpathlineto{\pgfqpoint{3.406064in}{0.944427in}}%
\pgfpathlineto{\pgfqpoint{3.380781in}{0.942137in}}%
\pgfpathlineto{\pgfqpoint{3.355497in}{0.940977in}}%
\pgfpathlineto{\pgfqpoint{3.330214in}{0.940784in}}%
\pgfpathlineto{\pgfqpoint{3.304931in}{0.941465in}}%
\pgfpathlineto{\pgfqpoint{3.279648in}{0.942980in}}%
\pgfpathlineto{\pgfqpoint{3.254365in}{0.945334in}}%
\pgfpathlineto{\pgfqpoint{3.229082in}{0.948576in}}%
\pgfpathlineto{\pgfqpoint{3.203798in}{0.952800in}}%
\pgfpathlineto{\pgfqpoint{3.200540in}{0.953471in}}%
\pgfpathclose%
\pgfusepath{fill}%
\end{pgfscope}%
\begin{pgfscope}%
\pgfpathrectangle{\pgfqpoint{2.470586in}{0.476543in}}{\pgfqpoint{1.820388in}{1.656000in}}%
\pgfusepath{clip}%
\pgfsetbuttcap%
\pgfsetroundjoin%
\definecolor{currentfill}{rgb}{0.999154,0.973779,0.709266}%
\pgfsetfillcolor{currentfill}%
\pgfsetlinewidth{0.000000pt}%
\definecolor{currentstroke}{rgb}{0.000000,0.000000,0.000000}%
\pgfsetstrokecolor{currentstroke}%
\pgfsetdash{}{0pt}%
\pgfpathmoveto{\pgfqpoint{3.203798in}{0.952800in}}%
\pgfpathlineto{\pgfqpoint{3.229082in}{0.948576in}}%
\pgfpathlineto{\pgfqpoint{3.254365in}{0.945334in}}%
\pgfpathlineto{\pgfqpoint{3.279648in}{0.942980in}}%
\pgfpathlineto{\pgfqpoint{3.304931in}{0.941465in}}%
\pgfpathlineto{\pgfqpoint{3.330214in}{0.940784in}}%
\pgfpathlineto{\pgfqpoint{3.355497in}{0.940977in}}%
\pgfpathlineto{\pgfqpoint{3.380781in}{0.942137in}}%
\pgfpathlineto{\pgfqpoint{3.406064in}{0.944427in}}%
\pgfpathlineto{\pgfqpoint{3.431347in}{0.948103in}}%
\pgfpathlineto{\pgfqpoint{3.456238in}{0.953471in}}%
\pgfpathlineto{\pgfqpoint{3.456630in}{0.953604in}}%
\pgfpathlineto{\pgfqpoint{3.481913in}{0.963885in}}%
\pgfpathlineto{\pgfqpoint{3.487837in}{0.966719in}}%
\pgfpathlineto{\pgfqpoint{3.505441in}{0.979967in}}%
\pgfpathlineto{\pgfqpoint{3.507196in}{0.980951in}}%
\pgfpathlineto{\pgfqpoint{3.526167in}{0.993215in}}%
\pgfpathlineto{\pgfqpoint{3.529831in}{1.006463in}}%
\pgfpathlineto{\pgfqpoint{3.532480in}{1.008173in}}%
\pgfpathlineto{\pgfqpoint{3.548114in}{1.019711in}}%
\pgfpathlineto{\pgfqpoint{3.544479in}{1.032959in}}%
\pgfpathlineto{\pgfqpoint{3.557763in}{1.042980in}}%
\pgfpathlineto{\pgfqpoint{3.561550in}{1.046207in}}%
\pgfpathlineto{\pgfqpoint{3.557763in}{1.055378in}}%
\pgfpathlineto{\pgfqpoint{3.556068in}{1.059455in}}%
\pgfpathlineto{\pgfqpoint{3.557763in}{1.060763in}}%
\pgfpathlineto{\pgfqpoint{3.571279in}{1.072703in}}%
\pgfpathlineto{\pgfqpoint{3.566702in}{1.085951in}}%
\pgfpathlineto{\pgfqpoint{3.575314in}{1.099199in}}%
\pgfpathlineto{\pgfqpoint{3.567745in}{1.112447in}}%
\pgfpathlineto{\pgfqpoint{3.573984in}{1.125695in}}%
\pgfpathlineto{\pgfqpoint{3.569581in}{1.138943in}}%
\pgfpathlineto{\pgfqpoint{3.572775in}{1.152191in}}%
\pgfpathlineto{\pgfqpoint{3.569750in}{1.165439in}}%
\pgfpathlineto{\pgfqpoint{3.569138in}{1.178687in}}%
\pgfpathlineto{\pgfqpoint{3.567422in}{1.191935in}}%
\pgfpathlineto{\pgfqpoint{3.565464in}{1.205183in}}%
\pgfpathlineto{\pgfqpoint{3.565374in}{1.218431in}}%
\pgfpathlineto{\pgfqpoint{3.560467in}{1.231679in}}%
\pgfpathlineto{\pgfqpoint{3.562065in}{1.244927in}}%
\pgfpathlineto{\pgfqpoint{3.557763in}{1.255041in}}%
\pgfpathlineto{\pgfqpoint{3.556475in}{1.258175in}}%
\pgfpathlineto{\pgfqpoint{3.557763in}{1.268086in}}%
\pgfpathlineto{\pgfqpoint{3.558185in}{1.271423in}}%
\pgfpathlineto{\pgfqpoint{3.557763in}{1.272177in}}%
\pgfpathlineto{\pgfqpoint{3.551018in}{1.284671in}}%
\pgfpathlineto{\pgfqpoint{3.551543in}{1.297919in}}%
\pgfpathlineto{\pgfqpoint{3.543456in}{1.311167in}}%
\pgfpathlineto{\pgfqpoint{3.545422in}{1.324415in}}%
\pgfpathlineto{\pgfqpoint{3.538834in}{1.337663in}}%
\pgfpathlineto{\pgfqpoint{3.539790in}{1.350911in}}%
\pgfpathlineto{\pgfqpoint{3.533089in}{1.364159in}}%
\pgfpathlineto{\pgfqpoint{3.533112in}{1.377407in}}%
\pgfpathlineto{\pgfqpoint{3.532480in}{1.378804in}}%
\pgfpathlineto{\pgfqpoint{3.527177in}{1.390655in}}%
\pgfpathlineto{\pgfqpoint{3.526415in}{1.403903in}}%
\pgfpathlineto{\pgfqpoint{3.521271in}{1.417151in}}%
\pgfpathlineto{\pgfqpoint{3.518753in}{1.430399in}}%
\pgfpathlineto{\pgfqpoint{3.515073in}{1.443647in}}%
\pgfpathlineto{\pgfqpoint{3.511791in}{1.456895in}}%
\pgfpathlineto{\pgfqpoint{3.508442in}{1.470143in}}%
\pgfpathlineto{\pgfqpoint{3.507196in}{1.473410in}}%
\pgfpathlineto{\pgfqpoint{3.503546in}{1.483391in}}%
\pgfpathlineto{\pgfqpoint{3.500942in}{1.496639in}}%
\pgfpathlineto{\pgfqpoint{3.495647in}{1.509887in}}%
\pgfpathlineto{\pgfqpoint{3.495073in}{1.523135in}}%
\pgfpathlineto{\pgfqpoint{3.489492in}{1.536383in}}%
\pgfpathlineto{\pgfqpoint{3.488382in}{1.549631in}}%
\pgfpathlineto{\pgfqpoint{3.481913in}{1.562078in}}%
\pgfpathlineto{\pgfqpoint{3.481494in}{1.562879in}}%
\pgfpathlineto{\pgfqpoint{3.480746in}{1.576127in}}%
\pgfpathlineto{\pgfqpoint{3.473540in}{1.589375in}}%
\pgfpathlineto{\pgfqpoint{3.472916in}{1.602623in}}%
\pgfpathlineto{\pgfqpoint{3.465562in}{1.615871in}}%
\pgfpathlineto{\pgfqpoint{3.464535in}{1.629119in}}%
\pgfpathlineto{\pgfqpoint{3.457891in}{1.642367in}}%
\pgfpathlineto{\pgfqpoint{3.456687in}{1.655615in}}%
\pgfpathlineto{\pgfqpoint{3.456630in}{1.655727in}}%
\pgfpathlineto{\pgfqpoint{3.449510in}{1.668863in}}%
\pgfpathlineto{\pgfqpoint{3.447370in}{1.682111in}}%
\pgfpathlineto{\pgfqpoint{3.441017in}{1.695359in}}%
\pgfpathlineto{\pgfqpoint{3.438723in}{1.708607in}}%
\pgfpathlineto{\pgfqpoint{3.433780in}{1.721855in}}%
\pgfpathlineto{\pgfqpoint{3.431347in}{1.730598in}}%
\pgfpathlineto{\pgfqpoint{3.429980in}{1.735103in}}%
\pgfpathlineto{\pgfqpoint{3.423649in}{1.748351in}}%
\pgfpathlineto{\pgfqpoint{3.418072in}{1.761599in}}%
\pgfpathlineto{\pgfqpoint{3.413039in}{1.774847in}}%
\pgfpathlineto{\pgfqpoint{3.406793in}{1.788095in}}%
\pgfpathlineto{\pgfqpoint{3.406064in}{1.790580in}}%
\pgfpathlineto{\pgfqpoint{3.401921in}{1.801343in}}%
\pgfpathlineto{\pgfqpoint{3.392721in}{1.814591in}}%
\pgfpathlineto{\pgfqpoint{3.388133in}{1.827839in}}%
\pgfpathlineto{\pgfqpoint{3.380781in}{1.838021in}}%
\pgfpathlineto{\pgfqpoint{3.377191in}{1.841087in}}%
\pgfpathlineto{\pgfqpoint{3.370712in}{1.854335in}}%
\pgfpathlineto{\pgfqpoint{3.355497in}{1.865326in}}%
\pgfpathlineto{\pgfqpoint{3.341646in}{1.867583in}}%
\pgfpathlineto{\pgfqpoint{3.330214in}{1.872108in}}%
\pgfpathlineto{\pgfqpoint{3.326091in}{1.867583in}}%
\pgfpathlineto{\pgfqpoint{3.304931in}{1.855451in}}%
\pgfpathlineto{\pgfqpoint{3.304104in}{1.854335in}}%
\pgfpathlineto{\pgfqpoint{3.295907in}{1.841087in}}%
\pgfpathlineto{\pgfqpoint{3.288155in}{1.827839in}}%
\pgfpathlineto{\pgfqpoint{3.279648in}{1.817957in}}%
\pgfpathlineto{\pgfqpoint{3.277519in}{1.814591in}}%
\pgfpathlineto{\pgfqpoint{3.273691in}{1.801343in}}%
\pgfpathlineto{\pgfqpoint{3.265078in}{1.788095in}}%
\pgfpathlineto{\pgfqpoint{3.262904in}{1.774847in}}%
\pgfpathlineto{\pgfqpoint{3.254365in}{1.763932in}}%
\pgfpathlineto{\pgfqpoint{3.252824in}{1.761599in}}%
\pgfpathlineto{\pgfqpoint{3.251088in}{1.748351in}}%
\pgfpathlineto{\pgfqpoint{3.242926in}{1.735103in}}%
\pgfpathlineto{\pgfqpoint{3.242245in}{1.721855in}}%
\pgfpathlineto{\pgfqpoint{3.234872in}{1.708607in}}%
\pgfpathlineto{\pgfqpoint{3.233439in}{1.695359in}}%
\pgfpathlineto{\pgfqpoint{3.229082in}{1.688114in}}%
\pgfpathlineto{\pgfqpoint{3.225832in}{1.682111in}}%
\pgfpathlineto{\pgfqpoint{3.223455in}{1.668863in}}%
\pgfpathlineto{\pgfqpoint{3.217688in}{1.655615in}}%
\pgfpathlineto{\pgfqpoint{3.214175in}{1.642367in}}%
\pgfpathlineto{\pgfqpoint{3.209394in}{1.629119in}}%
\pgfpathlineto{\pgfqpoint{3.205168in}{1.615871in}}%
\pgfpathlineto{\pgfqpoint{3.203798in}{1.609321in}}%
\pgfpathlineto{\pgfqpoint{3.202455in}{1.602623in}}%
\pgfpathlineto{\pgfqpoint{3.197655in}{1.589375in}}%
\pgfpathlineto{\pgfqpoint{3.195678in}{1.576127in}}%
\pgfpathlineto{\pgfqpoint{3.189030in}{1.562879in}}%
\pgfpathlineto{\pgfqpoint{3.187971in}{1.549631in}}%
\pgfpathlineto{\pgfqpoint{3.181617in}{1.536383in}}%
\pgfpathlineto{\pgfqpoint{3.181504in}{1.523135in}}%
\pgfpathlineto{\pgfqpoint{3.178515in}{1.517774in}}%
\pgfpathlineto{\pgfqpoint{3.174111in}{1.509887in}}%
\pgfpathlineto{\pgfqpoint{3.173663in}{1.496639in}}%
\pgfpathlineto{\pgfqpoint{3.165938in}{1.483391in}}%
\pgfpathlineto{\pgfqpoint{3.165851in}{1.470143in}}%
\pgfpathlineto{\pgfqpoint{3.159229in}{1.456895in}}%
\pgfpathlineto{\pgfqpoint{3.158663in}{1.443647in}}%
\pgfpathlineto{\pgfqpoint{3.153232in}{1.431423in}}%
\pgfpathlineto{\pgfqpoint{3.152777in}{1.430399in}}%
\pgfpathlineto{\pgfqpoint{3.152032in}{1.417151in}}%
\pgfpathlineto{\pgfqpoint{3.147521in}{1.403903in}}%
\pgfpathlineto{\pgfqpoint{3.145106in}{1.390655in}}%
\pgfpathlineto{\pgfqpoint{3.140743in}{1.377407in}}%
\pgfpathlineto{\pgfqpoint{3.136374in}{1.364159in}}%
\pgfpathlineto{\pgfqpoint{3.134339in}{1.350911in}}%
\pgfpathlineto{\pgfqpoint{3.130016in}{1.337663in}}%
\pgfpathlineto{\pgfqpoint{3.129218in}{1.324415in}}%
\pgfpathlineto{\pgfqpoint{3.127949in}{1.321773in}}%
\pgfpathlineto{\pgfqpoint{3.122647in}{1.311167in}}%
\pgfpathlineto{\pgfqpoint{3.123117in}{1.297919in}}%
\pgfpathlineto{\pgfqpoint{3.116096in}{1.284671in}}%
\pgfpathlineto{\pgfqpoint{3.118723in}{1.271423in}}%
\pgfpathlineto{\pgfqpoint{3.111213in}{1.258175in}}%
\pgfpathlineto{\pgfqpoint{3.114302in}{1.244927in}}%
\pgfpathlineto{\pgfqpoint{3.106257in}{1.231679in}}%
\pgfpathlineto{\pgfqpoint{3.110324in}{1.218431in}}%
\pgfpathlineto{\pgfqpoint{3.102666in}{1.206315in}}%
\pgfpathlineto{\pgfqpoint{3.101942in}{1.205183in}}%
\pgfpathlineto{\pgfqpoint{3.102666in}{1.201501in}}%
\pgfpathlineto{\pgfqpoint{3.104691in}{1.191935in}}%
\pgfpathlineto{\pgfqpoint{3.102666in}{1.188694in}}%
\pgfpathlineto{\pgfqpoint{3.096461in}{1.178687in}}%
\pgfpathlineto{\pgfqpoint{3.099952in}{1.165439in}}%
\pgfpathlineto{\pgfqpoint{3.095716in}{1.152191in}}%
\pgfpathlineto{\pgfqpoint{3.100117in}{1.138943in}}%
\pgfpathlineto{\pgfqpoint{3.096463in}{1.125695in}}%
\pgfpathlineto{\pgfqpoint{3.098334in}{1.112447in}}%
\pgfpathlineto{\pgfqpoint{3.097767in}{1.099199in}}%
\pgfpathlineto{\pgfqpoint{3.099648in}{1.085951in}}%
\pgfpathlineto{\pgfqpoint{3.102666in}{1.075781in}}%
\pgfpathlineto{\pgfqpoint{3.103642in}{1.072703in}}%
\pgfpathlineto{\pgfqpoint{3.104240in}{1.059455in}}%
\pgfpathlineto{\pgfqpoint{3.112346in}{1.046207in}}%
\pgfpathlineto{\pgfqpoint{3.111473in}{1.032959in}}%
\pgfpathlineto{\pgfqpoint{3.126441in}{1.019711in}}%
\pgfpathlineto{\pgfqpoint{3.124789in}{1.006463in}}%
\pgfpathlineto{\pgfqpoint{3.127949in}{1.004325in}}%
\pgfpathlineto{\pgfqpoint{3.146621in}{0.993215in}}%
\pgfpathlineto{\pgfqpoint{3.146664in}{0.979967in}}%
\pgfpathlineto{\pgfqpoint{3.153232in}{0.977551in}}%
\pgfpathlineto{\pgfqpoint{3.178515in}{0.969549in}}%
\pgfpathlineto{\pgfqpoint{3.189082in}{0.966719in}}%
\pgfpathlineto{\pgfqpoint{3.200540in}{0.953471in}}%
\pgfpathlineto{\pgfqpoint{3.203798in}{0.952800in}}%
\pgfpathclose%
\pgfpathmoveto{\pgfqpoint{3.275145in}{0.966719in}}%
\pgfpathlineto{\pgfqpoint{3.254365in}{0.969553in}}%
\pgfpathlineto{\pgfqpoint{3.229082in}{0.974088in}}%
\pgfpathlineto{\pgfqpoint{3.203798in}{0.979771in}}%
\pgfpathlineto{\pgfqpoint{3.203074in}{0.979967in}}%
\pgfpathlineto{\pgfqpoint{3.195594in}{0.993215in}}%
\pgfpathlineto{\pgfqpoint{3.178515in}{1.000636in}}%
\pgfpathlineto{\pgfqpoint{3.167136in}{1.006463in}}%
\pgfpathlineto{\pgfqpoint{3.163843in}{1.019711in}}%
\pgfpathlineto{\pgfqpoint{3.153232in}{1.026910in}}%
\pgfpathlineto{\pgfqpoint{3.145409in}{1.032959in}}%
\pgfpathlineto{\pgfqpoint{3.142742in}{1.046207in}}%
\pgfpathlineto{\pgfqpoint{3.133356in}{1.059455in}}%
\pgfpathlineto{\pgfqpoint{3.129394in}{1.072703in}}%
\pgfpathlineto{\pgfqpoint{3.127949in}{1.076423in}}%
\pgfpathlineto{\pgfqpoint{3.124555in}{1.085951in}}%
\pgfpathlineto{\pgfqpoint{3.120151in}{1.099199in}}%
\pgfpathlineto{\pgfqpoint{3.120456in}{1.112447in}}%
\pgfpathlineto{\pgfqpoint{3.116549in}{1.125695in}}%
\pgfpathlineto{\pgfqpoint{3.120088in}{1.138943in}}%
\pgfpathlineto{\pgfqpoint{3.113864in}{1.152191in}}%
\pgfpathlineto{\pgfqpoint{3.118088in}{1.165439in}}%
\pgfpathlineto{\pgfqpoint{3.113259in}{1.178687in}}%
\pgfpathlineto{\pgfqpoint{3.121442in}{1.191935in}}%
\pgfpathlineto{\pgfqpoint{3.117591in}{1.205183in}}%
\pgfpathlineto{\pgfqpoint{3.125695in}{1.218431in}}%
\pgfpathlineto{\pgfqpoint{3.120781in}{1.231679in}}%
\pgfpathlineto{\pgfqpoint{3.127949in}{1.243810in}}%
\pgfpathlineto{\pgfqpoint{3.128609in}{1.244927in}}%
\pgfpathlineto{\pgfqpoint{3.127949in}{1.247176in}}%
\pgfpathlineto{\pgfqpoint{3.124904in}{1.258175in}}%
\pgfpathlineto{\pgfqpoint{3.127949in}{1.263717in}}%
\pgfpathlineto{\pgfqpoint{3.132099in}{1.271423in}}%
\pgfpathlineto{\pgfqpoint{3.129017in}{1.284671in}}%
\pgfpathlineto{\pgfqpoint{3.135691in}{1.297919in}}%
\pgfpathlineto{\pgfqpoint{3.135043in}{1.311167in}}%
\pgfpathlineto{\pgfqpoint{3.141133in}{1.324415in}}%
\pgfpathlineto{\pgfqpoint{3.141855in}{1.337663in}}%
\pgfpathlineto{\pgfqpoint{3.145668in}{1.350911in}}%
\pgfpathlineto{\pgfqpoint{3.147768in}{1.364159in}}%
\pgfpathlineto{\pgfqpoint{3.151667in}{1.377407in}}%
\pgfpathlineto{\pgfqpoint{3.153232in}{1.382071in}}%
\pgfpathlineto{\pgfqpoint{3.156145in}{1.390655in}}%
\pgfpathlineto{\pgfqpoint{3.158012in}{1.403903in}}%
\pgfpathlineto{\pgfqpoint{3.162582in}{1.417151in}}%
\pgfpathlineto{\pgfqpoint{3.162841in}{1.430399in}}%
\pgfpathlineto{\pgfqpoint{3.168886in}{1.443647in}}%
\pgfpathlineto{\pgfqpoint{3.169041in}{1.456895in}}%
\pgfpathlineto{\pgfqpoint{3.175788in}{1.470143in}}%
\pgfpathlineto{\pgfqpoint{3.175554in}{1.483391in}}%
\pgfpathlineto{\pgfqpoint{3.178515in}{1.488417in}}%
\pgfpathlineto{\pgfqpoint{3.183463in}{1.496639in}}%
\pgfpathlineto{\pgfqpoint{3.183699in}{1.509887in}}%
\pgfpathlineto{\pgfqpoint{3.191059in}{1.523135in}}%
\pgfpathlineto{\pgfqpoint{3.191015in}{1.536383in}}%
\pgfpathlineto{\pgfqpoint{3.197287in}{1.549631in}}%
\pgfpathlineto{\pgfqpoint{3.198335in}{1.562879in}}%
\pgfpathlineto{\pgfqpoint{3.203798in}{1.573936in}}%
\pgfpathlineto{\pgfqpoint{3.204900in}{1.576127in}}%
\pgfpathlineto{\pgfqpoint{3.207002in}{1.589375in}}%
\pgfpathlineto{\pgfqpoint{3.211685in}{1.602623in}}%
\pgfpathlineto{\pgfqpoint{3.214633in}{1.615871in}}%
\pgfpathlineto{\pgfqpoint{3.218612in}{1.629119in}}%
\pgfpathlineto{\pgfqpoint{3.223660in}{1.642367in}}%
\pgfpathlineto{\pgfqpoint{3.226877in}{1.655615in}}%
\pgfpathlineto{\pgfqpoint{3.229082in}{1.660537in}}%
\pgfpathlineto{\pgfqpoint{3.233256in}{1.668863in}}%
\pgfpathlineto{\pgfqpoint{3.235514in}{1.682111in}}%
\pgfpathlineto{\pgfqpoint{3.243672in}{1.695359in}}%
\pgfpathlineto{\pgfqpoint{3.244791in}{1.708607in}}%
\pgfpathlineto{\pgfqpoint{3.252364in}{1.721855in}}%
\pgfpathlineto{\pgfqpoint{3.252852in}{1.735103in}}%
\pgfpathlineto{\pgfqpoint{3.254365in}{1.737536in}}%
\pgfpathlineto{\pgfqpoint{3.262508in}{1.748351in}}%
\pgfpathlineto{\pgfqpoint{3.264571in}{1.761599in}}%
\pgfpathlineto{\pgfqpoint{3.274909in}{1.774847in}}%
\pgfpathlineto{\pgfqpoint{3.277240in}{1.788095in}}%
\pgfpathlineto{\pgfqpoint{3.279648in}{1.791819in}}%
\pgfpathlineto{\pgfqpoint{3.288186in}{1.801343in}}%
\pgfpathlineto{\pgfqpoint{3.294353in}{1.814591in}}%
\pgfpathlineto{\pgfqpoint{3.304931in}{1.827635in}}%
\pgfpathlineto{\pgfqpoint{3.305263in}{1.827839in}}%
\pgfpathlineto{\pgfqpoint{3.324259in}{1.841087in}}%
\pgfpathlineto{\pgfqpoint{3.330214in}{1.846385in}}%
\pgfpathlineto{\pgfqpoint{3.346745in}{1.841087in}}%
\pgfpathlineto{\pgfqpoint{3.355497in}{1.839353in}}%
\pgfpathlineto{\pgfqpoint{3.368602in}{1.827839in}}%
\pgfpathlineto{\pgfqpoint{3.377393in}{1.814591in}}%
\pgfpathlineto{\pgfqpoint{3.380781in}{1.811442in}}%
\pgfpathlineto{\pgfqpoint{3.387601in}{1.801343in}}%
\pgfpathlineto{\pgfqpoint{3.393204in}{1.788095in}}%
\pgfpathlineto{\pgfqpoint{3.400712in}{1.774847in}}%
\pgfpathlineto{\pgfqpoint{3.406064in}{1.764768in}}%
\pgfpathlineto{\pgfqpoint{3.407324in}{1.761599in}}%
\pgfpathlineto{\pgfqpoint{3.412635in}{1.748351in}}%
\pgfpathlineto{\pgfqpoint{3.419318in}{1.735103in}}%
\pgfpathlineto{\pgfqpoint{3.423040in}{1.721855in}}%
\pgfpathlineto{\pgfqpoint{3.428913in}{1.708607in}}%
\pgfpathlineto{\pgfqpoint{3.431184in}{1.695359in}}%
\pgfpathlineto{\pgfqpoint{3.431347in}{1.695084in}}%
\pgfpathlineto{\pgfqpoint{3.437886in}{1.682111in}}%
\pgfpathlineto{\pgfqpoint{3.439747in}{1.668863in}}%
\pgfpathlineto{\pgfqpoint{3.447141in}{1.655615in}}%
\pgfpathlineto{\pgfqpoint{3.448148in}{1.642367in}}%
\pgfpathlineto{\pgfqpoint{3.455340in}{1.629119in}}%
\pgfpathlineto{\pgfqpoint{3.456243in}{1.615871in}}%
\pgfpathlineto{\pgfqpoint{3.456630in}{1.615238in}}%
\pgfpathlineto{\pgfqpoint{3.463794in}{1.602623in}}%
\pgfpathlineto{\pgfqpoint{3.464218in}{1.589375in}}%
\pgfpathlineto{\pgfqpoint{3.471440in}{1.576127in}}%
\pgfpathlineto{\pgfqpoint{3.472078in}{1.562879in}}%
\pgfpathlineto{\pgfqpoint{3.478989in}{1.549631in}}%
\pgfpathlineto{\pgfqpoint{3.480026in}{1.536383in}}%
\pgfpathlineto{\pgfqpoint{3.481913in}{1.531875in}}%
\pgfpathlineto{\pgfqpoint{3.485551in}{1.523135in}}%
\pgfpathlineto{\pgfqpoint{3.486075in}{1.509887in}}%
\pgfpathlineto{\pgfqpoint{3.491194in}{1.496639in}}%
\pgfpathlineto{\pgfqpoint{3.493895in}{1.483391in}}%
\pgfpathlineto{\pgfqpoint{3.498470in}{1.470143in}}%
\pgfpathlineto{\pgfqpoint{3.501846in}{1.456895in}}%
\pgfpathlineto{\pgfqpoint{3.504792in}{1.443647in}}%
\pgfpathlineto{\pgfqpoint{3.507196in}{1.435299in}}%
\pgfpathlineto{\pgfqpoint{3.508590in}{1.430399in}}%
\pgfpathlineto{\pgfqpoint{3.510714in}{1.417151in}}%
\pgfpathlineto{\pgfqpoint{3.515964in}{1.403903in}}%
\pgfpathlineto{\pgfqpoint{3.516193in}{1.390655in}}%
\pgfpathlineto{\pgfqpoint{3.522174in}{1.377407in}}%
\pgfpathlineto{\pgfqpoint{3.521671in}{1.364159in}}%
\pgfpathlineto{\pgfqpoint{3.528359in}{1.350911in}}%
\pgfpathlineto{\pgfqpoint{3.526880in}{1.337663in}}%
\pgfpathlineto{\pgfqpoint{3.532480in}{1.326394in}}%
\pgfpathlineto{\pgfqpoint{3.533464in}{1.324415in}}%
\pgfpathlineto{\pgfqpoint{3.532480in}{1.318970in}}%
\pgfpathlineto{\pgfqpoint{3.531003in}{1.311167in}}%
\pgfpathlineto{\pgfqpoint{3.532480in}{1.308708in}}%
\pgfpathlineto{\pgfqpoint{3.538983in}{1.297919in}}%
\pgfpathlineto{\pgfqpoint{3.538069in}{1.284671in}}%
\pgfpathlineto{\pgfqpoint{3.544840in}{1.271423in}}%
\pgfpathlineto{\pgfqpoint{3.542675in}{1.258175in}}%
\pgfpathlineto{\pgfqpoint{3.547736in}{1.244927in}}%
\pgfpathlineto{\pgfqpoint{3.545812in}{1.231679in}}%
\pgfpathlineto{\pgfqpoint{3.549938in}{1.218431in}}%
\pgfpathlineto{\pgfqpoint{3.549755in}{1.205183in}}%
\pgfpathlineto{\pgfqpoint{3.550675in}{1.191935in}}%
\pgfpathlineto{\pgfqpoint{3.552237in}{1.178687in}}%
\pgfpathlineto{\pgfqpoint{3.551674in}{1.165439in}}%
\pgfpathlineto{\pgfqpoint{3.554581in}{1.152191in}}%
\pgfpathlineto{\pgfqpoint{3.549609in}{1.138943in}}%
\pgfpathlineto{\pgfqpoint{3.553918in}{1.125695in}}%
\pgfpathlineto{\pgfqpoint{3.545604in}{1.112447in}}%
\pgfpathlineto{\pgfqpoint{3.552902in}{1.099199in}}%
\pgfpathlineto{\pgfqpoint{3.541791in}{1.085951in}}%
\pgfpathlineto{\pgfqpoint{3.545461in}{1.072703in}}%
\pgfpathlineto{\pgfqpoint{3.532480in}{1.063051in}}%
\pgfpathlineto{\pgfqpoint{3.526833in}{1.059455in}}%
\pgfpathlineto{\pgfqpoint{3.531292in}{1.046207in}}%
\pgfpathlineto{\pgfqpoint{3.510525in}{1.032959in}}%
\pgfpathlineto{\pgfqpoint{3.511010in}{1.019711in}}%
\pgfpathlineto{\pgfqpoint{3.507196in}{1.017277in}}%
\pgfpathlineto{\pgfqpoint{3.487493in}{1.006463in}}%
\pgfpathlineto{\pgfqpoint{3.481913in}{0.998338in}}%
\pgfpathlineto{\pgfqpoint{3.477325in}{0.993215in}}%
\pgfpathlineto{\pgfqpoint{3.456630in}{0.983351in}}%
\pgfpathlineto{\pgfqpoint{3.448091in}{0.979967in}}%
\pgfpathlineto{\pgfqpoint{3.431347in}{0.973544in}}%
\pgfpathlineto{\pgfqpoint{3.406064in}{0.967354in}}%
\pgfpathlineto{\pgfqpoint{3.401793in}{0.966719in}}%
\pgfpathlineto{\pgfqpoint{3.380781in}{0.962722in}}%
\pgfpathlineto{\pgfqpoint{3.355497in}{0.959774in}}%
\pgfpathlineto{\pgfqpoint{3.330214in}{0.958861in}}%
\pgfpathlineto{\pgfqpoint{3.304931in}{0.960410in}}%
\pgfpathlineto{\pgfqpoint{3.279648in}{0.965211in}}%
\pgfpathlineto{\pgfqpoint{3.275145in}{0.966719in}}%
\pgfpathclose%
\pgfusepath{fill}%
\end{pgfscope}%
\begin{pgfscope}%
\pgfpathrectangle{\pgfqpoint{2.470586in}{0.476543in}}{\pgfqpoint{1.820388in}{1.656000in}}%
\pgfusepath{clip}%
\pgfsetbuttcap%
\pgfsetroundjoin%
\definecolor{currentfill}{rgb}{0.997463,0.921338,0.629758}%
\pgfsetfillcolor{currentfill}%
\pgfsetlinewidth{0.000000pt}%
\definecolor{currentstroke}{rgb}{0.000000,0.000000,0.000000}%
\pgfsetstrokecolor{currentstroke}%
\pgfsetdash{}{0pt}%
\pgfpathmoveto{\pgfqpoint{3.279648in}{0.965211in}}%
\pgfpathlineto{\pgfqpoint{3.304931in}{0.960410in}}%
\pgfpathlineto{\pgfqpoint{3.330214in}{0.958861in}}%
\pgfpathlineto{\pgfqpoint{3.355497in}{0.959774in}}%
\pgfpathlineto{\pgfqpoint{3.380781in}{0.962722in}}%
\pgfpathlineto{\pgfqpoint{3.401793in}{0.966719in}}%
\pgfpathlineto{\pgfqpoint{3.406064in}{0.967354in}}%
\pgfpathlineto{\pgfqpoint{3.431347in}{0.973544in}}%
\pgfpathlineto{\pgfqpoint{3.448091in}{0.979967in}}%
\pgfpathlineto{\pgfqpoint{3.456630in}{0.983351in}}%
\pgfpathlineto{\pgfqpoint{3.477325in}{0.993215in}}%
\pgfpathlineto{\pgfqpoint{3.481913in}{0.998338in}}%
\pgfpathlineto{\pgfqpoint{3.487493in}{1.006463in}}%
\pgfpathlineto{\pgfqpoint{3.507196in}{1.017277in}}%
\pgfpathlineto{\pgfqpoint{3.511010in}{1.019711in}}%
\pgfpathlineto{\pgfqpoint{3.510525in}{1.032959in}}%
\pgfpathlineto{\pgfqpoint{3.531292in}{1.046207in}}%
\pgfpathlineto{\pgfqpoint{3.526833in}{1.059455in}}%
\pgfpathlineto{\pgfqpoint{3.532480in}{1.063051in}}%
\pgfpathlineto{\pgfqpoint{3.545461in}{1.072703in}}%
\pgfpathlineto{\pgfqpoint{3.541791in}{1.085951in}}%
\pgfpathlineto{\pgfqpoint{3.552902in}{1.099199in}}%
\pgfpathlineto{\pgfqpoint{3.545604in}{1.112447in}}%
\pgfpathlineto{\pgfqpoint{3.553918in}{1.125695in}}%
\pgfpathlineto{\pgfqpoint{3.549609in}{1.138943in}}%
\pgfpathlineto{\pgfqpoint{3.554581in}{1.152191in}}%
\pgfpathlineto{\pgfqpoint{3.551674in}{1.165439in}}%
\pgfpathlineto{\pgfqpoint{3.552237in}{1.178687in}}%
\pgfpathlineto{\pgfqpoint{3.550675in}{1.191935in}}%
\pgfpathlineto{\pgfqpoint{3.549755in}{1.205183in}}%
\pgfpathlineto{\pgfqpoint{3.549938in}{1.218431in}}%
\pgfpathlineto{\pgfqpoint{3.545812in}{1.231679in}}%
\pgfpathlineto{\pgfqpoint{3.547736in}{1.244927in}}%
\pgfpathlineto{\pgfqpoint{3.542675in}{1.258175in}}%
\pgfpathlineto{\pgfqpoint{3.544840in}{1.271423in}}%
\pgfpathlineto{\pgfqpoint{3.538069in}{1.284671in}}%
\pgfpathlineto{\pgfqpoint{3.538983in}{1.297919in}}%
\pgfpathlineto{\pgfqpoint{3.532480in}{1.308708in}}%
\pgfpathlineto{\pgfqpoint{3.531003in}{1.311167in}}%
\pgfpathlineto{\pgfqpoint{3.532480in}{1.318970in}}%
\pgfpathlineto{\pgfqpoint{3.533464in}{1.324415in}}%
\pgfpathlineto{\pgfqpoint{3.532480in}{1.326394in}}%
\pgfpathlineto{\pgfqpoint{3.526880in}{1.337663in}}%
\pgfpathlineto{\pgfqpoint{3.528359in}{1.350911in}}%
\pgfpathlineto{\pgfqpoint{3.521671in}{1.364159in}}%
\pgfpathlineto{\pgfqpoint{3.522174in}{1.377407in}}%
\pgfpathlineto{\pgfqpoint{3.516193in}{1.390655in}}%
\pgfpathlineto{\pgfqpoint{3.515964in}{1.403903in}}%
\pgfpathlineto{\pgfqpoint{3.510714in}{1.417151in}}%
\pgfpathlineto{\pgfqpoint{3.508590in}{1.430399in}}%
\pgfpathlineto{\pgfqpoint{3.507196in}{1.435299in}}%
\pgfpathlineto{\pgfqpoint{3.504792in}{1.443647in}}%
\pgfpathlineto{\pgfqpoint{3.501846in}{1.456895in}}%
\pgfpathlineto{\pgfqpoint{3.498470in}{1.470143in}}%
\pgfpathlineto{\pgfqpoint{3.493895in}{1.483391in}}%
\pgfpathlineto{\pgfqpoint{3.491194in}{1.496639in}}%
\pgfpathlineto{\pgfqpoint{3.486075in}{1.509887in}}%
\pgfpathlineto{\pgfqpoint{3.485551in}{1.523135in}}%
\pgfpathlineto{\pgfqpoint{3.481913in}{1.531875in}}%
\pgfpathlineto{\pgfqpoint{3.480026in}{1.536383in}}%
\pgfpathlineto{\pgfqpoint{3.478989in}{1.549631in}}%
\pgfpathlineto{\pgfqpoint{3.472078in}{1.562879in}}%
\pgfpathlineto{\pgfqpoint{3.471440in}{1.576127in}}%
\pgfpathlineto{\pgfqpoint{3.464218in}{1.589375in}}%
\pgfpathlineto{\pgfqpoint{3.463794in}{1.602623in}}%
\pgfpathlineto{\pgfqpoint{3.456630in}{1.615238in}}%
\pgfpathlineto{\pgfqpoint{3.456243in}{1.615871in}}%
\pgfpathlineto{\pgfqpoint{3.455340in}{1.629119in}}%
\pgfpathlineto{\pgfqpoint{3.448148in}{1.642367in}}%
\pgfpathlineto{\pgfqpoint{3.447141in}{1.655615in}}%
\pgfpathlineto{\pgfqpoint{3.439747in}{1.668863in}}%
\pgfpathlineto{\pgfqpoint{3.437886in}{1.682111in}}%
\pgfpathlineto{\pgfqpoint{3.431347in}{1.695084in}}%
\pgfpathlineto{\pgfqpoint{3.431184in}{1.695359in}}%
\pgfpathlineto{\pgfqpoint{3.428913in}{1.708607in}}%
\pgfpathlineto{\pgfqpoint{3.423040in}{1.721855in}}%
\pgfpathlineto{\pgfqpoint{3.419318in}{1.735103in}}%
\pgfpathlineto{\pgfqpoint{3.412635in}{1.748351in}}%
\pgfpathlineto{\pgfqpoint{3.407324in}{1.761599in}}%
\pgfpathlineto{\pgfqpoint{3.406064in}{1.764768in}}%
\pgfpathlineto{\pgfqpoint{3.400712in}{1.774847in}}%
\pgfpathlineto{\pgfqpoint{3.393204in}{1.788095in}}%
\pgfpathlineto{\pgfqpoint{3.387601in}{1.801343in}}%
\pgfpathlineto{\pgfqpoint{3.380781in}{1.811442in}}%
\pgfpathlineto{\pgfqpoint{3.377393in}{1.814591in}}%
\pgfpathlineto{\pgfqpoint{3.368602in}{1.827839in}}%
\pgfpathlineto{\pgfqpoint{3.355497in}{1.839353in}}%
\pgfpathlineto{\pgfqpoint{3.346745in}{1.841087in}}%
\pgfpathlineto{\pgfqpoint{3.330214in}{1.846385in}}%
\pgfpathlineto{\pgfqpoint{3.324259in}{1.841087in}}%
\pgfpathlineto{\pgfqpoint{3.305263in}{1.827839in}}%
\pgfpathlineto{\pgfqpoint{3.304931in}{1.827635in}}%
\pgfpathlineto{\pgfqpoint{3.294353in}{1.814591in}}%
\pgfpathlineto{\pgfqpoint{3.288186in}{1.801343in}}%
\pgfpathlineto{\pgfqpoint{3.279648in}{1.791819in}}%
\pgfpathlineto{\pgfqpoint{3.277240in}{1.788095in}}%
\pgfpathlineto{\pgfqpoint{3.274909in}{1.774847in}}%
\pgfpathlineto{\pgfqpoint{3.264571in}{1.761599in}}%
\pgfpathlineto{\pgfqpoint{3.262508in}{1.748351in}}%
\pgfpathlineto{\pgfqpoint{3.254365in}{1.737536in}}%
\pgfpathlineto{\pgfqpoint{3.252852in}{1.735103in}}%
\pgfpathlineto{\pgfqpoint{3.252364in}{1.721855in}}%
\pgfpathlineto{\pgfqpoint{3.244791in}{1.708607in}}%
\pgfpathlineto{\pgfqpoint{3.243672in}{1.695359in}}%
\pgfpathlineto{\pgfqpoint{3.235514in}{1.682111in}}%
\pgfpathlineto{\pgfqpoint{3.233256in}{1.668863in}}%
\pgfpathlineto{\pgfqpoint{3.229082in}{1.660537in}}%
\pgfpathlineto{\pgfqpoint{3.226877in}{1.655615in}}%
\pgfpathlineto{\pgfqpoint{3.223660in}{1.642367in}}%
\pgfpathlineto{\pgfqpoint{3.218612in}{1.629119in}}%
\pgfpathlineto{\pgfqpoint{3.214633in}{1.615871in}}%
\pgfpathlineto{\pgfqpoint{3.211685in}{1.602623in}}%
\pgfpathlineto{\pgfqpoint{3.207002in}{1.589375in}}%
\pgfpathlineto{\pgfqpoint{3.204900in}{1.576127in}}%
\pgfpathlineto{\pgfqpoint{3.203798in}{1.573936in}}%
\pgfpathlineto{\pgfqpoint{3.198335in}{1.562879in}}%
\pgfpathlineto{\pgfqpoint{3.197287in}{1.549631in}}%
\pgfpathlineto{\pgfqpoint{3.191015in}{1.536383in}}%
\pgfpathlineto{\pgfqpoint{3.191059in}{1.523135in}}%
\pgfpathlineto{\pgfqpoint{3.183699in}{1.509887in}}%
\pgfpathlineto{\pgfqpoint{3.183463in}{1.496639in}}%
\pgfpathlineto{\pgfqpoint{3.178515in}{1.488417in}}%
\pgfpathlineto{\pgfqpoint{3.175554in}{1.483391in}}%
\pgfpathlineto{\pgfqpoint{3.175788in}{1.470143in}}%
\pgfpathlineto{\pgfqpoint{3.169041in}{1.456895in}}%
\pgfpathlineto{\pgfqpoint{3.168886in}{1.443647in}}%
\pgfpathlineto{\pgfqpoint{3.162841in}{1.430399in}}%
\pgfpathlineto{\pgfqpoint{3.162582in}{1.417151in}}%
\pgfpathlineto{\pgfqpoint{3.158012in}{1.403903in}}%
\pgfpathlineto{\pgfqpoint{3.156145in}{1.390655in}}%
\pgfpathlineto{\pgfqpoint{3.153232in}{1.382071in}}%
\pgfpathlineto{\pgfqpoint{3.151667in}{1.377407in}}%
\pgfpathlineto{\pgfqpoint{3.147768in}{1.364159in}}%
\pgfpathlineto{\pgfqpoint{3.145668in}{1.350911in}}%
\pgfpathlineto{\pgfqpoint{3.141855in}{1.337663in}}%
\pgfpathlineto{\pgfqpoint{3.141133in}{1.324415in}}%
\pgfpathlineto{\pgfqpoint{3.135043in}{1.311167in}}%
\pgfpathlineto{\pgfqpoint{3.135691in}{1.297919in}}%
\pgfpathlineto{\pgfqpoint{3.129017in}{1.284671in}}%
\pgfpathlineto{\pgfqpoint{3.132099in}{1.271423in}}%
\pgfpathlineto{\pgfqpoint{3.127949in}{1.263717in}}%
\pgfpathlineto{\pgfqpoint{3.124904in}{1.258175in}}%
\pgfpathlineto{\pgfqpoint{3.127949in}{1.247176in}}%
\pgfpathlineto{\pgfqpoint{3.128609in}{1.244927in}}%
\pgfpathlineto{\pgfqpoint{3.127949in}{1.243810in}}%
\pgfpathlineto{\pgfqpoint{3.120781in}{1.231679in}}%
\pgfpathlineto{\pgfqpoint{3.125695in}{1.218431in}}%
\pgfpathlineto{\pgfqpoint{3.117591in}{1.205183in}}%
\pgfpathlineto{\pgfqpoint{3.121442in}{1.191935in}}%
\pgfpathlineto{\pgfqpoint{3.113259in}{1.178687in}}%
\pgfpathlineto{\pgfqpoint{3.118088in}{1.165439in}}%
\pgfpathlineto{\pgfqpoint{3.113864in}{1.152191in}}%
\pgfpathlineto{\pgfqpoint{3.120088in}{1.138943in}}%
\pgfpathlineto{\pgfqpoint{3.116549in}{1.125695in}}%
\pgfpathlineto{\pgfqpoint{3.120456in}{1.112447in}}%
\pgfpathlineto{\pgfqpoint{3.120151in}{1.099199in}}%
\pgfpathlineto{\pgfqpoint{3.124555in}{1.085951in}}%
\pgfpathlineto{\pgfqpoint{3.127949in}{1.076423in}}%
\pgfpathlineto{\pgfqpoint{3.129394in}{1.072703in}}%
\pgfpathlineto{\pgfqpoint{3.133356in}{1.059455in}}%
\pgfpathlineto{\pgfqpoint{3.142742in}{1.046207in}}%
\pgfpathlineto{\pgfqpoint{3.145409in}{1.032959in}}%
\pgfpathlineto{\pgfqpoint{3.153232in}{1.026910in}}%
\pgfpathlineto{\pgfqpoint{3.163843in}{1.019711in}}%
\pgfpathlineto{\pgfqpoint{3.167136in}{1.006463in}}%
\pgfpathlineto{\pgfqpoint{3.178515in}{1.000636in}}%
\pgfpathlineto{\pgfqpoint{3.195594in}{0.993215in}}%
\pgfpathlineto{\pgfqpoint{3.203074in}{0.979967in}}%
\pgfpathlineto{\pgfqpoint{3.203798in}{0.979771in}}%
\pgfpathlineto{\pgfqpoint{3.229082in}{0.974088in}}%
\pgfpathlineto{\pgfqpoint{3.254365in}{0.969553in}}%
\pgfpathlineto{\pgfqpoint{3.275145in}{0.966719in}}%
\pgfpathlineto{\pgfqpoint{3.279648in}{0.965211in}}%
\pgfpathclose%
\pgfpathmoveto{\pgfqpoint{3.268572in}{0.993215in}}%
\pgfpathlineto{\pgfqpoint{3.254365in}{0.996302in}}%
\pgfpathlineto{\pgfqpoint{3.229082in}{1.003590in}}%
\pgfpathlineto{\pgfqpoint{3.221143in}{1.006463in}}%
\pgfpathlineto{\pgfqpoint{3.208120in}{1.019711in}}%
\pgfpathlineto{\pgfqpoint{3.203798in}{1.021809in}}%
\pgfpathlineto{\pgfqpoint{3.184519in}{1.032959in}}%
\pgfpathlineto{\pgfqpoint{3.178515in}{1.041885in}}%
\pgfpathlineto{\pgfqpoint{3.176278in}{1.046207in}}%
\pgfpathlineto{\pgfqpoint{3.165427in}{1.059455in}}%
\pgfpathlineto{\pgfqpoint{3.157027in}{1.072703in}}%
\pgfpathlineto{\pgfqpoint{3.153232in}{1.080750in}}%
\pgfpathlineto{\pgfqpoint{3.151033in}{1.085951in}}%
\pgfpathlineto{\pgfqpoint{3.143588in}{1.099199in}}%
\pgfpathlineto{\pgfqpoint{3.143626in}{1.112447in}}%
\pgfpathlineto{\pgfqpoint{3.137297in}{1.125695in}}%
\pgfpathlineto{\pgfqpoint{3.140762in}{1.138943in}}%
\pgfpathlineto{\pgfqpoint{3.132372in}{1.152191in}}%
\pgfpathlineto{\pgfqpoint{3.136645in}{1.165439in}}%
\pgfpathlineto{\pgfqpoint{3.130254in}{1.178687in}}%
\pgfpathlineto{\pgfqpoint{3.138556in}{1.191935in}}%
\pgfpathlineto{\pgfqpoint{3.133424in}{1.205183in}}%
\pgfpathlineto{\pgfqpoint{3.141396in}{1.218431in}}%
\pgfpathlineto{\pgfqpoint{3.135495in}{1.231679in}}%
\pgfpathlineto{\pgfqpoint{3.143124in}{1.244927in}}%
\pgfpathlineto{\pgfqpoint{3.138776in}{1.258175in}}%
\pgfpathlineto{\pgfqpoint{3.145519in}{1.271423in}}%
\pgfpathlineto{\pgfqpoint{3.142026in}{1.284671in}}%
\pgfpathlineto{\pgfqpoint{3.148248in}{1.297919in}}%
\pgfpathlineto{\pgfqpoint{3.147432in}{1.311167in}}%
\pgfpathlineto{\pgfqpoint{3.153049in}{1.324415in}}%
\pgfpathlineto{\pgfqpoint{3.153232in}{1.328164in}}%
\pgfpathlineto{\pgfqpoint{3.153705in}{1.337663in}}%
\pgfpathlineto{\pgfqpoint{3.157048in}{1.350911in}}%
\pgfpathlineto{\pgfqpoint{3.159245in}{1.364159in}}%
\pgfpathlineto{\pgfqpoint{3.162628in}{1.377407in}}%
\pgfpathlineto{\pgfqpoint{3.167191in}{1.390655in}}%
\pgfpathlineto{\pgfqpoint{3.168459in}{1.403903in}}%
\pgfpathlineto{\pgfqpoint{3.173119in}{1.417151in}}%
\pgfpathlineto{\pgfqpoint{3.172898in}{1.430399in}}%
\pgfpathlineto{\pgfqpoint{3.178515in}{1.442401in}}%
\pgfpathlineto{\pgfqpoint{3.179123in}{1.443647in}}%
\pgfpathlineto{\pgfqpoint{3.178862in}{1.456895in}}%
\pgfpathlineto{\pgfqpoint{3.185831in}{1.470143in}}%
\pgfpathlineto{\pgfqpoint{3.185265in}{1.483391in}}%
\pgfpathlineto{\pgfqpoint{3.193280in}{1.496639in}}%
\pgfpathlineto{\pgfqpoint{3.193302in}{1.509887in}}%
\pgfpathlineto{\pgfqpoint{3.200614in}{1.523135in}}%
\pgfpathlineto{\pgfqpoint{3.200413in}{1.536383in}}%
\pgfpathlineto{\pgfqpoint{3.203798in}{1.543601in}}%
\pgfpathlineto{\pgfqpoint{3.206719in}{1.549631in}}%
\pgfpathlineto{\pgfqpoint{3.207818in}{1.562879in}}%
\pgfpathlineto{\pgfqpoint{3.214370in}{1.576127in}}%
\pgfpathlineto{\pgfqpoint{3.216554in}{1.589375in}}%
\pgfpathlineto{\pgfqpoint{3.220941in}{1.602623in}}%
\pgfpathlineto{\pgfqpoint{3.224099in}{1.615871in}}%
\pgfpathlineto{\pgfqpoint{3.227830in}{1.629119in}}%
\pgfpathlineto{\pgfqpoint{3.229082in}{1.632309in}}%
\pgfpathlineto{\pgfqpoint{3.233553in}{1.642367in}}%
\pgfpathlineto{\pgfqpoint{3.236691in}{1.655615in}}%
\pgfpathlineto{\pgfqpoint{3.243545in}{1.668863in}}%
\pgfpathlineto{\pgfqpoint{3.245453in}{1.682111in}}%
\pgfpathlineto{\pgfqpoint{3.253906in}{1.695359in}}%
\pgfpathlineto{\pgfqpoint{3.254365in}{1.702814in}}%
\pgfpathlineto{\pgfqpoint{3.254780in}{1.708607in}}%
\pgfpathlineto{\pgfqpoint{3.264019in}{1.721855in}}%
\pgfpathlineto{\pgfqpoint{3.264417in}{1.735103in}}%
\pgfpathlineto{\pgfqpoint{3.274521in}{1.748351in}}%
\pgfpathlineto{\pgfqpoint{3.276607in}{1.761599in}}%
\pgfpathlineto{\pgfqpoint{3.279648in}{1.765501in}}%
\pgfpathlineto{\pgfqpoint{3.289860in}{1.774847in}}%
\pgfpathlineto{\pgfqpoint{3.293636in}{1.788095in}}%
\pgfpathlineto{\pgfqpoint{3.304931in}{1.801206in}}%
\pgfpathlineto{\pgfqpoint{3.305174in}{1.801343in}}%
\pgfpathlineto{\pgfqpoint{3.321280in}{1.814591in}}%
\pgfpathlineto{\pgfqpoint{3.330214in}{1.820233in}}%
\pgfpathlineto{\pgfqpoint{3.351609in}{1.814591in}}%
\pgfpathlineto{\pgfqpoint{3.355497in}{1.813642in}}%
\pgfpathlineto{\pgfqpoint{3.368085in}{1.801343in}}%
\pgfpathlineto{\pgfqpoint{3.378559in}{1.788095in}}%
\pgfpathlineto{\pgfqpoint{3.380781in}{1.785421in}}%
\pgfpathlineto{\pgfqpoint{3.386365in}{1.774847in}}%
\pgfpathlineto{\pgfqpoint{3.393988in}{1.761599in}}%
\pgfpathlineto{\pgfqpoint{3.400306in}{1.748351in}}%
\pgfpathlineto{\pgfqpoint{3.406064in}{1.740084in}}%
\pgfpathlineto{\pgfqpoint{3.408655in}{1.735103in}}%
\pgfpathlineto{\pgfqpoint{3.411985in}{1.721855in}}%
\pgfpathlineto{\pgfqpoint{3.418161in}{1.708607in}}%
\pgfpathlineto{\pgfqpoint{3.420014in}{1.695359in}}%
\pgfpathlineto{\pgfqpoint{3.427994in}{1.682111in}}%
\pgfpathlineto{\pgfqpoint{3.429787in}{1.668863in}}%
\pgfpathlineto{\pgfqpoint{3.431347in}{1.666533in}}%
\pgfpathlineto{\pgfqpoint{3.437593in}{1.655615in}}%
\pgfpathlineto{\pgfqpoint{3.438345in}{1.642367in}}%
\pgfpathlineto{\pgfqpoint{3.445723in}{1.629119in}}%
\pgfpathlineto{\pgfqpoint{3.446429in}{1.615871in}}%
\pgfpathlineto{\pgfqpoint{3.454540in}{1.602623in}}%
\pgfpathlineto{\pgfqpoint{3.454779in}{1.589375in}}%
\pgfpathlineto{\pgfqpoint{3.456630in}{1.586205in}}%
\pgfpathlineto{\pgfqpoint{3.462134in}{1.576127in}}%
\pgfpathlineto{\pgfqpoint{3.462661in}{1.562879in}}%
\pgfpathlineto{\pgfqpoint{3.469498in}{1.549631in}}%
\pgfpathlineto{\pgfqpoint{3.470448in}{1.536383in}}%
\pgfpathlineto{\pgfqpoint{3.475862in}{1.523135in}}%
\pgfpathlineto{\pgfqpoint{3.476371in}{1.509887in}}%
\pgfpathlineto{\pgfqpoint{3.481428in}{1.496639in}}%
\pgfpathlineto{\pgfqpoint{3.481913in}{1.494441in}}%
\pgfpathlineto{\pgfqpoint{3.484244in}{1.483391in}}%
\pgfpathlineto{\pgfqpoint{3.488511in}{1.470143in}}%
\pgfpathlineto{\pgfqpoint{3.491941in}{1.456895in}}%
\pgfpathlineto{\pgfqpoint{3.494497in}{1.443647in}}%
\pgfpathlineto{\pgfqpoint{3.498404in}{1.430399in}}%
\pgfpathlineto{\pgfqpoint{3.500071in}{1.417151in}}%
\pgfpathlineto{\pgfqpoint{3.505485in}{1.403903in}}%
\pgfpathlineto{\pgfqpoint{3.505166in}{1.390655in}}%
\pgfpathlineto{\pgfqpoint{3.507196in}{1.386311in}}%
\pgfpathlineto{\pgfqpoint{3.511247in}{1.377407in}}%
\pgfpathlineto{\pgfqpoint{3.510259in}{1.364159in}}%
\pgfpathlineto{\pgfqpoint{3.516967in}{1.350911in}}%
\pgfpathlineto{\pgfqpoint{3.514940in}{1.337663in}}%
\pgfpathlineto{\pgfqpoint{3.521442in}{1.324415in}}%
\pgfpathlineto{\pgfqpoint{3.518491in}{1.311167in}}%
\pgfpathlineto{\pgfqpoint{3.526328in}{1.297919in}}%
\pgfpathlineto{\pgfqpoint{3.525005in}{1.284671in}}%
\pgfpathlineto{\pgfqpoint{3.531471in}{1.271423in}}%
\pgfpathlineto{\pgfqpoint{3.528787in}{1.258175in}}%
\pgfpathlineto{\pgfqpoint{3.532480in}{1.247512in}}%
\pgfpathlineto{\pgfqpoint{3.533388in}{1.244927in}}%
\pgfpathlineto{\pgfqpoint{3.532480in}{1.239483in}}%
\pgfpathlineto{\pgfqpoint{3.531107in}{1.231679in}}%
\pgfpathlineto{\pgfqpoint{3.532480in}{1.226065in}}%
\pgfpathlineto{\pgfqpoint{3.534396in}{1.218431in}}%
\pgfpathlineto{\pgfqpoint{3.533946in}{1.205183in}}%
\pgfpathlineto{\pgfqpoint{3.533725in}{1.191935in}}%
\pgfpathlineto{\pgfqpoint{3.535106in}{1.178687in}}%
\pgfpathlineto{\pgfqpoint{3.533261in}{1.165439in}}%
\pgfpathlineto{\pgfqpoint{3.535949in}{1.152191in}}%
\pgfpathlineto{\pgfqpoint{3.532480in}{1.145068in}}%
\pgfpathlineto{\pgfqpoint{3.529014in}{1.138943in}}%
\pgfpathlineto{\pgfqpoint{3.532480in}{1.128143in}}%
\pgfpathlineto{\pgfqpoint{3.533230in}{1.125695in}}%
\pgfpathlineto{\pgfqpoint{3.532480in}{1.124637in}}%
\pgfpathlineto{\pgfqpoint{3.522430in}{1.112447in}}%
\pgfpathlineto{\pgfqpoint{3.529472in}{1.099199in}}%
\pgfpathlineto{\pgfqpoint{3.515279in}{1.085951in}}%
\pgfpathlineto{\pgfqpoint{3.517839in}{1.072703in}}%
\pgfpathlineto{\pgfqpoint{3.507196in}{1.066036in}}%
\pgfpathlineto{\pgfqpoint{3.494661in}{1.059455in}}%
\pgfpathlineto{\pgfqpoint{3.497339in}{1.046207in}}%
\pgfpathlineto{\pgfqpoint{3.481913in}{1.037876in}}%
\pgfpathlineto{\pgfqpoint{3.471043in}{1.032959in}}%
\pgfpathlineto{\pgfqpoint{3.466340in}{1.019711in}}%
\pgfpathlineto{\pgfqpoint{3.456630in}{1.015214in}}%
\pgfpathlineto{\pgfqpoint{3.433626in}{1.006463in}}%
\pgfpathlineto{\pgfqpoint{3.431347in}{1.004924in}}%
\pgfpathlineto{\pgfqpoint{3.406064in}{0.993456in}}%
\pgfpathlineto{\pgfqpoint{3.405177in}{0.993215in}}%
\pgfpathlineto{\pgfqpoint{3.380781in}{0.987620in}}%
\pgfpathlineto{\pgfqpoint{3.355497in}{0.983842in}}%
\pgfpathlineto{\pgfqpoint{3.330214in}{0.982300in}}%
\pgfpathlineto{\pgfqpoint{3.304931in}{0.983543in}}%
\pgfpathlineto{\pgfqpoint{3.279648in}{0.988707in}}%
\pgfpathlineto{\pgfqpoint{3.268572in}{0.993215in}}%
\pgfpathclose%
\pgfusepath{fill}%
\end{pgfscope}%
\begin{pgfscope}%
\pgfpathrectangle{\pgfqpoint{2.470586in}{0.476543in}}{\pgfqpoint{1.820388in}{1.656000in}}%
\pgfusepath{clip}%
\pgfsetbuttcap%
\pgfsetroundjoin%
\definecolor{currentfill}{rgb}{0.995925,0.870742,0.557478}%
\pgfsetfillcolor{currentfill}%
\pgfsetlinewidth{0.000000pt}%
\definecolor{currentstroke}{rgb}{0.000000,0.000000,0.000000}%
\pgfsetstrokecolor{currentstroke}%
\pgfsetdash{}{0pt}%
\pgfpathmoveto{\pgfqpoint{3.279648in}{0.988707in}}%
\pgfpathlineto{\pgfqpoint{3.304931in}{0.983543in}}%
\pgfpathlineto{\pgfqpoint{3.330214in}{0.982300in}}%
\pgfpathlineto{\pgfqpoint{3.355497in}{0.983842in}}%
\pgfpathlineto{\pgfqpoint{3.380781in}{0.987620in}}%
\pgfpathlineto{\pgfqpoint{3.405177in}{0.993215in}}%
\pgfpathlineto{\pgfqpoint{3.406064in}{0.993456in}}%
\pgfpathlineto{\pgfqpoint{3.431347in}{1.004924in}}%
\pgfpathlineto{\pgfqpoint{3.433626in}{1.006463in}}%
\pgfpathlineto{\pgfqpoint{3.456630in}{1.015214in}}%
\pgfpathlineto{\pgfqpoint{3.466340in}{1.019711in}}%
\pgfpathlineto{\pgfqpoint{3.471043in}{1.032959in}}%
\pgfpathlineto{\pgfqpoint{3.481913in}{1.037876in}}%
\pgfpathlineto{\pgfqpoint{3.497339in}{1.046207in}}%
\pgfpathlineto{\pgfqpoint{3.494661in}{1.059455in}}%
\pgfpathlineto{\pgfqpoint{3.507196in}{1.066036in}}%
\pgfpathlineto{\pgfqpoint{3.517839in}{1.072703in}}%
\pgfpathlineto{\pgfqpoint{3.515279in}{1.085951in}}%
\pgfpathlineto{\pgfqpoint{3.529472in}{1.099199in}}%
\pgfpathlineto{\pgfqpoint{3.522430in}{1.112447in}}%
\pgfpathlineto{\pgfqpoint{3.532480in}{1.124637in}}%
\pgfpathlineto{\pgfqpoint{3.533230in}{1.125695in}}%
\pgfpathlineto{\pgfqpoint{3.532480in}{1.128143in}}%
\pgfpathlineto{\pgfqpoint{3.529014in}{1.138943in}}%
\pgfpathlineto{\pgfqpoint{3.532480in}{1.145068in}}%
\pgfpathlineto{\pgfqpoint{3.535949in}{1.152191in}}%
\pgfpathlineto{\pgfqpoint{3.533261in}{1.165439in}}%
\pgfpathlineto{\pgfqpoint{3.535106in}{1.178687in}}%
\pgfpathlineto{\pgfqpoint{3.533725in}{1.191935in}}%
\pgfpathlineto{\pgfqpoint{3.533946in}{1.205183in}}%
\pgfpathlineto{\pgfqpoint{3.534396in}{1.218431in}}%
\pgfpathlineto{\pgfqpoint{3.532480in}{1.226065in}}%
\pgfpathlineto{\pgfqpoint{3.531107in}{1.231679in}}%
\pgfpathlineto{\pgfqpoint{3.532480in}{1.239483in}}%
\pgfpathlineto{\pgfqpoint{3.533388in}{1.244927in}}%
\pgfpathlineto{\pgfqpoint{3.532480in}{1.247512in}}%
\pgfpathlineto{\pgfqpoint{3.528787in}{1.258175in}}%
\pgfpathlineto{\pgfqpoint{3.531471in}{1.271423in}}%
\pgfpathlineto{\pgfqpoint{3.525005in}{1.284671in}}%
\pgfpathlineto{\pgfqpoint{3.526328in}{1.297919in}}%
\pgfpathlineto{\pgfqpoint{3.518491in}{1.311167in}}%
\pgfpathlineto{\pgfqpoint{3.521442in}{1.324415in}}%
\pgfpathlineto{\pgfqpoint{3.514940in}{1.337663in}}%
\pgfpathlineto{\pgfqpoint{3.516967in}{1.350911in}}%
\pgfpathlineto{\pgfqpoint{3.510259in}{1.364159in}}%
\pgfpathlineto{\pgfqpoint{3.511247in}{1.377407in}}%
\pgfpathlineto{\pgfqpoint{3.507196in}{1.386311in}}%
\pgfpathlineto{\pgfqpoint{3.505166in}{1.390655in}}%
\pgfpathlineto{\pgfqpoint{3.505485in}{1.403903in}}%
\pgfpathlineto{\pgfqpoint{3.500071in}{1.417151in}}%
\pgfpathlineto{\pgfqpoint{3.498404in}{1.430399in}}%
\pgfpathlineto{\pgfqpoint{3.494497in}{1.443647in}}%
\pgfpathlineto{\pgfqpoint{3.491941in}{1.456895in}}%
\pgfpathlineto{\pgfqpoint{3.488511in}{1.470143in}}%
\pgfpathlineto{\pgfqpoint{3.484244in}{1.483391in}}%
\pgfpathlineto{\pgfqpoint{3.481913in}{1.494441in}}%
\pgfpathlineto{\pgfqpoint{3.481428in}{1.496639in}}%
\pgfpathlineto{\pgfqpoint{3.476371in}{1.509887in}}%
\pgfpathlineto{\pgfqpoint{3.475862in}{1.523135in}}%
\pgfpathlineto{\pgfqpoint{3.470448in}{1.536383in}}%
\pgfpathlineto{\pgfqpoint{3.469498in}{1.549631in}}%
\pgfpathlineto{\pgfqpoint{3.462661in}{1.562879in}}%
\pgfpathlineto{\pgfqpoint{3.462134in}{1.576127in}}%
\pgfpathlineto{\pgfqpoint{3.456630in}{1.586205in}}%
\pgfpathlineto{\pgfqpoint{3.454779in}{1.589375in}}%
\pgfpathlineto{\pgfqpoint{3.454540in}{1.602623in}}%
\pgfpathlineto{\pgfqpoint{3.446429in}{1.615871in}}%
\pgfpathlineto{\pgfqpoint{3.445723in}{1.629119in}}%
\pgfpathlineto{\pgfqpoint{3.438345in}{1.642367in}}%
\pgfpathlineto{\pgfqpoint{3.437593in}{1.655615in}}%
\pgfpathlineto{\pgfqpoint{3.431347in}{1.666533in}}%
\pgfpathlineto{\pgfqpoint{3.429787in}{1.668863in}}%
\pgfpathlineto{\pgfqpoint{3.427994in}{1.682111in}}%
\pgfpathlineto{\pgfqpoint{3.420014in}{1.695359in}}%
\pgfpathlineto{\pgfqpoint{3.418161in}{1.708607in}}%
\pgfpathlineto{\pgfqpoint{3.411985in}{1.721855in}}%
\pgfpathlineto{\pgfqpoint{3.408655in}{1.735103in}}%
\pgfpathlineto{\pgfqpoint{3.406064in}{1.740084in}}%
\pgfpathlineto{\pgfqpoint{3.400306in}{1.748351in}}%
\pgfpathlineto{\pgfqpoint{3.393988in}{1.761599in}}%
\pgfpathlineto{\pgfqpoint{3.386365in}{1.774847in}}%
\pgfpathlineto{\pgfqpoint{3.380781in}{1.785421in}}%
\pgfpathlineto{\pgfqpoint{3.378559in}{1.788095in}}%
\pgfpathlineto{\pgfqpoint{3.368085in}{1.801343in}}%
\pgfpathlineto{\pgfqpoint{3.355497in}{1.813642in}}%
\pgfpathlineto{\pgfqpoint{3.351609in}{1.814591in}}%
\pgfpathlineto{\pgfqpoint{3.330214in}{1.820233in}}%
\pgfpathlineto{\pgfqpoint{3.321280in}{1.814591in}}%
\pgfpathlineto{\pgfqpoint{3.305174in}{1.801343in}}%
\pgfpathlineto{\pgfqpoint{3.304931in}{1.801206in}}%
\pgfpathlineto{\pgfqpoint{3.293636in}{1.788095in}}%
\pgfpathlineto{\pgfqpoint{3.289860in}{1.774847in}}%
\pgfpathlineto{\pgfqpoint{3.279648in}{1.765501in}}%
\pgfpathlineto{\pgfqpoint{3.276607in}{1.761599in}}%
\pgfpathlineto{\pgfqpoint{3.274521in}{1.748351in}}%
\pgfpathlineto{\pgfqpoint{3.264417in}{1.735103in}}%
\pgfpathlineto{\pgfqpoint{3.264019in}{1.721855in}}%
\pgfpathlineto{\pgfqpoint{3.254780in}{1.708607in}}%
\pgfpathlineto{\pgfqpoint{3.254365in}{1.702814in}}%
\pgfpathlineto{\pgfqpoint{3.253906in}{1.695359in}}%
\pgfpathlineto{\pgfqpoint{3.245453in}{1.682111in}}%
\pgfpathlineto{\pgfqpoint{3.243545in}{1.668863in}}%
\pgfpathlineto{\pgfqpoint{3.236691in}{1.655615in}}%
\pgfpathlineto{\pgfqpoint{3.233553in}{1.642367in}}%
\pgfpathlineto{\pgfqpoint{3.229082in}{1.632309in}}%
\pgfpathlineto{\pgfqpoint{3.227830in}{1.629119in}}%
\pgfpathlineto{\pgfqpoint{3.224099in}{1.615871in}}%
\pgfpathlineto{\pgfqpoint{3.220941in}{1.602623in}}%
\pgfpathlineto{\pgfqpoint{3.216554in}{1.589375in}}%
\pgfpathlineto{\pgfqpoint{3.214370in}{1.576127in}}%
\pgfpathlineto{\pgfqpoint{3.207818in}{1.562879in}}%
\pgfpathlineto{\pgfqpoint{3.206719in}{1.549631in}}%
\pgfpathlineto{\pgfqpoint{3.203798in}{1.543601in}}%
\pgfpathlineto{\pgfqpoint{3.200413in}{1.536383in}}%
\pgfpathlineto{\pgfqpoint{3.200614in}{1.523135in}}%
\pgfpathlineto{\pgfqpoint{3.193302in}{1.509887in}}%
\pgfpathlineto{\pgfqpoint{3.193280in}{1.496639in}}%
\pgfpathlineto{\pgfqpoint{3.185265in}{1.483391in}}%
\pgfpathlineto{\pgfqpoint{3.185831in}{1.470143in}}%
\pgfpathlineto{\pgfqpoint{3.178862in}{1.456895in}}%
\pgfpathlineto{\pgfqpoint{3.179123in}{1.443647in}}%
\pgfpathlineto{\pgfqpoint{3.178515in}{1.442401in}}%
\pgfpathlineto{\pgfqpoint{3.172898in}{1.430399in}}%
\pgfpathlineto{\pgfqpoint{3.173119in}{1.417151in}}%
\pgfpathlineto{\pgfqpoint{3.168459in}{1.403903in}}%
\pgfpathlineto{\pgfqpoint{3.167191in}{1.390655in}}%
\pgfpathlineto{\pgfqpoint{3.162628in}{1.377407in}}%
\pgfpathlineto{\pgfqpoint{3.159245in}{1.364159in}}%
\pgfpathlineto{\pgfqpoint{3.157048in}{1.350911in}}%
\pgfpathlineto{\pgfqpoint{3.153705in}{1.337663in}}%
\pgfpathlineto{\pgfqpoint{3.153232in}{1.328164in}}%
\pgfpathlineto{\pgfqpoint{3.153049in}{1.324415in}}%
\pgfpathlineto{\pgfqpoint{3.147432in}{1.311167in}}%
\pgfpathlineto{\pgfqpoint{3.148248in}{1.297919in}}%
\pgfpathlineto{\pgfqpoint{3.142026in}{1.284671in}}%
\pgfpathlineto{\pgfqpoint{3.145519in}{1.271423in}}%
\pgfpathlineto{\pgfqpoint{3.138776in}{1.258175in}}%
\pgfpathlineto{\pgfqpoint{3.143124in}{1.244927in}}%
\pgfpathlineto{\pgfqpoint{3.135495in}{1.231679in}}%
\pgfpathlineto{\pgfqpoint{3.141396in}{1.218431in}}%
\pgfpathlineto{\pgfqpoint{3.133424in}{1.205183in}}%
\pgfpathlineto{\pgfqpoint{3.138556in}{1.191935in}}%
\pgfpathlineto{\pgfqpoint{3.130254in}{1.178687in}}%
\pgfpathlineto{\pgfqpoint{3.136645in}{1.165439in}}%
\pgfpathlineto{\pgfqpoint{3.132372in}{1.152191in}}%
\pgfpathlineto{\pgfqpoint{3.140762in}{1.138943in}}%
\pgfpathlineto{\pgfqpoint{3.137297in}{1.125695in}}%
\pgfpathlineto{\pgfqpoint{3.143626in}{1.112447in}}%
\pgfpathlineto{\pgfqpoint{3.143588in}{1.099199in}}%
\pgfpathlineto{\pgfqpoint{3.151033in}{1.085951in}}%
\pgfpathlineto{\pgfqpoint{3.153232in}{1.080750in}}%
\pgfpathlineto{\pgfqpoint{3.157027in}{1.072703in}}%
\pgfpathlineto{\pgfqpoint{3.165427in}{1.059455in}}%
\pgfpathlineto{\pgfqpoint{3.176278in}{1.046207in}}%
\pgfpathlineto{\pgfqpoint{3.178515in}{1.041885in}}%
\pgfpathlineto{\pgfqpoint{3.184519in}{1.032959in}}%
\pgfpathlineto{\pgfqpoint{3.203798in}{1.021809in}}%
\pgfpathlineto{\pgfqpoint{3.208120in}{1.019711in}}%
\pgfpathlineto{\pgfqpoint{3.221143in}{1.006463in}}%
\pgfpathlineto{\pgfqpoint{3.229082in}{1.003590in}}%
\pgfpathlineto{\pgfqpoint{3.254365in}{0.996302in}}%
\pgfpathlineto{\pgfqpoint{3.268572in}{0.993215in}}%
\pgfpathlineto{\pgfqpoint{3.279648in}{0.988707in}}%
\pgfpathclose%
\pgfpathmoveto{\pgfqpoint{3.274310in}{1.019711in}}%
\pgfpathlineto{\pgfqpoint{3.254365in}{1.025701in}}%
\pgfpathlineto{\pgfqpoint{3.235794in}{1.032959in}}%
\pgfpathlineto{\pgfqpoint{3.229082in}{1.036543in}}%
\pgfpathlineto{\pgfqpoint{3.216894in}{1.046207in}}%
\pgfpathlineto{\pgfqpoint{3.203798in}{1.058037in}}%
\pgfpathlineto{\pgfqpoint{3.202463in}{1.059455in}}%
\pgfpathlineto{\pgfqpoint{3.187843in}{1.072703in}}%
\pgfpathlineto{\pgfqpoint{3.180226in}{1.085951in}}%
\pgfpathlineto{\pgfqpoint{3.178515in}{1.087555in}}%
\pgfpathlineto{\pgfqpoint{3.168721in}{1.099199in}}%
\pgfpathlineto{\pgfqpoint{3.168488in}{1.112447in}}%
\pgfpathlineto{\pgfqpoint{3.159104in}{1.125695in}}%
\pgfpathlineto{\pgfqpoint{3.162540in}{1.138943in}}%
\pgfpathlineto{\pgfqpoint{3.153232in}{1.149855in}}%
\pgfpathlineto{\pgfqpoint{3.151554in}{1.152191in}}%
\pgfpathlineto{\pgfqpoint{3.153232in}{1.157600in}}%
\pgfpathlineto{\pgfqpoint{3.155818in}{1.165439in}}%
\pgfpathlineto{\pgfqpoint{3.153232in}{1.169342in}}%
\pgfpathlineto{\pgfqpoint{3.147881in}{1.178687in}}%
\pgfpathlineto{\pgfqpoint{3.153232in}{1.187479in}}%
\pgfpathlineto{\pgfqpoint{3.156076in}{1.191935in}}%
\pgfpathlineto{\pgfqpoint{3.153232in}{1.197330in}}%
\pgfpathlineto{\pgfqpoint{3.149607in}{1.205183in}}%
\pgfpathlineto{\pgfqpoint{3.153232in}{1.211460in}}%
\pgfpathlineto{\pgfqpoint{3.157372in}{1.218431in}}%
\pgfpathlineto{\pgfqpoint{3.153232in}{1.225935in}}%
\pgfpathlineto{\pgfqpoint{3.150394in}{1.231679in}}%
\pgfpathlineto{\pgfqpoint{3.153232in}{1.236786in}}%
\pgfpathlineto{\pgfqpoint{3.157849in}{1.244927in}}%
\pgfpathlineto{\pgfqpoint{3.153232in}{1.256691in}}%
\pgfpathlineto{\pgfqpoint{3.152698in}{1.258175in}}%
\pgfpathlineto{\pgfqpoint{3.153232in}{1.259270in}}%
\pgfpathlineto{\pgfqpoint{3.159159in}{1.271423in}}%
\pgfpathlineto{\pgfqpoint{3.155111in}{1.284671in}}%
\pgfpathlineto{\pgfqpoint{3.161049in}{1.297919in}}%
\pgfpathlineto{\pgfqpoint{3.160054in}{1.311167in}}%
\pgfpathlineto{\pgfqpoint{3.165234in}{1.324415in}}%
\pgfpathlineto{\pgfqpoint{3.165836in}{1.337663in}}%
\pgfpathlineto{\pgfqpoint{3.168530in}{1.350911in}}%
\pgfpathlineto{\pgfqpoint{3.170800in}{1.364159in}}%
\pgfpathlineto{\pgfqpoint{3.173594in}{1.377407in}}%
\pgfpathlineto{\pgfqpoint{3.178238in}{1.390655in}}%
\pgfpathlineto{\pgfqpoint{3.178515in}{1.395970in}}%
\pgfpathlineto{\pgfqpoint{3.178922in}{1.403903in}}%
\pgfpathlineto{\pgfqpoint{3.183846in}{1.417151in}}%
\pgfpathlineto{\pgfqpoint{3.183101in}{1.430399in}}%
\pgfpathlineto{\pgfqpoint{3.189614in}{1.443647in}}%
\pgfpathlineto{\pgfqpoint{3.188906in}{1.456895in}}%
\pgfpathlineto{\pgfqpoint{3.195913in}{1.470143in}}%
\pgfpathlineto{\pgfqpoint{3.195019in}{1.483391in}}%
\pgfpathlineto{\pgfqpoint{3.203098in}{1.496639in}}%
\pgfpathlineto{\pgfqpoint{3.202905in}{1.509887in}}%
\pgfpathlineto{\pgfqpoint{3.203798in}{1.511510in}}%
\pgfpathlineto{\pgfqpoint{3.210506in}{1.523135in}}%
\pgfpathlineto{\pgfqpoint{3.210150in}{1.536383in}}%
\pgfpathlineto{\pgfqpoint{3.216425in}{1.549631in}}%
\pgfpathlineto{\pgfqpoint{3.217552in}{1.562879in}}%
\pgfpathlineto{\pgfqpoint{3.223839in}{1.576127in}}%
\pgfpathlineto{\pgfqpoint{3.226106in}{1.589375in}}%
\pgfpathlineto{\pgfqpoint{3.229082in}{1.598927in}}%
\pgfpathlineto{\pgfqpoint{3.230317in}{1.602623in}}%
\pgfpathlineto{\pgfqpoint{3.234074in}{1.615871in}}%
\pgfpathlineto{\pgfqpoint{3.237849in}{1.629119in}}%
\pgfpathlineto{\pgfqpoint{3.243992in}{1.642367in}}%
\pgfpathlineto{\pgfqpoint{3.246701in}{1.655615in}}%
\pgfpathlineto{\pgfqpoint{3.253834in}{1.668863in}}%
\pgfpathlineto{\pgfqpoint{3.254365in}{1.673271in}}%
\pgfpathlineto{\pgfqpoint{3.255611in}{1.682111in}}%
\pgfpathlineto{\pgfqpoint{3.266140in}{1.695359in}}%
\pgfpathlineto{\pgfqpoint{3.266703in}{1.708607in}}%
\pgfpathlineto{\pgfqpoint{3.276053in}{1.721855in}}%
\pgfpathlineto{\pgfqpoint{3.276277in}{1.735103in}}%
\pgfpathlineto{\pgfqpoint{3.279648in}{1.739495in}}%
\pgfpathlineto{\pgfqpoint{3.289422in}{1.748351in}}%
\pgfpathlineto{\pgfqpoint{3.292621in}{1.761599in}}%
\pgfpathlineto{\pgfqpoint{3.304931in}{1.773115in}}%
\pgfpathlineto{\pgfqpoint{3.308755in}{1.774847in}}%
\pgfpathlineto{\pgfqpoint{3.318963in}{1.788095in}}%
\pgfpathlineto{\pgfqpoint{3.330214in}{1.794519in}}%
\pgfpathlineto{\pgfqpoint{3.355497in}{1.788679in}}%
\pgfpathlineto{\pgfqpoint{3.355997in}{1.788095in}}%
\pgfpathlineto{\pgfqpoint{3.365844in}{1.774847in}}%
\pgfpathlineto{\pgfqpoint{3.380001in}{1.761599in}}%
\pgfpathlineto{\pgfqpoint{3.380781in}{1.760461in}}%
\pgfpathlineto{\pgfqpoint{3.386034in}{1.748351in}}%
\pgfpathlineto{\pgfqpoint{3.395666in}{1.735103in}}%
\pgfpathlineto{\pgfqpoint{3.399340in}{1.721855in}}%
\pgfpathlineto{\pgfqpoint{3.406064in}{1.711420in}}%
\pgfpathlineto{\pgfqpoint{3.407409in}{1.708607in}}%
\pgfpathlineto{\pgfqpoint{3.408844in}{1.695359in}}%
\pgfpathlineto{\pgfqpoint{3.417194in}{1.682111in}}%
\pgfpathlineto{\pgfqpoint{3.418610in}{1.668863in}}%
\pgfpathlineto{\pgfqpoint{3.427546in}{1.655615in}}%
\pgfpathlineto{\pgfqpoint{3.428115in}{1.642367in}}%
\pgfpathlineto{\pgfqpoint{3.431347in}{1.637513in}}%
\pgfpathlineto{\pgfqpoint{3.436105in}{1.629119in}}%
\pgfpathlineto{\pgfqpoint{3.436615in}{1.615871in}}%
\pgfpathlineto{\pgfqpoint{3.444799in}{1.602623in}}%
\pgfpathlineto{\pgfqpoint{3.444833in}{1.589375in}}%
\pgfpathlineto{\pgfqpoint{3.452522in}{1.576127in}}%
\pgfpathlineto{\pgfqpoint{3.452986in}{1.562879in}}%
\pgfpathlineto{\pgfqpoint{3.456630in}{1.556221in}}%
\pgfpathlineto{\pgfqpoint{3.460007in}{1.549631in}}%
\pgfpathlineto{\pgfqpoint{3.460870in}{1.536383in}}%
\pgfpathlineto{\pgfqpoint{3.466070in}{1.523135in}}%
\pgfpathlineto{\pgfqpoint{3.466567in}{1.509887in}}%
\pgfpathlineto{\pgfqpoint{3.471321in}{1.496639in}}%
\pgfpathlineto{\pgfqpoint{3.474343in}{1.483391in}}%
\pgfpathlineto{\pgfqpoint{3.478397in}{1.470143in}}%
\pgfpathlineto{\pgfqpoint{3.481913in}{1.457364in}}%
\pgfpathlineto{\pgfqpoint{3.482036in}{1.456895in}}%
\pgfpathlineto{\pgfqpoint{3.484202in}{1.443647in}}%
\pgfpathlineto{\pgfqpoint{3.488215in}{1.430399in}}%
\pgfpathlineto{\pgfqpoint{3.489387in}{1.417151in}}%
\pgfpathlineto{\pgfqpoint{3.494865in}{1.403903in}}%
\pgfpathlineto{\pgfqpoint{3.493943in}{1.390655in}}%
\pgfpathlineto{\pgfqpoint{3.500139in}{1.377407in}}%
\pgfpathlineto{\pgfqpoint{3.498606in}{1.364159in}}%
\pgfpathlineto{\pgfqpoint{3.505518in}{1.350911in}}%
\pgfpathlineto{\pgfqpoint{3.502844in}{1.337663in}}%
\pgfpathlineto{\pgfqpoint{3.507196in}{1.328973in}}%
\pgfpathlineto{\pgfqpoint{3.509414in}{1.324415in}}%
\pgfpathlineto{\pgfqpoint{3.507196in}{1.315745in}}%
\pgfpathlineto{\pgfqpoint{3.505926in}{1.311167in}}%
\pgfpathlineto{\pgfqpoint{3.507196in}{1.309008in}}%
\pgfpathlineto{\pgfqpoint{3.513572in}{1.297919in}}%
\pgfpathlineto{\pgfqpoint{3.511854in}{1.284671in}}%
\pgfpathlineto{\pgfqpoint{3.517776in}{1.271423in}}%
\pgfpathlineto{\pgfqpoint{3.514644in}{1.258175in}}%
\pgfpathlineto{\pgfqpoint{3.518562in}{1.244927in}}%
\pgfpathlineto{\pgfqpoint{3.515972in}{1.231679in}}%
\pgfpathlineto{\pgfqpoint{3.518264in}{1.218431in}}%
\pgfpathlineto{\pgfqpoint{3.517547in}{1.205183in}}%
\pgfpathlineto{\pgfqpoint{3.516006in}{1.191935in}}%
\pgfpathlineto{\pgfqpoint{3.517290in}{1.178687in}}%
\pgfpathlineto{\pgfqpoint{3.513880in}{1.165439in}}%
\pgfpathlineto{\pgfqpoint{3.516477in}{1.152191in}}%
\pgfpathlineto{\pgfqpoint{3.507363in}{1.138943in}}%
\pgfpathlineto{\pgfqpoint{3.511271in}{1.125695in}}%
\pgfpathlineto{\pgfqpoint{3.507196in}{1.121247in}}%
\pgfpathlineto{\pgfqpoint{3.497553in}{1.112447in}}%
\pgfpathlineto{\pgfqpoint{3.504318in}{1.099199in}}%
\pgfpathlineto{\pgfqpoint{3.486154in}{1.085951in}}%
\pgfpathlineto{\pgfqpoint{3.487248in}{1.072703in}}%
\pgfpathlineto{\pgfqpoint{3.481913in}{1.069903in}}%
\pgfpathlineto{\pgfqpoint{3.457553in}{1.059455in}}%
\pgfpathlineto{\pgfqpoint{3.457490in}{1.046207in}}%
\pgfpathlineto{\pgfqpoint{3.456630in}{1.045816in}}%
\pgfpathlineto{\pgfqpoint{3.431347in}{1.036398in}}%
\pgfpathlineto{\pgfqpoint{3.419582in}{1.032959in}}%
\pgfpathlineto{\pgfqpoint{3.406064in}{1.022389in}}%
\pgfpathlineto{\pgfqpoint{3.400240in}{1.019711in}}%
\pgfpathlineto{\pgfqpoint{3.380781in}{1.015507in}}%
\pgfpathlineto{\pgfqpoint{3.355497in}{1.012119in}}%
\pgfpathlineto{\pgfqpoint{3.330214in}{1.010947in}}%
\pgfpathlineto{\pgfqpoint{3.304931in}{1.012470in}}%
\pgfpathlineto{\pgfqpoint{3.279648in}{1.017689in}}%
\pgfpathlineto{\pgfqpoint{3.274310in}{1.019711in}}%
\pgfpathclose%
\pgfusepath{fill}%
\end{pgfscope}%
\begin{pgfscope}%
\pgfpathrectangle{\pgfqpoint{2.470586in}{0.476543in}}{\pgfqpoint{1.820388in}{1.656000in}}%
\pgfusepath{clip}%
\pgfsetbuttcap%
\pgfsetroundjoin%
\definecolor{currentfill}{rgb}{0.994233,0.786159,0.477970}%
\pgfsetfillcolor{currentfill}%
\pgfsetlinewidth{0.000000pt}%
\definecolor{currentstroke}{rgb}{0.000000,0.000000,0.000000}%
\pgfsetstrokecolor{currentstroke}%
\pgfsetdash{}{0pt}%
\pgfpathmoveto{\pgfqpoint{3.279648in}{1.017689in}}%
\pgfpathlineto{\pgfqpoint{3.304931in}{1.012470in}}%
\pgfpathlineto{\pgfqpoint{3.330214in}{1.010947in}}%
\pgfpathlineto{\pgfqpoint{3.355497in}{1.012119in}}%
\pgfpathlineto{\pgfqpoint{3.380781in}{1.015507in}}%
\pgfpathlineto{\pgfqpoint{3.400240in}{1.019711in}}%
\pgfpathlineto{\pgfqpoint{3.406064in}{1.022389in}}%
\pgfpathlineto{\pgfqpoint{3.419582in}{1.032959in}}%
\pgfpathlineto{\pgfqpoint{3.431347in}{1.036398in}}%
\pgfpathlineto{\pgfqpoint{3.456630in}{1.045816in}}%
\pgfpathlineto{\pgfqpoint{3.457490in}{1.046207in}}%
\pgfpathlineto{\pgfqpoint{3.457553in}{1.059455in}}%
\pgfpathlineto{\pgfqpoint{3.481913in}{1.069903in}}%
\pgfpathlineto{\pgfqpoint{3.487248in}{1.072703in}}%
\pgfpathlineto{\pgfqpoint{3.486154in}{1.085951in}}%
\pgfpathlineto{\pgfqpoint{3.504318in}{1.099199in}}%
\pgfpathlineto{\pgfqpoint{3.497553in}{1.112447in}}%
\pgfpathlineto{\pgfqpoint{3.507196in}{1.121247in}}%
\pgfpathlineto{\pgfqpoint{3.511271in}{1.125695in}}%
\pgfpathlineto{\pgfqpoint{3.507363in}{1.138943in}}%
\pgfpathlineto{\pgfqpoint{3.516477in}{1.152191in}}%
\pgfpathlineto{\pgfqpoint{3.513880in}{1.165439in}}%
\pgfpathlineto{\pgfqpoint{3.517290in}{1.178687in}}%
\pgfpathlineto{\pgfqpoint{3.516006in}{1.191935in}}%
\pgfpathlineto{\pgfqpoint{3.517547in}{1.205183in}}%
\pgfpathlineto{\pgfqpoint{3.518264in}{1.218431in}}%
\pgfpathlineto{\pgfqpoint{3.515972in}{1.231679in}}%
\pgfpathlineto{\pgfqpoint{3.518562in}{1.244927in}}%
\pgfpathlineto{\pgfqpoint{3.514644in}{1.258175in}}%
\pgfpathlineto{\pgfqpoint{3.517776in}{1.271423in}}%
\pgfpathlineto{\pgfqpoint{3.511854in}{1.284671in}}%
\pgfpathlineto{\pgfqpoint{3.513572in}{1.297919in}}%
\pgfpathlineto{\pgfqpoint{3.507196in}{1.309008in}}%
\pgfpathlineto{\pgfqpoint{3.505926in}{1.311167in}}%
\pgfpathlineto{\pgfqpoint{3.507196in}{1.315745in}}%
\pgfpathlineto{\pgfqpoint{3.509414in}{1.324415in}}%
\pgfpathlineto{\pgfqpoint{3.507196in}{1.328973in}}%
\pgfpathlineto{\pgfqpoint{3.502844in}{1.337663in}}%
\pgfpathlineto{\pgfqpoint{3.505518in}{1.350911in}}%
\pgfpathlineto{\pgfqpoint{3.498606in}{1.364159in}}%
\pgfpathlineto{\pgfqpoint{3.500139in}{1.377407in}}%
\pgfpathlineto{\pgfqpoint{3.493943in}{1.390655in}}%
\pgfpathlineto{\pgfqpoint{3.494865in}{1.403903in}}%
\pgfpathlineto{\pgfqpoint{3.489387in}{1.417151in}}%
\pgfpathlineto{\pgfqpoint{3.488215in}{1.430399in}}%
\pgfpathlineto{\pgfqpoint{3.484202in}{1.443647in}}%
\pgfpathlineto{\pgfqpoint{3.482036in}{1.456895in}}%
\pgfpathlineto{\pgfqpoint{3.481913in}{1.457364in}}%
\pgfpathlineto{\pgfqpoint{3.478397in}{1.470143in}}%
\pgfpathlineto{\pgfqpoint{3.474343in}{1.483391in}}%
\pgfpathlineto{\pgfqpoint{3.471321in}{1.496639in}}%
\pgfpathlineto{\pgfqpoint{3.466567in}{1.509887in}}%
\pgfpathlineto{\pgfqpoint{3.466070in}{1.523135in}}%
\pgfpathlineto{\pgfqpoint{3.460870in}{1.536383in}}%
\pgfpathlineto{\pgfqpoint{3.460007in}{1.549631in}}%
\pgfpathlineto{\pgfqpoint{3.456630in}{1.556221in}}%
\pgfpathlineto{\pgfqpoint{3.452986in}{1.562879in}}%
\pgfpathlineto{\pgfqpoint{3.452522in}{1.576127in}}%
\pgfpathlineto{\pgfqpoint{3.444833in}{1.589375in}}%
\pgfpathlineto{\pgfqpoint{3.444799in}{1.602623in}}%
\pgfpathlineto{\pgfqpoint{3.436615in}{1.615871in}}%
\pgfpathlineto{\pgfqpoint{3.436105in}{1.629119in}}%
\pgfpathlineto{\pgfqpoint{3.431347in}{1.637513in}}%
\pgfpathlineto{\pgfqpoint{3.428115in}{1.642367in}}%
\pgfpathlineto{\pgfqpoint{3.427546in}{1.655615in}}%
\pgfpathlineto{\pgfqpoint{3.418610in}{1.668863in}}%
\pgfpathlineto{\pgfqpoint{3.417194in}{1.682111in}}%
\pgfpathlineto{\pgfqpoint{3.408844in}{1.695359in}}%
\pgfpathlineto{\pgfqpoint{3.407409in}{1.708607in}}%
\pgfpathlineto{\pgfqpoint{3.406064in}{1.711420in}}%
\pgfpathlineto{\pgfqpoint{3.399340in}{1.721855in}}%
\pgfpathlineto{\pgfqpoint{3.395666in}{1.735103in}}%
\pgfpathlineto{\pgfqpoint{3.386034in}{1.748351in}}%
\pgfpathlineto{\pgfqpoint{3.380781in}{1.760461in}}%
\pgfpathlineto{\pgfqpoint{3.380001in}{1.761599in}}%
\pgfpathlineto{\pgfqpoint{3.365844in}{1.774847in}}%
\pgfpathlineto{\pgfqpoint{3.355997in}{1.788095in}}%
\pgfpathlineto{\pgfqpoint{3.355497in}{1.788679in}}%
\pgfpathlineto{\pgfqpoint{3.330214in}{1.794519in}}%
\pgfpathlineto{\pgfqpoint{3.318963in}{1.788095in}}%
\pgfpathlineto{\pgfqpoint{3.308755in}{1.774847in}}%
\pgfpathlineto{\pgfqpoint{3.304931in}{1.773115in}}%
\pgfpathlineto{\pgfqpoint{3.292621in}{1.761599in}}%
\pgfpathlineto{\pgfqpoint{3.289422in}{1.748351in}}%
\pgfpathlineto{\pgfqpoint{3.279648in}{1.739495in}}%
\pgfpathlineto{\pgfqpoint{3.276277in}{1.735103in}}%
\pgfpathlineto{\pgfqpoint{3.276053in}{1.721855in}}%
\pgfpathlineto{\pgfqpoint{3.266703in}{1.708607in}}%
\pgfpathlineto{\pgfqpoint{3.266140in}{1.695359in}}%
\pgfpathlineto{\pgfqpoint{3.255611in}{1.682111in}}%
\pgfpathlineto{\pgfqpoint{3.254365in}{1.673271in}}%
\pgfpathlineto{\pgfqpoint{3.253834in}{1.668863in}}%
\pgfpathlineto{\pgfqpoint{3.246701in}{1.655615in}}%
\pgfpathlineto{\pgfqpoint{3.243992in}{1.642367in}}%
\pgfpathlineto{\pgfqpoint{3.237849in}{1.629119in}}%
\pgfpathlineto{\pgfqpoint{3.234074in}{1.615871in}}%
\pgfpathlineto{\pgfqpoint{3.230317in}{1.602623in}}%
\pgfpathlineto{\pgfqpoint{3.229082in}{1.598927in}}%
\pgfpathlineto{\pgfqpoint{3.226106in}{1.589375in}}%
\pgfpathlineto{\pgfqpoint{3.223839in}{1.576127in}}%
\pgfpathlineto{\pgfqpoint{3.217552in}{1.562879in}}%
\pgfpathlineto{\pgfqpoint{3.216425in}{1.549631in}}%
\pgfpathlineto{\pgfqpoint{3.210150in}{1.536383in}}%
\pgfpathlineto{\pgfqpoint{3.210506in}{1.523135in}}%
\pgfpathlineto{\pgfqpoint{3.203798in}{1.511510in}}%
\pgfpathlineto{\pgfqpoint{3.202905in}{1.509887in}}%
\pgfpathlineto{\pgfqpoint{3.203098in}{1.496639in}}%
\pgfpathlineto{\pgfqpoint{3.195019in}{1.483391in}}%
\pgfpathlineto{\pgfqpoint{3.195913in}{1.470143in}}%
\pgfpathlineto{\pgfqpoint{3.188906in}{1.456895in}}%
\pgfpathlineto{\pgfqpoint{3.189614in}{1.443647in}}%
\pgfpathlineto{\pgfqpoint{3.183101in}{1.430399in}}%
\pgfpathlineto{\pgfqpoint{3.183846in}{1.417151in}}%
\pgfpathlineto{\pgfqpoint{3.178922in}{1.403903in}}%
\pgfpathlineto{\pgfqpoint{3.178515in}{1.395970in}}%
\pgfpathlineto{\pgfqpoint{3.178238in}{1.390655in}}%
\pgfpathlineto{\pgfqpoint{3.173594in}{1.377407in}}%
\pgfpathlineto{\pgfqpoint{3.170800in}{1.364159in}}%
\pgfpathlineto{\pgfqpoint{3.168530in}{1.350911in}}%
\pgfpathlineto{\pgfqpoint{3.165836in}{1.337663in}}%
\pgfpathlineto{\pgfqpoint{3.165234in}{1.324415in}}%
\pgfpathlineto{\pgfqpoint{3.160054in}{1.311167in}}%
\pgfpathlineto{\pgfqpoint{3.161049in}{1.297919in}}%
\pgfpathlineto{\pgfqpoint{3.155111in}{1.284671in}}%
\pgfpathlineto{\pgfqpoint{3.159159in}{1.271423in}}%
\pgfpathlineto{\pgfqpoint{3.153232in}{1.259270in}}%
\pgfpathlineto{\pgfqpoint{3.152698in}{1.258175in}}%
\pgfpathlineto{\pgfqpoint{3.153232in}{1.256691in}}%
\pgfpathlineto{\pgfqpoint{3.157849in}{1.244927in}}%
\pgfpathlineto{\pgfqpoint{3.153232in}{1.236786in}}%
\pgfpathlineto{\pgfqpoint{3.150394in}{1.231679in}}%
\pgfpathlineto{\pgfqpoint{3.153232in}{1.225935in}}%
\pgfpathlineto{\pgfqpoint{3.157372in}{1.218431in}}%
\pgfpathlineto{\pgfqpoint{3.153232in}{1.211460in}}%
\pgfpathlineto{\pgfqpoint{3.149607in}{1.205183in}}%
\pgfpathlineto{\pgfqpoint{3.153232in}{1.197330in}}%
\pgfpathlineto{\pgfqpoint{3.156076in}{1.191935in}}%
\pgfpathlineto{\pgfqpoint{3.153232in}{1.187479in}}%
\pgfpathlineto{\pgfqpoint{3.147881in}{1.178687in}}%
\pgfpathlineto{\pgfqpoint{3.153232in}{1.169342in}}%
\pgfpathlineto{\pgfqpoint{3.155818in}{1.165439in}}%
\pgfpathlineto{\pgfqpoint{3.153232in}{1.157600in}}%
\pgfpathlineto{\pgfqpoint{3.151554in}{1.152191in}}%
\pgfpathlineto{\pgfqpoint{3.153232in}{1.149855in}}%
\pgfpathlineto{\pgfqpoint{3.162540in}{1.138943in}}%
\pgfpathlineto{\pgfqpoint{3.159104in}{1.125695in}}%
\pgfpathlineto{\pgfqpoint{3.168488in}{1.112447in}}%
\pgfpathlineto{\pgfqpoint{3.168721in}{1.099199in}}%
\pgfpathlineto{\pgfqpoint{3.178515in}{1.087555in}}%
\pgfpathlineto{\pgfqpoint{3.180226in}{1.085951in}}%
\pgfpathlineto{\pgfqpoint{3.187843in}{1.072703in}}%
\pgfpathlineto{\pgfqpoint{3.202463in}{1.059455in}}%
\pgfpathlineto{\pgfqpoint{3.203798in}{1.058037in}}%
\pgfpathlineto{\pgfqpoint{3.216894in}{1.046207in}}%
\pgfpathlineto{\pgfqpoint{3.229082in}{1.036543in}}%
\pgfpathlineto{\pgfqpoint{3.235794in}{1.032959in}}%
\pgfpathlineto{\pgfqpoint{3.254365in}{1.025701in}}%
\pgfpathlineto{\pgfqpoint{3.274310in}{1.019711in}}%
\pgfpathlineto{\pgfqpoint{3.279648in}{1.017689in}}%
\pgfpathclose%
\pgfpathmoveto{\pgfqpoint{3.275688in}{1.046207in}}%
\pgfpathlineto{\pgfqpoint{3.254365in}{1.058534in}}%
\pgfpathlineto{\pgfqpoint{3.253138in}{1.059455in}}%
\pgfpathlineto{\pgfqpoint{3.229082in}{1.069647in}}%
\pgfpathlineto{\pgfqpoint{3.224210in}{1.072703in}}%
\pgfpathlineto{\pgfqpoint{3.214788in}{1.085951in}}%
\pgfpathlineto{\pgfqpoint{3.203798in}{1.093049in}}%
\pgfpathlineto{\pgfqpoint{3.196632in}{1.099199in}}%
\pgfpathlineto{\pgfqpoint{3.196134in}{1.112447in}}%
\pgfpathlineto{\pgfqpoint{3.182655in}{1.125695in}}%
\pgfpathlineto{\pgfqpoint{3.186113in}{1.138943in}}%
\pgfpathlineto{\pgfqpoint{3.178515in}{1.145439in}}%
\pgfpathlineto{\pgfqpoint{3.172094in}{1.152191in}}%
\pgfpathlineto{\pgfqpoint{3.176219in}{1.165439in}}%
\pgfpathlineto{\pgfqpoint{3.166414in}{1.178687in}}%
\pgfpathlineto{\pgfqpoint{3.174558in}{1.191935in}}%
\pgfpathlineto{\pgfqpoint{3.166611in}{1.205183in}}%
\pgfpathlineto{\pgfqpoint{3.174013in}{1.218431in}}%
\pgfpathlineto{\pgfqpoint{3.166001in}{1.231679in}}%
\pgfpathlineto{\pgfqpoint{3.173053in}{1.244927in}}%
\pgfpathlineto{\pgfqpoint{3.167301in}{1.258175in}}%
\pgfpathlineto{\pgfqpoint{3.173097in}{1.271423in}}%
\pgfpathlineto{\pgfqpoint{3.168675in}{1.284671in}}%
\pgfpathlineto{\pgfqpoint{3.174010in}{1.297919in}}%
\pgfpathlineto{\pgfqpoint{3.172882in}{1.311167in}}%
\pgfpathlineto{\pgfqpoint{3.177424in}{1.324415in}}%
\pgfpathlineto{\pgfqpoint{3.177967in}{1.337663in}}%
\pgfpathlineto{\pgfqpoint{3.178515in}{1.341070in}}%
\pgfpathlineto{\pgfqpoint{3.180099in}{1.350911in}}%
\pgfpathlineto{\pgfqpoint{3.182581in}{1.364159in}}%
\pgfpathlineto{\pgfqpoint{3.184859in}{1.377407in}}%
\pgfpathlineto{\pgfqpoint{3.189792in}{1.390655in}}%
\pgfpathlineto{\pgfqpoint{3.189782in}{1.403903in}}%
\pgfpathlineto{\pgfqpoint{3.194772in}{1.417151in}}%
\pgfpathlineto{\pgfqpoint{3.193487in}{1.430399in}}%
\pgfpathlineto{\pgfqpoint{3.200105in}{1.443647in}}%
\pgfpathlineto{\pgfqpoint{3.198951in}{1.456895in}}%
\pgfpathlineto{\pgfqpoint{3.203798in}{1.466022in}}%
\pgfpathlineto{\pgfqpoint{3.206161in}{1.470143in}}%
\pgfpathlineto{\pgfqpoint{3.204848in}{1.483391in}}%
\pgfpathlineto{\pgfqpoint{3.213516in}{1.496639in}}%
\pgfpathlineto{\pgfqpoint{3.213099in}{1.509887in}}%
\pgfpathlineto{\pgfqpoint{3.220564in}{1.523135in}}%
\pgfpathlineto{\pgfqpoint{3.220078in}{1.536383in}}%
\pgfpathlineto{\pgfqpoint{3.226131in}{1.549631in}}%
\pgfpathlineto{\pgfqpoint{3.227286in}{1.562879in}}%
\pgfpathlineto{\pgfqpoint{3.229082in}{1.566753in}}%
\pgfpathlineto{\pgfqpoint{3.233796in}{1.576127in}}%
\pgfpathlineto{\pgfqpoint{3.236459in}{1.589375in}}%
\pgfpathlineto{\pgfqpoint{3.240566in}{1.602623in}}%
\pgfpathlineto{\pgfqpoint{3.244616in}{1.615871in}}%
\pgfpathlineto{\pgfqpoint{3.247995in}{1.629119in}}%
\pgfpathlineto{\pgfqpoint{3.254365in}{1.642235in}}%
\pgfpathlineto{\pgfqpoint{3.254446in}{1.642367in}}%
\pgfpathlineto{\pgfqpoint{3.257226in}{1.655615in}}%
\pgfpathlineto{\pgfqpoint{3.266256in}{1.668863in}}%
\pgfpathlineto{\pgfqpoint{3.267661in}{1.682111in}}%
\pgfpathlineto{\pgfqpoint{3.278468in}{1.695359in}}%
\pgfpathlineto{\pgfqpoint{3.278627in}{1.708607in}}%
\pgfpathlineto{\pgfqpoint{3.279648in}{1.710049in}}%
\pgfpathlineto{\pgfqpoint{3.291650in}{1.721855in}}%
\pgfpathlineto{\pgfqpoint{3.291885in}{1.735103in}}%
\pgfpathlineto{\pgfqpoint{3.304931in}{1.746948in}}%
\pgfpathlineto{\pgfqpoint{3.308255in}{1.748351in}}%
\pgfpathlineto{\pgfqpoint{3.316475in}{1.761599in}}%
\pgfpathlineto{\pgfqpoint{3.330214in}{1.767735in}}%
\pgfpathlineto{\pgfqpoint{3.355497in}{1.763343in}}%
\pgfpathlineto{\pgfqpoint{3.357467in}{1.761599in}}%
\pgfpathlineto{\pgfqpoint{3.365435in}{1.748351in}}%
\pgfpathlineto{\pgfqpoint{3.380781in}{1.736655in}}%
\pgfpathlineto{\pgfqpoint{3.381932in}{1.735103in}}%
\pgfpathlineto{\pgfqpoint{3.384860in}{1.721855in}}%
\pgfpathlineto{\pgfqpoint{3.393822in}{1.708607in}}%
\pgfpathlineto{\pgfqpoint{3.395053in}{1.695359in}}%
\pgfpathlineto{\pgfqpoint{3.406064in}{1.682631in}}%
\pgfpathlineto{\pgfqpoint{3.406395in}{1.682111in}}%
\pgfpathlineto{\pgfqpoint{3.407434in}{1.668863in}}%
\pgfpathlineto{\pgfqpoint{3.416555in}{1.655615in}}%
\pgfpathlineto{\pgfqpoint{3.416820in}{1.642367in}}%
\pgfpathlineto{\pgfqpoint{3.425713in}{1.629119in}}%
\pgfpathlineto{\pgfqpoint{3.426085in}{1.615871in}}%
\pgfpathlineto{\pgfqpoint{3.431347in}{1.608603in}}%
\pgfpathlineto{\pgfqpoint{3.435059in}{1.602623in}}%
\pgfpathlineto{\pgfqpoint{3.434886in}{1.589375in}}%
\pgfpathlineto{\pgfqpoint{3.442466in}{1.576127in}}%
\pgfpathlineto{\pgfqpoint{3.442848in}{1.562879in}}%
\pgfpathlineto{\pgfqpoint{3.449963in}{1.549631in}}%
\pgfpathlineto{\pgfqpoint{3.450828in}{1.536383in}}%
\pgfpathlineto{\pgfqpoint{3.456243in}{1.523135in}}%
\pgfpathlineto{\pgfqpoint{3.456630in}{1.513513in}}%
\pgfpathlineto{\pgfqpoint{3.456762in}{1.509887in}}%
\pgfpathlineto{\pgfqpoint{3.461214in}{1.496639in}}%
\pgfpathlineto{\pgfqpoint{3.464362in}{1.483391in}}%
\pgfpathlineto{\pgfqpoint{3.467979in}{1.470143in}}%
\pgfpathlineto{\pgfqpoint{3.471685in}{1.456895in}}%
\pgfpathlineto{\pgfqpoint{3.473464in}{1.443647in}}%
\pgfpathlineto{\pgfqpoint{3.477813in}{1.430399in}}%
\pgfpathlineto{\pgfqpoint{3.478498in}{1.417151in}}%
\pgfpathlineto{\pgfqpoint{3.481913in}{1.409496in}}%
\pgfpathlineto{\pgfqpoint{3.484245in}{1.403903in}}%
\pgfpathlineto{\pgfqpoint{3.482720in}{1.390655in}}%
\pgfpathlineto{\pgfqpoint{3.488923in}{1.377407in}}%
\pgfpathlineto{\pgfqpoint{3.486863in}{1.364159in}}%
\pgfpathlineto{\pgfqpoint{3.493722in}{1.350911in}}%
\pgfpathlineto{\pgfqpoint{3.490461in}{1.337663in}}%
\pgfpathlineto{\pgfqpoint{3.496941in}{1.324415in}}%
\pgfpathlineto{\pgfqpoint{3.492877in}{1.311167in}}%
\pgfpathlineto{\pgfqpoint{3.500474in}{1.297919in}}%
\pgfpathlineto{\pgfqpoint{3.498263in}{1.284671in}}%
\pgfpathlineto{\pgfqpoint{3.503883in}{1.271423in}}%
\pgfpathlineto{\pgfqpoint{3.500101in}{1.258175in}}%
\pgfpathlineto{\pgfqpoint{3.503456in}{1.244927in}}%
\pgfpathlineto{\pgfqpoint{3.500417in}{1.231679in}}%
\pgfpathlineto{\pgfqpoint{3.501650in}{1.218431in}}%
\pgfpathlineto{\pgfqpoint{3.500635in}{1.205183in}}%
\pgfpathlineto{\pgfqpoint{3.497491in}{1.191935in}}%
\pgfpathlineto{\pgfqpoint{3.498734in}{1.178687in}}%
\pgfpathlineto{\pgfqpoint{3.493355in}{1.165439in}}%
\pgfpathlineto{\pgfqpoint{3.495916in}{1.152191in}}%
\pgfpathlineto{\pgfqpoint{3.483734in}{1.138943in}}%
\pgfpathlineto{\pgfqpoint{3.487574in}{1.125695in}}%
\pgfpathlineto{\pgfqpoint{3.481913in}{1.120837in}}%
\pgfpathlineto{\pgfqpoint{3.469812in}{1.112447in}}%
\pgfpathlineto{\pgfqpoint{3.476209in}{1.099199in}}%
\pgfpathlineto{\pgfqpoint{3.456630in}{1.088023in}}%
\pgfpathlineto{\pgfqpoint{3.451931in}{1.085951in}}%
\pgfpathlineto{\pgfqpoint{3.450667in}{1.072703in}}%
\pgfpathlineto{\pgfqpoint{3.431347in}{1.065949in}}%
\pgfpathlineto{\pgfqpoint{3.406821in}{1.059455in}}%
\pgfpathlineto{\pgfqpoint{3.406064in}{1.057766in}}%
\pgfpathlineto{\pgfqpoint{3.397763in}{1.046207in}}%
\pgfpathlineto{\pgfqpoint{3.380781in}{1.042656in}}%
\pgfpathlineto{\pgfqpoint{3.355497in}{1.039445in}}%
\pgfpathlineto{\pgfqpoint{3.330214in}{1.038406in}}%
\pgfpathlineto{\pgfqpoint{3.304931in}{1.039934in}}%
\pgfpathlineto{\pgfqpoint{3.279648in}{1.044842in}}%
\pgfpathlineto{\pgfqpoint{3.275688in}{1.046207in}}%
\pgfpathclose%
\pgfusepath{fill}%
\end{pgfscope}%
\begin{pgfscope}%
\pgfpathrectangle{\pgfqpoint{2.470586in}{0.476543in}}{\pgfqpoint{1.820388in}{1.656000in}}%
\pgfusepath{clip}%
\pgfsetbuttcap%
\pgfsetroundjoin%
\definecolor{currentfill}{rgb}{0.992541,0.701576,0.398462}%
\pgfsetfillcolor{currentfill}%
\pgfsetlinewidth{0.000000pt}%
\definecolor{currentstroke}{rgb}{0.000000,0.000000,0.000000}%
\pgfsetstrokecolor{currentstroke}%
\pgfsetdash{}{0pt}%
\pgfpathmoveto{\pgfqpoint{3.279648in}{1.044842in}}%
\pgfpathlineto{\pgfqpoint{3.304931in}{1.039934in}}%
\pgfpathlineto{\pgfqpoint{3.330214in}{1.038406in}}%
\pgfpathlineto{\pgfqpoint{3.355497in}{1.039445in}}%
\pgfpathlineto{\pgfqpoint{3.380781in}{1.042656in}}%
\pgfpathlineto{\pgfqpoint{3.397763in}{1.046207in}}%
\pgfpathlineto{\pgfqpoint{3.406064in}{1.057766in}}%
\pgfpathlineto{\pgfqpoint{3.406821in}{1.059455in}}%
\pgfpathlineto{\pgfqpoint{3.431347in}{1.065949in}}%
\pgfpathlineto{\pgfqpoint{3.450667in}{1.072703in}}%
\pgfpathlineto{\pgfqpoint{3.451931in}{1.085951in}}%
\pgfpathlineto{\pgfqpoint{3.456630in}{1.088023in}}%
\pgfpathlineto{\pgfqpoint{3.476209in}{1.099199in}}%
\pgfpathlineto{\pgfqpoint{3.469812in}{1.112447in}}%
\pgfpathlineto{\pgfqpoint{3.481913in}{1.120837in}}%
\pgfpathlineto{\pgfqpoint{3.487574in}{1.125695in}}%
\pgfpathlineto{\pgfqpoint{3.483734in}{1.138943in}}%
\pgfpathlineto{\pgfqpoint{3.495916in}{1.152191in}}%
\pgfpathlineto{\pgfqpoint{3.493355in}{1.165439in}}%
\pgfpathlineto{\pgfqpoint{3.498734in}{1.178687in}}%
\pgfpathlineto{\pgfqpoint{3.497491in}{1.191935in}}%
\pgfpathlineto{\pgfqpoint{3.500635in}{1.205183in}}%
\pgfpathlineto{\pgfqpoint{3.501650in}{1.218431in}}%
\pgfpathlineto{\pgfqpoint{3.500417in}{1.231679in}}%
\pgfpathlineto{\pgfqpoint{3.503456in}{1.244927in}}%
\pgfpathlineto{\pgfqpoint{3.500101in}{1.258175in}}%
\pgfpathlineto{\pgfqpoint{3.503883in}{1.271423in}}%
\pgfpathlineto{\pgfqpoint{3.498263in}{1.284671in}}%
\pgfpathlineto{\pgfqpoint{3.500474in}{1.297919in}}%
\pgfpathlineto{\pgfqpoint{3.492877in}{1.311167in}}%
\pgfpathlineto{\pgfqpoint{3.496941in}{1.324415in}}%
\pgfpathlineto{\pgfqpoint{3.490461in}{1.337663in}}%
\pgfpathlineto{\pgfqpoint{3.493722in}{1.350911in}}%
\pgfpathlineto{\pgfqpoint{3.486863in}{1.364159in}}%
\pgfpathlineto{\pgfqpoint{3.488923in}{1.377407in}}%
\pgfpathlineto{\pgfqpoint{3.482720in}{1.390655in}}%
\pgfpathlineto{\pgfqpoint{3.484245in}{1.403903in}}%
\pgfpathlineto{\pgfqpoint{3.481913in}{1.409496in}}%
\pgfpathlineto{\pgfqpoint{3.478498in}{1.417151in}}%
\pgfpathlineto{\pgfqpoint{3.477813in}{1.430399in}}%
\pgfpathlineto{\pgfqpoint{3.473464in}{1.443647in}}%
\pgfpathlineto{\pgfqpoint{3.471685in}{1.456895in}}%
\pgfpathlineto{\pgfqpoint{3.467979in}{1.470143in}}%
\pgfpathlineto{\pgfqpoint{3.464362in}{1.483391in}}%
\pgfpathlineto{\pgfqpoint{3.461214in}{1.496639in}}%
\pgfpathlineto{\pgfqpoint{3.456762in}{1.509887in}}%
\pgfpathlineto{\pgfqpoint{3.456630in}{1.513513in}}%
\pgfpathlineto{\pgfqpoint{3.456243in}{1.523135in}}%
\pgfpathlineto{\pgfqpoint{3.450828in}{1.536383in}}%
\pgfpathlineto{\pgfqpoint{3.449963in}{1.549631in}}%
\pgfpathlineto{\pgfqpoint{3.442848in}{1.562879in}}%
\pgfpathlineto{\pgfqpoint{3.442466in}{1.576127in}}%
\pgfpathlineto{\pgfqpoint{3.434886in}{1.589375in}}%
\pgfpathlineto{\pgfqpoint{3.435059in}{1.602623in}}%
\pgfpathlineto{\pgfqpoint{3.431347in}{1.608603in}}%
\pgfpathlineto{\pgfqpoint{3.426085in}{1.615871in}}%
\pgfpathlineto{\pgfqpoint{3.425713in}{1.629119in}}%
\pgfpathlineto{\pgfqpoint{3.416820in}{1.642367in}}%
\pgfpathlineto{\pgfqpoint{3.416555in}{1.655615in}}%
\pgfpathlineto{\pgfqpoint{3.407434in}{1.668863in}}%
\pgfpathlineto{\pgfqpoint{3.406395in}{1.682111in}}%
\pgfpathlineto{\pgfqpoint{3.406064in}{1.682631in}}%
\pgfpathlineto{\pgfqpoint{3.395053in}{1.695359in}}%
\pgfpathlineto{\pgfqpoint{3.393822in}{1.708607in}}%
\pgfpathlineto{\pgfqpoint{3.384860in}{1.721855in}}%
\pgfpathlineto{\pgfqpoint{3.381932in}{1.735103in}}%
\pgfpathlineto{\pgfqpoint{3.380781in}{1.736655in}}%
\pgfpathlineto{\pgfqpoint{3.365435in}{1.748351in}}%
\pgfpathlineto{\pgfqpoint{3.357467in}{1.761599in}}%
\pgfpathlineto{\pgfqpoint{3.355497in}{1.763343in}}%
\pgfpathlineto{\pgfqpoint{3.330214in}{1.767735in}}%
\pgfpathlineto{\pgfqpoint{3.316475in}{1.761599in}}%
\pgfpathlineto{\pgfqpoint{3.308255in}{1.748351in}}%
\pgfpathlineto{\pgfqpoint{3.304931in}{1.746948in}}%
\pgfpathlineto{\pgfqpoint{3.291885in}{1.735103in}}%
\pgfpathlineto{\pgfqpoint{3.291650in}{1.721855in}}%
\pgfpathlineto{\pgfqpoint{3.279648in}{1.710049in}}%
\pgfpathlineto{\pgfqpoint{3.278627in}{1.708607in}}%
\pgfpathlineto{\pgfqpoint{3.278468in}{1.695359in}}%
\pgfpathlineto{\pgfqpoint{3.267661in}{1.682111in}}%
\pgfpathlineto{\pgfqpoint{3.266256in}{1.668863in}}%
\pgfpathlineto{\pgfqpoint{3.257226in}{1.655615in}}%
\pgfpathlineto{\pgfqpoint{3.254446in}{1.642367in}}%
\pgfpathlineto{\pgfqpoint{3.254365in}{1.642235in}}%
\pgfpathlineto{\pgfqpoint{3.247995in}{1.629119in}}%
\pgfpathlineto{\pgfqpoint{3.244616in}{1.615871in}}%
\pgfpathlineto{\pgfqpoint{3.240566in}{1.602623in}}%
\pgfpathlineto{\pgfqpoint{3.236459in}{1.589375in}}%
\pgfpathlineto{\pgfqpoint{3.233796in}{1.576127in}}%
\pgfpathlineto{\pgfqpoint{3.229082in}{1.566753in}}%
\pgfpathlineto{\pgfqpoint{3.227286in}{1.562879in}}%
\pgfpathlineto{\pgfqpoint{3.226131in}{1.549631in}}%
\pgfpathlineto{\pgfqpoint{3.220078in}{1.536383in}}%
\pgfpathlineto{\pgfqpoint{3.220564in}{1.523135in}}%
\pgfpathlineto{\pgfqpoint{3.213099in}{1.509887in}}%
\pgfpathlineto{\pgfqpoint{3.213516in}{1.496639in}}%
\pgfpathlineto{\pgfqpoint{3.204848in}{1.483391in}}%
\pgfpathlineto{\pgfqpoint{3.206161in}{1.470143in}}%
\pgfpathlineto{\pgfqpoint{3.203798in}{1.466022in}}%
\pgfpathlineto{\pgfqpoint{3.198951in}{1.456895in}}%
\pgfpathlineto{\pgfqpoint{3.200105in}{1.443647in}}%
\pgfpathlineto{\pgfqpoint{3.193487in}{1.430399in}}%
\pgfpathlineto{\pgfqpoint{3.194772in}{1.417151in}}%
\pgfpathlineto{\pgfqpoint{3.189782in}{1.403903in}}%
\pgfpathlineto{\pgfqpoint{3.189792in}{1.390655in}}%
\pgfpathlineto{\pgfqpoint{3.184859in}{1.377407in}}%
\pgfpathlineto{\pgfqpoint{3.182581in}{1.364159in}}%
\pgfpathlineto{\pgfqpoint{3.180099in}{1.350911in}}%
\pgfpathlineto{\pgfqpoint{3.178515in}{1.341070in}}%
\pgfpathlineto{\pgfqpoint{3.177967in}{1.337663in}}%
\pgfpathlineto{\pgfqpoint{3.177424in}{1.324415in}}%
\pgfpathlineto{\pgfqpoint{3.172882in}{1.311167in}}%
\pgfpathlineto{\pgfqpoint{3.174010in}{1.297919in}}%
\pgfpathlineto{\pgfqpoint{3.168675in}{1.284671in}}%
\pgfpathlineto{\pgfqpoint{3.173097in}{1.271423in}}%
\pgfpathlineto{\pgfqpoint{3.167301in}{1.258175in}}%
\pgfpathlineto{\pgfqpoint{3.173053in}{1.244927in}}%
\pgfpathlineto{\pgfqpoint{3.166001in}{1.231679in}}%
\pgfpathlineto{\pgfqpoint{3.174013in}{1.218431in}}%
\pgfpathlineto{\pgfqpoint{3.166611in}{1.205183in}}%
\pgfpathlineto{\pgfqpoint{3.174558in}{1.191935in}}%
\pgfpathlineto{\pgfqpoint{3.166414in}{1.178687in}}%
\pgfpathlineto{\pgfqpoint{3.176219in}{1.165439in}}%
\pgfpathlineto{\pgfqpoint{3.172094in}{1.152191in}}%
\pgfpathlineto{\pgfqpoint{3.178515in}{1.145439in}}%
\pgfpathlineto{\pgfqpoint{3.186113in}{1.138943in}}%
\pgfpathlineto{\pgfqpoint{3.182655in}{1.125695in}}%
\pgfpathlineto{\pgfqpoint{3.196134in}{1.112447in}}%
\pgfpathlineto{\pgfqpoint{3.196632in}{1.099199in}}%
\pgfpathlineto{\pgfqpoint{3.203798in}{1.093049in}}%
\pgfpathlineto{\pgfqpoint{3.214788in}{1.085951in}}%
\pgfpathlineto{\pgfqpoint{3.224210in}{1.072703in}}%
\pgfpathlineto{\pgfqpoint{3.229082in}{1.069647in}}%
\pgfpathlineto{\pgfqpoint{3.253138in}{1.059455in}}%
\pgfpathlineto{\pgfqpoint{3.254365in}{1.058534in}}%
\pgfpathlineto{\pgfqpoint{3.275688in}{1.046207in}}%
\pgfpathlineto{\pgfqpoint{3.279648in}{1.044842in}}%
\pgfpathclose%
\pgfpathmoveto{\pgfqpoint{3.276485in}{1.072703in}}%
\pgfpathlineto{\pgfqpoint{3.261004in}{1.085951in}}%
\pgfpathlineto{\pgfqpoint{3.254365in}{1.087883in}}%
\pgfpathlineto{\pgfqpoint{3.229381in}{1.099199in}}%
\pgfpathlineto{\pgfqpoint{3.229082in}{1.103772in}}%
\pgfpathlineto{\pgfqpoint{3.228620in}{1.112447in}}%
\pgfpathlineto{\pgfqpoint{3.209237in}{1.125695in}}%
\pgfpathlineto{\pgfqpoint{3.212787in}{1.138943in}}%
\pgfpathlineto{\pgfqpoint{3.203798in}{1.144594in}}%
\pgfpathlineto{\pgfqpoint{3.194381in}{1.152191in}}%
\pgfpathlineto{\pgfqpoint{3.198609in}{1.165439in}}%
\pgfpathlineto{\pgfqpoint{3.186100in}{1.178687in}}%
\pgfpathlineto{\pgfqpoint{3.194524in}{1.191935in}}%
\pgfpathlineto{\pgfqpoint{3.184402in}{1.205183in}}%
\pgfpathlineto{\pgfqpoint{3.191783in}{1.218431in}}%
\pgfpathlineto{\pgfqpoint{3.182093in}{1.231679in}}%
\pgfpathlineto{\pgfqpoint{3.189090in}{1.244927in}}%
\pgfpathlineto{\pgfqpoint{3.182241in}{1.258175in}}%
\pgfpathlineto{\pgfqpoint{3.187696in}{1.271423in}}%
\pgfpathlineto{\pgfqpoint{3.182551in}{1.284671in}}%
\pgfpathlineto{\pgfqpoint{3.187587in}{1.297919in}}%
\pgfpathlineto{\pgfqpoint{3.186273in}{1.311167in}}%
\pgfpathlineto{\pgfqpoint{3.190337in}{1.324415in}}%
\pgfpathlineto{\pgfqpoint{3.190890in}{1.337663in}}%
\pgfpathlineto{\pgfqpoint{3.192241in}{1.350911in}}%
\pgfpathlineto{\pgfqpoint{3.194819in}{1.364159in}}%
\pgfpathlineto{\pgfqpoint{3.196366in}{1.377407in}}%
\pgfpathlineto{\pgfqpoint{3.201360in}{1.390655in}}%
\pgfpathlineto{\pgfqpoint{3.200642in}{1.403903in}}%
\pgfpathlineto{\pgfqpoint{3.203798in}{1.412192in}}%
\pgfpathlineto{\pgfqpoint{3.205876in}{1.417151in}}%
\pgfpathlineto{\pgfqpoint{3.203880in}{1.430399in}}%
\pgfpathlineto{\pgfqpoint{3.211173in}{1.443647in}}%
\pgfpathlineto{\pgfqpoint{3.209434in}{1.456895in}}%
\pgfpathlineto{\pgfqpoint{3.217001in}{1.470143in}}%
\pgfpathlineto{\pgfqpoint{3.215356in}{1.483391in}}%
\pgfpathlineto{\pgfqpoint{3.223980in}{1.496639in}}%
\pgfpathlineto{\pgfqpoint{3.223355in}{1.509887in}}%
\pgfpathlineto{\pgfqpoint{3.229082in}{1.520282in}}%
\pgfpathlineto{\pgfqpoint{3.230828in}{1.523135in}}%
\pgfpathlineto{\pgfqpoint{3.230135in}{1.536383in}}%
\pgfpathlineto{\pgfqpoint{3.236678in}{1.549631in}}%
\pgfpathlineto{\pgfqpoint{3.238072in}{1.562879in}}%
\pgfpathlineto{\pgfqpoint{3.244358in}{1.576127in}}%
\pgfpathlineto{\pgfqpoint{3.247174in}{1.589375in}}%
\pgfpathlineto{\pgfqpoint{3.250815in}{1.602623in}}%
\pgfpathlineto{\pgfqpoint{3.254365in}{1.613506in}}%
\pgfpathlineto{\pgfqpoint{3.255352in}{1.615871in}}%
\pgfpathlineto{\pgfqpoint{3.259004in}{1.629119in}}%
\pgfpathlineto{\pgfqpoint{3.267299in}{1.642367in}}%
\pgfpathlineto{\pgfqpoint{3.269429in}{1.655615in}}%
\pgfpathlineto{\pgfqpoint{3.278794in}{1.668863in}}%
\pgfpathlineto{\pgfqpoint{3.279648in}{1.681159in}}%
\pgfpathlineto{\pgfqpoint{3.279740in}{1.682111in}}%
\pgfpathlineto{\pgfqpoint{3.295701in}{1.695359in}}%
\pgfpathlineto{\pgfqpoint{3.295379in}{1.708607in}}%
\pgfpathlineto{\pgfqpoint{3.304931in}{1.718047in}}%
\pgfpathlineto{\pgfqpoint{3.313122in}{1.721855in}}%
\pgfpathlineto{\pgfqpoint{3.314026in}{1.735103in}}%
\pgfpathlineto{\pgfqpoint{3.330214in}{1.741880in}}%
\pgfpathlineto{\pgfqpoint{3.355497in}{1.738273in}}%
\pgfpathlineto{\pgfqpoint{3.359760in}{1.735103in}}%
\pgfpathlineto{\pgfqpoint{3.362882in}{1.721855in}}%
\pgfpathlineto{\pgfqpoint{3.379173in}{1.708607in}}%
\pgfpathlineto{\pgfqpoint{3.380116in}{1.695359in}}%
\pgfpathlineto{\pgfqpoint{3.380781in}{1.694944in}}%
\pgfpathlineto{\pgfqpoint{3.392354in}{1.682111in}}%
\pgfpathlineto{\pgfqpoint{3.393180in}{1.668863in}}%
\pgfpathlineto{\pgfqpoint{3.405403in}{1.655615in}}%
\pgfpathlineto{\pgfqpoint{3.405355in}{1.642367in}}%
\pgfpathlineto{\pgfqpoint{3.406064in}{1.641587in}}%
\pgfpathlineto{\pgfqpoint{3.414562in}{1.629119in}}%
\pgfpathlineto{\pgfqpoint{3.414725in}{1.615871in}}%
\pgfpathlineto{\pgfqpoint{3.424277in}{1.602623in}}%
\pgfpathlineto{\pgfqpoint{3.423866in}{1.589375in}}%
\pgfpathlineto{\pgfqpoint{3.431347in}{1.577994in}}%
\pgfpathlineto{\pgfqpoint{3.432410in}{1.576127in}}%
\pgfpathlineto{\pgfqpoint{3.432710in}{1.562879in}}%
\pgfpathlineto{\pgfqpoint{3.439613in}{1.549631in}}%
\pgfpathlineto{\pgfqpoint{3.440420in}{1.536383in}}%
\pgfpathlineto{\pgfqpoint{3.445461in}{1.523135in}}%
\pgfpathlineto{\pgfqpoint{3.446052in}{1.509887in}}%
\pgfpathlineto{\pgfqpoint{3.450518in}{1.496639in}}%
\pgfpathlineto{\pgfqpoint{3.454153in}{1.483391in}}%
\pgfpathlineto{\pgfqpoint{3.456630in}{1.473906in}}%
\pgfpathlineto{\pgfqpoint{3.457561in}{1.470143in}}%
\pgfpathlineto{\pgfqpoint{3.461330in}{1.456895in}}%
\pgfpathlineto{\pgfqpoint{3.462598in}{1.443647in}}%
\pgfpathlineto{\pgfqpoint{3.467063in}{1.430399in}}%
\pgfpathlineto{\pgfqpoint{3.467133in}{1.417151in}}%
\pgfpathlineto{\pgfqpoint{3.473066in}{1.403903in}}%
\pgfpathlineto{\pgfqpoint{3.470751in}{1.390655in}}%
\pgfpathlineto{\pgfqpoint{3.477389in}{1.377407in}}%
\pgfpathlineto{\pgfqpoint{3.474598in}{1.364159in}}%
\pgfpathlineto{\pgfqpoint{3.481913in}{1.350937in}}%
\pgfpathlineto{\pgfqpoint{3.481927in}{1.350911in}}%
\pgfpathlineto{\pgfqpoint{3.481913in}{1.350863in}}%
\pgfpathlineto{\pgfqpoint{3.477760in}{1.337663in}}%
\pgfpathlineto{\pgfqpoint{3.481913in}{1.329538in}}%
\pgfpathlineto{\pgfqpoint{3.484368in}{1.324415in}}%
\pgfpathlineto{\pgfqpoint{3.481913in}{1.317130in}}%
\pgfpathlineto{\pgfqpoint{3.479647in}{1.311167in}}%
\pgfpathlineto{\pgfqpoint{3.481913in}{1.307254in}}%
\pgfpathlineto{\pgfqpoint{3.487035in}{1.297919in}}%
\pgfpathlineto{\pgfqpoint{3.484430in}{1.284671in}}%
\pgfpathlineto{\pgfqpoint{3.489319in}{1.271423in}}%
\pgfpathlineto{\pgfqpoint{3.485112in}{1.258175in}}%
\pgfpathlineto{\pgfqpoint{3.487541in}{1.244927in}}%
\pgfpathlineto{\pgfqpoint{3.484283in}{1.231679in}}%
\pgfpathlineto{\pgfqpoint{3.484177in}{1.218431in}}%
\pgfpathlineto{\pgfqpoint{3.482957in}{1.205183in}}%
\pgfpathlineto{\pgfqpoint{3.481913in}{1.202045in}}%
\pgfpathlineto{\pgfqpoint{3.477808in}{1.191935in}}%
\pgfpathlineto{\pgfqpoint{3.479086in}{1.178687in}}%
\pgfpathlineto{\pgfqpoint{3.471038in}{1.165439in}}%
\pgfpathlineto{\pgfqpoint{3.473634in}{1.152191in}}%
\pgfpathlineto{\pgfqpoint{3.457241in}{1.138943in}}%
\pgfpathlineto{\pgfqpoint{3.461053in}{1.125695in}}%
\pgfpathlineto{\pgfqpoint{3.456630in}{1.122711in}}%
\pgfpathlineto{\pgfqpoint{3.436976in}{1.112447in}}%
\pgfpathlineto{\pgfqpoint{3.442749in}{1.099199in}}%
\pgfpathlineto{\pgfqpoint{3.431347in}{1.094108in}}%
\pgfpathlineto{\pgfqpoint{3.406300in}{1.085951in}}%
\pgfpathlineto{\pgfqpoint{3.406064in}{1.085381in}}%
\pgfpathlineto{\pgfqpoint{3.397884in}{1.072703in}}%
\pgfpathlineto{\pgfqpoint{3.380781in}{1.069469in}}%
\pgfpathlineto{\pgfqpoint{3.355497in}{1.066689in}}%
\pgfpathlineto{\pgfqpoint{3.330214in}{1.065933in}}%
\pgfpathlineto{\pgfqpoint{3.304931in}{1.067445in}}%
\pgfpathlineto{\pgfqpoint{3.279648in}{1.071777in}}%
\pgfpathlineto{\pgfqpoint{3.276485in}{1.072703in}}%
\pgfpathclose%
\pgfusepath{fill}%
\end{pgfscope}%
\begin{pgfscope}%
\pgfpathrectangle{\pgfqpoint{2.470586in}{0.476543in}}{\pgfqpoint{1.820388in}{1.656000in}}%
\pgfusepath{clip}%
\pgfsetbuttcap%
\pgfsetroundjoin%
\definecolor{currentfill}{rgb}{0.981776,0.607382,0.345790}%
\pgfsetfillcolor{currentfill}%
\pgfsetlinewidth{0.000000pt}%
\definecolor{currentstroke}{rgb}{0.000000,0.000000,0.000000}%
\pgfsetstrokecolor{currentstroke}%
\pgfsetdash{}{0pt}%
\pgfpathmoveto{\pgfqpoint{3.279648in}{1.071777in}}%
\pgfpathlineto{\pgfqpoint{3.304931in}{1.067445in}}%
\pgfpathlineto{\pgfqpoint{3.330214in}{1.065933in}}%
\pgfpathlineto{\pgfqpoint{3.355497in}{1.066689in}}%
\pgfpathlineto{\pgfqpoint{3.380781in}{1.069469in}}%
\pgfpathlineto{\pgfqpoint{3.397884in}{1.072703in}}%
\pgfpathlineto{\pgfqpoint{3.406064in}{1.085381in}}%
\pgfpathlineto{\pgfqpoint{3.406300in}{1.085951in}}%
\pgfpathlineto{\pgfqpoint{3.431347in}{1.094108in}}%
\pgfpathlineto{\pgfqpoint{3.442749in}{1.099199in}}%
\pgfpathlineto{\pgfqpoint{3.436976in}{1.112447in}}%
\pgfpathlineto{\pgfqpoint{3.456630in}{1.122711in}}%
\pgfpathlineto{\pgfqpoint{3.461053in}{1.125695in}}%
\pgfpathlineto{\pgfqpoint{3.457241in}{1.138943in}}%
\pgfpathlineto{\pgfqpoint{3.473634in}{1.152191in}}%
\pgfpathlineto{\pgfqpoint{3.471038in}{1.165439in}}%
\pgfpathlineto{\pgfqpoint{3.479086in}{1.178687in}}%
\pgfpathlineto{\pgfqpoint{3.477808in}{1.191935in}}%
\pgfpathlineto{\pgfqpoint{3.481913in}{1.202045in}}%
\pgfpathlineto{\pgfqpoint{3.482957in}{1.205183in}}%
\pgfpathlineto{\pgfqpoint{3.484177in}{1.218431in}}%
\pgfpathlineto{\pgfqpoint{3.484283in}{1.231679in}}%
\pgfpathlineto{\pgfqpoint{3.487541in}{1.244927in}}%
\pgfpathlineto{\pgfqpoint{3.485112in}{1.258175in}}%
\pgfpathlineto{\pgfqpoint{3.489319in}{1.271423in}}%
\pgfpathlineto{\pgfqpoint{3.484430in}{1.284671in}}%
\pgfpathlineto{\pgfqpoint{3.487035in}{1.297919in}}%
\pgfpathlineto{\pgfqpoint{3.481913in}{1.307254in}}%
\pgfpathlineto{\pgfqpoint{3.479647in}{1.311167in}}%
\pgfpathlineto{\pgfqpoint{3.481913in}{1.317130in}}%
\pgfpathlineto{\pgfqpoint{3.484368in}{1.324415in}}%
\pgfpathlineto{\pgfqpoint{3.481913in}{1.329538in}}%
\pgfpathlineto{\pgfqpoint{3.477760in}{1.337663in}}%
\pgfpathlineto{\pgfqpoint{3.481913in}{1.350863in}}%
\pgfpathlineto{\pgfqpoint{3.481927in}{1.350911in}}%
\pgfpathlineto{\pgfqpoint{3.481913in}{1.350937in}}%
\pgfpathlineto{\pgfqpoint{3.474598in}{1.364159in}}%
\pgfpathlineto{\pgfqpoint{3.477389in}{1.377407in}}%
\pgfpathlineto{\pgfqpoint{3.470751in}{1.390655in}}%
\pgfpathlineto{\pgfqpoint{3.473066in}{1.403903in}}%
\pgfpathlineto{\pgfqpoint{3.467133in}{1.417151in}}%
\pgfpathlineto{\pgfqpoint{3.467063in}{1.430399in}}%
\pgfpathlineto{\pgfqpoint{3.462598in}{1.443647in}}%
\pgfpathlineto{\pgfqpoint{3.461330in}{1.456895in}}%
\pgfpathlineto{\pgfqpoint{3.457561in}{1.470143in}}%
\pgfpathlineto{\pgfqpoint{3.456630in}{1.473906in}}%
\pgfpathlineto{\pgfqpoint{3.454153in}{1.483391in}}%
\pgfpathlineto{\pgfqpoint{3.450518in}{1.496639in}}%
\pgfpathlineto{\pgfqpoint{3.446052in}{1.509887in}}%
\pgfpathlineto{\pgfqpoint{3.445461in}{1.523135in}}%
\pgfpathlineto{\pgfqpoint{3.440420in}{1.536383in}}%
\pgfpathlineto{\pgfqpoint{3.439613in}{1.549631in}}%
\pgfpathlineto{\pgfqpoint{3.432710in}{1.562879in}}%
\pgfpathlineto{\pgfqpoint{3.432410in}{1.576127in}}%
\pgfpathlineto{\pgfqpoint{3.431347in}{1.577994in}}%
\pgfpathlineto{\pgfqpoint{3.423866in}{1.589375in}}%
\pgfpathlineto{\pgfqpoint{3.424277in}{1.602623in}}%
\pgfpathlineto{\pgfqpoint{3.414725in}{1.615871in}}%
\pgfpathlineto{\pgfqpoint{3.414562in}{1.629119in}}%
\pgfpathlineto{\pgfqpoint{3.406064in}{1.641587in}}%
\pgfpathlineto{\pgfqpoint{3.405355in}{1.642367in}}%
\pgfpathlineto{\pgfqpoint{3.405403in}{1.655615in}}%
\pgfpathlineto{\pgfqpoint{3.393180in}{1.668863in}}%
\pgfpathlineto{\pgfqpoint{3.392354in}{1.682111in}}%
\pgfpathlineto{\pgfqpoint{3.380781in}{1.694944in}}%
\pgfpathlineto{\pgfqpoint{3.380116in}{1.695359in}}%
\pgfpathlineto{\pgfqpoint{3.379173in}{1.708607in}}%
\pgfpathlineto{\pgfqpoint{3.362882in}{1.721855in}}%
\pgfpathlineto{\pgfqpoint{3.359760in}{1.735103in}}%
\pgfpathlineto{\pgfqpoint{3.355497in}{1.738273in}}%
\pgfpathlineto{\pgfqpoint{3.330214in}{1.741880in}}%
\pgfpathlineto{\pgfqpoint{3.314026in}{1.735103in}}%
\pgfpathlineto{\pgfqpoint{3.313122in}{1.721855in}}%
\pgfpathlineto{\pgfqpoint{3.304931in}{1.718047in}}%
\pgfpathlineto{\pgfqpoint{3.295379in}{1.708607in}}%
\pgfpathlineto{\pgfqpoint{3.295701in}{1.695359in}}%
\pgfpathlineto{\pgfqpoint{3.279740in}{1.682111in}}%
\pgfpathlineto{\pgfqpoint{3.279648in}{1.681159in}}%
\pgfpathlineto{\pgfqpoint{3.278794in}{1.668863in}}%
\pgfpathlineto{\pgfqpoint{3.269429in}{1.655615in}}%
\pgfpathlineto{\pgfqpoint{3.267299in}{1.642367in}}%
\pgfpathlineto{\pgfqpoint{3.259004in}{1.629119in}}%
\pgfpathlineto{\pgfqpoint{3.255352in}{1.615871in}}%
\pgfpathlineto{\pgfqpoint{3.254365in}{1.613506in}}%
\pgfpathlineto{\pgfqpoint{3.250815in}{1.602623in}}%
\pgfpathlineto{\pgfqpoint{3.247174in}{1.589375in}}%
\pgfpathlineto{\pgfqpoint{3.244358in}{1.576127in}}%
\pgfpathlineto{\pgfqpoint{3.238072in}{1.562879in}}%
\pgfpathlineto{\pgfqpoint{3.236678in}{1.549631in}}%
\pgfpathlineto{\pgfqpoint{3.230135in}{1.536383in}}%
\pgfpathlineto{\pgfqpoint{3.230828in}{1.523135in}}%
\pgfpathlineto{\pgfqpoint{3.229082in}{1.520282in}}%
\pgfpathlineto{\pgfqpoint{3.223355in}{1.509887in}}%
\pgfpathlineto{\pgfqpoint{3.223980in}{1.496639in}}%
\pgfpathlineto{\pgfqpoint{3.215356in}{1.483391in}}%
\pgfpathlineto{\pgfqpoint{3.217001in}{1.470143in}}%
\pgfpathlineto{\pgfqpoint{3.209434in}{1.456895in}}%
\pgfpathlineto{\pgfqpoint{3.211173in}{1.443647in}}%
\pgfpathlineto{\pgfqpoint{3.203880in}{1.430399in}}%
\pgfpathlineto{\pgfqpoint{3.205876in}{1.417151in}}%
\pgfpathlineto{\pgfqpoint{3.203798in}{1.412192in}}%
\pgfpathlineto{\pgfqpoint{3.200642in}{1.403903in}}%
\pgfpathlineto{\pgfqpoint{3.201360in}{1.390655in}}%
\pgfpathlineto{\pgfqpoint{3.196366in}{1.377407in}}%
\pgfpathlineto{\pgfqpoint{3.194819in}{1.364159in}}%
\pgfpathlineto{\pgfqpoint{3.192241in}{1.350911in}}%
\pgfpathlineto{\pgfqpoint{3.190890in}{1.337663in}}%
\pgfpathlineto{\pgfqpoint{3.190337in}{1.324415in}}%
\pgfpathlineto{\pgfqpoint{3.186273in}{1.311167in}}%
\pgfpathlineto{\pgfqpoint{3.187587in}{1.297919in}}%
\pgfpathlineto{\pgfqpoint{3.182551in}{1.284671in}}%
\pgfpathlineto{\pgfqpoint{3.187696in}{1.271423in}}%
\pgfpathlineto{\pgfqpoint{3.182241in}{1.258175in}}%
\pgfpathlineto{\pgfqpoint{3.189090in}{1.244927in}}%
\pgfpathlineto{\pgfqpoint{3.182093in}{1.231679in}}%
\pgfpathlineto{\pgfqpoint{3.191783in}{1.218431in}}%
\pgfpathlineto{\pgfqpoint{3.184402in}{1.205183in}}%
\pgfpathlineto{\pgfqpoint{3.194524in}{1.191935in}}%
\pgfpathlineto{\pgfqpoint{3.186100in}{1.178687in}}%
\pgfpathlineto{\pgfqpoint{3.198609in}{1.165439in}}%
\pgfpathlineto{\pgfqpoint{3.194381in}{1.152191in}}%
\pgfpathlineto{\pgfqpoint{3.203798in}{1.144594in}}%
\pgfpathlineto{\pgfqpoint{3.212787in}{1.138943in}}%
\pgfpathlineto{\pgfqpoint{3.209237in}{1.125695in}}%
\pgfpathlineto{\pgfqpoint{3.228620in}{1.112447in}}%
\pgfpathlineto{\pgfqpoint{3.229082in}{1.103772in}}%
\pgfpathlineto{\pgfqpoint{3.229381in}{1.099199in}}%
\pgfpathlineto{\pgfqpoint{3.254365in}{1.087883in}}%
\pgfpathlineto{\pgfqpoint{3.261004in}{1.085951in}}%
\pgfpathlineto{\pgfqpoint{3.276485in}{1.072703in}}%
\pgfpathlineto{\pgfqpoint{3.279648in}{1.071777in}}%
\pgfpathclose%
\pgfpathmoveto{\pgfqpoint{3.276511in}{1.099199in}}%
\pgfpathlineto{\pgfqpoint{3.275919in}{1.112447in}}%
\pgfpathlineto{\pgfqpoint{3.254365in}{1.119390in}}%
\pgfpathlineto{\pgfqpoint{3.241583in}{1.125695in}}%
\pgfpathlineto{\pgfqpoint{3.245347in}{1.138943in}}%
\pgfpathlineto{\pgfqpoint{3.229082in}{1.146333in}}%
\pgfpathlineto{\pgfqpoint{3.219585in}{1.152191in}}%
\pgfpathlineto{\pgfqpoint{3.223986in}{1.165439in}}%
\pgfpathlineto{\pgfqpoint{3.207666in}{1.178687in}}%
\pgfpathlineto{\pgfqpoint{3.216548in}{1.191935in}}%
\pgfpathlineto{\pgfqpoint{3.203798in}{1.204720in}}%
\pgfpathlineto{\pgfqpoint{3.203423in}{1.205183in}}%
\pgfpathlineto{\pgfqpoint{3.203798in}{1.205911in}}%
\pgfpathlineto{\pgfqpoint{3.210837in}{1.218431in}}%
\pgfpathlineto{\pgfqpoint{3.203798in}{1.226016in}}%
\pgfpathlineto{\pgfqpoint{3.199409in}{1.231679in}}%
\pgfpathlineto{\pgfqpoint{3.203798in}{1.240951in}}%
\pgfpathlineto{\pgfqpoint{3.205831in}{1.244927in}}%
\pgfpathlineto{\pgfqpoint{3.203798in}{1.248062in}}%
\pgfpathlineto{\pgfqpoint{3.198200in}{1.258175in}}%
\pgfpathlineto{\pgfqpoint{3.202716in}{1.271423in}}%
\pgfpathlineto{\pgfqpoint{3.197253in}{1.284671in}}%
\pgfpathlineto{\pgfqpoint{3.201494in}{1.297919in}}%
\pgfpathlineto{\pgfqpoint{3.200104in}{1.311167in}}%
\pgfpathlineto{\pgfqpoint{3.203321in}{1.324415in}}%
\pgfpathlineto{\pgfqpoint{3.203798in}{1.336355in}}%
\pgfpathlineto{\pgfqpoint{3.203857in}{1.337663in}}%
\pgfpathlineto{\pgfqpoint{3.204448in}{1.350911in}}%
\pgfpathlineto{\pgfqpoint{3.207427in}{1.364159in}}%
\pgfpathlineto{\pgfqpoint{3.208299in}{1.377407in}}%
\pgfpathlineto{\pgfqpoint{3.213865in}{1.390655in}}%
\pgfpathlineto{\pgfqpoint{3.212242in}{1.403903in}}%
\pgfpathlineto{\pgfqpoint{3.217830in}{1.417151in}}%
\pgfpathlineto{\pgfqpoint{3.215218in}{1.430399in}}%
\pgfpathlineto{\pgfqpoint{3.222554in}{1.443647in}}%
\pgfpathlineto{\pgfqpoint{3.220326in}{1.456895in}}%
\pgfpathlineto{\pgfqpoint{3.227840in}{1.470143in}}%
\pgfpathlineto{\pgfqpoint{3.225864in}{1.483391in}}%
\pgfpathlineto{\pgfqpoint{3.229082in}{1.488346in}}%
\pgfpathlineto{\pgfqpoint{3.235223in}{1.496639in}}%
\pgfpathlineto{\pgfqpoint{3.234295in}{1.509887in}}%
\pgfpathlineto{\pgfqpoint{3.242223in}{1.523135in}}%
\pgfpathlineto{\pgfqpoint{3.241457in}{1.536383in}}%
\pgfpathlineto{\pgfqpoint{3.247592in}{1.549631in}}%
\pgfpathlineto{\pgfqpoint{3.249095in}{1.562879in}}%
\pgfpathlineto{\pgfqpoint{3.254365in}{1.574815in}}%
\pgfpathlineto{\pgfqpoint{3.255051in}{1.576127in}}%
\pgfpathlineto{\pgfqpoint{3.258768in}{1.589375in}}%
\pgfpathlineto{\pgfqpoint{3.262613in}{1.602623in}}%
\pgfpathlineto{\pgfqpoint{3.268466in}{1.615871in}}%
\pgfpathlineto{\pgfqpoint{3.271469in}{1.629119in}}%
\pgfpathlineto{\pgfqpoint{3.279648in}{1.641621in}}%
\pgfpathlineto{\pgfqpoint{3.280389in}{1.642367in}}%
\pgfpathlineto{\pgfqpoint{3.282532in}{1.655615in}}%
\pgfpathlineto{\pgfqpoint{3.296674in}{1.668863in}}%
\pgfpathlineto{\pgfqpoint{3.297232in}{1.682111in}}%
\pgfpathlineto{\pgfqpoint{3.304931in}{1.688448in}}%
\pgfpathlineto{\pgfqpoint{3.323563in}{1.695359in}}%
\pgfpathlineto{\pgfqpoint{3.321916in}{1.708607in}}%
\pgfpathlineto{\pgfqpoint{3.330214in}{1.712446in}}%
\pgfpathlineto{\pgfqpoint{3.355497in}{1.708634in}}%
\pgfpathlineto{\pgfqpoint{3.355531in}{1.708607in}}%
\pgfpathlineto{\pgfqpoint{3.355497in}{1.707596in}}%
\pgfpathlineto{\pgfqpoint{3.353159in}{1.695359in}}%
\pgfpathlineto{\pgfqpoint{3.355497in}{1.695100in}}%
\pgfpathlineto{\pgfqpoint{3.376424in}{1.682111in}}%
\pgfpathlineto{\pgfqpoint{3.376905in}{1.668863in}}%
\pgfpathlineto{\pgfqpoint{3.380781in}{1.666440in}}%
\pgfpathlineto{\pgfqpoint{3.390880in}{1.655615in}}%
\pgfpathlineto{\pgfqpoint{3.390479in}{1.642367in}}%
\pgfpathlineto{\pgfqpoint{3.402544in}{1.629119in}}%
\pgfpathlineto{\pgfqpoint{3.402508in}{1.615871in}}%
\pgfpathlineto{\pgfqpoint{3.406064in}{1.612089in}}%
\pgfpathlineto{\pgfqpoint{3.412854in}{1.602623in}}%
\pgfpathlineto{\pgfqpoint{3.412254in}{1.589375in}}%
\pgfpathlineto{\pgfqpoint{3.420698in}{1.576127in}}%
\pgfpathlineto{\pgfqpoint{3.421042in}{1.562879in}}%
\pgfpathlineto{\pgfqpoint{3.428863in}{1.549631in}}%
\pgfpathlineto{\pgfqpoint{3.429767in}{1.536383in}}%
\pgfpathlineto{\pgfqpoint{3.431347in}{1.532495in}}%
\pgfpathlineto{\pgfqpoint{3.434678in}{1.523135in}}%
\pgfpathlineto{\pgfqpoint{3.435329in}{1.509887in}}%
\pgfpathlineto{\pgfqpoint{3.439331in}{1.496639in}}%
\pgfpathlineto{\pgfqpoint{3.443157in}{1.483391in}}%
\pgfpathlineto{\pgfqpoint{3.446059in}{1.470143in}}%
\pgfpathlineto{\pgfqpoint{3.450342in}{1.456895in}}%
\pgfpathlineto{\pgfqpoint{3.451136in}{1.443647in}}%
\pgfpathlineto{\pgfqpoint{3.456276in}{1.430399in}}%
\pgfpathlineto{\pgfqpoint{3.455657in}{1.417151in}}%
\pgfpathlineto{\pgfqpoint{3.456630in}{1.415239in}}%
\pgfpathlineto{\pgfqpoint{3.461729in}{1.403903in}}%
\pgfpathlineto{\pgfqpoint{3.458724in}{1.390655in}}%
\pgfpathlineto{\pgfqpoint{3.465324in}{1.377407in}}%
\pgfpathlineto{\pgfqpoint{3.461953in}{1.364159in}}%
\pgfpathlineto{\pgfqpoint{3.469147in}{1.350911in}}%
\pgfpathlineto{\pgfqpoint{3.464348in}{1.337663in}}%
\pgfpathlineto{\pgfqpoint{3.470863in}{1.324415in}}%
\pgfpathlineto{\pgfqpoint{3.465460in}{1.311167in}}%
\pgfpathlineto{\pgfqpoint{3.472774in}{1.297919in}}%
\pgfpathlineto{\pgfqpoint{3.469520in}{1.284671in}}%
\pgfpathlineto{\pgfqpoint{3.473980in}{1.271423in}}%
\pgfpathlineto{\pgfqpoint{3.468919in}{1.258175in}}%
\pgfpathlineto{\pgfqpoint{3.470444in}{1.244927in}}%
\pgfpathlineto{\pgfqpoint{3.466676in}{1.231679in}}%
\pgfpathlineto{\pgfqpoint{3.464878in}{1.218431in}}%
\pgfpathlineto{\pgfqpoint{3.463380in}{1.205183in}}%
\pgfpathlineto{\pgfqpoint{3.456630in}{1.192637in}}%
\pgfpathlineto{\pgfqpoint{3.456143in}{1.191935in}}%
\pgfpathlineto{\pgfqpoint{3.456630in}{1.187536in}}%
\pgfpathlineto{\pgfqpoint{3.457463in}{1.178687in}}%
\pgfpathlineto{\pgfqpoint{3.456630in}{1.177506in}}%
\pgfpathlineto{\pgfqpoint{3.445611in}{1.165439in}}%
\pgfpathlineto{\pgfqpoint{3.448358in}{1.152191in}}%
\pgfpathlineto{\pgfqpoint{3.431347in}{1.141958in}}%
\pgfpathlineto{\pgfqpoint{3.424289in}{1.138943in}}%
\pgfpathlineto{\pgfqpoint{3.429105in}{1.125695in}}%
\pgfpathlineto{\pgfqpoint{3.406064in}{1.116723in}}%
\pgfpathlineto{\pgfqpoint{3.389519in}{1.112447in}}%
\pgfpathlineto{\pgfqpoint{3.395460in}{1.099199in}}%
\pgfpathlineto{\pgfqpoint{3.380781in}{1.095811in}}%
\pgfpathlineto{\pgfqpoint{3.355497in}{1.092552in}}%
\pgfpathlineto{\pgfqpoint{3.330214in}{1.091758in}}%
\pgfpathlineto{\pgfqpoint{3.304931in}{1.093516in}}%
\pgfpathlineto{\pgfqpoint{3.279648in}{1.098224in}}%
\pgfpathlineto{\pgfqpoint{3.276511in}{1.099199in}}%
\pgfpathclose%
\pgfusepath{fill}%
\end{pgfscope}%
\begin{pgfscope}%
\pgfpathrectangle{\pgfqpoint{2.470586in}{0.476543in}}{\pgfqpoint{1.820388in}{1.656000in}}%
\pgfusepath{clip}%
\pgfsetbuttcap%
\pgfsetroundjoin%
\definecolor{currentfill}{rgb}{0.966551,0.497424,0.295040}%
\pgfsetfillcolor{currentfill}%
\pgfsetlinewidth{0.000000pt}%
\definecolor{currentstroke}{rgb}{0.000000,0.000000,0.000000}%
\pgfsetstrokecolor{currentstroke}%
\pgfsetdash{}{0pt}%
\pgfpathmoveto{\pgfqpoint{3.279648in}{1.098224in}}%
\pgfpathlineto{\pgfqpoint{3.304931in}{1.093516in}}%
\pgfpathlineto{\pgfqpoint{3.330214in}{1.091758in}}%
\pgfpathlineto{\pgfqpoint{3.355497in}{1.092552in}}%
\pgfpathlineto{\pgfqpoint{3.380781in}{1.095811in}}%
\pgfpathlineto{\pgfqpoint{3.395460in}{1.099199in}}%
\pgfpathlineto{\pgfqpoint{3.389519in}{1.112447in}}%
\pgfpathlineto{\pgfqpoint{3.406064in}{1.116723in}}%
\pgfpathlineto{\pgfqpoint{3.429105in}{1.125695in}}%
\pgfpathlineto{\pgfqpoint{3.424289in}{1.138943in}}%
\pgfpathlineto{\pgfqpoint{3.431347in}{1.141958in}}%
\pgfpathlineto{\pgfqpoint{3.448358in}{1.152191in}}%
\pgfpathlineto{\pgfqpoint{3.445611in}{1.165439in}}%
\pgfpathlineto{\pgfqpoint{3.456630in}{1.177506in}}%
\pgfpathlineto{\pgfqpoint{3.457463in}{1.178687in}}%
\pgfpathlineto{\pgfqpoint{3.456630in}{1.187536in}}%
\pgfpathlineto{\pgfqpoint{3.456143in}{1.191935in}}%
\pgfpathlineto{\pgfqpoint{3.456630in}{1.192637in}}%
\pgfpathlineto{\pgfqpoint{3.463380in}{1.205183in}}%
\pgfpathlineto{\pgfqpoint{3.464878in}{1.218431in}}%
\pgfpathlineto{\pgfqpoint{3.466676in}{1.231679in}}%
\pgfpathlineto{\pgfqpoint{3.470444in}{1.244927in}}%
\pgfpathlineto{\pgfqpoint{3.468919in}{1.258175in}}%
\pgfpathlineto{\pgfqpoint{3.473980in}{1.271423in}}%
\pgfpathlineto{\pgfqpoint{3.469520in}{1.284671in}}%
\pgfpathlineto{\pgfqpoint{3.472774in}{1.297919in}}%
\pgfpathlineto{\pgfqpoint{3.465460in}{1.311167in}}%
\pgfpathlineto{\pgfqpoint{3.470863in}{1.324415in}}%
\pgfpathlineto{\pgfqpoint{3.464348in}{1.337663in}}%
\pgfpathlineto{\pgfqpoint{3.469147in}{1.350911in}}%
\pgfpathlineto{\pgfqpoint{3.461953in}{1.364159in}}%
\pgfpathlineto{\pgfqpoint{3.465324in}{1.377407in}}%
\pgfpathlineto{\pgfqpoint{3.458724in}{1.390655in}}%
\pgfpathlineto{\pgfqpoint{3.461729in}{1.403903in}}%
\pgfpathlineto{\pgfqpoint{3.456630in}{1.415239in}}%
\pgfpathlineto{\pgfqpoint{3.455657in}{1.417151in}}%
\pgfpathlineto{\pgfqpoint{3.456276in}{1.430399in}}%
\pgfpathlineto{\pgfqpoint{3.451136in}{1.443647in}}%
\pgfpathlineto{\pgfqpoint{3.450342in}{1.456895in}}%
\pgfpathlineto{\pgfqpoint{3.446059in}{1.470143in}}%
\pgfpathlineto{\pgfqpoint{3.443157in}{1.483391in}}%
\pgfpathlineto{\pgfqpoint{3.439331in}{1.496639in}}%
\pgfpathlineto{\pgfqpoint{3.435329in}{1.509887in}}%
\pgfpathlineto{\pgfqpoint{3.434678in}{1.523135in}}%
\pgfpathlineto{\pgfqpoint{3.431347in}{1.532495in}}%
\pgfpathlineto{\pgfqpoint{3.429767in}{1.536383in}}%
\pgfpathlineto{\pgfqpoint{3.428863in}{1.549631in}}%
\pgfpathlineto{\pgfqpoint{3.421042in}{1.562879in}}%
\pgfpathlineto{\pgfqpoint{3.420698in}{1.576127in}}%
\pgfpathlineto{\pgfqpoint{3.412254in}{1.589375in}}%
\pgfpathlineto{\pgfqpoint{3.412854in}{1.602623in}}%
\pgfpathlineto{\pgfqpoint{3.406064in}{1.612089in}}%
\pgfpathlineto{\pgfqpoint{3.402508in}{1.615871in}}%
\pgfpathlineto{\pgfqpoint{3.402544in}{1.629119in}}%
\pgfpathlineto{\pgfqpoint{3.390479in}{1.642367in}}%
\pgfpathlineto{\pgfqpoint{3.390880in}{1.655615in}}%
\pgfpathlineto{\pgfqpoint{3.380781in}{1.666440in}}%
\pgfpathlineto{\pgfqpoint{3.376905in}{1.668863in}}%
\pgfpathlineto{\pgfqpoint{3.376424in}{1.682111in}}%
\pgfpathlineto{\pgfqpoint{3.355497in}{1.695100in}}%
\pgfpathlineto{\pgfqpoint{3.353159in}{1.695359in}}%
\pgfpathlineto{\pgfqpoint{3.355497in}{1.707596in}}%
\pgfpathlineto{\pgfqpoint{3.355531in}{1.708607in}}%
\pgfpathlineto{\pgfqpoint{3.355497in}{1.708634in}}%
\pgfpathlineto{\pgfqpoint{3.330214in}{1.712446in}}%
\pgfpathlineto{\pgfqpoint{3.321916in}{1.708607in}}%
\pgfpathlineto{\pgfqpoint{3.323563in}{1.695359in}}%
\pgfpathlineto{\pgfqpoint{3.304931in}{1.688448in}}%
\pgfpathlineto{\pgfqpoint{3.297232in}{1.682111in}}%
\pgfpathlineto{\pgfqpoint{3.296674in}{1.668863in}}%
\pgfpathlineto{\pgfqpoint{3.282532in}{1.655615in}}%
\pgfpathlineto{\pgfqpoint{3.280389in}{1.642367in}}%
\pgfpathlineto{\pgfqpoint{3.279648in}{1.641621in}}%
\pgfpathlineto{\pgfqpoint{3.271469in}{1.629119in}}%
\pgfpathlineto{\pgfqpoint{3.268466in}{1.615871in}}%
\pgfpathlineto{\pgfqpoint{3.262613in}{1.602623in}}%
\pgfpathlineto{\pgfqpoint{3.258768in}{1.589375in}}%
\pgfpathlineto{\pgfqpoint{3.255051in}{1.576127in}}%
\pgfpathlineto{\pgfqpoint{3.254365in}{1.574815in}}%
\pgfpathlineto{\pgfqpoint{3.249095in}{1.562879in}}%
\pgfpathlineto{\pgfqpoint{3.247592in}{1.549631in}}%
\pgfpathlineto{\pgfqpoint{3.241457in}{1.536383in}}%
\pgfpathlineto{\pgfqpoint{3.242223in}{1.523135in}}%
\pgfpathlineto{\pgfqpoint{3.234295in}{1.509887in}}%
\pgfpathlineto{\pgfqpoint{3.235223in}{1.496639in}}%
\pgfpathlineto{\pgfqpoint{3.229082in}{1.488346in}}%
\pgfpathlineto{\pgfqpoint{3.225864in}{1.483391in}}%
\pgfpathlineto{\pgfqpoint{3.227840in}{1.470143in}}%
\pgfpathlineto{\pgfqpoint{3.220326in}{1.456895in}}%
\pgfpathlineto{\pgfqpoint{3.222554in}{1.443647in}}%
\pgfpathlineto{\pgfqpoint{3.215218in}{1.430399in}}%
\pgfpathlineto{\pgfqpoint{3.217830in}{1.417151in}}%
\pgfpathlineto{\pgfqpoint{3.212242in}{1.403903in}}%
\pgfpathlineto{\pgfqpoint{3.213865in}{1.390655in}}%
\pgfpathlineto{\pgfqpoint{3.208299in}{1.377407in}}%
\pgfpathlineto{\pgfqpoint{3.207427in}{1.364159in}}%
\pgfpathlineto{\pgfqpoint{3.204448in}{1.350911in}}%
\pgfpathlineto{\pgfqpoint{3.203857in}{1.337663in}}%
\pgfpathlineto{\pgfqpoint{3.203798in}{1.336355in}}%
\pgfpathlineto{\pgfqpoint{3.203321in}{1.324415in}}%
\pgfpathlineto{\pgfqpoint{3.200104in}{1.311167in}}%
\pgfpathlineto{\pgfqpoint{3.201494in}{1.297919in}}%
\pgfpathlineto{\pgfqpoint{3.197253in}{1.284671in}}%
\pgfpathlineto{\pgfqpoint{3.202716in}{1.271423in}}%
\pgfpathlineto{\pgfqpoint{3.198200in}{1.258175in}}%
\pgfpathlineto{\pgfqpoint{3.203798in}{1.248062in}}%
\pgfpathlineto{\pgfqpoint{3.205831in}{1.244927in}}%
\pgfpathlineto{\pgfqpoint{3.203798in}{1.240951in}}%
\pgfpathlineto{\pgfqpoint{3.199409in}{1.231679in}}%
\pgfpathlineto{\pgfqpoint{3.203798in}{1.226016in}}%
\pgfpathlineto{\pgfqpoint{3.210837in}{1.218431in}}%
\pgfpathlineto{\pgfqpoint{3.203798in}{1.205911in}}%
\pgfpathlineto{\pgfqpoint{3.203423in}{1.205183in}}%
\pgfpathlineto{\pgfqpoint{3.203798in}{1.204720in}}%
\pgfpathlineto{\pgfqpoint{3.216548in}{1.191935in}}%
\pgfpathlineto{\pgfqpoint{3.207666in}{1.178687in}}%
\pgfpathlineto{\pgfqpoint{3.223986in}{1.165439in}}%
\pgfpathlineto{\pgfqpoint{3.219585in}{1.152191in}}%
\pgfpathlineto{\pgfqpoint{3.229082in}{1.146333in}}%
\pgfpathlineto{\pgfqpoint{3.245347in}{1.138943in}}%
\pgfpathlineto{\pgfqpoint{3.241583in}{1.125695in}}%
\pgfpathlineto{\pgfqpoint{3.254365in}{1.119390in}}%
\pgfpathlineto{\pgfqpoint{3.275919in}{1.112447in}}%
\pgfpathlineto{\pgfqpoint{3.276511in}{1.099199in}}%
\pgfpathlineto{\pgfqpoint{3.279648in}{1.098224in}}%
\pgfpathclose%
\pgfpathmoveto{\pgfqpoint{3.290483in}{1.125695in}}%
\pgfpathlineto{\pgfqpoint{3.297161in}{1.138943in}}%
\pgfpathlineto{\pgfqpoint{3.279648in}{1.142223in}}%
\pgfpathlineto{\pgfqpoint{3.254365in}{1.150294in}}%
\pgfpathlineto{\pgfqpoint{3.250241in}{1.152191in}}%
\pgfpathlineto{\pgfqpoint{3.254365in}{1.163907in}}%
\pgfpathlineto{\pgfqpoint{3.255067in}{1.165439in}}%
\pgfpathlineto{\pgfqpoint{3.254365in}{1.165735in}}%
\pgfpathlineto{\pgfqpoint{3.232754in}{1.178687in}}%
\pgfpathlineto{\pgfqpoint{3.242298in}{1.191935in}}%
\pgfpathlineto{\pgfqpoint{3.229082in}{1.201584in}}%
\pgfpathlineto{\pgfqpoint{3.225284in}{1.205183in}}%
\pgfpathlineto{\pgfqpoint{3.229082in}{1.213007in}}%
\pgfpathlineto{\pgfqpoint{3.232085in}{1.218431in}}%
\pgfpathlineto{\pgfqpoint{3.229082in}{1.220898in}}%
\pgfpathlineto{\pgfqpoint{3.218641in}{1.231679in}}%
\pgfpathlineto{\pgfqpoint{3.224530in}{1.244927in}}%
\pgfpathlineto{\pgfqpoint{3.215629in}{1.258175in}}%
\pgfpathlineto{\pgfqpoint{3.219499in}{1.271423in}}%
\pgfpathlineto{\pgfqpoint{3.213066in}{1.284671in}}%
\pgfpathlineto{\pgfqpoint{3.216835in}{1.297919in}}%
\pgfpathlineto{\pgfqpoint{3.215273in}{1.311167in}}%
\pgfpathlineto{\pgfqpoint{3.217770in}{1.324415in}}%
\pgfpathlineto{\pgfqpoint{3.218402in}{1.337663in}}%
\pgfpathlineto{\pgfqpoint{3.217936in}{1.350911in}}%
\pgfpathlineto{\pgfqpoint{3.221059in}{1.364159in}}%
\pgfpathlineto{\pgfqpoint{3.221006in}{1.377407in}}%
\pgfpathlineto{\pgfqpoint{3.226621in}{1.390655in}}%
\pgfpathlineto{\pgfqpoint{3.224144in}{1.403903in}}%
\pgfpathlineto{\pgfqpoint{3.229082in}{1.415510in}}%
\pgfpathlineto{\pgfqpoint{3.229902in}{1.417151in}}%
\pgfpathlineto{\pgfqpoint{3.229082in}{1.419913in}}%
\pgfpathlineto{\pgfqpoint{3.226555in}{1.430399in}}%
\pgfpathlineto{\pgfqpoint{3.229082in}{1.434946in}}%
\pgfpathlineto{\pgfqpoint{3.234717in}{1.443647in}}%
\pgfpathlineto{\pgfqpoint{3.231568in}{1.456895in}}%
\pgfpathlineto{\pgfqpoint{3.240147in}{1.470143in}}%
\pgfpathlineto{\pgfqpoint{3.237530in}{1.483391in}}%
\pgfpathlineto{\pgfqpoint{3.247206in}{1.496639in}}%
\pgfpathlineto{\pgfqpoint{3.246099in}{1.509887in}}%
\pgfpathlineto{\pgfqpoint{3.253617in}{1.523135in}}%
\pgfpathlineto{\pgfqpoint{3.252780in}{1.536383in}}%
\pgfpathlineto{\pgfqpoint{3.254365in}{1.539954in}}%
\pgfpathlineto{\pgfqpoint{3.259514in}{1.549631in}}%
\pgfpathlineto{\pgfqpoint{3.261609in}{1.562879in}}%
\pgfpathlineto{\pgfqpoint{3.268106in}{1.576127in}}%
\pgfpathlineto{\pgfqpoint{3.272155in}{1.589375in}}%
\pgfpathlineto{\pgfqpoint{3.275230in}{1.602623in}}%
\pgfpathlineto{\pgfqpoint{3.279648in}{1.611946in}}%
\pgfpathlineto{\pgfqpoint{3.282526in}{1.615871in}}%
\pgfpathlineto{\pgfqpoint{3.285916in}{1.629119in}}%
\pgfpathlineto{\pgfqpoint{3.299310in}{1.642367in}}%
\pgfpathlineto{\pgfqpoint{3.300267in}{1.655615in}}%
\pgfpathlineto{\pgfqpoint{3.304931in}{1.659911in}}%
\pgfpathlineto{\pgfqpoint{3.327382in}{1.668863in}}%
\pgfpathlineto{\pgfqpoint{3.326781in}{1.682111in}}%
\pgfpathlineto{\pgfqpoint{3.330214in}{1.683383in}}%
\pgfpathlineto{\pgfqpoint{3.340567in}{1.682111in}}%
\pgfpathlineto{\pgfqpoint{3.338558in}{1.668863in}}%
\pgfpathlineto{\pgfqpoint{3.355497in}{1.666683in}}%
\pgfpathlineto{\pgfqpoint{3.373209in}{1.655615in}}%
\pgfpathlineto{\pgfqpoint{3.372042in}{1.642367in}}%
\pgfpathlineto{\pgfqpoint{3.380781in}{1.636737in}}%
\pgfpathlineto{\pgfqpoint{3.387747in}{1.629119in}}%
\pgfpathlineto{\pgfqpoint{3.387539in}{1.615871in}}%
\pgfpathlineto{\pgfqpoint{3.399848in}{1.602623in}}%
\pgfpathlineto{\pgfqpoint{3.398870in}{1.589375in}}%
\pgfpathlineto{\pgfqpoint{3.406064in}{1.580488in}}%
\pgfpathlineto{\pgfqpoint{3.408791in}{1.576127in}}%
\pgfpathlineto{\pgfqpoint{3.409136in}{1.562879in}}%
\pgfpathlineto{\pgfqpoint{3.416528in}{1.549631in}}%
\pgfpathlineto{\pgfqpoint{3.417455in}{1.536383in}}%
\pgfpathlineto{\pgfqpoint{3.422398in}{1.523135in}}%
\pgfpathlineto{\pgfqpoint{3.423346in}{1.509887in}}%
\pgfpathlineto{\pgfqpoint{3.427494in}{1.496639in}}%
\pgfpathlineto{\pgfqpoint{3.431347in}{1.486114in}}%
\pgfpathlineto{\pgfqpoint{3.432161in}{1.483391in}}%
\pgfpathlineto{\pgfqpoint{3.434450in}{1.470143in}}%
\pgfpathlineto{\pgfqpoint{3.438830in}{1.456895in}}%
\pgfpathlineto{\pgfqpoint{3.438947in}{1.443647in}}%
\pgfpathlineto{\pgfqpoint{3.444240in}{1.430399in}}%
\pgfpathlineto{\pgfqpoint{3.442835in}{1.417151in}}%
\pgfpathlineto{\pgfqpoint{3.449571in}{1.403903in}}%
\pgfpathlineto{\pgfqpoint{3.445362in}{1.390655in}}%
\pgfpathlineto{\pgfqpoint{3.452791in}{1.377407in}}%
\pgfpathlineto{\pgfqpoint{3.448304in}{1.364159in}}%
\pgfpathlineto{\pgfqpoint{3.456327in}{1.350911in}}%
\pgfpathlineto{\pgfqpoint{3.450130in}{1.337663in}}%
\pgfpathlineto{\pgfqpoint{3.456630in}{1.325465in}}%
\pgfpathlineto{\pgfqpoint{3.457132in}{1.324415in}}%
\pgfpathlineto{\pgfqpoint{3.456630in}{1.323246in}}%
\pgfpathlineto{\pgfqpoint{3.450504in}{1.311167in}}%
\pgfpathlineto{\pgfqpoint{3.456630in}{1.300539in}}%
\pgfpathlineto{\pgfqpoint{3.458005in}{1.297919in}}%
\pgfpathlineto{\pgfqpoint{3.456630in}{1.292828in}}%
\pgfpathlineto{\pgfqpoint{3.454030in}{1.284671in}}%
\pgfpathlineto{\pgfqpoint{3.456630in}{1.275855in}}%
\pgfpathlineto{\pgfqpoint{3.457840in}{1.271423in}}%
\pgfpathlineto{\pgfqpoint{3.456630in}{1.268423in}}%
\pgfpathlineto{\pgfqpoint{3.451737in}{1.258175in}}%
\pgfpathlineto{\pgfqpoint{3.452023in}{1.244927in}}%
\pgfpathlineto{\pgfqpoint{3.447559in}{1.231679in}}%
\pgfpathlineto{\pgfqpoint{3.443307in}{1.218431in}}%
\pgfpathlineto{\pgfqpoint{3.441517in}{1.205183in}}%
\pgfpathlineto{\pgfqpoint{3.431347in}{1.192625in}}%
\pgfpathlineto{\pgfqpoint{3.430570in}{1.191935in}}%
\pgfpathlineto{\pgfqpoint{3.431347in}{1.186842in}}%
\pgfpathlineto{\pgfqpoint{3.432330in}{1.178687in}}%
\pgfpathlineto{\pgfqpoint{3.431347in}{1.177714in}}%
\pgfpathlineto{\pgfqpoint{3.414077in}{1.165439in}}%
\pgfpathlineto{\pgfqpoint{3.417218in}{1.152191in}}%
\pgfpathlineto{\pgfqpoint{3.406064in}{1.147328in}}%
\pgfpathlineto{\pgfqpoint{3.380781in}{1.140219in}}%
\pgfpathlineto{\pgfqpoint{3.372273in}{1.138943in}}%
\pgfpathlineto{\pgfqpoint{3.380781in}{1.126776in}}%
\pgfpathlineto{\pgfqpoint{3.381228in}{1.125695in}}%
\pgfpathlineto{\pgfqpoint{3.380781in}{1.125575in}}%
\pgfpathlineto{\pgfqpoint{3.355497in}{1.121732in}}%
\pgfpathlineto{\pgfqpoint{3.330214in}{1.120774in}}%
\pgfpathlineto{\pgfqpoint{3.304931in}{1.122742in}}%
\pgfpathlineto{\pgfqpoint{3.290483in}{1.125695in}}%
\pgfpathclose%
\pgfusepath{fill}%
\end{pgfscope}%
\begin{pgfscope}%
\pgfpathrectangle{\pgfqpoint{2.470586in}{0.476543in}}{\pgfqpoint{1.820388in}{1.656000in}}%
\pgfusepath{clip}%
\pgfsetbuttcap%
\pgfsetroundjoin%
\definecolor{currentfill}{rgb}{0.939023,0.389927,0.245521}%
\pgfsetfillcolor{currentfill}%
\pgfsetlinewidth{0.000000pt}%
\definecolor{currentstroke}{rgb}{0.000000,0.000000,0.000000}%
\pgfsetstrokecolor{currentstroke}%
\pgfsetdash{}{0pt}%
\pgfpathmoveto{\pgfqpoint{3.304931in}{1.122742in}}%
\pgfpathlineto{\pgfqpoint{3.330214in}{1.120774in}}%
\pgfpathlineto{\pgfqpoint{3.355497in}{1.121732in}}%
\pgfpathlineto{\pgfqpoint{3.380781in}{1.125575in}}%
\pgfpathlineto{\pgfqpoint{3.381228in}{1.125695in}}%
\pgfpathlineto{\pgfqpoint{3.380781in}{1.126776in}}%
\pgfpathlineto{\pgfqpoint{3.372273in}{1.138943in}}%
\pgfpathlineto{\pgfqpoint{3.380781in}{1.140219in}}%
\pgfpathlineto{\pgfqpoint{3.406064in}{1.147328in}}%
\pgfpathlineto{\pgfqpoint{3.417218in}{1.152191in}}%
\pgfpathlineto{\pgfqpoint{3.414077in}{1.165439in}}%
\pgfpathlineto{\pgfqpoint{3.431347in}{1.177714in}}%
\pgfpathlineto{\pgfqpoint{3.432330in}{1.178687in}}%
\pgfpathlineto{\pgfqpoint{3.431347in}{1.186842in}}%
\pgfpathlineto{\pgfqpoint{3.430570in}{1.191935in}}%
\pgfpathlineto{\pgfqpoint{3.431347in}{1.192625in}}%
\pgfpathlineto{\pgfqpoint{3.441517in}{1.205183in}}%
\pgfpathlineto{\pgfqpoint{3.443307in}{1.218431in}}%
\pgfpathlineto{\pgfqpoint{3.447559in}{1.231679in}}%
\pgfpathlineto{\pgfqpoint{3.452023in}{1.244927in}}%
\pgfpathlineto{\pgfqpoint{3.451737in}{1.258175in}}%
\pgfpathlineto{\pgfqpoint{3.456630in}{1.268423in}}%
\pgfpathlineto{\pgfqpoint{3.457840in}{1.271423in}}%
\pgfpathlineto{\pgfqpoint{3.456630in}{1.275855in}}%
\pgfpathlineto{\pgfqpoint{3.454030in}{1.284671in}}%
\pgfpathlineto{\pgfqpoint{3.456630in}{1.292828in}}%
\pgfpathlineto{\pgfqpoint{3.458005in}{1.297919in}}%
\pgfpathlineto{\pgfqpoint{3.456630in}{1.300539in}}%
\pgfpathlineto{\pgfqpoint{3.450504in}{1.311167in}}%
\pgfpathlineto{\pgfqpoint{3.456630in}{1.323246in}}%
\pgfpathlineto{\pgfqpoint{3.457132in}{1.324415in}}%
\pgfpathlineto{\pgfqpoint{3.456630in}{1.325465in}}%
\pgfpathlineto{\pgfqpoint{3.450130in}{1.337663in}}%
\pgfpathlineto{\pgfqpoint{3.456327in}{1.350911in}}%
\pgfpathlineto{\pgfqpoint{3.448304in}{1.364159in}}%
\pgfpathlineto{\pgfqpoint{3.452791in}{1.377407in}}%
\pgfpathlineto{\pgfqpoint{3.445362in}{1.390655in}}%
\pgfpathlineto{\pgfqpoint{3.449571in}{1.403903in}}%
\pgfpathlineto{\pgfqpoint{3.442835in}{1.417151in}}%
\pgfpathlineto{\pgfqpoint{3.444240in}{1.430399in}}%
\pgfpathlineto{\pgfqpoint{3.438947in}{1.443647in}}%
\pgfpathlineto{\pgfqpoint{3.438830in}{1.456895in}}%
\pgfpathlineto{\pgfqpoint{3.434450in}{1.470143in}}%
\pgfpathlineto{\pgfqpoint{3.432161in}{1.483391in}}%
\pgfpathlineto{\pgfqpoint{3.431347in}{1.486114in}}%
\pgfpathlineto{\pgfqpoint{3.427494in}{1.496639in}}%
\pgfpathlineto{\pgfqpoint{3.423346in}{1.509887in}}%
\pgfpathlineto{\pgfqpoint{3.422398in}{1.523135in}}%
\pgfpathlineto{\pgfqpoint{3.417455in}{1.536383in}}%
\pgfpathlineto{\pgfqpoint{3.416528in}{1.549631in}}%
\pgfpathlineto{\pgfqpoint{3.409136in}{1.562879in}}%
\pgfpathlineto{\pgfqpoint{3.408791in}{1.576127in}}%
\pgfpathlineto{\pgfqpoint{3.406064in}{1.580488in}}%
\pgfpathlineto{\pgfqpoint{3.398870in}{1.589375in}}%
\pgfpathlineto{\pgfqpoint{3.399848in}{1.602623in}}%
\pgfpathlineto{\pgfqpoint{3.387539in}{1.615871in}}%
\pgfpathlineto{\pgfqpoint{3.387747in}{1.629119in}}%
\pgfpathlineto{\pgfqpoint{3.380781in}{1.636737in}}%
\pgfpathlineto{\pgfqpoint{3.372042in}{1.642367in}}%
\pgfpathlineto{\pgfqpoint{3.373209in}{1.655615in}}%
\pgfpathlineto{\pgfqpoint{3.355497in}{1.666683in}}%
\pgfpathlineto{\pgfqpoint{3.338558in}{1.668863in}}%
\pgfpathlineto{\pgfqpoint{3.340567in}{1.682111in}}%
\pgfpathlineto{\pgfqpoint{3.330214in}{1.683383in}}%
\pgfpathlineto{\pgfqpoint{3.326781in}{1.682111in}}%
\pgfpathlineto{\pgfqpoint{3.327382in}{1.668863in}}%
\pgfpathlineto{\pgfqpoint{3.304931in}{1.659911in}}%
\pgfpathlineto{\pgfqpoint{3.300267in}{1.655615in}}%
\pgfpathlineto{\pgfqpoint{3.299310in}{1.642367in}}%
\pgfpathlineto{\pgfqpoint{3.285916in}{1.629119in}}%
\pgfpathlineto{\pgfqpoint{3.282526in}{1.615871in}}%
\pgfpathlineto{\pgfqpoint{3.279648in}{1.611946in}}%
\pgfpathlineto{\pgfqpoint{3.275230in}{1.602623in}}%
\pgfpathlineto{\pgfqpoint{3.272155in}{1.589375in}}%
\pgfpathlineto{\pgfqpoint{3.268106in}{1.576127in}}%
\pgfpathlineto{\pgfqpoint{3.261609in}{1.562879in}}%
\pgfpathlineto{\pgfqpoint{3.259514in}{1.549631in}}%
\pgfpathlineto{\pgfqpoint{3.254365in}{1.539954in}}%
\pgfpathlineto{\pgfqpoint{3.252780in}{1.536383in}}%
\pgfpathlineto{\pgfqpoint{3.253617in}{1.523135in}}%
\pgfpathlineto{\pgfqpoint{3.246099in}{1.509887in}}%
\pgfpathlineto{\pgfqpoint{3.247206in}{1.496639in}}%
\pgfpathlineto{\pgfqpoint{3.237530in}{1.483391in}}%
\pgfpathlineto{\pgfqpoint{3.240147in}{1.470143in}}%
\pgfpathlineto{\pgfqpoint{3.231568in}{1.456895in}}%
\pgfpathlineto{\pgfqpoint{3.234717in}{1.443647in}}%
\pgfpathlineto{\pgfqpoint{3.229082in}{1.434946in}}%
\pgfpathlineto{\pgfqpoint{3.226555in}{1.430399in}}%
\pgfpathlineto{\pgfqpoint{3.229082in}{1.419913in}}%
\pgfpathlineto{\pgfqpoint{3.229902in}{1.417151in}}%
\pgfpathlineto{\pgfqpoint{3.229082in}{1.415510in}}%
\pgfpathlineto{\pgfqpoint{3.224144in}{1.403903in}}%
\pgfpathlineto{\pgfqpoint{3.226621in}{1.390655in}}%
\pgfpathlineto{\pgfqpoint{3.221006in}{1.377407in}}%
\pgfpathlineto{\pgfqpoint{3.221059in}{1.364159in}}%
\pgfpathlineto{\pgfqpoint{3.217936in}{1.350911in}}%
\pgfpathlineto{\pgfqpoint{3.218402in}{1.337663in}}%
\pgfpathlineto{\pgfqpoint{3.217770in}{1.324415in}}%
\pgfpathlineto{\pgfqpoint{3.215273in}{1.311167in}}%
\pgfpathlineto{\pgfqpoint{3.216835in}{1.297919in}}%
\pgfpathlineto{\pgfqpoint{3.213066in}{1.284671in}}%
\pgfpathlineto{\pgfqpoint{3.219499in}{1.271423in}}%
\pgfpathlineto{\pgfqpoint{3.215629in}{1.258175in}}%
\pgfpathlineto{\pgfqpoint{3.224530in}{1.244927in}}%
\pgfpathlineto{\pgfqpoint{3.218641in}{1.231679in}}%
\pgfpathlineto{\pgfqpoint{3.229082in}{1.220898in}}%
\pgfpathlineto{\pgfqpoint{3.232085in}{1.218431in}}%
\pgfpathlineto{\pgfqpoint{3.229082in}{1.213007in}}%
\pgfpathlineto{\pgfqpoint{3.225284in}{1.205183in}}%
\pgfpathlineto{\pgfqpoint{3.229082in}{1.201584in}}%
\pgfpathlineto{\pgfqpoint{3.242298in}{1.191935in}}%
\pgfpathlineto{\pgfqpoint{3.232754in}{1.178687in}}%
\pgfpathlineto{\pgfqpoint{3.254365in}{1.165735in}}%
\pgfpathlineto{\pgfqpoint{3.255067in}{1.165439in}}%
\pgfpathlineto{\pgfqpoint{3.254365in}{1.163907in}}%
\pgfpathlineto{\pgfqpoint{3.250241in}{1.152191in}}%
\pgfpathlineto{\pgfqpoint{3.254365in}{1.150294in}}%
\pgfpathlineto{\pgfqpoint{3.279648in}{1.142223in}}%
\pgfpathlineto{\pgfqpoint{3.297161in}{1.138943in}}%
\pgfpathlineto{\pgfqpoint{3.290483in}{1.125695in}}%
\pgfpathlineto{\pgfqpoint{3.304931in}{1.122742in}}%
\pgfpathclose%
\pgfpathmoveto{\pgfqpoint{3.300492in}{1.152191in}}%
\pgfpathlineto{\pgfqpoint{3.304931in}{1.159919in}}%
\pgfpathlineto{\pgfqpoint{3.312336in}{1.165439in}}%
\pgfpathlineto{\pgfqpoint{3.304931in}{1.166214in}}%
\pgfpathlineto{\pgfqpoint{3.279648in}{1.172751in}}%
\pgfpathlineto{\pgfqpoint{3.265701in}{1.178687in}}%
\pgfpathlineto{\pgfqpoint{3.276035in}{1.191935in}}%
\pgfpathlineto{\pgfqpoint{3.254365in}{1.202943in}}%
\pgfpathlineto{\pgfqpoint{3.251226in}{1.205183in}}%
\pgfpathlineto{\pgfqpoint{3.254365in}{1.212057in}}%
\pgfpathlineto{\pgfqpoint{3.257922in}{1.218431in}}%
\pgfpathlineto{\pgfqpoint{3.254365in}{1.220559in}}%
\pgfpathlineto{\pgfqpoint{3.240594in}{1.231679in}}%
\pgfpathlineto{\pgfqpoint{3.246045in}{1.244927in}}%
\pgfpathlineto{\pgfqpoint{3.234870in}{1.258175in}}%
\pgfpathlineto{\pgfqpoint{3.237838in}{1.271423in}}%
\pgfpathlineto{\pgfqpoint{3.229913in}{1.284671in}}%
\pgfpathlineto{\pgfqpoint{3.233113in}{1.297919in}}%
\pgfpathlineto{\pgfqpoint{3.231309in}{1.311167in}}%
\pgfpathlineto{\pgfqpoint{3.232871in}{1.324415in}}%
\pgfpathlineto{\pgfqpoint{3.233706in}{1.337663in}}%
\pgfpathlineto{\pgfqpoint{3.231852in}{1.350911in}}%
\pgfpathlineto{\pgfqpoint{3.235748in}{1.364159in}}%
\pgfpathlineto{\pgfqpoint{3.234533in}{1.377407in}}%
\pgfpathlineto{\pgfqpoint{3.241198in}{1.390655in}}%
\pgfpathlineto{\pgfqpoint{3.237232in}{1.403903in}}%
\pgfpathlineto{\pgfqpoint{3.243877in}{1.417151in}}%
\pgfpathlineto{\pgfqpoint{3.239382in}{1.430399in}}%
\pgfpathlineto{\pgfqpoint{3.247932in}{1.443647in}}%
\pgfpathlineto{\pgfqpoint{3.244241in}{1.456895in}}%
\pgfpathlineto{\pgfqpoint{3.252644in}{1.470143in}}%
\pgfpathlineto{\pgfqpoint{3.249705in}{1.483391in}}%
\pgfpathlineto{\pgfqpoint{3.254365in}{1.489847in}}%
\pgfpathlineto{\pgfqpoint{3.260456in}{1.496639in}}%
\pgfpathlineto{\pgfqpoint{3.258877in}{1.509887in}}%
\pgfpathlineto{\pgfqpoint{3.267666in}{1.523135in}}%
\pgfpathlineto{\pgfqpoint{3.266681in}{1.536383in}}%
\pgfpathlineto{\pgfqpoint{3.273084in}{1.549631in}}%
\pgfpathlineto{\pgfqpoint{3.275487in}{1.562879in}}%
\pgfpathlineto{\pgfqpoint{3.279648in}{1.572435in}}%
\pgfpathlineto{\pgfqpoint{3.281851in}{1.576127in}}%
\pgfpathlineto{\pgfqpoint{3.288433in}{1.589375in}}%
\pgfpathlineto{\pgfqpoint{3.291596in}{1.602623in}}%
\pgfpathlineto{\pgfqpoint{3.302052in}{1.615871in}}%
\pgfpathlineto{\pgfqpoint{3.304144in}{1.629119in}}%
\pgfpathlineto{\pgfqpoint{3.304931in}{1.629886in}}%
\pgfpathlineto{\pgfqpoint{3.330214in}{1.640357in}}%
\pgfpathlineto{\pgfqpoint{3.355497in}{1.636789in}}%
\pgfpathlineto{\pgfqpoint{3.367375in}{1.629119in}}%
\pgfpathlineto{\pgfqpoint{3.367022in}{1.615871in}}%
\pgfpathlineto{\pgfqpoint{3.380781in}{1.606703in}}%
\pgfpathlineto{\pgfqpoint{3.384521in}{1.602623in}}%
\pgfpathlineto{\pgfqpoint{3.383464in}{1.589375in}}%
\pgfpathlineto{\pgfqpoint{3.393640in}{1.576127in}}%
\pgfpathlineto{\pgfqpoint{3.394309in}{1.562879in}}%
\pgfpathlineto{\pgfqpoint{3.403519in}{1.549631in}}%
\pgfpathlineto{\pgfqpoint{3.404832in}{1.536383in}}%
\pgfpathlineto{\pgfqpoint{3.406064in}{1.533436in}}%
\pgfpathlineto{\pgfqpoint{3.409447in}{1.523135in}}%
\pgfpathlineto{\pgfqpoint{3.410617in}{1.509887in}}%
\pgfpathlineto{\pgfqpoint{3.414035in}{1.496639in}}%
\pgfpathlineto{\pgfqpoint{3.419197in}{1.483391in}}%
\pgfpathlineto{\pgfqpoint{3.421063in}{1.470143in}}%
\pgfpathlineto{\pgfqpoint{3.426501in}{1.456895in}}%
\pgfpathlineto{\pgfqpoint{3.425774in}{1.443647in}}%
\pgfpathlineto{\pgfqpoint{3.431347in}{1.432505in}}%
\pgfpathlineto{\pgfqpoint{3.432203in}{1.430399in}}%
\pgfpathlineto{\pgfqpoint{3.431347in}{1.425021in}}%
\pgfpathlineto{\pgfqpoint{3.429720in}{1.417151in}}%
\pgfpathlineto{\pgfqpoint{3.431347in}{1.414523in}}%
\pgfpathlineto{\pgfqpoint{3.436742in}{1.403903in}}%
\pgfpathlineto{\pgfqpoint{3.431718in}{1.390655in}}%
\pgfpathlineto{\pgfqpoint{3.439054in}{1.377407in}}%
\pgfpathlineto{\pgfqpoint{3.433925in}{1.364159in}}%
\pgfpathlineto{\pgfqpoint{3.441692in}{1.350911in}}%
\pgfpathlineto{\pgfqpoint{3.434817in}{1.337663in}}%
\pgfpathlineto{\pgfqpoint{3.441382in}{1.324415in}}%
\pgfpathlineto{\pgfqpoint{3.434279in}{1.311167in}}%
\pgfpathlineto{\pgfqpoint{3.441116in}{1.297919in}}%
\pgfpathlineto{\pgfqpoint{3.436591in}{1.284671in}}%
\pgfpathlineto{\pgfqpoint{3.439198in}{1.271423in}}%
\pgfpathlineto{\pgfqpoint{3.432625in}{1.258175in}}%
\pgfpathlineto{\pgfqpoint{3.431347in}{1.246087in}}%
\pgfpathlineto{\pgfqpoint{3.431179in}{1.244927in}}%
\pgfpathlineto{\pgfqpoint{3.425749in}{1.231679in}}%
\pgfpathlineto{\pgfqpoint{3.417382in}{1.218431in}}%
\pgfpathlineto{\pgfqpoint{3.415381in}{1.205183in}}%
\pgfpathlineto{\pgfqpoint{3.406064in}{1.197473in}}%
\pgfpathlineto{\pgfqpoint{3.395579in}{1.191935in}}%
\pgfpathlineto{\pgfqpoint{3.399268in}{1.178687in}}%
\pgfpathlineto{\pgfqpoint{3.380781in}{1.170425in}}%
\pgfpathlineto{\pgfqpoint{3.358801in}{1.165439in}}%
\pgfpathlineto{\pgfqpoint{3.366316in}{1.152191in}}%
\pgfpathlineto{\pgfqpoint{3.355497in}{1.150407in}}%
\pgfpathlineto{\pgfqpoint{3.330214in}{1.149351in}}%
\pgfpathlineto{\pgfqpoint{3.304931in}{1.151311in}}%
\pgfpathlineto{\pgfqpoint{3.300492in}{1.152191in}}%
\pgfpathclose%
\pgfusepath{fill}%
\end{pgfscope}%
\begin{pgfscope}%
\pgfpathrectangle{\pgfqpoint{2.470586in}{0.476543in}}{\pgfqpoint{1.820388in}{1.656000in}}%
\pgfusepath{clip}%
\pgfsetbuttcap%
\pgfsetroundjoin%
\definecolor{currentfill}{rgb}{0.894425,0.296117,0.202461}%
\pgfsetfillcolor{currentfill}%
\pgfsetlinewidth{0.000000pt}%
\definecolor{currentstroke}{rgb}{0.000000,0.000000,0.000000}%
\pgfsetstrokecolor{currentstroke}%
\pgfsetdash{}{0pt}%
\pgfpathmoveto{\pgfqpoint{3.304931in}{1.151311in}}%
\pgfpathlineto{\pgfqpoint{3.330214in}{1.149351in}}%
\pgfpathlineto{\pgfqpoint{3.355497in}{1.150407in}}%
\pgfpathlineto{\pgfqpoint{3.366316in}{1.152191in}}%
\pgfpathlineto{\pgfqpoint{3.358801in}{1.165439in}}%
\pgfpathlineto{\pgfqpoint{3.380781in}{1.170425in}}%
\pgfpathlineto{\pgfqpoint{3.399268in}{1.178687in}}%
\pgfpathlineto{\pgfqpoint{3.395579in}{1.191935in}}%
\pgfpathlineto{\pgfqpoint{3.406064in}{1.197473in}}%
\pgfpathlineto{\pgfqpoint{3.415381in}{1.205183in}}%
\pgfpathlineto{\pgfqpoint{3.417382in}{1.218431in}}%
\pgfpathlineto{\pgfqpoint{3.425749in}{1.231679in}}%
\pgfpathlineto{\pgfqpoint{3.431179in}{1.244927in}}%
\pgfpathlineto{\pgfqpoint{3.431347in}{1.246087in}}%
\pgfpathlineto{\pgfqpoint{3.432625in}{1.258175in}}%
\pgfpathlineto{\pgfqpoint{3.439198in}{1.271423in}}%
\pgfpathlineto{\pgfqpoint{3.436591in}{1.284671in}}%
\pgfpathlineto{\pgfqpoint{3.441116in}{1.297919in}}%
\pgfpathlineto{\pgfqpoint{3.434279in}{1.311167in}}%
\pgfpathlineto{\pgfqpoint{3.441382in}{1.324415in}}%
\pgfpathlineto{\pgfqpoint{3.434817in}{1.337663in}}%
\pgfpathlineto{\pgfqpoint{3.441692in}{1.350911in}}%
\pgfpathlineto{\pgfqpoint{3.433925in}{1.364159in}}%
\pgfpathlineto{\pgfqpoint{3.439054in}{1.377407in}}%
\pgfpathlineto{\pgfqpoint{3.431718in}{1.390655in}}%
\pgfpathlineto{\pgfqpoint{3.436742in}{1.403903in}}%
\pgfpathlineto{\pgfqpoint{3.431347in}{1.414523in}}%
\pgfpathlineto{\pgfqpoint{3.429720in}{1.417151in}}%
\pgfpathlineto{\pgfqpoint{3.431347in}{1.425021in}}%
\pgfpathlineto{\pgfqpoint{3.432203in}{1.430399in}}%
\pgfpathlineto{\pgfqpoint{3.431347in}{1.432505in}}%
\pgfpathlineto{\pgfqpoint{3.425774in}{1.443647in}}%
\pgfpathlineto{\pgfqpoint{3.426501in}{1.456895in}}%
\pgfpathlineto{\pgfqpoint{3.421063in}{1.470143in}}%
\pgfpathlineto{\pgfqpoint{3.419197in}{1.483391in}}%
\pgfpathlineto{\pgfqpoint{3.414035in}{1.496639in}}%
\pgfpathlineto{\pgfqpoint{3.410617in}{1.509887in}}%
\pgfpathlineto{\pgfqpoint{3.409447in}{1.523135in}}%
\pgfpathlineto{\pgfqpoint{3.406064in}{1.533436in}}%
\pgfpathlineto{\pgfqpoint{3.404832in}{1.536383in}}%
\pgfpathlineto{\pgfqpoint{3.403519in}{1.549631in}}%
\pgfpathlineto{\pgfqpoint{3.394309in}{1.562879in}}%
\pgfpathlineto{\pgfqpoint{3.393640in}{1.576127in}}%
\pgfpathlineto{\pgfqpoint{3.383464in}{1.589375in}}%
\pgfpathlineto{\pgfqpoint{3.384521in}{1.602623in}}%
\pgfpathlineto{\pgfqpoint{3.380781in}{1.606703in}}%
\pgfpathlineto{\pgfqpoint{3.367022in}{1.615871in}}%
\pgfpathlineto{\pgfqpoint{3.367375in}{1.629119in}}%
\pgfpathlineto{\pgfqpoint{3.355497in}{1.636789in}}%
\pgfpathlineto{\pgfqpoint{3.330214in}{1.640357in}}%
\pgfpathlineto{\pgfqpoint{3.304931in}{1.629886in}}%
\pgfpathlineto{\pgfqpoint{3.304144in}{1.629119in}}%
\pgfpathlineto{\pgfqpoint{3.302052in}{1.615871in}}%
\pgfpathlineto{\pgfqpoint{3.291596in}{1.602623in}}%
\pgfpathlineto{\pgfqpoint{3.288433in}{1.589375in}}%
\pgfpathlineto{\pgfqpoint{3.281851in}{1.576127in}}%
\pgfpathlineto{\pgfqpoint{3.279648in}{1.572435in}}%
\pgfpathlineto{\pgfqpoint{3.275487in}{1.562879in}}%
\pgfpathlineto{\pgfqpoint{3.273084in}{1.549631in}}%
\pgfpathlineto{\pgfqpoint{3.266681in}{1.536383in}}%
\pgfpathlineto{\pgfqpoint{3.267666in}{1.523135in}}%
\pgfpathlineto{\pgfqpoint{3.258877in}{1.509887in}}%
\pgfpathlineto{\pgfqpoint{3.260456in}{1.496639in}}%
\pgfpathlineto{\pgfqpoint{3.254365in}{1.489847in}}%
\pgfpathlineto{\pgfqpoint{3.249705in}{1.483391in}}%
\pgfpathlineto{\pgfqpoint{3.252644in}{1.470143in}}%
\pgfpathlineto{\pgfqpoint{3.244241in}{1.456895in}}%
\pgfpathlineto{\pgfqpoint{3.247932in}{1.443647in}}%
\pgfpathlineto{\pgfqpoint{3.239382in}{1.430399in}}%
\pgfpathlineto{\pgfqpoint{3.243877in}{1.417151in}}%
\pgfpathlineto{\pgfqpoint{3.237232in}{1.403903in}}%
\pgfpathlineto{\pgfqpoint{3.241198in}{1.390655in}}%
\pgfpathlineto{\pgfqpoint{3.234533in}{1.377407in}}%
\pgfpathlineto{\pgfqpoint{3.235748in}{1.364159in}}%
\pgfpathlineto{\pgfqpoint{3.231852in}{1.350911in}}%
\pgfpathlineto{\pgfqpoint{3.233706in}{1.337663in}}%
\pgfpathlineto{\pgfqpoint{3.232871in}{1.324415in}}%
\pgfpathlineto{\pgfqpoint{3.231309in}{1.311167in}}%
\pgfpathlineto{\pgfqpoint{3.233113in}{1.297919in}}%
\pgfpathlineto{\pgfqpoint{3.229913in}{1.284671in}}%
\pgfpathlineto{\pgfqpoint{3.237838in}{1.271423in}}%
\pgfpathlineto{\pgfqpoint{3.234870in}{1.258175in}}%
\pgfpathlineto{\pgfqpoint{3.246045in}{1.244927in}}%
\pgfpathlineto{\pgfqpoint{3.240594in}{1.231679in}}%
\pgfpathlineto{\pgfqpoint{3.254365in}{1.220559in}}%
\pgfpathlineto{\pgfqpoint{3.257922in}{1.218431in}}%
\pgfpathlineto{\pgfqpoint{3.254365in}{1.212057in}}%
\pgfpathlineto{\pgfqpoint{3.251226in}{1.205183in}}%
\pgfpathlineto{\pgfqpoint{3.254365in}{1.202943in}}%
\pgfpathlineto{\pgfqpoint{3.276035in}{1.191935in}}%
\pgfpathlineto{\pgfqpoint{3.265701in}{1.178687in}}%
\pgfpathlineto{\pgfqpoint{3.279648in}{1.172751in}}%
\pgfpathlineto{\pgfqpoint{3.304931in}{1.166214in}}%
\pgfpathlineto{\pgfqpoint{3.312336in}{1.165439in}}%
\pgfpathlineto{\pgfqpoint{3.304931in}{1.159919in}}%
\pgfpathlineto{\pgfqpoint{3.300492in}{1.152191in}}%
\pgfpathlineto{\pgfqpoint{3.304931in}{1.151311in}}%
\pgfpathclose%
\pgfpathmoveto{\pgfqpoint{3.289636in}{1.205183in}}%
\pgfpathlineto{\pgfqpoint{3.296344in}{1.218431in}}%
\pgfpathlineto{\pgfqpoint{3.279648in}{1.224937in}}%
\pgfpathlineto{\pgfqpoint{3.268360in}{1.231679in}}%
\pgfpathlineto{\pgfqpoint{3.272735in}{1.244927in}}%
\pgfpathlineto{\pgfqpoint{3.257844in}{1.258175in}}%
\pgfpathlineto{\pgfqpoint{3.259122in}{1.271423in}}%
\pgfpathlineto{\pgfqpoint{3.254365in}{1.277550in}}%
\pgfpathlineto{\pgfqpoint{3.250104in}{1.284671in}}%
\pgfpathlineto{\pgfqpoint{3.251753in}{1.297919in}}%
\pgfpathlineto{\pgfqpoint{3.250194in}{1.311167in}}%
\pgfpathlineto{\pgfqpoint{3.250090in}{1.324415in}}%
\pgfpathlineto{\pgfqpoint{3.251106in}{1.337663in}}%
\pgfpathlineto{\pgfqpoint{3.247804in}{1.350911in}}%
\pgfpathlineto{\pgfqpoint{3.251949in}{1.364159in}}%
\pgfpathlineto{\pgfqpoint{3.249493in}{1.377407in}}%
\pgfpathlineto{\pgfqpoint{3.254365in}{1.387025in}}%
\pgfpathlineto{\pgfqpoint{3.256748in}{1.390655in}}%
\pgfpathlineto{\pgfqpoint{3.254365in}{1.395270in}}%
\pgfpathlineto{\pgfqpoint{3.251161in}{1.403903in}}%
\pgfpathlineto{\pgfqpoint{3.254365in}{1.410257in}}%
\pgfpathlineto{\pgfqpoint{3.258843in}{1.417151in}}%
\pgfpathlineto{\pgfqpoint{3.254365in}{1.425850in}}%
\pgfpathlineto{\pgfqpoint{3.252635in}{1.430399in}}%
\pgfpathlineto{\pgfqpoint{3.254365in}{1.433085in}}%
\pgfpathlineto{\pgfqpoint{3.263022in}{1.443647in}}%
\pgfpathlineto{\pgfqpoint{3.257638in}{1.456895in}}%
\pgfpathlineto{\pgfqpoint{3.268039in}{1.470143in}}%
\pgfpathlineto{\pgfqpoint{3.264000in}{1.483391in}}%
\pgfpathlineto{\pgfqpoint{3.275584in}{1.496639in}}%
\pgfpathlineto{\pgfqpoint{3.273932in}{1.509887in}}%
\pgfpathlineto{\pgfqpoint{3.279648in}{1.519236in}}%
\pgfpathlineto{\pgfqpoint{3.282951in}{1.523135in}}%
\pgfpathlineto{\pgfqpoint{3.281684in}{1.536383in}}%
\pgfpathlineto{\pgfqpoint{3.289895in}{1.549631in}}%
\pgfpathlineto{\pgfqpoint{3.294235in}{1.562879in}}%
\pgfpathlineto{\pgfqpoint{3.300866in}{1.576127in}}%
\pgfpathlineto{\pgfqpoint{3.304931in}{1.583447in}}%
\pgfpathlineto{\pgfqpoint{3.312986in}{1.589375in}}%
\pgfpathlineto{\pgfqpoint{3.315924in}{1.602623in}}%
\pgfpathlineto{\pgfqpoint{3.330214in}{1.609807in}}%
\pgfpathlineto{\pgfqpoint{3.355497in}{1.606092in}}%
\pgfpathlineto{\pgfqpoint{3.360612in}{1.602623in}}%
\pgfpathlineto{\pgfqpoint{3.359307in}{1.589375in}}%
\pgfpathlineto{\pgfqpoint{3.375051in}{1.576127in}}%
\pgfpathlineto{\pgfqpoint{3.376880in}{1.562879in}}%
\pgfpathlineto{\pgfqpoint{3.380781in}{1.559024in}}%
\pgfpathlineto{\pgfqpoint{3.386751in}{1.549631in}}%
\pgfpathlineto{\pgfqpoint{3.388357in}{1.536383in}}%
\pgfpathlineto{\pgfqpoint{3.392979in}{1.523135in}}%
\pgfpathlineto{\pgfqpoint{3.395121in}{1.509887in}}%
\pgfpathlineto{\pgfqpoint{3.398579in}{1.496639in}}%
\pgfpathlineto{\pgfqpoint{3.406064in}{1.483418in}}%
\pgfpathlineto{\pgfqpoint{3.406075in}{1.483391in}}%
\pgfpathlineto{\pgfqpoint{3.407029in}{1.470143in}}%
\pgfpathlineto{\pgfqpoint{3.412656in}{1.456895in}}%
\pgfpathlineto{\pgfqpoint{3.410970in}{1.443647in}}%
\pgfpathlineto{\pgfqpoint{3.417821in}{1.430399in}}%
\pgfpathlineto{\pgfqpoint{3.414088in}{1.417151in}}%
\pgfpathlineto{\pgfqpoint{3.422259in}{1.403903in}}%
\pgfpathlineto{\pgfqpoint{3.415110in}{1.390655in}}%
\pgfpathlineto{\pgfqpoint{3.423944in}{1.377407in}}%
\pgfpathlineto{\pgfqpoint{3.416918in}{1.364159in}}%
\pgfpathlineto{\pgfqpoint{3.426056in}{1.350911in}}%
\pgfpathlineto{\pgfqpoint{3.416836in}{1.337663in}}%
\pgfpathlineto{\pgfqpoint{3.424164in}{1.324415in}}%
\pgfpathlineto{\pgfqpoint{3.415071in}{1.311167in}}%
\pgfpathlineto{\pgfqpoint{3.422213in}{1.297919in}}%
\pgfpathlineto{\pgfqpoint{3.416327in}{1.284671in}}%
\pgfpathlineto{\pgfqpoint{3.417546in}{1.271423in}}%
\pgfpathlineto{\pgfqpoint{3.409292in}{1.258175in}}%
\pgfpathlineto{\pgfqpoint{3.406064in}{1.248032in}}%
\pgfpathlineto{\pgfqpoint{3.404513in}{1.244927in}}%
\pgfpathlineto{\pgfqpoint{3.397830in}{1.231679in}}%
\pgfpathlineto{\pgfqpoint{3.380984in}{1.218431in}}%
\pgfpathlineto{\pgfqpoint{3.380781in}{1.216775in}}%
\pgfpathlineto{\pgfqpoint{3.378455in}{1.205183in}}%
\pgfpathlineto{\pgfqpoint{3.355497in}{1.198629in}}%
\pgfpathlineto{\pgfqpoint{3.330214in}{1.197070in}}%
\pgfpathlineto{\pgfqpoint{3.304931in}{1.200333in}}%
\pgfpathlineto{\pgfqpoint{3.289636in}{1.205183in}}%
\pgfpathclose%
\pgfusepath{fill}%
\end{pgfscope}%
\begin{pgfscope}%
\pgfpathrectangle{\pgfqpoint{2.470586in}{0.476543in}}{\pgfqpoint{1.820388in}{1.656000in}}%
\pgfusepath{clip}%
\pgfsetbuttcap%
\pgfsetroundjoin%
\definecolor{currentfill}{rgb}{0.845367,0.192926,0.155094}%
\pgfsetfillcolor{currentfill}%
\pgfsetlinewidth{0.000000pt}%
\definecolor{currentstroke}{rgb}{0.000000,0.000000,0.000000}%
\pgfsetstrokecolor{currentstroke}%
\pgfsetdash{}{0pt}%
\pgfpathmoveto{\pgfqpoint{3.304931in}{1.200333in}}%
\pgfpathlineto{\pgfqpoint{3.330214in}{1.197070in}}%
\pgfpathlineto{\pgfqpoint{3.355497in}{1.198629in}}%
\pgfpathlineto{\pgfqpoint{3.378455in}{1.205183in}}%
\pgfpathlineto{\pgfqpoint{3.380781in}{1.216775in}}%
\pgfpathlineto{\pgfqpoint{3.380984in}{1.218431in}}%
\pgfpathlineto{\pgfqpoint{3.397830in}{1.231679in}}%
\pgfpathlineto{\pgfqpoint{3.404513in}{1.244927in}}%
\pgfpathlineto{\pgfqpoint{3.406064in}{1.248032in}}%
\pgfpathlineto{\pgfqpoint{3.409292in}{1.258175in}}%
\pgfpathlineto{\pgfqpoint{3.417546in}{1.271423in}}%
\pgfpathlineto{\pgfqpoint{3.416327in}{1.284671in}}%
\pgfpathlineto{\pgfqpoint{3.422213in}{1.297919in}}%
\pgfpathlineto{\pgfqpoint{3.415071in}{1.311167in}}%
\pgfpathlineto{\pgfqpoint{3.424164in}{1.324415in}}%
\pgfpathlineto{\pgfqpoint{3.416836in}{1.337663in}}%
\pgfpathlineto{\pgfqpoint{3.426056in}{1.350911in}}%
\pgfpathlineto{\pgfqpoint{3.416918in}{1.364159in}}%
\pgfpathlineto{\pgfqpoint{3.423944in}{1.377407in}}%
\pgfpathlineto{\pgfqpoint{3.415110in}{1.390655in}}%
\pgfpathlineto{\pgfqpoint{3.422259in}{1.403903in}}%
\pgfpathlineto{\pgfqpoint{3.414088in}{1.417151in}}%
\pgfpathlineto{\pgfqpoint{3.417821in}{1.430399in}}%
\pgfpathlineto{\pgfqpoint{3.410970in}{1.443647in}}%
\pgfpathlineto{\pgfqpoint{3.412656in}{1.456895in}}%
\pgfpathlineto{\pgfqpoint{3.407029in}{1.470143in}}%
\pgfpathlineto{\pgfqpoint{3.406075in}{1.483391in}}%
\pgfpathlineto{\pgfqpoint{3.406064in}{1.483418in}}%
\pgfpathlineto{\pgfqpoint{3.398579in}{1.496639in}}%
\pgfpathlineto{\pgfqpoint{3.395121in}{1.509887in}}%
\pgfpathlineto{\pgfqpoint{3.392979in}{1.523135in}}%
\pgfpathlineto{\pgfqpoint{3.388357in}{1.536383in}}%
\pgfpathlineto{\pgfqpoint{3.386751in}{1.549631in}}%
\pgfpathlineto{\pgfqpoint{3.380781in}{1.559024in}}%
\pgfpathlineto{\pgfqpoint{3.376880in}{1.562879in}}%
\pgfpathlineto{\pgfqpoint{3.375051in}{1.576127in}}%
\pgfpathlineto{\pgfqpoint{3.359307in}{1.589375in}}%
\pgfpathlineto{\pgfqpoint{3.360612in}{1.602623in}}%
\pgfpathlineto{\pgfqpoint{3.355497in}{1.606092in}}%
\pgfpathlineto{\pgfqpoint{3.330214in}{1.609807in}}%
\pgfpathlineto{\pgfqpoint{3.315924in}{1.602623in}}%
\pgfpathlineto{\pgfqpoint{3.312986in}{1.589375in}}%
\pgfpathlineto{\pgfqpoint{3.304931in}{1.583447in}}%
\pgfpathlineto{\pgfqpoint{3.300866in}{1.576127in}}%
\pgfpathlineto{\pgfqpoint{3.294235in}{1.562879in}}%
\pgfpathlineto{\pgfqpoint{3.289895in}{1.549631in}}%
\pgfpathlineto{\pgfqpoint{3.281684in}{1.536383in}}%
\pgfpathlineto{\pgfqpoint{3.282951in}{1.523135in}}%
\pgfpathlineto{\pgfqpoint{3.279648in}{1.519236in}}%
\pgfpathlineto{\pgfqpoint{3.273932in}{1.509887in}}%
\pgfpathlineto{\pgfqpoint{3.275584in}{1.496639in}}%
\pgfpathlineto{\pgfqpoint{3.264000in}{1.483391in}}%
\pgfpathlineto{\pgfqpoint{3.268039in}{1.470143in}}%
\pgfpathlineto{\pgfqpoint{3.257638in}{1.456895in}}%
\pgfpathlineto{\pgfqpoint{3.263022in}{1.443647in}}%
\pgfpathlineto{\pgfqpoint{3.254365in}{1.433085in}}%
\pgfpathlineto{\pgfqpoint{3.252635in}{1.430399in}}%
\pgfpathlineto{\pgfqpoint{3.254365in}{1.425850in}}%
\pgfpathlineto{\pgfqpoint{3.258843in}{1.417151in}}%
\pgfpathlineto{\pgfqpoint{3.254365in}{1.410257in}}%
\pgfpathlineto{\pgfqpoint{3.251161in}{1.403903in}}%
\pgfpathlineto{\pgfqpoint{3.254365in}{1.395270in}}%
\pgfpathlineto{\pgfqpoint{3.256748in}{1.390655in}}%
\pgfpathlineto{\pgfqpoint{3.254365in}{1.387025in}}%
\pgfpathlineto{\pgfqpoint{3.249493in}{1.377407in}}%
\pgfpathlineto{\pgfqpoint{3.251949in}{1.364159in}}%
\pgfpathlineto{\pgfqpoint{3.247804in}{1.350911in}}%
\pgfpathlineto{\pgfqpoint{3.251106in}{1.337663in}}%
\pgfpathlineto{\pgfqpoint{3.250090in}{1.324415in}}%
\pgfpathlineto{\pgfqpoint{3.250194in}{1.311167in}}%
\pgfpathlineto{\pgfqpoint{3.251753in}{1.297919in}}%
\pgfpathlineto{\pgfqpoint{3.250104in}{1.284671in}}%
\pgfpathlineto{\pgfqpoint{3.254365in}{1.277550in}}%
\pgfpathlineto{\pgfqpoint{3.259122in}{1.271423in}}%
\pgfpathlineto{\pgfqpoint{3.257844in}{1.258175in}}%
\pgfpathlineto{\pgfqpoint{3.272735in}{1.244927in}}%
\pgfpathlineto{\pgfqpoint{3.268360in}{1.231679in}}%
\pgfpathlineto{\pgfqpoint{3.279648in}{1.224937in}}%
\pgfpathlineto{\pgfqpoint{3.296344in}{1.218431in}}%
\pgfpathlineto{\pgfqpoint{3.289636in}{1.205183in}}%
\pgfpathlineto{\pgfqpoint{3.304931in}{1.200333in}}%
\pgfpathclose%
\pgfpathmoveto{\pgfqpoint{3.326171in}{1.231679in}}%
\pgfpathlineto{\pgfqpoint{3.323138in}{1.244927in}}%
\pgfpathlineto{\pgfqpoint{3.304931in}{1.250181in}}%
\pgfpathlineto{\pgfqpoint{3.291604in}{1.258175in}}%
\pgfpathlineto{\pgfqpoint{3.288181in}{1.271423in}}%
\pgfpathlineto{\pgfqpoint{3.279648in}{1.279146in}}%
\pgfpathlineto{\pgfqpoint{3.275495in}{1.284671in}}%
\pgfpathlineto{\pgfqpoint{3.275221in}{1.297919in}}%
\pgfpathlineto{\pgfqpoint{3.273884in}{1.311167in}}%
\pgfpathlineto{\pgfqpoint{3.271148in}{1.324415in}}%
\pgfpathlineto{\pgfqpoint{3.272939in}{1.337663in}}%
\pgfpathlineto{\pgfqpoint{3.266516in}{1.350911in}}%
\pgfpathlineto{\pgfqpoint{3.272359in}{1.364159in}}%
\pgfpathlineto{\pgfqpoint{3.267378in}{1.377407in}}%
\pgfpathlineto{\pgfqpoint{3.276134in}{1.390655in}}%
\pgfpathlineto{\pgfqpoint{3.268134in}{1.403903in}}%
\pgfpathlineto{\pgfqpoint{3.276787in}{1.417151in}}%
\pgfpathlineto{\pgfqpoint{3.269199in}{1.430399in}}%
\pgfpathlineto{\pgfqpoint{3.279648in}{1.443343in}}%
\pgfpathlineto{\pgfqpoint{3.280012in}{1.443647in}}%
\pgfpathlineto{\pgfqpoint{3.279648in}{1.444166in}}%
\pgfpathlineto{\pgfqpoint{3.273907in}{1.456895in}}%
\pgfpathlineto{\pgfqpoint{3.279648in}{1.464425in}}%
\pgfpathlineto{\pgfqpoint{3.285986in}{1.470143in}}%
\pgfpathlineto{\pgfqpoint{3.279648in}{1.483269in}}%
\pgfpathlineto{\pgfqpoint{3.279609in}{1.483391in}}%
\pgfpathlineto{\pgfqpoint{3.279648in}{1.483436in}}%
\pgfpathlineto{\pgfqpoint{3.296076in}{1.496639in}}%
\pgfpathlineto{\pgfqpoint{3.293832in}{1.509887in}}%
\pgfpathlineto{\pgfqpoint{3.303814in}{1.523135in}}%
\pgfpathlineto{\pgfqpoint{3.303225in}{1.536383in}}%
\pgfpathlineto{\pgfqpoint{3.304931in}{1.539646in}}%
\pgfpathlineto{\pgfqpoint{3.315314in}{1.549631in}}%
\pgfpathlineto{\pgfqpoint{3.328893in}{1.562879in}}%
\pgfpathlineto{\pgfqpoint{3.330214in}{1.564701in}}%
\pgfpathlineto{\pgfqpoint{3.333218in}{1.562879in}}%
\pgfpathlineto{\pgfqpoint{3.355497in}{1.556443in}}%
\pgfpathlineto{\pgfqpoint{3.361727in}{1.549631in}}%
\pgfpathlineto{\pgfqpoint{3.365710in}{1.536383in}}%
\pgfpathlineto{\pgfqpoint{3.370988in}{1.523135in}}%
\pgfpathlineto{\pgfqpoint{3.376244in}{1.509887in}}%
\pgfpathlineto{\pgfqpoint{3.379802in}{1.496639in}}%
\pgfpathlineto{\pgfqpoint{3.380781in}{1.495774in}}%
\pgfpathlineto{\pgfqpoint{3.388456in}{1.483391in}}%
\pgfpathlineto{\pgfqpoint{3.388188in}{1.470143in}}%
\pgfpathlineto{\pgfqpoint{3.396249in}{1.456895in}}%
\pgfpathlineto{\pgfqpoint{3.392495in}{1.443647in}}%
\pgfpathlineto{\pgfqpoint{3.402249in}{1.430399in}}%
\pgfpathlineto{\pgfqpoint{3.395619in}{1.417151in}}%
\pgfpathlineto{\pgfqpoint{3.406064in}{1.404734in}}%
\pgfpathlineto{\pgfqpoint{3.406574in}{1.403903in}}%
\pgfpathlineto{\pgfqpoint{3.406064in}{1.403024in}}%
\pgfpathlineto{\pgfqpoint{3.395555in}{1.390655in}}%
\pgfpathlineto{\pgfqpoint{3.406064in}{1.378943in}}%
\pgfpathlineto{\pgfqpoint{3.407077in}{1.377407in}}%
\pgfpathlineto{\pgfqpoint{3.406064in}{1.375610in}}%
\pgfpathlineto{\pgfqpoint{3.396868in}{1.364159in}}%
\pgfpathlineto{\pgfqpoint{3.406064in}{1.353818in}}%
\pgfpathlineto{\pgfqpoint{3.408005in}{1.350911in}}%
\pgfpathlineto{\pgfqpoint{3.406064in}{1.348241in}}%
\pgfpathlineto{\pgfqpoint{3.395142in}{1.337663in}}%
\pgfpathlineto{\pgfqpoint{3.403930in}{1.324415in}}%
\pgfpathlineto{\pgfqpoint{3.391298in}{1.311167in}}%
\pgfpathlineto{\pgfqpoint{3.398875in}{1.297919in}}%
\pgfpathlineto{\pgfqpoint{3.390717in}{1.284671in}}%
\pgfpathlineto{\pgfqpoint{3.388871in}{1.271423in}}%
\pgfpathlineto{\pgfqpoint{3.380781in}{1.261598in}}%
\pgfpathlineto{\pgfqpoint{3.376150in}{1.258175in}}%
\pgfpathlineto{\pgfqpoint{3.356578in}{1.244927in}}%
\pgfpathlineto{\pgfqpoint{3.355497in}{1.242928in}}%
\pgfpathlineto{\pgfqpoint{3.336551in}{1.231679in}}%
\pgfpathlineto{\pgfqpoint{3.330214in}{1.231041in}}%
\pgfpathlineto{\pgfqpoint{3.326171in}{1.231679in}}%
\pgfpathclose%
\pgfusepath{fill}%
\end{pgfscope}%
\begin{pgfscope}%
\pgfpathrectangle{\pgfqpoint{2.470586in}{0.476543in}}{\pgfqpoint{1.820388in}{1.656000in}}%
\pgfusepath{clip}%
\pgfsetbuttcap%
\pgfsetroundjoin%
\definecolor{currentfill}{rgb}{0.762399,0.110727,0.151326}%
\pgfsetfillcolor{currentfill}%
\pgfsetlinewidth{0.000000pt}%
\definecolor{currentstroke}{rgb}{0.000000,0.000000,0.000000}%
\pgfsetstrokecolor{currentstroke}%
\pgfsetdash{}{0pt}%
\pgfpathmoveto{\pgfqpoint{3.330214in}{1.231041in}}%
\pgfpathlineto{\pgfqpoint{3.336551in}{1.231679in}}%
\pgfpathlineto{\pgfqpoint{3.355497in}{1.242928in}}%
\pgfpathlineto{\pgfqpoint{3.356578in}{1.244927in}}%
\pgfpathlineto{\pgfqpoint{3.376150in}{1.258175in}}%
\pgfpathlineto{\pgfqpoint{3.380781in}{1.261598in}}%
\pgfpathlineto{\pgfqpoint{3.388871in}{1.271423in}}%
\pgfpathlineto{\pgfqpoint{3.390717in}{1.284671in}}%
\pgfpathlineto{\pgfqpoint{3.398875in}{1.297919in}}%
\pgfpathlineto{\pgfqpoint{3.391298in}{1.311167in}}%
\pgfpathlineto{\pgfqpoint{3.403930in}{1.324415in}}%
\pgfpathlineto{\pgfqpoint{3.395142in}{1.337663in}}%
\pgfpathlineto{\pgfqpoint{3.406064in}{1.348241in}}%
\pgfpathlineto{\pgfqpoint{3.408005in}{1.350911in}}%
\pgfpathlineto{\pgfqpoint{3.406064in}{1.353818in}}%
\pgfpathlineto{\pgfqpoint{3.396868in}{1.364159in}}%
\pgfpathlineto{\pgfqpoint{3.406064in}{1.375610in}}%
\pgfpathlineto{\pgfqpoint{3.407077in}{1.377407in}}%
\pgfpathlineto{\pgfqpoint{3.406064in}{1.378943in}}%
\pgfpathlineto{\pgfqpoint{3.395555in}{1.390655in}}%
\pgfpathlineto{\pgfqpoint{3.406064in}{1.403024in}}%
\pgfpathlineto{\pgfqpoint{3.406574in}{1.403903in}}%
\pgfpathlineto{\pgfqpoint{3.406064in}{1.404734in}}%
\pgfpathlineto{\pgfqpoint{3.395619in}{1.417151in}}%
\pgfpathlineto{\pgfqpoint{3.402249in}{1.430399in}}%
\pgfpathlineto{\pgfqpoint{3.392495in}{1.443647in}}%
\pgfpathlineto{\pgfqpoint{3.396249in}{1.456895in}}%
\pgfpathlineto{\pgfqpoint{3.388188in}{1.470143in}}%
\pgfpathlineto{\pgfqpoint{3.388456in}{1.483391in}}%
\pgfpathlineto{\pgfqpoint{3.380781in}{1.495774in}}%
\pgfpathlineto{\pgfqpoint{3.379802in}{1.496639in}}%
\pgfpathlineto{\pgfqpoint{3.376244in}{1.509887in}}%
\pgfpathlineto{\pgfqpoint{3.370988in}{1.523135in}}%
\pgfpathlineto{\pgfqpoint{3.365710in}{1.536383in}}%
\pgfpathlineto{\pgfqpoint{3.361727in}{1.549631in}}%
\pgfpathlineto{\pgfqpoint{3.355497in}{1.556443in}}%
\pgfpathlineto{\pgfqpoint{3.333218in}{1.562879in}}%
\pgfpathlineto{\pgfqpoint{3.330214in}{1.564701in}}%
\pgfpathlineto{\pgfqpoint{3.328893in}{1.562879in}}%
\pgfpathlineto{\pgfqpoint{3.315314in}{1.549631in}}%
\pgfpathlineto{\pgfqpoint{3.304931in}{1.539646in}}%
\pgfpathlineto{\pgfqpoint{3.303225in}{1.536383in}}%
\pgfpathlineto{\pgfqpoint{3.303814in}{1.523135in}}%
\pgfpathlineto{\pgfqpoint{3.293832in}{1.509887in}}%
\pgfpathlineto{\pgfqpoint{3.296076in}{1.496639in}}%
\pgfpathlineto{\pgfqpoint{3.279648in}{1.483436in}}%
\pgfpathlineto{\pgfqpoint{3.279609in}{1.483391in}}%
\pgfpathlineto{\pgfqpoint{3.279648in}{1.483269in}}%
\pgfpathlineto{\pgfqpoint{3.285986in}{1.470143in}}%
\pgfpathlineto{\pgfqpoint{3.279648in}{1.464425in}}%
\pgfpathlineto{\pgfqpoint{3.273907in}{1.456895in}}%
\pgfpathlineto{\pgfqpoint{3.279648in}{1.444166in}}%
\pgfpathlineto{\pgfqpoint{3.280012in}{1.443647in}}%
\pgfpathlineto{\pgfqpoint{3.279648in}{1.443343in}}%
\pgfpathlineto{\pgfqpoint{3.269199in}{1.430399in}}%
\pgfpathlineto{\pgfqpoint{3.276787in}{1.417151in}}%
\pgfpathlineto{\pgfqpoint{3.268134in}{1.403903in}}%
\pgfpathlineto{\pgfqpoint{3.276134in}{1.390655in}}%
\pgfpathlineto{\pgfqpoint{3.267378in}{1.377407in}}%
\pgfpathlineto{\pgfqpoint{3.272359in}{1.364159in}}%
\pgfpathlineto{\pgfqpoint{3.266516in}{1.350911in}}%
\pgfpathlineto{\pgfqpoint{3.272939in}{1.337663in}}%
\pgfpathlineto{\pgfqpoint{3.271148in}{1.324415in}}%
\pgfpathlineto{\pgfqpoint{3.273884in}{1.311167in}}%
\pgfpathlineto{\pgfqpoint{3.275221in}{1.297919in}}%
\pgfpathlineto{\pgfqpoint{3.275495in}{1.284671in}}%
\pgfpathlineto{\pgfqpoint{3.279648in}{1.279146in}}%
\pgfpathlineto{\pgfqpoint{3.288181in}{1.271423in}}%
\pgfpathlineto{\pgfqpoint{3.291604in}{1.258175in}}%
\pgfpathlineto{\pgfqpoint{3.304931in}{1.250181in}}%
\pgfpathlineto{\pgfqpoint{3.323138in}{1.244927in}}%
\pgfpathlineto{\pgfqpoint{3.326171in}{1.231679in}}%
\pgfpathlineto{\pgfqpoint{3.330214in}{1.231041in}}%
\pgfpathclose%
\pgfpathmoveto{\pgfqpoint{3.315687in}{1.297919in}}%
\pgfpathlineto{\pgfqpoint{3.318531in}{1.311167in}}%
\pgfpathlineto{\pgfqpoint{3.304931in}{1.317820in}}%
\pgfpathlineto{\pgfqpoint{3.300554in}{1.324415in}}%
\pgfpathlineto{\pgfqpoint{3.304741in}{1.337663in}}%
\pgfpathlineto{\pgfqpoint{3.291009in}{1.350911in}}%
\pgfpathlineto{\pgfqpoint{3.301020in}{1.364159in}}%
\pgfpathlineto{\pgfqpoint{3.290263in}{1.377407in}}%
\pgfpathlineto{\pgfqpoint{3.303731in}{1.390655in}}%
\pgfpathlineto{\pgfqpoint{3.289233in}{1.403903in}}%
\pgfpathlineto{\pgfqpoint{3.302393in}{1.417151in}}%
\pgfpathlineto{\pgfqpoint{3.289703in}{1.430399in}}%
\pgfpathlineto{\pgfqpoint{3.304931in}{1.443324in}}%
\pgfpathlineto{\pgfqpoint{3.305770in}{1.443647in}}%
\pgfpathlineto{\pgfqpoint{3.304931in}{1.444147in}}%
\pgfpathlineto{\pgfqpoint{3.295655in}{1.456895in}}%
\pgfpathlineto{\pgfqpoint{3.304931in}{1.465552in}}%
\pgfpathlineto{\pgfqpoint{3.315440in}{1.470143in}}%
\pgfpathlineto{\pgfqpoint{3.304931in}{1.480146in}}%
\pgfpathlineto{\pgfqpoint{3.303395in}{1.483391in}}%
\pgfpathlineto{\pgfqpoint{3.304931in}{1.484667in}}%
\pgfpathlineto{\pgfqpoint{3.330214in}{1.494460in}}%
\pgfpathlineto{\pgfqpoint{3.355497in}{1.490563in}}%
\pgfpathlineto{\pgfqpoint{3.364031in}{1.483391in}}%
\pgfpathlineto{\pgfqpoint{3.360145in}{1.470143in}}%
\pgfpathlineto{\pgfqpoint{3.375214in}{1.456895in}}%
\pgfpathlineto{\pgfqpoint{3.365782in}{1.443647in}}%
\pgfpathlineto{\pgfqpoint{3.380781in}{1.432592in}}%
\pgfpathlineto{\pgfqpoint{3.382441in}{1.430399in}}%
\pgfpathlineto{\pgfqpoint{3.380781in}{1.427569in}}%
\pgfpathlineto{\pgfqpoint{3.369230in}{1.417151in}}%
\pgfpathlineto{\pgfqpoint{3.380781in}{1.409156in}}%
\pgfpathlineto{\pgfqpoint{3.385162in}{1.403903in}}%
\pgfpathlineto{\pgfqpoint{3.380781in}{1.399096in}}%
\pgfpathlineto{\pgfqpoint{3.366585in}{1.390655in}}%
\pgfpathlineto{\pgfqpoint{3.380781in}{1.381248in}}%
\pgfpathlineto{\pgfqpoint{3.384167in}{1.377407in}}%
\pgfpathlineto{\pgfqpoint{3.380781in}{1.373362in}}%
\pgfpathlineto{\pgfqpoint{3.367224in}{1.364159in}}%
\pgfpathlineto{\pgfqpoint{3.380781in}{1.354396in}}%
\pgfpathlineto{\pgfqpoint{3.383734in}{1.350911in}}%
\pgfpathlineto{\pgfqpoint{3.380781in}{1.348109in}}%
\pgfpathlineto{\pgfqpoint{3.361640in}{1.337663in}}%
\pgfpathlineto{\pgfqpoint{3.373320in}{1.324415in}}%
\pgfpathlineto{\pgfqpoint{3.355497in}{1.312752in}}%
\pgfpathlineto{\pgfqpoint{3.347720in}{1.311167in}}%
\pgfpathlineto{\pgfqpoint{3.355497in}{1.306338in}}%
\pgfpathlineto{\pgfqpoint{3.360093in}{1.297919in}}%
\pgfpathlineto{\pgfqpoint{3.355497in}{1.293085in}}%
\pgfpathlineto{\pgfqpoint{3.330214in}{1.285416in}}%
\pgfpathlineto{\pgfqpoint{3.315687in}{1.297919in}}%
\pgfpathclose%
\pgfusepath{fill}%
\end{pgfscope}%
\begin{pgfscope}%
\pgfpathrectangle{\pgfqpoint{2.470586in}{0.476543in}}{\pgfqpoint{1.820388in}{1.656000in}}%
\pgfusepath{clip}%
\pgfsetbuttcap%
\pgfsetroundjoin%
\definecolor{currentfill}{rgb}{0.685506,0.036909,0.149789}%
\pgfsetfillcolor{currentfill}%
\pgfsetlinewidth{0.000000pt}%
\definecolor{currentstroke}{rgb}{0.000000,0.000000,0.000000}%
\pgfsetstrokecolor{currentstroke}%
\pgfsetdash{}{0pt}%
\pgfpathmoveto{\pgfqpoint{3.330214in}{1.285416in}}%
\pgfpathlineto{\pgfqpoint{3.355497in}{1.293085in}}%
\pgfpathlineto{\pgfqpoint{3.360093in}{1.297919in}}%
\pgfpathlineto{\pgfqpoint{3.355497in}{1.306338in}}%
\pgfpathlineto{\pgfqpoint{3.347720in}{1.311167in}}%
\pgfpathlineto{\pgfqpoint{3.355497in}{1.312752in}}%
\pgfpathlineto{\pgfqpoint{3.373320in}{1.324415in}}%
\pgfpathlineto{\pgfqpoint{3.361640in}{1.337663in}}%
\pgfpathlineto{\pgfqpoint{3.380781in}{1.348109in}}%
\pgfpathlineto{\pgfqpoint{3.383734in}{1.350911in}}%
\pgfpathlineto{\pgfqpoint{3.380781in}{1.354396in}}%
\pgfpathlineto{\pgfqpoint{3.367224in}{1.364159in}}%
\pgfpathlineto{\pgfqpoint{3.380781in}{1.373362in}}%
\pgfpathlineto{\pgfqpoint{3.384167in}{1.377407in}}%
\pgfpathlineto{\pgfqpoint{3.380781in}{1.381248in}}%
\pgfpathlineto{\pgfqpoint{3.366585in}{1.390655in}}%
\pgfpathlineto{\pgfqpoint{3.380781in}{1.399096in}}%
\pgfpathlineto{\pgfqpoint{3.385162in}{1.403903in}}%
\pgfpathlineto{\pgfqpoint{3.380781in}{1.409156in}}%
\pgfpathlineto{\pgfqpoint{3.369230in}{1.417151in}}%
\pgfpathlineto{\pgfqpoint{3.380781in}{1.427569in}}%
\pgfpathlineto{\pgfqpoint{3.382441in}{1.430399in}}%
\pgfpathlineto{\pgfqpoint{3.380781in}{1.432592in}}%
\pgfpathlineto{\pgfqpoint{3.365782in}{1.443647in}}%
\pgfpathlineto{\pgfqpoint{3.375214in}{1.456895in}}%
\pgfpathlineto{\pgfqpoint{3.360145in}{1.470143in}}%
\pgfpathlineto{\pgfqpoint{3.364031in}{1.483391in}}%
\pgfpathlineto{\pgfqpoint{3.355497in}{1.490563in}}%
\pgfpathlineto{\pgfqpoint{3.330214in}{1.494460in}}%
\pgfpathlineto{\pgfqpoint{3.304931in}{1.484667in}}%
\pgfpathlineto{\pgfqpoint{3.303395in}{1.483391in}}%
\pgfpathlineto{\pgfqpoint{3.304931in}{1.480146in}}%
\pgfpathlineto{\pgfqpoint{3.315440in}{1.470143in}}%
\pgfpathlineto{\pgfqpoint{3.304931in}{1.465552in}}%
\pgfpathlineto{\pgfqpoint{3.295655in}{1.456895in}}%
\pgfpathlineto{\pgfqpoint{3.304931in}{1.444147in}}%
\pgfpathlineto{\pgfqpoint{3.305770in}{1.443647in}}%
\pgfpathlineto{\pgfqpoint{3.304931in}{1.443324in}}%
\pgfpathlineto{\pgfqpoint{3.289703in}{1.430399in}}%
\pgfpathlineto{\pgfqpoint{3.302393in}{1.417151in}}%
\pgfpathlineto{\pgfqpoint{3.289233in}{1.403903in}}%
\pgfpathlineto{\pgfqpoint{3.303731in}{1.390655in}}%
\pgfpathlineto{\pgfqpoint{3.290263in}{1.377407in}}%
\pgfpathlineto{\pgfqpoint{3.301020in}{1.364159in}}%
\pgfpathlineto{\pgfqpoint{3.291009in}{1.350911in}}%
\pgfpathlineto{\pgfqpoint{3.304741in}{1.337663in}}%
\pgfpathlineto{\pgfqpoint{3.300554in}{1.324415in}}%
\pgfpathlineto{\pgfqpoint{3.304931in}{1.317820in}}%
\pgfpathlineto{\pgfqpoint{3.318531in}{1.311167in}}%
\pgfpathlineto{\pgfqpoint{3.315687in}{1.297919in}}%
\pgfpathlineto{\pgfqpoint{3.330214in}{1.285416in}}%
\pgfpathclose%
\pgfusepath{fill}%
\end{pgfscope}%
\begin{pgfscope}%
\pgfsetbuttcap%
\pgfsetroundjoin%
\definecolor{currentfill}{rgb}{0.000000,0.000000,0.000000}%
\pgfsetfillcolor{currentfill}%
\pgfsetlinewidth{0.803000pt}%
\definecolor{currentstroke}{rgb}{0.000000,0.000000,0.000000}%
\pgfsetstrokecolor{currentstroke}%
\pgfsetdash{}{0pt}%
\pgfsys@defobject{currentmarker}{\pgfqpoint{0.000000in}{-0.048611in}}{\pgfqpoint{0.000000in}{0.000000in}}{%
\pgfpathmoveto{\pgfqpoint{0.000000in}{0.000000in}}%
\pgfpathlineto{\pgfqpoint{0.000000in}{-0.048611in}}%
\pgfusepath{stroke,fill}%
}%
\begin{pgfscope}%
\pgfsys@transformshift{2.571719in}{0.476543in}%
\pgfsys@useobject{currentmarker}{}%
\end{pgfscope}%
\end{pgfscope}%
\begin{pgfscope}%
\definecolor{textcolor}{rgb}{0.000000,0.000000,0.000000}%
\pgfsetstrokecolor{textcolor}%
\pgfsetfillcolor{textcolor}%
\pgftext[x=2.571719in,y=0.379321in,,top]{\color{textcolor}{\rmfamily\fontsize{8.000000}{9.600000}\selectfont\catcode`\^=\active\def^{\ifmmode\sp\else\^{}\fi}\catcode`\%=\active\def%{\%}\ensuremath{-}0.10}}%
\end{pgfscope}%
\begin{pgfscope}%
\pgfsetbuttcap%
\pgfsetroundjoin%
\definecolor{currentfill}{rgb}{0.000000,0.000000,0.000000}%
\pgfsetfillcolor{currentfill}%
\pgfsetlinewidth{0.803000pt}%
\definecolor{currentstroke}{rgb}{0.000000,0.000000,0.000000}%
\pgfsetstrokecolor{currentstroke}%
\pgfsetdash{}{0pt}%
\pgfsys@defobject{currentmarker}{\pgfqpoint{0.000000in}{-0.048611in}}{\pgfqpoint{0.000000in}{0.000000in}}{%
\pgfpathmoveto{\pgfqpoint{0.000000in}{0.000000in}}%
\pgfpathlineto{\pgfqpoint{0.000000in}{-0.048611in}}%
\pgfusepath{stroke,fill}%
}%
\begin{pgfscope}%
\pgfsys@transformshift{3.077382in}{0.476543in}%
\pgfsys@useobject{currentmarker}{}%
\end{pgfscope}%
\end{pgfscope}%
\begin{pgfscope}%
\definecolor{textcolor}{rgb}{0.000000,0.000000,0.000000}%
\pgfsetstrokecolor{textcolor}%
\pgfsetfillcolor{textcolor}%
\pgftext[x=3.077382in,y=0.379321in,,top]{\color{textcolor}{\rmfamily\fontsize{8.000000}{9.600000}\selectfont\catcode`\^=\active\def^{\ifmmode\sp\else\^{}\fi}\catcode`\%=\active\def%{\%}\ensuremath{-}0.05}}%
\end{pgfscope}%
\begin{pgfscope}%
\pgfsetbuttcap%
\pgfsetroundjoin%
\definecolor{currentfill}{rgb}{0.000000,0.000000,0.000000}%
\pgfsetfillcolor{currentfill}%
\pgfsetlinewidth{0.803000pt}%
\definecolor{currentstroke}{rgb}{0.000000,0.000000,0.000000}%
\pgfsetstrokecolor{currentstroke}%
\pgfsetdash{}{0pt}%
\pgfsys@defobject{currentmarker}{\pgfqpoint{0.000000in}{-0.048611in}}{\pgfqpoint{0.000000in}{0.000000in}}{%
\pgfpathmoveto{\pgfqpoint{0.000000in}{0.000000in}}%
\pgfpathlineto{\pgfqpoint{0.000000in}{-0.048611in}}%
\pgfusepath{stroke,fill}%
}%
\begin{pgfscope}%
\pgfsys@transformshift{3.583046in}{0.476543in}%
\pgfsys@useobject{currentmarker}{}%
\end{pgfscope}%
\end{pgfscope}%
\begin{pgfscope}%
\definecolor{textcolor}{rgb}{0.000000,0.000000,0.000000}%
\pgfsetstrokecolor{textcolor}%
\pgfsetfillcolor{textcolor}%
\pgftext[x=3.583046in,y=0.379321in,,top]{\color{textcolor}{\rmfamily\fontsize{8.000000}{9.600000}\selectfont\catcode`\^=\active\def^{\ifmmode\sp\else\^{}\fi}\catcode`\%=\active\def%{\%}0.00}}%
\end{pgfscope}%
\begin{pgfscope}%
\pgfsetbuttcap%
\pgfsetroundjoin%
\definecolor{currentfill}{rgb}{0.000000,0.000000,0.000000}%
\pgfsetfillcolor{currentfill}%
\pgfsetlinewidth{0.803000pt}%
\definecolor{currentstroke}{rgb}{0.000000,0.000000,0.000000}%
\pgfsetstrokecolor{currentstroke}%
\pgfsetdash{}{0pt}%
\pgfsys@defobject{currentmarker}{\pgfqpoint{0.000000in}{-0.048611in}}{\pgfqpoint{0.000000in}{0.000000in}}{%
\pgfpathmoveto{\pgfqpoint{0.000000in}{0.000000in}}%
\pgfpathlineto{\pgfqpoint{0.000000in}{-0.048611in}}%
\pgfusepath{stroke,fill}%
}%
\begin{pgfscope}%
\pgfsys@transformshift{4.088709in}{0.476543in}%
\pgfsys@useobject{currentmarker}{}%
\end{pgfscope}%
\end{pgfscope}%
\begin{pgfscope}%
\definecolor{textcolor}{rgb}{0.000000,0.000000,0.000000}%
\pgfsetstrokecolor{textcolor}%
\pgfsetfillcolor{textcolor}%
\pgftext[x=4.088709in,y=0.379321in,,top]{\color{textcolor}{\rmfamily\fontsize{8.000000}{9.600000}\selectfont\catcode`\^=\active\def^{\ifmmode\sp\else\^{}\fi}\catcode`\%=\active\def%{\%}0.05}}%
\end{pgfscope}%
\begin{pgfscope}%
\definecolor{textcolor}{rgb}{0.000000,0.000000,0.000000}%
\pgfsetstrokecolor{textcolor}%
\pgfsetfillcolor{textcolor}%
\pgftext[x=3.380781in,y=0.225000in,,top]{\color{textcolor}{\rmfamily\fontsize{9.000000}{10.800000}\selectfont\catcode`\^=\active\def^{\ifmmode\sp\else\^{}\fi}\catcode`\%=\active\def%{\%}$\delta_2/\omega_b$}}%
\end{pgfscope}%
\begin{pgfscope}%
\pgfpathrectangle{\pgfqpoint{2.470586in}{0.476543in}}{\pgfqpoint{1.820388in}{1.656000in}}%
\pgfusepath{clip}%
\pgfsetbuttcap%
\pgfsetroundjoin%
\pgfsetlinewidth{0.501875pt}%
\definecolor{currentstroke}{rgb}{1.000000,1.000000,1.000000}%
\pgfsetstrokecolor{currentstroke}%
\pgfsetstrokeopacity{0.800000}%
\pgfsetdash{}{0pt}%
\pgfpathmoveto{\pgfqpoint{2.968676in}{2.132543in}}%
\pgfpathlineto{\pgfqpoint{2.964032in}{2.119295in}}%
\pgfpathlineto{\pgfqpoint{2.962878in}{2.106047in}}%
\pgfpathlineto{\pgfqpoint{2.957987in}{2.092799in}}%
\pgfpathlineto{\pgfqpoint{2.957687in}{2.079551in}}%
\pgfpathlineto{\pgfqpoint{2.951033in}{2.066303in}}%
\pgfpathlineto{\pgfqpoint{2.951566in}{2.053055in}}%
\pgfpathlineto{\pgfqpoint{2.950967in}{2.051991in}}%
\pgfpathlineto{\pgfqpoint{2.944977in}{2.039807in}}%
\pgfpathlineto{\pgfqpoint{2.946194in}{2.026559in}}%
\pgfpathlineto{\pgfqpoint{2.938723in}{2.013311in}}%
\pgfpathlineto{\pgfqpoint{2.940242in}{2.000063in}}%
\pgfpathlineto{\pgfqpoint{2.933291in}{1.986815in}}%
\pgfpathlineto{\pgfqpoint{2.935098in}{1.973567in}}%
\pgfpathlineto{\pgfqpoint{2.926677in}{1.960319in}}%
\pgfpathlineto{\pgfqpoint{2.925971in}{1.947071in}}%
\pgfpathlineto{\pgfqpoint{2.925683in}{1.946642in}}%
\pgfpathlineto{\pgfqpoint{2.919312in}{1.933823in}}%
\pgfpathlineto{\pgfqpoint{2.919172in}{1.920575in}}%
\pgfpathlineto{\pgfqpoint{2.915877in}{1.907327in}}%
\pgfpathlineto{\pgfqpoint{2.915670in}{1.894079in}}%
\pgfpathlineto{\pgfqpoint{2.912767in}{1.880831in}}%
\pgfpathlineto{\pgfqpoint{2.909557in}{1.867583in}}%
\pgfpathlineto{\pgfqpoint{2.908989in}{1.854335in}}%
\pgfpathlineto{\pgfqpoint{2.902615in}{1.841087in}}%
\pgfpathlineto{\pgfqpoint{2.905226in}{1.827839in}}%
\pgfpathlineto{\pgfqpoint{2.900400in}{1.822364in}}%
\pgfpathlineto{\pgfqpoint{2.895853in}{1.814591in}}%
\pgfpathlineto{\pgfqpoint{2.900400in}{1.804583in}}%
\pgfpathlineto{\pgfqpoint{2.904365in}{1.801343in}}%
\pgfpathlineto{\pgfqpoint{2.900400in}{1.798080in}}%
\pgfpathlineto{\pgfqpoint{2.894471in}{1.788095in}}%
\pgfpathlineto{\pgfqpoint{2.900400in}{1.781197in}}%
\pgfpathlineto{\pgfqpoint{2.925683in}{1.776912in}}%
\pgfpathlineto{\pgfqpoint{2.946279in}{1.788095in}}%
\pgfpathlineto{\pgfqpoint{2.950967in}{1.795919in}}%
\pgfpathlineto{\pgfqpoint{2.952644in}{1.801343in}}%
\pgfpathlineto{\pgfqpoint{2.963478in}{1.814591in}}%
\pgfpathlineto{\pgfqpoint{2.973186in}{1.827839in}}%
\pgfpathlineto{\pgfqpoint{2.976250in}{1.832010in}}%
\pgfpathlineto{\pgfqpoint{2.979800in}{1.841087in}}%
\pgfpathlineto{\pgfqpoint{2.990145in}{1.854335in}}%
\pgfpathlineto{\pgfqpoint{2.994015in}{1.867583in}}%
\pgfpathlineto{\pgfqpoint{3.001533in}{1.874590in}}%
\pgfpathlineto{\pgfqpoint{3.006157in}{1.880831in}}%
\pgfpathlineto{\pgfqpoint{3.007961in}{1.894079in}}%
\pgfpathlineto{\pgfqpoint{3.021806in}{1.907327in}}%
\pgfpathlineto{\pgfqpoint{3.023790in}{1.920575in}}%
\pgfpathlineto{\pgfqpoint{3.026816in}{1.923453in}}%
\pgfpathlineto{\pgfqpoint{3.034537in}{1.933823in}}%
\pgfpathlineto{\pgfqpoint{3.034523in}{1.947071in}}%
\pgfpathlineto{\pgfqpoint{3.044185in}{1.960319in}}%
\pgfpathlineto{\pgfqpoint{3.046184in}{1.973567in}}%
\pgfpathlineto{\pgfqpoint{3.052099in}{1.980296in}}%
\pgfpathlineto{\pgfqpoint{3.056390in}{1.986815in}}%
\pgfpathlineto{\pgfqpoint{3.059487in}{2.000063in}}%
\pgfpathlineto{\pgfqpoint{3.067081in}{2.013311in}}%
\pgfpathlineto{\pgfqpoint{3.072136in}{2.026559in}}%
\pgfpathlineto{\pgfqpoint{3.077382in}{2.037186in}}%
\pgfpathlineto{\pgfqpoint{3.078430in}{2.039807in}}%
\pgfpathlineto{\pgfqpoint{3.083412in}{2.053055in}}%
\pgfpathlineto{\pgfqpoint{3.087037in}{2.066303in}}%
\pgfpathlineto{\pgfqpoint{3.093368in}{2.079551in}}%
\pgfpathlineto{\pgfqpoint{3.096592in}{2.092799in}}%
\pgfpathlineto{\pgfqpoint{3.102666in}{2.102291in}}%
\pgfpathlineto{\pgfqpoint{3.104835in}{2.106047in}}%
\pgfpathlineto{\pgfqpoint{3.106594in}{2.119295in}}%
\pgfpathlineto{\pgfqpoint{3.113674in}{2.132543in}}%
\pgfpathmoveto{\pgfqpoint{3.562436in}{2.132543in}}%
\pgfpathlineto{\pgfqpoint{3.567036in}{2.119295in}}%
\pgfpathlineto{\pgfqpoint{3.571858in}{2.106047in}}%
\pgfpathlineto{\pgfqpoint{3.576709in}{2.092799in}}%
\pgfpathlineto{\pgfqpoint{3.582418in}{2.079551in}}%
\pgfpathlineto{\pgfqpoint{3.583046in}{2.077118in}}%
\pgfpathlineto{\pgfqpoint{3.586347in}{2.066303in}}%
\pgfpathlineto{\pgfqpoint{3.593770in}{2.053055in}}%
\pgfpathlineto{\pgfqpoint{3.596749in}{2.039807in}}%
\pgfpathlineto{\pgfqpoint{3.604467in}{2.026559in}}%
\pgfpathlineto{\pgfqpoint{3.607003in}{2.013311in}}%
\pgfpathlineto{\pgfqpoint{3.608329in}{2.011124in}}%
\pgfpathlineto{\pgfqpoint{3.616498in}{2.000063in}}%
\pgfpathlineto{\pgfqpoint{3.619127in}{1.986815in}}%
\pgfpathlineto{\pgfqpoint{3.629632in}{1.973567in}}%
\pgfpathlineto{\pgfqpoint{3.631166in}{1.960319in}}%
\pgfpathlineto{\pgfqpoint{3.633612in}{1.956617in}}%
\pgfpathlineto{\pgfqpoint{3.642224in}{1.947071in}}%
\pgfpathlineto{\pgfqpoint{3.642591in}{1.933823in}}%
\pgfpathlineto{\pgfqpoint{3.652825in}{1.920575in}}%
\pgfpathlineto{\pgfqpoint{3.654970in}{1.907327in}}%
\pgfpathlineto{\pgfqpoint{3.658895in}{1.902605in}}%
\pgfpathlineto{\pgfqpoint{3.669391in}{1.894079in}}%
\pgfpathlineto{\pgfqpoint{3.671429in}{1.880831in}}%
\pgfpathlineto{\pgfqpoint{3.683001in}{1.867583in}}%
\pgfpathlineto{\pgfqpoint{3.684179in}{1.861819in}}%
\pgfpathlineto{\pgfqpoint{3.686552in}{1.854335in}}%
\pgfpathlineto{\pgfqpoint{3.700865in}{1.841087in}}%
\pgfpathlineto{\pgfqpoint{3.705106in}{1.827839in}}%
\pgfpathlineto{\pgfqpoint{3.709462in}{1.822358in}}%
\pgfpathlineto{\pgfqpoint{3.723661in}{1.814591in}}%
\pgfpathlineto{\pgfqpoint{3.728744in}{1.801343in}}%
\pgfpathlineto{\pgfqpoint{3.734745in}{1.795422in}}%
\pgfpathlineto{\pgfqpoint{3.760028in}{1.795187in}}%
\pgfpathlineto{\pgfqpoint{3.764503in}{1.801343in}}%
\pgfpathlineto{\pgfqpoint{3.760028in}{1.812772in}}%
\pgfpathlineto{\pgfqpoint{3.757859in}{1.814591in}}%
\pgfpathlineto{\pgfqpoint{3.760028in}{1.815555in}}%
\pgfpathlineto{\pgfqpoint{3.766414in}{1.827839in}}%
\pgfpathlineto{\pgfqpoint{3.760270in}{1.841087in}}%
\pgfpathlineto{\pgfqpoint{3.764910in}{1.854335in}}%
\pgfpathlineto{\pgfqpoint{3.760028in}{1.861782in}}%
\pgfpathlineto{\pgfqpoint{3.754243in}{1.867583in}}%
\pgfpathlineto{\pgfqpoint{3.760028in}{1.879518in}}%
\pgfpathlineto{\pgfqpoint{3.760370in}{1.880831in}}%
\pgfpathlineto{\pgfqpoint{3.760028in}{1.881496in}}%
\pgfpathlineto{\pgfqpoint{3.751047in}{1.894079in}}%
\pgfpathlineto{\pgfqpoint{3.755227in}{1.907327in}}%
\pgfpathlineto{\pgfqpoint{3.747755in}{1.920575in}}%
\pgfpathlineto{\pgfqpoint{3.748416in}{1.933823in}}%
\pgfpathlineto{\pgfqpoint{3.740975in}{1.947071in}}%
\pgfpathlineto{\pgfqpoint{3.741384in}{1.960319in}}%
\pgfpathlineto{\pgfqpoint{3.737500in}{1.973567in}}%
\pgfpathlineto{\pgfqpoint{3.736398in}{1.986815in}}%
\pgfpathlineto{\pgfqpoint{3.734745in}{1.994776in}}%
\pgfpathlineto{\pgfqpoint{3.733392in}{2.000063in}}%
\pgfpathlineto{\pgfqpoint{3.730293in}{2.013311in}}%
\pgfpathlineto{\pgfqpoint{3.727700in}{2.026559in}}%
\pgfpathlineto{\pgfqpoint{3.723595in}{2.039807in}}%
\pgfpathlineto{\pgfqpoint{3.721715in}{2.053055in}}%
\pgfpathlineto{\pgfqpoint{3.716208in}{2.066303in}}%
\pgfpathlineto{\pgfqpoint{3.715884in}{2.079551in}}%
\pgfpathlineto{\pgfqpoint{3.711045in}{2.092799in}}%
\pgfpathlineto{\pgfqpoint{3.710867in}{2.106047in}}%
\pgfpathlineto{\pgfqpoint{3.709462in}{2.109225in}}%
\pgfpathlineto{\pgfqpoint{3.704636in}{2.119295in}}%
\pgfpathlineto{\pgfqpoint{3.704265in}{2.132543in}}%
\pgfpathmoveto{\pgfqpoint{3.178515in}{0.766793in}}%
\pgfpathlineto{\pgfqpoint{3.203798in}{0.765046in}}%
\pgfpathlineto{\pgfqpoint{3.229082in}{0.763733in}}%
\pgfpathlineto{\pgfqpoint{3.254365in}{0.762786in}}%
\pgfpathlineto{\pgfqpoint{3.279648in}{0.762158in}}%
\pgfpathlineto{\pgfqpoint{3.304931in}{0.761809in}}%
\pgfpathlineto{\pgfqpoint{3.330214in}{0.761711in}}%
\pgfpathlineto{\pgfqpoint{3.355497in}{0.761844in}}%
\pgfpathlineto{\pgfqpoint{3.380781in}{0.762193in}}%
\pgfpathlineto{\pgfqpoint{3.406064in}{0.762750in}}%
\pgfpathlineto{\pgfqpoint{3.431347in}{0.763508in}}%
\pgfpathlineto{\pgfqpoint{3.456630in}{0.764468in}}%
\pgfpathlineto{\pgfqpoint{3.481913in}{0.765632in}}%
\pgfpathlineto{\pgfqpoint{3.507196in}{0.767007in}}%
\pgfpathlineto{\pgfqpoint{3.522938in}{0.767999in}}%
\pgfpathlineto{\pgfqpoint{3.532480in}{0.769563in}}%
\pgfpathlineto{\pgfqpoint{3.557763in}{0.774917in}}%
\pgfpathlineto{\pgfqpoint{3.579993in}{0.781247in}}%
\pgfpathlineto{\pgfqpoint{3.583046in}{0.781531in}}%
\pgfpathlineto{\pgfqpoint{3.608329in}{0.784160in}}%
\pgfpathlineto{\pgfqpoint{3.633612in}{0.787175in}}%
\pgfpathlineto{\pgfqpoint{3.658895in}{0.790641in}}%
\pgfpathlineto{\pgfqpoint{3.683305in}{0.794495in}}%
\pgfpathlineto{\pgfqpoint{3.684179in}{0.796543in}}%
\pgfpathlineto{\pgfqpoint{3.688314in}{0.807743in}}%
\pgfpathlineto{\pgfqpoint{3.709462in}{0.812533in}}%
\pgfpathlineto{\pgfqpoint{3.734745in}{0.819389in}}%
\pgfpathlineto{\pgfqpoint{3.739942in}{0.820991in}}%
\pgfpathlineto{\pgfqpoint{3.734745in}{0.830380in}}%
\pgfpathlineto{\pgfqpoint{3.732735in}{0.834239in}}%
\pgfpathlineto{\pgfqpoint{3.734745in}{0.835002in}}%
\pgfpathlineto{\pgfqpoint{3.760028in}{0.846319in}}%
\pgfpathlineto{\pgfqpoint{3.762256in}{0.847487in}}%
\pgfpathlineto{\pgfqpoint{3.760028in}{0.852917in}}%
\pgfpathlineto{\pgfqpoint{3.756903in}{0.860735in}}%
\pgfpathlineto{\pgfqpoint{3.760028in}{0.862797in}}%
\pgfpathlineto{\pgfqpoint{3.774593in}{0.873983in}}%
\pgfpathlineto{\pgfqpoint{3.778123in}{0.887231in}}%
\pgfpathlineto{\pgfqpoint{3.782025in}{0.900479in}}%
\pgfpathlineto{\pgfqpoint{3.780512in}{0.913727in}}%
\pgfpathlineto{\pgfqpoint{3.776278in}{0.926975in}}%
\pgfpathlineto{\pgfqpoint{3.779568in}{0.940223in}}%
\pgfpathlineto{\pgfqpoint{3.771264in}{0.953471in}}%
\pgfpathlineto{\pgfqpoint{3.776252in}{0.966719in}}%
\pgfpathlineto{\pgfqpoint{3.764403in}{0.979967in}}%
\pgfpathlineto{\pgfqpoint{3.767518in}{0.993215in}}%
\pgfpathlineto{\pgfqpoint{3.760028in}{1.000760in}}%
\pgfpathlineto{\pgfqpoint{3.754758in}{1.006463in}}%
\pgfpathlineto{\pgfqpoint{3.757434in}{1.019711in}}%
\pgfpathlineto{\pgfqpoint{3.743400in}{1.032959in}}%
\pgfpathlineto{\pgfqpoint{3.746165in}{1.046207in}}%
\pgfpathlineto{\pgfqpoint{3.735829in}{1.059455in}}%
\pgfpathlineto{\pgfqpoint{3.737193in}{1.072703in}}%
\pgfpathlineto{\pgfqpoint{3.734745in}{1.076038in}}%
\pgfpathlineto{\pgfqpoint{3.727972in}{1.085951in}}%
\pgfpathlineto{\pgfqpoint{3.725216in}{1.099199in}}%
\pgfpathlineto{\pgfqpoint{3.715179in}{1.112447in}}%
\pgfpathlineto{\pgfqpoint{3.711129in}{1.125695in}}%
\pgfpathlineto{\pgfqpoint{3.709462in}{1.129424in}}%
\pgfpathlineto{\pgfqpoint{3.705522in}{1.138943in}}%
\pgfpathlineto{\pgfqpoint{3.700077in}{1.152191in}}%
\pgfpathlineto{\pgfqpoint{3.695653in}{1.165439in}}%
\pgfpathlineto{\pgfqpoint{3.688619in}{1.178687in}}%
\pgfpathlineto{\pgfqpoint{3.685098in}{1.191935in}}%
\pgfpathlineto{\pgfqpoint{3.684179in}{1.193325in}}%
\pgfpathlineto{\pgfqpoint{3.677231in}{1.205183in}}%
\pgfpathlineto{\pgfqpoint{3.675095in}{1.218431in}}%
\pgfpathlineto{\pgfqpoint{3.665710in}{1.231679in}}%
\pgfpathlineto{\pgfqpoint{3.664924in}{1.244927in}}%
\pgfpathlineto{\pgfqpoint{3.658895in}{1.253981in}}%
\pgfpathlineto{\pgfqpoint{3.656377in}{1.258175in}}%
\pgfpathlineto{\pgfqpoint{3.655113in}{1.271423in}}%
\pgfpathlineto{\pgfqpoint{3.645635in}{1.284671in}}%
\pgfpathlineto{\pgfqpoint{3.643498in}{1.297919in}}%
\pgfpathlineto{\pgfqpoint{3.633842in}{1.311167in}}%
\pgfpathlineto{\pgfqpoint{3.633612in}{1.313400in}}%
\pgfpathlineto{\pgfqpoint{3.632652in}{1.324415in}}%
\pgfpathlineto{\pgfqpoint{3.625955in}{1.337663in}}%
\pgfpathlineto{\pgfqpoint{3.623718in}{1.350911in}}%
\pgfpathlineto{\pgfqpoint{3.616828in}{1.364159in}}%
\pgfpathlineto{\pgfqpoint{3.613729in}{1.377407in}}%
\pgfpathlineto{\pgfqpoint{3.608329in}{1.388412in}}%
\pgfpathlineto{\pgfqpoint{3.607353in}{1.390655in}}%
\pgfpathlineto{\pgfqpoint{3.603401in}{1.403903in}}%
\pgfpathlineto{\pgfqpoint{3.598625in}{1.417151in}}%
\pgfpathlineto{\pgfqpoint{3.593557in}{1.430399in}}%
\pgfpathlineto{\pgfqpoint{3.589989in}{1.443647in}}%
\pgfpathlineto{\pgfqpoint{3.584251in}{1.456895in}}%
\pgfpathlineto{\pgfqpoint{3.583046in}{1.461154in}}%
\pgfpathlineto{\pgfqpoint{3.580868in}{1.470143in}}%
\pgfpathlineto{\pgfqpoint{3.574326in}{1.483391in}}%
\pgfpathlineto{\pgfqpoint{3.571912in}{1.496639in}}%
\pgfpathlineto{\pgfqpoint{3.565305in}{1.509887in}}%
\pgfpathlineto{\pgfqpoint{3.564067in}{1.523135in}}%
\pgfpathlineto{\pgfqpoint{3.557763in}{1.535380in}}%
\pgfpathlineto{\pgfqpoint{3.557313in}{1.536383in}}%
\pgfpathlineto{\pgfqpoint{3.555527in}{1.549631in}}%
\pgfpathlineto{\pgfqpoint{3.548653in}{1.562879in}}%
\pgfpathlineto{\pgfqpoint{3.546971in}{1.576127in}}%
\pgfpathlineto{\pgfqpoint{3.539830in}{1.589375in}}%
\pgfpathlineto{\pgfqpoint{3.537719in}{1.602623in}}%
\pgfpathlineto{\pgfqpoint{3.532480in}{1.612734in}}%
\pgfpathlineto{\pgfqpoint{3.531041in}{1.615871in}}%
\pgfpathlineto{\pgfqpoint{3.529028in}{1.629119in}}%
\pgfpathlineto{\pgfqpoint{3.523725in}{1.642367in}}%
\pgfpathlineto{\pgfqpoint{3.521183in}{1.655615in}}%
\pgfpathlineto{\pgfqpoint{3.515452in}{1.668863in}}%
\pgfpathlineto{\pgfqpoint{3.511691in}{1.682111in}}%
\pgfpathlineto{\pgfqpoint{3.507196in}{1.694640in}}%
\pgfpathlineto{\pgfqpoint{3.506967in}{1.695359in}}%
\pgfpathlineto{\pgfqpoint{3.503512in}{1.708607in}}%
\pgfpathlineto{\pgfqpoint{3.500122in}{1.721855in}}%
\pgfpathlineto{\pgfqpoint{3.495607in}{1.735103in}}%
\pgfpathlineto{\pgfqpoint{3.491639in}{1.748351in}}%
\pgfpathlineto{\pgfqpoint{3.485758in}{1.761599in}}%
\pgfpathlineto{\pgfqpoint{3.482688in}{1.774847in}}%
\pgfpathlineto{\pgfqpoint{3.481913in}{1.776431in}}%
\pgfpathlineto{\pgfqpoint{3.476845in}{1.788095in}}%
\pgfpathlineto{\pgfqpoint{3.474613in}{1.801343in}}%
\pgfpathlineto{\pgfqpoint{3.468541in}{1.814591in}}%
\pgfpathlineto{\pgfqpoint{3.466564in}{1.827839in}}%
\pgfpathlineto{\pgfqpoint{3.459903in}{1.841087in}}%
\pgfpathlineto{\pgfqpoint{3.457726in}{1.854335in}}%
\pgfpathlineto{\pgfqpoint{3.456630in}{1.856209in}}%
\pgfpathlineto{\pgfqpoint{3.450293in}{1.867583in}}%
\pgfpathlineto{\pgfqpoint{3.448144in}{1.880831in}}%
\pgfpathlineto{\pgfqpoint{3.441206in}{1.894079in}}%
\pgfpathlineto{\pgfqpoint{3.439123in}{1.907327in}}%
\pgfpathlineto{\pgfqpoint{3.431725in}{1.920575in}}%
\pgfpathlineto{\pgfqpoint{3.431347in}{1.922117in}}%
\pgfpathlineto{\pgfqpoint{3.428489in}{1.933823in}}%
\pgfpathlineto{\pgfqpoint{3.419872in}{1.947071in}}%
\pgfpathlineto{\pgfqpoint{3.416382in}{1.960319in}}%
\pgfpathlineto{\pgfqpoint{3.408941in}{1.973567in}}%
\pgfpathlineto{\pgfqpoint{3.406064in}{1.982990in}}%
\pgfpathlineto{\pgfqpoint{3.404714in}{1.986815in}}%
\pgfpathlineto{\pgfqpoint{3.396201in}{2.000063in}}%
\pgfpathlineto{\pgfqpoint{3.390529in}{2.013311in}}%
\pgfpathlineto{\pgfqpoint{3.381443in}{2.026559in}}%
\pgfpathlineto{\pgfqpoint{3.380781in}{2.027821in}}%
\pgfpathlineto{\pgfqpoint{3.371174in}{2.039807in}}%
\pgfpathlineto{\pgfqpoint{3.356175in}{2.053055in}}%
\pgfpathlineto{\pgfqpoint{3.355497in}{2.053713in}}%
\pgfpathlineto{\pgfqpoint{3.330214in}{2.058512in}}%
\pgfpathlineto{\pgfqpoint{3.323427in}{2.053055in}}%
\pgfpathlineto{\pgfqpoint{3.304931in}{2.043740in}}%
\pgfpathlineto{\pgfqpoint{3.301068in}{2.039807in}}%
\pgfpathlineto{\pgfqpoint{3.296318in}{2.026559in}}%
\pgfpathlineto{\pgfqpoint{3.282564in}{2.013311in}}%
\pgfpathlineto{\pgfqpoint{3.279648in}{2.003895in}}%
\pgfpathlineto{\pgfqpoint{3.278743in}{2.000063in}}%
\pgfpathlineto{\pgfqpoint{3.269309in}{1.986815in}}%
\pgfpathlineto{\pgfqpoint{3.266679in}{1.973567in}}%
\pgfpathlineto{\pgfqpoint{3.257309in}{1.960319in}}%
\pgfpathlineto{\pgfqpoint{3.254365in}{1.950924in}}%
\pgfpathlineto{\pgfqpoint{3.253222in}{1.947071in}}%
\pgfpathlineto{\pgfqpoint{3.245518in}{1.933823in}}%
\pgfpathlineto{\pgfqpoint{3.241536in}{1.920575in}}%
\pgfpathlineto{\pgfqpoint{3.236483in}{1.907327in}}%
\pgfpathlineto{\pgfqpoint{3.232322in}{1.894079in}}%
\pgfpathlineto{\pgfqpoint{3.229082in}{1.884568in}}%
\pgfpathlineto{\pgfqpoint{3.227801in}{1.880831in}}%
\pgfpathlineto{\pgfqpoint{3.222253in}{1.867583in}}%
\pgfpathlineto{\pgfqpoint{3.218982in}{1.854335in}}%
\pgfpathlineto{\pgfqpoint{3.212625in}{1.841087in}}%
\pgfpathlineto{\pgfqpoint{3.209916in}{1.827839in}}%
\pgfpathlineto{\pgfqpoint{3.203798in}{1.815620in}}%
\pgfpathlineto{\pgfqpoint{3.203237in}{1.814591in}}%
\pgfpathlineto{\pgfqpoint{3.201264in}{1.801343in}}%
\pgfpathlineto{\pgfqpoint{3.194609in}{1.788095in}}%
\pgfpathlineto{\pgfqpoint{3.192942in}{1.774847in}}%
\pgfpathlineto{\pgfqpoint{3.185814in}{1.761599in}}%
\pgfpathlineto{\pgfqpoint{3.183731in}{1.748351in}}%
\pgfpathlineto{\pgfqpoint{3.178515in}{1.736203in}}%
\pgfpathlineto{\pgfqpoint{3.177987in}{1.735103in}}%
\pgfpathlineto{\pgfqpoint{3.175848in}{1.721855in}}%
\pgfpathlineto{\pgfqpoint{3.170199in}{1.708607in}}%
\pgfpathlineto{\pgfqpoint{3.167184in}{1.695359in}}%
\pgfpathlineto{\pgfqpoint{3.161870in}{1.682111in}}%
\pgfpathlineto{\pgfqpoint{3.157979in}{1.668863in}}%
\pgfpathlineto{\pgfqpoint{3.154015in}{1.655615in}}%
\pgfpathlineto{\pgfqpoint{3.153232in}{1.653462in}}%
\pgfpathlineto{\pgfqpoint{3.148481in}{1.642367in}}%
\pgfpathlineto{\pgfqpoint{3.145014in}{1.629119in}}%
\pgfpathlineto{\pgfqpoint{3.139712in}{1.615871in}}%
\pgfpathlineto{\pgfqpoint{3.137857in}{1.602623in}}%
\pgfpathlineto{\pgfqpoint{3.131984in}{1.589375in}}%
\pgfpathlineto{\pgfqpoint{3.130124in}{1.576127in}}%
\pgfpathlineto{\pgfqpoint{3.127949in}{1.572102in}}%
\pgfpathlineto{\pgfqpoint{3.122218in}{1.562879in}}%
\pgfpathlineto{\pgfqpoint{3.120642in}{1.549631in}}%
\pgfpathlineto{\pgfqpoint{3.113769in}{1.536383in}}%
\pgfpathlineto{\pgfqpoint{3.112389in}{1.523135in}}%
\pgfpathlineto{\pgfqpoint{3.105313in}{1.509887in}}%
\pgfpathlineto{\pgfqpoint{3.103398in}{1.496639in}}%
\pgfpathlineto{\pgfqpoint{3.102666in}{1.495233in}}%
\pgfpathlineto{\pgfqpoint{3.095748in}{1.483391in}}%
\pgfpathlineto{\pgfqpoint{3.093213in}{1.470143in}}%
\pgfpathlineto{\pgfqpoint{3.087050in}{1.456895in}}%
\pgfpathlineto{\pgfqpoint{3.083783in}{1.443647in}}%
\pgfpathlineto{\pgfqpoint{3.079017in}{1.430399in}}%
\pgfpathlineto{\pgfqpoint{3.077382in}{1.424114in}}%
\pgfpathlineto{\pgfqpoint{3.075225in}{1.417151in}}%
\pgfpathlineto{\pgfqpoint{3.070764in}{1.403903in}}%
\pgfpathlineto{\pgfqpoint{3.064566in}{1.390655in}}%
\pgfpathlineto{\pgfqpoint{3.060519in}{1.377407in}}%
\pgfpathlineto{\pgfqpoint{3.053172in}{1.364159in}}%
\pgfpathlineto{\pgfqpoint{3.052099in}{1.356397in}}%
\pgfpathlineto{\pgfqpoint{3.051228in}{1.350911in}}%
\pgfpathlineto{\pgfqpoint{3.043486in}{1.337663in}}%
\pgfpathlineto{\pgfqpoint{3.041835in}{1.324415in}}%
\pgfpathlineto{\pgfqpoint{3.032750in}{1.311167in}}%
\pgfpathlineto{\pgfqpoint{3.031728in}{1.297919in}}%
\pgfpathlineto{\pgfqpoint{3.026816in}{1.290744in}}%
\pgfpathlineto{\pgfqpoint{3.022178in}{1.284671in}}%
\pgfpathlineto{\pgfqpoint{3.021687in}{1.271423in}}%
\pgfpathlineto{\pgfqpoint{3.011987in}{1.258175in}}%
\pgfpathlineto{\pgfqpoint{3.011041in}{1.244927in}}%
\pgfpathlineto{\pgfqpoint{3.002214in}{1.231679in}}%
\pgfpathlineto{\pgfqpoint{3.001533in}{1.224808in}}%
\pgfpathlineto{\pgfqpoint{3.000796in}{1.218431in}}%
\pgfpathlineto{\pgfqpoint{2.990734in}{1.205183in}}%
\pgfpathlineto{\pgfqpoint{2.986573in}{1.191935in}}%
\pgfpathlineto{\pgfqpoint{2.978462in}{1.178687in}}%
\pgfpathlineto{\pgfqpoint{2.976250in}{1.171344in}}%
\pgfpathlineto{\pgfqpoint{2.974204in}{1.165439in}}%
\pgfpathlineto{\pgfqpoint{2.969696in}{1.152191in}}%
\pgfpathlineto{\pgfqpoint{2.964513in}{1.138943in}}%
\pgfpathlineto{\pgfqpoint{2.959808in}{1.125695in}}%
\pgfpathlineto{\pgfqpoint{2.951369in}{1.112447in}}%
\pgfpathlineto{\pgfqpoint{2.950967in}{1.110131in}}%
\pgfpathlineto{\pgfqpoint{2.948904in}{1.099199in}}%
\pgfpathlineto{\pgfqpoint{2.938181in}{1.085951in}}%
\pgfpathlineto{\pgfqpoint{2.937947in}{1.072703in}}%
\pgfpathlineto{\pgfqpoint{2.925748in}{1.059455in}}%
\pgfpathlineto{\pgfqpoint{2.927587in}{1.046207in}}%
\pgfpathlineto{\pgfqpoint{2.925683in}{1.044217in}}%
\pgfpathlineto{\pgfqpoint{2.914106in}{1.032959in}}%
\pgfpathlineto{\pgfqpoint{2.917558in}{1.019711in}}%
\pgfpathlineto{\pgfqpoint{2.901061in}{1.006463in}}%
\pgfpathlineto{\pgfqpoint{2.905234in}{0.993215in}}%
\pgfpathlineto{\pgfqpoint{2.900400in}{0.988676in}}%
\pgfpathlineto{\pgfqpoint{2.890677in}{0.979967in}}%
\pgfpathlineto{\pgfqpoint{2.899461in}{0.966719in}}%
\pgfpathlineto{\pgfqpoint{2.887673in}{0.953471in}}%
\pgfpathlineto{\pgfqpoint{2.895541in}{0.940223in}}%
\pgfpathlineto{\pgfqpoint{2.879898in}{0.926975in}}%
\pgfpathlineto{\pgfqpoint{2.886517in}{0.913727in}}%
\pgfpathlineto{\pgfqpoint{2.875773in}{0.900479in}}%
\pgfpathlineto{\pgfqpoint{2.885735in}{0.887231in}}%
\pgfpathlineto{\pgfqpoint{2.880757in}{0.873983in}}%
\pgfpathlineto{\pgfqpoint{2.891545in}{0.860735in}}%
\pgfpathlineto{\pgfqpoint{2.895418in}{0.847487in}}%
\pgfpathlineto{\pgfqpoint{2.900400in}{0.842605in}}%
\pgfpathlineto{\pgfqpoint{2.910685in}{0.834239in}}%
\pgfpathlineto{\pgfqpoint{2.925683in}{0.821265in}}%
\pgfpathlineto{\pgfqpoint{2.926026in}{0.820991in}}%
\pgfpathlineto{\pgfqpoint{2.948790in}{0.807743in}}%
\pgfpathlineto{\pgfqpoint{2.950967in}{0.806900in}}%
\pgfpathlineto{\pgfqpoint{2.976250in}{0.797764in}}%
\pgfpathlineto{\pgfqpoint{2.986058in}{0.794495in}}%
\pgfpathlineto{\pgfqpoint{3.001533in}{0.790008in}}%
\pgfpathlineto{\pgfqpoint{3.026816in}{0.784292in}}%
\pgfpathlineto{\pgfqpoint{3.044174in}{0.781247in}}%
\pgfpathlineto{\pgfqpoint{3.052099in}{0.780150in}}%
\pgfpathlineto{\pgfqpoint{3.077382in}{0.776942in}}%
\pgfpathlineto{\pgfqpoint{3.102666in}{0.774022in}}%
\pgfpathlineto{\pgfqpoint{3.127949in}{0.771371in}}%
\pgfpathlineto{\pgfqpoint{3.153232in}{0.768972in}}%
\pgfpathlineto{\pgfqpoint{3.164697in}{0.767999in}}%
\pgfpathlineto{\pgfqpoint{3.178515in}{0.766793in}}%
\pgfpathclose%
\pgfpathmoveto{\pgfqpoint{2.900400in}{1.757219in}}%
\pgfpathlineto{\pgfqpoint{2.925683in}{1.760469in}}%
\pgfpathlineto{\pgfqpoint{2.927144in}{1.761599in}}%
\pgfpathlineto{\pgfqpoint{2.925683in}{1.766459in}}%
\pgfpathlineto{\pgfqpoint{2.900400in}{1.767326in}}%
\pgfpathlineto{\pgfqpoint{2.893950in}{1.761599in}}%
\pgfpathlineto{\pgfqpoint{2.900400in}{1.757219in}}%
\pgfpathclose%
\pgfusepath{stroke}%
\end{pgfscope}%
\begin{pgfscope}%
\pgfpathrectangle{\pgfqpoint{2.470586in}{0.476543in}}{\pgfqpoint{1.820388in}{1.656000in}}%
\pgfusepath{clip}%
\pgfsetbuttcap%
\pgfsetroundjoin%
\pgfsetlinewidth{0.501875pt}%
\definecolor{currentstroke}{rgb}{1.000000,1.000000,1.000000}%
\pgfsetstrokecolor{currentstroke}%
\pgfsetstrokeopacity{0.800000}%
\pgfsetdash{}{0pt}%
\pgfpathmoveto{\pgfqpoint{3.279648in}{0.899860in}}%
\pgfpathlineto{\pgfqpoint{3.304931in}{0.898971in}}%
\pgfpathlineto{\pgfqpoint{3.330214in}{0.898711in}}%
\pgfpathlineto{\pgfqpoint{3.355497in}{0.899084in}}%
\pgfpathlineto{\pgfqpoint{3.380781in}{0.900125in}}%
\pgfpathlineto{\pgfqpoint{3.385969in}{0.900479in}}%
\pgfpathlineto{\pgfqpoint{3.406064in}{0.906613in}}%
\pgfpathlineto{\pgfqpoint{3.423217in}{0.913727in}}%
\pgfpathlineto{\pgfqpoint{3.431347in}{0.914651in}}%
\pgfpathlineto{\pgfqpoint{3.456630in}{0.918639in}}%
\pgfpathlineto{\pgfqpoint{3.481913in}{0.924171in}}%
\pgfpathlineto{\pgfqpoint{3.491841in}{0.926975in}}%
\pgfpathlineto{\pgfqpoint{3.507196in}{0.937169in}}%
\pgfpathlineto{\pgfqpoint{3.511234in}{0.940223in}}%
\pgfpathlineto{\pgfqpoint{3.532480in}{0.949562in}}%
\pgfpathlineto{\pgfqpoint{3.539310in}{0.953471in}}%
\pgfpathlineto{\pgfqpoint{3.557763in}{0.965182in}}%
\pgfpathlineto{\pgfqpoint{3.559918in}{0.966719in}}%
\pgfpathlineto{\pgfqpoint{3.565308in}{0.979967in}}%
\pgfpathlineto{\pgfqpoint{3.580569in}{0.993215in}}%
\pgfpathlineto{\pgfqpoint{3.578963in}{1.006463in}}%
\pgfpathlineto{\pgfqpoint{3.583046in}{1.010075in}}%
\pgfpathlineto{\pgfqpoint{3.592848in}{1.019711in}}%
\pgfpathlineto{\pgfqpoint{3.586317in}{1.032959in}}%
\pgfpathlineto{\pgfqpoint{3.599466in}{1.046207in}}%
\pgfpathlineto{\pgfqpoint{3.592721in}{1.059455in}}%
\pgfpathlineto{\pgfqpoint{3.604467in}{1.072703in}}%
\pgfpathlineto{\pgfqpoint{3.598801in}{1.085951in}}%
\pgfpathlineto{\pgfqpoint{3.604659in}{1.099199in}}%
\pgfpathlineto{\pgfqpoint{3.596643in}{1.112447in}}%
\pgfpathlineto{\pgfqpoint{3.600536in}{1.125695in}}%
\pgfpathlineto{\pgfqpoint{3.595926in}{1.138943in}}%
\pgfpathlineto{\pgfqpoint{3.597073in}{1.152191in}}%
\pgfpathlineto{\pgfqpoint{3.593821in}{1.165439in}}%
\pgfpathlineto{\pgfqpoint{3.591785in}{1.178687in}}%
\pgfpathlineto{\pgfqpoint{3.589805in}{1.191935in}}%
\pgfpathlineto{\pgfqpoint{3.586523in}{1.205183in}}%
\pgfpathlineto{\pgfqpoint{3.586056in}{1.218431in}}%
\pgfpathlineto{\pgfqpoint{3.583046in}{1.224990in}}%
\pgfpathlineto{\pgfqpoint{3.580156in}{1.231679in}}%
\pgfpathlineto{\pgfqpoint{3.581298in}{1.244927in}}%
\pgfpathlineto{\pgfqpoint{3.575160in}{1.258175in}}%
\pgfpathlineto{\pgfqpoint{3.576259in}{1.271423in}}%
\pgfpathlineto{\pgfqpoint{3.568639in}{1.284671in}}%
\pgfpathlineto{\pgfqpoint{3.568655in}{1.297919in}}%
\pgfpathlineto{\pgfqpoint{3.560285in}{1.311167in}}%
\pgfpathlineto{\pgfqpoint{3.561639in}{1.324415in}}%
\pgfpathlineto{\pgfqpoint{3.557763in}{1.331995in}}%
\pgfpathlineto{\pgfqpoint{3.554948in}{1.337663in}}%
\pgfpathlineto{\pgfqpoint{3.555215in}{1.350911in}}%
\pgfpathlineto{\pgfqpoint{3.548641in}{1.364159in}}%
\pgfpathlineto{\pgfqpoint{3.548070in}{1.377407in}}%
\pgfpathlineto{\pgfqpoint{3.542175in}{1.390655in}}%
\pgfpathlineto{\pgfqpoint{3.540724in}{1.403903in}}%
\pgfpathlineto{\pgfqpoint{3.535592in}{1.417151in}}%
\pgfpathlineto{\pgfqpoint{3.532480in}{1.430217in}}%
\pgfpathlineto{\pgfqpoint{3.532439in}{1.430399in}}%
\pgfpathlineto{\pgfqpoint{3.528913in}{1.443647in}}%
\pgfpathlineto{\pgfqpoint{3.525246in}{1.456895in}}%
\pgfpathlineto{\pgfqpoint{3.521986in}{1.470143in}}%
\pgfpathlineto{\pgfqpoint{3.516752in}{1.483391in}}%
\pgfpathlineto{\pgfqpoint{3.514221in}{1.496639in}}%
\pgfpathlineto{\pgfqpoint{3.508584in}{1.509887in}}%
\pgfpathlineto{\pgfqpoint{3.507920in}{1.523135in}}%
\pgfpathlineto{\pgfqpoint{3.507196in}{1.524752in}}%
\pgfpathlineto{\pgfqpoint{3.502203in}{1.536383in}}%
\pgfpathlineto{\pgfqpoint{3.500970in}{1.549631in}}%
\pgfpathlineto{\pgfqpoint{3.494145in}{1.562879in}}%
\pgfpathlineto{\pgfqpoint{3.493242in}{1.576127in}}%
\pgfpathlineto{\pgfqpoint{3.486135in}{1.589375in}}%
\pgfpathlineto{\pgfqpoint{3.485240in}{1.602623in}}%
\pgfpathlineto{\pgfqpoint{3.481913in}{1.608805in}}%
\pgfpathlineto{\pgfqpoint{3.478084in}{1.615871in}}%
\pgfpathlineto{\pgfqpoint{3.476828in}{1.629119in}}%
\pgfpathlineto{\pgfqpoint{3.470503in}{1.642367in}}%
\pgfpathlineto{\pgfqpoint{3.469022in}{1.655615in}}%
\pgfpathlineto{\pgfqpoint{3.462450in}{1.668863in}}%
\pgfpathlineto{\pgfqpoint{3.460047in}{1.682111in}}%
\pgfpathlineto{\pgfqpoint{3.456630in}{1.690092in}}%
\pgfpathlineto{\pgfqpoint{3.454232in}{1.695359in}}%
\pgfpathlineto{\pgfqpoint{3.451561in}{1.708607in}}%
\pgfpathlineto{\pgfqpoint{3.446956in}{1.721855in}}%
\pgfpathlineto{\pgfqpoint{3.443016in}{1.735103in}}%
\pgfpathlineto{\pgfqpoint{3.437745in}{1.748351in}}%
\pgfpathlineto{\pgfqpoint{3.432437in}{1.761599in}}%
\pgfpathlineto{\pgfqpoint{3.431347in}{1.765194in}}%
\pgfpathlineto{\pgfqpoint{3.428019in}{1.774847in}}%
\pgfpathlineto{\pgfqpoint{3.421503in}{1.788095in}}%
\pgfpathlineto{\pgfqpoint{3.417914in}{1.801343in}}%
\pgfpathlineto{\pgfqpoint{3.410497in}{1.814591in}}%
\pgfpathlineto{\pgfqpoint{3.407376in}{1.827839in}}%
\pgfpathlineto{\pgfqpoint{3.406064in}{1.830084in}}%
\pgfpathlineto{\pgfqpoint{3.397986in}{1.841087in}}%
\pgfpathlineto{\pgfqpoint{3.394606in}{1.854335in}}%
\pgfpathlineto{\pgfqpoint{3.383150in}{1.867583in}}%
\pgfpathlineto{\pgfqpoint{3.380781in}{1.877391in}}%
\pgfpathlineto{\pgfqpoint{3.379385in}{1.880831in}}%
\pgfpathlineto{\pgfqpoint{3.360013in}{1.894079in}}%
\pgfpathlineto{\pgfqpoint{3.355963in}{1.907327in}}%
\pgfpathlineto{\pgfqpoint{3.355497in}{1.907650in}}%
\pgfpathlineto{\pgfqpoint{3.330214in}{1.910929in}}%
\pgfpathlineto{\pgfqpoint{3.322721in}{1.907327in}}%
\pgfpathlineto{\pgfqpoint{3.314430in}{1.894079in}}%
\pgfpathlineto{\pgfqpoint{3.304931in}{1.889470in}}%
\pgfpathlineto{\pgfqpoint{3.297235in}{1.880831in}}%
\pgfpathlineto{\pgfqpoint{3.290287in}{1.867583in}}%
\pgfpathlineto{\pgfqpoint{3.281376in}{1.854335in}}%
\pgfpathlineto{\pgfqpoint{3.279648in}{1.851791in}}%
\pgfpathlineto{\pgfqpoint{3.274199in}{1.841087in}}%
\pgfpathlineto{\pgfqpoint{3.269306in}{1.827839in}}%
\pgfpathlineto{\pgfqpoint{3.260770in}{1.814591in}}%
\pgfpathlineto{\pgfqpoint{3.257438in}{1.801343in}}%
\pgfpathlineto{\pgfqpoint{3.254365in}{1.796661in}}%
\pgfpathlineto{\pgfqpoint{3.249552in}{1.788095in}}%
\pgfpathlineto{\pgfqpoint{3.247829in}{1.774847in}}%
\pgfpathlineto{\pgfqpoint{3.239185in}{1.761599in}}%
\pgfpathlineto{\pgfqpoint{3.237389in}{1.748351in}}%
\pgfpathlineto{\pgfqpoint{3.229560in}{1.735103in}}%
\pgfpathlineto{\pgfqpoint{3.229082in}{1.728308in}}%
\pgfpathlineto{\pgfqpoint{3.228643in}{1.721855in}}%
\pgfpathlineto{\pgfqpoint{3.221984in}{1.708607in}}%
\pgfpathlineto{\pgfqpoint{3.220290in}{1.695359in}}%
\pgfpathlineto{\pgfqpoint{3.213430in}{1.682111in}}%
\pgfpathlineto{\pgfqpoint{3.210705in}{1.668863in}}%
\pgfpathlineto{\pgfqpoint{3.205313in}{1.655615in}}%
\pgfpathlineto{\pgfqpoint{3.203798in}{1.650383in}}%
\pgfpathlineto{\pgfqpoint{3.201421in}{1.642367in}}%
\pgfpathlineto{\pgfqpoint{3.197045in}{1.629119in}}%
\pgfpathlineto{\pgfqpoint{3.192675in}{1.615871in}}%
\pgfpathlineto{\pgfqpoint{3.190228in}{1.602623in}}%
\pgfpathlineto{\pgfqpoint{3.185208in}{1.589375in}}%
\pgfpathlineto{\pgfqpoint{3.183305in}{1.576127in}}%
\pgfpathlineto{\pgfqpoint{3.178515in}{1.566768in}}%
\pgfpathlineto{\pgfqpoint{3.176452in}{1.562879in}}%
\pgfpathlineto{\pgfqpoint{3.175348in}{1.549631in}}%
\pgfpathlineto{\pgfqpoint{3.168859in}{1.536383in}}%
\pgfpathlineto{\pgfqpoint{3.168522in}{1.523135in}}%
\pgfpathlineto{\pgfqpoint{3.161221in}{1.509887in}}%
\pgfpathlineto{\pgfqpoint{3.160490in}{1.496639in}}%
\pgfpathlineto{\pgfqpoint{3.153232in}{1.483826in}}%
\pgfpathlineto{\pgfqpoint{3.152980in}{1.483391in}}%
\pgfpathlineto{\pgfqpoint{3.152441in}{1.470143in}}%
\pgfpathlineto{\pgfqpoint{3.145806in}{1.456895in}}%
\pgfpathlineto{\pgfqpoint{3.144696in}{1.443647in}}%
\pgfpathlineto{\pgfqpoint{3.138999in}{1.430399in}}%
\pgfpathlineto{\pgfqpoint{3.137684in}{1.417151in}}%
\pgfpathlineto{\pgfqpoint{3.133343in}{1.403903in}}%
\pgfpathlineto{\pgfqpoint{3.130244in}{1.390655in}}%
\pgfpathlineto{\pgfqpoint{3.127949in}{1.383469in}}%
\pgfpathlineto{\pgfqpoint{3.125959in}{1.377407in}}%
\pgfpathlineto{\pgfqpoint{3.120850in}{1.364159in}}%
\pgfpathlineto{\pgfqpoint{3.118854in}{1.350911in}}%
\pgfpathlineto{\pgfqpoint{3.113886in}{1.337663in}}%
\pgfpathlineto{\pgfqpoint{3.112964in}{1.324415in}}%
\pgfpathlineto{\pgfqpoint{3.105942in}{1.311167in}}%
\pgfpathlineto{\pgfqpoint{3.106152in}{1.297919in}}%
\pgfpathlineto{\pgfqpoint{3.102666in}{1.291634in}}%
\pgfpathlineto{\pgfqpoint{3.098607in}{1.284671in}}%
\pgfpathlineto{\pgfqpoint{3.100689in}{1.271423in}}%
\pgfpathlineto{\pgfqpoint{3.092642in}{1.258175in}}%
\pgfpathlineto{\pgfqpoint{3.094944in}{1.244927in}}%
\pgfpathlineto{\pgfqpoint{3.086650in}{1.231679in}}%
\pgfpathlineto{\pgfqpoint{3.089591in}{1.218431in}}%
\pgfpathlineto{\pgfqpoint{3.080995in}{1.205183in}}%
\pgfpathlineto{\pgfqpoint{3.082314in}{1.191935in}}%
\pgfpathlineto{\pgfqpoint{3.077382in}{1.184004in}}%
\pgfpathlineto{\pgfqpoint{3.074046in}{1.178687in}}%
\pgfpathlineto{\pgfqpoint{3.075913in}{1.165439in}}%
\pgfpathlineto{\pgfqpoint{3.071609in}{1.152191in}}%
\pgfpathlineto{\pgfqpoint{3.073873in}{1.138943in}}%
\pgfpathlineto{\pgfqpoint{3.070013in}{1.125695in}}%
\pgfpathlineto{\pgfqpoint{3.069508in}{1.112447in}}%
\pgfpathlineto{\pgfqpoint{3.068584in}{1.099199in}}%
\pgfpathlineto{\pgfqpoint{3.067570in}{1.085951in}}%
\pgfpathlineto{\pgfqpoint{3.070600in}{1.072703in}}%
\pgfpathlineto{\pgfqpoint{3.067767in}{1.059455in}}%
\pgfpathlineto{\pgfqpoint{3.074253in}{1.046207in}}%
\pgfpathlineto{\pgfqpoint{3.069916in}{1.032959in}}%
\pgfpathlineto{\pgfqpoint{3.077382in}{1.024439in}}%
\pgfpathlineto{\pgfqpoint{3.081899in}{1.019711in}}%
\pgfpathlineto{\pgfqpoint{3.077382in}{1.009837in}}%
\pgfpathlineto{\pgfqpoint{3.075761in}{1.006463in}}%
\pgfpathlineto{\pgfqpoint{3.077382in}{1.005041in}}%
\pgfpathlineto{\pgfqpoint{3.092317in}{0.993215in}}%
\pgfpathlineto{\pgfqpoint{3.087131in}{0.979967in}}%
\pgfpathlineto{\pgfqpoint{3.102666in}{0.972592in}}%
\pgfpathlineto{\pgfqpoint{3.116860in}{0.966719in}}%
\pgfpathlineto{\pgfqpoint{3.117653in}{0.953471in}}%
\pgfpathlineto{\pgfqpoint{3.127949in}{0.950027in}}%
\pgfpathlineto{\pgfqpoint{3.153232in}{0.942924in}}%
\pgfpathlineto{\pgfqpoint{3.164807in}{0.940223in}}%
\pgfpathlineto{\pgfqpoint{3.163210in}{0.926975in}}%
\pgfpathlineto{\pgfqpoint{3.178515in}{0.923579in}}%
\pgfpathlineto{\pgfqpoint{3.203798in}{0.919035in}}%
\pgfpathlineto{\pgfqpoint{3.229082in}{0.915473in}}%
\pgfpathlineto{\pgfqpoint{3.245257in}{0.913727in}}%
\pgfpathlineto{\pgfqpoint{3.254365in}{0.907180in}}%
\pgfpathlineto{\pgfqpoint{3.269367in}{0.900479in}}%
\pgfpathlineto{\pgfqpoint{3.279648in}{0.899860in}}%
\pgfpathclose%
\pgfusepath{stroke}%
\end{pgfscope}%
\begin{pgfscope}%
\pgfpathrectangle{\pgfqpoint{2.470586in}{0.476543in}}{\pgfqpoint{1.820388in}{1.656000in}}%
\pgfusepath{clip}%
\pgfsetbuttcap%
\pgfsetroundjoin%
\pgfsetlinewidth{0.501875pt}%
\definecolor{currentstroke}{rgb}{1.000000,1.000000,1.000000}%
\pgfsetstrokecolor{currentstroke}%
\pgfsetstrokeopacity{0.800000}%
\pgfsetdash{}{0pt}%
\pgfpathmoveto{\pgfqpoint{3.279648in}{1.044327in}}%
\pgfpathlineto{\pgfqpoint{3.304931in}{1.039474in}}%
\pgfpathlineto{\pgfqpoint{3.330214in}{1.037983in}}%
\pgfpathlineto{\pgfqpoint{3.355497in}{1.039046in}}%
\pgfpathlineto{\pgfqpoint{3.380781in}{1.042273in}}%
\pgfpathlineto{\pgfqpoint{3.399595in}{1.046207in}}%
\pgfpathlineto{\pgfqpoint{3.406064in}{1.055215in}}%
\pgfpathlineto{\pgfqpoint{3.407963in}{1.059455in}}%
\pgfpathlineto{\pgfqpoint{3.431347in}{1.065646in}}%
\pgfpathlineto{\pgfqpoint{3.451532in}{1.072703in}}%
\pgfpathlineto{\pgfqpoint{3.452734in}{1.085951in}}%
\pgfpathlineto{\pgfqpoint{3.456630in}{1.087669in}}%
\pgfpathlineto{\pgfqpoint{3.476829in}{1.099199in}}%
\pgfpathlineto{\pgfqpoint{3.470404in}{1.112447in}}%
\pgfpathlineto{\pgfqpoint{3.481913in}{1.120427in}}%
\pgfpathlineto{\pgfqpoint{3.488052in}{1.125695in}}%
\pgfpathlineto{\pgfqpoint{3.484204in}{1.138943in}}%
\pgfpathlineto{\pgfqpoint{3.496341in}{1.152191in}}%
\pgfpathlineto{\pgfqpoint{3.493775in}{1.165439in}}%
\pgfpathlineto{\pgfqpoint{3.499118in}{1.178687in}}%
\pgfpathlineto{\pgfqpoint{3.497873in}{1.191935in}}%
\pgfpathlineto{\pgfqpoint{3.500986in}{1.205183in}}%
\pgfpathlineto{\pgfqpoint{3.501997in}{1.218431in}}%
\pgfpathlineto{\pgfqpoint{3.500737in}{1.231679in}}%
\pgfpathlineto{\pgfqpoint{3.503772in}{1.244927in}}%
\pgfpathlineto{\pgfqpoint{3.500398in}{1.258175in}}%
\pgfpathlineto{\pgfqpoint{3.504172in}{1.271423in}}%
\pgfpathlineto{\pgfqpoint{3.498538in}{1.284671in}}%
\pgfpathlineto{\pgfqpoint{3.500741in}{1.297919in}}%
\pgfpathlineto{\pgfqpoint{3.493136in}{1.311167in}}%
\pgfpathlineto{\pgfqpoint{3.497191in}{1.324415in}}%
\pgfpathlineto{\pgfqpoint{3.490707in}{1.337663in}}%
\pgfpathlineto{\pgfqpoint{3.493957in}{1.350911in}}%
\pgfpathlineto{\pgfqpoint{3.487096in}{1.364159in}}%
\pgfpathlineto{\pgfqpoint{3.489146in}{1.377407in}}%
\pgfpathlineto{\pgfqpoint{3.482943in}{1.390655in}}%
\pgfpathlineto{\pgfqpoint{3.484456in}{1.403903in}}%
\pgfpathlineto{\pgfqpoint{3.481913in}{1.410002in}}%
\pgfpathlineto{\pgfqpoint{3.478724in}{1.417151in}}%
\pgfpathlineto{\pgfqpoint{3.478026in}{1.430399in}}%
\pgfpathlineto{\pgfqpoint{3.473679in}{1.443647in}}%
\pgfpathlineto{\pgfqpoint{3.471891in}{1.456895in}}%
\pgfpathlineto{\pgfqpoint{3.468186in}{1.470143in}}%
\pgfpathlineto{\pgfqpoint{3.464561in}{1.483391in}}%
\pgfpathlineto{\pgfqpoint{3.461415in}{1.496639in}}%
\pgfpathlineto{\pgfqpoint{3.456957in}{1.509887in}}%
\pgfpathlineto{\pgfqpoint{3.456630in}{1.518842in}}%
\pgfpathlineto{\pgfqpoint{3.456458in}{1.523135in}}%
\pgfpathlineto{\pgfqpoint{3.451035in}{1.536383in}}%
\pgfpathlineto{\pgfqpoint{3.450168in}{1.549631in}}%
\pgfpathlineto{\pgfqpoint{3.443049in}{1.562879in}}%
\pgfpathlineto{\pgfqpoint{3.442666in}{1.576127in}}%
\pgfpathlineto{\pgfqpoint{3.435083in}{1.589375in}}%
\pgfpathlineto{\pgfqpoint{3.435252in}{1.602623in}}%
\pgfpathlineto{\pgfqpoint{3.431347in}{1.608915in}}%
\pgfpathlineto{\pgfqpoint{3.426311in}{1.615871in}}%
\pgfpathlineto{\pgfqpoint{3.425934in}{1.629119in}}%
\pgfpathlineto{\pgfqpoint{3.417045in}{1.642367in}}%
\pgfpathlineto{\pgfqpoint{3.416773in}{1.655615in}}%
\pgfpathlineto{\pgfqpoint{3.407656in}{1.668863in}}%
\pgfpathlineto{\pgfqpoint{3.406609in}{1.682111in}}%
\pgfpathlineto{\pgfqpoint{3.406064in}{1.682968in}}%
\pgfpathlineto{\pgfqpoint{3.395344in}{1.695359in}}%
\pgfpathlineto{\pgfqpoint{3.394100in}{1.708607in}}%
\pgfpathlineto{\pgfqpoint{3.385147in}{1.721855in}}%
\pgfpathlineto{\pgfqpoint{3.382204in}{1.735103in}}%
\pgfpathlineto{\pgfqpoint{3.380781in}{1.737022in}}%
\pgfpathlineto{\pgfqpoint{3.365917in}{1.748351in}}%
\pgfpathlineto{\pgfqpoint{3.357914in}{1.761599in}}%
\pgfpathlineto{\pgfqpoint{3.355497in}{1.763739in}}%
\pgfpathlineto{\pgfqpoint{3.330214in}{1.768087in}}%
\pgfpathlineto{\pgfqpoint{3.315688in}{1.761599in}}%
\pgfpathlineto{\pgfqpoint{3.307525in}{1.748351in}}%
\pgfpathlineto{\pgfqpoint{3.304931in}{1.747256in}}%
\pgfpathlineto{\pgfqpoint{3.291545in}{1.735103in}}%
\pgfpathlineto{\pgfqpoint{3.291310in}{1.721855in}}%
\pgfpathlineto{\pgfqpoint{3.279648in}{1.710384in}}%
\pgfpathlineto{\pgfqpoint{3.278390in}{1.708607in}}%
\pgfpathlineto{\pgfqpoint{3.278223in}{1.695359in}}%
\pgfpathlineto{\pgfqpoint{3.267422in}{1.682111in}}%
\pgfpathlineto{\pgfqpoint{3.266007in}{1.668863in}}%
\pgfpathlineto{\pgfqpoint{3.256983in}{1.655615in}}%
\pgfpathlineto{\pgfqpoint{3.254365in}{1.643153in}}%
\pgfpathlineto{\pgfqpoint{3.254223in}{1.642367in}}%
\pgfpathlineto{\pgfqpoint{3.247794in}{1.629119in}}%
\pgfpathlineto{\pgfqpoint{3.244407in}{1.615871in}}%
\pgfpathlineto{\pgfqpoint{3.240362in}{1.602623in}}%
\pgfpathlineto{\pgfqpoint{3.236246in}{1.589375in}}%
\pgfpathlineto{\pgfqpoint{3.233587in}{1.576127in}}%
\pgfpathlineto{\pgfqpoint{3.229082in}{1.567170in}}%
\pgfpathlineto{\pgfqpoint{3.227093in}{1.562879in}}%
\pgfpathlineto{\pgfqpoint{3.225938in}{1.549631in}}%
\pgfpathlineto{\pgfqpoint{3.219881in}{1.536383in}}%
\pgfpathlineto{\pgfqpoint{3.220365in}{1.523135in}}%
\pgfpathlineto{\pgfqpoint{3.212896in}{1.509887in}}%
\pgfpathlineto{\pgfqpoint{3.213308in}{1.496639in}}%
\pgfpathlineto{\pgfqpoint{3.204640in}{1.483391in}}%
\pgfpathlineto{\pgfqpoint{3.205946in}{1.470143in}}%
\pgfpathlineto{\pgfqpoint{3.203798in}{1.466397in}}%
\pgfpathlineto{\pgfqpoint{3.198752in}{1.456895in}}%
\pgfpathlineto{\pgfqpoint{3.199897in}{1.443647in}}%
\pgfpathlineto{\pgfqpoint{3.193281in}{1.430399in}}%
\pgfpathlineto{\pgfqpoint{3.194555in}{1.417151in}}%
\pgfpathlineto{\pgfqpoint{3.189566in}{1.403903in}}%
\pgfpathlineto{\pgfqpoint{3.189563in}{1.390655in}}%
\pgfpathlineto{\pgfqpoint{3.184630in}{1.377407in}}%
\pgfpathlineto{\pgfqpoint{3.182338in}{1.364159in}}%
\pgfpathlineto{\pgfqpoint{3.179858in}{1.350911in}}%
\pgfpathlineto{\pgfqpoint{3.178515in}{1.342568in}}%
\pgfpathlineto{\pgfqpoint{3.177727in}{1.337663in}}%
\pgfpathlineto{\pgfqpoint{3.177182in}{1.324415in}}%
\pgfpathlineto{\pgfqpoint{3.172627in}{1.311167in}}%
\pgfpathlineto{\pgfqpoint{3.173752in}{1.297919in}}%
\pgfpathlineto{\pgfqpoint{3.168406in}{1.284671in}}%
\pgfpathlineto{\pgfqpoint{3.172820in}{1.271423in}}%
\pgfpathlineto{\pgfqpoint{3.167011in}{1.258175in}}%
\pgfpathlineto{\pgfqpoint{3.172752in}{1.244927in}}%
\pgfpathlineto{\pgfqpoint{3.165688in}{1.231679in}}%
\pgfpathlineto{\pgfqpoint{3.173683in}{1.218431in}}%
\pgfpathlineto{\pgfqpoint{3.166268in}{1.205183in}}%
\pgfpathlineto{\pgfqpoint{3.174191in}{1.191935in}}%
\pgfpathlineto{\pgfqpoint{3.166038in}{1.178687in}}%
\pgfpathlineto{\pgfqpoint{3.175814in}{1.165439in}}%
\pgfpathlineto{\pgfqpoint{3.171683in}{1.152191in}}%
\pgfpathlineto{\pgfqpoint{3.178515in}{1.145007in}}%
\pgfpathlineto{\pgfqpoint{3.185608in}{1.138943in}}%
\pgfpathlineto{\pgfqpoint{3.182142in}{1.125695in}}%
\pgfpathlineto{\pgfqpoint{3.195560in}{1.112447in}}%
\pgfpathlineto{\pgfqpoint{3.196054in}{1.099199in}}%
\pgfpathlineto{\pgfqpoint{3.203798in}{1.092554in}}%
\pgfpathlineto{\pgfqpoint{3.214021in}{1.085951in}}%
\pgfpathlineto{\pgfqpoint{3.223436in}{1.072703in}}%
\pgfpathlineto{\pgfqpoint{3.229082in}{1.069161in}}%
\pgfpathlineto{\pgfqpoint{3.251991in}{1.059455in}}%
\pgfpathlineto{\pgfqpoint{3.254365in}{1.057672in}}%
\pgfpathlineto{\pgfqpoint{3.274196in}{1.046207in}}%
\pgfpathlineto{\pgfqpoint{3.279648in}{1.044327in}}%
\pgfpathclose%
\pgfusepath{stroke}%
\end{pgfscope}%
\begin{pgfscope}%
\pgfpathrectangle{\pgfqpoint{2.470586in}{0.476543in}}{\pgfqpoint{1.820388in}{1.656000in}}%
\pgfusepath{clip}%
\pgfsetbuttcap%
\pgfsetroundjoin%
\pgfsetlinewidth{0.501875pt}%
\definecolor{currentstroke}{rgb}{1.000000,1.000000,1.000000}%
\pgfsetstrokecolor{currentstroke}%
\pgfsetstrokeopacity{0.800000}%
\pgfsetdash{}{0pt}%
\pgfpathmoveto{\pgfqpoint{3.330214in}{1.202661in}}%
\pgfpathlineto{\pgfqpoint{3.355497in}{1.204502in}}%
\pgfpathlineto{\pgfqpoint{3.357885in}{1.205183in}}%
\pgfpathlineto{\pgfqpoint{3.358421in}{1.218431in}}%
\pgfpathlineto{\pgfqpoint{3.380781in}{1.226788in}}%
\pgfpathlineto{\pgfqpoint{3.387011in}{1.231679in}}%
\pgfpathlineto{\pgfqpoint{3.393135in}{1.244927in}}%
\pgfpathlineto{\pgfqpoint{3.400538in}{1.258175in}}%
\pgfpathlineto{\pgfqpoint{3.406064in}{1.264643in}}%
\pgfpathlineto{\pgfqpoint{3.410287in}{1.271423in}}%
\pgfpathlineto{\pgfqpoint{3.409737in}{1.284671in}}%
\pgfpathlineto{\pgfqpoint{3.415679in}{1.297919in}}%
\pgfpathlineto{\pgfqpoint{3.408975in}{1.311167in}}%
\pgfpathlineto{\pgfqpoint{3.418140in}{1.324415in}}%
\pgfpathlineto{\pgfqpoint{3.411078in}{1.337663in}}%
\pgfpathlineto{\pgfqpoint{3.420517in}{1.350911in}}%
\pgfpathlineto{\pgfqpoint{3.411524in}{1.364159in}}%
\pgfpathlineto{\pgfqpoint{3.418768in}{1.377407in}}%
\pgfpathlineto{\pgfqpoint{3.409988in}{1.390655in}}%
\pgfpathlineto{\pgfqpoint{3.417446in}{1.403903in}}%
\pgfpathlineto{\pgfqpoint{3.409292in}{1.417151in}}%
\pgfpathlineto{\pgfqpoint{3.413353in}{1.430399in}}%
\pgfpathlineto{\pgfqpoint{3.406428in}{1.443647in}}%
\pgfpathlineto{\pgfqpoint{3.408408in}{1.456895in}}%
\pgfpathlineto{\pgfqpoint{3.406064in}{1.462401in}}%
\pgfpathlineto{\pgfqpoint{3.401494in}{1.470143in}}%
\pgfpathlineto{\pgfqpoint{3.400671in}{1.483391in}}%
\pgfpathlineto{\pgfqpoint{3.392946in}{1.496639in}}%
\pgfpathlineto{\pgfqpoint{3.389893in}{1.509887in}}%
\pgfpathlineto{\pgfqpoint{3.387543in}{1.523135in}}%
\pgfpathlineto{\pgfqpoint{3.383302in}{1.536383in}}%
\pgfpathlineto{\pgfqpoint{3.381605in}{1.549631in}}%
\pgfpathlineto{\pgfqpoint{3.380781in}{1.550929in}}%
\pgfpathlineto{\pgfqpoint{3.368689in}{1.562879in}}%
\pgfpathlineto{\pgfqpoint{3.366346in}{1.576127in}}%
\pgfpathlineto{\pgfqpoint{3.355497in}{1.585461in}}%
\pgfpathlineto{\pgfqpoint{3.337388in}{1.589375in}}%
\pgfpathlineto{\pgfqpoint{3.337881in}{1.602623in}}%
\pgfpathlineto{\pgfqpoint{3.330214in}{1.603632in}}%
\pgfpathlineto{\pgfqpoint{3.328208in}{1.602623in}}%
\pgfpathlineto{\pgfqpoint{3.327259in}{1.589375in}}%
\pgfpathlineto{\pgfqpoint{3.308739in}{1.576127in}}%
\pgfpathlineto{\pgfqpoint{3.304931in}{1.572012in}}%
\pgfpathlineto{\pgfqpoint{3.300628in}{1.562879in}}%
\pgfpathlineto{\pgfqpoint{3.295985in}{1.549631in}}%
\pgfpathlineto{\pgfqpoint{3.288294in}{1.536383in}}%
\pgfpathlineto{\pgfqpoint{3.289353in}{1.523135in}}%
\pgfpathlineto{\pgfqpoint{3.279648in}{1.511680in}}%
\pgfpathlineto{\pgfqpoint{3.278552in}{1.509887in}}%
\pgfpathlineto{\pgfqpoint{3.279648in}{1.501201in}}%
\pgfpathlineto{\pgfqpoint{3.280507in}{1.496639in}}%
\pgfpathlineto{\pgfqpoint{3.279648in}{1.495949in}}%
\pgfpathlineto{\pgfqpoint{3.268790in}{1.483391in}}%
\pgfpathlineto{\pgfqpoint{3.272906in}{1.470143in}}%
\pgfpathlineto{\pgfqpoint{3.262630in}{1.456895in}}%
\pgfpathlineto{\pgfqpoint{3.268198in}{1.443647in}}%
\pgfpathlineto{\pgfqpoint{3.257373in}{1.430399in}}%
\pgfpathlineto{\pgfqpoint{3.264350in}{1.417151in}}%
\pgfpathlineto{\pgfqpoint{3.255739in}{1.403903in}}%
\pgfpathlineto{\pgfqpoint{3.262697in}{1.390655in}}%
\pgfpathlineto{\pgfqpoint{3.254365in}{1.377962in}}%
\pgfpathlineto{\pgfqpoint{3.254084in}{1.377407in}}%
\pgfpathlineto{\pgfqpoint{3.254365in}{1.375998in}}%
\pgfpathlineto{\pgfqpoint{3.257701in}{1.364159in}}%
\pgfpathlineto{\pgfqpoint{3.254365in}{1.356188in}}%
\pgfpathlineto{\pgfqpoint{3.252699in}{1.350911in}}%
\pgfpathlineto{\pgfqpoint{3.254365in}{1.344735in}}%
\pgfpathlineto{\pgfqpoint{3.257097in}{1.337663in}}%
\pgfpathlineto{\pgfqpoint{3.255673in}{1.324415in}}%
\pgfpathlineto{\pgfqpoint{3.256520in}{1.311167in}}%
\pgfpathlineto{\pgfqpoint{3.258408in}{1.297919in}}%
\pgfpathlineto{\pgfqpoint{3.256932in}{1.284671in}}%
\pgfpathlineto{\pgfqpoint{3.267163in}{1.271423in}}%
\pgfpathlineto{\pgfqpoint{3.266867in}{1.258175in}}%
\pgfpathlineto{\pgfqpoint{3.279648in}{1.246769in}}%
\pgfpathlineto{\pgfqpoint{3.282740in}{1.244927in}}%
\pgfpathlineto{\pgfqpoint{3.279648in}{1.236608in}}%
\pgfpathlineto{\pgfqpoint{3.278300in}{1.231679in}}%
\pgfpathlineto{\pgfqpoint{3.279648in}{1.230875in}}%
\pgfpathlineto{\pgfqpoint{3.304931in}{1.221086in}}%
\pgfpathlineto{\pgfqpoint{3.319654in}{1.218431in}}%
\pgfpathlineto{\pgfqpoint{3.310185in}{1.205183in}}%
\pgfpathlineto{\pgfqpoint{3.330214in}{1.202661in}}%
\pgfpathclose%
\pgfusepath{stroke}%
\end{pgfscope}%
\begin{pgfscope}%
\pgfpathrectangle{\pgfqpoint{2.470586in}{0.476543in}}{\pgfqpoint{1.820388in}{1.656000in}}%
\pgfusepath{clip}%
\pgfsetbuttcap%
\pgfsetbeveljoin%
\definecolor{currentfill}{rgb}{1.000000,1.000000,1.000000}%
\pgfsetfillcolor{currentfill}%
\pgfsetlinewidth{0.501875pt}%
\definecolor{currentstroke}{rgb}{0.000000,0.000000,0.000000}%
\pgfsetstrokecolor{currentstroke}%
\pgfsetdash{}{0pt}%
\pgfsys@defobject{currentmarker}{\pgfqpoint{-0.052836in}{-0.044945in}}{\pgfqpoint{0.052836in}{0.055556in}}{%
\pgfpathmoveto{\pgfqpoint{0.000000in}{0.055556in}}%
\pgfpathlineto{\pgfqpoint{-0.012473in}{0.017168in}}%
\pgfpathlineto{\pgfqpoint{-0.052836in}{0.017168in}}%
\pgfpathlineto{\pgfqpoint{-0.020182in}{-0.006557in}}%
\pgfpathlineto{\pgfqpoint{-0.032655in}{-0.044945in}}%
\pgfpathlineto{\pgfqpoint{-0.000000in}{-0.021220in}}%
\pgfpathlineto{\pgfqpoint{0.032655in}{-0.044945in}}%
\pgfpathlineto{\pgfqpoint{0.020182in}{-0.006557in}}%
\pgfpathlineto{\pgfqpoint{0.052836in}{0.017168in}}%
\pgfpathlineto{\pgfqpoint{0.012473in}{0.017168in}}%
\pgfpathlineto{\pgfqpoint{0.000000in}{0.055556in}}%
\pgfpathclose%
\pgfusepath{stroke,fill}%
}%
\begin{pgfscope}%
\pgfsys@transformshift{3.330214in}{1.470143in}%
\pgfsys@useobject{currentmarker}{}%
\end{pgfscope}%
\end{pgfscope}%
\begin{pgfscope}%
\pgfsetrectcap%
\pgfsetmiterjoin%
\pgfsetlinewidth{0.803000pt}%
\definecolor{currentstroke}{rgb}{0.000000,0.000000,0.000000}%
\pgfsetstrokecolor{currentstroke}%
\pgfsetdash{}{0pt}%
\pgfpathmoveto{\pgfqpoint{2.470586in}{0.476543in}}%
\pgfpathlineto{\pgfqpoint{2.470586in}{2.132543in}}%
\pgfusepath{stroke}%
\end{pgfscope}%
\begin{pgfscope}%
\pgfsetrectcap%
\pgfsetmiterjoin%
\pgfsetlinewidth{0.803000pt}%
\definecolor{currentstroke}{rgb}{0.000000,0.000000,0.000000}%
\pgfsetstrokecolor{currentstroke}%
\pgfsetdash{}{0pt}%
\pgfpathmoveto{\pgfqpoint{4.290975in}{0.476543in}}%
\pgfpathlineto{\pgfqpoint{4.290975in}{2.132543in}}%
\pgfusepath{stroke}%
\end{pgfscope}%
\begin{pgfscope}%
\pgfsetrectcap%
\pgfsetmiterjoin%
\pgfsetlinewidth{0.803000pt}%
\definecolor{currentstroke}{rgb}{0.000000,0.000000,0.000000}%
\pgfsetstrokecolor{currentstroke}%
\pgfsetdash{}{0pt}%
\pgfpathmoveto{\pgfqpoint{2.470586in}{0.476543in}}%
\pgfpathlineto{\pgfqpoint{4.290975in}{0.476543in}}%
\pgfusepath{stroke}%
\end{pgfscope}%
\begin{pgfscope}%
\pgfsetrectcap%
\pgfsetmiterjoin%
\pgfsetlinewidth{0.803000pt}%
\definecolor{currentstroke}{rgb}{0.000000,0.000000,0.000000}%
\pgfsetstrokecolor{currentstroke}%
\pgfsetdash{}{0pt}%
\pgfpathmoveto{\pgfqpoint{2.470586in}{2.132543in}}%
\pgfpathlineto{\pgfqpoint{4.290975in}{2.132543in}}%
\pgfusepath{stroke}%
\end{pgfscope}%
\begin{pgfscope}%
\definecolor{textcolor}{rgb}{1.000000,1.000000,1.000000}%
\pgfsetstrokecolor{textcolor}%
\pgfsetfillcolor{textcolor}%
\pgftext[x=2.525198in,y=2.049743in,left,top]{\color{textcolor}{\rmfamily\fontsize{9.000000}{10.800000}\selectfont\catcode`\^=\active\def^{\ifmmode\sp\else\^{}\fi}\catcode`\%=\active\def%{\%}$(b)$}}%
\end{pgfscope}%
\begin{pgfscope}%
\definecolor{textcolor}{rgb}{0.200000,0.200000,0.200000}%
\pgfsetstrokecolor{textcolor}%
\pgfsetfillcolor{textcolor}%
\pgftext[x=3.380781in,y=2.182223in,,bottom]{\color{textcolor}{\rmfamily\fontsize{8.000000}{9.600000}\selectfont\catcode`\^=\active\def^{\ifmmode\sp\else\^{}\fi}\catcode`\%=\active\def%{\%}$D_2$}}%
\end{pgfscope}%
\begin{pgfscope}%
\pgfsetbuttcap%
\pgfsetmiterjoin%
\definecolor{currentfill}{rgb}{1.000000,1.000000,1.000000}%
\pgfsetfillcolor{currentfill}%
\pgfsetlinewidth{0.000000pt}%
\definecolor{currentstroke}{rgb}{0.000000,0.000000,0.000000}%
\pgfsetstrokecolor{currentstroke}%
\pgfsetstrokeopacity{0.000000}%
\pgfsetdash{}{0pt}%
\pgfpathmoveto{\pgfqpoint{4.390975in}{0.476543in}}%
\pgfpathlineto{\pgfqpoint{4.490975in}{0.476543in}}%
\pgfpathlineto{\pgfqpoint{4.490975in}{2.132543in}}%
\pgfpathlineto{\pgfqpoint{4.390975in}{2.132543in}}%
\pgfpathlineto{\pgfqpoint{4.390975in}{0.476543in}}%
\pgfpathclose%
\pgfusepath{fill}%
\end{pgfscope}%
\begin{pgfscope}%
\pgfpathrectangle{\pgfqpoint{4.390975in}{0.476543in}}{\pgfqpoint{0.100000in}{1.656000in}}%
\pgfusepath{clip}%
\pgfsetbuttcap%
\pgfsetroundjoin%
\definecolor{currentfill}{rgb}{0.207536,0.260208,0.608151}%
\pgfsetfillcolor{currentfill}%
\pgfsetlinewidth{0.000000pt}%
\definecolor{currentstroke}{rgb}{0.000000,0.000000,0.000000}%
\pgfsetstrokecolor{currentstroke}%
\pgfsetdash{}{0pt}%
\pgfpathmoveto{\pgfqpoint{4.390975in}{0.476543in}}%
\pgfpathlineto{\pgfqpoint{4.490975in}{0.476543in}}%
\pgfpathlineto{\pgfqpoint{4.490975in}{0.545543in}}%
\pgfpathlineto{\pgfqpoint{4.390975in}{0.545543in}}%
\pgfpathlineto{\pgfqpoint{4.390975in}{0.476543in}}%
\pgfusepath{fill}%
\end{pgfscope}%
\begin{pgfscope}%
\pgfpathrectangle{\pgfqpoint{4.390975in}{0.476543in}}{\pgfqpoint{0.100000in}{1.656000in}}%
\pgfusepath{clip}%
\pgfsetbuttcap%
\pgfsetroundjoin%
\definecolor{currentfill}{rgb}{0.241369,0.366782,0.660592}%
\pgfsetfillcolor{currentfill}%
\pgfsetlinewidth{0.000000pt}%
\definecolor{currentstroke}{rgb}{0.000000,0.000000,0.000000}%
\pgfsetstrokecolor{currentstroke}%
\pgfsetdash{}{0pt}%
\pgfpathmoveto{\pgfqpoint{4.390975in}{0.545543in}}%
\pgfpathlineto{\pgfqpoint{4.490975in}{0.545543in}}%
\pgfpathlineto{\pgfqpoint{4.490975in}{0.614543in}}%
\pgfpathlineto{\pgfqpoint{4.390975in}{0.614543in}}%
\pgfpathlineto{\pgfqpoint{4.390975in}{0.545543in}}%
\pgfusepath{fill}%
\end{pgfscope}%
\begin{pgfscope}%
\pgfpathrectangle{\pgfqpoint{4.390975in}{0.476543in}}{\pgfqpoint{0.100000in}{1.656000in}}%
\pgfusepath{clip}%
\pgfsetbuttcap%
\pgfsetroundjoin%
\definecolor{currentfill}{rgb}{0.274202,0.463130,0.708112}%
\pgfsetfillcolor{currentfill}%
\pgfsetlinewidth{0.000000pt}%
\definecolor{currentstroke}{rgb}{0.000000,0.000000,0.000000}%
\pgfsetstrokecolor{currentstroke}%
\pgfsetdash{}{0pt}%
\pgfpathmoveto{\pgfqpoint{4.390975in}{0.614543in}}%
\pgfpathlineto{\pgfqpoint{4.490975in}{0.614543in}}%
\pgfpathlineto{\pgfqpoint{4.490975in}{0.683543in}}%
\pgfpathlineto{\pgfqpoint{4.390975in}{0.683543in}}%
\pgfpathlineto{\pgfqpoint{4.390975in}{0.614543in}}%
\pgfusepath{fill}%
\end{pgfscope}%
\begin{pgfscope}%
\pgfpathrectangle{\pgfqpoint{4.390975in}{0.476543in}}{\pgfqpoint{0.100000in}{1.656000in}}%
\pgfusepath{clip}%
\pgfsetbuttcap%
\pgfsetroundjoin%
\definecolor{currentfill}{rgb}{0.353710,0.557862,0.757170}%
\pgfsetfillcolor{currentfill}%
\pgfsetlinewidth{0.000000pt}%
\definecolor{currentstroke}{rgb}{0.000000,0.000000,0.000000}%
\pgfsetstrokecolor{currentstroke}%
\pgfsetdash{}{0pt}%
\pgfpathmoveto{\pgfqpoint{4.390975in}{0.683543in}}%
\pgfpathlineto{\pgfqpoint{4.490975in}{0.683543in}}%
\pgfpathlineto{\pgfqpoint{4.490975in}{0.752543in}}%
\pgfpathlineto{\pgfqpoint{4.390975in}{0.752543in}}%
\pgfpathlineto{\pgfqpoint{4.390975in}{0.683543in}}%
\pgfusepath{fill}%
\end{pgfscope}%
\begin{pgfscope}%
\pgfpathrectangle{\pgfqpoint{4.390975in}{0.476543in}}{\pgfqpoint{0.100000in}{1.656000in}}%
\pgfusepath{clip}%
\pgfsetbuttcap%
\pgfsetroundjoin%
\definecolor{currentfill}{rgb}{0.433218,0.652595,0.806228}%
\pgfsetfillcolor{currentfill}%
\pgfsetlinewidth{0.000000pt}%
\definecolor{currentstroke}{rgb}{0.000000,0.000000,0.000000}%
\pgfsetstrokecolor{currentstroke}%
\pgfsetdash{}{0pt}%
\pgfpathmoveto{\pgfqpoint{4.390975in}{0.752543in}}%
\pgfpathlineto{\pgfqpoint{4.490975in}{0.752543in}}%
\pgfpathlineto{\pgfqpoint{4.490975in}{0.821543in}}%
\pgfpathlineto{\pgfqpoint{4.390975in}{0.821543in}}%
\pgfpathlineto{\pgfqpoint{4.390975in}{0.752543in}}%
\pgfusepath{fill}%
\end{pgfscope}%
\begin{pgfscope}%
\pgfpathrectangle{\pgfqpoint{4.390975in}{0.476543in}}{\pgfqpoint{0.100000in}{1.656000in}}%
\pgfusepath{clip}%
\pgfsetbuttcap%
\pgfsetroundjoin%
\definecolor{currentfill}{rgb}{0.514110,0.725798,0.845444}%
\pgfsetfillcolor{currentfill}%
\pgfsetlinewidth{0.000000pt}%
\definecolor{currentstroke}{rgb}{0.000000,0.000000,0.000000}%
\pgfsetstrokecolor{currentstroke}%
\pgfsetdash{}{0pt}%
\pgfpathmoveto{\pgfqpoint{4.390975in}{0.821543in}}%
\pgfpathlineto{\pgfqpoint{4.490975in}{0.821543in}}%
\pgfpathlineto{\pgfqpoint{4.490975in}{0.890543in}}%
\pgfpathlineto{\pgfqpoint{4.390975in}{0.890543in}}%
\pgfpathlineto{\pgfqpoint{4.390975in}{0.821543in}}%
\pgfusepath{fill}%
\end{pgfscope}%
\begin{pgfscope}%
\pgfpathrectangle{\pgfqpoint{4.390975in}{0.476543in}}{\pgfqpoint{0.100000in}{1.656000in}}%
\pgfusepath{clip}%
\pgfsetbuttcap%
\pgfsetroundjoin%
\definecolor{currentfill}{rgb}{0.607151,0.800231,0.886044}%
\pgfsetfillcolor{currentfill}%
\pgfsetlinewidth{0.000000pt}%
\definecolor{currentstroke}{rgb}{0.000000,0.000000,0.000000}%
\pgfsetstrokecolor{currentstroke}%
\pgfsetdash{}{0pt}%
\pgfpathmoveto{\pgfqpoint{4.390975in}{0.890543in}}%
\pgfpathlineto{\pgfqpoint{4.490975in}{0.890543in}}%
\pgfpathlineto{\pgfqpoint{4.490975in}{0.959543in}}%
\pgfpathlineto{\pgfqpoint{4.390975in}{0.959543in}}%
\pgfpathlineto{\pgfqpoint{4.390975in}{0.890543in}}%
\pgfusepath{fill}%
\end{pgfscope}%
\begin{pgfscope}%
\pgfpathrectangle{\pgfqpoint{4.390975in}{0.476543in}}{\pgfqpoint{0.100000in}{1.656000in}}%
\pgfusepath{clip}%
\pgfsetbuttcap%
\pgfsetroundjoin%
\definecolor{currentfill}{rgb}{0.699116,0.864975,0.921799}%
\pgfsetfillcolor{currentfill}%
\pgfsetlinewidth{0.000000pt}%
\definecolor{currentstroke}{rgb}{0.000000,0.000000,0.000000}%
\pgfsetstrokecolor{currentstroke}%
\pgfsetdash{}{0pt}%
\pgfpathmoveto{\pgfqpoint{4.390975in}{0.959543in}}%
\pgfpathlineto{\pgfqpoint{4.490975in}{0.959543in}}%
\pgfpathlineto{\pgfqpoint{4.490975in}{1.028543in}}%
\pgfpathlineto{\pgfqpoint{4.390975in}{1.028543in}}%
\pgfpathlineto{\pgfqpoint{4.390975in}{0.959543in}}%
\pgfusepath{fill}%
\end{pgfscope}%
\begin{pgfscope}%
\pgfpathrectangle{\pgfqpoint{4.390975in}{0.476543in}}{\pgfqpoint{0.100000in}{1.656000in}}%
\pgfusepath{clip}%
\pgfsetbuttcap%
\pgfsetroundjoin%
\definecolor{currentfill}{rgb}{0.780623,0.904960,0.944867}%
\pgfsetfillcolor{currentfill}%
\pgfsetlinewidth{0.000000pt}%
\definecolor{currentstroke}{rgb}{0.000000,0.000000,0.000000}%
\pgfsetstrokecolor{currentstroke}%
\pgfsetdash{}{0pt}%
\pgfpathmoveto{\pgfqpoint{4.390975in}{1.028543in}}%
\pgfpathlineto{\pgfqpoint{4.490975in}{1.028543in}}%
\pgfpathlineto{\pgfqpoint{4.490975in}{1.097543in}}%
\pgfpathlineto{\pgfqpoint{4.390975in}{1.097543in}}%
\pgfpathlineto{\pgfqpoint{4.390975in}{1.028543in}}%
\pgfusepath{fill}%
\end{pgfscope}%
\begin{pgfscope}%
\pgfpathrectangle{\pgfqpoint{4.390975in}{0.476543in}}{\pgfqpoint{0.100000in}{1.656000in}}%
\pgfusepath{clip}%
\pgfsetbuttcap%
\pgfsetroundjoin%
\definecolor{currentfill}{rgb}{0.870281,0.948943,0.970242}%
\pgfsetfillcolor{currentfill}%
\pgfsetlinewidth{0.000000pt}%
\definecolor{currentstroke}{rgb}{0.000000,0.000000,0.000000}%
\pgfsetstrokecolor{currentstroke}%
\pgfsetdash{}{0pt}%
\pgfpathmoveto{\pgfqpoint{4.390975in}{1.097543in}}%
\pgfpathlineto{\pgfqpoint{4.490975in}{1.097543in}}%
\pgfpathlineto{\pgfqpoint{4.490975in}{1.166543in}}%
\pgfpathlineto{\pgfqpoint{4.390975in}{1.166543in}}%
\pgfpathlineto{\pgfqpoint{4.390975in}{1.097543in}}%
\pgfusepath{fill}%
\end{pgfscope}%
\begin{pgfscope}%
\pgfpathrectangle{\pgfqpoint{4.390975in}{0.476543in}}{\pgfqpoint{0.100000in}{1.656000in}}%
\pgfusepath{clip}%
\pgfsetbuttcap%
\pgfsetroundjoin%
\definecolor{currentfill}{rgb}{0.926105,0.971396,0.884890}%
\pgfsetfillcolor{currentfill}%
\pgfsetlinewidth{0.000000pt}%
\definecolor{currentstroke}{rgb}{0.000000,0.000000,0.000000}%
\pgfsetstrokecolor{currentstroke}%
\pgfsetdash{}{0pt}%
\pgfpathmoveto{\pgfqpoint{4.390975in}{1.166543in}}%
\pgfpathlineto{\pgfqpoint{4.490975in}{1.166543in}}%
\pgfpathlineto{\pgfqpoint{4.490975in}{1.235543in}}%
\pgfpathlineto{\pgfqpoint{4.390975in}{1.235543in}}%
\pgfpathlineto{\pgfqpoint{4.390975in}{1.166543in}}%
\pgfusepath{fill}%
\end{pgfscope}%
\begin{pgfscope}%
\pgfpathrectangle{\pgfqpoint{4.390975in}{0.476543in}}{\pgfqpoint{0.100000in}{1.656000in}}%
\pgfusepath{clip}%
\pgfsetbuttcap%
\pgfsetroundjoin%
\definecolor{currentfill}{rgb}{0.973779,0.989850,0.797232}%
\pgfsetfillcolor{currentfill}%
\pgfsetlinewidth{0.000000pt}%
\definecolor{currentstroke}{rgb}{0.000000,0.000000,0.000000}%
\pgfsetstrokecolor{currentstroke}%
\pgfsetdash{}{0pt}%
\pgfpathmoveto{\pgfqpoint{4.390975in}{1.235543in}}%
\pgfpathlineto{\pgfqpoint{4.490975in}{1.235543in}}%
\pgfpathlineto{\pgfqpoint{4.490975in}{1.304543in}}%
\pgfpathlineto{\pgfqpoint{4.390975in}{1.304543in}}%
\pgfpathlineto{\pgfqpoint{4.390975in}{1.235543in}}%
\pgfusepath{fill}%
\end{pgfscope}%
\begin{pgfscope}%
\pgfpathrectangle{\pgfqpoint{4.390975in}{0.476543in}}{\pgfqpoint{0.100000in}{1.656000in}}%
\pgfusepath{clip}%
\pgfsetbuttcap%
\pgfsetroundjoin%
\definecolor{currentfill}{rgb}{0.999154,0.973779,0.709266}%
\pgfsetfillcolor{currentfill}%
\pgfsetlinewidth{0.000000pt}%
\definecolor{currentstroke}{rgb}{0.000000,0.000000,0.000000}%
\pgfsetstrokecolor{currentstroke}%
\pgfsetdash{}{0pt}%
\pgfpathmoveto{\pgfqpoint{4.390975in}{1.304543in}}%
\pgfpathlineto{\pgfqpoint{4.490975in}{1.304543in}}%
\pgfpathlineto{\pgfqpoint{4.490975in}{1.373543in}}%
\pgfpathlineto{\pgfqpoint{4.390975in}{1.373543in}}%
\pgfpathlineto{\pgfqpoint{4.390975in}{1.304543in}}%
\pgfusepath{fill}%
\end{pgfscope}%
\begin{pgfscope}%
\pgfpathrectangle{\pgfqpoint{4.390975in}{0.476543in}}{\pgfqpoint{0.100000in}{1.656000in}}%
\pgfusepath{clip}%
\pgfsetbuttcap%
\pgfsetroundjoin%
\definecolor{currentfill}{rgb}{0.997463,0.921338,0.629758}%
\pgfsetfillcolor{currentfill}%
\pgfsetlinewidth{0.000000pt}%
\definecolor{currentstroke}{rgb}{0.000000,0.000000,0.000000}%
\pgfsetstrokecolor{currentstroke}%
\pgfsetdash{}{0pt}%
\pgfpathmoveto{\pgfqpoint{4.390975in}{1.373543in}}%
\pgfpathlineto{\pgfqpoint{4.490975in}{1.373543in}}%
\pgfpathlineto{\pgfqpoint{4.490975in}{1.442543in}}%
\pgfpathlineto{\pgfqpoint{4.390975in}{1.442543in}}%
\pgfpathlineto{\pgfqpoint{4.390975in}{1.373543in}}%
\pgfusepath{fill}%
\end{pgfscope}%
\begin{pgfscope}%
\pgfpathrectangle{\pgfqpoint{4.390975in}{0.476543in}}{\pgfqpoint{0.100000in}{1.656000in}}%
\pgfusepath{clip}%
\pgfsetbuttcap%
\pgfsetroundjoin%
\definecolor{currentfill}{rgb}{0.995925,0.870742,0.557478}%
\pgfsetfillcolor{currentfill}%
\pgfsetlinewidth{0.000000pt}%
\definecolor{currentstroke}{rgb}{0.000000,0.000000,0.000000}%
\pgfsetstrokecolor{currentstroke}%
\pgfsetdash{}{0pt}%
\pgfpathmoveto{\pgfqpoint{4.390975in}{1.442543in}}%
\pgfpathlineto{\pgfqpoint{4.490975in}{1.442543in}}%
\pgfpathlineto{\pgfqpoint{4.490975in}{1.511543in}}%
\pgfpathlineto{\pgfqpoint{4.390975in}{1.511543in}}%
\pgfpathlineto{\pgfqpoint{4.390975in}{1.442543in}}%
\pgfusepath{fill}%
\end{pgfscope}%
\begin{pgfscope}%
\pgfpathrectangle{\pgfqpoint{4.390975in}{0.476543in}}{\pgfqpoint{0.100000in}{1.656000in}}%
\pgfusepath{clip}%
\pgfsetbuttcap%
\pgfsetroundjoin%
\definecolor{currentfill}{rgb}{0.994233,0.786159,0.477970}%
\pgfsetfillcolor{currentfill}%
\pgfsetlinewidth{0.000000pt}%
\definecolor{currentstroke}{rgb}{0.000000,0.000000,0.000000}%
\pgfsetstrokecolor{currentstroke}%
\pgfsetdash{}{0pt}%
\pgfpathmoveto{\pgfqpoint{4.390975in}{1.511543in}}%
\pgfpathlineto{\pgfqpoint{4.490975in}{1.511543in}}%
\pgfpathlineto{\pgfqpoint{4.490975in}{1.580543in}}%
\pgfpathlineto{\pgfqpoint{4.390975in}{1.580543in}}%
\pgfpathlineto{\pgfqpoint{4.390975in}{1.511543in}}%
\pgfusepath{fill}%
\end{pgfscope}%
\begin{pgfscope}%
\pgfpathrectangle{\pgfqpoint{4.390975in}{0.476543in}}{\pgfqpoint{0.100000in}{1.656000in}}%
\pgfusepath{clip}%
\pgfsetbuttcap%
\pgfsetroundjoin%
\definecolor{currentfill}{rgb}{0.992541,0.701576,0.398462}%
\pgfsetfillcolor{currentfill}%
\pgfsetlinewidth{0.000000pt}%
\definecolor{currentstroke}{rgb}{0.000000,0.000000,0.000000}%
\pgfsetstrokecolor{currentstroke}%
\pgfsetdash{}{0pt}%
\pgfpathmoveto{\pgfqpoint{4.390975in}{1.580543in}}%
\pgfpathlineto{\pgfqpoint{4.490975in}{1.580543in}}%
\pgfpathlineto{\pgfqpoint{4.490975in}{1.649543in}}%
\pgfpathlineto{\pgfqpoint{4.390975in}{1.649543in}}%
\pgfpathlineto{\pgfqpoint{4.390975in}{1.580543in}}%
\pgfusepath{fill}%
\end{pgfscope}%
\begin{pgfscope}%
\pgfpathrectangle{\pgfqpoint{4.390975in}{0.476543in}}{\pgfqpoint{0.100000in}{1.656000in}}%
\pgfusepath{clip}%
\pgfsetbuttcap%
\pgfsetroundjoin%
\definecolor{currentfill}{rgb}{0.981776,0.607382,0.345790}%
\pgfsetfillcolor{currentfill}%
\pgfsetlinewidth{0.000000pt}%
\definecolor{currentstroke}{rgb}{0.000000,0.000000,0.000000}%
\pgfsetstrokecolor{currentstroke}%
\pgfsetdash{}{0pt}%
\pgfpathmoveto{\pgfqpoint{4.390975in}{1.649543in}}%
\pgfpathlineto{\pgfqpoint{4.490975in}{1.649543in}}%
\pgfpathlineto{\pgfqpoint{4.490975in}{1.718543in}}%
\pgfpathlineto{\pgfqpoint{4.390975in}{1.718543in}}%
\pgfpathlineto{\pgfqpoint{4.390975in}{1.649543in}}%
\pgfusepath{fill}%
\end{pgfscope}%
\begin{pgfscope}%
\pgfpathrectangle{\pgfqpoint{4.390975in}{0.476543in}}{\pgfqpoint{0.100000in}{1.656000in}}%
\pgfusepath{clip}%
\pgfsetbuttcap%
\pgfsetroundjoin%
\definecolor{currentfill}{rgb}{0.966551,0.497424,0.295040}%
\pgfsetfillcolor{currentfill}%
\pgfsetlinewidth{0.000000pt}%
\definecolor{currentstroke}{rgb}{0.000000,0.000000,0.000000}%
\pgfsetstrokecolor{currentstroke}%
\pgfsetdash{}{0pt}%
\pgfpathmoveto{\pgfqpoint{4.390975in}{1.718543in}}%
\pgfpathlineto{\pgfqpoint{4.490975in}{1.718543in}}%
\pgfpathlineto{\pgfqpoint{4.490975in}{1.787543in}}%
\pgfpathlineto{\pgfqpoint{4.390975in}{1.787543in}}%
\pgfpathlineto{\pgfqpoint{4.390975in}{1.718543in}}%
\pgfusepath{fill}%
\end{pgfscope}%
\begin{pgfscope}%
\pgfpathrectangle{\pgfqpoint{4.390975in}{0.476543in}}{\pgfqpoint{0.100000in}{1.656000in}}%
\pgfusepath{clip}%
\pgfsetbuttcap%
\pgfsetroundjoin%
\definecolor{currentfill}{rgb}{0.939023,0.389927,0.245521}%
\pgfsetfillcolor{currentfill}%
\pgfsetlinewidth{0.000000pt}%
\definecolor{currentstroke}{rgb}{0.000000,0.000000,0.000000}%
\pgfsetstrokecolor{currentstroke}%
\pgfsetdash{}{0pt}%
\pgfpathmoveto{\pgfqpoint{4.390975in}{1.787543in}}%
\pgfpathlineto{\pgfqpoint{4.490975in}{1.787543in}}%
\pgfpathlineto{\pgfqpoint{4.490975in}{1.856543in}}%
\pgfpathlineto{\pgfqpoint{4.390975in}{1.856543in}}%
\pgfpathlineto{\pgfqpoint{4.390975in}{1.787543in}}%
\pgfusepath{fill}%
\end{pgfscope}%
\begin{pgfscope}%
\pgfpathrectangle{\pgfqpoint{4.390975in}{0.476543in}}{\pgfqpoint{0.100000in}{1.656000in}}%
\pgfusepath{clip}%
\pgfsetbuttcap%
\pgfsetroundjoin%
\definecolor{currentfill}{rgb}{0.894425,0.296117,0.202461}%
\pgfsetfillcolor{currentfill}%
\pgfsetlinewidth{0.000000pt}%
\definecolor{currentstroke}{rgb}{0.000000,0.000000,0.000000}%
\pgfsetstrokecolor{currentstroke}%
\pgfsetdash{}{0pt}%
\pgfpathmoveto{\pgfqpoint{4.390975in}{1.856543in}}%
\pgfpathlineto{\pgfqpoint{4.490975in}{1.856543in}}%
\pgfpathlineto{\pgfqpoint{4.490975in}{1.925543in}}%
\pgfpathlineto{\pgfqpoint{4.390975in}{1.925543in}}%
\pgfpathlineto{\pgfqpoint{4.390975in}{1.856543in}}%
\pgfusepath{fill}%
\end{pgfscope}%
\begin{pgfscope}%
\pgfpathrectangle{\pgfqpoint{4.390975in}{0.476543in}}{\pgfqpoint{0.100000in}{1.656000in}}%
\pgfusepath{clip}%
\pgfsetbuttcap%
\pgfsetroundjoin%
\definecolor{currentfill}{rgb}{0.845367,0.192926,0.155094}%
\pgfsetfillcolor{currentfill}%
\pgfsetlinewidth{0.000000pt}%
\definecolor{currentstroke}{rgb}{0.000000,0.000000,0.000000}%
\pgfsetstrokecolor{currentstroke}%
\pgfsetdash{}{0pt}%
\pgfpathmoveto{\pgfqpoint{4.390975in}{1.925543in}}%
\pgfpathlineto{\pgfqpoint{4.490975in}{1.925543in}}%
\pgfpathlineto{\pgfqpoint{4.490975in}{1.994543in}}%
\pgfpathlineto{\pgfqpoint{4.390975in}{1.994543in}}%
\pgfpathlineto{\pgfqpoint{4.390975in}{1.925543in}}%
\pgfusepath{fill}%
\end{pgfscope}%
\begin{pgfscope}%
\pgfpathrectangle{\pgfqpoint{4.390975in}{0.476543in}}{\pgfqpoint{0.100000in}{1.656000in}}%
\pgfusepath{clip}%
\pgfsetbuttcap%
\pgfsetroundjoin%
\definecolor{currentfill}{rgb}{0.762399,0.110727,0.151326}%
\pgfsetfillcolor{currentfill}%
\pgfsetlinewidth{0.000000pt}%
\definecolor{currentstroke}{rgb}{0.000000,0.000000,0.000000}%
\pgfsetstrokecolor{currentstroke}%
\pgfsetdash{}{0pt}%
\pgfpathmoveto{\pgfqpoint{4.390975in}{1.994543in}}%
\pgfpathlineto{\pgfqpoint{4.490975in}{1.994543in}}%
\pgfpathlineto{\pgfqpoint{4.490975in}{2.063543in}}%
\pgfpathlineto{\pgfqpoint{4.390975in}{2.063543in}}%
\pgfpathlineto{\pgfqpoint{4.390975in}{1.994543in}}%
\pgfusepath{fill}%
\end{pgfscope}%
\begin{pgfscope}%
\pgfpathrectangle{\pgfqpoint{4.390975in}{0.476543in}}{\pgfqpoint{0.100000in}{1.656000in}}%
\pgfusepath{clip}%
\pgfsetbuttcap%
\pgfsetroundjoin%
\definecolor{currentfill}{rgb}{0.685506,0.036909,0.149789}%
\pgfsetfillcolor{currentfill}%
\pgfsetlinewidth{0.000000pt}%
\definecolor{currentstroke}{rgb}{0.000000,0.000000,0.000000}%
\pgfsetstrokecolor{currentstroke}%
\pgfsetdash{}{0pt}%
\pgfpathmoveto{\pgfqpoint{4.390975in}{2.063543in}}%
\pgfpathlineto{\pgfqpoint{4.490975in}{2.063543in}}%
\pgfpathlineto{\pgfqpoint{4.490975in}{2.132543in}}%
\pgfpathlineto{\pgfqpoint{4.390975in}{2.132543in}}%
\pgfpathlineto{\pgfqpoint{4.390975in}{2.063543in}}%
\pgfusepath{fill}%
\end{pgfscope}%
\begin{pgfscope}%
\pgfsetbuttcap%
\pgfsetroundjoin%
\definecolor{currentfill}{rgb}{0.000000,0.000000,0.000000}%
\pgfsetfillcolor{currentfill}%
\pgfsetlinewidth{0.803000pt}%
\definecolor{currentstroke}{rgb}{0.000000,0.000000,0.000000}%
\pgfsetstrokecolor{currentstroke}%
\pgfsetdash{}{0pt}%
\pgfsys@defobject{currentmarker}{\pgfqpoint{0.000000in}{0.000000in}}{\pgfqpoint{0.048611in}{0.000000in}}{%
\pgfpathmoveto{\pgfqpoint{0.000000in}{0.000000in}}%
\pgfpathlineto{\pgfqpoint{0.048611in}{0.000000in}}%
\pgfusepath{stroke,fill}%
}%
\begin{pgfscope}%
\pgfsys@transformshift{4.490975in}{0.476543in}%
\pgfsys@useobject{currentmarker}{}%
\end{pgfscope}%
\end{pgfscope}%
\begin{pgfscope}%
\definecolor{textcolor}{rgb}{0.000000,0.000000,0.000000}%
\pgfsetstrokecolor{textcolor}%
\pgfsetfillcolor{textcolor}%
\pgftext[x=4.588197in, y=0.437963in, left, base]{\color{textcolor}{\rmfamily\fontsize{8.000000}{9.600000}\selectfont\catcode`\^=\active\def^{\ifmmode\sp\else\^{}\fi}\catcode`\%=\active\def%{\%}0.0}}%
\end{pgfscope}%
\begin{pgfscope}%
\pgfsetbuttcap%
\pgfsetroundjoin%
\definecolor{currentfill}{rgb}{0.000000,0.000000,0.000000}%
\pgfsetfillcolor{currentfill}%
\pgfsetlinewidth{0.803000pt}%
\definecolor{currentstroke}{rgb}{0.000000,0.000000,0.000000}%
\pgfsetstrokecolor{currentstroke}%
\pgfsetdash{}{0pt}%
\pgfsys@defobject{currentmarker}{\pgfqpoint{0.000000in}{0.000000in}}{\pgfqpoint{0.048611in}{0.000000in}}{%
\pgfpathmoveto{\pgfqpoint{0.000000in}{0.000000in}}%
\pgfpathlineto{\pgfqpoint{0.048611in}{0.000000in}}%
\pgfusepath{stroke,fill}%
}%
\begin{pgfscope}%
\pgfsys@transformshift{4.490975in}{1.211630in}%
\pgfsys@useobject{currentmarker}{}%
\end{pgfscope}%
\end{pgfscope}%
\begin{pgfscope}%
\definecolor{textcolor}{rgb}{0.000000,0.000000,0.000000}%
\pgfsetstrokecolor{textcolor}%
\pgfsetfillcolor{textcolor}%
\pgftext[x=4.588197in, y=1.173050in, left, base]{\color{textcolor}{\rmfamily\fontsize{8.000000}{9.600000}\selectfont\catcode`\^=\active\def^{\ifmmode\sp\else\^{}\fi}\catcode`\%=\active\def%{\%}0.4}}%
\end{pgfscope}%
\begin{pgfscope}%
\pgfsetbuttcap%
\pgfsetroundjoin%
\definecolor{currentfill}{rgb}{0.000000,0.000000,0.000000}%
\pgfsetfillcolor{currentfill}%
\pgfsetlinewidth{0.803000pt}%
\definecolor{currentstroke}{rgb}{0.000000,0.000000,0.000000}%
\pgfsetstrokecolor{currentstroke}%
\pgfsetdash{}{0pt}%
\pgfsys@defobject{currentmarker}{\pgfqpoint{0.000000in}{0.000000in}}{\pgfqpoint{0.048611in}{0.000000in}}{%
\pgfpathmoveto{\pgfqpoint{0.000000in}{0.000000in}}%
\pgfpathlineto{\pgfqpoint{0.048611in}{0.000000in}}%
\pgfusepath{stroke,fill}%
}%
\begin{pgfscope}%
\pgfsys@transformshift{4.490975in}{1.946717in}%
\pgfsys@useobject{currentmarker}{}%
\end{pgfscope}%
\end{pgfscope}%
\begin{pgfscope}%
\definecolor{textcolor}{rgb}{0.000000,0.000000,0.000000}%
\pgfsetstrokecolor{textcolor}%
\pgfsetfillcolor{textcolor}%
\pgftext[x=4.588197in, y=1.908136in, left, base]{\color{textcolor}{\rmfamily\fontsize{8.000000}{9.600000}\selectfont\catcode`\^=\active\def^{\ifmmode\sp\else\^{}\fi}\catcode`\%=\active\def%{\%}0.8}}%
\end{pgfscope}%
\begin{pgfscope}%
\definecolor{textcolor}{rgb}{0.000000,0.000000,0.000000}%
\pgfsetstrokecolor{textcolor}%
\pgfsetfillcolor{textcolor}%
\pgftext[x=4.794603in,y=1.304543in,,top,rotate=90.000000]{\color{textcolor}{\rmfamily\fontsize{9.000000}{10.800000}\selectfont\catcode`\^=\active\def^{\ifmmode\sp\else\^{}\fi}\catcode`\%=\active\def%{\%}$P_{D_s}$}}%
\end{pgfscope}%
\begin{pgfscope}%
\pgfsetrectcap%
\pgfsetmiterjoin%
\pgfsetlinewidth{0.803000pt}%
\definecolor{currentstroke}{rgb}{0.000000,0.000000,0.000000}%
\pgfsetstrokecolor{currentstroke}%
\pgfsetdash{}{0pt}%
\pgfpathmoveto{\pgfqpoint{4.390975in}{0.476543in}}%
\pgfpathlineto{\pgfqpoint{4.440975in}{0.476543in}}%
\pgfpathlineto{\pgfqpoint{4.490975in}{0.476543in}}%
\pgfpathlineto{\pgfqpoint{4.490975in}{2.132543in}}%
\pgfpathlineto{\pgfqpoint{4.440975in}{2.132543in}}%
\pgfpathlineto{\pgfqpoint{4.390975in}{2.132543in}}%
\pgfpathlineto{\pgfqpoint{4.390975in}{0.476543in}}%
\pgfpathclose%
\pgfusepath{stroke}%
\end{pgfscope}%
\end{pgfpicture}%
\makeatother%
\endgroup%
}
	\caption{(En color en l\'inea) Poblaci\'on del estado oscuro $D_s$ tras un pulso Raman desde $\ket{G,1}$ en funci\'on de la desintonizaci\'on de dos fotones $\delta_2/\omega_b$ y de la duraci\'on del pulso $t\omega_b$, para (a) $D_1$ con $\Omega/\omega_b=0.197$ y (b) $D_2$ con $\Omega/\omega_b=0.290$. Las estrellas marcan los pulsos optimizados num\'ericamente, $(\delta_2/\omega_b,\,d\,\omega_b)=(-0.020,\,170)$ y $(-0.025,\,150)$. Modelo reducido con p\'erdidas y par\'ametros del sistema: $G_1/\omega_b=0.35$ y $G_2/\omega_b=0.60$, que dan $\varepsilon_i/\omega_b=(-0.400,-0.275,-0.275)$, $J_{12}/\omega_b=J_{13}/\omega_b=0.200$ y $J_{23}/\omega_b=0.075$; $\lambda_i/\omega_b=(0.20,0.14,0)$, que dan $g_1/\omega_b=0.0572$ y $g_2/\omega_b=0.0613$; $\kappa/\omega_b=10^{-3}$, $\gamma/\omega_b=2\times10^{-4}$, $\gamma_\varphi/\omega_b=4\times10^{-4}$, $\chi=1$.}
	\label{fig:raman_map}
\end{figure}

La Figura~\ref{fig:raman_map} muestra la poblaci\'on del estado oscuro objetivo tras un pulso de escritura desde $\ket{G,1}$, en funci\'on de la desintonizaci\'on de dos fotones y de la duraci\'on del pulso. Cada mapa tiene la forma de la oscilaci\'on de Rabi de la Ec.~\eqref{eq:heff}: en la resonancia la poblaci\'on crece hasta un m\'aximo cerca de $t_\pi$ y vuelve a caer; al desintonizar, las oscilaciones se aceleran y su amplitud disminuye. La poblaci\'on m\'axima en funci\'on de la desintonizaci\'on es una lorentziana de ancho total $4\Omega_{\rm eff,s}$, de $0.035\,\omega_b$ para $D_1$ y $0.045\,\omega_b$ para $D_2$ en los mapas, y hay un solo m\'aximo por ranura en el intervalo explorado, de modo que la frecuencia del l\'aser selecciona la ranura. Ese m\'aximo no est\'a en $\delta_2=0$ sino en $\delta_2^\star$, lo que obliga a usar la Ec.~\eqref{eq:stark}: el desplazamiento de primer orden $(g_s^2-\Omega^2)/\Delta_I$ falla la resonancia de $D_2$ por un factor dos y la de $D_1$ en cerca del $30\%$, mientras que la Ec.~\eqref{eq:stark} reproduce la de ambas ranuras dentro de $10^{-3}\,\omega_b$. La poblaci\'on m\'axima, cerca de $0.9$, queda por debajo de la unidad porque el fon\'on decae mientras act\'ua el pulso, Ec.~\eqref{eq:etaw}, y por la mezcla del estado brillante; la duraci\'on \'optima, marcada con estrellas, coincide con $t_\pi$ dentro del $10\%$. La eficiencia de la Ec.~\eqref{eq:etaw} extendida a un pulso finito de duraci\'on $d$, $\eta_w\simeq\sin^2(\Omega_{\rm eff}d)\,\bigl[1-\kappa\bigl(d/2+\sin 2\Omega_{\rm eff}d/4\Omega_{\rm eff}\bigr)\bigr]$, coincide con la simulaci\'on para $D_1$ y la supera para $D_2$ en una cantidad del orden de la mezcla del estado brillante, $(\Omega/\Delta_I)^2$, que se pierde cuando el pulso se apaga abruptamente.

La Figura~\ref{fig:raman_dyn} es un corte del mapa en la desintonizaci\'on \'optima. La poblaci\'on de la ranura objetivo sigue la oscilaci\'on de la Ec.~\eqref{eq:heff} con un peque\~no rizado r\'apido, y la ranura no direccionada permanece cerca de un por ciento, con un m\'aximo simulado de $0.011$ y $0.009$, a un factor dos de la estimaci\'on de la Ec.~\eqref{eq:leak}.

\begin{figure}[t]
\centering
\resizebox{\columnwidth}{!}{%% Creator: Matplotlib, PGF backend
%%
%% To include the figure in your LaTeX document, write
%%   \input{<filename>.pgf}
%%
%% Make sure the required packages are loaded in your preamble
%%   \usepackage{pgf}
%%
%% Also ensure that all the required font packages are loaded; for instance,
%% the lmodern package is sometimes necessary when using math font.
%%   \usepackage{lmodern}
%%
%% Figures using additional raster images can only be included by \input if
%% they are in the same directory as the main LaTeX file. For loading figures
%% from other directories you can use the `import` package
%%   \usepackage{import}
%%
%% and then include the figures with
%%   \import{<path to file>}{<filename>.pgf}
%%
%% Matplotlib used the following preamble
%%   \def\mathdefault#1{#1}
%%   \everymath=\expandafter{\the\everymath\displaystyle}
%%   \IfFileExists{scrextend.sty}{
%%     \usepackage[fontsize=8.000000pt]{scrextend}
%%   }{
%%     \renewcommand{\normalsize}{\fontsize{8.000000}{9.600000}\selectfont}
%%     \normalsize
%%   }
%%   \usepackage{amsmath,amssymb}
%%   \makeatletter\@ifpackageloaded{underscore}{}{\usepackage[strings]{underscore}}\makeatother
%%
\begingroup%
\makeatletter%
\begin{pgfpicture}%
\pgfpathrectangle{\pgfpointorigin}{\pgfqpoint{4.721286in}{1.831877in}}%
\pgfusepath{use as bounding box, clip}%
\begin{pgfscope}%
\pgfsetbuttcap%
\pgfsetmiterjoin%
\definecolor{currentfill}{rgb}{1.000000,1.000000,1.000000}%
\pgfsetfillcolor{currentfill}%
\pgfsetlinewidth{0.000000pt}%
\definecolor{currentstroke}{rgb}{1.000000,1.000000,1.000000}%
\pgfsetstrokecolor{currentstroke}%
\pgfsetdash{}{0pt}%
\pgfpathmoveto{\pgfqpoint{0.000000in}{0.000000in}}%
\pgfpathlineto{\pgfqpoint{4.721286in}{0.000000in}}%
\pgfpathlineto{\pgfqpoint{4.721286in}{1.831877in}}%
\pgfpathlineto{\pgfqpoint{0.000000in}{1.831877in}}%
\pgfpathlineto{\pgfqpoint{0.000000in}{0.000000in}}%
\pgfpathclose%
\pgfusepath{fill}%
\end{pgfscope}%
\begin{pgfscope}%
\pgfsetbuttcap%
\pgfsetmiterjoin%
\definecolor{currentfill}{rgb}{1.000000,1.000000,1.000000}%
\pgfsetfillcolor{currentfill}%
\pgfsetlinewidth{0.000000pt}%
\definecolor{currentstroke}{rgb}{0.000000,0.000000,0.000000}%
\pgfsetstrokecolor{currentstroke}%
\pgfsetstrokeopacity{0.000000}%
\pgfsetdash{}{0pt}%
\pgfpathmoveto{\pgfqpoint{0.514740in}{0.462654in}}%
\pgfpathlineto{\pgfqpoint{2.424714in}{0.462654in}}%
\pgfpathlineto{\pgfqpoint{2.424714in}{1.731877in}}%
\pgfpathlineto{\pgfqpoint{0.514740in}{1.731877in}}%
\pgfpathlineto{\pgfqpoint{0.514740in}{0.462654in}}%
\pgfpathclose%
\pgfusepath{fill}%
\end{pgfscope}%
\begin{pgfscope}%
\pgfsetbuttcap%
\pgfsetroundjoin%
\definecolor{currentfill}{rgb}{0.000000,0.000000,0.000000}%
\pgfsetfillcolor{currentfill}%
\pgfsetlinewidth{0.803000pt}%
\definecolor{currentstroke}{rgb}{0.000000,0.000000,0.000000}%
\pgfsetstrokecolor{currentstroke}%
\pgfsetdash{}{0pt}%
\pgfsys@defobject{currentmarker}{\pgfqpoint{0.000000in}{-0.048611in}}{\pgfqpoint{0.000000in}{0.000000in}}{%
\pgfpathmoveto{\pgfqpoint{0.000000in}{0.000000in}}%
\pgfpathlineto{\pgfqpoint{0.000000in}{-0.048611in}}%
\pgfusepath{stroke,fill}%
}%
\begin{pgfscope}%
\pgfsys@transformshift{0.514740in}{0.462654in}%
\pgfsys@useobject{currentmarker}{}%
\end{pgfscope}%
\end{pgfscope}%
\begin{pgfscope}%
\definecolor{textcolor}{rgb}{0.000000,0.000000,0.000000}%
\pgfsetstrokecolor{textcolor}%
\pgfsetfillcolor{textcolor}%
\pgftext[x=0.514740in,y=0.365432in,,top]{\color{textcolor}{\rmfamily\fontsize{8.000000}{9.600000}\selectfont\catcode`\^=\active\def^{\ifmmode\sp\else\^{}\fi}\catcode`\%=\active\def%{\%}0}}%
\end{pgfscope}%
\begin{pgfscope}%
\pgfsetbuttcap%
\pgfsetroundjoin%
\definecolor{currentfill}{rgb}{0.000000,0.000000,0.000000}%
\pgfsetfillcolor{currentfill}%
\pgfsetlinewidth{0.803000pt}%
\definecolor{currentstroke}{rgb}{0.000000,0.000000,0.000000}%
\pgfsetstrokecolor{currentstroke}%
\pgfsetdash{}{0pt}%
\pgfsys@defobject{currentmarker}{\pgfqpoint{0.000000in}{-0.048611in}}{\pgfqpoint{0.000000in}{0.000000in}}{%
\pgfpathmoveto{\pgfqpoint{0.000000in}{0.000000in}}%
\pgfpathlineto{\pgfqpoint{0.000000in}{-0.048611in}}%
\pgfusepath{stroke,fill}%
}%
\begin{pgfscope}%
\pgfsys@transformshift{0.896735in}{0.462654in}%
\pgfsys@useobject{currentmarker}{}%
\end{pgfscope}%
\end{pgfscope}%
\begin{pgfscope}%
\definecolor{textcolor}{rgb}{0.000000,0.000000,0.000000}%
\pgfsetstrokecolor{textcolor}%
\pgfsetfillcolor{textcolor}%
\pgftext[x=0.896735in,y=0.365432in,,top]{\color{textcolor}{\rmfamily\fontsize{8.000000}{9.600000}\selectfont\catcode`\^=\active\def^{\ifmmode\sp\else\^{}\fi}\catcode`\%=\active\def%{\%}50}}%
\end{pgfscope}%
\begin{pgfscope}%
\pgfsetbuttcap%
\pgfsetroundjoin%
\definecolor{currentfill}{rgb}{0.000000,0.000000,0.000000}%
\pgfsetfillcolor{currentfill}%
\pgfsetlinewidth{0.803000pt}%
\definecolor{currentstroke}{rgb}{0.000000,0.000000,0.000000}%
\pgfsetstrokecolor{currentstroke}%
\pgfsetdash{}{0pt}%
\pgfsys@defobject{currentmarker}{\pgfqpoint{0.000000in}{-0.048611in}}{\pgfqpoint{0.000000in}{0.000000in}}{%
\pgfpathmoveto{\pgfqpoint{0.000000in}{0.000000in}}%
\pgfpathlineto{\pgfqpoint{0.000000in}{-0.048611in}}%
\pgfusepath{stroke,fill}%
}%
\begin{pgfscope}%
\pgfsys@transformshift{1.278729in}{0.462654in}%
\pgfsys@useobject{currentmarker}{}%
\end{pgfscope}%
\end{pgfscope}%
\begin{pgfscope}%
\definecolor{textcolor}{rgb}{0.000000,0.000000,0.000000}%
\pgfsetstrokecolor{textcolor}%
\pgfsetfillcolor{textcolor}%
\pgftext[x=1.278729in,y=0.365432in,,top]{\color{textcolor}{\rmfamily\fontsize{8.000000}{9.600000}\selectfont\catcode`\^=\active\def^{\ifmmode\sp\else\^{}\fi}\catcode`\%=\active\def%{\%}100}}%
\end{pgfscope}%
\begin{pgfscope}%
\pgfsetbuttcap%
\pgfsetroundjoin%
\definecolor{currentfill}{rgb}{0.000000,0.000000,0.000000}%
\pgfsetfillcolor{currentfill}%
\pgfsetlinewidth{0.803000pt}%
\definecolor{currentstroke}{rgb}{0.000000,0.000000,0.000000}%
\pgfsetstrokecolor{currentstroke}%
\pgfsetdash{}{0pt}%
\pgfsys@defobject{currentmarker}{\pgfqpoint{0.000000in}{-0.048611in}}{\pgfqpoint{0.000000in}{0.000000in}}{%
\pgfpathmoveto{\pgfqpoint{0.000000in}{0.000000in}}%
\pgfpathlineto{\pgfqpoint{0.000000in}{-0.048611in}}%
\pgfusepath{stroke,fill}%
}%
\begin{pgfscope}%
\pgfsys@transformshift{1.660724in}{0.462654in}%
\pgfsys@useobject{currentmarker}{}%
\end{pgfscope}%
\end{pgfscope}%
\begin{pgfscope}%
\definecolor{textcolor}{rgb}{0.000000,0.000000,0.000000}%
\pgfsetstrokecolor{textcolor}%
\pgfsetfillcolor{textcolor}%
\pgftext[x=1.660724in,y=0.365432in,,top]{\color{textcolor}{\rmfamily\fontsize{8.000000}{9.600000}\selectfont\catcode`\^=\active\def^{\ifmmode\sp\else\^{}\fi}\catcode`\%=\active\def%{\%}150}}%
\end{pgfscope}%
\begin{pgfscope}%
\pgfsetbuttcap%
\pgfsetroundjoin%
\definecolor{currentfill}{rgb}{0.000000,0.000000,0.000000}%
\pgfsetfillcolor{currentfill}%
\pgfsetlinewidth{0.803000pt}%
\definecolor{currentstroke}{rgb}{0.000000,0.000000,0.000000}%
\pgfsetstrokecolor{currentstroke}%
\pgfsetdash{}{0pt}%
\pgfsys@defobject{currentmarker}{\pgfqpoint{0.000000in}{-0.048611in}}{\pgfqpoint{0.000000in}{0.000000in}}{%
\pgfpathmoveto{\pgfqpoint{0.000000in}{0.000000in}}%
\pgfpathlineto{\pgfqpoint{0.000000in}{-0.048611in}}%
\pgfusepath{stroke,fill}%
}%
\begin{pgfscope}%
\pgfsys@transformshift{2.042719in}{0.462654in}%
\pgfsys@useobject{currentmarker}{}%
\end{pgfscope}%
\end{pgfscope}%
\begin{pgfscope}%
\definecolor{textcolor}{rgb}{0.000000,0.000000,0.000000}%
\pgfsetstrokecolor{textcolor}%
\pgfsetfillcolor{textcolor}%
\pgftext[x=2.042719in,y=0.365432in,,top]{\color{textcolor}{\rmfamily\fontsize{8.000000}{9.600000}\selectfont\catcode`\^=\active\def^{\ifmmode\sp\else\^{}\fi}\catcode`\%=\active\def%{\%}200}}%
\end{pgfscope}%
\begin{pgfscope}%
\pgfsetbuttcap%
\pgfsetroundjoin%
\definecolor{currentfill}{rgb}{0.000000,0.000000,0.000000}%
\pgfsetfillcolor{currentfill}%
\pgfsetlinewidth{0.803000pt}%
\definecolor{currentstroke}{rgb}{0.000000,0.000000,0.000000}%
\pgfsetstrokecolor{currentstroke}%
\pgfsetdash{}{0pt}%
\pgfsys@defobject{currentmarker}{\pgfqpoint{0.000000in}{-0.048611in}}{\pgfqpoint{0.000000in}{0.000000in}}{%
\pgfpathmoveto{\pgfqpoint{0.000000in}{0.000000in}}%
\pgfpathlineto{\pgfqpoint{0.000000in}{-0.048611in}}%
\pgfusepath{stroke,fill}%
}%
\begin{pgfscope}%
\pgfsys@transformshift{2.424714in}{0.462654in}%
\pgfsys@useobject{currentmarker}{}%
\end{pgfscope}%
\end{pgfscope}%
\begin{pgfscope}%
\definecolor{textcolor}{rgb}{0.000000,0.000000,0.000000}%
\pgfsetstrokecolor{textcolor}%
\pgfsetfillcolor{textcolor}%
\pgftext[x=2.424714in,y=0.365432in,,top]{\color{textcolor}{\rmfamily\fontsize{8.000000}{9.600000}\selectfont\catcode`\^=\active\def^{\ifmmode\sp\else\^{}\fi}\catcode`\%=\active\def%{\%}250}}%
\end{pgfscope}%
\begin{pgfscope}%
\definecolor{textcolor}{rgb}{0.000000,0.000000,0.000000}%
\pgfsetstrokecolor{textcolor}%
\pgfsetfillcolor{textcolor}%
\pgftext[x=1.469727in,y=0.211111in,,top]{\color{textcolor}{\rmfamily\fontsize{9.000000}{10.800000}\selectfont\catcode`\^=\active\def^{\ifmmode\sp\else\^{}\fi}\catcode`\%=\active\def%{\%}$\omega_b\,t$}}%
\end{pgfscope}%
\begin{pgfscope}%
\pgfsetbuttcap%
\pgfsetroundjoin%
\definecolor{currentfill}{rgb}{0.000000,0.000000,0.000000}%
\pgfsetfillcolor{currentfill}%
\pgfsetlinewidth{0.803000pt}%
\definecolor{currentstroke}{rgb}{0.000000,0.000000,0.000000}%
\pgfsetstrokecolor{currentstroke}%
\pgfsetdash{}{0pt}%
\pgfsys@defobject{currentmarker}{\pgfqpoint{-0.048611in}{0.000000in}}{\pgfqpoint{-0.000000in}{0.000000in}}{%
\pgfpathmoveto{\pgfqpoint{-0.000000in}{0.000000in}}%
\pgfpathlineto{\pgfqpoint{-0.048611in}{0.000000in}}%
\pgfusepath{stroke,fill}%
}%
\begin{pgfscope}%
\pgfsys@transformshift{0.514740in}{0.462654in}%
\pgfsys@useobject{currentmarker}{}%
\end{pgfscope}%
\end{pgfscope}%
\begin{pgfscope}%
\definecolor{textcolor}{rgb}{0.000000,0.000000,0.000000}%
\pgfsetstrokecolor{textcolor}%
\pgfsetfillcolor{textcolor}%
\pgftext[x=0.358489in, y=0.424074in, left, base]{\color{textcolor}{\rmfamily\fontsize{8.000000}{9.600000}\selectfont\catcode`\^=\active\def^{\ifmmode\sp\else\^{}\fi}\catcode`\%=\active\def%{\%}$0$}}%
\end{pgfscope}%
\begin{pgfscope}%
\pgfsetbuttcap%
\pgfsetroundjoin%
\definecolor{currentfill}{rgb}{0.000000,0.000000,0.000000}%
\pgfsetfillcolor{currentfill}%
\pgfsetlinewidth{0.803000pt}%
\definecolor{currentstroke}{rgb}{0.000000,0.000000,0.000000}%
\pgfsetstrokecolor{currentstroke}%
\pgfsetdash{}{0pt}%
\pgfsys@defobject{currentmarker}{\pgfqpoint{-0.048611in}{0.000000in}}{\pgfqpoint{-0.000000in}{0.000000in}}{%
\pgfpathmoveto{\pgfqpoint{-0.000000in}{0.000000in}}%
\pgfpathlineto{\pgfqpoint{-0.048611in}{0.000000in}}%
\pgfusepath{stroke,fill}%
}%
\begin{pgfscope}%
\pgfsys@transformshift{0.514740in}{1.067046in}%
\pgfsys@useobject{currentmarker}{}%
\end{pgfscope}%
\end{pgfscope}%
\begin{pgfscope}%
\definecolor{textcolor}{rgb}{0.000000,0.000000,0.000000}%
\pgfsetstrokecolor{textcolor}%
\pgfsetfillcolor{textcolor}%
\pgftext[x=0.266667in, y=1.028466in, left, base]{\color{textcolor}{\rmfamily\fontsize{8.000000}{9.600000}\selectfont\catcode`\^=\active\def^{\ifmmode\sp\else\^{}\fi}\catcode`\%=\active\def%{\%}$0.5$}}%
\end{pgfscope}%
\begin{pgfscope}%
\pgfsetbuttcap%
\pgfsetroundjoin%
\definecolor{currentfill}{rgb}{0.000000,0.000000,0.000000}%
\pgfsetfillcolor{currentfill}%
\pgfsetlinewidth{0.803000pt}%
\definecolor{currentstroke}{rgb}{0.000000,0.000000,0.000000}%
\pgfsetstrokecolor{currentstroke}%
\pgfsetdash{}{0pt}%
\pgfsys@defobject{currentmarker}{\pgfqpoint{-0.048611in}{0.000000in}}{\pgfqpoint{-0.000000in}{0.000000in}}{%
\pgfpathmoveto{\pgfqpoint{-0.000000in}{0.000000in}}%
\pgfpathlineto{\pgfqpoint{-0.048611in}{0.000000in}}%
\pgfusepath{stroke,fill}%
}%
\begin{pgfscope}%
\pgfsys@transformshift{0.514740in}{1.671438in}%
\pgfsys@useobject{currentmarker}{}%
\end{pgfscope}%
\end{pgfscope}%
\begin{pgfscope}%
\definecolor{textcolor}{rgb}{0.000000,0.000000,0.000000}%
\pgfsetstrokecolor{textcolor}%
\pgfsetfillcolor{textcolor}%
\pgftext[x=0.358489in, y=1.632857in, left, base]{\color{textcolor}{\rmfamily\fontsize{8.000000}{9.600000}\selectfont\catcode`\^=\active\def^{\ifmmode\sp\else\^{}\fi}\catcode`\%=\active\def%{\%}$1$}}%
\end{pgfscope}%
\begin{pgfscope}%
\definecolor{textcolor}{rgb}{0.000000,0.000000,0.000000}%
\pgfsetstrokecolor{textcolor}%
\pgfsetfillcolor{textcolor}%
\pgftext[x=0.211111in,y=1.097266in,,bottom,rotate=90.000000]{\color{textcolor}{\rmfamily\fontsize{9.000000}{10.800000}\selectfont\catcode`\^=\active\def^{\ifmmode\sp\else\^{}\fi}\catcode`\%=\active\def%{\%}Populations}}%
\end{pgfscope}%
\begin{pgfscope}%
\pgfpathrectangle{\pgfqpoint{0.514740in}{0.462654in}}{\pgfqpoint{1.909974in}{1.269222in}}%
\pgfusepath{clip}%
\pgfsetrectcap%
\pgfsetroundjoin%
\pgfsetlinewidth{0.903375pt}%
\definecolor{currentstroke}{rgb}{0.121569,0.466667,0.705882}%
\pgfsetstrokecolor{currentstroke}%
\pgfsetdash{}{0pt}%
\pgfpathmoveto{\pgfqpoint{0.514740in}{0.462654in}}%
\pgfpathlineto{\pgfqpoint{0.530020in}{0.463634in}}%
\pgfpathlineto{\pgfqpoint{0.545299in}{0.465156in}}%
\pgfpathlineto{\pgfqpoint{0.560579in}{0.466251in}}%
\pgfpathlineto{\pgfqpoint{0.575859in}{0.469529in}}%
\pgfpathlineto{\pgfqpoint{0.591139in}{0.470776in}}%
\pgfpathlineto{\pgfqpoint{0.606419in}{0.477530in}}%
\pgfpathlineto{\pgfqpoint{0.621698in}{0.480684in}}%
\pgfpathlineto{\pgfqpoint{0.636978in}{0.486156in}}%
\pgfpathlineto{\pgfqpoint{0.652258in}{0.492854in}}%
\pgfpathlineto{\pgfqpoint{0.667538in}{0.496092in}}%
\pgfpathlineto{\pgfqpoint{0.682818in}{0.508475in}}%
\pgfpathlineto{\pgfqpoint{0.698097in}{0.511225in}}%
\pgfpathlineto{\pgfqpoint{0.713377in}{0.523838in}}%
\pgfpathlineto{\pgfqpoint{0.728657in}{0.529502in}}%
\pgfpathlineto{\pgfqpoint{0.743937in}{0.540822in}}%
\pgfpathlineto{\pgfqpoint{0.759217in}{0.553200in}}%
\pgfpathlineto{\pgfqpoint{0.774496in}{0.559487in}}%
\pgfpathlineto{\pgfqpoint{0.789776in}{0.577175in}}%
\pgfpathlineto{\pgfqpoint{0.805056in}{0.580810in}}%
\pgfpathlineto{\pgfqpoint{0.820336in}{0.602098in}}%
\pgfpathlineto{\pgfqpoint{0.835615in}{0.609558in}}%
\pgfpathlineto{\pgfqpoint{0.850895in}{0.625701in}}%
\pgfpathlineto{\pgfqpoint{0.866175in}{0.640471in}}%
\pgfpathlineto{\pgfqpoint{0.881455in}{0.649397in}}%
\pgfpathlineto{\pgfqpoint{0.896735in}{0.673148in}}%
\pgfpathlineto{\pgfqpoint{0.912014in}{0.677169in}}%
\pgfpathlineto{\pgfqpoint{0.927294in}{0.702804in}}%
\pgfpathlineto{\pgfqpoint{0.942574in}{0.709764in}}%
\pgfpathlineto{\pgfqpoint{0.957854in}{0.731936in}}%
\pgfpathlineto{\pgfqpoint{0.973134in}{0.749783in}}%
\pgfpathlineto{\pgfqpoint{0.988413in}{0.762721in}}%
\pgfpathlineto{\pgfqpoint{1.003693in}{0.789464in}}%
\pgfpathlineto{\pgfqpoint{1.018973in}{0.793837in}}%
\pgfpathlineto{\pgfqpoint{1.034253in}{0.825422in}}%
\pgfpathlineto{\pgfqpoint{1.049533in}{0.830594in}}%
\pgfpathlineto{\pgfqpoint{1.064812in}{0.855587in}}%
\pgfpathlineto{\pgfqpoint{1.080092in}{0.871039in}}%
\pgfpathlineto{\pgfqpoint{1.095372in}{0.885187in}}%
\pgfpathlineto{\pgfqpoint{1.110652in}{0.915072in}}%
\pgfpathlineto{\pgfqpoint{1.125932in}{0.921253in}}%
\pgfpathlineto{\pgfqpoint{1.141211in}{0.956000in}}%
\pgfpathlineto{\pgfqpoint{1.156491in}{0.960098in}}%
\pgfpathlineto{\pgfqpoint{1.171771in}{0.990304in}}%
\pgfpathlineto{\pgfqpoint{1.187051in}{1.003834in}}%
\pgfpathlineto{\pgfqpoint{1.202330in}{1.019923in}}%
\pgfpathlineto{\pgfqpoint{1.217610in}{1.047372in}}%
\pgfpathlineto{\pgfqpoint{1.232890in}{1.050553in}}%
\pgfpathlineto{\pgfqpoint{1.248170in}{1.087392in}}%
\pgfpathlineto{\pgfqpoint{1.263450in}{1.089493in}}%
\pgfpathlineto{\pgfqpoint{1.278729in}{1.121595in}}%
\pgfpathlineto{\pgfqpoint{1.294009in}{1.131680in}}%
\pgfpathlineto{\pgfqpoint{1.309289in}{1.149770in}}%
\pgfpathlineto{\pgfqpoint{1.324569in}{1.175533in}}%
\pgfpathlineto{\pgfqpoint{1.339849in}{1.179495in}}%
\pgfpathlineto{\pgfqpoint{1.355128in}{1.215525in}}%
\pgfpathlineto{\pgfqpoint{1.370408in}{1.213381in}}%
\pgfpathlineto{\pgfqpoint{1.385688in}{1.247458in}}%
\pgfpathlineto{\pgfqpoint{1.400968in}{1.253365in}}%
\pgfpathlineto{\pgfqpoint{1.416248in}{1.271952in}}%
\pgfpathlineto{\pgfqpoint{1.431527in}{1.291748in}}%
\pgfpathlineto{\pgfqpoint{1.446807in}{1.292962in}}%
\pgfpathlineto{\pgfqpoint{1.462087in}{1.326134in}}%
\pgfpathlineto{\pgfqpoint{1.477367in}{1.320978in}}%
\pgfpathlineto{\pgfqpoint{1.492647in}{1.354588in}}%
\pgfpathlineto{\pgfqpoint{1.507926in}{1.356033in}}%
\pgfpathlineto{\pgfqpoint{1.523206in}{1.376908in}}%
\pgfpathlineto{\pgfqpoint{1.538486in}{1.392984in}}%
\pgfpathlineto{\pgfqpoint{1.553766in}{1.394413in}}%
\pgfpathlineto{\pgfqpoint{1.569045in}{1.422006in}}%
\pgfpathlineto{\pgfqpoint{1.584325in}{1.410923in}}%
\pgfpathlineto{\pgfqpoint{1.599605in}{1.441045in}}%
\pgfpathlineto{\pgfqpoint{1.614885in}{1.434464in}}%
\pgfpathlineto{\pgfqpoint{1.630165in}{1.453701in}}%
\pgfpathlineto{\pgfqpoint{1.645444in}{1.463276in}}%
\pgfpathlineto{\pgfqpoint{1.660724in}{1.464798in}}%
\pgfpathlineto{\pgfqpoint{1.676004in}{1.489602in}}%
\pgfpathlineto{\pgfqpoint{1.691284in}{1.477507in}}%
\pgfpathlineto{\pgfqpoint{1.706564in}{1.504671in}}%
\pgfpathlineto{\pgfqpoint{1.721843in}{1.491498in}}%
\pgfpathlineto{\pgfqpoint{1.737123in}{1.508902in}}%
\pgfpathlineto{\pgfqpoint{1.752403in}{1.509114in}}%
\pgfpathlineto{\pgfqpoint{1.767683in}{1.508332in}}%
\pgfpathlineto{\pgfqpoint{1.782963in}{1.525177in}}%
\pgfpathlineto{\pgfqpoint{1.798242in}{1.508998in}}%
\pgfpathlineto{\pgfqpoint{1.813522in}{1.533295in}}%
\pgfpathlineto{\pgfqpoint{1.828802in}{1.515166in}}%
\pgfpathlineto{\pgfqpoint{1.844082in}{1.530610in}}%
\pgfpathlineto{\pgfqpoint{1.859362in}{1.523834in}}%
\pgfpathlineto{\pgfqpoint{1.874641in}{1.522286in}}%
\pgfpathlineto{\pgfqpoint{1.889921in}{1.532449in}}%
\pgfpathlineto{\pgfqpoint{1.905201in}{1.514830in}}%
\pgfpathlineto{\pgfqpoint{1.920481in}{1.533625in}}%
\pgfpathlineto{\pgfqpoint{1.935760in}{1.509339in}}%
\pgfpathlineto{\pgfqpoint{1.951040in}{1.522177in}}%
\pgfpathlineto{\pgfqpoint{1.966320in}{1.507128in}}%
\pgfpathlineto{\pgfqpoint{1.981600in}{1.502793in}}%
\pgfpathlineto{\pgfqpoint{1.996880in}{1.504114in}}%
\pgfpathlineto{\pgfqpoint{2.012159in}{1.483630in}}%
\pgfpathlineto{\pgfqpoint{2.027439in}{1.498098in}}%
\pgfpathlineto{\pgfqpoint{2.042719in}{1.472364in}}%
\pgfpathlineto{\pgfqpoint{2.057999in}{1.483912in}}%
\pgfpathlineto{\pgfqpoint{2.073279in}{1.463490in}}%
\pgfpathlineto{\pgfqpoint{2.088558in}{1.458416in}}%
\pgfpathlineto{\pgfqpoint{2.103838in}{1.452439in}}%
\pgfpathlineto{\pgfqpoint{2.119118in}{1.429431in}}%
\pgfpathlineto{\pgfqpoint{2.134398in}{1.434767in}}%
\pgfpathlineto{\pgfqpoint{2.149678in}{1.403344in}}%
\pgfpathlineto{\pgfqpoint{2.164957in}{1.411323in}}%
\pgfpathlineto{\pgfqpoint{2.180237in}{1.386747in}}%
\pgfpathlineto{\pgfqpoint{2.195517in}{1.382618in}}%
\pgfpathlineto{\pgfqpoint{2.210797in}{1.371813in}}%
\pgfpathlineto{\pgfqpoint{2.226077in}{1.349166in}}%
\pgfpathlineto{\pgfqpoint{2.241356in}{1.350604in}}%
\pgfpathlineto{\pgfqpoint{2.256636in}{1.318298in}}%
\pgfpathlineto{\pgfqpoint{2.271916in}{1.320800in}}%
\pgfpathlineto{\pgfqpoint{2.287196in}{1.290593in}}%
\pgfpathlineto{\pgfqpoint{2.302475in}{1.284491in}}%
\pgfpathlineto{\pgfqpoint{2.317755in}{1.267601in}}%
\pgfpathlineto{\pgfqpoint{2.333035in}{1.246293in}}%
\pgfpathlineto{\pgfqpoint{2.348315in}{1.242392in}}%
\pgfpathlineto{\pgfqpoint{2.363595in}{1.209134in}}%
\pgfpathlineto{\pgfqpoint{2.378874in}{1.209679in}}%
\pgfpathlineto{\pgfqpoint{2.394154in}{1.178759in}}%
\pgfpathlineto{\pgfqpoint{2.409434in}{1.171870in}}%
\pgfpathlineto{\pgfqpoint{2.424714in}{1.151518in}}%
\pgfusepath{stroke}%
\end{pgfscope}%
\begin{pgfscope}%
\pgfpathrectangle{\pgfqpoint{0.514740in}{0.462654in}}{\pgfqpoint{1.909974in}{1.269222in}}%
\pgfusepath{clip}%
\pgfsetbuttcap%
\pgfsetroundjoin%
\pgfsetlinewidth{0.903375pt}%
\definecolor{currentstroke}{rgb}{1.000000,0.000000,0.000000}%
\pgfsetstrokecolor{currentstroke}%
\pgfsetdash{{3.330000pt}{1.440000pt}}{0.000000pt}%
\pgfpathmoveto{\pgfqpoint{0.514740in}{0.462654in}}%
\pgfpathlineto{\pgfqpoint{0.530020in}{0.463849in}}%
\pgfpathlineto{\pgfqpoint{0.545299in}{0.465677in}}%
\pgfpathlineto{\pgfqpoint{0.560579in}{0.467700in}}%
\pgfpathlineto{\pgfqpoint{0.575859in}{0.469517in}}%
\pgfpathlineto{\pgfqpoint{0.591139in}{0.468854in}}%
\pgfpathlineto{\pgfqpoint{0.606419in}{0.468930in}}%
\pgfpathlineto{\pgfqpoint{0.621698in}{0.466647in}}%
\pgfpathlineto{\pgfqpoint{0.636978in}{0.467741in}}%
\pgfpathlineto{\pgfqpoint{0.652258in}{0.465903in}}%
\pgfpathlineto{\pgfqpoint{0.667538in}{0.466907in}}%
\pgfpathlineto{\pgfqpoint{0.682818in}{0.465038in}}%
\pgfpathlineto{\pgfqpoint{0.698097in}{0.464587in}}%
\pgfpathlineto{\pgfqpoint{0.713377in}{0.465051in}}%
\pgfpathlineto{\pgfqpoint{0.728657in}{0.465763in}}%
\pgfpathlineto{\pgfqpoint{0.743937in}{0.469036in}}%
\pgfpathlineto{\pgfqpoint{0.759217in}{0.469119in}}%
\pgfpathlineto{\pgfqpoint{0.774496in}{0.471193in}}%
\pgfpathlineto{\pgfqpoint{0.789776in}{0.468820in}}%
\pgfpathlineto{\pgfqpoint{0.805056in}{0.467970in}}%
\pgfpathlineto{\pgfqpoint{0.820336in}{0.466212in}}%
\pgfpathlineto{\pgfqpoint{0.835615in}{0.465228in}}%
\pgfpathlineto{\pgfqpoint{0.850895in}{0.466663in}}%
\pgfpathlineto{\pgfqpoint{0.866175in}{0.466363in}}%
\pgfpathlineto{\pgfqpoint{0.881455in}{0.467910in}}%
\pgfpathlineto{\pgfqpoint{0.896735in}{0.467285in}}%
\pgfpathlineto{\pgfqpoint{0.912014in}{0.467281in}}%
\pgfpathlineto{\pgfqpoint{0.927294in}{0.468441in}}%
\pgfpathlineto{\pgfqpoint{0.942574in}{0.467967in}}%
\pgfpathlineto{\pgfqpoint{0.957854in}{0.471218in}}%
\pgfpathlineto{\pgfqpoint{0.973134in}{0.469436in}}%
\pgfpathlineto{\pgfqpoint{0.988413in}{0.470854in}}%
\pgfpathlineto{\pgfqpoint{1.003693in}{0.467252in}}%
\pgfpathlineto{\pgfqpoint{1.018973in}{0.466779in}}%
\pgfpathlineto{\pgfqpoint{1.034253in}{0.465581in}}%
\pgfpathlineto{\pgfqpoint{1.049533in}{0.465832in}}%
\pgfpathlineto{\pgfqpoint{1.064812in}{0.468184in}}%
\pgfpathlineto{\pgfqpoint{1.080092in}{0.468385in}}%
\pgfpathlineto{\pgfqpoint{1.095372in}{0.470376in}}%
\pgfpathlineto{\pgfqpoint{1.110652in}{0.470069in}}%
\pgfpathlineto{\pgfqpoint{1.125932in}{0.470479in}}%
\pgfpathlineto{\pgfqpoint{1.141211in}{0.470624in}}%
\pgfpathlineto{\pgfqpoint{1.156491in}{0.469737in}}%
\pgfpathlineto{\pgfqpoint{1.171771in}{0.469998in}}%
\pgfpathlineto{\pgfqpoint{1.187051in}{0.467829in}}%
\pgfpathlineto{\pgfqpoint{1.202330in}{0.467529in}}%
\pgfpathlineto{\pgfqpoint{1.217610in}{0.466032in}}%
\pgfpathlineto{\pgfqpoint{1.232890in}{0.466728in}}%
\pgfpathlineto{\pgfqpoint{1.248170in}{0.469091in}}%
\pgfpathlineto{\pgfqpoint{1.263450in}{0.470698in}}%
\pgfpathlineto{\pgfqpoint{1.278729in}{0.473320in}}%
\pgfpathlineto{\pgfqpoint{1.294009in}{0.471978in}}%
\pgfpathlineto{\pgfqpoint{1.309289in}{0.471880in}}%
\pgfpathlineto{\pgfqpoint{1.324569in}{0.468708in}}%
\pgfpathlineto{\pgfqpoint{1.339849in}{0.468431in}}%
\pgfpathlineto{\pgfqpoint{1.355128in}{0.467559in}}%
\pgfpathlineto{\pgfqpoint{1.370408in}{0.469012in}}%
\pgfpathlineto{\pgfqpoint{1.385688in}{0.469580in}}%
\pgfpathlineto{\pgfqpoint{1.400968in}{0.470427in}}%
\pgfpathlineto{\pgfqpoint{1.416248in}{0.470290in}}%
\pgfpathlineto{\pgfqpoint{1.431527in}{0.470277in}}%
\pgfpathlineto{\pgfqpoint{1.446807in}{0.470335in}}%
\pgfpathlineto{\pgfqpoint{1.462087in}{0.470929in}}%
\pgfpathlineto{\pgfqpoint{1.477367in}{0.471352in}}%
\pgfpathlineto{\pgfqpoint{1.492647in}{0.471924in}}%
\pgfpathlineto{\pgfqpoint{1.507926in}{0.470951in}}%
\pgfpathlineto{\pgfqpoint{1.523206in}{0.470361in}}%
\pgfpathlineto{\pgfqpoint{1.538486in}{0.469154in}}%
\pgfpathlineto{\pgfqpoint{1.553766in}{0.469566in}}%
\pgfpathlineto{\pgfqpoint{1.569045in}{0.469624in}}%
\pgfpathlineto{\pgfqpoint{1.584325in}{0.470811in}}%
\pgfpathlineto{\pgfqpoint{1.599605in}{0.470785in}}%
\pgfpathlineto{\pgfqpoint{1.614885in}{0.471082in}}%
\pgfpathlineto{\pgfqpoint{1.630165in}{0.471368in}}%
\pgfpathlineto{\pgfqpoint{1.645444in}{0.471496in}}%
\pgfpathlineto{\pgfqpoint{1.660724in}{0.472989in}}%
\pgfpathlineto{\pgfqpoint{1.676004in}{0.472176in}}%
\pgfpathlineto{\pgfqpoint{1.691284in}{0.472430in}}%
\pgfpathlineto{\pgfqpoint{1.706564in}{0.470253in}}%
\pgfpathlineto{\pgfqpoint{1.721843in}{0.469601in}}%
\pgfpathlineto{\pgfqpoint{1.737123in}{0.468016in}}%
\pgfpathlineto{\pgfqpoint{1.752403in}{0.469188in}}%
\pgfpathlineto{\pgfqpoint{1.767683in}{0.470695in}}%
\pgfpathlineto{\pgfqpoint{1.782963in}{0.473081in}}%
\pgfpathlineto{\pgfqpoint{1.798242in}{0.474913in}}%
\pgfpathlineto{\pgfqpoint{1.813522in}{0.474388in}}%
\pgfpathlineto{\pgfqpoint{1.828802in}{0.472611in}}%
\pgfpathlineto{\pgfqpoint{1.844082in}{0.471132in}}%
\pgfpathlineto{\pgfqpoint{1.859362in}{0.469154in}}%
\pgfpathlineto{\pgfqpoint{1.874641in}{0.469830in}}%
\pgfpathlineto{\pgfqpoint{1.889921in}{0.469817in}}%
\pgfpathlineto{\pgfqpoint{1.905201in}{0.471401in}}%
\pgfpathlineto{\pgfqpoint{1.920481in}{0.471275in}}%
\pgfpathlineto{\pgfqpoint{1.935760in}{0.472809in}}%
\pgfpathlineto{\pgfqpoint{1.951040in}{0.472552in}}%
\pgfpathlineto{\pgfqpoint{1.966320in}{0.472944in}}%
\pgfpathlineto{\pgfqpoint{1.981600in}{0.473023in}}%
\pgfpathlineto{\pgfqpoint{1.996880in}{0.472480in}}%
\pgfpathlineto{\pgfqpoint{2.012159in}{0.472032in}}%
\pgfpathlineto{\pgfqpoint{2.027439in}{0.470926in}}%
\pgfpathlineto{\pgfqpoint{2.042719in}{0.470418in}}%
\pgfpathlineto{\pgfqpoint{2.057999in}{0.470226in}}%
\pgfpathlineto{\pgfqpoint{2.073279in}{0.471401in}}%
\pgfpathlineto{\pgfqpoint{2.088558in}{0.472061in}}%
\pgfpathlineto{\pgfqpoint{2.103838in}{0.472951in}}%
\pgfpathlineto{\pgfqpoint{2.119118in}{0.472434in}}%
\pgfpathlineto{\pgfqpoint{2.134398in}{0.472470in}}%
\pgfpathlineto{\pgfqpoint{2.149678in}{0.471610in}}%
\pgfpathlineto{\pgfqpoint{2.164957in}{0.472136in}}%
\pgfpathlineto{\pgfqpoint{2.180237in}{0.472622in}}%
\pgfpathlineto{\pgfqpoint{2.195517in}{0.473202in}}%
\pgfpathlineto{\pgfqpoint{2.210797in}{0.472541in}}%
\pgfpathlineto{\pgfqpoint{2.226077in}{0.471546in}}%
\pgfpathlineto{\pgfqpoint{2.241356in}{0.469914in}}%
\pgfpathlineto{\pgfqpoint{2.256636in}{0.469822in}}%
\pgfpathlineto{\pgfqpoint{2.271916in}{0.470004in}}%
\pgfpathlineto{\pgfqpoint{2.287196in}{0.472526in}}%
\pgfpathlineto{\pgfqpoint{2.302475in}{0.472960in}}%
\pgfpathlineto{\pgfqpoint{2.317755in}{0.475405in}}%
\pgfpathlineto{\pgfqpoint{2.333035in}{0.474072in}}%
\pgfpathlineto{\pgfqpoint{2.348315in}{0.473546in}}%
\pgfpathlineto{\pgfqpoint{2.363595in}{0.471430in}}%
\pgfpathlineto{\pgfqpoint{2.378874in}{0.470261in}}%
\pgfpathlineto{\pgfqpoint{2.394154in}{0.470208in}}%
\pgfpathlineto{\pgfqpoint{2.409434in}{0.469808in}}%
\pgfpathlineto{\pgfqpoint{2.424714in}{0.471277in}}%
\pgfusepath{stroke}%
\end{pgfscope}%
\begin{pgfscope}%
\pgfpathrectangle{\pgfqpoint{0.514740in}{0.462654in}}{\pgfqpoint{1.909974in}{1.269222in}}%
\pgfusepath{clip}%
\pgfsetbuttcap%
\pgfsetroundjoin%
\pgfsetlinewidth{0.602250pt}%
\definecolor{currentstroke}{rgb}{0.501961,0.501961,0.501961}%
\pgfsetstrokecolor{currentstroke}%
\pgfsetdash{{0.600000pt}{0.990000pt}}{0.000000pt}%
\pgfpathmoveto{\pgfqpoint{1.813522in}{0.462654in}}%
\pgfpathlineto{\pgfqpoint{1.813522in}{1.731877in}}%
\pgfusepath{stroke}%
\end{pgfscope}%
\begin{pgfscope}%
\pgfsetrectcap%
\pgfsetmiterjoin%
\pgfsetlinewidth{0.803000pt}%
\definecolor{currentstroke}{rgb}{0.000000,0.000000,0.000000}%
\pgfsetstrokecolor{currentstroke}%
\pgfsetdash{}{0pt}%
\pgfpathmoveto{\pgfqpoint{0.514740in}{0.462654in}}%
\pgfpathlineto{\pgfqpoint{0.514740in}{1.731877in}}%
\pgfusepath{stroke}%
\end{pgfscope}%
\begin{pgfscope}%
\pgfsetrectcap%
\pgfsetmiterjoin%
\pgfsetlinewidth{0.803000pt}%
\definecolor{currentstroke}{rgb}{0.000000,0.000000,0.000000}%
\pgfsetstrokecolor{currentstroke}%
\pgfsetdash{}{0pt}%
\pgfpathmoveto{\pgfqpoint{2.424714in}{0.462654in}}%
\pgfpathlineto{\pgfqpoint{2.424714in}{1.731877in}}%
\pgfusepath{stroke}%
\end{pgfscope}%
\begin{pgfscope}%
\pgfsetrectcap%
\pgfsetmiterjoin%
\pgfsetlinewidth{0.803000pt}%
\definecolor{currentstroke}{rgb}{0.000000,0.000000,0.000000}%
\pgfsetstrokecolor{currentstroke}%
\pgfsetdash{}{0pt}%
\pgfpathmoveto{\pgfqpoint{0.514740in}{0.462654in}}%
\pgfpathlineto{\pgfqpoint{2.424714in}{0.462654in}}%
\pgfusepath{stroke}%
\end{pgfscope}%
\begin{pgfscope}%
\pgfsetrectcap%
\pgfsetmiterjoin%
\pgfsetlinewidth{0.803000pt}%
\definecolor{currentstroke}{rgb}{0.000000,0.000000,0.000000}%
\pgfsetstrokecolor{currentstroke}%
\pgfsetdash{}{0pt}%
\pgfpathmoveto{\pgfqpoint{0.514740in}{1.731877in}}%
\pgfpathlineto{\pgfqpoint{2.424714in}{1.731877in}}%
\pgfusepath{stroke}%
\end{pgfscope}%
\begin{pgfscope}%
\pgfpathrectangle{\pgfqpoint{0.514740in}{0.462654in}}{\pgfqpoint{1.909974in}{1.269222in}}%
\pgfusepath{clip}%
\pgfsetbuttcap%
\pgfsetroundjoin%
\definecolor{currentfill}{rgb}{0.000000,0.000000,0.000000}%
\pgfsetfillcolor{currentfill}%
\pgfsetlinewidth{1.003750pt}%
\definecolor{currentstroke}{rgb}{0.000000,0.000000,0.000000}%
\pgfsetstrokecolor{currentstroke}%
\pgfsetdash{}{0pt}%
\pgfsys@defobject{currentmarker}{\pgfqpoint{-0.020833in}{-0.020833in}}{\pgfqpoint{0.020833in}{0.020833in}}{%
\pgfpathmoveto{\pgfqpoint{0.000000in}{-0.020833in}}%
\pgfpathcurveto{\pgfqpoint{0.005525in}{-0.020833in}}{\pgfqpoint{0.010825in}{-0.018638in}}{\pgfqpoint{0.014731in}{-0.014731in}}%
\pgfpathcurveto{\pgfqpoint{0.018638in}{-0.010825in}}{\pgfqpoint{0.020833in}{-0.005525in}}{\pgfqpoint{0.020833in}{0.000000in}}%
\pgfpathcurveto{\pgfqpoint{0.020833in}{0.005525in}}{\pgfqpoint{0.018638in}{0.010825in}}{\pgfqpoint{0.014731in}{0.014731in}}%
\pgfpathcurveto{\pgfqpoint{0.010825in}{0.018638in}}{\pgfqpoint{0.005525in}{0.020833in}}{\pgfqpoint{0.000000in}{0.020833in}}%
\pgfpathcurveto{\pgfqpoint{-0.005525in}{0.020833in}}{\pgfqpoint{-0.010825in}{0.018638in}}{\pgfqpoint{-0.014731in}{0.014731in}}%
\pgfpathcurveto{\pgfqpoint{-0.018638in}{0.010825in}}{\pgfqpoint{-0.020833in}{0.005525in}}{\pgfqpoint{-0.020833in}{0.000000in}}%
\pgfpathcurveto{\pgfqpoint{-0.020833in}{-0.005525in}}{\pgfqpoint{-0.018638in}{-0.010825in}}{\pgfqpoint{-0.014731in}{-0.014731in}}%
\pgfpathcurveto{\pgfqpoint{-0.010825in}{-0.018638in}}{\pgfqpoint{-0.005525in}{-0.020833in}}{\pgfqpoint{0.000000in}{-0.020833in}}%
\pgfpathlineto{\pgfqpoint{0.000000in}{-0.020833in}}%
\pgfpathclose%
\pgfusepath{stroke,fill}%
}%
\begin{pgfscope}%
\pgfsys@transformshift{1.813522in}{1.533295in}%
\pgfsys@useobject{currentmarker}{}%
\end{pgfscope}%
\end{pgfscope}%
\begin{pgfscope}%
\definecolor{textcolor}{rgb}{0.000000,0.000000,0.000000}%
\pgfsetstrokecolor{textcolor}%
\pgfsetfillcolor{textcolor}%
\pgftext[x=2.367415in,y=1.668416in,right,top]{\color{textcolor}{\rmfamily\fontsize{9.000000}{10.800000}\selectfont\catcode`\^=\active\def^{\ifmmode\sp\else\^{}\fi}\catcode`\%=\active\def%{\%}$(a)$}}%
\end{pgfscope}%
\begin{pgfscope}%
\pgfsetrectcap%
\pgfsetroundjoin%
\pgfsetlinewidth{0.903375pt}%
\definecolor{currentstroke}{rgb}{0.121569,0.466667,0.705882}%
\pgfsetstrokecolor{currentstroke}%
\pgfsetdash{}{0pt}%
\pgfpathmoveto{\pgfqpoint{0.592518in}{1.602518in}}%
\pgfpathlineto{\pgfqpoint{0.681407in}{1.602518in}}%
\pgfpathlineto{\pgfqpoint{0.770295in}{1.602518in}}%
\pgfusepath{stroke}%
\end{pgfscope}%
\begin{pgfscope}%
\definecolor{textcolor}{rgb}{0.000000,0.000000,0.000000}%
\pgfsetstrokecolor{textcolor}%
\pgfsetfillcolor{textcolor}%
\pgftext[x=0.859184in,y=1.563629in,left,base]{\color{textcolor}{\rmfamily\fontsize{8.000000}{9.600000}\selectfont\catcode`\^=\active\def^{\ifmmode\sp\else\^{}\fi}\catcode`\%=\active\def%{\%}$P_{D_1}$}}%
\end{pgfscope}%
\begin{pgfscope}%
\pgfsetbuttcap%
\pgfsetroundjoin%
\pgfsetlinewidth{0.903375pt}%
\definecolor{currentstroke}{rgb}{1.000000,0.000000,0.000000}%
\pgfsetstrokecolor{currentstroke}%
\pgfsetdash{{3.330000pt}{1.440000pt}}{0.000000pt}%
\pgfpathmoveto{\pgfqpoint{0.592518in}{1.462240in}}%
\pgfpathlineto{\pgfqpoint{0.681407in}{1.462240in}}%
\pgfpathlineto{\pgfqpoint{0.770295in}{1.462240in}}%
\pgfusepath{stroke}%
\end{pgfscope}%
\begin{pgfscope}%
\definecolor{textcolor}{rgb}{0.000000,0.000000,0.000000}%
\pgfsetstrokecolor{textcolor}%
\pgfsetfillcolor{textcolor}%
\pgftext[x=0.859184in,y=1.423351in,left,base]{\color{textcolor}{\rmfamily\fontsize{8.000000}{9.600000}\selectfont\catcode`\^=\active\def^{\ifmmode\sp\else\^{}\fi}\catcode`\%=\active\def%{\%}$P_{D_2}$}}%
\end{pgfscope}%
\begin{pgfscope}%
\pgfsetbuttcap%
\pgfsetmiterjoin%
\definecolor{currentfill}{rgb}{1.000000,1.000000,1.000000}%
\pgfsetfillcolor{currentfill}%
\pgfsetlinewidth{0.000000pt}%
\definecolor{currentstroke}{rgb}{0.000000,0.000000,0.000000}%
\pgfsetstrokecolor{currentstroke}%
\pgfsetstrokeopacity{0.000000}%
\pgfsetdash{}{0pt}%
\pgfpathmoveto{\pgfqpoint{2.622769in}{0.462654in}}%
\pgfpathlineto{\pgfqpoint{4.532743in}{0.462654in}}%
\pgfpathlineto{\pgfqpoint{4.532743in}{1.731877in}}%
\pgfpathlineto{\pgfqpoint{2.622769in}{1.731877in}}%
\pgfpathlineto{\pgfqpoint{2.622769in}{0.462654in}}%
\pgfpathclose%
\pgfusepath{fill}%
\end{pgfscope}%
\begin{pgfscope}%
\pgfsetbuttcap%
\pgfsetroundjoin%
\definecolor{currentfill}{rgb}{0.000000,0.000000,0.000000}%
\pgfsetfillcolor{currentfill}%
\pgfsetlinewidth{0.803000pt}%
\definecolor{currentstroke}{rgb}{0.000000,0.000000,0.000000}%
\pgfsetstrokecolor{currentstroke}%
\pgfsetdash{}{0pt}%
\pgfsys@defobject{currentmarker}{\pgfqpoint{0.000000in}{-0.048611in}}{\pgfqpoint{0.000000in}{0.000000in}}{%
\pgfpathmoveto{\pgfqpoint{0.000000in}{0.000000in}}%
\pgfpathlineto{\pgfqpoint{0.000000in}{-0.048611in}}%
\pgfusepath{stroke,fill}%
}%
\begin{pgfscope}%
\pgfsys@transformshift{2.622769in}{0.462654in}%
\pgfsys@useobject{currentmarker}{}%
\end{pgfscope}%
\end{pgfscope}%
\begin{pgfscope}%
\definecolor{textcolor}{rgb}{0.000000,0.000000,0.000000}%
\pgfsetstrokecolor{textcolor}%
\pgfsetfillcolor{textcolor}%
\pgftext[x=2.622769in,y=0.365432in,,top]{\color{textcolor}{\rmfamily\fontsize{8.000000}{9.600000}\selectfont\catcode`\^=\active\def^{\ifmmode\sp\else\^{}\fi}\catcode`\%=\active\def%{\%}0}}%
\end{pgfscope}%
\begin{pgfscope}%
\pgfsetbuttcap%
\pgfsetroundjoin%
\definecolor{currentfill}{rgb}{0.000000,0.000000,0.000000}%
\pgfsetfillcolor{currentfill}%
\pgfsetlinewidth{0.803000pt}%
\definecolor{currentstroke}{rgb}{0.000000,0.000000,0.000000}%
\pgfsetstrokecolor{currentstroke}%
\pgfsetdash{}{0pt}%
\pgfsys@defobject{currentmarker}{\pgfqpoint{0.000000in}{-0.048611in}}{\pgfqpoint{0.000000in}{0.000000in}}{%
\pgfpathmoveto{\pgfqpoint{0.000000in}{0.000000in}}%
\pgfpathlineto{\pgfqpoint{0.000000in}{-0.048611in}}%
\pgfusepath{stroke,fill}%
}%
\begin{pgfscope}%
\pgfsys@transformshift{3.004764in}{0.462654in}%
\pgfsys@useobject{currentmarker}{}%
\end{pgfscope}%
\end{pgfscope}%
\begin{pgfscope}%
\definecolor{textcolor}{rgb}{0.000000,0.000000,0.000000}%
\pgfsetstrokecolor{textcolor}%
\pgfsetfillcolor{textcolor}%
\pgftext[x=3.004764in,y=0.365432in,,top]{\color{textcolor}{\rmfamily\fontsize{8.000000}{9.600000}\selectfont\catcode`\^=\active\def^{\ifmmode\sp\else\^{}\fi}\catcode`\%=\active\def%{\%}50}}%
\end{pgfscope}%
\begin{pgfscope}%
\pgfsetbuttcap%
\pgfsetroundjoin%
\definecolor{currentfill}{rgb}{0.000000,0.000000,0.000000}%
\pgfsetfillcolor{currentfill}%
\pgfsetlinewidth{0.803000pt}%
\definecolor{currentstroke}{rgb}{0.000000,0.000000,0.000000}%
\pgfsetstrokecolor{currentstroke}%
\pgfsetdash{}{0pt}%
\pgfsys@defobject{currentmarker}{\pgfqpoint{0.000000in}{-0.048611in}}{\pgfqpoint{0.000000in}{0.000000in}}{%
\pgfpathmoveto{\pgfqpoint{0.000000in}{0.000000in}}%
\pgfpathlineto{\pgfqpoint{0.000000in}{-0.048611in}}%
\pgfusepath{stroke,fill}%
}%
\begin{pgfscope}%
\pgfsys@transformshift{3.386759in}{0.462654in}%
\pgfsys@useobject{currentmarker}{}%
\end{pgfscope}%
\end{pgfscope}%
\begin{pgfscope}%
\definecolor{textcolor}{rgb}{0.000000,0.000000,0.000000}%
\pgfsetstrokecolor{textcolor}%
\pgfsetfillcolor{textcolor}%
\pgftext[x=3.386759in,y=0.365432in,,top]{\color{textcolor}{\rmfamily\fontsize{8.000000}{9.600000}\selectfont\catcode`\^=\active\def^{\ifmmode\sp\else\^{}\fi}\catcode`\%=\active\def%{\%}100}}%
\end{pgfscope}%
\begin{pgfscope}%
\pgfsetbuttcap%
\pgfsetroundjoin%
\definecolor{currentfill}{rgb}{0.000000,0.000000,0.000000}%
\pgfsetfillcolor{currentfill}%
\pgfsetlinewidth{0.803000pt}%
\definecolor{currentstroke}{rgb}{0.000000,0.000000,0.000000}%
\pgfsetstrokecolor{currentstroke}%
\pgfsetdash{}{0pt}%
\pgfsys@defobject{currentmarker}{\pgfqpoint{0.000000in}{-0.048611in}}{\pgfqpoint{0.000000in}{0.000000in}}{%
\pgfpathmoveto{\pgfqpoint{0.000000in}{0.000000in}}%
\pgfpathlineto{\pgfqpoint{0.000000in}{-0.048611in}}%
\pgfusepath{stroke,fill}%
}%
\begin{pgfscope}%
\pgfsys@transformshift{3.768754in}{0.462654in}%
\pgfsys@useobject{currentmarker}{}%
\end{pgfscope}%
\end{pgfscope}%
\begin{pgfscope}%
\definecolor{textcolor}{rgb}{0.000000,0.000000,0.000000}%
\pgfsetstrokecolor{textcolor}%
\pgfsetfillcolor{textcolor}%
\pgftext[x=3.768754in,y=0.365432in,,top]{\color{textcolor}{\rmfamily\fontsize{8.000000}{9.600000}\selectfont\catcode`\^=\active\def^{\ifmmode\sp\else\^{}\fi}\catcode`\%=\active\def%{\%}150}}%
\end{pgfscope}%
\begin{pgfscope}%
\pgfsetbuttcap%
\pgfsetroundjoin%
\definecolor{currentfill}{rgb}{0.000000,0.000000,0.000000}%
\pgfsetfillcolor{currentfill}%
\pgfsetlinewidth{0.803000pt}%
\definecolor{currentstroke}{rgb}{0.000000,0.000000,0.000000}%
\pgfsetstrokecolor{currentstroke}%
\pgfsetdash{}{0pt}%
\pgfsys@defobject{currentmarker}{\pgfqpoint{0.000000in}{-0.048611in}}{\pgfqpoint{0.000000in}{0.000000in}}{%
\pgfpathmoveto{\pgfqpoint{0.000000in}{0.000000in}}%
\pgfpathlineto{\pgfqpoint{0.000000in}{-0.048611in}}%
\pgfusepath{stroke,fill}%
}%
\begin{pgfscope}%
\pgfsys@transformshift{4.150749in}{0.462654in}%
\pgfsys@useobject{currentmarker}{}%
\end{pgfscope}%
\end{pgfscope}%
\begin{pgfscope}%
\definecolor{textcolor}{rgb}{0.000000,0.000000,0.000000}%
\pgfsetstrokecolor{textcolor}%
\pgfsetfillcolor{textcolor}%
\pgftext[x=4.150749in,y=0.365432in,,top]{\color{textcolor}{\rmfamily\fontsize{8.000000}{9.600000}\selectfont\catcode`\^=\active\def^{\ifmmode\sp\else\^{}\fi}\catcode`\%=\active\def%{\%}200}}%
\end{pgfscope}%
\begin{pgfscope}%
\pgfsetbuttcap%
\pgfsetroundjoin%
\definecolor{currentfill}{rgb}{0.000000,0.000000,0.000000}%
\pgfsetfillcolor{currentfill}%
\pgfsetlinewidth{0.803000pt}%
\definecolor{currentstroke}{rgb}{0.000000,0.000000,0.000000}%
\pgfsetstrokecolor{currentstroke}%
\pgfsetdash{}{0pt}%
\pgfsys@defobject{currentmarker}{\pgfqpoint{0.000000in}{-0.048611in}}{\pgfqpoint{0.000000in}{0.000000in}}{%
\pgfpathmoveto{\pgfqpoint{0.000000in}{0.000000in}}%
\pgfpathlineto{\pgfqpoint{0.000000in}{-0.048611in}}%
\pgfusepath{stroke,fill}%
}%
\begin{pgfscope}%
\pgfsys@transformshift{4.532743in}{0.462654in}%
\pgfsys@useobject{currentmarker}{}%
\end{pgfscope}%
\end{pgfscope}%
\begin{pgfscope}%
\definecolor{textcolor}{rgb}{0.000000,0.000000,0.000000}%
\pgfsetstrokecolor{textcolor}%
\pgfsetfillcolor{textcolor}%
\pgftext[x=4.532743in,y=0.365432in,,top]{\color{textcolor}{\rmfamily\fontsize{8.000000}{9.600000}\selectfont\catcode`\^=\active\def^{\ifmmode\sp\else\^{}\fi}\catcode`\%=\active\def%{\%}250}}%
\end{pgfscope}%
\begin{pgfscope}%
\definecolor{textcolor}{rgb}{0.000000,0.000000,0.000000}%
\pgfsetstrokecolor{textcolor}%
\pgfsetfillcolor{textcolor}%
\pgftext[x=3.577756in,y=0.211111in,,top]{\color{textcolor}{\rmfamily\fontsize{9.000000}{10.800000}\selectfont\catcode`\^=\active\def^{\ifmmode\sp\else\^{}\fi}\catcode`\%=\active\def%{\%}$\omega_b\,t$}}%
\end{pgfscope}%
\begin{pgfscope}%
\pgfpathrectangle{\pgfqpoint{2.622769in}{0.462654in}}{\pgfqpoint{1.909974in}{1.269222in}}%
\pgfusepath{clip}%
\pgfsetrectcap%
\pgfsetroundjoin%
\pgfsetlinewidth{0.903375pt}%
\definecolor{currentstroke}{rgb}{0.121569,0.466667,0.705882}%
\pgfsetstrokecolor{currentstroke}%
\pgfsetdash{}{0pt}%
\pgfpathmoveto{\pgfqpoint{2.622769in}{0.462654in}}%
\pgfpathlineto{\pgfqpoint{2.638049in}{0.464817in}}%
\pgfpathlineto{\pgfqpoint{2.653329in}{0.466835in}}%
\pgfpathlineto{\pgfqpoint{2.668609in}{0.473650in}}%
\pgfpathlineto{\pgfqpoint{2.683889in}{0.475118in}}%
\pgfpathlineto{\pgfqpoint{2.699168in}{0.481468in}}%
\pgfpathlineto{\pgfqpoint{2.714448in}{0.482661in}}%
\pgfpathlineto{\pgfqpoint{2.729728in}{0.492711in}}%
\pgfpathlineto{\pgfqpoint{2.745008in}{0.497022in}}%
\pgfpathlineto{\pgfqpoint{2.760287in}{0.514123in}}%
\pgfpathlineto{\pgfqpoint{2.775567in}{0.519037in}}%
\pgfpathlineto{\pgfqpoint{2.790847in}{0.536992in}}%
\pgfpathlineto{\pgfqpoint{2.806127in}{0.543079in}}%
\pgfpathlineto{\pgfqpoint{2.821407in}{0.561791in}}%
\pgfpathlineto{\pgfqpoint{2.836686in}{0.568516in}}%
\pgfpathlineto{\pgfqpoint{2.851966in}{0.590217in}}%
\pgfpathlineto{\pgfqpoint{2.867246in}{0.601653in}}%
\pgfpathlineto{\pgfqpoint{2.882526in}{0.624201in}}%
\pgfpathlineto{\pgfqpoint{2.897806in}{0.639505in}}%
\pgfpathlineto{\pgfqpoint{2.913085in}{0.657330in}}%
\pgfpathlineto{\pgfqpoint{2.928365in}{0.673266in}}%
\pgfpathlineto{\pgfqpoint{2.943645in}{0.691147in}}%
\pgfpathlineto{\pgfqpoint{2.958925in}{0.716397in}}%
\pgfpathlineto{\pgfqpoint{2.974205in}{0.733764in}}%
\pgfpathlineto{\pgfqpoint{2.989484in}{0.764780in}}%
\pgfpathlineto{\pgfqpoint{3.004764in}{0.776620in}}%
\pgfpathlineto{\pgfqpoint{3.020044in}{0.809510in}}%
\pgfpathlineto{\pgfqpoint{3.035324in}{0.820099in}}%
\pgfpathlineto{\pgfqpoint{3.050604in}{0.856077in}}%
\pgfpathlineto{\pgfqpoint{3.065883in}{0.861972in}}%
\pgfpathlineto{\pgfqpoint{3.081163in}{0.902162in}}%
\pgfpathlineto{\pgfqpoint{3.096443in}{0.910645in}}%
\pgfpathlineto{\pgfqpoint{3.111723in}{0.951641in}}%
\pgfpathlineto{\pgfqpoint{3.127002in}{0.959443in}}%
\pgfpathlineto{\pgfqpoint{3.142282in}{0.996820in}}%
\pgfpathlineto{\pgfqpoint{3.157562in}{1.005592in}}%
\pgfpathlineto{\pgfqpoint{3.172842in}{1.045552in}}%
\pgfpathlineto{\pgfqpoint{3.188122in}{1.063044in}}%
\pgfpathlineto{\pgfqpoint{3.203401in}{1.094925in}}%
\pgfpathlineto{\pgfqpoint{3.218681in}{1.112692in}}%
\pgfpathlineto{\pgfqpoint{3.233961in}{1.135822in}}%
\pgfpathlineto{\pgfqpoint{3.249241in}{1.158330in}}%
\pgfpathlineto{\pgfqpoint{3.264521in}{1.177227in}}%
\pgfpathlineto{\pgfqpoint{3.279800in}{1.205437in}}%
\pgfpathlineto{\pgfqpoint{3.295080in}{1.215840in}}%
\pgfpathlineto{\pgfqpoint{3.310360in}{1.253409in}}%
\pgfpathlineto{\pgfqpoint{3.325640in}{1.262686in}}%
\pgfpathlineto{\pgfqpoint{3.340920in}{1.301269in}}%
\pgfpathlineto{\pgfqpoint{3.356199in}{1.300190in}}%
\pgfpathlineto{\pgfqpoint{3.371479in}{1.339635in}}%
\pgfpathlineto{\pgfqpoint{3.386759in}{1.335789in}}%
\pgfpathlineto{\pgfqpoint{3.402039in}{1.379805in}}%
\pgfpathlineto{\pgfqpoint{3.417319in}{1.375370in}}%
\pgfpathlineto{\pgfqpoint{3.432598in}{1.410332in}}%
\pgfpathlineto{\pgfqpoint{3.447878in}{1.402962in}}%
\pgfpathlineto{\pgfqpoint{3.463158in}{1.435965in}}%
\pgfpathlineto{\pgfqpoint{3.478438in}{1.434040in}}%
\pgfpathlineto{\pgfqpoint{3.493717in}{1.462512in}}%
\pgfpathlineto{\pgfqpoint{3.508997in}{1.464479in}}%
\pgfpathlineto{\pgfqpoint{3.524277in}{1.483926in}}%
\pgfpathlineto{\pgfqpoint{3.539557in}{1.493864in}}%
\pgfpathlineto{\pgfqpoint{3.554837in}{1.506051in}}%
\pgfpathlineto{\pgfqpoint{3.570116in}{1.514492in}}%
\pgfpathlineto{\pgfqpoint{3.585396in}{1.511626in}}%
\pgfpathlineto{\pgfqpoint{3.600676in}{1.527553in}}%
\pgfpathlineto{\pgfqpoint{3.615956in}{1.519777in}}%
\pgfpathlineto{\pgfqpoint{3.631236in}{1.542976in}}%
\pgfpathlineto{\pgfqpoint{3.646515in}{1.526435in}}%
\pgfpathlineto{\pgfqpoint{3.661795in}{1.546708in}}%
\pgfpathlineto{\pgfqpoint{3.677075in}{1.525255in}}%
\pgfpathlineto{\pgfqpoint{3.692355in}{1.551879in}}%
\pgfpathlineto{\pgfqpoint{3.707635in}{1.529382in}}%
\pgfpathlineto{\pgfqpoint{3.722914in}{1.551911in}}%
\pgfpathlineto{\pgfqpoint{3.738194in}{1.525977in}}%
\pgfpathlineto{\pgfqpoint{3.753474in}{1.543061in}}%
\pgfpathlineto{\pgfqpoint{3.768754in}{1.519218in}}%
\pgfpathlineto{\pgfqpoint{3.784034in}{1.529308in}}%
\pgfpathlineto{\pgfqpoint{3.799313in}{1.502039in}}%
\pgfpathlineto{\pgfqpoint{3.814593in}{1.503963in}}%
\pgfpathlineto{\pgfqpoint{3.829873in}{1.489143in}}%
\pgfpathlineto{\pgfqpoint{3.845153in}{1.487364in}}%
\pgfpathlineto{\pgfqpoint{3.860432in}{1.476928in}}%
\pgfpathlineto{\pgfqpoint{3.875712in}{1.462357in}}%
\pgfpathlineto{\pgfqpoint{3.890992in}{1.453504in}}%
\pgfpathlineto{\pgfqpoint{3.906272in}{1.432564in}}%
\pgfpathlineto{\pgfqpoint{3.921552in}{1.431955in}}%
\pgfpathlineto{\pgfqpoint{3.936831in}{1.402078in}}%
\pgfpathlineto{\pgfqpoint{3.952111in}{1.400589in}}%
\pgfpathlineto{\pgfqpoint{3.967391in}{1.364963in}}%
\pgfpathlineto{\pgfqpoint{3.982671in}{1.366207in}}%
\pgfpathlineto{\pgfqpoint{3.997951in}{1.328120in}}%
\pgfpathlineto{\pgfqpoint{4.013230in}{1.328975in}}%
\pgfpathlineto{\pgfqpoint{4.028510in}{1.287374in}}%
\pgfpathlineto{\pgfqpoint{4.043790in}{1.289460in}}%
\pgfpathlineto{\pgfqpoint{4.059070in}{1.255425in}}%
\pgfpathlineto{\pgfqpoint{4.074350in}{1.253351in}}%
\pgfpathlineto{\pgfqpoint{4.089629in}{1.215935in}}%
\pgfpathlineto{\pgfqpoint{4.104909in}{1.204538in}}%
\pgfpathlineto{\pgfqpoint{4.120189in}{1.170614in}}%
\pgfpathlineto{\pgfqpoint{4.135469in}{1.156265in}}%
\pgfpathlineto{\pgfqpoint{4.150749in}{1.129812in}}%
\pgfpathlineto{\pgfqpoint{4.166028in}{1.108068in}}%
\pgfpathlineto{\pgfqpoint{4.181308in}{1.084553in}}%
\pgfpathlineto{\pgfqpoint{4.196588in}{1.059544in}}%
\pgfpathlineto{\pgfqpoint{4.211868in}{1.042350in}}%
\pgfpathlineto{\pgfqpoint{4.227148in}{1.011996in}}%
\pgfpathlineto{\pgfqpoint{4.242427in}{0.998195in}}%
\pgfpathlineto{\pgfqpoint{4.257707in}{0.964112in}}%
\pgfpathlineto{\pgfqpoint{4.272987in}{0.955333in}}%
\pgfpathlineto{\pgfqpoint{4.288267in}{0.921621in}}%
\pgfpathlineto{\pgfqpoint{4.303546in}{0.910658in}}%
\pgfpathlineto{\pgfqpoint{4.318826in}{0.871110in}}%
\pgfpathlineto{\pgfqpoint{4.334106in}{0.860609in}}%
\pgfpathlineto{\pgfqpoint{4.349386in}{0.826207in}}%
\pgfpathlineto{\pgfqpoint{4.364666in}{0.817775in}}%
\pgfpathlineto{\pgfqpoint{4.379945in}{0.787475in}}%
\pgfpathlineto{\pgfqpoint{4.395225in}{0.776429in}}%
\pgfpathlineto{\pgfqpoint{4.410505in}{0.747811in}}%
\pgfpathlineto{\pgfqpoint{4.425785in}{0.735783in}}%
\pgfpathlineto{\pgfqpoint{4.441065in}{0.711001in}}%
\pgfpathlineto{\pgfqpoint{4.456344in}{0.695000in}}%
\pgfpathlineto{\pgfqpoint{4.471624in}{0.675046in}}%
\pgfpathlineto{\pgfqpoint{4.486904in}{0.658639in}}%
\pgfpathlineto{\pgfqpoint{4.502184in}{0.643252in}}%
\pgfpathlineto{\pgfqpoint{4.517464in}{0.625608in}}%
\pgfpathlineto{\pgfqpoint{4.532743in}{0.611485in}}%
\pgfusepath{stroke}%
\end{pgfscope}%
\begin{pgfscope}%
\pgfpathrectangle{\pgfqpoint{2.622769in}{0.462654in}}{\pgfqpoint{1.909974in}{1.269222in}}%
\pgfusepath{clip}%
\pgfsetbuttcap%
\pgfsetroundjoin%
\pgfsetlinewidth{0.903375pt}%
\definecolor{currentstroke}{rgb}{1.000000,0.000000,0.000000}%
\pgfsetstrokecolor{currentstroke}%
\pgfsetdash{{3.330000pt}{1.440000pt}}{0.000000pt}%
\pgfpathmoveto{\pgfqpoint{2.622769in}{0.462654in}}%
\pgfpathlineto{\pgfqpoint{2.638049in}{0.464410in}}%
\pgfpathlineto{\pgfqpoint{2.653329in}{0.466011in}}%
\pgfpathlineto{\pgfqpoint{2.668609in}{0.468657in}}%
\pgfpathlineto{\pgfqpoint{2.683889in}{0.470287in}}%
\pgfpathlineto{\pgfqpoint{2.699168in}{0.470407in}}%
\pgfpathlineto{\pgfqpoint{2.714448in}{0.472842in}}%
\pgfpathlineto{\pgfqpoint{2.729728in}{0.471618in}}%
\pgfpathlineto{\pgfqpoint{2.745008in}{0.471691in}}%
\pgfpathlineto{\pgfqpoint{2.760287in}{0.466559in}}%
\pgfpathlineto{\pgfqpoint{2.775567in}{0.465886in}}%
\pgfpathlineto{\pgfqpoint{2.790847in}{0.464285in}}%
\pgfpathlineto{\pgfqpoint{2.806127in}{0.464673in}}%
\pgfpathlineto{\pgfqpoint{2.821407in}{0.464757in}}%
\pgfpathlineto{\pgfqpoint{2.836686in}{0.464895in}}%
\pgfpathlineto{\pgfqpoint{2.851966in}{0.468424in}}%
\pgfpathlineto{\pgfqpoint{2.867246in}{0.469266in}}%
\pgfpathlineto{\pgfqpoint{2.882526in}{0.473013in}}%
\pgfpathlineto{\pgfqpoint{2.897806in}{0.471763in}}%
\pgfpathlineto{\pgfqpoint{2.913085in}{0.471763in}}%
\pgfpathlineto{\pgfqpoint{2.928365in}{0.469123in}}%
\pgfpathlineto{\pgfqpoint{2.943645in}{0.468564in}}%
\pgfpathlineto{\pgfqpoint{2.958925in}{0.467128in}}%
\pgfpathlineto{\pgfqpoint{2.974205in}{0.464745in}}%
\pgfpathlineto{\pgfqpoint{2.989484in}{0.463697in}}%
\pgfpathlineto{\pgfqpoint{3.004764in}{0.463757in}}%
\pgfpathlineto{\pgfqpoint{3.020044in}{0.466903in}}%
\pgfpathlineto{\pgfqpoint{3.035324in}{0.468525in}}%
\pgfpathlineto{\pgfqpoint{3.050604in}{0.469547in}}%
\pgfpathlineto{\pgfqpoint{3.065883in}{0.470979in}}%
\pgfpathlineto{\pgfqpoint{3.081163in}{0.471327in}}%
\pgfpathlineto{\pgfqpoint{3.096443in}{0.472761in}}%
\pgfpathlineto{\pgfqpoint{3.111723in}{0.469906in}}%
\pgfpathlineto{\pgfqpoint{3.127002in}{0.469196in}}%
\pgfpathlineto{\pgfqpoint{3.142282in}{0.465080in}}%
\pgfpathlineto{\pgfqpoint{3.157562in}{0.465510in}}%
\pgfpathlineto{\pgfqpoint{3.172842in}{0.464942in}}%
\pgfpathlineto{\pgfqpoint{3.188122in}{0.465484in}}%
\pgfpathlineto{\pgfqpoint{3.203401in}{0.466395in}}%
\pgfpathlineto{\pgfqpoint{3.218681in}{0.466338in}}%
\pgfpathlineto{\pgfqpoint{3.233961in}{0.470455in}}%
\pgfpathlineto{\pgfqpoint{3.249241in}{0.470705in}}%
\pgfpathlineto{\pgfqpoint{3.264521in}{0.472734in}}%
\pgfpathlineto{\pgfqpoint{3.279800in}{0.470080in}}%
\pgfpathlineto{\pgfqpoint{3.295080in}{0.470408in}}%
\pgfpathlineto{\pgfqpoint{3.310360in}{0.468878in}}%
\pgfpathlineto{\pgfqpoint{3.325640in}{0.467724in}}%
\pgfpathlineto{\pgfqpoint{3.340920in}{0.466339in}}%
\pgfpathlineto{\pgfqpoint{3.356199in}{0.464314in}}%
\pgfpathlineto{\pgfqpoint{3.371479in}{0.464848in}}%
\pgfpathlineto{\pgfqpoint{3.386759in}{0.465942in}}%
\pgfpathlineto{\pgfqpoint{3.402039in}{0.468365in}}%
\pgfpathlineto{\pgfqpoint{3.417319in}{0.469441in}}%
\pgfpathlineto{\pgfqpoint{3.432598in}{0.469957in}}%
\pgfpathlineto{\pgfqpoint{3.447878in}{0.471191in}}%
\pgfpathlineto{\pgfqpoint{3.463158in}{0.471098in}}%
\pgfpathlineto{\pgfqpoint{3.478438in}{0.471969in}}%
\pgfpathlineto{\pgfqpoint{3.493717in}{0.468358in}}%
\pgfpathlineto{\pgfqpoint{3.508997in}{0.467612in}}%
\pgfpathlineto{\pgfqpoint{3.524277in}{0.465469in}}%
\pgfpathlineto{\pgfqpoint{3.539557in}{0.466621in}}%
\pgfpathlineto{\pgfqpoint{3.554837in}{0.465845in}}%
\pgfpathlineto{\pgfqpoint{3.570116in}{0.465777in}}%
\pgfpathlineto{\pgfqpoint{3.585396in}{0.467042in}}%
\pgfpathlineto{\pgfqpoint{3.600676in}{0.467512in}}%
\pgfpathlineto{\pgfqpoint{3.615956in}{0.471667in}}%
\pgfpathlineto{\pgfqpoint{3.631236in}{0.470542in}}%
\pgfpathlineto{\pgfqpoint{3.646515in}{0.471835in}}%
\pgfpathlineto{\pgfqpoint{3.661795in}{0.469566in}}%
\pgfpathlineto{\pgfqpoint{3.677075in}{0.470361in}}%
\pgfpathlineto{\pgfqpoint{3.692355in}{0.468843in}}%
\pgfpathlineto{\pgfqpoint{3.707635in}{0.467049in}}%
\pgfpathlineto{\pgfqpoint{3.722914in}{0.465591in}}%
\pgfpathlineto{\pgfqpoint{3.738194in}{0.464389in}}%
\pgfpathlineto{\pgfqpoint{3.753474in}{0.466352in}}%
\pgfpathlineto{\pgfqpoint{3.768754in}{0.467249in}}%
\pgfpathlineto{\pgfqpoint{3.784034in}{0.468367in}}%
\pgfpathlineto{\pgfqpoint{3.799313in}{0.469444in}}%
\pgfpathlineto{\pgfqpoint{3.814593in}{0.470088in}}%
\pgfpathlineto{\pgfqpoint{3.829873in}{0.471977in}}%
\pgfpathlineto{\pgfqpoint{3.845153in}{0.471327in}}%
\pgfpathlineto{\pgfqpoint{3.860432in}{0.470946in}}%
\pgfpathlineto{\pgfqpoint{3.875712in}{0.467595in}}%
\pgfpathlineto{\pgfqpoint{3.890992in}{0.467518in}}%
\pgfpathlineto{\pgfqpoint{3.906272in}{0.466216in}}%
\pgfpathlineto{\pgfqpoint{3.921552in}{0.466062in}}%
\pgfpathlineto{\pgfqpoint{3.936831in}{0.465415in}}%
\pgfpathlineto{\pgfqpoint{3.952111in}{0.465309in}}%
\pgfpathlineto{\pgfqpoint{3.967391in}{0.468045in}}%
\pgfpathlineto{\pgfqpoint{3.982671in}{0.469097in}}%
\pgfpathlineto{\pgfqpoint{3.997951in}{0.472085in}}%
\pgfpathlineto{\pgfqpoint{4.013230in}{0.470280in}}%
\pgfpathlineto{\pgfqpoint{4.028510in}{0.471951in}}%
\pgfpathlineto{\pgfqpoint{4.043790in}{0.469931in}}%
\pgfpathlineto{\pgfqpoint{4.059070in}{0.470203in}}%
\pgfpathlineto{\pgfqpoint{4.074350in}{0.468204in}}%
\pgfpathlineto{\pgfqpoint{4.089629in}{0.465565in}}%
\pgfpathlineto{\pgfqpoint{4.104909in}{0.464845in}}%
\pgfpathlineto{\pgfqpoint{4.120189in}{0.464784in}}%
\pgfpathlineto{\pgfqpoint{4.135469in}{0.467034in}}%
\pgfpathlineto{\pgfqpoint{4.150749in}{0.467084in}}%
\pgfpathlineto{\pgfqpoint{4.166028in}{0.468619in}}%
\pgfpathlineto{\pgfqpoint{4.181308in}{0.470340in}}%
\pgfpathlineto{\pgfqpoint{4.196588in}{0.471649in}}%
\pgfpathlineto{\pgfqpoint{4.211868in}{0.473087in}}%
\pgfpathlineto{\pgfqpoint{4.227148in}{0.470508in}}%
\pgfpathlineto{\pgfqpoint{4.242427in}{0.469708in}}%
\pgfpathlineto{\pgfqpoint{4.257707in}{0.467106in}}%
\pgfpathlineto{\pgfqpoint{4.272987in}{0.466748in}}%
\pgfpathlineto{\pgfqpoint{4.288267in}{0.465268in}}%
\pgfpathlineto{\pgfqpoint{4.303546in}{0.464933in}}%
\pgfpathlineto{\pgfqpoint{4.318826in}{0.465287in}}%
\pgfpathlineto{\pgfqpoint{4.334106in}{0.466102in}}%
\pgfpathlineto{\pgfqpoint{4.349386in}{0.470104in}}%
\pgfpathlineto{\pgfqpoint{4.364666in}{0.470451in}}%
\pgfpathlineto{\pgfqpoint{4.379945in}{0.472347in}}%
\pgfpathlineto{\pgfqpoint{4.395225in}{0.470663in}}%
\pgfpathlineto{\pgfqpoint{4.410505in}{0.472044in}}%
\pgfpathlineto{\pgfqpoint{4.425785in}{0.469986in}}%
\pgfpathlineto{\pgfqpoint{4.441065in}{0.468584in}}%
\pgfpathlineto{\pgfqpoint{4.456344in}{0.465765in}}%
\pgfpathlineto{\pgfqpoint{4.471624in}{0.464189in}}%
\pgfpathlineto{\pgfqpoint{4.486904in}{0.465127in}}%
\pgfpathlineto{\pgfqpoint{4.502184in}{0.465224in}}%
\pgfpathlineto{\pgfqpoint{4.517464in}{0.467378in}}%
\pgfpathlineto{\pgfqpoint{4.532743in}{0.468320in}}%
\pgfusepath{stroke}%
\end{pgfscope}%
\begin{pgfscope}%
\pgfpathrectangle{\pgfqpoint{2.622769in}{0.462654in}}{\pgfqpoint{1.909974in}{1.269222in}}%
\pgfusepath{clip}%
\pgfsetbuttcap%
\pgfsetroundjoin%
\pgfsetlinewidth{0.602250pt}%
\definecolor{currentstroke}{rgb}{0.501961,0.501961,0.501961}%
\pgfsetstrokecolor{currentstroke}%
\pgfsetdash{{0.600000pt}{0.990000pt}}{0.000000pt}%
\pgfpathmoveto{\pgfqpoint{3.768754in}{0.462654in}}%
\pgfpathlineto{\pgfqpoint{3.768754in}{1.731877in}}%
\pgfusepath{stroke}%
\end{pgfscope}%
\begin{pgfscope}%
\pgfsetrectcap%
\pgfsetmiterjoin%
\pgfsetlinewidth{0.803000pt}%
\definecolor{currentstroke}{rgb}{0.000000,0.000000,0.000000}%
\pgfsetstrokecolor{currentstroke}%
\pgfsetdash{}{0pt}%
\pgfpathmoveto{\pgfqpoint{2.622769in}{0.462654in}}%
\pgfpathlineto{\pgfqpoint{2.622769in}{1.731877in}}%
\pgfusepath{stroke}%
\end{pgfscope}%
\begin{pgfscope}%
\pgfsetrectcap%
\pgfsetmiterjoin%
\pgfsetlinewidth{0.803000pt}%
\definecolor{currentstroke}{rgb}{0.000000,0.000000,0.000000}%
\pgfsetstrokecolor{currentstroke}%
\pgfsetdash{}{0pt}%
\pgfpathmoveto{\pgfqpoint{4.532743in}{0.462654in}}%
\pgfpathlineto{\pgfqpoint{4.532743in}{1.731877in}}%
\pgfusepath{stroke}%
\end{pgfscope}%
\begin{pgfscope}%
\pgfsetrectcap%
\pgfsetmiterjoin%
\pgfsetlinewidth{0.803000pt}%
\definecolor{currentstroke}{rgb}{0.000000,0.000000,0.000000}%
\pgfsetstrokecolor{currentstroke}%
\pgfsetdash{}{0pt}%
\pgfpathmoveto{\pgfqpoint{2.622769in}{0.462654in}}%
\pgfpathlineto{\pgfqpoint{4.532743in}{0.462654in}}%
\pgfusepath{stroke}%
\end{pgfscope}%
\begin{pgfscope}%
\pgfsetrectcap%
\pgfsetmiterjoin%
\pgfsetlinewidth{0.803000pt}%
\definecolor{currentstroke}{rgb}{0.000000,0.000000,0.000000}%
\pgfsetstrokecolor{currentstroke}%
\pgfsetdash{}{0pt}%
\pgfpathmoveto{\pgfqpoint{2.622769in}{1.731877in}}%
\pgfpathlineto{\pgfqpoint{4.532743in}{1.731877in}}%
\pgfusepath{stroke}%
\end{pgfscope}%
\begin{pgfscope}%
\pgfpathrectangle{\pgfqpoint{2.622769in}{0.462654in}}{\pgfqpoint{1.909974in}{1.269222in}}%
\pgfusepath{clip}%
\pgfsetbuttcap%
\pgfsetroundjoin%
\definecolor{currentfill}{rgb}{0.000000,0.000000,0.000000}%
\pgfsetfillcolor{currentfill}%
\pgfsetlinewidth{1.003750pt}%
\definecolor{currentstroke}{rgb}{0.000000,0.000000,0.000000}%
\pgfsetstrokecolor{currentstroke}%
\pgfsetdash{}{0pt}%
\pgfsys@defobject{currentmarker}{\pgfqpoint{-0.020833in}{-0.020833in}}{\pgfqpoint{0.020833in}{0.020833in}}{%
\pgfpathmoveto{\pgfqpoint{0.000000in}{-0.020833in}}%
\pgfpathcurveto{\pgfqpoint{0.005525in}{-0.020833in}}{\pgfqpoint{0.010825in}{-0.018638in}}{\pgfqpoint{0.014731in}{-0.014731in}}%
\pgfpathcurveto{\pgfqpoint{0.018638in}{-0.010825in}}{\pgfqpoint{0.020833in}{-0.005525in}}{\pgfqpoint{0.020833in}{0.000000in}}%
\pgfpathcurveto{\pgfqpoint{0.020833in}{0.005525in}}{\pgfqpoint{0.018638in}{0.010825in}}{\pgfqpoint{0.014731in}{0.014731in}}%
\pgfpathcurveto{\pgfqpoint{0.010825in}{0.018638in}}{\pgfqpoint{0.005525in}{0.020833in}}{\pgfqpoint{0.000000in}{0.020833in}}%
\pgfpathcurveto{\pgfqpoint{-0.005525in}{0.020833in}}{\pgfqpoint{-0.010825in}{0.018638in}}{\pgfqpoint{-0.014731in}{0.014731in}}%
\pgfpathcurveto{\pgfqpoint{-0.018638in}{0.010825in}}{\pgfqpoint{-0.020833in}{0.005525in}}{\pgfqpoint{-0.020833in}{0.000000in}}%
\pgfpathcurveto{\pgfqpoint{-0.020833in}{-0.005525in}}{\pgfqpoint{-0.018638in}{-0.010825in}}{\pgfqpoint{-0.014731in}{-0.014731in}}%
\pgfpathcurveto{\pgfqpoint{-0.010825in}{-0.018638in}}{\pgfqpoint{-0.005525in}{-0.020833in}}{\pgfqpoint{0.000000in}{-0.020833in}}%
\pgfpathlineto{\pgfqpoint{0.000000in}{-0.020833in}}%
\pgfpathclose%
\pgfusepath{stroke,fill}%
}%
\begin{pgfscope}%
\pgfsys@transformshift{3.768754in}{1.519218in}%
\pgfsys@useobject{currentmarker}{}%
\end{pgfscope}%
\end{pgfscope}%
\begin{pgfscope}%
\definecolor{textcolor}{rgb}{0.000000,0.000000,0.000000}%
\pgfsetstrokecolor{textcolor}%
\pgfsetfillcolor{textcolor}%
\pgftext[x=4.475444in,y=1.668416in,right,top]{\color{textcolor}{\rmfamily\fontsize{9.000000}{10.800000}\selectfont\catcode`\^=\active\def^{\ifmmode\sp\else\^{}\fi}\catcode`\%=\active\def%{\%}$(b)$}}%
\end{pgfscope}%
\begin{pgfscope}%
\pgfsetrectcap%
\pgfsetroundjoin%
\pgfsetlinewidth{0.903375pt}%
\definecolor{currentstroke}{rgb}{0.121569,0.466667,0.705882}%
\pgfsetstrokecolor{currentstroke}%
\pgfsetdash{}{0pt}%
\pgfpathmoveto{\pgfqpoint{2.700547in}{1.602518in}}%
\pgfpathlineto{\pgfqpoint{2.789436in}{1.602518in}}%
\pgfpathlineto{\pgfqpoint{2.878325in}{1.602518in}}%
\pgfusepath{stroke}%
\end{pgfscope}%
\begin{pgfscope}%
\definecolor{textcolor}{rgb}{0.000000,0.000000,0.000000}%
\pgfsetstrokecolor{textcolor}%
\pgfsetfillcolor{textcolor}%
\pgftext[x=2.967214in,y=1.563629in,left,base]{\color{textcolor}{\rmfamily\fontsize{8.000000}{9.600000}\selectfont\catcode`\^=\active\def^{\ifmmode\sp\else\^{}\fi}\catcode`\%=\active\def%{\%}$P_{D_2}$}}%
\end{pgfscope}%
\begin{pgfscope}%
\pgfsetbuttcap%
\pgfsetroundjoin%
\pgfsetlinewidth{0.903375pt}%
\definecolor{currentstroke}{rgb}{1.000000,0.000000,0.000000}%
\pgfsetstrokecolor{currentstroke}%
\pgfsetdash{{3.330000pt}{1.440000pt}}{0.000000pt}%
\pgfpathmoveto{\pgfqpoint{2.700547in}{1.462240in}}%
\pgfpathlineto{\pgfqpoint{2.789436in}{1.462240in}}%
\pgfpathlineto{\pgfqpoint{2.878325in}{1.462240in}}%
\pgfusepath{stroke}%
\end{pgfscope}%
\begin{pgfscope}%
\definecolor{textcolor}{rgb}{0.000000,0.000000,0.000000}%
\pgfsetstrokecolor{textcolor}%
\pgfsetfillcolor{textcolor}%
\pgftext[x=2.967214in,y=1.423351in,left,base]{\color{textcolor}{\rmfamily\fontsize{8.000000}{9.600000}\selectfont\catcode`\^=\active\def^{\ifmmode\sp\else\^{}\fi}\catcode`\%=\active\def%{\%}$P_{D_1}$}}%
\end{pgfscope}%
\end{pgfpicture}%
\makeatother%
\endgroup%
}
\caption{(En color en l\'inea) Poblaci\'on de la ranura oscura objetivo (azul) y de la otra ranura (rojo, discontinuo) durante un pulso de escritura desde $\ket{G,1}$ a la desintonizaci\'on optimizada, para (a) la ranura $D_1$ y (b) la ranura $D_2$. Las l\'ineas de puntos y los puntos marcan las duraciones optimizadas de la Fig.~\ref{fig:raman_map}. Par\'ametros del sistema como en la Fig.~\ref{fig:raman_map}.}
\label{fig:raman_dyn}
\end{figure}

\section{P\'erdidas, ruido y puertos}
\label{sec:noise}

Tres mecanismos act\'uan sobre el qubit almacenado: la p\'erdida de la excitaci\'on ac\'ustica, el desfase de las energ\'ias de los excitones y la relajaci\'on por el continuo de fonones. Afectan de manera distinta a las dos ranuras oscuras, y solo el desfase depende de c\'omo est\'e correlacionado el ruido entre los puntos.

La din\'amica abierta sigue la ecuaci\'on maestra de Lindblad
\begin{equation}
\begin{aligned}\dot\rho={}&-i[H,\rho]+\mathcal D[L_\kappa]\rho+\mathcal D[L_\gamma]\rho\\&+\sum_\alpha\mathcal D[L_{\varphi,\alpha}]\rho+\sum_{a\neq b}\mathcal D[L_{a\to b}]\rho,\end{aligned}\label{eq:master}
\end{equation}
con $\mathcal D[L]\rho=L\rho L^\dagger-\tfrac12\{L^\dagger L,\rho\}$ y $H$ el hamiltoniano de la Ec.~\eqref{eq:H} incluido el campo de excitaci\'on. Los operadores de salto son la p\'erdida ac\'ustica $L_\kappa=\sqrt\kappa\,b$, el canal radiativo colectivo $L_\gamma$ de los excitones, los operadores de desfase $L_{\varphi,\alpha}$ y la relajaci\'on $L_{a\to b}=\sqrt{\Gamma_{a\to b}}\ket b\bra a$ entre autoestados excit\'onicos por el continuo de fonones, que definimos m\'as abajo. El canal radiativo es
\begin{align}
L_\gamma&=\sqrt{\frac{2\gamma}{3}}\sum_i\sigma_i,\nonumber\\
\|L_\gamma\ket B\|^2&=2\gamma,\qquad L_\gamma\ket{D_s}=0.\label{eq:radiative}
\end{align}
La \'ultima igualdad requiere dipolos \'opticos paralelos, fases equivalentes y un canal radiativo colectivo; las p\'erdidas radiativas independientes romper\'ian la protecci\'on~\cite{dicke1954,mokhlespour2012,carlson2019}.

Modelamos las fluctuaciones de las energ\'ias de sitio como ruido gaussiano blanco, $\langle\xi_i(t)\xi_j(t')\rangle=\gamma_\varphi K_{ij}\delta(t-t')$, con $K$ semidefinida positiva y $K_{ii}=1$. Si $K=\sum_\alpha\mu_\alpha\bm u_\alpha\bm u_\alpha^T$, los operadores de Lindblad son
\begin{equation}
L_{\varphi,\alpha}=\sqrt{\gamma_\varphi\mu_\alpha}\sum_iu_{i\alpha}\,n_i.\label{eq:noisejumps}
\end{equation}
Tomamos $K=(1-\chi)I+\chi\bm1\bm1^T$, con $\chi=0$ para ruido local y $\chi=1$ para ruido perfectamente com\'un. Este \'ultimo es el l\'imite de ruido colectivo en el que se basan las codificaciones libres de decoherencia~\cite{lidar1998,friesen2017}, no una correlaci\'on medida para la estructura de tres puntos propuesta. En la base de autoestados $\{B,D_1,D_2\}$ y bajo la aproximaci\'on secular $|E_a-E_b|\gg\gamma_\varphi$, la tasa de saltos inducidos por ruido entre autoestados es $W_{ab}=\gamma_\varphi\sum_{ij}K_{ij}(a_ib_i)(a_jb_j)$, y para las dos ranuras
\begin{equation}
W_{D_1D_2}=\frac{\gamma_\varphi}{6}(1-\chi),\qquad
\Gamma_{D_1D_2}^{(\varphi)}=\frac{5\gamma_\varphi}{6}(1-\chi),\label{eq:gamma_dark}
\end{equation}
donde $\Gamma_{D_1D_2}^{(\varphi)}$ es la tasa de decaimiento de la coherencia entre $D_1$ y $D_2$. Para $\chi=1$ el operador de ruido es proporcional al n\'umero total de excitones y act\'ua como la identidad sobre el subespacio de almacenamiento $\mathcal H_L$, de modo que la coherencia entre las ranuras queda protegida. La coherencia entre el vac\'io y un estado de un ex\'citon decae a $\gamma_\varphi/2$ para cualquier $\chi$.

La Figura~\ref{fig:storage} muestra la retenci\'on de cada ranura y de la coherencia durante la espera. En (a) ambas ranuras conservan casi toda su poblaci\'on durante esperas en las que un fon\'on pasivo ha deca\'ido como $e^{-\kappa t}$ hasta desaparecer; $D_1$ pierde algo m\'as que $D_2$, sobre todo por una lenta transferencia a la otra ranura. En (b), para $\chi=1$ la coherencia decae lentamente, para $\chi=0$ sigue la ley secular de la Ec.~\eqref{eq:gamma_dark}, y una correlaci\'on de $\chi=0.9$ ya sit\'ua la coherencia claramente por debajo de la curva ideal a esperas largas. La protecci\'on requiere por tanto una correlaci\'on cercana a la unidad, no una moderada.

La protecci\'on no cubre la escritura y la lectura secuenciales: entre dos ventanas una rama contiene un ex\'citon y la otra ninguno. Para tiempos sin protecci\'on $\tau_u^w$ y $\tau_u^r$ y una espera protegida $T_{\rm prot}$, la visibilidad de la superposici\'on almacenada es
\begin{equation}
\mathcal V\simeq\mathcal V_0\exp\!\Bigl[-\frac{\gamma_\varphi}{2}(\tau_u^w+\tau_u^r)-\Gamma_{D_1D_2}^{(\varphi)}T_{\rm prot}\Bigr].\label{eq:visibility}
\end{equation}
Variando solo la separaci\'on entre las dos ventanas en una simulaci\'on con carga ideal y $\chi=1$, la pendiente de $\ln\mathcal V$ concuerda aproximadamente con $-\gamma_\varphi/2$. El intervalo sin protecci\'on es por tanto el origen del costo de cargar lentamente el segundo fon\'on.

\begin{figure}[t]
	\centering
	\resizebox{\columnwidth}{!}{%% Creator: Matplotlib, PGF backend
%%
%% To include the figure in your LaTeX document, write
%%   \input{<filename>.pgf}
%%
%% Make sure the required packages are loaded in your preamble
%%   \usepackage{pgf}
%%
%% Also ensure that all the required font packages are loaded; for instance,
%% the lmodern package is sometimes necessary when using math font.
%%   \usepackage{lmodern}
%%
%% Figures using additional raster images can only be included by \input if
%% they are in the same directory as the main LaTeX file. For loading figures
%% from other directories you can use the `import` package
%%   \usepackage{import}
%%
%% and then include the figures with
%%   \import{<path to file>}{<filename>.pgf}
%%
%% Matplotlib used the following preamble
%%   \def\mathdefault#1{#1}
%%   \everymath=\expandafter{\the\everymath\displaystyle}
%%   \IfFileExists{scrextend.sty}{
%%     \usepackage[fontsize=8.000000pt]{scrextend}
%%   }{
%%     \renewcommand{\normalsize}{\fontsize{8.000000}{9.600000}\selectfont}
%%     \normalsize
%%   }
%%   \usepackage{amsmath,amssymb}
%%   \makeatletter\@ifpackageloaded{underscore}{}{\usepackage[strings]{underscore}}\makeatother
%%
\begingroup%
\makeatletter%
\begin{pgfpicture}%
\pgfpathrectangle{\pgfpointorigin}{\pgfqpoint{4.733508in}{1.836515in}}%
\pgfusepath{use as bounding box, clip}%
\begin{pgfscope}%
\pgfsetbuttcap%
\pgfsetmiterjoin%
\definecolor{currentfill}{rgb}{1.000000,1.000000,1.000000}%
\pgfsetfillcolor{currentfill}%
\pgfsetlinewidth{0.000000pt}%
\definecolor{currentstroke}{rgb}{1.000000,1.000000,1.000000}%
\pgfsetstrokecolor{currentstroke}%
\pgfsetdash{}{0pt}%
\pgfpathmoveto{\pgfqpoint{0.000000in}{0.000000in}}%
\pgfpathlineto{\pgfqpoint{4.733508in}{0.000000in}}%
\pgfpathlineto{\pgfqpoint{4.733508in}{1.836515in}}%
\pgfpathlineto{\pgfqpoint{0.000000in}{1.836515in}}%
\pgfpathlineto{\pgfqpoint{0.000000in}{0.000000in}}%
\pgfpathclose%
\pgfusepath{fill}%
\end{pgfscope}%
\begin{pgfscope}%
\pgfsetbuttcap%
\pgfsetmiterjoin%
\definecolor{currentfill}{rgb}{1.000000,1.000000,1.000000}%
\pgfsetfillcolor{currentfill}%
\pgfsetlinewidth{0.000000pt}%
\definecolor{currentstroke}{rgb}{0.000000,0.000000,0.000000}%
\pgfsetstrokecolor{currentstroke}%
\pgfsetstrokeopacity{0.000000}%
\pgfsetdash{}{0pt}%
\pgfpathmoveto{\pgfqpoint{0.514740in}{0.496542in}}%
\pgfpathlineto{\pgfqpoint{2.264249in}{0.496542in}}%
\pgfpathlineto{\pgfqpoint{2.264249in}{1.736515in}}%
\pgfpathlineto{\pgfqpoint{0.514740in}{1.736515in}}%
\pgfpathlineto{\pgfqpoint{0.514740in}{0.496542in}}%
\pgfpathclose%
\pgfusepath{fill}%
\end{pgfscope}%
\begin{pgfscope}%
\pgfsetbuttcap%
\pgfsetroundjoin%
\definecolor{currentfill}{rgb}{0.000000,0.000000,0.000000}%
\pgfsetfillcolor{currentfill}%
\pgfsetlinewidth{0.803000pt}%
\definecolor{currentstroke}{rgb}{0.000000,0.000000,0.000000}%
\pgfsetstrokecolor{currentstroke}%
\pgfsetdash{}{0pt}%
\pgfsys@defobject{currentmarker}{\pgfqpoint{0.000000in}{-0.048611in}}{\pgfqpoint{0.000000in}{0.000000in}}{%
\pgfpathmoveto{\pgfqpoint{0.000000in}{0.000000in}}%
\pgfpathlineto{\pgfqpoint{0.000000in}{-0.048611in}}%
\pgfusepath{stroke,fill}%
}%
\begin{pgfscope}%
\pgfsys@transformshift{0.514740in}{0.496542in}%
\pgfsys@useobject{currentmarker}{}%
\end{pgfscope}%
\end{pgfscope}%
\begin{pgfscope}%
\definecolor{textcolor}{rgb}{0.000000,0.000000,0.000000}%
\pgfsetstrokecolor{textcolor}%
\pgfsetfillcolor{textcolor}%
\pgftext[x=0.514740in,y=0.399320in,,top]{\color{textcolor}{\rmfamily\fontsize{8.000000}{9.600000}\selectfont\catcode`\^=\active\def^{\ifmmode\sp\else\^{}\fi}\catcode`\%=\active\def%{\%}0}}%
\end{pgfscope}%
\begin{pgfscope}%
\pgfsetbuttcap%
\pgfsetroundjoin%
\definecolor{currentfill}{rgb}{0.000000,0.000000,0.000000}%
\pgfsetfillcolor{currentfill}%
\pgfsetlinewidth{0.803000pt}%
\definecolor{currentstroke}{rgb}{0.000000,0.000000,0.000000}%
\pgfsetstrokecolor{currentstroke}%
\pgfsetdash{}{0pt}%
\pgfsys@defobject{currentmarker}{\pgfqpoint{0.000000in}{-0.048611in}}{\pgfqpoint{0.000000in}{0.000000in}}{%
\pgfpathmoveto{\pgfqpoint{0.000000in}{0.000000in}}%
\pgfpathlineto{\pgfqpoint{0.000000in}{-0.048611in}}%
\pgfusepath{stroke,fill}%
}%
\begin{pgfscope}%
\pgfsys@transformshift{0.952117in}{0.496542in}%
\pgfsys@useobject{currentmarker}{}%
\end{pgfscope}%
\end{pgfscope}%
\begin{pgfscope}%
\definecolor{textcolor}{rgb}{0.000000,0.000000,0.000000}%
\pgfsetstrokecolor{textcolor}%
\pgfsetfillcolor{textcolor}%
\pgftext[x=0.952117in,y=0.399320in,,top]{\color{textcolor}{\rmfamily\fontsize{8.000000}{9.600000}\selectfont\catcode`\^=\active\def^{\ifmmode\sp\else\^{}\fi}\catcode`\%=\active\def%{\%}2}}%
\end{pgfscope}%
\begin{pgfscope}%
\pgfsetbuttcap%
\pgfsetroundjoin%
\definecolor{currentfill}{rgb}{0.000000,0.000000,0.000000}%
\pgfsetfillcolor{currentfill}%
\pgfsetlinewidth{0.803000pt}%
\definecolor{currentstroke}{rgb}{0.000000,0.000000,0.000000}%
\pgfsetstrokecolor{currentstroke}%
\pgfsetdash{}{0pt}%
\pgfsys@defobject{currentmarker}{\pgfqpoint{0.000000in}{-0.048611in}}{\pgfqpoint{0.000000in}{0.000000in}}{%
\pgfpathmoveto{\pgfqpoint{0.000000in}{0.000000in}}%
\pgfpathlineto{\pgfqpoint{0.000000in}{-0.048611in}}%
\pgfusepath{stroke,fill}%
}%
\begin{pgfscope}%
\pgfsys@transformshift{1.389494in}{0.496542in}%
\pgfsys@useobject{currentmarker}{}%
\end{pgfscope}%
\end{pgfscope}%
\begin{pgfscope}%
\definecolor{textcolor}{rgb}{0.000000,0.000000,0.000000}%
\pgfsetstrokecolor{textcolor}%
\pgfsetfillcolor{textcolor}%
\pgftext[x=1.389494in,y=0.399320in,,top]{\color{textcolor}{\rmfamily\fontsize{8.000000}{9.600000}\selectfont\catcode`\^=\active\def^{\ifmmode\sp\else\^{}\fi}\catcode`\%=\active\def%{\%}4}}%
\end{pgfscope}%
\begin{pgfscope}%
\pgfsetbuttcap%
\pgfsetroundjoin%
\definecolor{currentfill}{rgb}{0.000000,0.000000,0.000000}%
\pgfsetfillcolor{currentfill}%
\pgfsetlinewidth{0.803000pt}%
\definecolor{currentstroke}{rgb}{0.000000,0.000000,0.000000}%
\pgfsetstrokecolor{currentstroke}%
\pgfsetdash{}{0pt}%
\pgfsys@defobject{currentmarker}{\pgfqpoint{0.000000in}{-0.048611in}}{\pgfqpoint{0.000000in}{0.000000in}}{%
\pgfpathmoveto{\pgfqpoint{0.000000in}{0.000000in}}%
\pgfpathlineto{\pgfqpoint{0.000000in}{-0.048611in}}%
\pgfusepath{stroke,fill}%
}%
\begin{pgfscope}%
\pgfsys@transformshift{1.826872in}{0.496542in}%
\pgfsys@useobject{currentmarker}{}%
\end{pgfscope}%
\end{pgfscope}%
\begin{pgfscope}%
\definecolor{textcolor}{rgb}{0.000000,0.000000,0.000000}%
\pgfsetstrokecolor{textcolor}%
\pgfsetfillcolor{textcolor}%
\pgftext[x=1.826872in,y=0.399320in,,top]{\color{textcolor}{\rmfamily\fontsize{8.000000}{9.600000}\selectfont\catcode`\^=\active\def^{\ifmmode\sp\else\^{}\fi}\catcode`\%=\active\def%{\%}6}}%
\end{pgfscope}%
\begin{pgfscope}%
\pgfsetbuttcap%
\pgfsetroundjoin%
\definecolor{currentfill}{rgb}{0.000000,0.000000,0.000000}%
\pgfsetfillcolor{currentfill}%
\pgfsetlinewidth{0.803000pt}%
\definecolor{currentstroke}{rgb}{0.000000,0.000000,0.000000}%
\pgfsetstrokecolor{currentstroke}%
\pgfsetdash{}{0pt}%
\pgfsys@defobject{currentmarker}{\pgfqpoint{0.000000in}{-0.048611in}}{\pgfqpoint{0.000000in}{0.000000in}}{%
\pgfpathmoveto{\pgfqpoint{0.000000in}{0.000000in}}%
\pgfpathlineto{\pgfqpoint{0.000000in}{-0.048611in}}%
\pgfusepath{stroke,fill}%
}%
\begin{pgfscope}%
\pgfsys@transformshift{2.264249in}{0.496542in}%
\pgfsys@useobject{currentmarker}{}%
\end{pgfscope}%
\end{pgfscope}%
\begin{pgfscope}%
\definecolor{textcolor}{rgb}{0.000000,0.000000,0.000000}%
\pgfsetstrokecolor{textcolor}%
\pgfsetfillcolor{textcolor}%
\pgftext[x=2.264249in,y=0.399320in,,top]{\color{textcolor}{\rmfamily\fontsize{8.000000}{9.600000}\selectfont\catcode`\^=\active\def^{\ifmmode\sp\else\^{}\fi}\catcode`\%=\active\def%{\%}8}}%
\end{pgfscope}%
\begin{pgfscope}%
\definecolor{textcolor}{rgb}{0.000000,0.000000,0.000000}%
\pgfsetstrokecolor{textcolor}%
\pgfsetfillcolor{textcolor}%
\pgftext[x=1.389494in,y=0.244999in,,top]{\color{textcolor}{\rmfamily\fontsize{9.000000}{10.800000}\selectfont\catcode`\^=\active\def^{\ifmmode\sp\else\^{}\fi}\catcode`\%=\active\def%{\%}$T\ (10^3/\omega_b)$}}%
\end{pgfscope}%
\begin{pgfscope}%
\pgfsetbuttcap%
\pgfsetroundjoin%
\definecolor{currentfill}{rgb}{0.000000,0.000000,0.000000}%
\pgfsetfillcolor{currentfill}%
\pgfsetlinewidth{0.803000pt}%
\definecolor{currentstroke}{rgb}{0.000000,0.000000,0.000000}%
\pgfsetstrokecolor{currentstroke}%
\pgfsetdash{}{0pt}%
\pgfsys@defobject{currentmarker}{\pgfqpoint{-0.048611in}{0.000000in}}{\pgfqpoint{-0.000000in}{0.000000in}}{%
\pgfpathmoveto{\pgfqpoint{-0.000000in}{0.000000in}}%
\pgfpathlineto{\pgfqpoint{-0.048611in}{0.000000in}}%
\pgfusepath{stroke,fill}%
}%
\begin{pgfscope}%
\pgfsys@transformshift{0.514740in}{0.496542in}%
\pgfsys@useobject{currentmarker}{}%
\end{pgfscope}%
\end{pgfscope}%
\begin{pgfscope}%
\definecolor{textcolor}{rgb}{0.000000,0.000000,0.000000}%
\pgfsetstrokecolor{textcolor}%
\pgfsetfillcolor{textcolor}%
\pgftext[x=0.358489in, y=0.457962in, left, base]{\color{textcolor}{\rmfamily\fontsize{8.000000}{9.600000}\selectfont\catcode`\^=\active\def^{\ifmmode\sp\else\^{}\fi}\catcode`\%=\active\def%{\%}$0$}}%
\end{pgfscope}%
\begin{pgfscope}%
\pgfsetbuttcap%
\pgfsetroundjoin%
\definecolor{currentfill}{rgb}{0.000000,0.000000,0.000000}%
\pgfsetfillcolor{currentfill}%
\pgfsetlinewidth{0.803000pt}%
\definecolor{currentstroke}{rgb}{0.000000,0.000000,0.000000}%
\pgfsetstrokecolor{currentstroke}%
\pgfsetdash{}{0pt}%
\pgfsys@defobject{currentmarker}{\pgfqpoint{-0.048611in}{0.000000in}}{\pgfqpoint{-0.000000in}{0.000000in}}{%
\pgfpathmoveto{\pgfqpoint{-0.000000in}{0.000000in}}%
\pgfpathlineto{\pgfqpoint{-0.048611in}{0.000000in}}%
\pgfusepath{stroke,fill}%
}%
\begin{pgfscope}%
\pgfsys@transformshift{0.514740in}{0.909867in}%
\pgfsys@useobject{currentmarker}{}%
\end{pgfscope}%
\end{pgfscope}%
\begin{pgfscope}%
\definecolor{textcolor}{rgb}{0.000000,0.000000,0.000000}%
\pgfsetstrokecolor{textcolor}%
\pgfsetfillcolor{textcolor}%
\pgftext[x=0.266667in, y=0.871286in, left, base]{\color{textcolor}{\rmfamily\fontsize{8.000000}{9.600000}\selectfont\catcode`\^=\active\def^{\ifmmode\sp\else\^{}\fi}\catcode`\%=\active\def%{\%}$0.5$}}%
\end{pgfscope}%
\begin{pgfscope}%
\pgfsetbuttcap%
\pgfsetroundjoin%
\definecolor{currentfill}{rgb}{0.000000,0.000000,0.000000}%
\pgfsetfillcolor{currentfill}%
\pgfsetlinewidth{0.803000pt}%
\definecolor{currentstroke}{rgb}{0.000000,0.000000,0.000000}%
\pgfsetstrokecolor{currentstroke}%
\pgfsetdash{}{0pt}%
\pgfsys@defobject{currentmarker}{\pgfqpoint{-0.048611in}{0.000000in}}{\pgfqpoint{-0.000000in}{0.000000in}}{%
\pgfpathmoveto{\pgfqpoint{-0.000000in}{0.000000in}}%
\pgfpathlineto{\pgfqpoint{-0.048611in}{0.000000in}}%
\pgfusepath{stroke,fill}%
}%
\begin{pgfscope}%
\pgfsys@transformshift{0.514740in}{1.323191in}%
\pgfsys@useobject{currentmarker}{}%
\end{pgfscope}%
\end{pgfscope}%
\begin{pgfscope}%
\definecolor{textcolor}{rgb}{0.000000,0.000000,0.000000}%
\pgfsetstrokecolor{textcolor}%
\pgfsetfillcolor{textcolor}%
\pgftext[x=0.358489in, y=1.284610in, left, base]{\color{textcolor}{\rmfamily\fontsize{8.000000}{9.600000}\selectfont\catcode`\^=\active\def^{\ifmmode\sp\else\^{}\fi}\catcode`\%=\active\def%{\%}$1$}}%
\end{pgfscope}%
\begin{pgfscope}%
\definecolor{textcolor}{rgb}{0.000000,0.000000,0.000000}%
\pgfsetstrokecolor{textcolor}%
\pgfsetfillcolor{textcolor}%
\pgftext[x=0.211111in,y=1.116529in,,bottom,rotate=90.000000]{\color{textcolor}{\rmfamily\fontsize{9.000000}{10.800000}\selectfont\catcode`\^=\active\def^{\ifmmode\sp\else\^{}\fi}\catcode`\%=\active\def%{\%}Retained population}}%
\end{pgfscope}%
\begin{pgfscope}%
\pgfpathrectangle{\pgfqpoint{0.514740in}{0.496542in}}{\pgfqpoint{1.749509in}{1.239972in}}%
\pgfusepath{clip}%
\pgfsetrectcap%
\pgfsetroundjoin%
\pgfsetlinewidth{0.903375pt}%
\definecolor{currentstroke}{rgb}{0.121569,0.466667,0.705882}%
\pgfsetstrokecolor{currentstroke}%
\pgfsetdash{}{0pt}%
\pgfpathmoveto{\pgfqpoint{0.514740in}{1.323191in}}%
\pgfpathlineto{\pgfqpoint{0.515396in}{1.310235in}}%
\pgfpathlineto{\pgfqpoint{0.516271in}{1.312120in}}%
\pgfpathlineto{\pgfqpoint{0.516708in}{1.319728in}}%
\pgfpathlineto{\pgfqpoint{0.517145in}{1.309293in}}%
\pgfpathlineto{\pgfqpoint{0.517364in}{1.304979in}}%
\pgfpathlineto{\pgfqpoint{0.517802in}{1.312283in}}%
\pgfpathlineto{\pgfqpoint{0.518676in}{1.318651in}}%
\pgfpathlineto{\pgfqpoint{0.518895in}{1.317558in}}%
\pgfpathlineto{\pgfqpoint{0.519332in}{1.307616in}}%
\pgfpathlineto{\pgfqpoint{0.519770in}{1.315607in}}%
\pgfpathlineto{\pgfqpoint{0.519988in}{1.319275in}}%
\pgfpathlineto{\pgfqpoint{0.520644in}{1.315348in}}%
\pgfpathlineto{\pgfqpoint{0.520863in}{1.315076in}}%
\pgfpathlineto{\pgfqpoint{0.521301in}{1.306592in}}%
\pgfpathlineto{\pgfqpoint{0.521738in}{1.313279in}}%
\pgfpathlineto{\pgfqpoint{0.521957in}{1.318928in}}%
\pgfpathlineto{\pgfqpoint{0.522613in}{1.310651in}}%
\pgfpathlineto{\pgfqpoint{0.523050in}{1.310871in}}%
\pgfpathlineto{\pgfqpoint{0.523269in}{1.309273in}}%
\pgfpathlineto{\pgfqpoint{0.523487in}{1.309868in}}%
\pgfpathlineto{\pgfqpoint{0.524143in}{1.320458in}}%
\pgfpathlineto{\pgfqpoint{0.524581in}{1.309481in}}%
\pgfpathlineto{\pgfqpoint{0.524800in}{1.309871in}}%
\pgfpathlineto{\pgfqpoint{0.526112in}{1.319767in}}%
\pgfpathlineto{\pgfqpoint{0.526768in}{1.306314in}}%
\pgfpathlineto{\pgfqpoint{0.527424in}{1.314903in}}%
\pgfpathlineto{\pgfqpoint{0.527642in}{1.314417in}}%
\pgfpathlineto{\pgfqpoint{0.527861in}{1.315494in}}%
\pgfpathlineto{\pgfqpoint{0.528080in}{1.316178in}}%
\pgfpathlineto{\pgfqpoint{0.528736in}{1.305319in}}%
\pgfpathlineto{\pgfqpoint{0.528955in}{1.310714in}}%
\pgfpathlineto{\pgfqpoint{0.529392in}{1.319174in}}%
\pgfpathlineto{\pgfqpoint{0.530048in}{1.314646in}}%
\pgfpathlineto{\pgfqpoint{0.530485in}{1.310503in}}%
\pgfpathlineto{\pgfqpoint{0.530704in}{1.308233in}}%
\pgfpathlineto{\pgfqpoint{0.530923in}{1.311738in}}%
\pgfpathlineto{\pgfqpoint{0.531360in}{1.320868in}}%
\pgfpathlineto{\pgfqpoint{0.531798in}{1.311977in}}%
\pgfpathlineto{\pgfqpoint{0.532672in}{1.307656in}}%
\pgfpathlineto{\pgfqpoint{0.532891in}{1.310036in}}%
\pgfpathlineto{\pgfqpoint{0.533328in}{1.319657in}}%
\pgfpathlineto{\pgfqpoint{0.533766in}{1.310108in}}%
\pgfpathlineto{\pgfqpoint{0.533984in}{1.306832in}}%
\pgfpathlineto{\pgfqpoint{0.534422in}{1.312883in}}%
\pgfpathlineto{\pgfqpoint{0.535297in}{1.320455in}}%
\pgfpathlineto{\pgfqpoint{0.535515in}{1.318510in}}%
\pgfpathlineto{\pgfqpoint{0.535953in}{1.307219in}}%
\pgfpathlineto{\pgfqpoint{0.536609in}{1.315273in}}%
\pgfpathlineto{\pgfqpoint{0.538796in}{1.317419in}}%
\pgfpathlineto{\pgfqpoint{0.571351in}{1.307868in}}%
\pgfpathlineto{\pgfqpoint{0.603907in}{1.314067in}}%
\pgfpathlineto{\pgfqpoint{0.614759in}{1.312136in}}%
\pgfpathlineto{\pgfqpoint{0.625611in}{1.311580in}}%
\pgfpathlineto{\pgfqpoint{0.636463in}{1.306072in}}%
\pgfpathlineto{\pgfqpoint{0.647315in}{1.308146in}}%
\pgfpathlineto{\pgfqpoint{0.658167in}{1.307545in}}%
\pgfpathlineto{\pgfqpoint{0.669018in}{1.311692in}}%
\pgfpathlineto{\pgfqpoint{0.679870in}{1.310255in}}%
\pgfpathlineto{\pgfqpoint{0.690722in}{1.310152in}}%
\pgfpathlineto{\pgfqpoint{0.701574in}{1.305012in}}%
\pgfpathlineto{\pgfqpoint{0.712426in}{1.305668in}}%
\pgfpathlineto{\pgfqpoint{0.723278in}{1.303536in}}%
\pgfpathlineto{\pgfqpoint{0.734130in}{1.308702in}}%
\pgfpathlineto{\pgfqpoint{0.744982in}{1.308104in}}%
\pgfpathlineto{\pgfqpoint{0.755834in}{1.309316in}}%
\pgfpathlineto{\pgfqpoint{0.766686in}{1.305076in}}%
\pgfpathlineto{\pgfqpoint{0.777538in}{1.303749in}}%
\pgfpathlineto{\pgfqpoint{0.788389in}{1.300682in}}%
\pgfpathlineto{\pgfqpoint{0.799241in}{1.304319in}}%
\pgfpathlineto{\pgfqpoint{0.810093in}{1.304726in}}%
\pgfpathlineto{\pgfqpoint{0.820945in}{1.308081in}}%
\pgfpathlineto{\pgfqpoint{0.853501in}{1.299756in}}%
\pgfpathlineto{\pgfqpoint{0.875205in}{1.300135in}}%
\pgfpathlineto{\pgfqpoint{0.886057in}{1.303823in}}%
\pgfpathlineto{\pgfqpoint{0.907760in}{1.303520in}}%
\pgfpathlineto{\pgfqpoint{0.918612in}{1.300442in}}%
\pgfpathlineto{\pgfqpoint{0.940316in}{1.297781in}}%
\pgfpathlineto{\pgfqpoint{0.972872in}{1.301144in}}%
\pgfpathlineto{\pgfqpoint{1.005428in}{1.296807in}}%
\pgfpathlineto{\pgfqpoint{1.016280in}{1.296270in}}%
\pgfpathlineto{\pgfqpoint{1.037983in}{1.297500in}}%
\pgfpathlineto{\pgfqpoint{1.070539in}{1.295702in}}%
\pgfpathlineto{\pgfqpoint{1.103095in}{1.295135in}}%
\pgfpathlineto{\pgfqpoint{1.113947in}{1.295082in}}%
\pgfpathlineto{\pgfqpoint{1.135651in}{1.292966in}}%
\pgfpathlineto{\pgfqpoint{1.157354in}{1.293228in}}%
\pgfpathlineto{\pgfqpoint{1.179058in}{1.293910in}}%
\pgfpathlineto{\pgfqpoint{1.211614in}{1.289936in}}%
\pgfpathlineto{\pgfqpoint{1.233318in}{1.290783in}}%
\pgfpathlineto{\pgfqpoint{1.244170in}{1.291832in}}%
\pgfpathlineto{\pgfqpoint{1.265873in}{1.290730in}}%
\pgfpathlineto{\pgfqpoint{1.276725in}{1.288499in}}%
\pgfpathlineto{\pgfqpoint{1.298429in}{1.287210in}}%
\pgfpathlineto{\pgfqpoint{1.330985in}{1.289462in}}%
\pgfpathlineto{\pgfqpoint{1.374392in}{1.285032in}}%
\pgfpathlineto{\pgfqpoint{1.417800in}{1.285965in}}%
\pgfpathlineto{\pgfqpoint{1.450356in}{1.283032in}}%
\pgfpathlineto{\pgfqpoint{1.515467in}{1.282035in}}%
\pgfpathlineto{\pgfqpoint{1.569727in}{1.280082in}}%
\pgfpathlineto{\pgfqpoint{1.602283in}{1.279831in}}%
\pgfpathlineto{\pgfqpoint{1.656542in}{1.277368in}}%
\pgfpathlineto{\pgfqpoint{1.678246in}{1.277933in}}%
\pgfpathlineto{\pgfqpoint{1.732505in}{1.274657in}}%
\pgfpathlineto{\pgfqpoint{1.765061in}{1.275324in}}%
\pgfpathlineto{\pgfqpoint{1.808469in}{1.271862in}}%
\pgfpathlineto{\pgfqpoint{1.841024in}{1.272880in}}%
\pgfpathlineto{\pgfqpoint{1.916988in}{1.270090in}}%
\pgfpathlineto{\pgfqpoint{1.938692in}{1.269467in}}%
\pgfpathlineto{\pgfqpoint{1.982099in}{1.266965in}}%
\pgfpathlineto{\pgfqpoint{2.025507in}{1.266889in}}%
\pgfpathlineto{\pgfqpoint{2.068915in}{1.264211in}}%
\pgfpathlineto{\pgfqpoint{2.112322in}{1.264338in}}%
\pgfpathlineto{\pgfqpoint{2.155730in}{1.261567in}}%
\pgfpathlineto{\pgfqpoint{2.188286in}{1.262239in}}%
\pgfpathlineto{\pgfqpoint{2.253397in}{1.259502in}}%
\pgfpathlineto{\pgfqpoint{2.264249in}{1.259813in}}%
\pgfpathlineto{\pgfqpoint{2.264249in}{1.259813in}}%
\pgfusepath{stroke}%
\end{pgfscope}%
\begin{pgfscope}%
\pgfpathrectangle{\pgfqpoint{0.514740in}{0.496542in}}{\pgfqpoint{1.749509in}{1.239972in}}%
\pgfusepath{clip}%
\pgfsetbuttcap%
\pgfsetroundjoin%
\pgfsetlinewidth{0.903375pt}%
\definecolor{currentstroke}{rgb}{1.000000,0.000000,0.000000}%
\pgfsetstrokecolor{currentstroke}%
\pgfsetdash{{5.760000pt}{1.440000pt}{0.900000pt}{1.440000pt}}{0.000000pt}%
\pgfpathmoveto{\pgfqpoint{0.514740in}{1.323191in}}%
\pgfpathlineto{\pgfqpoint{0.516052in}{1.311384in}}%
\pgfpathlineto{\pgfqpoint{0.516271in}{1.312275in}}%
\pgfpathlineto{\pgfqpoint{0.517145in}{1.316246in}}%
\pgfpathlineto{\pgfqpoint{0.517364in}{1.315975in}}%
\pgfpathlineto{\pgfqpoint{0.517583in}{1.313123in}}%
\pgfpathlineto{\pgfqpoint{0.518020in}{1.318599in}}%
\pgfpathlineto{\pgfqpoint{0.518458in}{1.314333in}}%
\pgfpathlineto{\pgfqpoint{0.518676in}{1.314459in}}%
\pgfpathlineto{\pgfqpoint{0.519114in}{1.320025in}}%
\pgfpathlineto{\pgfqpoint{0.519551in}{1.313620in}}%
\pgfpathlineto{\pgfqpoint{0.519770in}{1.315812in}}%
\pgfpathlineto{\pgfqpoint{0.519988in}{1.317117in}}%
\pgfpathlineto{\pgfqpoint{0.520207in}{1.313914in}}%
\pgfpathlineto{\pgfqpoint{0.520426in}{1.312835in}}%
\pgfpathlineto{\pgfqpoint{0.520644in}{1.315930in}}%
\pgfpathlineto{\pgfqpoint{0.521519in}{1.318808in}}%
\pgfpathlineto{\pgfqpoint{0.521738in}{1.317511in}}%
\pgfpathlineto{\pgfqpoint{0.522831in}{1.308487in}}%
\pgfpathlineto{\pgfqpoint{0.523050in}{1.310271in}}%
\pgfpathlineto{\pgfqpoint{0.524362in}{1.319317in}}%
\pgfpathlineto{\pgfqpoint{0.525456in}{1.315739in}}%
\pgfpathlineto{\pgfqpoint{0.525674in}{1.316925in}}%
\pgfpathlineto{\pgfqpoint{0.525893in}{1.318246in}}%
\pgfpathlineto{\pgfqpoint{0.526112in}{1.315505in}}%
\pgfpathlineto{\pgfqpoint{0.526549in}{1.315647in}}%
\pgfpathlineto{\pgfqpoint{0.527205in}{1.310062in}}%
\pgfpathlineto{\pgfqpoint{0.528517in}{1.319505in}}%
\pgfpathlineto{\pgfqpoint{0.528736in}{1.318193in}}%
\pgfpathlineto{\pgfqpoint{0.529611in}{1.313499in}}%
\pgfpathlineto{\pgfqpoint{0.529829in}{1.313894in}}%
\pgfpathlineto{\pgfqpoint{0.531142in}{1.319410in}}%
\pgfpathlineto{\pgfqpoint{0.532454in}{1.311911in}}%
\pgfpathlineto{\pgfqpoint{0.532891in}{1.317490in}}%
\pgfpathlineto{\pgfqpoint{0.533547in}{1.313998in}}%
\pgfpathlineto{\pgfqpoint{0.533766in}{1.313976in}}%
\pgfpathlineto{\pgfqpoint{0.533984in}{1.312764in}}%
\pgfpathlineto{\pgfqpoint{0.534203in}{1.314703in}}%
\pgfpathlineto{\pgfqpoint{0.535297in}{1.322093in}}%
\pgfpathlineto{\pgfqpoint{0.535515in}{1.320023in}}%
\pgfpathlineto{\pgfqpoint{0.536609in}{1.310491in}}%
\pgfpathlineto{\pgfqpoint{0.538796in}{1.317475in}}%
\pgfpathlineto{\pgfqpoint{0.549648in}{1.313790in}}%
\pgfpathlineto{\pgfqpoint{0.560499in}{1.317114in}}%
\pgfpathlineto{\pgfqpoint{0.571351in}{1.317210in}}%
\pgfpathlineto{\pgfqpoint{0.582203in}{1.313465in}}%
\pgfpathlineto{\pgfqpoint{0.593055in}{1.314858in}}%
\pgfpathlineto{\pgfqpoint{0.603907in}{1.313192in}}%
\pgfpathlineto{\pgfqpoint{0.614759in}{1.313880in}}%
\pgfpathlineto{\pgfqpoint{0.625611in}{1.315811in}}%
\pgfpathlineto{\pgfqpoint{0.636463in}{1.314186in}}%
\pgfpathlineto{\pgfqpoint{0.647315in}{1.317643in}}%
\pgfpathlineto{\pgfqpoint{0.669018in}{1.314792in}}%
\pgfpathlineto{\pgfqpoint{0.679870in}{1.316392in}}%
\pgfpathlineto{\pgfqpoint{0.690722in}{1.315183in}}%
\pgfpathlineto{\pgfqpoint{0.723278in}{1.314978in}}%
\pgfpathlineto{\pgfqpoint{0.734130in}{1.318168in}}%
\pgfpathlineto{\pgfqpoint{0.755834in}{1.313431in}}%
\pgfpathlineto{\pgfqpoint{0.777538in}{1.312891in}}%
\pgfpathlineto{\pgfqpoint{0.788389in}{1.314182in}}%
\pgfpathlineto{\pgfqpoint{0.799241in}{1.312563in}}%
\pgfpathlineto{\pgfqpoint{0.810093in}{1.312733in}}%
\pgfpathlineto{\pgfqpoint{0.820945in}{1.315382in}}%
\pgfpathlineto{\pgfqpoint{0.831797in}{1.314679in}}%
\pgfpathlineto{\pgfqpoint{0.853501in}{1.310925in}}%
\pgfpathlineto{\pgfqpoint{0.864353in}{1.311287in}}%
\pgfpathlineto{\pgfqpoint{0.875205in}{1.313853in}}%
\pgfpathlineto{\pgfqpoint{0.896909in}{1.313159in}}%
\pgfpathlineto{\pgfqpoint{0.907760in}{1.314504in}}%
\pgfpathlineto{\pgfqpoint{0.929464in}{1.313521in}}%
\pgfpathlineto{\pgfqpoint{0.940316in}{1.311536in}}%
\pgfpathlineto{\pgfqpoint{0.951168in}{1.311695in}}%
\pgfpathlineto{\pgfqpoint{0.962020in}{1.313714in}}%
\pgfpathlineto{\pgfqpoint{1.005428in}{1.313162in}}%
\pgfpathlineto{\pgfqpoint{1.016280in}{1.312605in}}%
\pgfpathlineto{\pgfqpoint{1.027131in}{1.310865in}}%
\pgfpathlineto{\pgfqpoint{1.037983in}{1.310838in}}%
\pgfpathlineto{\pgfqpoint{1.059687in}{1.312301in}}%
\pgfpathlineto{\pgfqpoint{1.092243in}{1.312180in}}%
\pgfpathlineto{\pgfqpoint{1.124799in}{1.310417in}}%
\pgfpathlineto{\pgfqpoint{1.146502in}{1.312073in}}%
\pgfpathlineto{\pgfqpoint{1.179058in}{1.312004in}}%
\pgfpathlineto{\pgfqpoint{1.211614in}{1.309710in}}%
\pgfpathlineto{\pgfqpoint{1.233318in}{1.311552in}}%
\pgfpathlineto{\pgfqpoint{1.265873in}{1.310552in}}%
\pgfpathlineto{\pgfqpoint{1.298429in}{1.308521in}}%
\pgfpathlineto{\pgfqpoint{1.330985in}{1.310762in}}%
\pgfpathlineto{\pgfqpoint{1.374392in}{1.308298in}}%
\pgfpathlineto{\pgfqpoint{1.396096in}{1.309232in}}%
\pgfpathlineto{\pgfqpoint{1.417800in}{1.310474in}}%
\pgfpathlineto{\pgfqpoint{1.461208in}{1.307889in}}%
\pgfpathlineto{\pgfqpoint{1.482912in}{1.308572in}}%
\pgfpathlineto{\pgfqpoint{1.504615in}{1.309547in}}%
\pgfpathlineto{\pgfqpoint{1.548023in}{1.307009in}}%
\pgfpathlineto{\pgfqpoint{1.569727in}{1.307787in}}%
\pgfpathlineto{\pgfqpoint{1.591431in}{1.308799in}}%
\pgfpathlineto{\pgfqpoint{1.634838in}{1.306441in}}%
\pgfpathlineto{\pgfqpoint{1.656542in}{1.307282in}}%
\pgfpathlineto{\pgfqpoint{1.678246in}{1.308173in}}%
\pgfpathlineto{\pgfqpoint{1.721654in}{1.305652in}}%
\pgfpathlineto{\pgfqpoint{1.775913in}{1.306382in}}%
\pgfpathlineto{\pgfqpoint{1.797617in}{1.304754in}}%
\pgfpathlineto{\pgfqpoint{1.830173in}{1.305979in}}%
\pgfpathlineto{\pgfqpoint{1.851876in}{1.306502in}}%
\pgfpathlineto{\pgfqpoint{1.895284in}{1.304345in}}%
\pgfpathlineto{\pgfqpoint{1.938692in}{1.305725in}}%
\pgfpathlineto{\pgfqpoint{1.971247in}{1.303364in}}%
\pgfpathlineto{\pgfqpoint{2.036359in}{1.304040in}}%
\pgfpathlineto{\pgfqpoint{2.058063in}{1.302600in}}%
\pgfpathlineto{\pgfqpoint{2.112322in}{1.304081in}}%
\pgfpathlineto{\pgfqpoint{2.144878in}{1.301907in}}%
\pgfpathlineto{\pgfqpoint{2.199137in}{1.303228in}}%
\pgfpathlineto{\pgfqpoint{2.231693in}{1.301156in}}%
\pgfpathlineto{\pgfqpoint{2.264249in}{1.302887in}}%
\pgfpathlineto{\pgfqpoint{2.264249in}{1.302887in}}%
\pgfusepath{stroke}%
\end{pgfscope}%
\begin{pgfscope}%
\pgfpathrectangle{\pgfqpoint{0.514740in}{0.496542in}}{\pgfqpoint{1.749509in}{1.239972in}}%
\pgfusepath{clip}%
\pgfsetbuttcap%
\pgfsetroundjoin%
\pgfsetlinewidth{0.903375pt}%
\definecolor{currentstroke}{rgb}{0.000000,0.000000,0.000000}%
\pgfsetstrokecolor{currentstroke}%
\pgfsetdash{{3.330000pt}{1.440000pt}}{0.000000pt}%
\pgfpathmoveto{\pgfqpoint{0.514740in}{1.323191in}}%
\pgfpathlineto{\pgfqpoint{0.536609in}{1.244525in}}%
\pgfpathlineto{\pgfqpoint{0.560499in}{1.167116in}}%
\pgfpathlineto{\pgfqpoint{0.582203in}{1.103761in}}%
\pgfpathlineto{\pgfqpoint{0.603907in}{1.046391in}}%
\pgfpathlineto{\pgfqpoint{0.625611in}{0.994442in}}%
\pgfpathlineto{\pgfqpoint{0.647315in}{0.947401in}}%
\pgfpathlineto{\pgfqpoint{0.669018in}{0.904804in}}%
\pgfpathlineto{\pgfqpoint{0.690722in}{0.866232in}}%
\pgfpathlineto{\pgfqpoint{0.712426in}{0.831304in}}%
\pgfpathlineto{\pgfqpoint{0.734130in}{0.799675in}}%
\pgfpathlineto{\pgfqpoint{0.755834in}{0.771036in}}%
\pgfpathlineto{\pgfqpoint{0.777538in}{0.745102in}}%
\pgfpathlineto{\pgfqpoint{0.799241in}{0.721618in}}%
\pgfpathlineto{\pgfqpoint{0.831797in}{0.690486in}}%
\pgfpathlineto{\pgfqpoint{0.864353in}{0.663660in}}%
\pgfpathlineto{\pgfqpoint{0.896909in}{0.640545in}}%
\pgfpathlineto{\pgfqpoint{0.929464in}{0.620627in}}%
\pgfpathlineto{\pgfqpoint{0.962020in}{0.603464in}}%
\pgfpathlineto{\pgfqpoint{0.994576in}{0.588675in}}%
\pgfpathlineto{\pgfqpoint{1.027131in}{0.575932in}}%
\pgfpathlineto{\pgfqpoint{1.070539in}{0.561639in}}%
\pgfpathlineto{\pgfqpoint{1.113947in}{0.549919in}}%
\pgfpathlineto{\pgfqpoint{1.157354in}{0.540310in}}%
\pgfpathlineto{\pgfqpoint{1.211614in}{0.530693in}}%
\pgfpathlineto{\pgfqpoint{1.265873in}{0.523189in}}%
\pgfpathlineto{\pgfqpoint{1.330985in}{0.516328in}}%
\pgfpathlineto{\pgfqpoint{1.406948in}{0.510522in}}%
\pgfpathlineto{\pgfqpoint{1.504615in}{0.505486in}}%
\pgfpathlineto{\pgfqpoint{1.623986in}{0.501724in}}%
\pgfpathlineto{\pgfqpoint{1.786765in}{0.499004in}}%
\pgfpathlineto{\pgfqpoint{2.036359in}{0.497329in}}%
\pgfpathlineto{\pgfqpoint{2.264249in}{0.496820in}}%
\pgfpathlineto{\pgfqpoint{2.264249in}{0.496820in}}%
\pgfusepath{stroke}%
\end{pgfscope}%
\begin{pgfscope}%
\pgfsetrectcap%
\pgfsetmiterjoin%
\pgfsetlinewidth{0.803000pt}%
\definecolor{currentstroke}{rgb}{0.000000,0.000000,0.000000}%
\pgfsetstrokecolor{currentstroke}%
\pgfsetdash{}{0pt}%
\pgfpathmoveto{\pgfqpoint{0.514740in}{0.496542in}}%
\pgfpathlineto{\pgfqpoint{0.514740in}{1.736515in}}%
\pgfusepath{stroke}%
\end{pgfscope}%
\begin{pgfscope}%
\pgfsetrectcap%
\pgfsetmiterjoin%
\pgfsetlinewidth{0.803000pt}%
\definecolor{currentstroke}{rgb}{0.000000,0.000000,0.000000}%
\pgfsetstrokecolor{currentstroke}%
\pgfsetdash{}{0pt}%
\pgfpathmoveto{\pgfqpoint{2.264249in}{0.496542in}}%
\pgfpathlineto{\pgfqpoint{2.264249in}{1.736515in}}%
\pgfusepath{stroke}%
\end{pgfscope}%
\begin{pgfscope}%
\pgfsetrectcap%
\pgfsetmiterjoin%
\pgfsetlinewidth{0.803000pt}%
\definecolor{currentstroke}{rgb}{0.000000,0.000000,0.000000}%
\pgfsetstrokecolor{currentstroke}%
\pgfsetdash{}{0pt}%
\pgfpathmoveto{\pgfqpoint{0.514740in}{0.496542in}}%
\pgfpathlineto{\pgfqpoint{2.264249in}{0.496542in}}%
\pgfusepath{stroke}%
\end{pgfscope}%
\begin{pgfscope}%
\pgfsetrectcap%
\pgfsetmiterjoin%
\pgfsetlinewidth{0.803000pt}%
\definecolor{currentstroke}{rgb}{0.000000,0.000000,0.000000}%
\pgfsetstrokecolor{currentstroke}%
\pgfsetdash{}{0pt}%
\pgfpathmoveto{\pgfqpoint{0.514740in}{1.736515in}}%
\pgfpathlineto{\pgfqpoint{2.264249in}{1.736515in}}%
\pgfusepath{stroke}%
\end{pgfscope}%
\begin{pgfscope}%
\definecolor{textcolor}{rgb}{0.000000,0.000000,0.000000}%
\pgfsetstrokecolor{textcolor}%
\pgfsetfillcolor{textcolor}%
\pgftext[x=2.211764in,y=1.699316in,right,top]{\color{textcolor}{\rmfamily\fontsize{9.000000}{10.800000}\selectfont\catcode`\^=\active\def^{\ifmmode\sp\else\^{}\fi}\catcode`\%=\active\def%{\%}$(a)$}}%
\end{pgfscope}%
\begin{pgfscope}%
\pgfsetrectcap%
\pgfsetroundjoin%
\pgfsetlinewidth{0.903375pt}%
\definecolor{currentstroke}{rgb}{0.121569,0.466667,0.705882}%
\pgfsetstrokecolor{currentstroke}%
\pgfsetdash{}{0pt}%
\pgfpathmoveto{\pgfqpoint{1.677637in}{0.988348in}}%
\pgfpathlineto{\pgfqpoint{1.766526in}{0.988348in}}%
\pgfpathlineto{\pgfqpoint{1.855414in}{0.988348in}}%
\pgfusepath{stroke}%
\end{pgfscope}%
\begin{pgfscope}%
\definecolor{textcolor}{rgb}{0.000000,0.000000,0.000000}%
\pgfsetstrokecolor{textcolor}%
\pgfsetfillcolor{textcolor}%
\pgftext[x=1.944303in,y=0.949459in,left,base]{\color{textcolor}{\rmfamily\fontsize{8.000000}{9.600000}\selectfont\catcode`\^=\active\def^{\ifmmode\sp\else\^{}\fi}\catcode`\%=\active\def%{\%}$D_1$}}%
\end{pgfscope}%
\begin{pgfscope}%
\pgfsetbuttcap%
\pgfsetroundjoin%
\pgfsetlinewidth{0.903375pt}%
\definecolor{currentstroke}{rgb}{1.000000,0.000000,0.000000}%
\pgfsetstrokecolor{currentstroke}%
\pgfsetdash{{5.760000pt}{1.440000pt}{0.900000pt}{1.440000pt}}{0.000000pt}%
\pgfpathmoveto{\pgfqpoint{1.677637in}{0.861188in}}%
\pgfpathlineto{\pgfqpoint{1.766526in}{0.861188in}}%
\pgfpathlineto{\pgfqpoint{1.855414in}{0.861188in}}%
\pgfusepath{stroke}%
\end{pgfscope}%
\begin{pgfscope}%
\definecolor{textcolor}{rgb}{0.000000,0.000000,0.000000}%
\pgfsetstrokecolor{textcolor}%
\pgfsetfillcolor{textcolor}%
\pgftext[x=1.944303in,y=0.822299in,left,base]{\color{textcolor}{\rmfamily\fontsize{8.000000}{9.600000}\selectfont\catcode`\^=\active\def^{\ifmmode\sp\else\^{}\fi}\catcode`\%=\active\def%{\%}$D_2$}}%
\end{pgfscope}%
\begin{pgfscope}%
\pgfsetbuttcap%
\pgfsetroundjoin%
\pgfsetlinewidth{0.903375pt}%
\definecolor{currentstroke}{rgb}{0.000000,0.000000,0.000000}%
\pgfsetstrokecolor{currentstroke}%
\pgfsetdash{{3.330000pt}{1.440000pt}}{0.000000pt}%
\pgfpathmoveto{\pgfqpoint{1.677637in}{0.734012in}}%
\pgfpathlineto{\pgfqpoint{1.766526in}{0.734012in}}%
\pgfpathlineto{\pgfqpoint{1.855414in}{0.734012in}}%
\pgfusepath{stroke}%
\end{pgfscope}%
\begin{pgfscope}%
\definecolor{textcolor}{rgb}{0.000000,0.000000,0.000000}%
\pgfsetstrokecolor{textcolor}%
\pgfsetfillcolor{textcolor}%
\pgftext[x=1.944303in,y=0.695123in,left,base]{\color{textcolor}{\rmfamily\fontsize{8.000000}{9.600000}\selectfont\catcode`\^=\active\def^{\ifmmode\sp\else\^{}\fi}\catcode`\%=\active\def%{\%}$e^{-\kappa t}$}}%
\end{pgfscope}%
\begin{pgfscope}%
\pgfsetbuttcap%
\pgfsetmiterjoin%
\definecolor{currentfill}{rgb}{1.000000,1.000000,1.000000}%
\pgfsetfillcolor{currentfill}%
\pgfsetlinewidth{0.000000pt}%
\definecolor{currentstroke}{rgb}{0.000000,0.000000,0.000000}%
\pgfsetstrokecolor{currentstroke}%
\pgfsetstrokeopacity{0.000000}%
\pgfsetdash{}{0pt}%
\pgfpathmoveto{\pgfqpoint{2.854484in}{0.496542in}}%
\pgfpathlineto{\pgfqpoint{4.603993in}{0.496542in}}%
\pgfpathlineto{\pgfqpoint{4.603993in}{1.736515in}}%
\pgfpathlineto{\pgfqpoint{2.854484in}{1.736515in}}%
\pgfpathlineto{\pgfqpoint{2.854484in}{0.496542in}}%
\pgfpathclose%
\pgfusepath{fill}%
\end{pgfscope}%
\begin{pgfscope}%
\pgfsetbuttcap%
\pgfsetroundjoin%
\definecolor{currentfill}{rgb}{0.000000,0.000000,0.000000}%
\pgfsetfillcolor{currentfill}%
\pgfsetlinewidth{0.803000pt}%
\definecolor{currentstroke}{rgb}{0.000000,0.000000,0.000000}%
\pgfsetstrokecolor{currentstroke}%
\pgfsetdash{}{0pt}%
\pgfsys@defobject{currentmarker}{\pgfqpoint{0.000000in}{-0.048611in}}{\pgfqpoint{0.000000in}{0.000000in}}{%
\pgfpathmoveto{\pgfqpoint{0.000000in}{0.000000in}}%
\pgfpathlineto{\pgfqpoint{0.000000in}{-0.048611in}}%
\pgfusepath{stroke,fill}%
}%
\begin{pgfscope}%
\pgfsys@transformshift{2.854484in}{0.496542in}%
\pgfsys@useobject{currentmarker}{}%
\end{pgfscope}%
\end{pgfscope}%
\begin{pgfscope}%
\definecolor{textcolor}{rgb}{0.000000,0.000000,0.000000}%
\pgfsetstrokecolor{textcolor}%
\pgfsetfillcolor{textcolor}%
\pgftext[x=2.854484in,y=0.399320in,,top]{\color{textcolor}{\rmfamily\fontsize{8.000000}{9.600000}\selectfont\catcode`\^=\active\def^{\ifmmode\sp\else\^{}\fi}\catcode`\%=\active\def%{\%}0}}%
\end{pgfscope}%
\begin{pgfscope}%
\pgfsetbuttcap%
\pgfsetroundjoin%
\definecolor{currentfill}{rgb}{0.000000,0.000000,0.000000}%
\pgfsetfillcolor{currentfill}%
\pgfsetlinewidth{0.803000pt}%
\definecolor{currentstroke}{rgb}{0.000000,0.000000,0.000000}%
\pgfsetstrokecolor{currentstroke}%
\pgfsetdash{}{0pt}%
\pgfsys@defobject{currentmarker}{\pgfqpoint{0.000000in}{-0.048611in}}{\pgfqpoint{0.000000in}{0.000000in}}{%
\pgfpathmoveto{\pgfqpoint{0.000000in}{0.000000in}}%
\pgfpathlineto{\pgfqpoint{0.000000in}{-0.048611in}}%
\pgfusepath{stroke,fill}%
}%
\begin{pgfscope}%
\pgfsys@transformshift{3.291862in}{0.496542in}%
\pgfsys@useobject{currentmarker}{}%
\end{pgfscope}%
\end{pgfscope}%
\begin{pgfscope}%
\definecolor{textcolor}{rgb}{0.000000,0.000000,0.000000}%
\pgfsetstrokecolor{textcolor}%
\pgfsetfillcolor{textcolor}%
\pgftext[x=3.291862in,y=0.399320in,,top]{\color{textcolor}{\rmfamily\fontsize{8.000000}{9.600000}\selectfont\catcode`\^=\active\def^{\ifmmode\sp\else\^{}\fi}\catcode`\%=\active\def%{\%}2}}%
\end{pgfscope}%
\begin{pgfscope}%
\pgfsetbuttcap%
\pgfsetroundjoin%
\definecolor{currentfill}{rgb}{0.000000,0.000000,0.000000}%
\pgfsetfillcolor{currentfill}%
\pgfsetlinewidth{0.803000pt}%
\definecolor{currentstroke}{rgb}{0.000000,0.000000,0.000000}%
\pgfsetstrokecolor{currentstroke}%
\pgfsetdash{}{0pt}%
\pgfsys@defobject{currentmarker}{\pgfqpoint{0.000000in}{-0.048611in}}{\pgfqpoint{0.000000in}{0.000000in}}{%
\pgfpathmoveto{\pgfqpoint{0.000000in}{0.000000in}}%
\pgfpathlineto{\pgfqpoint{0.000000in}{-0.048611in}}%
\pgfusepath{stroke,fill}%
}%
\begin{pgfscope}%
\pgfsys@transformshift{3.729239in}{0.496542in}%
\pgfsys@useobject{currentmarker}{}%
\end{pgfscope}%
\end{pgfscope}%
\begin{pgfscope}%
\definecolor{textcolor}{rgb}{0.000000,0.000000,0.000000}%
\pgfsetstrokecolor{textcolor}%
\pgfsetfillcolor{textcolor}%
\pgftext[x=3.729239in,y=0.399320in,,top]{\color{textcolor}{\rmfamily\fontsize{8.000000}{9.600000}\selectfont\catcode`\^=\active\def^{\ifmmode\sp\else\^{}\fi}\catcode`\%=\active\def%{\%}4}}%
\end{pgfscope}%
\begin{pgfscope}%
\pgfsetbuttcap%
\pgfsetroundjoin%
\definecolor{currentfill}{rgb}{0.000000,0.000000,0.000000}%
\pgfsetfillcolor{currentfill}%
\pgfsetlinewidth{0.803000pt}%
\definecolor{currentstroke}{rgb}{0.000000,0.000000,0.000000}%
\pgfsetstrokecolor{currentstroke}%
\pgfsetdash{}{0pt}%
\pgfsys@defobject{currentmarker}{\pgfqpoint{0.000000in}{-0.048611in}}{\pgfqpoint{0.000000in}{0.000000in}}{%
\pgfpathmoveto{\pgfqpoint{0.000000in}{0.000000in}}%
\pgfpathlineto{\pgfqpoint{0.000000in}{-0.048611in}}%
\pgfusepath{stroke,fill}%
}%
\begin{pgfscope}%
\pgfsys@transformshift{4.166616in}{0.496542in}%
\pgfsys@useobject{currentmarker}{}%
\end{pgfscope}%
\end{pgfscope}%
\begin{pgfscope}%
\definecolor{textcolor}{rgb}{0.000000,0.000000,0.000000}%
\pgfsetstrokecolor{textcolor}%
\pgfsetfillcolor{textcolor}%
\pgftext[x=4.166616in,y=0.399320in,,top]{\color{textcolor}{\rmfamily\fontsize{8.000000}{9.600000}\selectfont\catcode`\^=\active\def^{\ifmmode\sp\else\^{}\fi}\catcode`\%=\active\def%{\%}6}}%
\end{pgfscope}%
\begin{pgfscope}%
\pgfsetbuttcap%
\pgfsetroundjoin%
\definecolor{currentfill}{rgb}{0.000000,0.000000,0.000000}%
\pgfsetfillcolor{currentfill}%
\pgfsetlinewidth{0.803000pt}%
\definecolor{currentstroke}{rgb}{0.000000,0.000000,0.000000}%
\pgfsetstrokecolor{currentstroke}%
\pgfsetdash{}{0pt}%
\pgfsys@defobject{currentmarker}{\pgfqpoint{0.000000in}{-0.048611in}}{\pgfqpoint{0.000000in}{0.000000in}}{%
\pgfpathmoveto{\pgfqpoint{0.000000in}{0.000000in}}%
\pgfpathlineto{\pgfqpoint{0.000000in}{-0.048611in}}%
\pgfusepath{stroke,fill}%
}%
\begin{pgfscope}%
\pgfsys@transformshift{4.603993in}{0.496542in}%
\pgfsys@useobject{currentmarker}{}%
\end{pgfscope}%
\end{pgfscope}%
\begin{pgfscope}%
\definecolor{textcolor}{rgb}{0.000000,0.000000,0.000000}%
\pgfsetstrokecolor{textcolor}%
\pgfsetfillcolor{textcolor}%
\pgftext[x=4.603993in,y=0.399320in,,top]{\color{textcolor}{\rmfamily\fontsize{8.000000}{9.600000}\selectfont\catcode`\^=\active\def^{\ifmmode\sp\else\^{}\fi}\catcode`\%=\active\def%{\%}8}}%
\end{pgfscope}%
\begin{pgfscope}%
\definecolor{textcolor}{rgb}{0.000000,0.000000,0.000000}%
\pgfsetstrokecolor{textcolor}%
\pgfsetfillcolor{textcolor}%
\pgftext[x=3.729239in,y=0.244999in,,top]{\color{textcolor}{\rmfamily\fontsize{9.000000}{10.800000}\selectfont\catcode`\^=\active\def^{\ifmmode\sp\else\^{}\fi}\catcode`\%=\active\def%{\%}$T\ (10^3/\omega_b)$}}%
\end{pgfscope}%
\begin{pgfscope}%
\pgfsetbuttcap%
\pgfsetroundjoin%
\definecolor{currentfill}{rgb}{0.000000,0.000000,0.000000}%
\pgfsetfillcolor{currentfill}%
\pgfsetlinewidth{0.803000pt}%
\definecolor{currentstroke}{rgb}{0.000000,0.000000,0.000000}%
\pgfsetstrokecolor{currentstroke}%
\pgfsetdash{}{0pt}%
\pgfsys@defobject{currentmarker}{\pgfqpoint{-0.048611in}{0.000000in}}{\pgfqpoint{-0.000000in}{0.000000in}}{%
\pgfpathmoveto{\pgfqpoint{-0.000000in}{0.000000in}}%
\pgfpathlineto{\pgfqpoint{-0.048611in}{0.000000in}}%
\pgfusepath{stroke,fill}%
}%
\begin{pgfscope}%
\pgfsys@transformshift{2.854484in}{0.496542in}%
\pgfsys@useobject{currentmarker}{}%
\end{pgfscope}%
\end{pgfscope}%
\begin{pgfscope}%
\definecolor{textcolor}{rgb}{0.000000,0.000000,0.000000}%
\pgfsetstrokecolor{textcolor}%
\pgfsetfillcolor{textcolor}%
\pgftext[x=2.698233in, y=0.457962in, left, base]{\color{textcolor}{\rmfamily\fontsize{8.000000}{9.600000}\selectfont\catcode`\^=\active\def^{\ifmmode\sp\else\^{}\fi}\catcode`\%=\active\def%{\%}$0$}}%
\end{pgfscope}%
\begin{pgfscope}%
\pgfsetbuttcap%
\pgfsetroundjoin%
\definecolor{currentfill}{rgb}{0.000000,0.000000,0.000000}%
\pgfsetfillcolor{currentfill}%
\pgfsetlinewidth{0.803000pt}%
\definecolor{currentstroke}{rgb}{0.000000,0.000000,0.000000}%
\pgfsetstrokecolor{currentstroke}%
\pgfsetdash{}{0pt}%
\pgfsys@defobject{currentmarker}{\pgfqpoint{-0.048611in}{0.000000in}}{\pgfqpoint{-0.000000in}{0.000000in}}{%
\pgfpathmoveto{\pgfqpoint{-0.000000in}{0.000000in}}%
\pgfpathlineto{\pgfqpoint{-0.048611in}{0.000000in}}%
\pgfusepath{stroke,fill}%
}%
\begin{pgfscope}%
\pgfsys@transformshift{2.854484in}{0.909867in}%
\pgfsys@useobject{currentmarker}{}%
\end{pgfscope}%
\end{pgfscope}%
\begin{pgfscope}%
\definecolor{textcolor}{rgb}{0.000000,0.000000,0.000000}%
\pgfsetstrokecolor{textcolor}%
\pgfsetfillcolor{textcolor}%
\pgftext[x=2.606411in, y=0.871286in, left, base]{\color{textcolor}{\rmfamily\fontsize{8.000000}{9.600000}\selectfont\catcode`\^=\active\def^{\ifmmode\sp\else\^{}\fi}\catcode`\%=\active\def%{\%}$0.5$}}%
\end{pgfscope}%
\begin{pgfscope}%
\pgfsetbuttcap%
\pgfsetroundjoin%
\definecolor{currentfill}{rgb}{0.000000,0.000000,0.000000}%
\pgfsetfillcolor{currentfill}%
\pgfsetlinewidth{0.803000pt}%
\definecolor{currentstroke}{rgb}{0.000000,0.000000,0.000000}%
\pgfsetstrokecolor{currentstroke}%
\pgfsetdash{}{0pt}%
\pgfsys@defobject{currentmarker}{\pgfqpoint{-0.048611in}{0.000000in}}{\pgfqpoint{-0.000000in}{0.000000in}}{%
\pgfpathmoveto{\pgfqpoint{-0.000000in}{0.000000in}}%
\pgfpathlineto{\pgfqpoint{-0.048611in}{0.000000in}}%
\pgfusepath{stroke,fill}%
}%
\begin{pgfscope}%
\pgfsys@transformshift{2.854484in}{1.323191in}%
\pgfsys@useobject{currentmarker}{}%
\end{pgfscope}%
\end{pgfscope}%
\begin{pgfscope}%
\definecolor{textcolor}{rgb}{0.000000,0.000000,0.000000}%
\pgfsetstrokecolor{textcolor}%
\pgfsetfillcolor{textcolor}%
\pgftext[x=2.698233in, y=1.284610in, left, base]{\color{textcolor}{\rmfamily\fontsize{8.000000}{9.600000}\selectfont\catcode`\^=\active\def^{\ifmmode\sp\else\^{}\fi}\catcode`\%=\active\def%{\%}$1$}}%
\end{pgfscope}%
\begin{pgfscope}%
\definecolor{textcolor}{rgb}{0.000000,0.000000,0.000000}%
\pgfsetstrokecolor{textcolor}%
\pgfsetfillcolor{textcolor}%
\pgftext[x=2.550856in,y=1.116529in,,bottom,rotate=90.000000]{\color{textcolor}{\rmfamily\fontsize{9.000000}{10.800000}\selectfont\catcode`\^=\active\def^{\ifmmode\sp\else\^{}\fi}\catcode`\%=\active\def%{\%}$2|\rho_{D_1D_2}|$}}%
\end{pgfscope}%
\begin{pgfscope}%
\pgfpathrectangle{\pgfqpoint{2.854484in}{0.496542in}}{\pgfqpoint{1.749509in}{1.239972in}}%
\pgfusepath{clip}%
\pgfsetrectcap%
\pgfsetroundjoin%
\pgfsetlinewidth{0.903375pt}%
\definecolor{currentstroke}{rgb}{0.121569,0.466667,0.705882}%
\pgfsetstrokecolor{currentstroke}%
\pgfsetdash{}{0pt}%
\pgfpathmoveto{\pgfqpoint{2.854484in}{1.323191in}}%
\pgfpathlineto{\pgfqpoint{2.854922in}{1.305138in}}%
\pgfpathlineto{\pgfqpoint{2.856015in}{1.305343in}}%
\pgfpathlineto{\pgfqpoint{2.856234in}{1.304996in}}%
\pgfpathlineto{\pgfqpoint{2.856671in}{1.305907in}}%
\pgfpathlineto{\pgfqpoint{2.857109in}{1.312342in}}%
\pgfpathlineto{\pgfqpoint{2.857546in}{1.307826in}}%
\pgfpathlineto{\pgfqpoint{2.857983in}{1.300743in}}%
\pgfpathlineto{\pgfqpoint{2.858421in}{1.310432in}}%
\pgfpathlineto{\pgfqpoint{2.858639in}{1.319440in}}%
\pgfpathlineto{\pgfqpoint{2.859295in}{1.306359in}}%
\pgfpathlineto{\pgfqpoint{2.859514in}{1.313591in}}%
\pgfpathlineto{\pgfqpoint{2.859733in}{1.317200in}}%
\pgfpathlineto{\pgfqpoint{2.859952in}{1.309636in}}%
\pgfpathlineto{\pgfqpoint{2.860170in}{1.302706in}}%
\pgfpathlineto{\pgfqpoint{2.860826in}{1.310110in}}%
\pgfpathlineto{\pgfqpoint{2.861264in}{1.311918in}}%
\pgfpathlineto{\pgfqpoint{2.861482in}{1.309901in}}%
\pgfpathlineto{\pgfqpoint{2.861920in}{1.301696in}}%
\pgfpathlineto{\pgfqpoint{2.862576in}{1.309446in}}%
\pgfpathlineto{\pgfqpoint{2.862794in}{1.311998in}}%
\pgfpathlineto{\pgfqpoint{2.863013in}{1.309659in}}%
\pgfpathlineto{\pgfqpoint{2.863451in}{1.300289in}}%
\pgfpathlineto{\pgfqpoint{2.863888in}{1.316571in}}%
\pgfpathlineto{\pgfqpoint{2.866075in}{1.303887in}}%
\pgfpathlineto{\pgfqpoint{2.866512in}{1.311799in}}%
\pgfpathlineto{\pgfqpoint{2.866950in}{1.305501in}}%
\pgfpathlineto{\pgfqpoint{2.867387in}{1.300579in}}%
\pgfpathlineto{\pgfqpoint{2.867824in}{1.308222in}}%
\pgfpathlineto{\pgfqpoint{2.868043in}{1.315647in}}%
\pgfpathlineto{\pgfqpoint{2.868699in}{1.302115in}}%
\pgfpathlineto{\pgfqpoint{2.870011in}{1.311711in}}%
\pgfpathlineto{\pgfqpoint{2.870230in}{1.310945in}}%
\pgfpathlineto{\pgfqpoint{2.870449in}{1.311222in}}%
\pgfpathlineto{\pgfqpoint{2.870667in}{1.314624in}}%
\pgfpathlineto{\pgfqpoint{2.871105in}{1.304498in}}%
\pgfpathlineto{\pgfqpoint{2.871323in}{1.297736in}}%
\pgfpathlineto{\pgfqpoint{2.871979in}{1.311697in}}%
\pgfpathlineto{\pgfqpoint{2.872417in}{1.308122in}}%
\pgfpathlineto{\pgfqpoint{2.872854in}{1.300485in}}%
\pgfpathlineto{\pgfqpoint{2.873292in}{1.310346in}}%
\pgfpathlineto{\pgfqpoint{2.873510in}{1.311314in}}%
\pgfpathlineto{\pgfqpoint{2.873948in}{1.305254in}}%
\pgfpathlineto{\pgfqpoint{2.874604in}{1.308888in}}%
\pgfpathlineto{\pgfqpoint{2.875041in}{1.317537in}}%
\pgfpathlineto{\pgfqpoint{2.875478in}{1.304532in}}%
\pgfpathlineto{\pgfqpoint{2.875697in}{1.305006in}}%
\pgfpathlineto{\pgfqpoint{2.876134in}{1.313843in}}%
\pgfpathlineto{\pgfqpoint{2.876353in}{1.305041in}}%
\pgfpathlineto{\pgfqpoint{2.878540in}{1.303084in}}%
\pgfpathlineto{\pgfqpoint{2.889392in}{1.299479in}}%
\pgfpathlineto{\pgfqpoint{2.900244in}{1.307294in}}%
\pgfpathlineto{\pgfqpoint{2.911096in}{1.306617in}}%
\pgfpathlineto{\pgfqpoint{2.921948in}{1.298606in}}%
\pgfpathlineto{\pgfqpoint{2.954503in}{1.288694in}}%
\pgfpathlineto{\pgfqpoint{2.965355in}{1.293573in}}%
\pgfpathlineto{\pgfqpoint{2.976207in}{1.294479in}}%
\pgfpathlineto{\pgfqpoint{2.987059in}{1.303266in}}%
\pgfpathlineto{\pgfqpoint{2.997911in}{1.296553in}}%
\pgfpathlineto{\pgfqpoint{3.008763in}{1.294311in}}%
\pgfpathlineto{\pgfqpoint{3.019615in}{1.294682in}}%
\pgfpathlineto{\pgfqpoint{3.030467in}{1.298345in}}%
\pgfpathlineto{\pgfqpoint{3.041319in}{1.298332in}}%
\pgfpathlineto{\pgfqpoint{3.052171in}{1.299606in}}%
\pgfpathlineto{\pgfqpoint{3.063022in}{1.297461in}}%
\pgfpathlineto{\pgfqpoint{3.073874in}{1.302592in}}%
\pgfpathlineto{\pgfqpoint{3.084726in}{1.294754in}}%
\pgfpathlineto{\pgfqpoint{3.095578in}{1.291596in}}%
\pgfpathlineto{\pgfqpoint{3.106430in}{1.291165in}}%
\pgfpathlineto{\pgfqpoint{3.117282in}{1.296463in}}%
\pgfpathlineto{\pgfqpoint{3.128134in}{1.299959in}}%
\pgfpathlineto{\pgfqpoint{3.138986in}{1.301058in}}%
\pgfpathlineto{\pgfqpoint{3.149838in}{1.297718in}}%
\pgfpathlineto{\pgfqpoint{3.160690in}{1.301035in}}%
\pgfpathlineto{\pgfqpoint{3.171542in}{1.297991in}}%
\pgfpathlineto{\pgfqpoint{3.182393in}{1.297160in}}%
\pgfpathlineto{\pgfqpoint{3.193245in}{1.294548in}}%
\pgfpathlineto{\pgfqpoint{3.236653in}{1.292581in}}%
\pgfpathlineto{\pgfqpoint{3.247505in}{1.295722in}}%
\pgfpathlineto{\pgfqpoint{3.290913in}{1.289149in}}%
\pgfpathlineto{\pgfqpoint{3.301764in}{1.286623in}}%
\pgfpathlineto{\pgfqpoint{3.312616in}{1.285244in}}%
\pgfpathlineto{\pgfqpoint{3.323468in}{1.286372in}}%
\pgfpathlineto{\pgfqpoint{3.334320in}{1.288919in}}%
\pgfpathlineto{\pgfqpoint{3.345172in}{1.287563in}}%
\pgfpathlineto{\pgfqpoint{3.366876in}{1.282002in}}%
\pgfpathlineto{\pgfqpoint{3.377728in}{1.282121in}}%
\pgfpathlineto{\pgfqpoint{3.410284in}{1.285501in}}%
\pgfpathlineto{\pgfqpoint{3.431987in}{1.283591in}}%
\pgfpathlineto{\pgfqpoint{3.453691in}{1.279842in}}%
\pgfpathlineto{\pgfqpoint{3.464543in}{1.279165in}}%
\pgfpathlineto{\pgfqpoint{3.497099in}{1.282485in}}%
\pgfpathlineto{\pgfqpoint{3.518803in}{1.278556in}}%
\pgfpathlineto{\pgfqpoint{3.529654in}{1.278068in}}%
\pgfpathlineto{\pgfqpoint{3.562210in}{1.280131in}}%
\pgfpathlineto{\pgfqpoint{3.583914in}{1.280792in}}%
\pgfpathlineto{\pgfqpoint{3.605618in}{1.277536in}}%
\pgfpathlineto{\pgfqpoint{3.627322in}{1.276534in}}%
\pgfpathlineto{\pgfqpoint{3.649025in}{1.276515in}}%
\pgfpathlineto{\pgfqpoint{3.670729in}{1.277824in}}%
\pgfpathlineto{\pgfqpoint{3.692433in}{1.275420in}}%
\pgfpathlineto{\pgfqpoint{3.714137in}{1.273516in}}%
\pgfpathlineto{\pgfqpoint{3.724989in}{1.271456in}}%
\pgfpathlineto{\pgfqpoint{3.735841in}{1.270787in}}%
\pgfpathlineto{\pgfqpoint{3.757545in}{1.273436in}}%
\pgfpathlineto{\pgfqpoint{3.779248in}{1.270290in}}%
\pgfpathlineto{\pgfqpoint{3.833508in}{1.269317in}}%
\pgfpathlineto{\pgfqpoint{3.844360in}{1.269944in}}%
\pgfpathlineto{\pgfqpoint{3.866064in}{1.266132in}}%
\pgfpathlineto{\pgfqpoint{3.909471in}{1.265083in}}%
\pgfpathlineto{\pgfqpoint{3.920323in}{1.266701in}}%
\pgfpathlineto{\pgfqpoint{3.931175in}{1.266603in}}%
\pgfpathlineto{\pgfqpoint{3.952879in}{1.263056in}}%
\pgfpathlineto{\pgfqpoint{3.974583in}{1.263524in}}%
\pgfpathlineto{\pgfqpoint{3.996287in}{1.262433in}}%
\pgfpathlineto{\pgfqpoint{4.007138in}{1.263690in}}%
\pgfpathlineto{\pgfqpoint{4.017990in}{1.263381in}}%
\pgfpathlineto{\pgfqpoint{4.039694in}{1.260050in}}%
\pgfpathlineto{\pgfqpoint{4.061398in}{1.260991in}}%
\pgfpathlineto{\pgfqpoint{4.083102in}{1.259830in}}%
\pgfpathlineto{\pgfqpoint{4.093954in}{1.260729in}}%
\pgfpathlineto{\pgfqpoint{4.104806in}{1.260510in}}%
\pgfpathlineto{\pgfqpoint{4.126509in}{1.257712in}}%
\pgfpathlineto{\pgfqpoint{4.148213in}{1.257356in}}%
\pgfpathlineto{\pgfqpoint{4.159065in}{1.256219in}}%
\pgfpathlineto{\pgfqpoint{4.202473in}{1.255290in}}%
\pgfpathlineto{\pgfqpoint{4.245880in}{1.252678in}}%
\pgfpathlineto{\pgfqpoint{4.278436in}{1.253169in}}%
\pgfpathlineto{\pgfqpoint{4.300140in}{1.250899in}}%
\pgfpathlineto{\pgfqpoint{4.321844in}{1.250696in}}%
\pgfpathlineto{\pgfqpoint{4.343548in}{1.250090in}}%
\pgfpathlineto{\pgfqpoint{4.354399in}{1.250555in}}%
\pgfpathlineto{\pgfqpoint{4.419511in}{1.246853in}}%
\pgfpathlineto{\pgfqpoint{4.441215in}{1.247366in}}%
\pgfpathlineto{\pgfqpoint{4.462919in}{1.245166in}}%
\pgfpathlineto{\pgfqpoint{4.495474in}{1.244767in}}%
\pgfpathlineto{\pgfqpoint{4.517178in}{1.244248in}}%
\pgfpathlineto{\pgfqpoint{4.538882in}{1.243174in}}%
\pgfpathlineto{\pgfqpoint{4.549734in}{1.242248in}}%
\pgfpathlineto{\pgfqpoint{4.571438in}{1.242583in}}%
\pgfpathlineto{\pgfqpoint{4.593141in}{1.240831in}}%
\pgfpathlineto{\pgfqpoint{4.603993in}{1.241194in}}%
\pgfpathlineto{\pgfqpoint{4.603993in}{1.241194in}}%
\pgfusepath{stroke}%
\end{pgfscope}%
\begin{pgfscope}%
\pgfpathrectangle{\pgfqpoint{2.854484in}{0.496542in}}{\pgfqpoint{1.749509in}{1.239972in}}%
\pgfusepath{clip}%
\pgfsetbuttcap%
\pgfsetroundjoin%
\pgfsetlinewidth{0.903375pt}%
\definecolor{currentstroke}{rgb}{0.768627,0.580392,0.164706}%
\pgfsetstrokecolor{currentstroke}%
\pgfsetdash{{3.330000pt}{1.440000pt}}{0.000000pt}%
\pgfpathmoveto{\pgfqpoint{2.854484in}{1.323191in}}%
\pgfpathlineto{\pgfqpoint{2.854922in}{1.305094in}}%
\pgfpathlineto{\pgfqpoint{2.856015in}{1.305146in}}%
\pgfpathlineto{\pgfqpoint{2.856234in}{1.304769in}}%
\pgfpathlineto{\pgfqpoint{2.856671in}{1.305630in}}%
\pgfpathlineto{\pgfqpoint{2.857109in}{1.312020in}}%
\pgfpathlineto{\pgfqpoint{2.857546in}{1.307458in}}%
\pgfpathlineto{\pgfqpoint{2.857983in}{1.300318in}}%
\pgfpathlineto{\pgfqpoint{2.858421in}{1.309941in}}%
\pgfpathlineto{\pgfqpoint{2.858639in}{1.318914in}}%
\pgfpathlineto{\pgfqpoint{2.859295in}{1.305768in}}%
\pgfpathlineto{\pgfqpoint{2.859514in}{1.312974in}}%
\pgfpathlineto{\pgfqpoint{2.859733in}{1.316557in}}%
\pgfpathlineto{\pgfqpoint{2.859952in}{1.308971in}}%
\pgfpathlineto{\pgfqpoint{2.860170in}{1.302017in}}%
\pgfpathlineto{\pgfqpoint{2.860826in}{1.309321in}}%
\pgfpathlineto{\pgfqpoint{2.861264in}{1.311065in}}%
\pgfpathlineto{\pgfqpoint{2.861482in}{1.309021in}}%
\pgfpathlineto{\pgfqpoint{2.861920in}{1.300771in}}%
\pgfpathlineto{\pgfqpoint{2.862576in}{1.308432in}}%
\pgfpathlineto{\pgfqpoint{2.862794in}{1.310952in}}%
\pgfpathlineto{\pgfqpoint{2.863013in}{1.308589in}}%
\pgfpathlineto{\pgfqpoint{2.863451in}{1.299179in}}%
\pgfpathlineto{\pgfqpoint{2.863888in}{1.315396in}}%
\pgfpathlineto{\pgfqpoint{2.866075in}{1.302441in}}%
\pgfpathlineto{\pgfqpoint{2.866512in}{1.310280in}}%
\pgfpathlineto{\pgfqpoint{2.866950in}{1.303926in}}%
\pgfpathlineto{\pgfqpoint{2.867387in}{1.298959in}}%
\pgfpathlineto{\pgfqpoint{2.867824in}{1.306539in}}%
\pgfpathlineto{\pgfqpoint{2.868043in}{1.313928in}}%
\pgfpathlineto{\pgfqpoint{2.868699in}{1.300346in}}%
\pgfpathlineto{\pgfqpoint{2.870011in}{1.309771in}}%
\pgfpathlineto{\pgfqpoint{2.870230in}{1.308983in}}%
\pgfpathlineto{\pgfqpoint{2.870449in}{1.309232in}}%
\pgfpathlineto{\pgfqpoint{2.870667in}{1.312596in}}%
\pgfpathlineto{\pgfqpoint{2.871105in}{1.302419in}}%
\pgfpathlineto{\pgfqpoint{2.871323in}{1.295638in}}%
\pgfpathlineto{\pgfqpoint{2.871979in}{1.309467in}}%
\pgfpathlineto{\pgfqpoint{2.872417in}{1.305850in}}%
\pgfpathlineto{\pgfqpoint{2.872854in}{1.298192in}}%
\pgfpathlineto{\pgfqpoint{2.873292in}{1.307984in}}%
\pgfpathlineto{\pgfqpoint{2.873510in}{1.308924in}}%
\pgfpathlineto{\pgfqpoint{2.873948in}{1.302829in}}%
\pgfpathlineto{\pgfqpoint{2.874604in}{1.306378in}}%
\pgfpathlineto{\pgfqpoint{2.875041in}{1.314963in}}%
\pgfpathlineto{\pgfqpoint{2.875478in}{1.301945in}}%
\pgfpathlineto{\pgfqpoint{2.875697in}{1.302388in}}%
\pgfpathlineto{\pgfqpoint{2.876134in}{1.311120in}}%
\pgfpathlineto{\pgfqpoint{2.876353in}{1.302301in}}%
\pgfpathlineto{\pgfqpoint{2.878540in}{1.300093in}}%
\pgfpathlineto{\pgfqpoint{2.889392in}{1.295167in}}%
\pgfpathlineto{\pgfqpoint{2.900244in}{1.301638in}}%
\pgfpathlineto{\pgfqpoint{2.911096in}{1.299627in}}%
\pgfpathlineto{\pgfqpoint{2.921948in}{1.290333in}}%
\pgfpathlineto{\pgfqpoint{2.954503in}{1.276483in}}%
\pgfpathlineto{\pgfqpoint{2.965355in}{1.280054in}}%
\pgfpathlineto{\pgfqpoint{2.976207in}{1.279683in}}%
\pgfpathlineto{\pgfqpoint{2.987059in}{1.287050in}}%
\pgfpathlineto{\pgfqpoint{2.997911in}{1.279092in}}%
\pgfpathlineto{\pgfqpoint{3.008763in}{1.275516in}}%
\pgfpathlineto{\pgfqpoint{3.019615in}{1.274503in}}%
\pgfpathlineto{\pgfqpoint{3.030467in}{1.276778in}}%
\pgfpathlineto{\pgfqpoint{3.041319in}{1.275493in}}%
\pgfpathlineto{\pgfqpoint{3.052171in}{1.275502in}}%
\pgfpathlineto{\pgfqpoint{3.063022in}{1.272153in}}%
\pgfpathlineto{\pgfqpoint{3.073874in}{1.275878in}}%
\pgfpathlineto{\pgfqpoint{3.084726in}{1.266915in}}%
\pgfpathlineto{\pgfqpoint{3.095578in}{1.262501in}}%
\pgfpathlineto{\pgfqpoint{3.106430in}{1.260747in}}%
\pgfpathlineto{\pgfqpoint{3.117282in}{1.264622in}}%
\pgfpathlineto{\pgfqpoint{3.128134in}{1.266755in}}%
\pgfpathlineto{\pgfqpoint{3.138986in}{1.266562in}}%
\pgfpathlineto{\pgfqpoint{3.149838in}{1.262041in}}%
\pgfpathlineto{\pgfqpoint{3.160690in}{1.263958in}}%
\pgfpathlineto{\pgfqpoint{3.171542in}{1.259741in}}%
\pgfpathlineto{\pgfqpoint{3.182393in}{1.257673in}}%
\pgfpathlineto{\pgfqpoint{3.193245in}{1.253885in}}%
\pgfpathlineto{\pgfqpoint{3.236653in}{1.246949in}}%
\pgfpathlineto{\pgfqpoint{3.247505in}{1.248697in}}%
\pgfpathlineto{\pgfqpoint{3.290913in}{1.237562in}}%
\pgfpathlineto{\pgfqpoint{3.312616in}{1.231313in}}%
\pgfpathlineto{\pgfqpoint{3.323468in}{1.231113in}}%
\pgfpathlineto{\pgfqpoint{3.334320in}{1.232332in}}%
\pgfpathlineto{\pgfqpoint{3.345172in}{1.229907in}}%
\pgfpathlineto{\pgfqpoint{3.366876in}{1.222357in}}%
\pgfpathlineto{\pgfqpoint{3.388580in}{1.221198in}}%
\pgfpathlineto{\pgfqpoint{3.421135in}{1.218931in}}%
\pgfpathlineto{\pgfqpoint{3.464543in}{1.209021in}}%
\pgfpathlineto{\pgfqpoint{3.497099in}{1.208079in}}%
\pgfpathlineto{\pgfqpoint{3.518803in}{1.202228in}}%
\pgfpathlineto{\pgfqpoint{3.529654in}{1.200818in}}%
\pgfpathlineto{\pgfqpoint{3.551358in}{1.200018in}}%
\pgfpathlineto{\pgfqpoint{3.583914in}{1.196968in}}%
\pgfpathlineto{\pgfqpoint{3.605618in}{1.191986in}}%
\pgfpathlineto{\pgfqpoint{3.638174in}{1.187735in}}%
\pgfpathlineto{\pgfqpoint{3.692433in}{1.180941in}}%
\pgfpathlineto{\pgfqpoint{3.714137in}{1.177473in}}%
\pgfpathlineto{\pgfqpoint{3.735841in}{1.172291in}}%
\pgfpathlineto{\pgfqpoint{3.757545in}{1.171851in}}%
\pgfpathlineto{\pgfqpoint{3.779248in}{1.167420in}}%
\pgfpathlineto{\pgfqpoint{3.800952in}{1.165094in}}%
\pgfpathlineto{\pgfqpoint{3.822656in}{1.160640in}}%
\pgfpathlineto{\pgfqpoint{3.844360in}{1.159769in}}%
\pgfpathlineto{\pgfqpoint{3.866064in}{1.154991in}}%
\pgfpathlineto{\pgfqpoint{3.898619in}{1.150892in}}%
\pgfpathlineto{\pgfqpoint{3.909471in}{1.149513in}}%
\pgfpathlineto{\pgfqpoint{3.920323in}{1.149326in}}%
\pgfpathlineto{\pgfqpoint{3.931175in}{1.148002in}}%
\pgfpathlineto{\pgfqpoint{3.952879in}{1.143678in}}%
\pgfpathlineto{\pgfqpoint{3.974583in}{1.142550in}}%
\pgfpathlineto{\pgfqpoint{3.996287in}{1.138495in}}%
\pgfpathlineto{\pgfqpoint{4.017990in}{1.136545in}}%
\pgfpathlineto{\pgfqpoint{4.028842in}{1.134003in}}%
\pgfpathlineto{\pgfqpoint{4.039694in}{1.132664in}}%
\pgfpathlineto{\pgfqpoint{4.061398in}{1.131949in}}%
\pgfpathlineto{\pgfqpoint{4.083102in}{1.127669in}}%
\pgfpathlineto{\pgfqpoint{4.104806in}{1.125584in}}%
\pgfpathlineto{\pgfqpoint{4.126509in}{1.122446in}}%
\pgfpathlineto{\pgfqpoint{4.148213in}{1.120571in}}%
\pgfpathlineto{\pgfqpoint{4.169917in}{1.116170in}}%
\pgfpathlineto{\pgfqpoint{4.191621in}{1.114239in}}%
\pgfpathlineto{\pgfqpoint{4.213325in}{1.111704in}}%
\pgfpathlineto{\pgfqpoint{4.224177in}{1.111307in}}%
\pgfpathlineto{\pgfqpoint{4.289288in}{1.101346in}}%
\pgfpathlineto{\pgfqpoint{4.321844in}{1.098855in}}%
\pgfpathlineto{\pgfqpoint{4.343548in}{1.094625in}}%
\pgfpathlineto{\pgfqpoint{4.386955in}{1.090604in}}%
\pgfpathlineto{\pgfqpoint{4.397807in}{1.090230in}}%
\pgfpathlineto{\pgfqpoint{4.452067in}{1.081528in}}%
\pgfpathlineto{\pgfqpoint{4.473770in}{1.080787in}}%
\pgfpathlineto{\pgfqpoint{4.484622in}{1.080421in}}%
\pgfpathlineto{\pgfqpoint{4.528030in}{1.072610in}}%
\pgfpathlineto{\pgfqpoint{4.549734in}{1.070762in}}%
\pgfpathlineto{\pgfqpoint{4.571438in}{1.070424in}}%
\pgfpathlineto{\pgfqpoint{4.603993in}{1.063734in}}%
\pgfpathlineto{\pgfqpoint{4.603993in}{1.063734in}}%
\pgfusepath{stroke}%
\end{pgfscope}%
\begin{pgfscope}%
\pgfpathrectangle{\pgfqpoint{2.854484in}{0.496542in}}{\pgfqpoint{1.749509in}{1.239972in}}%
\pgfusepath{clip}%
\pgfsetbuttcap%
\pgfsetroundjoin%
\pgfsetlinewidth{0.903375pt}%
\definecolor{currentstroke}{rgb}{1.000000,0.000000,0.000000}%
\pgfsetstrokecolor{currentstroke}%
\pgfsetdash{{5.760000pt}{1.440000pt}{0.900000pt}{1.440000pt}}{0.000000pt}%
\pgfpathmoveto{\pgfqpoint{2.854484in}{1.323191in}}%
\pgfpathlineto{\pgfqpoint{2.856015in}{1.303374in}}%
\pgfpathlineto{\pgfqpoint{2.856453in}{1.302635in}}%
\pgfpathlineto{\pgfqpoint{2.856671in}{1.303143in}}%
\pgfpathlineto{\pgfqpoint{2.857109in}{1.309133in}}%
\pgfpathlineto{\pgfqpoint{2.857546in}{1.304152in}}%
\pgfpathlineto{\pgfqpoint{2.857983in}{1.296502in}}%
\pgfpathlineto{\pgfqpoint{2.858421in}{1.305530in}}%
\pgfpathlineto{\pgfqpoint{2.858639in}{1.314196in}}%
\pgfpathlineto{\pgfqpoint{2.859295in}{1.300468in}}%
\pgfpathlineto{\pgfqpoint{2.859514in}{1.307446in}}%
\pgfpathlineto{\pgfqpoint{2.859733in}{1.310796in}}%
\pgfpathlineto{\pgfqpoint{2.859952in}{1.303018in}}%
\pgfpathlineto{\pgfqpoint{2.860170in}{1.295847in}}%
\pgfpathlineto{\pgfqpoint{2.861045in}{1.302811in}}%
\pgfpathlineto{\pgfqpoint{2.861264in}{1.303434in}}%
\pgfpathlineto{\pgfqpoint{2.861482in}{1.301140in}}%
\pgfpathlineto{\pgfqpoint{2.861920in}{1.292494in}}%
\pgfpathlineto{\pgfqpoint{2.862576in}{1.299361in}}%
\pgfpathlineto{\pgfqpoint{2.862794in}{1.301609in}}%
\pgfpathlineto{\pgfqpoint{2.863013in}{1.299020in}}%
\pgfpathlineto{\pgfqpoint{2.863451in}{1.289263in}}%
\pgfpathlineto{\pgfqpoint{2.863888in}{1.304895in}}%
\pgfpathlineto{\pgfqpoint{2.866075in}{1.289544in}}%
\pgfpathlineto{\pgfqpoint{2.866512in}{1.296737in}}%
\pgfpathlineto{\pgfqpoint{2.866950in}{1.289892in}}%
\pgfpathlineto{\pgfqpoint{2.867387in}{1.284523in}}%
\pgfpathlineto{\pgfqpoint{2.867824in}{1.291553in}}%
\pgfpathlineto{\pgfqpoint{2.868043in}{1.298622in}}%
\pgfpathlineto{\pgfqpoint{2.868699in}{1.284602in}}%
\pgfpathlineto{\pgfqpoint{2.870011in}{1.292522in}}%
\pgfpathlineto{\pgfqpoint{2.870230in}{1.291536in}}%
\pgfpathlineto{\pgfqpoint{2.870449in}{1.291541in}}%
\pgfpathlineto{\pgfqpoint{2.870667in}{1.294566in}}%
\pgfpathlineto{\pgfqpoint{2.870886in}{1.293391in}}%
\pgfpathlineto{\pgfqpoint{2.871323in}{1.277003in}}%
\pgfpathlineto{\pgfqpoint{2.871979in}{1.289669in}}%
\pgfpathlineto{\pgfqpoint{2.872635in}{1.281135in}}%
\pgfpathlineto{\pgfqpoint{2.872854in}{1.277842in}}%
\pgfpathlineto{\pgfqpoint{2.873292in}{1.287032in}}%
\pgfpathlineto{\pgfqpoint{2.873510in}{1.287727in}}%
\pgfpathlineto{\pgfqpoint{2.873948in}{1.281327in}}%
\pgfpathlineto{\pgfqpoint{2.874604in}{1.284134in}}%
\pgfpathlineto{\pgfqpoint{2.875041in}{1.292151in}}%
\pgfpathlineto{\pgfqpoint{2.875478in}{1.279036in}}%
\pgfpathlineto{\pgfqpoint{2.875697in}{1.279201in}}%
\pgfpathlineto{\pgfqpoint{2.876134in}{1.287017in}}%
\pgfpathlineto{\pgfqpoint{2.876353in}{1.278054in}}%
\pgfpathlineto{\pgfqpoint{2.878540in}{1.273676in}}%
\pgfpathlineto{\pgfqpoint{2.889392in}{1.257391in}}%
\pgfpathlineto{\pgfqpoint{2.900244in}{1.252491in}}%
\pgfpathlineto{\pgfqpoint{2.911096in}{1.239381in}}%
\pgfpathlineto{\pgfqpoint{2.921948in}{1.219616in}}%
\pgfpathlineto{\pgfqpoint{2.943651in}{1.188688in}}%
\pgfpathlineto{\pgfqpoint{2.954503in}{1.174669in}}%
\pgfpathlineto{\pgfqpoint{2.965355in}{1.168250in}}%
\pgfpathlineto{\pgfqpoint{2.976207in}{1.158318in}}%
\pgfpathlineto{\pgfqpoint{2.987059in}{1.155092in}}%
\pgfpathlineto{\pgfqpoint{2.997911in}{1.138177in}}%
\pgfpathlineto{\pgfqpoint{3.008763in}{1.125076in}}%
\pgfpathlineto{\pgfqpoint{3.019615in}{1.114295in}}%
\pgfpathlineto{\pgfqpoint{3.030467in}{1.106899in}}%
\pgfpathlineto{\pgfqpoint{3.041319in}{1.097027in}}%
\pgfpathlineto{\pgfqpoint{3.052171in}{1.088661in}}%
\pgfpathlineto{\pgfqpoint{3.063022in}{1.077583in}}%
\pgfpathlineto{\pgfqpoint{3.073874in}{1.072101in}}%
\pgfpathlineto{\pgfqpoint{3.084726in}{1.056301in}}%
\pgfpathlineto{\pgfqpoint{3.095578in}{1.044133in}}%
\pgfpathlineto{\pgfqpoint{3.106430in}{1.034256in}}%
\pgfpathlineto{\pgfqpoint{3.128134in}{1.023277in}}%
\pgfpathlineto{\pgfqpoint{3.138986in}{1.015631in}}%
\pgfpathlineto{\pgfqpoint{3.149838in}{1.004684in}}%
\pgfpathlineto{\pgfqpoint{3.160690in}{0.998573in}}%
\pgfpathlineto{\pgfqpoint{3.171542in}{0.988099in}}%
\pgfpathlineto{\pgfqpoint{3.225801in}{0.945521in}}%
\pgfpathlineto{\pgfqpoint{3.236653in}{0.937744in}}%
\pgfpathlineto{\pgfqpoint{3.247505in}{0.932447in}}%
\pgfpathlineto{\pgfqpoint{3.312616in}{0.884919in}}%
\pgfpathlineto{\pgfqpoint{3.323468in}{0.878969in}}%
\pgfpathlineto{\pgfqpoint{3.334320in}{0.874318in}}%
\pgfpathlineto{\pgfqpoint{3.356024in}{0.860408in}}%
\pgfpathlineto{\pgfqpoint{3.366876in}{0.853100in}}%
\pgfpathlineto{\pgfqpoint{3.453691in}{0.808002in}}%
\pgfpathlineto{\pgfqpoint{3.518803in}{0.776435in}}%
\pgfpathlineto{\pgfqpoint{3.551358in}{0.765058in}}%
\pgfpathlineto{\pgfqpoint{3.594766in}{0.744655in}}%
\pgfpathlineto{\pgfqpoint{3.638174in}{0.730728in}}%
\pgfpathlineto{\pgfqpoint{3.670729in}{0.716212in}}%
\pgfpathlineto{\pgfqpoint{3.681581in}{0.712522in}}%
\pgfpathlineto{\pgfqpoint{3.714137in}{0.704918in}}%
\pgfpathlineto{\pgfqpoint{3.757545in}{0.687053in}}%
\pgfpathlineto{\pgfqpoint{3.768396in}{0.684146in}}%
\pgfpathlineto{\pgfqpoint{3.800952in}{0.678437in}}%
\pgfpathlineto{\pgfqpoint{3.844360in}{0.661846in}}%
\pgfpathlineto{\pgfqpoint{3.866064in}{0.658473in}}%
\pgfpathlineto{\pgfqpoint{3.876916in}{0.657496in}}%
\pgfpathlineto{\pgfqpoint{3.887767in}{0.655378in}}%
\pgfpathlineto{\pgfqpoint{3.931175in}{0.639888in}}%
\pgfpathlineto{\pgfqpoint{3.942027in}{0.638216in}}%
\pgfpathlineto{\pgfqpoint{3.963731in}{0.637551in}}%
\pgfpathlineto{\pgfqpoint{3.974583in}{0.635701in}}%
\pgfpathlineto{\pgfqpoint{4.017990in}{0.620843in}}%
\pgfpathlineto{\pgfqpoint{4.028842in}{0.619877in}}%
\pgfpathlineto{\pgfqpoint{4.050546in}{0.620122in}}%
\pgfpathlineto{\pgfqpoint{4.061398in}{0.618397in}}%
\pgfpathlineto{\pgfqpoint{4.104806in}{0.604380in}}%
\pgfpathlineto{\pgfqpoint{4.148213in}{0.603037in}}%
\pgfpathlineto{\pgfqpoint{4.180769in}{0.591385in}}%
\pgfpathlineto{\pgfqpoint{4.191621in}{0.590013in}}%
\pgfpathlineto{\pgfqpoint{4.224177in}{0.591657in}}%
\pgfpathlineto{\pgfqpoint{4.235028in}{0.589652in}}%
\pgfpathlineto{\pgfqpoint{4.267584in}{0.578424in}}%
\pgfpathlineto{\pgfqpoint{4.278436in}{0.577559in}}%
\pgfpathlineto{\pgfqpoint{4.310992in}{0.580000in}}%
\pgfpathlineto{\pgfqpoint{4.321844in}{0.577988in}}%
\pgfpathlineto{\pgfqpoint{4.354399in}{0.567200in}}%
\pgfpathlineto{\pgfqpoint{4.365251in}{0.566786in}}%
\pgfpathlineto{\pgfqpoint{4.397807in}{0.569906in}}%
\pgfpathlineto{\pgfqpoint{4.408659in}{0.567764in}}%
\pgfpathlineto{\pgfqpoint{4.430363in}{0.559840in}}%
\pgfpathlineto{\pgfqpoint{4.441215in}{0.557436in}}%
\pgfpathlineto{\pgfqpoint{4.452067in}{0.557485in}}%
\pgfpathlineto{\pgfqpoint{4.473770in}{0.561046in}}%
\pgfpathlineto{\pgfqpoint{4.484622in}{0.561114in}}%
\pgfpathlineto{\pgfqpoint{4.495474in}{0.558807in}}%
\pgfpathlineto{\pgfqpoint{4.517178in}{0.550971in}}%
\pgfpathlineto{\pgfqpoint{4.528030in}{0.548977in}}%
\pgfpathlineto{\pgfqpoint{4.538882in}{0.549541in}}%
\pgfpathlineto{\pgfqpoint{4.560586in}{0.553472in}}%
\pgfpathlineto{\pgfqpoint{4.571438in}{0.553366in}}%
\pgfpathlineto{\pgfqpoint{4.582290in}{0.550890in}}%
\pgfpathlineto{\pgfqpoint{4.603993in}{0.543235in}}%
\pgfpathlineto{\pgfqpoint{4.603993in}{0.543235in}}%
\pgfusepath{stroke}%
\end{pgfscope}%
\begin{pgfscope}%
\pgfpathrectangle{\pgfqpoint{2.854484in}{0.496542in}}{\pgfqpoint{1.749509in}{1.239972in}}%
\pgfusepath{clip}%
\pgfsetbuttcap%
\pgfsetroundjoin%
\pgfsetlinewidth{0.903375pt}%
\definecolor{currentstroke}{rgb}{0.000000,0.000000,0.000000}%
\pgfsetstrokecolor{currentstroke}%
\pgfsetdash{{0.900000pt}{1.485000pt}}{0.000000pt}%
\pgfpathmoveto{\pgfqpoint{2.854484in}{1.323191in}}%
\pgfpathlineto{\pgfqpoint{2.889392in}{1.280356in}}%
\pgfpathlineto{\pgfqpoint{2.932800in}{1.230175in}}%
\pgfpathlineto{\pgfqpoint{2.976207in}{1.183206in}}%
\pgfpathlineto{\pgfqpoint{3.019615in}{1.139244in}}%
\pgfpathlineto{\pgfqpoint{3.063022in}{1.098097in}}%
\pgfpathlineto{\pgfqpoint{3.106430in}{1.059584in}}%
\pgfpathlineto{\pgfqpoint{3.149838in}{1.023537in}}%
\pgfpathlineto{\pgfqpoint{3.204097in}{0.981706in}}%
\pgfpathlineto{\pgfqpoint{3.258357in}{0.943195in}}%
\pgfpathlineto{\pgfqpoint{3.312616in}{0.907741in}}%
\pgfpathlineto{\pgfqpoint{3.366876in}{0.875102in}}%
\pgfpathlineto{\pgfqpoint{3.421135in}{0.845053in}}%
\pgfpathlineto{\pgfqpoint{3.475395in}{0.817389in}}%
\pgfpathlineto{\pgfqpoint{3.529654in}{0.791922in}}%
\pgfpathlineto{\pgfqpoint{3.594766in}{0.764014in}}%
\pgfpathlineto{\pgfqpoint{3.659877in}{0.738744in}}%
\pgfpathlineto{\pgfqpoint{3.724989in}{0.715861in}}%
\pgfpathlineto{\pgfqpoint{3.790100in}{0.695140in}}%
\pgfpathlineto{\pgfqpoint{3.866064in}{0.673426in}}%
\pgfpathlineto{\pgfqpoint{3.942027in}{0.654087in}}%
\pgfpathlineto{\pgfqpoint{4.017990in}{0.636862in}}%
\pgfpathlineto{\pgfqpoint{4.104806in}{0.619470in}}%
\pgfpathlineto{\pgfqpoint{4.191621in}{0.604233in}}%
\pgfpathlineto{\pgfqpoint{4.289288in}{0.589338in}}%
\pgfpathlineto{\pgfqpoint{4.397807in}{0.575191in}}%
\pgfpathlineto{\pgfqpoint{4.517178in}{0.562107in}}%
\pgfpathlineto{\pgfqpoint{4.603993in}{0.553981in}}%
\pgfpathlineto{\pgfqpoint{4.603993in}{0.553981in}}%
\pgfusepath{stroke}%
\end{pgfscope}%
\begin{pgfscope}%
\pgfsetrectcap%
\pgfsetmiterjoin%
\pgfsetlinewidth{0.803000pt}%
\definecolor{currentstroke}{rgb}{0.000000,0.000000,0.000000}%
\pgfsetstrokecolor{currentstroke}%
\pgfsetdash{}{0pt}%
\pgfpathmoveto{\pgfqpoint{2.854484in}{0.496542in}}%
\pgfpathlineto{\pgfqpoint{2.854484in}{1.736515in}}%
\pgfusepath{stroke}%
\end{pgfscope}%
\begin{pgfscope}%
\pgfsetrectcap%
\pgfsetmiterjoin%
\pgfsetlinewidth{0.803000pt}%
\definecolor{currentstroke}{rgb}{0.000000,0.000000,0.000000}%
\pgfsetstrokecolor{currentstroke}%
\pgfsetdash{}{0pt}%
\pgfpathmoveto{\pgfqpoint{4.603993in}{0.496542in}}%
\pgfpathlineto{\pgfqpoint{4.603993in}{1.736515in}}%
\pgfusepath{stroke}%
\end{pgfscope}%
\begin{pgfscope}%
\pgfsetrectcap%
\pgfsetmiterjoin%
\pgfsetlinewidth{0.803000pt}%
\definecolor{currentstroke}{rgb}{0.000000,0.000000,0.000000}%
\pgfsetstrokecolor{currentstroke}%
\pgfsetdash{}{0pt}%
\pgfpathmoveto{\pgfqpoint{2.854484in}{0.496542in}}%
\pgfpathlineto{\pgfqpoint{4.603993in}{0.496542in}}%
\pgfusepath{stroke}%
\end{pgfscope}%
\begin{pgfscope}%
\pgfsetrectcap%
\pgfsetmiterjoin%
\pgfsetlinewidth{0.803000pt}%
\definecolor{currentstroke}{rgb}{0.000000,0.000000,0.000000}%
\pgfsetstrokecolor{currentstroke}%
\pgfsetdash{}{0pt}%
\pgfpathmoveto{\pgfqpoint{2.854484in}{1.736515in}}%
\pgfpathlineto{\pgfqpoint{4.603993in}{1.736515in}}%
\pgfusepath{stroke}%
\end{pgfscope}%
\begin{pgfscope}%
\definecolor{textcolor}{rgb}{0.000000,0.000000,0.000000}%
\pgfsetstrokecolor{textcolor}%
\pgfsetfillcolor{textcolor}%
\pgftext[x=4.551508in,y=1.699316in,right,top]{\color{textcolor}{\rmfamily\fontsize{9.000000}{10.800000}\selectfont\catcode`\^=\active\def^{\ifmmode\sp\else\^{}\fi}\catcode`\%=\active\def%{\%}$(b)$}}%
\end{pgfscope}%
\begin{pgfscope}%
\pgfsetrectcap%
\pgfsetroundjoin%
\pgfsetlinewidth{0.903375pt}%
\definecolor{currentstroke}{rgb}{0.121569,0.466667,0.705882}%
\pgfsetstrokecolor{currentstroke}%
\pgfsetdash{}{0pt}%
\pgfpathmoveto{\pgfqpoint{2.922540in}{1.622032in}}%
\pgfpathlineto{\pgfqpoint{2.990595in}{1.622032in}}%
\pgfpathlineto{\pgfqpoint{3.058651in}{1.622032in}}%
\pgfusepath{stroke}%
\end{pgfscope}%
\begin{pgfscope}%
\definecolor{textcolor}{rgb}{0.000000,0.000000,0.000000}%
\pgfsetstrokecolor{textcolor}%
\pgfsetfillcolor{textcolor}%
\pgftext[x=3.136429in,y=1.588004in,left,base]{\color{textcolor}{\rmfamily\fontsize{7.000000}{8.400000}\selectfont\catcode`\^=\active\def^{\ifmmode\sp\else\^{}\fi}\catcode`\%=\active\def%{\%}$\chi=1$}}%
\end{pgfscope}%
\begin{pgfscope}%
\pgfsetbuttcap%
\pgfsetroundjoin%
\pgfsetlinewidth{0.903375pt}%
\definecolor{currentstroke}{rgb}{0.768627,0.580392,0.164706}%
\pgfsetstrokecolor{currentstroke}%
\pgfsetdash{{3.330000pt}{1.440000pt}}{0.000000pt}%
\pgfpathmoveto{\pgfqpoint{2.922540in}{1.510766in}}%
\pgfpathlineto{\pgfqpoint{2.990595in}{1.510766in}}%
\pgfpathlineto{\pgfqpoint{3.058651in}{1.510766in}}%
\pgfusepath{stroke}%
\end{pgfscope}%
\begin{pgfscope}%
\definecolor{textcolor}{rgb}{0.000000,0.000000,0.000000}%
\pgfsetstrokecolor{textcolor}%
\pgfsetfillcolor{textcolor}%
\pgftext[x=3.136429in,y=1.476739in,left,base]{\color{textcolor}{\rmfamily\fontsize{7.000000}{8.400000}\selectfont\catcode`\^=\active\def^{\ifmmode\sp\else\^{}\fi}\catcode`\%=\active\def%{\%}$\chi=0.9$}}%
\end{pgfscope}%
\begin{pgfscope}%
\pgfsetbuttcap%
\pgfsetroundjoin%
\pgfsetlinewidth{0.903375pt}%
\definecolor{currentstroke}{rgb}{1.000000,0.000000,0.000000}%
\pgfsetstrokecolor{currentstroke}%
\pgfsetdash{{5.760000pt}{1.440000pt}{0.900000pt}{1.440000pt}}{0.000000pt}%
\pgfpathmoveto{\pgfqpoint{3.575806in}{1.622032in}}%
\pgfpathlineto{\pgfqpoint{3.643862in}{1.622032in}}%
\pgfpathlineto{\pgfqpoint{3.711917in}{1.622032in}}%
\pgfusepath{stroke}%
\end{pgfscope}%
\begin{pgfscope}%
\definecolor{textcolor}{rgb}{0.000000,0.000000,0.000000}%
\pgfsetstrokecolor{textcolor}%
\pgfsetfillcolor{textcolor}%
\pgftext[x=3.789695in,y=1.588004in,left,base]{\color{textcolor}{\rmfamily\fontsize{7.000000}{8.400000}\selectfont\catcode`\^=\active\def^{\ifmmode\sp\else\^{}\fi}\catcode`\%=\active\def%{\%}$\chi=0$}}%
\end{pgfscope}%
\begin{pgfscope}%
\pgfsetbuttcap%
\pgfsetroundjoin%
\pgfsetlinewidth{0.903375pt}%
\definecolor{currentstroke}{rgb}{0.000000,0.000000,0.000000}%
\pgfsetstrokecolor{currentstroke}%
\pgfsetdash{{0.900000pt}{1.485000pt}}{0.000000pt}%
\pgfpathmoveto{\pgfqpoint{3.575806in}{1.510766in}}%
\pgfpathlineto{\pgfqpoint{3.643862in}{1.510766in}}%
\pgfpathlineto{\pgfqpoint{3.711917in}{1.510766in}}%
\pgfusepath{stroke}%
\end{pgfscope}%
\begin{pgfscope}%
\definecolor{textcolor}{rgb}{0.000000,0.000000,0.000000}%
\pgfsetstrokecolor{textcolor}%
\pgfsetfillcolor{textcolor}%
\pgftext[x=3.789695in,y=1.476739in,left,base]{\color{textcolor}{\rmfamily\fontsize{7.000000}{8.400000}\selectfont\catcode`\^=\active\def^{\ifmmode\sp\else\^{}\fi}\catcode`\%=\active\def%{\%}secular}}%
\end{pgfscope}%
\end{pgfpicture}%
\makeatother%
\endgroup%
}
	\caption{(En color en l\'inea) Retenci\'on de las ranuras oscuras durante la espera con el l\'aser apagado. (a) Poblaci\'on del estado preparado $D_1$ (azul, continua) y $D_2$ (rojo, punto-raya) y decaimiento $e^{-\kappa t}$ de un fon\'on pasivo (negro, discontinuo). (b) Coherencia $2|\rho_{D_1D_2}|$ para $\chi=1$ (azul, continua), $0.9$ (dorado, discontinua) y $0$ (rojo, punto-raya); la l\'inea de puntos es la ley secular de la Ec.~\eqref{eq:gamma_dark} para $\chi=0$. Par\'ametros del sistema como en la Fig.~\ref{fig:raman_map}, m\'as las tasas de relajaci\'on por el continuo de fonones a temperatura cero.}
	\label{fig:storage}
\end{figure}

Adem\'as del modo discreto $b$, los puntos se acoplan al continuo de fonones longitudinales. Calculamos la tasa de relajaci\'on a temperatura cero $\Gamma_{a\to b}$ entre autoestados con la regla de oro de Fermi para el acoplo de potencial de deformaci\'on en volumen, con densidad espectral $J_{\rm LA}(\omega)\propto\omega^3$ a baja frecuencia~\cite{krummheuer2002} y un factor de estructura que depende del solapamiento de los dos estados y de las distancias $r_{ij}$ entre puntos. Para $qr_{ij}\ll1$, con $q=(E_a-E_b)/\hbar c$, la tasa escala como $(E_a-E_b)^5$, de modo que las separaciones peque\~nas entre ranuras resultan favorecidas. El confinamiento y las superficies libres cercanas pueden modificar este modelo de continuo~\cite{krummheuer2005}, y para las constantes del material y la geometr\'ia estimadas aqu\'i la tasa entre las ranuras es despreciable frente a los otros canales, por lo que no extrapolamos un valor absoluto. Las tasas entran en las simulaciones como saltos proyectados $L_{a\to b}=\sqrt{\Gamma_{a\to b}}\ket b\bra a$, que suponen un continuo markoviano y separaciones que resuelven las transiciones; el ba\~no est\'a a temperatura cero.

Un fon\'on itinerante de envolvente normalizada $u(t)$ se representa con una cavidad fuente $a_s$ en cascada con el modo $b$, siguiendo el formalismo de sistemas en cascada y de paquetes de onda de un solo cuanto~\cite{gardiner1993,carmichael1993,gough2012},
\begin{align}
L_{\rm cas}&=\sqrt{\gamma_s(t)}\,a_s+\sqrt{\kappa_{\rm in}}\,b,\nonumber\\
H_{\rm cas}&=-\frac i2\sqrt{\gamma_s(t)\kappa_{\rm in}}\,(b^\dagger a_s-a_s^\dagger b),\nonumber\\
\gamma_s(t)&=\frac{|u(t)|^2}{1-\int_0^t|u(t')|^2\mathrm dt'}.\label{eq:cascade}
\end{align}
El signo del t\'ermino coherente fija la direcci\'on de la cascada. La carga usa un paquete exponencial creciente truncado tras una ventana $T_L$, motivado por la emisi\'on de una cavidad invertida en el tiempo y la transferencia con adaptaci\'on de impedancia~\cite{cirac1997}, de modo que queda una fracci\'on $\epsilon$ en la fuente al final. La lectura invierte los pulsos Raman y abre un puerto de salida $\kappa_{\rm out}$ que captura un paquete exponencial decreciente en una cavidad sumidero, regularizado con el mismo $\epsilon$. Las eficiencias de estos puertos aislados, $\eta_L$ y $\eta_R$, son cerca de $0.96$ para $\kappa_{\rm in}T_L=\kappa_{\rm out}t=4$ y $\epsilon=0.02$; el canal completo incluye adem\'as la escritura Raman, la espera y la lectura Raman. Durante la espera usamos $\kappa=10^{-3}\,\omega_b$, y durante la carga y la extracci\'on $\kappa_{\rm in}=\kappa_{\rm out}=10\kappa$, con una conmutaci\'on instant\'anea del puerto.

La Figura~\ref{fig:protocol} muestra la secuencia completa para una espera $T\omega_b=2000$. La ocupaci\'on del modo sube durante cada carga y cae cuando el pulso Raman transfiere el fon\'on a su ranura; la poblaci\'on de $D_1$ condicionada a $R=0$ aparece en la primera escritura y la de $D_2$ condicionada a $R=1$ en la segunda. La coherencia entre las dos componentes existe solo mientras ambas est\'an almacenadas: aparece tras la segunda escritura, se mantiene casi constante durante la espera, de acuerdo con la Fig.~\ref{fig:storage}(b), y desaparece cuando la lectura vac\'ia las ranuras. Su valor es menor del que permitir\'ian las poblaciones por s\'i solas: la visibilidad se pierde sobre todo mientras la segunda componente se carga y se escribe, cuando una rama contiene un ex\'citon y la otra no, y apenas durante la espera.

\begin{figure}[t]
\centering
\resizebox{\columnwidth}{!}{%% Creator: Matplotlib, PGF backend
%%
%% To include the figure in your LaTeX document, write
%%   \input{<filename>.pgf}
%%
%% Make sure the required packages are loaded in your preamble
%%   \usepackage{pgf}
%%
%% Also ensure that all the required font packages are loaded; for instance,
%% the lmodern package is sometimes necessary when using math font.
%%   \usepackage{lmodern}
%%
%% Figures using additional raster images can only be included by \input if
%% they are in the same directory as the main LaTeX file. For loading figures
%% from other directories you can use the `import` package
%%   \usepackage{import}
%%
%% and then include the figures with
%%   \import{<path to file>}{<filename>.pgf}
%%
%% Matplotlib used the following preamble
%%   \def\mathdefault#1{#1}
%%   \everymath=\expandafter{\the\everymath\displaystyle}
%%   \IfFileExists{scrextend.sty}{
%%     \usepackage[fontsize=8.000000pt]{scrextend}
%%   }{
%%     \renewcommand{\normalsize}{\fontsize{8.000000}{9.600000}\selectfont}
%%     \normalsize
%%   }
%%   \usepackage{amsmath,amssymb}
%%   \makeatletter\@ifpackageloaded{underscore}{}{\usepackage[strings]{underscore}}\makeatother
%%
\begingroup%
\makeatletter%
\begin{pgfpicture}%
\pgfpathrectangle{\pgfpointorigin}{\pgfqpoint{4.635868in}{2.729451in}}%
\pgfusepath{use as bounding box, clip}%
\begin{pgfscope}%
\pgfsetbuttcap%
\pgfsetmiterjoin%
\definecolor{currentfill}{rgb}{1.000000,1.000000,1.000000}%
\pgfsetfillcolor{currentfill}%
\pgfsetlinewidth{0.000000pt}%
\definecolor{currentstroke}{rgb}{1.000000,1.000000,1.000000}%
\pgfsetstrokecolor{currentstroke}%
\pgfsetdash{}{0pt}%
\pgfpathmoveto{\pgfqpoint{0.000000in}{0.000000in}}%
\pgfpathlineto{\pgfqpoint{4.635868in}{0.000000in}}%
\pgfpathlineto{\pgfqpoint{4.635868in}{2.729451in}}%
\pgfpathlineto{\pgfqpoint{0.000000in}{2.729451in}}%
\pgfpathlineto{\pgfqpoint{0.000000in}{0.000000in}}%
\pgfpathclose%
\pgfusepath{fill}%
\end{pgfscope}%
\begin{pgfscope}%
\pgfsetbuttcap%
\pgfsetmiterjoin%
\definecolor{currentfill}{rgb}{1.000000,1.000000,1.000000}%
\pgfsetfillcolor{currentfill}%
\pgfsetlinewidth{0.000000pt}%
\definecolor{currentstroke}{rgb}{0.000000,0.000000,0.000000}%
\pgfsetstrokecolor{currentstroke}%
\pgfsetstrokeopacity{0.000000}%
\pgfsetdash{}{0pt}%
\pgfpathmoveto{\pgfqpoint{0.514740in}{1.372588in}}%
\pgfpathlineto{\pgfqpoint{4.535868in}{1.372588in}}%
\pgfpathlineto{\pgfqpoint{4.535868in}{2.486143in}}%
\pgfpathlineto{\pgfqpoint{0.514740in}{2.486143in}}%
\pgfpathlineto{\pgfqpoint{0.514740in}{1.372588in}}%
\pgfpathclose%
\pgfusepath{fill}%
\end{pgfscope}%
\begin{pgfscope}%
\pgfpathrectangle{\pgfqpoint{0.514740in}{1.372588in}}{\pgfqpoint{4.021128in}{1.113555in}}%
\pgfusepath{clip}%
\pgfsetbuttcap%
\pgfsetmiterjoin%
\definecolor{currentfill}{rgb}{0.901961,0.937255,0.972549}%
\pgfsetfillcolor{currentfill}%
\pgfsetfillopacity{0.700000}%
\pgfsetlinewidth{0.000000pt}%
\definecolor{currentstroke}{rgb}{0.901961,0.937255,0.972549}%
\pgfsetstrokecolor{currentstroke}%
\pgfsetstrokeopacity{0.700000}%
\pgfsetdash{}{0pt}%
\pgfpathmoveto{\pgfqpoint{0.514740in}{1.372588in}}%
\pgfpathlineto{\pgfqpoint{0.883823in}{1.372588in}}%
\pgfpathlineto{\pgfqpoint{0.883823in}{2.486143in}}%
\pgfpathlineto{\pgfqpoint{0.514740in}{2.486143in}}%
\pgfpathlineto{\pgfqpoint{0.514740in}{1.372588in}}%
\pgfpathclose%
\pgfusepath{fill}%
\end{pgfscope}%
\begin{pgfscope}%
\pgfpathrectangle{\pgfqpoint{0.514740in}{1.372588in}}{\pgfqpoint{4.021128in}{1.113555in}}%
\pgfusepath{clip}%
\pgfsetbuttcap%
\pgfsetmiterjoin%
\definecolor{currentfill}{rgb}{0.901961,0.937255,0.972549}%
\pgfsetfillcolor{currentfill}%
\pgfsetfillopacity{0.700000}%
\pgfsetlinewidth{0.000000pt}%
\definecolor{currentstroke}{rgb}{0.901961,0.937255,0.972549}%
\pgfsetstrokecolor{currentstroke}%
\pgfsetstrokeopacity{0.700000}%
\pgfsetdash{}{0pt}%
\pgfpathmoveto{\pgfqpoint{1.134218in}{1.372588in}}%
\pgfpathlineto{\pgfqpoint{1.503301in}{1.372588in}}%
\pgfpathlineto{\pgfqpoint{1.503301in}{2.486143in}}%
\pgfpathlineto{\pgfqpoint{1.134218in}{2.486143in}}%
\pgfpathlineto{\pgfqpoint{1.134218in}{1.372588in}}%
\pgfpathclose%
\pgfusepath{fill}%
\end{pgfscope}%
\begin{pgfscope}%
\pgfpathrectangle{\pgfqpoint{0.514740in}{1.372588in}}{\pgfqpoint{4.021128in}{1.113555in}}%
\pgfusepath{clip}%
\pgfsetbuttcap%
\pgfsetmiterjoin%
\definecolor{currentfill}{rgb}{0.921569,0.960784,0.937255}%
\pgfsetfillcolor{currentfill}%
\pgfsetfillopacity{0.750000}%
\pgfsetlinewidth{0.000000pt}%
\definecolor{currentstroke}{rgb}{0.921569,0.960784,0.937255}%
\pgfsetstrokecolor{currentstroke}%
\pgfsetstrokeopacity{0.750000}%
\pgfsetdash{}{0pt}%
\pgfpathmoveto{\pgfqpoint{1.642409in}{1.372588in}}%
\pgfpathlineto{\pgfqpoint{3.497191in}{1.372588in}}%
\pgfpathlineto{\pgfqpoint{3.497191in}{2.486143in}}%
\pgfpathlineto{\pgfqpoint{1.642409in}{2.486143in}}%
\pgfpathlineto{\pgfqpoint{1.642409in}{1.372588in}}%
\pgfpathclose%
\pgfusepath{fill}%
\end{pgfscope}%
\begin{pgfscope}%
\pgfpathrectangle{\pgfqpoint{0.514740in}{1.372588in}}{\pgfqpoint{4.021128in}{1.113555in}}%
\pgfusepath{clip}%
\pgfsetbuttcap%
\pgfsetmiterjoin%
\definecolor{currentfill}{rgb}{0.901961,0.937255,0.972549}%
\pgfsetfillcolor{currentfill}%
\pgfsetfillopacity{0.700000}%
\pgfsetlinewidth{0.000000pt}%
\definecolor{currentstroke}{rgb}{0.901961,0.937255,0.972549}%
\pgfsetstrokecolor{currentstroke}%
\pgfsetstrokeopacity{0.700000}%
\pgfsetdash{}{0pt}%
\pgfpathmoveto{\pgfqpoint{3.636299in}{1.372588in}}%
\pgfpathlineto{\pgfqpoint{4.007256in}{1.372588in}}%
\pgfpathlineto{\pgfqpoint{4.007256in}{2.486143in}}%
\pgfpathlineto{\pgfqpoint{3.636299in}{2.486143in}}%
\pgfpathlineto{\pgfqpoint{3.636299in}{1.372588in}}%
\pgfpathclose%
\pgfusepath{fill}%
\end{pgfscope}%
\begin{pgfscope}%
\pgfpathrectangle{\pgfqpoint{0.514740in}{1.372588in}}{\pgfqpoint{4.021128in}{1.113555in}}%
\pgfusepath{clip}%
\pgfsetbuttcap%
\pgfsetmiterjoin%
\definecolor{currentfill}{rgb}{0.901961,0.937255,0.972549}%
\pgfsetfillcolor{currentfill}%
\pgfsetfillopacity{0.700000}%
\pgfsetlinewidth{0.000000pt}%
\definecolor{currentstroke}{rgb}{0.901961,0.937255,0.972549}%
\pgfsetstrokecolor{currentstroke}%
\pgfsetstrokeopacity{0.700000}%
\pgfsetdash{}{0pt}%
\pgfpathmoveto{\pgfqpoint{4.164912in}{1.372588in}}%
\pgfpathlineto{\pgfqpoint{4.535868in}{1.372588in}}%
\pgfpathlineto{\pgfqpoint{4.535868in}{2.486143in}}%
\pgfpathlineto{\pgfqpoint{4.164912in}{2.486143in}}%
\pgfpathlineto{\pgfqpoint{4.164912in}{1.372588in}}%
\pgfpathclose%
\pgfusepath{fill}%
\end{pgfscope}%
\begin{pgfscope}%
\pgfsetbuttcap%
\pgfsetroundjoin%
\definecolor{currentfill}{rgb}{0.000000,0.000000,0.000000}%
\pgfsetfillcolor{currentfill}%
\pgfsetlinewidth{0.803000pt}%
\definecolor{currentstroke}{rgb}{0.000000,0.000000,0.000000}%
\pgfsetstrokecolor{currentstroke}%
\pgfsetdash{}{0pt}%
\pgfsys@defobject{currentmarker}{\pgfqpoint{0.000000in}{-0.048611in}}{\pgfqpoint{0.000000in}{0.000000in}}{%
\pgfpathmoveto{\pgfqpoint{0.000000in}{0.000000in}}%
\pgfpathlineto{\pgfqpoint{0.000000in}{-0.048611in}}%
\pgfusepath{stroke,fill}%
}%
\begin{pgfscope}%
\pgfsys@transformshift{0.514740in}{1.372588in}%
\pgfsys@useobject{currentmarker}{}%
\end{pgfscope}%
\end{pgfscope}%
\begin{pgfscope}%
\pgfsetbuttcap%
\pgfsetroundjoin%
\definecolor{currentfill}{rgb}{0.000000,0.000000,0.000000}%
\pgfsetfillcolor{currentfill}%
\pgfsetlinewidth{0.803000pt}%
\definecolor{currentstroke}{rgb}{0.000000,0.000000,0.000000}%
\pgfsetstrokecolor{currentstroke}%
\pgfsetdash{}{0pt}%
\pgfsys@defobject{currentmarker}{\pgfqpoint{0.000000in}{-0.048611in}}{\pgfqpoint{0.000000in}{0.000000in}}{%
\pgfpathmoveto{\pgfqpoint{0.000000in}{0.000000in}}%
\pgfpathlineto{\pgfqpoint{0.000000in}{-0.048611in}}%
\pgfusepath{stroke,fill}%
}%
\begin{pgfscope}%
\pgfsys@transformshift{0.978435in}{1.372588in}%
\pgfsys@useobject{currentmarker}{}%
\end{pgfscope}%
\end{pgfscope}%
\begin{pgfscope}%
\pgfsetbuttcap%
\pgfsetroundjoin%
\definecolor{currentfill}{rgb}{0.000000,0.000000,0.000000}%
\pgfsetfillcolor{currentfill}%
\pgfsetlinewidth{0.803000pt}%
\definecolor{currentstroke}{rgb}{0.000000,0.000000,0.000000}%
\pgfsetstrokecolor{currentstroke}%
\pgfsetdash{}{0pt}%
\pgfsys@defobject{currentmarker}{\pgfqpoint{0.000000in}{-0.048611in}}{\pgfqpoint{0.000000in}{0.000000in}}{%
\pgfpathmoveto{\pgfqpoint{0.000000in}{0.000000in}}%
\pgfpathlineto{\pgfqpoint{0.000000in}{-0.048611in}}%
\pgfusepath{stroke,fill}%
}%
\begin{pgfscope}%
\pgfsys@transformshift{1.442131in}{1.372588in}%
\pgfsys@useobject{currentmarker}{}%
\end{pgfscope}%
\end{pgfscope}%
\begin{pgfscope}%
\pgfsetbuttcap%
\pgfsetroundjoin%
\definecolor{currentfill}{rgb}{0.000000,0.000000,0.000000}%
\pgfsetfillcolor{currentfill}%
\pgfsetlinewidth{0.803000pt}%
\definecolor{currentstroke}{rgb}{0.000000,0.000000,0.000000}%
\pgfsetstrokecolor{currentstroke}%
\pgfsetdash{}{0pt}%
\pgfsys@defobject{currentmarker}{\pgfqpoint{0.000000in}{-0.048611in}}{\pgfqpoint{0.000000in}{0.000000in}}{%
\pgfpathmoveto{\pgfqpoint{0.000000in}{0.000000in}}%
\pgfpathlineto{\pgfqpoint{0.000000in}{-0.048611in}}%
\pgfusepath{stroke,fill}%
}%
\begin{pgfscope}%
\pgfsys@transformshift{1.905826in}{1.372588in}%
\pgfsys@useobject{currentmarker}{}%
\end{pgfscope}%
\end{pgfscope}%
\begin{pgfscope}%
\pgfsetbuttcap%
\pgfsetroundjoin%
\definecolor{currentfill}{rgb}{0.000000,0.000000,0.000000}%
\pgfsetfillcolor{currentfill}%
\pgfsetlinewidth{0.803000pt}%
\definecolor{currentstroke}{rgb}{0.000000,0.000000,0.000000}%
\pgfsetstrokecolor{currentstroke}%
\pgfsetdash{}{0pt}%
\pgfsys@defobject{currentmarker}{\pgfqpoint{0.000000in}{-0.048611in}}{\pgfqpoint{0.000000in}{0.000000in}}{%
\pgfpathmoveto{\pgfqpoint{0.000000in}{0.000000in}}%
\pgfpathlineto{\pgfqpoint{0.000000in}{-0.048611in}}%
\pgfusepath{stroke,fill}%
}%
\begin{pgfscope}%
\pgfsys@transformshift{2.369521in}{1.372588in}%
\pgfsys@useobject{currentmarker}{}%
\end{pgfscope}%
\end{pgfscope}%
\begin{pgfscope}%
\pgfsetbuttcap%
\pgfsetroundjoin%
\definecolor{currentfill}{rgb}{0.000000,0.000000,0.000000}%
\pgfsetfillcolor{currentfill}%
\pgfsetlinewidth{0.803000pt}%
\definecolor{currentstroke}{rgb}{0.000000,0.000000,0.000000}%
\pgfsetstrokecolor{currentstroke}%
\pgfsetdash{}{0pt}%
\pgfsys@defobject{currentmarker}{\pgfqpoint{0.000000in}{-0.048611in}}{\pgfqpoint{0.000000in}{0.000000in}}{%
\pgfpathmoveto{\pgfqpoint{0.000000in}{0.000000in}}%
\pgfpathlineto{\pgfqpoint{0.000000in}{-0.048611in}}%
\pgfusepath{stroke,fill}%
}%
\begin{pgfscope}%
\pgfsys@transformshift{2.833217in}{1.372588in}%
\pgfsys@useobject{currentmarker}{}%
\end{pgfscope}%
\end{pgfscope}%
\begin{pgfscope}%
\pgfsetbuttcap%
\pgfsetroundjoin%
\definecolor{currentfill}{rgb}{0.000000,0.000000,0.000000}%
\pgfsetfillcolor{currentfill}%
\pgfsetlinewidth{0.803000pt}%
\definecolor{currentstroke}{rgb}{0.000000,0.000000,0.000000}%
\pgfsetstrokecolor{currentstroke}%
\pgfsetdash{}{0pt}%
\pgfsys@defobject{currentmarker}{\pgfqpoint{0.000000in}{-0.048611in}}{\pgfqpoint{0.000000in}{0.000000in}}{%
\pgfpathmoveto{\pgfqpoint{0.000000in}{0.000000in}}%
\pgfpathlineto{\pgfqpoint{0.000000in}{-0.048611in}}%
\pgfusepath{stroke,fill}%
}%
\begin{pgfscope}%
\pgfsys@transformshift{3.296912in}{1.372588in}%
\pgfsys@useobject{currentmarker}{}%
\end{pgfscope}%
\end{pgfscope}%
\begin{pgfscope}%
\pgfsetbuttcap%
\pgfsetroundjoin%
\definecolor{currentfill}{rgb}{0.000000,0.000000,0.000000}%
\pgfsetfillcolor{currentfill}%
\pgfsetlinewidth{0.803000pt}%
\definecolor{currentstroke}{rgb}{0.000000,0.000000,0.000000}%
\pgfsetstrokecolor{currentstroke}%
\pgfsetdash{}{0pt}%
\pgfsys@defobject{currentmarker}{\pgfqpoint{0.000000in}{-0.048611in}}{\pgfqpoint{0.000000in}{0.000000in}}{%
\pgfpathmoveto{\pgfqpoint{0.000000in}{0.000000in}}%
\pgfpathlineto{\pgfqpoint{0.000000in}{-0.048611in}}%
\pgfusepath{stroke,fill}%
}%
\begin{pgfscope}%
\pgfsys@transformshift{3.760607in}{1.372588in}%
\pgfsys@useobject{currentmarker}{}%
\end{pgfscope}%
\end{pgfscope}%
\begin{pgfscope}%
\pgfsetbuttcap%
\pgfsetroundjoin%
\definecolor{currentfill}{rgb}{0.000000,0.000000,0.000000}%
\pgfsetfillcolor{currentfill}%
\pgfsetlinewidth{0.803000pt}%
\definecolor{currentstroke}{rgb}{0.000000,0.000000,0.000000}%
\pgfsetstrokecolor{currentstroke}%
\pgfsetdash{}{0pt}%
\pgfsys@defobject{currentmarker}{\pgfqpoint{0.000000in}{-0.048611in}}{\pgfqpoint{0.000000in}{0.000000in}}{%
\pgfpathmoveto{\pgfqpoint{0.000000in}{0.000000in}}%
\pgfpathlineto{\pgfqpoint{0.000000in}{-0.048611in}}%
\pgfusepath{stroke,fill}%
}%
\begin{pgfscope}%
\pgfsys@transformshift{4.224303in}{1.372588in}%
\pgfsys@useobject{currentmarker}{}%
\end{pgfscope}%
\end{pgfscope}%
\begin{pgfscope}%
\pgfsetbuttcap%
\pgfsetroundjoin%
\definecolor{currentfill}{rgb}{0.000000,0.000000,0.000000}%
\pgfsetfillcolor{currentfill}%
\pgfsetlinewidth{0.803000pt}%
\definecolor{currentstroke}{rgb}{0.000000,0.000000,0.000000}%
\pgfsetstrokecolor{currentstroke}%
\pgfsetdash{}{0pt}%
\pgfsys@defobject{currentmarker}{\pgfqpoint{-0.048611in}{0.000000in}}{\pgfqpoint{-0.000000in}{0.000000in}}{%
\pgfpathmoveto{\pgfqpoint{-0.000000in}{0.000000in}}%
\pgfpathlineto{\pgfqpoint{-0.048611in}{0.000000in}}%
\pgfusepath{stroke,fill}%
}%
\begin{pgfscope}%
\pgfsys@transformshift{0.514740in}{1.399313in}%
\pgfsys@useobject{currentmarker}{}%
\end{pgfscope}%
\end{pgfscope}%
\begin{pgfscope}%
\definecolor{textcolor}{rgb}{0.000000,0.000000,0.000000}%
\pgfsetstrokecolor{textcolor}%
\pgfsetfillcolor{textcolor}%
\pgftext[x=0.358489in, y=1.360733in, left, base]{\color{textcolor}{\rmfamily\fontsize{8.000000}{9.600000}\selectfont\catcode`\^=\active\def^{\ifmmode\sp\else\^{}\fi}\catcode`\%=\active\def%{\%}$0$}}%
\end{pgfscope}%
\begin{pgfscope}%
\pgfsetbuttcap%
\pgfsetroundjoin%
\definecolor{currentfill}{rgb}{0.000000,0.000000,0.000000}%
\pgfsetfillcolor{currentfill}%
\pgfsetlinewidth{0.803000pt}%
\definecolor{currentstroke}{rgb}{0.000000,0.000000,0.000000}%
\pgfsetstrokecolor{currentstroke}%
\pgfsetdash{}{0pt}%
\pgfsys@defobject{currentmarker}{\pgfqpoint{-0.048611in}{0.000000in}}{\pgfqpoint{-0.000000in}{0.000000in}}{%
\pgfpathmoveto{\pgfqpoint{-0.000000in}{0.000000in}}%
\pgfpathlineto{\pgfqpoint{-0.048611in}{0.000000in}}%
\pgfusepath{stroke,fill}%
}%
\begin{pgfscope}%
\pgfsys@transformshift{0.514740in}{1.844736in}%
\pgfsys@useobject{currentmarker}{}%
\end{pgfscope}%
\end{pgfscope}%
\begin{pgfscope}%
\definecolor{textcolor}{rgb}{0.000000,0.000000,0.000000}%
\pgfsetstrokecolor{textcolor}%
\pgfsetfillcolor{textcolor}%
\pgftext[x=0.266667in, y=1.806155in, left, base]{\color{textcolor}{\rmfamily\fontsize{8.000000}{9.600000}\selectfont\catcode`\^=\active\def^{\ifmmode\sp\else\^{}\fi}\catcode`\%=\active\def%{\%}$0.5$}}%
\end{pgfscope}%
\begin{pgfscope}%
\pgfsetbuttcap%
\pgfsetroundjoin%
\definecolor{currentfill}{rgb}{0.000000,0.000000,0.000000}%
\pgfsetfillcolor{currentfill}%
\pgfsetlinewidth{0.803000pt}%
\definecolor{currentstroke}{rgb}{0.000000,0.000000,0.000000}%
\pgfsetstrokecolor{currentstroke}%
\pgfsetdash{}{0pt}%
\pgfsys@defobject{currentmarker}{\pgfqpoint{-0.048611in}{0.000000in}}{\pgfqpoint{-0.000000in}{0.000000in}}{%
\pgfpathmoveto{\pgfqpoint{-0.000000in}{0.000000in}}%
\pgfpathlineto{\pgfqpoint{-0.048611in}{0.000000in}}%
\pgfusepath{stroke,fill}%
}%
\begin{pgfscope}%
\pgfsys@transformshift{0.514740in}{2.290158in}%
\pgfsys@useobject{currentmarker}{}%
\end{pgfscope}%
\end{pgfscope}%
\begin{pgfscope}%
\definecolor{textcolor}{rgb}{0.000000,0.000000,0.000000}%
\pgfsetstrokecolor{textcolor}%
\pgfsetfillcolor{textcolor}%
\pgftext[x=0.358489in, y=2.251577in, left, base]{\color{textcolor}{\rmfamily\fontsize{8.000000}{9.600000}\selectfont\catcode`\^=\active\def^{\ifmmode\sp\else\^{}\fi}\catcode`\%=\active\def%{\%}$1$}}%
\end{pgfscope}%
\begin{pgfscope}%
\definecolor{textcolor}{rgb}{0.000000,0.000000,0.000000}%
\pgfsetstrokecolor{textcolor}%
\pgfsetfillcolor{textcolor}%
\pgftext[x=0.211111in,y=1.929366in,,bottom,rotate=90.000000]{\color{textcolor}{\rmfamily\fontsize{9.000000}{10.800000}\selectfont\catcode`\^=\active\def^{\ifmmode\sp\else\^{}\fi}\catcode`\%=\active\def%{\%}Populations}}%
\end{pgfscope}%
\begin{pgfscope}%
\pgfpathrectangle{\pgfqpoint{0.514740in}{1.372588in}}{\pgfqpoint{4.021128in}{1.113555in}}%
\pgfusepath{clip}%
\pgfsetbuttcap%
\pgfsetroundjoin%
\pgfsetlinewidth{0.903375pt}%
\definecolor{currentstroke}{rgb}{0.000000,0.000000,0.000000}%
\pgfsetstrokecolor{currentstroke}%
\pgfsetdash{{0.900000pt}{1.485000pt}}{0.000000pt}%
\pgfpathmoveto{\pgfqpoint{0.514740in}{1.399313in}}%
\pgfpathlineto{\pgfqpoint{0.551648in}{1.400647in}}%
\pgfpathlineto{\pgfqpoint{0.588556in}{1.404862in}}%
\pgfpathlineto{\pgfqpoint{0.613162in}{1.409586in}}%
\pgfpathlineto{\pgfqpoint{0.637767in}{1.416213in}}%
\pgfpathlineto{\pgfqpoint{0.662373in}{1.425214in}}%
\pgfpathlineto{\pgfqpoint{0.686978in}{1.437227in}}%
\pgfpathlineto{\pgfqpoint{0.711584in}{1.453100in}}%
\pgfpathlineto{\pgfqpoint{0.736189in}{1.473959in}}%
\pgfpathlineto{\pgfqpoint{0.748492in}{1.486703in}}%
\pgfpathlineto{\pgfqpoint{0.760795in}{1.501280in}}%
\pgfpathlineto{\pgfqpoint{0.773098in}{1.517948in}}%
\pgfpathlineto{\pgfqpoint{0.785401in}{1.536999in}}%
\pgfpathlineto{\pgfqpoint{0.797703in}{1.558770in}}%
\pgfpathlineto{\pgfqpoint{0.810006in}{1.583643in}}%
\pgfpathlineto{\pgfqpoint{0.822309in}{1.612059in}}%
\pgfpathlineto{\pgfqpoint{0.834612in}{1.644517in}}%
\pgfpathlineto{\pgfqpoint{0.846914in}{1.681590in}}%
\pgfpathlineto{\pgfqpoint{0.859217in}{1.723930in}}%
\pgfpathlineto{\pgfqpoint{0.871520in}{1.772285in}}%
\pgfpathlineto{\pgfqpoint{0.883823in}{1.827506in}}%
\pgfpathlineto{\pgfqpoint{0.883823in}{1.827506in}}%
\pgfpathlineto{\pgfqpoint{0.899588in}{1.810594in}}%
\pgfpathlineto{\pgfqpoint{0.915354in}{1.777800in}}%
\pgfpathlineto{\pgfqpoint{0.925864in}{1.747573in}}%
\pgfpathlineto{\pgfqpoint{0.931119in}{1.734678in}}%
\pgfpathlineto{\pgfqpoint{0.936375in}{1.711293in}}%
\pgfpathlineto{\pgfqpoint{0.941630in}{1.693718in}}%
\pgfpathlineto{\pgfqpoint{0.946885in}{1.681611in}}%
\pgfpathlineto{\pgfqpoint{0.957396in}{1.634298in}}%
\pgfpathlineto{\pgfqpoint{0.967906in}{1.603523in}}%
\pgfpathlineto{\pgfqpoint{0.978416in}{1.557000in}}%
\pgfpathlineto{\pgfqpoint{0.983672in}{1.543299in}}%
\pgfpathlineto{\pgfqpoint{0.988927in}{1.526344in}}%
\pgfpathlineto{\pgfqpoint{0.999437in}{1.487080in}}%
\pgfpathlineto{\pgfqpoint{1.004692in}{1.477493in}}%
\pgfpathlineto{\pgfqpoint{1.009948in}{1.465519in}}%
\pgfpathlineto{\pgfqpoint{1.015203in}{1.447163in}}%
\pgfpathlineto{\pgfqpoint{1.020458in}{1.437327in}}%
\pgfpathlineto{\pgfqpoint{1.025713in}{1.435827in}}%
\pgfpathlineto{\pgfqpoint{1.036224in}{1.415250in}}%
\pgfpathlineto{\pgfqpoint{1.041479in}{1.416091in}}%
\pgfpathlineto{\pgfqpoint{1.041479in}{1.416091in}}%
\pgfpathlineto{\pgfqpoint{1.044570in}{1.411413in}}%
\pgfpathlineto{\pgfqpoint{1.047662in}{1.417758in}}%
\pgfpathlineto{\pgfqpoint{1.050753in}{1.419911in}}%
\pgfpathlineto{\pgfqpoint{1.053844in}{1.414409in}}%
\pgfpathlineto{\pgfqpoint{1.056935in}{1.412535in}}%
\pgfpathlineto{\pgfqpoint{1.060027in}{1.417943in}}%
\pgfpathlineto{\pgfqpoint{1.063118in}{1.419375in}}%
\pgfpathlineto{\pgfqpoint{1.066209in}{1.412323in}}%
\pgfpathlineto{\pgfqpoint{1.069301in}{1.414264in}}%
\pgfpathlineto{\pgfqpoint{1.072392in}{1.419079in}}%
\pgfpathlineto{\pgfqpoint{1.075483in}{1.416917in}}%
\pgfpathlineto{\pgfqpoint{1.078575in}{1.412023in}}%
\pgfpathlineto{\pgfqpoint{1.084757in}{1.419803in}}%
\pgfpathlineto{\pgfqpoint{1.087849in}{1.414539in}}%
\pgfpathlineto{\pgfqpoint{1.090940in}{1.412643in}}%
\pgfpathlineto{\pgfqpoint{1.097122in}{1.418629in}}%
\pgfpathlineto{\pgfqpoint{1.100214in}{1.413467in}}%
\pgfpathlineto{\pgfqpoint{1.103305in}{1.412813in}}%
\pgfpathlineto{\pgfqpoint{1.106396in}{1.418680in}}%
\pgfpathlineto{\pgfqpoint{1.112579in}{1.412738in}}%
\pgfpathlineto{\pgfqpoint{1.115670in}{1.414589in}}%
\pgfpathlineto{\pgfqpoint{1.118762in}{1.418960in}}%
\pgfpathlineto{\pgfqpoint{1.124944in}{1.411409in}}%
\pgfpathlineto{\pgfqpoint{1.128035in}{1.416268in}}%
\pgfpathlineto{\pgfqpoint{1.131127in}{1.417752in}}%
\pgfpathlineto{\pgfqpoint{1.134218in}{1.413749in}}%
\pgfpathlineto{\pgfqpoint{1.134218in}{1.413749in}}%
\pgfpathlineto{\pgfqpoint{1.146521in}{1.410366in}}%
\pgfpathlineto{\pgfqpoint{1.171126in}{1.409686in}}%
\pgfpathlineto{\pgfqpoint{1.195732in}{1.413634in}}%
\pgfpathlineto{\pgfqpoint{1.208035in}{1.416730in}}%
\pgfpathlineto{\pgfqpoint{1.232640in}{1.418796in}}%
\pgfpathlineto{\pgfqpoint{1.257246in}{1.422893in}}%
\pgfpathlineto{\pgfqpoint{1.281851in}{1.429934in}}%
\pgfpathlineto{\pgfqpoint{1.306457in}{1.441634in}}%
\pgfpathlineto{\pgfqpoint{1.331062in}{1.458434in}}%
\pgfpathlineto{\pgfqpoint{1.355668in}{1.479080in}}%
\pgfpathlineto{\pgfqpoint{1.367970in}{1.491040in}}%
\pgfpathlineto{\pgfqpoint{1.392576in}{1.520970in}}%
\pgfpathlineto{\pgfqpoint{1.404879in}{1.538923in}}%
\pgfpathlineto{\pgfqpoint{1.417181in}{1.560444in}}%
\pgfpathlineto{\pgfqpoint{1.429484in}{1.584846in}}%
\pgfpathlineto{\pgfqpoint{1.441787in}{1.613126in}}%
\pgfpathlineto{\pgfqpoint{1.454090in}{1.645259in}}%
\pgfpathlineto{\pgfqpoint{1.466392in}{1.682357in}}%
\pgfpathlineto{\pgfqpoint{1.478695in}{1.723672in}}%
\pgfpathlineto{\pgfqpoint{1.490998in}{1.771304in}}%
\pgfpathlineto{\pgfqpoint{1.503301in}{1.826000in}}%
\pgfpathlineto{\pgfqpoint{1.503301in}{1.826000in}}%
\pgfpathlineto{\pgfqpoint{1.512575in}{1.816305in}}%
\pgfpathlineto{\pgfqpoint{1.517212in}{1.806271in}}%
\pgfpathlineto{\pgfqpoint{1.521849in}{1.800064in}}%
\pgfpathlineto{\pgfqpoint{1.526486in}{1.785109in}}%
\pgfpathlineto{\pgfqpoint{1.531122in}{1.775959in}}%
\pgfpathlineto{\pgfqpoint{1.545033in}{1.727503in}}%
\pgfpathlineto{\pgfqpoint{1.549670in}{1.709163in}}%
\pgfpathlineto{\pgfqpoint{1.554307in}{1.684888in}}%
\pgfpathlineto{\pgfqpoint{1.558944in}{1.671281in}}%
\pgfpathlineto{\pgfqpoint{1.568218in}{1.627030in}}%
\pgfpathlineto{\pgfqpoint{1.572855in}{1.611820in}}%
\pgfpathlineto{\pgfqpoint{1.582129in}{1.566866in}}%
\pgfpathlineto{\pgfqpoint{1.586766in}{1.553683in}}%
\pgfpathlineto{\pgfqpoint{1.591403in}{1.528822in}}%
\pgfpathlineto{\pgfqpoint{1.596040in}{1.513848in}}%
\pgfpathlineto{\pgfqpoint{1.600677in}{1.503579in}}%
\pgfpathlineto{\pgfqpoint{1.605314in}{1.483099in}}%
\pgfpathlineto{\pgfqpoint{1.609951in}{1.476140in}}%
\pgfpathlineto{\pgfqpoint{1.614588in}{1.465669in}}%
\pgfpathlineto{\pgfqpoint{1.619225in}{1.447933in}}%
\pgfpathlineto{\pgfqpoint{1.623862in}{1.447634in}}%
\pgfpathlineto{\pgfqpoint{1.628498in}{1.444828in}}%
\pgfpathlineto{\pgfqpoint{1.633135in}{1.435538in}}%
\pgfpathlineto{\pgfqpoint{1.637772in}{1.443146in}}%
\pgfpathlineto{\pgfqpoint{1.642409in}{1.439007in}}%
\pgfpathlineto{\pgfqpoint{1.673322in}{1.441819in}}%
\pgfpathlineto{\pgfqpoint{1.735148in}{1.432726in}}%
\pgfpathlineto{\pgfqpoint{1.766061in}{1.440022in}}%
\pgfpathlineto{\pgfqpoint{1.796974in}{1.437009in}}%
\pgfpathlineto{\pgfqpoint{1.827887in}{1.431187in}}%
\pgfpathlineto{\pgfqpoint{1.858801in}{1.432886in}}%
\pgfpathlineto{\pgfqpoint{1.889714in}{1.436587in}}%
\pgfpathlineto{\pgfqpoint{1.920627in}{1.433959in}}%
\pgfpathlineto{\pgfqpoint{1.951540in}{1.430033in}}%
\pgfpathlineto{\pgfqpoint{1.982453in}{1.429928in}}%
\pgfpathlineto{\pgfqpoint{2.013366in}{1.434085in}}%
\pgfpathlineto{\pgfqpoint{2.075192in}{1.427014in}}%
\pgfpathlineto{\pgfqpoint{2.106105in}{1.431145in}}%
\pgfpathlineto{\pgfqpoint{2.137018in}{1.430725in}}%
\pgfpathlineto{\pgfqpoint{2.167931in}{1.426532in}}%
\pgfpathlineto{\pgfqpoint{2.198844in}{1.426963in}}%
\pgfpathlineto{\pgfqpoint{2.229757in}{1.429606in}}%
\pgfpathlineto{\pgfqpoint{2.260670in}{1.428254in}}%
\pgfpathlineto{\pgfqpoint{2.291583in}{1.424549in}}%
\pgfpathlineto{\pgfqpoint{2.322496in}{1.424956in}}%
\pgfpathlineto{\pgfqpoint{2.353409in}{1.429462in}}%
\pgfpathlineto{\pgfqpoint{2.415235in}{1.422322in}}%
\pgfpathlineto{\pgfqpoint{2.446148in}{1.425489in}}%
\pgfpathlineto{\pgfqpoint{2.477061in}{1.426244in}}%
\pgfpathlineto{\pgfqpoint{2.507974in}{1.423419in}}%
\pgfpathlineto{\pgfqpoint{2.662539in}{1.422045in}}%
\pgfpathlineto{\pgfqpoint{2.693452in}{1.424721in}}%
\pgfpathlineto{\pgfqpoint{2.724365in}{1.421762in}}%
\pgfpathlineto{\pgfqpoint{2.755278in}{1.419947in}}%
\pgfpathlineto{\pgfqpoint{2.786191in}{1.421945in}}%
\pgfpathlineto{\pgfqpoint{2.817104in}{1.422528in}}%
\pgfpathlineto{\pgfqpoint{2.878930in}{1.420117in}}%
\pgfpathlineto{\pgfqpoint{2.940756in}{1.421535in}}%
\pgfpathlineto{\pgfqpoint{2.971669in}{1.419524in}}%
\pgfpathlineto{\pgfqpoint{3.002582in}{1.419362in}}%
\pgfpathlineto{\pgfqpoint{3.033495in}{1.421254in}}%
\pgfpathlineto{\pgfqpoint{3.095321in}{1.418222in}}%
\pgfpathlineto{\pgfqpoint{3.157147in}{1.420186in}}%
\pgfpathlineto{\pgfqpoint{3.218974in}{1.417717in}}%
\pgfpathlineto{\pgfqpoint{3.280800in}{1.419943in}}%
\pgfpathlineto{\pgfqpoint{3.311713in}{1.417949in}}%
\pgfpathlineto{\pgfqpoint{3.342626in}{1.417397in}}%
\pgfpathlineto{\pgfqpoint{3.373539in}{1.418724in}}%
\pgfpathlineto{\pgfqpoint{3.435365in}{1.417530in}}%
\pgfpathlineto{\pgfqpoint{3.501828in}{1.419566in}}%
\pgfpathlineto{\pgfqpoint{3.515739in}{1.434710in}}%
\pgfpathlineto{\pgfqpoint{3.525012in}{1.456677in}}%
\pgfpathlineto{\pgfqpoint{3.538923in}{1.497355in}}%
\pgfpathlineto{\pgfqpoint{3.543560in}{1.513115in}}%
\pgfpathlineto{\pgfqpoint{3.552834in}{1.549510in}}%
\pgfpathlineto{\pgfqpoint{3.562108in}{1.585989in}}%
\pgfpathlineto{\pgfqpoint{3.576019in}{1.642652in}}%
\pgfpathlineto{\pgfqpoint{3.585293in}{1.675171in}}%
\pgfpathlineto{\pgfqpoint{3.594567in}{1.706460in}}%
\pgfpathlineto{\pgfqpoint{3.608478in}{1.740348in}}%
\pgfpathlineto{\pgfqpoint{3.613115in}{1.749521in}}%
\pgfpathlineto{\pgfqpoint{3.622388in}{1.760112in}}%
\pgfpathlineto{\pgfqpoint{3.631662in}{1.764437in}}%
\pgfpathlineto{\pgfqpoint{3.636299in}{1.762339in}}%
\pgfpathlineto{\pgfqpoint{3.636299in}{1.762339in}}%
\pgfpathlineto{\pgfqpoint{3.648665in}{1.717672in}}%
\pgfpathlineto{\pgfqpoint{3.661030in}{1.677645in}}%
\pgfpathlineto{\pgfqpoint{3.673395in}{1.644055in}}%
\pgfpathlineto{\pgfqpoint{3.685760in}{1.614142in}}%
\pgfpathlineto{\pgfqpoint{3.710491in}{1.565056in}}%
\pgfpathlineto{\pgfqpoint{3.722856in}{1.544951in}}%
\pgfpathlineto{\pgfqpoint{3.735221in}{1.527326in}}%
\pgfpathlineto{\pgfqpoint{3.747586in}{1.512111in}}%
\pgfpathlineto{\pgfqpoint{3.772317in}{1.486530in}}%
\pgfpathlineto{\pgfqpoint{3.797047in}{1.467294in}}%
\pgfpathlineto{\pgfqpoint{3.834143in}{1.446535in}}%
\pgfpathlineto{\pgfqpoint{3.858873in}{1.436516in}}%
\pgfpathlineto{\pgfqpoint{3.908334in}{1.423016in}}%
\pgfpathlineto{\pgfqpoint{3.945430in}{1.416590in}}%
\pgfpathlineto{\pgfqpoint{4.007256in}{1.410528in}}%
\pgfpathlineto{\pgfqpoint{4.007256in}{1.410528in}}%
\pgfpathlineto{\pgfqpoint{4.017766in}{1.415279in}}%
\pgfpathlineto{\pgfqpoint{4.023021in}{1.418695in}}%
\pgfpathlineto{\pgfqpoint{4.033532in}{1.431490in}}%
\pgfpathlineto{\pgfqpoint{4.049297in}{1.460568in}}%
\pgfpathlineto{\pgfqpoint{4.059808in}{1.486877in}}%
\pgfpathlineto{\pgfqpoint{4.080829in}{1.545724in}}%
\pgfpathlineto{\pgfqpoint{4.117615in}{1.649435in}}%
\pgfpathlineto{\pgfqpoint{4.128126in}{1.676243in}}%
\pgfpathlineto{\pgfqpoint{4.143891in}{1.704380in}}%
\pgfpathlineto{\pgfqpoint{4.149146in}{1.712648in}}%
\pgfpathlineto{\pgfqpoint{4.154402in}{1.717395in}}%
\pgfpathlineto{\pgfqpoint{4.164912in}{1.724785in}}%
\pgfpathlineto{\pgfqpoint{4.164912in}{1.724785in}}%
\pgfpathlineto{\pgfqpoint{4.177277in}{1.683941in}}%
\pgfpathlineto{\pgfqpoint{4.189642in}{1.647992in}}%
\pgfpathlineto{\pgfqpoint{4.202008in}{1.617239in}}%
\pgfpathlineto{\pgfqpoint{4.214373in}{1.590637in}}%
\pgfpathlineto{\pgfqpoint{4.226738in}{1.567295in}}%
\pgfpathlineto{\pgfqpoint{4.239103in}{1.546535in}}%
\pgfpathlineto{\pgfqpoint{4.251469in}{1.528243in}}%
\pgfpathlineto{\pgfqpoint{4.263834in}{1.512673in}}%
\pgfpathlineto{\pgfqpoint{4.288564in}{1.487018in}}%
\pgfpathlineto{\pgfqpoint{4.313295in}{1.466412in}}%
\pgfpathlineto{\pgfqpoint{4.338025in}{1.451163in}}%
\pgfpathlineto{\pgfqpoint{4.362755in}{1.439058in}}%
\pgfpathlineto{\pgfqpoint{4.387486in}{1.430305in}}%
\pgfpathlineto{\pgfqpoint{4.424581in}{1.420732in}}%
\pgfpathlineto{\pgfqpoint{4.461677in}{1.414337in}}%
\pgfpathlineto{\pgfqpoint{4.498773in}{1.409789in}}%
\pgfpathlineto{\pgfqpoint{4.535868in}{1.406785in}}%
\pgfpathlineto{\pgfqpoint{4.535868in}{1.406785in}}%
\pgfusepath{stroke}%
\end{pgfscope}%
\begin{pgfscope}%
\pgfpathrectangle{\pgfqpoint{0.514740in}{1.372588in}}{\pgfqpoint{4.021128in}{1.113555in}}%
\pgfusepath{clip}%
\pgfsetrectcap%
\pgfsetroundjoin%
\pgfsetlinewidth{0.903375pt}%
\definecolor{currentstroke}{rgb}{0.121569,0.466667,0.705882}%
\pgfsetstrokecolor{currentstroke}%
\pgfsetdash{}{0pt}%
\pgfpathmoveto{\pgfqpoint{0.514740in}{1.399313in}}%
\pgfpathlineto{\pgfqpoint{0.883823in}{1.399313in}}%
\pgfpathlineto{\pgfqpoint{0.883823in}{1.399313in}}%
\pgfpathlineto{\pgfqpoint{0.889078in}{1.401065in}}%
\pgfpathlineto{\pgfqpoint{0.894333in}{1.407750in}}%
\pgfpathlineto{\pgfqpoint{0.904843in}{1.429930in}}%
\pgfpathlineto{\pgfqpoint{0.910099in}{1.445865in}}%
\pgfpathlineto{\pgfqpoint{0.931119in}{1.544478in}}%
\pgfpathlineto{\pgfqpoint{0.941630in}{1.614289in}}%
\pgfpathlineto{\pgfqpoint{0.946885in}{1.642364in}}%
\pgfpathlineto{\pgfqpoint{0.983672in}{1.901193in}}%
\pgfpathlineto{\pgfqpoint{0.988927in}{1.932973in}}%
\pgfpathlineto{\pgfqpoint{0.994182in}{1.975386in}}%
\pgfpathlineto{\pgfqpoint{0.999437in}{2.006077in}}%
\pgfpathlineto{\pgfqpoint{1.004692in}{2.025593in}}%
\pgfpathlineto{\pgfqpoint{1.015203in}{2.085684in}}%
\pgfpathlineto{\pgfqpoint{1.020458in}{2.102703in}}%
\pgfpathlineto{\pgfqpoint{1.025713in}{2.110745in}}%
\pgfpathlineto{\pgfqpoint{1.036224in}{2.149730in}}%
\pgfpathlineto{\pgfqpoint{1.041479in}{2.148108in}}%
\pgfpathlineto{\pgfqpoint{1.041479in}{2.148108in}}%
\pgfpathlineto{\pgfqpoint{1.044570in}{2.155281in}}%
\pgfpathlineto{\pgfqpoint{1.047662in}{2.145237in}}%
\pgfpathlineto{\pgfqpoint{1.050753in}{2.142267in}}%
\pgfpathlineto{\pgfqpoint{1.053844in}{2.151724in}}%
\pgfpathlineto{\pgfqpoint{1.056935in}{2.155310in}}%
\pgfpathlineto{\pgfqpoint{1.060027in}{2.143560in}}%
\pgfpathlineto{\pgfqpoint{1.063118in}{2.141740in}}%
\pgfpathlineto{\pgfqpoint{1.066209in}{2.152649in}}%
\pgfpathlineto{\pgfqpoint{1.069301in}{2.150138in}}%
\pgfpathlineto{\pgfqpoint{1.072392in}{2.143941in}}%
\pgfpathlineto{\pgfqpoint{1.075483in}{2.148108in}}%
\pgfpathlineto{\pgfqpoint{1.078575in}{2.155199in}}%
\pgfpathlineto{\pgfqpoint{1.084757in}{2.138414in}}%
\pgfpathlineto{\pgfqpoint{1.087849in}{2.149664in}}%
\pgfpathlineto{\pgfqpoint{1.090940in}{2.154906in}}%
\pgfpathlineto{\pgfqpoint{1.094031in}{2.147171in}}%
\pgfpathlineto{\pgfqpoint{1.097122in}{2.143166in}}%
\pgfpathlineto{\pgfqpoint{1.100214in}{2.151032in}}%
\pgfpathlineto{\pgfqpoint{1.103305in}{2.150770in}}%
\pgfpathlineto{\pgfqpoint{1.106396in}{2.140891in}}%
\pgfpathlineto{\pgfqpoint{1.109488in}{2.146277in}}%
\pgfpathlineto{\pgfqpoint{1.112579in}{2.154613in}}%
\pgfpathlineto{\pgfqpoint{1.115670in}{2.149855in}}%
\pgfpathlineto{\pgfqpoint{1.118762in}{2.140927in}}%
\pgfpathlineto{\pgfqpoint{1.124944in}{2.154689in}}%
\pgfpathlineto{\pgfqpoint{1.128035in}{2.145262in}}%
\pgfpathlineto{\pgfqpoint{1.131127in}{2.142179in}}%
\pgfpathlineto{\pgfqpoint{1.134218in}{2.150041in}}%
\pgfpathlineto{\pgfqpoint{1.134218in}{2.150041in}}%
\pgfpathlineto{\pgfqpoint{1.146521in}{2.153431in}}%
\pgfpathlineto{\pgfqpoint{1.158824in}{2.153352in}}%
\pgfpathlineto{\pgfqpoint{1.171126in}{2.147649in}}%
\pgfpathlineto{\pgfqpoint{1.183429in}{2.149016in}}%
\pgfpathlineto{\pgfqpoint{1.195732in}{2.144354in}}%
\pgfpathlineto{\pgfqpoint{1.208035in}{2.140958in}}%
\pgfpathlineto{\pgfqpoint{1.232640in}{2.141506in}}%
\pgfpathlineto{\pgfqpoint{1.244943in}{2.142676in}}%
\pgfpathlineto{\pgfqpoint{1.257246in}{2.141900in}}%
\pgfpathlineto{\pgfqpoint{1.269548in}{2.146441in}}%
\pgfpathlineto{\pgfqpoint{1.281851in}{2.144725in}}%
\pgfpathlineto{\pgfqpoint{1.306457in}{2.143222in}}%
\pgfpathlineto{\pgfqpoint{1.331062in}{2.140501in}}%
\pgfpathlineto{\pgfqpoint{1.343365in}{2.137174in}}%
\pgfpathlineto{\pgfqpoint{1.355668in}{2.140021in}}%
\pgfpathlineto{\pgfqpoint{1.367970in}{2.137828in}}%
\pgfpathlineto{\pgfqpoint{1.417181in}{2.139076in}}%
\pgfpathlineto{\pgfqpoint{1.429484in}{2.137107in}}%
\pgfpathlineto{\pgfqpoint{1.441787in}{2.139770in}}%
\pgfpathlineto{\pgfqpoint{1.454090in}{2.135727in}}%
\pgfpathlineto{\pgfqpoint{1.503301in}{2.133953in}}%
\pgfpathlineto{\pgfqpoint{1.503301in}{2.133953in}}%
\pgfpathlineto{\pgfqpoint{1.507938in}{2.133305in}}%
\pgfpathlineto{\pgfqpoint{1.512575in}{2.128367in}}%
\pgfpathlineto{\pgfqpoint{1.517212in}{2.126861in}}%
\pgfpathlineto{\pgfqpoint{1.521849in}{2.132667in}}%
\pgfpathlineto{\pgfqpoint{1.526486in}{2.136289in}}%
\pgfpathlineto{\pgfqpoint{1.535759in}{2.127627in}}%
\pgfpathlineto{\pgfqpoint{1.540396in}{2.127755in}}%
\pgfpathlineto{\pgfqpoint{1.549670in}{2.134095in}}%
\pgfpathlineto{\pgfqpoint{1.554307in}{2.132589in}}%
\pgfpathlineto{\pgfqpoint{1.558944in}{2.128564in}}%
\pgfpathlineto{\pgfqpoint{1.563581in}{2.128157in}}%
\pgfpathlineto{\pgfqpoint{1.568218in}{2.132626in}}%
\pgfpathlineto{\pgfqpoint{1.572855in}{2.134150in}}%
\pgfpathlineto{\pgfqpoint{1.577492in}{2.129598in}}%
\pgfpathlineto{\pgfqpoint{1.582129in}{2.127318in}}%
\pgfpathlineto{\pgfqpoint{1.596040in}{2.134145in}}%
\pgfpathlineto{\pgfqpoint{1.600677in}{2.131297in}}%
\pgfpathlineto{\pgfqpoint{1.605314in}{2.126555in}}%
\pgfpathlineto{\pgfqpoint{1.609951in}{2.127769in}}%
\pgfpathlineto{\pgfqpoint{1.614588in}{2.132684in}}%
\pgfpathlineto{\pgfqpoint{1.619225in}{2.133622in}}%
\pgfpathlineto{\pgfqpoint{1.623862in}{2.129912in}}%
\pgfpathlineto{\pgfqpoint{1.628498in}{2.127854in}}%
\pgfpathlineto{\pgfqpoint{1.642409in}{2.132469in}}%
\pgfpathlineto{\pgfqpoint{1.642409in}{2.132469in}}%
\pgfpathlineto{\pgfqpoint{1.673322in}{2.128909in}}%
\pgfpathlineto{\pgfqpoint{1.704235in}{2.133697in}}%
\pgfpathlineto{\pgfqpoint{1.735148in}{2.130125in}}%
\pgfpathlineto{\pgfqpoint{1.766061in}{2.132879in}}%
\pgfpathlineto{\pgfqpoint{1.796974in}{2.130161in}}%
\pgfpathlineto{\pgfqpoint{1.827887in}{2.129489in}}%
\pgfpathlineto{\pgfqpoint{1.858801in}{2.132023in}}%
\pgfpathlineto{\pgfqpoint{1.889714in}{2.129705in}}%
\pgfpathlineto{\pgfqpoint{1.920627in}{2.132319in}}%
\pgfpathlineto{\pgfqpoint{1.951540in}{2.127137in}}%
\pgfpathlineto{\pgfqpoint{1.982453in}{2.130672in}}%
\pgfpathlineto{\pgfqpoint{2.013366in}{2.130625in}}%
\pgfpathlineto{\pgfqpoint{2.044279in}{2.129053in}}%
\pgfpathlineto{\pgfqpoint{2.075192in}{2.129967in}}%
\pgfpathlineto{\pgfqpoint{2.106105in}{2.126843in}}%
\pgfpathlineto{\pgfqpoint{2.137018in}{2.131244in}}%
\pgfpathlineto{\pgfqpoint{2.167931in}{2.127641in}}%
\pgfpathlineto{\pgfqpoint{2.198844in}{2.129407in}}%
\pgfpathlineto{\pgfqpoint{2.260670in}{2.126797in}}%
\pgfpathlineto{\pgfqpoint{2.291583in}{2.130472in}}%
\pgfpathlineto{\pgfqpoint{2.322496in}{2.126164in}}%
\pgfpathlineto{\pgfqpoint{2.353409in}{2.128669in}}%
\pgfpathlineto{\pgfqpoint{2.384322in}{2.125737in}}%
\pgfpathlineto{\pgfqpoint{2.415235in}{2.128270in}}%
\pgfpathlineto{\pgfqpoint{2.446148in}{2.127851in}}%
\pgfpathlineto{\pgfqpoint{2.477061in}{2.125127in}}%
\pgfpathlineto{\pgfqpoint{2.507974in}{2.128109in}}%
\pgfpathlineto{\pgfqpoint{2.538887in}{2.124643in}}%
\pgfpathlineto{\pgfqpoint{2.569800in}{2.128278in}}%
\pgfpathlineto{\pgfqpoint{2.600713in}{2.125080in}}%
\pgfpathlineto{\pgfqpoint{2.662539in}{2.126152in}}%
\pgfpathlineto{\pgfqpoint{2.693452in}{2.124441in}}%
\pgfpathlineto{\pgfqpoint{2.724365in}{2.127601in}}%
\pgfpathlineto{\pgfqpoint{2.755278in}{2.122595in}}%
\pgfpathlineto{\pgfqpoint{2.786191in}{2.126019in}}%
\pgfpathlineto{\pgfqpoint{2.817104in}{2.124448in}}%
\pgfpathlineto{\pgfqpoint{2.878930in}{2.125055in}}%
\pgfpathlineto{\pgfqpoint{2.909843in}{2.122033in}}%
\pgfpathlineto{\pgfqpoint{2.940756in}{2.125990in}}%
\pgfpathlineto{\pgfqpoint{2.971669in}{2.122321in}}%
\pgfpathlineto{\pgfqpoint{3.002582in}{2.124908in}}%
\pgfpathlineto{\pgfqpoint{3.033495in}{2.122478in}}%
\pgfpathlineto{\pgfqpoint{3.064408in}{2.122352in}}%
\pgfpathlineto{\pgfqpoint{3.095321in}{2.124652in}}%
\pgfpathlineto{\pgfqpoint{3.126234in}{2.121365in}}%
\pgfpathlineto{\pgfqpoint{3.157147in}{2.124164in}}%
\pgfpathlineto{\pgfqpoint{3.188061in}{2.120167in}}%
\pgfpathlineto{\pgfqpoint{3.218974in}{2.123449in}}%
\pgfpathlineto{\pgfqpoint{3.280800in}{2.120862in}}%
\pgfpathlineto{\pgfqpoint{3.311713in}{2.122694in}}%
\pgfpathlineto{\pgfqpoint{3.342626in}{2.119345in}}%
\pgfpathlineto{\pgfqpoint{3.373539in}{2.123444in}}%
\pgfpathlineto{\pgfqpoint{3.404452in}{2.119825in}}%
\pgfpathlineto{\pgfqpoint{3.435365in}{2.121219in}}%
\pgfpathlineto{\pgfqpoint{3.501828in}{2.119624in}}%
\pgfpathlineto{\pgfqpoint{3.511102in}{2.116078in}}%
\pgfpathlineto{\pgfqpoint{3.515739in}{2.117026in}}%
\pgfpathlineto{\pgfqpoint{3.520376in}{2.119337in}}%
\pgfpathlineto{\pgfqpoint{3.525012in}{2.119104in}}%
\pgfpathlineto{\pgfqpoint{3.534286in}{2.115861in}}%
\pgfpathlineto{\pgfqpoint{3.543560in}{2.118942in}}%
\pgfpathlineto{\pgfqpoint{3.548197in}{2.118698in}}%
\pgfpathlineto{\pgfqpoint{3.557471in}{2.116108in}}%
\pgfpathlineto{\pgfqpoint{3.562108in}{2.116529in}}%
\pgfpathlineto{\pgfqpoint{3.566745in}{2.118595in}}%
\pgfpathlineto{\pgfqpoint{3.571382in}{2.118883in}}%
\pgfpathlineto{\pgfqpoint{3.580656in}{2.115821in}}%
\pgfpathlineto{\pgfqpoint{3.594567in}{2.118636in}}%
\pgfpathlineto{\pgfqpoint{3.603841in}{2.115212in}}%
\pgfpathlineto{\pgfqpoint{3.608478in}{2.116409in}}%
\pgfpathlineto{\pgfqpoint{3.613115in}{2.119037in}}%
\pgfpathlineto{\pgfqpoint{3.617752in}{2.118935in}}%
\pgfpathlineto{\pgfqpoint{3.622388in}{2.115894in}}%
\pgfpathlineto{\pgfqpoint{3.627025in}{2.114556in}}%
\pgfpathlineto{\pgfqpoint{3.636299in}{2.119062in}}%
\pgfpathlineto{\pgfqpoint{3.636299in}{2.119062in}}%
\pgfpathlineto{\pgfqpoint{3.661030in}{2.116470in}}%
\pgfpathlineto{\pgfqpoint{3.673395in}{2.118146in}}%
\pgfpathlineto{\pgfqpoint{3.685760in}{2.114659in}}%
\pgfpathlineto{\pgfqpoint{3.698125in}{2.117874in}}%
\pgfpathlineto{\pgfqpoint{3.710491in}{2.114475in}}%
\pgfpathlineto{\pgfqpoint{3.722856in}{2.116201in}}%
\pgfpathlineto{\pgfqpoint{3.747586in}{2.114410in}}%
\pgfpathlineto{\pgfqpoint{3.759951in}{2.115835in}}%
\pgfpathlineto{\pgfqpoint{3.772317in}{2.112757in}}%
\pgfpathlineto{\pgfqpoint{3.784682in}{2.115603in}}%
\pgfpathlineto{\pgfqpoint{3.797047in}{2.111808in}}%
\pgfpathlineto{\pgfqpoint{3.809412in}{2.113551in}}%
\pgfpathlineto{\pgfqpoint{3.834143in}{2.111247in}}%
\pgfpathlineto{\pgfqpoint{3.846508in}{2.112496in}}%
\pgfpathlineto{\pgfqpoint{3.858873in}{2.109745in}}%
\pgfpathlineto{\pgfqpoint{3.871238in}{2.112447in}}%
\pgfpathlineto{\pgfqpoint{3.883604in}{2.108954in}}%
\pgfpathlineto{\pgfqpoint{3.895969in}{2.111011in}}%
\pgfpathlineto{\pgfqpoint{3.920699in}{2.108793in}}%
\pgfpathlineto{\pgfqpoint{3.933064in}{2.109976in}}%
\pgfpathlineto{\pgfqpoint{3.945430in}{2.107201in}}%
\pgfpathlineto{\pgfqpoint{3.957795in}{2.109682in}}%
\pgfpathlineto{\pgfqpoint{3.970160in}{2.106154in}}%
\pgfpathlineto{\pgfqpoint{3.982525in}{2.108195in}}%
\pgfpathlineto{\pgfqpoint{4.012511in}{2.103449in}}%
\pgfpathlineto{\pgfqpoint{4.017766in}{2.099702in}}%
\pgfpathlineto{\pgfqpoint{4.023021in}{2.092047in}}%
\pgfpathlineto{\pgfqpoint{4.033532in}{2.062477in}}%
\pgfpathlineto{\pgfqpoint{4.044042in}{2.026854in}}%
\pgfpathlineto{\pgfqpoint{4.059808in}{1.949311in}}%
\pgfpathlineto{\pgfqpoint{4.070318in}{1.889570in}}%
\pgfpathlineto{\pgfqpoint{4.117615in}{1.594817in}}%
\pgfpathlineto{\pgfqpoint{4.128126in}{1.539866in}}%
\pgfpathlineto{\pgfqpoint{4.138636in}{1.491037in}}%
\pgfpathlineto{\pgfqpoint{4.149146in}{1.453091in}}%
\pgfpathlineto{\pgfqpoint{4.159657in}{1.425545in}}%
\pgfpathlineto{\pgfqpoint{4.164912in}{1.416200in}}%
\pgfpathlineto{\pgfqpoint{4.164912in}{1.416200in}}%
\pgfpathlineto{\pgfqpoint{4.177277in}{1.417231in}}%
\pgfpathlineto{\pgfqpoint{4.300929in}{1.416687in}}%
\pgfpathlineto{\pgfqpoint{4.535868in}{1.417283in}}%
\pgfpathlineto{\pgfqpoint{4.535868in}{1.417283in}}%
\pgfusepath{stroke}%
\end{pgfscope}%
\begin{pgfscope}%
\pgfpathrectangle{\pgfqpoint{0.514740in}{1.372588in}}{\pgfqpoint{4.021128in}{1.113555in}}%
\pgfusepath{clip}%
\pgfsetbuttcap%
\pgfsetroundjoin%
\pgfsetlinewidth{0.903375pt}%
\definecolor{currentstroke}{rgb}{1.000000,0.000000,0.000000}%
\pgfsetstrokecolor{currentstroke}%
\pgfsetdash{{5.760000pt}{1.440000pt}{0.900000pt}{1.440000pt}}{0.000000pt}%
\pgfpathmoveto{\pgfqpoint{0.514740in}{1.399313in}}%
\pgfpathlineto{\pgfqpoint{0.899588in}{1.400020in}}%
\pgfpathlineto{\pgfqpoint{0.904843in}{1.400517in}}%
\pgfpathlineto{\pgfqpoint{0.910099in}{1.399838in}}%
\pgfpathlineto{\pgfqpoint{0.920609in}{1.400705in}}%
\pgfpathlineto{\pgfqpoint{0.931119in}{1.400867in}}%
\pgfpathlineto{\pgfqpoint{0.941630in}{1.401051in}}%
\pgfpathlineto{\pgfqpoint{1.020458in}{1.401535in}}%
\pgfpathlineto{\pgfqpoint{1.025713in}{1.402415in}}%
\pgfpathlineto{\pgfqpoint{1.036224in}{1.402420in}}%
\pgfpathlineto{\pgfqpoint{1.047662in}{1.402997in}}%
\pgfpathlineto{\pgfqpoint{1.050753in}{1.403428in}}%
\pgfpathlineto{\pgfqpoint{1.056935in}{1.402165in}}%
\pgfpathlineto{\pgfqpoint{1.063118in}{1.403473in}}%
\pgfpathlineto{\pgfqpoint{1.069301in}{1.402136in}}%
\pgfpathlineto{\pgfqpoint{1.075483in}{1.403570in}}%
\pgfpathlineto{\pgfqpoint{1.081666in}{1.402240in}}%
\pgfpathlineto{\pgfqpoint{1.087849in}{1.403587in}}%
\pgfpathlineto{\pgfqpoint{1.097122in}{1.402758in}}%
\pgfpathlineto{\pgfqpoint{1.103305in}{1.403310in}}%
\pgfpathlineto{\pgfqpoint{1.106396in}{1.402422in}}%
\pgfpathlineto{\pgfqpoint{1.115670in}{1.403391in}}%
\pgfpathlineto{\pgfqpoint{1.121853in}{1.402730in}}%
\pgfpathlineto{\pgfqpoint{1.128035in}{1.403496in}}%
\pgfpathlineto{\pgfqpoint{1.134218in}{1.402714in}}%
\pgfpathlineto{\pgfqpoint{1.134218in}{1.402714in}}%
\pgfpathlineto{\pgfqpoint{1.257246in}{1.403775in}}%
\pgfpathlineto{\pgfqpoint{1.503301in}{1.404962in}}%
\pgfpathlineto{\pgfqpoint{1.503301in}{1.404962in}}%
\pgfpathlineto{\pgfqpoint{1.507938in}{1.407405in}}%
\pgfpathlineto{\pgfqpoint{1.517212in}{1.425189in}}%
\pgfpathlineto{\pgfqpoint{1.521849in}{1.439634in}}%
\pgfpathlineto{\pgfqpoint{1.526486in}{1.464204in}}%
\pgfpathlineto{\pgfqpoint{1.531122in}{1.482772in}}%
\pgfpathlineto{\pgfqpoint{1.545033in}{1.570393in}}%
\pgfpathlineto{\pgfqpoint{1.572855in}{1.799464in}}%
\pgfpathlineto{\pgfqpoint{1.582129in}{1.885804in}}%
\pgfpathlineto{\pgfqpoint{1.586766in}{1.911413in}}%
\pgfpathlineto{\pgfqpoint{1.591403in}{1.963824in}}%
\pgfpathlineto{\pgfqpoint{1.609951in}{2.074032in}}%
\pgfpathlineto{\pgfqpoint{1.614588in}{2.093562in}}%
\pgfpathlineto{\pgfqpoint{1.619225in}{2.125925in}}%
\pgfpathlineto{\pgfqpoint{1.623862in}{2.126194in}}%
\pgfpathlineto{\pgfqpoint{1.628498in}{2.135156in}}%
\pgfpathlineto{\pgfqpoint{1.633135in}{2.155984in}}%
\pgfpathlineto{\pgfqpoint{1.637772in}{2.144129in}}%
\pgfpathlineto{\pgfqpoint{1.642409in}{2.147230in}}%
\pgfpathlineto{\pgfqpoint{1.642409in}{2.147230in}}%
\pgfpathlineto{\pgfqpoint{1.673322in}{2.139002in}}%
\pgfpathlineto{\pgfqpoint{1.704235in}{2.150758in}}%
\pgfpathlineto{\pgfqpoint{1.735148in}{2.150321in}}%
\pgfpathlineto{\pgfqpoint{1.766061in}{2.143942in}}%
\pgfpathlineto{\pgfqpoint{1.796974in}{2.143344in}}%
\pgfpathlineto{\pgfqpoint{1.827887in}{2.154689in}}%
\pgfpathlineto{\pgfqpoint{1.858801in}{2.152657in}}%
\pgfpathlineto{\pgfqpoint{1.889714in}{2.140206in}}%
\pgfpathlineto{\pgfqpoint{1.951540in}{2.154330in}}%
\pgfpathlineto{\pgfqpoint{1.982453in}{2.151956in}}%
\pgfpathlineto{\pgfqpoint{2.013366in}{2.144028in}}%
\pgfpathlineto{\pgfqpoint{2.075192in}{2.155404in}}%
\pgfpathlineto{\pgfqpoint{2.106105in}{2.147289in}}%
\pgfpathlineto{\pgfqpoint{2.137018in}{2.148786in}}%
\pgfpathlineto{\pgfqpoint{2.167931in}{2.152672in}}%
\pgfpathlineto{\pgfqpoint{2.198844in}{2.154144in}}%
\pgfpathlineto{\pgfqpoint{2.229757in}{2.146986in}}%
\pgfpathlineto{\pgfqpoint{2.260670in}{2.149305in}}%
\pgfpathlineto{\pgfqpoint{2.291583in}{2.157414in}}%
\pgfpathlineto{\pgfqpoint{2.322496in}{2.150946in}}%
\pgfpathlineto{\pgfqpoint{2.353409in}{2.148507in}}%
\pgfpathlineto{\pgfqpoint{2.384322in}{2.150750in}}%
\pgfpathlineto{\pgfqpoint{2.415235in}{2.157170in}}%
\pgfpathlineto{\pgfqpoint{2.477061in}{2.148214in}}%
\pgfpathlineto{\pgfqpoint{2.507974in}{2.154539in}}%
\pgfpathlineto{\pgfqpoint{2.600713in}{2.151317in}}%
\pgfpathlineto{\pgfqpoint{2.631626in}{2.156031in}}%
\pgfpathlineto{\pgfqpoint{2.662539in}{2.153777in}}%
\pgfpathlineto{\pgfqpoint{2.693452in}{2.148913in}}%
\pgfpathlineto{\pgfqpoint{2.724365in}{2.154798in}}%
\pgfpathlineto{\pgfqpoint{2.755278in}{2.155764in}}%
\pgfpathlineto{\pgfqpoint{2.786191in}{2.154895in}}%
\pgfpathlineto{\pgfqpoint{2.817104in}{2.149927in}}%
\pgfpathlineto{\pgfqpoint{2.878930in}{2.156791in}}%
\pgfpathlineto{\pgfqpoint{2.909843in}{2.152219in}}%
\pgfpathlineto{\pgfqpoint{2.940756in}{2.154219in}}%
\pgfpathlineto{\pgfqpoint{2.971669in}{2.153930in}}%
\pgfpathlineto{\pgfqpoint{3.002582in}{2.155783in}}%
\pgfpathlineto{\pgfqpoint{3.033495in}{2.152178in}}%
\pgfpathlineto{\pgfqpoint{3.064408in}{2.153960in}}%
\pgfpathlineto{\pgfqpoint{3.095321in}{2.157294in}}%
\pgfpathlineto{\pgfqpoint{3.126234in}{2.152862in}}%
\pgfpathlineto{\pgfqpoint{3.188061in}{2.153821in}}%
\pgfpathlineto{\pgfqpoint{3.218974in}{2.157557in}}%
\pgfpathlineto{\pgfqpoint{3.249887in}{2.153827in}}%
\pgfpathlineto{\pgfqpoint{3.280800in}{2.152340in}}%
\pgfpathlineto{\pgfqpoint{3.311713in}{2.156306in}}%
\pgfpathlineto{\pgfqpoint{3.342626in}{2.154670in}}%
\pgfpathlineto{\pgfqpoint{3.373539in}{2.155773in}}%
\pgfpathlineto{\pgfqpoint{3.404452in}{2.153152in}}%
\pgfpathlineto{\pgfqpoint{3.435365in}{2.156114in}}%
\pgfpathlineto{\pgfqpoint{3.497191in}{2.153117in}}%
\pgfpathlineto{\pgfqpoint{3.497191in}{2.153117in}}%
\pgfpathlineto{\pgfqpoint{3.501828in}{2.155007in}}%
\pgfpathlineto{\pgfqpoint{3.506465in}{2.147814in}}%
\pgfpathlineto{\pgfqpoint{3.520376in}{2.101138in}}%
\pgfpathlineto{\pgfqpoint{3.529649in}{2.054875in}}%
\pgfpathlineto{\pgfqpoint{3.548197in}{1.924313in}}%
\pgfpathlineto{\pgfqpoint{3.594567in}{1.554564in}}%
\pgfpathlineto{\pgfqpoint{3.603841in}{1.497675in}}%
\pgfpathlineto{\pgfqpoint{3.613115in}{1.455572in}}%
\pgfpathlineto{\pgfqpoint{3.617752in}{1.438838in}}%
\pgfpathlineto{\pgfqpoint{3.622388in}{1.426367in}}%
\pgfpathlineto{\pgfqpoint{3.627025in}{1.417345in}}%
\pgfpathlineto{\pgfqpoint{3.631662in}{1.412785in}}%
\pgfpathlineto{\pgfqpoint{3.636299in}{1.411801in}}%
\pgfpathlineto{\pgfqpoint{3.636299in}{1.411801in}}%
\pgfpathlineto{\pgfqpoint{3.759951in}{1.415407in}}%
\pgfpathlineto{\pgfqpoint{3.895969in}{1.418257in}}%
\pgfpathlineto{\pgfqpoint{4.101849in}{1.422798in}}%
\pgfpathlineto{\pgfqpoint{4.107105in}{1.423987in}}%
\pgfpathlineto{\pgfqpoint{4.128126in}{1.423696in}}%
\pgfpathlineto{\pgfqpoint{4.138636in}{1.424525in}}%
\pgfpathlineto{\pgfqpoint{4.143891in}{1.423819in}}%
\pgfpathlineto{\pgfqpoint{4.149146in}{1.424966in}}%
\pgfpathlineto{\pgfqpoint{4.164912in}{1.425070in}}%
\pgfpathlineto{\pgfqpoint{4.164912in}{1.425070in}}%
\pgfpathlineto{\pgfqpoint{4.276199in}{1.427261in}}%
\pgfpathlineto{\pgfqpoint{4.375121in}{1.428305in}}%
\pgfpathlineto{\pgfqpoint{4.535868in}{1.430298in}}%
\pgfpathlineto{\pgfqpoint{4.535868in}{1.430298in}}%
\pgfusepath{stroke}%
\end{pgfscope}%
\begin{pgfscope}%
\pgfsetrectcap%
\pgfsetmiterjoin%
\pgfsetlinewidth{0.803000pt}%
\definecolor{currentstroke}{rgb}{0.000000,0.000000,0.000000}%
\pgfsetstrokecolor{currentstroke}%
\pgfsetdash{}{0pt}%
\pgfpathmoveto{\pgfqpoint{0.514740in}{1.372588in}}%
\pgfpathlineto{\pgfqpoint{0.514740in}{2.486143in}}%
\pgfusepath{stroke}%
\end{pgfscope}%
\begin{pgfscope}%
\pgfsetrectcap%
\pgfsetmiterjoin%
\pgfsetlinewidth{0.803000pt}%
\definecolor{currentstroke}{rgb}{0.000000,0.000000,0.000000}%
\pgfsetstrokecolor{currentstroke}%
\pgfsetdash{}{0pt}%
\pgfpathmoveto{\pgfqpoint{4.535868in}{1.372588in}}%
\pgfpathlineto{\pgfqpoint{4.535868in}{2.486143in}}%
\pgfusepath{stroke}%
\end{pgfscope}%
\begin{pgfscope}%
\pgfsetrectcap%
\pgfsetmiterjoin%
\pgfsetlinewidth{0.803000pt}%
\definecolor{currentstroke}{rgb}{0.000000,0.000000,0.000000}%
\pgfsetstrokecolor{currentstroke}%
\pgfsetdash{}{0pt}%
\pgfpathmoveto{\pgfqpoint{0.514740in}{1.372588in}}%
\pgfpathlineto{\pgfqpoint{4.535868in}{1.372588in}}%
\pgfusepath{stroke}%
\end{pgfscope}%
\begin{pgfscope}%
\pgfsetrectcap%
\pgfsetmiterjoin%
\pgfsetlinewidth{0.803000pt}%
\definecolor{currentstroke}{rgb}{0.000000,0.000000,0.000000}%
\pgfsetstrokecolor{currentstroke}%
\pgfsetdash{}{0pt}%
\pgfpathmoveto{\pgfqpoint{0.514740in}{2.486143in}}%
\pgfpathlineto{\pgfqpoint{4.535868in}{2.486143in}}%
\pgfusepath{stroke}%
\end{pgfscope}%
\begin{pgfscope}%
\definecolor{textcolor}{rgb}{0.000000,0.000000,0.000000}%
\pgfsetstrokecolor{textcolor}%
\pgfsetfillcolor{textcolor}%
\pgftext[x=0.699281in,y=2.530686in,,bottom]{\color{textcolor}{\rmfamily\fontsize{8.000000}{9.600000}\selectfont\catcode`\^=\active\def^{\ifmmode\sp\else\^{}\fi}\catcode`\%=\active\def%{\%}L1}}%
\end{pgfscope}%
\begin{pgfscope}%
\definecolor{textcolor}{rgb}{0.000000,0.000000,0.000000}%
\pgfsetstrokecolor{textcolor}%
\pgfsetfillcolor{textcolor}%
\pgftext[x=0.962651in,y=2.530686in,,bottom]{\color{textcolor}{\rmfamily\fontsize{8.000000}{9.600000}\selectfont\catcode`\^=\active\def^{\ifmmode\sp\else\^{}\fi}\catcode`\%=\active\def%{\%}W1}}%
\end{pgfscope}%
\begin{pgfscope}%
\definecolor{textcolor}{rgb}{0.000000,0.000000,0.000000}%
\pgfsetstrokecolor{textcolor}%
\pgfsetfillcolor{textcolor}%
\pgftext[x=1.318759in,y=2.530686in,,bottom]{\color{textcolor}{\rmfamily\fontsize{8.000000}{9.600000}\selectfont\catcode`\^=\active\def^{\ifmmode\sp\else\^{}\fi}\catcode`\%=\active\def%{\%}L2}}%
\end{pgfscope}%
\begin{pgfscope}%
\definecolor{textcolor}{rgb}{0.000000,0.000000,0.000000}%
\pgfsetstrokecolor{textcolor}%
\pgfsetfillcolor{textcolor}%
\pgftext[x=1.572855in,y=2.530686in,,bottom]{\color{textcolor}{\rmfamily\fontsize{8.000000}{9.600000}\selectfont\catcode`\^=\active\def^{\ifmmode\sp\else\^{}\fi}\catcode`\%=\active\def%{\%}W2}}%
\end{pgfscope}%
\begin{pgfscope}%
\definecolor{textcolor}{rgb}{0.000000,0.000000,0.000000}%
\pgfsetstrokecolor{textcolor}%
\pgfsetfillcolor{textcolor}%
\pgftext[x=2.569800in,y=2.530686in,,bottom]{\color{textcolor}{\rmfamily\fontsize{8.000000}{9.600000}\selectfont\catcode`\^=\active\def^{\ifmmode\sp\else\^{}\fi}\catcode`\%=\active\def%{\%}storage}}%
\end{pgfscope}%
\begin{pgfscope}%
\definecolor{textcolor}{rgb}{0.000000,0.000000,0.000000}%
\pgfsetstrokecolor{textcolor}%
\pgfsetfillcolor{textcolor}%
\pgftext[x=3.566745in,y=2.530686in,,bottom]{\color{textcolor}{\rmfamily\fontsize{8.000000}{9.600000}\selectfont\catcode`\^=\active\def^{\ifmmode\sp\else\^{}\fi}\catcode`\%=\active\def%{\%}R2}}%
\end{pgfscope}%
\begin{pgfscope}%
\definecolor{textcolor}{rgb}{0.000000,0.000000,0.000000}%
\pgfsetstrokecolor{textcolor}%
\pgfsetfillcolor{textcolor}%
\pgftext[x=3.821777in,y=2.530686in,,bottom]{\color{textcolor}{\rmfamily\fontsize{8.000000}{9.600000}\selectfont\catcode`\^=\active\def^{\ifmmode\sp\else\^{}\fi}\catcode`\%=\active\def%{\%}O2}}%
\end{pgfscope}%
\begin{pgfscope}%
\definecolor{textcolor}{rgb}{0.000000,0.000000,0.000000}%
\pgfsetstrokecolor{textcolor}%
\pgfsetfillcolor{textcolor}%
\pgftext[x=4.086084in,y=2.530686in,,bottom]{\color{textcolor}{\rmfamily\fontsize{8.000000}{9.600000}\selectfont\catcode`\^=\active\def^{\ifmmode\sp\else\^{}\fi}\catcode`\%=\active\def%{\%}R1}}%
\end{pgfscope}%
\begin{pgfscope}%
\definecolor{textcolor}{rgb}{0.000000,0.000000,0.000000}%
\pgfsetstrokecolor{textcolor}%
\pgfsetfillcolor{textcolor}%
\pgftext[x=4.350390in,y=2.530686in,,bottom]{\color{textcolor}{\rmfamily\fontsize{8.000000}{9.600000}\selectfont\catcode`\^=\active\def^{\ifmmode\sp\else\^{}\fi}\catcode`\%=\active\def%{\%}O1}}%
\end{pgfscope}%
\begin{pgfscope}%
\definecolor{textcolor}{rgb}{0.000000,0.000000,0.000000}%
\pgfsetstrokecolor{textcolor}%
\pgfsetfillcolor{textcolor}%
\pgftext[x=4.515763in,y=2.452737in,right,top]{\color{textcolor}{\rmfamily\fontsize{9.000000}{10.800000}\selectfont\catcode`\^=\active\def^{\ifmmode\sp\else\^{}\fi}\catcode`\%=\active\def%{\%}$(a)$}}%
\end{pgfscope}%
\begin{pgfscope}%
\pgfsetbuttcap%
\pgfsetroundjoin%
\pgfsetlinewidth{0.903375pt}%
\definecolor{currentstroke}{rgb}{0.000000,0.000000,0.000000}%
\pgfsetstrokecolor{currentstroke}%
\pgfsetdash{{0.900000pt}{1.485000pt}}{0.000000pt}%
\pgfpathmoveto{\pgfqpoint{2.075341in}{2.013341in}}%
\pgfpathlineto{\pgfqpoint{2.175341in}{2.013341in}}%
\pgfpathlineto{\pgfqpoint{2.275341in}{2.013341in}}%
\pgfusepath{stroke}%
\end{pgfscope}%
\begin{pgfscope}%
\definecolor{textcolor}{rgb}{0.000000,0.000000,0.000000}%
\pgfsetstrokecolor{textcolor}%
\pgfsetfillcolor{textcolor}%
\pgftext[x=2.364230in,y=1.974452in,left,base]{\color{textcolor}{\rmfamily\fontsize{8.000000}{9.600000}\selectfont\catcode`\^=\active\def^{\ifmmode\sp\else\^{}\fi}\catcode`\%=\active\def%{\%}$\langle b^\dagger b\rangle$}}%
\end{pgfscope}%
\begin{pgfscope}%
\pgfsetrectcap%
\pgfsetroundjoin%
\pgfsetlinewidth{0.903375pt}%
\definecolor{currentstroke}{rgb}{0.121569,0.466667,0.705882}%
\pgfsetstrokecolor{currentstroke}%
\pgfsetdash{}{0pt}%
\pgfpathmoveto{\pgfqpoint{2.075341in}{1.874452in}}%
\pgfpathlineto{\pgfqpoint{2.175341in}{1.874452in}}%
\pgfpathlineto{\pgfqpoint{2.275341in}{1.874452in}}%
\pgfusepath{stroke}%
\end{pgfscope}%
\begin{pgfscope}%
\definecolor{textcolor}{rgb}{0.000000,0.000000,0.000000}%
\pgfsetstrokecolor{textcolor}%
\pgfsetfillcolor{textcolor}%
\pgftext[x=2.364230in,y=1.835563in,left,base]{\color{textcolor}{\rmfamily\fontsize{8.000000}{9.600000}\selectfont\catcode`\^=\active\def^{\ifmmode\sp\else\^{}\fi}\catcode`\%=\active\def%{\%}$P(D_1|R{=}0)$}}%
\end{pgfscope}%
\begin{pgfscope}%
\pgfsetbuttcap%
\pgfsetroundjoin%
\pgfsetlinewidth{0.903375pt}%
\definecolor{currentstroke}{rgb}{1.000000,0.000000,0.000000}%
\pgfsetstrokecolor{currentstroke}%
\pgfsetdash{{5.760000pt}{1.440000pt}{0.900000pt}{1.440000pt}}{0.000000pt}%
\pgfpathmoveto{\pgfqpoint{2.075341in}{1.735563in}}%
\pgfpathlineto{\pgfqpoint{2.175341in}{1.735563in}}%
\pgfpathlineto{\pgfqpoint{2.275341in}{1.735563in}}%
\pgfusepath{stroke}%
\end{pgfscope}%
\begin{pgfscope}%
\definecolor{textcolor}{rgb}{0.000000,0.000000,0.000000}%
\pgfsetstrokecolor{textcolor}%
\pgfsetfillcolor{textcolor}%
\pgftext[x=2.364230in,y=1.696674in,left,base]{\color{textcolor}{\rmfamily\fontsize{8.000000}{9.600000}\selectfont\catcode`\^=\active\def^{\ifmmode\sp\else\^{}\fi}\catcode`\%=\active\def%{\%}$P(D_2|R{=}1)$}}%
\end{pgfscope}%
\begin{pgfscope}%
\pgfsetbuttcap%
\pgfsetmiterjoin%
\definecolor{currentfill}{rgb}{1.000000,1.000000,1.000000}%
\pgfsetfillcolor{currentfill}%
\pgfsetlinewidth{0.000000pt}%
\definecolor{currentstroke}{rgb}{0.000000,0.000000,0.000000}%
\pgfsetstrokecolor{currentstroke}%
\pgfsetstrokeopacity{0.000000}%
\pgfsetdash{}{0pt}%
\pgfpathmoveto{\pgfqpoint{0.514740in}{0.462654in}}%
\pgfpathlineto{\pgfqpoint{4.535868in}{0.462654in}}%
\pgfpathlineto{\pgfqpoint{4.535868in}{1.258051in}}%
\pgfpathlineto{\pgfqpoint{0.514740in}{1.258051in}}%
\pgfpathlineto{\pgfqpoint{0.514740in}{0.462654in}}%
\pgfpathclose%
\pgfusepath{fill}%
\end{pgfscope}%
\begin{pgfscope}%
\pgfpathrectangle{\pgfqpoint{0.514740in}{0.462654in}}{\pgfqpoint{4.021128in}{0.795397in}}%
\pgfusepath{clip}%
\pgfsetbuttcap%
\pgfsetmiterjoin%
\definecolor{currentfill}{rgb}{0.901961,0.937255,0.972549}%
\pgfsetfillcolor{currentfill}%
\pgfsetfillopacity{0.700000}%
\pgfsetlinewidth{0.000000pt}%
\definecolor{currentstroke}{rgb}{0.901961,0.937255,0.972549}%
\pgfsetstrokecolor{currentstroke}%
\pgfsetstrokeopacity{0.700000}%
\pgfsetdash{}{0pt}%
\pgfpathmoveto{\pgfqpoint{0.514740in}{0.462654in}}%
\pgfpathlineto{\pgfqpoint{0.883823in}{0.462654in}}%
\pgfpathlineto{\pgfqpoint{0.883823in}{1.258051in}}%
\pgfpathlineto{\pgfqpoint{0.514740in}{1.258051in}}%
\pgfpathlineto{\pgfqpoint{0.514740in}{0.462654in}}%
\pgfpathclose%
\pgfusepath{fill}%
\end{pgfscope}%
\begin{pgfscope}%
\pgfpathrectangle{\pgfqpoint{0.514740in}{0.462654in}}{\pgfqpoint{4.021128in}{0.795397in}}%
\pgfusepath{clip}%
\pgfsetbuttcap%
\pgfsetmiterjoin%
\definecolor{currentfill}{rgb}{0.901961,0.937255,0.972549}%
\pgfsetfillcolor{currentfill}%
\pgfsetfillopacity{0.700000}%
\pgfsetlinewidth{0.000000pt}%
\definecolor{currentstroke}{rgb}{0.901961,0.937255,0.972549}%
\pgfsetstrokecolor{currentstroke}%
\pgfsetstrokeopacity{0.700000}%
\pgfsetdash{}{0pt}%
\pgfpathmoveto{\pgfqpoint{1.134218in}{0.462654in}}%
\pgfpathlineto{\pgfqpoint{1.503301in}{0.462654in}}%
\pgfpathlineto{\pgfqpoint{1.503301in}{1.258051in}}%
\pgfpathlineto{\pgfqpoint{1.134218in}{1.258051in}}%
\pgfpathlineto{\pgfqpoint{1.134218in}{0.462654in}}%
\pgfpathclose%
\pgfusepath{fill}%
\end{pgfscope}%
\begin{pgfscope}%
\pgfpathrectangle{\pgfqpoint{0.514740in}{0.462654in}}{\pgfqpoint{4.021128in}{0.795397in}}%
\pgfusepath{clip}%
\pgfsetbuttcap%
\pgfsetmiterjoin%
\definecolor{currentfill}{rgb}{0.921569,0.960784,0.937255}%
\pgfsetfillcolor{currentfill}%
\pgfsetfillopacity{0.750000}%
\pgfsetlinewidth{0.000000pt}%
\definecolor{currentstroke}{rgb}{0.921569,0.960784,0.937255}%
\pgfsetstrokecolor{currentstroke}%
\pgfsetstrokeopacity{0.750000}%
\pgfsetdash{}{0pt}%
\pgfpathmoveto{\pgfqpoint{1.642409in}{0.462654in}}%
\pgfpathlineto{\pgfqpoint{3.497191in}{0.462654in}}%
\pgfpathlineto{\pgfqpoint{3.497191in}{1.258051in}}%
\pgfpathlineto{\pgfqpoint{1.642409in}{1.258051in}}%
\pgfpathlineto{\pgfqpoint{1.642409in}{0.462654in}}%
\pgfpathclose%
\pgfusepath{fill}%
\end{pgfscope}%
\begin{pgfscope}%
\pgfpathrectangle{\pgfqpoint{0.514740in}{0.462654in}}{\pgfqpoint{4.021128in}{0.795397in}}%
\pgfusepath{clip}%
\pgfsetbuttcap%
\pgfsetmiterjoin%
\definecolor{currentfill}{rgb}{0.901961,0.937255,0.972549}%
\pgfsetfillcolor{currentfill}%
\pgfsetfillopacity{0.700000}%
\pgfsetlinewidth{0.000000pt}%
\definecolor{currentstroke}{rgb}{0.901961,0.937255,0.972549}%
\pgfsetstrokecolor{currentstroke}%
\pgfsetstrokeopacity{0.700000}%
\pgfsetdash{}{0pt}%
\pgfpathmoveto{\pgfqpoint{3.636299in}{0.462654in}}%
\pgfpathlineto{\pgfqpoint{4.007256in}{0.462654in}}%
\pgfpathlineto{\pgfqpoint{4.007256in}{1.258051in}}%
\pgfpathlineto{\pgfqpoint{3.636299in}{1.258051in}}%
\pgfpathlineto{\pgfqpoint{3.636299in}{0.462654in}}%
\pgfpathclose%
\pgfusepath{fill}%
\end{pgfscope}%
\begin{pgfscope}%
\pgfpathrectangle{\pgfqpoint{0.514740in}{0.462654in}}{\pgfqpoint{4.021128in}{0.795397in}}%
\pgfusepath{clip}%
\pgfsetbuttcap%
\pgfsetmiterjoin%
\definecolor{currentfill}{rgb}{0.901961,0.937255,0.972549}%
\pgfsetfillcolor{currentfill}%
\pgfsetfillopacity{0.700000}%
\pgfsetlinewidth{0.000000pt}%
\definecolor{currentstroke}{rgb}{0.901961,0.937255,0.972549}%
\pgfsetstrokecolor{currentstroke}%
\pgfsetstrokeopacity{0.700000}%
\pgfsetdash{}{0pt}%
\pgfpathmoveto{\pgfqpoint{4.164912in}{0.462654in}}%
\pgfpathlineto{\pgfqpoint{4.535868in}{0.462654in}}%
\pgfpathlineto{\pgfqpoint{4.535868in}{1.258051in}}%
\pgfpathlineto{\pgfqpoint{4.164912in}{1.258051in}}%
\pgfpathlineto{\pgfqpoint{4.164912in}{0.462654in}}%
\pgfpathclose%
\pgfusepath{fill}%
\end{pgfscope}%
\begin{pgfscope}%
\pgfsetbuttcap%
\pgfsetroundjoin%
\definecolor{currentfill}{rgb}{0.000000,0.000000,0.000000}%
\pgfsetfillcolor{currentfill}%
\pgfsetlinewidth{0.803000pt}%
\definecolor{currentstroke}{rgb}{0.000000,0.000000,0.000000}%
\pgfsetstrokecolor{currentstroke}%
\pgfsetdash{}{0pt}%
\pgfsys@defobject{currentmarker}{\pgfqpoint{0.000000in}{-0.048611in}}{\pgfqpoint{0.000000in}{0.000000in}}{%
\pgfpathmoveto{\pgfqpoint{0.000000in}{0.000000in}}%
\pgfpathlineto{\pgfqpoint{0.000000in}{-0.048611in}}%
\pgfusepath{stroke,fill}%
}%
\begin{pgfscope}%
\pgfsys@transformshift{0.514740in}{0.462654in}%
\pgfsys@useobject{currentmarker}{}%
\end{pgfscope}%
\end{pgfscope}%
\begin{pgfscope}%
\definecolor{textcolor}{rgb}{0.000000,0.000000,0.000000}%
\pgfsetstrokecolor{textcolor}%
\pgfsetfillcolor{textcolor}%
\pgftext[x=0.514740in,y=0.365432in,,top]{\color{textcolor}{\rmfamily\fontsize{8.000000}{9.600000}\selectfont\catcode`\^=\active\def^{\ifmmode\sp\else\^{}\fi}\catcode`\%=\active\def%{\%}0}}%
\end{pgfscope}%
\begin{pgfscope}%
\pgfsetbuttcap%
\pgfsetroundjoin%
\definecolor{currentfill}{rgb}{0.000000,0.000000,0.000000}%
\pgfsetfillcolor{currentfill}%
\pgfsetlinewidth{0.803000pt}%
\definecolor{currentstroke}{rgb}{0.000000,0.000000,0.000000}%
\pgfsetstrokecolor{currentstroke}%
\pgfsetdash{}{0pt}%
\pgfsys@defobject{currentmarker}{\pgfqpoint{0.000000in}{-0.048611in}}{\pgfqpoint{0.000000in}{0.000000in}}{%
\pgfpathmoveto{\pgfqpoint{0.000000in}{0.000000in}}%
\pgfpathlineto{\pgfqpoint{0.000000in}{-0.048611in}}%
\pgfusepath{stroke,fill}%
}%
\begin{pgfscope}%
\pgfsys@transformshift{0.978435in}{0.462654in}%
\pgfsys@useobject{currentmarker}{}%
\end{pgfscope}%
\end{pgfscope}%
\begin{pgfscope}%
\definecolor{textcolor}{rgb}{0.000000,0.000000,0.000000}%
\pgfsetstrokecolor{textcolor}%
\pgfsetfillcolor{textcolor}%
\pgftext[x=0.978435in,y=0.365432in,,top]{\color{textcolor}{\rmfamily\fontsize{8.000000}{9.600000}\selectfont\catcode`\^=\active\def^{\ifmmode\sp\else\^{}\fi}\catcode`\%=\active\def%{\%}500}}%
\end{pgfscope}%
\begin{pgfscope}%
\pgfsetbuttcap%
\pgfsetroundjoin%
\definecolor{currentfill}{rgb}{0.000000,0.000000,0.000000}%
\pgfsetfillcolor{currentfill}%
\pgfsetlinewidth{0.803000pt}%
\definecolor{currentstroke}{rgb}{0.000000,0.000000,0.000000}%
\pgfsetstrokecolor{currentstroke}%
\pgfsetdash{}{0pt}%
\pgfsys@defobject{currentmarker}{\pgfqpoint{0.000000in}{-0.048611in}}{\pgfqpoint{0.000000in}{0.000000in}}{%
\pgfpathmoveto{\pgfqpoint{0.000000in}{0.000000in}}%
\pgfpathlineto{\pgfqpoint{0.000000in}{-0.048611in}}%
\pgfusepath{stroke,fill}%
}%
\begin{pgfscope}%
\pgfsys@transformshift{1.442131in}{0.462654in}%
\pgfsys@useobject{currentmarker}{}%
\end{pgfscope}%
\end{pgfscope}%
\begin{pgfscope}%
\definecolor{textcolor}{rgb}{0.000000,0.000000,0.000000}%
\pgfsetstrokecolor{textcolor}%
\pgfsetfillcolor{textcolor}%
\pgftext[x=1.442131in,y=0.365432in,,top]{\color{textcolor}{\rmfamily\fontsize{8.000000}{9.600000}\selectfont\catcode`\^=\active\def^{\ifmmode\sp\else\^{}\fi}\catcode`\%=\active\def%{\%}1000}}%
\end{pgfscope}%
\begin{pgfscope}%
\pgfsetbuttcap%
\pgfsetroundjoin%
\definecolor{currentfill}{rgb}{0.000000,0.000000,0.000000}%
\pgfsetfillcolor{currentfill}%
\pgfsetlinewidth{0.803000pt}%
\definecolor{currentstroke}{rgb}{0.000000,0.000000,0.000000}%
\pgfsetstrokecolor{currentstroke}%
\pgfsetdash{}{0pt}%
\pgfsys@defobject{currentmarker}{\pgfqpoint{0.000000in}{-0.048611in}}{\pgfqpoint{0.000000in}{0.000000in}}{%
\pgfpathmoveto{\pgfqpoint{0.000000in}{0.000000in}}%
\pgfpathlineto{\pgfqpoint{0.000000in}{-0.048611in}}%
\pgfusepath{stroke,fill}%
}%
\begin{pgfscope}%
\pgfsys@transformshift{1.905826in}{0.462654in}%
\pgfsys@useobject{currentmarker}{}%
\end{pgfscope}%
\end{pgfscope}%
\begin{pgfscope}%
\definecolor{textcolor}{rgb}{0.000000,0.000000,0.000000}%
\pgfsetstrokecolor{textcolor}%
\pgfsetfillcolor{textcolor}%
\pgftext[x=1.905826in,y=0.365432in,,top]{\color{textcolor}{\rmfamily\fontsize{8.000000}{9.600000}\selectfont\catcode`\^=\active\def^{\ifmmode\sp\else\^{}\fi}\catcode`\%=\active\def%{\%}1500}}%
\end{pgfscope}%
\begin{pgfscope}%
\pgfsetbuttcap%
\pgfsetroundjoin%
\definecolor{currentfill}{rgb}{0.000000,0.000000,0.000000}%
\pgfsetfillcolor{currentfill}%
\pgfsetlinewidth{0.803000pt}%
\definecolor{currentstroke}{rgb}{0.000000,0.000000,0.000000}%
\pgfsetstrokecolor{currentstroke}%
\pgfsetdash{}{0pt}%
\pgfsys@defobject{currentmarker}{\pgfqpoint{0.000000in}{-0.048611in}}{\pgfqpoint{0.000000in}{0.000000in}}{%
\pgfpathmoveto{\pgfqpoint{0.000000in}{0.000000in}}%
\pgfpathlineto{\pgfqpoint{0.000000in}{-0.048611in}}%
\pgfusepath{stroke,fill}%
}%
\begin{pgfscope}%
\pgfsys@transformshift{2.369521in}{0.462654in}%
\pgfsys@useobject{currentmarker}{}%
\end{pgfscope}%
\end{pgfscope}%
\begin{pgfscope}%
\definecolor{textcolor}{rgb}{0.000000,0.000000,0.000000}%
\pgfsetstrokecolor{textcolor}%
\pgfsetfillcolor{textcolor}%
\pgftext[x=2.369521in,y=0.365432in,,top]{\color{textcolor}{\rmfamily\fontsize{8.000000}{9.600000}\selectfont\catcode`\^=\active\def^{\ifmmode\sp\else\^{}\fi}\catcode`\%=\active\def%{\%}2000}}%
\end{pgfscope}%
\begin{pgfscope}%
\pgfsetbuttcap%
\pgfsetroundjoin%
\definecolor{currentfill}{rgb}{0.000000,0.000000,0.000000}%
\pgfsetfillcolor{currentfill}%
\pgfsetlinewidth{0.803000pt}%
\definecolor{currentstroke}{rgb}{0.000000,0.000000,0.000000}%
\pgfsetstrokecolor{currentstroke}%
\pgfsetdash{}{0pt}%
\pgfsys@defobject{currentmarker}{\pgfqpoint{0.000000in}{-0.048611in}}{\pgfqpoint{0.000000in}{0.000000in}}{%
\pgfpathmoveto{\pgfqpoint{0.000000in}{0.000000in}}%
\pgfpathlineto{\pgfqpoint{0.000000in}{-0.048611in}}%
\pgfusepath{stroke,fill}%
}%
\begin{pgfscope}%
\pgfsys@transformshift{2.833217in}{0.462654in}%
\pgfsys@useobject{currentmarker}{}%
\end{pgfscope}%
\end{pgfscope}%
\begin{pgfscope}%
\definecolor{textcolor}{rgb}{0.000000,0.000000,0.000000}%
\pgfsetstrokecolor{textcolor}%
\pgfsetfillcolor{textcolor}%
\pgftext[x=2.833217in,y=0.365432in,,top]{\color{textcolor}{\rmfamily\fontsize{8.000000}{9.600000}\selectfont\catcode`\^=\active\def^{\ifmmode\sp\else\^{}\fi}\catcode`\%=\active\def%{\%}2500}}%
\end{pgfscope}%
\begin{pgfscope}%
\pgfsetbuttcap%
\pgfsetroundjoin%
\definecolor{currentfill}{rgb}{0.000000,0.000000,0.000000}%
\pgfsetfillcolor{currentfill}%
\pgfsetlinewidth{0.803000pt}%
\definecolor{currentstroke}{rgb}{0.000000,0.000000,0.000000}%
\pgfsetstrokecolor{currentstroke}%
\pgfsetdash{}{0pt}%
\pgfsys@defobject{currentmarker}{\pgfqpoint{0.000000in}{-0.048611in}}{\pgfqpoint{0.000000in}{0.000000in}}{%
\pgfpathmoveto{\pgfqpoint{0.000000in}{0.000000in}}%
\pgfpathlineto{\pgfqpoint{0.000000in}{-0.048611in}}%
\pgfusepath{stroke,fill}%
}%
\begin{pgfscope}%
\pgfsys@transformshift{3.296912in}{0.462654in}%
\pgfsys@useobject{currentmarker}{}%
\end{pgfscope}%
\end{pgfscope}%
\begin{pgfscope}%
\definecolor{textcolor}{rgb}{0.000000,0.000000,0.000000}%
\pgfsetstrokecolor{textcolor}%
\pgfsetfillcolor{textcolor}%
\pgftext[x=3.296912in,y=0.365432in,,top]{\color{textcolor}{\rmfamily\fontsize{8.000000}{9.600000}\selectfont\catcode`\^=\active\def^{\ifmmode\sp\else\^{}\fi}\catcode`\%=\active\def%{\%}3000}}%
\end{pgfscope}%
\begin{pgfscope}%
\pgfsetbuttcap%
\pgfsetroundjoin%
\definecolor{currentfill}{rgb}{0.000000,0.000000,0.000000}%
\pgfsetfillcolor{currentfill}%
\pgfsetlinewidth{0.803000pt}%
\definecolor{currentstroke}{rgb}{0.000000,0.000000,0.000000}%
\pgfsetstrokecolor{currentstroke}%
\pgfsetdash{}{0pt}%
\pgfsys@defobject{currentmarker}{\pgfqpoint{0.000000in}{-0.048611in}}{\pgfqpoint{0.000000in}{0.000000in}}{%
\pgfpathmoveto{\pgfqpoint{0.000000in}{0.000000in}}%
\pgfpathlineto{\pgfqpoint{0.000000in}{-0.048611in}}%
\pgfusepath{stroke,fill}%
}%
\begin{pgfscope}%
\pgfsys@transformshift{3.760607in}{0.462654in}%
\pgfsys@useobject{currentmarker}{}%
\end{pgfscope}%
\end{pgfscope}%
\begin{pgfscope}%
\definecolor{textcolor}{rgb}{0.000000,0.000000,0.000000}%
\pgfsetstrokecolor{textcolor}%
\pgfsetfillcolor{textcolor}%
\pgftext[x=3.760607in,y=0.365432in,,top]{\color{textcolor}{\rmfamily\fontsize{8.000000}{9.600000}\selectfont\catcode`\^=\active\def^{\ifmmode\sp\else\^{}\fi}\catcode`\%=\active\def%{\%}3500}}%
\end{pgfscope}%
\begin{pgfscope}%
\pgfsetbuttcap%
\pgfsetroundjoin%
\definecolor{currentfill}{rgb}{0.000000,0.000000,0.000000}%
\pgfsetfillcolor{currentfill}%
\pgfsetlinewidth{0.803000pt}%
\definecolor{currentstroke}{rgb}{0.000000,0.000000,0.000000}%
\pgfsetstrokecolor{currentstroke}%
\pgfsetdash{}{0pt}%
\pgfsys@defobject{currentmarker}{\pgfqpoint{0.000000in}{-0.048611in}}{\pgfqpoint{0.000000in}{0.000000in}}{%
\pgfpathmoveto{\pgfqpoint{0.000000in}{0.000000in}}%
\pgfpathlineto{\pgfqpoint{0.000000in}{-0.048611in}}%
\pgfusepath{stroke,fill}%
}%
\begin{pgfscope}%
\pgfsys@transformshift{4.224303in}{0.462654in}%
\pgfsys@useobject{currentmarker}{}%
\end{pgfscope}%
\end{pgfscope}%
\begin{pgfscope}%
\definecolor{textcolor}{rgb}{0.000000,0.000000,0.000000}%
\pgfsetstrokecolor{textcolor}%
\pgfsetfillcolor{textcolor}%
\pgftext[x=4.224303in,y=0.365432in,,top]{\color{textcolor}{\rmfamily\fontsize{8.000000}{9.600000}\selectfont\catcode`\^=\active\def^{\ifmmode\sp\else\^{}\fi}\catcode`\%=\active\def%{\%}4000}}%
\end{pgfscope}%
\begin{pgfscope}%
\definecolor{textcolor}{rgb}{0.000000,0.000000,0.000000}%
\pgfsetstrokecolor{textcolor}%
\pgfsetfillcolor{textcolor}%
\pgftext[x=2.525304in,y=0.211111in,,top]{\color{textcolor}{\rmfamily\fontsize{9.000000}{10.800000}\selectfont\catcode`\^=\active\def^{\ifmmode\sp\else\^{}\fi}\catcode`\%=\active\def%{\%}$\omega_b\,t$}}%
\end{pgfscope}%
\begin{pgfscope}%
\pgfsetbuttcap%
\pgfsetroundjoin%
\definecolor{currentfill}{rgb}{0.000000,0.000000,0.000000}%
\pgfsetfillcolor{currentfill}%
\pgfsetlinewidth{0.803000pt}%
\definecolor{currentstroke}{rgb}{0.000000,0.000000,0.000000}%
\pgfsetstrokecolor{currentstroke}%
\pgfsetdash{}{0pt}%
\pgfsys@defobject{currentmarker}{\pgfqpoint{-0.048611in}{0.000000in}}{\pgfqpoint{-0.000000in}{0.000000in}}{%
\pgfpathmoveto{\pgfqpoint{-0.000000in}{0.000000in}}%
\pgfpathlineto{\pgfqpoint{-0.048611in}{0.000000in}}%
\pgfusepath{stroke,fill}%
}%
\begin{pgfscope}%
\pgfsys@transformshift{0.514740in}{0.478250in}%
\pgfsys@useobject{currentmarker}{}%
\end{pgfscope}%
\end{pgfscope}%
\begin{pgfscope}%
\definecolor{textcolor}{rgb}{0.000000,0.000000,0.000000}%
\pgfsetstrokecolor{textcolor}%
\pgfsetfillcolor{textcolor}%
\pgftext[x=0.358489in, y=0.439670in, left, base]{\color{textcolor}{\rmfamily\fontsize{8.000000}{9.600000}\selectfont\catcode`\^=\active\def^{\ifmmode\sp\else\^{}\fi}\catcode`\%=\active\def%{\%}$0$}}%
\end{pgfscope}%
\begin{pgfscope}%
\pgfsetbuttcap%
\pgfsetroundjoin%
\definecolor{currentfill}{rgb}{0.000000,0.000000,0.000000}%
\pgfsetfillcolor{currentfill}%
\pgfsetlinewidth{0.803000pt}%
\definecolor{currentstroke}{rgb}{0.000000,0.000000,0.000000}%
\pgfsetstrokecolor{currentstroke}%
\pgfsetdash{}{0pt}%
\pgfsys@defobject{currentmarker}{\pgfqpoint{-0.048611in}{0.000000in}}{\pgfqpoint{-0.000000in}{0.000000in}}{%
\pgfpathmoveto{\pgfqpoint{-0.000000in}{0.000000in}}%
\pgfpathlineto{\pgfqpoint{-0.048611in}{0.000000in}}%
\pgfusepath{stroke,fill}%
}%
\begin{pgfscope}%
\pgfsys@transformshift{0.514740in}{0.790171in}%
\pgfsys@useobject{currentmarker}{}%
\end{pgfscope}%
\end{pgfscope}%
\begin{pgfscope}%
\definecolor{textcolor}{rgb}{0.000000,0.000000,0.000000}%
\pgfsetstrokecolor{textcolor}%
\pgfsetfillcolor{textcolor}%
\pgftext[x=0.266667in, y=0.751590in, left, base]{\color{textcolor}{\rmfamily\fontsize{8.000000}{9.600000}\selectfont\catcode`\^=\active\def^{\ifmmode\sp\else\^{}\fi}\catcode`\%=\active\def%{\%}$0.4$}}%
\end{pgfscope}%
\begin{pgfscope}%
\pgfsetbuttcap%
\pgfsetroundjoin%
\definecolor{currentfill}{rgb}{0.000000,0.000000,0.000000}%
\pgfsetfillcolor{currentfill}%
\pgfsetlinewidth{0.803000pt}%
\definecolor{currentstroke}{rgb}{0.000000,0.000000,0.000000}%
\pgfsetstrokecolor{currentstroke}%
\pgfsetdash{}{0pt}%
\pgfsys@defobject{currentmarker}{\pgfqpoint{-0.048611in}{0.000000in}}{\pgfqpoint{-0.000000in}{0.000000in}}{%
\pgfpathmoveto{\pgfqpoint{-0.000000in}{0.000000in}}%
\pgfpathlineto{\pgfqpoint{-0.048611in}{0.000000in}}%
\pgfusepath{stroke,fill}%
}%
\begin{pgfscope}%
\pgfsys@transformshift{0.514740in}{1.102091in}%
\pgfsys@useobject{currentmarker}{}%
\end{pgfscope}%
\end{pgfscope}%
\begin{pgfscope}%
\definecolor{textcolor}{rgb}{0.000000,0.000000,0.000000}%
\pgfsetstrokecolor{textcolor}%
\pgfsetfillcolor{textcolor}%
\pgftext[x=0.266667in, y=1.063511in, left, base]{\color{textcolor}{\rmfamily\fontsize{8.000000}{9.600000}\selectfont\catcode`\^=\active\def^{\ifmmode\sp\else\^{}\fi}\catcode`\%=\active\def%{\%}$0.8$}}%
\end{pgfscope}%
\begin{pgfscope}%
\definecolor{textcolor}{rgb}{0.000000,0.000000,0.000000}%
\pgfsetstrokecolor{textcolor}%
\pgfsetfillcolor{textcolor}%
\pgftext[x=0.211111in,y=0.860353in,,bottom,rotate=90.000000]{\color{textcolor}{\rmfamily\fontsize{9.000000}{10.800000}\selectfont\catcode`\^=\active\def^{\ifmmode\sp\else\^{}\fi}\catcode`\%=\active\def%{\%}Coherence}}%
\end{pgfscope}%
\begin{pgfscope}%
\pgfpathrectangle{\pgfqpoint{0.514740in}{0.462654in}}{\pgfqpoint{4.021128in}{0.795397in}}%
\pgfusepath{clip}%
\pgfsetrectcap%
\pgfsetroundjoin%
\pgfsetlinewidth{0.903375pt}%
\definecolor{currentstroke}{rgb}{0.768627,0.580392,0.164706}%
\pgfsetstrokecolor{currentstroke}%
\pgfsetdash{}{0pt}%
\pgfpathmoveto{\pgfqpoint{0.514740in}{0.478250in}}%
\pgfpathlineto{\pgfqpoint{0.899588in}{0.478881in}}%
\pgfpathlineto{\pgfqpoint{0.915354in}{0.480022in}}%
\pgfpathlineto{\pgfqpoint{0.920609in}{0.478816in}}%
\pgfpathlineto{\pgfqpoint{0.925864in}{0.480773in}}%
\pgfpathlineto{\pgfqpoint{0.931119in}{0.479929in}}%
\pgfpathlineto{\pgfqpoint{0.936375in}{0.480884in}}%
\pgfpathlineto{\pgfqpoint{0.946885in}{0.479369in}}%
\pgfpathlineto{\pgfqpoint{0.952140in}{0.482249in}}%
\pgfpathlineto{\pgfqpoint{0.957396in}{0.479297in}}%
\pgfpathlineto{\pgfqpoint{0.962651in}{0.481625in}}%
\pgfpathlineto{\pgfqpoint{0.967906in}{0.480500in}}%
\pgfpathlineto{\pgfqpoint{0.973161in}{0.482766in}}%
\pgfpathlineto{\pgfqpoint{0.978416in}{0.481543in}}%
\pgfpathlineto{\pgfqpoint{0.988927in}{0.483189in}}%
\pgfpathlineto{\pgfqpoint{0.994182in}{0.478866in}}%
\pgfpathlineto{\pgfqpoint{0.999437in}{0.486124in}}%
\pgfpathlineto{\pgfqpoint{1.004692in}{0.479780in}}%
\pgfpathlineto{\pgfqpoint{1.009948in}{0.483011in}}%
\pgfpathlineto{\pgfqpoint{1.015203in}{0.484302in}}%
\pgfpathlineto{\pgfqpoint{1.020458in}{0.481869in}}%
\pgfpathlineto{\pgfqpoint{1.025713in}{0.485105in}}%
\pgfpathlineto{\pgfqpoint{1.030969in}{0.480739in}}%
\pgfpathlineto{\pgfqpoint{1.036224in}{0.483805in}}%
\pgfpathlineto{\pgfqpoint{1.041479in}{0.481653in}}%
\pgfpathlineto{\pgfqpoint{1.041479in}{0.481653in}}%
\pgfpathlineto{\pgfqpoint{1.047662in}{0.485297in}}%
\pgfpathlineto{\pgfqpoint{1.050753in}{0.478822in}}%
\pgfpathlineto{\pgfqpoint{1.053844in}{0.485028in}}%
\pgfpathlineto{\pgfqpoint{1.056935in}{0.483992in}}%
\pgfpathlineto{\pgfqpoint{1.060027in}{0.481345in}}%
\pgfpathlineto{\pgfqpoint{1.063118in}{0.484628in}}%
\pgfpathlineto{\pgfqpoint{1.066209in}{0.484329in}}%
\pgfpathlineto{\pgfqpoint{1.069301in}{0.481840in}}%
\pgfpathlineto{\pgfqpoint{1.072392in}{0.481649in}}%
\pgfpathlineto{\pgfqpoint{1.075483in}{0.486836in}}%
\pgfpathlineto{\pgfqpoint{1.078575in}{0.479850in}}%
\pgfpathlineto{\pgfqpoint{1.084757in}{0.485641in}}%
\pgfpathlineto{\pgfqpoint{1.087849in}{0.482435in}}%
\pgfpathlineto{\pgfqpoint{1.094031in}{0.483784in}}%
\pgfpathlineto{\pgfqpoint{1.097122in}{0.483933in}}%
\pgfpathlineto{\pgfqpoint{1.100214in}{0.480520in}}%
\pgfpathlineto{\pgfqpoint{1.103305in}{0.486503in}}%
\pgfpathlineto{\pgfqpoint{1.106396in}{0.482579in}}%
\pgfpathlineto{\pgfqpoint{1.109488in}{0.480814in}}%
\pgfpathlineto{\pgfqpoint{1.112579in}{0.485865in}}%
\pgfpathlineto{\pgfqpoint{1.118762in}{0.481783in}}%
\pgfpathlineto{\pgfqpoint{1.121853in}{0.483044in}}%
\pgfpathlineto{\pgfqpoint{1.124944in}{0.486357in}}%
\pgfpathlineto{\pgfqpoint{1.128035in}{0.479809in}}%
\pgfpathlineto{\pgfqpoint{1.131127in}{0.484387in}}%
\pgfpathlineto{\pgfqpoint{1.134218in}{0.484539in}}%
\pgfpathlineto{\pgfqpoint{1.134218in}{0.478250in}}%
\pgfpathlineto{\pgfqpoint{1.503301in}{0.478841in}}%
\pgfpathlineto{\pgfqpoint{1.503301in}{0.478841in}}%
\pgfpathlineto{\pgfqpoint{1.507938in}{0.507769in}}%
\pgfpathlineto{\pgfqpoint{1.512575in}{0.544574in}}%
\pgfpathlineto{\pgfqpoint{1.521849in}{0.599863in}}%
\pgfpathlineto{\pgfqpoint{1.526486in}{0.639042in}}%
\pgfpathlineto{\pgfqpoint{1.531122in}{0.661590in}}%
\pgfpathlineto{\pgfqpoint{1.554307in}{0.805764in}}%
\pgfpathlineto{\pgfqpoint{1.558944in}{0.820125in}}%
\pgfpathlineto{\pgfqpoint{1.568218in}{0.873266in}}%
\pgfpathlineto{\pgfqpoint{1.572855in}{0.891103in}}%
\pgfpathlineto{\pgfqpoint{1.582129in}{0.932062in}}%
\pgfpathlineto{\pgfqpoint{1.586766in}{0.943856in}}%
\pgfpathlineto{\pgfqpoint{1.591403in}{0.969406in}}%
\pgfpathlineto{\pgfqpoint{1.596040in}{0.981637in}}%
\pgfpathlineto{\pgfqpoint{1.600677in}{0.989278in}}%
\pgfpathlineto{\pgfqpoint{1.605314in}{1.004611in}}%
\pgfpathlineto{\pgfqpoint{1.614588in}{1.022010in}}%
\pgfpathlineto{\pgfqpoint{1.619225in}{1.035173in}}%
\pgfpathlineto{\pgfqpoint{1.623862in}{1.032874in}}%
\pgfpathlineto{\pgfqpoint{1.628498in}{1.034125in}}%
\pgfpathlineto{\pgfqpoint{1.633135in}{1.043571in}}%
\pgfpathlineto{\pgfqpoint{1.637772in}{1.039861in}}%
\pgfpathlineto{\pgfqpoint{1.642409in}{1.041491in}}%
\pgfpathlineto{\pgfqpoint{1.642409in}{1.041491in}}%
\pgfpathlineto{\pgfqpoint{1.673322in}{1.035838in}}%
\pgfpathlineto{\pgfqpoint{1.704235in}{1.042289in}}%
\pgfpathlineto{\pgfqpoint{1.766061in}{1.038068in}}%
\pgfpathlineto{\pgfqpoint{1.796974in}{1.037434in}}%
\pgfpathlineto{\pgfqpoint{1.827887in}{1.041726in}}%
\pgfpathlineto{\pgfqpoint{1.858801in}{1.040879in}}%
\pgfpathlineto{\pgfqpoint{1.889714in}{1.034374in}}%
\pgfpathlineto{\pgfqpoint{1.920627in}{1.038998in}}%
\pgfpathlineto{\pgfqpoint{1.982453in}{1.039964in}}%
\pgfpathlineto{\pgfqpoint{2.013366in}{1.035700in}}%
\pgfpathlineto{\pgfqpoint{2.044279in}{1.036922in}}%
\pgfpathlineto{\pgfqpoint{2.075192in}{1.040464in}}%
\pgfpathlineto{\pgfqpoint{2.106105in}{1.034868in}}%
\pgfpathlineto{\pgfqpoint{2.137018in}{1.037728in}}%
\pgfpathlineto{\pgfqpoint{2.198844in}{1.037986in}}%
\pgfpathlineto{\pgfqpoint{2.229757in}{1.035057in}}%
\pgfpathlineto{\pgfqpoint{2.260670in}{1.034983in}}%
\pgfpathlineto{\pgfqpoint{2.291583in}{1.039733in}}%
\pgfpathlineto{\pgfqpoint{2.322496in}{1.035186in}}%
\pgfpathlineto{\pgfqpoint{2.384322in}{1.034470in}}%
\pgfpathlineto{\pgfqpoint{2.415235in}{1.038186in}}%
\pgfpathlineto{\pgfqpoint{2.477061in}{1.032803in}}%
\pgfpathlineto{\pgfqpoint{2.507974in}{1.036533in}}%
\pgfpathlineto{\pgfqpoint{2.538887in}{1.033827in}}%
\pgfpathlineto{\pgfqpoint{2.569800in}{1.035604in}}%
\pgfpathlineto{\pgfqpoint{2.600713in}{1.033374in}}%
\pgfpathlineto{\pgfqpoint{2.631626in}{1.035436in}}%
\pgfpathlineto{\pgfqpoint{2.662539in}{1.034635in}}%
\pgfpathlineto{\pgfqpoint{2.693452in}{1.031011in}}%
\pgfpathlineto{\pgfqpoint{2.724365in}{1.035445in}}%
\pgfpathlineto{\pgfqpoint{2.755278in}{1.033673in}}%
\pgfpathlineto{\pgfqpoint{2.786191in}{1.034316in}}%
\pgfpathlineto{\pgfqpoint{2.817104in}{1.031383in}}%
\pgfpathlineto{\pgfqpoint{2.878930in}{1.034358in}}%
\pgfpathlineto{\pgfqpoint{2.909843in}{1.031325in}}%
\pgfpathlineto{\pgfqpoint{2.940756in}{1.033383in}}%
\pgfpathlineto{\pgfqpoint{2.971669in}{1.031239in}}%
\pgfpathlineto{\pgfqpoint{3.002582in}{1.033374in}}%
\pgfpathlineto{\pgfqpoint{3.033495in}{1.030564in}}%
\pgfpathlineto{\pgfqpoint{3.095321in}{1.033516in}}%
\pgfpathlineto{\pgfqpoint{3.126234in}{1.029604in}}%
\pgfpathlineto{\pgfqpoint{3.157147in}{1.031766in}}%
\pgfpathlineto{\pgfqpoint{3.188061in}{1.029680in}}%
\pgfpathlineto{\pgfqpoint{3.218974in}{1.032684in}}%
\pgfpathlineto{\pgfqpoint{3.280800in}{1.028765in}}%
\pgfpathlineto{\pgfqpoint{3.311713in}{1.031553in}}%
\pgfpathlineto{\pgfqpoint{3.342626in}{1.029129in}}%
\pgfpathlineto{\pgfqpoint{3.373539in}{1.031251in}}%
\pgfpathlineto{\pgfqpoint{3.404452in}{1.028513in}}%
\pgfpathlineto{\pgfqpoint{3.435365in}{1.030233in}}%
\pgfpathlineto{\pgfqpoint{3.497191in}{1.028043in}}%
\pgfpathlineto{\pgfqpoint{3.497191in}{1.028043in}}%
\pgfpathlineto{\pgfqpoint{3.501828in}{1.028992in}}%
\pgfpathlineto{\pgfqpoint{3.506465in}{1.025379in}}%
\pgfpathlineto{\pgfqpoint{3.511102in}{1.018752in}}%
\pgfpathlineto{\pgfqpoint{3.515739in}{1.014532in}}%
\pgfpathlineto{\pgfqpoint{3.525012in}{1.000293in}}%
\pgfpathlineto{\pgfqpoint{3.534286in}{0.976397in}}%
\pgfpathlineto{\pgfqpoint{3.548197in}{0.935324in}}%
\pgfpathlineto{\pgfqpoint{3.562108in}{0.880327in}}%
\pgfpathlineto{\pgfqpoint{3.571382in}{0.841213in}}%
\pgfpathlineto{\pgfqpoint{3.617752in}{0.586897in}}%
\pgfpathlineto{\pgfqpoint{3.631662in}{0.504272in}}%
\pgfpathlineto{\pgfqpoint{3.636299in}{0.485780in}}%
\pgfpathlineto{\pgfqpoint{3.636299in}{0.485780in}}%
\pgfpathlineto{\pgfqpoint{3.648665in}{0.482844in}}%
\pgfpathlineto{\pgfqpoint{3.661030in}{0.484158in}}%
\pgfpathlineto{\pgfqpoint{3.673395in}{0.481994in}}%
\pgfpathlineto{\pgfqpoint{3.685760in}{0.485386in}}%
\pgfpathlineto{\pgfqpoint{3.698125in}{0.482042in}}%
\pgfpathlineto{\pgfqpoint{3.710491in}{0.484135in}}%
\pgfpathlineto{\pgfqpoint{3.722856in}{0.482322in}}%
\pgfpathlineto{\pgfqpoint{3.735221in}{0.483960in}}%
\pgfpathlineto{\pgfqpoint{3.747586in}{0.481222in}}%
\pgfpathlineto{\pgfqpoint{3.759951in}{0.483232in}}%
\pgfpathlineto{\pgfqpoint{3.772317in}{0.483746in}}%
\pgfpathlineto{\pgfqpoint{3.784682in}{0.481940in}}%
\pgfpathlineto{\pgfqpoint{3.809412in}{0.483089in}}%
\pgfpathlineto{\pgfqpoint{3.821777in}{0.484174in}}%
\pgfpathlineto{\pgfqpoint{3.834143in}{0.481556in}}%
\pgfpathlineto{\pgfqpoint{3.871238in}{0.482895in}}%
\pgfpathlineto{\pgfqpoint{3.883604in}{0.482030in}}%
\pgfpathlineto{\pgfqpoint{3.920699in}{0.482119in}}%
\pgfpathlineto{\pgfqpoint{3.933064in}{0.481649in}}%
\pgfpathlineto{\pgfqpoint{3.945430in}{0.483694in}}%
\pgfpathlineto{\pgfqpoint{3.957795in}{0.482177in}}%
\pgfpathlineto{\pgfqpoint{3.970160in}{0.482570in}}%
\pgfpathlineto{\pgfqpoint{3.982525in}{0.481748in}}%
\pgfpathlineto{\pgfqpoint{3.994890in}{0.484062in}}%
\pgfpathlineto{\pgfqpoint{4.007256in}{0.481656in}}%
\pgfpathlineto{\pgfqpoint{4.007256in}{0.481656in}}%
\pgfpathlineto{\pgfqpoint{4.017766in}{0.483593in}}%
\pgfpathlineto{\pgfqpoint{4.023021in}{0.484863in}}%
\pgfpathlineto{\pgfqpoint{4.028276in}{0.481753in}}%
\pgfpathlineto{\pgfqpoint{4.033532in}{0.481383in}}%
\pgfpathlineto{\pgfqpoint{4.038787in}{0.484330in}}%
\pgfpathlineto{\pgfqpoint{4.049297in}{0.482584in}}%
\pgfpathlineto{\pgfqpoint{4.054553in}{0.480482in}}%
\pgfpathlineto{\pgfqpoint{4.059808in}{0.483819in}}%
\pgfpathlineto{\pgfqpoint{4.075573in}{0.480930in}}%
\pgfpathlineto{\pgfqpoint{4.080829in}{0.482402in}}%
\pgfpathlineto{\pgfqpoint{4.086084in}{0.482572in}}%
\pgfpathlineto{\pgfqpoint{4.091339in}{0.480966in}}%
\pgfpathlineto{\pgfqpoint{4.096594in}{0.480681in}}%
\pgfpathlineto{\pgfqpoint{4.101849in}{0.481872in}}%
\pgfpathlineto{\pgfqpoint{4.138636in}{0.479363in}}%
\pgfpathlineto{\pgfqpoint{4.143891in}{0.480169in}}%
\pgfpathlineto{\pgfqpoint{4.159657in}{0.478966in}}%
\pgfpathlineto{\pgfqpoint{4.535868in}{0.479074in}}%
\pgfpathlineto{\pgfqpoint{4.535868in}{0.479074in}}%
\pgfusepath{stroke}%
\end{pgfscope}%
\begin{pgfscope}%
\pgfsetrectcap%
\pgfsetmiterjoin%
\pgfsetlinewidth{0.803000pt}%
\definecolor{currentstroke}{rgb}{0.000000,0.000000,0.000000}%
\pgfsetstrokecolor{currentstroke}%
\pgfsetdash{}{0pt}%
\pgfpathmoveto{\pgfqpoint{0.514740in}{0.462654in}}%
\pgfpathlineto{\pgfqpoint{0.514740in}{1.258051in}}%
\pgfusepath{stroke}%
\end{pgfscope}%
\begin{pgfscope}%
\pgfsetrectcap%
\pgfsetmiterjoin%
\pgfsetlinewidth{0.803000pt}%
\definecolor{currentstroke}{rgb}{0.000000,0.000000,0.000000}%
\pgfsetstrokecolor{currentstroke}%
\pgfsetdash{}{0pt}%
\pgfpathmoveto{\pgfqpoint{4.535868in}{0.462654in}}%
\pgfpathlineto{\pgfqpoint{4.535868in}{1.258051in}}%
\pgfusepath{stroke}%
\end{pgfscope}%
\begin{pgfscope}%
\pgfsetrectcap%
\pgfsetmiterjoin%
\pgfsetlinewidth{0.803000pt}%
\definecolor{currentstroke}{rgb}{0.000000,0.000000,0.000000}%
\pgfsetstrokecolor{currentstroke}%
\pgfsetdash{}{0pt}%
\pgfpathmoveto{\pgfqpoint{0.514740in}{0.462654in}}%
\pgfpathlineto{\pgfqpoint{4.535868in}{0.462654in}}%
\pgfusepath{stroke}%
\end{pgfscope}%
\begin{pgfscope}%
\pgfsetrectcap%
\pgfsetmiterjoin%
\pgfsetlinewidth{0.803000pt}%
\definecolor{currentstroke}{rgb}{0.000000,0.000000,0.000000}%
\pgfsetstrokecolor{currentstroke}%
\pgfsetdash{}{0pt}%
\pgfpathmoveto{\pgfqpoint{0.514740in}{1.258051in}}%
\pgfpathlineto{\pgfqpoint{4.535868in}{1.258051in}}%
\pgfusepath{stroke}%
\end{pgfscope}%
\begin{pgfscope}%
\definecolor{textcolor}{rgb}{0.000000,0.000000,0.000000}%
\pgfsetstrokecolor{textcolor}%
\pgfsetfillcolor{textcolor}%
\pgftext[x=4.515763in,y=1.234189in,right,top]{\color{textcolor}{\rmfamily\fontsize{9.000000}{10.800000}\selectfont\catcode`\^=\active\def^{\ifmmode\sp\else\^{}\fi}\catcode`\%=\active\def%{\%}$(b)$}}%
\end{pgfscope}%
\begin{pgfscope}%
\pgfsetrectcap%
\pgfsetroundjoin%
\pgfsetlinewidth{0.903375pt}%
\definecolor{currentstroke}{rgb}{0.768627,0.580392,0.164706}%
\pgfsetstrokecolor{currentstroke}%
\pgfsetdash{}{0pt}%
\pgfpathmoveto{\pgfqpoint{2.003563in}{0.755704in}}%
\pgfpathlineto{\pgfqpoint{2.103563in}{0.755704in}}%
\pgfpathlineto{\pgfqpoint{2.203563in}{0.755704in}}%
\pgfusepath{stroke}%
\end{pgfscope}%
\begin{pgfscope}%
\definecolor{textcolor}{rgb}{0.000000,0.000000,0.000000}%
\pgfsetstrokecolor{textcolor}%
\pgfsetfillcolor{textcolor}%
\pgftext[x=2.292451in,y=0.716815in,left,base]{\color{textcolor}{\rmfamily\fontsize{8.000000}{9.600000}\selectfont\catcode`\^=\active\def^{\ifmmode\sp\else\^{}\fi}\catcode`\%=\active\def%{\%}$2|\rho_{R0D_1,R1D_2}|$}}%
\end{pgfscope}%
\end{pgfpicture}%
\makeatother%
\endgroup%
}
\caption{(En color en l\'inea) Protocolo completo para $T\omega_b=2000$: ocupaci\'on del modo $\langle b^\dagger b\rangle$ (negro, punteado), y poblaciones de $D_1$ condicionada a la rama $R=0$ (azul, continua) y de $D_2$ condicionada a $R=1$ (rojo, punto-raya), y coherencia almacenada $2|\rho_{R0D_1,R1D_2}|$ (dorado). Las etiquetas marcan las etapas: carga (L1, L2), escritura Raman (W1, W2), almacenamiento, lectura Raman (R1, R2) y salida (O1, O2); el n\'umero es el intervalo temporal. El sombreado azul marca las ventanas de carga y emisi\'on; el verde marca la espera. La coherencia se define solo mientras ambas componentes est\'an almacenadas. Par\'ametros del sistema: como en la Fig.~\ref{fig:storage}, con $\kappa_{\rm in}/\kappa=\kappa_{\rm out}/\kappa=10$ y una separaci\'on de $100/\omega_b$ entre los dos intervalos temporales.}
\label{fig:protocol}
\end{figure}

Sea $\rho_{R,{\rm out}}$ el estado conjunto de la referencia y las dos ventanas de salida, incluyendo el vac\'io, y $\Pi_1$ el proyector sobre exactamente un fon\'on en la salida. Definimos
\begin{align}
s&=\mathrm{Tr}(\Pi_1\rho_{R,{\rm out}}),\qquad
F_h=F/s,\nonumber\\
F&=\bra{\Psi_{\rm id}}\rho_{R,{\rm out}}\ket{\Psi_{\rm id}},\label{eq:metrics}
\end{align}
con $\ket{\Psi_{\rm id}}$ el estado de salida ideal de la Ec.~\eqref{eq:input}; la fase determinista del qubit l\'ogico se calibra antes de evaluar $F$. La fidelidad condicional $F_h$ no es una eficiencia y no implica que cada intento produjo el qubit. Es el solapamiento de la salida normalizada de un fon\'on y su referencia con un estado de dos qubits m\'aximamente entrelazado, una medida de fidelidad de entrelazamiento apropiada para una memoria de un solo cuanto~\cite{surmacz2006}, y un valor $F_h>1/2$ certifica que este estado condicional est\'a entrelazado, entrelazamiento que es destilable en principio~\cite{bennett1996,horodecki1997}. Para probabilidades de recuperaci\'on $P_0,P_1$ por rama y visibilidad residual $\mathcal V$, sin errores de poblaci\'on adicionales, $F=(P_0+P_1)/4+\sqrt{P_0P_1}\,\mathcal V/2$ y $s=(P_0+P_1)/2$, de modo que $F_h=(1+\mathcal V)/2$ para $P_0=P_1$. Esta identidad separa la p\'erdida de excitaciones de la p\'erdida de coherencia: un aumento de $F$ puede provenir de una mayor probabilidad de \'exito aun si $F_h$ disminuye.

La referencia pasiva se simula como dos modos ac\'usticos independientes con la misma p\'erdida intr\'inseca $\kappa$ y los mismos puertos, uno por intervalo temporal. Cada modo tiene a lo sumo una excitaci\'on y solo decae por amplitud, de modo que el resultado se escribe con las probabilidades de recuperaci\'on de cada intervalo y la coherencia entre ellos:
\begin{align}
p_1&=\eta_{L1}\eta_R\,e^{-\kappa(\tau_{12}+t_{L2}+T+t_E)},\qquad
p_2=\eta_{L2}\eta_R\,e^{-\kappa T},\nonumber\\
s_{\rm pas}&=\frac{p_1+p_2}{2},\qquad
F_{\rm pas}=\frac{p_1+p_2}{4}+\frac{\sqrt{p_1p_2}}{2},\label{eq:passive}
\end{align}
con $\tau_{12}$ la separaci\'on entre las dos ventanas de carga y $t_{L2}$ y $t_E$ las duraciones de la segunda carga y de la primera extracci\'on: el primer intervalo espera durante todo ese tiempo y el segundo solo durante $T$.

\section{Resultados y referencia pasiva}
\label{sec:results}

Los resultados de esta secci\'on provienen de la din\'amica abierta del modelo reducido con los par\'ametros de la Fig.~\ref{fig:raman_map} y los pulsos de la Sec.~\ref{sec:raman} optimizados una sola vez; los par\'ametros se var\'ian de uno en uno con los pulsos fijos, salvo al escalar el acoplo por deformaci\'on. La tolerancia relativa de integraci\'on es $10^{-9}$, ya que con $10^{-6}$ el modelo completo devolvi\'o matrices de densidad con autovalores negativos; todas las matrices reportadas tienen autovalores no negativos y cumplen $|\rho_{01}|^2\le\rho_{00}\rho_{11}$. Truncar el modo de fonones en $N_{\rm cut}=3$ en lugar de $2$ cambia $F_h$ en cerca de $0.016$ y $s$ en cerca de $0.012$ a $T\omega_b=2000$, una incertidumbre que se aplica a la segunda cifra decimal de los resultados; las corridas completas usan $N_{\rm cut}=2$.

\begin{figure}[h]
	\centering
	\resizebox{\columnwidth}{!}{%% Creator: Matplotlib, PGF backend
%%
%% To include the figure in your LaTeX document, write
%%   \input{<filename>.pgf}
%%
%% Make sure the required packages are loaded in your preamble
%%   \usepackage{pgf}
%%
%% Also ensure that all the required font packages are loaded; for instance,
%% the lmodern package is sometimes necessary when using math font.
%%   \usepackage{lmodern}
%%
%% Figures using additional raster images can only be included by \input if
%% they are in the same directory as the main LaTeX file. For loading figures
%% from other directories you can use the `import` package
%%   \usepackage{import}
%%
%% and then include the figures with
%%   \import{<path to file>}{<filename>.pgf}
%%
%% Matplotlib used the following preamble
%%   \def\mathdefault#1{#1}
%%   \everymath=\expandafter{\the\everymath\displaystyle}
%%   \IfFileExists{scrextend.sty}{
%%     \usepackage[fontsize=8.000000pt]{scrextend}
%%   }{
%%     \renewcommand{\normalsize}{\fontsize{8.000000}{9.600000}\selectfont}
%%     \normalsize
%%   }
%%   \usepackage{amsmath,amssymb}
%%   \makeatletter\@ifpackageloaded{underscore}{}{\usepackage[strings]{underscore}}\makeatother
%%
\begingroup%
\makeatletter%
\begin{pgfpicture}%
\pgfpathrectangle{\pgfpointorigin}{\pgfqpoint{4.934702in}{2.200342in}}%
\pgfusepath{use as bounding box, clip}%
\begin{pgfscope}%
\pgfsetbuttcap%
\pgfsetmiterjoin%
\definecolor{currentfill}{rgb}{1.000000,1.000000,1.000000}%
\pgfsetfillcolor{currentfill}%
\pgfsetlinewidth{0.000000pt}%
\definecolor{currentstroke}{rgb}{1.000000,1.000000,1.000000}%
\pgfsetstrokecolor{currentstroke}%
\pgfsetdash{}{0pt}%
\pgfpathmoveto{\pgfqpoint{0.000000in}{0.000000in}}%
\pgfpathlineto{\pgfqpoint{4.934702in}{0.000000in}}%
\pgfpathlineto{\pgfqpoint{4.934702in}{2.200342in}}%
\pgfpathlineto{\pgfqpoint{0.000000in}{2.200342in}}%
\pgfpathlineto{\pgfqpoint{0.000000in}{0.000000in}}%
\pgfpathclose%
\pgfusepath{fill}%
\end{pgfscope}%
\begin{pgfscope}%
\pgfsetbuttcap%
\pgfsetmiterjoin%
\definecolor{currentfill}{rgb}{1.000000,1.000000,1.000000}%
\pgfsetfillcolor{currentfill}%
\pgfsetlinewidth{0.000000pt}%
\definecolor{currentstroke}{rgb}{0.000000,0.000000,0.000000}%
\pgfsetstrokecolor{currentstroke}%
\pgfsetstrokeopacity{0.000000}%
\pgfsetdash{}{0pt}%
\pgfpathmoveto{\pgfqpoint{0.500851in}{0.496542in}}%
\pgfpathlineto{\pgfqpoint{1.543842in}{0.496542in}}%
\pgfpathlineto{\pgfqpoint{1.543842in}{1.736515in}}%
\pgfpathlineto{\pgfqpoint{0.500851in}{1.736515in}}%
\pgfpathlineto{\pgfqpoint{0.500851in}{0.496542in}}%
\pgfpathclose%
\pgfusepath{fill}%
\end{pgfscope}%
\begin{pgfscope}%
\pgfsetbuttcap%
\pgfsetroundjoin%
\definecolor{currentfill}{rgb}{0.000000,0.000000,0.000000}%
\pgfsetfillcolor{currentfill}%
\pgfsetlinewidth{0.803000pt}%
\definecolor{currentstroke}{rgb}{0.000000,0.000000,0.000000}%
\pgfsetstrokecolor{currentstroke}%
\pgfsetdash{}{0pt}%
\pgfsys@defobject{currentmarker}{\pgfqpoint{0.000000in}{-0.048611in}}{\pgfqpoint{0.000000in}{0.000000in}}{%
\pgfpathmoveto{\pgfqpoint{0.000000in}{0.000000in}}%
\pgfpathlineto{\pgfqpoint{0.000000in}{-0.048611in}}%
\pgfusepath{stroke,fill}%
}%
\begin{pgfscope}%
\pgfsys@transformshift{0.500851in}{0.496542in}%
\pgfsys@useobject{currentmarker}{}%
\end{pgfscope}%
\end{pgfscope}%
\begin{pgfscope}%
\definecolor{textcolor}{rgb}{0.000000,0.000000,0.000000}%
\pgfsetstrokecolor{textcolor}%
\pgfsetfillcolor{textcolor}%
\pgftext[x=0.500851in,y=0.399320in,,top]{\color{textcolor}{\rmfamily\fontsize{8.000000}{9.600000}\selectfont\catcode`\^=\active\def^{\ifmmode\sp\else\^{}\fi}\catcode`\%=\active\def%{\%}0}}%
\end{pgfscope}%
\begin{pgfscope}%
\pgfsetbuttcap%
\pgfsetroundjoin%
\definecolor{currentfill}{rgb}{0.000000,0.000000,0.000000}%
\pgfsetfillcolor{currentfill}%
\pgfsetlinewidth{0.803000pt}%
\definecolor{currentstroke}{rgb}{0.000000,0.000000,0.000000}%
\pgfsetstrokecolor{currentstroke}%
\pgfsetdash{}{0pt}%
\pgfsys@defobject{currentmarker}{\pgfqpoint{0.000000in}{-0.048611in}}{\pgfqpoint{0.000000in}{0.000000in}}{%
\pgfpathmoveto{\pgfqpoint{0.000000in}{0.000000in}}%
\pgfpathlineto{\pgfqpoint{0.000000in}{-0.048611in}}%
\pgfusepath{stroke,fill}%
}%
\begin{pgfscope}%
\pgfsys@transformshift{1.015908in}{0.496542in}%
\pgfsys@useobject{currentmarker}{}%
\end{pgfscope}%
\end{pgfscope}%
\begin{pgfscope}%
\definecolor{textcolor}{rgb}{0.000000,0.000000,0.000000}%
\pgfsetstrokecolor{textcolor}%
\pgfsetfillcolor{textcolor}%
\pgftext[x=1.015908in,y=0.399320in,,top]{\color{textcolor}{\rmfamily\fontsize{8.000000}{9.600000}\selectfont\catcode`\^=\active\def^{\ifmmode\sp\else\^{}\fi}\catcode`\%=\active\def%{\%}4}}%
\end{pgfscope}%
\begin{pgfscope}%
\pgfsetbuttcap%
\pgfsetroundjoin%
\definecolor{currentfill}{rgb}{0.000000,0.000000,0.000000}%
\pgfsetfillcolor{currentfill}%
\pgfsetlinewidth{0.803000pt}%
\definecolor{currentstroke}{rgb}{0.000000,0.000000,0.000000}%
\pgfsetstrokecolor{currentstroke}%
\pgfsetdash{}{0pt}%
\pgfsys@defobject{currentmarker}{\pgfqpoint{0.000000in}{-0.048611in}}{\pgfqpoint{0.000000in}{0.000000in}}{%
\pgfpathmoveto{\pgfqpoint{0.000000in}{0.000000in}}%
\pgfpathlineto{\pgfqpoint{0.000000in}{-0.048611in}}%
\pgfusepath{stroke,fill}%
}%
\begin{pgfscope}%
\pgfsys@transformshift{1.530965in}{0.496542in}%
\pgfsys@useobject{currentmarker}{}%
\end{pgfscope}%
\end{pgfscope}%
\begin{pgfscope}%
\definecolor{textcolor}{rgb}{0.000000,0.000000,0.000000}%
\pgfsetstrokecolor{textcolor}%
\pgfsetfillcolor{textcolor}%
\pgftext[x=1.530965in,y=0.399320in,,top]{\color{textcolor}{\rmfamily\fontsize{8.000000}{9.600000}\selectfont\catcode`\^=\active\def^{\ifmmode\sp\else\^{}\fi}\catcode`\%=\active\def%{\%}8}}%
\end{pgfscope}%
\begin{pgfscope}%
\pgfsetbuttcap%
\pgfsetroundjoin%
\definecolor{currentfill}{rgb}{0.000000,0.000000,0.000000}%
\pgfsetfillcolor{currentfill}%
\pgfsetlinewidth{0.803000pt}%
\definecolor{currentstroke}{rgb}{0.000000,0.000000,0.000000}%
\pgfsetstrokecolor{currentstroke}%
\pgfsetdash{}{0pt}%
\pgfsys@defobject{currentmarker}{\pgfqpoint{-0.048611in}{0.000000in}}{\pgfqpoint{-0.000000in}{0.000000in}}{%
\pgfpathmoveto{\pgfqpoint{-0.000000in}{0.000000in}}%
\pgfpathlineto{\pgfqpoint{-0.048611in}{0.000000in}}%
\pgfusepath{stroke,fill}%
}%
\begin{pgfscope}%
\pgfsys@transformshift{0.500851in}{0.496542in}%
\pgfsys@useobject{currentmarker}{}%
\end{pgfscope}%
\end{pgfscope}%
\begin{pgfscope}%
\definecolor{textcolor}{rgb}{0.000000,0.000000,0.000000}%
\pgfsetstrokecolor{textcolor}%
\pgfsetfillcolor{textcolor}%
\pgftext[x=0.344600in, y=0.457962in, left, base]{\color{textcolor}{\rmfamily\fontsize{8.000000}{9.600000}\selectfont\catcode`\^=\active\def^{\ifmmode\sp\else\^{}\fi}\catcode`\%=\active\def%{\%}$0$}}%
\end{pgfscope}%
\begin{pgfscope}%
\pgfsetbuttcap%
\pgfsetroundjoin%
\definecolor{currentfill}{rgb}{0.000000,0.000000,0.000000}%
\pgfsetfillcolor{currentfill}%
\pgfsetlinewidth{0.803000pt}%
\definecolor{currentstroke}{rgb}{0.000000,0.000000,0.000000}%
\pgfsetstrokecolor{currentstroke}%
\pgfsetdash{}{0pt}%
\pgfsys@defobject{currentmarker}{\pgfqpoint{-0.048611in}{0.000000in}}{\pgfqpoint{-0.000000in}{0.000000in}}{%
\pgfpathmoveto{\pgfqpoint{-0.000000in}{0.000000in}}%
\pgfpathlineto{\pgfqpoint{-0.048611in}{0.000000in}}%
\pgfusepath{stroke,fill}%
}%
\begin{pgfscope}%
\pgfsys@transformshift{0.500851in}{1.087005in}%
\pgfsys@useobject{currentmarker}{}%
\end{pgfscope}%
\end{pgfscope}%
\begin{pgfscope}%
\definecolor{textcolor}{rgb}{0.000000,0.000000,0.000000}%
\pgfsetstrokecolor{textcolor}%
\pgfsetfillcolor{textcolor}%
\pgftext[x=0.252778in, y=1.048425in, left, base]{\color{textcolor}{\rmfamily\fontsize{8.000000}{9.600000}\selectfont\catcode`\^=\active\def^{\ifmmode\sp\else\^{}\fi}\catcode`\%=\active\def%{\%}$0.5$}}%
\end{pgfscope}%
\begin{pgfscope}%
\pgfsetbuttcap%
\pgfsetroundjoin%
\definecolor{currentfill}{rgb}{0.000000,0.000000,0.000000}%
\pgfsetfillcolor{currentfill}%
\pgfsetlinewidth{0.803000pt}%
\definecolor{currentstroke}{rgb}{0.000000,0.000000,0.000000}%
\pgfsetstrokecolor{currentstroke}%
\pgfsetdash{}{0pt}%
\pgfsys@defobject{currentmarker}{\pgfqpoint{-0.048611in}{0.000000in}}{\pgfqpoint{-0.000000in}{0.000000in}}{%
\pgfpathmoveto{\pgfqpoint{-0.000000in}{0.000000in}}%
\pgfpathlineto{\pgfqpoint{-0.048611in}{0.000000in}}%
\pgfusepath{stroke,fill}%
}%
\begin{pgfscope}%
\pgfsys@transformshift{0.500851in}{1.677468in}%
\pgfsys@useobject{currentmarker}{}%
\end{pgfscope}%
\end{pgfscope}%
\begin{pgfscope}%
\definecolor{textcolor}{rgb}{0.000000,0.000000,0.000000}%
\pgfsetstrokecolor{textcolor}%
\pgfsetfillcolor{textcolor}%
\pgftext[x=0.344600in, y=1.638888in, left, base]{\color{textcolor}{\rmfamily\fontsize{8.000000}{9.600000}\selectfont\catcode`\^=\active\def^{\ifmmode\sp\else\^{}\fi}\catcode`\%=\active\def%{\%}$1$}}%
\end{pgfscope}%
\begin{pgfscope}%
\definecolor{textcolor}{rgb}{0.000000,0.000000,0.000000}%
\pgfsetstrokecolor{textcolor}%
\pgfsetfillcolor{textcolor}%
\pgftext[x=0.211111in,y=1.116529in,,bottom,rotate=90.000000]{\color{textcolor}{\rmfamily\fontsize{9.000000}{10.800000}\selectfont\catcode`\^=\active\def^{\ifmmode\sp\else\^{}\fi}\catcode`\%=\active\def%{\%}$F$}}%
\end{pgfscope}%
\begin{pgfscope}%
\pgfpathrectangle{\pgfqpoint{0.500851in}{0.496542in}}{\pgfqpoint{1.042991in}{1.239972in}}%
\pgfusepath{clip}%
\pgfsetrectcap%
\pgfsetroundjoin%
\pgfsetlinewidth{0.903375pt}%
\definecolor{currentstroke}{rgb}{1.000000,0.000000,0.000000}%
\pgfsetstrokecolor{currentstroke}%
\pgfsetdash{}{0pt}%
\pgfpathmoveto{\pgfqpoint{0.500851in}{1.209589in}}%
\pgfpathlineto{\pgfqpoint{0.521557in}{1.221326in}}%
\pgfpathlineto{\pgfqpoint{0.537086in}{1.227986in}}%
\pgfpathlineto{\pgfqpoint{0.547439in}{1.230768in}}%
\pgfpathlineto{\pgfqpoint{0.568145in}{1.233042in}}%
\pgfpathlineto{\pgfqpoint{0.599204in}{1.234664in}}%
\pgfpathlineto{\pgfqpoint{0.609556in}{1.234378in}}%
\pgfpathlineto{\pgfqpoint{0.625086in}{1.231201in}}%
\pgfpathlineto{\pgfqpoint{0.650968in}{1.225067in}}%
\pgfpathlineto{\pgfqpoint{0.682027in}{1.223952in}}%
\pgfpathlineto{\pgfqpoint{0.723438in}{1.224053in}}%
\pgfpathlineto{\pgfqpoint{0.816615in}{1.227225in}}%
\pgfpathlineto{\pgfqpoint{0.847673in}{1.225129in}}%
\pgfpathlineto{\pgfqpoint{0.940849in}{1.218238in}}%
\pgfpathlineto{\pgfqpoint{1.023673in}{1.213571in}}%
\pgfpathlineto{\pgfqpoint{1.111672in}{1.208357in}}%
\pgfpathlineto{\pgfqpoint{1.189319in}{1.207401in}}%
\pgfpathlineto{\pgfqpoint{1.354966in}{1.204861in}}%
\pgfpathlineto{\pgfqpoint{1.448142in}{1.199813in}}%
\pgfpathlineto{\pgfqpoint{1.530965in}{1.194041in}}%
\pgfpathlineto{\pgfqpoint{1.530965in}{1.194041in}}%
\pgfusepath{stroke}%
\end{pgfscope}%
\begin{pgfscope}%
\pgfpathrectangle{\pgfqpoint{0.500851in}{0.496542in}}{\pgfqpoint{1.042991in}{1.239972in}}%
\pgfusepath{clip}%
\pgfsetbuttcap%
\pgfsetroundjoin%
\definecolor{currentfill}{rgb}{0.000000,0.000000,0.000000}%
\pgfsetfillcolor{currentfill}%
\pgfsetfillopacity{0.000000}%
\pgfsetlinewidth{0.803000pt}%
\definecolor{currentstroke}{rgb}{1.000000,0.000000,0.000000}%
\pgfsetstrokecolor{currentstroke}%
\pgfsetdash{}{0pt}%
\pgfsys@defobject{currentmarker}{\pgfqpoint{-0.017500in}{-0.017500in}}{\pgfqpoint{0.017500in}{0.017500in}}{%
\pgfpathmoveto{\pgfqpoint{0.000000in}{-0.017500in}}%
\pgfpathcurveto{\pgfqpoint{0.004641in}{-0.017500in}}{\pgfqpoint{0.009093in}{-0.015656in}}{\pgfqpoint{0.012374in}{-0.012374in}}%
\pgfpathcurveto{\pgfqpoint{0.015656in}{-0.009093in}}{\pgfqpoint{0.017500in}{-0.004641in}}{\pgfqpoint{0.017500in}{0.000000in}}%
\pgfpathcurveto{\pgfqpoint{0.017500in}{0.004641in}}{\pgfqpoint{0.015656in}{0.009093in}}{\pgfqpoint{0.012374in}{0.012374in}}%
\pgfpathcurveto{\pgfqpoint{0.009093in}{0.015656in}}{\pgfqpoint{0.004641in}{0.017500in}}{\pgfqpoint{0.000000in}{0.017500in}}%
\pgfpathcurveto{\pgfqpoint{-0.004641in}{0.017500in}}{\pgfqpoint{-0.009093in}{0.015656in}}{\pgfqpoint{-0.012374in}{0.012374in}}%
\pgfpathcurveto{\pgfqpoint{-0.015656in}{0.009093in}}{\pgfqpoint{-0.017500in}{0.004641in}}{\pgfqpoint{-0.017500in}{0.000000in}}%
\pgfpathcurveto{\pgfqpoint{-0.017500in}{-0.004641in}}{\pgfqpoint{-0.015656in}{-0.009093in}}{\pgfqpoint{-0.012374in}{-0.012374in}}%
\pgfpathcurveto{\pgfqpoint{-0.009093in}{-0.015656in}}{\pgfqpoint{-0.004641in}{-0.017500in}}{\pgfqpoint{0.000000in}{-0.017500in}}%
\pgfpathlineto{\pgfqpoint{0.000000in}{-0.017500in}}%
\pgfpathclose%
\pgfusepath{stroke,fill}%
}%
\begin{pgfscope}%
\pgfsys@transformshift{0.500851in}{1.209589in}%
\pgfsys@useobject{currentmarker}{}%
\end{pgfscope}%
\begin{pgfscope}%
\pgfsys@transformshift{0.552357in}{1.231506in}%
\pgfsys@useobject{currentmarker}{}%
\end{pgfscope}%
\begin{pgfscope}%
\pgfsys@transformshift{0.603862in}{1.234696in}%
\pgfsys@useobject{currentmarker}{}%
\end{pgfscope}%
\begin{pgfscope}%
\pgfsys@transformshift{0.655368in}{1.224727in}%
\pgfsys@useobject{currentmarker}{}%
\end{pgfscope}%
\begin{pgfscope}%
\pgfsys@transformshift{0.706874in}{1.223686in}%
\pgfsys@useobject{currentmarker}{}%
\end{pgfscope}%
\begin{pgfscope}%
\pgfsys@transformshift{0.758380in}{1.225709in}%
\pgfsys@useobject{currentmarker}{}%
\end{pgfscope}%
\begin{pgfscope}%
\pgfsys@transformshift{0.809885in}{1.227324in}%
\pgfsys@useobject{currentmarker}{}%
\end{pgfscope}%
\begin{pgfscope}%
\pgfsys@transformshift{0.861391in}{1.224008in}%
\pgfsys@useobject{currentmarker}{}%
\end{pgfscope}%
\begin{pgfscope}%
\pgfsys@transformshift{0.912897in}{1.219643in}%
\pgfsys@useobject{currentmarker}{}%
\end{pgfscope}%
\begin{pgfscope}%
\pgfsys@transformshift{0.964402in}{1.217136in}%
\pgfsys@useobject{currentmarker}{}%
\end{pgfscope}%
\begin{pgfscope}%
\pgfsys@transformshift{1.015908in}{1.214085in}%
\pgfsys@useobject{currentmarker}{}%
\end{pgfscope}%
\begin{pgfscope}%
\pgfsys@transformshift{1.067414in}{1.210580in}%
\pgfsys@useobject{currentmarker}{}%
\end{pgfscope}%
\begin{pgfscope}%
\pgfsys@transformshift{1.118919in}{1.208177in}%
\pgfsys@useobject{currentmarker}{}%
\end{pgfscope}%
\begin{pgfscope}%
\pgfsys@transformshift{1.170425in}{1.207582in}%
\pgfsys@useobject{currentmarker}{}%
\end{pgfscope}%
\begin{pgfscope}%
\pgfsys@transformshift{1.221931in}{1.207099in}%
\pgfsys@useobject{currentmarker}{}%
\end{pgfscope}%
\begin{pgfscope}%
\pgfsys@transformshift{1.273437in}{1.206532in}%
\pgfsys@useobject{currentmarker}{}%
\end{pgfscope}%
\begin{pgfscope}%
\pgfsys@transformshift{1.324942in}{1.205851in}%
\pgfsys@useobject{currentmarker}{}%
\end{pgfscope}%
\begin{pgfscope}%
\pgfsys@transformshift{1.376448in}{1.203905in}%
\pgfsys@useobject{currentmarker}{}%
\end{pgfscope}%
\begin{pgfscope}%
\pgfsys@transformshift{1.427954in}{1.201109in}%
\pgfsys@useobject{currentmarker}{}%
\end{pgfscope}%
\begin{pgfscope}%
\pgfsys@transformshift{1.479459in}{1.197650in}%
\pgfsys@useobject{currentmarker}{}%
\end{pgfscope}%
\begin{pgfscope}%
\pgfsys@transformshift{1.530965in}{1.194041in}%
\pgfsys@useobject{currentmarker}{}%
\end{pgfscope}%
\end{pgfscope}%
\begin{pgfscope}%
\pgfpathrectangle{\pgfqpoint{0.500851in}{0.496542in}}{\pgfqpoint{1.042991in}{1.239972in}}%
\pgfusepath{clip}%
\pgfsetbuttcap%
\pgfsetroundjoin%
\pgfsetlinewidth{0.903375pt}%
\definecolor{currentstroke}{rgb}{0.000000,0.000000,0.000000}%
\pgfsetstrokecolor{currentstroke}%
\pgfsetdash{{3.330000pt}{1.440000pt}}{0.000000pt}%
\pgfpathmoveto{\pgfqpoint{0.500851in}{1.229392in}}%
\pgfpathlineto{\pgfqpoint{0.521557in}{1.121930in}}%
\pgfpathlineto{\pgfqpoint{0.537086in}{1.049958in}}%
\pgfpathlineto{\pgfqpoint{0.552616in}{0.986813in}}%
\pgfpathlineto{\pgfqpoint{0.573321in}{0.914150in}}%
\pgfpathlineto{\pgfqpoint{0.588851in}{0.866434in}}%
\pgfpathlineto{\pgfqpoint{0.604380in}{0.824529in}}%
\pgfpathlineto{\pgfqpoint{0.625086in}{0.775908in}}%
\pgfpathlineto{\pgfqpoint{0.645792in}{0.734214in}}%
\pgfpathlineto{\pgfqpoint{0.666497in}{0.699115in}}%
\pgfpathlineto{\pgfqpoint{0.687203in}{0.668903in}}%
\pgfpathlineto{\pgfqpoint{0.707909in}{0.643332in}}%
\pgfpathlineto{\pgfqpoint{0.728615in}{0.621563in}}%
\pgfpathlineto{\pgfqpoint{0.749321in}{0.602907in}}%
\pgfpathlineto{\pgfqpoint{0.770027in}{0.587200in}}%
\pgfpathlineto{\pgfqpoint{0.795909in}{0.570622in}}%
\pgfpathlineto{\pgfqpoint{0.821791in}{0.557190in}}%
\pgfpathlineto{\pgfqpoint{0.847673in}{0.546099in}}%
\pgfpathlineto{\pgfqpoint{0.878732in}{0.535513in}}%
\pgfpathlineto{\pgfqpoint{0.909791in}{0.527137in}}%
\pgfpathlineto{\pgfqpoint{0.951202in}{0.518720in}}%
\pgfpathlineto{\pgfqpoint{0.992614in}{0.512628in}}%
\pgfpathlineto{\pgfqpoint{1.044379in}{0.507303in}}%
\pgfpathlineto{\pgfqpoint{1.106496in}{0.503182in}}%
\pgfpathlineto{\pgfqpoint{1.189319in}{0.500035in}}%
\pgfpathlineto{\pgfqpoint{1.308378in}{0.497927in}}%
\pgfpathlineto{\pgfqpoint{1.515436in}{0.496819in}}%
\pgfpathlineto{\pgfqpoint{1.530965in}{0.496788in}}%
\pgfpathlineto{\pgfqpoint{1.530965in}{0.496788in}}%
\pgfusepath{stroke}%
\end{pgfscope}%
\begin{pgfscope}%
\pgfpathrectangle{\pgfqpoint{0.500851in}{0.496542in}}{\pgfqpoint{1.042991in}{1.239972in}}%
\pgfusepath{clip}%
\pgfsetbuttcap%
\pgfsetroundjoin%
\pgfsetlinewidth{0.602250pt}%
\definecolor{currentstroke}{rgb}{0.788235,0.807843,0.819608}%
\pgfsetstrokecolor{currentstroke}%
\pgfsetdash{{0.600000pt}{0.990000pt}}{0.000000pt}%
\pgfpathmoveto{\pgfqpoint{0.758380in}{0.496542in}}%
\pgfpathlineto{\pgfqpoint{0.758380in}{1.736515in}}%
\pgfusepath{stroke}%
\end{pgfscope}%
\begin{pgfscope}%
\pgfsetrectcap%
\pgfsetmiterjoin%
\pgfsetlinewidth{0.803000pt}%
\definecolor{currentstroke}{rgb}{0.000000,0.000000,0.000000}%
\pgfsetstrokecolor{currentstroke}%
\pgfsetdash{}{0pt}%
\pgfpathmoveto{\pgfqpoint{0.500851in}{0.496542in}}%
\pgfpathlineto{\pgfqpoint{0.500851in}{1.736515in}}%
\pgfusepath{stroke}%
\end{pgfscope}%
\begin{pgfscope}%
\pgfsetrectcap%
\pgfsetmiterjoin%
\pgfsetlinewidth{0.803000pt}%
\definecolor{currentstroke}{rgb}{0.000000,0.000000,0.000000}%
\pgfsetstrokecolor{currentstroke}%
\pgfsetdash{}{0pt}%
\pgfpathmoveto{\pgfqpoint{1.543842in}{0.496542in}}%
\pgfpathlineto{\pgfqpoint{1.543842in}{1.736515in}}%
\pgfusepath{stroke}%
\end{pgfscope}%
\begin{pgfscope}%
\pgfsetrectcap%
\pgfsetmiterjoin%
\pgfsetlinewidth{0.803000pt}%
\definecolor{currentstroke}{rgb}{0.000000,0.000000,0.000000}%
\pgfsetstrokecolor{currentstroke}%
\pgfsetdash{}{0pt}%
\pgfpathmoveto{\pgfqpoint{0.500851in}{0.496542in}}%
\pgfpathlineto{\pgfqpoint{1.543842in}{0.496542in}}%
\pgfusepath{stroke}%
\end{pgfscope}%
\begin{pgfscope}%
\pgfsetrectcap%
\pgfsetmiterjoin%
\pgfsetlinewidth{0.803000pt}%
\definecolor{currentstroke}{rgb}{0.000000,0.000000,0.000000}%
\pgfsetstrokecolor{currentstroke}%
\pgfsetdash{}{0pt}%
\pgfpathmoveto{\pgfqpoint{0.500851in}{1.736515in}}%
\pgfpathlineto{\pgfqpoint{1.543842in}{1.736515in}}%
\pgfusepath{stroke}%
\end{pgfscope}%
\begin{pgfscope}%
\definecolor{textcolor}{rgb}{0.000000,0.000000,0.000000}%
\pgfsetstrokecolor{textcolor}%
\pgfsetfillcolor{textcolor}%
\pgftext[x=1.512552in,y=1.674516in,right,top]{\color{textcolor}{\rmfamily\fontsize{9.000000}{10.800000}\selectfont\catcode`\^=\active\def^{\ifmmode\sp\else\^{}\fi}\catcode`\%=\active\def%{\%}$(a)$}}%
\end{pgfscope}%
\begin{pgfscope}%
\pgfsetbuttcap%
\pgfsetmiterjoin%
\definecolor{currentfill}{rgb}{1.000000,1.000000,1.000000}%
\pgfsetfillcolor{currentfill}%
\pgfsetlinewidth{0.000000pt}%
\definecolor{currentstroke}{rgb}{0.000000,0.000000,0.000000}%
\pgfsetstrokecolor{currentstroke}%
\pgfsetstrokeopacity{0.000000}%
\pgfsetdash{}{0pt}%
\pgfpathmoveto{\pgfqpoint{2.137962in}{0.496542in}}%
\pgfpathlineto{\pgfqpoint{3.180953in}{0.496542in}}%
\pgfpathlineto{\pgfqpoint{3.180953in}{1.736515in}}%
\pgfpathlineto{\pgfqpoint{2.137962in}{1.736515in}}%
\pgfpathlineto{\pgfqpoint{2.137962in}{0.496542in}}%
\pgfpathclose%
\pgfusepath{fill}%
\end{pgfscope}%
\begin{pgfscope}%
\pgfsetbuttcap%
\pgfsetroundjoin%
\definecolor{currentfill}{rgb}{0.000000,0.000000,0.000000}%
\pgfsetfillcolor{currentfill}%
\pgfsetlinewidth{0.803000pt}%
\definecolor{currentstroke}{rgb}{0.000000,0.000000,0.000000}%
\pgfsetstrokecolor{currentstroke}%
\pgfsetdash{}{0pt}%
\pgfsys@defobject{currentmarker}{\pgfqpoint{0.000000in}{-0.048611in}}{\pgfqpoint{0.000000in}{0.000000in}}{%
\pgfpathmoveto{\pgfqpoint{0.000000in}{0.000000in}}%
\pgfpathlineto{\pgfqpoint{0.000000in}{-0.048611in}}%
\pgfusepath{stroke,fill}%
}%
\begin{pgfscope}%
\pgfsys@transformshift{2.137962in}{0.496542in}%
\pgfsys@useobject{currentmarker}{}%
\end{pgfscope}%
\end{pgfscope}%
\begin{pgfscope}%
\definecolor{textcolor}{rgb}{0.000000,0.000000,0.000000}%
\pgfsetstrokecolor{textcolor}%
\pgfsetfillcolor{textcolor}%
\pgftext[x=2.137962in,y=0.399320in,,top]{\color{textcolor}{\rmfamily\fontsize{8.000000}{9.600000}\selectfont\catcode`\^=\active\def^{\ifmmode\sp\else\^{}\fi}\catcode`\%=\active\def%{\%}0}}%
\end{pgfscope}%
\begin{pgfscope}%
\pgfsetbuttcap%
\pgfsetroundjoin%
\definecolor{currentfill}{rgb}{0.000000,0.000000,0.000000}%
\pgfsetfillcolor{currentfill}%
\pgfsetlinewidth{0.803000pt}%
\definecolor{currentstroke}{rgb}{0.000000,0.000000,0.000000}%
\pgfsetstrokecolor{currentstroke}%
\pgfsetdash{}{0pt}%
\pgfsys@defobject{currentmarker}{\pgfqpoint{0.000000in}{-0.048611in}}{\pgfqpoint{0.000000in}{0.000000in}}{%
\pgfpathmoveto{\pgfqpoint{0.000000in}{0.000000in}}%
\pgfpathlineto{\pgfqpoint{0.000000in}{-0.048611in}}%
\pgfusepath{stroke,fill}%
}%
\begin{pgfscope}%
\pgfsys@transformshift{2.653019in}{0.496542in}%
\pgfsys@useobject{currentmarker}{}%
\end{pgfscope}%
\end{pgfscope}%
\begin{pgfscope}%
\definecolor{textcolor}{rgb}{0.000000,0.000000,0.000000}%
\pgfsetstrokecolor{textcolor}%
\pgfsetfillcolor{textcolor}%
\pgftext[x=2.653019in,y=0.399320in,,top]{\color{textcolor}{\rmfamily\fontsize{8.000000}{9.600000}\selectfont\catcode`\^=\active\def^{\ifmmode\sp\else\^{}\fi}\catcode`\%=\active\def%{\%}4}}%
\end{pgfscope}%
\begin{pgfscope}%
\pgfsetbuttcap%
\pgfsetroundjoin%
\definecolor{currentfill}{rgb}{0.000000,0.000000,0.000000}%
\pgfsetfillcolor{currentfill}%
\pgfsetlinewidth{0.803000pt}%
\definecolor{currentstroke}{rgb}{0.000000,0.000000,0.000000}%
\pgfsetstrokecolor{currentstroke}%
\pgfsetdash{}{0pt}%
\pgfsys@defobject{currentmarker}{\pgfqpoint{0.000000in}{-0.048611in}}{\pgfqpoint{0.000000in}{0.000000in}}{%
\pgfpathmoveto{\pgfqpoint{0.000000in}{0.000000in}}%
\pgfpathlineto{\pgfqpoint{0.000000in}{-0.048611in}}%
\pgfusepath{stroke,fill}%
}%
\begin{pgfscope}%
\pgfsys@transformshift{3.168076in}{0.496542in}%
\pgfsys@useobject{currentmarker}{}%
\end{pgfscope}%
\end{pgfscope}%
\begin{pgfscope}%
\definecolor{textcolor}{rgb}{0.000000,0.000000,0.000000}%
\pgfsetstrokecolor{textcolor}%
\pgfsetfillcolor{textcolor}%
\pgftext[x=3.168076in,y=0.399320in,,top]{\color{textcolor}{\rmfamily\fontsize{8.000000}{9.600000}\selectfont\catcode`\^=\active\def^{\ifmmode\sp\else\^{}\fi}\catcode`\%=\active\def%{\%}8}}%
\end{pgfscope}%
\begin{pgfscope}%
\definecolor{textcolor}{rgb}{0.000000,0.000000,0.000000}%
\pgfsetstrokecolor{textcolor}%
\pgfsetfillcolor{textcolor}%
\pgftext[x=2.659457in,y=0.244999in,,top]{\color{textcolor}{\rmfamily\fontsize{9.000000}{10.800000}\selectfont\catcode`\^=\active\def^{\ifmmode\sp\else\^{}\fi}\catcode`\%=\active\def%{\%}$T\ (10^3/\omega_b)$}}%
\end{pgfscope}%
\begin{pgfscope}%
\pgfsetbuttcap%
\pgfsetroundjoin%
\definecolor{currentfill}{rgb}{0.000000,0.000000,0.000000}%
\pgfsetfillcolor{currentfill}%
\pgfsetlinewidth{0.803000pt}%
\definecolor{currentstroke}{rgb}{0.000000,0.000000,0.000000}%
\pgfsetstrokecolor{currentstroke}%
\pgfsetdash{}{0pt}%
\pgfsys@defobject{currentmarker}{\pgfqpoint{-0.048611in}{0.000000in}}{\pgfqpoint{-0.000000in}{0.000000in}}{%
\pgfpathmoveto{\pgfqpoint{-0.000000in}{0.000000in}}%
\pgfpathlineto{\pgfqpoint{-0.048611in}{0.000000in}}%
\pgfusepath{stroke,fill}%
}%
\begin{pgfscope}%
\pgfsys@transformshift{2.137962in}{0.496542in}%
\pgfsys@useobject{currentmarker}{}%
\end{pgfscope}%
\end{pgfscope}%
\begin{pgfscope}%
\definecolor{textcolor}{rgb}{0.000000,0.000000,0.000000}%
\pgfsetstrokecolor{textcolor}%
\pgfsetfillcolor{textcolor}%
\pgftext[x=1.981711in, y=0.457962in, left, base]{\color{textcolor}{\rmfamily\fontsize{8.000000}{9.600000}\selectfont\catcode`\^=\active\def^{\ifmmode\sp\else\^{}\fi}\catcode`\%=\active\def%{\%}$0$}}%
\end{pgfscope}%
\begin{pgfscope}%
\pgfsetbuttcap%
\pgfsetroundjoin%
\definecolor{currentfill}{rgb}{0.000000,0.000000,0.000000}%
\pgfsetfillcolor{currentfill}%
\pgfsetlinewidth{0.803000pt}%
\definecolor{currentstroke}{rgb}{0.000000,0.000000,0.000000}%
\pgfsetstrokecolor{currentstroke}%
\pgfsetdash{}{0pt}%
\pgfsys@defobject{currentmarker}{\pgfqpoint{-0.048611in}{0.000000in}}{\pgfqpoint{-0.000000in}{0.000000in}}{%
\pgfpathmoveto{\pgfqpoint{-0.000000in}{0.000000in}}%
\pgfpathlineto{\pgfqpoint{-0.048611in}{0.000000in}}%
\pgfusepath{stroke,fill}%
}%
\begin{pgfscope}%
\pgfsys@transformshift{2.137962in}{1.087005in}%
\pgfsys@useobject{currentmarker}{}%
\end{pgfscope}%
\end{pgfscope}%
\begin{pgfscope}%
\definecolor{textcolor}{rgb}{0.000000,0.000000,0.000000}%
\pgfsetstrokecolor{textcolor}%
\pgfsetfillcolor{textcolor}%
\pgftext[x=1.889889in, y=1.048425in, left, base]{\color{textcolor}{\rmfamily\fontsize{8.000000}{9.600000}\selectfont\catcode`\^=\active\def^{\ifmmode\sp\else\^{}\fi}\catcode`\%=\active\def%{\%}$0.5$}}%
\end{pgfscope}%
\begin{pgfscope}%
\pgfsetbuttcap%
\pgfsetroundjoin%
\definecolor{currentfill}{rgb}{0.000000,0.000000,0.000000}%
\pgfsetfillcolor{currentfill}%
\pgfsetlinewidth{0.803000pt}%
\definecolor{currentstroke}{rgb}{0.000000,0.000000,0.000000}%
\pgfsetstrokecolor{currentstroke}%
\pgfsetdash{}{0pt}%
\pgfsys@defobject{currentmarker}{\pgfqpoint{-0.048611in}{0.000000in}}{\pgfqpoint{-0.000000in}{0.000000in}}{%
\pgfpathmoveto{\pgfqpoint{-0.000000in}{0.000000in}}%
\pgfpathlineto{\pgfqpoint{-0.048611in}{0.000000in}}%
\pgfusepath{stroke,fill}%
}%
\begin{pgfscope}%
\pgfsys@transformshift{2.137962in}{1.677468in}%
\pgfsys@useobject{currentmarker}{}%
\end{pgfscope}%
\end{pgfscope}%
\begin{pgfscope}%
\definecolor{textcolor}{rgb}{0.000000,0.000000,0.000000}%
\pgfsetstrokecolor{textcolor}%
\pgfsetfillcolor{textcolor}%
\pgftext[x=1.981711in, y=1.638888in, left, base]{\color{textcolor}{\rmfamily\fontsize{8.000000}{9.600000}\selectfont\catcode`\^=\active\def^{\ifmmode\sp\else\^{}\fi}\catcode`\%=\active\def%{\%}$1$}}%
\end{pgfscope}%
\begin{pgfscope}%
\definecolor{textcolor}{rgb}{0.000000,0.000000,0.000000}%
\pgfsetstrokecolor{textcolor}%
\pgfsetfillcolor{textcolor}%
\pgftext[x=1.848222in,y=1.116529in,,bottom,rotate=90.000000]{\color{textcolor}{\rmfamily\fontsize{9.000000}{10.800000}\selectfont\catcode`\^=\active\def^{\ifmmode\sp\else\^{}\fi}\catcode`\%=\active\def%{\%}$s$}}%
\end{pgfscope}%
\begin{pgfscope}%
\pgfpathrectangle{\pgfqpoint{2.137962in}{0.496542in}}{\pgfqpoint{1.042991in}{1.239972in}}%
\pgfusepath{clip}%
\pgfsetrectcap%
\pgfsetroundjoin%
\pgfsetlinewidth{0.903375pt}%
\definecolor{currentstroke}{rgb}{1.000000,0.000000,0.000000}%
\pgfsetstrokecolor{currentstroke}%
\pgfsetdash{}{0pt}%
\pgfpathmoveto{\pgfqpoint{2.137962in}{1.311439in}}%
\pgfpathlineto{\pgfqpoint{2.158668in}{1.325931in}}%
\pgfpathlineto{\pgfqpoint{2.174197in}{1.334444in}}%
\pgfpathlineto{\pgfqpoint{2.189727in}{1.339819in}}%
\pgfpathlineto{\pgfqpoint{2.220785in}{1.345738in}}%
\pgfpathlineto{\pgfqpoint{2.241491in}{1.347159in}}%
\pgfpathlineto{\pgfqpoint{2.251844in}{1.346146in}}%
\pgfpathlineto{\pgfqpoint{2.293256in}{1.338208in}}%
\pgfpathlineto{\pgfqpoint{2.365726in}{1.338865in}}%
\pgfpathlineto{\pgfqpoint{2.448549in}{1.343683in}}%
\pgfpathlineto{\pgfqpoint{2.500314in}{1.342313in}}%
\pgfpathlineto{\pgfqpoint{2.567608in}{1.341020in}}%
\pgfpathlineto{\pgfqpoint{2.665960in}{1.342077in}}%
\pgfpathlineto{\pgfqpoint{2.955842in}{1.339373in}}%
\pgfpathlineto{\pgfqpoint{3.168076in}{1.337375in}}%
\pgfpathlineto{\pgfqpoint{3.168076in}{1.337375in}}%
\pgfusepath{stroke}%
\end{pgfscope}%
\begin{pgfscope}%
\pgfpathrectangle{\pgfqpoint{2.137962in}{0.496542in}}{\pgfqpoint{1.042991in}{1.239972in}}%
\pgfusepath{clip}%
\pgfsetbuttcap%
\pgfsetroundjoin%
\definecolor{currentfill}{rgb}{0.000000,0.000000,0.000000}%
\pgfsetfillcolor{currentfill}%
\pgfsetfillopacity{0.000000}%
\pgfsetlinewidth{0.803000pt}%
\definecolor{currentstroke}{rgb}{1.000000,0.000000,0.000000}%
\pgfsetstrokecolor{currentstroke}%
\pgfsetdash{}{0pt}%
\pgfsys@defobject{currentmarker}{\pgfqpoint{-0.017500in}{-0.017500in}}{\pgfqpoint{0.017500in}{0.017500in}}{%
\pgfpathmoveto{\pgfqpoint{0.000000in}{-0.017500in}}%
\pgfpathcurveto{\pgfqpoint{0.004641in}{-0.017500in}}{\pgfqpoint{0.009093in}{-0.015656in}}{\pgfqpoint{0.012374in}{-0.012374in}}%
\pgfpathcurveto{\pgfqpoint{0.015656in}{-0.009093in}}{\pgfqpoint{0.017500in}{-0.004641in}}{\pgfqpoint{0.017500in}{0.000000in}}%
\pgfpathcurveto{\pgfqpoint{0.017500in}{0.004641in}}{\pgfqpoint{0.015656in}{0.009093in}}{\pgfqpoint{0.012374in}{0.012374in}}%
\pgfpathcurveto{\pgfqpoint{0.009093in}{0.015656in}}{\pgfqpoint{0.004641in}{0.017500in}}{\pgfqpoint{0.000000in}{0.017500in}}%
\pgfpathcurveto{\pgfqpoint{-0.004641in}{0.017500in}}{\pgfqpoint{-0.009093in}{0.015656in}}{\pgfqpoint{-0.012374in}{0.012374in}}%
\pgfpathcurveto{\pgfqpoint{-0.015656in}{0.009093in}}{\pgfqpoint{-0.017500in}{0.004641in}}{\pgfqpoint{-0.017500in}{0.000000in}}%
\pgfpathcurveto{\pgfqpoint{-0.017500in}{-0.004641in}}{\pgfqpoint{-0.015656in}{-0.009093in}}{\pgfqpoint{-0.012374in}{-0.012374in}}%
\pgfpathcurveto{\pgfqpoint{-0.009093in}{-0.015656in}}{\pgfqpoint{-0.004641in}{-0.017500in}}{\pgfqpoint{0.000000in}{-0.017500in}}%
\pgfpathlineto{\pgfqpoint{0.000000in}{-0.017500in}}%
\pgfpathclose%
\pgfusepath{stroke,fill}%
}%
\begin{pgfscope}%
\pgfsys@transformshift{2.137962in}{1.311439in}%
\pgfsys@useobject{currentmarker}{}%
\end{pgfscope}%
\begin{pgfscope}%
\pgfsys@transformshift{2.189468in}{1.339760in}%
\pgfsys@useobject{currentmarker}{}%
\end{pgfscope}%
\begin{pgfscope}%
\pgfsys@transformshift{2.240973in}{1.347162in}%
\pgfsys@useobject{currentmarker}{}%
\end{pgfscope}%
\begin{pgfscope}%
\pgfsys@transformshift{2.292479in}{1.338213in}%
\pgfsys@useobject{currentmarker}{}%
\end{pgfscope}%
\begin{pgfscope}%
\pgfsys@transformshift{2.343985in}{1.338040in}%
\pgfsys@useobject{currentmarker}{}%
\end{pgfscope}%
\begin{pgfscope}%
\pgfsys@transformshift{2.395491in}{1.340944in}%
\pgfsys@useobject{currentmarker}{}%
\end{pgfscope}%
\begin{pgfscope}%
\pgfsys@transformshift{2.446996in}{1.343685in}%
\pgfsys@useobject{currentmarker}{}%
\end{pgfscope}%
\begin{pgfscope}%
\pgfsys@transformshift{2.498502in}{1.342365in}%
\pgfsys@useobject{currentmarker}{}%
\end{pgfscope}%
\begin{pgfscope}%
\pgfsys@transformshift{2.550008in}{1.340823in}%
\pgfsys@useobject{currentmarker}{}%
\end{pgfscope}%
\begin{pgfscope}%
\pgfsys@transformshift{2.601513in}{1.341719in}%
\pgfsys@useobject{currentmarker}{}%
\end{pgfscope}%
\begin{pgfscope}%
\pgfsys@transformshift{2.653019in}{1.342157in}%
\pgfsys@useobject{currentmarker}{}%
\end{pgfscope}%
\begin{pgfscope}%
\pgfsys@transformshift{2.704525in}{1.341434in}%
\pgfsys@useobject{currentmarker}{}%
\end{pgfscope}%
\begin{pgfscope}%
\pgfsys@transformshift{2.756031in}{1.340708in}%
\pgfsys@useobject{currentmarker}{}%
\end{pgfscope}%
\begin{pgfscope}%
\pgfsys@transformshift{2.807536in}{1.340851in}%
\pgfsys@useobject{currentmarker}{}%
\end{pgfscope}%
\begin{pgfscope}%
\pgfsys@transformshift{2.859042in}{1.340350in}%
\pgfsys@useobject{currentmarker}{}%
\end{pgfscope}%
\begin{pgfscope}%
\pgfsys@transformshift{2.910548in}{1.339591in}%
\pgfsys@useobject{currentmarker}{}%
\end{pgfscope}%
\begin{pgfscope}%
\pgfsys@transformshift{2.962053in}{1.339333in}%
\pgfsys@useobject{currentmarker}{}%
\end{pgfscope}%
\begin{pgfscope}%
\pgfsys@transformshift{3.013559in}{1.338714in}%
\pgfsys@useobject{currentmarker}{}%
\end{pgfscope}%
\begin{pgfscope}%
\pgfsys@transformshift{3.065065in}{1.338223in}%
\pgfsys@useobject{currentmarker}{}%
\end{pgfscope}%
\begin{pgfscope}%
\pgfsys@transformshift{3.116571in}{1.337808in}%
\pgfsys@useobject{currentmarker}{}%
\end{pgfscope}%
\begin{pgfscope}%
\pgfsys@transformshift{3.168076in}{1.337375in}%
\pgfsys@useobject{currentmarker}{}%
\end{pgfscope}%
\end{pgfscope}%
\begin{pgfscope}%
\pgfpathrectangle{\pgfqpoint{2.137962in}{0.496542in}}{\pgfqpoint{1.042991in}{1.239972in}}%
\pgfusepath{clip}%
\pgfsetbuttcap%
\pgfsetroundjoin%
\pgfsetlinewidth{0.903375pt}%
\definecolor{currentstroke}{rgb}{0.000000,0.000000,0.000000}%
\pgfsetstrokecolor{currentstroke}%
\pgfsetdash{{3.330000pt}{1.440000pt}}{0.000000pt}%
\pgfpathmoveto{\pgfqpoint{2.137962in}{1.265120in}}%
\pgfpathlineto{\pgfqpoint{2.158668in}{1.152419in}}%
\pgfpathlineto{\pgfqpoint{2.174197in}{1.076938in}}%
\pgfpathlineto{\pgfqpoint{2.189727in}{1.010714in}}%
\pgfpathlineto{\pgfqpoint{2.210432in}{0.934510in}}%
\pgfpathlineto{\pgfqpoint{2.225962in}{0.884467in}}%
\pgfpathlineto{\pgfqpoint{2.241491in}{0.840519in}}%
\pgfpathlineto{\pgfqpoint{2.262197in}{0.789528in}}%
\pgfpathlineto{\pgfqpoint{2.282903in}{0.745801in}}%
\pgfpathlineto{\pgfqpoint{2.303609in}{0.708991in}}%
\pgfpathlineto{\pgfqpoint{2.324314in}{0.677306in}}%
\pgfpathlineto{\pgfqpoint{2.345020in}{0.650488in}}%
\pgfpathlineto{\pgfqpoint{2.365726in}{0.627658in}}%
\pgfpathlineto{\pgfqpoint{2.386432in}{0.608092in}}%
\pgfpathlineto{\pgfqpoint{2.407138in}{0.591619in}}%
\pgfpathlineto{\pgfqpoint{2.433020in}{0.574233in}}%
\pgfpathlineto{\pgfqpoint{2.458902in}{0.560146in}}%
\pgfpathlineto{\pgfqpoint{2.484784in}{0.548515in}}%
\pgfpathlineto{\pgfqpoint{2.510667in}{0.539092in}}%
\pgfpathlineto{\pgfqpoint{2.541725in}{0.529939in}}%
\pgfpathlineto{\pgfqpoint{2.577961in}{0.521760in}}%
\pgfpathlineto{\pgfqpoint{2.619372in}{0.514833in}}%
\pgfpathlineto{\pgfqpoint{2.665960in}{0.509281in}}%
\pgfpathlineto{\pgfqpoint{2.722901in}{0.504728in}}%
\pgfpathlineto{\pgfqpoint{2.795372in}{0.501201in}}%
\pgfpathlineto{\pgfqpoint{2.898901in}{0.498627in}}%
\pgfpathlineto{\pgfqpoint{3.059371in}{0.497142in}}%
\pgfpathlineto{\pgfqpoint{3.168076in}{0.496800in}}%
\pgfpathlineto{\pgfqpoint{3.168076in}{0.496800in}}%
\pgfusepath{stroke}%
\end{pgfscope}%
\begin{pgfscope}%
\pgfpathrectangle{\pgfqpoint{2.137962in}{0.496542in}}{\pgfqpoint{1.042991in}{1.239972in}}%
\pgfusepath{clip}%
\pgfsetbuttcap%
\pgfsetroundjoin%
\pgfsetlinewidth{0.602250pt}%
\definecolor{currentstroke}{rgb}{0.788235,0.807843,0.819608}%
\pgfsetstrokecolor{currentstroke}%
\pgfsetdash{{0.600000pt}{0.990000pt}}{0.000000pt}%
\pgfpathmoveto{\pgfqpoint{2.395491in}{0.496542in}}%
\pgfpathlineto{\pgfqpoint{2.395491in}{1.736515in}}%
\pgfusepath{stroke}%
\end{pgfscope}%
\begin{pgfscope}%
\pgfsetrectcap%
\pgfsetmiterjoin%
\pgfsetlinewidth{0.803000pt}%
\definecolor{currentstroke}{rgb}{0.000000,0.000000,0.000000}%
\pgfsetstrokecolor{currentstroke}%
\pgfsetdash{}{0pt}%
\pgfpathmoveto{\pgfqpoint{2.137962in}{0.496542in}}%
\pgfpathlineto{\pgfqpoint{2.137962in}{1.736515in}}%
\pgfusepath{stroke}%
\end{pgfscope}%
\begin{pgfscope}%
\pgfsetrectcap%
\pgfsetmiterjoin%
\pgfsetlinewidth{0.803000pt}%
\definecolor{currentstroke}{rgb}{0.000000,0.000000,0.000000}%
\pgfsetstrokecolor{currentstroke}%
\pgfsetdash{}{0pt}%
\pgfpathmoveto{\pgfqpoint{3.180953in}{0.496542in}}%
\pgfpathlineto{\pgfqpoint{3.180953in}{1.736515in}}%
\pgfusepath{stroke}%
\end{pgfscope}%
\begin{pgfscope}%
\pgfsetrectcap%
\pgfsetmiterjoin%
\pgfsetlinewidth{0.803000pt}%
\definecolor{currentstroke}{rgb}{0.000000,0.000000,0.000000}%
\pgfsetstrokecolor{currentstroke}%
\pgfsetdash{}{0pt}%
\pgfpathmoveto{\pgfqpoint{2.137962in}{0.496542in}}%
\pgfpathlineto{\pgfqpoint{3.180953in}{0.496542in}}%
\pgfusepath{stroke}%
\end{pgfscope}%
\begin{pgfscope}%
\pgfsetrectcap%
\pgfsetmiterjoin%
\pgfsetlinewidth{0.803000pt}%
\definecolor{currentstroke}{rgb}{0.000000,0.000000,0.000000}%
\pgfsetstrokecolor{currentstroke}%
\pgfsetdash{}{0pt}%
\pgfpathmoveto{\pgfqpoint{2.137962in}{1.736515in}}%
\pgfpathlineto{\pgfqpoint{3.180953in}{1.736515in}}%
\pgfusepath{stroke}%
\end{pgfscope}%
\begin{pgfscope}%
\definecolor{textcolor}{rgb}{0.000000,0.000000,0.000000}%
\pgfsetstrokecolor{textcolor}%
\pgfsetfillcolor{textcolor}%
\pgftext[x=3.149663in,y=1.674516in,right,top]{\color{textcolor}{\rmfamily\fontsize{9.000000}{10.800000}\selectfont\catcode`\^=\active\def^{\ifmmode\sp\else\^{}\fi}\catcode`\%=\active\def%{\%}$(b)$}}%
\end{pgfscope}%
\begin{pgfscope}%
\pgfsetbuttcap%
\pgfsetmiterjoin%
\definecolor{currentfill}{rgb}{1.000000,1.000000,1.000000}%
\pgfsetfillcolor{currentfill}%
\pgfsetlinewidth{0.000000pt}%
\definecolor{currentstroke}{rgb}{0.000000,0.000000,0.000000}%
\pgfsetstrokecolor{currentstroke}%
\pgfsetstrokeopacity{0.000000}%
\pgfsetdash{}{0pt}%
\pgfpathmoveto{\pgfqpoint{3.775073in}{0.496542in}}%
\pgfpathlineto{\pgfqpoint{4.818064in}{0.496542in}}%
\pgfpathlineto{\pgfqpoint{4.818064in}{1.736515in}}%
\pgfpathlineto{\pgfqpoint{3.775073in}{1.736515in}}%
\pgfpathlineto{\pgfqpoint{3.775073in}{0.496542in}}%
\pgfpathclose%
\pgfusepath{fill}%
\end{pgfscope}%
\begin{pgfscope}%
\pgfsetbuttcap%
\pgfsetroundjoin%
\definecolor{currentfill}{rgb}{0.000000,0.000000,0.000000}%
\pgfsetfillcolor{currentfill}%
\pgfsetlinewidth{0.803000pt}%
\definecolor{currentstroke}{rgb}{0.000000,0.000000,0.000000}%
\pgfsetstrokecolor{currentstroke}%
\pgfsetdash{}{0pt}%
\pgfsys@defobject{currentmarker}{\pgfqpoint{0.000000in}{-0.048611in}}{\pgfqpoint{0.000000in}{0.000000in}}{%
\pgfpathmoveto{\pgfqpoint{0.000000in}{0.000000in}}%
\pgfpathlineto{\pgfqpoint{0.000000in}{-0.048611in}}%
\pgfusepath{stroke,fill}%
}%
\begin{pgfscope}%
\pgfsys@transformshift{3.775073in}{0.496542in}%
\pgfsys@useobject{currentmarker}{}%
\end{pgfscope}%
\end{pgfscope}%
\begin{pgfscope}%
\definecolor{textcolor}{rgb}{0.000000,0.000000,0.000000}%
\pgfsetstrokecolor{textcolor}%
\pgfsetfillcolor{textcolor}%
\pgftext[x=3.775073in,y=0.399320in,,top]{\color{textcolor}{\rmfamily\fontsize{8.000000}{9.600000}\selectfont\catcode`\^=\active\def^{\ifmmode\sp\else\^{}\fi}\catcode`\%=\active\def%{\%}0}}%
\end{pgfscope}%
\begin{pgfscope}%
\pgfsetbuttcap%
\pgfsetroundjoin%
\definecolor{currentfill}{rgb}{0.000000,0.000000,0.000000}%
\pgfsetfillcolor{currentfill}%
\pgfsetlinewidth{0.803000pt}%
\definecolor{currentstroke}{rgb}{0.000000,0.000000,0.000000}%
\pgfsetstrokecolor{currentstroke}%
\pgfsetdash{}{0pt}%
\pgfsys@defobject{currentmarker}{\pgfqpoint{0.000000in}{-0.048611in}}{\pgfqpoint{0.000000in}{0.000000in}}{%
\pgfpathmoveto{\pgfqpoint{0.000000in}{0.000000in}}%
\pgfpathlineto{\pgfqpoint{0.000000in}{-0.048611in}}%
\pgfusepath{stroke,fill}%
}%
\begin{pgfscope}%
\pgfsys@transformshift{4.290130in}{0.496542in}%
\pgfsys@useobject{currentmarker}{}%
\end{pgfscope}%
\end{pgfscope}%
\begin{pgfscope}%
\definecolor{textcolor}{rgb}{0.000000,0.000000,0.000000}%
\pgfsetstrokecolor{textcolor}%
\pgfsetfillcolor{textcolor}%
\pgftext[x=4.290130in,y=0.399320in,,top]{\color{textcolor}{\rmfamily\fontsize{8.000000}{9.600000}\selectfont\catcode`\^=\active\def^{\ifmmode\sp\else\^{}\fi}\catcode`\%=\active\def%{\%}4}}%
\end{pgfscope}%
\begin{pgfscope}%
\pgfsetbuttcap%
\pgfsetroundjoin%
\definecolor{currentfill}{rgb}{0.000000,0.000000,0.000000}%
\pgfsetfillcolor{currentfill}%
\pgfsetlinewidth{0.803000pt}%
\definecolor{currentstroke}{rgb}{0.000000,0.000000,0.000000}%
\pgfsetstrokecolor{currentstroke}%
\pgfsetdash{}{0pt}%
\pgfsys@defobject{currentmarker}{\pgfqpoint{0.000000in}{-0.048611in}}{\pgfqpoint{0.000000in}{0.000000in}}{%
\pgfpathmoveto{\pgfqpoint{0.000000in}{0.000000in}}%
\pgfpathlineto{\pgfqpoint{0.000000in}{-0.048611in}}%
\pgfusepath{stroke,fill}%
}%
\begin{pgfscope}%
\pgfsys@transformshift{4.805187in}{0.496542in}%
\pgfsys@useobject{currentmarker}{}%
\end{pgfscope}%
\end{pgfscope}%
\begin{pgfscope}%
\definecolor{textcolor}{rgb}{0.000000,0.000000,0.000000}%
\pgfsetstrokecolor{textcolor}%
\pgfsetfillcolor{textcolor}%
\pgftext[x=4.805187in,y=0.399320in,,top]{\color{textcolor}{\rmfamily\fontsize{8.000000}{9.600000}\selectfont\catcode`\^=\active\def^{\ifmmode\sp\else\^{}\fi}\catcode`\%=\active\def%{\%}8}}%
\end{pgfscope}%
\begin{pgfscope}%
\pgfsetbuttcap%
\pgfsetroundjoin%
\definecolor{currentfill}{rgb}{0.000000,0.000000,0.000000}%
\pgfsetfillcolor{currentfill}%
\pgfsetlinewidth{0.803000pt}%
\definecolor{currentstroke}{rgb}{0.000000,0.000000,0.000000}%
\pgfsetstrokecolor{currentstroke}%
\pgfsetdash{}{0pt}%
\pgfsys@defobject{currentmarker}{\pgfqpoint{-0.048611in}{0.000000in}}{\pgfqpoint{-0.000000in}{0.000000in}}{%
\pgfpathmoveto{\pgfqpoint{-0.000000in}{0.000000in}}%
\pgfpathlineto{\pgfqpoint{-0.048611in}{0.000000in}}%
\pgfusepath{stroke,fill}%
}%
\begin{pgfscope}%
\pgfsys@transformshift{3.775073in}{0.496542in}%
\pgfsys@useobject{currentmarker}{}%
\end{pgfscope}%
\end{pgfscope}%
\begin{pgfscope}%
\definecolor{textcolor}{rgb}{0.000000,0.000000,0.000000}%
\pgfsetstrokecolor{textcolor}%
\pgfsetfillcolor{textcolor}%
\pgftext[x=3.618822in, y=0.457962in, left, base]{\color{textcolor}{\rmfamily\fontsize{8.000000}{9.600000}\selectfont\catcode`\^=\active\def^{\ifmmode\sp\else\^{}\fi}\catcode`\%=\active\def%{\%}$0$}}%
\end{pgfscope}%
\begin{pgfscope}%
\pgfsetbuttcap%
\pgfsetroundjoin%
\definecolor{currentfill}{rgb}{0.000000,0.000000,0.000000}%
\pgfsetfillcolor{currentfill}%
\pgfsetlinewidth{0.803000pt}%
\definecolor{currentstroke}{rgb}{0.000000,0.000000,0.000000}%
\pgfsetstrokecolor{currentstroke}%
\pgfsetdash{}{0pt}%
\pgfsys@defobject{currentmarker}{\pgfqpoint{-0.048611in}{0.000000in}}{\pgfqpoint{-0.000000in}{0.000000in}}{%
\pgfpathmoveto{\pgfqpoint{-0.000000in}{0.000000in}}%
\pgfpathlineto{\pgfqpoint{-0.048611in}{0.000000in}}%
\pgfusepath{stroke,fill}%
}%
\begin{pgfscope}%
\pgfsys@transformshift{3.775073in}{1.013198in}%
\pgfsys@useobject{currentmarker}{}%
\end{pgfscope}%
\end{pgfscope}%
\begin{pgfscope}%
\definecolor{textcolor}{rgb}{0.000000,0.000000,0.000000}%
\pgfsetstrokecolor{textcolor}%
\pgfsetfillcolor{textcolor}%
\pgftext[x=3.527000in, y=0.974617in, left, base]{\color{textcolor}{\rmfamily\fontsize{8.000000}{9.600000}\selectfont\catcode`\^=\active\def^{\ifmmode\sp\else\^{}\fi}\catcode`\%=\active\def%{\%}$0.5$}}%
\end{pgfscope}%
\begin{pgfscope}%
\pgfsetbuttcap%
\pgfsetroundjoin%
\definecolor{currentfill}{rgb}{0.000000,0.000000,0.000000}%
\pgfsetfillcolor{currentfill}%
\pgfsetlinewidth{0.803000pt}%
\definecolor{currentstroke}{rgb}{0.000000,0.000000,0.000000}%
\pgfsetstrokecolor{currentstroke}%
\pgfsetdash{}{0pt}%
\pgfsys@defobject{currentmarker}{\pgfqpoint{-0.048611in}{0.000000in}}{\pgfqpoint{-0.000000in}{0.000000in}}{%
\pgfpathmoveto{\pgfqpoint{-0.000000in}{0.000000in}}%
\pgfpathlineto{\pgfqpoint{-0.048611in}{0.000000in}}%
\pgfusepath{stroke,fill}%
}%
\begin{pgfscope}%
\pgfsys@transformshift{3.775073in}{1.529853in}%
\pgfsys@useobject{currentmarker}{}%
\end{pgfscope}%
\end{pgfscope}%
\begin{pgfscope}%
\definecolor{textcolor}{rgb}{0.000000,0.000000,0.000000}%
\pgfsetstrokecolor{textcolor}%
\pgfsetfillcolor{textcolor}%
\pgftext[x=3.618822in, y=1.491272in, left, base]{\color{textcolor}{\rmfamily\fontsize{8.000000}{9.600000}\selectfont\catcode`\^=\active\def^{\ifmmode\sp\else\^{}\fi}\catcode`\%=\active\def%{\%}$1$}}%
\end{pgfscope}%
\begin{pgfscope}%
\definecolor{textcolor}{rgb}{0.000000,0.000000,0.000000}%
\pgfsetstrokecolor{textcolor}%
\pgfsetfillcolor{textcolor}%
\pgftext[x=3.485333in,y=1.116529in,,bottom,rotate=90.000000]{\color{textcolor}{\rmfamily\fontsize{9.000000}{10.800000}\selectfont\catcode`\^=\active\def^{\ifmmode\sp\else\^{}\fi}\catcode`\%=\active\def%{\%}$F_h$}}%
\end{pgfscope}%
\begin{pgfscope}%
\pgfpathrectangle{\pgfqpoint{3.775073in}{0.496542in}}{\pgfqpoint{1.042991in}{1.239972in}}%
\pgfusepath{clip}%
\pgfsetrectcap%
\pgfsetroundjoin%
\pgfsetlinewidth{0.903375pt}%
\definecolor{currentstroke}{rgb}{1.000000,0.000000,0.000000}%
\pgfsetstrokecolor{currentstroke}%
\pgfsetdash{}{0pt}%
\pgfpathmoveto{\pgfqpoint{3.775073in}{1.400703in}}%
\pgfpathlineto{\pgfqpoint{3.961426in}{1.389761in}}%
\pgfpathlineto{\pgfqpoint{4.127072in}{1.385782in}}%
\pgfpathlineto{\pgfqpoint{4.220248in}{1.378991in}}%
\pgfpathlineto{\pgfqpoint{4.365189in}{1.368640in}}%
\pgfpathlineto{\pgfqpoint{4.427306in}{1.366909in}}%
\pgfpathlineto{\pgfqpoint{4.520482in}{1.366717in}}%
\pgfpathlineto{\pgfqpoint{4.613659in}{1.365858in}}%
\pgfpathlineto{\pgfqpoint{4.686129in}{1.362540in}}%
\pgfpathlineto{\pgfqpoint{4.784482in}{1.355333in}}%
\pgfpathlineto{\pgfqpoint{4.805187in}{1.353708in}}%
\pgfpathlineto{\pgfqpoint{4.805187in}{1.353708in}}%
\pgfusepath{stroke}%
\end{pgfscope}%
\begin{pgfscope}%
\pgfpathrectangle{\pgfqpoint{3.775073in}{0.496542in}}{\pgfqpoint{1.042991in}{1.239972in}}%
\pgfusepath{clip}%
\pgfsetbuttcap%
\pgfsetroundjoin%
\definecolor{currentfill}{rgb}{0.000000,0.000000,0.000000}%
\pgfsetfillcolor{currentfill}%
\pgfsetfillopacity{0.000000}%
\pgfsetlinewidth{0.803000pt}%
\definecolor{currentstroke}{rgb}{1.000000,0.000000,0.000000}%
\pgfsetstrokecolor{currentstroke}%
\pgfsetdash{}{0pt}%
\pgfsys@defobject{currentmarker}{\pgfqpoint{-0.017500in}{-0.017500in}}{\pgfqpoint{0.017500in}{0.017500in}}{%
\pgfpathmoveto{\pgfqpoint{0.000000in}{-0.017500in}}%
\pgfpathcurveto{\pgfqpoint{0.004641in}{-0.017500in}}{\pgfqpoint{0.009093in}{-0.015656in}}{\pgfqpoint{0.012374in}{-0.012374in}}%
\pgfpathcurveto{\pgfqpoint{0.015656in}{-0.009093in}}{\pgfqpoint{0.017500in}{-0.004641in}}{\pgfqpoint{0.017500in}{0.000000in}}%
\pgfpathcurveto{\pgfqpoint{0.017500in}{0.004641in}}{\pgfqpoint{0.015656in}{0.009093in}}{\pgfqpoint{0.012374in}{0.012374in}}%
\pgfpathcurveto{\pgfqpoint{0.009093in}{0.015656in}}{\pgfqpoint{0.004641in}{0.017500in}}{\pgfqpoint{0.000000in}{0.017500in}}%
\pgfpathcurveto{\pgfqpoint{-0.004641in}{0.017500in}}{\pgfqpoint{-0.009093in}{0.015656in}}{\pgfqpoint{-0.012374in}{0.012374in}}%
\pgfpathcurveto{\pgfqpoint{-0.015656in}{0.009093in}}{\pgfqpoint{-0.017500in}{0.004641in}}{\pgfqpoint{-0.017500in}{0.000000in}}%
\pgfpathcurveto{\pgfqpoint{-0.017500in}{-0.004641in}}{\pgfqpoint{-0.015656in}{-0.009093in}}{\pgfqpoint{-0.012374in}{-0.012374in}}%
\pgfpathcurveto{\pgfqpoint{-0.009093in}{-0.015656in}}{\pgfqpoint{-0.004641in}{-0.017500in}}{\pgfqpoint{0.000000in}{-0.017500in}}%
\pgfpathlineto{\pgfqpoint{0.000000in}{-0.017500in}}%
\pgfpathclose%
\pgfusepath{stroke,fill}%
}%
\begin{pgfscope}%
\pgfsys@transformshift{3.775073in}{1.400703in}%
\pgfsys@useobject{currentmarker}{}%
\end{pgfscope}%
\begin{pgfscope}%
\pgfsys@transformshift{3.826579in}{1.397194in}%
\pgfsys@useobject{currentmarker}{}%
\end{pgfscope}%
\begin{pgfscope}%
\pgfsys@transformshift{3.878085in}{1.393232in}%
\pgfsys@useobject{currentmarker}{}%
\end{pgfscope}%
\begin{pgfscope}%
\pgfsys@transformshift{3.929590in}{1.390527in}%
\pgfsys@useobject{currentmarker}{}%
\end{pgfscope}%
\begin{pgfscope}%
\pgfsys@transformshift{3.981096in}{1.389433in}%
\pgfsys@useobject{currentmarker}{}%
\end{pgfscope}%
\begin{pgfscope}%
\pgfsys@transformshift{4.032602in}{1.388838in}%
\pgfsys@useobject{currentmarker}{}%
\end{pgfscope}%
\begin{pgfscope}%
\pgfsys@transformshift{4.084107in}{1.387920in}%
\pgfsys@useobject{currentmarker}{}%
\end{pgfscope}%
\begin{pgfscope}%
\pgfsys@transformshift{4.135613in}{1.385260in}%
\pgfsys@useobject{currentmarker}{}%
\end{pgfscope}%
\begin{pgfscope}%
\pgfsys@transformshift{4.187119in}{1.381541in}%
\pgfsys@useobject{currentmarker}{}%
\end{pgfscope}%
\begin{pgfscope}%
\pgfsys@transformshift{4.238625in}{1.377538in}%
\pgfsys@useobject{currentmarker}{}%
\end{pgfscope}%
\begin{pgfscope}%
\pgfsys@transformshift{4.290130in}{1.373354in}%
\pgfsys@useobject{currentmarker}{}%
\end{pgfscope}%
\begin{pgfscope}%
\pgfsys@transformshift{4.341636in}{1.369817in}%
\pgfsys@useobject{currentmarker}{}%
\end{pgfscope}%
\begin{pgfscope}%
\pgfsys@transformshift{4.393142in}{1.367627in}%
\pgfsys@useobject{currentmarker}{}%
\end{pgfscope}%
\begin{pgfscope}%
\pgfsys@transformshift{4.444647in}{1.366751in}%
\pgfsys@useobject{currentmarker}{}%
\end{pgfscope}%
\begin{pgfscope}%
\pgfsys@transformshift{4.496153in}{1.366677in}%
\pgfsys@useobject{currentmarker}{}%
\end{pgfscope}%
\begin{pgfscope}%
\pgfsys@transformshift{4.547659in}{1.366765in}%
\pgfsys@useobject{currentmarker}{}%
\end{pgfscope}%
\begin{pgfscope}%
\pgfsys@transformshift{4.599165in}{1.366196in}%
\pgfsys@useobject{currentmarker}{}%
\end{pgfscope}%
\begin{pgfscope}%
\pgfsys@transformshift{4.650670in}{1.364447in}%
\pgfsys@useobject{currentmarker}{}%
\end{pgfscope}%
\begin{pgfscope}%
\pgfsys@transformshift{4.702176in}{1.361521in}%
\pgfsys@useobject{currentmarker}{}%
\end{pgfscope}%
\begin{pgfscope}%
\pgfsys@transformshift{4.753682in}{1.357699in}%
\pgfsys@useobject{currentmarker}{}%
\end{pgfscope}%
\begin{pgfscope}%
\pgfsys@transformshift{4.805187in}{1.353708in}%
\pgfsys@useobject{currentmarker}{}%
\end{pgfscope}%
\end{pgfscope}%
\begin{pgfscope}%
\pgfpathrectangle{\pgfqpoint{3.775073in}{0.496542in}}{\pgfqpoint{1.042991in}{1.239972in}}%
\pgfusepath{clip}%
\pgfsetbuttcap%
\pgfsetroundjoin%
\pgfsetlinewidth{0.903375pt}%
\definecolor{currentstroke}{rgb}{0.000000,0.000000,0.000000}%
\pgfsetstrokecolor{currentstroke}%
\pgfsetdash{{3.330000pt}{1.440000pt}}{0.000000pt}%
\pgfpathmoveto{\pgfqpoint{3.775073in}{1.481819in}}%
\pgfpathlineto{\pgfqpoint{4.805187in}{1.481817in}}%
\pgfpathlineto{\pgfqpoint{4.805187in}{1.481817in}}%
\pgfusepath{stroke}%
\end{pgfscope}%
\begin{pgfscope}%
\pgfpathrectangle{\pgfqpoint{3.775073in}{0.496542in}}{\pgfqpoint{1.042991in}{1.239972in}}%
\pgfusepath{clip}%
\pgfsetbuttcap%
\pgfsetroundjoin%
\pgfsetlinewidth{0.602250pt}%
\definecolor{currentstroke}{rgb}{0.788235,0.807843,0.819608}%
\pgfsetstrokecolor{currentstroke}%
\pgfsetdash{{0.600000pt}{0.990000pt}}{0.000000pt}%
\pgfpathmoveto{\pgfqpoint{4.032602in}{0.496542in}}%
\pgfpathlineto{\pgfqpoint{4.032602in}{1.736515in}}%
\pgfusepath{stroke}%
\end{pgfscope}%
\begin{pgfscope}%
\pgfsetrectcap%
\pgfsetmiterjoin%
\pgfsetlinewidth{0.803000pt}%
\definecolor{currentstroke}{rgb}{0.000000,0.000000,0.000000}%
\pgfsetstrokecolor{currentstroke}%
\pgfsetdash{}{0pt}%
\pgfpathmoveto{\pgfqpoint{3.775073in}{0.496542in}}%
\pgfpathlineto{\pgfqpoint{3.775073in}{1.736515in}}%
\pgfusepath{stroke}%
\end{pgfscope}%
\begin{pgfscope}%
\pgfsetrectcap%
\pgfsetmiterjoin%
\pgfsetlinewidth{0.803000pt}%
\definecolor{currentstroke}{rgb}{0.000000,0.000000,0.000000}%
\pgfsetstrokecolor{currentstroke}%
\pgfsetdash{}{0pt}%
\pgfpathmoveto{\pgfqpoint{4.818064in}{0.496542in}}%
\pgfpathlineto{\pgfqpoint{4.818064in}{1.736515in}}%
\pgfusepath{stroke}%
\end{pgfscope}%
\begin{pgfscope}%
\pgfsetrectcap%
\pgfsetmiterjoin%
\pgfsetlinewidth{0.803000pt}%
\definecolor{currentstroke}{rgb}{0.000000,0.000000,0.000000}%
\pgfsetstrokecolor{currentstroke}%
\pgfsetdash{}{0pt}%
\pgfpathmoveto{\pgfqpoint{3.775073in}{0.496542in}}%
\pgfpathlineto{\pgfqpoint{4.818064in}{0.496542in}}%
\pgfusepath{stroke}%
\end{pgfscope}%
\begin{pgfscope}%
\pgfsetrectcap%
\pgfsetmiterjoin%
\pgfsetlinewidth{0.803000pt}%
\definecolor{currentstroke}{rgb}{0.000000,0.000000,0.000000}%
\pgfsetstrokecolor{currentstroke}%
\pgfsetdash{}{0pt}%
\pgfpathmoveto{\pgfqpoint{3.775073in}{1.736515in}}%
\pgfpathlineto{\pgfqpoint{4.818064in}{1.736515in}}%
\pgfusepath{stroke}%
\end{pgfscope}%
\begin{pgfscope}%
\definecolor{textcolor}{rgb}{0.000000,0.000000,0.000000}%
\pgfsetstrokecolor{textcolor}%
\pgfsetfillcolor{textcolor}%
\pgftext[x=4.786774in,y=1.674516in,right,top]{\color{textcolor}{\rmfamily\fontsize{9.000000}{10.800000}\selectfont\catcode`\^=\active\def^{\ifmmode\sp\else\^{}\fi}\catcode`\%=\active\def%{\%}$(c)$}}%
\end{pgfscope}%
\begin{pgfscope}%
\pgfsetrectcap%
\pgfsetroundjoin%
\pgfsetlinewidth{0.903375pt}%
\definecolor{currentstroke}{rgb}{1.000000,0.000000,0.000000}%
\pgfsetstrokecolor{currentstroke}%
\pgfsetdash{}{0pt}%
\pgfpathmoveto{\pgfqpoint{1.698087in}{2.017009in}}%
\pgfpathlineto{\pgfqpoint{1.798087in}{2.017009in}}%
\pgfpathlineto{\pgfqpoint{1.898087in}{2.017009in}}%
\pgfusepath{stroke}%
\end{pgfscope}%
\begin{pgfscope}%
\pgfsetbuttcap%
\pgfsetroundjoin%
\definecolor{currentfill}{rgb}{0.000000,0.000000,0.000000}%
\pgfsetfillcolor{currentfill}%
\pgfsetfillopacity{0.000000}%
\pgfsetlinewidth{0.803000pt}%
\definecolor{currentstroke}{rgb}{1.000000,0.000000,0.000000}%
\pgfsetstrokecolor{currentstroke}%
\pgfsetdash{}{0pt}%
\pgfsys@defobject{currentmarker}{\pgfqpoint{-0.017500in}{-0.017500in}}{\pgfqpoint{0.017500in}{0.017500in}}{%
\pgfpathmoveto{\pgfqpoint{0.000000in}{-0.017500in}}%
\pgfpathcurveto{\pgfqpoint{0.004641in}{-0.017500in}}{\pgfqpoint{0.009093in}{-0.015656in}}{\pgfqpoint{0.012374in}{-0.012374in}}%
\pgfpathcurveto{\pgfqpoint{0.015656in}{-0.009093in}}{\pgfqpoint{0.017500in}{-0.004641in}}{\pgfqpoint{0.017500in}{0.000000in}}%
\pgfpathcurveto{\pgfqpoint{0.017500in}{0.004641in}}{\pgfqpoint{0.015656in}{0.009093in}}{\pgfqpoint{0.012374in}{0.012374in}}%
\pgfpathcurveto{\pgfqpoint{0.009093in}{0.015656in}}{\pgfqpoint{0.004641in}{0.017500in}}{\pgfqpoint{0.000000in}{0.017500in}}%
\pgfpathcurveto{\pgfqpoint{-0.004641in}{0.017500in}}{\pgfqpoint{-0.009093in}{0.015656in}}{\pgfqpoint{-0.012374in}{0.012374in}}%
\pgfpathcurveto{\pgfqpoint{-0.015656in}{0.009093in}}{\pgfqpoint{-0.017500in}{0.004641in}}{\pgfqpoint{-0.017500in}{0.000000in}}%
\pgfpathcurveto{\pgfqpoint{-0.017500in}{-0.004641in}}{\pgfqpoint{-0.015656in}{-0.009093in}}{\pgfqpoint{-0.012374in}{-0.012374in}}%
\pgfpathcurveto{\pgfqpoint{-0.009093in}{-0.015656in}}{\pgfqpoint{-0.004641in}{-0.017500in}}{\pgfqpoint{0.000000in}{-0.017500in}}%
\pgfpathlineto{\pgfqpoint{0.000000in}{-0.017500in}}%
\pgfpathclose%
\pgfusepath{stroke,fill}%
}%
\begin{pgfscope}%
\pgfsys@transformshift{1.798087in}{2.017009in}%
\pgfsys@useobject{currentmarker}{}%
\end{pgfscope}%
\end{pgfscope}%
\begin{pgfscope}%
\definecolor{textcolor}{rgb}{0.000000,0.000000,0.000000}%
\pgfsetstrokecolor{textcolor}%
\pgfsetfillcolor{textcolor}%
\pgftext[x=1.986976in,y=1.978120in,left,base]{\color{textcolor}{\rmfamily\fontsize{8.000000}{9.600000}\selectfont\catcode`\^=\active\def^{\ifmmode\sp\else\^{}\fi}\catcode`\%=\active\def%{\%}memory}}%
\end{pgfscope}%
\begin{pgfscope}%
\pgfsetbuttcap%
\pgfsetroundjoin%
\pgfsetlinewidth{0.903375pt}%
\definecolor{currentstroke}{rgb}{0.000000,0.000000,0.000000}%
\pgfsetstrokecolor{currentstroke}%
\pgfsetdash{{3.330000pt}{1.440000pt}}{0.000000pt}%
\pgfpathmoveto{\pgfqpoint{2.570219in}{2.017009in}}%
\pgfpathlineto{\pgfqpoint{2.670219in}{2.017009in}}%
\pgfpathlineto{\pgfqpoint{2.770219in}{2.017009in}}%
\pgfusepath{stroke}%
\end{pgfscope}%
\begin{pgfscope}%
\definecolor{textcolor}{rgb}{0.000000,0.000000,0.000000}%
\pgfsetstrokecolor{textcolor}%
\pgfsetfillcolor{textcolor}%
\pgftext[x=2.859107in,y=1.978120in,left,base]{\color{textcolor}{\rmfamily\fontsize{8.000000}{9.600000}\selectfont\catcode`\^=\active\def^{\ifmmode\sp\else\^{}\fi}\catcode`\%=\active\def%{\%}passive}}%
\end{pgfscope}%
\end{pgfpicture}%
\makeatother%
\endgroup%
}
	\caption{(En color en l\'inea) (a) Fidelidad incondicional $F$, (b) probabilidad de recuperaci\'on $s$ y (c) fidelidad condicional $F_h=F/s$ en funci\'on de la espera $T\omega_b$, con carga y lectura reales, para la memoria de estados oscuros (rojo; los c\'irculos abiertos son los puntos calculados a esperas equiespaciadas) y para dos modos ac\'usticos pasivos con los mismos puertos (negro, discontinuo). Las l\'ineas verticales de puntos marcan $T\omega_b=2000$. Par\'ametros del sistema: como en la Fig.~\ref{fig:storage}, con $\kappa_{\rm in}/\kappa=\kappa_{\rm out}/\kappa=10$, $\chi=1$, una separaci\'on de $100/\omega_b$ entre los dos intervalos temporales y los pulsos optimizados de la Fig.~\ref{fig:raman_map}.}
	\label{fig:sweep}
\end{figure}

La Figura~\ref{fig:sweep} compara la memoria con la referencia pasiva en funci\'on de la espera $T$. La fidelidad incondicional $F$ y la probabilidad de recuperaci\'on $s$ de los modos pasivos decaen con el tiempo de vida de la cavidad, ya que la excitaci\'on se mantiene como fon\'on durante toda la espera. En la memoria el fon\'on se convierte en un ex\'citon oscuro, de modo que el decaimiento de la cavidad ya no limita la amplitud almacenada y $F$ y $s$ son casi independientes de $T$. Para $F$ las dos curvas se cruzan a una espera muy corta, donde los modos pasivos a\'un conservan casi toda la excitaci\'on, y $s$ de la memoria ya es mayor a espera cero; m\'as all\'a la ventaja crece con $T$. La meseta queda por debajo de la unidad porque $s$ est\'a acotada por el producto de las eficiencias de carga, escritura, lectura y emisi\'on, ninguna de las cuales depende de $T$.

El panel (c) muestra lo que la memoria no hace. La fidelidad condicional, que mide la coherencia del qubit una vez recuperado, queda por debajo del valor pasivo incluso a espera cero y disminuye lentamente con $T$. Una cavidad pasiva pierde amplitud pero conserva la fase, de modo que su fidelidad condicional est\'a limitada solo por los puertos. La brecha a espera cero no la causa el almacenamiento: proviene de la escritura y la lectura secuenciales, durante las cuales una rama contiene un ex\'citon y la otra ninguno, y de las transferencias Raman imperfectas (Fig.~\ref{fig:protocol}). La disminuci\'on lenta con $T$ es el costo adicional de la relajaci\'on y el desfase residuales de las ranuras durante la espera. La ventaja de la memoria es por tanto la probabilidad de recuperar el qubit tras una espera en la que los modos pasivos pierden la excitaci\'on, no una mayor fidelidad por qubit recuperado.

La Figura~\ref{fig:ports}(a),(b) var\'ia la velocidad de los puertos. La carga es el paso cr\'itico. Mientras el primer intervalo temporal est\'a almacenado y el segundo a\'un se est\'a cargando, una rama contiene un ex\'citon y la otra ninguno, y este intervalo sin protecci\'on no est\'a cubierto por la cancelaci\'on de modo com\'un (Ec.~\eqref{eq:visibility}). Un puerto de entrada m\'as r\'apido lo acorta, y $F_h$ aumenta con $\kappa_{\rm in}/\kappa$. El puerto de salida tiene la misma estructura, pues la primera componente sigue almacenada mientras se emite la segunda, y tiene un efecto comparable, aunque su ganancia satura antes. La dependencia con las tasas de los puertos es casi la misma para todas las esperas y las curvas difieren por un desplazamiento vertical, de modo que la espera solo a\~nade la lenta decoherencia de la Fig.~\ref{fig:sweep}(c) sobre lo que dejan la carga y la lectura.

El acoplo por deformaci\'on fija el acoplo Raman efectivo $g_s$ y con ello el tiempo necesario para escribir y leer. Para comprobar el requisito de la Ec.~\eqref{eq:gmin}, escalamos todos los acoplos por deformaci\'on por un factor $S$, lo que escala $g_s$ de la misma manera, y reoptimizamos la amplitud, la duraci\'on y la desintonizaci\'on de ambos pulsos para cada $S$ con puertos ideales y amplitud acotada; donde se hicieron dos optimizaciones se conservan los pulsos con mayor $F_h$ con puertos reales. La Figura~\ref{fig:ports}(d) muestra que $F_h$ es casi independiente de $S$ para acoplos grandes y disminuye por debajo de un umbral ($g_s\simeq0.03\,\omega_b$). La disminuci\'on con puertos reales ocurre en escalones, que corresponden a distintos reg\'imenes de amplitud encontrados por el optimizador de pulsos, de modo que la posici\'on de cada escal\'on no debe sobreinterpretarse; lo robusto es que por debajo del umbral la fidelidad condicional cae. La raz\'on es la anticipada en la Sec.~\ref{sec:raman}: un $g_s$ menor requiere un pulso m\'as intenso o m\'as largo, y una vez alcanzado el l\'imite de amplitud, la poblaci\'on virtual del estado brillante y el desplazamiento por luz crecen y limitan la transferencia aun cuando la memoria siga superando a los modos pasivos en $F$ y $s$. La curva con puertos ideales queda por encima de la de puertos reales por el costo de la carga y la emisi\'on, que es peque\~no para acoplos grandes y crece cuando los pulsos se alargan. La curva con ruido no correlacionado es baja y plana: un acoplo mayor acorta los pulsos pero no puede reemplazar la protecci\'on contra el ruido.

\begin{figure}[t]
\centering
\resizebox{\columnwidth}{!}{%% Creator: Matplotlib, PGF backend
%%
%% To include the figure in your LaTeX document, write
%%   \input{<filename>.pgf}
%%
%% Make sure the required packages are loaded in your preamble
%%   \usepackage{pgf}
%%
%% Also ensure that all the required font packages are loaded; for instance,
%% the lmodern package is sometimes necessary when using math font.
%%   \usepackage{lmodern}
%%
%% Figures using additional raster images can only be included by \input if
%% they are in the same directory as the main LaTeX file. For loading figures
%% from other directories you can use the `import` package
%%   \usepackage{import}
%%
%% and then include the figures with
%%   \import{<path to file>}{<filename>.pgf}
%%
%% Matplotlib used the following preamble
%%   \def\mathdefault#1{#1}
%%   \everymath=\expandafter{\the\everymath\displaystyle}
%%   \IfFileExists{scrextend.sty}{
%%     \usepackage[fontsize=8.000000pt]{scrextend}
%%   }{
%%     \renewcommand{\normalsize}{\fontsize{8.000000}{9.600000}\selectfont}
%%     \normalsize
%%   }
%%   \usepackage{amsmath,amssymb}
%%   \makeatletter\@ifpackageloaded{underscore}{}{\usepackage[strings]{underscore}}\makeatother
%%
\begingroup%
\makeatletter%
\begin{pgfpicture}%
\pgfpathrectangle{\pgfpointorigin}{\pgfqpoint{4.321618in}{3.783704in}}%
\pgfusepath{use as bounding box, clip}%
\begin{pgfscope}%
\pgfsetbuttcap%
\pgfsetmiterjoin%
\definecolor{currentfill}{rgb}{1.000000,1.000000,1.000000}%
\pgfsetfillcolor{currentfill}%
\pgfsetlinewidth{0.000000pt}%
\definecolor{currentstroke}{rgb}{1.000000,1.000000,1.000000}%
\pgfsetstrokecolor{currentstroke}%
\pgfsetdash{}{0pt}%
\pgfpathmoveto{\pgfqpoint{0.000000in}{0.000000in}}%
\pgfpathlineto{\pgfqpoint{4.321618in}{0.000000in}}%
\pgfpathlineto{\pgfqpoint{4.321618in}{3.783704in}}%
\pgfpathlineto{\pgfqpoint{0.000000in}{3.783704in}}%
\pgfpathlineto{\pgfqpoint{0.000000in}{0.000000in}}%
\pgfpathclose%
\pgfusepath{fill}%
\end{pgfscope}%
\begin{pgfscope}%
\pgfsetbuttcap%
\pgfsetmiterjoin%
\definecolor{currentfill}{rgb}{1.000000,1.000000,1.000000}%
\pgfsetfillcolor{currentfill}%
\pgfsetlinewidth{0.000000pt}%
\definecolor{currentstroke}{rgb}{0.000000,0.000000,0.000000}%
\pgfsetstrokecolor{currentstroke}%
\pgfsetstrokeopacity{0.000000}%
\pgfsetdash{}{0pt}%
\pgfpathmoveto{\pgfqpoint{0.473073in}{2.151654in}}%
\pgfpathlineto{\pgfqpoint{2.178118in}{2.151654in}}%
\pgfpathlineto{\pgfqpoint{2.178118in}{3.319877in}}%
\pgfpathlineto{\pgfqpoint{0.473073in}{3.319877in}}%
\pgfpathlineto{\pgfqpoint{0.473073in}{2.151654in}}%
\pgfpathclose%
\pgfusepath{fill}%
\end{pgfscope}%
\begin{pgfscope}%
\pgfsetbuttcap%
\pgfsetroundjoin%
\definecolor{currentfill}{rgb}{0.000000,0.000000,0.000000}%
\pgfsetfillcolor{currentfill}%
\pgfsetlinewidth{0.803000pt}%
\definecolor{currentstroke}{rgb}{0.000000,0.000000,0.000000}%
\pgfsetstrokecolor{currentstroke}%
\pgfsetdash{}{0pt}%
\pgfsys@defobject{currentmarker}{\pgfqpoint{0.000000in}{-0.048611in}}{\pgfqpoint{0.000000in}{0.000000in}}{%
\pgfpathmoveto{\pgfqpoint{0.000000in}{0.000000in}}%
\pgfpathlineto{\pgfqpoint{0.000000in}{-0.048611in}}%
\pgfusepath{stroke,fill}%
}%
\begin{pgfscope}%
\pgfsys@transformshift{0.473073in}{2.151654in}%
\pgfsys@useobject{currentmarker}{}%
\end{pgfscope}%
\end{pgfscope}%
\begin{pgfscope}%
\definecolor{textcolor}{rgb}{0.000000,0.000000,0.000000}%
\pgfsetstrokecolor{textcolor}%
\pgfsetfillcolor{textcolor}%
\pgftext[x=0.473073in,y=2.054432in,,top]{\color{textcolor}{\rmfamily\fontsize{8.000000}{9.600000}\selectfont\catcode`\^=\active\def^{\ifmmode\sp\else\^{}\fi}\catcode`\%=\active\def%{\%}$0$}}%
\end{pgfscope}%
\begin{pgfscope}%
\pgfsetbuttcap%
\pgfsetroundjoin%
\definecolor{currentfill}{rgb}{0.000000,0.000000,0.000000}%
\pgfsetfillcolor{currentfill}%
\pgfsetlinewidth{0.803000pt}%
\definecolor{currentstroke}{rgb}{0.000000,0.000000,0.000000}%
\pgfsetstrokecolor{currentstroke}%
\pgfsetdash{}{0pt}%
\pgfsys@defobject{currentmarker}{\pgfqpoint{0.000000in}{-0.048611in}}{\pgfqpoint{0.000000in}{0.000000in}}{%
\pgfpathmoveto{\pgfqpoint{0.000000in}{0.000000in}}%
\pgfpathlineto{\pgfqpoint{0.000000in}{-0.048611in}}%
\pgfusepath{stroke,fill}%
}%
\begin{pgfscope}%
\pgfsys@transformshift{1.076095in}{2.151654in}%
\pgfsys@useobject{currentmarker}{}%
\end{pgfscope}%
\end{pgfscope}%
\begin{pgfscope}%
\definecolor{textcolor}{rgb}{0.000000,0.000000,0.000000}%
\pgfsetstrokecolor{textcolor}%
\pgfsetfillcolor{textcolor}%
\pgftext[x=1.076095in,y=2.054432in,,top]{\color{textcolor}{\rmfamily\fontsize{8.000000}{9.600000}\selectfont\catcode`\^=\active\def^{\ifmmode\sp\else\^{}\fi}\catcode`\%=\active\def%{\%}$10$}}%
\end{pgfscope}%
\begin{pgfscope}%
\pgfsetbuttcap%
\pgfsetroundjoin%
\definecolor{currentfill}{rgb}{0.000000,0.000000,0.000000}%
\pgfsetfillcolor{currentfill}%
\pgfsetlinewidth{0.803000pt}%
\definecolor{currentstroke}{rgb}{0.000000,0.000000,0.000000}%
\pgfsetstrokecolor{currentstroke}%
\pgfsetdash{}{0pt}%
\pgfsys@defobject{currentmarker}{\pgfqpoint{0.000000in}{-0.048611in}}{\pgfqpoint{0.000000in}{0.000000in}}{%
\pgfpathmoveto{\pgfqpoint{0.000000in}{0.000000in}}%
\pgfpathlineto{\pgfqpoint{0.000000in}{-0.048611in}}%
\pgfusepath{stroke,fill}%
}%
\begin{pgfscope}%
\pgfsys@transformshift{1.679118in}{2.151654in}%
\pgfsys@useobject{currentmarker}{}%
\end{pgfscope}%
\end{pgfscope}%
\begin{pgfscope}%
\definecolor{textcolor}{rgb}{0.000000,0.000000,0.000000}%
\pgfsetstrokecolor{textcolor}%
\pgfsetfillcolor{textcolor}%
\pgftext[x=1.679118in,y=2.054432in,,top]{\color{textcolor}{\rmfamily\fontsize{8.000000}{9.600000}\selectfont\catcode`\^=\active\def^{\ifmmode\sp\else\^{}\fi}\catcode`\%=\active\def%{\%}$20$}}%
\end{pgfscope}%
\begin{pgfscope}%
\definecolor{textcolor}{rgb}{0.000000,0.000000,0.000000}%
\pgfsetstrokecolor{textcolor}%
\pgfsetfillcolor{textcolor}%
\pgftext[x=1.325596in,y=1.900111in,,top]{\color{textcolor}{\rmfamily\fontsize{9.000000}{10.800000}\selectfont\catcode`\^=\active\def^{\ifmmode\sp\else\^{}\fi}\catcode`\%=\active\def%{\%}$\kappa_{\rm in}/\kappa$}}%
\end{pgfscope}%
\begin{pgfscope}%
\pgfsetbuttcap%
\pgfsetroundjoin%
\definecolor{currentfill}{rgb}{0.000000,0.000000,0.000000}%
\pgfsetfillcolor{currentfill}%
\pgfsetlinewidth{0.803000pt}%
\definecolor{currentstroke}{rgb}{0.000000,0.000000,0.000000}%
\pgfsetstrokecolor{currentstroke}%
\pgfsetdash{}{0pt}%
\pgfsys@defobject{currentmarker}{\pgfqpoint{-0.048611in}{0.000000in}}{\pgfqpoint{-0.000000in}{0.000000in}}{%
\pgfpathmoveto{\pgfqpoint{-0.000000in}{0.000000in}}%
\pgfpathlineto{\pgfqpoint{-0.048611in}{0.000000in}}%
\pgfusepath{stroke,fill}%
}%
\begin{pgfscope}%
\pgfsys@transformshift{0.473073in}{2.151654in}%
\pgfsys@useobject{currentmarker}{}%
\end{pgfscope}%
\end{pgfscope}%
\begin{pgfscope}%
\definecolor{textcolor}{rgb}{0.000000,0.000000,0.000000}%
\pgfsetstrokecolor{textcolor}%
\pgfsetfillcolor{textcolor}%
\pgftext[x=0.316822in, y=2.113074in, left, base]{\color{textcolor}{\rmfamily\fontsize{8.000000}{9.600000}\selectfont\catcode`\^=\active\def^{\ifmmode\sp\else\^{}\fi}\catcode`\%=\active\def%{\%}$0$}}%
\end{pgfscope}%
\begin{pgfscope}%
\pgfsetbuttcap%
\pgfsetroundjoin%
\definecolor{currentfill}{rgb}{0.000000,0.000000,0.000000}%
\pgfsetfillcolor{currentfill}%
\pgfsetlinewidth{0.803000pt}%
\definecolor{currentstroke}{rgb}{0.000000,0.000000,0.000000}%
\pgfsetstrokecolor{currentstroke}%
\pgfsetdash{}{0pt}%
\pgfsys@defobject{currentmarker}{\pgfqpoint{-0.048611in}{0.000000in}}{\pgfqpoint{-0.000000in}{0.000000in}}{%
\pgfpathmoveto{\pgfqpoint{-0.000000in}{0.000000in}}%
\pgfpathlineto{\pgfqpoint{-0.048611in}{0.000000in}}%
\pgfusepath{stroke,fill}%
}%
\begin{pgfscope}%
\pgfsys@transformshift{0.473073in}{2.638414in}%
\pgfsys@useobject{currentmarker}{}%
\end{pgfscope}%
\end{pgfscope}%
\begin{pgfscope}%
\definecolor{textcolor}{rgb}{0.000000,0.000000,0.000000}%
\pgfsetstrokecolor{textcolor}%
\pgfsetfillcolor{textcolor}%
\pgftext[x=0.225000in, y=2.599833in, left, base]{\color{textcolor}{\rmfamily\fontsize{8.000000}{9.600000}\selectfont\catcode`\^=\active\def^{\ifmmode\sp\else\^{}\fi}\catcode`\%=\active\def%{\%}$0.5$}}%
\end{pgfscope}%
\begin{pgfscope}%
\pgfsetbuttcap%
\pgfsetroundjoin%
\definecolor{currentfill}{rgb}{0.000000,0.000000,0.000000}%
\pgfsetfillcolor{currentfill}%
\pgfsetlinewidth{0.803000pt}%
\definecolor{currentstroke}{rgb}{0.000000,0.000000,0.000000}%
\pgfsetstrokecolor{currentstroke}%
\pgfsetdash{}{0pt}%
\pgfsys@defobject{currentmarker}{\pgfqpoint{-0.048611in}{0.000000in}}{\pgfqpoint{-0.000000in}{0.000000in}}{%
\pgfpathmoveto{\pgfqpoint{-0.000000in}{0.000000in}}%
\pgfpathlineto{\pgfqpoint{-0.048611in}{0.000000in}}%
\pgfusepath{stroke,fill}%
}%
\begin{pgfscope}%
\pgfsys@transformshift{0.473073in}{3.125173in}%
\pgfsys@useobject{currentmarker}{}%
\end{pgfscope}%
\end{pgfscope}%
\begin{pgfscope}%
\definecolor{textcolor}{rgb}{0.000000,0.000000,0.000000}%
\pgfsetstrokecolor{textcolor}%
\pgfsetfillcolor{textcolor}%
\pgftext[x=0.316822in, y=3.086593in, left, base]{\color{textcolor}{\rmfamily\fontsize{8.000000}{9.600000}\selectfont\catcode`\^=\active\def^{\ifmmode\sp\else\^{}\fi}\catcode`\%=\active\def%{\%}$1$}}%
\end{pgfscope}%
\begin{pgfscope}%
\definecolor{textcolor}{rgb}{0.000000,0.000000,0.000000}%
\pgfsetstrokecolor{textcolor}%
\pgfsetfillcolor{textcolor}%
\pgftext[x=0.211111in,y=2.735766in,,bottom,rotate=90.000000]{\color{textcolor}{\rmfamily\fontsize{9.000000}{10.800000}\selectfont\catcode`\^=\active\def^{\ifmmode\sp\else\^{}\fi}\catcode`\%=\active\def%{\%}$F_h$}}%
\end{pgfscope}%
\begin{pgfscope}%
\pgfpathrectangle{\pgfqpoint{0.473073in}{2.151654in}}{\pgfqpoint{1.705045in}{1.168222in}}%
\pgfusepath{clip}%
\pgfsetrectcap%
\pgfsetroundjoin%
\pgfsetlinewidth{0.903375pt}%
\definecolor{currentstroke}{rgb}{0.121569,0.466667,0.705882}%
\pgfsetstrokecolor{currentstroke}%
\pgfsetdash{}{0pt}%
\pgfpathmoveto{\pgfqpoint{0.563527in}{2.884123in}}%
\pgfpathlineto{\pgfqpoint{0.594435in}{2.915964in}}%
\pgfpathlineto{\pgfqpoint{0.617617in}{2.936511in}}%
\pgfpathlineto{\pgfqpoint{0.633071in}{2.948108in}}%
\pgfpathlineto{\pgfqpoint{0.648525in}{2.957701in}}%
\pgfpathlineto{\pgfqpoint{0.671707in}{2.968826in}}%
\pgfpathlineto{\pgfqpoint{0.702616in}{2.980774in}}%
\pgfpathlineto{\pgfqpoint{0.733524in}{2.989832in}}%
\pgfpathlineto{\pgfqpoint{0.764433in}{2.996378in}}%
\pgfpathlineto{\pgfqpoint{0.810796in}{3.003508in}}%
\pgfpathlineto{\pgfqpoint{0.942158in}{3.020662in}}%
\pgfpathlineto{\pgfqpoint{1.003975in}{3.026072in}}%
\pgfpathlineto{\pgfqpoint{1.065792in}{3.028747in}}%
\pgfpathlineto{\pgfqpoint{1.158518in}{3.032792in}}%
\pgfpathlineto{\pgfqpoint{1.228063in}{3.035229in}}%
\pgfpathlineto{\pgfqpoint{1.513968in}{3.040265in}}%
\pgfpathlineto{\pgfqpoint{1.598967in}{3.041098in}}%
\pgfpathlineto{\pgfqpoint{1.738056in}{3.045558in}}%
\pgfpathlineto{\pgfqpoint{1.861691in}{3.046148in}}%
\pgfpathlineto{\pgfqpoint{1.923508in}{3.049080in}}%
\pgfpathlineto{\pgfqpoint{1.962144in}{3.047608in}}%
\pgfpathlineto{\pgfqpoint{2.008507in}{3.045753in}}%
\pgfpathlineto{\pgfqpoint{2.078052in}{3.046752in}}%
\pgfpathlineto{\pgfqpoint{2.101233in}{3.047738in}}%
\pgfpathlineto{\pgfqpoint{2.101233in}{3.047738in}}%
\pgfusepath{stroke}%
\end{pgfscope}%
\begin{pgfscope}%
\pgfpathrectangle{\pgfqpoint{0.473073in}{2.151654in}}{\pgfqpoint{1.705045in}{1.168222in}}%
\pgfusepath{clip}%
\pgfsetbuttcap%
\pgfsetroundjoin%
\definecolor{currentfill}{rgb}{0.121569,0.466667,0.705882}%
\pgfsetfillcolor{currentfill}%
\pgfsetlinewidth{0.803000pt}%
\definecolor{currentstroke}{rgb}{0.121569,0.466667,0.705882}%
\pgfsetstrokecolor{currentstroke}%
\pgfsetdash{}{0pt}%
\pgfsys@defobject{currentmarker}{\pgfqpoint{-0.019444in}{-0.019444in}}{\pgfqpoint{0.019444in}{0.019444in}}{%
\pgfpathmoveto{\pgfqpoint{-0.019444in}{-0.019444in}}%
\pgfpathlineto{\pgfqpoint{0.019444in}{0.019444in}}%
\pgfpathmoveto{\pgfqpoint{-0.019444in}{0.019444in}}%
\pgfpathlineto{\pgfqpoint{0.019444in}{-0.019444in}}%
\pgfusepath{stroke,fill}%
}%
\begin{pgfscope}%
\pgfsys@transformshift{0.563527in}{2.884123in}%
\pgfsys@useobject{currentmarker}{}%
\end{pgfscope}%
\begin{pgfscope}%
\pgfsys@transformshift{0.653980in}{2.960558in}%
\pgfsys@useobject{currentmarker}{}%
\end{pgfscope}%
\begin{pgfscope}%
\pgfsys@transformshift{0.744433in}{2.992393in}%
\pgfsys@useobject{currentmarker}{}%
\end{pgfscope}%
\begin{pgfscope}%
\pgfsys@transformshift{0.834886in}{3.006893in}%
\pgfsys@useobject{currentmarker}{}%
\end{pgfscope}%
\begin{pgfscope}%
\pgfsys@transformshift{0.925340in}{3.018911in}%
\pgfsys@useobject{currentmarker}{}%
\end{pgfscope}%
\begin{pgfscope}%
\pgfsys@transformshift{1.015793in}{3.026768in}%
\pgfsys@useobject{currentmarker}{}%
\end{pgfscope}%
\begin{pgfscope}%
\pgfsys@transformshift{1.106246in}{3.030172in}%
\pgfsys@useobject{currentmarker}{}%
\end{pgfscope}%
\begin{pgfscope}%
\pgfsys@transformshift{1.196700in}{3.034425in}%
\pgfsys@useobject{currentmarker}{}%
\end{pgfscope}%
\begin{pgfscope}%
\pgfsys@transformshift{1.287153in}{3.036469in}%
\pgfsys@useobject{currentmarker}{}%
\end{pgfscope}%
\begin{pgfscope}%
\pgfsys@transformshift{1.377606in}{3.038798in}%
\pgfsys@useobject{currentmarker}{}%
\end{pgfscope}%
\begin{pgfscope}%
\pgfsys@transformshift{1.468060in}{3.040342in}%
\pgfsys@useobject{currentmarker}{}%
\end{pgfscope}%
\begin{pgfscope}%
\pgfsys@transformshift{1.558513in}{3.040192in}%
\pgfsys@useobject{currentmarker}{}%
\end{pgfscope}%
\begin{pgfscope}%
\pgfsys@transformshift{1.648966in}{3.043046in}%
\pgfsys@useobject{currentmarker}{}%
\end{pgfscope}%
\begin{pgfscope}%
\pgfsys@transformshift{1.739420in}{3.045560in}%
\pgfsys@useobject{currentmarker}{}%
\end{pgfscope}%
\begin{pgfscope}%
\pgfsys@transformshift{1.829873in}{3.044979in}%
\pgfsys@useobject{currentmarker}{}%
\end{pgfscope}%
\begin{pgfscope}%
\pgfsys@transformshift{1.920326in}{3.049092in}%
\pgfsys@useobject{currentmarker}{}%
\end{pgfscope}%
\begin{pgfscope}%
\pgfsys@transformshift{2.010780in}{3.045747in}%
\pgfsys@useobject{currentmarker}{}%
\end{pgfscope}%
\begin{pgfscope}%
\pgfsys@transformshift{2.101233in}{3.047738in}%
\pgfsys@useobject{currentmarker}{}%
\end{pgfscope}%
\end{pgfscope}%
\begin{pgfscope}%
\pgfpathrectangle{\pgfqpoint{0.473073in}{2.151654in}}{\pgfqpoint{1.705045in}{1.168222in}}%
\pgfusepath{clip}%
\pgfsetrectcap%
\pgfsetroundjoin%
\pgfsetlinewidth{0.903375pt}%
\definecolor{currentstroke}{rgb}{1.000000,0.000000,0.000000}%
\pgfsetstrokecolor{currentstroke}%
\pgfsetdash{}{0pt}%
\pgfpathmoveto{\pgfqpoint{0.563527in}{2.879869in}}%
\pgfpathlineto{\pgfqpoint{0.594435in}{2.911736in}}%
\pgfpathlineto{\pgfqpoint{0.617617in}{2.932234in}}%
\pgfpathlineto{\pgfqpoint{0.633071in}{2.943731in}}%
\pgfpathlineto{\pgfqpoint{0.648525in}{2.953154in}}%
\pgfpathlineto{\pgfqpoint{0.671707in}{2.963903in}}%
\pgfpathlineto{\pgfqpoint{0.702616in}{2.975320in}}%
\pgfpathlineto{\pgfqpoint{0.733524in}{2.983948in}}%
\pgfpathlineto{\pgfqpoint{0.764433in}{2.990256in}}%
\pgfpathlineto{\pgfqpoint{0.810796in}{2.997158in}}%
\pgfpathlineto{\pgfqpoint{0.942158in}{3.013941in}}%
\pgfpathlineto{\pgfqpoint{1.003975in}{3.019134in}}%
\pgfpathlineto{\pgfqpoint{1.058065in}{3.020905in}}%
\pgfpathlineto{\pgfqpoint{1.135337in}{3.023463in}}%
\pgfpathlineto{\pgfqpoint{1.235790in}{3.028119in}}%
\pgfpathlineto{\pgfqpoint{1.343971in}{3.030442in}}%
\pgfpathlineto{\pgfqpoint{1.629876in}{3.035607in}}%
\pgfpathlineto{\pgfqpoint{1.714875in}{3.037172in}}%
\pgfpathlineto{\pgfqpoint{1.877145in}{3.039592in}}%
\pgfpathlineto{\pgfqpoint{1.923508in}{3.041330in}}%
\pgfpathlineto{\pgfqpoint{1.962144in}{3.039877in}}%
\pgfpathlineto{\pgfqpoint{2.008507in}{3.038046in}}%
\pgfpathlineto{\pgfqpoint{2.078052in}{3.039053in}}%
\pgfpathlineto{\pgfqpoint{2.101233in}{3.040038in}}%
\pgfpathlineto{\pgfqpoint{2.101233in}{3.040038in}}%
\pgfusepath{stroke}%
\end{pgfscope}%
\begin{pgfscope}%
\pgfpathrectangle{\pgfqpoint{0.473073in}{2.151654in}}{\pgfqpoint{1.705045in}{1.168222in}}%
\pgfusepath{clip}%
\pgfsetbuttcap%
\pgfsetroundjoin%
\definecolor{currentfill}{rgb}{0.000000,0.000000,0.000000}%
\pgfsetfillcolor{currentfill}%
\pgfsetfillopacity{0.000000}%
\pgfsetlinewidth{0.803000pt}%
\definecolor{currentstroke}{rgb}{1.000000,0.000000,0.000000}%
\pgfsetstrokecolor{currentstroke}%
\pgfsetdash{}{0pt}%
\pgfsys@defobject{currentmarker}{\pgfqpoint{-0.017500in}{-0.017500in}}{\pgfqpoint{0.017500in}{0.017500in}}{%
\pgfpathmoveto{\pgfqpoint{0.000000in}{-0.017500in}}%
\pgfpathcurveto{\pgfqpoint{0.004641in}{-0.017500in}}{\pgfqpoint{0.009093in}{-0.015656in}}{\pgfqpoint{0.012374in}{-0.012374in}}%
\pgfpathcurveto{\pgfqpoint{0.015656in}{-0.009093in}}{\pgfqpoint{0.017500in}{-0.004641in}}{\pgfqpoint{0.017500in}{0.000000in}}%
\pgfpathcurveto{\pgfqpoint{0.017500in}{0.004641in}}{\pgfqpoint{0.015656in}{0.009093in}}{\pgfqpoint{0.012374in}{0.012374in}}%
\pgfpathcurveto{\pgfqpoint{0.009093in}{0.015656in}}{\pgfqpoint{0.004641in}{0.017500in}}{\pgfqpoint{0.000000in}{0.017500in}}%
\pgfpathcurveto{\pgfqpoint{-0.004641in}{0.017500in}}{\pgfqpoint{-0.009093in}{0.015656in}}{\pgfqpoint{-0.012374in}{0.012374in}}%
\pgfpathcurveto{\pgfqpoint{-0.015656in}{0.009093in}}{\pgfqpoint{-0.017500in}{0.004641in}}{\pgfqpoint{-0.017500in}{0.000000in}}%
\pgfpathcurveto{\pgfqpoint{-0.017500in}{-0.004641in}}{\pgfqpoint{-0.015656in}{-0.009093in}}{\pgfqpoint{-0.012374in}{-0.012374in}}%
\pgfpathcurveto{\pgfqpoint{-0.009093in}{-0.015656in}}{\pgfqpoint{-0.004641in}{-0.017500in}}{\pgfqpoint{0.000000in}{-0.017500in}}%
\pgfpathlineto{\pgfqpoint{0.000000in}{-0.017500in}}%
\pgfpathclose%
\pgfusepath{stroke,fill}%
}%
\begin{pgfscope}%
\pgfsys@transformshift{0.563527in}{2.879869in}%
\pgfsys@useobject{currentmarker}{}%
\end{pgfscope}%
\begin{pgfscope}%
\pgfsys@transformshift{0.653980in}{2.955930in}%
\pgfsys@useobject{currentmarker}{}%
\end{pgfscope}%
\begin{pgfscope}%
\pgfsys@transformshift{0.744433in}{2.986408in}%
\pgfsys@useobject{currentmarker}{}%
\end{pgfscope}%
\begin{pgfscope}%
\pgfsys@transformshift{0.834886in}{3.000453in}%
\pgfsys@useobject{currentmarker}{}%
\end{pgfscope}%
\begin{pgfscope}%
\pgfsys@transformshift{0.925340in}{3.012238in}%
\pgfsys@useobject{currentmarker}{}%
\end{pgfscope}%
\begin{pgfscope}%
\pgfsys@transformshift{1.015793in}{3.019707in}%
\pgfsys@useobject{currentmarker}{}%
\end{pgfscope}%
\begin{pgfscope}%
\pgfsys@transformshift{1.106246in}{3.022109in}%
\pgfsys@useobject{currentmarker}{}%
\end{pgfscope}%
\begin{pgfscope}%
\pgfsys@transformshift{1.196700in}{3.026737in}%
\pgfsys@useobject{currentmarker}{}%
\end{pgfscope}%
\begin{pgfscope}%
\pgfsys@transformshift{1.287153in}{3.029497in}%
\pgfsys@useobject{currentmarker}{}%
\end{pgfscope}%
\begin{pgfscope}%
\pgfsys@transformshift{1.377606in}{3.030935in}%
\pgfsys@useobject{currentmarker}{}%
\end{pgfscope}%
\begin{pgfscope}%
\pgfsys@transformshift{1.468060in}{3.032352in}%
\pgfsys@useobject{currentmarker}{}%
\end{pgfscope}%
\begin{pgfscope}%
\pgfsys@transformshift{1.558513in}{3.033291in}%
\pgfsys@useobject{currentmarker}{}%
\end{pgfscope}%
\begin{pgfscope}%
\pgfsys@transformshift{1.648966in}{3.036107in}%
\pgfsys@useobject{currentmarker}{}%
\end{pgfscope}%
\begin{pgfscope}%
\pgfsys@transformshift{1.739420in}{3.037345in}%
\pgfsys@useobject{currentmarker}{}%
\end{pgfscope}%
\begin{pgfscope}%
\pgfsys@transformshift{1.829873in}{3.037508in}%
\pgfsys@useobject{currentmarker}{}%
\end{pgfscope}%
\begin{pgfscope}%
\pgfsys@transformshift{1.920326in}{3.041342in}%
\pgfsys@useobject{currentmarker}{}%
\end{pgfscope}%
\begin{pgfscope}%
\pgfsys@transformshift{2.010780in}{3.038040in}%
\pgfsys@useobject{currentmarker}{}%
\end{pgfscope}%
\begin{pgfscope}%
\pgfsys@transformshift{2.101233in}{3.040038in}%
\pgfsys@useobject{currentmarker}{}%
\end{pgfscope}%
\end{pgfscope}%
\begin{pgfscope}%
\pgfpathrectangle{\pgfqpoint{0.473073in}{2.151654in}}{\pgfqpoint{1.705045in}{1.168222in}}%
\pgfusepath{clip}%
\pgfsetrectcap%
\pgfsetroundjoin%
\pgfsetlinewidth{0.903375pt}%
\definecolor{currentstroke}{rgb}{0.172549,0.627451,0.172549}%
\pgfsetstrokecolor{currentstroke}%
\pgfsetdash{}{0pt}%
\pgfpathmoveto{\pgfqpoint{0.563527in}{2.854226in}}%
\pgfpathlineto{\pgfqpoint{0.594435in}{2.884602in}}%
\pgfpathlineto{\pgfqpoint{0.617617in}{2.904133in}}%
\pgfpathlineto{\pgfqpoint{0.633071in}{2.915081in}}%
\pgfpathlineto{\pgfqpoint{0.648525in}{2.924045in}}%
\pgfpathlineto{\pgfqpoint{0.671707in}{2.934296in}}%
\pgfpathlineto{\pgfqpoint{0.702616in}{2.945275in}}%
\pgfpathlineto{\pgfqpoint{0.733524in}{2.953395in}}%
\pgfpathlineto{\pgfqpoint{0.764433in}{2.958866in}}%
\pgfpathlineto{\pgfqpoint{0.818523in}{2.965394in}}%
\pgfpathlineto{\pgfqpoint{0.888068in}{2.974728in}}%
\pgfpathlineto{\pgfqpoint{0.942158in}{2.980938in}}%
\pgfpathlineto{\pgfqpoint{0.996248in}{2.985770in}}%
\pgfpathlineto{\pgfqpoint{1.034884in}{2.986875in}}%
\pgfpathlineto{\pgfqpoint{1.127610in}{2.988663in}}%
\pgfpathlineto{\pgfqpoint{1.212609in}{2.994098in}}%
\pgfpathlineto{\pgfqpoint{1.529423in}{2.998732in}}%
\pgfpathlineto{\pgfqpoint{1.622149in}{3.000163in}}%
\pgfpathlineto{\pgfqpoint{1.768965in}{3.003083in}}%
\pgfpathlineto{\pgfqpoint{1.846236in}{3.002712in}}%
\pgfpathlineto{\pgfqpoint{1.938962in}{3.006160in}}%
\pgfpathlineto{\pgfqpoint{2.031688in}{3.003462in}}%
\pgfpathlineto{\pgfqpoint{2.101233in}{3.004529in}}%
\pgfpathlineto{\pgfqpoint{2.101233in}{3.004529in}}%
\pgfusepath{stroke}%
\end{pgfscope}%
\begin{pgfscope}%
\pgfpathrectangle{\pgfqpoint{0.473073in}{2.151654in}}{\pgfqpoint{1.705045in}{1.168222in}}%
\pgfusepath{clip}%
\pgfsetbuttcap%
\pgfsetmiterjoin%
\definecolor{currentfill}{rgb}{0.000000,0.000000,0.000000}%
\pgfsetfillcolor{currentfill}%
\pgfsetfillopacity{0.000000}%
\pgfsetlinewidth{0.803000pt}%
\definecolor{currentstroke}{rgb}{0.172549,0.627451,0.172549}%
\pgfsetstrokecolor{currentstroke}%
\pgfsetdash{}{0pt}%
\pgfsys@defobject{currentmarker}{\pgfqpoint{-0.017500in}{-0.017500in}}{\pgfqpoint{0.017500in}{0.017500in}}{%
\pgfpathmoveto{\pgfqpoint{-0.017500in}{-0.017500in}}%
\pgfpathlineto{\pgfqpoint{0.017500in}{-0.017500in}}%
\pgfpathlineto{\pgfqpoint{0.017500in}{0.017500in}}%
\pgfpathlineto{\pgfqpoint{-0.017500in}{0.017500in}}%
\pgfpathlineto{\pgfqpoint{-0.017500in}{-0.017500in}}%
\pgfpathclose%
\pgfusepath{stroke,fill}%
}%
\begin{pgfscope}%
\pgfsys@transformshift{0.563527in}{2.854226in}%
\pgfsys@useobject{currentmarker}{}%
\end{pgfscope}%
\begin{pgfscope}%
\pgfsys@transformshift{0.653980in}{2.926682in}%
\pgfsys@useobject{currentmarker}{}%
\end{pgfscope}%
\begin{pgfscope}%
\pgfsys@transformshift{0.744433in}{2.955590in}%
\pgfsys@useobject{currentmarker}{}%
\end{pgfscope}%
\begin{pgfscope}%
\pgfsys@transformshift{0.834886in}{2.967397in}%
\pgfsys@useobject{currentmarker}{}%
\end{pgfscope}%
\begin{pgfscope}%
\pgfsys@transformshift{0.925340in}{2.979204in}%
\pgfsys@useobject{currentmarker}{}%
\end{pgfscope}%
\begin{pgfscope}%
\pgfsys@transformshift{1.015793in}{2.986570in}%
\pgfsys@useobject{currentmarker}{}%
\end{pgfscope}%
\begin{pgfscope}%
\pgfsys@transformshift{1.106246in}{2.987644in}%
\pgfsys@useobject{currentmarker}{}%
\end{pgfscope}%
\begin{pgfscope}%
\pgfsys@transformshift{1.196700in}{2.993755in}%
\pgfsys@useobject{currentmarker}{}%
\end{pgfscope}%
\begin{pgfscope}%
\pgfsys@transformshift{1.287153in}{2.995095in}%
\pgfsys@useobject{currentmarker}{}%
\end{pgfscope}%
\begin{pgfscope}%
\pgfsys@transformshift{1.377606in}{2.996104in}%
\pgfsys@useobject{currentmarker}{}%
\end{pgfscope}%
\begin{pgfscope}%
\pgfsys@transformshift{1.468060in}{2.998414in}%
\pgfsys@useobject{currentmarker}{}%
\end{pgfscope}%
\begin{pgfscope}%
\pgfsys@transformshift{1.558513in}{2.998914in}%
\pgfsys@useobject{currentmarker}{}%
\end{pgfscope}%
\begin{pgfscope}%
\pgfsys@transformshift{1.648966in}{3.000840in}%
\pgfsys@useobject{currentmarker}{}%
\end{pgfscope}%
\begin{pgfscope}%
\pgfsys@transformshift{1.739420in}{3.003326in}%
\pgfsys@useobject{currentmarker}{}%
\end{pgfscope}%
\begin{pgfscope}%
\pgfsys@transformshift{1.829873in}{3.002355in}%
\pgfsys@useobject{currentmarker}{}%
\end{pgfscope}%
\begin{pgfscope}%
\pgfsys@transformshift{1.920326in}{3.006496in}%
\pgfsys@useobject{currentmarker}{}%
\end{pgfscope}%
\begin{pgfscope}%
\pgfsys@transformshift{2.010780in}{3.003445in}%
\pgfsys@useobject{currentmarker}{}%
\end{pgfscope}%
\begin{pgfscope}%
\pgfsys@transformshift{2.101233in}{3.004529in}%
\pgfsys@useobject{currentmarker}{}%
\end{pgfscope}%
\end{pgfscope}%
\begin{pgfscope}%
\pgfsetrectcap%
\pgfsetmiterjoin%
\pgfsetlinewidth{0.803000pt}%
\definecolor{currentstroke}{rgb}{0.000000,0.000000,0.000000}%
\pgfsetstrokecolor{currentstroke}%
\pgfsetdash{}{0pt}%
\pgfpathmoveto{\pgfqpoint{0.473073in}{2.151654in}}%
\pgfpathlineto{\pgfqpoint{0.473073in}{3.319877in}}%
\pgfusepath{stroke}%
\end{pgfscope}%
\begin{pgfscope}%
\pgfsetrectcap%
\pgfsetmiterjoin%
\pgfsetlinewidth{0.803000pt}%
\definecolor{currentstroke}{rgb}{0.000000,0.000000,0.000000}%
\pgfsetstrokecolor{currentstroke}%
\pgfsetdash{}{0pt}%
\pgfpathmoveto{\pgfqpoint{2.178118in}{2.151654in}}%
\pgfpathlineto{\pgfqpoint{2.178118in}{3.319877in}}%
\pgfusepath{stroke}%
\end{pgfscope}%
\begin{pgfscope}%
\pgfsetrectcap%
\pgfsetmiterjoin%
\pgfsetlinewidth{0.803000pt}%
\definecolor{currentstroke}{rgb}{0.000000,0.000000,0.000000}%
\pgfsetstrokecolor{currentstroke}%
\pgfsetdash{}{0pt}%
\pgfpathmoveto{\pgfqpoint{0.473073in}{2.151654in}}%
\pgfpathlineto{\pgfqpoint{2.178118in}{2.151654in}}%
\pgfusepath{stroke}%
\end{pgfscope}%
\begin{pgfscope}%
\pgfsetrectcap%
\pgfsetmiterjoin%
\pgfsetlinewidth{0.803000pt}%
\definecolor{currentstroke}{rgb}{0.000000,0.000000,0.000000}%
\pgfsetstrokecolor{currentstroke}%
\pgfsetdash{}{0pt}%
\pgfpathmoveto{\pgfqpoint{0.473073in}{3.319877in}}%
\pgfpathlineto{\pgfqpoint{2.178118in}{3.319877in}}%
\pgfusepath{stroke}%
\end{pgfscope}%
\begin{pgfscope}%
\definecolor{textcolor}{rgb}{0.000000,0.000000,0.000000}%
\pgfsetstrokecolor{textcolor}%
\pgfsetfillcolor{textcolor}%
\pgftext[x=0.524225in,y=3.284830in,left,top]{\color{textcolor}{\rmfamily\fontsize{9.000000}{10.800000}\selectfont\catcode`\^=\active\def^{\ifmmode\sp\else\^{}\fi}\catcode`\%=\active\def%{\%}$(a)$}}%
\end{pgfscope}%
\begin{pgfscope}%
\pgfsetbuttcap%
\pgfsetmiterjoin%
\definecolor{currentfill}{rgb}{1.000000,1.000000,1.000000}%
\pgfsetfillcolor{currentfill}%
\pgfsetlinewidth{0.000000pt}%
\definecolor{currentstroke}{rgb}{0.000000,0.000000,0.000000}%
\pgfsetstrokecolor{currentstroke}%
\pgfsetstrokeopacity{0.000000}%
\pgfsetdash{}{0pt}%
\pgfpathmoveto{\pgfqpoint{2.516573in}{2.151654in}}%
\pgfpathlineto{\pgfqpoint{4.221618in}{2.151654in}}%
\pgfpathlineto{\pgfqpoint{4.221618in}{3.319877in}}%
\pgfpathlineto{\pgfqpoint{2.516573in}{3.319877in}}%
\pgfpathlineto{\pgfqpoint{2.516573in}{2.151654in}}%
\pgfpathclose%
\pgfusepath{fill}%
\end{pgfscope}%
\begin{pgfscope}%
\pgfsetbuttcap%
\pgfsetroundjoin%
\definecolor{currentfill}{rgb}{0.000000,0.000000,0.000000}%
\pgfsetfillcolor{currentfill}%
\pgfsetlinewidth{0.803000pt}%
\definecolor{currentstroke}{rgb}{0.000000,0.000000,0.000000}%
\pgfsetstrokecolor{currentstroke}%
\pgfsetdash{}{0pt}%
\pgfsys@defobject{currentmarker}{\pgfqpoint{0.000000in}{-0.048611in}}{\pgfqpoint{0.000000in}{0.000000in}}{%
\pgfpathmoveto{\pgfqpoint{0.000000in}{0.000000in}}%
\pgfpathlineto{\pgfqpoint{0.000000in}{-0.048611in}}%
\pgfusepath{stroke,fill}%
}%
\begin{pgfscope}%
\pgfsys@transformshift{2.516573in}{2.151654in}%
\pgfsys@useobject{currentmarker}{}%
\end{pgfscope}%
\end{pgfscope}%
\begin{pgfscope}%
\definecolor{textcolor}{rgb}{0.000000,0.000000,0.000000}%
\pgfsetstrokecolor{textcolor}%
\pgfsetfillcolor{textcolor}%
\pgftext[x=2.516573in,y=2.054432in,,top]{\color{textcolor}{\rmfamily\fontsize{8.000000}{9.600000}\selectfont\catcode`\^=\active\def^{\ifmmode\sp\else\^{}\fi}\catcode`\%=\active\def%{\%}$0$}}%
\end{pgfscope}%
\begin{pgfscope}%
\pgfsetbuttcap%
\pgfsetroundjoin%
\definecolor{currentfill}{rgb}{0.000000,0.000000,0.000000}%
\pgfsetfillcolor{currentfill}%
\pgfsetlinewidth{0.803000pt}%
\definecolor{currentstroke}{rgb}{0.000000,0.000000,0.000000}%
\pgfsetstrokecolor{currentstroke}%
\pgfsetdash{}{0pt}%
\pgfsys@defobject{currentmarker}{\pgfqpoint{0.000000in}{-0.048611in}}{\pgfqpoint{0.000000in}{0.000000in}}{%
\pgfpathmoveto{\pgfqpoint{0.000000in}{0.000000in}}%
\pgfpathlineto{\pgfqpoint{0.000000in}{-0.048611in}}%
\pgfusepath{stroke,fill}%
}%
\begin{pgfscope}%
\pgfsys@transformshift{3.119595in}{2.151654in}%
\pgfsys@useobject{currentmarker}{}%
\end{pgfscope}%
\end{pgfscope}%
\begin{pgfscope}%
\definecolor{textcolor}{rgb}{0.000000,0.000000,0.000000}%
\pgfsetstrokecolor{textcolor}%
\pgfsetfillcolor{textcolor}%
\pgftext[x=3.119595in,y=2.054432in,,top]{\color{textcolor}{\rmfamily\fontsize{8.000000}{9.600000}\selectfont\catcode`\^=\active\def^{\ifmmode\sp\else\^{}\fi}\catcode`\%=\active\def%{\%}$10$}}%
\end{pgfscope}%
\begin{pgfscope}%
\pgfsetbuttcap%
\pgfsetroundjoin%
\definecolor{currentfill}{rgb}{0.000000,0.000000,0.000000}%
\pgfsetfillcolor{currentfill}%
\pgfsetlinewidth{0.803000pt}%
\definecolor{currentstroke}{rgb}{0.000000,0.000000,0.000000}%
\pgfsetstrokecolor{currentstroke}%
\pgfsetdash{}{0pt}%
\pgfsys@defobject{currentmarker}{\pgfqpoint{0.000000in}{-0.048611in}}{\pgfqpoint{0.000000in}{0.000000in}}{%
\pgfpathmoveto{\pgfqpoint{0.000000in}{0.000000in}}%
\pgfpathlineto{\pgfqpoint{0.000000in}{-0.048611in}}%
\pgfusepath{stroke,fill}%
}%
\begin{pgfscope}%
\pgfsys@transformshift{3.722618in}{2.151654in}%
\pgfsys@useobject{currentmarker}{}%
\end{pgfscope}%
\end{pgfscope}%
\begin{pgfscope}%
\definecolor{textcolor}{rgb}{0.000000,0.000000,0.000000}%
\pgfsetstrokecolor{textcolor}%
\pgfsetfillcolor{textcolor}%
\pgftext[x=3.722618in,y=2.054432in,,top]{\color{textcolor}{\rmfamily\fontsize{8.000000}{9.600000}\selectfont\catcode`\^=\active\def^{\ifmmode\sp\else\^{}\fi}\catcode`\%=\active\def%{\%}$20$}}%
\end{pgfscope}%
\begin{pgfscope}%
\definecolor{textcolor}{rgb}{0.000000,0.000000,0.000000}%
\pgfsetstrokecolor{textcolor}%
\pgfsetfillcolor{textcolor}%
\pgftext[x=3.369096in,y=1.900111in,,top]{\color{textcolor}{\rmfamily\fontsize{9.000000}{10.800000}\selectfont\catcode`\^=\active\def^{\ifmmode\sp\else\^{}\fi}\catcode`\%=\active\def%{\%}$\kappa_{\rm out}/\kappa$}}%
\end{pgfscope}%
\begin{pgfscope}%
\pgfsetbuttcap%
\pgfsetroundjoin%
\definecolor{currentfill}{rgb}{0.000000,0.000000,0.000000}%
\pgfsetfillcolor{currentfill}%
\pgfsetlinewidth{0.803000pt}%
\definecolor{currentstroke}{rgb}{0.000000,0.000000,0.000000}%
\pgfsetstrokecolor{currentstroke}%
\pgfsetdash{}{0pt}%
\pgfsys@defobject{currentmarker}{\pgfqpoint{-0.048611in}{0.000000in}}{\pgfqpoint{-0.000000in}{0.000000in}}{%
\pgfpathmoveto{\pgfqpoint{-0.000000in}{0.000000in}}%
\pgfpathlineto{\pgfqpoint{-0.048611in}{0.000000in}}%
\pgfusepath{stroke,fill}%
}%
\begin{pgfscope}%
\pgfsys@transformshift{2.516573in}{2.151654in}%
\pgfsys@useobject{currentmarker}{}%
\end{pgfscope}%
\end{pgfscope}%
\begin{pgfscope}%
\definecolor{textcolor}{rgb}{0.000000,0.000000,0.000000}%
\pgfsetstrokecolor{textcolor}%
\pgfsetfillcolor{textcolor}%
\pgftext[x=2.360322in, y=2.113074in, left, base]{\color{textcolor}{\rmfamily\fontsize{8.000000}{9.600000}\selectfont\catcode`\^=\active\def^{\ifmmode\sp\else\^{}\fi}\catcode`\%=\active\def%{\%}$0$}}%
\end{pgfscope}%
\begin{pgfscope}%
\pgfsetbuttcap%
\pgfsetroundjoin%
\definecolor{currentfill}{rgb}{0.000000,0.000000,0.000000}%
\pgfsetfillcolor{currentfill}%
\pgfsetlinewidth{0.803000pt}%
\definecolor{currentstroke}{rgb}{0.000000,0.000000,0.000000}%
\pgfsetstrokecolor{currentstroke}%
\pgfsetdash{}{0pt}%
\pgfsys@defobject{currentmarker}{\pgfqpoint{-0.048611in}{0.000000in}}{\pgfqpoint{-0.000000in}{0.000000in}}{%
\pgfpathmoveto{\pgfqpoint{-0.000000in}{0.000000in}}%
\pgfpathlineto{\pgfqpoint{-0.048611in}{0.000000in}}%
\pgfusepath{stroke,fill}%
}%
\begin{pgfscope}%
\pgfsys@transformshift{2.516573in}{2.638414in}%
\pgfsys@useobject{currentmarker}{}%
\end{pgfscope}%
\end{pgfscope}%
\begin{pgfscope}%
\definecolor{textcolor}{rgb}{0.000000,0.000000,0.000000}%
\pgfsetstrokecolor{textcolor}%
\pgfsetfillcolor{textcolor}%
\pgftext[x=2.268500in, y=2.599833in, left, base]{\color{textcolor}{\rmfamily\fontsize{8.000000}{9.600000}\selectfont\catcode`\^=\active\def^{\ifmmode\sp\else\^{}\fi}\catcode`\%=\active\def%{\%}$0.5$}}%
\end{pgfscope}%
\begin{pgfscope}%
\pgfsetbuttcap%
\pgfsetroundjoin%
\definecolor{currentfill}{rgb}{0.000000,0.000000,0.000000}%
\pgfsetfillcolor{currentfill}%
\pgfsetlinewidth{0.803000pt}%
\definecolor{currentstroke}{rgb}{0.000000,0.000000,0.000000}%
\pgfsetstrokecolor{currentstroke}%
\pgfsetdash{}{0pt}%
\pgfsys@defobject{currentmarker}{\pgfqpoint{-0.048611in}{0.000000in}}{\pgfqpoint{-0.000000in}{0.000000in}}{%
\pgfpathmoveto{\pgfqpoint{-0.000000in}{0.000000in}}%
\pgfpathlineto{\pgfqpoint{-0.048611in}{0.000000in}}%
\pgfusepath{stroke,fill}%
}%
\begin{pgfscope}%
\pgfsys@transformshift{2.516573in}{3.125173in}%
\pgfsys@useobject{currentmarker}{}%
\end{pgfscope}%
\end{pgfscope}%
\begin{pgfscope}%
\definecolor{textcolor}{rgb}{0.000000,0.000000,0.000000}%
\pgfsetstrokecolor{textcolor}%
\pgfsetfillcolor{textcolor}%
\pgftext[x=2.360322in, y=3.086593in, left, base]{\color{textcolor}{\rmfamily\fontsize{8.000000}{9.600000}\selectfont\catcode`\^=\active\def^{\ifmmode\sp\else\^{}\fi}\catcode`\%=\active\def%{\%}$1$}}%
\end{pgfscope}%
\begin{pgfscope}%
\pgfpathrectangle{\pgfqpoint{2.516573in}{2.151654in}}{\pgfqpoint{1.705045in}{1.168222in}}%
\pgfusepath{clip}%
\pgfsetrectcap%
\pgfsetroundjoin%
\pgfsetlinewidth{0.903375pt}%
\definecolor{currentstroke}{rgb}{0.121569,0.466667,0.705882}%
\pgfsetstrokecolor{currentstroke}%
\pgfsetdash{}{0pt}%
\pgfpathmoveto{\pgfqpoint{2.607027in}{2.863389in}}%
\pgfpathlineto{\pgfqpoint{2.637935in}{2.893541in}}%
\pgfpathlineto{\pgfqpoint{2.661117in}{2.912912in}}%
\pgfpathlineto{\pgfqpoint{2.676571in}{2.923755in}}%
\pgfpathlineto{\pgfqpoint{2.692025in}{2.932611in}}%
\pgfpathlineto{\pgfqpoint{2.715207in}{2.942578in}}%
\pgfpathlineto{\pgfqpoint{2.746116in}{2.952955in}}%
\pgfpathlineto{\pgfqpoint{2.784751in}{2.962869in}}%
\pgfpathlineto{\pgfqpoint{2.831114in}{2.972235in}}%
\pgfpathlineto{\pgfqpoint{2.877477in}{2.979346in}}%
\pgfpathlineto{\pgfqpoint{2.939295in}{2.986270in}}%
\pgfpathlineto{\pgfqpoint{3.024294in}{2.993185in}}%
\pgfpathlineto{\pgfqpoint{3.086111in}{2.996133in}}%
\pgfpathlineto{\pgfqpoint{3.217473in}{3.001128in}}%
\pgfpathlineto{\pgfqpoint{3.279290in}{3.002645in}}%
\pgfpathlineto{\pgfqpoint{3.356562in}{3.004707in}}%
\pgfpathlineto{\pgfqpoint{3.449288in}{3.009109in}}%
\pgfpathlineto{\pgfqpoint{3.665649in}{3.013300in}}%
\pgfpathlineto{\pgfqpoint{4.021098in}{3.015278in}}%
\pgfpathlineto{\pgfqpoint{4.067461in}{3.016040in}}%
\pgfpathlineto{\pgfqpoint{4.144733in}{3.014959in}}%
\pgfpathlineto{\pgfqpoint{4.144733in}{3.014959in}}%
\pgfusepath{stroke}%
\end{pgfscope}%
\begin{pgfscope}%
\pgfpathrectangle{\pgfqpoint{2.516573in}{2.151654in}}{\pgfqpoint{1.705045in}{1.168222in}}%
\pgfusepath{clip}%
\pgfsetbuttcap%
\pgfsetroundjoin%
\definecolor{currentfill}{rgb}{0.121569,0.466667,0.705882}%
\pgfsetfillcolor{currentfill}%
\pgfsetlinewidth{0.803000pt}%
\definecolor{currentstroke}{rgb}{0.121569,0.466667,0.705882}%
\pgfsetstrokecolor{currentstroke}%
\pgfsetdash{}{0pt}%
\pgfsys@defobject{currentmarker}{\pgfqpoint{-0.019444in}{-0.019444in}}{\pgfqpoint{0.019444in}{0.019444in}}{%
\pgfpathmoveto{\pgfqpoint{-0.019444in}{-0.019444in}}%
\pgfpathlineto{\pgfqpoint{0.019444in}{0.019444in}}%
\pgfpathmoveto{\pgfqpoint{-0.019444in}{0.019444in}}%
\pgfpathlineto{\pgfqpoint{0.019444in}{-0.019444in}}%
\pgfusepath{stroke,fill}%
}%
\begin{pgfscope}%
\pgfsys@transformshift{2.607027in}{2.863389in}%
\pgfsys@useobject{currentmarker}{}%
\end{pgfscope}%
\begin{pgfscope}%
\pgfsys@transformshift{2.697480in}{2.935210in}%
\pgfsys@useobject{currentmarker}{}%
\end{pgfscope}%
\begin{pgfscope}%
\pgfsys@transformshift{2.787933in}{2.963590in}%
\pgfsys@useobject{currentmarker}{}%
\end{pgfscope}%
\begin{pgfscope}%
\pgfsys@transformshift{2.878386in}{2.979465in}%
\pgfsys@useobject{currentmarker}{}%
\end{pgfscope}%
\begin{pgfscope}%
\pgfsys@transformshift{2.968840in}{2.988895in}%
\pgfsys@useobject{currentmarker}{}%
\end{pgfscope}%
\begin{pgfscope}%
\pgfsys@transformshift{3.059293in}{2.995153in}%
\pgfsys@useobject{currentmarker}{}%
\end{pgfscope}%
\begin{pgfscope}%
\pgfsys@transformshift{3.149746in}{2.998022in}%
\pgfsys@useobject{currentmarker}{}%
\end{pgfscope}%
\begin{pgfscope}%
\pgfsys@transformshift{3.240200in}{3.001915in}%
\pgfsys@useobject{currentmarker}{}%
\end{pgfscope}%
\begin{pgfscope}%
\pgfsys@transformshift{3.330653in}{3.003586in}%
\pgfsys@useobject{currentmarker}{}%
\end{pgfscope}%
\begin{pgfscope}%
\pgfsys@transformshift{3.421106in}{3.008301in}%
\pgfsys@useobject{currentmarker}{}%
\end{pgfscope}%
\begin{pgfscope}%
\pgfsys@transformshift{3.511560in}{3.010532in}%
\pgfsys@useobject{currentmarker}{}%
\end{pgfscope}%
\begin{pgfscope}%
\pgfsys@transformshift{3.602013in}{3.012920in}%
\pgfsys@useobject{currentmarker}{}%
\end{pgfscope}%
\begin{pgfscope}%
\pgfsys@transformshift{3.692466in}{3.013346in}%
\pgfsys@useobject{currentmarker}{}%
\end{pgfscope}%
\begin{pgfscope}%
\pgfsys@transformshift{3.782920in}{3.013207in}%
\pgfsys@useobject{currentmarker}{}%
\end{pgfscope}%
\begin{pgfscope}%
\pgfsys@transformshift{3.873373in}{3.014328in}%
\pgfsys@useobject{currentmarker}{}%
\end{pgfscope}%
\begin{pgfscope}%
\pgfsys@transformshift{3.963826in}{3.013517in}%
\pgfsys@useobject{currentmarker}{}%
\end{pgfscope}%
\begin{pgfscope}%
\pgfsys@transformshift{4.054280in}{3.016050in}%
\pgfsys@useobject{currentmarker}{}%
\end{pgfscope}%
\begin{pgfscope}%
\pgfsys@transformshift{4.144733in}{3.014959in}%
\pgfsys@useobject{currentmarker}{}%
\end{pgfscope}%
\end{pgfscope}%
\begin{pgfscope}%
\pgfpathrectangle{\pgfqpoint{2.516573in}{2.151654in}}{\pgfqpoint{1.705045in}{1.168222in}}%
\pgfusepath{clip}%
\pgfsetrectcap%
\pgfsetroundjoin%
\pgfsetlinewidth{0.903375pt}%
\definecolor{currentstroke}{rgb}{1.000000,0.000000,0.000000}%
\pgfsetstrokecolor{currentstroke}%
\pgfsetdash{}{0pt}%
\pgfpathmoveto{\pgfqpoint{2.607027in}{2.858289in}}%
\pgfpathlineto{\pgfqpoint{2.637935in}{2.887812in}}%
\pgfpathlineto{\pgfqpoint{2.661117in}{2.906803in}}%
\pgfpathlineto{\pgfqpoint{2.676571in}{2.917455in}}%
\pgfpathlineto{\pgfqpoint{2.692025in}{2.926185in}}%
\pgfpathlineto{\pgfqpoint{2.715207in}{2.936086in}}%
\pgfpathlineto{\pgfqpoint{2.746116in}{2.946470in}}%
\pgfpathlineto{\pgfqpoint{2.777024in}{2.954542in}}%
\pgfpathlineto{\pgfqpoint{2.815660in}{2.962527in}}%
\pgfpathlineto{\pgfqpoint{2.862023in}{2.969921in}}%
\pgfpathlineto{\pgfqpoint{2.931568in}{2.978417in}}%
\pgfpathlineto{\pgfqpoint{3.001112in}{2.984528in}}%
\pgfpathlineto{\pgfqpoint{3.070657in}{2.988703in}}%
\pgfpathlineto{\pgfqpoint{3.279290in}{2.996205in}}%
\pgfpathlineto{\pgfqpoint{3.364289in}{2.998744in}}%
\pgfpathlineto{\pgfqpoint{3.449288in}{3.002114in}}%
\pgfpathlineto{\pgfqpoint{3.650194in}{3.005883in}}%
\pgfpathlineto{\pgfqpoint{3.959281in}{3.007409in}}%
\pgfpathlineto{\pgfqpoint{4.005644in}{3.008480in}}%
\pgfpathlineto{\pgfqpoint{4.059734in}{3.009826in}}%
\pgfpathlineto{\pgfqpoint{4.144733in}{3.008888in}}%
\pgfpathlineto{\pgfqpoint{4.144733in}{3.008888in}}%
\pgfusepath{stroke}%
\end{pgfscope}%
\begin{pgfscope}%
\pgfpathrectangle{\pgfqpoint{2.516573in}{2.151654in}}{\pgfqpoint{1.705045in}{1.168222in}}%
\pgfusepath{clip}%
\pgfsetbuttcap%
\pgfsetroundjoin%
\definecolor{currentfill}{rgb}{0.000000,0.000000,0.000000}%
\pgfsetfillcolor{currentfill}%
\pgfsetfillopacity{0.000000}%
\pgfsetlinewidth{0.803000pt}%
\definecolor{currentstroke}{rgb}{1.000000,0.000000,0.000000}%
\pgfsetstrokecolor{currentstroke}%
\pgfsetdash{}{0pt}%
\pgfsys@defobject{currentmarker}{\pgfqpoint{-0.017500in}{-0.017500in}}{\pgfqpoint{0.017500in}{0.017500in}}{%
\pgfpathmoveto{\pgfqpoint{0.000000in}{-0.017500in}}%
\pgfpathcurveto{\pgfqpoint{0.004641in}{-0.017500in}}{\pgfqpoint{0.009093in}{-0.015656in}}{\pgfqpoint{0.012374in}{-0.012374in}}%
\pgfpathcurveto{\pgfqpoint{0.015656in}{-0.009093in}}{\pgfqpoint{0.017500in}{-0.004641in}}{\pgfqpoint{0.017500in}{0.000000in}}%
\pgfpathcurveto{\pgfqpoint{0.017500in}{0.004641in}}{\pgfqpoint{0.015656in}{0.009093in}}{\pgfqpoint{0.012374in}{0.012374in}}%
\pgfpathcurveto{\pgfqpoint{0.009093in}{0.015656in}}{\pgfqpoint{0.004641in}{0.017500in}}{\pgfqpoint{0.000000in}{0.017500in}}%
\pgfpathcurveto{\pgfqpoint{-0.004641in}{0.017500in}}{\pgfqpoint{-0.009093in}{0.015656in}}{\pgfqpoint{-0.012374in}{0.012374in}}%
\pgfpathcurveto{\pgfqpoint{-0.015656in}{0.009093in}}{\pgfqpoint{-0.017500in}{0.004641in}}{\pgfqpoint{-0.017500in}{0.000000in}}%
\pgfpathcurveto{\pgfqpoint{-0.017500in}{-0.004641in}}{\pgfqpoint{-0.015656in}{-0.009093in}}{\pgfqpoint{-0.012374in}{-0.012374in}}%
\pgfpathcurveto{\pgfqpoint{-0.009093in}{-0.015656in}}{\pgfqpoint{-0.004641in}{-0.017500in}}{\pgfqpoint{0.000000in}{-0.017500in}}%
\pgfpathlineto{\pgfqpoint{0.000000in}{-0.017500in}}%
\pgfpathclose%
\pgfusepath{stroke,fill}%
}%
\begin{pgfscope}%
\pgfsys@transformshift{2.607027in}{2.858289in}%
\pgfsys@useobject{currentmarker}{}%
\end{pgfscope}%
\begin{pgfscope}%
\pgfsys@transformshift{2.697480in}{2.928757in}%
\pgfsys@useobject{currentmarker}{}%
\end{pgfscope}%
\begin{pgfscope}%
\pgfsys@transformshift{2.787933in}{2.956998in}%
\pgfsys@useobject{currentmarker}{}%
\end{pgfscope}%
\begin{pgfscope}%
\pgfsys@transformshift{2.878386in}{2.972143in}%
\pgfsys@useobject{currentmarker}{}%
\end{pgfscope}%
\begin{pgfscope}%
\pgfsys@transformshift{2.968840in}{2.981935in}%
\pgfsys@useobject{currentmarker}{}%
\end{pgfscope}%
\begin{pgfscope}%
\pgfsys@transformshift{3.059293in}{2.988192in}%
\pgfsys@useobject{currentmarker}{}%
\end{pgfscope}%
\begin{pgfscope}%
\pgfsys@transformshift{3.149746in}{2.991501in}%
\pgfsys@useobject{currentmarker}{}%
\end{pgfscope}%
\begin{pgfscope}%
\pgfsys@transformshift{3.240200in}{2.995312in}%
\pgfsys@useobject{currentmarker}{}%
\end{pgfscope}%
\begin{pgfscope}%
\pgfsys@transformshift{3.330653in}{2.997341in}%
\pgfsys@useobject{currentmarker}{}%
\end{pgfscope}%
\begin{pgfscope}%
\pgfsys@transformshift{3.421106in}{3.001345in}%
\pgfsys@useobject{currentmarker}{}%
\end{pgfscope}%
\begin{pgfscope}%
\pgfsys@transformshift{3.511560in}{3.003478in}%
\pgfsys@useobject{currentmarker}{}%
\end{pgfscope}%
\begin{pgfscope}%
\pgfsys@transformshift{3.602013in}{3.005374in}%
\pgfsys@useobject{currentmarker}{}%
\end{pgfscope}%
\begin{pgfscope}%
\pgfsys@transformshift{3.692466in}{3.006208in}%
\pgfsys@useobject{currentmarker}{}%
\end{pgfscope}%
\begin{pgfscope}%
\pgfsys@transformshift{3.782920in}{3.006718in}%
\pgfsys@useobject{currentmarker}{}%
\end{pgfscope}%
\begin{pgfscope}%
\pgfsys@transformshift{3.873373in}{3.007773in}%
\pgfsys@useobject{currentmarker}{}%
\end{pgfscope}%
\begin{pgfscope}%
\pgfsys@transformshift{3.963826in}{3.007406in}%
\pgfsys@useobject{currentmarker}{}%
\end{pgfscope}%
\begin{pgfscope}%
\pgfsys@transformshift{4.054280in}{3.009827in}%
\pgfsys@useobject{currentmarker}{}%
\end{pgfscope}%
\begin{pgfscope}%
\pgfsys@transformshift{4.144733in}{3.008888in}%
\pgfsys@useobject{currentmarker}{}%
\end{pgfscope}%
\end{pgfscope}%
\begin{pgfscope}%
\pgfpathrectangle{\pgfqpoint{2.516573in}{2.151654in}}{\pgfqpoint{1.705045in}{1.168222in}}%
\pgfusepath{clip}%
\pgfsetrectcap%
\pgfsetroundjoin%
\pgfsetlinewidth{0.903375pt}%
\definecolor{currentstroke}{rgb}{0.172549,0.627451,0.172549}%
\pgfsetstrokecolor{currentstroke}%
\pgfsetdash{}{0pt}%
\pgfpathmoveto{\pgfqpoint{2.607027in}{2.832103in}}%
\pgfpathlineto{\pgfqpoint{2.637935in}{2.860594in}}%
\pgfpathlineto{\pgfqpoint{2.661117in}{2.878832in}}%
\pgfpathlineto{\pgfqpoint{2.676571in}{2.888968in}}%
\pgfpathlineto{\pgfqpoint{2.692025in}{2.897155in}}%
\pgfpathlineto{\pgfqpoint{2.715207in}{2.906148in}}%
\pgfpathlineto{\pgfqpoint{2.746116in}{2.915305in}}%
\pgfpathlineto{\pgfqpoint{2.792479in}{2.925945in}}%
\pgfpathlineto{\pgfqpoint{2.846569in}{2.936417in}}%
\pgfpathlineto{\pgfqpoint{2.885205in}{2.941673in}}%
\pgfpathlineto{\pgfqpoint{2.947022in}{2.946990in}}%
\pgfpathlineto{\pgfqpoint{3.086111in}{2.956402in}}%
\pgfpathlineto{\pgfqpoint{3.178837in}{2.959730in}}%
\pgfpathlineto{\pgfqpoint{3.341108in}{2.964551in}}%
\pgfpathlineto{\pgfqpoint{3.449288in}{2.968204in}}%
\pgfpathlineto{\pgfqpoint{3.572923in}{2.970472in}}%
\pgfpathlineto{\pgfqpoint{3.696557in}{2.973097in}}%
\pgfpathlineto{\pgfqpoint{3.835646in}{2.973182in}}%
\pgfpathlineto{\pgfqpoint{3.889736in}{2.973992in}}%
\pgfpathlineto{\pgfqpoint{4.028825in}{2.974420in}}%
\pgfpathlineto{\pgfqpoint{4.144733in}{2.975958in}}%
\pgfpathlineto{\pgfqpoint{4.144733in}{2.975958in}}%
\pgfusepath{stroke}%
\end{pgfscope}%
\begin{pgfscope}%
\pgfpathrectangle{\pgfqpoint{2.516573in}{2.151654in}}{\pgfqpoint{1.705045in}{1.168222in}}%
\pgfusepath{clip}%
\pgfsetbuttcap%
\pgfsetmiterjoin%
\definecolor{currentfill}{rgb}{0.000000,0.000000,0.000000}%
\pgfsetfillcolor{currentfill}%
\pgfsetfillopacity{0.000000}%
\pgfsetlinewidth{0.803000pt}%
\definecolor{currentstroke}{rgb}{0.172549,0.627451,0.172549}%
\pgfsetstrokecolor{currentstroke}%
\pgfsetdash{}{0pt}%
\pgfsys@defobject{currentmarker}{\pgfqpoint{-0.017500in}{-0.017500in}}{\pgfqpoint{0.017500in}{0.017500in}}{%
\pgfpathmoveto{\pgfqpoint{-0.017500in}{-0.017500in}}%
\pgfpathlineto{\pgfqpoint{0.017500in}{-0.017500in}}%
\pgfpathlineto{\pgfqpoint{0.017500in}{0.017500in}}%
\pgfpathlineto{\pgfqpoint{-0.017500in}{0.017500in}}%
\pgfpathlineto{\pgfqpoint{-0.017500in}{-0.017500in}}%
\pgfpathclose%
\pgfusepath{stroke,fill}%
}%
\begin{pgfscope}%
\pgfsys@transformshift{2.607027in}{2.832103in}%
\pgfsys@useobject{currentmarker}{}%
\end{pgfscope}%
\begin{pgfscope}%
\pgfsys@transformshift{2.697480in}{2.899526in}%
\pgfsys@useobject{currentmarker}{}%
\end{pgfscope}%
\begin{pgfscope}%
\pgfsys@transformshift{2.787933in}{2.924964in}%
\pgfsys@useobject{currentmarker}{}%
\end{pgfscope}%
\begin{pgfscope}%
\pgfsys@transformshift{2.878386in}{2.940914in}%
\pgfsys@useobject{currentmarker}{}%
\end{pgfscope}%
\begin{pgfscope}%
\pgfsys@transformshift{2.968840in}{2.948634in}%
\pgfsys@useobject{currentmarker}{}%
\end{pgfscope}%
\begin{pgfscope}%
\pgfsys@transformshift{3.059293in}{2.955045in}%
\pgfsys@useobject{currentmarker}{}%
\end{pgfscope}%
\begin{pgfscope}%
\pgfsys@transformshift{3.149746in}{2.958935in}%
\pgfsys@useobject{currentmarker}{}%
\end{pgfscope}%
\begin{pgfscope}%
\pgfsys@transformshift{3.240200in}{2.961196in}%
\pgfsys@useobject{currentmarker}{}%
\end{pgfscope}%
\begin{pgfscope}%
\pgfsys@transformshift{3.330653in}{2.964141in}%
\pgfsys@useobject{currentmarker}{}%
\end{pgfscope}%
\begin{pgfscope}%
\pgfsys@transformshift{3.421106in}{2.967801in}%
\pgfsys@useobject{currentmarker}{}%
\end{pgfscope}%
\begin{pgfscope}%
\pgfsys@transformshift{3.511560in}{2.968901in}%
\pgfsys@useobject{currentmarker}{}%
\end{pgfscope}%
\begin{pgfscope}%
\pgfsys@transformshift{3.602013in}{2.971245in}%
\pgfsys@useobject{currentmarker}{}%
\end{pgfscope}%
\begin{pgfscope}%
\pgfsys@transformshift{3.692466in}{2.973104in}%
\pgfsys@useobject{currentmarker}{}%
\end{pgfscope}%
\begin{pgfscope}%
\pgfsys@transformshift{3.782920in}{2.971782in}%
\pgfsys@useobject{currentmarker}{}%
\end{pgfscope}%
\begin{pgfscope}%
\pgfsys@transformshift{3.873373in}{2.974028in}%
\pgfsys@useobject{currentmarker}{}%
\end{pgfscope}%
\begin{pgfscope}%
\pgfsys@transformshift{3.963826in}{2.973612in}%
\pgfsys@useobject{currentmarker}{}%
\end{pgfscope}%
\begin{pgfscope}%
\pgfsys@transformshift{4.054280in}{2.974824in}%
\pgfsys@useobject{currentmarker}{}%
\end{pgfscope}%
\begin{pgfscope}%
\pgfsys@transformshift{4.144733in}{2.975958in}%
\pgfsys@useobject{currentmarker}{}%
\end{pgfscope}%
\end{pgfscope}%
\begin{pgfscope}%
\pgfpathrectangle{\pgfqpoint{2.516573in}{2.151654in}}{\pgfqpoint{1.705045in}{1.168222in}}%
\pgfusepath{clip}%
\pgfsetrectcap%
\pgfsetroundjoin%
\pgfsetlinewidth{0.903375pt}%
\definecolor{currentstroke}{rgb}{0.000000,0.000000,0.000000}%
\pgfsetstrokecolor{currentstroke}%
\pgfsetdash{}{0pt}%
\pgfpathmoveto{\pgfqpoint{2.516573in}{3.079918in}}%
\pgfpathlineto{\pgfqpoint{4.221618in}{3.079918in}}%
\pgfusepath{stroke}%
\end{pgfscope}%
\begin{pgfscope}%
\pgfsetrectcap%
\pgfsetmiterjoin%
\pgfsetlinewidth{0.803000pt}%
\definecolor{currentstroke}{rgb}{0.000000,0.000000,0.000000}%
\pgfsetstrokecolor{currentstroke}%
\pgfsetdash{}{0pt}%
\pgfpathmoveto{\pgfqpoint{2.516573in}{2.151654in}}%
\pgfpathlineto{\pgfqpoint{2.516573in}{3.319877in}}%
\pgfusepath{stroke}%
\end{pgfscope}%
\begin{pgfscope}%
\pgfsetrectcap%
\pgfsetmiterjoin%
\pgfsetlinewidth{0.803000pt}%
\definecolor{currentstroke}{rgb}{0.000000,0.000000,0.000000}%
\pgfsetstrokecolor{currentstroke}%
\pgfsetdash{}{0pt}%
\pgfpathmoveto{\pgfqpoint{4.221618in}{2.151654in}}%
\pgfpathlineto{\pgfqpoint{4.221618in}{3.319877in}}%
\pgfusepath{stroke}%
\end{pgfscope}%
\begin{pgfscope}%
\pgfsetrectcap%
\pgfsetmiterjoin%
\pgfsetlinewidth{0.803000pt}%
\definecolor{currentstroke}{rgb}{0.000000,0.000000,0.000000}%
\pgfsetstrokecolor{currentstroke}%
\pgfsetdash{}{0pt}%
\pgfpathmoveto{\pgfqpoint{2.516573in}{2.151654in}}%
\pgfpathlineto{\pgfqpoint{4.221618in}{2.151654in}}%
\pgfusepath{stroke}%
\end{pgfscope}%
\begin{pgfscope}%
\pgfsetrectcap%
\pgfsetmiterjoin%
\pgfsetlinewidth{0.803000pt}%
\definecolor{currentstroke}{rgb}{0.000000,0.000000,0.000000}%
\pgfsetstrokecolor{currentstroke}%
\pgfsetdash{}{0pt}%
\pgfpathmoveto{\pgfqpoint{2.516573in}{3.319877in}}%
\pgfpathlineto{\pgfqpoint{4.221618in}{3.319877in}}%
\pgfusepath{stroke}%
\end{pgfscope}%
\begin{pgfscope}%
\definecolor{textcolor}{rgb}{0.000000,0.000000,0.000000}%
\pgfsetstrokecolor{textcolor}%
\pgfsetfillcolor{textcolor}%
\pgftext[x=2.567725in,y=3.284830in,left,top]{\color{textcolor}{\rmfamily\fontsize{9.000000}{10.800000}\selectfont\catcode`\^=\active\def^{\ifmmode\sp\else\^{}\fi}\catcode`\%=\active\def%{\%}$(b)$}}%
\end{pgfscope}%
\begin{pgfscope}%
\pgfsetbuttcap%
\pgfsetmiterjoin%
\definecolor{currentfill}{rgb}{1.000000,1.000000,1.000000}%
\pgfsetfillcolor{currentfill}%
\pgfsetlinewidth{0.000000pt}%
\definecolor{currentstroke}{rgb}{0.000000,0.000000,0.000000}%
\pgfsetstrokecolor{currentstroke}%
\pgfsetstrokeopacity{0.000000}%
\pgfsetdash{}{0pt}%
\pgfpathmoveto{\pgfqpoint{0.473073in}{0.462654in}}%
\pgfpathlineto{\pgfqpoint{2.178118in}{0.462654in}}%
\pgfpathlineto{\pgfqpoint{2.178118in}{1.630877in}}%
\pgfpathlineto{\pgfqpoint{0.473073in}{1.630877in}}%
\pgfpathlineto{\pgfqpoint{0.473073in}{0.462654in}}%
\pgfpathclose%
\pgfusepath{fill}%
\end{pgfscope}%
\begin{pgfscope}%
\pgfsetbuttcap%
\pgfsetroundjoin%
\definecolor{currentfill}{rgb}{0.000000,0.000000,0.000000}%
\pgfsetfillcolor{currentfill}%
\pgfsetlinewidth{0.803000pt}%
\definecolor{currentstroke}{rgb}{0.000000,0.000000,0.000000}%
\pgfsetstrokecolor{currentstroke}%
\pgfsetdash{}{0pt}%
\pgfsys@defobject{currentmarker}{\pgfqpoint{0.000000in}{-0.048611in}}{\pgfqpoint{0.000000in}{0.000000in}}{%
\pgfpathmoveto{\pgfqpoint{0.000000in}{0.000000in}}%
\pgfpathlineto{\pgfqpoint{0.000000in}{-0.048611in}}%
\pgfusepath{stroke,fill}%
}%
\begin{pgfscope}%
\pgfsys@transformshift{0.550575in}{0.462654in}%
\pgfsys@useobject{currentmarker}{}%
\end{pgfscope}%
\end{pgfscope}%
\begin{pgfscope}%
\definecolor{textcolor}{rgb}{0.000000,0.000000,0.000000}%
\pgfsetstrokecolor{textcolor}%
\pgfsetfillcolor{textcolor}%
\pgftext[x=0.550575in,y=0.365432in,,top]{\color{textcolor}{\rmfamily\fontsize{8.000000}{9.600000}\selectfont\catcode`\^=\active\def^{\ifmmode\sp\else\^{}\fi}\catcode`\%=\active\def%{\%}$0$}}%
\end{pgfscope}%
\begin{pgfscope}%
\pgfsetbuttcap%
\pgfsetroundjoin%
\definecolor{currentfill}{rgb}{0.000000,0.000000,0.000000}%
\pgfsetfillcolor{currentfill}%
\pgfsetlinewidth{0.803000pt}%
\definecolor{currentstroke}{rgb}{0.000000,0.000000,0.000000}%
\pgfsetstrokecolor{currentstroke}%
\pgfsetdash{}{0pt}%
\pgfsys@defobject{currentmarker}{\pgfqpoint{0.000000in}{-0.048611in}}{\pgfqpoint{0.000000in}{0.000000in}}{%
\pgfpathmoveto{\pgfqpoint{0.000000in}{0.000000in}}%
\pgfpathlineto{\pgfqpoint{0.000000in}{-0.048611in}}%
\pgfusepath{stroke,fill}%
}%
\begin{pgfscope}%
\pgfsys@transformshift{1.325596in}{0.462654in}%
\pgfsys@useobject{currentmarker}{}%
\end{pgfscope}%
\end{pgfscope}%
\begin{pgfscope}%
\definecolor{textcolor}{rgb}{0.000000,0.000000,0.000000}%
\pgfsetstrokecolor{textcolor}%
\pgfsetfillcolor{textcolor}%
\pgftext[x=1.325596in,y=0.365432in,,top]{\color{textcolor}{\rmfamily\fontsize{8.000000}{9.600000}\selectfont\catcode`\^=\active\def^{\ifmmode\sp\else\^{}\fi}\catcode`\%=\active\def%{\%}$0.5$}}%
\end{pgfscope}%
\begin{pgfscope}%
\pgfsetbuttcap%
\pgfsetroundjoin%
\definecolor{currentfill}{rgb}{0.000000,0.000000,0.000000}%
\pgfsetfillcolor{currentfill}%
\pgfsetlinewidth{0.803000pt}%
\definecolor{currentstroke}{rgb}{0.000000,0.000000,0.000000}%
\pgfsetstrokecolor{currentstroke}%
\pgfsetdash{}{0pt}%
\pgfsys@defobject{currentmarker}{\pgfqpoint{0.000000in}{-0.048611in}}{\pgfqpoint{0.000000in}{0.000000in}}{%
\pgfpathmoveto{\pgfqpoint{0.000000in}{0.000000in}}%
\pgfpathlineto{\pgfqpoint{0.000000in}{-0.048611in}}%
\pgfusepath{stroke,fill}%
}%
\begin{pgfscope}%
\pgfsys@transformshift{2.100616in}{0.462654in}%
\pgfsys@useobject{currentmarker}{}%
\end{pgfscope}%
\end{pgfscope}%
\begin{pgfscope}%
\definecolor{textcolor}{rgb}{0.000000,0.000000,0.000000}%
\pgfsetstrokecolor{textcolor}%
\pgfsetfillcolor{textcolor}%
\pgftext[x=2.100616in,y=0.365432in,,top]{\color{textcolor}{\rmfamily\fontsize{8.000000}{9.600000}\selectfont\catcode`\^=\active\def^{\ifmmode\sp\else\^{}\fi}\catcode`\%=\active\def%{\%}$1$}}%
\end{pgfscope}%
\begin{pgfscope}%
\definecolor{textcolor}{rgb}{0.000000,0.000000,0.000000}%
\pgfsetstrokecolor{textcolor}%
\pgfsetfillcolor{textcolor}%
\pgftext[x=1.325596in,y=0.211111in,,top]{\color{textcolor}{\rmfamily\fontsize{9.000000}{10.800000}\selectfont\catcode`\^=\active\def^{\ifmmode\sp\else\^{}\fi}\catcode`\%=\active\def%{\%}$\chi$}}%
\end{pgfscope}%
\begin{pgfscope}%
\pgfsetbuttcap%
\pgfsetroundjoin%
\definecolor{currentfill}{rgb}{0.000000,0.000000,0.000000}%
\pgfsetfillcolor{currentfill}%
\pgfsetlinewidth{0.803000pt}%
\definecolor{currentstroke}{rgb}{0.000000,0.000000,0.000000}%
\pgfsetstrokecolor{currentstroke}%
\pgfsetdash{}{0pt}%
\pgfsys@defobject{currentmarker}{\pgfqpoint{-0.048611in}{0.000000in}}{\pgfqpoint{-0.000000in}{0.000000in}}{%
\pgfpathmoveto{\pgfqpoint{-0.000000in}{0.000000in}}%
\pgfpathlineto{\pgfqpoint{-0.048611in}{0.000000in}}%
\pgfusepath{stroke,fill}%
}%
\begin{pgfscope}%
\pgfsys@transformshift{0.473073in}{0.462654in}%
\pgfsys@useobject{currentmarker}{}%
\end{pgfscope}%
\end{pgfscope}%
\begin{pgfscope}%
\definecolor{textcolor}{rgb}{0.000000,0.000000,0.000000}%
\pgfsetstrokecolor{textcolor}%
\pgfsetfillcolor{textcolor}%
\pgftext[x=0.316822in, y=0.424074in, left, base]{\color{textcolor}{\rmfamily\fontsize{8.000000}{9.600000}\selectfont\catcode`\^=\active\def^{\ifmmode\sp\else\^{}\fi}\catcode`\%=\active\def%{\%}$0$}}%
\end{pgfscope}%
\begin{pgfscope}%
\pgfsetbuttcap%
\pgfsetroundjoin%
\definecolor{currentfill}{rgb}{0.000000,0.000000,0.000000}%
\pgfsetfillcolor{currentfill}%
\pgfsetlinewidth{0.803000pt}%
\definecolor{currentstroke}{rgb}{0.000000,0.000000,0.000000}%
\pgfsetstrokecolor{currentstroke}%
\pgfsetdash{}{0pt}%
\pgfsys@defobject{currentmarker}{\pgfqpoint{-0.048611in}{0.000000in}}{\pgfqpoint{-0.000000in}{0.000000in}}{%
\pgfpathmoveto{\pgfqpoint{-0.000000in}{0.000000in}}%
\pgfpathlineto{\pgfqpoint{-0.048611in}{0.000000in}}%
\pgfusepath{stroke,fill}%
}%
\begin{pgfscope}%
\pgfsys@transformshift{0.473073in}{0.949414in}%
\pgfsys@useobject{currentmarker}{}%
\end{pgfscope}%
\end{pgfscope}%
\begin{pgfscope}%
\definecolor{textcolor}{rgb}{0.000000,0.000000,0.000000}%
\pgfsetstrokecolor{textcolor}%
\pgfsetfillcolor{textcolor}%
\pgftext[x=0.225000in, y=0.910833in, left, base]{\color{textcolor}{\rmfamily\fontsize{8.000000}{9.600000}\selectfont\catcode`\^=\active\def^{\ifmmode\sp\else\^{}\fi}\catcode`\%=\active\def%{\%}$0.5$}}%
\end{pgfscope}%
\begin{pgfscope}%
\pgfsetbuttcap%
\pgfsetroundjoin%
\definecolor{currentfill}{rgb}{0.000000,0.000000,0.000000}%
\pgfsetfillcolor{currentfill}%
\pgfsetlinewidth{0.803000pt}%
\definecolor{currentstroke}{rgb}{0.000000,0.000000,0.000000}%
\pgfsetstrokecolor{currentstroke}%
\pgfsetdash{}{0pt}%
\pgfsys@defobject{currentmarker}{\pgfqpoint{-0.048611in}{0.000000in}}{\pgfqpoint{-0.000000in}{0.000000in}}{%
\pgfpathmoveto{\pgfqpoint{-0.000000in}{0.000000in}}%
\pgfpathlineto{\pgfqpoint{-0.048611in}{0.000000in}}%
\pgfusepath{stroke,fill}%
}%
\begin{pgfscope}%
\pgfsys@transformshift{0.473073in}{1.436173in}%
\pgfsys@useobject{currentmarker}{}%
\end{pgfscope}%
\end{pgfscope}%
\begin{pgfscope}%
\definecolor{textcolor}{rgb}{0.000000,0.000000,0.000000}%
\pgfsetstrokecolor{textcolor}%
\pgfsetfillcolor{textcolor}%
\pgftext[x=0.316822in, y=1.397593in, left, base]{\color{textcolor}{\rmfamily\fontsize{8.000000}{9.600000}\selectfont\catcode`\^=\active\def^{\ifmmode\sp\else\^{}\fi}\catcode`\%=\active\def%{\%}$1$}}%
\end{pgfscope}%
\begin{pgfscope}%
\definecolor{textcolor}{rgb}{0.000000,0.000000,0.000000}%
\pgfsetstrokecolor{textcolor}%
\pgfsetfillcolor{textcolor}%
\pgftext[x=0.211111in,y=1.046766in,,bottom,rotate=90.000000]{\color{textcolor}{\rmfamily\fontsize{9.000000}{10.800000}\selectfont\catcode`\^=\active\def^{\ifmmode\sp\else\^{}\fi}\catcode`\%=\active\def%{\%}$F_h$}}%
\end{pgfscope}%
\begin{pgfscope}%
\pgfpathrectangle{\pgfqpoint{0.473073in}{0.462654in}}{\pgfqpoint{1.705045in}{1.168222in}}%
\pgfusepath{clip}%
\pgfsetrectcap%
\pgfsetroundjoin%
\pgfsetlinewidth{0.903375pt}%
\definecolor{currentstroke}{rgb}{0.121569,0.466667,0.705882}%
\pgfsetstrokecolor{currentstroke}%
\pgfsetdash{}{0pt}%
\pgfpathmoveto{\pgfqpoint{0.550575in}{1.252512in}}%
\pgfpathlineto{\pgfqpoint{1.321701in}{1.283277in}}%
\pgfpathlineto{\pgfqpoint{1.937044in}{1.305648in}}%
\pgfpathlineto{\pgfqpoint{2.100616in}{1.311162in}}%
\pgfpathlineto{\pgfqpoint{2.100616in}{1.311162in}}%
\pgfusepath{stroke}%
\end{pgfscope}%
\begin{pgfscope}%
\pgfpathrectangle{\pgfqpoint{0.473073in}{0.462654in}}{\pgfqpoint{1.705045in}{1.168222in}}%
\pgfusepath{clip}%
\pgfsetbuttcap%
\pgfsetroundjoin%
\definecolor{currentfill}{rgb}{0.121569,0.466667,0.705882}%
\pgfsetfillcolor{currentfill}%
\pgfsetlinewidth{0.803000pt}%
\definecolor{currentstroke}{rgb}{0.121569,0.466667,0.705882}%
\pgfsetstrokecolor{currentstroke}%
\pgfsetdash{}{0pt}%
\pgfsys@defobject{currentmarker}{\pgfqpoint{-0.019444in}{-0.019444in}}{\pgfqpoint{0.019444in}{0.019444in}}{%
\pgfpathmoveto{\pgfqpoint{-0.019444in}{-0.019444in}}%
\pgfpathlineto{\pgfqpoint{0.019444in}{0.019444in}}%
\pgfpathmoveto{\pgfqpoint{-0.019444in}{0.019444in}}%
\pgfpathlineto{\pgfqpoint{0.019444in}{-0.019444in}}%
\pgfusepath{stroke,fill}%
}%
\begin{pgfscope}%
\pgfsys@transformshift{0.550575in}{1.252512in}%
\pgfsys@useobject{currentmarker}{}%
\end{pgfscope}%
\begin{pgfscope}%
\pgfsys@transformshift{0.641754in}{1.256267in}%
\pgfsys@useobject{currentmarker}{}%
\end{pgfscope}%
\begin{pgfscope}%
\pgfsys@transformshift{0.732933in}{1.259994in}%
\pgfsys@useobject{currentmarker}{}%
\end{pgfscope}%
\begin{pgfscope}%
\pgfsys@transformshift{0.824112in}{1.263699in}%
\pgfsys@useobject{currentmarker}{}%
\end{pgfscope}%
\begin{pgfscope}%
\pgfsys@transformshift{0.915291in}{1.267366in}%
\pgfsys@useobject{currentmarker}{}%
\end{pgfscope}%
\begin{pgfscope}%
\pgfsys@transformshift{1.006470in}{1.270999in}%
\pgfsys@useobject{currentmarker}{}%
\end{pgfscope}%
\begin{pgfscope}%
\pgfsys@transformshift{1.097649in}{1.274597in}%
\pgfsys@useobject{currentmarker}{}%
\end{pgfscope}%
\begin{pgfscope}%
\pgfsys@transformshift{1.188827in}{1.278163in}%
\pgfsys@useobject{currentmarker}{}%
\end{pgfscope}%
\begin{pgfscope}%
\pgfsys@transformshift{1.280006in}{1.281682in}%
\pgfsys@useobject{currentmarker}{}%
\end{pgfscope}%
\begin{pgfscope}%
\pgfsys@transformshift{1.371185in}{1.285159in}%
\pgfsys@useobject{currentmarker}{}%
\end{pgfscope}%
\begin{pgfscope}%
\pgfsys@transformshift{1.462364in}{1.288591in}%
\pgfsys@useobject{currentmarker}{}%
\end{pgfscope}%
\begin{pgfscope}%
\pgfsys@transformshift{1.553543in}{1.291981in}%
\pgfsys@useobject{currentmarker}{}%
\end{pgfscope}%
\begin{pgfscope}%
\pgfsys@transformshift{1.644722in}{1.295316in}%
\pgfsys@useobject{currentmarker}{}%
\end{pgfscope}%
\begin{pgfscope}%
\pgfsys@transformshift{1.735901in}{1.298598in}%
\pgfsys@useobject{currentmarker}{}%
\end{pgfscope}%
\begin{pgfscope}%
\pgfsys@transformshift{1.827080in}{1.301826in}%
\pgfsys@useobject{currentmarker}{}%
\end{pgfscope}%
\begin{pgfscope}%
\pgfsys@transformshift{1.918259in}{1.305002in}%
\pgfsys@useobject{currentmarker}{}%
\end{pgfscope}%
\begin{pgfscope}%
\pgfsys@transformshift{2.009437in}{1.308113in}%
\pgfsys@useobject{currentmarker}{}%
\end{pgfscope}%
\begin{pgfscope}%
\pgfsys@transformshift{2.100616in}{1.311162in}%
\pgfsys@useobject{currentmarker}{}%
\end{pgfscope}%
\end{pgfscope}%
\begin{pgfscope}%
\pgfpathrectangle{\pgfqpoint{0.473073in}{0.462654in}}{\pgfqpoint{1.705045in}{1.168222in}}%
\pgfusepath{clip}%
\pgfsetrectcap%
\pgfsetroundjoin%
\pgfsetlinewidth{0.903375pt}%
\definecolor{currentstroke}{rgb}{1.000000,0.000000,0.000000}%
\pgfsetstrokecolor{currentstroke}%
\pgfsetdash{}{0pt}%
\pgfpathmoveto{\pgfqpoint{0.550575in}{1.076090in}}%
\pgfpathlineto{\pgfqpoint{0.815406in}{1.109889in}}%
\pgfpathlineto{\pgfqpoint{1.088027in}{1.146998in}}%
\pgfpathlineto{\pgfqpoint{1.360647in}{1.186411in}}%
\pgfpathlineto{\pgfqpoint{1.641056in}{1.229161in}}%
\pgfpathlineto{\pgfqpoint{1.944833in}{1.277695in}}%
\pgfpathlineto{\pgfqpoint{2.100616in}{1.303318in}}%
\pgfpathlineto{\pgfqpoint{2.100616in}{1.303318in}}%
\pgfusepath{stroke}%
\end{pgfscope}%
\begin{pgfscope}%
\pgfpathrectangle{\pgfqpoint{0.473073in}{0.462654in}}{\pgfqpoint{1.705045in}{1.168222in}}%
\pgfusepath{clip}%
\pgfsetbuttcap%
\pgfsetroundjoin%
\definecolor{currentfill}{rgb}{0.000000,0.000000,0.000000}%
\pgfsetfillcolor{currentfill}%
\pgfsetfillopacity{0.000000}%
\pgfsetlinewidth{0.803000pt}%
\definecolor{currentstroke}{rgb}{1.000000,0.000000,0.000000}%
\pgfsetstrokecolor{currentstroke}%
\pgfsetdash{}{0pt}%
\pgfsys@defobject{currentmarker}{\pgfqpoint{-0.017500in}{-0.017500in}}{\pgfqpoint{0.017500in}{0.017500in}}{%
\pgfpathmoveto{\pgfqpoint{0.000000in}{-0.017500in}}%
\pgfpathcurveto{\pgfqpoint{0.004641in}{-0.017500in}}{\pgfqpoint{0.009093in}{-0.015656in}}{\pgfqpoint{0.012374in}{-0.012374in}}%
\pgfpathcurveto{\pgfqpoint{0.015656in}{-0.009093in}}{\pgfqpoint{0.017500in}{-0.004641in}}{\pgfqpoint{0.017500in}{0.000000in}}%
\pgfpathcurveto{\pgfqpoint{0.017500in}{0.004641in}}{\pgfqpoint{0.015656in}{0.009093in}}{\pgfqpoint{0.012374in}{0.012374in}}%
\pgfpathcurveto{\pgfqpoint{0.009093in}{0.015656in}}{\pgfqpoint{0.004641in}{0.017500in}}{\pgfqpoint{0.000000in}{0.017500in}}%
\pgfpathcurveto{\pgfqpoint{-0.004641in}{0.017500in}}{\pgfqpoint{-0.009093in}{0.015656in}}{\pgfqpoint{-0.012374in}{0.012374in}}%
\pgfpathcurveto{\pgfqpoint{-0.015656in}{0.009093in}}{\pgfqpoint{-0.017500in}{0.004641in}}{\pgfqpoint{-0.017500in}{0.000000in}}%
\pgfpathcurveto{\pgfqpoint{-0.017500in}{-0.004641in}}{\pgfqpoint{-0.015656in}{-0.009093in}}{\pgfqpoint{-0.012374in}{-0.012374in}}%
\pgfpathcurveto{\pgfqpoint{-0.009093in}{-0.015656in}}{\pgfqpoint{-0.004641in}{-0.017500in}}{\pgfqpoint{0.000000in}{-0.017500in}}%
\pgfpathlineto{\pgfqpoint{0.000000in}{-0.017500in}}%
\pgfpathclose%
\pgfusepath{stroke,fill}%
}%
\begin{pgfscope}%
\pgfsys@transformshift{0.550575in}{1.076090in}%
\pgfsys@useobject{currentmarker}{}%
\end{pgfscope}%
\begin{pgfscope}%
\pgfsys@transformshift{0.641754in}{1.087465in}%
\pgfsys@useobject{currentmarker}{}%
\end{pgfscope}%
\begin{pgfscope}%
\pgfsys@transformshift{0.732933in}{1.099108in}%
\pgfsys@useobject{currentmarker}{}%
\end{pgfscope}%
\begin{pgfscope}%
\pgfsys@transformshift{0.824112in}{1.111038in}%
\pgfsys@useobject{currentmarker}{}%
\end{pgfscope}%
\begin{pgfscope}%
\pgfsys@transformshift{0.915291in}{1.123214in}%
\pgfsys@useobject{currentmarker}{}%
\end{pgfscope}%
\begin{pgfscope}%
\pgfsys@transformshift{1.006470in}{1.135653in}%
\pgfsys@useobject{currentmarker}{}%
\end{pgfscope}%
\begin{pgfscope}%
\pgfsys@transformshift{1.097649in}{1.148351in}%
\pgfsys@useobject{currentmarker}{}%
\end{pgfscope}%
\begin{pgfscope}%
\pgfsys@transformshift{1.188827in}{1.161326in}%
\pgfsys@useobject{currentmarker}{}%
\end{pgfscope}%
\begin{pgfscope}%
\pgfsys@transformshift{1.280006in}{1.174530in}%
\pgfsys@useobject{currentmarker}{}%
\end{pgfscope}%
\begin{pgfscope}%
\pgfsys@transformshift{1.371185in}{1.187978in}%
\pgfsys@useobject{currentmarker}{}%
\end{pgfscope}%
\begin{pgfscope}%
\pgfsys@transformshift{1.462364in}{1.201662in}%
\pgfsys@useobject{currentmarker}{}%
\end{pgfscope}%
\begin{pgfscope}%
\pgfsys@transformshift{1.553543in}{1.215600in}%
\pgfsys@useobject{currentmarker}{}%
\end{pgfscope}%
\begin{pgfscope}%
\pgfsys@transformshift{1.644722in}{1.229734in}%
\pgfsys@useobject{currentmarker}{}%
\end{pgfscope}%
\begin{pgfscope}%
\pgfsys@transformshift{1.735901in}{1.244077in}%
\pgfsys@useobject{currentmarker}{}%
\end{pgfscope}%
\begin{pgfscope}%
\pgfsys@transformshift{1.827080in}{1.258618in}%
\pgfsys@useobject{currentmarker}{}%
\end{pgfscope}%
\begin{pgfscope}%
\pgfsys@transformshift{1.918259in}{1.273368in}%
\pgfsys@useobject{currentmarker}{}%
\end{pgfscope}%
\begin{pgfscope}%
\pgfsys@transformshift{2.009437in}{1.288265in}%
\pgfsys@useobject{currentmarker}{}%
\end{pgfscope}%
\begin{pgfscope}%
\pgfsys@transformshift{2.100616in}{1.303318in}%
\pgfsys@useobject{currentmarker}{}%
\end{pgfscope}%
\end{pgfscope}%
\begin{pgfscope}%
\pgfpathrectangle{\pgfqpoint{0.473073in}{0.462654in}}{\pgfqpoint{1.705045in}{1.168222in}}%
\pgfusepath{clip}%
\pgfsetrectcap%
\pgfsetroundjoin%
\pgfsetlinewidth{0.903375pt}%
\definecolor{currentstroke}{rgb}{0.172549,0.627451,0.172549}%
\pgfsetstrokecolor{currentstroke}%
\pgfsetdash{}{0pt}%
\pgfpathmoveto{\pgfqpoint{0.550575in}{0.782897in}}%
\pgfpathlineto{\pgfqpoint{0.675202in}{0.796306in}}%
\pgfpathlineto{\pgfqpoint{0.792039in}{0.811114in}}%
\pgfpathlineto{\pgfqpoint{0.901087in}{0.827173in}}%
\pgfpathlineto{\pgfqpoint{1.002346in}{0.844301in}}%
\pgfpathlineto{\pgfqpoint{1.095816in}{0.862263in}}%
\pgfpathlineto{\pgfqpoint{1.181496in}{0.880796in}}%
\pgfpathlineto{\pgfqpoint{1.267177in}{0.901459in}}%
\pgfpathlineto{\pgfqpoint{1.352858in}{0.924505in}}%
\pgfpathlineto{\pgfqpoint{1.430749in}{0.947729in}}%
\pgfpathlineto{\pgfqpoint{1.508641in}{0.973361in}}%
\pgfpathlineto{\pgfqpoint{1.586532in}{1.001584in}}%
\pgfpathlineto{\pgfqpoint{1.656635in}{1.029355in}}%
\pgfpathlineto{\pgfqpoint{1.726737in}{1.059592in}}%
\pgfpathlineto{\pgfqpoint{1.789050in}{1.088664in}}%
\pgfpathlineto{\pgfqpoint{1.851363in}{1.119975in}}%
\pgfpathlineto{\pgfqpoint{1.913677in}{1.153697in}}%
\pgfpathlineto{\pgfqpoint{1.975990in}{1.189857in}}%
\pgfpathlineto{\pgfqpoint{2.038303in}{1.228640in}}%
\pgfpathlineto{\pgfqpoint{2.100616in}{1.270221in}}%
\pgfpathlineto{\pgfqpoint{2.100616in}{1.270221in}}%
\pgfusepath{stroke}%
\end{pgfscope}%
\begin{pgfscope}%
\pgfpathrectangle{\pgfqpoint{0.473073in}{0.462654in}}{\pgfqpoint{1.705045in}{1.168222in}}%
\pgfusepath{clip}%
\pgfsetbuttcap%
\pgfsetmiterjoin%
\definecolor{currentfill}{rgb}{0.000000,0.000000,0.000000}%
\pgfsetfillcolor{currentfill}%
\pgfsetfillopacity{0.000000}%
\pgfsetlinewidth{0.803000pt}%
\definecolor{currentstroke}{rgb}{0.172549,0.627451,0.172549}%
\pgfsetstrokecolor{currentstroke}%
\pgfsetdash{}{0pt}%
\pgfsys@defobject{currentmarker}{\pgfqpoint{-0.017500in}{-0.017500in}}{\pgfqpoint{0.017500in}{0.017500in}}{%
\pgfpathmoveto{\pgfqpoint{-0.017500in}{-0.017500in}}%
\pgfpathlineto{\pgfqpoint{0.017500in}{-0.017500in}}%
\pgfpathlineto{\pgfqpoint{0.017500in}{0.017500in}}%
\pgfpathlineto{\pgfqpoint{-0.017500in}{0.017500in}}%
\pgfpathlineto{\pgfqpoint{-0.017500in}{-0.017500in}}%
\pgfpathclose%
\pgfusepath{stroke,fill}%
}%
\begin{pgfscope}%
\pgfsys@transformshift{0.550575in}{0.782897in}%
\pgfsys@useobject{currentmarker}{}%
\end{pgfscope}%
\begin{pgfscope}%
\pgfsys@transformshift{0.641754in}{0.792484in}%
\pgfsys@useobject{currentmarker}{}%
\end{pgfscope}%
\begin{pgfscope}%
\pgfsys@transformshift{0.732933in}{0.803326in}%
\pgfsys@useobject{currentmarker}{}%
\end{pgfscope}%
\begin{pgfscope}%
\pgfsys@transformshift{0.824112in}{0.815601in}%
\pgfsys@useobject{currentmarker}{}%
\end{pgfscope}%
\begin{pgfscope}%
\pgfsys@transformshift{0.915291in}{0.829438in}%
\pgfsys@useobject{currentmarker}{}%
\end{pgfscope}%
\begin{pgfscope}%
\pgfsys@transformshift{1.006470in}{0.845047in}%
\pgfsys@useobject{currentmarker}{}%
\end{pgfscope}%
\begin{pgfscope}%
\pgfsys@transformshift{1.097649in}{0.862638in}%
\pgfsys@useobject{currentmarker}{}%
\end{pgfscope}%
\begin{pgfscope}%
\pgfsys@transformshift{1.188827in}{0.882476in}%
\pgfsys@useobject{currentmarker}{}%
\end{pgfscope}%
\begin{pgfscope}%
\pgfsys@transformshift{1.280006in}{0.904748in}%
\pgfsys@useobject{currentmarker}{}%
\end{pgfscope}%
\begin{pgfscope}%
\pgfsys@transformshift{1.371185in}{0.929764in}%
\pgfsys@useobject{currentmarker}{}%
\end{pgfscope}%
\begin{pgfscope}%
\pgfsys@transformshift{1.462364in}{0.957823in}%
\pgfsys@useobject{currentmarker}{}%
\end{pgfscope}%
\begin{pgfscope}%
\pgfsys@transformshift{1.553543in}{0.989310in}%
\pgfsys@useobject{currentmarker}{}%
\end{pgfscope}%
\begin{pgfscope}%
\pgfsys@transformshift{1.644722in}{1.024470in}%
\pgfsys@useobject{currentmarker}{}%
\end{pgfscope}%
\begin{pgfscope}%
\pgfsys@transformshift{1.735901in}{1.063731in}%
\pgfsys@useobject{currentmarker}{}%
\end{pgfscope}%
\begin{pgfscope}%
\pgfsys@transformshift{1.827080in}{1.107494in}%
\pgfsys@useobject{currentmarker}{}%
\end{pgfscope}%
\begin{pgfscope}%
\pgfsys@transformshift{1.918259in}{1.156268in}%
\pgfsys@useobject{currentmarker}{}%
\end{pgfscope}%
\begin{pgfscope}%
\pgfsys@transformshift{2.009437in}{1.210332in}%
\pgfsys@useobject{currentmarker}{}%
\end{pgfscope}%
\begin{pgfscope}%
\pgfsys@transformshift{2.100616in}{1.270221in}%
\pgfsys@useobject{currentmarker}{}%
\end{pgfscope}%
\end{pgfscope}%
\begin{pgfscope}%
\pgfpathrectangle{\pgfqpoint{0.473073in}{0.462654in}}{\pgfqpoint{1.705045in}{1.168222in}}%
\pgfusepath{clip}%
\pgfsetrectcap%
\pgfsetroundjoin%
\pgfsetlinewidth{0.903375pt}%
\definecolor{currentstroke}{rgb}{0.000000,0.000000,0.000000}%
\pgfsetstrokecolor{currentstroke}%
\pgfsetdash{}{0pt}%
\pgfpathmoveto{\pgfqpoint{0.473073in}{1.390918in}}%
\pgfpathlineto{\pgfqpoint{2.178118in}{1.390918in}}%
\pgfusepath{stroke}%
\end{pgfscope}%
\begin{pgfscope}%
\pgfsetrectcap%
\pgfsetmiterjoin%
\pgfsetlinewidth{0.803000pt}%
\definecolor{currentstroke}{rgb}{0.000000,0.000000,0.000000}%
\pgfsetstrokecolor{currentstroke}%
\pgfsetdash{}{0pt}%
\pgfpathmoveto{\pgfqpoint{0.473073in}{0.462654in}}%
\pgfpathlineto{\pgfqpoint{0.473073in}{1.630877in}}%
\pgfusepath{stroke}%
\end{pgfscope}%
\begin{pgfscope}%
\pgfsetrectcap%
\pgfsetmiterjoin%
\pgfsetlinewidth{0.803000pt}%
\definecolor{currentstroke}{rgb}{0.000000,0.000000,0.000000}%
\pgfsetstrokecolor{currentstroke}%
\pgfsetdash{}{0pt}%
\pgfpathmoveto{\pgfqpoint{2.178118in}{0.462654in}}%
\pgfpathlineto{\pgfqpoint{2.178118in}{1.630877in}}%
\pgfusepath{stroke}%
\end{pgfscope}%
\begin{pgfscope}%
\pgfsetrectcap%
\pgfsetmiterjoin%
\pgfsetlinewidth{0.803000pt}%
\definecolor{currentstroke}{rgb}{0.000000,0.000000,0.000000}%
\pgfsetstrokecolor{currentstroke}%
\pgfsetdash{}{0pt}%
\pgfpathmoveto{\pgfqpoint{0.473073in}{0.462654in}}%
\pgfpathlineto{\pgfqpoint{2.178118in}{0.462654in}}%
\pgfusepath{stroke}%
\end{pgfscope}%
\begin{pgfscope}%
\pgfsetrectcap%
\pgfsetmiterjoin%
\pgfsetlinewidth{0.803000pt}%
\definecolor{currentstroke}{rgb}{0.000000,0.000000,0.000000}%
\pgfsetstrokecolor{currentstroke}%
\pgfsetdash{}{0pt}%
\pgfpathmoveto{\pgfqpoint{0.473073in}{1.630877in}}%
\pgfpathlineto{\pgfqpoint{2.178118in}{1.630877in}}%
\pgfusepath{stroke}%
\end{pgfscope}%
\begin{pgfscope}%
\definecolor{textcolor}{rgb}{0.000000,0.000000,0.000000}%
\pgfsetstrokecolor{textcolor}%
\pgfsetfillcolor{textcolor}%
\pgftext[x=0.524225in,y=1.595830in,left,top]{\color{textcolor}{\rmfamily\fontsize{9.000000}{10.800000}\selectfont\catcode`\^=\active\def^{\ifmmode\sp\else\^{}\fi}\catcode`\%=\active\def%{\%}$(c)$}}%
\end{pgfscope}%
\begin{pgfscope}%
\pgfsetbuttcap%
\pgfsetmiterjoin%
\definecolor{currentfill}{rgb}{1.000000,1.000000,1.000000}%
\pgfsetfillcolor{currentfill}%
\pgfsetlinewidth{0.000000pt}%
\definecolor{currentstroke}{rgb}{0.000000,0.000000,0.000000}%
\pgfsetstrokecolor{currentstroke}%
\pgfsetstrokeopacity{0.000000}%
\pgfsetdash{}{0pt}%
\pgfpathmoveto{\pgfqpoint{2.516573in}{0.462654in}}%
\pgfpathlineto{\pgfqpoint{4.221618in}{0.462654in}}%
\pgfpathlineto{\pgfqpoint{4.221618in}{1.630877in}}%
\pgfpathlineto{\pgfqpoint{2.516573in}{1.630877in}}%
\pgfpathlineto{\pgfqpoint{2.516573in}{0.462654in}}%
\pgfpathclose%
\pgfusepath{fill}%
\end{pgfscope}%
\begin{pgfscope}%
\pgfsetbuttcap%
\pgfsetroundjoin%
\definecolor{currentfill}{rgb}{0.000000,0.000000,0.000000}%
\pgfsetfillcolor{currentfill}%
\pgfsetlinewidth{0.803000pt}%
\definecolor{currentstroke}{rgb}{0.000000,0.000000,0.000000}%
\pgfsetstrokecolor{currentstroke}%
\pgfsetdash{}{0pt}%
\pgfsys@defobject{currentmarker}{\pgfqpoint{0.000000in}{-0.048611in}}{\pgfqpoint{0.000000in}{0.000000in}}{%
\pgfpathmoveto{\pgfqpoint{0.000000in}{0.000000in}}%
\pgfpathlineto{\pgfqpoint{0.000000in}{-0.048611in}}%
\pgfusepath{stroke,fill}%
}%
\begin{pgfscope}%
\pgfsys@transformshift{2.516573in}{0.462654in}%
\pgfsys@useobject{currentmarker}{}%
\end{pgfscope}%
\end{pgfscope}%
\begin{pgfscope}%
\definecolor{textcolor}{rgb}{0.000000,0.000000,0.000000}%
\pgfsetstrokecolor{textcolor}%
\pgfsetfillcolor{textcolor}%
\pgftext[x=2.516573in,y=0.365432in,,top]{\color{textcolor}{\rmfamily\fontsize{8.000000}{9.600000}\selectfont\catcode`\^=\active\def^{\ifmmode\sp\else\^{}\fi}\catcode`\%=\active\def%{\%}$0$}}%
\end{pgfscope}%
\begin{pgfscope}%
\pgfsetbuttcap%
\pgfsetroundjoin%
\definecolor{currentfill}{rgb}{0.000000,0.000000,0.000000}%
\pgfsetfillcolor{currentfill}%
\pgfsetlinewidth{0.803000pt}%
\definecolor{currentstroke}{rgb}{0.000000,0.000000,0.000000}%
\pgfsetstrokecolor{currentstroke}%
\pgfsetdash{}{0pt}%
\pgfsys@defobject{currentmarker}{\pgfqpoint{0.000000in}{-0.048611in}}{\pgfqpoint{0.000000in}{0.000000in}}{%
\pgfpathmoveto{\pgfqpoint{0.000000in}{0.000000in}}%
\pgfpathlineto{\pgfqpoint{0.000000in}{-0.048611in}}%
\pgfusepath{stroke,fill}%
}%
\begin{pgfscope}%
\pgfsys@transformshift{3.328499in}{0.462654in}%
\pgfsys@useobject{currentmarker}{}%
\end{pgfscope}%
\end{pgfscope}%
\begin{pgfscope}%
\definecolor{textcolor}{rgb}{0.000000,0.000000,0.000000}%
\pgfsetstrokecolor{textcolor}%
\pgfsetfillcolor{textcolor}%
\pgftext[x=3.328499in,y=0.365432in,,top]{\color{textcolor}{\rmfamily\fontsize{8.000000}{9.600000}\selectfont\catcode`\^=\active\def^{\ifmmode\sp\else\^{}\fi}\catcode`\%=\active\def%{\%}$0.5$}}%
\end{pgfscope}%
\begin{pgfscope}%
\pgfsetbuttcap%
\pgfsetroundjoin%
\definecolor{currentfill}{rgb}{0.000000,0.000000,0.000000}%
\pgfsetfillcolor{currentfill}%
\pgfsetlinewidth{0.803000pt}%
\definecolor{currentstroke}{rgb}{0.000000,0.000000,0.000000}%
\pgfsetstrokecolor{currentstroke}%
\pgfsetdash{}{0pt}%
\pgfsys@defobject{currentmarker}{\pgfqpoint{0.000000in}{-0.048611in}}{\pgfqpoint{0.000000in}{0.000000in}}{%
\pgfpathmoveto{\pgfqpoint{0.000000in}{0.000000in}}%
\pgfpathlineto{\pgfqpoint{0.000000in}{-0.048611in}}%
\pgfusepath{stroke,fill}%
}%
\begin{pgfscope}%
\pgfsys@transformshift{4.140426in}{0.462654in}%
\pgfsys@useobject{currentmarker}{}%
\end{pgfscope}%
\end{pgfscope}%
\begin{pgfscope}%
\definecolor{textcolor}{rgb}{0.000000,0.000000,0.000000}%
\pgfsetstrokecolor{textcolor}%
\pgfsetfillcolor{textcolor}%
\pgftext[x=4.140426in,y=0.365432in,,top]{\color{textcolor}{\rmfamily\fontsize{8.000000}{9.600000}\selectfont\catcode`\^=\active\def^{\ifmmode\sp\else\^{}\fi}\catcode`\%=\active\def%{\%}$1$}}%
\end{pgfscope}%
\begin{pgfscope}%
\definecolor{textcolor}{rgb}{0.000000,0.000000,0.000000}%
\pgfsetstrokecolor{textcolor}%
\pgfsetfillcolor{textcolor}%
\pgftext[x=3.369096in,y=0.211111in,,top]{\color{textcolor}{\rmfamily\fontsize{9.000000}{10.800000}\selectfont\catcode`\^=\active\def^{\ifmmode\sp\else\^{}\fi}\catcode`\%=\active\def%{\%}$S$}}%
\end{pgfscope}%
\begin{pgfscope}%
\pgfsetbuttcap%
\pgfsetroundjoin%
\definecolor{currentfill}{rgb}{0.000000,0.000000,0.000000}%
\pgfsetfillcolor{currentfill}%
\pgfsetlinewidth{0.803000pt}%
\definecolor{currentstroke}{rgb}{0.000000,0.000000,0.000000}%
\pgfsetstrokecolor{currentstroke}%
\pgfsetdash{}{0pt}%
\pgfsys@defobject{currentmarker}{\pgfqpoint{-0.048611in}{0.000000in}}{\pgfqpoint{-0.000000in}{0.000000in}}{%
\pgfpathmoveto{\pgfqpoint{-0.000000in}{0.000000in}}%
\pgfpathlineto{\pgfqpoint{-0.048611in}{0.000000in}}%
\pgfusepath{stroke,fill}%
}%
\begin{pgfscope}%
\pgfsys@transformshift{2.516573in}{0.462654in}%
\pgfsys@useobject{currentmarker}{}%
\end{pgfscope}%
\end{pgfscope}%
\begin{pgfscope}%
\definecolor{textcolor}{rgb}{0.000000,0.000000,0.000000}%
\pgfsetstrokecolor{textcolor}%
\pgfsetfillcolor{textcolor}%
\pgftext[x=2.360322in, y=0.424074in, left, base]{\color{textcolor}{\rmfamily\fontsize{8.000000}{9.600000}\selectfont\catcode`\^=\active\def^{\ifmmode\sp\else\^{}\fi}\catcode`\%=\active\def%{\%}$0$}}%
\end{pgfscope}%
\begin{pgfscope}%
\pgfsetbuttcap%
\pgfsetroundjoin%
\definecolor{currentfill}{rgb}{0.000000,0.000000,0.000000}%
\pgfsetfillcolor{currentfill}%
\pgfsetlinewidth{0.803000pt}%
\definecolor{currentstroke}{rgb}{0.000000,0.000000,0.000000}%
\pgfsetstrokecolor{currentstroke}%
\pgfsetdash{}{0pt}%
\pgfsys@defobject{currentmarker}{\pgfqpoint{-0.048611in}{0.000000in}}{\pgfqpoint{-0.000000in}{0.000000in}}{%
\pgfpathmoveto{\pgfqpoint{-0.000000in}{0.000000in}}%
\pgfpathlineto{\pgfqpoint{-0.048611in}{0.000000in}}%
\pgfusepath{stroke,fill}%
}%
\begin{pgfscope}%
\pgfsys@transformshift{2.516573in}{0.949414in}%
\pgfsys@useobject{currentmarker}{}%
\end{pgfscope}%
\end{pgfscope}%
\begin{pgfscope}%
\definecolor{textcolor}{rgb}{0.000000,0.000000,0.000000}%
\pgfsetstrokecolor{textcolor}%
\pgfsetfillcolor{textcolor}%
\pgftext[x=2.268500in, y=0.910833in, left, base]{\color{textcolor}{\rmfamily\fontsize{8.000000}{9.600000}\selectfont\catcode`\^=\active\def^{\ifmmode\sp\else\^{}\fi}\catcode`\%=\active\def%{\%}$0.5$}}%
\end{pgfscope}%
\begin{pgfscope}%
\pgfsetbuttcap%
\pgfsetroundjoin%
\definecolor{currentfill}{rgb}{0.000000,0.000000,0.000000}%
\pgfsetfillcolor{currentfill}%
\pgfsetlinewidth{0.803000pt}%
\definecolor{currentstroke}{rgb}{0.000000,0.000000,0.000000}%
\pgfsetstrokecolor{currentstroke}%
\pgfsetdash{}{0pt}%
\pgfsys@defobject{currentmarker}{\pgfqpoint{-0.048611in}{0.000000in}}{\pgfqpoint{-0.000000in}{0.000000in}}{%
\pgfpathmoveto{\pgfqpoint{-0.000000in}{0.000000in}}%
\pgfpathlineto{\pgfqpoint{-0.048611in}{0.000000in}}%
\pgfusepath{stroke,fill}%
}%
\begin{pgfscope}%
\pgfsys@transformshift{2.516573in}{1.436173in}%
\pgfsys@useobject{currentmarker}{}%
\end{pgfscope}%
\end{pgfscope}%
\begin{pgfscope}%
\definecolor{textcolor}{rgb}{0.000000,0.000000,0.000000}%
\pgfsetstrokecolor{textcolor}%
\pgfsetfillcolor{textcolor}%
\pgftext[x=2.360322in, y=1.397593in, left, base]{\color{textcolor}{\rmfamily\fontsize{8.000000}{9.600000}\selectfont\catcode`\^=\active\def^{\ifmmode\sp\else\^{}\fi}\catcode`\%=\active\def%{\%}$1$}}%
\end{pgfscope}%
\begin{pgfscope}%
\pgfpathrectangle{\pgfqpoint{2.516573in}{0.462654in}}{\pgfqpoint{1.705045in}{1.168222in}}%
\pgfusepath{clip}%
\pgfsetrectcap%
\pgfsetroundjoin%
\pgfsetlinewidth{0.903375pt}%
\definecolor{currentstroke}{rgb}{0.121569,0.466667,0.705882}%
\pgfsetstrokecolor{currentstroke}%
\pgfsetdash{}{0pt}%
\pgfpathmoveto{\pgfqpoint{2.760151in}{1.212149in}}%
\pgfpathlineto{\pgfqpoint{2.794831in}{1.213358in}}%
\pgfpathlineto{\pgfqpoint{2.829512in}{1.216671in}}%
\pgfpathlineto{\pgfqpoint{2.843384in}{1.218558in}}%
\pgfpathlineto{\pgfqpoint{2.850320in}{1.220906in}}%
\pgfpathlineto{\pgfqpoint{2.864192in}{1.229920in}}%
\pgfpathlineto{\pgfqpoint{2.878064in}{1.242855in}}%
\pgfpathlineto{\pgfqpoint{2.912744in}{1.277891in}}%
\pgfpathlineto{\pgfqpoint{2.926616in}{1.287484in}}%
\pgfpathlineto{\pgfqpoint{2.968233in}{1.313121in}}%
\pgfpathlineto{\pgfqpoint{2.982105in}{1.319317in}}%
\pgfpathlineto{\pgfqpoint{2.995977in}{1.322934in}}%
\pgfpathlineto{\pgfqpoint{3.009849in}{1.323506in}}%
\pgfpathlineto{\pgfqpoint{3.079210in}{1.321996in}}%
\pgfpathlineto{\pgfqpoint{3.093082in}{1.321025in}}%
\pgfpathlineto{\pgfqpoint{3.106954in}{1.318145in}}%
\pgfpathlineto{\pgfqpoint{3.155506in}{1.304988in}}%
\pgfpathlineto{\pgfqpoint{3.169378in}{1.304326in}}%
\pgfpathlineto{\pgfqpoint{3.183250in}{1.307012in}}%
\pgfpathlineto{\pgfqpoint{3.197122in}{1.312384in}}%
\pgfpathlineto{\pgfqpoint{3.245675in}{1.335142in}}%
\pgfpathlineto{\pgfqpoint{3.301163in}{1.350718in}}%
\pgfpathlineto{\pgfqpoint{3.321971in}{1.353579in}}%
\pgfpathlineto{\pgfqpoint{3.356652in}{1.353577in}}%
\pgfpathlineto{\pgfqpoint{3.412140in}{1.352173in}}%
\pgfpathlineto{\pgfqpoint{3.426012in}{1.350325in}}%
\pgfpathlineto{\pgfqpoint{3.453756in}{1.343291in}}%
\pgfpathlineto{\pgfqpoint{3.474565in}{1.338406in}}%
\pgfpathlineto{\pgfqpoint{3.488437in}{1.336914in}}%
\pgfpathlineto{\pgfqpoint{3.502309in}{1.338063in}}%
\pgfpathlineto{\pgfqpoint{3.516181in}{1.341988in}}%
\pgfpathlineto{\pgfqpoint{3.564733in}{1.359283in}}%
\pgfpathlineto{\pgfqpoint{3.585541in}{1.360659in}}%
\pgfpathlineto{\pgfqpoint{3.647966in}{1.362066in}}%
\pgfpathlineto{\pgfqpoint{3.724263in}{1.361172in}}%
\pgfpathlineto{\pgfqpoint{3.752007in}{1.359775in}}%
\pgfpathlineto{\pgfqpoint{3.814431in}{1.353481in}}%
\pgfpathlineto{\pgfqpoint{3.828303in}{1.354551in}}%
\pgfpathlineto{\pgfqpoint{3.849111in}{1.359561in}}%
\pgfpathlineto{\pgfqpoint{3.876856in}{1.367442in}}%
\pgfpathlineto{\pgfqpoint{3.890728in}{1.369680in}}%
\pgfpathlineto{\pgfqpoint{3.918472in}{1.369690in}}%
\pgfpathlineto{\pgfqpoint{3.994769in}{1.368885in}}%
\pgfpathlineto{\pgfqpoint{4.078001in}{1.372160in}}%
\pgfpathlineto{\pgfqpoint{4.140426in}{1.371493in}}%
\pgfpathlineto{\pgfqpoint{4.140426in}{1.371493in}}%
\pgfusepath{stroke}%
\end{pgfscope}%
\begin{pgfscope}%
\pgfpathrectangle{\pgfqpoint{2.516573in}{0.462654in}}{\pgfqpoint{1.705045in}{1.168222in}}%
\pgfusepath{clip}%
\pgfsetbuttcap%
\pgfsetmiterjoin%
\definecolor{currentfill}{rgb}{0.000000,0.000000,0.000000}%
\pgfsetfillcolor{currentfill}%
\pgfsetfillopacity{0.000000}%
\pgfsetlinewidth{0.803000pt}%
\definecolor{currentstroke}{rgb}{0.121569,0.466667,0.705882}%
\pgfsetstrokecolor{currentstroke}%
\pgfsetdash{}{0pt}%
\pgfsys@defobject{currentmarker}{\pgfqpoint{-0.017500in}{-0.017500in}}{\pgfqpoint{0.017500in}{0.017500in}}{%
\pgfpathmoveto{\pgfqpoint{-0.017500in}{-0.017500in}}%
\pgfpathlineto{\pgfqpoint{0.017500in}{-0.017500in}}%
\pgfpathlineto{\pgfqpoint{0.017500in}{0.017500in}}%
\pgfpathlineto{\pgfqpoint{-0.017500in}{0.017500in}}%
\pgfpathlineto{\pgfqpoint{-0.017500in}{-0.017500in}}%
\pgfpathclose%
\pgfusepath{stroke,fill}%
}%
\begin{pgfscope}%
\pgfsys@transformshift{2.760151in}{1.212149in}%
\pgfsys@useobject{currentmarker}{}%
\end{pgfscope}%
\begin{pgfscope}%
\pgfsys@transformshift{2.841344in}{1.218198in}%
\pgfsys@useobject{currentmarker}{}%
\end{pgfscope}%
\begin{pgfscope}%
\pgfsys@transformshift{2.922536in}{1.284987in}%
\pgfsys@useobject{currentmarker}{}%
\end{pgfscope}%
\begin{pgfscope}%
\pgfsys@transformshift{3.003729in}{1.323518in}%
\pgfsys@useobject{currentmarker}{}%
\end{pgfscope}%
\begin{pgfscope}%
\pgfsys@transformshift{3.084922in}{1.321779in}%
\pgfsys@useobject{currentmarker}{}%
\end{pgfscope}%
\begin{pgfscope}%
\pgfsys@transformshift{3.166114in}{1.304214in}%
\pgfsys@useobject{currentmarker}{}%
\end{pgfscope}%
\begin{pgfscope}%
\pgfsys@transformshift{3.247307in}{1.335624in}%
\pgfsys@useobject{currentmarker}{}%
\end{pgfscope}%
\begin{pgfscope}%
\pgfsys@transformshift{3.328499in}{1.353776in}%
\pgfsys@useobject{currentmarker}{}%
\end{pgfscope}%
\begin{pgfscope}%
\pgfsys@transformshift{3.409692in}{1.352291in}%
\pgfsys@useobject{currentmarker}{}%
\end{pgfscope}%
\begin{pgfscope}%
\pgfsys@transformshift{3.490885in}{1.336876in}%
\pgfsys@useobject{currentmarker}{}%
\end{pgfscope}%
\begin{pgfscope}%
\pgfsys@transformshift{3.572077in}{1.360092in}%
\pgfsys@useobject{currentmarker}{}%
\end{pgfscope}%
\begin{pgfscope}%
\pgfsys@transformshift{3.653270in}{1.362076in}%
\pgfsys@useobject{currentmarker}{}%
\end{pgfscope}%
\begin{pgfscope}%
\pgfsys@transformshift{3.734463in}{1.360932in}%
\pgfsys@useobject{currentmarker}{}%
\end{pgfscope}%
\begin{pgfscope}%
\pgfsys@transformshift{3.815655in}{1.353476in}%
\pgfsys@useobject{currentmarker}{}%
\end{pgfscope}%
\begin{pgfscope}%
\pgfsys@transformshift{3.896848in}{1.369946in}%
\pgfsys@useobject{currentmarker}{}%
\end{pgfscope}%
\begin{pgfscope}%
\pgfsys@transformshift{3.978040in}{1.368480in}%
\pgfsys@useobject{currentmarker}{}%
\end{pgfscope}%
\begin{pgfscope}%
\pgfsys@transformshift{4.059233in}{1.372168in}%
\pgfsys@useobject{currentmarker}{}%
\end{pgfscope}%
\begin{pgfscope}%
\pgfsys@transformshift{4.140426in}{1.371493in}%
\pgfsys@useobject{currentmarker}{}%
\end{pgfscope}%
\end{pgfscope}%
\begin{pgfscope}%
\pgfpathrectangle{\pgfqpoint{2.516573in}{0.462654in}}{\pgfqpoint{1.705045in}{1.168222in}}%
\pgfusepath{clip}%
\pgfsetrectcap%
\pgfsetroundjoin%
\pgfsetlinewidth{0.903375pt}%
\definecolor{currentstroke}{rgb}{1.000000,0.000000,0.000000}%
\pgfsetstrokecolor{currentstroke}%
\pgfsetdash{}{0pt}%
\pgfpathmoveto{\pgfqpoint{2.760151in}{1.011753in}}%
\pgfpathlineto{\pgfqpoint{2.780959in}{1.012931in}}%
\pgfpathlineto{\pgfqpoint{2.801767in}{1.016245in}}%
\pgfpathlineto{\pgfqpoint{2.829512in}{1.023428in}}%
\pgfpathlineto{\pgfqpoint{2.843384in}{1.028183in}}%
\pgfpathlineto{\pgfqpoint{2.850320in}{1.034041in}}%
\pgfpathlineto{\pgfqpoint{2.857256in}{1.043589in}}%
\pgfpathlineto{\pgfqpoint{2.871128in}{1.070478in}}%
\pgfpathlineto{\pgfqpoint{2.898872in}{1.132482in}}%
\pgfpathlineto{\pgfqpoint{2.905808in}{1.144917in}}%
\pgfpathlineto{\pgfqpoint{2.912744in}{1.154486in}}%
\pgfpathlineto{\pgfqpoint{2.919680in}{1.160371in}}%
\pgfpathlineto{\pgfqpoint{2.933552in}{1.164267in}}%
\pgfpathlineto{\pgfqpoint{2.961297in}{1.169493in}}%
\pgfpathlineto{\pgfqpoint{2.989041in}{1.172506in}}%
\pgfpathlineto{\pgfqpoint{3.023721in}{1.173764in}}%
\pgfpathlineto{\pgfqpoint{3.086146in}{1.174954in}}%
\pgfpathlineto{\pgfqpoint{3.093082in}{1.177457in}}%
\pgfpathlineto{\pgfqpoint{3.100018in}{1.182828in}}%
\pgfpathlineto{\pgfqpoint{3.106954in}{1.190495in}}%
\pgfpathlineto{\pgfqpoint{3.120826in}{1.210440in}}%
\pgfpathlineto{\pgfqpoint{3.141634in}{1.243336in}}%
\pgfpathlineto{\pgfqpoint{3.155506in}{1.260602in}}%
\pgfpathlineto{\pgfqpoint{3.162442in}{1.266125in}}%
\pgfpathlineto{\pgfqpoint{3.176314in}{1.271587in}}%
\pgfpathlineto{\pgfqpoint{3.204058in}{1.279481in}}%
\pgfpathlineto{\pgfqpoint{3.231803in}{1.284590in}}%
\pgfpathlineto{\pgfqpoint{3.266483in}{1.287896in}}%
\pgfpathlineto{\pgfqpoint{3.315035in}{1.290500in}}%
\pgfpathlineto{\pgfqpoint{3.377460in}{1.290585in}}%
\pgfpathlineto{\pgfqpoint{3.439884in}{1.289319in}}%
\pgfpathlineto{\pgfqpoint{3.488437in}{1.286971in}}%
\pgfpathlineto{\pgfqpoint{3.502309in}{1.287678in}}%
\pgfpathlineto{\pgfqpoint{3.523117in}{1.291584in}}%
\pgfpathlineto{\pgfqpoint{3.557797in}{1.299167in}}%
\pgfpathlineto{\pgfqpoint{3.571669in}{1.300256in}}%
\pgfpathlineto{\pgfqpoint{3.592477in}{1.298993in}}%
\pgfpathlineto{\pgfqpoint{3.654902in}{1.292240in}}%
\pgfpathlineto{\pgfqpoint{3.745071in}{1.293627in}}%
\pgfpathlineto{\pgfqpoint{3.828303in}{1.295336in}}%
\pgfpathlineto{\pgfqpoint{3.862984in}{1.299298in}}%
\pgfpathlineto{\pgfqpoint{3.897664in}{1.302596in}}%
\pgfpathlineto{\pgfqpoint{3.973960in}{1.304801in}}%
\pgfpathlineto{\pgfqpoint{4.078001in}{1.303920in}}%
\pgfpathlineto{\pgfqpoint{4.140426in}{1.301766in}}%
\pgfpathlineto{\pgfqpoint{4.140426in}{1.301766in}}%
\pgfusepath{stroke}%
\end{pgfscope}%
\begin{pgfscope}%
\pgfpathrectangle{\pgfqpoint{2.516573in}{0.462654in}}{\pgfqpoint{1.705045in}{1.168222in}}%
\pgfusepath{clip}%
\pgfsetbuttcap%
\pgfsetroundjoin%
\definecolor{currentfill}{rgb}{0.000000,0.000000,0.000000}%
\pgfsetfillcolor{currentfill}%
\pgfsetfillopacity{0.000000}%
\pgfsetlinewidth{0.803000pt}%
\definecolor{currentstroke}{rgb}{1.000000,0.000000,0.000000}%
\pgfsetstrokecolor{currentstroke}%
\pgfsetdash{}{0pt}%
\pgfsys@defobject{currentmarker}{\pgfqpoint{-0.017500in}{-0.017500in}}{\pgfqpoint{0.017500in}{0.017500in}}{%
\pgfpathmoveto{\pgfqpoint{0.000000in}{-0.017500in}}%
\pgfpathcurveto{\pgfqpoint{0.004641in}{-0.017500in}}{\pgfqpoint{0.009093in}{-0.015656in}}{\pgfqpoint{0.012374in}{-0.012374in}}%
\pgfpathcurveto{\pgfqpoint{0.015656in}{-0.009093in}}{\pgfqpoint{0.017500in}{-0.004641in}}{\pgfqpoint{0.017500in}{0.000000in}}%
\pgfpathcurveto{\pgfqpoint{0.017500in}{0.004641in}}{\pgfqpoint{0.015656in}{0.009093in}}{\pgfqpoint{0.012374in}{0.012374in}}%
\pgfpathcurveto{\pgfqpoint{0.009093in}{0.015656in}}{\pgfqpoint{0.004641in}{0.017500in}}{\pgfqpoint{0.000000in}{0.017500in}}%
\pgfpathcurveto{\pgfqpoint{-0.004641in}{0.017500in}}{\pgfqpoint{-0.009093in}{0.015656in}}{\pgfqpoint{-0.012374in}{0.012374in}}%
\pgfpathcurveto{\pgfqpoint{-0.015656in}{0.009093in}}{\pgfqpoint{-0.017500in}{0.004641in}}{\pgfqpoint{-0.017500in}{0.000000in}}%
\pgfpathcurveto{\pgfqpoint{-0.017500in}{-0.004641in}}{\pgfqpoint{-0.015656in}{-0.009093in}}{\pgfqpoint{-0.012374in}{-0.012374in}}%
\pgfpathcurveto{\pgfqpoint{-0.009093in}{-0.015656in}}{\pgfqpoint{-0.004641in}{-0.017500in}}{\pgfqpoint{0.000000in}{-0.017500in}}%
\pgfpathlineto{\pgfqpoint{0.000000in}{-0.017500in}}%
\pgfpathclose%
\pgfusepath{stroke,fill}%
}%
\begin{pgfscope}%
\pgfsys@transformshift{2.760151in}{1.011753in}%
\pgfsys@useobject{currentmarker}{}%
\end{pgfscope}%
\begin{pgfscope}%
\pgfsys@transformshift{2.841344in}{1.027282in}%
\pgfsys@useobject{currentmarker}{}%
\end{pgfscope}%
\begin{pgfscope}%
\pgfsys@transformshift{2.922536in}{1.161532in}%
\pgfsys@useobject{currentmarker}{}%
\end{pgfscope}%
\begin{pgfscope}%
\pgfsys@transformshift{3.003729in}{1.173278in}%
\pgfsys@useobject{currentmarker}{}%
\end{pgfscope}%
\begin{pgfscope}%
\pgfsys@transformshift{3.084922in}{1.174853in}%
\pgfsys@useobject{currentmarker}{}%
\end{pgfscope}%
\begin{pgfscope}%
\pgfsys@transformshift{3.166114in}{1.267938in}%
\pgfsys@useobject{currentmarker}{}%
\end{pgfscope}%
\begin{pgfscope}%
\pgfsys@transformshift{3.247307in}{1.286316in}%
\pgfsys@useobject{currentmarker}{}%
\end{pgfscope}%
\begin{pgfscope}%
\pgfsys@transformshift{3.328499in}{1.290657in}%
\pgfsys@useobject{currentmarker}{}%
\end{pgfscope}%
\begin{pgfscope}%
\pgfsys@transformshift{3.409692in}{1.290465in}%
\pgfsys@useobject{currentmarker}{}%
\end{pgfscope}%
\begin{pgfscope}%
\pgfsys@transformshift{3.490885in}{1.286962in}%
\pgfsys@useobject{currentmarker}{}%
\end{pgfscope}%
\begin{pgfscope}%
\pgfsys@transformshift{3.572077in}{1.300257in}%
\pgfsys@useobject{currentmarker}{}%
\end{pgfscope}%
\begin{pgfscope}%
\pgfsys@transformshift{3.653270in}{1.292239in}%
\pgfsys@useobject{currentmarker}{}%
\end{pgfscope}%
\begin{pgfscope}%
\pgfsys@transformshift{3.734463in}{1.293482in}%
\pgfsys@useobject{currentmarker}{}%
\end{pgfscope}%
\begin{pgfscope}%
\pgfsys@transformshift{3.815655in}{1.294664in}%
\pgfsys@useobject{currentmarker}{}%
\end{pgfscope}%
\begin{pgfscope}%
\pgfsys@transformshift{3.896848in}{1.302561in}%
\pgfsys@useobject{currentmarker}{}%
\end{pgfscope}%
\begin{pgfscope}%
\pgfsys@transformshift{3.978040in}{1.304809in}%
\pgfsys@useobject{currentmarker}{}%
\end{pgfscope}%
\begin{pgfscope}%
\pgfsys@transformshift{4.059233in}{1.304241in}%
\pgfsys@useobject{currentmarker}{}%
\end{pgfscope}%
\begin{pgfscope}%
\pgfsys@transformshift{4.140426in}{1.301766in}%
\pgfsys@useobject{currentmarker}{}%
\end{pgfscope}%
\end{pgfscope}%
\begin{pgfscope}%
\pgfpathrectangle{\pgfqpoint{2.516573in}{0.462654in}}{\pgfqpoint{1.705045in}{1.168222in}}%
\pgfusepath{clip}%
\pgfsetrectcap%
\pgfsetroundjoin%
\pgfsetlinewidth{0.903375pt}%
\definecolor{currentstroke}{rgb}{0.172549,0.627451,0.172549}%
\pgfsetstrokecolor{currentstroke}%
\pgfsetdash{}{0pt}%
\pgfpathmoveto{\pgfqpoint{2.760151in}{0.973623in}}%
\pgfpathlineto{\pgfqpoint{2.780959in}{0.974639in}}%
\pgfpathlineto{\pgfqpoint{2.808703in}{0.978584in}}%
\pgfpathlineto{\pgfqpoint{2.843384in}{0.986148in}}%
\pgfpathlineto{\pgfqpoint{2.857256in}{0.991894in}}%
\pgfpathlineto{\pgfqpoint{2.871128in}{1.000442in}}%
\pgfpathlineto{\pgfqpoint{2.919680in}{1.033845in}}%
\pgfpathlineto{\pgfqpoint{2.975169in}{1.058598in}}%
\pgfpathlineto{\pgfqpoint{2.989041in}{1.062532in}}%
\pgfpathlineto{\pgfqpoint{3.002913in}{1.064497in}}%
\pgfpathlineto{\pgfqpoint{3.065337in}{1.066449in}}%
\pgfpathlineto{\pgfqpoint{3.093082in}{1.066216in}}%
\pgfpathlineto{\pgfqpoint{3.113890in}{1.062844in}}%
\pgfpathlineto{\pgfqpoint{3.155506in}{1.054347in}}%
\pgfpathlineto{\pgfqpoint{3.169378in}{1.053831in}}%
\pgfpathlineto{\pgfqpoint{3.183250in}{1.055822in}}%
\pgfpathlineto{\pgfqpoint{3.204058in}{1.062249in}}%
\pgfpathlineto{\pgfqpoint{3.238739in}{1.074690in}}%
\pgfpathlineto{\pgfqpoint{3.266483in}{1.080618in}}%
\pgfpathlineto{\pgfqpoint{3.301163in}{1.086677in}}%
\pgfpathlineto{\pgfqpoint{3.328907in}{1.089279in}}%
\pgfpathlineto{\pgfqpoint{3.398268in}{1.091200in}}%
\pgfpathlineto{\pgfqpoint{3.426012in}{1.090334in}}%
\pgfpathlineto{\pgfqpoint{3.460692in}{1.085255in}}%
\pgfpathlineto{\pgfqpoint{3.481501in}{1.082859in}}%
\pgfpathlineto{\pgfqpoint{3.502309in}{1.083044in}}%
\pgfpathlineto{\pgfqpoint{3.523117in}{1.085889in}}%
\pgfpathlineto{\pgfqpoint{3.571669in}{1.093618in}}%
\pgfpathlineto{\pgfqpoint{3.634094in}{1.096919in}}%
\pgfpathlineto{\pgfqpoint{3.689582in}{1.097186in}}%
\pgfpathlineto{\pgfqpoint{3.758943in}{1.096377in}}%
\pgfpathlineto{\pgfqpoint{3.821367in}{1.093634in}}%
\pgfpathlineto{\pgfqpoint{3.849111in}{1.096080in}}%
\pgfpathlineto{\pgfqpoint{3.890728in}{1.100311in}}%
\pgfpathlineto{\pgfqpoint{3.953152in}{1.100303in}}%
\pgfpathlineto{\pgfqpoint{3.994769in}{1.100864in}}%
\pgfpathlineto{\pgfqpoint{4.071065in}{1.105667in}}%
\pgfpathlineto{\pgfqpoint{4.126554in}{1.104389in}}%
\pgfpathlineto{\pgfqpoint{4.140426in}{1.103522in}}%
\pgfpathlineto{\pgfqpoint{4.140426in}{1.103522in}}%
\pgfusepath{stroke}%
\end{pgfscope}%
\begin{pgfscope}%
\pgfpathrectangle{\pgfqpoint{2.516573in}{0.462654in}}{\pgfqpoint{1.705045in}{1.168222in}}%
\pgfusepath{clip}%
\pgfsetbuttcap%
\pgfsetroundjoin%
\definecolor{currentfill}{rgb}{0.172549,0.627451,0.172549}%
\pgfsetfillcolor{currentfill}%
\pgfsetlinewidth{0.803000pt}%
\definecolor{currentstroke}{rgb}{0.172549,0.627451,0.172549}%
\pgfsetstrokecolor{currentstroke}%
\pgfsetdash{}{0pt}%
\pgfsys@defobject{currentmarker}{\pgfqpoint{-0.019444in}{-0.019444in}}{\pgfqpoint{0.019444in}{0.019444in}}{%
\pgfpathmoveto{\pgfqpoint{-0.019444in}{-0.019444in}}%
\pgfpathlineto{\pgfqpoint{0.019444in}{0.019444in}}%
\pgfpathmoveto{\pgfqpoint{-0.019444in}{0.019444in}}%
\pgfpathlineto{\pgfqpoint{0.019444in}{-0.019444in}}%
\pgfusepath{stroke,fill}%
}%
\begin{pgfscope}%
\pgfsys@transformshift{2.760151in}{0.973623in}%
\pgfsys@useobject{currentmarker}{}%
\end{pgfscope}%
\begin{pgfscope}%
\pgfsys@transformshift{2.841344in}{0.985617in}%
\pgfsys@useobject{currentmarker}{}%
\end{pgfscope}%
\begin{pgfscope}%
\pgfsys@transformshift{2.922536in}{1.035208in}%
\pgfsys@useobject{currentmarker}{}%
\end{pgfscope}%
\begin{pgfscope}%
\pgfsys@transformshift{3.003729in}{1.064540in}%
\pgfsys@useobject{currentmarker}{}%
\end{pgfscope}%
\begin{pgfscope}%
\pgfsys@transformshift{3.084922in}{1.066579in}%
\pgfsys@useobject{currentmarker}{}%
\end{pgfscope}%
\begin{pgfscope}%
\pgfsys@transformshift{3.166114in}{1.053747in}%
\pgfsys@useobject{currentmarker}{}%
\end{pgfscope}%
\begin{pgfscope}%
\pgfsys@transformshift{3.247307in}{1.076772in}%
\pgfsys@useobject{currentmarker}{}%
\end{pgfscope}%
\begin{pgfscope}%
\pgfsys@transformshift{3.328499in}{1.089261in}%
\pgfsys@useobject{currentmarker}{}%
\end{pgfscope}%
\begin{pgfscope}%
\pgfsys@transformshift{3.409692in}{1.091249in}%
\pgfsys@useobject{currentmarker}{}%
\end{pgfscope}%
\begin{pgfscope}%
\pgfsys@transformshift{3.490885in}{1.082537in}%
\pgfsys@useobject{currentmarker}{}%
\end{pgfscope}%
\begin{pgfscope}%
\pgfsys@transformshift{3.572077in}{1.093646in}%
\pgfsys@useobject{currentmarker}{}%
\end{pgfscope}%
\begin{pgfscope}%
\pgfsys@transformshift{3.653270in}{1.097189in}%
\pgfsys@useobject{currentmarker}{}%
\end{pgfscope}%
\begin{pgfscope}%
\pgfsys@transformshift{3.734463in}{1.097174in}%
\pgfsys@useobject{currentmarker}{}%
\end{pgfscope}%
\begin{pgfscope}%
\pgfsys@transformshift{3.815655in}{1.093536in}%
\pgfsys@useobject{currentmarker}{}%
\end{pgfscope}%
\begin{pgfscope}%
\pgfsys@transformshift{3.896848in}{1.100422in}%
\pgfsys@useobject{currentmarker}{}%
\end{pgfscope}%
\begin{pgfscope}%
\pgfsys@transformshift{3.978040in}{1.100269in}%
\pgfsys@useobject{currentmarker}{}%
\end{pgfscope}%
\begin{pgfscope}%
\pgfsys@transformshift{4.059233in}{1.105684in}%
\pgfsys@useobject{currentmarker}{}%
\end{pgfscope}%
\begin{pgfscope}%
\pgfsys@transformshift{4.140426in}{1.103522in}%
\pgfsys@useobject{currentmarker}{}%
\end{pgfscope}%
\end{pgfscope}%
\begin{pgfscope}%
\pgfpathrectangle{\pgfqpoint{2.516573in}{0.462654in}}{\pgfqpoint{1.705045in}{1.168222in}}%
\pgfusepath{clip}%
\pgfsetrectcap%
\pgfsetroundjoin%
\pgfsetlinewidth{0.903375pt}%
\definecolor{currentstroke}{rgb}{0.000000,0.000000,0.000000}%
\pgfsetstrokecolor{currentstroke}%
\pgfsetdash{}{0pt}%
\pgfpathmoveto{\pgfqpoint{2.516573in}{1.390918in}}%
\pgfpathlineto{\pgfqpoint{4.221618in}{1.390918in}}%
\pgfusepath{stroke}%
\end{pgfscope}%
\begin{pgfscope}%
\pgfsetrectcap%
\pgfsetmiterjoin%
\pgfsetlinewidth{0.803000pt}%
\definecolor{currentstroke}{rgb}{0.000000,0.000000,0.000000}%
\pgfsetstrokecolor{currentstroke}%
\pgfsetdash{}{0pt}%
\pgfpathmoveto{\pgfqpoint{2.516573in}{0.462654in}}%
\pgfpathlineto{\pgfqpoint{2.516573in}{1.630877in}}%
\pgfusepath{stroke}%
\end{pgfscope}%
\begin{pgfscope}%
\pgfsetrectcap%
\pgfsetmiterjoin%
\pgfsetlinewidth{0.803000pt}%
\definecolor{currentstroke}{rgb}{0.000000,0.000000,0.000000}%
\pgfsetstrokecolor{currentstroke}%
\pgfsetdash{}{0pt}%
\pgfpathmoveto{\pgfqpoint{4.221618in}{0.462654in}}%
\pgfpathlineto{\pgfqpoint{4.221618in}{1.630877in}}%
\pgfusepath{stroke}%
\end{pgfscope}%
\begin{pgfscope}%
\pgfsetrectcap%
\pgfsetmiterjoin%
\pgfsetlinewidth{0.803000pt}%
\definecolor{currentstroke}{rgb}{0.000000,0.000000,0.000000}%
\pgfsetstrokecolor{currentstroke}%
\pgfsetdash{}{0pt}%
\pgfpathmoveto{\pgfqpoint{2.516573in}{0.462654in}}%
\pgfpathlineto{\pgfqpoint{4.221618in}{0.462654in}}%
\pgfusepath{stroke}%
\end{pgfscope}%
\begin{pgfscope}%
\pgfsetrectcap%
\pgfsetmiterjoin%
\pgfsetlinewidth{0.803000pt}%
\definecolor{currentstroke}{rgb}{0.000000,0.000000,0.000000}%
\pgfsetstrokecolor{currentstroke}%
\pgfsetdash{}{0pt}%
\pgfpathmoveto{\pgfqpoint{2.516573in}{1.630877in}}%
\pgfpathlineto{\pgfqpoint{4.221618in}{1.630877in}}%
\pgfusepath{stroke}%
\end{pgfscope}%
\begin{pgfscope}%
\definecolor{textcolor}{rgb}{0.000000,0.000000,0.000000}%
\pgfsetstrokecolor{textcolor}%
\pgfsetfillcolor{textcolor}%
\pgftext[x=2.567725in,y=1.595830in,left,top]{\color{textcolor}{\rmfamily\fontsize{9.000000}{10.800000}\selectfont\catcode`\^=\active\def^{\ifmmode\sp\else\^{}\fi}\catcode`\%=\active\def%{\%}$(d)$}}%
\end{pgfscope}%
\begin{pgfscope}%
\pgfsetbuttcap%
\pgfsetmiterjoin%
\definecolor{currentfill}{rgb}{0.000000,0.000000,0.000000}%
\pgfsetfillcolor{currentfill}%
\pgfsetfillopacity{0.000000}%
\pgfsetlinewidth{0.803000pt}%
\definecolor{currentstroke}{rgb}{0.121569,0.466667,0.705882}%
\pgfsetstrokecolor{currentstroke}%
\pgfsetdash{}{0pt}%
\pgfsys@defobject{currentmarker}{\pgfqpoint{-0.017500in}{-0.017500in}}{\pgfqpoint{0.017500in}{0.017500in}}{%
\pgfpathmoveto{\pgfqpoint{-0.017500in}{-0.017500in}}%
\pgfpathlineto{\pgfqpoint{0.017500in}{-0.017500in}}%
\pgfpathlineto{\pgfqpoint{0.017500in}{0.017500in}}%
\pgfpathlineto{\pgfqpoint{-0.017500in}{0.017500in}}%
\pgfpathlineto{\pgfqpoint{-0.017500in}{-0.017500in}}%
\pgfpathclose%
\pgfusepath{stroke,fill}%
}%
\begin{pgfscope}%
\pgfsys@transformshift{3.441499in}{0.806172in}%
\pgfsys@useobject{currentmarker}{}%
\end{pgfscope}%
\end{pgfscope}%
\begin{pgfscope}%
\definecolor{textcolor}{rgb}{0.000000,0.000000,0.000000}%
\pgfsetstrokecolor{textcolor}%
\pgfsetfillcolor{textcolor}%
\pgftext[x=3.582471in,y=0.772145in,left,base]{\color{textcolor}{\rmfamily\fontsize{7.000000}{8.400000}\selectfont\catcode`\^=\active\def^{\ifmmode\sp\else\^{}\fi}\catcode`\%=\active\def%{\%}ideal ports}}%
\end{pgfscope}%
\begin{pgfscope}%
\pgfsetbuttcap%
\pgfsetroundjoin%
\definecolor{currentfill}{rgb}{0.000000,0.000000,0.000000}%
\pgfsetfillcolor{currentfill}%
\pgfsetfillopacity{0.000000}%
\pgfsetlinewidth{0.803000pt}%
\definecolor{currentstroke}{rgb}{1.000000,0.000000,0.000000}%
\pgfsetstrokecolor{currentstroke}%
\pgfsetdash{}{0pt}%
\pgfsys@defobject{currentmarker}{\pgfqpoint{-0.017500in}{-0.017500in}}{\pgfqpoint{0.017500in}{0.017500in}}{%
\pgfpathmoveto{\pgfqpoint{0.000000in}{-0.017500in}}%
\pgfpathcurveto{\pgfqpoint{0.004641in}{-0.017500in}}{\pgfqpoint{0.009093in}{-0.015656in}}{\pgfqpoint{0.012374in}{-0.012374in}}%
\pgfpathcurveto{\pgfqpoint{0.015656in}{-0.009093in}}{\pgfqpoint{0.017500in}{-0.004641in}}{\pgfqpoint{0.017500in}{0.000000in}}%
\pgfpathcurveto{\pgfqpoint{0.017500in}{0.004641in}}{\pgfqpoint{0.015656in}{0.009093in}}{\pgfqpoint{0.012374in}{0.012374in}}%
\pgfpathcurveto{\pgfqpoint{0.009093in}{0.015656in}}{\pgfqpoint{0.004641in}{0.017500in}}{\pgfqpoint{0.000000in}{0.017500in}}%
\pgfpathcurveto{\pgfqpoint{-0.004641in}{0.017500in}}{\pgfqpoint{-0.009093in}{0.015656in}}{\pgfqpoint{-0.012374in}{0.012374in}}%
\pgfpathcurveto{\pgfqpoint{-0.015656in}{0.009093in}}{\pgfqpoint{-0.017500in}{0.004641in}}{\pgfqpoint{-0.017500in}{0.000000in}}%
\pgfpathcurveto{\pgfqpoint{-0.017500in}{-0.004641in}}{\pgfqpoint{-0.015656in}{-0.009093in}}{\pgfqpoint{-0.012374in}{-0.012374in}}%
\pgfpathcurveto{\pgfqpoint{-0.009093in}{-0.015656in}}{\pgfqpoint{-0.004641in}{-0.017500in}}{\pgfqpoint{0.000000in}{-0.017500in}}%
\pgfpathlineto{\pgfqpoint{0.000000in}{-0.017500in}}%
\pgfpathclose%
\pgfusepath{stroke,fill}%
}%
\begin{pgfscope}%
\pgfsys@transformshift{3.441499in}{0.694907in}%
\pgfsys@useobject{currentmarker}{}%
\end{pgfscope}%
\end{pgfscope}%
\begin{pgfscope}%
\definecolor{textcolor}{rgb}{0.000000,0.000000,0.000000}%
\pgfsetstrokecolor{textcolor}%
\pgfsetfillcolor{textcolor}%
\pgftext[x=3.582471in,y=0.660879in,left,base]{\color{textcolor}{\rmfamily\fontsize{7.000000}{8.400000}\selectfont\catcode`\^=\active\def^{\ifmmode\sp\else\^{}\fi}\catcode`\%=\active\def%{\%}real ports}}%
\end{pgfscope}%
\begin{pgfscope}%
\pgfsetbuttcap%
\pgfsetroundjoin%
\definecolor{currentfill}{rgb}{0.172549,0.627451,0.172549}%
\pgfsetfillcolor{currentfill}%
\pgfsetlinewidth{0.803000pt}%
\definecolor{currentstroke}{rgb}{0.172549,0.627451,0.172549}%
\pgfsetstrokecolor{currentstroke}%
\pgfsetdash{}{0pt}%
\pgfsys@defobject{currentmarker}{\pgfqpoint{-0.019444in}{-0.019444in}}{\pgfqpoint{0.019444in}{0.019444in}}{%
\pgfpathmoveto{\pgfqpoint{-0.019444in}{-0.019444in}}%
\pgfpathlineto{\pgfqpoint{0.019444in}{0.019444in}}%
\pgfpathmoveto{\pgfqpoint{-0.019444in}{0.019444in}}%
\pgfpathlineto{\pgfqpoint{0.019444in}{-0.019444in}}%
\pgfusepath{stroke,fill}%
}%
\begin{pgfscope}%
\pgfsys@transformshift{3.441499in}{0.583642in}%
\pgfsys@useobject{currentmarker}{}%
\end{pgfscope}%
\end{pgfscope}%
\begin{pgfscope}%
\definecolor{textcolor}{rgb}{0.000000,0.000000,0.000000}%
\pgfsetstrokecolor{textcolor}%
\pgfsetfillcolor{textcolor}%
\pgftext[x=3.582471in,y=0.549614in,left,base]{\color{textcolor}{\rmfamily\fontsize{7.000000}{8.400000}\selectfont\catcode`\^=\active\def^{\ifmmode\sp\else\^{}\fi}\catcode`\%=\active\def%{\%}ideal, $\chi=0$}}%
\end{pgfscope}%
\begin{pgfscope}%
\pgfsetbuttcap%
\pgfsetroundjoin%
\definecolor{currentfill}{rgb}{0.121569,0.466667,0.705882}%
\pgfsetfillcolor{currentfill}%
\pgfsetlinewidth{0.803000pt}%
\definecolor{currentstroke}{rgb}{0.121569,0.466667,0.705882}%
\pgfsetstrokecolor{currentstroke}%
\pgfsetdash{}{0pt}%
\pgfsys@defobject{currentmarker}{\pgfqpoint{-0.019444in}{-0.019444in}}{\pgfqpoint{0.019444in}{0.019444in}}{%
\pgfpathmoveto{\pgfqpoint{-0.019444in}{-0.019444in}}%
\pgfpathlineto{\pgfqpoint{0.019444in}{0.019444in}}%
\pgfpathmoveto{\pgfqpoint{-0.019444in}{0.019444in}}%
\pgfpathlineto{\pgfqpoint{0.019444in}{-0.019444in}}%
\pgfusepath{stroke,fill}%
}%
\begin{pgfscope}%
\pgfsys@transformshift{0.505367in}{3.600371in}%
\pgfsys@useobject{currentmarker}{}%
\end{pgfscope}%
\end{pgfscope}%
\begin{pgfscope}%
\definecolor{textcolor}{rgb}{0.000000,0.000000,0.000000}%
\pgfsetstrokecolor{textcolor}%
\pgfsetfillcolor{textcolor}%
\pgftext[x=0.672034in,y=3.561482in,left,base]{\color{textcolor}{\rmfamily\fontsize{8.000000}{9.600000}\selectfont\catcode`\^=\active\def^{\ifmmode\sp\else\^{}\fi}\catcode`\%=\active\def%{\%}$T\omega_b=500$}}%
\end{pgfscope}%
\begin{pgfscope}%
\pgfsetbuttcap%
\pgfsetroundjoin%
\definecolor{currentfill}{rgb}{0.000000,0.000000,0.000000}%
\pgfsetfillcolor{currentfill}%
\pgfsetfillopacity{0.000000}%
\pgfsetlinewidth{0.803000pt}%
\definecolor{currentstroke}{rgb}{1.000000,0.000000,0.000000}%
\pgfsetstrokecolor{currentstroke}%
\pgfsetdash{}{0pt}%
\pgfsys@defobject{currentmarker}{\pgfqpoint{-0.017500in}{-0.017500in}}{\pgfqpoint{0.017500in}{0.017500in}}{%
\pgfpathmoveto{\pgfqpoint{0.000000in}{-0.017500in}}%
\pgfpathcurveto{\pgfqpoint{0.004641in}{-0.017500in}}{\pgfqpoint{0.009093in}{-0.015656in}}{\pgfqpoint{0.012374in}{-0.012374in}}%
\pgfpathcurveto{\pgfqpoint{0.015656in}{-0.009093in}}{\pgfqpoint{0.017500in}{-0.004641in}}{\pgfqpoint{0.017500in}{0.000000in}}%
\pgfpathcurveto{\pgfqpoint{0.017500in}{0.004641in}}{\pgfqpoint{0.015656in}{0.009093in}}{\pgfqpoint{0.012374in}{0.012374in}}%
\pgfpathcurveto{\pgfqpoint{0.009093in}{0.015656in}}{\pgfqpoint{0.004641in}{0.017500in}}{\pgfqpoint{0.000000in}{0.017500in}}%
\pgfpathcurveto{\pgfqpoint{-0.004641in}{0.017500in}}{\pgfqpoint{-0.009093in}{0.015656in}}{\pgfqpoint{-0.012374in}{0.012374in}}%
\pgfpathcurveto{\pgfqpoint{-0.015656in}{0.009093in}}{\pgfqpoint{-0.017500in}{0.004641in}}{\pgfqpoint{-0.017500in}{0.000000in}}%
\pgfpathcurveto{\pgfqpoint{-0.017500in}{-0.004641in}}{\pgfqpoint{-0.015656in}{-0.009093in}}{\pgfqpoint{-0.012374in}{-0.012374in}}%
\pgfpathcurveto{\pgfqpoint{-0.009093in}{-0.015656in}}{\pgfqpoint{-0.004641in}{-0.017500in}}{\pgfqpoint{0.000000in}{-0.017500in}}%
\pgfpathlineto{\pgfqpoint{0.000000in}{-0.017500in}}%
\pgfpathclose%
\pgfusepath{stroke,fill}%
}%
\begin{pgfscope}%
\pgfsys@transformshift{1.406539in}{3.600371in}%
\pgfsys@useobject{currentmarker}{}%
\end{pgfscope}%
\end{pgfscope}%
\begin{pgfscope}%
\definecolor{textcolor}{rgb}{0.000000,0.000000,0.000000}%
\pgfsetstrokecolor{textcolor}%
\pgfsetfillcolor{textcolor}%
\pgftext[x=1.573205in,y=3.561482in,left,base]{\color{textcolor}{\rmfamily\fontsize{8.000000}{9.600000}\selectfont\catcode`\^=\active\def^{\ifmmode\sp\else\^{}\fi}\catcode`\%=\active\def%{\%}$T\omega_b=2000$}}%
\end{pgfscope}%
\begin{pgfscope}%
\pgfsetbuttcap%
\pgfsetmiterjoin%
\definecolor{currentfill}{rgb}{0.000000,0.000000,0.000000}%
\pgfsetfillcolor{currentfill}%
\pgfsetfillopacity{0.000000}%
\pgfsetlinewidth{0.803000pt}%
\definecolor{currentstroke}{rgb}{0.172549,0.627451,0.172549}%
\pgfsetstrokecolor{currentstroke}%
\pgfsetdash{}{0pt}%
\pgfsys@defobject{currentmarker}{\pgfqpoint{-0.017500in}{-0.017500in}}{\pgfqpoint{0.017500in}{0.017500in}}{%
\pgfpathmoveto{\pgfqpoint{-0.017500in}{-0.017500in}}%
\pgfpathlineto{\pgfqpoint{0.017500in}{-0.017500in}}%
\pgfpathlineto{\pgfqpoint{0.017500in}{0.017500in}}%
\pgfpathlineto{\pgfqpoint{-0.017500in}{0.017500in}}%
\pgfpathlineto{\pgfqpoint{-0.017500in}{-0.017500in}}%
\pgfpathclose%
\pgfusepath{stroke,fill}%
}%
\begin{pgfscope}%
\pgfsys@transformshift{2.366738in}{3.600371in}%
\pgfsys@useobject{currentmarker}{}%
\end{pgfscope}%
\end{pgfscope}%
\begin{pgfscope}%
\definecolor{textcolor}{rgb}{0.000000,0.000000,0.000000}%
\pgfsetstrokecolor{textcolor}%
\pgfsetfillcolor{textcolor}%
\pgftext[x=2.533405in,y=3.561482in,left,base]{\color{textcolor}{\rmfamily\fontsize{8.000000}{9.600000}\selectfont\catcode`\^=\active\def^{\ifmmode\sp\else\^{}\fi}\catcode`\%=\active\def%{\%}$T\omega_b=8000$}}%
\end{pgfscope}%
\begin{pgfscope}%
\pgfsetrectcap%
\pgfsetroundjoin%
\pgfsetlinewidth{0.903375pt}%
\definecolor{currentstroke}{rgb}{0.000000,0.000000,0.000000}%
\pgfsetstrokecolor{currentstroke}%
\pgfsetdash{}{0pt}%
\pgfpathmoveto{\pgfqpoint{3.249160in}{3.600371in}}%
\pgfpathlineto{\pgfqpoint{3.326938in}{3.600371in}}%
\pgfpathlineto{\pgfqpoint{3.404716in}{3.600371in}}%
\pgfusepath{stroke}%
\end{pgfscope}%
\begin{pgfscope}%
\definecolor{textcolor}{rgb}{0.000000,0.000000,0.000000}%
\pgfsetstrokecolor{textcolor}%
\pgfsetfillcolor{textcolor}%
\pgftext[x=3.493605in,y=3.561482in,left,base]{\color{textcolor}{\rmfamily\fontsize{8.000000}{9.600000}\selectfont\catcode`\^=\active\def^{\ifmmode\sp\else\^{}\fi}\catcode`\%=\active\def%{\%}passive}}%
\end{pgfscope}%
\end{pgfpicture}%
\makeatother%
\endgroup%
}
\caption{(En color en l\'inea) Fidelidad condicional $F_h$ de la memoria en funci\'on de (a) la tasa del puerto de entrada $\kappa_{\rm in}/\kappa$ con lectura ideal, (b) la tasa del puerto de salida $\kappa_{\rm out}/\kappa$, (c) el factor de correlaci\'on del ruido $\chi$ (Sec.~\ref{sec:noise}) y (d) la escala $S$ de los acoplos por deformaci\'on $\lambda_i\to S\lambda_i$. En (a)--(c) los s\'imbolos corresponden a las esperas $T\omega_b=500$, $2000$ y $8000$; en (d), a $T\omega_b=2000$ con los pulsos reoptimizados para cada $S$, las curvas son para puertos ideales (cuadrados), puertos reales (c\'irculos) y puertos ideales con ruido no correlacionado, $\chi=0$ (cruces). Las l\'ineas negras son la referencia pasiva. Las l\'ineas son interpolaciones mon\'otonas de los puntos calculados. Par\'ametros del sistema: modelo reducido, los mismos que en la Fig.~\ref{fig:sweep} salvo los que se var\'ian en cada panel ($S=1$ por defecto); tolerancia de integraci\'on $10^{-9}$.}
\label{fig:ports}
\end{figure}

La Figura~\ref{fig:ports}(c) y la Fig.~\ref{fig:requirements}(a) muestran qu\'e supuestos sobre el entorno determinan la fidelidad condicional. La Fig.~\ref{fig:ports}(c) var\'ia el grado $\chi$ de correlaci\'on entre el ruido que act\'ua sobre las dos ranuras: $F_h$ crece mon\'otonamente con $\chi$, porque las fluctuaciones de modo com\'un desplazan ambas ranuras por igual y dejan invariante la fase relativa almacenada, mientras que la parte no correlacionada la desfasa. La protecci\'on de la coherencia descansa por tanto en la correlaci\'on del ruido, que es una condici\'on del modelo y no una propiedad derivada (Sec.~\ref{sec:noise}). Para ruido no correlacionado y esperas largas la memoria aun supera a los modos pasivos en $F$ y $s$, pero su fidelidad condicional est\'a muy por debajo del valor pasivo, lo que muestra que la protecci\'on por ruido com\'un, y no solo el car\'acter oscuro de los estados, es lo que preserva el qubit. La dependencia con $\chi$ se acent\'ua con la espera: tras una espera corta el ruido no correlacionado a\'un no ha desfasado la fase almacenada y $F_h$ es casi plana en $\chi$, mientras que tras una espera larga es la correlaci\'on del ruido la que decide cu\'anta coherencia sobrevive, como se espera de la tasa de desfase entre las ranuras, proporcional a $1-\chi$ (Ec.~\eqref{eq:gamma_dark}).

La Figura~\ref{fig:requirements}(a) representa $F_h$ en el plano de $\chi$ y la espera, con los contornos $F_h=0.8$ y $F_h=1/2$, por encima del cual el estado condicional es entrelazado. Los contornos se curvan hacia $\chi=1$ al crecer la espera: cuanto m\'as largo es el almacenamiento, m\'as cerca de un ruido perfectamente com\'un debe estar el entorno, y para las esperas m\'as largas la regi\'on de entrelazamiento certificado se reduce a $\chi$ grandes.

Un segundo supuesto es que el ruido se acopla con igual intensidad a todos los puntos. La Figura~\ref{fig:requirements}(b) lo relaja tomando el operador de ruido como $\sum_i a_i n_i$ con respuestas Stark desiguales $a_i$ y ruido perfectamente correlacionado. Cuando un punto se desv\'ia de los dem\'as (patr\'on A), las dos ranuras adquieren desplazamientos de ruido ligeramente distintos, un efecto de segundo orden en $\delta$, y la p\'erdida de $F_h$ es peque\~na. Cuando dos puntos se desv\'ian en sentidos opuestos (patr\'on B), el ruido adquiere un elemento de matriz entre las dos ranuras e induce transiciones entre ellas, que se acumulan con la espera y dan una p\'erdida mayor. En ambos casos las desviaciones del orden del diez por ciento cuestan poco. El requisito es por tanto la correlaci\'on espacial del campo que act\'ua sobre los puntos, no una igualdad exacta de sus respuestas Stark.

\begin{figure}[t]
\centering
\resizebox{\columnwidth}{!}{%% Creator: Matplotlib, PGF backend
%%
%% To include the figure in your LaTeX document, write
%%   \input{<filename>.pgf}
%%
%% Make sure the required packages are loaded in your preamble
%%   \usepackage{pgf}
%%
%% Also ensure that all the required font packages are loaded; for instance,
%% the lmodern package is sometimes necessary when using math font.
%%   \usepackage{lmodern}
%%
%% Figures using additional raster images can only be included by \input if
%% they are in the same directory as the main LaTeX file. For loading figures
%% from other directories you can use the `import` package
%%   \usepackage{import}
%%
%% and then include the figures with
%%   \import{<path to file>}{<filename>.pgf}
%%
%% Matplotlib used the following preamble
%%   \def\mathdefault#1{#1}
%%   \everymath=\expandafter{\the\everymath\displaystyle}
%%   \IfFileExists{scrextend.sty}{
%%     \usepackage[fontsize=8.000000pt]{scrextend}
%%   }{
%%     \renewcommand{\normalsize}{\fontsize{8.000000}{9.600000}\selectfont}
%%     \normalsize
%%   }
%%   \usepackage{amsmath,amssymb}
%%   \makeatletter\@ifpackageloaded{underscore}{}{\usepackage[strings]{underscore}}\makeatother
%%
\begingroup%
\makeatletter%
\begin{pgfpicture}%
\pgfpathrectangle{\pgfpointorigin}{\pgfqpoint{4.910872in}{2.038457in}}%
\pgfusepath{use as bounding box, clip}%
\begin{pgfscope}%
\pgfsetbuttcap%
\pgfsetmiterjoin%
\definecolor{currentfill}{rgb}{1.000000,1.000000,1.000000}%
\pgfsetfillcolor{currentfill}%
\pgfsetlinewidth{0.000000pt}%
\definecolor{currentstroke}{rgb}{1.000000,1.000000,1.000000}%
\pgfsetstrokecolor{currentstroke}%
\pgfsetdash{}{0pt}%
\pgfpathmoveto{\pgfqpoint{0.000000in}{0.000000in}}%
\pgfpathlineto{\pgfqpoint{4.910872in}{0.000000in}}%
\pgfpathlineto{\pgfqpoint{4.910872in}{2.038457in}}%
\pgfpathlineto{\pgfqpoint{0.000000in}{2.038457in}}%
\pgfpathlineto{\pgfqpoint{0.000000in}{0.000000in}}%
\pgfpathclose%
\pgfusepath{fill}%
\end{pgfscope}%
\begin{pgfscope}%
\pgfsetbuttcap%
\pgfsetmiterjoin%
\definecolor{currentfill}{rgb}{1.000000,1.000000,1.000000}%
\pgfsetfillcolor{currentfill}%
\pgfsetlinewidth{0.000000pt}%
\definecolor{currentstroke}{rgb}{0.000000,0.000000,0.000000}%
\pgfsetstrokecolor{currentstroke}%
\pgfsetstrokeopacity{0.000000}%
\pgfsetdash{}{0pt}%
\pgfpathmoveto{\pgfqpoint{0.415139in}{0.462654in}}%
\pgfpathlineto{\pgfqpoint{2.056734in}{0.462654in}}%
\pgfpathlineto{\pgfqpoint{2.056734in}{1.899877in}}%
\pgfpathlineto{\pgfqpoint{0.415139in}{1.899877in}}%
\pgfpathlineto{\pgfqpoint{0.415139in}{0.462654in}}%
\pgfpathclose%
\pgfusepath{fill}%
\end{pgfscope}%
\begin{pgfscope}%
\pgfpathrectangle{\pgfqpoint{0.415139in}{0.462654in}}{\pgfqpoint{1.641596in}{1.437222in}}%
\pgfusepath{clip}%
\pgfsetbuttcap%
\pgfsetroundjoin%
\definecolor{currentfill}{rgb}{0.956863,0.427451,0.262745}%
\pgfsetfillcolor{currentfill}%
\pgfsetlinewidth{0.000000pt}%
\definecolor{currentstroke}{rgb}{0.000000,0.000000,0.000000}%
\pgfsetstrokecolor{currentstroke}%
\pgfsetdash{}{0pt}%
\pgfpathmoveto{\pgfqpoint{0.415577in}{1.430432in}}%
\pgfpathlineto{\pgfqpoint{0.423388in}{1.435557in}}%
\pgfpathlineto{\pgfqpoint{0.426575in}{1.437654in}}%
\pgfpathlineto{\pgfqpoint{0.431637in}{1.441009in}}%
\pgfpathlineto{\pgfqpoint{0.437456in}{1.444877in}}%
\pgfpathlineto{\pgfqpoint{0.439887in}{1.446503in}}%
\pgfpathlineto{\pgfqpoint{0.448136in}{1.452042in}}%
\pgfpathlineto{\pgfqpoint{0.448221in}{1.452099in}}%
\pgfpathlineto{\pgfqpoint{0.456385in}{1.457641in}}%
\pgfpathlineto{\pgfqpoint{0.458852in}{1.459321in}}%
\pgfpathlineto{\pgfqpoint{0.464634in}{1.463286in}}%
\pgfpathlineto{\pgfqpoint{0.469369in}{1.466543in}}%
\pgfpathlineto{\pgfqpoint{0.472883in}{1.468977in}}%
\pgfpathlineto{\pgfqpoint{0.479773in}{1.473766in}}%
\pgfpathlineto{\pgfqpoint{0.481133in}{1.474717in}}%
\pgfpathlineto{\pgfqpoint{0.489382in}{1.480510in}}%
\pgfpathlineto{\pgfqpoint{0.490059in}{1.480988in}}%
\pgfpathlineto{\pgfqpoint{0.497631in}{1.486362in}}%
\pgfpathlineto{\pgfqpoint{0.500224in}{1.488210in}}%
\pgfpathlineto{\pgfqpoint{0.505880in}{1.492265in}}%
\pgfpathlineto{\pgfqpoint{0.510280in}{1.495432in}}%
\pgfpathlineto{\pgfqpoint{0.514130in}{1.498220in}}%
\pgfpathlineto{\pgfqpoint{0.520227in}{1.502654in}}%
\pgfpathlineto{\pgfqpoint{0.522379in}{1.504228in}}%
\pgfpathlineto{\pgfqpoint{0.530066in}{1.509877in}}%
\pgfpathlineto{\pgfqpoint{0.530628in}{1.510292in}}%
\pgfpathlineto{\pgfqpoint{0.538877in}{1.516416in}}%
\pgfpathlineto{\pgfqpoint{0.539792in}{1.517099in}}%
\pgfpathlineto{\pgfqpoint{0.547126in}{1.522601in}}%
\pgfpathlineto{\pgfqpoint{0.549408in}{1.524321in}}%
\pgfpathlineto{\pgfqpoint{0.555376in}{1.528844in}}%
\pgfpathlineto{\pgfqpoint{0.558920in}{1.531543in}}%
\pgfpathlineto{\pgfqpoint{0.563625in}{1.535145in}}%
\pgfpathlineto{\pgfqpoint{0.568330in}{1.538766in}}%
\pgfpathlineto{\pgfqpoint{0.571874in}{1.541506in}}%
\pgfpathlineto{\pgfqpoint{0.577639in}{1.545988in}}%
\pgfpathlineto{\pgfqpoint{0.580123in}{1.547929in}}%
\pgfpathlineto{\pgfqpoint{0.586848in}{1.553210in}}%
\pgfpathlineto{\pgfqpoint{0.588373in}{1.554413in}}%
\pgfpathlineto{\pgfqpoint{0.595957in}{1.560432in}}%
\pgfpathlineto{\pgfqpoint{0.596622in}{1.560962in}}%
\pgfpathlineto{\pgfqpoint{0.604871in}{1.567576in}}%
\pgfpathlineto{\pgfqpoint{0.604969in}{1.567654in}}%
\pgfpathlineto{\pgfqpoint{0.613120in}{1.574259in}}%
\pgfpathlineto{\pgfqpoint{0.613878in}{1.574877in}}%
\pgfpathlineto{\pgfqpoint{0.621369in}{1.581009in}}%
\pgfpathlineto{\pgfqpoint{0.622694in}{1.582099in}}%
\pgfpathlineto{\pgfqpoint{0.629619in}{1.587825in}}%
\pgfpathlineto{\pgfqpoint{0.631417in}{1.589321in}}%
\pgfpathlineto{\pgfqpoint{0.637868in}{1.594710in}}%
\pgfpathlineto{\pgfqpoint{0.640050in}{1.596543in}}%
\pgfpathlineto{\pgfqpoint{0.646117in}{1.601662in}}%
\pgfpathlineto{\pgfqpoint{0.648596in}{1.603766in}}%
\pgfpathlineto{\pgfqpoint{0.654366in}{1.608684in}}%
\pgfpathlineto{\pgfqpoint{0.657054in}{1.610988in}}%
\pgfpathlineto{\pgfqpoint{0.662616in}{1.615775in}}%
\pgfpathlineto{\pgfqpoint{0.665428in}{1.618210in}}%
\pgfpathlineto{\pgfqpoint{0.670865in}{1.622936in}}%
\pgfpathlineto{\pgfqpoint{0.673719in}{1.625432in}}%
\pgfpathlineto{\pgfqpoint{0.679114in}{1.630168in}}%
\pgfpathlineto{\pgfqpoint{0.681929in}{1.632654in}}%
\pgfpathlineto{\pgfqpoint{0.687363in}{1.637472in}}%
\pgfpathlineto{\pgfqpoint{0.690059in}{1.639877in}}%
\pgfpathlineto{\pgfqpoint{0.695612in}{1.644849in}}%
\pgfpathlineto{\pgfqpoint{0.698111in}{1.647099in}}%
\pgfpathlineto{\pgfqpoint{0.703862in}{1.652298in}}%
\pgfpathlineto{\pgfqpoint{0.706086in}{1.654321in}}%
\pgfpathlineto{\pgfqpoint{0.712111in}{1.659821in}}%
\pgfpathlineto{\pgfqpoint{0.713986in}{1.661543in}}%
\pgfpathlineto{\pgfqpoint{0.720360in}{1.667419in}}%
\pgfpathlineto{\pgfqpoint{0.721812in}{1.668766in}}%
\pgfpathlineto{\pgfqpoint{0.728609in}{1.675093in}}%
\pgfpathlineto{\pgfqpoint{0.729564in}{1.675988in}}%
\pgfpathlineto{\pgfqpoint{0.736859in}{1.682845in}}%
\pgfpathlineto{\pgfqpoint{0.737244in}{1.683210in}}%
\pgfpathlineto{\pgfqpoint{0.744850in}{1.690432in}}%
\pgfpathlineto{\pgfqpoint{0.745108in}{1.690677in}}%
\pgfpathlineto{\pgfqpoint{0.752382in}{1.697654in}}%
\pgfpathlineto{\pgfqpoint{0.753357in}{1.698593in}}%
\pgfpathlineto{\pgfqpoint{0.759843in}{1.704877in}}%
\pgfpathlineto{\pgfqpoint{0.761606in}{1.706590in}}%
\pgfpathlineto{\pgfqpoint{0.767235in}{1.712099in}}%
\pgfpathlineto{\pgfqpoint{0.769856in}{1.714671in}}%
\pgfpathlineto{\pgfqpoint{0.774559in}{1.719321in}}%
\pgfpathlineto{\pgfqpoint{0.778105in}{1.722837in}}%
\pgfpathlineto{\pgfqpoint{0.781815in}{1.726543in}}%
\pgfpathlineto{\pgfqpoint{0.786354in}{1.731089in}}%
\pgfpathlineto{\pgfqpoint{0.789006in}{1.733766in}}%
\pgfpathlineto{\pgfqpoint{0.794603in}{1.739429in}}%
\pgfpathlineto{\pgfqpoint{0.796131in}{1.740988in}}%
\pgfpathlineto{\pgfqpoint{0.802852in}{1.747859in}}%
\pgfpathlineto{\pgfqpoint{0.803193in}{1.748210in}}%
\pgfpathlineto{\pgfqpoint{0.810183in}{1.755432in}}%
\pgfpathlineto{\pgfqpoint{0.811102in}{1.756383in}}%
\pgfpathlineto{\pgfqpoint{0.817108in}{1.762654in}}%
\pgfpathlineto{\pgfqpoint{0.819351in}{1.765002in}}%
\pgfpathlineto{\pgfqpoint{0.823972in}{1.769877in}}%
\pgfpathlineto{\pgfqpoint{0.827600in}{1.773713in}}%
\pgfpathlineto{\pgfqpoint{0.830777in}{1.777099in}}%
\pgfpathlineto{\pgfqpoint{0.835849in}{1.782518in}}%
\pgfpathlineto{\pgfqpoint{0.837524in}{1.784321in}}%
\pgfpathlineto{\pgfqpoint{0.844099in}{1.791418in}}%
\pgfpathlineto{\pgfqpoint{0.844213in}{1.791543in}}%
\pgfpathlineto{\pgfqpoint{0.850833in}{1.798766in}}%
\pgfpathlineto{\pgfqpoint{0.852348in}{1.800422in}}%
\pgfpathlineto{\pgfqpoint{0.857396in}{1.805988in}}%
\pgfpathlineto{\pgfqpoint{0.860597in}{1.809524in}}%
\pgfpathlineto{\pgfqpoint{0.863906in}{1.813210in}}%
\pgfpathlineto{\pgfqpoint{0.868846in}{1.818725in}}%
\pgfpathlineto{\pgfqpoint{0.870363in}{1.820432in}}%
\pgfpathlineto{\pgfqpoint{0.876766in}{1.827654in}}%
\pgfpathlineto{\pgfqpoint{0.877095in}{1.828027in}}%
\pgfpathlineto{\pgfqpoint{0.883103in}{1.834877in}}%
\pgfpathlineto{\pgfqpoint{0.885345in}{1.837437in}}%
\pgfpathlineto{\pgfqpoint{0.889391in}{1.842099in}}%
\pgfpathlineto{\pgfqpoint{0.893594in}{1.846950in}}%
\pgfpathlineto{\pgfqpoint{0.895631in}{1.849321in}}%
\pgfpathlineto{\pgfqpoint{0.901823in}{1.856543in}}%
\pgfpathlineto{\pgfqpoint{0.901843in}{1.856567in}}%
\pgfpathlineto{\pgfqpoint{0.907949in}{1.863766in}}%
\pgfpathlineto{\pgfqpoint{0.910092in}{1.866297in}}%
\pgfpathlineto{\pgfqpoint{0.914029in}{1.870988in}}%
\pgfpathlineto{\pgfqpoint{0.918342in}{1.876134in}}%
\pgfpathlineto{\pgfqpoint{0.920065in}{1.878210in}}%
\pgfpathlineto{\pgfqpoint{0.926053in}{1.885432in}}%
\pgfpathlineto{\pgfqpoint{0.926591in}{1.886083in}}%
\pgfpathlineto{\pgfqpoint{0.931979in}{1.892654in}}%
\pgfpathlineto{\pgfqpoint{0.934840in}{1.896149in}}%
\pgfpathlineto{\pgfqpoint{0.937864in}{1.899877in}}%
\pgfpathlineto{\pgfqpoint{0.934840in}{1.899877in}}%
\pgfpathlineto{\pgfqpoint{0.926591in}{1.899877in}}%
\pgfpathlineto{\pgfqpoint{0.918342in}{1.899877in}}%
\pgfpathlineto{\pgfqpoint{0.910092in}{1.899877in}}%
\pgfpathlineto{\pgfqpoint{0.901843in}{1.899877in}}%
\pgfpathlineto{\pgfqpoint{0.893594in}{1.899877in}}%
\pgfpathlineto{\pgfqpoint{0.885345in}{1.899877in}}%
\pgfpathlineto{\pgfqpoint{0.877095in}{1.899877in}}%
\pgfpathlineto{\pgfqpoint{0.868846in}{1.899877in}}%
\pgfpathlineto{\pgfqpoint{0.860597in}{1.899877in}}%
\pgfpathlineto{\pgfqpoint{0.852348in}{1.899877in}}%
\pgfpathlineto{\pgfqpoint{0.844099in}{1.899877in}}%
\pgfpathlineto{\pgfqpoint{0.835849in}{1.899877in}}%
\pgfpathlineto{\pgfqpoint{0.827600in}{1.899877in}}%
\pgfpathlineto{\pgfqpoint{0.819351in}{1.899877in}}%
\pgfpathlineto{\pgfqpoint{0.811102in}{1.899877in}}%
\pgfpathlineto{\pgfqpoint{0.802852in}{1.899877in}}%
\pgfpathlineto{\pgfqpoint{0.794603in}{1.899877in}}%
\pgfpathlineto{\pgfqpoint{0.786354in}{1.899877in}}%
\pgfpathlineto{\pgfqpoint{0.778105in}{1.899877in}}%
\pgfpathlineto{\pgfqpoint{0.769856in}{1.899877in}}%
\pgfpathlineto{\pgfqpoint{0.761606in}{1.899877in}}%
\pgfpathlineto{\pgfqpoint{0.753357in}{1.899877in}}%
\pgfpathlineto{\pgfqpoint{0.745108in}{1.899877in}}%
\pgfpathlineto{\pgfqpoint{0.736859in}{1.899877in}}%
\pgfpathlineto{\pgfqpoint{0.728609in}{1.899877in}}%
\pgfpathlineto{\pgfqpoint{0.720360in}{1.899877in}}%
\pgfpathlineto{\pgfqpoint{0.712111in}{1.899877in}}%
\pgfpathlineto{\pgfqpoint{0.703862in}{1.899877in}}%
\pgfpathlineto{\pgfqpoint{0.695612in}{1.899877in}}%
\pgfpathlineto{\pgfqpoint{0.687363in}{1.899877in}}%
\pgfpathlineto{\pgfqpoint{0.679114in}{1.899877in}}%
\pgfpathlineto{\pgfqpoint{0.670865in}{1.899877in}}%
\pgfpathlineto{\pgfqpoint{0.662616in}{1.899877in}}%
\pgfpathlineto{\pgfqpoint{0.654366in}{1.899877in}}%
\pgfpathlineto{\pgfqpoint{0.646117in}{1.899877in}}%
\pgfpathlineto{\pgfqpoint{0.637868in}{1.899877in}}%
\pgfpathlineto{\pgfqpoint{0.629619in}{1.899877in}}%
\pgfpathlineto{\pgfqpoint{0.621369in}{1.899877in}}%
\pgfpathlineto{\pgfqpoint{0.613120in}{1.899877in}}%
\pgfpathlineto{\pgfqpoint{0.604871in}{1.899877in}}%
\pgfpathlineto{\pgfqpoint{0.596622in}{1.899877in}}%
\pgfpathlineto{\pgfqpoint{0.588373in}{1.899877in}}%
\pgfpathlineto{\pgfqpoint{0.580123in}{1.899877in}}%
\pgfpathlineto{\pgfqpoint{0.571874in}{1.899877in}}%
\pgfpathlineto{\pgfqpoint{0.563625in}{1.899877in}}%
\pgfpathlineto{\pgfqpoint{0.555376in}{1.899877in}}%
\pgfpathlineto{\pgfqpoint{0.547126in}{1.899877in}}%
\pgfpathlineto{\pgfqpoint{0.538877in}{1.899877in}}%
\pgfpathlineto{\pgfqpoint{0.530628in}{1.899877in}}%
\pgfpathlineto{\pgfqpoint{0.522379in}{1.899877in}}%
\pgfpathlineto{\pgfqpoint{0.514130in}{1.899877in}}%
\pgfpathlineto{\pgfqpoint{0.505880in}{1.899877in}}%
\pgfpathlineto{\pgfqpoint{0.497631in}{1.899877in}}%
\pgfpathlineto{\pgfqpoint{0.489382in}{1.899877in}}%
\pgfpathlineto{\pgfqpoint{0.481133in}{1.899877in}}%
\pgfpathlineto{\pgfqpoint{0.472883in}{1.899877in}}%
\pgfpathlineto{\pgfqpoint{0.464634in}{1.899877in}}%
\pgfpathlineto{\pgfqpoint{0.456385in}{1.899877in}}%
\pgfpathlineto{\pgfqpoint{0.448136in}{1.899877in}}%
\pgfpathlineto{\pgfqpoint{0.439887in}{1.899877in}}%
\pgfpathlineto{\pgfqpoint{0.431637in}{1.899877in}}%
\pgfpathlineto{\pgfqpoint{0.423388in}{1.899877in}}%
\pgfpathlineto{\pgfqpoint{0.415139in}{1.899877in}}%
\pgfpathlineto{\pgfqpoint{0.415139in}{1.892654in}}%
\pgfpathlineto{\pgfqpoint{0.415139in}{1.885432in}}%
\pgfpathlineto{\pgfqpoint{0.415139in}{1.878210in}}%
\pgfpathlineto{\pgfqpoint{0.415139in}{1.870988in}}%
\pgfpathlineto{\pgfqpoint{0.415139in}{1.863766in}}%
\pgfpathlineto{\pgfqpoint{0.415139in}{1.856543in}}%
\pgfpathlineto{\pgfqpoint{0.415139in}{1.849321in}}%
\pgfpathlineto{\pgfqpoint{0.415139in}{1.842099in}}%
\pgfpathlineto{\pgfqpoint{0.415139in}{1.834877in}}%
\pgfpathlineto{\pgfqpoint{0.415139in}{1.827654in}}%
\pgfpathlineto{\pgfqpoint{0.415139in}{1.820432in}}%
\pgfpathlineto{\pgfqpoint{0.415139in}{1.813210in}}%
\pgfpathlineto{\pgfqpoint{0.415139in}{1.805988in}}%
\pgfpathlineto{\pgfqpoint{0.415139in}{1.798766in}}%
\pgfpathlineto{\pgfqpoint{0.415139in}{1.791543in}}%
\pgfpathlineto{\pgfqpoint{0.415139in}{1.784321in}}%
\pgfpathlineto{\pgfqpoint{0.415139in}{1.777099in}}%
\pgfpathlineto{\pgfqpoint{0.415139in}{1.769877in}}%
\pgfpathlineto{\pgfqpoint{0.415139in}{1.762654in}}%
\pgfpathlineto{\pgfqpoint{0.415139in}{1.755432in}}%
\pgfpathlineto{\pgfqpoint{0.415139in}{1.748210in}}%
\pgfpathlineto{\pgfqpoint{0.415139in}{1.740988in}}%
\pgfpathlineto{\pgfqpoint{0.415139in}{1.733766in}}%
\pgfpathlineto{\pgfqpoint{0.415139in}{1.726543in}}%
\pgfpathlineto{\pgfqpoint{0.415139in}{1.719321in}}%
\pgfpathlineto{\pgfqpoint{0.415139in}{1.712099in}}%
\pgfpathlineto{\pgfqpoint{0.415139in}{1.704877in}}%
\pgfpathlineto{\pgfqpoint{0.415139in}{1.697654in}}%
\pgfpathlineto{\pgfqpoint{0.415139in}{1.690432in}}%
\pgfpathlineto{\pgfqpoint{0.415139in}{1.683210in}}%
\pgfpathlineto{\pgfqpoint{0.415139in}{1.675988in}}%
\pgfpathlineto{\pgfqpoint{0.415139in}{1.668766in}}%
\pgfpathlineto{\pgfqpoint{0.415139in}{1.661543in}}%
\pgfpathlineto{\pgfqpoint{0.415139in}{1.654321in}}%
\pgfpathlineto{\pgfqpoint{0.415139in}{1.647099in}}%
\pgfpathlineto{\pgfqpoint{0.415139in}{1.639877in}}%
\pgfpathlineto{\pgfqpoint{0.415139in}{1.632654in}}%
\pgfpathlineto{\pgfqpoint{0.415139in}{1.625432in}}%
\pgfpathlineto{\pgfqpoint{0.415139in}{1.618210in}}%
\pgfpathlineto{\pgfqpoint{0.415139in}{1.610988in}}%
\pgfpathlineto{\pgfqpoint{0.415139in}{1.603766in}}%
\pgfpathlineto{\pgfqpoint{0.415139in}{1.596543in}}%
\pgfpathlineto{\pgfqpoint{0.415139in}{1.589321in}}%
\pgfpathlineto{\pgfqpoint{0.415139in}{1.582099in}}%
\pgfpathlineto{\pgfqpoint{0.415139in}{1.574877in}}%
\pgfpathlineto{\pgfqpoint{0.415139in}{1.567654in}}%
\pgfpathlineto{\pgfqpoint{0.415139in}{1.560432in}}%
\pgfpathlineto{\pgfqpoint{0.415139in}{1.553210in}}%
\pgfpathlineto{\pgfqpoint{0.415139in}{1.545988in}}%
\pgfpathlineto{\pgfqpoint{0.415139in}{1.538766in}}%
\pgfpathlineto{\pgfqpoint{0.415139in}{1.531543in}}%
\pgfpathlineto{\pgfqpoint{0.415139in}{1.524321in}}%
\pgfpathlineto{\pgfqpoint{0.415139in}{1.517099in}}%
\pgfpathlineto{\pgfqpoint{0.415139in}{1.509877in}}%
\pgfpathlineto{\pgfqpoint{0.415139in}{1.502654in}}%
\pgfpathlineto{\pgfqpoint{0.415139in}{1.495432in}}%
\pgfpathlineto{\pgfqpoint{0.415139in}{1.488210in}}%
\pgfpathlineto{\pgfqpoint{0.415139in}{1.480988in}}%
\pgfpathlineto{\pgfqpoint{0.415139in}{1.473766in}}%
\pgfpathlineto{\pgfqpoint{0.415139in}{1.466543in}}%
\pgfpathlineto{\pgfqpoint{0.415139in}{1.459321in}}%
\pgfpathlineto{\pgfqpoint{0.415139in}{1.452099in}}%
\pgfpathlineto{\pgfqpoint{0.415139in}{1.444877in}}%
\pgfpathlineto{\pgfqpoint{0.415139in}{1.437654in}}%
\pgfpathlineto{\pgfqpoint{0.415139in}{1.430432in}}%
\pgfpathlineto{\pgfqpoint{0.415139in}{1.430147in}}%
\pgfpathlineto{\pgfqpoint{0.415577in}{1.430432in}}%
\pgfpathclose%
\pgfusepath{fill}%
\end{pgfscope}%
\begin{pgfscope}%
\pgfpathrectangle{\pgfqpoint{0.415139in}{0.462654in}}{\pgfqpoint{1.641596in}{1.437222in}}%
\pgfusepath{clip}%
\pgfsetbuttcap%
\pgfsetroundjoin%
\definecolor{currentfill}{rgb}{0.998078,0.940408,0.658670}%
\pgfsetfillcolor{currentfill}%
\pgfsetlinewidth{0.000000pt}%
\definecolor{currentstroke}{rgb}{0.000000,0.000000,0.000000}%
\pgfsetstrokecolor{currentstroke}%
\pgfsetdash{}{0pt}%
\pgfpathmoveto{\pgfqpoint{0.418442in}{1.047654in}}%
\pgfpathlineto{\pgfqpoint{0.423388in}{1.049693in}}%
\pgfpathlineto{\pgfqpoint{0.431637in}{1.053107in}}%
\pgfpathlineto{\pgfqpoint{0.435896in}{1.054877in}}%
\pgfpathlineto{\pgfqpoint{0.439887in}{1.056554in}}%
\pgfpathlineto{\pgfqpoint{0.448136in}{1.060036in}}%
\pgfpathlineto{\pgfqpoint{0.453001in}{1.062099in}}%
\pgfpathlineto{\pgfqpoint{0.456385in}{1.063550in}}%
\pgfpathlineto{\pgfqpoint{0.464634in}{1.067102in}}%
\pgfpathlineto{\pgfqpoint{0.469763in}{1.069321in}}%
\pgfpathlineto{\pgfqpoint{0.472883in}{1.070686in}}%
\pgfpathlineto{\pgfqpoint{0.481133in}{1.074310in}}%
\pgfpathlineto{\pgfqpoint{0.486192in}{1.076543in}}%
\pgfpathlineto{\pgfqpoint{0.489382in}{1.077967in}}%
\pgfpathlineto{\pgfqpoint{0.497631in}{1.081665in}}%
\pgfpathlineto{\pgfqpoint{0.502295in}{1.083766in}}%
\pgfpathlineto{\pgfqpoint{0.505880in}{1.085398in}}%
\pgfpathlineto{\pgfqpoint{0.514130in}{1.089172in}}%
\pgfpathlineto{\pgfqpoint{0.518081in}{1.090988in}}%
\pgfpathlineto{\pgfqpoint{0.522379in}{1.092985in}}%
\pgfpathlineto{\pgfqpoint{0.530628in}{1.096836in}}%
\pgfpathlineto{\pgfqpoint{0.533558in}{1.098210in}}%
\pgfpathlineto{\pgfqpoint{0.538877in}{1.100732in}}%
\pgfpathlineto{\pgfqpoint{0.547126in}{1.104663in}}%
\pgfpathlineto{\pgfqpoint{0.548734in}{1.105432in}}%
\pgfpathlineto{\pgfqpoint{0.555376in}{1.108646in}}%
\pgfpathlineto{\pgfqpoint{0.563617in}{1.112654in}}%
\pgfpathlineto{\pgfqpoint{0.563625in}{1.112658in}}%
\pgfpathlineto{\pgfqpoint{0.571874in}{1.116733in}}%
\pgfpathlineto{\pgfqpoint{0.578207in}{1.119877in}}%
\pgfpathlineto{\pgfqpoint{0.580123in}{1.120838in}}%
\pgfpathlineto{\pgfqpoint{0.588373in}{1.124997in}}%
\pgfpathlineto{\pgfqpoint{0.592519in}{1.127099in}}%
\pgfpathlineto{\pgfqpoint{0.596622in}{1.129200in}}%
\pgfpathlineto{\pgfqpoint{0.604871in}{1.133446in}}%
\pgfpathlineto{\pgfqpoint{0.606562in}{1.134321in}}%
\pgfpathlineto{\pgfqpoint{0.613120in}{1.137749in}}%
\pgfpathlineto{\pgfqpoint{0.620339in}{1.141543in}}%
\pgfpathlineto{\pgfqpoint{0.621369in}{1.142090in}}%
\pgfpathlineto{\pgfqpoint{0.629619in}{1.146492in}}%
\pgfpathlineto{\pgfqpoint{0.633856in}{1.148766in}}%
\pgfpathlineto{\pgfqpoint{0.637868in}{1.150939in}}%
\pgfpathlineto{\pgfqpoint{0.646117in}{1.155433in}}%
\pgfpathlineto{\pgfqpoint{0.647131in}{1.155988in}}%
\pgfpathlineto{\pgfqpoint{0.654366in}{1.159988in}}%
\pgfpathlineto{\pgfqpoint{0.660163in}{1.163210in}}%
\pgfpathlineto{\pgfqpoint{0.662616in}{1.164586in}}%
\pgfpathlineto{\pgfqpoint{0.670865in}{1.169240in}}%
\pgfpathlineto{\pgfqpoint{0.672967in}{1.170432in}}%
\pgfpathlineto{\pgfqpoint{0.679114in}{1.173952in}}%
\pgfpathlineto{\pgfqpoint{0.685548in}{1.177654in}}%
\pgfpathlineto{\pgfqpoint{0.687363in}{1.178709in}}%
\pgfpathlineto{\pgfqpoint{0.695612in}{1.183527in}}%
\pgfpathlineto{\pgfqpoint{0.697912in}{1.184877in}}%
\pgfpathlineto{\pgfqpoint{0.703862in}{1.188402in}}%
\pgfpathlineto{\pgfqpoint{0.710069in}{1.192099in}}%
\pgfpathlineto{\pgfqpoint{0.712111in}{1.193327in}}%
\pgfpathlineto{\pgfqpoint{0.720360in}{1.198312in}}%
\pgfpathlineto{\pgfqpoint{0.722022in}{1.199321in}}%
\pgfpathlineto{\pgfqpoint{0.728609in}{1.203360in}}%
\pgfpathlineto{\pgfqpoint{0.733774in}{1.206543in}}%
\pgfpathlineto{\pgfqpoint{0.736859in}{1.208462in}}%
\pgfpathlineto{\pgfqpoint{0.745108in}{1.213621in}}%
\pgfpathlineto{\pgfqpoint{0.745338in}{1.213766in}}%
\pgfpathlineto{\pgfqpoint{0.753357in}{1.218854in}}%
\pgfpathlineto{\pgfqpoint{0.756701in}{1.220988in}}%
\pgfpathlineto{\pgfqpoint{0.761606in}{1.224146in}}%
\pgfpathlineto{\pgfqpoint{0.767883in}{1.228210in}}%
\pgfpathlineto{\pgfqpoint{0.769856in}{1.229499in}}%
\pgfpathlineto{\pgfqpoint{0.778105in}{1.234918in}}%
\pgfpathlineto{\pgfqpoint{0.778882in}{1.235432in}}%
\pgfpathlineto{\pgfqpoint{0.786354in}{1.240413in}}%
\pgfpathlineto{\pgfqpoint{0.789697in}{1.242654in}}%
\pgfpathlineto{\pgfqpoint{0.794603in}{1.245973in}}%
\pgfpathlineto{\pgfqpoint{0.800340in}{1.249877in}}%
\pgfpathlineto{\pgfqpoint{0.802852in}{1.251601in}}%
\pgfpathlineto{\pgfqpoint{0.810816in}{1.257099in}}%
\pgfpathlineto{\pgfqpoint{0.811102in}{1.257298in}}%
\pgfpathlineto{\pgfqpoint{0.819351in}{1.263077in}}%
\pgfpathlineto{\pgfqpoint{0.821116in}{1.264321in}}%
\pgfpathlineto{\pgfqpoint{0.827600in}{1.268930in}}%
\pgfpathlineto{\pgfqpoint{0.831254in}{1.271543in}}%
\pgfpathlineto{\pgfqpoint{0.835849in}{1.274857in}}%
\pgfpathlineto{\pgfqpoint{0.841236in}{1.278766in}}%
\pgfpathlineto{\pgfqpoint{0.844099in}{1.280860in}}%
\pgfpathlineto{\pgfqpoint{0.851065in}{1.285988in}}%
\pgfpathlineto{\pgfqpoint{0.852348in}{1.286940in}}%
\pgfpathlineto{\pgfqpoint{0.860597in}{1.293100in}}%
\pgfpathlineto{\pgfqpoint{0.860743in}{1.293210in}}%
\pgfpathlineto{\pgfqpoint{0.868846in}{1.299349in}}%
\pgfpathlineto{\pgfqpoint{0.870268in}{1.300432in}}%
\pgfpathlineto{\pgfqpoint{0.877095in}{1.305679in}}%
\pgfpathlineto{\pgfqpoint{0.879650in}{1.307654in}}%
\pgfpathlineto{\pgfqpoint{0.885345in}{1.312094in}}%
\pgfpathlineto{\pgfqpoint{0.888892in}{1.314877in}}%
\pgfpathlineto{\pgfqpoint{0.893594in}{1.318594in}}%
\pgfpathlineto{\pgfqpoint{0.897999in}{1.322099in}}%
\pgfpathlineto{\pgfqpoint{0.901843in}{1.325182in}}%
\pgfpathlineto{\pgfqpoint{0.906972in}{1.329321in}}%
\pgfpathlineto{\pgfqpoint{0.910092in}{1.331860in}}%
\pgfpathlineto{\pgfqpoint{0.915813in}{1.336543in}}%
\pgfpathlineto{\pgfqpoint{0.918342in}{1.338629in}}%
\pgfpathlineto{\pgfqpoint{0.924526in}{1.343766in}}%
\pgfpathlineto{\pgfqpoint{0.926591in}{1.345494in}}%
\pgfpathlineto{\pgfqpoint{0.933112in}{1.350988in}}%
\pgfpathlineto{\pgfqpoint{0.934840in}{1.352455in}}%
\pgfpathlineto{\pgfqpoint{0.941574in}{1.358210in}}%
\pgfpathlineto{\pgfqpoint{0.943089in}{1.359515in}}%
\pgfpathlineto{\pgfqpoint{0.949914in}{1.365432in}}%
\pgfpathlineto{\pgfqpoint{0.951338in}{1.366676in}}%
\pgfpathlineto{\pgfqpoint{0.958135in}{1.372654in}}%
\pgfpathlineto{\pgfqpoint{0.959588in}{1.373942in}}%
\pgfpathlineto{\pgfqpoint{0.966238in}{1.379877in}}%
\pgfpathlineto{\pgfqpoint{0.967837in}{1.381315in}}%
\pgfpathlineto{\pgfqpoint{0.974225in}{1.387099in}}%
\pgfpathlineto{\pgfqpoint{0.976086in}{1.388796in}}%
\pgfpathlineto{\pgfqpoint{0.982100in}{1.394321in}}%
\pgfpathlineto{\pgfqpoint{0.984335in}{1.396390in}}%
\pgfpathlineto{\pgfqpoint{0.989864in}{1.401543in}}%
\pgfpathlineto{\pgfqpoint{0.992585in}{1.404097in}}%
\pgfpathlineto{\pgfqpoint{0.997520in}{1.408766in}}%
\pgfpathlineto{\pgfqpoint{1.000834in}{1.411922in}}%
\pgfpathlineto{\pgfqpoint{1.005071in}{1.415988in}}%
\pgfpathlineto{\pgfqpoint{1.009083in}{1.419866in}}%
\pgfpathlineto{\pgfqpoint{1.012518in}{1.423210in}}%
\pgfpathlineto{\pgfqpoint{1.017332in}{1.427931in}}%
\pgfpathlineto{\pgfqpoint{1.019865in}{1.430432in}}%
\pgfpathlineto{\pgfqpoint{1.025581in}{1.436119in}}%
\pgfpathlineto{\pgfqpoint{1.027114in}{1.437654in}}%
\pgfpathlineto{\pgfqpoint{1.033831in}{1.444432in}}%
\pgfpathlineto{\pgfqpoint{1.034268in}{1.444877in}}%
\pgfpathlineto{\pgfqpoint{1.041324in}{1.452099in}}%
\pgfpathlineto{\pgfqpoint{1.042080in}{1.452878in}}%
\pgfpathlineto{\pgfqpoint{1.048286in}{1.459321in}}%
\pgfpathlineto{\pgfqpoint{1.050329in}{1.461456in}}%
\pgfpathlineto{\pgfqpoint{1.055160in}{1.466543in}}%
\pgfpathlineto{\pgfqpoint{1.058578in}{1.470167in}}%
\pgfpathlineto{\pgfqpoint{1.061948in}{1.473766in}}%
\pgfpathlineto{\pgfqpoint{1.066828in}{1.479011in}}%
\pgfpathlineto{\pgfqpoint{1.068652in}{1.480988in}}%
\pgfpathlineto{\pgfqpoint{1.075077in}{1.487992in}}%
\pgfpathlineto{\pgfqpoint{1.075275in}{1.488210in}}%
\pgfpathlineto{\pgfqpoint{1.081808in}{1.495432in}}%
\pgfpathlineto{\pgfqpoint{1.083326in}{1.497122in}}%
\pgfpathlineto{\pgfqpoint{1.088261in}{1.502654in}}%
\pgfpathlineto{\pgfqpoint{1.091575in}{1.506394in}}%
\pgfpathlineto{\pgfqpoint{1.094639in}{1.509877in}}%
\pgfpathlineto{\pgfqpoint{1.099824in}{1.515809in}}%
\pgfpathlineto{\pgfqpoint{1.100943in}{1.517099in}}%
\pgfpathlineto{\pgfqpoint{1.107168in}{1.524321in}}%
\pgfpathlineto{\pgfqpoint{1.108074in}{1.525378in}}%
\pgfpathlineto{\pgfqpoint{1.113315in}{1.531543in}}%
\pgfpathlineto{\pgfqpoint{1.116323in}{1.535105in}}%
\pgfpathlineto{\pgfqpoint{1.119391in}{1.538766in}}%
\pgfpathlineto{\pgfqpoint{1.124572in}{1.544985in}}%
\pgfpathlineto{\pgfqpoint{1.125401in}{1.545988in}}%
\pgfpathlineto{\pgfqpoint{1.131334in}{1.553210in}}%
\pgfpathlineto{\pgfqpoint{1.132821in}{1.555032in}}%
\pgfpathlineto{\pgfqpoint{1.137195in}{1.560432in}}%
\pgfpathlineto{\pgfqpoint{1.141071in}{1.565246in}}%
\pgfpathlineto{\pgfqpoint{1.142994in}{1.567654in}}%
\pgfpathlineto{\pgfqpoint{1.148727in}{1.574877in}}%
\pgfpathlineto{\pgfqpoint{1.149320in}{1.575628in}}%
\pgfpathlineto{\pgfqpoint{1.154384in}{1.582099in}}%
\pgfpathlineto{\pgfqpoint{1.157569in}{1.586192in}}%
\pgfpathlineto{\pgfqpoint{1.159984in}{1.589321in}}%
\pgfpathlineto{\pgfqpoint{1.165524in}{1.596543in}}%
\pgfpathlineto{\pgfqpoint{1.165818in}{1.596929in}}%
\pgfpathlineto{\pgfqpoint{1.170989in}{1.603766in}}%
\pgfpathlineto{\pgfqpoint{1.174067in}{1.607859in}}%
\pgfpathlineto{\pgfqpoint{1.176400in}{1.610988in}}%
\pgfpathlineto{\pgfqpoint{1.181754in}{1.618210in}}%
\pgfpathlineto{\pgfqpoint{1.182317in}{1.618973in}}%
\pgfpathlineto{\pgfqpoint{1.187038in}{1.625432in}}%
\pgfpathlineto{\pgfqpoint{1.190566in}{1.630286in}}%
\pgfpathlineto{\pgfqpoint{1.192272in}{1.632654in}}%
\pgfpathlineto{\pgfqpoint{1.197447in}{1.639877in}}%
\pgfpathlineto{\pgfqpoint{1.198815in}{1.641796in}}%
\pgfpathlineto{\pgfqpoint{1.202561in}{1.647099in}}%
\pgfpathlineto{\pgfqpoint{1.207064in}{1.653509in}}%
\pgfpathlineto{\pgfqpoint{1.207630in}{1.654321in}}%
\pgfpathlineto{\pgfqpoint{1.212632in}{1.661543in}}%
\pgfpathlineto{\pgfqpoint{1.215314in}{1.665435in}}%
\pgfpathlineto{\pgfqpoint{1.217588in}{1.668766in}}%
\pgfpathlineto{\pgfqpoint{1.222493in}{1.675988in}}%
\pgfpathlineto{\pgfqpoint{1.223563in}{1.677571in}}%
\pgfpathlineto{\pgfqpoint{1.227340in}{1.683210in}}%
\pgfpathlineto{\pgfqpoint{1.231812in}{1.689924in}}%
\pgfpathlineto{\pgfqpoint{1.232148in}{1.690432in}}%
\pgfpathlineto{\pgfqpoint{1.236892in}{1.697654in}}%
\pgfpathlineto{\pgfqpoint{1.240061in}{1.702504in}}%
\pgfpathlineto{\pgfqpoint{1.241598in}{1.704877in}}%
\pgfpathlineto{\pgfqpoint{1.246252in}{1.712099in}}%
\pgfpathlineto{\pgfqpoint{1.248310in}{1.715310in}}%
\pgfpathlineto{\pgfqpoint{1.250859in}{1.719321in}}%
\pgfpathlineto{\pgfqpoint{1.255425in}{1.726543in}}%
\pgfpathlineto{\pgfqpoint{1.256560in}{1.728347in}}%
\pgfpathlineto{\pgfqpoint{1.259937in}{1.733766in}}%
\pgfpathlineto{\pgfqpoint{1.264417in}{1.740988in}}%
\pgfpathlineto{\pgfqpoint{1.264809in}{1.741622in}}%
\pgfpathlineto{\pgfqpoint{1.268839in}{1.748210in}}%
\pgfpathlineto{\pgfqpoint{1.273058in}{1.755141in}}%
\pgfpathlineto{\pgfqpoint{1.273234in}{1.755432in}}%
\pgfpathlineto{\pgfqpoint{1.277570in}{1.762654in}}%
\pgfpathlineto{\pgfqpoint{1.281307in}{1.768909in}}%
\pgfpathlineto{\pgfqpoint{1.281880in}{1.769877in}}%
\pgfpathlineto{\pgfqpoint{1.286136in}{1.777099in}}%
\pgfpathlineto{\pgfqpoint{1.289557in}{1.782930in}}%
\pgfpathlineto{\pgfqpoint{1.290365in}{1.784321in}}%
\pgfpathlineto{\pgfqpoint{1.294543in}{1.791543in}}%
\pgfpathlineto{\pgfqpoint{1.297806in}{1.797210in}}%
\pgfpathlineto{\pgfqpoint{1.298693in}{1.798766in}}%
\pgfpathlineto{\pgfqpoint{1.302794in}{1.805988in}}%
\pgfpathlineto{\pgfqpoint{1.306055in}{1.811755in}}%
\pgfpathlineto{\pgfqpoint{1.306870in}{1.813210in}}%
\pgfpathlineto{\pgfqpoint{1.310897in}{1.820432in}}%
\pgfpathlineto{\pgfqpoint{1.314304in}{1.826571in}}%
\pgfpathlineto{\pgfqpoint{1.314900in}{1.827654in}}%
\pgfpathlineto{\pgfqpoint{1.318854in}{1.834877in}}%
\pgfpathlineto{\pgfqpoint{1.322553in}{1.841664in}}%
\pgfpathlineto{\pgfqpoint{1.322788in}{1.842099in}}%
\pgfpathlineto{\pgfqpoint{1.326671in}{1.849321in}}%
\pgfpathlineto{\pgfqpoint{1.330537in}{1.856543in}}%
\pgfpathlineto{\pgfqpoint{1.330803in}{1.857042in}}%
\pgfpathlineto{\pgfqpoint{1.334352in}{1.863766in}}%
\pgfpathlineto{\pgfqpoint{1.338149in}{1.870988in}}%
\pgfpathlineto{\pgfqpoint{1.339052in}{1.872711in}}%
\pgfpathlineto{\pgfqpoint{1.341903in}{1.878210in}}%
\pgfpathlineto{\pgfqpoint{1.345633in}{1.885432in}}%
\pgfpathlineto{\pgfqpoint{1.347301in}{1.888675in}}%
\pgfpathlineto{\pgfqpoint{1.349327in}{1.892654in}}%
\pgfpathlineto{\pgfqpoint{1.352991in}{1.899877in}}%
\pgfpathlineto{\pgfqpoint{1.347301in}{1.899877in}}%
\pgfpathlineto{\pgfqpoint{1.339052in}{1.899877in}}%
\pgfpathlineto{\pgfqpoint{1.330803in}{1.899877in}}%
\pgfpathlineto{\pgfqpoint{1.322553in}{1.899877in}}%
\pgfpathlineto{\pgfqpoint{1.314304in}{1.899877in}}%
\pgfpathlineto{\pgfqpoint{1.306055in}{1.899877in}}%
\pgfpathlineto{\pgfqpoint{1.297806in}{1.899877in}}%
\pgfpathlineto{\pgfqpoint{1.289557in}{1.899877in}}%
\pgfpathlineto{\pgfqpoint{1.281307in}{1.899877in}}%
\pgfpathlineto{\pgfqpoint{1.273058in}{1.899877in}}%
\pgfpathlineto{\pgfqpoint{1.264809in}{1.899877in}}%
\pgfpathlineto{\pgfqpoint{1.256560in}{1.899877in}}%
\pgfpathlineto{\pgfqpoint{1.248310in}{1.899877in}}%
\pgfpathlineto{\pgfqpoint{1.240061in}{1.899877in}}%
\pgfpathlineto{\pgfqpoint{1.231812in}{1.899877in}}%
\pgfpathlineto{\pgfqpoint{1.223563in}{1.899877in}}%
\pgfpathlineto{\pgfqpoint{1.215314in}{1.899877in}}%
\pgfpathlineto{\pgfqpoint{1.207064in}{1.899877in}}%
\pgfpathlineto{\pgfqpoint{1.198815in}{1.899877in}}%
\pgfpathlineto{\pgfqpoint{1.190566in}{1.899877in}}%
\pgfpathlineto{\pgfqpoint{1.182317in}{1.899877in}}%
\pgfpathlineto{\pgfqpoint{1.174067in}{1.899877in}}%
\pgfpathlineto{\pgfqpoint{1.165818in}{1.899877in}}%
\pgfpathlineto{\pgfqpoint{1.157569in}{1.899877in}}%
\pgfpathlineto{\pgfqpoint{1.149320in}{1.899877in}}%
\pgfpathlineto{\pgfqpoint{1.141071in}{1.899877in}}%
\pgfpathlineto{\pgfqpoint{1.132821in}{1.899877in}}%
\pgfpathlineto{\pgfqpoint{1.124572in}{1.899877in}}%
\pgfpathlineto{\pgfqpoint{1.116323in}{1.899877in}}%
\pgfpathlineto{\pgfqpoint{1.108074in}{1.899877in}}%
\pgfpathlineto{\pgfqpoint{1.099824in}{1.899877in}}%
\pgfpathlineto{\pgfqpoint{1.091575in}{1.899877in}}%
\pgfpathlineto{\pgfqpoint{1.083326in}{1.899877in}}%
\pgfpathlineto{\pgfqpoint{1.075077in}{1.899877in}}%
\pgfpathlineto{\pgfqpoint{1.066828in}{1.899877in}}%
\pgfpathlineto{\pgfqpoint{1.058578in}{1.899877in}}%
\pgfpathlineto{\pgfqpoint{1.050329in}{1.899877in}}%
\pgfpathlineto{\pgfqpoint{1.042080in}{1.899877in}}%
\pgfpathlineto{\pgfqpoint{1.033831in}{1.899877in}}%
\pgfpathlineto{\pgfqpoint{1.025581in}{1.899877in}}%
\pgfpathlineto{\pgfqpoint{1.017332in}{1.899877in}}%
\pgfpathlineto{\pgfqpoint{1.009083in}{1.899877in}}%
\pgfpathlineto{\pgfqpoint{1.000834in}{1.899877in}}%
\pgfpathlineto{\pgfqpoint{0.992585in}{1.899877in}}%
\pgfpathlineto{\pgfqpoint{0.984335in}{1.899877in}}%
\pgfpathlineto{\pgfqpoint{0.976086in}{1.899877in}}%
\pgfpathlineto{\pgfqpoint{0.967837in}{1.899877in}}%
\pgfpathlineto{\pgfqpoint{0.959588in}{1.899877in}}%
\pgfpathlineto{\pgfqpoint{0.951338in}{1.899877in}}%
\pgfpathlineto{\pgfqpoint{0.943089in}{1.899877in}}%
\pgfpathlineto{\pgfqpoint{0.937864in}{1.899877in}}%
\pgfpathlineto{\pgfqpoint{0.934840in}{1.896149in}}%
\pgfpathlineto{\pgfqpoint{0.931979in}{1.892654in}}%
\pgfpathlineto{\pgfqpoint{0.926591in}{1.886083in}}%
\pgfpathlineto{\pgfqpoint{0.926053in}{1.885432in}}%
\pgfpathlineto{\pgfqpoint{0.920065in}{1.878210in}}%
\pgfpathlineto{\pgfqpoint{0.918342in}{1.876134in}}%
\pgfpathlineto{\pgfqpoint{0.914029in}{1.870988in}}%
\pgfpathlineto{\pgfqpoint{0.910092in}{1.866297in}}%
\pgfpathlineto{\pgfqpoint{0.907949in}{1.863766in}}%
\pgfpathlineto{\pgfqpoint{0.901843in}{1.856567in}}%
\pgfpathlineto{\pgfqpoint{0.901823in}{1.856543in}}%
\pgfpathlineto{\pgfqpoint{0.895631in}{1.849321in}}%
\pgfpathlineto{\pgfqpoint{0.893594in}{1.846950in}}%
\pgfpathlineto{\pgfqpoint{0.889391in}{1.842099in}}%
\pgfpathlineto{\pgfqpoint{0.885345in}{1.837437in}}%
\pgfpathlineto{\pgfqpoint{0.883103in}{1.834877in}}%
\pgfpathlineto{\pgfqpoint{0.877095in}{1.828027in}}%
\pgfpathlineto{\pgfqpoint{0.876766in}{1.827654in}}%
\pgfpathlineto{\pgfqpoint{0.870363in}{1.820432in}}%
\pgfpathlineto{\pgfqpoint{0.868846in}{1.818725in}}%
\pgfpathlineto{\pgfqpoint{0.863906in}{1.813210in}}%
\pgfpathlineto{\pgfqpoint{0.860597in}{1.809524in}}%
\pgfpathlineto{\pgfqpoint{0.857396in}{1.805988in}}%
\pgfpathlineto{\pgfqpoint{0.852348in}{1.800422in}}%
\pgfpathlineto{\pgfqpoint{0.850833in}{1.798766in}}%
\pgfpathlineto{\pgfqpoint{0.844213in}{1.791543in}}%
\pgfpathlineto{\pgfqpoint{0.844099in}{1.791418in}}%
\pgfpathlineto{\pgfqpoint{0.837524in}{1.784321in}}%
\pgfpathlineto{\pgfqpoint{0.835849in}{1.782518in}}%
\pgfpathlineto{\pgfqpoint{0.830777in}{1.777099in}}%
\pgfpathlineto{\pgfqpoint{0.827600in}{1.773713in}}%
\pgfpathlineto{\pgfqpoint{0.823972in}{1.769877in}}%
\pgfpathlineto{\pgfqpoint{0.819351in}{1.765002in}}%
\pgfpathlineto{\pgfqpoint{0.817108in}{1.762654in}}%
\pgfpathlineto{\pgfqpoint{0.811102in}{1.756383in}}%
\pgfpathlineto{\pgfqpoint{0.810183in}{1.755432in}}%
\pgfpathlineto{\pgfqpoint{0.803193in}{1.748210in}}%
\pgfpathlineto{\pgfqpoint{0.802852in}{1.747859in}}%
\pgfpathlineto{\pgfqpoint{0.796131in}{1.740988in}}%
\pgfpathlineto{\pgfqpoint{0.794603in}{1.739429in}}%
\pgfpathlineto{\pgfqpoint{0.789006in}{1.733766in}}%
\pgfpathlineto{\pgfqpoint{0.786354in}{1.731089in}}%
\pgfpathlineto{\pgfqpoint{0.781815in}{1.726543in}}%
\pgfpathlineto{\pgfqpoint{0.778105in}{1.722837in}}%
\pgfpathlineto{\pgfqpoint{0.774559in}{1.719321in}}%
\pgfpathlineto{\pgfqpoint{0.769856in}{1.714671in}}%
\pgfpathlineto{\pgfqpoint{0.767235in}{1.712099in}}%
\pgfpathlineto{\pgfqpoint{0.761606in}{1.706590in}}%
\pgfpathlineto{\pgfqpoint{0.759843in}{1.704877in}}%
\pgfpathlineto{\pgfqpoint{0.753357in}{1.698593in}}%
\pgfpathlineto{\pgfqpoint{0.752382in}{1.697654in}}%
\pgfpathlineto{\pgfqpoint{0.745108in}{1.690677in}}%
\pgfpathlineto{\pgfqpoint{0.744850in}{1.690432in}}%
\pgfpathlineto{\pgfqpoint{0.737244in}{1.683210in}}%
\pgfpathlineto{\pgfqpoint{0.736859in}{1.682845in}}%
\pgfpathlineto{\pgfqpoint{0.729564in}{1.675988in}}%
\pgfpathlineto{\pgfqpoint{0.728609in}{1.675093in}}%
\pgfpathlineto{\pgfqpoint{0.721812in}{1.668766in}}%
\pgfpathlineto{\pgfqpoint{0.720360in}{1.667419in}}%
\pgfpathlineto{\pgfqpoint{0.713986in}{1.661543in}}%
\pgfpathlineto{\pgfqpoint{0.712111in}{1.659821in}}%
\pgfpathlineto{\pgfqpoint{0.706086in}{1.654321in}}%
\pgfpathlineto{\pgfqpoint{0.703862in}{1.652298in}}%
\pgfpathlineto{\pgfqpoint{0.698111in}{1.647099in}}%
\pgfpathlineto{\pgfqpoint{0.695612in}{1.644849in}}%
\pgfpathlineto{\pgfqpoint{0.690059in}{1.639877in}}%
\pgfpathlineto{\pgfqpoint{0.687363in}{1.637472in}}%
\pgfpathlineto{\pgfqpoint{0.681929in}{1.632654in}}%
\pgfpathlineto{\pgfqpoint{0.679114in}{1.630168in}}%
\pgfpathlineto{\pgfqpoint{0.673719in}{1.625432in}}%
\pgfpathlineto{\pgfqpoint{0.670865in}{1.622936in}}%
\pgfpathlineto{\pgfqpoint{0.665428in}{1.618210in}}%
\pgfpathlineto{\pgfqpoint{0.662616in}{1.615775in}}%
\pgfpathlineto{\pgfqpoint{0.657054in}{1.610988in}}%
\pgfpathlineto{\pgfqpoint{0.654366in}{1.608684in}}%
\pgfpathlineto{\pgfqpoint{0.648596in}{1.603766in}}%
\pgfpathlineto{\pgfqpoint{0.646117in}{1.601662in}}%
\pgfpathlineto{\pgfqpoint{0.640050in}{1.596543in}}%
\pgfpathlineto{\pgfqpoint{0.637868in}{1.594710in}}%
\pgfpathlineto{\pgfqpoint{0.631417in}{1.589321in}}%
\pgfpathlineto{\pgfqpoint{0.629619in}{1.587825in}}%
\pgfpathlineto{\pgfqpoint{0.622694in}{1.582099in}}%
\pgfpathlineto{\pgfqpoint{0.621369in}{1.581009in}}%
\pgfpathlineto{\pgfqpoint{0.613878in}{1.574877in}}%
\pgfpathlineto{\pgfqpoint{0.613120in}{1.574259in}}%
\pgfpathlineto{\pgfqpoint{0.604969in}{1.567654in}}%
\pgfpathlineto{\pgfqpoint{0.604871in}{1.567576in}}%
\pgfpathlineto{\pgfqpoint{0.596622in}{1.560962in}}%
\pgfpathlineto{\pgfqpoint{0.595957in}{1.560432in}}%
\pgfpathlineto{\pgfqpoint{0.588373in}{1.554413in}}%
\pgfpathlineto{\pgfqpoint{0.586848in}{1.553210in}}%
\pgfpathlineto{\pgfqpoint{0.580123in}{1.547929in}}%
\pgfpathlineto{\pgfqpoint{0.577639in}{1.545988in}}%
\pgfpathlineto{\pgfqpoint{0.571874in}{1.541506in}}%
\pgfpathlineto{\pgfqpoint{0.568330in}{1.538766in}}%
\pgfpathlineto{\pgfqpoint{0.563625in}{1.535145in}}%
\pgfpathlineto{\pgfqpoint{0.558920in}{1.531543in}}%
\pgfpathlineto{\pgfqpoint{0.555376in}{1.528844in}}%
\pgfpathlineto{\pgfqpoint{0.549408in}{1.524321in}}%
\pgfpathlineto{\pgfqpoint{0.547126in}{1.522601in}}%
\pgfpathlineto{\pgfqpoint{0.539792in}{1.517099in}}%
\pgfpathlineto{\pgfqpoint{0.538877in}{1.516416in}}%
\pgfpathlineto{\pgfqpoint{0.530628in}{1.510292in}}%
\pgfpathlineto{\pgfqpoint{0.530066in}{1.509877in}}%
\pgfpathlineto{\pgfqpoint{0.522379in}{1.504228in}}%
\pgfpathlineto{\pgfqpoint{0.520227in}{1.502654in}}%
\pgfpathlineto{\pgfqpoint{0.514130in}{1.498220in}}%
\pgfpathlineto{\pgfqpoint{0.510280in}{1.495432in}}%
\pgfpathlineto{\pgfqpoint{0.505880in}{1.492265in}}%
\pgfpathlineto{\pgfqpoint{0.500224in}{1.488210in}}%
\pgfpathlineto{\pgfqpoint{0.497631in}{1.486362in}}%
\pgfpathlineto{\pgfqpoint{0.490059in}{1.480988in}}%
\pgfpathlineto{\pgfqpoint{0.489382in}{1.480510in}}%
\pgfpathlineto{\pgfqpoint{0.481133in}{1.474717in}}%
\pgfpathlineto{\pgfqpoint{0.479773in}{1.473766in}}%
\pgfpathlineto{\pgfqpoint{0.472883in}{1.468977in}}%
\pgfpathlineto{\pgfqpoint{0.469369in}{1.466543in}}%
\pgfpathlineto{\pgfqpoint{0.464634in}{1.463286in}}%
\pgfpathlineto{\pgfqpoint{0.458852in}{1.459321in}}%
\pgfpathlineto{\pgfqpoint{0.456385in}{1.457641in}}%
\pgfpathlineto{\pgfqpoint{0.448221in}{1.452099in}}%
\pgfpathlineto{\pgfqpoint{0.448136in}{1.452042in}}%
\pgfpathlineto{\pgfqpoint{0.439887in}{1.446503in}}%
\pgfpathlineto{\pgfqpoint{0.437456in}{1.444877in}}%
\pgfpathlineto{\pgfqpoint{0.431637in}{1.441009in}}%
\pgfpathlineto{\pgfqpoint{0.426575in}{1.437654in}}%
\pgfpathlineto{\pgfqpoint{0.423388in}{1.435557in}}%
\pgfpathlineto{\pgfqpoint{0.415577in}{1.430432in}}%
\pgfpathlineto{\pgfqpoint{0.415139in}{1.430147in}}%
\pgfpathlineto{\pgfqpoint{0.415139in}{1.423210in}}%
\pgfpathlineto{\pgfqpoint{0.415139in}{1.415988in}}%
\pgfpathlineto{\pgfqpoint{0.415139in}{1.408766in}}%
\pgfpathlineto{\pgfqpoint{0.415139in}{1.401543in}}%
\pgfpathlineto{\pgfqpoint{0.415139in}{1.394321in}}%
\pgfpathlineto{\pgfqpoint{0.415139in}{1.387099in}}%
\pgfpathlineto{\pgfqpoint{0.415139in}{1.379877in}}%
\pgfpathlineto{\pgfqpoint{0.415139in}{1.372654in}}%
\pgfpathlineto{\pgfqpoint{0.415139in}{1.365432in}}%
\pgfpathlineto{\pgfqpoint{0.415139in}{1.358210in}}%
\pgfpathlineto{\pgfqpoint{0.415139in}{1.350988in}}%
\pgfpathlineto{\pgfqpoint{0.415139in}{1.343766in}}%
\pgfpathlineto{\pgfqpoint{0.415139in}{1.336543in}}%
\pgfpathlineto{\pgfqpoint{0.415139in}{1.329321in}}%
\pgfpathlineto{\pgfqpoint{0.415139in}{1.322099in}}%
\pgfpathlineto{\pgfqpoint{0.415139in}{1.314877in}}%
\pgfpathlineto{\pgfqpoint{0.415139in}{1.307654in}}%
\pgfpathlineto{\pgfqpoint{0.415139in}{1.300432in}}%
\pgfpathlineto{\pgfqpoint{0.415139in}{1.293210in}}%
\pgfpathlineto{\pgfqpoint{0.415139in}{1.285988in}}%
\pgfpathlineto{\pgfqpoint{0.415139in}{1.278766in}}%
\pgfpathlineto{\pgfqpoint{0.415139in}{1.271543in}}%
\pgfpathlineto{\pgfqpoint{0.415139in}{1.264321in}}%
\pgfpathlineto{\pgfqpoint{0.415139in}{1.257099in}}%
\pgfpathlineto{\pgfqpoint{0.415139in}{1.249877in}}%
\pgfpathlineto{\pgfqpoint{0.415139in}{1.242654in}}%
\pgfpathlineto{\pgfqpoint{0.415139in}{1.235432in}}%
\pgfpathlineto{\pgfqpoint{0.415139in}{1.228210in}}%
\pgfpathlineto{\pgfqpoint{0.415139in}{1.220988in}}%
\pgfpathlineto{\pgfqpoint{0.415139in}{1.213766in}}%
\pgfpathlineto{\pgfqpoint{0.415139in}{1.206543in}}%
\pgfpathlineto{\pgfqpoint{0.415139in}{1.199321in}}%
\pgfpathlineto{\pgfqpoint{0.415139in}{1.192099in}}%
\pgfpathlineto{\pgfqpoint{0.415139in}{1.184877in}}%
\pgfpathlineto{\pgfqpoint{0.415139in}{1.177654in}}%
\pgfpathlineto{\pgfqpoint{0.415139in}{1.170432in}}%
\pgfpathlineto{\pgfqpoint{0.415139in}{1.163210in}}%
\pgfpathlineto{\pgfqpoint{0.415139in}{1.155988in}}%
\pgfpathlineto{\pgfqpoint{0.415139in}{1.148766in}}%
\pgfpathlineto{\pgfqpoint{0.415139in}{1.141543in}}%
\pgfpathlineto{\pgfqpoint{0.415139in}{1.134321in}}%
\pgfpathlineto{\pgfqpoint{0.415139in}{1.127099in}}%
\pgfpathlineto{\pgfqpoint{0.415139in}{1.119877in}}%
\pgfpathlineto{\pgfqpoint{0.415139in}{1.112654in}}%
\pgfpathlineto{\pgfqpoint{0.415139in}{1.105432in}}%
\pgfpathlineto{\pgfqpoint{0.415139in}{1.098210in}}%
\pgfpathlineto{\pgfqpoint{0.415139in}{1.090988in}}%
\pgfpathlineto{\pgfqpoint{0.415139in}{1.083766in}}%
\pgfpathlineto{\pgfqpoint{0.415139in}{1.076543in}}%
\pgfpathlineto{\pgfqpoint{0.415139in}{1.069321in}}%
\pgfpathlineto{\pgfqpoint{0.415139in}{1.062099in}}%
\pgfpathlineto{\pgfqpoint{0.415139in}{1.054877in}}%
\pgfpathlineto{\pgfqpoint{0.415139in}{1.047654in}}%
\pgfpathlineto{\pgfqpoint{0.415139in}{1.046309in}}%
\pgfpathlineto{\pgfqpoint{0.418442in}{1.047654in}}%
\pgfpathclose%
\pgfusepath{fill}%
\end{pgfscope}%
\begin{pgfscope}%
\pgfpathrectangle{\pgfqpoint{0.415139in}{0.462654in}}{\pgfqpoint{1.641596in}{1.437222in}}%
\pgfusepath{clip}%
\pgfsetbuttcap%
\pgfsetroundjoin%
\definecolor{currentfill}{rgb}{0.940408,0.976932,0.858593}%
\pgfsetfillcolor{currentfill}%
\pgfsetlinewidth{0.000000pt}%
\definecolor{currentstroke}{rgb}{0.000000,0.000000,0.000000}%
\pgfsetstrokecolor{currentstroke}%
\pgfsetdash{}{0pt}%
\pgfpathmoveto{\pgfqpoint{0.419871in}{0.809321in}}%
\pgfpathlineto{\pgfqpoint{0.423388in}{0.810234in}}%
\pgfpathlineto{\pgfqpoint{0.431637in}{0.812384in}}%
\pgfpathlineto{\pgfqpoint{0.439887in}{0.814549in}}%
\pgfpathlineto{\pgfqpoint{0.447443in}{0.816543in}}%
\pgfpathlineto{\pgfqpoint{0.448136in}{0.816729in}}%
\pgfpathlineto{\pgfqpoint{0.456385in}{0.818943in}}%
\pgfpathlineto{\pgfqpoint{0.464634in}{0.821172in}}%
\pgfpathlineto{\pgfqpoint{0.472883in}{0.823416in}}%
\pgfpathlineto{\pgfqpoint{0.474163in}{0.823766in}}%
\pgfpathlineto{\pgfqpoint{0.481133in}{0.825694in}}%
\pgfpathlineto{\pgfqpoint{0.489382in}{0.827990in}}%
\pgfpathlineto{\pgfqpoint{0.497631in}{0.830302in}}%
\pgfpathlineto{\pgfqpoint{0.500063in}{0.830988in}}%
\pgfpathlineto{\pgfqpoint{0.505880in}{0.832647in}}%
\pgfpathlineto{\pgfqpoint{0.514130in}{0.835014in}}%
\pgfpathlineto{\pgfqpoint{0.522379in}{0.837398in}}%
\pgfpathlineto{\pgfqpoint{0.525175in}{0.838210in}}%
\pgfpathlineto{\pgfqpoint{0.530628in}{0.839814in}}%
\pgfpathlineto{\pgfqpoint{0.538877in}{0.842255in}}%
\pgfpathlineto{\pgfqpoint{0.547126in}{0.844713in}}%
\pgfpathlineto{\pgfqpoint{0.549526in}{0.845432in}}%
\pgfpathlineto{\pgfqpoint{0.555376in}{0.847206in}}%
\pgfpathlineto{\pgfqpoint{0.563625in}{0.849725in}}%
\pgfpathlineto{\pgfqpoint{0.571874in}{0.852261in}}%
\pgfpathlineto{\pgfqpoint{0.573146in}{0.852654in}}%
\pgfpathlineto{\pgfqpoint{0.580123in}{0.854838in}}%
\pgfpathlineto{\pgfqpoint{0.588373in}{0.857437in}}%
\pgfpathlineto{\pgfqpoint{0.596060in}{0.859877in}}%
\pgfpathlineto{\pgfqpoint{0.596622in}{0.860057in}}%
\pgfpathlineto{\pgfqpoint{0.604871in}{0.862721in}}%
\pgfpathlineto{\pgfqpoint{0.613120in}{0.865405in}}%
\pgfpathlineto{\pgfqpoint{0.618292in}{0.867099in}}%
\pgfpathlineto{\pgfqpoint{0.621369in}{0.868119in}}%
\pgfpathlineto{\pgfqpoint{0.629619in}{0.870869in}}%
\pgfpathlineto{\pgfqpoint{0.637868in}{0.873641in}}%
\pgfpathlineto{\pgfqpoint{0.639881in}{0.874321in}}%
\pgfpathlineto{\pgfqpoint{0.646117in}{0.876453in}}%
\pgfpathlineto{\pgfqpoint{0.654366in}{0.879293in}}%
\pgfpathlineto{\pgfqpoint{0.660856in}{0.881543in}}%
\pgfpathlineto{\pgfqpoint{0.662616in}{0.882160in}}%
\pgfpathlineto{\pgfqpoint{0.670865in}{0.885070in}}%
\pgfpathlineto{\pgfqpoint{0.679114in}{0.888002in}}%
\pgfpathlineto{\pgfqpoint{0.681251in}{0.888766in}}%
\pgfpathlineto{\pgfqpoint{0.687363in}{0.890976in}}%
\pgfpathlineto{\pgfqpoint{0.695612in}{0.893979in}}%
\pgfpathlineto{\pgfqpoint{0.701095in}{0.895988in}}%
\pgfpathlineto{\pgfqpoint{0.703862in}{0.897013in}}%
\pgfpathlineto{\pgfqpoint{0.712111in}{0.900089in}}%
\pgfpathlineto{\pgfqpoint{0.720360in}{0.903188in}}%
\pgfpathlineto{\pgfqpoint{0.720417in}{0.903210in}}%
\pgfpathlineto{\pgfqpoint{0.728609in}{0.906339in}}%
\pgfpathlineto{\pgfqpoint{0.736859in}{0.909514in}}%
\pgfpathlineto{\pgfqpoint{0.739230in}{0.910432in}}%
\pgfpathlineto{\pgfqpoint{0.745108in}{0.912734in}}%
\pgfpathlineto{\pgfqpoint{0.753357in}{0.915988in}}%
\pgfpathlineto{\pgfqpoint{0.757555in}{0.917654in}}%
\pgfpathlineto{\pgfqpoint{0.761606in}{0.919281in}}%
\pgfpathlineto{\pgfqpoint{0.769856in}{0.922616in}}%
\pgfpathlineto{\pgfqpoint{0.775406in}{0.924877in}}%
\pgfpathlineto{\pgfqpoint{0.778105in}{0.925988in}}%
\pgfpathlineto{\pgfqpoint{0.786354in}{0.929407in}}%
\pgfpathlineto{\pgfqpoint{0.792799in}{0.932099in}}%
\pgfpathlineto{\pgfqpoint{0.794603in}{0.932861in}}%
\pgfpathlineto{\pgfqpoint{0.802852in}{0.936367in}}%
\pgfpathlineto{\pgfqpoint{0.809747in}{0.939321in}}%
\pgfpathlineto{\pgfqpoint{0.811102in}{0.939908in}}%
\pgfpathlineto{\pgfqpoint{0.819351in}{0.943504in}}%
\pgfpathlineto{\pgfqpoint{0.826267in}{0.946543in}}%
\pgfpathlineto{\pgfqpoint{0.827600in}{0.947136in}}%
\pgfpathlineto{\pgfqpoint{0.835849in}{0.950824in}}%
\pgfpathlineto{\pgfqpoint{0.842375in}{0.953766in}}%
\pgfpathlineto{\pgfqpoint{0.844099in}{0.954551in}}%
\pgfpathlineto{\pgfqpoint{0.852348in}{0.958334in}}%
\pgfpathlineto{\pgfqpoint{0.858088in}{0.960988in}}%
\pgfpathlineto{\pgfqpoint{0.860597in}{0.962160in}}%
\pgfpathlineto{\pgfqpoint{0.868846in}{0.966041in}}%
\pgfpathlineto{\pgfqpoint{0.873420in}{0.968210in}}%
\pgfpathlineto{\pgfqpoint{0.877095in}{0.969971in}}%
\pgfpathlineto{\pgfqpoint{0.885345in}{0.973953in}}%
\pgfpathlineto{\pgfqpoint{0.888388in}{0.975432in}}%
\pgfpathlineto{\pgfqpoint{0.893594in}{0.977990in}}%
\pgfpathlineto{\pgfqpoint{0.901843in}{0.982075in}}%
\pgfpathlineto{\pgfqpoint{0.903005in}{0.982654in}}%
\pgfpathlineto{\pgfqpoint{0.910092in}{0.986225in}}%
\pgfpathlineto{\pgfqpoint{0.917279in}{0.989877in}}%
\pgfpathlineto{\pgfqpoint{0.918342in}{0.990422in}}%
\pgfpathlineto{\pgfqpoint{0.926591in}{0.994686in}}%
\pgfpathlineto{\pgfqpoint{0.931222in}{0.997099in}}%
\pgfpathlineto{\pgfqpoint{0.934840in}{0.999004in}}%
\pgfpathlineto{\pgfqpoint{0.943089in}{1.003382in}}%
\pgfpathlineto{\pgfqpoint{0.944846in}{1.004321in}}%
\pgfpathlineto{\pgfqpoint{0.951338in}{1.007828in}}%
\pgfpathlineto{\pgfqpoint{0.958159in}{1.011543in}}%
\pgfpathlineto{\pgfqpoint{0.959588in}{1.012330in}}%
\pgfpathlineto{\pgfqpoint{0.967837in}{1.016904in}}%
\pgfpathlineto{\pgfqpoint{0.971167in}{1.018766in}}%
\pgfpathlineto{\pgfqpoint{0.976086in}{1.021544in}}%
\pgfpathlineto{\pgfqpoint{0.983887in}{1.025988in}}%
\pgfpathlineto{\pgfqpoint{0.984335in}{1.026246in}}%
\pgfpathlineto{\pgfqpoint{0.992585in}{1.031028in}}%
\pgfpathlineto{\pgfqpoint{0.996315in}{1.033210in}}%
\pgfpathlineto{\pgfqpoint{1.000834in}{1.035879in}}%
\pgfpathlineto{\pgfqpoint{1.008474in}{1.040432in}}%
\pgfpathlineto{\pgfqpoint{1.009083in}{1.040799in}}%
\pgfpathlineto{\pgfqpoint{1.017332in}{1.045802in}}%
\pgfpathlineto{\pgfqpoint{1.020362in}{1.047654in}}%
\pgfpathlineto{\pgfqpoint{1.025581in}{1.050879in}}%
\pgfpathlineto{\pgfqpoint{1.031997in}{1.054877in}}%
\pgfpathlineto{\pgfqpoint{1.033831in}{1.056031in}}%
\pgfpathlineto{\pgfqpoint{1.042080in}{1.061264in}}%
\pgfpathlineto{\pgfqpoint{1.043385in}{1.062099in}}%
\pgfpathlineto{\pgfqpoint{1.050329in}{1.066583in}}%
\pgfpathlineto{\pgfqpoint{1.054533in}{1.069321in}}%
\pgfpathlineto{\pgfqpoint{1.058578in}{1.071982in}}%
\pgfpathlineto{\pgfqpoint{1.065455in}{1.076543in}}%
\pgfpathlineto{\pgfqpoint{1.066828in}{1.077463in}}%
\pgfpathlineto{\pgfqpoint{1.075077in}{1.083032in}}%
\pgfpathlineto{\pgfqpoint{1.076155in}{1.083766in}}%
\pgfpathlineto{\pgfqpoint{1.083326in}{1.088694in}}%
\pgfpathlineto{\pgfqpoint{1.086637in}{1.090988in}}%
\pgfpathlineto{\pgfqpoint{1.091575in}{1.094443in}}%
\pgfpathlineto{\pgfqpoint{1.096914in}{1.098210in}}%
\pgfpathlineto{\pgfqpoint{1.099824in}{1.100284in}}%
\pgfpathlineto{\pgfqpoint{1.106992in}{1.105432in}}%
\pgfpathlineto{\pgfqpoint{1.108074in}{1.106217in}}%
\pgfpathlineto{\pgfqpoint{1.116323in}{1.112249in}}%
\pgfpathlineto{\pgfqpoint{1.116873in}{1.112654in}}%
\pgfpathlineto{\pgfqpoint{1.124572in}{1.118385in}}%
\pgfpathlineto{\pgfqpoint{1.126560in}{1.119877in}}%
\pgfpathlineto{\pgfqpoint{1.132821in}{1.124623in}}%
\pgfpathlineto{\pgfqpoint{1.136061in}{1.127099in}}%
\pgfpathlineto{\pgfqpoint{1.141071in}{1.130965in}}%
\pgfpathlineto{\pgfqpoint{1.145382in}{1.134321in}}%
\pgfpathlineto{\pgfqpoint{1.149320in}{1.137416in}}%
\pgfpathlineto{\pgfqpoint{1.154527in}{1.141543in}}%
\pgfpathlineto{\pgfqpoint{1.157569in}{1.143979in}}%
\pgfpathlineto{\pgfqpoint{1.163499in}{1.148766in}}%
\pgfpathlineto{\pgfqpoint{1.165818in}{1.150656in}}%
\pgfpathlineto{\pgfqpoint{1.172303in}{1.155988in}}%
\pgfpathlineto{\pgfqpoint{1.174067in}{1.157453in}}%
\pgfpathlineto{\pgfqpoint{1.180944in}{1.163210in}}%
\pgfpathlineto{\pgfqpoint{1.182317in}{1.164371in}}%
\pgfpathlineto{\pgfqpoint{1.189424in}{1.170432in}}%
\pgfpathlineto{\pgfqpoint{1.190566in}{1.171416in}}%
\pgfpathlineto{\pgfqpoint{1.197749in}{1.177654in}}%
\pgfpathlineto{\pgfqpoint{1.198815in}{1.178590in}}%
\pgfpathlineto{\pgfqpoint{1.205922in}{1.184877in}}%
\pgfpathlineto{\pgfqpoint{1.207064in}{1.185897in}}%
\pgfpathlineto{\pgfqpoint{1.213947in}{1.192099in}}%
\pgfpathlineto{\pgfqpoint{1.215314in}{1.193342in}}%
\pgfpathlineto{\pgfqpoint{1.221830in}{1.199321in}}%
\pgfpathlineto{\pgfqpoint{1.223563in}{1.200927in}}%
\pgfpathlineto{\pgfqpoint{1.229573in}{1.206543in}}%
\pgfpathlineto{\pgfqpoint{1.231812in}{1.208656in}}%
\pgfpathlineto{\pgfqpoint{1.237180in}{1.213766in}}%
\pgfpathlineto{\pgfqpoint{1.240061in}{1.216534in}}%
\pgfpathlineto{\pgfqpoint{1.244656in}{1.220988in}}%
\pgfpathlineto{\pgfqpoint{1.248310in}{1.224564in}}%
\pgfpathlineto{\pgfqpoint{1.252003in}{1.228210in}}%
\pgfpathlineto{\pgfqpoint{1.256560in}{1.232751in}}%
\pgfpathlineto{\pgfqpoint{1.259226in}{1.235432in}}%
\pgfpathlineto{\pgfqpoint{1.264809in}{1.241099in}}%
\pgfpathlineto{\pgfqpoint{1.266328in}{1.242654in}}%
\pgfpathlineto{\pgfqpoint{1.273058in}{1.249612in}}%
\pgfpathlineto{\pgfqpoint{1.273312in}{1.249877in}}%
\pgfpathlineto{\pgfqpoint{1.280176in}{1.257099in}}%
\pgfpathlineto{\pgfqpoint{1.281307in}{1.258301in}}%
\pgfpathlineto{\pgfqpoint{1.286926in}{1.264321in}}%
\pgfpathlineto{\pgfqpoint{1.289557in}{1.267167in}}%
\pgfpathlineto{\pgfqpoint{1.293568in}{1.271543in}}%
\pgfpathlineto{\pgfqpoint{1.297806in}{1.276212in}}%
\pgfpathlineto{\pgfqpoint{1.300103in}{1.278766in}}%
\pgfpathlineto{\pgfqpoint{1.306055in}{1.285444in}}%
\pgfpathlineto{\pgfqpoint{1.306535in}{1.285988in}}%
\pgfpathlineto{\pgfqpoint{1.312859in}{1.293210in}}%
\pgfpathlineto{\pgfqpoint{1.314304in}{1.294876in}}%
\pgfpathlineto{\pgfqpoint{1.319082in}{1.300432in}}%
\pgfpathlineto{\pgfqpoint{1.322553in}{1.304508in}}%
\pgfpathlineto{\pgfqpoint{1.325209in}{1.307654in}}%
\pgfpathlineto{\pgfqpoint{1.330803in}{1.314345in}}%
\pgfpathlineto{\pgfqpoint{1.331243in}{1.314877in}}%
\pgfpathlineto{\pgfqpoint{1.337176in}{1.322099in}}%
\pgfpathlineto{\pgfqpoint{1.339052in}{1.324403in}}%
\pgfpathlineto{\pgfqpoint{1.343019in}{1.329321in}}%
\pgfpathlineto{\pgfqpoint{1.347301in}{1.334682in}}%
\pgfpathlineto{\pgfqpoint{1.348775in}{1.336543in}}%
\pgfpathlineto{\pgfqpoint{1.354441in}{1.343766in}}%
\pgfpathlineto{\pgfqpoint{1.355550in}{1.345192in}}%
\pgfpathlineto{\pgfqpoint{1.360017in}{1.350988in}}%
\pgfpathlineto{\pgfqpoint{1.363800in}{1.355943in}}%
\pgfpathlineto{\pgfqpoint{1.365514in}{1.358210in}}%
\pgfpathlineto{\pgfqpoint{1.370928in}{1.365432in}}%
\pgfpathlineto{\pgfqpoint{1.372049in}{1.366941in}}%
\pgfpathlineto{\pgfqpoint{1.376256in}{1.372654in}}%
\pgfpathlineto{\pgfqpoint{1.380298in}{1.378197in}}%
\pgfpathlineto{\pgfqpoint{1.381511in}{1.379877in}}%
\pgfpathlineto{\pgfqpoint{1.386685in}{1.387099in}}%
\pgfpathlineto{\pgfqpoint{1.388547in}{1.389721in}}%
\pgfpathlineto{\pgfqpoint{1.391784in}{1.394321in}}%
\pgfpathlineto{\pgfqpoint{1.396797in}{1.401515in}}%
\pgfpathlineto{\pgfqpoint{1.396816in}{1.401543in}}%
\pgfpathlineto{\pgfqpoint{1.401764in}{1.408766in}}%
\pgfpathlineto{\pgfqpoint{1.405046in}{1.413601in}}%
\pgfpathlineto{\pgfqpoint{1.406650in}{1.415988in}}%
\pgfpathlineto{\pgfqpoint{1.411465in}{1.423210in}}%
\pgfpathlineto{\pgfqpoint{1.413295in}{1.425978in}}%
\pgfpathlineto{\pgfqpoint{1.416212in}{1.430432in}}%
\pgfpathlineto{\pgfqpoint{1.420899in}{1.437654in}}%
\pgfpathlineto{\pgfqpoint{1.421544in}{1.438656in}}%
\pgfpathlineto{\pgfqpoint{1.425513in}{1.444877in}}%
\pgfpathlineto{\pgfqpoint{1.429793in}{1.451650in}}%
\pgfpathlineto{\pgfqpoint{1.430075in}{1.452099in}}%
\pgfpathlineto{\pgfqpoint{1.434564in}{1.459321in}}%
\pgfpathlineto{\pgfqpoint{1.438043in}{1.464970in}}%
\pgfpathlineto{\pgfqpoint{1.439002in}{1.466543in}}%
\pgfpathlineto{\pgfqpoint{1.443375in}{1.473766in}}%
\pgfpathlineto{\pgfqpoint{1.446292in}{1.478625in}}%
\pgfpathlineto{\pgfqpoint{1.447697in}{1.480988in}}%
\pgfpathlineto{\pgfqpoint{1.451958in}{1.488210in}}%
\pgfpathlineto{\pgfqpoint{1.454541in}{1.492626in}}%
\pgfpathlineto{\pgfqpoint{1.456167in}{1.495432in}}%
\pgfpathlineto{\pgfqpoint{1.460320in}{1.502654in}}%
\pgfpathlineto{\pgfqpoint{1.462790in}{1.506985in}}%
\pgfpathlineto{\pgfqpoint{1.464424in}{1.509877in}}%
\pgfpathlineto{\pgfqpoint{1.468473in}{1.517099in}}%
\pgfpathlineto{\pgfqpoint{1.471040in}{1.521715in}}%
\pgfpathlineto{\pgfqpoint{1.472475in}{1.524321in}}%
\pgfpathlineto{\pgfqpoint{1.476423in}{1.531543in}}%
\pgfpathlineto{\pgfqpoint{1.479289in}{1.536829in}}%
\pgfpathlineto{\pgfqpoint{1.480329in}{1.538766in}}%
\pgfpathlineto{\pgfqpoint{1.484180in}{1.545988in}}%
\pgfpathlineto{\pgfqpoint{1.487538in}{1.552338in}}%
\pgfpathlineto{\pgfqpoint{1.487994in}{1.553210in}}%
\pgfpathlineto{\pgfqpoint{1.491752in}{1.560432in}}%
\pgfpathlineto{\pgfqpoint{1.495477in}{1.567654in}}%
\pgfpathlineto{\pgfqpoint{1.495787in}{1.568260in}}%
\pgfpathlineto{\pgfqpoint{1.499145in}{1.574877in}}%
\pgfpathlineto{\pgfqpoint{1.502780in}{1.582099in}}%
\pgfpathlineto{\pgfqpoint{1.504036in}{1.584612in}}%
\pgfpathlineto{\pgfqpoint{1.506367in}{1.589321in}}%
\pgfpathlineto{\pgfqpoint{1.509915in}{1.596543in}}%
\pgfpathlineto{\pgfqpoint{1.512286in}{1.601406in}}%
\pgfpathlineto{\pgfqpoint{1.513425in}{1.603766in}}%
\pgfpathlineto{\pgfqpoint{1.516888in}{1.610988in}}%
\pgfpathlineto{\pgfqpoint{1.520323in}{1.618210in}}%
\pgfpathlineto{\pgfqpoint{1.520535in}{1.618657in}}%
\pgfpathlineto{\pgfqpoint{1.523707in}{1.625432in}}%
\pgfpathlineto{\pgfqpoint{1.527061in}{1.632654in}}%
\pgfpathlineto{\pgfqpoint{1.528784in}{1.636389in}}%
\pgfpathlineto{\pgfqpoint{1.530377in}{1.639877in}}%
\pgfpathlineto{\pgfqpoint{1.533653in}{1.647099in}}%
\pgfpathlineto{\pgfqpoint{1.536903in}{1.654321in}}%
\pgfpathlineto{\pgfqpoint{1.537033in}{1.654612in}}%
\pgfpathlineto{\pgfqpoint{1.540104in}{1.661543in}}%
\pgfpathlineto{\pgfqpoint{1.543280in}{1.668766in}}%
\pgfpathlineto{\pgfqpoint{1.545283in}{1.673351in}}%
\pgfpathlineto{\pgfqpoint{1.546422in}{1.675988in}}%
\pgfpathlineto{\pgfqpoint{1.549525in}{1.683210in}}%
\pgfpathlineto{\pgfqpoint{1.552605in}{1.690432in}}%
\pgfpathlineto{\pgfqpoint{1.553532in}{1.692619in}}%
\pgfpathlineto{\pgfqpoint{1.555644in}{1.697654in}}%
\pgfpathlineto{\pgfqpoint{1.558655in}{1.704877in}}%
\pgfpathlineto{\pgfqpoint{1.561643in}{1.712099in}}%
\pgfpathlineto{\pgfqpoint{1.561781in}{1.712434in}}%
\pgfpathlineto{\pgfqpoint{1.564586in}{1.719321in}}%
\pgfpathlineto{\pgfqpoint{1.567508in}{1.726543in}}%
\pgfpathlineto{\pgfqpoint{1.570030in}{1.732820in}}%
\pgfpathlineto{\pgfqpoint{1.570406in}{1.733766in}}%
\pgfpathlineto{\pgfqpoint{1.573263in}{1.740988in}}%
\pgfpathlineto{\pgfqpoint{1.576100in}{1.748210in}}%
\pgfpathlineto{\pgfqpoint{1.578279in}{1.753793in}}%
\pgfpathlineto{\pgfqpoint{1.578912in}{1.755432in}}%
\pgfpathlineto{\pgfqpoint{1.581688in}{1.762654in}}%
\pgfpathlineto{\pgfqpoint{1.584444in}{1.769877in}}%
\pgfpathlineto{\pgfqpoint{1.586529in}{1.775371in}}%
\pgfpathlineto{\pgfqpoint{1.587177in}{1.777099in}}%
\pgfpathlineto{\pgfqpoint{1.589874in}{1.784321in}}%
\pgfpathlineto{\pgfqpoint{1.592554in}{1.791543in}}%
\pgfpathlineto{\pgfqpoint{1.594778in}{1.797571in}}%
\pgfpathlineto{\pgfqpoint{1.595214in}{1.798766in}}%
\pgfpathlineto{\pgfqpoint{1.597837in}{1.805988in}}%
\pgfpathlineto{\pgfqpoint{1.600444in}{1.813210in}}%
\pgfpathlineto{\pgfqpoint{1.603027in}{1.820412in}}%
\pgfpathlineto{\pgfqpoint{1.603034in}{1.820432in}}%
\pgfpathlineto{\pgfqpoint{1.605587in}{1.827654in}}%
\pgfpathlineto{\pgfqpoint{1.608124in}{1.834877in}}%
\pgfpathlineto{\pgfqpoint{1.610646in}{1.842099in}}%
\pgfpathlineto{\pgfqpoint{1.611276in}{1.843911in}}%
\pgfpathlineto{\pgfqpoint{1.613137in}{1.849321in}}%
\pgfpathlineto{\pgfqpoint{1.615607in}{1.856543in}}%
\pgfpathlineto{\pgfqpoint{1.618064in}{1.863766in}}%
\pgfpathlineto{\pgfqpoint{1.619526in}{1.868083in}}%
\pgfpathlineto{\pgfqpoint{1.620498in}{1.870988in}}%
\pgfpathlineto{\pgfqpoint{1.622905in}{1.878210in}}%
\pgfpathlineto{\pgfqpoint{1.625299in}{1.885432in}}%
\pgfpathlineto{\pgfqpoint{1.627681in}{1.892654in}}%
\pgfpathlineto{\pgfqpoint{1.627775in}{1.892941in}}%
\pgfpathlineto{\pgfqpoint{1.630028in}{1.899877in}}%
\pgfpathlineto{\pgfqpoint{1.627775in}{1.899877in}}%
\pgfpathlineto{\pgfqpoint{1.619526in}{1.899877in}}%
\pgfpathlineto{\pgfqpoint{1.611276in}{1.899877in}}%
\pgfpathlineto{\pgfqpoint{1.603027in}{1.899877in}}%
\pgfpathlineto{\pgfqpoint{1.594778in}{1.899877in}}%
\pgfpathlineto{\pgfqpoint{1.586529in}{1.899877in}}%
\pgfpathlineto{\pgfqpoint{1.578279in}{1.899877in}}%
\pgfpathlineto{\pgfqpoint{1.570030in}{1.899877in}}%
\pgfpathlineto{\pgfqpoint{1.561781in}{1.899877in}}%
\pgfpathlineto{\pgfqpoint{1.553532in}{1.899877in}}%
\pgfpathlineto{\pgfqpoint{1.545283in}{1.899877in}}%
\pgfpathlineto{\pgfqpoint{1.537033in}{1.899877in}}%
\pgfpathlineto{\pgfqpoint{1.528784in}{1.899877in}}%
\pgfpathlineto{\pgfqpoint{1.520535in}{1.899877in}}%
\pgfpathlineto{\pgfqpoint{1.512286in}{1.899877in}}%
\pgfpathlineto{\pgfqpoint{1.504036in}{1.899877in}}%
\pgfpathlineto{\pgfqpoint{1.495787in}{1.899877in}}%
\pgfpathlineto{\pgfqpoint{1.487538in}{1.899877in}}%
\pgfpathlineto{\pgfqpoint{1.479289in}{1.899877in}}%
\pgfpathlineto{\pgfqpoint{1.471040in}{1.899877in}}%
\pgfpathlineto{\pgfqpoint{1.462790in}{1.899877in}}%
\pgfpathlineto{\pgfqpoint{1.454541in}{1.899877in}}%
\pgfpathlineto{\pgfqpoint{1.446292in}{1.899877in}}%
\pgfpathlineto{\pgfqpoint{1.438043in}{1.899877in}}%
\pgfpathlineto{\pgfqpoint{1.429793in}{1.899877in}}%
\pgfpathlineto{\pgfqpoint{1.421544in}{1.899877in}}%
\pgfpathlineto{\pgfqpoint{1.413295in}{1.899877in}}%
\pgfpathlineto{\pgfqpoint{1.405046in}{1.899877in}}%
\pgfpathlineto{\pgfqpoint{1.396797in}{1.899877in}}%
\pgfpathlineto{\pgfqpoint{1.388547in}{1.899877in}}%
\pgfpathlineto{\pgfqpoint{1.380298in}{1.899877in}}%
\pgfpathlineto{\pgfqpoint{1.372049in}{1.899877in}}%
\pgfpathlineto{\pgfqpoint{1.363800in}{1.899877in}}%
\pgfpathlineto{\pgfqpoint{1.355550in}{1.899877in}}%
\pgfpathlineto{\pgfqpoint{1.352991in}{1.899877in}}%
\pgfpathlineto{\pgfqpoint{1.349327in}{1.892654in}}%
\pgfpathlineto{\pgfqpoint{1.347301in}{1.888675in}}%
\pgfpathlineto{\pgfqpoint{1.345633in}{1.885432in}}%
\pgfpathlineto{\pgfqpoint{1.341903in}{1.878210in}}%
\pgfpathlineto{\pgfqpoint{1.339052in}{1.872711in}}%
\pgfpathlineto{\pgfqpoint{1.338149in}{1.870988in}}%
\pgfpathlineto{\pgfqpoint{1.334352in}{1.863766in}}%
\pgfpathlineto{\pgfqpoint{1.330803in}{1.857042in}}%
\pgfpathlineto{\pgfqpoint{1.330537in}{1.856543in}}%
\pgfpathlineto{\pgfqpoint{1.326671in}{1.849321in}}%
\pgfpathlineto{\pgfqpoint{1.322788in}{1.842099in}}%
\pgfpathlineto{\pgfqpoint{1.322553in}{1.841664in}}%
\pgfpathlineto{\pgfqpoint{1.318854in}{1.834877in}}%
\pgfpathlineto{\pgfqpoint{1.314900in}{1.827654in}}%
\pgfpathlineto{\pgfqpoint{1.314304in}{1.826571in}}%
\pgfpathlineto{\pgfqpoint{1.310897in}{1.820432in}}%
\pgfpathlineto{\pgfqpoint{1.306870in}{1.813210in}}%
\pgfpathlineto{\pgfqpoint{1.306055in}{1.811755in}}%
\pgfpathlineto{\pgfqpoint{1.302794in}{1.805988in}}%
\pgfpathlineto{\pgfqpoint{1.298693in}{1.798766in}}%
\pgfpathlineto{\pgfqpoint{1.297806in}{1.797210in}}%
\pgfpathlineto{\pgfqpoint{1.294543in}{1.791543in}}%
\pgfpathlineto{\pgfqpoint{1.290365in}{1.784321in}}%
\pgfpathlineto{\pgfqpoint{1.289557in}{1.782930in}}%
\pgfpathlineto{\pgfqpoint{1.286136in}{1.777099in}}%
\pgfpathlineto{\pgfqpoint{1.281880in}{1.769877in}}%
\pgfpathlineto{\pgfqpoint{1.281307in}{1.768909in}}%
\pgfpathlineto{\pgfqpoint{1.277570in}{1.762654in}}%
\pgfpathlineto{\pgfqpoint{1.273234in}{1.755432in}}%
\pgfpathlineto{\pgfqpoint{1.273058in}{1.755141in}}%
\pgfpathlineto{\pgfqpoint{1.268839in}{1.748210in}}%
\pgfpathlineto{\pgfqpoint{1.264809in}{1.741622in}}%
\pgfpathlineto{\pgfqpoint{1.264417in}{1.740988in}}%
\pgfpathlineto{\pgfqpoint{1.259937in}{1.733766in}}%
\pgfpathlineto{\pgfqpoint{1.256560in}{1.728347in}}%
\pgfpathlineto{\pgfqpoint{1.255425in}{1.726543in}}%
\pgfpathlineto{\pgfqpoint{1.250859in}{1.719321in}}%
\pgfpathlineto{\pgfqpoint{1.248310in}{1.715310in}}%
\pgfpathlineto{\pgfqpoint{1.246252in}{1.712099in}}%
\pgfpathlineto{\pgfqpoint{1.241598in}{1.704877in}}%
\pgfpathlineto{\pgfqpoint{1.240061in}{1.702504in}}%
\pgfpathlineto{\pgfqpoint{1.236892in}{1.697654in}}%
\pgfpathlineto{\pgfqpoint{1.232148in}{1.690432in}}%
\pgfpathlineto{\pgfqpoint{1.231812in}{1.689924in}}%
\pgfpathlineto{\pgfqpoint{1.227340in}{1.683210in}}%
\pgfpathlineto{\pgfqpoint{1.223563in}{1.677571in}}%
\pgfpathlineto{\pgfqpoint{1.222493in}{1.675988in}}%
\pgfpathlineto{\pgfqpoint{1.217588in}{1.668766in}}%
\pgfpathlineto{\pgfqpoint{1.215314in}{1.665435in}}%
\pgfpathlineto{\pgfqpoint{1.212632in}{1.661543in}}%
\pgfpathlineto{\pgfqpoint{1.207630in}{1.654321in}}%
\pgfpathlineto{\pgfqpoint{1.207064in}{1.653509in}}%
\pgfpathlineto{\pgfqpoint{1.202561in}{1.647099in}}%
\pgfpathlineto{\pgfqpoint{1.198815in}{1.641796in}}%
\pgfpathlineto{\pgfqpoint{1.197447in}{1.639877in}}%
\pgfpathlineto{\pgfqpoint{1.192272in}{1.632654in}}%
\pgfpathlineto{\pgfqpoint{1.190566in}{1.630286in}}%
\pgfpathlineto{\pgfqpoint{1.187038in}{1.625432in}}%
\pgfpathlineto{\pgfqpoint{1.182317in}{1.618973in}}%
\pgfpathlineto{\pgfqpoint{1.181754in}{1.618210in}}%
\pgfpathlineto{\pgfqpoint{1.176400in}{1.610988in}}%
\pgfpathlineto{\pgfqpoint{1.174067in}{1.607859in}}%
\pgfpathlineto{\pgfqpoint{1.170989in}{1.603766in}}%
\pgfpathlineto{\pgfqpoint{1.165818in}{1.596929in}}%
\pgfpathlineto{\pgfqpoint{1.165524in}{1.596543in}}%
\pgfpathlineto{\pgfqpoint{1.159984in}{1.589321in}}%
\pgfpathlineto{\pgfqpoint{1.157569in}{1.586192in}}%
\pgfpathlineto{\pgfqpoint{1.154384in}{1.582099in}}%
\pgfpathlineto{\pgfqpoint{1.149320in}{1.575628in}}%
\pgfpathlineto{\pgfqpoint{1.148727in}{1.574877in}}%
\pgfpathlineto{\pgfqpoint{1.142994in}{1.567654in}}%
\pgfpathlineto{\pgfqpoint{1.141071in}{1.565246in}}%
\pgfpathlineto{\pgfqpoint{1.137195in}{1.560432in}}%
\pgfpathlineto{\pgfqpoint{1.132821in}{1.555032in}}%
\pgfpathlineto{\pgfqpoint{1.131334in}{1.553210in}}%
\pgfpathlineto{\pgfqpoint{1.125401in}{1.545988in}}%
\pgfpathlineto{\pgfqpoint{1.124572in}{1.544985in}}%
\pgfpathlineto{\pgfqpoint{1.119391in}{1.538766in}}%
\pgfpathlineto{\pgfqpoint{1.116323in}{1.535105in}}%
\pgfpathlineto{\pgfqpoint{1.113315in}{1.531543in}}%
\pgfpathlineto{\pgfqpoint{1.108074in}{1.525378in}}%
\pgfpathlineto{\pgfqpoint{1.107168in}{1.524321in}}%
\pgfpathlineto{\pgfqpoint{1.100943in}{1.517099in}}%
\pgfpathlineto{\pgfqpoint{1.099824in}{1.515809in}}%
\pgfpathlineto{\pgfqpoint{1.094639in}{1.509877in}}%
\pgfpathlineto{\pgfqpoint{1.091575in}{1.506394in}}%
\pgfpathlineto{\pgfqpoint{1.088261in}{1.502654in}}%
\pgfpathlineto{\pgfqpoint{1.083326in}{1.497122in}}%
\pgfpathlineto{\pgfqpoint{1.081808in}{1.495432in}}%
\pgfpathlineto{\pgfqpoint{1.075275in}{1.488210in}}%
\pgfpathlineto{\pgfqpoint{1.075077in}{1.487992in}}%
\pgfpathlineto{\pgfqpoint{1.068652in}{1.480988in}}%
\pgfpathlineto{\pgfqpoint{1.066828in}{1.479011in}}%
\pgfpathlineto{\pgfqpoint{1.061948in}{1.473766in}}%
\pgfpathlineto{\pgfqpoint{1.058578in}{1.470167in}}%
\pgfpathlineto{\pgfqpoint{1.055160in}{1.466543in}}%
\pgfpathlineto{\pgfqpoint{1.050329in}{1.461456in}}%
\pgfpathlineto{\pgfqpoint{1.048286in}{1.459321in}}%
\pgfpathlineto{\pgfqpoint{1.042080in}{1.452878in}}%
\pgfpathlineto{\pgfqpoint{1.041324in}{1.452099in}}%
\pgfpathlineto{\pgfqpoint{1.034268in}{1.444877in}}%
\pgfpathlineto{\pgfqpoint{1.033831in}{1.444432in}}%
\pgfpathlineto{\pgfqpoint{1.027114in}{1.437654in}}%
\pgfpathlineto{\pgfqpoint{1.025581in}{1.436119in}}%
\pgfpathlineto{\pgfqpoint{1.019865in}{1.430432in}}%
\pgfpathlineto{\pgfqpoint{1.017332in}{1.427931in}}%
\pgfpathlineto{\pgfqpoint{1.012518in}{1.423210in}}%
\pgfpathlineto{\pgfqpoint{1.009083in}{1.419866in}}%
\pgfpathlineto{\pgfqpoint{1.005071in}{1.415988in}}%
\pgfpathlineto{\pgfqpoint{1.000834in}{1.411922in}}%
\pgfpathlineto{\pgfqpoint{0.997520in}{1.408766in}}%
\pgfpathlineto{\pgfqpoint{0.992585in}{1.404097in}}%
\pgfpathlineto{\pgfqpoint{0.989864in}{1.401543in}}%
\pgfpathlineto{\pgfqpoint{0.984335in}{1.396390in}}%
\pgfpathlineto{\pgfqpoint{0.982100in}{1.394321in}}%
\pgfpathlineto{\pgfqpoint{0.976086in}{1.388796in}}%
\pgfpathlineto{\pgfqpoint{0.974225in}{1.387099in}}%
\pgfpathlineto{\pgfqpoint{0.967837in}{1.381315in}}%
\pgfpathlineto{\pgfqpoint{0.966238in}{1.379877in}}%
\pgfpathlineto{\pgfqpoint{0.959588in}{1.373942in}}%
\pgfpathlineto{\pgfqpoint{0.958135in}{1.372654in}}%
\pgfpathlineto{\pgfqpoint{0.951338in}{1.366676in}}%
\pgfpathlineto{\pgfqpoint{0.949914in}{1.365432in}}%
\pgfpathlineto{\pgfqpoint{0.943089in}{1.359515in}}%
\pgfpathlineto{\pgfqpoint{0.941574in}{1.358210in}}%
\pgfpathlineto{\pgfqpoint{0.934840in}{1.352455in}}%
\pgfpathlineto{\pgfqpoint{0.933112in}{1.350988in}}%
\pgfpathlineto{\pgfqpoint{0.926591in}{1.345494in}}%
\pgfpathlineto{\pgfqpoint{0.924526in}{1.343766in}}%
\pgfpathlineto{\pgfqpoint{0.918342in}{1.338629in}}%
\pgfpathlineto{\pgfqpoint{0.915813in}{1.336543in}}%
\pgfpathlineto{\pgfqpoint{0.910092in}{1.331860in}}%
\pgfpathlineto{\pgfqpoint{0.906972in}{1.329321in}}%
\pgfpathlineto{\pgfqpoint{0.901843in}{1.325182in}}%
\pgfpathlineto{\pgfqpoint{0.897999in}{1.322099in}}%
\pgfpathlineto{\pgfqpoint{0.893594in}{1.318594in}}%
\pgfpathlineto{\pgfqpoint{0.888892in}{1.314877in}}%
\pgfpathlineto{\pgfqpoint{0.885345in}{1.312094in}}%
\pgfpathlineto{\pgfqpoint{0.879650in}{1.307654in}}%
\pgfpathlineto{\pgfqpoint{0.877095in}{1.305679in}}%
\pgfpathlineto{\pgfqpoint{0.870268in}{1.300432in}}%
\pgfpathlineto{\pgfqpoint{0.868846in}{1.299349in}}%
\pgfpathlineto{\pgfqpoint{0.860743in}{1.293210in}}%
\pgfpathlineto{\pgfqpoint{0.860597in}{1.293100in}}%
\pgfpathlineto{\pgfqpoint{0.852348in}{1.286940in}}%
\pgfpathlineto{\pgfqpoint{0.851065in}{1.285988in}}%
\pgfpathlineto{\pgfqpoint{0.844099in}{1.280860in}}%
\pgfpathlineto{\pgfqpoint{0.841236in}{1.278766in}}%
\pgfpathlineto{\pgfqpoint{0.835849in}{1.274857in}}%
\pgfpathlineto{\pgfqpoint{0.831254in}{1.271543in}}%
\pgfpathlineto{\pgfqpoint{0.827600in}{1.268930in}}%
\pgfpathlineto{\pgfqpoint{0.821116in}{1.264321in}}%
\pgfpathlineto{\pgfqpoint{0.819351in}{1.263077in}}%
\pgfpathlineto{\pgfqpoint{0.811102in}{1.257298in}}%
\pgfpathlineto{\pgfqpoint{0.810816in}{1.257099in}}%
\pgfpathlineto{\pgfqpoint{0.802852in}{1.251601in}}%
\pgfpathlineto{\pgfqpoint{0.800340in}{1.249877in}}%
\pgfpathlineto{\pgfqpoint{0.794603in}{1.245973in}}%
\pgfpathlineto{\pgfqpoint{0.789697in}{1.242654in}}%
\pgfpathlineto{\pgfqpoint{0.786354in}{1.240413in}}%
\pgfpathlineto{\pgfqpoint{0.778882in}{1.235432in}}%
\pgfpathlineto{\pgfqpoint{0.778105in}{1.234918in}}%
\pgfpathlineto{\pgfqpoint{0.769856in}{1.229499in}}%
\pgfpathlineto{\pgfqpoint{0.767883in}{1.228210in}}%
\pgfpathlineto{\pgfqpoint{0.761606in}{1.224146in}}%
\pgfpathlineto{\pgfqpoint{0.756701in}{1.220988in}}%
\pgfpathlineto{\pgfqpoint{0.753357in}{1.218854in}}%
\pgfpathlineto{\pgfqpoint{0.745338in}{1.213766in}}%
\pgfpathlineto{\pgfqpoint{0.745108in}{1.213621in}}%
\pgfpathlineto{\pgfqpoint{0.736859in}{1.208462in}}%
\pgfpathlineto{\pgfqpoint{0.733774in}{1.206543in}}%
\pgfpathlineto{\pgfqpoint{0.728609in}{1.203360in}}%
\pgfpathlineto{\pgfqpoint{0.722022in}{1.199321in}}%
\pgfpathlineto{\pgfqpoint{0.720360in}{1.198312in}}%
\pgfpathlineto{\pgfqpoint{0.712111in}{1.193327in}}%
\pgfpathlineto{\pgfqpoint{0.710069in}{1.192099in}}%
\pgfpathlineto{\pgfqpoint{0.703862in}{1.188402in}}%
\pgfpathlineto{\pgfqpoint{0.697912in}{1.184877in}}%
\pgfpathlineto{\pgfqpoint{0.695612in}{1.183527in}}%
\pgfpathlineto{\pgfqpoint{0.687363in}{1.178709in}}%
\pgfpathlineto{\pgfqpoint{0.685548in}{1.177654in}}%
\pgfpathlineto{\pgfqpoint{0.679114in}{1.173952in}}%
\pgfpathlineto{\pgfqpoint{0.672967in}{1.170432in}}%
\pgfpathlineto{\pgfqpoint{0.670865in}{1.169240in}}%
\pgfpathlineto{\pgfqpoint{0.662616in}{1.164586in}}%
\pgfpathlineto{\pgfqpoint{0.660163in}{1.163210in}}%
\pgfpathlineto{\pgfqpoint{0.654366in}{1.159988in}}%
\pgfpathlineto{\pgfqpoint{0.647131in}{1.155988in}}%
\pgfpathlineto{\pgfqpoint{0.646117in}{1.155433in}}%
\pgfpathlineto{\pgfqpoint{0.637868in}{1.150939in}}%
\pgfpathlineto{\pgfqpoint{0.633856in}{1.148766in}}%
\pgfpathlineto{\pgfqpoint{0.629619in}{1.146492in}}%
\pgfpathlineto{\pgfqpoint{0.621369in}{1.142090in}}%
\pgfpathlineto{\pgfqpoint{0.620339in}{1.141543in}}%
\pgfpathlineto{\pgfqpoint{0.613120in}{1.137749in}}%
\pgfpathlineto{\pgfqpoint{0.606562in}{1.134321in}}%
\pgfpathlineto{\pgfqpoint{0.604871in}{1.133446in}}%
\pgfpathlineto{\pgfqpoint{0.596622in}{1.129200in}}%
\pgfpathlineto{\pgfqpoint{0.592519in}{1.127099in}}%
\pgfpathlineto{\pgfqpoint{0.588373in}{1.124997in}}%
\pgfpathlineto{\pgfqpoint{0.580123in}{1.120838in}}%
\pgfpathlineto{\pgfqpoint{0.578207in}{1.119877in}}%
\pgfpathlineto{\pgfqpoint{0.571874in}{1.116733in}}%
\pgfpathlineto{\pgfqpoint{0.563625in}{1.112658in}}%
\pgfpathlineto{\pgfqpoint{0.563617in}{1.112654in}}%
\pgfpathlineto{\pgfqpoint{0.555376in}{1.108646in}}%
\pgfpathlineto{\pgfqpoint{0.548734in}{1.105432in}}%
\pgfpathlineto{\pgfqpoint{0.547126in}{1.104663in}}%
\pgfpathlineto{\pgfqpoint{0.538877in}{1.100732in}}%
\pgfpathlineto{\pgfqpoint{0.533558in}{1.098210in}}%
\pgfpathlineto{\pgfqpoint{0.530628in}{1.096836in}}%
\pgfpathlineto{\pgfqpoint{0.522379in}{1.092985in}}%
\pgfpathlineto{\pgfqpoint{0.518081in}{1.090988in}}%
\pgfpathlineto{\pgfqpoint{0.514130in}{1.089172in}}%
\pgfpathlineto{\pgfqpoint{0.505880in}{1.085398in}}%
\pgfpathlineto{\pgfqpoint{0.502295in}{1.083766in}}%
\pgfpathlineto{\pgfqpoint{0.497631in}{1.081665in}}%
\pgfpathlineto{\pgfqpoint{0.489382in}{1.077967in}}%
\pgfpathlineto{\pgfqpoint{0.486192in}{1.076543in}}%
\pgfpathlineto{\pgfqpoint{0.481133in}{1.074310in}}%
\pgfpathlineto{\pgfqpoint{0.472883in}{1.070686in}}%
\pgfpathlineto{\pgfqpoint{0.469763in}{1.069321in}}%
\pgfpathlineto{\pgfqpoint{0.464634in}{1.067102in}}%
\pgfpathlineto{\pgfqpoint{0.456385in}{1.063550in}}%
\pgfpathlineto{\pgfqpoint{0.453001in}{1.062099in}}%
\pgfpathlineto{\pgfqpoint{0.448136in}{1.060036in}}%
\pgfpathlineto{\pgfqpoint{0.439887in}{1.056554in}}%
\pgfpathlineto{\pgfqpoint{0.435896in}{1.054877in}}%
\pgfpathlineto{\pgfqpoint{0.431637in}{1.053107in}}%
\pgfpathlineto{\pgfqpoint{0.423388in}{1.049693in}}%
\pgfpathlineto{\pgfqpoint{0.418442in}{1.047654in}}%
\pgfpathlineto{\pgfqpoint{0.415139in}{1.046309in}}%
\pgfpathlineto{\pgfqpoint{0.415139in}{1.040432in}}%
\pgfpathlineto{\pgfqpoint{0.415139in}{1.033210in}}%
\pgfpathlineto{\pgfqpoint{0.415139in}{1.025988in}}%
\pgfpathlineto{\pgfqpoint{0.415139in}{1.018766in}}%
\pgfpathlineto{\pgfqpoint{0.415139in}{1.011543in}}%
\pgfpathlineto{\pgfqpoint{0.415139in}{1.004321in}}%
\pgfpathlineto{\pgfqpoint{0.415139in}{0.997099in}}%
\pgfpathlineto{\pgfqpoint{0.415139in}{0.989877in}}%
\pgfpathlineto{\pgfqpoint{0.415139in}{0.982654in}}%
\pgfpathlineto{\pgfqpoint{0.415139in}{0.975432in}}%
\pgfpathlineto{\pgfqpoint{0.415139in}{0.968210in}}%
\pgfpathlineto{\pgfqpoint{0.415139in}{0.960988in}}%
\pgfpathlineto{\pgfqpoint{0.415139in}{0.953766in}}%
\pgfpathlineto{\pgfqpoint{0.415139in}{0.946543in}}%
\pgfpathlineto{\pgfqpoint{0.415139in}{0.939321in}}%
\pgfpathlineto{\pgfqpoint{0.415139in}{0.932099in}}%
\pgfpathlineto{\pgfqpoint{0.415139in}{0.924877in}}%
\pgfpathlineto{\pgfqpoint{0.415139in}{0.917654in}}%
\pgfpathlineto{\pgfqpoint{0.415139in}{0.910432in}}%
\pgfpathlineto{\pgfqpoint{0.415139in}{0.903210in}}%
\pgfpathlineto{\pgfqpoint{0.415139in}{0.895988in}}%
\pgfpathlineto{\pgfqpoint{0.415139in}{0.888766in}}%
\pgfpathlineto{\pgfqpoint{0.415139in}{0.881543in}}%
\pgfpathlineto{\pgfqpoint{0.415139in}{0.874321in}}%
\pgfpathlineto{\pgfqpoint{0.415139in}{0.867099in}}%
\pgfpathlineto{\pgfqpoint{0.415139in}{0.859877in}}%
\pgfpathlineto{\pgfqpoint{0.415139in}{0.852654in}}%
\pgfpathlineto{\pgfqpoint{0.415139in}{0.845432in}}%
\pgfpathlineto{\pgfqpoint{0.415139in}{0.838210in}}%
\pgfpathlineto{\pgfqpoint{0.415139in}{0.830988in}}%
\pgfpathlineto{\pgfqpoint{0.415139in}{0.823766in}}%
\pgfpathlineto{\pgfqpoint{0.415139in}{0.816543in}}%
\pgfpathlineto{\pgfqpoint{0.415139in}{0.809321in}}%
\pgfpathlineto{\pgfqpoint{0.415139in}{0.808109in}}%
\pgfpathlineto{\pgfqpoint{0.419871in}{0.809321in}}%
\pgfpathclose%
\pgfusepath{fill}%
\end{pgfscope}%
\begin{pgfscope}%
\pgfpathrectangle{\pgfqpoint{0.415139in}{0.462654in}}{\pgfqpoint{1.641596in}{1.437222in}}%
\pgfusepath{clip}%
\pgfsetbuttcap%
\pgfsetroundjoin%
\definecolor{currentfill}{rgb}{0.772472,0.900961,0.942561}%
\pgfsetfillcolor{currentfill}%
\pgfsetlinewidth{0.000000pt}%
\definecolor{currentstroke}{rgb}{0.000000,0.000000,0.000000}%
\pgfsetstrokecolor{currentstroke}%
\pgfsetdash{}{0pt}%
\pgfpathmoveto{\pgfqpoint{0.423388in}{0.631473in}}%
\pgfpathlineto{\pgfqpoint{0.431637in}{0.632760in}}%
\pgfpathlineto{\pgfqpoint{0.439887in}{0.634057in}}%
\pgfpathlineto{\pgfqpoint{0.448136in}{0.635364in}}%
\pgfpathlineto{\pgfqpoint{0.452052in}{0.635988in}}%
\pgfpathlineto{\pgfqpoint{0.456385in}{0.636685in}}%
\pgfpathlineto{\pgfqpoint{0.464634in}{0.638021in}}%
\pgfpathlineto{\pgfqpoint{0.472883in}{0.639367in}}%
\pgfpathlineto{\pgfqpoint{0.481133in}{0.640723in}}%
\pgfpathlineto{\pgfqpoint{0.489382in}{0.642090in}}%
\pgfpathlineto{\pgfqpoint{0.496088in}{0.643210in}}%
\pgfpathlineto{\pgfqpoint{0.497631in}{0.643470in}}%
\pgfpathlineto{\pgfqpoint{0.505880in}{0.644869in}}%
\pgfpathlineto{\pgfqpoint{0.514130in}{0.646278in}}%
\pgfpathlineto{\pgfqpoint{0.522379in}{0.647699in}}%
\pgfpathlineto{\pgfqpoint{0.530628in}{0.649132in}}%
\pgfpathlineto{\pgfqpoint{0.538056in}{0.650432in}}%
\pgfpathlineto{\pgfqpoint{0.538877in}{0.650577in}}%
\pgfpathlineto{\pgfqpoint{0.547126in}{0.652043in}}%
\pgfpathlineto{\pgfqpoint{0.555376in}{0.653522in}}%
\pgfpathlineto{\pgfqpoint{0.563625in}{0.655013in}}%
\pgfpathlineto{\pgfqpoint{0.571874in}{0.656516in}}%
\pgfpathlineto{\pgfqpoint{0.578070in}{0.657654in}}%
\pgfpathlineto{\pgfqpoint{0.580123in}{0.658035in}}%
\pgfpathlineto{\pgfqpoint{0.588373in}{0.659575in}}%
\pgfpathlineto{\pgfqpoint{0.596622in}{0.661128in}}%
\pgfpathlineto{\pgfqpoint{0.604871in}{0.662695in}}%
\pgfpathlineto{\pgfqpoint{0.613120in}{0.664275in}}%
\pgfpathlineto{\pgfqpoint{0.616237in}{0.664877in}}%
\pgfpathlineto{\pgfqpoint{0.621369in}{0.665877in}}%
\pgfpathlineto{\pgfqpoint{0.629619in}{0.667496in}}%
\pgfpathlineto{\pgfqpoint{0.637868in}{0.669129in}}%
\pgfpathlineto{\pgfqpoint{0.646117in}{0.670777in}}%
\pgfpathlineto{\pgfqpoint{0.652678in}{0.672099in}}%
\pgfpathlineto{\pgfqpoint{0.654366in}{0.672442in}}%
\pgfpathlineto{\pgfqpoint{0.662616in}{0.674131in}}%
\pgfpathlineto{\pgfqpoint{0.670865in}{0.675834in}}%
\pgfpathlineto{\pgfqpoint{0.679114in}{0.677552in}}%
\pgfpathlineto{\pgfqpoint{0.687363in}{0.679285in}}%
\pgfpathlineto{\pgfqpoint{0.687532in}{0.679321in}}%
\pgfpathlineto{\pgfqpoint{0.695612in}{0.681045in}}%
\pgfpathlineto{\pgfqpoint{0.703862in}{0.682820in}}%
\pgfpathlineto{\pgfqpoint{0.712111in}{0.684612in}}%
\pgfpathlineto{\pgfqpoint{0.720360in}{0.686419in}}%
\pgfpathlineto{\pgfqpoint{0.720926in}{0.686543in}}%
\pgfpathlineto{\pgfqpoint{0.728609in}{0.688253in}}%
\pgfpathlineto{\pgfqpoint{0.736859in}{0.690105in}}%
\pgfpathlineto{\pgfqpoint{0.745108in}{0.691973in}}%
\pgfpathlineto{\pgfqpoint{0.752948in}{0.693766in}}%
\pgfpathlineto{\pgfqpoint{0.753357in}{0.693860in}}%
\pgfpathlineto{\pgfqpoint{0.761606in}{0.695776in}}%
\pgfpathlineto{\pgfqpoint{0.769856in}{0.697709in}}%
\pgfpathlineto{\pgfqpoint{0.778105in}{0.699662in}}%
\pgfpathlineto{\pgfqpoint{0.783661in}{0.700988in}}%
\pgfpathlineto{\pgfqpoint{0.786354in}{0.701637in}}%
\pgfpathlineto{\pgfqpoint{0.794603in}{0.703640in}}%
\pgfpathlineto{\pgfqpoint{0.802852in}{0.705663in}}%
\pgfpathlineto{\pgfqpoint{0.811102in}{0.707706in}}%
\pgfpathlineto{\pgfqpoint{0.813123in}{0.708210in}}%
\pgfpathlineto{\pgfqpoint{0.819351in}{0.709779in}}%
\pgfpathlineto{\pgfqpoint{0.827600in}{0.711876in}}%
\pgfpathlineto{\pgfqpoint{0.835849in}{0.713995in}}%
\pgfpathlineto{\pgfqpoint{0.841399in}{0.715432in}}%
\pgfpathlineto{\pgfqpoint{0.844099in}{0.716139in}}%
\pgfpathlineto{\pgfqpoint{0.852348in}{0.718315in}}%
\pgfpathlineto{\pgfqpoint{0.860597in}{0.720512in}}%
\pgfpathlineto{\pgfqpoint{0.868557in}{0.722654in}}%
\pgfpathlineto{\pgfqpoint{0.868846in}{0.722733in}}%
\pgfpathlineto{\pgfqpoint{0.877095in}{0.724991in}}%
\pgfpathlineto{\pgfqpoint{0.885345in}{0.727272in}}%
\pgfpathlineto{\pgfqpoint{0.893594in}{0.729577in}}%
\pgfpathlineto{\pgfqpoint{0.894659in}{0.729877in}}%
\pgfpathlineto{\pgfqpoint{0.901843in}{0.731920in}}%
\pgfpathlineto{\pgfqpoint{0.910092in}{0.734289in}}%
\pgfpathlineto{\pgfqpoint{0.918342in}{0.736684in}}%
\pgfpathlineto{\pgfqpoint{0.919762in}{0.737099in}}%
\pgfpathlineto{\pgfqpoint{0.926591in}{0.739118in}}%
\pgfpathlineto{\pgfqpoint{0.934840in}{0.741581in}}%
\pgfpathlineto{\pgfqpoint{0.943089in}{0.744072in}}%
\pgfpathlineto{\pgfqpoint{0.943910in}{0.744321in}}%
\pgfpathlineto{\pgfqpoint{0.951338in}{0.746606in}}%
\pgfpathlineto{\pgfqpoint{0.959588in}{0.749170in}}%
\pgfpathlineto{\pgfqpoint{0.967140in}{0.751543in}}%
\pgfpathlineto{\pgfqpoint{0.967837in}{0.751765in}}%
\pgfpathlineto{\pgfqpoint{0.976086in}{0.754406in}}%
\pgfpathlineto{\pgfqpoint{0.984335in}{0.757079in}}%
\pgfpathlineto{\pgfqpoint{0.989491in}{0.758766in}}%
\pgfpathlineto{\pgfqpoint{0.992585in}{0.759790in}}%
\pgfpathlineto{\pgfqpoint{1.000834in}{0.762544in}}%
\pgfpathlineto{\pgfqpoint{1.009083in}{0.765332in}}%
\pgfpathlineto{\pgfqpoint{1.011008in}{0.765988in}}%
\pgfpathlineto{\pgfqpoint{1.017332in}{0.768168in}}%
\pgfpathlineto{\pgfqpoint{1.025581in}{0.771042in}}%
\pgfpathlineto{\pgfqpoint{1.031737in}{0.773210in}}%
\pgfpathlineto{\pgfqpoint{1.033831in}{0.773956in}}%
\pgfpathlineto{\pgfqpoint{1.042080in}{0.776920in}}%
\pgfpathlineto{\pgfqpoint{1.050329in}{0.779920in}}%
\pgfpathlineto{\pgfqpoint{1.051727in}{0.780432in}}%
\pgfpathlineto{\pgfqpoint{1.058578in}{0.782973in}}%
\pgfpathlineto{\pgfqpoint{1.066828in}{0.786067in}}%
\pgfpathlineto{\pgfqpoint{1.071021in}{0.787654in}}%
\pgfpathlineto{\pgfqpoint{1.075077in}{0.789208in}}%
\pgfpathlineto{\pgfqpoint{1.083326in}{0.792399in}}%
\pgfpathlineto{\pgfqpoint{1.089661in}{0.794877in}}%
\pgfpathlineto{\pgfqpoint{1.091575in}{0.795634in}}%
\pgfpathlineto{\pgfqpoint{1.099824in}{0.798925in}}%
\pgfpathlineto{\pgfqpoint{1.107682in}{0.802099in}}%
\pgfpathlineto{\pgfqpoint{1.108074in}{0.802259in}}%
\pgfpathlineto{\pgfqpoint{1.116323in}{0.805655in}}%
\pgfpathlineto{\pgfqpoint{1.124572in}{0.809095in}}%
\pgfpathlineto{\pgfqpoint{1.125111in}{0.809321in}}%
\pgfpathlineto{\pgfqpoint{1.132821in}{0.812600in}}%
\pgfpathlineto{\pgfqpoint{1.141071in}{0.816152in}}%
\pgfpathlineto{\pgfqpoint{1.141973in}{0.816543in}}%
\pgfpathlineto{\pgfqpoint{1.149320in}{0.819771in}}%
\pgfpathlineto{\pgfqpoint{1.157569in}{0.823441in}}%
\pgfpathlineto{\pgfqpoint{1.158292in}{0.823766in}}%
\pgfpathlineto{\pgfqpoint{1.165818in}{0.827182in}}%
\pgfpathlineto{\pgfqpoint{1.174067in}{0.830976in}}%
\pgfpathlineto{\pgfqpoint{1.174092in}{0.830988in}}%
\pgfpathlineto{\pgfqpoint{1.182317in}{0.834847in}}%
\pgfpathlineto{\pgfqpoint{1.189392in}{0.838210in}}%
\pgfpathlineto{\pgfqpoint{1.190566in}{0.838775in}}%
\pgfpathlineto{\pgfqpoint{1.198815in}{0.842779in}}%
\pgfpathlineto{\pgfqpoint{1.204215in}{0.845432in}}%
\pgfpathlineto{\pgfqpoint{1.207064in}{0.846849in}}%
\pgfpathlineto{\pgfqpoint{1.215314in}{0.850993in}}%
\pgfpathlineto{\pgfqpoint{1.218584in}{0.852654in}}%
\pgfpathlineto{\pgfqpoint{1.223563in}{0.855213in}}%
\pgfpathlineto{\pgfqpoint{1.231812in}{0.859505in}}%
\pgfpathlineto{\pgfqpoint{1.232520in}{0.859877in}}%
\pgfpathlineto{\pgfqpoint{1.240061in}{0.863883in}}%
\pgfpathlineto{\pgfqpoint{1.246039in}{0.867099in}}%
\pgfpathlineto{\pgfqpoint{1.248310in}{0.868336in}}%
\pgfpathlineto{\pgfqpoint{1.256560in}{0.872874in}}%
\pgfpathlineto{\pgfqpoint{1.259162in}{0.874321in}}%
\pgfpathlineto{\pgfqpoint{1.264809in}{0.877499in}}%
\pgfpathlineto{\pgfqpoint{1.271906in}{0.881543in}}%
\pgfpathlineto{\pgfqpoint{1.273058in}{0.882208in}}%
\pgfpathlineto{\pgfqpoint{1.281307in}{0.887011in}}%
\pgfpathlineto{\pgfqpoint{1.284286in}{0.888766in}}%
\pgfpathlineto{\pgfqpoint{1.289557in}{0.891908in}}%
\pgfpathlineto{\pgfqpoint{1.296319in}{0.895988in}}%
\pgfpathlineto{\pgfqpoint{1.297806in}{0.896896in}}%
\pgfpathlineto{\pgfqpoint{1.306055in}{0.901987in}}%
\pgfpathlineto{\pgfqpoint{1.308016in}{0.903210in}}%
\pgfpathlineto{\pgfqpoint{1.314304in}{0.907182in}}%
\pgfpathlineto{\pgfqpoint{1.319389in}{0.910432in}}%
\pgfpathlineto{\pgfqpoint{1.322553in}{0.912480in}}%
\pgfpathlineto{\pgfqpoint{1.330455in}{0.917654in}}%
\pgfpathlineto{\pgfqpoint{1.330803in}{0.917885in}}%
\pgfpathlineto{\pgfqpoint{1.339052in}{0.923409in}}%
\pgfpathlineto{\pgfqpoint{1.341218in}{0.924877in}}%
\pgfpathlineto{\pgfqpoint{1.347301in}{0.929050in}}%
\pgfpathlineto{\pgfqpoint{1.351693in}{0.932099in}}%
\pgfpathlineto{\pgfqpoint{1.355550in}{0.934810in}}%
\pgfpathlineto{\pgfqpoint{1.361892in}{0.939321in}}%
\pgfpathlineto{\pgfqpoint{1.363800in}{0.940695in}}%
\pgfpathlineto{\pgfqpoint{1.371823in}{0.946543in}}%
\pgfpathlineto{\pgfqpoint{1.372049in}{0.946710in}}%
\pgfpathlineto{\pgfqpoint{1.380298in}{0.952864in}}%
\pgfpathlineto{\pgfqpoint{1.381493in}{0.953766in}}%
\pgfpathlineto{\pgfqpoint{1.388547in}{0.959158in}}%
\pgfpathlineto{\pgfqpoint{1.390914in}{0.960988in}}%
\pgfpathlineto{\pgfqpoint{1.396797in}{0.965596in}}%
\pgfpathlineto{\pgfqpoint{1.400095in}{0.968210in}}%
\pgfpathlineto{\pgfqpoint{1.405046in}{0.972183in}}%
\pgfpathlineto{\pgfqpoint{1.409047in}{0.975432in}}%
\pgfpathlineto{\pgfqpoint{1.413295in}{0.978926in}}%
\pgfpathlineto{\pgfqpoint{1.417777in}{0.982654in}}%
\pgfpathlineto{\pgfqpoint{1.421544in}{0.985828in}}%
\pgfpathlineto{\pgfqpoint{1.426295in}{0.989877in}}%
\pgfpathlineto{\pgfqpoint{1.429793in}{0.992897in}}%
\pgfpathlineto{\pgfqpoint{1.434607in}{0.997099in}}%
\pgfpathlineto{\pgfqpoint{1.438043in}{1.000137in}}%
\pgfpathlineto{\pgfqpoint{1.442721in}{1.004321in}}%
\pgfpathlineto{\pgfqpoint{1.446292in}{1.007556in}}%
\pgfpathlineto{\pgfqpoint{1.450644in}{1.011543in}}%
\pgfpathlineto{\pgfqpoint{1.454541in}{1.015160in}}%
\pgfpathlineto{\pgfqpoint{1.458384in}{1.018766in}}%
\pgfpathlineto{\pgfqpoint{1.462790in}{1.022955in}}%
\pgfpathlineto{\pgfqpoint{1.465946in}{1.025988in}}%
\pgfpathlineto{\pgfqpoint{1.471040in}{1.030948in}}%
\pgfpathlineto{\pgfqpoint{1.473337in}{1.033210in}}%
\pgfpathlineto{\pgfqpoint{1.479289in}{1.039147in}}%
\pgfpathlineto{\pgfqpoint{1.480563in}{1.040432in}}%
\pgfpathlineto{\pgfqpoint{1.487538in}{1.047560in}}%
\pgfpathlineto{\pgfqpoint{1.487629in}{1.047654in}}%
\pgfpathlineto{\pgfqpoint{1.494537in}{1.054877in}}%
\pgfpathlineto{\pgfqpoint{1.495787in}{1.056201in}}%
\pgfpathlineto{\pgfqpoint{1.501294in}{1.062099in}}%
\pgfpathlineto{\pgfqpoint{1.504036in}{1.065076in}}%
\pgfpathlineto{\pgfqpoint{1.507906in}{1.069321in}}%
\pgfpathlineto{\pgfqpoint{1.512286in}{1.074192in}}%
\pgfpathlineto{\pgfqpoint{1.514377in}{1.076543in}}%
\pgfpathlineto{\pgfqpoint{1.520535in}{1.083561in}}%
\pgfpathlineto{\pgfqpoint{1.520712in}{1.083766in}}%
\pgfpathlineto{\pgfqpoint{1.526909in}{1.090988in}}%
\pgfpathlineto{\pgfqpoint{1.528784in}{1.093204in}}%
\pgfpathlineto{\pgfqpoint{1.532976in}{1.098210in}}%
\pgfpathlineto{\pgfqpoint{1.537033in}{1.103122in}}%
\pgfpathlineto{\pgfqpoint{1.538921in}{1.105432in}}%
\pgfpathlineto{\pgfqpoint{1.544743in}{1.112654in}}%
\pgfpathlineto{\pgfqpoint{1.545283in}{1.113332in}}%
\pgfpathlineto{\pgfqpoint{1.550442in}{1.119877in}}%
\pgfpathlineto{\pgfqpoint{1.553532in}{1.123852in}}%
\pgfpathlineto{\pgfqpoint{1.556028in}{1.127099in}}%
\pgfpathlineto{\pgfqpoint{1.561505in}{1.134321in}}%
\pgfpathlineto{\pgfqpoint{1.561781in}{1.134689in}}%
\pgfpathlineto{\pgfqpoint{1.566866in}{1.141543in}}%
\pgfpathlineto{\pgfqpoint{1.570030in}{1.145869in}}%
\pgfpathlineto{\pgfqpoint{1.572126in}{1.148766in}}%
\pgfpathlineto{\pgfqpoint{1.577283in}{1.155988in}}%
\pgfpathlineto{\pgfqpoint{1.578279in}{1.157400in}}%
\pgfpathlineto{\pgfqpoint{1.582335in}{1.163210in}}%
\pgfpathlineto{\pgfqpoint{1.586529in}{1.169305in}}%
\pgfpathlineto{\pgfqpoint{1.587296in}{1.170432in}}%
\pgfpathlineto{\pgfqpoint{1.592156in}{1.177654in}}%
\pgfpathlineto{\pgfqpoint{1.594778in}{1.181603in}}%
\pgfpathlineto{\pgfqpoint{1.596928in}{1.184877in}}%
\pgfpathlineto{\pgfqpoint{1.601610in}{1.192099in}}%
\pgfpathlineto{\pgfqpoint{1.603027in}{1.194310in}}%
\pgfpathlineto{\pgfqpoint{1.606203in}{1.199321in}}%
\pgfpathlineto{\pgfqpoint{1.610716in}{1.206543in}}%
\pgfpathlineto{\pgfqpoint{1.611276in}{1.207450in}}%
\pgfpathlineto{\pgfqpoint{1.615140in}{1.213766in}}%
\pgfpathlineto{\pgfqpoint{1.619491in}{1.220988in}}%
\pgfpathlineto{\pgfqpoint{1.619526in}{1.221045in}}%
\pgfpathlineto{\pgfqpoint{1.623756in}{1.228210in}}%
\pgfpathlineto{\pgfqpoint{1.627775in}{1.235122in}}%
\pgfpathlineto{\pgfqpoint{1.627953in}{1.235432in}}%
\pgfpathlineto{\pgfqpoint{1.632068in}{1.242654in}}%
\pgfpathlineto{\pgfqpoint{1.636024in}{1.249704in}}%
\pgfpathlineto{\pgfqpoint{1.636120in}{1.249877in}}%
\pgfpathlineto{\pgfqpoint{1.640092in}{1.257099in}}%
\pgfpathlineto{\pgfqpoint{1.644004in}{1.264321in}}%
\pgfpathlineto{\pgfqpoint{1.644273in}{1.264822in}}%
\pgfpathlineto{\pgfqpoint{1.647842in}{1.271543in}}%
\pgfpathlineto{\pgfqpoint{1.651621in}{1.278766in}}%
\pgfpathlineto{\pgfqpoint{1.652522in}{1.280505in}}%
\pgfpathlineto{\pgfqpoint{1.655332in}{1.285988in}}%
\pgfpathlineto{\pgfqpoint{1.658985in}{1.293210in}}%
\pgfpathlineto{\pgfqpoint{1.660772in}{1.296785in}}%
\pgfpathlineto{\pgfqpoint{1.662576in}{1.300432in}}%
\pgfpathlineto{\pgfqpoint{1.666107in}{1.307654in}}%
\pgfpathlineto{\pgfqpoint{1.669021in}{1.313693in}}%
\pgfpathlineto{\pgfqpoint{1.669586in}{1.314877in}}%
\pgfpathlineto{\pgfqpoint{1.673001in}{1.322099in}}%
\pgfpathlineto{\pgfqpoint{1.676369in}{1.329321in}}%
\pgfpathlineto{\pgfqpoint{1.677270in}{1.331272in}}%
\pgfpathlineto{\pgfqpoint{1.679679in}{1.336543in}}%
\pgfpathlineto{\pgfqpoint{1.682938in}{1.343766in}}%
\pgfpathlineto{\pgfqpoint{1.685519in}{1.349558in}}%
\pgfpathlineto{\pgfqpoint{1.686150in}{1.350988in}}%
\pgfpathlineto{\pgfqpoint{1.689305in}{1.358210in}}%
\pgfpathlineto{\pgfqpoint{1.692418in}{1.365432in}}%
\pgfpathlineto{\pgfqpoint{1.693769in}{1.368599in}}%
\pgfpathlineto{\pgfqpoint{1.695480in}{1.372654in}}%
\pgfpathlineto{\pgfqpoint{1.698495in}{1.379877in}}%
\pgfpathlineto{\pgfqpoint{1.701471in}{1.387099in}}%
\pgfpathlineto{\pgfqpoint{1.702018in}{1.388437in}}%
\pgfpathlineto{\pgfqpoint{1.704395in}{1.394321in}}%
\pgfpathlineto{\pgfqpoint{1.707279in}{1.401543in}}%
\pgfpathlineto{\pgfqpoint{1.710126in}{1.408766in}}%
\pgfpathlineto{\pgfqpoint{1.710267in}{1.409126in}}%
\pgfpathlineto{\pgfqpoint{1.712922in}{1.415988in}}%
\pgfpathlineto{\pgfqpoint{1.715682in}{1.423210in}}%
\pgfpathlineto{\pgfqpoint{1.718409in}{1.430432in}}%
\pgfpathlineto{\pgfqpoint{1.718516in}{1.430719in}}%
\pgfpathlineto{\pgfqpoint{1.721086in}{1.437654in}}%
\pgfpathlineto{\pgfqpoint{1.723730in}{1.444877in}}%
\pgfpathlineto{\pgfqpoint{1.726343in}{1.452099in}}%
\pgfpathlineto{\pgfqpoint{1.726765in}{1.453276in}}%
\pgfpathlineto{\pgfqpoint{1.728911in}{1.459321in}}%
\pgfpathlineto{\pgfqpoint{1.731446in}{1.466543in}}%
\pgfpathlineto{\pgfqpoint{1.733952in}{1.473766in}}%
\pgfpathlineto{\pgfqpoint{1.735015in}{1.476856in}}%
\pgfpathlineto{\pgfqpoint{1.736420in}{1.480988in}}%
\pgfpathlineto{\pgfqpoint{1.738852in}{1.488210in}}%
\pgfpathlineto{\pgfqpoint{1.741257in}{1.495432in}}%
\pgfpathlineto{\pgfqpoint{1.743264in}{1.501523in}}%
\pgfpathlineto{\pgfqpoint{1.743633in}{1.502654in}}%
\pgfpathlineto{\pgfqpoint{1.745968in}{1.509877in}}%
\pgfpathlineto{\pgfqpoint{1.748278in}{1.517099in}}%
\pgfpathlineto{\pgfqpoint{1.750563in}{1.524321in}}%
\pgfpathlineto{\pgfqpoint{1.751513in}{1.527350in}}%
\pgfpathlineto{\pgfqpoint{1.752814in}{1.531543in}}%
\pgfpathlineto{\pgfqpoint{1.755034in}{1.538766in}}%
\pgfpathlineto{\pgfqpoint{1.757231in}{1.545988in}}%
\pgfpathlineto{\pgfqpoint{1.759405in}{1.553210in}}%
\pgfpathlineto{\pgfqpoint{1.759762in}{1.554407in}}%
\pgfpathlineto{\pgfqpoint{1.761543in}{1.560432in}}%
\pgfpathlineto{\pgfqpoint{1.763657in}{1.567654in}}%
\pgfpathlineto{\pgfqpoint{1.765749in}{1.574877in}}%
\pgfpathlineto{\pgfqpoint{1.767820in}{1.582099in}}%
\pgfpathlineto{\pgfqpoint{1.768012in}{1.582772in}}%
\pgfpathlineto{\pgfqpoint{1.769857in}{1.589321in}}%
\pgfpathlineto{\pgfqpoint{1.771872in}{1.596543in}}%
\pgfpathlineto{\pgfqpoint{1.773867in}{1.603766in}}%
\pgfpathlineto{\pgfqpoint{1.775843in}{1.610988in}}%
\pgfpathlineto{\pgfqpoint{1.776261in}{1.612524in}}%
\pgfpathlineto{\pgfqpoint{1.777789in}{1.618210in}}%
\pgfpathlineto{\pgfqpoint{1.779713in}{1.625432in}}%
\pgfpathlineto{\pgfqpoint{1.781619in}{1.632654in}}%
\pgfpathlineto{\pgfqpoint{1.783507in}{1.639877in}}%
\pgfpathlineto{\pgfqpoint{1.784510in}{1.643741in}}%
\pgfpathlineto{\pgfqpoint{1.785371in}{1.647099in}}%
\pgfpathlineto{\pgfqpoint{1.787211in}{1.654321in}}%
\pgfpathlineto{\pgfqpoint{1.789034in}{1.661543in}}%
\pgfpathlineto{\pgfqpoint{1.790841in}{1.668766in}}%
\pgfpathlineto{\pgfqpoint{1.792633in}{1.675988in}}%
\pgfpathlineto{\pgfqpoint{1.792759in}{1.676501in}}%
\pgfpathlineto{\pgfqpoint{1.794396in}{1.683210in}}%
\pgfpathlineto{\pgfqpoint{1.796142in}{1.690432in}}%
\pgfpathlineto{\pgfqpoint{1.797874in}{1.697654in}}%
\pgfpathlineto{\pgfqpoint{1.799591in}{1.704877in}}%
\pgfpathlineto{\pgfqpoint{1.801008in}{1.710882in}}%
\pgfpathlineto{\pgfqpoint{1.801292in}{1.712099in}}%
\pgfpathlineto{\pgfqpoint{1.802968in}{1.719321in}}%
\pgfpathlineto{\pgfqpoint{1.804630in}{1.726543in}}%
\pgfpathlineto{\pgfqpoint{1.806279in}{1.733766in}}%
\pgfpathlineto{\pgfqpoint{1.807915in}{1.740988in}}%
\pgfpathlineto{\pgfqpoint{1.809258in}{1.746957in}}%
\pgfpathlineto{\pgfqpoint{1.809536in}{1.748210in}}%
\pgfpathlineto{\pgfqpoint{1.811134in}{1.755432in}}%
\pgfpathlineto{\pgfqpoint{1.812719in}{1.762654in}}%
\pgfpathlineto{\pgfqpoint{1.814293in}{1.769877in}}%
\pgfpathlineto{\pgfqpoint{1.815855in}{1.777099in}}%
\pgfpathlineto{\pgfqpoint{1.817406in}{1.784321in}}%
\pgfpathlineto{\pgfqpoint{1.817507in}{1.784794in}}%
\pgfpathlineto{\pgfqpoint{1.818934in}{1.791543in}}%
\pgfpathlineto{\pgfqpoint{1.820450in}{1.798766in}}%
\pgfpathlineto{\pgfqpoint{1.821955in}{1.805988in}}%
\pgfpathlineto{\pgfqpoint{1.823450in}{1.813210in}}%
\pgfpathlineto{\pgfqpoint{1.824935in}{1.820432in}}%
\pgfpathlineto{\pgfqpoint{1.825756in}{1.824449in}}%
\pgfpathlineto{\pgfqpoint{1.826404in}{1.827654in}}%
\pgfpathlineto{\pgfqpoint{1.827857in}{1.834877in}}%
\pgfpathlineto{\pgfqpoint{1.829300in}{1.842099in}}%
\pgfpathlineto{\pgfqpoint{1.830734in}{1.849321in}}%
\pgfpathlineto{\pgfqpoint{1.832159in}{1.856543in}}%
\pgfpathlineto{\pgfqpoint{1.833575in}{1.863766in}}%
\pgfpathlineto{\pgfqpoint{1.834005in}{1.865967in}}%
\pgfpathlineto{\pgfqpoint{1.834974in}{1.870988in}}%
\pgfpathlineto{\pgfqpoint{1.836361in}{1.878210in}}%
\pgfpathlineto{\pgfqpoint{1.837740in}{1.885432in}}%
\pgfpathlineto{\pgfqpoint{1.839111in}{1.892654in}}%
\pgfpathlineto{\pgfqpoint{1.840474in}{1.899877in}}%
\pgfpathlineto{\pgfqpoint{1.834005in}{1.899877in}}%
\pgfpathlineto{\pgfqpoint{1.825756in}{1.899877in}}%
\pgfpathlineto{\pgfqpoint{1.817507in}{1.899877in}}%
\pgfpathlineto{\pgfqpoint{1.809258in}{1.899877in}}%
\pgfpathlineto{\pgfqpoint{1.801008in}{1.899877in}}%
\pgfpathlineto{\pgfqpoint{1.792759in}{1.899877in}}%
\pgfpathlineto{\pgfqpoint{1.784510in}{1.899877in}}%
\pgfpathlineto{\pgfqpoint{1.776261in}{1.899877in}}%
\pgfpathlineto{\pgfqpoint{1.768012in}{1.899877in}}%
\pgfpathlineto{\pgfqpoint{1.759762in}{1.899877in}}%
\pgfpathlineto{\pgfqpoint{1.751513in}{1.899877in}}%
\pgfpathlineto{\pgfqpoint{1.743264in}{1.899877in}}%
\pgfpathlineto{\pgfqpoint{1.735015in}{1.899877in}}%
\pgfpathlineto{\pgfqpoint{1.726765in}{1.899877in}}%
\pgfpathlineto{\pgfqpoint{1.718516in}{1.899877in}}%
\pgfpathlineto{\pgfqpoint{1.710267in}{1.899877in}}%
\pgfpathlineto{\pgfqpoint{1.702018in}{1.899877in}}%
\pgfpathlineto{\pgfqpoint{1.693769in}{1.899877in}}%
\pgfpathlineto{\pgfqpoint{1.685519in}{1.899877in}}%
\pgfpathlineto{\pgfqpoint{1.677270in}{1.899877in}}%
\pgfpathlineto{\pgfqpoint{1.669021in}{1.899877in}}%
\pgfpathlineto{\pgfqpoint{1.660772in}{1.899877in}}%
\pgfpathlineto{\pgfqpoint{1.652522in}{1.899877in}}%
\pgfpathlineto{\pgfqpoint{1.644273in}{1.899877in}}%
\pgfpathlineto{\pgfqpoint{1.636024in}{1.899877in}}%
\pgfpathlineto{\pgfqpoint{1.630028in}{1.899877in}}%
\pgfpathlineto{\pgfqpoint{1.627775in}{1.892941in}}%
\pgfpathlineto{\pgfqpoint{1.627681in}{1.892654in}}%
\pgfpathlineto{\pgfqpoint{1.625299in}{1.885432in}}%
\pgfpathlineto{\pgfqpoint{1.622905in}{1.878210in}}%
\pgfpathlineto{\pgfqpoint{1.620498in}{1.870988in}}%
\pgfpathlineto{\pgfqpoint{1.619526in}{1.868083in}}%
\pgfpathlineto{\pgfqpoint{1.618064in}{1.863766in}}%
\pgfpathlineto{\pgfqpoint{1.615607in}{1.856543in}}%
\pgfpathlineto{\pgfqpoint{1.613137in}{1.849321in}}%
\pgfpathlineto{\pgfqpoint{1.611276in}{1.843911in}}%
\pgfpathlineto{\pgfqpoint{1.610646in}{1.842099in}}%
\pgfpathlineto{\pgfqpoint{1.608124in}{1.834877in}}%
\pgfpathlineto{\pgfqpoint{1.605587in}{1.827654in}}%
\pgfpathlineto{\pgfqpoint{1.603034in}{1.820432in}}%
\pgfpathlineto{\pgfqpoint{1.603027in}{1.820412in}}%
\pgfpathlineto{\pgfqpoint{1.600444in}{1.813210in}}%
\pgfpathlineto{\pgfqpoint{1.597837in}{1.805988in}}%
\pgfpathlineto{\pgfqpoint{1.595214in}{1.798766in}}%
\pgfpathlineto{\pgfqpoint{1.594778in}{1.797571in}}%
\pgfpathlineto{\pgfqpoint{1.592554in}{1.791543in}}%
\pgfpathlineto{\pgfqpoint{1.589874in}{1.784321in}}%
\pgfpathlineto{\pgfqpoint{1.587177in}{1.777099in}}%
\pgfpathlineto{\pgfqpoint{1.586529in}{1.775371in}}%
\pgfpathlineto{\pgfqpoint{1.584444in}{1.769877in}}%
\pgfpathlineto{\pgfqpoint{1.581688in}{1.762654in}}%
\pgfpathlineto{\pgfqpoint{1.578912in}{1.755432in}}%
\pgfpathlineto{\pgfqpoint{1.578279in}{1.753793in}}%
\pgfpathlineto{\pgfqpoint{1.576100in}{1.748210in}}%
\pgfpathlineto{\pgfqpoint{1.573263in}{1.740988in}}%
\pgfpathlineto{\pgfqpoint{1.570406in}{1.733766in}}%
\pgfpathlineto{\pgfqpoint{1.570030in}{1.732820in}}%
\pgfpathlineto{\pgfqpoint{1.567508in}{1.726543in}}%
\pgfpathlineto{\pgfqpoint{1.564586in}{1.719321in}}%
\pgfpathlineto{\pgfqpoint{1.561781in}{1.712434in}}%
\pgfpathlineto{\pgfqpoint{1.561643in}{1.712099in}}%
\pgfpathlineto{\pgfqpoint{1.558655in}{1.704877in}}%
\pgfpathlineto{\pgfqpoint{1.555644in}{1.697654in}}%
\pgfpathlineto{\pgfqpoint{1.553532in}{1.692619in}}%
\pgfpathlineto{\pgfqpoint{1.552605in}{1.690432in}}%
\pgfpathlineto{\pgfqpoint{1.549525in}{1.683210in}}%
\pgfpathlineto{\pgfqpoint{1.546422in}{1.675988in}}%
\pgfpathlineto{\pgfqpoint{1.545283in}{1.673351in}}%
\pgfpathlineto{\pgfqpoint{1.543280in}{1.668766in}}%
\pgfpathlineto{\pgfqpoint{1.540104in}{1.661543in}}%
\pgfpathlineto{\pgfqpoint{1.537033in}{1.654612in}}%
\pgfpathlineto{\pgfqpoint{1.536903in}{1.654321in}}%
\pgfpathlineto{\pgfqpoint{1.533653in}{1.647099in}}%
\pgfpathlineto{\pgfqpoint{1.530377in}{1.639877in}}%
\pgfpathlineto{\pgfqpoint{1.528784in}{1.636389in}}%
\pgfpathlineto{\pgfqpoint{1.527061in}{1.632654in}}%
\pgfpathlineto{\pgfqpoint{1.523707in}{1.625432in}}%
\pgfpathlineto{\pgfqpoint{1.520535in}{1.618657in}}%
\pgfpathlineto{\pgfqpoint{1.520323in}{1.618210in}}%
\pgfpathlineto{\pgfqpoint{1.516888in}{1.610988in}}%
\pgfpathlineto{\pgfqpoint{1.513425in}{1.603766in}}%
\pgfpathlineto{\pgfqpoint{1.512286in}{1.601406in}}%
\pgfpathlineto{\pgfqpoint{1.509915in}{1.596543in}}%
\pgfpathlineto{\pgfqpoint{1.506367in}{1.589321in}}%
\pgfpathlineto{\pgfqpoint{1.504036in}{1.584612in}}%
\pgfpathlineto{\pgfqpoint{1.502780in}{1.582099in}}%
\pgfpathlineto{\pgfqpoint{1.499145in}{1.574877in}}%
\pgfpathlineto{\pgfqpoint{1.495787in}{1.568260in}}%
\pgfpathlineto{\pgfqpoint{1.495477in}{1.567654in}}%
\pgfpathlineto{\pgfqpoint{1.491752in}{1.560432in}}%
\pgfpathlineto{\pgfqpoint{1.487994in}{1.553210in}}%
\pgfpathlineto{\pgfqpoint{1.487538in}{1.552338in}}%
\pgfpathlineto{\pgfqpoint{1.484180in}{1.545988in}}%
\pgfpathlineto{\pgfqpoint{1.480329in}{1.538766in}}%
\pgfpathlineto{\pgfqpoint{1.479289in}{1.536829in}}%
\pgfpathlineto{\pgfqpoint{1.476423in}{1.531543in}}%
\pgfpathlineto{\pgfqpoint{1.472475in}{1.524321in}}%
\pgfpathlineto{\pgfqpoint{1.471040in}{1.521715in}}%
\pgfpathlineto{\pgfqpoint{1.468473in}{1.517099in}}%
\pgfpathlineto{\pgfqpoint{1.464424in}{1.509877in}}%
\pgfpathlineto{\pgfqpoint{1.462790in}{1.506985in}}%
\pgfpathlineto{\pgfqpoint{1.460320in}{1.502654in}}%
\pgfpathlineto{\pgfqpoint{1.456167in}{1.495432in}}%
\pgfpathlineto{\pgfqpoint{1.454541in}{1.492626in}}%
\pgfpathlineto{\pgfqpoint{1.451958in}{1.488210in}}%
\pgfpathlineto{\pgfqpoint{1.447697in}{1.480988in}}%
\pgfpathlineto{\pgfqpoint{1.446292in}{1.478625in}}%
\pgfpathlineto{\pgfqpoint{1.443375in}{1.473766in}}%
\pgfpathlineto{\pgfqpoint{1.439002in}{1.466543in}}%
\pgfpathlineto{\pgfqpoint{1.438043in}{1.464970in}}%
\pgfpathlineto{\pgfqpoint{1.434564in}{1.459321in}}%
\pgfpathlineto{\pgfqpoint{1.430075in}{1.452099in}}%
\pgfpathlineto{\pgfqpoint{1.429793in}{1.451650in}}%
\pgfpathlineto{\pgfqpoint{1.425513in}{1.444877in}}%
\pgfpathlineto{\pgfqpoint{1.421544in}{1.438656in}}%
\pgfpathlineto{\pgfqpoint{1.420899in}{1.437654in}}%
\pgfpathlineto{\pgfqpoint{1.416212in}{1.430432in}}%
\pgfpathlineto{\pgfqpoint{1.413295in}{1.425978in}}%
\pgfpathlineto{\pgfqpoint{1.411465in}{1.423210in}}%
\pgfpathlineto{\pgfqpoint{1.406650in}{1.415988in}}%
\pgfpathlineto{\pgfqpoint{1.405046in}{1.413601in}}%
\pgfpathlineto{\pgfqpoint{1.401764in}{1.408766in}}%
\pgfpathlineto{\pgfqpoint{1.396816in}{1.401543in}}%
\pgfpathlineto{\pgfqpoint{1.396797in}{1.401515in}}%
\pgfpathlineto{\pgfqpoint{1.391784in}{1.394321in}}%
\pgfpathlineto{\pgfqpoint{1.388547in}{1.389721in}}%
\pgfpathlineto{\pgfqpoint{1.386685in}{1.387099in}}%
\pgfpathlineto{\pgfqpoint{1.381511in}{1.379877in}}%
\pgfpathlineto{\pgfqpoint{1.380298in}{1.378197in}}%
\pgfpathlineto{\pgfqpoint{1.376256in}{1.372654in}}%
\pgfpathlineto{\pgfqpoint{1.372049in}{1.366941in}}%
\pgfpathlineto{\pgfqpoint{1.370928in}{1.365432in}}%
\pgfpathlineto{\pgfqpoint{1.365514in}{1.358210in}}%
\pgfpathlineto{\pgfqpoint{1.363800in}{1.355943in}}%
\pgfpathlineto{\pgfqpoint{1.360017in}{1.350988in}}%
\pgfpathlineto{\pgfqpoint{1.355550in}{1.345192in}}%
\pgfpathlineto{\pgfqpoint{1.354441in}{1.343766in}}%
\pgfpathlineto{\pgfqpoint{1.348775in}{1.336543in}}%
\pgfpathlineto{\pgfqpoint{1.347301in}{1.334682in}}%
\pgfpathlineto{\pgfqpoint{1.343019in}{1.329321in}}%
\pgfpathlineto{\pgfqpoint{1.339052in}{1.324403in}}%
\pgfpathlineto{\pgfqpoint{1.337176in}{1.322099in}}%
\pgfpathlineto{\pgfqpoint{1.331243in}{1.314877in}}%
\pgfpathlineto{\pgfqpoint{1.330803in}{1.314345in}}%
\pgfpathlineto{\pgfqpoint{1.325209in}{1.307654in}}%
\pgfpathlineto{\pgfqpoint{1.322553in}{1.304508in}}%
\pgfpathlineto{\pgfqpoint{1.319082in}{1.300432in}}%
\pgfpathlineto{\pgfqpoint{1.314304in}{1.294876in}}%
\pgfpathlineto{\pgfqpoint{1.312859in}{1.293210in}}%
\pgfpathlineto{\pgfqpoint{1.306535in}{1.285988in}}%
\pgfpathlineto{\pgfqpoint{1.306055in}{1.285444in}}%
\pgfpathlineto{\pgfqpoint{1.300103in}{1.278766in}}%
\pgfpathlineto{\pgfqpoint{1.297806in}{1.276212in}}%
\pgfpathlineto{\pgfqpoint{1.293568in}{1.271543in}}%
\pgfpathlineto{\pgfqpoint{1.289557in}{1.267167in}}%
\pgfpathlineto{\pgfqpoint{1.286926in}{1.264321in}}%
\pgfpathlineto{\pgfqpoint{1.281307in}{1.258301in}}%
\pgfpathlineto{\pgfqpoint{1.280176in}{1.257099in}}%
\pgfpathlineto{\pgfqpoint{1.273312in}{1.249877in}}%
\pgfpathlineto{\pgfqpoint{1.273058in}{1.249612in}}%
\pgfpathlineto{\pgfqpoint{1.266328in}{1.242654in}}%
\pgfpathlineto{\pgfqpoint{1.264809in}{1.241099in}}%
\pgfpathlineto{\pgfqpoint{1.259226in}{1.235432in}}%
\pgfpathlineto{\pgfqpoint{1.256560in}{1.232751in}}%
\pgfpathlineto{\pgfqpoint{1.252003in}{1.228210in}}%
\pgfpathlineto{\pgfqpoint{1.248310in}{1.224564in}}%
\pgfpathlineto{\pgfqpoint{1.244656in}{1.220988in}}%
\pgfpathlineto{\pgfqpoint{1.240061in}{1.216534in}}%
\pgfpathlineto{\pgfqpoint{1.237180in}{1.213766in}}%
\pgfpathlineto{\pgfqpoint{1.231812in}{1.208656in}}%
\pgfpathlineto{\pgfqpoint{1.229573in}{1.206543in}}%
\pgfpathlineto{\pgfqpoint{1.223563in}{1.200927in}}%
\pgfpathlineto{\pgfqpoint{1.221830in}{1.199321in}}%
\pgfpathlineto{\pgfqpoint{1.215314in}{1.193342in}}%
\pgfpathlineto{\pgfqpoint{1.213947in}{1.192099in}}%
\pgfpathlineto{\pgfqpoint{1.207064in}{1.185897in}}%
\pgfpathlineto{\pgfqpoint{1.205922in}{1.184877in}}%
\pgfpathlineto{\pgfqpoint{1.198815in}{1.178590in}}%
\pgfpathlineto{\pgfqpoint{1.197749in}{1.177654in}}%
\pgfpathlineto{\pgfqpoint{1.190566in}{1.171416in}}%
\pgfpathlineto{\pgfqpoint{1.189424in}{1.170432in}}%
\pgfpathlineto{\pgfqpoint{1.182317in}{1.164371in}}%
\pgfpathlineto{\pgfqpoint{1.180944in}{1.163210in}}%
\pgfpathlineto{\pgfqpoint{1.174067in}{1.157453in}}%
\pgfpathlineto{\pgfqpoint{1.172303in}{1.155988in}}%
\pgfpathlineto{\pgfqpoint{1.165818in}{1.150656in}}%
\pgfpathlineto{\pgfqpoint{1.163499in}{1.148766in}}%
\pgfpathlineto{\pgfqpoint{1.157569in}{1.143979in}}%
\pgfpathlineto{\pgfqpoint{1.154527in}{1.141543in}}%
\pgfpathlineto{\pgfqpoint{1.149320in}{1.137416in}}%
\pgfpathlineto{\pgfqpoint{1.145382in}{1.134321in}}%
\pgfpathlineto{\pgfqpoint{1.141071in}{1.130965in}}%
\pgfpathlineto{\pgfqpoint{1.136061in}{1.127099in}}%
\pgfpathlineto{\pgfqpoint{1.132821in}{1.124623in}}%
\pgfpathlineto{\pgfqpoint{1.126560in}{1.119877in}}%
\pgfpathlineto{\pgfqpoint{1.124572in}{1.118385in}}%
\pgfpathlineto{\pgfqpoint{1.116873in}{1.112654in}}%
\pgfpathlineto{\pgfqpoint{1.116323in}{1.112249in}}%
\pgfpathlineto{\pgfqpoint{1.108074in}{1.106217in}}%
\pgfpathlineto{\pgfqpoint{1.106992in}{1.105432in}}%
\pgfpathlineto{\pgfqpoint{1.099824in}{1.100284in}}%
\pgfpathlineto{\pgfqpoint{1.096914in}{1.098210in}}%
\pgfpathlineto{\pgfqpoint{1.091575in}{1.094443in}}%
\pgfpathlineto{\pgfqpoint{1.086637in}{1.090988in}}%
\pgfpathlineto{\pgfqpoint{1.083326in}{1.088694in}}%
\pgfpathlineto{\pgfqpoint{1.076155in}{1.083766in}}%
\pgfpathlineto{\pgfqpoint{1.075077in}{1.083032in}}%
\pgfpathlineto{\pgfqpoint{1.066828in}{1.077463in}}%
\pgfpathlineto{\pgfqpoint{1.065455in}{1.076543in}}%
\pgfpathlineto{\pgfqpoint{1.058578in}{1.071982in}}%
\pgfpathlineto{\pgfqpoint{1.054533in}{1.069321in}}%
\pgfpathlineto{\pgfqpoint{1.050329in}{1.066583in}}%
\pgfpathlineto{\pgfqpoint{1.043385in}{1.062099in}}%
\pgfpathlineto{\pgfqpoint{1.042080in}{1.061264in}}%
\pgfpathlineto{\pgfqpoint{1.033831in}{1.056031in}}%
\pgfpathlineto{\pgfqpoint{1.031997in}{1.054877in}}%
\pgfpathlineto{\pgfqpoint{1.025581in}{1.050879in}}%
\pgfpathlineto{\pgfqpoint{1.020362in}{1.047654in}}%
\pgfpathlineto{\pgfqpoint{1.017332in}{1.045802in}}%
\pgfpathlineto{\pgfqpoint{1.009083in}{1.040799in}}%
\pgfpathlineto{\pgfqpoint{1.008474in}{1.040432in}}%
\pgfpathlineto{\pgfqpoint{1.000834in}{1.035879in}}%
\pgfpathlineto{\pgfqpoint{0.996315in}{1.033210in}}%
\pgfpathlineto{\pgfqpoint{0.992585in}{1.031028in}}%
\pgfpathlineto{\pgfqpoint{0.984335in}{1.026246in}}%
\pgfpathlineto{\pgfqpoint{0.983887in}{1.025988in}}%
\pgfpathlineto{\pgfqpoint{0.976086in}{1.021544in}}%
\pgfpathlineto{\pgfqpoint{0.971167in}{1.018766in}}%
\pgfpathlineto{\pgfqpoint{0.967837in}{1.016904in}}%
\pgfpathlineto{\pgfqpoint{0.959588in}{1.012330in}}%
\pgfpathlineto{\pgfqpoint{0.958159in}{1.011543in}}%
\pgfpathlineto{\pgfqpoint{0.951338in}{1.007828in}}%
\pgfpathlineto{\pgfqpoint{0.944846in}{1.004321in}}%
\pgfpathlineto{\pgfqpoint{0.943089in}{1.003382in}}%
\pgfpathlineto{\pgfqpoint{0.934840in}{0.999004in}}%
\pgfpathlineto{\pgfqpoint{0.931222in}{0.997099in}}%
\pgfpathlineto{\pgfqpoint{0.926591in}{0.994686in}}%
\pgfpathlineto{\pgfqpoint{0.918342in}{0.990422in}}%
\pgfpathlineto{\pgfqpoint{0.917279in}{0.989877in}}%
\pgfpathlineto{\pgfqpoint{0.910092in}{0.986225in}}%
\pgfpathlineto{\pgfqpoint{0.903005in}{0.982654in}}%
\pgfpathlineto{\pgfqpoint{0.901843in}{0.982075in}}%
\pgfpathlineto{\pgfqpoint{0.893594in}{0.977990in}}%
\pgfpathlineto{\pgfqpoint{0.888388in}{0.975432in}}%
\pgfpathlineto{\pgfqpoint{0.885345in}{0.973953in}}%
\pgfpathlineto{\pgfqpoint{0.877095in}{0.969971in}}%
\pgfpathlineto{\pgfqpoint{0.873420in}{0.968210in}}%
\pgfpathlineto{\pgfqpoint{0.868846in}{0.966041in}}%
\pgfpathlineto{\pgfqpoint{0.860597in}{0.962160in}}%
\pgfpathlineto{\pgfqpoint{0.858088in}{0.960988in}}%
\pgfpathlineto{\pgfqpoint{0.852348in}{0.958334in}}%
\pgfpathlineto{\pgfqpoint{0.844099in}{0.954551in}}%
\pgfpathlineto{\pgfqpoint{0.842375in}{0.953766in}}%
\pgfpathlineto{\pgfqpoint{0.835849in}{0.950824in}}%
\pgfpathlineto{\pgfqpoint{0.827600in}{0.947136in}}%
\pgfpathlineto{\pgfqpoint{0.826267in}{0.946543in}}%
\pgfpathlineto{\pgfqpoint{0.819351in}{0.943504in}}%
\pgfpathlineto{\pgfqpoint{0.811102in}{0.939908in}}%
\pgfpathlineto{\pgfqpoint{0.809747in}{0.939321in}}%
\pgfpathlineto{\pgfqpoint{0.802852in}{0.936367in}}%
\pgfpathlineto{\pgfqpoint{0.794603in}{0.932861in}}%
\pgfpathlineto{\pgfqpoint{0.792799in}{0.932099in}}%
\pgfpathlineto{\pgfqpoint{0.786354in}{0.929407in}}%
\pgfpathlineto{\pgfqpoint{0.778105in}{0.925988in}}%
\pgfpathlineto{\pgfqpoint{0.775406in}{0.924877in}}%
\pgfpathlineto{\pgfqpoint{0.769856in}{0.922616in}}%
\pgfpathlineto{\pgfqpoint{0.761606in}{0.919281in}}%
\pgfpathlineto{\pgfqpoint{0.757555in}{0.917654in}}%
\pgfpathlineto{\pgfqpoint{0.753357in}{0.915988in}}%
\pgfpathlineto{\pgfqpoint{0.745108in}{0.912734in}}%
\pgfpathlineto{\pgfqpoint{0.739230in}{0.910432in}}%
\pgfpathlineto{\pgfqpoint{0.736859in}{0.909514in}}%
\pgfpathlineto{\pgfqpoint{0.728609in}{0.906339in}}%
\pgfpathlineto{\pgfqpoint{0.720417in}{0.903210in}}%
\pgfpathlineto{\pgfqpoint{0.720360in}{0.903188in}}%
\pgfpathlineto{\pgfqpoint{0.712111in}{0.900089in}}%
\pgfpathlineto{\pgfqpoint{0.703862in}{0.897013in}}%
\pgfpathlineto{\pgfqpoint{0.701095in}{0.895988in}}%
\pgfpathlineto{\pgfqpoint{0.695612in}{0.893979in}}%
\pgfpathlineto{\pgfqpoint{0.687363in}{0.890976in}}%
\pgfpathlineto{\pgfqpoint{0.681251in}{0.888766in}}%
\pgfpathlineto{\pgfqpoint{0.679114in}{0.888002in}}%
\pgfpathlineto{\pgfqpoint{0.670865in}{0.885070in}}%
\pgfpathlineto{\pgfqpoint{0.662616in}{0.882160in}}%
\pgfpathlineto{\pgfqpoint{0.660856in}{0.881543in}}%
\pgfpathlineto{\pgfqpoint{0.654366in}{0.879293in}}%
\pgfpathlineto{\pgfqpoint{0.646117in}{0.876453in}}%
\pgfpathlineto{\pgfqpoint{0.639881in}{0.874321in}}%
\pgfpathlineto{\pgfqpoint{0.637868in}{0.873641in}}%
\pgfpathlineto{\pgfqpoint{0.629619in}{0.870869in}}%
\pgfpathlineto{\pgfqpoint{0.621369in}{0.868119in}}%
\pgfpathlineto{\pgfqpoint{0.618292in}{0.867099in}}%
\pgfpathlineto{\pgfqpoint{0.613120in}{0.865405in}}%
\pgfpathlineto{\pgfqpoint{0.604871in}{0.862721in}}%
\pgfpathlineto{\pgfqpoint{0.596622in}{0.860057in}}%
\pgfpathlineto{\pgfqpoint{0.596060in}{0.859877in}}%
\pgfpathlineto{\pgfqpoint{0.588373in}{0.857437in}}%
\pgfpathlineto{\pgfqpoint{0.580123in}{0.854838in}}%
\pgfpathlineto{\pgfqpoint{0.573146in}{0.852654in}}%
\pgfpathlineto{\pgfqpoint{0.571874in}{0.852261in}}%
\pgfpathlineto{\pgfqpoint{0.563625in}{0.849725in}}%
\pgfpathlineto{\pgfqpoint{0.555376in}{0.847206in}}%
\pgfpathlineto{\pgfqpoint{0.549526in}{0.845432in}}%
\pgfpathlineto{\pgfqpoint{0.547126in}{0.844713in}}%
\pgfpathlineto{\pgfqpoint{0.538877in}{0.842255in}}%
\pgfpathlineto{\pgfqpoint{0.530628in}{0.839814in}}%
\pgfpathlineto{\pgfqpoint{0.525175in}{0.838210in}}%
\pgfpathlineto{\pgfqpoint{0.522379in}{0.837398in}}%
\pgfpathlineto{\pgfqpoint{0.514130in}{0.835014in}}%
\pgfpathlineto{\pgfqpoint{0.505880in}{0.832647in}}%
\pgfpathlineto{\pgfqpoint{0.500063in}{0.830988in}}%
\pgfpathlineto{\pgfqpoint{0.497631in}{0.830302in}}%
\pgfpathlineto{\pgfqpoint{0.489382in}{0.827990in}}%
\pgfpathlineto{\pgfqpoint{0.481133in}{0.825694in}}%
\pgfpathlineto{\pgfqpoint{0.474163in}{0.823766in}}%
\pgfpathlineto{\pgfqpoint{0.472883in}{0.823416in}}%
\pgfpathlineto{\pgfqpoint{0.464634in}{0.821172in}}%
\pgfpathlineto{\pgfqpoint{0.456385in}{0.818943in}}%
\pgfpathlineto{\pgfqpoint{0.448136in}{0.816729in}}%
\pgfpathlineto{\pgfqpoint{0.447443in}{0.816543in}}%
\pgfpathlineto{\pgfqpoint{0.439887in}{0.814549in}}%
\pgfpathlineto{\pgfqpoint{0.431637in}{0.812384in}}%
\pgfpathlineto{\pgfqpoint{0.423388in}{0.810234in}}%
\pgfpathlineto{\pgfqpoint{0.419871in}{0.809321in}}%
\pgfpathlineto{\pgfqpoint{0.415139in}{0.808109in}}%
\pgfpathlineto{\pgfqpoint{0.415139in}{0.802099in}}%
\pgfpathlineto{\pgfqpoint{0.415139in}{0.794877in}}%
\pgfpathlineto{\pgfqpoint{0.415139in}{0.787654in}}%
\pgfpathlineto{\pgfqpoint{0.415139in}{0.780432in}}%
\pgfpathlineto{\pgfqpoint{0.415139in}{0.773210in}}%
\pgfpathlineto{\pgfqpoint{0.415139in}{0.765988in}}%
\pgfpathlineto{\pgfqpoint{0.415139in}{0.758766in}}%
\pgfpathlineto{\pgfqpoint{0.415139in}{0.751543in}}%
\pgfpathlineto{\pgfqpoint{0.415139in}{0.744321in}}%
\pgfpathlineto{\pgfqpoint{0.415139in}{0.737099in}}%
\pgfpathlineto{\pgfqpoint{0.415139in}{0.729877in}}%
\pgfpathlineto{\pgfqpoint{0.415139in}{0.722654in}}%
\pgfpathlineto{\pgfqpoint{0.415139in}{0.715432in}}%
\pgfpathlineto{\pgfqpoint{0.415139in}{0.708210in}}%
\pgfpathlineto{\pgfqpoint{0.415139in}{0.700988in}}%
\pgfpathlineto{\pgfqpoint{0.415139in}{0.693766in}}%
\pgfpathlineto{\pgfqpoint{0.415139in}{0.686543in}}%
\pgfpathlineto{\pgfqpoint{0.415139in}{0.679321in}}%
\pgfpathlineto{\pgfqpoint{0.415139in}{0.672099in}}%
\pgfpathlineto{\pgfqpoint{0.415139in}{0.664877in}}%
\pgfpathlineto{\pgfqpoint{0.415139in}{0.657654in}}%
\pgfpathlineto{\pgfqpoint{0.415139in}{0.650432in}}%
\pgfpathlineto{\pgfqpoint{0.415139in}{0.643210in}}%
\pgfpathlineto{\pgfqpoint{0.415139in}{0.635988in}}%
\pgfpathlineto{\pgfqpoint{0.415139in}{0.630196in}}%
\pgfpathlineto{\pgfqpoint{0.423388in}{0.631473in}}%
\pgfpathclose%
\pgfusepath{fill}%
\end{pgfscope}%
\begin{pgfscope}%
\pgfpathrectangle{\pgfqpoint{0.415139in}{0.462654in}}{\pgfqpoint{1.641596in}{1.437222in}}%
\pgfusepath{clip}%
\pgfsetbuttcap%
\pgfsetroundjoin%
\definecolor{currentfill}{rgb}{0.556401,0.759631,0.863899}%
\pgfsetfillcolor{currentfill}%
\pgfsetlinewidth{0.000000pt}%
\definecolor{currentstroke}{rgb}{0.000000,0.000000,0.000000}%
\pgfsetstrokecolor{currentstroke}%
\pgfsetdash{}{0pt}%
\pgfpathmoveto{\pgfqpoint{0.423388in}{0.479460in}}%
\pgfpathlineto{\pgfqpoint{0.431637in}{0.480024in}}%
\pgfpathlineto{\pgfqpoint{0.439887in}{0.480593in}}%
\pgfpathlineto{\pgfqpoint{0.448136in}{0.481168in}}%
\pgfpathlineto{\pgfqpoint{0.456385in}{0.481747in}}%
\pgfpathlineto{\pgfqpoint{0.464634in}{0.482333in}}%
\pgfpathlineto{\pgfqpoint{0.472883in}{0.482923in}}%
\pgfpathlineto{\pgfqpoint{0.481133in}{0.483520in}}%
\pgfpathlineto{\pgfqpoint{0.489382in}{0.484122in}}%
\pgfpathlineto{\pgfqpoint{0.492092in}{0.484321in}}%
\pgfpathlineto{\pgfqpoint{0.497631in}{0.484731in}}%
\pgfpathlineto{\pgfqpoint{0.505880in}{0.485347in}}%
\pgfpathlineto{\pgfqpoint{0.514130in}{0.485969in}}%
\pgfpathlineto{\pgfqpoint{0.522379in}{0.486597in}}%
\pgfpathlineto{\pgfqpoint{0.530628in}{0.487232in}}%
\pgfpathlineto{\pgfqpoint{0.538877in}{0.487872in}}%
\pgfpathlineto{\pgfqpoint{0.547126in}{0.488519in}}%
\pgfpathlineto{\pgfqpoint{0.555376in}{0.489173in}}%
\pgfpathlineto{\pgfqpoint{0.563625in}{0.489834in}}%
\pgfpathlineto{\pgfqpoint{0.571874in}{0.490501in}}%
\pgfpathlineto{\pgfqpoint{0.580123in}{0.491175in}}%
\pgfpathlineto{\pgfqpoint{0.584589in}{0.491543in}}%
\pgfpathlineto{\pgfqpoint{0.588373in}{0.491858in}}%
\pgfpathlineto{\pgfqpoint{0.596622in}{0.492548in}}%
\pgfpathlineto{\pgfqpoint{0.604871in}{0.493246in}}%
\pgfpathlineto{\pgfqpoint{0.613120in}{0.493951in}}%
\pgfpathlineto{\pgfqpoint{0.621369in}{0.494664in}}%
\pgfpathlineto{\pgfqpoint{0.629619in}{0.495384in}}%
\pgfpathlineto{\pgfqpoint{0.637868in}{0.496113in}}%
\pgfpathlineto{\pgfqpoint{0.646117in}{0.496849in}}%
\pgfpathlineto{\pgfqpoint{0.654366in}{0.497593in}}%
\pgfpathlineto{\pgfqpoint{0.662616in}{0.498346in}}%
\pgfpathlineto{\pgfqpoint{0.667177in}{0.498766in}}%
\pgfpathlineto{\pgfqpoint{0.670865in}{0.499107in}}%
\pgfpathlineto{\pgfqpoint{0.679114in}{0.499878in}}%
\pgfpathlineto{\pgfqpoint{0.687363in}{0.500656in}}%
\pgfpathlineto{\pgfqpoint{0.695612in}{0.501443in}}%
\pgfpathlineto{\pgfqpoint{0.703862in}{0.502239in}}%
\pgfpathlineto{\pgfqpoint{0.712111in}{0.503043in}}%
\pgfpathlineto{\pgfqpoint{0.720360in}{0.503856in}}%
\pgfpathlineto{\pgfqpoint{0.728609in}{0.504679in}}%
\pgfpathlineto{\pgfqpoint{0.736859in}{0.505510in}}%
\pgfpathlineto{\pgfqpoint{0.741557in}{0.505988in}}%
\pgfpathlineto{\pgfqpoint{0.745108in}{0.506351in}}%
\pgfpathlineto{\pgfqpoint{0.753357in}{0.507203in}}%
\pgfpathlineto{\pgfqpoint{0.761606in}{0.508064in}}%
\pgfpathlineto{\pgfqpoint{0.769856in}{0.508935in}}%
\pgfpathlineto{\pgfqpoint{0.778105in}{0.509816in}}%
\pgfpathlineto{\pgfqpoint{0.786354in}{0.510708in}}%
\pgfpathlineto{\pgfqpoint{0.794603in}{0.511610in}}%
\pgfpathlineto{\pgfqpoint{0.802852in}{0.512524in}}%
\pgfpathlineto{\pgfqpoint{0.808986in}{0.513210in}}%
\pgfpathlineto{\pgfqpoint{0.811102in}{0.513448in}}%
\pgfpathlineto{\pgfqpoint{0.819351in}{0.514384in}}%
\pgfpathlineto{\pgfqpoint{0.827600in}{0.515332in}}%
\pgfpathlineto{\pgfqpoint{0.835849in}{0.516291in}}%
\pgfpathlineto{\pgfqpoint{0.844099in}{0.517262in}}%
\pgfpathlineto{\pgfqpoint{0.852348in}{0.518245in}}%
\pgfpathlineto{\pgfqpoint{0.860597in}{0.519241in}}%
\pgfpathlineto{\pgfqpoint{0.868846in}{0.520249in}}%
\pgfpathlineto{\pgfqpoint{0.870336in}{0.520432in}}%
\pgfpathlineto{\pgfqpoint{0.877095in}{0.521270in}}%
\pgfpathlineto{\pgfqpoint{0.885345in}{0.522303in}}%
\pgfpathlineto{\pgfqpoint{0.893594in}{0.523350in}}%
\pgfpathlineto{\pgfqpoint{0.901843in}{0.524410in}}%
\pgfpathlineto{\pgfqpoint{0.910092in}{0.525484in}}%
\pgfpathlineto{\pgfqpoint{0.918342in}{0.526571in}}%
\pgfpathlineto{\pgfqpoint{0.926450in}{0.527654in}}%
\pgfpathlineto{\pgfqpoint{0.926591in}{0.527673in}}%
\pgfpathlineto{\pgfqpoint{0.934840in}{0.528790in}}%
\pgfpathlineto{\pgfqpoint{0.943089in}{0.529921in}}%
\pgfpathlineto{\pgfqpoint{0.951338in}{0.531067in}}%
\pgfpathlineto{\pgfqpoint{0.959588in}{0.532228in}}%
\pgfpathlineto{\pgfqpoint{0.967837in}{0.533406in}}%
\pgfpathlineto{\pgfqpoint{0.976086in}{0.534600in}}%
\pgfpathlineto{\pgfqpoint{0.977978in}{0.534877in}}%
\pgfpathlineto{\pgfqpoint{0.984335in}{0.535811in}}%
\pgfpathlineto{\pgfqpoint{0.992585in}{0.537038in}}%
\pgfpathlineto{\pgfqpoint{1.000834in}{0.538282in}}%
\pgfpathlineto{\pgfqpoint{1.009083in}{0.539544in}}%
\pgfpathlineto{\pgfqpoint{1.017332in}{0.540825in}}%
\pgfpathlineto{\pgfqpoint{1.025423in}{0.542099in}}%
\pgfpathlineto{\pgfqpoint{1.025581in}{0.542124in}}%
\pgfpathlineto{\pgfqpoint{1.033831in}{0.543440in}}%
\pgfpathlineto{\pgfqpoint{1.042080in}{0.544776in}}%
\pgfpathlineto{\pgfqpoint{1.050329in}{0.546130in}}%
\pgfpathlineto{\pgfqpoint{1.058578in}{0.547505in}}%
\pgfpathlineto{\pgfqpoint{1.066828in}{0.548900in}}%
\pgfpathlineto{\pgfqpoint{1.069297in}{0.549321in}}%
\pgfpathlineto{\pgfqpoint{1.075077in}{0.550313in}}%
\pgfpathlineto{\pgfqpoint{1.083326in}{0.551747in}}%
\pgfpathlineto{\pgfqpoint{1.091575in}{0.553203in}}%
\pgfpathlineto{\pgfqpoint{1.099824in}{0.554679in}}%
\pgfpathlineto{\pgfqpoint{1.108074in}{0.556178in}}%
\pgfpathlineto{\pgfqpoint{1.110064in}{0.556543in}}%
\pgfpathlineto{\pgfqpoint{1.116323in}{0.557698in}}%
\pgfpathlineto{\pgfqpoint{1.124572in}{0.559240in}}%
\pgfpathlineto{\pgfqpoint{1.132821in}{0.560806in}}%
\pgfpathlineto{\pgfqpoint{1.141071in}{0.562397in}}%
\pgfpathlineto{\pgfqpoint{1.148068in}{0.563766in}}%
\pgfpathlineto{\pgfqpoint{1.149320in}{0.564012in}}%
\pgfpathlineto{\pgfqpoint{1.157569in}{0.565650in}}%
\pgfpathlineto{\pgfqpoint{1.165818in}{0.567314in}}%
\pgfpathlineto{\pgfqpoint{1.174067in}{0.569005in}}%
\pgfpathlineto{\pgfqpoint{1.182317in}{0.570724in}}%
\pgfpathlineto{\pgfqpoint{1.183569in}{0.570988in}}%
\pgfpathlineto{\pgfqpoint{1.190566in}{0.572469in}}%
\pgfpathlineto{\pgfqpoint{1.198815in}{0.574242in}}%
\pgfpathlineto{\pgfqpoint{1.207064in}{0.576044in}}%
\pgfpathlineto{\pgfqpoint{1.215314in}{0.577877in}}%
\pgfpathlineto{\pgfqpoint{1.216797in}{0.578210in}}%
\pgfpathlineto{\pgfqpoint{1.223563in}{0.579738in}}%
\pgfpathlineto{\pgfqpoint{1.231812in}{0.581630in}}%
\pgfpathlineto{\pgfqpoint{1.240061in}{0.583555in}}%
\pgfpathlineto{\pgfqpoint{1.247974in}{0.585432in}}%
\pgfpathlineto{\pgfqpoint{1.248310in}{0.585513in}}%
\pgfpathlineto{\pgfqpoint{1.256560in}{0.587501in}}%
\pgfpathlineto{\pgfqpoint{1.264809in}{0.589524in}}%
\pgfpathlineto{\pgfqpoint{1.273058in}{0.591582in}}%
\pgfpathlineto{\pgfqpoint{1.277297in}{0.592654in}}%
\pgfpathlineto{\pgfqpoint{1.281307in}{0.593675in}}%
\pgfpathlineto{\pgfqpoint{1.289557in}{0.595804in}}%
\pgfpathlineto{\pgfqpoint{1.297806in}{0.597971in}}%
\pgfpathlineto{\pgfqpoint{1.304939in}{0.599877in}}%
\pgfpathlineto{\pgfqpoint{1.306055in}{0.600177in}}%
\pgfpathlineto{\pgfqpoint{1.314304in}{0.602419in}}%
\pgfpathlineto{\pgfqpoint{1.322553in}{0.604704in}}%
\pgfpathlineto{\pgfqpoint{1.330803in}{0.607032in}}%
\pgfpathlineto{\pgfqpoint{1.331039in}{0.607099in}}%
\pgfpathlineto{\pgfqpoint{1.339052in}{0.609399in}}%
\pgfpathlineto{\pgfqpoint{1.347301in}{0.611812in}}%
\pgfpathlineto{\pgfqpoint{1.355550in}{0.614272in}}%
\pgfpathlineto{\pgfqpoint{1.355714in}{0.614321in}}%
\pgfpathlineto{\pgfqpoint{1.363800in}{0.616775in}}%
\pgfpathlineto{\pgfqpoint{1.372049in}{0.619328in}}%
\pgfpathlineto{\pgfqpoint{1.379074in}{0.621543in}}%
\pgfpathlineto{\pgfqpoint{1.380298in}{0.621932in}}%
\pgfpathlineto{\pgfqpoint{1.388547in}{0.624585in}}%
\pgfpathlineto{\pgfqpoint{1.396797in}{0.627292in}}%
\pgfpathlineto{\pgfqpoint{1.401216in}{0.628766in}}%
\pgfpathlineto{\pgfqpoint{1.405046in}{0.630053in}}%
\pgfpathlineto{\pgfqpoint{1.413295in}{0.632868in}}%
\pgfpathlineto{\pgfqpoint{1.421544in}{0.635743in}}%
\pgfpathlineto{\pgfqpoint{1.422240in}{0.635988in}}%
\pgfpathlineto{\pgfqpoint{1.429793in}{0.638673in}}%
\pgfpathlineto{\pgfqpoint{1.438043in}{0.641666in}}%
\pgfpathlineto{\pgfqpoint{1.442232in}{0.643210in}}%
\pgfpathlineto{\pgfqpoint{1.446292in}{0.644720in}}%
\pgfpathlineto{\pgfqpoint{1.454541in}{0.647838in}}%
\pgfpathlineto{\pgfqpoint{1.461273in}{0.650432in}}%
\pgfpathlineto{\pgfqpoint{1.462790in}{0.651022in}}%
\pgfpathlineto{\pgfqpoint{1.471040in}{0.654273in}}%
\pgfpathlineto{\pgfqpoint{1.479289in}{0.657596in}}%
\pgfpathlineto{\pgfqpoint{1.479432in}{0.657654in}}%
\pgfpathlineto{\pgfqpoint{1.487538in}{0.660989in}}%
\pgfpathlineto{\pgfqpoint{1.495787in}{0.664459in}}%
\pgfpathlineto{\pgfqpoint{1.496768in}{0.664877in}}%
\pgfpathlineto{\pgfqpoint{1.504036in}{0.668005in}}%
\pgfpathlineto{\pgfqpoint{1.512286in}{0.671633in}}%
\pgfpathlineto{\pgfqpoint{1.513330in}{0.672099in}}%
\pgfpathlineto{\pgfqpoint{1.520535in}{0.675345in}}%
\pgfpathlineto{\pgfqpoint{1.528784in}{0.679144in}}%
\pgfpathlineto{\pgfqpoint{1.529165in}{0.679321in}}%
\pgfpathlineto{\pgfqpoint{1.537033in}{0.683033in}}%
\pgfpathlineto{\pgfqpoint{1.544311in}{0.686543in}}%
\pgfpathlineto{\pgfqpoint{1.545283in}{0.687018in}}%
\pgfpathlineto{\pgfqpoint{1.553532in}{0.691101in}}%
\pgfpathlineto{\pgfqpoint{1.558806in}{0.693766in}}%
\pgfpathlineto{\pgfqpoint{1.561781in}{0.695287in}}%
\pgfpathlineto{\pgfqpoint{1.570030in}{0.699581in}}%
\pgfpathlineto{\pgfqpoint{1.572687in}{0.700988in}}%
\pgfpathlineto{\pgfqpoint{1.578279in}{0.703989in}}%
\pgfpathlineto{\pgfqpoint{1.585983in}{0.708210in}}%
\pgfpathlineto{\pgfqpoint{1.586529in}{0.708513in}}%
\pgfpathlineto{\pgfqpoint{1.594778in}{0.713163in}}%
\pgfpathlineto{\pgfqpoint{1.598725in}{0.715432in}}%
\pgfpathlineto{\pgfqpoint{1.603027in}{0.717942in}}%
\pgfpathlineto{\pgfqpoint{1.610944in}{0.722654in}}%
\pgfpathlineto{\pgfqpoint{1.611276in}{0.722855in}}%
\pgfpathlineto{\pgfqpoint{1.619526in}{0.727915in}}%
\pgfpathlineto{\pgfqpoint{1.622663in}{0.729877in}}%
\pgfpathlineto{\pgfqpoint{1.627775in}{0.733125in}}%
\pgfpathlineto{\pgfqpoint{1.633908in}{0.737099in}}%
\pgfpathlineto{\pgfqpoint{1.636024in}{0.738493in}}%
\pgfpathlineto{\pgfqpoint{1.644273in}{0.744030in}}%
\pgfpathlineto{\pgfqpoint{1.644700in}{0.744321in}}%
\pgfpathlineto{\pgfqpoint{1.652522in}{0.749750in}}%
\pgfpathlineto{\pgfqpoint{1.655058in}{0.751543in}}%
\pgfpathlineto{\pgfqpoint{1.660772in}{0.755659in}}%
\pgfpathlineto{\pgfqpoint{1.665005in}{0.758766in}}%
\pgfpathlineto{\pgfqpoint{1.669021in}{0.761768in}}%
\pgfpathlineto{\pgfqpoint{1.674562in}{0.765988in}}%
\pgfpathlineto{\pgfqpoint{1.677270in}{0.768089in}}%
\pgfpathlineto{\pgfqpoint{1.683751in}{0.773210in}}%
\pgfpathlineto{\pgfqpoint{1.685519in}{0.774634in}}%
\pgfpathlineto{\pgfqpoint{1.692591in}{0.780432in}}%
\pgfpathlineto{\pgfqpoint{1.693769in}{0.781417in}}%
\pgfpathlineto{\pgfqpoint{1.701099in}{0.787654in}}%
\pgfpathlineto{\pgfqpoint{1.702018in}{0.788451in}}%
\pgfpathlineto{\pgfqpoint{1.709293in}{0.794877in}}%
\pgfpathlineto{\pgfqpoint{1.710267in}{0.795754in}}%
\pgfpathlineto{\pgfqpoint{1.717189in}{0.802099in}}%
\pgfpathlineto{\pgfqpoint{1.718516in}{0.803341in}}%
\pgfpathlineto{\pgfqpoint{1.724800in}{0.809321in}}%
\pgfpathlineto{\pgfqpoint{1.726765in}{0.811231in}}%
\pgfpathlineto{\pgfqpoint{1.732142in}{0.816543in}}%
\pgfpathlineto{\pgfqpoint{1.735015in}{0.819443in}}%
\pgfpathlineto{\pgfqpoint{1.739226in}{0.823766in}}%
\pgfpathlineto{\pgfqpoint{1.743264in}{0.828000in}}%
\pgfpathlineto{\pgfqpoint{1.746067in}{0.830988in}}%
\pgfpathlineto{\pgfqpoint{1.751513in}{0.836924in}}%
\pgfpathlineto{\pgfqpoint{1.752674in}{0.838210in}}%
\pgfpathlineto{\pgfqpoint{1.759059in}{0.845432in}}%
\pgfpathlineto{\pgfqpoint{1.759762in}{0.846244in}}%
\pgfpathlineto{\pgfqpoint{1.765230in}{0.852654in}}%
\pgfpathlineto{\pgfqpoint{1.768012in}{0.855991in}}%
\pgfpathlineto{\pgfqpoint{1.771199in}{0.859877in}}%
\pgfpathlineto{\pgfqpoint{1.776261in}{0.866192in}}%
\pgfpathlineto{\pgfqpoint{1.776977in}{0.867099in}}%
\pgfpathlineto{\pgfqpoint{1.782567in}{0.874321in}}%
\pgfpathlineto{\pgfqpoint{1.784510in}{0.876888in}}%
\pgfpathlineto{\pgfqpoint{1.787981in}{0.881543in}}%
\pgfpathlineto{\pgfqpoint{1.792759in}{0.888111in}}%
\pgfpathlineto{\pgfqpoint{1.793229in}{0.888766in}}%
\pgfpathlineto{\pgfqpoint{1.798311in}{0.895988in}}%
\pgfpathlineto{\pgfqpoint{1.801008in}{0.899912in}}%
\pgfpathlineto{\pgfqpoint{1.803241in}{0.903210in}}%
\pgfpathlineto{\pgfqpoint{1.808024in}{0.910432in}}%
\pgfpathlineto{\pgfqpoint{1.809258in}{0.912333in}}%
\pgfpathlineto{\pgfqpoint{1.812662in}{0.917654in}}%
\pgfpathlineto{\pgfqpoint{1.817168in}{0.924877in}}%
\pgfpathlineto{\pgfqpoint{1.817507in}{0.925430in}}%
\pgfpathlineto{\pgfqpoint{1.821539in}{0.932099in}}%
\pgfpathlineto{\pgfqpoint{1.825756in}{0.939265in}}%
\pgfpathlineto{\pgfqpoint{1.825789in}{0.939321in}}%
\pgfpathlineto{\pgfqpoint{1.829914in}{0.946543in}}%
\pgfpathlineto{\pgfqpoint{1.833927in}{0.953766in}}%
\pgfpathlineto{\pgfqpoint{1.834005in}{0.953908in}}%
\pgfpathlineto{\pgfqpoint{1.837825in}{0.960988in}}%
\pgfpathlineto{\pgfqpoint{1.841620in}{0.968210in}}%
\pgfpathlineto{\pgfqpoint{1.842255in}{0.969438in}}%
\pgfpathlineto{\pgfqpoint{1.845309in}{0.975432in}}%
\pgfpathlineto{\pgfqpoint{1.848900in}{0.982654in}}%
\pgfpathlineto{\pgfqpoint{1.850504in}{0.985944in}}%
\pgfpathlineto{\pgfqpoint{1.852396in}{0.989877in}}%
\pgfpathlineto{\pgfqpoint{1.855799in}{0.997099in}}%
\pgfpathlineto{\pgfqpoint{1.858753in}{1.003523in}}%
\pgfpathlineto{\pgfqpoint{1.859115in}{1.004321in}}%
\pgfpathlineto{\pgfqpoint{1.862342in}{1.011543in}}%
\pgfpathlineto{\pgfqpoint{1.865490in}{1.018766in}}%
\pgfpathlineto{\pgfqpoint{1.867002in}{1.022303in}}%
\pgfpathlineto{\pgfqpoint{1.868557in}{1.025988in}}%
\pgfpathlineto{\pgfqpoint{1.871546in}{1.033210in}}%
\pgfpathlineto{\pgfqpoint{1.874464in}{1.040432in}}%
\pgfpathlineto{\pgfqpoint{1.875251in}{1.042415in}}%
\pgfpathlineto{\pgfqpoint{1.877307in}{1.047654in}}%
\pgfpathlineto{\pgfqpoint{1.880081in}{1.054877in}}%
\pgfpathlineto{\pgfqpoint{1.882790in}{1.062099in}}%
\pgfpathlineto{\pgfqpoint{1.883501in}{1.064023in}}%
\pgfpathlineto{\pgfqpoint{1.885432in}{1.069321in}}%
\pgfpathlineto{\pgfqpoint{1.888011in}{1.076543in}}%
\pgfpathlineto{\pgfqpoint{1.890532in}{1.083766in}}%
\pgfpathlineto{\pgfqpoint{1.891750in}{1.087318in}}%
\pgfpathlineto{\pgfqpoint{1.892993in}{1.090988in}}%
\pgfpathlineto{\pgfqpoint{1.895395in}{1.098210in}}%
\pgfpathlineto{\pgfqpoint{1.897744in}{1.105432in}}%
\pgfpathlineto{\pgfqpoint{1.899999in}{1.112519in}}%
\pgfpathlineto{\pgfqpoint{1.900042in}{1.112654in}}%
\pgfpathlineto{\pgfqpoint{1.902283in}{1.119877in}}%
\pgfpathlineto{\pgfqpoint{1.904476in}{1.127099in}}%
\pgfpathlineto{\pgfqpoint{1.906622in}{1.134321in}}%
\pgfpathlineto{\pgfqpoint{1.908248in}{1.139901in}}%
\pgfpathlineto{\pgfqpoint{1.908721in}{1.141543in}}%
\pgfpathlineto{\pgfqpoint{1.910771in}{1.148766in}}%
\pgfpathlineto{\pgfqpoint{1.912779in}{1.155988in}}%
\pgfpathlineto{\pgfqpoint{1.914745in}{1.163210in}}%
\pgfpathlineto{\pgfqpoint{1.916498in}{1.169776in}}%
\pgfpathlineto{\pgfqpoint{1.916671in}{1.170432in}}%
\pgfpathlineto{\pgfqpoint{1.918551in}{1.177654in}}%
\pgfpathlineto{\pgfqpoint{1.920394in}{1.184877in}}%
\pgfpathlineto{\pgfqpoint{1.922201in}{1.192099in}}%
\pgfpathlineto{\pgfqpoint{1.923972in}{1.199321in}}%
\pgfpathlineto{\pgfqpoint{1.924747in}{1.202531in}}%
\pgfpathlineto{\pgfqpoint{1.925704in}{1.206543in}}%
\pgfpathlineto{\pgfqpoint{1.927401in}{1.213766in}}%
\pgfpathlineto{\pgfqpoint{1.929064in}{1.220988in}}%
\pgfpathlineto{\pgfqpoint{1.930697in}{1.228210in}}%
\pgfpathlineto{\pgfqpoint{1.932298in}{1.235432in}}%
\pgfpathlineto{\pgfqpoint{1.932996in}{1.238624in}}%
\pgfpathlineto{\pgfqpoint{1.933867in}{1.242654in}}%
\pgfpathlineto{\pgfqpoint{1.935403in}{1.249877in}}%
\pgfpathlineto{\pgfqpoint{1.936911in}{1.257099in}}%
\pgfpathlineto{\pgfqpoint{1.938392in}{1.264321in}}%
\pgfpathlineto{\pgfqpoint{1.939846in}{1.271543in}}%
\pgfpathlineto{\pgfqpoint{1.941245in}{1.278616in}}%
\pgfpathlineto{\pgfqpoint{1.941275in}{1.278766in}}%
\pgfpathlineto{\pgfqpoint{1.942672in}{1.285988in}}%
\pgfpathlineto{\pgfqpoint{1.944044in}{1.293210in}}%
\pgfpathlineto{\pgfqpoint{1.945393in}{1.300432in}}%
\pgfpathlineto{\pgfqpoint{1.946719in}{1.307654in}}%
\pgfpathlineto{\pgfqpoint{1.948022in}{1.314877in}}%
\pgfpathlineto{\pgfqpoint{1.949303in}{1.322099in}}%
\pgfpathlineto{\pgfqpoint{1.949494in}{1.323190in}}%
\pgfpathlineto{\pgfqpoint{1.950558in}{1.329321in}}%
\pgfpathlineto{\pgfqpoint{1.951791in}{1.336543in}}%
\pgfpathlineto{\pgfqpoint{1.953004in}{1.343766in}}%
\pgfpathlineto{\pgfqpoint{1.954197in}{1.350988in}}%
\pgfpathlineto{\pgfqpoint{1.955372in}{1.358210in}}%
\pgfpathlineto{\pgfqpoint{1.956527in}{1.365432in}}%
\pgfpathlineto{\pgfqpoint{1.957665in}{1.372654in}}%
\pgfpathlineto{\pgfqpoint{1.957744in}{1.373159in}}%
\pgfpathlineto{\pgfqpoint{1.958780in}{1.379877in}}%
\pgfpathlineto{\pgfqpoint{1.959877in}{1.387099in}}%
\pgfpathlineto{\pgfqpoint{1.960957in}{1.394321in}}%
\pgfpathlineto{\pgfqpoint{1.962021in}{1.401543in}}%
\pgfpathlineto{\pgfqpoint{1.963069in}{1.408766in}}%
\pgfpathlineto{\pgfqpoint{1.964102in}{1.415988in}}%
\pgfpathlineto{\pgfqpoint{1.965120in}{1.423210in}}%
\pgfpathlineto{\pgfqpoint{1.965993in}{1.429489in}}%
\pgfpathlineto{\pgfqpoint{1.966123in}{1.430432in}}%
\pgfpathlineto{\pgfqpoint{1.967106in}{1.437654in}}%
\pgfpathlineto{\pgfqpoint{1.968076in}{1.444877in}}%
\pgfpathlineto{\pgfqpoint{1.969032in}{1.452099in}}%
\pgfpathlineto{\pgfqpoint{1.969975in}{1.459321in}}%
\pgfpathlineto{\pgfqpoint{1.970905in}{1.466543in}}%
\pgfpathlineto{\pgfqpoint{1.971823in}{1.473766in}}%
\pgfpathlineto{\pgfqpoint{1.972728in}{1.480988in}}%
\pgfpathlineto{\pgfqpoint{1.973622in}{1.488210in}}%
\pgfpathlineto{\pgfqpoint{1.974242in}{1.493279in}}%
\pgfpathlineto{\pgfqpoint{1.974502in}{1.495432in}}%
\pgfpathlineto{\pgfqpoint{1.975368in}{1.502654in}}%
\pgfpathlineto{\pgfqpoint{1.976223in}{1.509877in}}%
\pgfpathlineto{\pgfqpoint{1.977067in}{1.517099in}}%
\pgfpathlineto{\pgfqpoint{1.977901in}{1.524321in}}%
\pgfpathlineto{\pgfqpoint{1.978725in}{1.531543in}}%
\pgfpathlineto{\pgfqpoint{1.979539in}{1.538766in}}%
\pgfpathlineto{\pgfqpoint{1.980343in}{1.545988in}}%
\pgfpathlineto{\pgfqpoint{1.981138in}{1.553210in}}%
\pgfpathlineto{\pgfqpoint{1.981925in}{1.560432in}}%
\pgfpathlineto{\pgfqpoint{1.982491in}{1.565691in}}%
\pgfpathlineto{\pgfqpoint{1.982700in}{1.567654in}}%
\pgfpathlineto{\pgfqpoint{1.983464in}{1.574877in}}%
\pgfpathlineto{\pgfqpoint{1.984219in}{1.582099in}}%
\pgfpathlineto{\pgfqpoint{1.984967in}{1.589321in}}%
\pgfpathlineto{\pgfqpoint{1.985706in}{1.596543in}}%
\pgfpathlineto{\pgfqpoint{1.986438in}{1.603766in}}%
\pgfpathlineto{\pgfqpoint{1.987162in}{1.610988in}}%
\pgfpathlineto{\pgfqpoint{1.987879in}{1.618210in}}%
\pgfpathlineto{\pgfqpoint{1.988588in}{1.625432in}}%
\pgfpathlineto{\pgfqpoint{1.989291in}{1.632654in}}%
\pgfpathlineto{\pgfqpoint{1.989987in}{1.639877in}}%
\pgfpathlineto{\pgfqpoint{1.990677in}{1.647099in}}%
\pgfpathlineto{\pgfqpoint{1.990741in}{1.647767in}}%
\pgfpathlineto{\pgfqpoint{1.991356in}{1.654321in}}%
\pgfpathlineto{\pgfqpoint{1.992029in}{1.661543in}}%
\pgfpathlineto{\pgfqpoint{1.992695in}{1.668766in}}%
\pgfpathlineto{\pgfqpoint{1.993356in}{1.675988in}}%
\pgfpathlineto{\pgfqpoint{1.994011in}{1.683210in}}%
\pgfpathlineto{\pgfqpoint{1.994660in}{1.690432in}}%
\pgfpathlineto{\pgfqpoint{1.995304in}{1.697654in}}%
\pgfpathlineto{\pgfqpoint{1.995943in}{1.704877in}}%
\pgfpathlineto{\pgfqpoint{1.996578in}{1.712099in}}%
\pgfpathlineto{\pgfqpoint{1.997207in}{1.719321in}}%
\pgfpathlineto{\pgfqpoint{1.997831in}{1.726543in}}%
\pgfpathlineto{\pgfqpoint{1.998451in}{1.733766in}}%
\pgfpathlineto{\pgfqpoint{1.998990in}{1.740081in}}%
\pgfpathlineto{\pgfqpoint{1.999066in}{1.740988in}}%
\pgfpathlineto{\pgfqpoint{1.999673in}{1.748210in}}%
\pgfpathlineto{\pgfqpoint{2.000275in}{1.755432in}}%
\pgfpathlineto{\pgfqpoint{2.000874in}{1.762654in}}%
\pgfpathlineto{\pgfqpoint{2.001468in}{1.769877in}}%
\pgfpathlineto{\pgfqpoint{2.002059in}{1.777099in}}%
\pgfpathlineto{\pgfqpoint{2.002646in}{1.784321in}}%
\pgfpathlineto{\pgfqpoint{2.003229in}{1.791543in}}%
\pgfpathlineto{\pgfqpoint{2.003809in}{1.798766in}}%
\pgfpathlineto{\pgfqpoint{2.004386in}{1.805988in}}%
\pgfpathlineto{\pgfqpoint{2.004960in}{1.813210in}}%
\pgfpathlineto{\pgfqpoint{2.005531in}{1.820432in}}%
\pgfpathlineto{\pgfqpoint{2.006098in}{1.827654in}}%
\pgfpathlineto{\pgfqpoint{2.006663in}{1.834877in}}%
\pgfpathlineto{\pgfqpoint{2.007226in}{1.842099in}}%
\pgfpathlineto{\pgfqpoint{2.007239in}{1.842273in}}%
\pgfpathlineto{\pgfqpoint{2.007781in}{1.849321in}}%
\pgfpathlineto{\pgfqpoint{2.008333in}{1.856543in}}%
\pgfpathlineto{\pgfqpoint{2.008883in}{1.863766in}}%
\pgfpathlineto{\pgfqpoint{2.009430in}{1.870988in}}%
\pgfpathlineto{\pgfqpoint{2.009976in}{1.878210in}}%
\pgfpathlineto{\pgfqpoint{2.010519in}{1.885432in}}%
\pgfpathlineto{\pgfqpoint{2.011060in}{1.892654in}}%
\pgfpathlineto{\pgfqpoint{2.011599in}{1.899877in}}%
\pgfpathlineto{\pgfqpoint{2.007239in}{1.899877in}}%
\pgfpathlineto{\pgfqpoint{1.998990in}{1.899877in}}%
\pgfpathlineto{\pgfqpoint{1.990741in}{1.899877in}}%
\pgfpathlineto{\pgfqpoint{1.982491in}{1.899877in}}%
\pgfpathlineto{\pgfqpoint{1.974242in}{1.899877in}}%
\pgfpathlineto{\pgfqpoint{1.965993in}{1.899877in}}%
\pgfpathlineto{\pgfqpoint{1.957744in}{1.899877in}}%
\pgfpathlineto{\pgfqpoint{1.949494in}{1.899877in}}%
\pgfpathlineto{\pgfqpoint{1.941245in}{1.899877in}}%
\pgfpathlineto{\pgfqpoint{1.932996in}{1.899877in}}%
\pgfpathlineto{\pgfqpoint{1.924747in}{1.899877in}}%
\pgfpathlineto{\pgfqpoint{1.916498in}{1.899877in}}%
\pgfpathlineto{\pgfqpoint{1.908248in}{1.899877in}}%
\pgfpathlineto{\pgfqpoint{1.899999in}{1.899877in}}%
\pgfpathlineto{\pgfqpoint{1.891750in}{1.899877in}}%
\pgfpathlineto{\pgfqpoint{1.883501in}{1.899877in}}%
\pgfpathlineto{\pgfqpoint{1.875251in}{1.899877in}}%
\pgfpathlineto{\pgfqpoint{1.867002in}{1.899877in}}%
\pgfpathlineto{\pgfqpoint{1.858753in}{1.899877in}}%
\pgfpathlineto{\pgfqpoint{1.850504in}{1.899877in}}%
\pgfpathlineto{\pgfqpoint{1.842255in}{1.899877in}}%
\pgfpathlineto{\pgfqpoint{1.840474in}{1.899877in}}%
\pgfpathlineto{\pgfqpoint{1.839111in}{1.892654in}}%
\pgfpathlineto{\pgfqpoint{1.837740in}{1.885432in}}%
\pgfpathlineto{\pgfqpoint{1.836361in}{1.878210in}}%
\pgfpathlineto{\pgfqpoint{1.834974in}{1.870988in}}%
\pgfpathlineto{\pgfqpoint{1.834005in}{1.865967in}}%
\pgfpathlineto{\pgfqpoint{1.833575in}{1.863766in}}%
\pgfpathlineto{\pgfqpoint{1.832159in}{1.856543in}}%
\pgfpathlineto{\pgfqpoint{1.830734in}{1.849321in}}%
\pgfpathlineto{\pgfqpoint{1.829300in}{1.842099in}}%
\pgfpathlineto{\pgfqpoint{1.827857in}{1.834877in}}%
\pgfpathlineto{\pgfqpoint{1.826404in}{1.827654in}}%
\pgfpathlineto{\pgfqpoint{1.825756in}{1.824449in}}%
\pgfpathlineto{\pgfqpoint{1.824935in}{1.820432in}}%
\pgfpathlineto{\pgfqpoint{1.823450in}{1.813210in}}%
\pgfpathlineto{\pgfqpoint{1.821955in}{1.805988in}}%
\pgfpathlineto{\pgfqpoint{1.820450in}{1.798766in}}%
\pgfpathlineto{\pgfqpoint{1.818934in}{1.791543in}}%
\pgfpathlineto{\pgfqpoint{1.817507in}{1.784794in}}%
\pgfpathlineto{\pgfqpoint{1.817406in}{1.784321in}}%
\pgfpathlineto{\pgfqpoint{1.815855in}{1.777099in}}%
\pgfpathlineto{\pgfqpoint{1.814293in}{1.769877in}}%
\pgfpathlineto{\pgfqpoint{1.812719in}{1.762654in}}%
\pgfpathlineto{\pgfqpoint{1.811134in}{1.755432in}}%
\pgfpathlineto{\pgfqpoint{1.809536in}{1.748210in}}%
\pgfpathlineto{\pgfqpoint{1.809258in}{1.746957in}}%
\pgfpathlineto{\pgfqpoint{1.807915in}{1.740988in}}%
\pgfpathlineto{\pgfqpoint{1.806279in}{1.733766in}}%
\pgfpathlineto{\pgfqpoint{1.804630in}{1.726543in}}%
\pgfpathlineto{\pgfqpoint{1.802968in}{1.719321in}}%
\pgfpathlineto{\pgfqpoint{1.801292in}{1.712099in}}%
\pgfpathlineto{\pgfqpoint{1.801008in}{1.710882in}}%
\pgfpathlineto{\pgfqpoint{1.799591in}{1.704877in}}%
\pgfpathlineto{\pgfqpoint{1.797874in}{1.697654in}}%
\pgfpathlineto{\pgfqpoint{1.796142in}{1.690432in}}%
\pgfpathlineto{\pgfqpoint{1.794396in}{1.683210in}}%
\pgfpathlineto{\pgfqpoint{1.792759in}{1.676501in}}%
\pgfpathlineto{\pgfqpoint{1.792633in}{1.675988in}}%
\pgfpathlineto{\pgfqpoint{1.790841in}{1.668766in}}%
\pgfpathlineto{\pgfqpoint{1.789034in}{1.661543in}}%
\pgfpathlineto{\pgfqpoint{1.787211in}{1.654321in}}%
\pgfpathlineto{\pgfqpoint{1.785371in}{1.647099in}}%
\pgfpathlineto{\pgfqpoint{1.784510in}{1.643741in}}%
\pgfpathlineto{\pgfqpoint{1.783507in}{1.639877in}}%
\pgfpathlineto{\pgfqpoint{1.781619in}{1.632654in}}%
\pgfpathlineto{\pgfqpoint{1.779713in}{1.625432in}}%
\pgfpathlineto{\pgfqpoint{1.777789in}{1.618210in}}%
\pgfpathlineto{\pgfqpoint{1.776261in}{1.612524in}}%
\pgfpathlineto{\pgfqpoint{1.775843in}{1.610988in}}%
\pgfpathlineto{\pgfqpoint{1.773867in}{1.603766in}}%
\pgfpathlineto{\pgfqpoint{1.771872in}{1.596543in}}%
\pgfpathlineto{\pgfqpoint{1.769857in}{1.589321in}}%
\pgfpathlineto{\pgfqpoint{1.768012in}{1.582772in}}%
\pgfpathlineto{\pgfqpoint{1.767820in}{1.582099in}}%
\pgfpathlineto{\pgfqpoint{1.765749in}{1.574877in}}%
\pgfpathlineto{\pgfqpoint{1.763657in}{1.567654in}}%
\pgfpathlineto{\pgfqpoint{1.761543in}{1.560432in}}%
\pgfpathlineto{\pgfqpoint{1.759762in}{1.554407in}}%
\pgfpathlineto{\pgfqpoint{1.759405in}{1.553210in}}%
\pgfpathlineto{\pgfqpoint{1.757231in}{1.545988in}}%
\pgfpathlineto{\pgfqpoint{1.755034in}{1.538766in}}%
\pgfpathlineto{\pgfqpoint{1.752814in}{1.531543in}}%
\pgfpathlineto{\pgfqpoint{1.751513in}{1.527350in}}%
\pgfpathlineto{\pgfqpoint{1.750563in}{1.524321in}}%
\pgfpathlineto{\pgfqpoint{1.748278in}{1.517099in}}%
\pgfpathlineto{\pgfqpoint{1.745968in}{1.509877in}}%
\pgfpathlineto{\pgfqpoint{1.743633in}{1.502654in}}%
\pgfpathlineto{\pgfqpoint{1.743264in}{1.501523in}}%
\pgfpathlineto{\pgfqpoint{1.741257in}{1.495432in}}%
\pgfpathlineto{\pgfqpoint{1.738852in}{1.488210in}}%
\pgfpathlineto{\pgfqpoint{1.736420in}{1.480988in}}%
\pgfpathlineto{\pgfqpoint{1.735015in}{1.476856in}}%
\pgfpathlineto{\pgfqpoint{1.733952in}{1.473766in}}%
\pgfpathlineto{\pgfqpoint{1.731446in}{1.466543in}}%
\pgfpathlineto{\pgfqpoint{1.728911in}{1.459321in}}%
\pgfpathlineto{\pgfqpoint{1.726765in}{1.453276in}}%
\pgfpathlineto{\pgfqpoint{1.726343in}{1.452099in}}%
\pgfpathlineto{\pgfqpoint{1.723730in}{1.444877in}}%
\pgfpathlineto{\pgfqpoint{1.721086in}{1.437654in}}%
\pgfpathlineto{\pgfqpoint{1.718516in}{1.430719in}}%
\pgfpathlineto{\pgfqpoint{1.718409in}{1.430432in}}%
\pgfpathlineto{\pgfqpoint{1.715682in}{1.423210in}}%
\pgfpathlineto{\pgfqpoint{1.712922in}{1.415988in}}%
\pgfpathlineto{\pgfqpoint{1.710267in}{1.409126in}}%
\pgfpathlineto{\pgfqpoint{1.710126in}{1.408766in}}%
\pgfpathlineto{\pgfqpoint{1.707279in}{1.401543in}}%
\pgfpathlineto{\pgfqpoint{1.704395in}{1.394321in}}%
\pgfpathlineto{\pgfqpoint{1.702018in}{1.388437in}}%
\pgfpathlineto{\pgfqpoint{1.701471in}{1.387099in}}%
\pgfpathlineto{\pgfqpoint{1.698495in}{1.379877in}}%
\pgfpathlineto{\pgfqpoint{1.695480in}{1.372654in}}%
\pgfpathlineto{\pgfqpoint{1.693769in}{1.368599in}}%
\pgfpathlineto{\pgfqpoint{1.692418in}{1.365432in}}%
\pgfpathlineto{\pgfqpoint{1.689305in}{1.358210in}}%
\pgfpathlineto{\pgfqpoint{1.686150in}{1.350988in}}%
\pgfpathlineto{\pgfqpoint{1.685519in}{1.349558in}}%
\pgfpathlineto{\pgfqpoint{1.682938in}{1.343766in}}%
\pgfpathlineto{\pgfqpoint{1.679679in}{1.336543in}}%
\pgfpathlineto{\pgfqpoint{1.677270in}{1.331272in}}%
\pgfpathlineto{\pgfqpoint{1.676369in}{1.329321in}}%
\pgfpathlineto{\pgfqpoint{1.673001in}{1.322099in}}%
\pgfpathlineto{\pgfqpoint{1.669586in}{1.314877in}}%
\pgfpathlineto{\pgfqpoint{1.669021in}{1.313693in}}%
\pgfpathlineto{\pgfqpoint{1.666107in}{1.307654in}}%
\pgfpathlineto{\pgfqpoint{1.662576in}{1.300432in}}%
\pgfpathlineto{\pgfqpoint{1.660772in}{1.296785in}}%
\pgfpathlineto{\pgfqpoint{1.658985in}{1.293210in}}%
\pgfpathlineto{\pgfqpoint{1.655332in}{1.285988in}}%
\pgfpathlineto{\pgfqpoint{1.652522in}{1.280505in}}%
\pgfpathlineto{\pgfqpoint{1.651621in}{1.278766in}}%
\pgfpathlineto{\pgfqpoint{1.647842in}{1.271543in}}%
\pgfpathlineto{\pgfqpoint{1.644273in}{1.264822in}}%
\pgfpathlineto{\pgfqpoint{1.644004in}{1.264321in}}%
\pgfpathlineto{\pgfqpoint{1.640092in}{1.257099in}}%
\pgfpathlineto{\pgfqpoint{1.636120in}{1.249877in}}%
\pgfpathlineto{\pgfqpoint{1.636024in}{1.249704in}}%
\pgfpathlineto{\pgfqpoint{1.632068in}{1.242654in}}%
\pgfpathlineto{\pgfqpoint{1.627953in}{1.235432in}}%
\pgfpathlineto{\pgfqpoint{1.627775in}{1.235122in}}%
\pgfpathlineto{\pgfqpoint{1.623756in}{1.228210in}}%
\pgfpathlineto{\pgfqpoint{1.619526in}{1.221045in}}%
\pgfpathlineto{\pgfqpoint{1.619491in}{1.220988in}}%
\pgfpathlineto{\pgfqpoint{1.615140in}{1.213766in}}%
\pgfpathlineto{\pgfqpoint{1.611276in}{1.207450in}}%
\pgfpathlineto{\pgfqpoint{1.610716in}{1.206543in}}%
\pgfpathlineto{\pgfqpoint{1.606203in}{1.199321in}}%
\pgfpathlineto{\pgfqpoint{1.603027in}{1.194310in}}%
\pgfpathlineto{\pgfqpoint{1.601610in}{1.192099in}}%
\pgfpathlineto{\pgfqpoint{1.596928in}{1.184877in}}%
\pgfpathlineto{\pgfqpoint{1.594778in}{1.181603in}}%
\pgfpathlineto{\pgfqpoint{1.592156in}{1.177654in}}%
\pgfpathlineto{\pgfqpoint{1.587296in}{1.170432in}}%
\pgfpathlineto{\pgfqpoint{1.586529in}{1.169305in}}%
\pgfpathlineto{\pgfqpoint{1.582335in}{1.163210in}}%
\pgfpathlineto{\pgfqpoint{1.578279in}{1.157400in}}%
\pgfpathlineto{\pgfqpoint{1.577283in}{1.155988in}}%
\pgfpathlineto{\pgfqpoint{1.572126in}{1.148766in}}%
\pgfpathlineto{\pgfqpoint{1.570030in}{1.145869in}}%
\pgfpathlineto{\pgfqpoint{1.566866in}{1.141543in}}%
\pgfpathlineto{\pgfqpoint{1.561781in}{1.134689in}}%
\pgfpathlineto{\pgfqpoint{1.561505in}{1.134321in}}%
\pgfpathlineto{\pgfqpoint{1.556028in}{1.127099in}}%
\pgfpathlineto{\pgfqpoint{1.553532in}{1.123852in}}%
\pgfpathlineto{\pgfqpoint{1.550442in}{1.119877in}}%
\pgfpathlineto{\pgfqpoint{1.545283in}{1.113332in}}%
\pgfpathlineto{\pgfqpoint{1.544743in}{1.112654in}}%
\pgfpathlineto{\pgfqpoint{1.538921in}{1.105432in}}%
\pgfpathlineto{\pgfqpoint{1.537033in}{1.103122in}}%
\pgfpathlineto{\pgfqpoint{1.532976in}{1.098210in}}%
\pgfpathlineto{\pgfqpoint{1.528784in}{1.093204in}}%
\pgfpathlineto{\pgfqpoint{1.526909in}{1.090988in}}%
\pgfpathlineto{\pgfqpoint{1.520712in}{1.083766in}}%
\pgfpathlineto{\pgfqpoint{1.520535in}{1.083561in}}%
\pgfpathlineto{\pgfqpoint{1.514377in}{1.076543in}}%
\pgfpathlineto{\pgfqpoint{1.512286in}{1.074192in}}%
\pgfpathlineto{\pgfqpoint{1.507906in}{1.069321in}}%
\pgfpathlineto{\pgfqpoint{1.504036in}{1.065076in}}%
\pgfpathlineto{\pgfqpoint{1.501294in}{1.062099in}}%
\pgfpathlineto{\pgfqpoint{1.495787in}{1.056201in}}%
\pgfpathlineto{\pgfqpoint{1.494537in}{1.054877in}}%
\pgfpathlineto{\pgfqpoint{1.487629in}{1.047654in}}%
\pgfpathlineto{\pgfqpoint{1.487538in}{1.047560in}}%
\pgfpathlineto{\pgfqpoint{1.480563in}{1.040432in}}%
\pgfpathlineto{\pgfqpoint{1.479289in}{1.039147in}}%
\pgfpathlineto{\pgfqpoint{1.473337in}{1.033210in}}%
\pgfpathlineto{\pgfqpoint{1.471040in}{1.030948in}}%
\pgfpathlineto{\pgfqpoint{1.465946in}{1.025988in}}%
\pgfpathlineto{\pgfqpoint{1.462790in}{1.022955in}}%
\pgfpathlineto{\pgfqpoint{1.458384in}{1.018766in}}%
\pgfpathlineto{\pgfqpoint{1.454541in}{1.015160in}}%
\pgfpathlineto{\pgfqpoint{1.450644in}{1.011543in}}%
\pgfpathlineto{\pgfqpoint{1.446292in}{1.007556in}}%
\pgfpathlineto{\pgfqpoint{1.442721in}{1.004321in}}%
\pgfpathlineto{\pgfqpoint{1.438043in}{1.000137in}}%
\pgfpathlineto{\pgfqpoint{1.434607in}{0.997099in}}%
\pgfpathlineto{\pgfqpoint{1.429793in}{0.992897in}}%
\pgfpathlineto{\pgfqpoint{1.426295in}{0.989877in}}%
\pgfpathlineto{\pgfqpoint{1.421544in}{0.985828in}}%
\pgfpathlineto{\pgfqpoint{1.417777in}{0.982654in}}%
\pgfpathlineto{\pgfqpoint{1.413295in}{0.978926in}}%
\pgfpathlineto{\pgfqpoint{1.409047in}{0.975432in}}%
\pgfpathlineto{\pgfqpoint{1.405046in}{0.972183in}}%
\pgfpathlineto{\pgfqpoint{1.400095in}{0.968210in}}%
\pgfpathlineto{\pgfqpoint{1.396797in}{0.965596in}}%
\pgfpathlineto{\pgfqpoint{1.390914in}{0.960988in}}%
\pgfpathlineto{\pgfqpoint{1.388547in}{0.959158in}}%
\pgfpathlineto{\pgfqpoint{1.381493in}{0.953766in}}%
\pgfpathlineto{\pgfqpoint{1.380298in}{0.952864in}}%
\pgfpathlineto{\pgfqpoint{1.372049in}{0.946710in}}%
\pgfpathlineto{\pgfqpoint{1.371823in}{0.946543in}}%
\pgfpathlineto{\pgfqpoint{1.363800in}{0.940695in}}%
\pgfpathlineto{\pgfqpoint{1.361892in}{0.939321in}}%
\pgfpathlineto{\pgfqpoint{1.355550in}{0.934810in}}%
\pgfpathlineto{\pgfqpoint{1.351693in}{0.932099in}}%
\pgfpathlineto{\pgfqpoint{1.347301in}{0.929050in}}%
\pgfpathlineto{\pgfqpoint{1.341218in}{0.924877in}}%
\pgfpathlineto{\pgfqpoint{1.339052in}{0.923409in}}%
\pgfpathlineto{\pgfqpoint{1.330803in}{0.917885in}}%
\pgfpathlineto{\pgfqpoint{1.330455in}{0.917654in}}%
\pgfpathlineto{\pgfqpoint{1.322553in}{0.912480in}}%
\pgfpathlineto{\pgfqpoint{1.319389in}{0.910432in}}%
\pgfpathlineto{\pgfqpoint{1.314304in}{0.907182in}}%
\pgfpathlineto{\pgfqpoint{1.308016in}{0.903210in}}%
\pgfpathlineto{\pgfqpoint{1.306055in}{0.901987in}}%
\pgfpathlineto{\pgfqpoint{1.297806in}{0.896896in}}%
\pgfpathlineto{\pgfqpoint{1.296319in}{0.895988in}}%
\pgfpathlineto{\pgfqpoint{1.289557in}{0.891908in}}%
\pgfpathlineto{\pgfqpoint{1.284286in}{0.888766in}}%
\pgfpathlineto{\pgfqpoint{1.281307in}{0.887011in}}%
\pgfpathlineto{\pgfqpoint{1.273058in}{0.882208in}}%
\pgfpathlineto{\pgfqpoint{1.271906in}{0.881543in}}%
\pgfpathlineto{\pgfqpoint{1.264809in}{0.877499in}}%
\pgfpathlineto{\pgfqpoint{1.259162in}{0.874321in}}%
\pgfpathlineto{\pgfqpoint{1.256560in}{0.872874in}}%
\pgfpathlineto{\pgfqpoint{1.248310in}{0.868336in}}%
\pgfpathlineto{\pgfqpoint{1.246039in}{0.867099in}}%
\pgfpathlineto{\pgfqpoint{1.240061in}{0.863883in}}%
\pgfpathlineto{\pgfqpoint{1.232520in}{0.859877in}}%
\pgfpathlineto{\pgfqpoint{1.231812in}{0.859505in}}%
\pgfpathlineto{\pgfqpoint{1.223563in}{0.855213in}}%
\pgfpathlineto{\pgfqpoint{1.218584in}{0.852654in}}%
\pgfpathlineto{\pgfqpoint{1.215314in}{0.850993in}}%
\pgfpathlineto{\pgfqpoint{1.207064in}{0.846849in}}%
\pgfpathlineto{\pgfqpoint{1.204215in}{0.845432in}}%
\pgfpathlineto{\pgfqpoint{1.198815in}{0.842779in}}%
\pgfpathlineto{\pgfqpoint{1.190566in}{0.838775in}}%
\pgfpathlineto{\pgfqpoint{1.189392in}{0.838210in}}%
\pgfpathlineto{\pgfqpoint{1.182317in}{0.834847in}}%
\pgfpathlineto{\pgfqpoint{1.174092in}{0.830988in}}%
\pgfpathlineto{\pgfqpoint{1.174067in}{0.830976in}}%
\pgfpathlineto{\pgfqpoint{1.165818in}{0.827182in}}%
\pgfpathlineto{\pgfqpoint{1.158292in}{0.823766in}}%
\pgfpathlineto{\pgfqpoint{1.157569in}{0.823441in}}%
\pgfpathlineto{\pgfqpoint{1.149320in}{0.819771in}}%
\pgfpathlineto{\pgfqpoint{1.141973in}{0.816543in}}%
\pgfpathlineto{\pgfqpoint{1.141071in}{0.816152in}}%
\pgfpathlineto{\pgfqpoint{1.132821in}{0.812600in}}%
\pgfpathlineto{\pgfqpoint{1.125111in}{0.809321in}}%
\pgfpathlineto{\pgfqpoint{1.124572in}{0.809095in}}%
\pgfpathlineto{\pgfqpoint{1.116323in}{0.805655in}}%
\pgfpathlineto{\pgfqpoint{1.108074in}{0.802259in}}%
\pgfpathlineto{\pgfqpoint{1.107682in}{0.802099in}}%
\pgfpathlineto{\pgfqpoint{1.099824in}{0.798925in}}%
\pgfpathlineto{\pgfqpoint{1.091575in}{0.795634in}}%
\pgfpathlineto{\pgfqpoint{1.089661in}{0.794877in}}%
\pgfpathlineto{\pgfqpoint{1.083326in}{0.792399in}}%
\pgfpathlineto{\pgfqpoint{1.075077in}{0.789208in}}%
\pgfpathlineto{\pgfqpoint{1.071021in}{0.787654in}}%
\pgfpathlineto{\pgfqpoint{1.066828in}{0.786067in}}%
\pgfpathlineto{\pgfqpoint{1.058578in}{0.782973in}}%
\pgfpathlineto{\pgfqpoint{1.051727in}{0.780432in}}%
\pgfpathlineto{\pgfqpoint{1.050329in}{0.779920in}}%
\pgfpathlineto{\pgfqpoint{1.042080in}{0.776920in}}%
\pgfpathlineto{\pgfqpoint{1.033831in}{0.773956in}}%
\pgfpathlineto{\pgfqpoint{1.031737in}{0.773210in}}%
\pgfpathlineto{\pgfqpoint{1.025581in}{0.771042in}}%
\pgfpathlineto{\pgfqpoint{1.017332in}{0.768168in}}%
\pgfpathlineto{\pgfqpoint{1.011008in}{0.765988in}}%
\pgfpathlineto{\pgfqpoint{1.009083in}{0.765332in}}%
\pgfpathlineto{\pgfqpoint{1.000834in}{0.762544in}}%
\pgfpathlineto{\pgfqpoint{0.992585in}{0.759790in}}%
\pgfpathlineto{\pgfqpoint{0.989491in}{0.758766in}}%
\pgfpathlineto{\pgfqpoint{0.984335in}{0.757079in}}%
\pgfpathlineto{\pgfqpoint{0.976086in}{0.754406in}}%
\pgfpathlineto{\pgfqpoint{0.967837in}{0.751765in}}%
\pgfpathlineto{\pgfqpoint{0.967140in}{0.751543in}}%
\pgfpathlineto{\pgfqpoint{0.959588in}{0.749170in}}%
\pgfpathlineto{\pgfqpoint{0.951338in}{0.746606in}}%
\pgfpathlineto{\pgfqpoint{0.943910in}{0.744321in}}%
\pgfpathlineto{\pgfqpoint{0.943089in}{0.744072in}}%
\pgfpathlineto{\pgfqpoint{0.934840in}{0.741581in}}%
\pgfpathlineto{\pgfqpoint{0.926591in}{0.739118in}}%
\pgfpathlineto{\pgfqpoint{0.919762in}{0.737099in}}%
\pgfpathlineto{\pgfqpoint{0.918342in}{0.736684in}}%
\pgfpathlineto{\pgfqpoint{0.910092in}{0.734289in}}%
\pgfpathlineto{\pgfqpoint{0.901843in}{0.731920in}}%
\pgfpathlineto{\pgfqpoint{0.894659in}{0.729877in}}%
\pgfpathlineto{\pgfqpoint{0.893594in}{0.729577in}}%
\pgfpathlineto{\pgfqpoint{0.885345in}{0.727272in}}%
\pgfpathlineto{\pgfqpoint{0.877095in}{0.724991in}}%
\pgfpathlineto{\pgfqpoint{0.868846in}{0.722733in}}%
\pgfpathlineto{\pgfqpoint{0.868557in}{0.722654in}}%
\pgfpathlineto{\pgfqpoint{0.860597in}{0.720512in}}%
\pgfpathlineto{\pgfqpoint{0.852348in}{0.718315in}}%
\pgfpathlineto{\pgfqpoint{0.844099in}{0.716139in}}%
\pgfpathlineto{\pgfqpoint{0.841399in}{0.715432in}}%
\pgfpathlineto{\pgfqpoint{0.835849in}{0.713995in}}%
\pgfpathlineto{\pgfqpoint{0.827600in}{0.711876in}}%
\pgfpathlineto{\pgfqpoint{0.819351in}{0.709779in}}%
\pgfpathlineto{\pgfqpoint{0.813123in}{0.708210in}}%
\pgfpathlineto{\pgfqpoint{0.811102in}{0.707706in}}%
\pgfpathlineto{\pgfqpoint{0.802852in}{0.705663in}}%
\pgfpathlineto{\pgfqpoint{0.794603in}{0.703640in}}%
\pgfpathlineto{\pgfqpoint{0.786354in}{0.701637in}}%
\pgfpathlineto{\pgfqpoint{0.783661in}{0.700988in}}%
\pgfpathlineto{\pgfqpoint{0.778105in}{0.699662in}}%
\pgfpathlineto{\pgfqpoint{0.769856in}{0.697709in}}%
\pgfpathlineto{\pgfqpoint{0.761606in}{0.695776in}}%
\pgfpathlineto{\pgfqpoint{0.753357in}{0.693860in}}%
\pgfpathlineto{\pgfqpoint{0.752948in}{0.693766in}}%
\pgfpathlineto{\pgfqpoint{0.745108in}{0.691973in}}%
\pgfpathlineto{\pgfqpoint{0.736859in}{0.690105in}}%
\pgfpathlineto{\pgfqpoint{0.728609in}{0.688253in}}%
\pgfpathlineto{\pgfqpoint{0.720926in}{0.686543in}}%
\pgfpathlineto{\pgfqpoint{0.720360in}{0.686419in}}%
\pgfpathlineto{\pgfqpoint{0.712111in}{0.684612in}}%
\pgfpathlineto{\pgfqpoint{0.703862in}{0.682820in}}%
\pgfpathlineto{\pgfqpoint{0.695612in}{0.681045in}}%
\pgfpathlineto{\pgfqpoint{0.687532in}{0.679321in}}%
\pgfpathlineto{\pgfqpoint{0.687363in}{0.679285in}}%
\pgfpathlineto{\pgfqpoint{0.679114in}{0.677552in}}%
\pgfpathlineto{\pgfqpoint{0.670865in}{0.675834in}}%
\pgfpathlineto{\pgfqpoint{0.662616in}{0.674131in}}%
\pgfpathlineto{\pgfqpoint{0.654366in}{0.672442in}}%
\pgfpathlineto{\pgfqpoint{0.652678in}{0.672099in}}%
\pgfpathlineto{\pgfqpoint{0.646117in}{0.670777in}}%
\pgfpathlineto{\pgfqpoint{0.637868in}{0.669129in}}%
\pgfpathlineto{\pgfqpoint{0.629619in}{0.667496in}}%
\pgfpathlineto{\pgfqpoint{0.621369in}{0.665877in}}%
\pgfpathlineto{\pgfqpoint{0.616237in}{0.664877in}}%
\pgfpathlineto{\pgfqpoint{0.613120in}{0.664275in}}%
\pgfpathlineto{\pgfqpoint{0.604871in}{0.662695in}}%
\pgfpathlineto{\pgfqpoint{0.596622in}{0.661128in}}%
\pgfpathlineto{\pgfqpoint{0.588373in}{0.659575in}}%
\pgfpathlineto{\pgfqpoint{0.580123in}{0.658035in}}%
\pgfpathlineto{\pgfqpoint{0.578070in}{0.657654in}}%
\pgfpathlineto{\pgfqpoint{0.571874in}{0.656516in}}%
\pgfpathlineto{\pgfqpoint{0.563625in}{0.655013in}}%
\pgfpathlineto{\pgfqpoint{0.555376in}{0.653522in}}%
\pgfpathlineto{\pgfqpoint{0.547126in}{0.652043in}}%
\pgfpathlineto{\pgfqpoint{0.538877in}{0.650577in}}%
\pgfpathlineto{\pgfqpoint{0.538056in}{0.650432in}}%
\pgfpathlineto{\pgfqpoint{0.530628in}{0.649132in}}%
\pgfpathlineto{\pgfqpoint{0.522379in}{0.647699in}}%
\pgfpathlineto{\pgfqpoint{0.514130in}{0.646278in}}%
\pgfpathlineto{\pgfqpoint{0.505880in}{0.644869in}}%
\pgfpathlineto{\pgfqpoint{0.497631in}{0.643470in}}%
\pgfpathlineto{\pgfqpoint{0.496088in}{0.643210in}}%
\pgfpathlineto{\pgfqpoint{0.489382in}{0.642090in}}%
\pgfpathlineto{\pgfqpoint{0.481133in}{0.640723in}}%
\pgfpathlineto{\pgfqpoint{0.472883in}{0.639367in}}%
\pgfpathlineto{\pgfqpoint{0.464634in}{0.638021in}}%
\pgfpathlineto{\pgfqpoint{0.456385in}{0.636685in}}%
\pgfpathlineto{\pgfqpoint{0.452052in}{0.635988in}}%
\pgfpathlineto{\pgfqpoint{0.448136in}{0.635364in}}%
\pgfpathlineto{\pgfqpoint{0.439887in}{0.634057in}}%
\pgfpathlineto{\pgfqpoint{0.431637in}{0.632760in}}%
\pgfpathlineto{\pgfqpoint{0.423388in}{0.631473in}}%
\pgfpathlineto{\pgfqpoint{0.415139in}{0.630196in}}%
\pgfpathlineto{\pgfqpoint{0.415139in}{0.628766in}}%
\pgfpathlineto{\pgfqpoint{0.415139in}{0.621543in}}%
\pgfpathlineto{\pgfqpoint{0.415139in}{0.614321in}}%
\pgfpathlineto{\pgfqpoint{0.415139in}{0.607099in}}%
\pgfpathlineto{\pgfqpoint{0.415139in}{0.599877in}}%
\pgfpathlineto{\pgfqpoint{0.415139in}{0.592654in}}%
\pgfpathlineto{\pgfqpoint{0.415139in}{0.585432in}}%
\pgfpathlineto{\pgfqpoint{0.415139in}{0.578210in}}%
\pgfpathlineto{\pgfqpoint{0.415139in}{0.570988in}}%
\pgfpathlineto{\pgfqpoint{0.415139in}{0.563766in}}%
\pgfpathlineto{\pgfqpoint{0.415139in}{0.556543in}}%
\pgfpathlineto{\pgfqpoint{0.415139in}{0.549321in}}%
\pgfpathlineto{\pgfqpoint{0.415139in}{0.542099in}}%
\pgfpathlineto{\pgfqpoint{0.415139in}{0.534877in}}%
\pgfpathlineto{\pgfqpoint{0.415139in}{0.527654in}}%
\pgfpathlineto{\pgfqpoint{0.415139in}{0.520432in}}%
\pgfpathlineto{\pgfqpoint{0.415139in}{0.513210in}}%
\pgfpathlineto{\pgfqpoint{0.415139in}{0.505988in}}%
\pgfpathlineto{\pgfqpoint{0.415139in}{0.498766in}}%
\pgfpathlineto{\pgfqpoint{0.415139in}{0.491543in}}%
\pgfpathlineto{\pgfqpoint{0.415139in}{0.484321in}}%
\pgfpathlineto{\pgfqpoint{0.415139in}{0.478901in}}%
\pgfpathlineto{\pgfqpoint{0.423388in}{0.479460in}}%
\pgfpathclose%
\pgfusepath{fill}%
\end{pgfscope}%
\begin{pgfscope}%
\pgfpathrectangle{\pgfqpoint{0.415139in}{0.462654in}}{\pgfqpoint{1.641596in}{1.437222in}}%
\pgfusepath{clip}%
\pgfsetbuttcap%
\pgfsetroundjoin%
\definecolor{currentfill}{rgb}{0.360938,0.566474,0.761630}%
\pgfsetfillcolor{currentfill}%
\pgfsetlinewidth{0.000000pt}%
\definecolor{currentstroke}{rgb}{0.000000,0.000000,0.000000}%
\pgfsetstrokecolor{currentstroke}%
\pgfsetdash{}{0pt}%
\pgfpathmoveto{\pgfqpoint{0.423388in}{0.462654in}}%
\pgfpathlineto{\pgfqpoint{0.431637in}{0.462654in}}%
\pgfpathlineto{\pgfqpoint{0.439887in}{0.462654in}}%
\pgfpathlineto{\pgfqpoint{0.448136in}{0.462654in}}%
\pgfpathlineto{\pgfqpoint{0.456385in}{0.462654in}}%
\pgfpathlineto{\pgfqpoint{0.464634in}{0.462654in}}%
\pgfpathlineto{\pgfqpoint{0.472883in}{0.462654in}}%
\pgfpathlineto{\pgfqpoint{0.481133in}{0.462654in}}%
\pgfpathlineto{\pgfqpoint{0.489382in}{0.462654in}}%
\pgfpathlineto{\pgfqpoint{0.497631in}{0.462654in}}%
\pgfpathlineto{\pgfqpoint{0.505880in}{0.462654in}}%
\pgfpathlineto{\pgfqpoint{0.514130in}{0.462654in}}%
\pgfpathlineto{\pgfqpoint{0.522379in}{0.462654in}}%
\pgfpathlineto{\pgfqpoint{0.530628in}{0.462654in}}%
\pgfpathlineto{\pgfqpoint{0.538877in}{0.462654in}}%
\pgfpathlineto{\pgfqpoint{0.547126in}{0.462654in}}%
\pgfpathlineto{\pgfqpoint{0.555376in}{0.462654in}}%
\pgfpathlineto{\pgfqpoint{0.563625in}{0.462654in}}%
\pgfpathlineto{\pgfqpoint{0.571874in}{0.462654in}}%
\pgfpathlineto{\pgfqpoint{0.580123in}{0.462654in}}%
\pgfpathlineto{\pgfqpoint{0.588373in}{0.462654in}}%
\pgfpathlineto{\pgfqpoint{0.596622in}{0.462654in}}%
\pgfpathlineto{\pgfqpoint{0.604871in}{0.462654in}}%
\pgfpathlineto{\pgfqpoint{0.613120in}{0.462654in}}%
\pgfpathlineto{\pgfqpoint{0.621369in}{0.462654in}}%
\pgfpathlineto{\pgfqpoint{0.629619in}{0.462654in}}%
\pgfpathlineto{\pgfqpoint{0.637868in}{0.462654in}}%
\pgfpathlineto{\pgfqpoint{0.646117in}{0.462654in}}%
\pgfpathlineto{\pgfqpoint{0.654366in}{0.462654in}}%
\pgfpathlineto{\pgfqpoint{0.662616in}{0.462654in}}%
\pgfpathlineto{\pgfqpoint{0.670865in}{0.462654in}}%
\pgfpathlineto{\pgfqpoint{0.679114in}{0.462654in}}%
\pgfpathlineto{\pgfqpoint{0.687363in}{0.462654in}}%
\pgfpathlineto{\pgfqpoint{0.695612in}{0.462654in}}%
\pgfpathlineto{\pgfqpoint{0.703862in}{0.462654in}}%
\pgfpathlineto{\pgfqpoint{0.712111in}{0.462654in}}%
\pgfpathlineto{\pgfqpoint{0.720360in}{0.462654in}}%
\pgfpathlineto{\pgfqpoint{0.728609in}{0.462654in}}%
\pgfpathlineto{\pgfqpoint{0.736859in}{0.462654in}}%
\pgfpathlineto{\pgfqpoint{0.745108in}{0.462654in}}%
\pgfpathlineto{\pgfqpoint{0.753357in}{0.462654in}}%
\pgfpathlineto{\pgfqpoint{0.761606in}{0.462654in}}%
\pgfpathlineto{\pgfqpoint{0.769856in}{0.462654in}}%
\pgfpathlineto{\pgfqpoint{0.778105in}{0.462654in}}%
\pgfpathlineto{\pgfqpoint{0.786354in}{0.462654in}}%
\pgfpathlineto{\pgfqpoint{0.794603in}{0.462654in}}%
\pgfpathlineto{\pgfqpoint{0.802852in}{0.462654in}}%
\pgfpathlineto{\pgfqpoint{0.811102in}{0.462654in}}%
\pgfpathlineto{\pgfqpoint{0.819351in}{0.462654in}}%
\pgfpathlineto{\pgfqpoint{0.827600in}{0.462654in}}%
\pgfpathlineto{\pgfqpoint{0.835849in}{0.462654in}}%
\pgfpathlineto{\pgfqpoint{0.844099in}{0.462654in}}%
\pgfpathlineto{\pgfqpoint{0.852348in}{0.462654in}}%
\pgfpathlineto{\pgfqpoint{0.860597in}{0.462654in}}%
\pgfpathlineto{\pgfqpoint{0.868846in}{0.462654in}}%
\pgfpathlineto{\pgfqpoint{0.877095in}{0.462654in}}%
\pgfpathlineto{\pgfqpoint{0.885345in}{0.462654in}}%
\pgfpathlineto{\pgfqpoint{0.893594in}{0.462654in}}%
\pgfpathlineto{\pgfqpoint{0.901843in}{0.462654in}}%
\pgfpathlineto{\pgfqpoint{0.910092in}{0.462654in}}%
\pgfpathlineto{\pgfqpoint{0.918342in}{0.462654in}}%
\pgfpathlineto{\pgfqpoint{0.926591in}{0.462654in}}%
\pgfpathlineto{\pgfqpoint{0.934840in}{0.462654in}}%
\pgfpathlineto{\pgfqpoint{0.943089in}{0.462654in}}%
\pgfpathlineto{\pgfqpoint{0.951338in}{0.462654in}}%
\pgfpathlineto{\pgfqpoint{0.959588in}{0.462654in}}%
\pgfpathlineto{\pgfqpoint{0.967837in}{0.462654in}}%
\pgfpathlineto{\pgfqpoint{0.976086in}{0.462654in}}%
\pgfpathlineto{\pgfqpoint{0.984335in}{0.462654in}}%
\pgfpathlineto{\pgfqpoint{0.992585in}{0.462654in}}%
\pgfpathlineto{\pgfqpoint{1.000834in}{0.462654in}}%
\pgfpathlineto{\pgfqpoint{1.009083in}{0.462654in}}%
\pgfpathlineto{\pgfqpoint{1.017332in}{0.462654in}}%
\pgfpathlineto{\pgfqpoint{1.025581in}{0.462654in}}%
\pgfpathlineto{\pgfqpoint{1.033831in}{0.462654in}}%
\pgfpathlineto{\pgfqpoint{1.042080in}{0.462654in}}%
\pgfpathlineto{\pgfqpoint{1.050329in}{0.462654in}}%
\pgfpathlineto{\pgfqpoint{1.058578in}{0.462654in}}%
\pgfpathlineto{\pgfqpoint{1.066828in}{0.462654in}}%
\pgfpathlineto{\pgfqpoint{1.075077in}{0.462654in}}%
\pgfpathlineto{\pgfqpoint{1.083326in}{0.462654in}}%
\pgfpathlineto{\pgfqpoint{1.091575in}{0.462654in}}%
\pgfpathlineto{\pgfqpoint{1.099824in}{0.462654in}}%
\pgfpathlineto{\pgfqpoint{1.108074in}{0.462654in}}%
\pgfpathlineto{\pgfqpoint{1.116323in}{0.462654in}}%
\pgfpathlineto{\pgfqpoint{1.124572in}{0.462654in}}%
\pgfpathlineto{\pgfqpoint{1.132821in}{0.462654in}}%
\pgfpathlineto{\pgfqpoint{1.141071in}{0.462654in}}%
\pgfpathlineto{\pgfqpoint{1.149320in}{0.462654in}}%
\pgfpathlineto{\pgfqpoint{1.157569in}{0.462654in}}%
\pgfpathlineto{\pgfqpoint{1.165818in}{0.462654in}}%
\pgfpathlineto{\pgfqpoint{1.174067in}{0.462654in}}%
\pgfpathlineto{\pgfqpoint{1.182317in}{0.462654in}}%
\pgfpathlineto{\pgfqpoint{1.190566in}{0.462654in}}%
\pgfpathlineto{\pgfqpoint{1.198815in}{0.462654in}}%
\pgfpathlineto{\pgfqpoint{1.207064in}{0.462654in}}%
\pgfpathlineto{\pgfqpoint{1.215314in}{0.462654in}}%
\pgfpathlineto{\pgfqpoint{1.223563in}{0.462654in}}%
\pgfpathlineto{\pgfqpoint{1.231812in}{0.462654in}}%
\pgfpathlineto{\pgfqpoint{1.240061in}{0.462654in}}%
\pgfpathlineto{\pgfqpoint{1.248310in}{0.462654in}}%
\pgfpathlineto{\pgfqpoint{1.256560in}{0.462654in}}%
\pgfpathlineto{\pgfqpoint{1.264809in}{0.462654in}}%
\pgfpathlineto{\pgfqpoint{1.273058in}{0.462654in}}%
\pgfpathlineto{\pgfqpoint{1.281307in}{0.462654in}}%
\pgfpathlineto{\pgfqpoint{1.289557in}{0.462654in}}%
\pgfpathlineto{\pgfqpoint{1.297806in}{0.462654in}}%
\pgfpathlineto{\pgfqpoint{1.306055in}{0.462654in}}%
\pgfpathlineto{\pgfqpoint{1.314304in}{0.462654in}}%
\pgfpathlineto{\pgfqpoint{1.322553in}{0.462654in}}%
\pgfpathlineto{\pgfqpoint{1.330803in}{0.462654in}}%
\pgfpathlineto{\pgfqpoint{1.339052in}{0.462654in}}%
\pgfpathlineto{\pgfqpoint{1.347301in}{0.462654in}}%
\pgfpathlineto{\pgfqpoint{1.355550in}{0.462654in}}%
\pgfpathlineto{\pgfqpoint{1.363800in}{0.462654in}}%
\pgfpathlineto{\pgfqpoint{1.372049in}{0.462654in}}%
\pgfpathlineto{\pgfqpoint{1.380298in}{0.462654in}}%
\pgfpathlineto{\pgfqpoint{1.388547in}{0.462654in}}%
\pgfpathlineto{\pgfqpoint{1.396797in}{0.462654in}}%
\pgfpathlineto{\pgfqpoint{1.405046in}{0.462654in}}%
\pgfpathlineto{\pgfqpoint{1.413295in}{0.462654in}}%
\pgfpathlineto{\pgfqpoint{1.421544in}{0.462654in}}%
\pgfpathlineto{\pgfqpoint{1.429793in}{0.462654in}}%
\pgfpathlineto{\pgfqpoint{1.438043in}{0.462654in}}%
\pgfpathlineto{\pgfqpoint{1.446292in}{0.462654in}}%
\pgfpathlineto{\pgfqpoint{1.454541in}{0.462654in}}%
\pgfpathlineto{\pgfqpoint{1.462790in}{0.462654in}}%
\pgfpathlineto{\pgfqpoint{1.471040in}{0.462654in}}%
\pgfpathlineto{\pgfqpoint{1.479289in}{0.462654in}}%
\pgfpathlineto{\pgfqpoint{1.487538in}{0.462654in}}%
\pgfpathlineto{\pgfqpoint{1.495787in}{0.462654in}}%
\pgfpathlineto{\pgfqpoint{1.504036in}{0.462654in}}%
\pgfpathlineto{\pgfqpoint{1.512286in}{0.462654in}}%
\pgfpathlineto{\pgfqpoint{1.520535in}{0.462654in}}%
\pgfpathlineto{\pgfqpoint{1.528784in}{0.462654in}}%
\pgfpathlineto{\pgfqpoint{1.537033in}{0.462654in}}%
\pgfpathlineto{\pgfqpoint{1.545283in}{0.462654in}}%
\pgfpathlineto{\pgfqpoint{1.553532in}{0.462654in}}%
\pgfpathlineto{\pgfqpoint{1.561781in}{0.462654in}}%
\pgfpathlineto{\pgfqpoint{1.570030in}{0.462654in}}%
\pgfpathlineto{\pgfqpoint{1.578279in}{0.462654in}}%
\pgfpathlineto{\pgfqpoint{1.586529in}{0.462654in}}%
\pgfpathlineto{\pgfqpoint{1.594778in}{0.462654in}}%
\pgfpathlineto{\pgfqpoint{1.603027in}{0.462654in}}%
\pgfpathlineto{\pgfqpoint{1.611276in}{0.462654in}}%
\pgfpathlineto{\pgfqpoint{1.619526in}{0.462654in}}%
\pgfpathlineto{\pgfqpoint{1.627775in}{0.462654in}}%
\pgfpathlineto{\pgfqpoint{1.636024in}{0.462654in}}%
\pgfpathlineto{\pgfqpoint{1.644273in}{0.462654in}}%
\pgfpathlineto{\pgfqpoint{1.652522in}{0.462654in}}%
\pgfpathlineto{\pgfqpoint{1.660772in}{0.462654in}}%
\pgfpathlineto{\pgfqpoint{1.669021in}{0.462654in}}%
\pgfpathlineto{\pgfqpoint{1.677270in}{0.462654in}}%
\pgfpathlineto{\pgfqpoint{1.685519in}{0.462654in}}%
\pgfpathlineto{\pgfqpoint{1.693769in}{0.462654in}}%
\pgfpathlineto{\pgfqpoint{1.702018in}{0.462654in}}%
\pgfpathlineto{\pgfqpoint{1.710267in}{0.462654in}}%
\pgfpathlineto{\pgfqpoint{1.718516in}{0.462654in}}%
\pgfpathlineto{\pgfqpoint{1.726765in}{0.462654in}}%
\pgfpathlineto{\pgfqpoint{1.735015in}{0.462654in}}%
\pgfpathlineto{\pgfqpoint{1.743264in}{0.462654in}}%
\pgfpathlineto{\pgfqpoint{1.751513in}{0.462654in}}%
\pgfpathlineto{\pgfqpoint{1.759762in}{0.462654in}}%
\pgfpathlineto{\pgfqpoint{1.768012in}{0.462654in}}%
\pgfpathlineto{\pgfqpoint{1.776261in}{0.462654in}}%
\pgfpathlineto{\pgfqpoint{1.784510in}{0.462654in}}%
\pgfpathlineto{\pgfqpoint{1.792759in}{0.462654in}}%
\pgfpathlineto{\pgfqpoint{1.801008in}{0.462654in}}%
\pgfpathlineto{\pgfqpoint{1.809258in}{0.462654in}}%
\pgfpathlineto{\pgfqpoint{1.817507in}{0.462654in}}%
\pgfpathlineto{\pgfqpoint{1.825756in}{0.462654in}}%
\pgfpathlineto{\pgfqpoint{1.834005in}{0.462654in}}%
\pgfpathlineto{\pgfqpoint{1.842255in}{0.462654in}}%
\pgfpathlineto{\pgfqpoint{1.850504in}{0.462654in}}%
\pgfpathlineto{\pgfqpoint{1.858753in}{0.462654in}}%
\pgfpathlineto{\pgfqpoint{1.867002in}{0.462654in}}%
\pgfpathlineto{\pgfqpoint{1.875251in}{0.462654in}}%
\pgfpathlineto{\pgfqpoint{1.883501in}{0.462654in}}%
\pgfpathlineto{\pgfqpoint{1.891750in}{0.462654in}}%
\pgfpathlineto{\pgfqpoint{1.899999in}{0.462654in}}%
\pgfpathlineto{\pgfqpoint{1.908248in}{0.462654in}}%
\pgfpathlineto{\pgfqpoint{1.916498in}{0.462654in}}%
\pgfpathlineto{\pgfqpoint{1.924747in}{0.462654in}}%
\pgfpathlineto{\pgfqpoint{1.932996in}{0.462654in}}%
\pgfpathlineto{\pgfqpoint{1.941245in}{0.462654in}}%
\pgfpathlineto{\pgfqpoint{1.949494in}{0.462654in}}%
\pgfpathlineto{\pgfqpoint{1.957744in}{0.462654in}}%
\pgfpathlineto{\pgfqpoint{1.965993in}{0.462654in}}%
\pgfpathlineto{\pgfqpoint{1.974242in}{0.462654in}}%
\pgfpathlineto{\pgfqpoint{1.982491in}{0.462654in}}%
\pgfpathlineto{\pgfqpoint{1.990741in}{0.462654in}}%
\pgfpathlineto{\pgfqpoint{1.998990in}{0.462654in}}%
\pgfpathlineto{\pgfqpoint{2.007239in}{0.462654in}}%
\pgfpathlineto{\pgfqpoint{2.015488in}{0.462654in}}%
\pgfpathlineto{\pgfqpoint{2.023737in}{0.462654in}}%
\pgfpathlineto{\pgfqpoint{2.031987in}{0.462654in}}%
\pgfpathlineto{\pgfqpoint{2.040236in}{0.462654in}}%
\pgfpathlineto{\pgfqpoint{2.048485in}{0.462654in}}%
\pgfpathlineto{\pgfqpoint{2.056734in}{0.462654in}}%
\pgfpathlineto{\pgfqpoint{2.056734in}{0.469877in}}%
\pgfpathlineto{\pgfqpoint{2.056734in}{0.477099in}}%
\pgfpathlineto{\pgfqpoint{2.056734in}{0.484321in}}%
\pgfpathlineto{\pgfqpoint{2.056734in}{0.491543in}}%
\pgfpathlineto{\pgfqpoint{2.056734in}{0.498766in}}%
\pgfpathlineto{\pgfqpoint{2.056734in}{0.505988in}}%
\pgfpathlineto{\pgfqpoint{2.056734in}{0.513210in}}%
\pgfpathlineto{\pgfqpoint{2.056734in}{0.520432in}}%
\pgfpathlineto{\pgfqpoint{2.056734in}{0.527654in}}%
\pgfpathlineto{\pgfqpoint{2.056734in}{0.534877in}}%
\pgfpathlineto{\pgfqpoint{2.056734in}{0.542099in}}%
\pgfpathlineto{\pgfqpoint{2.056734in}{0.549321in}}%
\pgfpathlineto{\pgfqpoint{2.056734in}{0.556543in}}%
\pgfpathlineto{\pgfqpoint{2.056734in}{0.563766in}}%
\pgfpathlineto{\pgfqpoint{2.056734in}{0.570988in}}%
\pgfpathlineto{\pgfqpoint{2.056734in}{0.578210in}}%
\pgfpathlineto{\pgfqpoint{2.056734in}{0.585432in}}%
\pgfpathlineto{\pgfqpoint{2.056734in}{0.592654in}}%
\pgfpathlineto{\pgfqpoint{2.056734in}{0.599877in}}%
\pgfpathlineto{\pgfqpoint{2.056734in}{0.607099in}}%
\pgfpathlineto{\pgfqpoint{2.056734in}{0.614321in}}%
\pgfpathlineto{\pgfqpoint{2.056734in}{0.621543in}}%
\pgfpathlineto{\pgfqpoint{2.056734in}{0.628766in}}%
\pgfpathlineto{\pgfqpoint{2.056734in}{0.635988in}}%
\pgfpathlineto{\pgfqpoint{2.056734in}{0.643210in}}%
\pgfpathlineto{\pgfqpoint{2.056734in}{0.650432in}}%
\pgfpathlineto{\pgfqpoint{2.056734in}{0.657654in}}%
\pgfpathlineto{\pgfqpoint{2.056734in}{0.664877in}}%
\pgfpathlineto{\pgfqpoint{2.056734in}{0.672099in}}%
\pgfpathlineto{\pgfqpoint{2.056734in}{0.679321in}}%
\pgfpathlineto{\pgfqpoint{2.056734in}{0.686543in}}%
\pgfpathlineto{\pgfqpoint{2.056734in}{0.693766in}}%
\pgfpathlineto{\pgfqpoint{2.056734in}{0.700988in}}%
\pgfpathlineto{\pgfqpoint{2.056734in}{0.708210in}}%
\pgfpathlineto{\pgfqpoint{2.056734in}{0.715432in}}%
\pgfpathlineto{\pgfqpoint{2.056734in}{0.722654in}}%
\pgfpathlineto{\pgfqpoint{2.056734in}{0.729877in}}%
\pgfpathlineto{\pgfqpoint{2.056734in}{0.737099in}}%
\pgfpathlineto{\pgfqpoint{2.056734in}{0.744321in}}%
\pgfpathlineto{\pgfqpoint{2.056734in}{0.751543in}}%
\pgfpathlineto{\pgfqpoint{2.056734in}{0.758766in}}%
\pgfpathlineto{\pgfqpoint{2.056734in}{0.765988in}}%
\pgfpathlineto{\pgfqpoint{2.056734in}{0.773210in}}%
\pgfpathlineto{\pgfqpoint{2.056734in}{0.780432in}}%
\pgfpathlineto{\pgfqpoint{2.056734in}{0.787654in}}%
\pgfpathlineto{\pgfqpoint{2.056734in}{0.794877in}}%
\pgfpathlineto{\pgfqpoint{2.056734in}{0.802099in}}%
\pgfpathlineto{\pgfqpoint{2.056734in}{0.809321in}}%
\pgfpathlineto{\pgfqpoint{2.056734in}{0.816543in}}%
\pgfpathlineto{\pgfqpoint{2.056734in}{0.823766in}}%
\pgfpathlineto{\pgfqpoint{2.056734in}{0.830988in}}%
\pgfpathlineto{\pgfqpoint{2.056734in}{0.838210in}}%
\pgfpathlineto{\pgfqpoint{2.056734in}{0.845432in}}%
\pgfpathlineto{\pgfqpoint{2.056734in}{0.852654in}}%
\pgfpathlineto{\pgfqpoint{2.056734in}{0.859877in}}%
\pgfpathlineto{\pgfqpoint{2.056734in}{0.867099in}}%
\pgfpathlineto{\pgfqpoint{2.056734in}{0.874321in}}%
\pgfpathlineto{\pgfqpoint{2.056734in}{0.881543in}}%
\pgfpathlineto{\pgfqpoint{2.056734in}{0.888766in}}%
\pgfpathlineto{\pgfqpoint{2.056734in}{0.895988in}}%
\pgfpathlineto{\pgfqpoint{2.056734in}{0.903210in}}%
\pgfpathlineto{\pgfqpoint{2.056734in}{0.910432in}}%
\pgfpathlineto{\pgfqpoint{2.056734in}{0.917654in}}%
\pgfpathlineto{\pgfqpoint{2.056734in}{0.924877in}}%
\pgfpathlineto{\pgfqpoint{2.056734in}{0.932099in}}%
\pgfpathlineto{\pgfqpoint{2.056734in}{0.939321in}}%
\pgfpathlineto{\pgfqpoint{2.056734in}{0.946543in}}%
\pgfpathlineto{\pgfqpoint{2.056734in}{0.953766in}}%
\pgfpathlineto{\pgfqpoint{2.056734in}{0.960988in}}%
\pgfpathlineto{\pgfqpoint{2.056734in}{0.968210in}}%
\pgfpathlineto{\pgfqpoint{2.056734in}{0.975432in}}%
\pgfpathlineto{\pgfqpoint{2.056734in}{0.982654in}}%
\pgfpathlineto{\pgfqpoint{2.056734in}{0.989877in}}%
\pgfpathlineto{\pgfqpoint{2.056734in}{0.997099in}}%
\pgfpathlineto{\pgfqpoint{2.056734in}{1.004321in}}%
\pgfpathlineto{\pgfqpoint{2.056734in}{1.011543in}}%
\pgfpathlineto{\pgfqpoint{2.056734in}{1.018766in}}%
\pgfpathlineto{\pgfqpoint{2.056734in}{1.025988in}}%
\pgfpathlineto{\pgfqpoint{2.056734in}{1.033210in}}%
\pgfpathlineto{\pgfqpoint{2.056734in}{1.040432in}}%
\pgfpathlineto{\pgfqpoint{2.056734in}{1.047654in}}%
\pgfpathlineto{\pgfqpoint{2.056734in}{1.054877in}}%
\pgfpathlineto{\pgfqpoint{2.056734in}{1.062099in}}%
\pgfpathlineto{\pgfqpoint{2.056734in}{1.069321in}}%
\pgfpathlineto{\pgfqpoint{2.056734in}{1.076543in}}%
\pgfpathlineto{\pgfqpoint{2.056734in}{1.083766in}}%
\pgfpathlineto{\pgfqpoint{2.056734in}{1.090988in}}%
\pgfpathlineto{\pgfqpoint{2.056734in}{1.098210in}}%
\pgfpathlineto{\pgfqpoint{2.056734in}{1.105432in}}%
\pgfpathlineto{\pgfqpoint{2.056734in}{1.112654in}}%
\pgfpathlineto{\pgfqpoint{2.056734in}{1.119877in}}%
\pgfpathlineto{\pgfqpoint{2.056734in}{1.127099in}}%
\pgfpathlineto{\pgfqpoint{2.056734in}{1.134321in}}%
\pgfpathlineto{\pgfqpoint{2.056734in}{1.141543in}}%
\pgfpathlineto{\pgfqpoint{2.056734in}{1.148766in}}%
\pgfpathlineto{\pgfqpoint{2.056734in}{1.155988in}}%
\pgfpathlineto{\pgfqpoint{2.056734in}{1.163210in}}%
\pgfpathlineto{\pgfqpoint{2.056734in}{1.170432in}}%
\pgfpathlineto{\pgfqpoint{2.056734in}{1.177654in}}%
\pgfpathlineto{\pgfqpoint{2.056734in}{1.184877in}}%
\pgfpathlineto{\pgfqpoint{2.056734in}{1.192099in}}%
\pgfpathlineto{\pgfqpoint{2.056734in}{1.199321in}}%
\pgfpathlineto{\pgfqpoint{2.056734in}{1.206543in}}%
\pgfpathlineto{\pgfqpoint{2.056734in}{1.213766in}}%
\pgfpathlineto{\pgfqpoint{2.056734in}{1.220988in}}%
\pgfpathlineto{\pgfqpoint{2.056734in}{1.228210in}}%
\pgfpathlineto{\pgfqpoint{2.056734in}{1.235432in}}%
\pgfpathlineto{\pgfqpoint{2.056734in}{1.242654in}}%
\pgfpathlineto{\pgfqpoint{2.056734in}{1.249877in}}%
\pgfpathlineto{\pgfqpoint{2.056734in}{1.257099in}}%
\pgfpathlineto{\pgfqpoint{2.056734in}{1.264321in}}%
\pgfpathlineto{\pgfqpoint{2.056734in}{1.271543in}}%
\pgfpathlineto{\pgfqpoint{2.056734in}{1.278766in}}%
\pgfpathlineto{\pgfqpoint{2.056734in}{1.285988in}}%
\pgfpathlineto{\pgfqpoint{2.056734in}{1.293210in}}%
\pgfpathlineto{\pgfqpoint{2.056734in}{1.300432in}}%
\pgfpathlineto{\pgfqpoint{2.056734in}{1.307654in}}%
\pgfpathlineto{\pgfqpoint{2.056734in}{1.314877in}}%
\pgfpathlineto{\pgfqpoint{2.056734in}{1.322099in}}%
\pgfpathlineto{\pgfqpoint{2.056734in}{1.329321in}}%
\pgfpathlineto{\pgfqpoint{2.056734in}{1.336543in}}%
\pgfpathlineto{\pgfqpoint{2.056734in}{1.343766in}}%
\pgfpathlineto{\pgfqpoint{2.056734in}{1.350988in}}%
\pgfpathlineto{\pgfqpoint{2.056734in}{1.358210in}}%
\pgfpathlineto{\pgfqpoint{2.056734in}{1.365432in}}%
\pgfpathlineto{\pgfqpoint{2.056734in}{1.372654in}}%
\pgfpathlineto{\pgfqpoint{2.056734in}{1.379877in}}%
\pgfpathlineto{\pgfqpoint{2.056734in}{1.387099in}}%
\pgfpathlineto{\pgfqpoint{2.056734in}{1.394321in}}%
\pgfpathlineto{\pgfqpoint{2.056734in}{1.401543in}}%
\pgfpathlineto{\pgfqpoint{2.056734in}{1.408766in}}%
\pgfpathlineto{\pgfqpoint{2.056734in}{1.415988in}}%
\pgfpathlineto{\pgfqpoint{2.056734in}{1.423210in}}%
\pgfpathlineto{\pgfqpoint{2.056734in}{1.430432in}}%
\pgfpathlineto{\pgfqpoint{2.056734in}{1.437654in}}%
\pgfpathlineto{\pgfqpoint{2.056734in}{1.444877in}}%
\pgfpathlineto{\pgfqpoint{2.056734in}{1.452099in}}%
\pgfpathlineto{\pgfqpoint{2.056734in}{1.459321in}}%
\pgfpathlineto{\pgfqpoint{2.056734in}{1.466543in}}%
\pgfpathlineto{\pgfqpoint{2.056734in}{1.473766in}}%
\pgfpathlineto{\pgfqpoint{2.056734in}{1.480988in}}%
\pgfpathlineto{\pgfqpoint{2.056734in}{1.488210in}}%
\pgfpathlineto{\pgfqpoint{2.056734in}{1.495432in}}%
\pgfpathlineto{\pgfqpoint{2.056734in}{1.502654in}}%
\pgfpathlineto{\pgfqpoint{2.056734in}{1.509877in}}%
\pgfpathlineto{\pgfqpoint{2.056734in}{1.517099in}}%
\pgfpathlineto{\pgfqpoint{2.056734in}{1.524321in}}%
\pgfpathlineto{\pgfqpoint{2.056734in}{1.531543in}}%
\pgfpathlineto{\pgfqpoint{2.056734in}{1.538766in}}%
\pgfpathlineto{\pgfqpoint{2.056734in}{1.545988in}}%
\pgfpathlineto{\pgfqpoint{2.056734in}{1.553210in}}%
\pgfpathlineto{\pgfqpoint{2.056734in}{1.560432in}}%
\pgfpathlineto{\pgfqpoint{2.056734in}{1.567654in}}%
\pgfpathlineto{\pgfqpoint{2.056734in}{1.574877in}}%
\pgfpathlineto{\pgfqpoint{2.056734in}{1.582099in}}%
\pgfpathlineto{\pgfqpoint{2.056734in}{1.589321in}}%
\pgfpathlineto{\pgfqpoint{2.056734in}{1.596543in}}%
\pgfpathlineto{\pgfqpoint{2.056734in}{1.603766in}}%
\pgfpathlineto{\pgfqpoint{2.056734in}{1.610988in}}%
\pgfpathlineto{\pgfqpoint{2.056734in}{1.618210in}}%
\pgfpathlineto{\pgfqpoint{2.056734in}{1.625432in}}%
\pgfpathlineto{\pgfqpoint{2.056734in}{1.632654in}}%
\pgfpathlineto{\pgfqpoint{2.056734in}{1.639877in}}%
\pgfpathlineto{\pgfqpoint{2.056734in}{1.647099in}}%
\pgfpathlineto{\pgfqpoint{2.056734in}{1.654321in}}%
\pgfpathlineto{\pgfqpoint{2.056734in}{1.661543in}}%
\pgfpathlineto{\pgfqpoint{2.056734in}{1.668766in}}%
\pgfpathlineto{\pgfqpoint{2.056734in}{1.675988in}}%
\pgfpathlineto{\pgfqpoint{2.056734in}{1.683210in}}%
\pgfpathlineto{\pgfqpoint{2.056734in}{1.690432in}}%
\pgfpathlineto{\pgfqpoint{2.056734in}{1.697654in}}%
\pgfpathlineto{\pgfqpoint{2.056734in}{1.704877in}}%
\pgfpathlineto{\pgfqpoint{2.056734in}{1.712099in}}%
\pgfpathlineto{\pgfqpoint{2.056734in}{1.719321in}}%
\pgfpathlineto{\pgfqpoint{2.056734in}{1.726543in}}%
\pgfpathlineto{\pgfqpoint{2.056734in}{1.733766in}}%
\pgfpathlineto{\pgfqpoint{2.056734in}{1.740988in}}%
\pgfpathlineto{\pgfqpoint{2.056734in}{1.748210in}}%
\pgfpathlineto{\pgfqpoint{2.056734in}{1.755432in}}%
\pgfpathlineto{\pgfqpoint{2.056734in}{1.762654in}}%
\pgfpathlineto{\pgfqpoint{2.056734in}{1.769877in}}%
\pgfpathlineto{\pgfqpoint{2.056734in}{1.777099in}}%
\pgfpathlineto{\pgfqpoint{2.056734in}{1.784321in}}%
\pgfpathlineto{\pgfqpoint{2.056734in}{1.791543in}}%
\pgfpathlineto{\pgfqpoint{2.056734in}{1.798766in}}%
\pgfpathlineto{\pgfqpoint{2.056734in}{1.805988in}}%
\pgfpathlineto{\pgfqpoint{2.056734in}{1.813210in}}%
\pgfpathlineto{\pgfqpoint{2.056734in}{1.820432in}}%
\pgfpathlineto{\pgfqpoint{2.056734in}{1.827654in}}%
\pgfpathlineto{\pgfqpoint{2.056734in}{1.834877in}}%
\pgfpathlineto{\pgfqpoint{2.056734in}{1.842099in}}%
\pgfpathlineto{\pgfqpoint{2.056734in}{1.849321in}}%
\pgfpathlineto{\pgfqpoint{2.056734in}{1.856543in}}%
\pgfpathlineto{\pgfqpoint{2.056734in}{1.863766in}}%
\pgfpathlineto{\pgfqpoint{2.056734in}{1.870988in}}%
\pgfpathlineto{\pgfqpoint{2.056734in}{1.878210in}}%
\pgfpathlineto{\pgfqpoint{2.056734in}{1.885432in}}%
\pgfpathlineto{\pgfqpoint{2.056734in}{1.892654in}}%
\pgfpathlineto{\pgfqpoint{2.056734in}{1.899877in}}%
\pgfpathlineto{\pgfqpoint{2.048485in}{1.899877in}}%
\pgfpathlineto{\pgfqpoint{2.040236in}{1.899877in}}%
\pgfpathlineto{\pgfqpoint{2.031987in}{1.899877in}}%
\pgfpathlineto{\pgfqpoint{2.023737in}{1.899877in}}%
\pgfpathlineto{\pgfqpoint{2.015488in}{1.899877in}}%
\pgfpathlineto{\pgfqpoint{2.011599in}{1.899877in}}%
\pgfpathlineto{\pgfqpoint{2.011060in}{1.892654in}}%
\pgfpathlineto{\pgfqpoint{2.010519in}{1.885432in}}%
\pgfpathlineto{\pgfqpoint{2.009976in}{1.878210in}}%
\pgfpathlineto{\pgfqpoint{2.009430in}{1.870988in}}%
\pgfpathlineto{\pgfqpoint{2.008883in}{1.863766in}}%
\pgfpathlineto{\pgfqpoint{2.008333in}{1.856543in}}%
\pgfpathlineto{\pgfqpoint{2.007781in}{1.849321in}}%
\pgfpathlineto{\pgfqpoint{2.007239in}{1.842273in}}%
\pgfpathlineto{\pgfqpoint{2.007226in}{1.842099in}}%
\pgfpathlineto{\pgfqpoint{2.006663in}{1.834877in}}%
\pgfpathlineto{\pgfqpoint{2.006098in}{1.827654in}}%
\pgfpathlineto{\pgfqpoint{2.005531in}{1.820432in}}%
\pgfpathlineto{\pgfqpoint{2.004960in}{1.813210in}}%
\pgfpathlineto{\pgfqpoint{2.004386in}{1.805988in}}%
\pgfpathlineto{\pgfqpoint{2.003809in}{1.798766in}}%
\pgfpathlineto{\pgfqpoint{2.003229in}{1.791543in}}%
\pgfpathlineto{\pgfqpoint{2.002646in}{1.784321in}}%
\pgfpathlineto{\pgfqpoint{2.002059in}{1.777099in}}%
\pgfpathlineto{\pgfqpoint{2.001468in}{1.769877in}}%
\pgfpathlineto{\pgfqpoint{2.000874in}{1.762654in}}%
\pgfpathlineto{\pgfqpoint{2.000275in}{1.755432in}}%
\pgfpathlineto{\pgfqpoint{1.999673in}{1.748210in}}%
\pgfpathlineto{\pgfqpoint{1.999066in}{1.740988in}}%
\pgfpathlineto{\pgfqpoint{1.998990in}{1.740081in}}%
\pgfpathlineto{\pgfqpoint{1.998451in}{1.733766in}}%
\pgfpathlineto{\pgfqpoint{1.997831in}{1.726543in}}%
\pgfpathlineto{\pgfqpoint{1.997207in}{1.719321in}}%
\pgfpathlineto{\pgfqpoint{1.996578in}{1.712099in}}%
\pgfpathlineto{\pgfqpoint{1.995943in}{1.704877in}}%
\pgfpathlineto{\pgfqpoint{1.995304in}{1.697654in}}%
\pgfpathlineto{\pgfqpoint{1.994660in}{1.690432in}}%
\pgfpathlineto{\pgfqpoint{1.994011in}{1.683210in}}%
\pgfpathlineto{\pgfqpoint{1.993356in}{1.675988in}}%
\pgfpathlineto{\pgfqpoint{1.992695in}{1.668766in}}%
\pgfpathlineto{\pgfqpoint{1.992029in}{1.661543in}}%
\pgfpathlineto{\pgfqpoint{1.991356in}{1.654321in}}%
\pgfpathlineto{\pgfqpoint{1.990741in}{1.647767in}}%
\pgfpathlineto{\pgfqpoint{1.990677in}{1.647099in}}%
\pgfpathlineto{\pgfqpoint{1.989987in}{1.639877in}}%
\pgfpathlineto{\pgfqpoint{1.989291in}{1.632654in}}%
\pgfpathlineto{\pgfqpoint{1.988588in}{1.625432in}}%
\pgfpathlineto{\pgfqpoint{1.987879in}{1.618210in}}%
\pgfpathlineto{\pgfqpoint{1.987162in}{1.610988in}}%
\pgfpathlineto{\pgfqpoint{1.986438in}{1.603766in}}%
\pgfpathlineto{\pgfqpoint{1.985706in}{1.596543in}}%
\pgfpathlineto{\pgfqpoint{1.984967in}{1.589321in}}%
\pgfpathlineto{\pgfqpoint{1.984219in}{1.582099in}}%
\pgfpathlineto{\pgfqpoint{1.983464in}{1.574877in}}%
\pgfpathlineto{\pgfqpoint{1.982700in}{1.567654in}}%
\pgfpathlineto{\pgfqpoint{1.982491in}{1.565691in}}%
\pgfpathlineto{\pgfqpoint{1.981925in}{1.560432in}}%
\pgfpathlineto{\pgfqpoint{1.981138in}{1.553210in}}%
\pgfpathlineto{\pgfqpoint{1.980343in}{1.545988in}}%
\pgfpathlineto{\pgfqpoint{1.979539in}{1.538766in}}%
\pgfpathlineto{\pgfqpoint{1.978725in}{1.531543in}}%
\pgfpathlineto{\pgfqpoint{1.977901in}{1.524321in}}%
\pgfpathlineto{\pgfqpoint{1.977067in}{1.517099in}}%
\pgfpathlineto{\pgfqpoint{1.976223in}{1.509877in}}%
\pgfpathlineto{\pgfqpoint{1.975368in}{1.502654in}}%
\pgfpathlineto{\pgfqpoint{1.974502in}{1.495432in}}%
\pgfpathlineto{\pgfqpoint{1.974242in}{1.493279in}}%
\pgfpathlineto{\pgfqpoint{1.973622in}{1.488210in}}%
\pgfpathlineto{\pgfqpoint{1.972728in}{1.480988in}}%
\pgfpathlineto{\pgfqpoint{1.971823in}{1.473766in}}%
\pgfpathlineto{\pgfqpoint{1.970905in}{1.466543in}}%
\pgfpathlineto{\pgfqpoint{1.969975in}{1.459321in}}%
\pgfpathlineto{\pgfqpoint{1.969032in}{1.452099in}}%
\pgfpathlineto{\pgfqpoint{1.968076in}{1.444877in}}%
\pgfpathlineto{\pgfqpoint{1.967106in}{1.437654in}}%
\pgfpathlineto{\pgfqpoint{1.966123in}{1.430432in}}%
\pgfpathlineto{\pgfqpoint{1.965993in}{1.429489in}}%
\pgfpathlineto{\pgfqpoint{1.965120in}{1.423210in}}%
\pgfpathlineto{\pgfqpoint{1.964102in}{1.415988in}}%
\pgfpathlineto{\pgfqpoint{1.963069in}{1.408766in}}%
\pgfpathlineto{\pgfqpoint{1.962021in}{1.401543in}}%
\pgfpathlineto{\pgfqpoint{1.960957in}{1.394321in}}%
\pgfpathlineto{\pgfqpoint{1.959877in}{1.387099in}}%
\pgfpathlineto{\pgfqpoint{1.958780in}{1.379877in}}%
\pgfpathlineto{\pgfqpoint{1.957744in}{1.373159in}}%
\pgfpathlineto{\pgfqpoint{1.957665in}{1.372654in}}%
\pgfpathlineto{\pgfqpoint{1.956527in}{1.365432in}}%
\pgfpathlineto{\pgfqpoint{1.955372in}{1.358210in}}%
\pgfpathlineto{\pgfqpoint{1.954197in}{1.350988in}}%
\pgfpathlineto{\pgfqpoint{1.953004in}{1.343766in}}%
\pgfpathlineto{\pgfqpoint{1.951791in}{1.336543in}}%
\pgfpathlineto{\pgfqpoint{1.950558in}{1.329321in}}%
\pgfpathlineto{\pgfqpoint{1.949494in}{1.323190in}}%
\pgfpathlineto{\pgfqpoint{1.949303in}{1.322099in}}%
\pgfpathlineto{\pgfqpoint{1.948022in}{1.314877in}}%
\pgfpathlineto{\pgfqpoint{1.946719in}{1.307654in}}%
\pgfpathlineto{\pgfqpoint{1.945393in}{1.300432in}}%
\pgfpathlineto{\pgfqpoint{1.944044in}{1.293210in}}%
\pgfpathlineto{\pgfqpoint{1.942672in}{1.285988in}}%
\pgfpathlineto{\pgfqpoint{1.941275in}{1.278766in}}%
\pgfpathlineto{\pgfqpoint{1.941245in}{1.278616in}}%
\pgfpathlineto{\pgfqpoint{1.939846in}{1.271543in}}%
\pgfpathlineto{\pgfqpoint{1.938392in}{1.264321in}}%
\pgfpathlineto{\pgfqpoint{1.936911in}{1.257099in}}%
\pgfpathlineto{\pgfqpoint{1.935403in}{1.249877in}}%
\pgfpathlineto{\pgfqpoint{1.933867in}{1.242654in}}%
\pgfpathlineto{\pgfqpoint{1.932996in}{1.238624in}}%
\pgfpathlineto{\pgfqpoint{1.932298in}{1.235432in}}%
\pgfpathlineto{\pgfqpoint{1.930697in}{1.228210in}}%
\pgfpathlineto{\pgfqpoint{1.929064in}{1.220988in}}%
\pgfpathlineto{\pgfqpoint{1.927401in}{1.213766in}}%
\pgfpathlineto{\pgfqpoint{1.925704in}{1.206543in}}%
\pgfpathlineto{\pgfqpoint{1.924747in}{1.202531in}}%
\pgfpathlineto{\pgfqpoint{1.923972in}{1.199321in}}%
\pgfpathlineto{\pgfqpoint{1.922201in}{1.192099in}}%
\pgfpathlineto{\pgfqpoint{1.920394in}{1.184877in}}%
\pgfpathlineto{\pgfqpoint{1.918551in}{1.177654in}}%
\pgfpathlineto{\pgfqpoint{1.916671in}{1.170432in}}%
\pgfpathlineto{\pgfqpoint{1.916498in}{1.169776in}}%
\pgfpathlineto{\pgfqpoint{1.914745in}{1.163210in}}%
\pgfpathlineto{\pgfqpoint{1.912779in}{1.155988in}}%
\pgfpathlineto{\pgfqpoint{1.910771in}{1.148766in}}%
\pgfpathlineto{\pgfqpoint{1.908721in}{1.141543in}}%
\pgfpathlineto{\pgfqpoint{1.908248in}{1.139901in}}%
\pgfpathlineto{\pgfqpoint{1.906622in}{1.134321in}}%
\pgfpathlineto{\pgfqpoint{1.904476in}{1.127099in}}%
\pgfpathlineto{\pgfqpoint{1.902283in}{1.119877in}}%
\pgfpathlineto{\pgfqpoint{1.900042in}{1.112654in}}%
\pgfpathlineto{\pgfqpoint{1.899999in}{1.112519in}}%
\pgfpathlineto{\pgfqpoint{1.897744in}{1.105432in}}%
\pgfpathlineto{\pgfqpoint{1.895395in}{1.098210in}}%
\pgfpathlineto{\pgfqpoint{1.892993in}{1.090988in}}%
\pgfpathlineto{\pgfqpoint{1.891750in}{1.087318in}}%
\pgfpathlineto{\pgfqpoint{1.890532in}{1.083766in}}%
\pgfpathlineto{\pgfqpoint{1.888011in}{1.076543in}}%
\pgfpathlineto{\pgfqpoint{1.885432in}{1.069321in}}%
\pgfpathlineto{\pgfqpoint{1.883501in}{1.064023in}}%
\pgfpathlineto{\pgfqpoint{1.882790in}{1.062099in}}%
\pgfpathlineto{\pgfqpoint{1.880081in}{1.054877in}}%
\pgfpathlineto{\pgfqpoint{1.877307in}{1.047654in}}%
\pgfpathlineto{\pgfqpoint{1.875251in}{1.042415in}}%
\pgfpathlineto{\pgfqpoint{1.874464in}{1.040432in}}%
\pgfpathlineto{\pgfqpoint{1.871546in}{1.033210in}}%
\pgfpathlineto{\pgfqpoint{1.868557in}{1.025988in}}%
\pgfpathlineto{\pgfqpoint{1.867002in}{1.022303in}}%
\pgfpathlineto{\pgfqpoint{1.865490in}{1.018766in}}%
\pgfpathlineto{\pgfqpoint{1.862342in}{1.011543in}}%
\pgfpathlineto{\pgfqpoint{1.859115in}{1.004321in}}%
\pgfpathlineto{\pgfqpoint{1.858753in}{1.003523in}}%
\pgfpathlineto{\pgfqpoint{1.855799in}{0.997099in}}%
\pgfpathlineto{\pgfqpoint{1.852396in}{0.989877in}}%
\pgfpathlineto{\pgfqpoint{1.850504in}{0.985944in}}%
\pgfpathlineto{\pgfqpoint{1.848900in}{0.982654in}}%
\pgfpathlineto{\pgfqpoint{1.845309in}{0.975432in}}%
\pgfpathlineto{\pgfqpoint{1.842255in}{0.969438in}}%
\pgfpathlineto{\pgfqpoint{1.841620in}{0.968210in}}%
\pgfpathlineto{\pgfqpoint{1.837825in}{0.960988in}}%
\pgfpathlineto{\pgfqpoint{1.834005in}{0.953908in}}%
\pgfpathlineto{\pgfqpoint{1.833927in}{0.953766in}}%
\pgfpathlineto{\pgfqpoint{1.829914in}{0.946543in}}%
\pgfpathlineto{\pgfqpoint{1.825789in}{0.939321in}}%
\pgfpathlineto{\pgfqpoint{1.825756in}{0.939265in}}%
\pgfpathlineto{\pgfqpoint{1.821539in}{0.932099in}}%
\pgfpathlineto{\pgfqpoint{1.817507in}{0.925430in}}%
\pgfpathlineto{\pgfqpoint{1.817168in}{0.924877in}}%
\pgfpathlineto{\pgfqpoint{1.812662in}{0.917654in}}%
\pgfpathlineto{\pgfqpoint{1.809258in}{0.912333in}}%
\pgfpathlineto{\pgfqpoint{1.808024in}{0.910432in}}%
\pgfpathlineto{\pgfqpoint{1.803241in}{0.903210in}}%
\pgfpathlineto{\pgfqpoint{1.801008in}{0.899912in}}%
\pgfpathlineto{\pgfqpoint{1.798311in}{0.895988in}}%
\pgfpathlineto{\pgfqpoint{1.793229in}{0.888766in}}%
\pgfpathlineto{\pgfqpoint{1.792759in}{0.888111in}}%
\pgfpathlineto{\pgfqpoint{1.787981in}{0.881543in}}%
\pgfpathlineto{\pgfqpoint{1.784510in}{0.876888in}}%
\pgfpathlineto{\pgfqpoint{1.782567in}{0.874321in}}%
\pgfpathlineto{\pgfqpoint{1.776977in}{0.867099in}}%
\pgfpathlineto{\pgfqpoint{1.776261in}{0.866192in}}%
\pgfpathlineto{\pgfqpoint{1.771199in}{0.859877in}}%
\pgfpathlineto{\pgfqpoint{1.768012in}{0.855991in}}%
\pgfpathlineto{\pgfqpoint{1.765230in}{0.852654in}}%
\pgfpathlineto{\pgfqpoint{1.759762in}{0.846244in}}%
\pgfpathlineto{\pgfqpoint{1.759059in}{0.845432in}}%
\pgfpathlineto{\pgfqpoint{1.752674in}{0.838210in}}%
\pgfpathlineto{\pgfqpoint{1.751513in}{0.836924in}}%
\pgfpathlineto{\pgfqpoint{1.746067in}{0.830988in}}%
\pgfpathlineto{\pgfqpoint{1.743264in}{0.828000in}}%
\pgfpathlineto{\pgfqpoint{1.739226in}{0.823766in}}%
\pgfpathlineto{\pgfqpoint{1.735015in}{0.819443in}}%
\pgfpathlineto{\pgfqpoint{1.732142in}{0.816543in}}%
\pgfpathlineto{\pgfqpoint{1.726765in}{0.811231in}}%
\pgfpathlineto{\pgfqpoint{1.724800in}{0.809321in}}%
\pgfpathlineto{\pgfqpoint{1.718516in}{0.803341in}}%
\pgfpathlineto{\pgfqpoint{1.717189in}{0.802099in}}%
\pgfpathlineto{\pgfqpoint{1.710267in}{0.795754in}}%
\pgfpathlineto{\pgfqpoint{1.709293in}{0.794877in}}%
\pgfpathlineto{\pgfqpoint{1.702018in}{0.788451in}}%
\pgfpathlineto{\pgfqpoint{1.701099in}{0.787654in}}%
\pgfpathlineto{\pgfqpoint{1.693769in}{0.781417in}}%
\pgfpathlineto{\pgfqpoint{1.692591in}{0.780432in}}%
\pgfpathlineto{\pgfqpoint{1.685519in}{0.774634in}}%
\pgfpathlineto{\pgfqpoint{1.683751in}{0.773210in}}%
\pgfpathlineto{\pgfqpoint{1.677270in}{0.768089in}}%
\pgfpathlineto{\pgfqpoint{1.674562in}{0.765988in}}%
\pgfpathlineto{\pgfqpoint{1.669021in}{0.761768in}}%
\pgfpathlineto{\pgfqpoint{1.665005in}{0.758766in}}%
\pgfpathlineto{\pgfqpoint{1.660772in}{0.755659in}}%
\pgfpathlineto{\pgfqpoint{1.655058in}{0.751543in}}%
\pgfpathlineto{\pgfqpoint{1.652522in}{0.749750in}}%
\pgfpathlineto{\pgfqpoint{1.644700in}{0.744321in}}%
\pgfpathlineto{\pgfqpoint{1.644273in}{0.744030in}}%
\pgfpathlineto{\pgfqpoint{1.636024in}{0.738493in}}%
\pgfpathlineto{\pgfqpoint{1.633908in}{0.737099in}}%
\pgfpathlineto{\pgfqpoint{1.627775in}{0.733125in}}%
\pgfpathlineto{\pgfqpoint{1.622663in}{0.729877in}}%
\pgfpathlineto{\pgfqpoint{1.619526in}{0.727915in}}%
\pgfpathlineto{\pgfqpoint{1.611276in}{0.722855in}}%
\pgfpathlineto{\pgfqpoint{1.610944in}{0.722654in}}%
\pgfpathlineto{\pgfqpoint{1.603027in}{0.717942in}}%
\pgfpathlineto{\pgfqpoint{1.598725in}{0.715432in}}%
\pgfpathlineto{\pgfqpoint{1.594778in}{0.713163in}}%
\pgfpathlineto{\pgfqpoint{1.586529in}{0.708513in}}%
\pgfpathlineto{\pgfqpoint{1.585983in}{0.708210in}}%
\pgfpathlineto{\pgfqpoint{1.578279in}{0.703989in}}%
\pgfpathlineto{\pgfqpoint{1.572687in}{0.700988in}}%
\pgfpathlineto{\pgfqpoint{1.570030in}{0.699581in}}%
\pgfpathlineto{\pgfqpoint{1.561781in}{0.695287in}}%
\pgfpathlineto{\pgfqpoint{1.558806in}{0.693766in}}%
\pgfpathlineto{\pgfqpoint{1.553532in}{0.691101in}}%
\pgfpathlineto{\pgfqpoint{1.545283in}{0.687018in}}%
\pgfpathlineto{\pgfqpoint{1.544311in}{0.686543in}}%
\pgfpathlineto{\pgfqpoint{1.537033in}{0.683033in}}%
\pgfpathlineto{\pgfqpoint{1.529165in}{0.679321in}}%
\pgfpathlineto{\pgfqpoint{1.528784in}{0.679144in}}%
\pgfpathlineto{\pgfqpoint{1.520535in}{0.675345in}}%
\pgfpathlineto{\pgfqpoint{1.513330in}{0.672099in}}%
\pgfpathlineto{\pgfqpoint{1.512286in}{0.671633in}}%
\pgfpathlineto{\pgfqpoint{1.504036in}{0.668005in}}%
\pgfpathlineto{\pgfqpoint{1.496768in}{0.664877in}}%
\pgfpathlineto{\pgfqpoint{1.495787in}{0.664459in}}%
\pgfpathlineto{\pgfqpoint{1.487538in}{0.660989in}}%
\pgfpathlineto{\pgfqpoint{1.479432in}{0.657654in}}%
\pgfpathlineto{\pgfqpoint{1.479289in}{0.657596in}}%
\pgfpathlineto{\pgfqpoint{1.471040in}{0.654273in}}%
\pgfpathlineto{\pgfqpoint{1.462790in}{0.651022in}}%
\pgfpathlineto{\pgfqpoint{1.461273in}{0.650432in}}%
\pgfpathlineto{\pgfqpoint{1.454541in}{0.647838in}}%
\pgfpathlineto{\pgfqpoint{1.446292in}{0.644720in}}%
\pgfpathlineto{\pgfqpoint{1.442232in}{0.643210in}}%
\pgfpathlineto{\pgfqpoint{1.438043in}{0.641666in}}%
\pgfpathlineto{\pgfqpoint{1.429793in}{0.638673in}}%
\pgfpathlineto{\pgfqpoint{1.422240in}{0.635988in}}%
\pgfpathlineto{\pgfqpoint{1.421544in}{0.635743in}}%
\pgfpathlineto{\pgfqpoint{1.413295in}{0.632868in}}%
\pgfpathlineto{\pgfqpoint{1.405046in}{0.630053in}}%
\pgfpathlineto{\pgfqpoint{1.401216in}{0.628766in}}%
\pgfpathlineto{\pgfqpoint{1.396797in}{0.627292in}}%
\pgfpathlineto{\pgfqpoint{1.388547in}{0.624585in}}%
\pgfpathlineto{\pgfqpoint{1.380298in}{0.621932in}}%
\pgfpathlineto{\pgfqpoint{1.379074in}{0.621543in}}%
\pgfpathlineto{\pgfqpoint{1.372049in}{0.619328in}}%
\pgfpathlineto{\pgfqpoint{1.363800in}{0.616775in}}%
\pgfpathlineto{\pgfqpoint{1.355714in}{0.614321in}}%
\pgfpathlineto{\pgfqpoint{1.355550in}{0.614272in}}%
\pgfpathlineto{\pgfqpoint{1.347301in}{0.611812in}}%
\pgfpathlineto{\pgfqpoint{1.339052in}{0.609399in}}%
\pgfpathlineto{\pgfqpoint{1.331039in}{0.607099in}}%
\pgfpathlineto{\pgfqpoint{1.330803in}{0.607032in}}%
\pgfpathlineto{\pgfqpoint{1.322553in}{0.604704in}}%
\pgfpathlineto{\pgfqpoint{1.314304in}{0.602419in}}%
\pgfpathlineto{\pgfqpoint{1.306055in}{0.600177in}}%
\pgfpathlineto{\pgfqpoint{1.304939in}{0.599877in}}%
\pgfpathlineto{\pgfqpoint{1.297806in}{0.597971in}}%
\pgfpathlineto{\pgfqpoint{1.289557in}{0.595804in}}%
\pgfpathlineto{\pgfqpoint{1.281307in}{0.593675in}}%
\pgfpathlineto{\pgfqpoint{1.277297in}{0.592654in}}%
\pgfpathlineto{\pgfqpoint{1.273058in}{0.591582in}}%
\pgfpathlineto{\pgfqpoint{1.264809in}{0.589524in}}%
\pgfpathlineto{\pgfqpoint{1.256560in}{0.587501in}}%
\pgfpathlineto{\pgfqpoint{1.248310in}{0.585513in}}%
\pgfpathlineto{\pgfqpoint{1.247974in}{0.585432in}}%
\pgfpathlineto{\pgfqpoint{1.240061in}{0.583555in}}%
\pgfpathlineto{\pgfqpoint{1.231812in}{0.581630in}}%
\pgfpathlineto{\pgfqpoint{1.223563in}{0.579738in}}%
\pgfpathlineto{\pgfqpoint{1.216797in}{0.578210in}}%
\pgfpathlineto{\pgfqpoint{1.215314in}{0.577877in}}%
\pgfpathlineto{\pgfqpoint{1.207064in}{0.576044in}}%
\pgfpathlineto{\pgfqpoint{1.198815in}{0.574242in}}%
\pgfpathlineto{\pgfqpoint{1.190566in}{0.572469in}}%
\pgfpathlineto{\pgfqpoint{1.183569in}{0.570988in}}%
\pgfpathlineto{\pgfqpoint{1.182317in}{0.570724in}}%
\pgfpathlineto{\pgfqpoint{1.174067in}{0.569005in}}%
\pgfpathlineto{\pgfqpoint{1.165818in}{0.567314in}}%
\pgfpathlineto{\pgfqpoint{1.157569in}{0.565650in}}%
\pgfpathlineto{\pgfqpoint{1.149320in}{0.564012in}}%
\pgfpathlineto{\pgfqpoint{1.148068in}{0.563766in}}%
\pgfpathlineto{\pgfqpoint{1.141071in}{0.562397in}}%
\pgfpathlineto{\pgfqpoint{1.132821in}{0.560806in}}%
\pgfpathlineto{\pgfqpoint{1.124572in}{0.559240in}}%
\pgfpathlineto{\pgfqpoint{1.116323in}{0.557698in}}%
\pgfpathlineto{\pgfqpoint{1.110064in}{0.556543in}}%
\pgfpathlineto{\pgfqpoint{1.108074in}{0.556178in}}%
\pgfpathlineto{\pgfqpoint{1.099824in}{0.554679in}}%
\pgfpathlineto{\pgfqpoint{1.091575in}{0.553203in}}%
\pgfpathlineto{\pgfqpoint{1.083326in}{0.551747in}}%
\pgfpathlineto{\pgfqpoint{1.075077in}{0.550313in}}%
\pgfpathlineto{\pgfqpoint{1.069297in}{0.549321in}}%
\pgfpathlineto{\pgfqpoint{1.066828in}{0.548900in}}%
\pgfpathlineto{\pgfqpoint{1.058578in}{0.547505in}}%
\pgfpathlineto{\pgfqpoint{1.050329in}{0.546130in}}%
\pgfpathlineto{\pgfqpoint{1.042080in}{0.544776in}}%
\pgfpathlineto{\pgfqpoint{1.033831in}{0.543440in}}%
\pgfpathlineto{\pgfqpoint{1.025581in}{0.542124in}}%
\pgfpathlineto{\pgfqpoint{1.025423in}{0.542099in}}%
\pgfpathlineto{\pgfqpoint{1.017332in}{0.540825in}}%
\pgfpathlineto{\pgfqpoint{1.009083in}{0.539544in}}%
\pgfpathlineto{\pgfqpoint{1.000834in}{0.538282in}}%
\pgfpathlineto{\pgfqpoint{0.992585in}{0.537038in}}%
\pgfpathlineto{\pgfqpoint{0.984335in}{0.535811in}}%
\pgfpathlineto{\pgfqpoint{0.977978in}{0.534877in}}%
\pgfpathlineto{\pgfqpoint{0.976086in}{0.534600in}}%
\pgfpathlineto{\pgfqpoint{0.967837in}{0.533406in}}%
\pgfpathlineto{\pgfqpoint{0.959588in}{0.532228in}}%
\pgfpathlineto{\pgfqpoint{0.951338in}{0.531067in}}%
\pgfpathlineto{\pgfqpoint{0.943089in}{0.529921in}}%
\pgfpathlineto{\pgfqpoint{0.934840in}{0.528790in}}%
\pgfpathlineto{\pgfqpoint{0.926591in}{0.527673in}}%
\pgfpathlineto{\pgfqpoint{0.926450in}{0.527654in}}%
\pgfpathlineto{\pgfqpoint{0.918342in}{0.526571in}}%
\pgfpathlineto{\pgfqpoint{0.910092in}{0.525484in}}%
\pgfpathlineto{\pgfqpoint{0.901843in}{0.524410in}}%
\pgfpathlineto{\pgfqpoint{0.893594in}{0.523350in}}%
\pgfpathlineto{\pgfqpoint{0.885345in}{0.522303in}}%
\pgfpathlineto{\pgfqpoint{0.877095in}{0.521270in}}%
\pgfpathlineto{\pgfqpoint{0.870336in}{0.520432in}}%
\pgfpathlineto{\pgfqpoint{0.868846in}{0.520249in}}%
\pgfpathlineto{\pgfqpoint{0.860597in}{0.519241in}}%
\pgfpathlineto{\pgfqpoint{0.852348in}{0.518245in}}%
\pgfpathlineto{\pgfqpoint{0.844099in}{0.517262in}}%
\pgfpathlineto{\pgfqpoint{0.835849in}{0.516291in}}%
\pgfpathlineto{\pgfqpoint{0.827600in}{0.515332in}}%
\pgfpathlineto{\pgfqpoint{0.819351in}{0.514384in}}%
\pgfpathlineto{\pgfqpoint{0.811102in}{0.513448in}}%
\pgfpathlineto{\pgfqpoint{0.808986in}{0.513210in}}%
\pgfpathlineto{\pgfqpoint{0.802852in}{0.512524in}}%
\pgfpathlineto{\pgfqpoint{0.794603in}{0.511610in}}%
\pgfpathlineto{\pgfqpoint{0.786354in}{0.510708in}}%
\pgfpathlineto{\pgfqpoint{0.778105in}{0.509816in}}%
\pgfpathlineto{\pgfqpoint{0.769856in}{0.508935in}}%
\pgfpathlineto{\pgfqpoint{0.761606in}{0.508064in}}%
\pgfpathlineto{\pgfqpoint{0.753357in}{0.507203in}}%
\pgfpathlineto{\pgfqpoint{0.745108in}{0.506351in}}%
\pgfpathlineto{\pgfqpoint{0.741557in}{0.505988in}}%
\pgfpathlineto{\pgfqpoint{0.736859in}{0.505510in}}%
\pgfpathlineto{\pgfqpoint{0.728609in}{0.504679in}}%
\pgfpathlineto{\pgfqpoint{0.720360in}{0.503856in}}%
\pgfpathlineto{\pgfqpoint{0.712111in}{0.503043in}}%
\pgfpathlineto{\pgfqpoint{0.703862in}{0.502239in}}%
\pgfpathlineto{\pgfqpoint{0.695612in}{0.501443in}}%
\pgfpathlineto{\pgfqpoint{0.687363in}{0.500656in}}%
\pgfpathlineto{\pgfqpoint{0.679114in}{0.499878in}}%
\pgfpathlineto{\pgfqpoint{0.670865in}{0.499107in}}%
\pgfpathlineto{\pgfqpoint{0.667177in}{0.498766in}}%
\pgfpathlineto{\pgfqpoint{0.662616in}{0.498346in}}%
\pgfpathlineto{\pgfqpoint{0.654366in}{0.497593in}}%
\pgfpathlineto{\pgfqpoint{0.646117in}{0.496849in}}%
\pgfpathlineto{\pgfqpoint{0.637868in}{0.496113in}}%
\pgfpathlineto{\pgfqpoint{0.629619in}{0.495384in}}%
\pgfpathlineto{\pgfqpoint{0.621369in}{0.494664in}}%
\pgfpathlineto{\pgfqpoint{0.613120in}{0.493951in}}%
\pgfpathlineto{\pgfqpoint{0.604871in}{0.493246in}}%
\pgfpathlineto{\pgfqpoint{0.596622in}{0.492548in}}%
\pgfpathlineto{\pgfqpoint{0.588373in}{0.491858in}}%
\pgfpathlineto{\pgfqpoint{0.584589in}{0.491543in}}%
\pgfpathlineto{\pgfqpoint{0.580123in}{0.491175in}}%
\pgfpathlineto{\pgfqpoint{0.571874in}{0.490501in}}%
\pgfpathlineto{\pgfqpoint{0.563625in}{0.489834in}}%
\pgfpathlineto{\pgfqpoint{0.555376in}{0.489173in}}%
\pgfpathlineto{\pgfqpoint{0.547126in}{0.488519in}}%
\pgfpathlineto{\pgfqpoint{0.538877in}{0.487872in}}%
\pgfpathlineto{\pgfqpoint{0.530628in}{0.487232in}}%
\pgfpathlineto{\pgfqpoint{0.522379in}{0.486597in}}%
\pgfpathlineto{\pgfqpoint{0.514130in}{0.485969in}}%
\pgfpathlineto{\pgfqpoint{0.505880in}{0.485347in}}%
\pgfpathlineto{\pgfqpoint{0.497631in}{0.484731in}}%
\pgfpathlineto{\pgfqpoint{0.492092in}{0.484321in}}%
\pgfpathlineto{\pgfqpoint{0.489382in}{0.484122in}}%
\pgfpathlineto{\pgfqpoint{0.481133in}{0.483520in}}%
\pgfpathlineto{\pgfqpoint{0.472883in}{0.482923in}}%
\pgfpathlineto{\pgfqpoint{0.464634in}{0.482333in}}%
\pgfpathlineto{\pgfqpoint{0.456385in}{0.481747in}}%
\pgfpathlineto{\pgfqpoint{0.448136in}{0.481168in}}%
\pgfpathlineto{\pgfqpoint{0.439887in}{0.480593in}}%
\pgfpathlineto{\pgfqpoint{0.431637in}{0.480024in}}%
\pgfpathlineto{\pgfqpoint{0.423388in}{0.479460in}}%
\pgfpathlineto{\pgfqpoint{0.415139in}{0.478901in}}%
\pgfpathlineto{\pgfqpoint{0.415139in}{0.477099in}}%
\pgfpathlineto{\pgfqpoint{0.415139in}{0.469877in}}%
\pgfpathlineto{\pgfqpoint{0.415139in}{0.462654in}}%
\pgfpathlineto{\pgfqpoint{0.423388in}{0.462654in}}%
\pgfpathclose%
\pgfusepath{fill}%
\end{pgfscope}%
\begin{pgfscope}%
\pgfpathrectangle{\pgfqpoint{0.415139in}{0.462654in}}{\pgfqpoint{1.641596in}{1.437222in}}%
\pgfusepath{clip}%
\pgfsetbuttcap%
\pgfsetroundjoin%
\definecolor{currentfill}{rgb}{0.229066,0.328028,0.641522}%
\pgfsetfillcolor{currentfill}%
\pgfsetlinewidth{0.000000pt}%
\definecolor{currentstroke}{rgb}{0.000000,0.000000,0.000000}%
\pgfsetstrokecolor{currentstroke}%
\pgfsetdash{}{0pt}%
\pgfusepath{fill}%
\end{pgfscope}%
\begin{pgfscope}%
\pgfsetbuttcap%
\pgfsetroundjoin%
\definecolor{currentfill}{rgb}{0.000000,0.000000,0.000000}%
\pgfsetfillcolor{currentfill}%
\pgfsetlinewidth{0.803000pt}%
\definecolor{currentstroke}{rgb}{0.000000,0.000000,0.000000}%
\pgfsetstrokecolor{currentstroke}%
\pgfsetdash{}{0pt}%
\pgfsys@defobject{currentmarker}{\pgfqpoint{0.000000in}{-0.048611in}}{\pgfqpoint{0.000000in}{0.000000in}}{%
\pgfpathmoveto{\pgfqpoint{0.000000in}{0.000000in}}%
\pgfpathlineto{\pgfqpoint{0.000000in}{-0.048611in}}%
\pgfusepath{stroke,fill}%
}%
\begin{pgfscope}%
\pgfsys@transformshift{0.415139in}{0.462654in}%
\pgfsys@useobject{currentmarker}{}%
\end{pgfscope}%
\end{pgfscope}%
\begin{pgfscope}%
\definecolor{textcolor}{rgb}{0.000000,0.000000,0.000000}%
\pgfsetstrokecolor{textcolor}%
\pgfsetfillcolor{textcolor}%
\pgftext[x=0.415139in,y=0.365432in,,top]{\color{textcolor}{\rmfamily\fontsize{8.000000}{9.600000}\selectfont\catcode`\^=\active\def^{\ifmmode\sp\else\^{}\fi}\catcode`\%=\active\def%{\%}0.00}}%
\end{pgfscope}%
\begin{pgfscope}%
\pgfsetbuttcap%
\pgfsetroundjoin%
\definecolor{currentfill}{rgb}{0.000000,0.000000,0.000000}%
\pgfsetfillcolor{currentfill}%
\pgfsetlinewidth{0.803000pt}%
\definecolor{currentstroke}{rgb}{0.000000,0.000000,0.000000}%
\pgfsetstrokecolor{currentstroke}%
\pgfsetdash{}{0pt}%
\pgfsys@defobject{currentmarker}{\pgfqpoint{0.000000in}{-0.048611in}}{\pgfqpoint{0.000000in}{0.000000in}}{%
\pgfpathmoveto{\pgfqpoint{0.000000in}{0.000000in}}%
\pgfpathlineto{\pgfqpoint{0.000000in}{-0.048611in}}%
\pgfusepath{stroke,fill}%
}%
\begin{pgfscope}%
\pgfsys@transformshift{0.825538in}{0.462654in}%
\pgfsys@useobject{currentmarker}{}%
\end{pgfscope}%
\end{pgfscope}%
\begin{pgfscope}%
\definecolor{textcolor}{rgb}{0.000000,0.000000,0.000000}%
\pgfsetstrokecolor{textcolor}%
\pgfsetfillcolor{textcolor}%
\pgftext[x=0.825538in,y=0.365432in,,top]{\color{textcolor}{\rmfamily\fontsize{8.000000}{9.600000}\selectfont\catcode`\^=\active\def^{\ifmmode\sp\else\^{}\fi}\catcode`\%=\active\def%{\%}0.25}}%
\end{pgfscope}%
\begin{pgfscope}%
\pgfsetbuttcap%
\pgfsetroundjoin%
\definecolor{currentfill}{rgb}{0.000000,0.000000,0.000000}%
\pgfsetfillcolor{currentfill}%
\pgfsetlinewidth{0.803000pt}%
\definecolor{currentstroke}{rgb}{0.000000,0.000000,0.000000}%
\pgfsetstrokecolor{currentstroke}%
\pgfsetdash{}{0pt}%
\pgfsys@defobject{currentmarker}{\pgfqpoint{0.000000in}{-0.048611in}}{\pgfqpoint{0.000000in}{0.000000in}}{%
\pgfpathmoveto{\pgfqpoint{0.000000in}{0.000000in}}%
\pgfpathlineto{\pgfqpoint{0.000000in}{-0.048611in}}%
\pgfusepath{stroke,fill}%
}%
\begin{pgfscope}%
\pgfsys@transformshift{1.235937in}{0.462654in}%
\pgfsys@useobject{currentmarker}{}%
\end{pgfscope}%
\end{pgfscope}%
\begin{pgfscope}%
\definecolor{textcolor}{rgb}{0.000000,0.000000,0.000000}%
\pgfsetstrokecolor{textcolor}%
\pgfsetfillcolor{textcolor}%
\pgftext[x=1.235937in,y=0.365432in,,top]{\color{textcolor}{\rmfamily\fontsize{8.000000}{9.600000}\selectfont\catcode`\^=\active\def^{\ifmmode\sp\else\^{}\fi}\catcode`\%=\active\def%{\%}0.50}}%
\end{pgfscope}%
\begin{pgfscope}%
\pgfsetbuttcap%
\pgfsetroundjoin%
\definecolor{currentfill}{rgb}{0.000000,0.000000,0.000000}%
\pgfsetfillcolor{currentfill}%
\pgfsetlinewidth{0.803000pt}%
\definecolor{currentstroke}{rgb}{0.000000,0.000000,0.000000}%
\pgfsetstrokecolor{currentstroke}%
\pgfsetdash{}{0pt}%
\pgfsys@defobject{currentmarker}{\pgfqpoint{0.000000in}{-0.048611in}}{\pgfqpoint{0.000000in}{0.000000in}}{%
\pgfpathmoveto{\pgfqpoint{0.000000in}{0.000000in}}%
\pgfpathlineto{\pgfqpoint{0.000000in}{-0.048611in}}%
\pgfusepath{stroke,fill}%
}%
\begin{pgfscope}%
\pgfsys@transformshift{1.646336in}{0.462654in}%
\pgfsys@useobject{currentmarker}{}%
\end{pgfscope}%
\end{pgfscope}%
\begin{pgfscope}%
\definecolor{textcolor}{rgb}{0.000000,0.000000,0.000000}%
\pgfsetstrokecolor{textcolor}%
\pgfsetfillcolor{textcolor}%
\pgftext[x=1.646336in,y=0.365432in,,top]{\color{textcolor}{\rmfamily\fontsize{8.000000}{9.600000}\selectfont\catcode`\^=\active\def^{\ifmmode\sp\else\^{}\fi}\catcode`\%=\active\def%{\%}0.75}}%
\end{pgfscope}%
\begin{pgfscope}%
\pgfsetbuttcap%
\pgfsetroundjoin%
\definecolor{currentfill}{rgb}{0.000000,0.000000,0.000000}%
\pgfsetfillcolor{currentfill}%
\pgfsetlinewidth{0.803000pt}%
\definecolor{currentstroke}{rgb}{0.000000,0.000000,0.000000}%
\pgfsetstrokecolor{currentstroke}%
\pgfsetdash{}{0pt}%
\pgfsys@defobject{currentmarker}{\pgfqpoint{0.000000in}{-0.048611in}}{\pgfqpoint{0.000000in}{0.000000in}}{%
\pgfpathmoveto{\pgfqpoint{0.000000in}{0.000000in}}%
\pgfpathlineto{\pgfqpoint{0.000000in}{-0.048611in}}%
\pgfusepath{stroke,fill}%
}%
\begin{pgfscope}%
\pgfsys@transformshift{2.056734in}{0.462654in}%
\pgfsys@useobject{currentmarker}{}%
\end{pgfscope}%
\end{pgfscope}%
\begin{pgfscope}%
\definecolor{textcolor}{rgb}{0.000000,0.000000,0.000000}%
\pgfsetstrokecolor{textcolor}%
\pgfsetfillcolor{textcolor}%
\pgftext[x=2.056734in,y=0.365432in,,top]{\color{textcolor}{\rmfamily\fontsize{8.000000}{9.600000}\selectfont\catcode`\^=\active\def^{\ifmmode\sp\else\^{}\fi}\catcode`\%=\active\def%{\%}1.00}}%
\end{pgfscope}%
\begin{pgfscope}%
\definecolor{textcolor}{rgb}{0.000000,0.000000,0.000000}%
\pgfsetstrokecolor{textcolor}%
\pgfsetfillcolor{textcolor}%
\pgftext[x=1.235937in,y=0.211111in,,top]{\color{textcolor}{\rmfamily\fontsize{9.000000}{10.800000}\selectfont\catcode`\^=\active\def^{\ifmmode\sp\else\^{}\fi}\catcode`\%=\active\def%{\%}$\chi$}}%
\end{pgfscope}%
\begin{pgfscope}%
\pgfsetbuttcap%
\pgfsetroundjoin%
\definecolor{currentfill}{rgb}{0.000000,0.000000,0.000000}%
\pgfsetfillcolor{currentfill}%
\pgfsetlinewidth{0.803000pt}%
\definecolor{currentstroke}{rgb}{0.000000,0.000000,0.000000}%
\pgfsetstrokecolor{currentstroke}%
\pgfsetdash{}{0pt}%
\pgfsys@defobject{currentmarker}{\pgfqpoint{-0.048611in}{0.000000in}}{\pgfqpoint{-0.000000in}{0.000000in}}{%
\pgfpathmoveto{\pgfqpoint{-0.000000in}{0.000000in}}%
\pgfpathlineto{\pgfqpoint{-0.048611in}{0.000000in}}%
\pgfusepath{stroke,fill}%
}%
\begin{pgfscope}%
\pgfsys@transformshift{0.415139in}{0.750099in}%
\pgfsys@useobject{currentmarker}{}%
\end{pgfscope}%
\end{pgfscope}%
\begin{pgfscope}%
\definecolor{textcolor}{rgb}{0.000000,0.000000,0.000000}%
\pgfsetstrokecolor{textcolor}%
\pgfsetfillcolor{textcolor}%
\pgftext[x=0.258888in, y=0.711519in, left, base]{\color{textcolor}{\rmfamily\fontsize{8.000000}{9.600000}\selectfont\catcode`\^=\active\def^{\ifmmode\sp\else\^{}\fi}\catcode`\%=\active\def%{\%}2}}%
\end{pgfscope}%
\begin{pgfscope}%
\pgfsetbuttcap%
\pgfsetroundjoin%
\definecolor{currentfill}{rgb}{0.000000,0.000000,0.000000}%
\pgfsetfillcolor{currentfill}%
\pgfsetlinewidth{0.803000pt}%
\definecolor{currentstroke}{rgb}{0.000000,0.000000,0.000000}%
\pgfsetstrokecolor{currentstroke}%
\pgfsetdash{}{0pt}%
\pgfsys@defobject{currentmarker}{\pgfqpoint{-0.048611in}{0.000000in}}{\pgfqpoint{-0.000000in}{0.000000in}}{%
\pgfpathmoveto{\pgfqpoint{-0.000000in}{0.000000in}}%
\pgfpathlineto{\pgfqpoint{-0.048611in}{0.000000in}}%
\pgfusepath{stroke,fill}%
}%
\begin{pgfscope}%
\pgfsys@transformshift{0.415139in}{1.133358in}%
\pgfsys@useobject{currentmarker}{}%
\end{pgfscope}%
\end{pgfscope}%
\begin{pgfscope}%
\definecolor{textcolor}{rgb}{0.000000,0.000000,0.000000}%
\pgfsetstrokecolor{textcolor}%
\pgfsetfillcolor{textcolor}%
\pgftext[x=0.258888in, y=1.094778in, left, base]{\color{textcolor}{\rmfamily\fontsize{8.000000}{9.600000}\selectfont\catcode`\^=\active\def^{\ifmmode\sp\else\^{}\fi}\catcode`\%=\active\def%{\%}4}}%
\end{pgfscope}%
\begin{pgfscope}%
\pgfsetbuttcap%
\pgfsetroundjoin%
\definecolor{currentfill}{rgb}{0.000000,0.000000,0.000000}%
\pgfsetfillcolor{currentfill}%
\pgfsetlinewidth{0.803000pt}%
\definecolor{currentstroke}{rgb}{0.000000,0.000000,0.000000}%
\pgfsetstrokecolor{currentstroke}%
\pgfsetdash{}{0pt}%
\pgfsys@defobject{currentmarker}{\pgfqpoint{-0.048611in}{0.000000in}}{\pgfqpoint{-0.000000in}{0.000000in}}{%
\pgfpathmoveto{\pgfqpoint{-0.000000in}{0.000000in}}%
\pgfpathlineto{\pgfqpoint{-0.048611in}{0.000000in}}%
\pgfusepath{stroke,fill}%
}%
\begin{pgfscope}%
\pgfsys@transformshift{0.415139in}{1.516617in}%
\pgfsys@useobject{currentmarker}{}%
\end{pgfscope}%
\end{pgfscope}%
\begin{pgfscope}%
\definecolor{textcolor}{rgb}{0.000000,0.000000,0.000000}%
\pgfsetstrokecolor{textcolor}%
\pgfsetfillcolor{textcolor}%
\pgftext[x=0.258888in, y=1.478037in, left, base]{\color{textcolor}{\rmfamily\fontsize{8.000000}{9.600000}\selectfont\catcode`\^=\active\def^{\ifmmode\sp\else\^{}\fi}\catcode`\%=\active\def%{\%}6}}%
\end{pgfscope}%
\begin{pgfscope}%
\pgfsetbuttcap%
\pgfsetroundjoin%
\definecolor{currentfill}{rgb}{0.000000,0.000000,0.000000}%
\pgfsetfillcolor{currentfill}%
\pgfsetlinewidth{0.803000pt}%
\definecolor{currentstroke}{rgb}{0.000000,0.000000,0.000000}%
\pgfsetstrokecolor{currentstroke}%
\pgfsetdash{}{0pt}%
\pgfsys@defobject{currentmarker}{\pgfqpoint{-0.048611in}{0.000000in}}{\pgfqpoint{-0.000000in}{0.000000in}}{%
\pgfpathmoveto{\pgfqpoint{-0.000000in}{0.000000in}}%
\pgfpathlineto{\pgfqpoint{-0.048611in}{0.000000in}}%
\pgfusepath{stroke,fill}%
}%
\begin{pgfscope}%
\pgfsys@transformshift{0.415139in}{1.899877in}%
\pgfsys@useobject{currentmarker}{}%
\end{pgfscope}%
\end{pgfscope}%
\begin{pgfscope}%
\definecolor{textcolor}{rgb}{0.000000,0.000000,0.000000}%
\pgfsetstrokecolor{textcolor}%
\pgfsetfillcolor{textcolor}%
\pgftext[x=0.258888in, y=1.861296in, left, base]{\color{textcolor}{\rmfamily\fontsize{8.000000}{9.600000}\selectfont\catcode`\^=\active\def^{\ifmmode\sp\else\^{}\fi}\catcode`\%=\active\def%{\%}8}}%
\end{pgfscope}%
\begin{pgfscope}%
\definecolor{textcolor}{rgb}{0.000000,0.000000,0.000000}%
\pgfsetstrokecolor{textcolor}%
\pgfsetfillcolor{textcolor}%
\pgftext[x=0.244999in,y=1.181266in,,bottom,rotate=90.000000]{\color{textcolor}{\rmfamily\fontsize{9.000000}{10.800000}\selectfont\catcode`\^=\active\def^{\ifmmode\sp\else\^{}\fi}\catcode`\%=\active\def%{\%}$T\ (10^3/\omega_b)$}}%
\end{pgfscope}%
\begin{pgfscope}%
\pgfpathrectangle{\pgfqpoint{0.415139in}{0.462654in}}{\pgfqpoint{1.641596in}{1.437222in}}%
\pgfusepath{clip}%
\pgfsetbuttcap%
\pgfsetroundjoin%
\pgfsetlinewidth{0.803000pt}%
\definecolor{currentstroke}{rgb}{0.000000,0.000000,0.000000}%
\pgfsetstrokecolor{currentstroke}%
\pgfsetdash{}{0pt}%
\pgfpathmoveto{\pgfqpoint{1.352991in}{1.899877in}}%
\pgfpathlineto{\pgfqpoint{1.318854in}{1.834877in}}%
\pgfpathlineto{\pgfqpoint{1.286136in}{1.777099in}}%
\pgfpathlineto{\pgfqpoint{1.250859in}{1.719321in}}%
\pgfpathlineto{\pgfqpoint{1.215314in}{1.665435in}}%
\pgfpathlineto{\pgfqpoint{1.181754in}{1.618210in}}%
\pgfpathlineto{\pgfqpoint{1.142994in}{1.567654in}}%
\pgfpathlineto{\pgfqpoint{1.107168in}{1.524321in}}%
\pgfpathlineto{\pgfqpoint{1.066828in}{1.479011in}}%
\pgfpathlineto{\pgfqpoint{1.025581in}{1.436119in}}%
\pgfpathlineto{\pgfqpoint{0.982100in}{1.394321in}}%
\pgfpathlineto{\pgfqpoint{0.968214in}{1.381656in}}%
\pgfpathmoveto{\pgfqpoint{0.798823in}{1.248844in}}%
\pgfpathlineto{\pgfqpoint{0.745108in}{1.213621in}}%
\pgfpathlineto{\pgfqpoint{0.685548in}{1.177654in}}%
\pgfpathlineto{\pgfqpoint{0.629619in}{1.146492in}}%
\pgfpathlineto{\pgfqpoint{0.571874in}{1.116733in}}%
\pgfpathlineto{\pgfqpoint{0.514130in}{1.089172in}}%
\pgfpathlineto{\pgfqpoint{0.448136in}{1.060036in}}%
\pgfpathlineto{\pgfqpoint{0.415139in}{1.046309in}}%
\pgfpathlineto{\pgfqpoint{0.415139in}{1.046309in}}%
\pgfusepath{stroke}%
\end{pgfscope}%
\begin{pgfscope}%
\pgfpathrectangle{\pgfqpoint{0.415139in}{0.462654in}}{\pgfqpoint{1.641596in}{1.437222in}}%
\pgfusepath{clip}%
\pgfsetbuttcap%
\pgfsetroundjoin%
\pgfsetlinewidth{0.803000pt}%
\definecolor{currentstroke}{rgb}{0.000000,0.000000,0.000000}%
\pgfsetstrokecolor{currentstroke}%
\pgfsetdash{}{0pt}%
\pgfpathmoveto{\pgfqpoint{2.011599in}{1.899877in}}%
\pgfpathlineto{\pgfqpoint{1.995943in}{1.704877in}}%
\pgfpathlineto{\pgfqpoint{1.982700in}{1.567654in}}%
\pgfpathlineto{\pgfqpoint{1.970905in}{1.466543in}}%
\pgfpathlineto{\pgfqpoint{1.958780in}{1.379877in}}%
\pgfpathlineto{\pgfqpoint{1.946719in}{1.307654in}}%
\pgfpathlineto{\pgfqpoint{1.932996in}{1.238624in}}%
\pgfpathlineto{\pgfqpoint{1.918551in}{1.177654in}}%
\pgfpathlineto{\pgfqpoint{1.904476in}{1.127099in}}%
\pgfpathlineto{\pgfqpoint{1.888011in}{1.076543in}}%
\pgfpathlineto{\pgfqpoint{1.871546in}{1.033210in}}%
\pgfpathlineto{\pgfqpoint{1.855799in}{0.997099in}}%
\pgfpathlineto{\pgfqpoint{1.837825in}{0.960988in}}%
\pgfpathlineto{\pgfqpoint{1.817168in}{0.924877in}}%
\pgfpathlineto{\pgfqpoint{1.798311in}{0.895988in}}%
\pgfpathlineto{\pgfqpoint{1.776261in}{0.866192in}}%
\pgfpathlineto{\pgfqpoint{1.751513in}{0.836924in}}%
\pgfpathlineto{\pgfqpoint{1.748179in}{0.833290in}}%
\pgfpathmoveto{\pgfqpoint{1.577239in}{0.703430in}}%
\pgfpathlineto{\pgfqpoint{1.537033in}{0.683033in}}%
\pgfpathlineto{\pgfqpoint{1.495787in}{0.664459in}}%
\pgfpathlineto{\pgfqpoint{1.446292in}{0.644720in}}%
\pgfpathlineto{\pgfqpoint{1.396797in}{0.627292in}}%
\pgfpathlineto{\pgfqpoint{1.339052in}{0.609399in}}%
\pgfpathlineto{\pgfqpoint{1.273058in}{0.591582in}}%
\pgfpathlineto{\pgfqpoint{1.207064in}{0.576044in}}%
\pgfpathlineto{\pgfqpoint{1.132821in}{0.560806in}}%
\pgfpathlineto{\pgfqpoint{1.050329in}{0.546130in}}%
\pgfpathlineto{\pgfqpoint{0.959588in}{0.532228in}}%
\pgfpathlineto{\pgfqpoint{0.852348in}{0.518245in}}%
\pgfpathlineto{\pgfqpoint{0.736859in}{0.505510in}}%
\pgfpathlineto{\pgfqpoint{0.604871in}{0.493246in}}%
\pgfpathlineto{\pgfqpoint{0.456385in}{0.481747in}}%
\pgfpathlineto{\pgfqpoint{0.415139in}{0.478901in}}%
\pgfpathlineto{\pgfqpoint{0.415139in}{0.478901in}}%
\pgfusepath{stroke}%
\end{pgfscope}%
\begin{pgfscope}%
\pgfsetrectcap%
\pgfsetmiterjoin%
\pgfsetlinewidth{0.803000pt}%
\definecolor{currentstroke}{rgb}{0.000000,0.000000,0.000000}%
\pgfsetstrokecolor{currentstroke}%
\pgfsetdash{}{0pt}%
\pgfpathmoveto{\pgfqpoint{0.415139in}{0.462654in}}%
\pgfpathlineto{\pgfqpoint{0.415139in}{1.899877in}}%
\pgfusepath{stroke}%
\end{pgfscope}%
\begin{pgfscope}%
\pgfsetrectcap%
\pgfsetmiterjoin%
\pgfsetlinewidth{0.803000pt}%
\definecolor{currentstroke}{rgb}{0.000000,0.000000,0.000000}%
\pgfsetstrokecolor{currentstroke}%
\pgfsetdash{}{0pt}%
\pgfpathmoveto{\pgfqpoint{2.056734in}{0.462654in}}%
\pgfpathlineto{\pgfqpoint{2.056734in}{1.899877in}}%
\pgfusepath{stroke}%
\end{pgfscope}%
\begin{pgfscope}%
\pgfsetrectcap%
\pgfsetmiterjoin%
\pgfsetlinewidth{0.803000pt}%
\definecolor{currentstroke}{rgb}{0.000000,0.000000,0.000000}%
\pgfsetstrokecolor{currentstroke}%
\pgfsetdash{}{0pt}%
\pgfpathmoveto{\pgfqpoint{0.415139in}{0.462654in}}%
\pgfpathlineto{\pgfqpoint{2.056734in}{0.462654in}}%
\pgfusepath{stroke}%
\end{pgfscope}%
\begin{pgfscope}%
\pgfsetrectcap%
\pgfsetmiterjoin%
\pgfsetlinewidth{0.803000pt}%
\definecolor{currentstroke}{rgb}{0.000000,0.000000,0.000000}%
\pgfsetstrokecolor{currentstroke}%
\pgfsetdash{}{0pt}%
\pgfpathmoveto{\pgfqpoint{0.415139in}{1.899877in}}%
\pgfpathlineto{\pgfqpoint{2.056734in}{1.899877in}}%
\pgfusepath{stroke}%
\end{pgfscope}%
\begin{pgfscope}%
\definecolor{textcolor}{rgb}{0.000000,0.000000,0.000000}%
\pgfsetstrokecolor{textcolor}%
\pgfsetfillcolor{textcolor}%
\pgftext[x=0.464387in,y=1.856760in,left,top]{\color{textcolor}{\rmfamily\fontsize{9.000000}{10.800000}\selectfont\catcode`\^=\active\def^{\ifmmode\sp\else\^{}\fi}\catcode`\%=\active\def%{\%}$(a)$}}%
\end{pgfscope}%
\begin{pgfscope}%
\pgfpathrectangle{\pgfqpoint{0.415139in}{0.462654in}}{\pgfqpoint{1.641596in}{1.437222in}}%
\pgfusepath{clip}%
\definecolor{textcolor}{rgb}{0.000000,0.000000,0.000000}%
\pgfsetstrokecolor{textcolor}%
\pgfsetfillcolor{textcolor}%
\pgftext[x=0.845788in, y=1.248118in, left, base,rotate=39.406269]{\color{textcolor}{\rmfamily\fontsize{7.000000}{8.400000}\selectfont\catcode`\^=\active\def^{\ifmmode\sp\else\^{}\fi}\catcode`\%=\active\def%{\%}$0.5$}}%
\end{pgfscope}%
\begin{pgfscope}%
\pgfpathrectangle{\pgfqpoint{0.415139in}{0.462654in}}{\pgfqpoint{1.641596in}{1.437222in}}%
\pgfusepath{clip}%
\definecolor{textcolor}{rgb}{0.000000,0.000000,0.000000}%
\pgfsetstrokecolor{textcolor}%
\pgfsetfillcolor{textcolor}%
\pgftext[x=1.627262in, y=0.697773in, left, base,rotate=38.340704]{\color{textcolor}{\rmfamily\fontsize{7.000000}{8.400000}\selectfont\catcode`\^=\active\def^{\ifmmode\sp\else\^{}\fi}\catcode`\%=\active\def%{\%}$0.8$}}%
\end{pgfscope}%
\begin{pgfscope}%
\pgfsetbuttcap%
\pgfsetmiterjoin%
\definecolor{currentfill}{rgb}{1.000000,1.000000,1.000000}%
\pgfsetfillcolor{currentfill}%
\pgfsetlinewidth{0.000000pt}%
\definecolor{currentstroke}{rgb}{0.000000,0.000000,0.000000}%
\pgfsetstrokecolor{currentstroke}%
\pgfsetstrokeopacity{0.000000}%
\pgfsetdash{}{0pt}%
\pgfpathmoveto{\pgfqpoint{2.987869in}{0.462654in}}%
\pgfpathlineto{\pgfqpoint{4.791820in}{0.462654in}}%
\pgfpathlineto{\pgfqpoint{4.791820in}{1.899877in}}%
\pgfpathlineto{\pgfqpoint{2.987869in}{1.899877in}}%
\pgfpathlineto{\pgfqpoint{2.987869in}{0.462654in}}%
\pgfpathclose%
\pgfusepath{fill}%
\end{pgfscope}%
\begin{pgfscope}%
\pgfsetbuttcap%
\pgfsetroundjoin%
\definecolor{currentfill}{rgb}{0.000000,0.000000,0.000000}%
\pgfsetfillcolor{currentfill}%
\pgfsetlinewidth{0.803000pt}%
\definecolor{currentstroke}{rgb}{0.000000,0.000000,0.000000}%
\pgfsetstrokecolor{currentstroke}%
\pgfsetdash{}{0pt}%
\pgfsys@defobject{currentmarker}{\pgfqpoint{0.000000in}{-0.048611in}}{\pgfqpoint{0.000000in}{0.000000in}}{%
\pgfpathmoveto{\pgfqpoint{0.000000in}{0.000000in}}%
\pgfpathlineto{\pgfqpoint{0.000000in}{-0.048611in}}%
\pgfusepath{stroke,fill}%
}%
\begin{pgfscope}%
\pgfsys@transformshift{3.044243in}{0.462654in}%
\pgfsys@useobject{currentmarker}{}%
\end{pgfscope}%
\end{pgfscope}%
\begin{pgfscope}%
\definecolor{textcolor}{rgb}{0.000000,0.000000,0.000000}%
\pgfsetstrokecolor{textcolor}%
\pgfsetfillcolor{textcolor}%
\pgftext[x=3.044243in,y=0.365432in,,top]{\color{textcolor}{\rmfamily\fontsize{8.000000}{9.600000}\selectfont\catcode`\^=\active\def^{\ifmmode\sp\else\^{}\fi}\catcode`\%=\active\def%{\%}0.0}}%
\end{pgfscope}%
\begin{pgfscope}%
\pgfsetbuttcap%
\pgfsetroundjoin%
\definecolor{currentfill}{rgb}{0.000000,0.000000,0.000000}%
\pgfsetfillcolor{currentfill}%
\pgfsetlinewidth{0.803000pt}%
\definecolor{currentstroke}{rgb}{0.000000,0.000000,0.000000}%
\pgfsetstrokecolor{currentstroke}%
\pgfsetdash{}{0pt}%
\pgfsys@defobject{currentmarker}{\pgfqpoint{0.000000in}{-0.048611in}}{\pgfqpoint{0.000000in}{0.000000in}}{%
\pgfpathmoveto{\pgfqpoint{0.000000in}{0.000000in}}%
\pgfpathlineto{\pgfqpoint{0.000000in}{-0.048611in}}%
\pgfusepath{stroke,fill}%
}%
\begin{pgfscope}%
\pgfsys@transformshift{3.607977in}{0.462654in}%
\pgfsys@useobject{currentmarker}{}%
\end{pgfscope}%
\end{pgfscope}%
\begin{pgfscope}%
\definecolor{textcolor}{rgb}{0.000000,0.000000,0.000000}%
\pgfsetstrokecolor{textcolor}%
\pgfsetfillcolor{textcolor}%
\pgftext[x=3.607977in,y=0.365432in,,top]{\color{textcolor}{\rmfamily\fontsize{8.000000}{9.600000}\selectfont\catcode`\^=\active\def^{\ifmmode\sp\else\^{}\fi}\catcode`\%=\active\def%{\%}0.1}}%
\end{pgfscope}%
\begin{pgfscope}%
\pgfsetbuttcap%
\pgfsetroundjoin%
\definecolor{currentfill}{rgb}{0.000000,0.000000,0.000000}%
\pgfsetfillcolor{currentfill}%
\pgfsetlinewidth{0.803000pt}%
\definecolor{currentstroke}{rgb}{0.000000,0.000000,0.000000}%
\pgfsetstrokecolor{currentstroke}%
\pgfsetdash{}{0pt}%
\pgfsys@defobject{currentmarker}{\pgfqpoint{0.000000in}{-0.048611in}}{\pgfqpoint{0.000000in}{0.000000in}}{%
\pgfpathmoveto{\pgfqpoint{0.000000in}{0.000000in}}%
\pgfpathlineto{\pgfqpoint{0.000000in}{-0.048611in}}%
\pgfusepath{stroke,fill}%
}%
\begin{pgfscope}%
\pgfsys@transformshift{4.171712in}{0.462654in}%
\pgfsys@useobject{currentmarker}{}%
\end{pgfscope}%
\end{pgfscope}%
\begin{pgfscope}%
\definecolor{textcolor}{rgb}{0.000000,0.000000,0.000000}%
\pgfsetstrokecolor{textcolor}%
\pgfsetfillcolor{textcolor}%
\pgftext[x=4.171712in,y=0.365432in,,top]{\color{textcolor}{\rmfamily\fontsize{8.000000}{9.600000}\selectfont\catcode`\^=\active\def^{\ifmmode\sp\else\^{}\fi}\catcode`\%=\active\def%{\%}0.2}}%
\end{pgfscope}%
\begin{pgfscope}%
\pgfsetbuttcap%
\pgfsetroundjoin%
\definecolor{currentfill}{rgb}{0.000000,0.000000,0.000000}%
\pgfsetfillcolor{currentfill}%
\pgfsetlinewidth{0.803000pt}%
\definecolor{currentstroke}{rgb}{0.000000,0.000000,0.000000}%
\pgfsetstrokecolor{currentstroke}%
\pgfsetdash{}{0pt}%
\pgfsys@defobject{currentmarker}{\pgfqpoint{0.000000in}{-0.048611in}}{\pgfqpoint{0.000000in}{0.000000in}}{%
\pgfpathmoveto{\pgfqpoint{0.000000in}{0.000000in}}%
\pgfpathlineto{\pgfqpoint{0.000000in}{-0.048611in}}%
\pgfusepath{stroke,fill}%
}%
\begin{pgfscope}%
\pgfsys@transformshift{4.735447in}{0.462654in}%
\pgfsys@useobject{currentmarker}{}%
\end{pgfscope}%
\end{pgfscope}%
\begin{pgfscope}%
\definecolor{textcolor}{rgb}{0.000000,0.000000,0.000000}%
\pgfsetstrokecolor{textcolor}%
\pgfsetfillcolor{textcolor}%
\pgftext[x=4.735447in,y=0.365432in,,top]{\color{textcolor}{\rmfamily\fontsize{8.000000}{9.600000}\selectfont\catcode`\^=\active\def^{\ifmmode\sp\else\^{}\fi}\catcode`\%=\active\def%{\%}0.3}}%
\end{pgfscope}%
\begin{pgfscope}%
\definecolor{textcolor}{rgb}{0.000000,0.000000,0.000000}%
\pgfsetstrokecolor{textcolor}%
\pgfsetfillcolor{textcolor}%
\pgftext[x=3.889845in,y=0.211111in,,top]{\color{textcolor}{\rmfamily\fontsize{9.000000}{10.800000}\selectfont\catcode`\^=\active\def^{\ifmmode\sp\else\^{}\fi}\catcode`\%=\active\def%{\%}$\delta$}}%
\end{pgfscope}%
\begin{pgfscope}%
\pgfsetbuttcap%
\pgfsetroundjoin%
\definecolor{currentfill}{rgb}{0.000000,0.000000,0.000000}%
\pgfsetfillcolor{currentfill}%
\pgfsetlinewidth{0.803000pt}%
\definecolor{currentstroke}{rgb}{0.000000,0.000000,0.000000}%
\pgfsetstrokecolor{currentstroke}%
\pgfsetdash{}{0pt}%
\pgfsys@defobject{currentmarker}{\pgfqpoint{-0.048611in}{0.000000in}}{\pgfqpoint{-0.000000in}{0.000000in}}{%
\pgfpathmoveto{\pgfqpoint{-0.000000in}{0.000000in}}%
\pgfpathlineto{\pgfqpoint{-0.048611in}{0.000000in}}%
\pgfusepath{stroke,fill}%
}%
\begin{pgfscope}%
\pgfsys@transformshift{2.987869in}{0.462654in}%
\pgfsys@useobject{currentmarker}{}%
\end{pgfscope}%
\end{pgfscope}%
\begin{pgfscope}%
\definecolor{textcolor}{rgb}{0.000000,0.000000,0.000000}%
\pgfsetstrokecolor{textcolor}%
\pgfsetfillcolor{textcolor}%
\pgftext[x=2.739796in, y=0.424074in, left, base]{\color{textcolor}{\rmfamily\fontsize{8.000000}{9.600000}\selectfont\catcode`\^=\active\def^{\ifmmode\sp\else\^{}\fi}\catcode`\%=\active\def%{\%}0.5}}%
\end{pgfscope}%
\begin{pgfscope}%
\pgfsetbuttcap%
\pgfsetroundjoin%
\definecolor{currentfill}{rgb}{0.000000,0.000000,0.000000}%
\pgfsetfillcolor{currentfill}%
\pgfsetlinewidth{0.803000pt}%
\definecolor{currentstroke}{rgb}{0.000000,0.000000,0.000000}%
\pgfsetstrokecolor{currentstroke}%
\pgfsetdash{}{0pt}%
\pgfsys@defobject{currentmarker}{\pgfqpoint{-0.048611in}{0.000000in}}{\pgfqpoint{-0.000000in}{0.000000in}}{%
\pgfpathmoveto{\pgfqpoint{-0.000000in}{0.000000in}}%
\pgfpathlineto{\pgfqpoint{-0.048611in}{0.000000in}}%
\pgfusepath{stroke,fill}%
}%
\begin{pgfscope}%
\pgfsys@transformshift{2.987869in}{0.750099in}%
\pgfsys@useobject{currentmarker}{}%
\end{pgfscope}%
\end{pgfscope}%
\begin{pgfscope}%
\definecolor{textcolor}{rgb}{0.000000,0.000000,0.000000}%
\pgfsetstrokecolor{textcolor}%
\pgfsetfillcolor{textcolor}%
\pgftext[x=2.739796in, y=0.711519in, left, base]{\color{textcolor}{\rmfamily\fontsize{8.000000}{9.600000}\selectfont\catcode`\^=\active\def^{\ifmmode\sp\else\^{}\fi}\catcode`\%=\active\def%{\%}0.6}}%
\end{pgfscope}%
\begin{pgfscope}%
\pgfsetbuttcap%
\pgfsetroundjoin%
\definecolor{currentfill}{rgb}{0.000000,0.000000,0.000000}%
\pgfsetfillcolor{currentfill}%
\pgfsetlinewidth{0.803000pt}%
\definecolor{currentstroke}{rgb}{0.000000,0.000000,0.000000}%
\pgfsetstrokecolor{currentstroke}%
\pgfsetdash{}{0pt}%
\pgfsys@defobject{currentmarker}{\pgfqpoint{-0.048611in}{0.000000in}}{\pgfqpoint{-0.000000in}{0.000000in}}{%
\pgfpathmoveto{\pgfqpoint{-0.000000in}{0.000000in}}%
\pgfpathlineto{\pgfqpoint{-0.048611in}{0.000000in}}%
\pgfusepath{stroke,fill}%
}%
\begin{pgfscope}%
\pgfsys@transformshift{2.987869in}{1.037543in}%
\pgfsys@useobject{currentmarker}{}%
\end{pgfscope}%
\end{pgfscope}%
\begin{pgfscope}%
\definecolor{textcolor}{rgb}{0.000000,0.000000,0.000000}%
\pgfsetstrokecolor{textcolor}%
\pgfsetfillcolor{textcolor}%
\pgftext[x=2.739796in, y=0.998963in, left, base]{\color{textcolor}{\rmfamily\fontsize{8.000000}{9.600000}\selectfont\catcode`\^=\active\def^{\ifmmode\sp\else\^{}\fi}\catcode`\%=\active\def%{\%}0.7}}%
\end{pgfscope}%
\begin{pgfscope}%
\pgfsetbuttcap%
\pgfsetroundjoin%
\definecolor{currentfill}{rgb}{0.000000,0.000000,0.000000}%
\pgfsetfillcolor{currentfill}%
\pgfsetlinewidth{0.803000pt}%
\definecolor{currentstroke}{rgb}{0.000000,0.000000,0.000000}%
\pgfsetstrokecolor{currentstroke}%
\pgfsetdash{}{0pt}%
\pgfsys@defobject{currentmarker}{\pgfqpoint{-0.048611in}{0.000000in}}{\pgfqpoint{-0.000000in}{0.000000in}}{%
\pgfpathmoveto{\pgfqpoint{-0.000000in}{0.000000in}}%
\pgfpathlineto{\pgfqpoint{-0.048611in}{0.000000in}}%
\pgfusepath{stroke,fill}%
}%
\begin{pgfscope}%
\pgfsys@transformshift{2.987869in}{1.324988in}%
\pgfsys@useobject{currentmarker}{}%
\end{pgfscope}%
\end{pgfscope}%
\begin{pgfscope}%
\definecolor{textcolor}{rgb}{0.000000,0.000000,0.000000}%
\pgfsetstrokecolor{textcolor}%
\pgfsetfillcolor{textcolor}%
\pgftext[x=2.739796in, y=1.286407in, left, base]{\color{textcolor}{\rmfamily\fontsize{8.000000}{9.600000}\selectfont\catcode`\^=\active\def^{\ifmmode\sp\else\^{}\fi}\catcode`\%=\active\def%{\%}0.8}}%
\end{pgfscope}%
\begin{pgfscope}%
\pgfsetbuttcap%
\pgfsetroundjoin%
\definecolor{currentfill}{rgb}{0.000000,0.000000,0.000000}%
\pgfsetfillcolor{currentfill}%
\pgfsetlinewidth{0.803000pt}%
\definecolor{currentstroke}{rgb}{0.000000,0.000000,0.000000}%
\pgfsetstrokecolor{currentstroke}%
\pgfsetdash{}{0pt}%
\pgfsys@defobject{currentmarker}{\pgfqpoint{-0.048611in}{0.000000in}}{\pgfqpoint{-0.000000in}{0.000000in}}{%
\pgfpathmoveto{\pgfqpoint{-0.000000in}{0.000000in}}%
\pgfpathlineto{\pgfqpoint{-0.048611in}{0.000000in}}%
\pgfusepath{stroke,fill}%
}%
\begin{pgfscope}%
\pgfsys@transformshift{2.987869in}{1.612432in}%
\pgfsys@useobject{currentmarker}{}%
\end{pgfscope}%
\end{pgfscope}%
\begin{pgfscope}%
\definecolor{textcolor}{rgb}{0.000000,0.000000,0.000000}%
\pgfsetstrokecolor{textcolor}%
\pgfsetfillcolor{textcolor}%
\pgftext[x=2.739796in, y=1.573852in, left, base]{\color{textcolor}{\rmfamily\fontsize{8.000000}{9.600000}\selectfont\catcode`\^=\active\def^{\ifmmode\sp\else\^{}\fi}\catcode`\%=\active\def%{\%}0.9}}%
\end{pgfscope}%
\begin{pgfscope}%
\pgfsetbuttcap%
\pgfsetroundjoin%
\definecolor{currentfill}{rgb}{0.000000,0.000000,0.000000}%
\pgfsetfillcolor{currentfill}%
\pgfsetlinewidth{0.803000pt}%
\definecolor{currentstroke}{rgb}{0.000000,0.000000,0.000000}%
\pgfsetstrokecolor{currentstroke}%
\pgfsetdash{}{0pt}%
\pgfsys@defobject{currentmarker}{\pgfqpoint{-0.048611in}{0.000000in}}{\pgfqpoint{-0.000000in}{0.000000in}}{%
\pgfpathmoveto{\pgfqpoint{-0.000000in}{0.000000in}}%
\pgfpathlineto{\pgfqpoint{-0.048611in}{0.000000in}}%
\pgfusepath{stroke,fill}%
}%
\begin{pgfscope}%
\pgfsys@transformshift{2.987869in}{1.899877in}%
\pgfsys@useobject{currentmarker}{}%
\end{pgfscope}%
\end{pgfscope}%
\begin{pgfscope}%
\definecolor{textcolor}{rgb}{0.000000,0.000000,0.000000}%
\pgfsetstrokecolor{textcolor}%
\pgfsetfillcolor{textcolor}%
\pgftext[x=2.739796in, y=1.861296in, left, base]{\color{textcolor}{\rmfamily\fontsize{8.000000}{9.600000}\selectfont\catcode`\^=\active\def^{\ifmmode\sp\else\^{}\fi}\catcode`\%=\active\def%{\%}1.0}}%
\end{pgfscope}%
\begin{pgfscope}%
\pgfpathrectangle{\pgfqpoint{2.987869in}{0.462654in}}{\pgfqpoint{1.803951in}{1.437222in}}%
\pgfusepath{clip}%
\pgfsetrectcap%
\pgfsetroundjoin%
\pgfsetlinewidth{0.903375pt}%
\definecolor{currentstroke}{rgb}{1.000000,0.000000,0.000000}%
\pgfsetstrokecolor{currentstroke}%
\pgfsetdash{}{0pt}%
\pgfpathmoveto{\pgfqpoint{3.044243in}{1.507604in}}%
\pgfpathlineto{\pgfqpoint{3.452171in}{1.506861in}}%
\pgfpathlineto{\pgfqpoint{3.817607in}{1.503947in}}%
\pgfpathlineto{\pgfqpoint{4.200041in}{1.498626in}}%
\pgfpathlineto{\pgfqpoint{4.573975in}{1.491199in}}%
\pgfpathlineto{\pgfqpoint{4.735447in}{1.487321in}}%
\pgfpathlineto{\pgfqpoint{4.735447in}{1.487321in}}%
\pgfusepath{stroke}%
\end{pgfscope}%
\begin{pgfscope}%
\pgfpathrectangle{\pgfqpoint{2.987869in}{0.462654in}}{\pgfqpoint{1.803951in}{1.437222in}}%
\pgfusepath{clip}%
\pgfsetbuttcap%
\pgfsetmiterjoin%
\definecolor{currentfill}{rgb}{0.000000,0.000000,0.000000}%
\pgfsetfillcolor{currentfill}%
\pgfsetfillopacity{0.000000}%
\pgfsetlinewidth{0.803000pt}%
\definecolor{currentstroke}{rgb}{1.000000,0.000000,0.000000}%
\pgfsetstrokecolor{currentstroke}%
\pgfsetdash{}{0pt}%
\pgfsys@defobject{currentmarker}{\pgfqpoint{-0.017500in}{-0.017500in}}{\pgfqpoint{0.017500in}{0.017500in}}{%
\pgfpathmoveto{\pgfqpoint{-0.017500in}{-0.017500in}}%
\pgfpathlineto{\pgfqpoint{0.017500in}{-0.017500in}}%
\pgfpathlineto{\pgfqpoint{0.017500in}{0.017500in}}%
\pgfpathlineto{\pgfqpoint{-0.017500in}{0.017500in}}%
\pgfpathlineto{\pgfqpoint{-0.017500in}{-0.017500in}}%
\pgfpathclose%
\pgfusepath{stroke,fill}%
}%
\begin{pgfscope}%
\pgfsys@transformshift{3.044243in}{1.507604in}%
\pgfsys@useobject{currentmarker}{}%
\end{pgfscope}%
\begin{pgfscope}%
\pgfsys@transformshift{3.185176in}{1.507648in}%
\pgfsys@useobject{currentmarker}{}%
\end{pgfscope}%
\begin{pgfscope}%
\pgfsys@transformshift{3.326110in}{1.507376in}%
\pgfsys@useobject{currentmarker}{}%
\end{pgfscope}%
\begin{pgfscope}%
\pgfsys@transformshift{3.467044in}{1.506786in}%
\pgfsys@useobject{currentmarker}{}%
\end{pgfscope}%
\begin{pgfscope}%
\pgfsys@transformshift{3.607977in}{1.505879in}%
\pgfsys@useobject{currentmarker}{}%
\end{pgfscope}%
\begin{pgfscope}%
\pgfsys@transformshift{3.748911in}{1.504657in}%
\pgfsys@useobject{currentmarker}{}%
\end{pgfscope}%
\begin{pgfscope}%
\pgfsys@transformshift{3.889845in}{1.503118in}%
\pgfsys@useobject{currentmarker}{}%
\end{pgfscope}%
\begin{pgfscope}%
\pgfsys@transformshift{4.030779in}{1.501265in}%
\pgfsys@useobject{currentmarker}{}%
\end{pgfscope}%
\begin{pgfscope}%
\pgfsys@transformshift{4.171712in}{1.499098in}%
\pgfsys@useobject{currentmarker}{}%
\end{pgfscope}%
\begin{pgfscope}%
\pgfsys@transformshift{4.312646in}{1.496619in}%
\pgfsys@useobject{currentmarker}{}%
\end{pgfscope}%
\begin{pgfscope}%
\pgfsys@transformshift{4.453580in}{1.493828in}%
\pgfsys@useobject{currentmarker}{}%
\end{pgfscope}%
\begin{pgfscope}%
\pgfsys@transformshift{4.594513in}{1.490728in}%
\pgfsys@useobject{currentmarker}{}%
\end{pgfscope}%
\begin{pgfscope}%
\pgfsys@transformshift{4.735447in}{1.487321in}%
\pgfsys@useobject{currentmarker}{}%
\end{pgfscope}%
\end{pgfscope}%
\begin{pgfscope}%
\pgfpathrectangle{\pgfqpoint{2.987869in}{0.462654in}}{\pgfqpoint{1.803951in}{1.437222in}}%
\pgfusepath{clip}%
\pgfsetrectcap%
\pgfsetroundjoin%
\pgfsetlinewidth{0.903375pt}%
\definecolor{currentstroke}{rgb}{1.000000,0.000000,0.000000}%
\pgfsetstrokecolor{currentstroke}%
\pgfsetdash{}{0pt}%
\pgfpathmoveto{\pgfqpoint{3.044243in}{1.507604in}}%
\pgfpathlineto{\pgfqpoint{3.239709in}{1.506440in}}%
\pgfpathlineto{\pgfqpoint{3.435174in}{1.502816in}}%
\pgfpathlineto{\pgfqpoint{3.622142in}{1.497168in}}%
\pgfpathlineto{\pgfqpoint{3.809109in}{1.489329in}}%
\pgfpathlineto{\pgfqpoint{4.004575in}{1.478817in}}%
\pgfpathlineto{\pgfqpoint{4.191542in}{1.466581in}}%
\pgfpathlineto{\pgfqpoint{4.378509in}{1.452210in}}%
\pgfpathlineto{\pgfqpoint{4.573975in}{1.434941in}}%
\pgfpathlineto{\pgfqpoint{4.735447in}{1.418979in}}%
\pgfpathlineto{\pgfqpoint{4.735447in}{1.418979in}}%
\pgfusepath{stroke}%
\end{pgfscope}%
\begin{pgfscope}%
\pgfpathrectangle{\pgfqpoint{2.987869in}{0.462654in}}{\pgfqpoint{1.803951in}{1.437222in}}%
\pgfusepath{clip}%
\pgfsetbuttcap%
\pgfsetroundjoin%
\definecolor{currentfill}{rgb}{1.000000,0.000000,0.000000}%
\pgfsetfillcolor{currentfill}%
\pgfsetlinewidth{0.803000pt}%
\definecolor{currentstroke}{rgb}{1.000000,0.000000,0.000000}%
\pgfsetstrokecolor{currentstroke}%
\pgfsetdash{}{0pt}%
\pgfsys@defobject{currentmarker}{\pgfqpoint{-0.019444in}{-0.019444in}}{\pgfqpoint{0.019444in}{0.019444in}}{%
\pgfpathmoveto{\pgfqpoint{-0.019444in}{-0.019444in}}%
\pgfpathlineto{\pgfqpoint{0.019444in}{0.019444in}}%
\pgfpathmoveto{\pgfqpoint{-0.019444in}{0.019444in}}%
\pgfpathlineto{\pgfqpoint{0.019444in}{-0.019444in}}%
\pgfusepath{stroke,fill}%
}%
\begin{pgfscope}%
\pgfsys@transformshift{3.044243in}{1.507604in}%
\pgfsys@useobject{currentmarker}{}%
\end{pgfscope}%
\begin{pgfscope}%
\pgfsys@transformshift{3.185176in}{1.506983in}%
\pgfsys@useobject{currentmarker}{}%
\end{pgfscope}%
\begin{pgfscope}%
\pgfsys@transformshift{3.326110in}{1.505119in}%
\pgfsys@useobject{currentmarker}{}%
\end{pgfscope}%
\begin{pgfscope}%
\pgfsys@transformshift{3.467044in}{1.502013in}%
\pgfsys@useobject{currentmarker}{}%
\end{pgfscope}%
\begin{pgfscope}%
\pgfsys@transformshift{3.607977in}{1.497667in}%
\pgfsys@useobject{currentmarker}{}%
\end{pgfscope}%
\begin{pgfscope}%
\pgfsys@transformshift{3.748911in}{1.492086in}%
\pgfsys@useobject{currentmarker}{}%
\end{pgfscope}%
\begin{pgfscope}%
\pgfsys@transformshift{3.889845in}{1.485273in}%
\pgfsys@useobject{currentmarker}{}%
\end{pgfscope}%
\begin{pgfscope}%
\pgfsys@transformshift{4.030779in}{1.477234in}%
\pgfsys@useobject{currentmarker}{}%
\end{pgfscope}%
\begin{pgfscope}%
\pgfsys@transformshift{4.171712in}{1.467977in}%
\pgfsys@useobject{currentmarker}{}%
\end{pgfscope}%
\begin{pgfscope}%
\pgfsys@transformshift{4.312646in}{1.457511in}%
\pgfsys@useobject{currentmarker}{}%
\end{pgfscope}%
\begin{pgfscope}%
\pgfsys@transformshift{4.453580in}{1.445847in}%
\pgfsys@useobject{currentmarker}{}%
\end{pgfscope}%
\begin{pgfscope}%
\pgfsys@transformshift{4.594513in}{1.432998in}%
\pgfsys@useobject{currentmarker}{}%
\end{pgfscope}%
\begin{pgfscope}%
\pgfsys@transformshift{4.735447in}{1.418979in}%
\pgfsys@useobject{currentmarker}{}%
\end{pgfscope}%
\end{pgfscope}%
\begin{pgfscope}%
\pgfpathrectangle{\pgfqpoint{2.987869in}{0.462654in}}{\pgfqpoint{1.803951in}{1.437222in}}%
\pgfusepath{clip}%
\pgfsetrectcap%
\pgfsetroundjoin%
\pgfsetlinewidth{0.903375pt}%
\definecolor{currentstroke}{rgb}{0.172549,0.627451,0.172549}%
\pgfsetstrokecolor{currentstroke}%
\pgfsetdash{}{0pt}%
\pgfpathmoveto{\pgfqpoint{3.044243in}{1.409880in}}%
\pgfpathlineto{\pgfqpoint{3.265204in}{1.408731in}}%
\pgfpathlineto{\pgfqpoint{3.477667in}{1.404842in}}%
\pgfpathlineto{\pgfqpoint{3.664634in}{1.399207in}}%
\pgfpathlineto{\pgfqpoint{3.860100in}{1.391120in}}%
\pgfpathlineto{\pgfqpoint{4.055566in}{1.380859in}}%
\pgfpathlineto{\pgfqpoint{4.259530in}{1.367867in}}%
\pgfpathlineto{\pgfqpoint{4.463494in}{1.352642in}}%
\pgfpathlineto{\pgfqpoint{4.667459in}{1.335272in}}%
\pgfpathlineto{\pgfqpoint{4.735447in}{1.329018in}}%
\pgfpathlineto{\pgfqpoint{4.735447in}{1.329018in}}%
\pgfusepath{stroke}%
\end{pgfscope}%
\begin{pgfscope}%
\pgfpathrectangle{\pgfqpoint{2.987869in}{0.462654in}}{\pgfqpoint{1.803951in}{1.437222in}}%
\pgfusepath{clip}%
\pgfsetbuttcap%
\pgfsetmiterjoin%
\definecolor{currentfill}{rgb}{0.000000,0.000000,0.000000}%
\pgfsetfillcolor{currentfill}%
\pgfsetfillopacity{0.000000}%
\pgfsetlinewidth{0.803000pt}%
\definecolor{currentstroke}{rgb}{0.172549,0.627451,0.172549}%
\pgfsetstrokecolor{currentstroke}%
\pgfsetdash{}{0pt}%
\pgfsys@defobject{currentmarker}{\pgfqpoint{-0.017500in}{-0.017500in}}{\pgfqpoint{0.017500in}{0.017500in}}{%
\pgfpathmoveto{\pgfqpoint{-0.017500in}{-0.017500in}}%
\pgfpathlineto{\pgfqpoint{0.017500in}{-0.017500in}}%
\pgfpathlineto{\pgfqpoint{0.017500in}{0.017500in}}%
\pgfpathlineto{\pgfqpoint{-0.017500in}{0.017500in}}%
\pgfpathlineto{\pgfqpoint{-0.017500in}{-0.017500in}}%
\pgfpathclose%
\pgfusepath{stroke,fill}%
}%
\begin{pgfscope}%
\pgfsys@transformshift{3.044243in}{1.409880in}%
\pgfsys@useobject{currentmarker}{}%
\end{pgfscope}%
\begin{pgfscope}%
\pgfsys@transformshift{3.185176in}{1.409477in}%
\pgfsys@useobject{currentmarker}{}%
\end{pgfscope}%
\begin{pgfscope}%
\pgfsys@transformshift{3.326110in}{1.407879in}%
\pgfsys@useobject{currentmarker}{}%
\end{pgfscope}%
\begin{pgfscope}%
\pgfsys@transformshift{3.467044in}{1.405093in}%
\pgfsys@useobject{currentmarker}{}%
\end{pgfscope}%
\begin{pgfscope}%
\pgfsys@transformshift{3.607977in}{1.401127in}%
\pgfsys@useobject{currentmarker}{}%
\end{pgfscope}%
\begin{pgfscope}%
\pgfsys@transformshift{3.748911in}{1.395991in}%
\pgfsys@useobject{currentmarker}{}%
\end{pgfscope}%
\begin{pgfscope}%
\pgfsys@transformshift{3.889845in}{1.389701in}%
\pgfsys@useobject{currentmarker}{}%
\end{pgfscope}%
\begin{pgfscope}%
\pgfsys@transformshift{4.030779in}{1.382277in}%
\pgfsys@useobject{currentmarker}{}%
\end{pgfscope}%
\begin{pgfscope}%
\pgfsys@transformshift{4.171712in}{1.373740in}%
\pgfsys@useobject{currentmarker}{}%
\end{pgfscope}%
\begin{pgfscope}%
\pgfsys@transformshift{4.312646in}{1.364115in}%
\pgfsys@useobject{currentmarker}{}%
\end{pgfscope}%
\begin{pgfscope}%
\pgfsys@transformshift{4.453580in}{1.353431in}%
\pgfsys@useobject{currentmarker}{}%
\end{pgfscope}%
\begin{pgfscope}%
\pgfsys@transformshift{4.594513in}{1.341721in}%
\pgfsys@useobject{currentmarker}{}%
\end{pgfscope}%
\begin{pgfscope}%
\pgfsys@transformshift{4.735447in}{1.329018in}%
\pgfsys@useobject{currentmarker}{}%
\end{pgfscope}%
\end{pgfscope}%
\begin{pgfscope}%
\pgfpathrectangle{\pgfqpoint{2.987869in}{0.462654in}}{\pgfqpoint{1.803951in}{1.437222in}}%
\pgfusepath{clip}%
\pgfsetrectcap%
\pgfsetroundjoin%
\pgfsetlinewidth{0.903375pt}%
\definecolor{currentstroke}{rgb}{0.172549,0.627451,0.172549}%
\pgfsetstrokecolor{currentstroke}%
\pgfsetdash{}{0pt}%
\pgfpathmoveto{\pgfqpoint{3.044243in}{1.409880in}}%
\pgfpathlineto{\pgfqpoint{3.137726in}{1.408700in}}%
\pgfpathlineto{\pgfqpoint{3.248207in}{1.405170in}}%
\pgfpathlineto{\pgfqpoint{3.358688in}{1.398541in}}%
\pgfpathlineto{\pgfqpoint{3.452171in}{1.390702in}}%
\pgfpathlineto{\pgfqpoint{3.545655in}{1.381022in}}%
\pgfpathlineto{\pgfqpoint{3.656136in}{1.367090in}}%
\pgfpathlineto{\pgfqpoint{3.758118in}{1.351864in}}%
\pgfpathlineto{\pgfqpoint{3.860100in}{1.334447in}}%
\pgfpathlineto{\pgfqpoint{3.970581in}{1.313250in}}%
\pgfpathlineto{\pgfqpoint{4.081061in}{1.289683in}}%
\pgfpathlineto{\pgfqpoint{4.191542in}{1.263844in}}%
\pgfpathlineto{\pgfqpoint{4.302023in}{1.235853in}}%
\pgfpathlineto{\pgfqpoint{4.421002in}{1.203495in}}%
\pgfpathlineto{\pgfqpoint{4.548480in}{1.166464in}}%
\pgfpathlineto{\pgfqpoint{4.684456in}{1.124554in}}%
\pgfpathlineto{\pgfqpoint{4.735447in}{1.108212in}}%
\pgfpathlineto{\pgfqpoint{4.735447in}{1.108212in}}%
\pgfusepath{stroke}%
\end{pgfscope}%
\begin{pgfscope}%
\pgfpathrectangle{\pgfqpoint{2.987869in}{0.462654in}}{\pgfqpoint{1.803951in}{1.437222in}}%
\pgfusepath{clip}%
\pgfsetbuttcap%
\pgfsetroundjoin%
\definecolor{currentfill}{rgb}{0.172549,0.627451,0.172549}%
\pgfsetfillcolor{currentfill}%
\pgfsetlinewidth{0.803000pt}%
\definecolor{currentstroke}{rgb}{0.172549,0.627451,0.172549}%
\pgfsetstrokecolor{currentstroke}%
\pgfsetdash{}{0pt}%
\pgfsys@defobject{currentmarker}{\pgfqpoint{-0.019444in}{-0.019444in}}{\pgfqpoint{0.019444in}{0.019444in}}{%
\pgfpathmoveto{\pgfqpoint{-0.019444in}{-0.019444in}}%
\pgfpathlineto{\pgfqpoint{0.019444in}{0.019444in}}%
\pgfpathmoveto{\pgfqpoint{-0.019444in}{0.019444in}}%
\pgfpathlineto{\pgfqpoint{0.019444in}{-0.019444in}}%
\pgfusepath{stroke,fill}%
}%
\begin{pgfscope}%
\pgfsys@transformshift{3.044243in}{1.409880in}%
\pgfsys@useobject{currentmarker}{}%
\end{pgfscope}%
\begin{pgfscope}%
\pgfsys@transformshift{3.185176in}{1.407586in}%
\pgfsys@useobject{currentmarker}{}%
\end{pgfscope}%
\begin{pgfscope}%
\pgfsys@transformshift{3.326110in}{1.400713in}%
\pgfsys@useobject{currentmarker}{}%
\end{pgfscope}%
\begin{pgfscope}%
\pgfsys@transformshift{3.467044in}{1.389314in}%
\pgfsys@useobject{currentmarker}{}%
\end{pgfscope}%
\begin{pgfscope}%
\pgfsys@transformshift{3.607977in}{1.373475in}%
\pgfsys@useobject{currentmarker}{}%
\end{pgfscope}%
\begin{pgfscope}%
\pgfsys@transformshift{3.748911in}{1.353316in}%
\pgfsys@useobject{currentmarker}{}%
\end{pgfscope}%
\begin{pgfscope}%
\pgfsys@transformshift{3.889845in}{1.328992in}%
\pgfsys@useobject{currentmarker}{}%
\end{pgfscope}%
\begin{pgfscope}%
\pgfsys@transformshift{4.030779in}{1.300690in}%
\pgfsys@useobject{currentmarker}{}%
\end{pgfscope}%
\begin{pgfscope}%
\pgfsys@transformshift{4.171712in}{1.268631in}%
\pgfsys@useobject{currentmarker}{}%
\end{pgfscope}%
\begin{pgfscope}%
\pgfsys@transformshift{4.312646in}{1.233063in}%
\pgfsys@useobject{currentmarker}{}%
\end{pgfscope}%
\begin{pgfscope}%
\pgfsys@transformshift{4.453580in}{1.194264in}%
\pgfsys@useobject{currentmarker}{}%
\end{pgfscope}%
\begin{pgfscope}%
\pgfsys@transformshift{4.594513in}{1.152539in}%
\pgfsys@useobject{currentmarker}{}%
\end{pgfscope}%
\begin{pgfscope}%
\pgfsys@transformshift{4.735447in}{1.108212in}%
\pgfsys@useobject{currentmarker}{}%
\end{pgfscope}%
\end{pgfscope}%
\begin{pgfscope}%
\pgfpathrectangle{\pgfqpoint{2.987869in}{0.462654in}}{\pgfqpoint{1.803951in}{1.437222in}}%
\pgfusepath{clip}%
\pgfsetrectcap%
\pgfsetroundjoin%
\pgfsetlinewidth{0.903375pt}%
\definecolor{currentstroke}{rgb}{0.000000,0.000000,0.000000}%
\pgfsetstrokecolor{currentstroke}%
\pgfsetdash{}{0pt}%
\pgfpathmoveto{\pgfqpoint{2.987869in}{1.766257in}}%
\pgfpathlineto{\pgfqpoint{4.791820in}{1.766257in}}%
\pgfusepath{stroke}%
\end{pgfscope}%
\begin{pgfscope}%
\pgfsetrectcap%
\pgfsetmiterjoin%
\pgfsetlinewidth{0.803000pt}%
\definecolor{currentstroke}{rgb}{0.000000,0.000000,0.000000}%
\pgfsetstrokecolor{currentstroke}%
\pgfsetdash{}{0pt}%
\pgfpathmoveto{\pgfqpoint{2.987869in}{0.462654in}}%
\pgfpathlineto{\pgfqpoint{2.987869in}{1.899877in}}%
\pgfusepath{stroke}%
\end{pgfscope}%
\begin{pgfscope}%
\pgfsetrectcap%
\pgfsetmiterjoin%
\pgfsetlinewidth{0.803000pt}%
\definecolor{currentstroke}{rgb}{0.000000,0.000000,0.000000}%
\pgfsetstrokecolor{currentstroke}%
\pgfsetdash{}{0pt}%
\pgfpathmoveto{\pgfqpoint{4.791820in}{0.462654in}}%
\pgfpathlineto{\pgfqpoint{4.791820in}{1.899877in}}%
\pgfusepath{stroke}%
\end{pgfscope}%
\begin{pgfscope}%
\pgfsetrectcap%
\pgfsetmiterjoin%
\pgfsetlinewidth{0.803000pt}%
\definecolor{currentstroke}{rgb}{0.000000,0.000000,0.000000}%
\pgfsetstrokecolor{currentstroke}%
\pgfsetdash{}{0pt}%
\pgfpathmoveto{\pgfqpoint{2.987869in}{0.462654in}}%
\pgfpathlineto{\pgfqpoint{4.791820in}{0.462654in}}%
\pgfusepath{stroke}%
\end{pgfscope}%
\begin{pgfscope}%
\pgfsetrectcap%
\pgfsetmiterjoin%
\pgfsetlinewidth{0.803000pt}%
\definecolor{currentstroke}{rgb}{0.000000,0.000000,0.000000}%
\pgfsetstrokecolor{currentstroke}%
\pgfsetdash{}{0pt}%
\pgfpathmoveto{\pgfqpoint{2.987869in}{1.899877in}}%
\pgfpathlineto{\pgfqpoint{4.791820in}{1.899877in}}%
\pgfusepath{stroke}%
\end{pgfscope}%
\begin{pgfscope}%
\definecolor{textcolor}{rgb}{0.000000,0.000000,0.000000}%
\pgfsetstrokecolor{textcolor}%
\pgfsetfillcolor{textcolor}%
\pgftext[x=4.737702in,y=0.635121in,right,top]{\color{textcolor}{\rmfamily\fontsize{9.000000}{10.800000}\selectfont\catcode`\^=\active\def^{\ifmmode\sp\else\^{}\fi}\catcode`\%=\active\def%{\%}$(b)$}}%
\end{pgfscope}%
\begin{pgfscope}%
\pgfsetbuttcap%
\pgfsetmiterjoin%
\definecolor{currentfill}{rgb}{0.000000,0.000000,0.000000}%
\pgfsetfillcolor{currentfill}%
\pgfsetfillopacity{0.000000}%
\pgfsetlinewidth{0.803000pt}%
\definecolor{currentstroke}{rgb}{1.000000,0.000000,0.000000}%
\pgfsetstrokecolor{currentstroke}%
\pgfsetdash{}{0pt}%
\pgfsys@defobject{currentmarker}{\pgfqpoint{-0.017500in}{-0.017500in}}{\pgfqpoint{0.017500in}{0.017500in}}{%
\pgfpathmoveto{\pgfqpoint{-0.017500in}{-0.017500in}}%
\pgfpathlineto{\pgfqpoint{0.017500in}{-0.017500in}}%
\pgfpathlineto{\pgfqpoint{0.017500in}{0.017500in}}%
\pgfpathlineto{\pgfqpoint{-0.017500in}{0.017500in}}%
\pgfpathlineto{\pgfqpoint{-0.017500in}{-0.017500in}}%
\pgfpathclose%
\pgfusepath{stroke,fill}%
}%
\begin{pgfscope}%
\pgfsys@transformshift{3.132282in}{0.870255in}%
\pgfsys@useobject{currentmarker}{}%
\end{pgfscope}%
\end{pgfscope}%
\begin{pgfscope}%
\definecolor{textcolor}{rgb}{0.000000,0.000000,0.000000}%
\pgfsetstrokecolor{textcolor}%
\pgfsetfillcolor{textcolor}%
\pgftext[x=3.258671in,y=0.838658in,left,base]{\color{textcolor}{\rmfamily\fontsize{6.500000}{7.800000}\selectfont\catcode`\^=\active\def^{\ifmmode\sp\else\^{}\fi}\catcode`\%=\active\def%{\%}A, $T\omega_b=2000$}}%
\end{pgfscope}%
\begin{pgfscope}%
\pgfsetbuttcap%
\pgfsetroundjoin%
\definecolor{currentfill}{rgb}{1.000000,0.000000,0.000000}%
\pgfsetfillcolor{currentfill}%
\pgfsetlinewidth{0.803000pt}%
\definecolor{currentstroke}{rgb}{1.000000,0.000000,0.000000}%
\pgfsetstrokecolor{currentstroke}%
\pgfsetdash{}{0pt}%
\pgfsys@defobject{currentmarker}{\pgfqpoint{-0.019444in}{-0.019444in}}{\pgfqpoint{0.019444in}{0.019444in}}{%
\pgfpathmoveto{\pgfqpoint{-0.019444in}{-0.019444in}}%
\pgfpathlineto{\pgfqpoint{0.019444in}{0.019444in}}%
\pgfpathmoveto{\pgfqpoint{-0.019444in}{0.019444in}}%
\pgfpathlineto{\pgfqpoint{0.019444in}{-0.019444in}}%
\pgfusepath{stroke,fill}%
}%
\begin{pgfscope}%
\pgfsys@transformshift{3.132282in}{0.769483in}%
\pgfsys@useobject{currentmarker}{}%
\end{pgfscope}%
\end{pgfscope}%
\begin{pgfscope}%
\definecolor{textcolor}{rgb}{0.000000,0.000000,0.000000}%
\pgfsetstrokecolor{textcolor}%
\pgfsetfillcolor{textcolor}%
\pgftext[x=3.258671in,y=0.737886in,left,base]{\color{textcolor}{\rmfamily\fontsize{6.500000}{7.800000}\selectfont\catcode`\^=\active\def^{\ifmmode\sp\else\^{}\fi}\catcode`\%=\active\def%{\%}B, $T\omega_b=2000$}}%
\end{pgfscope}%
\begin{pgfscope}%
\pgfsetbuttcap%
\pgfsetmiterjoin%
\definecolor{currentfill}{rgb}{0.000000,0.000000,0.000000}%
\pgfsetfillcolor{currentfill}%
\pgfsetfillopacity{0.000000}%
\pgfsetlinewidth{0.803000pt}%
\definecolor{currentstroke}{rgb}{0.172549,0.627451,0.172549}%
\pgfsetstrokecolor{currentstroke}%
\pgfsetdash{}{0pt}%
\pgfsys@defobject{currentmarker}{\pgfqpoint{-0.017500in}{-0.017500in}}{\pgfqpoint{0.017500in}{0.017500in}}{%
\pgfpathmoveto{\pgfqpoint{-0.017500in}{-0.017500in}}%
\pgfpathlineto{\pgfqpoint{0.017500in}{-0.017500in}}%
\pgfpathlineto{\pgfqpoint{0.017500in}{0.017500in}}%
\pgfpathlineto{\pgfqpoint{-0.017500in}{0.017500in}}%
\pgfpathlineto{\pgfqpoint{-0.017500in}{-0.017500in}}%
\pgfpathclose%
\pgfusepath{stroke,fill}%
}%
\begin{pgfscope}%
\pgfsys@transformshift{3.132282in}{0.668712in}%
\pgfsys@useobject{currentmarker}{}%
\end{pgfscope}%
\end{pgfscope}%
\begin{pgfscope}%
\definecolor{textcolor}{rgb}{0.000000,0.000000,0.000000}%
\pgfsetstrokecolor{textcolor}%
\pgfsetfillcolor{textcolor}%
\pgftext[x=3.258671in,y=0.637114in,left,base]{\color{textcolor}{\rmfamily\fontsize{6.500000}{7.800000}\selectfont\catcode`\^=\active\def^{\ifmmode\sp\else\^{}\fi}\catcode`\%=\active\def%{\%}A, $T\omega_b=8000$}}%
\end{pgfscope}%
\begin{pgfscope}%
\pgfsetbuttcap%
\pgfsetroundjoin%
\definecolor{currentfill}{rgb}{0.172549,0.627451,0.172549}%
\pgfsetfillcolor{currentfill}%
\pgfsetlinewidth{0.803000pt}%
\definecolor{currentstroke}{rgb}{0.172549,0.627451,0.172549}%
\pgfsetstrokecolor{currentstroke}%
\pgfsetdash{}{0pt}%
\pgfsys@defobject{currentmarker}{\pgfqpoint{-0.019444in}{-0.019444in}}{\pgfqpoint{0.019444in}{0.019444in}}{%
\pgfpathmoveto{\pgfqpoint{-0.019444in}{-0.019444in}}%
\pgfpathlineto{\pgfqpoint{0.019444in}{0.019444in}}%
\pgfpathmoveto{\pgfqpoint{-0.019444in}{0.019444in}}%
\pgfpathlineto{\pgfqpoint{0.019444in}{-0.019444in}}%
\pgfusepath{stroke,fill}%
}%
\begin{pgfscope}%
\pgfsys@transformshift{3.132282in}{0.567940in}%
\pgfsys@useobject{currentmarker}{}%
\end{pgfscope}%
\end{pgfscope}%
\begin{pgfscope}%
\definecolor{textcolor}{rgb}{0.000000,0.000000,0.000000}%
\pgfsetstrokecolor{textcolor}%
\pgfsetfillcolor{textcolor}%
\pgftext[x=3.258671in,y=0.536343in,left,base]{\color{textcolor}{\rmfamily\fontsize{6.500000}{7.800000}\selectfont\catcode`\^=\active\def^{\ifmmode\sp\else\^{}\fi}\catcode`\%=\active\def%{\%}B, $T\omega_b=8000$}}%
\end{pgfscope}%
\begin{pgfscope}%
\pgfsetbuttcap%
\pgfsetmiterjoin%
\definecolor{currentfill}{rgb}{1.000000,1.000000,1.000000}%
\pgfsetfillcolor{currentfill}%
\pgfsetlinewidth{0.000000pt}%
\definecolor{currentstroke}{rgb}{0.000000,0.000000,0.000000}%
\pgfsetstrokecolor{currentstroke}%
\pgfsetstrokeopacity{0.000000}%
\pgfsetdash{}{0pt}%
\pgfpathmoveto{\pgfqpoint{2.110853in}{0.462654in}}%
\pgfpathlineto{\pgfqpoint{2.182714in}{0.462654in}}%
\pgfpathlineto{\pgfqpoint{2.182714in}{1.899877in}}%
\pgfpathlineto{\pgfqpoint{2.110853in}{1.899877in}}%
\pgfpathlineto{\pgfqpoint{2.110853in}{0.462654in}}%
\pgfpathclose%
\pgfusepath{fill}%
\end{pgfscope}%
\begin{pgfscope}%
\pgfpathrectangle{\pgfqpoint{2.110853in}{0.462654in}}{\pgfqpoint{0.071861in}{1.437222in}}%
\pgfusepath{clip}%
\pgfsetbuttcap%
\pgfsetroundjoin%
\definecolor{currentfill}{rgb}{0.956863,0.427451,0.262745}%
\pgfsetfillcolor{currentfill}%
\pgfsetlinewidth{0.000000pt}%
\definecolor{currentstroke}{rgb}{0.000000,0.000000,0.000000}%
\pgfsetstrokecolor{currentstroke}%
\pgfsetdash{}{0pt}%
\pgfpathmoveto{\pgfqpoint{2.110853in}{0.462654in}}%
\pgfpathlineto{\pgfqpoint{2.182714in}{0.462654in}}%
\pgfpathlineto{\pgfqpoint{2.182714in}{0.667972in}}%
\pgfpathlineto{\pgfqpoint{2.110853in}{0.667972in}}%
\pgfpathlineto{\pgfqpoint{2.110853in}{0.462654in}}%
\pgfusepath{fill}%
\end{pgfscope}%
\begin{pgfscope}%
\pgfpathrectangle{\pgfqpoint{2.110853in}{0.462654in}}{\pgfqpoint{0.071861in}{1.437222in}}%
\pgfusepath{clip}%
\pgfsetbuttcap%
\pgfsetroundjoin%
\definecolor{currentfill}{rgb}{0.998078,0.940408,0.658670}%
\pgfsetfillcolor{currentfill}%
\pgfsetlinewidth{0.000000pt}%
\definecolor{currentstroke}{rgb}{0.000000,0.000000,0.000000}%
\pgfsetstrokecolor{currentstroke}%
\pgfsetdash{}{0pt}%
\pgfpathmoveto{\pgfqpoint{2.110853in}{0.667972in}}%
\pgfpathlineto{\pgfqpoint{2.182714in}{0.667972in}}%
\pgfpathlineto{\pgfqpoint{2.182714in}{0.873289in}}%
\pgfpathlineto{\pgfqpoint{2.110853in}{0.873289in}}%
\pgfpathlineto{\pgfqpoint{2.110853in}{0.667972in}}%
\pgfusepath{fill}%
\end{pgfscope}%
\begin{pgfscope}%
\pgfpathrectangle{\pgfqpoint{2.110853in}{0.462654in}}{\pgfqpoint{0.071861in}{1.437222in}}%
\pgfusepath{clip}%
\pgfsetbuttcap%
\pgfsetroundjoin%
\definecolor{currentfill}{rgb}{0.940408,0.976932,0.858593}%
\pgfsetfillcolor{currentfill}%
\pgfsetlinewidth{0.000000pt}%
\definecolor{currentstroke}{rgb}{0.000000,0.000000,0.000000}%
\pgfsetstrokecolor{currentstroke}%
\pgfsetdash{}{0pt}%
\pgfpathmoveto{\pgfqpoint{2.110853in}{0.873289in}}%
\pgfpathlineto{\pgfqpoint{2.182714in}{0.873289in}}%
\pgfpathlineto{\pgfqpoint{2.182714in}{1.078607in}}%
\pgfpathlineto{\pgfqpoint{2.110853in}{1.078607in}}%
\pgfpathlineto{\pgfqpoint{2.110853in}{0.873289in}}%
\pgfusepath{fill}%
\end{pgfscope}%
\begin{pgfscope}%
\pgfpathrectangle{\pgfqpoint{2.110853in}{0.462654in}}{\pgfqpoint{0.071861in}{1.437222in}}%
\pgfusepath{clip}%
\pgfsetbuttcap%
\pgfsetroundjoin%
\definecolor{currentfill}{rgb}{0.772472,0.900961,0.942561}%
\pgfsetfillcolor{currentfill}%
\pgfsetlinewidth{0.000000pt}%
\definecolor{currentstroke}{rgb}{0.000000,0.000000,0.000000}%
\pgfsetstrokecolor{currentstroke}%
\pgfsetdash{}{0pt}%
\pgfpathmoveto{\pgfqpoint{2.110853in}{1.078607in}}%
\pgfpathlineto{\pgfqpoint{2.182714in}{1.078607in}}%
\pgfpathlineto{\pgfqpoint{2.182714in}{1.283924in}}%
\pgfpathlineto{\pgfqpoint{2.110853in}{1.283924in}}%
\pgfpathlineto{\pgfqpoint{2.110853in}{1.078607in}}%
\pgfusepath{fill}%
\end{pgfscope}%
\begin{pgfscope}%
\pgfpathrectangle{\pgfqpoint{2.110853in}{0.462654in}}{\pgfqpoint{0.071861in}{1.437222in}}%
\pgfusepath{clip}%
\pgfsetbuttcap%
\pgfsetroundjoin%
\definecolor{currentfill}{rgb}{0.556401,0.759631,0.863899}%
\pgfsetfillcolor{currentfill}%
\pgfsetlinewidth{0.000000pt}%
\definecolor{currentstroke}{rgb}{0.000000,0.000000,0.000000}%
\pgfsetstrokecolor{currentstroke}%
\pgfsetdash{}{0pt}%
\pgfpathmoveto{\pgfqpoint{2.110853in}{1.283924in}}%
\pgfpathlineto{\pgfqpoint{2.182714in}{1.283924in}}%
\pgfpathlineto{\pgfqpoint{2.182714in}{1.489242in}}%
\pgfpathlineto{\pgfqpoint{2.110853in}{1.489242in}}%
\pgfpathlineto{\pgfqpoint{2.110853in}{1.283924in}}%
\pgfusepath{fill}%
\end{pgfscope}%
\begin{pgfscope}%
\pgfpathrectangle{\pgfqpoint{2.110853in}{0.462654in}}{\pgfqpoint{0.071861in}{1.437222in}}%
\pgfusepath{clip}%
\pgfsetbuttcap%
\pgfsetroundjoin%
\definecolor{currentfill}{rgb}{0.360938,0.566474,0.761630}%
\pgfsetfillcolor{currentfill}%
\pgfsetlinewidth{0.000000pt}%
\definecolor{currentstroke}{rgb}{0.000000,0.000000,0.000000}%
\pgfsetstrokecolor{currentstroke}%
\pgfsetdash{}{0pt}%
\pgfpathmoveto{\pgfqpoint{2.110853in}{1.489242in}}%
\pgfpathlineto{\pgfqpoint{2.182714in}{1.489242in}}%
\pgfpathlineto{\pgfqpoint{2.182714in}{1.694559in}}%
\pgfpathlineto{\pgfqpoint{2.110853in}{1.694559in}}%
\pgfpathlineto{\pgfqpoint{2.110853in}{1.489242in}}%
\pgfusepath{fill}%
\end{pgfscope}%
\begin{pgfscope}%
\pgfpathrectangle{\pgfqpoint{2.110853in}{0.462654in}}{\pgfqpoint{0.071861in}{1.437222in}}%
\pgfusepath{clip}%
\pgfsetbuttcap%
\pgfsetroundjoin%
\definecolor{currentfill}{rgb}{0.229066,0.328028,0.641522}%
\pgfsetfillcolor{currentfill}%
\pgfsetlinewidth{0.000000pt}%
\definecolor{currentstroke}{rgb}{0.000000,0.000000,0.000000}%
\pgfsetstrokecolor{currentstroke}%
\pgfsetdash{}{0pt}%
\pgfpathmoveto{\pgfqpoint{2.110853in}{1.694559in}}%
\pgfpathlineto{\pgfqpoint{2.182714in}{1.694559in}}%
\pgfpathlineto{\pgfqpoint{2.182714in}{1.899877in}}%
\pgfpathlineto{\pgfqpoint{2.110853in}{1.899877in}}%
\pgfpathlineto{\pgfqpoint{2.110853in}{1.694559in}}%
\pgfusepath{fill}%
\end{pgfscope}%
\begin{pgfscope}%
\pgfsetbuttcap%
\pgfsetroundjoin%
\definecolor{currentfill}{rgb}{0.000000,0.000000,0.000000}%
\pgfsetfillcolor{currentfill}%
\pgfsetlinewidth{0.803000pt}%
\definecolor{currentstroke}{rgb}{0.000000,0.000000,0.000000}%
\pgfsetstrokecolor{currentstroke}%
\pgfsetdash{}{0pt}%
\pgfsys@defobject{currentmarker}{\pgfqpoint{0.000000in}{0.000000in}}{\pgfqpoint{0.048611in}{0.000000in}}{%
\pgfpathmoveto{\pgfqpoint{0.000000in}{0.000000in}}%
\pgfpathlineto{\pgfqpoint{0.048611in}{0.000000in}}%
\pgfusepath{stroke,fill}%
}%
\begin{pgfscope}%
\pgfsys@transformshift{2.182714in}{0.462654in}%
\pgfsys@useobject{currentmarker}{}%
\end{pgfscope}%
\end{pgfscope}%
\begin{pgfscope}%
\definecolor{textcolor}{rgb}{0.000000,0.000000,0.000000}%
\pgfsetstrokecolor{textcolor}%
\pgfsetfillcolor{textcolor}%
\pgftext[x=2.279936in, y=0.428897in, left, base]{\color{textcolor}{\rmfamily\fontsize{7.000000}{8.400000}\selectfont\catcode`\^=\active\def^{\ifmmode\sp\else\^{}\fi}\catcode`\%=\active\def%{\%}0.0}}%
\end{pgfscope}%
\begin{pgfscope}%
\pgfsetbuttcap%
\pgfsetroundjoin%
\definecolor{currentfill}{rgb}{0.000000,0.000000,0.000000}%
\pgfsetfillcolor{currentfill}%
\pgfsetlinewidth{0.803000pt}%
\definecolor{currentstroke}{rgb}{0.000000,0.000000,0.000000}%
\pgfsetstrokecolor{currentstroke}%
\pgfsetdash{}{0pt}%
\pgfsys@defobject{currentmarker}{\pgfqpoint{0.000000in}{0.000000in}}{\pgfqpoint{0.048611in}{0.000000in}}{%
\pgfpathmoveto{\pgfqpoint{0.000000in}{0.000000in}}%
\pgfpathlineto{\pgfqpoint{0.048611in}{0.000000in}}%
\pgfusepath{stroke,fill}%
}%
\begin{pgfscope}%
\pgfsys@transformshift{2.182714in}{0.667972in}%
\pgfsys@useobject{currentmarker}{}%
\end{pgfscope}%
\end{pgfscope}%
\begin{pgfscope}%
\definecolor{textcolor}{rgb}{0.000000,0.000000,0.000000}%
\pgfsetstrokecolor{textcolor}%
\pgfsetfillcolor{textcolor}%
\pgftext[x=2.279936in, y=0.634214in, left, base]{\color{textcolor}{\rmfamily\fontsize{7.000000}{8.400000}\selectfont\catcode`\^=\active\def^{\ifmmode\sp\else\^{}\fi}\catcode`\%=\active\def%{\%}0.4}}%
\end{pgfscope}%
\begin{pgfscope}%
\pgfsetbuttcap%
\pgfsetroundjoin%
\definecolor{currentfill}{rgb}{0.000000,0.000000,0.000000}%
\pgfsetfillcolor{currentfill}%
\pgfsetlinewidth{0.803000pt}%
\definecolor{currentstroke}{rgb}{0.000000,0.000000,0.000000}%
\pgfsetstrokecolor{currentstroke}%
\pgfsetdash{}{0pt}%
\pgfsys@defobject{currentmarker}{\pgfqpoint{0.000000in}{0.000000in}}{\pgfqpoint{0.048611in}{0.000000in}}{%
\pgfpathmoveto{\pgfqpoint{0.000000in}{0.000000in}}%
\pgfpathlineto{\pgfqpoint{0.048611in}{0.000000in}}%
\pgfusepath{stroke,fill}%
}%
\begin{pgfscope}%
\pgfsys@transformshift{2.182714in}{0.873289in}%
\pgfsys@useobject{currentmarker}{}%
\end{pgfscope}%
\end{pgfscope}%
\begin{pgfscope}%
\definecolor{textcolor}{rgb}{0.000000,0.000000,0.000000}%
\pgfsetstrokecolor{textcolor}%
\pgfsetfillcolor{textcolor}%
\pgftext[x=2.279936in, y=0.839532in, left, base]{\color{textcolor}{\rmfamily\fontsize{7.000000}{8.400000}\selectfont\catcode`\^=\active\def^{\ifmmode\sp\else\^{}\fi}\catcode`\%=\active\def%{\%}0.5}}%
\end{pgfscope}%
\begin{pgfscope}%
\pgfsetbuttcap%
\pgfsetroundjoin%
\definecolor{currentfill}{rgb}{0.000000,0.000000,0.000000}%
\pgfsetfillcolor{currentfill}%
\pgfsetlinewidth{0.803000pt}%
\definecolor{currentstroke}{rgb}{0.000000,0.000000,0.000000}%
\pgfsetstrokecolor{currentstroke}%
\pgfsetdash{}{0pt}%
\pgfsys@defobject{currentmarker}{\pgfqpoint{0.000000in}{0.000000in}}{\pgfqpoint{0.048611in}{0.000000in}}{%
\pgfpathmoveto{\pgfqpoint{0.000000in}{0.000000in}}%
\pgfpathlineto{\pgfqpoint{0.048611in}{0.000000in}}%
\pgfusepath{stroke,fill}%
}%
\begin{pgfscope}%
\pgfsys@transformshift{2.182714in}{1.078607in}%
\pgfsys@useobject{currentmarker}{}%
\end{pgfscope}%
\end{pgfscope}%
\begin{pgfscope}%
\definecolor{textcolor}{rgb}{0.000000,0.000000,0.000000}%
\pgfsetstrokecolor{textcolor}%
\pgfsetfillcolor{textcolor}%
\pgftext[x=2.279936in, y=1.044849in, left, base]{\color{textcolor}{\rmfamily\fontsize{7.000000}{8.400000}\selectfont\catcode`\^=\active\def^{\ifmmode\sp\else\^{}\fi}\catcode`\%=\active\def%{\%}0.6}}%
\end{pgfscope}%
\begin{pgfscope}%
\pgfsetbuttcap%
\pgfsetroundjoin%
\definecolor{currentfill}{rgb}{0.000000,0.000000,0.000000}%
\pgfsetfillcolor{currentfill}%
\pgfsetlinewidth{0.803000pt}%
\definecolor{currentstroke}{rgb}{0.000000,0.000000,0.000000}%
\pgfsetstrokecolor{currentstroke}%
\pgfsetdash{}{0pt}%
\pgfsys@defobject{currentmarker}{\pgfqpoint{0.000000in}{0.000000in}}{\pgfqpoint{0.048611in}{0.000000in}}{%
\pgfpathmoveto{\pgfqpoint{0.000000in}{0.000000in}}%
\pgfpathlineto{\pgfqpoint{0.048611in}{0.000000in}}%
\pgfusepath{stroke,fill}%
}%
\begin{pgfscope}%
\pgfsys@transformshift{2.182714in}{1.283924in}%
\pgfsys@useobject{currentmarker}{}%
\end{pgfscope}%
\end{pgfscope}%
\begin{pgfscope}%
\definecolor{textcolor}{rgb}{0.000000,0.000000,0.000000}%
\pgfsetstrokecolor{textcolor}%
\pgfsetfillcolor{textcolor}%
\pgftext[x=2.279936in, y=1.250167in, left, base]{\color{textcolor}{\rmfamily\fontsize{7.000000}{8.400000}\selectfont\catcode`\^=\active\def^{\ifmmode\sp\else\^{}\fi}\catcode`\%=\active\def%{\%}0.7}}%
\end{pgfscope}%
\begin{pgfscope}%
\pgfsetbuttcap%
\pgfsetroundjoin%
\definecolor{currentfill}{rgb}{0.000000,0.000000,0.000000}%
\pgfsetfillcolor{currentfill}%
\pgfsetlinewidth{0.803000pt}%
\definecolor{currentstroke}{rgb}{0.000000,0.000000,0.000000}%
\pgfsetstrokecolor{currentstroke}%
\pgfsetdash{}{0pt}%
\pgfsys@defobject{currentmarker}{\pgfqpoint{0.000000in}{0.000000in}}{\pgfqpoint{0.048611in}{0.000000in}}{%
\pgfpathmoveto{\pgfqpoint{0.000000in}{0.000000in}}%
\pgfpathlineto{\pgfqpoint{0.048611in}{0.000000in}}%
\pgfusepath{stroke,fill}%
}%
\begin{pgfscope}%
\pgfsys@transformshift{2.182714in}{1.489242in}%
\pgfsys@useobject{currentmarker}{}%
\end{pgfscope}%
\end{pgfscope}%
\begin{pgfscope}%
\definecolor{textcolor}{rgb}{0.000000,0.000000,0.000000}%
\pgfsetstrokecolor{textcolor}%
\pgfsetfillcolor{textcolor}%
\pgftext[x=2.279936in, y=1.455484in, left, base]{\color{textcolor}{\rmfamily\fontsize{7.000000}{8.400000}\selectfont\catcode`\^=\active\def^{\ifmmode\sp\else\^{}\fi}\catcode`\%=\active\def%{\%}0.8}}%
\end{pgfscope}%
\begin{pgfscope}%
\pgfsetbuttcap%
\pgfsetroundjoin%
\definecolor{currentfill}{rgb}{0.000000,0.000000,0.000000}%
\pgfsetfillcolor{currentfill}%
\pgfsetlinewidth{0.803000pt}%
\definecolor{currentstroke}{rgb}{0.000000,0.000000,0.000000}%
\pgfsetstrokecolor{currentstroke}%
\pgfsetdash{}{0pt}%
\pgfsys@defobject{currentmarker}{\pgfqpoint{0.000000in}{0.000000in}}{\pgfqpoint{0.048611in}{0.000000in}}{%
\pgfpathmoveto{\pgfqpoint{0.000000in}{0.000000in}}%
\pgfpathlineto{\pgfqpoint{0.048611in}{0.000000in}}%
\pgfusepath{stroke,fill}%
}%
\begin{pgfscope}%
\pgfsys@transformshift{2.182714in}{1.694559in}%
\pgfsys@useobject{currentmarker}{}%
\end{pgfscope}%
\end{pgfscope}%
\begin{pgfscope}%
\definecolor{textcolor}{rgb}{0.000000,0.000000,0.000000}%
\pgfsetstrokecolor{textcolor}%
\pgfsetfillcolor{textcolor}%
\pgftext[x=2.279936in, y=1.660802in, left, base]{\color{textcolor}{\rmfamily\fontsize{7.000000}{8.400000}\selectfont\catcode`\^=\active\def^{\ifmmode\sp\else\^{}\fi}\catcode`\%=\active\def%{\%}0.9}}%
\end{pgfscope}%
\begin{pgfscope}%
\pgfsetbuttcap%
\pgfsetroundjoin%
\definecolor{currentfill}{rgb}{0.000000,0.000000,0.000000}%
\pgfsetfillcolor{currentfill}%
\pgfsetlinewidth{0.803000pt}%
\definecolor{currentstroke}{rgb}{0.000000,0.000000,0.000000}%
\pgfsetstrokecolor{currentstroke}%
\pgfsetdash{}{0pt}%
\pgfsys@defobject{currentmarker}{\pgfqpoint{0.000000in}{0.000000in}}{\pgfqpoint{0.048611in}{0.000000in}}{%
\pgfpathmoveto{\pgfqpoint{0.000000in}{0.000000in}}%
\pgfpathlineto{\pgfqpoint{0.048611in}{0.000000in}}%
\pgfusepath{stroke,fill}%
}%
\begin{pgfscope}%
\pgfsys@transformshift{2.182714in}{1.899877in}%
\pgfsys@useobject{currentmarker}{}%
\end{pgfscope}%
\end{pgfscope}%
\begin{pgfscope}%
\definecolor{textcolor}{rgb}{0.000000,0.000000,0.000000}%
\pgfsetstrokecolor{textcolor}%
\pgfsetfillcolor{textcolor}%
\pgftext[x=2.279936in, y=1.866119in, left, base]{\color{textcolor}{\rmfamily\fontsize{7.000000}{8.400000}\selectfont\catcode`\^=\active\def^{\ifmmode\sp\else\^{}\fi}\catcode`\%=\active\def%{\%}1.0}}%
\end{pgfscope}%
\begin{pgfscope}%
\pgfsetrectcap%
\pgfsetmiterjoin%
\pgfsetlinewidth{0.803000pt}%
\definecolor{currentstroke}{rgb}{0.000000,0.000000,0.000000}%
\pgfsetstrokecolor{currentstroke}%
\pgfsetdash{}{0pt}%
\pgfpathmoveto{\pgfqpoint{2.110853in}{0.462654in}}%
\pgfpathlineto{\pgfqpoint{2.146783in}{0.462654in}}%
\pgfpathlineto{\pgfqpoint{2.182714in}{0.462654in}}%
\pgfpathlineto{\pgfqpoint{2.182714in}{1.899877in}}%
\pgfpathlineto{\pgfqpoint{2.146783in}{1.899877in}}%
\pgfpathlineto{\pgfqpoint{2.110853in}{1.899877in}}%
\pgfpathlineto{\pgfqpoint{2.110853in}{0.462654in}}%
\pgfpathclose%
\pgfusepath{stroke}%
\end{pgfscope}%
\begin{pgfscope}%
\definecolor{textcolor}{rgb}{0.000000,0.000000,0.000000}%
\pgfsetstrokecolor{textcolor}%
\pgfsetfillcolor{textcolor}%
\pgftext[x=2.605292in, y=1.106541in, left, base,rotate=90.000000]{\color{textcolor}{\rmfamily\fontsize{9.000000}{10.800000}\selectfont\catcode`\^=\active\def^{\ifmmode\sp\else\^{}\fi}\catcode`\%=\active\def%{\%}$F_h$}}%
\end{pgfscope}%
\end{pgfpicture}%
\makeatother%
\endgroup%
}
\caption{(En color en l\'inea) (a) Fidelidad condicional $F_h$ en funci\'on del factor de correlaci\'on del ruido $\chi$ y de la espera $T\omega_b$; los contornos son $F_h=0.5$ y $0.8$, interpolados entre los puntos calculados. (b) $F_h$ en funci\'on de la desviaci\'on $\delta$ de las respuestas Stark, con operador de ruido $\sum_i a_i n_i$ y $\chi=1$: patr\'on A, $a=(1+\delta,1,1)$ (cuadrados), y patr\'on B, $a=(1,1+\delta,1-\delta)$ (cruces), para $T\omega_b=2000$ (rojo) y $8000$ (verde); la l\'inea negra es la referencia pasiva. Las l\'ineas en (b) son interpolaciones mon\'otonas. Par\'ametros del sistema: modelo reducido, los mismos que en la Fig.~\ref{fig:sweep}; tolerancia de integraci\'on $10^{-9}$.}
\label{fig:requirements}
\end{figure}

Cuando una rama ya contiene un ex\'citon, el siguiente pulso puede alcanzar la variedad de dos excitones, que el modelo reducido omite. La Figura~\ref{fig:biexciton} muestra el efecto sobre el protocolo completo con puertos reales para valores seleccionados del desplazamiento de dos excitones entre puntos $\Delta_{xx}$. Para desplazamientos grandes la fidelidad condicional se aproxima al valor del modelo reducido, mientras que la probabilidad de recuperaci\'on disminuye antes a medida que la variedad adicional se hace accesible: el pulso que debe escribir el segundo intervalo temporal empieza a excitar estados de dos excitones de la otra rama y retira poblaci\'on de la amplitud almacenada. El biexcit\'on reduce por tanto la probabilidad de \'exito antes de reducir la calidad del qubit recuperado.

El modelo reducido es v\'alido para desplazamientos grandes, y la separaci\'on requerida depende de los par\'ametros del pulso adem\'as de $\Delta_{xx}$. El modelo completo con un desplazamiento grande tambi\'en reproduce la dependencia con la espera de la Fig.~\ref{fig:sweep} (comprobado a tres esperas), y proyectado sobre el subespacio de un ex\'citon reproduce los resultados reducidos.

\begin{figure}[h]
\centering
\resizebox{0.6\columnwidth}{!}{%% Creator: Matplotlib, PGF backend
%%
%% To include the figure in your LaTeX document, write
%%   \input{<filename>.pgf}
%%
%% Make sure the required packages are loaded in your preamble
%%   \usepackage{pgf}
%%
%% Also ensure that all the required font packages are loaded; for instance,
%% the lmodern package is sometimes necessary when using math font.
%%   \usepackage{lmodern}
%%
%% Figures using additional raster images can only be included by \input if
%% they are in the same directory as the main LaTeX file. For loading figures
%% from other directories you can use the `import` package
%%   \usepackage{import}
%%
%% and then include the figures with
%%   \import{<path to file>}{<filename>.pgf}
%%
%% Matplotlib used the following preamble
%%   \def\mathdefault#1{#1}
%%   \everymath=\expandafter{\the\everymath\displaystyle}
%%   \IfFileExists{scrextend.sty}{
%%     \usepackage[fontsize=8.000000pt]{scrextend}
%%   }{
%%     \renewcommand{\normalsize}{\fontsize{8.000000}{9.600000}\selectfont}
%%     \normalsize
%%   }
%%   \usepackage{amsmath,amssymb}
%%   \makeatletter\@ifpackageloaded{underscore}{}{\usepackage[strings]{underscore}}\makeatother
%%
\begingroup%
\makeatletter%
\begin{pgfpicture}%
\pgfpathrectangle{\pgfpointorigin}{\pgfqpoint{2.976702in}{1.989099in}}%
\pgfusepath{use as bounding box, clip}%
\begin{pgfscope}%
\pgfsetbuttcap%
\pgfsetmiterjoin%
\definecolor{currentfill}{rgb}{1.000000,1.000000,1.000000}%
\pgfsetfillcolor{currentfill}%
\pgfsetlinewidth{0.000000pt}%
\definecolor{currentstroke}{rgb}{1.000000,1.000000,1.000000}%
\pgfsetstrokecolor{currentstroke}%
\pgfsetdash{}{0pt}%
\pgfpathmoveto{\pgfqpoint{0.000000in}{0.000000in}}%
\pgfpathlineto{\pgfqpoint{2.976702in}{0.000000in}}%
\pgfpathlineto{\pgfqpoint{2.976702in}{1.989099in}}%
\pgfpathlineto{\pgfqpoint{0.000000in}{1.989099in}}%
\pgfpathlineto{\pgfqpoint{0.000000in}{0.000000in}}%
\pgfpathclose%
\pgfusepath{fill}%
\end{pgfscope}%
\begin{pgfscope}%
\pgfsetbuttcap%
\pgfsetmiterjoin%
\definecolor{currentfill}{rgb}{1.000000,1.000000,1.000000}%
\pgfsetfillcolor{currentfill}%
\pgfsetlinewidth{0.000000pt}%
\definecolor{currentstroke}{rgb}{0.000000,0.000000,0.000000}%
\pgfsetstrokecolor{currentstroke}%
\pgfsetstrokeopacity{0.000000}%
\pgfsetdash{}{0pt}%
\pgfpathmoveto{\pgfqpoint{0.514740in}{0.476543in}}%
\pgfpathlineto{\pgfqpoint{2.876702in}{0.476543in}}%
\pgfpathlineto{\pgfqpoint{2.876702in}{1.889099in}}%
\pgfpathlineto{\pgfqpoint{0.514740in}{1.889099in}}%
\pgfpathlineto{\pgfqpoint{0.514740in}{0.476543in}}%
\pgfpathclose%
\pgfusepath{fill}%
\end{pgfscope}%
\begin{pgfscope}%
\pgfsetbuttcap%
\pgfsetroundjoin%
\definecolor{currentfill}{rgb}{0.000000,0.000000,0.000000}%
\pgfsetfillcolor{currentfill}%
\pgfsetlinewidth{0.803000pt}%
\definecolor{currentstroke}{rgb}{0.000000,0.000000,0.000000}%
\pgfsetstrokecolor{currentstroke}%
\pgfsetdash{}{0pt}%
\pgfsys@defobject{currentmarker}{\pgfqpoint{0.000000in}{-0.048611in}}{\pgfqpoint{0.000000in}{0.000000in}}{%
\pgfpathmoveto{\pgfqpoint{0.000000in}{0.000000in}}%
\pgfpathlineto{\pgfqpoint{0.000000in}{-0.048611in}}%
\pgfusepath{stroke,fill}%
}%
\begin{pgfscope}%
\pgfsys@transformshift{0.892654in}{0.476543in}%
\pgfsys@useobject{currentmarker}{}%
\end{pgfscope}%
\end{pgfscope}%
\begin{pgfscope}%
\definecolor{textcolor}{rgb}{0.000000,0.000000,0.000000}%
\pgfsetstrokecolor{textcolor}%
\pgfsetfillcolor{textcolor}%
\pgftext[x=0.892654in,y=0.379321in,,top]{\color{textcolor}{\rmfamily\fontsize{8.000000}{9.600000}\selectfont\catcode`\^=\active\def^{\ifmmode\sp\else\^{}\fi}\catcode`\%=\active\def%{\%}$-12$}}%
\end{pgfscope}%
\begin{pgfscope}%
\pgfsetbuttcap%
\pgfsetroundjoin%
\definecolor{currentfill}{rgb}{0.000000,0.000000,0.000000}%
\pgfsetfillcolor{currentfill}%
\pgfsetlinewidth{0.803000pt}%
\definecolor{currentstroke}{rgb}{0.000000,0.000000,0.000000}%
\pgfsetstrokecolor{currentstroke}%
\pgfsetdash{}{0pt}%
\pgfsys@defobject{currentmarker}{\pgfqpoint{0.000000in}{-0.048611in}}{\pgfqpoint{0.000000in}{0.000000in}}{%
\pgfpathmoveto{\pgfqpoint{0.000000in}{0.000000in}}%
\pgfpathlineto{\pgfqpoint{0.000000in}{-0.048611in}}%
\pgfusepath{stroke,fill}%
}%
\begin{pgfscope}%
\pgfsys@transformshift{1.648482in}{0.476543in}%
\pgfsys@useobject{currentmarker}{}%
\end{pgfscope}%
\end{pgfscope}%
\begin{pgfscope}%
\definecolor{textcolor}{rgb}{0.000000,0.000000,0.000000}%
\pgfsetstrokecolor{textcolor}%
\pgfsetfillcolor{textcolor}%
\pgftext[x=1.648482in,y=0.379321in,,top]{\color{textcolor}{\rmfamily\fontsize{8.000000}{9.600000}\selectfont\catcode`\^=\active\def^{\ifmmode\sp\else\^{}\fi}\catcode`\%=\active\def%{\%}$-8$}}%
\end{pgfscope}%
\begin{pgfscope}%
\pgfsetbuttcap%
\pgfsetroundjoin%
\definecolor{currentfill}{rgb}{0.000000,0.000000,0.000000}%
\pgfsetfillcolor{currentfill}%
\pgfsetlinewidth{0.803000pt}%
\definecolor{currentstroke}{rgb}{0.000000,0.000000,0.000000}%
\pgfsetstrokecolor{currentstroke}%
\pgfsetdash{}{0pt}%
\pgfsys@defobject{currentmarker}{\pgfqpoint{0.000000in}{-0.048611in}}{\pgfqpoint{0.000000in}{0.000000in}}{%
\pgfpathmoveto{\pgfqpoint{0.000000in}{0.000000in}}%
\pgfpathlineto{\pgfqpoint{0.000000in}{-0.048611in}}%
\pgfusepath{stroke,fill}%
}%
\begin{pgfscope}%
\pgfsys@transformshift{2.404309in}{0.476543in}%
\pgfsys@useobject{currentmarker}{}%
\end{pgfscope}%
\end{pgfscope}%
\begin{pgfscope}%
\definecolor{textcolor}{rgb}{0.000000,0.000000,0.000000}%
\pgfsetstrokecolor{textcolor}%
\pgfsetfillcolor{textcolor}%
\pgftext[x=2.404309in,y=0.379321in,,top]{\color{textcolor}{\rmfamily\fontsize{8.000000}{9.600000}\selectfont\catcode`\^=\active\def^{\ifmmode\sp\else\^{}\fi}\catcode`\%=\active\def%{\%}$-4$}}%
\end{pgfscope}%
\begin{pgfscope}%
\definecolor{textcolor}{rgb}{0.000000,0.000000,0.000000}%
\pgfsetstrokecolor{textcolor}%
\pgfsetfillcolor{textcolor}%
\pgftext[x=1.695721in,y=0.225000in,,top]{\color{textcolor}{\rmfamily\fontsize{9.000000}{10.800000}\selectfont\catcode`\^=\active\def^{\ifmmode\sp\else\^{}\fi}\catcode`\%=\active\def%{\%}$\Delta_{xx}/\omega_b$}}%
\end{pgfscope}%
\begin{pgfscope}%
\pgfsetbuttcap%
\pgfsetroundjoin%
\definecolor{currentfill}{rgb}{0.000000,0.000000,0.000000}%
\pgfsetfillcolor{currentfill}%
\pgfsetlinewidth{0.803000pt}%
\definecolor{currentstroke}{rgb}{0.000000,0.000000,0.000000}%
\pgfsetstrokecolor{currentstroke}%
\pgfsetdash{}{0pt}%
\pgfsys@defobject{currentmarker}{\pgfqpoint{-0.048611in}{0.000000in}}{\pgfqpoint{-0.000000in}{0.000000in}}{%
\pgfpathmoveto{\pgfqpoint{-0.000000in}{0.000000in}}%
\pgfpathlineto{\pgfqpoint{-0.048611in}{0.000000in}}%
\pgfusepath{stroke,fill}%
}%
\begin{pgfscope}%
\pgfsys@transformshift{0.514740in}{0.476543in}%
\pgfsys@useobject{currentmarker}{}%
\end{pgfscope}%
\end{pgfscope}%
\begin{pgfscope}%
\definecolor{textcolor}{rgb}{0.000000,0.000000,0.000000}%
\pgfsetstrokecolor{textcolor}%
\pgfsetfillcolor{textcolor}%
\pgftext[x=0.358489in, y=0.437963in, left, base]{\color{textcolor}{\rmfamily\fontsize{8.000000}{9.600000}\selectfont\catcode`\^=\active\def^{\ifmmode\sp\else\^{}\fi}\catcode`\%=\active\def%{\%}$0$}}%
\end{pgfscope}%
\begin{pgfscope}%
\pgfsetbuttcap%
\pgfsetroundjoin%
\definecolor{currentfill}{rgb}{0.000000,0.000000,0.000000}%
\pgfsetfillcolor{currentfill}%
\pgfsetlinewidth{0.803000pt}%
\definecolor{currentstroke}{rgb}{0.000000,0.000000,0.000000}%
\pgfsetstrokecolor{currentstroke}%
\pgfsetdash{}{0pt}%
\pgfsys@defobject{currentmarker}{\pgfqpoint{-0.048611in}{0.000000in}}{\pgfqpoint{-0.000000in}{0.000000in}}{%
\pgfpathmoveto{\pgfqpoint{-0.000000in}{0.000000in}}%
\pgfpathlineto{\pgfqpoint{-0.048611in}{0.000000in}}%
\pgfusepath{stroke,fill}%
}%
\begin{pgfscope}%
\pgfsys@transformshift{0.514740in}{1.118614in}%
\pgfsys@useobject{currentmarker}{}%
\end{pgfscope}%
\end{pgfscope}%
\begin{pgfscope}%
\definecolor{textcolor}{rgb}{0.000000,0.000000,0.000000}%
\pgfsetstrokecolor{textcolor}%
\pgfsetfillcolor{textcolor}%
\pgftext[x=0.266667in, y=1.080034in, left, base]{\color{textcolor}{\rmfamily\fontsize{8.000000}{9.600000}\selectfont\catcode`\^=\active\def^{\ifmmode\sp\else\^{}\fi}\catcode`\%=\active\def%{\%}$0.5$}}%
\end{pgfscope}%
\begin{pgfscope}%
\pgfsetbuttcap%
\pgfsetroundjoin%
\definecolor{currentfill}{rgb}{0.000000,0.000000,0.000000}%
\pgfsetfillcolor{currentfill}%
\pgfsetlinewidth{0.803000pt}%
\definecolor{currentstroke}{rgb}{0.000000,0.000000,0.000000}%
\pgfsetstrokecolor{currentstroke}%
\pgfsetdash{}{0pt}%
\pgfsys@defobject{currentmarker}{\pgfqpoint{-0.048611in}{0.000000in}}{\pgfqpoint{-0.000000in}{0.000000in}}{%
\pgfpathmoveto{\pgfqpoint{-0.000000in}{0.000000in}}%
\pgfpathlineto{\pgfqpoint{-0.048611in}{0.000000in}}%
\pgfusepath{stroke,fill}%
}%
\begin{pgfscope}%
\pgfsys@transformshift{0.514740in}{1.760685in}%
\pgfsys@useobject{currentmarker}{}%
\end{pgfscope}%
\end{pgfscope}%
\begin{pgfscope}%
\definecolor{textcolor}{rgb}{0.000000,0.000000,0.000000}%
\pgfsetstrokecolor{textcolor}%
\pgfsetfillcolor{textcolor}%
\pgftext[x=0.358489in, y=1.722104in, left, base]{\color{textcolor}{\rmfamily\fontsize{8.000000}{9.600000}\selectfont\catcode`\^=\active\def^{\ifmmode\sp\else\^{}\fi}\catcode`\%=\active\def%{\%}$1$}}%
\end{pgfscope}%
\begin{pgfscope}%
\definecolor{textcolor}{rgb}{0.000000,0.000000,0.000000}%
\pgfsetstrokecolor{textcolor}%
\pgfsetfillcolor{textcolor}%
\pgftext[x=0.211111in,y=1.182821in,,bottom,rotate=90.000000]{\color{textcolor}{\rmfamily\fontsize{9.000000}{10.800000}\selectfont\catcode`\^=\active\def^{\ifmmode\sp\else\^{}\fi}\catcode`\%=\active\def%{\%}$s,F_h$}}%
\end{pgfscope}%
\begin{pgfscope}%
\pgfpathrectangle{\pgfqpoint{0.514740in}{0.476543in}}{\pgfqpoint{2.361962in}{1.412556in}}%
\pgfusepath{clip}%
\pgfsetrectcap%
\pgfsetroundjoin%
\pgfsetlinewidth{0.903375pt}%
\definecolor{currentstroke}{rgb}{0.121569,0.466667,0.705882}%
\pgfsetstrokecolor{currentstroke}%
\pgfsetdash{}{0pt}%
\pgfpathmoveto{\pgfqpoint{0.892654in}{1.386080in}}%
\pgfpathlineto{\pgfqpoint{1.139532in}{1.382387in}}%
\pgfpathlineto{\pgfqpoint{1.471869in}{1.375023in}}%
\pgfpathlineto{\pgfqpoint{1.652280in}{1.369444in}}%
\pgfpathlineto{\pgfqpoint{1.756728in}{1.364111in}}%
\pgfpathlineto{\pgfqpoint{1.899158in}{1.354411in}}%
\pgfpathlineto{\pgfqpoint{1.994111in}{1.346034in}}%
\pgfpathlineto{\pgfqpoint{2.070074in}{1.337244in}}%
\pgfpathlineto{\pgfqpoint{2.136541in}{1.327251in}}%
\pgfpathlineto{\pgfqpoint{2.212504in}{1.313537in}}%
\pgfpathlineto{\pgfqpoint{2.316952in}{1.292320in}}%
\pgfpathlineto{\pgfqpoint{2.354934in}{1.281975in}}%
\pgfpathlineto{\pgfqpoint{2.383420in}{1.272330in}}%
\pgfpathlineto{\pgfqpoint{2.411906in}{1.260320in}}%
\pgfpathlineto{\pgfqpoint{2.430896in}{1.248098in}}%
\pgfpathlineto{\pgfqpoint{2.449887in}{1.231805in}}%
\pgfpathlineto{\pgfqpoint{2.468878in}{1.211921in}}%
\pgfpathlineto{\pgfqpoint{2.487868in}{1.188927in}}%
\pgfpathlineto{\pgfqpoint{2.516354in}{1.149657in}}%
\pgfpathlineto{\pgfqpoint{2.544840in}{1.106095in}}%
\pgfpathlineto{\pgfqpoint{2.592317in}{1.028366in}}%
\pgfpathlineto{\pgfqpoint{2.611307in}{0.995115in}}%
\pgfpathlineto{\pgfqpoint{2.630298in}{0.958288in}}%
\pgfpathlineto{\pgfqpoint{2.658784in}{0.896826in}}%
\pgfpathlineto{\pgfqpoint{2.687270in}{0.828573in}}%
\pgfpathlineto{\pgfqpoint{2.715756in}{0.754290in}}%
\pgfpathlineto{\pgfqpoint{2.753737in}{0.647177in}}%
\pgfpathlineto{\pgfqpoint{2.782223in}{0.561777in}}%
\pgfpathlineto{\pgfqpoint{2.782223in}{0.561777in}}%
\pgfusepath{stroke}%
\end{pgfscope}%
\begin{pgfscope}%
\pgfpathrectangle{\pgfqpoint{0.514740in}{0.476543in}}{\pgfqpoint{2.361962in}{1.412556in}}%
\pgfusepath{clip}%
\pgfsetbuttcap%
\pgfsetmiterjoin%
\definecolor{currentfill}{rgb}{0.000000,0.000000,0.000000}%
\pgfsetfillcolor{currentfill}%
\pgfsetfillopacity{0.000000}%
\pgfsetlinewidth{0.803000pt}%
\definecolor{currentstroke}{rgb}{0.121569,0.466667,0.705882}%
\pgfsetstrokecolor{currentstroke}%
\pgfsetdash{}{0pt}%
\pgfsys@defobject{currentmarker}{\pgfqpoint{-0.025000in}{-0.025000in}}{\pgfqpoint{0.025000in}{0.025000in}}{%
\pgfpathmoveto{\pgfqpoint{-0.025000in}{-0.025000in}}%
\pgfpathlineto{\pgfqpoint{0.025000in}{-0.025000in}}%
\pgfpathlineto{\pgfqpoint{0.025000in}{0.025000in}}%
\pgfpathlineto{\pgfqpoint{-0.025000in}{0.025000in}}%
\pgfpathlineto{\pgfqpoint{-0.025000in}{-0.025000in}}%
\pgfpathclose%
\pgfusepath{stroke,fill}%
}%
\begin{pgfscope}%
\pgfsys@transformshift{0.892654in}{1.386080in}%
\pgfsys@useobject{currentmarker}{}%
\end{pgfscope}%
\begin{pgfscope}%
\pgfsys@transformshift{1.459525in}{1.375325in}%
\pgfsys@useobject{currentmarker}{}%
\end{pgfscope}%
\begin{pgfscope}%
\pgfsys@transformshift{1.648482in}{1.369596in}%
\pgfsys@useobject{currentmarker}{}%
\end{pgfscope}%
\begin{pgfscope}%
\pgfsys@transformshift{1.837438in}{1.358889in}%
\pgfsys@useobject{currentmarker}{}%
\end{pgfscope}%
\begin{pgfscope}%
\pgfsys@transformshift{2.026395in}{1.342616in}%
\pgfsys@useobject{currentmarker}{}%
\end{pgfscope}%
\begin{pgfscope}%
\pgfsys@transformshift{2.215352in}{1.312982in}%
\pgfsys@useobject{currentmarker}{}%
\end{pgfscope}%
\begin{pgfscope}%
\pgfsys@transformshift{2.404309in}{1.263961in}%
\pgfsys@useobject{currentmarker}{}%
\end{pgfscope}%
\begin{pgfscope}%
\pgfsys@transformshift{2.593266in}{1.026787in}%
\pgfsys@useobject{currentmarker}{}%
\end{pgfscope}%
\begin{pgfscope}%
\pgfsys@transformshift{2.782223in}{0.561777in}%
\pgfsys@useobject{currentmarker}{}%
\end{pgfscope}%
\end{pgfscope}%
\begin{pgfscope}%
\pgfpathrectangle{\pgfqpoint{0.514740in}{0.476543in}}{\pgfqpoint{2.361962in}{1.412556in}}%
\pgfusepath{clip}%
\pgfsetrectcap%
\pgfsetroundjoin%
\pgfsetlinewidth{0.903375pt}%
\definecolor{currentstroke}{rgb}{1.000000,0.000000,0.000000}%
\pgfsetstrokecolor{currentstroke}%
\pgfsetdash{}{0pt}%
\pgfpathmoveto{\pgfqpoint{0.892654in}{1.588429in}}%
\pgfpathlineto{\pgfqpoint{1.509850in}{1.588848in}}%
\pgfpathlineto{\pgfqpoint{1.823195in}{1.587037in}}%
\pgfpathlineto{\pgfqpoint{1.965625in}{1.584419in}}%
\pgfpathlineto{\pgfqpoint{2.117551in}{1.579456in}}%
\pgfpathlineto{\pgfqpoint{2.316952in}{1.569827in}}%
\pgfpathlineto{\pgfqpoint{2.383420in}{1.564723in}}%
\pgfpathlineto{\pgfqpoint{2.411906in}{1.561211in}}%
\pgfpathlineto{\pgfqpoint{2.430896in}{1.555677in}}%
\pgfpathlineto{\pgfqpoint{2.449887in}{1.546951in}}%
\pgfpathlineto{\pgfqpoint{2.468878in}{1.535506in}}%
\pgfpathlineto{\pgfqpoint{2.487868in}{1.521813in}}%
\pgfpathlineto{\pgfqpoint{2.516354in}{1.498092in}}%
\pgfpathlineto{\pgfqpoint{2.563831in}{1.454005in}}%
\pgfpathlineto{\pgfqpoint{2.630298in}{1.390178in}}%
\pgfpathlineto{\pgfqpoint{2.658784in}{1.358692in}}%
\pgfpathlineto{\pgfqpoint{2.687270in}{1.324448in}}%
\pgfpathlineto{\pgfqpoint{2.725251in}{1.274948in}}%
\pgfpathlineto{\pgfqpoint{2.763233in}{1.221630in}}%
\pgfpathlineto{\pgfqpoint{2.782223in}{1.193733in}}%
\pgfpathlineto{\pgfqpoint{2.782223in}{1.193733in}}%
\pgfusepath{stroke}%
\end{pgfscope}%
\begin{pgfscope}%
\pgfpathrectangle{\pgfqpoint{0.514740in}{0.476543in}}{\pgfqpoint{2.361962in}{1.412556in}}%
\pgfusepath{clip}%
\pgfsetbuttcap%
\pgfsetroundjoin%
\definecolor{currentfill}{rgb}{0.000000,0.000000,0.000000}%
\pgfsetfillcolor{currentfill}%
\pgfsetfillopacity{0.000000}%
\pgfsetlinewidth{0.803000pt}%
\definecolor{currentstroke}{rgb}{1.000000,0.000000,0.000000}%
\pgfsetstrokecolor{currentstroke}%
\pgfsetdash{}{0pt}%
\pgfsys@defobject{currentmarker}{\pgfqpoint{-0.025000in}{-0.025000in}}{\pgfqpoint{0.025000in}{0.025000in}}{%
\pgfpathmoveto{\pgfqpoint{0.000000in}{-0.025000in}}%
\pgfpathcurveto{\pgfqpoint{0.006630in}{-0.025000in}}{\pgfqpoint{0.012989in}{-0.022366in}}{\pgfqpoint{0.017678in}{-0.017678in}}%
\pgfpathcurveto{\pgfqpoint{0.022366in}{-0.012989in}}{\pgfqpoint{0.025000in}{-0.006630in}}{\pgfqpoint{0.025000in}{0.000000in}}%
\pgfpathcurveto{\pgfqpoint{0.025000in}{0.006630in}}{\pgfqpoint{0.022366in}{0.012989in}}{\pgfqpoint{0.017678in}{0.017678in}}%
\pgfpathcurveto{\pgfqpoint{0.012989in}{0.022366in}}{\pgfqpoint{0.006630in}{0.025000in}}{\pgfqpoint{0.000000in}{0.025000in}}%
\pgfpathcurveto{\pgfqpoint{-0.006630in}{0.025000in}}{\pgfqpoint{-0.012989in}{0.022366in}}{\pgfqpoint{-0.017678in}{0.017678in}}%
\pgfpathcurveto{\pgfqpoint{-0.022366in}{0.012989in}}{\pgfqpoint{-0.025000in}{0.006630in}}{\pgfqpoint{-0.025000in}{0.000000in}}%
\pgfpathcurveto{\pgfqpoint{-0.025000in}{-0.006630in}}{\pgfqpoint{-0.022366in}{-0.012989in}}{\pgfqpoint{-0.017678in}{-0.017678in}}%
\pgfpathcurveto{\pgfqpoint{-0.012989in}{-0.022366in}}{\pgfqpoint{-0.006630in}{-0.025000in}}{\pgfqpoint{0.000000in}{-0.025000in}}%
\pgfpathlineto{\pgfqpoint{0.000000in}{-0.025000in}}%
\pgfpathclose%
\pgfusepath{stroke,fill}%
}%
\begin{pgfscope}%
\pgfsys@transformshift{0.892654in}{1.588429in}%
\pgfsys@useobject{currentmarker}{}%
\end{pgfscope}%
\begin{pgfscope}%
\pgfsys@transformshift{1.459525in}{1.588903in}%
\pgfsys@useobject{currentmarker}{}%
\end{pgfscope}%
\begin{pgfscope}%
\pgfsys@transformshift{1.648482in}{1.588348in}%
\pgfsys@useobject{currentmarker}{}%
\end{pgfscope}%
\begin{pgfscope}%
\pgfsys@transformshift{1.837438in}{1.586878in}%
\pgfsys@useobject{currentmarker}{}%
\end{pgfscope}%
\begin{pgfscope}%
\pgfsys@transformshift{2.026395in}{1.582766in}%
\pgfsys@useobject{currentmarker}{}%
\end{pgfscope}%
\begin{pgfscope}%
\pgfsys@transformshift{2.215352in}{1.574793in}%
\pgfsys@useobject{currentmarker}{}%
\end{pgfscope}%
\begin{pgfscope}%
\pgfsys@transformshift{2.404309in}{1.562426in}%
\pgfsys@useobject{currentmarker}{}%
\end{pgfscope}%
\begin{pgfscope}%
\pgfsys@transformshift{2.593266in}{1.426553in}%
\pgfsys@useobject{currentmarker}{}%
\end{pgfscope}%
\begin{pgfscope}%
\pgfsys@transformshift{2.782223in}{1.193733in}%
\pgfsys@useobject{currentmarker}{}%
\end{pgfscope}%
\end{pgfscope}%
\begin{pgfscope}%
\pgfpathrectangle{\pgfqpoint{0.514740in}{0.476543in}}{\pgfqpoint{2.361962in}{1.412556in}}%
\pgfusepath{clip}%
\pgfsetbuttcap%
\pgfsetroundjoin%
\pgfsetlinewidth{0.602250pt}%
\definecolor{currentstroke}{rgb}{0.121569,0.466667,0.705882}%
\pgfsetstrokecolor{currentstroke}%
\pgfsetdash{{2.220000pt}{0.960000pt}}{0.000000pt}%
\pgfpathmoveto{\pgfqpoint{0.514740in}{1.394747in}}%
\pgfpathlineto{\pgfqpoint{2.876702in}{1.394747in}}%
\pgfusepath{stroke}%
\end{pgfscope}%
\begin{pgfscope}%
\pgfpathrectangle{\pgfqpoint{0.514740in}{0.476543in}}{\pgfqpoint{2.361962in}{1.412556in}}%
\pgfusepath{clip}%
\pgfsetbuttcap%
\pgfsetroundjoin%
\pgfsetlinewidth{0.602250pt}%
\definecolor{currentstroke}{rgb}{1.000000,0.000000,0.000000}%
\pgfsetstrokecolor{currentstroke}%
\pgfsetdash{{2.220000pt}{0.960000pt}}{0.000000pt}%
\pgfpathmoveto{\pgfqpoint{0.514740in}{1.585439in}}%
\pgfpathlineto{\pgfqpoint{2.876702in}{1.585439in}}%
\pgfusepath{stroke}%
\end{pgfscope}%
\begin{pgfscope}%
\pgfsetrectcap%
\pgfsetmiterjoin%
\pgfsetlinewidth{0.803000pt}%
\definecolor{currentstroke}{rgb}{0.000000,0.000000,0.000000}%
\pgfsetstrokecolor{currentstroke}%
\pgfsetdash{}{0pt}%
\pgfpathmoveto{\pgfqpoint{0.514740in}{0.476543in}}%
\pgfpathlineto{\pgfqpoint{0.514740in}{1.889099in}}%
\pgfusepath{stroke}%
\end{pgfscope}%
\begin{pgfscope}%
\pgfsetrectcap%
\pgfsetmiterjoin%
\pgfsetlinewidth{0.803000pt}%
\definecolor{currentstroke}{rgb}{0.000000,0.000000,0.000000}%
\pgfsetstrokecolor{currentstroke}%
\pgfsetdash{}{0pt}%
\pgfpathmoveto{\pgfqpoint{2.876702in}{0.476543in}}%
\pgfpathlineto{\pgfqpoint{2.876702in}{1.889099in}}%
\pgfusepath{stroke}%
\end{pgfscope}%
\begin{pgfscope}%
\pgfsetrectcap%
\pgfsetmiterjoin%
\pgfsetlinewidth{0.803000pt}%
\definecolor{currentstroke}{rgb}{0.000000,0.000000,0.000000}%
\pgfsetstrokecolor{currentstroke}%
\pgfsetdash{}{0pt}%
\pgfpathmoveto{\pgfqpoint{0.514740in}{0.476543in}}%
\pgfpathlineto{\pgfqpoint{2.876702in}{0.476543in}}%
\pgfusepath{stroke}%
\end{pgfscope}%
\begin{pgfscope}%
\pgfsetrectcap%
\pgfsetmiterjoin%
\pgfsetlinewidth{0.803000pt}%
\definecolor{currentstroke}{rgb}{0.000000,0.000000,0.000000}%
\pgfsetstrokecolor{currentstroke}%
\pgfsetdash{}{0pt}%
\pgfpathmoveto{\pgfqpoint{0.514740in}{1.889099in}}%
\pgfpathlineto{\pgfqpoint{2.876702in}{1.889099in}}%
\pgfusepath{stroke}%
\end{pgfscope}%
\begin{pgfscope}%
\pgfsetbuttcap%
\pgfsetmiterjoin%
\definecolor{currentfill}{rgb}{0.000000,0.000000,0.000000}%
\pgfsetfillcolor{currentfill}%
\pgfsetfillopacity{0.000000}%
\pgfsetlinewidth{0.803000pt}%
\definecolor{currentstroke}{rgb}{0.121569,0.466667,0.705882}%
\pgfsetstrokecolor{currentstroke}%
\pgfsetdash{}{0pt}%
\pgfsys@defobject{currentmarker}{\pgfqpoint{-0.025000in}{-0.025000in}}{\pgfqpoint{0.025000in}{0.025000in}}{%
\pgfpathmoveto{\pgfqpoint{-0.025000in}{-0.025000in}}%
\pgfpathlineto{\pgfqpoint{0.025000in}{-0.025000in}}%
\pgfpathlineto{\pgfqpoint{0.025000in}{0.025000in}}%
\pgfpathlineto{\pgfqpoint{-0.025000in}{0.025000in}}%
\pgfpathlineto{\pgfqpoint{-0.025000in}{-0.025000in}}%
\pgfpathclose%
\pgfusepath{stroke,fill}%
}%
\begin{pgfscope}%
\pgfsys@transformshift{0.648073in}{0.614815in}%
\pgfsys@useobject{currentmarker}{}%
\end{pgfscope}%
\end{pgfscope}%
\begin{pgfscope}%
\definecolor{textcolor}{rgb}{0.000000,0.000000,0.000000}%
\pgfsetstrokecolor{textcolor}%
\pgfsetfillcolor{textcolor}%
\pgftext[x=0.792518in,y=0.575926in,left,base]{\color{textcolor}{\rmfamily\fontsize{8.000000}{9.600000}\selectfont\catcode`\^=\active\def^{\ifmmode\sp\else\^{}\fi}\catcode`\%=\active\def%{\%}$s$}}%
\end{pgfscope}%
\begin{pgfscope}%
\pgfsetbuttcap%
\pgfsetroundjoin%
\definecolor{currentfill}{rgb}{0.000000,0.000000,0.000000}%
\pgfsetfillcolor{currentfill}%
\pgfsetfillopacity{0.000000}%
\pgfsetlinewidth{0.803000pt}%
\definecolor{currentstroke}{rgb}{1.000000,0.000000,0.000000}%
\pgfsetstrokecolor{currentstroke}%
\pgfsetdash{}{0pt}%
\pgfsys@defobject{currentmarker}{\pgfqpoint{-0.025000in}{-0.025000in}}{\pgfqpoint{0.025000in}{0.025000in}}{%
\pgfpathmoveto{\pgfqpoint{0.000000in}{-0.025000in}}%
\pgfpathcurveto{\pgfqpoint{0.006630in}{-0.025000in}}{\pgfqpoint{0.012989in}{-0.022366in}}{\pgfqpoint{0.017678in}{-0.017678in}}%
\pgfpathcurveto{\pgfqpoint{0.022366in}{-0.012989in}}{\pgfqpoint{0.025000in}{-0.006630in}}{\pgfqpoint{0.025000in}{0.000000in}}%
\pgfpathcurveto{\pgfqpoint{0.025000in}{0.006630in}}{\pgfqpoint{0.022366in}{0.012989in}}{\pgfqpoint{0.017678in}{0.017678in}}%
\pgfpathcurveto{\pgfqpoint{0.012989in}{0.022366in}}{\pgfqpoint{0.006630in}{0.025000in}}{\pgfqpoint{0.000000in}{0.025000in}}%
\pgfpathcurveto{\pgfqpoint{-0.006630in}{0.025000in}}{\pgfqpoint{-0.012989in}{0.022366in}}{\pgfqpoint{-0.017678in}{0.017678in}}%
\pgfpathcurveto{\pgfqpoint{-0.022366in}{0.012989in}}{\pgfqpoint{-0.025000in}{0.006630in}}{\pgfqpoint{-0.025000in}{0.000000in}}%
\pgfpathcurveto{\pgfqpoint{-0.025000in}{-0.006630in}}{\pgfqpoint{-0.022366in}{-0.012989in}}{\pgfqpoint{-0.017678in}{-0.017678in}}%
\pgfpathcurveto{\pgfqpoint{-0.012989in}{-0.022366in}}{\pgfqpoint{-0.006630in}{-0.025000in}}{\pgfqpoint{0.000000in}{-0.025000in}}%
\pgfpathlineto{\pgfqpoint{0.000000in}{-0.025000in}}%
\pgfpathclose%
\pgfusepath{stroke,fill}%
}%
\begin{pgfscope}%
\pgfsys@transformshift{0.991554in}{0.614815in}%
\pgfsys@useobject{currentmarker}{}%
\end{pgfscope}%
\end{pgfscope}%
\begin{pgfscope}%
\definecolor{textcolor}{rgb}{0.000000,0.000000,0.000000}%
\pgfsetstrokecolor{textcolor}%
\pgfsetfillcolor{textcolor}%
\pgftext[x=1.135998in,y=0.575926in,left,base]{\color{textcolor}{\rmfamily\fontsize{8.000000}{9.600000}\selectfont\catcode`\^=\active\def^{\ifmmode\sp\else\^{}\fi}\catcode`\%=\active\def%{\%}$F_h$}}%
\end{pgfscope}%
\end{pgfpicture}%
\makeatother%
\endgroup%
}
\caption{(En color en l\'inea) Probabilidad de recuperaci\'on $s$ (cuadrados) y fidelidad condicional $F_h$ (c\'irculos) del modelo completo de ocho estados con puertos reales a $T\omega_b=2000$ en funci\'on del desplazamiento de dos excitones entre puntos $\Delta_{xx}$; las l\'ineas discontinuas son el modelo reducido. Otros par\'ametros: los mismos que en la Fig.~\ref{fig:sweep}, con $N_{\rm cut}=2$ y tolerancia de integraci\'on $10^{-9}$.}
\label{fig:biexciton}
\end{figure}

\section{Requisitos de plataforma}
\label{sec:platform}

Los resultados valen bajo cinco condiciones, independientes entre s\'i. Damos las conversiones de los par\'ametros del modelo a $f_b=\omega_b/2\pi=100$~GHz ($\hbar\omega_b=0.414$~meV) y a 10~GHz ($0.0414$~meV), que se siguen de las definiciones y no son propiedades medidas de un dispositivo, junto con los precedentes parciales que conocemos. La primera es el factor de calidad fon\'onico, $Q=\omega_b/\kappa=10^3$. Se han reportado resonancias ac\'usticas con $Q\sim10^3$ cerca de 100~GHz en microcavidades de Bragg de GaAs/AlAs~\cite{czerniuk2014}, y productos frecuencia--factor de calidad altos en otros materiales~\cite{zalalutdinov2021}; ninguno concierne a una estructura de tres puntos. La segunda es el acoplo por deformaci\'on, $g_s\gtrsim0.03\,\omega_b$ (Fig.~\ref{fig:ports}(d)), es decir $\hbar g_s\gtrsim12.4~\mu$eV a 100~GHz y $1.24~\mu$eV a 10~GHz; el mecanismo de un punto acoplado a un modo de cristal fon\'onico se ha mostrado a 1.466~GHz~\cite{spinnler2024}, pero el umbral en s\'i es una estimaci\'on nuestra. La tercera son los puertos conmutables r\'apidos, $\kappa_{\rm in},\kappa_{\rm out}\gtrsim10\kappa$ (Fig.~\ref{fig:ports}); existen acopladores conmutables de este tipo en circuitos superconductores y de ondas ac\'usticas superficiales a unos pocos GHz, no para fonones a 10--100~GHz.

Las otras dos condiciones dependen de la estructura y de su entorno. La cuarta es el desplazamiento de dos excitones entre puntos, $|\Delta_{xx}|\gtrsim6\,\omega_b$ (Fig.~\ref{fig:biexciton}), es decir, cerca de 2.5~meV a 100~GHz y 0.25~meV a 10~GHz, y no conocemos ninguna medida de esta cantidad para puntos apilados en el rango requerido. La quinta es la correlaci\'on del ruido, $\chi$ cercano a uno, con el valor exacto dependiendo de la espera (Fig.~\ref{fig:requirements}) y tolerante a desajustes moderados de las respuestas Stark; depende de la fuente de ruido y de la geometr\'ia, y no tenemos ninguna medida para esta estructura. Si una estructura puede cumplir las cinco a la vez es una pregunta abierta que este trabajo no aborda.

La mayor probabilidad de recuperaci\'on sugiere la sincronizaci\'on local de operaciones fon\'onicas probabil\'isticas como un posible uso, con purificaci\'on solo si se dispone de varios pares recuperados y de compuertas adecuadas; las memorias cu\'anticas pueden aumentar la tasa de eventos probabil\'isticos sincronizados, mientras que los errores de la memoria limitan el beneficio de la purificaci\'on~\cite{nunn2013,briegel1998,razavi2009}. Nuestro c\'alculo no implementa esas operaciones ni aporta una bandera no destructiva para la proyecci\'on de un fon\'on que define $F_h$. Los tiempos de almacenamiento cortos considerados aqu\'i tampoco establecen una aplicaci\'on a repetidores de larga distancia. Por tanto afirmamos un recurso de entrelazamiento condicional y una ventaja en la probabilidad de recuperaci\'on dentro del modelo, no una memoria tolerante a fallos ni un protocolo operacional de purificaci\'on.

\section{Conclusiones}
\label{sec:conclusions}
Hemos mostrado que tres puntos cu\'anticos acoplados por F\"orster pueden almacenar las dos componentes de un qubit fon\'onico de intervalos temporales en dos excitones oscuros colectivos. El tercer punto es lo que hace posible el esquema: coloca ambas componentes en la misma variedad de un ex\'citon, de modo que las fluctuaciones de energ\'ia de modo com\'un dejan invariante su fase relativa, y las energ\'ias distintas de los dos excitones oscuros permiten direccionarlos por frecuencia con pulsos Raman, sin conmutar ninguna interacci\'on. La deformaci\'on no uniforme, que conecta los estados oscuros con el brillante, es la que permite escribir y leer.

Frente a dos modos ac\'usticos pasivos con los mismos puertos, la memoria mantiene una probabilidad de recuperaci\'on cercana a $0.7$ y una fidelidad incondicional casi constante durante esperas en las que la referencia pasiva pierde casi toda la excitaci\'on. La fidelidad condicionada a la recuperaci\'on, de $0.83$ a $0.88$ para la memoria frente a $0.95$ para la referencia, no mejora: la limitan la escritura y la lectura secuenciales, no el almacenamiento. La ganancia es por tanto la probabilidad de recuperar el qubit, y no una mayor calidad por qubit recuperado.

Esta ventaja depende de condiciones que el c\'alculo hace expl\'icitas: un ruido casi perfectamente correlacionado entre los puntos, pues con ruido local la fidelidad condicional cae por debajo de $0.35$ a la espera m\'as larga; puertos de entrada y salida unas diez veces m\'as r\'apidos que la p\'erdida del modo; un acoplo por deformaci\'on por encima de un umbral; y un desplazamiento de dos excitones entre puntos grande (Sec.~\ref{sec:platform}). Si una sola estructura puede reunirlas queda abierto. El paso siguiente m\'as directo es derivar la correlaci\'on del ruido para una fuente y una geometr\'ia concretas en lugar de tratarla como un par\'ametro, y dotar a la lectura de una bandera no destructiva para la proyecci\'on de un fon\'on, sin la cual una fidelidad condicional por encima de un medio, que certifica entrelazamiento condicional, a\'un no es un recurso operacional.

\bibliography{../referencias_intro}% el estilo apsrev4-2 lo elige la opci\'on aps

\end{document}
